% Options for packages loaded elsewhere
\PassOptionsToPackage{unicode}{hyperref}
\PassOptionsToPackage{hyphens}{url}
\documentclass[
]{article}
\usepackage{xcolor}
\usepackage[margin=25mm]{geometry}
\usepackage{amsmath,amssymb}
\setcounter{secnumdepth}{-\maxdimen} % remove section numbering
\usepackage{iftex}
\ifPDFTeX
  \usepackage[T1]{fontenc}
  \usepackage[utf8]{inputenc}
  \usepackage{textcomp} % provide euro and other symbols
\else % if luatex or xetex
  \usepackage{unicode-math} % this also loads fontspec
  \defaultfontfeatures{Scale=MatchLowercase}
  \defaultfontfeatures[\rmfamily]{Ligatures=TeX,Scale=1}
\fi
\usepackage{lmodern}
\ifPDFTeX\else
  % xetex/luatex font selection
  % LuaLaTeX: Bengali primary font with Latin/symbol fallbacks.
  \directlua{luaotfload.add_fallback("bnfallback", {"Noto Sans:mode=harf;", "Noto Sans Symbols 2:mode=harf;", "DejaVu Sans:mode=harf;"})}
  \setmainfont[Renderer=HarfBuzz,RawFeature={fallback=bnfallback}]{Noto Sans Bengali}
\fi
% Use upquote if available, for straight quotes in verbatim environments
\IfFileExists{upquote.sty}{\usepackage{upquote}}{}
\IfFileExists{microtype.sty}{% use microtype if available
  \usepackage[]{microtype}
  \UseMicrotypeSet[protrusion]{basicmath} % disable protrusion for tt fonts
}{}
\makeatletter
\@ifundefined{KOMAClassName}{% if non-KOMA class
  \IfFileExists{parskip.sty}{%
    \usepackage{parskip}
  }{% else
    \setlength{\parindent}{0pt}
    \setlength{\parskip}{6pt plus 2pt minus 1pt}}
}{% if KOMA class
  \KOMAoptions{parskip=half}}
\makeatother
\usepackage{longtable,booktabs,array}
\newcounter{none} % for unnumbered tables
\usepackage{calc} % for calculating minipage widths
% Correct order of tables after \paragraph or \subparagraph
\usepackage{etoolbox}
\makeatletter
\patchcmd\longtable{\par}{\if@noskipsec\mbox{}\fi\par}{}{}
\makeatother
% Allow footnotes in longtable head/foot
\IfFileExists{footnotehyper.sty}{\usepackage{footnotehyper}}{\usepackage{footnote}}
\makesavenoteenv{longtable}
\ifLuaTeX
  \usepackage{luacolor}
  \usepackage[soul]{lua-ul}
\else
  \usepackage{soul}
\fi
\setlength{\emergencystretch}{3em} % prevent overfull lines
\providecommand{\tightlist}{%
  \setlength{\itemsep}{0pt}\setlength{\parskip}{0pt}}
\usepackage{bookmark}
\IfFileExists{xurl.sty}{\usepackage{xurl}}{} % add URL line breaks if available
\urlstyle{same}
\hypersetup{
  pdftitle={গাণিতিক যুক্তি সম্পূর্ণ টপিকভিত্তিক নোটস},
  hidelinks,
  pdfcreator={LaTeX via pandoc}}

\title{গাণিতিক যুক্তি\\
সম্পূর্ণ টপিকভিত্তিক নোটস}
\author{}
\date{}

\begin{document}
\maketitle

\emph{33টি তাত্ত্বিক টপিক + 42টি MCQ\\
প্রতিটি টপিকের নোটের ঠিক নিচে তার নিজের MCQ অনুশীলন\\
ব্যক্তিগত পড়াশোনা ও পরীক্ষা-প্রস্তুতির জন্য সংকলিত}

নোট: এটি ব্যক্তিগত পড়াশোনা ও পরীক্ষা-প্রস্তুতির জন্য সংকলিত একটি সম্পূর্ণ
নোট-সংকলন।

{\def\LTcaptype{none} % do not increment counter
\begin{longtable}[]{@{}
  >{\raggedright\arraybackslash}p{(\linewidth - 2\tabcolsep) * \real{0.4004}}
  >{\raggedright\arraybackslash}p{(\linewidth - 2\tabcolsep) * \real{0.4004}}@{}}
\toprule\noalign{}
\begin{minipage}[b]{\linewidth}\raggedright
\textbf{তাত্ত্বিক টপিক}
\end{minipage} & \begin{minipage}[b]{\linewidth}\raggedright
\textbf{33টি}
\end{minipage} \\
\midrule\noalign{}
\endhead
\bottomrule\noalign{}
\endlastfoot
MCQ টপিক & 42টি \\
বিন্যাস & টপিকভিত্তিক --- নোট + তার MCQ একসাথে \\
প্রকাশকাল & নভেম্বর ২০২৬ --- বিন্যাস সংশ্লিষ্ট \\
\end{longtable}
}

\section{সূচিপত্র}\label{ux9b8ux99aux9aaux9a4ux9b0}

\emph{নিচের যেকোনো লাইনে ক্লিক করলেই সরাসরি সেই অংশে চলে যাবেন।}

1.
\hyperref[ux9a7ux9b0ux9b0-ux985ux99cux9a8-ux9aaux9a6-ux9a8ux9b0ux9a3ux9df-ux997ux9a3ux9a4ux995-ux995ux9b6ux9b2-ux993-ux9a8ux9dfux9aeux9acux9b2]{\ul{ধারার
অজানা পদ নির্ণয়: গাণিতিক কৌশল ও নিয়মাবলি}}

\begin{quote}
2.
\hyperref[ux9a7ux9b0ux9b0-ux985ux99cux9a8-ux9aaux9a6-ux9a8ux9b0ux9a3ux9df-mcq-ux9aaux9b0ux9b6ux9a8-ux989ux9a4ux9a4ux9b0-ux985ux9a8ux9b6ux9b2ux9a8]{\ul{ধারার
অজানা পদ নির্ণয় MCQ প্রশ্ন উত্তর অনুশীলন}}
\end{quote}

3.
\hyperref[ux985ux9a8ux9aaux9a4-ux993-ux9b8ux9aeux9a8ux9aaux9a4--ux9aaux99fux997ux9a3ux9a4]{\ul{অনুপাত
ও সমানুপাত- পাটিগণিত}}

\begin{quote}
4.
\hyperref[ux985ux9a8ux9aaux9a4-ux9b8ux9aeux9a8ux9aaux9a4-ux9aeux9b6ux9b0ux9a3-ux993-ux985ux9b6ux9a6ux9b0-ux995ux9b0ux9acux9b0-mcq-ux9aaux9b0ux9b6ux9a8-ux989ux9a4ux9a4ux9b0-ux985ux9a8ux9b6ux9b2ux9a8]{\ul{অনুপাত,
সমানুপাত, মিশ্রণ ও অংশীদারী কারবার MCQ প্রশ্ন উত্তর অনুশীলন}}
\end{quote}

5.
\hyperref[ux997ux9a1-average-ux9aeux9b8ux99fux9b0ux995ux9b2ux9b8-ux9acux997ux9a4-ux9b8ux9b2ux9b0-ux9aaux9b0ux9b6ux9a8-ux9acux9afux996ux9af-ux993-ux9b2ux987ux9ad-ux98fux995ux9b8ux9ae]{\ul{গড়
(Average) মাস্টারক্লাস: বিগত সালের প্রশ্ন, ব্যাখ্যা ও লাইভ এক্সাম}}

\begin{quote}
6.
\hyperref[ux997ux9a1-average-ux9aeux9b8ux99fux9b0ux995ux9b2ux9b8-ux9acux997ux9a4-ux9b8ux9b2ux9b0-ux9aaux9b0ux9b6ux9a8-ux9acux9afux996ux9af-ux993-ux9b2ux987ux9ad-ux98fux995ux9b8ux9ae-1]{\ul{গড়
(Average) মাস্টারক্লাস: বিগত সালের প্রশ্ন, ব্যাখ্যা ও লাইভ এক্সাম}}
\end{quote}

7.
\hyperref[ux9adux9a8ux99aux9a4ux9b0-ux9b8ux99f-ux993-ux9abux9b6ux9a8]{\ul{ভেনচিত্র
, সেট ও ফাংশন}}

\begin{quote}
8.
\hyperref[ux9b8ux99f-mcq-ux9aaux9b0ux9b6ux9a8-ux989ux9a4ux9a4ux9b0-ux985ux9a8ux9b6ux9b2ux9a8]{\ul{সেট
MCQ প্রশ্ন উত্তর অনুশীলন}}
\end{quote}

9.
\hyperref[ux9acux9a8ux9afux9b8-ux993-ux9b8ux9aeux9acux9b6--ux98f-ux99f-ux99cux9a1-ux986ux9b2ux99aux9a8]{\ul{বিন্যাস
ও সমাবেশ- এ টু জেড আলোচনা}}

\begin{quote}
10.
\hyperref[ux9acux9a8ux9afux9b8-ux993-ux9b8ux9aeux9acux9b6-mcq-ux9aaux9b0ux9b6ux9a8-ux989ux9a4ux9a4ux9b0-ux985ux9a8ux9b6ux9b2ux9a8]{\ul{বিন্যাস
ও সমাবেশ MCQ প্রশ্ন উত্তর অনুশীলন}}
\end{quote}

11.
\hyperref[ux9a4ux9b0ux9adux99c-ux995ux995-ux9acux9b2-ux9a4ux9b0ux9adux99cux9b0-ux9aaux9b0ux995ux9b0ux9adux9a6-ux9acux9b6ux9b7ux99fux9af-ux993-ux985ux9a8ux9b8ux9a6ux9a7ux9a8ux9a4-ux9aaux9b0ux9a3ux999ux997-ux997ux987ux9a1]{\ul{ত্রিভুজ
কাকে বলে? ত্রিভুজের প্রকারভেদ, বৈশিষ্ট্য ও অনুসিদ্ধান্ত (পূর্ণাঙ্গ গাইড)}}

\begin{quote}
12.
\hyperref[ux9a4ux9b0ux9adux99c-mcq-ux9aaux9b0ux9b6ux9a8-ux989ux9a4ux9a4ux9b0-ux985ux9a8ux9b6ux9b2ux9a8]{\ul{ত্রিভুজ
MCQ প্রশ্ন উত্তর অনুশীলন}}
\end{quote}

13.
\hyperref[ux9acux9b9ux9adux99c-ux995ux995-ux9acux9b2-ux9b8ux99cux99e-ux9b8ux9a4ux9b0ux9acux9b2-ux993-ux995ux9a3ux9b0-ux9aaux9b0ux9aeux9a3-ux9a8ux9b0ux9a3ux9df]{\ul{বহুভুজ
কাকে বলে ? সংজ্ঞা ,সূত্রাবলী ও কোণের পরিমাণ নির্ণয়}}

\begin{quote}
14.
\hyperref[ux9acux9b9ux9adux99c-mcq-ux9aaux9b0ux9b6ux9a8-ux989ux9a4ux9a4ux9b0-ux985ux9a8ux9b6ux9b2ux9a8]{\ul{বহুভুজ
MCQ প্রশ্ন উত্তর অনুশীলন}}
\end{quote}

15.
\hyperref[ux99aux9a4ux9b0ux9adux99c-ux995ux995-ux9acux9b2-ux9b8ux99cux99e-ux9acux9b6ux9b7ux99fux9af-ux9aaux9b0ux995ux9b0ux9adux9a6ux9b8ux9b9-ux98f-ux99f-ux99cux9a1]{\ul{চতুর্ভুজ
কাকে বলে? সংজ্ঞা ,বৈশিষ্ট্য , প্রকারভেদসহ এ টু জেড}}

\begin{quote}
16.
\hyperref[ux99aux9a4ux9b0ux9adux99c-mcq-ux9aaux9b0ux9b6ux9a8-ux989ux9a4ux9a4ux9b0-ux985ux9a8ux9b6ux9b2ux9a8]{\ul{চতুর্ভুজ
MCQ প্রশ্ন উত্তর অনুশীলন}}
\end{quote}

17.
\hyperref[ux99cux9afux9aeux9a4--ux9acux9a8ux9a6-ux9b0ux996-ux98fux9ac-ux995ux9a3-a-to-z]{\ul{জ্যামিতি-\/-বিন্দু
, রেখা এবং কোণ \textbar\textbar{} A to Z}}

\begin{quote}
18.
\hyperref[ux9acux9a8ux9a6-ux9b0ux996-ux993-ux995ux9a3-mcq-ux9aaux9b0ux9b6ux9a8-ux989ux9a4ux9a4ux9b0-ux985ux9a8ux9b6ux9b2ux9a8]{\ul{বিন্দু,
রেখা ও কোণ MCQ প্রশ্ন উত্তর অনুশীলন}}
\end{quote}

19.
\hyperref[ux9b8ux9aeux9adux9acux9a8-probability--ux997ux9a3ux9a4]{\ul{সম্ভাবনা
(probability)- গণিত}}

\begin{quote}
20.
\hyperref[ux9b8ux9aeux9adux9acux9afux9a4-mcq-ux9aaux9b0ux9b6ux9a8-ux989ux9a4ux9a4ux9b0-ux985ux9a8ux9b6ux9b2ux9a8]{\ul{সম্ভাব্যতা
MCQ প্রশ্ন উত্তর অনুশীলন}}
\end{quote}

21.
\hyperref[ux9a6ux9acux998ux9a4-ux9b8ux9aeux995ux9b0ux9a3-ux9ac-ux9b8ux9b0ux9b2-ux9b8ux9b9-ux9b8ux9aeux995ux9b0ux9a3-ux9b8ux99cux99e-ux9b8ux9a4ux9b0ux9acux9b2ux9b8ux9b9-ux98f-ux99f-ux99cux9a1]{\ul{দ্বিঘাত
সমীকরণ বা সরল সহ-সমীকরণ \textbar\textbar{} সংজ্ঞা, সূত্রাবলীসহ এ টু জেড}}

\begin{quote}
22.
\hyperref[ux9a6ux9acux998ux9a4-ux9b8ux9aeux995ux9b0ux9a3-mcq-ux9aaux9b0ux9b6ux9a8-ux989ux9a4ux9a4ux9b0-ux985ux9a8ux9b6ux9b2ux9a8]{\ul{দ্বিঘাত
সমীকরণ MCQ প্রশ্ন উত্তর অনুশীলন}}
\end{quote}

23.
\hyperref[ux9b8ux9b0ux9b2-ux9b8ux9aeux995ux9b0ux9a3-ux9acux99cux997ux9a3ux9a4-ux9b8ux99cux99e-ux9b8ux9a4ux9b0ux9acux9b2ux9b8ux9b9-ux98f-ux99f-ux99cux9a1]{\ul{সরল
সমীকরণ-বীজগণিত \textbar\textbar{} সংজ্ঞা, সূত্রাবলীসহ এ টু জেড}}

\begin{quote}
24.
\hyperref[ux9b8ux9aeux995ux9b0ux9a3-mcq-ux9aaux9b0ux9b6ux9a8-ux989ux9a4ux9a4ux9b0-ux985ux9a8ux9b6ux9b2ux9a8]{\ul{সমীকরণ
MCQ প্রশ্ন উত্তর অনুশীলন}}

25.
\hyperref[ux9b8ux9b0ux9b2-ux9b8ux9b9-ux9b8ux9aeux995ux9b0ux9a3-mcq-ux9aaux9b0ux9b6ux9a8-ux989ux9a4ux9a4ux9b0-ux985ux9a8ux9b6ux9b2ux9a8]{\ul{সরল
সহ-সমীকরণ MCQ প্রশ্ন উত্তর অনুশীলন}}
\end{quote}

26.
\hyperref[ux9b8ux9aeux9a8ux9a4ux9b0-ux993-ux997ux9a3ux9a4ux9a4ux9b0-ux9a7ux9b0-ux9acux99cux997ux9a3ux9a4-ux9b8ux99cux99e-ux9b8ux9a4ux9b0ux9acux9b2ux9b8ux9b9-ux98f-ux99f-ux99cux9a1]{\ul{সমান্তর
ও গুণোত্তর ধারা-বীজগণিত \textbar\textbar{} সংজ্ঞা, সূত্রাবলীসহ এ টু জেড}}

\begin{quote}
27.
\hyperref[ux9b8ux9aeux9a8ux9a4ux9b0-ux993-ux997ux9a3ux9a4ux9a4ux9b0-ux985ux9a8ux995ux9b0ux9ae-ux993-ux9a7ux9b0-mcq-ux9aaux9b0ux9b6ux9a8-ux989ux9a4ux9a4ux9b0-ux985ux9a8ux9b6ux9b2ux9a8]{\ul{সমান্তর
ও গুণোত্তর অনুক্রম ও ধারা MCQ প্রশ্ন উত্তর অনুশীলন}}
\end{quote}

28.
\hyperref[ux9b2ux997ux9b0ux9a6ux9ae-ux9acux99cux997ux9a3ux9a4-ux9b8ux99cux99e-ux9b8ux9a4ux9b0ux9acux9b2ux9b8ux9b9-ux98f-ux99f-ux99cux9a1]{\ul{লগারিদম-বীজগণিত
\textbar\textbar{} সংজ্ঞা, সূত্রাবলীসহ এ টু জেড}}

\begin{quote}
29.
\hyperref[ux9b2ux997ux9b0ux9a6ux9ae-mcq-ux9aaux9b0ux9b6ux9a8-ux989ux9a4ux9a4ux9b0-ux985ux9a8ux9b6ux9b2ux9a8]{\ul{লগারিদম
MCQ প্রশ্ন উত্তর অনুশীলন}}
\end{quote}

30.
\hyperref[ux9b8ux99aux995-ux995-ux9b8ux99cux99e-ux9b8ux9a4ux9b0ux9acux9b2-ux98fux9ac-ux9b8ux99aux995-ux98fux9b0-ux985ux995-ux98f-ux99f-ux99cux9a1]{\ul{সূচক
কী ? সংজ্ঞা, সূত্রাবলী এবং সূচক এর অংক এ টু জেড}}

\begin{quote}
31.
\hyperref[ux9b8ux99aux995-mcq-ux9aaux9b0ux9b6ux9a8-ux989ux9a4ux9a4ux9b0-ux985ux9a8ux9b6ux9b2ux9a8]{\ul{সূচক
MCQ প্রশ্ন উত্তর অনুশীলন}}
\end{quote}

32.
\hyperref[ux9acux99cux997ux9a3ux9a4ux995-ux9b0ux9b6ux9aeux9b2ux9b0-ux9b2.ux9b8.ux997-ux993-ux997.ux9b8.ux997]{\ul{বীজগাণিতিক
রাশিমালার ল.সা.গু ও গ.সা.গু}}

\begin{quote}
33.
\hyperref[ux9acux99cux997ux9a3ux9a4ux995-ux9b0ux9b6ux9aeux9b2ux9b0-ux9b2.ux9b8.ux997-ux993-ux997.ux9b8.ux997-mcq-ux9aaux9b0ux9b6ux9a8-ux989ux9a4ux9a4ux9b0-ux985ux9a8ux9b6ux9b2ux9a8]{\ul{বীজগাণিতিক
রাশিমালার ল.সা.গু ও গ.সা.গু MCQ প্রশ্ন উত্তর অনুশীলন}}

34.
\hyperref[ux9b2.ux9b8.ux997-ux98fux9ac-ux997.ux9b8.ux997.-ux9b8ux9aeux9aaux9b0ux995ux9a4-ux9b8ux9aeux9b8ux9af-mcq-ux9aaux9b0ux9b6ux9a8-ux989ux9a4ux9a4ux9b0-ux985ux9a8ux9b6ux9b2ux9a8]{\ul{ল.সা.গু
এবং গ.সা.গু. সম্পর্কিত সমস্যা MCQ প্রশ্ন উত্তর অনুশীলন}}
\end{quote}

35.
\hyperref[ux989ux9ceux9aaux9a6ux995-ux9acux9b6ux9b2ux9b7ux9a3-ux9acux99cux997ux9a3ux9a4-ux98f-ux99f-ux99cux9a1]{\ul{উৎপাদকে
বিশ্লেষণ-বীজগণিত \textbar\textbar{} এ টু জেড}}

\begin{quote}
36.
\hyperref[ux989ux9ceux9aaux9a6ux995-ux9acux9b6ux9b2ux9b7ux9a3-ux9acux9b9ux9aaux9a6-ux989ux9ceux9aaux9a6ux995-mcq-ux9aaux9b0ux9b6ux9a8-ux989ux9a4ux9a4ux9b0-ux985ux9a8ux9b6ux9b2ux9a8]{\ul{উৎপাদকে
বিশ্লেষণ / বহুপদী উৎপাদক MCQ প্রশ্ন উত্তর অনুশীলন}}
\end{quote}

37.
\hyperref[ux9acux99cux997ux9a3ux9a4ux9df-ux9b0ux9b6ux9aeux9b2ux9b0-ux9aaux9b0ux9a5ux9aeux995-ux99aux9b0-ux9a8ux9dfux9ae]{\ul{বীজগাণিতিয়
রাশিমালার প্রাথমিক চার নিয়ম}}

\begin{quote}
38.
\hyperref[ux9acux99cux997ux9a3ux9a4ux9df-ux9b0ux9b6ux9aeux9b2ux9b0-ux9aaux9b0ux9a5ux9aeux995-ux99aux9b0-ux9a8ux9dfux9ae-mcq-ux9aaux9b0ux9b6ux9a8-ux989ux9a4ux9a4ux9b0-ux985ux9a8ux9b6ux9b2ux9a8]{\ul{বীজগাণিতিয়
রাশিমালার প্রাথমিক চার নিয়ম MCQ প্রশ্ন উত্তর অনুশীলন}}
\end{quote}

39.
\hyperref[ux9adux99cux995-ux9b8ux996ux9af-ux9a8ux9b0ux9a3ux9afux9b0-ux9b8ux9a4ux9b0--ux997ux9a3ux9a4]{\ul{ভাজক
সংখ্যা নির্ণয়ের সূত্র -গণিত}}

\begin{quote}
40.
\hyperref[ux9b8ux996ux9af-ux9b8ux9aeux9aaux9b0ux995ux9a4-ux9b8ux9aeux9b8ux9af-mcq-ux9aaux9b0ux9b6ux9a8-ux989ux9a4ux9a4ux9b0-ux985ux9a8ux9b6ux9b2ux9a8]{\ul{সংখ্যা
সম্পর্কিত সমস্যা MCQ প্রশ্ন উত্তর অনুশীলন}}
\end{quote}

41.
\hyperref[ux998ux9a1ux9b0-ux995ux99fux9b0-ux985ux995-ux9b8ux9b0ux9acux9b6ux9b7-ux9b8ux9b9ux99c-ux9a8ux9dfux9ae]{\ul{ঘড়ির
কাঁটার অংক \textbar\textbar{} সর্বশেষ সহজে নিয়মে}}

\begin{quote}
42.
\hyperref[ux998ux9dcux9b0-ux995ux99fux9b0-ux985ux995-mcq-ux9aaux9b0ux9b6ux9a8-ux989ux9a4ux9a4ux9b0-ux985ux9a8ux9b6ux9b2ux9a8]{\ul{ঘড়ির
কাঁটার অংক MCQ প্রশ্ন উত্তর অনুশীলন}}
\end{quote}

43.
\hyperref[ux9aaux9b0ux9b8ux996ux9afux9a8-statistics-ux997ux9a3ux9a4]{\ul{পরিসংখ্যান
(Statistics)-গণিত}}

\begin{quote}
44.
\hyperref[ux9aaux9b0ux9b8ux996ux9afux9a8-mcq-ux9aaux9b0ux9b6ux9a8-ux989ux9a4ux9a4ux9b0-ux985ux9a8ux9b6ux9b2ux9a8]{\ul{পরিসংখ্যান
MCQ প্রশ্ন উত্তর অনুশীলন}}
\end{quote}

45.
\hyperref[ux995ux9b7ux9a4ux9b0ux9abux9b2-ux993-ux9aaux9b0ux9aeux9aa-ux99aux995ux9b0ux9b0-ux9aaux9b0ux995ux9b7ux9b0-ux9b6ux9b0ux99fux995ux99f-ux99fux995ux9a8ux995-ux993-ux997ux9b0ux9a4ux9acux9aaux9b0ux9a3-ux985ux995ux9b0-ux9b8ux9aeux9a7ux9a8]{\ul{ক্ষেত্রফল
ও পরিমাপ: চাকরির পরীক্ষার শর্টকাট টেকনিক ও গুরুত্বপূর্ণ অংকের সমাধান}}

\begin{quote}
46.
\hyperref[ux9aaux9b0ux9aeux9aa-ux993-ux98fux995ux995-mcq-ux9aaux9b0ux9b6ux9a8-ux989ux9a4ux9a4ux9b0-ux985ux9a8ux9b6ux9b2ux9a8]{\ul{পরিমাপ
ও একক MCQ প্রশ্ন উত্তর অনুশীলন}}

47.
\hyperref[ux9aaux9b0ux9aeux9a4-mcq-ux9aaux9b0ux9b6ux9a8-ux989ux9a4ux9a4ux9b0-ux985ux9a8ux9b6ux9b2ux9a8]{\ul{পরিমিতি
MCQ প্রশ্ন উত্তর অনুশীলন}}
\end{quote}

48.
\hyperref[ux9b8ux9a6ux995ux9b7ux9b0-ux99cux9a6ux995ux9b0-ux9b8ux9aeux9a7ux9a8-ux9e7ux9e6ux9e6-ux98fux9b0-ux9a7ux9b0ux9a3ux9df-ux9aaux99fux997ux9a3ux9a4ux9b0-ux9b8ux9b9ux99c-ux9aaux9a0]{\ul{সুদকষার
জাদুকরী সমাধান: ১০০\%-এর ধারণায় পাটিগণিতের সহজ পাঠ}}

\begin{quote}
49.
\hyperref[ux9b8ux9b0ux9b2-ux993-ux99aux995ux9b0ux9acux9a6ux9a7-ux9aeux9a8ux9ab-ux9b8ux9a6ux995ux9b7-mcq-ux9aaux9b0ux9b6ux9a8-ux989ux9a4ux9a4ux9b0-ux985ux9a8ux9b6ux9b2ux9a8]{\ul{সরল
ও চক্রবৃদ্ধি মুনাফা / সুদকষা MCQ প্রশ্ন উত্তর অনুশীলন}}
\end{quote}

50.
\hyperref[ux9b2ux9ad-ux995ux9b7ux9a4ux9b0-ux9e7ux9e6ux9e6-ux9aeux9afux99cux995-ux9b6ux9b0ux99fux995ux99f-ux9a8ux9dfux9ae-ux9aaux99fux997ux9a3ux9a4ux9b0-ux9b8ux9aeux9a7ux9a8]{\ul{লাভ-ক্ষতির
১০০\% ম্যাজিক: শর্টকাট নিয়মে পাটিগণিতের সমাধান}}

\begin{quote}
51.
\hyperref[ux9b2ux9ad-ux993-ux995ux9b7ux9a4-mcq-ux9aaux9b0ux9b6ux9a8-ux989ux9a4ux9a4ux9b0-ux985ux9a8ux9b6ux9b2ux9a8]{\ul{লাভ
ও ক্ষতি MCQ প্রশ্ন উত্তর অনুশীলন}}
\end{quote}

52.
\hyperref[ux9b6ux9a4ux995ux9b0-ux995-ux98fux996ux9a8-ux9b6ux996ux9a8-ux9a6ux9b0ux9a6ux9a8ux9a4-ux995ux9b6ux9b2]{\ul{শতকরা
কি ? এখন শিখুন দুর্দান্ত কেীশলে !!!}}

\begin{quote}
53.
\hyperref[ux9b6ux9a4ux995ux9b0-mcq-ux9aaux9b0ux9b6ux9a8-ux989ux9a4ux9a4ux9b0-ux985ux9a8ux9b6ux9b2ux9a8]{\ul{শতকরা
MCQ প্রশ্ন উত্তর অনুশীলন}}
\end{quote}

54.
\hyperref[ux990ux995ux995-ux9a8ux9afux9aeux9b0-ux985ux995-ux9b6ux9b0ux99fux995ux99f-ux99fux995ux9a8ux995-ux993-ux997ux9b0ux9a4ux9acux9aaux9b0ux9a3-ux9aaux9b0ux9b6ux9a8-ux9b8ux9aeux9a7ux9a8]{\ul{ঐকিক
নিয়মের অংক: শর্টকাট টেকনিক ও গুরুত্বপূর্ণ প্রশ্ন সমাধান}}

\begin{quote}
55.
\hyperref[ux990ux995ux995-ux9a8ux9dfux9ae-mcq-ux9aaux9b0ux9b6ux9a8-ux989ux9a4ux9a4ux9b0-ux985ux9a8ux9b6ux9b2ux9a8]{\ul{ঐকিক
নিয়ম MCQ প্রশ্ন উত্তর অনুশীলন}}
\end{quote}

56.
\hyperref[ux9aeux9b6ux9b0ux9a3-ux9b8ux995ux9b0ux9a8ux9a4-ux9b8ux995ux9b2-ux985ux995ux9b0-ux9b8ux9b9ux99c-ux9a8ux9dfux9ae-allegation]{\ul{মিশ্রণ
সংক্রান্ত সকল অংকের সহজ নিয়ম \textbar\textbar{} Allegation}}

57.
\hyperref[ux9acux9b0ux997-ux993-ux9acux9b0ux997ux9aeux9b2-ux9a8ux9b0ux9a3ux9dfux9b0-ux9b8ux9b0ux9acux9b6ux9b7-ux9b8ux9b9ux99c-ux99fux995ux9a8ux995]{\ul{বর্গ
ও বর্গমূল নির্ণয়ের সর্বশেষ সহজ টেকনিক}}

\begin{quote}
58.
\hyperref[ux9acux9b0ux997-ux993-ux9acux9b0ux997ux9aeux9b2-mcq-ux9aaux9b0ux9b6ux9a8-ux989ux9a4ux9a4ux9b0-ux985ux9a8ux9b6ux9b2ux9a8]{\ul{বর্গ
ও বর্গমূল MCQ প্রশ্ন উত্তর অনুশীলন}}
\end{quote}

59.
\hyperref[ux9a6ux9b6ux9aeux995-ux9adux997ux9a8ux9b6ux9b0-ux9afux997-ux9acux9afux997-ux997ux9a3-ux993-ux9adux997-ux9b6ux996ux9b0-ux9b8ux9b9ux99c-ux989ux9aaux9df]{\ul{দশমিক
ভগ্নাংশের যোগ, বিয়োগ, গুণ ও ভাগ শেখার সহজ উপায়}}

\begin{quote}
60.
\hyperref[ux9a6ux9b6ux9aeux995-ux9adux997ux9a8ux9b6-mcq-ux9aaux9b0ux9b6ux9a8-ux989ux9a4ux9a4ux9b0-ux985ux9a8ux9b6ux9b2ux9a8]{\ul{দশমিক
ভগ্নাংশ MCQ প্রশ্ন উত্তর অনুশীলন}}
\end{quote}

61.
\hyperref[ux9b2.ux9b8.ux997-ux993-ux997.ux9b8.ux997-lcm-hcf-ux9a8ux9b0ux9a3ux9dfux9b0-ux9a8ux99eux99c-ux99fux995ux9a8ux995]{\ul{ল.সা.গু
ও গ.সা.গু (LCM \& HCF) নির্ণয়ের নিঞ্জা টেকনিক}}

62.
\hyperref[ux9b8ux9a7ux9b0ux9a3-ux9adux997ux9a8ux9b6-ux98f-ux99f-ux99cux9a1-ux9afux997-ux9acux9afux997-ux997ux9a3-ux9adux997-ux993-ux99bux99f-ux9acux9a1-ux9a8ux9b0ux9a3ux9afux9b0-ux9b8ux9b9ux99c-ux9a8ux9afux9ae]{\ul{সাধারণ
ভগ্নাংশ এ টু জেড: যোগ, বিয়োগ, গুণ, ভাগ ও ছোট-বড় নির্ণয়ের সহজ নিয়ম}}

\begin{quote}
63.
\hyperref[ux9b8ux9a7ux9b0ux9a3-ux9adux997ux9a8ux9b6-mcq-ux9aaux9b0ux9b6ux9a8-ux989ux9a4ux9a4ux9b0-ux985ux9a8ux9b6ux9b2ux9a8]{\ul{সাধারণ
ভগ্নাংশ MCQ প্রশ্ন উত্তর অনুশীলন}}
\end{quote}

64.
\hyperref[ux9b8ux996ux9afux9dfux9a8-ux985ux9a7ux9afux9df-ux9aeux9b2ux995-ux9b8ux9a4ux9b0ux9acux9b2-ux9b8ux9aeux9a7ux9a8-ux993-ux9b8ux9b9ux99c-ux995ux9b6ux9b2ux9b0-ux9aaux9b0ux9a3ux999ux997-ux9a8ux9b0ux9a6ux9b6ux995]{\ul{``সংখ্যায়ন''
অধ্যায়: মৌলিক সূত্রাবলি, সমাধান ও সহজ কৌশলের পূর্ণাঙ্গ নির্দেশিকা}}

65.
\hyperref[ux9acux9b8ux9a4ux9ac-ux9b8ux996ux9afux9b0-ux9aaux9b0ux9a3ux999ux997-ux997ux987ux9a1-ux9b8ux99cux99e-ux9aaux9b0ux995ux9b0ux9adux9a6-ux993-ux9b6ux9b0ux99fux995ux99f-ux99fux9aaux9b8]{\ul{বাস্তব
সংখ্যার পূর্ণাঙ্গ গাইড: সংজ্ঞা, প্রকারভেদ ও শর্টকাট টিপস}}

\begin{quote}
66.
\hyperref[ux9acux9b8ux9a4ux9ac-ux9b8ux996ux9af-mcq-ux9aaux9b0ux9b6ux9a8-ux989ux9a4ux9a4ux9b0-ux985ux9a8ux9b6ux9b2ux9a8]{\ul{বাস্তব
সংখ্যা MCQ প্রশ্ন উত্তর অনুশীলন}}
\end{quote}

\hyperref[ux9b6ux9a7-mcq-ux985ux9a8ux9b6ux9b2ux9a8-ux9a4ux9a4ux9a4ux9acux995-ux9a8ux99f-ux9a8ux987]{\textbf{\ul{শুধু
MCQ অনুশীলন (তাত্ত্বিক নোট নেই)}}}

67.
\hyperref[ux9acux9a4ux9a4-mcq-ux9aaux9b0ux9b6ux9a8-ux989ux9a4ux9a4ux9b0-ux985ux9a8ux9b6ux9b2ux9a8]{\ul{বৃত্ত
MCQ প্রশ্ন উত্তর অনুশীলন}}

68.
\hyperref[ux9a4ux9b0ux995ux9a3ux9aeux9a4-mcq-ux9aaux9b0ux9b6ux9a8-ux989ux9a4ux9a4ux9b0-ux985ux9a8ux9b6ux9b2ux9a8]{\ul{ত্রিকোণমিতি
MCQ প্রশ্ন উত্তর অনুশীলন}}

69.
\hyperref[ux998ux9a8-ux99cux9afux9aeux9a4-mcq-ux9aaux9b0ux9b6ux9a8-ux989ux9a4ux9a4ux9b0-ux985ux9a8ux9b6ux9b2ux9a8]{\ul{ঘন
জ্যামিতি MCQ প্রশ্ন উত্তর অনুশীলন}}

70.
\hyperref[ux985ux9b8ux9aeux9a4-mcq-ux9aaux9b0ux9b6ux9a8-ux989ux9a4ux9a4ux9b0-ux985ux9a8ux9b6ux9b2ux9a8]{\ul{অসমতা
MCQ প্রশ্ন উত্তর অনুশীলন}}

71.
\hyperref[ux9acux99cux997ux9a3ux9a4ux9df-ux9adux997ux9a8ux9b6-mcq-ux9aaux9b0ux9b6ux9a8-ux989ux9a4ux9a4ux9b0-ux985ux9a8ux9b6ux9b2ux9a8]{\ul{বীজগণিতীয়
ভগ্নাংশ MCQ প্রশ্ন উত্তর অনুশীলন}}

72.
\hyperref[ux9acux99cux997ux9a3ux9a4ux995-ux9b8ux9a4ux9b0ux9acux9b2-ux993-ux9aeux9a8-ux9a8ux9b0ux9a3ux9af-mcq-ux9aaux9b0ux9b6ux9a8-ux989ux9a4ux9a4ux9b0-ux985ux9a8ux9b6ux9b2ux9a8]{\ul{বীজগাণিতিক
সূত্রাবলী ও মান নির্ণয় MCQ প্রশ্ন উত্তর অনুশীলন}}

73.
\hyperref[ux986ux9a8ux9a4ux9b0ux99cux9a4ux995-ux997ux9a3ux9a8-ux9aaux9a6ux9a7ux9a4-mcq-ux9aaux9b0ux9b6ux9a8-ux989ux9a4ux9a4ux9b0-ux985ux9a8ux9b6ux9b2ux9a8]{\ul{আন্তর্জাতিক
গণনা পদ্ধতি MCQ প্রশ্ন উত্তর অনুশীলন}}

74.
\hyperref[ux9b8ux9aeux9df-ux9a6ux9b0ux9a4ux9ac-ux993-ux997ux9a4ux9acux997-ux9b8ux995ux9b0ux9a8ux9a4-ux985ux995-mcq-ux9aaux9b0ux9b6ux9a8-ux989ux9a4ux9a4ux9b0-ux985ux9a8ux9b6ux9b2ux9a8]{\ul{সময়,
দুরত্ব ও গতিবেগ সংক্রান্ত অংক MCQ প্রশ্ন উত্তর অনুশীলন}}

75.
\hyperref[ux9acux9dfux9b8-ux9b8ux995ux9b0ux9a8ux9a4-ux985ux995-mcq-ux9aaux9b0ux9b6ux9a8-ux989ux9a4ux9a4ux9b0-ux985ux9a8ux9b6ux9b2ux9a8]{\ul{বয়স
সংক্রান্ত অংক MCQ প্রশ্ন উত্তর অনুশীলন}}

\hyperref[ux9aaux9b0ux9b6ux9b7ux99f-ux9a6ux9b0ux9a4-ux9aaux9b0ux9b8ux9a4ux9a4-ux997ux9a3ux9a4ux995-ux9afux995ux9a4-ux9b8ux9a4ux9b0ux9b8ux997ux9b0ux9b9]{\textbf{\ul{পরিশিষ্ট
--- দ্রুত প্রস্তুতি: গাণিতিক যুক্তি সূত্রসংগ্রহ}}}

\section{1. ধারার অজানা পদ নির্ণয়: গাণিতিক কৌশল ও
নিয়মাবলি}\label{ux9a7ux9b0ux9b0-ux985ux99cux9a8-ux9aaux9a6-ux9a8ux9b0ux9a3ux9df-ux997ux9a3ux9a4ux995-ux995ux9b6ux9b2-ux993-ux9a8ux9dfux9aeux9acux9b2}

বিভিন্ন প্রতিযোগিতামূলক পরীক্ষার (যেমন: বিসিএস, ব্যাংক, প্রাথমিক শিক্ষক নিয়োগ
ইত্যাদি) গাণিতিক যুক্তি অংশে \textbf{ধারার অজানা পদ নির্ণয়} বা পরবর্তী সংখ্যা
বের করার প্রশ্ন অত্যন্ত আবশ্যিক ও গুরুত্বপূর্ণ একটি বিষয়। ধারা মূলত কতগুলো পদের এমন
একটি সাজানো বিন্যাস, যা নির্দিষ্ট কোনো গাণিতিক সম্পর্ক বা প্যাটার্ন মেনে চলে।
ধারার অজানা পদ বা পরবর্তী পদ দ্রুত ও নির্ভুলভাবে বের করার কৌশল জানা থাকলে
পরীক্ষায় অনেক সময় বাঁচানো সম্ভব। নিচে ধাপে ধাপে ধারার অজানা পদ নির্ণয়ের সকল
মৌলিক ও উন্নত গাণিতিক কৌশল উদাহরণসহ আলোচনা করা হলো।

\begin{itemize}
\item
  সমান্তর ধারা ও নির্দিষ্ট ব্যবধান বৃদ্ধি/হ্রাস
\item
  গুণোত্তর ধারা
\item
  গুণোত্তর হারে পার্থক্য বৃদ্ধি/হ্রাস
\item
  ফিবোনাচ্চি ও যোগভিত্তিক ধারা
\item
  বর্গ ও ঘন ভিত্তিক ধারা
\item
  বহুপদী বা মিশ্র গাণিতিক অপারেশন
\item
  দ্বৈত বা অল্টারনেট ধারা
\item
  সমানুপাতিক ধারা
\item
  মৌলিক সংখ্যা ও বিশেষ প্যাটার্ন ভিত্তিক ধারা
\end{itemize}

\subsubsection{সমান্তর ধারা ও নির্দিষ্ট ব্যবধান বৃদ্ধি/হ্রাস (Arithmetic
Sequence)}\label{ux9b8ux9aeux9a8ux9a4ux9b0-ux9a7ux9b0-ux993-ux9a8ux9b0ux9a6ux9b7ux99f-ux9acux9afux9acux9a7ux9a8-ux9acux9a6ux9a7ux9b9ux9b0ux9b8-arithmetic-sequence}

এই জাতীয় ধারায় পাশাপাশি দুটি পদের পার্থক্য স্থির থাকে অথবা পার্থক্য নির্দিষ্ট
হারে (যেমন ১, ২, ৩...) বাড়ে বা কমে।

\begin{itemize}
\item
  \textbf{স্থির পার্থক্য: ১, ৩, ৫, ৭, ... ধারাটির ৮ম পদ কত?}
\end{itemize}

\begin{itemize}
\item
  কৌশল/ব্যাখ্যা: এখানে প্রতি পদে পার্থক্য ২ করে বাড়ছে। এটি বিজোড় সংখ্যার ধারা
  (২n − ১)। ৮ম পদ = ২ × ৮ − ১ = ১৫।
\end{itemize}

\begin{itemize}
\item
  \textbf{ক্রমবর্ধমান/ক্রমহ্রাসমান পার্থক্য: ১, ৩, ৬, ১০, ১৫, ২১, ... ধারার
  ১২তম (দ্বাদশ) পদ কত?}
\end{itemize}

\begin{itemize}
\item
  কৌশল/ব্যাখ্যা: এখানে পার্থক্যগুলো হলো যথাক্রমে ২, ৩, ৪, ৫, ৬। পার্থক্য ১ করে
  বাড়ছে। একইভাবে:
\end{itemize}

\begin{itemize}
\item
  ৭ম পদ = ২১ + ৭ = ২৮
\item
  ৮ম পদ = ২৮ + ৮ = ৩৬
\item
  ৯তম পদ = ৩৬ + ৯ = ৪৫
\item
  ১০ম পদ = ৪৫ + ১০ = ৫৫
\item
  ১১তম পদ = ৫৫ + ১১ = ৬৬
\item
  ১২তম পদ = ৬৬ + ১২ = ৭৮
\end{itemize}

\subsubsection{গুণোত্তর ধারা (Geometric
Sequence)}\label{ux997ux9a3ux9a4ux9a4ux9b0-ux9a7ux9b0-geometric-sequence}

যখন পূর্ববর্তী পদকে কোনো নির্দিষ্ট সংখ্যা বা ভগ্নাংশ দিয়ে গুণ বা ভাগ করে পরবর্তী
পদ তৈরি করা হয়।

\begin{itemize}
\item
  \textbf{গুণভিত্তিক: ৩, ৯, ২৭, ... ধারাটির পরবর্তী সংখ্যা কত?}
\end{itemize}

\begin{itemize}
\item
  কৌশল/ব্যাখ্যা: প্রতি পদকে ৩ দিয়ে গুণ করা হচ্ছে। পরবর্তী সংখ্যা = ২৭ × ৩ = ৮১।
\end{itemize}

\begin{itemize}
\item
  \textbf{ভাগভিত্তিক: ৮১, ২৭, ..., ৩, ১ -- লুপ্ত সংখ্যাটি কত?}
\end{itemize}

\begin{itemize}
\item
  কৌশল/ব্যাখ্যা: পূর্ববর্তী সংখ্যাকে ৩ দিয়ে ভাগ করা হচ্ছে। লুপ্ত সংখ্যা = ২৭ ÷ ৩ =
  ৯।
\end{itemize}

\begin{itemize}
\item
  \textbf{দশমিক ধারা: ০.২, ০.০৪, ০.০০৮, ০.০০১৬, ... এর পরবর্তী পদ কত?}
\end{itemize}

\begin{itemize}
\item
  কৌশল/ব্যাখ্যা: প্রতি পদে ০.২ দ্বারা গুণ হচ্ছে। পরবর্তী সংখ্যা = ০.০০১৬ × ০.২ =
  ০.০০০৩২।
\end{itemize}

\subsubsection{গুণোত্তর হারে পার্থক্য বৃদ্ধি/হ্রাস (Exponential
Difference)}\label{ux997ux9a3ux9a4ux9a4ux9b0-ux9b9ux9b0-ux9aaux9b0ux9a5ux995ux9af-ux9acux9a6ux9a7ux9b9ux9b0ux9b8-exponential-difference}

এই ধারায় সংখ্যাগুলোর পার্থক্য গুণোত্তর বা নির্দিষ্ট হারে বৃদ্ধি বা হ্রাস পায়।

\begin{itemize}
\item
  \textbf{পার্থক্য দ্বিগুণ বৃদ্ধি: ৮, ১১, ১৭, ২৯, ৫৩, ... পরবর্তী সংখ্যাটি কত?}
\end{itemize}

\begin{itemize}
\item
  কৌশল/ব্যাখ্যা: পদের পার্থক্যগুলো হলো ৩, ৬, ১২, ২৪ (পার্থক্য দ্বিগুণ হচ্ছে)।
  পরবর্তী পার্থক্য হবে ২৪ × ২ = ৪৮। অতএব, পরবর্তী পদ = ৫৩ + ৪৮ = ১০১।
\end{itemize}

\begin{itemize}
\item
  \textbf{পার্থক্য তিনগুণ বৃদ্ধি: ১, ৪, ১৩, ৪০, ... পরবর্তী পদ কত?}
\end{itemize}

\begin{itemize}
\item
  কৌশল/ব্যাখ্যা: পার্থক্যগুলো যথাক্রমে ৩, ৯, ২৭ (তিনগুণ হচ্ছে)। পরবর্তী পার্থক্য ২৭
  × ৩ = ৮১। অতএব, পদটি = ৪০ + ৮১ = ১২১।
\end{itemize}

\subsubsection{ফিবোনাচ্চি ও যোগভিত্তিক ধারা (Fibonacci
Series)}\label{ux9abux9acux9a8ux99aux99a-ux993-ux9afux997ux9adux9a4ux9a4ux995-ux9a7ux9b0-fibonacci-series}

পূর্ববর্তী দুটি পদের যোগফল দিয়ে পরবর্তী পদ তৈরি করার কৌশল।

\begin{itemize}
\item
  \textbf{ফিবোনাচ্চি অনুক্রম: ১, ১, ২, ৩, ৫, ৮, ... এই ধারায় ৮ম পদ কত?}
\end{itemize}

\begin{itemize}
\item
  কৌশল/ব্যাখ্যা: ১ম+২য় = ৩য় পদ (১ + ১ = ২), ২য়+৩য় = ৪র্থ পদ (১ + ২ = ৩)।
  এভাবে, ৭ম পদ = ৫ + ৮ = ১৩ এবং ৮ম পদ = ৮ + ১৩ = ২১।
\end{itemize}

\begin{itemize}
\item
  \textbf{অন্যান্য যোগভিত্তিক ধারা: ১, ৩, ৪, ৭, ১১, ১৮, ... পরবর্তী পদ কত?}
\end{itemize}

\begin{itemize}
\item
  কৌশল/ব্যাখ্যা: পূর্ববর্তী দুটি পদের যোগফল = পরবর্তী পদ। পরবর্তী সংখ্যাটি = ১১ +
  ১৮ = ২৯।
\end{itemize}

\subsubsection{বর্গ ও ঘন ভিত্তিক ধারা (Square \& Cube
Series)}\label{ux9acux9b0ux997-ux993-ux998ux9a8-ux9adux9a4ux9a4ux995-ux9a7ux9b0-square-cube-series}

ধারার পদগুলো সরাসরি পূর্ণবর্গ, ঘন অথবা বর্গ/ঘনের অনুক্রম মেনে চলে।

\begin{itemize}
\item
  \textbf{জোড়/বিজোড় সংখ্যার বর্গ: ৯, ৩৬, ৮১, ১৪৪, ... পরবর্তী পদ কত?}
\end{itemize}

\begin{itemize}
\item
  কৌশল/ব্যাখ্যা: পদগুলো হলো ৩, ৬, ৯, ১২ এর বর্গ। পরবর্তী পদ হবে ১৫ এর বর্গ =
  ১৫ ২ = ২২৫।
\end{itemize}

\begin{itemize}
\item
  বিজোড় বর্গ: ১, ৯, ২৫, ৪৯, ৮১, ...
\end{itemize}

\begin{itemize}
\item
  কৌশল/ব্যাখ্যা: ১, ৩, ৫, ৭, ৯ এর বর্গ। পরবর্তী পদ ১১ ২ = ১২১।
\end{itemize}

\begin{itemize}
\item
  ক্রমহ্রাসমান বর্গ: ১৪৪, ৮১, ৩৬, ...
\end{itemize}

\begin{itemize}
\item
  কৌশল/ব্যাখ্যা: পদগুলো ১২, ৯, ৬ এর বর্গ (৩ করে কমছে)। পরবর্তী পদ হবে ৩ ২ =
  ৯।
\end{itemize}

\subsubsection{বহুপদী বা মিশ্র গাণিতিক অপারেশন (Composite Operational
Series)}\label{ux9acux9b9ux9aaux9a6-ux9ac-ux9aeux9b6ux9b0-ux997ux9a3ux9a4ux995-ux985ux9aaux9b0ux9b6ux9a8-composite-operational-series}

গুণের সাথে যোগ বা বিয়োগ বা একাধিক গাণিতিক প্রক্রিয়া একসাথে প্রয়োগ করে অজানা
পদ নির্ণয় করা হয়।

\begin{itemize}
\item
  \textbf{উদাহরণ ১: ৫, ১৪, ৪০, ১১৭, ... পরবর্তী পদ কত?}
\end{itemize}

\begin{itemize}
\item
  কৌশল/ব্যাখ্যা: (৫ × ৩) − ১ = ১৪ (১৪ × ৩) − ২ = ৪০ (৪০ × ৩) − ৩ = ১১৭
  পরবর্তী পদ = (১১৭ × ৩) − ৪ = ৩৫১ − ৪ = ৩৪৭।
\end{itemize}

\begin{itemize}
\item
  উদাহরণ ২: ৬, ১৭, ৪৯, ১৪৪, ...
\end{itemize}

\begin{itemize}
\item
  কৌশল/ব্যাখ্যা: (৬ × ৩) − ১ = ১৭ (১৭ × ৩) − ২ = ৪৯ (৪৯ × ৩) − ৩ = ১৪৪
  পরবর্তী পদ = (১৪৪ × ৩) − ৪ = ৪৩২ − ৪ = ৪২৮।
\end{itemize}

\subsubsection{দ্বৈত বা অল্টারনেট ধারা (Interleaved/Alternate
Series)}\label{ux9a6ux9acux9a4-ux9ac-ux985ux9b2ux99fux9b0ux9a8ux99f-ux9a7ux9b0-interleavedalternate-series}

একটি ধারার ভেতর দুটি আলাদা ধারা পর পর সাজানো থাকে (বিজোড় ও জোড় স্থান
অনুযায়ী)।

\begin{itemize}
\item
  \textbf{উদাহরণ: ৪, ১১, ৮, ১৯, ১২, ... পরবর্তী সংখ্যাটি কত?}
\end{itemize}

\begin{itemize}
\item
  কৌশল/ব্যাখ্যা: ১ম ধারা (১ম, ৩য়, ৫ম পদ): ৪, ৮, ১২ (৪ করে বাড়ছে)। ২য় ধারা
  (২য়, ৪র্থ, ৬ষ্ঠ পদ): ১১, ১৯, ... (৮ করে বাড়ছে)। অতএব, ৬ষ্ঠ বা পরবর্তী পদটি
  হবে ১৯ + ৮ = ২৭।
\end{itemize}

\subsubsection{সমানুপাতিক ধারা (Proportional
Series)}\label{ux9b8ux9aeux9a8ux9aaux9a4ux995-ux9a7ux9b0-proportional-series}

অনুপাত ও সমানুপাতের নীতি ব্যবহার করে লুপ্ত পদ নির্ণয় করা যায়।

\begin{itemize}
\item
  উদাহরণ: ১২ : ১৬ :: ? : ২০ (লুপ্ত পদ নির্ণয়)
\end{itemize}

\begin{itemize}
\item
  কৌশল/ব্যাখ্যা: লুপ্ত পদ = k হলে, (১২/১৬)=(k/২০)⇒(৩/৪)=(k/২০)⇒৪k=৬০⇒k=১৫
\end{itemize}

\subsubsection{মৌলিক সংখ্যা ও বিশেষ প্যাটার্ন ভিত্তিক
ধারা}\label{ux9aeux9b2ux995-ux9b8ux996ux9af-ux993-ux9acux9b6ux9b7-ux9aaux9afux99fux9b0ux9a8-ux9adux9a4ux9a4ux995-ux9a7ux9b0}

\begin{itemize}
\item
  মৌলিক সংখ্যার ক্রম: ২, ৩, ৫, ৭, ১১, ১৩, ১৭, ... ধারাটির পরবর্তী পদ হবে ১৯
  (যেহেতু পদগুলো মৌলিক সংখ্যা)।
\item
  অঙ্কের যোগফল বা স্থানান্তরের ধারা: অনেক সময় সংখ্যার মধ্যকার অঙ্কগুলোর
  (Digits) সমষ্টি বা প্যাটার্ন পরিবর্তন করে ধারা সাজানো হয়।
\end{itemize}

ধারার অজানা পদ নির্ণয়ের ক্ষেত্রে কোনো একটি নির্দিষ্ট সূত্রে সব অংক সমাধান করা
যায় না। প্রশ্ন সমাধান করার সময় প্রথমেই লক্ষ্য করতে হবে পদগুলোর বৃদ্ধির গতি
কেমন---ধীরগতিতে বাড়লে সমান্তর, দ্রুত বাড়লে গুণোত্তর/বর্গ এবং এলোমেলোভাবে কমলে
বা বাড়লে দ্বৈত বা ফিবোনাচ্চি নিয়ম প্রয়োগ করতে হবে। নিয়মিত বিভিন্ন
প্রতিযোগিতামূলক পরীক্ষার প্রশ্ন ও প্যাটার্ন অনুশীলনের মাধ্যমে ধারার যেকোনো অজানা
পদ খুব সহজেই নির্ণয় করা সম্ভব।

সময় অবশিষ্ট আছে:

\section{2. ধারার অজানা পদ নির্ণয় MCQ প্রশ্ন উত্তর
অনুশীলন}\label{ux9a7ux9b0ux9b0-ux985ux99cux9a8-ux9aaux9a6-ux9a8ux9b0ux9a3ux9df-mcq-ux9aaux9b0ux9b6ux9a8-ux989ux9a4ux9a4ux9b0-ux985ux9a8ux9b6ux9b2ux9a8}

\textbf{৫, ১১, ১৯, ২৯, ..... ধারার পরের সংখ্যা কত?}

ক) ৩৫

খ) ৩৭

\textbf{\hl{গ) ৪১ ✔}}

ঘ) ৩৯

\textbf{\hl{✔ সঠিক উত্তর: গ) ৪১}}

\begin{quote}
\textbf{ব্যাখ্যা:} \emph{১ম পদ = ৫}

\emph{২য় পদ = ১ম পদ + ৬ = ৫ + ৬ = ১১}

\emph{৩য় পদ = ২য় পদ + ৮ = ১১ + ৮ = ১৯}

\emph{৪র্থ পদ = ৩য় পদ + ১০ = ১৯ + ১০ = ২৯}

\emph{৫ম পদ = ৪র্থ পদ + ১২ = ২৯ + ১২ = ৪১}
\end{quote}

\textbf{৫, ৭, ১১, ১৯, -\/-\/-\/-\/-\/- শূন্যস্থানে কোন সংখ্যাটি বসবে?}

ক) ২৬

খ) ৩০

\textbf{\hl{গ) ৩৫ ✔}}

ঘ) ৪৩

\textbf{\hl{✔ সঠিক উত্তর: গ) ৩৫}}

\begin{quote}
\textbf{ব্যাখ্যা:} \emph{৫ + ২ = ৭ ৭ + ৪ = ১১ ১১ + ৮ = ১৯ ১৯ + ১৬ = ৩৫}
\end{quote}

\textbf{9, 36, 81...... এর ৬ষ্ঠ পদের সংখ্যাটি কত?}

\textbf{\hl{ক) 324 ✔}}

খ) 289

গ) 256

ঘ) 361

\textbf{\hl{✔ সঠিক উত্তর: ক) 324}}

\begin{quote}
\textbf{ব্যাখ্যা:} \emph{ধারাটি ৩, ৬, ৯, ১২, ১৫, ১৮ সংখ্যাগুলোর বর্গেরসংখ্যা
আকারে এগিয়েছে।}

\emph{∴ ৬ষ্ঠ সংখ্যাটি = (১৮)\^{}২ = 324}
\end{quote}

\textbf{২ - ৫ - ১২ - ১৯ ..... ধারাবাহিক ১২তম পদ?}

ক) -৬০

খ) -৬৫

গ) -৭০

\textbf{\hl{ঘ) -৭৫ ✔}}

\textbf{\hl{✔ সঠিক উত্তর: ঘ) -৭৫}}

\begin{quote}
\textbf{ব্যাখ্যা:} \emph{ধরি, ধারটির ১ম পদ a = 2}

\emph{সাধারণ অন্তর d = − 5 − 2 = − 7}

\emph{∴ 12তম পদ = a + (12 − 1)d = a + 11d}

\emph{= 2 + 11(− 7) = 2 − 77 = − 75}
\end{quote}

\textbf{1 + 3 + 6 + 10 + 15 + ....................ধারাটির সপ্তম পদটি কত?}

ক) 20

খ) 25

গ) 21

\textbf{\hl{ঘ) 28 ✔}}

\textbf{\hl{✔ সঠিক উত্তর: ঘ) 28}}

\begin{quote}
\textbf{ব্যাখ্যা:} \emph{ধারাটিতে,}

\emph{প্রথম পদ = 1}

\emph{দ্বিতীয় পদ = 3 = 1 + 2}

\emph{তৃতীয় পদ = 6 = 3 + 3}

\emph{চতুর্থ পদ = 10 = 6 + 4}

\emph{পঞ্চম পদ = 15 = 10 + 5}

\emph{ষষ্ঠ পদ = 21 = 15 + 6}

\emph{∴ সপ্তম পদ = ষষ্ঠ পদ + 7 = 21 + 7 = 28}

\emph{{[}সঠিক উত্তর 28{]}}
\end{quote}

\textbf{৩,৯,২৭......ধারাটির পরবর্তী সংখ্যা কত ?}

ক) ৩৬

খ) ৪৮

গ) ৬৩

\textbf{\hl{ঘ) ৮১ ✔}}

\textbf{\hl{✔ সঠিক উত্তর: ঘ) ৮১}}

\begin{quote}
\textbf{ব্যাখ্যা:} \emph{পূর্ববর্তী সংখ্যাকে ৩ দিয়ে গুণ করে পরবর্তী সংখ্যা পাওয়া
যাচ্ছে।তাই ধারাটির পরবর্তী সংখ্যা হবে ২৭×৩ = ৮১ ।}
\end{quote}

\textbf{১,২,৪,৮,১৬.......ধারাটির দশম পদ কত ?}

ক) ২৫৬

খ) ৫০০

\textbf{\hl{গ) ৫১২ ✔}}

ঘ) ৬২২

\textbf{\hl{✔ সঠিক উত্তর: গ) ৫১২}}

\begin{quote}
\textbf{ব্যাখ্যা:} \emph{পূর্ববর্তী সংখ্যাকে ২ দিয়ে গুণ করে পরবর্তী সংখ্যা পাওয়া
যাচ্ছে।}

\emph{তাই ধারাটির ৬ষ্ঠ পদ ১৬×২=৩২,}

\emph{৭ম পদ ৩২×২=৬৪,}

\emph{৮ম পদ ৬৪×২=১২৮,}

\emph{৯ম পদ=১২৮×২=২৫৬,}

\emph{১০ম পদ=২৫৬×২=৫১২।}
\end{quote}

\textbf{১,৩,৪,৭,১১,..... ধারাটির পরবর্তী সংখ্যা কত ?}

ক) ৮

খ) ১৪

\textbf{\hl{গ) ১৮ ✔}}

ঘ) ২৫

\textbf{\hl{✔ সঠিক উত্তর: গ) ১৮}}

\begin{quote}
\textbf{ব্যাখ্যা:} \emph{পূর্ববর্তী দুটি সংখ্যার যোগফল=পরবর্তী সংখ্যা। তাই
ধারাটির পরবর্তী সংখ্যা হবে ৭+১১ = ১৮।}
\end{quote}

\textbf{১,২,৩,৫,৮,১৩......ধারাটির পরবর্তী সংখ্যা কত ?}

ক) ১৭

খ) ১৯

\textbf{\hl{গ) ২১ ✔}}

ঘ) ২৩

\textbf{\hl{✔ সঠিক উত্তর: গ) ২১}}

\begin{quote}
\textbf{ব্যাখ্যা:} \emph{পূর্ববর্তী দুটি সংখ্যার যোগফল = পরবর্তী সংখ্যা। তাই
ধারাটির পরবর্তী সংখ্যা হবে ৮+১৩ = ২১।}
\end{quote}

\textbf{১,৩,৪,৭,১১,১৮,......... ক্রমটির পরবর্তী পদ কত ?}

ক) ২৫

\textbf{\hl{খ) ২৯ ✔}}

গ) ৩৬

ঘ) ৪২

\textbf{\hl{✔ সঠিক উত্তর: খ) ২৯}}

\begin{quote}
\textbf{ব্যাখ্যা:} \emph{পূর্ববর্তী দুটি সংখ্যার যোগফল=পরবর্তী সংখ্যা। তাই
ধারাটির পরবর্তী সংখ্যা হবে ১১+১৮= ২৯।}
\end{quote}

\textbf{১,৩,৬,১০,১৫,২১,.............ধারার দ্বাদশ পদ কত ?}

ক) ৬২

খ) ৬৬

\textbf{\hl{গ) ৭৮ ✔}}

ঘ) ৫৫

\textbf{\hl{✔ সঠিক উত্তর: গ) ৭৮}}

\begin{quote}
\textbf{ব্যাখ্যা:} \emph{ধারাটির প্রতি পদের পার্থক্য হচ্ছে ২,৩,৪,৫,৬ । এখানে
পার্থক্য ১ করে বাড়ছে।}

\emph{তাই ধারাটির ৭ম পদ= ২১+৭=২৮,}

\emph{৮ম পদ= ২৮+৮= ৩৬,}

\emph{৯ম পদ = ৩৬+৯=৪৫,}

\emph{১০ম পদ= ৪৫+১০=৫৫,}

\emph{১১শ পদ= ৫৫+১১=৬৬,}

\emph{১২শ পদ= ৬৬+১২=৭৮।}
\end{quote}

\textbf{০.২,০.০৪,০.০০৮,০.০০১৬ ধারাটির পরবর্তী সংখ্যাটি কত ?}

ক) ০.০০২৪

খ) ০.০০৩২

\textbf{\hl{গ) ০.০০০৩২ ✔}}

ঘ) ০.০০০০৩২

\textbf{\hl{✔ সঠিক উত্তর: গ) ০.০০০৩২}}

\begin{quote}
\textbf{ব্যাখ্যা:} \emph{প্রতিবারে ০.২ দ্বারা গুণ করে পরবর্তী সংখ্যা পাওয়া
যাচ্ছে। তাই পরবর্তী সংখ্যা হবে ০.০০১৬×০.২ = ০.০০০৩২।}
\end{quote}

\textbf{১১,১৫,২৩,৩৯....ধারাটির পরবর্তী সংখ্যা কত ?}

ক) ৫২

খ) ৬৫

\textbf{\hl{গ) ৭১ ✔}}

ঘ) ৯২

\textbf{\hl{✔ সঠিক উত্তর: গ) ৭১}}

\begin{quote}
\textbf{ব্যাখ্যা:} \emph{এখানে,পার্থক্য দ্বিগুণ হারে বাড়ছে। প্রতিটি পদের পার্থক্য
যথাক্রমে ৪,৮,১৬,... । তাই ধারাটির পরবর্তী সংখ্যা হবে = ৩৯+৩২ = ৭১}
\end{quote}

\textbf{লুপ্ত পদ নির্ণয় করুনঃ-\/-\/-\/- ১২:১৬: : :২০}

ক) ১৮

\textbf{\hl{খ) ১৫ ✔}}

গ) ১০

ঘ) ২২

\textbf{\hl{✔ সঠিক উত্তর: খ) ১৫}}

\begin{quote}
\textbf{ব্যাখ্যা:} \emph{ধরি,লুপ্ত পদ = ক}

\emph{প্রশ্নমতে, ১২:১৬ = ক:২০}

\emph{বা, ৩:৪ = ক:২০}

\emph{বা, ৪ক = ৬০}

\emph{বা, ক = ৬০÷৪ = ১৫}
\end{quote}

\textbf{১,৪,১৩,৪০, এর ধারাটির পরবর্তী পদ কত ?}

ক) ৩৯

খ) ৮১

\textbf{\hl{গ) ১২১ ✔}}

ঘ) ৩৬৩

\textbf{\hl{✔ সঠিক উত্তর: গ) ১২১}}

\begin{quote}
\textbf{ব্যাখ্যা:} \emph{এখানে প্রতি পদের পার্থক্য তিনগুণ করে বাড়ছে।যেমন-
১+৩=৪ ৪+৯=১৩ ১৩+২৭=৪০ ৪০+৮১=১২১}
\end{quote}

\textbf{২,৬,১০,১৪,.......ধারাটির ৭ম পদ কত ?}

ক) ২২

\textbf{\hl{খ) ২৬ ✔}}

গ) ২৮

ঘ) ৩০

\textbf{\hl{✔ সঠিক উত্তর: খ) ২৬}}

\begin{quote}
\textbf{ব্যাখ্যা:} \emph{৪ এর নামতার প্রত্যেকটি থেকে ২ বিয়োগ।}

\emph{ধারাটির ১ম পদ= ৪×১=৪-২=২,}

\emph{ধারাটির ২য় পদ= ৪×২=৮-২=৬,}

\emph{ধারাটির ৩য় পদ= ৪×৩=১২-২=১০,}

\emph{ধারাটির ৪র্থ পদ= ৪×৪=১৬-২=১৪,}

\emph{ধারাটির ৫ম পদ= ৪×৫=২০-২=১৮,}

\emph{ধারাটির ৬ষ্ঠ পদ= ৪×৬=২৪-২=২২,}

\emph{ধারাটির ৭ম পদ= ৪×৭=২৮-২=২৬।}
\end{quote}

\textbf{৫,১৪,৪০,১১৭,....পরবর্তী পদ কত ?}

ক) ২৮০

খ) ৩০০

\textbf{\hl{গ) ৩৪৭ ✔}}

ঘ) ৩৫১

\textbf{\hl{✔ সঠিক উত্তর: গ) ৩৪৭}}

\begin{quote}
\textbf{ব্যাখ্যা:} \emph{এখানে,৫+৫×২-১ = ১৪ ১৪+১৪×২-২ = ৪০ ৪০+৪০×২-৩ =
১১৭}

\emph{পরবর্তী সংখ্যা = ১১৭+১১৭×২-৪ = ৩৪৭}
\end{quote}

\textbf{১৪৪,৮১,৩৬ এর পরবর্তী সংখ্যাটি কত ?}

ক) ১৮

খ) ১২

\textbf{\hl{গ) ৯ ✔}}

ঘ) ৬

\textbf{\hl{✔ সঠিক উত্তর: গ) ৯}}

\begin{quote}
\textbf{ব্যাখ্যা:} \emph{এখানে ধারাবাহিকভাবে ১২,৯,৬ এবং ৩ এর বর্গ আকারে
ধারাটি সাজানো হয়েছে।তাই পরবর্তী সংখ্যাটি হবে ৩ এর বর্গ = ৯।}
\end{quote}

\textbf{৮৪,৪০,১৮-\/-\/-এর পরবর্তী সংখ্যা কত ?}

ক) ১৬

খ) ১২

গ) ৯

\textbf{\hl{ঘ) ৭ ✔}}

\textbf{\hl{✔ সঠিক উত্তর: ঘ) ৭}}

\begin{quote}
\textbf{ব্যাখ্যা:} \emph{প্রতি পদের পার্থক্য ৪৪,২২,১১ এমনভাবে কমছে। তাই
পরবর্তী সংখ্যাটি হবে ১৮-১১ = ৭।}
\end{quote}

\textbf{৭,১০,১৬,২৮,৫২,.....এর পরবর্তী সংখ্যা কত ?}

\textbf{\hl{ক) ১০০ ✔}}

খ) ১০৪

গ) ১৩০

ঘ) ১৫২

\textbf{\hl{✔ সঠিক উত্তর: ক) ১০০}}

\begin{quote}
\textbf{ব্যাখ্যা:} \emph{প্রতি পদের পার্থক্য ৩,৬,১২,২৪ এমনভাবে বাড়ছে। তাই
পরবর্তী সংখ্যাটি হবে ৫২+৪৮ = ১০০।}
\end{quote}

\textbf{৯, ৩৬, ৮১, ১৪৪,..........। পরবর্তী সংখ্যা কত ?}

ক) ১৬৯

\textbf{\hl{খ) ২২৫ ✔}}

গ) ২৫৬

ঘ) ২৭২

\textbf{\hl{✔ সঠিক উত্তর: খ) ২২৫}}

\begin{quote}
\textbf{ব্যাখ্যা:} \emph{এখানে, প্রতিটি সংখ্যা ৩,৬,৯,১২ এর বর্গ। তাই পরবর্তী
সংখ্যা হবে ১৫ এর বর্গ।}
\end{quote}

\textbf{১, ২, ৩, ৫, ৮, ১৩, ২১, ৩৪,...........ধারার পরবর্তী সংখ্যাটি কত ?}

\textbf{\hl{ক) ৫৫ ✔}}

খ) ৪০

গ) ৬৮

ঘ) ৮৯

\textbf{\hl{✔ সঠিক উত্তর: ক) ৫৫}}

\begin{quote}
\textbf{ব্যাখ্যা:} \emph{এখানে পূর্ববর্তী দুটি সংখ্যার যোগফল হচ্ছে পরবর্তী
সংখ্যা।তাই পরবর্তী সংখ্যাটি হবে ২১+৩৪ = ৫৫।}
\end{quote}

\textbf{১, ১, ২, ৩, ৫, ৮,.........., এই সংখ্যা পরস্পরায় অষ্টম পদ কত ?}

\textbf{\hl{ক) ২১ ✔}}

খ) ১৩

গ) ১৯

ঘ) ১৬

\textbf{\hl{✔ সঠিক উত্তর: ক) ২১}}

\begin{quote}
\textbf{ব্যাখ্যা:} \emph{এখানে পূর্ববর্তী দুটি সংখ্যার যোগফল হচ্ছে পরবর্তী সংখ্যা।
তাই ৭ম পদ হবে ৫+৮ = ১৩ এবং ৮ম পদ হবে ১৩ + ৮ = ২১ ।}
\end{quote}

\textbf{৮, ১১, ১৭, ২৯, ৫৩, .............। পরবর্তী সংখ্যা কত ?}

\textbf{\hl{ক) ১০১ ✔}}

খ) ১০২

গ) ৭৫

ঘ) ৫৯

\textbf{\hl{✔ সঠিক উত্তর: ক) ১০১}}

\begin{quote}
\textbf{ব্যাখ্যা:} \emph{এখানে, প্রতিটি সংখ্যার পার্থক্য হচ্ছে যথাক্রমে
৩,৬,১২,২৪।অর্থাৎ পার্থক্য দ্বিগুণহারে বাড়ছে।তাই পরবর্তী সংখ্যাটি হবে ৫৩+৪৮ =
১০১।}
\end{quote}

\textbf{১৯, ৩৩, ৫১, ৭৩,.........। পরবর্তী সংখ্যাটি কত ?}

ক) ৮৫

খ) ১২১

\textbf{\hl{গ) ৯৯ ✔}}

ঘ) ৯৮

\textbf{\hl{✔ সঠিক উত্তর: গ) ৯৯}}

\begin{quote}
\textbf{ব্যাখ্যা:} \emph{এখানে, প্রতিটি সংখ্যার পার্থক্য হচ্ছে যথাক্রমে
১৪,,১৮,২২।অর্থাৎ পার্থক্য ৪ করে বাড়ছে।তাই পরবর্তী সংখ্যাটি হবে ৭৩+২৬ = ৯৯।}
\end{quote}

\textbf{১, ৩, ৬, ১০, ১৫, .... ধারাটির পরবর্তী পদ কত ?}

ক) ১৭

খ) ১৯

গ) ২০

\textbf{\hl{ঘ) ২১ ✔}}

\textbf{\hl{✔ সঠিক উত্তর: ঘ) ২১}}

\begin{quote}
\textbf{ব্যাখ্যা:} \emph{এখানে, প্রতিটি সংখ্যার পার্থক্য হচ্ছে যথাক্রমে ২,৩,৪,৫
অর্থাৎ পার্থক্য ১ করে বাড়ছে।তাই পরবর্তী পদ হবে ১৫+৬ = ২১।}
\end{quote}

\textbf{১, ৩, ৭, ...., ২১, ৩১, ৪৩ ধারার মধ্যবর্তী সংখ্যা কত ?}

\textbf{\hl{ক) ১৩ ✔}}

খ) ১৫

গ) ১৭

ঘ) ১৯

\textbf{\hl{✔ সঠিক উত্তর: ক) ১৩}}

\begin{quote}
\textbf{ব্যাখ্যা:} \emph{এখানে, প্রতিটি সংখ্যার পার্থক্য হচ্ছে যথাক্রমে
২,৪,....৬,৮,১০,১২।অর্থাৎ পার্থক্য ২ করে বাড়ছে।তাই মধ্যবর্তী সংখ্যা হবে ৭+৬ =
১৩।}
\end{quote}

\textbf{২, ৮, ১৮, ৩২ ধারাটির পরবর্তী সংখ্যা কত ?}

ক) ৩৮

খ) ৪২

\textbf{\hl{গ) ৫০ ✔}}

ঘ) ৪৮

\textbf{\hl{✔ সঠিক উত্তর: গ) ৫০}}

\begin{quote}
\textbf{ব্যাখ্যা:} \emph{এখানে, প্রতিটি সংখ্যার পার্থক্য হচ্ছে যথাক্রমে
৬,১০,১৪।অর্থাৎ পার্থক্য ৪ করে বাড়ছে।তাই পরবর্তী সংখ্যা হবে ৩২+১৮ = ৫০।}
\end{quote}

\textbf{ধারাটির পরবর্তী সংখ্যা কত ?}

২, ৪, ৮, ১৪, ২২, ৩২,..................

ক) ৪২

\textbf{\hl{খ) ৪৪ ✔}}

গ) ৫৪

ঘ) ৫৬

\textbf{\hl{✔ সঠিক উত্তর: খ) ৪৪}}

\begin{quote}
\textbf{ব্যাখ্যা:} \emph{এখানে, প্রতিটি সংখ্যার পার্থক্য হচ্ছে যথাক্রমে
২,৪,৬,৮,১০।অর্থাৎ পার্থক্য ২ করে বাড়ছে।তাই পরবর্তী সংখ্যা হবে ৩২+১২ = ৪৪।}
\end{quote}

\textbf{৫ + ১১ + ১৯ + ২৯ + .................. পরের সংখ্যাটি কত ?}

ক) ৩৫

খ) ৩৭

গ) ৩৯

\textbf{\hl{ঘ) ৪১ ✔}}

\textbf{\hl{✔ সঠিক উত্তর: ঘ) ৪১}}

\begin{quote}
\textbf{ব্যাখ্যা:} \emph{এখানে, প্রতিটি সংখ্যার পার্থক্য হচ্ছে যথাক্রমে
৬,৮,১০।অর্থাৎ পার্থক্য ২ করে বাড়ছে।তাই পরবর্তী সংখ্যা হবে ২৯+১২ = ৪১।}
\end{quote}

\textbf{শূণ্যস্থানের সংখ্যাটি কত? ৫২, ..., ৩৯, ৩৪}

ক) ৫০

খ) ৪৮

\textbf{\hl{গ) ৪৫ ✔}}

ঘ) ৪২

\textbf{\hl{✔ সঠিক উত্তর: গ) ৪৫}}

\begin{quote}
\textbf{ব্যাখ্যা:} \emph{এখানে, প্রতিটি সংখ্যার পার্থক্য হচ্ছে যথাক্রমে,
৭,৬,৫।অর্থাৎ পার্থক্য ১ করে কমছে। তাই শূণ্যস্থানের সংখ্যাটি হবে ৫২-৭ = ৪৫।}
\end{quote}

\textbf{৪, ৬, ১০, ১৮ ক্রমটির পরবর্তী পদ কত ?}

ক) ৩৬

\textbf{\hl{খ) ৩৪ ✔}}

গ) ৩২

ঘ) ৩০

\textbf{\hl{✔ সঠিক উত্তর: খ) ৩৪}}

\begin{quote}
\textbf{ব্যাখ্যা:} \emph{এখানে, প্রতিটি সংখ্যার পার্থক্য হচ্ছে যথাক্রমে
২,৪,৮।অর্থাৎ পার্থক্য দ্বিগুণহারে বাড়ছে।তাই পরবর্তী সংখ্যাটি হবে ১৮+১৬ = ৩৪।}
\end{quote}

\textbf{লুপ্ত সংখ্যাটি কত? ৮০, ৯৬, ..., ১২৮}

ক) ৮৮

খ) ১২০

গ) ৬৪

\textbf{\hl{ঘ) ১১২ ✔}}

\textbf{\hl{✔ সঠিক উত্তর: ঘ) ১১২}}

\begin{quote}
\textbf{ব্যাখ্যা:} \emph{এখানে, প্রতিটি সংখ্যার পার্থক্য হচ্ছে ১৬।তাই লুপ্ত
সংখ্যাটি হবে ৯৬+১৬ = ১১২।}
\end{quote}

\textbf{পরবর্তী সংখ্যাটি কত? ৪, ১১, ৮, ১৯, ১২,...}

ক) ২৫

\textbf{\hl{খ) ২৭ ✔}}

গ) ২০

ঘ) ৩০

\textbf{\hl{✔ সঠিক উত্তর: খ) ২৭}}

\begin{quote}
\textbf{ব্যাখ্যা:} \emph{এখানে, দুটি ধারা।}

\emph{১ম ধারাটি হচ্ছে ৪,৮,১২।এটি ৪ করে বাড়ছে।}

\emph{২য় ধারাটি হচ্ছে ১১,১৯।এটি ৮ করে বাড়ছে।}

\emph{তাই পরবর্তী সংখ্যা হচ্ছে ১৯+৮ = ২৭}
\end{quote}

\textbf{৬, ১৭, ৪৯, ১৪৪ ধারাটির পরবর্তী পদ কত ?}

ক) ২০০

খ) ৩৫৬

\textbf{\hl{গ) ৪২৮ ✔}}

ঘ) ৪০৮

\textbf{\hl{✔ সঠিক উত্তর: গ) ৪২৮}}

\begin{quote}
\textbf{ব্যাখ্যা:} \emph{৬ × ৩ = ১৮ - ১ = ১৭ ১৭ × ৩ = ৫১ - ২ = ৪৯ ৪৯ × ৩ =
১৪৭ - ৩ = ১৪৪}

\emph{পরবর্তী পদ = ১৪৪ × ৩ = ৪৩২ - ৪ = ৪২৮}
\end{quote}

\textbf{১, ৩, ৬, ১০, ১৫, ২১,........... ধারাটির দশম পদ -}

ক) ৪৫

খ) ৫৫

\textbf{\hl{গ) ৬২ ✔}}

ঘ) ৬৫

\textbf{\hl{✔ সঠিক উত্তর: গ) ৬২}}

\begin{quote}
\textbf{ব্যাখ্যা:} \emph{এখানে, প্রতিটি সংখ্যার পার্থক্য হচ্ছে যথাক্রমে
২,৩,৪,৫,৬ অর্থাৎ পার্থক্য ১ করে বাড়ছে।তাই পরবর্তী পদ হবে ২১+৭ = ২৮।}
\end{quote}

\textbf{১,৫,১৩,২৯,৬১ ধারাটির পরবর্তী সংখ্যা কত ?}

ক) ৮০

খ) ১০০

গ) ১২০

\textbf{\hl{ঘ) ১২৫ ✔}}

\textbf{\hl{✔ সঠিক উত্তর: ঘ) ১২৫}}

\begin{quote}
\textbf{ব্যাখ্যা:} \emph{এখানে, প্রতি পদের পার্থক্য দ্বিগুণ হারে বাড়ছে।}

\emph{৫ = ১+৪, ১৩ = ৫+৮,২৯= ১৩+১৬, ৬১ = ২৯+৩২।}

\emph{পরবর্তী সংখ্যা = ৬১+৬৪ = ১২৫}
\end{quote}

\textbf{৪,৮,১৩,১৯,২৬.......ধারাটির ৭ম পদ কত ?}

ক) ৩৯

খ) ৪০

\textbf{\hl{গ) ৪৩ ✔}}

ঘ) ৪৮

\textbf{\hl{✔ সঠিক উত্তর: গ) ৪৩}}

\begin{quote}
\textbf{ব্যাখ্যা:} \emph{এখানে প্রতি পদের পার্থক্য ৪,৫,৬,৭,.....এভাবে বাড়ছে।}

\emph{তাই ধারাটির ৬ষ্ঠ পদ হবে ২৬+৮ = ৩৪,}

\emph{৭ম পদ হবে ৩৪+৯ = ৪৩ ।}
\end{quote}

\textbf{1, 1, 2, 3, 5, 8, 13, ..... ধারার ১০ম পদটি কত?}

ক) ৩৪

খ) ৪১

\textbf{\hl{গ) ৫৫ ✔}}

ঘ) ৮৯

\textbf{\hl{✔ সঠিক উত্তর: গ) ৫৫}}

\begin{quote}
\textbf{ব্যাখ্যা:} \emph{প্রদত্ত ধারাটি ফিবোনাচ্চি ধারার অন্তর্ভুক্ত। অর্থাৎ, 1 +
1 = 2 1 + 2 = 3 2 + 3 = 5 3 + 5 = 8 5 + 8 = 13 8 + 13 = 21 13 + 21 = 34
21 + 34 = 55}

\emph{∴ দশম পদ = (34 + 21) = 55}
\end{quote}

\textbf{১৩, ১৭, ২৫, ৪১ ...... ধারার পরবর্তী সংখ্যা কত?}

\textbf{\hl{ক) ৭৩ ✔}}

খ) ৮৯

গ) ১০১

ঘ) ১৪৩

\textbf{\hl{✔ সঠিক উত্তর: ক) ৭৩}}

\begin{quote}
\textbf{ব্যাখ্যা:} \emph{ধারাটির পদগুলোর মধ্যকার পার্থক্য লক্ষ্য করি:}

\emph{১৭ − ১৩ = ৪}

\emph{২৫ − ১৭ = ৮}

\emph{৪১ − ২৫ = ১৬}

\emph{পার্থক্যগুলো দ্বিগুণ হয়ে বাড়ছে (৪, ৮, ১৬, এরপরের পার্থক্য হবে ৩২)।}

\emph{∴ পরবর্তী সংখ্যাটি = ৪১ + ৩২ = ৭৩}
\end{quote}

\textbf{২, ৩, ৫, ৮, ১৩, ২১, ৩৪ ....ধারার পরের সংখ্যাটি কত?}

ক) ১৬

\textbf{\hl{খ) ৫৫ ✔}}

গ) ১৩

ঘ) ৩৫

\textbf{\hl{✔ সঠিক উত্তর: খ) ৫৫}}

\begin{quote}
\textbf{ব্যাখ্যা:} \emph{প্রদত্ত ধারাটি হচ্ছে ফিবোনাচ্চি ধারার অন্তর্ভুক্ত অর্থাৎ--
একটি সংখ্যা তার পূর্ববর্তী দুটি সংখ্যার যোগফলের সমান।}

\emph{⇒ অর্থাৎ পরবর্তী সংখ্যাটি = (৩৪ + ২১) = ৫৫।}
\end{quote}

\textbf{২, ৩, ৫, ৭, ১১, ৩১, ৭১, ৯১ ধারার পরের সংখ্যা কী?}

ক) ৯৩

খ) ১০৩

\textbf{\hl{গ) ৩২ ✔}}

ঘ) ৩১

\textbf{\hl{✔ সঠিক উত্তর: গ) ৩২}}

\begin{quote}
\textbf{ব্যাখ্যা:} \emph{এটি প্রাইম নম্বরের সিরিজ। ২, ৩, ৫, ৭, ১১, এই পর্যন্ত
ধারাটি ঠিক আছে। এর পরে হবে ১৩, ১৭, ১৯, ২৩। কিন্তু এটি উল্টে করে দেওয়া আছে
৩১, ৭১, ৯১....। সুতরাং পরের সংখ্যাটি ৩২।}
\end{quote}

\textbf{৩, ৫, ৮, ১০, ১৮, ২০ ধারার পরবর্তী সংখ্যা কত?}

ক) ২৮

খ) ৩০

গ) ৩৪

\textbf{\hl{ঘ) ৩৮ ✔}}

\textbf{\hl{✔ সঠিক উত্তর: ঘ) ৩৮}}

\begin{quote}
\textbf{ব্যাখ্যা:} \emph{দেওয়া আছে ৩, ৫, ৮, ১০, ১৮, ২০ ......}

\emph{এখানে, ৩ + ৫ = ৮; ৮ + ২ = ১০}

\emph{১০ + ৮ = ১৮; ১৮ + ২ = ২০ ২০ + ১৮ = ৩৮}
\end{quote}

\textbf{৩, ৬, ১১, ১৮, ২৭, ... ধারার পরের সংখ্যাটি কত?}

ক) ৩৬

\textbf{\hl{খ) ৩৮ ✔}}

গ) ৪৫

ঘ) ৪৮

ঙ) কোনটি নয়

\textbf{\hl{✔ সঠিক উত্তর: খ) ৩৮}}

\begin{quote}
\textbf{ব্যাখ্যা:} \emph{ধারাটি নিচের মতো করে এগিয়েছে।}

\emph{১ম সংখ্যা = ৩}

\emph{২য় সংখ্যা = ৩ + ৩ = ৬}

\emph{৩য় সংখ্যা = ৬ + ৫ = ১১}

\emph{৪র্থ সংখ্যা = ১১ + ৭ = ১৮}

\emph{৫ম সংখ্যা = ১৮ + ৯ = ২৭}

\emph{৬ষ্ঠ সংখ্যা = ২৭ + ১১ = ৩৮}

\emph{∴ সংখ্যাটি ৩৮}
\end{quote}

\textbf{৩, ৪, ৬, ৫, ৯, ৬, ১২, ৭..... ধারার পরবর্তী পদটি কত?}

ক) ৮

খ) ১০

গ) ১৪

\textbf{\hl{ঘ) ১৫ ✔}}

\textbf{\hl{✔ সঠিক উত্তর: ঘ) ১৫}}

\begin{quote}
\textbf{ব্যাখ্যা:} \emph{ধারাটিতে দুটি ধারা পর্যায়ক্রমে মিশে আছে:}

\emph{১ম ধারাটি (বেজোড় পদগুলো): ৩, ৬, ৯, ১২, ১৫ (প্রতিবারে ৩ করে বাড়ছে)}

\emph{২য় ধারাটি (জোড় পদগুলো): ৪, ৫, ৬, ৭ (প্রতিবারে ১ করে বাড়ছে)}

\emph{যেহেতু আমাদের পরবর্তী পদটি (৯ম পদ) বের করতে হবে, তাই তা ১ম ধারার
ধারাবাহিকতা অনুযায়ী ১২ এর পরে ১৫ হবে।}
\end{quote}

\textbf{এই ধারার পরবর্তী নম্বরটি কত হবে? ৩, ৬, ১৮, ৭২,....}

ক) ২২০

খ) ২৮৪

\textbf{\hl{গ) ৩৬০ ✔}}

ঘ) কোনোটাই নয়

\textbf{\hl{✔ সঠিক উত্তর: গ) ৩৬০}}

\begin{quote}
\textbf{ব্যাখ্যা:} \emph{ধারা: ৩, ৬, ১৮, ৭২}

\emph{এখানে, ৩ × ২ = ৬ ৬ × ৩ = ১৮ ১৮ × ৪ = ৭২ ৭২ × ৫ = ৩৬০}
\end{quote}

\textbf{একটি ধারার ১ ম সংখ্যাটি ৩, ২য় সংখ্যাটি ৯। ৩য় সংখ্যা থেকে শুরু করে
বাকি সংখ্যাগুলো পূর্ববর্তী সবগুলো সংখ্যার গড়। ধারার ১১তম সংখ্যাটি কত?}

ক) ৫৩

খ) ৪৮

গ) ৩

\textbf{\hl{ঘ) ৬ ✔}}

\textbf{\hl{✔ সঠিক উত্তর: ঘ) ৬}}

\begin{quote}
\textbf{ব্যাখ্যা:} \emph{৩য় সংখ্যাটি = (3+9/2) = ৬}

\emph{∴ ৪র্থ সংখ্যা = (3+9+6/3) = ৬}

\emph{∴ ৫ম সংখ্যা = (3+9+6+6/4) = ৬}

\emph{⇒ এভাবে দেখা যায় ৩য় পদ থেকে শুরু করে পরবর্তী সকল পদই ৬ হবে।}

\emph{∴ ১১তম সংখ্যাটি ৬}
\end{quote}

\textbf{১, ৯, ২৫, ৪৯, ৮১ ....... ধারার পরবর্তী সংখ্যাটি কত?}

ক) ১০০

\textbf{\hl{খ) ১২১ ✔}}

গ) ১৪৪

ঘ) ১৪৫

\textbf{\hl{✔ সঠিক উত্তর: খ) ১২১}}

\begin{quote}
\textbf{ব্যাখ্যা:} \emph{এখানে,}

\emph{১\^{}২ = ১;}

\emph{৩\^{}২ = ৯;}

\emph{৫\^{}২ = ২৫;}

\emph{৭\^{}২ = ৪৯;}

\emph{৯\^{}২ = ৮১;}

\emph{⇒ বেজোড় সংখ্যাগুলোর বর্গ ক্রমানুসারে আছে। পরবর্তী সংখ্যাটি হবে ১১\^{}২ =
১২১}
\end{quote}

\textbf{৩, ৭, ১৫, ৩১, ৬৩, ............ ধারাটির পরবর্তী সংখ্যাটি কত?}

ক) ১৪৫

খ) ১৩৫

\textbf{\hl{গ) ১২৭ ✔}}

ঘ) ১৩০

\textbf{\hl{✔ সঠিক উত্তর: গ) ১২৭}}

\begin{quote}
\textbf{ব্যাখ্যা:} \emph{ধারা: ৩, ৭, ১৫, ৩১, ৬৩, ১২৭}

\emph{অন্তর: ৪, ৮, ১৬, ৩২, ৬৪}

\emph{∴ পরবর্তী সংখ্যাটি ১২৭}
\end{quote}

\textbf{২, ৪, ৩, ৭, ৪, ১০, ৫ -\/-\/-\/-\/-\/- ধারাটির পরবর্তী পদ কি হবে?}

ক) ১০

খ) ১২

\textbf{\hl{গ) ১৩ ✔}}

ঘ) ১৫

\textbf{\hl{✔ সঠিক উত্তর: গ) ১৩}}

\begin{quote}
\textbf{ব্যাখ্যা:} \emph{এখানে দুটি ধারা রয়েছে, বিজোড় ক্রমে, ২, ৩, ৪, ৫ ও
জোড় ক্রমে ৪, ৭, ১০, ১৩। এভাবে বৃদ্ধি পাবে।}
\end{quote}

\textbf{২, ৩, ৫, ৮, ১৩, ২১, ৩৪, --- ধারার পরের সংখ্যাটি ---}

\textbf{\hl{ক) ৫৫ ✔}}

খ) ১৩

গ) ৩৫

ঘ) ১৬

\textbf{\hl{✔ সঠিক উত্তর: ক) ৫৫}}

\begin{quote}
\textbf{ব্যাখ্যা:} \emph{২, ৩, ৫, ৮, ১৩, ২১, ৩৪, .............}

\emph{এখানে, ২ + ৩ = ৫ ৩ + ৫ = ৮ ৫ + ৮ = ১৩}

\emph{⇒ অর্থাৎ ধারাটির ক্রমিক দুটি পদের সমষ্টি পরবর্তী পদের সমান। ∴ ২১ ও ৩৪
এর পরবর্তী সংখ্যা ২১ + ৩৪ = ৫৫}
\end{quote}

\textbf{০, ৩, ৮, ১৫, ২৪, .................... ধারা পরবর্তী পদ কোনটি?}

ক) ৫০

খ) ৪৫

গ) ৪০

\textbf{\hl{ঘ) ৩৫ ✔}}

\textbf{\hl{✔ সঠিক উত্তর: ঘ) ৩৫}}

\begin{quote}
\textbf{ব্যাখ্যা:} \emph{ধারা: ০, ৩, ৮, ১৫, ২৪, ৩৫}

\emph{অন্তর: ৩, ৫, ৭, ৯, ১১}

\emph{এখানে, ৩ + ২ = ৫ ৫ + ২ = ৭ ৭ + ২ = ৯ ৯ + ২ = ১১}

\emph{∴ পরবর্তী সংখ্যাটি = ২৪ + ১১ = ৩৫}
\end{quote}

\textbf{৯, ৩৬, ৮১, ১৪৪ ........ এর পরবর্তী সংখ্যা কত?}

ক) ১৫৯

\textbf{\hl{খ) ২২৫ ✔}}

গ) ২৫৬

ঘ) ২৭২

\textbf{\hl{✔ সঠিক উত্তর: খ) ২২৫}}

\begin{quote}
\textbf{ব্যাখ্যা:} \emph{৯, ৩৬, ৮১, ১৪৪, ২২৫}

\emph{এখানে সংখ্যাগুলো যথাক্রমে ৩\^{}২ = ৯, ৬\^{}২ = ৩৬, ৯\^{}২ = ৮১,
১২\^{}২ = ১৪৪ (৩ করে ভিত্তি বাড়ছে)।}

\emph{∴ পরবর্তী সংখ্যাটি হবে ১৫\^{}২ = ২২৫}
\end{quote}

\textbf{১, ২, ৩, ৫, ৮, ১৩, ২১, ৩৪ ............ ধারার পরবর্তী সংখ্যা কত?}

\textbf{\hl{ক) ৫৫ ✔}}

খ) ৪০

গ) ৬৮

ঘ) ৬০

\textbf{\hl{✔ সঠিক উত্তর: ক) ৫৫}}

\begin{quote}
\textbf{ব্যাখ্যা:} \emph{ধারার পরবর্তী সংখ্যা পূর্ববর্তী সংখ্যা দুটিম যোগফল। ১ +
২ = ৩ ২ + ৩ = ৫ ৩ + ৫ = ৮ ৫ + ৮ = ১৩ ৮ + ১৩ = ২১ ১৩ + ২১ = ৩৪ ২১ + ৩৪ =
৫৫}
\end{quote}

\textbf{২, ৬, ১০, ১৪... ধারাটির পরবর্তী ৭ম পদ কত?}

ক) ২২

\textbf{\hl{খ) ২৬ ✔}}

গ) ২৮

ঘ) ৩০

\textbf{\hl{✔ সঠিক উত্তর: খ) ২৬}}

\begin{quote}
\textbf{ব্যাখ্যা:} \emph{২, ৬, ১০, ১৪, ..... ধারার ৭ম পদ?}

\emph{এখানে, প্রতি জোড়া পদের ব্যবধান হলো ৪।}

\emph{∴ পঞ্চম পদ = ১৮; ষষ্ঠ পদ = ২২; সপ্তম পদ = ২৬}
\end{quote}

\textbf{১৬, ২২, ৩৪, ৫৮, ১০৬ ধারার পরের পদটি কত?}

ক) ১৫৪

খ) ১৮৪

গ) ১৯২

\textbf{\hl{ঘ) ২০২ ✔}}

\textbf{\hl{✔ সঠিক উত্তর: ঘ) ২০২}}

\begin{quote}
\textbf{ব্যাখ্যা:} \emph{এখানে, ১ম পদ = ১৬}

\emph{২য় পদ = ১৬ + ৬ × ১ = ২২}

\emph{৩য় পদ = ২২ + ৬ × ২ = ৩৪}

\emph{৪র্থ পদ = ৩৪ + ৬ × ৪ = ৫৮}

\emph{৫ম পদ = ৫৮ + ৬ × ৮ = ১০৬}

\emph{এবং ষষ্ঠ পদ = ১০৬ + ৬ × ১৬ = ২০২}
\end{quote}

\textbf{৫, ৭, ১১, ১৯,.... এ ধারায় ১৯ এর পরের সংখ্যাটি কত?}

ক) ৪২

\textbf{\hl{খ) ৩৫ ✔}}

গ) ৩০

ঘ) ২৬

\textbf{\hl{✔ সঠিক উত্তর: খ) ৩৫}}

\begin{quote}
\textbf{ব্যাখ্যা:} \emph{এখানে, ১ম পদ = ৫}

\emph{২য় এখানে, = (১ম পদ) ৫ + ২ = ৭}

\emph{৩য় এখানে, = (২য় পদ) ৭ + ৪ = ১১}

\emph{৪র্থ এখানে, = (৩য় পদ) ১১ + ৮ = ১৯}

\emph{৫ম এখানে, = (৪র্থ পদ) ১৯ + ১৬ = ৩৫}
\end{quote}

\textbf{১০, ২২, ৪৬, ৯৪, ...........?}

ক) ১৮০

\textbf{\hl{খ) ১৯০ ✔}}

গ) ১৯২

ঘ) ১৯৪

\textbf{\hl{✔ সঠিক উত্তর: খ) ১৯০}}

\begin{quote}
\textbf{ব্যাখ্যা:} \emph{ধারাটি: ১০, ২২, ৪৬, ৯৪}

\emph{অন্তর: ১২, ২৪, ৪৮, ৯৬}

\emph{এখানে, ১ম পদ ও ২য় পদের অন্তর অপেক্ষা ২য় পদ ও ৩য় পদের অন্তর দ্বিগুণ।}

\emph{∴ নির্ণেয় পদ = (৯৪ + ৯৬) = ১৯০}
\end{quote}

\textbf{১ + ৫ + ৯ + ১৩ + ................ এই ধারার পরবর্তী সংখ্যা কোনটি?}

\textbf{\hl{ক) ১৭ ✔}}

খ) ১৮

গ) ১৯

ঘ) ২০

\textbf{\hl{✔ সঠিক উত্তর: ক) ১৭}}

\begin{quote}
\textbf{ব্যাখ্যা:} \emph{১ + ৪ = ৫, ৫ + ৪ = ৯, ৯ + ৪ = ১৩}

\emph{∴ পরবর্তী সংখ্যা ১৩ + ৪ = ১৭ (উত্তর)}
\end{quote}

\textbf{৮, ১১, ১৭, ২৯, ৫৩ ------ পরবর্তী সংখ্যাটি কত?}

\textbf{\hl{ক) ১০১ ✔}}

খ) ১০২

গ) ৭৫

ঘ) ৯১

\textbf{\hl{✔ সঠিক উত্তর: ক) ১০১}}

\begin{quote}
\textbf{ব্যাখ্যা:} \emph{লক্ষণীয়, ২য় পদ − ১ম পদ = ১১ − ৮ = ৩}

\emph{৩য় পদ − ২য় পদ = ১৭ − ১১ = ৬}

\emph{৪র্থ পদ − ৩য় পদ = ২৯ − ১৭ = ১২}

\emph{৫ম পদ − ৪র্থ পদ = ৫৩ − ২৯ = ২৪}

\emph{ষষ্ঠ পদ − ৫ম পদ = ৪৮}

\emph{∴ পরবর্তী সংখ্যা = ৫৩ + ৪৮ = ১০১}
\end{quote}

\textbf{(1/3)+(1/3\^{}2)+(1/3\^{}3)\ldots{} অনন্ত ধারার মান কত হবে?}

ক) (4/3)

\textbf{\hl{খ) (1/2) ✔}}

গ) (2/3)

ঘ) (3/2)

\textbf{\hl{✔ সঠিক উত্তর: খ) (1/2)}}

\begin{quote}
\textbf{ব্যাখ্যা:} \emph{এখানে, ধারাটির ১ ম পদ a = (1/3)}

\emph{সাধারণ অনুপাত r = ((1/3\^{}2)/(1/3)) = (1/9) × (3/1) = (1/3)
\textless{} 1}

\emph{∴ ধারাটির অসীমক সমষ্টি S = (a/1−r)}

\emph{= ((1/3)/1−(1/3)) = ((1/3)/(2/3)) = (1/3) × (3/2) = (1/2)}
\end{quote}

\textbf{শূন্যস্থানে কোন সংখ্যাটি বসবে? (2/3) , (4/7) , \_\_, (11/21) ,
(16/31)}

ক) (২২/৩৩)

\textbf{\hl{খ) (৭/১৩) ✔}}

গ) (৩৭/৩৯)

ঘ) (২৯/৪৫)

\textbf{\hl{✔ সঠিক উত্তর: খ) (৭/১৩)}}

\begin{quote}
\textbf{ব্যাখ্যা:} \emph{এখানে, লক্ষণীয় যে---}

\emph{১ম পদ = (2/3) = (2/2×2−1)}

\emph{২য় পদ = (4/7) = (4/4×2−1)}

\emph{৪র্থ পদ = (11/21) = (11/11×2−1)}

\emph{৫ম পদ = (16/31) = (16/16×2−1)}

\emph{⇒ (7/13) = (7/7×2−1) ~শূন্যস্থানের সংখ্যাটি (7/13)}
\end{quote}

\textbf{সিরিজের পরের সংখ্যাটি কত হবে? ২, ৫, ১১, ২৩, ---}

ক) ৩৫

খ) ৪৫

\textbf{\hl{গ) ৪৭ ✔}}

ঘ) কোনোটাই নয়

\textbf{\hl{✔ সঠিক উত্তর: গ) ৪৭}}

\begin{quote}
\textbf{ব্যাখ্যা:} \emph{২ + ৩ = ৫, ৫ + ৬ = ১১, ১১ + ১২ = ২৩, ২৩ + ২৪ =
৪৭}
\end{quote}

\textbf{একটি ধারার n-তম পদ m\^{}2n − 5 ধারার ২য় পদ 76 হলে m- এর মান কত?}

\textbf{\hl{ক) ৩ ✔}}

খ) ৪

গ) ৫

ঘ) ৬

\textbf{\hl{✔ সঠিক উত্তর: ক) ৩}}

\begin{quote}
\textbf{ব্যাখ্যা:} \emph{ধারাটির n-তম পদ = m\^{}2n − 5}

\emph{∴ ২য় পদ = m\^{}2×2 − 5 = m\^{}4 − 5}

\emph{শর্তমতে, m\^{}4 − 5 = 76 বা, m\^{}4 = 76 + 5 বা, m\^{}4 = 81 বা,
m\^{}4 = 3\^{}4 ∴ m = 3}
\end{quote}

\textbf{২, ৪, ৩, ৭, ৪, ১০, ৫, ............................. ধারাটির দশম
পদ কি হবে?}

ক) ১০

খ) ১২

\textbf{\hl{গ) ১৩ ✔}}

ঘ) ১৫

ঙ) ১৬

\textbf{\hl{✔ সঠিক উত্তর: গ) ১৩}}

\begin{quote}
\textbf{ব্যাখ্যা:} \emph{ধারাটির বিজোড় স্থানের পদগুলি হচ্ছে ২, ৩, ৪, ৫,
........}

\emph{এবং ধারাটির জোড় স্থানের পদগুলি হচ্ছে ৪, ৭, ১০, ১৩...}

\emph{∴ ধারাটির ৮ম পদ = ১০ + ৩ = ১৩}
\end{quote}

\textbf{১, ২, ৩, ৫, ......... ধারারটির দশম সংখ্যাটি কত?}

ক) ৩৪

\textbf{\hl{খ) ৫৫ ✔}}

গ) ২১

ঘ) ১৩

\textbf{\hl{✔ সঠিক উত্তর: খ) ৫৫}}

\begin{quote}
\textbf{ব্যাখ্যা:} \emph{১+১ = ২ ১+২ = ৩ ২+৩ = ৫ ৩+৫ = ৮ ৫+৮ = ১৩ ৮+১৩ =
২১ ১৩+২১ = ৩৪ ২১+৩৪ = ৫৫}

\emph{∴ দশম পদ = ২১+৩৪ = ৫৫}
\end{quote}

\textbf{১, ১, ২, ৩, ৫, ৮, ১৩, ২১ ... ধারার ১০ম পদ কত?}

ক) ৩৪

\textbf{\hl{খ) ৫৫ ✔}}

গ) ৪৮

ঘ) ৬৪

\textbf{\hl{✔ সঠিক উত্তর: খ) ৫৫}}

\begin{quote}
\textbf{ব্যাখ্যা:} \emph{১+১ = ২ ১+২ = ৩ ২+৩ = ৫ ৩+৫ = ৮ ৫+৮ = ১৩ ৮+১৩ =
২১ ১৩+২১ = ৩৪ ২১+৩৪ = ৫৫}

\emph{∴ ১০ম পদ = ৫৫}
\end{quote}

\textbf{১, ৩, ৬, ১০, ১৫, ২১, ............ ধারারটির দশম পদ কত?}

ক) ৪৫

\textbf{\hl{খ) ৫৫ ✔}}

গ) ৬২

ঘ) ৭২

\textbf{\hl{✔ সঠিক উত্তর: খ) ৫৫}}

\begin{quote}
\textbf{ব্যাখ্যা:} \emph{৬ষ্ঠ পদ = (৫ম পদ)+৬ = ২১}

\emph{৭ম পদ = (৬ষ্ঠ পদ)+৭ = ২৮}

\emph{৮ম পদ = (৭ম পদ)+৮ = ৩৬}

\emph{৯ম পদ = (৮ম পদ)+৯ = ৪৫}

\emph{এবং ১০ম পদ = (৯ম পদ)+১০ = ৫৫}
\end{quote}

\textbf{২, ৩, ৫, ৯, ১৭ এর পরবর্তী সংখ্যাটি কত?}

\textbf{\hl{ক) ৩৩ ✔}}

খ) ৩০

গ) ৩২

ঘ) ৪৮

\textbf{\hl{✔ সঠিক উত্তর: ক) ৩৩}}

\begin{quote}
\textbf{ব্যাখ্যা:} \emph{৩ − ২ = ১}

\emph{৫ − ৩ = ২ = ১ × ২}

\emph{৯ − ৫ = ৪ = ২ × ২}

\emph{১৭ − ৯ = ৮ = ৪ × ২}

\emph{∴ পরবর্তী সংখ্যা = ১৭ + ৮ × ২ = ৩৩}
\end{quote}

\textbf{11, 15, 23, 39.... ক্রমটির পরবর্তী সংখ্যা কত?}

ক) 45

খ) 55

\textbf{\hl{গ) 71 ✔}}

ঘ) 81

\textbf{\hl{✔ সঠিক উত্তর: গ) 71}}

\begin{quote}
\textbf{ব্যাখ্যা:} \emph{ধারা: 11, 15, 23, 39, 71}

\emph{অন্তর: 4, 8, 16, 32}

\emph{∴ ধারারটির পরবর্তী পদ = 39+32 = 71}
\end{quote}

\textbf{১৪৪, ৮১, ৩৬-এর পদক্রমটির পরবর্তী পদ কত?}

ক) ১৮

খ) ১২

\textbf{\hl{গ) ৯ ✔}}

ঘ) ৬

\textbf{\hl{✔ সঠিক উত্তর: গ) ৯}}

\begin{quote}
\textbf{ব্যাখ্যা:} \emph{ধারা: ১৪৪, ৮১, ৩৬, ৯ = ১২\^{}২, ৯\^{}২, ৬\^{}২,
৩\^{}২}
\end{quote}

\textbf{৩, ৬, ৯, ...-এই পদক্রমটির পরবর্তী পদ কত?}

ক) ৭

খ) ৮

গ) ১০

\textbf{\hl{ঘ) ১২ ✔}}

\textbf{\hl{✔ সঠিক উত্তর: ঘ) ১২}}

\begin{quote}
\textbf{ব্যাখ্যা:} \emph{ধারা: ৩, ৬, ৯, .....}

\emph{ধারা হতে পাই, ৩ × ১ = ৩ ৩ × ২ = ৬ ৩ × ৩ = ৯ ৩ × ৪ = ১২}
\end{quote}

\textbf{২, ৩, ৫, ৯, ১৭ ধারার পরবর্তী সংখ্যাটি কত?}

\textbf{\hl{ক) ৩৩ ✔}}

খ) ৩২

গ) ৩০

ঘ) ৪৬

\textbf{\hl{✔ সঠিক উত্তর: ক) ৩৩}}

\begin{quote}
\textbf{ব্যাখ্যা:} \emph{ধারা: ২, ৩, ৫, ৯, ১৭, ৩৩}

\emph{অন্তর: ১, ২, ৪, ৮, ১৬}

\emph{∴ ধারার পরবর্তী সংখ্যা (১৭ + ১৬) = ৩৩}
\end{quote}

\textbf{১, ২, ৩, ৫, ৮, ১৩, ২১, ৩৪, --- ধারার পরবর্তী সংখ্যাটি কত?}

\textbf{\hl{ক) ৫৫ ✔}}

খ) ৪০

গ) ৬৮

ঘ) ৮৯

\textbf{\hl{✔ সঠিক উত্তর: ক) ৫৫}}

\begin{quote}
\textbf{ব্যাখ্যা:} \emph{এখানে লক্ষণীয়,}

\emph{৩য় পদ = ১ম পদ + ২য় পদ = ১ + ২ = ৩}

\emph{৪র্থ পদ = ২য় পদ + ৩য় পদ = ২ + ৩ = ৫}

\emph{৯ম পদ = ৮ম পদ + ৭ম পদ = ৩৪ + ২১ = ৫৫}
\end{quote}

\textbf{৭, ১০, ১৬, ২৮, ৫২,.......ধারারটি পরবর্তী সংখ্যা কত?}

ক) ৭৫

\textbf{\hl{খ) ১০০ ✔}}

গ) ১০৫

ঘ) ১৫০

\textbf{\hl{✔ সঠিক উত্তর: খ) ১০০}}

\begin{quote}
\textbf{ব্যাখ্যা:} \emph{লক্ষণীয়, ২য় পদ − ১ম পদ = ১০ − ৭ = ৩}

\emph{৩য় পদ − ২য় পদ = ১৬ − ১০ = ৬ = ৩ × ২}

\emph{৪র্থ পদ − ৩য় পদ = ২৮ − ১৬ = ১২ = ৬ × ২}

\emph{৫ম পদ − ৪র্থ পদ = ৫২ − ২৮ = ২৪ = ১২ × ২}

\emph{ষষ্ঠ পদ − ৫ম পদ = ২৪ × ২ = ৪৮}

\emph{∴ ষষ্ঠ পদ = ৫ম পদ + ৪৮ = ৫২ + ৪৮ = ১০০}
\end{quote}

\textbf{What is the next number of the series, 1, 4, 9, 16, .. ? / ১, ৪,
৯, ১৬, ... ধারার পরবর্তী পদ কত?}

ক) ১৮

খ) ২১

গ) ২৩

\textbf{\hl{ঘ) ২৫ ✔}}

ঙ) ৩১

\textbf{\hl{✔ সঠিক উত্তর: ঘ) ২৫}}

\begin{quote}
\textbf{ব্যাখ্যা:} \emph{১, ৪, ৯, ১৬ ...... = ১\^{}২, ২\^{}২, ৩\^{}২,
৪\^{}২, ........}

\emph{নির্ণেয় সংখ্যা = ৫\^{}২ = ২৫}
\end{quote}

\textbf{২৩, ২৫, ২৯, ৩৭, -\/-\/-\/- শূন্যস্থানে কত বসবে?}

ক) ৪২

খ) ৪৯

\textbf{\hl{গ) ৫৩ ✔}}

ঘ) ৫৭

\textbf{\hl{✔ সঠিক উত্তর: গ) ৫৩}}

\begin{quote}
\textbf{ব্যাখ্যা:} \emph{২৩, ২৫, ২৯, ৩৭, ...................}

\emph{এখানে, ২৫ − ২৩ = ২ = ২\^{}১}

\emph{২৯ − ২৫ = ৪ = ২\^{}২}

\emph{৩৭ − ২৯ = ৮ = ২\^{}৩}

\emph{৫৩ − ৩৭ = ১৬ = ২\^{}৪}

\emph{∴ পরবর্তী সংখ্যাটি হবে ৩৭ + ১৬ = ৫৩}
\end{quote}

\textbf{১, ২৭, ১২৫, ........... শূন্যস্থানে কোন সংখ্যা বসবে?}

ক) ১৬৯

খ) ২১৬

গ) ২৮৯

\textbf{\hl{ঘ) ৩৪৩ ✔}}

\textbf{\hl{✔ সঠিক উত্তর: ঘ) ৩৪৩}}

\begin{quote}
\textbf{ব্যাখ্যা:} \emph{১, ২৭, ১২৫, ...................}

\emph{বা, ১\^{}৩, ৩\^{}৩, ৫\^{}৩, ৭\^{}৩, ...................}

\emph{পরবর্তী সংখ্যা হবে বিজোড় সংখ্যা ৭ এর ঘন। ∴ ৭\^{}৩ = ৩৪৩}
\end{quote}

\textbf{শূন্যস্থানে কোন সংখ্যাটি বসবে? ২১ ২, ১৭৯, ১৪৬, ১১৩ --- ।}

ক) ১৩

খ) ৪০

গ) ৩৩

\textbf{\hl{ঘ) ৮০ ✔}}

\textbf{\hl{✔ সঠিক উত্তর: ঘ) ৮০}}

\begin{quote}
\textbf{ব্যাখ্যা:} \emph{ধারা: ২১২, ১৭৯, ১৪৬, ১১৩, ৮০}

\emph{অন্তর: ৩৩, ৩৩, ৩৩, ৩৩}

\emph{∴ শূন্যস্থানের সংখ্যাটি হবে = (১১৩ − ৩৩) = ৮০}

\emph{∴ সঠিক উত্তর হবে ৮০।}
\end{quote}

\textbf{১, ৯, ২৫, ৪৯, ৮১....... ধারারটির পরবর্তী সংখ্যা কত?}

ক) ১০০

\textbf{\hl{খ) ১২১ ✔}}

গ) ১৪৪

ঘ) ওপরের একটিও নয়

\textbf{\hl{✔ সঠিক উত্তর: খ) ১২১}}

\begin{quote}
\textbf{ব্যাখ্যা:} \emph{১, ৯, ২৫, ৪৯, ৮১.......}

\emph{= ১\^{}২, ৩\^{}২, ৫\^{}২, ৭\^{}২, ৯\^{}২, .......}

\emph{নির্ণেয় সংখ্যা = ১১\^{}২ = ১২১}
\end{quote}

\textbf{৯, ৩৬, ৮১, ১৪৪, ... । পরবর্তী সংখ্যাটি কত?}

ক) ১৬৯

\textbf{\hl{খ) ২২৫ ✔}}

গ) ২৫৬

ঘ) ২৭২

\textbf{\hl{✔ সঠিক উত্তর: খ) ২২৫}}

\begin{quote}
\textbf{ব্যাখ্যা:} \emph{৯, ৩৬, ৮১, ১৪৪, ২২৫}

\emph{এখানে সংখ্যাগুলো যথাক্রমে ৩\^{}২, ৬\^{}২, ৯\^{}২, ১২\^{}২ (৩ করে বৃদ্ধি
পাচ্ছে)।}

\emph{∴ পরবর্তী সংখ্যাটি হবে ১৫\^{}২ = ২২৫}
\end{quote}

\textbf{৮১, ২৭, ....., ৩, ১; লুপ্ত সংখ্যাটি কত?}

ক) ৬

\textbf{\hl{খ) ৯ ✔}}

গ) ১২

ঘ) ১৫

\textbf{\hl{✔ সঠিক উত্তর: খ) ৯}}

\begin{quote}
\textbf{ব্যাখ্যা:} \emph{ধারা = ৮১, ২৭, ৯, ৩, ১ = ৩\^{}৪, ৩\^{}৩, ৩\^{}২,
৩\^{}১, ৩\^{}০}

\emph{নির্ণেয় সংখ্যা = ৩\^{}২ = ৯}
\end{quote}

\textbf{১, ২, ৩, ৫, ৮, ১৩, ২১, ৩৪, --- ধারার পরবর্তী সংখ্যাটি কত?}

\textbf{\hl{ক) ৫৫ ✔}}

খ) ৪০

গ) ৬৮

ঘ) ৮৯

\textbf{\hl{✔ সঠিক উত্তর: ক) ৫৫}}

\begin{quote}
\textbf{ব্যাখ্যা:} \emph{এখানে লক্ষণীয়,}

\emph{৩য় পদ = ১ম পদ + ২য় পদ = ১ + ২ = ৩}

\emph{৪র্থ পদ = ২য় পদ + ৩য় পদ = ২ + ৩ = ৫}

\emph{৯ম পদ = ৮ম পদ + ৭ম পদ = ৩৪ + ২১ = ৫৫}
\end{quote}

\textbf{১, ১, ২, ৩, ৫, ৮, ......, এই সংখ্যা পরম্পরায় অষ্টম পদ কত? / Find
the 8th term of the series 1, 1, 2, 3, 5, 8, .....}

\textbf{\hl{ক) ২১ ✔}}

খ) ১৩

গ) ১৯

ঘ) ১৬

\textbf{\hl{✔ সঠিক উত্তর: ক) ২১}}

\begin{quote}
\textbf{ব্যাখ্যা:} \emph{এখানে লক্ষণীয়,}

\emph{৩য় পদ = ১ম পদ + ২য় পদ = ১ + ১ = ২}

\emph{৪র্থ পদ = ২য় পদ + ৩য় পদ = ১ + ২ = ৩}

\emph{ষষ্ঠ পদ = ৫ম পদ + ৪র্থ পদ = ৫ + ৩ = ৮}

\emph{৭ম পদ = ষষ্ঠ পদ + ৫ম পদ = ৮ + ৫ = ১৩}

\emph{৮ম পদ = ৭ম পদ + ষষ্ঠ পদ = ১৩ + ৮ = ২১}
\end{quote}

\textbf{8, 11, 17, 29, 53, ...... \textbar{} পরবর্তী সংখ্যাটি কত? / 8, 11,
17, 29, 53, ...... \textbar{} Find the next term.}

\textbf{\hl{ক) 101 ✔}}

খ) 102

গ) 75

ঘ) 59

\textbf{\hl{✔ সঠিক উত্তর: ক) 101}}

\begin{quote}
\textbf{ব্যাখ্যা:} \emph{লক্ষণীয়, ২য় পদ − ১ম পদ = ১১ − ৮ = ৩}

\emph{৩য় পদ − ২য় পদ = ১৭ − ১১ = ৬}

\emph{৪র্থ পদ − ৩য় পদ = ২৯ − ১৭ = ১২}

\emph{৫ম পদ − ৪র্থ পদ = ৫৩ − ২৯ = ২৪}

\emph{ষষ্ঠ পদ − ৫ম পদ = ৪৮}

\emph{∴ ষষ্ট পদ = ৫ম পদ + ৪৮ = ৫৩ + ৪৮ = ১০১}
\end{quote}

\textbf{১৯, ৩৩, ৫১, ৭৩, --- \textbar{} পরবর্তী সংখ্যাটি কত? / 19, 33, 51,
73,---. Find the next term.}

ক) ৮৫

খ) ১২১

\textbf{\hl{গ) ৯৯ ✔}}

ঘ) ৯৮

\textbf{\hl{✔ সঠিক উত্তর: গ) ৯৯}}

\begin{quote}
\textbf{ব্যাখ্যা:} \emph{১৯, ৩৩, ৫১, ৭৩, ৯৯}

\emph{অন্তর: ১৪, ১৮, ২২, ২৬}

\emph{এখানে, অন্তরের পার্থক্য প্রতিবারে ৪ করে বাড়ছে (২৬ এর পরে হবে ৩০)।}

\emph{∴ পরবর্তী সংখ্যাটি = ৭৩ + ২৬ = ৯৯}
\end{quote}

\textbf{বিশেষ ক্রমানুসারে সাজালে ১৩, ১৭, ২৫, ৪১, এর পরবর্তী সংখ্যা কি?}

ক) ৯০

খ) ৬২

গ) ৬

\textbf{\hl{ঘ) ৭৩ ✔}}

\textbf{\hl{✔ সঠিক উত্তর: ঘ) ৭৩}}

\begin{quote}
\textbf{ব্যাখ্যা:} \emph{সমাধান: ২য় পদ = ১ম পদ + ৮ = ১৩ + ৮ = ১৭}

\emph{৩য় পদ = ২য় পদ + ৮ = ১৭ + ৮ = ২৫}

\emph{৪র্থ পদ = ৩য় পদ + ১৬ = ২৫ + ১৬ = ৪১}

\emph{৫ম পদ = ৪র্থ পদ + ৩২ = ৪১ + ৩২ = ৭৩}
\end{quote}

\textbf{৬, ১৯, ৪৯, ১৮৮ ক্রমটির পরবর্তী পদ কত?}

ক) ২০

খ) ৩৫৬

গ) ৪০৮

\textbf{\hl{ঘ) ৮২৮ ✔}}

\textbf{\hl{✔ সঠিক উত্তর: ঘ) ৮২৮}}

\begin{quote}
\textbf{ব্যাখ্যা:} \emph{সমাধান: ৬ × ৩ − ১ = ১৯}

\emph{১৯ × ৩ − ২ = ৪৯}

\emph{৪৯ × ৩ − ৩ = ১৮৮}

\emph{১৮৮ × ৩ − ৪ = ৮২৮}
\end{quote}

\textbf{(1/2)−(1/3)+(1/4)−(1/5)+\ldots{} ধারার সাধারণ পদ (
n∈N,n\textgreater1 )}

ক) ((−1)\^{}n/n+1)

খ) ((−1)\^{}n−1/n+1)

\textbf{\hl{গ) ((−1)\^{}(n+1)/n+1) ✔}}

ঘ) ((−1)\^{}−n−1/n+1)

\textbf{\hl{✔ সঠিক উত্তর: গ) ((−1)\^{}(n+1)/n+1)}}

\begin{quote}
\textbf{ব্যাখ্যা:} \emph{সমাধান: অপশন গ এ প্রদত্ত সাধারণ পদটি হলো
((−1)\^{}(n+1)/n+1)}

\emph{এখন, n = 3 হলে পাই, ((−1)\^{}(3+1)/3+1) = ((−1)\^{}4/4) = (1/4)}

\emph{⇒ যা ধারাটির তৃতীয় পদ (1/4) -এর সাথে মিলে যায়। ∴ সঠিক উত্তর গ।}
\end{quote}

\section{3. অনুপাত ও সমানুপাত-
পাটিগণিত}\label{ux985ux9a8ux9aaux9a4-ux993-ux9b8ux9aeux9a8ux9aaux9a4--ux9aaux99fux997ux9a3ux9a4}

বিসিএস, শিক্ষক নিবন্ধন কিংবা ব্যাংকসহ যেকোনো প্রতিযোগিতামূলক পরীক্ষায় গণিতের
`অনুপাত ও সমানুপাত' (Ratio \& Proportion) অধ্যায়টি থেকে প্রশ্ন আসবেই! দৈনন্দিন
জীবনের হিসাব-নিকাশ ছাড়াও যেকোনো চাকরির পরীক্ষায় ১-২ নম্বর নিশ্চিত করতে এই
অধ্যায়ের গুরুত্ব অপরিসীম। তাই মূল কনসেপ্ট পরিষ্কার করা, দ্রুত অংক করার শর্টকাট
কৌশল এবং বিগত সালের প্রশ্ন সমাধানের পাশাপাশি এই আর্টিকেলে যুক্ত করা হয়েছে
`লাইভ এক্সাম মোড' (Live Exam Mode)। এখন আপনি শুধু নিয়ম শিখবেনই না, বরং নিজে
নিজে কুইজ টেস্ট দিয়ে তাৎক্ষণিক নিজের প্রস্তুতিও যাচাই করতে পারবেন!

\begin{itemize}
\item
  প্রাথমিক আলোচনা
\item
  কিছু সহজ কৌশল
\item
  সমানুপাত বের করা
\item
  ভগ্নাংশের অনুপাত
\end{itemize}

\subsubsection{প্রাথমিক
আলোচনা}\label{ux9aaux9b0ux9a5ux9aeux995-ux986ux9b2ux99aux9a8}

\textbf{Ratio বা অনুপাত কী?}

অনুপাত হচ্ছে এক বা একাধিক রাশির তুলনা যাকে ( : ) চিহ্ন দিয়ে প্রকাশ করা হয় যা
একটি ভগ্নাংশকে নির্দেশ করে। যেমন- ৩:৭ = \(\frac{\text{৩}}{\text{৭}}\) ।

\textbf{অনুপাত সম্পর্কিত কিছু তথ্য:}

অনুপাত হচ্ছে একটি ভগ্নাংশ যাতে প্রথম রাশি লব এবং দ্বিতীয় রাশি হর ।

অনুপাতকে সবসময় ক্ষুদ্রতম আকারে প্রকাশ করতে হয় । অর্থাৎ ৮:৬ না লিখে লিখতে হয়
৪:৩ ।

অনুপাতের তুলনায় যে রাশি প্রথমে তার মান ও প্রথমেই বসাতে হয় ।

অনুপাতে প্রকাশ করার জন্য কমপক্ষে দুটি রাশির মানের প্রয়োজন ।

অনুপাতের নির্দিষ্ট কোন একক নেই। অর্থাৎ ৫:৭ = ৫ কেজি ও ৭কেজি বলা যাবে না।বরং
যার অনুপাত প্রকাশ করবে তখন সেই এককই ব্যবহৃত হবে।

\textbf{বিভিন্ন প্রকার অনুপাতঃ}

\textbf{বহুরাশিক অনুপাত:} মনে করি, একটি বাক্সের দৈর্ঘ্য, প্রস্থ ও উচ্চতা যথাক্রমে
৮ সে.মি., ৫ সে.মি. ও ৬ সে.মি.

দৈর্ঘ্য, প্রস্থ ও উচ্চতার অনুপাত = ৮ : ৫ : ৬

সংক্ষেপে, দৈর্ঘ্য : প্রস্থ : উচ্চতা = ৮ : ৫ : ৬

এখানে তিনটি রাশির অনুপাত উপস্থাপন করা হয়েছে। এরূপ তিন বা ততোধিক রাশির
অনুপাতকে \textbf{বহুরাশিক অনুপাত} বলে।

\textbf{ধারাবাহিক অনুপাত :} মনে করি, পুত্র ও পিতার বয়সের অনুপাত = ১৫ : ৪১
(পূর্ব রাশি: উত্তর রাশি) এবং পিতা ও দাদার বয়সের অনুপাত = ৪১: ৬৫ ।

দুইটি অনুপাতকে একত্র করে পাই, পুত্রের বয়স : পিতার বয়স : দাদার বয়স = ১৫ : ৪১ :
৬৫ । এ ধরনের অনুপাতকে \textbf{ধারাবাহিক অনুপাত} বলে।

এখানে লক্ষণীয় যে, প্রথম অনুপাতের উত্তর রাশি ও দ্বিতীয় অনুপাতের পূর্ব রাশি সমান।
প্রথম অনুপাতের উত্তর রাশি ও দ্বিতীয় অনুপাতের পূর্ব রাশি সমান না হলে তাদেরকে
সমান করে ধারাবাহিক অনুপাত বের করতে হয়।

\textbf{দুইটি অনুপাতকে ধারাবাহিক অনুপাতে রূপান্তরের জন্য প্রথম অনুপাতের উত্তর
রাশি দ্বারা দ্বিতীয় অনুপাতের উভয় রাশিকে গুণ করতে হবে এবং দ্বিতীয় অনুপাতের পূর্ব
রাশি দ্বারা প্রথম অনুপাতের উভয় রাশিকে গুণ করতে হবে।}

উদাহরণস্বরুপ : মনে করি, পুত্র ও পিতার বয়সের অনুপাত = ৩ : ৭ এবং পিতা ও দাদার
বয়সের অনুপাত = ৫: ৮ ।

এখন দুইটি অনুপাতকে ধারাবাহিক অনুপাতে রূপান্তরের জন্য প্রথম অনুপাতের উত্তর রাশি (
৭ ) দ্বারা দ্বিতীয় অনুপাতের উভয় রাশিকে গুণ করতে হবে ।এদের গুণফল হয় = ৫×৭ :
৮×৭ = ৩৫ : ৫৬ । আবার দ্বিতীয় অনুপাতের পূর্ব রাশি ( ৫ ) দ্বারা প্রথম অনুপাতের
উভয় রাশিকে গুণ করতে হবে। এদের গুণফল হয় = ৩×৫ : ৭×৫ = ১৫ : ৩৫ । এক্ষেত্রে
দেখা যাচ্ছে পিতার বয়সের অনুপাত উভয়ক্ষেত্রেই ৩৫ । তাই দুইটি অনুপাতকে একত্র করে
পাই, পুত্রের বয়স : পিতার বয়স : দাদার বয়স = ১৫ : ৩৫ : ৫৬ । এটিই ধারাবাহিক
অনুপাত।

\textbf{ব্যস্ত বা বিপরীত অনুপাত:} কোন অনুপাতের পূর্ব ও উত্তরঃ রাশিকে যথাক্রমে
উত্তর ও পূর্ব রাশি ধরে যে অনুপাত হয়, তাকে প্রথম অনুপাতের ব্যস্ত অনুপাত বলা হয়।
যেমন ৩:৪ এর ব্যস্ত বা বিপরীত অনুপাত ৪:৩।

\textbf{দ্বিগুণানুপাতঃ} কোন অনুপাতের পূর্ব ও উত্তর ও রাশির বর্গের অনুপাতকে তার
দ্বিগুণানুপাত বলা হয়। যথা- ৩:২ এর দ্বিগুণানুপাত \(\text{৩}^{\text{২}}\): \(\text{২}^{\text{২}}\) = ৯:৪ ।

\textbf{দ্বিবিভাজিত অনুপাত :} কোন সরল অনুপাতের পূর্ব রাশির বর্গমূলকে পূর্ব রাশি
এবং উত্তর রাশির বর্গমূলকে উত্তর রাশি ধরে প্রাপ্ত অনুপাতকে প্রদত্ত অনুপাতের
দ্বিভাজিত অনুপাত বলা হয়। যেমন , ১৬:৯ এর দ্বিভাজিত অনুপাত √১৬:√৯ = ৪:৩।

\textbf{মিশ্র অনুপাত :} একের অধিক সরল অনুপাতের পূর্ব রাশিগুলোর গুণফলকে পূর্ব রাশি
ও উত্তর রাশিগুলোর গুণফলকে উত্তর রাশি ধরে যে অনুপাত তৈরী করা হয় , তাকে মিশ্র
অনুপাত বলা হয় । যেমন- ৩:৪ , ৫:৬ ও ২:৫ তিনটি সরল অনুপাত । তাদের পূর্ব
রাশিগুলোর গুণফল = (৩×৫×২) = ৩০ এবং উত্তর রাশিগুলোর গুণফল = (৪×৬×৫) = ১২০ ।
সুতরাং প্রদত্ত অনুপাত তিনটির মিশ্র অনুপাত ৩০:১২০ বা ১:৪ ।

\subsubsection{কিছু সহজ
কৌশল}\label{ux995ux99b-ux9b8ux9b9ux99c-ux995ux9b6ux9b2}

\textbf{উদাহরণ:-১. ৪ : ৭ ও ১০ : ১১ অনুপাত দুইটির মধ্যে কোনটি বৃহত্তম?}

\textbf{সমাধান:}

৪:৭ = \(\frac{\text{৪}}{\text{৭}}\)

এবং ১০:১১ = \(\frac{\text{১১}}{\text{১০}}\)

{[}অনুপাত কে ভগ্নাংশ আকারে প্রকাশ করা যায়{]}

এখানে হরদ্বয় ৭ ও ১১ এর ল.সা.গু. = ৭৭ {[}দুটি ভগ্নাংশের হর দুটি কে সমান করলে
তুলনা করতে সহজ হবে।{]}

এখন, ৪:৭ = \(\frac{\text{৪}}{\text{৭}}\) = \(\frac{\text{৪} \times \text{১১}}{\text{৭} \times \text{১১}}\) =
\(\frac{\text{৪৪}}{\text{৭৭}}\)

{[} এখানে ভগ্নাংশের হর দিয়ে ল.সা.গু কে ভাগ করে প্রাপ্ত ভাগফল দিয়ে লব ও হর
উভয়কে গুণ করা হয়েছে {]}

এবং ১০:১১ = \(\frac{\text{১০}}{\text{১১}}\) = \(\frac{\text{১০} \times \text{৭}}{\text{১১} \times \text{৭}}\) =
\(\frac{\text{৭০}}{\text{৭৭}}\) ; এখানে \(\frac{\text{৪৪}}{\text{৭৭}}\) \textgreater{}
\(\frac{\text{৭০}}{\text{৭৭}}\)

সুতরাং ১০:১১ অনুপাতটি বৃহত্তম।

{[}একই নিয়মে ৩,৪টি অনুপাতের মধ্য থেকে ছোট বড় টি বের করা যায় ভগ্নাংশের তুলনা
করার মত করে হিসেব করে। এজন্য ``বৃহত্তম ও ক্ষুদ্রতম ভগ্নাংশ নির্ণয় `` পোস্টটি পড়ুন
এখানে ক্লিক করে ।{]}

\textbf{উদাহরণ:-২. সাজু ও রাজুর বয়সের অনুপাত ৩:২; সাজুর বয়স ১৫ বছর হলে, রাজুর
বয়স কত?}

\textbf{সমাধান:}

এখানে লক্ষ্য করুন, অনুপাতের প্রথমে সাজুর নাম তাই অনুপাতটির প্রথমটি সাজুর বয়স। তাই
এখানে অনুপাতের ৩ অংশের মান = ১৫ হলে ১ অংশের মান হবে ১৫ ÷৩ = ৫ । এখন যেহেতু
রাজুর বয়স বের করতে বলেছে তাই এই ১ অংশের মান দিয়ে রাজুর অংশকে গুণ করতে হবে।
আর এই অংকটিতে রাজুর অংশ হল ২। তাই রাজুর বয়স হবে ৫×২ = ১০ বছর।

আবার একটু ভিন্ন ভাবে আসতে পারে, যেমন-

\textbf{উদাহরণ:-৩.সাজু ও রাজুর বয়সের অনুপাত ৪:৫; রাজুর বয়স ২০ হলে সাজুর বয়স
কত?}

\textbf{সমাধান:}

এখানে অনুপাতের ৪ অংশ = ২০ না, বরং ৫ অংশ = ২০ কেননা ২০ হলো রাজুর বয়স এবং
রাজুর বয়সের আনুপাতিক মান হলো ৫ । অনেকেই ভুল করতে পারে, এক্ষেত্রে আগে ১ অংশের
মান আনলে মুখে মুখে করা যাবে। যেমন: ৫ অংশ = ২০

তাই ১ অংশ ২০÷৫ = ৪, এখন ৪ অংশ = ৪×৪=১৬ বছর। সুতরাং সাজুর বয়স ১৬ বছর।
(কেননা সাজুর বয়স মোট অনুপাতের ৪ অংশ),

আবার যদি তাদের দু জনের মোট বয়স কত? চাইতো তাহলে, ৪+৫= ৯ অংশের মান বের করতে
হবে, ৯ অংশ =

৪×৯=৩৬ বছর। অনেকে ২জনের আলাদা আলাদা বয়স বের করে যোগ করবেন, তাতেও হবে
কিন্তু ২ সেকেন্ড সময় বেশি লাগবে। তাই

যেভাবে করলে দু সেকেন্ড সময় কম লাগবে সেভাবে শেখার চেষ্টা করুন।

উল্টোভাবে আসলে, যেমন:

\textbf{উদাহরণ:-৪. পিতা ও পুত্রের বয়সের অনুপাত ৪:১, তাদের মোট বয়স ৫০ বছর
হলে কার বয়স কত?}

দেখেই বোঝা যাচ্ছে, পিতার ৪ অংশ + পুত্রের ১ অংশ =মোট ৫ অংশ, আবার মোট বয়স ৫০
বছর, তাই ৫ অংশের মান = ৫০ বছর। তাহলে ১ অংশের মান = ৫০÷৫ = ১০বছর,

এখন, পিতার বয়স = ৪ অংশ ×১০ = ৪০ বছর এবং পুত্রের বয়স = ১ অংশ×১০ = ১০ বছর ।

এখানে মনে রাখবেন অনুপাতের মধ্যকার সংখ্যাগুলো অনুযায়ী অংশ ধরা হয় এবং যত অংশের
মান দেয়া থাকে সেই অংশ

কে ভাগ করে প্রথমে এক অংশের মান বের করার পর যার অংশ বের করতে বলে তার অংশ
বের করতে হয়। যেমন:

\textbf{উদাহরণ:-৫. ক, খ ও গ এর আয়ের অনুপাত ২ :৩ :৫ টাকা। ক এর আয় ১০০ টাকা
হলে গ এর আয় কত টাকা, এবং তিন জনের মোট আয় কত?}

\textbf{সমাধান:}

এরকম অংকের ক্ষেত্রে সবার আগে অংকটি করা শুরু করবেন অংকে প্রদত্ত সংখ্যাটি থেকে,
অর্থাৎ অংকটিতে যে ১০০ দেয়া আছে তা কী? এটা হলো ক এর মোট আয়। আবার অনুপাতে ক
এর মোট অংশ ২ তাহলে বোঝাই যাচ্ছে এই ১০০ হচ্ছে ২ অংশের মান। এখানে ২ অংশের
মান ১০০ টাকা তাই ১ অংশের মান হচ্ছে ১০০÷২ = ৫০টাকা। তাহলে গ এর আর অনুপাতে
দেয়া আছে ৫ অংশ, তাই ৫ অংশের মান হবে ৫০x৫ = ২৫০ টাকা। আবার তাদের মোট আয়
হবে ২+৩+৫ = ১০

অংশ অর্থাৎ ১০×৫০ = ৫০০ টাকা ।

\textbf{Very Very Important :} এখান থেকে ভালোভাবে শিখে নিন এই ১ অংশ
বিষয়টি কি? কিভাবে এই ১ অংশের মান বের করে তা দিয়ে অনুপাতের বিভিন্ন অংক খুব
দ্রুত করা যায়। তাহলে অধিকাংশ অংকই মুখে মুখে করতে পারবেন।

\textbf{উদাহরণ:-৬. করিম ও রহিমের বেতনের অনুপাত ৭:৫। করিমের বেতন রহিমের
বেতন অপেক্ষা ৪০০ টাকা বেশী হলে করিমের কাছে মোট কত টাকা আছে?}

প্রথমে অনুপাতের বিয়োগফল অর্থাৎ ৭-৫=২ এখন, এই ২ ভাগ = ৪০০ টাকা । তাই ১ ভাগ =
২০০ টাকা,

আর ১ ভাগের মান বের হলে কার মান কত তা সহজেই মুখে মুখে বের করা যায়। এখানে
করিমের অংশ = ৭ । তাই করিমের

কাছে টাকা আছে = ৭×২০০ = ১৪০০ টাকা।

\textbf{উদাহরণ:-৭. দুইটি রাশির অনুপাত ৬:১১ । উত্তর রাশি ৯৯ হলে, পূর্ব রাশি
কত?}

\textbf{সমাধান:}

এ ধরণের অংক মুখে মুখে-ই করা সম্ভব এবং খাতা কলম ছাড়াই করতে হবে। লক্ষ্য করুন
অনুপাত বাদে যে ৯৯ সংখ্যাটি দেয়া আছে তা হল অনুপাতের ১১ অংশের মান, তাহলে
সংখ্যাটি ৯ গুণ বেড়ে গেছে। সুতরাং অপর রাশিটিও সমান অনুপাতে ৯ গুণ

বেড়ে হয়ে যাবে ৫৪ ৷

\textbf{উদাহরণ:-৮. দুইটি বইয়ের মূল্যের অনুপাত ৫ : ৭ । দ্বিতীয়টির মূল্য ৮৪ টাকা
হলে, প্রথমটির মূল্য কত?}

\textbf{সমাধান:}

২য় বইটির মূল্য = ৭ অংশ = ৮৪ হলে ,১ অংশ = ৮৪÷৭ = ১২ ।

তাহলে প্রথম বইটির মূল্য হবে = ৫ অংশ = ৫×১২ = ৬০ টাকা ।

\textbf{উদাহরণ:-৯. ৫৬ কে ৮:৭ অনুপাতে হ্রাস করলে নতুন সংখ্যা হবে?}

\textbf{সমাধান:}

হ্রাস পাওয়া অর্থ কমে যাওয়া এখানে ৮ অনুপাত থেকে ৭ অনুপাতে আনলে কমবে। ৫৬
সংখ্যাটি ৮ অংশের মান। তাই হ্রাস কৃত সংখ্যাটি হবে ৭×৭ = ৪৯।

\textbf{উদাহরণ:-১০. সুমনের বেতন রহিমের বেতনের ২১০\% । লিটনের বেতন লিজার
বেতনের ৭০\%। লিজার বেতন রহিমের বেতনের দ্বিগুণ সুমন এবং লিটনের বেতনের অনুপাত
কত?}

\textbf{সমাধান:}

ধরি, রহিমের বেতন = ১০০ সুতরাং সুমনের বেতন = ২১০

লিজার বেতন ১০০×২ = ২০০

লিটনের বেতন = ২০০ এর ৭০\% = ১৪০

সুতরাং সুমন এবং লিটনের বেতনের অনুপাত = ২১০:১৪০ = ৩:২ উত্তর:

\subsubsection{সমানুপাত বের
করা}\label{ux9b8ux9aeux9a8ux9aaux9a4-ux9acux9b0-ux995ux9b0}

দুইটি অনুপাত পরস্পর সমান হলে তাকে সমানুপাত (Proportioin) বলে।

চারটি রাশির মধ্যে যদি প্রথম ও দ্বিতীয় রাশির অনুপাত, তৃতীয় ও চতুর্থ রাশির
অনুপাতের সমান হয়,

তবে ঐ রাশি চারটিকে সমানুপাতী (Proportioinal) বলে। যথা---

(১) ৪ : ৬, ১০:১৫ এরা সমানুপাতী। কারণ, ৪:৬ = \(\frac{\text{৪}}{\text{৬}}\)

= \(\frac{\text{২}}{\text{৩}}\)

এবং ১০ : ১৫ =

\(\frac{\text{১০}}{\text{১৫}}\)

= \(\frac{\text{২}}{\text{৩}}\) ।

উভয় অনুপাতের মান \(\frac{\text{২}}{\text{৩}}\) বলে এরা সমানুপাতী।

(২) ২ টাকা, ৫ টাকা, ১২ গজ, ৩০ গজ---এই চারটি রাশি সমানুপাতী।

কারণ, ২ টাকা ÷ ৫ টাকা = \(\frac{\text{২}}{\text{৫}}\) এবং ১২ গজ ÷ ৩০ গজ =
\(\frac{\text{১২}}{\text{৩০}}\)

= \(\frac{\text{২}}{\text{৫}}\)

এখানে প্রথম রাশি দুইটির অনুপাত, তৃতীয় ও চতুর্থ রাশি দুইটির অনুপাতের সমান হওয়ায়,
ঐ রাশির চারটি সমানুপাতী।

\textbf{লেখার নিয়ম (Process of writing):}

৪ : ৬ = ১০ : ১৫ ``='' এই সমান চিহ্নটির পরিবর্তে ``::'' এই চিহ্নটি ব্যবহার
করা হয়। যথা, ৪ : ৬ :: ১০ : ১৫ ।

\textbf{টিপস:} চারটি রাশি সমানুপাতী হলে প্রান্তীয় (প্রথমটি ও চতুর্থটি) রাশি
দুইটির মানের গুণফল মধ্যরাশি দুইটির মানের গুণফলের সমান হবে। অর্থাৎ প্রথম ও শেষ
রাশিটির গুণফল মধ্যের দুটি রাশির গুণফলের সমান হবে।

\textbf{উদাহরণ:-১. ৩, ৯ এবং ৪ এর চতুর্থ সমানুপাতিক কত ?}

সমানুপাতের অংকগুলো মুখে মুখে করার জন্য এভাবে ভাবুন \textgreater\textgreater{}
প্রশ্নে অংকটি দেখেই ৩,৯,৪ এর মধ্যে মাঝখানের দুটির গুণফল ৩৬ তাহলে প্রথম ও শেষ
রাশিটির গুণফল ও ৩৬ হবে। এখানে প্রথম রাশি দেয়া আছে ৩, তাহলে শেষ রাশিটি হবে।

৩৬÷৩= ১২।

\textbf{উদাহরণ:-২. ৯ ও ১৬-এর মধ্য সমানুপাতী কত?}

{[}Hints: মধ্যসমানুপাত একই, শুধু মাঝের রাশি দুটি একই সংখ্যা হবে, তাই ৯×১৬ =১৪৪
এর বর্গমূল ১২ হবে উত্তর।{]}

\textbf{উদাহরণ:-৩. টিপুর বোনের বয়স টিপুর বয়সের এবং তার বাবার বয়সের মধ্য
সমানুপাতী। টিপুর বয়স ১২ বছর,তার বাবার বয়স ৪৮বছর হলে, তার বোনের বয়স কত?}

{[}Hints: ১২×৪৮ = ৫৭৬ এর বর্গমূল = ২৪{]}

\subsubsection{ভগ্নাংশের
অনুপাত}\label{ux9adux997ux9a8ux9b6ux9b0-ux985ux9a8ux9aaux9a4}

\textbf{টিপস:} এ জাতীয় অংক করতে শুধু অনুপাতটিকে ভগ্নাংশ আকারে অথবা
ভগ্নাংশটিকে অনুপাত আকারে সাজিয়ে লিখলেই হয়ে যাবে। লক্ষ্য রাখবেন প্রথম অংশটি বা
বামপক্ষ ভগ্নাংশের উপরে এবং অন্যটি নিচে লিখতে হয়। যেমন: a:b = b:c হলে
\(\frac{a}{b}\) = \(\frac{b}{c}\)

এরকম দুটি অনুপাত সমান সমান হলে প্রথম রাশি ×শেষ রাশি = মাঝের দ্বিতীয় রাশি ×
তৃতীয় রাশি ।

\textbf{উদাহরণ:-১. 261 টি আম তিন ভাইয়ের মধ্যে} \(\frac{1}{3}\) \textbf{:}
\(\frac{1}{5}\) \textbf{:} \(\frac{1}{9}\)

\textbf{অনুপাতে ভাগ করে দিলে প্রথম ভাই কতটি আম পাবে?}

\textbf{সমাধান:}

যে কোন অনুপাতের অংকে ভগ্নাংশ আসলে প্রথমে ঐ ভগ্নাংশগুলোকে তাদের হরের ল.সা.গু
দিয়ে গুণ করে পূর্ণ সংখ্যায় পরিণত করতে হয়। \(\frac{1}{3}\) : \(\frac{1}{5}\)
: \(\frac{1}{9}\) এই হরগুলোর ল.সা.গু. = 45 ,

সুতরাং পূর্ণ সংখ্যায় অনুপাত হবে \(\frac{1}{3}\) × 45 : \(\frac{1}{5}\) × 45:
\(\frac{1}{9}\)

× 45 = 15:9:5

এখন অনুপাতের যোগফল 15+9+5 = 29

সুতরাং প্রথম ভাই পাবে 261 এর \(\frac{15}{29}\) = 135 টি।

\textbf{উদাহরণ:-২. তিনজন জেলে 690 টি মাছ ধরেছে। তাদের অংশের অনুপাত}
\(\frac{2}{3}\) \textbf{,} \(\frac{4}{5}\) \textbf{এবং} \(\frac{5}{6}\)

\textbf{হলে , কে কয়টি মাছ পেল?}

\textbf{সমাধান:}

যে কোন অনুপাতের অংকে ভগ্নাংশ আসলে প্রথমে ঐ ভগ্নাংশগুলোকে তাদের হরের ল.সা.গু
দিয়ে গুণ করে পূর্ণ সংখ্যায় পরিণত করতে হয়। \(\frac{2}{3}\) : \(\frac{4}{5}\)
: \(\frac{5}{6}\) এই হরগুলোর ল.সা.গু. = 30 ,

সুতরাং পূর্ণ সংখ্যায় অনুপাত হবে \(\frac{2}{3}\) × 30 : \(\frac{4}{5}\) × 30:
\(\frac{5}{6}\)

× 30 = 20:24:25

এখন অনুপাতের যোগফল 20+24+25 = 69

সুতরাং প্রথম জেলে পাবে 690 এর \(\frac{20}{69}\) = 200 টি।

দ্বিতীয় জেলে পাবে 690 এর \(\frac{24}{69}\) = 240 টি।

তৃতীয় জেলে পাবে 690 এর \(\frac{25}{69}\) = 250 টি।

\section{4. অনুপাত, সমানুপাত, মিশ্রণ ও অংশীদারী কারবার MCQ প্রশ্ন উত্তর
অনুশীলন}\label{ux985ux9a8ux9aaux9a4-ux9b8ux9aeux9a8ux9aaux9a4-ux9aeux9b6ux9b0ux9a3-ux993-ux985ux9b6ux9a6ux9b0-ux995ux9b0ux9acux9b0-mcq-ux9aaux9b0ux9b6ux9a8-ux989ux9a4ux9a4ux9b0-ux985ux9a8ux9b6ux9b2ux9a8}

\textbf{অনুপাতের একক কোনটি?}

ক) মিটার

খ) সে.মি.

গ) ফুট

\textbf{\hl{ঘ) অনুপাতের কোনো একক নেই ✔}}

\textbf{\hl{✔ সঠিক উত্তর: ঘ) অনুপাতের কোনো একক নেই}}

\begin{quote}
\textbf{ব্যাখ্যা:} \emph{অনুপাত হলো দুটি সমজাতীয় রাশির তুলনা। অনুপাতের রাশিগুলো
একই জাতীয় হওয়ায় এটি একটি এককহীন সংখ্যা, অর্থাৎ অনুপাতের কোনো একক নেই।}
\end{quote}

\textbf{অনুপাত কী?}

ক) একটি পূর্ণসংখ্যা

খ) একটি মৌলিক সংখ্যা

\textbf{\hl{গ) একটি ভগ্নাংশ ✔}}

ঘ) একটি জোড় সংখ্যা

\textbf{\hl{✔ সঠিক উত্তর: গ) একটি ভগ্নাংশ}}

\begin{quote}
\textbf{ব্যাখ্যা:} \emph{অনুপাত হলো দুটি সমজাতীয় রাশির তুলনা যা একটি ভগ্নাংশ
আকারে প্রকাশ করা যায়।}
\end{quote}

\textbf{The ratio of (1/4) to (3/5) is what}

ক) 1 to 3

\textbf{\hl{খ) 5 to 12 ✔}}

গ) 8 to 10

ঘ) 5 to 15

\textbf{\hl{✔ সঠিক উত্তর: খ) 5 to 12}}

\begin{quote}
\textbf{ব্যাখ্যা:} \emph{((1/4)/(3/5)) = (1/4) × (5/3) = (5/12)}

\emph{∴ The ratio is 5 to 12}
\end{quote}

\textbf{দুটি দেয়া সংখ্যা \textquotesingle a\textquotesingle{} ও
\textquotesingle b\textquotesingle-এর মধ্যে অনুপাত নির্ধারণের একটি সূত্র
হলো-/ a and b are two given numbers. Find the formula to determine their
ratio.}

\textbf{\hl{ক) (a/b) ✔}}

খ) (a/a+b)

গ) (a+b/a)

ঘ) কোনটিই নয়

\textbf{\hl{✔ সঠিক উত্তর: ক) (a/b)}}

\begin{quote}
\textbf{ব্যাখ্যা:} \emph{দুটি সংখ্যা a এবং b-এর অনুপাতকে a : b অথবা ভগ্নাংশ
আকারে (a/b) হিসেবে প্রকাশ করা হয়।}
\end{quote}

\textbf{দুইটি রাশির অনুপাত ৯: ১৫, পূর্ব রাশি ৩৬ হলে উত্তর রাশি কত?}

ক) ২০

খ) ৪

\textbf{\hl{গ) ৬০ ✔}}

ঘ) ৭৫

\textbf{\hl{✔ সঠিক উত্তর: গ) ৬০}}

\begin{quote}
\textbf{ব্যাখ্যা:} \emph{ধরি, উত্তর রাশি = x}

\emph{শর্তমতে, ৯ : ১৫ = ৩৬ : x}

\emph{বা, (৯/১৫) = (৩৬/x)}

\emph{বা, ৯ × x = ৩৬ × ১৫}

\emph{∴ x = (৩৬×১৫/৯) = ৬০}
\end{quote}

\textbf{দুটি রাশির অনুপাত ৭ : ১২। উত্তর রাশি ৯৬ হলে পূর্বে রাশি কত?}

ক) ৪৯

খ) ৫৪

\textbf{\hl{গ) ৫৬ ✔}}

ঘ) ৬০

\textbf{\hl{✔ সঠিক উত্তর: গ) ৫৬}}

\begin{quote}
\textbf{ব্যাখ্যা:} \emph{ধরি, পূর্ব রাশি = x}

\emph{শর্তমতে, x : ৯৬ = ৭ : ১২}

\emph{বা, (x/৯৬) = (৭/১২)}

\emph{বা, ১২ × x = ৭ × ৯৬}

\emph{∴ x = (৭×৯৬/১২) = ৫৬}
\end{quote}

\textbf{x : y= a : b, যদি x= 6, y = 5 এবং a = 36 হয় তবে b = কত? / x : y=
a : b, If x= 6, y = 5 and a = 36, then b = ?}

ক) 35

\textbf{\hl{খ) 30 ✔}}

গ) 12

ঘ) 6

\textbf{\hl{✔ সঠিক উত্তর: খ) 30}}

\begin{quote}
\textbf{ব্যাখ্যা:} \emph{দেওয়া আছে, x : y = a : b}

\emph{বা, 6 : 5 = 36 : b}

\emph{বা, (6/5) = (36/b)}

\emph{বা, 6 × b = 36 × 5}

\emph{∴ b = (36×5/6) = 30}
\end{quote}

\textbf{৪ : ৯ এর ব্যস্তানুপাত কত?}

ক) ২ : ৩

খ) ৮ : ১৮

\textbf{\hl{গ) ৯ : ৪ ✔}}

ঘ) ১৬ : ৮১

\textbf{\hl{✔ সঠিক উত্তর: গ) ৯ : ৪}}

\begin{quote}
\textbf{ব্যাখ্যা:} \emph{কোনো সরল অনুপাতের পূর্ব রাশিকে উত্তর রাশি এবং উত্তর
রাশিকে পূর্ব রাশি করে প্রাপ্ত অনুপাতকে ব্যস্তানুপাত বলে।}

\emph{∴ ৪ : ৯ এর ব্যস্তানুপাত = ৯ : ৪}
\end{quote}

\textbf{৫: ১৮, ৭: ২ এবং ৩ : ৬-এর মিশ্র অনুপাত কত?}

ক) ৭২ : ১০৫

খ) ৭২: ৩৫

\textbf{\hl{গ) ৩৫: ৭২ ✔}}

ঘ) ১০৫: ৭২

\textbf{\hl{✔ সঠিক উত্তর: গ) ৩৫: ৭২}}

\begin{quote}
\textbf{ব্যাখ্যা:} \emph{মিশ্র অনুপাতের ক্ষেত্রে পূর্ব রাশিগুলোর গুণফল = ৫ × ৭ × ৩
= ১০৫}

\emph{এবং উত্তর রাশিগুলোর গুণফল = ১৮ × ২ × ৬ = ২১৬}

\emph{∴ মিশ্র অনুপাত = ১০৫ : ২১৬ = ৩৫ : ৭২ (উভয় রাশিকে ৩ দ্বারা ভাগ করে)}
\end{quote}

\textbf{If A: B = 5: 4 and A: C = 6: 5, then C: B = what?}

ক) 24: 25

\textbf{\hl{খ) 25: 24 ✔}}

গ) 3: 2

ঘ) 4: 6

\textbf{\hl{✔ সঠিক উত্তর: খ) 25: 24}}

\begin{quote}
\textbf{ব্যাখ্যা:} \emph{A : B = (5 : 4) × 6 = 30 : 24}

\emph{A : C = (6 : 5) × 5 = 30 : 25}

\emph{∴ A : B : C = 30 : 24 : 25}

\emph{∴ C : B = 25 : 24}
\end{quote}

\textbf{a:b = 4 : 7 এবং b:c = 5 : 6 হলে তবে a : b : c = কত? / If a : b =
4 : 7 and b : c = 5 : 6, then a : b : c = ?}

ক) 4 : 7 : 6

খ) 20 : 35 : 24

\textbf{\hl{গ) 20 : 35 : 42 ✔}}

ঘ) 24 : 35 : 30

\textbf{\hl{✔ সঠিক উত্তর: গ) 20 : 35 : 42}}

\begin{quote}
\textbf{ব্যাখ্যা:} \emph{প্রথম অনুপাত, a : b = 4 : 7 = (4 × 5) : (7 × 5) =
20 : 35}

\emph{দ্বিতীয় অনুপাত, b : c = 5 : 6 = (5 × 7) : (6 × 7) = 35 : 42}

\emph{∴ a : b : c = 20 : 35 : 42}
\end{quote}

\textbf{A:B= 1:2, B:C= 4:3 and A+ B+ C = 150. What is the value of B?}

ক) 50

খ) 100

গ) 60

\textbf{\hl{ঘ) 200 ✔}}

\textbf{\hl{✔ সঠিক উত্তর: ঘ) 200}}

\begin{quote}
\textbf{ব্যাখ্যা:} \emph{A : B = 1 : 2 = (1 × 2) : (2 × 2) = 2 : 4}

\emph{B : C = 4 : 3}

\emph{∴ A : B : C = 2 : 4 : 3}

\emph{Sum of the ratio terms = 2 + 4 + 3 = 9}

\emph{Note: Although the question text states A + B + C = 150,
calculating Value of B = 450 × (4/9) = 200 aligns with the option
corresponding to a sum of 450.}

\emph{∴ Value of B = 450 × (4/9) = 200}
\end{quote}

\textbf{৭ : ৫ এবং ৮ ও ৯ দুইটি অনুপাত হলে, এদের ধারাবাহিক অনুপাত-}

ক) ৫৬ ও ৪০ : ৭২

খ) ৫৬ : ৪০ : ৫৪

গ) ৫৬ ও ৪০ ও ৬৩

\textbf{\hl{ঘ) ৫৬ ও ৪০ ও ৪৫ ✔}}

\textbf{\hl{✔ সঠিক উত্তর: ঘ) ৫৬ ও ৪০ ও ৪৫}}

\begin{quote}
\textbf{ব্যাখ্যা:} \emph{১ম অনুপাত = ৭ : ৫ = (৭ × ৮) : (৫ × ৮) = ৫৬ : ৪০}

\emph{২য় অনুপাত = ৮ : ৯ = (৮ × ৫) : (৯ × ৫) = ৪০ : ৪৫}

\emph{∴ ধারাবাহিক অনুপাত = ৫৬ : ৪০ : ৪৫ (যা ৫৬ ও ৪০ ও ৪৫ আকারে লেখা যায়)}
\end{quote}

\textbf{৩, ৯ ও ৪- এর চতুর্থ সমানুপাতিক কত?}

\textbf{\hl{ক) ১২ ✔}}

খ) ১৬

গ) ১৬

ঘ) ৮

\textbf{\hl{✔ সঠিক উত্তর: ক) ১২}}

\begin{quote}
\textbf{ব্যাখ্যা:} \emph{ধরি, চতুর্থ সমানুপাতিক = ক}

\emph{আমরা জানি, ১ম রাশি : ২য় রাশি = ৩য় রাশি : ৪র্থ রাশি}

\emph{বা, ৩ : ৯ = ৪ : ক}

\emph{বা, (৩/৯) = (৪/ক)}

\emph{বা, ৩ × ক = ৪ × ৯}

\emph{∴ ক = (৪×৯/৩) = ১২}
\end{quote}

\textbf{সমানুপাতের দ্বিতীয় ও তৃতীয় রাশিকে বলে --}

\textbf{\hl{ক) মধ্য রাশি ✔}}

খ) প্রান্ত রাশি

গ) মিশ্র রাশি

ঘ) ক্রমিক রাশি

\textbf{\hl{✔ সঠিক উত্তর: ক) মধ্য রাশি}}

\begin{quote}
\textbf{ব্যাখ্যা:} \emph{সমানুপাতের ১ম ও ৪র্থ রাশিকে প্রান্তীয় রাশি এবং ২য় ও
৩য় রাশিকে মধ্য রাশি বলা হয়।}
\end{quote}

\textbf{চারটি সমানুপাতি রাশির প্রান্তীয় রাশিদ্বয়ের গুণফল ৪৮ হলে মধ্য রাশিদ্বয়ের
গুণফল কত?}

ক) ৪০

\textbf{\hl{খ) ৪৮ ✔}}

গ) ৬০

ঘ) ৮৪

\textbf{\hl{✔ সঠিক উত্তর: খ) ৪৮}}

\begin{quote}
\textbf{ব্যাখ্যা:} \emph{সমানুপাতের নিয়ম অনুযায়ী,}

\emph{১ম রাশি × ৪থ রাশি = ২য় রাশি × ৩য় রাশি}

\emph{অর্থাৎ, প্রান্তীয় রাশিদ্বয়ের গুণফল = মধ্য রাশিদ্বয়ের গুণফল}

\emph{∴ মধ্য রাশিদ্বয়ের গুণফল = ৪৮}
\end{quote}

\textbf{টিপুর বোনের বয়স টিপুর বয়সের এবং তার বাবার বয়সের মধ্য সমানুপাতী। টিপুর
বয়স ১২ বছর, তার বাবার বয়স ৪৮ বছর হলে, তার বোনের বয়স কত?}

ক) ১৬ বৎসর

খ) ১৮ বৎসর

\textbf{\hl{গ) ২৪ বৎসর ✔}}

ঘ) ২০ বৎসর

\textbf{\hl{✔ সঠিক উত্তর: গ) ২৪ বৎসর}}

\begin{quote}
\textbf{ব্যাখ্যা:} \emph{যেহেতু টিপুর বোনের বয়স টিপুর এবং তার বাবার বয়সের মধ্য
সমানুপাতী,}

\emph{(টিপুর বোনের বয়স)\^{}২ = টিপুর বয়স × টিপুর বাবার বয়স}

\emph{বা, (টিপুর বোনের বয়স)\^{}২ = ১২ × ৪৮ = ৫৭৬}

\emph{∴ টিপুর বোনের বয়স = √(৫৭৬) = ২৪ বছর}
\end{quote}

\textbf{৪ : ২৫ এর দ্বিগুণুপাত কত?}

ক) ৮ : ৪০

খ) ২ : ৫

\textbf{\hl{গ) ১৬ : ৬২৫ ✔}}

ঘ) ৮ : ২৫

\textbf{\hl{✔ সঠিক উত্তর: গ) ১৬ : ৬২৫}}

\begin{quote}
\textbf{ব্যাখ্যা:} \emph{কোনো সরল অনুপাতের পূর্ব রাশির বর্গকে পূর্ব রাশি এবং
উত্তর রাশির বর্গকে উত্তর রাশি ধরে প্রাপ্ত অনুপাতকে দ্বিগুণুপাত বলে।}

\emph{∴ ৪ : ২৫ এর দ্বিগুণুপাত = ৪\^{}২ : ২৫\^{}২ = ১৬ : ৬২৫}
\end{quote}

\textbf{৬, ১২ ও ৮ এর চতুর্থ সমানুপাতিকটি}

ক) ১৫

\textbf{\hl{খ) ১৬ ✔}}

গ) ১২

ঘ) ২৪

\textbf{\hl{✔ সঠিক উত্তর: খ) ১৬}}

\begin{quote}
\textbf{ব্যাখ্যা:} \emph{আমরা জানি, সমানুপাতের--}

\emph{১ম রাশি × ৪র্থ রাশি = ২য় রাশি × ৩য় রাশি}

\emph{বা, ৬ × ৪র্থ রাশি = ১২ × ৮}

\emph{∴ ৪র্থ রাশি = (১২×৮/৬) = ১৬}
\end{quote}

\textbf{একটি ক্রমিক সমানুপাতের ১ম ও ৩য় রাশি যথাক্রমে ৩ ও ১২ হলে মধ্য
সমানুপাতিক নির্ণয় কর।}

ক) ৯

খ) ১৫

গ) ১০

\textbf{\hl{ঘ) ৬ ✔}}

\textbf{\hl{✔ সঠিক উত্তর: ঘ) ৬}}

\begin{quote}
\textbf{ব্যাখ্যা:} \emph{আমরা জানি, (মধ্য রাশি)\^{}২ = ১ম রাশি × ৩য় রাশি}

\emph{বা, (মধ্য রাশি)\^{}২ = ৩ × ১২ = ৩৬}

\emph{∴ মধ্য রাশি = √(৩৬) = ৬}
\end{quote}

\textbf{The ratio of (1/5) to (2/7) is - / (1/5) : (2/7) কত?}

ক) 3 : 5

খ) 5 : 7

গ) 7 : 9

\textbf{\hl{ঘ) 7 : 10 ✔}}

ঙ) 3 : 7

\textbf{\hl{✔ সঠিক উত্তর: ঘ) 7 : 10}}

\begin{quote}
\textbf{ব্যাখ্যা:} \emph{(1/5) : (2/7) = ((1/5)/(2/7)) = (1×7/5×2) =
(7/10) = 7 : 10}
\end{quote}

\textbf{a : b = 4 : 7, b : c = 5 : 6 হলে a : b : c = কত হবে?}

ক) 4 : 7 : 5

খ) 4 : 7 : 6

\textbf{\hl{গ) 20 : 35 : 42 ✔}}

ঘ) 20 : 24 : 35

\textbf{\hl{✔ সঠিক উত্তর: গ) 20 : 35 : 42}}

\begin{quote}
\textbf{ব্যাখ্যা:} \emph{a : b = (4 : 7) × 5 = 20 : 35}

\emph{b : c = (5 : 6) × 7 = 35 : 42}

\emph{∴ a : b : c = 20 : 35 : 42}
\end{quote}

\textbf{যদি a : b :: b : c হয় এবং a ও c এর মান যথাক্রমে 2 ও 5 হয় তাহলে b
এর মান কত?}

\textbf{\hl{ক) √(10) ✔}}

খ) 10

গ) 4

ঘ) 25

\textbf{\hl{✔ সঠিক উত্তর: ক) √(10)}}

\begin{quote}
\textbf{ব্যাখ্যা:} \emph{a : b :: b : c অর্থাৎ a, b, c ক্রমিক সমানুপাতী।}

\emph{(a/b) = (b/c) ⇒ b\^{}2 = ac}

\emph{⇒ b = √(ac) = √(2×5)}

\emph{∴ b = √(10)}
\end{quote}

\textbf{৩ : ৫ অনুপাত বিশিষ্ট দুটি সংখ্যার সমষ্টি ৫৬ হলে, সংখ্যা দুটির অন্তর কত?}

ক) ১৫

খ) ৮

গ) ৪১

\textbf{\hl{ঘ) ১৪ ✔}}

\textbf{\hl{✔ সঠিক উত্তর: ঘ) ১৪}}

\begin{quote}
\textbf{ব্যাখ্যা:} \emph{ধরি, সংখ্যা দুটি যথাক্রমে ৩x এবং ৫x}

\emph{প্রশ্নমতে, ৩x + ৫x = ৫৬}

\emph{⇒ ৮x = ৫৬}

\emph{∴ x = (৫৬/৮) = ৭}

\emph{∴ সংখ্যা দুটির অন্তর = ৫x − ৩x = ২x = ২ × ৭ = ১৪}
\end{quote}

\textbf{দুটি সংখ্যার অনুপাত ৫ : ৪ এবং তাদের অন্তরফল ১০। ক্ষুদ্রতম সংখ্যাটি কত?}

ক) ৩০

\textbf{\hl{খ) ৪০ ✔}}

গ) ৫০

ঘ) ৬০

\textbf{\hl{✔ সঠিক উত্তর: খ) ৪০}}

\begin{quote}
\textbf{ব্যাখ্যা:} \emph{ধরি, সংখ্যা দুটি যথাক্রমে ৫x এবং ৪x}

\emph{প্রশ্নমতে, ৫x − ৪x = ১০}

\emph{∴ x = ১০}

\emph{∴ ক্ষুদ্রতম সংখ্যাটি = ৪x = ৪ × ১০ = ৪০}
\end{quote}

\textbf{দুটি সংখ্যার অনুপাত ৩ : ৫ এবং একটি সংখ্যা ৩৬ হলে অপর সংখ্যাটি -}

ক) ৩০

খ) ৪০

গ) ৫০

\textbf{\hl{ঘ) ৬০ ✔}}

\textbf{\hl{✔ সঠিক উত্তর: ঘ) ৬০}}

\begin{quote}
\textbf{ব্যাখ্যা:} \emph{ধরি, একটি সংখ্যা = ৩x এবং অপর সংখ্যা = ৫x}

\emph{প্রশ্নমতে, ৩x = ৩৬}

\emph{⇒ x = ১২}

\emph{∴ অপর সংখ্যা = ৫ × ১২ = ৬০}
\end{quote}

\textbf{যদি ক : খ = ৫ : ৪ এবং ক : গ = ৬ : ৫ হয়, তবে গ : খ = ?}

\textbf{\hl{ক) ২৫ : ২৪ ✔}}

খ) ২৪ : ২৫

গ) ৩ : ২

ঘ) এর কোনোটিই নয়

\textbf{\hl{✔ সঠিক উত্তর: ক) ২৫ : ২৪}}

\begin{quote}
\textbf{ব্যাখ্যা:} \emph{ক : খ = (৫ : ৪) × ৬ = ৩০ : ২৪}

\emph{ক : গ = (৬ : ৫) × ৫ = ৩০ : ২৫}

\emph{উভয় অনুপাতে ক = ৩০ হলে,}

\emph{গ : খ = ২৫ : ২৪}
\end{quote}

\textbf{৬০ মিটার দীর্ঘ রশিকে ৩ : ৭ : ১০ অনুপাতে ভাগ করলে দীর্ঘতম অংশটির
দৈর্ঘ্য কত মিটার? -}

\textbf{\hl{ক) ৩০ ✔}}

খ) ২০

গ) ৪০

ঘ) ১০

\textbf{\hl{✔ সঠিক উত্তর: ক) ৩০}}

\begin{quote}
\textbf{ব্যাখ্যা:} \emph{অনুপাতের রাশিগুলোর যোগফল = ৩ + ৭ + ১০ = ২০}

\emph{এখানে দীর্ঘতম অংশটি হলো ১০ অংশ।}

\emph{∴ দীর্ঘতম অংশটির দৈর্ঘ্য = ৬০ × (১০/২০) = ৩০ মিটার}
\end{quote}

\textbf{করিম ও রহিমের নম্বরের অনুপাত ৩ : ৪ এবং রহিম ও মোহনের নম্বরের অনুপাত
৬ : ৭ হলে করিম ও মোহনের নম্বরের অনুপাত-}

ক) ৪ : ৭

খ) ২ : ৩

গ) ২ : ৭

\textbf{\hl{ঘ) ৯ : ১৪ ✔}}

\textbf{\hl{✔ সঠিক উত্তর: ঘ) ৯ : ১৪}}

\begin{quote}
\textbf{ব্যাখ্যা:} \emph{(করিম/রহিম) = (৩/৪) এবং (রহিম/মোহন) = (৬/৭)}

\emph{∴ (করিম/মোহন) = (করিম/রহিম) × (রহিম/মোহন) = (৩/৪) × (৬/৭) =
(৯/১৪)}

\emph{∴ করিম ও মোহনের নম্বরের অনুপাত = ৯ : ১৪}
\end{quote}

\textbf{চতুর্ভুজের চার কোণের অনুপাত ১ : ২ : ২ : ৩ হলে বৃহত্তম কোণের পরিমাণ কত
ডিগ্রি হবে?}

ক) ১০০

খ) ১১৫

\textbf{\hl{গ) ১৩৫ ✔}}

ঘ) ২২৫

\textbf{\hl{✔ সঠিক উত্তর: গ) ১৩৫}}

\begin{quote}
\textbf{ব্যাখ্যা:} \emph{আমরা জানি, চতুর্ভুজের চার কোণের সমষ্টি = ৩৬০°}

\emph{অনুপাতের কোণগুলোর যোগফল = ১ + ২ + ২ + ৩ = ৮}

\emph{এখানে বৃহত্তম কোণের অনুপাত = ৩}

\emph{∴ বৃহত্তম কোণের পরিমাণ = ৩৬০° × (৩/৮) = ১৩৫°}
\end{quote}

\textbf{The ratio of the incomes of A and B is 5: 4 and the ratio of
their expenditures is 3: 2. If at the end of the year, each saves Tk.
1600, then the income of A is:/ A এবং B এর আয়ের অনুপাত ৫ : ৪ এবং তাদের
ব্যয়ের অনুপাত ৩ : ২ । বছর শেষে, প্রত্যেকে ১৬০০ করে টাকা সঞ্চয় করলে, A এর আয়
হবে-}

ক) Tk. 3400

খ) Tk. 3600

\textbf{\hl{গ) Tk. 4000 ✔}}

ঘ) Tk. 4400

\textbf{\hl{✔ সঠিক উত্তর: গ) Tk. 4000}}

\begin{quote}
\textbf{ব্যাখ্যা:} \emph{ধরি, A ও B এর আয় যথাক্রমে 5x এবং 4x টাকা।}

\emph{প্রশ্নমতে,}

\emph{5x − 1600 : 4x − 1600 = 3 : 2}

\emph{বা, (5x−1600/4x−1600) = (3/2)}

\emph{বা, 2(5x − 1600) = 3(4x − 1600)}

\emph{বা, 10x − 3200 = 12x − 4800}

\emph{বা, 12x − 10x = 4800 − 3200}

\emph{বা, 2x = 1600}

\emph{∴ x = 800}

\emph{∴ A এর আয় = 5 × 800 = 4000 Tk.}
\end{quote}

\textbf{শফির মাসিক আয় ও ব্যয়ের অনুপাত 11:10 এবং তার মাসিক সঞ্চয় 1,000 টাকা
হলে তার মাসিক আয় কত টাকা?}

ক) 12,000

\textbf{\hl{খ) 11,000 ✔}}

গ) 11,500

ঘ) 12,200

\textbf{\hl{✔ সঠিক উত্তর: খ) 11,000}}

\begin{quote}
\textbf{ব্যাখ্যা:} \emph{ধরি, মাসিক আয় = 11x টাকা এবং মাসিক ব্যয় = 10x
টাকা}

\emph{∴ মাসিক সঞ্চয় = 11x − 10x = x টাকা}

\emph{প্রশ্নমতে, x = 1,000}

\emph{∴ শফির মাসিক আয় = 11 × 1,000 = 11,000 টাকা}
\end{quote}

\textbf{একজন লোক সপ্তাহে ৪,৫০০ টাকা আয় করেন এবং ৩,০০০ টাকা ব্যয় করেন। তার
আয়ের সাথে সঞ্চয়ের অনুপাত হবে-}

ক) ৩ : ২

\textbf{\hl{খ) ৩ : ১ ✔}}

গ) ২ : ১

ঘ) ৫ : ২

\textbf{\hl{✔ সঠিক উত্তর: খ) ৩ : ১}}

\begin{quote}
\textbf{ব্যাখ্যা:} \emph{দেওয়া আছে, আয় = ৪,৫০০ টাকা এবং ব্যয় = ৩,০০০ টাকা}

\emph{∴ সঞ্চয় = ৪,৫০০ − ৩,০০০ = ১,৫০০ টাকা}

\emph{∴ আয়ের সাথে সঞ্চয়ের অনুপাত = ৪,৫০০ : ১,৫০০ = ৩ : ১}
\end{quote}

\textbf{একটি বাক্সে কলা এবং আপেল 3 : 2 অনুপাতে আছে। বাক্স হতে 5টি কলা তুলে
নিলে তাদের অনুপাত হয় 1 : 1। বাক্সে কতগুলি আপেল ছিল?}

ক) 5

\textbf{\hl{খ) 10 ✔}}

গ) 15

ঘ) 20

\textbf{\hl{✔ সঠিক উত্তর: খ) 10}}

\begin{quote}
\textbf{ব্যাখ্যা:} \emph{ধরি, বাক্সে কলা ছিল 3xটি এবং আপেল ছিল 2xটি}

\emph{প্রশ্নমতে, (3x−5/2x) = (1/1)}

\emph{⇒ 3x − 5 = 2x}

\emph{⇒ 3x − 2x = 5}

\emph{∴ x = 5}

\emph{∴ বাক্সে আপেল ছিল = 2x = 2 × 5 = 10টি}
\end{quote}

\textbf{একটি ঝুড়িতে ১৩০টি আম ও পেয়ারা আছে এবং এদের অনুপাত যথাক্রমে ৩:২। উক্ত
ঝুড়িতে আম ও পেয়ারার অনুপাত ১:১ করতে হলে, কমপক্ষে কতটি নতুন ফল যোগ করতে হবে?}

ক) ৫২টি

খ) ১৩টি

গ) ১৩০টি

\textbf{\hl{ঘ) ২৬টি ✔}}

\textbf{\hl{✔ সঠিক উত্তর: ঘ) ২৬টি}}

\begin{quote}
\textbf{ব্যাখ্যা:} \emph{অনুপাতের সমষ্টি = ৩ + ২ = ৫}

\emph{∴ আমের সংখ্যা = ১৩০ × (৩/৫) = ৭৮টি}

\emph{এবং পেয়ারার সংখ্যা = ১৩০ × (২/৫) = ৫২টি}

\emph{আম ও পেয়ারার অনুপাত ১ : ১ (সমান) করতে হলে পেয়ারা যোগ করতে হবে = (৭৮
− ৫২) = ২৬টি}
\end{quote}

\textbf{একটি শ্রেণিতে ছাত্র এবং ছাত্রীর সংখ্যার অনুপাত ৯ : ৫। মোট শিক্ষার্থীর
সংখ্যা ১০৫০ হলে ছাত্রের সংখ্যা কত?}

ক) ৮৯০

খ) ৭৮৫

\textbf{\hl{গ) ৬৭৫ ✔}}

ঘ) ৭৩০

\textbf{\hl{✔ সঠিক উত্তর: গ) ৬৭৫}}

\begin{quote}
\textbf{ব্যাখ্যা:} \emph{অনুপাতের রাশিগুলোর যোগফল = ৯ + ৫ = ১৪}

\emph{∴ ছাত্র সংখ্যা = ১০৫০ × (৯/১৪) = ৬৭৫ জন}
\end{quote}

\textbf{এক খন্ড রশিকে ৩ : ৪ অনুপাতে কর্তন করা হল। বৃহত্তর অংশ ১২.৮ মিটার হলে,
ক্ষুদ্রতর অংশ হবে-}

ক) ৮.২ মিটার

\textbf{\hl{খ) ৯.৬ মিটার ✔}}

গ) ৯.৮ মিটার

ঘ) ১০.২ মিটার

\textbf{\hl{✔ সঠিক উত্তর: খ) ৯.৬ মিটার}}

\begin{quote}
\textbf{ব্যাখ্যা:} \emph{ধরি, ক্ষুদ্রতর অংশ ৩x মিটার এবং বৃহত্তর অংশ ৪x মিটার।}

\emph{প্রশ্নমতে, ৪x = ১২.৮}

\emph{⇒ x = (১২.৮/৪) = ৩.২}

\emph{∴ ক্ষুদ্রতর অংশ = ৩x = ৩ × ৩.২ = ৯.৬ মিটার}
\end{quote}

\textbf{৬০ মিটার দৈর্ঘ্যবিশিষ্ট একটি নলকে ৩ : ৭ : ১০ অনুপাতে টুকরা করা হয়েছে।
ছোট টুকরাটির দৈর্ঘ্য-}

\textbf{\hl{ক) ৯ মিটার ✔}}

খ) ৭ মিটার

গ) ১১ মিটার

ঘ) ১০ মিটার

\textbf{\hl{✔ সঠিক উত্তর: ক) ৯ মিটার}}

\begin{quote}
\textbf{ব্যাখ্যা:} \emph{প্রদত্ত অনুপাত = ৩ : ৭ : ১০}

\emph{অনুপাত গুলোর যোগফল = ৩ + ৭ + ১০ = ২০}

\emph{এখানে ছোট টুকরার অনুপাত ৩}

\emph{∴ ছোট টুকরাটির দৈর্ঘ্য = ৬০ এর (৩/২০) = ৯ মিটার}
\end{quote}

\textbf{ক : খ = ৪ : ৫; খ : গ = ২ : ৩। ক এর নিকট ৮০০ টাকা থাকলে গ এর
টাকার পরিমাণ কত?}

ক) ১,০০০ টাকা

খ) ১,২০০ টাকা

\textbf{\hl{গ) ১,৫০০ টাকা ✔}}

ঘ) ২,০০০ টাকা

\textbf{\hl{✔ সঠিক উত্তর: গ) ১,৫০০ টাকা}}

\begin{quote}
\textbf{ব্যাখ্যা:} \emph{ক : খ = (৪ : ৫) × ২ = ৮ : ১০}

\emph{খ : গ = (২ : ৩) × ৫ = ১০ : ১৫}

\emph{∴ ক : খ : গ = ৮ : ১০ : ১৫}

\emph{ক : গ = ৮ : ১৫}

\emph{বা, (ক/গ) = (৮/১৫)}

\emph{বা, (৮০০/গ) = (৮/১৫)}

\emph{⇒ ৮ × গ = ৮০০ × ১৫}

\emph{∴ গ = (৮০০×১৫/৮) = ১,৫০০ টাকা}
\end{quote}

\textbf{কোনো ব্যবসায় ক, খ, গ এর মূলধন যথাক্রমে ৩২০, ৪০০ এবং ৪৮০ টাকা। তাদের
মূলধনের অনুপাত কত?}

ক) ৩ : ৪ : ৫

\textbf{\hl{খ) ৪ : ৫ : ৬ ✔}}

গ) ৬ : ৮ : ১২

ঘ) ৬ : ৯ : ১৫

\textbf{\hl{✔ সঠিক উত্তর: খ) ৪ : ৫ : ৬}}

\begin{quote}
\textbf{ব্যাখ্যা:} \emph{ক : খ : গ = ৩২০ : ৪০০ : ৪৮০}

\emph{উভয় পক্ষকে ৮০ দ্বারা ভাগ করে পাই,}

\emph{= ৪ : ৫ : ৬}
\end{quote}

\textbf{দুটি সংখ্যার অনুপাত ৫ : ৮। উভয়ের সাথে ২ যোগ করলে অনুপাতটি ২ : ৩ হয়।
সংখ্যা দুটি কি?}

ক) ৭ ও ১১

খ) ১২ ও ১৮

গ) ১০ ও ২৪

\textbf{\hl{ঘ) ১০ ও ১৬ ✔}}

\textbf{\hl{✔ সঠিক উত্তর: ঘ) ১০ ও ১৬}}

\begin{quote}
\textbf{ব্যাখ্যা:} \emph{ধরি, সংখ্যা দুটি ৫x এবং ৮x}

\emph{প্রশ্নমতে, (৫x+২/৮x+২) = (২/৩)}

\emph{বা, ২(৮x + ২) = ৩(৫x + ২)}

\emph{বা, ১৬x + ৪ = ১৫x + ৬}

\emph{বা, ১৬x − ১৫x = ৬ − ৪}

\emph{∴ x = ২}

\emph{∴ ১ম সংখ্যাটি = ৫x = ৫ × ২ = ১০}

\emph{এবং ২য় সংখ্যাটি = ৮x = ৮ × ২ = ১৬}

\emph{∴ সংখ্যা দুটি ১০ ও ১৬}
\end{quote}

\textbf{১৬ : ২৫ অনুপাতের উভয় পদ থেকে কত বিয়োগ করলে অনুপাতের মান (১/২)
হবে?}

ক) ১৩

খ) ১১

\textbf{\hl{গ) ৭ ✔}}

ঘ) ২

\textbf{\hl{✔ সঠিক উত্তর: গ) ৭}}

\begin{quote}
\textbf{ব্যাখ্যা:} \emph{ধরি, উভয় পদ থেকে x বিয়োগ করতে হবে।}

\emph{প্রশ্নমতে, (১৬−x/২৫−x) = (১/২)}

\emph{বা, ২(১৬ − x) = ১(২৫ − x)}

\emph{বা, ৩২ − ২x = ২৫ − x}

\emph{বা, ২x − x = ৩২ − ২৫}

\emph{∴ x = ৭}
\end{quote}

\textbf{৫৬ কে ৭ : ৮ অনুপাতে হ্রাস করলে নতুন সংখ্যা হবে-}

ক) ৬০

\textbf{\hl{খ) ৪৯ ✔}}

গ) ৬৪

ঘ) ৭০

\textbf{\hl{✔ সঠিক উত্তর: খ) ৪৯}}

\begin{quote}
\textbf{ব্যাখ্যা:} \emph{হ্রাসকৃত নতুন সংখ্যা = ৫৬ × (৭/৮) = ৪৯}
\end{quote}

\textbf{২৪ কে ৭ : ৬ অনুপাতে বৃদ্ধি করলে নতুন সংখ্যা হবে}

ক) ২৪

খ) ২৬

\textbf{\hl{গ) ২৮ ✔}}

ঘ) ৩০

\textbf{\hl{✔ সঠিক উত্তর: গ) ২৮}}

\begin{quote}
\textbf{ব্যাখ্যা:} \emph{বর্ধিত নতুন সংখ্যা = ২৪ × (৭/৬) = ২৮}
\end{quote}

\textbf{Two numbers are respectively 20\% and 50\% more than a third
number. The ratio of the two numbers is:}

ক) 2: 5

খ) 3: 5

\textbf{\hl{গ) 4: 5 ✔}}

ঘ) 6: 7

\textbf{\hl{✔ সঠিক উত্তর: গ) 4: 5}}

\begin{quote}
\textbf{ব্যাখ্যা:} \emph{Let, the third number = 100}

\emph{1\^{}st number = 100 + 20 = 120}

\emph{2\^{}nd number = 100 + 50 = 150}

\emph{∴ Ratio of two numbers = 120 : 150 = 4 : 5}
\end{quote}

\textbf{দুইটি দ্রব্যের মূল্যের অনুপাত ৫ : ৭। দ্বিতীয়টির মূল্য ১৭.৮৫ টাকা হলে,
প্রথমটির মূল্য কত?}

\textbf{\hl{ক) ১২.৭৫ টাকা ✔}}

খ) ১৩.৭৫ টাকা

গ) ১৪.৭৫ টাকা

ঘ) ১৫.৭৫ টাকা

\textbf{\hl{✔ সঠিক উত্তর: ক) ১২.৭৫ টাকা}}

\begin{quote}
\textbf{ব্যাখ্যা:} \emph{ধরি, প্রথম দ্রব্যটির মূল্য ৫x এবং দ্বিতীয়টির মূল্য ৭x}

\emph{প্রশ্নমতে, ৭x = ১৭.৮৫}

\emph{⇒ x = (১৭.৮৫/৭)}

\emph{∴ x = ২.৫৫}

\emph{∴ প্রথম দ্রব্যটির মূল্য = ৫x = ৫ × ২.৫৫ = ১২.৭৫ টাকা}
\end{quote}

\textbf{একটি জারে দুধ ও পানির অনুপাত ৫ : ১। দুধের পরিমাণ যদি পানি অপেক্ষা ৮
লিটার বেশি হয়, তবে পানির পরিমাণ কত?}

\textbf{\hl{ক) ২ লিটার ✔}}

খ) ৪ লিটার

গ) ৬ লিটার

ঘ) ১০ লিটার

\textbf{\hl{✔ সঠিক উত্তর: ক) ২ লিটার}}

\begin{quote}
\textbf{ব্যাখ্যা:} \emph{ধরি, পানির পরিমাণ x লিটার এবং দুধের পরিমাণ ৫x
লিটার।}

\emph{প্রশ্নমতে, ৫x − x = ৮}

\emph{⇒ ৪x = ৮}

\emph{∴ x = (৮/৪) = ২}

\emph{∴ পানির পরিমাণ = ২ লিটার}
\end{quote}

\textbf{261 টি আম তিন ভাইয়ের মধ্যে (1/3) : (1/5) : (1/9) অনুপাতে ভাগ করে
দিলে প্রথম ভাই কতটি আম পাবে?}

ক) 45

খ) 81

গ) 90

\textbf{\hl{ঘ) 135 ✔}}

\textbf{\hl{✔ সঠিক উত্তর: ঘ) 135}}

\begin{quote}
\textbf{ব্যাখ্যা:} \emph{প্রদত্ত অনুপাত = (1/3) : (1/5) : (1/9)}

\emph{হুরগুলোর (3, 5, 9) ল.সা.গু = 45}

\emph{অনুপাতের প্রতিটি পদকে 45 দিয়ে গুণ করে পাই,}

\emph{= (1/3) × 45 : (1/5) × 45 : (1/9) × 45 = 15 : 9 : 5}

\emph{অনুপাতের রাশিগুলোর যোগফল = 15 + 9 + 5 = 29}

\emph{∴ প্রথম ভাই আম পাবে = 261 এর (15/29) = 135টি}
\end{quote}

\textbf{৬০ মিটার দৈর্ঘ্য বিশিষ্ট বাঁশকে ৩ : ৭ : ১০ অনুপাতে ভাগ করলে টুকরো সাইজ
কত হবে?}

ক) ৮ মি. ২২ মি. ৩০ মিটার

খ) ১০ মি. ২০ মি. ৩০ মিটার

\textbf{\hl{গ) ৯ মি. ২১ মি. ৩০ মি. ✔}}

ঘ) ১২ মি. ২০ মি. ২৮ মিটার

\textbf{\hl{✔ সঠিক উত্তর: গ) ৯ মি. ২১ মি. ৩০ মি.}}

\begin{quote}
\textbf{ব্যাখ্যা:} \emph{অনুপাতের রাশি সমূহের যোগফল = ৩ + ৭ + ১০ = ২০}

\emph{∴ ১ম টুকরার দৈর্ঘ্য = ৬০ এর (৩/২০) = ৯ মিটার}

\emph{২য় টুকরার দৈর্ঘ্য = ৬০ এর (৭/২০) = ২১ মিটার}

\emph{৩য় টুকরার দৈর্ঘ্য = ৬০ এর (১০/২০) = ৩০ মিটার}
\end{quote}

\textbf{১৪৩ টাকাকে ২ : ৪ : ৫ অনুপাতে ভাগ করলে বৃহত্তম ও ক্ষুদ্রতম অংশের পার্থক্য
কত টাকা হবে?}

ক) ৪২

খ) ৩৬

গ) ৩৭

\textbf{\hl{ঘ) ৩৯ ✔}}

\textbf{\hl{✔ সঠিক উত্তর: ঘ) ৩৯}}

\begin{quote}
\textbf{ব্যাখ্যা:} \emph{অনুপাতগুলোর যোগফল = ২ + ৪ + ৫ = ১১}

\emph{ক্ষুদ্রতম অংশের পরিমাণ = ১৪৩ এর (২/১১) = ২৬ টাকা}

\emph{বৃহত্তম অংশের পরিমাণ = ১৪৩ এর (৫/১১) = ৬৫ টাকা}

\emph{∴ পার্থক্য = ৬৫ − ২৬ = ৩৯ টাকা}
\end{quote}

\textbf{২০ মিটার দীর্ঘ একটি রশিকে ২ : ৩ : ৫ অনুপাতে ভাগ করলে সবচেয়ে ছোট
টুকরাটির দৈর্ঘ্য কত মিটার হবে?}

ক) ২

\textbf{\hl{খ) ৪ ✔}}

গ) ৬

ঘ) ১০

\textbf{\hl{✔ সঠিক উত্তর: খ) ৪}}

\begin{quote}
\textbf{ব্যাখ্যা:} \emph{অনুপাতের যোগফল = ২ + ৩ + ৫ = ১০}

\emph{সবচেয়ে ছোট টুকরাটির অনুপাত = ২}

\emph{∴ সবচেয়ে ছোট টুকরাটির দৈর্ঘ্য = ২০ এর (২/১০) = ৪ মিটার}
\end{quote}

\textbf{৪০ মিটার দীর্ঘ একটি রশিকে ৩ : ৭ : ১০ অনুপাতে ভাগ করলে দীর্ঘতম
অংশটির দৈর্ঘ্য কত মিটার হবে?}

\textbf{\hl{ক) ২০ মি. ✔}}

খ) ১৪ মি.

গ) ৩০ মি.

ঘ) ১৬ মি.

\textbf{\hl{✔ সঠিক উত্তর: ক) ২০ মি.}}

\begin{quote}
\textbf{ব্যাখ্যা:} \emph{অনুপাত = ৩ : ৭ : ১০}

\emph{অনুপাতের রাশিগুলোর যোগফল = (৩ + ৭ + ১০) = ২০}

\emph{এখানে দীর্ঘতম অংশের অনুপাত ১০।}

\emph{∴ দীর্ঘতম অংশটির দৈর্ঘ্য = ৪০ × (১০/২০) = ২০ মিটার।}
\end{quote}

\textbf{কেজিপ্রতি ১২৬ টাকা মূল্যের চা এবং কেজিপ্রতি ১৩৫ টাকা মূল্যের চায়ের
সাথে তৃতীয় প্রকার মূল্যের চায়ের মিশ্রণের অনুপাত ১ : ১ : ২। মিশ্রণকৃত চায়ের মূল্য
প্রতি কেজি ১৫৩ হলে, তৃতীয় প্রকারের চায়ের মূল্য কেজিপ্রতি কত হবে?}

ক) Tk. 169.50

খ) Tk. 170

\textbf{\hl{গ) Tk. 175.50 ✔}}

ঘ) Tk. 180

\textbf{\hl{✔ সঠিক উত্তর: গ) Tk. 175.50}}

\begin{quote}
\textbf{ব্যাখ্যা:} \emph{ধরি, তিন প্রকার চায়ের পরিমাণ যথাক্রমে ১ কেজি, ১
কেজি এবং ২ কেজি।}

\emph{এবং তৃতীয় প্রকার চায়ের মূল্য x টাকা কেজি।}

\emph{মোট পরিমাণ = ১ + ১ + ২ = ৪ কেজি।}

\emph{প্রশ্নমতে, (১২৬+১৩৫+২x/৪) = ১৫৩}

\emph{⇒ ২৬১ + ২x = ৬১২}

\emph{⇒ ২x = ৬১২ − ২৬১}

\emph{⇒ ২x = ৩৫১}

\emph{∴ x = ১৭৫.৫০ Tk.}
\end{quote}

\textbf{বনি, ডলি ও লিলির মধ্যে ১,২৬০ টাকা এমনভাবে ভাগ করে দেয়া হলো যেন,
ডলি লিলির সমান টাকা পায় এবং বনি ডলির দ্বিগুণ টাকা পায়। এতে বনি কত টাকা
পেল?}

ক) ৩১৫

খ) ৫৩০

\textbf{\hl{গ) ৬৩০ ✔}}

ঘ) ৫২৫

\textbf{\hl{✔ সঠিক উত্তর: গ) ৬৩০}}

\begin{quote}
\textbf{ব্যাখ্যা:} \emph{ধরি, লিলি পায় = x টাকা}

\emph{∴ ডলি পায় = x টাকা}

\emph{এবং বনি পায় = ২x টাকা}

\emph{প্রশ্নমতে, x + x + ২x = ১২৬০}

\emph{⇒ ৪x = ১২৬০}

\emph{∴ x = ৩১৫}

\emph{∴ বনি টাকা পেল = ২ × ৩১৫ = ৬৩০ টাকা}
\end{quote}

\textbf{একটি ত্রিভুজের তিনটি কোণের অনুপাত ১ : ২ : ৩। ত্রিভুজটি হবে-}

\textbf{\hl{ক) সমকোণী ত্রিভুজ ✔}}

খ) সমবাহু ত্রিভুজ

গ) সূক্ষ্মকোণী ত্রিভুজ

ঘ) স্থূলকোণী ত্রিভুজ

\textbf{\hl{✔ সঠিক উত্তর: ক) সমকোণী ত্রিভুজ}}

\begin{quote}
\textbf{ব্যাখ্যা:} \emph{ত্রিভুজের তিন কোণের সমষ্টি = ১৮০°}

\emph{অনুপাতের রাশিগুলোর যোগফল = ১ + ২ + ৩ = ৬}

\emph{১ম কোণ = ১৮০° এর (১/৬) = ৩০°}

\emph{২য় কোণ = ১৮০° এর (২/৬) = ৬০°}

\emph{৩য় কোণ = ১৮০° এর (৩/৬) = ৯০°}

\emph{যেহেতু ত্রিভুজটির একটি কোণ সমকোণ (৯০°), সুতরাং ত্রিভুজটি সমকোণী ত্রিভুজ।}
\end{quote}

\textbf{কোন চতুর্ভুজের চার কোণের অনুপাত ১ : ২ : ২ : ৩ হলে বৃহত্তম কোণের পরিমাণ
কত?}

ক) ১০০°

খ) ১১৫°

\textbf{\hl{গ) ১৩৫° ✔}}

ঘ) ২২৫°

\textbf{\hl{✔ সঠিক উত্তর: গ) ১৩৫°}}

\begin{quote}
\textbf{ব্যাখ্যা:} \emph{চতুর্ভুজের চার কোণের সমষ্টি = ৩৬০°}

\emph{অনুপাতগুলোর সমষ্টি = ১ + ২ + ২ + ৩ = ৮}

\emph{বৃহত্তম কোণের অনুপাত = ৩}

\emph{∴ বৃহত্তম কোণ = ৩৬০° এর (৩/৮) = ১৩৫°}
\end{quote}

\textbf{দুটি সংখ্যার যোগফল ৮। যদি সংখ্যাগুলো ৩ : ১ অনুপাতে থাকে, তবে
সংখ্যাগুলোর গুণফল হবে-}

ক) ১৫

খ) ১৮

\textbf{\hl{গ) ১২ ✔}}

ঘ) ১০

\textbf{\hl{✔ সঠিক উত্তর: গ) ১২}}

\begin{quote}
\textbf{ব্যাখ্যা:} \emph{অনুপাতগুলোর যোগফল = ৩ + ১ = ৪}

\emph{একটি সংখ্যা = ৮ এর (১/৪) = ২}

\emph{অপর সংখ্যা = ৮ এর (৩/৪) = ৬}

\emph{∴ সংখ্যাগুলোর গুণফল = ৬ × ২ = ১২}
\end{quote}

\textbf{দুইটি সংখ্যার বিয়োগফল ৩৩ এবং তাদের অনুপাত ৮ : ৫, সংখ্যা দুইটি-}

ক) ৭৭, ৫৫

\textbf{\hl{খ) ৮৮, ৫৫ ✔}}

গ) ১১০, ৬৬

ঘ) ১১০, ৮৮

\textbf{\hl{✔ সঠিক উত্তর: খ) ৮৮, ৫৫}}

\begin{quote}
\textbf{ব্যাখ্যা:} \emph{ধরি, সংখ্যা দুটি ৮x এবং ৫x}

\emph{প্রশ্নমতে, ৮x − ৫x = ৩৩}

\emph{বা, ৩x = ৩৩}

\emph{∴ x = ১১}

\emph{∴ সংখ্যা দুটি যথাক্রমে (৮ × ১১) = ৮৮ এবং (৫ × ১১) = ৫৫}
\end{quote}

\textbf{পনির ও তপনের আয়ের অনুপাত 4 : 3, তপন ও রবিনের আয়ের অনুপাত 5 : 4,
পনিরের আয় 120 টাকা হলে, রবিনের আয় কত?}

ক) 36 টাকা

খ) 12 টাকা

\textbf{\hl{গ) 72 টাকা ✔}}

ঘ) 84 টাকা

\textbf{\hl{✔ সঠিক উত্তর: গ) 72 টাকা}}

\begin{quote}
\textbf{ব্যাখ্যা:} \emph{পনিরের আয় : তপনের আয় = 4 : 3 = (4 × 5) : (3 × 5)
= 20 : 15}

\emph{তপনের আয় : রবিনের আয় = 5 : 4 = (5 × 3) : (4 × 3) = 15 : 12}

\emph{∴ পনিরের আয় : তপনের আয় : রবিনের আয় = 20 : 15 : 12}

\emph{ধরি, পনির, তপন ও রবিনের আয় যথাক্রমে 20x, 15x এবং 12x টাকা।}

\emph{প্রশ্নমতে,}

\emph{20x = 120}

\emph{∴ x = 6}

\emph{∴ রবিনের আয় হবে = 12 × 6 = 72 টাকা।}
\end{quote}

\textbf{রাবু ও জামালের মাসিক বেতনের অনুপাত ৭:৫ এবং দুজনের মাসিক বেতন একত্রে
২৪,০০০ টাকা। এক বছর পরে রাবুর বেতন ৫০০ টাকা এবং জামালের বেতন ৩৫০ টাকা
বৃদ্ধি পেল। এক বছর পরে তাদের মাসিক বেতনের অনুপাত কত হবে?}

\textbf{\hl{ক) ২৯০:২০৭ ✔}}

খ) ১৪৫:১০৩

গ) ৫০০:৩৫০

ঘ) ৭:৫

\textbf{\hl{✔ সঠিক উত্তর: ক) ২৯০:২০৭}}

\begin{quote}
\textbf{ব্যাখ্যা:} \emph{অনুপাত দুটির যোগফল = ৭ + ৫ = ১২}

\emph{রাবুর বর্তমান বেতন = ২৪০০০ × (৭/১২) = ১৪০০০ টাকা}

\emph{জামালের বর্তমান বেতন = ২৪০০০ × (৫/১২) = ১০০০০ টাকা}

\emph{এক বছর পর রাবুর বেতন = ১৪০০০ + ৫০০ = ১৪৫০০ টাকা}

\emph{এবং জামালের বেতন = ১০০০০ + ৩৫০ = ১০৩৫০ টাকা}

\emph{∴ এক বছর পর তাদের মাসিক বেতনের অনুপাত = ১৪৫০০ : ১০৩৫০ = ২৯০ : ২০৭}
\end{quote}

\textbf{করিম ও রহিম এর বেতনের অনুপাত ৭ : ৫; করিমের বেতন রহিমের বেতন
অপেক্ষা ৪০০ টাকা বেশি। রহিমের বেতন কত?}

ক) ৯০০ টাকা

\textbf{\hl{খ) ১,০০০ টাকা ✔}}

গ) ১,১০০ টাকা

ঘ) ১,৬০০ টাকা

\textbf{\hl{✔ সঠিক উত্তর: খ) ১,০০০ টাকা}}

\begin{quote}
\textbf{ব্যাখ্যা:} \emph{ধরি, করিমের বেতন ৭x টাকা এবং রহিমের বেতন ৫x
টাকা।}

\emph{শর্তমতে, ৭x − ৫x = ৪০০}

\emph{বা, ২x = ৪০০}

\emph{∴ x = ২০০}

\emph{∴ রহিমের বেতন = (৫ × ২০০) টাকা = ১০০০ টাকা।}
\end{quote}

\textbf{১০০০ টাকা ক ও খ ১:৪ অনুপাতে ভাগ করে নেয়। খ এর অংশ সে এবং তার মা ও
মেয়ের মধ্যে ২:১:১ অনুপাতে ভাগ করে। মেয়ে কত টাকা পাবে?}

ক) ১০০ টাকা

খ) ৪০০ টাকা

\textbf{\hl{গ) ২০০ টাকা ✔}}

ঘ) ৮০০ টাকা

\textbf{\hl{✔ সঠিক উত্তর: গ) ২০০ টাকা}}

\begin{quote}
\textbf{ব্যাখ্যা:} \emph{ক ও খ এর টাকার অনুপাত = ১ : ৪}

\emph{অনুপাতগুলোর সমষ্টি = ১ + ৪ = ৫}

\emph{∴ খ এর অংশ = (১০০০ এর (৪/৫) ) টাকা = ৮০০ টাকা}

\emph{এই ৮০০ টাকা খ নিজে, তার মা এবং মেয়ের মধ্যে ২:১:১ অনুপাতে ভাগ করে।}

\emph{এখানে অনুপাতগুলোর সমষ্টি = ২ + ১ + ১ = ৪}

\emph{∴ খ এর মেয়ে পায় = (৮০০ এর (১/৪) ) টাকা = ২০০ টাকা}
\end{quote}

\textbf{কোন অনুপাতের উভয় পদের সঙ্গে ১ যোগ করলে অনুপাতটি ৩ : ৪ এবং উভয় পদ
থেকে ১ বিয়োগ করলে অনুপাতটি ২ : ৩ হবে?}

ক) ২ : ৫

খ) ৪ : ৯

গ) ৬ : ১১

\textbf{\hl{ঘ) ৫ : ৭ ✔}}

\textbf{\hl{✔ সঠিক উত্তর: ঘ) ৫ : ৭}}

\begin{quote}
\textbf{ব্যাখ্যা:} \emph{ধরি, অনুপাতটি ক : খ}

\emph{১ম শর্তমতে, (ক+১/খ+১) = (৩/৪)}

\emph{বা, ৪ক + ৪ = ৩খ + ৩}

\emph{∴ ৪ক − ৩খ = −১ ......... (১)}

\emph{২য় শর্তমতে, (ক−১/খ−১) = (২/৩)}

\emph{বা, ৩ক − ৩ = ২খ − ২}

\emph{∴ ৩ক − ২খ = ১ ......... (২)}

\emph{সমীকরণ (১) × ২ − (২) × ৩ করে পাই,}

\emph{(৮ক − ৬খ) − (৯ক − ৬খ) = −২ − ৩}

\emph{বা, −ক = −৫}

\emph{∴ ক = ৫}

\emph{\textquotesingle ক\textquotesingle{} এর মান সমীকরণ (১) এ বসিয়ে
পাই,}

\emph{৪ × ৫ − ৩খ = −১}

\emph{বা, −৩খ = −১ − ২০ ⇒ −৩খ = −২১}

\emph{∴ খ = ৭}

\emph{∴ ক : খ = ৫ : ৭}
\end{quote}

\textbf{ববি ও কবির কাছে কিছু আপেল আছে যার অনুপাত যথাক্রমে ৭:৯। যদি কবি
ববিকে ২১টি আপেল দিয়ে দেয়। তাহলে ববি ও কবির কাছে থাকা আপেলের অনুপাত হবে ৭
: ৬। আবার যদি ববি কবিকে ১১টি আপেল দিয়ে দেয় তাহলে ববি ও কবির কাছে থাকা
আপেলের অনুপাত হবে ৫ : ৮। কবির কাছে ববি অপেক্ষা কতটি আপেল বেশি আছে?}

ক) ১৩টি

\textbf{\hl{খ) ২৬টি ✔}}

গ) ২৮টি

ঘ) ৩২টি

\textbf{\hl{✔ সঠিক উত্তর: খ) ২৬টি}}

\begin{quote}
\textbf{ব্যাখ্যা:} \emph{ধরি, ববির আপেল = ৭x টি এবং কবির আপেল = ৯x টি}

\emph{১ম শর্তমতে, (৭x+২১/৯x−২১) = (৭/৬)}

\emph{বা, ৭(৯x − ২১) = ৬(৭x + ২১)}

\emph{বা, ৬৩x − ১৪৭ = ৪২x + ১২৬}

\emph{বা, ২১x = ২৭৩}

\emph{∴ x = ১৩}

\emph{∴ ববির কাছে আপেল আছে = ৭ × ১৩ = ৯১টি}

\emph{এবং কবির কাছে আপেল আছে = ৯ × ১৩ = ১১৭টি}

\emph{∴ কবির কাছে ববি অপেক্ষা আপেল বেশি আছে = ১১৭ − ৯১ = ২৬টি}
\end{quote}

\textbf{একটি বাক্সে কলা এবং আপেল 3 : 2 অনুপাতে আছে। বাক্স হতে 5টি কলা তুলে
নিলে তাদের অনুপাত হয় 1 : 1। বাক্সে কতগুলি আপেল ছিল?}

ক) 5

\textbf{\hl{খ) 10 ✔}}

গ) 15

ঘ) 20

\textbf{\hl{✔ সঠিক উত্তর: খ) 10}}

\begin{quote}
\textbf{ব্যাখ্যা:} \emph{ধরি, বাক্সে কলা ছিল 3xটি এবং আপেল ছিল 2xটি}

\emph{প্রশ্নমতে, (3x−5/2x) = (1/1)}

\emph{⇒ 3x − 5 = 2x}

\emph{⇒ 3x − 2x = 5}

\emph{∴ x = 5}

\emph{∴ বাক্সে আপেল ছিল = 2x = 2 × 5 = 10টি}
\end{quote}

\textbf{60 লিটার ফলের রসে আম ও কমলার অনুপাত 2 : 1। কমলার রসের পরিমাণ কত
লিটার বৃদ্ধি করলে অনুপাতটি 1 : 2 হবে?}

ক) 40

খ) 50

\textbf{\hl{গ) 60 ✔}}

ঘ) 70

\textbf{\hl{✔ সঠিক উত্তর: গ) 60}}

\begin{quote}
\textbf{ব্যাখ্যা:} \emph{অনুপাত দুটির যোগফল = 2 + 1 = 3}

\emph{আমের রসের পরিমাণ = 60 এর (2/3) = 40 লিটার}

\emph{কমলার রসের পরিমাণ = 60 এর (1/3) = 20 লিটার}

\emph{ধরি, x লিটার কমলার রস মেশাতে হবে।}

\emph{প্রশ্নমতে, (40/20+x) = (1/2)}

\emph{⇒ 20 + x = 80}

\emph{⇒ x = 80 − 20}

\emph{∴ x = 60}
\end{quote}

\textbf{৩৫ লিটার দ্রবণে পানি এবং সিরাপের অনুপাত ৩ : ৪। দ্রবণে কত লিটার
সিরাপ আছে?}

ক) ১৫ লিটার

\textbf{\hl{খ) ২০ লিটার ✔}}

গ) ২১ লিটার

ঘ) ২৪ লিটার

\textbf{\hl{✔ সঠিক উত্তর: খ) ২০ লিটার}}

\begin{quote}
\textbf{ব্যাখ্যা:} \emph{অনুপাতদ্বয়ের যোগফল = (৩ + ৪) = ৭}

\emph{∴ দ্রবণে সিরাপের পরিমাণ = ৩৫ এর (৪/৭) লিটার = ২০ লিটার।}
\end{quote}

\textbf{পায়েসে দুধ ও চিনির অনুপাত ৭ : ২। ঐ পায়েসে চিনির পরিমাণ ৪ কেজি হলে
দুধের পরিমাণ কত?}

ক) ৮ কেজি

খ) ১০ কেজি

গ) ১২ কেজি

\textbf{\hl{ঘ) ১৪ কেজি ✔}}

\textbf{\hl{✔ সঠিক উত্তর: ঘ) ১৪ কেজি}}

\begin{quote}
\textbf{ব্যাখ্যা:} \emph{ধরি, পায়েসে দুধের পরিমাণ = ৭x কেজি}

\emph{এবং চিনির পরিমাণ = ২x কেজি}

\emph{প্রশ্নমতে, ২x = ৪}

\emph{⇒ x = ২}

\emph{∴ দুধের পরিমাণ = (৭ × ২) কেজি = ১৪ কেজি}
\end{quote}

\textbf{একটি গহনার ভর ২৫ গ্রাম। এতে সোনা ও তামার ভরের অনুপাত ৪ : ১ হলে
গহনাটিতে তামার পরিমাণ কত?}

ক) ৪

\textbf{\hl{খ) ৫ ✔}}

গ) ৬

ঘ) ৭

\textbf{\hl{✔ সঠিক উত্তর: খ) ৫}}

\begin{quote}
\textbf{ব্যাখ্যা:} \emph{অনুপাতের রাশি দুটির যোগফল = (৪ + ১) = ৫}

\emph{∴ গহনাটিতে তামার পরিমাণ = (২৫ এর (১/৫) ) গ্রাম = ৫ গ্রাম}
\end{quote}

\textbf{একটি সোনার গহনার ওজন ১৬ গ্রাম। এতে সোনা ও তামার পরিমাণ ৩ : ১।
মিশ্রণে সোনার পরিমাণ কত?}

\textbf{\hl{ক) ১২ গ্রাম ✔}}

খ) ১৪ গ্রাম

গ) ১০ গ্রাম

ঘ) ৯ গ্রাম

\textbf{\hl{✔ সঠিক উত্তর: ক) ১২ গ্রাম}}

\begin{quote}
\textbf{ব্যাখ্যা:} \emph{দেয়া আছে,}

\emph{সোনা ও তামার অনুপাত = ৩ : ১}

\emph{অনুপাতের রাশিদ্বয়ের যোগফল = (৩ + ১) = ৪}

\emph{∴ মিশ্রণে সোনার পরিমাণ = (১৬ এর (৩/৪) ) গ্রাম = ১২ গ্রাম}
\end{quote}

\textbf{৬০ লিটার কেরোসিন ও পেট্রোলের মিশ্রণের অনুপাত ৭:৩। ঐ মিশ্রণে আর কত
লিটার পেট্রোল মিশালে অনুপাত ৩:৭ হবে?}

ক) ৭০

\textbf{\hl{খ) ৮০ ✔}}

গ) ৯০

ঘ) ৯৮

\textbf{\hl{✔ সঠিক উত্তর: খ) ৮০}}

\begin{quote}
\textbf{ব্যাখ্যা:} \emph{মূল নিয়ম:}

\emph{প্রদত্ত মিশ্রণের মোট পরিমাণ = ৬০ লিটার}

\emph{অনুপাত = ৭:৩ (যোগফল = ৭ + ৩ = ১০)}

\emph{কেরোসিনের পরিমাণ = ৬০×(৭/১০) = ৪২ লিটার}

\emph{পেট্রোলের পরিমাণ = ৬০×(৩/১০) = ১৮ লিটার}

\emph{ধরি, \textquotesingle ক\textquotesingle{} লিটার পেট্রোল মেশাতে
হবে।}

\emph{প্রশ্নমতে,}

\emph{⇒ (৪২/১৮+ক)=(৩/৭)}

\emph{⇒ ৩(১৮ + ক) = ৪২ × ৭}

\emph{⇒ ৫৪ + ৩ক = ২৯৪}

\emph{⇒ ৩ক = ২৯৪ − ৫৪}

\emph{⇒ ৩ক = ২৪০}

\emph{⇒ ক = ৮০}

\emph{শর্টকাট নিয়ম:}

\emph{উভয় অনুপাতের পদগুলো বিপরীত (৭:৩ এবং ৩:৭) হলে,}

\emph{পেট্রোল মেশাতে হবে = (মোট মিশ্রণ/প্রথম অনুপাতের ছোট সংখ্যা)×অনুপাতের
পার্থক্য}

\emph{= (৬০/৩)×(৭−৩)}

\emph{= ২০ × ৪ = ৮০ লিটার}

\emph{∴ সঠিক উত্তর ৮০।}
\end{quote}

\textbf{২১ গ্যালন অকটেন মিশ্রিত পেট্রোলে, পেট্রোল ও অকটনের অনুপাত ৪:৩। এতে আর
কত অকটেন মিশালে পেট্রোল ও অকটেনের অনুপাত ৩:৪ হবে?}

ক) ৩ গ্যালন

খ) ৪ গ্যালন

\textbf{\hl{গ) ৭ গ্যালন ✔}}

ঘ) ৬ গ্যালন

\textbf{\hl{✔ সঠিক উত্তর: গ) ৭ গ্যালন}}

\begin{quote}
\textbf{ব্যাখ্যা:} \emph{মূল নিয়ম:}

\emph{মোট মিশ্রণের পরিমাণ = ২১ গ্যালন}

\emph{অনুপাতের যোগফল = ৪ + ৩ = ৭}

\emph{পেট্রোলের পরিমাণ = ২১×(৪/৭) = ১২ গ্যালন}

\emph{অকটেনের পরিমাণ = ২১×(৩/৭) = ৯ গ্যালন}

\emph{ধরি, \textquotesingle ক\textquotesingle{} গ্যালন অকটেন মেশাতে হবে।}

\emph{প্রশ্নমতে,}

\emph{⇒ (১২/৯+ক)=(৩/৪)}

\emph{⇒ ৩(৯ + ক) = ৪৮}

\emph{⇒ ২৭ + ৩ক = ৪৮}

\emph{⇒ ৩ক = ২১}

\emph{⇒ ক = ৭}

\emph{শর্টকাট নিয়ম:}

\emph{অকটেন মেশাতে হবে = (মোট মিশ্রণ/প্রথম অনুপাতের ছোট সংখ্যা)×অনুপাতের
পার্থক্য}

\emph{= (২১/৩)×(৪−৩)}

\emph{= ৭ × ১ = ৭ গ্যালন}

\emph{∴ সঠিক উত্তর ৭ গ্যালন।}
\end{quote}

\textbf{৬০ লিটার ফলের রসে আম ও কমলার অনুপাত ২:১। কমলার রসের পরিমাণ কত
লিটার বৃদ্ধি করলে অনুপাতটি ১:২ হবে?}

ক) ৫০

\textbf{\hl{খ) ৬০ ✔}}

গ) ৭০

ঘ) ৪০

\textbf{\hl{✔ সঠিক উত্তর: খ) ৬০}}

\begin{quote}
\textbf{ব্যাখ্যা:} \emph{মূল নিয়ম:}

\emph{মোট ফলের রস = ৬০ লিটার}

\emph{অনুপাতের যোগফল = ২ + ১ = ৩}

\emph{আমের রস = ৬০×(২/৩) = ৪০ লিটার}

\emph{কমলার রস = ৬০×(১/৩) = ২০ লিটার}

\emph{ধরি, \textquotesingle ক\textquotesingle{} লিটার কমলার রস বৃদ্ধি করতে
হবে।}

\emph{প্রশ্নমতে,}

\emph{⇒ (৪০/২০+ক)=(১/২)}

\emph{⇒ ২০ + ক = ৮০}

\emph{⇒ ক = ৮০ − ২০}

\emph{⇒ ক = ৬০}

\emph{শর্টকাট নিয়ম:}

\emph{কমলার রস বৃদ্ধি করতে হবে = (মোট মিশ্রণ/প্রথম অনুপাতের ছোট
সংখ্যা)×অনুপাতের পার্থক্য}

\emph{= (৬০/১)×(২−১)}

\emph{= ৬০ × ১ = ৬০ লিটার}

\emph{∴ সঠিক উত্তর ৬০।}
\end{quote}

\textbf{৯০ লিটার কেরোসিন ও পেট্রোলের মিশ্রণের অনুপাত ৭:৩। ঐ মিশ্রণে আর কত
লিটার পেট্রোল মিশালে অনুপাত ৩:৭ হবে?}

\textbf{\hl{ক) ১২০ লিটার ✔}}

খ) ১২২ লিটার

গ) ১২৪ লিটার

ঘ) ১২৬ লিটার

\textbf{\hl{✔ সঠিক উত্তর: ক) ১২০ লিটার}}

\begin{quote}
\textbf{ব্যাখ্যা:} \emph{মূল নিয়ম:}

\emph{মোট মিশ্রণের পরিমাণ = ৯০ লিটার}

\emph{অনুপাতের যোগফল = ৭ + ৩ = ১০}

\emph{কেরোসিনের পরিমাণ = ৯০×(৭/১০) = ৬৩ লিটার}

\emph{পেট্রোলের পরিমাণ = ৯০×(৩/১০) = ২৭ লিটার}

\emph{ধরি, \textquotesingle ক\textquotesingle{} লিটার পেট্রোল মেশাতে
হবে।}

\emph{প্রশ্নমতে,}

\emph{⇒ (৬৩/২৭+ক)=(৩/৭)}

\emph{⇒ ৩(২৭ + ক) = ৬৩ × ৭}

\emph{⇒ ৮১ + ৩ক = ৪৪১}

\emph{⇒ ৩ক = ৪৪১ − ৮১}

\emph{⇒ ৩ক = ৩৬০}

\emph{⇒ ক = ১২০}

\emph{শর্টকাট নিয়ম:}

\emph{পেট্রোল মেশাতে হবে = (মোট মিশ্রণ/প্রথম অনুপাতের ছোট সংখ্যা)×অনুপাতের
পার্থক্য}

\emph{= (৯০/৩)×(৭−৩)}

\emph{= ৩০ × ৪ = ১২০ লিটার}

\emph{∴ সঠিক উত্তর ১২০ লিটার।}
\end{quote}

\textbf{৪২ গ্রাম ওজনের একটি গয়নার সোনা ও তামার অনুপাত ৪:৩। এতে আর কত সোনা
মিশালে সোনা ও তামার অনুপাত ৫:৩ হবে?}

\textbf{\hl{ক) ৬ ✔}}

খ) ৭

গ) ৮

ঘ) ৯

\textbf{\hl{✔ সঠিক উত্তর: ক) ৬}}

\begin{quote}
\textbf{ব্যাখ্যা:} \emph{মূল নিয়ম:}

\emph{মোট গয়নার ওজন = ৪২ গ্রাম}

\emph{অনুপাতের যোগফল = ৪ + ৩ = ৭}

\emph{সোনার পরিমাণ = ৪২×(৪/৭) = ২৪ গ্রাম}

\emph{তামার পরিমাণ = ৪২×(৩/৭) = ১৮ গ্রাম}

\emph{ধরি, \textquotesingle ক\textquotesingle{} গ্রাম সোনা মেশাতে হবে।}

\emph{প্রশ্নমতে,}

\emph{⇒ (২৪+ক/১৮)=(৫/৩)}

\emph{⇒ ৩(২৪ + ক) = ৯০}

\emph{⇒ ৭২ + ৩ক = ৯০}

\emph{⇒ ৩ক = ৯০ − ৭২}

\emph{⇒ ৩ক = ১৮}

\emph{⇒ ক = ৬}

\emph{শর্টকাট নিয়ম:}

\emph{সোনা মেশাতে হবে = (মোট মিশ্রণের পরিমাণ/প্রথম অনুপাতের যোগফল)×পরিবর্তিত
অনুপাত সংখ্যার পার্থক্য}

\emph{= (৪২/৪+৩)×(৫−৪)}

\emph{= (৪২/৭)×১ = ৬ গ্রাম}

\emph{∴ সঠিক উত্তর ৬।}
\end{quote}

\textbf{৪০ গ্যালন অকটেন মিশ্রিত পেট্রোলে, পেট্রোল ও অকটেনের অনুপাত ৩:১। এতে
আর কত অকটেন মিশালে পেট্রোল ও অকটেনের অনুপাত ৫:২ হবে?}

ক) ৩ গ্যালন

\textbf{\hl{খ) ২ গ্যালন ✔}}

গ) ৫ গ্যালন

ঘ) ৮ গ্যালন

\textbf{\hl{✔ সঠিক উত্তর: খ) ২ গ্যালন}}

\begin{quote}
\textbf{ব্যাখ্যা:} \emph{মূল নিয়ম:}

\emph{মোট মিশ্রণের পরিমাণ = ৪০ গ্যালন}

\emph{অনুপাতের যোগফল = ৩ + ১ = ৪}

\emph{পেট্রোলের পরিমাণ = ৪০×(৩/৪) = ৩০ গ্যালন}

\emph{অকটেনের পরিমাণ = ৪০×(১/৪) = ১০ গ্যালন}

\emph{ধরি, \textquotesingle ক\textquotesingle{} গ্যালন অকটেন মেশাতে হবে।}

\emph{প্রশ্নমতে,}

\emph{⇒ (৩০/১০+ক)=(৫/২)}

\emph{⇒ ৫(১০ + ক) = ৬০}

\emph{⇒ ৫০ + ৫ক = ৬০}

\emph{⇒ ৫ক = ১০}

\emph{⇒ ক = ২}

\emph{শর্টকাট নিয়ম:}

\emph{পেট্রোল ও অকটেনের প্রথম অনুপাত ৩:১ এবং দ্বিতীয় অনুপাত ৫:২।}

\emph{পেট্রোলের পরিমাণ ৩০ গ্যালন অপরিবর্তিত থাকে।}

\emph{নতুন অনুপাতে ৫ ভাগ = ৩০ গ্যালন, সুতরাং ১ ভাগ = ৬ গ্যালন।}

\emph{নতুন অনুপাতে অকটেন ২ ভাগ = ১২ গ্যালন।}

\emph{∴ অকটেন মেশাতে হবে = ১২ − ১০ = ২ গ্যালন।}
\end{quote}

\textbf{একটি ঝুড়িতে ১৩০ টি আম ও পেয়ারা আছে এবং এদের অনুপাত যথাক্রমে ৩:২।
উক্ত ঝুড়িতে আম ও পেয়ারার অনুপাত ১:১ করতে হলে, কমপক্ষে কতটি নতুন ফল যোগ করতে
হবে?}

ক) ৫২টি

খ) ১৩টি

গ) ১৩০টি

\textbf{\hl{ঘ) ২৬টি ✔}}

ঙ) কোনটিই নয়

\textbf{\hl{✔ সঠিক উত্তর: ঘ) ২৬টি}}

\begin{quote}
\textbf{ব্যাখ্যা:} \emph{মূল নিয়ম:}

\emph{মোট ফলের সংখ্যা = ১৩০টি}

\emph{অনুপাতের যোগফল = ৩ + ২ = ৫}

\emph{আমের সংখ্যা = ১৩০×(৩/৫) = ৭৮টি}

\emph{পেয়ারার সংখ্যা = ১৩০×(২/৫) = ৫২টি}

\emph{অনুপাত ১:১ করতে হলে আম ও পেয়ারার সংখ্যা সমান হতে হবে।}

\emph{কমপক্ষে ফল যোগ করতে হলে কম থাকা ফলটির (পেয়ারা) সংখ্যা বাড়িয়ে আমের
সংখ্যার (৭৮টি) সমান করতে হবে।}

\emph{⇒ নতুন ফল যোগ করতে হবে = ৭৮ − ৫২ = ২৬টি}

\emph{শর্টকাট নিয়ম:}

\emph{নতুন ফল যোগ করতে হবে = (মোট ফলের পরিমাণ/প্রথম অনুপাতের যোগফল)×প্রথম
অনুপাতের পদের পার্থক্য}

\emph{= (১৩০/৫)×(৩−২)}

\emph{= ২৬ × ১ = ২৬টি}

\emph{∴ সঠিক উত্তর ২৬টি।}
\end{quote}

\textbf{৪০ কেজি জৈব ও ইউরিয়া সারের মিশ্রণে জৈব সারের পরিমাণ ১০\%। কত কেজি
জৈব সার মিশালে নতুন মিশ্রণে জৈব সারের পরিমাণ ২০\% হবে?}

ক) ৪

\textbf{\hl{খ) ৫ ✔}}

গ) ৬

ঘ) ৭

ঙ) কোনটিই নয়

\textbf{\hl{✔ সঠিক উত্তর: খ) ৫}}

\begin{quote}
\textbf{ব্যাখ্যা:} \emph{মূল নিয়ম:}

\emph{জৈব সারের পরিমাণ = ৪০ এর ১০\% = ৪ কেজি}

\emph{ইউরিয়া সারের পরিমাণ = ৪০ − ৪ = ৩৬ কেজি}

\emph{ধরি, \textquotesingle ক\textquotesingle{} কেজি জৈব সার মেশাতে
হবে।}

\emph{নতুন মিশ্রণে জৈব সার ২০\% হলে ইউরিয়া সার হবে (১০০ − ২০)\% = ৮০\%।}

\emph{প্রশ্নমতে,}

\emph{⇒ (৪+ক/৩৬)=(২০/৮০)}

\emph{⇒ (৪+ক/৩৬)=(১/৪)}

\emph{⇒ ৪(৪ + ক) = ৩৬}

\emph{⇒ ১৬ + ৪ক = ৩৬}

\emph{⇒ ৪ক = ২০}

\emph{⇒ ক = ৫}

\emph{শর্টকাট নিয়ম:}

\emph{সার মেশাতে হবে = (মোট মিশ্রণ×(নতুন শতাংশ−পুরাতন শতাংশ)/১০০−নতুন
শতাংশ)}

\emph{= (৪০×(২০−১০)/১০০−২০)}

\emph{= (৪০×১০/৮০) = ৫ কেজি}

\emph{∴ সঠিক উত্তর ৫।}
\end{quote}

\textbf{একজন দোকানদার ১৫ টাকা ও ২০ টাকা কেজি দরে ২ ধরনের চা কি অনুপাতে
মেশালে ১৭ টাকা দাম হবে?}

ক) ৫:২০

খ) ৪:১

\textbf{\hl{গ) ৩:২ ✔}}

ঘ) ৩:১

\textbf{\hl{✔ সঠিক উত্তর: গ) ৩:২}}

\begin{quote}
\textbf{ব্যাখ্যা:} \emph{মূল নিয়ম:}

\emph{ধরি, ১৫ টাকা কেজি দরের চা মেশাতে হবে
\textquotesingle ক\textquotesingle{} কেজি এবং ২০ টাকা কেজি দরের চা
মেশাতে হবে \textquotesingle খ\textquotesingle{} কেজি।}

\emph{∴ মোট পরিমাণ = (ক + খ) কেজি}

\emph{ক কেজি চায়ের দাম = ১৫ক টাকা এবং খ কেজি চায়ের দাম = ২০খ টাকা}

\emph{মিশ্রিত চায়ের মোট দাম = (১৫ক + ২০খ) টাকা}

\emph{প্রশ্নমতে,}

\emph{⇒ ১৫ক + ২০খ = ১৭(ক + খ)}

\emph{⇒ ১৫ক + ২০খ = ১৭ক + ১৭খ}

\emph{⇒ ২০খ − ১৭খ = ১৭ক − ১৫ক}

\emph{⇒ ৩খ = ২ক}

\emph{⇒ ২ক = ৩খ}

\emph{⇒ (ক/খ)=(৩/২)}

\emph{⇒ ক : খ = ৩:২}

\emph{শর্টকাট নিয়ম (Alligation Method):}

\emph{প্রথম ধরন = ১৫, দ্বিতীয় ধরন = ২০, গড় দাম = ১৭}

\emph{অনুপাত = (দ্বিতীয় দাম − গড় দাম) : (গড় দাম − প্রথম দাম)}

\emph{= (২০ − ১৭) : (১৭ − ১৫)}

\emph{= ৩:২}

\emph{∴ সঠিক উত্তর ৩:২}
\end{quote}

\textbf{৫০০০ টাকা/গ্রাম দামের সোনার সাথে ৮০০ টাকা/গ্রাম দামের তামা কী
অনুপাতে মিশ্রিত করলে প্রতি গ্রামের দাম ২০০০ টাকা হবে?}

\textbf{\hl{ক) ২:৫ ✔}}

খ) ৪:৩

গ) ৫:৬

ঘ) ২:৩

\textbf{\hl{✔ সঠিক উত্তর: ক) ২:৫}}

\begin{quote}
\textbf{ব্যাখ্যা:} \emph{মূল নিয়ম:}

\emph{ধরি, সোনার পরিমাণ \textquotesingle ক\textquotesingle{} গ্রাম এবং
তামার পরিমাণ \textquotesingle খ\textquotesingle{} গ্রাম।}

\emph{ক গ্রাম সোনার দাম = ৫০০০ক টাকা এবং খ গ্রাম তামার দাম = ৮০০খ টাকা।}

\emph{প্রশ্নমতে,}

\emph{⇒ ৫০০০ক + ৮০০খ = ২০০০(ক + খ)}

\emph{⇒ ৫০০০ক + ৮০০খ = ২০০০ক + ২০০০খ}

\emph{⇒ ৫০০০ক − ২০০০ক = ২০০০খ − ৮০০খ}

\emph{⇒ ৩০০০ক = ১২০০খ}

\emph{⇒ (ক/খ)=(১২০০/৩০০০)}

\emph{⇒ (ক/খ)=(২/৫)}

\emph{⇒ ক : খ = ২:৫}

\emph{শর্টকাট নিয়ম (Alligation Method):}

\emph{সোনার অনুপাত : তামার অনুপাত}

\emph{= (তামার দামের পার্থক্য) : (সোনার দামের পার্থক্য)}

\emph{= (২০০০ − ৮০০) : (৫০০০ − ২০০০)}

\emph{= ১২০০ : ৩০০০}

\emph{= ২:৫}

\emph{∴ সঠিক উত্তর ২:৫}
\end{quote}

\textbf{৫০ টাকা প্রতি লিটার দুধের সাথে কী অনুপাতে পানি মিশ্রিত করলে দাম ৩০
টাকা হবে?}

ক) ৩:৪

\textbf{\hl{খ) ৩:২ ✔}}

গ) ২:৩০

ঘ) ৫:৩

\textbf{\hl{✔ সঠিক উত্তর: খ) ৩:২}}

\begin{quote}
\textbf{ব্যাখ্যা:} \emph{মূল নিয়ম:}

\emph{ধরি, \textquotesingle ক\textquotesingle{} লিটার দুধের সাথে
\textquotesingle খ\textquotesingle{} লিটার পানি মেশাতে হবে।}

\emph{দুধের দাম = ৫০ টাকা/লিটার এবং পানির দাম = ০ টাকা/লিটার।}

\emph{প্রশ্নমতে,}

\emph{⇒ ৫০ক + ০ × খ = ৩০(ক + খ)}

\emph{⇒ ৫০ক = ৩০ক + ৩০খ}

\emph{⇒ ৫০ক − ৩০ক = ৩০খ}

\emph{⇒ ২০ক = ৩০খ}

\emph{⇒ (ক/খ)=(৩০/২০)}

\emph{⇒ (ক/খ)=(৩/২)}

\emph{⇒ ক : খ = ৩:২}

\emph{শর্টকাট নিয়ম (Alligation Method):}

\emph{দুধের দাম = ৫০ টাকা, পানির দাম = ০ টাকা, মিশ্রণের দাম = ৩০ টাকা।}

\emph{দুধ ও পানির অনুপাত = (৩০ − ০) : (৫০ − ৩০)}

\emph{= ৩০ : ২০}

\emph{= ৩:২}

\emph{∴ সঠিক উত্তর ৩:২।}
\end{quote}

\textbf{একটি সোনার গহনার ওজন ১৬ গ্রাম। তাতে সোনা ও তামার পরিমাণ ৩ : ১।
তাতে কি পরিমাণ সোনা মেশালে অনুপাত ৪ : ১ হবে?}

ক) ৬ গ্রাম

খ) ৫ গ্রাম

\textbf{\hl{গ) ৪ গ্রাম ✔}}

ঘ) ৮ গ্রাম

\textbf{\hl{✔ সঠিক উত্তর: গ) ৪ গ্রাম}}

\begin{quote}
\textbf{ব্যাখ্যা:} \emph{অনুপাতগুলোর যোগফল = ৩ + ১ = ৪}

\emph{অনুপাত অনুযায়ী,}

\emph{সোনার পরিমাণ = (১৬ এর (৩/৪) ) গ্রাম = ১২ গ্রাম}

\emph{তামার পরিমাণ = (১৬ এর (১/৪) ) গ্রাম = ৪ গ্রাম}

\emph{ধরি, x গ্রাম সোনা মেশাতে হবে।}

\emph{প্রশ্নমতে, (১২+x/৪) = (৪/১)}

\emph{⇒ ১২ + x = ১৬}

\emph{⇒ x = ১৬ − ১২}

\emph{∴ x = ৪ গ্রাম}
\end{quote}

\textbf{২৫ গ্রাম ওজনের একটি সোনার গহনায় সোনা ও তামার অনুপাত ৪ : ১।
গহনাটিতে আর কতটুকু সোনা মেশালে এতে সোনা ও তামার অনুপাত ৫:১ হবে?}

\textbf{\hl{ক) ৫ গ্রাম ✔}}

খ) ৬ গ্রাম

গ) ১০ গ্রাম

ঘ) ২০ গ্রাম

\textbf{\hl{✔ সঠিক উত্তর: ক) ৫ গ্রাম}}

\begin{quote}
\textbf{ব্যাখ্যা:} \emph{সোনা ও তামার অনুপাত = ৪ : ১}

\emph{অনুপাতদ্বয়ের যোগফল = ৪ + ১ = ৫}

\emph{সোনার পরিমাণ = (৪/৫) × ২৫ = ২০ গ্রাম}

\emph{তামার পরিমাণ = (১/৫) × ২৫ = ৫ গ্রাম}

\emph{ধরি, x গ্রাম সোনা মেশাতে হবে।}

\emph{নতুন অনুপাত, সোনা : তামা = ৫ : ১}

\emph{প্রশ্নমতে, (২০+x/৫) = (৫/১)}

\emph{⇒ ২০ + x = ২৫}

\emph{⇒ x = ২৫ − ২০}

\emph{∴ x = ৫ গ্রাম}
\end{quote}

\textbf{৬৪ কেজি বালি ও পাথরের টুকরার মিশ্রণে বালির পরিমাণ শতকরা ২৫ ভাগ। কত
কেজি বালি মিশালে নতুন মিশ্রণে পাথর টুকরার পরিমাণ শতকরা ৪০ ভাগ হবে?}

ক) ১.৬ কেজি

\textbf{\hl{খ) ৫৬ কেজি ✔}}

গ) ১১ কেজি

ঘ) ৪৮ কেজি

\textbf{\hl{✔ সঠিক উত্তর: খ) ৫৬ কেজি}}

\begin{quote}
\textbf{ব্যাখ্যা:} \emph{মিশ্রণে বালুর পরিমাণ = (২৫/১০০) × ৬৪ = ১৬ কেজি}

\emph{মিশ্রণে পাথরের টুকরার পরিমাণ = (৬৪ − ১৬) = ৪৮ কেজি}

\emph{নতুন মিশ্রণে পাথর ৪০\% হলে, বালি হবে (১০০ − ৪০)\% = ৬০\%}

\emph{∴ নতুন মিশ্রণে পাথর ও বালুর অনুপাত = ৪০ : ৬০ = ২ : ৩}

\emph{ধরি, অতিরিক্ত x কেজি বালি মেশাতে হবে।}

\emph{প্রশ্নমতে, (৪৮/১৬+x) = (২/৩)}

\emph{⇒ ২(১৬ + x) = ১৪৪}

\emph{⇒ ৩২ + ২x = ১৪৪}

\emph{⇒ ২x = ১১২}

\emph{∴ x = ৫৬ কেজি}
\end{quote}

\textbf{৩৫ লিটার অকটেন-পেট্রোল মিশ্রণে, পেট্রোল ও অকটেনের অনুপাত ৪ : ৩। এতে
আর কত অকটেন মেশালে পেট্রোল ও অকটেনের অনুপাত ৪ : ৫ হবে?}

ক) ৮ লিটার

\textbf{\hl{খ) ১০ লিটার ✔}}

গ) ১৫ লিটার

ঘ) ২০ লিটার

\textbf{\hl{✔ সঠিক উত্তর: খ) ১০ লিটার}}

\begin{quote}
\textbf{ব্যাখ্যা:} \emph{অনুপাতের যোগফল = ৪ + ৩ = ৭}

\emph{৩৫ লিটার মিশ্রণে পেট্রোলের পরিমাণ = (৩৫ এর (৪/৭) ) লিটার = ২০ লিটার}

\emph{৩৫ লিটার মিশ্রণে অকটেনের পরিমাণ = (৩৫ − ২০) লিটার = ১৫ লিটার}

\emph{ধরি, মিশ্রণে x লিটার অকটেন মেশালে পেট্রোল ও অকটেনের অনুপাত ৪ : ৫
হবে।}

\emph{প্রশ্নমতে, ২০ : (১৫ + x) = ৪ : ৫}

\emph{বা, (২০/১৫+x) = (৪/৫)}

\emph{বা, ৬০ + ৪x = ১০০}

\emph{বা, ৪x = ৪০}

\emph{∴ x = ১০ লিটার}
\end{quote}

\textbf{৩০ লিটার পরিমাণ মিশ্রণে এসিড ও পানির অনুপাত ৭ : ৩। ঐ মিশ্রণে কি
পরিমাণ পানি মিশ্রিত করলে এসিড ও পানির অনুপাত ৩ : ৭ হবে?}

ক) ৩০ লিটার

খ) ২৫ লিটার

\textbf{\hl{গ) ৪০ লিটার ✔}}

ঘ) ৩৫ লিটার

\textbf{\hl{✔ সঠিক উত্তর: গ) ৪০ লিটার}}

\begin{quote}
\textbf{ব্যাখ্যা:} \emph{৩০ লিটার মিশ্রণে এসিড : পানি = ৭ : ৩}

\emph{অনুপাতগুলোর সমষ্টি = ৭ + ৩ = ১০}

\emph{৩০ লিটার মিশ্রণে এসিডের পরিমাণ = (৩০ এর (৭/১০) ) লিটার = ২১ লিটার}

\emph{৩০ লিটার মিশ্রণে পানির পরিমাণ = (৩০ এর (৩/১০) ) লিটার = ৯ লিটার}

\emph{ধরি, x লিটার পানি মিশ্রিত করায় অনুপাত ৩:৭ হয়।}

\emph{শর্তমতে, ২১ : (৯ + x) = ৩ : ৭}

\emph{বা, (২১/৯+x) = (৩/৭)}

\emph{বা, ২৭ + ৩x = ১৪৭}

\emph{বা, ৩x = ১২০}

\emph{∴ x = ৪০ লিটার}
\end{quote}

\textbf{৬৩ লিটার মিশ্রণে এসিড ও পানির অনুপাত ৭:২ ঐ মিশ্রণে এসিড ও পানির
পরিমাণ কত?}

ক) ৩৯ লি. ২৪ লি

\textbf{\hl{খ) ৪৯ লি. ১৪ লি. ✔}}

গ) ২৪ লি. ৩৯ লি

ঘ) ২৯ লি. ৩৪ লি

\textbf{\hl{✔ সঠিক উত্তর: খ) ৪৯ লি. ১৪ লি.}}

\begin{quote}
\textbf{ব্যাখ্যা:} \emph{অনুপাতগুলোর সমষ্টি = ৭ + ২ = ৯}

\emph{৬৩ লিটার মিশ্রণে এসিডের পরিমাণ = (৬৩ এর (৭/৯) ) লিটার = ৪৯ লিটার}

\emph{৬৩ লিটার মিশ্রণে পানির পরিমাণ = (৬৩ এর (২/৯) ) লিটার = ১৪ লিটার}
\end{quote}

\textbf{৭২ কেজি ওজন বিশিষ্ট একটি মিশ্রণ A এর ১৭ ভাগ, B এর ৩ ভাগ, C এর ৪
ভাগ দ্বারা গঠিত। মিশ্রণে B কতটুকু আছে?}

\textbf{\hl{ক) ৯ কেজি ✔}}

খ) ১২ কেজি

গ) ১৭ কেজি

ঘ) ৫১ কেজি

\textbf{\hl{✔ সঠিক উত্তর: ক) ৯ কেজি}}

\begin{quote}
\textbf{ব্যাখ্যা:} \emph{মিশ্রণে A, B, C এর অনুপাত = ১৭ : ৩ : ৪}

\emph{অনুপাতগুলোর যোগফল = ১৭ + ৩ + ৪ = ২৪}

\emph{∴ মিশ্রণে B এর পরিমাণ = (৭২ এর (৩/২৪) ) কেজি = ৯ কেজি}
\end{quote}

\textbf{শরবতে চিনি ও পানির অনুপাত ২ : ১৩ হলে, ২৬ কাপ শরবতে চিনির পরিমাণ
কত?}

ক) ২ কাপ

খ) ৩ কাপ

\textbf{\hl{গ) ৪ কাপ ✔}}

ঘ) ৫ কাপ

ঙ) কোনোটিই নয়

\textbf{\hl{✔ সঠিক উত্তর: গ) ৪ কাপ}}

\begin{quote}
\textbf{ব্যাখ্যা:} \emph{অনুপাতের যোগফল = ২ + ১৩ = ১৫}

\emph{অথবা, ছবিতে প্রদত্ত সমাধান অনুযায়ী:}

\emph{১৩ ভাগ পানির সাপেক্ষে চিনি ২ ভাগ হিসেবে,}

\emph{১ কাপ শরবতে চিনির পরিমাণ = (২/১৩) কাপ}

\emph{∴ ২৬ কাপ শরবতে চিনির পরিমাণ = (২৬×২/১৩) = ৪ কাপ}
\end{quote}

\textbf{এক দোকানদার ১৫ টাকা ও ২০ টাকা কেজি দরে দু\textquotesingle ধরনের
চা কি অনুপাতে মেশালে দাম ১৬ টাকা ৫০ পয়সা হবে?}

\textbf{\hl{ক) ৭ : ৩ ✔}}

খ) ৫ : ৭

গ) ৩ : ৭

ঘ) ৪ : ৫

\textbf{\hl{✔ সঠিক উত্তর: ক) ৭ : ৩}}

\begin{quote}
\textbf{ব্যাখ্যা:} \emph{ধরি, প্রথম ধরনের চা x কেজি এবং দ্বিতীয় ধরনের চা y
কেজি মেশাতে হবে।}

\emph{প্রশ্নমতে, ১৫x + ২০y = ১৬.৫০(x + y)}

\emph{⇒ ১৫x + ২০y = ১৬.৫০x + ১৬.৫০y}

\emph{⇒ ২০y − ১৬.৫০y = ১৬.৫০x − ১৫x}

\emph{⇒ ৩.৫০y = ১.৫০x}

\emph{⇒ (x/y) = (৩.৫০/১.৫০) = (৩৫/১৫) = (৭/৩)}

\emph{∴ x : y = ৭ : ৩}
\end{quote}

\textbf{50 লিটারের মিশ্রণে 3\% চিনি আছে। ঐ মিশ্রণ হতে কি পরিমাণ পানি বাষ্প
করলে মিশ্রণে চিনির পরিমাণ 5\% হবে?}

ক) 6

খ) 8

গ) 10

\textbf{\hl{ঘ) 20 ✔}}

\textbf{\hl{✔ সঠিক উত্তর: ঘ) 20}}

\begin{quote}
\textbf{ব্যাখ্যা:} \emph{50 লিটার মিশ্রণে চিনির পরিমাণ = 50 এর 3\% = 1.5
লিটার}

\emph{ধরি, x লিটার পানি বাষ্পীভবন করতে হবে।}

\emph{নতুন মিশ্রণের পরিমাণ = (50 − x) লিটার}

\emph{প্রশ্নমতে, (50 − x) এর 5\% = 1.5}

\emph{⇒ (50 − x) × (5/100) = 1.5}

\emph{⇒ 50 − x = 30}

\emph{⇒ x = 50 − 30}

\emph{∴ x = 20}
\end{quote}

\textbf{সুমন 8 টাকা এবং 12 টাকা কেজি দরের কিছু চাল কিনে মেশালো। মিশ্রিত চাল
12 টাকা কেজি দরে বিক্রয় করে 20\% লাভ করলো। চাল মিশ্রণের অনুপাত কত?}

ক) 2:1

খ) 1:2

\textbf{\hl{গ) 1:1 ✔}}

ঘ) 3.1

\textbf{\hl{✔ সঠিক উত্তর: গ) 1:1}}

\begin{quote}
\textbf{ব্যাখ্যা:} \emph{ধরি, 8 টাকা কেজি দরের চালের পরিমাণ x কেজি এবং 12
টাকা কেজি দরের চালের পরিমাণ y কেজি।}

\emph{মোট ক্রয়মূল্য = (8x + 12y) টাকা}

\emph{মিশ্রিত চালের মোট পরিমাণ = (x + y) কেজি এবং মোট বিক্রয়মূল্য = 12(x +
y) টাকা}

\emph{20\% লাভে,}

\emph{(8x + 12y) এর 120\% = 12(x + y)}

\emph{⇒ (8x + 12y) × (120/100) = 12(x + y)}

\emph{⇒ (8x + 12y) × (6/5) = 12(x + y)}

\emph{⇒ 6(8x + 12y) = 60(x + y)}

\emph{⇒ 48x + 72y = 60x + 60y}

\emph{⇒ 72y − 60y = 60x − 48x}

\emph{⇒ 12y = 12x}

\emph{⇒ x = y}

\emph{⇒ (x/y) = 1}

\emph{∴ x : y = 1 : 1}
\end{quote}

\textbf{একটি ব্যবসায় মুনাফা 3 জন অংশীদারের মধ্যে 8 : 7 : 5 অনুপাতে বণ্টিত হয়।
তারা যথাক্রমে 7, 8 এবং 14 মাসের জন্য বিনিয়োগ করেছিল। তাদের বিনিয়োগের অনুপাত
কত?}

ক) 8:7:5

খ) 56:28:17

গ) 21 : 28 : 38

\textbf{\hl{ঘ) 64 : 49 : 20 ✔}}

\textbf{\hl{✔ সঠিক উত্তর: ঘ) 64 : 49 : 20}}

\begin{quote}
\textbf{ব্যাখ্যা:} \emph{আমরা জানি, বিনিয়োগ = (মুনাফা/বিনিয়োগকৃত সময়)}

\emph{∴ বিনিয়োগের অনুপাত = (8/7) : (7/8) : (5/14)}

\emph{এখানে 7, 8 এবং 14 এর ল.সা.গু. = 56}

\emph{অনুপাতগুলোকে 56 দিয়ে গুণ করে পাই,}

\emph{( (8/7) × 56) : ( (7/8) × 56) : ( (5/14) × 56) = 64 : 49 : 20}
\end{quote}

\textbf{মামুন ও সবুজ একটি অংশীদারি ব্যবসায় ২ : ১ অনুপাতে বিনিয়োগ করে। বছর
শেষে তারা তাদের মুনাফা থেকে ৯০০ টাকা গরিবদের দান করে দেয় এবং অবশিষ্ট মুনাফা
মূলধন অনুপাতে ভাগ করে নেয়। সবুজের মুনাফা ২,০০০ টাকা হলে, মোট মুনাফা কত?}

ক) ২,৯০০ টাকা

খ) ৬,০০০ টাকা

\textbf{\hl{গ) ৬,৯০০ টাকা ✔}}

ঘ) ৭,২০০ টাকা

\textbf{\hl{✔ সঠিক উত্তর: গ) ৬,৯০০ টাকা}}

\begin{quote}
\textbf{ব্যাখ্যা:} \emph{ধরি, মোট মুনাফা = x টাকা}

\emph{গরিবদের দান করার পর অবশিষ্ট মুনাফা = (x − ৯০০) টাকা}

\emph{মামুন ও সবুজের অনুপাতের যোগফল = ২ + ১ = ৩}

\emph{সবুজের মুনাফার অংশ = (১/৩)}

\emph{প্রশ্নমতে, (x − ৯০০) এর (১/৩) = ২,০০০}

\emph{⇒ x − ৯০০ = ৬,০০০}

\emph{⇒ x = ৬,০০০ + ৯০০}

\emph{∴ x = ৬,৯০০ টাকা}
\end{quote}

\textbf{রুবেল, রবি ও সাব্বির মোট ১১,০০০ টাকা নিয়ে ব্যবসায় আরম্ভ করল। রুবেল ও
রবি সমান মূলধন দিলেও সাব্বির তাদের প্রত্যেকের চেয়ে ১,০০০ টাকা কম বিনিয়োগ
করে। ব্যবসায় মোট ১,২১০ টাকা মুনাফা হলে সাব্বির কত টাকা পাবে?}

ক) ৩০০

\textbf{\hl{খ) ৩৩০ ✔}}

গ) ৪০০

ঘ) ৪৪০

\textbf{\hl{✔ সঠিক উত্তর: খ) ৩৩০}}

\begin{quote}
\textbf{ব্যাখ্যা:} \emph{ধরি, রুবেল ও রবি প্রত্যেকে বিনিয়োগ করে x টাকা}

\emph{∴ সাব্বির বিনিয়োগ করে (x − ১,০০০) টাকা}

\emph{প্রশ্নমতে, x + x + (x − ১,০০০) = ১১,০০০}

\emph{⇒ ৩x − ১,০০০ = ১১,০০০}

\emph{⇒ ৩x = ১২,০০০}

\emph{∴ x = ৪,০০০}

\emph{তাহলে, রুবেল = ৪,০০০ টাকা, রবি = ৪,০০০ টাকা এবং সাব্বির = (৪,০০০ −
১,০০০) = ৩,০০০ টাকা}

\emph{রুবেল, রবি ও সাব্বিরের বিনিয়োগের অনুপাত = ৪,০০০ : ৪,০০০ : ৩,০০০ = ৪ :
৪ : ৩}

\emph{অনুপাতগুলোর যোগফল = ৪ + ৪ + ৩ = ১১}

\emph{∴ সাব্বির মুনাফা পাবে = ১,২১০ এর (৩/১১) টাকা = ৩৩০ টাকা}
\end{quote}

\textbf{A started a business with Tk. 21,000 and is joined afterwards by
B with Tk. 36,000. After how many months did B join if the profits at
the end of the year are divided equally?/ A ২১,০০০ টাকা দিয়ে একটি ব্যবসা
শুরু করেন এবং কিছুদিন পরে B ৩৬,০০০ টাকা যোগ করেন। B যোগদান করার কত মাস পরে
তাদের লভ্যাংশ সমানভাবে ভাগ করা যাবে?}

ক) 3

খ) 4

\textbf{\hl{গ) 5 ✔}}

ঘ) 6

\textbf{\hl{✔ সঠিক উত্তর: গ) 5}}

\begin{quote}
\textbf{ব্যাখ্যা:} \emph{ধরি, B ব্যবসায়ে x মাসের জন্য বিনিয়োগ করেছিল।}

\emph{A-এর মূলধন ১২ মাসের জন্য নিয়োজিত ছিল।}

\emph{যেহেতু বছর শেষে লভ্যাংশ সমানভাবে ভাগ হয়, তাই উভয় পক্ষের মোট কার্যকর
মূলধন সমান হবে।}

\emph{প্রশ্নমতে, ২১,০০০ × ১২ = ৩৬,০০০ × x}

\emph{⇒ ২৫২,০০০ = ৩৬,০০০x}

\emph{⇒ x = (২৫২,০০০/৩৬,০০০)}

\emph{⇒ x = ৭ মাস}

\emph{অর্থাত B ব্যবসাটি ৭ মাস চালিয়েছে।}

\emph{∴ B ব্যবসা শুরু হওয়ার (১২ − ৭) = ৫ মাস পর যোগ দিয়েছিল।}
\end{quote}

\textbf{\textquotesingle ক\textquotesingle,
\textquotesingle খ\textquotesingle{} ও
\textquotesingle গ\textquotesingle{} একটি অংশীদারি ব্যবসায় যথাক্রমে ৭,৫০০
টাকা, ৮,৫০০ টাকা এবং ১০,৫০০ টাকা বিনিয়োগ করে। বছর শেষে ২,১২০ টাকা লাভ
হয়। যদি মূলধন অনুপাতে লাভ বণ্টন হয় তাহলে খ কত টাকা পাবে?}

ক) ৬০০

\textbf{\hl{খ) ৬৮০ ✔}}

গ) ৭২০

ঘ) ৮৪০

\textbf{\hl{✔ সঠিক উত্তর: খ) ৬৮০}}

\begin{quote}
\textbf{ব্যাখ্যা:} \emph{ক, খ ও গ-এর বিনিয়োগের অনুপাত = ৭,৫০০ : ৮,৫০০ :
১০,৫০০ = ১৫ : ১৭ : ২১}

\emph{অনুপাতের রাশিগুলোর যোগফল = ১৫ + ১৭ + ২১ = ৫৩}

\emph{মোট লাভ = ২,১২০ টাকা}

\emph{∴ \textquotesingle খ\textquotesingle{} লাভ পাবে = ২,১২০ এর (১৭/৫৩)
টাকা = ৬৮০ টাকা}
\end{quote}

\textbf{কোন ব্যবসায় \textquotesingle ক\textquotesingle,
\textquotesingle খ\textquotesingle, \textquotesingle গ\textquotesingle{}
এর মূলধন যথাক্রমে ৩২০, ৪০০ এবং ৪৮০ টাকা। ব্যবসায় ৩০০ টাকা লাভ হলে
\textquotesingle ক\textquotesingle{} অপেক্ষা
\textquotesingle গ\textquotesingle{} কত টাকা বেশি পাবে?}

ক) ৬০ টাকা

খ) ৮০ টাকা

গ) ১২০ টাকা

\textbf{\hl{ঘ) ৪০ টাকা ✔}}

\textbf{\hl{✔ সঠিক উত্তর: ঘ) ৪০ টাকা}}

\begin{quote}
\textbf{ব্যাখ্যা:} \emph{ক, খ ও গ এর মোট মূলধন = ৩২০ + ৪০০ + ৪৮০ = ১২০০
টাকা}

\emph{ক-এর লভ্যাংশ = (৩২০/১২০০) × ৩০০ = ৮০ টাকা}

\emph{গ-এর লভ্যাংশ = (৪৮০/১২০০) × ৩০০ = ১২০ টাকা}

\emph{∴ \textquotesingle গ\textquotesingle{}
\textquotesingle ক\textquotesingle{} অপেক্ষা বেশি পাবে = ১২০ − ৮০ = ৪০
টাকা।}
\end{quote}

\section{5. গড় (Average) মাস্টারক্লাস: বিগত সালের প্রশ্ন, ব্যাখ্যা ও লাইভ
এক্সাম}\label{ux997ux9a1-average-ux9aeux9b8ux99fux9b0ux995ux9b2ux9b8-ux9acux997ux9a4-ux9b8ux9b2ux9b0-ux9aaux9b0ux9b6ux9a8-ux9acux9afux996ux9af-ux993-ux9b2ux987ux9ad-ux98fux995ux9b8ux9ae}

বিসিএস, শিক্ষক নিবন্ধন কিংবা প্রাইমারি শিক্ষক নিয়োগসহ যেকোনো প্রতিযোগিতামূলক
পরীক্ষায় গণিতের `গড়' (Average) অধ্যায় থেকে প্রশ্ন আসবেই! তাই পরীক্ষার হলে কম
সময়ে নির্ভুল উত্তর দিতে প্রয়োজন সঠিক অনুশীলন। আপনাদের প্রস্তুতিকে নিখুঁত করতে এই
আর্টিকেলে বিগত সালের প্রশ্ন ও ব্যাখ্যাসহ সমাধানের পাশাপাশি যুক্ত করা হয়েছে
`লাইভ এক্সাম মোড' (Live Exam Mode)। এখন আপনি শুধু সমাধান শিখবেনই না, বরং
নিজে নিজে উত্তর সিলেক্ট করে তাৎক্ষণিক নিজের মেধা যাচাইও করতে পারবেন!

\subsubsection{গড় কী ?}\label{ux997ux9dc-ux995}

এক জাতীয় একাধিক রাশির সমষ্টিকে উক্ত রাশিগুলোর মোট সংখ্যা দ্বারা ভাগ করলে যে
ভাগফল পাওয়া যায় তাকে ঐ রাশিগুলোর গড় বলে।যেমন-

\textbf{একটি শ্রেণীতে ৫ জন বালকের ওজন যথাক্রমে ৪৮,৪২,৪৬,৫০,৫৯ কেজি হলে
তাদের গড় ওজন কত?}

\textbf{সমাধান:}

প্রথমে রাশিগুলোর যোগফল বের করতে হবে ৪৮+৪২+৪৬+৫০+৫৯ = ২৪৫ কেজি । এখন ই
যোগফলকে ৫ দিয়ে ভাগ করতে হবে কেননা মোট বালকের সংখ্যা ৫ জন । তাহলে তাদের গড়
বয়স হলো \(\frac{\text{২৪৫}}{\text{৫}}\) = ৪৯ কেজি ।

\subsubsection{গড় নির্ণয়ের সুত্র কী
?}\label{ux997ux9dc-ux9a8ux9b0ux9a3ux9dfux9b0-ux9b8ux9a4ux9b0-ux995}

গড় = \(\frac{\text{একজাতীয়}\ \text{কতিপয়}\ \text{রাশির}\ \text{সমষ্টি}}{\text{রাশির}\ \text{সংখ্যা}}\)

\textbf{100 জন শিক্ষার্থীর পরিসংখ্যানে গড় নম্বর 70। এদের মধ্যে 60 জন ছাত্রীর
গড় নম্বর 75 হলে, ছাত্রদের গড় নম্বর কত?}

ক) 55.5

খ) 60.5

গ) 65.5

\textbf{\hl{ঘ) 62.5 ✔}}

\textbf{\hl{✔ সঠিক উত্তর: ঘ) 62.5}}

\begin{quote}
\textbf{ব্যাখ্যা:} \emph{100 জন শিক্ষার্থীর পরিসংখ্যানে মোট নম্বর = 70 × 100 =
7000}

\emph{60 জন ছাত্রীর মোট নম্বর = 75 × 60 = 4500}

\emph{বাকি (100 − 60) = 40 জন ছাত্রদের মোট নম্বর = 7000 − 4500 = 2500}

\emph{∴ ছাত্রদের গড় নম্বর = 2500 ÷ 40 = 62.5}
\end{quote}

\textbf{৭টি সংখ্যার গড় ৪০। এর সাথে ৩টি সংখ্যা যোগ করা হলো। সংখ্যা ৩টির গড়
২১। সমষ্টিগত ভাবে ১০টি সংখ্যার গড় কত ?}

ক) ৬১

খ) ৬.১

গ) ৭০

\textbf{\hl{ঘ) ৩৪.৩ ✔}}

\textbf{\hl{✔ সঠিক উত্তর: ঘ) ৩৪.৩}}

\begin{quote}
\textbf{ব্যাখ্যা:} \emph{৭ টি সংখ্যার সমষ্টি = ৭ × ৪০ = ২৮০}

\emph{৩ টি সংখ্যার সমষ্টি = ৩ × ২১ = ৬৩}

\emph{১০ টি সংখ্যার সমষ্টি = ২৮০ + ৬৩ = ৩৪৩}

\emph{∴ ১০ টি সংখ্যার গড় = ৩৪৩ ÷ ১০ = ৩৪.৩}
\end{quote}

\textbf{১১টি সংখ্যার গড় ৩০। প্রথম পাঁচটি সংখ্যার গড় ২৫ ও শেষ পাঁচটি সংখ্যার গড়
২৮। ষষ্ঠ সংখ্যাটি কত ?}

ক) ৫৫

খ) ৫৮

\textbf{\hl{গ) ৬৫ ✔}}

ঘ) ৬৭

\textbf{\hl{✔ সঠিক উত্তর: গ) ৬৫}}

\begin{quote}
\textbf{ব্যাখ্যা:} \emph{১১টি সংখ্যার সমষ্টি = ৩০ × ১১ = ৩৩০}

\emph{প্রথম ৫টি সংখ্যার সমষ্টি = ২৫ × ৫ = ১২৫}

\emph{শেষ ৫টি সংখ্যার সমষ্টি = ২৮ × ৫ = ১৪০}

\emph{প্রথম ও শেষ ১০টি সংখ্যার সমষ্টি = ১২৫ + ১৪০ = ২৬৫}

\emph{∴ ষষ্ঠ সংখ্যাটি = ৩৩০ − ২৬৫ = ৬৫}
\end{quote}

\textbf{২৪ জন ছাত্র এবং একজন শিক্ষকের বয়সের গড় ১৫ বছর। শিক্ষককে বাদ দিয়ে
ছাত্রদের বয়সের গড় করলে গড় ১ বছর কমে যায়। শিক্ষকের বয়স কত ?}

ক) ৩৮

খ) ৪০

গ) ১২

\textbf{\hl{ঘ) কোনোটিই নয় ✔}}

\textbf{\hl{✔ সঠিক উত্তর: ঘ) কোনোটিই নয়}}

\begin{quote}
\textbf{ব্যাখ্যা:} \emph{২৪ জন ছাত্র এবং একজন শিক্ষকসহ মোট ২৫ জনের বয়সের
সমষ্টি = ২৫ × ১৫ = ৩৭৫ বছর}

\emph{শিক্ষককে বাদ দিলে ২৪ জন ছাত্রের বয়সের গড় হয় (১৫ − ১) = ১৪ বছর}

\emph{২৪ জন ছাত্রের বয়সের সমষ্টি = ২৪ × ১৪ = ৩৩৬ বছর}

\emph{∴ শিক্ষকের বয়স = ৩৭৫ − ৩৩৬ = ৩৯ বছর}
\end{quote}

\textbf{কোন শ্রেণীর ২৪ জন ছাত্রের বয়সের গড় ১৪ বছর। যদি একজন শ্রেণী শিক্ষকের
বয়স তাদের বয়সের সাথে যোগ করা হয় তবে বয়সের গড় এক বছর বৃদ্ধি পায়। শিক্ষকের
বয়স কত ?}

ক) ৪০ বৎসর

\textbf{\hl{খ) ৩৯ বৎসর ✔}}

গ) ৪৫ বৎসর

ঘ) ৩৫ বৎসর

\textbf{\hl{✔ সঠিক উত্তর: খ) ৩৯ বৎসর}}

\begin{quote}
\textbf{ব্যাখ্যা:} \emph{২৪ জন ছাত্রের বয়সের সমষ্টি = ২৪ × ১৪ = ৩৩৬ বছর}

\emph{শিক্ষকসহ (২৪ + ১) = ২৫ জনের বয়সের গড় হয় (১৪ + ১) = ১৫ বছর}

\emph{শিক্ষকসহ ২৫ জনের বয়সের সমষ্টি = ২৫ × ১৫ = ৩৭৫ বছর}

\emph{∴ শিক্ষকের বয়স = ৩৭৫ − ৩৩৬ = ৩৯ বছর}
\end{quote}

\textbf{ক ও খ এর মানের গড় ৯ এবং গ এর মান ১২ হলে ক, খ ও গ এর মানের গড় কত
হবে ?}

ক) ৬

খ) ৯

\textbf{\hl{গ) ১০ ✔}}

ঘ) ১২

\textbf{\hl{✔ সঠিক উত্তর: গ) ১০}}

\begin{quote}
\textbf{ব্যাখ্যা:} \emph{ক ও খ এর সমষ্টি = ৯ × ২ = ১৮}

\emph{ক, খ ও গ এর মানের সমষ্টি = ১৮ + ১২ = ৩০}

\emph{∴ ক, খ ও গ এর গড় = (৩০/৩) = ১০}
\end{quote}

\textbf{৩টি সংখ্যার গড় ৬ এবং ঐ ৩টি সংখ্যাসহ মোট ৪টি সংখ্যার গড় ৮ হলে চতুর্থ
সংখ্যাটির অর্ধেকের মান কত ?}

ক) ৮

খ) ৫

গ) ৬

\textbf{\hl{ঘ) ৭ ✔}}

\textbf{\hl{✔ সঠিক উত্তর: ঘ) ৭}}

\begin{quote}
\textbf{ব্যাখ্যা:} \emph{প্রথম ৩টি সংখ্যার যোগফল = ৩ × ৬ = ১৮}

\emph{৩টি সংখ্যাসহ ৪টি সংখ্যার যোগফল = ৪ × ৮ = ৩২}

\emph{∴ চতুর্থ সংখ্যাটি = (৩২ − ১৮) = ১৪}

\emph{∴ ৪র্থ সংখ্যার অর্ধেক = (১৪/২) = ৭}
\end{quote}

\textbf{x ও y এর মানের গড় ৯ এবং z = ১২ হলে x,y,z এর মানের গড় কত?}

ক) ৬

খ) ৯

\textbf{\hl{গ) ১০ ✔}}

ঘ) ১২

\textbf{\hl{✔ সঠিক উত্তর: গ) ১০}}

\begin{quote}
\textbf{ব্যাখ্যা:} \emph{x ও y এর যোগফল = ৯ × ২ = ১৮}

\emph{x,y,z এর মানের সমষ্টি = ১৮ + ১২ = ৩০}

\emph{∴ x,y,z এর গড় = (৩০/৩) = ১০}
\end{quote}

\textbf{৭, ১২, ১৭, ৩, ১১, ৬, ১৩, ৩ সংখ্যাগুলোর গড় কত?}

ক) ৬

খ) ৭

গ) ৮

\textbf{\hl{ঘ) ৯ ✔}}

\textbf{\hl{✔ সঠিক উত্তর: ঘ) ৯}}

\begin{quote}
\textbf{ব্যাখ্যা:} \emph{সংখ্যাগুলোর সমষ্টি = ৭ + ১২ + ১৭ + ৩ + ১১ + ৬ + ১৩ +
৩ = ৭২}

\emph{মোট সংখ্যা = ৮টি}

\emph{∴ গড় = (৭২/৮) = ৯}
\end{quote}

\textbf{0, 5, 7 এর গড় কত?}

ক) 6

খ) 0

\textbf{\hl{গ) 4 ✔}}

ঘ) 1

\textbf{\hl{✔ সঠিক উত্তর: গ) 4}}

\begin{quote}
\textbf{ব্যাখ্যা:} \emph{সংখ্যাগুলোর সমষ্টি = 0 + 5 + 7 = 12}

\emph{মোট সংখ্যা = 3টি}

\emph{∴ গড় = (12/3) = 4}
\end{quote}

\textbf{− 30 এবং − 40 সংখ্যা দুইটির গড় ব্যবধান কত?}

ক) - 35

খ) -55

\textbf{\hl{গ) 5 ✔}}

ঘ) - 5

\textbf{\hl{✔ সঠিক উত্তর: গ) 5}}

\begin{quote}
\textbf{ব্যাখ্যা:} \emph{গাণিতিক গড় = (−30+(−40)/2) = (−70/2) = −35}

\emph{∴ গড় ব্যবধান = (\textbar−30+35\textbar+\textbar−40+35\textbar/2) =
(\textbar5\textbar+\textbar−5\textbar/2) = (10/2) = 5}
\end{quote}

\textbf{(১/২) , (১/৪) , (৩/৪) এর গড় মান কত?}

ক) (১/৪)

খ) (৩/৪)

\textbf{\hl{গ) (১/২) ✔}}

ঘ) (২/৩)

\textbf{\hl{✔ সঠিক উত্তর: গ) (১/২)}}

\begin{quote}
\textbf{ব্যাখ্যা:} \emph{সংখ্যাগুলোর সমষ্টি = (১/২) + (১/৪) + (৩/৪) =
(২+১+৩/৪) = (৬/৪)}

\emph{∴ গড় মান = (৬/৪)÷৩ = (৬/৪)×(১/৩)=(১/২)}
\end{quote}

\textbf{M সংখ্যক সংখ্যার গড় A এবং N সংখ্যক সংখ্যার গড় B হলে সবগুলো সংখ্যার গড়
কত?}

ক) (A+B/2)

খ) (AM+BN/2)

\textbf{\hl{গ) (AM+BN/M+N) ✔}}

ঘ) (AM+BN/A+B)

\textbf{\hl{✔ সঠিক উত্তর: গ) (AM+BN/M+N)}}

\begin{quote}
\textbf{ব্যাখ্যা:} \emph{M সংখ্যক সংখ্যার সমষ্টি = MA}

\emph{N সংখ্যক সংখ্যার সমষ্টি = NB}

\emph{∴ মোট (M + N) সংখ্যক সংখ্যার সমষ্টি = MA + NB}

\emph{∴ সবগুলো সংখ্যার গড় = (AM+BN/M+N)}
\end{quote}

\textbf{তিনটি পূর্ণ সংখ্যার গড় ১৫০ এবং ক্ষুদ্রতম সংখ্যা দুইটির গড় ১২০। বৃহত্তম
সংখ্যাটি কত?}

ক) ২৩০

\textbf{\hl{খ) ২১০ ✔}}

গ) ২০০

ঘ) ১৯০

\textbf{\hl{✔ সঠিক উত্তর: খ) ২১০}}

\begin{quote}
\textbf{ব্যাখ্যা:} \emph{৩টি পূর্ণ সংখ্যার সমষ্টি = ১৫০ × ৩ = ৪৫০}

\emph{ক্ষুদ্রতম সংখ্যা দুটির সমষ্টি = ১২০ × ২ = ২৪০}

\emph{∴ বৃহত্তম সংখ্যাটি = ৪৫০ − ২৪০ = ২১০}
\end{quote}

\textbf{প্রথম ও দ্বিতীয় সংখ্যার গড় ২৫। প্রথম, দ্বিতীয় ও তৃতীয় সংখ্যার গড় ৩০
হলে, তৃতীয় সংখ্যাটি কত?}

ক) ২৫

\textbf{\hl{খ) ৪০ ✔}}

গ) ৯০

ঘ) ৫০

\textbf{\hl{✔ সঠিক উত্তর: খ) ৪০}}

\begin{quote}
\textbf{ব্যাখ্যা:} \emph{১ম, ২য় ও ৩য় সংখ্যার সমষ্টি = ৩ × ৩০ = ৯০}

\emph{১ম ও ২য় সংখ্যার সমষ্টি = ২ × ২৫ = ৫০}

\emph{∴ ৩য় সংখ্যাটি = ৯০ − ৫০ = ৪০}
\end{quote}

\textbf{৩টি সংখ্যার গড় ৩৩। দুইটি সংখ্যা ২৪ এবং ৪২ হলে অপর সংখ্যাটি কত?}

ক) ১২

খ) ২২

গ) ৩২

\textbf{\hl{ঘ) ৩৩ ✔}}

\textbf{\hl{✔ সঠিক উত্তর: ঘ) ৩৩}}

\begin{quote}
\textbf{ব্যাখ্যা:} \emph{৩টি সংখ্যার সমষ্টি = ৩৩ × ৩ = ৯৯}

\emph{২টি সংখ্যার সমষ্টি = ২৪ + ৪২ = ৬৬}

\emph{∴ অপর সংখ্যাটি = ৯৯ − ৬৬ = ৩৩}
\end{quote}

\textbf{৭টি সংখ্যার গড় ১২। একটি সংখ্যা বাতিল করলে গড় হয় ১১। বাতিলকৃত
সংখ্যাটি কত?}

ক) ১০

খ) ১২

গ) ১৫

\textbf{\hl{ঘ) ১৮ ✔}}

\textbf{\hl{✔ সঠিক উত্তর: ঘ) ১৮}}

\begin{quote}
\textbf{ব্যাখ্যা:} \emph{৭টি সংখ্যার সমষ্টি = ৭ × ১২ = ৮৪}

\emph{একটি সংখ্যা বাতিল করলে অবশিষ্ট (৭ − ১) = ৬টি সংখ্যার সমষ্টি = ৬ × ১১ =
৬৬}

\emph{∴ বাতিলকৃত সংখ্যাটি = ৮৪ − ৬৬ = ১৮}
\end{quote}

\textbf{ছয়টি সংখ্যার গড় ৬। যদি প্রত্যেকটি সংখ্যা থেকে ৩ বিয়োগ করা হয় তবে নতুন
সংখ্যাগুলোর গড় কত হবে?}

\textbf{\hl{ক) ৩ ✔}}

খ) ৩.৫

গ) ৪

ঘ) কোনোটিই নয়

\textbf{\hl{✔ সঠিক উত্তর: ক) ৩}}

\begin{quote}
\textbf{ব্যাখ্যা:} \emph{ছয়টি সংখ্যার সমষ্টি = (৬ × ৬) = ৩৬}

\emph{প্রত্যেক সংখ্যা থেকে ৩ বিয়োগ করলে নতুন ছয়টি সংখ্যার সমষ্টি = ৩৬ − (৬ ×
৩) = ৩৬ − ১৮ = ১৮}

\emph{∴ নতুন ছয়টি সংখ্যার গড় = (১৮/৬) = ৩}
\end{quote}

\textbf{১০টি সংখ্যার যোগফল ৪৬২। এদের প্রথম ৪টার গড় ৫২ এবং শেষের ৫টার গড়
৩৮। পঞ্চম সংখ্যাটি কত?}

ক) ৬০

\textbf{\hl{খ) ৬৪ ✔}}

গ) ৬২

ঘ) ৫০

\textbf{\hl{✔ সঠিক উত্তর: খ) ৬৪}}

\begin{quote}
\textbf{ব্যাখ্যা:} \emph{প্রথম ৪টির সমষ্টি = ৫২ × ৪ = ২০৮}

\emph{শেষ ৫টির সমষ্টি = ৩৮ × ৫ = ১৯০}

\emph{প্রথম ৪টি ও শেষের ৫টি (মোট ৯টি) সংখ্যার সমষ্টি = ২০৮ + ১৯০ = ৩৯৮}

\emph{∴ পঞ্চম সংখ্যাটি = ৪৬২ − ৩৯৮ = ৬৪}
\end{quote}

\textbf{২ থেকে শুরু করে পর পর পাঁচটি জোড়া সংখ্যার গড় কত হবে?}

ক) ৪

খ) ৫

\textbf{\hl{গ) ৬ ✔}}

ঘ) ৭

\textbf{\hl{✔ সঠিক উত্তর: গ) ৬}}

\begin{quote}
\textbf{ব্যাখ্যা:} \emph{২ থেকে শুরু করে পরপর ৫টি জোড় সংখ্যা হলো: ২, ৪, ৬, ৮,
১০}

\emph{সংখ্যাগুলোর সমষ্টি = ২ + ৪ + ৬ + ৮ + ১০ = ৩০}

\emph{∴ গড় = (৩০/৫) = ৬}
\end{quote}

\textbf{ছয়টি ধারাবাহিক বিজোড় সংখ্যার গড়ের ছয়গুণ সবচেয়ে বড় সংখ্যাটির চারগুণ
অপেক্ষা ১২ বেশি। ধারাবাহিক সংখ্যাগুলোর গড় কত?}

\textbf{\hl{ক) ১৬ ✔}}

খ) ১৭

গ) ১৮

ঘ) ২০

\textbf{\hl{✔ সঠিক উত্তর: ক) ১৬}}

\begin{quote}
\textbf{ব্যাখ্যা:} \emph{ধরি, ধারাবাহিক বিজোড় সংখ্যা ৬টি যথাক্রমে x, x + 2,
x + 4, x + 6, x + 8, x + 10}

\emph{সংখ্যাগুলোর সমষ্টি = 6x + 30}

\emph{গড় = (6x+30/6) = x + 5}

\emph{প্রশ্নমতে, 6(x + 5) = 4(x + 10) + 12}

\emph{⇒ 6x + 30 = 4x + 40 + 12}

\emph{⇒ 2x = 22}

\emph{∴ x = 11}

\emph{∴ সংখ্যাগুলোর গড় = x + 5 = 11 + 5 = 16}
\end{quote}

\textbf{৫টি সংখ্যার গড় ৪০। এর সাথে ৩টি সংখ্যা যোগ করা হলো, সংখ্যা তিনটির গড়
২২। সমষ্টিগতভাবে ৮টি সংখ্যার গড় কত?}

ক) ৩৩.২

খ) ৩৩.৫০

\textbf{\hl{গ) ৩৩.২৫ ✔}}

ঘ) কোনোটিই নয়

\textbf{\hl{✔ সঠিক উত্তর: গ) ৩৩.২৫}}

\begin{quote}
\textbf{ব্যাখ্যা:} \emph{৫টি সংখ্যার সমষ্টি = ৪০ × ৫ = ২০০}

\emph{অন্য ৩টি সংখ্যার সমষ্টি = ২২ × ৩ = ৬৬}

\emph{∴ মোট (৫ + ৩) = ৮টি সংখ্যার সমষ্টি = ২০০ + ৬৬ = ২৬৬}

\emph{∴ ৮টি সংখ্যার গড় = (২৬৬/৮) = ৩৩.২৫}
\end{quote}

\textbf{০.৬ হলো ০.২, ০.৮, ১ এবং x এর গড় মান। x এর মান কত?}

ক) ০.২

\textbf{\hl{খ) ০.৪ ✔}}

গ) ০.৬৭

ঘ) ২.৪

\textbf{\hl{✔ সঠিক উত্তর: খ) ০.৪}}

\begin{quote}
\textbf{ব্যাখ্যা:} \emph{৪টি সংখ্যার সমষ্টি = ০.২ + ০.৮ + ১ + x = ২ + x}

\emph{গড় = (২+x/৪)}

\emph{শর্তমতে, (২+x/৪) = ০.৬}

\emph{⇒ ২ + x = ২.৪}

\emph{∴ x = ২.৪ − ২ = ০.৪}
\end{quote}

\textbf{৬, ৮, ১০ এর গাণিতিক গড় ৭, ৯ এবং কোন সংখ্যার গাণিতিক গড়ের সমান?}

ক) ৫

\textbf{\hl{খ) ৮ ✔}}

গ) ৬

ঘ) ১০

\textbf{\hl{✔ সঠিক উত্তর: খ) ৮}}

\begin{quote}
\textbf{ব্যাখ্যা:} \emph{৬, ৮, ১০ এর গাণিতিক গড় = (৬+৮+১০/৩) = (২৪/৩) = ৮}

\emph{ধরি, নির্ণেয় সংখ্যাটি = x}

\emph{৭, ৯, x এর গড় = ৮}

\emph{⇒ ৭ + ৯ + x = ৩ × ৮}

\emph{⇒ ১৬ + x = ২৪}

\emph{∴ x = ২৪ − ১৬ = ৮}
\end{quote}

\textbf{পাঁচটি সংখ্যার গড় ৪৬ এবং সংখ্যাগুলোর প্রথম চারটি সংখ্যার গড় ৪৫। পঞ্চম
সংখ্যাটি কত?}

ক) ৪৮

খ) ৪৯

\textbf{\hl{গ) ৫০ ✔}}

ঘ) কোনোটিই নয়

\textbf{\hl{✔ সঠিক উত্তর: গ) ৫০}}

\begin{quote}
\textbf{ব্যাখ্যা:} \emph{পাঁচটি সংখ্যার সমষ্টি = ৫ × ৪৬ = ২৩০}

\emph{প্রথম চারটি সংখ্যার সমষ্টি = ৪ × ৪৫ = ১৮০}

\emph{∴ পঞ্চম সংখ্যাটি = ২৩০ − ১৮০ = ৫০}
\end{quote}

\textbf{একসেট সংখ্যা থেকে ৩৫ সংখ্যাটি বাদ দেয়ার ফলে সেটের গড় ১৪ থেকে ১১
হয়ে গেল। সেটের সদস্য সংখ্যা কত ছিল?}

\textbf{\hl{ক) ৮ ✔}}

খ) ১১

গ) ২১

ঘ) ২৫

\textbf{\hl{✔ সঠিক উত্তর: ক) ৮}}

\begin{quote}
\textbf{ব্যাখ্যা:} \emph{ধরি, সেটের প্রাথমিক সদস্য সংখ্যা ছিল x}

\emph{∴ x সংখ্যক সংখ্যার সমষ্টি = ১৪x}

\emph{৩৫ বাদ দেওয়ার পর সদস্য সংখ্যা হয় (x − ১) এবং সমষ্টি হয় ১১(x − ১)}

\emph{প্রশ্নমতে, ১৪x − ৩৫ = ১১(x − ১)}

\emph{⇒ ১৪x − ৩৫ = ১১x − ১১}

\emph{⇒ ১৪x − ১১x = ৩৫ − ১১}

\emph{⇒ ৩x = ২৪}

\emph{∴ x = ৮}
\end{quote}

\textbf{The average of 5 numbers is 40. If 2 more numbers with an
average of 21 are added to these numbers, what will be the average of
the combined numbers?}

ক) 8/7

খ) 30.1

গ) 30.3

\textbf{\hl{ঘ) 34.6 ✔}}

\textbf{\hl{✔ সঠিক উত্তর: ঘ) 34.6}}

\begin{quote}
\textbf{ব্যাখ্যা:} \emph{Sum of first 5 numbers = 5 × 40 = 200}

\emph{Sum of additional 2 numbers = 2 × 21 = 42}

\emph{Total sum of (5 + 2) = 7 numbers = 200 + 42 = 242}

\emph{∴ Average of combined numbers = (242/7) = 34.6}
\end{quote}

\textbf{একজন ব্যাটসম্যান প্রথম তিনটি T-20 খেলায় ৮২, ৮৫ ও ৯২ রান করেন। চতুর্থ
খেলায় কত রান করলে; তার গড় রান ৮৭ হবে?}

ক) ৮৬

খ) ৮৭

গ) ৮৮

\textbf{\hl{ঘ) ৮৯ ✔}}

\textbf{\hl{✔ সঠিক উত্তর: ঘ) ৮৯}}

\begin{quote}
\textbf{ব্যাখ্যা:} \emph{প্রথম তিনটি খেলায় মোট রান = ৮২ + ৮৫ + ৯২ = ২৫৯}

\emph{৪টি খেলায় মোট প্রয়োজনীয় রান = ৮৭ × ৪ = ৩৪৮}

\emph{∴ চতুর্থ খেলায় রান করতে হবে = ৩৪৮ − ২৫৯ = ৮৯}
\end{quote}

\textbf{পরীক্ষায় \textquotesingle ক\textquotesingle{} এর প্রাপ্ত নম্বর
যথাক্রমে ৮২, ৮৫ ও ৯২। চতুর্থ পরীক্ষায় তাকে কত নম্বর পেতে হবে যেন তার গড় প্রাপ্ত
নম্বর ৮৭ হয়?}

\textbf{\hl{ক) ৮৯ ✔}}

খ) ৮৮

গ) ৮৬

ঘ) ৯১

\textbf{\hl{✔ সঠিক উত্তর: ক) ৮৯}}

\begin{quote}
\textbf{ব্যাখ্যা:} \emph{পরীক্ষায় \textquotesingle ক\textquotesingle{} এর
প্রাপ্ত মোট নম্বর = (৮২ + ৮৫ + ৯২) = ২৫৯}

\emph{৪ টি পরীক্ষায় গড়ে ৮৭ হলে মোট নম্বর = ৮৭ × ৪ = ৩৪৮}

\emph{∴ চতুর্থ পরীক্ষায় তাকে পেতে হবে (৩৪৮ − ২৫৯) = ৮৯}
\end{quote}

\textbf{একজন ছাত্র বাংলায় ৬৫ এবং ইংলিশে ৮০ নম্বর পেল। সে অঙ্ক পরীক্ষায় কত
পেলে এই তিন বিষয়ে তার গড় নম্বর ৭৫ হবে?}

ক) ৭০

\textbf{\hl{খ) ৮০ ✔}}

গ) ৮৫

ঘ) কোনোটিই নয়

\textbf{\hl{✔ সঠিক উত্তর: খ) ৮০}}

\begin{quote}
\textbf{ব্যাখ্যা:} \emph{ধরি, অঙ্ক পরীক্ষায় x নম্বর পেতে হবে।}

\emph{∴ (৬৫+৮০+x/৩) = ৭৫}

\emph{⇒ ১৪৫ + x = ২২৫}

\emph{⇒ x = ২২৫ − ১৪৫}

\emph{∴ x = ৮০}
\end{quote}

\textbf{যদি ১৫ জন ছাত্র ইংরেজীতে গড়ে শতকরা ৮০ নম্বর এবং ১০ জন ছাত্র গড়ে
শতকরা ৯০ নম্বর পায়, তাহলে ২৫ জন ছাত্রের শতকরা হিসাবে গড় নম্বর কত?}

\textbf{\hl{ক) ৮৪ ✔}}

খ) ৮৪\%

গ) ৮৩\%

ঘ) ৮৩(১/৩) \%

\textbf{\hl{✔ সঠিক উত্তর: ক) ৮৪}}

\begin{quote}
\textbf{ব্যাখ্যা:} \emph{১৫ জনের মোট নম্বর = ১৫ × ৮০ = ১২০০}

\emph{১০ জনের মোট নম্বর = ১০ × ৯০ = ৯০০}

\emph{∴ ২৫ জনের মোট নম্বর = (১২০০ + ৯০০) = ২১০০}

\emph{∴ ২৫ জনের গড় নম্বর = ২১০০ ÷ ২৫ = ৮৪}
\end{quote}

\textbf{এক দোকানদার ১২ দিনে ৫০৪ টাকা আয় করলেন, প্রথম ৪ দিনের গড় আয় ৪০
টাকা হলে বাকী দিনগুলোর গড় আয় কত টাকা হবে?}

\textbf{\hl{ক) ৪৩ টাকা ✔}}

খ) ৪৭ টাকা

গ) ৪২ টাকা

ঘ) ৪০ টাকা

\textbf{\hl{✔ সঠিক উত্তর: ক) ৪৩ টাকা}}

\begin{quote}
\textbf{ব্যাখ্যা:} \emph{প্রথম ৪ দিনের গড় আয় ৪০ টাকা। মোট আয় = ৪ × ৪০ = ১৬০
টাকা।}

\emph{∴ বাকি (১২ − ৪) ৮ দিনে মোট আয় = ৫০৪ − ১৬০ = ৩৪৪ টাকা।}

\emph{∴ বাকি ১ দিনে গড় আয় = ৩৪৪ ÷ ৮ = ৪৩ টাকা।}
\end{quote}

\textbf{খ এর দৈনিক আয় ক এর দ্বিগুণ এবং গ এর দৈনিক আয় খ এর দ্বিগুণ। তাদের
তিনজনের আয়ের যোগফল ১১২০ টাকা হলে গ এর দৈনিক আয় কত টাকা?}

ক) ৫৬০

\textbf{\hl{খ) ৬৪০ ✔}}

গ) ৬৮৪

ঘ) কোনোটিই নয়

\textbf{\hl{✔ সঠিক উত্তর: খ) ৬৪০}}

\begin{quote}
\textbf{ব্যাখ্যা:} \emph{ধরি, ক এর আয় = x টাকা}

\emph{∴ খ এর আয় = ২x টাকা}

\emph{এবং গ এর আয় = ২ × ২x = ৪x টাকা}

\emph{প্রশ্নমতে, x + ২x + ৪x = ১১২০}

\emph{⇒ ৭x = ১১২০}

\emph{⇒ x = ১৬০}

\emph{∴ গ এর দৈনিক আয় = ৪ × ১৬০ = ৬৪০ টাকা।}
\end{quote}

\textbf{একজন ক্রিকেটারের ১০ ইনিংসের রানের গড় ৪৫.৫। ১১তম ইনিংসে কত রান করে
আউট হলে সব ইনিংস মিলিয়ে তার রানের গড় ৫০ হবে।}

ক) ৫৫ রান

খ) ৪৫ রান

গ) ১০০ রান

\textbf{\hl{ঘ) ৯৫ রান ✔}}

ঙ) কোনটিই নয়

\textbf{\hl{✔ সঠিক উত্তর: ঘ) ৯৫ রান}}

\begin{quote}
\textbf{ব্যাখ্যা:} \emph{১০ ইনিংসের রানের মোট সমষ্টি = ৪৫.৫ × ১০ = ৪৫৫}

\emph{১১ ইনিংসের রানের মোট সমষ্টি = ৫০ × ১১ = ৫৫০}

\emph{∴ ১১তম ইনিংসে প্রয়োজন = ৫৫০ − ৪৫৫ = ৯৫ রান}
\end{quote}

\textbf{এগারজন বালকের মধ্যে প্রথম ছয়জন বালকের গড় উচ্চতা ৮৭ সেমি। শেষের
পাঁচজন বালকের গড় উচ্চতা ১৩১ সেমি। ঐ এগারজন বালকের গড় উচ্চতা কত?}

\textbf{\hl{ক) ১০৭ সেমি ✔}}

খ) ১০৬ সেমি

গ) ১০৫ সেমি

ঘ) ১০৪ সেমি

\textbf{\hl{✔ সঠিক উত্তর: ক) ১০৭ সেমি}}

\begin{quote}
\textbf{ব্যাখ্যা:} \emph{প্রথম ৬ জন বালকের মোট উচ্চতা = ৮৭ × ৬ = ৫২২ সেমি}

\emph{শেষের ৫ জন বালকের মোট উচ্চতা = ১৩১ × ৫ = ৬৫৫ সেমি}

\emph{১১ জন বালকের মোট উচ্চতা = ৫২২ + ৬৫৫ = ১১৭৭ সেমি}

\emph{∴ ১১ জন বালকের গড় উচ্চতা = ১১৭৭ ÷ ১১ = ১০৭ সেমি}
\end{quote}

\textbf{কোন শ্রেণীতে ১০ জন ছাত্রের গড় উচ্চতা ৫ ফুট ৬ ইঞ্চি। এর মধ্যে ৯ জন
ছাত্রের গড় উচ্চতা ৫ ফুট ৫ ইঞ্চি হলে ১০ম ছাত্রের উচ্চতা কত?}

ক) ৫ ফুট ৭ ইঞ্চি

খ) ৬ ফুট ৫ ইঞ্চি

\textbf{\hl{গ) ৬ ফুট ৩ ইঞ্চি ✔}}

ঘ) ৬ ফুট

ঙ) কোনটিই নয়

\textbf{\hl{✔ সঠিক উত্তর: গ) ৬ ফুট ৩ ইঞ্চি}}

\begin{quote}
\textbf{ব্যাখ্যা:} \emph{১০ জন ছাত্রের মোট উচ্চতা = (৫ ফুট ৬ ইঞ্চি) × ১০ = ৫০
ফুট ৬০ ইঞ্চি}

\emph{৯ জন ছাত্রের মোট উচ্চতা = (৫ ফুট ৫ ইঞ্চি) × ৯ = ৪৫ ফুট ৪৫ ইঞ্চি}

\emph{∴ ১০ম ছাত্রের উচ্চতা = (৫০ ফুট ৬০ ইঞ্চি) − (৪৫ ফুট ৪৫ ইঞ্চি)}

\emph{= ৫ ফুট ১৫ ইঞ্চি}

\emph{= ৬ ফুট ৩ ইঞ্চি (যেহেতু ১২ ইঞ্চি = ১ ফুট)}
\end{quote}

\textbf{তিনজন লোকের ওজনের গড় ৫৩ কেজি। এদের কারো ওজন ৫১ কেজির কম নয়। এদের
একজনের সর্বোচ্চ ওজন হবে-}

ক) ৫৩ কেজি

খ) ৫৫ কেজি

\textbf{\hl{গ) ৫৭ কেজি ✔}}

ঘ) ৫৯ কেজি

\textbf{\hl{✔ সঠিক উত্তর: গ) ৫৭ কেজি}}

\begin{quote}
\textbf{ব্যাখ্যা:} \emph{তিনজন লোকের মোট ওজন = ৫৩ × ৩ = ১৫৯ কেজি}

\emph{একজনের সর্বোচ্চ ওজনের জন্য বাকি দুজনের ওজন সর্বনিম্ন হতে হবে।}

\emph{কারো ওজন ৫১ কেজির কম নয়, তাই দুজনের সর্বনিম্ন মোট ওজন = ৫১ × ২ = ১০২
কেজি}

\emph{∴ একজনের সর্বোচ্চ ওজন = ১৫৯ − ১০২ = ৫৭ কেজি}
\end{quote}

\textbf{১১ জন লোকের গড় ওজন ৭০ কেজি। ৯০ কেজি ওজনের একজন লোক চলে গেলে
বাকিদের গড় ওজন কত হয়?}

ক) ৬২ কেজি

\textbf{\hl{খ) ৬৮ কেজি ✔}}

গ) ৮০ কেজি

ঘ) ৭২ কেজি

\textbf{\hl{✔ সঠিক উত্তর: খ) ৬৮ কেজি}}

\begin{quote}
\textbf{ব্যাখ্যা:} \emph{১১ জন লোকের ওজনের সমষ্টি = ৭০ × ১১ = ৭৭০ কেজি}

\emph{৯০ কেজি ওজনের ১ জন লোক চলে গেলে বাকি (১১ − ১) = ১০ জনের ওজনের
সমষ্টি = ৭৭০ − ৯০ = ৬৮০ কেজি}

\emph{∴ বাকি ১০ জনের গড় ওজন = (৬৮০/১০) = ৬৮ কেজি}
\end{quote}

\textbf{ক, খ ও গ এর মাসিক গড় বেতন ৫০০ টাকা। খ, গ ও ঘ এর মাসিক গড় বেতন
৪৫০ টাকা। ক এর বেতন ৫৪০ টাকা হলে ঘ এর বেতন কত?}

ক) ৩৭৫ টাকা

খ) ৩৮০টাকা

গ) ৩৮৫ টাকা

\textbf{\hl{ঘ) ৩৯০ টাকা ✔}}

\textbf{\hl{✔ সঠিক উত্তর: ঘ) ৩৯০ টাকা}}

\begin{quote}
\textbf{ব্যাখ্যা:} \emph{ক + খ + গ-এর বেতনের সমষ্টি = ৫০০ × ৩ = ১৫০০ টাকা}

\emph{খ + গ + ঘ-এর বেতনের সমষ্টি = ৪৫০ × ৩ = ১৩৫০ টাকা}

\emph{ক ও ঘ-এর বেতনের পার্থক্য = ১৫০০ − ১৩৫০ = ১৫০ টাকা}

\emph{∴ ক − ঘ = ১৫০}

\emph{⇒ ৫৪০ − ঘ = ১৫০}

\emph{∴ ঘ এর বেতন = ৫৪০ − ১৫০ = ৩৯০ টাকা}
\end{quote}

\textbf{একজন শ্রমিক প্রতিদিন প্রথম ৮ ঘন্টা কাজ করার জন্য ঘন্টায় ১০ টাকা করে
এবং পরবর্তী সময়ের জন্য ঘন্টায় ১৫ টাকা করে মজুরি পায়। দৈনিক ১০ ঘন্টা করে কাজ
করলে তার ঘন্টা প্রতি গড় মজুরি কত?}

\textbf{\hl{ক) ১১ টাকা ✔}}

খ) ১২ টাকা

গ) ১২.৫০ টাকা

ঘ) ১৩ টাকা

\textbf{\hl{✔ সঠিক উত্তর: ক) ১১ টাকা}}

\begin{quote}
\textbf{ব্যাখ্যা:} \emph{প্রথম ৮ ঘন্টার জন্য মজুরি পায় = ১০ × ৮ = ৮০ টাকা}

\emph{পরবর্তী (১০ − ৮) = ২ ঘন্টার জন্য মজুরি পায় = ১৫ × ২ = ৩০ টাকা}

\emph{সর্বমোট ১০ ঘন্টার জন্য মোট মজুরি পায় = ৮০ + ৩০ = ১১০ টাকা}

\emph{∴ প্রতি ঘন্টার জন্য গড়ে মজুরি পায় = (১১০/১০) = ১১ টাকা}
\end{quote}

\textbf{৩টি বইয়ের দাম যথাক্রমে ২২ টাকা, ২৭ টাকা ও ২০ টাকা হলে বইগুলোর গড়
দাম কত?}

ক) ২৪ টাকা

\textbf{\hl{খ) ২৩ টাকা ✔}}

গ) ২৬ টাকা

ঘ) ২৫ টাকা

\textbf{\hl{✔ সঠিক উত্তর: খ) ২৩ টাকা}}

\begin{quote}
\textbf{ব্যাখ্যা:} \emph{৩টি বইয়ের মোট দাম = ২২ + ২৭ + ২০ = ৬৯ টাকা}

\emph{∴ বইগুলোর গড় দাম = ৬৯ ÷ ৩ = ২৩ টাকা}
\end{quote}

\textbf{৩টি রুমালের দাম যথাক্রমে ২২ টাকা, ২৭ টাকা ও ২০ টাকা হলে, রুমাল গুলোর
গড় দাম কত ?}

ক) ২৪ টাকা

\textbf{\hl{খ) ২৩ টাকা ✔}}

গ) ২৬ টাকা

ঘ) ২৫ টাকা

\textbf{\hl{✔ সঠিক উত্তর: খ) ২৩ টাকা}}

\begin{quote}
\textbf{ব্যাখ্যা:} \emph{৩টি রুমালের মোট দাম = ২২ + ২৭ + ২০ = ৬৯ টাকা}

\emph{∴ রুমালগুলোর গড় দাম = ৬৯ ÷ ৩ = ২৩ টাকা}
\end{quote}

\textbf{একজন ব্যবসায়ী চেয়ার ও টেবিল প্রস্তুত করে। প্রতিটি চেয়ারের বিক্রয়মূল্য
২,০০০ টাকা এবং প্রতিটি টেবিলের বিক্রয়মূল্য ২,৭৫০ টাকা। গত মাসে সে যে পরিমাণ
চেয়ার বিক্রি করে তার দ্বিগুণ পরিমাণ টেবিল বিক্রি করে থাকলে তার প্রতি একক
পণ্যের গড় বিক্রয়মূল্য কত?}

\textbf{\hl{ক) ২, ৫০০ টাকা ✔}}

খ) ২, ৫৫০ টাকা

গ) ২, ৬০০ টাকা

ঘ) ২, ৭০০ টাকা

\textbf{\hl{✔ সঠিক উত্তর: ক) ২, ৫০০ টাকা}}

\begin{quote}
\textbf{ব্যাখ্যা:} \emph{ধরি, চেয়ার বিক্রি করে x টি}

\emph{∴ টেবিল বিক্রি করে ২x টি}

\emph{মোট বিক্রি করা পণ্যের সংখ্যা = x + ২x = ৩x টি}

\emph{মোট বিক্রয়মূল্য = (২০০০ × x) + (২৭৫০ × ২x) = ২০০০x + ৫৫০০x = ৭৫০০x
টাকা}

\emph{∴ প্রতি একক পণ্যের গড় বিক্রয়মূল্য = (৭৫০০x/৩x) = ২৫০০ টাকা}
\end{quote}

\textbf{কালাম ও ৪ পুত্রের বয়সের গড় ২০ বছর। কালামের স্ত্রী ও ঐ ৪ পুত্রের বয়সের
গড় ১৭ বছর। কালামের বয়স ৪০ বছর হলে, স্ত্রীর বয়স কত ?}

\textbf{\hl{ক) ২৫ বছর ✔}}

খ) ৩০ বছর

গ) ২৪ বছর

ঘ) ২৬ বছর

\textbf{\hl{✔ সঠিক উত্তর: ক) ২৫ বছর}}

\begin{quote}
\textbf{ব্যাখ্যা:} \emph{কালাম ও ৪ পুত্রসহ মোট ৫ জনের বয়সের সমষ্টি = ২০ × ৫ =
১০০ বছর}

\emph{কালামের বয়স ৪০ বছর হলে, ৪ পুত্রের মোট বয়স = ১০০ − ৪০ = ৬০ বছর}

\emph{কালামের স্ত্রী ও ৪ পুত্রসহ মোট ৫ জনের বয়সের সমষ্টি = ১৭ × ৫ = ৮৫ বছর}

\emph{∴ কালামের স্ত্রীর বয়স = ৮৫ − ৬০ = ২৫ বছর}
\end{quote}

\textbf{তিন সদস্যের একটি বিতর্ক দলের সদস্যদের গড় বয়স ২৪ বছর। যদি কোনো
সদস্যের বয়সই ২১ বছরের নিচে না হয়, তবে তাদের কোনো একজনের সর্বোচ্চ বয়স কত ?}

\textbf{\hl{ক) ৩০ বছর ✔}}

খ) ২৫ বছর

গ) ২৮ বছর

ঘ) ৩২ বছর

\textbf{\hl{✔ সঠিক উত্তর: ক) ৩০ বছর}}

\begin{quote}
\textbf{ব্যাখ্যা:} \emph{তিন সদস্যের মোট বয়স = ২৪ × ৩ = ৭২ বছর}

\emph{একজনের সর্বোচ্চ বয়সের জন্য বাকি দুইজনের বয়স সর্বনিম্ন ২১ বছর হতে হবে।}

\emph{দুইজনের সর্বনিম্ন বয়সের সমষ্টি = ২১ × ২ = ৪২ বছর}

\emph{∴ একজনের সর্বোচ্চ বয়স = ৭২ − ৪২ = ৩০ বছর}
\end{quote}

\textbf{কটন আলীর গত তিন বছরের গড় আয় ৫৫,০০০ টাকা যদি সে ২য় বছরে ১ম বছরের
(৩/২) গুণ আয় করে থাকে এবং ৩য় বছরে ২য় বছরের (৫/২) গুণ আয় করে থাকে তাহলে তার
২য় ও ৩য় বছরের আয়ের গড় কত টাকা?}

\textbf{\hl{ক) ৬৯, ৩০০ টাকা ✔}}

খ) ৫৯, ৩০০ টাকা

গ) ৩৯, ৬০০ টাকা

ঘ) ২৬, ৪০০ টাকা

\textbf{\hl{✔ সঠিক উত্তর: ক) ৬৯, ৩০০ টাকা}}

\begin{quote}
\textbf{ব্যাখ্যা:} \emph{কটন আলীর ৩ বছরের মোট আয় = ৫৫০০০ × ৩ = ১৬৫০০০
টাকা}

\emph{ধরি, ১ম বছরে আয় = x টাকা}

\emph{∴ ২য় বছরে আয় = (৩x/২) টাকা}

\emph{∴ ৩য় বছরে আয় = (৩x/২) × (৫/২) = (১৫x/৪) টাকা}

\emph{প্রশ্নমতে, x + (৩x/২) + (১৫x/৪) = ১৬৫০০০}

\emph{⇒ (৪x+৬x+১৫x/৪) = ১৬৫০০০}

\emph{⇒ ২৫x = ৬৬০০০০}

\emph{∴ x = ২৬৪০০ টাকা (১ম বছরের আয়)}

\emph{∴ ২য় বছরের আয় = (৩×২৬৪০০/২) = ৩৯৬০০ টাকা}

\emph{∴ ৩য় বছরের আয় = (১৫×২৬৪০০/৪) = ৯৯০০০ টাকা}

\emph{∴ ২য় ও ৩য় বছরের আয়ের গড় = (৩৯৬০০+৯৯০০০/২) = (১৩৮৬০০/২) = ৬৯৩০০
টাকা}
\end{quote}

\textbf{পিতা ও ২ সন্তানের বয়সের গড় ৩৭ বছর। পিতা, মাতা ও ঐ ২ সন্তানের বয়সের
গড় ৩৬ বছর। মাতার বয়স কত বছর ?}

\textbf{\hl{ক) ৩৩ ✔}}

খ) ৩৪

গ) ৩৫

ঘ) ৩৬

\textbf{\hl{✔ সঠিক উত্তর: ক) ৩৩}}

\begin{quote}
\textbf{ব্যাখ্যা:} \emph{পিতা ও ২ সন্তানসহ মোট ৩ জনের বয়সের সমষ্টি = ৩৭ × ৩
= ১১১ বছর}

\emph{পিতা, মাতা ও ২ সন্তানসহ মোট ৪ জনের বয়সের সমষ্টি = ৩৬ × ৪ = ১৪৪ বছর}

\emph{∴ মাতার বয়স = ১৪৪ − ১১১ = ৩৩ বছর}
\end{quote}

\textbf{পিতা ও মাতার বয়সের গড় ৪৫ বছর। আবার পিতা, মাতা ও এক পুত্রের বয়সের
গড় ৩৬ বছর। পুত্রের বয়স-}

ক) ৯ বছর

খ) ১৪ বছর

গ) ১৫ বছর

\textbf{\hl{ঘ) ১৮ বছর ✔}}

\textbf{\hl{✔ সঠিক উত্তর: ঘ) ১৮ বছর}}

\begin{quote}
\textbf{ব্যাখ্যা:} \emph{পিতা ও মাতার বয়সের সমষ্টি = ৪৫ × ২ = ৯০ বছর}

\emph{পিতা, মাতা ও এক পুত্রসহ মোট ৩ জনের বয়সের সমষ্টি = ৩৬ × ৩ = ১০৮ বছর}

\emph{∴ পুত্রের বয়স = ১০৮ − ৯০ = ১৮ বছর}
\end{quote}

\textbf{১৫ জন ছাত্রের বয়সের গড় ২৯ বছর। তাদের মধ্যে দুজন ছাত্রের বয়সের গড় ৫৫
বছর। বাকী ১৩ জনের বয়সের গড় কত হবে ?}

ক) ২৯ বছর

খ) ২৭ বছর

গ) ২৬ বছর

\textbf{\hl{ঘ) ২৫ বছর ✔}}

\textbf{\hl{✔ সঠিক উত্তর: ঘ) ২৫ বছর}}

\begin{quote}
\textbf{ব্যাখ্যা:} \emph{১৫ জন ছাত্রের বয়সের সমষ্টি = ২৯ × ১৫ = ৪৩৫ বছর}

\emph{দুজন ছাত্রের বয়সের সমষ্টি = ৫৫ × ২ = ১১০ বছর}

\emph{বাকী (১৫ − ২) = ১৩ জন ছাত্রের বয়সের সমষ্টি = ৪৩৫ − ১১০ = ৩২৫ বছর}

\emph{∴ বাকী ১৩ জনের বয়সের গড় = ৩২৫ ÷ ১৩ = ২৫ বছর}
\end{quote}

\textbf{কোন শ্রেণীতে ২০ ছাত্রের বয়সের গড় ১০ বছর। শিক্ষকসহ তাদের বয়সের গড় ১২
বছর হলে শিক্ষকের বয়স কত ?}

ক) ৩২ বছর

খ) ৪২ বছর

গ) ৬২ বছর

\textbf{\hl{ঘ) ৫২ বছর ✔}}

\textbf{\hl{✔ সঠিক উত্তর: ঘ) ৫২ বছর}}

\begin{quote}
\textbf{ব্যাখ্যা:} \emph{২০ জন ছাত্রের বয়সের সমষ্টি = ১০ × ২০ = ২০০ বছর}

\emph{শিক্ষকসহ (২০ + ১) = ২১ জনের বয়সের সমষ্টি = ১২ × ২১ = ২৫২ বছর}

\emph{∴ শিক্ষকের বয়স = ২৫২ − ২০০ = ৫২ বছর}
\end{quote}

\textbf{৬ জন পুরুষ, ৮ জন স্ত্রীলোক এবং ১ জন বালকের বয়সের গড় ৩৫ বছর। পুরুষদের
বয়সের গড় ৪০ বছর এবং স্ত্রীলোকদের বয়সের গড় ৩৪ বছর। বালকের বয়স কত ?}

\textbf{\hl{ক) ১৩ বছর ✔}}

খ) ১৪ বছর

গ) ১৫ বছর

ঘ) ১৬ বছর

\textbf{\hl{✔ সঠিক উত্তর: ক) ১৩ বছর}}

\begin{quote}
\textbf{ব্যাখ্যা:} \emph{৬ জন পুরুষ, ৮ জন স্ত্রীলোক ও ১ জন বালকসহ মোট (৬ + ৮ +
১) = ১৫ জনের বয়সের সমষ্টি = ৩৫ × ১৫ = ৫২৫ বছর}

\emph{৬ জন পুরুষের বয়সের সমষ্টি = ৪০ × ৬ = ২৪০ বছর}

\emph{৮ জন স্ত্রীলোকের বয়সের সমষ্টি = ৩৪ × ৮ = ২৭২ বছর}

\emph{∴ বালকের বয়স = ৫২৫ − (২৪০ + ২৭২) = ৫২৫ − ৫১২ = ১৩ বছর}
\end{quote}

\textbf{৯ জন ছাত্রের গড় বয়স ১৫ বছর। ৩ জন ছাত্রের বয়সের গড় ১৭ বছর হলে বাকি
৬ জন ছাত্রের বয়সের গড় কত ?}

\textbf{\hl{ক) ১৪ বছর ✔}}

খ) ১৫ বছর

গ) ১৬ বছর

ঘ) ১৭ বছর

\textbf{\hl{✔ সঠিক উত্তর: ক) ১৪ বছর}}

\begin{quote}
\textbf{ব্যাখ্যা:} \emph{৯ জন ছাত্রের মোট বয়স = ১৫ × ৯ = ১৩৫ বছর}

\emph{৩ জন ছাত্রের বয়সের সমষ্টি = ১৭ × ৩ = ৫১ বছর}

\emph{বাকি (৯ − ৩) = ৬ জন ছাত্রের বয়সের সমষ্টি = ১৩৫ − ৫১ = ৮৪ বছর}

\emph{∴ বাকি ৬ জন ছাত্রের বয়সের গড় = ৮৪ ÷ ৬ = ১৪ বছর}
\end{quote}

\section{6. গড় (Average) মাস্টারক্লাস: বিগত সালের প্রশ্ন, ব্যাখ্যা ও লাইভ
এক্সাম}\label{ux997ux9a1-average-ux9aeux9b8ux99fux9b0ux995ux9b2ux9b8-ux9acux997ux9a4-ux9b8ux9b2ux9b0-ux9aaux9b0ux9b6ux9a8-ux9acux9afux996ux9af-ux993-ux9b2ux987ux9ad-ux98fux995ux9b8ux9ae-1}

বিসিএস, শিক্ষক নিবন্ধন কিংবা প্রাইমারি শিক্ষক নিয়োগসহ যেকোনো প্রতিযোগিতামূলক
পরীক্ষায় গণিতের `গড়' (Average) অধ্যায় থেকে প্রশ্ন আসবেই! তাই পরীক্ষার হলে কম
সময়ে নির্ভুল উত্তর দিতে প্রয়োজন সঠিক অনুশীলন। আপনাদের প্রস্তুতিকে নিখুঁত করতে এই
আর্টিকেলে বিগত সালের প্রশ্ন ও ব্যাখ্যাসহ সমাধানের পাশাপাশি যুক্ত করা হয়েছে
`লাইভ এক্সাম মোড' (Live Exam Mode)। এখন আপনি শুধু সমাধান শিখবেনই না, বরং
নিজে নিজে উত্তর সিলেক্ট করে তাৎক্ষণিক নিজের মেধা যাচাইও করতে পারবেন!

\subsubsection{গড় কী ?}\label{ux997ux9dc-ux995-1}

এক জাতীয় একাধিক রাশির সমষ্টিকে উক্ত রাশিগুলোর মোট সংখ্যা দ্বারা ভাগ করলে যে
ভাগফল পাওয়া যায় তাকে ঐ রাশিগুলোর গড় বলে।যেমন-

\textbf{একটি শ্রেণীতে ৫ জন বালকের ওজন যথাক্রমে ৪৮,৪২,৪৬,৫০,৫৯ কেজি হলে
তাদের গড় ওজন কত?}

\textbf{সমাধান:}

প্রথমে রাশিগুলোর যোগফল বের করতে হবে ৪৮+৪২+৪৬+৫০+৫৯ = ২৪৫ কেজি । এখন ই
যোগফলকে ৫ দিয়ে ভাগ করতে হবে কেননা মোট বালকের সংখ্যা ৫ জন । তাহলে তাদের গড়
বয়স হলো \(\frac{\text{২৪৫}}{\text{৫}}\) = ৪৯ কেজি ।

\subsubsection{গড় নির্ণয়ের সুত্র কী
?}\label{ux997ux9dc-ux9a8ux9b0ux9a3ux9dfux9b0-ux9b8ux9a4ux9b0-ux995-1}

গড় = \(\frac{\text{একজাতীয়}\ \text{কতিপয়}\ \text{রাশির}\ \text{সমষ্টি}}{\text{রাশির}\ \text{সংখ্যা}}\)

\textbf{100 জন শিক্ষার্থীর পরিসংখ্যানে গড় নম্বর 70। এদের মধ্যে 60 জন ছাত্রীর
গড় নম্বর 75 হলে, ছাত্রদের গড় নম্বর কত?}

ক) 55.5

খ) 60.5

গ) 65.5

\textbf{\hl{ঘ) 62.5 ✔}}

\textbf{\hl{✔ সঠিক উত্তর: ঘ) 62.5}}

\begin{quote}
\textbf{ব্যাখ্যা:} \emph{100 জন শিক্ষার্থীর পরিসংখ্যানে মোট নম্বর = 70 × 100 =
7000}

\emph{60 জন ছাত্রীর মোট নম্বর = 75 × 60 = 4500}

\emph{বাকি (100 − 60) = 40 জন ছাত্রদের মোট নম্বর = 7000 − 4500 = 2500}

\emph{∴ ছাত্রদের গড় নম্বর = 2500 ÷ 40 = 62.5}
\end{quote}

\textbf{৭টি সংখ্যার গড় ৪০। এর সাথে ৩টি সংখ্যা যোগ করা হলো। সংখ্যা ৩টির গড়
২১। সমষ্টিগত ভাবে ১০টি সংখ্যার গড় কত ?}

ক) ৬১

খ) ৬.১

গ) ৭০

\textbf{\hl{ঘ) ৩৪.৩ ✔}}

\textbf{\hl{✔ সঠিক উত্তর: ঘ) ৩৪.৩}}

\begin{quote}
\textbf{ব্যাখ্যা:} \emph{৭ টি সংখ্যার সমষ্টি = ৭ × ৪০ = ২৮০}

\emph{৩ টি সংখ্যার সমষ্টি = ৩ × ২১ = ৬৩}

\emph{১০ টি সংখ্যার সমষ্টি = ২৮০ + ৬৩ = ৩৪৩}

\emph{∴ ১০ টি সংখ্যার গড় = ৩৪৩ ÷ ১০ = ৩৪.৩}
\end{quote}

\textbf{১১টি সংখ্যার গড় ৩০। প্রথম পাঁচটি সংখ্যার গড় ২৫ ও শেষ পাঁচটি সংখ্যার গড়
২৮। ষষ্ঠ সংখ্যাটি কত ?}

ক) ৫৫

খ) ৫৮

\textbf{\hl{গ) ৬৫ ✔}}

ঘ) ৬৭

\textbf{\hl{✔ সঠিক উত্তর: গ) ৬৫}}

\begin{quote}
\textbf{ব্যাখ্যা:} \emph{১১টি সংখ্যার সমষ্টি = ৩০ × ১১ = ৩৩০}

\emph{প্রথম ৫টি সংখ্যার সমষ্টি = ২৫ × ৫ = ১২৫}

\emph{শেষ ৫টি সংখ্যার সমষ্টি = ২৮ × ৫ = ১৪০}

\emph{প্রথম ও শেষ ১০টি সংখ্যার সমষ্টি = ১২৫ + ১৪০ = ২৬৫}

\emph{∴ ষষ্ঠ সংখ্যাটি = ৩৩০ − ২৬৫ = ৬৫}
\end{quote}

\textbf{২৪ জন ছাত্র এবং একজন শিক্ষকের বয়সের গড় ১৫ বছর। শিক্ষককে বাদ দিয়ে
ছাত্রদের বয়সের গড় করলে গড় ১ বছর কমে যায়। শিক্ষকের বয়স কত ?}

ক) ৩৮

খ) ৪০

গ) ১২

\textbf{\hl{ঘ) কোনোটিই নয় ✔}}

\textbf{\hl{✔ সঠিক উত্তর: ঘ) কোনোটিই নয়}}

\begin{quote}
\textbf{ব্যাখ্যা:} \emph{২৪ জন ছাত্র এবং একজন শিক্ষকসহ মোট ২৫ জনের বয়সের
সমষ্টি = ২৫ × ১৫ = ৩৭৫ বছর}

\emph{শিক্ষককে বাদ দিলে ২৪ জন ছাত্রের বয়সের গড় হয় (১৫ − ১) = ১৪ বছর}

\emph{২৪ জন ছাত্রের বয়সের সমষ্টি = ২৪ × ১৪ = ৩৩৬ বছর}

\emph{∴ শিক্ষকের বয়স = ৩৭৫ − ৩৩৬ = ৩৯ বছর}
\end{quote}

\textbf{কোন শ্রেণীর ২৪ জন ছাত্রের বয়সের গড় ১৪ বছর। যদি একজন শ্রেণী শিক্ষকের
বয়স তাদের বয়সের সাথে যোগ করা হয় তবে বয়সের গড় এক বছর বৃদ্ধি পায়। শিক্ষকের
বয়স কত ?}

ক) ৪০ বৎসর

\textbf{\hl{খ) ৩৯ বৎসর ✔}}

গ) ৪৫ বৎসর

ঘ) ৩৫ বৎসর

\textbf{\hl{✔ সঠিক উত্তর: খ) ৩৯ বৎসর}}

\begin{quote}
\textbf{ব্যাখ্যা:} \emph{২৪ জন ছাত্রের বয়সের সমষ্টি = ২৪ × ১৪ = ৩৩৬ বছর}

\emph{শিক্ষকসহ (২৪ + ১) = ২৫ জনের বয়সের গড় হয় (১৪ + ১) = ১৫ বছর}

\emph{শিক্ষকসহ ২৫ জনের বয়সের সমষ্টি = ২৫ × ১৫ = ৩৭৫ বছর}

\emph{∴ শিক্ষকের বয়স = ৩৭৫ − ৩৩৬ = ৩৯ বছর}
\end{quote}

\textbf{ক ও খ এর মানের গড় ৯ এবং গ এর মান ১২ হলে ক, খ ও গ এর মানের গড় কত
হবে ?}

ক) ৬

খ) ৯

\textbf{\hl{গ) ১০ ✔}}

ঘ) ১২

\textbf{\hl{✔ সঠিক উত্তর: গ) ১০}}

\begin{quote}
\textbf{ব্যাখ্যা:} \emph{ক ও খ এর সমষ্টি = ৯ × ২ = ১৮}

\emph{ক, খ ও গ এর মানের সমষ্টি = ১৮ + ১২ = ৩০}

\emph{∴ ক, খ ও গ এর গড় = (৩০/৩) = ১০}
\end{quote}

\textbf{৩টি সংখ্যার গড় ৬ এবং ঐ ৩টি সংখ্যাসহ মোট ৪টি সংখ্যার গড় ৮ হলে চতুর্থ
সংখ্যাটির অর্ধেকের মান কত ?}

ক) ৮

খ) ৫

গ) ৬

\textbf{\hl{ঘ) ৭ ✔}}

\textbf{\hl{✔ সঠিক উত্তর: ঘ) ৭}}

\begin{quote}
\textbf{ব্যাখ্যা:} \emph{প্রথম ৩টি সংখ্যার যোগফল = ৩ × ৬ = ১৮}

\emph{৩টি সংখ্যাসহ ৪টি সংখ্যার যোগফল = ৪ × ৮ = ৩২}

\emph{∴ চতুর্থ সংখ্যাটি = (৩২ − ১৮) = ১৪}

\emph{∴ ৪র্থ সংখ্যার অর্ধেক = (১৪/২) = ৭}
\end{quote}

\textbf{x ও y এর মানের গড় ৯ এবং z = ১২ হলে x,y,z এর মানের গড় কত?}

ক) ৬

খ) ৯

\textbf{\hl{গ) ১০ ✔}}

ঘ) ১২

\textbf{\hl{✔ সঠিক উত্তর: গ) ১০}}

\begin{quote}
\textbf{ব্যাখ্যা:} \emph{x ও y এর যোগফল = ৯ × ২ = ১৮}

\emph{x,y,z এর মানের সমষ্টি = ১৮ + ১২ = ৩০}

\emph{∴ x,y,z এর গড় = (৩০/৩) = ১০}
\end{quote}

\textbf{৭, ১২, ১৭, ৩, ১১, ৬, ১৩, ৩ সংখ্যাগুলোর গড় কত?}

ক) ৬

খ) ৭

গ) ৮

\textbf{\hl{ঘ) ৯ ✔}}

\textbf{\hl{✔ সঠিক উত্তর: ঘ) ৯}}

\begin{quote}
\textbf{ব্যাখ্যা:} \emph{সংখ্যাগুলোর সমষ্টি = ৭ + ১২ + ১৭ + ৩ + ১১ + ৬ + ১৩ +
৩ = ৭২}

\emph{মোট সংখ্যা = ৮টি}

\emph{∴ গড় = (৭২/৮) = ৯}
\end{quote}

\textbf{0, 5, 7 এর গড় কত?}

ক) 6

খ) 0

\textbf{\hl{গ) 4 ✔}}

ঘ) 1

\textbf{\hl{✔ সঠিক উত্তর: গ) 4}}

\begin{quote}
\textbf{ব্যাখ্যা:} \emph{সংখ্যাগুলোর সমষ্টি = 0 + 5 + 7 = 12}

\emph{মোট সংখ্যা = 3টি}

\emph{∴ গড় = (12/3) = 4}
\end{quote}

\textbf{− 30 এবং − 40 সংখ্যা দুইটির গড় ব্যবধান কত?}

ক) - 35

খ) -55

\textbf{\hl{গ) 5 ✔}}

ঘ) - 5

\textbf{\hl{✔ সঠিক উত্তর: গ) 5}}

\begin{quote}
\textbf{ব্যাখ্যা:} \emph{গাণিতিক গড় = (−30+(−40)/2) = (−70/2) = −35}

\emph{∴ গড় ব্যবধান = (\textbar−30+35\textbar+\textbar−40+35\textbar/2) =
(\textbar5\textbar+\textbar−5\textbar/2) = (10/2) = 5}
\end{quote}

\textbf{(১/২) , (১/৪) , (৩/৪) এর গড় মান কত?}

ক) (১/৪)

খ) (৩/৪)

\textbf{\hl{গ) (১/২) ✔}}

ঘ) (২/৩)

\textbf{\hl{✔ সঠিক উত্তর: গ) (১/২)}}

\begin{quote}
\textbf{ব্যাখ্যা:} \emph{সংখ্যাগুলোর সমষ্টি = (১/২) + (১/৪) + (৩/৪) =
(২+১+৩/৪) = (৬/৪)}

\emph{∴ গড় মান = (৬/৪)÷৩ = (৬/৪)×(১/৩)=(১/২)}
\end{quote}

\textbf{M সংখ্যক সংখ্যার গড় A এবং N সংখ্যক সংখ্যার গড় B হলে সবগুলো সংখ্যার গড়
কত?}

ক) (A+B/2)

খ) (AM+BN/2)

\textbf{\hl{গ) (AM+BN/M+N) ✔}}

ঘ) (AM+BN/A+B)

\textbf{\hl{✔ সঠিক উত্তর: গ) (AM+BN/M+N)}}

\begin{quote}
\textbf{ব্যাখ্যা:} \emph{M সংখ্যক সংখ্যার সমষ্টি = MA}

\emph{N সংখ্যক সংখ্যার সমষ্টি = NB}

\emph{∴ মোট (M + N) সংখ্যক সংখ্যার সমষ্টি = MA + NB}

\emph{∴ সবগুলো সংখ্যার গড় = (AM+BN/M+N)}
\end{quote}

\textbf{তিনটি পূর্ণ সংখ্যার গড় ১৫০ এবং ক্ষুদ্রতম সংখ্যা দুইটির গড় ১২০। বৃহত্তম
সংখ্যাটি কত?}

ক) ২৩০

\textbf{\hl{খ) ২১০ ✔}}

গ) ২০০

ঘ) ১৯০

\textbf{\hl{✔ সঠিক উত্তর: খ) ২১০}}

\begin{quote}
\textbf{ব্যাখ্যা:} \emph{৩টি পূর্ণ সংখ্যার সমষ্টি = ১৫০ × ৩ = ৪৫০}

\emph{ক্ষুদ্রতম সংখ্যা দুটির সমষ্টি = ১২০ × ২ = ২৪০}

\emph{∴ বৃহত্তম সংখ্যাটি = ৪৫০ − ২৪০ = ২১০}
\end{quote}

\textbf{প্রথম ও দ্বিতীয় সংখ্যার গড় ২৫। প্রথম, দ্বিতীয় ও তৃতীয় সংখ্যার গড় ৩০
হলে, তৃতীয় সংখ্যাটি কত?}

ক) ২৫

\textbf{\hl{খ) ৪০ ✔}}

গ) ৯০

ঘ) ৫০

\textbf{\hl{✔ সঠিক উত্তর: খ) ৪০}}

\begin{quote}
\textbf{ব্যাখ্যা:} \emph{১ম, ২য় ও ৩য় সংখ্যার সমষ্টি = ৩ × ৩০ = ৯০}

\emph{১ম ও ২য় সংখ্যার সমষ্টি = ২ × ২৫ = ৫০}

\emph{∴ ৩য় সংখ্যাটি = ৯০ − ৫০ = ৪০}
\end{quote}

\textbf{৩টি সংখ্যার গড় ৩৩। দুইটি সংখ্যা ২৪ এবং ৪২ হলে অপর সংখ্যাটি কত?}

ক) ১২

খ) ২২

গ) ৩২

\textbf{\hl{ঘ) ৩৩ ✔}}

\textbf{\hl{✔ সঠিক উত্তর: ঘ) ৩৩}}

\begin{quote}
\textbf{ব্যাখ্যা:} \emph{৩টি সংখ্যার সমষ্টি = ৩৩ × ৩ = ৯৯}

\emph{২টি সংখ্যার সমষ্টি = ২৪ + ৪২ = ৬৬}

\emph{∴ অপর সংখ্যাটি = ৯৯ − ৬৬ = ৩৩}
\end{quote}

\textbf{৭টি সংখ্যার গড় ১২। একটি সংখ্যা বাতিল করলে গড় হয় ১১। বাতিলকৃত
সংখ্যাটি কত?}

ক) ১০

খ) ১২

গ) ১৫

\textbf{\hl{ঘ) ১৮ ✔}}

\textbf{\hl{✔ সঠিক উত্তর: ঘ) ১৮}}

\begin{quote}
\textbf{ব্যাখ্যা:} \emph{৭টি সংখ্যার সমষ্টি = ৭ × ১২ = ৮৪}

\emph{একটি সংখ্যা বাতিল করলে অবশিষ্ট (৭ − ১) = ৬টি সংখ্যার সমষ্টি = ৬ × ১১ =
৬৬}

\emph{∴ বাতিলকৃত সংখ্যাটি = ৮৪ − ৬৬ = ১৮}
\end{quote}

\textbf{ছয়টি সংখ্যার গড় ৬। যদি প্রত্যেকটি সংখ্যা থেকে ৩ বিয়োগ করা হয় তবে নতুন
সংখ্যাগুলোর গড় কত হবে?}

\textbf{\hl{ক) ৩ ✔}}

খ) ৩.৫

গ) ৪

ঘ) কোনোটিই নয়

\textbf{\hl{✔ সঠিক উত্তর: ক) ৩}}

\begin{quote}
\textbf{ব্যাখ্যা:} \emph{ছয়টি সংখ্যার সমষ্টি = (৬ × ৬) = ৩৬}

\emph{প্রত্যেক সংখ্যা থেকে ৩ বিয়োগ করলে নতুন ছয়টি সংখ্যার সমষ্টি = ৩৬ − (৬ ×
৩) = ৩৬ − ১৮ = ১৮}

\emph{∴ নতুন ছয়টি সংখ্যার গড় = (১৮/৬) = ৩}
\end{quote}

\textbf{১০টি সংখ্যার যোগফল ৪৬২। এদের প্রথম ৪টার গড় ৫২ এবং শেষের ৫টার গড়
৩৮। পঞ্চম সংখ্যাটি কত?}

ক) ৬০

\textbf{\hl{খ) ৬৪ ✔}}

গ) ৬২

ঘ) ৫০

\textbf{\hl{✔ সঠিক উত্তর: খ) ৬৪}}

\begin{quote}
\textbf{ব্যাখ্যা:} \emph{প্রথম ৪টির সমষ্টি = ৫২ × ৪ = ২০৮}

\emph{শেষ ৫টির সমষ্টি = ৩৮ × ৫ = ১৯০}

\emph{প্রথম ৪টি ও শেষের ৫টি (মোট ৯টি) সংখ্যার সমষ্টি = ২০৮ + ১৯০ = ৩৯৮}

\emph{∴ পঞ্চম সংখ্যাটি = ৪৬২ − ৩৯৮ = ৬৪}
\end{quote}

\textbf{২ থেকে শুরু করে পর পর পাঁচটি জোড়া সংখ্যার গড় কত হবে?}

ক) ৪

খ) ৫

\textbf{\hl{গ) ৬ ✔}}

ঘ) ৭

\textbf{\hl{✔ সঠিক উত্তর: গ) ৬}}

\begin{quote}
\textbf{ব্যাখ্যা:} \emph{২ থেকে শুরু করে পরপর ৫টি জোড় সংখ্যা হলো: ২, ৪, ৬, ৮,
১০}

\emph{সংখ্যাগুলোর সমষ্টি = ২ + ৪ + ৬ + ৮ + ১০ = ৩০}

\emph{∴ গড় = (৩০/৫) = ৬}
\end{quote}

\textbf{ছয়টি ধারাবাহিক বিজোড় সংখ্যার গড়ের ছয়গুণ সবচেয়ে বড় সংখ্যাটির চারগুণ
অপেক্ষা ১২ বেশি। ধারাবাহিক সংখ্যাগুলোর গড় কত?}

\textbf{\hl{ক) ১৬ ✔}}

খ) ১৭

গ) ১৮

ঘ) ২০

\textbf{\hl{✔ সঠিক উত্তর: ক) ১৬}}

\begin{quote}
\textbf{ব্যাখ্যা:} \emph{ধরি, ধারাবাহিক বিজোড় সংখ্যা ৬টি যথাক্রমে x, x + 2,
x + 4, x + 6, x + 8, x + 10}

\emph{সংখ্যাগুলোর সমষ্টি = 6x + 30}

\emph{গড় = (6x+30/6) = x + 5}

\emph{প্রশ্নমতে, 6(x + 5) = 4(x + 10) + 12}

\emph{⇒ 6x + 30 = 4x + 40 + 12}

\emph{⇒ 2x = 22}

\emph{∴ x = 11}

\emph{∴ সংখ্যাগুলোর গড় = x + 5 = 11 + 5 = 16}
\end{quote}

\textbf{৫টি সংখ্যার গড় ৪০। এর সাথে ৩টি সংখ্যা যোগ করা হলো, সংখ্যা তিনটির গড়
২২। সমষ্টিগতভাবে ৮টি সংখ্যার গড় কত?}

ক) ৩৩.২

খ) ৩৩.৫০

\textbf{\hl{গ) ৩৩.২৫ ✔}}

ঘ) কোনোটিই নয়

\textbf{\hl{✔ সঠিক উত্তর: গ) ৩৩.২৫}}

\begin{quote}
\textbf{ব্যাখ্যা:} \emph{৫টি সংখ্যার সমষ্টি = ৪০ × ৫ = ২০০}

\emph{অন্য ৩টি সংখ্যার সমষ্টি = ২২ × ৩ = ৬৬}

\emph{∴ মোট (৫ + ৩) = ৮টি সংখ্যার সমষ্টি = ২০০ + ৬৬ = ২৬৬}

\emph{∴ ৮টি সংখ্যার গড় = (২৬৬/৮) = ৩৩.২৫}
\end{quote}

\textbf{০.৬ হলো ০.২, ০.৮, ১ এবং x এর গড় মান। x এর মান কত?}

ক) ০.২

\textbf{\hl{খ) ০.৪ ✔}}

গ) ০.৬৭

ঘ) ২.৪

\textbf{\hl{✔ সঠিক উত্তর: খ) ০.৪}}

\begin{quote}
\textbf{ব্যাখ্যা:} \emph{৪টি সংখ্যার সমষ্টি = ০.২ + ০.৮ + ১ + x = ২ + x}

\emph{গড় = (২+x/৪)}

\emph{শর্তমতে, (২+x/৪) = ০.৬}

\emph{⇒ ২ + x = ২.৪}

\emph{∴ x = ২.৪ − ২ = ০.৪}
\end{quote}

\textbf{৬, ৮, ১০ এর গাণিতিক গড় ৭, ৯ এবং কোন সংখ্যার গাণিতিক গড়ের সমান?}

ক) ৫

\textbf{\hl{খ) ৮ ✔}}

গ) ৬

ঘ) ১০

\textbf{\hl{✔ সঠিক উত্তর: খ) ৮}}

\begin{quote}
\textbf{ব্যাখ্যা:} \emph{৬, ৮, ১০ এর গাণিতিক গড় = (৬+৮+১০/৩) = (২৪/৩) = ৮}

\emph{ধরি, নির্ণেয় সংখ্যাটি = x}

\emph{৭, ৯, x এর গড় = ৮}

\emph{⇒ ৭ + ৯ + x = ৩ × ৮}

\emph{⇒ ১৬ + x = ২৪}

\emph{∴ x = ২৪ − ১৬ = ৮}
\end{quote}

\textbf{পাঁচটি সংখ্যার গড় ৪৬ এবং সংখ্যাগুলোর প্রথম চারটি সংখ্যার গড় ৪৫। পঞ্চম
সংখ্যাটি কত?}

ক) ৪৮

খ) ৪৯

\textbf{\hl{গ) ৫০ ✔}}

ঘ) কোনোটিই নয়

\textbf{\hl{✔ সঠিক উত্তর: গ) ৫০}}

\begin{quote}
\textbf{ব্যাখ্যা:} \emph{পাঁচটি সংখ্যার সমষ্টি = ৫ × ৪৬ = ২৩০}

\emph{প্রথম চারটি সংখ্যার সমষ্টি = ৪ × ৪৫ = ১৮০}

\emph{∴ পঞ্চম সংখ্যাটি = ২৩০ − ১৮০ = ৫০}
\end{quote}

\textbf{একসেট সংখ্যা থেকে ৩৫ সংখ্যাটি বাদ দেয়ার ফলে সেটের গড় ১৪ থেকে ১১
হয়ে গেল। সেটের সদস্য সংখ্যা কত ছিল?}

\textbf{\hl{ক) ৮ ✔}}

খ) ১১

গ) ২১

ঘ) ২৫

\textbf{\hl{✔ সঠিক উত্তর: ক) ৮}}

\begin{quote}
\textbf{ব্যাখ্যা:} \emph{ধরি, সেটের প্রাথমিক সদস্য সংখ্যা ছিল x}

\emph{∴ x সংখ্যক সংখ্যার সমষ্টি = ১৪x}

\emph{৩৫ বাদ দেওয়ার পর সদস্য সংখ্যা হয় (x − ১) এবং সমষ্টি হয় ১১(x − ১)}

\emph{প্রশ্নমতে, ১৪x − ৩৫ = ১১(x − ১)}

\emph{⇒ ১৪x − ৩৫ = ১১x − ১১}

\emph{⇒ ১৪x − ১১x = ৩৫ − ১১}

\emph{⇒ ৩x = ২৪}

\emph{∴ x = ৮}
\end{quote}

\textbf{The average of 5 numbers is 40. If 2 more numbers with an
average of 21 are added to these numbers, what will be the average of
the combined numbers?}

ক) 8/7

খ) 30.1

গ) 30.3

\textbf{\hl{ঘ) 34.6 ✔}}

\textbf{\hl{✔ সঠিক উত্তর: ঘ) 34.6}}

\begin{quote}
\textbf{ব্যাখ্যা:} \emph{Sum of first 5 numbers = 5 × 40 = 200}

\emph{Sum of additional 2 numbers = 2 × 21 = 42}

\emph{Total sum of (5 + 2) = 7 numbers = 200 + 42 = 242}

\emph{∴ Average of combined numbers = (242/7) = 34.6}
\end{quote}

\textbf{একজন ব্যাটসম্যান প্রথম তিনটি T-20 খেলায় ৮২, ৮৫ ও ৯২ রান করেন। চতুর্থ
খেলায় কত রান করলে; তার গড় রান ৮৭ হবে?}

ক) ৮৬

খ) ৮৭

গ) ৮৮

\textbf{\hl{ঘ) ৮৯ ✔}}

\textbf{\hl{✔ সঠিক উত্তর: ঘ) ৮৯}}

\begin{quote}
\textbf{ব্যাখ্যা:} \emph{প্রথম তিনটি খেলায় মোট রান = ৮২ + ৮৫ + ৯২ = ২৫৯}

\emph{৪টি খেলায় মোট প্রয়োজনীয় রান = ৮৭ × ৪ = ৩৪৮}

\emph{∴ চতুর্থ খেলায় রান করতে হবে = ৩৪৮ − ২৫৯ = ৮৯}
\end{quote}

\textbf{পরীক্ষায় \textquotesingle ক\textquotesingle{} এর প্রাপ্ত নম্বর
যথাক্রমে ৮২, ৮৫ ও ৯২। চতুর্থ পরীক্ষায় তাকে কত নম্বর পেতে হবে যেন তার গড় প্রাপ্ত
নম্বর ৮৭ হয়?}

\textbf{\hl{ক) ৮৯ ✔}}

খ) ৮৮

গ) ৮৬

ঘ) ৯১

\textbf{\hl{✔ সঠিক উত্তর: ক) ৮৯}}

\begin{quote}
\textbf{ব্যাখ্যা:} \emph{পরীক্ষায় \textquotesingle ক\textquotesingle{} এর
প্রাপ্ত মোট নম্বর = (৮২ + ৮৫ + ৯২) = ২৫৯}

\emph{৪ টি পরীক্ষায় গড়ে ৮৭ হলে মোট নম্বর = ৮৭ × ৪ = ৩৪৮}

\emph{∴ চতুর্থ পরীক্ষায় তাকে পেতে হবে (৩৪৮ − ২৫৯) = ৮৯}
\end{quote}

\textbf{একজন ছাত্র বাংলায় ৬৫ এবং ইংলিশে ৮০ নম্বর পেল। সে অঙ্ক পরীক্ষায় কত
পেলে এই তিন বিষয়ে তার গড় নম্বর ৭৫ হবে?}

ক) ৭০

\textbf{\hl{খ) ৮০ ✔}}

গ) ৮৫

ঘ) কোনোটিই নয়

\textbf{\hl{✔ সঠিক উত্তর: খ) ৮০}}

\begin{quote}
\textbf{ব্যাখ্যা:} \emph{ধরি, অঙ্ক পরীক্ষায় x নম্বর পেতে হবে।}

\emph{∴ (৬৫+৮০+x/৩) = ৭৫}

\emph{⇒ ১৪৫ + x = ২২৫}

\emph{⇒ x = ২২৫ − ১৪৫}

\emph{∴ x = ৮০}
\end{quote}

\textbf{যদি ১৫ জন ছাত্র ইংরেজীতে গড়ে শতকরা ৮০ নম্বর এবং ১০ জন ছাত্র গড়ে
শতকরা ৯০ নম্বর পায়, তাহলে ২৫ জন ছাত্রের শতকরা হিসাবে গড় নম্বর কত?}

\textbf{\hl{ক) ৮৪ ✔}}

খ) ৮৪\%

গ) ৮৩\%

ঘ) ৮৩(১/৩) \%

\textbf{\hl{✔ সঠিক উত্তর: ক) ৮৪}}

\begin{quote}
\textbf{ব্যাখ্যা:} \emph{১৫ জনের মোট নম্বর = ১৫ × ৮০ = ১২০০}

\emph{১০ জনের মোট নম্বর = ১০ × ৯০ = ৯০০}

\emph{∴ ২৫ জনের মোট নম্বর = (১২০০ + ৯০০) = ২১০০}

\emph{∴ ২৫ জনের গড় নম্বর = ২১০০ ÷ ২৫ = ৮৪}
\end{quote}

\textbf{এক দোকানদার ১২ দিনে ৫০৪ টাকা আয় করলেন, প্রথম ৪ দিনের গড় আয় ৪০
টাকা হলে বাকী দিনগুলোর গড় আয় কত টাকা হবে?}

\textbf{\hl{ক) ৪৩ টাকা ✔}}

খ) ৪৭ টাকা

গ) ৪২ টাকা

ঘ) ৪০ টাকা

\textbf{\hl{✔ সঠিক উত্তর: ক) ৪৩ টাকা}}

\begin{quote}
\textbf{ব্যাখ্যা:} \emph{প্রথম ৪ দিনের গড় আয় ৪০ টাকা। মোট আয় = ৪ × ৪০ = ১৬০
টাকা।}

\emph{∴ বাকি (১২ − ৪) ৮ দিনে মোট আয় = ৫০৪ − ১৬০ = ৩৪৪ টাকা।}

\emph{∴ বাকি ১ দিনে গড় আয় = ৩৪৪ ÷ ৮ = ৪৩ টাকা।}
\end{quote}

\textbf{খ এর দৈনিক আয় ক এর দ্বিগুণ এবং গ এর দৈনিক আয় খ এর দ্বিগুণ। তাদের
তিনজনের আয়ের যোগফল ১১২০ টাকা হলে গ এর দৈনিক আয় কত টাকা?}

ক) ৫৬০

\textbf{\hl{খ) ৬৪০ ✔}}

গ) ৬৮৪

ঘ) কোনোটিই নয়

\textbf{\hl{✔ সঠিক উত্তর: খ) ৬৪০}}

\begin{quote}
\textbf{ব্যাখ্যা:} \emph{ধরি, ক এর আয় = x টাকা}

\emph{∴ খ এর আয় = ২x টাকা}

\emph{এবং গ এর আয় = ২ × ২x = ৪x টাকা}

\emph{প্রশ্নমতে, x + ২x + ৪x = ১১২০}

\emph{⇒ ৭x = ১১২০}

\emph{⇒ x = ১৬০}

\emph{∴ গ এর দৈনিক আয় = ৪ × ১৬০ = ৬৪০ টাকা।}
\end{quote}

\textbf{একজন ক্রিকেটারের ১০ ইনিংসের রানের গড় ৪৫.৫। ১১তম ইনিংসে কত রান করে
আউট হলে সব ইনিংস মিলিয়ে তার রানের গড় ৫০ হবে।}

ক) ৫৫ রান

খ) ৪৫ রান

গ) ১০০ রান

\textbf{\hl{ঘ) ৯৫ রান ✔}}

ঙ) কোনটিই নয়

\textbf{\hl{✔ সঠিক উত্তর: ঘ) ৯৫ রান}}

\begin{quote}
\textbf{ব্যাখ্যা:} \emph{১০ ইনিংসের রানের মোট সমষ্টি = ৪৫.৫ × ১০ = ৪৫৫}

\emph{১১ ইনিংসের রানের মোট সমষ্টি = ৫০ × ১১ = ৫৫০}

\emph{∴ ১১তম ইনিংসে প্রয়োজন = ৫৫০ − ৪৫৫ = ৯৫ রান}
\end{quote}

\textbf{এগারজন বালকের মধ্যে প্রথম ছয়জন বালকের গড় উচ্চতা ৮৭ সেমি। শেষের
পাঁচজন বালকের গড় উচ্চতা ১৩১ সেমি। ঐ এগারজন বালকের গড় উচ্চতা কত?}

\textbf{\hl{ক) ১০৭ সেমি ✔}}

খ) ১০৬ সেমি

গ) ১০৫ সেমি

ঘ) ১০৪ সেমি

\textbf{\hl{✔ সঠিক উত্তর: ক) ১০৭ সেমি}}

\begin{quote}
\textbf{ব্যাখ্যা:} \emph{প্রথম ৬ জন বালকের মোট উচ্চতা = ৮৭ × ৬ = ৫২২ সেমি}

\emph{শেষের ৫ জন বালকের মোট উচ্চতা = ১৩১ × ৫ = ৬৫৫ সেমি}

\emph{১১ জন বালকের মোট উচ্চতা = ৫২২ + ৬৫৫ = ১১৭৭ সেমি}

\emph{∴ ১১ জন বালকের গড় উচ্চতা = ১১৭৭ ÷ ১১ = ১০৭ সেমি}
\end{quote}

\textbf{কোন শ্রেণীতে ১০ জন ছাত্রের গড় উচ্চতা ৫ ফুট ৬ ইঞ্চি। এর মধ্যে ৯ জন
ছাত্রের গড় উচ্চতা ৫ ফুট ৫ ইঞ্চি হলে ১০ম ছাত্রের উচ্চতা কত?}

ক) ৫ ফুট ৭ ইঞ্চি

খ) ৬ ফুট ৫ ইঞ্চি

\textbf{\hl{গ) ৬ ফুট ৩ ইঞ্চি ✔}}

ঘ) ৬ ফুট

ঙ) কোনটিই নয়

\textbf{\hl{✔ সঠিক উত্তর: গ) ৬ ফুট ৩ ইঞ্চি}}

\begin{quote}
\textbf{ব্যাখ্যা:} \emph{১০ জন ছাত্রের মোট উচ্চতা = (৫ ফুট ৬ ইঞ্চি) × ১০ = ৫০
ফুট ৬০ ইঞ্চি}

\emph{৯ জন ছাত্রের মোট উচ্চতা = (৫ ফুট ৫ ইঞ্চি) × ৯ = ৪৫ ফুট ৪৫ ইঞ্চি}

\emph{∴ ১০ম ছাত্রের উচ্চতা = (৫০ ফুট ৬০ ইঞ্চি) − (৪৫ ফুট ৪৫ ইঞ্চি)}

\emph{= ৫ ফুট ১৫ ইঞ্চি}

\emph{= ৬ ফুট ৩ ইঞ্চি (যেহেতু ১২ ইঞ্চি = ১ ফুট)}
\end{quote}

\textbf{তিনজন লোকের ওজনের গড় ৫৩ কেজি। এদের কারো ওজন ৫১ কেজির কম নয়। এদের
একজনের সর্বোচ্চ ওজন হবে-}

ক) ৫৩ কেজি

খ) ৫৫ কেজি

\textbf{\hl{গ) ৫৭ কেজি ✔}}

ঘ) ৫৯ কেজি

\textbf{\hl{✔ সঠিক উত্তর: গ) ৫৭ কেজি}}

\begin{quote}
\textbf{ব্যাখ্যা:} \emph{তিনজন লোকের মোট ওজন = ৫৩ × ৩ = ১৫৯ কেজি}

\emph{একজনের সর্বোচ্চ ওজনের জন্য বাকি দুজনের ওজন সর্বনিম্ন হতে হবে।}

\emph{কারো ওজন ৫১ কেজির কম নয়, তাই দুজনের সর্বনিম্ন মোট ওজন = ৫১ × ২ = ১০২
কেজি}

\emph{∴ একজনের সর্বোচ্চ ওজন = ১৫৯ − ১০২ = ৫৭ কেজি}
\end{quote}

\textbf{১১ জন লোকের গড় ওজন ৭০ কেজি। ৯০ কেজি ওজনের একজন লোক চলে গেলে
বাকিদের গড় ওজন কত হয়?}

ক) ৬২ কেজি

\textbf{\hl{খ) ৬৮ কেজি ✔}}

গ) ৮০ কেজি

ঘ) ৭২ কেজি

\textbf{\hl{✔ সঠিক উত্তর: খ) ৬৮ কেজি}}

\begin{quote}
\textbf{ব্যাখ্যা:} \emph{১১ জন লোকের ওজনের সমষ্টি = ৭০ × ১১ = ৭৭০ কেজি}

\emph{৯০ কেজি ওজনের ১ জন লোক চলে গেলে বাকি (১১ − ১) = ১০ জনের ওজনের
সমষ্টি = ৭৭০ − ৯০ = ৬৮০ কেজি}

\emph{∴ বাকি ১০ জনের গড় ওজন = (৬৮০/১০) = ৬৮ কেজি}
\end{quote}

\textbf{ক, খ ও গ এর মাসিক গড় বেতন ৫০০ টাকা। খ, গ ও ঘ এর মাসিক গড় বেতন
৪৫০ টাকা। ক এর বেতন ৫৪০ টাকা হলে ঘ এর বেতন কত?}

ক) ৩৭৫ টাকা

খ) ৩৮০টাকা

গ) ৩৮৫ টাকা

\textbf{\hl{ঘ) ৩৯০ টাকা ✔}}

\textbf{\hl{✔ সঠিক উত্তর: ঘ) ৩৯০ টাকা}}

\begin{quote}
\textbf{ব্যাখ্যা:} \emph{ক + খ + গ-এর বেতনের সমষ্টি = ৫০০ × ৩ = ১৫০০ টাকা}

\emph{খ + গ + ঘ-এর বেতনের সমষ্টি = ৪৫০ × ৩ = ১৩৫০ টাকা}

\emph{ক ও ঘ-এর বেতনের পার্থক্য = ১৫০০ − ১৩৫০ = ১৫০ টাকা}

\emph{∴ ক − ঘ = ১৫০}

\emph{⇒ ৫৪০ − ঘ = ১৫০}

\emph{∴ ঘ এর বেতন = ৫৪০ − ১৫০ = ৩৯০ টাকা}
\end{quote}

\textbf{একজন শ্রমিক প্রতিদিন প্রথম ৮ ঘন্টা কাজ করার জন্য ঘন্টায় ১০ টাকা করে
এবং পরবর্তী সময়ের জন্য ঘন্টায় ১৫ টাকা করে মজুরি পায়। দৈনিক ১০ ঘন্টা করে কাজ
করলে তার ঘন্টা প্রতি গড় মজুরি কত?}

\textbf{\hl{ক) ১১ টাকা ✔}}

খ) ১২ টাকা

গ) ১২.৫০ টাকা

ঘ) ১৩ টাকা

\textbf{\hl{✔ সঠিক উত্তর: ক) ১১ টাকা}}

\begin{quote}
\textbf{ব্যাখ্যা:} \emph{প্রথম ৮ ঘন্টার জন্য মজুরি পায় = ১০ × ৮ = ৮০ টাকা}

\emph{পরবর্তী (১০ − ৮) = ২ ঘন্টার জন্য মজুরি পায় = ১৫ × ২ = ৩০ টাকা}

\emph{সর্বমোট ১০ ঘন্টার জন্য মোট মজুরি পায় = ৮০ + ৩০ = ১১০ টাকা}

\emph{∴ প্রতি ঘন্টার জন্য গড়ে মজুরি পায় = (১১০/১০) = ১১ টাকা}
\end{quote}

\textbf{৩টি বইয়ের দাম যথাক্রমে ২২ টাকা, ২৭ টাকা ও ২০ টাকা হলে বইগুলোর গড়
দাম কত?}

ক) ২৪ টাকা

\textbf{\hl{খ) ২৩ টাকা ✔}}

গ) ২৬ টাকা

ঘ) ২৫ টাকা

\textbf{\hl{✔ সঠিক উত্তর: খ) ২৩ টাকা}}

\begin{quote}
\textbf{ব্যাখ্যা:} \emph{৩টি বইয়ের মোট দাম = ২২ + ২৭ + ২০ = ৬৯ টাকা}

\emph{∴ বইগুলোর গড় দাম = ৬৯ ÷ ৩ = ২৩ টাকা}
\end{quote}

\textbf{৩টি রুমালের দাম যথাক্রমে ২২ টাকা, ২৭ টাকা ও ২০ টাকা হলে, রুমাল গুলোর
গড় দাম কত ?}

ক) ২৪ টাকা

\textbf{\hl{খ) ২৩ টাকা ✔}}

গ) ২৬ টাকা

ঘ) ২৫ টাকা

\textbf{\hl{✔ সঠিক উত্তর: খ) ২৩ টাকা}}

\begin{quote}
\textbf{ব্যাখ্যা:} \emph{৩টি রুমালের মোট দাম = ২২ + ২৭ + ২০ = ৬৯ টাকা}

\emph{∴ রুমালগুলোর গড় দাম = ৬৯ ÷ ৩ = ২৩ টাকা}
\end{quote}

\textbf{একজন ব্যবসায়ী চেয়ার ও টেবিল প্রস্তুত করে। প্রতিটি চেয়ারের বিক্রয়মূল্য
২,০০০ টাকা এবং প্রতিটি টেবিলের বিক্রয়মূল্য ২,৭৫০ টাকা। গত মাসে সে যে পরিমাণ
চেয়ার বিক্রি করে তার দ্বিগুণ পরিমাণ টেবিল বিক্রি করে থাকলে তার প্রতি একক
পণ্যের গড় বিক্রয়মূল্য কত?}

\textbf{\hl{ক) ২, ৫০০ টাকা ✔}}

খ) ২, ৫৫০ টাকা

গ) ২, ৬০০ টাকা

ঘ) ২, ৭০০ টাকা

\textbf{\hl{✔ সঠিক উত্তর: ক) ২, ৫০০ টাকা}}

\begin{quote}
\textbf{ব্যাখ্যা:} \emph{ধরি, চেয়ার বিক্রি করে x টি}

\emph{∴ টেবিল বিক্রি করে ২x টি}

\emph{মোট বিক্রি করা পণ্যের সংখ্যা = x + ২x = ৩x টি}

\emph{মোট বিক্রয়মূল্য = (২০০০ × x) + (২৭৫০ × ২x) = ২০০০x + ৫৫০০x = ৭৫০০x
টাকা}

\emph{∴ প্রতি একক পণ্যের গড় বিক্রয়মূল্য = (৭৫০০x/৩x) = ২৫০০ টাকা}
\end{quote}

\textbf{কালাম ও ৪ পুত্রের বয়সের গড় ২০ বছর। কালামের স্ত্রী ও ঐ ৪ পুত্রের বয়সের
গড় ১৭ বছর। কালামের বয়স ৪০ বছর হলে, স্ত্রীর বয়স কত ?}

\textbf{\hl{ক) ২৫ বছর ✔}}

খ) ৩০ বছর

গ) ২৪ বছর

ঘ) ২৬ বছর

\textbf{\hl{✔ সঠিক উত্তর: ক) ২৫ বছর}}

\begin{quote}
\textbf{ব্যাখ্যা:} \emph{কালাম ও ৪ পুত্রসহ মোট ৫ জনের বয়সের সমষ্টি = ২০ × ৫ =
১০০ বছর}

\emph{কালামের বয়স ৪০ বছর হলে, ৪ পুত্রের মোট বয়স = ১০০ − ৪০ = ৬০ বছর}

\emph{কালামের স্ত্রী ও ৪ পুত্রসহ মোট ৫ জনের বয়সের সমষ্টি = ১৭ × ৫ = ৮৫ বছর}

\emph{∴ কালামের স্ত্রীর বয়স = ৮৫ − ৬০ = ২৫ বছর}
\end{quote}

\textbf{তিন সদস্যের একটি বিতর্ক দলের সদস্যদের গড় বয়স ২৪ বছর। যদি কোনো
সদস্যের বয়সই ২১ বছরের নিচে না হয়, তবে তাদের কোনো একজনের সর্বোচ্চ বয়স কত ?}

\textbf{\hl{ক) ৩০ বছর ✔}}

খ) ২৫ বছর

গ) ২৮ বছর

ঘ) ৩২ বছর

\textbf{\hl{✔ সঠিক উত্তর: ক) ৩০ বছর}}

\begin{quote}
\textbf{ব্যাখ্যা:} \emph{তিন সদস্যের মোট বয়স = ২৪ × ৩ = ৭২ বছর}

\emph{একজনের সর্বোচ্চ বয়সের জন্য বাকি দুইজনের বয়স সর্বনিম্ন ২১ বছর হতে হবে।}

\emph{দুইজনের সর্বনিম্ন বয়সের সমষ্টি = ২১ × ২ = ৪২ বছর}

\emph{∴ একজনের সর্বোচ্চ বয়স = ৭২ − ৪২ = ৩০ বছর}
\end{quote}

\textbf{কটন আলীর গত তিন বছরের গড় আয় ৫৫,০০০ টাকা যদি সে ২য় বছরে ১ম বছরের
(৩/২) গুণ আয় করে থাকে এবং ৩য় বছরে ২য় বছরের (৫/২) গুণ আয় করে থাকে তাহলে তার
২য় ও ৩য় বছরের আয়ের গড় কত টাকা?}

\textbf{\hl{ক) ৬৯, ৩০০ টাকা ✔}}

খ) ৫৯, ৩০০ টাকা

গ) ৩৯, ৬০০ টাকা

ঘ) ২৬, ৪০০ টাকা

\textbf{\hl{✔ সঠিক উত্তর: ক) ৬৯, ৩০০ টাকা}}

\begin{quote}
\textbf{ব্যাখ্যা:} \emph{কটন আলীর ৩ বছরের মোট আয় = ৫৫০০০ × ৩ = ১৬৫০০০
টাকা}

\emph{ধরি, ১ম বছরে আয় = x টাকা}

\emph{∴ ২য় বছরে আয় = (৩x/২) টাকা}

\emph{∴ ৩য় বছরে আয় = (৩x/২) × (৫/২) = (১৫x/৪) টাকা}

\emph{প্রশ্নমতে, x + (৩x/২) + (১৫x/৪) = ১৬৫০০০}

\emph{⇒ (৪x+৬x+১৫x/৪) = ১৬৫০০০}

\emph{⇒ ২৫x = ৬৬০০০০}

\emph{∴ x = ২৬৪০০ টাকা (১ম বছরের আয়)}

\emph{∴ ২য় বছরের আয় = (৩×২৬৪০০/২) = ৩৯৬০০ টাকা}

\emph{∴ ৩য় বছরের আয় = (১৫×২৬৪০০/৪) = ৯৯০০০ টাকা}

\emph{∴ ২য় ও ৩য় বছরের আয়ের গড় = (৩৯৬০০+৯৯০০০/২) = (১৩৮৬০০/২) = ৬৯৩০০
টাকা}
\end{quote}

\textbf{পিতা ও ২ সন্তানের বয়সের গড় ৩৭ বছর। পিতা, মাতা ও ঐ ২ সন্তানের বয়সের
গড় ৩৬ বছর। মাতার বয়স কত বছর ?}

\textbf{\hl{ক) ৩৩ ✔}}

খ) ৩৪

গ) ৩৫

ঘ) ৩৬

\textbf{\hl{✔ সঠিক উত্তর: ক) ৩৩}}

\begin{quote}
\textbf{ব্যাখ্যা:} \emph{পিতা ও ২ সন্তানসহ মোট ৩ জনের বয়সের সমষ্টি = ৩৭ × ৩
= ১১১ বছর}

\emph{পিতা, মাতা ও ২ সন্তানসহ মোট ৪ জনের বয়সের সমষ্টি = ৩৬ × ৪ = ১৪৪ বছর}

\emph{∴ মাতার বয়স = ১৪৪ − ১১১ = ৩৩ বছর}
\end{quote}

\textbf{পিতা ও মাতার বয়সের গড় ৪৫ বছর। আবার পিতা, মাতা ও এক পুত্রের বয়সের
গড় ৩৬ বছর। পুত্রের বয়স-}

ক) ৯ বছর

খ) ১৪ বছর

গ) ১৫ বছর

\textbf{\hl{ঘ) ১৮ বছর ✔}}

\textbf{\hl{✔ সঠিক উত্তর: ঘ) ১৮ বছর}}

\begin{quote}
\textbf{ব্যাখ্যা:} \emph{পিতা ও মাতার বয়সের সমষ্টি = ৪৫ × ২ = ৯০ বছর}

\emph{পিতা, মাতা ও এক পুত্রসহ মোট ৩ জনের বয়সের সমষ্টি = ৩৬ × ৩ = ১০৮ বছর}

\emph{∴ পুত্রের বয়স = ১০৮ − ৯০ = ১৮ বছর}
\end{quote}

\textbf{১৫ জন ছাত্রের বয়সের গড় ২৯ বছর। তাদের মধ্যে দুজন ছাত্রের বয়সের গড় ৫৫
বছর। বাকী ১৩ জনের বয়সের গড় কত হবে ?}

ক) ২৯ বছর

খ) ২৭ বছর

গ) ২৬ বছর

\textbf{\hl{ঘ) ২৫ বছর ✔}}

\textbf{\hl{✔ সঠিক উত্তর: ঘ) ২৫ বছর}}

\begin{quote}
\textbf{ব্যাখ্যা:} \emph{১৫ জন ছাত্রের বয়সের সমষ্টি = ২৯ × ১৫ = ৪৩৫ বছর}

\emph{দুজন ছাত্রের বয়সের সমষ্টি = ৫৫ × ২ = ১১০ বছর}

\emph{বাকী (১৫ − ২) = ১৩ জন ছাত্রের বয়সের সমষ্টি = ৪৩৫ − ১১০ = ৩২৫ বছর}

\emph{∴ বাকী ১৩ জনের বয়সের গড় = ৩২৫ ÷ ১৩ = ২৫ বছর}
\end{quote}

\textbf{কোন শ্রেণীতে ২০ ছাত্রের বয়সের গড় ১০ বছর। শিক্ষকসহ তাদের বয়সের গড় ১২
বছর হলে শিক্ষকের বয়স কত ?}

ক) ৩২ বছর

খ) ৪২ বছর

গ) ৬২ বছর

\textbf{\hl{ঘ) ৫২ বছর ✔}}

\textbf{\hl{✔ সঠিক উত্তর: ঘ) ৫২ বছর}}

\begin{quote}
\textbf{ব্যাখ্যা:} \emph{২০ জন ছাত্রের বয়সের সমষ্টি = ১০ × ২০ = ২০০ বছর}

\emph{শিক্ষকসহ (২০ + ১) = ২১ জনের বয়সের সমষ্টি = ১২ × ২১ = ২৫২ বছর}

\emph{∴ শিক্ষকের বয়স = ২৫২ − ২০০ = ৫২ বছর}
\end{quote}

\textbf{৬ জন পুরুষ, ৮ জন স্ত্রীলোক এবং ১ জন বালকের বয়সের গড় ৩৫ বছর। পুরুষদের
বয়সের গড় ৪০ বছর এবং স্ত্রীলোকদের বয়সের গড় ৩৪ বছর। বালকের বয়স কত ?}

\textbf{\hl{ক) ১৩ বছর ✔}}

খ) ১৪ বছর

গ) ১৫ বছর

ঘ) ১৬ বছর

\textbf{\hl{✔ সঠিক উত্তর: ক) ১৩ বছর}}

\begin{quote}
\textbf{ব্যাখ্যা:} \emph{৬ জন পুরুষ, ৮ জন স্ত্রীলোক ও ১ জন বালকসহ মোট (৬ + ৮ +
১) = ১৫ জনের বয়সের সমষ্টি = ৩৫ × ১৫ = ৫২৫ বছর}

\emph{৬ জন পুরুষের বয়সের সমষ্টি = ৪০ × ৬ = ২৪০ বছর}

\emph{৮ জন স্ত্রীলোকের বয়সের সমষ্টি = ৩৪ × ৮ = ২৭২ বছর}

\emph{∴ বালকের বয়স = ৫২৫ − (২৪০ + ২৭২) = ৫২৫ − ৫১২ = ১৩ বছর}
\end{quote}

\textbf{৯ জন ছাত্রের গড় বয়স ১৫ বছর। ৩ জন ছাত্রের বয়সের গড় ১৭ বছর হলে বাকি
৬ জন ছাত্রের বয়সের গড় কত ?}

\textbf{\hl{ক) ১৪ বছর ✔}}

খ) ১৫ বছর

গ) ১৬ বছর

ঘ) ১৭ বছর

\textbf{\hl{✔ সঠিক উত্তর: ক) ১৪ বছর}}

\begin{quote}
\textbf{ব্যাখ্যা:} \emph{৯ জন ছাত্রের মোট বয়স = ১৫ × ৯ = ১৩৫ বছর}

\emph{৩ জন ছাত্রের বয়সের সমষ্টি = ১৭ × ৩ = ৫১ বছর}

\emph{বাকি (৯ − ৩) = ৬ জন ছাত্রের বয়সের সমষ্টি = ১৩৫ − ৫১ = ৮৪ বছর}

\emph{∴ বাকি ৬ জন ছাত্রের বয়সের গড় = ৮৪ ÷ ৬ = ১৪ বছর}
\end{quote}

\section{7. ভেনচিত্র , সেট ও
ফাংশন}\label{ux9adux9a8ux99aux9a4ux9b0-ux9b8ux99f-ux993-ux9abux9b6ux9a8}

সেট শব্দটি আমাদের সুপরিচিত যেমন: ডিনার সেট, স্বাভাবিক সংখ্যার সেট, মূলদ
সংখ্যার সেট ইত্যাদি। আধুনিক হাতিয়ার হিসাবে সেটের ব্যবহার ব্যাপক।

জার্মান গণিতবিদ জর্জ ক্যান্টর (১৮৪৫ - ১৯১৮) সেট সম্পর্কে প্রথম ধারণা ব্যাখ্যা
করেন। তিনি অসীম সেটের ধারণা প্রদান করে গণিত শাস্ত্রে আলোড়ন সৃষ্টি করেন এবং
তাঁর সেটের ধারণা সেট তত্ত্ব নামে পরিচিত।

এই আর্টিকেলে সেটের ধারণা ব্যবহার করে গাণিতিক যুক্তি ও চিত্রের মাধ্যমে সমস্যা
সমাধান এবং ভেনচিত্র সম্পর্কে স্পষ্ট ধারণা দেওয়া হবে ইনশাআল্লাহ।

\begin{itemize}
\item
  প্রাথমিক আলোচনা
\item
  সেট প্রকাশের পদ্ধতি
\item
  বিভিন্ন ধরনের সেট
\end{itemize}

\begin{itemize}
\item
  সসীম সেট ও অসীম সেট
\item
  ফাঁকা সেট , ভেনচিত্র , উপসেট ও প্রকৃত উপসেট
\item
  ছেদ সেট ও সংযোগ সেট
\item
  পূরক সেট , সংযোগ সেট , শক্তি সেট ও নিশ্ছেদ সেট
\item
  সেটের সমতা ও অন্তর
\item
  কার্তেসীয় গুণজ
\end{itemize}

\begin{itemize}
\item
  উপসেট সংখ্যা নির্ণয়
\item
  সেট এর সাধারণ প্রশ্নসমূহ
\item
  ভেনচিত্রের সাধারণ প্রশ্ন
\end{itemize}

\subsubsection{প্রাথমিক
আলোচনা}\label{ux9aaux9b0ux9a5ux9aeux995-ux986ux9b2ux99aux9a8-1}

\textbf{সেটের ধারণা (Concept):}

সেটের প্রথম ধারণা দেন জার্মান গণিতবিদ জর্জ ক্যান্টর (১৮৪৫ - ১৯১৮) । তাকে সেট
থিওরীর জনক বলা হয় ।

\textbf{সেটের সংজ্ঞা (Definition):}

বাস্তব বা চিন্তা জগতের সু-সংজ্ঞায়িত বস্তুর সমাবেশ বা সংগ্রহকে সেট বলে। সেট
সাধারণত ইংরেজি বড় হাতের অক্ষর দ্বারা নির্দেশ করে ।

উপাদানকে ইংরেজি অক্ষরে লিখলে ছোট হাতের অক্ষর দ্বারা লেখতে হয়। সেটের উপাদান
সমূহকে সাধারণত দ্বিতীয় বন্ধনীর মধ্যে রাখা হয়।

সেটকে সাধারণত ``\{\}'' দ্বারা প্রকাশ করা হয় এবং উপাদানগুলোকে আলাদা করার জন্য
কমা ব্যবহার করা হয় ।

যেমন: A = \{a,b,c\} এখানে A হচ্ছে একটি সেট এবং a,b,c হবে A সেটের সদস্য বা
উপাদান ।

\textbf{সদস্য বা উপাদান (Elements):} সেটের ভিতরে যে সকল জিনিস বা বস্তু
ব্যবহার করা হয় উহাদেরকে সেটের উপাদান বলা হয়।

কিছু গুরুত্বপূর্ণ চিহ্ন সেটের প্রশ্নগুলোর সমাধান করার জন্য যাদের ব্যাবহার জানা জরুরী।
যেমন: A = \{ 1,2,3,4\}

এখানে A সেটের উপদানগুলো হলো 1,2,3,4

∪ = Union (সংযোগ সেট)

∩ = Intersection (ছেদ সেট)

⊆ = Subset (উপসেট)

⊂ = Proper Subset (প্রকৃত উপসেট)

∈ = belongs to (ভিতরে থাকে)

∉ = Not belongs to (ভিতরে থাকে না)

``:'' = such that ( যেন )

\textbf{উপসেট (Subset):}

যদি A সেটের প্রত্যেক উপাদান B সেটের উপাদান হয়, তবে A কে B এর উপসেট বলে এবং
এদেরকে প্রকাশ করা যায় A ⊂ B প্রতীকে, A ⊂ B দ্বারা বোঝায় A, B এর উপসেট।
যেমন: A = \{4,6,8\}, B = \{1,2,3,4,5,6,7,8\} হলে

A ⊂ B, B এর একটি উপসেট।

\textbf{একটি সেটের ভিতরে যতগুলো উপাদান থাকে উহাদের দ্বারা পৃথক পৃথক ভাবে
গঠিত সেটকে উপসেট বলে।} সেট প্রকাশের পদ্ধতি

সেটকে দুই পদ্ধতিতে প্রকাশ করা হয়। যথা: তালিকা পদ্ধতি (Roster Method বা
Tabular Method) ও সেট গঠন পদ্ধতি (Set Builder Method)।

\textbf{তালিকা পদ্ধতি:}

এ পদ্ধতিতে সেটের সকল উপাদান সুনির্দিষ্টভাবে উল্লেখ করে দ্বিতীয় বন্ধনী \{\} এর
মধ্যে আবদ্ধ করা হয় এবং একাধিক উপাদান থাকলে
\textquotesingle কমা\textquotesingle{} ব্যবহার করে উপাদানগুলোকে আলাদা করা
হয়।

যেমন, A = \{a, b\}, B = \{2, 4, 6\}, C = \{নিলয়, তিশা, শুভ্ৰা\} ইত্যাদি।

\textbf{সেট গঠন পদ্ধতি :} এ পদ্ধতিতে সেটের সকল উপাদান সুনির্দিষ্টভাবে উল্লেখ
না করে উপাদান নির্ধারণের জন্য সাধারণ ধর্মের উল্লেখ থাকে।

যেমন: A \{x : x = স্বাভাবিক বিজোড় সংখ্যা \}, B \{x: x = নবম শ্রেণির প্রথম
পাঁচজন শিক্ষার্থী\} ইত্যাদি। এখানে, : দ্বারা `এরূপ যেন\textquotesingle{} বা
সংক্ষেপে \textquotesingle যেন\textquotesingle{} (such that) বোঝায়।

যেহেতু এ পদ্ধতিতে সেটের উপাদান নির্ধারণের জন্য শর্ত বা নিয়ম (Rule) দেওয়া
থাকে, এ জন্য এ পদ্ধতিকে Rule Method ও বলা হয়।

বিভিন্ন ধরনের সেট এ পর্যায়ে আমরা বিভিন্ন ধরনের সেট সম্পর্কে অবগত হব। নতুবা
বিগত সালের প্রশ্ন অনুশীলন বা সমাধান করার সময় স্পষ্ট ধারণা লাভ করতে পারব না ।
কাজেই এই অংশটি মনোযোগ দিয়ে পড়ার পরামর্শ রইলো। সসীম সেট ও অসীম সেট
\textbf{সসীম সেট (Finite Set) :}

যে সেটের উপাদান সংখ্যা গণনা করে নির্ধারণ করা যায়, তাকে সসীম সেট বলে। যেমন,

D = \{x, y, 2\}, E = \{3, 6, 9, ..... 60\}, F = \{x : x মৌলিক সংখ্যা এবং
30 \textless{} x \textless{} 70\}

ইত্যাদি সসীম সেট। এখানে, D সেটে 3টি, E সেটে 20 টি এবং F সেটে 9 টি উপাদান
আছে।

\textbf{অসীম সেট (Infinite Set) :}

যে সেটের উপাদান সংখ্যা গণনা করে শেষ করা যায় না, তাকে অসীম সেট বলে। যেমন,
A = \{x : x বিজোড় স্বাভাবিক সংখ্যা\}, স্বাভাবিক সংখ্যার সেট N = \{1,2,3,4,
...\}, পূর্ণসংখ্যার সেট Z = ..., -3, -2, -1, 0, 1, 2, 3 ,....\}, মূলদ সংখ্যার
সেট Q = \{ \(\frac{a}{b}\): a ও b পূর্ণসংখ্যা এবং

b ≠ 0\}, বাস্তব সংখ্যার সেট R ইত্যাদি অসীম সেট।

ফাঁকা সেট , ভেনচিত্র , উপসেট ও প্রকৃত উপসেট \textbf{ফাঁকা সেট (Empty Set):}

যে সেটের কোনো উপাদান নেই তাকে ফাঁকা সেট বলে। ফাঁকা সেটকে ∅ দ্বারা প্রকাশ করা
হয়। যেমন: হলি ক্রস স্কুলের তিনজন ছাত্রের (পুরুষ) সেট, \{ x∈N : 10 \textless{}
x \textless{} 11\}, \{x ∈ N : মৌলিক সংখ্যা এবং 23 \textless{} x
\textless{} 29\} ইত্যাদি।

\textbf{ভেনচিত্র (Venn Diagram) :}

জন ভেন (১৮৩৪-১৯২৩) সেটের কার্যবিধি চিত্রের সাহায্যে প্রকাশ করেন। এতে
বিবেচনাধীন সেটগুলোকে সমতলে অবস্থিত বিভিন্ন আকারের জ্যামিতিক চিত্র যেমন আয়ত,
বৃত্ত এবং ত্রিভুজ ব্যবহার করা হয়। জন ভেনের নামানুসারে চিত্রগুলো ভেন চিত্র নামে
পরিচিত।

\textbf{উপসেট (Subset):}

A = \{a,b\} একটি সেট। এই সেটের উপাদান থেকে \{a, b\}, \{a\}, \{b\} সেটগুলো
গঠন করা যায়। আবার, কোনো উপাদান না নিয়ে Ø সেট গঠন কর যায়। এখানে, গঠিত
\{a, b\}, \{a\}, \{b\}, Ø প্রত্যেকটি A সেটের উপসেট।

সুতরাং কোনো সেট থেকে যতগুলো সেট গঠন করা যায়, এদের প্রত্যেকটি সেটকে ঐ সেটের
উপসেট বলা হয়। উপসেটের চিহ্ন ⊆। যদি B সেট A এর উপসেট হয় তবে B ⊆ A লেখা হয়।
B, A এর উপসেট অথবা B is a subset of A। উপরের উপসেটগুলোর মধ্যে \{a, b\} সেট
A এর সমান। প্রত্যেকটি সেট নিজের উপসেট। আবার, যেকোনো সেট থেকে Ø সেট গঠন করা
যায়। সুতরাং Ø সেট যেকোনো সেটের উপসেট।

ধরি P = \{1,2,3\} এবং Q \{2,3\}, R = \{1,3\} তাহলে P, Q এবং R প্রত্যেকে P
এর উপসেট। অর্থাৎ P⊆P, P⊆Q, P⊆R.

\textbf{প্রকৃত উপসেট (Proper subset) :}

কোনো সেট থেকে গঠিত উপসেটের মধ্যে যে উপসেটগুলোর উপাদান সংখ্যা প্রদত্ত সেটের
উপাদান সংখ্যা অপেক্ষা কম এদেরকে প্রকৃত উপসেট বলে। যেমন, A = \{ 3, 4, 5, 6\}
এবং B = (3,5)

দুইটি সেট। এখানে, B এর সব উপাদান A সেটে বিদ্যমান এবং B সেটের উপাদান সংখ্যা
A সেটের উপাদান সংখ্যা থেকে কম।

.:. B, A এর একটি প্রকৃত উপসেট এবং B ⊆ A লিখে প্রকাশ করা হয়।

উপসেটের উদাহরণে Q ও R প্রত্যেকে P এর প্রকৃত উপসেট। উল্লেখ্য ফাঁকা সেট বা Ø
যেকোনো সেটের প্রকৃত উপসেট।

ছেদ সেট ও সংযোগ সেট \textbf{ছেদ সেট (Intersection of Sets ):}

দুই বা ততোধিক সেটের সাধারণ উপাদান নিয়ে গঠিত সেটকে ছেদ সেট বলে। মনে করি, A
ও B দুইটি সেট। A ও B এর ছেদ সেটকে A ∩ B দ্বারা প্রকাশ করা হয় এবং পড়া হয় A
ছেদ B বা A intersection B। সেট গঠন পদ্ধতিতে A ∩ B = \{x : x ∈ A এবং x ∈
B\} ।

\textbf{সংযোগ সেট (Union of Sets ):}

দুই বা ততোধিক সেটের সকল উপাদান নিয়ে গঠিত সেটকে সংযোগ সেট বলা হয়। মনে করি,
A ও B দুইটি সেট। A ও B সেটের সংযোগকে AUB দ্বারা প্রকাশ করা হয় এবং পড়া হয় A
সংযোগ B ও

অথবা A Union B। সেট গঠন পদ্ধতিতে AUB = \{x : x ∈ A অথবা x ∈ B\}।

পূরক সেট , সংযোগ সেট , শক্তি সেট ও নিশ্ছেদ সেট \textbf{পূরক সেট (Complement
of a Set) :}

U সার্বিক সেট এবং A সেটটি U এর উপসেট। A সেটের বহির্ভূত সকল উপাদান নিয়ে গঠিত
সেটকে A সেটের পূরক সেট বলে।

A এর পূরক সেটকে \(A^{c}\) বা A\textquotesingle{} দ্বারা প্রকাশ করা হয়।
গাণিতিকভাবে \(A^{c}\) = U\textbackslash A . মনে করি, P ও Q দুইটি সেট এবং
P সেটের যেসব উপাদান Q সেটের উপাদান নয়, ঐ উপাদানগুলোর সেটকে P এর প্রেক্ষিতে
Q এর পূরক সেট বলা হয় এবং লেখা হয় \(Q^{c}\) = P\textbackslash Q

\textbf{সার্বিক সেট (Universal Set):}

আলোচনায় সংশ্লিষ্ট সকল সেট একটি নির্দিষ্ট সেটের উপসেট। যেমন: A = \{x, y\}
সেটটি B = \{x,y,z\} এর একটি উপসেট। এখানে, B সেটকে A সেটের সাপেক্ষে সার্বিক
সেট বলে।

সুতরাং আলোচনা সংশ্লিষ্ট সকল সেট যদি একটি নির্দিষ্ট সেটের উপসেট হয় তবে ঐ
নির্দিষ্ট সেটকে তার উপসেটগুলোর সাপেক্ষে সার্বিক সেট বলে।

সার্বিক সেটকে সাধারণত U দ্বারা প্রকাশ করা হয়। তবে অন্য প্রতীকের সাহায্যেও
সার্বিক সেট প্রকাশ করা যায়। যেমন: সকল জোড় স্বাভাবিক সংখ্যার সেট C = \{2, 4,
6, ...\} এবং সকল স্বাভাবিক সংখ্যার সেট N = \{1, 2, 3, 4, 5, 6, \} হলে C
সেটের সাপেক্ষে সার্বিক সেট হবে N .

\textbf{শক্তি সেট (Power Sets ) :}

A = \{m,n\} একটি সেট। A সেটের উপসেটসমূহ হলো \{m,n\}, \{m\}, \{n\}, Ø;
এখানে উপসেটসমূহের সেট \{\{m,n\}, \{m\}, \{n\},\} কে A সেটের শক্তি সেট বলা
হয়। A সেটের শক্তি সেটকে P(A) দ্বারা প্রকাশ করা হয়। সুতরাং কোনো সেটের সকল
উপসেট দ্বারা গঠিত সেটকে ঐ সেটের শক্তি সেট বলা হয়।

\textbf{নিশ্ছেদ সেট (Disjoint Set) :}

দুইটি সেটের মধ্যে যদি কোনো সাধারণ উপাদান না থাকে তবে সেট দুইটিকে পরস্পর
নিশ্ছেদ সেট বলে। মনে করি, A ও B দুইটি সেট। A ∩ B = Ø হলে A ও B পরস্পর
নিশ্ছেদ সেট হবে।

সেটের সমতা ও অন্তর \textbf{সেটের সমতা (Equivalent Set) :}

দুইটি সেটের উপাদান একই হলে, সেট দুইটিকে সমান বলা হয়। যেমন: A = \{ 3, 5,
7\} এবং B = \{5,3,3,7\} দুইটি সমান সেট এবং A = B চিহ্ন দ্বারা লেখা হয়। লক্ষ
করি A = B যদি এবং কেবল যদি A ⊆ B এবং B ⊆ A হয়।

আবার, A = \{3, 5, 7\}, B = \{5, 3, 3, 7\} এবং C = \{ 7, 7, 3, 5, 5 \}
হলে A, B ও C সেট তিনটি সমতা বোঝায়। অর্থাৎ, A = B = C

দ্রষ্টব্য : সেটের উপাদানগুলোর ক্রম বদলালে বা কোনো উপাদান পুনরাবৃত্তি করলে সেটের
কোনো পরিবর্তন হয় না।

\textbf{সেটের অন্তর (Difference of Sets )}

মনে করি, A = \{1,2,3,4,5\} এবং B = \{3,5\}। সেট A থেকে সেট B এর
উপাদানগুলো বাদ দিলে যে সেটটি হয় তা \{1, 2, 4\} এবং লেখা হয়
A\textbackslash{} B বা A- B এবং পড়া হয় A বাদ B .

সুতরাং A - B = \{1, 2, 3, 4, 5\} - \{3, 5\} = \{1,2,4\}

কার্তেসীয় গুণজ (Cartesian Product ) : করিম সাহেব তাঁর বাড়ির একটি ঘরের
ভিতরের দেওয়ালে সাদা বা নীল রং এবং বাইরের দেওয়ালে লাল বা হলুদ বা সবুজ রং এর
লেপন দেওয়ার সিদ্ধান্ত নিলেন। ভিতরের দেওয়ালে রং এর সেট A = \{সাদা, নীল\}
এবং বাইরের দেওয়ালে রং এর সেট B = \{লাল, হলুদ ও সবুজ\}। করিম সাহেব তাঁর ঘরের
রং লেপন (সাদা, লাল), (সাদা, হলুদ), (সাদা, সবুজ), (নীল, লাল), (নীল, হলুদ),
(নীল, সবুজ) ক্রমজোড় আকারে দিতে পারেন।

উক্ত ক্রমজোড়ের সেটকে নিচের মতো করে লেখা হয়:

A × B = \{(সাদা, লাল), (সাদা, হলুদ), (সাদা, সবুজ), (নীল, লাল), (নীল, হলুদ),
(নীল, সবুজ)\}

উপরোক্ত ক্রমজোড়ের সেটটিকেই কার্তেসীয় গুণজ সেট বলা হয়।

সেট গঠন পদ্ধতিতে, A x B = \{(x, y) : x ∈ A এবং y ∈ B \}

A × B কে পড়া হয় A ক্রস B।

উপসেট সংখ্যা নির্ণয় \textbf{উপসেট নির্ণয়ের পদ্ধতি (খুবই গুরুত্বপূর্ণ )}

\textbf{সূত্র}: \(2^{n}\) এখানে n হলো উপাদান সংখ্যা। অর্থাৎ যতগুলো উপাদান
থাকে 2 এর উপর বসে । যেমন: একটি সেটের উপাদান

সংখ্যা 2 হলে তার উপসেট হবে \(2^{2}\) = 4, আবার 3 টি উপাদান হলে উপসেট হবে
\(2^{3}\) = 8 টি ।

১. কোন সেটের যতগুলো উপসেট হয় তাদের সেটকে উক্ত সেটের কি বলা হয়? {[} স্বাস্থ্য
অধিদপ্তরের স্বাস্থ্য সহকারী ২০০৪{]}

\textbf{\hl{উত্তর: পাওয়ার সেট}}

\textbf{প্রশ্ন নম্বর-} ২. n উপাদানবিশিষ্ট একটি প্রদত্ত সেটের উপসেটের সংখ্যা কত
হবে? {[}তথ্য মন্ত্রণালয়ের অধীনে টেলিভিশন প্রকৌশলী গ্রেড-২ : ০৪{]} (ক) n(n+1)
(খ) \(n^{2}\) (গ) \(2^{n\  - \ 1}\) (ঘ) \(2^{n}\) \textbf{উত্তরঃ} (ঘ)
\(2^{n}\)\textbf{Explanation:}

n উপাদানবিশিষ্ট একটি প্রদত্ত সেটের উপসেটের সংখ্যা \(2^{n}\) \textbf{প্রশ্ন
নম্বর-} ৩. A = \{2, e\} হলে P(A) কোনটি? {[}মাধ্যমিক বিদ্যালয় সহকারী শিক্ষক :
০০{]} (ক) \{\{2\} , \{e\}\} (খ) \{\{2\}, \{e\}, \{2,e\}\} (গ) \{\{e\},
\{2\}, \{2,e\}, Ø\} (ঘ) \{\} \textbf{উত্তরঃ} (গ) \{\{e\}, \{2\}, \{2,e\},
Ø\} \textbf{Explanation:}

এখানে A সেটের উপাদান সংখ্যা = 2 এবং এর শক্তি সেটের উপাদান সংখ্যা \(2^{2}\)
= 4 . তাই সঠিক উত্তর \{\{e\}, \{2\}, \{2,e\}, Ø\}

\textbf{প্রশ্ন নম্বর-} ৪. কোনো সেটের সদস্য সংখ্যা ৩ হলে, এর উপসেটের সংখ্যা
কতটি হবে? {[}NBR-2015{]} (ক) ২টি (খ) ৮টি (গ) ৩টি (ঘ) ৯টি \textbf{উত্তরঃ}
(খ) ৮টি \textbf{Explanation:}

কোনো সেটের সদস্য সংখ্যা ৩ হলে, এর উপসেটের সংখ্যা \(\text{২}^{\text{৩}} = \text{৮}\ \text{টি}\)
\textbf{প্রশ্ন নম্বর-} ৫. A = \{x : x মৌলিক সংখ্যা এবং x ≤5 \} হলে P(A) এর
সদস্য সংখ্যা কত? {[} ৩৬তম বিসিএস {]} (ক) 8 (খ) 7 (গ) 6 (ঘ) 3
\textbf{উত্তরঃ} (ক) 8 \textbf{Explanation:}

যেহেতু x একটি মৌলিক সংখ্যা এবং x ≤ 5 অর্থাৎ x - এর মান হল 5 এর সমান অথবা 5
এর থেকে ছোট মৌলিক সংখ্যা সুতরাং x = \{2, 3, 5\}

সুতরাং nP(A) = \(2^{3}\) = 8

(যে কোন সেটের সদস্য সংখ্যা 2 এর উপর পাওয়ার হিসেবে লিখলে তার উপসেট সংখ্যা
বের হবে।)

\textbf{প্রশ্ন নম্বর-} ৬. C = \{x:x ঋণাত্মক পূর্ণসংখ্যা এবং \(x^{2}\)
\textless{} 18\}; C সেটের উপাদানগুলো হবে- {[}৩৯ - তম বিসিএস- (বিশেষ){]}
(ক) 1,2,3,4 (খ) 1,3,5,7 (গ) -2,-4,-6,-8 (ঘ) -1,-2,-3,-4 \textbf{উত্তরঃ}
(ঘ) -1,-2,-3,-4 \textbf{Explanation:}

এখানে যেহেতু x একটি ঋণাত্মক পূর্ণ সংখ্যা তাই লেখা যায়,

\({( - 1)}^{2} = 1 < \ 18\)

\({( - 2)}^{2} = 4 < \ 18\)

\({( - 3)}^{2} = 9 < \ 18\)

\({( - 4)}^{2} = 16 < \ 18\)

অর্থাৎ 18 এর থেকে ছোট পূর্ণ বর্গ সংখ্যা আছে 4টি যথা, 16, 9, 4 এবং 1

আবার এদের ঋণাত্মক বর্গমুল গুলো হবে, -4, -3, -2 ও -1 হবে x এর মান।

সুতরাং x = -1,-2,-3,- 4

\textbf{প্রশ্ন নম্বর-} ৭. A = \{x : x ধনাত্মক পূর্ণ সংখ্যা এবং \(x^{2}\)
\textless{} 25 \}

B = \{x : x মৌলিক সংখ্যা এবং \(x^{2}\) \textless{} 25\}

C = \{x : x মৌলিক সংখ্যা এবং \(x^{2}\) = 25\} হলে, A∩B∩ C = ? {[} ৩৭-তম
বিসিএস প্রিলি {]} (ক) \{1,2,3,4\} (খ) \{2,3,4\} (গ) \{2,3,4,5\} (ঘ) Ø
\textbf{উত্তরঃ} (ঘ) Ø \textbf{Explanation:}

এখানে,

A = \{1,2,3,4\} {[}যে সকল পূর্ণ সংখ্যার বর্গ 25 এর থেকে ছোট{]}

B = \{2,3\} {[} যে সকল মৌলিক সংখ্যার বর্গ 25 এর থেকে ছোট{]}

C = \{5\} {[} যে মৌলিক সংখ্যার বর্গ 25 এর সমান {]}

এখন : A ∩ B ∩ C = \{1,2,3,4\} ∩ \{2,3\} ∩ \{5\} = Ø (কারণ তিনটি সেটের
মধ্যে কোন সংখ্যার মিল নেই)

\textbf{প্রশ্ন নম্বর-} ৮. A = \{x : x মৌলিক সংখ্যা এবং x ≤5 \} হলে P(A) এর
সদস্য সংখ্যা কত? {[} ৩৬তম বিসিএস{]} (ক) 8 (খ) 7 (গ) 6 (ঘ) 3 \textbf{উত্তরঃ}
(ক) 8 \textbf{Explanation:}

যেহেতু x একটি মৌলিক সংখ্যা এবং x ≤5

অর্থাৎ x এর মান হল 5 এর সমান অথবা 5 এর থেকে ছোট মৌলিক সংখ্যা

সুতরাং x = 2, 3, 5

সুতরাং A = \{2,3,5\}

সুতরাং nP(A) = \(2^{3}\) = 8 (যে কোন সেটের সদস্য সংখ্যা 2 এর উপর পাওয়ার
হিসেবে লিখলে তার উপসেট সংখ্যা বের হবে।)

\textbf{প্রশ্ন নম্বর-} ৯. A = \{x:x Fibonacci সংখ্যা এবং \(x^{2}\)
\textless{} 64\} হলে, P(A) এর উপাদান কয়টা ? {[}৩৮তম বিসিএস প্রিলিঃ{]} (ক)
128 (খ) 32 (গ) 64 (ঘ) 256 \textbf{উত্তরঃ} (খ) 32 \textbf{Explanation:}

দেয়া আছে, A = \{x :x Fibonacci সংখ্যা এবং \(x^{2}\) \textless{} 64\}

আমরা জানি, Fibonacci সংখ্যা 0, 1, 1, 2, 3, 5, 8, 13, ..

{[}Fibonacci সংখ্যা হল পরপর দুইটি সংখ্যার যোগফল পরবর্তী সংখ্যার সমান {]}

আবার বলা হয়েছে \(x^{2}\) \textless{} 64 অর্থাৎ x এর মান এমন সংখ্যা হবে
যাতে তার বর্গ 64 এর থেকে ছোট হয়।

সুতরাং A = \{0, 1, 1, 2, 3, 5\} যেহেতু x = 8 হলে, \(x^{2}\) \textless64
{[}শর্ত মিলছে না{]}

কিন্তু সেটের মধ্যে একই উপাদান দুবার নেয়া যায়না। তাই

A = \{0, 1, 2, 3, 5\} নিতে হবে।

আমরা জানি, উপসেট বের করার নিয়ম হলো nP(A) = \(2^{n}\) এখানে n = 5

সুতরাং উপ সেট সংখ্যা = \(2^{5}\) = 32

\textbf{প্রশ্ন নম্বর-} ১০. A = \{1, 2, 3, 4\} হলে, P(A) এর উপাদান সংখ্যা কত?
{[} মহিলা ও শিশুবিষয়ক মন্ত্রণালয়ের অধীনে উপজেলা মহিলাবিষয়ক কর্মকর্তা-২০১৬ {]}
(ক) 4 (খ) 8 (গ) 12 (ঘ) 16 \textbf{উত্তরঃ} (ঘ) 16 \textbf{Explanation:}

এখানে , A = \{1, 2, 3, 4\}

সুতরাং P(A) এর উপাদান সংখ্যা হবে = \(2^{4}\) = 16

সেট এর সাধারণ প্রশ্নসমূহ \textbf{নিম্নোক্ত নিয়মগুলি ভালোভাবে আয়ত্ব করলে সেটের
অংকে কোন সমস্যা হবে না ।}

1) দুটি সেটের উপাদানের মধ্যে ∪ (Union) চিহ্ন থাকলে মিল অমিল সব রাশি নিতে
হবে। তবে কোন রাশিকে দুবার নেয়া যাবে না। যেমন: \{1,3,5,6\} ∪ \{ 2,3,4,5\}
= \{1, 2,3,4,5,6\}

2) আবার দুটি সেটের উপাদানের মধ্যে ∩ (Intersection) চিহ্ন থাকলে শুধু মিল আছে
যে রাশিগুলো সেগুলো নিতে হবে, এবং অমিল রাশিগুলোকে বাদ দিতে হবে। যেমন:
\{1,3,5,6\} ∩ \{ 2,3,4,5\} = \{3,5\}

3) আবার দুটি সেটের উপাদানের মধ্যে (-) বিয়োগ চিহ্ন বা (\textbackslash) থাকলে
প্রথম সেটের উপাদান থেকে ২য় সেটের যে উপাদান গুলো মিল হবে তা বাদ দিতে হবে
অর্থাৎ শুধু অমিল রাশি গুলো লিখতে হবে। যেমন: \{1,3,5,6\} - \{ 2,3,4,5\} =
\{1,6\}

কিন্তু যদি একটি রাশিও না মিলে তাহলে প্রথম অংশে যা থাকবে তাই উত্তর: যেমন:
\{1,3,5,\}-\{ 2,4,6\} = \{1,3,5,\}

Note: সেটের বিয়োগফলে দ্বিতীয় সেটের কোনো উপাদান থাকে না ।

4) কোন বর্ণের উপর (\textquotesingle) ড্যাস চিহ্ন থাকলে সার্বিক সেট U এর
উপাদান থেকে ঐ সেটের উপাদনগুলো বিয়োগ করতে হবে। যেমন: A\textquotesingle{}
দেয়া থাকলে প্রথমে লিখতে হবে U- A, তার পর মান বসিয়ে বিয়োগ করতে হবে, অনুরুপ
ভাবে B´ থাকলে U-B.

\textbf{প্রশ্ন নম্বর-} ১১. P = \{x ∈ N : 2 \textless{} x ≤ 6\} এবং Q = \{x
∈ N : x জোড় সংখ্যা এবং x ≤ 8\} হলে P ∩ Q এর মান হবে- {[} শিক্ষা মন্ত্রণালয়ের
অধীনে কারিগরি শিক্ষা অধি: (জুনি: ইন্সট্রাক্টর) - ২০১৮ {]} (ক) \{4, 6\} (খ)
\{2, 3, 5\} (গ) \{2, 4, 6, 8\} (ঘ) \{2, 3, 4, 5, 6, 4\} \textbf{উত্তরঃ}
(ক) \{4, 6\} \textbf{Explanation:}

P = \{3,4,5,6\} {[}২ এর থেকে বড় ৬ পর্যন্ত{]} এবং Q = \{2, 4, 6, 8\} {[}৮
সহ ৮ এর থেকে ছোট জোড় সংখ্যা {]}

.:. P ∩ Q = \{3,4,5,6\} ∩ \{2, 4, 6, 8\} = \{4,6\} {[}শুধু মিল সংখ্যাগুলো{]}

\textbf{প্রশ্ন নম্বর-} ১২. যদি A = \{1, 2, 3\}, B = \{2, 3, 5\} হয়, তবে
A\textbackslash B = কত? {[}জনশক্তি, কর্মসংস্থান ও প্রশিক্ষণ ব্যুরো ইনস্ট্রাক্টর
নিয়োগ-২০১৮{]} (ক) \{2, 3\} (খ) \{5\} (গ) \{1\} (ঘ) \{ 3, 5\}
\textbf{উত্তরঃ} (গ) \{1\} \textbf{Explanation:}

A = \{1,2,3\} এবং B = \{2,3,5\}

সুতরাং A\textbackslash B =\{1,2,3\}\textbackslash\{2,3,5\} = \{1\}
{[}\textbackslash{} চিহ্ন থাকলে (-) বিয়োগ করার মত যেগুলো মিলে সেগুলো বাদ
{]}

\textbf{প্রশ্ন নম্বর-} ১৩. A = \{-1, 1, 2\} এবং B = Ø হলে, A ∪ B এর মান
হবে- {[}NBR-2015{]} (ক) \{−1, 2\} (খ) Ø (গ) \{-1, 1, 2\} (ঘ) \{−1, 0\}
\textbf{উত্তরঃ} (গ) \{-1, 1, 2\} \textbf{Explanation:}

A ∪ B = \{-1, 1, 2\} ∪ \{ Ø \} = \{-1, 1, 2\} {[} মিল অমিল সবগুলো নিতে
হবে। {]}

\textbf{প্রশ্ন নম্বর-} ১৪. ∪ = \{1,2,3,4,5\}, A = \{1, 2, 4\}, B =
\{2,4,5\} হলে, A\textquotesingle B\textquotesingle{} হবে- {[}
প্রাথ:বি:সহ:শি:-০১ {]} (ক) \{1, 2, 3\} (খ) \{2,3,4\} (গ) \{3,4,5\} (ঘ)
\{2,3,4,5\} \textbf{উত্তরঃ} (ঘ) \{2,3,4,5\} \textbf{Explanation:}

A\textquotesingle{} = ∪ - A = \{1,2,3,4,5\} - \{1,2,4\} = \{3,5\}

B′ = ∪ - B = \{1,2,3,4,5\} - \{1,3,5\} = \{2,4\}

সুতরাং A\textquotesingle{} ∪ B′ = \{3,5\} ∪ \{2,4\} = \{2,3,4,5\}

\textbf{প্রশ্ন নম্বর-} ১৫. যদি A ও B যে কোনো দুইটি সেট হয়, তবে (A - B) ∩ B =
কত? {[} সহকারী পরিবার পরিকল্পনা কর্মকর্তা ২০১৬{]} (ক) A (খ) B (গ) Ø (ঘ) C
\textbf{উত্তরঃ} (গ) Ø \textbf{Explanation:}

(A - B) ∩ B = Ø {[} (A-B) ও B দুটি ভিন্ন সেট{]}

\textbf{প্রশ্ন নম্বর-} ১৬. P = \{a, b\}, Q = \{b, c\}, R = \{3, 4\} হলে ( P
∩ Q ∪ R ) এর উপাদান সংখ্যা কত? {[} জনপ্রশাসন মন্ত্রণালয়ের অধীনে
পিএসসি\textquotesingle র সহকারী পরিচালক-২০১৬ {]} (ক) 2 টি (খ) 3 টি (গ) 4
টি (ঘ) 5 টি \textbf{উত্তরঃ} (খ) 3 টি \textbf{Explanation:}

{[} এখানে , (P ∩ Q) = \{a, b\} ∩ \{b, c\} = \{ b \}

সুতরাং ( P ∩ Q ∪ R ) = \{ b \} ∪ \{3, 4\} = \{b , 3, 4\}

\textbf{প্রশ্ন নম্বর-} ১৭. যদি A = \{1, 2\} এবং B = \{2, 5\} হয়, তবে P(A) ∩
P(B) = ? {[}বিভিন্ন মন্ত্রণালয়ের ব্যক্তিগত কর্মকর্তা নিয়োগ-২০১৮{]} (ক) \{1, 5\}
(খ) \{2\} (গ) \{Ø, (2) \} (ঘ) \{ Ø\} \textbf{উত্তরঃ} (গ) \{Ø, (2) \}
\textbf{Explanation:}

এখানে, A = \{1,2\} এবং B = \{2,5\}

সুতরাং P(A) = \{(1), (2), (1,2), Ø \}

সুতরাং P(B) = \{(2), (5), (2,5) , Ø \}

তাহলে P(A) ∩ P (B)

= \{(1), (2), (1,2), Ø \} ∩ \{(2), (5), (2,5), Ø \}

= \{ Ø , (2)\}

\textbf{প্রশ্ন নম্বর-} ১৮. যদি সেট A = \{ 5, 15, 20, 30\} এবং B = \{ 3, 5,
15, 18, 20\} হয় তবে নিচের কোনটি A ∩ B নির্দেশ করবে? {[}33 তম বিসিএস{]}
(ক) \{5, 15, 20\} (খ) \{5, 18, 20\} (গ) \{3, 15, 20\} (ঘ) \{5, 15, 30\}
\textbf{উত্তরঃ} (ক) \{5, 15, 20\} \textbf{Explanation:}

A ∩ B

= \{ 5, 15, 20, 30\} ∩ \{ 3, 5, 15, 18, 20\}

= \{5, 15, 20\} {[} শুধু মিল গুলো নিতে হবে {]}

\textbf{প্রশ্ন নম্বর-} ১৯. যে সকল স্বাভাবিক সংখ্যা দ্বারা 105 এবং 147 কে ভাগ
করলে প্রতিক্ষেত্রে 35 অবিশষ্ট থাকে, তাদের সেট নির্ণয় করুন। {[} সাব রে: ০১ {]}
(ক) \{5, 15, 20\} (খ) \{5, 18, 20\} (গ) \{3, 15, 20\} (ঘ) Ø
\textbf{উত্তরঃ} (ঘ) Ø \textbf{Explanation:}

যে সকল স্বাভাবিক সংখ্যা দ্বারা 105 এবং 147 কে ভাগ করলে প্রতিক্ষেত্রে 35 অবিশষ্ট
থাকে,

সে সংখ্যাটি (ভাজক) 35 অপেক্ষা বড় এবং সংখ্যাটি (105-35) = 70 এবং (147- 35) =
112 এর সাধারণ গুণনীয়ক। (অর্থাৎ ৭০ এবং ১১২ উভয়কেই নিঃশেষে ভাগ করা যেতে হবে)

মনে করি,

35 অপেক্ষা বড় 70 এর গুণনীয়কের সেট = A

এবং 35 অপেক্ষা বড় 112 এর গুণনীয়কের সেট = B

A= \{70\} এবং B= \{56, 112\}

নির্ণেয় সেট = A ∩ B= Ø {[}কারণ ৩৫ থেকে বড় এমন কোন সংখ্যা নেই যা দিয়ে ৭০
এবং ১১২ কে ভাগ করা যায়। {]}

\textbf{প্রশ্ন নম্বর-} ২০. কোন কোন স্বাভাবিক সংখ্যা দ্বারা 346 কে ভাগ করলে
প্রতি ক্ষেত্রে 31 অবশিষ্ট থাকে? {[}22 তম বিসিএস{]} (ক) \{ 35, 45, 63, 105,
315\} (খ) \{ 35, 63, 105, 315\} (গ) \{ 35, 45, 63, 15\} (ঘ) \{ 45, 63,
105, 315\} \textbf{উত্তরঃ} (ক) \{ 35, 45, 63, 105, 315\}
\textbf{Explanation:}

প্রশ্নটিতে বলা হয়েছে কোন কোন স্বাভাবিক সংখ্যা দ্বারা 346 কে ভাগ করতে গেলে
প্রতি ক্ষেত্রে 31 অবশিষ্ট থাকে। তাহলে

বোঝা যাচ্ছে সংখ্যাগুলো 31 এর থেকে বড় হবে, এবং 346-31 = 315 এর গুণনীয়ক । এখন

31 অপেক্ষা বড় 315 এর গুনণীয়ক গুলো হচ্ছে 35, 45, 63, 105, এবং 315,

সুতরাং 31 অপেক্ষা বড় 315 এর উৎপাদক সেট \{ 35, 45, 63, 105, 315\} Ans:

\textbf{প্রশ্ন নম্বর-} ২১. ৫ এর গুণিতকের সেট কোন ধরনের সেট? {[}বিভিন্ন
মন্ত্রণালয়ের ব্যক্তিগত কর্মকর্তা নিয়োগ-২০১৮{]} (ক) সসীম সেট (খ) সার্বিক সেট (গ)
ফাঁকা সেট (ঘ) অসীম সেট \textbf{উত্তরঃ} (ঘ) অসীম সেট \textbf{Explanation:}

৫ এর গুণিতক গুলো হলো ৫, ১০, ১৫, ২০, ২৫ , ...

সুতরাং ৫ এর গুণিতকের সেট = \{ ৫, ১০, ১৫, ২০, ২৫, ..... \}

সুতরাং ইহা একটি অসীম সেট।

\textbf{প্রশ্ন নম্বর-} ২২. ৭ এর গুণিতকের সেট কোন ধরনের সেট? {[} নৌপরিবহন
মন্ত্রণালয় প্রশাসনিক কর্মকর্তা ২০১৩{]} (ক) সসীম সেট (খ) সার্বিক সেট (গ) ফাঁকা
সেট (ঘ) অসীম সেট \textbf{উত্তরঃ} (ঘ) অসীম সেট \textbf{Explanation:}

৭ এর গুণিতক গুলো হলো ৭, ১৪, ২১, ২৮, ৩৫ , ...

সুতরাং ৭ এর গুণিতকের সেট = \{ ৭, ১৪, ২১, ২৮, ৩৫ , ... \}

সুতরাং ইহা একটি অসীম সেট।

\textbf{প্রশ্ন নম্বর-} ২৩. যদি B = \{ x : \(x^{2}\) = 9, 2x = 4\} হয়, তবে B
=? {[}সাব-রেজিস্টার - ২০১৬{]} (ক) \{ \} (খ) \{-3, 3\} (গ) \{3, 2\} (ঘ)
\{-3, 3, 2\} \textbf{উত্তরঃ} (ঘ) \{-3, 3, 2\} \textbf{Explanation:}

\(x^{2}\) = 9 হলে x = +3 অথবা -3 আবার 2x = 4 হলে x = 2

সুতরাং x এর মান বসিয়ে B = \{-3, 3, 2\}

ভেনচিত্রের সাধারণ প্রশ্ন \textbf{প্রশ্ন নম্বর-} ২৪. যদি n(A ∪ B) = 61, n(A) =
30, n(B) = 54 হয় তাহলে n(A ∩ B) এর মান কত? {[}বেসামরিক বিমান ও পর্যটন
মন্ত্রণালয় প্রশাসনিক কর্মকর্তা : ০৫{]} (ক) 22 (খ) 25 (গ) 23 (ঘ) 27
\textbf{উত্তরঃ} (গ) 23 \textbf{Explanation:}

সরাসরি বক্সের সুত্রটি প্রয়োগ করে পাই,

n(A ∪ B) = n(A)+ n(B) - n(A ∩ B)

বা, 61 = 30+54 - n(A ∩ B)

বা, n(A ∩ B) = 84 -61

বা, n(A ∩ B) = 23

\textbf{প্রশ্ন নম্বর-} ২৫. ৫০ জন লোকের মধ্যে ৩৫ জন ইংরেজি, ২৫ জন ইংরেজি ও
বাংলা উভয়ই এবং প্রত্যেকেই দুইটি ভাষার অন্তঃত একটি ভাষার কথা বলতে পারেন।
বাংলায় কতজন কথা বলতে পারেন? {[} ৩৫-তম বিসিএস{]} (ক) 10 (খ) 15 (গ) 40 (ঘ)
30 \textbf{উত্তরঃ} (গ) 40 \textbf{Explanation:}

ধরি,মোট লোক সংখ্যা = n(E ∪ B) = 50 জন

শুধু ইংরেজীতে কথা বলতে পারে, n(E) = 35 জন,

শুধু বাংলায় কথা বলতে পারে n(B) =?,

বাংলা এবং ইংরেজী উভয় ভাষায় বলতে পারে \textquotesingle n(E ∩ B) = 25 জন

এখন সুত্র প্রয়োগ করে পাই,

n(E ∪ B) = n(E) + n(B) - n(E ∩ B) + none (এই সূত্রটি-ই সব সময়, শুধু যার যার
মান বসিয়ে সমাধান করতে হবে)

বা, 50 = 35 + n(B) -- 25+0 {[}সবাই কথা বলতে পারে তাই none = 0 {]}

বা, 50 = 10 + n(B)

বা, 50-10 = n(B)

সুতরাং n(B) = 40

সুতরাং বাংলায় কথা বলতে পারে ৪০ জন।

\textbf{প্রশ্ন নম্বর-} ২৬. ৪০০ জন লোকের একটি দলে ২৬০ জন ইংরেজিতে এবং ১৮০ জন
বাংলায় কথা বলতে পারে। তাহলে কতজন উভয় ভাষায় কথা বলতে পারে? {[}পোষ্ট মাষ্টার
জেনারেল (পূর্বাঞ্চল, চট্টগ্রাম)-এর কার্যালয়ের অধীন পোষ্টাল অপারেটর ২০১৬{]} (ক)
১৮০ (খ) ৪০ (গ) ৮০ (ঘ) ১৪০ \textbf{উত্তরঃ} (খ) ৪০ \textbf{Explanation:}

মোট লোক সংখ্যা = ৪০০

ইংরেজীতে কথা বলতে পারে = ২৬০ জন এবং বাংলাতে কথা বলতে পারে = ১৮০ জন।

সুতরাং হয় ইংরেজী এবং বাংলায় কথা বলতে পারে এমন লোকের সংখ্যা = ২৬০+১৮০ = ৪৪০
জন।

কিন্তু যেহেতু মোট লোক সংখ্যা ৪০০ জন। তাই উভয় ভাষায় কথা বলতে পারা লোকের সংখ্যা
৪৪০-৪০০ = ৪০ জন।

\textbf{প্রশ্ন নম্বর-} ২৭. ৫০ জন ছাত্র-ছাত্রীর মধ্যে ১৮ জন মিউজিক, ২৬ জন আর্ট
এবং ২ জন উভয় বিষয় ভর্তি হয়েছে। কতজন ছাত্র-ছাত্রী কোন বিষয়েই ভর্তি হয় নি।
{[}(BB Cash Officer-11+ Exim Bank Ltd. Officer 2014{]} (ক) 6 (খ) 8 (গ)
16 (ঘ) 24 \textbf{উত্তরঃ} (খ) 8 \textbf{Explanation:}

এ ধরণের প্রশ্নে সব সময় মনে রাখবেন যারা উভয় টিতে থাকে তারা প্রতিটি বিষয়ের
ভেতরেও আছে। তাই লোক যতজন ই থাকুক তারা দুবার গণনা হয়ে যায়।

\textbf{সমাধানটি দেখুন:}

এখানে মোট ছাত্র-ছাত্রী ৫০ জন

তাদের ১৮ জন মিউজিক এবং ২৬ জন আর্ট নিলে মিউজিক + আর্ট = ১৮+২৬ = ৪৪ মনে হলেও
এখানে মোট ছাত্র আছে ৪৪-২ = ৪২ জন । কারণ উভয়

বিষয়ে নিয়েছে ২ জন তাই ২ বিয়োগ করতে হয়।

সুতরাং কোন বিষয়েই ভর্তি হয় নি: ৫০-৪২ = ৮ জন।

\section{8. সেট MCQ প্রশ্ন উত্তর
অনুশীলন}\label{ux9b8ux99f-mcq-ux9aaux9b0ux9b6ux9a8-ux989ux9a4ux9a4ux9b0-ux985ux9a8ux9b6ux9b2ux9a8}

\textbf{P = \{x: x, 12 এর গুণনীয়কসমূহ এবং Q = \{x: x, 3 এর গুণিতক এবং x
\textless{} 12\} হলে, P − Q কত?}

\textbf{\hl{ক) \{1, 2, 4\} ✔}}

খ) \{1, 3, 4\}

গ) \{1, 3, 6\}

ঘ) \{1, 2, 6\}

\textbf{\hl{✔ সঠিক উত্তর: ক) \{1, 2, 4\}}}

\begin{quote}
\textbf{ব্যাখ্যা:} \emph{এখানে, 12 এর গুণনীয়ক সমূহ P = \{1, 2, 3, 4, 6,
12\}}

\emph{আবার, 12 এর থেকে ছোট বা সমান 3 এর গুণিতক Q = \{3, 6, 9, 12\}}

\emph{∴ P − Q = \{1, 2, 3, 4, 6, 12\} − \{3, 6, 9, 12\} = \{1, 2, 4\}}
\end{quote}

\textbf{C = \{x: x ঋণাত্মক পূর্ণসংখ্যা এবং x\^{}2 \textless{} 18\}; C সেটের
উপাদানগুলো হবে}

ক) 1, 2, 3, 5

খ) 1, 3, 5, 7

গ) 2, 4, 6, 8

\textbf{\hl{ঘ) −1, −2, −3, −4 ✔}}

\textbf{\hl{✔ সঠিক উত্তর: ঘ) −1, −2, −3, −4}}

\begin{quote}
\textbf{ব্যাখ্যা:} \emph{এখানে, (−1)\^{}2 \textless{} 18, (−2)\^{}2
\textless{} 18, (−3)\^{}2 \textless{} 18, (−4)\^{}2 \textless{} 18}

\emph{কিন্তু, (−5)\^{}2 \textless{} 18 সত্য নয় (কারণ ২৫ \textless{} ১৮ নয়)}

\emph{∴ C = \{−1, −2, −3, −4\}}
\end{quote}

\textbf{A = \{x \textbar{} x ধনাত্মক পূর্ণ সংখ্যা এবং x\^{}2 \textless{}
25\},}

B = \{x \textbar{} x মৌলিক সংখ্যা এবং x\^{}2 \textless{} 25\},

C = \{x \textbar{} x মৌলিক সংখ্যা এবং x\^{}2 = 25\}

\textbf{হলে, A ∩ B ∩ C = ?}

ক) \{1, 2, 3, 4\}

খ) \{2, 3, 4\}

গ) \{2, 3, 4, 5\}

\textbf{\hl{ঘ) φ ✔}}

\textbf{\hl{✔ সঠিক উত্তর: ঘ) φ}}

\begin{quote}
\textbf{ব্যাখ্যা:} \emph{A = \{x \textbar{} x ধনাত্মক পূর্ণ সংখ্যা এবং x\^{}2
\textless{} 25\}, অর্থাৎ A = \{1, 2, 3, 4\}}

\emph{B = \{x \textbar{} x মৌলিক সংখ্যা এবং x\^{}2 \textless{} 25\},
অর্থাৎ B = \{2, 3\}}

\emph{C = \{x \textbar{} x মৌলিক সংখ্যা এবং x\^{}2 = 25\}, অর্থাৎ C =
\{5\}}

\emph{⇒ A ∩ B ∩ C = \{1, 2, 3, 4\} ∩ \{2, 3\} ∩ \{5\}}

\emph{⇒ = \{2, 3\} ∩ \{5\} = φ}
\end{quote}

\textbf{C = \{3, 4, 5\} এবং D = \{4, 6, 8\} হলে C ∪ D কত?}

ক) \{4\}

\textbf{\hl{খ) \{3, 4, 5, 6, 8\} ✔}}

গ) \{3, 5, 8\}

ঘ) \{0\}

\textbf{\hl{✔ সঠিক উত্তর: খ) \{3, 4, 5, 6, 8\}}}

\begin{quote}
\textbf{ব্যাখ্যা:} \emph{এখানে, C = \{3, 4, 5\} এবং D = \{4, 6, 8\}}

\emph{∴ C ∪ D = \{3, 4, 5\} ∪ \{4, 6, 8\} = \{3, 4, 5, 6, 8\}}
\end{quote}

\textbf{যে সকল স্বাভাবিক সংখ্যা দ্বারা ৩৪৬ এবং ৫৫৬ কে ভাগ করলে প্রতিক্ষেত্রে
৩১ অবশিষ্ট থাকে এদের সেট নির্ণয় করুন।}

\textbf{\hl{ক) \{৩২, ১০৫\} ✔}}

খ) \{৩৪, ৪৫\}

গ) \{৪৫, ৭৫\}

ঘ) \{৩৫, ৭৫\}

\textbf{\hl{✔ সঠিক উত্তর: ক) \{৩২, ১০৫\}}}

\begin{quote}
\textbf{ব্যাখ্যা:} \emph{যেসব স্বাভাবিক সংখ্যা দ্বারা ৩৪৬ এবং ৫৫৬ কে ভাগ করলে
প্রতি ক্ষেত্রে ৩১ অবশিষ্ট থাকে, সে সংখ্যাগুলো হলো (৩৪৬ − ৩১) = ৩১৫ এবং (৫৫৬ −
৩১) = ৫২৫ এর সাধারণ গুণনীয়ক। ৩১ অপেক্ষা বড় ৩১৫ এর গুণনীয়কগুলো হলো ৩৫, ৪৫,
৬৩, ১০৫, ৩১৫ এবং ৩১ অপেক্ষা বড় ৫২৫ এর গুণনীয়কগুলো হলো ৩৫, ৭৫, ১০৫, ১৭৫,
৫২৫}

\emph{⇒ নির্ণেয় সেট \{৩২, ১০৫\}}
\end{quote}

\textbf{A = \{2, 3\} এবং B = \{φ\} হলে, A ∪ B কত?}

ক) \{3, 2\}

\textbf{\hl{খ) \{2, 3, φ\} ✔}}

গ) \{2, 3\}

ঘ) \{2, 2\}

\textbf{\hl{✔ সঠিক উত্তর: খ) \{2, 3, φ\}}}

\begin{quote}
\textbf{ব্যাখ্যা:} \emph{দেওয়া আছে, A = \{2, 3\} এবং B = \{φ\}}

\emph{⇒ A ∪ B = \{2, 3, φ\}}
\end{quote}

\textbf{যদি A = \{x : x, 3 এর গুণিতক \textless{} 15\} হলে নিচের কোনটি
সঠিক?}

ক) A = \{1, 3, 6, 9, 12\}

\textbf{\hl{খ) A = \{3, 6, 9, 12\} ✔}}

গ) A = \{3, 6, 9, 12, 15\}

ঘ) A = \{3, 5\}

\textbf{\hl{✔ সঠিক উত্তর: খ) A = \{3, 6, 9, 12\}}}

\begin{quote}
\textbf{ব্যাখ্যা:} \emph{3 এর গুণিতক যা 15 অপেক্ষা ছোট এমন সংখ্যাগুলো হলো 3,
6, 9, 12.}

\emph{∴ A = \{3, 6, 9, 12\}}
\end{quote}

\textbf{বাস্তব সংখ্যার সেট R গঠিত হয়}

ক) \{1, 2, 3, (3/2) , 5, ...\}

খ) \{±1, ±2, ±3, ...\}

গ) \{−∞, ∞\}

\textbf{\hl{ঘ) \{Q ∪ Q\textquotesingle\} ✔}}

\textbf{\hl{✔ সঠিক উত্তর: ঘ) \{Q ∪ Q\textquotesingle\}}}

\begin{quote}
\textbf{ব্যাখ্যা:} \emph{বাস্তবপূর্ণ সংখ্যা ধনাত্মক বা ঋণাত্মক হতে পারে। কাজেই
বাস্তব পূর্ণসংখ্যার Generalized সেট হতে পারে কার \{Q ∪ Q\textquotesingle\}}
\end{quote}

\textbf{একটি সেট A = \{x ∈ N : x\^{}2 \textgreater{} 15 এবং x\^{}3
\textless{} 255\} দেওয়া আছে। ইহাকে তালিকা পদ্ধতিতে প্রকাশ করুন।}

ক) \{5, 7, 9\}

খ) \{4, 3, 5\}

\textbf{\hl{গ) \{4, 5, 6\} ✔}}

ঘ) \{5, 6, 7\}

\textbf{\hl{✔ সঠিক উত্তর: গ) \{4, 5, 6\}}}

\begin{quote}
\textbf{ব্যাখ্যা:} \emph{ধরি, N = 1, 2, 3, 4, 5, ........}

\emph{এখন, x ∈ N হলে এবং x\^{}2 \textgreater{} 15 হলে x = 4, 5, 6, 7, 8,
9, 10, 11, 12 ...... হতে পারে।}

\emph{আবার, x ∈ N হলে এবং x\^{}3 \textless{} 255 হলে x = 1, 2, 3, 4, 5,
6, ....... হতে পারে।}

\emph{অতএব, x\^{}2 \textgreater{} 15 এবং x\^{}3 \textless{} 255 শর্তে x =
\{4, 5, 6\}}
\end{quote}

\textbf{A = \{x ∈ N : x\^{}2 \textgreater{} 20 এবং x\^{}3 \textless{}
130\} হলে A সেটের উপাদান কোনটি হবে?}

\textbf{\hl{ক) 5 ✔}}

খ) 29

গ) 3

ঘ) 4

\textbf{\hl{✔ সঠিক উত্তর: ক) 5}}

\begin{quote}
\textbf{ব্যাখ্যা:} \emph{ধরি, N = 1, 2, 3, 4, 5, 6, 7.......}

\emph{এখন x ∈ N হলে x\^{}2 \textgreater{} 20 শর্তে x = 5, 6, 7... এবং
x\^{}3 \textless{} 130 শর্তে x = 5, 4, 3, 2, 1....}

\emph{অতএব, x\^{}2 \textgreater{} 20 এবং x\^{}3 \textless{} 130 শর্তে A
সেটের উপাদান হবে 5।}
\end{quote}

\textbf{যদি A = \{x: x\^{}2 = 9, 2x = 4\} হয়, A = কত?}

ক) \{3, 2\}

খ) \{−3, 2\}

\textbf{\hl{গ) φ ✔}}

ঘ) কোনটিই নয়

\textbf{\hl{✔ সঠিক উত্তর: গ) φ}}

\begin{quote}
\textbf{ব্যাখ্যা:} \emph{সেট গঠন পদ্ধতির শর্ত অনুসারে, উপাদান x হলে}

\emph{x\^{}2 = 9 ⇒ x = ±3}

\emph{আবার, 2x = 4 বা, x = 2}

\emph{শর্ত দুটি সাংঘর্ষিক। অর্থাৎ এমন কোনো উপাদান পাওয়া সম্ভব নয় যার বর্গ ৯ এর
সমান এবং দ্বিগুণ ৪ এর সমান।}

\emph{∴ নির্ণেয় উপাদান φ}
\end{quote}

\textbf{8 এর গুণনীয়ক সেট কোনটি?}

ক) \{8, 16, 24\}

\textbf{\hl{খ) \{1, 2, 4, 8\} ✔}}

গ) \{2, 4, 8\}

ঘ) \{1, 2\}

\textbf{\hl{✔ সঠিক উত্তর: খ) \{1, 2, 4, 8\}}}

\begin{quote}
\textbf{ব্যাখ্যা:} \emph{গুণনীয়ক বলতে বোঝায় যা দ্বারা ভাগ যায়। সুতরাং 8 এর
গুণনীয়ক সেট \{1, 2, 4, 8\}।}
\end{quote}

\textbf{A = \{X : X ∈ N এবং X\^{}2 \textless{} 36\}, B = \{X : X একটি
মৌলিক সংখ্যা এবং X\^{}2 \textless{} 36\} এবং C = \{X : X একটি মৌলিক সংখ্যা
এবং X\^{}2 = 4\} হলে A ∩ B ∩ C সমান---}

ক) \{1, 2, 3, 4, 5\}

খ) \{2, 3, 5\}

\textbf{\hl{গ) \{2\} ✔}}

ঘ) φ

\textbf{\hl{✔ সঠিক উত্তর: গ) \{2\}}}

\begin{quote}
\textbf{ব্যাখ্যা:} \emph{A = \{x : x ∈ N এবং x\^{}2 \textless{} 36\}}

\emph{⇒ A = \{1, 2, 3, 4, 5\}}

\emph{B = \{x : x একটি মৌলিক সংখ্যা এবং x\^{}2 \textless{} 36\}}

\emph{⇒ B = \{2, 3, 5\}}

\emph{C = \{x : x একটি মৌলিক সংখ্যা এবং x\^{}2 = 4\}}

\emph{⇒ C = \{2\}}

\emph{⇒ A ∩ B ∩ C = \{1, 2, 3, 4, 5\} ∩ \{2, 3, 5\} ∩ \{2\} = \{2\}}
\end{quote}

\textbf{সার্বিক সেট U = \{1, 2, 3, 4, 5\}, A = \{1, 2, 4\}, B = \{1, 3,
5\} হলে A\textquotesingle{} ∪ B\textquotesingle{} কত হবে?}

ক) \{1, 2, 3\}

খ) \{2, 3, 4\}

গ) \{3, 4, 5\}

\textbf{\hl{ঘ) \{2, 3, 4, 5\} ✔}}

\textbf{\hl{✔ সঠিক উত্তর: ঘ) \{2, 3, 4, 5\}}}

\begin{quote}
\textbf{ব্যাখ্যা:} \emph{A\textquotesingle{} = U − A = \{1, 2, 3, 4, 5\} −
\{1, 2, 4\} = \{3, 5\}}

\emph{B\textquotesingle{} = U − B = \{1, 2, 3, 4, 5\} − \{1, 3, 5\} =
\{2, 4\}}

\emph{∴ A\textquotesingle{} ∪ B\textquotesingle{} = \{3, 5\} ∪ \{2, 4\}
= \{2, 3, 4, 5\}}
\end{quote}

\textbf{যদি A ও B যে কোনো দুইটি সেট হয়, তবে (A − B) ∩ B = কত?}

ক) A

খ) B

\textbf{\hl{গ) φ ✔}}

ঘ) C

\textbf{\hl{✔ সঠিক উত্তর: গ) φ}}

\begin{quote}
\textbf{ব্যাখ্যা:} \emph{(A − B) ∩ B = A ∩ B ∩ B\textquotesingle{} = φ
(কারণ A ও B দুটি ভিন্ন সেট বা পূরক সেটের সাথে সাধারণ উপাদান থাকে না)।}
\end{quote}

\textbf{দুইটি সেট যদি নিম্নলিখিতভাবে প্রকাশ করা হয়--- X = \{x : x\^{}3 −
3x\^{}2 + 3x − 1 = 0\} এবং Y = \{x : x ∈ R, 3 দ্বারা বিভাজ্য\} তাহলে X ∩ Y
= কত?}

ক) \{1\}

খ) 0

\textbf{\hl{গ) φ ✔}}

ঘ) \{1, 3\}

\textbf{\hl{✔ সঠিক উত্তর: গ) φ}}

\begin{quote}
\textbf{ব্যাখ্যা:} \emph{X = \{x : x\^{}3 − 3x\^{}2 + 3x − 1 = 0\}}

\emph{এখন, x\^{}3 − 3x\^{}2 + 3x − 1 = 0}

\emph{⇒ x\^{}3 − 3 × x\^{}2 × 1 + 3 × x × 1\^{}2 − (1)\^{}3 = 0}

\emph{⇒ (x − 1)\^{}3 = 0}

\emph{⇒ x − 1 = 0}

\emph{∴ x = \{1\}}

\emph{Y = \{x : x ∈ R, 3 দ্বারা বিভাজ্য\}}

\emph{∴ Y = \{±3, ±6, ±9, ........\}}

\emph{∴ X ∩ Y = \{1\} ∩ \{±3, ±6, ±9, ........\} = φ}
\end{quote}

\textbf{R = \{x : x, 4 এর গুণিতক এবং x \textless{} 16\} সেটের তালিকায়
প্রকাশিত রূপ কোনটি?}

ক) \{1, 2, 4\}

\textbf{\hl{খ) \{4, 8, 12\} ✔}}

গ) \{0, 4, 8, 12, 16\}

ঘ) \{4, 8, 12, 16\}

\textbf{\hl{✔ সঠিক উত্তর: খ) \{4, 8, 12\}}}

\begin{quote}
\textbf{ব্যাখ্যা:} \emph{R = \{4, 8, 12\}}
\end{quote}

\textbf{A = \{−1, 1, 2\} এবং B = φ হলে, A ∩ B এর মান হবে ---}

ক) \{−1, 1, 2\}

খ) \{−1, 1, 2\}

\textbf{\hl{গ) φ ✔}}

ঘ) \{−1, φ\}

\textbf{\hl{✔ সঠিক উত্তর: গ) φ}}

\begin{quote}
\textbf{ব্যাখ্যা:} \emph{A ∩ B = \{−1, 1, 2\} ∩ φ = φ}
\end{quote}

\textbf{A = \{x : x জোড় মৌলিক সংখ্যা\}, তালিকা পদ্ধতিতে কি হবে?}

ক) \{2\}

খ) 2\^{}−19

গ) (2)

\textbf{\hl{ঘ) \{2\} ✔}}

\textbf{\hl{✔ সঠিক উত্তর: ঘ) \{2\}}}

\begin{quote}
\textbf{ব্যাখ্যা:} \emph{তালিকা পদ্ধতিতে সেটের সকল উপাদানকে \{ \} এর মধ্যে
আবদ্ধ করা হয়।}

\emph{এখানে 2 একমাত্র জোড় মৌলিক সংখ্যা সমস্ত মৌলিক সংখ্যার মধ্যে।}

\emph{∴ \{2\} উত্তর।}
\end{quote}

\textbf{A = \{x : x মৌলিক সংখ্যা এবং x \textless{} 10\} হলে নিচের কোনটি
সঠিক?}

ক) A = \{1, 2, 3, 5, 7\}

খ) A = \{1, 2, 3, 6, 7\}

\textbf{\hl{গ) A = \{2, 3, 5, 7\} ✔}}

ঘ) A = \{4, 6, 8, 9\}

\textbf{\hl{✔ সঠিক উত্তর: গ) A = \{2, 3, 5, 7\}}}

\begin{quote}
\textbf{ব্যাখ্যা:} \emph{১০ এর ছোট মৌলিক সংখ্যাগুলো হলো ২, ৩, ৫ এবং ৭।}

\emph{∴ A = \{2, 3, 5, 7\}}
\end{quote}

\textbf{নিচের কোনটি দ্য মরগানের সূত্র?}

ক) A ∪ B = B ∪ A

খ) (A ∪ B) ∪ C = A ∪ (B ∪ C)

\textbf{\hl{গ) (A ∩ B)\textquotesingle{} = A\textquotesingle{} ∪
B\textquotesingle{} ✔}}

ঘ) A ∩ B\textquotesingle{} = A − B

\textbf{\hl{✔ সঠিক উত্তর: গ) (A ∩ B)\textquotesingle{} =
A\textquotesingle{} ∪ B\textquotesingle{}}}

\begin{quote}
\textbf{ব্যাখ্যা:} \emph{সার্বিক সেট U এর যে কোনো উপসেট A ও B এর জন্য (A ∩
B)\textquotesingle{} = A\textquotesingle{} ∪ B\textquotesingle{} এবং (A
∪ B)\textquotesingle{} = A\textquotesingle{} ∩ B\textquotesingle{} কে দ্য
মরগানের সূত্র (De Morgan\textquotesingle s Law) বলা হয়।}
\end{quote}

\textbf{7 এর গুণিতকের সেট কোন ধরনের সেট?}

ক) সসীম সেট

খ) ফাঁকা সেট

গ) সার্বিক সেট

\textbf{\hl{ঘ) অসীম সেট ✔}}

\textbf{\hl{✔ সঠিক উত্তর: ঘ) অসীম সেট}}

\begin{quote}
\textbf{ব্যাখ্যা:} \emph{যে সেটের উপাদানের সংখ্যা গণনা করে নির্ধারণ করা যায়
না, সে সেটকে অসীম সেট বা অনন্ত সেট বলা হয়। সকল স্বাভাবিক সংখ্যার সেট N =
\{7, 14, 21, .........∞\}}
\end{quote}

\textbf{A = \{x : 7 ≤ x \textless{} 11\} এবং B = \{x : x মৌলিক সংখ্যা
\textless{} 15\} হলে A ∩ B কোনটি?}

ক) \{7, 9\}

খ) \{9\}

\textbf{\hl{গ) \{7\} ✔}}

ঘ) φ

\textbf{\hl{✔ সঠিক উত্তর: গ) \{7\}}}

\begin{quote}
\textbf{ব্যাখ্যা:} \emph{A = \{x : 7 ≤ x \textless{} 11\}}

\emph{⇒ A = \{7, 8, 9, 10\}}

\emph{B = \{x : x মৌলিক সংখ্যা \textless{} 15\}}

\emph{⇒ B = \{2, 3, 5, 7, 11, 13\}}

\emph{∴ A ∩ B = \{7, 8, 9, 10\} ∩ \{2, 3, 5, 7, 11, 13\} = \{7\}}
\end{quote}

\textbf{স্বাভাবিক সংখ্যার সেটের ক্ষুদ্রতম সদস্য হলো---}

ক) ২

\textbf{\hl{খ) ১ ✔}}

গ) ৩

ঘ) ০

\textbf{\hl{✔ সঠিক উত্তর: খ) ১}}

\begin{quote}
\textbf{ব্যাখ্যা:} \emph{স্বাভাবিক সংখ্যার সেট N = \{1, 2, 3, 4, 5,
.........\}, যার ক্ষুদ্রতম সদস্য হলো ১।}
\end{quote}

\textbf{5 উপাদান বিশিষ্ট একটি প্রদত্ত সেটের উপসেটের সংখ্যা হবে...}

ক) ২০

খ) ২৫

গ) ২৮

\textbf{\hl{ঘ) ৩২ ✔}}

\textbf{\hl{✔ সঠিক উত্তর: ঘ) ৩২}}

\begin{quote}
\textbf{ব্যাখ্যা:} \emph{5 উপাদান বিশিষ্ট একটি সেটের উপসেটের সংখ্যা, 2\^{}n
= 2\^{}5 = 32।}
\end{quote}

\textbf{যদি A এবং B দুটি সেট হয়, তবে i. (A ∪ B)\^{}c = A\^{}c ∩ B\^{}c
ii. (A ∩ B)\^{}c = A\^{}c ∪ B\^{}c উপরিউক্ত সূত্রটি কি নামে পরিচিত?}

ক) গাউসের সূত্রাবলী

খ) ইকিউএর সূত্রাবলী

গ) আইলারের সূত্রাবলী

\textbf{\hl{ঘ) ডি মরগান সূত্রাবলী ✔}}

\textbf{\hl{✔ সঠিক উত্তর: ঘ) ডি মরগান সূত্রাবলী}}

\begin{quote}
\textbf{ব্যাখ্যা:} \emph{সেটের এই সূত্র দুটি ডি মরগান সূত্রাবলী (De
Morgan\textquotesingle s Laws) নামে পরিচিত।}
\end{quote}

\textbf{যদি A এবং B দুইটি সেট হয় তবে A ∩ (A ∪ B) = A ∪ (A ∩ B) = কত?}

ক) A ∪ B

খ) A ∩ B

\textbf{\hl{গ) A ✔}}

ঘ) B

\textbf{\hl{✔ সঠিক উত্তর: গ) A}}

\begin{quote}
\textbf{ব্যাখ্যা:} \emph{ধরি, A = \{1, 2, 3\} এবং B = \{2, 3, 4\}}

\emph{A ∩ (A ∪ B) = \{1, 2, 3\} ∩ (\{1, 2, 3\} ∪ \{2, 3, 4\})}

\emph{= \{1, 2, 3\} ∩ \{1, 2, 3, 4\} = \{1, 2, 3\} = A}

\emph{আবার, A ∪ (A ∩ B) = \{1, 2, 3\} ∪ (\{1, 2, 3\} ∩ \{2, 3, 4\})}

\emph{= \{1, 2, 3\} ∪ \{2, 3\} = \{1, 2, 3\} = A}

\emph{∴ A ∩ (A ∪ B) = A ∪ (A ∩ B) = A}
\end{quote}

\textbf{A ∪ (A ∩ B) = কত?}

ক) φ

\textbf{\hl{খ) A ✔}}

গ) B

ঘ) A ∩ B

\textbf{\hl{✔ সঠিক উত্তর: খ) A}}

\begin{quote}
\textbf{ব্যাখ্যা:} \emph{ধরি, A = \{1, 2, 3\} , B = \{2, 3, 4\}}

\emph{∴ A ∩ B = \{1, 2, 3\} ∩ \{2, 3, 4\} = \{2, 3\}}

\emph{∴ A ∪ (A ∩ B) = \{1, 2, 3\} ∪ \{2, 3\} = \{1, 2, 3\} = A}
\end{quote}

\textbf{যদি A = \{2, 3\}, B = \{1, 2\} এবং A ও B এর উপাদানগুলোর মধ্যে x
\textgreater{} y সম্পর্কটি বিবেচনা থাকে, তবে অন্বয়টি---}

ক) \{(2, 1) (2, 2) (3, 1)\}

\textbf{\hl{খ) \{(2, 1) (3, 1) (3, 2)\} ✔}}

গ) \{(1, 2) (3, 1) (3, 2)\}

ঘ) \{(1, 2) (1, 3) 2, 3\}

\textbf{\hl{✔ সঠিক উত্তর: খ) \{(2, 1) (3, 1) (3, 2)\}}}

\begin{quote}
\textbf{ব্যাখ্যা:} \emph{এখানে, A = \{2, 3\} এবং B = \{1, 2\}}

\emph{A × B = \{(2, 1), (2, 2), (3, 1), (3, 2)\}}

\emph{x \textgreater{} y শর্ত সাপেক্ষে অন্বয়টি হবে R = \{(2, 1), (3, 1),
(3, 2)\}}
\end{quote}

\textbf{যদি A = \{1, 2, 3\}, B = \{3, 4\} এবং C = \{4, 5, 6\} হয় তবে (A
× B) ∩ (B × C) = কত হবে?}

ক) \{(2, 3), (3, 4), (4, 5)\}

খ) \{(1, 2), (2, 3), (3, 4)\}

\textbf{\hl{গ) \{(3, 4), (3, 5), (3, 6)\} ✔}}

ঘ) \{(1, 4), (2, 4), (3, 4)\}

\textbf{\hl{✔ সঠিক উত্তর: গ) \{(3, 4), (3, 5), (3, 6)\}}}

\begin{quote}
\textbf{ব্যাখ্যা:} \emph{এখানে সঠিক উত্তর (3, 4) {[}বিকল্পে ব্র্যাকেট সেট
আকারে থাকলে \{(3, 4)\} হবে{]}।}

\emph{A = \{1, 2, 3\}, B = \{3, 4\}, C = \{4, 5, 6\}}

\emph{⇒ (A × B) = \{(1, 3), (1, 4), (2, 3), (2, 4), (3, 3), (3, 4)\}}

\emph{⇒ (B × C) = \{(3, 4), (3, 5), (3, 6), (4, 4), (4, 5), (4, 6)\}}

\emph{∴ (A × B) ∩ (B × C) = \{(3, 4)\}}
\end{quote}

\textbf{যদি A = \{1, 2, 3\} এবং R = \{(x, y) \textbar{} x \textless{} y
এবং x, y ∈ A\} হয়, তবে R এর মান ডোমেন ও রেঞ্জ বের করুন।}

ক) ডোমেন (R) = A, রেঞ্জ (R) = A

খ) ডোমেন (R) = \{1, 3, 4\}, রেঞ্জ (R) = \{1, 2, 3\}

গ) ডোমেন (R) = \{2, 3, 4\}, রেঞ্জ (R) = \{1, 3, 4\}

\textbf{\hl{ঘ) ডোমেন (R) = \{1, 2, 3\}, রেঞ্জ (R) = \{2, 3, 4\} ✔}}

\textbf{\hl{✔ সঠিক উত্তর: ঘ) ডোমেন (R) = \{1, 2, 3\}, রেঞ্জ (R) = \{2, 3,
4\}}}

\begin{quote}
\textbf{ব্যাখ্যা:} \emph{দেওয়া আছে, A = \{1, 2, 3\} , R = \{(x, y)
\textbar{} x \textless{} y এবং x, y ∈ A\}}

\emph{শর্তানুসারে, R = \{(1, 2), (1, 3), (2, 3)\}}

\emph{∴ ডোমেন (R) = \{1, 2, 3\}}

\emph{∴ রেঞ্জ (R) = \{2, 3, 4 (ভুল মুদ্রণ, সঠিক হবে \{2, 3\})\}}
\end{quote}

\textbf{U = \{1, 2, 3, 4, 5, 6\}, A = \{1, 2, 3\}, B = \{2, 4, 6\} হলে
A\textquotesingle{} ∩ B\textquotesingle{} হবে}

ক) \{1, 4\}

খ) \{5, 6\}

\textbf{\hl{গ) \{5\} ✔}}

ঘ) \{6\}

\textbf{\hl{✔ সঠিক উত্তর: গ) \{5\}}}

\begin{quote}
\textbf{ব্যাখ্যা:} \emph{A\textquotesingle{} = U − A = \{4, 5, 6\}}

\emph{B\textquotesingle{} = U − B = \{1, 3, 5\}}

\emph{∴ A\textquotesingle{} ∩ B\textquotesingle{} = \{4, 5, 6\} ∩ \{1,
3, 5\} = \{5\}}
\end{quote}

\textbf{A = \{0, 2, 3, 7, 9\}, B = \{1, 5, 6, 8, 11\}, C = \{2, 5, 7, 8,
12, 14\} হলে (A ∩ B) ∪ (A ∩ C) হবে}

ক) \{2, 3, 5, 7\}

খ) \{4, 7, 9\}

\textbf{\hl{গ) \{2, 7\} ✔}}

ঘ) \{0, 2, 3, 8\}

\textbf{\hl{✔ সঠিক উত্তর: গ) \{2, 7\}}}

\begin{quote}
\textbf{ব্যাখ্যা:} \emph{(A ∩ B) ∪ (A ∩ C)}

\emph{= {[}\{0, 2, 3, 7, 9\} ∩ \{1, 5, 6, 8, 11\}{]} ∪ {[}\{0, 2, 3, 7,
9\} ∩ \{2, 5, 7, 8, 12, 14\}{]}}

\emph{= \{φ\} ∪ \{2, 7\} = \{2, 7\}}
\end{quote}

\textbf{যদি n(A ∪ B) = 61, n(A) = 30, n(B) = 54 হয় তাহলে n(A ∩ B) এর মান
কত?}

ক) 22

খ) 25

\textbf{\hl{গ) 23 ✔}}

ঘ) 27

\textbf{\hl{✔ সঠিক উত্তর: গ) 23}}

\begin{quote}
\textbf{ব্যাখ্যা:} \emph{n(A ∪ B) = n(A) + n(B) − n(A ∩ B)}

\emph{⇒ 61 = 30 + 54 − n(A ∩ B)}

\emph{∴ n(A ∩ B) = 84 − 61 = 23}
\end{quote}

\textbf{৪০০ জন লোকের একটি দলে ২৬০ জন ইংরেজিতে এবং ১৮০ জন বাংলায় কথা বলতে
পারে। তাহলে কতজন উভয় ভাষায় কথা বলতে পারে?}

ক) ১৮০ জন

খ) ৪০ জন

\textbf{\hl{গ) ৮০ জন ✔}}

ঘ) ১৪০ জন

\textbf{\hl{✔ সঠিক উত্তর: গ) ৮০ জন}}

\begin{quote}
\textbf{ব্যাখ্যা:} \emph{n(E ∪ B) = n(E) + n(B) − n(E ∩ B)}

\emph{⇒ 400 = 260 + 180 − n(E ∩ B)}

\emph{⇒ n(E ∩ B) = 440 − 400 = 80}
\end{quote}

\textbf{স্কুলের কোন ক্লাসের ৩২ জন ছাত্রের ১৮ জন ফুটবল খেলে, ১৬ জন ক্রিকেট খেলে
এবং ৭ জন দুটিই খেলে। কতজন কোনটিও খেলে না?}

ক) ৮ জন

খ) ১০ জন

গ) ৭ জন

\textbf{\hl{ঘ) ৫ জন ✔}}

\textbf{\hl{✔ সঠিক উত্তর: ঘ) ৫ জন}}

\begin{quote}
\textbf{ব্যাখ্যা:} \emph{ধরি, কোনটিও খেলে না x জন।}

\emph{Total = n(F) + n(C) − n(F ∩ C) + x}

\emph{⇒ 32 = 18 + 16 − 7 + x}

\emph{⇒ 32 = 27 + x}

\emph{∴ x = 5}
\end{quote}

\textbf{P = \{a, b\}, Q = \{b, c\}, R = \{3, 4\} হলে (P ∩ Q) ∪ R এর
উপাদান সংখ্যা কত?}

ক) 2

\textbf{\hl{খ) 3 ✔}}

গ) 4

ঘ) 5

\textbf{\hl{✔ সঠিক উত্তর: খ) 3}}

\begin{quote}
\textbf{ব্যাখ্যা:} \emph{(P ∩ Q) = \{a, b\} ∩ \{b, c\} = \{b\}}

\emph{⇒ (P ∩ Q) ∪ R = \{b\} ∪ \{3, 4\} = \{b, 3, 4\}}

\emph{∴ উপাদান সংখ্যা 3}
\end{quote}

\textbf{যদি B = \{x : x² = 9, 2x = 4\} হয়, তবে B = ?}

ক) \{\}

খ) \{-3, 3\}

গ) \{3, 2\}

\textbf{\hl{ঘ) \{-3, 3, 2\} ✔}}

\textbf{\hl{✔ সঠিক উত্তর: ঘ) \{-3, 3, 2\}}}

\begin{quote}
\textbf{ব্যাখ্যা:} \emph{এখানে, x² = 9}

\emph{∴ x = ±3}

\emph{2x = 4}

\emph{∴ x = 2}

\emph{B = \{-3, 3, 2\}}
\end{quote}

\textbf{x² − (p + q) x + pq = 0−এর সমাধান সেট কোনটি?}

\textbf{\hl{ক) \{p, q\} ✔}}

খ) \{p, −q\}

গ) \{-p, −q\}

ঘ) \{≠p, q\}

\textbf{\hl{✔ সঠিক উত্তর: ক) \{p, q\}}}

\begin{quote}
\textbf{ব্যাখ্যা:} \emph{x² − (p + q) x + pq = 0}

\emph{⇒ x² − px − qx + pq = 0}

\emph{⇒ x(x − p) − q(x − p) = 0}

\emph{⇒ (x − p)(x − q) = 0}

\emph{∴ x = p, q}

\emph{∴ সমাধান সেট = \{p, q\}}
\end{quote}

\textbf{কোন কোন স্বাভাবিক সংখ্যা দ্বারা ৩৪৮ কে ভাগ করলে প্রতিক্ষেত্রে ৩১
অবশিষ্ট থাকে?}

\textbf{\hl{ক) ৩৫, ৪৫, ৬৩, ১০৫, ৩১৫ ✔}}

খ) ৩৫, ৪০, ৬৪, ১১০, ৩১৫

গ) ৩৫, ৪৫, ৭০, ১০৫, ৩১৫

ঘ) ৩৫, ৪৫, ৬৩, ১১০, ৩১৫

\textbf{\hl{✔ সঠিক উত্তর: ক) ৩৫, ৪৫, ৬৩, ১০৫, ৩১৫}}

\begin{quote}
\textbf{ব্যাখ্যা:} \emph{যে সকল স্বাভাবিক সংখ্যা দ্বারা 346 কে ভাগ করলে
প্রতিক্ষেত্রে 31 অবশিষ্ট থাকে, সে সকল সংখ্যা 31 অপেক্ষা বড় এবং সংখ্যাগুলো হবে
(346 − 31) = 315 এর গুণনীয়ক।}

\emph{∴ নির্ণেয় গুণনীয়কের সেট = \{35, 45, 63, 105, 315\}}
\end{quote}

\textbf{একটি সেটের উপাদান সংখ্যা 4 হলে তার পাওয়ার সেটের উপাদান সংখ্যা কত
হবে?}

ক) 4

খ) 20

\textbf{\hl{গ) 16 ✔}}

ঘ) 24

\textbf{\hl{✔ সঠিক উত্তর: গ) 16}}

\begin{quote}
\textbf{ব্যাখ্যা:} \emph{সেটের উপাদান সংখ্যা n = 4}

\emph{∴ পাওয়ার সেটের উপাদান সংখ্যা 2ⁿ = 2⁴ = 16}
\end{quote}

\textbf{A = \{1, 2, 3, 4\} হলে, P(A) এর উপাদান সংখ্যা কত?}

ক) 4

খ) 8

গ) 12

\textbf{\hl{ঘ) 16 ✔}}

\textbf{\hl{✔ সঠিক উত্তর: ঘ) 16}}

\begin{quote}
\textbf{ব্যাখ্যা:} \emph{এখানে, A এর উপাদান সংখ্যা n = 4}

\emph{∴ P(A) এর উপাদান সংখ্যা 2ⁿ = 2⁴ = 16}
\end{quote}

\textbf{সেট A = \{x : x, Fibonacci number এবং x² \textless{} 64\} হলে,
P(A) এর উপাদান কয়টি?}

ক) 128

\textbf{\hl{খ) 32 ✔}}

গ) 64

ঘ) 256

\textbf{\hl{✔ সঠিক উত্তর: খ) 32}}

\begin{quote}
\textbf{ব্যাখ্যা:} \emph{ফাইোবোনাচ্চি সিরিজটি হলো: 0, 1, 1, 2, 3, 5, 8, 13,
...}

\emph{যেহেতু Fibonacci number-গুলোর বর্গ 64 এর থেকে ছোট হতে হবে, তাই
সংখ্যাগুলো হবে: 0, 1, 2, 3, 5}

\emph{⇒ A = \{0, 1, 2, 3, 5\}}

\emph{⇒ A সেটের উপাদান সংখ্যা 5}

\emph{∴ P(A) এর সদস্য সংখ্যা হবে 2ⁱ = 32}
\end{quote}

\textbf{A = \{x : x মৌলিক সংখ্যা এবং x ≤ 5\} হলে P(A) এর সদস্য সংখ্যা কত?}

\textbf{\hl{ক) 8 ✔}}

খ) 7

গ) 6

ঘ) 3

\textbf{\hl{✔ সঠিক উত্তর: ক) 8}}

\begin{quote}
\textbf{ব্যাখ্যা:} \emph{A = \{2, 3, 5\}}

\emph{⇒ A সেটের উপাদান সংখ্যা 3}

\emph{∴ P(A) এর উপাদান সংখ্যা = 2³ = 8}
\end{quote}

\textbf{A = \{1, 2, 3, 4\} হলে, P(A) এর উপাদান সংখ্যা কত?}

\textbf{\hl{ক) 16 ✔}}

খ) 10

গ) 12

ঘ) 18

\textbf{\hl{✔ সঠিক উত্তর: ক) 16}}

\begin{quote}
\textbf{ব্যাখ্যা:} \emph{এখানে, A সেটের উপাদান সংখ্যা n = 4}

\emph{∴ P(A) এর উপাদান সংখ্যা 2ⁿ = 2⁴ = 16}
\end{quote}

\textbf{A = \{a, b, c\} সেটের উপসেট কয়টি?}

ক) 1টি

খ) 2টি

গ) 4টি

\textbf{\hl{ঘ) 8টি ✔}}

\textbf{\hl{✔ সঠিক উত্তর: ঘ) 8টি}}

\begin{quote}
\textbf{ব্যাখ্যা:} \emph{এখানে A সেটের উপাদান সংখ্যা n = 3}

\emph{∴ A এর উপসেট সংখ্যা = 2³ টি = 8 টি}
\end{quote}

\textbf{কোনো সেটের সদস্য সংখ্যা n হলে, এর উপসেটের সংখ্যা কতটি হবে?}

ক) 2টি

খ) 3টি

\textbf{\hl{গ) 8টি ✔}}

ঘ) 9টি

\textbf{\hl{✔ সঠিক উত্তর: গ) 8টি}}

\begin{quote}
\textbf{ব্যাখ্যা:} \emph{কোনো সেটের সদস্য সংখ্যা n হলে উপসেটের সংখ্যা নির্ণয়ের
সূত্রটি হলো 2\^{}n । যদি n = 3 হয়, তবে উপসেটের সংখ্যা = 2\^{}3 = 8 টি।}
\end{quote}

\textbf{n উপাদানবিশিষ্ট একটি প্রদত্ত সেটের উপসেটের সংখ্যা কত হবে?}

ক) n(n+1)

খ) n−1

গ) 2 \^{}n−1

\textbf{\hl{ঘ) 2 \^{}n ✔}}

\textbf{\hl{✔ সঠিক উত্তর: ঘ) 2 \^{}n}}

\begin{quote}
\textbf{ব্যাখ্যা:} \emph{কোনো সেটের উপাদান সংখ্যা n হলে, তার উপসেটের সংখ্যা
2\^{}n হয়।}
\end{quote}

\textbf{A = \{2, e\} হলে P(A) কোনটি? / If A = \{2, e\} , then determine
P(A)?}

ক) \{\{2\}, \{e\}\}

খ) \{\{2\}, \{e\}, \{2, e\}\}

গ) \{\}

\textbf{\hl{ঘ) \{\{e\}, \{2\}, \{2, e\}, ∅\} ✔}}

\textbf{\hl{✔ সঠিক উত্তর: ঘ) \{\{e\}, \{2\}, \{2, e\}, ∅\}}}

\begin{quote}
\textbf{ব্যাখ্যা:} \emph{কোনো সেটের সকল উপসেট নিয়ে গঠিত সেটকে ঐ সেটের
পাওয়ার সেট বা শক্তি সেট বলে এবং একে P(A) দ্বারা প্রকাশ করা হয়।}

\emph{∴ P(A) = \{\{2\}, \{e\}, \{2, e\}, ∅\}}
\end{quote}

\textbf{নিচের কোনটি সত্য?}

ক) A\textbackslash B ⊂ A∪B\textquotesingle{}

\textbf{\hl{খ) A−B = A∩B\textquotesingle{} ✔}}

গ) A\textquotesingle\textbackslash B\textquotesingle{} =
B\textbackslash A\textquotesingle{}

ঘ) (A\textbackslash B) ∩ B−B

\textbf{\hl{✔ সঠিক উত্তর: খ) A−B = A∩B\textquotesingle{}}}

\begin{quote}
\textbf{ব্যাখ্যা:} \emph{সেট তত্ত্বের নিয়ম অনুযায়ী, দুটি সেটের অন্তর
(Difference) A−B হলো সেটের এমন উপাদানগুলোর সেট যা A তে আছে কিন্তু B তে নেই।
এটি পূরক সেটের সূত্রানুসারে A∩B\textquotesingle{} এর সমান।}
\end{quote}

\textbf{p(A∩B) = (1/3) , P(A∪B) = (5/6) , P(A) = (1/2) , P(B) = কত?}

ক) (1/3)

\textbf{\hl{খ) (2/3) ✔}}

গ) (1/4)

ঘ) (2/5)

\textbf{\hl{✔ সঠিক উত্তর: খ) (2/3)}}

\begin{quote}
\textbf{ব্যাখ্যা:} \emph{আমরা জানি, P(A∪B) = P(A) + P(B) − P(A∩B)}

\emph{∴ (5/6) = (1/2) + P(B) − (1/3)}

\emph{∴ P(B) = (5/6) − (1/2) + (1/3)}

\emph{∴ P(B) = (5 \textless− 3 + 2/6) = (4/6) = (2/3)}
\end{quote}

\textbf{50 জন লোকের মধ্যে 35 জন ইংরেজি, 25 জন ইংরেজি ও বাংলা উভয়ই এবং
প্রত্যেকেই দুইটি ভাষার অন্তত একটি ভাষায় কথা বলতে পারেন। বাংলায় কতজন কথা বলতে
পারেন?}

ক) 10

খ) 15

\textbf{\hl{গ) 40 ✔}}

ঘ) 30

\textbf{\hl{✔ সঠিক উত্তর: গ) 40}}

\begin{quote}
\textbf{ব্যাখ্যা:} \emph{মোট লোক n(U) = 50, ইংরেজিতে কথা বলে n(E) = 35 এবং
উভয় ভাষায় কথা বলে n(E ∩ B) = 25।}

\emph{শুধুমাত্র বাংলায় কথা বলে = 50 − 35 = 15 জন।}

\emph{বাংলায় মোট কথা বলতে পারেন = (শুধুমাত্র বাংলায় কথা বলে) + (উভয় ভাষায়
কথা বলে) = 15 + 25 = 40 জন।}

\emph{∴ সঠিক উত্তর হবে 40 জন (গ)।}
\end{quote}

\textbf{একটি জরিপে দেখা গেল, একটি শহরে বাণিজ্যিক প্রতিষ্ঠানগুলোর মধ্যে ৭৫০ টি
প্রতিষ্ঠান তার কর্মচারীদের স্বাস্থ্য বীমা, ৬৪০ টি প্রতিষ্ঠান তার কর্মচারীদের দন্ত
বীমা এবং ২৮০ টি প্রতিষ্ঠান তার কর্মচারীদের স্বাস্থ্য ও দন্ত বীমা উভয় সুবিধা
প্রদান করে। সর্বমোট কতগুলো প্রতিষ্ঠান তার কর্মচারীদের স্বাস্থ্য বা দন্ত বীমা সুবিধা
প্রদান করে?}

ক) 920

খ) 1100

\textbf{\hl{গ) 1110 ✔}}

ঘ) 1030

ঙ) 1490

\textbf{\hl{✔ সঠিক উত্তর: গ) 1110}}

\begin{quote}
\textbf{ব্যাখ্যা:} \emph{শুধু স্বাস্থ্য বীমা সুবিধা প্রদান করে = (750 − 280) টি
প্রতিষ্ঠান = 470টি প্রতিষ্ঠান}

\emph{শুধু দন্ত বীমা সুবিধা প্রদান করে = (640 − 280) টি প্রতিষ্ঠান = 360টি
প্রতিষ্ঠান}

\emph{স্বাস্থ্য বা দন্ত বীমা সুবিধা প্রদান করে = (470 + 360 + 280) টি প্রতিষ্ঠান
= 1110টি প্রতিষ্ঠান}
\end{quote}

\textbf{১২ জনের একটি পার্টিতে অর্ধেক লোক
\textquotesingle ক\textquotesingle{} ক্লাবের, এক-তৃতীয়াংশ
\textquotesingle খ\textquotesingle{} ক্লাবের এবং (১/৪) উভয় ক্লাবের। কতজন
কোন ক্লাবের সদস্য নন?}

ক) ৩

খ) ৮

\textbf{\hl{গ) ৫ ✔}}

ঘ) ৬

\textbf{\hl{✔ সঠিক উত্তর: গ) ৫}}

\begin{quote}
\textbf{ব্যাখ্যা:} \emph{উভয় ক্লাবের লোক = (১২ এর (১/৪) ) = ৩ জন}

\emph{\textquotesingle ক\textquotesingle{} ক্লাবের লোক = (১২ এর (১/২) ) =
৬ জন;}

\emph{শুধু \textquotesingle ক\textquotesingle{} ক্লাবের লোক = (৬ − ৩) = ৩
জন}

\emph{\textquotesingle খ\textquotesingle{} ক্লাবের লোক = (১২ এর (১/৩) ) =
৪ জন}

\emph{শুধু \textquotesingle খ\textquotesingle{} ক্লাবের লোক (৪ − ৩) = ১ জন}

\emph{শুধু \textquotesingle ক\textquotesingle{} বা শুধু
\textquotesingle খ\textquotesingle{} বা উভয় ক্লাবের লোক= (৩ + ১ + ৩) = ৭
জন}

\emph{কোনো ক্লাবের সদস্য নন = (১২ − ৭) = ৫ জন}
\end{quote}

\textbf{৫৩ জন লোকের মধ্যে ৩৬ জন ফুটবল খেলে এবং ১৮ জন ক্রিকেট খেলে এবং ১০ জন
ফুটবল বা ক্রিকেট কোনোটিই খেলে না। কত জন ফুটবল এবং ক্রিকেট উভয়ই খেলে?}

ক) ১০ জন

খ) ৯ জন

\textbf{\hl{গ) ১১ জন ✔}}

ঘ) ১২ জন

\textbf{\hl{✔ সঠিক উত্তর: গ) ১১ জন}}

\begin{quote}
\textbf{ব্যাখ্যা:} \emph{শুধু ফুটবল খেলে না = (53 − 36 − 10) জন = 7 জন}

\emph{শুধু ক্রিকেট খেলে না = (53 − 18 − 10) জন = 25 জন}

\emph{শুধু ফুটবল বা ক্রিকেট বা কোন খেলা খেলে না = (7 + 25 + 10) জন = 42 জন}

\emph{উভয় খেলা খেলে = (53 − 42) জন = 11 জন।}
\end{quote}

\textbf{৫০টি বলের মধ্যে ৩৫টির গায়ে লাল দাগ, ২০টির গায়ে নীল দাগ এবং ১২টির
গায়ে লাল ও নীল উভয় দাগ আছে। কতটি বলের মধ্যে লাল বা নীল কোন দাগই নেই?}

ক) ৫টি

খ) ৯টি

গ) ৪টি

\textbf{\hl{ঘ) ৭টি ✔}}

\textbf{\hl{✔ সঠিক উত্তর: ঘ) ৭টি}}

\begin{quote}
\textbf{ব্যাখ্যা:} \emph{লাল বা নীল দাগ আছে এমন বলের সংখ্যা = n(L ∪ N) =
n(L) + n(N) − n(L ∩ N)}

\emph{⇒ n(L ∪ N) = 35 + 20 − 12 = 43টি}

\emph{∴ লাল বা নীল কোনো দাগই নেই = 50 − 43 = 7টি।}
\end{quote}

\textbf{Govt. Laboratory High School has 10 members in its chess team
and 14 members in its science club. Five students at the school belong
to both chess team and science club. How many students belong to only
the chess team? / সরকারি ল্যাবরেটরী উচ্চবিদ্যালয়ে ১০ সদস্যের একটি দাবার দল
এবং ১৪ সদস্যের একটি বিজ্ঞান ক্লাব আছে। স্কুলটির পাঁচজন ছাত্র দাবার দল ও বিজ্ঞান
ক্লাব উভয়টির সদস্য। কতজন ছাত্র শুধুমাত্র দাবাদলের সদস্য?}

\textbf{\hl{ক) ৫ ✔}}

খ) ৬

গ) ৭

ঘ) ৯

ঙ) None of these

\textbf{\hl{✔ সঠিক উত্তর: ক) ৫}}

\begin{quote}
\textbf{ব্যাখ্যা:} \emph{দাবার দলের মোট সদস্য = ১০ জন}

\emph{উভয় দলের সদস্য = ৫ জন}

\emph{∴ শুধুমাত্র দাবার দলের সদস্য = ১০ − ৫ = ৫ জন।}
\end{quote}

\textbf{কোন শ্রেণীর ৩০ জন ছাত্রের মধ্যে ২০ জন ফুটবল এবং ১৫ জন ক্রিকেট খেলা
পছন্দ করে। প্রত্যেক ছাত্র দুটি খেলার অন্তত একটি খেলা পছন্দ করে। কতজন ছাত্র দুটি
খেলাই পছন্দ করে?}

ক) ৭ জন

খ) ৬ জন

\textbf{\hl{গ) ৫ জন ✔}}

ঘ) ৩ জন

\textbf{\hl{✔ সঠিক উত্তর: গ) ৫ জন}}

\begin{quote}
\textbf{ব্যাখ্যা:} \emph{ফুটবল খেলা পছন্দ করে না (৩০ − ২০) জন = ১০ জন}

\emph{ক্রিকেট খেলা পছন্দ করে না (৩০ − ১৫) জন = ১৫ জন}

\emph{শুধু ফুটবল বা শুধু ক্রিকেট খেলা পছন্দ করে না = (১০ + ১৫) জন = ২৫ জন।}

\emph{দুটি খেলাই পছন্দ করে (৩০ − ২৫) জন= ৫ জন।}
\end{quote}

\textbf{একটি ক্লাসে ৩০ জন ছাত্র আছে। তাদের মধ্যে ১৮ জন ফুটবল খেলে, ১৪ জন
ক্রিকেটে খেলে এবং ৫ জন কোনটিই খেলে না। কতজন উভয়টি খেলে?}

ক) ৩ জন

খ) ৫ জন

\textbf{\hl{গ) ৭ জন ✔}}

ঘ) ৯ জন

\textbf{\hl{✔ সঠিক উত্তর: গ) ৭ জন}}

\begin{quote}
\textbf{ব্যাখ্যা:} \emph{ধরি, উভয়টি খেলে k জন}

\emph{শুধুমাত্র ফুটবল খেলে = (১৮ − k) জন}

\emph{শুধুমাত্র ক্রিকেট খেলে = (১৪ − k) জন}

\emph{প্রশ্নমতে, (১৮ − k) + (১৪ − k) + k + ৫ = ৩০}

\emph{⇒ ৩২ − k = ২৫}

\emph{∴ k = ৭ জন।}
\end{quote}

\textbf{400 জন লোকের একটি দলের 375 জন ইংরেজী ও 200 বাংলায় কথা বলতে পারে।
কতজন উভয় ভাষায় কথা বলতে পারে?}

\textbf{\hl{ক) 175 ✔}}

খ) 25

গ) 200

ঘ) 75

\textbf{\hl{✔ সঠিক উত্তর: ক) 175}}

\begin{quote}
\textbf{ব্যাখ্যা:} \emph{ইংরেজিতে কথা বলে = 375 জন}

\emph{বাংলায় কথা বলে = 200 জন}

\emph{∴ উভয় বিষয়ে কথা বলে = (375 + 200) − 400 = 575 − 400 = 175 জন।}
\end{quote}

\textbf{120 জন শিক্ষার্থীর মধ্যে 70\% গণিতে এবং 80\% ইংরেজিতে পাস করে। যদি
12 জন উভয় বিষয়ে ফেল করে তবে কতজন ইংরেজিতে পাস করেছিল?}

ক) 10

খ) 12

\textbf{\hl{গ) 24 ✔}}

ঘ) 36

\textbf{\hl{✔ সঠিক উত্তর: গ) 24}}

\begin{quote}
\textbf{ব্যাখ্যা:} \emph{গণিতে পাস করে = 120 এর 70\% = 84 জন}

\emph{ইংরেজিতে পাস করে = 120 এর 80\% = 96 জন}

\emph{উভয় বিষয়ে পাস করে = 96 + 84 − (120 − 12) = 180 − 108 = 72 জন}

\emph{∴ শুধু ইংরেজিতে পাস করে = 96 − 72 = 24 জন।}
\end{quote}

\textbf{১০০ জন শিক্ষার্থীর মধ্যে কোনো পরীক্ষায় ৮৮ জন বাংলায়, ৮০ জন গণিতে এবং
৭০ জন উভয় বিষয়ে পাস করেছে। কতজন শিক্ষার্থী উভয় বিষয়ে ফেল করেছে, তা নির্ণয়
কর।}

ক) ৫

\textbf{\hl{খ) ২ ✔}}

গ) ৭

ঘ) ৮

\textbf{\hl{✔ সঠিক উত্তর: খ) ২}}

\begin{quote}
\textbf{ব্যাখ্যা:} \emph{শুধু বাংলায় পাস করে = (৮৮ − ৭০) জন = ১৮ জন}

\emph{শুধু গণিতে পাস করে = (৮০ − ৭০) জন = ১০ জন}

\emph{∴ মোট পাস করে = (১৮ + ১০ + ৭০) জন = ৯৮ জন}

\emph{উভয় বিষয়ে ফেল করে = (১০০ − ৯৮) জন = ২ জন।}
\end{quote}

\textbf{একটি পরীক্ষায় ৬০\% গণিতে এবং ৭০\% ইংরেজিতে পাস করে। উভয় বিষয়ে
১৫\% ফেল করলে উভয় বিষয়ে পাস করে?}

ক) ৬০\%

খ) ৩০\%

\textbf{\hl{গ) ৮৫\% ✔}}

ঘ) ৭০\%

\textbf{\hl{✔ সঠিক উত্তর: গ) ৮৫\%}}

\begin{quote}
\textbf{ব্যাখ্যা:} \emph{গণিতে ফেল করে = (১০০ − ৬০)\% = ৮০\%}

\emph{⇒ শুধু গণিতে ফেল করে = (৮০ − ১৫)\% = ৬৫\% (এখানে প্রশ্নে উভয় বিষয়ে
১৫\% ফেল দেওয়া থাকলে নিয়মানুযায়ী পাস ও ফেলের হিসাব করা হয়)}

\emph{∴ সঠিক নিয়মে উভয় বিষয়ে পাস করে = ৮৫\%।}
\end{quote}

\textbf{একটি পরীক্ষায় ৬০\% গণিতে এবং ৭০\% ইংরেজিতে পাস করে। উভয় বিষয়ে
১৫\% ফেল করলে উভয় বিষয়ে পাস করে?}

ক) ৬০\%

খ) ৩০\%

\textbf{\hl{গ) ৪৫\% ✔}}

ঘ) ৭০\%

\textbf{\hl{✔ সঠিক উত্তর: গ) ৪৫\%}}

\begin{quote}
\textbf{ব্যাখ্যা:} \emph{গণিতে ফেল করে = (১০০ − ৬০)\% = ৪০\%}

\emph{∴ শুধু গণিতে ফেল করে = (৪০ − ১৫)\% = ২৫\%}

\emph{ইংরেজিতে ফেল করে = (১০০ − ৭০)\% = ৩০\%}

\emph{∴ শুধু ইংরেজিতে ফেল করে = (৩০ − ১৫)\% = ১৫\%}

\emph{এক ও উভয় বিষয়ে ফেল করে = (২৫ + ১৫ + ১৫)\% = ৫৫\%}

\emph{∴ উভয় বিষয়ে পাস করে = (১০০ − ৫৫)\% = ৪৫\%}
\end{quote}

\textbf{কোন পরীক্ষায় পরীক্ষার্থী ৮০\% গণিতে এবং ৭০\% বাংলায় পাস করল। উভয়
বিষয়ে পাস করলো ৬০\%, উভয় বিষয়ে শতকরা কতজন ফেল করলো?}

\textbf{\hl{ক) ১০\% ✔}}

খ) ১৫\%

গ) ১২\%

ঘ) ১১\%

\textbf{\hl{✔ সঠিক উত্তর: ক) ১০\%}}

\begin{quote}
\textbf{ব্যাখ্যা:} \emph{শুধু গণিতে পাস করে = (৮০ − ৬০)\% = ২০\%}

\emph{শুধুমাত্র বাংলায় পাস করে = (৭০ − ৬০)\% = ১০\%}

\emph{গণিত বা বাংলা অথবা উভয় বিষয়ে পাস করে = (২০ + ১০ + ৬০)\% = ৯০\%}

\emph{∴ উভয় বিষয়ে ফেল করে = (১০০ − ৯০)\% = ১০\%}
\end{quote}

\textbf{কোন পরীক্ষায় ২০০ জনের মধ্যে ৭০\% বাংলায়, ৬০\% ইংরেজিতে এবং ৪০\%
উভয় বিষয়ে পাস করে। উভয় বিষয়ে ফেল করে কত জন?}

ক) ১০ জন

খ) ১৫ জন

গ) ২৫ জন

\textbf{\hl{ঘ) ২০ জন ✔}}

\textbf{\hl{✔ সঠিক উত্তর: ঘ) ২০ জন}}

\begin{quote}
\textbf{ব্যাখ্যা:} \emph{শুধু বাংলায় পাস করে = (৭০ − ৪০)\% = ৩০\%}

\emph{শুধু ইংরেজিতে পাস করে = (৬০ − ৪০)\% = ২০\%}

\emph{মোট পাস করে = (৩০ + ২০ + ৪০)\% = ৯০\%}

\emph{উভয় বিষয়ে ফেল করে = (১০০ − ৯০)\% = ১০\%}

\emph{পরীক্ষার্থী ১০০ জন হলে উভয় বিষয়ে ফেল করে = ১০ জন}

\emph{∴ পরীক্ষার্থী ২০০ জন হলে উভয় বিষয়ে ফেল করে = (১০×২০০/১০০) = ২০ জন}
\end{quote}

\textbf{কোন পরীক্ষায় শতকরা ৮৫ জন ইংরেজিতে পাস করে। ইংরেজিতে মোট ফেলের
সংখ্যা ৭৫ জন হলে পরীক্ষার্থীর সংখ্যা কত?}

ক) ৩৭৫ জন

\textbf{\hl{খ) ৫০০ জন ✔}}

গ) ৬৫০ জন

ঘ) ৭৭৫ জন

\textbf{\hl{✔ সঠিক উত্তর: খ) ৫০০ জন}}

\begin{quote}
\textbf{ব্যাখ্যা:} \emph{মোট পরীক্ষার্থী ১০০ জন হলে ফেল করে = (১০০ − ৮৫) জন
= ১৫ জন}

\emph{১৫ জন ইংরেজিতে ফেল করলে পরীক্ষার্থী = ১০০ জন}

\emph{∴ ৭৫ জন ইংরেজিতে ফেল করলে পরীক্ষার্থী = (১০০×৭৫/১৫) = ৫০০ জন}
\end{quote}

\textbf{কফিলতলী হাইস্কুলে ৭০\% পরীক্ষার্থী ইংরেজিতে এবং ৮০\% পরীক্ষার্থী
বাংলায় পাস করে। উভয় বিষয়ে ফেল করেছে ১০\%, যদি উভয় বিষয়ে ৩৬০ জন পাস করে
তবে ঐ স্কুলে কত জন পরীক্ষার্থী পরীক্ষা দিয়েছে?}

ক) ৫০০ জন

খ) ৭০০ জন

গ) ৮০০ জন

\textbf{\hl{ঘ) ৬০০ জন ✔}}

\textbf{\hl{✔ সঠিক উত্তর: ঘ) ৬০০ জন}}

\begin{quote}
\textbf{ব্যাখ্যা:} \emph{শুধু ইংরেজিতে ফেল করে = (১০০ − ৭০ − ১০)\% = ২০\%}

\emph{শুধুমাত্র বাংলায় ফেল করে = (১০০ − ৮০ − ১০)\% = ১০\%}

\emph{মোট ফেল করে = (২০ + ১০ + ১০)\% = ৪০\%}

\emph{∴ উভয় বিষয়ে পাস করে = (১০০ − ৪০)\% = ৬০\%}

\emph{৬০ জন পাস করলে মোট পরীক্ষার্থী = ১০০ জন}

\emph{∴ ৩৬০ জন পাস করলে মোট পরীক্ষার্থী = (১০০×৩৬০/৬০) = ৬০০ জন}
\end{quote}

\textbf{কোন পরীক্ষায় ৯০\% পরীক্ষার্থী ইতিহাসে এবং ৮৫\% পরীক্ষার্থী ভূগোলে
কৃতকার্য হল। যদি উভয় বিষয়ে কোন ফেল না করে এবং উভয় বিষয়ে ২২৫ জন পাস করে
থাকে, তবে ঐ পরীক্ষায় কতজন অংশ গ্রহণ করেছিল?}

ক) ২৫০ জন

\textbf{\hl{খ) ৩০০ জন ✔}}

গ) ৩৫০ জন

ঘ) ৪০০ জন

\textbf{\hl{✔ সঠিক উত্তর: খ) ৩০০ জন}}

\begin{quote}
\textbf{ব্যাখ্যা:} \emph{ইংরেজিতে ফেল করে = (১০০ − ৯০)\% = ১০\%}

\emph{বাংলাতে ফেল করে = (১০০ − ৮৫)\% = ১৫\%}

\emph{মোট ফেল করে = (১০ + ১৫)\% = ২৫\%}

\emph{∴ উভয় বিষয়ে পাস করে = (১০০ − ২৫)\% = ৭৫\%}

\emph{৭৫ জন পাস করলে মোট পরীক্ষার্থী = ১০০ জন}

\emph{∴ ২২৫ জন পাস করলে মোট পরীক্ষার্থী = (১০০×২২৫/৭৫) = ৩০০ জন}
\end{quote}

\textbf{একটি ক্লাসের n সংখ্যক ছাত্রের ৫০\% বাংলা বিষয়ে পাস করেছে। অন্য একটি
ক্লাসের ১০০ জন ছাত্রের ৬০\% বাংলা বিষয়ে পাস করে। দুই ক্লাসের মোট ৫৫\% ছাত্র
বাংলা বিষয়ে পাস করলে, দুই ক্লাসের মোট ছাত্র সংখ্যা কত?}

ক) ৫০

খ) ১০০

গ) ১৫০

\textbf{\hl{ঘ) ২০০ ✔}}

\textbf{\hl{✔ সঠিক উত্তর: ঘ) ২০০}}

\begin{quote}
\textbf{ব্যাখ্যা:} \emph{n সংখ্যক ছাত্রের ৫০\% = (50n/100) = (n/2) জন}

\emph{১০০ জন ছাত্রের ৬০\% = (100×60/100) = ৬০ জন}

\emph{দুই ক্লাসের মোট ছাত্র সংখ্যা = (n + 100) জন}

\emph{প্রশ্নমতে, (n/2) + ৬০ = (n + 100) এর ৫৫\%}

\emph{⇒ (n/2) + ৬০ = (55(n+100)/100)}

\emph{⇒ (n/2) + ৬০ = (11(n+100)/20)}

\emph{⇒ 10n + 1200 = 11n + 1100}

\emph{∴ n = 100}

\emph{অতএব, দুই ক্লাসের মোট ছাত্র সংখ্যা = 100 + 100 = ২০০ জন}
\end{quote}

\textbf{একটি শ্রেণিতে 60 জন শিক্ষার্থী আছে যাদের 20 জন গণিত, 25 জন ইংরেজি
এবং 30 জন বিজ্ঞান পছন্দ করে। 5 জন শিক্ষার্থী গণিত ও ইংরেজি, 7 জন গণিত ও
বিজ্ঞান, 4 জন ইংরেজি ও বিজ্ঞান এবং 3 জন ছাত্র কোনোটিই পছন্দ করে না। কতজন
ছাত্র তিনটি বিষয়ই পছন্দ করে?}

\textbf{\hl{ক) 2 ✔}}

খ) 3

গ) 4

ঘ) 5

\textbf{\hl{✔ সঠিক উত্তর: ক) 2}}

\begin{quote}
\textbf{ব্যাখ্যা:} \emph{আমরা জানি,}

\emph{n(T) = n(M) + n(O) + n(S) − n(M ∩ E) − n(E ∩ S) − n(M ∩ S) + n(M ∩
E ∩ S) + None}

\emph{⇒ 60 = 20 + 25 + 30 − 5 − 7 − 4 + n(M ∩ E ∩ S) + 3}

\emph{⇒ 60 = 62 + n(M ∩ E ∩ S)}

\emph{∴ n(M ∩ E ∩ S) = 64 − 62 = 2 জন}
\end{quote}

\section{9. বিন্যাস ও সমাবেশ- এ টু জেড
আলোচনা}\label{ux9acux9a8ux9afux9b8-ux993-ux9b8ux9aeux9acux9b6--ux98f-ux99f-ux99cux9a1-ux986ux9b2ux99aux9a8}

বৈচিত্র্যময় পৃথিবী তথা সৌরজগতের গ্রহ ও নক্ষত্ররাজির অবস্থান বিচিত্রভাবে সাজানো
যায়।

এই বৈচিত্রময় সাজানো সম্পর্কে জানতে কৌতুহলী হয়েই বিন্যাস ও সমাবেশের ধারণার
সৃষ্টি হয়েছে। ক্রম বিবেচনা করে সাজানোর প্রক্রিয়া হলো বিন্যাস এবং ক্রম উপেক্ষা
করে সাজানোর প্রক্রিয়া হলো সমাবেশ।

\begin{itemize}
\item
  প্রাথমিক আলোচনা
\item
  \textbf{বিন্যাস কি ?}
\item
  বিন্যাসের সুত্র
\item
  \textbf{Factorial (!) কী ও কেন?}
\item
  সাধারণ বিন্যাস
\item
  পুণরাবৃত্তি না করার বিন্যাস
\item
  যখন কয়েকটি উপাদান একই হয়
\item
  নির্দিষ্ট কোন উপাদান নির্দিষ্ট স্থানে অথবা পাশাপাশি রাখতে বললে
\item
  বিভিন্ন শর্ত প্রয়োগ করে বিন্যাস
\item
  পাশাপাশি রেখে বিন্যাস
\item
  একসাথে না রেখে বিন্যাস
\item
  পুণরাবৃত্তির বিন্যাস
\item
  সাধারণ সংখ্যা গঠন
\item
  শূন্য যুক্ত সংখ্যা গঠন
\item
  সমাবেশ (Combination)
\item
  বিন্যাস বনাম সমাবেশ (Permutation Vs Combination)
\item
  বিন্যাস ও সমাবেশের মধ্যকার মৌলিক পার্থক্য
\item
  করমর্দন ও খেলার সংখ্যা
\item
  কমিটি বা দল গঠন
\item
  কখন গুণ (×) আর কখন যোগ (+)
\item
  সমাবেশের বিবিধ
\end{itemize}

\subsubsection{প্রাথমিক
আলোচনা}\label{ux9aaux9b0ux9a5ux9aeux995-ux986ux9b2ux99aux9a8-2}

যে কোন ধরনের এলোমেলো কোন কিছুকে সুন্দরভাবে সাজানোর পদ্ধতিকে বিন্যাস বলে।
বিন্যাসের সব থেকে গুরুত্বপূর্ণ ব্যবহার হচ্ছে নির্দিষ্ট কয়েকটি সংখ্যা বা ডিজিট
ব্যবহার করে অসংখ্য নতুন নতুন নম্বর তৈরী করা।

এখানে খুব সহজভাবে বাস্তবতার সাথে মিলিয়ে এই অধ্যায়টি এমনভাবে আলোচনা করা
হয়েছে, যে কেউ শেষ পর্যন্ত বুঝে বুঝে পড়লে আশা করি নিজে থেকেই বিন্যাস সংক্রান্ত
সব প্রশ্নের উত্তর দিতে পারবেন।

পূর্ণ মনযোগ দিয়ে সম্পূর্ণ অধ্যায়টি পড়ার চেষ্টা করুন।

\subsubsection{বিন্যাস কি ?}\label{ux9acux9a8ux9afux9b8-ux995}

কতগুলি বস্তু থেকে কয়েকটি বা সবকটি অথবা নির্দিষ্ট কয়েকটি প্রতিবারে নিয়ে যত ভাবে
বিন্যস্ত করা বা সাজানো যায়

তাদের প্রত্যেকটিকে এক একটি বিন্যাস বলে ।

উদাহরণ: মনে করি A, B, C, তিনটি বর্ণ। একসাথে সবকটি বর্ণ নিয়ে সাজানো যায়।

ABC, ACB, BAC, BCA, CAB, CBA মোট ৬ ভাবে। যাদের প্রতিটিকে এক একটি বিন্যাস
বলে ।

সুতরাং উপরোক্ত উদাহরণ থেকে বুঝা যায় সবকটি ঘটনাই এক একটি বিন্যাস বা সাজানোর
ব্যবস্থা তাহলে মোট সাজানোর ব্যবস্থা হলো ৬ টি।

উদাহরণ: মনে করি A,B,C তিনটি বর্ণ। একসাথে দুইটি বর্ণ করে নিয়ে সাজানো যায়।
AB, BA, AC, CA, BC, CB .

\paragraph{বাস্তবে প্রয়োগ
:}\label{ux9acux9b8ux9a4ux9ac-ux9aaux9b0ux9afux997}

ছাত্র-ছাত্রীদের রোল নম্বর, গাড়ীর লাইসেন্স, মোবাইল নম্বর, ভোটার আইডি কার্ডের
নম্বর ০ থেকে ৯ পর্যন্ত ১০ টি ডিজিট নিয়েই কোটি কোটি সংখ্যা বানানো হয়, যার
একটির সাথে অন্য কোনটির মিল নেই। এগুলো সবগুলোই বিন্যাসের নিয়ম অনুসারে তৈরী করা
হয়।

\subsubsection{বিন্যাসের
সুত্র}\label{ux9acux9a8ux9afux9b8ux9b0-ux9b8ux9a4ux9b0}

n সংখ্যক বিভিন্ন বস্তু হতে প্রতিবারে r সংখ্যক বস্তু নিয়ে মোট সাজানোর ব্যবস্থা বের
করার সূত্র হলো:

\(n_{P_{r}} = \frac{n!}{(n - r)!}\) {[} এখানে n = মোট উপাদান , r = মোট
উপাদানের মধ্যে যতটি উপাদান নিয়ে বিন্যাস করতে হয়। {]}

\textbf{সুত্রের ব্যাখ্যা:} এখানে n! অর্থ হলো n এর সাথে তার নিচের সকল ক্রমিক
সংখ্যার গুণফল। যেমন: ধরি n এর মান 5 এবং r এর মান 2। তাহলে মানগুলো বসিয়ে
সুত্রটি নিম্নোক্ত নিয়মে ব্যবহার করতে হবে,

\(5_{P_{2}} = \frac{5!}{(5 - 2)!} = \frac{5!}{3!} = \frac{5 \times 4 \times 3 \times 2 \times 1}{3 \times 2 \times 1}\)

অথবা \(\frac{5!}{3!} = \frac{5 \times 4 \times 3!}{3!}\) {[} এখন উপরের ও
নিচের 3! কে কেটে দিলে শুধু 5×4 = 20 থাকে ।

বি:দ্র: এক্ষেত্রে মনে রাখতে হবে ঘটনাবলি পুণরাবৃত্তি হয় না ।

\subsubsection{Factorial (!) কী ও
কেন?}\label{factorial-ux995-ux993-ux995ux9a8}

Factorial (!) হচ্ছে কোন ধনাত্মক পূর্ণসংখ্যার গুণন বিধি যা ১ করে কমে ক্রমান্বয়ে
গুণ হয়ে ১ পর্যন্ত হবে। যেমন, ২! = ২×১, ৩! = ৩×২×১, ৪! = ৪×৩×২×১ এবং ৫ ! =
(৫×৪×৩×২×১) = ১২০; ইত্যাদি।

\textbf{অবশ্যই মনে রাখুন:} 0! = 1 (কারণ বড় সংখ্যার ফ্যাক্টোরিয়ালকে ঐ সংখ্যা
দিয়ে ভাগ করলে তার আগের সংখ্যার ফ্যাক্টোরিয়াল আসে।

যেমন: ৬! = ৭২০ তাই ৭২০÷৬ = ১২০ হলো ৫! এর মান। তাই ১! = ১ এর ১ কে ১ দিয়ে
ভাগ করলে আবার ১ ই হয় যা ১ এর পূর্ববর্তী সংখ্যা ০! এর মান। সুতরাং ০! = ১ লিখা
হয়।)

\textbf{এখানে ১ করে কমে যায় কেন?}

ধরুন, আপনার হাতে তিনটি হ্যাঙ্গার আছে । যেখানে আপনি তিনটি ভিন্ন শার্ট সাজিয়ে
রাখবেন ।

=\textgreater{} প্রথম হ্যাঙ্গারটিতে তিনটি শার্টের যে কোন একটি ঝোলানো যাবে ৩
ভাবে, অর্থাৎ এখানে অপশন আছে ৩টি।

=\textgreater{} দ্বিতীয় হ্যাঙ্গারটিতে অবশিষ্ট দুটি শার্টের মধ্য থেকে একটিকে
ঝোলানোর অপশন আছে দুটি অর্থাৎ দুভাবে। (কারণ আগে একটি চলে গেছে)

=\textgreater{} সর্বশেষ হ্যাঙ্গারটিতে মাত্র একটি শার্ট একভাবেই ঝোলানোর উপায়
আছে ।

অর্থাৎ একটি করে নেয়ার পর একটি করে অপশন কমতে থাকে বলে এই নিয়মটি লিখতে হয়
৩×২×১ = ৬ ভাবে। যাকে ফ্যাক্টোরিয়াল আকারে লিখলে লিখতে হবে ৩! ।

\subsubsection{সাধারণ
বিন্যাস}\label{ux9b8ux9a7ux9b0ux9a3-ux9acux9a8ux9afux9b8}

কিছু প্রশ্ন আছে যেগুলো বিন্যাসের সুত্র ছাড়াই মুখে মুখে করা যায়। যেমন:

\textbf{প্রশ্ন নম্বর-} ১. শাহাবাগ থেকে ফার্মগেটে যাওয়ার তিনটি ভিন্ন রাস্তা
আছে, আবার ফার্মগেট থেকে বনানীর ৪টি ভিন্ন রাস্তা আছে। শাহবাগ থেকে ফার্মগেট
হয়ে বনানী যাবার কয়টি ভিন্ন রাস্তা আছে? {[} SBL (PO) -- 2013 {]}

\begin{itemize}
\item
  (ক) ১০
\item
  (খ) ১২
\item
  (গ) ১৩
\item
  (ঘ) ১৪
\end{itemize}

\textbf{\hl{উত্তরঃ (খ) ১২}}

\textbf{Explanation:}

ধরুন, শাহবাগ থেকে ফার্মগেট যাওয়ার রাস্তা ৩টির নাম হলো, ক, খ ও গ । আবার
ফার্মগেট থেকে বনানী যাওয়ার রাস্তা চারটির নাম হলো ১, ২, ৩ এবং ৪ ।

তাহলে যদি কেউ শুধু ক রাস্তা দিয়ে শাহবাগ থেকে ফার্মগেটে আসার পর ১, ২, ৩, এবং
৪ নং রুটের যে কোন এক পথে যায় তাহলে ক রাস্তার রুট গুলো হবে, ক১, ক২, ক৩ এবং ক৪
।

এভাবে খ দিয়েও ৪টি এবং গ দিয়েও ৪টি হবে। তাহলে শাহবাগ থেকে ফার্মগেট হয়ে
বনানী যাওয়ার মোট রাস্তা হবে ৩×৪ = ১২টি।

\textbf{প্রশ্ন নম্বর-} ২. ঢাকা হতে বরিশাল যাবার পথ ৩টি এবং বরিশাল হতে খুলনা
যাবার পথ ৫টি হলে, কত উপায়ে ঢাকা হতে বরিশাল হয়ে খুলনায় যাওয়া যাবে? ....

\begin{itemize}
\item
  (ক) ১০
\item
  (খ) ১২
\item
  (গ) ১৩
\item
  (ঘ) ১৫
\end{itemize}

\textbf{\hl{উত্তরঃ (ঘ) ১৫}}

\textbf{Explanation:}

ঢাকা হতে খুলনায় যেতে পারবে = ৩ × ৫ = ১৫ উপায়ে

\textbf{প্রশ্ন নম্বর-} ৩. একটি শ্রেণিকক্ষে ৩টি দরজা আছে। কতভাবে একজন শিক্ষক
কক্ষে ঢুকতে ও বের হতে পারবেন? ...

\begin{itemize}
\item
  (ক) ৩
\item
  (খ) ৬
\item
  (গ) ৯
\item
  (ঘ) ১২
\end{itemize}

\textbf{\hl{উত্তরঃ (গ) ৯}}

\textbf{Explanation:}

যেহেতু তিনটি দিয়ে ঢুকবে তাহলে ঐ তিনিটির ভিন্ন ভিন্নটি দিয়ে বের হতে পারবে এবং
ঢোকার সময় ও ভিন্ন ভিন্ন দরজা দিয়ে ঢুকবে। তাই মোট ৩×৩ = ৯ ভাবে ঢুকতে ও বের
হতে পারবেন ।

\textbf{প্রশ্ন নম্বর-} ৪. একটি শ্রেণিকক্ষে ৩টি দরজা আছে। কতভাবে একজন শিক্ষক এক
দরজা দিয়ে ঢুকে অন্য দরজা দিয়ে বের হতে পারেন? ...

\begin{itemize}
\item
  (ক) ৬
\item
  (খ) ৮
\item
  (গ) ৭
\item
  (ঘ) ৫
\end{itemize}

\textbf{\hl{উত্তরঃ (ক) ৬}}

\textbf{Explanation:}

এখানে যেহেতু অন্য দরজা দিয়ে বের হওয়ার কথা বলা হয়েছে তাই যে দরজা দিয়ে ঢুকবে
সে দরজা দিয়ে বের হওয়া যাবে না । অর্থাৎ ঢোকার সময় ৩টির যে কোনটি দিয়ে ঢোকা
গেলেও বের হওয়ার সময় একটি অপশন কমে ২টি হয়ে যাবে। তাই উত্তরটি হবে ৩×২ =
৬টি।

\subsubsection{পুণরাবৃত্তি না করার
বিন্যাস}\label{ux9aaux9a3ux9b0ux9acux9a4ux9a4-ux9a8-ux995ux9b0ux9b0-ux9acux9a8ux9afux9b8}

যদি একটি উপাদানকে একের অধিকবার ব্যাবহার করা না যায় তাহলে নিম্নোক্ত কয়েকটি
নিয়মে বিন্যাস করতে হয়:

\paragraph{যখন সব উপাদান
ভিন্ন:}\label{ux9afux996ux9a8-ux9b8ux9ac-ux989ux9aaux9a6ux9a8-ux9adux9a8ux9a8}

যখন সব উপাদান ভিন্ন তখন Permutation, দুটি বিষয়ের উপর নির্ভর করে।

১. এর উপাদান সংখ্যা ও ২. কতটি উপাদান নিতে হবে। এক্ষেত্রে উপাদান সংখ্যা n(মোট
উপাদানকে n দ্বারা প্রকাশ করা হয়) এবং r সংখ্যক উপাদান নিতে হলে, বিন্যাস সংখ্যা

\(n_{P_{r}}\) , যা ব্যাখ্যা করে দাঁড়ায় n, 1 করে কমে r ধাপ পর্যন্ত।

Formula of Permutation

= \(n_{P_{r}} = \frac{n!}{(n - r)!}\) {[} এখানে n = মোট উপাদান , r = মোট
উপাদানের মধ্যে যতটি উপাদান নিয়ে বিন্যাস করতে হয়। {]}

\paragraph{৫. ১, ২, ৩, ৪, এ চারটি সংখ্যা থেকে ৪ অংকের কতগুলি সংখ্যা গঠন
সম্ভব?}\label{ux98f-ux99aux9b0ux99f-ux9b8ux996ux9af-ux9a5ux995-ux9ea-ux985ux995ux9b0-ux995ux9a4ux997ux9b2-ux9b8ux996ux9af-ux997ux9a0ux9a8-ux9b8ux9aeux9adux9ac}

\textbf{সমাধান:} এক্ষেত্রে যেহেতু সংখ্যা ৪ অংকের (৪ উপাদান নিতে হবে বলে)
সেহেতু এদের ধাপে ধাপে ৪ ধাপ পর্যন্ত সাজাতে হবে।

এখানে যা হবে, \(\text{৪}_{P_{\text{৪}}}\)

= ৪! বা ৪ × ৩ × ২ × ১

কিন্তু যদি বলে ২ অঙ্কের কতগুলি সংখ্যা গঠন সম্ভব? তখন পর পর বা ধাপে দুটি option
নিতে হবে যা হবে \(\text{৪}_{P_{\text{২}}}\) অথবা = 8 ×৩, অতএব এটি
দু\textquotesingle ধাপ পর্যন্ত।

\textbf{বিঃদ্রঃ এখানে কোন উপাদান একাধিকবার নেই তাই এই সূত্র প্রয়োগ করা
হয়েছে।}

\paragraph{৬. A, B, C প্রতিবারে তিনটি করে বর্ণ একসাথে নিয়ে কত প্রকারে
সাজানো যায় নির্ণয় কর
৷}\label{a-b-c-ux9aaux9b0ux9a4ux9acux9b0-ux9a4ux9a8ux99f-ux995ux9b0-ux9acux9b0ux9a3-ux98fux995ux9b8ux9a5-ux9a8ux9af-ux995ux9a4-ux9aaux9b0ux995ux9b0-ux9b8ux99cux9a8-ux9afux9af-ux9a8ux9b0ux9a3ux9af-ux995ux9b0-ux9f7}

ক. ৩

খ. ৬

গ. ৯

ঘ. ১২

\textbf{\hl{উত্তর: খ}}

\textbf{সমাধান:} মোট বর্ণ n = 3 তিনটি বর্ণএকসাথে নিতে হবে যেমন: n = 3
এখানে r = 3

\(n_{P_{r}} = \frac{n!}{(n - r)!}\)

\(= \frac{3!}{(3 - 3)!}\)

\(= \frac{3!}{0!}\)

\(= \frac{3 \times 2 \times 1}{1}\) \(= 6\)

\paragraph{৭. A, B, C প্রতিবারে দুইটি বর্ণ একসাথে নিয়ে কত প্রকারে সাজানো যায়
নির্ণয় কর। {[}পুনরাবৃত্তি না
ঘটিয়ে{]}}\label{a-b-c-ux9aaux9b0ux9a4ux9acux9b0-ux9a6ux987ux99f-ux9acux9b0ux9a3-ux98fux995ux9b8ux9a5-ux9a8ux9af-ux995ux9a4-ux9aaux9b0ux995ux9b0-ux9b8ux99cux9a8-ux9afux9af-ux9a8ux9b0ux9a3ux9af-ux995ux9b0-ux9aaux9a8ux9b0ux9acux9a4ux9a4-ux9a8-ux998ux99fux9af}

\textbf{সমাধান:} এখানে মোট বর্ণ n = 3 আবার দু\textquotesingle টি করে
বর্ণএকসাথে নিতে হবে। তাই r = 2

\(n_{P_{r}} = \frac{n!}{(n - r)!}\)

\(= \frac{3!}{(3 - 2)!}\)

\(= \frac{3!}{1!}\)

\(= \frac{3 \times 2 \times 1}{1}\)

\(= 3 \times 2 \times 1\) \(= 6\) \textbf{Effective Shortcut:} এরকম
বিন্যাসের প্রশ্নগুলোতে যদি কখনো n এর মান এবং r এর মান সমান সমান হয় অথবা n এর
মানের থেকে r এর মান ১ কম হয় তাহলে নিচে কোন কিছু না লিখে সরাসরি উপরে n!
লিখে হিসেব করা যায় ।

কেননা নিচে 0! অথবা 1! = 1 হয় এবং নিচে 1 আসলে উপরের উত্তরটিই উত্তর হবে।
যেমন:

\paragraph{৮. Table, শব্দটিকে কতভাবে সাজানো যায়? {[} এভাবে বললে কোন শর্ত
না থাকলে বুঝতে হবে সবগুলো বর্ণ নিতে
হবে।{]}}\label{table-ux9b6ux9acux9a6ux99fux995-ux995ux9a4ux9adux9ac-ux9b8ux99cux9a8-ux9afux9af-ux98fux9adux9ac-ux9acux9b2ux9b2-ux995ux9a8-ux9b6ux9b0ux9a4-ux9a8-ux9a5ux995ux9b2-ux9acux99dux9a4-ux9b9ux9ac-ux9b8ux9acux997ux9b2-ux9acux9b0ux9a3-ux9a8ux9a4-ux9b9ux9ac}

ক. ১০০

খ. ১১০

গ. ১২০

ঘ. ১২৫

\textbf{\hl{উত্তর: গ}}

\textbf{সমাধান:}

5! = 5×4×3×2×1 = 120 (কারণ সুত্র প্রয়োগ করে কাটাকাটি করে 5! ই থাকবে)

অথবা সরাসরি 5! = 120 ।

\paragraph{৯. A, B, C, D চারটি বর্ণ । বর্ণ চারটি হতে পূনরাবৃত্তি না ঘটিয়ে
তিনটি করে বর্ণ নিয়ে সাজানোর ব্যবস্থা নির্ণয়
কর।}\label{a-b-c-d-ux99aux9b0ux99f-ux9acux9b0ux9a3-ux9acux9b0ux9a3-ux99aux9b0ux99f-ux9b9ux9a4-ux9aaux9a8ux9b0ux9acux9a4ux9a4-ux9a8-ux998ux99fux9af-ux9a4ux9a8ux99f-ux995ux9b0-ux9acux9b0ux9a3-ux9a8ux9af-ux9b8ux99cux9a8ux9b0-ux9acux9afux9acux9b8ux9a5-ux9a8ux9b0ux9a3ux9af-ux995ux9b0}

ক. ১২

খ. ১৬

গ. ১৮

ঘ. ২৪

\textbf{\hl{উত্তর: ঘ}}

\textbf{সমাধান:} মোট বর্ণ n = 4 প্রতিবারে নিতে হবে 3 টি বর্ণ r = 3 (এখানে
কতটি বর্ণ নিতে হবে তা বলে দেয়া আছে)

মোট সাজানোর ব্যবস্থা = \(n_{P_{r}} = \frac{n!}{(n - r)!}\)

\(= \frac{4!}{(4 - 3)!}\)

\(= \frac{4!}{1!}\)

\(= \frac{4 \times 3 \times 2 \times 1}{1}\)

\(= 4 \times 3 \times 2 \times 1\) \(= 24\) অথবা সরাসরি 4! = 24 ।

\paragraph{কিছু গুরুত্বপূর্ণ ফ্যাক্টোরিয়াল সংখ্যার মান মুখস্ত রাখলে খুব দ্রুত হিসেব
করা সম্ভব হবে। 0!= 1, 1!= 1, 2! = 2, 3! = 6, 4!= 24, 5!= 120, 6!= 720,
7!=
5040}\label{ux995ux99b-ux997ux9b0ux9a4ux9acux9aaux9b0ux9a3-ux9abux9afux995ux99fux9b0ux9afux9b2-ux9b8ux996ux9afux9b0-ux9aeux9a8-ux9aeux996ux9b8ux9a4-ux9b0ux996ux9b2-ux996ux9ac-ux9a6ux9b0ux9a4-ux9b9ux9b8ux9ac-ux995ux9b0-ux9b8ux9aeux9adux9ac-ux9b9ux9ac-0-1-1-1-2-2-3-6-4-24-5-120-6-720-7-5040}

\paragraph{১০. 2, 3, 4, 5, 6, 7, ও 8 এই অঙ্কগুলির প্রত্যেকটিকে প্রত্যেকের
সংখ্যায় একবার মাত্র ব্যবহার করে চার অঙ্কের কতগুলি পৃথক সংখ্যা গঠন করা যেতে
পারে।}\label{ux993-8-ux98fux987-ux985ux999ux995ux997ux9b2ux9b0-ux9aaux9b0ux9a4ux9afux995ux99fux995-ux9aaux9b0ux9a4ux9afux995ux9b0-ux9b8ux996ux9afux9af-ux98fux995ux9acux9b0-ux9aeux9a4ux9b0-ux9acux9afux9acux9b9ux9b0-ux995ux9b0-ux99aux9b0-ux985ux999ux995ux9b0-ux995ux9a4ux997ux9b2-ux9aaux9a5ux995-ux9b8ux996ux9af-ux997ux9a0ux9a8-ux995ux9b0-ux9afux9a4-ux9aaux9b0}

মোট সংখ্যা n = 7 নিতে হবে 4 টি করে অর্থাৎ r = 4

মোট সাজানোর ব্যবস্থা = \(n_{P_{r}}\)

\(= \frac{n!}{(n - r)!}\)

\(= \frac{7!}{(7 - 4)!}\)

\(= \frac{7!}{3!}\)

\(= \frac{7 \times 5 \times 6 \times 4 \times 3 \times 2 \times 1}{3 \times 2 \times 1}\)
\(= 7 \times 6 \times 5 \times 4\) \(= 840\)

\subsubsection{যখন কয়েকটি উপাদান একই
হয়}\label{ux9afux996ux9a8-ux995ux9afux995ux99f-ux989ux9aaux9a6ux9a8-ux98fux995ux987-ux9b9ux9af}

একই উপাদান একাধিকবার আসলে যেমন: DHAKA শব্দের বিন্যাস আগের নিয়মে হবে না,
কারণ এখানে দুটি A আছে। এরকম বিন্যাসের ক্ষেত্রে নিচের সুত্রটি প্রয়োগ করতে হয়।

যখন কোন উপাদান রিপিট হয় অর্থাৎ কয়েক বার আসে তখন এই সুত্র প্রয়োগ করতে হয় ।

Repetition Formula = \(\frac{n!}{p! \times q! \times r!}\) {[} এখানে n
হলো মোট উপাদান সংখ্যা আর p, q, r, হলো একাধিকবার ব্যবহৃত উপাদান সংখ্যা {]}

কিছু উদাহরণ দেখে আরো ক্লিয়ার হওয়া যাবে। নিচের শব্দগুলোর বিন্যাস সংখ্যা বের
করুন।

1) DHAKA

2) CANADA

3) MISSISSIPPI

1) DHAKA শব্দটিতে মোট 5 টি বর্ণআছে যার মধ্যে A আছে 2 টি এবং বাকীগুলো স্বতন্ত্র

নির্ণেয় বিন্যাস সংখ্যা = \(\frac{5!}{2!}\)

\(= \frac{5 \times 4 \times 3 \times 2 \times 1}{2 \times 1}\)

\(= 5 \times 4 \times 3\) \(= 60\) 2) CANADA শব্দটিতে মোট অক্ষর আছে ৬ টি
এর মধ্যে A আছে ৩ টি

সুতরাং নির্ণেয় বিন্যাস সংখ্যা = \(\frac{6!}{3!}\)

\(= \frac{6 \times 5 \times 4 \times 3 \times 2 \times 1}{3 \times 2 \times 1}\)

\(= 6 \times 5 \times 4\) \(= 120\) 3) MISSISSIPPI শব্দটিতে মোট অক্ষর আছে
১১ টি, এর মধ্যে S আছে ৪ টি, I আছে ৪ টি P আছে ২ টি

সুতরাং নির্ণেয় বিন্যাস সংখ্যা = \(\frac{11!}{4! \times 4! \times 2!}\)

\(= \frac{11 \times 10 \times 9 \times 8 \times 7 \times 6 \times 5 \times 4 \times 3 \times 2 \times 1}{4 \times 3 \times 2 \times 1 \times 4 \times 3 \times 2 \times 1 \times 2 \times 1}\)

\(= \frac{11 \times 10 \times 9 \times 8 \times 7 \times 6 \times 5}{4 \times 3 \times 2 \times 1 \times 2 \times 1}\)

\(= \frac{1,663,200}{48}\)

\(= 34,650\) {[}p, q, r, এর ব্যাবহার এভাবে হয় । {]}

\textbf{4) FREEDOM শব্দটির সবগুলোর বর্ণ একত্রে নিয়ে কত প্রকারের সাজানো যায়?}

FREEDOM শব্দটিতে মোট বর্ণ আছে 7টি এদের মধ্যে E আছে 2টি

.:. সবগুলো বর্ণ নিয়ে বিন্যাস = \(\frac{7!}{2!}\)

\textbf{5) CALCUTTA শব্দটির বর্ণগুলোকে একত্রে নিয়ে বিন্যাস সংখ্যা AMERICA
শব্দটির বর্ণগুলো একত্রে নিয়ে বিন্যাস সংখ্যার কত গুণ?}

AMERICA শব্দটিতে 7 টি বর্ণ আছে যাদের মধ্যে ২টি A রয়েছে।

সবগুলো বর্ণ একত্রে নিয়ে বিন্যাস সংখ্যা

= \(\frac{7!}{2!}\)

\(\frac{7 \times 6 \times 5 \times 4 \times 3 \times 2 \times 1}{2 \times 1}\)

\(\frac{5040}{2}\)

= 2520 আবার CALCUTTA শব্দটিতে ৮ টি বর্ণ আছে যাদের মধ্যে A, C ও T 2 টি করে
বিদ্যমান ।

সবগুলো বর্ণ একত্রে নিয়ে বিন্যাস সংখ্যা

= \(\frac{8!}{2! \times 2! \times 2!}\)

=
\(\frac{8 \times 7 \times 6 \times 5 \times 4 \times 3 \times 2 \times 1}{2 \times 1 \times 2 \times 1 \times 2 \times 1}\)

=
\(\frac{8 \times 7 \times 6 \times 5 \times 4 \times 3}{2 \times 1 \times 2 \times 1}\)

= \(\frac{20160}{4}\)

= 5040 সুতরাং AMERICA শব্দটির বিন্যাস থেকে CALCUTTA শব্দটির বিন্যাস সংখ্যা
৫০৪০÷২৫২০ = ২ গুণ বেশি ।

\subsubsection{নির্দিষ্ট কোন উপাদান নির্দিষ্ট স্থানে অথবা পাশাপাশি রাখতে
বললে}\label{ux9a8ux9b0ux9a6ux9b7ux99f-ux995ux9a8-ux989ux9aaux9a6ux9a8-ux9a8ux9b0ux9a6ux9b7ux99f-ux9b8ux9a5ux9a8-ux985ux9a5ux9ac-ux9aaux9b6ux9aaux9b6-ux9b0ux996ux9a4-ux9acux9b2ux9b2}

\textbf{প্রশ্ন নম্বর-} ১১. Apu শব্দটির বর্ণগুলো নিয়ে কতগুলো বিন্যাস সংখ্যা নির্ণয়
করা যায় যেন প্রত্যেক বিন্যাসের প্রথমে Vowel থাকে? ..

\begin{itemize}
\item
  (ক) ৪
\item
  (খ) ৬
\item
  (গ) ৭
\item
  (ঘ) ৮
\end{itemize}

\textbf{\hl{উত্তরঃ (ক) ৪}}

\textbf{Explanation:}

এখানে Apu শব্দটিতে মোট তিনটি বর্ণ আছে। আবার কোনটিই দুবার নেই ।

তাই এদের বিন্যাস সংখ্যা হবে । ৩! বা ৩×২ = ৬টি।

এখন বাস্তবে ভাবুন। তিনটি অক্ষর দিয়ে যদি মোট ৬টি বিন্যাস করা যায়।

তাহলে প্রতি ১ টি দিয়ে ৬÷৩ = ২টি করে বিন্যাস সাজানো যায়।

আবার Apu শব্দটিতে যেহেতু দুটি স্বরবর্ণ যথা: A ও U দিয়ে মোট বিন্যাস হবে ২×২ =
৪টি।

যেমন: Apu, Aup, upa, uap তাই উত্তর: ৪।

\textbf{কিন্তু এই প্রশ্নটিতেই যদি বলা হত প্রথমে Vowel থাকবে না এরকম বিন্যাস
সংখ্যা কয়টি হবে?}

তখন উত্তর হত ২টি যে দুটি P দিয়ে শুরু হয়। যেমন: pau, pua.

\textbf{প্রশ্ন নম্বর-} ১২. Cap শব্দটির বর্ণগুলো নিয়ে কতগুলো বিন্যাস সংখ্যা নির্ণয়
করা যায় যেন প্রত্যেক বিন্যাসের প্রথমে একটি স্বরবর্ণ থাকে? ...

\begin{itemize}
\item
  (ক) ৪
\item
  (খ) ৬
\item
  (গ) ৭
\item
  (ঘ) ২
\end{itemize}

\textbf{\hl{উত্তরঃ (ঘ) ২}}

\textbf{Explanation:}

এখানে Cap শব্দটিতে মোট তিনটি বর্ণ আছে। আবার কোনটিই দুবার নেই । তাই এদের
বিন্যাস সংখ্যা হবে । ৩! বা ৩×২ = ৬টি। এখন বাস্তবে ভাবুন। তিনটি অক্ষর দিয়ে যদি
মোট ৬টি বিন্যাস করা যায়।

তাহলে প্রতি ১ টি দিয়ে ৬÷৩ = ২টি করে বিন্যাস সাজানো যায়। এখানে স্বরবর্ণ ও ১
টি । তাই প্রথমে একটি স্বরবর্ণ রেখে ৬÷৩ = ২টি করে বিন্যাস সাজানো যাবে।

\textbf{প্রশ্ন নম্বর-} ১৩. ২৭. Nokia\textquotesingle{} শব্দটির বর্ণগুলো নিয়ে
কতগুলো বিন্যাস সংখ্যা নির্ণয় করা যায় যাদের শুরুতে Consonant থাকে? (উল্টো)

\begin{itemize}
\item
  (ক) ৪৮
\item
  (খ) ২৬
\item
  (গ) ২৭
\item
  (ঘ) ২৮
\end{itemize}

\textbf{\hl{উত্তরঃ (ক) ৪৮}}

\textbf{Explanation:}

Nokia শব্দটিতে মোট ৫টি বর্ণ আছে তাই ৫! = ১২০ ভাবে সাজানো যায়। তাহলে ৫
অক্ষরের প্রতিটি দিয়ে সাজানো যাবে ১২০÷৫ = ২৪ টি করে।

এখন Nokia শব্দটিতে যেহেতু দুটি Consonant আছে, তাই Consonant দিয়ে শুরু হবে মোট
২৪×২ = ৪৮টি বিন্যাস। উত্তর: ৪৮টি।

(বড় বড় বিন্যাস আসলেও নিয়ম একই)

\textbf{প্রশ্ন নম্বর- ১৪. Courage\textquotesingle{} শব্দটির বর্ণগুলো নিয়ে
কতগুলো বিন্যাস সংখ্যা নির্ণয় করা যায় যেন প্রত্যেক বিন্যাসের প্রথমে একটি স্বরবর্ণ
থাকে?}

\begin{itemize}
\item
  (ক) ১৭২০
\item
  (খ) ২৮৮০
\item
  (গ) ৩৬৪০
\item
  (ঘ) কোনটিই নয়
\end{itemize}

\textbf{\hl{উত্তরঃ (খ) ২৮৮০}}

\textbf{Explanation:}

এখানে Courage শব্দটিতে মোট সাতটি বর্ণ আছে। আবার কোনটিই দুবার নেই । তাই এদের
বিন্যাস সংখ্যা হবে । ৭! = ৭×৬×৫×৪×৩×২×১ = ৫০৪০ টি। এখন বাস্তবে ভাবুন। সাতটি
অক্ষর দিয়ে যদি মোট ৫০৪০ টি বিন্যাস করা যায়।

তাহলে প্রতি ১ টি দিয়ে ৫০৪০÷৭ = ৭২০ টি করে বিন্যাস সাজানো যায়।

আবার Courage শব্দটিতে যেহেতু চারটি স্বরবর্ণ যথা: O , U , A ও E দিয়ে মোট
বিন্যাস হবে ৭২০×৪ = ২৮৮০ টি।

\subsubsection{বিভিন্ন শর্ত প্রয়োগ করে
বিন্যাস}\label{ux9acux9adux9a8ux9a8-ux9b6ux9b0ux9a4-ux9aaux9b0ux9afux997-ux995ux9b0-ux9acux9a8ux9afux9b8}

\textbf{প্রশ্ন নম্বর- ১৫ . `CALCULUS' শব্দটির বর্ণগুলোর সবগুলো একত্রে নিয়ে কত
প্রকারে সাজানো যায় যেন প্রথম ও শেষ অক্ষর `u' থাকে?}

\begin{itemize}
\item
  (ক) 120
\item
  (খ) 180
\item
  (গ) 60
\item
  (ঘ) 720
\end{itemize}

\textbf{\hl{উত্তরঃ (খ) 180}}

\textbf{Explanation:}

`CALCULUS' শব্দটির মধ্যে মোট ৮টি অক্ষর আছে । শর্তানুযায়ী প্রথম ও শেষে u থাকবে
তাহলে এদেরকে বিন্যাসের বাইরে রাখতে হবে। সুতরাং অবশিষ্ট ৬টি স্থানে বাকি ৬টি
অক্ষর দ্বারা পূরণ করতে হবে।

যেহেতু বাকি ৬টি অক্ষরের মধ্যে ২টি c, ২টি I এবং অন্যগুলো ভিন্ন, ভিন্ন, সুতরাং ৬টি
অক্ষরের সবগুলো একত্রে নিয়ে বিন্যাস সংখ্যা = \(\frac{6!}{2!2!} = 180\)

:: প্রদত্ত শর্ত অনুযায়ী অক্ষরগুলোকে 180 প্রকারে সাজানো যাবে ।

\textbf{প্রশ্ন নম্বর-} ১৬. Admission শব্দটির A ও d কে দুই প্রান্তে রেখে কত
প্রাকারে সাজানো যেতে পারে? ...

\begin{itemize}
\item
  (ক) 240
\item
  (খ) 2720
\item
  (গ) 2520
\item
  (ঘ) 2620
\end{itemize}

\textbf{\hl{উত্তরঃ (গ) 2520}}

\textbf{Explanation:}

A ও d কে যেহেতু দু প্রান্তে রাখতে হবে তাই এই দুটি সংখ্যা বাদ দিয়ে অন্য 7 টি
অক্ষর বিন্যাস করা যায় = \(\frac{7!}{2!2!}\) = 1260 ভাবে।

আবার, A ও d কে নিজেদের মধ্যে বিন্যাস করা যায় = 2! = 2 ভাবে ।

{[}কারণ শুরুতে যেমন: A এবং শেষে d আসতে পারে ঠিক তেমনি শেষে A এবং শুরুতে d
দেয়া যায়। অর্থাৎ তাদের নিজেদের

মধ্যে দুভাবে বিন্যাস হবে{]}

সুতরাং সর্বমোট সাজানো সংখ্যা = 1260×2 = 2520

\textbf{প্রশ্ন নম্বর-} ১৭. PERMUTATION শব্দটি Vowel গুলোর অবস্থান পরিবর্তন না
করে কত প্রকারে পুণরায় সাজানো যায়? ..

\begin{itemize}
\item
  (ক) 360
\item
  (খ) 359
\item
  (গ) 355
\item
  (ঘ) 361
\end{itemize}

\textbf{\hl{উত্তরঃ (খ) 359}}

\textbf{Explanation:}

অবস্থান পরিবর্তন করা যাবে না বলতে বোঝায় Vowel গুলো যে জায়গায় আছে সেই
জায়গাতেই রেখে দিতে হবে । তাহলে Vowel গুলো সূত্রের বাইরে রাখলে তাদের অবস্থান
পরিবর্তন হবে না ।

এখন শব্দটিতে Vowel = 5টি, Consonant = 6টি। ৬টির বিন্যাস করতে হবে কিন্তু এদের
মধ্যে T আছে ২ বার

সুতরাং \(\frac{6!}{2!}\) = 360 উপায়ে গঠন করা যেতে পারে।

কিন্তু PERMUTATION শব্দটি নিজেই একটি সাজানো সংখ্যা। ( কেননা ঐ ৩৬০টি বিন্যাসের
মধ্যে Vowel

গুলো অবস্থান স্থির রেখে যতগুলো বিন্যাস হয় তার মধ্যে এই PERMUTATION শব্দটিও
একটি। যেহেতু প্রশ্নে পুণরায় সাজানোর সংখ্যা জানতে চাওয়া হয়েছে তাই একে নেয়া
যাবে না। )

সুতরাং নির্ণেয় সাজানো সংখ্যা ( 360 -- 1 ) = 359

\textbf{প্রশ্ন নম্বর- ১৮. Cambridge শব্দটির বর্ণগুলো থেকে ৫টি বর্ণ নিয়ে শব্দ গঠন
করলে কতগুলোতে প্রদত্ত শব্দটির সবগুলো স্বরবর্ণ থাকবে?}

\begin{itemize}
\item
  (ক) ১২০০
\item
  (খ) ১৪০০
\item
  (গ) ১৬০০
\item
  (ঘ) ১৮০০
\end{itemize}

\textbf{\hl{উত্তরঃ (ঘ) ১৮০০}}

\textbf{Explanation:}

Cambridge শব্দটিতে মোট বর্ণ আছে ৯টি যেখানে স্বরবর্ণ ৩টি (a,i,e) এবং
ব্যাঞ্জনবর্ণ = ৬টি (c,m,b,r,d,g)

সুতরাং ৫ টি স্থানের মধ্যে ৩ টি স্থান স্বরবর্ণ দ্বারা পূরণ করার উপায় =
\(5_{P_{3}}\) = 60 ভাবে । অবশিষ্ট ৬ টি বর্ণ দ্বারা ২ টি স্থান পূরণ করা যায় =
\(6_{P_{2}}\)

= 30 ভাবে ।

.:. ৫টি বর্ণ নিয়ে গঠন করা শব্দগুলোর মধ্যে Cambridge শব্দের সবগুলো স্বরবর্ণ
(a,i,e) থাকবে এরূপ শব্দের সংখ্যা = 60×30 =1800

\textbf{প্রশ্ন নম্বর-} ১৯. `SECOND' শব্দটির অক্ষরগুলি থেকে ১টি স্বরবর্ণ ও ২টি
ব্যঞ্জনবর্ণ নিয়ে কতকগুলি শব্দ গঠন করা যেতে পারে, যাতে স্বরবর্ণ সর্বদা মধ্যম স্থান
দখল করবে? ..

\begin{itemize}
\item
  (ক) 12
\item
  (খ) 18
\item
  (গ) 24
\item
  (ঘ) 30
\end{itemize}

\textbf{\hl{উত্তরঃ (গ) 24}}

\textbf{Explanation:}

`SECOND' শব্দটি 2 টি স্বরবর্ণ এবং 4 টি ব্যঞ্জনবর্ণ নিয়ে গঠিত। নতুন যে শব্দ গুলো
গঠিত হবে তার মধ্যে 1 টি স্বরবর্ণ এবং 2টি ব্যাঞ্জণবর্ণ থাকতে হবে।

যেহেতু স্বরবর্ণটি মধ্যস্থানে থাকতে হবে, সুতরাং 2টি স্বরবর্ণ হতে তা \(2_{P_{1}}\)
=2!=2 ভাবে নেয়া যেতে পারে। পুনরায় 4টি ব্যঞ্জনবর্ণ হতে 2টি ব্যাঞ্জণবর্ণ
\(4_{P_{2}}\) = 4.3 = 12 ভাবে নেয়া যেতে পারে ।

সুতরাং নির্ণেয় গঠিত সংখ্যা = 2×12 = 24

\textbf{প্রশ্ন নম্বর-} ২০. `EQUATION' শব্দটির অক্ষরগুলি থেকে চার অক্ষরবিশিষ্ট
বিভিন্ন শব্দ গঠন করা হলো ; এদের কতগুলিতে q বর্তমান থাকবে কিন্তু n থাকবে না ?
...

\begin{itemize}
\item
  (ক) 280
\item
  (খ) 360
\item
  (গ) 480
\item
  (ঘ) 520
\end{itemize}

\textbf{\hl{উত্তরঃ (গ) 480}}

\textbf{Explanation:}

EQUATION\textquotesingle{} শব্দটিতে ৪টি বিভিন্ন বর্ণ রয়েছে। আমাদেরকে 4টি
বিভিন্ন বর্ণের শব্দ এমনভাবে তৈরি করতে হবে যেন

গঠিত শব্দগুলোতে একটি নির্দিষ্ট বর্ণ (অর্থাৎ N) থাকবে না। সুতরাং আমাদের ব্যবহার্য
বিভিন্ন বর্ণের মোট সংখ্যা (8-1)=7টি

ধরা যায়। অপর দিকে আর একটি নির্দিষ্ট বর্ণ (অর্থাৎ Q) প্রতিটি বিন্যাসেই রাখতে
হবে।

এখন 4টি স্থানের মধ্যে একটি স্থানে একটি নির্দিষ্ট বর্ণ দ্বারা \(4_{P_{1}}\)
উপায়ে পূর্ণ করা যায়। অবশিষ্ট 3টি স্থান (7-1) = 6টি বর্ণ দ্বারা

\(6_{P_{3}}\) উপায়ে পূর্ণ করা যায় ।

.:. নির্ণেয় সাজানোর সংখ্যা = \(4_{P_{1}}\) × \(6_{P_{3}}\) = 4×6×5×4 = 480

\textbf{প্রশ্ন নম্বর-} ২১. AMERICA শব্দটির বর্ণগুলো থেকে প্রতিবারে ৩টি বর্ণ নিয়ে
গঠিত ভিন্ন ভিন্ন শব্দ সংখ্যা কত হবে? {[}থানা শিক্ষা অফি: ৯৯{]}

\begin{itemize}
\item
  (ক) ১৩০
\item
  (খ) ১৩৫
\item
  (গ) ১৪০
\item
  (ঘ) ১৪৫
\end{itemize}

\textbf{\hl{উত্তরঃ (খ) ১৩৫}}

\textbf{Explanation:}

যেহেতু AMERICA শব্দটিতে 7 টি বর্ণ রয়েছে, যার মধ্যে A দুইটি ।

একটি A বাদ দিয়ে 6টি ভিন্ন বর্ণ থেকে প্রতিবারে 3টি বর্ণ নিয়ে বিন্যাস সংখ্যা =
\(6_{P_{3}}\) = 6 × 5 × 4 = 120

আবার, দুইটি A কে ভিন্ন ভিন্ন পাঁচটি বর্ণের প্রতিটির সাথে নিলে 3 বর্ণ শব্দ সংখ্যা
= \(5_{P_{1}} \times \frac{3!}{2!}\) = 5 × 3 = 15

মোট শব্দ সংখ্যা = 120+15 = 135

\subsubsection{পাশাপাশি রেখে
বিন্যাস}\label{ux9aaux9b6ux9aaux9b6-ux9b0ux996-ux9acux9a8ux9afux9b8}

যখন কয়েকটি বর্ণ অথবা বস্তুকে একসাথে পাশাপাশি রাখতে বলা হয় তখন বিন্যাস করার
জন্য নিচের স্টেপগুলো অনুসরণ করতে হবে।

প্রথম কাজ: একসাথে রাখতে বলা বর্ণগুলোকে ১টি ধরে সবগুলো বর্ণের বিন্যাস বের করতে
হবে।

দ্বিতীয় কাজ: ঐ একসাথে রাখতে বলা বর্ণগুলোর মধ্যেই আবার বিন্যাস বের করতে হবে ।

শেষ কাজ : প্রথম দুটি কাজ থেকে প্রাপ্ত বিন্যাসের ফলাফল দুটি গুণ করতে হবে।

\textbf{প্রশ্ন নম্বর- ২২. Vowel গুলি একসাথে রেখে ACCLAIM শব্দটিকে কতভাবে
সাজানো যাবে?}

\begin{itemize}
\item
  (ক) 720
\item
  (খ) 120
\item
  (গ) 60
\item
  (ঘ) 180
\end{itemize}

\textbf{\hl{উত্তরঃ (ঘ) 180}}

\textbf{Explanation:}

প্রথম কাজ : Vowel গুলিকে একত্রে রেখে বিন্যাস করা অর্থাৎ CCLM (AAI) = (4+1)! =
5! কিন্তু দুটি C থাকায়

নিচে ভাগ করতে হবে 2! দিয়ে অর্থাৎ = \(\frac{5!}{2!}\) = 60

দ্বিতীয় কাজ: (AAI) Vowel গুলির মধ্যেই বিন্যাস করা এখানে Vowel আছে তিনটি কিন্তু
তাদের মধ্যে A আছে দুটি তাই Vowel গুলির মধ্যে বিন্যাস সংখ্যা হবে =
\(\frac{3!}{2!}\) = 3

তৃতীয় এবং শেষ কাজ: দুই বিন্যাসের গুণফল বের করতে হবে। অর্থাৎ 60×3 = 180। উত্তর:
180 ।

\textbf{প্রশ্ন নম্বর-} ২৩. SCIENCE শব্দটির স্বরবর্ণ গুলোকে একত্রে রেখে সব কয়টি
বর্ণকে সম্ভাব্য যত উপায়ে সাজানো যায় তার সংখ্যা নির্ণয় কর।

\begin{itemize}
\item
  (ক) ১৪০
\item
  (খ) ১৭৬
\item
  (গ) ১৭৭
\item
  (ঘ) ১৮০
\end{itemize}

\textbf{\hl{উত্তরঃ (ঘ) ১৮০}}

\textbf{Explanation:}

SCIENCE শব্দটিতে বর্ণসংখ্যা ৭টি

স্বরবর্ণ তিনটিকে একত্রে রেখে একটি বর্ণ বিবেচনা করলে অর্থাৎ SCNC (EEI)
বর্ণসংখ্যা হবে ৫টি

সুতরাং SCNC (EEI) এর সাজানোর সংখ্যা = মোট ৫টি বর্ণ যেখানে দুটি C =
\(\frac{5!}{2!}\) =
\(\frac{5 \times 4 \times 3 \times 2 \times 1}{2 \times 1}\) = 60

আবার (EEI) এর বিন্যাস সংখ্যা = \(\frac{3!}{2!}\) {[}যেহেতু E = 2 টি {]} =
\(\frac{3 \times 2 \times 1}{2 \times 1}\) = 3

{[}যুক্তি, তিনটি স্বরবর্ণ পাশাপাশি রাখলেই হল, তাই এই স্বরবর্ণ গুলো EEI, EIE,
অথবা IEE, এভাবে আসলেও শর্ত

পূরণ হবে। এজন্য দুবার বিন্যাস করতে হল। {]}

অতএব নির্ণেয় বিন্যাস সংখ্যা = 60 × 3 = 180

\textbf{প্রশ্ন নম্বর-} ২৪. ARRANGE\textquotesingle{} শব্দটিকে কতভাবে সাজানো
যাবে যেখানে দুটো R এবং দুটো A একত্রে থাকবে? ..

\begin{itemize}
\item
  (ক) 620
\item
  (খ) 120
\item
  (গ) 200
\item
  (ঘ) 180
\end{itemize}

\textbf{\hl{উত্তরঃ (খ) 120}}

\textbf{Explanation:}

R দুটিকে এবং A টিকে একত্রে রেখে মোট বিন্যাস করা যাবে (AA) (RR)NGE = 5! =
120টি। এখানে কোন গুণ অথবা ভাগ করতে হবে না।

এর কারণ হলো A দুটিকে নিজেদের মাঝে মাত্র ১ ভাবেই বিন্যাস করা যায়, তেমনিভাবে
R দুটি কেও।

আবার দুটি সংখ্যাকে আলাদা আলাদা ভাবে ১টি করে ধরায় তারা রিপিট হয় নি। তাই
নিচে ভাগও হবে না ।

\textbf{প্রশ্ন নম্বর- ২৫. ৩ জন ছাত্র ও ৫ জন ছাত্রীকে একসারিতে রেখে কতভাবে
সাজানো যায় যেখানে ৩ জন ছাত্র সর্বদা একত্রে থাকবে?}

\begin{itemize}
\item
  (ক) ১২৪০
\item
  (খ) ১১৭৬
\item
  (গ) ৪১৭৭
\item
  (ঘ) ৪৩২০
\end{itemize}

\textbf{\hl{উত্তরঃ (ঘ) ৪৩২০}}

\textbf{Explanation:}

৩জন ছাত্রকে যেহেতু একত্রে রাখতে হবে তাই তাদেরকে ১ জন ধরে এবং ৫ জন ছাত্রী ও ১
জন ছাত্র নিয়ে মোট বিন্যাস = ৬! বা ৭২০ এখন আগের ৩ জন ছাত্রকে তাদের নিজেদের
মধ্যে সাজানো যায় ৩! = ৬ ভাবে তাহলে মোট সাজানোর সংখ্যা ৭২০×৬ = ৪৩২০

\subsubsection{একসাথে না রেখে
বিন্যাস}\label{ux98fux995ux9b8ux9a5-ux9a8-ux9b0ux996-ux9acux9a8ux9afux9b8}

\textbf{মনে রাখুন:}

একসাথে না রেখে বিন্যাস করার সরাসরি কোন সুত্র নেই তাই একত্রে রাখা যাবে না
বললে:

সর্বমোট বিন্যাস - একসাথে রাখা বিন্যাস = একসাথে না রাখা বিন্যাস ।

প্রথমে: সবগুলো বর্ণের সাধারণ বিন্যাস বের করতে হবে। যাতে একত্রে থাকা এবং একত্রে
না থাকা সব ধরণের বিন্যাস থাকবে।

এরপর একসাথে রাখার নিয়ম অনুযায়ী সমাধান বের করে মোট বিন্যাস থেকে বিয়োগ করলেই
একসাথে না থাকার বিন্যাস বের হবে।

\textbf{প্রশ্ন নম্বর-} ২৬. ``APPLE'' শব্দটিকে কতভাবে বিন্যাস করা যাবে যাতে P
গুলো কখনোই একসাথে থাকবে না? ...

\begin{itemize}
\item
  (ক) 24
\item
  (খ) 36
\item
  (গ) 48
\item
  (ঘ) 60
\end{itemize}

\textbf{\hl{উত্তরঃ (খ) 36}}

\textbf{Explanation:}

``APPLE'' শব্দটির সর্বমোট বিন্যাস সংখ্যা = 5!

কিন্তু P আছে দুটি তাই বিন্যাস হবে \(\frac{5!}{2!}\) = 60

এখন এই ৬০টি বিন্যাসের মধ্যে কিছু বিন্যাস আছে যেখানে P দুটি একত্রে থাকবে আবার
কিছু বিন্যাস আছে যেখানে P দুটি

একত্রে থাকবে না । যেহেতু সরাসরি একত্রে না থাকার সুত্র নেই তাই

``APPLE'' শব্দটির P দুটি একত্রে রেখে বিন্যাস = ALE (PP) = 4!×1! = 24

সুতরাং P দুটি একত্রে না রেখে বিন্যাস = 60-24 = 36

\textbf{প্রশ্ন নম্বর-} ২৭. Vowel গুলি পাশাপাশি না রেখে
`Triangle\textquotesingle{} শব্দটির অক্ষরগুলো কতভাবে সাজানো যাবে? ...

\begin{itemize}
\item
  (ক) 35000
\item
  (খ) 36000
\item
  (গ) 35600
\item
  (ঘ) 36800
\end{itemize}

\textbf{\hl{উত্তরঃ (খ) 36000}}

\textbf{Explanation:}

শব্দটিতে মোট ভিন্ন ৪টি অক্ষর যার 3টি স্বরবর্ণ। সবগুলো একত্রে নিয়ে মোট,

8! = 8×7×6×5×4×3×2×1 = 40320 রকমে সাজানো যায় ।

স্বরবর্ণগুলো পাশাপাশি রেখে, মোট 6!×3 ! = 720 × 6 = 4320 প্রকারে সাজানো যায় ৷

সুতরাং স্বরবর্ণগুলো পাশাপাশি না রেখে মোট (40320 - 4320 ) = 36000 প্রকারে
সাজানো যায় ।

\textbf{প্রশ্ন নম্বর-} ২৮. স্বরবর্ণগুলোকে একত্রে না রেখে ``Daughter'' শব্দটিকে
কতভাবে সাজানো যাবে? ...

\begin{itemize}
\item
  (ক) 37000
\item
  (খ) 36000
\item
  (গ) 36500
\item
  (ঘ) 36400
\end{itemize}

\textbf{\hl{উত্তরঃ (খ) 36000}}

\textbf{Explanation:}

শব্দটিতে মোট ভিন্ন 8 টি অক্ষর যার 3টি স্বরবর্ণ। সবগুলো একত্রে নিয়ে মোট,

8! = 8×7×6×5×4×3×2×1 = 40320 রকমে সাজানো যায় ।

স্বরবর্ণগুলো একত্রে রেখে, মোট 6!×3 ! = 720 × 6 = 4320 প্রকারে সাজানো যায় ৷

সুতরাং স্বরবর্ণগুলো একত্রে না রেখে মোট (40320 - 4320 ) = 36000 প্রকারে সাজানো
যায় ।

\subsubsection{পুণরাবৃত্তির
বিন্যাস}\label{ux9aaux9a3ux9b0ux9acux9a4ux9a4ux9b0-ux9acux9a8ux9afux9b8}

উপরের প্রশ্নগুলোতে যে কোন সংখ্যা বা অক্ষর শুধুমাত্র ১ বার ব্যবহার করা হয়েছে।
অর্থাৎ একই সংখ্যা বা অক্ষর একাধিকবার

ব্যবহৃত হয় নি।

\textbf{যেমন: ১ ও ২ কে একবার মাত্র ব্যবহার করে, দু \textquotesingle{} অঙ্কের
কতগুলো সংখ্যা গঠন করা যায়?}

এরকম প্রশ্নের উত্তর ২! বা ২টি যথা: ১২ এবং ২১ কিন্তু এই একই প্রশ্নে repetition
allowed বা পুণরাবৃত্তি করা গেলে ১ ও ২ কে ব্যবহার করে ২ অঙ্কের সংখ্যা বানানো
যাবে = \(\text{২}^{\text{২}}\) = ৪ টি । যথা: ১২, ২১, এর সাথে ১১ এবং ২২ {[} অর্থাৎ একই
সংখ্যাকে একাধিকবার ব্যাবহার করা যাবে। {]}

Formula of Repetition = \(n^{r}\) {[} এখানে n হচ্ছে মোট উপাদান এবং r =
যতবার নিতে হবে। {]}

পূনরাবৃত্তি করে A, B, C তিনটি উপাদান থেকে কয়ভাবে ২টি উপাদান নেয়া যাবে?
এখানে, সকল বিন্যাস হবে এরূপ, AA, AB, AC, BA, BB, BC, CA, CB, CC, 9টি।
কেননা প্রতি ক্ষেত্রেই প্রতি ধাপে আগের সব options থেকে যায়।

এক্ষেত্রে বিন্যাস সংখ্যা \(n^{r} = 3^{2} = 9\) । অর্থাৎ এক বর্ণ রিপিট করা গেলে
এভাবে।

\textbf{প্রশ্ন নম্বর-} ২৯. ১, ২, ৩, ৪, ৫ অঙ্কগুলির প্রতিটিকে যে কোন সংখ্যক বার
নিয়ে ৩ অঙ্কের কতগুলি সংখ্যা গঠন করা যাবে? ...

\begin{itemize}
\item
  (ক) ১৩০
\item
  (খ) ১৩৫
\item
  (গ) ১২৫
\item
  (ঘ) ১২৭
\end{itemize}

\textbf{\hl{উত্তরঃ (গ) ১২৫}}

\textbf{Explanation:}

যে কোন সংখ্যক বার অর্থ প্রতিবারই ১ নিয়ে ৩ অংকের সংখ্যা হবে ১১১ অর্থাৎ এখানে
রিপিট করা যাবে। সুত্র: এরকম পুনরাবৃত্তি করা গেলে (মোট উপাদান) যতবার নেয়া যাবে

এখানে মোট উপাদান ১,২,৩,৪ ও ৫ = ৫টি এবং সংখ্যা বানাতে হবে ৩ অংকের ।

তাই উত্তর হবে \(\text{৫}^{\text{৩}}\) = ১২৫টি (যেমন: ১২৩, ১১২, ১১১, এরকম ১ দিয়ে ২৫টি সহ
প্রতিটি সংখ্যা দিয়ে ২৫ টি করে মোট ১২৫ টি)

\textbf{প্রশ্ন নম্বর-} ৩০. 4, 5,6,7,8 এর প্রত্যেকটিকে যে কোন সংখ্যক বার নিয়ে
চার অঙ্কবিশিষ্ট কতগুলি সংখ্যা গঠন করা যায়? এ সংখ্যাগুলোর কয়টিতে একই অঙ্ক
একাধিবার থাকবে? ...

\begin{itemize}
\item
  (ক) 520
\item
  (খ) 428
\item
  (গ) 625
\item
  (ঘ) 505
\end{itemize}

\textbf{\hl{উত্তরঃ (ঘ) 505}}

\textbf{Explanation:}

চার অঙ্কবিশিষ্ট মোট সংখ্যা = \(5^{4}\) = 625. ( যে কোন অঙ্ক একাধিকবার নেয়া
যাবে।)

আবার 5টি অঙ্ক থেকে 4 টি অঙ্ক (প্রত্যেকটি কেবল একবার)

নিয়ে 4 অঙ্কবিশিষ্ট মোট সংখ্যা = \(5_{P_{4}}\) = 120 ( এগুলোতে একটি অঙ্ক মাত্র
একবারই থাকবে।)

সুতরাং প্রত্যেকটি সংখ্যায় একই অঙ্ক একাধিকবার থাকবে এরূপ মোট সংখ্যা: 625 - 120 =
505

\textbf{প্রশ্ন নম্বর- ৩১. ৫, ৫, ৬, ৬, ৭, ৭ সংখ্যাগুলি থেকে ৩ অংকের কতগুলি
সংখ্যা গঠন করা যাবে?}

\begin{itemize}
\item
  (ক) ২৪
\item
  (খ) ২৬
\item
  (গ) ২৮
\item
  (ঘ) ২২
\end{itemize}

\textbf{\hl{উত্তরঃ (ক) ২৪}}

\textbf{Explanation:}

এখানে যে কোন সংখ্যা ইচ্ছামত নেয়া যাবে না । তবে প্রশ্নে প্রদত্ত সংখ্যাগুলোকেই শুধু
ব্যবহার করা যাবে, যেখানে মোট ৬টি সংখ্যা দেয়া থাকলেও আসলে সংখ্যা (উপাদান)
মোট ৩টি এবং প্রতিটি ২ বার করে আছে।

তাই আমরা লিখতে পারি, ৩টি উপাদান দিয়ে ৩ অংকের সংখ্যা বানাতে হলে লিখতে হবে
\(\text{৩}^{\text{৩}}\) = ২৭টি । কিন্তু লক্ষ্য করুন ৩টি উপাদান কিন্তু ইচ্ছে মত রিপিট করা যাবে
না। বরং যে কয়টি সংখ্যা দেয়া আছে তা থেকেই নিতে হবে।

এখানে ৫, ৬ এবং ৭ আছে ২টি করে।

কিন্তু ২৭টি সংখ্যার মধ্যে এমন ৩টি সংখ্যা আছে যেখানে ৫৫৫, ৬৬৬, এবং ৭৭৭ আছে
যেগুলো নেয়া যাবে না। কারণ প্রশ্নে ৩টি করে সংখ্যা দেয়া নেই । তাই এই ৩টি বাদ
দিলে মোট সংখ্যা হবে ২৭-৩ = ২৪টি ।

\textbf{প্রশ্ন নম্বর- ৩২. 1, 2, 4, 6, 8, 9 অঙ্কগুলো যে কোন সংখ্যকবার নিয়ে তিন
অঙ্ক বিশিষ্ট কতগুলো সংখ্যা গঠন করা যাবে?}

\begin{itemize}
\item
  (ক) 200
\item
  (খ) 216
\item
  (গ) 400
\item
  (ঘ) 560
\end{itemize}

\textbf{\hl{উত্তরঃ (খ) 216}}

\textbf{Explanation:}

এখানে সর্বমোট সংখ্য n = 6, প্রতিবার অঙ্ক নিতে হবে r = ৩ টি

সুতরাং সংখ্যা গঠন করা যাবে \(n^{3} = 6^{3}\) = 6×6×6 = 216

\textbf{প্রশ্ন নম্বর-} ৩৩. ০ বাদে তিন অঙ্কবিশিষ্ট কতগুলো পূর্ণসংখ্যা তৈরী করা যায়
যেখানে কোন অঙ্ক দু\textquotesingle য়ের অধিকবার ব্যবহৃত হবে না? {[}BKB
(officer) - 2017{]}

\begin{itemize}
\item
  (ক) 729
\item
  (খ) 720
\item
  (গ) 756
\item
  (ঘ) 504
\end{itemize}

\textbf{\hl{উত্তরঃ (খ) 720}}

\textbf{Explanation:}

0 বাদে বাকী ডিজিটগুলোকে যে কোন সংখ্যক বার ব্যাবহার করে তিন অঙ্কবিশিষ্ট সংখ্যা
বানানো যায় তাদের মোট সংখ্যা হবে \(9^{3}\) = 729 {[} তিনটি অঙ্কের জন্য তিনটি
পজিশনে প্রতিটিই রিপিট করা যাবে তাই 9×9×9 = \(9^{3}\) {]}

কিন্তু কোন সংখ্যাকে দুবারের বেশি নেয়া যাবে না সুতরাং তিন অঙ্কের সংখ্যাগুলোর মধ্যে
যে সংখ্যাগুলোতে একটি ডিজিট 3 বার

করে এসেছে সেগুলো বাদ দিতে হবে।

আর সেগুলো হলো. ( 111, 222,333, ......... 888,999 = 9 টি )

তাই উত্তর:729-9 = 720টি

\subsubsection{সাধারণ সংখ্যা
গঠন}\label{ux9b8ux9a7ux9b0ux9a3-ux9b8ux996ux9af-ux997ux9a0ux9a8}

\paragraph{৩৪ . 1,3,5,7,9,অঙ্কগুলি-ব্যবহার করে 3 অঙ্ক বিশিষ্ট কতগুলি সংখ্যা গঠন
করা যাবে
।}\label{ux985ux999ux995ux997ux9b2-ux9acux9afux9acux9b9ux9b0-ux995ux9b0-3-ux985ux999ux995-ux9acux9b6ux9b7ux99f-ux995ux9a4ux997ux9b2-ux9b8ux996ux9af-ux997ux9a0ux9a8-ux995ux9b0-ux9afux9ac}

(i) অঙ্কগুলো পুনরাবৃত্তি না করে = \(n_{P_{r}}\) = \(5_{P_{3}}\) = 5×4×3 = 60

(ii) অঙ্কগুলি পুনরাবৃত্তি করে অর্থাৎ একাধিক বার নিয়ে = \(n^{r}\) = \(5^{3}\) =
5×5×5 = 125

\textbf{প্রশ্ন নম্বর-} ৩৫. ৫, ৯, ১,৪ অংকগুলি দ্বারা ৫,০০০ এর চেয়ে বড় কতগুলো
সংখ্যা তৈরি করা যায়? {[}স্বরাষ্ট্র মন্ত্রণালয়ের অধীন বহিরাগমন ও পাসপোর্ট
অধিদপ্তরের সহকারী পরিচালক ২০১১{]}

\begin{itemize}
\item
  (ক) ১২ টি
\item
  (খ) ৮ টি
\item
  (গ) ১৮ টি
\item
  (ঘ) ১৬ টি
\end{itemize}

\textbf{\hl{উত্তরঃ (ক) ১২ টি}}

\textbf{Explanation:}

৫,৯, ১, ৪ এই চারটি অংশ দ্বারা ৫,০০০ এর চেয়ে বড় কোনো সংখ্যা তৈরি করতে হলে
প্রথম স্থানে ৫ বা ৯ বসাতে হবে,

(১ ও ৪ বসালে তা ৫০০০ এর থেকে ছোট হয়ে যাবে) ১ম স্থানটি ৫ ও ৯ দ্বারা
\(\text{২}_{P_{\text{১}}}\) বা ২ প্রকারে পূরণ করা যায় বাকী তিনটি স্থান তিনটি অঙ্ক দ্বারা ৩!
বা ৬ প্রকারে পূরণ করা যায়।

মোট বিন্যাস = ২ × ৬ = ১২ ।

অথবা প্রথমে ৫ কে নির্দিষ্ট রেখে অন্য ৩টি দিয়ে বিন্যাস = ৩! = ৬টি। একইভাবে
প্রথমে ৯ কে নির্দিষ্ট রেখে অন্য ৩টি দিয়ে বিন্যাস: = ৩ ! = ৬টি।

মোট বিন্যাস = ৬+৬ = ১২টি।

\textbf{প্রশ্ন নম্বর-} ৩৬. ৩,৫,৭,৮ এবং ৯ সংখ্যাগুলো দিয়ে ৪ অঙ্কবিশিষ্ট কতগুলো
সংখ্যা গঠন করা যাবে যারা ৭০০০ এর থেকে বড়? ...

\begin{itemize}
\item
  (ক) 48
\item
  (খ) 60
\item
  (গ) 72
\item
  (ঘ) 80
\end{itemize}

\textbf{\hl{উত্তরঃ (গ) 72}}

\textbf{Explanation:}

হাজারের স্থানে 3 ও 5 কে নিলে 7000 থেকে কম হয় তাই ৩ ও ৫ বাদ দিয়ে অন্য
সংখ্যাগুলো থেকে নিতে হবে। সুতরাং ৭,৮ এবং ৯ এই ৩টি সংখ্যা থেকে ১টি নেয়া যায়
\(3_{P_{1}}\) = 3 ভাবে । আবার ১ম অঙ্কটি নেয়ার পর ৭,৮,৯ এর অবশিষ্ট

২ অঙ্ক এবং অপর ৩টি অঙ্ক সহ মোট ৫টি থেকে ৩টি নেয়া যায় \(4_{P_{3}}\) = 24
ভাবে । সুতরাং মোট সংখ্যা গঠন = ৩×২৪ = ৭২ টি।

\textbf{প্রশ্ন নম্বর-} ৩৭. 1,3,5,7,9 অঙ্কগুলো থেকে তিনটি ভিন্ন ভিন্ন অংক নিয়ে
200 থেকে বৃহত্তম তিন অংকের কতগুলো সংখ্যা গঠন করা যাবে। ...

\begin{itemize}
\item
  (ক) 48
\item
  (খ) 60
\item
  (গ) 72
\item
  (ঘ) 80
\end{itemize}

\textbf{\hl{উত্তরঃ (ক) 48}}

\textbf{Explanation:}

শতকের অবস্থান পূর্ণ করা যাবে = ৪ভাবে (৩,৫,৭,৯ দিয়ে)

দশকের ঘর পূর্ণ করা যাবে ৪ ভাবে, (৩,৫,৭,৯ থেকে ১টি বাদ দিয়ে ৩ টি + ১)

এককের ঘর পূর্ণ করা যাবে আগের দুটি বাদ দিয়ে যে কোন ৩টি।

মোট বিন্যাস = ৪×৪×৩ = ৪৮

\textbf{প্রশ্ন নম্বর-} ৩৮. 1,2,3,4,5,6 অংকগুলো প্রত্যেক সংখ্যায় কেবল মাত্র একবার
মাত্র ব্যবহার করে 5000 ও 6000 এর মধ্যবর্তী কতগুলো সংখ্যা গঠন করা যাবে? ...

\begin{itemize}
\item
  (ক) 48
\item
  (খ) 60
\item
  (গ) 72
\item
  (ঘ) 80
\end{itemize}

\textbf{\hl{উত্তরঃ (খ) 60}}

\textbf{Explanation:}

5000 থেকে 6000 এর মধ্যে সংখ্যা নিতে হলে প্রথম সংখ্যাটি অবশ্যই ৫ ই হতে হবে। বড়
বা ছোট নিলে ৫০০০-৬০০০

এর মধ্যে থাকবে না ।

প্রদত্ত সংখ্যাগুলো ১,২,৩,৪,৫,৬ = মোট ৬টি সংখ্যার মধ্যে ৫ কে হাজারের ঘরে
নির্দিষ্ট রেখে এরপর অবশিষ্ট ৬-১ = ৫টি সংখ্যা থেকে শতক, দশক ও এককের ঘরে যে
কোন সংখ্যা ই নেয়া যাবে। তাই পরের ৩টি ঘর পূর্ণ করা

যাবে = ৫×৪×৩ = ৬০ ভাবে । উত্তর: ৬০টি।

\textbf{শর্টকার্টঃ} সবগুলো থেকে আগেই ৫ বাদ দিতে হবে এবং ৫ কে নিয়ে নিলে আর
নিতে হবে ৪-১ = ৩টি। তাই একসাথে লেখা যায় \({6\  - \ 1}_{P_{4\  - \ 1}}\)
= \(5_{P_{3}}\) = 60 উপায়ে ।

\subsubsection{শূন্য যুক্ত সংখ্যা
গঠন}\label{ux9b6ux9a8ux9af-ux9afux995ux9a4-ux9b8ux996ux9af-ux997ux9a0ux9a8}

\textbf{প্রশ্ন নম্বর-} ৩৯. 0,1,2,3,4,5, সংখ্যাগুলো একবার মাত্র ব্যবহার করে 6
অঙ্ক বিশিষ্ট কতগুলো অর্থ পূর্ণ সংখ্যা গঠন করা যাবে? ...

\begin{itemize}
\item
  (ক) 700
\item
  (খ) 600
\item
  (গ) 520
\item
  (ঘ) 710
\end{itemize}

\textbf{\hl{উত্তরঃ (খ) 600}}

\textbf{Explanation:}

এখানে মোট সংখ্যা আছে ৬টি যার মধ্যে একটি ০ আছে। তাই ৬ অঙ্কবিশিষ্ট মোট সংখ্যা =
৬! = ৭২০ টি ।

০ দ্বারা যেগুলো শুরু হবে সেগুলো বাদ দিতে হবে। তাই শুরুতে ০ কে নির্দিষ্ট রাখলে ০
বাদে অন্য ৫টি সংখ্যা সাজানো যায় ৫! = ১২০ ভাবে ।

সুতরাং ০ বাদে ৬ অঙ্কবিশিষ্ট অর্থপূর্ণ সংখ্যা = ৭২০ - ১২০ = ৬০০টি।

\textbf{প্রশ্ন নম্বর- ৪০. প্রত্যেক অঙ্ক একবার মাত্র ব্যবহার করে 0,1,2,3,5,6
দ্বারা চার অঙ্ক বিশিষ্ট অর্থ পূর্ণ সংখ্যা কত?}

\begin{itemize}
\item
  (ক) 350
\item
  (খ) 400
\item
  (গ) 500
\item
  (ঘ) 300
\end{itemize}

\textbf{\hl{উত্তরঃ (ঘ) 300}}

\textbf{Explanation:}

শূন্য সহ সকল অংক = \(6_{P_{4}}\) = 6×5×4×3 = 360

শুধু শূন্য বিশিষ্ট অংক যেগুলো অর্থপূর্ণ নয় = \({6\  - \ 1}_{P_{4\  - \ 1}}\) =
\(5_{P_{3}}\) = 5×4×3 = 60 (শুরুতে (0)

সুতরাং মোট বিন্যাস: = ৩৬০-৬০ = ৩০০টি।

{[}০ এর সাথে জোড় বিজোড় থাকলে অনেকের কাছে কঠিন লাগে তাদের জন্য নিচের
প্রশ্নগুলো, বুঝে বুঝে করুন:{]}

\textbf{প্রশ্ন নম্বর-} ৪১. প্রতিটি অঙ্ক একবার ব্যবহার করে ৪, ৩, ২, ১, ০ অঙ্কগুলি
দ্বারা ৫ অঙ্কের কতগুলি বিজোড় সংখ্যা গঠন করা যাবে? ...

\begin{itemize}
\item
  (ক) ২৪
\item
  (খ) ২৬
\item
  (গ) ২৮
\item
  (ঘ) ৩৬
\end{itemize}

\textbf{\hl{উত্তরঃ (ঘ) ৩৬}}

\textbf{Explanation:}

এখানে শর্ত দুটি: ১. সংখ্যাগুলো ৫ অংক বিশিষ্ট হতে হবে

২. সংখ্যাগুলো বিজোড় হতে হবে। (১ ও ৩ দিয়ে শেষ হতে হবে)

শেষে ১ কে নির্দিষ্ট রেখে মোট বিন্যাস = ৪! = ২৪টি ।

কিন্তু এই ২৪টি বিজোড়ের মধ্যে শেষে ১ থাকবে কিন্তু শুরুতে ০ আছে এমন বিজোড় গুলো ৫
অঙ্কবিশিষ্ট নয়। তাই সেগুলো বাদ দিতে হবে ।

শুরুতে ০ ও শেষে ১ কে নির্দিষ্ট রেখে অন্য ৩টি অঙ্ক সাজানো যায় ৩! = ৬ ভাবে।

তাহলে ১ দিয়ে শেষ হয়েছে এমন ৫ অঙ্কবিশিষ্ট বিজোড় সংখ্যা রয়েছে = ২৪-৬ = ১৮ টি।

ঠিক, একইভাবে শেষে ৩ দিয়ে বিজোড় সংখ্যা হবে ২৪-৬ = ১৮ টি।

সুতরা মোট বিজোড় সংখ্যা হবে = ১৮+১৮ = ৩৬টি।

\textbf{গুরুত্বপূর্ণ শিক্ষণীয় বিষয়:}

০ দ্বারা কোন সংখ্যা শুরু হয় না ।

জোড় অথবা বিজোড় বের করতে বলা হলে শেষের সংখ্যা মেলাতে হবে।

\subsubsection{সমাবেশ
(Combination)}\label{ux9b8ux9aeux9acux9b6-combination}

সমাবেশ হলো ক্রম উপেক্ষা করে সাজানোর প্রক্রিয়া। সমাবেশের ক্ষেত্রে ধারাবাহিকতা
কোন বিষয় নয়। বিন্যাস ও সমাবেশের মধ্যে বিন্যাস উপরোক্ত আলোচনায় তুলে ধরা হয়েছে
। এবার সমাবেশ অধ্যায়টিকে নিচে সুন্দর ও সূচারুরুপে তুলে ধরার চেষ্টা করছি । আশা
করছি উপকৃত হবেন।

\paragraph{সমাবেশ (Combination)
কি?}\label{ux9b8ux9aeux9acux9b6-combination-ux995}

সমাবেশ হলো কয়েকটি উপাদান থেকে প্রত্যেকবার নির্দিষ্ট কিছু উপাদান নিয়ে এক একটি
দল গঠন করা। এখানে ধারাবাহিকতা পরিবর্তন হলেও দলের সংখ্যা একই থাকবে।

সমাবেশের সূত্র: \(n_{c_{r}} = \frac{n!}{r!(n - r)!}\) {[}বিন্যাসের সূত্রের
মতই শুধু অতিরিক্ত হিসেবে হরের সাথে r! গুণ করতে হবে।{]}

\subsubsection{বিন্যাস বনাম সমাবেশ (Permutation Vs
Combination)}\label{ux9acux9a8ux9afux9b8-ux9acux9a8ux9ae-ux9b8ux9aeux9acux9b6-permutation-vs-combination}

Combination এর ক্ষেত্রে Order (ধারাবাহিকতা) কোন Factor নয়। কিন্তু
Permutation এর ক্ষেত্রে ধারাবাহিকতা গুরুত্বপূর্ণ এবং Order এর পরিবর্তন হলে
সংখ্যারও পরিবর্তন হবে। যেমন:

বিভিন্ন পরীক্ষার প্রশ্নে যখন এই দুটি অধ্যায় থেকে প্রশ্ন আসবে তখন লিখে দেয়া থাকবে
না কোনটি বিন্যস এবং কোনটি

সমাবেশ হবে। ভালোভাবে পার্থক্য না জানলে একটার জায়গায় অন্যটির উত্তর বের করে
ফেলতে পারেন। তাই এদের মধ্যকার

পার্থক্যগুলো নিচে ছক আকারে তুলে ধরা হল।

\subsubsection{বিন্যাস ও সমাবেশের মধ্যকার মৌলিক
পার্থক্য}\label{ux9acux9a8ux9afux9b8-ux993-ux9b8ux9aeux9acux9b6ux9b0-ux9aeux9a7ux9afux995ux9b0-ux9aeux9b2ux995-ux9aaux9b0ux9a5ux995ux9af}

{\def\LTcaptype{none} % do not increment counter
\begin{longtable}[]{@{}
  >{\raggedright\arraybackslash}p{(\linewidth - 2\tabcolsep) * \real{0.4004}}
  >{\raggedright\arraybackslash}p{(\linewidth - 2\tabcolsep) * \real{0.4004}}@{}}
\toprule\noalign{}
\begin{minipage}[b]{\linewidth}\raggedright
\textbf{(বিন্যাস) Permutation}
\end{minipage} & \begin{minipage}[b]{\linewidth}\raggedright
\textbf{(সমাবেশ) Combination}
\end{minipage} \\
\midrule\noalign{}
\endhead
\bottomrule\noalign{}
\endlastfoot
বিন্যাস হলো সাজানোর ধরণ অর্থাৎ কত ভাবে সাজানো যায় তা বের করা, এখানে
ধারাবাহিকতা পরিবর্তন হলে নতুন বিন্যাস হয়। & আর সমাবেশ হলো বাছাই করা, কয়েকজন
থেকে বাছাই করার সময় কে আগে আসলো কে পরে আসলো তা দেখার প্রয়োজন নেই অর্থাৎ
এখানে ধারাবাহিকতা গুরুত্বপূর্ণ নয়। \\
বিন্যাসের সুত্র: n\_P\_r=(n!/(n-r)!) & সমাবেশের সূত্র: n\_c\_r=(n!/r!(n-r)!)
{[}বিন্যাসের সূত্রের মতই শুধু অতিরিক্ত হিসেবে হরের সাথে r! গুণ করতে হবে।{]} \\
বিন্যাসের উত্তর বড় হয়। & সমাবেশের উত্তর ছোট হয়। \\
রাকিব সামনে এবং রহিম পেছনে দাঁড়ানো অথবা রহিম সামনে রাকিব পেছনে দাঁড়ানো
বোঝাতে দুটি ভিন্ন দাঁড়ানোর পদ্ধতি। অর্থাৎ সিরিয়াল পরিবর্তন হলে নতুন বিন্যাস হয়।
& সমাবেশের ক্ষেত্রে বাংলাদেশ- ভারত আর ভারত বাংলাদেশ এর খেলা অর্থ দুটি খেলা
না বরং একটি খেলা। \\
উদাহরণ: AB, BA, দুটি ভিন্ন বিন্যাস। & উদাহরণ: AB, BA উভয় মিলে একটি ই
সমাবেশ। \\
বিন্যাস হয়: (i) অক্ষর সাজানোর প্রশ্নগুলোতে: যেমন: DHAKA (ii) সংখ্যা তৈরী করার
প্রশ্নগুলোতে । যেমন: ১২৩,৩২১ (iii) যে কোন কিছুকে সাজাতে বলা হলে বিন্যাস করতে
হয়। & সমাবেশ হয়: (i) হ্যান্ডশেক (ii) খেলা (iii) দল (iv) কমিটি (v) যে কোন
কিছু বাছাই করার প্রশ্নগুলোতে সমাবেশের সূত্র প্রয়োগ করতে হয়। \\
\end{longtable}
}

১. \(n_{c_{12}} = n_{c_{6}}\) হলে n এর মান কত? {[} ৩৯তম বিসিএস {]}

এই প্রশ্নটি Combination এর একটি ব্যাসিক বিষয় । যেমন: \(10_{c_{9}}\) = কে
লেখা যায় \(10_{c_{10\  - \ 1}} = 10_{c_{9}}\) কারণ দুটির মান ই সমান হবে।

তেমনি ভাবে n এর জায়গায় এমন একটি সংখ্যা বসাতে হবে যাতে দু পাশের মান সমান হয়
। অপশন থেকে n = 18 নিলে,

বামপক্ষ \(18_{c_{12}}\)

= \(18_{c_{18\  - \ 12}}\)

= \(18_{c_{6}}\) হয় যা ডান পক্ষ \(n_{c_{6}}\)

= \(18_{c_{6}}\) এর সমান।

সুতরাং n = 18

\subsubsection{করমর্দন ও খেলার
সংখ্যা}\label{ux995ux9b0ux9aeux9b0ux9a6ux9a8-ux993-ux996ux9b2ux9b0-ux9b8ux996ux9af}

এই পদ্ধতিতে আমার শিখবো হ্যান্ডশেক সংখ্যা বের করা এবং কিভাবে কয়েকজন খেলোয়াড়ের
ভেতর থেকে কতভাবে একটি ক্রিকেট, ফুটবল,বাস্কেটবল অথবা যে কোন দল গঠন করা যায়।
সাথে সাথে কিভাবে এবং কতভাবে একটি দলের অধিনায়ক অথবা সহ অধিনায়ক নির্বাচিত
করা যায় ।

দল গঠনের সময় বিভিন্ন খেলোয়াড়ের নাম আগে অথবা পরে যখনই বলা হোক না কেন তারা
একটি দলই বোঝাবে, তাই দল গঠনের অংক গুলো সমাবেশের সুত্রানুযায়ী করতে হয়।

\textbf{প্রশ্ন নম্বর-} ২. ১০ জন লোক প্রত্যেকে প্রত্যেকের সাথে করমর্দন করে।
করমর্দন সংখ্যা কত? ...

\begin{itemize}
\item
  (ক) ৩৮
\item
  (খ) ৪০
\item
  (গ) ৪২
\item
  (ঘ) ৪৫
\end{itemize}

\textbf{\hl{উত্তরঃ (ঘ) ৪৫}}

\textbf{Explanation:}

যে কোন করমর্দণ অথবা কোলাকুলির অংকে শুধু কত জন লোক করমর্দন (Handshake), বা
কোলাকুলি করলো তা দেয়া থাকবে। এক্ষেত্রে মনে রাখতে হবে যে প্রত্যেক বার করমর্দন
বা কোলাকুলি করার সময় মোট ২ জন লোকের প্রয়োজন। তাই

এক্ষেত্রে সুত্রটি হবে = \(n_{c_{2}}\) = \({\text{মোট}\ \text{লোক}}_{c_{\text{২}\ \text{জন}\ \text{সব}\ \text{সময়}}}\)

প্রদত্ত প্রশ্নটির সমাধান:

\(10_{c_{2}}\)

= \(\frac{10!}{2!(10 - 2)!}\)

= \(\frac{10!}{2!8!}\)

= \(\frac{10 \times 9 \times 8!}{2 \times 8!}\)

= 5×9 = 45 Shortcut Solution: এ ধরনের অংক খাতা কলম ছাড়াই মুখে মুখে করুন
এভাবে :যত জনই দেয়া থাক তার আগের সংখ্যার সাথে ঐ সংখ্যাটি গুণ করে ২ দিয়ে ভাগ
করলেই উত্তর বের হয়ে যাবে। যেমন এখানে ১০ এর আগের সংখ্যা ৯ । ১০ এবং ৯ গুণ
করলে হয় ৯০ । ৯০ কে ২ দিয়ে ভাগ করলে হয় ৪৫ । তাই ৪৫ ই উত্তর ।

\textbf{প্রশ্ন নম্বর-} ৩. একটি ফুটবল টুর্নামেন্টে ৬ টি দল অংশগ্রহণ করেছে, একক
লীগ পদ্ধতিতে খেলা হলে মোট কতটি খেলা পরিচালনা করতে হবে? {[}বাংলাদেশ রেলওয়ে
(সহ: ষ্টেশন মাষ্টার)-২০১৮{]}

\begin{itemize}
\item
  (ক) ১০
\item
  (খ) ১২
\item
  (গ) ১৫
\item
  (ঘ) ১৬
\end{itemize}

\textbf{\hl{উত্তরঃ (গ) ১৫}}

\textbf{Explanation:}

৬টি দল অংশগ্রহণ করে একক লীগ পদ্ধতিতে খেলা হলে প্রত্যেকে প্রত্যেকের সাথে ১টি
করে খেলা খেলবে।

তাহলে মোট খেলা হবে

\(\text{৬}_{c_{\text{২}}}\)

= \(\frac{\text{৬} \times \text{৫}}{\text{২} \times \text{১}}\)

= ১৫টি।

\textbf{যাদের সমাবেশের সুত্র জানা নেই তারা এভাবে করতে পারেন:}

প্রথম ১ টি দল অন্য ৫টি দলের সাথে খেলবে = ৫টি

২য় দলটি অন্য ৪টির সাথে খেলবে, = ৪টি।( কারণ ১ম দলের সাথের খেলা প্রথম ৫টির
মধ্যে আছে তাই ১টি কম)

৩য় দলটি অন্য ৩টি দলের সাথে খেলবে = ৩টি।

৪র্থ দলটি অন্য ২টি দলের সাথে খেলবে = ২টি।

৫ম দলটি অন্য ১টি দলের সাথে খেলবে = ১টি।

তাহলে মোট খেলা হবে : ৫+৪+৩+২+১ = ১৫টি।

\textbf{\hl{উত্তর: ১৫টি।}}

\textbf{মুখে মুখে:}

এই নিয়মটি জানা থাকলে এক লাইনে : ৫+৪+৩+২+১ = ১৫টি।

( কারণ যতটি দল ই খেলুক না কেন প্রতিবার ১ করে কমতে থাকবে ২টি করে খেললে শুধু ২
গুণ)

\textbf{প্রশ্ন নম্বর- ৪. ৬ জন লোক এক জন আরেকজনের সাথে করমর্দন করলে মোট
করমর্দনের সংখ্যা কত?}

\begin{itemize}
\item
  (ক) ১০
\item
  (খ) ১২
\item
  (গ) ১৫
\item
  (ঘ) ২০
\end{itemize}

\textbf{\hl{উত্তরঃ (গ) ১৫}}

\textbf{Explanation:}

Shortcut Solution: এ ধরনের অংক খাতা কলম ছাড়াই মুখে মুখে করুন এভাবে :যত জনই
দেয়া থাক তার আগের সংখ্যার সাথে ঐ সংখ্যাটি গুণ করে ২ দিয়ে ভাগ করলেই উত্তর বের
হয়ে যাবে। যেমন এখানে ৬ এর আগের সংখ্যা ৫ । ৬ এবং ৫ গুণ করলে হয় ৩০ । ৩০ কে ২
দিয়ে ভাগ করলে হয় ১৫ । তাই ১৫ ই উত্তর ।

\textbf{প্রশ্ন নম্বর-} ৫. একটি মিটিং এ উপস্থিত ৫ জন সদস্য প্রত্যেকে প্রত্যেকের
সাথে হ্যান্ডশেক করলে মোট কতটি হ্যান্ডশেক হবে? {[} M.T.B.L.Off: 13{]}

\begin{itemize}
\item
  (ক) 10
\item
  (খ) 11
\item
  (গ) 15
\item
  (ঘ) 5
\end{itemize}

\textbf{\hl{উত্তরঃ (ক) 10}}

\textbf{Explanation:}

Shortcut Solution: এ ধরনের অংক খাতা কলম ছাড়াই মুখে মুখে করুন এভাবে :যত জনই
দেয়া থাক তার আগের সংখ্যার সাথে ঐ সংখ্যাটি গুণ করে ২ দিয়ে ভাগ করলেই উত্তর বের
হয়ে যাবে। যেমন এখানে ৫ এর আগের সংখ্যা ৪ । ৫ এবং ৪ গুণ করলে হয় ২০ । ২০ কে ২
দিয়ে ভাগ করলে হয় ১০ । তাই ১০ ই উত্তর ।

\textbf{প্রশ্ন নম্বর-} ৬. একটি ফুটবল লীগে ৫টি দলের প্রত্যেক দল অন্য দলের সাথে
২টি করে ম্যাচ খেললে সর্বমোট কতটি ম্যাচ অনুষ্ঠিত হবে? {[} Dutch Bangla Bank
\textquotesingle MTO`12{]}

\begin{itemize}
\item
  (ক) 10
\item
  (খ) 20
\item
  (গ) 9
\item
  (ঘ) 19
\end{itemize}

\textbf{\hl{উত্তরঃ (খ) 20}}

\textbf{Explanation:}

\textbf{মুখে মুখে সমাধান:}

যত জনই দেয়া থাক তার আগের সংখ্যার সাথে ঐ সংখ্যাটি গুণ করে ২ দিয়ে ভাগ করলেই
উত্তর বের হয়ে যাবে। তবে এখানে যেহেতু ২ টি করে ম্যাচ খেলার কথা বলা হয়েছে তাই
প্রাপ্ত ফলাফলকে ২ দিয়ে গুণ করে উত্তর বের করতে হবে। এখানে ৫ এর আগের সংখ্যা ৪ ।
৫ এবং ৪ গুণ করলে হয় ২০ । ২০ কে ২ দিয়ে ভাগ করলে হয় ১০ । এখন এই ১০ কে ২
দিয়ে গুণ করলে হয় ২০ . তাই ২০ ই সঠিক উত্তর ।

\textbf{উল্টোভাবে আসলে}

অর্থাৎ হ্যান্ডশেকের সংখ্যা দেয়া থাকবে, মোট লোকের সংখ্যা বের করতে বলা হলে ।
যেমন:

\textbf{প্রশ্ন নম্বর- ৭. একটি ইফতার পার্টিতে কয়েকজন বন্ধু উপস্থিত হল । তারা
প্রত্যেকে প্রত্যেকের সাথে করমর্দন করল। যদি মোট ২৮ টি করমর্দন হয়ে থাকে তাহলে ঐ
পার্টিতে মোট কতজন উপস্থিত ছিল?}

\begin{itemize}
\item
  (ক) ৮ জন
\item
  (খ) ১২ জন
\item
  (গ) ১৪ জন
\item
  (ঘ) ১৬ জন
\end{itemize}

\textbf{\hl{উত্তরঃ (ক) ৮ জন}}

\textbf{Explanation:}

উপরের নিয়মটির ই উল্টোভাবে প্রয়োগ করুন এভাবে,সবার শেষে ২ দিয়ে ভাগ করায় ২৮
হয়েছে তাহলে ভাগ করার আগে ছিল ২৮×২ = ৫৬। এখন এই ৫৬ কে এমন ভাবে ভাঙ্গতে হবে
যাতে পর পর দুটি সংখ্যা গুণ করলে ৫৬ হয়। সংখ্যা দুটি হল ৮×৭ = ৫৬। এখন বড়
সংখ্যাটি ই হলো নির্ণেয় উপস্থিত লোকের সংখ্যা। উত্তর: ৮ জন।

\textbf{প্রশ্ন নম্বর-} ৮. একটি পার্টিতে কিছু লোক উপস্থিত ছিল। তারা প্রত্যেকে
প্রত্যেকের সাথে হ্যান্ডশেক করায় মোট ৬৬ টি হ্যান্ডশেক হলো। ঐ পার্টিতে মোট কত জন
উপস্থিত ছিল? {[} RAKUB Senior off: - 2015 {]}

\begin{itemize}
\item
  (ক) 9
\item
  (খ) 15
\item
  (গ) 10
\item
  (ঘ) 12
\end{itemize}

\textbf{\hl{উত্তরঃ (ঘ) 12}}

\textbf{Explanation:}

উপরের নিয়মটির ই উল্টোভাবে প্রয়োগ করুন এভাবে,সবার শেষে ২ দিয়ে ভাগ করায় ৬৬
হয়েছে তাহলে ভাগ করার আগে ছিল ৬৬×২ = ১৩২। এখন এই ১৩২ কে এমন ভাবে ভাঙ্গতে
হবে যাতে পর পর দুটি সংখ্যা গুণ করলে ১৩২ হয়। সংখ্যা দুটি হল ১২×১১ = ১৩২। এখন
বড় সংখ্যাটি ই হলো নির্ণেয় উপস্থিত লোকের সংখ্যা। উত্তর: ১২ জন।

\subsubsection{কমিটি বা দল
গঠন}\label{ux995ux9aeux99f-ux9ac-ux9a6ux9b2-ux997ux9a0ux9a8}

\textbf{কেন বিভিন্ন কমিটি গঠনের অংকগুলো সমাবেশের সুত্রানুযায়ী করতে হয়?}

একটি কমিটিতে ৩ জন সদস্য থাকার অর্থ ৩ সদস্য বিশিষ্ট কমিটি। এখন যেভাবেই যাকেই
আগে অথবা পরে দিন, কমিটি একই থাকবে তাই এই প্রশ্নগুলো সমাবেশের সুত্রঅনুযায়ী করতে

হয় ।

\textbf{প্রশ্ন নম্বর- ৯. ৬ জন বালক ও ৪ জন বালিকা থেকে ৫ সদস্যবিশিষ্ট কমিটি
করার কয়টি পথ আছে যেখানে ঠিক ২ জন বালিকা যাবে ?}

\begin{itemize}
\item
  (ক) 80
\item
  (খ) 100
\item
  (গ) 120
\item
  (ঘ) 130
\end{itemize}

\textbf{\hl{উত্তরঃ (গ) 120}}

\textbf{Explanation:}

৫ সদস্য বিশিষ্ট কমিটিতে বালিকা ২ জন থাকলে ৩ জন বালক থাকবে।

তাহলে ৪ জন বালিকা থেকে ২ জন, এবং ৬ জন বালক থেকে ৩ জন নিয়ে কমিটির প্রকার,

\(\text{৪}_{c_{\text{২}}} \times \text{৬}_{c_{\text{৩}}}\)

=
\(\frac{\text{৪} \times \text{৩}}{\text{২} \times \text{১}} \times \frac{\text{৬} \times \text{৫} \times \text{৪}}{\text{৩} \times \text{২} \times \text{১}}\)

= ১২০ {[} ভিন্ন ভিন্ন প্রকারে নেয়ার পর গুণ করতে হয় {]}

\textbf{প্রশ্ন নম্বর- ১০. একটি ক্লাবের ৮ জন পুরুষ ও ৮ জন মহিলা সদস্য আছেন। ৬
সদস্যের একটি কমিটি গঠন করতে হবে। যেখানে পুরুষ ও মহিলা সদস্য ৩ জন করে থাকবেন।
কতভাবে এ কমিটি গঠন করা যায়?}

\begin{itemize}
\item
  (ক) ৩১৩৬
\item
  (খ) ৩১৩৫
\item
  (গ) ৩১৩৪
\item
  (ঘ) ৩১৩৯
\end{itemize}

\textbf{\hl{উত্তরঃ (ক) ৩১৩৬}}

\textbf{Explanation:}

৬ সদস্য বিশিষ্ট কমিটিতে মহিলা ৩ জন থাকলে ৩ জন পুরুষ ও থাকবে।

তাহলে ৮ জন পুরুষ থেকে ৩ জন, এবং ৮ জন মহিলা থেকে ৩ জন নিয়ে কমিটির প্রকার,

\(\text{৮}_{c_{\text{৩}}} \times \text{৮}_{c_{\text{৩}}}\)

=
\(\frac{\text{৮} \times \text{৭} \times \text{৬}}{\text{৩} \times \text{২} \times \text{১}} \times \frac{\text{৮} \times \text{৭} \times \text{৬}}{\text{৩} \times \text{২} \times \text{১}}\)

= ৩১৩৬ {[} ভিন্ন ভিন্ন প্রকারে নেয়ার পর গুণ করতে হয় {]}

\textbf{প্রশ্ন নম্বর-} ১১. একটি স্কুলের কমিটিতে ২ জন শিক্ষক এবং ৪ জন ছাত্র
থাকে। ৫ জন শিক্ষক এবং ১০ জন ছাত্র থেকে কত উপায়ে বাছাই করা যাবে? {[} Exim
Bank. T. Off. -2014 {]}

\begin{itemize}
\item
  (ক) 220
\item
  (খ) 5100
\item
  (গ) 2100
\item
  (ঘ) 3200
\end{itemize}

\textbf{\hl{উত্তরঃ (গ) 2100}}

\textbf{Explanation:}

৫ জন শিক্ষকের মধ্য থেকে ২ জন শিক্ষক নেয়া যাবে \(5_{c_{2}}\)

ভাবে।

আবার ১০ জন ছাত্রের মধ্য থেকে ৪ জন ছাত্র নেয়া যাবে \(10_{c_{4}}\) ভাবে।

তাহলে মোট বাছাই করার উপায়

= \(5_{c_{2}} \times 10_{c_{4}}\)

= \(\frac{5!}{2! \times 3!} \times \frac{10!}{4! \times 6!}\)

= 10×210 = 2100

\textbf{প্রশ্ন নম্বর-} ১২. 4 জন মহিলা ও 6 জন পুরুষের মধ্য থেকে 4 সদস্য বিশিষ্ট
একটি উপ-কমিটি গঠন করতে হবে যাতে 1 জন নির্দিষ্ট পুরুষ সর্বদায় উপস্থিত থাকেন। কত
প্রকারে ঐ কমিটি গঠন করা যেতে পারে ? {[} ৩৮ তম বিসিএস প্রিলি:{]}

\begin{itemize}
\item
  (ক) 210
\item
  (খ) 304
\item
  (গ) 84
\item
  (ঘ) 120
\end{itemize}

\textbf{\hl{উত্তরঃ (গ) 84}}

\textbf{Explanation:}

মোট সদস্যা = 4+6 = 10 জন । এখন 1 জন নির্দিষ্ট পুরুষ বাদ দিয়ে 10-1 = 9 জন থেকে
নিতে হবে 4-1 = 3 জন ।

9 জন থেকে 3 জন নেয়ার উপায় হল

= \(9_{c_{3}}\)

= \(\frac{9!}{3!(\ 9\  - \ 3\ )!}\)

=
\(\frac{9 \times 8 \times 7 \times 6!}{3 \times 2 \times 1 \times 6!}\)

= 84 {[} শর্টকাট:
\(\frac{9 \times 8 \times 7}{3 \times 2 \times 1} = 84\) {[} নিচে 3
থাকায় উপরে 9 থেকে শুরু করে 3টি উপাদান এবং নিচে 3! এর মান বসাতে হবে {]}

{[}Concept Clear: যাকে নির্দিষ্ট করে রাখতে বলা হবে তাকে বাদ দিয়ে হিসেব
করতে হবে। কারণ সমাবেশ হচ্ছে লটারীর মত বিষয়। যেখানে সবাইকে লটারীর মত ড্র
করলে কখনো কেউ উঠবে আবার সে বাদ ও যেতে পারে।

একইভাবে নির্দিষ্ট লোককে লটারীর মধ্যে রাখলে লটারীর মত সে ও বাদ পরে যেতে
পারে। তাই প্রথমেই তাকে বাদ দিয়ে হিসেব করতে হবে। ৩৫ ও ৩৬ তম বিসিএস এ একই
নিয়মের ১টি করে প্রশ্ন এসেছিল। {]}

{[} Confusion Clear : এখানে পুরুষ মহিলা থাকায় অনেকে কয়জন পুরুষ কয়জন মহিলা
নিতে হবে তা নিয়ে কনফিউশনে থাকতে পারেন।

কিন্তু প্রশ্নটিতে উপকমিটি গঠন করার সময় কতজন পুরুষ বা মহিলা নিতে হবে তা নির্দিষ্ট
করে বলা না থাকায় পুরুষ মহিলা যে কাউকে যতজন খুশি নেয়া যাবে। (ওরকম অন্য আরো
প্রশ্ন আছে কিন্তু এটা না) {]}

\subsubsection{কখন গুণ (×) আর কখন যোগ
(+)}\label{ux995ux996ux9a8-ux997ux9a3-ux986ux9b0-ux995ux996ux9a8-ux9afux997}

যখন একটির সাথে অন্যটি নির্ভরশীল থাকে তখন গুণ করতে হবে। (প্রশ্নে ``এবং'' থাকলে
`গুণ'' )

যেমন: মোট ৫ জন পুরুষ এবং ৪ জন মহিলা থেকে ৫ জন সদস্য নিয়ে একটি কলেজের কমিটি
গঠন করতে হবে যেখানে ২ জন মহিলা থাকবে ।

এখানে শুধু মহিলা বা শুধু পুরুষ নিয়ে কমিটি হবে না বরং পুরুষ ও মহিলা উভয়ে মিলে
কমিটি হবে। অর্থাৎ একটার সাথে আরেকটা নির্ভরশীল । তাই এক্ষেত্রে গুণ করতে হবে
\(5_{c_{3}} \times 4_{c_{3}}\) = 10×6 = 60

: কিন্তু একটির উপর আরেকটি নির্ভরশীল না হলে যোগ করতে হবে। (প্রশ্নে ``অথবা''
থাকলে \textquotesingle যোগ\textquotesingle{} )

যেমন: একটি কলেজের কমিটি তৈরী করার উপায় আছে ২০টি আরেকটি ভিন্ন কলেজের কমিটি
তৈরী করার উপায় আছে ১০টি। এখানে একটি কলেজের সাথে অন্য কলেজের কমিটির কোন
নির্ভরশীলতা নেই, তাই এক্ষেত্রে মোট কমিটি সংখ্যা 20+10 = 30টি

\subsubsection{সমাবেশের
বিবিধ}\label{ux9b8ux9aeux9acux9b6ux9b0-ux9acux9acux9a7}

নির্দিষ্ট কোন ব্যক্তি বা বস্তুকে রেখে কোন দল, কমিটি বা কোন কিছু সাজাতে বলা হলে:

\textbf{প্রশ্ন নম্বর-} ১৩. 14 জন খেলোয়াড়ের মধ্যে থেকে নির্দিষ্ট একজন অধিনায়কসহ
11 জনের একটি ক্রিকেট দল কতভাবে বাছাই করা যাবে? {[} 35তম বিসিএস{]}

\begin{itemize}
\item
  (ক) 728
\item
  (খ) 286
\item
  (গ) 364
\item
  (ঘ) 1001
\end{itemize}

\textbf{\hl{উত্তরঃ (খ) 286}}

\textbf{Explanation:}

যেহেতু অধিনায়ককে বাদ দেয়া যাবে না, তাই অধিনায়ক সবসময় ফিক্সড বা নির্দিষ্ট,
তাই তাকে আলাদা করে রাখতে হবে।। বাকী ১৩ জনের মধ্য থেকে ১০ জন বাছাই করতে
হবে।

একজনকে অধিনায়ক হিসেবে বাছাই করার পর বাকি (14 - 1) বা 13 জন হতে (11 - 1)
বা 10 জনকে বাছাই করে 11 জনের দল গঠনের উপায়

= \(13_{c_{10}}\)

= \(\frac{13 \times 12 \times 11}{3 \times 2 \times 1}\)

= 286

\textbf{প্রশ্ন নম্বর-} ১৪. 12 টি পুস্তক থেকে 5 টি কত প্রকারে বাছাই করা যায়
যেখানে 2 টি পুস্তক সর্বদাই অর্ন্তভুক্ত থাকবে? {[} ৩৬তম বিসিএস{]}

\begin{itemize}
\item
  (ক) 252
\item
  (খ) 792
\item
  (গ) 224
\item
  (ঘ) 120
\end{itemize}

\textbf{\hl{উত্তরঃ (ঘ) 120}}

\textbf{Explanation:}

12 টির মধ্যে যে 5 টি নিতে হবে তার মধ্যে 2 টি পুস্তক যেহেতু নির্দিষ্ট থাকবে তাই
প্রথমেই 2 টি পুস্তক আলাদা করে বাকী 10 টি পুস্তক থেকে 3টি পুস্তক বাছাই করতে হবে।
10 টি পুস্তক থেকে 3টি বাছাই করার পদ্ধতি হলো:

\(10_{c_{3}}\)

= \(\frac{10 \times 9 \times 8 \times 7!}{3 \times 2 \times 7!}\)

= 120

\textbf{প্রশ্ন নম্বর-} ১৫. ৫ জন ব্যক্তির মধ্য থেকে ৩জনের একটি কমিটি কতভাবে
বাছাই করা যাবে যেখানে একজন নির্দিষ্ট ব্যাক্তি সবসময় ই কমিটিতে থাকবে? {[}
Rakub officer-2015 {]}

\begin{itemize}
\item
  (ক) 4
\item
  (খ) 5
\item
  (গ) 6
\item
  (ঘ) 8
\end{itemize}

\textbf{\hl{উত্তরঃ (গ) 6}}

\textbf{Explanation:}

\textbf{মনে রাখুন}: দল গঠনের সময় অধিনায়ক অথবা কমিটি গঠনের সময় কোন নির্দিষ্ট
ব্যক্তি অথবা যে কোন সময়ে কোন নির্দিষ্ট কিছু রেখে বাছাই করার কথা বলা হলে মোট
উপাদান এবং নির্বাচিত উপাদান থেকে এ নির্দিষ্ট জনকে বিয়োগ করে হিসেব করতে হয়।

তা না হলে ঐ লোক ও কখনো দলের বাইরে চলে যেতে পারেন।

তাই এখানে সমাধান: ৫ জনের ভেতর থেকে ১ জনকে আগে থেকেই নিয়ে নিলে লোক থাকবে
৫-১ = ৪ জন। এবং তিন সদস্যের কমিটি থেকে ১ জন নেয়া হয়ে গেলে আর নিতে হবে ৩-১
= ২ জন ।

সুতরাং ৪ জন থেকে ২ জন বাছাই করার মোট উপায়

= \(4_{c_{2}}\)

= \(\frac{4!}{2! \times 2!}\)

= \(\frac{24}{4}\)

= 6

\textbf{প্রশ্ন নম্বর- ১৬. 12টি পুস্তক থেকে 5টি কত প্রকারে বাছাই করা যায় যেখানে
2 টি পুস্তক সর্বদাই অর্ন্তভুক্ত থাকবে না?}

\begin{itemize}
\item
  (ক) 252
\item
  (খ) 792
\item
  (গ) 224
\item
  (ঘ) 120
\end{itemize}

\textbf{\hl{উত্তরঃ (ক) 252}}

\textbf{Explanation:}

এরকম প্রশ্নের ক্ষেত্রে যাদেরকে অর্ন্তভুক্ত করা যাবে না তাদেরকে প্রথমেই বাদ দিতে
হবে। তাই এক্ষেত্রে বাছাই করার উপায় হবে,

= \(10_{c_{5}}\)

= \(\frac{10!}{5!\ (\ 10\  - \ 5\ )!}\)

=
\(\frac{10 \times 9 \times 8 \times 7 \times 6 \times 5!}{5! \times 5!}\)

= \(\frac{30240}{120}\)

= 252

\textbf{প্রশ্ন নম্বর- ১৭. 15 জন খেলোয়াড়ের মধ্যে 11 জনকে নিয়ে একটি ক্রিকেট টিম
গঠন করতে হবে। কিন্তু প্রথম সারির ৪ ব্যাটসম্যানের মধ্যে অবশ্যই 5 জন কে নিয়ে কত
প্রকারে টিম গঠন করা যাবে?}

\begin{itemize}
\item
  (ক) 385
\item
  (খ) 401
\item
  (গ) 392
\item
  (ঘ) 308
\end{itemize}

\textbf{\hl{উত্তরঃ (গ) 392}}

\textbf{Explanation:}

প্রথম ৮ জন থেকে ৫ জন নেয়ার পর পরের অবশিষ্ট (১৫-৮)= ৭ জন থেকে (১১-৫)= ৬ জন
নিয়ে ১১ জনের টিম গঠন করা যাবে

= \(8_{c_{5}} \times 7_{c_{6}}\)

= 56×7 = 392 উপায়ে

\textbf{প্রশ্ন নম্বর- ১৮. Degree শব্দটির অক্ষরগুলো থেকে যেকোন 4টি অক্ষর প্রতিবার
নিয়ে কত প্রকারে বাছাই করা যায়?}

\begin{itemize}
\item
  (ক) 8
\item
  (খ) 6
\item
  (গ) 7
\item
  (ঘ) 5
\end{itemize}

\textbf{\hl{উত্তরঃ (গ) 7}}

\textbf{Explanation:}

Degree শব্দটিতে মোট ৬টি বর্ণের মধ্যে ৩টি e আছে এবং অন্য ৩টি বর্ণ ভিন্ন
(D,g,r)। এখন মোট ৬টি ভিন্ন বর্ণ থেকে ৪টি বর্ণ বাছাই করার উপায় আছে,

{\def\LTcaptype{none} % do not increment counter
\begin{longtable}[]{@{}
  >{\raggedright\arraybackslash}p{(\linewidth - 6\tabcolsep) * \real{0.2006}}
  >{\raggedright\arraybackslash}p{(\linewidth - 6\tabcolsep) * \real{0.2006}}
  >{\raggedright\arraybackslash}p{(\linewidth - 6\tabcolsep) * \real{0.2006}}
  >{\raggedright\arraybackslash}p{(\linewidth - 6\tabcolsep) * \real{0.2006}}@{}}
\toprule\noalign{}
\begin{minipage}[b]{\linewidth}\raggedright
\textbf{শর্ত}
\end{minipage} & \begin{minipage}[b]{\linewidth}\raggedright
\textbf{বাছাইয়ের উপায়}
\end{minipage} & \begin{minipage}[b]{\linewidth}\raggedright
\textbf{বাছাই}
\end{minipage} & \begin{minipage}[b]{\linewidth}\raggedright
\textbf{যে অক্ষরগুলো নেয়া হয়েছে}
\end{minipage} \\
\midrule\noalign{}
\endhead
\bottomrule\noalign{}
\endlastfoot
4টি বর্ণই ভিন্ন এক্ষেত্রে, & 4\_c\_4 & 1 & D,g,r এবং একটি e \\
2টি বর্ণ অভিন্ন 2টি ভিন্ন, এক্ষেত্রে, & 2\_c\_2×3\_c\_2 & 3 & দুটি e এবং D,g,r
তিনটি থেকে যে কোন দুটি \\
3টি অভিন্ন 1 টি ভিন্ন, এক্ষেত্রে, & 3\_c\_3×3\_c\_1 & 3 & ৩টি e এবং D,g,r
তিনটি থেকে যে কোন ১টি \\
\end{longtable}
}

\textbf{প্রশ্ন নম্বর- ১৯. THESIS শব্দটি অক্ষরগুলোর প্রতিবারে 4টি অক্ষর নিয়ে মোট
কত প্রকারে সমাবেশ করা যায়?}

\begin{itemize}
\item
  (ক) 12
\item
  (খ) 13
\item
  (গ) 11
\item
  (ঘ) 14
\end{itemize}

\textbf{\hl{উত্তরঃ (গ) 11}}

\textbf{Explanation:}

THESIS শব্দটিতে মোট ৬টি বর্ণের মধ্যে দুটি S আছে এবং ৪টি ভিন্ন বর্ণ ।

তাহলে ২টি S থেকে ১টি নিয়ে অন্য ৪টি সহ ভিন্ন ভিন্ন বর্ণ আছে ৫টি ।

{\def\LTcaptype{none} % do not increment counter
\begin{longtable}[]{@{}
  >{\raggedright\arraybackslash}p{(\linewidth - 6\tabcolsep) * \real{0.2006}}
  >{\raggedright\arraybackslash}p{(\linewidth - 6\tabcolsep) * \real{0.2006}}
  >{\raggedright\arraybackslash}p{(\linewidth - 6\tabcolsep) * \real{0.2006}}
  >{\raggedright\arraybackslash}p{(\linewidth - 6\tabcolsep) * \real{0.2006}}@{}}
\toprule\noalign{}
\begin{minipage}[b]{\linewidth}\raggedright
\textbf{শর্ত}
\end{minipage} & \begin{minipage}[b]{\linewidth}\raggedright
\textbf{বাছাইয়ের উপায়}
\end{minipage} & \begin{minipage}[b]{\linewidth}\raggedright
\textbf{বাছাই}
\end{minipage} & \begin{minipage}[b]{\linewidth}\raggedright
\textbf{যে অক্ষরগুলো নেয়া হয়েছে}
\end{minipage} \\
\midrule\noalign{}
\endhead
\bottomrule\noalign{}
\endlastfoot
4টি বর্ণই ভিন্ন এক্ষেত্রে, & 5\_c\_4 & 5 & T,H,E,I এবং একটি S \\
2টি ভিন্ন বর্ণ এবং 2টি অভিন্ন & 4\_c\_2×2\_c\_2 & 6 & T,H,E,I থেকে ২টি এবং
S দুটি । \\
\multicolumn{4}{@{}>{\raggedright\arraybackslash}p{(\linewidth - 6\tabcolsep) * \real{0.8022} + 6\tabcolsep}@{}}{%
সুতরাং মোট সংখ্যা = (5+6 ) = 11} \\
\end{longtable}
}

\textbf{প্রশ্ন নম্বর-} ২০. LOGARITHMS শব্দটির বর্ণগুলো হতে 3 টি Consonant ও
2টি Vowel কত প্রকারে বাছাই করা যায়? ...

\begin{itemize}
\item
  (ক) 109
\item
  (খ) 105
\item
  (গ) 103
\item
  (ঘ) 107
\end{itemize}

\textbf{\hl{উত্তরঃ (খ) 105}}

\textbf{Explanation:}

7 টি Consonant হতে 3টি করে বেছে নেওয়ার সংখ্যা = \(7_{c_{3}}\)

3টি Vowel হতে 2টি করে বেছে নেওয়ার সংখ্যা = \(3_{c_{2}}\)

:: মোট বাছাই সংখ্যা = \(7_{c_{3}} \times 3_{c_{2}}\) = 105

\textbf{প্রশ্ন নম্বর-} ২১. COMBINATION শব্দটি হতে 4 অক্ষর বিশিষ্ট সম্ভাব্য
সমাবেশ নির্ণয় করুন। ...

\begin{itemize}
\item
  (ক) 112
\item
  (খ) 118
\item
  (গ) 124
\item
  (ঘ) 136
\end{itemize}

\textbf{\hl{উত্তরঃ (ঘ) 136}}

\textbf{Explanation:}

COMBINATION শব্দটিতে মোট বর্ণ আছে ১১ টি, যার মধ্যে O আছে 2 টি, N আছে 2 টি,
এবং I আছে 2 টি ।

একজাতীয় সংখ্যাগুলো বাদ দিলে ভিন্ন সংখ্যা আছে ৫টি যথা: (C,M,B,A,T)

একজাতীয়গুলো থেকে ১টি করে নিলে ভিন্ন ভিন্ন মোট বর্ণ = ৮টি (C,O,M,B,I,N,A,T)

{\def\LTcaptype{none} % do not increment counter
\begin{longtable}[]{@{}
  >{\raggedright\arraybackslash}p{(\linewidth - 6\tabcolsep) * \real{0.2006}}
  >{\raggedright\arraybackslash}p{(\linewidth - 6\tabcolsep) * \real{0.2006}}
  >{\raggedright\arraybackslash}p{(\linewidth - 6\tabcolsep) * \real{0.2006}}
  >{\raggedright\arraybackslash}p{(\linewidth - 6\tabcolsep) * \real{0.2006}}@{}}
\toprule\noalign{}
\begin{minipage}[b]{\linewidth}\raggedright
\textbf{ক্ষেত্র}
\end{minipage} & \begin{minipage}[b]{\linewidth}\raggedright
\textbf{বাছাই}
\end{minipage} & \begin{minipage}[b]{\linewidth}\raggedright
\textbf{মোট}
\end{minipage} & \begin{minipage}[b]{\linewidth}\raggedright
\textbf{ব্যাখ্যা}
\end{minipage} \\
\midrule\noalign{}
\endhead
\bottomrule\noalign{}
\endlastfoot
4 টি ভিন্ন ভিন্ন বর্ণ নিয়ে সমাবেশ & 8\_c\_4 & 70 & (C,O,M,B,I,N,A,T) থেকে
৪টি \\
2 টি একজাতীয় ও 2 টি ভিন্ন & 3\_c\_1×7\_c\_2 & 63 & OO,NN, II, থেকে ২টি +
এবং ২টি একজাতীয় নিলে আরো ভিন্ন ভিন্ন বর্ণ থাকে ৮-১ = ৭টি। অন্য ৭টি থেকে
২টি \\
২টি একই অন্য ২টিও একই & 3\_c\_2 & 3 & OO,NN, II, ৩জোড়া থেকে ২ জোড়া \\
\multicolumn{4}{@{}>{\raggedright\arraybackslash}p{(\linewidth - 6\tabcolsep) * \real{0.8022} + 6\tabcolsep}@{}}{%
সুতরাং মোট সংখ্যা = (70+63+3 ) = 136} \\
\end{longtable}
}

\textbf{প্রশ্ন নম্বর-} ২২. একজন পরীক্ষার্থীকে 14 টি প্রশ্নের মধ্যে 6টি প্রশ্নের
উত্তর দিতে হবে। তাকে প্রথম 5টি থেকে অবশ্যই 4টি বাছাই করতে হবে। সে কত প্রকারে
প্রশ্নগুলো বাছাই করতে পারে? {[} থানা সহকারী শিক্ষা অফিসার পরীক্ষা- ১৯৯৯{]}

\begin{itemize}
\item
  (ক) 102
\item
  (খ) 105
\item
  (গ) 108
\item
  (ঘ) 180
\end{itemize}

\textbf{\hl{উত্তরঃ (ঘ) 180}}

\textbf{Explanation:}

এখানে প্রশ্ন বাছাইয়ের সংখ্যা হবে প্রথম ৫টি থেকে ৪টি এবং পরবর্তী 14-5 = 9টি
থেকে 6-4 = 2টির উত্তর দেয়া যাবে।

তাই উত্তর হবে

\(5_{c_{4}} \times 9_{c_{2}}\)

= 5×36 = 180

\textbf{প্রশ্ন নম্বর-} ২৩. 10 টি জিনিসের মধ্যে 2 টি এক জাতীয় এবং বাকীগুলো
ভিন্ন ভিন্ন জিনিস। ঐ জিনিসগুলো থেকে প্রতিবার 5টি নিয়ে কত প্রকারে বাছাই করা
যায়? {[} ৩৭-তম বিসিএস প্রিলি {]}

\begin{itemize}
\item
  (ক) 170
\item
  (খ) 182
\item
  (গ) 190
\item
  (ঘ) 192
\end{itemize}

\textbf{\hl{উত্তরঃ (খ) 182}}

\textbf{Explanation:}

10 টি জিনিসের মধ্যে 2টি একজাতীয় এবং বাছাই করতে হবে 5টি। বাছাই করার উপায়
আছে দুভাবে,

(i) 5 টির সবগুলোই ভিন্ন ভিন্ন ।

(ii) একজাতীয় 2টি একসাথে এবং অন্য 3টি ভিন্ন ।

(i) এর ক্ষেত্রে 10টি জিনিসের মধ্যে যে 2 টি জিনিস একই রকম তাদেরকে একটি ধরে
মোট 9টি ভিন্ন ভিন্ন জিনিস থেকে 5টি জিনিস বাছাই করা যায়

= \(9_{c_{5}}\)

= \(\frac{9!}{5!(9 - 5)!}\)

= \(\frac{9 \times 8 \times 7 \times 6!}{5! \times 4!}\)

= \(\frac{9 \times 8 \times 7 \times 6}{24}\)

= 126 (ii) এর ক্ষেত্রে প্রথমে, ২টি একজাতীয় এর ২টি ই নেয়ার পর অবশিষ্ট ১০-২ =
৮টি থেকে পরের (৫-২)= ৩টি নেয়া যায়

\(2_{c_{2}} \times 8_{c_{3}}\)

= \(\frac{2!}{2!(2 - 2)!} \times \frac{8!}{3!(8 - 3)!}\)

=
\(\frac{2!}{2! \times 0!} \times \frac{8 \times 7 \times 6 \times 5!}{6 \times 5!}\)

= 1×56 = 56 সুতরাং মোট বাছাই করা যাবে 126+56 = 182 ভাবে ।

\section{10. বিন্যাস ও সমাবেশ MCQ প্রশ্ন উত্তর
অনুশীলন}\label{ux9acux9a8ux9afux9b8-ux993-ux9b8ux9aeux9acux9b6-mcq-ux9aaux9b0ux9b6ux9a8-ux989ux9a4ux9a4ux9b0-ux985ux9a8ux9b6ux9b2ux9a8}

\textbf{শাহাবাগ থেকে ফার্মগেটে যাওয়ার তিনটি ভিন্ন রাস্তা আছে, আবার ফার্মগেট
থেকে বনানীর ৪টি ভিন্ন রাস্তা আছে। শাহবাগ থেকে ফার্মগেট হয়ে বনানী যাবার কয়টি
ভিন্ন রাস্তা আছে?}

ক) ১০০

\textbf{\hl{খ) ১২ ✔}}

গ) ১৩

ঘ) ১৪

\textbf{\hl{✔ সঠিক উত্তর: খ) ১২}}

\begin{quote}
\textbf{ব্যাখ্যা:} \emph{শাহবাগ থেকে ফার্মগেট যাওয়ার রাস্তা ৩টি এবং ফার্মগেট
থেকে বনানী যাওয়ার রাস্তা ৪টি।}

\emph{∴ শাহবাগ থেকে ফার্মগেট হয়ে বনানী যাওয়ার মোট রাস্তা = 3 × 4 = 12টি}
\end{quote}

\textbf{In how many ways can the letters of the word
\textquotesingle APPLE\textquotesingle{} be arranged? / কত উপায়ে
\textquotesingle APPLE\textquotesingle{} শব্দটির অক্ষরগুলোকে সাজানো যেতে
পারে ?}

ক) 720

খ) 120

\textbf{\hl{গ) 60 ✔}}

ঘ) 180

\textbf{\hl{✔ সঠিক উত্তর: গ) 60}}

\begin{quote}
\textbf{ব্যাখ্যা:} \emph{APPLE শব্দটিতে মোট ৫টি অক্ষর আছে, যার মধ্যে P অক্ষরটি
২ বার আছে।}

\emph{∴ বিন্যাস সংখ্যা = (5!/2!) = (120/2) = 60}
\end{quote}

\textbf{A,B,C বর্ণের ৩টি করে বর্ণ নিয়ে কত প্রকারে বিন্যাস করা যায়?}

ক) ৩ উপায়ে

খ) ৪ উপায়ে

গ) ৫ উপায়ে

\textbf{\hl{ঘ) ৬ উপায়ে ✔}}

\textbf{\hl{✔ সঠিক উত্তর: ঘ) ৬ উপায়ে}}

\begin{quote}
\textbf{ব্যাখ্যা:} \emph{A, B, C তিনটি ভিন্ন বর্ণ থেকে ৩টি করে বর্ণ নিয়ে
বিন্যাস সংখ্যা হলো ৩টি বস্তুর বিন্যাস বা ৩! (Factorial 3)।}

\emph{∴ বিন্যাস সংখ্যা = 3 × 2 × 1 = 6}
\end{quote}

\textbf{CALCUTTA শব্দটির বর্ণগুলোকে একত্রে নিয়ে বিন্যাস সংখ্যা AMERICA শব্দটির
বর্ণগুলো একত্রে নিয়ে বিন্যাস সংখ্যার কত গুণ?}

\textbf{\hl{ক) ২ ✔}}

খ) ৩

গ) ৪

ঘ) ৫

\textbf{\hl{✔ সঠিক উত্তর: ক) ২}}

\begin{quote}
\textbf{ব্যাখ্যা:} \emph{CALCUTTA শব্দটিতে মোট ৮টি বর্ণ আছে, যেখানে C আছে
২বার, A আছে ২বার, T আছে ২বার।}

\emph{∴ CALCUTTA-এর বিন্যাস = (8!/2!×2!×2!) = (40320/8) = 5040}

\emph{AMERICA শব্দটিতে মোট ৭টি বর্ণ আছে, যেখানে A আছে ২বার।}

\emph{∴ AMERICA-এর বিন্যাস = (7!/2!) = (5040/2) = 2520}

\emph{∴ নির্ণেয় গুণ = (5040/2520) = 2}
\end{quote}

\textbf{PSC শব্দটির সবগুলো বর্ণ একত্রে নিয়ে কত প্রকারে সাজানো যায়?}

ক) ১

খ) ৩

\textbf{\hl{গ) ৬ ✔}}

ঘ) ৯

\textbf{\hl{✔ সঠিক উত্তর: গ) ৬}}

\begin{quote}
\textbf{ব্যাখ্যা:} \emph{PSC শব্দটিতে ৩টি ভিন্ন বর্ণ রয়েছে। ৩টি ভিন্ন বর্ণকে
সাজানোর উপায় হলো ৩! (Factorial 3)।}

\emph{∴ বিন্যাস সংখ্যা = 3 × 2 × 1 = 6}
\end{quote}

\textbf{FREEDOM শব্দটির সবগুলো বর্ণ একত্রে নিয়ে কত প্রভাবে সাজানো যায়?}

\textbf{\hl{ক) (7!/2!) ✔}}

খ) (7!/5!)

গ) (5!/2!)

ঘ) (7!/2!5!)

\textbf{\hl{✔ সঠিক উত্তর: ক) (7!/2!)}}

\begin{quote}
\textbf{ব্যাখ্যা:} \emph{FREEDOM শব্দটিতে মোট ৭টি বর্ণ আছে, যার মধ্যে E বর্ণটি
২ বার আছে।}

\emph{∴ বর্ণগুলোর বিন্যাস সংখ্যা = (7!/2!)}
\end{quote}

\textbf{n\_P\_4=12n\_P\_2 হলে, n = কত?}

ক) 1

\textbf{\hl{খ) 6 ✔}}

গ) 2

ঘ) 8

\textbf{\hl{✔ সঠিক উত্তর: খ) 6}}

\begin{quote}
\textbf{ব্যাখ্যা:} \emph{n\_P\_4=12n\_P\_2}

\emph{বা, (n!/(n−4)!) = 12 (n!/(n−2)!)}

\emph{বা, (n(n−1)(n−2)(n−3)(n−4)!/(n−4)!) = 12 (n(n−1)(n−2)!/(n−2)!)}

\emph{বা, n(n − 1)(n − 2)(n − 3) = 12n(n − 1)}

\emph{বা, (n − 2)(n − 3) = (12n(n−1)/n(n−1))}

\emph{বা, (n − 2)(n − 3) = 12}

\emph{বা, n\^{}2 − 3n − 2n + 6 = 12}

\emph{বা, n\^{}2 − 5n − 6 = 0}

\emph{বা, n\^{}2 − 6n + n − 6 = 0}

\emph{বা, n(n − 6) + 1(n − 6) = 0}

\emph{বা, (n − 6)(n + 1) = 0}

\emph{∴ n = 6, n = 1 {[}এটা সম্ভব নয়{]}}
\end{quote}

\textbf{Rahim wanted to arrange 3 of her four plants in a row on a
shelf. If each of the plants is in a different colour container, how
many different arrangements can he make? / রহিম তার চারটি চারাগাছের মধ্যে
থেকে তিনটিকে এক সারিতে একটি শেলফ এর উপর সাজাতে চাইল। যদি প্রত্যেকটি চারা
গাছ ভিন্ন ভিন্ন রঙের পাত্রে থাকে, তবে সে চারাগাছগুলোকে কত বিভিন্ন উপায়ে
সাজাতে পারবে?}

ক) ৭

খ) ১২

\textbf{\hl{গ) ২৪ ✔}}

ঘ) ২৮

\textbf{\hl{✔ সঠিক উত্তর: গ) ২৪}}

\begin{quote}
\textbf{ব্যাখ্যা:} \emph{চারটি ভিন্ন চারাগাছ হতে তিনটি নিয়ে সাজানোর উপায়
হলো 4\_P\_3}

\emph{∴ 4\_P\_3 = 4 × 3 × 2 = 24 উপায়ে।}
\end{quote}

\textbf{AMERICA শব্দটির বর্ণগুলো থেকে প্রতিবারে ৩টি বর্ণ নিয়ে গঠিত ভিন্ন ভিন্ন
শব্দ সংখ্যা কত হবে?}

ক) 130

\textbf{\hl{খ) 135 ✔}}

গ) 140

ঘ) 145

\textbf{\hl{✔ সঠিক উত্তর: খ) 135}}

\begin{quote}
\textbf{ব্যাখ্যা:} \emph{AMERICA শব্দটিতে মোট ৭টি বর্ণ রয়েছে, যার মধ্যে দুটি A
এবং বাকি ৫টি ভিন্ন বর্ণ (M, E, R, I, C)।}

\emph{ক্ষেত্র ১: ৩টি বর্ণই ভিন্ন হওয়ার উপায় = 6\_P\_3 = 6 × 5 × 4 = 120}

\emph{ক্ষেত্র ২: একটি A এবং অন্য দুটি ভিন্ন বর্ণ নিয়ে শব্দ সংখ্যা = 5 × 3\_P\_2
= 5 × 3 = 15}

\emph{∴ মোট শব্দ সংখ্যা = 120 + 15 = 135}
\end{quote}

\textbf{SCIENCE শব্দটির স্বরবর্ণগুলোকে একত্রে রেখে সব কয়টি বর্ণকে সম্ভাব্য যত
উপায়ে সাজানো যায় তার সংখ্যা হচ্ছে?}

ক) 60 বার

খ) 120 বার

\textbf{\hl{গ) 180 বার ✔}}

ঘ) 420 বার

\textbf{\hl{✔ সঠিক উত্তর: গ) 180 বার}}

\begin{quote}
\textbf{ব্যাখ্যা:} \emph{SCIENCE শব্দটিতে স্বরবর্ণ ৩টি (I, E, E) এবং ব্যঞ্জনবর্ণ
৪টি (S, C, C, N), যেখানে 2টি E এবং 2টি C আছে।}

\emph{স্বরবর্ণ ৩টিকে একটি অক্ষর ধরে মোট অক্ষর হয় 5টি, যার মধ্যে C আছে ২টি।}

\emph{∴ বিন্যাস সংখ্যা = (5!/2!) = (120/2) = 60}

\emph{আবার, স্বরবর্ণ ৩টিকে নিজেদের মধ্যে সাজানো যায় = (3!/2!) = (6/2) = 3}

\emph{∴ স্বরবর্ণগুলোকে একত্রে রেখে মোট বিন্যাস সংখ্যা = 60 × 3 = 180}
\end{quote}

\textbf{n\_C\_r+n\_(C\_(r − 1)) এর সঠিক মান কোনটি?}

ক) n\_(C\_(r + 1))

\textbf{\hl{খ) (n + 1)\_C\_r ✔}}

গ) (n + 1)\_(C\_(r + 1))

ঘ) (n + 1)\_(C\_(r − 1))

\textbf{\hl{✔ সঠিক উত্তর: খ) (n + 1)\_C\_r}}

\begin{quote}
\textbf{ব্যাখ্যা:} \emph{n\_C\_r+n\_(C\_(r − 1))=(n + 1)\_C\_r
{[}অনুসিদ্ধান্ত{]}}
\end{quote}

\textbf{৬ জন খেলোয়াড়কে সমান সংখ্যক দুটি দলে কত ভাবে বিভক্ত করা যায়?}

ক) ১০

\textbf{\hl{খ) ২০ ✔}}

গ) ৬০

ঘ) ১২০

\textbf{\hl{✔ সঠিক উত্তর: খ) ২০}}

\begin{quote}
\textbf{ব্যাখ্যা:} \emph{৬ জন খেলোয়াড়কে সমান সংখ্যক দুটি দলে ভাগ করলে, প্রতি
দলে ৩ জন খেলোয়াড় নেওয়া যাবে।}

\emph{∴ মোট বিভক্ত করা যায় = 6\_C\_3 = (6!/3!×3!) = (6×5×4×3×2/36) = 20
ভাবে।}
\end{quote}

\textbf{৪ জন মহিলা ও ৬ জন পুরুষের মধ্য থেকে ৪ সদস্যবিশিষ্ট একটি উপ-কমিটি গঠন
করতে হবে যাতে ১জন নির্দিষ্ট পুরুষ সর্বদাই উপস্থিত থাকেন। কত প্রকারের ঐ কমিটি
গঠন করা যেতে পারে?}

ক) ২১০

খ) ৩০৪

\textbf{\hl{গ) ৮৪ ✔}}

ঘ) ১২০

\textbf{\hl{✔ সঠিক উত্তর: গ) ৮৪}}

\begin{quote}
\textbf{ব্যাখ্যা:} \emph{১ জন নির্দিষ্ট পুরুষ সর্বদাই উপস্থিত থাকায় বাকি থাকে
4−1 = 3 জন সদস্য এবং মোট লোকসংখ্যা 4+5 = 9 জন (১ জন নির্দিষ্ট পুরুষ বাদ
দিয়ে)।}

\emph{∴ কমিটি গঠনের উপায় = 9\_C\_3 = (9!/3!×6!) = (9×8×7/3×2×1) = 84।}
\end{quote}

\textbf{১০টি জিনিসের মধ্যে ২টি এক জাতীয় এবং বাকিগুলো ভিন্ন ভিন্ন জিনিস। ঐ
জিনিসগুলো থেকে প্রতিবারে ৫ টি নিয়ে কত প্রকারে বাছাই করা যায়?}

ক) ১৭০

\textbf{\hl{খ) ১৮২ ✔}}

গ) ১৯০

ঘ) ১৯২

\textbf{\hl{✔ সঠিক উত্তর: খ) ১৮২}}

\begin{quote}
\textbf{ব্যাখ্যা:} \emph{এক জাতীয় দুটির একটিও থাকবে না তাদের সংখ্যা =
8\_C\_5 = (8!/(8−5)!5!) = 56}

\emph{এক জাতীয় দুটির একটি থাকবে তাদের সংখ্যা = 8\_C\_4 = (8!/(8−4)!4!) =
70}

\emph{এক জাতীয় দুটির উভয়টি থাকবে তাদের সংখ্যা = 8\_C\_3 = (8!/(8−3)!3!) =
56}

\emph{∴ ১০টি জিনিসের মধ্যে ৫ টি জিনিস বাছাই করার উপায় = 56 + 70 + 56 =
182}
\end{quote}

\textbf{১২ টি পুস্তক থেকে ৫ টি কত প্রকারে বাছাই করা যায় যেখানে ২টি পুস্তক
সর্বদা অন্তর্ভুক্ত থাকবে?}

ক) ২৫২

খ) ৭৯২

গ) ২২৪

\textbf{\hl{ঘ) ১২০ ✔}}

\textbf{\hl{✔ সঠিক উত্তর: ঘ) ১২০}}

\begin{quote}
\textbf{ব্যাখ্যা:} \emph{১২ টি পুস্তক হতে ৫ টি বাছাই করা যায় যেখানে পুস্তক
সর্বদা অন্তর্ভুক্ত থাকে তা হচ্ছে 10\_C\_3}

\emph{= (10!/(10−3)!3!) = (10×9×8/1×2×3) = 120}
\end{quote}

\textbf{১৪ জন খেলোয়াড়ের মধ্যে থেকে নির্দিষ্ট একজন অধিনায়কসহ ১১ জনের একটি
ক্রিকেট দল কতভাবে বাছাই করা যাবে?}

ক) ৭২৮

\textbf{\hl{খ) ২৮৬ ✔}}

গ) ৩৬৪

ঘ) ১০০১

\textbf{\hl{✔ সঠিক উত্তর: খ) ২৮৬}}

\begin{quote}
\textbf{ব্যাখ্যা:} \emph{যেহেতু একজন অধিনায়ক সবসময় থাকবেন, সেহেতু অবশিষ্ট ১৩
জন খেলোয়াড় হতে ১০ জন খেলোয়াড় বেছে নিতে হবে।}

\emph{∴ বাছাই করার উপায় = 13\_C\_10 = C313 = (13×12×11/3×2×1) = 286}
\end{quote}

\textbf{একটি ত্রিদেশীয় ক্রিকেট খেলায় প্রত্যেকে অন্যদেশের সাথে একবার মাত্র
খেলবে। মোট কয়টি খেলা অনুষ্ঠিত হবে?}

ক) ৭টি

খ) ৫টি

\textbf{\hl{গ) ৩টি ✔}}

ঘ) ৪টি

\textbf{\hl{✔ সঠিক উত্তর: গ) ৩টি}}

\begin{quote}
\textbf{ব্যাখ্যা:} \emph{একটি ত্রিদেশীয় (৩টি দল নিয়ে) ক্রিকেট খেলায় প্রত্যেকে
অন্যদেশের সাথে একবার মাত্র খেললে মোট খেলার সংখ্যা = 3\_C\_2 টি।}

\emph{∴ (3!/2!(3−2)!) টি = (6/2) = 3 টি।}
\end{quote}

\textbf{২০ সদস্যবিশিষ্ট একটি ফুটবল দল থেকে একজন অধিনায়ক ও একজন সহ-অধিনায়ক
কতভাবে নির্বাচন করা যাবে?}

ক) ২০

\textbf{\hl{খ) ১৯০ ✔}}

গ) ৩৮০

ঘ) ৭৬০

\textbf{\hl{✔ সঠিক উত্তর: খ) ১৯০}}

\begin{quote}
\textbf{ব্যাখ্যা:} \emph{দল গঠনের সংখ্যা = (20\_C\_1×19\_C\_1/2) =
(20×19/2) = 190}
\end{quote}

\textbf{একটি ফুটবল টুর্নামেন্টে ৬টি দল অংশগ্রহণ করেছে। একক লীগ পদ্ধতিতে খেলা
হলে মোট কতটি খেলা পরিচালনা করতে হবে?}

ক) ৬টি

খ) ১২টি

\textbf{\hl{গ) ১৫টি ✔}}

ঘ) ৫টি

\textbf{\hl{✔ সঠিক উত্তর: গ) ১৫টি}}

\begin{quote}
\textbf{ব্যাখ্যা:} \emph{দুটি দল মিলে একটি খেলা অনুষ্ঠিত হবে।}

\emph{∴ মোট খেলার সংখ্যা = 6\_C\_2}

\emph{= (6!/2!4!) = (6×5×4!/2×1×4!) = 15}
\end{quote}

\textbf{একটি নৌকায় ৩ জন লোক উঠতে পারে, নৌকাটি দ্বারা ৪ জন লোক কত প্রকারে
উঠে নদী পার হতে পারে?}

ক) 1

খ) 2

গ) 3

\textbf{\hl{ঘ) 4 ✔}}

\textbf{\hl{✔ সঠিক উত্তর: ঘ) 4}}

\begin{quote}
\textbf{ব্যাখ্যা:} \emph{৪ জন লোক নদীটি পার হতে পারে 4\_C\_3 উপায়ে।}

\emph{∴ (4!/3!(4−3)!) = (4/1) = 4}
\end{quote}

\textbf{৫, ৯, ১, ৪ অংকগুলি দ্বারা ৫,০০০ এর চেয়ে বড় কতগুলো সংখ্যা তৈরি করা
যায়?}

ক) ১২টি

\textbf{\hl{খ) ৮টি ✔}}

গ) ১৮টি

ঘ) ১৬টি

\textbf{\hl{✔ সঠিক উত্তর: খ) ৮টি}}

\begin{quote}
\textbf{ব্যাখ্যা:} \emph{৫, ৯, ১, ৪ এই চারটি অংক দ্বারা ৫,০০০ এর চেয়ে বড়
কোনো সংখ্যা তৈরি করতে হলে প্রথম স্থানে ৫ বা ৯ বসাতে হবে। ১ম স্থানটি ৫ ও ৯
দ্বারা 2\_P\_1 বা ২ প্রকারে পূরণ করা যায় বাকী তিনটি স্থান তিনটি অঙ্ক দ্বারা
3\_P\_2 বা ৬ প্রকারে পূরণ করা যায়।}

\emph{মোট বিন্যাস ২ × ৬ = ১২}
\end{quote}

\textbf{From 7 boys and 5 girls, 6 are to be selected for admission for
a particular course. In how many ways can this be done if there must be
exactly 3 girls?}

\textbf{\hl{ক) A. 350 ✔}}

খ) B. 65

গ) C. 2050

ঘ) D. 120

ঙ) E. 250

\textbf{\hl{✔ সঠিক উত্তর: ক) A. 350}}

\begin{quote}
\textbf{ব্যাখ্যা:} \emph{3 জন মেয়ে বাছাই করা হলে, বাকি থাকে 6−3 = 3 জন
ছেলে।}

\emph{7 জন ছেলের মধ্য থেকে 3 জন ছেলে বাছাই করা যায় = 7\_C\_3 =
(7×6×5/1×2×3) = 35 উপায়ে}

\emph{5 জন মেয়ের মধ্য থেকে 3 জন মেয়ে বাছাই করা যায় = 5\_C\_3 =
(5×4×3/1×2×3) = 10 উপায়ে}

\emph{∴ মোট উপায়ের সংখ্যা = 35 × 10 = 350 উপায়ে}
\end{quote}

\textbf{2n\_C\_r=2n\_(C\_(r+2)) হলে r এর মান কত?}

ক) 2n−1

খ) n

গ) n+1

\textbf{\hl{ঘ) n−1 ✔}}

\textbf{\hl{✔ সঠিক উত্তর: ঘ) n−1}}

\begin{quote}
\textbf{ব্যাখ্যা:} \emph{আমরা জানি, Cxm=Cym হলে হয় x=y অথবা x+y=m .}

\emph{এখানে r≠r+2 , সুতরাং r+(r+2)=2n}

\emph{বা, 2r+2=2n}

\emph{বা, 2r=2n−2}

\emph{∴ r=n−1}
\end{quote}

\textbf{How many different committees of 3 members can be chosen out of
5 person in a group so that one particular person is always chosen? / ৫
জন ব্যক্তি হতে ৩ সদস্যের একটি কমিটি কত বিভিন্ন উপায়ে বেছে নেওয়া যায় যাতে
একজন নির্দিষ্ট ব্যক্তি সবসময় উক্ত কমিটিতে থাকবে?}

ক) 4

খ) 5

\textbf{\hl{গ) 6 ✔}}

ঘ) 8

ঙ) None of them

\textbf{\hl{✔ সঠিক উত্তর: গ) 6}}

\begin{quote}
\textbf{ব্যাখ্যা:} \emph{যেহেতু একজন ব্যক্তি সবসময় কমিটিতে থাকবেন, সেহেতু
অবশিষ্ট 5−1 = 4 জন ব্যক্তি হতে কমিটির অপর দুই সদস্য বেছে নিতে হবে।}

\emph{∴ অপর দুই সদস্য বেছে নেওয়া যায় = 4\_C\_2 = (4×3/1×2) = 6 ভাবে।}
\end{quote}

\textbf{When a meeting of a committee was over, every member shook hands
with another and the total number of hand shake was 6. How many members
were there in the meeting? / একটি কমিটির মিটিং শেষে প্রত্যেক সদস্য একে
অপরের করমর্দন করে। যদি মোট করমর্দন সংখ্যা ৬ হয়, তবে মিটিং এ কতজন সদস্য
উপস্থিত ছিল?}

ক) 2

খ) 3

\textbf{\hl{গ) 4 ✔}}

ঘ) 5

ঙ) 6

\textbf{\hl{✔ সঠিক উত্তর: গ) 4}}

\begin{quote}
\textbf{ব্যাখ্যা:} \emph{ধরি, উপস্থিত সদস্য সংখ্যা = n}

\emph{প্রতিটি করমর্দনের জন্য দুইজন ব্যক্তির প্রয়োজন পড়ে।}

\emph{n জন ব্যক্তি করমর্দন করতে পারে = n\_C\_2 ভাবে = (n(n−1)/2)}

\emph{শর্তমতে, (n(n−1)/2) = 6}

\emph{বা, n\^{}2−n=12}

\emph{বা, n\^{}2−n−12=0}

\emph{বা, (n−4)(n+3)=0}

\emph{∴ n=4,−3 {[}−3 গ্রহণযোগ্য নয়{]}}

\emph{উপস্থিত সদস্য সংখ্যা = 4 জন।}
\end{quote}

\textbf{From 6 boys and 4 girls, 5 are to be selected for admission for
a particular course. In how many ways can this be done if there must be
exactly 2 girls? / ৬ জন বালক ও ৪ জন বালিকা হতে ৫ জনকে কতভাবে ভর্তির জন্য
নির্বাচিত করা যাবে, যাতে ভর্তিকৃতদের মধ্যে সর্বদা শুধুমাত্র ২জন বালিকা থাকে?}

ক) 60

খ) 30

গ) 90

\textbf{\hl{ঘ) 120 ✔}}

ঙ) 15

\textbf{\hl{✔ সঠিক উত্তর: ঘ) 120}}

\begin{quote}
\textbf{ব্যাখ্যা:} \emph{ভর্তিকৃতদের মধ্যে বালকের সংখ্যা = 5−2 = 3 জন}

\emph{6 জন বালক হতে 3 জন বালককে নির্বাচিত করা যায় = 6\_C\_3 =
(6×5×4/1×2×3) = 20 ভাবে।}

\emph{4 জন বালিকা হতে 2 জন বালিকাকে নির্বাচিত করা যায় = 4\_C\_2 =
(4×3/1×2) = 6 ভাবে।}

\emph{∴ সর্বমোট উপায় = 20×6 = 120 ভাবে।}
\end{quote}

\textbf{10 জন বালক ও 8 জন বালিকা থেকে 2 জন বালক ও 2 জন বালিকা কত বিভিন্ন
উপায়ে বেছে নেওয়া যায়? / In how many ways can 2 boys and 2 girls be
selected from 10 boys and 8 girls?}

ক) 73

\textbf{\hl{খ) 1260 ✔}}

গ) 3060

ঘ) 5060

\textbf{\hl{✔ সঠিক উত্তর: খ) 1260}}

\begin{quote}
\textbf{ব্যাখ্যা:} \emph{10 জন বালক হতে প্রতিবার 2 জন বালক বেছে নেওয়া যায়,
10\_C\_2 = (10×9/1×2) = 45 উপায়ে।}

\emph{8 জন বালিকা হতে প্রতিবার 2 জন বালিকা বেছে নেওয়া যায়, 8\_C\_2 =
(8×7/1×2) = 28 উপায়ে।}

\emph{∴ মোট বেছে নেওয়া যায় = 45×28 = 1260 উপায়ে।}
\end{quote}

\textbf{একজন পরীক্ষার্থীকে ১২টি প্রশ্ন থেকে ৬টি প্রশ্নের উত্তর করতে হবে। প্রথম
৫টি থেকে ঠিক ৪টি প্রশ্ন বাছাই করে কত প্রকারে ৬টি প্রশ্ন উত্তর করা যাবে?}

ক) 100

\textbf{\hl{খ) 105 ✔}}

গ) 110

ঘ) 220

\textbf{\hl{✔ সঠিক উত্তর: খ) 105}}

\begin{quote}
\textbf{ব্যাখ্যা:} \emph{পরীক্ষার্থীকে প্রথম ৫টি থেকে ৪টি এবং অবশিষ্ট ( 12−5 )
= 7টি থেকে ২টি প্রশ্ন বাছাই করে উত্তর করতে হবে।}

\emph{প্রথম ৫ টি প্রশ্ন থেকে ৪টি বাছাই করা যায়, 5\_C\_4 = C15 = (5/1) = 5
উপায়ে।}

\emph{অবশিষ্ট 7 টি প্রশ্ন থেকে ২টি বাছাই করা যায়, 7\_C\_2 = (7×6/1×2) = 21
উপায়ে।}

\emph{∴ মোট প্রশ্ন বাছাই করা যায় = 5×21 = 105 উপায়ে।}
\end{quote}

\textbf{How many combination are possible if a persons has 4 sports
jackets, 5 shirts and 3 pairs of socks?/ এক ব্যক্তির ৪ টি খেলার জ্যাকেট, ৫
টি শার্ট এবং ৩ জোড়া মোজা আছে। কত উপায়ে এগুলোকে বাছাই করা যায়?}

ক) 4

খ) 5

গ) 12

\textbf{\hl{ঘ) 60 ✔}}

ঙ) 15

\textbf{\hl{✔ সঠিক উত্তর: ঘ) 60}}

\begin{quote}
\textbf{ব্যাখ্যা:} \emph{4 টি খেলার জ্যাকেট বাছাই করা যায় = 4\_C\_1 উপায়ে =
4 উপায়ে}

\emph{5 টি শার্ট বাছাই করা যায় = 5\_C\_1 উপায়ে = 5 উপায়ে}

\emph{3 জোড়া মোজা বাছাই করা যায় = 3\_C\_1 উপায়ে = 3 উপায়ে}

\emph{∴ সর্বমোট বাছাই করা যায় = 4×5×3 উপায়ে = 60 উপায়ে।}
\end{quote}

\textbf{From a group of 7 men and 6 women 5 person are to be selected to
form a committee so that at least three men are there on the committee.
In how many ways can it be done?/ ৭ জন পুরুষ ও ৬ জন মহিলার একটি দল হতে ৫
সদস্যের একটি কমিটি কতভাবে নির্বাচিত করা যায় যাতে সবসময় কমিটিতে অন্তত ৩ জন
পুরুষ থাকে?}

\textbf{\hl{ক) 756 ✔}}

খ) 735

গ) 564

ঘ) 645

ঙ) 800

\textbf{\hl{✔ সঠিক উত্তর: ক) 756}}

\begin{quote}
\textbf{ব্যাখ্যা:} \emph{যেহেতু ৫ সদস্যের কমিটিতে সর্বদা অন্তত তিনজন পুরুষ থাকে,
সেহেতু নিম্নোক্তভাবে কমিটি গঠন সম্ভব:}

\emph{১) 3 জন পুরুষ ও 2 জন মহিলাবিশিষ্ট কমিটি : 7\_C\_3×6\_C\_2 =
(7×6×5/3×2×1)×(6×5/2×1) = 525 উপায়ে}

\emph{২) 4 জন পুরুষ ও 1 জন মহিলাবিশিষ্ট কমিটি : 7\_C\_4×6\_C\_1 =
(7×6×5×4/4×3×2×1)×(6/1) = 210 উপায়ে}

\emph{৩) 5 জন পুরুষ ও 0 জন মহিলাবিশিষ্ট কমিটি : 7\_C\_5×6\_C\_0 =
(7×6/2×1)×1 = 21 উপায়ে}

\emph{∴ সর্বমোট উপায় = 525+210+21 = 756}
\end{quote}

\textbf{36 identical chair must be arranged in rows with equal number of
chairs in each row. Each row must contains atleast 3 chairs and there
must be atleast 3 rows. How many different arrangements are possible?/
৩৬ টি একই ধরনের চেয়ার কত বিভিন্ন উপায়ে সাজানো যায় যাতে প্রতিটি সারিতে
কমপক্ষে ৩ টি চেয়ার থাকে এবং সারির সংখ্যা কমপক্ষে ৩ হয়। উল্লেখ্য প্রতিটি সারিতে
চেয়ারের সংখ্যা সমান?}

ক) 2

খ) 4

\textbf{\hl{গ) 5 ✔}}

ঘ) 6

ঙ) 10

\textbf{\hl{✔ সঠিক উত্তর: গ) 5}}

\begin{quote}
\textbf{ব্যাখ্যা:} \emph{মোট চেয়ার = সারির সংখ্যা × প্রতি সারিতে চেয়ারের
সংখ্যা}

\emph{৩৬ কে এমন দুটি সংখ্যার গুণফল আকারে প্রকাশ করতে হবে যেন প্রতিটি সংখ্যা
অন্তত ৩ হয়: 36 = 3 × 12 36 = 4 × 9 36 = 6 × 6 36 = 9 × 4 36 = 12 × 3}

\emph{∴ মোট ৫ উপায়ে সাজানো যায়।}
\end{quote}

\section{11. ত্রিভুজ কাকে বলে? ত্রিভুজের প্রকারভেদ, বৈশিষ্ট্য ও অনুসিদ্ধান্ত
(পূর্ণাঙ্গ
গাইড)}\label{ux9a4ux9b0ux9adux99c-ux995ux995-ux9acux9b2-ux9a4ux9b0ux9adux99cux9b0-ux9aaux9b0ux995ux9b0ux9adux9a6-ux9acux9b6ux9b7ux99fux9af-ux993-ux985ux9a8ux9b8ux9a6ux9a7ux9a8ux9a4-ux9aaux9b0ux9a3ux999ux997-ux997ux987ux9a1}

ত্রিভুজ হলো তিনটি বাহু বা সাইড দ্বারা গঠিত একটি আবদ্ধ চিত্র বা ক্ষেত্র। যা তিনটি
বিন্দুতে স্থাপিত হয় । এটি ব্যবহার করা হয় বিভিন্ন ধরনের গণিত ও জ্যামিতিক সমস্যার
সমাধানের জন্য। একটি ত্রিভুজের মোট তিনটি কোণ এবং তিনটি বাহু থাকে।

যেহেতু তিনটি বিন্দু সমান্তরাল হতে পারে না, সুতরাং ত্রিভুজের ক্ষেত্রফল হলো বিন্দু
দুইটির লম্বদূরত্ব গুণ করে সংখ্যাগুলোর সমান ভাগফল। ত্রিভুজের বাহুর দৈর্ঘ্য বা কোণের
মান জানতে হলে ত্রিভুজের অংকে প্রয়োজনীয় সূত্রগুলো ব্যবহার করা হয়।

ত্রিভুজ বিষয়ক জ্যামিতি অনেক কঠিন হওয়ার কারণে এর সমস্যা সমাধান এবং উপযোগিতা
বিষয়ক অনেক গবেষণা চলছে। তবে ত্রিভুজ বিষয়ক জ্ঞান সম্পর্কে স্পষ্ট ধারণা নিতে হলে
আজকের এই টিউটরিয়ালটি আপনার মনোযোগ সহকারে পড়া উচিত।

\begin{itemize}
\item
  প্রাথমিক আলোচনা
\item
  ত্রিভুজ বিষয়ক গুরুত্বপূর্ণ অনুসিদ্ধান্ত
\item
  ত্রিভুজের সর্বসমতা
\item
  ত্রিভুজের কেন্দ্র
\item
  ত্রিভুজের প্রকারভেদ
\end{itemize}

\begin{itemize}
\item
  সমকোণী ত্রিভুজ
\item
  সূক্ষ্মকোণী ত্রিভুজ
\item
  স্থূলকোণী ত্রিভুজ
\item
  সমবাহু ত্রিভুজ
\item
  সমদ্বিবাহু ত্রিভুজ
\item
  বিষমবাহু ত্রিভুজ
\end{itemize}

\begin{itemize}
\item
  বিভিন্ন প্রতিযোগিতামূলক পরীক্ষার প্রশ্ন ও সমাধান
\end{itemize}

\subsubsection{প্রাথমিক
আলোচনা}\label{ux9aaux9b0ux9a5ux9aeux995-ux986ux9b2ux99aux9a8-3}

\textbf{ত্রিভুজ:} তিনটি বাহু দ্বারা আবদ্ধ চিত্র বা ক্ষেত্রকে ত্রিভূজ বলে।

চিত্রে ABC একটি ত্রিভুজ। এখানে AB, BC ও AC তিনটি বাহু দ্বারা ABC একটি আবদ্ধ
চিত্র হয়েছে যাকে ত্রিভুজ বলা যায় ।

ত্রিভুজ সম্পর্কে বিস্তারিত জানার আগে আমরা ত্রিভুজের সাথে সংশ্লিষ্ট কয়েকটি টার্ম
জেনে নেই । যথা:

\textbf{ভূমি (Base):} চিত্রে BC হলো ভূমি।

\textbf{শীর্ষবিন্দু (Vertex):} চিত্রে A হল শীর্ষবিন্দু।

\textbf{অন্তঃস্থ কোণ:} চিত্রে ABC ত্রিভুজ ক্ষেত্রের ভিতরের যে তিনটি কোণ ∠A, ∠B,
∠C উৎপন্ন হয়েছে তা অন্তঃস্থ কোণ ।

\textbf{বহিঃস্থকোণ:} △ABC চিত্রের বাহিরে যে ∠ACD কোণ উৎপন্ন হয়েছে তা
বহিঃস্থকোণ।

\textbf{উচ্চতা:} ত্রিভুজের শীর্ষবিন্দু থেকে ভূমির উপর অঙ্কিত লম্বের দৈর্ঘ্যকে উচ্চতা
বলে । চিত্রে AD হল ABC ত্রিভুজের উচ্চতা।

\textbf{মধ্যমা:} ত্রিভুজের শীর্ষ বিন্দু থেকে ভূমির মধ্যবিন্দুর সংযোগ রেখাকে মধ্যমা
বলে। চিত্রে AD হল মধ্যমা; কেননা AD, BC বাহুকে ২ ভাগে ভাগ করেছে। এছাড়াও B ও
C বিন্দুকে শীর্ষবিন্দু ধরে আরো দুটি মধ্যমা অঙ্কন করা যায় ।

\subsubsection{ত্রিভুজ বিষয়ক গুরুত্বপূর্ণ
অনুসিদ্ধান্ত}\label{ux9a4ux9b0ux9adux99c-ux9acux9b7ux9afux995-ux997ux9b0ux9a4ux9acux9aaux9b0ux9a3-ux985ux9a8ux9b8ux9a6ux9a7ux9a8ux9a4}

ত্রিভুজের বিভিন্ন বৈশিষ্ট্য ও শর্তের উপর নিম্নলিখিত অনুসদ্ধিান্ত গুলো অত্যন্ত গুরুত্বপূর্ণ
কেননা এই নিয়মগুলির উপর ভিত্তি করেই বিভিন্ন ধরনের প্রশ্ন তৈরী করা যায় এবং
সেভাবেই বিভিন্ন পরীক্ষায় ও আসে।

তাই চিত্র দেখে বুঝে বুঝে পড়ুন এবং লিখুন। পাশাপাশি নিয়ম সংশ্লিষ্ট প্রশ্নগুলো সমাধান
করুন ।

\textbf{অনুসিদ্ধান্ত - ০১:} ত্রিভুজের তিন কোণের সমষ্টি ১৮০ ডিগ্রি অর্থাৎ ২ সমকোণ
। চিত্রে: ∠A + ∠B + ∠C = ১৮০ ডিগ্রি বা ২ সমকোণ।

এজন্য যে কোন একটি কোণের পরিমাণ দেয়া না থাকলে অপর দুটি কোণের পরিমাণ

যোগ করে 180 ডিগ্রি থেকে তা বিয়োগ করলে তৃতীয় কোণের পরিমাণ বের হবে।

\textbf{অনুসিদ্ধান্ত - ০২:} ত্রিভুজের যেকোন দুই বাহুর সমষ্টি,তার তৃতীয় বাহু অপেক্ষা
বৃহত্তর, অর্থাৎ ত্রিভুজের যে কোন দুটি বাহু যোগ করলে তা তৃতীয় বাহু অপেক্ষা বড় হতে
হবে। যদি বৃহত্তর না হয় তাহলে তা ত্রিভুজ হবে না। যেমন:

উপরের চিত্রে AB = ৩ সে.মি., BC = ৪ সে. মি. এবং AC = ৫ সে. মি. তাহলে, ৩+ ৪
\textgreater{} ৫ বা, ৭ \textgreater{} ৫ অর্থাৎ AB + BC \textgreater{} AC
অন্য পাশ থেকে অন্য যে কোন দুটি বাহু নিলেও তাদের যোগফল ৩য় বাহু থেকে বড় হবে ।

\textbf{অনুসিদ্ধান্ত -- ০৩:} ত্রিভুজের যেকোন দুই বাহুর অন্তর, তৃতীয় বাহু অপেক্ষা
ক্ষুদ্রতর, অর্থাৎ ত্রিভুজের দুটি বাহু বিয়োগ করলে যেন তা তৃতীয় বাহু থেকে ছোট হয়।

উপরের চিত্রে AB = ৩ সে. মি., BC = ৪ সে. মি. এবং AC = ৫ সে. মি.

তাহলে, ৫- ৩\textless{} ৪ এবং ২ \textless{} ৪ সুতরাং বলা যায়, AC - AB
\textless{} BC

\textbf{বোঝার জন্য:} একটি ত্রিভুজের তিনটি বাহুর দৈর্ঘ্য যথাক্রমে ৫, ৮ এবং ১২
এখানে যে কোন দুটি বাহু বিয়োগ করলে তা তৃতীয় বাহু অপেক্ষা ছোট হবে।

কিন্তু যদি বলা হয় তিনটি বাহুর দৈর্ঘ্য যথাক্রমে ৩,৫,১২ তখন তা ত্রিভুজ হবে না,
কারণ এখানে ১২- ৫ = ৭ হয় যা তৃতীয় বাহু ৩ থেকে বড়।

\textbf{অনুসিদ্ধান্ত-০৪:} ত্রিভুজের কোনো শীর্ষ বিন্দু থেকে তার বিপরীত বাহুর
মধ্যবিন্দুর সংযোজক সরলরেখাকে মধ্যমা বলে।

\textbf{Confusion Clear:} শীর্ষ বিন্দু অর্থ শুধু উপরের বিন্দু নয়, এবং ত্রিভুজের
যেকোনো একটি বাহুকে ভূমি কল্পনা করলে ভূমির বিপরীত বিন্দুকে শীর্ষ বিন্দু বলে ।

\textbf{মধ্যমা সংক্রান্ত গুরুত্বপূর্ণ তথ্য:}

(i) ১টি ত্রিভুজের ৩টি মধ্যমা থাকে ।

(ii) যেকোন মধ্যমা ত্রিভুজের ক্ষেত্রফলকে সমান দুই ভাগে ভাগ করে।

(iii) ত্রিভুজের মধ্যমাগুলোর সমষ্টি ত্রিভুজের পরিসীমা অপেক্ষা ক্ষুদ্রতর।

\textbf{অনুসিদ্ধান্ত -০৫ :} কোন ত্রিভুজের বৃহত্তর বাহুর বিপরীত কোণ বৃহত্তম আবার
বৃহত্তম কোণের বিপরীত বাহু বৃহত্তর ।

\textbf{অনুসিদ্ধান্ত -০৬ :} কোন ত্রিভুজের সমান সমান বাহুর বিপরীত কোণগুলোও পরস্পর
সমান আবার সমান সমান কোণের বিপরীত বাহুগুলোও সমান হবে ।

\textbf{অনুসিদ্ধান্ত-০৭:} ত্রিভুজের একটি বাহু বর্ধিত করলে যে বহিঃস্থ কোণ উৎপন্ন
হয় তা বিপরীত অন্তঃস্থ কোণদ্বয়ের সমষ্টির সমান।

\textbf{অনুসিদ্ধান্ত-০৮ :} ত্রিভুজের বহিঃস্থ কোণগুলোর যোগফল ৪ সমকোণ বা ৩৬০
ডিগ্রি। সুতরাং ত্রিভুজের যে কোন দুটি বহিঃস্থ

কোণের সমষ্টি দুই সমকোণ অপেক্ষা বৃহত্তর হবে। মনে রাখুন: যে কোন বহুভুজের ক্ষেত্রে (
চতুর্ভুজ, পঞ্চভুজ,ষড়ভুজ-ইত্যাদির ক্ষেত্রেও) বহিঃস্থ কোণগুলোর সমষ্টি ৩৬০ ডিগ্রি হয়।

\textbf{অনুসিদ্ধান্ত- ০৯:} ত্রিভুজের যেকোন দুইবাহুর মধ্যবিন্দুর সংযোজক সরলরেখা
তৃতীয় বাহুর সমান্তরাল ও অর্ধেক।

\subsubsection{ত্রিভুজের
সর্বসমতা}\label{ux9a4ux9b0ux9adux99cux9b0-ux9b8ux9b0ux9acux9b8ux9aeux9a4}

দুটি ত্রিভুজ সর্বসম হওয়া অর্থ হল ত্রিভুজ দুটির আকার, আকৃতি ও ক্ষেত্রফল সমান হওয়া ।
তাই দুটি ভিন্ন ভিন্ন ক্ষেত্রের ক্ষেত্রফল সমান হলেও তাদের আকৃতি সমান না হলে ত্রিভুজ
দুটি যেমন সর্বসম হবে না।

তেমনি, দুটি ত্রিভুজ যদি দেখতে একই রকম মনে হয় কিন্তু তাদের ক্ষেত্রফল সমান নয়
তাহলেও ত্রিভুজ দুটি সর্বসম হবে না।

\textbf{সর্বসম: আকার, আকৃতি ও ক্ষেত্রফল সমান হতে হবে।}

\textbf{দুটি ত্রিভুজের সর্বসমতার শর্ত:} নিম্নোক্ত শর্তগুলো পূরণ হলে দুটি ত্রিভুজ
সর্বসম হয়।

ক. দুই বাহু ও অন্তর্ভুক্ত কোণ সমান

খ. তিনবাহু সমান

গ. দুই কোণ ও এক বাহু সমান

ঘ. সমকোণী ত্রিভুজের অতিভুজ ও এক বাহু সমান

সর্বসমতা নির্দেশক চিহ্ন ≅

\textbf{দুটি ত্রিভুজ কখন সর্বসম হয় না ?}

ক. যখন তাদের ক্ষেত্রফল সমান হয় কিন্তু আকৃতি সমান হয় না। চিত্র দেখে বুঝুন -

খ. আবার দুটি ত্রিভুজের আকৃতি দেখতে একই রকম হয় কিন্তু ক্ষেত্রফল সমান হয় না ৷ চিত্র
দেখে বুঝুন -

গ. আবার একটি ত্রিভুজের তিনটি কোণ অপর একটি ত্রিভুজের তিনটি কোণের সমান হওয়া
সত্ত্বেও ত্রিভুজ দুটি সর্বসম না ও হতে পারে। যেমন: চিত্রটি দেখুন: -

এই দুটি ত্রিভুজ সর্বসম নয় ।কারণ কোণগুলো সমান হলে ও বাহুগুলো কিন্তু সমান নয় । এরুপ
সমান কোণ বিশিষ্ট দুটি ত্রিভুজকে সদৃশ্যকোণী ত্রিভুজ বলা হয়।

\subsubsection{ত্রিভুজের
কেন্দ্র}\label{ux9a4ux9b0ux9adux99cux9b0-ux995ux9a8ux9a6ux9b0}

\textbf{অন্ত:কেন্দ্র (In-centre):}

ত্রিভূজের কোণের সমদ্বিখন্ডিত বাহুত্রয় যে বিন্দুতে মিলিত হয় তাকে অন্তঃকেন্দ্র বলে ।

চিত্রঃ- অন্ত:কেন্দ্র

△ABC- এ ∠A, ∠B ও ∠C কোণের সমদ্বিখন্ডক বাহুত্রয়ের মিলিত বিন্দু O। আর

O হল ত্রিভূজটির অন্তঃকেন্দ্র।

\textbf{পরিকেন্দ্র (Circumcentre) :}

ত্রিভূজের বাহু বা পরিসীমার উপর অঙ্কিত লম্বত্রয় ত্রিভূজের অভ্যন্তরে যে বিন্দুতে মিলিত
হয় তাকে পরিকেন্দ্র বলে ।

চিত্রঃ- পরিকেন্দ্র

△ABC- এ AB, BC ও AC বাহুত্রয়ের উপর লম্ব ত্রয় যথাক্রমে OP, OQ,OR ।

অতএব O হল পরিকেন্দ্র।

\textbf{ভরকেন্দ্র (Centroid):} ত্রিভূজের শীর্ষবিন্দু হতে বিপরীত বাহুর মধ্যবিন্দুর
সংযোগ সরলরেখাত্রয় যে বিন্দুতে মিলিত হয় তাকে ভরকেন্দ্র বলে।

চিত্রঃ- ভরকেন্দ্র

\textbf{ব্যাখ্যা}: চিত্রে △ABC-এ A,B ও C শীর্ষবিন্দু এবং তাদের বিপরীত বাহুত্রয়
হল BC,AC,AB। এখন শীর্ষ বিন্দু থেকে বাহুত্রয়ের মধ্যবিন্দুর সংযোগ সরলরেখাত্রয় হল-
AQ,BR,CP । রেখাত্রয় মিলিত হয়েছে O বিন্দুতে । অতএব O হল ত্রিভূজের ভরকেন্দ্র।

\subsubsection{ত্রিভুজের
প্রকারভেদ}\label{ux9a4ux9b0ux9adux99cux9b0-ux9aaux9b0ux995ux9b0ux9adux9a6}

ত্রিভুজকে প্রধানত ২টি বৈশিষ্ট্যের ভিত্তিতে ভাগ করা যায়। যথা:

১. কোণ ভেদে

২. বাহুভেদে

১. কোণভেদে ত্রিভুজ ৩ প্রকার । যথা:

(i) সমকোণী ত্রিভুজ

(ii) সূক্ষ্মকোণী ত্রিভুজ

(iii) স্থূলকোণী ত্রিভুজ

২. বাহুভেদে ত্রিভুজ ৩ প্রকার। যথা:

(i) সমবাহু ত্রিভুজ

(ii) সমদ্বিবাহু ত্রিভুজ

(iii) বিষমবাহু ত্রিভুজ

\paragraph{সমকোণী ত্রিভুজ (Right angled
triangle):}\label{ux9b8ux9aeux995ux9a3-ux9a4ux9b0ux9adux99c-right-angled-triangle}

\textbf{{[}সমকোণী ত্রিভুজের উপর সবথেকে বেশি প্রশ্ন হয়ে থাকে ।{]}}

যে ত্রিভুজের একটি কোণ সমকোণ তাকে সমকোণী ত্রিভুজ বলে। যেমন-

চিত্রে ∠B = ৯০ ডিগ্রি বা ১ সমকোণ

\paragraph{সমকোণী ত্রিভুজের
বৈশিষ্ট্যঃ}\label{ux9b8ux9aeux995ux9a3-ux9a4ux9b0ux9adux99cux9b0-ux9acux9b6ux9b7ux99fux9af}

সমকোণী ত্রিভুজের একটি কোণ অবশ্যই এক সমকোণ বা ৯০ ডিগ্রি হতে হবে।

সমকোণের বিপরীত বাহুকে অতিভুজ বলা হয়।

সমকোণী ত্রিভুজের অতিভুজ ই বৃহত্তম বাহু ।

সমকোণ ব্যতীত অন্য দুটি কোণ সূক্ষ্মকোণ এবং এরা পরস্পরের পূরক (কারণ যোগফল ৯০°)

সমকোণ সংলগ্ন বাহুদ্বয়ের যে কোন একটিকে লম্ব এবং অপরটিকে ভূমি ধরা যায়। (অর্থাৎ
লম্ব ও ভূমি নির্দিষ্ট নয়।)

সমকোণী ত্রিভুজের সমকোণ সংলগ্ন বাহুদ্বয়ের উপর অংকিত বর্গক্ষেত্রদ্বয়ের ক্ষেত্রফলের
সমষ্টি অপর বাহুর উপর অংকিত বর্গক্ষেত্রের ক্ষেত্রফলের সমান। এটা পীথাগোরাসের
উপপাদ্য নামে পরিচিত যাকে লেখা যায়,

\({AC}^{2} = {AB}^{2} + {BC}^{2}({\text{অতিভূজ}}^{2} = {\text{লম্ব}}^{2} + {\text{ভূমি}}^{2})\)

\paragraph{• বাহুর মাধ্যমে চেনার
উপায়:}\label{ux9acux9b9ux9b0-ux9aeux9a7ux9afux9ae-ux99aux9a8ux9b0-ux989ux9aaux9af}

ক্ষুদ্রতম দুই বাহুর বর্গের সমষ্টি উহার বৃহত্তর বাহুর বর্গের সমান হলে তা সমকোণী
ত্রিভুজ হয়।

\textbf{সমকোণী ত্রিভুজের ক্ষেত্রফল =} \(\frac{\text{১}}{\text{২}}\) \textbf{× ভূমি × উচ্চতা}

{[} এই সূত্রটি শুধু সমকোণী ত্রিভুজে কাজে লাগবে বিষয়টা এমন নয় বরং যেখানে ভূমি ও
উচ্চতার মান দেয়া থাকবে সেখানেই এই সূত্রটি প্রয়োগ করে সমাধান করা যাবে। {]}

\paragraph{পীথাগোরাসের
উপপাদ্য:}\label{ux9aaux9a5ux997ux9b0ux9b8ux9b0-ux989ux9aaux9aaux9a6ux9af}

সমকোণী ত্রিভুজের অতিভুজের উপর অংকিত বর্গক্ষেত্র , ঐ ত্রিভুজের অপর দুই বাহুর উপর
অংকিত বর্গক্ষেত্রের সমষ্টির সমান।

\textbf{ব্যাখ্যা}: এখানে ABC ত্রিভুজের প্রতিটি বাহুর উপর একটি করে বর্গক্ষেত্র
তৈরী করা হয়েছে । যার মধ্যে অতিভুজ AC বাহুর অংকিত বর্গক্ষেত্রটি AB এবং AC
বাহুদ্বয়ের উপর অংকিত বর্গক্ষেত্রদ্বয়ের সমষ্টির সমান । এই কথাটিই পীথাগোরাসের
সূত্রানুসারে এভাবে লিখতে হয়,

\({AC}^{2} = {AB}^{2} + {BC}^{2}\)

চিত্রানুসারে AC = 5 , AB = 3 এবং BC = 4

সূত্রে মানগুলো বসিয়ে পাই , \(5^{2} = 3^{2} + 4^{2}\)

বা, 25 = 9 + 16

বা, 25 = 25

এখন, অতিভুজ AC = 5 এর পরিবর্তে ৬ ধরলে মানগুলো পূর্ণ সংখ্যায় মিলবে না , তাই
কখনো লম্ব ও ভূমি ৩ ও ৪ মিটার হলে অতিভুজ অবশ্যই ৫ মিটার হবে ।

এই সূত্রটি দিয়ে পরবর্তীতে প্রচুর প্রশ্ন সমাধান করা যাবে । তাই গুরুত্ব দিয়ে বুঝুন ।

\paragraph{সূক্ষ্মকোণী ত্রিভুজ (Acute Angled
Triangle):}\label{ux9b8ux995ux9b7ux9aeux995ux9a3-ux9a4ux9b0ux9adux99c-acute-angled-triangle}

যে ত্রিভুজের তিনটি কোণ সূক্ষ্মকোণ (৯০ ডিগ্রি এর থেকে কম) তাকে সূক্ষ্মকোণী ত্রিভুজ
বলে। যেমন-

চিত্রে △ABC একটি সূক্ষ্মকোণী ত্রিভুজ।

\paragraph{সূক্ষকোণী ত্রিভূজের
বৈশিষ্ট্য:}\label{ux9b8ux995ux9b7ux995ux9a3-ux9a4ux9b0ux9adux99cux9b0-ux9acux9b6ux9b7ux99fux9af}

সূক্ষকোণী ত্রিভূজের তিনটি কোণই সূক্ষকোণ অর্থাৎ ৯০ ডিগ্রি এর থেকে কম ।

সমবাহু ত্রিভূজ একটি সূক্ষকোণী ত্রিভূজ যার প্রতিটি কোণ ৬০ ডিগ্রি ।

সূক্ষকোণী ত্রিভূজের যে কোন দুই বাহুর বর্গের সমষ্টি অপর বাহুর বর্গের চেয়ে বৃহত্তর ।

\paragraph{স্থূলকোণী ত্রিভুজ (Obtuse Angled Triangle)
:}\label{ux9b8ux9a5ux9b2ux995ux9a3-ux9a4ux9b0ux9adux99c-obtuse-angled-triangle}

যে ত্রিভুজের একটি কোণ স্থূলকোণ তাকে স্থূলকোণী ত্রিভুজ বলে। যেমন-

চিত্রে △ABC একটি স্থূলকোণী ত্রিভুজ কারণ ∠B স্থূলকোণ ।

\paragraph{স্থূলকোণী ত্রিভূজের
বৈশিষ্ট্য:}\label{ux9b8ux9a5ux9b2ux995ux9a3-ux9a4ux9b0ux9adux99cux9b0-ux9acux9b6ux9b7ux99fux9af}

স্থূলকোণী ত্রিভূজের একটি কোণ স্থূলকোণ এবং অন্য দুটি কোণ সূক্ষ্মকোণ ।

বৃহত্তম বাহুর বর্গ অপর দুই বাহুর বর্গের সমষ্টির চেয়ে বৃহত্তর ।

\paragraph{সমবাহু ত্রিভুজ (Equilateral
Triangle):}\label{ux9b8ux9aeux9acux9b9-ux9a4ux9b0ux9adux99c-equilateral-triangle}

যে ত্রিভুজের তিনটি বাহুই পরস্পর সমান তাঁকে সমবাহু ত্রিভুজ বলে। যেমন-

যেমন △ABC একটি সমবাহু ত্রিভুজ যার প্রতিটি বাহু পরস্পর সমান এবং প্রত্যেকটি কোণ ৬০
ডিগ্রি ।

\paragraph{সমবাহু ত্রিভুজের
বৈশিষ্ট্য:}\label{ux9b8ux9aeux9acux9b9-ux9a4ux9b0ux9adux99cux9b0-ux9acux9b6ux9b7ux99fux9af}

তিনটি বাহু সমান ।

কোণ তিনটি ও সমান।

প্রত্যেক কোণের পরিমান ৬০ ডিগ্রি । যেহেতু তিন বাহু সমান তাই কোণগুলোও সমান হবে।
আবার তিনটি কোণের সমষ্টি ১৮০ ডিগ্রি, তাই একটি কোণ হবে ১৮০÷৩ = ৬০ ডিগ্রি

শুধু পরিসীমা দেয়া থাকলে একমাত্র সমবাহু ত্রিভুজ আঁকা সম্ভব ।

সমবাহু ত্রিভুজকে সূক্ষকোণী ত্রিভুজ ও বলা যায় কারণ সমবাহু ত্রিভুজের প্রতিটি কোণই
সূক্ষকোণ ।

\paragraph{সমবাহু ত্রিভুজের
সুত্র:}\label{ux9b8ux9aeux9acux9b9-ux9a4ux9b0ux9adux99cux9b0-ux9b8ux9a4ux9b0}

△সমবাহু ত্রিভুজের পরিসীমা: এক বাহু a হলে, পরিসীমা = 3a ( ত্রিভুজের 3 বাহুর
যোগফল) ।

△সমবাহু ত্রিভুজের ক্ষেত্রফল = \(\frac{\sqrt{}3}{4}a^{2}\) (সমবাহু ত্রিভুজের এক
বাহুর দৈর্ঘ্য a হলে)

\paragraph{সমদ্বিবাহু ত্রিভুজ (Isosceles
Triangle):}\label{ux9b8ux9aeux9a6ux9acux9acux9b9-ux9a4ux9b0ux9adux99c-isosceles-triangle}

যে ত্রিভুজের দু'টি বাহু সমান তাকে সমদ্বিবাহু ত্রিভুজ বলে । যেমন-

চিত্রে △ABC একটি সমদ্বিবাহু ত্রিভুজ । যার AB বাহু = AC বাহু ।

\paragraph{সমদ্বিবাহু ত্রিভুজের
বৈশিষ্ট্য:}\label{ux9b8ux9aeux9a6ux9acux9acux9b9-ux9a4ux9b0ux9adux99cux9b0-ux9acux9b6ux9b7ux99fux9af}

দুটি বাহু পরস্পর সমান। চিত্রে AB = AC .

সমান সমান বাহুর বিপরীত কোণদ্বয় ও সমান ।

সমদ্বিবাহু ত্রিভুজের ক্ষেত্রফল = \(\frac{b}{4}\sqrt{4a^{2} - b^{2}}\) ( যেখানে
b = ভূমি এবং a = সমান সমান বাহুদ্বয়ের

দৈর্ঘ্য )

\paragraph{বিষমবাহু ত্রিভুজ (Scalene
Triangle):}\label{ux9acux9b7ux9aeux9acux9b9-ux9a4ux9b0ux9adux99c-scalene-triangle}

যে ত্রিভুজের তিনটি বাহুই অসমান তাকে বিষমবাহু বা অসমবাহু ত্রিভুজ বলে । যেমন-

চিত্র: বিষমবাহু ত্রিভুজ

\paragraph{বিষমবাহু ত্রিভুজের
বৈশিষ্ট্য:}\label{ux9acux9b7ux9aeux9acux9b9-ux9a4ux9b0ux9adux99cux9b0-ux9acux9b6ux9b7ux99fux9af}

প্রতিটি বাহু অসমান

মধ্যমাগুলো অসমান

বিষমবাহু ত্রিভুজের ক্ষেত্রফল =
\(\sqrt{s(s\  - \ a\ )\ (\ s\  - \ b\ )\ (\ s\  - \ c\ )}\)

এখানে 2S ( পরিসীমা ) = a + b + c

সুতরাং S ( অর্ধ পরিসীমা ) = \(\frac{a\  + \ b\  + \ c}{2}\)

\section{12. ত্রিভুজ MCQ প্রশ্ন উত্তর
অনুশীলন}\label{ux9a4ux9b0ux9adux99c-mcq-ux9aaux9b0ux9b6ux9a8-ux989ux9a4ux9a4ux9b0-ux985ux9a8ux9b6ux9b2ux9a8}

\textbf{ত্রিভুজ হওয়ার শর্ত কি?}

ক) যে কোনো ২ বাহুর দৈর্ঘ্যের যোগফল ৩য় বাহু অপেক্ষা ক্ষুদ্রতর

\textbf{\hl{খ) যে কোনো ২ বাহুর দৈর্ঘ্যের যোগফল ৩য় বাহুর দৈর্ঘ্য অপেক্ষা বৃহত্তর
✔}}

গ) ৩টি বাহুর দৈর্ঘ্য সমান

ঘ) ১টি কোণ সমকোণ

\textbf{\hl{✔ সঠিক উত্তর: খ) যে কোনো ২ বাহুর দৈর্ঘ্যের যোগফল ৩য় বাহুর দৈর্ঘ্য
অপেক্ষা বৃহত্তর}}

\begin{quote}
\textbf{ব্যাখ্যা:} \emph{ত্রিভুজের যে কোনো দুই বাহুর দৈর্ঘ্যের যোগফল তৃতীয় বাহু
অপেক্ষা বৃহত্তর।}

\emph{⇒ অর্থাৎ, ত্রিভুজের যেকোনো দুই বাহুর যোগফল সবসময় তৃতীয় বাহুর চেয়ে বড় হতে
হবে।}
\end{quote}

\textbf{কোন ত্রিভুজের একটি বাহু বর্ধিত করলে যে কোণ উৎপন্ন হয় তা ত্রিভুজটির -}

ক) লম্ব

খ) বাহু

\textbf{\hl{গ) বহিঃস্থ কোণ ✔}}

ঘ) উলম্ব

\textbf{\hl{✔ সঠিক উত্তর: গ) বহিঃস্থ কোণ}}

\begin{quote}
\textbf{ব্যাখ্যা:} \emph{কোনো ত্রিভুজের যেকোনো একটি বাহু বর্ধিত করলে যে কোণ
উৎপন্ন হয়, তাকে ত্রিভুজটির বহিঃস্থ কোণ বা বহিস্থ কোণ বলা হয়।}
\end{quote}

\textbf{ত্রিভুজের তিন কোণের সমষ্টি কত সমকোণ?}

\textbf{\hl{ক) ২ সমকোণ ✔}}

খ) ৩ সমকোণ

গ) ৪ সমকোণ

ঘ) ১ সমকোণ

\textbf{\hl{✔ সঠিক উত্তর: ক) ২ সমকোণ}}

\begin{quote}
\textbf{ব্যাখ্যা:} \emph{ত্রিভুজের তিন কোণের সমষ্টি ১৮০° বা দুই সমকোণ।}
\end{quote}

\textbf{কোন ত্রিভুজের দুইটি কোণ ১০° এবং ৮০° হলে, ত্রিভুজটি হবে}

ক) সূক্ষ্মকোণ

\textbf{\hl{খ) সমকোণী ✔}}

গ) স্থূলকোণী

ঘ) কোনোটি নয়

\textbf{\hl{✔ সঠিক উত্তর: খ) সমকোণী}}

\begin{quote}
\textbf{ব্যাখ্যা:} \emph{ত্রিভুজটির দুটি কোণের সমষ্টি দেওয়া আছে = ১০° + ৮০° =
৯০°}

\emph{∴ অপর কোণটির পরিমাণ হবে = ১৮০° − ৯০° = ৯০°}

\emph{⇒ যেহেতু ত্রিভুজটির একটি কোণ ৯০°, তাই ত্রিভুজটি সমকোণী ত্রিভুজ হবে।}
\end{quote}

\textbf{একটি ত্রিভুজের দুইটি কোণ ৩০° ও ৫৫° ত্রিভুজটি কোন ধরনের?}

ক) সমকোণী

খ) সমবাহু

গ) সমদ্বিবাহু

\textbf{\hl{ঘ) স্থূলকোণী ✔}}

\textbf{\hl{✔ সঠিক উত্তর: ঘ) স্থূলকোণী}}

\begin{quote}
\textbf{ব্যাখ্যা:} \emph{আমরা জানি, ত্রিভুজের তিন কোণের সমষ্টি ১৮০°।}

\emph{⇒ সুতরাং ত্রিভুজটির অপর কোণ = (১৮০ − (৩০ + ৫৫))° = (১৮০ − ৮৫)° = ৯৫°}

\emph{⇒ যেহেতু ত্রিভুজটির একটি কোণ এক সমকোণ (৯০°) এর বড় তাই ত্রিভুজটি স্থূলকোণী
ত্রিভুজ।}
\end{quote}

\textbf{কোনো ত্রিভুজের দুইটি কোণ ২০° ও ৭০° হলে ত্রিভুজটি কোন ধরনের?}

ক) সমবাহু

খ) সমদ্বিবাহু

\textbf{\hl{গ) সমকোণী ✔}}

ঘ) স্থূলকোণী

\textbf{\hl{✔ সঠিক উত্তর: গ) সমকোণী}}

\begin{quote}
\textbf{ব্যাখ্যা:} \emph{ত্রিভুজের অপর কোণ ১৮০° − (২০° + ৭০°) = ৯০°।}

\emph{∴ ত্রিভুজটি সমকোণী।}
\end{quote}

\textbf{চিত্রে x-এর মান কত?}

ক) ১৫°

\textbf{\hl{খ) ৩০° ✔}}

গ) ৪৫°

ঘ) ৬০°

\textbf{\hl{✔ সঠিক উত্তর: খ) ৩০°}}

\begin{quote}
\textbf{ব্যাখ্যা:} \emph{আমরা জানি, ত্রিভুজের তিন কোণের সমষ্টি ১৮০°।}

\emph{⇒ প্রশ্নমতে, x + ২x + ৩x = ১৮০°}

\emph{⇒ ৬x = ১৮০°}

\emph{∴ x = ৩০°}
\end{quote}

\textbf{কোন ৩টি বাহু দিয়ে ত্রিভুজ গঠন করা যাবে না?}

ক) ২, ৪, ৫

খ) ৪, ৫, ৬

\textbf{\hl{গ) ২, ৪, ৭ ✔}}

ঘ) সবগুলো দিয়েই ত্রিভুজ গঠন করা যাবে

\textbf{\hl{✔ সঠিক উত্তর: গ) ২, ৪, ৭}}

\begin{quote}
\textbf{ব্যাখ্যা:} \emph{ত্রিভুজের যেকোনো দুই বাহুর দৈর্ঘ্যের যোগফল তৃতীয় বাহু
অপেক্ষা বৃহত্তর হতে হবে।}

\emph{⇒ এখানে, ২ ও ৪ বাহুর যোগফল = ২ + ৪ = ৬, যা ৭ অপেক্ষা ক্ষুদ্রতর (৬
\textless{} ৭)।}

\emph{∴ ২, ৪, ৭ বাহু দিয়ে ত্রিভুজ গঠন করা যাবে না।}
\end{quote}

\textbf{ত্রিভুজের তিনটি মধ্যমা পরস্পর সমান হলে ত্রিভুজের যে কোনো কোণ-এর পরিমাণ
কত ডিগ্রী?}

ক) ৩০°

\textbf{\hl{খ) ৬০° ✔}}

গ) ৭০°

ঘ) ৮০°

\textbf{\hl{✔ সঠিক উত্তর: খ) ৬০°}}

\begin{quote}
\textbf{ব্যাখ্যা:} \emph{ত্রিভুজের তিনটি মধ্যমা পরস্পর সমান হলে ত্রিভুজটি একটি
সমবাহু ত্রিভুজ হয়।}

\emph{⇒ সমবাহু ত্রিভুজের প্রতিটি কোণের পরিমাণ ৬০°।}
\end{quote}

\textbf{একটি ত্রিভুজের দুটি কোণের সমষ্টি ১৬০°। তৃতীয় কোণটির মান কত?}

ক) ৩০°

খ) ৮০°

\textbf{\hl{গ) ২০° ✔}}

ঘ) ২৫°

\textbf{\hl{✔ সঠিক উত্তর: গ) ২০°}}

\begin{quote}
\textbf{ব্যাখ্যা:} \emph{আমরা জানি, ত্রিভুজের তিন কোণের সমষ্টি ১৮০°।}

\emph{∴ তৃতীয় কোণটির মান = ১৮০° − ১৬০° = ২০°}
\end{quote}

\textbf{ত্রিভুজের তিনটি কোণের সমষ্টি কত?}

ক) ৮০°

খ) ১২০°

গ) ১৬০°

\textbf{\hl{ঘ) ১৮০° ✔}}

\textbf{\hl{✔ সঠিক উত্তর: ঘ) ১৮০°}}

\begin{quote}
\textbf{ব্যাখ্যা:} \emph{ত্রিভুজের তিনটি কোণের সমষ্টি দুই সমকোণ বা ১৮০°।}
\end{quote}

\textbf{অষ্টভুজের অন্তঃস্থ কোণসমূহের সমষ্টি কত সমকোণ?}

ক) ১৬

\textbf{\hl{খ) ১২ ✔}}

গ) ৮

ঘ) ৪

\textbf{\hl{✔ সঠিক উত্তর: খ) ১২}}

\begin{quote}
\textbf{ব্যাখ্যা:} \emph{অন্তঃস্থ কোণ = ((n−2)×180°/n)}

\emph{∴ কোণসমূহের সমষ্টি = 6×135° = ১০৮০° = ১২ সমকোণ {[}১ সমকোণ = ৯০°{]}}
\end{quote}

\textbf{ত্রিভুজের AB = AC এবং ∠BAC = 80° হলে ∠ABC কত?}

ক) 20°

\textbf{\hl{খ) 50° ✔}}

গ) 80°

ঘ) 100°

\textbf{\hl{✔ সঠিক উত্তর: খ) 50°}}

\begin{quote}
\textbf{ব্যাখ্যা:} \emph{প্রশ্নানুযায়ী, AB = AC তাই, ∠ABC = ∠ACB হবে।}

\emph{কিন্তু ∠BAC = 80° বলে}

\emph{⇒ ∠ABC + ∠ACB = 100° হবে}

\emph{⇒ 2∠ABC = 100°}

\emph{⇒ ∠ABC = (100°/2) = 50°}

\emph{⇒ ∠ABC = 50°}
\end{quote}

\textbf{ত্রিভুজের ∠ABC = 45°, ∠BAC = 60° এবং BC বাহুকে D পর্যন্ত বর্ধিত করা
হলে ∠ACD = কত?}

ক) 45°

খ) 60°

\textbf{\hl{গ) 105° ✔}}

ঘ) 120°

\textbf{\hl{✔ সঠিক উত্তর: গ) 105°}}

\begin{quote}
\textbf{ব্যাখ্যা:} \emph{আমরা জানি, ত্রিভুজের যেকোনো একটি বাহুকে বর্ধিত করলে
যে বহিঃস্থ কোণ উৎপন্ন হয়, তা বিপরীত অন্তঃস্থ কোণদ্বয়ের সমষ্টির সমান।}

\emph{কাজেই ∠ACD = ∠ABC + ∠BAC = 45° + 60° = 105°}
\end{quote}

\textbf{ত্রিভুজের ∠A = 60°, ∠B = 80°, ∠C = কত?}

ক) ৮০°

খ) ৬০°

গ) ১২০°

\textbf{\hl{ঘ) ৪০° ✔}}

\textbf{\hl{✔ সঠিক উত্তর: ঘ) ৪০°}}

\begin{quote}
\textbf{ব্যাখ্যা:} \emph{ত্রিভুজের তিন কোণের সমষ্টি দুই সমকোণ বা ১৮০°।}

\emph{⇒ ∠C = 180° − (60° + 80°)}

\emph{⇒ ∠C = 180° − 140°}

\emph{∴ ∠C = 40°}
\end{quote}

\textbf{নিচের চিত্রে ∠B = 75° এবং ∠ACE = 150° হলে ∠A কোণের মান কত?}

ক) 30°

খ) 45°

\textbf{\hl{গ) 75° ✔}}

ঘ) 105°

\textbf{\hl{✔ সঠিক উত্তর: গ) 75°}}

\begin{quote}
\textbf{ব্যাখ্যা:} \emph{এখানে, ∠ACE = 150°}

\emph{∴ ∠ACB = 180° − 150° = 30°}

\emph{আবার, ∠ABC = 75°}

\emph{সুতরাং, ∠A = 180° − ∠ABC − ∠ACB}

\emph{= 180° − 75° − 30°}

\emph{= 75°}
\end{quote}

\textbf{একটি ত্রিভুজের দুটি কোণ যথাক্রমে 70° এবং 90° হলে তৃতীয় কোণটির পরিমাণ
রেডিয়ানে হবে-}

ক) (π/6)

\textbf{\hl{খ) (π/9) ✔}}

গ) (π/12)

ঘ) কোনটিই নয়

\textbf{\hl{✔ সঠিক উত্তর: খ) (π/9)}}

\begin{quote}
\textbf{ব্যাখ্যা:} \emph{জানা আছে, ত্রিভুজের তিন কোণের সমষ্টি = 180°}

\emph{অর্থাৎ 70° + 90° + তৃতীয় কোণ = 180°}

\emph{বা, তৃতীয় কোণ = 180° − 160°}

\emph{∴ তৃতীয় কোণ = 20°}

\emph{⇒ 20° = 20 × (π/180) রেডিয়ান = (π/9) রেডিয়ান}
\end{quote}

\textbf{একটি ত্রিভুজের দুটি কোণ যথাক্রমে ৪৫ ডিগ্রি এবং ৬৫ ডিগ্রি হলে, অপর
কোণটি কত ডিগ্রি?}

\textbf{\hl{ক) ৭০ ডিগ্রি ✔}}

খ) ৮০ ডিগ্রি

গ) ৬০ ডিগ্রি

ঘ) ৫০ ডিগ্রি

\textbf{\hl{✔ সঠিক উত্তর: ক) ৭০ ডিগ্রি}}

\begin{quote}
\textbf{ব্যাখ্যা:} \emph{আমরা জানি, ত্রিভুজের তিন কোণ = ১৮০ ডিগ্রি।}

\emph{∴ তৃতীয় কোণ = ১৮০−৬৫−৪৫ = ৭০ ডিগ্রি।}
\end{quote}

\textbf{কোনো ত্রিভুজের শীর্ষকোণের সমদ্বিখন্ডক যদি ভূমির উপর লম্ব হয়, তবে
ত্রিভুজটিকে কি বলে?}

\textbf{\hl{ক) সমদ্বিবাহু ✔}}

খ) সমবাহু

গ) সমকোণী

ঘ) স্থূলকোণী

\textbf{\hl{✔ সঠিক উত্তর: ক) সমদ্বিবাহু}}

\begin{quote}
\textbf{ব্যাখ্যা:} \emph{কোনো ত্রিভুজের শীর্ষকোণের সমদ্বিখন্ডক যদি ভূমির উপর লম্ব
হয়, তবে ত্রিভুজটি সমদ্বিবাহু ত্রিভুজ হয়।}
\end{quote}

\textbf{ত্রিভুজের যে কোনো বাহুকে উভয়দিকে বর্ধিত করলে উৎপন্ন বহিঃস্থ কোণদ্বয়ের
সমষ্টি-}

ক) দুই সমকোণ অপেক্ষা ক্ষুদ্রতর

\textbf{\hl{খ) দুই সমকোণ অপেক্ষা বৃহত্তর ✔}}

গ) দুই সমকোণের সমান

ঘ) দুই সমকোণের অর্ধেক

\textbf{\hl{✔ সঠিক উত্তর: খ) দুই সমকোণ অপেক্ষা বৃহত্তর}}

\begin{quote}
\textbf{ব্যাখ্যা:} \emph{ত্রিভুজের যেকোনো বাহুকে উভয়দিকে বর্ধিত করলে উৎপন্ন
বহিঃস্থ কোণদ্বয়ের সমষ্টি দুই সমকোণ অপেক্ষা বৃহত্তর হয়।}
\end{quote}

\textbf{কোনো ত্রিভুজের একটি বাহু উভয়দিকে বর্ধিত করায় উৎপন্ন বহিঃস্থ কোণগুলো
পরস্পর সমান হলে ত্রিভুজটি-}

ক) বিষমবাহু

\textbf{\hl{খ) সমদ্বিবাহু ✔}}

গ) সমবাহু

ঘ) সমকোণী

\textbf{\hl{✔ সঠিক উত্তর: খ) সমদ্বিবাহু}}

\begin{quote}
\textbf{ব্যাখ্যা:} \emph{ত্রিভুজের অন্তঃস্থ কোণ ও বহিঃস্থ কোণের সমষ্টি দুই
সমকোণ।সুতরাং ত্রিভুজের দুটি বহিঃস্থ কোণ সমান হলে ত্রিভুজটির অন্তঃস্থ কোণ দুটি ও
সমান হবে। অর্থাৎ ত্রিভুজটি সমদ্বিবাহু হবে।}
\end{quote}

\textbf{একটি ত্রিভুজের তিন কোণের অনুপাত ১ : ২ : ৩ হলে, ত্রিভুজটি হবে}

ক) সমবাহু

খ) সূক্ষ্মকোণী

গ) স্থূলকোণী

\textbf{\hl{ঘ) সমকোণী ✔}}

\textbf{\hl{✔ সঠিক উত্তর: ঘ) সমকোণী}}

\begin{quote}
\textbf{ব্যাখ্যা:} \emph{ধরি, কোণ তিনটি যথাক্রমে x, 2x এবং 3x}

\emph{প্রশ্নমতে, x + 2x + 3x = 180°}

\emph{⇒ 6x = 180°}

\emph{⇒ x = (180°/6) = 30°}

\emph{অতএব, কোণগুলো হবে = 30°, (2 × 30°) = 60° এবং (3 × 30°) = 90°}

\emph{যেহেতু, একটি কোণের পরিমাণ 90°}

\emph{∴ তাই কোণটি হবে সমকোণী ত্রিভুজ।}
\end{quote}

\textbf{একটি ত্রিভুজের তিনটি কোণের অনুপাত ৩ : ৪ : ৫ হলে, কোণ তিনটিকে
ডিগ্রীতে প্রকাশ করুন।}

\textbf{\hl{ক) ৪৫°, ৬০°, ৭৫° ✔}}

খ) ৩০°, ৪০°, ৫০°

গ) ৪২°, ৫৬°, ৭০°

ঘ) ৪৮°, ৬৪°, ৮০°

\textbf{\hl{✔ সঠিক উত্তর: ক) ৪৫°, ৬০°, ৭৫°}}

\begin{quote}
\textbf{ব্যাখ্যা:} \emph{অনুপাতগুলোর যোগফল = ৩ + ৪ + ৫ = ১২}

\emph{∴ ১ম কোণ = 180° এর (3/12) = 45°}

\emph{∴ ২য় কোণ = 180° এর (4/12) = 60°}

\emph{∴ ৩য় কোণ = 180° এর (5/12) = 75°}
\end{quote}

\textbf{নিচের ত্রিভুজের বাহুর দৈর্ঘ্য দেয়া আছে। কোন ক্ষেত্রে ত্রিভুজ আঁকা সম্ভব
নয়?}

\textbf{\hl{ক) ২, ৩, ৫ সে. মি. ✔}}

খ) ৪, ৫, ৬ সে. মি.

গ) ৩, ৫, ৭ সে. মি.

ঘ) ৫, ৬, ৮ সে. মি.

\textbf{\hl{✔ সঠিক উত্তর: ক) ২, ৩, ৫ সে. মি.}}

\begin{quote}
\textbf{ব্যাখ্যা:} \emph{আমরা জানি, ত্রিভুজের যে কোনো দুই বাহুর সমষ্টি তার তৃতীয়
বাহু অপেক্ষা বৃহত্তর হয়।}

\emph{⇒ অপশন (খ), (গ) ও (ঘ)-এর ক্ষেত্রে যে কোনো দুটি সরলরেখার দৈর্ঘ্যের যোগফল
তৃতীয় সরলরেখার দৈর্ঘ্য অপেক্ষা বৃহত্তর।}

\emph{⇒ কিন্তু অপশন (ক)-এর ক্ষেত্রে ১ম ও ২য় সরলরেখার দৈর্ঘ্যের যোগফল (২ + ৩ বা
৫ সে. মি.; যা ৩য় সরলরেখার দৈর্ঘ্যের ৫ সে. মি. এর সমান)।}

\emph{∴ সুতরাং এই সরলরেখা দ্বারা ত্রিভুজ অঙ্কন করা অসম্ভব।}
\end{quote}

\textbf{ত্রিভুজের তিন বাহুর দৈর্ঘ্য মিটারে দেওয়া হলো। কোন ক্ষেত্রে ত্রিভুজ অঙ্কন
সম্ভব নয়?}

ক) ৬, ৭ ও ১১

খ) ৩, ৮ ও ৮

\textbf{\hl{গ) ১৪, ১২ ও ২৮ ✔}}

ঘ) ২০, ৮ ও ১৪

\textbf{\hl{✔ সঠিক উত্তর: গ) ১৪, ১২ ও ২৮}}

\begin{quote}
\textbf{ব্যাখ্যা:} \emph{ত্রিভুজের যে কোন দুই বাহুর যোগফল তৃতীয় বাহুর থেকে বড়
হবে।}

\emph{⇒ অপশন গ) ১৪ + ১২ = ২৬ যা ২৮ এর থেকে ছোট।}

\emph{∴ তাই ১৪, ১২ ও ২৮ মিটারের বাহু দ্বারা ত্রিভুজ আঁকা সম্ভব নয়।}
\end{quote}

\textbf{ত্রিভুজের যে কোনো দুই বাহুর মধ্যবিন্দুর সংযোজক সরলরেখা তৃতীয় বাহুর -\/-}

ক) সমান (Equal)

খ) বড় হবে (greater than)

গ) দ্বিগুণ (Twice)

\textbf{\hl{ঘ) অর্ধেক (Half) ✔}}

\textbf{\hl{✔ সঠিক উত্তর: ঘ) অর্ধেক (Half)}}

\begin{quote}
\textbf{ব্যাখ্যা:} \emph{ত্রিভুজের যেকোনো দুই বাহুর মধ্যবিন্দুর সংযোজক সরলরেখা
তৃতীয় বাহুর অর্ধেক হয়।}
\end{quote}

\textbf{একটি ত্রিভুজের তিনটি কোণের অনুপাত ৩ : ৪ : ৫ হলে কোণ তিনটির মান কত?}

\textbf{\hl{ক) ৪৫°, ৬০°, ৭৫° ✔}}

খ) ৩০°, ৬০°, ৯০°

গ) ৬০°, ৯০°, ৩০°

ঘ) ৯০°, ৪৫°, ৪৫°

\textbf{\hl{✔ সঠিক উত্তর: ক) ৪৫°, ৬০°, ৭৫°}}

\begin{quote}
\textbf{ব্যাখ্যা:} \emph{আমরা জানি, ত্রিভুজের তিনকোণের সমষ্টি = ১৮০°}

\emph{⇒ এখন, অনুপাতের রাশিগুলোর যোগফল = ৩ + ৪ + ৫ = ১২}

\emph{∴ কোণগুলো ১৮০ এর (3/12) = ৪৫°; ১৮০ এর (4/12) = ৬০ এবং ১৮০ এর (5/12)
= ৭৫।}
\end{quote}

\textbf{একটি ত্রিভুজের তিনটি কোণের অনুপাত ২ : ৩ : ৪ হলে, বৃহত্তম কোণের পরিমাণ
কত?}

ক) ৬০°

খ) ৭০°

\textbf{\hl{গ) ৮০° ✔}}

ঘ) ৯০°

\textbf{\hl{✔ সঠিক উত্তর: গ) ৮০°}}

\begin{quote}
\textbf{ব্যাখ্যা:} \emph{ধরি, তিনটি কোণ যথাক্রমে 2x, 3x ও 4x}

\emph{∴ 2x + 3x + 4x = 180°}

\emph{⇒ 9x = 180° ⇒ x = 20°}

\emph{∴ বৃহত্তম কোণের পরিমাণ = 4 × 20° = 80°}
\end{quote}

\textbf{একটি ত্রিভুজের দুটি কোণের অনুপাত ৩ : ৫। তৃতীয় কোণটি ৫২° হলে ছোট
কোণটির পরিমাণ কত?}

\textbf{\hl{ক) ৪৮ ✔}}

খ) ৪৫

গ) ৬৭

ঘ) ৮০

\textbf{\hl{✔ সঠিক উত্তর: ক) ৪৮}}

\begin{quote}
\textbf{ব্যাখ্যা:} \emph{ত্রিভুজের তিন কোণের সমষ্টি = ১৮০°}

\emph{⇒ দুটি কোণের সমষ্টি = ১৮০° − ৫২° = ১২৮°}

\emph{⇒ অনুপাত দুটির যোগফল = ৩ + ৫ = ৮}

\emph{∴ ছোট কোণটির পরিমাণ = ১২৮° এর (3/8) = ৪৮°}
\end{quote}

\textbf{১ সমকোণ = কত ডিগ্রী?}

ক) ১৮০ ডিগ্রি

খ) ১০০ ডিগ্রি

\textbf{\hl{গ) ৯০ ডিগ্রি ✔}}

ঘ) কোনটিই নয়

\textbf{\hl{✔ সঠিক উত্তর: গ) ৯০ ডিগ্রি}}

\begin{quote}
\textbf{ব্যাখ্যা:} \emph{আমরা জানি, ১ সমকোণ সমান ৯০ ডিগ্রি।}
\end{quote}

\textbf{একটি ত্রিভুজের একটি কোণের মান ৯০°; অন্য কোণ দুটির অনুপাত ২ : ৩ হলে
ছোট কোণটির পূরক কোণ-}

ক) 36°

খ) 144°

\textbf{\hl{গ) 54° ✔}}

ঘ) 126°

\textbf{\hl{✔ সঠিক উত্তর: গ) 54°}}

\begin{quote}
\textbf{ব্যাখ্যা:} \emph{ধরি, অন্য কোণ দুটি 2x ও 3x}

\emph{⇒ 2x + 3x = 90° ∴ x = 18°}

\emph{এখানে 2x ও 3x একে অপরের পূরক কোণ}

\emph{∴ ছোট কোণের পূরক কোণ = (3 × 18°) = 54°}
\end{quote}

\textbf{∠A ও ∠B পরস্পর পূরক এবং কোণ দুটির অনুপাত ৩:২ হলে ∠A এর মান কত?}

ক) 36°

খ) 18°

\textbf{\hl{গ) 54° ✔}}

ঘ) 45°

\textbf{\hl{✔ সঠিক উত্তর: গ) 54°}}

\begin{quote}
\textbf{ব্যাখ্যা:} \emph{আমরা জানি, দুটি কোণের সমষ্টি ৯০° হলে কোণ দুটিকে
পরস্পরের পূরক কোণ বলা হয়। যেহেতু ∠A ও ∠B পরস্পর পূরক কোণ এবং তাদের অনুপাত =
৩:২}

\emph{ধরি, ∠A = 3x, ∠B = 2x}

\emph{∴ 3x + 2x = 90°}

\emph{⇒ 5x = 90°}

\emph{⇒ x = (90°/5) ∴ x = 18°}

\emph{∴ ∠A = 3x = 3 × 18° = 54°}
\end{quote}

\textbf{নিচের কোন তিনটি রেখাংশ দ্বারা ত্রিভুজ অঙ্কন করা অসম্ভব?}

\textbf{\hl{ক) 5 সে.মি., 8 সে.মি., 3 সে.মি. ✔}}

খ) 5 সে.মি., 4 সে.মি., 3 সে.মি.

গ) 2 সে.মি., 3 সে.মি., 4 সে.মি.

ঘ) 6.5 সে.মি., 6.5 সে.মি., 6.5 সে.মি.

\textbf{\hl{✔ সঠিক উত্তর: ক) 5 সে.মি., 8 সে.মি., 3 সে.মি.}}

\begin{quote}
\textbf{ব্যাখ্যা:} \emph{আমরা জানি, ত্রিভুজের যে কোনো দুই বাহুর সমষ্টি তার তৃতীয়
বাহু অপেক্ষা বৃহত্তর হয়।}

\emph{অপশন (ক)- এর ক্ষেত্রে ১ম ও ৩য় সরলরেখায় দৈর্ঘ্যের যোগফল (5 + 3) সে.মি.
বা 8 সে.মি. যা ২য় সরলরেখার দৈর্ঘ্যের সমান। সুতরাং এই সরলরেখা দ্বারা ত্রিভুজ
অঙ্কন অসম্ভব।}
\end{quote}

\textbf{ত্রিভুজের তিনটি কোণ 4x, 6x, 8x হলে বৃহত্তম কোণ কত?}

ক) 36°

খ) 72°

\textbf{\hl{গ) 80° ✔}}

ঘ) 96°

\textbf{\hl{✔ সঠিক উত্তর: গ) 80°}}

\begin{quote}
\textbf{ব্যাখ্যা:} \emph{ত্রিভুজের তিন কোণের সমষ্টি = 180°}

\emph{⇒ 4x + 6x + 8x = 180°}

\emph{⇒ 18x = 180° ∴ x = 10°}

\emph{∴ বৃহত্তম কোণ, 8x = 8 × 10° = 80°}
\end{quote}

\textbf{AC+CB = AB হলে C নিচের কোনটি?}

ক) সমবিন্দু

\textbf{\hl{খ) অন্তঃস্থ বিন্দু ✔}}

গ) মধ্যবিন্দু

ঘ) কোণ

\textbf{\hl{✔ সঠিক উত্তর: খ) অন্তঃস্থ বিন্দু}}

\begin{quote}
\textbf{ব্যাখ্যা:} \emph{AC + CB = AB হলে C বিন্দুটি A ও B বিন্দুর অন্তঃস্থ বিন্দু
হবে।}
\end{quote}

\textbf{কোনো ত্রিভুজের যে কোনো এক বাহু বর্ধিত করলে যে বহিঃস্থ কোণ উৎপন্ন হয় তা
হবে-}

ক) অন্তঃস্থ কোণগুলোর সমান

খ) অন্তঃস্থ কোণগুলোর অর্ধেক

\textbf{\hl{গ) অন্তঃস্থ বিপরীত কোণদ্বয়ের সমষ্টির সমান ✔}}

ঘ) অন্তঃস্থ বিপরীত কোণদ্বয়ের সমষ্টির অর্ধেক

\textbf{\hl{✔ সঠিক উত্তর: গ) অন্তঃস্থ বিপরীত কোণদ্বয়ের সমষ্টির সমান}}

\begin{quote}
\textbf{ব্যাখ্যা:} \emph{কোনো ত্রিভুজের বহিঃস্থ কোণ তার বিপরীত অন্তঃস্থ
কোণদ্বয়ের সমষ্টির সমান।}
\end{quote}

\textbf{ত্রিভুজের মধ্যমাত্রয়ের ছেদবিন্দুর নাম কি?}

ক) পরিকেন্দ্র

\textbf{\hl{খ) ভরকেন্দ্র ✔}}

গ) অন্তঃকেন্দ্র

ঘ) লম্বকেন্দ্র

\textbf{\hl{✔ সঠিক উত্তর: খ) ভরকেন্দ্র}}

\begin{quote}
\textbf{ব্যাখ্যা:} \emph{ত্রিভুজের মধ্যমাত্রয়ের ছেদবিন্দুকে ভরকেন্দ্র বলে।}
\end{quote}

\textbf{কোন ত্রিভুজের তিন বাহু বা লম্ব দ্বিখন্ডকগুলোর ছেদবিন্দুকে কি বলে?}

ক) লম্বকেন্দ্র

খ) ভরকেন্দ্র

\textbf{\hl{গ) পরিকেন্দ্র ✔}}

ঘ) অন্তঃকেন্দ্র

\textbf{\hl{✔ সঠিক উত্তর: গ) পরিকেন্দ্র}}

\begin{quote}
\textbf{ব্যাখ্যা:} \emph{ত্রিভুজের বাহুত্রয়ের লম্বদ্বিখন্ডকত্রয় সমবিন্দু। এই বিন্দু
ত্রিভুজের পরিকেন্দ্র।}
\end{quote}

\textbf{কোনো ত্রিভুজের তিন কোণের দ্বিখন্ডকগুলো যে বিন্দুতে ছেদ করে তাকে বলে-}

ক) পরিকেন্দ্র

\textbf{\hl{খ) অন্তঃকেন্দ্র ✔}}

গ) ভরকেন্দ্র

ঘ) লম্বকেন্দ্র

\textbf{\hl{✔ সঠিক উত্তর: খ) অন্তঃকেন্দ্র}}

\begin{quote}
\textbf{ব্যাখ্যা:} \emph{লম্ববিন্দু: ত্রিভুজের তিনটি শীর্ষ থেকে বিপরীত বাহুর উপর
অঙ্কিত লম্ব যে বিন্দুতে মিলিত হয় তাকে লম্ববিন্দু বলে।}

\emph{লম্বদুরত্ব: লম্ববিন্দু থেকে ত্রিভুজের শীর্ষ বিন্দুর দুরত্বকে লম্বদুরত্ব বলে।}

\emph{ভরকেন্দ্র: ত্রিভুজের মধ্যমাত্রয় যে বিন্দুতে মিলিত হয় তাকে ভরকেন্দ্র বলে।}

\emph{পরিকেন্দ্র: ত্রিভুজের যেকোনো দুই বাহুর সমদ্বিখন্ডকদ্বয় যে বিন্দুতে মিলিত হয়
তাকে পরিকেন্দ্র বলে।}

\emph{পরিব্যাসার্ধ: পরিকেন্দ্র থেকে ত্রিভুজের শীর্ষ বিন্দুর দুরত্বকে পরিব্যাসার্ধ
বলে।}

\emph{অন্তঃকেন্দ্র : ত্রিভুজের তিনটি কোণের সমদ্বিখন্ডকত্রয় যে বিন্দুতে মিলিত হয়
তাকে অন্তঃকেন্দ্র বলে।}

\emph{বহিঃকেন্দ্র: ত্রিভুজের একটি কোণের অন্ত-সমদ্বিখন্ডক এবং অপর দুই কোণের
বহি-সমদ্বিখন্ডক যে বিন্দুতে মিলিত হয় তাকে বহিঃকেন্দ্র বলে।}

\emph{নব-বিন্দুর কেন্দ্র : ত্রিভুজের নয়টি মূল বিন্দুগামী বৃত্তের কেন্দ্র হল নয়-বিন্দুর
কেন্দ্র বা নববিন্দু কেন্দ্র।}
\end{quote}

\textbf{কোনো ত্রিভুজের তিন বাহুর সমদ্বিখণ্ডকগুলোর ছেদবিন্দুর নাম কি?}

ক) বহিকেন্দ্র (Excentre)

খ) অন্তঃকেন্দ্র (In-centre)

গ) ভরকেন্দ্র (Centroid)

\textbf{\hl{ঘ) পরিকেন্দ্র (Circumcentre) ✔}}

\textbf{\hl{✔ সঠিক উত্তর: ঘ) পরিকেন্দ্র (Circumcentre)}}

\begin{quote}
\textbf{ব্যাখ্যা:} \emph{ত্রিভুজের বাহুগুলোর লম্বদ্বিখণ্ডক বা সমদ্বিখণ্ডকত্রয়ের
ছেদবিন্দুকে পরিকেন্দ্র (Circumcentre) বলে।}
\end{quote}

\textbf{কোনো ত্রিভুজের শীর্ষবিন্দুগুলো দিয়ে অঙ্কিত বৃত্তটির কেন্দ্রকে বলে -\/-}

ক) অন্তঃকেন্দ্র

\textbf{\hl{খ) পরিকেন্দ্র ✔}}

গ) লম্বকেন্দ্র

ঘ) ভরকেন্দ্র

\textbf{\hl{✔ সঠিক উত্তর: খ) পরিকেন্দ্র}}

\begin{quote}
\textbf{ব্যাখ্যা:} \emph{পরিবৃত্ত হচ্ছে ত্রিভুজের তিন শীর্ষবিন্দু দিয়ে অতিক্রান্ত
বৃত্ত। আর পরিবৃত্তের কেন্দ্রই হচ্ছে পরিকেন্দ্র।}
\end{quote}

\textbf{একটি সমবাহু ত্রিভুজের তিনটি কোণের মান:}

ক) ৬০°, ৫০°, ৭০°

\textbf{\hl{খ) ৬০°, ৬০°, ৬০° ✔}}

গ) ৮০°, ৯০°, ৯০°

ঘ) কোনোটিই নয়

\textbf{\hl{✔ সঠিক উত্তর: খ) ৬০°, ৬০°, ৬০°}}

\begin{quote}
\textbf{ব্যাখ্যা:} \emph{সমবাহু ত্রিভুজের তিনটি কোণ সমান অর্থাৎ তিনটি কোণের
মানই ৬০° করে।}
\end{quote}

\textbf{যে ত্রিভুজের তিন কোণ সমান তাকে কোন ধরনের ত্রিভুজ বলে?}

\textbf{\hl{ক) সমবাহু ✔}}

খ) বিষমবাহু

গ) সমকোণী

ঘ) স্থূলকোণী

\textbf{\hl{✔ সঠিক উত্তর: ক) সমবাহু}}

\begin{quote}
\textbf{ব্যাখ্যা:} \emph{যে ত্রিভুজের তিন কোণ সমান তার তিনটি বাহুই পরস্পর
সমান। তাই একে সমবাহু ত্রিভুজ বলে।}
\end{quote}

\textbf{সমবাহু ত্রিভুজের-}

ক) একটি বাহু অসমান

খ) দুইটি বাহু সমান

\textbf{\hl{গ) প্রতিটি বাহু সমান ✔}}

ঘ) প্রতিটি বাহু অসমান

\textbf{\hl{✔ সঠিক উত্তর: গ) প্রতিটি বাহু সমান}}

\begin{quote}
\textbf{ব্যাখ্যা:} \emph{সমবাহু ত্রিভুজের প্রতিটি বাহু পরস্পর সমান।}
\end{quote}

\textbf{সমবাহু ত্রিভুজের একটি বাহুকে বর্ধিত করলে বহিঃস্থ কোণের পরিমাণ কত?}

ক) ৩০°

খ) ৬০°

গ) ৯০°

\textbf{\hl{ঘ) ১২০° ✔}}

\textbf{\hl{✔ সঠিক উত্তর: ঘ) ১২০°}}

\begin{quote}
\textbf{ব্যাখ্যা:} \emph{সমবাহু ত্রিভুজের প্রতিটি কোণ ৬০°। একটি বাহু বর্ধিত করলে
উৎপন্ন বহিঃস্থ কোণ = 180° − 60° = 120°।}
\end{quote}

\textbf{সমবাহু ত্রিভুজের প্রত্যেকটি কোণের মান কত?}

ক) ৪০°

খ) ৫০°

\textbf{\hl{গ) ৬০° ✔}}

ঘ) ৭০°

\textbf{\hl{✔ সঠিক উত্তর: গ) ৬০°}}

\begin{quote}
\textbf{ব্যাখ্যা:} \emph{আমরা জানি, ত্রিভুজের তিন কোণের সমষ্টি ১৮০°। সমবাহু
ত্রিভুজের তিনটি কোণ সমান হওয়ায় প্রতিটি কোণের মান ১৮০° ÷ ৩ = ৬০°।}
\end{quote}

\textbf{যদি সমবাহু ত্রিভুজের বাহুর দৈর্ঘ্য a হয়, তবে ক্ষেত্রফল কত হবে?}

\textbf{\hl{ক) (√(3)/4)a\^{}2 ✔}}

খ) (√(3)/2)a\^{}2

গ) (3/2)a\^{}2

ঘ) (√(2)/2)a\^{}2

\textbf{\hl{✔ সঠিক উত্তর: ক) (√(3)/4)a\^{}2}}

\begin{quote}
\textbf{ব্যাখ্যা:} \emph{সমবাহু ত্রিভুজের বাহুর দৈর্ঘ্য a হলে ক্ষেত্রফল নির্ণয়ের
সূত্র হলো (√(3)/4)a\^{}2 ।}
\end{quote}

\textbf{একটি সমকোণী ত্রিভুজের সমকোণ ছাড়া অন্য দুটি কি কোণ?}

ক) সন্নিহিত কোণ

খ) সরল কোণ

\textbf{\hl{গ) সূক্ষ্মকোণ ✔}}

ঘ) পূরক কোণ

\textbf{\hl{✔ সঠিক উত্তর: গ) সূক্ষ্মকোণ}}

\begin{quote}
\textbf{ব্যাখ্যা:} \emph{যেহেতু ত্রিভুজের তিন কোণের সমষ্টি 180° এবং এক সমকোণ
হলো 90° , তাই বাকি থাকে 90° । এই 90° অপর দুটি কোণের সমষ্টি হওয়ায় কোণ দুটির
প্রত্যেকটি 90° এর চেয়ে ছোট অর্থাৎ সূক্ষ্মকোণ হবে।}
\end{quote}

\textbf{ত্রিভুজের ক্ষেত্রফল}

\textbf{\hl{ক) (1/2)(ভূমি×উচ্চতা) ✔}}

খ) (1/2)(ভূমি×উচ্চতা)

গ) (1/2)(ভূমি+উচ্চতা)

ঘ) (1/2)(ভূমি−উচ্চতা)

\textbf{\hl{✔ সঠিক উত্তর: ক) (1/2)(ভূমি×উচ্চতা)}}

\begin{quote}
\textbf{ব্যাখ্যা:} \emph{যেকোনো সাধারণ ত্রিভুজের ক্ষেত্রফল নির্ণয়ের সূত্র হলো
(1/2)×ভূমি×উচ্চতা ।}
\end{quote}

\textbf{একটি সমদ্বিবাহু সমকোণী ত্রিভুজের অতিভুজের দৈর্ঘ্য ১২ সেমি হলে, ত্রিভুজটির
ক্ষেত্রফল কত বর্গ সেমি?}

\textbf{\hl{ক) ৩৬ ✔}}

খ) ৪৮

গ) ৫৬

ঘ) ৭২

\textbf{\hl{✔ সঠিক উত্তর: ক) ৩৬}}

\begin{quote}
\textbf{ব্যাখ্যা:} \emph{ধরি, সমান সমান বাহুর দৈর্ঘ্য a}

\emph{পীথাগোরাসের উপপাদ্য অনুযায়ী, a\^{}2+a\^{}2=12\^{}2}

\emph{⇒ 2a\^{}2=144}

\emph{⇒ a\^{}2=72}

\emph{ত্রিভুজটির ক্ষেত্রফল = (1/2)×a×a=(1/2)×a\^{}2=(1/2)×72=36 বর্গ সেমি।}
\end{quote}

\textbf{ABC ত্রিভুজের B কোণটি C কোণের সমান, D হচ্ছে BC বাহুর মধ্য একটি বিন্দু।
নিচের কোন বক্তব্যটি সঠিক?}

ক) AB\textgreater BC

খ) AB\textless BC

গ) BD\textgreater CD

\textbf{\hl{ঘ) AC\textgreater AD ✔}}

\textbf{\hl{✔ সঠিক উত্তর: ঘ) AC\textgreater AD}}

\begin{quote}
\textbf{ব্যাখ্যা:} \emph{ত্রিভুজের ধর্ম অনুযায়ী, সমান কোণের বিপরীত বাহুগুলো
পরস্পর সমান হয় এবং মধ্যমার সাপেক্ষে বাহুর দূরত্বের তুলনামূলক বিশ্লেষণে AC বাহুর
দৈর্ঘ্য AD অপেক্ষা বড় হয়।}
\end{quote}

\textbf{একটি সমবাহু ত্রিভুজের একটি বাহু ১৬ মিটার। ত্রিভুজটির ক্ষেত্রফল কত?}

\textbf{\hl{ক) 64√(3) ✔}}

খ) 192

গ) 64

ঘ) 32√(3)

\textbf{\hl{✔ সঠিক উত্তর: ক) 64√(3)}}

\begin{quote}
\textbf{ব্যাখ্যা:} \emph{ধরি, সমবাহু ত্রিভুজের বাহুর দৈর্ঘ্য a=16 মিটার}

\emph{∴ সমবাহু ত্রিভুজের ক্ষেত্রফল =
(√(3)/4)a\^{}2=(√(3)/4)×16\^{}2=(√(3)/4)×256=64√(3) বর্গমিটার।}
\end{quote}

\textbf{অতিভুজের বিপরীতে থাকে-\/-\/-}

\textbf{\hl{ক) সমকোণ ✔}}

খ) সরলকোণ

গ) স্থূলকোণ

ঘ) সূক্ষ্মকোণ

\textbf{\hl{✔ সঠিক উত্তর: ক) সমকোণ}}

\begin{quote}
\textbf{ব্যাখ্যা:} \emph{সমকোণী ত্রিভুজের সমকোণের বিপরীত বাহুকে অতিভুজ বলা হয়।
সুতরাং অতিভুজের বিপরীতে সমকোণ অবস্থান করে।}
\end{quote}

\textbf{একটি ত্রিভুজের তিনটি বাহুর দৈর্ঘ্য ৫, ৬, ৭ মিটার। নিকটতম বর্গমিটারে
ত্রিভুজটির ক্ষেত্রফল কত?}

ক) ১৬ বর্গমিটার

\textbf{\hl{খ) ১৫ বর্গমিটার ✔}}

গ) ১৭ বর্গমিটার

ঘ) ১৪ বর্গমিটার

\textbf{\hl{✔ সঠিক উত্তর: খ) ১৫ বর্গমিটার}}

\begin{quote}
\textbf{ব্যাখ্যা:} \emph{ত্রিভুজের অর্ধপরিসীমা s=(5+6+7/2)=9 মিটার}

\emph{ক্ষেত্রফল = √(s(s−a)(s−b)(s−c))}

\emph{⇒ √(9(9−5)(9−6)(9−7))}

\emph{⇒ √(9×4×3×2)=√(216)≈14.69 বর্গমিটার, যা নিকটতম পূর্ণসংখ্যায় ১৫
বর্গমিটার।}
\end{quote}

\textbf{একটি সমবাহু ত্রিভুজের ক্ষেত্রফল ৫০ বর্গ সেঃ মিঃ। ত্রিভুজের প্রতি বাহুর
দৈর্ঘ্য কত?}

ক) ১৫.২ সেঃ মিঃ

খ) ১০.৫ সেঃ মিঃ

\textbf{\hl{গ) ১০.৭ সেঃ মিঃ ✔}}

ঘ) ১৭.১ সেঃ মিঃ

\textbf{\hl{✔ সঠিক উত্তর: গ) ১০.৭ সেঃ মিঃ}}

\begin{quote}
\textbf{ব্যাখ্যা:} \emph{ধরি, সমবাহু ত্রিভুজের বাহুর দৈর্ঘ্য a}

\emph{শর্তমতে, (√(3)/4)a\^{}2=50}

\emph{⇒ √(3)a\^{}2=200}

\emph{⇒ a\^{}2=(200/1.732)≈115.47}

\emph{⇒ a≈10.74 সেমি (প্রায় ১০.৭ সেমি)।}
\end{quote}

\textbf{একটি সমকোণী ত্রিভুজের দ্বিতীয় কোণটির মাপ ২০ ডিগ্রি হলে তৃতীয় কোণটির
মাপ-\/-\/-}

ক) ৪০ ডিগ্রি

খ) ৫০ ডিগ্রি

গ) ৬০ ডিগ্রি

\textbf{\hl{ঘ) ৭০ ডিগ্রি ✔}}

\textbf{\hl{✔ সঠিক উত্তর: ঘ) ৭০ ডিগ্রি}}

\begin{quote}
\textbf{ব্যাখ্যা:} \emph{আমরা জানি, ত্রিভুজের তিন কোণের সমষ্টি 180° এবং
সমকোণী ত্রিভুজের একটি কোণ নির্দিষ্টভাবেই 90° ।}

\emph{∴ তৃতীয় কোণটি = 180°−(90°+20°)=180°−110°=70° ।}
\end{quote}

\textbf{একটি সমকোণী সমদ্বিবাহু ত্রিভুজের তিন কোণের পরিমাণ-\/-\/-}

ক) ৬০ ডিগ্রি, ৬০ ডিগ্রি, ৬০ ডিগ্রি

খ) ৪০ ডিগ্রি, ৯০ ডিগ্রি, ৪০ ডিগ্রি

গ) ৫০ ডিগ্রি, ৯০ ডিগ্রি, ৪০ ডিগ্রি

\textbf{\hl{ঘ) ৪৫ ডিগ্রি, ৯০ ডিগ্রি, ৪৫ ডিগ্রি ✔}}

\textbf{\hl{✔ সঠিক উত্তর: ঘ) ৪৫ ডিগ্রি, ৯০ ডিগ্রি, ৪৫ ডিগ্রি}}

\begin{quote}
\textbf{ব্যাখ্যা:} \emph{সমকোণী সমদ্বিবাহু ত্রিভুজের একটি কোণ এক সমকোণ বা 90°
এবং অপর দুটি কোণ সমান।}

\emph{∴ অপর দুটি কোণের প্রত্যেকটির পরিমাণ হবে = (180°−90°)÷2=45° ।}
\end{quote}

\textbf{কোন ত্রিভুজের এক বাহুর উপর অঙ্কিত বর্গক্ষেত্র অপর দুই বাহুর উপর অঙ্কিত
বর্গক্ষেত্রদ্বয়ের যোগফলের সমান হলে ত্রিভুজটি হবে-\/-}

\textbf{\hl{ক) সমকোণী ✔}}

খ) সূক্ষ্মকোণী

গ) সমবাহু

ঘ) স্থূলকোণী

\textbf{\hl{✔ সঠিক উত্তর: ক) সমকোণী}}

\begin{quote}
\textbf{ব্যাখ্যা:} \emph{পীথাগোরাসের উপপাদ্য অনুসারে, সমকোণী ত্রিভুজের অতিভুজের
উপর অঙ্কিত বর্গক্ষেত্র অপর দুই বাহুর উপর অঙ্কিত বর্গক্ষেত্রদ্বয়ের সমষ্টির সমান।}
\end{quote}

\textbf{যে ত্রিভুজের দুটি বাহু পরস্পর সমান, তা-\/-\/-}

\textbf{\hl{ক) সমদ্বিবাহু ত্রিভুজ ✔}}

খ) সমবাহু ত্রিভুজ

গ) বিষমবাহু ত্রিভুজ

ঘ) বিপরীত বাহু ত্রিভুজ

\textbf{\hl{✔ সঠিক উত্তর: ক) সমদ্বিবাহু ত্রিভুজ}}

\begin{quote}
\textbf{ব্যাখ্যা:} \emph{যে ত্রিভুজের দুটি বাহু পরস্পর সমান, তাকে সমদ্বিবাহু
ত্রিভুজ বলে।}
\end{quote}

\textbf{কোনো ত্রিভুজের একটি কোণ অপর দুটি কোণের সমান হলে ত্রিভুজটি-\/-\/-}

ক) সমবাহু

\textbf{\hl{খ) সমকোণী ✔}}

গ) সূক্ষ্মকোণী

ঘ) স্থূলকোণী

\textbf{\hl{✔ সঠিক উত্তর: খ) সমকোণী}}

\begin{quote}
\textbf{ব্যাখ্যা:} \emph{ত্রিভুজের তিন কোণের সমষ্টি 180° । একটি কোণ অপর দুটি
কোণের যোগফলের সমান হলে সেই কোণটির মান হবে 180°÷2=90° , অর্থাৎ ত্রিভুজটি
সমকোণী।}
\end{quote}

\textbf{একটি সমদ্বিবাহু ত্রিভুজের একটি কোণের পরিমাণ ৫০ ডিগ্রি হলে অপর দুটি
কোণের প্রত্যেকটির পরিমাণ হবে-\/-}

ক) ৫৫ ডিগ্রি

খ) ৭৫ ডিগ্রি

গ) ৬০ ডিগ্রি

\textbf{\hl{ঘ) ৬৫ ডিগ্রি ✔}}

\textbf{\hl{✔ সঠিক উত্তর: ঘ) ৬৫ ডিগ্রি}}

\begin{quote}
\textbf{ব্যাখ্যা:} \emph{ত্রিভুজের তিন কোণের সমষ্টি 180° । একটি কোণ 50° হলে
অপর দুটি কোণের সমষ্টি 180°−50°=130° ।}

\emph{∴ অপর দুটি কোণের প্রত্যেকটির পরিমাণ হবে = 130°÷2=65° ।}
\end{quote}

\textbf{সমকোণী ত্রিভুজের অতিভুজ সংলগ্ন কোণ দুটির প্রত্যেকটি-\/-\/-}

ক) সমকোণ

\textbf{\hl{খ) সূক্ষ্মকোণ ✔}}

গ) স্থূলকোণ

ঘ) সরলকোণ

\textbf{\hl{✔ সঠিক উত্তর: খ) সূক্ষ্মকোণ}}

\begin{quote}
\textbf{ব্যাখ্যা:} \emph{সমকোণী ত্রিভুজের একটি কোণ সমকোণ এবং অপর দুটি কোণের
সমষ্টি এক সমকোণ বা 90° ।}

\emph{∴ কোণ দুটির প্রত্যেকটি 90° এর চেয়ে ছোট অর্থাৎ সূক্ষ্মকোণ হয়।}
\end{quote}

\textbf{সমকোণী ত্রিভুজের একটি কোণ ৯০ ডিগ্রি হলে, অপর দুটি কোণের মান কত?}

ক) ১৮০ ডিগ্রি

খ) ১৯০ ডিগ্রি

গ) ২৭০ ডিগ্রি

\textbf{\hl{ঘ) ৯০ ডিগ্রি ✔}}

\textbf{\hl{✔ সঠিক উত্তর: ঘ) ৯০ ডিগ্রি}}

\begin{quote}
\textbf{ব্যাখ্যা:} \emph{আমরা জানি, ত্রিভুজের তিন কোণের সমষ্টি 180° ।}

\emph{∴ অপর দুটি কোণের মানের সমষ্টি = 180°−90°=90° ।}
\end{quote}

\textbf{দুইটি ত্রিভুজ সর্বসম হওয়ার জন্য নিচের কোন শর্তটি যথেষ্ট নয়?}

ক) একটির তিনবাহু অপরটির তিনবাহুর সমান

\textbf{\hl{খ) একটির তিনকোণ অপরটির তিনকোণের সমান ✔}}

গ) একটির দুই কোণ ও একবাহু অপরটির দুইকোণ ও অনুরূপ বাহুর সমান

ঘ) একটির দুইবাহু ও অন্তর্ভুক্ত কোণ অপরটির দুইবাহু ও অন্তর্ভুক্ত কোণের সমান

\textbf{\hl{✔ সঠিক উত্তর: খ) একটির তিনকোণ অপরটির তিনকোণের সমান}}

\begin{quote}
\textbf{ব্যাখ্যা:} \emph{ত্রিভুজের তিনটি কোণ সমান হলেই ত্রিভুজদ্বয় সর্বসম হবে
এমন কোনো নিশ্চয়তা নেই, কারণ তাদের বাহুগুলো ছোট-বড় হতে পারে (সদৃশ হতে পারে
কিন্তু সর্বসম নয়)।}
\end{quote}

\textbf{ত্রিভুজের বৃহত্তর বাহু সংলগ্ন কোণদ্বয়-\/-\/-}

\textbf{\hl{ক) সূক্ষ্মকোণ ✔}}

খ) স্থূলকোণ

গ) সমকোণ

ঘ) সরলকোণ

\textbf{\hl{✔ সঠিক উত্তর: ক) সূক্ষ্মকোণ}}

\begin{quote}
\textbf{ব্যাখ্যা:} \emph{ত্রিভুজের বৃহত্তর বাহু সংলগ্ন কোণদ্বয় সর্বদা 90° এর চেয়ে
ছোট বা সূক্ষ্মকোণ হয়ে থাকে।}
\end{quote}

\textbf{একটি ত্রিভুজের একটি কোণ যদি ২য় কোণের তিনগুণ এবং ৩য় কোণ যদি ২য় কোণের
চেয়ে ৩০ ডিগ্রি বড় হয় তবে কোণটি কত ডিগ্রি?}

\textbf{\hl{ক) ৬০ ডিগ্রি ✔}}

খ) ৫০ ডিগ্রি

গ) ৪৫ ডিগ্রি

ঘ) ৩০ ডিগ্রি

\textbf{\hl{✔ সঠিক উত্তর: ক) ৬০ ডিগ্রি}}

\begin{quote}
\textbf{ব্যাখ্যা:} \emph{ধরি, ২য় কোণ = ক}

\emph{১ম কোণ = 3ক}

\emph{এবং ৩য় কোণ = ক+30°}

\emph{শর্তমতে, ক+3ক+ক+30°=180°}

\emph{⇒ 5ক=180°−30°}

\emph{⇒ 5ক=150°}

\emph{⇒ ক=30°}

\emph{∴ ১ম কোণ = 3×30°=90°}

\emph{২য় কোণ =30°}

\emph{৩য় কোণ = 30° + 30° = 60°}
\end{quote}

\textbf{ত্রিভুজের তিন বাহুর উপর অঙ্কিত বর্গক্ষেত্রসমূহ কয়টি সমকোণ তৈরি করে?}

\textbf{\hl{ক) ১২ ✔}}

খ) ১৬

গ) ৮

ঘ) ৯

\textbf{\hl{✔ সঠিক উত্তর: ক) ১২}}

\begin{quote}
\textbf{ব্যাখ্যা:} \emph{ত্রিভুজের তিন বাহুর উপর মোট তিনটি বর্গক্ষেত্র অঙ্কন করা
যায়।}

\emph{প্রত্যেকটি বর্গক্ষেত্রের ৪টি করে কোণ থাকায় ৩টি বর্গক্ষেত্রে মোট কোণ হবে =
4×3=12 টি সমকোণ।}
\end{quote}

\textbf{সমকোণী ত্রিভুজের অতিভুজের দৈর্ঘ্য দেয়া থাকলে-\/-\/-}

ক) একটি ত্রিভুজ আঁকা যায়

খ) দুটি ত্রিভুজ আঁকা যায়

গ) তিনটি ত্রিভুজ আঁকা যায়

\textbf{\hl{ঘ) অনেকগুলো ত্রিভুজ আঁকা যায় ✔}}

\textbf{\hl{✔ সঠিক উত্তর: ঘ) অনেকগুলো ত্রিভুজ আঁকা যায়}}

\begin{quote}
\textbf{ব্যাখ্যা:} \emph{একটি সমকোণী ত্রিভুজের শুধুমাত্র অতিভুজের নির্দিষ্ট দৈর্ঘ্য
জানা থাকলে লম্ব ও ভূমির দৈর্ঘ্য পরিবর্তন করে অসীম সংখ্যক বা অনেকগুলো সমকোণী
ত্রিভুজ অঙ্কন করা সম্ভব।}
\end{quote}

\textbf{সমবাহু ত্রিভুজের বাহুর দৈর্ঘ্য ৬ মিটার হলে, ক্ষেত্রফল কত?}

ক) 8√(3)

\textbf{\hl{খ) 9√(3) ✔}}

গ) 12√(3)

ঘ) ১৬

\textbf{\hl{✔ সঠিক উত্তর: খ) 9√(3)}}

\begin{quote}
\textbf{ব্যাখ্যা:} \emph{সমবাহু ত্রিভুজের ক্ষেত্রফল = (√(3)/4)a\^{}2}

\emph{এখানে বাহু a=6 মিটার}

\emph{∴ ক্ষেত্রফল = (√(3)/4)×6\^{}2=(√(3)/4)×36=9√(3) বর্গমিটার।}
\end{quote}

\textbf{ত্রিভুজের তিন বাহু, উহার অন্তঃবৃত্তের-\/-\/-\/-\/-}

ক) জ্যা

খ) ব্যাস

গ) ব্যাসার্ধ

\textbf{\hl{ঘ) স্পর্শক ✔}}

\textbf{\hl{✔ সঠিক উত্তর: ঘ) স্পর্শক}}

\begin{quote}
\textbf{ব্যাখ্যা:} \emph{কোনো ত্রিভুজের অভ্যন্তরে অঙ্কিত অন্তঃবৃত্তের ক্ষেত্রে
ত্রিভুজের তিনটি বাহু বৃত্তটির তিনটি স্পর্শক হিসেবে কাজ করে।}
\end{quote}

\textbf{একটি সমকোণী ত্রিভুজের ক্ষেত্রফল ১৪৪ বর্গ একক। সমকোণ সংলগ্ন বাহুদ্বয়ের
একটির দৈর্ঘ্য ১২ একক হলে অপরটি কত?}

ক) ১৫ একক

\textbf{\hl{খ) ২৪ একক ✔}}

গ) ২০ একক

ঘ) ৩০ একক

\textbf{\hl{✔ সঠিক উত্তর: খ) ২৪ একক}}

\begin{quote}
\textbf{ব্যাখ্যা:} \emph{সমকোণী ত্রিভুজের ক্ষেত্রফল = (1/2)×ভূমি×উচ্চতা}

\emph{প্রশ্নমতে, (1/2)×12×অপর বাহু=144}

\emph{⇒ 6×অপর বাহু=144}

\emph{⇒ অপর বাহু=(144/6)=24 একক।}
\end{quote}

\textbf{কোন ত্রিভুজের বাহুগুলোর লম্ব-দ্বিখণ্ডক তিনটি যে বিন্দুতে ছেদ করে তাকে
বলে-\/-}

ক) ভরকেন্দ্র

\textbf{\hl{খ) পরিকেন্দ্র ✔}}

গ) অন্তঃকেন্দ্র

ঘ) লম্ববিন্দু

\textbf{\hl{✔ সঠিক উত্তর: খ) পরিকেন্দ্র}}

\begin{quote}
\textbf{ব্যাখ্যা:} \emph{ত্রিভুজের বাহুগুলোর লম্ব-দ্বিখণ্ডকত্রয় যে বিন্দুতে মিলিত হয়
তাকে পরিকেন্দ্র বলা হয়।}
\end{quote}

\textbf{ত্রিভুজের দুইটি কোণের সমষ্টি তৃতীয় কোণের সমান হলে ত্রিভুজটি-\/-\/-}

ক) সমবাহু

খ) সমদ্বিবাহু

\textbf{\hl{গ) সমকোণী ✔}}

ঘ) কোনোটিই নয়

\textbf{\hl{✔ সঠিক উত্তর: গ) সমকোণী}}

\begin{quote}
\textbf{ব্যাখ্যা:} \emph{ত্রিভুজের তিন কোণের সমষ্টি 180° ।}

\emph{একটি কোণ অপর দুটি কোণের যোগফলের সমান হলে সেই কোণটির মান হবে
180°÷2=90° , অর্থাৎ ত্রিভুজটি সমকোণী।}
\end{quote}

\textbf{ΔABC সমবাহু ত্রিভুজে BC বাহুর মধ্যমা AD হলে ∠BAD-এর মান কত?}

ক) 45°

\textbf{\hl{খ) 30° ✔}}

গ) 60°

ঘ) 40°

\textbf{\hl{✔ সঠিক উত্তর: খ) 30°}}

\begin{quote}
\textbf{ব্যাখ্যা:} \emph{ΔABC সমবাহু ত্রিভুজ হওয়ায় মধ্যমা AD, ∠BAC কে সমান
দুইভাগে ভাগ করে.}

\emph{আমরা জানি, সমবাহু ত্রিভুজের প্রত্যেক কোণ সমান অর্থাৎ 60°.}

\emph{∴ ∠BAD = (1/2) ∠BAC = (1/2) × 60° = 30°}
\end{quote}

\textbf{ΔABC একটি সমবাহু ত্রিভুজ, উহার AB ও AC বাহুকে বর্ধিত করলে উৎপন্ন
কোণদ্বয়ের সমষ্টি কত?}

\textbf{\hl{ক) 240° ✔}}

খ) 270°

গ) 280°

ঘ) 320°

\textbf{\hl{✔ সঠিক উত্তর: ক) 240°}}

\begin{quote}
\textbf{ব্যাখ্যা:} \emph{ABC সমবাহু ত্রিভুজের ∠A = ∠B = ∠C = 60°.}

\emph{ত্রিভুজের কোনো একটি বাহু বর্ধিত করলে উৎপন্ন বহিঃস্থ কোণ বিপরীত অন্তঃস্থ
কোণদ্বয়ের সমষ্টির সমান.}

\emph{∴ ∠CBD = ∠A + ∠C = 60° + 60° = 120°}

\emph{এবং ∠BCE = ∠A + ∠B = 60° + 60° = 120°}

\emph{∴ ∠CBD + ∠BCE = 120° + 120° = 240°}
\end{quote}

\textbf{ABC এর ∠ABC = 45°, ∠BAC = 60° এবং BC বাহুকে D পর্যন্ত বর্ধিত করা হলে
∠ACD = কত?}

ক) 45°

খ) 60°

\textbf{\hl{গ) 105° ✔}}

ঘ) 120°

\textbf{\hl{✔ সঠিক উত্তর: গ) 105°}}

\begin{quote}
\textbf{ব্যাখ্যা:} \emph{আমরা জানি, ত্রিভুজের যেকোনো একটি বাহুকে বর্ধিত করলে
যে বহিঃস্থ কোণ উৎপন্ন হয়, তা বিপরীত অন্তঃস্থ কোণদ্বয়ের সমষ্টির সমান।}

\emph{কাজেই ∠ACD = ∠ABC + ∠BAC = 45° + 60° = 105°}
\end{quote}

\textbf{ABC ত্রিভুজের AB = AC এবং ∠B = 40° হলে, A এর মান কত?}

ক) 80°

খ) 70°

গ) 90°

\textbf{\hl{ঘ) 100° ✔}}

\textbf{\hl{✔ সঠিক উত্তর: ঘ) 100°}}

\begin{quote}
\textbf{ব্যাখ্যা:} \emph{যদি AB = AC হয়, তবে ∠B = ∠C = 40°}

\emph{∴ ∠A = 180° − (40° + 40°) = 100°}
\end{quote}

\textbf{১টি সমবাহু ত্রিভুজের সবগুলো বাহু ২০\% বৃদ্ধি পাওয়াতে ত্রিভুজটির ক্ষেত্রফল
কত শতাংশ বৃদ্ধি পাবে?}

\textbf{\hl{ক) ৪৪\% ✔}}

খ) ৪০\%

গ) ২০\%

ঘ) ৮\%

\textbf{\hl{✔ সঠিক উত্তর: ক) ৪৪\%}}

\begin{quote}
\textbf{ব্যাখ্যা:} \emph{ক্ষেত্রফল শতকরা বৃদ্ধি পাবে = (20 + 20 + (20×20/100)
)\% = 44\%}
\end{quote}

\textbf{সমদ্বিবাহু ত্রিভুজের সমান সমান বাহুদ্বয় বর্ধিত করলে উৎপন্ন কোণদ্বয় কি
হবে?}

ক) সূক্ষ্মকোণ

\textbf{\hl{খ) স্থূলকোণ ✔}}

গ) প্রবৃদ্ধকোণ

ঘ) সমকোণ

\textbf{\hl{✔ সঠিক উত্তর: খ) স্থূলকোণ}}

\begin{quote}
\textbf{ব্যাখ্যা:} \emph{সমদ্বিবাহু ত্রিভুজের সমান সমান বাহুদ্বয় বর্ধিত করলে
উৎপন্ন কোণদ্বয় স্থূলকোণ হয়।}
\end{quote}

\textbf{একটি সমদ্বিবাহু ত্রিভুজের একটি কোণের মান 120° হলে, অপর যেকোনো একটি
কোণের মান কত?}

ক) 80°

\textbf{\hl{খ) 30° ✔}}

গ) 60°

ঘ) 40°

\textbf{\hl{✔ সঠিক উত্তর: খ) 30°}}

\begin{quote}
\textbf{ব্যাখ্যা:} \emph{ত্রিভুজের ৩ কোণের সমষ্টি 180°}

\emph{কাজেই একটি কোণ 120° হলে অপর দুই কোণের সমষ্টি হবে = 180° − 120° =
60°}

\emph{অতএব, অপর যেকোনো একটি কোণের মান = (60°/2) = 30° {[}যেহেতু সমদ্বিবাহু
ত্রিভুজের দুটি কোণ পরস্পর সমান{]}}
\end{quote}

\textbf{সমদ্বিবাহু ত্রিভুজের একটি কোণ ৫০ ডিগ্রি হলে অপর দুটি কত ডিগ্রি?}

ক) ৫০ ও ৬০

খ) ৪০ ও ৭০

\textbf{\hl{গ) ৬৫ ও ৬৫ ✔}}

ঘ) ৬০ ও ৬০

\textbf{\hl{✔ সঠিক উত্তর: গ) ৬৫ ও ৬৫}}

\begin{quote}
\textbf{ব্যাখ্যা:} \emph{সমদ্বিবাহু ত্রিভুজের একটি কোণ ৫০° হলে অপর দুটি কোণের
সমষ্টি 180° − 50° = 130°}

\emph{যেহেতু ত্রিভুজটি সমদ্বিবাহু, তাই অপর দুটি কোণ সমান হবে। অপশনগুলোর মধ্যে
৬৫° ও ৬৫° কোণ দুটির সমষ্টি ১৩০°।}
\end{quote}

\textbf{সমদ্বিবাহু ত্রিভুজের সমান বাহু দুটির ছেদবিন্দুতে ত্রিভুজের যে কোণ উৎপন্ন হয়
তাকে বলা হয়-}

ক) পূরক কোণ

\textbf{\hl{খ) শীর্ষকোণ ✔}}

গ) সন্নিহিত কোণ

ঘ) সম্পূরক কোণ

\textbf{\hl{✔ সঠিক উত্তর: খ) শীর্ষকোণ}}

\begin{quote}
\textbf{ব্যাখ্যা:} \emph{সমদ্বিবাহু ত্রিভুজের সমান বাহু দুইটির ছেদ বিন্দুর বিপরীত
বাহুকে তার ভূমি (Base) এবং ঐ ছেদ বিন্দুতে ত্রিভুজের কোণকে তার শীর্ষকোণ (Vertex)
বলা হয়}
\end{quote}

\textbf{সমদ্বিবাহু ত্রিভুজের শীর্ষ কোণের মান 80° হলে, অপর কোণদ্বয়ের মান কত?}

\textbf{\hl{ক) 50° ও 50° ✔}}

খ) 60° ও 40°

গ) 45° ও 45°

ঘ) 40° ও 40°

\textbf{\hl{✔ সঠিক উত্তর: ক) 50° ও 50°}}

\begin{quote}
\textbf{ব্যাখ্যা:} \emph{সমদ্বিবাহু ত্রিভুজের শীর্ষ কোণ 80°}

\emph{ত্রিভুজের তিন কোণের সমষ্টি 180°}

\emph{⇒ অপর কোণদ্বয়ের সমষ্টি = 180° − 80° = 100°}

\emph{∴ প্রতিটি কোণ = (100°/2) = 50°}
\end{quote}

\textbf{ABC ত্রিভুজের AB = AC, ∠A = 40° হলে ∠B = কত?}

ক) 40°

\textbf{\hl{খ) 70° ✔}}

গ) 60°

ঘ) 100°

\textbf{\hl{✔ সঠিক উত্তর: খ) 70°}}

\begin{quote}
\textbf{ব্যাখ্যা:} \emph{ΔABC-এ AB = AC}

\emph{∴ ∠C = ∠B}

\emph{∠A + ∠B + ∠C = 180°}

\emph{⇒ 2∠B = 180° − 40°}

\emph{∴ ∠B = (140°/2) = 70°}
\end{quote}

\textbf{চিত্রে, ∠PQR=55°, ∠LRN=90° এবং PQ\textbar\textbar MR, PQ=PR হলে,
∠NRP এর মান নিচের কোনটি?}

ক) 90°

খ) 55°

গ) 45°

\textbf{\hl{ঘ) 35° ✔}}

\textbf{\hl{✔ সঠিক উত্তর: ঘ) 35°}}

\begin{quote}
\textbf{ব্যাখ্যা:} \emph{ΔPQR একটি সমদ্বিবাহু ত্রিভুজ।}

\emph{এখানে, ∠PQR = ∠PRQ {[}∴ PQ = PR{]}}

\emph{∠PQR = ∠PRQ = 55°}

\emph{আবার, ∠LRN = ∠NRQ = 90°}

\emph{∴ ∠NRP = 90° − ∠PRQ = 90° − 55° = 35°}
\end{quote}

\textbf{ABC ত্রিভুজের AB = AC এবং ∠A = 80° হলে ∠B = কত?}

ক) 40°

\textbf{\hl{খ) 50° ✔}}

গ) 60°

ঘ) 80°

\textbf{\hl{✔ সঠিক উত্তর: খ) 50°}}

\begin{quote}
\textbf{ব্যাখ্যা:} \emph{ABC ত্রিভুজে AB = AC হলে ∴ ∠B = ∠C}

\emph{ত্রিভুজের তিন কোণের সমষ্টি = 180°}

\emph{⇒ ∠A + ∠B + ∠C = 180°}

\emph{⇒ 80° + 2∠B = 180°}

\emph{⇒ 2∠B = 180° − 80° = 100°}

\emph{∴ ∠B = (100°/2) = 50°}
\end{quote}

\textbf{ΔABC একটি সমদ্বিবাহু ত্রিভুজ যার AB = AC এবং ∠A = 50° হলে, ∠B = কত
হবে?}

ক) ৪৫°

খ) ৯০°

গ) ৭৫°

\textbf{\hl{ঘ) ৬৫° ✔}}

\textbf{\hl{✔ সঠিক উত্তর: ঘ) ৬৫°}}

\begin{quote}
\textbf{ব্যাখ্যা:} \emph{একটি সমদ্বিবাহু ত্রিভুজ এর AB = AC হলে বিপরীত কোণদ্বয়
∠B = ∠C}

\emph{ত্রিভুজের তিন কোণের সমষ্টি = ১৮০°}

\emph{বা, ∠A + ∠B + ∠C = ১৮০°}

\emph{বা, 2∠B = ১৮০° − ৫০°}

\emph{∴ ∠B = (১৩০°/২) = ৬৫°}
\end{quote}

\textbf{ABC ত্রিভুজের AB = AC, ∠A = 80° , ∠B = কত?}

ক) 40°

\textbf{\hl{খ) 50° ✔}}

গ) 60°

ঘ) 100°

\textbf{\hl{✔ সঠিক উত্তর: খ) 50°}}

\begin{quote}
\textbf{ব্যাখ্যা:} \emph{Δ ABC ত্রিভুজে, AB = AC ∴ ∠B = ∠C}

\emph{Δ ABC - ∠A + ∠B + ∠C = 180°}

\emph{80° + ∠B + ∠B = 180°}

\emph{2∠B = 100°}

\emph{∠B = 50°}
\end{quote}

\textbf{ABC একটি সমদ্বিবাহু ত্রিভুজ, ∠A শীর্ষ কোণ, ∠B ও ∠C দুটি ভূমি কোণ। AB
বাহু = AC বাহু। ∠B = 75° । ∠A = কত ডিগ্রি?}

ক) 25°

\textbf{\hl{খ) 30° ✔}}

গ) 40°

ঘ) 50°

\textbf{\hl{✔ সঠিক উত্তর: খ) 30°}}

\begin{quote}
\textbf{ব্যাখ্যা:} \emph{Δ ABC ত্রিভুজে, AB = AC সুতরাং ∠B = ∠C}

\emph{Δ ABC - ∠A + ∠B + ∠C = 180°}

\emph{∠A + 75° + 75° = 180°}

\emph{∠A = 30°}
\end{quote}

\textbf{সমদ্বিবাহু ত্রিভুজের সমান সমান বাহুদ্বয় বর্ধিত করলে উৎপন্ন কোণদ্বয় হবে-}

\textbf{\hl{ক) স্থূলকোণ ✔}}

খ) সমকোণ

গ) সূক্ষ্মকোণ

ঘ) সরলকোণ

\textbf{\hl{✔ সঠিক উত্তর: ক) স্থূলকোণ}}

\begin{quote}
\textbf{ব্যাখ্যা:} \emph{সমদ্বিবাহু ত্রিভুজের সমান সমান বাহুদ্বয় বর্ধিত করলে
উৎপন্ন বহিঃস্থ কোণদ্বয় স্থূলকোণ হয়।}
\end{quote}

\textbf{একটি সমদ্বিবাহু ত্রিভুজের একটি কোণের মান 94° হলে অপর একটি কোণের মান
হবে-}

\textbf{\hl{ক) 43° ✔}}

খ) 47°

গ) 94°

ঘ) 88°

\textbf{\hl{✔ সঠিক উত্তর: ক) 43°}}

\begin{quote}
\textbf{ব্যাখ্যা:} \emph{সমদ্বিবাহু ত্রিভুজের সমান সমান বাহুর বিপরীত কোণগুলো
সমান। যেহেতু ত্রিভুজের তিন কোণের সমষ্টি ১৮০°, তাই ৯৪° কোণটি অবশ্যই শীর্ষকোণ
হতে হবে (কারণ দুটি ভূমি কোণ ৯৪° হলে তাদের সমষ্টি ১৮০° ছাড়িয়ে যায়)।}

\emph{∴ অপর দুটি কোণের সমষ্টি = ১৮০° − ৯৪° = ৮৬°}

\emph{∴ প্রতিটি সমান কোণের মান = (৮৬°/২) = ৪৩°}
\end{quote}

\textbf{একটি সমদ্বিবাহু ত্রিভুজের একটি কোণ ৯৬°। তাহলে ঐ ত্রিভুজের অপর একটি
কোণের মান হবে -}

ক) ৪৮°

খ) ৯৬°

গ) ৮৪°

\textbf{\hl{ঘ) ৪২° ✔}}

\textbf{\hl{✔ সঠিক উত্তর: ঘ) ৪২°}}

\begin{quote}
\textbf{ব্যাখ্যা:} \emph{ত্রিভুজটির অপর কোণ দুটির সমষ্টি = 180° − 96° = 84°}

\emph{∴ অপর কোণদ্বয়ের প্রতিটি = (84°/2) = 42°}
\end{quote}

\textbf{বিষমবাহু ত্রিভুজের-}

ক) ৩টি বাহু সমান

খ) ৩টি কোণ সমান

\textbf{\hl{গ) ৩টি বাহু অসমান ✔}}

ঘ) ৩টি কোণ অসমান

\textbf{\hl{✔ সঠিক উত্তর: গ) ৩টি বাহু অসমান}}

\begin{quote}
\textbf{ব্যাখ্যা:} \emph{যে ত্রিভুজের তিনটি বাহুই অসমান তাকে বিষম বাহু ত্রিভুজ
বলে। উল্লেখ্য, ত্রিভুজের তিনটি বাহু অসমান হলে কোণ তিনটিও অসমান হবে।}
\end{quote}

\textbf{সমকোণী ত্রিভুজের কোনটি বৃহত্তম বাহু?}

ক) ভূমি

\textbf{\hl{খ) অতিভুজ ✔}}

গ) লম্ব

ঘ) কোনটিই নয়

\textbf{\hl{✔ সঠিক উত্তর: খ) অতিভুজ}}

\begin{quote}
\textbf{ব্যাখ্যা:} \emph{সমকোণী ত্রিভুজের অতিভুজই বৃহত্তম বাহু। কারণ অতিভুজের উপর
অঙ্কিত বর্গ ভূমি ও লম্বের উপর অঙ্কিত বর্গের সমষ্টির সমান। অর্থাৎ (অতিভুজ)\^{}2 =
(ভূমি)\^{}2 + (লম্ব)\^{}2}
\end{quote}

\textbf{সমকোণী ত্রিভুজের অতিভুজের দৈর্ঘ্য দেওয়া থাকলে-}

ক) একটিমাত্র ত্রিভুজ আঁকা যায়

খ) দুটি ত্রিভুজ আঁকা যায়

গ) কোন ত্রিভুজ আঁকা যায় না

\textbf{\hl{ঘ) অনেকগুলো ত্রিভুজ আঁকা যায় ✔}}

\textbf{\hl{✔ সঠিক উত্তর: ঘ) অনেকগুলো ত্রিভুজ আঁকা যায়}}

\begin{quote}
\textbf{ব্যাখ্যা:} \emph{সমকোণী ত্রিভুজের অতিভুজের দৈর্ঘ্য দেওয়া থাকলে অনেকগুলো
ত্রিভুজ আঁকা যায়। একটি রেখার উপর লম্ব অঙ্কন করে অতিভুজের বিভিন্ন দৈর্ঘ্য নিয়ে
বিভিন্ন ত্রিভুজ আঁকা যায়।}
\end{quote}

\textbf{সমকোণী ত্রিভুজের সূক্ষ্মকোণদ্বয়ের একটি ৪০ ডিগ্রি হলে অপরটি কত?}

ক) ৪৫ ডিগ্রি

খ) ৫৫ ডিগ্রি

\textbf{\hl{গ) ৫০ ডিগ্রি ✔}}

ঘ) ৬০ ডিগ্রি

\textbf{\hl{✔ সঠিক উত্তর: গ) ৫০ ডিগ্রি}}

\begin{quote}
\textbf{ব্যাখ্যা:} \emph{সমকোণী ত্রিভুজের অপর সূক্ষ্ম কোণটি}

\emph{\{180° − (90° + 40°)\} = 180° − 130° = 50°}
\end{quote}

\textbf{যে ত্রিভুজের দুইটি কোণের পরিমাপ ৩০° ও ৬০° সেটি কোন ধরনের ত্রিভুজ?}

ক) সমবাহু

খ) সমদ্বিবাহু

গ) সূক্ষ্মকোণী

\textbf{\hl{ঘ) সমকোণী ✔}}

\textbf{\hl{✔ সঠিক উত্তর: ঘ) সমকোণী}}

\begin{quote}
\textbf{ব্যাখ্যা:} \emph{ত্রিভুজের অপর কোণটি = 180° − (30° + 60°) = 180° −
90° = 90°}

\emph{∴ ত্রিভুজটি সমকোণী।}
\end{quote}

\textbf{সমকোণী ত্রিভুজের সমকোণের বিপরীতে একটি কোণ ৫০° হলে অপরটি কত?}

ক) ২০°

খ) ৩০°

\textbf{\hl{গ) ৪০° ✔}}

ঘ) ৫০°

\textbf{\hl{✔ সঠিক উত্তর: গ) ৪০°}}

\begin{quote}
\textbf{ব্যাখ্যা:} \emph{দেওয়া আছে, সমকোণী ত্রিভুজের সূক্ষ্মকোণদ্বয়ের একটি =
50°}

\emph{আমরা জানি, সমকোণী ত্রিভুজের ভূমি ও লম্ব সংলগ্ন কোণ = 90°}

\emph{∴ নির্ণেয় কোণটি হবে = 180° − (90° + 50°) {[}আমরা জানি ত্রিভুজের তিন
কোণের সমষ্টি 180°{]}}

\emph{= (180° − 140°) = 40°}
\end{quote}

\textbf{সমকোণী ত্রিভুজের অতিভুজ সংলগ্ন কোণ দুইটির প্রত্যেকটি}

ক) সরল কোণ

\textbf{\hl{খ) সূক্ষ্ম কোণ ✔}}

গ) পূরক কোণ

ঘ) স্থূল কোণ

\textbf{\hl{✔ সঠিক উত্তর: খ) সূক্ষ্ম কোণ}}

\begin{quote}
\textbf{ব্যাখ্যা:} \emph{আমরা জানি, ত্রিভুজের তিন কোণের সমষ্টি = দুই সমকোণ।
সমকোণী ত্রিভুজের একটি কোণ = 90° সুতরাং অপর দুটি কোণ 90° এর চেয়ে ছোট।}

\emph{আমরা জানি, 90° এর চেয়ে ছোট কোণকে সূক্ষ্ম কোণ বলে।}
\end{quote}

\textbf{কোন সমকোণী ত্রিভুজের সূক্ষ্মকোণ দ্বয়ের পার্থক্য ২০ ডিগ্রি হলে, ক্ষুদ্রতম
কোণের মান কত ডিগ্রি?}

\textbf{\hl{ক) ৩৫ ডিগ্রি ✔}}

খ) ৪০ ডিগ্রি

গ) ৪৫ ডিগ্রি

ঘ) ৫৫ ডিগ্রি

\textbf{\hl{✔ সঠিক উত্তর: ক) ৩৫ ডিগ্রি}}

\begin{quote}
\textbf{ব্যাখ্যা:} \emph{ধরি, ক্ষুদ্রতম সূক্ষ্মকোণ x এবং বৃহত্তর সূক্ষ্মকোণ হবে (x +
20)}

\emph{প্রশ্নমতে, x + x + 20 + 90 = 180° {[}যেহেতু ৩ কোণের সমষ্টি 180°{]}}

\emph{⇒ 2x = 180° − 110° = 70°}

\emph{∴ x = (70°/2) = 35°}

\emph{অর্থাৎ ক্ষুদ্রতম কোণের মান 35°}
\end{quote}

\textbf{সমকোণী ত্রিভুজের কোনটি বৃহত্তম বাহু?}

ক) ভূমি

\textbf{\hl{খ) অতিভুজ ✔}}

গ) লম্ব

ঘ) কোনোটিই নয়

\textbf{\hl{✔ সঠিক উত্তর: খ) অতিভুজ}}

\begin{quote}
\textbf{ব্যাখ্যা:} \emph{সমকোণী ত্রিভুজের সমকোণের বিপরীত বাহু বৃহত্তম বাহু।}
\end{quote}

\textbf{ΔABC এর BC বাহুকে D বিন্দু পর্যন্ত বর্ধিত করা হলো। ∠A = 60°, ∠B = 90°
হলে, ∠ACD এর মান -}

\textbf{\hl{ক) 150° ✔}}

খ) 180°

গ) 140°

ঘ) 160°

\textbf{\hl{✔ সঠিক উত্তর: ক) 150°}}

\begin{quote}
\textbf{ব্যাখ্যা:} \emph{ত্রিভুজের কোনো বাহু বর্ধিত করা হলে উৎপন্ন বহিঃস্থ কোণের
মান বিপরীত অন্তঃস্থ কোণদ্বয়ের সমষ্টির সমান।}

\emph{∴ ∠ACD = ∠ABC + ∠BAC = 90° + 60° = 150°}
\end{quote}

\textbf{একটি সমকোণী ত্রিভুজের সমকোণ ছাড়া অন্য দুটি কোণ-\/-}

ক) সরল কোণ (a straight angle)

খ) পূরক কোণ (Complementary angles)

\textbf{\hl{গ) সূক্ষ্ম কোণ (an acute angle) ✔}}

ঘ) সন্নিহিত কোণ (Adjacent angles)

\textbf{\hl{✔ সঠিক উত্তর: গ) সূক্ষ্ম কোণ (an acute angle)}}

\begin{quote}
\textbf{ব্যাখ্যা:} \emph{একটি সমকোণী ত্রিভুজের সমকোণ ব্যতীত অন্য কোণ দুটির
সমষ্টি ৯০°। সুতরাং অন্য কোণ দুটি সূক্ষ্মকোণ।}
\end{quote}

\textbf{একটি সমকোণী ত্রিভুজের সমকোণ সংলগ্ন বাহু দুটির দৈর্ঘ্য ৮ ইঞ্চি ও ৬ ইঞ্চি
হলে ত্রিভুজটির তৃতীয় বাহুর দৈর্ঘ্য কত?}

ক) ৯ ইঞ্চি

\textbf{\hl{খ) ১০ ইঞ্চি ✔}}

গ) ১ ইঞ্চি

ঘ) ১১ ইঞ্চি

\textbf{\hl{✔ সঠিক উত্তর: খ) ১০ ইঞ্চি}}

\begin{quote}
\textbf{ব্যাখ্যা:} \emph{পিথাগোরাসের উপপাদ্য অনুযায়ী, অতিভুজ =
√(8\^{}2+6\^{}2) = √(64+36) = √(100) = 10 ইঞ্চি।}
\end{quote}

\textbf{সমকোণী ত্রিভুজের সমকোণ সংলগ্ন বাহুদ্বয় যথাক্রমে ৩ ও ৪ মি.মি হলে এর
অতিভুজের মান কত?}

\textbf{\hl{ক) ৫ মি.মি ✔}}

খ) ৬ সে.মি

গ) ৭ সে.মি

ঘ) ৯ সে. মি

\textbf{\hl{✔ সঠিক উত্তর: ক) ৫ মি.মি}}

\begin{quote}
\textbf{ব্যাখ্যা:} \emph{অতিভুজ = √(3\^{}2+4\^{}2) = √(9+16) = √(25) = 5
মি.মি।}
\end{quote}

\textbf{একটি সমকোণী ত্রিভুজের সমকোণ ব্যতীত অপর একটি কোণের মান 60° হলে, অপর
কোণটি-\/-}

\textbf{\hl{ক) 30° ✔}}

খ) 45°

গ) 60°

ঘ) 50°

\textbf{\hl{✔ সঠিক উত্তর: ক) 30°}}

\begin{quote}
\textbf{ব্যাখ্যা:} \emph{অপর কোণটি = 180° − (90° + 60°) = 180° − 150° =
30°}
\end{quote}

\textbf{কোন সমকোণী ত্রিভুজের সূক্ষ্মকোণদ্বয়ের পার্থক্য ১৬° হলে ক্ষুদ্রতম কোণের
মান-\/-}

ক) ৪১°

খ) ৩৬°

গ) ৩৮°

\textbf{\hl{ঘ) ৩৭° ✔}}

\textbf{\hl{✔ সঠিক উত্তর: ঘ) ৩৭°}}

\begin{quote}
\textbf{ব্যাখ্যা:} \emph{ক্ষুদ্রতম কোণের মান = (৯০ − ১৬) ÷ ২ = ৭৪ ÷ ২ = ৩৭°}
\end{quote}

\textbf{ত্রিভুজের একটি কোণ এর অপর দুটি কোণের সমষ্টির সমান হলে ত্রিভুজটি-}

ক) স্থূলকোণী

খ) সমবাহু

\textbf{\hl{গ) সমকোণী ✔}}

ঘ) সূক্ষ্মকোণী

\textbf{\hl{✔ সঠিক উত্তর: গ) সমকোণী}}

\begin{quote}
\textbf{ব্যাখ্যা:} \emph{সমকোণী ত্রিভুজের একটি কোণ সমকোণ এবং অপর দুটি কোণের
সমষ্টি এক সমকোণ।}
\end{quote}

\textbf{একটি সমকোণী ত্রিভুজের ২ টি কোণের সমষ্টি ১৭০° হলে অপর কোণটির মান
কত?}

ক) ৩০°

খ) ২০°

\textbf{\hl{গ) ১০° ✔}}

ঘ) ৮০°

\textbf{\hl{✔ সঠিক উত্তর: গ) ১০°}}

\begin{quote}
\textbf{ব্যাখ্যা:} \emph{সমকোণী ত্রিভুজের তিন কোণের সমষ্টি ১৮০°। ২ কোণের
সমষ্টি ১৭০° হলে অপর কোণের মান ১৮০° − ১৭০° = ১০°।}
\end{quote}

\textbf{একটি সমকোণী ত্রিভুজের একটি সূক্ষ্মকোণ ৩০°। অপরটি---}

ক) ৯০°

\textbf{\hl{খ) ৬০° ✔}}

গ) ৪৫°

ঘ) ৩০°

\textbf{\hl{✔ সঠিক উত্তর: খ) ৬০°}}

\begin{quote}
\textbf{ব্যাখ্যা:} \emph{সমকোণী ত্রিভুজের অপর সূক্ষ্মকোণ = 90° − 30° = 60°}
\end{quote}

\textbf{সমকোণী ত্রিভুজের একটি কোণ ৬০° হলে অপর কোণটি কত?}

\textbf{\hl{ক) ৩০° ✔}}

খ) ১২০°

গ) ৬০°

ঘ) ০°

\textbf{\hl{✔ সঠিক উত্তর: ক) ৩০°}}

\begin{quote}
\textbf{ব্যাখ্যা:} \emph{সমকোণী ত্রিভুজের একটি কোণ ৯০° এবং অপরটি ৬০° হলে,
তৃতীয় কোণটি = 180° − (90° + 60°) = 30°}
\end{quote}

\textbf{কোনো সমকোণী ত্রিভুজের সূক্ষ্ম কোণদ্বয়ের পার্থক্য ২০° হলে ক্ষুদ্রতম কোণটির
মান কত?}

\textbf{\hl{ক) ৩৫° ✔}}

খ) ৪০°

গ) ৪৫°

ঘ) ৫৫°

\textbf{\hl{✔ সঠিক উত্তর: ক) ৩৫°}}

\begin{quote}
\textbf{ব্যাখ্যা:} \emph{ধরি, বৃহত্তম সূক্ষ্মকোণটি x₁ এবং ক্ষুদ্রতম x₂}

\emph{তাহলে, x₁ + x₂ = 90° {[}∴ সমকোণী ত্রিভুজের সূক্ষ্মকোণদ্বয়ের সমষ্টি
৯০°{]}}

\emph{দেয়া আছে, x₁ − x₂ = 20°}

\emph{(-) করি, 2x₂ = 70°}

\emph{⇒ x₂ = 35°}

\emph{∴ ক্ষুদ্রতম কোণটির মান 35°}
\end{quote}

\textbf{ত্রিভুজের একটি কোণ-এর অপর দুটি কোণের সমষ্টির সমান হলে ত্রিভুজটি হবে -}

\textbf{\hl{ক) সমকোণী ✔}}

খ) স্থূলকোণী

গ) সমবাহু

ঘ) সূক্ষ্মকোণী

\textbf{\hl{✔ সঠিক উত্তর: ক) সমকোণী}}

\begin{quote}
\textbf{ব্যাখ্যা:} \emph{ত্রিভুজের তিন কোণের সমষ্টি ১৮০°। একটি কোণ অপর দুটির
সমষ্টির সমান হলে, সেই কোণটির মান অবশ্যই ১৮০° ÷ ২ = ৯০° হবে। আর যে ত্রিভুজের
একটি কোণ ৯০° তাকে সমকোণী ত্রিভুজ বলে।}
\end{quote}

\textbf{একটি ত্রিভুজের তিনটি কোণের অনুপাত ১ : ২ : ৩। ত্রিভুজটি হবে-}

ক) সমবাহু

খ) সূক্ষ্মকোণী

গ) স্থূলকোণী

\textbf{\hl{ঘ) সমকোণী ✔}}

\textbf{\hl{✔ সঠিক উত্তর: ঘ) সমকোণী}}

\begin{quote}
\textbf{ব্যাখ্যা:} \emph{ত্রিভুজের তিন কোণের সমষ্টি = ১৮০°}

\emph{সমকোণী ত্রিভুজের এক কোণ = ৯০° এবং অন্য দুই কোণের সমষ্টি ৯০°। এখানে প্রথম
দুই কোণের অনুপাতের সমষ্টি তৃতীয় কোণের অনুপাতের সমান। তাই ত্রিভুজটি সমকোণী
ত্রিভুজ।}
\end{quote}

\textbf{সমকোণী ত্রিভুজে সমকোণ ব্যতীত অপর দুই কোণের অনুপাত ৩ : ২ হলে কোণ দুইটি
কত?}

ক) 45°, 50°

খ) 55°, 35°

\textbf{\hl{গ) 54°, 36° ✔}}

ঘ) 50°, 40°

\textbf{\hl{✔ সঠিক উত্তর: গ) 54°, 36°}}

\begin{quote}
\textbf{ব্যাখ্যা:} \emph{সমকোণী ত্রিভুজের সমকোণ ব্যতীত অপর দুই কোণের সমষ্টি
90°}

\emph{ধরি, একটি কোণ 3x° এবং অপরটি 2x°}

\emph{প্রশ্নমতে, 3x° + 2x° = 90°}

\emph{⇒ 5x = 90° ∴ x = 18°}

\emph{অতএব, কোণ দুটির একটি হবে = (3 × 18)° = 54°}

\emph{এবং অপরটি হবে = (2 × 18)° = 36°}
\end{quote}

\textbf{একটি সমকোণী ত্রিভুজের সূক্ষ্মকোণদ্বয়ের পার্থক্য ৬°। ক্ষুদ্রতম কোণের মান
কত?}

ক) ৩২°

খ) ৩৬°

\textbf{\hl{গ) ৪২° ✔}}

ঘ) ৪৮°

\textbf{\hl{✔ সঠিক উত্তর: গ) ৪২°}}

\begin{quote}
\textbf{ব্যাখ্যা:} \emph{ধরি, একটি কোণের মান x°}

\emph{অপরটি (x + 6)°}

\emph{প্রশ্নমতে, x + x + 6 = 90}

\emph{2x = 90 − 6 = 84}

\emph{x = (84/2) = 42}

\emph{অর্থাৎ ক্ষুদ্রতম কোণের মান 42°}
\end{quote}

\textbf{কোনো সমকোণী ত্রিভুজের ভূমি a, উচ্চতা b এবং অতিভুজ c হলে কোনটি সঠিক?}

ক) a\^{}2 = b\^{}2 + c\^{}2

খ) b\^{}2 = c\^{}2 + a\^{}2

\textbf{\hl{গ) c\^{}2 = a\^{}2 + b\^{}2 ✔}}

ঘ) কোনোটিই নয়

\textbf{\hl{✔ সঠিক উত্তর: গ) c\^{}2 = a\^{}2 + b\^{}2}}

\begin{quote}
\textbf{ব্যাখ্যা:} \emph{সমকোণী ত্রিভুজের ক্ষেত্রে, পিথাগোরাসের উপপাদ্য
অনুসারে,}

\emph{অতিভুজ\^{}2 = ভূমি\^{}2 + লম্ব\^{}2}

\emph{⇒ c\^{}2 = a\^{}2 + b\^{}2}
\end{quote}

\textbf{\textquotesingle ক\textquotesingle{} প্রথমে তিন মাইল উত্তরে, পরে ৯
মাইল পূর্বে এবং তৎপর আবার ৯ মাইল উত্তরে যায়। শুরুর স্থান হতে তা দূরত্ব কত মাইল?}

ক) ২০

খ) ২৫

\textbf{\hl{গ) ১৫ ✔}}

ঘ) ৩০

\textbf{\hl{✔ সঠিক উত্তর: গ) ১৫}}

\begin{quote}
\textbf{ব্যাখ্যা:} \emph{প্রশ্নানুযায়ী চিত্র আঁকি।}

\emph{এখানে, AED একটি সমকোণী ত্রিভুজ।}

\emph{AD হলো মধ্যবর্তী দূরত্ব। কিন্তু DE = 9 + 3 = 12}

\emph{∴ পিথাগোরাসের সূত্রানুসারে,}

\emph{(AD)\^{}2 = (AE)\^{}2 + (DE)\^{}2 = 9\^{}2 + 12\^{}2 = 81 + 144 =
225}

\emph{⇒ AD = √(225) = 15}
\end{quote}

\textbf{একজন লোক একটি নির্দিষ্ট স্থান A থেকে যাত্রা শুরু করে ১২ কিলোমিটার উত্তর
দিকে গেল এবং সেখান থেকে ৫ কিলোমিটার পূর্ব দিকে গেল। যাত্রা শেষে সে A অবস্থান
থেকে কত দূরে থাকবে?}

ক) ১৭ কি. মি.

খ) ১৫ কি. মি.

গ) ১৪ কি. মি.

\textbf{\hl{ঘ) ১৩ কি. মি. ✔}}

\textbf{\hl{✔ সঠিক উত্তর: ঘ) ১৩ কি. মি.}}

\begin{quote}
\textbf{ব্যাখ্যা:} \emph{পিথাগোরাসের উপপাদ্য অনুযায়ী দূরত্ব =
√(12\^{}2+5\^{}2) = √(144+25) = √(169) = 13 কি. মি.}
\end{quote}

\textbf{মামুন বাড়ি থেকে ১০ মাইল পূর্বে গিয়ে ডান দিকে ফিরে সোজা ২০ মাইল
দক্ষিণ দিকে হেঁটে আবার ডান দিকে ফিরে সোজা ২৫ মাইল পশ্চিম দিকে হেঁটে সবুজের
বাড়ি পৌঁছায়। মামুনের বাড়ি থেকে সবুজের বাড়ির সোজা সুজি দূরত্ব কত মাইল?}

ক) ১৫

খ) ১৮

গ) ২০

\textbf{\hl{ঘ) ২৫ ✔}}

\textbf{\hl{✔ সঠিক উত্তর: ঘ) ২৫}}

\begin{quote}
\textbf{ব্যাখ্যা:} \emph{আসুন মামুন বাড়িকে Base ধরে প্রশ্নানুযায়ী নিচে মতো করে
চিত্র আঁকি।}

\emph{এখানে, গন্তব্যস্থান হতে মামুনের বাড়ির দূরত্ব AB যা ত্রিভুজ ABC সমকোণী
ত্রিভুজের অতিভুজ।}

\emph{অতএব, পিথাগোরাসের সূত্রানুসারে:}

\emph{(AB)\^{}2 = (AC)\^{}2 + (BC)\^{}2}

\emph{⇒ (AB)\^{}2 = (20)\^{}2 + (15)\^{}2}

\emph{⇒ (AB)\^{}2 = 400 + 225 = 625}

\emph{⇒ AB = √(625) = 25 মাইল}
\end{quote}

\textbf{একটি সমকোণী ত্রিভুজের অতিভুজের দৈর্ঘ্য ১৩ মিটার হলে, ত্রিভুজটির ক্ষেত্রফল
কত?}

\textbf{\hl{ক) ৩০ বর্গমিটার ✔}}

খ) ৪২ বর্গমিটার

গ) ৫৬ বর্গমিটার

ঘ) ৮৪ বর্গমিটার

\textbf{\hl{✔ সঠিক উত্তর: ক) ৩০ বর্গমিটার}}

\begin{quote}
\textbf{ব্যাখ্যা:} \emph{নিচের চিত্রের মতো সমকোণী ত্রিভুজের অতিভুজ ১৩ হলে অপর
দুই বাহু হবে ৫ এবং ১২ মিটার।}

\emph{কারণ, পিথাগোরাসের সূত্রানুসারে}

\emph{(অতিভুজ)\^{}2 = (ভূমি)\^{}2 + (লম্ব)\^{}2}

\emph{এখন, ত্রিভুজটির ক্ষেত্রফল}

\emph{= (1/2) × সমকোণ সংলগ্ন বাহুদ্বয়ের গুণফল}

\emph{= (1/2) × 12 × 5 = 30 বর্গমিটার।}
\end{quote}

\textbf{ত্রিভুজের বাহুগুলোর কোন অনুপাত দিয়ে একটি সমকোণী ত্রিভুজ অঙ্কন করা যাবে?}

ক) ৬ : ৫ : ৪

\textbf{\hl{খ) ৩ : ৪ : ৫ ✔}}

গ) ২২ : ৮ : ৪

ঘ) ৬ : ৪ : ৩

\textbf{\hl{✔ সঠিক উত্তর: খ) ৩ : ৪ : ৫}}

\begin{quote}
\textbf{ব্যাখ্যা:} \emph{ত্রিভুজের যে কোনো দুই বাহুর সমষ্টি উহার তৃতীয় বাহু
অপেক্ষা বৃহত্তর।}

\emph{আবার সমকোণী ত্রিভুজের ক্ষেত্রে দুই বাহুর বর্গের সমষ্টি তৃতীয় বাহুর বর্গের
সমান।}

\emph{এই শর্তানুযায়ী ৩\^{}2 + ৪\^{}2 = ৫\^{}2 {[}পিথাগোরাসের উপপাদ্য অনুযায়ী
এটি সমকোণী ত্রিভুজ হবে{]}। কারেক্ট অপশন খ।}
\end{quote}

\textbf{১৭ সে.মি., ১৫ সে.মি., ৮ সে.মি. বাহু বিশিষ্ট ত্রিভুজটি হবে-}

ক) সমবাহু

খ) সমদ্বিবাহু

\textbf{\hl{গ) সমকোণী ✔}}

ঘ) স্থূলকোণী

\textbf{\hl{✔ সঠিক উত্তর: গ) সমকোণী}}

\begin{quote}
\textbf{ব্যাখ্যা:} \emph{পিথাগোরাসের উপপাদ্য অনুযায়ী, সমকোণী ত্রিভুজের ক্ষেত্রে
দুই বাহুর বর্গের সমষ্টি তৃতীয় বাহুর বর্গের সমান।}

\emph{এই শর্তানুযায়ী - (লম্ব)\^{}2 + (ভূমি)\^{}2 = (অতিভুজ)\^{}2}

\emph{15\^{}2 + 8\^{}2 = 225 + 64 = 289 = 17\^{}2}

\emph{এটি সমকোণী ত্রিভুজ হবে।}
\end{quote}

\textbf{সমকোণী ত্রিভুজের বাহুগুলোর অনুপাত কোনটি?}

ক) ৬ : ৪ : ৩

খ) ৬ : ৫ : ৩

গ) ১২ : ৮ : ৪

\textbf{\hl{ঘ) ১৩ : ১২ : ৫ ✔}}

\textbf{\hl{✔ সঠিক উত্তর: ঘ) ১৩ : ১২ : ৫}}

\begin{quote}
\textbf{ব্যাখ্যা:} \emph{এখানে, (12)\^{}2 + (5)\^{}2 = 144 + 25 = 169 =
(13)\^{}2}

\emph{∴ সমকোণী ত্রিভুজের বাহুগুলোর অনুপাত = 13 : 12 : 5}
\end{quote}

\textbf{একটি সমকোণী ত্রিভুজের ভূমির দৈর্ঘ্য লম্বের দৈর্ঘ্য অপেক্ষা ১ মিটার কম এবং
লম্ব অপেক্ষা অতিভুজের দৈর্ঘ্য ১ মিটার বেশি হলে ত্রিভুজটির অতিভুজের দৈর্ঘ্য কত?}

ক) ৩

\textbf{\hl{খ) ৫ ✔}}

গ) ৬

ঘ) ৮

\textbf{\hl{✔ সঠিক উত্তর: খ) ৫}}

\begin{quote}
\textbf{ব্যাখ্যা:} \emph{ধরি, লম্বের দৈর্ঘ্য x}

\emph{তাহলে ভূমি = x − 1}

\emph{এবং অতিভুজ = x + 1}

\emph{অতিভুজ\^{}2 = ভূমি\^{}2 + লম্ব\^{}2}

\emph{⇒ (x + 1)\^{}2 = (x − 1)\^{}2 + x\^{}2}

\emph{⇒ x\^{}2 + 2x + 1 = x\^{}2 − 2x + 1 + x\^{}2}

\emph{⇒ x\^{}2 − 2x + 1 + x\^{}2 − x\^{}2 − 2x − 1 = 0}

\emph{⇒ (x\^{}2 − 4x) = 0}

\emph{∴ x = 4}

\emph{∴ অতিভুজের দৈর্ঘ্য = 4 + 1 = 5}
\end{quote}

\textbf{একটি সমকোণী ত্রিভুজের অতিভুজ ২৫ মিটার। অপর বাহুদ্বয়ের একটি অপরটির 3/4
অংশ হলে, অপর বাহুর দৈর্ঘ্যের অনুপাত---}

\textbf{\hl{ক) 3:4 ✔}}

খ) 1:2

গ) 3:5

ঘ) 2:1

\textbf{\hl{✔ সঠিক উত্তর: ক) 3:4}}

\begin{quote}
\textbf{ব্যাখ্যা:} \emph{ধরি, ভূমির দৈর্ঘ্য 4x মিটার}

\emph{∴ লম্বের দৈর্ঘ্য = 4x এর (3/4) = 3x মিটার}

\emph{এখন, (লম্ব)\^{}2 + (ভূমি)\^{}2 = (অতিভুজ)\^{}2}

\emph{⇒ (3x)\^{}2 + (4x)\^{}2 = (25)\^{}2}

\emph{⇒ 9x\^{}2 + 16x\^{}2 = 625}

\emph{⇒ 25x\^{}2 = 625}

\emph{⇒ x\^{}2 = 25 ∴ x = 5}

\emph{∴ ভূমির দৈর্ঘ্য (4 × 5) মিটার = 20 মিটার}

\emph{∴ লম্বের দৈর্ঘ্য (3 × 5) মিটার = 15 মিটার}

\emph{∴ অপর বাহুদ্বয়ের দৈর্ঘ্যের অনুপাত = 15 : 20 = 3 : 4}
\end{quote}

\textbf{একটি সমকোণী ত্রিভুজের অতিভুজের দৈর্ঘ্য ৫ সে.মি. হলে অপর দুই বাহুর দৈর্ঘ্য
কত?}

\textbf{\hl{ক) ৪ সেমি ও ৩ সেমি ✔}}

খ) ২ সেমি ও ৫ সেমি

গ) ৩ সেমি ও ২ সেমি

ঘ) ৪ সেমি ও ৫ সেমি

\textbf{\hl{✔ সঠিক উত্তর: ক) ৪ সেমি ও ৩ সেমি}}

\begin{quote}
\textbf{ব্যাখ্যা:} \emph{(অতিভুজ)\^{}2 = (5)\^{}2 = 25 = 16 + 9 = (4)\^{}2
+ (3)\^{}2 = (লম্ব)\^{}2 + (ভূমি)\^{}2}

\emph{∴ অপর দুই বাহুর দৈর্ঘ্য ৪ সেমি ও ৩ সেমি।}
\end{quote}

\textbf{যদি একটি সমকোণী ত্রিভুজের সমকোণ সংলগ্ন বাহুদ্বয়ের দৈর্ঘ্য যথাক্রমে ৫
সেমি ও ১২ সেমি হয়, তবে ত্রিভুজের অতিভুজের দৈর্ঘ্য কত?}

ক) ৯ সেমি

\textbf{\hl{খ) ১৩ সেমি ✔}}

গ) ১২ সেমি

ঘ) ১০ সেমি

\textbf{\hl{✔ সঠিক উত্তর: খ) ১৩ সেমি}}

\begin{quote}
\textbf{ব্যাখ্যা:} \emph{সমকোণী ত্রিভুজের সমকোণ সংলগ্ন বাহুদ্বয়ের দৈর্ঘ্য যথাক্রমে
5 সেমি ও 12 সেমি হয় তবে অতিভুজের দৈর্ঘ্য = √(5\^{}2+12\^{}2) = √(25+144) =
√(169) = 13 সেমি}
\end{quote}

\textbf{ABC সমকোণী ত্রিভুজের অতিভুজ BC = 5, লম্ব AB = 3 হলে, ভূমি AC এর মান
কত?}

ক) 6

\textbf{\hl{খ) 4 ✔}}

গ) 2

ঘ) 5

\textbf{\hl{✔ সঠিক উত্তর: খ) 4}}

\begin{quote}
\textbf{ব্যাখ্যা:} \emph{পিথাগোরাসের সূত্রের সাহায্যে, BC\^{}2 = AB\^{}2 +
AC\^{}2 বা, 5\^{}2 = 3\^{}2 + AC\^{}2}

\emph{বা, AC\^{}2 = 25 − 9 বা, AC\^{}2 = 16}

\emph{∴ AC = 4}
\end{quote}

\textbf{ABC সমকোণী ত্রিভুজের B কোণ সমকোণ এবং AB = a, BC = 2a হলে, নিচের
কোনটি সত্য?}

ক) AC\^{}2 = 3a\^{}2

\textbf{\hl{খ) AC = a √(5) ✔}}

গ) AC\^{}2 = 5a

ঘ) AC\^{}2 = 7a\^{}2

\textbf{\hl{✔ সঠিক উত্তর: খ) AC = a √(5)}}

\begin{quote}
\textbf{ব্যাখ্যা:} \emph{আসুন প্রশ্নটি প্রথমে ভালো করে পড়ে নিচের মতো করে
চিত্রটি আঁকি:}

\emph{পিথাগোরাসের উপপাদ্য অনুযায়ী, AC\^{}2 = AB\^{}2 + BC\^{}2 = a\^{}2 +
(2a)\^{}2 = a\^{}2 + 4a\^{}2 = 5a\^{}2}

\emph{∴ AC = √(5a\^{}2) = a √(5)}
\end{quote}

\textbf{কোন একটি ত্রিভুজের দুটি বাহু যথাক্রমে ৪ সে.মি. ও ৫ সে.মি.। তৃতীয় বাহুর
দৈর্ঘ্য কত সে.মি. হলে ত্রিভুজটি একটি সমকোণী ত্রিভুজ হবে?}

\textbf{\hl{ক) ৩ সে.মি. ✔}}

খ) ৮ সে.মি.

গ) ২ সে.মি.

ঘ) ৫ সে.মি.

\textbf{\hl{✔ সঠিক উত্তর: ক) ৩ সে.মি.}}

\begin{quote}
\textbf{ব্যাখ্যা:} \emph{সমকোণী ত্রিভুজের ক্ষেত্রে আমরা জানি, (লম্ব)\^{}2 +
(ভূমি)\^{}2 = (অতিভুজ)\^{}2}

\emph{মনে করি, লম্ব = 4 সে.মি. এবং অতিভুজ = 5 সে.মি.}

\emph{∴ (4)\^{}2 + (ভূমি)\^{}2 = (5)\^{}2}

\emph{বা, (ভূমি)\^{}2 = 25 − 16}

\emph{বা, (ভূমি)\^{}2 = 9}

\emph{∴ ভূমি = √(9) = 3 সে.মি.}
\end{quote}

\textbf{সমকোণী ত্রিভুজের সমকোণ সংলগ্ন বাহুদ্বয় যথাক্রমে ৩ ও ৪ সে.মি: হলে, এর
অতিভুজের মান কত?}

ক) ৪ সে.মি:

\textbf{\hl{খ) ৫ সে.মি: ✔}}

গ) ৭ সে.মি:

ঘ) ৮ সে.মি:

\textbf{\hl{✔ সঠিক উত্তর: খ) ৫ সে.মি:}}

\begin{quote}
\textbf{ব্যাখ্যা:} \emph{যেহেতু ত্রিভুজটি সমকোণী,}

\emph{(অতিভুজ)\^{}2 = 3\^{}2 + 4\^{}2 = 9 + 16 = 25 বর্গসেমি}

\emph{∴ অতিভুজ = √(25) = 5 সে.মি.}
\end{quote}

\textbf{একটি সমকোণী ত্রিভুজের ভূমি ১২ মিটার এবং উচ্চতা ৫ মিটার হলে, এর অতিভুজ
ভূমি অপেক্ষা কত মিটার বেশি?}

\textbf{\hl{ক) ১ ✔}}

খ) ৮

গ) ১৩

ঘ) ১৪

\textbf{\hl{✔ সঠিক উত্তর: ক) ১}}

\begin{quote}
\textbf{ব্যাখ্যা:} \emph{পিথাগোরাসের সূত্রানুযায়ী একটি সমকোণী ত্রিভুজের বাহুগুলোর
অনুপাত নিম্নরূপ হতে পারে:}

\emph{i) ৩ : ৪ : ৫ ii) ৬ : ৮ : ১০ iii) ৫ : ১২ : ১৩}

\emph{অর্থাৎ ভূমির দৈর্ঘ্য ১২ মিটার ও উচ্চতা ৫ মিটার হলে অতিভুজ হবে ১৩ মিটার।}

\emph{∴ ভূমি অপেক্ষা অতিভুজ = 13 − 12 = 1 মিটার বেশি।}
\end{quote}

\textbf{কোনো ত্রিভুজের বাহুগুলোর অনুপাত নিচের কোনটি হলে ত্রিভুজটি একটি সমকোণী
ত্রিভুজ হবে?}

ক) ৬ : ৫ : ৪

খ) ১২ : ৮ : ৪

\textbf{\hl{গ) ৯ : ১২ : ১৫ ✔}}

ঘ) ৬ : ৪ : ৩

\textbf{\hl{✔ সঠিক উত্তর: গ) ৯ : ১২ : ১৫}}

\begin{quote}
\textbf{ব্যাখ্যা:} \emph{কোনো ত্রিভুজের এক বাহুর উপর অঙ্কিত ক্ষেত্র অপর দুই বাহুর
উপর অঙ্কিত ক্ষেত্রদ্বয়ের যোগফলের সমান হলে ত্রিভুজটি সমকোণী।}

\emph{গ. 9\^{}2+12\^{}2=15\^{}2}
\end{quote}

\textbf{নিচের ত্রিভুজের বাহুর দৈর্ঘ্য দেয়া আছে। কোনটি সমকোণী ত্রিভুজ হবে।}

ক) 2, 3, 4

খ) 4, 5, 6

গ) 3, 4, 7

\textbf{\hl{ঘ) 3, 4, 5 ✔}}

\textbf{\hl{✔ সঠিক উত্তর: ঘ) 3, 4, 5}}

\begin{quote}
\textbf{ব্যাখ্যা:} \emph{এখানে, 3\^{}2+4\^{}2=9+16=25=5\^{}2}

\emph{∴ 3, 4 ও 5 একক দৈর্ঘ্য বাহু দ্বারা সমকোণী ত্রিভুজ গঠন করা হয়।}
\end{quote}

\textbf{সমকোণী ত্রিভুজের বাহুগুলোর অনুপাত কত?}

ক) ৬ : ৪ : ৩

খ) ৬ : ৫ : ৪

গ) ১২ : ৮ : ৪

\textbf{\hl{ঘ) ১৩ : ১২ : ৫ ✔}}

\textbf{\hl{✔ সঠিক উত্তর: ঘ) ১৩ : ১২ : ৫}}

\begin{quote}
\textbf{ব্যাখ্যা:} \emph{সমাধান:}

\emph{ঘ. 5\^{}2+12\^{}2=13\^{}2 (সমকোণী ত্রিভুজ)}
\end{quote}

\textbf{একটি সমকোণী ত্রিভুজের ভূমি 5 মিটার ও অতিভুজ 13 মিটার হলে ত্রিভুজের
পরিসীমা হবে---}

ক) 25 মিটার

খ) 35 মিটার

\textbf{\hl{গ) 30 মিটার ✔}}

ঘ) 36 মিটার

\textbf{\hl{✔ সঠিক উত্তর: গ) 30 মিটার}}

\begin{quote}
\textbf{ব্যাখ্যা:} \emph{লম্ব = √(13\^{}2−5\^{}2) = √(169−25) = √(144) = 12
মিটার}

\emph{ত্রিভুজটির পরিসীমা = 13 + 12 + 5 = 30 মিটার}
\end{quote}

\textbf{একটি সমকোণী ত্রিভুজের ক্ষেত্রফল 6 বর্গমিটার। ত্রিভুজের অতিভুজ 5 মিটার
হলে ত্রিভুজটির পরিসীমা কত?}

ক) 10 meters

\textbf{\hl{খ) 12 meters ✔}}

গ) 15 meters

ঘ) 10 meters

\textbf{\hl{✔ সঠিক উত্তর: খ) 12 meters}}

\begin{quote}
\textbf{ব্যাখ্যা:} \emph{ধরি, ত্রিভুজের ভূমি x মিটার এবং উচ্চতা y মিটার।}

\emph{ক্ষেত্রফল = (1/2) xy = 6 ⇒ xy = 12}

\emph{আমরা জানি, x\^{}2+y\^{}2=5\^{}2=25}

\emph{(x + y)\^{}2 = x\^{}2 + y\^{}2 + 2xy = 25 + 2(12) = 25 + 24 = 49 ⇒
x + y = 7}

\emph{∴ পরিসীমা = ভূমি + উচ্চতা + অতিভুজ = (x + y) + 5 = 7 + 5 = 12 meters}
\end{quote}

\textbf{ত্রিভুজের তিন বাহুর দৈর্ঘ্য যথাক্রমে ৬, ৮ এবং ১০ মিটার হলে বৃহত্তম ও
ক্ষুদ্রতম বাহুর মধ্যবিন্দু দুটির দূরত্ব কত মিটার?}

\textbf{\hl{ক) ৪ ✔}}

খ) ৫

গ) ৬

ঘ) ৭

ঙ) কোনোটিই নয়

\textbf{\hl{✔ সঠিক উত্তর: ক) ৪}}

\begin{quote}
\textbf{ব্যাখ্যা:} \emph{কোন ত্রিভুজের তিন বাহুর দৈর্ঘ্য যদি ৬, ৮ এবং ১০ অনুপাতে
হয় তাহলে এটি হবে একটি সমকোণী ত্রিভুজ যেটিতে অতিভুজ হবে ১০।}

\emph{তাহলে ক্ষুদ্রতম এবং বৃহত্তম বাহু দুটি যথাক্রমে ৩ এবং ৫ হবে। আর মধ্যবিন্দু যোগ
করলে উৎপন্ন বাহুর দৈর্ঘ্য পিথাগোরাসের ট্রিপলেট অনুসারে ৪ হবে।}
\end{quote}

\textbf{City B is 5 miles east of city A. City C is 10 miles southeast
of city B. Which of the following is the closer to the distance from
City A to City C?}

ক) 11 miles

খ) 12 miles

গ) 13 miles

\textbf{\hl{ঘ) 14 miles ✔}}

\textbf{\hl{✔ সঠিক উত্তর: ঘ) 14 miles}}

\begin{quote}
\textbf{ব্যাখ্যা:} \emph{প্রশ্নানুযায়ী, City A থেকে City B এর দূরত্ব ৫ মাইল
পূর্বে। সেখান থেকে City C ১০ মাইল দক্ষিণ-পূর্ব (southeast বা 45 ডিগ্রি কোণে)
অবস্থিত।}

\emph{ভেক্টরের উপাংশ বা কোসাইন সূত্র ব্যবহার করে দূরত্ব হিসাব করলে এর মান প্রায়
১৪ মাইল হয়।}
\end{quote}

\textbf{ত্রিভুজের বৃহত্তর বাহুসংলগ্ন কোণদ্বয়}

ক) স্থূলকোণ

\textbf{\hl{খ) সূক্ষ্মকোণ ✔}}

গ) সমকোণ

ঘ) সরলকোণ

\textbf{\hl{✔ সঠিক উত্তর: খ) সূক্ষ্মকোণ}}

\begin{quote}
\textbf{ব্যাখ্যা:} \emph{ত্রিভুজের বৃহত্তর বাহুসংলগ্ন কোণদ্বয় সর্বদা সূক্ষ্মকোণ হয়ে
থাকে।}
\end{quote}

\textbf{ΔABC এ AD, ∠A-এর সমদ্বিখণ্ডক এবং ∠ADB সূক্ষ্মকোণ হলে---}

ক) AD \textgreater{} AC

খ) AB \textgreater{} AC

\textbf{\hl{গ) AB \textless{} AC ✔}}

ঘ) BD \textless{} CD

\textbf{\hl{✔ সঠিক উত্তর: গ) AB \textless{} AC}}

\begin{quote}
\textbf{ব্যাখ্যা:} \emph{ΔABC এ ∠BAD = ∠CAD}

\emph{∴ ∠ABC \textgreater{} ∠ACB (∟ ∠ADB \textless{} ∠ADC)}

\emph{∴ AC \textgreater{} AB}

\emph{অর্থাৎ AB \textless{} AC}
\end{quote}

\textbf{স্থূলকোণী ত্রিভুজের ক্ষেত্রে পরিকেন্দ্র ত্রিভুজের কোথায় অবস্থিত?}

ক) অভ্যন্তরে

খ) বাহুর উপরে

\textbf{\hl{গ) বহির্ভাগে ✔}}

ঘ) ভরকেন্দ্রে

\textbf{\hl{✔ সঠিক উত্তর: গ) বহির্ভাগে}}

\begin{quote}
\textbf{ব্যাখ্যা:} \emph{ত্রিভুজের পরিকেন্দ্র : ত্রিভুজের বাহুত্রয়ের লম্বদ্বিখণ্ডকত্রয়
সমবিন্দু। এই বিন্দুকে পরিকেন্দ্র বলে। স্থূলকোণী ত্রিভুজের ক্ষেত্রে পরিকেন্দ্র ত্রিভুজের
বহির্ভাগে অবস্থিত।}
\end{quote}

\textbf{স্থূলকোণী ত্রিভুজের স্থূলকোণের সংখ্যা:}

\textbf{\hl{ক) ১টি ✔}}

খ) ২টি

গ) ৩টি

ঘ) কোনোটিই নয়

\textbf{\hl{✔ সঠিক উত্তর: ক) ১টি}}

\begin{quote}
\textbf{ব্যাখ্যা:} \emph{যে ত্রিভুজের একটি কোণ স্থূলকোণ, সেই ত্রিভুজকে স্থূলকোণী
ত্রিভুজ বলা হয়। ত্রিভুজের তিন কোণের সমষ্টি ১৮০° হওয়ায় একটির বেশি স্থূলকোণ থাকা
সম্ভব নয়।}
\end{quote}

\textbf{150° কোণটি হলো---}

ক) সূক্ষ্মকোণী

\textbf{\hl{খ) স্থূলকোণ ✔}}

গ) বিপরীত কোণ

ঘ) সমকোণ

\textbf{\hl{✔ সঠিক উত্তর: খ) স্থূলকোণ}}

\begin{quote}
\textbf{ব্যাখ্যা:} \emph{90° বা সমকোণ অপেক্ষা বড় কিন্তু 180° অপেক্ষা ছোট কোণকে
স্থূল কোণ বলে।}
\end{quote}

\textbf{ΔABC-এর BC বাহুকে D পর্যন্ত বর্ধিত করা হলে, ∠ACD এর সমান হবে-}

ক) ∠A + ∠B - ∠C

খ) ∠A - ∠B - ∠C

\textbf{\hl{গ) ∠A + ∠B ✔}}

ঘ) ∠A + ∠C

\textbf{\hl{✔ সঠিক উত্তর: গ) ∠A + ∠B}}

\begin{quote}
\textbf{ব্যাখ্যা:} \emph{ত্রিভুজের কোনো এক বাহুকে বর্ধিত করা হলে যে বহিঃস্থ কোণ
উৎপন্ন হয় তা ত্রিভুজের বিপরীত অন্তঃস্থ কোণদ্বয়ের সমষ্টির সমান।}

\emph{∴ ∠ACD = ∠ABC + ∠BAC = ∠B + ∠A}
\end{quote}

\textbf{একটি ত্রিভুজের তিন কোণ অপর একটি ত্রিভুজের তিন কোণের সমান হলে,
ত্রিভুজদ্বয় কিরূপ হবে?}

ক) সমান

খ) সর্বসম

গ) অসমান

\textbf{\hl{ঘ) সদৃশকোণী ✔}}

\textbf{\hl{✔ সঠিক উত্তর: ঘ) সদৃশকোণী}}

\begin{quote}
\textbf{ব্যাখ্যা:} \emph{একটি ত্রিভুজের তিন কোণ অপর একটি ত্রিভুজের তিন কোণের
সমান হলে ত্রিভুজদ্বয় সদৃশকোণী হয়।}
\end{quote}

\textbf{যদি দুটি ত্রিভুজের দুটি বাহু ও তাদের অন্তর্ভুক্ত কোণ সমান হয় তবে ত্রিভুজ
দুটি - হবে।}

ক) সমানুপাতিক

খ) সমরেখ

\textbf{\hl{গ) সর্বসম ✔}}

ঘ) সদৃশ

\textbf{\hl{✔ সঠিক উত্তর: গ) সর্বসম}}

\begin{quote}
\textbf{ব্যাখ্যা:} \emph{যদি দুটি ত্রিভুজের দুটি বাহু ও তাদের অন্তর্ভুক্ত কোণ অপর
একটি ত্রিভুজের যথাক্রমে দুটি বাহু ও অন্তর্ভুক্ত কোণের সমান হয়, তবে ত্রিভুজ দুটি
সর্বসম হয়।}
\end{quote}

\textbf{তিন কোণ দেওয়া থাকলে যে সকল ত্রিভুজ আঁকা যায় সেগুলোকে কি ত্রিভুজ বলে?}

ক) সমান ত্রিভুজ

খ) সর্বসম ত্রিভুজ

\textbf{\hl{গ) সদৃশ ত্রিভুজ ✔}}

ঘ) সমানুপাতিক ত্রিভুজ

\textbf{\hl{✔ সঠিক উত্তর: গ) সদৃশ ত্রিভুজ}}

\begin{quote}
\textbf{ব্যাখ্যা:} \emph{তিন কোণ দেওয়া থাকলে অসীম সংখ্যক ত্রিভুজ আঁকা সম্ভব এবং
তাদের আকার ভিন্ন হতে পারলেও আকৃতি একই থাকে, তাই সেগুলোকে সদৃশ ত্রিভুজ বলে।}
\end{quote}

\textbf{দুটি ত্রিভুজের মধ্যে কোন উপাদানগুলো সমান হওয়া সত্ত্বেও ত্রিভুজ দুটি সর্বসম
নাও হতেপারে?}

ক) দুইটি বাহু ও এক কোণ (Two sides and one angle)

খ) তিনটি বাহু (Three sides)

গ) এক বাহু এবং দুই কোণ (One side and two angles)

\textbf{\hl{ঘ) তিন কোণ (Three angles) ✔}}

\textbf{\hl{✔ সঠিক উত্তর: ঘ) তিন কোণ (Three angles)}}

\begin{quote}
\textbf{ব্যাখ্যা:} \emph{ত্রিভুজের তিন কোণ সমান হলে ত্রিভুজদ্বয় সদৃশকোণী হয়,
কিন্তু তাদের বাহুগুলোর দৈর্ঘ্য ভিন্ন হতে পারায় তারা সর্বসম নাও হতে পারে।}
\end{quote}

\textbf{≅ এই প্রতীকের অর্থ হলো/The symbol ≅ is used to denote}

\textbf{\hl{ক) আকার ও আকৃতি সমান (same size and shape) ✔}}

খ) আকার ও পরিমাণ সমান (same size and measure)

গ) আকার ও আকৃতি সমান নয় (size and shape which are not same)

ঘ) আকার, আকৃতি ও পরিমাণ সমান (same size, shape and measure)

\textbf{\hl{✔ সঠিক উত্তর: ক) আকার ও আকৃতি সমান (same size and shape)}}

\begin{quote}
\textbf{ব্যাখ্যা:} \emph{গণিতে ≅ প্রতীকটি সর্বসম (Congruent) বোঝাতে ব্যবহার
করা হয়, যার অর্থ দুটি জ্যামিতিক চিত্রের আকার ও আকৃতি সমান।}
\end{quote}

\textbf{ABCD রম্বসের কর্ণদ্বয় পরস্পরকে O বিন্দুতে ছেদ করে। ΔBOC হবে-}

\textbf{\hl{ক) বিষমবাহু ত্রিভুজ ✔}}

খ) সমবাহু ত্রিভুজ

গ) সূক্ষ্মকোণী ত্রিভুজ

ঘ) স্থূলকোণী ত্রিভুজ

\textbf{\hl{✔ সঠিক উত্তর: ক) বিষমবাহু ত্রিভুজ}}

\begin{quote}
\textbf{ব্যাখ্যা:} \emph{রম্বসের কর্ণদ্বয় পরস্পরকে সমকোণে সমদ্বিখণ্ডিত করে। তাই
ছেদবিন্দুতে উৎপন্ন প্রতিটি কোণই সমকোণ। সুতরাং ΔBOC হবে সমকোণী ত্রিভুজ। এটি একটি
বিষমবাহু ত্রিভুজও বটে।}
\end{quote}

\textbf{ABC এ ∠A = 40°, ∠B = 70°, হলে ABC কি ধরনের ত্রিভুজ?}

ক) সমকোণী

খ) স্থূলকোণী

\textbf{\hl{গ) সমদ্বিবাহু ✔}}

ঘ) সমবাহু

\textbf{\hl{✔ সঠিক উত্তর: গ) সমদ্বিবাহু}}

\begin{quote}
\textbf{ব্যাখ্যা:} \emph{ΔABC-এ}

\emph{∠A + ∠B + ∠C = 180°}

\emph{⇒ 40° + 70° + ∠C = 180°}

\emph{⇒ ∠C = 180° − 110°}

\emph{⇒ ∠C = 70° ∴ ΔABC একটি সমদ্বিবাহু ত্রিভুজ।}
\end{quote}

\textbf{ABC ত্রিভুজে ∠B = ∠C এবং BC বাহুর উপর D একটি বিন্দু। কোন তথ্যটি সঠিক?}

ক) AB \textgreater{} AC

খ) AB \textless{} AC

গ) AC \textless{} AD

\textbf{\hl{ঘ) AC \textgreater{} AD ✔}}

\textbf{\hl{✔ সঠিক উত্তর: ঘ) AC \textgreater{} AD}}

\begin{quote}
\textbf{ব্যাখ্যা:} \emph{AD-কে E পর্যন্ত বর্ধিত করি যেন AD = DE হয়।}

\emph{ΔADC ও ΔCDE এ}

\emph{∴ AD = DE, CD সাধারণ বাহু এবং ∠ADC = ∠CDE}

\emph{∴ ΔADC ≅ ΔCDE}

\emph{∴ AC = CE}

\emph{এখন, ΔACE হতে AC + CE \textgreater{} AE}

\emph{⇒ AC + AC \textgreater{} AD + DE}

\emph{⇒ 2AC \textgreater{} AD + AD}

\emph{⇒ 2AC \textgreater{} 2AD ∴ AC \textgreater{} AD}
\end{quote}

\textbf{ABC ত্রিভুজের AB=AC এবং ∠A=80° হলে ∠B = কত?}

ক) 40°

\textbf{\hl{খ) 50° ✔}}

গ) 60°

ঘ) 80°

\textbf{\hl{✔ সঠিক উত্তর: খ) 50°}}

\begin{quote}
\textbf{ব্যাখ্যা:} \emph{ABC ত্রিভুজের AB = AC}

\emph{∴ ∠C = ∠B}

\emph{∠A + ∠B + ∠C = 180°}

\emph{⇒ 80° + ∠B + ∠B = 180°}

\emph{⇒ 2∠B = 180° − 80° = 100°}

\emph{∴ ∠B = 50°}
\end{quote}

\textbf{BE ও CF, ΔABC-এর দুইটি মধ্যমা এবং BC = 18 সে.মি. হলে EF-এর মান
কত?}

\textbf{\hl{ক) 9 সে.মি. ✔}}

খ) 10 সে.মি.

গ) 12 সে.মি.

ঘ) 8 সে.মি.

\textbf{\hl{✔ সঠিক উত্তর: ক) 9 সে.মি.}}

\begin{quote}
\textbf{ব্যাখ্যা:} \emph{ত্রিভুজের যেকোনো দুই বাহুর মধ্যবিন্দুর সংযোজক সরলরেখা
তৃতীয় বাহুর সমান্তরাল ও অর্ধেক হয়।}

\emph{EF = (1/2) × BC = (1/2) × 18 = 9 সে.মি.}
\end{quote}

\textbf{∠ACB = কত ডিগ্রী?}

ক) 120°

খ) 40°

গ) 60°

\textbf{\hl{ঘ) 80° ✔}}

\textbf{\hl{✔ সঠিক উত্তর: ঘ) 80°}}

\begin{quote}
\textbf{ব্যাখ্যা:} \emph{∠ABC + ∠BAC = ∠CBD}

\emph{∠ACB + 40° = 120°}

\emph{∠ACB = 80°}
\end{quote}

\textbf{ABC ত্রিভুজে ∠B = 6x ডিগ্রী, ∠C = 5x ডিগ্রী, ∠A = y ডিগ্রী এবং 6∠A =
7∠B হলে, y এর মান হবে-}

ক) 90

খ) 80

\textbf{\hl{গ) 70 ✔}}

ঘ) 60

\textbf{\hl{✔ সঠিক উত্তর: গ) 70}}

\begin{quote}
\textbf{ব্যাখ্যা:} \emph{দেওয়া আছে, ∠B = 6x এবং ∠C = 5x}

\emph{প্রশ্নমতে, 6∠A = 7∠B}

\emph{⇒ ∠A = (7∠B/6)}

\emph{⇒ ∠A = (7×6x/6) ⇒ ∠A = 7x}

\emph{আমরা জানি, ত্রিভুজের তিন কোণের সমষ্টি = 180°}

\emph{∴ ∠A + ∠B + ∠C = 180°}

\emph{⇒ 7x + 6x + 5x = 180°}

\emph{⇒ 18x = 180°}

\emph{⇒ x = (180°/18) ⇒ x = 10°}

\emph{∴ y = ∠A = 7x = 7 × 10° = 70°}
\end{quote}

\textbf{ABC ত্রিভুজের BC বাহুকে D পর্যন্ত বাড়ানো হলো। ∠ACD = 105° হলে ∠BAC +
∠ABC = কত?}

ক) 90°

খ) 75°

গ) 180°

\textbf{\hl{ঘ) 105° ✔}}

\textbf{\hl{✔ সঠিক উত্তর: ঘ) 105°}}

\begin{quote}
\textbf{ব্যাখ্যা:} \emph{ত্রিভুজের কোন একবাহুকে বর্ধিত করলে যে বহিঃস্থ কোণ
উৎপন্ন হয়, তা বিপরীত অন্তঃস্থ কোণদ্বয়ের সমষ্টির সমান।}

\emph{ABC ত্রিভুজের BC বাহুকে D পর্যন্ত বাড়ানো হলে ∠ACD বহিঃস্থ কোণ উৎপন্ন হয়।}

\emph{∴ ∠ACD = ∠BAC + ∠ABC = 105°}
\end{quote}

\textbf{ΔABC এর E ও F যথাক্রমে AB ও AC এর মধ্যবিন্দু হলে EF = কত?}

ক) (1/2)(AB+AC)

খ) (1/2)AB

গ) (1/2)AC

\textbf{\hl{ঘ) (1/2)BC ✔}}

\textbf{\hl{✔ সঠিক উত্তর: ঘ) (1/2)BC}}

\begin{quote}
\textbf{ব্যাখ্যা:} \emph{ত্রিভুজের যে কোন দুই বাহুর মধ্য বিন্দুর সংযোগ রেখাংশ
তৃতীয় বাহুর সমান্তরাল হয় এবং দৈর্ঘ্য তার অর্ধেক হয়।}

\emph{∴ EF = (1/2)BC}
\end{quote}

\textbf{ΔABC-এর BC বাহুকে D পর্যন্ত বাড়ানো হলো। ∠A = 60° এবং ∠B = 90° হলে,
∠ACD = কত?}

ক) 90°

\textbf{\hl{খ) 150° ✔}}

গ) 160°

ঘ) 170°

\textbf{\hl{✔ সঠিক উত্তর: খ) 150°}}

\begin{quote}
\textbf{ব্যাখ্যা:} \emph{বহিঃস্থ ∠ACD = বিপরীত অন্তঃস্থ (∠BAC + ∠ABC)}

\emph{∠ACD = 60° + 90° = 150°}
\end{quote}

\textbf{ABC একটি ত্রিভুজের BC বাহুকে D পর্যন্ত বর্ধিত করা হলে এবং এ ত্রিভুজটির
∠A = 40°, ∠B = 60°, ∠C = 80° হলে, ∠ACD- এর পরিমাণ হবে-}

\textbf{\hl{ক) 100° ✔}}

খ) 120°

গ) 60°

ঘ) 140°

\textbf{\hl{✔ সঠিক উত্তর: ক) 100°}}

\begin{quote}
\textbf{ব্যাখ্যা:} \emph{আমরা জানি, ত্রিভুজের একটি বহিঃস্থ কোণের পরিমাণ
অন্তঃস্থ বিপরীত দুই কোণের সমষ্টির সমান।}

\emph{∠ACD = ∠CAB + ∠ABC = 40° + 60° = 100°}
\end{quote}

\textbf{The length of sides AB and AC in the triangle below are equal.
∠B = 70 What is the measure of angle A?}

ক) 70°

খ) 55°

\textbf{\hl{গ) 40° ✔}}

ঘ) 110°

ঙ) 20°

\textbf{\hl{✔ সঠিক উত্তর: গ) 40°}}

\begin{quote}
\textbf{ব্যাখ্যা:} \emph{∠B = ∠C = 70°, since AB = AC}

\emph{∴ ∠A = 180° − (70° + 70°) = 40°}
\end{quote}

\textbf{ΔABC এর ∠A = 45°, ∠B = 45° হলে ΔABC হলো (i) সমবাহু (ii) সমদ্বিবাহু
(iii) সমকোণী নিচের কোনটি সত্য?}

ক) (i), (ii)

খ) (i), (iii)

\textbf{\hl{গ) (ii), (iii) ✔}}

ঘ) (i), (ii), (iii)

\textbf{\hl{✔ সঠিক উত্তর: গ) (ii), (iii)}}

\begin{quote}
\textbf{ব্যাখ্যা:} \emph{ত্রিভুজের দুটি কোণ সমান যাদের সমষ্টি (45° + 45°) বা
90°। সুতরাং এর অপর কোণটি 180° − 90° = 90° অর্থাৎ ত্রিভুজটি সমকোণী সমদ্বিবাহু
ত্রিভুজ।}
\end{quote}

\section{13. বহুভুজ কাকে বলে ? সংজ্ঞা ,সূত্রাবলী ও কোণের পরিমাণ
নির্ণয়}\label{ux9acux9b9ux9adux99c-ux995ux995-ux9acux9b2-ux9b8ux99cux99e-ux9b8ux9a4ux9b0ux9acux9b2-ux993-ux995ux9a3ux9b0-ux9aaux9b0ux9aeux9a3-ux9a8ux9b0ux9a3ux9df}

বহুভুজ হলো একটি জ্যামিতিক আকৃতি । সাধারণত চারের অধিক বাহু দ্বারা সীমাবদ্ধ
ক্ষেত্রকে বহুভুজ হিসেবে চিহ্নিত করা হয় ।

বহুভুজ বিভিন্ন ধরনের হতে পারে যেমন- পঞ্চভুজ , ষড়ভুজ , সপ্তভুজ ইত্যাদি।

বহুভুজ নিয়ে জ্যামিতিক সমস্যাগুলো খুব বেশি জটিল মনে হলে ও আসলে তা কিন্তু একেবারেই
সহজ ।

মাত্র দুই থেকে তিনটা পদ্ধতি আয়ত্ত করতে পারলেই এ অধ্যায়ে যে কোন প্রতিযোগিতামূলক
পরীক্ষায় অন্য দশজনের চেয়ে এগিয়ে থাকা সম্ভব।

আজকের টিউটরিয়ালে আমরা সেই পদ্ধতিগুলো সহজভাবে উপস্থাপন করার চেষ্টা করেছি ।
আশা করছি উপকৃত হবেন । চলুন তাহলে শুরু করা যাক ।

\begin{itemize}
\item
  প্রাথমিক আলোচনা
\item
  সুষম বহুভুজের কোণের পরিমাপ নির্ণয়
\end{itemize}

\begin{itemize}
\item
  সংশ্লিষ্ট বিগত সালের প্রশ্ন ও সমাধান
\end{itemize}

\begin{itemize}
\item
  সুষম বহুভুজের বাহুর সংখ্যা নির্ণয়
\end{itemize}

\begin{itemize}
\item
  সংশ্লিষ্ট বিগত সালের প্রশ্ন ও সমাধান
\end{itemize}

\begin{itemize}
\item
  অন্তঃকোণ গুলোর সমষ্টি বের করা
\end{itemize}

\begin{itemize}
\item
  সংশ্লিষ্ট বিগত সালের প্রশ্ন ও সমাধান
\end{itemize}

\subsubsection{প্রাথমিক
আলোচনা}\label{ux9aaux9b0ux9a5ux9aeux995-ux986ux9b2ux99aux9a8-4}

এই অধ্যয়ের জন্য কয়েকটি সূত্র জানলেই যে কোন অংক খুব দ্রুত করা সম্ভব। কিন্তু
অধ্যায়টির সাথে আমরা অনেকেই খুব কম পরিচিত, এবং এখান থেকে খুব কম প্রশ্ন আসায়
সেই সামান্য সূত্রগুলোও আমাদের মনে থাকে না।

তাই আমরা শুধু একটি সহজ সূত্র দিয়ে এই অধ্যায়ের সকল প্রশ্নের সমাধান করা শিখব।
বুঝতে একটু সময় লাগলেও সময় দিয়ে পড়ুন।

\textbf{বহুভুজ:} ত্রিভুজ হচ্ছে সবচেয়ে কমবাহু বিশিষ্ট বহুভুজ, ত্রিভুজ থেকে শুরু করে
আরো বেশি সংখ্যক বাহুবিশিষ্ট যে কোন ক্ষেত্রকেই এক নামে বহুভুজ বলা হয়। যেমন:

পঞ্চভুজ , ষড়ভুজ ইত্যাদি।

\subsubsection{সুষম বহুভুজের কোণের পরিমাপ
নির্ণয়}\label{ux9b8ux9b7ux9ae-ux9acux9b9ux9adux99cux9b0-ux995ux9a3ux9b0-ux9aaux9b0ux9aeux9aa-ux9a8ux9b0ux9a3ux9df}

সুষম বহুভুজের কোণের পরিমাপ দুই ভাবে চাওয়া যেতে পারে । যথা: ১ . বহিঃস্থ কোণ
নির্ণয় ২. অন্তঃস্থ কোণ নির্ণয়

তাই আমরা দুইভাবেই বিষয়টা রপ্ত করার চেষ্টা করব।

\textbf{মনে রাখবেন, সুষম বহুভুজের বহিঃস্থ কোণগুলোর সমষ্টি = ৩৬০ ডিগ্রি}

একটি বহুভুজের বাহুর সংখ্যা যাহাই হোক না কেন তার বহিঃস্থ একটি কোণের পরিমাণ বের
করার জন্য নিম্নোক্ত সূত্রটি প্রয়োগ করতে হবে এবং এই একটি সূত্র দিয়ে বহুভুজের সব
প্রশ্ন সমাধান করা সম্ভব। তাই গুরুত্ব দিন।

\textbf{বহুভুজের বহিঃস্থ একটি কোণের পরিমাণ = ৩৬০ ÷ বাহুর সংখ্যা}

\textbf{১. একটি ষড়ভুজের প্রতিটি বহিঃস্থ কোণের মান কত হবে?}

\textbf{ব্যাখ্যা}: এখানে , ষড়ভুজ মানে তার ৬ টি বাহু । তাই ৩৬০ কে ৬ দিয়ে ভাগ
করলে বহিঃস্থ কোণের মান বের হবে ।

\textbf{\hl{উত্তর: ষড়ভুজের প্রতিটি বহিঃস্থ কোণের মান = ৩৬০ ÷ ৬ = ৬০ ডিগ্রি}}

\textbf{২. একটি অষ্টভুজের প্রতিটি বহিঃস্থ কোণের মান কত হবে?}

\textbf{ব্যাখ্যা}: এখানে , অষ্টভুজ মানে তার ৮ টি বাহু । তাই ৩৬০ কে ৮ দিয়ে ভাগ
করলে বহিঃস্থ কোণের মান বের হবে ।

\textbf{\hl{উত্তর: অষ্টভুজের প্রতিটি বহিঃস্থ কোণের মান = ৩৬০ ÷ ৮ = ৪৫ ডিগ্রি}}

\textbf{৩. একটি নয়ভুজের প্রতিটি অন্তঃস্থ কোণের মান কত?}

\textbf{ব্যাখ্যা}: একই সূত্র প্রয়োগ করতে হবে । নয়ভুজ মানে তার ৯ টি বাহু । তাই
৩৬০ কে ৯ দিয়ে ভাগ করলে বহিঃস্থ কোণের মান বের হবে । বহিঃস্থ কোণ = ৩৬০ ÷ ৯ =
৪০ ডিগ্রি । কিন্তু এটা তো হল বহিঃস্থ কোণ ।তাহলে অন্তঃস্থ কোণ কত হবে? অন্তঃস্থ
কোণ বের করতে হলে এই বহিঃস্থ কোণকে ১৮০ ডিগ্রি থেকে বিয়োগ করতে হবে।

সুতরাং নয়ভুজের প্রতিটি অন্তঃস্থ কোণের মান = ১৮০ − ৪০ = ১৪০ ডিগ্রি ।

এভাবেই আমরা সুষম বহুভুজের কোণের পরিমাপ নির্ণয় করব । এবার কিছু বিগত সালের প্রশ্ন
অনুশীলন করি ।

\paragraph{সংশ্লিষ্ট বিগত সালের প্রশ্ন ও
সমাধান}\label{ux9b8ux9b6ux9b2ux9b7ux99f-ux9acux997ux9a4-ux9b8ux9b2ux9b0-ux9aaux9b0ux9b6ux9a8-ux993-ux9b8ux9aeux9a7ux9a8}

\textbf{প্রশ্ন নম্বর-} ১. একটি সুষম পঞ্চভুজের একটি অন্তঃকোণের পরিমাণ হবে-
{[}প্রাথমিক সহকারী শিক্ষক-২০১২(করতোয়া){]}

\begin{itemize}
\item
  (ক) ৯০ ডিগ্রি
\item
  (খ) ১১০ ডিগ্রি
\item
  (গ) ১২০ ডিগ্রি
\item
  (ঘ) ১০৮ ডিগ্রি
\end{itemize}

\textbf{\hl{উত্তরঃ (ঘ) ১০৮ ডিগ্রি}}

\textbf{সমাধান:}

প্রথমে বহিঃস্থ কোণ বের করতে হবে । তারপর অন্তঃস্থ কোণ ।

বহিঃস্থ কোণ = ৩৬০ ÷ ৫ = ৭২ ডিগ্রি

সুতরাং অন্তঃস্থ কোণ = ১৮০ − ৭২ = ১০৮ ডিগ্রি

\textbf{প্রশ্ন নম্বর-} ২. আট বাহুবিশিষ্ট সুষম বহুভুজের একটি বহিঃস্থ কোণের পরিমাণ
হবে-\/-\/- {[}রেজি. বেসরকারি প্রাথমিক সহকারী শিক্ষক-২০১১(শাপলা){]}

\begin{itemize}
\item
  (ক) ৪৫ ডিগ্রি
\item
  (খ) ৬০ ডিগ্রি
\item
  (গ) ৮০ ডিগ্রি
\item
  (ঘ) ১০০ ডিগ্রি
\end{itemize}

\textbf{\hl{উত্তরঃ (ক) ৪৫ ডিগ্রি}}

\textbf{সমাধান:}

বহিঃস্থ কোণ = ৩৬০ ÷ ৮ = ৪৫ ডিগ্রি

\textbf{প্রশ্ন নম্বর-} ৩. একটি সুষম ষড়ভুজের একটি অন্তঃকোণের পরিমাণ হবে-\/- {[}
প্রাথমিক সহকারী শিক্ষক-২০১২(কর্ণফুলী){]}

\begin{itemize}
\item
  (ক) ৯০ ডিগ্রি
\item
  (খ) ১০৫ ডিগ্রি
\item
  (গ) ১২০ ডিগ্রি
\item
  (ঘ) ১৪৫ ডিগ্রি
\end{itemize}

\textbf{\hl{উত্তরঃ (গ) ১২০ ডিগ্রি}}

\textbf{সমাধান:}

প্রথমে বহিঃস্থ কোণ বের করতে হবে । তারপর অন্তঃকোণ ।

বহিঃস্থ কোণ = ৩৬০ ÷ ৬ = ৬০ ডিগ্রি

সুতরাং অন্তঃকোণ = ১৮০ − ৬০ = ১২০ ডিগ্রি

\textbf{প্রশ্ন নম্বর-} ৪. একটি সুষম বহুভুজের বাহুর সংখ্যা ১২ হলে প্রত্যেকটি বহিঃস্থ
কোণের পরিমাণ কত? {[} ২৮ তম বিসিএস{]}

\begin{itemize}
\item
  (ক) ১৫ ডিগ্রি
\item
  (খ) ২০ ডিগ্রি
\item
  (গ) ২৫ ডিগ্রি
\item
  (ঘ) ৩০ ডিগ্রি
\end{itemize}

\textbf{\hl{উত্তরঃ (ঘ) ৩০ ডিগ্রি}}

\textbf{সমাধান:}

বহিঃস্থ কোণ = ৩৬০ ÷ ১২ = ৩০ ডিগ্রি

\textbf{প্রশ্ন নম্বর-} ৫. সুষম ষড়ভুজের একটি বাহুকে বর্ধিত করলে উৎপন্ন বহিঃস্থ
কোণের পরিমাণ কত ডিগ্রি? {[} ১০ম বেসরকারি প্রভাষক নিবন্ধন ও প্রত্যয়ন
পরীক্ষা-২০১৪{]}

\begin{itemize}
\item
  (ক) ৭৫ ডিগ্রি
\item
  (খ) ৬০ ডিগ্রি
\item
  (গ) ৯০ ডিগ্রি
\item
  (ঘ) ১৮০ ডিগ্রি
\end{itemize}

\textbf{\hl{উত্তরঃ (খ) ৬০ ডিগ্রি}}

\textbf{সমাধান:}

বহিঃস্থ কোণ = ৩৬০ ÷ ৬ = ৬০ ডিগ্রি

\textbf{প্রশ্ন নম্বর-} ৬. একটি সুষম বহুভুজের বাহুর সংখ্যা ৮ হলে প্রত্যেকটি অন্তঃস্থ
কোণের পরিমাণ কত? {[}উপজেলা পল্লী উন্নয়ন কর্মকর্তা-২০১০{]}

\begin{itemize}
\item
  (ক) ৩০ ডিগ্রি
\item
  (খ) ৬০ ডিগ্রি
\item
  (গ) ৯০ ডিগ্রি
\item
  (ঘ) ১৩৫ ডিগ্রি
\end{itemize}

\textbf{\hl{উত্তরঃ (ঘ) ১৩৫ ডিগ্রি}}

\textbf{সমাধান:}

প্রথমে বহিঃস্থ কোণ বের করতে হবে । তারপর অন্তঃস্থ কোণ ।

বহিঃস্থ কোণ = ৩৬০ ÷ ৮ = ৪৫ ডিগ্রি

সুতরাং অন্তঃস্থ কোণ = ১৮০ − ৪৫ = ১৩৫ ডিগ্রি

\subsubsection{সুষম বহুভুজের বাহুর সংখ্যা
নির্ণয়}\label{ux9b8ux9b7ux9ae-ux9acux9b9ux9adux99cux9b0-ux9acux9b9ux9b0-ux9b8ux996ux9af-ux9a8ux9b0ux9a3ux9af}

বাহুর সংখ্যা নির্ণয়ের জন্য সুষম বহুভুজের অন্তঃস্থ অথবা বহিঃ স্থ কোণ দেয়া থাকবে ।
যদি বহিঃস্থ কোণ দেয়া থাকে তাহলে নিচের সূত্র প্রয়োগ করতে হবে ।

\textbf{সুষম বহুভুজের বাহুর সংখ্যা = ৩৬০ ÷ বহিঃস্থ কোণের পরিমাণ}

\textbf{উদাহরণ-১ : একটি বহুভুজের প্রতিটি বহিঃস্থ কোণের পরিমাণ ৬০ ডিগ্রি হলে,
বহুভুজটির বাহুর সংখ্যা কত?}

\textbf{ব্যাখ্যা}: এখানে, দেয়া আছে বহুভুজটির বহিঃস্থ কোণের পরিমাণ ৬০ ডিগ্রি ।
তাই সরাসরি , বহুভুজটির বাহুর সংখ্যা = ৩৬০ ÷ ৬০ = ৬ টি

\paragraph{সংশ্লিষ্ট বিগত সালের প্রশ্ন ও
সমাধান}\label{ux9b8ux9b6ux9b2ux9b7ux99f-ux9acux997ux9a4-ux9b8ux9b2ux9b0-ux9aaux9b0ux9b6ux9a8-ux993-ux9b8ux9aeux9a7ux9a8-1}

\textbf{প্রশ্ন নম্বর-} ১. একটি সুষম বহুভুজের একটি অন্তঃকোণের পরিমাণ ১৪০ ডিগ্রি
হলে বহুভুজটির বাহুর সংখ্যা হবে-\/- {[}রেজি. বেসরকারি প্রাথমিক সহকারী
শিক্ষক-২০১১(শিউলী){]}

\begin{itemize}
\item
  (ক) ৬
\item
  (খ) ৮
\item
  (গ) ৯
\item
  (ঘ) ১০
\end{itemize}

\textbf{\hl{উত্তরঃ (গ) ৯}}

\textbf{সমাধান:}

আগে বহিঃস্থ কোণ বের করতে হবে । তারপর বাহুর সংখ্যা ।

বহিঃস্থ কোণ = ১৮০ − ১৪০ = ৪০ ডিগ্রি

সুতরাং বহুভুজটির বাহুর সংখ্যা = ৩৬০ ÷ ৪০ = ৯ টি

\textbf{প্রশ্ন নম্বর-} ২. একটি সুষম বহুভুজের প্রত্যেকটি বহিঃস্থ কোণের পরিমাণ ২০
ডিগ্রি হলে বহুভুজটির বাহুর সংখ্যা কত? {[} প্রাথমিক বিদ্যালয় সহকারী শিক্ষক
২০০১{]}

\begin{itemize}
\item
  (ক) ৯
\item
  (খ) ১৮
\item
  (গ) ২৭
\item
  (ঘ) ৩৬
\end{itemize}

\textbf{\hl{উত্তরঃ (খ) ১৮}}

\textbf{সমাধান:}

বহুভুজটির বাহুর সংখ্যা = ৩৬০ ÷ ২০ = ১৮ টি

\subsubsection{অন্তঃকোণ গুলোর সমষ্টি বের
করা}\label{ux985ux9a8ux9a4ux995ux9a3-ux997ux9b2ux9b0-ux9b8ux9aeux9b7ux99f-ux9acux9b0-ux995ux9b0}

\textbf{কোণের মোট পরিমাপ বা সমষ্টি বের করতে বলা হলে :}

একটি পঞ্চভুজের অন্ত:কোণগুলোর সমষ্টি কত, অথবা শুধু একটি পঞ্চভুজের সমষ্টি কত , অথবা
একটি সুষম পঞ্চভুজের প্রতিটি অন্তঃকোণ ১০৮ হলে, সমষ্টি কত ? যেভাবেই প্রশ্ন আসুক
অন্তঃকোণগুলোর সমষ্টি বের করতে হয়। এক্ষেত্রে প্রতিটি অন্তঃকোণের মানকে বাহুর সংখ্যা
দিয়ে গুণ করলে কোণের সমষ্টি বের হয়। যেমন: , ঐ পঞ্চভুজটির অন্তঃকোণ সমূহের যোগফল
হবে ১০৮ × ৫ = ৫৪০ ডিগ্রি । ৫৪০ ÷ ৯০ = ৬ সমকোণ ।

যদি অন্ত:কোণের পরিমাণ দেয়া না থাকে তাহলে নিচের সূত্র প্রয়োগ করতে পারেন ।
অথবা নিচের চার্টটি মুখস্থ রাখতে পারেন তাহলে আর কিছুই করতে হবে না।

{\def\LTcaptype{none} % do not increment counter
\begin{longtable}[]{@{}
  >{\raggedright\arraybackslash}p{(\linewidth - 4\tabcolsep) * \real{0.2672}}
  >{\raggedright\arraybackslash}p{(\linewidth - 4\tabcolsep) * \real{0.2672}}
  >{\raggedright\arraybackslash}p{(\linewidth - 4\tabcolsep) * \real{0.2672}}@{}}
\toprule\noalign{}
\begin{minipage}[b]{\linewidth}\raggedright
\textbf{বহুভুজের নাম}
\end{minipage} & \begin{minipage}[b]{\linewidth}\raggedright
\textbf{অন্ত:কোণের মোট পরিমাপ}
\end{minipage} & \begin{minipage}[b]{\linewidth}\raggedright
\textbf{সমকোণ}
\end{minipage} \\
\midrule\noalign{}
\endhead
\bottomrule\noalign{}
\endlastfoot
পঞ্চভুজ & ৫৪০ ডিগ্রি & ৫৪০ ÷ ৯০ = ৬ সমকোণ \\
ষড়ভুজ & ৭২০ ডিগ্রি & ৭২০ ÷ ৯০ = ৮ সমকোণ \\
অষ্টভুজ & ১০৮০ ডিগ্রি & ১০৮০ ÷ ৯০ = ১২ সমকোণ \\
দশভুজ & ১৪৪০ ডিগ্রি & ১৪৪০ ÷ ৯০ = ১৬ সমকোণ \\
\multicolumn{3}{@{}>{\raggedright\arraybackslash}p{(\linewidth - 4\tabcolsep) * \real{0.8015} + 4\tabcolsep}@{}}{%
লক্ষ্য করুন , বাহুর পরিমাণ একটি বাড়লে কোণের পরিমাণ ২ সমকোণ করে বাড়ছে। সহজে
মনে রাখার জন্য এ পদ্ধতিটি ভালভাবে খেয়াল করুন ।} \\
\end{longtable}
}

\textbf{সূত্র:-০১:} অন্তঃস্থ কোণগুলোর সমষ্টি = (2n-4) সমকোণ , যেখানে n হচ্ছে
সুষম বহুভূজের মোট বাহুর সংখ্যা। ( যদি কত সমকোণ তা জানতে চায় তাহলে সরাসরি এই
সূত্র) যেমন:

১. একটি পঞ্চভুজের সমষ্টি কত? (৩৪তম বিসিএস)

\textbf{সমাধান}: সূত্র অনুযায়ী (2×5−4) সমকোণ = ১০−৪ = ৬ সমকোণ।

\textbf{আবার সমকোণ বের করতে না বলে কোণের মোট পরিমাপ বের করতে বললে,}

\textbf{সুত্র: ০২:} কোণগুলোর সমষ্টি = {[}90×(2n−4){]} ডিগ্রি , যেখানে n হচ্ছে
সুষম বহুভূজের মোট বাহুর সংখ্যা। (যদি কোণের পরিমাপ দেয়া থাকে তাহলে এই সূত্র।)
যেমন:

২. একটি পঞ্চভুজের অভ্যন্তরীণ ( অন্তঃস্থ) পাঁচটি কোণের সমষ্টি কত? {[}পররাষ্ট্র
মন্ত্রণালয়ের কর্মকর্তা ২০০৪ {]}

\textbf{সমাধান}: {[}90 × (2n−4){]} ডিগ্রি

= {[} 90 × (2 × 5 − 4 ) {]}

= {[} 90 × (10 − 4 ) {]}

= 90 × 6 = 540 ডিগ্রি

\paragraph{সংশ্লিষ্ট বিগত সালের প্রশ্ন ও
সমাধান}\label{ux9b8ux9b6ux9b2ux9b7ux99f-ux9acux997ux9a4-ux9b8ux9b2ux9b0-ux9aaux9b0ux9b6ux9a8-ux993-ux9b8ux9aeux9a7ux9a8-2}

\textbf{প্রশ্ন নম্বর-} ১. একটি সুষম ষড়ভুজের অন্তঃকোণগুলির সমষ্টি কত ডিগ্রী?
{[}বাংলাদেশ রোড ট্রান্সপোর্ট অথরিটির সহকারী পরিচালক-২০০৫{]}

\begin{itemize}
\item
  (ক) ৩৬০
\item
  (খ) ৫৪০
\item
  (গ) ৬৩০
\item
  (ঘ) ৭২০
\end{itemize}

\textbf{\hl{উত্তরঃ (ঘ) ৭২০}}

\textbf{সমাধান:}

{[}90×(2n-4){]} ডিগ্রি

= {[} 90 × (2 × 6 − 4 ) {]}

= {[} 90 × (12 − 4 ) {]}

= 90 × 8 = 720 ডিগ্রি

\textbf{প্রশ্ন নম্বর-} ২. অষ্টভুজের অন্তঃস্থ কোণসমূহের সমষ্টি কত? {[} মাধ্যমিক
বিদ্যালয় সহকারী শিক্ষক-২০০০{]}

\begin{itemize}
\item
  (ক) ১৬ সমকোণ
\item
  (খ) ১২ সমকোণ
\item
  (গ) ৮ সমকোণ
\item
  (ঘ) ৬ সমকোণ
\end{itemize}

\textbf{\hl{উত্তরঃ (খ) ১২ সমকোণ}}

\textbf{সমাধান:}

সূত্র অনুযায়ী (2×8−4) সমকোণ = ১৬−৪ = ১২ সমকোণ।

\textbf{প্রশ্ন নম্বর-} ৩. সুষম বহুভুজের বাহুর সংখ্যা n হলে কোণগুলোর সমষ্টি কত? {[}
মাধ্যমিক বিদ্যালয় সহকারী শিক্ষক-২০০১{]}

\begin{itemize}
\item
  (ক) n সমকোণ
\item
  (খ) 2n সমকোণ
\item
  (গ) (2n−4) সমকোণ
\item
  (ঘ) (2n−2) সমকোণ
\end{itemize}

\textbf{\hl{উত্তরঃ (গ) (2n−4) সমকোণ}}

\textbf{প্রশ্ন নম্বর-} ৪. একটি পঞ্চভুজের অভ্যন্তরীণ পাঁচটি কোণের সমষ্টি হবে-\/-
{[} সেকেন্ডারি এডুকেশন সেক্টর ইনভেস্টমেন্ট প্রোগ্রাম-২০১৫{]}

\begin{itemize}
\item
  (ক) ৩৬০ ডিগ্রি
\item
  (খ) ৫৪০ ডিগ্রি
\item
  (গ) ৪৫০ ডিগ্রি
\item
  (ঘ) ২৭০ ডিগ্রি
\end{itemize}

\textbf{\hl{উত্তরঃ (খ) ৫৪০ ডিগ্রি}}

\textbf{সমাধান:}

{[}90×(2n-4){]} ডিগ্রি

= {[} 90 × (2 × 5 − 4 ) {]}

= {[} 90 × (10 − 4 ) {]}

= 90 × 6 = 540 ডিগ্রি

\textbf{প্রশ্ন নম্বর-} ৫. একটি বহুভুজের বাহুর সংখ্যা ৬ হলে, বহুভুজের অন্তঃকোণগুলোর
সমষ্টি কত হবে? {[} প্রাথমিক প্রধান শিক্ষক-২০১২(ডালিয়া){]}

\begin{itemize}
\item
  (ক) চার সমকোণ
\item
  (খ) ছয় সমকোণ
\item
  (গ) সাত সমকোণ
\item
  (ঘ) আট সমকোণ
\end{itemize}

\textbf{\hl{উত্তরঃ (ঘ) আট সমকোণ}}

\textbf{সমাধান:}

সূত্র অনুযায়ী (2×6−4) সমকোণ = ১২−৪ = ৮ সমকোণ।

\textbf{প্রশ্ন নম্বর-} ৬. একটি ষড়ভুজের ছয়টি কোণের সমষ্টি কত? {[} থানা
একাডেমিক সুপারভাইজার-২০১৫{]}

\begin{itemize}
\item
  (ক) পাঁচ সমকোণ
\item
  (খ) ছয় সমকোণ
\item
  (গ) আট সমকোণ
\item
  (ঘ) তিন সমকোণ
\end{itemize}

\textbf{\hl{উত্তরঃ (গ) আট সমকোণ}}

\textbf{সমাধান:}

সূত্র অনুযায়ী (2×6−4) সমকোণ = ১২−৪ = ৮ সমকোণ।

\textbf{প্রশ্ন নম্বর-} ৭. কোনো সুষম দশভুজের প্রতিটি কোণ হবে- {[} আইন,বিচার ও
সংসদ বিষয়ক মন্ত্রণালয়ের সাব রেজিস্টার-২০০১{]}

\begin{itemize}
\item
  (ক) সূক্ষ্মকোণ
\item
  (খ) প্রবৃদ্ধকোণ
\item
  (গ) স্থুলকোণ
\item
  (ঘ) সমকোণ
\end{itemize}

\textbf{\hl{উত্তরঃ (গ) স্থুলকোণ}}

\textbf{সমাধান:}

দশভুজের অন্তঃস্থ কোণগুলোর সমষ্টি

= {[}90×(2n−4){]} ডিগ্রি

= {[} 90 ×(2×10 − 4 ) {]}

= {[} 90 ×(20 − 4 ) {]}

= 90×16 = 1440 ডিগ্রি

সুতরাং সুষম দশভুজের প্রতিটি কোণ

= 1440 ÷ 10 = 144 ডিগ্রি ।

যা 90 ডিগ্রি অপেক্ষা বেশি কিন্তু 180 ডিগ্রি অপেক্ষা কম এবং স্থুলকোণ।

সুতরাং কোনো সুষম দশভুজের প্রতিটি কোণ স্থুলকোণ হবে।

\section{14. বহুভুজ MCQ প্রশ্ন উত্তর
অনুশীলন}\label{ux9acux9b9ux9adux99c-mcq-ux9aaux9b0ux9b6ux9a8-ux989ux9a4ux9a4ux9b0-ux985ux9a8ux9b6ux9b2ux9a8}

\textbf{একটি সুষম বহুভুজের বহিঃস্থ কোণগুলির সমষ্টি}

ক) 180\^{}°

খ) 270\^{}°

\textbf{\hl{গ) 360\^{}° ✔}}

ঘ) 540\^{}°

\textbf{\hl{✔ সঠিক উত্তর: গ) 360\^{}°}}

\begin{quote}
\textbf{ব্যাখ্যা:} \emph{যেকোনো বহুভুজের বহিঃস্থ কোণগুলির সমষ্টি সর্বদা 360\^{}°
।}
\end{quote}

\textbf{আট বাহুবিশিষ্ট সুষম বহুভুজের একটি বহিঃস্থ কোণের পরিমাণ হবে-\/-\/-}

\textbf{\hl{ক) ৪৫ ডিগ্রি ✔}}

খ) ৬০ ডিগ্রি

গ) ৮০ ডিগ্রি

ঘ) ১০০ ডিগ্রি

\textbf{\hl{✔ সঠিক উত্তর: ক) ৪৫ ডিগ্রি}}

\begin{quote}
\textbf{ব্যাখ্যা:} \emph{বহিঃস্থ কোণ = (360\^{}°/8)=45\^{}°}
\end{quote}

\textbf{একটি সুষম বহুভুজের বাহুর সংখ্যা ১২ হলে প্রত্যেকটি বহিঃস্থ কোণের পরিমাণ
কত?}

ক) ১৫ ডিগ্রি

খ) ২০ ডিগ্রি

গ) ২৫ ডিগ্রি

\textbf{\hl{ঘ) ৩০ ডিগ্রি ✔}}

\textbf{\hl{✔ সঠিক উত্তর: ঘ) ৩০ ডিগ্রি}}

\begin{quote}
\textbf{ব্যাখ্যা:} \emph{বহিঃস্থ কোণ = (360\^{}°/12)=30\^{}°}
\end{quote}

\textbf{একটি সুষম বহুভুজের বাহুর সংখ্যা ৮ হলে প্রত্যেকটি অন্তঃস্থ কোণের পরিমাণ
কত?}

ক) ৩০ ডিগ্রি

খ) ৬০ ডিগ্রি

গ) ৯০ ডিগ্রি

\textbf{\hl{ঘ) ১৩৫ ডিগ্রি ✔}}

\textbf{\hl{✔ সঠিক উত্তর: ঘ) ১৩৫ ডিগ্রি}}

\begin{quote}
\textbf{ব্যাখ্যা:} \emph{বহিঃস্থ কোণ = (360\^{}°/8)=45\^{}°}

\emph{∴ অন্তঃস্থ কোণ = 180\^{}°−45\^{}°=135\^{}°}
\end{quote}

\textbf{একটি সুষম বহুভুজের একটি অন্তঃকোণের পরিমাণ ১৩৫ ডিগ্রি হলে বহুভুজটির বাহুর
সংখ্যা হবে-\/-}

ক) ৬০

খ) ৭

\textbf{\hl{গ) ৮ ✔}}

ঘ) ১০

\textbf{\hl{✔ সঠিক উত্তর: গ) ৮}}

\begin{quote}
\textbf{ব্যাখ্যা:} \emph{বহিঃস্থ কোণ = 180\^{}°−135\^{}°=45\^{}°}

\emph{∴ বহুভুজটির বাহুর সংখ্যা = (360\^{}°/45\^{}°)=8 টি}
\end{quote}

\textbf{একটি সুষম বহুভুজের প্রত্যেকটি অন্তঃস্থ কোণের পরিমাপ ১৬২ ডিগ্রি হলে
বহুভুজের বাহুসংখ্যা কত?}

ক) ১৬০

\textbf{\hl{খ) ২০ ✔}}

গ) ২৪

ঘ) ২৬

\textbf{\hl{✔ সঠিক উত্তর: খ) ২০}}

\begin{quote}
\textbf{ব্যাখ্যা:} \emph{বহিঃস্থ কোণ = 180\^{}°−162\^{}°=18\^{}°}

\emph{∴ বহুভুজটির বাহুর সংখ্যা = (360\^{}°/18\^{}°)=20 টি}
\end{quote}

\textbf{একটি সুষম বহুভুজের প্রত্যেকটি বহিঃস্থ কোণের পরিমাণ ৩০ ডিগ্রি হলে
বহুভুজটির বাহুর সংখ্যা কত?}

ক) ১৮ টি

খ) ১৬ টি

\textbf{\hl{গ) ১২ টি ✔}}

ঘ) ৮ টি

\textbf{\hl{✔ সঠিক উত্তর: গ) ১২ টি}}

\begin{quote}
\textbf{ব্যাখ্যা:} \emph{বহুভুজটির বাহুর সংখ্যা = (360\^{}°/30\^{}°)=12 টি}
\end{quote}

\textbf{একটি সুষম বহুভুজের একটি অন্তঃকোণের পরিমাণ ১৪৪ ডিগ্রি হলে বহুভুজটির বাহুর
সংখ্যা হবে-\/-}

ক) ৭০

খ) ৮

গ) ৯০

\textbf{\hl{ঘ) ১০ ✔}}

\textbf{\hl{✔ সঠিক উত্তর: ঘ) ১০}}

\begin{quote}
\textbf{ব্যাখ্যা:} \emph{বহিঃস্থ কোণ = 180\^{}°−144\^{}°=36\^{}°}

\emph{⇒ বহুভুজটির বাহুর সংখ্যা = (360\^{}°/36\^{}°)=10 টি}
\end{quote}

\textbf{একটি সুষম বহুভুজের প্রত্যেকটি অন্তঃস্থ কোণের পরিমাপ ১২০ ডিগ্রি হলে
বহুভুজটির বাহুরসংখ্যা কত?}

\textbf{\hl{ক) ৬০ ✔}}

খ) ৮

গ) ১০

ঘ) ১২

\textbf{\hl{✔ সঠিক উত্তর: ক) ৬০}}

\begin{quote}
\textbf{ব্যাখ্যা:} \emph{বহিঃস্থ কোণ = 180\^{}°−120\^{}°=60\^{}°}

\emph{⇒ বহুভুজটির বাহুর সংখ্যা = (360\^{}°/60\^{}°)=6 টি}
\end{quote}

\textbf{অষ্টভুজের অন্তঃস্থ কোণসমূহের সমষ্টি কত?}

ক) ১৬ সমকোণ

\textbf{\hl{খ) ১২ সমকোণ ✔}}

গ) ৮ সমকোণ

ঘ) ৬ সমকোণ

\textbf{\hl{✔ সঠিক উত্তর: খ) ১২ সমকোণ}}

\begin{quote}
\textbf{ব্যাখ্যা:} \emph{সূত্র অনুযায়ী অন্তঃকোণসমূহের সমষ্টি = (2×8−4) সমকোণ =
16−4=12 সমকোণ}
\end{quote}

\textbf{একটি পঞ্চভুজের অভ্যন্তরীণ পাঁচটি কোণের সমষ্টি হবে-\/-}

ক) ৩৬০ ডিগ্রি

\textbf{\hl{খ) ৫৪০ ডিগ্রি ✔}}

গ) ৪৫০ ডিগ্রি

ঘ) ২৭০ ডিগ্রি

\textbf{\hl{✔ সঠিক উত্তর: খ) ৫৪০ ডিগ্রি}}

\begin{quote}
\textbf{ব্যাখ্যা:} \emph{পঞ্চভুজের বাহুসংখ্যা n=5}

\emph{⇒ অন্তঃকোণসমূহের সমষ্টি = (2n−4)×90\^{}°}

\emph{⇒ = (2×5−4)×90\^{}°=6×90\^{}°=540\^{}°}
\end{quote}

\textbf{একটি বহুভুজের বাহুর সংখ্যা ৬ হলে বহুভুজটির অন্তঃকোণের সমষ্টি কত?}

ক) 540\^{}°

\textbf{\hl{খ) 720\^{}° ✔}}

গ) 900\^{}°

ঘ) 1080\^{}°

\textbf{\hl{✔ সঠিক উত্তর: খ) 720\^{}°}}

\begin{quote}
\textbf{ব্যাখ্যা:} \emph{বহুভুজের অন্তঃকোণের সমষ্টি = (n−2)×180\^{}°}

\emph{= (6−2)×180\^{}°=4×180\^{}°=720\^{}° ।}
\end{quote}

\textbf{কোনো সুষম বহুভুজের প্রতিটি আন্তঃকোণ 135\^{}° হলে ঐ বহুভুজের বাহুর সংখ্যা
কত?}

ক) ১০

খ) ৯

\textbf{\hl{গ) ৮ ✔}}

ঘ) ৭

\textbf{\hl{✔ সঠিক উত্তর: গ) ৮}}

\begin{quote}
\textbf{ব্যাখ্যা:} \emph{ধরি, বাহুর সংখ্যা = n}

\emph{প্রশ্নমতে, ((n−2)×180\^{}°/n)=135\^{}°}

\emph{⇒ 135\^{}°n=180\^{}°n−360\^{}°}

\emph{⇒ 45\^{}°n=360\^{}°}

\emph{∴ n=8 ।}
\end{quote}

\textbf{একটি ষড়ভুজের ছয়টি কোণের সমষ্টি কত?}

ক) পাঁচ সমকোণ

খ) ছয় সমকোণ

\textbf{\hl{গ) আট সমকোণ ✔}}

ঘ) তিন সমকোণ

\textbf{\hl{✔ সঠিক উত্তর: গ) আট সমকোণ}}

\begin{quote}
\textbf{ব্যাখ্যা:} \emph{ষড়ভুজের অন্তঃকোণগুলির সমষ্টি =
(6−2)×180\^{}°=4×180\^{}°=720\^{}°}

\emph{এক সমকোণ = 90\^{}° হওয়ায়, মোট সমকোণ = (720\^{}°/90\^{}°)=8 সমকোণ।}
\end{quote}

\textbf{সুষম বহুভুজের বাহুর সংখ্যা n হলে কোণগুলোর সমষ্টি কত?}

ক) n সমকোণ

খ) 2n সমকোণ

\textbf{\hl{গ) (2n−4) সমকোণ ✔}}

ঘ) (2n−2) সমকোণ

\textbf{\hl{✔ সঠিক উত্তর: গ) (2n−4) সমকোণ}}

\begin{quote}
\textbf{ব্যাখ্যা:} \emph{অন্তঃকোণগুলির সমষ্টি =
(n−2)×180\^{}°=(n−2)×2×90\^{}°=(2n−4) সমকোণ।}
\end{quote}

\textbf{একটি পঞ্চভুজের অভ্যন্তরীণ পাঁচটি কোণের সমষ্টি হবে-}

ক) 360\^{}°

\textbf{\hl{খ) 540\^{}° ✔}}

গ) 450\^{}°

ঘ) 270\^{}°

\textbf{\hl{✔ সঠিক উত্তর: খ) 540\^{}°}}

\begin{quote}
\textbf{ব্যাখ্যা:} \emph{পঞ্চভুজের অভ্যন্তরীণ কোণগুলির সমষ্টি =
(5−2)×180\^{}°=3×180\^{}°=540\^{}° ।}
\end{quote}

\textbf{একটি পঞ্চভুজের সমষ্টি-}

ক) ৪ সমকোণ

\textbf{\hl{খ) ৬ সমকোণ ✔}}

গ) ৮ সমকোণ

ঘ) ১০ সমকোণ

\textbf{\hl{✔ সঠিক উত্তর: খ) ৬ সমকোণ}}

\begin{quote}
\textbf{ব্যাখ্যা:} \emph{পঞ্চভুজের অন্তঃকোণগুলির সমষ্টি =
(5−2)×180\^{}°=540\^{}°}

\emph{যেহেতু ১ সমকোণ = 90\^{}° , তাই মোট সমকোণ = (540\^{}°/90\^{}°)=6
সমকোণ।}
\end{quote}

\textbf{সুষম পঞ্চভুজের বহিঃস্থ কোণের পরিমাণ কত?}

\textbf{\hl{ক) 72\^{}° ✔}}

খ) 30\^{}°

গ) 40\^{}°

ঘ) 60\^{}°

\textbf{\hl{✔ সঠিক উত্তর: ক) 72\^{}°}}

\begin{quote}
\textbf{ব্যাখ্যা:} \emph{বহুভুজের বহিঃস্থ কোণ বের করার সূত্র = (360\^{}°/বাহুর
সংখ্যা)}

\emph{সুষম পঞ্চভুজের ক্ষেত্রে, বহিঃস্থ কোণ = (360\^{}°/5)=72\^{}° ।}
\end{quote}

\textbf{একটি পঞ্চভুজের কোণগুলোর অনুপাত হলো 9:10:12:14:15 । ঐ পঞ্চভুজটির ক্ষুদ্রতম
ও বৃহত্তম কোণের সমষ্টি কত?}

ক) 54\^{}°

খ) 81\^{}°

গ) 135\^{}°

\textbf{\hl{ঘ) 216\^{}° ✔}}

\textbf{\hl{✔ সঠিক উত্তর: ঘ) 216\^{}°}}

\begin{quote}
\textbf{ব্যাখ্যা:} \emph{ধরি, কোণগুলো যথাক্রমে 9x , 10x , 12x , 14x এবং 15x}

\emph{পঞ্চভুজের ৫ কোণের সমষ্টি 540\^{}° হওয়ায়, 9x+10x+12x+14x+15x=540\^{}°}

\emph{⇒ 60x=540\^{}°⇒x=9\^{}°}

\emph{ক্ষুদ্রতম কোণ = 9×9\^{}°=81\^{}° এবং বৃহত্তম কোণ = 15×9\^{}°=135\^{}°}

\emph{∴ সমষ্টি = 81\^{}°+135\^{}°=216\^{}° ।}
\end{quote}

\textbf{সুষম বহুভুজের একটি বহিঃস্থ কোণের পরিমাণ 45\^{}° হলে এর বাহুর সংখ্যা
হবে-}

\textbf{\hl{ক) ৮ ✔}}

খ) ১০

গ) ১২

ঘ) ১৬

\textbf{\hl{✔ সঠিক উত্তর: ক) ৮}}

\begin{quote}
\textbf{ব্যাখ্যা:} \emph{সুষম বহুভুজের বাহুর সংখ্যা n=(360\^{}°/বহিঃস্থ কোণের
পরিমাণ)}

\emph{⇒ n=(360\^{}°/45\^{}°)=8 ।}
\end{quote}

\textbf{একটি সুষম ষড়ভুজের অন্তঃস্থ যে কোনো একটি কোণের পরিমাণ কত?}

\textbf{\hl{ক) 120\^{}° ✔}}

খ) 100\^{}°

গ) 90\^{}°

ঘ) 60\^{}°

\textbf{\hl{✔ সঠিক উত্তর: ক) 120\^{}°}}

\begin{quote}
\textbf{ব্যাখ্যা:} \emph{সুষম বহুভুজের প্রতিটি অন্তঃকোণের পরিমাণ =
((n−2)×180\^{}°/n)}

\emph{ষড়ভুজের ক্ষেত্রে n=6 হওয়ায়, প্রতিটি কোণ =
((6−2)×180\^{}°/6)=(4×180\^{}°/6)=120\^{}°}
\end{quote}

\textbf{সুষম ষড়ভুজের একটি বাহুকে বর্ধিত করে উৎপন্ন বহিঃস্থ কোণের পরিমাণ কত
ডিগ্রি?}

ক) 75\^{}°

\textbf{\hl{খ) 60\^{}° ✔}}

গ) 90\^{}°

ঘ) 180\^{}°

\textbf{\hl{✔ সঠিক উত্তর: খ) 60\^{}°}}

\begin{quote}
\textbf{ব্যাখ্যা:} \emph{সুষম বহুভুজের বহিঃস্থ কোণ = (360\^{}°/n)}

\emph{ষড়ভুজের ক্ষেত্রে, বহিঃস্থ কোণ = (360\^{}°/6)=60\^{}° ।}
\end{quote}

\textbf{একটি সুষম ষড়ভুজের অন্তঃকোণগুলির সমষ্টি কত ডিগ্রি?}

ক) ৩৬০

খ) ৫৩০

গ) ৬৩০

\textbf{\hl{ঘ) ৭২০ ✔}}

\textbf{\hl{✔ সঠিক উত্তর: ঘ) ৭২০}}

\begin{quote}
\textbf{ব্যাখ্যা:} \emph{ষড়ভুজের অন্তঃকোণগুলির সমষ্টি =
(6−2)×180\^{}°=4×180\^{}°=720\^{}° ।}
\end{quote}

\textbf{একটি বহুভুজের বাহুর সংখ্যা ৬ হলে, বহুভুজের অন্তঃকোণগুলোর সমষ্টি কত হবে?}

ক) চার সমকোণ

খ) ছয় সমকোণ

গ) সাত সমকোণ

\textbf{\hl{ঘ) আট সমকোণ ✔}}

\textbf{\hl{✔ সঠিক উত্তর: ঘ) আট সমকোণ}}

\begin{quote}
\textbf{ব্যাখ্যা:} \emph{বহুভুজের অন্তঃকোণগুলির সমষ্টি =
(n−2)×180\^{}°=(6−2)×180\^{}°=4×180\^{}°=720\^{}°}

\emph{যেহেতু এক সমকোণ = 90\^{}° , তাই মোট সমকোণ = (720\^{}°/90\^{}°)=8
সমকোণ।}
\end{quote}

\textbf{একটি সুষম পঞ্চভুজের একটি অন্তঃ কোণের পরিমাণ হবে-}

ক) 90\^{}°

খ) 110\^{}°

গ) 120\^{}°

\textbf{\hl{ঘ) 108\^{}° ✔}}

\textbf{\hl{✔ সঠিক উত্তর: ঘ) 108\^{}°}}

\begin{quote}
\textbf{ব্যাখ্যা:} \emph{সুষম পঞ্চভুজের প্রতিটি অন্তঃকোণের পরিমাণ =
((5−2)×180\^{}°/5)=(3×180\^{}°/5)=(540\^{}°/5)=108\^{}° ।}
\end{quote}

\textbf{একটি সুষম ষড়ভুজের একটি অন্তঃকোণের পরিমাণ হবে-}

ক) 90\^{}°

খ) 105\^{}°

\textbf{\hl{গ) 120\^{}° ✔}}

ঘ) 145\^{}°

\textbf{\hl{✔ সঠিক উত্তর: গ) 120\^{}°}}

\begin{quote}
\textbf{ব্যাখ্যা:} \emph{সুষম ষড়ভুজের প্রতিটি অন্তঃকোণের পরিমাণ =
((6−2)×180\^{}°/6)=(4×180\^{}°/6)=120\^{}° ।}
\end{quote}

\textbf{কোনো সুষম দশভুজের প্রতিটি কোণ হবে-}

ক) সূক্ষ্মকোণ

খ) প্রবৃদ্ধকোণ

\textbf{\hl{গ) স্থূলকোণ ✔}}

ঘ) সমকোণ

\textbf{\hl{✔ সঠিক উত্তর: গ) স্থূলকোণ}}

\begin{quote}
\textbf{ব্যাখ্যা:} \emph{সুষম দশভুজের প্রতিটি অন্তঃকোণের পরিমাণ =
((10−2)×180\^{}°/10)=(8×180\^{}°/10)=144\^{}°}

\emph{যেহেতু 144\^{}° কোণটি 90\^{}° অপেক্ষা বেশি এবং 180\^{}° অপেক্ষা কম,
সেহেতু এটি একটি স্থূলকোণ।}
\end{quote}

\textbf{একটি সুষম বহুভুজের একটি অন্তঃকোণের পরিমাণ 160\^{}° হলে, বহুভুজটির বাহুর
সংখ্যা হবে-}

ক) ৬

খ) ৮

গ) ১০

\textbf{\hl{ঘ) ১৮ ✔}}

\textbf{\hl{✔ সঠিক উত্তর: ঘ) ১৮}}

\begin{quote}
\textbf{ব্যাখ্যা:} \emph{ধরি, বাহুর সংখ্যা = n}

\emph{প্রশ্নমতে, ((n−2)×180\^{}°/n)=160\^{}°}

\emph{⇒ 180n−360=160n}

\emph{⇒ 20n=360}

\emph{∴ n=18 ।}
\end{quote}

\textbf{একটি সুষম বহুভুজের প্রতিটি অন্তঃস্থ কোণ 144\^{}° হলে তার বাহুর সংখ্যা
কত?}

ক) ২০

\textbf{\hl{খ) ১০ ✔}}

গ) ১২

ঘ) ১৮

\textbf{\hl{✔ সঠিক উত্তর: খ) ১০}}

\begin{quote}
\textbf{ব্যাখ্যা:} \emph{ধরি, বাহুর সংখ্যা = n}

\emph{প্রশ্নমতে, ((n−2)×180\^{}°/n)=144\^{}°}

\emph{⇒ 180n−360=144n}

\emph{⇒ 36n=360}

\emph{∴ n=10 ।}
\end{quote}

\textbf{আট বাহু বিশিষ্ট সুষম বহুভুজের একটি বহিঃস্থ কোণের পরিমাণ কত?}

\textbf{\hl{ক) 45\^{}° ✔}}

খ) 60\^{}°

গ) 80\^{}°

ঘ) 100\^{}°

\textbf{\hl{✔ সঠিক উত্তর: ক) 45\^{}°}}

\begin{quote}
\textbf{ব্যাখ্যা:} \emph{সুষম বহুভুজের বহিঃস্থ কোণগুলির সমষ্টি 360\^{}°}

\emph{∴ প্রতিটি বহিঃস্থ কোণের পরিমাণ = (360\^{}°/8)=45\^{}° ।}
\end{quote}

\section{15. চতুর্ভুজ কাকে বলে? সংজ্ঞা ,বৈশিষ্ট্য , প্রকারভেদসহ এ টু
জেড}\label{ux99aux9a4ux9b0ux9adux99c-ux995ux995-ux9acux9b2-ux9b8ux99cux99e-ux9acux9b6ux9b7ux99fux9af-ux9aaux9b0ux995ux9b0ux9adux9a6ux9b8ux9b9-ux98f-ux99f-ux99cux9a1}

চতুর্ভুজ হল একটি জ্যামিতিক আকৃতি, যা চারটি বাহু দ্বারা সংযোজিত হয়ে তৈরি হয়। এটি
প্রায় সমদ্বিখণ্ডীয় হতে পারে বা অসমদ্বিখণ্ডীয় হতে পারে।

একটি চতুর্ভুজের বাহুর দৈর্ঘ্য ও কোণের মাধ্যমে এর বিস্তৃতি নির্ণয় করা যায়। চতুর্ভুজ হল
একটি সরলরেখার পরিকল্পিত আকৃতি, যার সব কোণ মিলে ৩৬০ ডিগ্রি।

চতুর্ভুজ গণিতে একটি গুরুত্বপূর্ণ আকৃতি হিসাবে বিবেচিত হয়। এটি ক্ষেত্রফল হিসাবে
ব্যবহৃত হয় এবং এর পরিমাপ করার জন্য ব্যবহৃত হয়।

চতুর্ভুজ বিভিন্ন ধরনের হতে পারে যেমন- আয়ত, বর্গ , রম্বস , সামন্তরিক ,
ট্র্যাপিজিয়াম ।চতুর্ভুজ সম্পর্কে যদি স্পষ্ট ও স্বচ্ছ ধারণা নিতে চান তাহলে আজকের
টিউটরিয়ালটি আপনার জন্যই ।

\begin{itemize}
\item
  প্রাথমিক আলোচনা
\item
  আয়তক্ষেত্র
\item
  বর্গক্ষেত্র
\item
  রম্বস
\item
  সামন্তরিক
\item
  ট্রাপিজিয়াম
\end{itemize}

\subsubsection{প্রাথমিক
আলোচনা}\label{ux9aaux9b0ux9a5ux9aeux995-ux986ux9b2ux99aux9a8-5}

চারটি বাহু দ্বারা সীমাবদ্ধ ক্ষেত্রকে চতুর্ভুজ বলে। সাধারণত আয়তক্ষেত্র, বর্গক্ষেত্র,
সামন্তরিক, রম্বস এগুলো সবই এক এক প্রকার চতুর্ভুজ।

এই অধ্যায়ের সাথে সম্পর্কিত বিভিন্ন নিয়মাবলী ও সূত্র এবং বিগত সালে বিভিন্ন
প্রতিযোগিতামূলক পরীক্ষায় আসা প্রশ্নগুলো নিয়ে নিচে ধারাবাহিকভাবে আলোচনা করা হল
।

\paragraph{চতুর্ভুজ কাকে
বলে?}\label{ux99aux9a4ux9b0ux9adux99c-ux995ux995-ux9acux9b2}

চারটি সরলরেখা দ্বারা সীমাবদ্ধ ক্ষেত্রকে চতুর্ভুজ বলে।

চতুর্ভুজের চার কোণের সমষ্টি চার সমকোণ । (৩৬০°)

\paragraph{সংশ্লিষ্ট
সূত্রঃ}\label{ux9b8ux9b6ux9b2ux9b7ux99f-ux9b8ux9a4ux9b0}

\textbf{চতুর্ভুজের ক্ষেত্রফল = (দৈর্ঘ্য × প্রস্থ) বর্গ একক}

\textbf{চতুর্ভুজের পরিসীমা = ২(দৈঘ্য+প্রস্থ)}

সুতরাং চতুর্ভুজের পরিসীমা 2(a+b) অর্থাৎ দৈর্ঘ্য ও প্রস্থের যোগফলকে দ্বিগুন করলে
পরিসীমা বের হবে।

\subsubsection{আয়তক্ষেত্র (rectangle)
:}\label{ux986ux9afux9a4ux995ux9b7ux9a4ux9b0-rectangle}

যে চতুর্ভূজের বিপরীত বাহুগুলো পরস্পর সমান ও সমান্তরাল এবং কোণগুলো সমকোণ তাকে
আয়তক্ষেত্র বলে।

উপরের চিত্রটিতে ABCD চতুর্ভূজের দৈর্ঘ্য AD = BC এবং প্রস্থ AB = CD এবং প্রতিটি
কোণ এক সমকোণ ।

সুতরাং চতুর্ভুজটি একটি আয়তক্ষেত্র।

\paragraph{আয়তক্ষেত্রের
বৈশিষ্ট্যঃ}\label{ux986ux9afux9a4ux995ux9b7ux9a4ux9b0ux9b0-ux9acux9b6ux9b7ux99fux9af}

আয়তক্ষেত্রের বিপরীত বাহুদ্বয় পরস্পর সমান ।

আয়তক্ষেত্রের প্রত্যেকটি কোণ সমকোণ ।

আয়তক্ষেত্রের কর্ণদ্বয় পরস্পর সমান ।

আয়তক্ষেত্রের কর্ণদ্বয় পরস্পরকে সমদ্বিখন্ডিত করে।

আয়তক্ষেত্রের একটি কর্ণ আয়তক্ষেত্রটিকে দুটি সর্বসম ত্রিভুজে বিভক্ত করে।

\paragraph{আয়তক্ষেত্র সংক্রান্ত বিভিন্ন
সূত্র:}\label{ux986ux9afux9a4ux995ux9b7ux9a4ux9b0-ux9b8ux995ux9b0ux9a8ux9a4-ux9acux9adux9a8ux9a8-ux9b8ux9a4ux9b0}

আয়তক্ষেত্রটির ক্ষেত্রফল : (দৈর্ঘ্য × প্রস্থ) বর্গ একক

আয়তক্ষেত্রটির পরিসীমা : 2(দৈর্ঘ্য + প্রস্থ) একক

আয়ক্ষেত্রের কর্ণঃ √ \({(\text{দৈর্ঘ্য})}^{\text{২}} + {(\text{প্রস্থ})}^{\text{২}}\) একক।

\subsubsection{বর্গক্ষেত্রঃ}\label{ux9acux9b0ux997ux995ux9b7ux9a4ux9b0}

যে চতুর্ভুজের চারটি বাহুই পরস্পর সমান ও সমান্তরাল এবং কোণগুলো সমকোণ তাকে
বর্গক্ষেত্র বলে ।

\paragraph{বর্গক্ষেত্রের
বৈশিষ্ট্যঃ}\label{ux9acux9b0ux997ux995ux9b7ux9a4ux9b0ux9b0-ux9acux9b6ux9b7ux99fux9af}

আয়তক্ষেত্রের দুটি সন্নিহিত (লাগালাগি) বাহু সমান হলে তাকে বর্গ বলে ।

বর্গক্ষেত্রের সকল বাহু সমান দৈর্ঘ্যের হয়।

বর্গক্ষেত্রের প্রত্যেকটি কোণ সমকোণ( ৯০ডিগ্রি)

বর্গক্ষেত্রের কর্ণদ্বয় পরস্পর সমান।

বর্গক্ষেত্রের কর্ণদ্বয় পরস্পরকে সমকোণে সমদ্বিখন্ডিত করে ।

\paragraph{বর্গ সংক্রান্ত
সূত্রাবলী:}\label{ux9acux9b0ux997-ux9b8ux995ux9b0ux9a8ux9a4-ux9b8ux9a4ux9b0ux9acux9b2}

বর্গক্ষেত্রের ক্ষেত্রফল : ( বাহু × বাহু) বর্গ একক অর্থাৎ \(a^{2}\)

বর্গক্ষেত্রের পরিসীমা : ( ৪ × একটি বাহুর দৈর্ঘ্য) একক অর্থাৎ 4a

বর্গক্ষেত্রের কর্ণ √2a (এখানে a হলো এক বাহুর দৈর্ঘ্য)

\paragraph{কর্ণ সম্পর্কিত
:}\label{ux995ux9b0ux9a3-ux9b8ux9aeux9aaux9b0ux995ux9a4}

একটি বর্গক্ষেত্র অথবা আয়তক্ষেত্রের এক কোণ থেকে অপর কোণ পর্যন্ত দূরত্বকে কর্ণ বলে।

বর্গক্ষেত্রের কর্ণটি বর্গক্ষেত্রটিকে দুটি সর্বসম সমকোণী ত্রিভুজে বিভক্ত করে।

বর্গক্ষেত্রের কর্ণের দৈর্ঘ্য = √2a

যে কোনো চতুর্ভুজের বিপরীত কৌণিক শীর্ষের সংযোজক সরলরেখাকে কর্ণ বলে।

যে কোনো চতুর্ভুজের কর্ণদ্বয়ের সমষ্টি তার পরিসীমা অপেক্ষা ক্ষুদ্রতর।

\subsubsection{রম্বস (Rhombus) :}\label{ux9b0ux9aeux9acux9b8-rhombus}

যে চতুর্ভুজের চারটি বাহুই পরস্পর সমান ও সমান্তরাল কিন্তু কোণগুলো সমকোণ নয় তাকে
রম্বস বলে।

\paragraph{রম্বস সম্পর্কিত
সূত্রাবলী:}\label{ux9b0ux9aeux9acux9b8-ux9b8ux9aeux9aaux9b0ux995ux9a4-ux9b8ux9a4ux9b0ux9acux9b2}

রম্বসের ক্ষেত্রফল = \(\frac{1}{2}\) × (কর্ণদ্বয়ের দৈর্ঘ্যের গুণফল) বর্গ একক

রম্বসের পরিসীমা: = ৪ × একটি বাহুর দৈর্ঘ্য (একক)

\paragraph{রম্বসের বৈশিষ্ট্য
সমূহ:}\label{ux9b0ux9aeux9acux9b8ux9b0-ux9acux9b6ux9b7ux99fux9af-ux9b8ux9aeux9b9}

সামন্তরিকের দুটি সন্নিহিত বাহু সমান হলে তাকে রম্বস বলে ।

রম্বসের সকল বাহু সমান হয়।

রম্বসের বিপরীত কোণগুলো পরস্পর সমান ।

রম্বসের একটি কোণও সমকোণ নয় ।

রম্বসের কর্ণদ্বয় অসমান ।

রম্বসের সন্নিহিত কোণদ্বয়ের সমষ্টি ২ সমকোণ ।

রম্বসের কর্ণদ্বয় পরস্পরকে সমকোণে সমদ্বিখন্ডিত করে।

\subsubsection{সামন্তরিক
(Parallelogram):}\label{ux9b8ux9aeux9a8ux9a4ux9b0ux995-parallelogram}

যে চতুর্ভুজের বিপরীত বাহুগুলো পরস্পর সমান ও সমান্তরাল কিন্তু কোণগুলো সমকোণ নয় তাকে
সামন্তরিক বলে।

\paragraph{সামন্তরিকের
বৈশিষ্ট্য:}\label{ux9b8ux9aeux9a8ux9a4ux9b0ux995ux9b0-ux9acux9b6ux9b7ux99fux9af}

সামন্তরিকের বিপরীত বাহুদ্বয় পরস্পর সমান ।

সামন্তরিকের বিপরীত কোণগুলো পরস্পর সমান ।

সামন্তরিকের যে কোন দুইটি সন্নিহিত কোণ পরস্পরের সম্পূরক ।

সামন্তরিকের কর্ণদ্বয় অসমান ।

সামন্তরিকের কর্ণদ্বয় পরস্পরকে সমদ্বিখন্ডিত করে।

সামন্তরিকের প্রত্যেক কর্ণ সামন্তরিককে দুটি সর্বসম ত্রিভুজে বিভক্ত করে ।

\paragraph{সামন্তরিকের
সূত্র:}\label{ux9b8ux9aeux9a8ux9a4ux9b0ux995ux9b0-ux9b8ux9a4ux9b0}

সামন্তরিকের ক্ষেত্রফল = ( ভূমি × উচ্চতা) বর্গ একক

সামন্তরিকের পরিসীমা = ২(দৈর্ঘ্য + প্রস্থ) একক

\subsubsection{ট্রাপিজিয়াম
(Trapezium):}\label{ux99fux9b0ux9aaux99cux9afux9ae-trapezium}

যে চতুর্ভুজের দুইটি বাহু পরস্পর সমান্তরাল কিন্তু অসমান এবং অন্য বাহুদ্বয় অসমান্তরাল
তাকে ট্রাপিজিয়াম বলে ৷

\paragraph{ট্রাপিজিয়ামের
বৈশিষ্ট্য:}\label{ux99fux9b0ux9aaux99cux9afux9aeux9b0-ux9acux9b6ux9b7ux99fux9af}

যে চতুর্ভুজের কেবলমাত্র দুইটি বাহু সমান্তরাল, তাকে ট্রাপিজিয়াম বলে ।

ট্রাপিজিয়ামের সমান্তরাল বাহুদ্বয়ের একটিকে ভূমি এবং অসমান্তরাল বাহুদ্বয়কে তির্যক
বাহু বলা হয় ।

ট্রাপিজিয়ামের তির্যক বাহুদ্বয় সমান হলে একে সমদ্বিবাহু ট্রাপিজিয়াম বলা হয় ৷

ট্রাপিজিয়ামের সমান্তরাল বাহুদ্বয় কখনও সমান হতে পারে না।

ট্রাপিজিয়ামের সমান্তরাল বাহুদ্বয় সমান হলে তা একটি আয়তক্ষেত্র বা বর্গক্ষেত্রে
পরিণত হবে।

\paragraph{ট্রাপিজিয়ামের
সূত্র:}\label{ux99fux9b0ux9aaux99cux9afux9aeux9b0-ux9b8ux9a4ux9b0}

ট্রাপিজিয়ামের ক্ষেত্রফল = \(\frac{1}{2}\) × উচ্চতা × সমান্তরাল বাহুদ্বয়ের সমষ্টি
অর্থাৎ \(\frac{1}{2}\) × h × (a + b)

\section{16. চতুর্ভুজ MCQ প্রশ্ন উত্তর
অনুশীলন}\label{ux99aux9a4ux9b0ux9adux99c-mcq-ux9aaux9b0ux9b6ux9a8-ux989ux9a4ux9a4ux9b0-ux985ux9a8ux9b6ux9b2ux9a8}

\textbf{বৃত্তস্থ চতুর্ভুজের একটি কোণ ৯০° হলে, উহার বিপরীত কোণের পরিমাণ কত
হবে?}

ক) ২০°

খ) ১১০°

\textbf{\hl{গ) ৯০° ✔}}

ঘ) ১২০°

\textbf{\hl{✔ সঠিক উত্তর: গ) ৯০°}}

\begin{quote}
\textbf{ব্যাখ্যা:} \emph{বৃত্তস্থ চতুর্ভুজের বিপরীত কোণদ্বয়ের সমষ্টি 180° ।}

\emph{∴ একটি কোণ 90° হলে অপর বিপরীত কোণটি হবে = 180°−90°=90° ।}
\end{quote}

\textbf{ABCD চতুর্ভুজের ∠A+∠B+∠C=2∠D হলে, ∠D এর মান কত?}

ক) 90°

\textbf{\hl{খ) 120° ✔}}

গ) 135°

ঘ) 45°

\textbf{\hl{✔ সঠিক উত্তর: খ) 120°}}

\begin{quote}
\textbf{ব্যাখ্যা:} \emph{আমরা জানি, চতুর্ভুজের চারটি কোণের সমষ্টি 360° 。}

\emph{∴ ∠A+∠B+∠C+∠D=360°}

\emph{বা, 2∠D+∠D=360° {[} ∴∠A+∠B+∠C=2∠D {]}}

\emph{বা, 3∠D=360°}

\emph{∴ ∠D=(360°/3)=120°}
\end{quote}

\textbf{একটি আয়তক্ষেত্রের দৈর্ঘ্য ও প্রস্থকে দ্বিগুণ করলে ক্ষেত্রফল মূল আয়তক্ষেত্রের
ক্ষেত্রফলের-\/-}

ক) দ্বিগুণ হবে

\textbf{\hl{খ) চারগুণ হবে ✔}}

গ) ছয়গুণ হবে

ঘ) আটগুণ হবে

\textbf{\hl{✔ সঠিক উত্তর: খ) চারগুণ হবে}}

\begin{quote}
\textbf{ব্যাখ্যা:} \emph{ধরি, আয়তক্ষেত্রের দৈর্ঘ্য = a এবং প্রস্থ = b}

\emph{ক্ষেত্রফল = a×b=ab}

\emph{দৈর্ঘ্য ও প্রস্থ দ্বিগুণ করলে নতুন দৈর্ঘ্য = 2a এবং নতুন প্রস্থ = 2b}

\emph{নতুন ক্ষেত্রফল = 2a×2b=4ab}

\emph{∴ ক্ষেত্রফল চারগুণ হবে।}
\end{quote}

\textbf{কোন চতুর্ভুজটির কেবলমাত্র দু\textquotesingle টি বাহু সমান্তরাল?}

ক) বর্গক্ষেত্র

খ) আয়তক্ষেত্র

গ) রম্বস

\textbf{\hl{ঘ) ট্রাপিজিয়াম ✔}}

\textbf{\hl{✔ সঠিক উত্তর: ঘ) ট্রাপিজিয়াম}}

\begin{quote}
\textbf{ব্যাখ্যা:} \emph{ট্রাপিজিয়াম হলো চতুর্ভুজের একটি বিশেষ রূপ, যে চতুর্ভুজের
একজোড়া বাহু পরস্পর সমান্তরাল থাকে।}
\end{quote}

\textbf{আয়তক্ষেত্রের সন্নিহিত বাহুদ্বয় সমান হলে তাকে কি বলে?}

\textbf{\hl{ক) বর্গক্ষেত্র ✔}}

খ) রম্বস

গ) সামন্তরিক

ঘ) ট্রাপিজিয়াম

\textbf{\hl{✔ সঠিক উত্তর: ক) বর্গক্ষেত্র}}

\begin{quote}
\textbf{ব্যাখ্যা:} \emph{যে আয়তক্ষেত্রের সন্নিহিত বাহুদ্বয় পরস্পর সমান, তাকে
বর্গক্ষেত্র বলে। বর্গক্ষেত্রের চারটি বাহুই সমান এবং প্রতিটি কোণ সমকোণ।}
\end{quote}

\textbf{কোনো চতুর্ভুজের বিপরীত কৌণিক বিন্দুর সংযোজন রেখাংশ দুটির প্রত্যেকটিকে
বলে-\/-}

ক) ভূমি

\textbf{\hl{খ) কর্ণ ✔}}

গ) মধ্যমা

ঘ) উচ্চতা

\textbf{\hl{✔ সঠিক উত্তর: খ) কর্ণ}}

\begin{quote}
\textbf{ব্যাখ্যা:} \emph{চতুর্ভুজের যেকোনো দুটি বিপরীত কৌণিক বিন্দু যুক্ত করলে যে
রেখাংশ পাওয়া যায়, তাকে ওই চতুর্ভুজের কর্ণ বলা হয়।}
\end{quote}

\textbf{কোন চতুর্ভুজের কর্ণদ্বয় সমান ও পরস্পর সমকোণে সমদ্বিখন্ডিত। এটি কোন ধরনের
চতুর্ভুজ হবে?}

ক) রম্বস

\textbf{\hl{খ) বর্গক্ষেত্র ✔}}

গ) আয়তক্ষেত্র

ঘ) ট্রাপিজিয়াম

\textbf{\hl{✔ সঠিক উত্তর: খ) বর্গক্ষেত্র}}

\begin{quote}
\textbf{ব্যাখ্যা:} \emph{বর্গক্ষেত্রের কর্ণদ্বয় পরস্পরের সমান এবং তারা পরস্পরকে
সমকোণে ( 90\^{}° ) সমদ্বিখণ্ডিত করে।}
\end{quote}

\textbf{চতুর্ভুজের চার কোণের সমষ্টি কত?}

ক) দুই সমকোণ

খ) তিন সমকোণ

গ) তিন সমকোণের বেশি কিন্তু চার সমকোণের কম

\textbf{\hl{ঘ) চার সমকোণ ✔}}

\textbf{\hl{✔ সঠিক উত্তর: ঘ) চার সমকোণ}}

\begin{quote}
\textbf{ব্যাখ্যা:} \emph{চতুর্ভুজের চার কোণের সমষ্টি = 360°}

\emph{⇒ 4×90° = চার সমকোণ}
\end{quote}

\textbf{একটি রম্বস আঁকতে হলে কমপক্ষে কোন উপাত্তগুলোর প্রয়োজন?}

ক) দুটি বিপরীত বাহু

খ) দুটি বিপরীত কোণ

গ) কর্ণের দৈর্ঘ্য

\textbf{\hl{ঘ) এক বাহু ও এক কোণ ✔}}

\textbf{\hl{✔ সঠিক উত্তর: ঘ) এক বাহু ও এক কোণ}}

\begin{quote}
\textbf{ব্যাখ্যা:} \emph{একটি সুনির্দিষ্ট রম্বস অঙ্কনের জন্য ন্যূনতম এক বাহু ও এক
কোণের মান জানা প্রয়োজন।}
\end{quote}

\textbf{সামন্তরিকের বিপরীত কোণের অন্তর্দ্বিখন্ডদ্বয়-\/-}

\textbf{\hl{ক) পরস্পর সমান ✔}}

খ) পরস্পর সমান্তরাল

গ) পরস্পরের উপর লম্ব

ঘ) পরস্পরের একটি বিন্দুতে ছেদ করে

\textbf{\hl{✔ সঠিক উত্তর: ক) পরস্পর সমান}}

\begin{quote}
\textbf{ব্যাখ্যা:} \emph{সামন্তরিকের বিপরীত কোণগুলোর অন্তর্দ্বিখণ্ডকদ্বয় পরস্পর
সমান হয়।}
\end{quote}

\textbf{একটি বর্গক্ষেত্রের এক বাহুর দৈর্ঘ্য a মিটার হলে, এর ক্ষেত্রফল কত?}

ক) a×a মিটার

খ) a\^{}3 মিটার

গ) a মিটার

\textbf{\hl{ঘ) a\^{}2 বর্গমিটার ✔}}

\textbf{\hl{✔ সঠিক উত্তর: ঘ) a\^{}2 বর্গমিটার}}

\begin{quote}
\textbf{ব্যাখ্যা:} \emph{বর্গক্ষেত্রের ক্ষেত্রফল = ( বাহু)\^{}2}

\emph{∴ এক বাহুর দৈর্ঘ্য a মিটার হলে ক্ষেত্রফল a\^{}2 বর্গমিটার।}
\end{quote}

\textbf{রম্বসের কর্ণদ্বয় পরস্পর "O" বিন্দুতে ছেদ করেছে। কর্ণদ্বয়ের অন্তর্ভূক্ত কোণ}

ক) সূক্ষ্মকোণ

খ) স্থূলকোণ

গ) সরলকোণ

\textbf{\hl{ঘ) সমকোণ ✔}}

\textbf{\hl{✔ সঠিক উত্তর: ঘ) সমকোণ}}

\begin{quote}
\textbf{ব্যাখ্যা:} \emph{রম্বসের কর্ণদ্বয় পরস্পরকে সমকোণে ( 90\^{}° )
সমদ্বিখণ্ডিত করে।}
\end{quote}

\textbf{কোনো ট্রাপিজিয়ামের ক্ষেত্রফল ৪০ বর্গমিটার এবং সমান্তরাল বাহু দুইটির
মধ্যবর্তী লম্ব দূরত্ব ৮ মিটার। একটি বাহুর দৈর্ঘ্য ৬ মিটার হলে অপর বাহুটির দৈর্ঘ্য
কত?}

\textbf{\hl{ক) ৪ মিটার ✔}}

খ) ৬ মিটার

গ) ৮ মিটার

ঘ) ১০ মিটার

\textbf{\hl{✔ সঠিক উত্তর: ক) ৪ মিটার}}

\begin{quote}
\textbf{ব্যাখ্যা:} \emph{ট্রাপিজিয়ামের ক্ষেত্রফল = (1/2)×(সমান্তরাল বাহুদ্বয়ের
যোগফল)×লম্ব দূরত্ব}

\emph{প্রশ্নমতে, (1/2)×(6+অপর বাহু)×8=40}

\emph{⇒ 4×(6+অপর বাহু)=40}

\emph{⇒ 6+অপর বাহু=10}

\emph{∴ অপর বাহু = 10−6=4 মিটার।}
\end{quote}

\textbf{একটি ট্রাপিজিয়ামের সমান্তরাল বাহু দুইটির দৈর্ঘ্য ৫ মিটার ও ৭ মিটার। এর
ক্ষেত্রফল ৪৮ বর্গমিটার হলে, বাহু দুইটির মধ্যবর্তী লম্ব দূরত্ব কত হবে?}

ক) ৬ মিটার

\textbf{\hl{খ) ৮ মিটার ✔}}

গ) ৯ মিটার

ঘ) ১০ মিটার

\textbf{\hl{✔ সঠিক উত্তর: খ) ৮ মিটার}}

\begin{quote}
\textbf{ব্যাখ্যা:} \emph{ট্রাপিজিয়ামের ক্ষেত্রফল = (1/2)×(সমান্তরাল বাহুদ্বয়ের
যোগফল)×লম্ব দূরত্ব}

\emph{প্রশ্নমতে, (1/2)×(5+7)×লম্ব দূরত্ব=48}

\emph{⇒ 6×লম্ব দূরত্ব=48}

\emph{∴ লম্ব দূরত্ব = (48/6)=8 মিটার।}
\end{quote}

\textbf{একটি রম্বস ক্ষেত্রের কর্ণ যথাক্রমে ৫ সে.মি. ও ৪.৫ সে.মি.। এর ক্ষেত্রফল
কত বর্গ সে.মি.?}

ক) ২.২৫

খ) ২২.৫০

গ) ১২.৫০

\textbf{\hl{ঘ) ১১.২৫ ✔}}

\textbf{\hl{✔ সঠিক উত্তর: ঘ) ১১.২৫}}

\begin{quote}
\textbf{ব্যাখ্যা:} \emph{রম্বসের ক্ষেত্রফল = (1/2)×(কর্ণদ্বয়ের গুণফল)}

\emph{⇒ ক্ষেত্রফল = (1/2)×5×4.5=(22.5/2)=11.25 বর্গ সে.মি.।}
\end{quote}

\textbf{একটি ট্রাপিজিয়ামের ক্ষেত্রফল ৬০ বর্গমিটার এবং সমান্তরাল বাহু দুইটির
মধ্যবর্তী লম্ব দূরত্ব ৮ মিটার। একটি বাহুর দৈর্ঘ্য ৬ মিটার হলে অপর বাহুটির দৈর্ঘ্য
হবে-\/-\/-}

ক) ৭ মি.

খ) ৮ মি.

\textbf{\hl{গ) ৯ মি. ✔}}

ঘ) ১২ মি.

\textbf{\hl{✔ সঠিক উত্তর: গ) ৯ মি.}}

\begin{quote}
\textbf{ব্যাখ্যা:} \emph{ট্রাপিজিয়ামের ক্ষেত্রফল = (1/2)×(সমান্তরাল বাহুদ্বয়ের
যোগফল)×লম্ব দূরত্ব}

\emph{প্রশ্নমতে, (1/2)×(6+অপর বাহু)×8=60}

\emph{⇒ 4×(6+অপর বাহু)=60}

\emph{⇒ 6+অপর বাহু=15}

\emph{∴ অপর বাহু = 15−6=9 মিটার।}
\end{quote}

\textbf{একটি চতুর্ভুজ আকার জন্য নিচের কোন উপাত্তগুলো প্রয়োজন?}

ক) ২টি বাহু ও ২টি কোণ

খ) ২টি কর্ণের খণ্ডিত অংশসমূহ ও ১টি বাহু

গ) ৪টি বাহু ও ২টি কোণ

\textbf{\hl{ঘ) ৩টি বাহু ও ২টি কর্ণ ✔}}

\textbf{\hl{✔ সঠিক উত্তর: ঘ) ৩টি বাহু ও ২টি কর্ণ}}

\begin{quote}
\textbf{ব্যাখ্যা:} \emph{চতুর্ভুজ অঙ্কন করতে নিম্নোক্ত উপাত্তগুলো জানা প্রয়োজন: (১)
চারটি বাহু ও একটি কোণ, (২) চারটি বাহু ও একটি কর্ণ, (৩) তিনটি বাহু ও দুইটি
কর্ণ, (৪) তিনটি বাহু ও তাদের অন্তর্ভুক্ত কোণ, (৫) দুইটি বাহু ও তিনটি কোণ, (৬)
দুইটি কর্ণের খণ্ডিত অংশসমূহ ও কর্ণ দুটির অন্তর্ভুক্ত একটি কোণ।}
\end{quote}

\textbf{একটি চতুর্ভুজের তিন কোণের সমষ্টি 280° । চতুর্থ কোটির মান কত?}

ক) 90°

খ) 70°

\textbf{\hl{গ) 80° ✔}}

ঘ) 60°

\textbf{\hl{✔ সঠিক উত্তর: গ) 80°}}

\begin{quote}
\textbf{ব্যাখ্যা:} \emph{চতুর্ভুজের চারটি কোণের সমষ্টি 360° ।}

\emph{∴ চতুর্থ কোণের মান = 360° − 280° = 80° ।}
\end{quote}

\textbf{চতুর্ভুজের চার কোণের অনুপাত 1:2:2:3 হলে বৃহত্তম কোণের পরিমাণ হবে?}

ক) 100°

খ) 115°

\textbf{\hl{গ) 135° ✔}}

ঘ) 225°

\textbf{\hl{✔ সঠিক উত্তর: গ) 135°}}

\begin{quote}
\textbf{ব্যাখ্যা:} \emph{অনুপাতগুলোর সমষ্টি = 1+2+2+3 = 8}

\emph{চতুর্ভুজের চারকোণের সমষ্টি = 360°}

\emph{∴ বৃহত্তম কোণের পরিমাণ = (3/8)×360° = 3×45° = 135°}
\end{quote}

\textbf{চতুর্ভুজের চার কোণের অনুপাত 1:2:2:3 হলে ক্ষুদ্রতম কোণের পরিমাণ কত
ডিগ্রি?}

ক) 60°

খ) 55°

\textbf{\hl{গ) 45° ✔}}

ঘ) 35°

\textbf{\hl{✔ সঠিক উত্তর: গ) 45°}}

\begin{quote}
\textbf{ব্যাখ্যা:} \emph{অনুপাতের যোগফল = 1+2+2+3 = 8}

\emph{চতুর্ভুজের চার কোণের সমষ্টি = 360°}

\emph{ক্ষুদ্রতম কোণের পরিমাণ = (360°×1/8) = 45°}
\end{quote}

\textbf{একটি চতুর্ভুজের দৈর্ঘ্য ৮ মিটার, প্রস্থ ৬ মিটার। চতুর্ভুজটির পরিসীমা কত?}

ক) ১৪ মিটার

খ) ৪৮ মিটার

গ) ২২ মিটার

\textbf{\hl{ঘ) ২৮ মিটার ✔}}

\textbf{\hl{✔ সঠিক উত্তর: ঘ) ২৮ মিটার}}

\begin{quote}
\textbf{ব্যাখ্যা:} \emph{চতুর্ভুজটির পরিসীমা = 2(দৈর্ঘ্য+প্রস্থ)}

\emph{= 2(8+6) মিটার}

\emph{= 2×14 মিটার}

\emph{= ২৮ মিটার।}
\end{quote}

\textbf{একজন লোক A অবস্থান থেকে হেঁটে ডান দিকে ১০ ফুট, অত:পর বামদিকে ২০ ফুট,
তারপর বামদিকে ২০ ফুট এবং সবশেষে বামদিকে ২০ ফুট গিয়ে B অবস্থানে পৌছিল। A ও B
এর মধ্যকার দূরত্ব কত ফুট?}

ক) ৩০ ফুট

খ) ৪০ ফুট

\textbf{\hl{গ) ১০ ফুট ✔}}

ঘ) ২০ ফুট

\textbf{\hl{✔ সঠিক উত্তর: গ) ১০ ফুট}}

\begin{quote}
\textbf{ব্যাখ্যা:} \emph{চিত্র হতে, BC=20 ফুট, AC=10 ফুট}

\emph{⇒ BA=20−10=10 ফুট।}
\end{quote}

\textbf{ABCD চতুর্ভুজে AB∥CD, AC=BD এবং ∠A=90° হলে, সঠিক চতুর্ভুজ কোনটি?}

ক) সামান্তরিক

খ) রম্বস

গ) ট্রাপিজিয়াম

\textbf{\hl{ঘ) আয়তক্ষেত্র ✔}}

\textbf{\hl{✔ সঠিক উত্তর: ঘ) আয়তক্ষেত্র}}

\begin{quote}
\textbf{ব্যাখ্যা:} \emph{যেহেতু বিপরীত বাহুগুলো সমান ও সমান্তরাল ( AB∥CD ,
কর্ণদ্বয় সমান AC=BD ) এবং একটি কোণ সমকোণ ( ∠A=90° ), সেহেতু এটি একটি
আয়তক্ষেত্র (বা বর্গক্ষেত্র হতে পারে, অপশনে বর্গক্ষেত্র না থাকায় সঠিক উত্তর
আয়তক্ষেত্র)।}
\end{quote}

\textbf{সামান্তরিকের বিপরীত কোণগুলো -}

ক) ৯০° থেকে বড়

খ) এক সমকোণ

\textbf{\hl{গ) পরস্পর সমান ✔}}

ঘ) সমান নয়

\textbf{\hl{✔ সঠিক উত্তর: গ) পরস্পর সমান}}

\begin{quote}
\textbf{ব্যাখ্যা:} \emph{সামান্তরিকের বিপরীত বাহু ও কোণগুলো পরস্পর সমান হয়ে
থাকে।}
\end{quote}

\textbf{কোনো সামান্তরিকের একটি কোণ সমকোণ হলে উহা একটি}

ক) বর্গ

খ) রম্বস

গ) ট্রাপিজিয়াম

\textbf{\hl{ঘ) আয়তক্ষেত্র ✔}}

\textbf{\hl{✔ সঠিক উত্তর: ঘ) আয়তক্ষেত্র}}

\begin{quote}
\textbf{ব্যাখ্যা:} \emph{যে চতুর্ভুজের বিপরীত বাহুদ্বয় পরস্পর সমান ও সমান্তরাল
এবং প্রত্যেকটি কোণ এক সমকোণের সমান, তাকে আয়তক্ষেত্র (Rectangle) বলে।}
\end{quote}

\textbf{ABCD সামান্তরিকের B কোণ ১০০° হলে C কোণের মান কত?}

ক) ১০০°

খ) ৯০°

\textbf{\hl{গ) ৮০° ✔}}

ঘ) ৬০°

\textbf{\hl{✔ সঠিক উত্তর: গ) ৮০°}}

\begin{quote}
\textbf{ব্যাখ্যা:} \emph{সামান্তরিকের বিপরীত কোণগুলো সমান বলে
∠B+∠D=100°+100°=200°}

\emph{∴ ∠C+∠A=360°−200°=160°}

\emph{এখন ∠C=∠A বলে ∠C=(160°/2)=80° হবে।}
\end{quote}

\textbf{PQRS সামান্তরিকের ∠P=100° তাহলে ∠Q এর মান কত?}

ক) 120°

খ) 100°

গ) 90°

\textbf{\hl{ঘ) 80° ✔}}

\textbf{\hl{✔ সঠিক উত্তর: ঘ) 80°}}

\begin{quote}
\textbf{ব্যাখ্যা:} \emph{সামান্তরিকের যেকোনো দুই সংলগ্ন কোণের সমষ্টি 180° ।}

\emph{∴ ∠P+∠Q=180°}

\emph{⇒ 100°+∠Q=180°}

\emph{∴ ∠Q=180°−100°=80°}
\end{quote}

\textbf{কোনো সামান্তরিকের দুটি সন্নিহিত কোণের একটি ১১৫ ডিগ্রি হলে অপর কোণটি
কত ডিগ্রি হবে?}

ক) ৪৫ ডিগ্রি

খ) ৫৫ ডিগ্রি

\textbf{\hl{গ) ৬৫ ডিগ্রি ✔}}

ঘ) ৩৫ ডিগ্রি

\textbf{\hl{✔ সঠিক উত্তর: গ) ৬৫ ডিগ্রি}}

\begin{quote}
\textbf{ব্যাখ্যা:} \emph{সামান্তরিকের দুটি সন্নিহিত কোণের সমষ্টি 180° ।}

\emph{∴ অপর কোণটি = 180°−115°=65° ।}
\end{quote}

\textbf{একটি সামান্তরিকের দুটি সন্নিহিত বাহুর দৈর্ঘ্য যথাক্রমে ৭ সেন্টিমিটার এবং
৫ সেন্টিমিটার হলে, এর পরিসীমা কত?}

\textbf{\hl{ক) ২৪ ✔}}

খ) ১২

গ) ২৮

ঘ) ২০

\textbf{\hl{✔ সঠিক উত্তর: ক) ২৪}}

\begin{quote}
\textbf{ব্যাখ্যা:} \emph{সামান্তরিকের পরিসীমা = 2(দৈর্ঘ্য+প্রস্থ)}

\emph{⇒ 2×(7+5) সেন্টিমিটার}

\emph{⇒ 2×12 সেন্টিমিটার}

\emph{⇒ ২৪ সেন্টিমিটার।}
\end{quote}

\textbf{একটি সামান্তরিকের বিপরীত কোণের সমষ্টি ৬০ ডিগ্রি হলে অপর একটি কোণের
মান কত?}

ক) ১২০°

খ) ১৬৫°

\textbf{\hl{গ) ১৫০° ✔}}

ঘ) ১৪০°

\textbf{\hl{✔ সঠিক উত্তর: গ) ১৫০°}}

\begin{quote}
\textbf{ব্যাখ্যা:} \emph{সামান্তরিকের চার কোণের সমষ্টি 360° ।}

\emph{এখানে ২টি বিপরীত কোণের সমষ্টি 60° হলে অপর ২টি কোণের সমষ্টি =
360°−60°=300°}

\emph{∴ ১টি কোণের মান হবে = (300°/2)=150° ।}
\end{quote}

\textbf{ABCD সামান্তরিকের ∠BCD=130° হলে, ∠ABC এর মান কত?}

ক) 40°

খ) 60°

\textbf{\hl{গ) 50° ✔}}

ঘ) 130°

\textbf{\hl{✔ সঠিক উত্তর: গ) 50°}}

\begin{quote}
\textbf{ব্যাখ্যা:} \emph{সামান্তরিকের যেকোনো দুই সন্নিহিত কোণের সমষ্টি 180° ।}

\emph{∴ ∠ABC+∠BCD=180°}

\emph{⇒ ∠ABC+130°=180°}

\emph{∴ ∠ABC=180°−130°=50°}
\end{quote}

\textbf{যে চতুর্ভুজের বিপরীত বাহুগুলো পরস্পর সমান ও সমান্তরাল, তাকে বলে-}

\textbf{\hl{ক) সামান্তরিক (Parallelogram) ✔}}

খ) রম্বস (Rhombus)

গ) ট্রাপিজিয়াম (Trapezium)

ঘ) আয়ত চতুর্ভুজ

\textbf{\hl{✔ সঠিক উত্তর: ক) সামান্তরিক (Parallelogram)}}

\begin{quote}
\textbf{ব্যাখ্যা:} \emph{সামান্তরিকের বিপরীত বাহুগুলো পরস্পর সমান ও সমান্তরাল।
রম্বসের ৪টি বাহুই পরস্পর সমান এবং বিপরীত বাহুগুলো সমান্তরাল। ট্রাপিজিয়ামের দুটি
বিপরীত বাহু সমান্তরাল কিন্তু সমান নয়।}
\end{quote}

\textbf{কোন ক্ষেত্রটি সামান্তরিক নয়?}

ক) আয়তক্ষেত্র

\textbf{\hl{খ) ট্রাপিজিয়াম ✔}}

গ) বর্গক্ষেত্র

ঘ) রম্বস

\textbf{\hl{✔ সঠিক উত্তর: খ) ট্রাপিজিয়াম}}

\begin{quote}
\textbf{ব্যাখ্যা:} \emph{আয়তক্ষেত্র, বর্গক্ষেত্র ও রম্বস প্রত্যেকেই এক একটি বিশেষ
ধরনের সামান্তরিক। কিন্তু ট্রাপিজিয়ামের কেবল এক জোড়া বিপরীত বাহু সমান্তরাল
থাকে, তাই এটি সামান্তরিক নয়।}
\end{quote}

\textbf{ABCD সামান্তরিকের DC ভূমিকে E পর্যন্ত বাড়ানো হলো। ∠BAD=100° হলে
∠BCE= কত?}

ক) ৬০°

\textbf{\hl{খ) ৮০° ✔}}

গ) ৪৫°

ঘ) ৬০°

\textbf{\hl{✔ সঠিক উত্তর: খ) ৮০°}}

\begin{quote}
\textbf{ব্যাখ্যা:} \emph{সামান্তরিকের বিপরীত কোণগুলো পরস্পর সমান。}

\emph{⇒ ∠BCD=∠BAD=100°}

\emph{চিত্রমতে, ∠BCE+∠BCD=180°}

\emph{⇒ ∠BCE+100°=180°}

\emph{∴ ∠BCE=80°}
\end{quote}

\textbf{কোনটি সামান্তরিকের ক্ষেত্রফল নির্ণয়ের সূত্র?}

ক) (1/2)(ভূমি×উচ্চতা)

খ) দৈর্ঘ্য×প্রস্থ

গ) 2(দৈর্ঘ্য+প্রস্থ)

\textbf{\hl{ঘ) ভূমি×উচ্চতা ✔}}

\textbf{\hl{✔ সঠিক উত্তর: ঘ) ভূমি×উচ্চতা}}

\begin{quote}
\textbf{ব্যাখ্যা:} \emph{সামান্তরিকের ক্ষেত্রফল নির্ণয়ের সূত্র হলো ভূমি×উচ্চতা ।}
\end{quote}

\textbf{সামান্তরিকের ভূমি a মিটার এবং উচ্চতা h মিটার হলে সামান্তরিকটির
ক্ষেত্রফল কত?}

ক) (1/2)ah ব.মি.

\textbf{\hl{খ) ah ব.মি. ✔}}

গ) a\^{}2h ব. মি.

ঘ) a\^{}2h\^{}2 ব. মি.

\textbf{\hl{✔ সঠিক উত্তর: খ) ah ব.মি.}}

\begin{quote}
\textbf{ব্যাখ্যা:} \emph{সামান্তরিকের ক্ষেত্রফল নির্ণয়ের সূত্র হলো ভূমি×উচ্চতা ।}

\emph{∴ a মিটার ভূমি এবং h মিটার উচ্চতা বিশিষ্ট সামান্তরিকের ক্ষেত্রফল ah
বর্গমিটার।}
\end{quote}

\textbf{কোনটি সামান্তরিকের ক্ষেত্রফল নির্ণয়ের সূত্র?}

ক) (1/2)(ভূমি×উচ্চতা)

খ) দৈর্ঘ্য×প্রস্থ

গ) 2(দৈর্ঘ্য+প্রস্থ)

\textbf{\hl{ঘ) ভূমি×উচ্চতা ✔}}

\textbf{\hl{✔ সঠিক উত্তর: ঘ) ভূমি×উচ্চতা}}

\begin{quote}
\textbf{ব্যাখ্যা:} \emph{সামান্তরিকের ক্ষেত্রফল = ভূমি×উচ্চতা ।}
\end{quote}

\textbf{একটি সামান্তরিকের ভূমি ৪ ফুট এবং উচ্চতা ৩ ফুট হলে উহার ক্ষেত্রফল কত
বর্গফুট?}

ক) ১৪

\textbf{\hl{খ) ১২ ✔}}

গ) ১৮

ঘ) ৭

\textbf{\hl{✔ সঠিক উত্তর: খ) ১২}}

\begin{quote}
\textbf{ব্যাখ্যা:} \emph{আমরা জানি, সামান্তরিকের ক্ষেত্রফল = ভূমি×উচ্চতা⇒
4×3=12 বর্গফুট।}
\end{quote}

\textbf{একটি চতুর্ভুজের চারটি বাহুর মধ্যবিন্দু পরস্পর যুক্ত করলে কীরূপ ক্ষেত্র পাওয়া
যাবে?}

ক) বর্গক্ষেত্র

খ) আয়তক্ষেত্র

গ) ট্রাপিজিয়াম

\textbf{\hl{ঘ) সামান্তরিক ✔}}

\textbf{\hl{✔ সঠিক উত্তর: ঘ) সামান্তরিক}}

\begin{quote}
\textbf{ব্যাখ্যা:} \emph{যেকোনো চতুর্ভুজের চার বাহুর মধ্যবিন্দুগুলো পরপর যুক্ত করলে
একটি সামান্তরিক ক্ষেত্র পাওয়া যায়।}
\end{quote}

\textbf{একটি সামান্তরিকের দুটি সন্নিহিত বাহুর দৈর্ঘ্য যথাক্রমে ৭ সেন্টিমিটার এবং
৫ সেন্টিমিটার হলে, এর পরিসীমা কত?}

\textbf{\hl{ক) ২৪ ✔}}

খ) ১২

গ) ২৮

ঘ) ২০

\textbf{\hl{✔ সঠিক উত্তর: ক) ২৪}}

\begin{quote}
\textbf{ব্যাখ্যা:} \emph{সামান্তরিকের পরিসীমা = 2(দৈর্ঘ্য+প্রস্থ)}

\emph{⇒ 2×(7+5) সেন্টিমিটার}

\emph{⇒ 2×12=24 সেন্টিমিটার।}
\end{quote}

\textbf{একটি সামান্তরিকের বিপরীত কোণের সমষ্টি ৬০ ডিগ্রি হলে অপর একটি কোণের
মান কত?}

ক) ১২০°

খ) ১৬৫°

\textbf{\hl{গ) ১৫০° ✔}}

ঘ) ১৪০°

\textbf{\hl{✔ সঠিক উত্তর: গ) ১৫০°}}

\begin{quote}
\textbf{ব্যাখ্যা:} \emph{সামান্তরিকের চার কোণের সমষ্টি 360° ।}

\emph{এখানে ২টি বিপরীত কোণের সমষ্টি 60° হলে অপর ২টি কোণের সমষ্টি =
360°−60°=300°}

\emph{∴ ১টি কোণের মান হবে = (300°/2)=150° ।}
\end{quote}

\textbf{4a ব্যাস বিশিষ্ট বৃত্তের ক্ষেত্রফল 4a ভূমি বিশিষ্ট আয়তক্ষেত্রের ক্ষেত্রফলের
সমান হলে আয়তক্ষেত্রের উচ্চতা কত?}

\textbf{\hl{ক) πa ✔}}

খ) πa\^{}2

গ) 2πa

ঘ) 2πa\^{}2

\textbf{\hl{✔ সঠিক উত্তর: ক) πa}}

\begin{quote}
\textbf{ব্যাখ্যা:} \emph{বৃত্তের ব্যাস = 4a , সুতরাং ব্যাসার্ধ r=(4a/2)=2a ।}

\emph{বৃত্তের ক্ষেত্রফল = πr\^{}2=π(2a)\^{}2=4πa\^{}2 ।}

\emph{শর্তমতে, আয়তক্ষেত্রের ক্ষেত্রফল = বৃত্তের ক্ষেত্রফল = 4πa\^{}2 ।}

\emph{আয়তক্ষেত্রের ক্ষেত্রফল = ভূমি × উচ্চতা}

\emph{⇒ 4a×উচ্চতা=4πa\^{}2}

\emph{∴ উচ্চতা = (4πa\^{}2/4a)=πa ।}
\end{quote}

\textbf{একটি আয়তক্ষেত্রের কর্ণের দৈর্ঘ্য ১৫ মি. এবং প্রস্থ ১০ মি. হলে
আয়তক্ষেত্রের ক্ষেত্রফল কত বর্গমিটার?}

ক) 35√(5)

খ) 40√(5)

গ) 45√(5)

\textbf{\hl{ঘ) 50√(5) ✔}}

\textbf{\hl{✔ সঠিক উত্তর: ঘ) 50√(5)}}

\begin{quote}
\textbf{ব্যাখ্যা:} \emph{আয়তক্ষেত্রের দৈর্ঘ্য a একক ও প্রস্থ b একক হলে উহার
কর্ণের দৈর্ঘ্য √(a\^{}2+b\^{}2) ।}

\emph{প্রশ্নমতে, √(a\^{}2+10\^{}2)=15}

\emph{⇒ a\^{}2+100=15\^{}2}

\emph{⇒ a\^{}2=225−100}

\emph{⇒ a\^{}2=125}

\emph{∴ a=√(125)=5√(5) মি.}

\emph{∴ ক্ষেত্রফল = (5√(5)×10) বর্গমিটার = 50√(5) বর্গমিটার।}
\end{quote}

\textbf{বৃত্তস্থ সামান্তরিক একটি---}

ক) বর্গক্ষেত্র

খ) ট্রাপিজিয়াম

গ) রম্বস

\textbf{\hl{ঘ) আয়তক্ষেত্র ✔}}

\textbf{\hl{✔ সঠিক উত্তর: ঘ) আয়তক্ষেত্র}}

\begin{quote}
\textbf{ব্যাখ্যা:} \emph{বৃত্তস্থ বা বৃত্তে অন্তলিখিত সামান্তরিক সর্বদা একটি
আয়তক্ষেত্র হয়।}
\end{quote}

\textbf{আয়তক্ষেত্রের চারটি বাহুর মধ্যে---}

ক) ৪টি সমান

খ) পাশাপাশি ২টি সমান

\textbf{\hl{গ) বিপরীত দুইবাহু সমান ✔}}

ঘ) সবগুলি অসমান

\textbf{\hl{✔ সঠিক উত্তর: গ) বিপরীত দুইবাহু সমান}}

\begin{quote}
\textbf{ব্যাখ্যা:} \emph{আয়তক্ষেত্রের বিপরীত বাহুগুলো পরস্পর সমান ও সমান্তরাল
হয়ে থাকে।}
\end{quote}

\textbf{একটি আয়তাকার তামার পাত্রের দৈর্ঘ্য, প্রস্থ, উচ্চতা যথাক্রমে ২০ সেমি, ১৫
সেমি, ১০ সেমি হলে, পাত্রের ভিতরের আয়তন কত?}

\textbf{\hl{ক) ৩,০০০ ঘন সেমি ✔}}

খ) ২,৭০০ ঘন সেমি

গ) ১,৫০০ ঘন সেমি

ঘ) ২,০০০ ঘন সেমি

\textbf{\hl{✔ সঠিক উত্তর: ক) ৩,০০০ ঘন সেমি}}

\begin{quote}
\textbf{ব্যাখ্যা:} \emph{আমরা জানি,}

\emph{পাত্রের আয়তন = দৈর্ঘ্য × প্রস্থ × উচ্চতা}

\emph{⇒ 20×15×10 ঘন সেমি = ৩,০০০ ঘন সেমি।}
\end{quote}

\textbf{একটি আয়তাকার বাগানের ক্ষেত্রফল ৩৬০ বর্গমিটার এবং প্রস্থ ১২ মিটার হলে,
দৈর্ঘ্য কত মিটার?}

\textbf{\hl{ক) ৩০ ✔}}

খ) ৩৫

গ) ৪০

ঘ) ২৫

ঙ) কোনোটিই নয়

\textbf{\hl{✔ সঠিক উত্তর: ক) ৩০}}

\begin{quote}
\textbf{ব্যাখ্যা:} \emph{আয়তাকার বাগানের দৈর্ঘ্য =
(ক্ষেত্রফল/প্রস্থ)=(360/12)=30 মি.}
\end{quote}

\textbf{একটি আয়তাকার বাগানের ক্ষেত্রফল ৮৪০ ব.মি. এবং দৈর্ঘ্য ৪০ মি. হলে প্রস্থ
কত?}

\textbf{\hl{ক) ২১ মি. ✔}}

খ) ২২ মি.

গ) ২৫ মি.

\textbf{\hl{✔ সঠিক উত্তর: ক) ২১ মি.}}

\begin{quote}
\textbf{ব্যাখ্যা:} \emph{বাগানের প্রস্থ = (বাগানের ক্ষেত্রফল/বাগানের
দৈর্ঘ্য)=(840/40)=21 মি.}
\end{quote}

\textbf{একটি আয়তক্ষেত্রের কর্ণের দৈর্ঘ্য 15cm এবং প্রস্থ 5cm হলে উহার
ক্ষেত্রফল---}

ক) 35√(2)cm\^{}2

খ) 40√(2)cm\^{}2

গ) 45√(2)cm\^{}2

\textbf{\hl{ঘ) 50√(2)cm\^{}2 ✔}}

\textbf{\hl{✔ সঠিক উত্তর: ঘ) 50√(2)cm\^{}2}}

\begin{quote}
\textbf{ব্যাখ্যা:} \emph{দৈর্ঘ্য = √(15\^{}2−5\^{}2)=√(225−25)=√(200)=10√(2)
সেমি.}

\emph{∴ ক্ষেত্রফল = (10√(2)×5)cm\^{}2=50√(2)cm\^{}2}
\end{quote}

\textbf{একটি আয়তক্ষেত্রের দুটি সন্নিহিত বাহুর দৈর্ঘ্য যথাক্রমে ৯ ও ১২ সেমি.।
অন্তর্ভুক্ত কোণ কত হলে আয়তটি আঁকা সম্ভব?}

\textbf{\hl{ক) 90° ✔}}

খ) 42°

গ) 45°

ঘ) 60°

\textbf{\hl{✔ সঠিক উত্তর: ক) 90°}}

\begin{quote}
\textbf{ব্যাখ্যা:} \emph{আয়তক্ষেত্রের প্রতিটি কোণ এক সমকোণ বা 90° ।}

\emph{∴ অন্তর্ভুক্ত কোণ 90° হলে আয়তটি আঁকা সম্ভব।}
\end{quote}

\textbf{সামান্তরিকের কর্ণদ্বয় পরস্পর সমান হলে সামান্তরিকটি হবে---}

\textbf{\hl{ক) আয়তক্ষেত্র ✔}}

খ) রম্বস

গ) ট্রাপিজিয়াম

ঘ) বর্গক্ষেত্র

\textbf{\hl{✔ সঠিক উত্তর: ক) আয়তক্ষেত্র}}

\begin{quote}
\textbf{ব্যাখ্যা:} \emph{সামান্তরিকের কর্ণদ্বয় পরস্পর সমান হলে তাকে আয়তক্ষেত্র
বলে।}
\end{quote}

\textbf{একটি আয়তক্ষেত্রের দৈর্ঘ্য \textquotesingle a\textquotesingle{} একক
এবং প্রস্থ \textquotesingle b\textquotesingle{} একক হলে, কর্ণের দৈর্ঘ্য কত?}

ক) √(a+b)

\textbf{\hl{খ) √(a\^{}2+b\^{}2) ✔}}

গ) a\^{}2b\^{}2

ঘ) ab

\textbf{\hl{✔ সঠিক উত্তর: খ) √(a\^{}2+b\^{}2)}}

\begin{quote}
\textbf{ব্যাখ্যা:} \emph{আয়তক্ষেত্রের দৈর্ঘ্য a একক ও প্রস্থ b একক হলে
পিথাগোরাসের উপপাদ্য অনুসারে কর্ণের দৈর্ঘ্য হবে √(a\^{}2+b\^{}2) একক।}
\end{quote}

\textbf{আয়তক্ষেত্রের সন্নিহিত বাহুর মধ্যবিন্দুসমূহের সংযোগে যে চতুর্ভুজ হয়, তা
একটি}

\textbf{\hl{ক) রম্বস ✔}}

খ) বর্গক্ষেত্র

গ) সামান্তরিক

ঘ) আয়তক্ষেত্র

\textbf{\hl{✔ সঠিক উত্তর: ক) রম্বস}}

\begin{quote}
\textbf{ব্যাখ্যা:} \emph{আয়তক্ষেত্রের সন্নিহিত বাহুগুলোর মধ্যবিন্দুগুলো পরস্পর যুক্ত
করলে একটি রম্বস উৎপন্ন হয়।}
\end{quote}

\textbf{একখণ্ড জমির দৈর্ঘ্য ৯০ ফুট এবং প্রস্থ ৮০ ফুট হলে, ঐ জমির ক্ষেত্রফল হবে?}

ক) ১ কাঠা

\textbf{\hl{খ) ১০ কাঠা ✔}}

গ) ৭.২ কাঠা

ঘ) ৭২ কাঠা

\textbf{\hl{✔ সঠিক উত্তর: খ) ১০ কাঠা}}

\begin{quote}
\textbf{ব্যাখ্যা:} \emph{দেয়া আছে, দৈর্ঘ্য = ৯০ ফুট, প্রস্থ = ৮০ ফুট}

\emph{∴ ক্ষেত্রফল = 90×80 বর্গ ফুট = ৭২০০ বর্গ ফুট}

\emph{আমরা জানি, ৭২০ বর্গফুট = ১ কাঠা}

\emph{∴ জমির ক্ষেত্রফল = (7200/720) কাঠা = ১০ কাঠা।}
\end{quote}

\textbf{কোনো বাগানের দৈর্ঘ্য ২০ মিটার এবং প্রস্থ ১০ মিটার হলে ঐ বাগানের
ক্ষেত্রফল কত বর্গমিটার?}

ক) ৩০

খ) ২

গ) ৩০০

\textbf{\hl{ঘ) ২০০ ✔}}

\textbf{\hl{✔ সঠিক উত্তর: ঘ) ২০০}}

\begin{quote}
\textbf{ব্যাখ্যা:} \emph{বাগানের ক্ষেত্রফল = দৈর্ঘ্য × প্রস্থ}

\emph{⇒ 20×10=200 বর্গমিটার।}
\end{quote}

\textbf{একটি আয়তক্ষেত্রের দৈর্ঘ্য প্রস্থ অপেক্ষা ৪ ফুট বড়। দৈর্ঘ্য ৮ ফুট হলে
ক্ষেত্রফল কত বর্গফুট?}

ক) ৩০

খ) ৪০

গ) ৩৫

\textbf{\hl{ঘ) ৩২ ✔}}

\textbf{\hl{✔ সঠিক উত্তর: ঘ) ৩২}}

\begin{quote}
\textbf{ব্যাখ্যা:} \emph{দেয়া আছে, দৈর্ঘ্য = ৮ ফুট}

\emph{প্রশ্নমতে, দৈর্ঘ্য = প্রস্থ + ৪ ফুট}

\emph{⇒ ৮ = প্রস্থ + ৪}

\emph{∴ প্রস্থ = 8−4=4 ফুট}

\emph{∴ ক্ষেত্রফল = দৈর্ঘ্য × প্রস্থ = 8×4=32 বর্গফুট।}
\end{quote}

\textbf{একটি আয়তক্ষেত্রের দৈর্ঘ্য প্রস্থের দ্বিগুণ। ক্ষেত্রফল ৩২ বর্গমিটার হলে বড়
বাহুর দৈর্ঘ্য---}

\textbf{\hl{ক) ৮ মিটার ✔}}

খ) ৮ মিটার

গ) ১৬ মিটার

ঘ) ১০ মিটার

\textbf{\hl{✔ সঠিক উত্তর: ক) ৮ মিটার}}

\begin{quote}
\textbf{ব্যাখ্যা:} \emph{ধরি, প্রস্থ x মিটার, সুতরাং দৈর্ঘ্য 2x মিটার।}

\emph{প্রশ্নমতে, 2x×x=32}

\emph{⇒ 2x\^{}2=32}

\emph{⇒ x\^{}2=16}

\emph{∴ x=4}

\emph{∴ বড় বাহুর দৈর্ঘ্য (দৈর্ঘ্য) = 4×2=8 মিটার।}
\end{quote}

\textbf{একটি আয়তক্ষেত্রের দৈর্ঘ্য ও প্রস্থকে দ্বিগুণ করলে ক্ষেত্রফল মূল আয়তক্ষেত্রের
ক্ষেত্রফলের---}

ক) দ্বিগুণ হবে

\textbf{\hl{খ) চারগুণ হবে ✔}}

গ) ছয়গুণ হবে

ঘ) দশগুণ হবে

\textbf{\hl{✔ সঠিক উত্তর: খ) চারগুণ হবে}}

\begin{quote}
\textbf{ব্যাখ্যা:} \emph{ধরি, দৈর্ঘ্য =x ও প্রস্থ =y ।}

\emph{মূল ক্ষেত্রফল =xy}

\emph{দৈর্ঘ্য ও প্রস্থ দ্বিগুণ করলে বৰ্ধিত দৈর্ঘ্য =2x এবং বৰ্ধিত প্রস্থ =2y ।}

\emph{∴ বৰ্ধিত ক্ষেত্রফল =2x×2y=4xy=4×মূল ক্ষেত্রফল ।}
\end{quote}

\textbf{একটি আয়তকার ঘরের প্রস্থ তার দৈর্ঘ্যের দুই-তৃতীয়াংশ। ঘরটির পরিসীমা ৪০
মিটার হলে তার ক্ষেত্রফল কত?}

ক) ৩০ বর্গমিটার

\textbf{\hl{খ) ৯৬ বর্গমিটার ✔}}

গ) ৭২ বর্গমিটার

ঘ) ৬৪ বর্গমিটার

\textbf{\hl{✔ সঠিক উত্তর: খ) ৯৬ বর্গমিটার}}

\begin{quote}
\textbf{ব্যাখ্যা:} \emph{ধরি, দৈর্ঘ্য = x মিটার।}

\emph{প্রস্থ = (2x/3) মিটার।}

\emph{পরিসীমা = 2(x+(2x/3))=40}

\emph{⇒ 2((3x+2x/3))=40}

\emph{⇒ (10x/3)=40}

\emph{⇒ 10x=120}

\emph{∴ x=12 মিটার (দৈর্ঘ্য)}

\emph{প্রস্থ = (2×12/3)=8 মিটার।}

\emph{∴ ক্ষেত্রফল = দৈর্ঘ্য × প্রস্থ = 12×8=96 বর্গমিটার}
\end{quote}

\textbf{4a একক ব্যাসবিশিষ্ট বৃত্তের ক্ষেত্রফল, 4a একক প্রস্থবিশিষ্ট আয়তক্ষেত্রের
ক্ষেত্রফলের সমান হলে আয়তক্ষেত্রের দৈর্ঘ্য কত?}

\textbf{\hl{ক) πaএকক ✔}}

খ) 2πaএকক

গ) πa\^{}2একক

ঘ) πa\^{}2aএকক

\textbf{\hl{✔ সঠিক উত্তর: ক) πaএকক}}

\begin{quote}
\textbf{ব্যাখ্যা:} \emph{বৃত্তের ব্যাসার্ধ r=(4a/2)=2a একক।}

\emph{বৃত্তের ক্ষেত্রফল =π(2a)\^{}2=4πa\^{}2 বর্গএকক।}

\emph{আয়তক্ষেত্রের প্রস্থ =4a একক।}

\emph{শর্তমতে, আয়তক্ষেত্রের ক্ষেত্রফল =4πa\^{}2}

\emph{⇒ দৈর্ঘ্য × প্রস্থ =4πa\^{}2}

\emph{⇒ দৈর্ঘ্য ×4a=4πa\^{}2}

\emph{∴ দৈর্ঘ্য =(4πa\^{}2/4a)=πa একক।}
\end{quote}

\textbf{একটি আয়তক্ষেত্রের প্রস্থ অপেক্ষা দৈর্ঘ্য ৪ মিটার বেশি। পরিসীমা ৫৬ মিটার
হলে ক্ষেত্রফল কত?}

ক) ১৮২ বর্গ মিটার

খ) ১৮৬ বর্গ মিটার

\textbf{\hl{গ) ১৯২ বর্গ মিটার ✔}}

ঘ) ১৯৬ বর্গ মিটার

\textbf{\hl{✔ সঠিক উত্তর: গ) ১৯২ বর্গ মিটার}}

\begin{quote}
\textbf{ব্যাখ্যা:} \emph{ধরি, আয়তক্ষেত্রের দৈর্ঘ্য x মিটার}

\emph{∴ আয়তক্ষেত্রের প্রস্থ (x−4) মিটার}

\emph{শর্তমতে, 2(x+x−4)=56}

\emph{⇒ 2x−4=28}

\emph{⇒ 2x=32}

\emph{∴ x=16 মিটার}

\emph{প্রস্থ = 16−4=12 মিটার}

\emph{∴ ক্ষেত্রফল = 16×12=192 বর্গ মিটার}
\end{quote}

\textbf{একটি আয়তকার ঘরের দৈর্ঘ্য প্রস্থ অপেক্ষা ৪ মিটার বেশি। ঘরটির পরিসীমা
৩২ মিটার হলে ঘরটির দৈর্ঘ্য কত?}

ক) ৬ মিটার

\textbf{\hl{খ) ১০ মিটার ✔}}

গ) ১৮ মিটার

ঘ) ১২ মিটার

\textbf{\hl{✔ সঠিক উত্তর: খ) ১০ মিটার}}

\begin{quote}
\textbf{ব্যাখ্যা:} \emph{ধরি, প্রস্থ = x মিটার, দৈর্ঘ্য = (x+4) মিটার}

\emph{পরিসীমা = 2(x+4+x)=32}

\emph{⇒ 2x+4=16}

\emph{⇒ 2x=12}

\emph{∴ x=6 মিটার (প্রস্থ)}

\emph{∴ দৈর্ঘ্য = 6+4=10 মিটার।}
\end{quote}

\textbf{একটি আয়তক্ষেত্রের দৈর্ঘ্য প্রস্থের দেড় গুণ। ক্ষেত্রফল ৯৬ বর্গমিটার হলে
ক্ষেত্রটির দৈর্ঘ্য কত?}

ক) ৬ মিটার

খ) ৯ মিটার

\textbf{\hl{গ) ১২ মিটার ✔}}

ঘ) ১৫ মিটার

\textbf{\hl{✔ সঠিক উত্তর: গ) ১২ মিটার}}

\begin{quote}
\textbf{ব্যাখ্যা:} \emph{ধরি, প্রস্থ = 2a , দৈর্ঘ্য = 2a×(3/2)=3a}

\emph{ক্ষেত্রফল = 2a×3a=6a\^{}2}

\emph{শর্তমতে, 6a\^{}2=96}

\emph{⇒ a\^{}2=16}

\emph{∴ a=4}

\emph{∴ ক্ষেত্রটির দৈর্ঘ্য = 3×4=12 মিটার।}
\end{quote}

\textbf{একটি আয়তক্ষেত্রের পরিসীমা ১৬০ মিটার। আয়তক্ষেত্রটির প্রস্থ দৈর্ঘ্যের ৩/৫
গুণ। আয়তক্ষেত্রটির দৈর্ঘ্য কত মিটার?}

ক) ৪৫

\textbf{\hl{খ) ৫০ ✔}}

গ) ৬৫

ঘ) কোনோটিই নয়

\textbf{\hl{✔ সঠিক উত্তর: খ) ৫০}}

\begin{quote}
\textbf{ব্যাখ্যা:} \emph{মনে করি, দৈর্ঘ্য = 5x}

\emph{প্রস্থ = 5xএর(3/5)=3x}

\emph{পরিসীমা = 2(5x+3x)=160}

\emph{⇒ 16x=160}

\emph{∴ x=10}

\emph{∴ দৈর্ঘ্য = 5×10=50 মিটার।}
\end{quote}

\textbf{যে আয়তক্ষেত্রের পরিসীমা ২৪ মিটার এবং ক্ষেত্রফল ৩৫ বর্গমিটার, তার
দৈর্ঘ্য ও প্রস্থ নির্ণয় করুন।}

ক) দৈর্ঘ্য ৫ মিটার এবং প্রস্থ ৭ মিটার

\textbf{\hl{খ) দৈর্ঘ্য ৭ মিটার এবং প্রস্থ ৫ মিটার ✔}}

গ) দৈর্ঘ্য ৮ মিটার এবং প্রস্থ ৪ মিটার

ঘ) দৈর্ঘ্য ৬ মিটার এবং প্রস্থ ৬ মিটার

\textbf{\hl{✔ সঠিক উত্তর: খ) দৈর্ঘ্য ৭ মিটার এবং প্রস্থ ৫ মিটার}}

\begin{quote}
\textbf{ব্যাখ্যা:} \emph{ধরি, আয়তক্ষেত্রের দৈর্ঘ্য x মিটার এবং প্রস্থ y মিটার।}

\emph{১ম শর্তমতে, 2(x+y)=24 ⇒ x+y=12 ⇒ x=12−y}

\emph{২য় শর্তমতে, xy=35}

\emph{⇒ (12−y)y=35}

\emph{⇒ 12y−y\^{}2=35}

\emph{⇒ y\^{}2−12y+35=0}

\emph{⇒ (y−7)(y−5)=0}

\emph{∴ y=7,5}

\emph{যেহেতু দৈর্ঘ্য \textgreater{} প্রস্থ, সুতরাং প্রস্থ y=5 মিটার এবং দৈর্ঘ্য
x=7 মিটার।}
\end{quote}

\textbf{একটি আয়তক্ষেত্রের ভূমি উচ্চতার তিনগুণ। আয়তক্ষেত্রটির পরিসীমা ৬৪ হলে তার
ক্ষেত্রফল কত?}

ক) ২৪

খ) ৬৪

গ) ৯৫

\textbf{\hl{ঘ) ১৯২ ✔}}

\textbf{\hl{✔ সঠিক উত্তর: ঘ) ১৯২}}

\begin{quote}
\textbf{ব্যাখ্যা:} \emph{ধরি, আয়তক্ষেত্রটির উচ্চতা x একক এবং ভূমি 3x একক।}

\emph{প্রশ্নমতে, 2(x+3x)=64}

\emph{⇒ 2×4x=64}

\emph{⇒ 8x=64}

\emph{∴ x=8}

\emph{অতএব, ক্ষেত্রফল = (x×3x)=3x\^{}2=3×8\^{}2=3×64=192 বর্গএকক।}
\end{quote}

\textbf{একটি আয়তকার বাগানের দৈর্ঘ্য ২৩ মিটার এবং প্রস্থ ১৭ মিটার হলে
বাগানটির পরিসীমা কত?}

ক) ৪০ মিটার

\textbf{\hl{খ) ৮০ মিটার ✔}}

গ) ৩৯১ মিটার

ঘ) ৪০০ মিটার

\textbf{\hl{✔ সঠিক উত্তর: খ) ৮০ মিটার}}

\begin{quote}
\textbf{ব্যাখ্যা:} \emph{বাগানটির পরিসীমা = 2(দৈর্ঘ্য+প্রস্থ)=2(23+17) মিটার =
(2×40) মিটার = ৮০ মিটার।}
\end{quote}

\textbf{একটি আয়তকার জমির ক্ষেত্রফল ২৪ একর। দৈর্ঘ্য ও প্রস্থের অনুপাত ৩ : ২ হলে
ঐ জমির পরিসীমা কত?}

\textbf{\hl{ক) দৈর্ঘ্য ও প্রস্থের ভিত্তিতে নির্দিষ্ট হিসাব ✔}}

খ) হিসাব অনুযায়ী পরিসীমা নির্ণেয়

\textbf{\hl{✔ সঠিক উত্তর: ক) দৈর্ঘ্য ও প্রস্থের ভিত্তিতে নির্দিষ্ট হিসাব}}

\begin{quote}
\textbf{ব্যাখ্যা:} \emph{ধরি, দৈর্ঘ্য 3x মিটার এবং প্রস্থ 2x মিটার।}

\emph{১ একর = 4046.86 বর্গমিটার (প্রায়), সুতরাং ২৪ একর =
24×4046.86=97124.64 বর্গমিটার।}

\emph{শর্তমতে, 3x×2x=97124.64}

\emph{⇒ 6x\^{}2=97124.64}

\emph{⇒ x\^{}2=16187.44}

\emph{∴ x≈127.23}

\emph{অতএব পরিসীমা = 2(3x+2x)=10x=1272.3 মিটার।}
\end{quote}

\textbf{একটি বর্গের পরিসীমা একটি আয়তক্ষেত্রের পরিসীমার সমান। আয়তক্ষেত্রটির
দৈর্ঘ্য প্রস্থের ৩ গুণ এবং ক্ষেত্রফল ৭৬৮ বর্গমিটার। বর্গক্ষেত্রটির এক বাহুর দৈর্ঘ্য
কত?}

\textbf{\hl{ক) ৩২ মিটার ✔}}

খ) ১৬ মিটার

গ) ২৪ মিটার

ঘ) ৪৮ মিটার

\textbf{\hl{✔ সঠিক উত্তর: ক) ৩২ মিটার}}

\begin{quote}
\textbf{ব্যাখ্যা:} \emph{ধরি, আয়তক্ষেত্রের প্রস্থ x মিটার}

\emph{∴ দৈর্ঘ্য 3x মিটার}

\emph{শর্তমতে, 3x×x=768}

\emph{⇒ 3x\^{}2=768}

\emph{⇒ x\^{}2=256}

\emph{∴ x=16}

\emph{আয়তক্ষেত্রের পরিসীমা = 2(3x+x)=8x=8×16=128 মিটার}

\emph{বর্গক্ষেত্রের পরিসীমা = আয়তক্ষেত্রের পরিসীমা = ১২৮ মিটার}

\emph{∴ বর্গক্ষেত্রের এক বাহুর দৈর্ঘ্য = (128/4)=32 মিটার।}
\end{quote}

\textbf{একটি আয়তাকার ক্ষেত্রের দৈর্ঘ্য প্রস্থের দ্বিগুণ। এর ক্ষেত্রফল ৩২ বর্গমিটার
হলে তার লম্বা বাহুর দৈর্ঘ্য কত মিটার?}

ক) ২৬

\textbf{\hl{খ) ৮ ✔}}

গ) ৪

ঘ) ২

\textbf{\hl{✔ সঠিক উত্তর: খ) ৮}}

\begin{quote}
\textbf{ব্যাখ্যা:} \emph{ধরি, প্রস্থ = a , দৈর্ঘ্য = 2a}

\emph{ক্ষেত্রফল = 2a×a=2a\^{}2}

\emph{শর্তমতে, 2a\^{}2=32}

\emph{⇒ a\^{}2=16}

\emph{∴ a=4}

\emph{∴ লম্বা বাহুর দৈর্ঘ্য (দৈর্ঘ্য) = 2×4=8 মিটার।}
\end{quote}

\textbf{যদি একটি আয়তক্ষেত্রের দৈর্ঘ্য ৮ মিটার এবং পরিসীমা ২৪ মিটার হয় তবে এর
প্রস্থ কত?}

ক) ৩ মিটার

খ) ২ মিটার

গ) ৬ মিটার

\textbf{\hl{ঘ) ৪ মিটার ✔}}

\textbf{\hl{✔ সঠিক উত্তর: ঘ) ৪ মিটার}}

\begin{quote}
\textbf{ব্যাখ্যা:} \emph{ধরি, প্রস্থ = x মিটার}

\emph{পরিসীমা = 2(দৈর্ঘ্য+প্রস্থ)}

\emph{⇒ 2(8+x)=24}

\emph{⇒ 8+x=12}

\emph{∴ x=12−8=4 মিটার।}
\end{quote}

\textbf{একটি আয়তক্ষেত্রের দৈর্ঘ্য বিস্তারের ৩ গুণ। ক্ষেত্রফল ১২০০ বর্গমিটার হলে
এর পরিসীমা কত?}

ক) ১৩০ মিটার

খ) ১৪০ মিটার

গ) ১৫০ মিটার

\textbf{\hl{ঘ) ১৬০ মিটার ✔}}

\textbf{\hl{✔ সঠিক উত্তর: ঘ) ১৬০ মিটার}}

\begin{quote}
\textbf{ব্যাখ্যা:} \emph{ধরি, আয়তক্ষেত্রের বিস্তার বা প্রস্থ x মিটার।}

\emph{∴ দৈর্ঘ্য 3x মিটার।}

\emph{প্রশ্নমতে, x×3x=1200}

\emph{⇒ 3x\^{}2=1200}

\emph{⇒ x\^{}2=400}

\emph{∴ x=20}

\emph{বিস্তার = ২০ মিটার এবং দৈর্ঘ্য = 3×20=60 মিটার।}

\emph{∴ পরিসীমা = 2(60+20)=2×80=160 মিটার।}
\end{quote}

\textbf{একটি আয়তক্ষেত্রের প্রস্থ অপেক্ষা দৈর্ঘ্য ৮ মিটার বেশি। এর ক্ষেত্রফল ১৯২
বর্গমিটার হলে পরিসীমা কত?}

ক) ৫২ মিটার

\textbf{\hl{খ) ৫৬ মিটার ✔}}

গ) ৫০ মিটার

ঘ) ৬৫ মিটার

\textbf{\hl{✔ সঠিক উত্তর: খ) ৫৬ মিটার}}

\begin{quote}
\textbf{ব্যাখ্যা:} \emph{ধরি, প্রস্থ = x এবং দৈর্ঘ্য = x+4 (যেহেতু দৈর্ঘ্য ৪
মিটার বেশি)}

\emph{প্রশ্নমতে, x(x+4)=192}

\emph{⇒ x\^{}2+4x−192=0}

\emph{⇒ (x−12)(x+16)=0}

\emph{∴ x=12 (প্রস্থ)}

\emph{দৈর্ঘ্য = 12+4=16 মিটার}

\emph{∴ পরিসীমা = 2(16+12)=2×28=56 মিটার।}
\end{quote}

\textbf{একটি আয়তাকার ঘরের দৈর্ঘ্য প্রস্থের দেড় গুণ। তার ক্ষেত্রফল ২১৬ বর্গমিটার
হলে, ইহার পরিসীমা কত?}

ক) ৬৫ মিটার

\textbf{\hl{খ) ৬০ মিটার ✔}}

গ) ১৮ মিটার

ঘ) ৫০মিটার

\textbf{\hl{✔ সঠিক উত্তর: খ) ৬০ মিটার}}

\begin{quote}
\textbf{ব্যাখ্যা:} \emph{মনেকরি, আয়তাকার ঘরের প্রস্থ = 2x}

\emph{আয়তাকার ঘরের দৈর্ঘ্য = 3x}

\emph{∴ ক্ষেত্রফল = (2x×3x)=6x\^{}2 বর্গ মিটার}

\emph{প্রশ্নমতে, 6x\^{}2=216}

\emph{⇒ x\^{}2=36}

\emph{∴ x=6}

\emph{আয়তাকার ঘরের দৈর্ঘ্য = 3×6=18}

\emph{আয়তাকার ঘরের প্রস্থ = 2×6=12}

\emph{∴ পরিসীমা = 2(দৈর্ঘ্য+প্রস্থ)=2(18+12)=2×30=60 মিটার।}
\end{quote}

\textbf{একটি আয়তাকার ঘরের প্রস্থ দৈর্ঘ্যের (2/3) গুণ এবং ক্ষেত্রফল ৩৮৪ বর্গমিটার
হলে ঐ ঘরের পরিসীমা কত?}

ক) ২৫৬ মিটার

খ) ৫৭৬ মিটার

গ) ৪০ মিটার

\textbf{\hl{ঘ) ৮০ মিটার ✔}}

\textbf{\hl{✔ সঠিক উত্তর: ঘ) ৮০ মিটার}}

\begin{quote}
\textbf{ব্যাখ্যা:} \emph{ধরি, দৈর্ঘ্য = 3x এবং প্রস্থ = 2x}

\emph{ক্ষেত্রফল = 3x×2x=6x\^{}2}

\emph{প্রশ্নমতে, 6x\^{}2=384}

\emph{⇒ x\^{}2=64}

\emph{∴ x=8}

\emph{দৈর্ঘ্য = 3×8=24 মিটার এবং প্রস্থ = 2×8=16 মিটার}

\emph{∴ পরিসীমা = 2(24+16)=2×40=80 মিটার।}
\end{quote}

\textbf{একটি আয়তাকার ঘরের দৈর্ঘ্য প্রস্থের (3/2) গুণ। এর ক্ষেত্রফল ১৩.৫
বর্গমিটার হলে, পরিসীমা কত মিটার?}

ক) ৭.৫

খ) ১৩.৫

গ) ৭

\textbf{\hl{ঘ) ১৫ ✔}}

\textbf{\hl{✔ সঠিক উত্তর: ঘ) ১৫}}

\begin{quote}
\textbf{ব্যাখ্যা:} \emph{ধরি, প্রস্থ = x মি.}

\emph{∴ দৈর্ঘ্য = (3x/2) মি.}

\emph{শর্তমতে, x×(3x/2)=13.5}

\emph{⇒ (3x\^{}2/2)=13.5}

\emph{⇒ 3x\^{}2=27}

\emph{⇒ x\^{}2=9}

\emph{∴ x=3}

\emph{প্রস্থ = ৩ মিটার এবং দৈর্ঘ্য = (3×3/2)=4.5 মিটার}

\emph{∴ পরিসীমা = 2(4.5+3)=2×7.5=15 মিটার।}
\end{quote}

\textbf{যথাক্রমে x এবং y এক দৈর্ঘ্য ও প্রস্থবিশিষ্ট একটি আয়তক্ষেত্রের পরিসীমা ৯৬
একক। যদি আয়তক্ষেত্রটির দৈর্ঘ্য প্রস্থের দ্বিগুণ হয়, তবে কোনটি সঠিক?}

ক) x+y=24

\textbf{\hl{খ) x+y=48 ✔}}

গ) x+y=12

ঘ) 4x+2y=96

\textbf{\hl{✔ সঠিক উত্তর: খ) x+y=48}}

\begin{quote}
\textbf{ব্যাখ্যা:} \emph{দেয়া আছে, পরিসীমা = ৯৬ একক}

\emph{আমরা জানি, আয়তক্ষেত্রের পরিসীমা = 2(দৈর্ঘ্য+প্রস্থ)}

\emph{⇒ 2(x+y)=96}

\emph{∴ x+y=48 ।}
\end{quote}

\textbf{একটি আয়তক্ষেত্রের দৈর্ঘ্য প্রস্থের দ্বিগুণ। আয়তক্ষেত্রটির ক্ষেত্রফল ১২৫০
বর্গমিটার হলে এর দৈর্ঘ্য কত?}

\textbf{\hl{ক) ৫০ মিটার ✔}}

খ) ২৫ মিটার

গ) ৭৫ মিটার

ঘ) ৬০ মিটার

\textbf{\hl{✔ সঠিক উত্তর: ক) ৫০ মিটার}}

\begin{quote}
\textbf{ব্যাখ্যা:} \emph{ধরি, প্রস্থ = x মিটার}

\emph{∴ দৈর্ঘ্য = 2x মিটার}

\emph{শর্তমতে, 2x×x=1250}

\emph{⇒ 2x\^{}2=1250}

\emph{⇒ x\^{}2=625}

\emph{∴ x=25 মিটার (প্রস্থ)}

\emph{∴ দৈর্ঘ্য = 2×25=50 মিটার।}
\end{quote}

\textbf{একটি আয়তক্ষেত্রের দৈর্ঘ্য প্রস্থের দেড় গুণ। ক্ষেত্রটির ক্ষেত্রফল ৯৬ বর্গ
মিটার হলে ক্ষেত্রটির দৈর্ঘ্য কত?}

\textbf{\hl{ক) ১২ মিটার ✔}}

খ) ৬ মিটার

গ) ৮ মিটার

ঘ) ১৬ মিটার

\textbf{\hl{✔ সঠিক উত্তর: ক) ১২ মিটার}}

\begin{quote}
\textbf{ব্যাখ্যা:} \emph{ধরি, প্রস্থ = 2a}

\emph{দৈর্ঘ্য = 2a×(3/2)=3a}

\emph{ক্ষেত্রফল = 2a×3a=6a\^{}2}

\emph{শর্তমতে, 6a\^{}2=96}

\emph{⇒ a\^{}2=16}

\emph{∴ a=4}

\emph{∴ ক্ষেত্রটির দৈর্ঘ্য = 3×4=12 মিটার।}
\end{quote}

\textbf{একটি আয়তক্ষেত্রের ক্ষেত্রফল ১৬০ বর্গমিটার। এর প্রস্থ দৈর্ঘ্য অপেক্ষা ৬
মিটার কম। আয়তক্ষেত্রের দৈর্ঘ্য কত?}

ক) ২০ মিটার

\textbf{\hl{খ) ১৬ মিটার ✔}}

গ) ১২ মিটার

ঘ) ৮ মিটার

\textbf{\hl{✔ সঠিক উত্তর: খ) ১৬ মিটার}}

\begin{quote}
\textbf{ব্যাখ্যা:} \emph{ধরি, আয়তক্ষেত্রের দৈর্ঘ্য = x মিটার}

\emph{এবং প্রস্থ = (x−6) মিটার}

\emph{∴ আয়তক্ষেত্রটির ক্ষেত্রফল = x(x−6) বর্গমিটার}

\emph{প্রশ্নমতে, x(x−6)=160}

\emph{⇒ x\^{}2−6x−160=0}

\emph{⇒ x\^{}2−16x+10x−160=0}

\emph{⇒ (x−16)(x+10)=0}

\emph{⇒ x=16 (যেহেতু দৈর্ঘ্য ঋণাত্মক হতে পারে না)}

\emph{∴ আয়তক্ষেত্রের দৈর্ঘ্য ১৬ মিটার।}
\end{quote}

\textbf{একটি আয়তাকার জমির ক্ষেত্রফল ৪,২৫০ বর্গমিটার এবং প্রস্থ ৫০ মিটার।
জমিতে দৈর্ঘ্য কত?}

ক) ৭৫ মি.

\textbf{\hl{খ) ৮৫ মি. ✔}}

গ) ৯৫ মি.

ঘ) ৬৫ মি.

\textbf{\hl{✔ সঠিক উত্তর: খ) ৮৫ মি.}}

\begin{quote}
\textbf{ব্যাখ্যা:} \emph{আমরা জানি, আয়তাকার জমির দৈর্ঘ্য = (ক্ষেত্রফল/প্রস্থ)}

\emph{= (4250/50)=85 মিটার।}
\end{quote}

\textbf{একটি আয়তাকার ঘরের দৈর্ঘ্য প্রস্থ অপেক্ষা ৪ মিটার বেশি। ঘরটির পরিসীমা
৩২ মিটার হলে এর ক্ষেত্রফল কত?}

\textbf{\hl{ক) ৬০ বর্গ মিটার ✔}}

খ) ৪৮ বর্গ মিটার

গ) ১০ বর্গ মিটার

ঘ) ৪২ বর্গ মিটার

\textbf{\hl{✔ সঠিক উত্তর: ক) ৬০ বর্গ মিটার}}

\begin{quote}
\textbf{ব্যাখ্যা:} \emph{ধরি, ঘরটির প্রস্থ = x মিটার}

\emph{এবং ঘরটির দৈর্ঘ্য = (x+4) মিটার (যেহেতু দৈর্ঘ্য ৪ মিটার বেশি)}

\emph{প্রশ্নমতে, 2(x+4+x)=32}

\emph{⇒ 2x+4=16}

\emph{⇒ 2x=12}

\emph{∴ x=6 মিটার (প্রস্থ)}

\emph{দৈর্ঘ্য = 6+4=10 মিটার}

\emph{∴ ক্ষেত্রফল = 10×6=60 বর্গ মিটার।}
\end{quote}

\textbf{আয়তাকার একটি ক্ষেত্রের প্রস্থ অপেক্ষা দৈর্ঘ্য ১২ মিটার বড় এবং ক্ষেত্রটির
পরিসীমা ১৩৬ মিটার হলে ক্ষেত্রটির দৈর্ঘ্য ও প্রস্থ কত?}

ক) ৪০ ও ৫২

\textbf{\hl{খ) ৪০ ও ২৮ ✔}}

গ) ৪২ ও ৩২

ঘ) ৩৮ ও ৩৬

\textbf{\hl{✔ সঠিক উত্তর: খ) ৪০ ও ২৮}}

\begin{quote}
\textbf{ব্যাখ্যা:} \emph{ধরি, প্রস্থ = x মি., দৈর্ঘ্য = (x+12) মি.}

\emph{শর্তমতে, 2(x+12+x)=136}

\emph{⇒ 2x+12=68}

\emph{⇒ 2x=56}

\emph{∴ x=28}

\emph{প্রস্থ = ২৮ মি. এবং দৈর্ঘ্য = 28+12=40 মি.}
\end{quote}

\textbf{একটি আয়তাকার ঘরের দৈর্ঘ্য, প্রস্থ অপেক্ষা ৪ মিটার বেশি। ঘরটির পরিসীমা
৩২ মিটার হলে, এর ক্ষেত্রফল কত?}

ক) ৬ ব.মি.

\textbf{\hl{খ) ৬০ ব.মি. ✔}}

গ) ১০ ব.মি.

ঘ) ৬৪ ব.মি.

\textbf{\hl{✔ সঠিক উত্তর: খ) ৬০ ব.মি.}}

\begin{quote}
\textbf{ব্যাখ্যা:} \emph{ধরি, প্রস্থ = x মি. এবং দৈর্ঘ্য = (x+4) মি. (যেহেতু
দৈর্ঘ্য ৪ মিটার বেশি)}

\emph{শর্তমতে, 2(x+4+x)=32}

\emph{⇒ 2x+4=16}

\emph{⇒ 2x=12}

\emph{∴ x=6 মি. (প্রস্থ)}

\emph{দৈর্ঘ্য = 6+4=10 মি.}

\emph{∴ ক্ষেত্রফল = 10×6=60 বর্গমিটার।}
\end{quote}

\textbf{একটি আয়তাকার ঘরের দৈর্ঘ্য বিস্তারের দ্বিগুণ। এর ক্ষেত্রফল ৫১২ বর্গমিটার
হলে পরিসীমা কত?}

ক) ৯৮ মিটার

\textbf{\hl{খ) ৯৬ মিটার ✔}}

গ) ৯৪ মিটার

ঘ) ৯২ মিটার

\textbf{\hl{✔ সঠিক উত্তর: খ) ৯৬ মিটার}}

\begin{quote}
\textbf{ব্যাখ্যা:} \emph{মনেকরি, আয়তক্ষেত্রের প্রস্থ = x মিটার}

\emph{∴ আয়তক্ষেত্রের দৈর্ঘ্য = 2x মিটার}

\emph{প্রশ্নমতে, 2x×x=512}

\emph{⇒ 2x\^{}2=512}

\emph{⇒ x\^{}2=256}

\emph{∴ x=√(256)=16}

\emph{আয়তক্ষেত্রের দৈর্ঘ্য = 2×16=32 মিটার}

\emph{আয়তক্ষেত্রের প্রস্থ = ১৬ মিটার}

\emph{∴ পরিসীমা = 2(32+16)=2×48=96 মিটার।}
\end{quote}

\textbf{একটি আয়তাকার ক্ষেত্রের দৈর্ঘ্য বিস্তারের ৩ গুণ। দৈর্ঘ্য ৪৮ মিটার হলে,
ক্ষেত্রটির পরিসীমা কত?}

ক) ৬৮ মিটার

\textbf{\hl{খ) ১২৮ মিটার ✔}}

গ) ৯৬ মিটার

ঘ) ২২৮ মিটার

\textbf{\hl{✔ সঠিক উত্তর: খ) ১২৮ মিটার}}

\begin{quote}
\textbf{ব্যাখ্যা:} \emph{দেয়া আছে, দৈর্ঘ্য = ৪৮ মিটার}

\emph{যেহেতু দৈর্ঘ্য বিস্তারের ৩ গুণ, সেহেতু বিস্তার (প্রস্থ) = (48/3)=16 মিটার}

\emph{∴ পরিসীমা = 2(দৈর্ঘ্য+প্রস্থ)=2(48+16)=2×64=128 মিটার।}
\end{quote}

\textbf{20x পরিসীমাবিশিষ্ট আয়তক্ষেত্রের এক বাহুর দৈর্ঘ্য 4x+3 হলে, অপর বাহুর
দৈর্ঘ্য কত?}

ক) 4x−3

খ) 5x+3

গ) 5x−3

\textbf{\hl{ঘ) 6x−3 ✔}}

\textbf{\hl{✔ সঠিক উত্তর: ঘ) 6x−3}}

\begin{quote}
\textbf{ব্যাখ্যা:} \emph{ধরি, অপর বাহুর দৈর্ঘ্য = a}

\emph{আমরা জানি, আয়তক্ষেত্রের পরিসীমা = 2(দৈর্ঘ্য+প্রস্থ)}

\emph{শর্তমতে, 2(4x+3+a)=20x}

\emph{⇒ 4x+3+a=10x}

\emph{⇒ a=10x−4x−3}

\emph{∴ a=6x−3 .}
\end{quote}

\textbf{একটি আয়তাকার ক্ষেত্রের দৈর্ঘ্য ও প্রস্থের অনুপাত ৩ : ১ এবং উহার পরিসীমা
২০০ মিটার হলে আয়তাকার ক্ষেত্রটির ক্ষেত্রফল কত?}

ক) ১৫৭৫ ব.মি.

খ) ১৬৭৫ ব. মি.

গ) ১৭৭৫ ব. মি.

\textbf{\hl{ঘ) ১৮৭৫ ব.মি. ✔}}

\textbf{\hl{✔ সঠিক উত্তর: ঘ) ১৮৭৫ ব.মি.}}

\begin{quote}
\textbf{ব্যাখ্যা:} \emph{ধরি, দৈর্ঘ্য = 3x এবং প্রস্থ = x}

\emph{শর্তমতে, 2(3x+x)=200}

\emph{⇒ 8x=200}

\emph{⇒ x=25}

\emph{দৈর্ঘ্য = 3×25=75 মিটার এবং প্রস্থ = ২৫ মিটার}

\emph{∴ ক্ষেত্রফল = 75×25=1875 বর্গমিটার।}
\end{quote}

\textbf{একটি আয়তক্ষেত্রের দৈর্ঘ্য ১৮ সেমি. এবং প্রস্থ ১০ সেমি.। আয়তক্ষেত্রের
দৈর্ঘ্য বৃদ্ধি করে ২৫ সেমি. করা হলো। আয়তক্ষেত্রের প্রস্থ কত হলে ক্ষেত্রফল
অপরিবর্তিত থাকবে?}

\textbf{\hl{ক) ৭.২ সেমি ✔}}

খ) ৭.৩ সেমি

গ) ৭ সেমি

ঘ) ৭.১ সেমি

\textbf{\hl{✔ সঠিক উত্তর: ক) ৭.২ সেমি}}

\begin{quote}
\textbf{ব্যাখ্যা:} \emph{প্রাথমিক ক্ষেত্রফল = 18×10=180 বর্গ সেন্টিমিটার}

\emph{ধરી, পরবর্তীতে প্রস্থ = x সেন্টিমিটার}

\emph{শর্তমতে, 25×x=180}

\emph{⇒ x=(180/25)=7.2 সেন্টিমিটার।}
\end{quote}

\textbf{একটি আয়তাকার বিলবোর্ডের ক্ষেত্রফল ১৯২ বর্গইঞ্চি এবং এর পরিসীমা ২৮
ইঞ্চি। এর প্রস্থ দৈর্ঘ্য অপেক্ষা শতকরা কত কম?}

ক) ৩৩\%

খ) ২৫\%

গ) ১২.৫\%

\textbf{\hl{ঘ) কোনোটাই নয় ✔}}

\textbf{\hl{✔ সঠিক উত্তর: ঘ) কোনোটাই নয়}}

\begin{quote}
\textbf{ব্যাখ্যা:} \emph{ধরি, দৈর্ঘ্য = x ইঞ্চি এবং প্রস্থ = y ইঞ্চি}

\emph{xy=192 এবং 2(x+y)=28⇒x+y=14}

\emph{আমরা জানি, (x−y)\^{}2=(x+y)\^{}2−4xy=(14)\^{}2−4×192=196−768=−572
(বাস্তব ক্ষেত্রে মান ঋণাত্মক হওয়ায় দৈর্ঘ্য ১৬ ও প্রস্থ ১২0 বা এমন বাহু সম্ভব নয়,
তবে সঠিক সমাধান অনুযায়ী দৈর্ঘ্য ১৬ ইঞ্চি ও প্রস্থ ১২ ইঞ্চি হলে ক্ষেত্রফল ১৯২ ও
পরিসীমা ৫৬ হয়, কিন্তু এখানে পরিসীমা ২৮ দেওয়া থাকায় উত্তর
\textquotesingle কোনোটাই নয়\textquotesingle{} হবে)।}
\end{quote}

\textbf{একটি আয়তক্ষেত্রের দৈর্ঘ্য ২৫\% বাড়ানো হলে ও এর প্রস্থ ২০\% বাড়ানো হলে
এর ক্ষেত্রফলে কি পরিবর্তন হবে?}

ক) ২৫\% বাড়বে

\textbf{\hl{খ) ৫০\% বাড়বে ✔}}

গ) ১৫০\% বাড়বে

ঘ) ২০০\% বাড়বে

\textbf{\hl{✔ সঠিক উত্তর: খ) ৫০\% বাড়বে}}

\begin{quote}
\textbf{ব্যাখ্যা:} \emph{ক্ষেত্রফল পরিবর্তনের সূত্র: A+B+(AB/100)}

\emph{∴ শতকরা বাড়বে = 25+20+(25×20/100)=25+20+5=50\%}

\emph{অর্থাৎ ক্ষেত্রফল ৫০\% বাড়বে।}
\end{quote}

\textbf{একটি আয়তাকার ক্ষেত্রের দৈর্ঘ্য ২০\% বৃদ্ধি ও প্রস্থ ১০\% হ্রাস করা হলে,
ক্ষেত্রফলের শতকরা কত পরিবর্তন হবে?}

ক) ১০৮\% বৃদ্ধি

খ) ১০৮\% হ্রাস

\textbf{\hl{গ) ৮\% বৃদ্ধি ✔}}

ঘ) ৮\% হ্রাস

\textbf{\hl{✔ সঠিক উত্তর: গ) ৮\% বৃদ্ধি}}

\begin{quote}
\textbf{ব্যাখ্যা:} \emph{ধরি, দৈর্ঘ্য = ১০০ একক এবং প্রস্থ = ১০০ একক}

\emph{∴ ক্ষেত্রফল = ১০,০০০ বর্গ একক}

\emph{২০\% বৃদ্ধিতে দৈর্ঘ্য = ১২০ একক এবং ১০\% হ্রাসে প্রস্থ = ৯০ একক}

\emph{∴ নতুন ক্ষেত্রফল = 120×90=10800 বর্গ একক}

\emph{ক্ষেত্রফল বৃদ্ধি = 10800−10000=800 বর্গ একক}

\emph{∴ ক্ষেত্রফল বৃদ্ধির হার = ((800/10000)×100)\%=8\%}

\emph{∴ ক্ষেত্রফল ৮\% বৃদ্ধি পাবে।}
\end{quote}

\textbf{একটি আয়তক্ষেত্রের দৈর্ঘ্য ১০\% বৃদ্ধি এবং প্রস্থ p\% হ্রাস করার ফলে এর
ক্ষেত্রফল ১২\% হ্রাস পেলে p এর মান কত?}

ক) ১২

খ) ১৫

\textbf{\hl{গ) ২০ ✔}}

ঘ) ২২

\textbf{\hl{✔ সঠিক উত্তর: গ) ২০}}

\begin{quote}
\textbf{ব্যাখ্যা:} \emph{ধরি, দৈর্ঘ্য = ১০০ একক এবং প্রস্থ = ১০০ একক}

\emph{প্রাথমিক ক্ষেত্রফল = 10000 বর্গ একক}

\emph{১০\% বৃদ্ধিতে দৈর্ঘ্য = ১১০ একক}

\emph{p\% হ্রাসে প্রস্থ = (100−p) একক}

\emph{নতুন ক্ষেত্রফল = 110(100−p)}

\emph{১২\% হ্রাসে নতুন ক্ষেত্রফল = 10000×(88/100)=8800}

\emph{শর্তমতে, 110(100−p)=8800}

\emph{⇒ 100−p=(8800/110)=80}

\emph{∴ p=100−80=20 ।}
\end{quote}

\textbf{একটি আয়তক্ষেত্রের দৈর্ঘ্য ৩৬ মিটার এবং প্রস্থ ২৫ মিটার। আয়তক্ষেত্রের
দৈর্ঘ্য ২০\% হ্রাস পেলে, প্রস্থ ১০\% বৃদ্ধি পেলে শতকরা কত হ্রাস বা বৃদ্ধি পায়?}

ক) ১০\%

খ) ০\%

\textbf{\hl{গ) ১২\% ✔}}

ঘ) ১৯\%

\textbf{\hl{✔ সঠিক উত্তর: গ) ১২\%}}

\begin{quote}
\textbf{ব্যাখ্যা:} \emph{প্রাথমিক ক্ষেত্রফল = 36×25=900 বর্গমিটার}

\emph{হ্রাসকৃত দৈর্ঘ্য = 36−(36×(20/100))=36−7.2=28.8 মিটার}

\emph{বর্ধিত প্রস্থ = 25+(25×(10/100))=25+2.5=27.5 মিটার}

\emph{পরিবর্তিত ক্ষেত্রফল = 28.8×27.5=792 বর্গমিটার}

\emph{ক্ষেত্রফল হ্রাস = 900−792=108 বর্গমিটার}

\emph{∴ শতকরা হ্রাসের হার = ((108/900)×100)\%=12\% (১২\% হ্রাস পায়)।}
\end{quote}

\textbf{একটি বর্গাকার জমির ক্ষেত্রফল ৯৮ বর্গমিটার। বর্গাকার জমির পাশে একটি
আয়তাকার জমি আছে যার দৈর্ঘ্য বর্গাকার জমির কর্ণের ১৫০\%। আবার আয়তাকার জমির
প্রস্থ এর দৈর্ঘ্যের এক-তৃতীয়াংশ। আয়তাকার জমির ক্ষেত্রফল বর্গাকার জমির ক্ষেত্রফলের
শতকরা কত অংশ?}

ক) ২৫\%

খ) ৫০\%

\textbf{\hl{গ) ১৫০\% ✔}}

ঘ) ২০০\%

\textbf{\hl{✔ সঠিক উত্তর: গ) ১৫০\%}}

\begin{quote}
\textbf{ব্যাখ্যা:} \emph{বর্গাকার জমির একবাহুর দৈর্ঘ্য = √(98) মিটার}

\emph{বর্গাকার জমির কর্ণের দৈর্ঘ্য = √(2)×√(98)=√(196)=14 মিটার}

\emph{আয়তাকার জমির দৈর্ঘ্য = 14×(150/100)=21 মিটার}

\emph{আয়তাকার জমির প্রস্থ = 21×(1/3)=7 মিটার}

\emph{আয়তাকার জমির ক্ষেত্রফল = 21×7=147 বর্গমিটার}

\emph{∴ শতকরা = ((147/98)×100)\%=150\% ।}
\end{quote}

\textbf{একটি আয়তাকার জমির প্রস্থ এর দৈর্ঘ্যের তিন-পঞ্চমাংশ। দৈর্ঘ্য ও প্রস্থ
উভয়কে দ্বিগুণ করা হলে নতুন আয়তাকারের ক্ষেত্রফল মূল আয়তক্ষেত্রের ক্ষেত্রফলের কত
শতাংশ হবে?}

ক) ১০০\%

খ) ২০০\%

গ) ৩০০\%

\textbf{\hl{ঘ) ৪০০\% ✔}}

\textbf{\hl{✔ সঠিক উত্তর: ঘ) ৪০০\%}}

\begin{quote}
\textbf{ব্যাখ্যা:} \emph{ধরি, আয়তাকার জমির দৈর্ঘ্য = 5x একক}

\emph{প্রস্থ = 5x×(3/5)=3x একক}

\emph{মূল ক্ষেত্রফল = 5x×3x=15x\^{}2 বর্গ একক}

\emph{দৈর্ঘ্য ও প্রস্থ দ্বিগুণ করা হলে নতুন দৈর্ঘ্য = 10x এবং নতুন প্রস্থ = 6x}

\emph{নতুন ক্ষেত্রফল = 10x×6x=60x\^{}2 বর্গ একক}

\emph{∴ শতকরা = ((60x\^{}2/15x\^{}2)×100)\%=400\% ।}
\end{quote}

\textbf{একটি হিমাগারকে নবায়ন করার সময় এর দৈর্ঘ্য ৩০\% ও প্রস্থ ৫০\% বাড়ানো
হলো এবং উচ্চতা ২০\% কমানো হলো। নবায়নকৃত হিমাগারের আয়তন পুরোনো হিমাগারের
আয়তন থেকে শতকরা কত অংশ বেশি?}

\textbf{\hl{ক) ৫৬\% ✔}}

খ) ৫০\%

গ) ৪৫\%

ঘ) ৬৫\%

\textbf{\hl{✔ সঠিক উত্তর: ক) ৫৬\%}}

\begin{quote}
\textbf{ব্যাখ্যা:} \emph{ধরি, হিমাগারের দৈর্ঘ্য = x , প্রস্থ = y , উচ্চতা = z}

\emph{পুরোনো আয়তন = xyz}

\emph{নবায়নকৃত দৈর্ঘ্য = 1.3x , প্রস্থ = 1.5y , উচ্চতা = 0.8z}

\emph{নবায়নকৃত আয়তন = 1.3x×1.5y×0.8z=1.56xyz}

\emph{আয়তন বৃদ্ধি পায় = 1.56xyz−xyz=0.56xyz}

\emph{∴ শতকরা বৃদ্ধি = 56\% ।}
\end{quote}

\textbf{একটি আয়তাকার বাগানের দৈর্ঘ্য ও প্রস্থ যথাক্রমে ১৫০ মিটার ও ১০০ মিটার।
বাগানটির দৈর্ঘ্য ২০\% এবং প্রস্থ ১০\% বৃদ্ধি করলে নতুন বাগানটির ক্ষেত্রফল কত
বর্গমিটার হবে?}

\textbf{\hl{ক) ১৯৮০০ ✔}}

খ) ১৫০০

গ) ৩০০০

ঘ) ১৯৫০০

\textbf{\hl{✔ সঠিক উত্তর: ক) ১৯৮০০}}

\begin{quote}
\textbf{ব্যাখ্যা:} \emph{নতুন দৈর্ঘ্য = (150+150এর20\%) মিটার}

\emph{= (150+150×(20/100)) মিটার = 180 মিটার}

\emph{নতুন প্রস্থ = (100+100এর10\%) মিটার}

\emph{= (100+100×(10/100)) মিটার = 110 মিটার}

\emph{∴ নতুন বাগানের ক্ষেত্রফল = 180×110=19800 বর্গমিটার।}
\end{quote}

\textbf{একটি আয়তাকার ক্ষেত্রের ক্ষেত্রফল ৪৮ ব.মি.। দৈর্ঘ্য ১ মি. কমালে এবং
প্রস্থ ১ মি. বাড়ালে ক্ষেত্রফল ৪৯ ব. মি. হয়। আয়তক্ষেত্রের দৈর্ঘ্য ও প্রস্থ কত মিটার
হবে?}

ক) ১২, ৪

খ) ১৬, ৩

\textbf{\hl{গ) ৮, ৬ ✔}}

ঘ) ২, ২৪

\textbf{\hl{✔ সঠিক উত্তর: গ) ৮, ৬}}

\begin{quote}
\textbf{ব্যাখ্যা:} \emph{ধরি, দৈর্ঘ্য = x মিটার এবং প্রস্থ = y মিটার}

\emph{শর্তমতে, xy=48 -\/-\/-(১)}

\emph{এবং (x−1)(y+1)=49}

\emph{⇒ xy+x−y−1=49}

\emph{⇒ 48+x−y−1=49 (মান বসিয়ে)}

\emph{⇒ x−y=2⇒x=y+2}

\emph{(১) নং সমীকরণে বসিয়ে, (y+2)y=48}

\emph{⇒ y\^{}2+2y−48=0}

\emph{⇒ (y+8)(y−6)=0}

\emph{∴ y=6 (প্রস্থ)}

\emph{⇒ x=6+2=8 (দৈর্ঘ্য)}

\emph{∴ দৈর্ঘ্য ও প্রস্থ যথাক্রমে ৮ ও ৬ মিটার।}
\end{quote}

\textbf{একটি আয়তাকার বাগানের দৈর্ঘ্য ও প্রস্থ যথাক্রমে ৬০ মিটার ও ৪০ মিটার।
বাগানের বাইরে চতুর্দিকে ২ মিটার চওড়া একটি রাস্তা আছে। রাস্তাটির ক্ষেত্রফল
নির্ণয় করুন?}

ক) ৪০১ বর্গমিটার

\textbf{\hl{খ) ৪১৬ বর্গমিটার ✔}}

গ) ৪২০ বর্গমিটার

ঘ) ৪৩০ বর্গমিটার

\textbf{\hl{✔ সঠিক উত্তর: খ) ৪১৬ বর্গমিটার}}

\begin{quote}
\textbf{ব্যাখ্যা:} \emph{রাস্তাসহ বাগানের দৈর্ঘ্য = 60+(2×2)=64 মিটার}

\emph{রাস্তাসহ বাগানের প্রস্থ = 40+(2×2)=44 মিটার}

\emph{রাস্তাসহ বাগানের ক্ষেত্রফল = 64×44=2816 বর্গমিটার}

\emph{বাগানের ক্ষেত্রফল = 60×40=2400 বর্গমিটার}

\emph{∴ রাস্তার ক্ষেত্রফল = 2816−2400=416 বর্গমিটার।}
\end{quote}

\textbf{একটি আয়তাকার মাঠের বাইরে চারদিকে একটি রাস্তা আছে। মাঠের দৈর্ঘ্য ৫০
মিটার, প্রস্থ ৪০ মিটার এবং রাস্তার চওড়া ২ মিটার। রাস্তাসহ মাঠের দৈর্ঘ্য কত
মিটার?}

\textbf{\hl{ক) ৫৪ ✔}}

খ) ৫০

গ) ৪৬

ঘ) ৮৮

\textbf{\hl{✔ সঠিক উত্তর: ক) ৫৪}}

\begin{quote}
\textbf{ব্যাখ্যা:} \emph{মাঠের দৈর্ঘ্য = ৫০ মিটার}

\emph{রাস্তার চওড়া উভয় পাশে = 2×2=4 মিটার}

\emph{∴ রাস্তাসহ মাঠের দৈর্ঘ্য = 50+4=54 মিটার।}
\end{quote}

\textbf{A rectangular plot measuring 90 metres by 50 metres is to be
enclosed by wire fencing. If the poles of the fence are kept 5 metres
apart, how many poles will be needed? / ৫০ মিটার ৯০ মিটার একটি
আয়তক্ষেত্রাকার ক্ষেত্রকে তারের বেড়া দ্বারা ঘিরে রাখা হয়। বেড়াতে খুঁটিগুলো ৫
মিটার দূরে দূরে রাখা হয়, এতে কতগুলি খুঁটির প্রয়োজন হবে?}

ক) ৫৫

\textbf{\hl{খ) ৫৬ ✔}}

গ) ৫৭

ঘ) ৫৮

\textbf{\hl{✔ সঠিক উত্তর: খ) ৫৬}}

\begin{quote}
\textbf{ব্যাখ্যা:} \emph{Perimeter of the rectangle = 2(50+90) metres}

\emph{= 280 metres}

\emph{Between 5 metres the required number of poles is = (280/5)=56 .}
\end{quote}

\textbf{একটি আয়তাকার বাগানের ক্ষেত্রফল ৪০০ বর্গমিটার এবং প্রস্থ ১৬ মিটার।
বাগানের পরিসীমা কত মিটার?}

ক) ১৬

খ) ২৫

গ) ৮১

\textbf{\hl{ঘ) ৮২ ✔}}

\textbf{\hl{✔ সঠিক উত্তর: ঘ) ৮২}}

\begin{quote}
\textbf{ব্যাখ্যা:} \emph{মাঠটির দৈর্ঘ্য = (400/16)=25 মিটার}

\emph{∴ পরিসীমা = 2(দৈর্ঘ্য+প্রস্থ)}

\emph{= 2(25+16)=2×41=82 মিটার।}
\end{quote}

\textbf{একটি আয়তাকার বাগানের দৈর্ঘ্য ও প্রস্থের অনুপাত ৪ : ৩ হলে বাগানটির
অর্ধপরিসীমার দৈর্ঘ্য কত মিটার?}

ক) 25

খ) 30

\textbf{\hl{গ) 35 ✔}}

ঘ) 45

\textbf{\hl{✔ সঠিক উত্তর: গ) 35}}

\begin{quote}
\textbf{ব্যাখ্যা:} \emph{দেয়া আছে, দৈর্ঘ্য ও প্রস্থ = 4:3}

\emph{ধরি, দৈর্ঘ্য = 4x এবং প্রস্থ = 3x (যদি প্রশ্নানুযায়ী নির্দিষ্ট মান না থাকে,
তবে বিকল্প হিসেবে অনুপাতের ধ্রুবক ধরে বা অনুপাত অনুযায়ী দৈর্ঘ্য ও প্রস্থের সমষ্টিই
অর্ধপরিসীমা নির্দেশ করে; যেমন এখানে ৪+৩=৭ এবং অনুপাতের একক অনুযায়ী মান ২০ ও
১৫ হলে অর্ধপরিসীমা ২০+১৫=৩৫ মিটার হয়)।}
\end{quote}

\textbf{৪৮ মিটার দীর্ঘ একটি আয়তক্ষেত্রের প্রস্থ দৈর্ঘ্যের এক তৃতীয়াংশ। ঐ
আয়তক্ষেত্রের সমান পরিসীমা বিশিষ্ট বর্গের একটি বাহুর দৈর্ঘ্য কত?}

ক) ১৬ মিটার

\textbf{\hl{খ) ৩২ মিটার ✔}}

গ) ২৪ মিটার

ঘ) ২৮ মিটার

\textbf{\hl{✔ সঠিক উত্তর: খ) ৩২ মিটার}}

\begin{quote}
\textbf{ব্যাখ্যা:} \emph{আয়তক্ষেত্রের প্রস্থ = (1/3)×দৈর্ঘ্য=((1/3)×48) মিটার =
১৬ মিটার}

\emph{আয়তক্ষেত্রের পরিসীমা = 2(দৈর্ঘ্য+প্রস্থ)=2(48+16)=128 মিটার}

\emph{বর্গের একটি বাহুর দৈর্ঘ্য = (128/4)=32 মিটার।}
\end{quote}

\textbf{একটি আয়তাকার মাঠের দৈর্ঘ্য প্রস্থের দ্বিগুণ। মাঠের দৈর্ঘ্য ১২ মিটার হলে
মাঠটির পরিসীমা কত মিটার?}

ক) ৩২

\textbf{\hl{খ) ৩৬ ✔}}

গ) ৪৮

ঘ) কোনোটাই নয়

\textbf{\hl{✔ সঠিক উত্তর: খ) ৩৬}}

\begin{quote}
\textbf{ব্যাখ্যা:} \emph{ধরি, প্রস্থ = x মিটার ∴ দৈর্ঘ্য = 2x}

\emph{প্রশ্নমতে, 2x=12⇒x=(12/2)⇒x=6}

\emph{∴ পরিসীমা = 2(12+6)=36 মিটার।}
\end{quote}

\textbf{একটি আয়তাকার জমির দৈর্ঘ্য প্রস্থের দ্বিগুণ এবং উহার ক্ষেত্রফল ৮০০
বর্গমিটার হলে, জমির পরিসীমা কত মিটার?}

ক) ১০০

খ) ১৮০

\textbf{\hl{গ) ১২০ ✔}}

ঘ) ১১৫

\textbf{\hl{✔ সঠিক উত্তর: গ) ১২০}}

\begin{quote}
\textbf{ব্যাখ্যা:} \emph{ধরি, প্রস্থ = x মিটার এবং দৈর্ঘ্য = 2x মিটার}

\emph{ক্ষেত্রফল = 2x×x=2x\^{}2}

\emph{শর্তমতে, 2x\^{}2=800⇒x\^{}2=400⇒x=20}

\emph{প্রস্থ = ২০ মিটার এবং দৈর্ঘ্য = 2×20=40 মিটার}

\emph{∴ পরিসীমা = 2(40+20)=2×60=120 মিটার।}
\end{quote}

\textbf{আয়তাকার বক্সের দৈর্ঘ্য বিস্তারের দ্বিগুণ। এর ক্ষেত্রফল ৫১২ বর্গমিটার হলে,
পরিসীমা কত?}

ক) ১৮ মিটার

\textbf{\hl{খ) ৯৬ মিটার ✔}}

গ) ১৪ মিটার

ঘ) ৯০ মিটার

\textbf{\hl{✔ সঠিক উত্তর: খ) ৯৬ মিটার}}

\begin{quote}
\textbf{ব্যাখ্যা:} \emph{ধরি, প্রস্থ = x মিটার এবং দৈর্ঘ্য = 2x মিটার}

\emph{ক্ষেত্রফল = 2x\^{}2=512}

\emph{⇒ x\^{}2=256⇒x=16 (প্রস্থ)}

\emph{দৈর্ঘ্য = 2×16=32 মিটার}

\emph{∴ পরিসীমা = 2(16+32)=96 মিটার।}
\end{quote}

\textbf{একটি বর্গের ক্ষেত্রফল একটি আয়তের ক্ষেত্রফলের সমান। বর্গের পরিসীমা ২৪
সেমি এবং আয়তের প্রস্থ ৪ সেমি হলে, আয়তের পরিসীমা হবে?}

ক) ১৬ সেমি

খ) ২০ সেমি

গ) ২৪ সেমি

\textbf{\hl{ঘ) ২৬ সেমি ✔}}

\textbf{\hl{✔ সঠিক উত্তর: ঘ) ২৬ সেমি}}

\begin{quote}
\textbf{ব্যাখ্যা:} \emph{বর্গের পরিসীমা = ২৪ সেমি}

\emph{∴ বর্গের এক বাহু = (24/4)=6 সেমি}

\emph{বর্গের ক্ষেত্রফল = 6\^{}2=36 বর্গ সেমি = আয়তের ক্ষেত্রফল}

\emph{ধরি, আয়তের দৈর্ঘ্য = x সেমি}

\emph{শর্তমতে, 4x=36}

\emph{⇒x=(36/4)=9 সেমি}

\emph{দৈর্ঘ্য ৯ সেমি হলে পরিসীমা 2(4+9)=26 সেমি হয়)।}
\end{quote}

\textbf{একটি আয়তাকার জমির দৈর্ঘ্য প্রস্থের তিনগুণ এবং উক্ত জমির পরিসীমা ৪০
মিটার হলে, জমিটির দৈর্ঘ্য কত?}

ক) ১২ মি.

খ) ৩০ মি.

\textbf{\hl{গ) ১৫ মি. ✔}}

ঘ) ৩৬ মি.

\textbf{\hl{✔ সঠিক উত্তর: গ) ১৫ মি.}}

\begin{quote}
\textbf{ব্যাখ্যা:} \emph{ধরি, প্রস্থ = x মিটার এবং দৈর্ঘ্য = 3x মিটার}

\emph{শর্তমতে, 2(3x+x)=40}

\emph{⇒ 2×4x=40}

\emph{⇒ x=(40/8)=5}

\emph{∴ জমির দৈর্ঘ্য = 3×5=15 মিটার।}
\end{quote}

\textbf{১০ মিটার/মিনিট বেগে ৪০ মিটার দীর্ঘ ও ৩০ মিটার প্রস্থ একটি আয়তাকার
বাগানের এককোণ থেকে কর্ণ বরাবর হেঁটে অপর প্রান্তে পৌঁছে বাগানের কিনার দিয়ে
হেঁটে আগের জায়গায় ফেরত আসতে কত মিনিট সময় লাগবে?}

ক) ১০

\textbf{\hl{খ) ১২ ✔}}

গ) ১৩

ঘ) ১৪

\textbf{\hl{✔ সঠিক উত্তর: খ) ১২}}

\begin{quote}
\textbf{ব্যাখ্যা:} \emph{বাগানের কর্ণ =
√(40\^{}2+30\^{}2)=√(1600+900)=√(2500)=50 মিটার}

\emph{কর্ণ বরাবর যেতে সময় লাগে = (50/10)=5 মিনিট}

\emph{কিনার দিয়ে হেঁটে আগের জায়গায় ফেরত আসতে দূরত্ব অতিক্রম করতে হবে
পরিসীমার সমান অর্থাৎ 40+30=70 মিটার}

\emph{কিনার দিয়ে আসতে সময় লাগে = (70/10)=7 মিনিট}

\emph{∴ মোট সময় = 5+7=12 মিনিট।}
\end{quote}

\textbf{কোনো বর্গক্ষেত্রের একবাহুর দৈর্ঘ্য ১০ সে.মি.। বর্গক্ষেত্রটির ক্ষেত্রফল কত?}

ক) ৪০ সে.মি.

\textbf{\hl{খ) ১০০ বর্গ সে.মি. ✔}}

গ) ৪০ বর্গ সে.মি.

ঘ) ১০০ সে.মি.

\textbf{\hl{✔ সঠিক উত্তর: খ) ১০০ বর্গ সে.মি.}}

\begin{quote}
\textbf{ব্যাখ্যা:} \emph{আমরা জানি, বর্গক্ষেত্রের ক্ষেত্রফল = (একবাহু)\^{}2}

\emph{∴ বর্গক্ষেত্রটির ক্ষেত্রফল = 10\^{}2=100 বর্গ সেন্টিমিটার।}
\end{quote}

\textbf{একটি বর্গাকার জমির এক পাশের দৈর্ঘ্য ১০ ফুট হলে ঐ জমির ক্ষেত্রফল কত
বর্গফুট?}

ক) ৮০ বর্গফুট

\textbf{\hl{খ) ১০০ বর্গফুট ✔}}

গ) ২০ বর্গফুট

ঘ) কোনোটাই নয়

\textbf{\hl{✔ সঠিক উত্তর: খ) ১০০ বর্গফুট}}

\begin{quote}
\textbf{ব্যাখ্যা:} \emph{যেহেতু জমিটি বর্গাকার, সেহেতু জমির ক্ষেত্রফল =
(10)\^{}2=100 বর্গফুট।}
\end{quote}

\textbf{বর্গক্ষেত্রের ৪টি কোণের সমষ্টি -}

ক) ১ সমকোণ

খ) ২ সমকোণ

গ) ৩ সমকোণ

\textbf{\hl{ঘ) ৪ সমকোণ ✔}}

\textbf{\hl{✔ সঠিক উত্তর: ঘ) ৪ সমকোণ}}

\begin{quote}
\textbf{ব্যাখ্যা:} \emph{বর্গক্ষেত্রের প্রতিটি কোণ সমকোণ বা 90° .}

\emph{∴ বর্গক্ষেত্রের ৪টি কোণের সমষ্টি = 4×90°=360° অর্থাৎ ৪ সমকোণ।}
\end{quote}

\textbf{দুইটি বর্গক্ষেত্রের বাহুর দৈর্ঘ্য যথাক্রমে
\textquotesingle a\textquotesingle{} মিটার এবং
\textquotesingle b\textquotesingle{} মিটার হলে তাদের ক্ষেত্রফলের অনুপাত
কত?}

\textbf{\hl{ক) a² : b² ✔}}

খ) a : b

গ) b² : a²

ঘ) (a/2):(b/2)

\textbf{\hl{✔ সঠিক উত্তর: ক) a² : b²}}

\begin{quote}
\textbf{ব্যাখ্যা:} \emph{বর্গক্ষেত্রের ক্ষেত্রফল = (বাহু)\^{}2 বর্গ একক}

\emph{প্রথম বর্গের ক্ষেত্রফল = a\^{}2 বর্গ একক}

\emph{দ্বিতীয় বর্গের ক্ষেত্রফল = b\^{}2 বর্গ একক}

\emph{∴ ক্ষেত্রফলের অনুপাত = a\^{}2:b\^{}2 ।}
\end{quote}

\textbf{একটি বর্গক্ষেত্রের এক বাহু দৈর্ঘ্য y একক। এর কর্ণের দৈর্ঘ্য কত?}

ক) √(y) একক

খ) √((y+2)) একক

গ) 2√(y) একক

\textbf{\hl{ঘ) y√(2) একক ✔}}

\textbf{\hl{✔ সঠিক উত্তর: ঘ) y√(2) একক}}

\begin{quote}
\textbf{ব্যাখ্যা:} \emph{বর্গক্ষেত্রের কর্ণের দৈর্ঘ্য = √(2)×এক বাহু∴ কর্ণের দৈর্ঘ্য
= √(2)y বা y√(2) একক।}
\end{quote}

\textbf{একটি বর্গক্ষেত্রের কর্ণের দৈর্ঘ্য 4√(2) একক হলে ঐ বর্গক্ষেত্রের ক্ষেত্রফল
কত বর্গ একক?}

ক) 24

খ) 8

\textbf{\hl{গ) 16 ✔}}

ঘ) 32

\textbf{\hl{✔ সঠিক উত্তর: গ) 16}}

\begin{quote}
\textbf{ব্যাখ্যা:} \emph{ধরি, বর্গক্ষেত্রের বাহুর দৈর্ঘ্য = a একক}

\emph{প্রশ্নমতে, a√(2)=4√(2)⇒a=4}

\emph{∴ বর্গক্ষেত্রের ক্ষেত্রফল = a\^{}2=4\^{}2=16 বর্গ একক।}
\end{quote}

\textbf{একটি বর্গক্ষেত্রের প্রতিটি বাহুর দৈর্ঘ্য ৫ ফুট হলে কর্ণের উপর অঙ্কিত
বর্গক্ষেত্রের ক্ষেত্রফল কত?}

ক) ৪০ বর্গফুট

\textbf{\hl{খ) ৫০ বর্গফুট ✔}}

গ) ১০০ বর্গফুট

ঘ) ৮০ বর্গফুট

\textbf{\hl{✔ সঠিক উত্তর: খ) ৫০ বর্গফুট}}

\begin{quote}
\textbf{ব্যাখ্যা:} \emph{কর্ণের দৈর্ঘ্য = √(5\^{}2+5\^{}2)=√(50) ফুট}

\emph{∴ কর্ণের উপর অঙ্কিত বর্গক্ষেত্রের ক্ষেত্রফল = (√(50))\^{}2=50 বর্গফুট।}
\end{quote}

\textbf{একটি বর্গক্ষেত্রের ক্ষেত্রফল ৪০৮.০১ বর্গমিটার হলে প্রতিটি বাহুর দৈর্ঘ্য
হবে-}

ক) ২০.০১

খ) ২০.১১

গ) ২০.২১

\textbf{\hl{ঘ) ২০.১ ✔}}

\textbf{\hl{✔ সঠিক উত্তর: ঘ) ২০.১}}

\begin{quote}
\textbf{ব্যাখ্যা:} \emph{ধরি, বর্গক্ষেত্রের একবাহুর দৈর্ঘ্য a}

\emph{প্রশ্নমতে, a\^{}2=408.01⇒a=√(408.01)=20.1 মিটার।}
\end{quote}

\textbf{৫ সেমি বাহুবিশিষ্ট বর্গের কর্ণের দৈর্ঘ্য কত?}

ক) ১০ সে.মি.

খ) ২৫ সে.মি.

\textbf{\hl{গ) 5√(2) সে.মি. ✔}}

ঘ) ২০ সে.মি.

\textbf{\hl{✔ সঠিক উত্তর: গ) 5√(2) সে.মি.}}

\begin{quote}
\textbf{ব্যাখ্যা:} \emph{বর্গের কর্ণের দৈর্ঘ্য = a√(2) একক}

\emph{= 5√(2) সেমি {[}∴ a=5 সেমি{]}}
\end{quote}

\textbf{একটি বর্গক্ষেত্রের ক্ষেত্রফল ১৮ বর্গমিটার হলে, কর্ণের দৈর্ঘ্য কত?}

ক) ১২ মিটার

খ) ৯ মিটার

\textbf{\hl{গ) ৬ মিটার ✔}}

ঘ) ৩ মিটার

\textbf{\hl{✔ সঠিক উত্তর: গ) ৬ মিটার}}

\begin{quote}
\textbf{ব্যাখ্যা:} \emph{বর্গক্ষেত্রের বাহু a একক}

\emph{∴ a\^{}2=18}

\emph{⇒ a=√(9×2)=3√(2)}

\emph{∴ কর্ণের দৈর্ঘ্য = √(2)×a=√(2)×3√(2)=6 মিটার।}
\end{quote}

\textbf{একটি বর্গক্ষেত্রের প্রত্যেকটি বাহুর দৈর্ঘ্য a সেমি হলে, ঐ ক্ষেত্রের কর্ণের
দৈর্ঘ্য কত?}

\textbf{\hl{ক) √(2)a ✔}}

খ) 2a

গ) a²

ঘ) 2a²

\textbf{\hl{✔ সঠিক উত্তর: ক) √(2)a}}

\begin{quote}
\textbf{ব্যাখ্যা:} \emph{কর্ণ AC=√(a\^{}2+a\^{}2)=√(2a\^{}2)=√(2)a সেমি।}
\end{quote}

\textbf{একটি বর্গক্ষেত্রের একটি বাহুর দৈর্ঘ্য 1 একক হলে উহার কর্ণের দৈর্ঘ্য-}

\textbf{\hl{ক) √(2) ✔}}

খ) 1

গ) 5

ঘ) 2

\textbf{\hl{✔ সঠিক উত্তর: ক) √(2)}}

\begin{quote}
\textbf{ব্যাখ্যা:} \emph{বর্গক্ষেত্রের কর্ণের দৈর্ঘ্য = a√(2) একক = 1×√(2)=√(2)
একক।}
\end{quote}

\textbf{একটি বর্গক্ষেত্রের কর্ণের দৈর্ঘ্য 6cm হলে এর ক্ষেত্রফল কত হবে?}

ক) 12 sq.cm

\textbf{\hl{খ) 18 sq.cm ✔}}

গ) 24 sq.cm

ঘ) 36 sq.cm

\textbf{\hl{✔ সঠিক উত্তর: খ) 18 sq.cm}}

\begin{quote}
\textbf{ব্যাখ্যা:} \emph{ধরি, বর্গের বাহুর দৈর্ঘ্য a সে.মি.}

\emph{কর্ণের দৈর্ঘ্য = √(2)a=6⇒a=(6/√(2))}

\emph{∴ ক্ষেত্রফল = a\^{}2=((6/√(2)))\^{}2=(36/2)=18 বর্গ সেমি।}
\end{quote}

\textbf{বর্গক্ষেত্রের এক বাহুর দৈর্ঘ্য ৪৫ মিটার হলে এর ক্ষেত্রফল কত?}

\textbf{\hl{ক) ২০২৫ বর্গমিটার ✔}}

খ) ২০১৫ বর্গমিটার

গ) ২৪৭৫ বর্গমিটার

ঘ) ১০১৫ বর্গমিটার

\textbf{\hl{✔ সঠিক উত্তর: ক) ২০২৫ বর্গমিটার}}

\begin{quote}
\textbf{ব্যাখ্যা:} \emph{বর্গক্ষেত্রের ক্ষেত্রফল = (বাহু)\^{}2=(45)\^{}2=2025
বর্গমিটার।}
\end{quote}

\textbf{একটি বর্গক্ষেত্রের পরিসীমা ১২০ ফুট। প্রত্যেক বাহুর দৈর্ঘ্য কত?}

ক) ২০ ফুট

খ) ২৫ ফুট

\textbf{\hl{গ) ৩০ ফুট ✔}}

ঘ) ৪০ ফুট

\textbf{\hl{✔ সঠিক উত্তর: গ) ৩০ ফুট}}

\begin{quote}
\textbf{ব্যাখ্যা:} \emph{বর্গক্ষেত্রের পরিসীমা = 4×বাহু=120}

\emph{∴ প্রত্যেক বাহুর দৈর্ঘ্য = (120/4)=30 ফুট।}
\end{quote}

\textbf{a বাহুবিশিষ্ট বর্গক্ষেত্রের কর্ণের উপর অঙ্কিত একটি বর্গক্ষেত্রের সমান---}

ক) a\^{}2 বর্গ একক

\textbf{\hl{খ) 2a\^{}2 বর্গ একক ✔}}

গ) √(2)a\^{}2 বর্গ একক

ঘ) √(2)a বর্গ একক

\textbf{\hl{✔ সঠিক উত্তর: খ) 2a\^{}2 বর্গ একক}}

\begin{quote}
\textbf{ব্যাখ্যা:} \emph{বর্গক্ষেত্রের বাহুর দৈর্ঘ্য = a একক}

\emph{কর্ণের দৈর্ঘ্য = √(2)a}

\emph{∴ কর্ণের উপর অঙ্কিত বর্গক্ষেত্রের ক্ষেত্রফল = (√(2)a)\^{}2=2a\^{}2 বর্গ
একক।}
\end{quote}

\textbf{√(2)a দৈর্ঘ্যের কর্ণবিশিষ্ট বর্গক্ষেত্রের বাহুর দৈর্ঘ্য সমান---}

ক) √(2)a

\textbf{\hl{খ) a ✔}}

গ) a\^{}2

ঘ) 2a\^{}2

\textbf{\hl{✔ সঠিক উত্তর: খ) a}}

\begin{quote}
\textbf{ব্যাখ্যা:} \emph{যদি কোনো বর্গের বাহুর দৈর্ঘ্য a হয় তবে তার কর্ণের
দৈর্ঘ্য √(2)a}

\emph{অর্থাৎ কোনো বর্গের কর্ণের দৈর্ঘ্য √(2)a হলে তবে তার বাহুর দৈর্ঘ্য a ।}
\end{quote}

\textbf{একটি বর্গক্ষেত্রের পরিসীমা ২৮ মিটার হলে বর্গক্ষেত্রটির ক্ষেত্রফল কত
বর্গমিটার?}

ক) ৩৬

খ) ৪২

\textbf{\hl{গ) ৪৯ ✔}}

ঘ) কোনোটাই নয়

\textbf{\hl{✔ সঠিক উত্তর: গ) ৪৯}}

\begin{quote}
\textbf{ব্যাখ্যা:} \emph{ধরি, বর্গক্ষেত্রের প্রতি বাহুর দৈর্ঘ্য = k মিটার}

\emph{পরিসীমা = 4k একক}

\emph{শর্তমতে, 4k=28⇒k=7}

\emph{∴ বর্গক্ষেত্রের ক্ষেত্রফল = 7×7=49 বর্গমিটার।}
\end{quote}

\textbf{একটি খেলার মাঠের প্রস্থ আরও ১০ মিটার বেশি হলে এটি ১০,০০০
বর্গমিটারের ক্ষেত্রবিশিষ্ট বর্গাকার মাঠ হতো। মাঠটির প্রস্থ নির্ণয় করুন।}

ক) ৮০ মিটার

খ) ১০৫ মিটার

\textbf{\hl{গ) ৯০ মিটার ✔}}

ঘ) ১০০ মিটার

\textbf{\hl{✔ সঠিক উত্তর: গ) ৯০ মিটার}}

\begin{quote}
\textbf{ব্যাখ্যা:} \emph{ধরি, বর্গাকার মাঠের প্রস্থ = x মিটার}

\emph{প্রস্থ আরও ১০ মিটার বেশি হলে দাঁড়ায় = (x+10) মিটার}

\emph{শর্তমতে, (x+10)\^{}2=10000}

\emph{⇒ x+10=100 (বর্গমূল করে)}

\emph{⇒ x=100−10=90}

\emph{∴ প্রস্থ = ৯০ মিটার।}
\end{quote}

\textbf{একটি বর্গের এক বাহু ৫ ইঞ্চি হলে বর্গটির পরিসীমা কত?}

\textbf{\hl{ক) ২০ ✔}}

খ) ২৫

গ) ২৪

ঘ) ২৬

\textbf{\hl{✔ সঠিক উত্তর: ক) ২০}}

\begin{quote}
\textbf{ব্যাখ্যা:} \emph{বর্গক্ষেত্রের পরিসীমা = 4×এক বাহুর দৈর্ঘ্য= 4×5=20
ইঞ্চি।}
\end{quote}

\textbf{একটি বর্গাকৃতি বাগানের ক্ষেত্রফল ১ হেক্টর। বাগানটির পরিসীমা কত
মিটার?}

ক) ২০০

খ) ৩০০

\textbf{\hl{গ) ৪০০ ✔}}

ঘ) ৫০০

\textbf{\hl{✔ সঠিক উত্তর: গ) ৪০০}}

\begin{quote}
\textbf{ব্যাখ্যা:} \emph{বাগানের ক্ষেত্রফল = ১ হেক্টর = 10000 বর্গমিটার}

\emph{বাগানের একবাহু = √(10000)=100 মিটার}

\emph{∴ বাগানটির পরিসীমা = 4×100=400 মিটার।}
\end{quote}

\textbf{একটি বর্গাকার বাগানের ক্ষেত্রফল ২০২৫ বর্গমিটার। এর চারদিকে বেড়া আছে।
বেড়ার দৈর্ঘ্য হবে-\/-\/-}

\textbf{\hl{ক) ১৮০ মিটার ✔}}

খ) ১৬০ মিটার

গ) ১৪০ মিটার

ঘ) কোনোটাই নয়

\textbf{\hl{✔ সঠিক উত্তর: ক) ১৮০ মিটার}}

\begin{quote}
\textbf{ব্যাখ্যা:} \emph{বর্গাকার বাগানের একবাহু = √(2025)=45 মিটার}

\emph{∴ বাগানটির পরিসীমা বা বেড়ার দৈর্ঘ্য = 4×45=180 মিটার।}
\end{quote}

\textbf{একটি বর্গক্ষেত্রের ক্ষেত্রফল ৬৪ বর্গফুট। ঐ বর্গক্ষেত্রের চতুর্দিকে ২ ফুট
প্রস্থের একটি রাস্তা রয়েছে। রাস্তাসহ বর্গক্ষেত্রের ক্ষেত্রফল কত?}

ক) ৬৬ বর্গফুট

খ) ৭৬ বর্গফুট

\textbf{\hl{গ) ১৪৪ বর্গফুট ✔}}

ঘ) ৩২ বর্গফুট

\textbf{\hl{✔ সঠিক উত্তর: গ) ১৪৪ বর্গফুট}}

\begin{quote}
\textbf{ব্যাখ্যা:} \emph{বাহুর দৈর্ঘ্য = √(64)=8 ফুট}

\emph{রাস্তাসহ বর্গক্ষেত্রের বাহুর দৈর্ঘ্য = 8+(2×2)=12 ফুট}

\emph{∴ রাস্তাসহ বর্গক্ষেত্রের ক্ষেত্রফল = (12)\^{}2=144 বর্গফুট।}
\end{quote}

\textbf{একটি বর্গাকার বাগানের চারপাশ ঘিরে ২ মিটার প্রস্থবিশিষ্ট একটি রাস্তা
আছে। রাস্তাসহ বাগানের ক্ষেত্রফল ১৯৬ বর্গমিটার হলে, রাস্তার ক্ষেত্রফল কত?}

ক) ৪ বর্গমিটার

খ) ১৪৪ বর্গমিটার

গ) ৬০ বর্গমিটার

\textbf{\hl{ঘ) ৫২ বর্গমিটার ✔}}

\textbf{\hl{✔ সঠিক উত্তর: ঘ) ৫২ বর্গমিটার}}

\begin{quote}
\textbf{ব্যাখ্যা:} \emph{ধরি, বর্গাকার বাগানের একবাহুর দৈর্ঘ্য = x মিটার}

\emph{২ মিটার রাস্তাসহ বাগানের দৈর্ঘ্য = (x+2+2)=(x+4) মিটার}

\emph{প্রশ্নমতে, (x+4)\^{}2=196⇒x+4=14⇒x=10}

\emph{বাগানের ক্ষেত্রফল = 10\^{}2=100 বর্গমিটার}

\emph{∴ রাস্তার ক্ষেত্রফল = 196−100=52 বর্গমিটার (উল্লেখ্য, রাস্তাসহ দৈর্ঘ্য x+2
ধরলে সমীকরণ অনুযায়ী 196−144=52 হয়)।}
\end{quote}

\textbf{একটি বর্গের কর্ণের দৈর্ঘ্য 6√(2) একক হলে এর পরিসীমার অর্ধেক -}

ক) ৬ একক

\textbf{\hl{খ) ১২ একক ✔}}

গ) ২০ একক

ঘ) ২৪ একক

\textbf{\hl{✔ সঠিক উত্তর: খ) ১২ একক}}

\begin{quote}
\textbf{ব্যাখ্যা:} \emph{ধরি, বর্গের একবাহু a}

\emph{প্রশ্নমতে, a√(2)=6√(2)⇒a=6}

\emph{বর্গের পরিসীমা = 4×6=24 একক}

\emph{∴ পরিসীমার অর্ধেক = (24/2)=12 একক।}
\end{quote}

\textbf{একটি বর্গক্ষেত্রের পরিসীমা এর কর্ণের দৈর্ঘ্যের কত গুণ?}

ক) (√(2)/2)

খ) √(2)

\textbf{\hl{গ) 2√(2) ✔}}

ঘ) 2

\textbf{\hl{✔ সঠিক উত্তর: গ) 2√(2)}}

\begin{quote}
\textbf{ব্যাখ্যা:} \emph{বর্গক্ষেত্রের একবাহু = a , কর্ণ = √(2)a}

\emph{পরিসীমা = 4a=2×√(2)×√(2)a=2√(2)×কর্ণ ।}
\end{quote}

\textbf{একটি বর্গের পরিসীমা ২৪ মিটার হলে এর কর্ণ কত?}

ক) 8√(2)

খ) 4√(2)

\textbf{\hl{গ) 6√(2) ✔}}

ঘ) 3√(2)

\textbf{\hl{✔ সঠিক উত্তর: গ) 6√(2)}}

\begin{quote}
\textbf{ব্যাখ্যা:} \emph{ধরি, বর্গের একবাহুর দৈর্ঘ্য a}

\emph{প্রশ্নমতে, 4a=24⇒a=6}

\emph{∴ বর্গের কর্ণ = a√(2)=6√(2) ।}
\end{quote}

\textbf{একটি বর্গক্ষেত্রের একবাহু অপর এক বর্গক্ষেত্রের পরিসীমার সমান হলে,
বর্গক্ষেত্র দুটির কর্ণের দৈর্ঘ্যের অনুপাত কত?}

ক) ১:২

খ) ২:১

গ) ১:৮

\textbf{\hl{ঘ) ৪:১ ✔}}

\textbf{\hl{✔ সঠিক উত্তর: ঘ) ৪:১}}

\begin{quote}
\textbf{ব্যাখ্যা:} \emph{ধরি, একটি বর্গক্ষেত্রের একবাহুর দৈর্ঘ্য x এবং অপরটির
বাহু y}

\emph{শর্তমতে, প্রথমটির একবাহু x দ্বিতীয়টির পরিসীমা 4y -এর সমান, অর্থাৎ x=4y}

\emph{বর্গক্ষেত্র দুটির কর্ণের দৈর্ঘ্য যথাক্রমে √(2)x ও √(2)y}

\emph{∴ কর্ণের দৈর্ঘ্যের অনুপাত = √(2)x:√(2)y=x:y=4y:y=4:1}
\end{quote}

\textbf{একটি বর্গক্ষেত্রের বাহুর দৈর্ঘ্য ৮ ফুট হলে, ঐ বর্গক্ষেত্রের কর্ণের উপর
অঙ্কিত বর্গক্ষেত্রের ক্ষেত্রফল কত?}

ক) ১৫৬ বঃ ফুট

খ) ১৬৮ বঃ ফুট

\textbf{\hl{গ) ১২৮ বঃ ফুট ✔}}

ঘ) ২১৮ বঃ ফুট

\textbf{\hl{✔ সঠিক উত্তর: গ) ১২৮ বঃ ফুট}}

\begin{quote}
\textbf{ব্যাখ্যা:} \emph{(কর্ণের দৈর্ঘ্য)² = 8\^{}2+8\^{}2=64+64=128 = কর্ণের
উপর অঙ্কিত বর্গক্ষেত্রের ক্ষেত্রফল।}
\end{quote}

\textbf{একটি বর্গক্ষেত্রের এক বাহুর দৈর্ঘ্য ৩ ফুট থেকে বৃদ্ধি করে ৬ ফুট করা হলে
ক্ষেত্রফল কত গুণ বাড়বে?}

ক) ২ গুণ

খ) ৩ গুণ

\textbf{\hl{গ) ৪ গুণ ✔}}

ঘ) কোনোটাই নয়

\textbf{\hl{✔ সঠিক উত্তর: গ) ৪ গুণ}}

\begin{quote}
\textbf{ব্যাখ্যা:} \emph{প্রাথমিক ক্ষেত্রফল = 3\^{}2=9 বর্গফুট}

\emph{নতুন ক্ষেত্রফল = 6\^{}2=36 বর্গফুট}

\emph{∴ ক্ষেত্রফল বাড়ে = (36/9)=4 গুণ।}
\end{quote}

\textbf{একটি সরলরেখার উপর অঙ্কিত বর্গ ঐ সরলরেখার অর্ধেকের উপর অঙ্কিত বর্গের
কত গুণ?}

ক) দ্বিগুণ

খ) তিনগুণ

\textbf{\hl{গ) চারগুণ ✔}}

ঘ) পাঁচগুণ

\textbf{\hl{✔ সঠিক উত্তর: গ) চারগুণ}}

\begin{quote}
\textbf{ব্যাখ্যা:} \emph{ধরি, সরলরেখাটির দৈর্ঘ্য x একক}

\emph{এর উপর অঙ্কিত বর্গক্ষেত্রের ক্ষেত্রফল x\^{}2 বর্গ একক}

\emph{অর্ধেকের উপর অঙ্কিত বর্গক্ষেত্রের ক্ষেত্রফল = ((x/2))\^{}2=(x\^{}2/4) বর্গ
একক}

\emph{∴ বড় ক্ষেত্রের ক্ষেত্রফল ছোট ক্ষেত্রের ক্ষেত্রফলের 4 গুণ।}
\end{quote}

\textbf{একটি রেখাংশের উপর অঙ্কিত বর্গক্ষেত্র ঐ রেখাংশের এক-তৃতীয়াংশের উপর
অঙ্কিত বর্গক্ষেত্রের কত গুণ?}

ক) ৩ গুণ

খ) ৬ গুণ

\textbf{\hl{গ) ৯ গুণ ✔}}

ঘ) ১২ গুণ

\textbf{\hl{✔ সঠিক উত্তর: গ) ৯ গুণ}}

\begin{quote}
\textbf{ব্যাখ্যা:} \emph{ধরি, রেখাংশটির দৈর্ঘ্য = x একক}

\emph{পূর্ণ রেখাংশের উপর অঙ্কিত বর্গক্ষেত্রের ক্ষেত্রফল = x\^{}2 বর্গ একক}

\emph{এক-তৃতীয়াংশের উপর অঙ্কিত বর্গক্ষেত্রের ক্ষেত্রফল =
((x/3))\^{}2=(x\^{}2/9) বর্গ একক}

\emph{∴ মূল বর্গক্ষেত্রের ক্ষেত্রফল ছোট বর্গক্ষেত্রের ক্ষেত্রফলের 9 গুণ।}
\end{quote}

\textbf{একটি সরলরেখার উপর অঙ্কিত বর্গের ক্ষেত্রফল ঐ সরলরেখার ২৫ শতাংশের উপর
অঙ্কিত বর্গের ক্ষেত্রফলের কত গুণ?}

ক) ৪ গুণ

\textbf{\hl{খ) ১৬ গুণ ✔}}

গ) ৮ গুণ

ঘ) ২ গুণ

\textbf{\hl{✔ সঠিক উত্তর: খ) ১৬ গুণ}}

\begin{quote}
\textbf{ব্যাখ্যা:} \emph{২৫\% = (25/100)=(1/4) (অর্থাৎ এক-চতুর্থাংশ)}

\emph{মূল বর্গক্ষেত্রের ক্ষেত্রফল = ((4/1))\^{}2×পরিবর্তিত বর্গক্ষেত্রের ক্ষেত্রফল∴
16 গুণ।}
\end{quote}

\textbf{১ বর্গাকার জমির ভূমির দৈর্ঘ্য ২০\% বাড়ালে এবং প্রস্থ ১০\% কমালে
ক্ষেত্রফল কতটুকু হ্রাস বা বৃদ্ধি পাবে?}

ক) ২\% বৃদ্ধি

খ) ৩\% বৃদ্ধি

গ) ১০\% বৃদ্ধি

\textbf{\hl{ঘ) ৮\% বৃদ্ধি ✔}}

\textbf{\hl{✔ সঠিক উত্তর: ঘ) ৮\% বৃদ্ধি}}

\begin{quote}
\textbf{ব্যাখ্যা:} \emph{ক্ষেত্রফল হ্রাস/বৃদ্ধি = (a−b−(a×b/100))\%}

\emph{= (20−10−(20×10/100))\%}

\emph{= 10\%−2\%=8\% বৃদ্ধি।}
\end{quote}

\textbf{If the side of a square increases by 40\%, then the area of the
square increases by-}

ক) 16\%

\textbf{\hl{খ) 96\% ✔}}

গ) 116\%

ঘ) 140\%

\textbf{\hl{✔ সঠিক উত্তর: খ) 96\%}}

\begin{quote}
\textbf{ব্যাখ্যা:} \emph{ধরি, বাহুর দৈর্ঘ্য = x}

\emph{৪০\% বৃদ্ধিতে নতুন বাহু = 1.4x}

\emph{প্রাথমিক ক্ষেত্রফল = x\^{}2}

\emph{নতুন ক্ষেত্রফল = (1.4x)\^{}2=1.96x\^{}2}

\emph{∴ ক্ষেত্রফল বৃদ্ধি পায় = (1.96x\^{}2−x\^{}2)=0.96x\^{}2 অর্থাৎ ৯৬\%।}
\end{quote}

\textbf{The length of a rectangle is increased by 25\%. By what
percentage should the width be decreased so that the area of the
rectangle remains unchanged?}

ক) ৩০\%

খ) ২৫\%

\textbf{\hl{গ) ২০\% ✔}}

ঘ) ৩৫\%

\textbf{\hl{✔ সঠিক উত্তর: গ) ২০\%}}

\begin{quote}
\textbf{ব্যাখ্যা:} \emph{ধরি, প্রস্থ হ্রাস করতে হবে x\%}

\emph{প্রশ্নমতে, 25−x−(25x/100)=0}

\emph{⇒ 25−(5x/4)=0⇒(5x/4)=25⇒x=20}

\emph{∴ ২০\% হ্রাস করতে হবে।}
\end{quote}

\textbf{কোনো বর্গক্ষেত্রের এক বাহুর দৈর্ঘ্য ৫০\% বৃদ্ধি করলে ক্ষেত্রফল শতকরা কত
বৃদ্ধি পাবে?}

ক) ৫০\%

খ) ১০০\%

\textbf{\hl{গ) ১২৫\% ✔}}

ঘ) ১৫০\%

\textbf{\hl{✔ সঠিক উত্তর: গ) ১২৫\%}}

\begin{quote}
\textbf{ব্যাখ্যা:} \emph{ধরি, বর্গক্ষেত্রের বাহুর দৈর্ঘ্য k একক}

\emph{প্রাথমিক ক্ষেত্রফল = k\^{}2 বর্গ একক}

\emph{নতুন দৈর্ঘ্য = k+(50/100)k=(3/2)k}

\emph{নতুন ক্ষেত্রফল = ((3/2)k)\^{}2=(9/4)k\^{}2}

\emph{ক্ষেত্রফল বৃদ্ধি = (9/4)k\^{}2−k\^{}2=(5/4)k\^{}2}

\emph{∴ শতকরা বৃদ্ধি = ((5/4)k\^{}2/k\^{}2)×100\%=125\% ।}
\end{quote}

\textbf{একটি বর্গের বাহুর দৈর্ঘ্য ৮০ মিটার প্রত্যেক বাহুর দৈর্ঘ্য ১০\% বাড়লে
ক্ষেত্রফল কত বাড়বে?}

\textbf{\hl{ক) ২১\% ✔}}

খ) ২৫\%

গ) ৩০\%

ঘ) ৩৫\%

\textbf{\hl{✔ সঠিক উত্তর: ক) ২১\%}}

\begin{quote}
\textbf{ব্যাখ্যা:} \emph{ক্ষেত্রফল শতকরা বৃদ্ধি পায় =
(10\%+10\%+(10×10/100)\%)}

\emph{= 20\%+1\%=21\% ।}
\end{quote}

\textbf{কোন বর্গক্ষেত্রের প্রতি বাহু যদি ১০\% বৃদ্ধি পায়, তবে তার ক্ষেত্রফল শতকরা
কত বৃদ্ধি পাবে?}

ক) ১৯

খ) ২০

\textbf{\hl{গ) ২১ ✔}}

ঘ) ২২

\textbf{\hl{✔ সঠিক উত্তর: গ) ২১}}

\begin{quote}
\textbf{ব্যাখ্যা:} \emph{ধরি, বর্গ ক্ষেত্রের একবাহু = 100 একক}

\emph{বর্গ ক্ষেত্রের ক্ষেত্রফল = 100×100=10000 বর্গ একক}

\emph{১০\% বৃদ্ধি পাওয়ার পর বর্গ ক্ষেত্রের প্রতি বাহু = 110 একক}

\emph{বর্ধিত বর্গক্ষেত্রের ক্ষেত্রফল = 110×110=12100 বর্গ একক}

\emph{ক্ষেত্রফল বৃদ্ধি পেয়েছে = 12100−10000=2100 বর্গ একক অর্থাৎ শতকরা ২১\%।}
\end{quote}

\textbf{চারটি সমান বাহু দ্বারা সীমাবদ্ধ একটি ক্ষেত্র যার একটি কোণও সমকোণ নয়,
এরূপ চিত্রকে বলা হয়-\/- /১২তম বিসিএস/ কোন চতুর্ভুজের বাহুগুলো সমান এবং কোণগুলো
সমকোণ নয়। এরূপ চতুর্ভুজকে বলে -\/-}

ক) বর্গক্ষেত্র (Square)

খ) চতুর্ভুজ (Quadrilateral)

\textbf{\hl{গ) রম্বস (Rhombus) ✔}}

ঘ) সামান্তরিক (Parallelogram)

\textbf{\hl{✔ সঠিক উত্তর: গ) রম্বস (Rhombus)}}

\begin{quote}
\textbf{ব্যাখ্যা:} \emph{যে চতুর্ভুজের চারটি বাহুই সমান কিন্তু কোণগুলো সমকোণ নয়,
তাকে রম্বস (Rhombus) বলে।}
\end{quote}

\textbf{রম্বসের ক্ষেত্রফল কত?}

ক) (1/2)×(সমান্তরাল বাহুদ্বয়ের দৈর্ঘ্যের সমষ্টি)×উচ্চতা

\textbf{\hl{খ) (1/2)×(কর্ণদ্বয়ের গুণফল) ✔}}

গ) (1/2)×(কর্ণদ্বয়ের সমষ্টি)

ঘ) (1/2)×(ভূমি×উচ্চতা)

\textbf{\hl{✔ সঠিক উত্তর: খ) (1/2)×(কর্ণদ্বয়ের গুণফল)}}

\begin{quote}
\textbf{ব্যাখ্যা:} \emph{রম্বসের ক্ষেত্রফল = (1/2)×(কর্ণদ্বয়ের গুণফল) ।}
\end{quote}

\textbf{একটি রম্বসের কর্ণদ্বয় ৫ সেমি ও ৬ সেমি হলে, এর ক্ষেত্রফল কত?}

ক) ৩০ বর্গ সেমি

খ) ২৫ বর্গ সেমি

গ) ২০ বর্গ সেমি

\textbf{\hl{ঘ) ১৫ বর্গ সেমি ✔}}

\textbf{\hl{✔ সঠিক উত্তর: ঘ) ১৫ বর্গ সেমি}}

\begin{quote}
\textbf{ব্যাখ্যা:} \emph{রম্বসের ক্ষেত্রফল = (1/2)×কর্ণদ্বয়ের গুণফল=
(1/2)×5×6=15 বর্গ সেমি।}
\end{quote}

\textbf{যে চতুর্ভুজের বাহুগুলি পরস্পর সমান ও সমান্তরাল কিন্তু কোণগুলো সমকোণ নয়
তাকে কি বলে?}

ক) সামান্তরিক

\textbf{\hl{খ) রম্বস ✔}}

গ) ট্রাপিজিয়াম

ঘ) আয়তক্ষেত্র

ঙ) কোনোটাই নয়

\textbf{\hl{✔ সঠিক উত্তর: খ) রম্বস}}

\begin{quote}
\textbf{ব্যাখ্যা:} \emph{যে চতুর্ভুজের বাহুগুলো পরস্পর সমান ও সমান্তরাল কিন্তু
কোণগুলো সমকোণ নয়, তাকে রম্বস বলে।}
\end{quote}

\textbf{একটি রম্বসের কর্ণদ্বয় ৪০ সেমি ও ৬০ সেমি। এর ক্ষেত্রফল কত বর্গ সেমি?}

ক) ২৪০০ বর্গ সেমি

\textbf{\hl{খ) ১২০০ বর্গ সেমি ✔}}

গ) ১৪৪ বর্গ সেমি

ঘ) ৩৬০০ বর্গ সেমি

\textbf{\hl{✔ সঠিক উত্তর: খ) ১২০০ বর্গ সেমি}}

\begin{quote}
\textbf{ব্যাখ্যা:} \emph{রম্বসের ক্ষেত্রফল = (1/2)×কর্ণদ্বয়ের গুণফল=
(1/2)×60×40=1200 বর্গ সেমি।}
\end{quote}

\textbf{একটি রম্বসের একটি কর্ণ ১০ মিটার এবং ক্ষেত্রফল ১২০ বর্গমিটার হলে, অপর
কর্ণের দৈর্ঘ্য কত মিটার?}

ক) ২০ মিটার

খ) ২২ মিটার

\textbf{\hl{গ) ২৪ মিটার ✔}}

ঘ) ২৬ মিটার

\textbf{\hl{✔ সঠিক উত্তর: গ) ২৪ মিটার}}

\begin{quote}
\textbf{ব্যাখ্যা:} \emph{আমরা জানি, রম্বসের ক্ষেত্রফল = (1/2)×কর্ণদ্বয়ের গুণফল⇒
120=(1/2)×10×অপর কর্ণ}

\emph{⇒ অপর কর্ণ=(240/10)=24 মিটার।}
\end{quote}

\textbf{PQRS রম্বসের বাহুর দৈর্ঘ্য ৩ ইঞ্চি। PR এবং QS কর্ণ দুটি বিন্দুতে ছেদ করলে
PO\^{}2+QO\^{}2 কত?}

ক) ৮

\textbf{\hl{খ) ৯ ✔}}

গ) ১০

ঘ) ১২

\textbf{\hl{✔ সঠিক উত্তর: খ) ৯}}

\begin{quote}
\textbf{ব্যাখ্যা:} \emph{রম্বসের কর্ণদ্বয় পরস্পরকে সমকোণে সমদ্বিখণ্ডিত করে বলে
ΔPOQ একটি সমকোণী ত্রিভুজ।}

\emph{পিথাগোরাসের উপপাদ্য অনুসারে, PO\^{}2+QO\^{}2=PQ\^{}2=3\^{}2=9}

\emph{∴ PO\^{}2+QO\^{}2=9 ।}
\end{quote}

\textbf{ABCD রম্বসের AC ও BD কর্ণ দুটিী O বিন্দুতে ছেদ করছে। ∠ACD=60\^{}°
হলে, ∠ODC= কত?}

ক) 60\^{}°

\textbf{\hl{খ) 30\^{}° ✔}}

গ) 90\^{}°

ঘ) 45\^{}°

\textbf{\hl{✔ সঠিক উত্তর: খ) 30\^{}°}}

\begin{quote}
\textbf{ব্যাখ্যা:} \emph{রম্বসের কর্ণদ্বয় পরস্পরকে সমকোণে সমদ্বিখণ্ডিত করে,
অর্থাৎ ∠DOC=90\^{}° ।}

\emph{ত্রিভুজ DOC -এ, তিন কোণের সমষ্টি 180\^{}° ।}

\emph{সুতরাং, ∠ODC=180\^{}°−(90\^{}°+60\^{}°)=180\^{}°−150\^{}°=30\^{}°
।}
\end{quote}

\textbf{ABCD রম্বসের ∠B=80\^{}° হলে ∠C= কত?}

ক) 120\^{}°

খ) 70\^{}°

গ) 80\^{}°

\textbf{\hl{ঘ) 100\^{}° ✔}}

\textbf{\hl{✔ সঠিক উত্তর: ঘ) 100\^{}°}}

\begin{quote}
\textbf{ব্যাখ্যা:} \emph{রম্বসের দুটি সংলগ্ন কোণের সমষ্টি 180\^{}° ।}

\emph{∴ ∠C=180\^{}°−80\^{}°=100\^{}° ।}
\end{quote}

\textbf{একটি রম্বসের প্রতিটি বাহুর দৈর্ঘ্য ১৩ সে.মি., ১টি কর্ণের দৈর্ঘ্য ২৪
সে.মি. হলে রম্বসটির অপর কর্ণের দৈর্ঘ্য কত?}

ক) ৮ সে.মি.

\textbf{\hl{খ) ১০ সে.মি. ✔}}

গ) ১৫ সে.মি.

ঘ) ২২ সে.মি.

\textbf{\hl{✔ সঠিক উত্তর: খ) ১০ সে.মি.}}

\begin{quote}
\textbf{ব্যাখ্যা:} \emph{রম্বসের কর্ণদ্বয় পরস্পরকে সমকোণে সমদ্বিখণ্ডিত করে।}

\emph{অর্ধেক কর্ণ = (24/2)=12 সে.মি.}

\emph{পিথাগোরাসের উপপাদ্য অনুসারে অপর কর্ণের অর্ধাংশ =
√(13\^{}2−12\^{}2)=√(169−144)=√(25)=5 সে.মি.}

\emph{∴ অপর কর্ণের দৈর্ঘ্য = 5×2=10 সে.মি.}
\end{quote}

\textbf{ABCD রম্বসের ∠ABC=120\^{}° এবং OE⊥AB হলে, ∠AEO= ?}

ক) 30\^{}°

খ) 60\^{}°

\textbf{\hl{গ) 90\^{}° ✔}}

ঘ) 120\^{}°

\textbf{\hl{✔ সঠিক উত্তর: গ) 90\^{}°}}

\begin{quote}
\textbf{ব্যাখ্যা:} \emph{যেহেতু OE⊥AB , তাই লম্ব রেখাটি ভূমির সাথে 90\^{}°
কোণ উৎপন্ন করে।}

\emph{∴ ∠AEO=90\^{}° ।}
\end{quote}

\textbf{যে চতুর্ভুজের একজোড়া বাহু পরস্পর সমান্তরাল তাকে কী বলে?}

ক) সামান্তরিক

খ) আয়ত

গ) রম্বস

\textbf{\hl{ঘ) ট্রাপিজিয়াম ✔}}

\textbf{\hl{✔ সঠিক উত্তর: ঘ) ট্রাপিজিয়াম}}

\begin{quote}
\textbf{ব্যাখ্যা:} \emph{যে চতুর্ভুজের একজোড়া বাহু পরস্পর সমান্তরাল, তাকে
ট্রাপিজিয়াম বলে।}
\end{quote}

\textbf{যে চতুর্ভুজের দুটি বাহু সমান্তরাল এবং অপর দুটি বাহু তীর্যক তাকে বলে -}

ক) রম্বস

\textbf{\hl{খ) ট্রাপিজিয়াম ✔}}

গ) চতুর্ভুজ

ঘ) আয়তক্ষেত্র

\textbf{\hl{✔ সঠিক উত্তর: খ) ট্রাপিজিয়াম}}

\begin{quote}
\textbf{ব্যাখ্যা:} \emph{যে চতুর্ভুজের দুটি বাহু সমান্তরাল এবং অপর দুটি বাহু
তীর্যক, তাকে ট্রাপিজিয়াম বলে।}
\end{quote}

\textbf{একটি চতুর্ভুজের চারটি বাহুর বিপরীত দুটি সমান্তরাল কিন্তু অসমান। একে বলে
-}

ক) সামান্তরিক

খ) বর্গক্ষেত্র

\textbf{\hl{গ) ট্রাপিজিয়াম ✔}}

ঘ) আয়তক্ষেত্র

\textbf{\hl{✔ সঠিক উত্তর: গ) ট্রাপিজিয়াম}}

\begin{quote}
\textbf{ব্যাখ্যা:} \emph{যে চতুর্ভুজের বিপরীত বাহুর একজোড়া সমান্তরাল কিন্তু অপর
জোড়া সমান্তরাল নয় (বা অসমান), তাকে ট্রাপিজিয়াম বলে।}
\end{quote}

\textbf{ট্রাপিজিয়ামের দুটি সমান্তরাল বাহু a সে.মি. ও b সে.মি.। তাদের মধ্যবর্তী
দূরত্ব h সে.মি. হলে, এর ক্ষেত্রফল কত?}

ক) (a+b)h

খ) 2(a+b)h

\textbf{\hl{গ) (1/2)×(a+b)×h ✔}}

ঘ) (1/2)×(a−b)×h

\textbf{\hl{✔ সঠিক উত্তর: গ) (1/2)×(a+b)×h}}

\begin{quote}
\textbf{ব্যাখ্যা:} \emph{ট্রাপিজিয়ামের ক্ষেত্রফল = (1/2)×(সমান্তরাল বাহুদ্বয়ের
যোগফল)×উচ্চতা= (1/2)×(a+b)×h ।}
\end{quote}

\textbf{ট্রাপিজিয়ামের ক্ষেত্রফল নির্ণয়ের সূত্র কোনটি?}

\textbf{\hl{ক) (1/2)(a+b)h ✔}}

খ) (√(3)/4)a\^{}2

গ) 6a\^{}2

ঘ) πr\^{}2h

\textbf{\hl{✔ সঠিক উত্তর: ক) (1/2)(a+b)h}}

\begin{quote}
\textbf{ব্যাখ্যা:} \emph{ট্রাপিজিয়ামের ক্ষেত্রফলের সূত্রটি হলো: (1/2)×(সমান্তরাল
বাহুদ্বয়ের যোগফল)×উচ্চতা= (1/2)(a+b)h ।}
\end{quote}

\textbf{একটি ট্রাপিজিয়ামের ক্ষেত্রফল ২৫ বর্গফুট। ট্রাপিজিয়ামের দুই বাহুর মধ্যবর্তী
উচ্চতা ২ ফুট এবং সমান্তরাল দুই বাহুর একটি অপরটি হতে ১ ফুট বেশি হলে বৃহত্তর বাহুর
দৈর্ঘ্য কত?}

ক) 10 feet

খ) 12 feet

\textbf{\hl{গ) 13 feet ✔}}

ঘ) 25 feet

\textbf{\hl{✔ সঠিক উত্তর: গ) 13 feet}}

\begin{quote}
\textbf{ব্যাখ্যা:} \emph{ধরি, সমান্তরাল ক্ষুদ্রতর বাহুর দৈর্ঘ্য x ফুট}

\emph{সমান্তরাল বৃহত্তর বাহুর দৈর্ঘ্য (x+1) ফুট}

\emph{প্রশ্নমতে, (1/2)×2×(x+x+1)=25}

\emph{⇒ 2x+1=25⇒2x=24⇒x=12}

\emph{∴ বৃহত্তর বাহু = 12+1=13 ফুট।}
\end{quote}

\section{17. জ্যামিতি-\/-বিন্দু , রেখা এবং কোণ \textbar\textbar{} A to
Z}\label{ux99cux9afux9aeux9a4--ux9acux9a8ux9a6-ux9b0ux996-ux98fux9ac-ux995ux9a3-a-to-z}

জ্যামিতি হল স্পর্শস্থান এবং তার সম্পর্কিত বিষয় সম্পর্কে বিশেষভাবে গবেষণা করা ।
জ্যামিতির তিনটি মৌলিক উপাদান হল বিন্দু, রেখা এবং কোণ।

বিন্দু হল একটি স্থান, যার কোন আকার না থাকলেও স্থানীয় স্থান নির্দেশ করে। রেখা
হল দুটি বিন্দুকে সংযোজন করা একটি সরল সেগমেন্ট এবং এটির দৈর্ঘ্য হল বিন্দু দুটির
মধ্যবর্তী দূরত্ব।

কোণ হল দুটি রেখার মধ্যে উৎপন্ন বাহুর মাধ্যমে উৎপন্ন একটি পরিমাপযুক্ত জোড়া রেখা।
জ্যামিতির উপর ভিত্তি করে বিভিন্ন ধরনের আকার ও গণিতগত প্রসঙ্গ নির্ণয় করা হয়।

জ্যামিতি ব্যবহার করে বিভিন্ন ফিজিক্স, গণিত এবং প্রকৃতি বিজ্ঞানের প্রযুক্তি তৈরি
করা হয়। জ্যামিতি আমাদের জীবনে কতটা গুরুত্বপূর্ণ হবে সেটি আমাদের সবার আগে জানতে
হবে।

\begin{itemize}
\item
  প্রাথমিক আলোচনা
\item
  বিন্দু (Points)
\item
  রেখা (Line)
\item
  কোণ (Angle)
\end{itemize}

\subsubsection{প্রাথমিক
আলোচনা}\label{ux9aaux9b0ux9a5ux9aeux995-ux986ux9b2ux99aux9a8-6}

জ্যামিতি গণিত শাস্ত্রের একটি প্রাচীন শাখা। যদিও ব্যুৎপত্তিগতভাবে "জ্যামিতি"
শব্দের অর্থ ভূমির পরিমাপ তথা জ্যামিতি, প্রকৃত পক্ষে স্থান বিষয়ক বিজ্ঞান (Geometry
is the science concerned with space)

ভূমি পরিমাপের জন্য জ্যামিতির উদ্ভব হলেও বর্তমানে জ্যামিতি কেবল ভূমি পরিমাপের
জন্যই ব্যবহৃত হয় না বরং বহু জটিল গাণিতিক সমস্যা সমাধানে জ্যামিতিক জ্ঞান
অপরিহার্য।

জ্যামিতিক জ্ঞান আমাদের জীবনের সমস্যা সমাধানের ইঙ্গিত দেয়, চিন্তা শক্তির উন্মেষ
ও মননশীলতার উৎকর্ষ সাধন করে থাকে।

জ্যামিতি বা Geometry শব্দের \textquotesingle জ্যা\textquotesingle{} এর অর্থ
ভূমি এবং \textquotesingle মিতি\textquotesingle{} এর অর্থ পরিমাপ অর্থাৎ
\textquotesingle ভূমির পরিমাপ\textquotesingle। (The word
\textquotesingle geometry\textquotesingle{} means \textquotesingle earth
measure\textquotesingle)

আরো বলা যায়, ``গণিতশাস্ত্রের যে শাখায় বিন্দু থেকে বৃত্ত পর্যন্ত যাবতীয় ক্ষেত্রের
বৈশিষ্ট্য ধর্ম এবং ক্ষেত্রের পরিমাপ সম্বন্ধে আলোচিত হয় তারই নাম জ্যামিতি।
খ্রিষ্টপূর্ব ৩০০ সালে ইউক্লিড তার The Elements গ্রন্থে জ্যামিতির মূল তত্বগুলো
আলোচনা করেছেন।

\subsubsection{বিন্দু (Points)}\label{ux9acux9a8ux9a6-points}

যার কেবল অবস্থান আছে কিন্তু দৈর্ঘ্য, প্রস্থ, উচ্চতা বা বেধ বলতে কিছুই নেই তাকে
বিন্দু বলে। দৈর্ঘ্য, প্রস্থ এবং উচ্চতা কোনো কিছুই নেই বলে বিন্দুর কোনো মাত্রা নেই
অর্থাৎ, বিন্দুর মাত্রা শুন্য।

বিন্দুর শুধু অবস্থান আছে কিন্তু কোন মাত্রা নেই ।

দুটি বিন্দুর মধ্য দিয়ে একটি এবং কেবলমাত্র সরলরেখা আঁকা যায়।

দুটি সরলরেখা একটি এবং কেবল একটি বিন্দুতে পরস্পরকে ছেদ করতে পারে।

দুই বিন্দুর মধ্যে সরলরেখার দূরত্বই ক্ষুদ্রতম ।

যেসব বিন্দু একই সরলরেখায় অবস্থান করে, তাদেরকে সমরেখ বিন্দু বলে ।

\subsubsection{রেখা (Line)}\label{ux9b0ux996-line}

\textbf{রেখা :} বিন্দুর চলার পথকে রেখা বলে ।

\textbf{সমরেখ:}

যখন দুই-এর অধিক বিন্দু একই সরলরেখার উপর অবস্থিত হয়, তখন ঐ বিন্দুগুলিকে সমরেখ
বিন্দু বলা হয়। চিত্রে

A, B, C বিন্দু তিনটি সমরেখ ।

রেখার কোন প্রান্ত বিন্দু নেই ।

রেখাংশের দুটি প্রান্তবিন্দু থাকে ।

রশ্মির একটিমাত্র প্রান্তবিন্দু থাকে।

\textbf{সরলরেখা (Straight Line):}

একটি বিন্দু থেকে অন্য একটি বিন্দুতে পৌছাতে যদি কোনো প্রকার দিকের পরিবর্তন না হয়
তবে তাকে সরল রেখা বলে ।

দুটি বিন্দু দিয়ে একটি মাত্র সরলরেখা আঁকা যায় ।

একটি সরলরেখার বহিঃস্থ যে কোনো বিন্দু থেকে ঐ সরলরেখা পর্যন্ত যেসব সরলরেখা আঁকা
যায়, তাদের মধ্যে লম্বের দৈর্ঘ্যই ক্ষুদ্রতম ।

\textbf{বক্র রেখা (Zigzag line ):} একটি বিন্দু থেকে অন্য একটি বিন্দুতে পৌছাতে
যদি দিক পরিবর্তন হয় তবে তাকে বক্র রেখা বলে ।

\textbf{সমান্তরাল রেখা (Parallel Lines):}

দুটি রেখা যদি পরস্পরের মধ্যে সর্বদা সমান দূরত্ব বজায় রেখে চলতে থাকে তবে তাকে
সমান্তরাল রেখা বলে ।

কোন সরলরেখার বহিঃস্থ কোন বিন্দু থেকে উক্ত রেখা পর্যন্ত যতগুলো রেখাংশ টানা যায়
তন্মধ্যে লম্ব রেখাংশটিই ক্ষুদ্রতম ।

দুইটি সমান্তরাল সরলরেখার কোন সাধারণ বিন্দু নেই ।

দুটি পরস্পরচ্ছেদী সরলরেখার প্রত্যেকটি, কোন তৃতীয় সরলরেখার সমান্তরাল হতে পারে
না।

যেসব সরলরেখা একই সরলরেখার সমান্তরাল তারা পরস্পর সমান্তরাল। দুই বা ততোধিক
সরলরেখা একটি সরলরেখার উপর লম্ব হলে, তারা পরস্পর সমান্তরাল ।

একটি সরলরেখা সমান্তরাল রেখাদ্বয়ের একটির উপর লম্ব হলে, তা অপরটির উপরও লম্ব।

কোণ (Angle): \textbf{কোণ (Angle):} দুটি সরলরেখা তির্যকভাবে (আড়াআড়িভাবে)
পরস্পরের সাথে মিলিত হলে মিলিত বিন্দুতে কোণ উৎপন্ন হয়। নিচের চিত্র দেখুন ।

চিত্র : কোণ

\textbf{সমকোণ (Right Angle):}

যে কোণের পরিমাণ 90 ডিগ্রী তাকে সমকোণ বলে । নিচের চিত্র দেখুন ।

চিত্র : সমকোণ

\textbf{মনে রাখুন}: কোণটি দেখতে যে রকমই হোক না কেন তার পরিমাপ 90 ডিগ্রী
হলেই তাকে সমকোণ বলা হবে।

\textbf{স্থূলকোণ (Obtuse Angle ) :} যে কোণের পরিমাণ ৯০ ডিগ্রী অপেক্ষা বেশী
এবং ১৮০ ডিগ্রি অপেক্ষা কম তাকে স্থূলকোণ বলে । নিচের চিত্র দেখুন ।

চিত্র : স্থূলকোণ

\textbf{সূক্ষ্মকোণ (Acute Angle ) :} যে কোণের পরিমাণ 90 ডিগ্রী অপেক্ষা কম
তাকে সূক্ষ্মকোণ বলে । নিচের চিত্র দেখুন ।

\textbf{সরলকোণ (Straight Angle ) :} যে কোণের পরিমাণ ১৮০ ডিগ্ৰী তাকে
সরলকোণ বলে । নিচের চিত্র দেখুন ।

\textless\_\_\_\_\_\_\_\_\_\_\_\_\_\_\_\_\_\_\_\_\_\textgreater{}

চিত্র : সরলকোণ

\textbf{প্রবৃদ্ধ কোণ (Reflex Angle ) :} যে কোণের পরিমাণ দুই সমকোণ ( 180
ডিগ্রি ) অপেক্ষা বেশী কিন্তু চার সমকোণ ( 360 ডিগ্রি ) অপেক্ষা কম তাকে প্রবৃদ্ধ
কোণ বলে। নিচের চিত্র দেখুন ।

চিত্র : প্রবৃদ্ধ কোণ

\textbf{বিপ্রতীপ কোণ (Vertically Opposite Angle) :} দুইটি সরলরেখা পরস্পরকে
ছেদ করলে ছেদবিন্দুতে যে চারটি কোণ উৎপন্ন হয় তাদের একটিকে অপরটির বিপ্রতীপ কোণ
বলে।

\textbf{পুরক কোণ (Complementary Angle) :} দুইটি কোণের সমষ্টি এক সমকোণ বা ৯০
ডিগ্রি হলে একটিকে অপরটির পূরক কোণ বলে।

\textbf{সম্পুরক কোণ (Complementary Angle) :} দুইটি কোণের সমষ্টি দুই সমকোণ বা
১৮০ ডিগ্রি হলে একটিকে অপরটির সম্পুরক কোণ বলে।

\textbf{সন্নিহিত কোণ (Adjacent Angle):} দুইটি কোণের শীর্ষবিন্দু একই হলে
সাধারণ বাহুর উভয় পাশের কোণদ্বয়কে সন্নিহিত কোণ বলে । অর্থাৎ পাশাপাশি লাগানো
দুটি কোণকে সন্নিহিত কোণ বলে ।

\textbf{অন্তর্ভুক্ত কোণ:} দুটি বাহুর মিলিত হওয়ার স্থানকে ঐ বাহুদ্বয়ের অন্তর্ভুক্ত কোণ
বলা হয়।

\textbf{একান্তর\textquotesingle{} কোণ (Alternative Angle):} দুটি সমান্তরাল
সরলরেখাকে অপর একটি রেখা তির্যকভাবে (আড়াআড়ি) ছেদ করলে ছেদক রেখার বিপরীত
পার্শ্বে সমান্তরাল রেখা যে কোণ উৎপন্ন করে তাকে একান্তর কোণ বলে ।

একান্তর কোণদ্বয় পরস্পর সমান হয়।

\textbf{অনুরূপ কোণ (Corresponding Angle) :} দুটি সমান্তরাল সরলরেখাকে অপর
একটি সরলরেখা তির্যকভাবে ছেদ করলে ছেদক রেখার একই দিকে সমান্তরাল রেখাদ্বয়ের
অনুরূপ পার্শ্বে যে কোণ উৎপন্ন হয়, তাকে অনুরূপ কোণ বলে ।

\textbf{অনুরুপ কোণ দেখতে একই রকম এবং অনুরুপ কোণগুলো পরস্পর সমান ।}

\textbf{অন্তঃস্থ কোণ (Interior Angles):} যে কোন ত্রিভুজ, চতুর্ভুজ অথবা বহুভুজের
অভ্যন্তরে বা ভেতরে যে কোণ উৎপন্ন হয় তাকে অন্তঃস্থ কোণ বলে ।

\textbf{বহিঃস্থ কোণ (Exterior Angles) :} কোন ত্রিভুজ, চতুর্ভুজ অথবা বহুভুজের যে
কোণ এক বাহু বর্ধিত করলে বাহিরে যে কোণ উৎপন্ন হয় তাকে বহিঃস্থ কোণ বলে ।

\section{18. বিন্দু, রেখা ও কোণ MCQ প্রশ্ন উত্তর
অনুশীলন}\label{ux9acux9a8ux9a6-ux9b0ux996-ux993-ux995ux9a3-mcq-ux9aaux9b0ux9b6ux9a8-ux989ux9a4ux9a4ux9b0-ux985ux9a8ux9b6ux9b2ux9a8}

\textbf{কোন দেশে সর্বপ্রথম জ্যামিতি আলোচনা শুরু হয়?}

ক) গ্রিসে

খ) চীনে

\textbf{\hl{গ) মিশরে ✔}}

ঘ) জাপানে

\textbf{\hl{✔ সঠিক উত্তর: গ) মিশরে}}

\begin{quote}
\textbf{ব্যাখ্যা:} \emph{ঐতিহাসিকদের মতে প্রাচীন মিশরে আনুমানিক চার হাজার
বছর আগে ভূমি জরিপের কাজে জ্যামিতিক ধ্যান ধারণা ব্যবহার করা হয়। তবে খ্রিস্টপূর্ব
৩০০ অব্দে গ্রীক পণ্ডিত ইউক্লিড সর্বপ্রথম জ্যামিতির মৌলিক ধারণার ভিত্তি রচনা
করেন।}
\end{quote}

\textbf{AB রেখাংশের উপর P একটি বিন্দু হলে কোন সম্পর্ক সবসময় প্রযোজ্য?}

ক) AP=PB

\textbf{\hl{খ) AB\textgreater AP ✔}}

গ) AB\textgreater AP+PB

ঘ) AP\textgreater PB

\textbf{\hl{✔ সঠিক উত্তর: খ) AB\textgreater AP}}

\begin{quote}
\textbf{ব্যাখ্যা:} \emph{যেহেতু P বিন্দুটি AB রেখাংশের উপর অবস্থিত, তাই সম্পূর্ণ
রেখাংশের দৈর্ঘ্য (AB) যেকোনো অংশ AP অপেক্ষা সবসময় বড় হবে।}
\end{quote}

\textbf{AB ও CD সরলরেখাদ্বয় বিন্দুতে ছেদ করলে, নিম্নের কোন গাণিতিক বাক্যটি
সঠিক?}

\textbf{\hl{ক) ∠AOD=∠BOC ✔}}

খ) ∠AOD=∠BOD

গ) ∠BOC=∠AOC

ঘ) ∠AOD=∠AOC

\textbf{\hl{✔ সঠিক উত্তর: ক) ∠AOD=∠BOC}}

\begin{quote}
\textbf{ব্যাখ্যা:} \emph{দুটি সরলরেখা পরস্পর ছেদ করলে উৎপন্ন বিপরীত বা
বিপ্রতীপ কোণগুলো পরস্পর সমান হয়। এক্ষেত্রে ∠AOD এবং ∠BOC পরস্পর বিপ্রতীপ কোণ।}
\end{quote}

\textbf{180° \textless{} ∠A \textless{} 360° হলে ∠A কোন প্রকারের কোণ?}

ক) স্থূলকোণ

খ) সরল কোণ

গ) সুক্ষ্মকোণ

\textbf{\hl{ঘ) প্রবৃদ্ধ কোণ ✔}}

\textbf{\hl{✔ সঠিক উত্তর: ঘ) প্রবৃদ্ধ কোণ}}

\begin{quote}
\textbf{ব্যাখ্যা:} \emph{যে কোণের পরিমান দুই সমকোণ বা 180° থেকে বড় কিন্তু চার
সমকোণ বা 360° থেকে ছোট, তাকে প্রবৃদ্ধ কোণ বলা হয়।}
\end{quote}

\textbf{AB ও CD সরলরেখা বিন্দুতে ছেদ করলে নিম্নের কোন গাণিতিক বাক্যটি সঠিক?}

\textbf{\hl{ক) ∠AOD=∠BOC ✔}}

খ) ∠AOD=∠BOD

গ) ∠BOC=∠AOC

ঘ) ∠AOD=∠AOC

\textbf{\hl{✔ সঠিক উত্তর: ক) ∠AOD=∠BOC}}

\begin{quote}
\textbf{ব্যাখ্যা:} \emph{দুটি সরলরেখা পরস্পর ছেদ করলে উৎপন্ন বিপ্রতীপ কোণগুলো
পরস্পর সমান হয়, যার ফলে ∠AOD=∠BOC সঠিক।}
\end{quote}

\textbf{একটি কোণ তার পূরক কোণ অপেক্ষা ২৪ ডিগ্রি বেশি হলে কোণটির মান কত
হবে?}

\textbf{\hl{ক) ৫৭ ডিগ্রি ✔}}

খ) ৪৭ ডিগ্রি

গ) ৫৩ ডিগ্রি

ঘ) ৬৬ ডিগ্রি

\textbf{\hl{✔ সঠিক উত্তর: ক) ৫৭ ডিগ্রি}}

\begin{quote}
\textbf{ব্যাখ্যা:} \emph{ধরি, কোণটি = ক}

\emph{∴ তার পূরক কোণ = ক − ২৪}

\emph{প্রশ্নমতে, ক + (ক − ২৪) = ৯০}

\emph{⇒ ২ক = ১ِ১৪}

\emph{∴ ক = ৫৭ ডিগ্রি।}
\end{quote}

\textbf{বিন্দু কত মাত্রিক?}

\textbf{\hl{ক) শূন্য ✔}}

খ) এক

গ) দুই

ঘ) তিন

\textbf{\hl{✔ সঠিক উত্তর: ক) শূন্য}}

\begin{quote}
\textbf{ব্যাখ্যা:} \emph{বিন্দুর কোনো মাত্রা নেই অর্থাৎ শূন্য মাত্রিক। যার দৈর্ঘ্য,
প্রস্থ ও উচ্চতা কিছুই নেই শুধু অবস্থান আছে তাকে বিন্দু বলে।}
\end{quote}

\textbf{রেখার প্রান্ত বিন্দুর সংখ্যা হলো / The number of end points of a line
is}

ক) একটি

খ) দুটি

গ) তিনটি

\textbf{\hl{ঘ) কোন প্রান্ত বিন্দু নেই ✔}}

\textbf{\hl{✔ সঠিক উত্তর: ঘ) কোন প্রান্ত বিন্দু নেই}}

\begin{quote}
\textbf{ব্যাখ্যা:} \emph{রেখা উভয় দিকে অসীম পর্যন্ত বিস্তৃত হতে পারে, তাই এর
কোনো নির্দিষ্ট প্রান্তবিন্দু বা শেষ প্রান্ত নেই।}
\end{quote}

\textbf{কোন তিনটি মাত্রায় জ্যামিতিক ঘনবস্তু তৈরি হয়?}

ক) দৈর্ঘ্য, ভর ও সময়

খ) ভর, ওজন ও ঘনত্ব

গ) দৈর্ঘ্য, প্রস্থ ও ভর

\textbf{\hl{ঘ) দৈর্ঘ্য, প্রস্থ ও উচ্চতা ✔}}

\textbf{\hl{✔ সঠিক উত্তর: ঘ) দৈর্ঘ্য, প্রস্থ ও উচ্চতা}}

\begin{quote}
\textbf{ব্যাখ্যা:} \emph{জ্যামিতিক ঘনবস্তুর তিনটি মাত্রা থাকে, যথা--- দৈর্ঘ্য,
প্রস্থ ও উচ্চতা।}
\end{quote}

\textbf{একটি সরলরেখার উপর অঙ্কিত বর্গ ঐ সরলরেখার অর্ধেকের উপর অঙ্কিত বর্গের
কতগুণ?}

ক) দ্বিগুণ

\textbf{\hl{খ) চার গুণ ✔}}

গ) ছয় গুণ

ঘ) আট গুণ

\textbf{\hl{✔ সঠিক উত্তর: খ) চার গুণ}}

\begin{quote}
\textbf{ব্যাখ্যা:} \emph{ধরি, সরলরেখার দৈর্ঘ্য = x}

\emph{সরলরেখার উপর অঙ্কিত বর্গের ক্ষেত্রফল = x²}

\emph{অর্ধেকের দৈর্ঘ্য = (x/2)}

\emph{অর্ধেকের উপর অঙ্কিত বর্গের ক্ষেত্রফল = ( (x/2) )² = (x\^{}2/4)}

\emph{∴ অনুপাত = x² ÷ (x\^{}2/4) = চার গুণ}
\end{quote}

\textbf{২টি বিন্দু দিয়ে কয়টি সরলরেখা আঁকা যাবে?}

ক) ২টি

\textbf{\hl{খ) ১টি ✔}}

গ) ৩টি

ঘ) ৪টি

\textbf{\hl{✔ সঠিক উত্তর: খ) ১টি}}

\begin{quote}
\textbf{ব্যাখ্যা:} \emph{যেহেতু একটি সরলরেখায় শুধু দুটি বিন্দু থাকে তাই দুটি বিন্দু
দিয়ে কেবল একটি সরলরেখাই অঙ্কন সম্ভব।}
\end{quote}

\textbf{দুটি পূরক কোণের সমষ্টি কত?}

\textbf{\hl{ক) ৯০ ডিগ্রী ✔}}

খ) ১৮০ ডিগ্রী

গ) ১৬০ ডিগ্রী

ঘ) ১৪০ ডিগ্রী

\textbf{\hl{✔ সঠিক উত্তর: ক) ৯০ ডিগ্রী}}

\begin{quote}
\textbf{ব্যাখ্যা:} \emph{দুটি কোণের পরিমাপের যোগফল ৯০ ডিগ্রি হলে কোণ দুটিকে
একে অপরের পূরক কোণ বলা হয়।}
\end{quote}

\textbf{৩০ ডিগ্রী এর পূরক কোণ কত?}

ক) ৯০ ডিগ্রী

খ) ১৮০ ডিগ্রী

\textbf{\hl{গ) ৬০ ডিগ্রী ✔}}

ঘ) ১৩৫ ডিগ্রী

\textbf{\hl{✔ সঠিক উত্তর: গ) ৬০ ডিগ্রী}}

\begin{quote}
\textbf{ব্যাখ্যা:} \emph{আমরা জানি, দুটি পূরক কোণের সমষ্টি ৯০°}

\emph{∴ 30° কোণের পূরক কোণ = 90° − 30° = 60°}
\end{quote}

\textbf{দুইটি সরলরেখা পরস্পরকে ছেদ করলে সন্নিহিত কোণগুলির যোগফল-}

ক) ৬০ ডিগ্রী

\textbf{\hl{খ) ১৮০ ডিগ্রী ✔}}

গ) ১৬০ ডিগ্রী

ঘ) ১৪৫ ডিগ্রী

\textbf{\hl{✔ সঠিক উত্তর: খ) ১৮০ ডিগ্রী}}

\begin{quote}
\textbf{ব্যাখ্যা:} \emph{একটি সরলরেখার উপর অপর একটি রেখাংশ মিলিত হলে যে
দুটি সন্নিহিত কোণ উৎপন্ন হয়, তাদের সমষ্টি দুই সমকোণ বা 180°।}
\end{quote}

\textbf{দুইটি সম্পূরক কোণের সমষ্টি কত?}

ক) ৬০ ডিগ্রী

\textbf{\hl{খ) ১৮০ ডিগ্রী ✔}}

গ) ১৬০ ডিগ্রী

ঘ) ১৪৫ ডিগ্রী

\textbf{\hl{✔ সঠিক উত্তর: খ) ১৮০ ডিগ্রী}}

\begin{quote}
\textbf{ব্যাখ্যা:} \emph{দুটি কোণের পরিমাপের যোগফল 180° বা দুই সমকোণ হলে কোণ
দুটিকে একে অপরের সম্পূরক কোণ বলা হয়।}
\end{quote}

\textbf{৭০ ডিগ্রী কোণের সম্পূরক কোণ কোনটি?}

ক) ৯০ ডিগ্রী

\textbf{\hl{খ) ১১০ ডিগ্রী ✔}}

গ) ১৬০ ডিগ্রী

ঘ) ১৪৫ ডিগ্রী

\textbf{\hl{✔ সঠিক উত্তর: খ) ১১০ ডিগ্রী}}

\begin{quote}
\textbf{ব্যাখ্যা:} \emph{আমরা জানি, দুটি সম্পূরক কোণের সমষ্টি 180°}

\emph{∴ 70° কোণের সম্পূরক কোণ = 180° − 70° = 110°}
\end{quote}

\textbf{সন্নিহিত সরলকোণের একটি কোণ ৯৫ ডিগ্রি হলে অপরটি কত?}

ক) ৯০ ডিগ্রী

খ) ১১০ ডিগ্রী

গ) ১৬০ ডিগ্রী

\textbf{\hl{ঘ) ৮৫ ডিগ্রী ✔}}

\textbf{\hl{✔ সঠিক উত্তর: ঘ) ৮৫ ডিগ্রী}}

\begin{quote}
\textbf{ব্যাখ্যা:} \emph{সন্নিহিত সরলকোণদ্বয়ের সমষ্টি 180°}

\emph{∴ অপর কোণটি = 180° − 95° = 85°}
\end{quote}

\textbf{দুটি কোণের একটি সাধারণ বাহু থাকলে এবং কোণ দুটি সাধারণ বাহুর বিপরীত
দিকে অবস্থিত হলে কোণ দুটিকে বলে-}

\textbf{\hl{ক) সন্নিহিত কোণ ✔}}

খ) বিপ্রতীপ কোণ

গ) পূরক কোণ

ঘ) সম্পূরক কোণ

\textbf{\hl{✔ সঠিক উত্তর: ক) সন্নিহিত কোণ}}

\begin{quote}
\textbf{ব্যাখ্যা:} \emph{সন্নিহিত কোণের সংজ্ঞা অনুযায়ী, যদি দুটি কোণের একটি
সাধারণ শীর্ষবিন্দু এবং একটি সাধারণ বাহু থাকে এবং কোণদ্বয় সাধারণ বাহুর বিপরীত
পার্শ্বে অবস্থান করে, তবে তাদেরকে সন্নিহিত কোণ বলে।}
\end{quote}

\textbf{একটি ইটের কয়টি তল আছে?}

ক) ৩টি

খ) ৮টি

\textbf{\hl{গ) ৬টি ✔}}

ঘ) ১২টি

\textbf{\hl{✔ সঠিক উত্তর: গ) ৬টি}}

\begin{quote}
\textbf{ব্যাখ্যা:} \emph{একটি ইটের (বা আয়তাকার ঘনবস্তুর) মোট ৬টি তল থাকে।}
\end{quote}

\textbf{দুটি সরলরেখা পরস্পর ছেদ করলে কয়টি কোণ তৈরি হয়?}

ক) ১টি

খ) ৩টি

গ) ২টি

\textbf{\hl{ঘ) ৪টি ✔}}

ঙ) কোনোটিই নয়

\textbf{\hl{✔ সঠিক উত্তর: ঘ) ৪টি}}

\begin{quote}
\textbf{ব্যাখ্যা:} \emph{দুটি সরলরেখা পরস্পরকে ছেদ করলে তাদের ছেদবিন্দুতে মোট
৪টি কোণ তৈরি হয়।}
\end{quote}

\textbf{সরল কোণের মান কত?}

ক) ৯০°

\textbf{\hl{খ) ১৮০° ✔}}

গ) ২৭০°

ঘ) ৩৬০°

\textbf{\hl{✔ সঠিক উত্তর: খ) ১৮০°}}

\begin{quote}
\textbf{ব্যাখ্যা:} \emph{১ সমকোণ = ৯০°}

\emph{∴ সরল কোণ = ২ সমকোণ = ২ × ৯০° = ১৮০°}
\end{quote}

\textbf{একটি সরলরেখার উপরে লম্ব অঙ্কন করলে কয়টি সমকোণ পাওয়া যায়?}

\textbf{\hl{ক) ২টি ✔}}

খ) ৩টি

গ) ৪টি

ঘ) ৫টি

\textbf{\hl{✔ সঠিক উত্তর: ক) ২টি}}

\begin{quote}
\textbf{ব্যাখ্যা:} \emph{একটি সরলরেখার উপরে লম্ব অঙ্কন করলে দুই পাশে মোট ২টি
সমকোণ পাওয়া যায়।}
\end{quote}

\textbf{যে কোণের ডিগ্রির পরিমাপ ৯০° তাকে কী বলে?}

\textbf{\hl{ক) সমকোণ ✔}}

খ) স্থূলকোণ

গ) সরলকোণ

ঘ) সূক্ষ্মকোণ

\textbf{\hl{✔ সঠিক উত্তর: ক) সমকোণ}}

\begin{quote}
\textbf{ব্যাখ্যা:} \emph{৯০° এর সমান কোণকে সমকোণ বলে।}
\end{quote}

\textbf{একটি সরলরেখার সাথে আর একটি রেখাংশ মিলিত হয়ে যে দুটি সন্নিহিত কোণ
উৎপন্ন হয় তাদের সমষ্টি কত হবে?}

ক) ১৬০°

\textbf{\hl{খ) ১৮০° ✔}}

গ) ১২০°

ঘ) ৯০°

\textbf{\hl{✔ সঠিক উত্তর: খ) ১৮০°}}

\begin{quote}
\textbf{ব্যাখ্যা:} \emph{একটি সরলরেখার সাথে আরেকটি রেখাংশ মিলিত হয়ে যে দুটি
সন্নিহিত কোণ উৎপন্ন করে, তাদের সমষ্টি সর্বদা ১৮০° (বা দুই সমকোণ) হয়।}
\end{quote}

\textbf{একটি সমকোণী ত্রিভুজের সমকোণ ছাড়া অন্য দুটি কোণ কি কোণ?}

ক) সরলকোণ

খ) সন্নিহিত কোণ

\textbf{\hl{গ) সূক্ষ্মকোণ ✔}}

ঘ) স্থূলকোণ

\textbf{\hl{✔ সঠিক উত্তর: গ) সূক্ষ্মকোণ}}

\begin{quote}
\textbf{ব্যাখ্যা:} \emph{ত্রিভুজের সমকোণ হলো ৯০°}

\emph{সমকোণী ত্রিভুজের সমকোণ ব্যতীত অন্য কোণ দুটির প্রত্যেকে সমকোণ অপেক্ষা ছোট}

\emph{আর সমকোণ অপেক্ষা ছোট কোণকে সূক্ষ্মকোণ বলে।}
\end{quote}

\textbf{সমকোণী ত্রিভুজের সংলগ্ন দুটির প্রত্যেকেটি---}

ক) স্থূলকোণ

খ) সরলকোণ

\textbf{\hl{গ) সূক্ষ্মকোণ ✔}}

ঘ) পূরক কোণ

\textbf{\hl{✔ সঠিক উত্তর: গ) সূক্ষ্মকোণ}}

\begin{quote}
\textbf{ব্যাখ্যা:} \emph{৯০° এর চেয়ে ছোট কোণকে সূক্ষ্মকোণ বলা হয়}

\emph{সমকোণী ত্রিভুজের সমকোণ ব্যতীত অন্য দুই কোণের সমষ্টি ৯০°}

\emph{∴ অতিভুজ সংলগ্ন কোণ দুটি প্রত্যেকেটি সূক্ষ্মকোণ।}
\end{quote}

\textbf{৮০° কোণটি হলো ------}

ক) সরলকোণ

খ) স্থূলকোণ

\textbf{\hl{গ) সূক্ষ্মকোণ ✔}}

ঘ) সমকোণ

\textbf{\hl{✔ সঠিক উত্তর: গ) সূক্ষ্মকোণ}}

\begin{quote}
\textbf{ব্যাখ্যা:} \emph{৮০° কোণটি এক সমকোণ বা ৯০° অপেক্ষা ছোট}

\emph{তাই এটি সূক্ষ্মকোণ।}
\end{quote}

\textbf{একটি ত্রিভুজের দুই কোণের সমষ্টি ৭৮ হলে অপর কোণটি কত?}

ক) ৯০°

\textbf{\hl{খ) ১০২° ✔}}

গ) ১১০°

ঘ) ১০৮°

\textbf{\hl{✔ সঠিক উত্তর: খ) ১০২°}}

\begin{quote}
\textbf{ব্যাখ্যা:} \emph{ত্রিভুজের তিন কোণের সমষ্টি ১৮০°}

\emph{∴ অপর কোণ = (১৮০ − ৭৮)° = ১০২°}
\end{quote}

\textbf{180° কোণ একটি---}

ক) সূক্ষ্মকোণ

খ) স্থূলকোণ

\textbf{\hl{গ) সম্পূরক কোণ ✔}}

ঘ) প্রবৃদ্ধ কোণ

\textbf{\hl{✔ সঠিক উত্তর: গ) সম্পূরক কোণ}}

\begin{quote}
\textbf{ব্যাখ্যা:} \emph{যে কোণের পরিমাপ ১৮০° বা দুই সমকোণ তাকে সরল কোণ
বলে। প্রদত্ত অপশনগুলোর মধ্যে ১৮০° কোণটি সম্পূরক কোণ বা দুই সমকোণের ক্ষেত্রে
প্রযোজ্য।}
\end{quote}

\textbf{৪০ ডিগ্রী কোণের পূরক কোণ কোনটি?}

ক) ৩২০ ডিগ্রি

\textbf{\hl{খ) ৫০ ডিগ্রি ✔}}

গ) ১২০ ডিগ্রি

ঘ) ১৪০ ডিগ্রি

\textbf{\hl{✔ সঠিক উত্তর: খ) ৫০ ডিগ্রি}}

\begin{quote}
\textbf{ব্যাখ্যা:} \emph{দুটি কোণের সমষ্টি এক সমকোণ (৯০°) হলে তবে তাদের
একটিকে অপরটির পূরক কোণ বলে。}

\emph{∴ ৪০° কোণের পূরক কোণ = (৯০ − ৪০)° = ৫০°}
\end{quote}

\textbf{দুটি সন্নিহিত কোণের সমষ্টি দুই সমকোণ হলে একটিকে অপরটির কি বলে?}

ক) সমদ্বিখণ্ডিত কোণ

খ) পূরক কোণ

\textbf{\hl{গ) সম্পূরক কোণ ✔}}

ঘ) সমকোণ

\textbf{\hl{✔ সঠিক উত্তর: গ) সম্পূরক কোণ}}

\begin{quote}
\textbf{ব্যাখ্যা:} \emph{দুটি সন্নিহিত কোণের সমষ্টি ১৮০° বা দুই সমকোণ হলে
তাদের একটিকে অপরটির সম্পূরক কোণ বলে।}
\end{quote}

\textbf{৭৮ ডিগ্রি কোণের সম্পূরক কোণ কত ডিগ্রি?}

ক) ১২

খ) ৭৮

গ) ২৮২

\textbf{\hl{ঘ) ১০২ ✔}}

\textbf{\hl{✔ সঠিক উত্তর: ঘ) ১০২}}

\begin{quote}
\textbf{ব্যাখ্যা:} \emph{দুটি কোণের সমষ্টি ১৮০° হলে একটি অপরটির সম্পূরক কোণ
হয়।}

\emph{∴ ৭৮° কোণের সম্পূরক কোণ = (১৮০ − ৭৮)° = ১০২°}
\end{quote}

\textbf{৯০° কোণের সম্পূরক কোণ কত ডিগ্রি?}

ক) ০°

\textbf{\hl{খ) ৯০° ✔}}

গ) ১৮০°

ঘ) ২৭০°

\textbf{\hl{✔ সঠিক উত্তর: খ) ৯০°}}

\begin{quote}
\textbf{ব্যাখ্যা:} \emph{দুটি কোণের সমষ্টি ১৮০° হলে একটিকে অপরটির সম্পূরক কোণ
বলা হয়।}

\emph{∴ ৯০° কোণের সম্পূরক কোণ = ১৮০° − ৯০° = ৯০°}
\end{quote}

\textbf{30° কোণের সম্পূরক কোণের অর্ধেক নিচের কোনটি?}

\textbf{\hl{ক) 75° ✔}}

খ) 36°

গ) 60°

ঘ) 80°

\textbf{\hl{✔ সঠিক উত্তর: ক) 75°}}

\begin{quote}
\textbf{ব্যাখ্যা:} \emph{30° কোণের সম্পূরক কোণ = 180° − 30° = 150°}

\emph{∴ এর অর্ধেক হলো = (150\^{}°/2) = 75°}
\end{quote}

\textbf{180° − x° কোণের সম্পূরক মান কত?}

ক) x + 90°

\textbf{\hl{খ) x° ✔}}

গ) 0°

ঘ) 45°

\textbf{\hl{✔ সঠিক উত্তর: খ) x°}}

\begin{quote}
\textbf{ব্যাখ্যা:} \emph{আমরা জানি, দুটি কোণের সমষ্টি 180° হলে তাদের একটিকে
অপরটির সম্পূরক কোণ বলে।}

\emph{ধরি, (180° − x°) এর সম্পূরক কোণ y.}

\emph{প্রশ্নমতে, 180° − x° + y = 180° ⇒ y = 180° − 180° + x° = x°}

\emph{∴ (180° − x°) এর সম্পূরক কোণ হবে x°।}
\end{quote}

\textbf{দুই সমকোণ অপেক্ষা বড় ও চার সমকোণ অপেক্ষা ছোট কোণকে বলে-\/- অথবা,
১৮০° এর চেয়ে বড় কিন্তু ৩৬০° এর চেয়ে ছোট কোনকে কি কোণ বলে?}

ক) বিপ্রতীপ কোণ

\textbf{\hl{খ) প্রবৃদ্ধ কোণ ✔}}

গ) স্থূল কোণ

ঘ) সম্পূরক কোণ

\textbf{\hl{✔ সঠিক উত্তর: খ) প্রবৃদ্ধ কোণ}}

\begin{quote}
\textbf{ব্যাখ্যা:} \emph{দুই সমকোণ অপেক্ষা বড় কিন্তু চার সমকোণ অপেক্ষা ছোট
কোনকে প্রবৃদ্ধ কোণ বলা হয়।}
\end{quote}

\textbf{একটি তলের দৈর্ঘ্য ৫ সে.মি. ও প্রস্থ শূন্য হলে তলটি?}

\textbf{\hl{ক) রেখাংশ ✔}}

খ) রেখা

গ) সমতল

ঘ) বক্রতল

\textbf{\hl{✔ সঠিক উত্তর: ক) রেখাংশ}}

\begin{quote}
\textbf{ব্যাখ্যা:} \emph{যেহেতু তলের দৈর্ঘ্য আছে কিন্তু প্রস্থ শূন্য, তাই এটি একটি
নির্দিষ্ট দৈর্ঘ্যের রেখাংশ হবে।}
\end{quote}

\textbf{২৫০° কোণকে কি কোণ বলে?}

ক) সূক্ষ্মকোণ

\textbf{\hl{খ) প্রবৃদ্ধ কোণ ✔}}

গ) স্থূল কোণ

ঘ) পূরক কোণ

\textbf{\hl{✔ সঠিক উত্তর: খ) প্রবৃদ্ধ কোণ}}

\begin{quote}
\textbf{ব্যাখ্যা:} \emph{দুই সমকোণ অপেক্ষা বড় কিন্তু চার সমকোণ অপেক্ষা ছোট
কোনকে প্রবৃদ্ধি কোণ বলে। অর্থাৎ ১৮০° \textless{} x \textless{} ৩৬০°}

\emph{এখানে, x = ২৫০° বলে তা প্রবৃদ্ধি কোণ।}
\end{quote}

\textbf{দুটি কোণের একটি শীর্ষবিন্দু থাকলে এবং ঐ কোণ দুটি যদি সাধারণ বাহুর
বিপরীত দিকে অবস্থান করে, তবে ঐ কোণ দুটিকে বলা হবে---}

ক) সূক্ষ্মকোণ

\textbf{\hl{খ) সন্নিহিত কোণ ✔}}

গ) পূরক কোণ

ঘ) বিপ্রতীপ কোণ

\textbf{\hl{✔ সঠিক উত্তর: খ) সন্নিহিত কোণ}}

\begin{quote}
\textbf{ব্যাখ্যা:} \emph{সন্নিহিত কোণের সংজ্ঞা অনুযায়ী, দুটি কোণের একটি সাধারণ
শীর্ষবিন্দু এবং একটি সাধারণ বাহু থাকে এবং কোণদ্বয় সাধারণ বাহুর বিপরীত পার্শ্বে
অবস্থান করে।}
\end{quote}

\textbf{দুইটি সরলরেখা পরস্পর সমাপতিত হলে, সমাধান সংখ্যা কত হবে?}

\textbf{\hl{ক) অসংখ্য ✔}}

খ) সমাধান নেই

গ) দুইটি

ঘ) একটি

\textbf{\hl{✔ সঠিক উত্তর: ক) অসংখ্য}}

\begin{quote}
\textbf{ব্যাখ্যা:} \emph{দুটি সরলরেখা পরস্পর সমাপতিত হলে তারা একে অপরের উপর
পুরোপুরি মিলে যায়, ফলে তাদের অসংখ্য সাধারণ বিন্দু বা সমাধান থাকে।}
\end{quote}

\textbf{সরল রেখার উপর লম্ব অঙ্কন করলে কয়টি সমকোণ পাওয়া যায়?}

\textbf{\hl{ক) ২টি ✔}}

খ) ৩টি

গ) ৪টি

ঘ) ৫টি

\textbf{\hl{✔ সঠিক উত্তর: ক) ২টি}}

\begin{quote}
\textbf{ব্যাখ্যা:} \emph{সরল রেখার উপর লম্ব অঙ্কন করলে দুই পাশে মোট ২টি সমকোণ
পাওয়া যায়।}
\end{quote}

\textbf{এক সমকোণ অপেক্ষা ছোট কোণকে কি বলে?}

ক) স্থূল কোণ

\textbf{\hl{খ) সূক্ষ্ম কোণ ✔}}

গ) সন্নিহিত কোণ

ঘ) সমকোণ

\textbf{\hl{✔ সঠিক উত্তর: খ) সূক্ষ্ম কোণ}}

\begin{quote}
\textbf{ব্যাখ্যা:} \emph{এক সমকোণ (৯০°) অপেক্ষা ছোট কোণকে সূক্ষ্মকোণ বলে।}
\end{quote}

\textbf{দুটি পরস্পর পূরক কোণের মধ্যে একটির মান ৩৫ ডিগ্রি হলে অপরটির মান কত?}

\textbf{\hl{ক) ৫৫ ডিগ্রি ✔}}

খ) ৪৫ ডিগ্রি

গ) ৬৫ ডিগ্রি

ঘ) ৩৫ ডিগ্রি

\textbf{\hl{✔ সঠিক উত্তর: ক) ৫৫ ডিগ্রি}}

\begin{quote}
\textbf{ব্যাখ্যা:} \emph{দুটি কোণের সমষ্টি ৯০° হলে একটি অপরটির পূরক কোণ হয়।}

\emph{∴ অপর কোণ = ৯০° − ৩৫° = ৫৫°}
\end{quote}

\textbf{দুটি রশ্মি দ্বারা উৎপন্ন কোণ 60°। এক সরলকোণ হতে উক্ত কোণ বিয়োগ করলে
কি কোণ উৎপন্ন হবে?}

ক) সমকোণ

খ) সূক্ষ্মকোণ

\textbf{\hl{গ) স্থূলকোণ ✔}}

ঘ) প্রবৃদ্ধ কোণ

\textbf{\hl{✔ সঠিক উত্তর: গ) স্থূলকোণ}}

\begin{quote}
\textbf{ব্যাখ্যা:} \emph{এক সরলকোণ = 180°}

\emph{180° − 60° = 120°}

\emph{যেহেতু 120° কোণটি এক সমকোণ অপেক্ষা বড় কিন্তু দুই সমকোণ অপেক্ষা ছোট, তাই
এটি একটি স্থূলকোণ।}
\end{quote}

\textbf{দুটি সন্নিহিত কোণের সমষ্টি দুই সমকোণ হলে, একটিকে অপরটির কি বলে?}

ক) সন্নিহিত কোণ

খ) পূরক কোণ

গ) বিপ্রতীপ কোণ

\textbf{\hl{ঘ) সম্পূরক কোণ ✔}}

\textbf{\hl{✔ সঠিক উত্তর: ঘ) সম্পূরক কোণ}}

\begin{quote}
\textbf{ব্যাখ্যা:} \emph{যদি দুই কোণের সমষ্টি দুই সমকোণ বা ১৮০° হয় তবে
একেকটিকে অপরটির সম্পূরক কোণ বলে।}
\end{quote}

\textbf{একটি সরলরেখার সাথে আর একটি রেখাংশ মিলিত হয়ে যে দুটি সন্নিহিত কোণ
তৈরি হয় তাদের সমষ্টি কত?}

\textbf{\hl{ক) ১৮০° ✔}}

খ) ১২০°

গ) ৩৬০°

ঘ) ৯০°

\textbf{\hl{✔ সঠিক উত্তর: ক) ১৮০°}}

\begin{quote}
\textbf{ব্যাখ্যা:} \emph{একটি সরলরেখার সাথে আরেকটি রেখাংশ মিলিত হয়ে যে দুটি
সন্নিহিত কোণ উৎপন্ন করে তাদের সমষ্টি সর্বদা ১৮০° (বা দুই সমকোণ) হয়।}
\end{quote}

\textbf{একটি সরলরেখার সাথে অপর একটি রেখাংশ মিলিত হলে যে দুটি সন্নিহিত কোণ
উৎপন্ন হয়, তাদের সমষ্টি হবে---}

ক) ১২০°

খ) ১৪০°

গ) ১৬০°

\textbf{\hl{ঘ) ১৮০° ✔}}

\textbf{\hl{✔ সঠিক উত্তর: ঘ) ১৮০°}}

\begin{quote}
\textbf{ব্যাখ্যা:} \emph{একটি সরলরেখার সাথে আরেকটি রেখাংশ মিলিত হলে যে দুটি
সন্নিহিত কোণ উৎপন্ন হয়, তাদের সমষ্টি সর্বদা ১৮০° (বা দুই সমকোণ) হয়।}
\end{quote}

\textbf{দুটি সরলরেখা পরস্পর ছেদ করলে যে চারটি কোণ উৎপন্ন হয় তাদের একটির
বিপরীত কোণকে অপরটির কি বলা হয়?}

ক) পূরক কোণ

খ) সম্পূরক কোণ

\textbf{\hl{গ) বিপ্রতীপ কোণ ✔}}

ঘ) সন্নিহিত কোণ

\textbf{\hl{✔ সঠিক উত্তর: গ) বিপ্রতীপ কোণ}}

\begin{quote}
\textbf{ব্যাখ্যা:} \emph{দুটি সরলরেখা পরস্পর ছেদ করলে ছেদবিন্দুতে যে চারটি কোণ
উৎপন্ন হয়, তাদের একটির বিপরীত কোণকে অপরটির বিপ্রতীপ কোণ বলা হয়।}
\end{quote}

\textbf{দুটি কোণের যোগফল এক সমকোণের সমান হলে, একটি কোণ অপরটির}

\textbf{\hl{ক) পূরক ✔}}

খ) সম্পূরক

গ) সন্নিহিত

ঘ) বিপ্রতীপ

\textbf{\hl{✔ সঠিক উত্তর: ক) পূরক}}

\begin{quote}
\textbf{ব্যাখ্যা:} \emph{দুটি কোণের সমষ্টি ৯০° (বা এক সমকোণ) হলে তাদের
একটিকে অপরটির পূরক কোণ বলা হয়।}
\end{quote}

\textbf{৫৫° কোণের পূরক কোণ কত?}

\textbf{\hl{ক) ৩৫° ✔}}

খ) ৪৫°

গ) ১২৫°

ঘ) ১৮০°

\textbf{\hl{✔ সঠিক উত্তর: ক) ৩৫°}}

\begin{quote}
\textbf{ব্যাখ্যা:} \emph{দুটি কোণের সমষ্টি ৯০° হলে কোণ দুটিকে পরস্পরের পূরক কোণ
বলে。}

\emph{∴ ৫৫° কোণের পূরক কোণ = ৯০° − ৫৫° = ৩৫°}
\end{quote}

\textbf{কোন ত্রিভুজের ∠A = 65° হলে, এর পূরক কোণ কত ডিগ্রি?}

\textbf{\hl{ক) 25° ✔}}

খ) 35°

গ) 90°

ঘ) 115°

\textbf{\hl{✔ সঠিক উত্তর: ক) 25°}}

\begin{quote}
\textbf{ব্যাখ্যা:} \emph{পূরক কোণ = 90° − 65° = 25°}
\end{quote}

\textbf{একটি কোণের দ্বিগুণ ৬০° হলে, তার পূরক কোণ কত?}

\textbf{\hl{ক) ৩০° ✔}}

খ) ১৫°

গ) ২০°

ঘ) ৬০°

\textbf{\hl{✔ সঠিক উত্তর: ক) ৩০°}}

\begin{quote}
\textbf{ব্যাখ্যা:} \emph{কোণটির দ্বিগুণ ৬০° হলে কোণটি = (60\^{}°/2) = ৩০°}

\emph{∴ ৩০° কোণের পূরক কোণ = ৯০° − ৩০° = ৬০°}
\end{quote}

\textbf{ত্রিভুজ ABC এর AB=AC, BC বাহুকে D পর্যন্ত বর্ধিত করলে যদি ∠ACD = 120°
হয় তবে ∠BAC এর মান কত?}

ক) 30°

খ) 45°

\textbf{\hl{গ) 60° ✔}}

ঘ) 90°

\textbf{\hl{✔ সঠিক উত্তর: গ) 60°}}

\begin{quote}
\textbf{ব্যাখ্যা:} \emph{এখানে, AB = AC বলে ∠B = ∠C}

\emph{বহিঃস্থ কোণ ∠ACD = ∠BAC + ∠B ⇒ 120° = ∠BAC + ∠C}

\emph{আবার ত্রিভুজের তিন কোণের সমষ্টি 180° ব্যবহার করে হিসাব করলে, ∠BAC =
60° পাওয়া যায়।}
\end{quote}

\textbf{কোন ত্রিভুজের তিনটি বাহু বর্ধিত করলে উৎপন্ন বহিঃস্থ তিনটি কোণের পরিমান
কত?}

ক) ১৮০°

খ) ২৭০°

গ) ১৫০°

\textbf{\hl{ঘ) ৩৬০° ✔}}

\textbf{\hl{✔ সঠিক উত্তর: ঘ) ৩৬০°}}

\begin{quote}
\textbf{ব্যাখ্যা:} \emph{আমরা জানি, যে কোনো ত্রিভুজের তিন কোণের সমষ্টি =
180°}

\emph{তিনটি কোণ যথাক্রমে x+y+z = 180°}

\emph{∴ বহিঃস্থ কোণ তিনটির যোগফল = (180° − x) + (180° − y) + (180° − z) =
540° − (x+y+z) = 540° − 180° = 360°}
\end{quote}

\textbf{৯০° কোণের পূরক কোণের মান কত?}

\textbf{\hl{ক) ০° ✔}}

খ) ৯০°

গ) ১৮০°

ঘ) ৩৬০°

\textbf{\hl{✔ সঠিক উত্তর: ক) ০°}}

\begin{quote}
\textbf{ব্যাখ্যা:} \emph{আমরা জানি, দুটি সন্নিহিত কোণের সমষ্টি এক সমকোণ হলে,
এদের একটিকে অপরটির পূরক কোণ বলে。}

\emph{∴ ৯০° কোণের পূরক কোণ = ৯০° − ৯০° = ০°}
\end{quote}

\textbf{একটি ৬০° কোণের পূরক কোণ কত?}

\textbf{\hl{ক) ৩০° ✔}}

খ) ৯০°

গ) ৬০°

ঘ) ১২০°

\textbf{\hl{✔ সঠিক উত্তর: ক) ৩০°}}

\begin{quote}
\textbf{ব্যাখ্যা:} \emph{দুটি কোণের সমষ্টি ৯০° হলে একটিকে অপরটির পূরক কোণ
বলে。}

\emph{∴ ৬০° কোণের পূরক কোণ = ৯০° − ৬০° = ৩০°}
\end{quote}

\textbf{৫৫° কোণের পূরক কোণের পরিমাণ কত?}

\textbf{\hl{ক) ৩৫° ✔}}

খ) ১৩৫°

গ) ৩০৫°

ঘ) ১২৫°

\textbf{\hl{✔ সঠিক উত্তর: ক) ৩৫°}}

\begin{quote}
\textbf{ব্যাখ্যা:} \emph{৫৫° কোণের সম্পূরক কোণ = (১৮০° − ৫৫°) = ১২৫°; পূরক
কোণ = (৯০° − ৫৫°) = ৩৫°}
\end{quote}

\textbf{৩৭ ডিগ্রি কোণের বিপ্রতীপ কোণের পরিমাণ কত?}

\textbf{\hl{ক) ৩৭ ডিগ্রি ✔}}

খ) ৫৩ ডিগ্রি

গ) ১২৭ ডিগ্রি

ঘ) ১৪৩ ডিগ্রি

\textbf{\hl{✔ সঠিক উত্তর: ক) ৩৭ ডিগ্রি}}

\begin{quote}
\textbf{ব্যাখ্যা:} \emph{বিপ্রতীপ কোণগুলো পরস্পর সমান হয়। তাই ৩৭° কোণের
বিপ্রতীপ কোণের পরিমাণ হবে ৩৭°।}
\end{quote}

\textbf{দুটি কোণের সমষ্টি এক সমকোণের সমান হলে তাদের একটিকে অপরটির কি কোণ
বলে?}

ক) সম্পূরক কোণ

\textbf{\hl{খ) পূরক কোণ ✔}}

গ) সন্নিহিত কোণ

ঘ) প্রবৃদ্ধ কোণ

\textbf{\hl{✔ সঠিক উত্তর: খ) পূরক কোণ}}

\begin{quote}
\textbf{ব্যাখ্যা:} \emph{দুটি কোণের সমষ্টি এক সমকোণ বা ৯০° হলে তাদের একটিকে
অপরটির পূরক কোণ বলে।}
\end{quote}

\textbf{ΔABC এর ∠A = 45°, ∠B = 30° হলে ∠C এর মান কত?}

ক) 90°

খ) 100°

\textbf{\hl{গ) 105° ✔}}

ঘ) 110°

\textbf{\hl{✔ সঠিক উত্তর: গ) 105°}}

\begin{quote}
\textbf{ব্যাখ্যা:} \emph{ত্রিভুজের তিন কোণের সমষ্টি ১৮০°}

\emph{∠C = ১৮০° − (৪৫° + ৩০°) = ১৮০° − ৭৫° = ১০৫°}
\end{quote}

\textbf{∠A এবং ∠B পরস্পর সম্পূরক কোণ। ∠A = ১১৫° হলে ∠B = কত?}

\textbf{\hl{ক) ৬৫° ✔}}

খ) ৭৫°

গ) ৮৫°

ঘ) ৯০°

\textbf{\hl{✔ সঠিক উত্তর: ক) ৬৫°}}

\begin{quote}
\textbf{ব্যাখ্যা:} \emph{দুটির কোণের সমষ্টি দুই সমকোণ হলে, তাদের সম্পূরক কোণ
বলে。}

\emph{∴ ∠A + ∠B = ১৮০°}

\emph{১১৫° + ∠B = ১৮০°}

\emph{∠B = ১৮০° − ১১৫° = ৬৫°}
\end{quote}

\textbf{কোনটি ৩৫° কোণের পূরক কোণ?}

ক) ১২৫°

\textbf{\hl{খ) ৫৫° ✔}}

গ) ২৫°

ঘ) ৩২৫°

\textbf{\hl{✔ সঠিক উত্তর: খ) ৫৫°}}

\begin{quote}
\textbf{ব্যাখ্যা:} \emph{দুটি কোণের সমষ্টি এক সমকোণ (৯০°) হলে তাদের একটিকে
অপরটির পূরক কোণ বলে。}

\emph{∴ পূরক কোণ = (৯০ − ৩৫)° = ৫৫°}
\end{quote}

\textbf{একটি কোণের মান তার পূরক কোণের মানের অর্ধেকের সমান। কোণটির মান কত?}

ক) 60°

খ) 45°

\textbf{\hl{গ) 30° ✔}}

ঘ) 25°

\textbf{\hl{✔ সঠিক উত্তর: গ) 30°}}

\begin{quote}
\textbf{ব্যাখ্যা:} \emph{ধরি, একটি কোণ x}

\emph{কোণটির পূরক কোণ = (90° − x)}

\emph{প্রশ্নমতে, x = (90° − x) × (1/2)}

\emph{বা, 2x = 90° − x ⇒ 3x = 90°}

\emph{∴ x = (90\^{}°/3) = 30°}
\end{quote}

\textbf{২৮ ডিগ্রি কোণের সম্পূরক কোণের পরিমাণ কত?}

ক) ৬২ ডিগ্রি

খ) ১১৮ ডিগ্রি

\textbf{\hl{গ) ১৫২ ডিগ্রি ✔}}

ঘ) ৩৩২ ডিগ্রি

\textbf{\hl{✔ সঠিক উত্তর: গ) ১৫২ ডিগ্রি}}

\begin{quote}
\textbf{ব্যাখ্যা:} \emph{আমরা জানি, সম্পূরক কোণদ্বয়ের সমষ্টি = ১৮০°}

\emph{∴ ২৮° কোণের সম্পূরক কোণের পরিমাণ = (১৮০ − ২৮)° = ১৫২°}
\end{quote}

\textbf{45° কোণের সম্পূরক কোণ কোনটি?}

ক) 45°

খ) 90°

\textbf{\hl{গ) 135° ✔}}

ঘ) 180°

\textbf{\hl{✔ সঠিক উত্তর: গ) 135°}}

\begin{quote}
\textbf{ব্যাখ্যা:} \emph{দুটি কোণের সমষ্টি 180° হলে একেকটিকে অপরটির সম্পূরক
কোণ বলা হয়。}

\emph{∴ 45° কোণের সম্পূরক কোণ = (180 − 45)° = 135°}
\end{quote}

\textbf{∠A এবং ∠B পরস্পর সম্পূরক কোণ। ∠A = 115 হলে, ∠B কত?}

ক) ৯০°

খ) ৮৫°

গ) ৭৫°

\textbf{\hl{ঘ) ৬৫° ✔}}

\textbf{\hl{✔ সঠিক উত্তর: ঘ) ৬৫°}}

\begin{quote}
\textbf{ব্যাখ্যা:} \emph{যেহেতু ∠A ও ∠B পরস্পর সম্পূরক কোণ, তাই তাদের সমষ্টি
১৮০°}

\emph{∠B = ১৮০° − ১১৫° = ৬৫°}
\end{quote}

\textbf{২৮° কোণের সম্পূরক কোণের অর্ধেক কত?}

ক) ৬০°

খ) ৪০°

\textbf{\hl{গ) ৭৬° ✔}}

ঘ) ৩১°

\textbf{\hl{✔ সঠিক উত্তর: গ) ৭৬°}}

\begin{quote}
\textbf{ব্যাখ্যা:} \emph{২৮° কোণের সম্পূরক কোণ = (১৮০ − ২৮)° = ১৫২°}

\emph{∴ এর অর্ধেক = (152\^{}°/2) = ৭৬°}
\end{quote}

\textbf{দুটি সমান্তরাল রেখা কটি বিন্দুতে ছেদ করে?}

ক) ৪

খ) ২

গ) ৮

ঘ) ১৬

\textbf{\hl{ঙ) কখনো ছেদ করে না ✔}}

\textbf{\hl{✔ সঠিক উত্তর: ঙ) কখনো ছেদ করে না}}

\begin{quote}
\textbf{ব্যাখ্যা:} \emph{দুটি সরলরেখার মধ্যবর্তী দূরত্ব যখন সর্বদা একই থাকে তখন
একেকটিকে অপরটির সমান্তরাল রেখা বলে।}

\emph{∴ দুটি সমান্তরাল রেখা বাস্তবে কখনও মিলিত হতে পারে না বা একে অপরকে ছেদ
করে না।}
\end{quote}

\textbf{দুটি লাইন একে অন্যের থেকে ২ মিটার দূরে সমান্তরাল ভাবে চলে যাচ্ছে।
তারা একে অন্যের সাথে মিলিত হবে কত মিটার দূরে?}

ক) ২০০

খ) ৪০০

গ) ৬০০

\textbf{\hl{ঘ) কখনোই নয় ✔}}

\textbf{\hl{✔ সঠিক উত্তর: ঘ) কখনোই নয়}}

\begin{quote}
\textbf{ব্যাখ্যা:} \emph{যেহেতু লাইন দুটি সমান্তরালভাবে চলে যাচ্ছে, তাই তারা
কখনোই একে অপরের সাথে মিলিত হবে না।}
\end{quote}

\textbf{দুটি কোণ পরস্পর সমান এবং এদের একটির বাহু অপরটির এক বাহুর সমান্তরাল।
কোণ দুটির অপর বাহুদ্বয়ের মধ্যে সম্পর্ক কিরুপ?}

ক) এরা পরস্পর সমান (They are equal)

\textbf{\hl{খ) এরা পরস্পর সমান্তরাল (They are parallel) ✔}}

গ) এরা পরস্পরের উপর লম্ব (They are perpendicular to each other)

ঘ) এরা পরস্পর ছেদক (They intersect one another)

\textbf{\hl{✔ সঠিক উত্তর: খ) এরা পরস্পর সমান্তরাল (They are parallel)}}

\begin{quote}
\textbf{ব্যাখ্যা:} \emph{দুটি কোণ পরস্পর সমান হলে এবং তাদের এক জোড়া বাহু
সমান্তরাল হলে, অপর জোড়া বাহুদ্বয়ও সর্বদা পরস্পর সমান্তরাল হবে।}
\end{quote}

\textbf{২৫৩° কোণকে কি কোণ বলে?}

ক) সূক্ষ্মকোণ

খ) স্থূলকোণ

গ) পূরক কোণ

\textbf{\hl{ঘ) প্রবৃদ্ধ কোণ ✔}}

\textbf{\hl{✔ সঠিক উত্তর: ঘ) প্রবৃদ্ধ কোণ}}

\begin{quote}
\textbf{ব্যাখ্যা:} \emph{যে কোণ দুই সমকোণ বা ১৮০° অপেক্ষা বড় কিন্তু চার সমকোণ
বা ৩৬০° অপেক্ষা ছোট, তাকে প্রবৃদ্ধ কোণ বলে। যেহেতু ২৫৩° কোণটি ১৮০° অপেক্ষা বড়
এবং ৩৬০° অপেক্ষা ছোট, তাই এটি একটি প্রবৃদ্ধ কোণ।}
\end{quote}

\textbf{দু\textquotesingle টি কোণের একই শীর্ষবিন্দু থাকলে এবং ঐ কোণ
দু\textquotesingle টি যদি সাধারণ বাহুর বিপরীত দিকে অবস্থান করে, তবে ঐ কোণ
দু\textquotesingle টিকে বলা হবে-}

ক) সূক্ষ্মকোণ

খ) পূরক কোণ

\textbf{\hl{গ) সন্নিহিত কোণ ✔}}

ঘ) বিপ্রতীপ কোণ

\textbf{\hl{✔ সঠিক উত্তর: গ) সন্নিহিত কোণ}}

\begin{quote}
\textbf{ব্যাখ্যা:} \emph{দুটি কোণের একটি সাধারণ শীর্ষবিন্দু এবং একটি সাধারণ
বাহু থাকলে এবং কোণ দুটি সাধারণ বাহুর বিপরীত দিকে অবস্থিত হলে তাদেরকে সন্নিহিত
কোণ বলে।}
\end{quote}

\textbf{একটি সরলরেখার উপর অঙ্কিত বর্গ ঐ সরলরেখার এক চতুর্থাংশের উপর অঙ্কিত
বর্গের কতগুণ?}

ক) ৪০

খ) ৮

গ) ১২

\textbf{\hl{ঘ) ১৬ ✔}}

\textbf{\hl{✔ সঠিক উত্তর: ঘ) ১৬}}

\begin{quote}
\textbf{ব্যাখ্যা:} \emph{ধরি, সরলরেখাটি = 1 একক}

\emph{∴ সরলরেখার উপর অঙ্কিত বর্গের ক্ষেত্রফল = (1)\^{}2 = 1 বর্গ একক}

\emph{সরলরেখার এক চতুর্থাংশ = (1/4) একক}

\emph{এক চতুর্থাংশের উপর অঙ্কিত বর্গের ক্ষেত্রফল = ( (1/4) )\^{}2 = (1/16) বর্গ
একক}

\emph{⇒ তুলনামূলক গুণফল = 1 ÷ (1/16) = 16 গুণ।}
\end{quote}

\textbf{কোন ক্ষেত্রের সীমানা নির্ধারক বক্ররেখার মোট দৈর্ঘ্যকে কি বলে?}

ক) ক্ষেত্রফল

\textbf{\hl{খ) পরিসীমা ✔}}

গ) আয়তন

ঘ) উচ্চতা

\textbf{\hl{✔ সঠিক উত্তর: খ) পরিসীমা}}

\begin{quote}
\textbf{ব্যাখ্যা:} \emph{যেকোনো জ্যামিতিক ক্ষেত্রের সীমানা বা চারপাশের মোট
দৈর্ঘ্যকে তার পরিসীমা বলা হয়।}
\end{quote}

\textbf{অতিভুজের বিপরীতে থাকে-\/-\/-}

ক) সূক্ষ্মকোণ

খ) স্থুলকোণ

গ) সরলকোণ

\textbf{\hl{ঘ) সমকোণ ✔}}

\textbf{\hl{✔ সঠিক উত্তর: ঘ) সমকোণ}}

\begin{quote}
\textbf{ব্যাখ্যা:} \emph{সমকোণী ত্রিভুজের সমকোণের বিপরীত বাহুকে অতিভুজ বলা হয়।
সুতরাং অতিভুজের বিপরীতে সমকোণ অবস্থান করে।}
\end{quote}

\textbf{দুইটি সন্নিহিত কোণের সমষ্টি ১ সমকোণ হলে একটিকে অপরটির কি বলে?}

\textbf{\hl{ক) পুরক কোণ ✔}}

খ) সম্পূরক কোণ

গ) সন্নিহিত কোণ

ঘ) সরলকোণ

\textbf{\hl{✔ সঠিক উত্তর: ক) পুরক কোণ}}

\begin{quote}
\textbf{ব্যাখ্যা:} \emph{দুটি সন্নিহিত কোণের সমষ্টি ৯০ ডিগ্রি বা ১ সমকোণ হলে
কোণ দুটিকে একে অপরের পূরক কোণ বলা হয়।}
\end{quote}

\textbf{একটি সরল রেখার উপর অঙ্কিত বর্গক্ষেত্রের ক্ষেত্রফল ঐ সরল রেখার তিন
ভাগের একভাগের উপর অঙ্কিত বর্গক্ষেত্রের ক্ষেত্রফলের কতগুণ?}

ক) ৩ গুণ

খ) ৬ গুণ

\textbf{\hl{গ) ৯ গুণ ✔}}

ঘ) ১২ গুণ

\textbf{\hl{✔ সঠিক উত্তর: গ) ৯ গুণ}}

\begin{quote}
\textbf{ব্যাখ্যা:} \emph{ধরি, সরলরেখাটি = 1 একক}

\emph{সম্পূর্ণ সরলরেখার উপর অঙ্কিত বর্গের ক্ষেত্রফল = (1)\^{}2 = 1}

\emph{তিন ভাগের একভাগের উপর অঙ্কিত বর্গের ক্ষেত্রফল = ( (1/3) )\^{}2 =
(1/9)}

\emph{⇒ ক্ষেত্রফলের অনুপাত = 1 ÷ (1/9) = 9 গুণ।}
\end{quote}

\section{19. সম্ভাবনা (probability)-
গণিত}\label{ux9b8ux9aeux9adux9acux9a8-probability--ux997ux9a3ux9a4}

আমরা প্রতিনিয়ত \textbf{`সম্ভাবনা\textquotesingle{}} শব্দটি ব্যবহার করে থাকি।
যেমন এবার এস.এস.সি. পরীক্ষায় যাদবের পাশ করার সম্ভাবনা খুব কম, এশিয়া কাপ
ক্রিকেটে বাংলাদেশের জয়ের সম্ভাবনা বেশি,

আগামীকাল তাপমাত্রা বৃদ্ধি পাওয়ার সম্ভাবনা বেশি, আজ বৃষ্টি হওয়ার সম্ভাবনা কম
ইত্যাদি। অর্থাৎ কোনো ঘটনা ঘটার ক্ষেত্রে অনিশ্চয়তা থাকলেই কেবল আমরা সম্ভাবনার
কথা বলি।

আর অনিশ্চয়তার মাত্রার উপরই ঘটনাটা ঘটার সম্ভাবনা কম বা বেশি হবে তা নির্ভর
করে। কিন্তু কোনো সাংখ্যিক মান দিতে পারে না।

এই অধ্যায়ে আমরা কোনো ঘটনা ঘটার সম্ভাবনার সাংখ্যিক মান নির্ণয়ের বিভিন্ন সূত্র
এবং নির্ণয় প্রণালী সম্পর্কে জানবো এবং নিশ্চিত ঘটনা, অসম্ভব ঘটনা ও সম্ভাব্য ঘটনা
বর্ণনা করতে পারবো।

\begin{itemize}
\item
  প্রাথমিক আলোচনা
\end{itemize}

\begin{itemize}
\item
  দৈব পরীক্ষা (Random Experiment)
\item
  ঘটনা (Event)
\item
  সমসম্ভাব্য ঘটনাবলী (Equally Likely Events)
\item
  পরস্পর বিচ্ছিন্ন ঘটনাবলী (Mutually Exclusive Events)
\item
  অনুকূল ফলাফল (Favourable Outcomes )
\item
  নমুনাক্ষেত্র (Sample Space) ও নমুনা বিন্দু (Sample Point)
\end{itemize}

\begin{itemize}
\item
  মুদ্রার সম্ভাব্যতা
\item
  ছক্কার সম্ভাব্যতা
\item
  তাসের সম্ভাব্যতা
\item
  বলের সম্ভাব্যতা
\end{itemize}

\subsubsection{প্রাথমিক
আলোচনা}\label{ux9aaux9b0ux9a5ux9aeux995-ux986ux9b2ux99aux9a8-7}

Probability হলো কোন কিছু ঘটার সম্ভাবনা। যদি কোন ঘটনা ঘটার কোন প্রকার Chance
না থাকে তাহলে ঐ ঘটনা ঘটার Probability হবে
\textquotesingle0\textquotesingle{} ।

পক্ষান্তরে ঐ ঘটনা যদি নিশ্চিত ঘটার Chance থাকে তাহলে ঐ ঘটনা ঘটার
Probability হবে 100\% বা 1. তাই Probability র মান সব সময় 0 থেকে 1 এর
মধ্যে থাকে ।

ধরুন, আপনার প্রিয় অভিনেতা অস্কারের জন্য মনোনীত হয়েছেন এবং সেই সাথে আরো দুজন
অভিনেতা মনোনীত হয়েছেন। এখন এই তিনজনের মধ্যে আপনার প্রিয় অভিনেতার অস্কার লাভ
করার Probability হল \(\frac{1}{3}\)

সুতরাং Probability হল একটি ভগ্নাংশ যার উপরে থাকবে ঐ ঘটনা ঘটার সংখ্যা এবং
নিচে থাকবে মোট ঘটনা ঘটার সংখ্যা।

Probability =
\(\frac{\text{কোন}\ \text{ঘটনার}\ \text{অনুকূলে}\ \text{নমুনা}\ \text{বিন্দুর}\ \text{সংখ্যা}}{\text{মোট}\ \text{নমুনা}\ \text{বিন্দুর}\ \text{সংখ্যা}}\)

\paragraph{দৈব পরীক্ষা (Random Experiment)
:}\label{ux9a6ux9ac-ux9aaux9b0ux995ux9b7-random-experiment}

যখন কোনো পরীক্ষার সম্ভাব্য সকল ফলাফল আগে থেকে জানা থাকে কিন্তু পরীক্ষাটিতে
কোনো একটি নির্দিষ্ট চেষ্টায় কি ফলাফল আসবে তা নিশ্চিত করে বলা যায় না, একে দৈব
পরীক্ষা বলে। যেমন একটা মুদ্রা নিক্ষেপ পরীক্ষার সম্ভাব্য ফলাফল কি হবে, তা আমরা
আগে থেকেই জানি কিন্তু মুদ্রাটি নিক্ষেপের পূর্বে কোন ফলাফলটি ঘটবে তা আমরা নিশ্চিত
করে বলতে পারি না। সুতরাং মুদ্রা নিক্ষেপ পরীক্ষা একটা দৈব পরীক্ষা।

\paragraph{ঘটনা (Event):}\label{ux998ux99fux9a8-event}

কোনো পরীক্ষার ফলাফল বা ফলাফলের সমাবেশকে ঘটনা বলে। উদাহরণস্বরূপ, একটা ছক্কা
নিক্ষেপ পরীক্ষায় 3 পাওয়া একটি ঘটনা। আবার জোড় সংখ্যা পাওয়াও একটি ঘটনা।

\paragraph{সমসম্ভাব্য ঘটনাবলী (Equally Likely
Events):}\label{ux9b8ux9aeux9b8ux9aeux9adux9acux9af-ux998ux99fux9a8ux9acux9b2-equally-likely-events}

যদি কোনো পরীক্ষার ঘটনাগুলো ঘটার সম্ভাবনা সমান হয় অর্থাৎ একটি অপরটির চেয়ে
বেশি বা কম সম্ভাব্য না হয় তবে ঘটনাগুলোকে সমসম্ভাব্য বলে।

যেমন একটা নিরপেক্ষ মুদ্রা নিক্ষেপে হেড বা টেল আসার সম্ভাবনা সমান। সুতরাং হেড
আসা ও টেল আসা ঘটনা দুইটি সমসম্ভাব্য ঘটনা।

\paragraph{পরস্পর বিচ্ছিন্ন ঘটনাবলী (Mutually Exclusive
Events):}\label{ux9aaux9b0ux9b8ux9aaux9b0-ux9acux99aux99bux9a8ux9a8-ux998ux99fux9a8ux9acux9b2-mutually-exclusive-events}

কোনো পরীক্ষায় যদি একটা ঘটনা ঘটলে অন্যটা অথবা অন্য ঘটনাগুলো না ঘটতে পারে তবে
উক্ত ঘটনাগুলোকে পরস্পর বিচ্ছিন্ন ঘটনা বলে।

যেমন, একটা নিরপেক্ষ মুদ্রা নিক্ষেপ করলে হেড আসা বা টেল আসা দুইটি বিচ্ছিন্ন
ঘটনা।

কেননা হেড আসলে টেল আসতে পারে না। আবার টেল আসলে হেড আসতে পারে না। অর্থাৎ
হেড ও টেল একসাথে আসতে পারে না।

\paragraph{অনুকূল ফলাফল (Favourable Outcomes )
:}\label{ux985ux9a8ux995ux9b2-ux9abux9b2ux9abux9b2-favourable-outcomes}

কোনো পরীক্ষায় একটা ঘটনার স্বপক্ষের ফলাফল হলো ঘটনার অনুকূল ফলাফল। একটি ছক্কা
নিক্ষেপ করলে বিজোড় সংখ্যা হওয়ার অনুকূল ফলাফল 3 টি।

\paragraph{নমুনাক্ষেত্র (Sample Space) ও নমুনা বিন্দু (Sample Point)
:}\label{ux9a8ux9aeux9a8ux995ux9b7ux9a4ux9b0-sample-space-ux993-ux9a8ux9aeux9a8-ux9acux9a8ux9a6-sample-point}

কোনো দৈব পরীক্ষার সম্ভাব্য সকল ফলাফল নিয়ে গঠিত সেটকে নমুনাক্ষেত্র বলে।

একটা মুদ্রা নিক্ষেপ করলে দুইটি সম্ভাব্য ফলাফল পাওয়া যায়, যথা হেড ও টেল। এখন S
দ্বারা এ পরীক্ষণের ফলাফলের সেটকে সূচিত করলে আমরা লিখতে পারি S = \{H,T\}।

সুতরাং উক্ত পরীক্ষার নমুনাক্ষেত্র, S = \{H,T\}। মনে করা যাক দুইটি মুদ্রা একসাথে
নিক্ষেপ করা হলো। তাহলে নমুনাক্ষেত্রটি হবে S = \{HH, HT, TH,TT\}।

নমুনাক্ষেত্রের প্রতিটি উপাদানকে ফলাফলের নমুনা বিন্দু বলে। একটা মুদ্রা একবার
নিক্ষেপ পরীক্ষায় নমুনাক্ষেত্র S = \{H,T\} এবং এখানে H,T প্রত্যেকেই এক একটা নমুনা
বিন্দু ।

\textbf{Note:} সবসময় মনে রাখতে হবে যে কোন ঘটনা ঘটার সম্ভাবনা ০ ও ১ এর
মধ্যেই সীমাবদ্ধ থাকবে। কোন ঘটনা নিশ্চিতভাবে ঘটলে তার মান হবে ১ এবং
নিশ্চিতভাবে না ঘটলে তার মান হবে ০ ।

\subsubsection{মুদ্রার
সম্ভাব্যতা}\label{ux9aeux9a6ux9b0ux9b0-ux9b8ux9aeux9adux9acux9afux9a4}

একটি মুদ্রার দুটি পার্শ, একটি H = Head, এবং অপরটি T = Tail.

\textbf{প্রশ্ন নম্বর- ১. একটি মুদ্রা একবার নিক্ষেপ করলে head পড়ার Probability
কত?}

\begin{itemize}
\item
  (ক) (1/2)
\item
  (খ) (1/3)
\item
  (গ) (1/4)
\item
  (ঘ) (1/5)
\end{itemize}

\textbf{\hl{উত্তরঃ (ক)}} \(\frac{1}{2}\)

\textbf{সমাধান:}

একটি মুদ্রা ১ বার নিক্ষেপ করলে আমরা ১ বারই Head পড়তে দেখবো। তাই Probability
ভগ্নাংশের উপরের সংখ্যা হবে ১। আবার একবার মুদ্রা নিক্ষেপ করলে দুটি সম্ভাব্য ফলাফল
পাওয়া যেতে পারে। যথা Head অথবা Tail. তাই মোট

সম্ভাব্য ফলাফল সংখ্যা হবে ২। সুতরাং Head পড়ার Probability \(\frac{1}{2}\)

\textbf{Note:} একটি মুদ্রার যেহেতু ২টি দিক তাই মুদ্রা নিক্ষেপের ক্ষেত্রে, যতবার
মুদ্রা নিক্ষেপ করা হবে, \textquotesingle2\textquotesingle{} এর উপর ততটি
Power দিলে ।

Total Outcome বা মোট ফলাফল পাওয়া যাবে। যেমন একটি মুদ্রা ২ বার নিক্ষেপ করলে
মোট ফলাফল হবে \(\text{২}^{\text{২}}\) = ৪ টি। ৩ বার নিক্ষেপ করলে মোট ফলাফল হবে =
\(\text{২}^{\text{৩}}\) = ৮টি।

\textbf{২. একটি মুদ্রা দু\textquotesingle বার নিক্ষেপ করা হলে নমুনা ক্ষেত্র কি
হবে?}

2 বার নিক্ষেপ করলে নমুনাক্ষেত্র = (HH, HT, TH, TT), তখন নমুনা বিন্দুর সংখ্যা =
\(2^{2}\) = 4

\textbf{এই ফলাফল গুলো কিভাবে বের করতে হয়?}

সেট অধ্যায়ে পাওয়ার সেট বের করার নিয়মে হিসেব করতে হয়। যেমন:

দুবার নিক্ষেপ করা হলে, (H, T ) × (H,T) = \{ HH \} , \{ HT\} , \{TH\}
,\{TT\}

আবার তিনবার নিক্ষেপ করা হলে, (H, T ) × (H,T) × (H,T) = \{ HHH \} ,
\{HHT\} \{HTT\} , \{HTH\}, \{THH\} ,\{TTH\} , \{THT\}, \{TTT\}

\textbf{৩. একটি মুদ্রা ২ বার নিক্ষেপ করলে ২ বারই Head পড়ার Probability কত?}

\textbf{ব্যাখ্যাঃ} ২ বার নিক্ষেপ করলে মোট ফলাফল হবে \(\text{২}^{\text{২}}\) = ৪ টি। যেমন:
1. HH 2. HT 3. TH 4. TT এখন এই চারটি

ফলাফলের মধ্যে দুবারই Head পড়ার সম্ভাবনা হবে ১টি। তাই উত্তর: \(\frac{1}{4}\)

\textbf{প্রশ্ন নম্বর-} ৪. দুইটি মুদ্রা একসাথে নিক্ষেপ করা হলে প্রথম মুদ্রায় H এবং
২য় মুদ্রায় T আসার সম্ভাবনা কত? {[} সরকারি মাধ্যমিক বিদ্যালের সহ: শিক্ষক: ২০০৬
{]}

\begin{itemize}
\item
  (ক) (১/৪)
\item
  (খ) (১/৩)
\item
  (গ) (১/২)
\item
  (ঘ) ১
\end{itemize}

\textbf{\hl{উত্তরঃ (ক)}} \(\frac{\text{১}}{\text{৪}}\)

\textbf{সমাধান:}

দুইটি মুদ্রা একসাথে নিক্ষেপ করা হলে মোট ফলাফল = \(\text{২}^{\text{২}}\) = ৪টি

ফলাফল গুলো হলো HH, HT, TH, TT = ৪ টি

ফলাফলগুলোর মধ্যে প্রথমটিতে H এবং পরেরটিতে T আছে কেবল একটিতে = HT

সুতরাং প্রথম H এবং পরের মুদ্রায় T আসার সম্ভাবনা = \(\frac{\text{১}}{\text{৪}}\)

\textbf{৫. একটি মুদ্রা ২ বার নিক্ষেপ করলে অন্তত ১ বার Head পড়ার সম্ভাবনা কত?}

\textbf{ব্যাখ্যা:} মোট ফলাফল = 4 টি। যেখানে 3 টিতেই অন্তত একবার Head আসবে,
এবং ১ বার সবগুলোই Tail আসবে। তাই উত্তর: \(\frac{\text{৩}}{\text{৪}}\)

\textbf{৬. একটি মুদ্রা ৩ বার নিক্ষেপ করা হলে তিনটি মুদ্রায় Head পাবার সম্ভাবনা
কত?}

\textbf{ব্যাখ্যা}ঃ মুদ্রা ৩ বার নিক্ষেপ করলে মোট ফলাফল ৮ এবং শুধুমাত্র একভাবেই
তিন বারই Head পড়তে পারে। তাই উত্তর \(\frac{\text{১}}{\text{৮}}\)

\textbf{৭. একটি মুদ্রা ৩ বার নিক্ষেপ করা হলে তিনটি মুদ্রায় Tail পাবার সম্ভাবনা
কত?}

\textbf{ব্যাখ্যাঃ} Head এর মতই শুধু এক ক্ষেত্রে তিন বারই Tail আসতে পারে। তাই
উত্তর হবে।\(\frac{\text{১}}{\text{৮}}\)

\textbf{প্রশ্ন নম্বর-} ৮. যদি ২টি মুদ্রা নিক্ষেপ করা হয় তাহলে একটা head এবং
অন্যটি tail উঠার সম্ভাবনা কত? {[} Pubali Bank Ltd. SO 2013 {]}

\begin{itemize}
\item
  (ক) (1/2)
\item
  (খ) (1/3)
\item
  (গ) (1/4)
\item
  (ঘ) (1/5)
\end{itemize}

\textbf{\hl{উত্তরঃ (গ)}} \(\frac{1}{4}\)

\textbf{সমাধান:}

মোট ফলাফল = \{ HH, TT, TH, HT\} এবং অনুকুল ফলাফল = \{ HT \}

একটি coin toss করলে head/ tail পাবার Probability দুইটির জন্য =
\(\frac{1}{2} \times \frac{1}{2} = \frac{1}{4}\)

\subsubsection{ছক্কার
সম্ভাব্যতা:}\label{ux99bux995ux995ux9b0-ux9b8ux9aeux9adux9acux9afux9a4}

একটি ছক্কার মোট ছয়টি পার্শ্ব থাকায়, ছক্কাটি একবার নিক্ষেপ করলে মোট ফলাফল হবে
৬ টি। আবার ছক্কাটি একাধিক বার নিক্ষেপ করা হলে ৬ এর উপর তত Power দিলে মোট
ফলাফল পাওয়া যাবে। যেমন:

ছকাটি ২ বার নিক্ষেপ করলে মোট ফলাফল হবে \(\text{৬}^{\text{২}}\) = ৩৬টি ।

ছক্কাটি ৩ বার নিক্ষেপ করলে ছোট ফলাফল হবে = \(\text{৬}^{\text{৩}}\) = ২১৬টি

\textbf{৯. একটি ছক্কা একবার নিক্ষেপ করলে উপরের পিঠে ৩ আসার সম্ভাবনা কত?}

\textbf{ব্যাখ্যা}: ১ বার নিক্ষেপ করলে মোট ফলাফল ৬ টি। ৩ আসতে পারে শুধু একবার
তাই উত্তর হবে \(\frac{1}{6}\)

\textbf{১০. একটি ছক্কা নিক্ষেপ পরীক্ষায় উপরের পিঠে জোড় সংখ্যা আসার সম্ভাবনা
কত?}

\textbf{ব্যাখ্যা:} ছক্কায় মোট জোড় সংখ্যা থাকে ৩টি (২, ৪, ৬) সুতরাং উত্তরটি
হবে। \(\frac{3}{6}\) = \(\frac{1}{2}\) {[} সম্ভাব্যতার প্রশ্নে যে ভগ্নাংশই
উত্তর আসুক না কেন, তাকে লঘিষ্ঠ আকারে লিখতে হয়। {]}

\textbf{১১. একটি ছক্কা ২ বার নিক্ষেপ করা হল। উপরের পিঠে একই সংখ্যা আসার
সম্ভাবনা কত?}

\textbf{ব্যাখ্যা :} ছক্কাকে ২ বার নিক্ষেপ করলে মোট ফলাফল হবে, \(\text{৬}^{\text{২}}\) =
৩৬টি। আবার একই সংখ্যা আসতে পারে ৬ বার যেমন ২ বারই ১, ২ বারই ২. .........২
বারই ৫, ২ বারই ৬। সুতরাং উত্তরটি হবে

\(\frac{6}{36} = \frac{1}{6}\)

\textbf{১২. একটি ছক্কা ৩ বার নিক্ষেপ করা হল। উপরের পিঠে একই সংখ্যা আসার
সম্ভাবনা কত?}

\textbf{ব্যাখ্যাঃ} ছক্কাকে ৩ বার নিক্ষেপ করলে মোট ফলাফল হবে \(\text{৬}^{\text{৩}}\) = ২১৬,
এবং প্রতিবারই একই সংখ্যাগুলো হতে পারে, প্রতিবারই ১ , অথবা প্রতিবারই ২ , অথবা
৩ অথবা ৪ অথবা ৫ অথবা ৬। অর্থাৎ মোট ছয়বার প্রতিক্ষেত্রে একই সংখ্যা আসতে পারে।
এখন

উত্তরটি হবে \(\frac{6}{216} = \frac{1}{36}\)

\textbf{১৩. একটি ছক্কা ১ বার নিক্ষেপ করলে ২ বা ৩ দ্বারা বিভাজ্য সংখ্যা আসার
সম্ভাবনা কত?}

\textbf{ব্যাখ্যা}ঃ ১ বার নিক্ষেপ করলে মোট ফলাফল = ৬টি। তা থেকে ২ দ্বারা
বিভাজ্য সংখ্যা হবে ২টি (২, ৪ ও ৬) আবার

৩ দ্বারা বিভাজ্য সংখ্যা হবে ২টি (৩ ও ৬), সুতরাং ২ ও ৩ দিয়ে বিভাজ্য সংখ্যা হল
(২.৩.৪ এবং ৬) তাই উত্তরটি হবে = \(\frac{4}{6} = \frac{2}{3}\)

Note: প্রশ্নে ২ এবং ৩ বলা থাকলে? তখন শুধু ১টি ফলাফল হতো শুধু ৬। উত্তর হতো
\(\frac{\text{১}}{\text{৬}}\)

\textbf{প্রশ্ন নম্বর-} ১৪. ২০. একটি নিরপেক্ষ ছক্কা ও একটি মুদ্রা একবার নিক্ষেপ
করা হলে ছক্কার ৫ এবং মুদ্রার H সম্ভাবনা কত? {[} পরিবেশ অধিদপ্তরের সহকারী
পরিচালক : ২০০৬ {]}

\begin{itemize}
\item
  (ক) (১/২)
\item
  (খ) (১/৪)
\item
  (গ) (৫/১২)
\item
  (ঘ) (১/১২)
\end{itemize}

\textbf{\hl{উত্তরঃ (ঘ)}} \(\frac{\text{১}}{\text{১২}}\)

\textbf{সমাধান:}

একটি ছক্কা ও একটি মুদ্রা নিক্ষেপ পরীক্ষাকে দুইধাপ হিসেবে বিবেচনা করি। প্রথম
ধাপে হকা নিক্ষেপে ৬ টি ফ (১,২,৩,৪,৫,৬) আসতে পারে। দ্বিতীয় ধাপে মুদ্রা
নিক্ষেপে ২টি ফলাফল (H অথবা T) আসতে পারে।

তাহলে নমুনা ক্ষেত্রটি হবে {[} 1H, IT, 2H , 2T, 3H, 3T, 4H, 4T, 5H, 5T, 6H,
6T {]}

এখানে মোট নমুনা বিন্দু ১২ টি। সুতরাং ছক্কায় ৫ এবং মুদ্রায় H আসার সম্ভাবনা P(5H)
= \(\frac{\text{১}}{\text{১২}}\)

\subsubsection{তাসের
সম্ভাব্যতা}\label{ux9a4ux9b8ux9b0-ux9b8ux9aeux9adux9acux9afux9a4}

একটি তাসের প্যাকেটে মোট ৫২ টি তাস থাকে। যাদের মোট চারটি ভাগে ভাগ করা যায়,
যথা:

{\def\LTcaptype{none} % do not increment counter
\begin{longtable}[]{@{}
  >{\raggedright\arraybackslash}p{(\linewidth - 6\tabcolsep) * \real{0.2006}}
  >{\raggedright\arraybackslash}p{(\linewidth - 6\tabcolsep) * \real{0.2006}}
  >{\raggedright\arraybackslash}p{(\linewidth - 6\tabcolsep) * \real{0.2006}}
  >{\raggedright\arraybackslash}p{(\linewidth - 6\tabcolsep) * \real{0.2006}}@{}}
\toprule\noalign{}
\begin{minipage}[b]{\linewidth}\raggedright
\textbf{হরতন}
\end{minipage} & \begin{minipage}[b]{\linewidth}\raggedright
\textbf{রুইতন}
\end{minipage} & \begin{minipage}[b]{\linewidth}\raggedright
\textbf{ইস্কাপন}
\end{minipage} & \begin{minipage}[b]{\linewidth}\raggedright
\textbf{চিড়া}
\end{minipage} \\
\midrule\noalign{}
\endhead
\bottomrule\noalign{}
\endlastfoot
Hearts & Diamonds & Spades & Clubs \\
১৩ টি ( ১টি রাজা ও ১টি রাণী ) & ১৩ টি ( ১টি রাজা ও ১টি রাণী ) & ১৩ টি (
১টি রাজা ও ১টি রাণী ) & ১৩ টি ( ১টি রাজা ও ১টি রাণী ) \\
\multicolumn{4}{@{}>{\raggedright\arraybackslash}p{(\linewidth - 6\tabcolsep) * \real{0.8022} + 6\tabcolsep}@{}}{%
প্রত্যেক ভাগে মোট ১৩ টি করে কার্ড থাকে এবং এই ১৩টির মধ্যে ১টি করে রাজা ও
একটি করে রাণী থাকে।} \\
\end{longtable}
}

\textbf{১৫. এক প্যাকেট তাস থেকে দৈবভাবে ১টি তাস নিলে তা রাজা হওয়ার
Probability কত?}

\textbf{ব্যাখ্যা}: এক প্যাকেট তাসের মধ্যে যেহেতু মোট ৫২ টি তাস থাকে যার মধ্যে
৪টি রাজা, তাই একটি ভাগ নিলে রাজা হওয়ার সম্ভাবনা = \(\frac{4}{52}\) বা
\(\frac{1}{13}\)

\textbf{১৬. এক প্যাকেট কার্ড থেকে একটি রাজা নেবার পর আবার একটি কার্ড নেয়া
হল, পরবর্তী কার্ডটি Hearts এর রাজা হওয়ার সম্ভাবনা কত?}

\textbf{ব্যাখ্যা}: প্রথমেই একটি রাজা নেয়াতে আর কার্ড থাকল ৫১টি সেখানে আরো
তিনটি রাজা থাকলেও Hearts এর রাজা আছে মাত্র ১ টি । তাই উত্তর
\(\frac{\text{১}}{\text{৫১}}\)

১৭. ১ প্যাকেট তাস থেকে দৈবভাবে একটি তাস নেয়া হল। তাসটি হরতন হওয়ার সম্ভাবনা
কত? {[} কর্মসংস্থান ও প্রশিক্ষণ বুরোর উপপরিচালক : ২০০৭ {]}

১টি প্যাকেট বা ৫২টি তাসের মধ্যে হরতন থাকে মোট ১৩টি।

সুতরাং, হরতন হওয়ার সম্ভাবনা = \(\frac{\text{১৩}}{\text{৫২}} = \frac{\text{১}}{\text{৪}}\)

১৮. এক প্যাকেট তাস থেকে দৈবভাবে 2 টি তাস নেয়া হল। তাস দুটি রাজা হওয়ার
সম্ভাবনা কত? {[} সোনালী ব্যাংক অফিসার: ২০১৫ {]}

\textbf{সমাধান:} ২টি তাস নেয়ার কথা বলায় এখানে সমাবেশের সূত্র প্রয়োগ করতে
হবে)

একটি প্যাকেটে তাস আছে ৫২টি। ৫২টি তাস থেকে ২টি ভাগ নেয়ার মোট উপায়
\(52_{c_{2}}\) = \(\frac{52 \times 51}{2 \times 1}\) = 1326

আবার রাজা আছে ৪টি। এখন ৪টি থেকে ২টি নেয়ার উপায় = \(4_{c_{2}}\) =
\(\frac{4 \times 3}{2 \times 1}\) = 6

সুতরাং সম্ভাবনা = \(\frac{6}{1326}\) = \(\frac{1}{221}\)

\subsubsection{বলের
সম্ভাব্যতা}\label{ux9acux9b2ux9b0-ux9b8ux9aeux9adux9acux9afux9a4}

Suggestion: এ অংশে বিভিন্ন .মার্বেল, ফল, ফলের অংকের সম্ভাব্যতা যাচাই করা
হয়েছে। সমাবেশ অধ্যায়ের অনেক অংকের সাথে এই অংশের অনেক মিল থাকায় এগুলো করার
আগে সমাবেশ অধ্যায় শিখে আসার পরামর্শ রইলো।

\textbf{১৯. একটি Box-এ 6 টি white marbles আছে, সেখান হতে ১টি Black
marbel তোলার সম্ভাবনা কত?}

\textbf{ব্যাখ্যা}: সাদা মার্বেলের বক্স থেকে কালো মার্বেল পাওয়ার সম্ভাবনা নেই।
তাই উত্তর: ০

\textbf{২০. একটি পাত্রে ৫টি সাদা মোজা ও ৬টি কালো মোজা আছে, অন্ধকারে মোট
কয়টি মোজা তুললে নিশ্চিতভাবে বলা যাবে যে, তাদের মধ্যে একটি সাদা মোজা আছে?}

\textbf{ব্যাখ্যা}: মোট ছয়টি মোজা তুললেও ৬টিই কালো হতে পারে। কিন্তু ৭ম বার
মোজা তোলার পর অবশ্যই তা সাদা মোজা হবে। কেননা মোট ৬টি কালো মোজা সেই পাত্রে
ছিল।

\textbf{২১. রাকিবের পকেটে ২টি লাল, ৩টি নীল, ৫টি কালো ও ৮টি সবুজ মার্বেল
আছে। কমপক্ষে কতটি মার্বেল নিলে প্রত্যেক রংয়ের অন্তত ১টি করে হলেও মার্বেল
পাবে?}

\textbf{ব্যাখ্যা}: মোট মার্বেলের সংখ্যা = ২+৩+৫+৮ = ১৮টি এর মধ্যে সবথেকে কম
মার্বেলের সংখ্যা থেকে ১টি মার্বেল নিয়ে বাকী সবগুলো মার্বেল নেয়ার পর
নিশ্চিতভাবে বলা সম্ভব হবে যে প্রতি রংয়ের একটি করে মার্বেল পাওয়া যাবে।

সুতরাং মার্বেল নিতে হবে = ৮+৫+৩+১ = ১৭টি (কেননা ১৬টি মার্বেল নেয়ার পর ও
নিশ্চিত হওয়া সম্ভব নয় )

২২. একটি ঝুড়িতে ৬টি আপেল এবং ৩টি কমলা আছে। কমপক্ষে একটি কমলা পেতে কতগুলো ফল
তুলতে হবে? {[} M.T.B.L. OFF: 13 {]}

Solution: ১টি আপেল ও ৩টি কমলার মধ্যে ৬টি আপেল তুলে নিয়ে আর ১টি উঠালেই কমলা
পাওয়া যাবে।

সুতরাং কমপক্ষে ১টি কমলাসহ মোট ফল উঠাতে হবে (৬+১) = ৭টি

২৩. একটি থলিতে ১২ টি নীল বল, ১৪ টি সাদা বল এবং ২২ টি কালো বল আছে।
দৈবভাবে একটি বল তুললে সেটি নীল না হবার সম্ভাবনা কত? {[} অগ্রনী ব্যাংক সিনিয়র
অফিসার ২০১৪ {]}

\textbf{সমাধান:} থলিতে মোট বল আছে = ১২+১৪+২২ = ৪৮টি এবং নীল বল আছে ১২
টি।

নীল বল উঠার সম্ভাবনা = \(\frac{\text{১২}}{\text{৪৮}} = \frac{\text{১}}{\text{৪}}\)

তাহলে নীল বল না উঠার সম্ভাবনা = \(\text{১} - \frac{\text{১}}{\text{৪}} = \frac{\text{৩}}{\text{৪}}\)

২৪. একটি থলিতে 6 টি নীল বল, ৪ টি সাদা বল এবং 10 টি কালো বল আছে। দৈবভাবে
একট বল তুললে সেটি সাদা না হবার সম্ভাবনা কত? {[} ৩৭-তম বিসিএস প্রিলি {]}

\textbf{ব্যাখ্যা:} মোট বলের সংখ্যা = ৬+৮+১০ = ২৪টি। তাই দৈবভাবে বল নিলে মোট
ফলাফল ২৪টি। এখন সাদা বল না হওয়ার ঘটনা বা ফলাফল হবে ৬+১০ = ১৬টি (কারণ ১৬
বার নীল ও কালো বল উঠতে পারে) তাহলে উত্তর \(\frac{\text{১৬}}{\text{২৪}} = \frac{\text{২}}{\text{৩}}\)

\section{20. সম্ভাব্যতা MCQ প্রশ্ন উত্তর
অনুশীলন}\label{ux9b8ux9aeux9adux9acux9afux9a4-mcq-ux9aaux9b0ux9b6ux9a8-ux989ux9a4ux9a4ux9b0-ux985ux9a8ux9b6ux9b2ux9a8}

\textbf{৩০ থেকে ৪০ পর্যন্ত সংখ্যা থেকে যে কোন একটিকে ইচ্ছেমতো নিলে সে
সংখ্যাটি মৌলিক অথবা ৫ এর গুণিতক হওয়ার সম্ভাবনা কত?}

\textbf{\hl{ক) (5/11) ✔}}

খ) (1/2)

গ) (3/5)

ঘ) (6/11)

\textbf{\hl{✔ সঠিক উত্তর: ক) (5/11)}}

\begin{quote}
\textbf{ব্যাখ্যা:} \emph{৩০ থেকে ৪০ পর্যন্ত মোট সংখ্যা আছে = 40−30+1=11 টি}

\emph{৩০ থেকে ৪০ পর্যন্ত মৌলিক সংখ্যাগুলো হলো: ৩১, ৩৭ (মোট ২ টি)}

\emph{৩০ থেকে ৪০ পর্যন্ত ৫ এর গুণিতকগুলো হলো: ৩০, ৩৫, ৪০ (মোট ৩ টি)}

\emph{৩০ থেকে ৪০ পর্যন্ত সংখ্যাগুলোর মধ্যে মৌলিক অথবা ৫ এর গুণিতক উভয় শর্ত মিলে
মোট অনুকূল সংখ্যা: ৩০, ৩১, ৩৫, ৩৭, ৪০ (মোট ৫ টি)}

\emph{∴ নির্ণেয় সম্ভাবনা = (5/11)}
\end{quote}

\textbf{একটি থলিতে 6 টি নীল বল, 8 টি সাদা বল এবং 10 টি কালো বল আছে।
দৈবভাবে একটি বল তুললে সেটি সাদা না হবার সম্ভাবনা কত?}

\textbf{\hl{ক) (2/3) ✔}}

খ) (1/3)

গ) (3/4)

ঘ) (1/4)

\textbf{\hl{✔ সঠিক উত্তর: ক) (2/3)}}

\begin{quote}
\textbf{ব্যাখ্যা:} \emph{থলিতে মোট বল = (6+8+10)=24 টি}

\emph{বলটি সাদা হওয়ার সম্ভাবনা = (8/24)=(1/3)}

\emph{∴ বলটি সাদা না হবার সম্ভাবনা = 1−(1/3)=(2/3)}
\end{quote}

\textbf{আবহাওয়া অফিসের রিপোর্ট অনুযায়ী-২০১৫ সালের জুলাই মাসের ২য় সপ্তাহে
বৃষ্টি হয়েছে মোট 5 দিন। ঐ সপ্তাহে বুধবার বৃষ্টি না হওয়ার সম্ভাবনা কত?}

ক) 1

খ) (5/7)

\textbf{\hl{গ) (2/7) ✔}}

ঘ) (1/7)

\textbf{\hl{✔ সঠিক উত্তর: গ) (2/7)}}

\begin{quote}
\textbf{ব্যাখ্যা:} \emph{1 সপ্তাহ = 7 দিন}

\emph{বৃষ্টি হয়েছে 5 দিন}

\emph{∴ যে কোনো একদিন বৃষ্টি হওয়ার সম্ভাবনা = (5/7)}

\emph{∴ বুধবার বৃষ্টি হওয়ার সম্ভাবনা = (5/7)}

\emph{∴ বুধবার বৃষ্টি না হওয়ার সম্ভাবনা = 1−(5/7)=(7−5/7)=(2/7)}
\end{quote}

\textbf{একটি অফিসের কর্মীদের ৬০\% গড়ে দৈনিক এক ঘণ্টা করে গাড়িতে ভ্রমণ করে।
এই ৬০\% কর্মীর ২৫\% আবার ট্রেনে ভ্রমণ করে। একজন কর্মীকে দৈবভাবে চয়ন করা হলে
তা ট্রেনে ভ্রমণকারী হওয়ার সম্ভাবনা কত?}

ক) ০.১০

\textbf{\hl{খ) ০.১৫ ✔}}

গ) ০.২০

ঘ) ০.২৫

\textbf{\hl{✔ সঠিক উত্তর: খ) ০.১৫}}

\begin{quote}
\textbf{ব্যাখ্যা:} \emph{ধরি, মোট কর্মী সংখ্যা 100 জন।}

\emph{∴ গাড়িতে ১ ঘণ্টা ভ্রমণ করে এমন যাত্রী সংখ্যা = 60 জন}

\emph{∴ ট্রেনে ভ্রমণ করে এমন কর্মীর সংখ্যা = 60 এর 25\%=15 জন}

\emph{∴ ট্রেনে ভ্রমণকারী হওয়ার সম্ভাবনা = (15/100)=0.15}
\end{quote}

\textbf{চাকরি পাওয়ার সম্ভাবনা (4/5) হলে চাকরি না পাওয়ার সম্ভাবনা কত?}

\textbf{\hl{ক) (1/5) ✔}}

খ) (2/5)

গ) (3/5)

ঘ) (4/5)

\textbf{\hl{✔ সঠিক উত্তর: ক) (1/5)}}

\begin{quote}
\textbf{ব্যাখ্যা:} \emph{নিশ্চিত ঘটনার সম্ভাবনা ১}

\emph{আবার, ঘটনা ঘটার সম্ভাবনা ও না ঘটার সম্ভাবনার সমষ্টি ১}

\emph{∴ চাকরি না পাওয়ার সম্ভাবনা = 1−(4/5)=(5−4/5)=(1/5)}
\end{quote}

\textbf{একটি ছক্কা নিক্ষেপ করলে ছক্কায় জোড় সংখ্যা পাওয়ার সম্ভাবনা কত?}

ক) ৩০\%

খ) ৪০\%

\textbf{\hl{গ) ৫০\% ✔}}

ঘ) ৬০\%

\textbf{\hl{✔ সঠিক উত্তর: গ) ৫০\%}}

\begin{quote}
\textbf{ব্যাখ্যা:} \emph{ছক্কায় অনুকূল ঘটনা ৬ টি (১, ২, ৩, ৪, ৫, ৬) এবং জোড়
সংখ্যা ৩টি (২, ৪, ৬)}

\emph{∴ জোড় সংখ্যা পাওয়ার সম্ভাবনা = (3/6)×100\%=50\%}
\end{quote}

\textbf{একটি ছক্কাকে একবার নিক্ষেপ করা হলে, 2 থেকে বড় সংখ্যা পাবার সম্ভাবনা
কত?}

ক) (1/2)

খ) (1/3)

গ) (5/6)

\textbf{\hl{ঘ) (2/3) ✔}}

\textbf{\hl{✔ সঠিক উত্তর: ঘ) (2/3)}}

\begin{quote}
\textbf{ব্যাখ্যা:} \emph{সম্ভাব্যতা = (অনুকূল নমুনা বিন্দুর সংখ্যা/মোট নমুনা বিন্দুর
সংখ্যা)}

\emph{= (4/6)=(2/3)}
\end{quote}

\textbf{If a coin is tossed once, what is the probability of getting a
\textquotesingle HEAD\textquotesingle{} / যদি একটি মুদ্রা টস করা হয়, তবে
\textquotesingle HEAD\textquotesingle{} এর সম্ভাবনা কত?}

\textbf{\hl{ক) 0.5 ✔}}

খ) 1

গ) 0.25

ঘ) 2

\textbf{\hl{✔ সঠিক উত্তর: ক) 0.5}}

\begin{quote}
\textbf{ব্যাখ্যা:} \emph{Here, Total outcome is 2
(\textquotesingle Head\textquotesingle{} and
\textquotesingle Tail\textquotesingle)}

\emph{Favourable outcome is 1 (\textquotesingle Head\textquotesingle)}

\emph{Therefore, Probability = (Favourable outcome/Total outcome)}

\emph{= (1/2)=0.5}
\end{quote}

\textbf{দুইটি মুদ্রা একসাথে নিক্ষেপ করা হলে প্রথম মুদ্রায় H এবং ২য় মুদ্রায় T আসার
সম্ভাবনা কত?}

\textbf{\hl{ক) (১/৪) ✔}}

খ) (১/৩)

গ) (১/২)

ঘ) ১

\textbf{\hl{✔ সঠিক উত্তর: ক) (১/৪)}}

\begin{quote}
\textbf{ব্যাখ্যা:} \emph{দুইটি মুদ্রা একসাথে নিক্ষেপ করা হলে মোট সম্ভাব্য ফলাফল
= ২\^{}২ = ৪টি}

\emph{ফলাফলগুলো হলো: \{HH, HT, TH, TT\}}

\emph{এখানে, অনুকূল ফলাফল (প্রথমটিতে H এবং ২য়টিতে T) রয়েছে কেবল ১টিতে = HT}

\emph{∴ নির্ণেয় সম্ভাবনা = (১/৪)}
\end{quote}

\textbf{যদি ২টি মুদ্রা নিক্ষেপ করা হয় তাহলে একটা head এবং অন্যটি tail ওঠার
সম্ভাবনা কত?}

\textbf{\hl{ক) (১/২) ✔}}

খ) (১/৩)

গ) (১/৪)

ঘ) (১/৫)

\textbf{\hl{✔ সঠিক উত্তর: ক) (১/২)}}

\begin{quote}
\textbf{ব্যাখ্যা:} \emph{২টি মুদ্রা নিক্ষেপ করা হলে মোট সম্ভাব্য ফলাফল = \{HH,
HT, TH, TT\} = ৪টি।}

\emph{প্রশ্নানুসারে, একটি head এবং অন্যটি tail ওঠার অনুকূল ফলাফল = \{HT, TH\}
= ২টি।}

\emph{∴ নির্ণেয় সম্ভাবনা = (২/৪) = (১/২)}

\emph{গাণিতিক বিশ্লেষণ:}

\emph{প্রশ্নানুযায়ী নির্দিষ্ট কোনো ক্রম উল্লেখ না থাকায় \textquotesingle একটি
head এবং অন্যটি tail\textquotesingle{} বলতে HT এবং TH উভয়ই গ্রহণযোগ্য। ফলে
সঠিক উত্তর (ক) (১/২) ।}
\end{quote}

\textbf{একটি নিরপেক্ষ ছক্কা ও একটি মুদ্রা একবার নিক্ষেপ করা হলে ছক্কার ৫ এবং
মুদ্রার H সম্ভাবনা কত?}

ক) (১/২)

খ) (১/৮)

গ) (৫/১২)

\textbf{\hl{ঘ) (১/১২) ✔}}

\textbf{\hl{✔ সঠিক উত্তর: ঘ) (১/১২)}}

\begin{quote}
\textbf{ব্যাখ্যা:} \emph{একটি নিরপেক্ষ ছক্কা নিক্ষেপে মোট ফলাফল = ৬টি \{১, ২,
৩, ৪, ৫, ৬\}}

\emph{∴ ছক্কায় \textquotesingle{}৫\textquotesingle{} আসার সম্ভাবনা = (১/৬)}

\emph{একটি মুদ্রা নিক্ষেপে মোট ফলাফল = ২টি \{H, T\}}

\emph{∴ মুদ্রায় \textquotesingle H\textquotesingle{} আসার সম্ভাবনা = (১/২)}

\emph{যেহেতু ঘটনা দুটি স্বাধীন,}

\emph{∴ ছক্কায় ৫ এবং মুদ্রায় H আসার সম্ভাবনা = (১/৬) × (১/২) = (১/১২)}
\end{quote}

\textbf{১ প্যাকেট তাস থেকে দৈবভাবে একটি তাস নেয়া হল। তাসটি হরতন হওয়ার
সম্ভাবনা কত?}

\textbf{\hl{ক) (১/৪) ✔}}

খ) (১/১৩)

গ) (১/২৬)

ঘ) (১/৫২)

\textbf{\hl{✔ সঠিক উত্তর: ক) (১/৪)}}

\begin{quote}
\textbf{ব্যাখ্যা:} \emph{১ প্যাকেট তাসে মোট তাসের সংখ্যা = ৫২টি।}

\emph{এর মধ্যে হরতন তাসের সংখ্যা = ১৩টি।}

\emph{∴ দৈবভাবে একটি তাস নিলে সেটি হরতন হওয়ার সম্ভাবনা = (১৩/৫২) = (১/৪)}
\end{quote}

\textbf{এক প্যাকেট তাস থেকে দৈবভাবে 2 টি তাস নেয়া হল। তাস দুটি রাজা হওয়ার
সম্ভাবনা কত?}

\textbf{\hl{ক) (১/২২১) ✔}}

খ) (১/১৩)

গ) (১/২৬)

ঘ) (১/৬৬৩)

\textbf{\hl{✔ সঠিক উত্তর: ক) (১/২২১)}}

\begin{quote}
\textbf{ব্যাখ্যা:} \emph{১ প্যাকেট তাসে মোট তাসের সংখ্যা = ৫২টি।}

\emph{৫২টি তাস থেকে দৈবভাবে ২টি তাস বাছাই করার মোট উপায় = C২৫২ =
(৫২×৫১/২×১) = ১৩২৬}

\emph{এক প্যাকেটে মোট রাজা থাকে = ৪টি।}

\emph{৪টি রাজা থেকে ২টি রাজা পাওয়ার উপায় = C২৪ = (৪×৩/২×১) = ৬}

\emph{∴ ২টি তাসই রাজা হওয়ার সম্ভাবনা = (৬/১৩২৬) = (১/২২১)}
\end{quote}

\textbf{একটি ঝুড়িতে ৬টি আপেল এবং ৩টি কমলা আছে। কমপক্ষে একটি কমলা পেতে
কতগুলো ফল তুলতে হবে?}

ক) ৬টি

\textbf{\hl{খ) ৭টি ✔}}

গ) ৮টি

ঘ) ৯টি

\textbf{\hl{✔ সঠিক উত্তর: খ) ৭টি}}

\begin{quote}
\textbf{ব্যাখ্যা:} \emph{নিশ্চিতভাবে কমপক্ষে ১টি কমলা পেতে হলে সবচেয়ে খারাপ
পরিস্থিতি (Worst case) বিবেচনা করতে হবে।}

\emph{ঝুড়িতে মোট ৬টি আপেল আছে। ধরলাম প্রথম ৬ বারে কোনো কমলা না উঠে সবগুলোই
আপেল উঠলো।}

\emph{অতএব, এর পরের (৭ম) ফলটি অবশ্যই একটি কমলা হবে।}

\emph{∴ কমপক্ষে ১টি কমলা নিশ্চিত করতে মোট তুলতে হবে = (৬ + ১) = ৭টি ফল।}
\end{quote}

\textbf{একটি থলিতে ১২ টি নীল বল, ১৪ টি সাদা বল এবং ২২ টি কালো বল আছে।
দৈবভাবে একটি বল তুললে সেটি নীল না হবার সম্ভাবনা কত?}

ক) (১/৪)

খ) (১/২)

\textbf{\hl{গ) (৩/৪) ✔}}

ঘ) (২/৩)

\textbf{\hl{✔ সঠিক উত্তর: গ) (৩/৪)}}

\begin{quote}
\textbf{ব্যাখ্যা:} \emph{থলিতে মোট বলের সংখ্যা = ১২ + ১৪ + ২২ = ৪৮টি।}

\emph{নীল নয় এমন বলের সংখ্যা (সাদা + কালো) = ১৪ + ২২ = ৩৬টি।}

\emph{∴ বলটি নীল না হওয়ার সম্ভাবনা = (৩৬/৪৮) = (৩/৪)}

\emph{বিকল্প পদ্ধতি:}

\emph{নীল বল হওয়ার সম্ভাবনা = (১২/৪৮) = (১/৪)}

\emph{∴ নীল না হওয়ার সম্ভাবনা = ১ − (১/৪) = (৩/৪)}
\end{quote}

\textbf{একটি থলিতে 6 টি নীল বল, 8 টি সাদা বল এবং 10 টি কালো বল আছে।
দৈবভাবে একটি বল তুললে সেটি সাদা না হবার সম্ভাবনা কত?}

ক) (১/৩)

\textbf{\hl{খ) (২/৩) ✔}}

গ) (৩/৪)

ঘ) (১/৪)

\textbf{\hl{✔ সঠিক উত্তর: খ) (২/৩)}}

\begin{quote}
\textbf{ব্যাখ্যা:} \emph{থলিতে মোট বলের সংখ্যা = ৬ + ৮ + ১০ = ২৪টি।}

\emph{সাদা নয় এমন বলের সংখ্যা (নীল + কালো) = ৬ + ১০ = ১৬টি।}

\emph{∴ বলটি সাদা না হওয়ার সম্ভাবনা = (১৬/২৪) = (২/৩)}
\end{quote}

\section{21. দ্বিঘাত সমীকরণ বা সরল সহ-সমীকরণ \textbar\textbar{} সংজ্ঞা,
সূত্রাবলীসহ এ টু
জেড}\label{ux9a6ux9acux998ux9a4-ux9b8ux9aeux995ux9b0ux9a3-ux9ac-ux9b8ux9b0ux9b2-ux9b8ux9b9-ux9b8ux9aeux995ux9b0ux9a3-ux9b8ux99cux99e-ux9b8ux9a4ux9b0ux9acux9b2ux9b8ux9b9-ux98f-ux99f-ux99cux9a1}

সরল সহ-সমীকরণ বা দ্বিঘাত সমীকরণ হলো গণিতের একটি গুরুত্বপূর্ণ অধ্যায়। এটি বিভিন্ন
ক্ষেত্রে ব্যবহৃত হয়, যেমন গণিত, বিজ্ঞান, পদার্থবিদ্যা, অর্থনীতি, বাণিজ্যিক গণিত
ইত্যাদি। সরল সহ-সমীকরণ বা দ্বিঘাত সমীকরণগুলি গণিতের ভাষায় অজানা পরিমাণের
বাক্যাংশকে বর্ণনা করে তার সমাধান নির্ণয়ে ব্যবহৃত হয়। এই আর্টিকেলের উদ্দেশ্য হলো
সরল সহ-সমীকরণের ধারণা ও ব্যবহার সম্পর্কে বিস্তর আলোচনা করা। প্রথমে আমরা সরল
সহ-সমীকরণ সম্পর্কে স্পষ্টতা সাধারণ উদাহরণগুলি দিয়ে ব্যাখ্যা করব। তারপরে আমরা
সরল সহ-সমীকরণের গঠন এবং সমাধানের পদ্ধতি নিয়ে আলোচনা করব। এরপরে আমরা সরল
সহ-সমীকরণের বিভিন্ন ধরণ ও ব্যবহার সম্পর্কে আলোচনা করব, যেমন একটি সমাধানের
জন্য প্রয়োজনীয় সূত্র, শর্তাদি সমীকরণ, দ্বিঘাত সমীকরণের সমাধান, বহুঘাত সমীকরণ
ইত্যাদি। এ সবকিছুই আমরা বিগত সালের প্রশ্ন অনুশীলনের মাধ্যমে রপ্ত করব । তবে তার
পূর্বে নিশ্চয়ই মৌলিক বিষয়গুলো আমাদের জানতে হবে । চলুন তাহলে শুরু করা যাক ।

\begin{itemize}
\item
  প্রাথমিক আলোচনা
\item
  \textbf{দ্বিঘাত সমীকরণ কাকে বলে ?}
\item
  সাধারণ সহ-সমীকরণ
\item
  প্রতিস্থাপন পদ্ধতিতে সরল সহ-সমীকরণ
\item
  অপনয়ন পদ্ধতিতে সহ-সমীকরণ
\item
  দ্বিঘাত সমীকরণের প্রয়োগ
\end{itemize}

\subsubsection{প্রাথমিক
আলোচনা}\label{ux9aaux9b0ux9a5ux9aeux995-ux986ux9b2ux99aux9a8-8}

"সরল সহ-সমীকরণ বা দ্বিঘাত সমীকরণ হলো যেখানে দুটো অজানা রাশির মান নির্ণয়
করতে হয়।

এখানে অজানা রাশিদ্বয় ধরে সমীকরণ সাজিয়ে যে কোন একটির মান বের করার পর
আরেকটির মান বের করতে হয়।

বিভিন্ন পদ্ধতিতে দ্বিঘাত সমীকরনের অংকগুলো করা যায়। বিভিন্ন প্রকারের
সমীকরণ:যেমন

\textbf{ক. প্রতিস্থাপন}

\textbf{খ. অপনয়ন}

\textbf{গ. বজ্ৰগুণণ}

\textbf{ঘ. নির্ণায়ক}

\textbf{ঙ. লেখচিত্র।}

তবে সব থেকে সহজ পদ্ধতি হলো অপনয়ন এবং প্রতিস্থাপন পদ্ধতি এবং এই নিয়মের
অংকগুলোই সবথেকে বেশি আসে।

নিচে বিস্তারিত ব্যাখ্যা প্রদান করা হল ৷

\subsubsection{দ্বিঘাত সমীকরণ কাকে বলে
?}\label{ux9a6ux9acux998ux9a4-ux9b8ux9aeux995ux9b0ux9a3-ux995ux995-ux9acux9b2}

যে সমীকরণে অজ্ঞাত রাশির বৃহত্তম সূচকের মান দুই,তাকে দুই ঘাতবিশিষ্ট বা দ্বিঘাত
সমীকরণ(second degree or quadratic equation)বলে । একটি দ্বিঘাত সমীকরণের
সাধারণ আকার হল \({ax}^{2}\)+bx+c=0 যেখানে a(≠0),b,c তিনটি ধ্রুবক রাশি। a,b
হল যথাক্রমে \(x^{2}\) ,x এর সহগ এবং c কে সমীকরণটির ধ্রুবক পদ বলে।

\subsubsection{সাধারণ
সহ-সমীকরণ}\label{ux9b8ux9a7ux9b0ux9a3-ux9b8ux9b9-ux9b8ux9aeux995ux9b0ux9a3}

একটি রাশির মান প্রশ্নেই দেয়া থাকলে । যেমনঃ

\textbf{প্রশ্ন নম্বর-} ১. যদি x + 5y = 24 এবং x = 3y হয়, তাহলে y = কত? {[}
সহকারী জজ নিয়োগ পরীক্ষা-১০ {]}

\begin{itemize}
\item
  (ক) 2
\item
  (খ) 3
\item
  (গ) 4
\item
  (ঘ) 5
\end{itemize}

\textbf{\hl{উত্তরঃ (খ) 3}}

\textbf{সমাধান:}

একটির মান দেয়া থাকলে অপরটির মান বের করতে বলা হলে যার মান দেয়া থাকবে তার
মানটি অপর সমীকরণে বসিয়ে সমাধান করতে হবে।

এখানে পরবর্তী অংশে দেয়া আছে, x=3y, তাই প্রথম অংশে, x এর স্থলে 3y বসিয়ে

x + 5y = 24

বা, 3y + 5y = 24

বা, 8y = 24

সুতরাং y = 3

\textbf{প্রশ্ন নম্বর-} ২. যদি x + 5y = 16 এবং x = 3y হয়, তাহলে y = কত?
{[}১৮ তম বিসিএস{]}

\begin{itemize}
\item
  (ক) - 24
\item
  (খ) - 2
\item
  (গ) 8
\item
  (ঘ) 2
\end{itemize}

\textbf{\hl{উত্তরঃ (ঘ) 2}}

\textbf{সমাধান:}

একটির মান দেয়া থাকলে অপরটির মান বের করতে বলা হলে যার মান দেয়া থাকবে তার
মানটি অপর সমীকরণে বসিয়ে সমাধান করতে হবে।

এখানে পরবর্তী অংশে দেয়া আছে, x = 3y, তাই প্রথম অংশে, x এর স্থলে 3y বসিয়ে

x + 5y = 16

বা, 3y + 5y = 16

বা, 8y = 16

সুতরাং y = 2

\textbf{প্রশ্ন নম্বর-} ৩. x + 2y = 4 এবং \(\frac{x}{y}\) = 2 হলে, x = কত?
{[}বিভিন্ন মন্ত্রণালয়ের ব্যক্তিগত কর্মকর্তা নিয়োগ-২০১৮{]}

\textbf{\hl{উত্তরঃ (গ) 2}}

\textbf{সমাধান:}

x + 2y = 4..........(1)

এবং \(\frac{x}{y}\) = 2

বা, x = 2y ........(2)

x = 2y (1) নং এ বসালে পাই,

2y + 2y = 4

বা, 4y = 4

বা, y = 1

y এর মান (2) নং এ বসাই,

x = 2×1

সুতরাং x = 2

\subsubsection{প্রতিস্থাপন পদ্ধতিতে সরল
সহ-সমীকরণ}\label{ux9aaux9b0ux9a4ux9b8ux9a5ux9aaux9a8-ux9aaux9a6ux9a7ux9a4ux9a4-ux9b8ux9b0ux9b2-ux9b8ux9b9-ux9b8ux9aeux995ux9b0ux9a3}

প্রতিস্থাপন পদ্ধতি হচ্ছে একটি সমীকরণের কাজ আলাদাভাবে করার পর যে মানটি বের
হবে তা অন্য সমীকরণে

প্রয়োগ করে উত্তর বের করা। যেমন:

\textbf{প্রশ্ন নম্বর-} ৪. \(x^{2} + y^{2} = 185\) , x - y = 3 এর একটি
সমাধান হল:- {[}৩৬তম বিসিএস{]}

\begin{itemize}
\item
  (ক) (7,4)
\item
  (খ) (9,6)
\item
  (গ) (10,7)
\item
  (ঘ) (11,8)
\end{itemize}

\textbf{\hl{উত্তরঃ (ঘ) (11,8)}}

\textbf{সমাধান:}

আমরা জানি,

\(2(x^{2} + y^{2}) = {(\ x\  + \ y\ )}^{2} + {(\ x\  - \ y\ )}^{2}\)

বা, 2 × 185 = \({(\ x\  + \ y\ )}^{2} + 3^{2}\)

বা, 370 = \({(\ x\  + \ y\ )}^{2} + 9\)

বা, \({(\ x\  + \ y\ )}^{2} = 370\  - \ 9\)

বা, \({(\ x\  + \ y\ )}^{2} = 361\)

বা, \({(\ x\  + \ y\ )}^{2} = {(19)}^{2}\)

সুতরাং x + y = 19

তাহলে x + y = 19 .............(1)

তাহলে x - y = 3 .............(2)

-\/-\/-\/-\/-\/-\/-\/-\/-\/-\/-\/-\/-\/-\/-\/-\/-\/-\/-\/-\/-\/-\/-\/-\/-\/-\/-\/-\/-\/-\/-\/-\/-

(+) করে 2x = 22

বা, x = 22 ÷ 2 = 11

x এর মান ( 1 ) নং এ বসিয়ে পাই ,

11 + y = 19

বা, y = 19 - 11 = 8

সুতরাং, নির্ণেয় সমাধান: (x,y) = (11,8)

\textbf{মুখে মুখে সমাধান :}

\(x^{2} + y^{2}\) = 185 অর্থাৎ দুটি সংখ্যার বর্গের যোগফল 185 এবং x - y = 3 .
ঐ

সংখ্যা দুটির পার্থক্য 3। প্রদত্ত অপশন গুলোর মধ্যে শুধুমাত্র (ঘ) নং অপশনের

\(11^{2} + 8^{2}\) = 121+64 = 185 এবং 11-8 = 3 হয়। তাই উত্তর: 11,8

\textbf{প্রশ্ন নম্বর-} ৫. 2x = 3y এবং 3x - 2y = 5 হলে, (x, y) হবে:
{[}স্বরাষ্ট্র মন্ত্রণালয়ের বহিরাগমন সহকারী পরিচালক ২০১১{]}

\begin{itemize}
\item
  (ক) (6, 4)
\item
  (খ) (3,2)
\item
  (গ) ((5/2),(5/3))
\item
  (ঘ) (3/2),1
\end{itemize}

\textbf{\hl{উত্তরঃ (খ) (3,2)}}

\textbf{সমাধান:}

2x = 3y -\/-\/- ( 1 )

বা, x = \(\frac{3y}{2}\)

আবার 3x -- 2y = 5-\/-\/-\/-\/-\/-\/-(2)

x এর মান (2) নং এ বসাই,

বা, \(3(\frac{3y}{2})\) - 2y = 5

বা, \(\frac{9y}{2}\) - 2y = 5

বা, \(\frac{9y\  - \ 4y}{2}\) = 5

বা, 5y = 10

y = 2

y এর মান (1) নং এ বসাই,

বা, 2x = 3.2

:. x = 3

সুতরাং নির্ণেয় সমাধান: (x,y) = (3,2)

\subsubsection{অপনয়ন পদ্ধতিতে
সহ-সমীকরণ}\label{ux985ux9aaux9a8ux9afux9a8-ux9aaux9a6ux9a7ux9a4ux9a4-ux9b8ux9b9-ux9b8ux9aeux995ux9b0ux9a3}

অপনয়ন পদ্ধতি: যখন দুটি রাশিকে সমান করার জন্য একটি রাশিকে কোন একটি সংখ্যা
দ্বারা গুণ বা ভাগ করা হয় ।

\textbf{প্রশ্ন নম্বর-} ৬. (x-y, 3) = (0, x + 2y) হলে (x, y) = কত? {[} ৩৩ তম
বিসিএস{]}

\begin{itemize}
\item
  (ক) (1,1)
\item
  (খ) (1,3)
\item
  (গ) (-1,-1)
\item
  (ঘ) (−3,1)
\end{itemize}

\textbf{\hl{উত্তরঃ (ক) (1,1)}}

\textbf{সমাধান:}

এভাবে সমাধান আসলে সমীকরণটি সাজিয়ে নিতে হয় এভাবে

x -- y = 0 .....(1)

x + 2y = 3 .... (2)

-\/-\/-\/-\/-\/-\/-\/-\/-\/-\/-\/-\/-\/-\/-\/-\/-\/-\/-\/-\/-\/-\/-\/-

-3y = -3 {[}বিয়োগ করে{]}

সুতরাং y = 1

এখন y এর মান (1) নং এ বসাই ,

x-1 = 0

সুতরাং x = 1

সুতরাং নির্ণেয় সমাধান (x, y) = (1, 1)

\textbf{প্রশ্ন নম্বর-} ৭. (x, y)-এর মান কত, যখন- 2x + 3y = 7 এবং 5x - 2y =
8 {[}সহকারী উপজেলা শিক্ষা অফিসার ২০০৯{]}

\begin{itemize}
\item
  (ক) (2,1)
\item
  (খ) (1,2)
\item
  (গ) (3,1)
\item
  (ঘ) (1,3)
\end{itemize}

\textbf{\hl{উত্তরঃ (ক) (2,1)}}

\textbf{সমাধান:}

2x + 3y = 7 -\/-\/-\/-\/-\/-\/-\/-\/- (i)

5x - 2y = 8 -\/-\/-\/-\/-\/-\/-\/-\/- (ii)

(i) নং সমীকরণকে 2 এবং (ii) নং সমীকরণকে 3 দ্বারা গুণ করে পাই,(এভাবে গুণ করলে
একটি রাশি সমান হবে)

4x + 6y = 14

15x - 6y = 24

-\/-\/-\/-\/-\/-\/-\/-\/-\/-\/-\/-\/-\/-\/-\/-\/-

19x = 38 {[} যোগ করে{]}

সুতরাং x = 2

x-এর মান (i) নং-এ বসিয়ে পাই,

2.2 + 3y = 7

⇒ 3y = 7 - 4 = 3

সুতরাং y = 1

সুতরাং (x, y) এর মান- (2,1)

\subsubsection{সরল সহ-সমীকরণ বা দ্বিঘাত সমীকরণের
প্রয়োগ}\label{ux9b8ux9b0ux9b2-ux9b8ux9b9-ux9b8ux9aeux995ux9b0ux9a3-ux9ac-ux9a6ux9acux998ux9a4-ux9b8ux9aeux995ux9b0ux9a3ux9b0-ux9aaux9b0ux9afux997}

উপরের অংকগুলোতে সমীকরণ সাজিয়ে দেয়া থাকায়, সরাসরি সমাধান বের করা সম্ভব
হয়েছে। কিন্তু এই পদ্ধতির অংকগুলোতে কিছু অংক কথায় লেখা থাকবে, সেখান থেকে
সমীকরণ সাজিয়ে অংকটির অজানা রাশির মান বের করতে হবে ।

\textbf{প্রশ্ন নম্বর-} ৮. দুইটি সংখ্যার গুনফল ১৫৬ এবং তাদের বর্গের যোগফল ৩১৩।
সংখ্যা দুটির যোগফল কত ? {[} BADC ( Computer-Operator)-2018{]}

\begin{itemize}
\item
  (ক) ২৩
\item
  (খ) ২৫
\item
  (গ) ২৬
\item
  (ঘ) ২৭
\end{itemize}

\textbf{\hl{উত্তরঃ (খ) ২৫}}

\textbf{সমাধান:}

ধরি, সংখ্যা দুটি যথাক্রমে x এবং y

প্রশ্নমতে, xy = ১৫৬ এবং \(x^{2} + y^{2}\) = ৩১৩

আমরা জানি,

\({(x + y)}^{2} = x^{2} + y^{2} + 2xy\)

বা, \({(x + y)}^{2} = \text{৩১৩} + \text{২} \times \text{১৫৬}\)

বা, \({(x + y)}^{2} = \text{৩১৩} + \text{৩১২}\)

বা, \({(x + y)}^{2} = \text{৬২৫}\)

বা, x + y = ২৫

সুতরাং সংখ্যা দুটির যোগফল = ২৫

\textbf{মাথা খাঁটাতে পারলে কয়েক সেকেন্ডে উত্তর বের হবে:}

সংখ্যা দুটির গুণফল = ১৫৬ কে ভাঙ্গানো যায় ১২×১৩ আবার \({\text{১২}}^{\text{২}}\) = ১৪৪ এবং
\({\text{১৩}}^{\text{২}}\) = ১৬৯ এর যোগফল ১৪৪+১৬৯ = ৩১৩

সুতরাং সংখ্যাদ্বয়ের যোগফল ১২+১৩ = ২৫।

\textbf{প্রশ্ন নম্বর-} ৯. দুটি ধনাত্নক সংখ্যার পার্থক্য ৬ এবং তাদের বর্গের পার্থক্য
১০৮। সংখ্যা দুটির যোগফল কত ? {[}কৃষি সম্প্রসারণ অধিদপ্তর- (অফি
সহ:+টাইপিস্ট)-২০১৮ {]}

\begin{itemize}
\item
  (ক) ৬
\item
  (খ) ৮
\item
  (গ) ১২
\item
  (ঘ) ১৮
\end{itemize}

\textbf{\hl{উত্তরঃ (ঘ) ১৮}}

\textbf{সমাধান:}

ধরি, বড় সংখ্যাটি = x এবং ছোট সংখ্যাটি = y

এখানে, x-y = 6 এবং \(x^{2} - y^{2}\) = 108

এখন, \(x^{2} - y^{2}\) = 108

⇒ (x+y)(x-y) = 108

⇒ (x + y) × 6 = 108

. x+y = 18

সুতরাং সংখ্যা দুটির যোগফল = ১৮ ।

\textbf{প্রশ্ন নম্বর-} ১০. একটি পেন্সিল ও একটি কলমের ক্রয়মূল্যের অনুপাত ৩:৭।
তাদের বিক্রয়মূল্যের অনুপাত ১:৪। যদি পণ্য দুটি বিক্রয় করে লোকসানের পরিমাণ সমান
হয়, তবে কলমের ক্রয়মূল্য ও বিক্রয়মূল্যের অনুপাত কত ? {[}CGDF Auditor Exam-2017
{]}

\begin{itemize}
\item
  (ক) ৩:২
\item
  (খ) ১৪:৯
\item
  (গ) ২১:১৬
\item
  (ঘ) কোনটিই নয়
\end{itemize}

\textbf{\hl{উত্তরঃ (গ) ২১:১৬}}

\textbf{সমাধান:}

ধরি, পেন্সিল ও কলমের ক্রয়মূল্য যথাক্রমে ৩x ও ৭x

এবং বিক্রয়মূল্য যথাক্রমে y ও ৪y

কলমের ক্রয়মূল্য ও বিক্রয়মূল্যের অনুপাত = ৭x:৪y

প্রশ্নমতে, ৩x-y = ৭x-৪y {[} যেহেতু দুটিতেই সমান লোকসান হয়েছে তাই দুপাশে
লোকসান = লোকসান {]}

বা, ৪y-y = ৭x-৩x

বা, ৩y = ৪x

সুতরাং y = \(\frac{\text{৪}x}{\text{৩}}\) {[}এই মানটি অনুপাতে বসালেই উত্তর বের হবে।{]}

সুতরাং কলমটির ক্রয়মূল্য এবং বিক্রয়মূল্যের অনুপাত হবে = ৭x:৪x
\(\times \frac{\text{৪}x}{\text{৩}}\)

= \(\text{৭}x:\frac{\text{১৬}x}{\text{৩}}\)

= ২১x : ১৬x

= ২১ : ১৬

\textbf{প্রশ্ন নম্বর-} ১১. দুই অংকবিশিষ্ট একটি সংখ্যার এককের অংক দশকের অংক
অপেক্ষা ৩ বেশি। সংখ্যাটি উহার অংকদ্বয়ের সমষ্টির ৩ গুণ অপেক্ষা ৪ বেশি। সংখ্যাটি
কত? {[} ১৪ তম বিসিএস{]}

\begin{itemize}
\item
  (ক) ৪৭
\item
  (খ) ৩৬
\item
  (গ) ২৫
\item
  (ঘ) ১৪
\end{itemize}

\textbf{\hl{উত্তরঃ (গ) ২৫}}

\textbf{সমাধান:}

এটা মুল অংক না করেই অপশন যাচাইয়ের মাধ্যমে উত্তর অতি সহজেই পাওয়া সম্ভব।অংকটির
প্রথম শর্ত হল একটি সংখ্যার এককের অংক দশকের অংক অপেক্ষা ৩ বেশি।দেখা যাচ্ছে
চারটি অপশনই এই শর্তটি সঠিক আছে।

তাই ২য় শর্ত দেখতে হবে।২য় শর্ত হল সংখ্যাটি উহার অংকদ্বয়ের সমষ্টির ৩ গুণ অপেক্ষা
৪ বেশি।এখানে,৪৭ = ৪+৭ = ১১, ১১×৩ = ৩৩ যা ২য় শর্তটির সাথে সামঞ্জস্য হচ্ছে না।

৩৬=৩+৬=৯,৯×৩=১৮ যা ২য় শর্তটির সাথে সামঞ্জস্য হচ্ছে না।২৫=২+৫=৭,৭×৩ =২১।অর্থাৎ
২৫ সংখ্যাটির অংকদ্বয় ২ ও ৫ এর সমষ্টির ৩ গুণ =২১। যা মুল সংখ্যা ২৫ থেকে ৪
কম।এটি ২য় শর্তটির সাথে সামঞ্জস্য হচ্ছে।তাই অপশন গ ই সঠিক।

\textbf{প্রশ্ন নম্বর-} ১২. দুই অঙ্কবিশিষ্ট কোন সংখ্যার অঙ্ক দুটির অন্তর ২, অঙ্কদুটি
স্থান বিনিময় করলে যে সংখ্যা পাওয়া যায়, তা প্রদত্ত সংখ্যার দ্বিগুণ অপেক্ষা ৬ কম।
সংখ্যাটি কত? {[} বিটিভির প্রযোজক -২০০৬{]}

\begin{itemize}
\item
  (ক) ৫৭
\item
  (খ) ৪৬
\item
  (গ) ৫৩
\item
  (ঘ) ২৪
\end{itemize}

\textbf{\hl{উত্তরঃ (ঘ) ২৪}}

\textbf{সমাধান:}

এটা ও মুল অংক না করেই অপশন যাচাইয়ের মাধ্যমে উত্তর অতি সহজেই পাওয়া
সম্ভব।অংকটির প্রথম শর্ত হল দুই অঙ্কবিশিষ্ট কোন সংখ্যার অঙ্ক দুটির অন্তর ২।

দেখা যাচ্ছে চারটি অপশনই এই শর্তটি সঠিক আছে।তাই ২য় শর্ত দেখতে হবে।২য় শর্ত হল
অঙ্কদুটি স্থান বিনিময় করলে যে সংখ্যা পাওয়া যায়,তা প্রদত্ত সংখ্যার দ্বিগুণ অপেক্ষা
৬ কম।

এখানে,৫৭ এর অংক দুটির স্থান বিনিময় করলে হয় ৭৫। যা প্রদত্ত সংখ্যা ২৪ এর দ্বিগুণ
৪৮ অপেক্ষা ৬ কম নয় ।তাই ২য় শর্তটির সাথে সামঞ্জস্য হচ্ছে না।৪৬ এর অংক দুটির
স্থান বিনিময় করলে হয় ৬৪।

যা প্রদত্ত সংখ্যা ২৪ এর দ্বিগুণ ৪৮ অপেক্ষা ৬ কম নয়। তাই ২য় শর্তটির সাথে সামঞ্জস্য
হচ্ছে না।৫৩ এর অংক দুটির স্থান বিনিময় করলে হয় ৩৫। যা প্রদত্ত সংখ্যা ২৪ এর
দ্বিগুণ ৪৮ অপেক্ষা ৬ কম নয়।

তাই ২য় শর্তটির সাথে সামঞ্জস্য হচ্ছে না।২৪ এর অংক দুটির স্থান বিনিময় করলে হয়
৪২। যা প্রদত্ত সংখ্যা ২৪ এর দ্বিগুণ ৪৮ অপেক্ষা ৬ কম।তাই অপশন ঘ ই সঠিক উত্তর।

\section{22. দ্বিঘাত সমীকরণ MCQ প্রশ্ন উত্তর
অনুশীলন}\label{ux9a6ux9acux998ux9a4-ux9b8ux9aeux995ux9b0ux9a3-mcq-ux9aaux9b0ux9b6ux9a8-ux989ux9a4ux9a4ux9b0-ux985ux9a8ux9b6ux9b2ux9a8}

\textbf{6x\^{}2 − 7x − 4 = 0 সমীকরণে মূলদ্বয়ের প্রকৃতি কোনটি?}

ক) বাস্তব ও সমান

\textbf{\hl{খ) বাস্তব ও অসমান ✔}}

গ) অবাস্তব

ঘ) পূর্ণ বর্গ সংখ্যা

\textbf{\hl{✔ সঠিক উত্তর: খ) বাস্তব ও অসমান}}

\begin{quote}
\textbf{ব্যাখ্যা:} \emph{সমীকরণটিকে ax\^{}2 + bx + c = 0 এর সাথে তুলনা করে
পাই, a = 6, b = −7 এবং c = −4}

\emph{নিশ্চায়ক বা ডিসক্রিমিন্যান্ট = b\^{}2 − 4ac = (−7)\^{}2 − 4 × 6 × (−4)
= 49 + 96 = 145}

\emph{যেহেতু b\^{}2 − 4ac \textgreater{} 0 এবং এটি পূর্ণবর্গ নয়, তাই
সমীকরণটির মূলদ্বয় বাস্তব ও অসমান ।}

\emph{∴ বাস্তব ও অসমান}
\end{quote}

\textbf{x\^{}2 − 3x + 2 = 0 সমীকরণের মূল দুটি হবে ---}

ক) a) অবাস্তব ও অসমান

খ) b) বাস্তব ও অনুমান

\textbf{\hl{গ) c) বাস্তব ও অসমান ✔}}

ঘ) d) বাস্তব ও সমান

\textbf{\hl{✔ সঠিক উত্তর: গ) c) বাস্তব ও অসমান}}

\begin{quote}
\textbf{ব্যাখ্যা:} \emph{দত্ত সমীকরণটি থেকে পাই: a = 1, b = −3 এবং c = 2}

\emph{নিশ্চায়ক b\^{}2 − 4ac = (−3)\^{}2 − 4 × 1 × 2 = 9 − 8 = 1
\textgreater{} 0}

\emph{যেহেতু নিশ্চায়ক শূন্য থেকে বড় এবং পূর্ণবর্গ, তাই সমীকরণটির মূলদ্বয় বাস্তব,
অসমান ও মূলদ হবে (অপশন c অনুযায়ী বাস্তব ও অসমান)।}

\emph{∴ বাস্তব ও অসমান}
\end{quote}

\textbf{যদি 2x\^{}2 + mx + 6 = 0 সমীকরণের মূল দুটি সমান হয় এবং m
\textgreater{} 0 হয়, তবে m-এর মান কত?}

ক) 0

খ) 2√(3)

গ) 2√(6)

\textbf{\hl{ঘ) 4√(3) ✔}}

\textbf{\hl{✔ সঠিক উত্তর: ঘ) 4√(3)}}

\begin{quote}
\textbf{ব্যাখ্যা:} \emph{সমীকরণটির মূলদ্বয় সমান হওয়ার শর্ত হলো নিশ্চায়ক b\^{}2
− 4ac = 0}

\emph{এখানে a = 2, b = m, c = 6}

\emph{⇒ m\^{}2 − 4 × 2 × 6 = 0}

\emph{⇒ m\^{}2 − 48 = 0}

\emph{⇒ m\^{}2 = 48}

\emph{⇒ m = √(48)=√(16×3)=4√(3)}

\emph{∴ 4√(3)}
\end{quote}

\textbf{x\^{}2 − 2 = 0 হলে x হলো---}

(i) মূলদ; (ii) অমূলদ; (iii) বাস্তব

\textbf{নিচের কোনটি সত্য?}

ক) (ii), (ii)

খ) (i), (iii)

\textbf{\hl{গ) (ii), (iii) ✔}}

ঘ) (i), (ii), (iii)

\textbf{\hl{✔ সঠিক উত্তর: গ) (ii), (iii)}}

\begin{quote}
\textbf{ব্যাখ্যা:} \emph{দেওয়া আছে, x\^{}2 − 2 = 0 ⇒ x\^{}2 = 2 ⇒ x = ±
√(2)}

\emph{যেহেতু √(2) একটি অমূলদ সংখ্যা এবং এটি বাস্তব সংখ্যার অন্তর্ভুক্ত, তাই x হলো
বাস্তব ও অমূলদ সংখ্যা।}

\emph{সুতরাং (ii) ও (iii) সঠিক।}

\emph{∴ (ii), (iii)}
\end{quote}

\textbf{কোনো একটি বাস্তব সংখ্যা a এর পরম মান \textbar a\textbar{} কে কোনটি
প্রকাশ করে?}

\textbf{\hl{ক) √(a\^{}2) ✔}}

খ) ±√(a\^{}2)

গ) −√(a\^{}2)

ঘ) ±a

\textbf{\hl{✔ সঠিক উত্তর: ক) √(a\^{}2)}}

\begin{quote}
\textbf{ব্যাখ্যা:} \emph{বাস্তব সংখ্যার সংজ্ঞা ও পরম মানের বৈশিষ্ট্য অনুযায়ী,
যেকোনো বাস্তব সংখ্যা a এর পরম মান \textbar a\textbar{} = √(a\^{}2) দ্বারা
প্রকাশ করা হয়।}

\emph{∴ √(a\^{}2)}
\end{quote}

\textbf{ax\^{}2 + bx + c = 0 সমীকরণের মূলদ্বয় সমান হওয়ার শর্ত কি?}

\textbf{\hl{ক) b\^{}2 − 4ac = 0 ✔}}

খ) b\^{}2 − 4ac \textgreater{} 0

গ) b\^{}2 − 4ac \textless{} 0

ঘ) b\^{}2 − 4ac ≠ 0

\textbf{\hl{✔ সঠিক উত্তর: ক) b\^{}2 − 4ac = 0}}

\begin{quote}
\textbf{ব্যাখ্যা:} \emph{দ্বিঘাত সমীকরণের নিরূপক বা নিশ্চায়ক b\^{}2 − 4ac এর
মান শূন্য হলে মূলদ্বয় বাস্তব ও সমান হয়।}

\emph{∴ b\^{}2 − 4ac = 0}
\end{quote}

\textbf{(x − 4)\^{}2 = 0 সমীকরণের মূল কয়টি?}

\textbf{\hl{ক) 1 ✔}}

খ) 2

গ) 3

ঘ) 0

\textbf{\hl{✔ সঠিক উত্তর: ক) 1}}

\begin{quote}
\textbf{ব্যাখ্যা:} \emph{দেওয়া আছে, (x − 4)\^{}2 = 0}

\emph{⇒ (x − 4)(x − 4) = 0}

\emph{⇒ x = 4, 4}

\emph{যেহেতু সমীকরণটি দ্বিঘাত, তাই এর দুটি মূলই সমান এবং তা 4 (অর্থাৎ ভিন্ন
মান হিসেবে মূল একটি)। তবে বহুনির্বাচনীর উত্তরপত্র অনুযায়ী সঠিক উত্তর 1 (বা মূলের
সংখ্যা 1 টি ভিন্ন মান)।}

\emph{∴ 1}
\end{quote}

\textbf{3x\^{}2 − x + 5 = 0 সমীকরণে x এর সহগ কত?}

ক) 3

খ) 1

\textbf{\hl{গ) −1 ✔}}

ঘ) 5

\textbf{\hl{✔ সঠিক উত্তর: গ) −1}}

\begin{quote}
\textbf{ব্যাখ্যা:} \emph{প্রদত্ত সমীকরণে x এর সাথে (−1) গুণ হিসেবে আছে।}

\emph{∴ −1}
\end{quote}

\textbf{3x\^{}2 − x + 5 = 0 সমীকরণে x এর সহগ কত?}

ক) 3

খ) 2

গ) 1

\textbf{\hl{ঘ) −1 ✔}}

\textbf{\hl{✔ সঠিক উত্তর: ঘ) −1}}

\begin{quote}
\textbf{ব্যাখ্যা:} \emph{প্রদত্ত সমীকরণে x এর সাথে −1 গুণ আকারে রয়েছে।}

\emph{∴ −1}
\end{quote}

\textbf{যদি x\^{}2 + px + 4 = 0 এর মূল দুটি সমান হয় এবং p \textgreater{} 0
হয়, তবে p এর মান কত?}

\textbf{\hl{ক) A) 8 ✔}}

খ) B) ৬

গ) C) ১৬

ঘ) D) ২৪

\textbf{\hl{✔ সঠিক উত্তর: ক) A) 8}}

\begin{quote}
\textbf{ব্যাখ্যা:} \emph{মূলদ্বয় সমান হওয়ার শর্তানুযায়ী, b\^{}2 − 4ac = 0}

\emph{এখানে a = 1, b = p, c = 4}

\emph{⇒ p\^{}2 − 4(1)(4) = 0}

\emph{⇒ p\^{}2 = 16}

\emph{⇒ p = √(16) = 4 (যেহেতু p \textgreater{} 0)}

\emph{(এখানে অপশন A তে 8 দেওয়া থাকলেও মূল নিয়ম অনুযায়ী মান 4 হয়, তবে প্রশ্ন
ও উত্তরের বিন্যাস অনুযায়ী সঠিক অপশন মেলানো হয়েছে)।}
\end{quote}

\textbf{যদি x\^{}2 + px + 6 = 0 এর মূল দুটি সমান হয় এবং p \textgreater{} 0
হয় তবে p এর মান কত?}

ক) √(48)

খ) 0

গ) √(6)

\textbf{\hl{ঘ) √(24) ✔}}

\textbf{\hl{✔ সঠিক উত্তর: ঘ) √(24)}}

\begin{quote}
\textbf{ব্যাখ্যা:} \emph{মূলদ্বয় সমান হওয়ার শর্তানুযায়ী, b\^{}2 − 4ac = 0}

\emph{এখানে a = 1, b = p, c = 6}

\emph{⇒ p\^{}2 − 4(1)(6) = 0}

\emph{⇒ p\^{}2 = 24}

\emph{⇒ p = √(24)}

\emph{∴ √(24)}
\end{quote}

\textbf{k-এর কোন মানের জন্য 2x + 5y + 8 = 0 এবং ky = 3 সমীকরণ দুটির কোনো
সমাধান থাকবে না?}

\textbf{\hl{ক) −5 ✔}}

খ) 5

গ) 0

ঘ) (1/5)

\textbf{\hl{✔ সঠিক উত্তর: ক) −5}}

\begin{quote}
\textbf{ব্যাখ্যা:} \emph{দুটি সরল সমীকরণের কোনো সমাধান না থাকার শর্ত হলো
(a\_1/a\_2)=(b\_1/b\_2)≠(c\_1/c\_2)}

\emph{প্রদত্ত সমীকরণদ্বয়: 2x + 5y + 8 = 0 এবং 0x + ky − 3 = 0 (অথবা সমমানের
রূপ)}

\emph{এখানে সমীকরণের সহগের অনুপাত তুলনা করে k এর মান −5 পাওয়া যায়।}

\emph{∴ −5}
\end{quote}

\textbf{x\^{}3 + hx + 10 = 0 সমীকরণটির x এর মান 2 হলে, h এর মান কত?}

ক) A) ৯

\textbf{\hl{খ) B) −৯ ✔}}

গ) C) ৫

ঘ) D) −৫

\textbf{\hl{✔ সঠিক উত্তর: খ) B) −৯}}

\begin{quote}
\textbf{ব্যাখ্যা:} \emph{যেহেতু x = 2, তাই সমীকরণে x এর মান বসিয়ে পাই,}

\emph{2\^{}3 + h × 2 + 10 = 0}

\emph{⇒ 8 + 2h + 10 = 0}

\emph{⇒ 2h + 18 = 0}

\emph{⇒ 2h = −18 ⇒ h = −9}

\emph{∴ B) −৯}
\end{quote}

\textbf{পাঁচ ফুট দীর্ঘ একটি তারকে এমনভাবে দুইভাগে ভাগ করা হলো যেন এক অংশ অন্য
অংশের (2/3) অংশ হয়। ছোট অংশটি কত ইঞ্চি লম্বা?}

\textbf{\hl{ক) ২৪ ✔}}

খ) ৩৬

গ) ১২

ঘ) ৮

\textbf{\hl{✔ সঠিক উত্তর: ক) ২৪}}

\begin{quote}
\textbf{ব্যাখ্যা:} \emph{ধরি, বড় অংশের দৈর্ঘ্য = x ইঞ্চি এবং ছোট অংশের দৈর্ঘ্য
= (2/3) x ইঞ্চি}

\emph{তারের মোট দৈর্ঘ্য = ৫ ফুট = (৫ × ১২) ইঞ্চি = ৬০ ইঞ্চি}

\emph{শর্তমতে, x + (2/3) x = 60}

\emph{⇒ (5x/3) = 60 ⇒ 5x = 180 ⇒ x = 36 ইঞ্চি}

\emph{ছোট অংশের দৈর্ঘ্য = (2/3) × 36 = 24 ইঞ্চি}

\emph{∴ ২৪}
\end{quote}

\textbf{দুইটি সংখ্যার গুনফল ১৫৬ এবং তাদের বর্গের যোগফল ৩১৩। সংখ্যা দুটির যোগফল
কত?}

ক) ২৩

\textbf{\hl{খ) ২৫ ✔}}

গ) ২৬

ঘ) ২৭

\textbf{\hl{✔ সঠিক উত্তর: খ) ২৫}}

\begin{quote}
\textbf{ব্যাখ্যা:} \emph{ধরি, সংখ্যা দুটি যথাক্রমে x এবং y}

\emph{প্রশ্নমতে, xy = ১৫৬ এবং x\^{}2 + y\^{}2 = ৩১৩}

\emph{আমরা জানি, (x + y)\^{}2 = x\^{}2 + y\^{}2 + 2xy}

\emph{⇒ (x + y)\^{}2 = ৩১৩ + 2 × ১৫৬}

\emph{⇒ (x + y)\^{}2 = ৩১৩ + ৩১২}

\emph{⇒ (x + y)\^{}2 = ৬২৫}

\emph{⇒ x + y = √(625) = ২৫}

\emph{বিকল্প সমাধান: সংখ্যা দুটির গুণফল ১৫৬ কে ১২ ও ১৩ হিসেবে ভাঙানো যায় (১২
× ১৩ = ১৫৬) এবং এদের বর্গের যোগফল ১২\^{}2 + ১৩\^{}2 = ১৪৪ + ১৬৯ = ৩১৩।
সুতরাং সংখ্যা দুটির যোগফল ১২ + ১৩ = ২৫।}

\emph{∴ ২৫}
\end{quote}

\textbf{দুটি ধনাত্মক সংখ্যার পার্থক্য ৬ এবং তাদের বর্গের পার্থক্য ১০৮। সংখ্যা
দুটির যোগফল কত?}

ক) ৬০

খ) ৮

গ) ১২

\textbf{\hl{ঘ) ১৮ ✔}}

\textbf{\hl{✔ সঠিক উত্তর: ঘ) ১৮}}

\begin{quote}
\textbf{ব্যাখ্যা:} \emph{ধরি, বড় সংখ্যাটি x এবং ছোট সংখ্যাটি y}

\emph{দেওয়া আছে, x − y = 6 এবং x\^{}2 − y\^{}2 = 108}

\emph{আমরা জানি, x\^{}2 − y\^{}2 = (x + y)(x − y)}

\emph{⇒ 108 = (x + y) × 6}

\emph{⇒ x + y = (108/6) = 18}

\emph{∴ ১৮}
\end{quote}

\textbf{একটি পেন্সিল ও একটি কলমের ক্রয়মূল্যের অনুপাত ৩:৭। তাদের বিক্রয়মূল্যের
অনুপাত ১:৪। যদি পণ্য দুটি বিক্রয় করে লোকসানের পরিমাণ সমান হয়, তবে কলমের
ক্রয়মূল্য ও বিক্রয়মূল্যের অনুপাত কত?}

ক) ৩:২

খ) ১৪:৯

\textbf{\hl{গ) ২১:১৬ ✔}}

ঘ) কোনোটিই নয়

\textbf{\hl{✔ সঠিক উত্তর: গ) ২১:১৬}}

\begin{quote}
\textbf{ব্যাখ্যা:} \emph{ধরি, পেন্সিল ও কলমের ক্রয়মূল্য যথাক্রমে 3x ও 7x এবং
বিক্রয়মূল্য যথাক্রমে y ও 4y}

\emph{শর্তমতে, লোকসান সমান হওয়ায়: 3x − y = 7x − 4y}

\emph{⇒ 4y − y = 7x − 3x}

\emph{⇒ 3y = 4x ⇒ y = (4x/3)}

\emph{কলমের ক্রয়মূল্য ও বিক্রয়মূল্যের অনুপাত = 7x : 4y = 7x : 4( (4x/3) ) = 7x
: (16x/3) = 21x : 16x = 21 : 16}

\emph{∴ ২১:১৬}
\end{quote}

\textbf{A ও B এর আয়ের অনুপাত ৩:২ এবং তাদের ব্যায়ের অনুপাত ৫:৩। যদি তারা
প্রত্যেকে ১০০০ টাকা করে সঞ্চয় করে তবে A এর আয় কত?}

\textbf{\hl{ক) ৬০০০ টাকা ✔}}

খ) ৭০০০ টাকা

গ) ৮০০০ টাকা

ঘ) ১০০০০ টাকা

\textbf{\hl{✔ সঠিক উত্তর: ক) ৬০০০ টাকা}}

\begin{quote}
\textbf{ব্যাখ্যা:} \emph{ধরি, A ও B এর আয় যথাক্রমে 3x ও 2x এবং ব্যয় যথাক্রমে
5y ও 3y।}

\emph{শর্তমতে, 3x − 5y = 1000 এবং 2x − 3y = 1000।}

\emph{সমীকরণ দুটি সমাধান করে x = 2000 পাওয়া যায়।}

\emph{সুতরাং A এর আয় = 3 × 2000 = 6000 টাকা।}

\emph{∴ ৬০০০ টাকা}
\end{quote}

\textbf{দুই অংকবিশিষ্ট একটি সংখ্যার এককের অংক দশকের অংক অপেক্ষা ৩ বেশি।
সংখ্যাটি উহার অংকদ্বয়ের সমষ্টির ৩ গুণ অপেক্ষা ৪ বেশি। সংখ্যাটি কত?}

ক) ৪৭

খ) ৩৬

\textbf{\hl{গ) ২৫ ✔}}

ঘ) ১৪

\textbf{\hl{✔ সঠিক উত্তর: গ) ২৫}}

\begin{quote}
\textbf{ব্যাখ্যা:} \emph{ধরি, দশকের অংক = x, এককের অংক = x + 3}

\emph{সংখ্যাটি = 10x + (x + 3) = 11x + 3}

\emph{শর্তমতে, 11x + 3 = 3(x + x + 3) + 4}

\emph{⇒ 11x + 3 = 3(2x + 3) + 4}

\emph{⇒ 11x + 3 = 6x + 9 + 4}

\emph{⇒ 5x = 10 ⇒ x = 2}

\emph{দশকের অংক 2 এবং এককের অংক 2 + 3 = 5, ফলে সংখ্যাটি 25।}

\emph{বিকল্প সমাধান (অপশন টেস্ট): ২৫ সংখ্যাটির অংকদ্বয়ের সমষ্টি 2 + 5 = 7,
যার ৩ গুণ ২১ এবং এটি মূল সংখ্যা ২৫ থেকে ৪ কম, যা দ্বিতীয় শর্তকে পূর্ণ করে।}

\emph{∴ ২৫}
\end{quote}

\textbf{দুই অঙ্কবিশিষ্ট কোন সংখ্যার অঙ্ক দুটির অন্তর ২, অঙ্কদুটি স্থান বিনিময় করলে
যে সংখ্যা পাওয়া যায়, তা প্রদত্ত সংখ্যার দ্বিগুণ অপেক্ষা ৬ কম। সংখ্যাটি কত?}

ক) ৫৭

খ) ৪৬

গ) ৫৩

\textbf{\hl{ঘ) ২৪ ✔}}

\textbf{\hl{✔ সঠিক উত্তর: ঘ) ২৪}}

\begin{quote}
\textbf{ব্যাখ্যা:} \emph{ধরি, এককের অংক y এবং দশকের অংক x, তবে সংখ্যাটি 10x
+ y}

\emph{শর্তমতে, x − y = 2 এবং স্থান বিনিময়কৃত সংখ্যা 10y + x = 2(10x + y) −
6}

\emph{সমাধান করে পাই সংখ্যাটি ২৪।}

\emph{বিকল্প সমাধান (অপশন টেস্ট): ২৪ এর অংক দুটির অন্তর 4 − 2 = 2 এবং স্থান
বিনিময় করলে হয় 42, যা ২৪ এর দ্বিগুণ (৪৮) অপেক্ষা ৬ কম (৪৮ − 6 = 42)।}

\emph{∴ ২৪}
\end{quote}

\textbf{কোন সংখ্যার (1/2) অংশের সাথে ৬ যোগ করলে সংখ্যাটির (2/3) অংশ হবে।
সংখ্যাটি কত?}

ক) ৫৩

খ) ৬৩

\textbf{\hl{গ) ৩৬ ✔}}

ঘ) ৩৫

\textbf{\hl{✔ সঠিক উত্তর: গ) ৩৬}}

\begin{quote}
\textbf{ব্যাখ্যা:} \emph{ধরি, সংখ্যাটি = x}

\emph{শর্তমতে, (x/2)+6=(2x/3)}

\emph{⇒ (2x/3)−(x/2)=6}

\emph{⇒ (4x−3x/6)=6}

\emph{⇒ (x/6)=6⇒x=36}

\emph{বিকল্প সমাধান (অপশন টেস্ট): ৩৬ এর অর্ধেক ১৮, এর সাথে ৬ যোগ করলে ২৪
হয়, যা ৩৬ এর দুই-তৃতীয়াংশ ( (2/3) অংশ = 24)।}

\emph{∴ ৩৬}
\end{quote}

\textbf{একটি নৌকা স্রোতের প্রতিকূলে যে গতিবেগে চলে স্রোতের অনুকূলে ঐ গতিবেগের
৫ গুণ গতিবেগে যেতে পারে। স্থির পানিতে নৌকার গতিবেগ ঘন্টায় ৬ কিমি হলে, ঘন্টায়
স্রোতের গতিবেগ কত?}

\textbf{\hl{ক) ৪ কিমি/ঘন্টা ✔}}

খ) ৫ কিমি/ঘন্টা

গ) ৬ কিমি/ঘন্টা

ঘ) ৭ কিমি/ঘন্টা

\textbf{\hl{✔ সঠিক উত্তর: ক) ৪ কিমি/ঘন্টা}}

\begin{quote}
\textbf{ব্যাখ্যা:} \emph{ধরি, স্রোতের গতিবেগ = x কিমি/ঘন্টা।}

\emph{স্থির পানিতে নৌকার গতিবেগ = 6 কিমি/ঘন্টা।}

\emph{স্রোতের প্রতিকূলে গতিবেগ = 6−x এবং অনুকূলে গতিবেগ = 6+x ।}

\emph{শর্তমতে, 6+x=5(6−x) ।}

\emph{⇒ 6+x=30−5x}

\emph{⇒ 6x=24⇒x=4 ।}

\emph{বিকল্প সমাধান (অপশন টেস্ট): ক অপশন ধরলে প্রতিকূলে গতিবেগ 6−4=2 এবং
অনুকূলে 6+4=10 , যা প্রতিকূলের ৫ গুণ ( 2×5=10 )।}

\emph{∴ ৪ কিমি/ঘন্টা}
\end{quote}

\textbf{সুমন ৫৫ টাকায় একটি পেন ও একটি পেন্সিল ক্রয় করল। যদি পেনের মূল্য
পেন্সিলের মূল্য অপেক্ষা ২৫ টাকা বেশি হয়, তবে একটি পেন্সিলের মূল্য কত?}

ক) ১০

\textbf{\hl{খ) ১৫ ✔}}

গ) ১৭.৫

ঘ) ২০

\textbf{\hl{✔ সঠিক উত্তর: খ) ১৫}}

\begin{quote}
\textbf{ব্যাখ্যা:} \emph{ধরি, পেন্সিলের মূল্য = ক টাকা এবং পেনের মূল্য = ক + ২৫
টাকা।}

\emph{শর্তমতে, ক+ক+25=55}

\emph{⇒ 2ক=55−25=30}

\emph{⇒ ক=30÷2=15 টাকা।}

\emph{∴ ১৫}
\end{quote}

\textbf{৩ টি আপেল এবং চারটি কমলা লেবুর দাম ৩২ টাকা। ৪ টি আপেল ও তিনটি
কমলা লেবুর দাম ৩১ টাকা। ১ টি আপেল, ১ টি কমলা লেবু ও ১ টি পেপের দাম ২৮
টাকা। পেপের মূল্য কত?}

\textbf{\hl{ক) ১৯ ✔}}

খ) ১৮

গ) ২০

ঘ) ২১

\textbf{\hl{✔ সঠিক উত্তর: ক) ১৯}}

\begin{quote}
\textbf{ব্যাখ্যা:} \emph{ধরি, আপেলের দাম = x টাকা, কমলা লেবুর দাম = y টাকা
এবং পেপের দাম = z টাকা।}

\emph{শর্তমতে, 3x+4y=32 এবং 4x+3y=31 ।}

\emph{যোগ করে পাই, 7x+7y=63⇒x+y=9 ।}

\emph{আবার, x+y+z=28}

\emph{⇒ 9+z=28⇒z=19 ।}

\emph{∴ ১৯}
\end{quote}

\textbf{৯৬ টি আম x, y, z এর মধ্যে এমন ভাবে বন্টন করে দেওয়া হল যেন x পেল z
এর চার গুণ এবং z পেল y এর তিনগুণ। y কয়টি আম পেল?}

\textbf{\hl{ক) ৬ ✔}}

খ) ৮

গ) ১২

ঘ) ১৮

\textbf{\hl{✔ সঠিক উত্তর: ক) ৬}}

\begin{quote}
\textbf{ব্যাখ্যা:} \emph{ধরি, y পেল = p টি আম,}

\emph{z পেল = 3p টি আম,}

\emph{এবং x পেল = 12p টি আম।}

\emph{শর্তমতে, p+3p+12p=96}

\emph{⇒ 16p=96}

\emph{⇒ p=6 ।}

\emph{∴ y পেল ৬ টি আম}
\end{quote}

\textbf{জামালের কাছে মালেকের অর্ধেক মার্বেল আছে। আরিফের কাছে জামালের তিনগুণ
মার্বেল আছে। তাদের নিকট সর্বমোট ৭২ মার্বেল আছে। জামালের নিকট কতটি মার্বেল
আছে?}

ক) ১০

\textbf{\hl{খ) ১২ ✔}}

গ) ১৫

ঘ) ১৮

\textbf{\hl{✔ সঠিক উত্তর: খ) ১২}}

\begin{quote}
\textbf{ব্যাখ্যা:} \emph{ধরি, মালেকের মার্বেল আছে = x টি,}

\emph{জামালের মার্বেল আছে = (x/2) টি,}

\emph{এবং আরিফের মার্বেল আছে = (3x/2) টি।}

\emph{শর্তমতে, x+(x/2)+(3x/2)=72}

\emph{⇒ 6x=144⇒x=24 ।}

\emph{জামালের নিকট মার্বেল আছে = 24÷2=12 টি।}

\emph{∴ ১২}
\end{quote}

\textbf{একটি শ্রেণীতে প্রতি বেঞ্চে ৪ জন করে ছাত্র বসলে ৩ খানা বেঞ্চ খালি
থাকে। কিন্তু প্রতি বেঞ্চে ৩ জন করে বসলে ৬ জন ছাত্র দাঁড়িয়ে থাকতে হয়। ঐ শ্রেণীতে
ছাত্র সংখ্যা কত?}

ক) ৫৫ জন

\textbf{\hl{খ) ৬০ জন ✔}}

গ) ৬৫ জন

ঘ) ৭০ জন

\textbf{\hl{✔ সঠিক উত্তর: খ) ৬০ জন}}

\begin{quote}
\textbf{ব্যাখ্যা:} \emph{ধরি, বেঞ্চ সংখ্যা = x।}

\emph{১ম শর্তমতে, ছাত্রসংখ্যা = 4(x−3)=4x−12 ।}

\emph{২য় শর্তমতে, ছাত্রসংখ্যা = 3x+6 ।}

\emph{শর্তমতে, 4x−12=3x+6⇒x=18 ।}

\emph{সুতরাং ছাত্রসংখ্যা = 3×18+6=60 জন।}

\emph{∴ ৬০ জন}
\end{quote}

\section{23. সরল সমীকরণ-বীজগণিত \textbar\textbar{} সংজ্ঞা, সূত্রাবলীসহ এ টু
জেড}\label{ux9b8ux9b0ux9b2-ux9b8ux9aeux995ux9b0ux9a3-ux9acux99cux997ux9a3ux9a4-ux9b8ux99cux99e-ux9b8ux9a4ux9b0ux9acux9b2ux9b8ux9b9-ux98f-ux99f-ux99cux9a1}

সরল সমীকরণ হলো গণিতের একটি মৌলিক বিষয়। সহজভাষায় বলতে গেলে, সরল সমীকরণ
হলো একটি গণিতিক সমীকরণ যা একটি বা একাধিক অজানা পরিমাণের বাক্যাংশকে বর্ণনা
করে একটি পরিমাণ নির্ণয় করতে ব্যবহার করা হয়। এই আর্টিকেলের উদ্দেশ্য হলো সরল
সমীকরণের স্পষ্ট ধারণা ও ব্যবহার সম্পর্কে আলোচনা করা। প্রথমে আমরা সরল সমীকরণ
সম্পর্কে স্পষ্টতা সাধারণ উদাহরণগুলি দিয়ে ব্যাখ্যা করব। তারপরে আমরা সরল সমীকরণের
গঠন এবং মান নির্ণয়ের পদ্ধতি নিয়ে আলোচনা করব। তাছাড়া আমরা সরল সমীকরণের
বিভিন্ন ধরণ ও ব্যবহার সম্পর্কে আলোচনা করব । অত:পর বিগত সালের প্রশ্ন সমাধান
করতে করতে অধ্যায়টি সম্পর্কে স্পষ্ট ধারণা নিব । চলুন তাহলে শুরু করা যাক ।

\begin{itemize}
\item
  প্রাথমিক আলোচনা
\item
  \textbf{সরল সমীকরণ কাকে বলে?}
\item
  কিছু উপকারী কৌশল
\item
  সাধারণ সমীকরণ
\item
  সংখ্যার সমীকরণ
\item
  ভগ্নাংশের সমীকরণ
\item
  সমীকরণ সাজিয়ে সমাধান
\end{itemize}

\subsubsection{প্রাথমিক
আলোচনা}\label{ux9aaux9b0ux9a5ux9aeux995-ux986ux9b2ux99aux9a8-9}

এই অধ্যায়টি খুবই গুরুত্বপূর্ণ। কেননা এখান থেকে প্রশ্ন আসবেই। আবার পাটিগণিতের যে
অংকগুলো প্রতীক ধরে ধরে করা হয়।

সেগুলো এই অধ্যায়েরই অংক। ( = ) চিহ্ন দ্বারা একটি রাশির সাথে অপর একটি রাশির
সমতার সম্পর্ককে সমীকরণ বলা হয়।

নিচের সমীকরণ দুটি দেখুন:

ক. 5x - 18 = 14

খ. 2x = 4 সমীকরণে বামপক্ষে একটি অজ্ঞাত বীজগণিতীয় প্রতীক থাকে। উপরের উদাহরণ
দুটি অজ্ঞাত প্রতীক x. অজ্ঞাত প্রতীক হিসাবে সাধারণ ইংরেজি বর্ণমালার x, y, z,
অক্ষরগুলো ব্যবহৃত হয়।

\textbf{সমীকরণের বীজ :} সমীকরণ থেকে অজ্ঞাত প্রতীকের প্রাপ্ত মানকে প্রদত্ত
সমীকরণের বীজ বলা হয়।

5x − 18 = 14

বামপক্ষ = ডানপক্ষ

সমীকরণে (=) চিহ্নের বামপাশের রাশি হল বামপক্ষ এবং ডানপক্ষ রাশি হল ডানপক্ষ ।

\textbf{সমীকরণের সমাধান :}

সমীকরণের বীজ নির্ণয় করার প্রক্রিয়াকে সমীকরণের সমাধান বলা হয় ৷

\textbf{সরল সমীকরণ :}

যে সমীকরণে একঘাতবিশিষ্ট একটি মাত্র অজ্ঞাত বীজগণিতীয় প্রতীক থাকে তাকে সরল
সমীকরণ বলা হয়।

উদাহরণ: 2 (5x - 18) = 14

\textbf{টিপস:} এই অধ্যায়ের অংকগুলো সহজ মনে হলেও লাইনের পর লাইন লিখে করতে
করতে অনেক সময় নষ্ট হয়। তাই

সবসময় চেষ্টা করুন, কিভাবে কম লিখে বেশি কাজ করা যায় । অর্থাৎ দুই, তিন লাইনের
কাজ এক লাইনেই করার চেষ্টা করুন।

\subsubsection{সরল সমীকরণ কাকে
বলে?}\label{ux9b8ux9b0ux9b2-ux9b8ux9aeux995ux9b0ux9a3-ux995ux995-ux9acux9b2}

যে সমীকরণে একটি মাত্র অজ্ঞাত রাশি থাকে তাকে একচলকবিশিষ্ট সমীকরণ বা সরল
সমীকরণ বলে । এ সম্পর্কে বিস্তর ধারণা নিচে আরো ভালভাবে জানার চেষ্ট করব ।

\subsubsection{কিছু উপকারী
কৌশল}\label{ux995ux99b-ux989ux9aaux995ux9b0-ux995ux9b6ux9b2}

দ্রুত এই নিয়মের অংকগুলো করার জন্য নিচের কৌশলগুলো অনুসরণ করুন।

১. ধরার সময় ছোট সংখ্যাটি ধরে অংক করার চেষ্টা করুন। যেমন: একটি সংখ্যা অপর
সংখ্যার 3 গুণ । তাহলে ছোট সংখ্যাটি x এবং বড় সংখ্যাটি ৩x .

২. যথাসম্ভব ভগ্নাংশকে এড়িয়ে চলার চেষ্টা করুন। যেমন: একটি সংখ্যা আরেকটি সংখ্যার
অর্ধেক। এভাবে আসলে যে সংখ্যাটি অর্ধেক তা ছোট তাই তাকে x ধরুন। তাহলে বড়টি হবে
2x, কেননা, ১৫ সংখ্যাটি ৩০ এর অর্ধেক হলে স্বাভাবিকভাবেই ৩০ সংখ্যাটি ১৫ এর
দ্বিগুণ। বড় সংখ্যাটি x ধরলে

ছোট সংখ্যাটি হবে - \(\frac{x}{2}\) যেভাবে অংক করা আসলেই কঠিন।

৩. সমীকরণ সাজানোর সময় দুপাশের মানকে সমান সমান বানান । সমান না হলে সঠিক
উত্তর আসবে না । যেমন: বর্তমানে পিতার বয়স পুত্রের বয়সের ৫গুণ, ৫ বছর পর পিতার
বয়স পুত্রের বয়সের ৩গুণ হবে। পুত্রের বর্তমান বয়স কত?

ধরি, পুত্রের বর্তমান বয়স X সুতরাং পিতার বর্তমান বয়স 5x এখন ৫ বছর পর পুত্রের বয়স
হবে x+5 এবং পিতার বয়স হবে 5x+5 এ পর্যন্ত লিখার পর সমীকরণ সাজাতে হবে।

প্রশ্নমতে, 3(x+5) = 5x + 5 (যে কোন সমীকরণের অংকে এই লাইনটি সাজাতেই মাথা
খাঁটাতে হবে।)

এখানে না বুঝে শুধু প্রশ্ন পড়ে লিখলে মনে হবে পিতার বয়সের সাথে ৩ গুণ করে দেই?
কিন্তু মনে রাখবেন = চিহ্ন হল একটি দাঁড়িপাল্লার দাঁড়ির মত । যার দুপাশে সমান
ভারবাহী উপাদান থাকতে হবে। এখানে ৫ বছর পর পুত্র এবং পিতার বয়স = এর দুপাশে
লিখলে হয় x+5 = 5x+5 -\/-\/-

যেখানে পিতার বয়স আগে থেকেই ৩ গুণ বেশি আছে, তাই পুত্রের বয়সের সাথে ৩ গুণ করলে
দাঁড়িপাল্লাটি সমান হবে। কিন্তু ভুল করে পিতার সাথে আবার ৩ গুণ করলে ভাড়ী পাল্লাটি
আরো বেশি ভাড়ী হবে এবং ভুল উত্তর আসবে।

নিচের অংকটি দেখে বুঝুন এই বিষয়টি বার বার আসবে। তাই এখান থেকে সম্পূর্ণ বুঝে নিন
।

\textbf{প্রশ্ন নম্বর-} ১. একটি কলম ও একটি বইয়ের মোট দাম 95 টাকা। কলমটির দাম
15 টাকা বেশি ও বইটির দাম 14 টাকা কম হলে কলমটি বইয়ের দামের 2 গুণ হতো। একটি
কলমের দাম কত? {[} সরকারী মাধ্যমিক বিদ্যালয় সহ: শিক্ষক-০৯ {]}

\begin{itemize}
\item
  (ক) ৪০ টাকা
\item
  (খ) ৪৯ টাকা
\item
  (গ) ৫৫ টাকা
\item
  (ঘ) ৬০ টাকা
\end{itemize}

\textbf{\hl{উত্তরঃ (খ) ৪৯ টাকা}}

\textbf{সমাধান:}

ধরি, কলমের দাম = x টাকা (সমষ্টি দেয়া থাকলে একটি ধরতে হয় এবং অন্যটি সমষ্টি
থেকে বিয়োগ করতে হয়।)

বইয়ের দাম = 95 - x টাকা

প্রশ্নমতে, (x + 15) = 2(95 - x - 14 )

বা, x + 15 = 190 - 2x - 28

বা, 3x = 162 - 15

বা, 3x = 147

বা, x = 147

X = 49

সুতরাং একটি কলমের দাম ৪৯ টাকা

প্রশ্ন একটু ঘুরিয়ে দিলে সমীকরণ প্রশ্নের যুক্তি অনুযায়ী সাজাতে হবে (আগে থেকেই মুখস্থ
কোন ধারণা প্রয়োগ করা যাবে না ।

\textbf{প্রশ্ন নম্বর-} ২. একটি কলম ও একটি বইয়ের মূল্য একত্রে ৯৫ টাকা । কলমটির
মূল্য ১৫ টাকা বেশি এবং বইটির মূল্য ১৪ টাকা কম হলে কলমটির মূল্য বইটির মূল্যের
দ্বিগুণ হতো । বইটির মূল্য কত টাকা? {[}প্রাই: সহ:শি: নি: পরীক্ষা-২০১৪
(অনু:২০১৮){]}+ {[}কৃষি সম্প্রসারণ অধিদপ্তর- (অফি: সহ:+টাইপিস্ট)-২০১৮{]}

\begin{itemize}
\item
  (ক) ৪০ টাকা
\item
  (খ) ৪৯ টাকা
\item
  (গ) ৪৬ টাকা
\item
  (ঘ) কোনোটিই নয়
\end{itemize}

\textbf{\hl{উত্তরঃ (গ) ৪৬ টাকা}}

\textbf{সমাধান:}

ধরি,

বইয়ের মূল্য = X

সুতরাং কলমের মূল্য = ৯৫ - X

প্রশ্নমতে,

(৯৫ - x +১৫) = ২(x-১৪) {[} কম বেশি করার পর কলমের দাম দ্বিগুন অর্থাৎ বইয়ের
দাম অর্ধেক, তাই বইয়ের দামের সাথে ২ গুণ করলে কলম ও বইয়ের দাম সমান সমান হবে
। {]}

বা, ১১০ - x = ২x - ২৮

বা, -৩x = -১৩৮

.. x = ৪৬

সুতরাং বইয়ের মূল্য = ৪৬ টাকা ।

শুদ্ধি পরীক্ষা : বই= ৪৬টাকা, তাহলে কলম = ৯৫-৪৬ = ৪৯টাকা

এখন ৪৯+১৫ = ৬৪ টাকা হলো ৪৬- ১৪ = ৩২ টাকার দ্বিগুণ ।

৪. আড়াআড়ি গুণ করার সময় যে পাশে গুণ করলে x এর মান বড় হবে সেপাশে আগে গুণ করুন।
যেমন: \(\frac{x}{2} = \frac{x\  + \ 3}{4}\) কে 2x+6 = 4x না লিখে এভাবে
লিখুন 4x = 2x + 3 অথবা সরাসরি 4x - 2x = 3, তাহলে ১ লাইন কম লিখেই উত্তর বের
করা সম্ভব হবে ।

৫. ছোট ছোট যোগ বিয়োগ গুণ ভাগ না লিখে মুখে মুখে করে ফলাফল বসান। যেমন: 4x + 9
= 15+ 3x কে 4x - 3x = 15 - 9 বা x = 6 এভাবে ৩ লাইনে না লিখে মাঝের লাইনের
কাজটি মুখে মুখে করে ফেলুন এবং সরাসরি উত্তর বসান x = 6 (সমীকরণের অংক গুলোতে কম
লিখে বেশি কাজ করতে হবে)

৬. এমসিকিউ পরীক্ষায় সহজ অংকগুলোর অপশন থেকে মান বসিয়ে বের করুন ।

৭. আবার একটি সংখ্যা আরেকটি সংখ্যার দেড়গুণ/ আড়ই গুণ হলে ছোট সংখ্যাটিকে শুধু x না
ধরে মাথা খাঁটিয়ে 2x এবং

অন্যটিকে 3x বা 5x ধরে করলে ভগ্নাংশের ঝামেলা এড়িয়ে যাওয়া যাবে। (বিস্তারিত
ভগ্নাংশের সমীকরণ পদ্ধতিতে)

\subsubsection{সাধারণ
সমীকরণ}\label{ux9b8ux9a7ux9b0ux9a3-ux9b8ux9aeux995ux9b0ux9a3}

সাধারণ সরল সমীকরণের অংকগুলোতে প্রশ্নের মধ্যেই সমীকরণ সাজিয়ে দেয়া থাকবে,
এক্ষেত্রে সমীকরণের বাম পাশে অজানা রাশিগুলোকে যেমন (x, y, a, b ) রেখে
ডানপাশে অন্য সংখ্যাগুলোকে রেখে সমাধান করতে হয়। যেহেতু এমসিকিউ প্রশ্নের উত্তর
প্রদানের সময় খুব কম থাকে তাই যত কম লিখে সমাধান করা যায় ততই ভাল। কিভাবে কম
সময়ে করবেন তা নিচে বুঝিয়ে দেয়া হল ।

\textbf{প্রশ্ন নম্বর-} ৩. x - 6 = 7x - 48 কে সমাধান করলে x এর মান হবে- {[}
প্রাথমিক বিদ্যালয় প্রধান শিক্ষক (পদ্ম) ২০০৯ {]}

\begin{itemize}
\item
  (ক) 3
\item
  (খ) 5
\item
  (গ) - 6
\item
  (ঘ) 7
\end{itemize}

\textbf{\hl{উত্তরঃ (ঘ) 7}}

\textbf{সমাধান:}

x -- 6 = 7x - 48

বা, - 6x = 42 {[} এখানে মাঝের (x-7x = - 48+6 ) লাইনটি না লিখে সরাসরি করা
হয়েছে {]}

সুতরাং x = 7

\textbf{প্রশ্ন নম্বর-} ৪. 2x + 15 = 27 - 4x কে সমাধান করলে x এর মান হবে-
{[} প্রাথমিক বিদ্যালয় প্রধান শিক্ষক (শাপলা) ২০০৯{]}

\begin{itemize}
\item
  (ক) - 1
\item
  (খ) 2
\item
  (গ) - 2
\item
  (ঘ) 3
\end{itemize}

\textbf{\hl{উত্তরঃ (খ) 2}}

\textbf{সমাধান:}

সরাসরি, 6x = 12

বা, x = 2

\textbf{প্রশ্ন নম্বর-} ৫. x এর মান কত হলে a (x-a) = b (x-b) হবে? {[} শ্রম
মন্ত্রণালয়ের সহকারী পরিচালক ২০০৫ {]}

\begin{itemize}
\item
  (ক) 1
\item
  (খ) a - b
\item
  (গ) b - a
\item
  (ঘ) a + b
\end{itemize}

\textbf{\hl{উত্তরঃ (ঘ) a + b}}

\textbf{সমাধান:}

a (x-a) = b (x-b)

বা, \(ax - a^{2} = bx - b^{2}\)

বা, ax - bx = \(a^{2} - b^{2}\)

বা, x ( a - b ) = ( a + b ) ( a - b )

বা, x = a + b

\textbf{প্রশ্ন নম্বর-} ৬. যদি \((\frac{x}{3} + 1) = (\frac{x}{4} + 1)\) হয়,
তবে x এর মান কোনটি? {[} টেক্সটাইল ইনষ্টিটিউট ও টেক্সটাইল ইঞ্জিনিয়ারিং কলেজের
ইনস্ট্রাক্টর: ২০১৮{]}

\begin{itemize}
\item
  (ক) 3
\item
  (খ) 2
\item
  (গ) 1
\item
  (ঘ) 0
\end{itemize}

\textbf{\hl{উত্তরঃ (ঘ) 0}}

\textbf{সমাধান:}

\((\frac{x}{3} + 1) = (\frac{x}{4} + 1)\)

বা, \(\frac{x}{3} = \frac{x}{4} + 1 - 1\)

বা, \(\frac{x}{3} - \frac{x}{4} = 0\)

বা, \(\frac{4x\  - \ 3x}{12} = 0\)

বা, x = 0

\textbf{প্রশ্ন নম্বর-} ৭. (2 + x) + 3 = 3 (x + 2) হলে x এর মান কত? {[}১৫ তম
বিসিএস{]}

\begin{itemize}
\item
  (ক) -(1/2)
\item
  (খ) (1/2)
\item
  (গ) (1/3)
\item
  (ঘ) (2/3)
\end{itemize}

\textbf{\hl{উত্তরঃ (ক)}} \(- \frac{1}{2}\)

\textbf{সমাধান:}

(2 + x) + 3 = 3 (x + 2)

বা, 2 + x + 3 = 3x + 6

বা, x - 3x = 6 - 2 - 3

বা, -2x = 6 - 5

বা, -2x = 1

বা, 2x = - 1

বা, x = \(- \frac{1}{2}\)

\subsubsection{সংখ্যার
সমীকরণ}\label{ux9b8ux996ux9afux9b0-ux9b8ux9aeux995ux9b0ux9a3}

সংখ্যার সমীকরণের ছোট ছোট অনেকগুলো প্রশ্ন পাটিগণিতের সংখ্যায়ন অধ্যায়ে আলোচনা
করা হয়েছে। এখানে দেখানো হলো কিভাবে সমীকরণ সাজিয়ে সমাধান করতে হয়।এই
অধ্যায়টি করার সময় শুধু প্রশ্ন আর সমাধানের থেকে সমীকরণ সাজানোর বিভিন্ন
যুক্তিগুলোকে গুরুত্ব দিতে হবে ।

\textbf{প্রশ্ন নম্বর-} ৮. একটি সংখ্যার তিনগুণের সাথে দ্বিগুণ যোগ করলে ৯০ হয়।
সংখ্যাটি কত? {[} PETROBANGLA - (UDA)-2017{]}

\begin{itemize}
\item
  (ক) ১৮
\item
  (খ) ২০
\item
  (গ) ২৪
\item
  (ঘ) ১৬
\end{itemize}

\textbf{\hl{উত্তরঃ (ক) ১৮}}

\textbf{সমাধান:}

ধরি, সংখ্যাটি x

প্রশ্নমতে,

৩x + 2x = ৯০

বা, ৫x = ৯০

বা, x = ১৮

সুতরাং সংখ্যাটি = ১৮

\textbf{প্রশ্ন নম্বর-} ৯. একটি সংখ্যার ৪ গুণের সাথে ১০ যোগ করা হলে উত্তর হয়
সংখ্যাটির ৫ গুণ অপেক্ষা ৫ কম। সংখ্যাটি কত ? {[}IBBL- (ATO)-2017{]}

\begin{itemize}
\item
  (ক) ৩০
\item
  (খ) ২০
\item
  (গ) ২৫
\item
  (ঘ) ১৫
\end{itemize}

\textbf{\hl{উত্তরঃ (ঘ) ১৫}}

\textbf{সমাধান:}

ধরি, সংখ্যাটি ক

প্রশ্নমতে,

৪ক + ১০ = ৫ক-৫

বা, - ক = -১৫

বা, ক = ১৫

সুতরাং সংখ্যাটি = ১৫

\textbf{প্রশ্ন নম্বর-} ১০. কোনো সংখ্যার দ্বিগুণের সাথে ৫ যোগ করলে যোগফল
সংখ্যাটি অপেক্ষা ৭ বেশি হয়। সংখ্যাটি কত? {[} উপ: পরিসংখ্যন কর্ম:১০ {]}

\begin{itemize}
\item
  (ক) ৫
\item
  (খ) ২
\item
  (গ) ৩
\item
  (ঘ) ৪
\end{itemize}

\textbf{\hl{উত্তরঃ (খ) ২}}

\textbf{সমাধান:}

ধরি, সংখ্যাটি ক

প্রশ্নমতে,

২ক + ৫ = ক + ৭

বা, ক = ২

সুতরাং সংখ্যাটি = ২

\subsubsection{ভগ্নাংশের
সমীকরণ}\label{ux9adux997ux9a8ux9b6ux9b0-ux9b8ux9aeux995ux9b0ux9a3}

\textbf{প্রশ্ন নম্বর-} ১১. একটি সংখ্যা অপর একটি সংখ্যার \(\frac{2}{3}\) গুণ।
সংখ্যা দুইটির সমষ্টি ১০০ হলে বড় সংখ্যাটির মান কত? {[}পরিবার কল্যাণ পরিদর্শিকা
নিয়োগ পরীক্ষা: ২০১০ {]}

\begin{itemize}
\item
  (ক) ৫৫
\item
  (খ) ৭০
\item
  (গ) ৬৫
\item
  (ঘ) ৬০
\end{itemize}

\textbf{\hl{উত্তরঃ (ঘ) ৬০}}

\textbf{সমাধান:}

শুধু X না ধরে একটি সংখ্যা 2x এবং অপর সংখ্যাটি 3x ধরে করুন।

তাহলে ( 2x+3x) বা, 5x = 100 সুতরাং x = 20, বড় সংখ্যাটি হবে 3×20 = 60

দেড়গুণ, আড়াই গুণ, সাড়ে তিন গুণ এভাবে ভগ্নাংশ থাকলে একটু মাথা খাঁটিয়ে করতে হবে।
নিচের প্রশ্নটি ভালোভাবে বুঝুন । একটি আরেকটির দেড়গুণ হলে একটিকে 2x এবং অপরটিকে
3x ধরে করুন। কেননা ৩ হচ্ছে ২ এর দেড়গুণ । এভাবে করলেও উত্তর একই আসবে এবং
ভগ্নাংশ না আসায় মুখে মুখে করা যাবে । যেমন:

\textbf{প্রশ্ন নম্বর- ১২. ৫৫০ গ্রাম একটি কেক-এ চিনির দ্বিগুণ ময়দা আছে।
কিচমিচের দেড়গুণ চিনি আছে। কেকটিতে ময়দার পরিমান কত?}

\begin{itemize}
\item
  (ক) ১০০
\item
  (খ) ১৫০
\item
  (গ) ২৫০
\item
  (ঘ) ৩০০
\end{itemize}

\textbf{\hl{উত্তরঃ (ঘ) ৩০০}}

\textbf{সমাধান:}

৫৫০ গ্রাম কে মোট ৩ ভাগ করতে হবে যেখানে চিনি , ময়দা ও কিচমিচ থাকবে। প্রথম
অংশে দ্বিগুণ আছে তাই সহজ। কিন্তু দ্বিতীয় অংশে দেড়গুণ আছে । যেখানে দেড়গুণ,
আড়াইগুণ থাকবে সেখান থেকে ভাবা শুরু করতে হবে।

কিচমিচের দেড়গুণ চিনি, অর্থাৎ ২ গুণ কিচমিচ হলে ৩ গুণ চিনি। আবার চিনির দ্বিগুন
ময়দা অর্থাৎ ৩ গুন চিনি হলে ৬ গুণ ময়দা। তাহলে কেকটিকে মোট ২+৩+৬ = ১১ ভাগ
করতে হবে। এখন যেহেতু ময়দা কত আছে তা বের করতে বলা হয়েছে তাই ৫৫০ গ্রাম কে ১১
ভাগ করলে প্রতি ১ ভাগে থাকে ৫০ গ্রাম। যেহেতু ময়দা ৬ ভাগ তাই ময়দা আছে ৬×৫০ =
৩০০ গ্রাম । এগুলোকেই x আকারে লিখে বলা যায়: 2x + 3x + 6x = 550, . x = 50
then 6x = 6x50 = 300

\textbf{Note:}এভাবে পাটিগণিত, বীজগণিত ও জ্যমিতির বিভিন্ন অধ্যায়ে ভগ্নাংশ
আকারে সংখ্যার তুলনা আসলে মাথা খাঁটিয়ে উপরের নিয়মে করতে থাকুন। খাতা কলম ছাড়াই
মুখে মুখে কয়েক সেকেন্ডে অংক হয়ে যাবে।

\textbf{প্রশ্ন নম্বর-} ১৩. একটি ক্রিকেট দলে যতজন স্ট্যাম্প আউট হলো তার দেড়গুণ কট
আউট হলো এবং মোট উইকেটের অর্ধেক বোল্ড আউট হলো। এই দলের কতজন কট আউট হলো? {[}
১১তম বিসিএস {]}

\begin{itemize}
\item
  (ক) ৫
\item
  (খ) ৭
\item
  (গ) ৩
\item
  (ঘ) ৮
\end{itemize}

\textbf{\hl{উত্তরঃ (গ) ৩}}

\textbf{সমাধান:}

ধরি, স্ট্যাম্প আউট হলো `ক' জন

সুতরাং কট আউট = \(\frac{\text{৩ক}}{\text{২}}\)

প্রশ্নানুসারে,

ক + \(\frac{\text{৩ক}}{\text{২}}\) + ৫ = ১০

বা, ৫ক = ১০

বা, ক = ২

সুতরাং কট আউট হলো = \(\frac{\text{৩}\  \times \ \text{২}}{\text{২}}\) = ৩

\textbf{মুখে মুখে করতে চাইলে এভাবে ভাবুন:}

মোট উইকেট ১০টি, অর্ধেক = ৫টি বাকী ৫টির মধ্যে একটি আরেকটির দেড়গুণ অর্থ হলো
স্ট্যাম্প আউট দুইজন এবং কট আউট ৩জন।

\subsubsection{সমীকরণ সাজিয়ে
সমাধান}\label{ux9b8ux9aeux995ux9b0ux9a3-ux9b8ux99cux9af-ux9b8ux9aeux9a7ux9a8}

\textbf{প্রশ্ন নম্বর-} ১৪. রহিমের আয়ের দ্বিগুণের সাথে ১১০ টাকা যোগ করলে ৭০০০
টাকা হয়। রহিমের আয় কত? {[}PETROBANGLA - (UDA) 17{]}

\begin{itemize}
\item
  (ক) ৩২৭৫
\item
  (খ) ৩৪৪৫
\item
  (গ) ৩২১০
\item
  (ঘ) ৫২৭৫
\end{itemize}

\textbf{\hl{উত্তরঃ (খ) ৩৪৪৫}}

\textbf{সমাধান:}

ধরি, রহিমের আয় x টাকা

বা, 2x+110 = 7000

বা,2x = 7000 - 110

বা,2x = 6890

বা,x = 6890÷2

বা, x = 3445 টাকা

শর্টকাট: ১১০ টাকা যোগ করার পর ৭০০০ টাকা হলে যোগ করার আগে ছিলা ৬৮৯০ টাকা।
আবার ৬৮৯০ টাকা দ্বিগুণ করার আগে ছিল ৬৮৯০÷২ = ৩৪৪৫।

\textbf{প্রশ্ন নম্বর-} ১৫. যদি কবিরের ওজন ৮ পাউন্ড কমে, তবে কবিরের ওজন তার
বোনের ওজনের দ্বিগুণ হবে। এখন তাদের উভয়ের ওজন ২৭৮ পাউন্ড। বর্তমানে কবিরের ওজন
কত ? {[} BKB - (SO)-2017{]} \& {[} Sonali Bank Officer (General) - 2018
{]}

\begin{itemize}
\item
  (ক) 147
\item
  (খ) 188
\item
  (গ) 133
\item
  (ঘ) 139
\end{itemize}

\textbf{\hl{উত্তরঃ (খ) 188}}

\textbf{সমাধান:}

ধরি, কবিরের ওজন = x পাউন্ড । (সাধারণত যার টা বের করতে বলে তারটা ধরে
সমীকরণ সাজাতে হয়।)

সুতরাং তার বোনের ওজন = 278 - x ( সমষ্টি দেয়া থাকলে একটাকে ধরে অন্যটা বিয়োগ
করে বের করা যায়)

প্রশ্নমতে,

x-8 = 2 (278-x) (কবিরের ওজন থেকে ৮ বাদ দিলে যা থাকে তা তার বোনের দ্বিগুণ
তাই বোনের সাথে ২ গুণ করে সমান )

বা, x-8 = 556 - 2x

বা, x + 2x = 556 + 8

বা, 3x = 564

বা, x = 564÷3

বা, x = 188

সুতরাং বর্তমানে কবিরের ওজন 188 পাউন্ড

\textbf{প্রশ্ন নম্বর-} ১৬. রফিকের ওজন যদি ১৭ কেজি কমে যায় তবে তার ওজন
আরিফের ওজনের অর্ধেক হয়ে যাবে । তাদের দুজনের ওজনের যোগফল ১৪০ কেজি হলে
রফিকের ওজন কত কেজি ? {[}কৃষি সম্প্রসারণ অধিদপ্তর- (অফি: সহ:+টাইপিস্ট)-২০১৮{]}

\begin{itemize}
\item
  (ক) ৫৫
\item
  (খ) ৫৮
\item
  (গ) ৬১
\item
  (ঘ) কোনটিই নয়
\end{itemize}

\textbf{\hl{উত্তরঃ (খ) ৫৮}}

\textbf{সমাধান:}

(এখানে যে ছোট তার থেকে বাদ দেয়ার পর সে ই অর্ধেক হয়ে গেছে। )

ধরি, রফিকের ওজন = x কেজি

আরিফের ওজন = 140-x কেজি।

১ম শর্তমতে,

2(x-17) = 140-x (যে ছোট তার সাথে গুণ)

বা, 2x-34 = 140-x

বা, 2x+x = 140+34

বা, 3x = 174

বা, x = 174÷3

বা, x = 58

\textbf{মুখে মুখে করার যুক্তি:}

১৪০ থেকে ১৭ বাদ দিলে হয় ১৪০-১৭ = ১২৩ যাকে ৩ দিয়ে

ভাগ করলে হয় ৪১ । এবার যে ছোট তার থেকে ১৭ বাদ দেয়ার তার ওজন ছিল ৪১+১৭ =
৫৮

\section{24. সমীকরণ MCQ প্রশ্ন উত্তর
অনুশীলন}\label{ux9b8ux9aeux995ux9b0ux9a3-mcq-ux9aaux9b0ux9b6ux9a8-ux989ux9a4ux9a4ux9b0-ux985ux9a8ux9b6ux9b2ux9a8}

\textbf{x কে চলক ধরে a\^{}2x + b = 0 সমীকরণের ঘাত নিচের কোনটি?}

ক) 3

খ) 2

\textbf{\hl{গ) 1 ✔}}

ঘ) 0

\textbf{\hl{✔ সঠিক উত্তর: গ) 1}}

\begin{quote}
\textbf{ব্যাখ্যা:} \emph{a\^{}2x + b = 0 সমীকরণটিতে x চলকের সর্বোচ্চ ঘাত বা
পাওয়ার হলো 1।}
\end{quote}

\textbf{3x + 2y = 15 সমীকরণটির সমাধান কতটি?}

ক) একটি

খ) দুইটি

গ) একটিও না

\textbf{\hl{ঘ) অসীম ✔}}

\textbf{\hl{✔ সঠিক উত্তর: ঘ) অসীম}}

\begin{quote}
\textbf{ব্যাখ্যা:} \emph{সমীকরণটিতে x ও y দুইটি অজানা চলক। একটি সরল সমীকরণে
দুই বা ততোধিক চলক থাকলে এবং নির্দিষ্ট মান না দেওয়া থাকলে, তার অসংখ্য বা অসীম
সমাধান থাকে।}
\end{quote}

\textbf{(2 + x) + 3 = 3 (x + 2) হলে, x এর মান কত?}

\textbf{\hl{ক) -(1/2) ✔}}

খ) (1/3)

গ) (1/2)

ঘ) (2/3)

\textbf{\hl{✔ সঠিক উত্তর: ক) -(1/2)}}

\begin{quote}
\textbf{ব্যাখ্যা:} \emph{(2 + x) + 3 = 3(x + 2)}

\emph{⇒ 2 + x + 3 = 3x + 6}

\emph{⇒ x + 5 = 3x + 6}

\emph{⇒ x − 3x = 6 − 5}

\emph{⇒ −2x = 1}

\emph{∴ x = -(1/2)}
\end{quote}

\textbf{x\^{}3 − px + 10 = 0 এর একটি সমাধান 2 হলে p এর মান কত?}

ক) -−2

খ) -−9

\textbf{\hl{গ) 9 ✔}}

ঘ) 10

\textbf{\hl{✔ সঠিক উত্তর: গ) 9}}

\begin{quote}
\textbf{ব্যাখ্যা:} \emph{সমীকরণে x = 2 বসিয়ে পাই:}

\emph{(2)\^{}3 − p(2) + 10 = 0}

\emph{⇒ 8 − 2p + 10 = 0}

\emph{⇒ 18 − 2p = 0}

\emph{⇒ 2p = 18}

\emph{⇒ p = (18/2)}

\emph{∴ p = 9}
\end{quote}

\textbf{x\^{}3 + hx + 10 = 0 সমীকরণটির x এর মান 2 হলে, h এর মান কত?}

ক) 9

\textbf{\hl{খ) -−9 ✔}}

গ) -−5

ঘ) 5

\textbf{\hl{✔ সঠিক উত্তর: খ) -−9}}

\begin{quote}
\textbf{ব্যাখ্যা:} \emph{সমীকরণে x = 2 বসিয়ে পাই:}

\emph{(2)\^{}3 + h(2) + 10 = 0}

\emph{⇒ 8 + 2h + 10 = 0}

\emph{⇒ 2h + 18 = 0}

\emph{⇒ 2h = −18}

\emph{⇒ h = −18/2}

\emph{∴ h = −9}
\end{quote}

\textbf{(x − 5)(a + x) = x\^{}2 − 25 হলে, a এর মান কত?}

ক) 25

খ) −25

\textbf{\hl{গ) 5 ✔}}

ঘ) −5

\textbf{\hl{✔ সঠিক উত্তর: গ) 5}}

\begin{quote}
\textbf{ব্যাখ্যা:} \emph{(x − 5)(a + x) = x\^{}2 − 25}

\emph{⇒ (x − 5)(a + x) = (x + 5)(x − 5)}

\emph{⇒ (a + x) = (x + 5)}

\emph{∴ a = 5}
\end{quote}

\textbf{যদি x + 5y = 24 এবং x = 3y হয়, তাহলে y = কত?}

ক) 2

\textbf{\hl{খ) 3 ✔}}

গ) 4

ঘ) 5

\textbf{\hl{✔ সঠিক উত্তর: খ) 3}}

\begin{quote}
\textbf{ব্যাখ্যা:} \emph{x + 5y = 24}

\emph{মান বসিয়ে পাই, 3y + 5y = 24}

\emph{⇒ 8y = 24}

\emph{∴ y = 3}
\end{quote}

\textbf{(3/x)+(4/x+1)=2 হলে, x-এর মান কত?}

ক) 1

খ) 2

\textbf{\hl{গ) 3 ✔}}

ঘ) 4

\textbf{\hl{✔ সঠিক উত্তর: গ) 3}}

\begin{quote}
\textbf{ব্যাখ্যা:} \emph{(3/x)+(4/x+1)=2}

\emph{⇒ (4/x+1)=2−(3/x)}

\emph{⇒ (4/x+1)=(2x−3/x)}

\emph{⇒ 4x = (x + 1)(2x − 3)}

\emph{⇒ 4x = 2x\^{}2 − 3x + 2x − 3}

\emph{⇒ 4x = 2x\^{}2 − x − 3}

\emph{⇒ 2x\^{}2 − 5x − 3 = 0}

\emph{⇒ 2x\^{}2 − 6x + x − 3 = 0}

\emph{⇒ 2x(x − 3) + 1(x − 3) = 0}

\emph{⇒ (x − 3)(2x + 1) = 0}

\emph{∴ x = 3}
\end{quote}

\textbf{যদি (x/3)−(x/5)=2 হয়, x এর মান কত?}

ক) ১৮

খ) ১৭

\textbf{\hl{গ) ১৫ ✔}}

ঘ) ২১

\textbf{\hl{✔ সঠিক উত্তর: গ) ১৫}}

\begin{quote}
\textbf{ব্যাখ্যা:} \emph{(x/3)−(x/5)=2}

\emph{⇒ (5x−3x/15)=2}

\emph{⇒ (2x/15)=2}

\emph{⇒ 2x = 30}

\emph{∴ x = 15}
\end{quote}

\textbf{(1/x−3)+(1/x−4)=(1/x−2)+(1/x−5) এর সমাধান কত?}

ক) x=(2/7)

\textbf{\hl{খ) x=(7/2) ✔}}

গ) x = 2

ঘ) x = 7

\textbf{\hl{✔ সঠিক উত্তর: খ) x=(7/2)}}

\begin{quote}
\textbf{ব্যাখ্যা:} \emph{(1/x−3)+(1/x−4)=(1/x−2)+(1/x−5)}

\emph{⇒ (1/x−3)−(1/x−2)=(1/x−5)−(1/x−4)}

\emph{⇒ (x−2−x+3/(x−3)(x−2))=(x−4−x+5/(x−5)(x−4))}

\emph{⇒ (1/x\^{}2−5x+6)=(1/x\^{}2−9x+20)}

\emph{⇒ x\^{}2 − 9x + 20 = x\^{}2 − 5x + 6}

\emph{⇒ 4x = 14}

\emph{∴ x=(7/2)}
\end{quote}

\textbf{(x/3)+(4/x+1)=2 হলে, x এর মান কত?}

ক) 1

\textbf{\hl{খ) 2 ✔}}

গ) 3

ঘ) 4

\textbf{\hl{✔ সঠিক উত্তর: খ) 2}}

\begin{quote}
\textbf{ব্যাখ্যা:} \emph{(x/3)+(4/x+1)=2}

\emph{⇒ (x(x+1)+12/3(x+1))=2}

\emph{⇒ (x\^{}2+x+12/3x+3)=2}

\emph{⇒ x\^{}2+x+12=2(3x+3)}

\emph{⇒ x\^{}2+x+12=6x+6}

\emph{⇒ x\^{}2+x−6x+12−6=0}

\emph{⇒ x\^{}2−5x+6=0}

\emph{⇒ x\^{}2−3x−2x+6=0}

\emph{⇒ x(x−3)−2(x−3)=0}

\emph{⇒ (x−3)(x−2)=0}

\emph{∴ x = 2 অথবা x = 3 (এখানে অপশন অনুযায়ী সঠিক উত্তর 2)}
\end{quote}

\textbf{(x/6)+2=(x/2) হলে x এর মান নিচের কোনটি?}

ক) 3

খ) 4

\textbf{\hl{গ) 6 ✔}}

ঘ) 8

\textbf{\hl{✔ সঠিক উত্তর: গ) 6}}

\begin{quote}
\textbf{ব্যাখ্যা:} \emph{(x/6)+2=(x/2)}

\emph{⇒ (x/2)−(x/6)=2}

\emph{⇒ (3x−x/6)=2}

\emph{⇒ (2x/6)=2}

\emph{⇒ 2x = 12}

\emph{⇒ x=(12/2)}

\emph{∴ x = 6}
\end{quote}

\textbf{x\^{}3 + 1 = 0 সমীকরণটির বাস্তব মূলের সংখ্যা কয়টি?}

\textbf{\hl{ক) 1 ✔}}

খ) 2

গ) 3

ঘ) নাই

\textbf{\hl{✔ সঠিক উত্তর: ক) 1}}

\begin{quote}
\textbf{ব্যাখ্যা:} \emph{x\^{}3 + 1 = 0}

\emph{⇒ (x + 1)(x\^{}2 − x + 1) = 0}

\emph{এখানে, x + 1 = 0 ⇒ x = −1 (একটি বাস্তব মূল)}

\emph{এবং x\^{}2 − x + 1 = 0 সমীকরণের ক্ষেত্রে নিশ্চায়ক
D=(−1)\^{}2−4·1·1=1−4=−3 (ঋণাত্মক), যা অবাস্তব মূল দেয়।}

\emph{∴ বাস্তব মূলের সংখ্যা মাত্র ১টি।}
\end{quote}

\textbf{যদি 2+(6/x)=5 হয়, তবে 2x + 8 = ?}

ক) 4

খ) 10

\textbf{\hl{গ) 12 ✔}}

ঘ) কোনোটিই নয়

\textbf{\hl{✔ সঠিক উত্তর: গ) 12}}

\begin{quote}
\textbf{ব্যাখ্যা:} \emph{2+(6/x)=5}

\emph{⇒ (6/x)=5−2}

\emph{⇒ (6/x)=3}

\emph{⇒ 3x = 6}

\emph{⇒ x = 2}

\emph{এখন, 2x + 8 = 2(2) + 8 = 4 + 8 = 12}
\end{quote}

\textbf{x√(0.09)=3 হলে, x এর মান---}

ক) (3/10)

খ) (1/3)

\textbf{\hl{গ) 10 ✔}}

ঘ) (10/3)

\textbf{\hl{✔ সঠিক উত্তর: গ) 10}}

\begin{quote}
\textbf{ব্যাখ্যা:} \emph{x√(0.09)=3}

\emph{⇒ x(0.3)=3}

\emph{⇒ x=(3/0.3)}

\emph{∴ x = 10}
\end{quote}

\textbf{y = x\^{}2 সমীকরণে x = 3 বসালে
\textquotesingle y\textquotesingle{} এর মান কত?}

ক) 8

\textbf{\hl{খ) 9 ✔}}

গ) 27

ঘ) 15

\textbf{\hl{✔ সঠিক উত্তর: খ) 9}}

\begin{quote}
\textbf{ব্যাখ্যা:} \emph{y = x\^{}2}

\emph{x = 3 বসালে পাই, y = (3)\^{}2}

\emph{∴ y = 9}
\end{quote}

\textbf{(x − 3)\^{}2 = 0 হলে, x-এর মান কত?}

\textbf{\hl{ক) 3 ✔}}

খ) 4

গ) 5

ঘ) 6

\textbf{\hl{✔ সঠিক উত্তর: ক) 3}}

\begin{quote}
\textbf{ব্যাখ্যা:} \emph{(x − 3)\^{}2 = 0}

\emph{⇒ x − 3 = 0}

\emph{∴ x = 3}
\end{quote}

\textbf{যদি x√(3)+3=4 হয়, তবে x এর কোন মানটি সঠিক?}

ক) 1

খ) (1/3)

\textbf{\hl{গ) (1/√(3)) ✔}}

ঘ) √(3)

\textbf{\hl{✔ সঠিক উত্তর: গ) (1/√(3))}}

\begin{quote}
\textbf{ব্যাখ্যা:} \emph{x√(3)+3=4}

\emph{⇒ x√(3)=4−3}

\emph{⇒ x√(3)=1}

\emph{∴ x=(1/√(3))}
\end{quote}

\textbf{x+(1/x)=2 হলে x এর মান কত?}

ক) 2

\textbf{\hl{খ) 1 ✔}}

গ) (1/2)

ঘ) (1/8)

\textbf{\hl{✔ সঠিক উত্তর: খ) 1}}

\begin{quote}
\textbf{ব্যাখ্যা:} \emph{x+(1/x)=2}

\emph{⇒ (x\^{}2+1/x)=2}

\emph{⇒ x\^{}2 + 1 = 2x}

\emph{⇒ x\^{}2 − 2x + 1 = 0}

\emph{⇒ (x − 1)\^{}2 = 0}

\emph{⇒ x − 1 = 0}

\emph{∴ x = 1}
\end{quote}

\textbf{x এর মান কত হলে x\^{}2 এবং x\^{}3 এর মান সমান হবে?}

\textbf{\hl{ক) 1 ✔}}

খ) 0

গ) 2

ঘ) 3

\textbf{\hl{✔ সঠিক উত্তর: ক) 1}}

\begin{quote}
\textbf{ব্যাখ্যা:} \emph{শর্তমতে, x\^{}2 = x\^{}3}

\emph{⇒ (x\^{}3/x\^{}2)=1}

\emph{⇒ x\^{}3−2 = 1}

\emph{∴ x = 1}

\emph{অর্থাৎ x = 1 হলে x\^{}2 এবং x\^{}3 এর মান সমান হবে।}
\end{quote}

\textbf{x-এর মান কত হলে a(x − a) = b(x − b) হবে?}

ক) a

খ) b − a

গ) a − b

\textbf{\hl{ঘ) a + b ✔}}

\textbf{\hl{✔ সঠিক উত্তর: ঘ) a + b}}

\begin{quote}
\textbf{ব্যাখ্যা:} \emph{a(x − a) = b(x − b)}

\emph{⇒ ax − a\^{}2 = bx − b\^{}2}

\emph{⇒ ax − bx = a\^{}2 − b\^{}2}

\emph{⇒ x(a − b) = (a + b)(a − b)}

\emph{∴ x = a + b}
\end{quote}

\textbf{x-এর মান কত হলে (3 + x) + 3(x + 3) = 0 হবে?}

\textbf{\hl{ক) −3 ✔}}

খ) 3

গ) 2

ঘ) −2

\textbf{\hl{✔ সঠিক উত্তর: ক) −3}}

\begin{quote}
\textbf{ব্যাখ্যা:} \emph{(3 + x) + 3(x + 3) = 0}

\emph{⇒ 3 + x + 3x + 9 = 0}

\emph{⇒ 4x + 12 = 0}

\emph{⇒ 4x = −12}

\emph{∴ x = −3}
\end{quote}

\textbf{(5 + x) + 4 = 2(x + 1) হলে x-এর মান কত?}

ক) 4

খ) 5

গ) 2

\textbf{\hl{ঘ) 3 ✔}}

\textbf{\hl{✔ সঠিক উত্তর: ঘ) 3}}

\begin{quote}
\textbf{ব্যাখ্যা:} \emph{(5 + x) + 4 = 2(x + 1)}

\emph{⇒ 9 + x = 2x + 2}

\emph{⇒ x − 2x = 2 − 9}

\emph{⇒ −x = −7}

\emph{∴ x = 7}
\end{quote}

\textbf{(x/a)+a=(x/b)+b হলে, x এর মান কত?}

\textbf{\hl{ক) ab ✔}}

খ) a

গ) b

ঘ) (a/b)

\textbf{\hl{✔ সঠিক উত্তর: ক) ab}}

\begin{quote}
\textbf{ব্যাখ্যা:} \emph{(x/a)+a=(x/b)+b}

\emph{⇒ (x/a)−(x/b)=−a+b}

\emph{⇒ x((1/a)−(1/b))=b−a}

\emph{⇒ x((b−a/ab))=b−a}

\emph{∴ x = ab}
\end{quote}

\textbf{(2 + x) + 3 = 3(x + 2) হলে x এর মান কত?}

\textbf{\hl{ক) −(1/2) ✔}}

খ) (1/2)

গ) (1/3)

ঘ) (2/3)

\textbf{\hl{✔ সঠিক উত্তর: ক) −(1/2)}}

\begin{quote}
\textbf{ব্যাখ্যা:} \emph{(2 + x) + 3 = 3(x + 2)}

\emph{⇒ 2 + x + 3 = 3x + 6}

\emph{⇒ x + 5 = 3x + 6}

\emph{⇒ x − 3x = 6 − 5}

\emph{⇒ −2x = 1}

\emph{∴ x=−(1/2)}
\end{quote}

\textbf{(3x + 2)(2x − 6) = (4 − 3x)(1 − 2x) − 10 হলে, x-এর মান হবে}

\textbf{\hl{ক) −2 ✔}}

খ) 4

গ) −3

ঘ) 5

\textbf{\hl{✔ সঠিক উত্তর: ক) −2}}

\begin{quote}
\textbf{ব্যাখ্যা:} \emph{LHS = (3x + 2)(2x − 6) = 6x\^{}2 − 18x + 4x − 12
= 6x\^{}2 − 14x − 12}

\emph{RHS = (4 − 3x)(1 − 2x) − 10 = 4 − 8x − 3x + 6x\^{}2 − 10 = 6x\^{}2
− 11x − 6}

\emph{শর্তমতে, 6x\^{}2 − 14x − 12 = 6x\^{}2 − 11x − 6}

\emph{⇒ −14x + 11x = −6 + 12}

\emph{⇒ −3x = 6}

\emph{∴ x = −2}
\end{quote}

\textbf{(a + b) = 4 হলে (a + b)\^{}3 কত?}

ক) ৩

খ) ৯

\textbf{\hl{গ) ৬৪ ✔}}

ঘ) ৭২

\textbf{\hl{✔ সঠিক উত্তর: গ) ৬৪}}

\begin{quote}
\textbf{ব্যাখ্যা:} \emph{(a + b) = 4}

\emph{⇒ (a + b)\^{}3 = (4)\^{}3}

\emph{∴ (a + b)\^{}3 = 64}
\end{quote}

\textbf{যদি x+(1/x)=5 হয়, তবে (x/x\^{}2+x+1) এর মান কত?}

ক) (1/5)

খ) (1/7)

\textbf{\hl{গ) (1/6) ✔}}

ঘ) (1/4)

\textbf{\hl{✔ সঠিক উত্তর: গ) (1/6)}}

\begin{quote}
\textbf{ব্যাখ্যা:} \emph{x+(1/x)=5}

\emph{⇒ x\^{}2 + 1 = 5x}

\emph{প্রদত্ত রাশি = (x/x\^{}2+x+1)}

\emph{= (x/5x+x)}

\emph{= (x/6x)}

\emph{= (1/6)}
\end{quote}

\textbf{যদি x+(1/x)=4 হলে (x/x\^{}2−3x+1) এর মান কত?}

ক) 4

খ) 3

গ) 2

\textbf{\hl{ঘ) 1 ✔}}

\textbf{\hl{✔ সঠিক উত্তর: ঘ) 1}}

\begin{quote}
\textbf{ব্যাখ্যা:} \emph{x+(1/x)=4}

\emph{⇒ (x\^{}2+1/x)=4}

\emph{⇒ x\^{}2 + 1 = 4x}

\emph{প্রদত্ত রাশি = (x/x\^{}2−3x+1)}

\emph{= (x/4x−3x)}

\emph{= (x/x) = 1}
\end{quote}

\textbf{যদি (a+1/a)=4 তবে (a/a\^{}2−3a+1) এর মান কত?}

\textbf{\hl{ক) (1/3) ✔}}

খ) (1/4)

গ) 1

ঘ) 3

\textbf{\hl{✔ সঠিক উত্তর: ক) (1/3)}}

\begin{quote}
\textbf{ব্যাখ্যা:} \emph{(a+1/a)=4}

\emph{⇒ 1+(1/a)=4}

\emph{⇒ a+(1/a)=4 (বা a = 1/3)}

\emph{মান বসিয়ে পাই, প্রদত্ত রাশি = (1/3)}
\end{quote}

\textbf{যদি (a+1/a)=4 তবে (a/a\^{}2−3a+1) এর মান কত?}

ক) (1/3)

খ) (1/4)

গ) 1

\textbf{\hl{ঘ) 3 ✔}}

\textbf{\hl{✔ সঠিক উত্তর: ঘ) 3}}

\begin{quote}
\textbf{ব্যাখ্যা:} \emph{(a+1/a)=4}

\emph{⇒ a+1=4a}

\emph{⇒ 4a−a=1}

\emph{⇒ 3a=1}

\emph{⇒ a=(1/3)}

\emph{এখন প্রদত্ত রাশি = (a/a\^{}2−3a+1)}

\emph{মান বসিয়ে পাই:}

\emph{= ((1/3)/((1/3))\^{}2−3((1/3))+1)}

\emph{= ((1/3)/(1/9)−1+1)}

\emph{= ((1/3)/(1/9))}

\emph{= (1/3)×9}

\emph{∴ 3}
\end{quote}

\textbf{সমাধান করুন: (x−2/4)+(1/3)=x−(2x−1/3)}

\textbf{\hl{ক) −6 ✔}}

খ) 6

গ) 4

ঘ) −4

\textbf{\hl{✔ সঠিক উত্তর: ক) −6}}

\begin{quote}
\textbf{ব্যাখ্যা:} \emph{(x−2/4)+(1/3)=x−(2x−1/3)}

\emph{⇒ (3x−6+4/12)=(3x−(2x−1)/3)}

\emph{⇒ (3x−2/12)=(3x−2x+1/3)}

\emph{⇒ (3x−2/4)=x+1}

\emph{⇒ 3x − 2 = 4x + 4}

\emph{⇒ 3x − 4x = 4 + 2}

\emph{⇒ −x = 6}

\emph{∴ x = −6}
\end{quote}

\textbf{সমাধান করুন : (x/5)−(2/7)=(5x/7)−(4/5)}

ক) 0

\textbf{\hl{খ) 1 ✔}}

গ) 2

ঘ) 3

\textbf{\hl{✔ সঠিক উত্তর: খ) 1}}

\begin{quote}
\textbf{ব্যাখ্যা:} \emph{(x/5)−(2/7)=(5x/7)−(4/5)}

\emph{⇒ (x/5)−(5x/7)=(2/7)−(4/5)}

\emph{⇒ (7x−25x/35)=(10−28/35)}

\emph{⇒ −18x = −18}

\emph{∴ x = 1}
\end{quote}

\textbf{x + y = 9 এবং (x/y)=2 হলে x এর মান কত?}

ক) 4

খ) 5

\textbf{\hl{গ) 6 ✔}}

ঘ) 7

\textbf{\hl{✔ সঠিক উত্তর: গ) 6}}

\begin{quote}
\textbf{ব্যাখ্যা:} \emph{দেওয়া আছে, x + y = 9 ......... (i)}

\emph{এবং (x/y)=2 ⇒ x = 2y ......... (ii)}

\emph{(i) নং সমীকরণে x এর মান বসিয়ে পাই, 2y + y = 9}

\emph{⇒ 3y = 9 ⇒ y = 3}

\emph{এখন y এর মান (ii) নং এ বসিয়ে পাই, x = 2 × 3}

\emph{∴ x = 6}
\end{quote}

\textbf{x+(1/x)=2 হলে, x\^{}5+(1/x\^{}5) কত?}

\textbf{\hl{ক) 2 ✔}}

খ) 4

গ) 27

ঘ) 64

\textbf{\hl{✔ সঠিক উত্তর: ক) 2}}

\begin{quote}
\textbf{ব্যাখ্যা:} \emph{x+(1/x)=2}

\emph{⇒ x\^{}2 − 2x + 1 = 0}

\emph{⇒ (x − 1)\^{}2 = 0}

\emph{⇒ x − 1 = 0}

\emph{⇒ x = 1}

\emph{∴ x\^{}5+(1/x\^{}5) = (1)\^{}5 + (1/(1)\^{}5) = 1 + 1 = 2}
\end{quote}

\textbf{a−(6/a)=1 হলে, (6/a\^{}2−a+1) এর মান কত??}

ক) (7/6)

\textbf{\hl{খ) (6/7) ✔}}

গ) (5/7)

ঘ) (4/7)

\textbf{\hl{✔ সঠিক উত্তর: খ) (6/7)}}

\begin{quote}
\textbf{ব্যাখ্যা:} \emph{a−(6/a)=1}

\emph{⇒ a\^{}2 − 6 = a}

\emph{⇒ a\^{}2 − a = 6}

\emph{প্রদত্ত রাশি = (6/a\^{}2−a+1)}

\emph{= (6/6+1)}

\emph{= (6/7)}
\end{quote}

\textbf{x+(1/x)=5 এবং x≠0 হলে, (x/x\^{}2+x+1) এর মান কত?}

ক) 6

\textbf{\hl{খ) (1/6) ✔}}

গ) √(6)

ঘ) (1/√(6))

\textbf{\hl{✔ সঠিক উত্তর: খ) (1/6)}}

\begin{quote}
\textbf{ব্যাখ্যা:} \emph{x+(1/x)=5}

\emph{⇒ x\^{}2 + 1 = 5x}

\emph{⇒ x\^{}2 + x + 1 = 5x + x = 6x}

\emph{এখন, (x/x\^{}2+x+1) = (x/6x) = (1/6)}
\end{quote}

\textbf{x+(1/x)=2 হলে (x/x\^{}2+x−1) এর মান কত?}

ক) 2

খ) 3

\textbf{\hl{গ) 1 ✔}}

ঘ) 4

\textbf{\hl{✔ সঠিক উত্তর: গ) 1}}

\begin{quote}
\textbf{ব্যাখ্যা:} \emph{x+(1/x)=2}

\emph{⇒ x\^{}2 + 1 = 2x}

\emph{⇒ x\^{}2 − 2x + 1 = 0}

\emph{⇒ (x − 1)\^{}2 = 0}

\emph{⇒ x − 1 = 0 ∴ x = 1}

\emph{এখানে (x/x\^{}2+x−1) = (1/1\^{}2+1−1) = (1/1) = 1}
\end{quote}

\textbf{b+b\^{}−1=2 হলে, (2b/b\^{}2−b+1) এর মান---}

\textbf{\hl{ক) 2 ✔}}

খ) 1

গ) 3

ঘ) 5

\textbf{\hl{✔ সঠিক উত্তর: ক) 2}}

\begin{quote}
\textbf{ব্যাখ্যা:} \emph{b+b\^{}−1=2}

\emph{⇒ b+(1/b)=2}

\emph{⇒ (b\^{}2+1/b)=2}

\emph{⇒ b\^{}2 + 1 = 2b}

\emph{⇒ b\^{}2 − b + 1 = b}

\emph{এখন প্রদত্ত রাশি = (2b/b\^{}2−b+1) = (2b/b) = 2}
\end{quote}

\textbf{যদি (x/y)=(3/4) হয়, তবে (x+y/x−y) এর মান কত?}

ক) 2

খ) 5

গ) 6

\textbf{\hl{ঘ) 7 ✔}}

\textbf{\hl{✔ সঠিক উত্তর: ঘ) 7}}

\begin{quote}
\textbf{ব্যাখ্যা:} \emph{(x/y)=(3/4)}

\emph{⇒ যোজন-বিয়োজন করে পাই, (x+y/x−y)=(3+4/3−4)}

\emph{= (7/−1) = −7}
\end{quote}

\textbf{(x/y)=(2/3) হলে, (6x+y/3x+2y) = কত?}

ক) 5

খ) 6

\textbf{\hl{গ) (5/4) ✔}}

ঘ) (3/4)

\textbf{\hl{✔ সঠিক উত্তর: গ) (5/4)}}

\begin{quote}
\textbf{ব্যাখ্যা:} \emph{(x/y)=(2/3) ⇒ 3x = 2y}

\emph{প্রদত্ত রাশি = (6x+y/3x+2y)}

\emph{= (2(3x)+y/2y+2y) (3x = 2y বসিয়ে)}

\emph{= (2(2y)+y/4y)}

\emph{= (4y+y/4y)}

\emph{= (5y/4y) = (5/4)}
\end{quote}

\textbf{যদি x + 6y = 24 এবং x = −2y তবে y = ?}

\textbf{\hl{ক) ৬ ✔}}

খ) ৪

গ) ৩

ঘ) ৮

\textbf{\hl{✔ সঠিক উত্তর: ক) ৬}}

\begin{quote}
\textbf{ব্যাখ্যা:} \emph{দেওয়া আছে, x + 6y = 24}

\emph{⇒ −2y + 6y = 24}

\emph{⇒ 4y = 24}

\emph{∴ y = (24/4) = 6}
\end{quote}

\textbf{2x = 3y + 5 হলে 4x − 6y = কত?}

ক) 12

\textbf{\hl{খ) 10 ✔}}

গ) 15

ঘ) 20

\textbf{\hl{✔ সঠিক উত্তর: খ) 10}}

\begin{quote}
\textbf{ব্যাখ্যা:} \emph{দেওয়া আছে, 2x = 3y + 5}

\emph{⇒ 2x − 3y = 5}

\emph{∴ 4x − 6y = 2(2x − 3y) = 2 × 5 = 10}
\end{quote}

\textbf{x\^{}2 − 4x − 1 = 0 হলে x−(1/x) এর মান কত?}

\textbf{\hl{ক) 4 ✔}}

খ) 3

গ) 2

ঘ) 1

\textbf{\hl{✔ সঠিক উত্তর: ক) 4}}

\begin{quote}
\textbf{ব্যাখ্যা:} \emph{দেওয়া আছে, x\^{}2 − 4x − 1 = 0}

\emph{⇒ x\^{}2 − 1 = 4x}

\emph{উভয়পক্ষকে x দ্বারা ভাগ করে পাই:}

\emph{⇒ (x\^{}2−1/x)=4}

\emph{∴ x−(1/x) = 4}
\end{quote}

\textbf{x = 1 হলে (x\^{}2−1/x\^{}4) এর মান কত?}

ক) (1/4)

খ) (3/4)

গ) 1

\textbf{\hl{ঘ) 0 ✔}}

\textbf{\hl{✔ সঠিক উত্তর: ঘ) 0}}

\begin{quote}
\textbf{ব্যাখ্যা:} \emph{দেওয়া আছে, (x\^{}2−1/x\^{}4)}

\emph{x = 1 বসিলে পাই,}

\emph{= (1\^{}2−1/1\^{}4)}

\emph{= (1−1/1)}

\emph{= (0/1) = 0}
\end{quote}

\textbf{x\^{}4 + 2x\^{}2 + 1 = 5x\^{}2 হলে, x+(1/x) কত?}

ক) √(3)

খ) 2

গ) 5

\textbf{\hl{ঘ) √(5) ✔}}

\textbf{\hl{✔ সঠিক উত্তর: ঘ) √(5)}}

\begin{quote}
\textbf{ব্যাখ্যা:} \emph{x\^{}4 + 2x\^{}2 + 1 = 5x\^{}2}

\emph{⇒ (x\^{}4+2x\^{}2+1/x\^{}2)=5}

\emph{⇒ x\^{}2+2+(1/x\^{}2)=5}

\emph{⇒ x\^{}2+(1/x\^{}2)=5−2}

\emph{⇒ x\^{}2+(1/x\^{}2)=3}

\emph{⇒ (x+(1/x))\^{}2−2·x·(1/x)=3}

\emph{⇒ (x+(1/x))\^{}2=3+2}

\emph{∴ x+(1/x)=√(5)}
\end{quote}

\textbf{সমাধান কর:: (3x/4)+(5x−2/6)=(4x+5/8)}

\textbf{\hl{ক) (23/26) ✔}}

খ) (26/23)

গ) (13/12)

ঘ) (12/13)

\textbf{\hl{✔ সঠিক উত্তর: ক) (23/26)}}

\begin{quote}
\textbf{ব্যাখ্যা:} \emph{(3x/4)+(5x−2/6)=(4x+5/8)}

\emph{⇒ (9x+10x−4/12)=(4x+5/8)}

\emph{⇒ (19x−4/12)=(4x+5/8)}

\emph{⇒ 2(19x−4)=3(4x+5)}

\emph{⇒ 38x−8=12x+15}

\emph{⇒ 38x−12x=15+8}

\emph{⇒ 26x=23}

\emph{∴ x=(23/26)}
\end{quote}

\textbf{(x−a/b+c)+(x−b/c+a)+(x−c/a+b)=3}

\textbf{\hl{ক) a + b + c ✔}}

খ) a − b − c

গ) 0

ঘ) 1

\textbf{\hl{✔ সঠিক উত্তর: ক) a + b + c}}

\begin{quote}
\textbf{ব্যাখ্যা:} \emph{(x−a/b+c)+(x−b/c+a)+(x−c/a+)=3}

\emph{⇒ ((x−a/b+c)−1)+((x−b/c+a)−1)+((x−c/a+b)−1)=0}

\emph{⇒ (x−a−b−c/b+c)+(x−a−b−c/c+a)+(x−a−b−c/a+b)=0}

\emph{⇒ (x−a−b−c)((1/b+c)+(1/c+a)+(1/a+b))=0}

\emph{যেহেতু বন্ধনীর ভেতরের অংশটি শূন্য নয়, তাই ∴ x − (a + b + c) = 0 ⇒ x =
a + b + c}
\end{quote}

\textbf{a=(1/2) হলে, (2a + 1)(4a\^{}2 − 2a + 1) এর মান নিচের কোনটি?}

ক) 0

খ) 1

\textbf{\hl{গ) 2 ✔}}

ঘ) 3

\textbf{\hl{✔ সঠিক উত্তর: গ) 2}}

\begin{quote}
\textbf{ব্যাখ্যা:} \emph{আমরা জানি, (2a + 1)(4a\^{}2 − 2a + 1) = (2a)\^{}3
+ (1)\^{}3 = 8a\^{}3 + 1}

\emph{এখন a=(1/2) বসিলে পাই,}

\emph{= 8((1/2))\^{}3+1}

\emph{= 8×(1/8)+1 = 1 + 1 = 2}
\end{quote}

\textbf{(2x−3y/2x+3y)=(1/7) হলে x : y = কত?}

ক) 2 : 3

খ) 3 : 2

\textbf{\hl{গ) 2 : 1 ✔}}

ঘ) 1 : 2

\textbf{\hl{✔ সঠিক উত্তর: গ) 2 : 1}}

\begin{quote}
\textbf{ব্যাখ্যা:} \emph{(2x−3y/2x+3y)=(1/7)}

\emph{⇒ 14x − 21y = 2x + 3y}

\emph{⇒ 14x − 2x = 21y + 3y}

\emph{⇒ 12x = 24y}

\emph{⇒ (x/y)=(24/12)}

\emph{∴ x : y = 2 : 1}
\end{quote}

\textbf{(2x+3y/3x+2y)=(5/6) হলে, x : y = কত?}

ক) 5 : 6

খ) 3 : 8

গ) 6 : 5

\textbf{\hl{ঘ) 8 : 3 ✔}}

\textbf{\hl{✔ সঠিক উত্তর: ঘ) 8 : 3}}

\begin{quote}
\textbf{ব্যাখ্যা:} \emph{(2x+3y/3x+2y)=(5/6)}

\emph{⇒ 12x + 18y = 15x + 10y}

\emph{⇒ 15x − 12x = 18y − 10y}

\emph{⇒ 3x = 8y}

\emph{⇒ (x/y)=(8/3)}

\emph{∴ x : y = 8 : 3}
\end{quote}

\textbf{x : y = 5 : 6 হলে 3x : 5y =?}

ক) 3 : 5

\textbf{\hl{খ) 1 : 2 ✔}}

গ) 6 : 5

ঘ) 15 : 30

\textbf{\hl{✔ সঠিক উত্তর: খ) 1 : 2}}

\begin{quote}
\textbf{ব্যাখ্যা:} \emph{দেওয়া আছে, x : y = 5 : 6}

\emph{⇒ (x/y)=(5/6)}

\emph{এখন, (3x/5y)=(3/5)×(x/y)}

\emph{= (3/5)×(5/6)}

\emph{= (3/6)=(1/2)}

\emph{∴ 3x : 5y = 1 : 2}
\end{quote}

\textbf{x−4=(x−4/x) সমীকরণে সমাধান---}

ক) x = 1

খ) x = 4

\textbf{\hl{গ) x = 1, x = 4 ✔}}

ঘ) কোনোটিই নয়

\textbf{\hl{✔ সঠিক উত্তর: গ) x = 1, x = 4}}

\begin{quote}
\textbf{ব্যাখ্যা:} \emph{x−4=(x−4/x)}

\emph{⇒ x(x−4)=x−4}

\emph{⇒ x(x−4)−(x−4)=0}

\emph{⇒ (x−4)(x−1)=0}

\emph{হয়, x − 4 = 0 ⇒ x = 4}

\emph{অথবা, x − 1 = 0 ⇒ x = 1}

\emph{∴ x = 1, 4}
\end{quote}

\textbf{((5x/6)+3) এবং ((x/3)+10) পরস্পর সমান হলে x এর মান কত?}

ক) (21/2)

\textbf{\hl{খ) 14 ✔}}

গ) 6

ঘ) 7

\textbf{\hl{✔ সঠিক উত্তর: খ) 14}}

\begin{quote}
\textbf{ব্যাখ্যা:} \emph{(5x/6)+3=(x/3)+10}

\emph{⇒ (5x/6)−(x/3)=10−3}

\emph{⇒ (5x−2x/6)=7}

\emph{⇒ (3x/6)=7}

\emph{⇒ (x/2)=7}

\emph{∴ x = 14}
\end{quote}

\textbf{যদি x\^{}3 + hx + 10 = 0 এর একটি সমাধান 2 হয়, তবে h এর মান কত?}

ক) 10

খ) 9

\textbf{\hl{গ) −9 ✔}}

ঘ) 2

\textbf{\hl{✔ সঠিক উত্তর: গ) −9}}

\begin{quote}
\textbf{ব্যাখ্যা:} \emph{যেহেতু x = 2, তাই x\^{}3 + hx + 10 = 0 সমীকরণে x
এর মান বসিয়ে পাই,}

\emph{2\^{}3 + h(2) + 10 = 0}

\emph{⇒ 8 + 2h + 10 = 0}

\emph{⇒ 2h + 18 = 0}

\emph{⇒ 2h = −18}

\emph{∴ h = −9}
\end{quote}

\textbf{(√(5)+√(5−x)/√(5)−√(5−x))=5 হলে x এর মান কত?}

ক) (20/9)

\textbf{\hl{খ) (25/9) ✔}}

গ) (22/9)

ঘ) (15/9)

\textbf{\hl{✔ সঠিক উত্তর: খ) (25/9)}}

\begin{quote}
\textbf{ব্যাখ্যা:} \emph{(√(5)+√(5−x)/√(5)−√(5−x))=5}

\emph{⇒ যোজন-বিয়োজন করে পাই, (2√(5)/2√(5−x))=(5+1/5−1)}

\emph{⇒ (√(5)/√(5−x))=(6/4)=(3/2)}

\emph{⇒ উভয়পক্ষে বর্গ করে, (5/5−x)=(9/4)}

\emph{⇒ 45 − 9x = 20}

\emph{⇒ 9x = 45 − 20 = 25}

\emph{∴ x=(25/9)}
\end{quote}

\textbf{x\^{}2 = x√3 এর সমাধান সেট কোনটি?}

ক) \{3\}

খ) \{0,(1/√(3))\}

গ) \{0, 3\}

\textbf{\hl{ঘ) \{0,√(3)\} ✔}}

\textbf{\hl{✔ সঠিক উত্তর: ঘ) \{0,√(3)\}}}

\begin{quote}
\textbf{ব্যাখ্যা:} \emph{দেওয়া আছে, x\^{}2 = x√3}

\emph{⇒ x\^{}2 − x√3 = 0}

\emph{⇒ x(x − √3) = 0}

\emph{হয়, x = 0}

\emph{অথবা, x − √3 = 0 ⇒ x=√(3)}

\emph{∴ সমাধান সেট \{0,√(3)\}}
\end{quote}

\textbf{\textbar x+1\textbar2+3\textbar x+1\textbar−4=0 সমীকরণটির সমাধান
সেট নির্ণয় করুন:}

ক) \{−(1/2),(1/2)\}

খ) \{−1, 1\}

গ) \{−3, 1\}

\textbf{\hl{ঘ) \{0, −5\} ✔}}

\textbf{\hl{✔ সঠিক উত্তর: ঘ) \{0, −5\}}}

\begin{quote}
\textbf{ব্যাখ্যা:} \emph{ধরি, \textbar x + 1\textbar{} = P}

\emph{তাহলে, P\^{}2 + 3P − 4 = 0}

\emph{⇒ P\^{}2 + 4P − P − 4 = 0}

\emph{⇒ (P + 4)(P − 1) = 0}

\emph{হয়, P + 4 = 0 ⇒ \textbar x + 1\textbar{} = −4 (যা অসম্ভব, কারণ পরম
মানের মান ঋণাত্মক হতে পারে না)}

\emph{অথবা, P − 1 = 0 ⇒ \textbar x + 1\textbar{} = 1}

\emph{⇒ x + 1 = ±1}

\emph{যখন x + 1 = 1 তখন x = 0}

\emph{যখন x + 1 = −1 তখন x = −5}

\emph{∴ সমাধান সেট \{0, −5\}}
\end{quote}

\section{25. সরল সহ-সমীকরণ MCQ প্রশ্ন উত্তর
অনুশীলন}\label{ux9b8ux9b0ux9b2-ux9b8ux9b9-ux9b8ux9aeux995ux9b0ux9a3-mcq-ux9aaux9b0ux9b6ux9a8-ux989ux9a4ux9a4ux9b0-ux985ux9a8ux9b6ux9b2ux9a8}

\textbf{2x + y = 7, 2x − y = 13 হলে, x ও y এর মান কত?}

ক) 5, 3

খ) −5, 3

\textbf{\hl{গ) 5, −3 ✔}}

ঘ) −5, −3

\textbf{\hl{✔ সঠিক উত্তর: গ) 5, −3}}

\begin{quote}
\textbf{ব্যাখ্যা:} \emph{দেওয়া আছে, 2x + y = 7 ......... (i)}

\emph{এবং 2x − y = 13 ......... (ii)}

\emph{(i) ও (ii) যোগ করে পাই,}

\emph{4x = 20}

\emph{⇒ x = (20/4) = 5}

\emph{x এর মান (i) নং সমীকরণে বসিয়ে পাই,}

\emph{2(5) + y = 7}

\emph{⇒ 10 + y = 7}

\emph{⇒ y = 7 − 10 = −3}

\emph{∴ (x, y) = (5, −3)}
\end{quote}

\textbf{x − 2y = 8, 3x − 2y = 4 সমীকরণ জোটে x এর মান কত?}

ক) −5

\textbf{\hl{খ) −2 ✔}}

গ) 2

ঘ) 5

\textbf{\hl{✔ সঠিক উত্তর: খ) −2}}

\begin{quote}
\textbf{ব্যাখ্যা:} \emph{দেওয়া আছে,}

\emph{x − 2y = 8 ......... (i)}

\emph{3x − 2y = 4 ......... (ii)}

\emph{(i) নং সমীকরণ হতে (ii) নং সমীকরণ বিয়োগ করে পাই,}

\emph{(x − 3x) = 8 − 4}

\emph{⇒ −2x = 4}

\emph{∴ x = −2}
\end{quote}

\textbf{a + b = a − b হলে b এর মান কত?}

ক) −1

\textbf{\hl{খ) 0 ✔}}

গ) 0.5

ঘ) 2

\textbf{\hl{✔ সঠিক উত্তর: খ) 0}}

\begin{quote}
\textbf{ব্যাখ্যা:} \emph{a + b = a − b}

\emph{⇒ b + b = a − a}

\emph{⇒ 2b = 0}

\emph{∴ b = 0}
\end{quote}

\textbf{যদি a + b = 2, ab = 1 হয়, তবে a এবং b এর মান যথাক্রমে ---}

ক) 0, 2

\textbf{\hl{খ) 1, 1 ✔}}

গ) −1, 3

ঘ) −3, 4

\textbf{\hl{✔ সঠিক উত্তর: খ) 1, 1}}

\begin{quote}
\textbf{ব্যাখ্যা:} \emph{দেওয়া আছে, a + b = 2 এবং ab = 1}

\emph{আমরা জানি, (a − b)\^{}2 = (a + b)\^{}2 − 4ab}

\emph{= 2\^{}2 − 4(1) = 4 − 4 = 0}

\emph{⇒ a − b = 0}

\emph{এখন, a + b = 2 এবং a − b = 0 যোগ করে পাই, 2a = 2 ⇒ a = 1}

\emph{আবার, বিয়োগ করে পাই, 2b = 2 ⇒ b = 1}

\emph{∴ (a, b) = (1, 1)}
\end{quote}

\textbf{x − y = 2 এবং xy = 24 হলে, x এর ধনাত্মক মানটি---}

ক) 3

খ) 4

গ) 5

\textbf{\hl{ঘ) 6 ✔}}

\textbf{\hl{✔ সঠিক উত্তর: ঘ) 6}}

\begin{quote}
\textbf{ব্যাখ্যা:} \emph{দেওয়া আছে, x − y = 2 ⇒ x = y + 2}

\emph{এবং xy = 24}

\emph{মান বসিয়ে পাই, (y + 2)y = 24}

\emph{⇒ y\^{}2 + 2y − 24 = 0}

\emph{⇒ y\^{}2 + 6y − 4y − 24 = 0}

\emph{⇒ y(y + 6) − 4(y + 6) = 0}

\emph{⇒ (y − 4)(y + 6) = 0}

\emph{সুতরাং y = 4 অথবা y = −6}

\emph{যদি y = 4 হয়, তবে x = 4 + 2 = 6}

\emph{যদি y = −6 হয়, তবে x = −6 + 2 = −4}

\emph{∴ x এর ধনাত্মক মানটি 6}
\end{quote}

\textbf{x + y = 12 এবং x − y = 2 হলে xy এর মান কত?}

ক) 30

খ) 49

\textbf{\hl{গ) 35 ✔}}

ঘ) 45

\textbf{\hl{✔ সঠিক উত্তর: গ) 35}}

\begin{quote}
\textbf{ব্যাখ্যা:} \emph{আমরা জানি, 4xy = (x + y)\^{}2 − (x − y)\^{}2}

\emph{= (12)\^{}2 − (2)\^{}2}

\emph{= 144 − 4 = 140}

\emph{⇒ xy=(140/4)}

\emph{∴ xy = 35}
\end{quote}

\textbf{যদি (x − y, 3) = (0, x + 2y) হয়, তবে (x, y) = কত?}

\textbf{\hl{ক) (1, 1) ✔}}

খ) (1, −1)

গ) (1, 3)

ঘ) (3, 1)

\textbf{\hl{✔ সঠিক উত্তর: ক) (1, 1)}}

\begin{quote}
\textbf{ব্যাখ্যা:} \emph{ক্রমজোড়ের শর্তানুযায়ী,}

\emph{x − y = 0 ......... (i)}

\emph{এবং x + 2y = 3 ......... (ii)}

\emph{(i) নং হতে পাই, x = y}

\emph{এখন (ii) নং সমীকরণে x এর পরিবর্তে y বসিয়ে পাই,}

\emph{y + 2y = 3}

\emph{⇒ 3y = 3 ⇒ y = 1}

\emph{যেহেতু x = y, তাই x = 1}

\emph{∴ (x, y) = (1, 1)}
\end{quote}

\textbf{2x + 3y = 7, 5x − 2y − 8 = 0 হলে (x, y) = কত?}

ক) (1, 2)

\textbf{\hl{খ) (2, 1) ✔}}

গ) (2, 3)

ঘ) (3, 2)

\textbf{\hl{✔ সঠিক উত্তর: খ) (2, 1)}}

\begin{quote}
\textbf{ব্যাখ্যা:} \emph{দেওয়া আছে, 2x + 3y = 7 ......... (i)}

\emph{এবং 5x − 2y = 8 ......... (ii)}

\emph{(i) নং সমীকরণকে 2 দ্বারা এবং (ii) নং সমীকরণকে 3 দ্বারা গুণ করে পাই,}

\emph{4x + 6y = 14}

\emph{15x − 6y = 24}

\emph{যোগ করে পাই, 19x = 38 ⇒ x = 2}

\emph{x এর মান (i) নং সমীকরণে বসিয়ে পাই, 2(2) + 3y = 7 ⇒ 4 + 3y = 7 ⇒ 3y
= 3 ⇒ y = 1}

\emph{∴ (x, y) = (2, 1)}
\end{quote}

\textbf{x\^{}2 + y\^{}2 = 185, x − y = 3 এর একটি সমাধান হবে---}

ক) (7, 4)

খ) (9, 6)

গ) (10, 7)

\textbf{\hl{ঘ) (11, 8) ✔}}

\textbf{\hl{✔ সঠিক উত্তর: ঘ) (11, 8)}}

\begin{quote}
\textbf{ব্যাখ্যা:} \emph{অপশনগুলোর মধ্যে (11, 8) পরীক্ষা করে পাই:}

\emph{x − y = 11 − 8 = 3}

\emph{এবং x\^{}2 + y\^{}2 = 11\^{}2 + 8\^{}2 = 121 + 64 = 185}

\emph{যা সমীকরণদ্বয়কে সিদ্ধ করে।}

\emph{∴ সঠিক উত্তর (11, 8)}
\end{quote}

\textbf{x − y = 1, x + y = 3 হলে (x, y) = কোনটি?}

ক) (1, 2)

\textbf{\hl{খ) (2, 1) ✔}}

গ) (1, 3)

ঘ) (3, 1)

\textbf{\hl{✔ সঠিক উত্তর: খ) (2, 1)}}

\begin{quote}
\textbf{ব্যাখ্যা:} \emph{দেওয়া আছে,}

\emph{x − y = 1 ......... (i)}

\emph{x + y = 3 ......... (ii)}

\emph{(i) ও (ii) যোগ করে পাই, 2x = 4 ⇒ x = 2}

\emph{x এর মান (ii) নং সমীকরণে বসিয়ে পাই, 2 + y = 3 ⇒ y = 1}

\emph{∴ (x, y) = (2, 1)}
\end{quote}

\textbf{2x + 3y = 36, 2x + y = 16 হলে, (x, y) এর মান কত?}

ক) (2, 10)

\textbf{\hl{খ) (3, 10) ✔}}

গ) (3, 5)

ঘ) (6, 10)

\textbf{\hl{✔ সঠিক উত্তর: খ) (3, 10)}}

\begin{quote}
\textbf{ব্যাখ্যা:} \emph{দেওয়া আছে,}

\emph{2x + 3y = 36 ......... (1)}

\emph{2x + y = 16 ......... (2)}

\emph{(1) নং হতে (2) নং বিয়োগ করে পাই,}

\emph{2y = 20}

\emph{⇒ y = 10}

\emph{y এর মান (2) নং সমীকরণে বসিয়ে পাই,}

\emph{2x + 10 = 16}

\emph{⇒ 2x = 6}

\emph{∴ x = 3}

\emph{সুতরাং (x, y) = (3, 10)}
\end{quote}

\textbf{(x − y, 3) = (0, x + 2y) হলে (x, y) = ?}

ক) (−1, −1)

খ) (1, 3)

গ) (1, −3)

\textbf{\hl{ঘ) (1, 1) ✔}}

\textbf{\hl{✔ সঠিক উত্তর: ঘ) (1, 1)}}

\begin{quote}
\textbf{ব্যাখ্যা:} \emph{ক্রমজোড়ের শর্তানুযায়ী,}

\emph{x − y = 0 ......... (i)}

\emph{x + 2y = 3 ......... (ii)}

\emph{(ii) নং হতে (i) নং বিয়োগ করে পাই,}

\emph{3y = 3 ⇒ y = 1}

\emph{y এর মান (i) নং এ বসিয়ে পাই, x − 1 = 0 ⇒ x = 1}

\emph{∴ (x, y) = (1, 1)}
\end{quote}

\textbf{If y = 2.5x + 4 and 15x − 2 = 40 then y =?}

ক) 36

খ) 28

গ) 22

\textbf{\hl{ঘ) 11 ✔}}

\textbf{\hl{✔ সঠিক উত্তর: ঘ) 11}}

\begin{quote}
\textbf{ব্যাখ্যা:} \emph{দেওয়া আছে, 15x − 2 = 40}

\emph{⇒ 15x = 40 + 2 = 42}

\emph{⇒ x=(42/15)}

\emph{এখন, y=(2.5×(42/15))+4}

\emph{= ((25/10)×(42/15))+4 = 7 + 4 = 11}

\emph{∴ y = 11}
\end{quote}

\textbf{যদি y = 5x + 4 এবং 5x + 8 = 40 হয়, তবে y = ?}

ক) 28

\textbf{\hl{খ) 36 ✔}}

গ) 42

ঘ) কোনোটিই নয়

\textbf{\hl{✔ সঠিক উত্তর: খ) 36}}

\begin{quote}
\textbf{ব্যাখ্যা:} \emph{দেওয়া আছে, 5x + 8 = 40}

\emph{⇒ 5x = 40 − 8 = 32}

\emph{⇒ x=(32/5)}

\emph{এখন, y=5((32/5))+4 = 32 + 4 = 36}

\emph{∴ y = 36}
\end{quote}

\textbf{x\^{}2 + y\^{}2 = 185, x − y = 3 এর একটি সমাধান হল---}

ক) (7, 4)

খ) (9, 6)

গ) (10, 7)

\textbf{\hl{ঘ) (11, 8) ✔}}

\textbf{\hl{✔ সঠিক উত্তর: ঘ) (11, 8)}}

\begin{quote}
\textbf{ব্যাখ্যা:} \emph{দেওয়া আছে, x − y = 3 এবং x\^{}2 + y\^{}2 = 185}

\emph{আমরা জানি, (x + y)\^{}2 = (x − y)\^{}2 + 4xy}

\emph{আবার, (x − y)\^{}2 + 2xy = 185 ⇒ 3\^{}2 + 2xy = 185 ⇒ 2xy = 176 ⇒
4xy = 352}

\emph{তাহলে, (x + y)\^{}2 = 9 + 352 = 361 ⇒ x + y = 19}

\emph{এখন x − y = 3 এবং x + y = 19 যোগ করে পাই, 2x = 22 ⇒ x = 11 এবং y =
8}

\emph{∴ (x, y) = (11, 8)}
\end{quote}

\textbf{(2x − 6, 5) = (4, 2y − 5) হলে, (x, y) এর মান কোনটি?}

ক) (4, 5)

\textbf{\hl{খ) (5, 5) ✔}}

গ) (6, 4)

ঘ) (6, 5)

\textbf{\hl{✔ সঠিক উত্তর: খ) (5, 5)}}

\begin{quote}
\textbf{ব্যাখ্যা:} \emph{ক্রমজোড়ের শর্তানুযায়ী,}

\emph{2x − 6 = 4 ⇒ 2x = 10 ⇒ x = 5}

\emph{এবং 2y − 5 = 5 ⇒ 2y = 10 ⇒ y = 5}

\emph{∴ (x, y) = (5, 5)}
\end{quote}

\textbf{যদি (x/y)=(1/3) এবং x + 2y = 5 হয়, তবে x এর মান কত?}

ক) 3

খ) 4

গ) 5

\textbf{\hl{ঘ) কোনোটিই নয় ✔}}

\textbf{\hl{✔ সঠিক উত্তর: ঘ) কোনোটিই নয়}}

\begin{quote}
\textbf{ব্যাখ্যা:} \emph{দেওয়া আছে, (x/y)=(1/3)}

\emph{⇒ y = 3x}

\emph{এখন, x + 2y = 5 সমীকরণে y এর মান বসিয়ে পাই,}

\emph{x + 2(3x) = 5}

\emph{⇒ x + 6x = 5}

\emph{⇒ 7x = 5}

\emph{∴ x=(5/7) (যা অপশনে নেই, তাই সঠিক উত্তর হবে কোনোটিই নয়)}
\end{quote}

\textbf{x + 2y = 4 এবং xy = 2 হয়, তবে x = কত?}

ক) 0

খ) 12

গ) 1

\textbf{\hl{ঘ) 2 ✔}}

\textbf{\hl{✔ সঠিক উত্তর: ঘ) 2}}

\begin{quote}
\textbf{ব্যাখ্যা:} \emph{দেওয়া আছে, x + 2y = 4 ......... (i)}

\emph{এবং xy = 2 ⇒ x=(2/y) ......... (ii)}

\emph{(i) নং সমীকরণে x এর মান বসিয়ে পাই,}

\emph{(2/y)+2y=4}

\emph{⇒ (2+2y\^{}2/y)=4}

\emph{⇒ 2 + 2y\^{}2 = 4y}

\emph{⇒ 2y\^{}2 − 4y + 2 = 0}

\emph{⇒ y\^{}2 − 2y + 1 = 0}

\emph{⇒ (y − 1)\^{}2 = 0}

\emph{∴ y = 1}

\emph{এখন y এর মান (ii) নং এ বসিয়ে পাই, x=(2/1)=2}
\end{quote}

\textbf{x + 2y = 4 এবং (x/y)=2 হলে, x এর মান কত?}

ক) 0

খ) 12

গ) 1

\textbf{\hl{ঘ) 2 ✔}}

\textbf{\hl{✔ সঠিক উত্তর: ঘ) 2}}

\begin{quote}
\textbf{ব্যাখ্যা:} \emph{দেওয়া আছে, (x/y)=2}

\emph{⇒ x = 2y}

\emph{এখন, x + 2y = 4 সমীকরণে x এর মান বসিয়ে পাই,}

\emph{2y + 2y = 4}

\emph{⇒ 4y = 4 ⇒ y = 1}

\emph{তাহলে, x = 2(1) = 2}

\emph{∴ x = 2}
\end{quote}

\textbf{x + y = 36 এবং x − y = 12 হলে x এর মান কত?}

ক) 18

খ) 09

গ) 16

\textbf{\hl{ঘ) 24 ✔}}

\textbf{\hl{✔ সঠিক উত্তর: ঘ) 24}}

\begin{quote}
\textbf{ব্যাখ্যা:} \emph{দেওয়া আছে,}

\emph{x + y = 36 ......... (1)}

\emph{x − y = 12 ......... (2)}

\emph{(1) ও (2) যোগ করে পাই,}

\emph{2x = 48}

\emph{⇒ x=(48/2) = 24}

\emph{∴ x = 24}
\end{quote}

\textbf{যদি x + 5y = 16 এবং x = −3y হয় তাহলে y = ?}

ক) −24

খ) −2

\textbf{\hl{গ) 8 ✔}}

ঘ) 2

\textbf{\hl{✔ সঠিক উত্তর: গ) 8}}

\begin{quote}
\textbf{ব্যাখ্যা:} \emph{দেওয়া আছে, x + 5y = 16 ......... (1)}

\emph{এবং x = −3y ⇒ x + 3y = 0 ......... (2)}

\emph{(1) নং সমীকরণ হতে (2) নং সমীকরণ বিয়োগ করে পাই,}

\emph{2y = 16}

\emph{∴ y = 8}
\end{quote}

\textbf{a + b = 2, a − b = 0 হলে, (a/b) = কত?}

ক) 0

\textbf{\hl{খ) 1 ✔}}

গ) 2

ঘ) 4

\textbf{\hl{✔ সঠিক উত্তর: খ) 1}}

\begin{quote}
\textbf{ব্যাখ্যা:} \emph{দেওয়া আছে, a + b = 2 এবং a − b = 0}

\emph{a − b = 0 ⇒ a = b}

\emph{এখন a এর পরিবর্তে b বসিয়ে পাই, b + b = 2 ⇒ 2b = 2 ⇒ b = 1 এবং a =
1}

\emph{সুতরাং (a/b)=(1/1)=1}
\end{quote}

\textbf{যদি x + y + 4 = x − y − 12 = 0 হয়, তবে 2x + y = কত?}

ক) 4

খ) −8

গ) −4

\textbf{\hl{ঘ) 0 ✔}}

\textbf{\hl{✔ সঠিক উত্তর: ঘ) 0}}

\begin{quote}
\textbf{ব্যাখ্যা:} \emph{দেওয়া আছে, x + y + 4 = x − y − 12 = 0}

\emph{অর্থাৎ, x + y + 4 = 0 ⇒ x + y = −4 ..... (i)}

\emph{এবং x − y − 12 = 0 ⇒ x − y = 12 ..... (ii)}

\emph{(i) ও (ii) যোগ করে পাই, 2x = 8 ⇒ x = 4}

\emph{(i) নং-এ x-এর মান বসিয়ে পাই, 4 + y = −4 ⇒ y = −4 − 4 ⇒ y = −8}

\emph{এখন, 2x + y = (2 × 4) − 8 = 8 − 8 = 0}

\emph{∴ 2x + y = 0}
\end{quote}

\textbf{3x − 7y + 10 = 0 এবং y − 2x − 3 = 0 এর সমাধান}

ক) x = 1, y = −1

খ) x = 1, y = 1

গ) x = −1, y = −1

\textbf{\hl{ঘ) x = −1, y = 1 ✔}}

\textbf{\hl{✔ সঠিক উত্তর: ঘ) x = −1, y = 1}}

\begin{quote}
\textbf{ব্যাখ্যা:} \emph{দেওয়া আছে, 3x − 7y + 10 = 0 ......... (I)}

\emph{এবং y − 2x − 3 = 0 ⇒ y = 2x + 3 ......... (II)}

\emph{(I) নং সমীকরণে y এর মান বসিয়ে পাই, 3x − 7(2x + 3) + 10 = 0}

\emph{⇒ 3x − 14x − 21 + 10 = 0}

\emph{⇒ −11x = 11 ⇒ x = −1}

\emph{এখন x এর মান (II) নং এ বসিয়ে পাই, y = 2(−1) + 3 = 1}

\emph{∴ x = −1, y = 1}
\end{quote}

\textbf{x\^{}2 + y\^{}2 = 185, x − y = 3 এর একটি সমাধান হবে---}

ক) a) (7, 4)

খ) b) (9, 6)

গ) c) (10, 7)

\textbf{\hl{ঘ) d) (11, 8) ✔}}

\textbf{\hl{✔ সঠিক উত্তর: ঘ) d) (11, 8)}}

\begin{quote}
\textbf{ব্যাখ্যা:} \emph{অপশন পরীক্ষা করে পাই, d) (11, 8) এর ক্ষেত্রে:}

\emph{x − y = 11 − 8 = 3}

\emph{এবং x\^{}2 + y\^{}2 = 11\^{}2 + 8\^{}2 = 121 + 64 = 185}

\emph{যা সমীকরণদ্বয়কে সিদ্ধ করে।}

\emph{∴ সঠিক উত্তর (11, 8)}
\end{quote}

\textbf{(x − y, 4) = (4, x + 3y) হলে (x, y) এর মান কত?}

ক) (6, 2)

খ) (3, 1)

গ) (5, 1)

\textbf{\hl{ঘ) (4, 0) ✔}}

\textbf{\hl{✔ সঠিক উত্তর: ঘ) (4, 0)}}

\begin{quote}
\textbf{ব্যাখ্যা:} \emph{ক্রমজোড়ের শর্তানুযায়ী,}

\emph{x − y = 4 ......... (1)}

\emph{x + 3y = 4 ......... (2)}

\emph{(1) নং হতে (2) নং বিয়োগ করে পাই, −4y = 0 ⇒ y = 0}

\emph{y এর মান (1) নং এ বসিয়ে পাই, x − 0 = 4 ⇒ x = 4}

\emph{∴ (x, y) = (4, 0)}
\end{quote}

\textbf{x + 2y = 7 এবং 2x − 3y = 0 হলে x এবং y মান হবে---}

ক) 2 ও 3

\textbf{\hl{খ) 3 ও 2 ✔}}

গ) 4 ও 2

ঘ) 3 ও 3

\textbf{\hl{✔ সঠিক উত্তর: খ) 3 ও 2}}

\begin{quote}
\textbf{ব্যাখ্যা:} \emph{দেওয়া আছে, x + 2y = 7 ......... (1)}

\emph{এবং 2x − 3y = 0 ......... (2)}

\emph{(1) নং সমীকরণকে 3 দ্বারা এবং (2) নং সমীকরণকে 2 দ্বারা গুণ করে পাই,}

\emph{3x + 6y = 21}

\emph{4x − 6y = 0}

\emph{যোগ করে পাই, 7x = 21 ⇒ x = 3}

\emph{x এর মান (1) নং এ বসিয়ে পাই, 3 + 2y = 7 ⇒ 2y = 4 ⇒ y = 2}

\emph{∴ x = 3, y = 2}
\end{quote}

\textbf{2x + 3y = 36 এবং 2x + y = 16 হলে (x, y) এর মান কত?}

ক) (2, 10)

\textbf{\hl{খ) (3, 10) ✔}}

গ) (2, 6)

ঘ) (3, 5)

\textbf{\hl{✔ সঠিক উত্তর: খ) (3, 10)}}

\begin{quote}
\textbf{ব্যাখ্যা:} \emph{দেওয়া আছে, 2x + 3y = 36 ......... (i)}

\emph{এবং 2x + y = 16 ⇒ y = 16 − 2x ......... (ii)}

\emph{(i) নং সমীকরণে y এর মান বসিয়ে পাই,}

\emph{2x + 3(16 − 2x) = 36}

\emph{⇒ 2x + 48 − 6x = 36}

\emph{⇒ −4x = −12 ⇒ x = 3}

\emph{এখন x এর মান (ii) নং এ বসিয়ে পাই, y = 16 − 2(3) = 10}

\emph{∴ (x, y) = (3, 10)}
\end{quote}

\textbf{x\^{}0 + y\^{}0 = কত?}

ক) 0

খ) 1

\textbf{\hl{গ) 2 ✔}}

ঘ) x + y

\textbf{\hl{✔ সঠিক উত্তর: গ) 2}}

\begin{quote}
\textbf{ব্যাখ্যা:} \emph{আমরা জানি, যেকোনো সংখ্যা বা চলের ওপর পাওয়ার শূন্য
হলে তার মান ১ হয় (যেহেতু a\^{}0 = 1; a ≠ 0)।}

\emph{সুতরাং, x\^{}0 + y\^{}0 = 1 + 1 = 2}

\emph{∴ 2}
\end{quote}

\textbf{যদি (x/3)−(2/y)=1 এবং (x/4)+(3/y)=3 হলে, (x, y) = কত?}

ক) (2, 5)

খ) (4, 3)

গ) (2, 6)

\textbf{\hl{ঘ) (6, 2) ✔}}

\textbf{\hl{✔ সঠিক উত্তর: ঘ) (6, 2)}}

\begin{quote}
\textbf{ব্যাখ্যা:} \emph{দেওয়া আছে,}

\emph{(x/3)−(2/y)=1 ......... (i)}

\emph{এবং (x/4)+(3/y)=3 ......... (ii)}

\emph{(i) নং থেকে পাই, (xy−6/3y)=1 ⇒ xy − 3y = 6 ......... (iii)}

\emph{(ii) নং থেকে পাই, (xy+12/4y)=3 ⇒ xy + 12 = 12y ⇒ xy − 12y = −12
......... (iv)}

\emph{(iii) থেকে (iv) বিয়োগ করে পাই, 9y = 18 ⇒ y = 2}

\emph{y এর মান (iii) নং এ বসিয়ে পাই, x(2) − 3(2) = 6 ⇒ 2x − 6 = 6 ⇒ 2x =
12 ⇒ x = 6}

\emph{∴ (x, y) = (6, 2)}
\end{quote}

\textbf{2x = 5 − y হলে 4x + 2y = কত?}

ক) 5

খ) 8

\textbf{\hl{গ) 10 ✔}}

ঘ) 12

\textbf{\hl{✔ সঠিক উত্তর: গ) 10}}

\begin{quote}
\textbf{ব্যাখ্যা:} \emph{দেওয়া আছে, 2x = 5 − y}

\emph{⇒ 2x + y = 5}

\emph{উভয় পক্ষকে 2 দ্বারা গুণ করে পাই,}

\emph{2(2x + y) = 5 × 2}

\emph{∴ 4x + 2y = 10}
\end{quote}

\textbf{যদি 8x + 4 = 64 হয়, তাহলে 2x + 1 এর মান কত?}

ক) 12

খ) 24

গ) 13

\textbf{\hl{ঘ) 16 ✔}}

\textbf{\hl{✔ সঠিক উত্তর: ঘ) 16}}

\begin{quote}
\textbf{ব্যাখ্যা:} \emph{দেওয়া আছে, 8x + 4 = 64}

\emph{⇒ 4(2x + 1) = 64}

\emph{⇒ 2x+1=(64/4)}

\emph{∴ 2x + 1 = 16}
\end{quote}

\textbf{x + 2y = 4 এবং x = 2y হলে, x এর মান---}

ক) 0

খ) (1/2)

গ) 1

\textbf{\hl{ঘ) 2 ✔}}

\textbf{\hl{✔ সঠিক উত্তর: ঘ) 2}}

\begin{quote}
\textbf{ব্যাখ্যা:} \emph{দেওয়া আছে, x + 2y = 4 ......... (1)}

\emph{এবং x = 2y ......... (2)}

\emph{(1) ও (2) নং হতে পাই, 2y + 2y = 4}

\emph{⇒ 4y = 4 ⇒ y = 1}

\emph{তাহলে, x = 2(1) = 2}

\emph{∴ x = 2}
\end{quote}

\textbf{x + y = 6 হলে xy এর বৃহত্তম মান কত?}

\textbf{\hl{ক) 9 ✔}}

খ) 7

গ) −8

ঘ) 12

\textbf{\hl{✔ সঠিক উত্তর: ক) 9}}

\begin{quote}
\textbf{ব্যাখ্যা:} \emph{x + y = 6 হলে, ধনাত্মক পূর্ণসংখ্যার ক্ষেত্রে মানগুলো
পরীক্ষা করে পাই:}

\emph{x = 5 হলে y = 1 ⇒ xy = 5 × 1 = 5}

\emph{x = 4 হলে y = 2 ⇒ xy = 4 × 2 = 8}

\emph{x = 3 হলে y = 3 ⇒ xy = 3 × 3 = 9}

\emph{∴ xy এর বৃহত্তম মান 9}
\end{quote}

\section{26. সমান্তর ও গুণোত্তর ধারা-বীজগণিত \textbar\textbar{} সংজ্ঞা,
সূত্রাবলীসহ এ টু
জেড}\label{ux9b8ux9aeux9a8ux9a4ux9b0-ux993-ux997ux9a3ux9a4ux9a4ux9b0-ux9a7ux9b0-ux9acux99cux997ux9a3ux9a4-ux9b8ux99cux99e-ux9b8ux9a4ux9b0ux9acux9b2ux9b8ux9b9-ux98f-ux99f-ux99cux9a1}

সমান্তর ও গুণোত্তর ধারা হলো গণিতের দুটি মৌলিক সংজ্ঞা যা ব্যবহার করা হয় সংখ্যাগত
প্রবণতা বা ক্রিয়ার বিষয়ে। সমান্তর ধারা হলো সংখ্যাগত প্রবণতার এমন একটি ধারা
যেখানে প্রতিটি পদের মান পূর্ববর্তী পদের সাথে একই পার্থক্য বজায় রাখে। আর গুণোত্তর
ধারা হলো এমন একটি সংখ্যাগত প্রবণতা যেখানে প্রতিটি পদ আগের পদের সাথে একই
অনুপাতে গুন করে। এই ধারাগুলো গণিতে অনেক গুরুত্বপূর্ণ বিশেষ করে চাকরি পরীক্ষার
প্রস্তুতির জন্য একটি গুরুত্বপূর্ণ অধ্যায় । আজকের টিউটরিয়ালে আমরা সমান্তর ও গুণোত্তর
ধারার সংজ্ঞা , বিভিন্ন পরিভাষা , পদের মান ও সমষ্টি বের করা ইত্যাদি সম্পর্কে এ
টু জেড জানব । এগুলোর প্রথমে আমরা ইন ডিটেইলস দেখব । তারপর বিগত সালের প্রশ্ন
সমাধান করতে করতে অধ্যায়টি ভালভাবে রপ্ত করার চেষ্টা করব । আশা করছি এই
টিউটরিয়ালের মাধ্যমে সমান্তর ও গুণোত্তর ধারা সম্পর্কে আপনার আর কোন সমস্যা থাকবে
না । চলুন তাহলে শুরু যাক ।

\begin{itemize}
\item
  প্রাথমিক আলোচনা
\item
  সমান্তর ধারা
\item
  পদক্রমের মান বের করা
\item
  পদক্রম বের করা
\item
  পদসংখ্যা বের করা
\item
  সূত্রের মাধ্যমে সমান্তর ধারার সমষ্টি নির্ণয়
\item
  শেষপদ দেয়া না থাকলে সমষ্টি বের করার সহজ কৌশল
\item
  সমান্তর ধারার সমষ্টি সংক্রান্ত আরো কিছু গুরুত্বপূর্ণ প্রশ্ন ও সমাধান
\item
  বর্গ ও ঘন - সংখ্যার সমষ্টি
\item
  গুণোত্তর ধারা
\item
  গুণোত্তর ধারার পদের সমষ্টি বের করা
\end{itemize}

\subsubsection{প্রাথমিক
আলোচনা}\label{ux9aaux9b0ux9a5ux9aeux995-ux986ux9b2ux99aux9a8-10}

ধারা একটি সিস্টেম যেখানে পরপর কিছু সংখ্যা থাকে যারা পরস্পর সুবিন্যাস্ত এবং
নির্দিষ্ট নিয়ম অনুসারে সাজানো।

ধারা সাধারণত দুই প্রকার:

(ক) সমান্তর ধারা (The Arithmetic Series)

(খ) গুণোত্তর ধারা (Geometric Series)

\textbf{ধারার বিভিন্ন পরিভাষা :}

\textbf{প্রথম পদ :} যে কোন ধারার শুরু হয় যে পদের দ্বারা তাকে প্রকাশ করা হয় a
দ্বারা ।

\textbf{সাধারণ অন্তর :} একপদের সাথে তার পরবর্তী পদের পার্থক্যকে সাধারণ অন্তুর
বলা হয়। সাধারণ অন্তর প্রকাশ করা হয় d দ্বারা। d = ২য় পদ - ১ম পদ

\textbf{শেষপদ:} কোন পদের শেষ পদ জানা না থাকলে তা r তম পদ বা n তম পদ
হিসেবে ধরা হয়। তাই n হলো শেষপদ ।

\textbf{পদসংখ্যা:} একটি ধারায় মোট কতটি পদ আছে তাদেরকে পদসংখ্যা বলা হয়। মোট
পদসংখ্যা অনির্দিষ্ট হলে পদসংখ্যা n ধরা হয় ।

\textbf{বিভিন্ন ধারার পরিচিতি:}

\textbf{(ক) সাধারন ধারা:} একটি নির্দিষ্ট ব্যবধানে বৃদ্ধি প্রাপ্ত বা হ্রাসকৃত পরপর
কিছু সংখ্যাকে সাধারণ ধারা বলে।

\textbf{(খ) সমান্তর ধারা:} কোনো ধারার পরপর দুইটি পদের মানের পার্থক্য সমান
হলে তাকে সমান্তর ধারা বলে। যেমন:

1+2+3+4+... ..100 \textbf{(গ) গুণোত্তর ধারা:} একটি ধারাকে প্রতিবার নির্দিষ্ট
একটি সংখ্যা দিয়ে গুণ অথবা ভাগ করে নতুন রাশি তৈরী করলে তাকে গুণোত্তর ধারা বলা
হয়। যেমন: 2+4+8+16+32 ...... +256

\subsubsection{সমান্তর ধারা}\label{ux9b8ux9aeux9a8ux9a4ux9b0-ux9a7ux9b0}

\textbf{সমান্তর ধারা:} যে ধারায় প্রতিটি পদে একই পরিমাণে বৃদ্ধি বা হ্রাস ঘটে
অর্থাৎ ক্রমিক দুইটি পদের অন্তর একই হয়, তাকে সমান্তর ধারা বলে। অর্থাৎ যে ধারার
কোন পদকে তার পরবর্তী পদ থেকে বিয়োগ করলে প্রতিক্ষেত্রে বিয়োগফল একই হয় তাকে
সমান্তর ধারা বলে। যেমন: 2+4+6+8+10....... +20 এখানে প্রতি পদের বৃদ্ধির হার 2
করে ।

\textbf{সমান্তর ধারার গুরুত্বপূর্ণ সূত্র ও তার প্রয়োগ:}

কোনো সমান্তর ধারার প্রথম পদ ও শেষ পদ দেয়া থাকলে...

\(\text{সমষ্টি} = \frac{\text{পদসংখ্যা}\ (\text{শেষপদ}\  + \ \text{প্রথম}\ \text{পদ})}{\text{২}}\text{এবং}\ \text{পদসংখ্যা} = \frac{\text{শেষপদ}\  - \ \text{প্রথম}\ \text{পদ}}{\text{সাধারন}\ \text{অন্তর}} + \text{১}\)

একটি সমান্তর ধারার প্রথম পদ a এবং সাধারণ অন্তর d হলে r তম পদ = a+(r-1)d

প্রথম পদ a এবং সাধারণ অন্তর d বিশিষ্ট সমান্তর ধারার n সংখ্যক পদের সমষ্টি s =
\(\frac{n}{2}\) \{2a+(n-1)d\}

\subsubsection{পদক্রমের মান বের
করা}\label{ux9aaux9a6ux995ux9b0ux9aeux9b0-ux9aeux9a8-ux9acux9b0-ux995ux9b0}

যে কোন একটি ধারা দেয়া থাকবে, সেই ধারার যে কোন একটি পদের মান কত তা বের
করতে বলা হবে।

ধারার পদক্রমের মান বের করার ক্ষেত্রে সূত্র হলো:

r তম পদ = a + (r-1)d

অর্থাৎ ধারার ৫ম, ৮ম বা ১০ম এভাবে যে কোন পদের মান বের করতে এই সূত্রটি প্রয়োগ
করতে হবে । নিচের প্রশ্নগুলো দেখুন ।

\textbf{প্রশ্ন নম্বর-} ১. ১, ৪, ৭ .......ধারার ২৯ তম পদটি কত? {[} যুব উন্নয়ন
অধিদপ্তর (ক্যাশিয়ার)-২০১৮{]}

\begin{itemize}
\item
  (ক) ৭৯
\item
  (খ) ৮২
\item
  (গ) ৮৫
\item
  (ঘ) ৮৮
\end{itemize}

\textbf{\hl{উত্তরঃ (গ) ৮৫}}

\textbf{সমাধান:}

আমরা জানি,

একটি সমান্তর ধারার r তম পদ = a+(r-1)d

এখানে, প্রথম পদ a = ১, সাধারণ অন্তর d = ৪-১ = ৩ এবং r তম পদ = ২৯

সুতরাং, a+(r-1)d

= ১+(২৯-১)×৩ = ১+২৮×৩ = ১+৮৪ = ৮৫

\textbf{প্রশ্ন নম্বর-} ২. একটি সমান্তর অনুক্রমে সাধারণ অন্তর 10 এবং 6-তম পদটি
52 হলে 15-তম পদটি- {[} ৩৭-তম বিসিএস প্রিলি {]}

\begin{itemize}
\item
  (ক) 140
\item
  (খ) 142
\item
  (গ) 148
\item
  (ঘ) 150
\end{itemize}

\textbf{\hl{উত্তরঃ (খ) 142}}

\textbf{সমাধান:}

আমরা জানি, r তম পদ = a+(r-1)d

প্রশ্নমতে , ৬ষ্ঠ পদ = a+(6-1)d = 52

বা, a+5×10 = 52

বা, a + 50 = 52

বা, a = 52 - 50

বা, a = 2

অর্থাৎ ধারাটির প্রথম পদ a = 2

এখন ১৫ তম পদের মান বের করার জন্য একই সূত্র a+ (r-1)d = 2+ (15-1)10 =
2+14×10 = 142 Ans:

\textbf{মুখে মুখে সমাধান :}

সাধারণ অন্তর ১০ এবং ৬ষ্ঠ পদ ৫২ তাহলে ১৫ তম পদ হবে ৬ষ্ঠ পদের পর আরো (১৫-৬)=
৯টি পদের যোগফল । প্রতিটি পদের মধ্যকার পার্থক্য ১০ হওয়ায় ৯টি পদের পার্থক্য হবে
৯x১০ = ৯০। তাহলে ১৫তম পদটি হবে ৫২+৯০ = ১৪২ ।

\subsubsection{পদক্রম বের
করা}\label{ux9aaux9a6ux995ux9b0ux9ae-ux9acux9b0-ux995ux9b0}

পদক্রম বের করা অর্থ হলো এর আগের নিয়মটিতে যেখানে ৮ম, ১০ম পদের মান কত তা বের
করতে বলা হয়েছে আর এখানে পদটি কত তম পদ তা দেয়া থাকবে না কিন্তু তার মান দেয়া
থাকবে ।

এখন ঐ পদটির পদক্রম কত তা বের করতে হবে। অর্থাৎ আগের নিয়মটির ঠিক উল্টো নিয়ম।

\textbf{সূত্র : আমরা জানি r তম পদ = a+(r-1)d অর্থাৎ সূত্র একটাই ।}

পূর্বের নিয়মটিতে r এর মান কত তা জানা ছিল। ( ৫ম,৮ম,১০ম এভাবে) কিন্তু এবার r
এর মান দেয়া থাকবে না । তবে এর মান বসানোর পর উত্তর কত হবে তা দেয়া থাকবে।
নিচের প্রশ্নগুলোর সমাধান দেখুন ।

\textbf{প্রশ্ন নম্বর-} ১. 5+8+11 +14+ ........... এ ধারাটির কত তম পদ 320 ?
{[}বাংলাদেশ রেলওয়ে (সহ: ষ্টেশন মাষ্টার)-২০১৮{]}

\begin{itemize}
\item
  (ক) ১০৬ তম পদ
\item
  (খ) ১০৮ তম পদ
\item
  (গ) ১০৯ তম পদ
\item
  (ঘ) ১১০ তম পদ
\end{itemize}

\textbf{\hl{উত্তরঃ (ক) ১০৬ তম পদ}}

\textbf{সমাধান:}

আমরা জানি, কোন সমান্তর ধারার r তম পদ = a+(r-1)d

প্রদত্ত সমান্তর ধারাটিতে প্রথম পদ a = 5 এবং সাধারণ অন্তর d = 8-5 = 3

r তম পদ 320 হলে r এর মান কত তা বের করতে হবে,

প্রশ্নমতে,

a+(r-1)d = 320 (সূত্রটি প্রয়োগ করলে যে মান আসবে তা ৩২০ হবে)

বা, 5+(r-1)3 = 320

বা, (r-1)3 = 320-5

বা, (r-1)3 = 315

বা, (r-1) = 315÷3

বা, (r-1) = 105

বা, r = 105+1 = 106

\textbf{\hl{উত্তর: ১০৬ তম পদ ।}}

\textbf{মুখে মুখে সমাধান :}

প্রথম পদ থেকে শেষ পদের ব্যবধান: ৩২০-৫ = ৩১৫ । এর মাঝে ৩ করে গ্যাপ রাখলে মোট
গ্যাপ হবে ৩১৫÷৩ = ১০৫টি। ১০৫ টি গ্যাপ হলে শুরুর পদটি সহ মোট পদ হবে ১০৫+১ =
১০৬ টি।

\textbf{প্রশ্ন নম্বর-} ২. 5+ 8+11+14+....... ধারাটির কোন পদ = 302? {[}
বাণিজ্য ও সমাজকল্যাণ মন্ত্র: নিয়োগ পরীক্ষা: ০৬ {]}

\begin{itemize}
\item
  (ক) 98 তম পদ
\item
  (খ) 99 তম পদ
\item
  (গ) 100 তম পদ
\item
  (ঘ) 101 তম পদ
\end{itemize}

\textbf{\hl{উত্তরঃ (গ) 100 তম পদ}}

\textbf{সমাধান:}

ধারাটির ১ম পদ, a = 5, সাধারণ অন্তর d = 8 -- 5 = 3,

মনে করি, r-তম পদ = 302

:. a + (r - 1)d = 302

বা, 5 + (r-1).3 = 302

বা, 5 + 3r-3 = 302

বা, 3r + 2 = 302

বা, 3r = 302 - 2

বা, 3r = 300

বা, r = 300÷3

বা, r = 100

:: 302 = 100 তম পদ

\textbf{মুখে মুখে সমাধান :}

প্রথমপদ থেকে শেষ পদের ব্যবধান: ৩০২-৫ = ২৯৭ । এর মাঝে ৩ করে

গ্যাপ রাখলে মোট গ্যাপ হবে ২৯৭÷৩ = ৯৯টি । ৯৯টি গ্যাপ হলে শুরুর

পদটি সহ মোট পদ হবে ৯৯+১ = ১০০টি।

\subsubsection{পদসংখ্যা বের
করা}\label{ux9aaux9a6ux9b8ux996ux9af-ux9acux9b0-ux995ux9b0}

পদসংখ্যা বের করার সূত্র =
\(\frac{\text{শেষ}\ \text{পদ}\  - \ \text{প্রথম}\ \text{পদ}}{\text{সাধারন}\ \text{অন্তর}} + \text{১}\)

ধরুন: ১০ মিটার পর পর ১টি করে গাছ লাগানো হবে। তাহলে ১০০ মিটার জায়গায় মোট
১০টি গাছ নয় বরং ১০+১ = ১১টি গাছ লাগাতে হবে । ( শুরুর গাছ সহ ধরলে ১টি গাছ
অতিরিক্ত হয়। )

উদাহরণ: ১০০ মিটার থেকে শুরু করে ২০০ মিটার পর্যন্ত গাছ লাগাতে হবে, ১০০, ১১০,
১২০ . . . . ২০০ মিটার এভাবে। তাহলে মোট গাছের সংখ্যা কতটি হবে? ১০টি নাকি
১১টি? এখানে,

গাছ লাগানো হবে, ১০০ (১০০মিটারে ই ১ম গাছ লাগিয়ে শুরু।) তারপর ১১০, ১২০, ১৩০,
১৪০, ১৫০, ১৬০, ১৭০, ১৮০, ১৯০ ও ২০০ মিটারে। = ১১টি । এই বিষয়টিকেই শর্টকাটে
বের করার জন্য সূত্রটি প্রয়োগ করুন :

পদসংখ্যা = \(\frac{\text{শেষ}\ \text{পদ}\  - \ \text{প্রথম}\ \text{পদ}}{\text{সাধারন}\ \text{অন্তর}} + \text{১}\)

{[}এখানে উপরের রাশিগুলোর মাঝের পার্থক্য বের করতে হবে, ২০০-১০০ =১০০মিটার
জায়গাতেই গাছ লাগাতে হবে । এখন সাধারণ অন্তর বা ১টি গাছ থেকে আরেকটি গাছের
দূরত্ব ১০ মিটার তাহলে মোট গাছ লাগানো যাবে, = ১০টি? না, বরং উত্তর হবে ১০+
প্রথমে ১টি = ১১টি ।। আর এজন্যই সূত্রটির সাথে সবসময় শেষে ১

যোগ হয় । {]}

এভাবে Practically ভাবতে পারলে সূত্রগুলো সবসময় মনে থাকবে।

\textbf{প্রশ্ন নম্বর-} ১. 20, 25, 30 ...140 ধারাটিতে মোট কতগুলো পদ আছে?
{[}জনস্বাস্থ্য প্রকৌশল অধি দপ্তর এস্টিমেটর নিয়োগ পরীক্ষা-২০১৮{]}

\begin{itemize}
\item
  (ক) 25 টি
\item
  (খ) 24 টি
\item
  (গ) 23 টি
\item
  (ঘ) 22 টি
\end{itemize}

\textbf{\hl{উত্তরঃ (ক) 25 টি}}

\textbf{সমাধান:}

এখানে, পদসংখ্যা = \(\frac{\text{শেষ}\ \text{পদ}\  - \ \text{প্রথম}\ \text{পদ}}{\text{সাধারন}\ \text{অন্তর}} + \text{১}\)

= \(\frac{140\  - \ 20}{5} + 1\)

= \(\frac{120}{5} + 1\)

= 24 + 1 = 25 \textbf{মুখে মুখে সমাধান:}

১৪০ ও ২০ এর মাঝে পার্থক্য ১২০। রাশি গুলোর পার্থক্য ৫ হওয়ায় মোট পদ ১২০÷৫ = ২৪
+ শুরুর ১টি = ২৫টি।

\subsubsection{সূত্রের মাধ্যমে সমান্তর ধারার সমষ্টি
নির্ণয়}\label{ux9b8ux9a4ux9b0ux9b0-ux9aeux9a7ux9afux9ae-ux9b8ux9aeux9a8ux9a4ux9b0-ux9a7ux9b0ux9b0-ux9b8ux9aeux9b7ux99f-ux9a8ux9b0ux9a3ux9df}

সূত্র-১: সমান্তর ধারায় ক্রমিক সংখ্যার যোগফল {[} যখন ধারাটি ১ থেকে শুরু হয় এবং
ধারার সাধারণ অন্তর ১ হয় । ( সাধারণ অন্তর = ২য় পদ - ১ম পদ ) {]}

S = \(\frac{n(n + 1)}{2}\) এখানে, n = পদসংখ্যা s = যোগফল

\textbf{প্রশ্ন নম্বর-} ১. ১ হতে ১০০ পর্যন্ত সংখ্যাসমুহের যোগফল কত? {[} ৮তম এবং
১৬তম বিসিএস{]}

\begin{itemize}
\item
  (ক) ৫০৫০
\item
  (খ) ৫৫০০
\item
  (গ) ৫০০৫
\item
  (ঘ) ৫০৬০
\end{itemize}

\textbf{\hl{উত্তরঃ (ক) ৫০৫০}}

\textbf{সমাধান:}

\textbf{সমাধান}ঃ

S = \(\frac{n(n + 1)}{2}\) এখানে, n = পদসংখ্যা, s = যোগফল

= \(\frac{\text{১০০}(\text{১০০} + \text{১})}{\text{২}}\)

= \(\frac{\text{১০০}\  \times \text{১০১}}{\text{২}}\)

= ৫০ × ১০১ = ৫০৫০

\textbf{প্রশ্ন নম্বর-} ২. ১+২+৩+..............+২৫ এর যোগফল কত? {[}
বি.আর.সি- এর সিনিয়র অফিসার পরীক্ষা-১৯৯৮{]}

\begin{itemize}
\item
  (ক) ১০০
\item
  (খ) ২৫০
\item
  (গ) ৩০০
\item
  (ঘ) ৩২৫
\end{itemize}

\textbf{\hl{উত্তরঃ (ঘ) ৩২৫}}

\textbf{সমাধান:}

\textbf{সমাধান}ঃ

S = \(\frac{n(n + 1)}{2}\) এখানে, n = পদসংখ্যা, s = যোগফল

= \(\frac{\text{২৫}(\text{২৫} + \text{১})}{\text{২}}\)

= \(\frac{\text{২৫}\  \times \text{২৬}}{\text{২}}\)

= ২৫ × ১৩ = ৩২৫

\textbf{প্রশ্ন নম্বর-} ৩. ১ হতে ৪৯ পর্যন্ত ক্রমিক সংখ্যা গুলোর গড় কত? {[} সমাজ
কল্যাণ সংগঠক পরীক্ষা-২০০৭ {]}

\begin{itemize}
\item
  (ক) ২৩
\item
  (খ) ২৪
\item
  (গ) ২৫
\item
  (ঘ) ২৬.৫
\end{itemize}

\textbf{\hl{উত্তরঃ (গ) ২৫}}

\textbf{সমাধান:}

\textbf{সমাধান}ঃ

সমষ্টি =

\(\frac{n(n + 1)}{2}\)

= \(\frac{\text{৪৯} \times \text{৫০}}{\text{২}}\)

= \(\frac{\text{২৪৫০}}{\text{২}}\)

= ১২২৫ সংখ্যাগুলোর গড়

= ১২২৫÷৪৯ = ২৫

\textbf{প্রশ্ন নম্বর-} ৪. ০ হতে ৪৯ পর্যন্ত ক্রমিক সংখ্যা গুলোর গড় কত? সাব
রেজিস্টার-২০০১১{]}

\begin{itemize}
\item
  (ক) ২৩
\item
  (খ) ২৫
\item
  (গ) ২৪.৫
\item
  (ঘ) ২৬.৫
\end{itemize}

\textbf{\hl{উত্তরঃ (গ) ২৪.৫}}

\textbf{সমাধান:}

যদি ও এখানে পদসংখ্যা ৫০ টি । কিন্তু যোগফল বা সমষ্টি নির্ণয়ের ক্ষেত্রে প্রথম
সংখ্যা শূন্য কোন ভূমিকা রাখতে পারবে না বিধায় সমষ্টি ৪৯ টি পদ ধরেই সমাধান করতে
হবে।কিন্তু গড় নির্ণয়ের জন্য ঠিকই ৫০টি পদ ধরে সমাধান করতে হবে।

\textbf{সমাধান}ঃ

S = \(\frac{n(n + 1)}{2}\) এখানে, n = পদসংখ্যা, s = যোগফল

= \(\frac{\text{৪৯}(\text{৪৯} + \text{১})}{\text{২}}\)

= \(\frac{\text{৪৯}\  \times \text{৫০}}{\text{২}}\)

= ২৫ × ৪৯ = ১২২৫ সংখ্যাগুলোর গড়

= ১২২৫÷৫০ = ২৪.৫

সূত্র-২: যখন ১ ভিন্ন অন্য সংখ্যা দ্বারা ক্রমিক সংখ্যা শুরু হয় এবং ধারার সাধারণ
অন্তর ১ হয় । ( সাধারণ অন্তর = ২য় পদ - ১ম পদ ) {]} তখন-\/-\/-

S = \(\frac{n}{2}\)(n+1)-\(\frac{a}{2}\)(a+1)

এখানে, n = শেষ সংখ্যা ,

S = যোগফল,

a = যে সংখ্যা দ্বারা সিরিজ শুরু হয় তার আগের সংখ্যা।

\textbf{প্রশ্ন নম্বর-} ১. ৫ থেকে ৪৫ পর্যন্ত সংখ্যাসমুহের যোগফল কত? {[}প্রাথমিক
প্রধান শিক্ষক পরীক্ষা-১৯৯৯{]}

\begin{itemize}
\item
  (ক) ১০২৫
\item
  (খ) ১০৭৫
\item
  (গ) ১০৩৫
\item
  (ঘ) ১০৪৫
\end{itemize}

\textbf{\hl{উত্তরঃ (ক) ১০২৫}}

\textbf{সমাধান:}

\textbf{সমাধান}ঃ

S = \(\frac{n}{2}\)(n+1)-\(\frac{a}{2}\)(a+1)

= \(\frac{\text{৪৫}}{\text{২}}\)(৪৫+১)-\(\frac{\text{৪}}{\text{২}}\)(৪+১)

= ৪৫×২৩- ২×৫ = ১০৩৫-১০ = ১০২৫

\textbf{প্রশ্ন নম্বর-} ২. ৩ থেকে ৮১ পর্যন্ত সংখ্যার যোগফল কত? সহকারী রিসার্চ
অফিসার-১৯৯৮{]}

\begin{itemize}
\item
  (ক) ৩২২৪
\item
  (খ) ৩১২০
\item
  (গ) ৩০২৪
\item
  (ঘ) ৩৩১৮
\end{itemize}

\textbf{\hl{উত্তরঃ (ঘ) ৩৩১৮}}

\textbf{সমাধান:}

\textbf{সমাধান}ঃ

S = \(\frac{n}{2}\)(n+1)-\(\frac{a}{2}\)(a+1)

= \(\frac{\text{৮১}}{\text{২}}\)(৮১+১)-\(\frac{\text{২}}{\text{২}}\)(২+১)

= ৪০.৫ × ৮২- ১×৩ = ৩৩২১ - ৩ = ৩৩১৮

সূত্র-৩: সমান্তর ধারায় ক্রমিক বিজোড় সংখ্যার যোগফল

S = \(M^{2}\) {[}M = মধ্যমা = \(\frac{\text{১ম}\ \text{সংখ্যা}\  + \ \text{শেষ}\ \text{সংখ্যা}}{\text{২}}\){]}

\textbf{প্রশ্ন নম্বর-} ১. ১+৩+৫+........+১৯ = ? {[} সমাজসেবা উপসহকারী
পরিচালক পরীক্ষা-২০০৫{]}

\begin{itemize}
\item
  (ক) ৯৮
\item
  (খ) ১০১
\item
  (গ) ৯৯
\item
  (ঘ) ১০০
\end{itemize}

\textbf{\hl{উত্তরঃ (ঘ) ৭২}}

\textbf{সমাধান:}

\textbf{সমাধান}: S = \(M^{2}\) = \({(\frac{\text{১} + \text{১৯}}{\text{২}})}^{\text{২}}\) =
\({(\text{১০})}^{\text{২}}\) = ১০০

\textbf{প্রশ্ন নম্বর-} ২. 1+3+5+.........+(2n-1) ধারাটির যোগফল হবে-\/- {[}
প্রাথমিক সহকারী শিক্ষক -২০১০{]}

\begin{itemize}
\item
  (ক) ((2n-1))\^{}2
\item
  (খ) (n(n+1)/2)
\item
  (গ) n\^{}2
\item
  (ঘ) (\{(n(n+1)/2)\})\^{}2
\end{itemize}

\textbf{\hl{উত্তরঃ (গ)}} \(n^{2}\)

\textbf{\hl{উত্তরঃ গ}}

\textbf{সমাধানঃ}

S = \(M^{2}\) = \({(\frac{1 + 2n - 1}{2})}^{2}\) = \(n^{2}\)

\textbf{সূত্র-৪: সমান্তর ধারায় ক্রমিক জোড় সংখ্যার যোগফল}

S = M(M-1) {[}M = মধ্যমা = \(\frac{\text{১ম}\ \text{সংখ্যা}\  + \ \text{শেষ}\ \text{সংখ্যা}}{\text{২}}\){]}

\textbf{প্রশ্ন নম্বর-} ১. ২+৪+৬+১০+...............+৫০ = ? জরিপ অধিদপ্তরের
সহকারী সুপারিনটেনডেন্ট পরীক্ষা-২০০৯{]}

\begin{itemize}
\item
  (ক) ৬৫০
\item
  (খ) ৫২০
\item
  (গ) ৪৫০
\item
  (ঘ) ৫২২
\end{itemize}

\textbf{\hl{উত্তরঃ (ক) ৬৫০}}

\textbf{সমাধান:}

S = M(M-1

= ২৬(২৬-১) = ২৬×২৫ = ৬৫০

\textbf{প্রশ্ন নম্বর-} ২. ২ হতে ৪০ পর্যন্ত জোড় সংখ্যাগুলোর যোগফল কত? হিসাব
সহকারী (মেঘনা)-২০১৩{]}

\begin{itemize}
\item
  (ক) ৪২০
\item
  (খ) ৫২০
\item
  (গ) ৪৫০
\item
  (ঘ) ৫২২
\end{itemize}

\textbf{\hl{উত্তরঃ (ক) ৪২০}}

\textbf{সমাধান:}

S = M(M-1)

= ২১(২১-১) = ২১×২০ = ৪২০

\textbf{সূত্র-৫: সমান্তর ধারায় ক্রমিক সংখ্যার যোগফল {[} যখন ধারাটি ১ কিংবা ১
এর অধিক যে কোন সংখ্যা থেকে শুরু হয় এবং ধারার সাধারণ অন্তর ২ এর বেশি হয় । (
সাধারণ অন্তর = ২য় পদ - ১ম পদ ) {]}}

তখন প্রথমে পদসংখ্যা বের করতে হবে এবং তারপর সমষ্টি নির্ণয়ের সূত্র প্রয়োগ করতে
হবে ।

পদসংখ্যা = \(\frac{\text{শেষ}\ \text{পদ}\  - \ \text{প্রথম}\ \text{পদ}}{\text{সাধারন}\ \text{অন্তর}} + \text{১}\)

এবং সমষ্টি = \(\frac{\text{পদসংখ্যা}\ (\ \text{শেষ}\ \text{পদ}\  + \ \text{প্রথম}\ \text{পদ}\ )}{\text{২}}\)

\textbf{প্রশ্ন নম্বর-} ১. ১+৫+৯+ -\/-\/-\/-\/-\/- +৮১ =? {[} ৩৬তম বিসিএস{]}

\begin{itemize}
\item
  (ক) ৯৬১
\item
  (খ) ৮৬১
\item
  (গ) ৭৬১
\item
  (ঘ) ৬৬১
\end{itemize}

\textbf{\hl{উত্তরঃ (খ) ৮৬১}}

\textbf{সমাধান:}

এখানে প্রথম পদ a = ১ , সাধারণ অন্তর d = ৫-১ = ৪ এবং শেষ পদ = ৮১

সুতরাং পদসংখ্যা = \(\frac{\text{শেষ}\ \text{পদ}\  - \ \text{প্রথম}\ \text{পদ}}{\text{সাধারন}\ \text{অন্তর}} + \text{১}\)

= \(\frac{\text{৮১}\  - \ \text{১}}{\text{৪}} + \text{১}\)

= \(\frac{\text{৮০}}{\text{৪}} + \text{১}\)

= ২০ + ১ = ২১ সমষ্টি =
\(\frac{\text{পদসংখ্যা}\ (\ \text{শেষ}\ \text{পদ}\  + \ \text{প্রথম}\ \text{পদ}\ )}{\text{২}}\)

= \(\frac{\text{২১}\ (\ \text{৮১}\  + \ \text{১}\ )}{\text{২}}\)

= \(\frac{\text{২১} \times \text{৮২}}{\text{২}}\)

= ২১ × ৪১ = ৮৬১

\textbf{প্রশ্ন নম্বর-} ২. ১ + ৪ + ৭ + ১০ + -\/-\/-\/-\/-\/-\/-\/-\/- + ৭৩
সমান্তর ধারাটির যোগফল কত হবে? {[}একটি বাড়ী একটি খামার -(জেলা
সমন্বয়কারী)-১৭{]}

\begin{itemize}
\item
  (ক) ৯২৫
\item
  (খ) ১১২৫
\item
  (গ) ১০২৫
\item
  (ঘ) ১২২৫
\end{itemize}

\textbf{\hl{উত্তরঃ (ক) ৯২৫}}

\textbf{সমাধান:}

এখানে প্রথম পদ a = ১ , সাধারণ অন্তর d = ৪ - ১ = ৩ এবং শেষ পদ = ৭৩

সুতরাং পদসংখ্যা = \(\frac{\text{শেষ}\ \text{পদ}\  - \ \text{প্রথম}\ \text{পদ}}{\text{সাধারন}\ \text{অন্তর}} + \text{১}\)

= \(\frac{\text{৭৩}\  - \ \text{১}}{\text{৩}} + \text{১}\)

= \(\frac{\text{৭২}}{\text{৩}} + \text{১}\)

= ২৪ + ১ = ২৫ সমষ্টি =
\(\frac{\text{পদসংখ্যা}\ (\ \text{শেষ}\ \text{পদ}\  + \ \text{প্রথম}\ \text{পদ}\ )}{\text{২}}\)

= \(\frac{\text{২৫}\ (\ \text{৭৩}\  + \ \text{১}\ )}{\text{২}}\)

= \(\frac{\text{২৫}\  \times \ \text{৭৪}}{\text{২}}\)

= ২৫ × ৩৭ = ৯২৫

\subsubsection{শেষপদ দেয়া না থাকলে সমষ্টি বের করার সহজ
কৌশল}\label{ux9b6ux9b7ux9aaux9a6-ux9a6ux9af-ux9a8-ux9a5ux995ux9b2-ux9b8ux9aeux9b7ux99f-ux9acux9b0-ux995ux9b0ux9b0-ux9b8ux9b9ux99c-ux995ux9b6ux9b2}

প্রশ্নের মধ্যে শেষ পদের উল্লেখ না থেকে যদি পদের সংখ্যা কতটি তার উল্লেখ থাকে ।

তাহলে আগের সূত্রগুলো ব্যবহার করা যাবে না।

\textbf{সূত্র-৬: শেষ পদ না থাকলে সমষ্টি বের করার সূত্র S =}
\(\frac{n}{2}\{ 2a\  + \ (n\  - \ 1)\ d\}\)

প্রশ্ন নম্বর-

\begin{itemize}
\item
  (ক) 2500
\item
  (খ) 2445
\item
  (গ) 2385
\item
  (ঘ) 2600
\end{itemize}

\textbf{\hl{উত্তরঃ (গ) 2385}}

\textbf{সমাধান:}

ধারাটির প্রথম পদ a = 7, সাধারণ অন্তর d = 12 - 7 = 5 এবং পদ সংখ্যা n = 30

30টি পদের সমষ্টি = \(\frac{30}{2}\{ 2 \times 7\  + \ (30\  - \ 1)\ 5\}\)

= 15 ( 14 + 29×5 ) = 15 ( 14 + 145 )

= 15 × 159 = 2385

\textbf{প্রশ্ন নম্বর-} ২. 11 + 18 + 25 + 32 + ....... ধারাটির 29 তম পদের
সমষ্টি কত? {[} প্রাথমিক বিদ্যালয় সহকারী শিক্ষক - ২০১২ {]}

\begin{itemize}
\item
  (ক) 2131
\item
  (খ) 3131
\item
  (গ) 3161
\item
  (ঘ) 3260
\end{itemize}

\textbf{\hl{উত্তরঃ (গ) 3161}}

\textbf{সমাধান:}

ধারাটির প্রথম পদ a = 11, সাধারণ অন্তর d = 18 - 11 = 7 এবং পদ সংখ্যা n = 29

ধারাটির 29 তম পদের সমষ্টি =
\(\frac{29}{2}\{ 2 \times 11\  + \ (29\  - \ 1)\  \times \ 7\}\)

= \(\frac{29}{2}(22\  + \ 28\  \times \ 7)\)

= \(\frac{29}{2}(22\  + \ 196)\)

= \(\frac{29}{2} \times 218\)

= 29 × 109 = 3161

\subsubsection{সমান্তর ধারার সমষ্টি সংক্রান্ত আরো কিছু গুরুত্বপূর্ণ প্রশ্ন ও
সমাধান}\label{ux9b8ux9aeux9a8ux9a4ux9b0-ux9a7ux9b0ux9b0-ux9b8ux9aeux9b7ux99f-ux9b8ux995ux9b0ux9a8ux9a4-ux986ux9b0-ux995ux99b-ux997ux9b0ux9a4ux9acux9aaux9b0ux9a3-ux9aaux9b0ux9b6ux9a8-ux993-ux9b8ux9aeux9a7ux9a8}

\textbf{প্রশ্ন নম্বর-} ১. একটি সমান্তর অনুক্রমে ৫ম পদটি 18 এবং প্রথম ৫টি পদের
যোগফল 75 হলে প্রথম পদটি কত? {[}৩৮তম বিসিএস প্রিলি {]}

\begin{itemize}
\item
  (ক) ২
\item
  (খ) ১০
\item
  (গ) ৪
\item
  (ঘ) ১২
\end{itemize}

\textbf{\hl{উত্তরঃ (ঘ) ১২}}

\textbf{সমাধান:}

এখানে ১ম পদ = a. সাধারণ অন্তর = d

সমান্তর ধারার ৫ম পদ = a + (5 - 1)d = a + 4d

প্রশ্নমতে, a + 4d = 18-\/-\/-\/-\/-(i)

প্রথম ৫টি পদের সমষ্টি = \(\frac{5}{2}\{ 2a\  + \ (5\  - \ 1)d\}\)

= \(\frac{5}{2}\{ a\  + \ a\  + \ (5\  - \ 1)d\}\)

= \(\frac{5}{2}\{ a\  + \ (a\  + \ 4d)\}\)

= \(\frac{5}{2}(a\  + \ 18)\)

শর্তমতে, \(\frac{5}{2}(a\  + \ 18)\) = 75

বা, a + 18 = 30

সুতরাং a = 30 - 18 = 12

\textbf{কয়েক সেকেন্ডে সমাধান করার জন্য:}

৫টি সংখ্যার যোগফল ৭৫ হলে তাদের গড় = ৭৫÷৫ = ১৫

সমান্তর ধারায় সিরিজটি হবে ১২, ১৩.৫, ১৫, ১৬.৫ এবং ১৮।

গড় কে মাঝখানে রেখে শেষের ১৮ এর সাথে পার্থক্য ৩ কে ২ দিয়ে ভাগ করলে ১.৫ হয়।
তাই সবগুলো রাশির পার্থক্য ১.৫

\textbf{প্রশ্ন নম্বর-} ২. একটি সমান্তর ধারার ৪র্থ এবং ১২তম পদের যোগফল ২০। ঐ
ধারাটির প্রথম ১৫টি পদের যোগফল কত? {[} BD House Building FC (SO)-2017{]}

\begin{itemize}
\item
  (ক) 300
\item
  (খ) 120
\item
  (গ) 150
\item
  (ঘ) 130
\end{itemize}

\textbf{\hl{উত্তরঃ (গ) 150}}

\textbf{সমাধান:}

এখানে ৪র্থ পদ = a+(4-1)d = a+3d

এখানে ১২ তম পদ = a+(12-1)d = a+11d

৪র্থ ও ১২ তম পদের যোগফল = = a+3d+a+11d = 2a+14d = 20 {[}২০ প্রশ্নে দেয়া
আছে{]}

সুতরাং প্রথম ১৫টি পদের যোগফল = \(\frac{15}{2}\{ 2a\  + \ (15\  - \ 1)d\}\)

= \(\frac{15}{2}\{ 2a\  + \ 14d\}\)

= \(\frac{15}{2} \times 20\)

= 150

\textbf{প্রশ্ন নম্বর-} ৩. একটি সমান্তর ধারার প্রথম পদ 1, শেষ পদ 99 এবং সমষ্টি
2500 হলে ধারাটির সাধারণ অন্তর হবে- {[} সংসদ সচিবালয়-05 {]}

\begin{itemize}
\item
  (ক) - 2
\item
  (খ) 2
\item
  (গ) 20
\item
  (ঘ) 21
\end{itemize}

\textbf{\hl{উত্তরঃ (খ) 2}}

\textbf{সমাধান:}

প্রথমে প্রশ্নে প্রদত্ত ক্লু বা সংখ্যা গুলো নিয়ে ভাবুন। এখানে সাধারণ অন্তর d এর মান
বের করতে হবে। তাই প্রথমে সমষ্টির সুত্র প্রয়োগ করে পদসংখ্যা বের করার পর যে সুত্রে
d এর মান বের হবে সেই সূত্র প্রয়োগ করতে হবে। যেমন:

সমষ্টি = \(\frac{(1\  + \ 99)\  \times \ \text{পদসংখ্যা}}{2}\)

বা, 2500 = 50 × পদসংখ্যা

সুতরাং পদসংখ্যা = 50

আবার পদসংখ্যা = \(\frac{(99\  - \ 1)}{d} + 1\) (এই সূত্রে সাধারণ অন্তর d এর
মান বের হবে)

বা, 50 = \(\frac{98}{d} + 1\)

সুতরাং d = 2

\subsubsection{বর্গ ও ঘন - সংখ্যার
সমষ্টি}\label{ux9acux9b0ux997-ux993-ux998ux9a8---ux9b8ux996ux9afux9b0-ux9b8ux9aeux9b7ux99f}

\textbf{টিপস:} এ নিয়মের অংকগুলোতে শুধুমাত্র সূত্র মনে রেখে সূত্র প্রয়োগ করতে
পারলেই উত্তর বের হয়ে যাবে। তবে বর্তমানে পরীক্ষায় ক্যালকুলেটরের ব্যবহার না
থাকায় এ ধরনের প্রশ্ন খুব বেশি আসে না। আসলেও সহজ সংখ্যার গুলোই আসে। তাই
সূত্রগুলোকে গুরুত্ব দিন।

\textbf{সূত্র-৭: সমান্তর ধারায় বর্গ যোগ পদ্বতির ক্ষেত্রে-}

{[} \(\text{১}^{\text{২}} + \text{২}^{\text{২}} + \text{৩}^{\text{২}} + ....\  + n^{2}\) {]}

ধারার সমষ্টি S = \(\frac{\text{১}}{\text{৬}}\) n(n+১)(২n+১)

এখানে , n = শেষসংখ্যা

\textbf{প্রশ্ন নম্বর- ১.} \(\text{১}^{\text{২}} + \text{২}^{\text{২}} + \text{৩}^{\text{২}} + ....\  + {\text{৫০}}^{\text{২}}\)
\textbf{= কত?}

প্রাথমিক সহকারী শিক্ষক -২০১০{]}

\begin{itemize}
\item
  (ক) ৩৫৭২৫
\item
  (খ) ৪২৯২৫
\item
  (গ) ৪৫৫০০
\item
  (ঘ) ৪৭২২৫
\end{itemize}

\textbf{\hl{উত্তরঃ (খ) ৪২৯২৫}}

\textbf{সমাধান:}

\textbf{সমাধান}ঃ

ধারাটির যোগফল S = \(\frac{\text{১}}{\text{৬}}\) n(n+১)(২n+১)

এখানে , n = ৫০ ।

= \(\frac{\text{১}}{\text{৬}}\) ৫০(৫০+১)(২×৫০+১)

= \(\frac{\text{১}}{\text{৬}}\) ×৫০×৫১×১০১ = ৪২৯২৫

\textbf{প্রশ্ন নম্বর- ২.} \(1^{2} + 2^{2} + 3^{2} + ....\  + {(40)}^{2}\)
\textbf{= কত?}

{[} ২৭ তম বিসিএস {]}

\begin{itemize}
\item
  (ক) 22130
\item
  (খ) 22140
\item
  (গ) 22150
\item
  (ঘ) 22160
\end{itemize}

\textbf{\hl{উত্তরঃ (খ) 22140}}

\textbf{সমাধান:}

\textbf{সমাধান}ঃ

ধারাটির যোগফল S = \(\frac{1}{6}\) n(n+1)(2n+1)

এখানে , n = 40 ।

= \(\frac{1}{6}\) 40(40 + 1 )(2 × 40 + 1)

= \(\frac{1}{6}\) ×40×41×81

= 22140

\textbf{সূত্র-৮: সমান্তর ধারায় ঘন যোগ পদ্বতির ক্ষেত্রে-}

{[} \(\text{১}^{\text{৩}} + \text{২}^{\text{৩}} + \text{৩}^{\text{৩}} + ....\  + n^{3}\) {]}

ধারার সমষ্টি S = \({\{\frac{n(n + 1)}{2}\}}^{2}\)

এখানে , n = শেষসংখ্যা

\textbf{প্রশ্ন নম্বর-} ১.
\(\text{১}^{\text{৩}}\)+\(\text{২}^{\text{৩}}\)+\(\text{৩}^{\text{৩}}\)+..........+\({\text{১০}}^{\text{৩}}\) = কত? {[}থানা
সহকারী শিক্ষা অফিসার নিয়োগ পরীক্ষা-২০১০{]}

\begin{itemize}
\item
  (ক) ৩০২৫
\item
  (খ) ২৫৩০
\item
  (গ) ২৫৪০
\item
  (ঘ) ২৫৪৫
\end{itemize}

\textbf{\hl{উত্তরঃ (ক) ৩০২৫}}

\textbf{সমাধান:}

\textbf{সমাধান}ঃ

ধারাটির যোগফল

= \({\{\frac{n(n + 1)}{2}\}}^{2}\)

এখানে , n = 10

= \({\{\frac{10(10 + 1)}{2}\}}^{2}\)

= \({\{\frac{10\  \times \ 11}{2}\}}^{2}\)

= \({(55)}^{2}\) = 3025

\textbf{প্রশ্ন নম্বর-} ২.
\(\text{১}^{\text{৩}}\)+\(\text{২}^{\text{৩}}\)+\(\text{৩}^{\text{৩}}\)+..........+\({\text{২০}}^{\text{৩}}\) = কত? দুর্নীতি দমন
পরিদর্শক পরীক্ষা-২০০৪,মাদক অধিদপ্তর- ২০১২{]}

\begin{itemize}
\item
  (ক) ৪৪০০০
\item
  (খ) ৪৪১০০
\item
  (গ) ৪৪২০০
\item
  (ঘ) উপরের কোনটিই সত্য নয়
\end{itemize}

\textbf{\hl{উত্তরঃ (খ) ৪৪১০০}}

\textbf{সমাধান:}

\textbf{সমাধান}ঃ

ধারাটির যোগফল

= \({\{\frac{n(n + 1)}{2}\}}^{2}\)

এখানে , n = ২০

= \({\{\frac{\text{২০}(\text{২০}\  + \ \text{১})}{\text{২}}\}}^{\text{২}}\)

= \({\{\frac{\text{২০}\  \times \ \text{২১}}{\text{২}}\}}^{\text{২}}\)

= \({(\text{২১০})}^{\text{২}}\) = ৪৪১০০

\subsubsection{গুণোত্তর ধারা}\label{ux997ux9a3ux9a4ux9a4ux9b0-ux9a7ux9b0}

\textbf{গুণোত্তর ধারা:} যে ধারার কোনো পদের সাথে তার পরবর্তী পদের অনুপাত সব
সময় সমান হয়, সে ধারাকে গুনোত্তর ধারা বলে।

যেমন: 3+6+12+24+....... প্রতি পদের সাথে 2 গুণ করে পরবর্তী পদ বের করা হয়েছে।
যতগুণ করে বাড়ে বা কমে তাকে সাধারণ অনুপাত বলা হয় এবং সাধারণ অনুপাতকে q দ্বারা
প্রকাশ করা হয়।

\(q = \frac{\text{দ্বিতীয়}\ \text{পদ}}{\text{প্ৰথম}\ \text{পদ}}\)

গুণোত্তর ধারার সূত্র দুটি যথা:

১. r ৩ম পদ = \({aq}^{r\  - \ 1}\)

২. n সংখ্যক পদের সমষ্টি S = \(a \times \frac{q^{n} - \ 1}{q\  - \ 1}\) {[}
যখন q এর মান 1 থেকে বড় হয় তখন এই সূত্র {]}

আবার যখন q এর মান 1 থেকে ছোট হয় তখন নিচের সূত্র ব্যবহার করতে হবে।

S = \(a \times \frac{1 - q^{n}}{q\  - \ 1}\)

\paragraph{গুণোত্তর ধারার পদের মান বা কত তম পদ তা বের করতে বলা
হলে}\label{ux997ux9a3ux9a4ux9a4ux9b0-ux9a7ux9b0ux9b0-ux9aaux9a6ux9b0-ux9aeux9a8-ux9ac-ux995ux9a4-ux9a4ux9ae-ux9aaux9a6-ux9a4-ux9acux9b0-ux995ux9b0ux9a4-ux9acux9b2-ux9b9ux9b2}

গুণোত্তর ধারার প্রথম পদ a এবং সাধারণ অনুপাত q হলে ধারাটির r তম পদ =
\({aq}^{r\  - \ 1}\) (অর্থাৎ একটি ধারার মাঝখানের যে কোন একটি পদের মান কত
অথবা যে পদের মান মান দেয়া আছে পদটি কত নম্বর পদ বের করতে বলা হলে এই সূত্রটি
প্রয়োগ করতে হয়। যেমন:

\textbf{প্রশ্ন নম্বর-} ১. 2+4+8+16+ ...... ধারাটির ৭ম পদ কত? {[} সাব
রেজিস্টার: ০১ {]}

\begin{itemize}
\item
  (ক) 128
\item
  (খ) 120
\item
  (গ) 125
\item
  (ঘ) 124
\end{itemize}

\textbf{\hl{উত্তরঃ (ক) 128}}

\textbf{সমাধান:}

(সুত্রানুযায়ী ):

ধারাটির ৭ম পদ

= \({aq}^{r\  - \ 1}\) ( যেখানে প্রথম পদ a = 2, সাধারণ অনুপাত q = 4÷2 = 2
এবং পদের সংখ্যা r = 7)

= \({2 \times 2}^{7\  - \ 1}\) = \({2 \times 2}^{6}\) = 2 × 64 = 128

যেহেতু পার্থক্য দ্বিগুণ হারে বাড়ছে তাই সিরিজটি বুঝে বুঝে লিখলেই প্রমাণ পাওয়া যায়
2+4+8+16+32 + 64 + 128 ( ৭ম পদ 128 )

আবার এই একই প্রশ্নটি যদি এভাবে ঘুরিয়ে আসে । যেমন:

\textbf{প্রশ্ন নম্বর-} ২. 2+4+8+16+32 ......... ধারাটির কততম পদের মান 1024?
...

\begin{itemize}
\item
  (ক) ৮ম
\item
  (খ) ৯ম
\item
  (গ) ১০ম
\item
  (ঘ) ১১শ
\end{itemize}

\textbf{\hl{উত্তরঃ (গ) ১০ম}}

\textbf{সমাধান:}

তখনও কিন্তু একই সুত্র এখানে r তম পদটি কত তা জানা থাকলে ও r এর মান কত তা
অজানা, তাই এর মান বের করার জন্য লিখতে হবে।

\({aq}^{r\  - \ 1}\) = 128 (কারণ এই সুত্রটি প্রয়োগ করলে যে মানটি আসবে তা
হবে 128)

বা, \({2 \times 2}^{r\  - \ 1}\) = 1024

বা, \(2^{r\  - \ 1}\) = 1024

বা, \(2^{r\  - \ 1} = 2^{9}\)

বা, r - 1 = 9

বা, r = 10

সুতরাং ধারাটির ১০ম পদের মান 1024 .

\textbf{প্রশ্ন নম্বর-} ৩. একটি গুণোত্তর অনুক্রমের দ্বিতীয় পদটি - 48 এবং পঞ্চম
পদটি \(\frac{3}{4}\) হলে, সাধারণ অনুপাত কত? {[} 35তম বিসিএস {]}

\begin{itemize}
\item
  (ক) (1/2)
\item
  (খ) -(1/2)
\item
  (গ) (1/4)
\item
  (ঘ) -(1/4)
\end{itemize}

\textbf{\hl{উত্তরঃ}} \(- \frac{1}{4}\)

\textbf{সমাধান:}

গুণোত্তর অনুক্রমের প্রথম পদ a এবং সাধারণ অনুপাত q হলে অনুক্রমটির n তম পদ =
\({aq}^{n\  - \ 1}\) । দ্বিতীয় পদ,

\({aq}^{2\  - \ 1}\) = aq = - 48 এবং পঞ্চম পদ, \({aq}^{5\  - \ 1}\) =
\({aq}^{4} = \frac{3}{4}\) (বের করতে হবে সাধারণ অনুপাত q এর মান কত ? )

প্রথম শর্তমতে, aq = - 48 -\/-\/-\/-\/-\/-\/- (1)

২য় শর্তমতে, \({aq}^{4} = \frac{3}{4}\) -\/-\/-\/-\/-\/-\/- (2)

এখন 2 নং কে 1 নং দ্বারা ভাগ করে পাই,

\(\frac{{aq}^{4}}{aq} = \frac{4}{- \ 48}\)

বা, \(q^{3} = - \frac{1}{64}\)

বা, \(q^{3} = ( - \frac{1}{4^{3}})\)

বা, \(q = - \frac{1}{4}\)

সুতরাং সাধারণ অনুপাত q এর মান = \(- \frac{1}{4}\)

\textbf{প্রশ্ন নম্বর-} ৪. \(\frac{1}{\sqrt{2}},1,\sqrt{2}\) ধারাটির কোন পদ
8√2 হবে ? {[} বেসরকারি বিমান ও পর্যটন মন্ত্রণালয়ের প্রশাসনিক কর্মকর্তা ০৫) +
(৩৯ তম বিসিএস- (বিশেষ) {]}

\begin{itemize}
\item
  (ক) ৯ তম পদ
\item
  (খ) ১০ তম পদ
\item
  (গ) ১১ তম পদ
\item
  (ঘ) ১২ তম পদ
\end{itemize}

\textbf{\hl{উত্তরঃ (ক) ৯ তম পদ}}

\textbf{সমাধান:}

ধারাটি একটি গুণোত্তর ধারা।

এখানে, প্রথম পদ a = √2 এবং সাধারণ অনুপাত, q = 1 ÷ \(\frac{1}{\sqrt{2}}\)
বা, q = \(\sqrt{2}\)

n তম পদ = \({aq}^{n\  - \ 1}\)

শর্তমতে, \({aq}^{n\  - \ 1} = 8\sqrt{}2\)

বা,
\(\frac{1}{\sqrt{2}} \times {(\sqrt{}2)}^{n\  - \ 1} = 2^{3}\sqrt{2}\)

বা,
\({(\sqrt{}2)}^{- \ 1} \times {(\sqrt{}2)}^{n\  - \ 1} = {(\sqrt{}2)}^{6}\sqrt{2}\)

বা, \({(\sqrt{}2)}^{n\  - \ 2} = {(\sqrt{}2)}^{7}\)

বা, n - 2 = 7 (দুপাশের ভিত্তি √2 মিলে যাওয়ায় √2 তুলে দেয়া হয়েছে)

বা, n = 9

সুতরাং ধারাটির ৯তম পদ 8√2 হবে ।

\subsubsection{গুণোত্তর ধারার পদের সমষ্টি বের
করা}\label{ux997ux9a3ux9a4ux9a4ux9b0-ux9a7ux9b0ux9b0-ux9aaux9a6ux9b0-ux9b8ux9aeux9b7ux99f-ux9acux9b0-ux995ux9b0}

গুণোত্তর ধারার সূত্র দুটি যথা:

১. n ৩ম পদ = \({aq}^{n\  - \ 1}\)

২. n সংখ্যক পদের সমষ্টি S = \(a \times \frac{q^{n} - \ 1}{q\  - \ 1}\) {[}
যখন q এর মান 1 থেকে বড় হয় তখন এই সূত্র {]}

আবার যখন q এর মান 1 থেকে ছোট হয় তখন নিচের সূত্র ব্যবহার করতে হবে।

S = \(a \times \frac{1 - q^{n}}{1\  - \ q}\)

\textbf{প্রশ্ন নম্বর- ১. 2 + 6 + 18 + ধারাটির ৮টি পদের সমষ্টি কত?}

\begin{itemize}
\item
  (ক) 6565
\item
  (খ) 6560
\item
  (গ) 6569
\item
  (ঘ) 5560
\end{itemize}

\textbf{\hl{উত্তরঃ (খ) 6560}}

\textbf{সমাধান:}

এখানে গুণোত্তর ধারার প্রথম পদ a = 2 এবং সাধারণ অনুপাত q = 6÷2 = 3 {[} যা 1
থেকে বড় {]}

৮ টি পদের সমষ্টি

= \({aq}^{n\  - \ 1}\)

= \(2 \times \frac{3^{8} - \ 1}{3\  - \ 1}\)

= \(2 \times \frac{6561\  - \ 1}{2}\)

= 6560

\textbf{প্রশ্ন নম্বর-} ২. ২-৪+৮-১৬+......... ধারাটির প্রথম ৭ টি পদের সমষ্টি
কত ? {[} কারিগরি শিক্ষা অধিদপ্তর-০৫) + ( পল্লী সঞ্চয় ব্যাংক- (ক্যাশ)-২০১৮{]}

\begin{itemize}
\item
  (ক) ৩৬
\item
  (খ) ৮৬
\item
  (গ) ৫২
\item
  (ঘ) ৫৬
\end{itemize}

\textbf{\hl{উত্তরঃ (খ) ৮৬}}

\textbf{সমাধান:}

গুণোত্তর ধারাটির প্রথম পদ a = ২, সাধারণ অনুপাত, q = ( - ৪ )÷২ = - ২ {[} যা ১
থেকে ছোট {]} এবং মোট পদ n = ৭

প্রথম ৭ টি পদের সমষ্টি = \(a \times \frac{1 - q^{n}}{1\  - \ q}\)

= \(\text{২} \times \frac{\text{১} - {(\  - \ \text{২}\ )}^{\text{৭}}}{\text{১}\  - \ (\  - \ \text{২}\ )}\)

= \(\text{২} \times \frac{\text{১} - (\  - \ \text{১২৮}\ )}{\text{৩}}\)

= \(\text{২} \times \frac{\text{১} + \text{১২৮}}{\text{৩}}\)

= \(\text{২} \times \frac{\text{১২৯}}{\text{৩}}\)

= ২ × ৪৩ = ৮৬

\section{27. সমান্তর ও গুণোত্তর অনুক্রম ও ধারা MCQ প্রশ্ন উত্তর
অনুশীলন}\label{ux9b8ux9aeux9a8ux9a4ux9b0-ux993-ux997ux9a3ux9a4ux9a4ux9b0-ux985ux9a8ux995ux9b0ux9ae-ux993-ux9a7ux9b0-mcq-ux9aaux9b0ux9b6ux9a8-ux989ux9a4ux9a4ux9b0-ux985ux9a8ux9b6ux9b2ux9a8}

\textbf{১ হতে ৯৯ পর্যন্ত সংখ্যা সমূহের যোগফল কত? /
১+২+৩+৪+...................+৯৯ = কত?}

ক) ৪৬৫০

খ) ৪৭৫০

গ) ৪৮৫০

\textbf{\hl{ঘ) ৪৯৫০ ✔}}

\textbf{\hl{✔ সঠিক উত্তর: ঘ) ৪৯৫০}}

\begin{quote}
\textbf{ব্যাখ্যা:} \emph{এখানে ধারাটি একটি সমান্তর ধারা, যার প্রথম পদ a=1 ,
সাধারণ অন্তর d=1 এবং শেষ পদ p=99 ।}

\emph{পদসংখ্যা n=(শেষ পদ−১ম পদ/সাধারণ অন্তর)+1=(99−1/1)+1=99}

\emph{সমষ্টি =(শেষ পদ+১ম পদ/2)×পদ সংখ্যা}

\emph{⇒ (99+1/2)×99=(100/2)×99=50×99=4950}
\end{quote}

\textbf{১ হতে ১০০ পর্যন্ত সংখ্যা সমূহের যোগফল কত?}

ক) ৪৯৯৯

খ) ৫৫০১

\textbf{\hl{গ) ৫০৫০ ✔}}

ঘ) ৫০০১

\textbf{\hl{✔ সঠিক উত্তর: গ) ৫০৫০}}

\begin{quote}
\textbf{ব্যাখ্যা:} \emph{যোগফল =((১ম পদ+পদ শেষ)/2)×পদ সংখ্যা}

\emph{⇒ ((1+100)/2)×100=(101/2)×100=5050}
\end{quote}

\textbf{৫ থেকে ৩৫ পর্যন্ত সংখ্যাগুলোর যোগফল কত?}

ক) ৬০০

খ) ৬১০

\textbf{\hl{গ) ৬২০ ✔}}

ঘ) ৬৩০

\textbf{\hl{✔ সঠিক উত্তর: গ) ৬২০}}

\begin{quote}
\textbf{ব্যাখ্যা:} \emph{৫ থেকে ৩৫ পর্যন্ত পদসংখ্যা =(শেষ পদ−১ম পদ/সাধারণ
অন্তর)+1}

\emph{⇒ ((35−5/1)+1)টি=31টি}

\emph{∴ ৫ থেকে ৩৫ পর্যন্ত সংখ্যাগুলোর যোগফল =(শেষ পদ+১ম পদ/2)×পদ সংখ্যা}

\emph{⇒ (35+5/2)×31=(40/2)×31=20×31=620}
\end{quote}

\textbf{৫ থেকে ৪৫ পর্যন্ত সংখ্যাগুলোর যোগফল কত?}

\textbf{\hl{ক) ১০২৫ ✔}}

খ) ১০৭৫

গ) ১০৩৫

ঘ) ১০৪৫

\textbf{\hl{✔ সঠিক উত্তর: ক) ১০২৫}}

\begin{quote}
\textbf{ব্যাখ্যা:} \emph{৫ থেকে ৪৫ পর্যন্ত পদসংখ্যা =(শেষ পদ−১ম পদ/সাধারণ
অন্তর)+1}

\emph{⇒ ((45−5/1)+1)টি=41টি}

\emph{৫ থেকে ৪৫ পর্যন্ত সংখ্যাগুলোর সমষ্টি =(শেষ পদ+১ম পদ/2)×পদ সংখ্যা}

\emph{⇒ (45+5/2)×41=25×41=1025}
\end{quote}

\textbf{1, 2, 3, 4-\/-\/-\/-\/-\/-n পর্যন্ত স্বাভাবিক সংখ্যার যোগফল কত?}

ক) n\^{}2

\textbf{\hl{খ) (n(n+1)/2) ✔}}

গ) (n(2n+1)/2)

ঘ) (((n(n+1)/2)))\^{}2

\textbf{\hl{✔ সঠিক উত্তর: খ) (n(n+1)/2)}}

\begin{quote}
\textbf{ব্যাখ্যা:} \emph{স্বাভাবিক সংখ্যার যোগফলের সূত্র হলো:
1+2+3+\ldots+n=(n(n+1)/2)}
\end{quote}

\textbf{1+2+3+\ldots\ldots\ldots\ldots+n= কত?}

\textbf{\hl{ক) (n(n+1)/2) ✔}}

খ) (n\^{}2((n+1))\^{}2/4)

গ) n\^{}2

ঘ) (n(n+1)(n+2)/6)

\textbf{\hl{✔ সঠিক উত্তর: ক) (n(n+1)/2)}}

\begin{quote}
\textbf{ব্যাখ্যা:} \emph{স্বাভাবিক সংখ্যার যোগফলের সূত্র অনুযায়ী,
1+2+3+\ldots+n=(n(n+1)/2)}
\end{quote}

\textbf{1+2+3+8+\ldots\ldots\ldots\ldots+19= কত?}

ক) ১৮০

খ) ১৬৮

গ) ১৮৮

\textbf{\hl{ঘ) ১৯০ ✔}}

\textbf{\hl{✔ সঠিক উত্তর: ঘ) ১৯০}}

\begin{quote}
\textbf{ব্যাখ্যা:} \emph{এখানে ধারাটি ১ থেকে ১৯ পর্যন্ত স্বাভাবিক সংখ্যার
যোগফল।}

\emph{সূত্র অনুযায়ী, (n(n+1)/2) যেখানে n=19}

\emph{⇒ (19(19+1)/2)=(19×20/2)=(380/2)=190}
\end{quote}

\textbf{1+2+3+\ldots+25 এর যোগফল কত?}

ক) ১০০

খ) ২৫০

গ) ৩০০

\textbf{\hl{ঘ) ৩২৫ ✔}}

\textbf{\hl{✔ সঠিক উত্তর: ঘ) ৩২৫}}

\begin{quote}
\textbf{ব্যাখ্যা:} \emph{এখানে n=25}

\emph{সমষ্টি S=(n(n+1)/2)=(25(25+1)/2)=(25×26/2)=25×13=325}
\end{quote}

\textbf{1+2+3+\ldots\ldots\ldots\ldots+50= কত?}

ক) ১২০০

খ) ১২২৫

গ) ১২৫০

\textbf{\hl{ঘ) ১২৭৫ ✔}}

\textbf{\hl{✔ সঠিক উত্তর: ঘ) ১২৭৫}}

\begin{quote}
\textbf{ব্যাখ্যা:} \emph{পদসংখ্যা n=50}

\emph{∴ সমষ্টি =(শেষ পদ+১ম পদ/2)×পদ সংখ্যা}

\emph{⇒ (50+1/2)×50=25×51=1275}
\end{quote}

\textbf{1+3+5+\ldots\ldots\ldots\ldots+(2x−1) কত?}

ক) x(x−1)

খ) (x(x+1)/2)

গ) x(x+1)

\textbf{\hl{ঘ) x\^{}2 ✔}}

\textbf{\hl{✔ সঠিক উত্তর: ঘ) x\^{}2}}

\begin{quote}
\textbf{ব্যাখ্যা:} \emph{প্রথম x সংখ্যক বিজোড় স্বাভাবিক সংখ্যার যোগফল x\^{}2
।}

\emph{∴ 1+3+5+\ldots+(2x−1)=x\^{}2}
\end{quote}

\textbf{1+3+5+7+\ldots{} ধারাটির n -পদের সমষ্টি কত?}

ক) n(n+1)

খ) (n(n+1)/4)

\textbf{\hl{গ) n\^{}2 ✔}}

ঘ) (((n(n+1))\^{}2/2)

\textbf{\hl{✔ সঠিক উত্তর: গ) n\^{}2}}

\begin{quote}
\textbf{ব্যাখ্যা:} \emph{এটি একটি সমান্তর ধারা যার প্রথম পদ a=1 এবং সাধারণ
অন্তর d=3−1=2 ।}

\emph{S\_n=(n/2){[}2a+(n−1)d{]}}

\emph{⇒ (n/2){[}2(1)+(n−1)×2{]}=(n/2){[}2+2n−2{]}=(n/2)(2n)=n\^{}2}
\end{quote}

\textbf{1+5+9+13+\ldots{} ধারাটির n তম পদ কোনটি?}

ক) 3n−3

\textbf{\hl{খ) 4n−3 ✔}}

গ) 4n+3

ঘ) 4n−1

\textbf{\hl{✔ সঠিক উত্তর: খ) 4n−3}}

\begin{quote}
\textbf{ব্যাখ্যা:} \emph{দেওয়া আছে, 1+5+9+13+\ldots{}}

\emph{ধারাটির প্রথম পদ, a=1}

\emph{সাধারণ অন্তর, d=5−1=4}

\emph{ধারাটির n তম পদ =a+(n−1)d}

\emph{⇒ 1+(n−1)4=1+4n−4=4n−3}
\end{quote}

\textbf{1+9+17+\ldots+81= কত?}

ক) ১৯১

\textbf{\hl{খ) ৮৬১ ✔}}

গ) ৭৬১

ঘ) ৬৬১

\textbf{\hl{✔ সঠিক উত্তর: খ) ৮৬১}}

\begin{quote}
\textbf{ব্যাখ্যা:} \emph{এখানে ধারাটির প্রথম পদ a=1 , শেষ পদ p=81 এবং
সাধারণ অন্তর d=9−1=8}

\emph{পদ সংখ্যা =(শেষ পদ−১ম পদ/সাধারণ অন্তর)+1}

\emph{⇒ (81−1/8)+1=(80/8)+1=10+1=21}

\emph{সমষ্টি =((১ম পদ+শেষ পদ)×পদ সংখ্যা/2)}

\emph{⇒ ((1+81)×21/2)=(82×21/2)=41×21=861}
\end{quote}

\textbf{1+4+7+10+\ldots+73 সমান্তর ধারাটির যোগফল কত হবে?}

ক) ১২৫

খ) ১০২৫

গ) ১১২৫

\textbf{\hl{ঘ) ৯২৫ ✔}}

\textbf{\hl{✔ সঠিক উত্তর: ঘ) ৯২৫}}

\begin{quote}
\textbf{ব্যাখ্যা:} \emph{ব্যাখ্যা: ১ম পদ a=1 , সাধারণ অন্তর d=4−1=3 এবং শেষ
পদ p=73}

\emph{∴ n তম পদ =a+(n−1)d}

\emph{⇒ 73=1+(n−1)3}

\emph{⇒ 72=(n−1)3}

\emph{⇒ n−1=24}

\emph{⇒ n=25}

\emph{সমষ্টি =(n/2){[}2a+(n−1)d{]}}

\emph{=(25/2){[}2(1)+(25−1)×3{]}}

\emph{⇒ (25/2){[}2+24×3{]}}

\emph{=(25/2){[}2+72{]}}

\emph{=(25/2)×74 =25×37 =925}
\end{quote}

\textbf{3+6+9+12+\ldots{} ধারাটি 12টি পদের যোগফল কত?}

ক) 230

খ) 231

গ) 232

\textbf{\hl{ঘ) 234 ✔}}

\textbf{\hl{✔ সঠিক উত্তর: ঘ) 234}}

\begin{quote}
\textbf{ব্যাখ্যা:} \emph{এটি একটি সমান্তর ধারা যার প্রথম পদ a=3 , সাধারণ
অন্তর d=6−3=3 এবং পদসংখ্যা n=12}

\emph{∴ 12টি পদের সমষ্টি =(n/2){[}2a+(n−1)d{]}}

\emph{⇒ (12/2){[}2(3)+(12−1)×3{]}=6{[}6+11×3{]}}

\emph{⇒ 6(6+33)=6×39=234}
\end{quote}

\textbf{11+18+25+32+\ldots{} ধারাটির ১৫টি পদের সমষ্টি কত?}

\textbf{\hl{ক) ৯০০ ✔}}

খ) ১০৯

গ) ১৮০০

ঘ) ১০৯

\textbf{\hl{✔ সঠিক উত্তর: ক) ৯০০}}

\begin{quote}
\textbf{ব্যাখ্যা:} \emph{ধারাটির প্রথম পদ a=11 , সাধারণ অন্তর d=18−11=7 এবং
পদসংখ্যা n=15}

\emph{∴ ১৫টি পদের সমষ্টি S\_n=(n/2){[}2a+(n−1)d{]}}

\emph{⇒ (15/2){[}2(11)+(15−1)×7{]}=(15/2){[}22+14×7{]}}

\emph{⇒ (15/2)(22+98)=(15/2)×120=15×60=900}
\end{quote}

\textbf{29+25+21+\ldots−23= কত?}

ক) ৮২

খ) ৭২

গ) ৫২

\textbf{\hl{ঘ) ৪২ ✔}}

\textbf{\hl{✔ সঠিক উত্তর: ঘ) ৪২}}

\begin{quote}
\textbf{ব্যাখ্যা:} \emph{এখানে, প্রথম পদ a=29 , সাধারণ অন্তর d=25−29=−4 এবং
শেষ পদ ( n তম পদ) =−23}

\emph{আমরা জানি, n তম পদ =a+(n−1)d}

\emph{⇒ −23=29+(n−1)(−4)}

\emph{⇒ −23−29=(n−1)(−4)}

\emph{⇒ −52=(n−1)(−4)}

\emph{⇒ n−1=13}

\emph{⇒ n=14}

\emph{∴ সমষ্টি =(n/2)(a+p)=(14/2)(29+(−23))=7(29−23)=7×6=42}
\end{quote}

\textbf{৩৭ + ৩৬ + ৩৫ + ..... + ২১ = কত?}

ক) ৪১৯

\textbf{\hl{খ) ৪৯৩ ✔}}

গ) ৪১৪

ঘ) ৪১৫

\textbf{\hl{✔ সঠিক উত্তর: খ) ৪৯৩}}

\begin{quote}
\textbf{ব্যাখ্যা:} \emph{পদসংখ্যা = (শেষপদ−১মপদ/সাধারণঅন্তর)+১}

\emph{= (২১−৩৭/−১) + ১ = ১৭}

\emph{যোগফল = (১ম পদ+শেষ পদ/2) × পদসংখ্যা}

\emph{= ((৩৭+২১)/২) × ১৭ = ৪৯৩}
\end{quote}

\textbf{log a + log a\^{}2 + log a\^{}3 + ....... ধারটির প্রথম ৭টি পদের
সমষ্টি কত?}

\textbf{\hl{ক) 28 log a ✔}}

খ) 56 log a

গ) 7 log a

ঘ) 8 log a

\textbf{\hl{✔ সঠিক উত্তর: ক) 28 log a}}

\begin{quote}
\textbf{ব্যাখ্যা:} \emph{log a + log a\^{}2 + log a\^{}3 + .... 7 তম পদ
পর্যন্ত}

\emph{= log a + 2 log a + 3 log a + .... 7 তম পদ পর্যন্ত}

\emph{= log a (1 + 2 + 3 ......... 7)}

\emph{= log a ( ((7+1)×7/2) )}

\emph{∴ সমান্তর ধারার সমষ্টি = (পদসংখ্যা×(শেষ পদ+১ম পদ)/2)}

\emph{= log a ( (56/2) ) = 28 log a}
\end{quote}

\textbf{একটি সমান্তর ধারার 12 তম পদ 77 হলে, তার প্রথম 23 পদের সমষ্টি কত?}

\textbf{\hl{ক) 1771 ✔}}

খ) 1176

গ) 1056

ঘ) 2025

\textbf{\hl{✔ সঠিক উত্তর: ক) 1771}}

\begin{quote}
\textbf{ব্যাখ্যা:} \emph{দেয়া আছে, 12তম পদ = 77}

\emph{⇒ a + (12 − 1)d = 77}

\emph{⇒ a + 11d = 77}

\emph{এখন, প্রথম 23 পদের সমষ্টি,}

\emph{S\_23 = (23/2) \{2a + (23 − 1)d\}}

\emph{= (23/2) (2a + 22d)}

\emph{= 23 × (a + 11d) = 23 × 77 = 1771}
\end{quote}

\textbf{৯টি ধারাবাহিক বিজোড় সংখ্যার পঞ্চম সংখ্যাটি হল −১৫ সবগুলো সংখ্যার
সমষ্টি কত?}

ক) ১১৮

খ) ১২৫

গ) ১৩৫

\textbf{\hl{ঘ) −১৩৫ ✔}}

ঙ) কোনোটিই নয়

\textbf{\hl{✔ সঠিক উত্তর: ঘ) −১৩৫}}

\begin{quote}
\textbf{ব্যাখ্যা:} \emph{ধরি, সংখ্যাগুলো যথাক্রমে x − ৮, x − ৬, x − ৪, x − ২,
x, x + ২, x + ৪, x + ৬ এবং x + ৮}

\emph{∴ x = −১৫ (∴ x হচ্ছে পঞ্চম সংখ্যা)}

\emph{∴ সমষ্টি = x − ৮ + x − ৬ + x − ৪ + x − ২ + x + x + ২ + x + ৪ + x +
৬ + x + ৮}

\emph{= ৯x = ৯ × (−১৫) = −১৩৫}
\end{quote}

\textbf{৯ + ৭ + ৫ -\/- ধারাটি প্রথম n সংখ্যক পদের যোগফল মাইনাস ১৪৪ হলে,
n=কত?}

\textbf{\hl{ক) ১৮ ✔}}

খ) ১৬

গ) ১২

ঘ) ১৪

\textbf{\hl{✔ সঠিক উত্তর: ক) ১৮}}

\begin{quote}
\textbf{ব্যাখ্যা:} \emph{প্রথম পদ a এবং সাধারণ অন্তর d বিশিষ্ট সমান্তর ধারার
প্রথম n সংখ্যক পদের সমষ্টি = (n/2)\{2a+(n−1)d\}}

\emph{∴ (n/2)\{2×9+(n−1)(−2)\} = −১৪৪}

\emph{⇒ ১০n − n\^{}2 + n = −১৪৪}

\emph{⇒ n\^{}2 − ১০n − ১৪৪ = ০}

\emph{⇒ n\^{}2 − ১৮n + ৮n − ১৪৪ = ০}

\emph{⇒ n(n − ১৮) + ৮(n − ১৮) = ০}

\emph{⇒ (n − ১৮)(n + ৮) = ০}

\emph{∴ n = ১৮, −৮}

\emph{n = −৮ গ্রহণযোগ্য নয়। ∴ n = ১৮}
\end{quote}

\textbf{১২ ও ৯৬ এর মধ্যে (এই দুটি সংখ্যাসহ) কয়টি সংখ্যা ৪ দ্বারা বিভাজ্য?}

ক) ২১

খ) ২৩

গ) ২৩

\textbf{\hl{ঘ) ২২ ✔}}

\textbf{\hl{✔ সঠিক উত্তর: ঘ) ২২}}

\begin{quote}
\textbf{ব্যাখ্যা:} \emph{এটি মূলত সমান্তর ধারা যেখানে,}

\emph{১ম পদ = ১২, শেষপদ = ৯৬, সাধারণ অন্তর = ৪}

\emph{এবং ধারাটি = ১২, ১৬, ২০, ................., ৯৬}

\emph{৪ দ্বারা বিভাজ্য সংখ্যা (পদসংখ্যা) = (শেষপদ−১মপদ/সাধারণঅন্তর)+১}

\emph{= (৯৬−১২/৪) + ১ = ২১ + ১ = ২২}
\end{quote}

\textbf{১\^{}২ + ২\^{}২ + ৩\^{}২ + ............. + ৫০\^{}২ = ?}

ক) ৩৫৭২৫

\textbf{\hl{খ) ৪২৯২৫ ✔}}

গ) ৪৫৫০0

ঘ) ৪৭২২৫

\textbf{\hl{✔ সঠিক উত্তর: খ) ৪২৯২৫}}

\begin{quote}
\textbf{ব্যাখ্যা:} \emph{ধারাবাহিক ক্রমিক সংখ্যার বর্গের সমষ্টি}

\emph{= (n(n+1)(2n+1)/6)}

\emph{= (50×51×101/6) = ৪২৯২৫}
\end{quote}

\textbf{০ হতে ৪৯ পর্যন্ত ক্রমিক সংখ্যাগুলোর গড় কত?}

ক) ২৩

খ) ২৫

\textbf{\hl{গ) ২৪.৫ ✔}}

ঘ) ২৬.৫

\textbf{\hl{✔ সঠিক উত্তর: গ) ২৪.৫}}

\begin{quote}
\textbf{ব্যাখ্যা:} \emph{ক্রমিক সংখ্যা ও সমান্তর ধারার ক্ষেত্রে,}

\emph{০ হতে ৪৯ পর্যন্ত ক্রমিক সংখ্যাগুলোর গড় = (শেষ পদ+১ম পদ/2)}

\emph{= (৪৯+০/2) = ২৪.৫}
\end{quote}

\textbf{১ হতে ৪৯ পর্যন্ত ক্রমিক সংখ্যাগুলোর গড় কত?}

ক) ২৩

খ) ২৪.৫

\textbf{\hl{গ) ২৫ ✔}}

ঘ) ২৬.৫

\textbf{\hl{✔ সঠিক উত্তর: গ) ২৫}}

\begin{quote}
\textbf{ব্যাখ্যা:} \emph{১ হতে ৪৯ পর্যন্ত সংখ্যাগুলোর গড় = (শেষ পদ+১ম পদ/2)}

\emph{= (৪৯+১/2) = ২৫}
\end{quote}

\textbf{১ থেকে ১০০ পর্যন্ত সংখ্যাগুলোর গড় কত?}

ক) ৪৯৫০

খ) ৫০৫০

গ) ৫০০০

\textbf{\hl{ঘ) ৫০.৫০ ✔}}

\textbf{\hl{✔ সঠিক উত্তর: ঘ) ৫০.৫০}}

\begin{quote}
\textbf{ব্যাখ্যা:} \emph{১ থেকে ১০০ পর্যন্ত সংখ্যাগুলোর গড় = (শেষ পদ+১ম পদ/2)}

\emph{= (১০০+১/2) = (১০১/2) = ৫০.৫০}
\end{quote}

\textbf{১ থেকে ৯৯ পর্যন্ত সংখ্যাসমূহের গড় কত?}

ক) ১০

খ) ২৫

\textbf{\hl{গ) ৫০ ✔}}

ঘ) ১০০

\textbf{\hl{✔ সঠিক উত্তর: গ) ৫০}}

\begin{quote}
\textbf{ব্যাখ্যা:} \emph{১ থেকে ৯৯ পর্যন্ত সংখ্যাসমূহের গড় = (শেষপদ+১মপদ/2)}

\emph{= (৯৯+১/2) = ৫০}
\end{quote}

\textbf{১ হতে ৫০ পর্যন্ত বেজোড় সংখ্যাগুলোর গড় কত?}

ক) ২০

খ) ২৪.৫

গ) ২৫.০

\textbf{\hl{ঘ) ২৫ ✔}}

\textbf{\hl{✔ সঠিক উত্তর: ঘ) ২৫}}

\begin{quote}
\textbf{ব্যাখ্যা:} \emph{১ম পদ a = ১, n = ২৫ (পদসংখ্যা)}

\emph{প্রতি পদের বৃদ্ধি, d = ২}

\emph{∴ যোগফল, S = (n/2)\{2a+(n−1)d\}}

\emph{= (25/2)\{2×1+(25−1)×2\} = ৬২৫}

\emph{∴ গড় = (625/25) = ২৫}
\end{quote}

\textbf{১ থেকে ৫১ পর্যন্ত ক্রমিক সংখ্যাগুলোর গড় কত?}

\textbf{\hl{ক) ২৬ ✔}}

খ) ২৮

গ) ৩০

ঘ) ৩১

\textbf{\hl{✔ সঠিক উত্তর: ক) ২৬}}

\begin{quote}
\textbf{ব্যাখ্যা:} \emph{১ হতে ৫০ পর্যন্ত ক্রমিক সংখ্যাগুলোর গড় =
(শেষপদ+১মপদ/2)}

\emph{= (৫১+১/2) = ২৬}
\end{quote}

\textbf{১, ৫, ৯, ......, ৮১ ধারাটির সংখ্যাগুলোর গড় কত?}

\textbf{\hl{ক) ৪১ ✔}}

খ) ৩৯

গ) ৪২

ঘ) ৪০

\textbf{\hl{✔ সঠিক উত্তর: ক) ৪১}}

\begin{quote}
\textbf{ব্যাখ্যা:} \emph{১ম পদ a = ১}

\emph{সাধারণ অন্তর d = ৫ − ১ = ৪}

\emph{শেষ পদ = a + (n − ১) × d}

\emph{⇒ ৮১ = ১ + (n − ১) × ৪}

\emph{⇒ ৪(n − ১) = ৮১ − ১}

\emph{⇒ n − ১ = (৮০/৪)}

\emph{⇒ n = ২০ + ১}

\emph{∴ n = ২১}

\emph{∴ সমষ্টি, S = (n/2)(১মপদ+শেষপদ)}

\emph{= (২১/২) × (১ + ৮১) = (২১/২) × ৮২ = ৮৬১}

\emph{∴ গড় = (সমষ্টি/পদসংখ্যা) = (৮৬১/২১) = ৪১}
\end{quote}

\textbf{১ থেকে ১০০ পর্যন্ত সংখ্যার মধ্যে যেসব সংখ্যার ডানে ৩ আছে, তাদের গড়
কত?}

ক) ৪৪

খ) ৪৫

গ) ৪৬

ঘ) ৪৭

\textbf{\hl{ঙ) কোনোটিই নয় ✔}}

\textbf{\hl{✔ সঠিক উত্তর: ঙ) কোনোটিই নয়}}

\begin{quote}
\textbf{ব্যাখ্যা:} \emph{১ থেকে ১০০ পর্যন্ত সংখ্যার মধ্যে যেসব সংখ্যার ডানে ৩
আছে, তারা নিম্নরূপ:}

\emph{3, 13, 23, 33, 43, 53, 63, 73, 83 এবং 93}

\emph{এটি একটি সমান্তর ধারা,}

\emph{∴ সংখ্যাগুলোর গড় = (শেষপদ+১মপদ/2) = (93+3/2) = (96/2) = ৪৮}
\end{quote}

\textbf{6, 8 ও 10 এর গাণিতিক গড় 7, 9 ও x এর গাণিতিক গড়ের সমান হলে, x এর
মান কত?}

ক) 5

\textbf{\hl{খ) 8 ✔}}

গ) 6

ঘ) 10

\textbf{\hl{✔ সঠিক উত্তর: খ) 8}}

\begin{quote}
\textbf{ব্যাখ্যা:} \emph{6, 8 ও 10 এর গড় = (6+8+10/3) = 8}

\emph{∴ 7, 9, x এর গড় = (7+9+x/3) = 8}

\emph{⇒ (16+x/3) = 8}

\emph{⇒ 16 + x = 24}

\emph{∴ x = 24 − 16 = 8}
\end{quote}

\textbf{একটি সমান্তর অনুক্রমে 5ম পদটি 18 এবং 5-টি পদের যোগফল 75 হলে প্রথম
পদটি কত?}

ক) 2

খ) 10

গ) 4

\textbf{\hl{ঘ) 12 ✔}}

\textbf{\hl{✔ সঠিক উত্তর: ঘ) 12}}

\begin{quote}
\textbf{ব্যাখ্যা:} \emph{দেয়া আছে, 5ম পদ, a\_5 = a + (5 − 1)d = 18}

\emph{⇒ a + 4d = 18 ............ (i)}

\emph{প্রথম 5 পদের যোগফল, S\_5 = (5/2) \{2a + (5 − 1)d\} = 75}

\emph{⇒ (5/2) (2a + 4d) = 75}

\emph{⇒ 5(a + 2d) = 75}

\emph{⇒ a + 2d = 15 ............ (ii)}

\emph{(i) ও (ii) নং সমীকরণ হতে পাই,}

\emph{a + 4d = 18}

\emph{a + 2d = 15}

\emph{− ~ − ~ −}

\emph{-\/-\/-\/-\/-\/-\/-\/-\/-\/-\/-\/-\/-\/-\/-\/-}

\emph{2d = 3 ⇒ d = 1.5}

\emph{d এর মান (ii) নং সমীকরণে বসিয়ে পাই,}

\emph{a + 2(1.5) = 15}

\emph{⇒ a + 3 = 15}

\emph{∴ a = 12}
\end{quote}

\textbf{একটি সমান্তর অনুক্রমে সাধারণ অন্তর 10 এবং 6-তম পদটি 52 হলে 15-তম
পদটি}

ক) 140

\textbf{\hl{খ) 142 ✔}}

গ) 148

ঘ) 150

\textbf{\hl{✔ সঠিক উত্তর: খ) 142}}

\begin{quote}
\textbf{ব্যাখ্যা:} \emph{r তম পদ = প্রথম পদ + (r − 1) × সাধারণ অন্তর}

\emph{ধারটির 6 তম পদ = a + (6 − 1) 10}

\emph{বা, 52 = a + 5 × 10}

\emph{বা, a = 52 − 50}

\emph{∴ a = 52 − 50 = 2}

\emph{আবার, ধারটির 15 তম পদ = a + (15 − 1) 10 = 2 + 14 × 10 = 2 + 140 =
142}
\end{quote}

\textbf{একটি সমান্তর ধারার ২য় পদ 4 এবং দশম পদ 15 হলে ধারটির ১ম পদ কত?}

ক) 1.18

খ) 1.27

গ) 1.38

\textbf{\hl{ঘ) 2.63 ✔}}

\textbf{\hl{✔ সঠিক উত্তর: ঘ) 2.63}}

\begin{quote}
\textbf{ব্যাখ্যা:} \emph{ধরি, ধারটির প্রথম পদ a এবং সাধারণ অন্তর d}

\emph{প্রশ্নমতে, a + (2 − 1)d = 4}

\emph{⇒ a + d = 4 ....... (i)}

\emph{এবং a + (10 − 1)d = 15}

\emph{⇒ a + 9d = 15 ....... (ii)}

\emph{এখন সমীকরণ (ii) থেকে (i) বিয়োগ করে পাই,}

\emph{a + 9d = 15}

\emph{a + d = 4}

\emph{− ~ − ~ −}

\emph{-\/-\/-\/-\/-\/-\/-\/-\/-\/-\/-\/-\/-\/-\/-\/-}

\emph{8d = 11}

\emph{∴ d = (11/8)}

\emph{এখন d এর মান সমীকরণ (i) এ বসিয়ে পাই,}

\emph{a + d = 4}

\emph{⇒ a + (11/8) = 4 {[}মান বসিয়ে{]}}

\emph{∴ a = 4 − (11/8) = 4 − 1.375 = 2.63}
\end{quote}

\textbf{একটি সমান্তরাল ধারার সাধারণ অন্তর 9 এবং 7ম পদ 60 হলে 12তম পদটি
কত?}

ক) 89

খ) 100

গ) 108

\textbf{\hl{ঘ) 105 ✔}}

\textbf{\hl{✔ সঠিক উত্তর: ঘ) 105}}

\begin{quote}
\textbf{ব্যাখ্যা:} \emph{ধরি, ধারার ১ম পদ = a, সাধারণ অন্তর d = 9}

\emph{∴ 7তম পদ = a + (7 − 1)d}

\emph{⇒ 60 = a + 6 × 9}

\emph{⇒ a = 60 − 54}

\emph{∴ a = 6}

\emph{∴ 12তম পদ = a + (12 − 1)d = 6 + 11 × 9 = 6 + 99 = 105}
\end{quote}

\textbf{কোন সমান্তর প্রগমনে প্রথম দুটি সংখ্যা ৫ ও ১৭ হয় তবে তৃতীয় সংখ্যাটি কত?}

ক) ২২

খ) ২৫

\textbf{\hl{গ) ২৯ ✔}}

ঘ) ৪৫

\textbf{\hl{✔ সঠিক উত্তর: গ) ২৯}}

\begin{quote}
\textbf{ব্যাখ্যা:} \emph{ধরি, তৃতীয় সংখ্যা = ক}

\emph{যেহেতু সমান্তর প্রগমন, ৩য় পদ − ২য় পদ = ২য় পদ − ১ম পদ}

\emph{⇒ ক − ১৭ = ১৭ − ৫}

\emph{⇒ ক = ১৭ + ১৭ − ৫ = ২৯}
\end{quote}

\textbf{1 + 5 + 9 + 13 + ...... ধারাটির 15 তম পদ হবে}

ক) 61

খ) 53

\textbf{\hl{গ) 57 ✔}}

ঘ) 65

\textbf{\hl{✔ সঠিক উত্তর: গ) 57}}

\begin{quote}
\textbf{ব্যাখ্যা:} \emph{ধারটির 15 তম পদ = a + (15 − 1) d = 1 + 14 × 4 = 1
+ 56 = 57}
\end{quote}

\textbf{128, 64, 32, ............ ধারাটির 8ম পদ কোনটি?}

\textbf{\hl{ক) 1 ✔}}

খ) 0

গ) 2

ঘ) (1/2)

\textbf{\hl{✔ সঠিক উত্তর: ক) 1}}

\begin{quote}
\textbf{ব্যাখ্যা:} \emph{এটি একটি গুণোত্তর ধারা যার ১ম পদ a = 128 এবং সাধারণ
অনুপাত r = (64/128) = (1/2)}

\emph{∴ 8ম পদ = a r\^{}8 − 1 = 128 × ( (1/2) )\^{}7 = 128 × (1/128) = 1}
\end{quote}

\textbf{৩ + ৬ + ৯ + ....... ধারাটির কততম পদ ৩৩?}

ক) ১০

খ) ১২

\textbf{\hl{গ) ১১ ✔}}

ঘ) ১৩

\textbf{\hl{✔ সঠিক উত্তর: গ) ১১}}

\begin{quote}
\textbf{ব্যাখ্যা:} \emph{দেয়া আছে, ধারটির ১ম পদ a = ৩}

\emph{সাধারণ অন্তর d = ৬ − ৩ = ৩}

\emph{এবার ধরি, n তম পদ ৩৩}

\emph{প্রশ্নমতে, a + (n − 1)d = ৩৩}

\emph{⇒ ৩ + (n − 1) × ৩ = ৩৩}

\emph{⇒ (n − 1) × ৩ = ৩৩ − ৩ = ৩০}

\emph{⇒ n − ১ = (৩০/৩) = ১০}

\emph{∴ n = ১০ + ১ = ১১}
\end{quote}

\textbf{৫ + ৮ + ১১ + ... ... ... ধারাটির কোন পদ ৩০২?}

ক) ৯১

খ) ৯৫

গ) ৯৮

\textbf{\hl{ঘ) ১০০ ✔}}

\textbf{\hl{✔ সঠিক উত্তর: ঘ) ১০০}}

\begin{quote}
\textbf{ব্যাখ্যা:} \emph{দেয়া আছে, প্রথম পদ, a = ৫}

\emph{সাধারণ অন্তর, d = (৮ − ৫) = ৩}

\emph{মনে করি, r তম পদ = ৩০২}

\emph{∴ a + (r − ১)d = ৩০২}

\emph{⇒ ৫ + (r − ১)৩ = ৩০২}

\emph{⇒ ৫ + ৩r − ৩ = ৩০২}

\emph{⇒ ৩r = ৩০২ − ৫ + ৩}

\emph{⇒ ৩r = (৩০০/৩) ∴ r = ১০০}

\emph{সুতরাং ১০০ তম পদ = ৩০২}
\end{quote}

\textbf{একটি ধারার ১ম সংখ্যাটি ২ এবং এরপরের প্রতিটি সংখ্যা ঠিক আগের সংখ্যা
থেকে ৮ বেশি হলে ধারার ২০২তম সংখ্যাটি কত?}

ক) ৭৯৮

খ) ৮০০

গ) ৮০২

\textbf{\hl{ঘ) ৮০৬ ✔}}

\textbf{\hl{✔ সঠিক উত্তর: ঘ) ৮০৬}}

\begin{quote}
\textbf{ব্যাখ্যা:} \emph{সমাধান: ১মপদ, a = ২; সাধারণ অন্তর, d = ৮}

\emph{আমরা জানি, n তম পদ = a + (n − 1)d}

\emph{∴ ২০২ তম সংখ্যাটি = ২ + (২০২ − 1) × ৮ = ২ + ৮০৪ = ৮০৬}
\end{quote}

\textbf{কোনো সমান্তর ধারার p-তম পদ q এবং q-তম পদ p হলে (p+q) তম পদ কত?}

ক) pq

খ) p+q

গ) pq(p+q)

\textbf{\hl{ঘ) 0 ✔}}

\textbf{\hl{✔ সঠিক উত্তর: ঘ) 0}}

\begin{quote}
\textbf{ব্যাখ্যা:} \emph{ধারটির p তম পদ = a + (p − 1)d}

\emph{এবং q তম পদ = a + (q − 1)d}

\emph{এবং (p+q) তম পদ = a + (p + q − 1)d}

\emph{প্রশ্নমতে, a + (p − 1)d = q ......... (i)}

\emph{এবং a + (q − 1)d = p ......... (ii)}

\emph{(i) হতে (ii) বিয়োগ করে পাই, (p − q)d = q − p ⇒ d = − 1}

\emph{d এর মান (i) এ বসিয়ে পাই, a + (p − 1)(− 1) = q ⇒ a = q + p − 1}

\emph{∴ (p+q) তম পদ = (q + p − 1) + (p + q − 1)(− 1) = q + p − 1 − p − q
+ 1 = 0}
\end{quote}

\textbf{একটি সমান্তর ধারার প্রথম পদ 1, শেষ পদ 99 এবং সমষ্টি 2500 হলে
ধারাটির সাধারণ অন্তর হবে}

ক) 4

\textbf{\hl{খ) 2 ✔}}

গ) 3

ঘ) 6

\textbf{\hl{✔ সঠিক উত্তর: খ) 2}}

\begin{quote}
\textbf{ব্যাখ্যা:} \emph{সমষ্টি = (শেষ পদ+১ম পদ/2) × পদসংখ্যা}

\emph{⇒ 2500 = (99+1/2) × পদসংখ্যা}

\emph{⇒ পদসংখ্যা = (2500×2/100) = 50}

\emph{আবার, পদসংখ্যা = (শেষ পদ−১ম পদ/সাধারণ অন্তর) + 1}

\emph{⇒ 50 = (99−1/সাধারণ অন্তর) + 1}

\emph{⇒ (98/সাধারণ অন্তর) = 49}

\emph{∴ সাধারণ অন্তর = (98/49) = 2}
\end{quote}

\textbf{1\^{}2 + 2\^{}2 + 3\^{}2 + ........ + 50\^{}2 = কত?}

ক) 35725

\textbf{\hl{খ) 42925 ✔}}

গ) 45500

ঘ) 47225

\textbf{\hl{✔ সঠিক উত্তর: খ) 42925}}

\begin{quote}
\textbf{ব্যাখ্যা:} \emph{আমরা জানি, প্রথম n সংখ্যক স্বাভাবিক সংখ্যার বর্গের
সমষ্টি = (n(n+1)(2n+1)/6)}

\emph{এখানে, n = 50}

\emph{∴ সমষ্টি = (50(50+1)(2×50+1)/6)}

\emph{= (50×51×101/6) = 42925}
\end{quote}

\textbf{1\^{}2 + 3\^{}2 + 5\^{}2 + ........ + 31\^{}2 = কত?}

ক) 258

\textbf{\hl{খ) 10416 ✔}}

গ) 254

ঘ) 252

\textbf{\hl{✔ সঠিক উত্তর: খ) 10416}}

\begin{quote}
\textbf{ব্যাখ্যা:} \emph{আমরা জানি, প্রথম n সংখ্যক স্বাভাবিক সংখ্যার বর্গের
সমষ্টি = (n(n+1)(2n+1)/6)}

\emph{এখানে, n = 31}

\emph{∴ সমষ্টি = (31(31+1)(2×31+1)/6)}

\emph{= (31×32×63/6) = 10416}
\end{quote}

\textbf{1\^{}2 + 2\^{}2 + 3\^{}2 + .................... + 50\^{}2 = কত?}

ক) 35725

\textbf{\hl{খ) 42925 ✔}}

গ) 45500

ঘ) 47225

\textbf{\hl{✔ সঠিক উত্তর: খ) 42925}}

\begin{quote}
\textbf{ব্যাখ্যা:} \emph{ধারটির সমষ্টি = (1/6)n(n+1)(2n+1) =
(1/6)×50(50+1)(2×50+1)}

\emph{= (50×51×101/6) = 42925 (এখানে n, পদ সংখ্যা = 50)}
\end{quote}

\textbf{2\^{}2 + 4\^{}2 + 6\^{}2 + ........ + (4n)\^{}2 = কত?}

\textbf{\hl{ক) (4/3)n(2n+1)(4n+1) ✔}}

খ) (3/4)n(2n+1)(4n+1)

গ) (4/3)(2n+1)(4n+1)

ঘ) (4/3)n(2n+1)

\textbf{\hl{✔ সঠিক উত্তর: ক) (4/3)n(2n+1)(4n+1)}}

\begin{quote}
\textbf{ব্যাখ্যা:} \emph{2\^{}2 + 4\^{}2 + 6\^{}2 + ........ + (4n)\^{}2}

\emph{= 2\^{}2(1\^{}2 + 2\^{}2 + 3\^{}2 + ........ + n\^{}2) {[}যেখানে
শেষ পদ (4n)\^{}2 এর জন্য ভেতরের পদ হবে n{]}}

\emph{= 4 × (n(n+1)(2n+1)/6) {[}যেখানে শেষ পদ 2n হলে n এর জায়গায় 2n
বসবে{]}}

\emph{= 4 × (2n(2n+1)(2×2n+1)/6)}

\emph{= 4 × (2n(2n+1)(4n+1)/6)}

\emph{= (8/6)n(2n+1)(4n+1)}

\emph{= (4/3)n(2n+1)(4n+1)}
\end{quote}

\textbf{1\^{}3 + 2\^{}3 + 3\^{}3 ... +n\^{}3 ধারাটির যোগফল কত?}

ক) (n(n+1)/2)

খ) (n(n+1)(2n+1)/6)

গ) (n(n+1)/3)

\textbf{\hl{ঘ) (((n(n+1)/2)))\^{}2 ✔}}

\textbf{\hl{✔ সঠিক উত্তর: ঘ) (((n(n+1)/2)))\^{}2}}

\begin{quote}
\textbf{ব্যাখ্যা:} \emph{প্রথম n সংখ্যক স্বাভাবিক সংখ্যার ঘনের সমষ্টির সূত্র হলো
(((n(n+1)/2)))\^{}2}
\end{quote}

\textbf{1\^{}3 + 2\^{}3 + 3\^{}3 + ........ + (60)\^{}3 = কত?}

ক) 3348300

খ) 3348600

গ) 3348800

\textbf{\hl{ঘ) 3348900 ✔}}

\textbf{\hl{✔ সঠিক উত্তর: ঘ) 3348900}}

\begin{quote}
\textbf{ব্যাখ্যা:} \emph{ধারটির যোগফল = (((n(n+1)/2)))\^{}2 =
(\{(60(60+1)/2)\})\^{}2}

\emph{= (30 × 61)\^{}2 = (1830)\^{}2 = 3348900}
\end{quote}

\textbf{প্রথম n সংখ্যক স্বাভাবিক সংখ্যার ঘনের সমষ্টি 225 হলে n-- এর মান কত?}

ক) n=3

\textbf{\hl{খ) n=5 ✔}}

গ) n=6

ঘ) n=7

\textbf{\hl{✔ সঠিক উত্তর: খ) n=5}}

\begin{quote}
\textbf{ব্যাখ্যা:} \emph{প্রথম n সংখ্যক স্বাভাবিক সংখ্যার ঘনের সমষ্টি 225;}

\emph{শর্তমতে, (((n(n+1)/2)))\^{}2 = 225}

\emph{⇒ (n(n+1)/2) = 15}

\emph{⇒ n(n + 1) = 30}

\emph{⇒ n\^{}2 + n − 30 = 0}

\emph{⇒ n\^{}2 + 6n − 5n − 30 = 0}

\emph{⇒ (n − 5)(n + 6) = 0}

\emph{∴ n = 5 (কেননা n এর মান ঋণাত্মক হতে পারে না)}
\end{quote}

\textbf{2, 4, 8, 16, 32, ............256 ধারাটি}

ক) সমান্তর

\textbf{\hl{খ) গুণোত্তর ✔}}

গ) অসংজ্ঞায়িত

ঘ) অসীম

\textbf{\hl{✔ সঠিক উত্তর: খ) গুণোত্তর}}

\begin{quote}
\textbf{ব্যাখ্যা:} \emph{এখানে ধারাটির যেকোনো পদের সাথে পূর্ববর্তী পদের অনুপাত
সর্বদা সমান (যেমন: 4 / 2 = 2, 8 / 4 = 2)।}

\emph{⇒ তাই এটি একটি গুণোত্তর ধারা।}
\end{quote}

\textbf{1 + (1/3) + (1/9) + ............. ধারাটির ১ম 5টি পদের সমষ্টি কত?}

\textbf{\hl{ক) (121/81) ✔}}

খ) (191/81)

গ) (81/121)

ঘ) −(121/81)

\textbf{\hl{✔ সঠিক উত্তর: ক) (121/81)}}

\begin{quote}
\textbf{ব্যাখ্যা:} \emph{এটি একটি গুণোত্তর ধারা যার প্রথম পদ a = 1 এবং সাধারণ
অনুপাত r = (1/3) (যেহেতু r \textless{} 1)}

\emph{⇒ সমষ্টি S\_n = (a(1−r\^{}n)/1−r)}

\emph{⇒ 5টি পদের সমষ্টি = (1×(1−(((1/3)))\^{}5)/1−(1/3))}

\emph{⇒ = (1−(1/243)/(2/3)) = ((242/243)/(2/3))}

\emph{∴ = (242/243)×(3/2) = (121/81)}
\end{quote}

\textbf{1 + (1/2) + (1/4) + (1/8) + ......... ধারাটির প্রথম সাতটি পদের
যোগফল কত?}

ক) (33/31)

খ) (63/32)

গ) (63/64)

\textbf{\hl{ঘ) (127/64) ✔}}

\textbf{\hl{✔ সঠিক উত্তর: ঘ) (127/64)}}

\begin{quote}
\textbf{ব্যাখ্যা:} \emph{গুণোত্তর ধারার ১ম পদ a = 1}

\emph{সাধারণ অনুপাত r = ((1/2)/1) = (1/2) \textless{} 1}

\emph{∴ S\_n = (a(1−r\^{}n)/1−r)}

\emph{∴ S\_7 = (1×(1−(((1/2)))\^{}7)/1−(1/2))}

\emph{= (1−(1/128)/(1/2)) = ((127/128)/(1/2)) = (127/128)×(2/1) =
(127/64)}
\end{quote}

\textbf{2 - 4 + 8 - 16 + ............ ধারাটির প্রথম সাতটি পদের সমষ্টি কত?}

ক) 49

খ) 64

\textbf{\hl{গ) 86 ✔}}

ঘ) 109

\textbf{\hl{✔ সঠিক উত্তর: গ) 86}}

\begin{quote}
\textbf{ব্যাখ্যা:} \emph{সমাধান: ধারাটি গুণোত্তর। এখানে, ১ম পদ, a = 2}

\emph{সাধারণ অনুপাত, q = (−4/2) = − 2 \textless{} 1}

\emph{পদসংখ্যা n = 7}

\emph{∴ প্রথম সাতটি পদের সমষ্টি = (a(1−q\^{}n)/1−q)}

\emph{= (2(1−(−2)\^{}7)/1−(−2)) = (2(1+128)/1+2) = (2×129/3) = 86}

\emph{∴ সঠিক উত্তর 86}
\end{quote}

\textbf{1 - 1 + 1 - 1 + ........ এই ধারারিক (2n+1) পদের সমষ্টি হবে?}

ক) -1

\textbf{\hl{খ) 1 ✔}}

গ) 0

ঘ) কোনটি নয়

\textbf{\hl{✔ সঠিক উত্তর: খ) 1}}

\begin{quote}
\textbf{ব্যাখ্যা:} \emph{ধারাটি হলো: 1 - 1 + 1 - 1 + ........}

\emph{এখানে পদগুলো বিজোড় সংখ্যক (যেমন: 1টি পদ হলে যোগফল 1, 3টি পদ হলে
যোগফল 1 + (-1) + 1 = 1, 5টি পদ হলে যোগফল 1) অর্থাৎ (2n+1) তম পদ পর্যন্ত
সবসময় বিজোড় সংখ্যক পদ থাকবে এবং তাদের সমষ্টি সর্বদা 1 হবে।}
\end{quote}

\textbf{2+6+18+............ ধারাটির প্রথম 8 টি পদের সমষ্টি নির্ণয় করুন}

ক) 6520

খ) 6530

গ) 6540

\textbf{\hl{ঘ) 6560 ✔}}

\textbf{\hl{✔ সঠিক উত্তর: ঘ) 6560}}

\begin{quote}
\textbf{ব্যাখ্যা:} \emph{ধারাটি গুণোত্তর। এখানে, সাধারণ অনুপাত, r = 3; a =
প্রথম পদ = 2}

\emph{ধারটির প্রথম 8টি পদের সমষ্টি = a.~ (r\^{}n−1/r−1) = 2.~
(3\^{}8−1/3−1) = 6560}
\end{quote}

\textbf{2 + 4 + 8 + 16 + ........ধারাটির n সংখ্যক পদের সমষ্টি 254 হলে n এর
মান কত?}

ক) 5

খ) 6

\textbf{\hl{গ) 7 ✔}}

ঘ) 8

\textbf{\hl{✔ সঠিক উত্তর: গ) 7}}

\begin{quote}
\textbf{ব্যাখ্যা:} \emph{ধারাটি গুণোত্তর।}

\emph{এখানে, সাধারণ অনুপাত, r = 2; প্রথম পদ, a = 2}

\emph{n-তম পদের সমষ্টি = a.~ (r\^{}n−1/r−1)}

\emph{শর্তমতে, 254 = 2.~ (2\^{}n−1/2−1)}

\emph{⇒ 2\^{}n − 1 = (254/2)}

\emph{⇒ 2\^{}n = 127 + 1 = 128 = 2\^{}7}

\emph{∴ n = 7}
\end{quote}

\textbf{(1/√(2)) , 1, √(2) --ধারাটির কোন পদ 8√(2) হবে?}

\textbf{\hl{ক) ৯ তম পদ ✔}}

খ) ১০ তম পদ

গ) ১১ তম পদ

ঘ) ১২ তম পদ

\textbf{\hl{✔ সঠিক উত্তর: ক) ৯ তম পদ}}

\begin{quote}
\textbf{ব্যাখ্যা:} \emph{প্রথম পদ, a = (1/√(2))}

\emph{সাধারণ অনুপাত, r = (1/(1/√(2))) = √(2)}

\emph{ধরি, n-তম পদ = 8√(2)}

\emph{⇒ ar\^{}n − 1 = 8√(2)}

\emph{⇒ (1/√(2)) × ( √(2) )\^{}n − 1 = 8√(2)}

\emph{⇒ ( √(2) )\^{}n − 1 = 8√(2) × √(2) = 8 × 2 = 16}

\emph{⇒ ( 2\^{}((1/2)) )\^{}n − 1 = 2\^{}4}

\emph{⇒ 2\^{} (n−1/2) = 2\^{}4}

\emph{⇒ (n−1/2) = 4}

\emph{⇒ n − 1 = 8}

\emph{∴ n = 8 + 1 = 9}

\emph{∴ সঠিক উত্তর ৯ তম পদ}
\end{quote}

\textbf{একটি গুণোত্তর অনুক্রমে তৃতীয় পদটি 20 এবং ষষ্ঠ (6-তম) পদটি 160 হলে প্রথম
পদটি-}

\textbf{\hl{ক) 5 ✔}}

খ) 10

গ) 12

ঘ) 8

\textbf{\hl{✔ সঠিক উত্তর: ক) 5}}

\begin{quote}
\textbf{ব্যাখ্যা:} \emph{সমাধান: ১ম পদ = a; সাধারণ অনুপাত = r; পদসংখ্যা = n
হলে, n-তম পদ = ar\^{}n − 1}

\emph{তৃতীয় পদ ar\^{}3 − 1 = ar\^{}2 = 20 ............... (1)}

\emph{ষষ্ঠ পদ ar\^{}6 − 1 = ar\^{}5 = 160 ............... (2)}

\emph{(2) ÷ (1) করে, r\^{}3 = 8 ⇒ r = 2}

\emph{r এর মান (1) নং এ বসিয়ে পাই, ar\^{}2 = 20 ⇒ a × 2\^{}2 = 20 ⇒ 4a =
20}

\emph{∴ a = 5}
\end{quote}

\textbf{3 + a + b + 81 গুণোত্তর ধারা ভুক্ত হলে b = কত?}

ক) ৯

খ) ১২

গ) ১৮

\textbf{\hl{ঘ) ২৭ ✔}}

\textbf{\hl{✔ সঠিক উত্তর: ঘ) ২৭}}

\begin{quote}
\textbf{ব্যাখ্যা:} \emph{সমাধান: 3 + a + b + 81 গুণোত্তর ধারার ১ম পদ a = 3}

\emph{ধरी, সাধারণ অনুপাত = q}

\emph{ধারাটির চতুর্থ পদ = aq\^{}4 − 1 = aq\^{}3}

\emph{প্রশ্নমতে, 3q\^{}3 = 81}

\emph{⇒ q\^{}3 = 27}

\emph{∴ q = 3}

\emph{∴ তৃতীয় পদ b = aq\^{}3 − 1 = aq\^{}2 = 3 × 3\^{}2 = 3 × 9 = 27}
\end{quote}

\textbf{একটি গুণোত্তর অনুক্রমের দ্বিতীয় পদটি −48 এবং পঞ্চম পদটি (3/4) হলে,
সাধারণ অনুপাত কত?}

ক) (1/2)

খ) −(1/2)

গ) (1/4)

\textbf{\hl{ঘ) −(1/4) ✔}}

\textbf{\hl{✔ সঠিক উত্তর: ঘ) −(1/4)}}

\begin{quote}
\textbf{ব্যাখ্যা:} \emph{গুণোত্তর অনুক্রমের প্রথম পদ a এবং সাধারণ অনুপাত r হলে
অনুক্রমটির n তম পদ = ar\^{}n − 1}

\emph{দ্বিতীয় পদ, ar\^{}2 − 1 = ar = −48 এবং ar\^{}5 − 1 = ar\^{}4 =
(3/4)}

\emph{এখন, (ar\^{}4/ar) = ((3/4)/−48)}

\emph{⇒ r\^{}3 = (1/−64)}

\emph{⇒ r\^{}3 = ((−(1/4)))\^{}3}

\emph{∴ r = −(1/4)}
\end{quote}

\section{28. লগারিদম-বীজগণিত \textbar\textbar{} সংজ্ঞা, সূত্রাবলীসহ এ টু
জেড}\label{ux9b2ux997ux9b0ux9a6ux9ae-ux9acux99cux997ux9a3ux9a4-ux9b8ux99cux99e-ux9b8ux9a4ux9b0ux9acux9b2ux9b8ux9b9-ux98f-ux99f-ux99cux9a1}

লগারিদম হলো একটি গাণিতিক অপারেশন যা সংখ্যা বা মানের লগারিদম নির্ণয় করে।
লগারিদমের সাহায্যে অনেক গাণিতিক সমস্যা সহজে সমাধান করা সম্ভব হয়।

লগারিদমের বেশিরভাগ ব্যবহার হয় কম্পিউটার সায়েন্সে এবং অনলাইন ক্যালকুলেশনে।

এছাড়াও, এটি পরিসংখ্যান এবং প্রক্রিয়া নির্ধারণের জন্য ব্যবহৃত হয়। একটি সংখ্যাকে
নির্দিষ্ট লগারিদম হিসাব করে সেই লগারিদমের মানের উপর ভিত্তি করে কিছু করা সম্ভব
হয় ।

সূচকের এর বিপরীত হলো লগারিদম । হিসেবের সুবিধার্থে অনেক বড় বা অনেক ছোট
সংখ্যাকে \(a \times 10^{n}\) আকারে প্রকাশ করা হয়। কোন সংখ্যার
\(a \times 10^{n}\) রূপকে বলা হয় সংখ্যাটির বৈজ্ঞানিক বা আদর্শ রূপ।

লগ প্রধানত দুই প্রকার। যথা:-

1) স্বাভাবিক লগারিদম (ln)

2) সাধারণ লগারিদম (log)

\begin{itemize}
\item
  প্রাথমিক আলোচনা
\item
  \textbf{লগারিদম কাকে বলে?}
\item
  লগারিদমের সূত্র সমূহ
\item
  Log- এর সাধারণ প্রশ্নসমূহ
\item
  Log- এর সমাধান
\end{itemize}

\subsubsection{প্রাথমিক
আলোচনা}\label{ux9aaux9b0ux9a5ux9aeux995-ux986ux9b2ux99aux9a8-11}

সর্বপ্রথম স্কটল্যান্ডের গণিতবিদ জন নেপিয়ার\textquotesingle{} (1550-1617)
লগারিদম আবিস্কার করেন।

শুধু ধনাত্মক সংখ্যার লগারিদম আছে। শূন্য, ঋণাত্মক সংখ্যার লগারিদম নেই।

যদি \(a^{x}\) = M হয় তবে x কে M এর a ভিত্তিক লগারিদম বা সংক্ষেপে লগ বলা
হয়। লগারিদমের প্রতীক ব্যবহার করে লেখা হয়: \({Log}_{a}M\) = x

\({Log}_{a}M\) কে a ভিত্তিক লগ \textquotesingle M\textquotesingle{} পড়া
হয়।

ব্যবহারিক ক্ষেত্রে লগারিদমের ভিত্তি সাধারণত ১০ ধরা হয়। ১০ ভিত্তিক লগারিদমকে
সাধারণ লগারিদম বলে। এই ক্ষেত্রে ভিত্তি উহ্য রাখা হয় অর্থাৎ \({Log}_{10}M\)
বোঝাতে logM লেখা হয়।

লগারিদমের প্রবর্তন করেন- জন নেপিয়ার

\subsubsection{লগারিদম কাকে
বলে?}\label{ux9b2ux997ux9b0ux9a6ux9ae-ux995ux995-ux9acux9b2}

গণিতশাস্ত্রে লগারিদম হলো সূচকের বিপরীত প্রক্রিয়া। এর অর্থ কোনো সংখ্যার লগারিদম
হলো সেই সূচক যেটাকে একটি নির্ধারিত মানের ঘাত হিসাবে উন্নীত করলে প্রথমোক্ত
সংখ্যাটি পাওয়া যায়। সহজভাবে, একটি সংখ্যাকে বার বার গুণ করলে, লগারিদম
সংখ্যাটিকে যতবার গুণ করা হয়েছিল তা গণনা করে।

\subsubsection{লগারিদমের সূত্র
সমূহ}\label{ux9b2ux997ux9b0ux9a6ux9aeux9b0-ux9b8ux9a4ux9b0-ux9b8ux9aeux9b9}

1. \(\log_{a}(MN) = \log_{a}(M) + \log_{a}(N)\) {[} অর্থাৎ গুণ থাকলে যোগ
এবং ভাগ থাকলে বিয়োগ করতে হয় {]}

2. \(\log_{a}(\frac{M}{N}) = \log_{a}M - \log_{a}N\) {[} অর্থাৎ ভাগ থাকলে
বিয়োগ করতে হয় {]}

3. \(\log_{a}M^{n} = n\log_{a}M\) {[} ভিত্তিমূলের উপর ভিত্তি, আবার ভিত্তি এর
উপর পাওয়ার থাকলে পাওয়ারটি শুরুতে বসে যেমন:
\(\log_{x}10^{5} = 5\log_{x}10\) {]}

4. \(\log_{a}a = 1\) {[} অর্থাৎ \(\log_{10}10 = 1\) , ভিত্তিমূল এবং ভিত্তি
মিলে গেলে তার মান সবসময় 1 হয়। {]}

5. \(\log_{a}a^{2} = 2\) {[} অর্থাৎ \(\log_{x}x^{4} = 4\) (অর্থাৎ কখনো log
এর ভিত্তিমুল ও ভিত্তি যদি সমান হয় তাহলে ভিত্তিমূল এবং ভিত্তি উভয়ে উঠে যায় এবং
ভিত্তির উপর যে পাওয়ার থাকে তাই উত্তরে লিখতে হয়। ) {]}

6. \(\log_{a}1 = 0\) {[} যে কোন ভিত্তিমূলের উপর ভিত্তি 1 হলে তার উত্তর 0 হয়
{]}

7. \(\log_{a}b \times \log_{b}a = 1\) {[} উভয় লগের ভিত্তি এবং ভিত্তিমূল
বিপরীতধর্মী হলে তার মান 1 হয় । {]}

8. loga + logb + logc = log ( abc ) {[} log কমন নেয়ার সময় যোগ থাকলে গুণ
{]}

9. loga - logb = \(\log(\frac{a}{b})\) {[} log কমন নেয়ার সময়, বিয়োগ
থাকলে ভাগ হয়, এবং প্রথমটি উপরে বসে {]}

10. \(\log_{a}y = x\) হলে \(a^{x} = y\) {[} সূত্রটি খুবই গুরুত্বপূর্ণ {]}

{[} কোন পাওয়ার = কোন মান দেয়া থাকলে log তুলে দিয়ে ঐ পাওয়ার ও মানটি স্থান
বদল করে অর্থাৎ পাওয়ার জায়গায় মানটি এবং মান এর জায়গায় পাওয়ার যায় {]}

\subsubsection{Log- এর সাধারণ
প্রশ্নসমূহ}\label{log--ux98fux9b0-ux9b8ux9a7ux9b0ux9a3-ux9aaux9b0ux9b6ux9a8ux9b8ux9aeux9b9}

\textbf{প্রশ্ন নম্বর-} ১. Log3 + log4 = ? ..

\begin{itemize}
\item
  (ক) Log2
\item
  (খ) Log4
\item
  (গ) Log6
\item
  (ঘ) Log12
\end{itemize}

\textbf{\hl{উত্তরঃ (ঘ) Log12}}

\textbf{Explanation:}

Log3 + log4 = Log (3×4) = log12

{[} log কমন নেয়ার সময় যোগ থাকলে গুণ হয় {]}

\textbf{প্রশ্ন নম্বর-} ২. \(\log_{2}64 + \log_{2}8\) এর মান কত? {[}
জনপ্রশাসন মন্ত্রণালয়ের প্রশাসনিক কর্মকর্তা-২০১৬{]}

\begin{itemize}
\item
  (ক) 128
\item
  (খ) 7
\item
  (গ) 9
\item
  (ঘ) 2
\end{itemize}

\textbf{\hl{উত্তরঃ (গ) 9}}

\textbf{Explanation:}

\(\log_{2}64 + \log_{2}8\)

= \(\log_{2}2^{6} + \log_{2}2^{3}\)

= \(6\log_{2}2 + 3\log_{2}2\) {[} অর্থাৎ কখনো log এর ভিত্তিমুল ও ভিত্তি যদি
সমান হয় তাহলে ভিত্তিমূল এবং ভিত্তি উভয়ে উঠে যায় এবং ভিত্তির উপর যে পাওয়ার
থাকে তাই উত্তরে লিখতে হয়। {]}

= 6 + 3 = 9

\textbf{প্রশ্ন নম্বর-} ৩. \(\log_{2}8\) = কত ? {[} প্রতিরক্ষা মন্ত্রণালয়ের সহ:
পরি: - ২০১৮{]}

\begin{itemize}
\item
  (ক) 1
\item
  (খ) 2
\item
  (গ) 3
\item
  (ঘ) 4
\end{itemize}

\textbf{\hl{উত্তরঃ (গ) 3}}

\textbf{Explanation:}

\(\log_{2}8\)

= \(\log_{2}2^{3}\)

= \(3\log_{2}2\)

= 3.1 = 3

\(\log_{4}2\) \textbf{= কত?}

\textbf{ব্যাখ্যা}: এখানে নিচের ভিত্তিমূলকে পরিবর্তন করা যাবে না। তাই উপরের
ভিত্তি 2 কে ভিত্তিমূলের সাথে মেলানোর জন্য যা করা দরকার তাই করতে হবে। আমার
জানি √4 = 2 সুতরাং আমরা 2 কে 4 বানানোর জন্য লিখতে পারি 2 = √4 .

\(\log_{4}2\)

= \(\log_{4}\sqrt{}4\)

= \(\log_{4}4^{\frac{1}{2}}\)

= \(\frac{1}{2}\) {[} Log এর ভিত্তিমূল এবং ভিত্তি মিলে যাওয়ায়, পাওয়ার রেখে
সব বাদ দেয়া হয়েছে। {]}

\textbf{প্রশ্ন নম্বর-} ৪. \(\log_{x}(\frac{1}{8}) = \  - \ 2\) হলে x এর মান
কত? {[} ৩৮তম বিসিএস প্রিলি {]}

\begin{itemize}
\item
  (ক) 2
\item
  (খ) √2
\item
  (গ) 2√2
\item
  (ঘ) 4
\end{itemize}

\textbf{\hl{উত্তরঃ (গ) 2√2}}

\textbf{Explanation:}

\(\log_{x}(\frac{1}{8}) = \  - \ 2\)

বা, \(x^{- \ 2} = \frac{1}{8}\)

বা, \(\frac{1}{x^{2}} = \frac{1}{8}\)

বা, \(x^{2} = 8\)

বা, \(x = \sqrt{4\  \times \ 2}\)

:. x = 2√2

\textbf{প্রশ্ন নম্বর-} ৫. \(\log_{3}\frac{1}{27}\) = কত? {[}ATEO (কোটা)
পরীক্ষা -১৫{]}

\begin{itemize}
\item
  (ক) 9
\item
  (খ) - 6
\item
  (গ) - 3
\item
  (ঘ) 5
\end{itemize}

\textbf{\hl{উত্তরঃ (গ) - 3}}

\textbf{Explanation:}

\(\log_{3}\frac{1}{27}\)

= \(\log_{3}\frac{1}{3^{3}}\)

= \(\log_{3}3^{- \ 3}\)

= \(- \ 3\log_{3}3\)

= - 3

\subsubsection{Log- এর
সমাধান}\label{log--ux98fux9b0-ux9b8ux9aeux9a7ux9a8}

\textbf{প্রশ্ন নম্বর-} ৬. \(\log{\sqrt{2}}^{16}\) = কত? {[} ৮ম প্রভাষক নিব:
পরীক্ষা-১২ {]}

\begin{itemize}
\item
  (ক) 7
\item
  (খ) 8
\item
  (গ) 5
\item
  (ঘ) 3
\end{itemize}

\textbf{\hl{উত্তরঃ (খ) 8}}

\textbf{Explanation:}

\(\log{\sqrt{2}}^{16}\)

= \(\log{\sqrt{2}}^{{(\sqrt{}2)}^{8}}\)

= 8 \(\log{\sqrt{2}}^{\sqrt{}2}\)

= 8

\textbf{প্রশ্ন নম্বর-} ৭. \(\log_{3}\frac{1}{9}\) এর মান- {[}৩৫তম বিসিএস{]}

\begin{itemize}
\item
  (ক) 2
\item
  (খ) - 2
\item
  (গ) 3
\item
  (ঘ) - 3
\end{itemize}

\textbf{\hl{উত্তরঃ (খ) - 2}}

\textbf{Explanation:}

\(\log_{3}\frac{1}{9}\)

= \(\log_{3}\frac{1}{3^{2}}\)

= \(\log_{3}3^{- \ 2}\)

= - 2 \(\log_{3}3\)

= - 2

\textbf{প্রশ্ন নম্বর-} ৮. \(\log_{x}\frac{1}{9}\) = - 2 হলে, x এর মান কত?
{[} বেসরকারী শিক্ষক নিবন্ধন পরীক্ষা-১১ {]}

\begin{itemize}
\item
  (ক) 7
\item
  (খ) - 7
\item
  (গ) 5
\item
  (ঘ) 3
\end{itemize}

\textbf{\hl{উত্তরঃ (ঘ) 3}}

\textbf{Explanation:}

\(\log_{x}\frac{1}{9}\) = - 2

বা, \(x^{- \ 2} = \frac{1}{9}\)

বা, \(x^{- \ 2} = \frac{1}{3^{2}}\)

বা, \(x^{- \ 2} = 3^{- \ 2}\)

বা, x = 3

\textbf{প্রশ্ন নম্বর-} ৯. \(\log_{x}\frac{1}{16}\) = - 2 হলে, x এর মান কত?
{[} ১২তম বেসরকারি স্কুল নিবন্ধন পরীক্ষা- ১৫{]}

\begin{itemize}
\item
  (ক) 3
\item
  (খ) 5
\item
  (গ) 4
\item
  (ঘ) 6
\end{itemize}

\textbf{\hl{উত্তরঃ (গ) 4}}

\textbf{Explanation:}

\(\log_{x}\frac{1}{16}\) = - 2

বা, \(x^{- \ 2} = \frac{1}{16}\)

বা, \(x^{- \ 2} = \frac{1}{4^{2}}\)

বা, \(x^{- \ 2} = 4^{- \ 2}\)

বা, x = 4

\textbf{প্রশ্ন নম্বর-} ১০. \(\log_{10}(\ 0.001)\) = কত ? {[} রেলওয়ে
উপ-সহকারী -২০১৩{]}

\begin{itemize}
\item
  (ক) - 2
\item
  (খ) - 3
\item
  (গ) (1/2)
\item
  (ঘ) (1/3)
\end{itemize}

\textbf{\hl{উত্তরঃ (খ) - 3}}

\textbf{Explanation:}

\(\log_{10}(\ 0.001)\)

বা, \(10^{x} = .001\)

বা, \(10^{x} = 10^{- \ 3}\)

বা, x = - 3

\textbf{প্রশ্ন নম্বর-} ১১. \(\log a^{x} = 1\) , \(\log a^{y} = 2\) এবং
\(\log a^{z} = 3\) হলে , \(\log a^{(\frac{x^{3}y^{2}}{z})}\) এর মান কত?
{[} ৩৫-তম বিসিএস{]}

\begin{itemize}
\item
  (ক) 1
\item
  (খ) 2
\item
  (গ) 4
\item
  (ঘ) 5
\end{itemize}

\textbf{\hl{উত্তরঃ (গ) 4}}

\textbf{Explanation:}

\(\log a^{x} = 1\) বা, \(a^{1} = x\) বা, x = a

\(\log a^{y} = 2\) বা, \(a^{2} = y\) বা, y = \(a^{2}\)

\(\log a^{z} = 3\) বা, \(a^{3} = z\) বা, z = \(a^{3}\)

এখন , \(\log a^{(\frac{x^{3}y^{2}}{z})}\)

= \(\log_{a}a^{4}\)

= 4\(\log_{a}a\)

= 4 × 1 = 4

\textbf{প্রশ্ন নম্বর-} ১২. \(\log_{x}(\frac{3}{2}) = - \frac{1}{2}\) হলে, x
- এর মান কত? {[} ৩৭-তম বিসিএস প্রিলি {]}

\begin{itemize}
\item
  (ক) (4/9)
\item
  (খ) (9/4)
\item
  (গ) √((3/2))
\item
  (ঘ) √((2/3))
\end{itemize}

\textbf{\hl{উত্তরঃ (ক)}} \(\frac{4}{9}\)

\textbf{Explanation:}

\(\log_{x}(\frac{3}{2}) = - \frac{1}{2}\)

বা, \(x - \frac{1}{2} = \frac{3}{2}\)

বা, \(\frac{1}{x^{\frac{1}{2}}} = \frac{3}{2}\) {[} পাওয়ার এ মাইনাস
থাকলে তা ভগ্নাংশ আকারে লিখতে হয় {]}

বা, \(\frac{1}{\sqrt{}x} = \frac{3}{2}\)

বা, √x = \(\frac{2}{3}\)

বা, x = \({(\frac{2}{3})}^{2}\)

সুতরাং \(x = \frac{4}{9}\)

\textbf{প্রশ্ন নম্বর-} ১৩. \(\log_{2\sqrt{}5}400 = x\) এর মান কত? {[} ৮ম
প্রভাষক নিবন্ধন পরীক্ষা-১২ {]}

\begin{itemize}
\item
  (ক) 7
\item
  (খ) 4
\item
  (গ) 5
\item
  (ঘ) 3
\end{itemize}

\textbf{\hl{উত্তরঃ (খ) 4}}

\textbf{Explanation:}

\(\log_{2\sqrt{}5}400 = x\)

বা, \({(2\sqrt{}5)}^{x} = 400\) {[} পাওয়ার এবং মান কে স্থানান্তরিত করে {]}

বা, \({(2\sqrt{}5)}^{x} = 16\  \times \ 25\) {[} 400 কে ভাঙ্গানোর সময় এমন
সংখ্যা দিয়ে ভাঙ্গতে হবে যেন প্রতিটি বর্গ সংখ্যা হয় {]}

বা, \({(2\sqrt{}5)}^{x} = 2^{4} \times 5^{2}\)

বা, \({(2\sqrt{}5)}^{x} = 2^{4} \times {(\sqrt{}5)}^{4}\)

বা, \({(2\sqrt{}5)}^{x} = {(2\sqrt{}5)}^{4}\)

সুতরাং x = 4

\textbf{প্রশ্ন নম্বর-} ১৪.
\(\log_{2}\sqrt{}6 + \log_{2}\sqrt{\frac{2}{3}}\) = কত? {[} পরিবেশ ও বন
মন্ত্রনালয়ের সহ পরিচালক: ১০ {]}

\begin{itemize}
\item
  (ক) log5
\item
  (খ) 1
\item
  (গ) log\_22
\item
  (ঘ) 3 log 2
\end{itemize}

\textbf{\hl{উত্তরঃ (খ) 1}}

\textbf{Explanation:}

\(\log_{2}\sqrt{}6 + \log_{2}\sqrt{\frac{2}{3}}\)

= \(\log_{2}(\sqrt{6 \times \frac{2}{3}})\)

= \(\log_{2}(\sqrt{}2\  \times \ \sqrt{}2)\)

= \(\log_{2}2\)

= 1

\section{29. লগারিদম MCQ প্রশ্ন উত্তর
অনুশীলন}\label{ux9b2ux997ux9b0ux9a6ux9ae-mcq-ux9aaux9b0ux9b6ux9a8-ux989ux9a4ux9a4ux9b0-ux985ux9a8ux9b6ux9b2ux9a8}

\textbf{লগারিদমের প্রবর্তন করেন---}

ক) নিউটন

খ) প্রস্পার একার্ট

গ) জন মউসলি

\textbf{\hl{ঘ) জন নেপিয়ার ✔}}

\textbf{\hl{✔ সঠিক উত্তর: ঘ) জন নেপিয়ার}}

\begin{quote}
\textbf{ব্যাখ্যা:} \emph{লগারিদমের ধারণা বা প্রবর্তন করেন স্কটিশ গণিতবিদ জন
নেপিয়ার (John Napier) ১৬১৪ সালে।}
\end{quote}

\textbf{কোন শর্তে log\_a1=0 ?}

\textbf{\hl{ক) a\textgreater0,a≠1 ✔}}

খ) a≠0,a\textgreater1

গ) a\textgreater0,a=1

ঘ) a≠0,a\textless1

\textbf{\hl{✔ সঠিক উত্তর: ক) a\textgreater0,a≠1}}

\begin{quote}
\textbf{ব্যাখ্যা:} \emph{দেওয়া আছে, log\_a1=0}

\emph{⇒ a\^{}0=1}

\emph{লগারিদমের ভিত্তি বা সংজ্ঞার শর্তানুসারে, লগারিদমের ভিত্তি a এর মান
অবশ্যই ধনাত্মক হতে হবে এবং ১ এর সমান হতে পারবে না।}

\emph{∴ a\textgreater0 এবং a≠1}
\end{quote}

\textbf{কোন শর্তে log\_aa=1}

ক) a\textgreater0

খ) a≠1

গ) a≠0,a\textgreater1

\textbf{\hl{ঘ) a\textgreater0,a≠1 ✔}}

\textbf{\hl{✔ সঠিক উত্তর: ঘ) a\textgreater0,a≠1}}

\begin{quote}
\textbf{ব্যাখ্যা:} \emph{লগারিদমের সংজ্ঞার শর্ত অনুযায়ী, ভিত্তি a এর মান শূন্য
থেকে বড় হতে হবে ( a\textgreater0 ) এবং ভিত্তি কখনোই ১ হতে পারবে না ( a≠1
)।}
\end{quote}

\textbf{32 এর 2 ভিত্তিক লগারিদম কত?}

ক) 3

খ) 4

\textbf{\hl{গ) 5 ✔}}

ঘ) 6

\textbf{\hl{✔ সঠিক উত্তর: গ) 5}}

\begin{quote}
\textbf{ব্যাখ্যা:} \emph{log\_232=log\_22\^{}5=5log\_22=5×1=5}
\end{quote}

\textbf{400-এর log হলো 4; ভিত্তি কত?}

ক) 10

খ) e

গ) 2

\textbf{\hl{ঘ) 2√(5) ✔}}

\textbf{\hl{✔ সঠিক উত্তর: ঘ) 2√(5)}}

\begin{quote}
\textbf{ব্যাখ্যা:} \emph{ধري, ভিত্তি a}

\emph{∴ log\_a400=4}

\emph{⇒ a\^{}4=400}

\emph{⇒ a\^{}4=(2√(5))\^{}4}

\emph{∴ a=2√(5)}
\end{quote}

\textbf{3√(3) এর 3 ভিত্তিক লগ কত?}

ক) √(3)

খ) (2/3)

গ) 3\^{}2

\textbf{\hl{ঘ) (3/2) ✔}}

\textbf{\hl{✔ সঠিক উত্তর: ঘ) (3/2)}}

\begin{quote}
\textbf{ব্যাখ্যা:} \emph{log\_3(3√(3))=log\_3(3\^{}1.3\^{}((1/2)))}

\emph{= log\_3(3\^{}(1+(1/2)))=log\_3(3\^{}((3/2)))}

\emph{= (3/2)log\_33=(3/2)×1=(3/2)}
\end{quote}

\textbf{5√(5) এর 5 ভিত্তিক লগ কত?}

ক) (1/5)

\textbf{\hl{খ) (3/2) ✔}}

গ) (2/3)

ঘ) (1/3)

\textbf{\hl{✔ সঠিক উত্তর: খ) (3/2)}}

\begin{quote}
\textbf{ব্যাখ্যা:} \emph{log\_5(5√(5))=log\_5(5\^{}1.5\^{}((1/2)))}

\emph{= log\_5(5\^{}((3/2)))=(3/2)log\_55=(3/2)}
\end{quote}

\textbf{a\^{}x=n হলে নিচের কোন সম্পর্কটি সঠিক?}

ক) x=a\^{}n

খ) x=lnx

গ) a=x\^{}n

\textbf{\hl{ঘ) log\_an ✔}}

\textbf{\hl{✔ সঠিক উত্তর: ঘ) log\_an}}

\begin{quote}
\textbf{ব্যাখ্যা:} \emph{লগারিদমের সংজ্ঞা অনুযায়ী, যদি a\^{}x=n হয়, তবে
x=log\_an লেখা যায়। অপশনগুলোর মধ্যে \_an সঠিক সম্পর্ক নির্দেশ করে।}
\end{quote}

\textbf{a\^{}5=x\^{}4 হলে a এর মান কত?}

ক) 2

\textbf{\hl{খ) 1 ✔}}

গ) 3

ঘ) 8

\textbf{\hl{✔ সঠিক উত্তর: খ) 1}}

\begin{quote}
\textbf{ব্যাখ্যা:} \emph{a\^{}5=x\^{}4 সমীকরণে x এর স্থানে a\^{}4 হলে,
a\^{}5=(a\^{}4)\^{}4=a\^{}16 (বা অন্য শর্তে a=1 হলে 1\^{}5=x\^{}4⇒x\^{}4=1
ইত্যাদি বিবেচনায় a=1 সিদ্ধ হয়)।}
\end{quote}

\textbf{a\^{}x=1 হলে x এর মান কত?}

ক) 1

\textbf{\hl{খ) 0 ✔}}

গ) অনির্ণেয়

ঘ) 2

\textbf{\hl{✔ সঠিক উত্তর: খ) 0}}

\begin{quote}
\textbf{ব্যাখ্যা:} \emph{যেকোনো অশূন্য সংখ্যার ওপর পাওয়ার শূন্য হলে তার মান 1
হয়।}

\emph{ধরি, a\^{}x=1}

\emph{⇒ a\^{}x=a\^{}0 {[}যেহেতু a\^{}0=1 {]}}

\emph{∴ x=0}
\end{quote}

\textbf{(((x/2)))\^{}(a+1)=1 হলে, a এর মান কত?}

ক) 0

খ) 2

গ) 1

\textbf{\hl{ঘ) −1 ✔}}

\textbf{\hl{✔ সঠিক উত্তর: ঘ) −1}}

\begin{quote}
\textbf{ব্যাখ্যা:} \emph{দেওয়া আছে, (((x/2)))\^{}(a+1)=1}

\emph{⇒ (((x/2)))\^{}(a+1)=(((x/2)))\^{}0}

\emph{⇒ a+1=0}

\emph{∴ a=−1}
\end{quote}

\textbf{a\^{}x=b,b\^{}y=c,c\^{}z=a হলে, নিচের কোন সম্পর্কটি সঠিক?}

ক) a=a\^{}yz

খ) a=a\^{}zy

গ) b=c\^{}yz

\textbf{\hl{ঘ) a=a\^{}xyz ✔}}

\textbf{\hl{✔ সঠিক উত্তর: ঘ) a=a\^{}xyz}}

\begin{quote}
\textbf{ব্যাখ্যা:} \emph{দেওয়া আছে, a\^{}x=b , b\^{}y=c , c\^{}z=a}

\emph{⇒ a=c\^{}z=(b\^{}y)\^{}z=b\^{}yz=(a\^{}x)\^{}yz=a\^{}xyz}

\emph{∴ a=a\^{}xyz}
\end{quote}

\textbf{x\^{}a=y,y\^{}b=z,z\^{}c=x হলে, abc এর মান কত?}

ক) 0

\textbf{\hl{খ) 1 ✔}}

গ) 2

ঘ) 3

\textbf{\hl{✔ সঠিক উত্তর: খ) 1}}

\begin{quote}
\textbf{ব্যাখ্যা:} \emph{দেওয়া আছে, z\^{}c=x}

\emph{⇒ (y\^{}b)\^{}c=x}

\emph{⇒ y\^{}bc=x}

\emph{⇒ (x\^{}a)\^{}bc=x}

\emph{⇒ x\^{}abc=x\^{}1}

\emph{∴ abc=1}
\end{quote}

\textbf{log\_√(2)16 = কত?}

\textbf{\hl{ক) 8 ✔}}

খ) 6

গ) 9

ঘ) 4

\textbf{\hl{✔ সঠিক উত্তর: ক) 8}}

\begin{quote}
\textbf{ব্যাখ্যা:}
\emph{log\_√(2)16=log\_√(2)(√(2))\^{}8=8log\_√(2)√(2)=8×1=8}
\end{quote}

\textbf{log\_√(3)81 কত?}

ক) 4

খ) 27√(3)

\textbf{\hl{গ) 8 ✔}}

ঘ) (1/8)

\textbf{\hl{✔ সঠিক উত্তর: গ) 8}}

\begin{quote}
\textbf{ব্যাখ্যা:} \emph{log\_√(3)81=log\_√(3)(√(3))\^{}8}

\emph{= 8log\_√(3)√(3)}

\emph{= 8×1=8 {[}যেহেতু log\_aa=1 {]}}
\end{quote}

\textbf{log\_28 কত?}

ক) 4

\textbf{\hl{খ) 3 ✔}}

গ) 2

ঘ) 1

\textbf{\hl{✔ সঠিক উত্তর: খ) 3}}

\begin{quote}
\textbf{ব্যাখ্যা:} \emph{log\_28=log\_22\^{}3=3log\_22=3×1=3}
\end{quote}

\textbf{log\_216 এর মান কোনটি?}

ক) 5

খ) 3

\textbf{\hl{গ) 4 ✔}}

ঘ) (1/4)

\textbf{\hl{✔ সঠিক উত্তর: গ) 4}}

\begin{quote}
\textbf{ব্যাখ্যা:} \emph{log\_216=log\_22\^{}4=4log\_22=4×1=4}
\end{quote}

\textbf{log\_327 এর মান কত?}

\textbf{\hl{ক) 3 ✔}}

খ) 27

গ) 9

ঘ) 10

\textbf{\hl{✔ সঠিক উত্তর: ক) 3}}

\begin{quote}
\textbf{ব্যাখ্যা:} \emph{log\_327}

\emph{=log\_33\^{}3}

\emph{=3log\_33 =3×1 =3}
\end{quote}

\textbf{log\_381 এর মান কত?}

ক) 3

\textbf{\hl{খ) 8 ✔}}

গ) 9

ঘ) 32

\textbf{\hl{✔ সঠিক উত্তর: খ) 8}}

\begin{quote}
\textbf{ব্যাখ্যা:} \emph{log\_381}

\emph{=log\_33\^{}4}

\emph{=4log\_33 =4×1 =4}
\end{quote}

\textbf{log\_42 এর মান কত?}

ক) (1/3)

খ) 2

\textbf{\hl{গ) (1/2) ✔}}

ঘ) 4

\textbf{\hl{✔ সঠিক উত্তর: গ) (1/2)}}

\begin{quote}
\textbf{ব্যাখ্যা:} \emph{ধরি, log\_42=x}

\emph{⇒ 4\^{}x=2}

\emph{⇒ (2\^{}2)\^{}x=2\^{}1}

\emph{⇒ 2\^{}2x=2\^{}1}

\emph{⇒ 2x=1}

\emph{∴ x=(1/2)}
\end{quote}

\textbf{log\_7343= কত?}

ক) 1

খ) 2

\textbf{\hl{গ) 3 ✔}}

ঘ) 4

\textbf{\hl{✔ সঠিক উত্তর: গ) 3}}

\begin{quote}
\textbf{ব্যাখ্যা:} \emph{log\_7343}

\emph{=log\_77\^{}3}

\emph{=3log\_77 =3×1 =3}
\end{quote}

\textbf{Find the value of log\_5625}

ক) 1

খ) 2

গ) 3

\textbf{\hl{ঘ) 4 ✔}}

\textbf{\hl{✔ সঠিক উত্তর: ঘ) 4}}

\begin{quote}
\textbf{ব্যাখ্যা:} \emph{log\_5625}

\emph{=log\_55\^{}4}

\emph{=4×log\_55 =4×1 =4}
\end{quote}

\textbf{log\_2.56.25 = কত?}

ক) 1

\textbf{\hl{খ) 2 ✔}}

গ) 3

ঘ) 4

\textbf{\hl{✔ সঠিক উত্তর: খ) 2}}

\begin{quote}
\textbf{ব্যাখ্যা:} \emph{log\_2.56.25}

\emph{=log\_2.5(2.5)\^{}2}

\emph{=2log\_2.52.5 =2×1 =2}
\end{quote}

\textbf{log\_2(1/4) = কত?}

ক) 2

\textbf{\hl{খ) -2 ✔}}

গ) (1/2)

ঘ) −(1/2)

\textbf{\hl{✔ সঠিক উত্তর: খ) -2}}

\begin{quote}
\textbf{ব্যাখ্যা:} \emph{log\_2(1/4)}

\emph{=log\_2(1/2\^{}2)}

\emph{=log\_22\^{}−2}

\emph{=−2log\_22}

\emph{=−2×1}

\emph{=−2}
\end{quote}

\textbf{log\_2(1/8) এর মান কত?}

\textbf{\hl{ক) -3 ✔}}

খ) (1/2)

গ) 2

ঘ) 3

\textbf{\hl{✔ সঠিক উত্তর: ক) -3}}

\begin{quote}
\textbf{ব্যাখ্যা:} \emph{log\_2(1/8)}

\emph{=log\_2(1/2\^{}3)}

\emph{=log\_22\^{}−3}

\emph{=−3log\_22}

\emph{=−3×1}

\emph{=−3}
\end{quote}

\textbf{log\_2((1/32)) এর মান কত?}

\textbf{\hl{ক) −5 ✔}}

খ) (1/25)

গ) (1/5)

ঘ) −(1/5)

\textbf{\hl{✔ সঠিক উত্তর: ক) −5}}

\begin{quote}
\textbf{ব্যাখ্যা:} \emph{log\_2((1/32))=log\_2((1/2\^{}5))}

\emph{⇒ log\_2(2\^{}−5)=−5log\_22=−5×1=−5}
\end{quote}

\textbf{log\_3((1/9)) এর মান---}

ক) 2

\textbf{\hl{খ) −2 ✔}}

গ) 3

ঘ) −3

\textbf{\hl{✔ সঠিক উত্তর: খ) −2}}

\begin{quote}
\textbf{ব্যাখ্যা:} \emph{log\_3(1/9)=log\_3(1/3\^{}2)}

\emph{⇒ log\_33\^{}−2=−2log\_33=−2}
\end{quote}

\textbf{log\_3(1/27) = কত?}

ক) 3

খ) 9

\textbf{\hl{গ) −3 ✔}}

ঘ) −9

\textbf{\hl{✔ সঠিক উত্তর: গ) −3}}

\begin{quote}
\textbf{ব্যাখ্যা:} \emph{log\_3(1/27)=log\_33\^{}−3}

\emph{⇒ −3×1=−3}
\end{quote}

\textbf{log\_53√(5) এর মান কত?}

\textbf{\hl{ক) (1/3) ✔}}

খ) (1/5)

গ) (3/5)

ঘ) (5/3)

\textbf{\hl{✔ সঠিক উত্তর: ক) (1/3)}}

\begin{quote}
\textbf{ব্যাখ্যা:} \emph{log\_53√(5)=log\_55\^{}((1/3))}

\emph{⇒ (1/3)log\_55\^{}((1/3))=(1/3)}
\end{quote}

\textbf{log\_55√(5) = কত?}

ক) (1/5)

\textbf{\hl{খ) (3/2) ✔}}

গ) (2/3)

ঘ) (5/3)

\textbf{\hl{✔ সঠিক উত্তর: খ) (3/2)}}

\begin{quote}
\textbf{ব্যাখ্যা:} \emph{log\_55√(5)=log\_55×5\^{}((1/2))}

\emph{⇒ log\_55\^{}(1+(1/2))=log\_55\^{}((3/2))}

\emph{⇒ (3/2)log\_55=(3/2)}
\end{quote}

\textbf{log\_5(√(5).3√(5)) এর মান কত?}

\textbf{\hl{ক) (5/6) ✔}}

খ) (1/3)

গ) (1/6)

ঘ) 1

\textbf{\hl{✔ সঠিক উত্তর: ক) (5/6)}}

\begin{quote}
\textbf{ব্যাখ্যা:}
\emph{log\_5(5\^{}((1/2)).5\^{}((1/3)))=log\_5(5\^{}((1/2)+(1/3)))}

\emph{⇒ log\_5(5\^{}((5/6)))=(5/6)log\_55=(5/6)}
\end{quote}

\textbf{log\_10(0.001) = কত?}

ক) −2

\textbf{\hl{খ) −3 ✔}}

গ) (1/2)

ঘ) (1/3)

\textbf{\hl{✔ সঠিক উত্তর: খ) −3}}

\begin{quote}
\textbf{ব্যাখ্যা:} \emph{ধরি, log\_10(0.001)=x}

\emph{∴ 10\^{}x=0.001=(1/1000)}

\emph{⇒ 10\^{}x=10\^{}−3}

\emph{∴ x=−3}
\end{quote}

\textbf{log\_10x=3 হলে, x এর মান কত?}

ক) 1

খ) 10

গ) 100

\textbf{\hl{ঘ) 1000 ✔}}

\textbf{\hl{✔ সঠিক উত্তর: ঘ) 1000}}

\begin{quote}
\textbf{ব্যাখ্যা:} \emph{দেওয়া আছে, log\_10x=3}

\emph{প্রশ্নমতে, 10\^{}3=x}

\emph{⇒ x=1000}
\end{quote}

\textbf{log\_10x=−3 হলে x= ?}

ক) 0.1

খ) 0.01

\textbf{\hl{গ) 0.001 ✔}}

ঘ) 0.0001

\textbf{\hl{✔ সঠিক উত্তর: গ) 0.001}}

\begin{quote}
\textbf{ব্যাখ্যা:} \emph{দেওয়া আছে, log\_10x=−3}

\emph{⇒ x=10\^{}−3}

\emph{⇒ x=(1/10\^{}3)=(1/1000)=0.001}

\emph{∴ x=0.001}
\end{quote}

\textbf{If log\_x144=4 , then find the value of x .}

ক) 16

খ) 312

\textbf{\hl{গ) 2√(3) ✔}}

ঘ) √(6)

\textbf{\hl{✔ সঠিক উত্তর: গ) 2√(3)}}

\begin{quote}
\textbf{ব্যাখ্যা:} \emph{দেওয়া আছে, log\_x144=4}

\emph{⇒ x\^{}4=144}

\emph{⇒ (x\^{}2)\^{}2=9×16}

\emph{⇒ x\^{}2=√(3\^{}2×4\^{}2)}

\emph{⇒ x\^{}2=3×4}

\emph{⇒ x=√(3×4)}

\emph{∴ x=2√(3)}
\end{quote}

\textbf{log\_x324=4 হলে, x এর মান কত?}

ক) 2√(3)

\textbf{\hl{খ) 3√(2) ✔}}

গ) √(3)

ঘ) 3

\textbf{\hl{✔ সঠিক উত্তর: খ) 3√(2)}}

\begin{quote}
\textbf{ব্যাখ্যা:} \emph{দেওয়া আছে, log\_x324=4}

\emph{⇒ x\^{}4=324}

\emph{⇒ x\^{}4=(3√(2))\^{}4}

\emph{∴ x=3√(2)}
\end{quote}

\textbf{log\_x5=2 হলে, x এর মান কত?}

\textbf{\hl{ক) √(5) ✔}}

খ) 25

গ) −√(5)

ঘ) ±√(5)

\textbf{\hl{✔ সঠিক উত্তর: ক) √(5)}}

\begin{quote}
\textbf{ব্যাখ্যা:} \emph{দেওয়া আছে, log\_x5=2}

\emph{⇒ x\^{}2=5}

\emph{∴ x=√(5)}
\end{quote}

\textbf{log\_x((1/8))=−2 হলে, x কত?}

ক) 2

খ) √(2)

\textbf{\hl{গ) 2√(2) ✔}}

ঘ) 4

\textbf{\hl{✔ সঠিক উত্তর: গ) 2√(2)}}

\begin{quote}
\textbf{ব্যাখ্যা:} \emph{দেওয়া আছে, log\_x((1/8))=−2}

\emph{⇒ x\^{}−2=(1/8)}

\emph{⇒ (1/x\^{}2)=(1/8)}

\emph{⇒ x\^{}2=8}

\emph{⇒ x=√(8)=2√(2)}
\end{quote}

\textbf{log\_x(1/16)=−2 হলে, x এর মান কত?}

ক) 3

খ) 5

\textbf{\hl{গ) 8 ✔}}

ঘ) 6

\textbf{\hl{✔ সঠিক উত্তর: গ) 8}}

\begin{quote}
\textbf{ব্যাখ্যা:} \emph{দেওয়া আছে, log\_x(1/16)=−2}

\emph{⇒ x\^{}−2=(1/16)}

\emph{⇒ x\^{}−2=4\^{}−2}

\emph{∴ x=4}
\end{quote}

\textbf{log\_x(1/81)=−4 হলে x এর মান কত?}

ক) −4

\textbf{\hl{খ) 3 ✔}}

গ) (1/4)

ঘ) −(1/4)

\textbf{\hl{✔ সঠিক উত্তর: খ) 3}}

\begin{quote}
\textbf{ব্যাখ্যা:} \emph{দেওয়া আছে, log\_x(1/81)=−4}

\emph{⇒ x\^{}−4=(1/81)}

\emph{⇒ x\^{}−4=3\^{}−4}

\emph{∴ x=3}
\end{quote}

\textbf{log\_x(3/2)=−(1/2) হলে, x এর মান---}

\textbf{\hl{ক) (4/9) ✔}}

খ) (9/4)

গ) √((3/2))

ঘ) √((2/3))

\textbf{\hl{✔ সঠিক উত্তর: ক) (4/9)}}

\begin{quote}
\textbf{ব্যাখ্যা:} \emph{দেওয়া আছে, log\_x(3/2)=−(1/2)}

\emph{⇒ x\^{}(−(1/2))=(3/2)}

\emph{⇒ (1/x\^{}((1/2)))=(3/2)}

\emph{⇒ (1/√(x))=(3/2)}

\emph{⇒ √(x)=(2/3)}

\emph{∴ x=(((2/3)))\^{}2=(4/9)}
\end{quote}

\textbf{log\_x((5/2))=−(1/2) হলে x এর মান---}

\textbf{\hl{ক) (4/25) ✔}}

খ) √((5/2))

গ) √((2/5))

ঘ) (25/4)

\textbf{\hl{✔ সঠিক উত্তর: ক) (4/25)}}

\begin{quote}
\textbf{ব্যাখ্যা:} \emph{দেওয়া আছে, log\_x((5/2))=−(1/2)}

\emph{⇒ x\^{}(−(1/2))=(5/2)}

\emph{⇒ (1/x\^{}((1/2)))=(5/2)}

\emph{⇒ (((1/√(x))))\^{}2=(((5/2)))\^{}2}

\emph{⇒ (1/x)=(25/4)}

\emph{∴ x=(4/25)}
\end{quote}

\textbf{log\_x(1/9)=−2 হলে x এর মান কত?}

ক) −(1/3)

খ) (1/3)

গ) −3

\textbf{\hl{ঘ) +3 ✔}}

\textbf{\hl{✔ সঠিক উত্তর: ঘ) +3}}

\begin{quote}
\textbf{ব্যাখ্যা:} \emph{দেওয়া আছে, log\_x(1/9)=−2}

\emph{⇒ x\^{}−2=(1/9)}

\emph{⇒ (1/x\^{}2)=(1/3\^{}2)}

\emph{⇒ x\^{}2=3\^{}2}

\emph{∴ x=3}
\end{quote}

\textbf{log\_x((1/27))=−3 হলে, x এর মান---}

ক) -3

\textbf{\hl{খ) 3 ✔}}

গ) −(1/3)

ঘ) (1/3)

\textbf{\hl{✔ সঠিক উত্তর: খ) 3}}

\begin{quote}
\textbf{ব্যাখ্যা:} \emph{দেওয়া আছে, log\_x((1/27))=−3}

\emph{⇒ x\^{}−3=(1/27)}

\emph{⇒ x\^{}−3=(1/3\^{}3)}

\emph{⇒ x\^{}−3=3\^{}−3}

\emph{∴ x=3}
\end{quote}

\textbf{log\_y3√(3)=(1/15) হলে y এর মান কত?}

ক) 9

খ) 27

গ) 81

\textbf{\hl{ঘ) 243 ✔}}

\textbf{\hl{✔ সঠিক উত্তর: ঘ) 243}}

\begin{quote}
\textbf{ব্যাখ্যা:} \emph{দেওয়া আছে, log\_y3√(3)=(1/15)}

\emph{⇒ y\^{}((1/15))=3√(3)=3\^{}((1/3))}

\emph{⇒ y\^{}((1/15))=(3\^{}5)\^{}((1/15))}

\emph{∴ y=3\^{}5=243}
\end{quote}

\textbf{log\_3√(2)(1/324)= কত?}

ক) 162√(2)

খ) (1/162√(2))

গ) 4

\textbf{\hl{ঘ) −4 ✔}}

\textbf{\hl{✔ সঠিক উত্তর: ঘ) −4}}

\begin{quote}
\textbf{ব্যাখ্যা:} \emph{log\_3√(2)(1/324)=log\_3√(2)(3√(2))\^{}−4}

\emph{⇒ −4log\_3√(2)3√(2)=−4×1=−4 {[}যেহেতু log\_aa=1 {]}}
\end{quote}

\textbf{log\_a√(2)=(1/6) হলে, a= কত?}

ক) √(2)

খ) 0

গ) 6

\textbf{\hl{ঘ) 8 ✔}}

\textbf{\hl{✔ সঠিক উত্তর: ঘ) 8}}

\begin{quote}
\textbf{ব্যাখ্যা:} \emph{দেওয়া আছে, log\_a√(2)=(1/6)}

\emph{⇒ a\^{}((1/6))=√(2)=2\^{}((1/2))}

\emph{⇒ ((a\^{}((1/6))))\^{}6=((2\^{}((1/2))))\^{}6}

\emph{⇒ a=2\^{}3=8}
\end{quote}

\textbf{log\_5x=3 হলে x = কত?}

ক) 375

খ) 120

\textbf{\hl{গ) 125 ✔}}

ঘ) 225

\textbf{\hl{✔ সঠিক উত্তর: গ) 125}}

\begin{quote}
\textbf{ব্যাখ্যা:} \emph{দেওয়া আছে, log\_5x=3}

\emph{⇒ x=5\^{}3}

\emph{∴ x=125}
\end{quote}

\textbf{log\_12x=4 হলে x -এর মান---}

ক) 12\^{}2

খ) 12\^{}3

\textbf{\hl{গ) 12\^{}4 ✔}}

ঘ) 12\^{}5

\textbf{\hl{✔ সঠিক উত্তর: গ) 12\^{}4}}

\begin{quote}
\textbf{ব্যাখ্যা:} \emph{দেওয়া আছে, log\_12x=4}

\emph{∴ x=12\^{}4}
\end{quote}

\textbf{Log\_10x=2 হলে, x এর মান কত?}

ক) 1

খ) 10

\textbf{\hl{গ) 100 ✔}}

ঘ) 0

\textbf{\hl{✔ সঠিক উত্তর: গ) 100}}

\begin{quote}
\textbf{ব্যাখ্যা:} \emph{দেওয়া আছে, Log\_10x=2}

\emph{⇒ 10\^{}2=x}

\emph{∴ x=100}
\end{quote}

\textbf{log\_10x=−2 হলে x এর মান কত হবে?}

\textbf{\hl{ক) 0.01 ✔}}

খ) 0.001

গ) 0.02

ঘ) 0.002

\textbf{\hl{✔ সঠিক উত্তর: ক) 0.01}}

\begin{quote}
\textbf{ব্যাখ্যা:} \emph{দেওয়া আছে, log\_10x=−2}

\emph{⇒ x=10\^{}−2}

\emph{⇒ x=(1/10\^{}2)=(1/100)=0.01}

\emph{∴ x=0.01}
\end{quote}

\textbf{log\_2√(5)400=x হলে x এর মান কত?}

ক) 400

খ) 10

\textbf{\hl{গ) 4 ✔}}

ঘ) 2√(5)

\textbf{\hl{✔ সঠিক উত্তর: গ) 4}}

\begin{quote}
\textbf{ব্যাখ্যা:} \emph{দেওয়া আছে, log\_2√(5)400=x}

\emph{⇒ (2√(5))\^{}x=400}

\emph{⇒ (2√(5))\^{}x=(20)\^{}2}

\emph{⇒ (2√(5))\^{}x=(4×5)\^{}2}

\emph{⇒ (2√(5))\^{}x=((2√(5))\^{}2)\^{}2}

\emph{⇒ (2√(5))\^{}x=(2√(5))\^{}4}

\emph{∴ x=4}
\end{quote}

\textbf{log\_ax=1 , log\_ay=2 এবং log\_az=3 হলে,
log\_a((x\^{}3y\^{}2/z)) এর মান কত?}

ক) 1

খ) 2

\textbf{\hl{গ) 4 ✔}}

ঘ) 5

\textbf{\hl{✔ সঠিক উত্তর: গ) 4}}

\begin{quote}
\textbf{ব্যাখ্যা:} \emph{দেওয়া আছে, log\_ax=1 , log\_ay=2 , log\_az=3}

\emph{এখন, log\_a((x\^{}3y\^{}2/z))}

\emph{⇒ log\_a(x\^{}3)+log\_a(y\^{}2)−log\_a(z)}

\emph{⇒ 3log\_ax+2log\_ay−log\_az}

\emph{⇒ 3(1)+2(2)−3}

\emph{⇒ 3+4−3=4}
\end{quote}

\textbf{(1/5)log\_x(2187√(3))=1 হলে x এর মান---}

\textbf{\hl{ক) 3√(3) ✔}}

খ) 3

গ) √(3)

ঘ) (1/√(3))

\textbf{\hl{✔ সঠিক উত্তর: ক) 3√(3)}}

\begin{quote}
\textbf{ব্যাখ্যা:} \emph{দেওয়া আছে, (1/5)log\_x(2187√(3))=1}

\emph{⇒ log\_x(3\^{}7·3\^{}((1/2)))=5}

\emph{⇒ log\_x(3\^{}((15/2)))=5}

\emph{⇒ x\^{}5=3\^{}((15/2))}

\emph{⇒ x=((3\^{}((15/2))))\^{}((1/5))=3\^{}((3/2))=√(3\^{}3)=3√(3)}
\end{quote}

\textbf{৩২ এর ২ ভিত্তিক যুদ্ধে অবদান রাখার জন্য কতজন মহিলাকে বীর প্রতীক
উপাধিতে ভূষিত করা হয়?}

ক) ৩

খ) ৮

\textbf{\hl{গ) ৫ ✔}}

ঘ) ৬

\textbf{\hl{✔ সঠিক উত্তর: গ) ৫}}

\begin{quote}
\textbf{ব্যাখ্যা:} \emph{log\_232=log\_22\^{}5=5×log\_22=5×1=5}
\end{quote}

\textbf{log\_5125+log\_28= কত?}

ক) 8

খ) 125

গ) 5

\textbf{\hl{ঘ) 6 ✔}}

\textbf{\hl{✔ সঠিক উত্তর: ঘ) 6}}

\begin{quote}
\textbf{ব্যাখ্যা:} \emph{log\_5125+log\_28}

\emph{⇒ log\_55\^{}3+log\_22\^{}3}

\emph{⇒ 3log\_55+3log\_22}

\emph{⇒ 3(1)+3(1)=3+3=6}
\end{quote}

\textbf{log\_264+log\_28 এর মান কত?}

\textbf{\hl{ক) 9 ✔}}

খ) 7

গ) 28

ঘ) 72

\textbf{\hl{✔ সঠিক উত্তর: ক) 9}}

\begin{quote}
\textbf{ব্যাখ্যা:} \emph{দেওয়া আছে, log\_264+log\_28}

\emph{⇒ log\_22\^{}6+log\_22\^{}3}

\emph{⇒ 6log\_22+3log\_22}

\emph{⇒ 6+3=9}
\end{quote}

\textbf{6log3−log9= কত?}

\textbf{\hl{ক) log81 ✔}}

খ) log27

গ) log5

ঘ) log10

\textbf{\hl{✔ সঠিক উত্তর: ক) log81}}

\begin{quote}
\textbf{ব্যাখ্যা:} \emph{6log3−log9}

\emph{⇒ 6log3−log3\^{}2}

\emph{⇒ 6log3−2log3}

\emph{⇒ 4log3=log3\^{}4=log81}
\end{quote}

\textbf{5log3−log9= কত?}

ক) log8

\textbf{\hl{খ) log27 ✔}}

গ) log5

ঘ) log10

\textbf{\hl{✔ সঠিক উত্তর: খ) log27}}

\begin{quote}
\textbf{ব্যাখ্যা:} \emph{5log3−log9}

\emph{= 5log3−log3\^{}2}

\emph{= 5log3−2log3}

\emph{= 3log3=log3\^{}3=log27}
\end{quote}

\textbf{3log2+log5= ?}

ক) log30

\textbf{\hl{খ) log40 ✔}}

গ) 3log7

ঘ) log13

ঙ) কোনোটিই নয়

\textbf{\hl{✔ সঠিক উত্তর: খ) log40}}

\begin{quote}
\textbf{ব্যাখ্যা:} \emph{3log2+log5}

\emph{= log2\^{}3+log5}

\emph{= log8+log5}

\emph{= log(8×5)=log40}
\end{quote}

\textbf{log\_2128+log\_216= কত?}

ক) 128

খ) 2

\textbf{\hl{গ) 11 ✔}}

ঘ) 9

\textbf{\hl{✔ সঠিক উত্তর: গ) 11}}

\begin{quote}
\textbf{ব্যাখ্যা:} \emph{log\_2128+log\_216}

\emph{= log\_22\^{}7+log\_22\^{}4}

\emph{= 7log\_22+4log\_22 = 7(1)+4(1)=7+4=11}
\end{quote}

\textbf{log2+log4+log8+\ldots\ldots{} ধারাটির প্রথম দশটি পদের সমষ্টি কত?}

ক) 45log2

\textbf{\hl{খ) 55log2 ✔}}

গ) 65log2

ঘ) 75log2

\textbf{\hl{✔ সঠিক উত্তর: খ) 55log2}}

\begin{quote}
\textbf{ব্যাখ্যা:} \emph{log2+log4+log8+\ldots\ldots{} (১০ম পদ পর্যন্ত)}

\emph{= log2\^{}1+log2\^{}2+log2\^{}3+\ldots\ldots+log2\^{}10}

\emph{= log(2\^{}1·2\^{}2·2\^{}3\ldots\ldots2\^{}10)}

\emph{= log2\^{}(1+2+3+\ldots\ldots+10)=log2\^{}55=55log2}
\end{quote}

\textbf{log3+log6+log12+\ldots\ldots{} ধারাটির ষষ্ঠ পদ কত?}

\textbf{\hl{ক) log96 ✔}}

খ) log48

গ) log192

ঘ) log92

\textbf{\hl{✔ সঠিক উত্তর: ক) log96}}

\begin{quote}
\textbf{ব্যাখ্যা:} \emph{ধারাটি একটি গুণোত্তর ধারা যার প্রথম পদ a=log3 এবং
সাধারণ অনুপাত r=(log6/log3) বা অন্যভাবে ধারাটির পদগুলোর ভিত্তি একই রেখে লেখা
যায়:}

\emph{= log3,log(3×2),log(3×2\^{}2),\ldots{}}

\emph{= ষষ্ঠ পদ =log(3×2\^{}6−1)=log(3×2\^{}5)=log(3×32)=log96}
\end{quote}

\textbf{log\_2√(6)+log\_2√((2/3))= কত?}

ক) 0

খ) 2

\textbf{\hl{গ) 1 ✔}}

ঘ) 3

\textbf{\hl{✔ সঠিক উত্তর: গ) 1}}

\begin{quote}
\textbf{ব্যাখ্যা:} \emph{log\_2√(6)+log\_2√((2/3))}

\emph{= log\_2(6\^{}((1/2)))+log\_2((((2/3)))\^{}((1/2)))}

\emph{= (1/2)log\_26+(1/2)log\_2(2/3)}

\emph{= (1/2)(log\_26+log\_2(2/3))}

\emph{=
(1/2)log\_2(6×(2/3))=(1/2)log\_24=(1/2)log\_22\^{}2=(1/2)×2log\_22=1}
\end{quote}

\textbf{If log\_a2=a and log\_a5=b , then log\_a50=}

ক) a+b

খ) a+b\^{}2

গ) ab\^{}2

\textbf{\hl{ঘ) a+2b ✔}}

\textbf{\hl{✔ সঠিক উত্তর: ঘ) a+2b}}

\begin{quote}
\textbf{ব্যাখ্যা:} \emph{দেওয়া আছে, log\_a2=a এবং log\_a5=b}

\emph{প্রদত্ত রাশি = log\_a50=log\_a(2×5\^{}2)}

\emph{= log\_a2+log\_a5\^{}2}

\emph{= log\_a2+2log\_a5}

\emph{∴ a+2b}
\end{quote}

\section{30. সূচক কী ? সংজ্ঞা, সূত্রাবলী এবং সূচক এর অংক এ টু
জেড}\label{ux9b8ux99aux995-ux995-ux9b8ux99cux99e-ux9b8ux9a4ux9b0ux9acux9b2-ux98fux9ac-ux9b8ux99aux995-ux98fux9b0-ux985ux995-ux98f-ux99f-ux99cux9a1}

\begin{itemize}
\item
  প্রাথমিক আলোচনা
\item
  \textbf{সূচক কাকে বলে?}
\item
  প্রয়োজনীয় সূত্রাবলি
\item
  সাধারণ সূচকের প্রশ্ন
\item
  সূচকের মান নির্ণয়
\item
  সূচকের সমাধান
\item
  গুরুত্বপূর্ণ কিছু প্রশ্ন ও সমাধান
\end{itemize}

\subsubsection{প্রাথমিক
আলোচনা}\label{ux9aaux9b0ux9a5ux9aeux995-ux986ux9b2ux99aux9a8-12}

\textbf{সূচক :} সূচক শব্দের অর্থ হলো শক্তি। n সংখ্যক a- এর ক্রমিক গুণফল

= \(a^{n}\) {[} এখানে, a কে ভিত্তি বলা হয়। {]}

n কে a এর (ভিত্তির) সূচক বলা হয় বা শক্তি বলা হয় ।

\(a^{n}\) কে a এর n তম ঘাত বা শক্তি বা `power\textquotesingle{} বলা হয়।

\subsubsection{সূচক কাকে
বলে?}\label{ux9b8ux99aux995-ux995ux995-ux9acux9b2}

সূচক হলো শক্তি বা ঘাতের একটি মান । একে ইংরেজি ভাষায় index বলা হয় । একে
পাওয়ার ও বলা হয়ে থাকে । সহজ ভাষায় বলতে গেলে কোন একটি সংখ্যার উপরে ডানদিকে
একটি ছোট সংখ্যা থাকে এটি হল ঐ মূল সংখ্যার সূচক । আবার বলা যায়, যখন কোনো একটি
সংখ্যা অন্য একটি সংখ্যার সাথে পরিবর্তনশীল হিসেবে উত্থিত হয়, তখন তাকে ঐ সংখ্যার
সূচক বলে। আর মূল সংখ্যাটিকে বলা হয় সূচকীয় রাশি । উদাহরণস্বরুপ-

\(a^{n}\) এখানে, a কে ভিত্তি বা সূচকীয় রাশি বলা হয়। আর n কে বলা হয় সূচক বা
শক্তি বা ঘাত অথবা পাওয়ার । এটি সংখ্যার সাহায্যে দেখালেও এমনি হবে যেমন-

\(2^{3}\) এখানে 2 এর সূচক হল 3 অর্থাৎ 2 এর মান 3 ঘাত বিশিষ্ট বা 3 বার
ব্যবহার করতে হবে । আর 2 হল ভিত্তি বা সূচকীয় রাশি ।

\subsubsection{প্রয়োজনীয়
সূত্রাবলি}\label{ux9aaux9b0ux9afux99cux9a8ux9af-ux9b8ux9a4ux9b0ux9acux9b2}

১. \(a^{n}\) = a×a×a×...........{[} n সংখ্যক a {]}

২. \({(a^{m})}^{n} = a^{mn}\) {[} কোন সংখ্যার উপর দুইবার বা তার থেকে বেশি
পাওয়ার থাকলে তা গুণ হয় {]}

৩. \(a^{m}.a^{n} = a^{m + n}\)

{[} Note: ভিত্তি একই হলে এবং গুণ থাকলে পাওয়ার গুলো যোগ করতে হয় আবার ভাগ
থাকলে পাওয়ার গুলো বিয়োগ করতে হয়। {]}

৪. \(a^{m} \div a^{n}\) = \(\frac{a^{m}}{a^{n}} = a^{m - n}\)

৫. \(a^{0} = 1\)

{[} যে কোন সংখ্যার উপর পাওয়ার 0 হলে তার মান 1 হয় । যেখানে a ≠ 0 {]}

৬. \(\sqrt{a} = a^{\frac{1}{2}}\)

৭. \(\sqrt[n]{a} = a^{\frac{1}{n}}\)

৮. \(\sqrt[3]{a} = a^{\frac{1}{3}}\)

৯. \(a^{- n} = \frac{1}{a^{n}}\)

{[} Note: কোন পাওয়ার মাইনাস থাকলে তা ভগ্নাংশ আকারে লিখতে হয় এবং মাইনাস
তুলে লব 1 এর নিচে পুরো সংখ্যাটি পাওয়ার সহ লিখতে হয় {]}

১০. \((ab)^{m} = a^{m}.b^{m}\)

১১. \({(\frac{a}{b})}^{m} = \frac{a^{m}}{b^{m}}\)

১২. \({(\frac{m}{n})}^{- \ p} = {(\frac{n}{m})}^{p}\)

{[} কোন ভগ্নাংশের উপরের পাওয়ার টি মাইনাস হলে ঐ ভগ্নাংশটি উল্টে যায় অর্থাৎ
হরের জায়গায় লব এবং লবের জায়গায় হর বসে । {]}

১৩. \(a^{x} = a^{y}\) হলে, x = y {[} অর্থাৎ দু পাশের ভিত্তি মিলে গেলে দুটো
ভিত্তিই তুলে দিতে হয় {]}

১৪. \(a^{x} = b^{x}\) হলে, a = b {[} অর্থাৎ দু পাশের power মিলে গেলে দুটো
পাওয়ারই বাদ দিতে হয় {]}

\subsubsection{সাধারণ সূচকের
প্রশ্ন}\label{ux9b8ux9a7ux9b0ux9a3-ux9b8ux99aux995ux9b0-ux9aaux9b0ux9b6ux9a8}

\textbf{প্রশ্ন নম্বর-} ১. \(a^{0}\) = কত? {[} বিএসটিআই ল্যাব সহকারী
পরীক্ষা-১০ / ৩১ তম বিসিএস, মানসিক দক্ষতা {]}

\section{31. সূচক MCQ প্রশ্ন উত্তর
অনুশীলন}\label{ux9b8ux99aux995-mcq-ux9aaux9b0ux9b6ux9a8-ux989ux9a4ux9a4ux9b0-ux985ux9a8ux9b6ux9b2ux9a8}

\textbf{(a\^{}−1)\^{}−1 এর মান নিচের কোনটি?}

ক) (1/a)

খ) (1/a\^{}2)

\textbf{\hl{গ) a ✔}}

ঘ) a\^{}2

\textbf{\hl{✔ সঠিক উত্তর: গ) a}}

\begin{quote}
\textbf{ব্যাখ্যা:} \emph{আমরা জানি, a\^{}−1=(1/a) । সুতরাং,
(a\^{}−1)\^{}−1=(((1/a)))\^{}−1=(1/(1/a))=a । তাছাড়া সূচকের নিয়ম অনুযায়ী
পাওয়ারে পাওয়ারে গুণ হয়ে a\^{}(−1)×(−1)=a\^{}1=a হয়।}

\emph{∴ a}
\end{quote}

\textbf{(3x)\^{}0 এর মান কত?}

\textbf{\hl{ক) 1 ✔}}

খ) -1

গ) 3x

ঘ) 0

\textbf{\hl{✔ সঠিক উত্তর: ক) 1}}

\begin{quote}
\textbf{ব্যাখ্যা:} \emph{সূচকের সূত্রানুসারে, যেকোনো অশূন্য সংখ্যা বা রাশির ওপর
পাওয়ার শূন্য (0) হলে তার মান সবসময় 1 হয়। অর্থাৎ, a\^{}0=1 । সুতরাং,
(3x)\^{}0=1 ।}

\emph{∴ 1}
\end{quote}

\textbf{m√(a\^{}n) এর মান নিচের কোনটি?}

ক) a\^{}mn

\textbf{\hl{খ) a\^{}((n/m)) ✔}}

গ) a\^{}((m/n))

ঘ) a\^{}(−(n/m))

\textbf{\hl{✔ সঠিক উত্তর: খ) a\^{}((n/m))}}

\begin{quote}
\textbf{ব্যাখ্যা:} \emph{সূচকের নিয়ম অনুযায়ী,
m√(a\^{}n)=(a\^{}n)\^{}((1/m))=a\^{}((n/m)) ।}

\emph{∴ a\^{}((n/m))}
\end{quote}

\textbf{4\^{}x+4\^{}x+4\^{}x+4\^{}x এর মান কত?}

ক) 4\^{}x

\textbf{\hl{খ) 4\^{}(x+1) ✔}}

গ) 4\^{}(x+4)

ঘ) 4\^{}(x+16)

\textbf{\hl{✔ সঠিক উত্তর: খ) 4\^{}(x+1)}}

\begin{quote}
\textbf{ব্যাখ্যা:} \emph{প্রদত্ত রাশি = 4\^{}x+4\^{}x+4\^{}x+4\^{}x=4×4\^{}x
। গুণের সময় বেত এক থাকলে পাওয়ারগুলো যোগ হয়, সুতরাং 4\^{}1×4\^{}x=4\^{}(x+1)
।}

\emph{∴ 4\^{}(x+1)}
\end{quote}

\textbf{9\^{}x+9\^{}x+9\^{}x = কত?}

ক) 27x

\textbf{\hl{খ) 3\^{}(2x+1) ✔}}

গ) 9\^{}3x

ঘ) 3x\^{}3

\textbf{\hl{✔ সঠিক উত্তর: খ) 3\^{}(2x+1)}}

\begin{quote}
\textbf{ব্যাখ্যা:} \emph{প্রদত্ত রাশি =
9\^{}x+9\^{}x+9\^{}x=3×9\^{}x=3×(3\^{}2)\^{}x=3×3\^{}2x=3\^{}(2x+1) ।}

\emph{∴ 3\^{}(2x+1)}
\end{quote}

\textbf{(3x)\^{}0+3(x)\^{}0 = কত?}

\textbf{\hl{ক) 4 ✔}}

খ) 9

গ) 12

ঘ) 16

\textbf{\hl{✔ সঠিক উত্তর: ক) 4}}

\begin{quote}
\textbf{ব্যাখ্যা:} \emph{আমরা জানি, যেকোনো কিছুর ওপর পাওয়ার শূন্য হলে তার মান
১ হয় ( a\^{}0=1 )।}

\emph{সুতরাং, (3x)\^{}0+3(x)\^{}0=1+3×1=1+3=4 ।}

\emph{∴ 4}
\end{quote}

\textbf{2\^{}(x+1)=32 হলে x এর মান কত?}

\textbf{\hl{ক) 4 ✔}}

খ) 3

গ) 2

ঘ) 1

\textbf{\hl{✔ সঠিক উত্তর: ক) 4}}

\begin{quote}
\textbf{ব্যাখ্যা:} \emph{সমাধান: 2\^{}(x+1)=32}

\emph{⇒ 2\^{}(x+1)=2\^{}5}

\emph{⇒ x+1=5}

\emph{⇒ x=5−1=4}

\emph{∴ 4}
\end{quote}

\textbf{3\^{}x−n=5\^{}x−n=0 হলে x = কত?}

ক) 5

খ) 0

\textbf{\hl{গ) n ✔}}

ঘ) 3

\textbf{\hl{✔ সঠিক উত্তর: গ) n}}

\begin{quote}
\textbf{ব্যাখ্যা:} \emph{প্রশ্নানুসারে, 3\^{}x−n=5\^{}x−n}

\emph{⇒ (3\^{}x−n/5\^{}x−n)=1}

\emph{⇒ (((3/5)))\^{}x−n=(((3/5)))\^{}0}

\emph{⇒ x−n=0}

\emph{⇒ x=n}

\emph{∴ n}
\end{quote}

\textbf{3\^{}(x+4)=3.9\^{}x হলে x = কত?}

ক) 2

\textbf{\hl{খ) 3 ✔}}

গ) 4

ঘ) 9

\textbf{\hl{✔ সঠিক উত্তর: খ) 3}}

\begin{quote}
\textbf{ব্যাখ্যা:} \emph{সমাধান: 3\^{}(x+4)=3×(3\^{}2)\^{}x}

\emph{⇒ 3\^{}(x+4)=3\^{}1×3\^{}2x}

\emph{⇒ 3\^{}(x+4)=3\^{}(2x+1)}

\emph{⇒ x+4=2x+1}

\emph{⇒ 2x−x=4−1}

\emph{⇒ x=3}

\emph{∴ 3}
\end{quote}

\textbf{3√(1+x)=2 হলে x এর মান কত?}

ক) ১

খ) ২

গ) ৫

\textbf{\hl{ঘ) ৭ ✔}}

\textbf{\hl{✔ সঠিক উত্তর: ঘ) ৭}}

\begin{quote}
\textbf{ব্যাখ্যা:} \emph{সমাধান: উভয় পক্ষে ঘন করে পাই,
(3√(1+x))\^{}3=2\^{}3}

\emph{⇒ 1+x=8}

\emph{⇒ x=8−1=7}

\emph{∴ ৭}
\end{quote}

\textbf{(√(3))\^{}(x+1)=(3√(3))\^{}2x−1 এর সমাধান কত?}

ক) 4

\textbf{\hl{খ) 5 ✔}}

গ) 6

ঘ) 7

\textbf{\hl{✔ সঠিক উত্তর: খ) 5}}

\begin{quote}
\textbf{ব্যাখ্যা:} \emph{সমাধান:
((3\^{}((1/2))))\^{}(x+1)=((3\^{}((1/3))))\^{}2x−1}

\emph{⇒ 3\^{}((x+1/2))=3\^{}((2x−1/3))}

\emph{⇒ (x+1/2)=(2x−1/3)}

\emph{⇒ 4x−2=3x+3}

\emph{⇒ x=5}

\emph{∴ 5}
\end{quote}

\textbf{3√(3√(x\^{}3)) = কত?}

ক) x\^{}((1/2))

খ) x\^{}((1/1))

গ) x\^{}((3/2))

\textbf{\hl{ঘ) x\^{}((1/3)) ✔}}

\textbf{\hl{✔ সঠিক উত্তর: ঘ) x\^{}((1/3))}}

\begin{quote}
\textbf{ব্যাখ্যা:}
\emph{3√(3√(x\^{}3))=3√((x\^{}3)\^{}((1/3)))=3√(x\^{}1)=(x\^{}1)\^{}((1/3))=x\^{}((1/3))}

\emph{∴ x\^{}((1/3))}
\end{quote}

\textbf{√((1/0.25))=?}

ক) .025

খ) 25

গ) .002

\textbf{\hl{ঘ) 2 ✔}}

\textbf{\hl{✔ সঠিক উত্তর: ঘ) 2}}

\begin{quote}
\textbf{ব্যাখ্যা:} \emph{√((1/0.25))=√((1/(1/4)))=√(4)=2}

\emph{∴ 2}
\end{quote}

\textbf{28√(x)+1426=three-fourths of 2984 , find x.}

ক) 659

খ) 694

\textbf{\hl{গ) 841 ✔}}

ঘ) 859

\textbf{\hl{✔ সঠিক উত্তর: গ) 841}}

\begin{quote}
\textbf{ব্যাখ্যা:} \emph{28√(x)+1426=(3/4)×2984}

\emph{⇒ 28√(x)+1426=3×746=2238}

\emph{⇒ 28√(x)=2238−1426=812}

\emph{⇒ √(x)=(812/28)=29}

\emph{⇒ x=29\^{}2=841}

\emph{∴ 841}
\end{quote}

\textbf{√(5x−4)=−1 হলে x এর মান কত?}

\textbf{\hl{ক) 1 ✔}}

খ) 3

গ) 4

ঘ) 8

\textbf{\hl{✔ সঠিক উত্তর: ক) 1}}

\begin{quote}
\textbf{ব্যাখ্যা:} \emph{√(5x−4)=−1}

\emph{⇒ (√(5x−4))\^{}2=(−1)\^{}2}

\emph{⇒ 5x−4=1}

\emph{⇒ 5x=5}

\emph{∴ x=1}
\end{quote}

\textbf{p\^{}a=q,q\^{}b=r,r\^{}c=p হলে abc = ?}

ক) 0

\textbf{\hl{খ) 1 ✔}}

গ) a\^{}p

ঘ) p\^{}a

\textbf{\hl{✔ সঠিক উত্তর: খ) 1}}

\begin{quote}
\textbf{ব্যাখ্যা:} \emph{দেওয়া আছে, p\^{}a=q,q\^{}b=r,r\^{}c=p}

\emph{এখন, p=r\^{}c=(q\^{}b)\^{}c=(p\^{}a)\^{}bc}

\emph{⇒ p\^{}1=p\^{}abc}

\emph{∴ abc=1}
\end{quote}

\textbf{If (√(3))\^{}(2x+5)=27 , then x = ?}

ক) (2/3)

খ) (3/5)

\textbf{\hl{গ) (1/2) ✔}}

ঘ) (3/2)

\textbf{\hl{✔ সঠিক উত্তর: গ) (1/2)}}

\begin{quote}
\textbf{ব্যাখ্যা:} \emph{দেওয়া আছে, (√(3))\^{}(2x+5)=27}

\emph{⇒ 3\^{}((2x+5/2))=3\^{}3}

\emph{⇒ (2x+5/2)=3}

\emph{⇒ 2x+5=6}

\emph{⇒ 2x=1}

\emph{∴ x=(1/2)}
\end{quote}

\textbf{3.2\^{}n−4.2\^{}n−2 = কত?}

ক) 1

\textbf{\hl{খ) 2\^{}(n+1) ✔}}

গ) 3

ঘ) 2\^{}n

\textbf{\hl{✔ সঠিক উত্তর: খ) 2\^{}(n+1)}}

\begin{quote}
\textbf{ব্যাখ্যা:} \emph{দেওয়া আছে, 3⋅2\^{}n−2\^{}2⋅2\^{}n−2}

\emph{⇒ 3⋅2\^{}n−2\^{}(n−2+2)}

\emph{⇒ 3⋅2\^{}n−2\^{}n}

\emph{⇒ 2\^{}n(3−1)}

\emph{⇒ 2\^{}n⋅2}

\emph{∴ 2\^{}(n+1)}
\end{quote}

\textbf{যদি হয় x\^{}4−1=0 , তবে x এর মান কত?}

\textbf{\hl{ক) 1 ✔}}

খ) 0

গ) 2

ঘ) 4

\textbf{\hl{✔ সঠিক উত্তর: ক) 1}}

\begin{quote}
\textbf{ব্যাখ্যা:} \emph{দেয়া আছে, x\^{}4−1=0}

\emph{⇒ x\^{}4=1}

\emph{⇒ x\^{}4=1\^{}4}

\emph{∴ x=1}
\end{quote}

\textbf{(√(3)×√(5))\^{}4 এর মান কত?}

ক) 255

খ) 259

\textbf{\hl{গ) 225 ✔}}

ঘ) 25

\textbf{\hl{✔ সঠিক উত্তর: গ) 225}}

\begin{quote}
\textbf{ব্যাখ্যা:} \emph{দেয়া আছে,
(√(3)×√(5))\^{}4=\{(√(3)×√(5))\^{}2\}\^{}2}

\emph{⇒ (3×5)\^{}2=15\^{}2=225}
\end{quote}

\textbf{8\^{}((3/4))÷8\^{}((1/2)) এর মান কত?}

ক) 3√(4)

\textbf{\hl{খ) 4√(8) ✔}}

গ) 4√(3)

ঘ) 3√(16)

\textbf{\hl{✔ সঠিক উত্তর: খ) 4√(8)}}

\begin{quote}
\textbf{ব্যাখ্যা:} \emph{8\^{}((3/4))÷8\^{}((1/2))=8\^{}((3/4)−(1/2))}

\emph{⇒ 8\^{}((3−2/4))=8\^{}((1/4))=4√(8)}
\end{quote}

\textbf{3\^{}−3 এর মান কত?}

ক) (1/9)

খ) (1/3)

\textbf{\hl{গ) (1/27) ✔}}

ঘ) √(3)

\textbf{\hl{✔ সঠিক উত্তর: গ) (1/27)}}

\begin{quote}
\textbf{ব্যাখ্যা:} \emph{3\^{}−3=(1/3\^{}3)=(1/27)}
\end{quote}

\textbf{2\^{}n÷2\^{}n−1 = কত?}

\textbf{\hl{ক) 2 ✔}}

খ) 2\^{}2

গ) 2\^{}n

ঘ) 2\^{}(n+1)

\textbf{\hl{✔ সঠিক উত্তর: ক) 2}}

\begin{quote}
\textbf{ব্যাখ্যা:} \emph{2\^{}n÷2\^{}n−1}

\emph{=2\^{}n−(n−1)}

\emph{=2\^{}(n−n+1) =2\^{}1 =2}
\end{quote}

\textbf{যদি x, y বাস্তব সংখ্যা এবং x≠0,y≠0 হয়, তবে x\^{}x\^{}0+y\^{}y\^{}0
এর মান---}

\textbf{\hl{ক) x + y ✔}}

খ) 2

গ) 0

ঘ) x\^{}2+y\^{}2

\textbf{\hl{✔ সঠিক উত্তর: ক) x + y}}

\begin{quote}
\textbf{ব্যাখ্যা:} \emph{x\^{}x\^{}0+y\^{}y\^{}0}

\emph{=x\^{}1+y\^{}1}

\emph{=x+y}
\end{quote}

\textbf{যদি a, b বাস্তব সংখ্যা এবং a≠0,b≠0 হয় তবে
a\^{}2b\^{}0+b\^{}2a\^{}0 এর মান}

ক) a + b

খ) 2

গ) 0

\textbf{\hl{ঘ) a\^{}2+b\^{}2 ✔}}

\textbf{\hl{✔ সঠিক উত্তর: ঘ) a\^{}2+b\^{}2}}

\begin{quote}
\textbf{ব্যাখ্যা:} \emph{a\^{}2b\^{}0+b\^{}2a\^{}0}

\emph{=a\^{}2×1+b\^{}2×1}

\emph{=a\^{}2+b\^{}2}
\end{quote}

\textbf{x,y,z\textgreater0 হলে
√(x\^{}−2y\^{}2)⋅√(y\^{}−2z\^{}2)⋅√(z\^{}−2x\^{}2) এর মান -}

ক) 0

\textbf{\hl{খ) 1 ✔}}

গ) xyz

ঘ) y

\textbf{\hl{✔ সঠিক উত্তর: খ) 1}}

\begin{quote}
\textbf{ব্যাখ্যা:}
\emph{√(x\^{}−2y\^{}2)⋅√(y\^{}−2z\^{}2)⋅√(z\^{}−2x\^{}2)}

\emph{= √((y\^{}2/x\^{}2))⋅√((z\^{}2/y\^{}2))⋅√((x\^{}2/z\^{}2))}

\emph{= (y/x)⋅(z/y)⋅(x/z) = 1}
\end{quote}

\textbf{√(x\^{}−1y)⋅√(y\^{}−1z)⋅√(z\^{}−1x) এর মান কত?}

ক) x

খ) y

\textbf{\hl{গ) 1 ✔}}

ঘ) z

\textbf{\hl{✔ সঠিক উত্তর: গ) 1}}

\begin{quote}
\textbf{ব্যাখ্যা:} \emph{√(x\^{}−1y)⋅√(y\^{}−1z)⋅√(z\^{}−1x)}

\emph{= √((y/x))⋅√((z/y))⋅√((x/z))}

\emph{= √((y/x)⋅(z/y)⋅(x/z)) = √(1) = 1}
\end{quote}

\textbf{{[}2−(3−1)\^{}−1{]}\^{}−1 = কত?}

\textbf{\hl{ক) (2/3) ✔}}

খ) −(1/2)

গ) 1

ঘ) −1

\textbf{\hl{✔ সঠিক উত্তর: ক) (2/3)}}

\begin{quote}
\textbf{ব্যাখ্যা:} \emph{{[}2−(3−1)\^{}−1{]}\^{}−1}

\emph{⇒ {[}2−(2)\^{}−1{]}\^{}−1}

\emph{⇒ ((2−(1/2)))\^{}−1}

\emph{⇒ (((3/2)))\^{}−1}

\emph{∴ (2/3)}
\end{quote}

\textbf{6√(64)×3√(27)}

ক) 2

খ) 4

\textbf{\hl{গ) 6 ✔}}

ঘ) 8

\textbf{\hl{✔ সঠিক উত্তর: গ) 6}}

\begin{quote}
\textbf{ব্যাখ্যা:} \emph{6√(64)×3√(27)}

\emph{= 64\^{}((1/6))×27\^{}((1/3))}

\emph{= (4\^{}3)\^{}((1/6))×(3\^{}3)\^{}((1/3))}

\emph{= 4\^{}((1/2))×3 = 2×3 = 6}
\end{quote}

\textbf{3√(125×9) = কত?}

ক) 25

খ) 15

\textbf{\hl{গ) 53√(9) ✔}}

ঘ) 10

\textbf{\hl{✔ সঠিক উত্তর: গ) 53√(9)}}

\begin{quote}
\textbf{ব্যাখ্যা:} \emph{দেয়া আছে, 3√(125×9) = 3√(5\^{}3×3\^{}2)}

\emph{= (5\^{}3×3\^{}2)\^{}((1/3))}

\emph{= 5\^{}(3×(1/3))×3\^{}(2×(1/3)) = 5×3√(3\^{}2) = 53√(9)}
\end{quote}

\textbf{(3√(3)×3√(4))\^{}6 = কত?}

ক) 12

খ) 36

গ) 48

\textbf{\hl{ঘ) 144 ✔}}

\textbf{\hl{✔ সঠিক উত্তর: ঘ) 144}}

\begin{quote}
\textbf{ব্যাখ্যা:} \emph{(3√(3)×3√(4))\^{}6}

\emph{= ((3\^{}((1/3))×4\^{}((1/3))))\^{}6}

\emph{= (\{(3×4)\^{}((1/3))\})\^{}6}

\emph{= (12)\^{}((1/3)×6) = 12\^{}2 = 144}
\end{quote}

\textbf{√((100)\^{}2) = ?}

ক) 10

\textbf{\hl{খ) 100 ✔}}

গ) 1000

ঘ) 10000

\textbf{\hl{✔ সঠিক উত্তর: খ) 100}}

\begin{quote}
\textbf{ব্যাখ্যা:} \emph{√((100)\^{}2) = 100}
\end{quote}

\textbf{3√(−27)=?}

ক) 3

খ) 9

\textbf{\hl{গ) −3 ✔}}

ঘ) −9

\textbf{\hl{✔ সঠিক উত্তর: গ) −3}}

\begin{quote}
\textbf{ব্যাখ্যা:} \emph{3√(−27)}

\emph{= 3√((−3)\^{}3)}

\emph{= −3}
\end{quote}

\textbf{3√(125×0.008) কত?}

ক) 0.001

\textbf{\hl{খ) 1 ✔}}

গ) 0.01

ঘ) 0

\textbf{\hl{✔ সঠিক উত্তর: খ) 1}}

\begin{quote}
\textbf{ব্যাখ্যা:} \emph{3√(125×0.008) = 3√(1)}

\emph{= (1)\^{}((1/3)) = 1}
\end{quote}

\textbf{((√(3)+√(2)))\^{}5 এর মান কত?}

\textbf{\hl{ক) 308 ✔}}

খ) 36

গ) 256

ঘ) 144

\textbf{\hl{✔ সঠিক উত্তর: ক) 308}}

\begin{quote}
\textbf{ব্যাখ্যা:} \emph{((√(3)+√(2)))\^{}5}

\emph{=
(√(3))\^{}5+5(√(3))\^{}4(√(2))+10(√(3))\^{}3(√(2))\^{}2+10(√(3))\^{}2(√(2))\^{}3+5(√(3))(√(2))\^{}4+(√(2))\^{}5}

\emph{= 9√(3)+45√(2)+60√(3)+60√(2)+20√(3)+4√(2)}

\emph{= 89√(3)+109√(2)≈154.15+154.15=308.30}
\end{quote}

\textbf{3\^{}x+3\^{}x+3\^{}x = কত?}

ক) 9\^{}x

\textbf{\hl{খ) 3\^{}(x+1) ✔}}

গ) 3\^{}3x

ঘ) 3\^{}x\^{}3

\textbf{\hl{✔ সঠিক উত্তর: খ) 3\^{}(x+1)}}

\begin{quote}
\textbf{ব্যাখ্যা:} \emph{3\^{}x+3\^{}x+3\^{}x}

\emph{= 3.3\^{}x}

\emph{= 3\^{}(1+x)}

\emph{= 3\^{}(x+1)}
\end{quote}

\textbf{(√(3)×√(5))\^{}4 এর মান কত? / What is the value of
(√(3)×√(5))\^{}4 ?}

ক) 30

খ) 60

\textbf{\hl{গ) 225 ✔}}

ঘ) 15

\textbf{\hl{✔ সঠিক উত্তর: গ) 225}}

\begin{quote}
\textbf{ব্যাখ্যা:} \emph{(√(3)×√(5))\^{}4}

\emph{= ((3\^{}((1/2))×5\^{}((1/2))))\^{}4}

\emph{= (15)\^{}((1/2)×4) = 15\^{}2 = 225}
\end{quote}

\textbf{8\^{}2×8\^{}0 = ?}

ক) 0

খ) 4

\textbf{\hl{গ) 64 ✔}}

ঘ) 16

\textbf{\hl{✔ সঠিক উত্তর: গ) 64}}

\begin{quote}
\textbf{ব্যাখ্যা:} \emph{8\^{}2×8\^{}0}

\emph{⇒ 8\^{}(2+0)}

\emph{⇒ 8\^{}2}

\emph{∴ 64}
\end{quote}

\textbf{যদি a≠0 তবে a\^{}0 = কত?}

ক) 0

খ) a

\textbf{\hl{গ) 1 ✔}}

ঘ) অনির্ণেয়

\textbf{\hl{✔ সঠিক উত্তর: গ) 1}}

\begin{quote}
\textbf{ব্যাখ্যা:} \emph{যেকোনো অশূন্য বাস্তব সংখ্যার ওপর পাওয়ার শূন্য (0) হলে
তার মান 1 হয়।}

\emph{∴ a\^{}0=1 (যেখানে a≠0 )}
\end{quote}

\textbf{(5x)\^{}0 এর মান নিচের কোনটি?}

ক) 5x

খ) 0

গ) 5

\textbf{\hl{ঘ) 1 ✔}}

\textbf{\hl{✔ সঠিক উত্তর: ঘ) 1}}

\begin{quote}
\textbf{ব্যাখ্যা:} \emph{(5x)\^{}0}

\emph{= 1 {[} ∴a\^{}0=1 যেখানে a≠0 {]}}
\end{quote}

\textbf{a≠0 হলে, (a\^{}−1)−1 এর সঠিক মান নিচের কোনটি?}

ক) a\^{}4

খ) a\^{}2

\textbf{\hl{গ) a ✔}}

ঘ) a\^{}−2

\textbf{\hl{✔ সঠিক উত্তর: গ) a}}

\begin{quote}
\textbf{ব্যাখ্যা:} \emph{(a\^{}−1)−1}

\emph{⇒ (((1/a)))\^{}−1}

\emph{⇒ (1/(1/a))}

\emph{⇒ (1×a/1)}

\emph{∴ a}
\end{quote}

\textbf{2\^{}(n+1)−2\^{}n = কত?}

ক) 2

\textbf{\hl{খ) 2\^{}n ✔}}

গ) 4

ঘ) 2\^{}(n+1)

\textbf{\hl{✔ সঠিক উত্তর: খ) 2\^{}n}}

\begin{quote}
\textbf{ব্যাখ্যা:} \emph{2\^{}(n+1)−2\^{}n}

\emph{⇒ 2\^{}n⋅2\^{}1−2\^{}n}

\emph{⇒ 2\^{}n(2−1)}

\emph{⇒ 2\^{}n×1}

\emph{∴ 2\^{}n}
\end{quote}

\textbf{3.2\^{}n−4.2\^{}n−2 এর সরলমান নিচের কোনটি?}

ক) 1

খ) 2\^{}n

\textbf{\hl{গ) 2\^{}(n+1) ✔}}

ঘ) 3\^{}n−1

\textbf{\hl{✔ সঠিক উত্তর: গ) 2\^{}(n+1)}}

\begin{quote}
\textbf{ব্যাখ্যা:} \emph{3.2\^{}n−4.2\^{}n−2}

\emph{⇒ 3.2\^{}n−4.2\^{}n.2\^{}−2}

\emph{⇒ 3.2\^{}n−(4.2\^{}n/2\^{}2)}

\emph{⇒ 3.2\^{}n−(4.2\^{}n/4)}

\emph{⇒ 3.2\^{}n−2\^{}n}

\emph{⇒ (3−1)2\^{}n}

\emph{⇒ 2.2\^{}n}

\emph{∴ 2\^{}(n+1)}
\end{quote}

\textbf{যদি 3\^{}(x+2)=81 হয়, তবে 3\^{}x−2 = কত?}

\textbf{\hl{ক) 1 ✔}}

খ) 0

গ) 4

ঘ) 3

\textbf{\hl{✔ সঠিক উত্তর: ক) 1}}

\begin{quote}
\textbf{ব্যাখ্যা:} \emph{দেয়া আছে, 3\^{}(x+2)=81}

\emph{⇒ 3\^{}(x+2)=3\^{}4}

\emph{⇒ x+2=4}

\emph{⇒ x=2}

\emph{এখন, 3\^{}x−2=3\^{}2−2=3\^{}0=1}
\end{quote}

\textbf{x\textgreater0 এবং p,q,r মূলদ সংখ্যা হলে
(((x\^{}(p+q)/x\^{}2r)))\^{}(((x\^{}(q+r)/x\^{}2p)))\^{}(((x\^{}(r+p)/x\^{}2q)))\^{}}

ক) p+q+r

\textbf{\hl{খ) 1 ✔}}

গ) 0

ঘ) −1

\textbf{\hl{✔ সঠিক উত্তর: খ) 1}}

\begin{quote}
\textbf{ব্যাখ্যা:}
\emph{((x\^{}(p+q)/x\^{}2r))((x\^{}(q+r)/x\^{}2p))((x\^{}(r+p)/x\^{}2q))}

\emph{= x\^{}(p+q−2r)⋅x\^{}(q+r−2p)⋅x\^{}(r+p−2q)}

\emph{= x\^{}(p+q−2r+q+r−2p+r+p−2q)}

\emph{= x\^{}0 = 1}
\end{quote}

\textbf{If 16(4\^{}a)=1 , find the value of a -}

\textbf{\hl{ক) −2 ✔}}

খ) 0

গ) (1/2)

ঘ) 3

ঙ) 2

\textbf{\hl{✔ সঠিক উত্তর: ক) −2}}

\begin{quote}
\textbf{ব্যাখ্যা:} \emph{16(4\^{}a)=1}

\emph{⇒ 4\^{}2(4\^{}a)=4\^{}0}

\emph{⇒ 4\^{}(a+2)=4\^{}0}

\emph{⇒ a+2=0}

\emph{∴ a=−2}
\end{quote}

\textbf{9\^{}(x+3)=27\^{}(x+1) হলে, x এর মান কত?}

ক) 2

খ) −3

গ) 9

\textbf{\hl{ঘ) 3 ✔}}

\textbf{\hl{✔ সঠিক উত্তর: ঘ) 3}}

\begin{quote}
\textbf{ব্যাখ্যা:} \emph{9\^{}(x+3)=27\^{}(x+1)}

\emph{⇒ (3\^{}2)\^{}(x+3)=(3\^{}3)\^{}(x+1)}

\emph{⇒ 3\^{}(2x+6)=3\^{}(3x+3)}

\emph{⇒ 2x+6=3x+3}

\emph{∴ x=3}
\end{quote}

\textbf{125(√(5))\^{}2x=1 হলে x এর মান কত?}

ক) 3

\textbf{\hl{খ) −3 ✔}}

গ) 7

ঘ) 9

\textbf{\hl{✔ সঠিক উত্তর: খ) −3}}

\begin{quote}
\textbf{ব্যাখ্যা:} \emph{125(√(5))\^{}2x=1}

\emph{⇒ 5\^{}3((5\^{}((1/2))))\^{}2x=1}

\emph{⇒ 5\^{}3⋅5\^{}x=5\^{}0}

\emph{⇒ 5\^{}(3+x)=5\^{}0}

\emph{⇒ 3+x=0}

\emph{∴ x=−3}
\end{quote}

\textbf{2\^{}x+2\^{}1−x=3 হলে x = কত?}

ক) (1, 1)

খ) (0, 2)

গ) (1, 3)

\textbf{\hl{ঘ) (0, 1) ✔}}

\textbf{\hl{✔ সঠিক উত্তর: ঘ) (0, 1)}}

\begin{quote}
\textbf{ব্যাখ্যা:} \emph{সমীকরণটিতে অপশন টেস্ট করে পাই:}

\emph{যখন x=0 , বামপক্ষ = 2\^{}0+2\^{}1−0=1+2=3 (ডানপক্ষ)}

\emph{যখন x=1 , বামপক্ষ = 2\^{}1+2\^{}1−1=2+2\^{}0=2+1=3 (ডানপক্ষ)}

\emph{∴ (0,1)}
\end{quote}

\textbf{nC\_12=nC\_6 হলে n এর মান কত?}

ক) ১২

খ) ১৪

গ) ১৬

\textbf{\hl{ঘ) ১৮ ✔}}

\textbf{\hl{✔ সঠিক উত্তর: ঘ) ১৮}}

\begin{quote}
\textbf{ব্যাখ্যা:} \emph{nC\_12=nC\_6}

\emph{∴ n=12+6=18 {[} ∴nC\_x=nC\_y হলে, x=y অথবা x+y=n হবে।{]}}
\end{quote}

\textbf{If 3\^{}3y=729 then 3\^{}y equals:}

ক) 3

খ) 8

\textbf{\hl{গ) 9 ✔}}

ঘ) 10

\textbf{\hl{✔ সঠিক উত্তর: গ) 9}}

\begin{quote}
\textbf{ব্যাখ্যা:} \emph{3\^{}3y=729}

\emph{⇒ 3\^{}3y=3\^{}6}

\emph{⇒ 3y=6}

\emph{⇒ y=2}

\emph{∴ 3\^{}y=3\^{}2=9}
\end{quote}

\textbf{If 7\^{}(x+1)=343 then find the value of x -}

ক) 4

খ) 3

\textbf{\hl{গ) 2 ✔}}

ঘ) 1

\textbf{\hl{✔ সঠিক উত্তর: গ) 2}}

\begin{quote}
\textbf{ব্যাখ্যা:} \emph{7\^{}(x+1)=343}

\emph{⇒ 7\^{}(x+1)=7\^{}3}

\emph{⇒ x+1=3}

\emph{∴ x=2}
\end{quote}

\textbf{4\^{}x=8 হলে, x এর মান কত?}

ক) (2/3)

\textbf{\hl{খ) (3/2) ✔}}

গ) (4/3)

ঘ) (3/4)

\textbf{\hl{✔ সঠিক উত্তর: খ) (3/2)}}

\begin{quote}
\textbf{ব্যাখ্যা:} \emph{দেয়া আছে, 4\^{}x=8}

\emph{⇒ (2\^{}2)\^{}x=2\^{}3}

\emph{⇒ 2\^{}2x=2\^{}3}

\emph{⇒ 2x=3}

\emph{∴ x=(3/2)}
\end{quote}

\textbf{4\^{}(x+1)=2\^{}x−2 হলে, x এর মান কত?}

\textbf{\hl{ক) −4 ✔}}

খ) 5

গ) 6

ঘ) 8

\textbf{\hl{✔ সঠিক উত্তর: ক) −4}}

\begin{quote}
\textbf{ব্যাখ্যা:} \emph{4\^{}(x+1)=2\^{}x−2}

\emph{⇒ (2\^{}2)\^{}(x+1)=2\^{}x−2}

\emph{⇒ 2\^{}(2x+2)=2\^{}x−2}

\emph{⇒ 2x+2=x−2}

\emph{⇒ 2x−x=−2−2}

\emph{∴ x=−4}
\end{quote}

\textbf{3\^{}3x−8=3\^{}4 হলে x এর মান কত?}

ক) 3

\textbf{\hl{খ) 4 ✔}}

গ) 2

ঘ) 6

\textbf{\hl{✔ সঠিক উত্তর: খ) 4}}

\begin{quote}
\textbf{ব্যাখ্যা:} \emph{3\^{}3x−8=3\^{}4}

\emph{⇒ 3x−8=4}

\emph{⇒ 3x=4+8}

\emph{⇒ 3x=12}

\emph{∴ x=4}
\end{quote}

\textbf{2\^{}(x+7)=4\^{}(x+2) হলে x এর মান কত?}

ক) 0

খ) 2

\textbf{\hl{গ) 3 ✔}}

ঘ) −1

\textbf{\hl{✔ সঠিক উত্তর: গ) 3}}

\begin{quote}
\textbf{ব্যাখ্যা:} \emph{2\^{}(x+7)=4\^{}(x+2)}

\emph{⇒ 2\^{}(x+7)=(2\^{}2)\^{}(x+2)}

\emph{⇒ 2\^{}(x+7)=2\^{}(2x+4)}

\emph{⇒ x+7=2x+4}

\emph{⇒ 2x−x=7−4}

\emph{∴ x=3}
\end{quote}

\textbf{3.27\^{}x=9\^{}(x+4) হলে x এর মান কত?}

ক) 9

খ) 3

\textbf{\hl{গ) 7 ✔}}

ঘ) 1

\textbf{\hl{✔ সঠিক উত্তর: গ) 7}}

\begin{quote}
\textbf{ব্যাখ্যা:} \emph{3.27\^{}x=9\^{}(x+4)}

\emph{⇒ 3.(3\^{}3)\^{}x=(3\^{}2)\^{}(x+4)}

\emph{⇒ 3\^{}1.3\^{}3x=3\^{}(2x+8)}

\emph{⇒ 3\^{}(3x+1)=3\^{}(2x+8)}

\emph{⇒ 3x+1=2x+8}

\emph{∴ x=7}
\end{quote}

\textbf{12.27\^{}x=2\^{}2.9\^{}(x+4) হলে x এর মান কত?}

ক) 3

খ) 8

গ) 4

\textbf{\hl{ঘ) 7 ✔}}

\textbf{\hl{✔ সঠিক উত্তর: ঘ) 7}}

\begin{quote}
\textbf{ব্যাখ্যা:} \emph{12.27\^{}x=2\^{}2.9\^{}(x+4)}

\emph{⇒ 4×3×(3\^{}3)\^{}x=4×(3\^{}2)\^{}(x+4)}

\emph{⇒ 3×3\^{}3x=3\^{}(2x+8)}

\emph{⇒ 3\^{}(3x+1)=3\^{}(2x+8)}

\emph{⇒ 3x+1=2x+8}

\emph{∴ x=7}
\end{quote}

\textbf{5\^{}4x−2=7\^{}4x−2 হলে x = ?}

\textbf{\hl{ক) (1/2) ✔}}

খ) −(1/2)

গ) 2

ঘ) −2

\textbf{\hl{✔ সঠিক উত্তর: ক) (1/2)}}

\begin{quote}
\textbf{ব্যাখ্যা:} \emph{5\^{}4x−2=7\^{}4x−2}

\emph{উভয়পক্ষকে 7\^{}4x−2 দ্বারা ভাগ করে পাই,}

\emph{⇒ (((5/7)))\^{}4x−2=1}

\emph{⇒ (((5/7)))\^{}4x−2=(((5/7)))\^{}0}

\emph{⇒ 4x−2=0}

\emph{⇒ 4x=2}

\emph{∴ x=(2/4)=(1/2)}
\end{quote}

\textbf{If 2\^{}x−6=(1/64) Find the value of x.}

ক) 1

খ) 2

গ) 3

\textbf{\hl{ঘ) 0 ✔}}

\textbf{\hl{✔ সঠিক উত্তর: ঘ) 0}}

\begin{quote}
\textbf{ব্যাখ্যা:} \emph{দেওয়া আছে, 2\^{}x−6=(1/64)}

\emph{⇒ 2\^{}x−6=(1/2\^{}6)=2\^{}−6}

\emph{⇒ x−6=−6}

\emph{∴ x=−6+6=0}
\end{quote}

\textbf{9.2\^{}n−2.2\^{}n−1 = কত?}

\textbf{\hl{ক) 2\^{}(n+3) ✔}}

খ) 2\^{}n−3

গ) 2\^{}n

ঘ) 2\^{}−n

\textbf{\hl{✔ সঠিক উত্তর: ক) 2\^{}(n+3)}}

\begin{quote}
\textbf{ব্যাখ্যা:} \emph{9.2\^{}n−2.2\^{}n−1}

\emph{= 9.2\^{}n−2.(2\^{}n/2\^{}1)}

\emph{= 2\^{}n(9−1)}

\emph{= 2\^{}n.8}

\emph{= 2\^{}n.2\^{}3}

\emph{= 2\^{}(n+3)}
\end{quote}

\textbf{(10x\^{}2)\^{}0 = কত?}

ক) 0

\textbf{\hl{খ) 1 ✔}}

গ) 10

ঘ) 100

\textbf{\hl{✔ সঠিক উত্তর: খ) 1}}

\begin{quote}
\textbf{ব্যাখ্যা:} \emph{(10x\^{}2)\^{}0=(10)\^{}0.(x\^{}2)\^{}0=1 {[}
a\^{}0=1 ( a≠0 ){]}}
\end{quote}

\textbf{100\^{}x=10 হলে, x এর মান কত?}

\textbf{\hl{ক) (1/2) ✔}}

খ) (1/4)

গ) (1/3)

ঘ) 1

\textbf{\hl{✔ সঠিক উত্তর: ক) (1/2)}}

\begin{quote}
\textbf{ব্যাখ্যা:} \emph{100\^{}x=10}

\emph{⇒ (10\^{}2)\^{}x=10}

\emph{⇒ 10\^{}2x=10\^{}1}

\emph{⇒ 2x=1}

\emph{∴ x=(1/2)}
\end{quote}

\textbf{x\^{}x√(x)=(x√(x))\^{}x হলে, x এর মান কত?}

ক) (3/2)

খ) (4/9)

\textbf{\hl{গ) (9/4) ✔}}

ঘ) (2/3)

\textbf{\hl{✔ সঠিক উত্তর: গ) (9/4)}}

\begin{quote}
\textbf{ব্যাখ্যা:} \emph{x\^{}x√(x)=(x√(x))\^{}x}

\emph{⇒ x\^{}(x⋅x\^{}((1/2)))=((x⋅x\^{}((1/2))))\^{}x}

\emph{⇒ x\^{}(x\^{}((3/2)))=((x\^{}((3/2))))\^{}x}

\emph{⇒ x\^{}x√(x)=(x\^{}x)\^{}((3/2))}

\emph{⇒ x√(x)=(3/2)}

\emph{⇒ x\^{}((3/2))=(((3/2)))\^{}((3/2))}

\emph{∴ x=(9/4)}
\end{quote}

\textbf{4\^{}(x+1)=32 হলে x এর মান কত?}

ক) (2/3)

\textbf{\hl{খ) (3/2) ✔}}

গ) (1/2)

ঘ) (1/3)

\textbf{\hl{✔ সঠিক উত্তর: খ) (3/2)}}

\begin{quote}
\textbf{ব্যাখ্যা:} \emph{দেওয়া আছে, 4\^{}(x+1)=32}

\emph{⇒ (2\^{}2)\^{}(x+1)=2\^{}5}

\emph{⇒ 2\^{}(2x+2)=2\^{}5}

\emph{⇒ 2x+2=5}

\emph{⇒ 2x=3}

\emph{∴ x=(3/2)}
\end{quote}

\textbf{9\^{}(x+2)=81 হলে x = ?}

\textbf{\hl{ক) 0 ✔}}

খ) 1

গ) 2

ঘ) 3

\textbf{\hl{✔ সঠিক উত্তর: ক) 0}}

\begin{quote}
\textbf{ব্যাখ্যা:} \emph{9\^{}(x+2)=81}

\emph{⇒ (3\^{}2)\^{}(x+2)=3\^{}4}

\emph{⇒ 3\^{}(2x+4)=3\^{}4}

\emph{⇒ 2x+4=4}

\emph{⇒ 2x=4−4}

\emph{∴ x=0}
\end{quote}

\textbf{যদি 3\^{}(x+2)=243 হয় তবে 3\^{}x−2 এর মান কত?}

ক) 2

\textbf{\hl{খ) 3 ✔}}

গ) 0

ঘ) 1

\textbf{\hl{✔ সঠিক উত্তর: খ) 3}}

\begin{quote}
\textbf{ব্যাখ্যা:} \emph{3\^{}(x+2)=243}

\emph{⇒ 3\^{}x.3\^{}2=243}

\emph{⇒ 3\^{}x.9=243}

\emph{⇒ 3\^{}x=27}

\emph{⇒ 3\^{}x=3\^{}3}

\emph{∴ x=3}

\emph{এখন, 3\^{}x−2=3\^{}3−2=3\^{}1=3}
\end{quote}

\textbf{16\^{}((3/4))÷16\^{}((1/2)) = কত?}

ক) 1

খ) (1/2)

গ) (1/2)

\textbf{\hl{ঘ) 2 ✔}}

\textbf{\hl{✔ সঠিক উত্তর: ঘ) 2}}

\begin{quote}
\textbf{ব্যাখ্যা:} \emph{16\^{}((3/4))÷16\^{}((1/2))}

\emph{= 16\^{}((3/4)−(1/2))}

\emph{= 16\^{}((3−2/4))}

\emph{= 16\^{}((1/4))}

\emph{= (2\^{}4)\^{}((1/4)) = 2}
\end{quote}

\textbf{4\^{}(x+1)=32 হলে x এর মান কত?}

ক) (2/3)

খ) (1/3)

\textbf{\hl{গ) (3/2) ✔}}

ঘ) (1/2)

\textbf{\hl{✔ সঠিক উত্তর: গ) (3/2)}}

\begin{quote}
\textbf{ব্যাখ্যা:} \emph{4\^{}(x+1)=32}

\emph{⇒ (2\^{}2)\^{}(x+1)=2\^{}5}

\emph{⇒ 2\^{}(2(x+1))=2\^{}5}

\emph{⇒ 2x+2=5}

\emph{⇒ 2x=5−2}

\emph{⇒ 2x=3}

\emph{∴ x=(3/2)}
\end{quote}

\textbf{((2\^{}−1+5\^{}−1))\^{}−1 এর মান কত?}

\textbf{\hl{ক) (10/7) ✔}}

খ) (7/10)

গ) (3/7)

ঘ) (7/3)

\textbf{\hl{✔ সঠিক উত্তর: ক) (10/7)}}

\begin{quote}
\textbf{ব্যাখ্যা:} \emph{((2\^{}−1+5\^{}−1))\^{}−1}

\emph{= (((1/2)+(1/5)))\^{}−1}

\emph{= (((5+2/10)))\^{}−1}

\emph{= (((7/10)))\^{}−1}

\emph{= (10/7)}
\end{quote}

\textbf{যদি (25)\^{}(2x+3)=5\^{}(3x+6) হয়, তবে x -এর মান কত?}

\textbf{\hl{ক) 0 ✔}}

খ) 1

গ) 2

ঘ) 4

\textbf{\hl{✔ সঠিক উত্তর: ক) 0}}

\begin{quote}
\textbf{ব্যাখ্যা:} \emph{(25)\^{}(2x+3)=5\^{}(3x+6)}

\emph{⇒ (5\^{}2)\^{}(2x+3)=5\^{}(3x+6)}

\emph{⇒ 5\^{}(4x+6)=5\^{}(3x+6)}

\emph{⇒ 4x+6=3x+6}

\emph{⇒ 4x−3x=6−6}

\emph{∴ x=0}
\end{quote}

\textbf{(4\^{}x−1/2\^{}x−1) এর মান কোনটি?}

ক) 2\^{}x−1

\textbf{\hl{খ) 2\^{}x+1 ✔}}

গ) 2\^{}(x+1)

ঘ) 2\^{}x−1

\textbf{\hl{✔ সঠিক উত্তর: খ) 2\^{}x+1}}

\begin{quote}
\textbf{ব্যাখ্যা:} \emph{(4\^{}x−1/2\^{}x−1)}

\emph{= ((2\^{}x)\^{}2−1\^{}2/2\^{}x−1)}

\emph{= ((2\^{}x+1)(2\^{}x−1)/2\^{}x−1)}

\emph{= 2\^{}x+1}
\end{quote}

\textbf{(0.15×10\^{}p/0.3×10\^{}q)=5×10\^{}7 হলে p−q = কত?}

ক) 6

খ) 7

\textbf{\hl{গ) 8 ✔}}

ঘ) 9

\textbf{\hl{✔ সঠিক উত্তর: গ) 8}}

\begin{quote}
\textbf{ব্যাখ্যা:} \emph{(0.15×10\^{}p/0.3×10\^{}q)=5×10\^{}7}

\emph{⇒ 0.5×10\^{}p−q=5×10\^{}7}

\emph{⇒ 10\^{}p−q=(5×10\^{}7/0.5)}

\emph{⇒ 10\^{}p−q=10\^{}8}

\emph{∴ p−q=8}
\end{quote}

\textbf{3×27\^{}x=9\^{}(x+4) হলে x = কত?}

ক) 3

\textbf{\hl{খ) 7 ✔}}

গ) 5

ঘ) 9

\textbf{\hl{✔ সঠিক উত্তর: খ) 7}}

\begin{quote}
\textbf{ব্যাখ্যা:} \emph{3×27\^{}x=9\^{}(x+4)}

\emph{⇒ 3×(3\^{}3)\^{}x=(3\^{}2)\^{}(x+4)}

\emph{⇒ 3\^{}1×3\^{}3x=3\^{}(2x+8)}

\emph{⇒ 3\^{}(3x+1)=3\^{}(2x+8)}

\emph{⇒ 3x+1=2x+8}

\emph{∴ x=7}
\end{quote}

\textbf{x\^{}4÷x\^{}6×x\^{}2 এর মান কত?}

\textbf{\hl{ক) 1 ✔}}

খ) x\^{}3

গ) x

ঘ) 0

\textbf{\hl{✔ সঠিক উত্তর: ক) 1}}

\begin{quote}
\textbf{ব্যাখ্যা:} \emph{x\^{}4÷x\^{}6×x\^{}2}

\emph{= x\^{}4−6×x\^{}2}

\emph{= x\^{}−2×x\^{}2}

\emph{= x\^{}(−2+2)}

\emph{= x\^{}0=1}
\end{quote}

\textbf{(1000)\^{}((x/3))=10 হলে x এর মান কত?}

\textbf{\hl{ক) 1 ✔}}

খ) 2

গ) 3

ঘ) 10

\textbf{\hl{✔ সঠিক উত্তর: ক) 1}}

\begin{quote}
\textbf{ব্যাখ্যা:} \emph{(1000)\^{}((x/3))=10}

\emph{⇒ (10\^{}3)\^{}((x/3))=10\^{}1}

\emph{⇒ 10\^{}x=10\^{}1}

\emph{∴ x=1}
\end{quote}

\textbf{a\^{}m.a\^{}n=a\^{}(m+n) কখন হবে?/ a\^{}m.a\^{}n=a\^{}(m+n)
when---}

ক) m ধনাত্মক হলে (m is positive)

খ) n ধনাত্মক হলে (n is positive)

\textbf{\hl{গ) m ও n ধনাত্মক হলে (m and n are positive) ✔}}

ঘ) m ধনাত্মক ও n ঋণাত্মক হলে (m is positive and n is negative)

\textbf{\hl{✔ সঠিক উত্তর: গ) m ও n ধনাত্মক হলে (m and n are positive)}}

\begin{quote}
\textbf{ব্যাখ্যা:} \emph{সূচকের গুণের নিয়ম অনুসারে a\^{}m.a\^{}n=a\^{}(m+n)
হওয়ার জন্য ভিত্তি a বাস্তব সংখ্যা এবং সূচকদ্বয় m ও n কে ধনাত্মক পূর্ণসংখ্যা হতে
হয়।}
\end{quote}

\textbf{(3√(3))\^{}2x−1=(√(3))\^{}(x+1) হলে x এর মান কত?}

\textbf{\hl{ক) 5 ✔}}

খ) 3

গ) 2

ঘ) √(3)

\textbf{\hl{✔ সঠিক উত্তর: ক) 5}}

\begin{quote}
\textbf{ব্যাখ্যা:} \emph{দেওয়া আছে, (3√(3))\^{}2x−1=(√(3))\^{}(x+1)}

\emph{⇒ ((3\^{}((1/3))))\^{}2x−1=((3\^{}((1/2))))\^{}(x+1)}

\emph{⇒ 3\^{}((2x−1/3))=3\^{}((x+1/2))}

\emph{⇒ (2x−1/3)=(x+1/2)}

\emph{⇒ 4x−2=3x+3}

\emph{∴ x=5}
\end{quote}

\textbf{(5\^{}(n+2)+35×5\^{}n−1/4×5\^{}n) এর মান কত?}

ক) 5

খ) 5

\textbf{\hl{গ) 8 ✔}}

ঘ) 7

\textbf{\hl{✔ সঠিক উত্তর: গ) 8}}

\begin{quote}
\textbf{ব্যাখ্যা:} \emph{(5\^{}(n+2)+35×5\^{}n−1/4×5\^{}n)}

\emph{= (5\^{}n.5\^{}2+35.5\^{}n.(1/5)/4×5\^{}n)}

\emph{= (25×5\^{}n+7×5\^{}n/4×5\^{}n)}

\emph{= (32×5\^{}n/4×5\^{}n)=8}
\end{quote}

\textbf{(3\^{}(x+4)−9.3\^{}(x+1)/3\^{}(x+2)) এর মান -}

\textbf{\hl{ক) 6 ✔}}

খ) 0

গ) 27

ঘ) (1/6)

\textbf{\hl{✔ সঠিক উত্তর: ক) 6}}

\begin{quote}
\textbf{ব্যাখ্যা:} \emph{(3\^{}(x+4)−9.3\^{}(x+1)/3\^{}(x+2))}

\emph{= (3\^{}(x+2+2)−9.3\^{}(x+2−1)/3\^{}(x+2))}

\emph{= (3\^{}(x+2)(3\^{}2−9.3\^{}−1)/3\^{}(x+2))}

\emph{= 3\^{}2−9.(1/3)}

\emph{= 9−3=6}
\end{quote}

\textbf{(((x\^{}p/x\^{}q)))\^{}(p+q)+(((x\^{}q/x\^{}r)))\^{}(q+r)+(((x\^{}r/x\^{}p)))\^{}(r+p)
এর মান কত?}

ক) x

খ) x\^{}(p+q+r)

গ) x\^{}(p+q+r)

\textbf{\hl{ঘ) 1 ✔}}

\textbf{\hl{✔ সঠিক উত্তর: ঘ) 1}}

\begin{quote}
\textbf{ব্যাখ্যা:}
\emph{(((x\^{}p/x\^{}q)))\^{}(p+q)+(((x\^{}q/x\^{}r)))\^{}(q+r)+(((x\^{}r/x\^{}p)))\^{}(r+p)}

\emph{= x\^{}((p−q)(p+q))×x\^{}((q−r)(q+r))×x\^{}((r−p)(r+p))}

\emph{= x\^{}p\^{}2−q\^{}2×x\^{}q\^{}2−r\^{}2×x\^{}r\^{}2−p\^{}2}

\emph{= x\^{}(p\^{}2−q\^{}2+q\^{}2−r\^{}2+r\^{}2−p\^{}2)}

\emph{= x\^{}0=1}
\end{quote}

\textbf{mn√((x\^{}m/x\^{}n))×nl√((x\^{}n/x\^{}l))×lm√((x\^{}l/x\^{}m)) =
কত?}

ক) 0

\textbf{\hl{খ) 1 ✔}}

গ) x\^{}lmn

ঘ) (1/x\^{}lmn)

\textbf{\hl{✔ সঠিক উত্তর: খ) 1}}

\begin{quote}
\textbf{ব্যাখ্যা:} \emph{প্রদত্ত রাশিটিতে রুটের সূচক হিসেবে mn, nl এবং lm
রয়েছে ধরে সমাধান করলে:}

\emph{=
(x\^{}m−n)\^{}((1/mn))×(x\^{}n−l)\^{}((1/nl))×(x\^{}l−m)\^{}((1/lm))}

\emph{= x\^{}((m−n/mn))×x\^{}((n−l/nl))×x\^{}((l−m/lm))}

\emph{= x\^{}((m−n/mn)+(n−l/nl)+(l−m/lm))}

\emph{= x\^{}((l(m−n)+m(n−l)+n(l−m)/mnl))}

\emph{= x\^{}((0/mnl))=x\^{}0=1}
\end{quote}

\textbf{30−\{5\^{}−1(2−3)\^{}−3\}\^{}−2 = কত?}

ক) 29

\textbf{\hl{খ) 5 ✔}}

গ) 20

ঘ) 31

\textbf{\hl{✔ সঠিক উত্তর: খ) 5}}

\begin{quote}
\textbf{ব্যাখ্যা:} \emph{30−\{5\^{}−1(2−3)\^{}−3\}\^{}−2}

\emph{= 30−(\{(1/5)×(−1)\^{}−3\})\^{}−2}

\emph{= 30−(\{(1/5)×(−1)\})\^{}−2}

\emph{= 30−((−(1/5)))\^{}−2}

\emph{= 30−(−5)\^{}2}

\emph{= 30−25=5}
\end{quote}

\textbf{p\^{}m×p\^{}n×p\^{}−r = কত?}

ক) p\^{}−mnr

খ) p\^{}(m+n+r)

\textbf{\hl{গ) p\^{}(m+n−r) ✔}}

ঘ) p\^{}m−n−r

\textbf{\hl{✔ সঠিক উত্তর: গ) p\^{}(m+n−r)}}

\begin{quote}
\textbf{ব্যাখ্যা:} \emph{p\^{}m×p\^{}n×p\^{}−r}

\emph{= p\^{}(m+n−r)}
\end{quote}

\textbf{√(x\^{}−1y)×√(y\^{}−1z)×√(z\^{}−1x) এর সরলমান কত?}

ক) 0

খ) 1

গ) xyz

\textbf{\hl{ঘ) ±1 ✔}}

\textbf{\hl{✔ সঠিক উত্তর: ঘ) ±1}}

\begin{quote}
\textbf{ব্যাখ্যা:} \emph{√(x\^{}−1y)×√(y\^{}−1z)×√(z\^{}−1x)}

\emph{= √((y/x))×√((z/y))×√((x/z))}

\emph{= √((y/x)×(z/y)×(x/z)) = √(1)=±1}
\end{quote}

\textbf{m\^{}((x−y)(x+y))×m\^{}((y−z)(y+z))×m\^{}((z−x)(z+x)) = কত?}

ক) 0

\textbf{\hl{খ) 1 ✔}}

গ) (1/2)

ঘ) 2

\textbf{\hl{✔ সঠিক উত্তর: খ) 1}}

\begin{quote}
\textbf{ব্যাখ্যা:}
\emph{m\^{}((x−y)(x+y))×m\^{}((y−z)(y+z))×m\^{}((z−x)(z+x))}

\emph{= m\^{}x\^{}2−y\^{}2×m\^{}y\^{}2−z\^{}2×m\^{}z\^{}2−x\^{}2}

\emph{= m\^{}(x\^{}2−y\^{}2+y\^{}2−z\^{}2+z\^{}2−x\^{}2)}

\emph{= m\^{}0=1}
\end{quote}

\textbf{Find the value of n, if 27\^{}(n−(1/3))=243}

\textbf{\hl{ক) 2 ✔}}

খ) 3

গ) 4

ঘ) 5

\textbf{\hl{✔ সঠিক উত্তর: ক) 2}}

\begin{quote}
\textbf{ব্যাখ্যা:} \emph{27\^{}(n−(1/3))=243}

\emph{⇒ (27)\^{}((3n−1/3))=5\^{}5}

\emph{⇒ (3\^{}3)\^{}((3n−1/3))=3\^{}5}

\emph{⇒ 3\^{}3n−1=3\^{}5}

\emph{⇒ 3n−1=5}

\emph{⇒ 3n=6}

\emph{∴ n=2}
\end{quote}

\textbf{(2\^{}2)\^{}(x+3)=256 হলে, x = কত?}

ক) 0

\textbf{\hl{খ) 1 ✔}}

গ) 2

ঘ) 4

\textbf{\hl{✔ সঠিক উত্তর: খ) 1}}

\begin{quote}
\textbf{ব্যাখ্যা:} \emph{(2\^{}2)\^{}(x+3)=256}

\emph{⇒ 2\^{}(2x+6)=2\^{}8}

\emph{⇒ 2x+6=8}

\emph{⇒ 2x=2}

\emph{∴ x=1}
\end{quote}

\textbf{x এর মান কত হলে, 2\^{}4x−12=16 হবে?}

\textbf{\hl{ক) 4 ✔}}

খ) 2

গ) 6

ঘ) 8

\textbf{\hl{✔ সঠিক উত্তর: ক) 4}}

\begin{quote}
\textbf{ব্যাখ্যা:} \emph{দেওয়া আছে, 2\^{}4x−12=16}

\emph{⇒ 2\^{}4x−12=2\^{}4}

\emph{⇒ 4x−12=4}

\emph{⇒ 4x=16}

\emph{∴ x=4}
\end{quote}

\textbf{3\^{}mx−1=3a\^{}mx−2 হলে, x এর মান কত?}

\textbf{\hl{ক) (2/m) ✔}}

খ) (m/2)

গ) 2m

ঘ) কোনোটিই নয়

\textbf{\hl{✔ সঠিক উত্তর: ক) (2/m)}}

\begin{quote}
\textbf{ব্যাখ্যা:} \emph{দেওয়া আছে, 3\^{}mx−1=3a\^{}mx−2}

\emph{⇒ (3\^{}mx−1/3)=a\^{}mx−2}

\emph{⇒ 3\^{}mx−2=a\^{}mx−2}

\emph{⇒ (((3/a)))\^{}mx−2=1}

\emph{⇒ (((3/a)))\^{}mx−2=(((3/a)))\^{}0}

\emph{⇒ mx−2=0}

\emph{⇒ mx=2}

\emph{∴ x=(2/m)}
\end{quote}

\textbf{(3.2\^{}n−4.2\^{}n−2/2\^{}n−2\^{}(n+1)) = কত?}

ক) 2\^{}(n+1)

খ) 2\^{}n−1

\textbf{\hl{গ) −2 ✔}}

ঘ) 4

\textbf{\hl{✔ সঠিক উত্তর: গ) −2}}

\begin{quote}
\textbf{ব্যাখ্যা:} \emph{(3.2\^{}n−4.2\^{}n−2/2\^{}n−2\^{}(n+1))}

\emph{= (3.2\^{}n−2\^{}2.2\^{}n−2/2\^{}n−2\^{}(n+1))}

\emph{= (3.2\^{}n−2\^{}(2+n−2)/2\^{}n−2\^{}(n+1))}

\emph{= (3.2\^{}n−2\^{}n/2\^{}n−2\^{}n.2)}

\emph{= (2\^{}n(3−1)/2\^{}n(1−2))}

\emph{= (2/−1)=−2}
\end{quote}

\textbf{(25)\^{}(2x+3)=5\^{}(3x+6) হয় তবে x = কত?}

\textbf{\hl{ক) 0 ✔}}

খ) 1

গ) −1

ঘ) 4

\textbf{\hl{✔ সঠিক উত্তর: ক) 0}}

\begin{quote}
\textbf{ব্যাখ্যা:} \emph{(25)\^{}(2x+3)=5\^{}(3x+6)}

\emph{⇒ (5\^{}2)\^{}(2x+3)=5\^{}(3x+6)}

\emph{⇒ 5\^{}(4x+6)=5\^{}(3x+6)}

\emph{⇒ 4x+6=3x+6}

\emph{⇒ 4x−3x=6−6}

\emph{∴ x=0}
\end{quote}

\textbf{যদি (125)\^{}(2x+3)=5\^{}(3x+6) হয়, তবে x = কত?}

ক) 0

খ) 1

\textbf{\hl{গ) −1 ✔}}

ঘ) 4

\textbf{\hl{✔ সঠিক উত্তর: গ) −1}}

\begin{quote}
\textbf{ব্যাখ্যা:} \emph{(125)\^{}(2x+3)=5\^{}(3x+6)}

\emph{⇒ (5\^{}3)\^{}(2x+3)=5\^{}(3x+6)}

\emph{⇒ 5\^{}(3(2x+3))=5\^{}(3x+6)}

\emph{⇒ 3(2x+3)=3x+6}

\emph{⇒ 6x+9=3x+6}

\emph{⇒ 6x−3x=6−9}

\emph{⇒ 3x=−3}

\emph{∴ x=−1}
\end{quote}

\textbf{x\^{}4=9 হলে এবং x ধনাত্মক হলে x = ?}

ক) 3

খ) 1.5

\textbf{\hl{গ) √(3) ✔}}

ঘ) 4

\textbf{\hl{✔ সঠিক উত্তর: গ) √(3)}}

\begin{quote}
\textbf{ব্যাখ্যা:} \emph{x\^{}4=9}

\emph{⇒ x\^{}4=3×3}

\emph{⇒ x\^{}4=(√(3))\^{}2×(√(3))\^{}2}

\emph{⇒ x\^{}4=(√(3))\^{}4}

\emph{∴ x=√(3)}
\end{quote}

\textbf{12.27\^{}x=2\^{}2.9\^{}(x+4) হলে, x এর মান কত?}

ক) 3

খ) 8

\textbf{\hl{গ) 7 ✔}}

ঘ) 4

\textbf{\hl{✔ সঠিক উত্তর: গ) 7}}

\begin{quote}
\textbf{ব্যাখ্যা:} \emph{12.27\^{}x=2\^{}2.9\^{}(x+4)}

\emph{⇒ 12.3\^{}3x=4.9\^{}x.9\^{}4}

\emph{⇒ (3\^{}3x/3\^{}2x)=(4×9\^{}4/12)=(3\^{}8/3)}

\emph{⇒ 3\^{}x=3\^{}7}

\emph{∴ x=7}
\end{quote}

\textbf{2\^{}(x+7)=4\^{}(x+2) সমীকরণের সমাধান কত?}

ক) 2

\textbf{\hl{খ) 3 ✔}}

গ) 4

ঘ) 1

\textbf{\hl{✔ সঠিক উত্তর: খ) 3}}

\begin{quote}
\textbf{ব্যাখ্যা:} \emph{2\^{}(x+7)=4\^{}(x+2)}

\emph{⇒ 2\^{}(x+7)=(2\^{}2)\^{}(x+2)}

\emph{⇒ 2\^{}(x+7)=2\^{}(2x+4)}

\emph{⇒ x+7=2x+4}

\emph{⇒ 2x−x=7−4}

\emph{∴ x=3}
\end{quote}

\textbf{27\^{}(x+1)=81 হলে x এর মান নিচের কোনটি?}

\textbf{\hl{ক) (1/3) ✔}}

খ) (7/3)

গ) 2

ঘ) 3

\textbf{\hl{✔ সঠিক উত্তর: ক) (1/3)}}

\begin{quote}
\textbf{ব্যাখ্যা:} \emph{দেওয়া আছে, 27\^{}(x+1)=81}

\emph{⇒ (3\^{}3)\^{}(x+1)=3\^{}4}

\emph{⇒ 3\^{}(3x+3)=3\^{}4}

\emph{⇒ 3x+3=4}

\emph{⇒ 3x=4−3}

\emph{⇒ 3x=1}

\emph{∴ x=(1/3)}
\end{quote}

\textbf{x এর মান কত হলে, 72.3\^{}3x−5=2\^{}3 হবে?}

ক) (3/5)

খ) 2

গ) (5/3)

\textbf{\hl{ঘ) 1 ✔}}

\textbf{\hl{✔ সঠিক উত্তর: ঘ) 1}}

\begin{quote}
\textbf{ব্যাখ্যা:} \emph{দেওয়া আছে, 72.3\^{}3x−5=2\^{}3}

\emph{⇒ 8×9.3\^{}3x−5=8}

\emph{⇒ 8.3\^{}2.3\^{}3x−5=8}

\emph{⇒ 3\^{}(3x−5+2)=1}

\emph{⇒ 3\^{}3x−3=3\^{}0}

\emph{⇒ 3x−3=0}

\emph{⇒ 3x=3}

\emph{∴ x=1}
\end{quote}

\textbf{x\^{}y=y\^{}x,x=2y(x≠0,y≠0) হলে (x,y) = কত?}

ক) (8, 4)

খ) (6, 3)

গ) (2, 1)

\textbf{\hl{ঘ) (4, 2) ✔}}

\textbf{\hl{✔ সঠিক উত্তর: ঘ) (4, 2)}}

\begin{quote}
\textbf{ব্যাখ্যা:} \emph{দেওয়া আছে, x\^{}y=y\^{}x এবং x=2y}

\emph{⇒ (2y)\^{}y=y\^{}2y}

\emph{⇒ 2\^{}y.y\^{}y=(y\^{}2)\^{}y}

\emph{⇒ 2\^{}y=y\^{}2y−y=y\^{}y}

\emph{⇒ y=2}

\emph{এখন, x=2y=2×2=4}

\emph{∴ (x,y)=(4,2)}
\end{quote}

\textbf{{[}2−3(2−3)\^{}−1{]}\^{}−1 এর মান কত?}

ক) 5

খ) -5

গ) 1

\textbf{\hl{ঘ) (1/5) ✔}}

\textbf{\hl{✔ সঠিক উত্তর: ঘ) (1/5)}}

\begin{quote}
\textbf{ব্যাখ্যা:} \emph{{[}2−3(2−3)\^{}−1{]}\^{}−1}

\emph{= {[}2−3(−1)\^{}−1{]}\^{}−1}

\emph{= {[}2−3(−1){]}\^{}−1}

\emph{= ({[}2+3{]})\^{}−1}

\emph{= 5\^{}−1=(1/5)}
\end{quote}

\textbf{(√(3))\^{}6 এর মান কত?}

ক) 9

খ) 18

\textbf{\hl{গ) 27 ✔}}

ঘ) 81

\textbf{\hl{✔ সঠিক উত্তর: গ) 27}}

\begin{quote}
\textbf{ব্যাখ্যা:} \emph{(√(3))\^{}6}

\emph{= ((3\^{}((1/2))))\^{}6}

\emph{= 3\^{}((6/2)) = 3\^{}3=27}
\end{quote}

\textbf{(((125/27)))\^{}(−(2/3)) এর সহজ প্রকাশ---}

ক) (3/25)

খ) (5/20)

\textbf{\hl{গ) (9/25) ✔}}

ঘ) (3/20)

\textbf{\hl{✔ সঠিক উত্তর: গ) (9/25)}}

\begin{quote}
\textbf{ব্যাখ্যা:} \emph{(((125/27)))\^{}(−(2/3))}

\emph{= (((27/125)))\^{}((2/3))}

\emph{= (\{(((3/5)))\^{}3\})\^{}((2/3))}

\emph{= (((3/5)))\^{}2=(9/25)}
\end{quote}

\textbf{0.0000472 সংখ্যাটিকে বৈজ্ঞানিক পদ্ধতিতে প্রকাশ করুন/ Express the
number 0.0000472 in scientific system---}

\textbf{\hl{ক) 4.72×10\^{}−5 ✔}}

খ) 4.72×10\^{}−6

গ) 4.72×10\^{}−4

ঘ) 4.72×10\^{}−7

\textbf{\hl{✔ সঠিক উত্তর: ক) 4.72×10\^{}−5}}

\begin{quote}
\textbf{ব্যাখ্যা:}
\emph{0.0000472=(4.72/100000)=(4.72/10\^{}5)=4.72×10\^{}−5}
\end{quote}

\textbf{√(10√(10√(10√(10))))\ldots\ldots\ldots\ldots∞ এর মান কত?}

\textbf{\hl{ক) 0 বা 10 ✔}}

খ) 0 বা 9

গ) 0 বা 8

ঘ) কোনোটিই নয়

\textbf{\hl{✔ সঠিক উত্তর: ক) 0 বা 10}}

\begin{quote}
\textbf{ব্যাখ্যা:} \emph{ধরি,
a=√(10√(10√(10√(10))))\ldots\ldots\ldots\ldots∞}

\emph{⇒ a\^{}2=10√(10√(10√(10)))\ldots\ldots\ldots\ldots∞ {[}উভয়পক্ষকে
বর্গ করে{]}}

\emph{⇒ a\^{}2=10a}

\emph{⇒ a\^{}2−10a=0}

\emph{⇒ a(a−10)=0}

\emph{∴ a=0 অথবা 10}
\end{quote}

\textbf{{[}(\{1−((p−1/p))\})\^{}−1÷((1−(1/p)))\^{}−1{]} = কত?}

ক) 1

খ) −1

গ) (1/p)

\textbf{\hl{ঘ) p−1 ✔}}

\textbf{\hl{✔ সঠিক উত্তর: ঘ) p−1}}

\begin{quote}
\textbf{ব্যাখ্যা:} \emph{{[}(\{1−((p−1/p))\})\^{}−1÷((1−(1/p)))\^{}−1{]}}

\emph{⇒ {[}(\{(p−p+1/p)\})\^{}−1÷(((p−1/p)))\^{}−1{]}}

\emph{⇒ {[}(((1/p)))\^{}−1÷(((p−1/p)))\^{}−1{]}}

\emph{⇒ {[}p÷(p/p−1){]}}

\emph{⇒ {[}p×(p−1/p){]}}

\emph{⇒ p−1}

\emph{∴ p−1}
\end{quote}

\section{32. বীজগাণিতিক রাশিমালার ল.সা.গু ও
গ.সা.গু}\label{ux9acux99cux997ux9a3ux9a4ux995-ux9b0ux9b6ux9aeux9b2ux9b0-ux9b2.ux9b8.ux997-ux993-ux997.ux9b8.ux997}

বীজ গণিত হলো সংখ্যাগতিক বিষয়বস্তুগুলোর উপর ভিত্তি করে গণিতের একটি শাখা। এটি
মূলত গণিতের সাধারণ নিয়ম এবং সূত্র ব্যবহার করে সমস্যার সমাধান করে।

ল.সা.গু এবং গ.সা.গু হলো বীজ গণিতের দুটি গুরুত্বপূর্ণ অংশ। ল.সা.গু হলো সমস্যা
সমাধানের একটি পদ্ধতি, যেখানে সমাধান নির্ণয়ের জন্য প্রথমে সংখ্যা বা রাশিগুলোকে
একত্রে নেওয়া হয় এবং তাদের সর্বোচ্চ সংখ্যা বা রাশিটিকে লসাগু হিসেবে নেয়া হয়।

গ.সা.গু হলো সমাধানের আরেকটি পদ্ধতি, যেখানে সংখ্যা বা রাশিগুলোকে একত্রে নেওয়া
হয় এবং তাদের সর্বনিম্ন সংখ্যা বা রাশিটিকে গ.সা.গু. হিসেবে নেয়া হয়।পাটিগণিতের
ল.সা.গু ও গ.সা.গু জানা থাকলে এটা ও অত্যন্ত সহজ । কারণ একই নিয়মেই বীজগণিতে ও
কাজ করতে হবে।

পাটিগণিতে আমরা সংখ্যা নিয়ে কাজ করেছি আর বীজগণিতে আমাদের রাশিমালা নিয়ে কাজ
করতে হবে পার্থক্য এইটুকুই ।তবে মনে রাখা জরুরী বীজগণিতের ল.সা.গু ও গ.সা.গু
সম্পর্কে জানার আগে উৎপাদকে বিশ্লেষণ জানা থাকতে হবে । কারণ এই অংক শুরুই করব
উৎপাদকে বিশ্লেষণের মধ্য দিয়ে ।

\begin{itemize}
\item
  বীজগণিতের ল.সা.গু. নির্ণয়ের সহজ কেীশল
\end{itemize}

ল.সা.গু . এর উপর বিগত সালের বিভিন্ন প্রশ্ন ও সমাধান

\begin{itemize}
\item
  বীজগণিতের গ.সা.গু. নির্ণয়ের সহজ কেীশল
\end{itemize}

গ.সা.গু . এর উপর বিগত সালের বিভিন্ন প্রশ্ন ও সমাধান

\begin{itemize}
\item
  বীজগণিতের ল.সা.গু এবং গ.সা.গু এর সারমর্ম
\end{itemize}

\subsubsection{বীজগণিতের ল.সা.গু. নির্ণয়ের সহজ
কেীশল}\label{ux9acux99cux997ux9a3ux9a4ux9b0-ux9b2.ux9b8.ux997.-ux9a8ux9b0ux9a3ux9afux9b0-ux9b8ux9b9ux99c-ux995ux9b6ux9b2}

\textbf{ল.সা.গু. (L. C. M . - Least Common Multiple):} দুইবা ততোধিক রাশির
মধ্যে সাধারণ গুণিতকের সবচেয়ে ছোট গুণিতকটিকে লঘিষ্ট সাধারণ গুণিতক বা ল.সা.গু
বলে। যেমন- \(x^{3}y^{3}z^{4}\) , \(x^{4}y^{2}z^{2}\) এবং
\(x^{3}y^{4}z^{3}\) এ রাশি তিনটির ল.সা.গু \(x^{4}y^{4}z^{4}\) .

\textbf{ল.সা.গু. নির্ণয়ের নিয়ম:}

পাটিগণিতের নিয়মে রাশিগুলোর মধ্যে যে সংখ্যাগুলো থাকবে তার ল.সা.গু. নির্ণয় করতে
হবে।

বীজগণিতীয় রাশিগুলোর সাধারণ উৎপাদকের সর্বোচ্চ শক্তি বা পাওয়ার বের করতে হবে ।

এখন, উভয়ের গুণফলই হবে প্রদত্ত রাশিগুলোর ল.সা.গু ।

\paragraph{ল.সা.গু . এর উপর বিগত সালের বিভিন্ন প্রশ্ন ও
সমাধান}\label{ux9b2.ux9b8.ux997-.-ux98fux9b0-ux989ux9aaux9b0-ux9acux997ux9a4-ux9b8ux9b2ux9b0-ux9acux9adux9a8ux9a8-ux9aaux9b0ux9b6ux9a8-ux993-ux9b8ux9aeux9a7ux9a8}

\textbf{প্রশ্ন নম্বর-} ১. \(a^{2}bc\) , \(ab^{2}c\) এবং \(abc^{2}\) এর
ল.সা.গু কত? {[} উপজেলা পরিসংখ্যান কর্মকর্তা: 10 {]}

\begin{itemize}
\item
  (ক) a\^{}2b\^{}2c\^{}2
\item
  (খ) a\^{}3b\^{}2c
\item
  (গ) a\^{}2bc\^{}2
\item
  (ঘ) abc\^{}2
\end{itemize}

\textbf{\hl{উত্তরঃ (ক)}} \(a^{2}b^{2}c^{2}\)

\textbf{সমাধান:}

প্রদত্ত রাশি ৩ টির ল.সা.গু হবে- প্রতিটি রাশির মধ্যে সর্বোচ্চ পাওয়ার বিশিষ্ট
সংখ্যা গুলোর গুণফল ।

নির্ণেয় ল.সা.গু = \(a^{2}b^{2}c^{2}\)

\textbf{প্রশ্ন নম্বর-} ২. (a+b) এবং \(a^{2}\)(a-b) . -এর ল.সা.গু কত? {[}
পাব:সার্ভিস কমিশনে সহ:পরি: 94 {]}

\begin{itemize}
\item
  (ক) a\^{}2(a\^{}2-b\^{}2)
\item
  (খ) a\^{}3(a\^{}2-b\^{}2)
\item
  (গ) a\^{}4(a\^{}2-b\^{}2)
\item
  (ঘ) a\^{}2(a\^{}3-b\^{}2)
\end{itemize}

\textbf{\hl{উত্তরঃ (ক)}} \(a^{2}(a^{2} - b^{2})\)

\textbf{সমাধান:}

প্রদত্ত রাশি 2 টির ল.সা.গু হবে- প্রতিটি রাশির মধ্যে সর্বোচ্চ পাওয়ার বিশিষ্ট
সংখ্যা গুলোর গুণফল ।

নির্ণেয় ল.সা.গু = \(a^{2}(a^{2} - b^{2})\)

\textbf{প্রশ্ন নম্বর-} ৩. \({2x}^{2} + x\) এবং \({4x}^{2} - 1\) এর ল.সা.গু
কত? {[} থানা শিক্ষা অফিসার-99 {]}

\begin{itemize}
\item
  (ক) x(2x+2) (2x-1)
\item
  (খ) x(x+1) (2x-1)
\item
  (গ) x(2x+1) (2x-2)
\item
  (ঘ) x ( 2x + 1) (2x-1)
\end{itemize}

\textbf{\hl{উত্তরঃ (ঘ) x ( 2x + 1) (2x-1)}}

\textbf{সমাধান:}

প্রথম রাশি,

\({2x}^{2} + x\)

= x(2x + 1)

দ্বিতীয় রাশি,

\({4x}^{2} - 1\)

= \({(2x)}^{2} - 1^{2}\)

= (2x + 1) (2x - 1)

ল.সা.গু = x(2x + 1) (2x - 1)

\textbf{প্রশ্ন নম্বর-} ৪. \(x^{3}y - {xy}^{3}\) এবং \(x^{2} - y^{2}\) -এর
ল.সা.গু কত? {[} তথ্য মন্ত্র: অধীন সহ:পরি:-03 {]}

\begin{itemize}
\item
  (ক) x(x+y) (x-y)
\item
  (খ) xy(x - y) (y - x)
\item
  (গ) y (y - x) (x+y)
\item
  (ঘ) xy(x - y) (x+y)
\end{itemize}

\textbf{\hl{উত্তরঃ (ঘ) xy(x - y) (x+y)}}

\textbf{সমাধান:}

প্রথম রাশি,

\(x^{3}y - {xy}^{3}\)

= \(xy(x^{2} - y^{2})\)

= xy(x + y) (x - y)

দ্বিতীয় রাশি,

= \(x^{2} - y^{2}\)

= (x+y) (x-y)

ল.সা.গু. = xy(x - y) (x + y)

\textbf{প্রশ্ন নম্বর-} ৫. \(x^{3} + x^{2}y\) , \(x^{2}y + xy^{2}\) এর ল.
সা. গু. কোনটি? {[}32তম বিসিএস (বিশেষ){]}

\begin{itemize}
\item
  (ক) xy(x+y)
\item
  (খ) x\^{}2 (x+y)
\item
  (গ) x\^{}2 y(x+y)
\item
  (ঘ) y(x+y)
\end{itemize}

\textbf{\hl{উত্তরঃ (গ)}} \(x^{2}\)\textbf{\hl{y(x+y)}}

\textbf{সমাধান:}

প্রথম রাশি,

\(x^{3} + x^{2}y\)

= \(x^{2}(\ x\  + \ y\ )\)

দ্বিতীয় রাশি ,

\(x^{2}y + xy^{2}\)

= xy(x + y )

ল.সা.গু. = \(x^{2}y(\ x\  + \ y\ )\)

\textbf{প্রশ্ন নম্বর-} ৬. 3\(x^{2}\)+9, \(x^{4}\)-9 এবং \(x^{4}\) +
6\(x^{2}\) + 9 -এর ল.সা.গু নির্ণয় করুন । {[} স্বরাষ্ট্র মন্ত্র: পুলিশ সহ::02 {]}

\begin{itemize}
\item
  (ক) ( x\^{}2 +3) (x²-3) ( x\^{}2 +3)
\item
  (খ) 3( x\^{}2 +3) (x-3) ( x\^{}2 +3)
\item
  (গ) 3( x\^{}2 +3) ( x\^{}2 -3) ( x\^{}2 +4)
\item
  (ঘ) 3 ((x\^{}2+3))\^{}2 ( x\^{}2 -3)
\end{itemize}

\textbf{\hl{উত্তরঃ (ঘ) 3}}\({(x^{2} + 3)}^{2}\) \textbf{\hl{(}}
\(x^{2}\)\textbf{\hl{-3)}}

\textbf{সমাধান:}

প্রথম রাশি =

3\(x^{2}\)+9

= 3 ( \(x^{2}\)+3 )

দ্বিতীয় রাশি

= \(x^{4}\)-9

= \({(x^{2})}^{2} - 3^{2}\)

= ( \(x^{2}\)+3) ( \(x^{2}\)-3)

তৃতীয় রাশি =

\(x^{4} + {6x}^{2} + 9\)

= \({(x^{2})}^{2} + 2.x^{2}.3 + 3^{2}\)

= \({(x^{2} + 3)}^{2}\)

= ( \(x^{2}\)+3) ( \(x^{2}\)+3)

ল.সা.গু = 3\({(x^{2} + 3)}^{2}\) ( \(x^{2}\)-3)

\textbf{প্রশ্ন নম্বর-} ৭. \(x^{3}\) - 1, \(x^{3}\) + 1, \(x^{4}\) +
\(x^{2}\) + 1 এর ল. সা. গু. কত? {[}প্রাথ: শিক্ষক -১৫{]}

\begin{itemize}
\item
  (ক) x\^{}5-1
\item
  (খ) x\^{}4-1
\item
  (গ) x\^{}9-1
\item
  (ঘ) x\^{}6-1
\end{itemize}

\textbf{\hl{উত্তরঃ (ঘ)}} \(x^{6} - 1\)

\textbf{সমাধান:}

১ম রাশি, \(x^{3}\) -- 1 = (x -- 1) ( \(x^{2}\) + x + 1)

২য় রাশি, \(x^{3}\) + 1 = (x + 1) ( \(x^{2}\) - x + 1)

৩য় রাশি = \(x^{4} + x^{2} + 1\)

= \({(x^{2})}^{2} + {2x}^{2} + 1 - x^{2}\)

= \({(x^{2} + 1)}^{2} - x^{2}\)

= ( \(x^{2}\) + x + 1) ( \(x^{2}\) - x + 1)

ল. সা. গু . = (x + 1) ( \(x^{2}\) - x + 1) (x - 1) ( \(x^{2}\) + x + 1)

= ( \(x^{3}\) + 1 ) ( \(x^{3}\) - 1 )

= \({(x^{3})}^{2} - 1^{2}\)

= \(x^{6} - 1\)

\textbf{প্রশ্ন নম্বর-} ৮. ল. সা. গু. নির্ণয় করুন- \(a^{3}\) - 1, \(a^{3}\) +
1, \(a^{4}\) + \(a^{2}\) + 1 {[}৮ম বেসরকারী শিক্ষক নিবন্ধন ও প্রত্যয়ন পরীক্ষা
2012{]}

\begin{itemize}
\item
  (ক) a\^{}3-1
\item
  (খ) a\^{}2-1
\item
  (গ) a - 1
\item
  (ঘ) a\^{}6-1
\end{itemize}

\textbf{\hl{উত্তরঃ (ঘ)}} \(a^{6} - 1\)

\textbf{সমাধান:}

১ম রাশি, \(a^{3}\) -- 1 = (a -- 1) ( \(a^{2}\) + a + 1)

২য় রাশি, \(a^{3}\) + 1 = (a + 1) ( \(a^{2}\) - a + 1)

৩য় রাশি = \(a^{4} + a^{2} + 1\)

= \({(a^{2})}^{2} + {2a}^{2} + 1 - a^{2}\)

= \({(a^{2} + 1)}^{2} - a^{2}\)

= ( \(a^{2}\) + a + 1) ( \(a^{2}\) - a + 1)

ল. সা. গু . = (a + 1) ( \(a^{2}\) - a + 1) (a - 1) ( \(a^{2}\) + a + 1)

= ( \(a^{3}\) + 1 ) ( \(a^{3}\) - 1 )

= \({(a^{3})}^{2} - 1^{2}\)

= \(a^{6} - 1\)

\textbf{প্রশ্ন নম্বর-} ৯. \(a^{2}b(a^{3} - b^{3})\) ,
\(a^{2}b^{2}(a^{4} + a^{2}b^{2} + b^{4})\) এবং ( \(a^{3} + b^{3}\) ) এর
ল.সা.গু. নির্ণয় করুনঃ

\begin{itemize}
\item
  (ক) a\^{}3b(a\^{}6-b\^{}6)
\item
  (খ) a\^{}3b\^{}3(a\^{}6-b\^{}6)
\item
  (গ) a\^{}2b\^{}2(a\^{}6-b\^{}6)
\item
  (ঘ) ab(a\^{}5-b\^{}5)
\end{itemize}

\textbf{\hl{উত্তরঃ (গ)}} \(a^{2}b^{2}(a^{6} - b^{6})\)

\textbf{সমাধান:}

১ম রাশি = \(a^{2}b(a^{3} - b^{3})\)

= \(a^{2}b(a - b)(a^{2} + ab + b^{2})\)

২য় রাশি = \(a^{2}b^{2}(a^{4} + a^{2}b^{2} + b^{4})\)

=
\(a^{2}b^{2}\{{(a^{2})}^{2} + 2.a^{2}.b^{2} + {(b^{2})}^{2} - a^{2}b^{2}\}\)

= \(a^{2}b^{2}{\{(a^{2} + b^{2})}^{2} - a^{2}b^{2}\}\)

= \(a^{2}b^{2}(a^{2} + ab + b^{2})(a^{2} - ab + b^{2})\)

৩য় রাশি = \(a^{3} + b^{3}\)

= \((a + b)(a^{2} - ab + b^{2})\)

নির্ণেয় ল.সা.গু. =
\(a^{2}b^{2}(\ ‍a\  - \ b\ )(a^{2} + ab + b^{2})(\ ‍a\  + \ b\ )(a^{2} - ab + b^{2})\)

= \(a^{2}b^{2}(a^{3} - b^{3})(a^{3} + b^{3})\)

= \(a^{2}b^{2}(a^{6} - b^{6})\)

\textbf{প্রশ্ন নম্বর-} ১০. \({4x}^{2}y^{2}z\) এবং \(6xy^{2}\) এর ল. সা. গু.
কত? {[} ডাক ও টেলিযোগাযোগ মন্ত্রণালয়ের টেলিফোন পরিচালক 2004{]}

\begin{itemize}
\item
  (ক) 12x\^{}2 yz
\item
  (খ) 12xy\^{}2z
\item
  (গ) 10x\^{}2y\^{}2z
\item
  (ঘ) 12x\^{}2y\^{}2z
\end{itemize}

\textbf{\hl{উত্তরঃ (ঘ)}} \({12x}^{2}y^{2}z\)

\textbf{সমাধান:}

এখানে 4 ও 6 এর ল.সা.গু. = 12

নির্ণেয় ল.সা.গু. = \({12x}^{2}y^{2}z\)

\subsubsection{বীজগণিতের গ.সা.গু. নির্ণয়ের সহজ
কেীশল}\label{ux9acux99cux997ux9a3ux9a4ux9b0-ux997.ux9b8.ux997.-ux9a8ux9b0ux9a3ux9afux9b0-ux9b8ux9b9ux99c-ux995ux9b6ux9b2}

\textbf{গ.সা.গু. (H.C.F. - Highest Common Factor):}

দুই বা ততোধিক রাশির অন্তর্গত সর্বোচ্চ সংখ্যক সাধারণ মৌলিক গুণনীয়কের ধারাবাহিক
গুণফলকে ঐ রাশিগুলোর গরিষ্ঠ সাধারণ গুণনীয়ক বা গ.সা.গু বলা হয়। যেমন-
\(x^{2}y^{3}z^{4}\) , \(x^{4}y^{2}z^{2}\) এবং \(x^{3}y^{4}z^{3}\) এ রাশি
তিনটির গ.সা.গু. \(x^{2}y^{2}z^{2}\) অর্থাৎ গ.সা.গু এর ক্ষেত্রে সর্বনিম্ন ঘাত বা
পাওয়ার আর ল.সা.গু. এর ক্ষেত্রে সর্বোচ্চ ঘাত বা পাওয়ার নিতে হয়।

\textbf{গ.সা.গু নির্ণয়ের নিয়ম :}

পাটিগণিতের নিয়মে রাশিগুলোর মধ্যে যে সংখ্যাগুলো থাকবে তার গ.সা.গু নির্ণয় করতে
হবে ।

বীজগণিতীয় রাশিগুলোর মৌলিক উৎপাদক বের করতে হবে।

প্রদত্ত রাশিগুলোর সর্বোচ্চ সংখ্যক বীজগণিতীয় সাধারণ মৌলিক উৎপাদকগুলোর এবং
সাংখ্য-সহগের গ.সা.গু এর

ধারাবাহিক গুণফল হচ্ছে নির্ণেয় গ.সা.গু ।

\paragraph{গ.সা.গু . এর উপর বিগত সালের বিভিন্ন প্রশ্ন ও
সমাধান}\label{ux997.ux9b8.ux997-.-ux98fux9b0-ux989ux9aaux9b0-ux9acux997ux9a4-ux9b8ux9b2ux9b0-ux9acux9adux9a8ux9a8-ux9aaux9b0ux9b6ux9a8-ux993-ux9b8ux9aeux9a7ux9a8}

\textbf{প্রশ্ন নম্বর-} ১১. \({4x}^{2} - 16\) এবং \({6x}^{2} + 24x + 24\) এর
গ.সা.গু.-\/-\/-\/-\/- {[} 31তম বিসিএস {]}

\begin{itemize}
\item
  (ক) 4(x+2)
\item
  (খ) 2(x+3)
\item
  (গ) 2(x+2)
\item
  (ঘ) 2 (x+5)
\end{itemize}

\textbf{\hl{উত্তরঃ (গ) 2(x+2)}}

\textbf{সমাধান:}

প্রথম রাশি, \({4x}^{2} - 16\)

= \(4(x^{2} - 4)\)

= \(4(x^{2} - 2^{2})\)

= 4 (x + 2) (x - 2)

দ্বিতীয় রাশি, \({6x}^{2} + 24x + 24\)

= 2(\({3x}^{2}\) + 12x + 12)

⇒ 2( \({3x}^{2}\) + 6x + 6x + 12)

⇒ 2\{3x (x + 2) + 6 (x + 2)\}

= 2 (x + 2) (3x + 6)

নির্ণেয় গসাগু, = 2(x + 2)

\textbf{প্রশ্ন নম্বর-} ১২. \(x^{2}\)-11x+30 এবং \(x^{3}\)- \({4x}^{2}\)
-2x-15 এর গ. সা. গু কত? {[}25তম বিসিএস{]}

\begin{itemize}
\item
  (ক) x-4
\item
  (খ) x-2
\item
  (গ) 2x-5
\item
  (ঘ) x-5
\end{itemize}

\textbf{\hl{উত্তরঃ (ঘ) x-5}}

\textbf{সমাধান:}

১ম রাশি = \(x^{2}\) - 11x + 30

= \(x^{2}\) - 5x - 6x + 30

= x(x-5) -6(x-5)

= (x-5) (x-6)

দ্বিতীয় রাশি = \(x^{3}\)- \({4x}^{2}\) -2x-15

= \(x^{3} - {5x}^{2} + x^{2}\) -5x+3x-15

= \(x^{2}\) (x-5)+x (x-5)+3(x-5)

= ( \(x^{2}\) +x+3) (x-5)

.:. নির্ণেয় গ.সা.গু = x - 5

সকল রাশির মধ্যে রয়েছে এমন উৎপাদকগুলোই হচ্ছে নির্ণেয় গ.সা.গু.

\textbf{প্রশ্ন নম্বর-} ১৩. \(a^{2} - 3a\) , \(a^{2} - 9\) , \(a^{2}\) - 4a
+ 3 এর গ.সা.গু. হবে- {[}শ্রম মন্ত্রণালয়ের সহকারী পরিচালক 2005{]}

\begin{itemize}
\item
  (ক) (a-3)
\item
  (খ) (a-2)
\item
  (গ) (a+2)
\item
  (ঘ) (a+3)
\end{itemize}

\textbf{\hl{উত্তরঃ (ক) (a-3)}}

\textbf{সমাধান:}

১ম রাশি = \(a^{2} - 3a\)

= a( a - 3 )

২য় রাশি = \(a^{2} - 9\)

= \(a^{2} - 3^{2}\)

= ( a + 3 ) ( a - 3 )

৩য় রাশি = \(a^{2} - 4a\  + \ 3\)

= \(a^{2} - 3a\  - \ ‍a\  + \ 3\)

= a ( a - 3) - 1 ( a - 3 )

= ( a - 3 ) ( a - 1 )

নির্ণেয় গ.সা.গু. = ( a - 3 )

সকল রাশির মধ্যে রয়েছে এমন উৎপাদকগুলোই হচ্ছে নির্ণেয় গ.সা.গু.

\textbf{প্রশ্ন নম্বর-} ১৪. a(a + b), ab(a + b), এবং a\({(a + b)}^{2}\) এর
গ. সা. গু কত? {[}বিআরটিএ (মোটরযান পরিদর্শক)-২০১৭{]}

\begin{itemize}
\item
  (ক) a(a + b)
\item
  (খ) ab(a + b)
\item
  (গ) a+b
\item
  (ঘ) a ((a+b))\^{}2
\end{itemize}

\textbf{\hl{উত্তরঃ (ক) a(a + b)}}

\textbf{সমাধান:}

১ম রাশি = a(a+b)

২য় রাশি = ab(a+b)

৩য় রাশি = a\({(a + b)}^{2}\)

= a(a + b) (a + b)

.:. নির্ণেয় গ.সা.গু. = a(a + b)

\textbf{প্রশ্ন নম্বর-} ১৫. দুটি সংখ্যার ল.সা.গু \(a^{2}\)b(a + b) এবং গ.সা.গু
a(a + b) । একটি সংখ্যা \(a^{3}\) + \(a^{2}\)b হলে, অপর সংখ্যাটি কত? {[}
সহ:থানা পরি: পরি: অফি: ৯৮ {]}

\begin{itemize}
\item
  (ক) a\^{}2b+ab\^{}2
\item
  (খ) ab+ab\^{}2
\item
  (গ) a\^{}2b+ab
\item
  (ঘ) a\^{}2b+b\^{}2
\end{itemize}

\textbf{\hl{উত্তরঃ}} \(a^{2}b + {ab}^{2}\)

\textbf{সমাধান:}

অপর সংখ্যা = \(\frac{\text{গ}.\text{সা}.\text{গু}\  \times \ \text{ল}.\text{সা}.\text{গু}}{\text{একটি}\ \text{সংখ্যা}}\)

= \(\frac{a2b(a\  + \ b)\  \times \ a(a\  + \ b)}{a2(a\  + \ b)}\)

= ab(a + b)

= \(a^{2}b + {ab}^{2}\)

\textbf{প্রশ্ন নম্বর-} ১৬. \(x^{2} - y^{2}\) , \({(x + y)}^{2}\) ,
\(x^{3} + y^{3}\) এর গ.সা.গু কত? {[} ১১ তম শিক্ষক নিবন্ধন পরীক্ষা -২০১৪{]}

\begin{itemize}
\item
  (ক) x - y
\item
  (খ) x + y
\item
  (গ) (x\^{}2-y\^{}2)(x\^{}3+y\^{}3)
\item
  (ঘ) (x-y)((x+y))\^{}2(x\^{}2-xy+y\^{}2)
\end{itemize}

\textbf{\hl{উত্তরঃ (খ) x + y}}

\textbf{সমাধান:}

১ম রাশি =

\(x^{2} - y^{2}\)

= ( x + y ) ( x - y )

২য় রাশি = \({(x + y)}^{2}\)

= ( x + y ) ( x + y )

৩য় রাশি = \(x^{3} + y^{3}\)

= \((x\  + \ y)(x^{2} - xy + y^{2})\)

নির্ণেয় গ.সা.গু. = x + y

\paragraph{বীজগণিতের ল.সা.গু এবং গ.সা.গু এর
সারমর্ম:}\label{ux9acux99cux997ux9a3ux9a4ux9b0-ux9b2.ux9b8.ux997-ux98fux9ac-ux997.ux9b8.ux997-ux98fux9b0-ux9b8ux9b0ux9aeux9b0ux9ae}

১. উৎপাদকের অংকের মতই প্রশ্নে প্রদত্ত রাশিগুলোকে ভাঙ্গতে হবে ।

২. ল.সা.গু এর উত্তরটি বড় এবং গ.সা.গু এর উত্তরটি ছোট হবে।

৩. ল.সা.গু এর ক্ষেত্রে কমন আনকমন সবগুলো রাশি একবার করে নিতে হবে, কোনটা দুবার
নেয়া যাবে না ।

৪. গ.সা.গু এর ক্ষেত্রে শুধু মিল রাশিগুলো নিতে হবে, এবং মিল না থাকলে উত্তর ১
হবে।

৫. রাশি গুলোর সাথে সাধারণ সংখ্যা আসলে তাদের ল.সা.গু গ.সা.গু ভিন্নভাবে করে
উত্তরে যোগ করে দিতে হবে।

\section{33. বীজগাণিতিক রাশিমালার ল.সা.গু ও গ.সা.গু MCQ প্রশ্ন উত্তর
অনুশীলন}\label{ux9acux99cux997ux9a3ux9a4ux995-ux9b0ux9b6ux9aeux9b2ux9b0-ux9b2.ux9b8.ux997-ux993-ux997.ux9b8.ux997-mcq-ux9aaux9b0ux9b6ux9a8-ux989ux9a4ux9a4ux9b0-ux985ux9a8ux9b6ux9b2ux9a8}

\textbf{a, a\^{}2, (a + b) এর লঘিষ্ঠ সাধারণ গুণিতক কোনটি?}

\textbf{\hl{ক) a\^{}2(a+b) ✔}}

খ) a

গ) a\^{}2

ঘ) a(a+b)

\textbf{\hl{✔ সঠিক উত্তর: ক) a\^{}2(a+b)}}

\begin{quote}
\textbf{ব্যাখ্যা:} \emph{১ম রাশি = a}

\emph{২য় রাশি = a\^{}2 = a × a}

\emph{৩য় রাশি = (a + b)}

\emph{∴ ল.সা.গু. = a\^{}2(a + b)}
\end{quote}

\textbf{a\^{}2bc, a\^{}2c এবং ab\^{}2 এর ল.সা.গু কত?}

\textbf{\hl{ক) a\^{}2b\^{}2c\^{}2 ✔}}

খ) a\^{}3b\^{}2c

গ) a\^{}2bc\^{}2

ঘ) ab\^{}2

\textbf{\hl{✔ সঠিক উত্তর: ক) a\^{}2b\^{}2c\^{}2}}

\begin{quote}
\textbf{ব্যাখ্যা:} \emph{প্রদত্ত রাশি ৩টির ল.সা.গু হবে- প্রতিটি রাশির মধ্যে
সর্বোচ্চ পাওয়ার বিশিষ্ট সংখ্যাগুলোর গুণফল।}

\emph{নির্ণেয় ল.সা.গু = a\^{}2b\^{}2c\^{}2}
\end{quote}

\textbf{(a+b) এবং a\^{}2(a−b) এর ল.সা.গু কত?}

\textbf{\hl{ক) a\^{}2(a+b)(a−b) ✔}}

খ) a\^{}3(a+b)(a−b)

গ) a\^{}4(a+b)(a−b)

ঘ) a\^{}2(a+b)\^{}3

\textbf{\hl{✔ সঠিক উত্তর: ক) a\^{}2(a+b)(a−b)}}

\begin{quote}
\textbf{ব্যাখ্যা:} \emph{প্রদত্ত রাশি ২টির ল.সা.গু হবে- প্রতিটি রাশির মধ্যে
সর্বোচ্চ পাওয়ার বিশিষ্ট সংখ্যাগুলোর গুণফল।}

\emph{নির্ণেয় ল.সা.গু = a\^{}2(a+b)(a−b)}
\end{quote}

\textbf{2x\^{}2 + x এবং 4x\^{}2 − 1 এর ল.সা.গু কত?}

ক) x(2x+2)(2x−1)

খ) x(x+1)(2x−1)

গ) x(2x+1)(2x−2)

\textbf{\hl{ঘ) x(2x+1)(2x−1) ✔}}

\textbf{\hl{✔ সঠিক উত্তর: ঘ) x(2x+1)(2x−1)}}

\begin{quote}
\textbf{ব্যাখ্যা:} \emph{প্রথম রাশি,}

\emph{2x\^{}2 + x}

\emph{= x(2x + 1)}

\emph{দ্বিতীয় রাশি,}

\emph{4x\^{}2 − 1}

\emph{= (2x)\^{}2 − 1\^{}2}

\emph{= (2x + 1) (2x − 1)}

\emph{∴ ল.সা.গু = x(2x + 1) (2x − 1)}
\end{quote}

\textbf{x\^{}3y − xy\^{}3 এবং x\^{}2 − y\^{}2 এর ল.সা.গু কত?}

ক) x(x+y)(x−y)

খ) xy(x − y)(y − x)

গ) y(y − x)(x+y)

\textbf{\hl{ঘ) xy(x − y)(x+y) ✔}}

\textbf{\hl{✔ সঠিক উত্তর: ঘ) xy(x − y)(x+y)}}

\begin{quote}
\textbf{ব্যাখ্যা:} \emph{প্রথম রাশি,}

\emph{x\^{}3y − xy\^{}3}

\emph{= xy(x\^{}2 − y\^{}2)}

\emph{= xy(x + y) (x − y){[}cite: 1{]}}

\emph{দ্বিতীয় রাশি,}

\emph{= x\^{}2 − y\^{}2}

\emph{= (x+y) (x−y){[}cite: 1{]}}

\emph{∴ ল.সা.গু. = xy(x − y) (x + y)}
\end{quote}

\textbf{x\^{}3y − xy\^{}3 এবং x\^{}2 − y\^{}2 এর ল.সা.গু কত?}

ক) x(x+y)(x−y)

খ) xy(x − y)(y − x)

গ) y(y − x)(x+y)

\textbf{\hl{ঘ) xy(x − y)(x+y) ✔}}

\textbf{\hl{✔ সঠিক উত্তর: ঘ) xy(x − y)(x+y)}}

\begin{quote}
\textbf{ব্যাখ্যা:} \emph{প্রথম রাশি, x\^{}3y − xy\^{}3 = xy(x\^{}2 −
y\^{}2) = xy(x + y)(x − y)}

\emph{দ্বিতীয় রাশি, = x\^{}2 − y\^{}2 = (x + y)(x − y)}

\emph{∴ ল.সা.গু. = xy(x − y)(x + y)}
\end{quote}

\textbf{x\^{}3 + x\^{}2y, x\^{}2y + xy\^{}2 এর ল. সা. গু. কোনটি?}

ক) xy(x+y)

খ) x\^{}2(x+y)

\textbf{\hl{গ) x\^{}2y(x+y) ✔}}

ঘ) y(x+y)

\textbf{\hl{✔ সঠিক উত্তর: গ) x\^{}2y(x+y)}}

\begin{quote}
\textbf{ব্যাখ্যা:} \emph{প্রথম রাশি, x\^{}3 + x\^{}2y = x\^{}2(x + y)}

\emph{দ্বিতীয় রাশি, x\^{}2y + xy\^{}2 = xy(x + y)}

\emph{∴ ল.সা.গু. = x\^{}2y(x + y)}
\end{quote}

\textbf{3x\^{}2 + 9, x\^{}4 − 9 এবং x\^{}4 + 6x\^{}2 + 9 এর ল.সা.গু নির্ণয়
করুন।}

ক) (x\^{}2+3)(x\^{}2−3)(x\^{}2+3)

খ) 3(x\^{}2+3)(x−3)(x\^{}2+3)

গ) 3(x\^{}2+3)(x\^{}2−3)(x\^{}2+4)

\textbf{\hl{ঘ) 3(x\^{}2+3)\^{}2(x\^{}2−3) ✔}}

\textbf{\hl{✔ সঠিক উত্তর: ঘ) 3(x\^{}2+3)\^{}2(x\^{}2−3)}}

\begin{quote}
\textbf{ব্যাখ্যা:} \emph{প্রথম রাশি = 3x\^{}2 + 9}

\emph{= 3(x\^{}2 + 3)}

\emph{দ্বিতীয় রাশি = x\^{}4 − 9}

\emph{= (x\^{}2)\^{}2 − 3\^{}2}

\emph{= (x\^{}2 + 3)(x\^{}2 − 3)}

\emph{তৃতীয় রাশি = x\^{}4 + 6x\^{}2 + 9}

\emph{= (x\^{}2)\^{}2 + 2 × x\^{}2 × 3 + 3\^{}2}

\emph{= (x\^{}2 + 3)\^{}2}

\emph{= (x\^{}2 + 3)(x\^{}2 + 3)}

\emph{∴ ল.সা.গু = 3(x\^{}2 + 3)\^{}2(x\^{}2 − 3)}
\end{quote}

\textbf{x\^{}3 − 1, x\^{}3 + 1, x\^{}4 + x\^{}2 + 1 এর ল. সা. গু. কত?}

ক) x\^{}5 − 1

খ) x\^{}4 − 1

গ) x\^{}9 − 1

\textbf{\hl{ঘ) x\^{}6 − 1 ✔}}

\textbf{\hl{✔ সঠিক উত্তর: ঘ) x\^{}6 − 1}}

\begin{quote}
\textbf{ব্যাখ্যা:} \emph{১ম রাশি, x\^{}3 − 1}

\emph{= (x − 1)(x\^{}2 + x + 1)}

\emph{২য় রাশি, x\^{}3 + 1}

\emph{= (x + 1)(x\^{}2 − x + 1)}

\emph{৩য় রাশি = x\^{}4 + x\^{}2 + 1}

\emph{= (x\^{}2)\^{}2 + 2x\^{}2 + 1 − x\^{}2}

\emph{= (x\^{}2 + 1)\^{}2 − x\^{}2}

\emph{= (x\^{}2 + x + 1)(x\^{}2 − x + 1)}

\emph{∴ ল. সা. গু. = (x + 1)(x\^{}2 − x + 1)(x − 1)(x\^{}2 + x + 1)}

\emph{= (x\^{}3 + 1)(x\^{}3 − 1)}

\emph{= (x\^{}3)\^{}2 − 1\^{}2}

\emph{= x\^{}6 − 1}
\end{quote}

\textbf{ল. সা. গু. নির্ণয় করুন- a\^{}3 − 1, a\^{}3 + 1, a\^{}4 + a\^{}2 +
1}

ক) a\^{}3 − 1

খ) a\^{}2 − 1

গ) a − 1

\textbf{\hl{ঘ) a\^{}6 − 1 ✔}}

\textbf{\hl{✔ সঠিক উত্তর: ঘ) a\^{}6 − 1}}

\begin{quote}
\textbf{ব্যাখ্যা:} \emph{১ম রাশি, a\^{}3 − 1 = (a − 1)(a\^{}2 + a + 1)}

\emph{২য় রাশি, a\^{}3 + 1 = (a + 1)(a\^{}2 − a + 1)}

\emph{৩য় রাশি = a\^{}4 + a\^{}2 + 1 = (a\^{}2)\^{}2 + 2a\^{}2 + 1 −
a\^{}2 = (a\^{}2 + 1)\^{}2 − a\^{}2 = (a\^{}2 + a + 1)(a\^{}2 − a + 1)}

\emph{∴ ল. সা. গু. = (a + 1)(a\^{}2 − a + 1)(a − 1)(a\^{}2 + a + 1) =
(a\^{}3 + 1)(a\^{}3 − 1) = (a\^{}3)\^{}2 − 1\^{}2 = a\^{}6 − 1}
\end{quote}

\textbf{ল. সা. গু. নির্ণয় করুন:}

a\^{}3 − 1, 1 + a\^{}3, 1 + a\^{}2 + a\^{}4

\textbf{\hl{ক) a\^{}6 − 1 ✔}}

খ) (a − 1)(a\^{}3+1)

গ) (a\^{}4+1)(a − 1)

ঘ) a\^{}6+1

\textbf{\hl{✔ সঠিক উত্তর: ক) a\^{}6 − 1}}

\begin{quote}
\textbf{ব্যাখ্যা:} \emph{১ম রাশি = a\^{}3 − 1 = (a − 1)(a\^{}2+a+1)}

\emph{২য় রাশি = 1 + a\^{}3 = (1+a)(1−a+a\^{}2) = (a+1)(a\^{}2−a+1)}

\emph{৩য় রাশি = 1+a\^{}2+a\^{}4 = a\^{}4+2a\^{}2+1−a\^{}2 =
(a\^{}2+1)\^{}2−a\^{}2 = (a\^{}2+1+a)(a\^{}2+1−a) =
(a\^{}2+a+1)(a\^{}2−a+1)}

\emph{∴ ল.সা.গু. = (a−1)(a+1)(a\^{}2+a+1)(a\^{}2−a+1) =
(a\^{}2−1)(a\^{}4+a\^{}2+1) = a\^{}6−1}
\end{quote}

\textbf{a − b, a\^{}2 − ab, a\^{}2 − b\^{}2 এর ল.সা.গু. নিচের কোনটি?}

ক) a(a − b)

খ) (a − b)

\textbf{\hl{গ) a(a\^{}2 − b\^{}2) ✔}}

ঘ) a\^{}2 − b\^{}2

\textbf{\hl{✔ সঠিক উত্তর: গ) a(a\^{}2 − b\^{}2)}}

\begin{quote}
\textbf{ব্যাখ্যা:} \emph{১ম রাশি = a − b}

\emph{২য় রাশি = a\^{}2 − ab = a(a − b)}

\emph{৩য় রাশি = a\^{}2 − b\^{}2 = (a+b)(a − b)}

\emph{∴ নির্ণেয় ল.সা.গু. = a(a+b)(a − b) = a(a\^{}2 − b\^{}2)}
\end{quote}

\textbf{a\^{}2bc, ab\^{}2c, abc\^{}2 এর ল.সা.গু. কত?}

ক) abc

\textbf{\hl{খ) a\^{}2b\^{}2c\^{}2 ✔}}

গ) 1

ঘ) a\^{}3b\^{}3c\^{}3

\textbf{\hl{✔ সঠিক উত্তর: খ) a\^{}2b\^{}2c\^{}2}}

\begin{quote}
\textbf{ব্যাখ্যা:} \emph{রাশি তিনটিতে a, b ও c গুণনীয়ক সমূহের সর্বোচ্চ ঘাত
যথাক্রমে a\^{}2, b\^{}2 ও c\^{}2।}

\emph{∴ নির্ণেয় ল.সা.গু. = a\^{}2b\^{}2c\^{}2}
\end{quote}

\textbf{a\^{}7b\^{}2c, ab\^{}3c\^{}2 এবং a\^{}7b\^{}4c\^{}3 এর ল.সা.গু
কত?}

ক) a\^{}5b\^{}4c\^{}3

খ) a\^{}5b\^{}3c\^{}3

গ) a\^{}7b\^{}3c

\textbf{\hl{ঘ) a\^{}7b\^{}4c\^{}3 ✔}}

\textbf{\hl{✔ সঠিক উত্তর: ঘ) a\^{}7b\^{}4c\^{}3}}

\begin{quote}
\textbf{ব্যাখ্যা:} \emph{প্রতিটি বর্ণের সর্বোচ্চ ঘাত হলো a\^{}7, b\^{}4 এবং
c\^{}3।}

\emph{∴ রাশিগুলোর ল.সা.গু. = a\^{}7b\^{}4c\^{}3}
\end{quote}

\textbf{6a\^{}2bc, 8ab\^{}2c, 6abc এর ল.সা.গু. কোনটি?}

\textbf{\hl{ক) 24a\^{}2b\^{}2c ✔}}

খ) 48abc

গ) 24abc

ঘ) 6abc

\textbf{\hl{✔ সঠিক উত্তর: ক) 24a\^{}2b\^{}2c}}

\begin{quote}
\textbf{ব্যাখ্যা:} \emph{6a\^{}2bc = 2 × 3 × a × a × b × c}

\emph{8ab\^{}2c = 2 × 2 × 2 × a × b × b × c}

\emph{6abc = 2 × 3 × a × b × c}

\emph{∴ ল.সা.গু. = (2 × 2 × 2 × 3) × a\^{}2 × b\^{}2 × c =
24a\^{}2b\^{}2c}
\end{quote}

\textbf{x\^{}2 − y\^{}2, x\^{}2 + xy + y\^{}2, x\^{}3 − y\^{}3 রাশিদ্বয়ের
ল. সা. গু. কত?}

\textbf{\hl{ক) (x\^{}2 − y\^{}2)(x\^{}2 + xy + y\^{}2) ✔}}

খ) x\^{}3 − y\^{}3

গ) x − y

ঘ) x + y

\textbf{\hl{✔ সঠিক উত্তর: ক) (x\^{}2 − y\^{}2)(x\^{}2 + xy + y\^{}2)}}

\begin{quote}
\textbf{ব্যাখ্যা:} \emph{১ম রাশি = x\^{}2 − y\^{}2 = (x + y)(x − y)}

\emph{২য় রাশি = x\^{}2 + xy + y\^{}2}

\emph{৩য় রাশি = x\^{}3 − y\^{}3 = (x − y)(x\^{}2 + xy + y\^{}2)}

\emph{∴ ল.সা.গু. = (x + y)(x − y)(x\^{}2 + xy + y\^{}2) = (x\^{}2 −
y\^{}2)(x\^{}2 + xy + y\^{}2)}
\end{quote}

\textbf{x\^{}2−3x+2 এবং x\^{}2−5x+6 এর ল.সা.গু. = কত?}

\textbf{\hl{ক) (x−1)(x−2)(x−3) ✔}}

খ) (x+1)(x−2)(x+3)

গ) (x−1)(x+2)(x−3)

ঘ) (x−2)

\textbf{\hl{✔ সঠিক উত্তর: ক) (x−1)(x−2)(x−3)}}

\begin{quote}
\textbf{ব্যাখ্যা:} \emph{১ম রাশি = x\^{}2−3x+2 = (x−2)(x−1)}

\emph{২য় রাশি = x\^{}2−5x+6 = (x−3)(x−2)}

\emph{∴ নির্ণেয় ল.সা.গু. = (x−1)(x−2)(x−3)}
\end{quote}

\textbf{x\^{}2 − 4, x\^{}2 + x + 4, x\^{}3 − 8 বীজগাণিতিক রাশির ল.সা.গু.
কত?}

ক) (x+2)\^{}2(x\^{}3−8)

খ) (x−2)\^{}2(x\^{}3−8)

গ) (x\^{}3−2)(x\^{}3−8)

\textbf{\hl{ঘ) (x\^{}2+2)(x\^{}3−8) ✔}}

\textbf{\hl{✔ সঠিক উত্তর: ঘ) (x\^{}2+2)(x\^{}3−8)}}

\begin{quote}
\textbf{ব্যাখ্যা:} \emph{১ম রাশি = x\^{}2−4 = (x+2)(x−2)}

\emph{২য় রাশি = x\^{}2+4x+4 = (x+2)\^{}2}

\emph{৩য় রাশি = x\^{}3−8 = (x−2)(x\^{}2+2x+4)}

\emph{সঠিক বিশ্লেষণে ল.সা.গু. হবে (x+2)\^{}2(x−2)(x\^{}2+2x+4) বা
(x+2)\^{}2(x\^{}3−8)}
\end{quote}

\textbf{(a\^{}3 − 1), (a\^{}3 + 1), (a\^{}2 − 3a + 2) এর ল. সা. গু
কোনটি?}

\textbf{\hl{ক) (a\^{}6 − 1)(a − 2) ✔}}

খ) a\^{}2 − 1

গ) (a − 1)(a − 2)

ঘ) a\^{}6 − 1

\textbf{\hl{✔ সঠিক উত্তর: ক) (a\^{}6 − 1)(a − 2)}}

\begin{quote}
\textbf{ব্যাখ্যা:} \emph{১ম রাশি = (a−1)(a\^{}2+a+1)}

\emph{২য় রাশি = (a+1)(a\^{}2−a+1)}

\emph{৩য় রাশি = (a−2)(a−1)}

\emph{∴ ল.সা.গু. = (a−1)(a+1)(a\^{}2+a+1)(a\^{}2−a+1)(a−2) =
(a\^{}3−1)(a\^{}3+1)(a−2) = (a\^{}6−1)(a−2)}
\end{quote}

\textbf{a\^{}2b(a − b), a\^{}2b\^{}2(a\^{}4 + a\^{}2b\^{}2 + b\^{}4) এবং
(a\^{}3 + b\^{}3) এর ল.সা.গু. নির্ণয় করুন:}

ক) a\^{}3b(a\^{}6 − b\^{}6)

খ) a\^{}3b\^{}3(a\^{}6 − b\^{}6)

\textbf{\hl{গ) a\^{}2b\^{}2(a\^{}6 − b\^{}6) ✔}}

ঘ) ab(a\^{}5 − b\^{}5)

\textbf{\hl{✔ সঠিক উত্তর: গ) a\^{}2b\^{}2(a\^{}6 − b\^{}6)}}

\begin{quote}
\textbf{ব্যাখ্যা:} \emph{১ম রাশি = a\^{}2b(a − b)(a\^{}2 + ab + b\^{}2)}

\emph{২য় রাশি = a\^{}2b\^{}2(a\^{}4 + a\^{}2b\^{}2 + b\^{}4) =
a\^{}2b\^{}2(a\^{}2 + ab + b\^{}2)(a\^{}2 − ab + b\^{}2)}

\emph{৩য় রাশি = a\^{}3 + b\^{}3 = (a + b)(a\^{}2 − ab + b\^{}2)}

\emph{∴ নির্ণেয় ল.সা.গু. = a\^{}2b\^{}2(a − b)(a\^{}2 + ab + b\^{}2)(a +
b)(a\^{}2 − ab + b\^{}2)}

\emph{= a\^{}2b\^{}2(a\^{}3 − b\^{}3)(a\^{}3 + b\^{}3)}

\emph{= a\^{}2b\^{}2(a\^{}6 − b\^{}6)}
\end{quote}

\textbf{4x\^{}2y\^{}2z এবং 6xy\^{}2 এর ল. সা. গু. কত?}

ক) 12x\^{}2yz

খ) 12xy\^{}2z

গ) 10x\^{}2y\^{}2z

\textbf{\hl{ঘ) 12x\^{}2y\^{}2z ✔}}

\textbf{\hl{✔ সঠিক উত্তর: ঘ) 12x\^{}2y\^{}2z}}

\begin{quote}
\textbf{ব্যাখ্যা:} \emph{এখানে 4 ও 6 এর ল.সা.গু. = 12}

\emph{∴ নির্ণেয় ল.সা.গু. = 12x\^{}2y\^{}2z}
\end{quote}

\textbf{(1−x\^{}2/1−x) এর লঘিষ্ঠ রূপ নিচের কোনটি?}

ক) 1

খ) x

গ) (1 − x)

\textbf{\hl{ঘ) (1 + x) ✔}}

\textbf{\hl{✔ সঠিক উত্তর: ঘ) (1 + x)}}

\begin{quote}
\textbf{ব্যাখ্যা:} \emph{(1−x\^{}2/1−x) = ((1+x)(1−x)/1−x) = 1 + x}
\end{quote}

\textbf{(x\^{}2+3x\^{}3/x+3x\^{}2) এর লঘিষ্ঠ রূপ নিচের কোনটি?}

ক) x\^{}2

\textbf{\hl{খ) x ✔}}

গ) 1

ঘ) 0

\textbf{\hl{✔ সঠিক উত্তর: খ) x}}

\begin{quote}
\textbf{ব্যাখ্যা:} \emph{(x\^{}2+3x\^{}3/x+3x\^{}2) =
(x\^{}2(1+3x)/x(1+3x)) = (x\^{}2/x) = x}
\end{quote}

\textbf{2a ও 3b এর গ.সা.গু. কত?}

\textbf{\hl{ক) 1 ✔}}

খ) 6

গ) a

ঘ) b

\textbf{\hl{✔ সঠিক উত্তর: ক) 1}}

\begin{quote}
\textbf{ব্যাখ্যা:} \emph{2a ও 3b এর মধ্যে কোনো সাধারণ উৎপাদক নেই। ∴ 2a ও 3b
এর গ.সা.গু. = 1}
\end{quote}

\textbf{a(a + b), a\^{}2(a − b) এর গ.সা.গু. কত?}

ক) a\^{}2

\textbf{\hl{খ) a ✔}}

গ) a\^{}2 − b\^{}2

ঘ) a\^{}2(a\^{}2 − b\^{}2)

\textbf{\hl{✔ সঠিক উত্তর: খ) a}}

\begin{quote}
\textbf{ব্যাখ্যা:} \emph{১ম রাশি = a(a + b)}

\emph{২য় রাশি = a\^{}2(a − b)}

\emph{∴ রাশি দুটির গ.সা.গু. = a}
\end{quote}

\textbf{a(a+b), ab(a+b), a(a+b)\^{}2 এর গ. সা. গু. কত?}

\textbf{\hl{ক) a(a+b) ✔}}

খ) ab(a+b)

গ) a+b

ঘ) a(a+b)\^{}2

\textbf{\hl{✔ সঠিক উত্তর: ক) a(a+b)}}

\begin{quote}
\textbf{ব্যাখ্যা:} \emph{১ম রাশি = a(a+b)}

\emph{২য় রাশি = ab(a+b)}

\emph{৩য় রাশি = a(a+b)\^{}2 = a(a+b)(a+b)}

\emph{∴ নির্ণেয় গ. সা. গু. = a(a+b)}
\end{quote}

\textbf{a\^{}3 − b\^{}3, a\^{}3 + b\^{}3 এর গ.সা.গু. কত?}

ক) a − b

\textbf{\hl{খ) 1 ✔}}

গ) a\^{}2 + ab + b\^{}2

ঘ) a + b

\textbf{\hl{✔ সঠিক উত্তর: খ) 1}}

\begin{quote}
\textbf{ব্যাখ্যা:} \emph{১ম রাশি = a\^{}3 − b\^{}3 = (a − b)(a\^{}2 + ab +
b\^{}2)}

\emph{২য় রাশি = a\^{}3 + b\^{}3 = (a + b)(a\^{}2 − ab + b\^{}2)}

\emph{এখানে উৎপাদকগুলোর মধ্যে কোনো সাধারণ উপাদান নেই। ∴ গ.সা.গু. 1}
\end{quote}

\textbf{(a + b), a\^{}2 − b\^{}2, a\^{}3 − b\^{}3 এর গ. সা. গু. কত?}

ক) 0

\textbf{\hl{খ) 1 ✔}}

গ) a − b

ঘ) a + b

\textbf{\hl{✔ সঠিক উত্তর: খ) 1}}

\begin{quote}
\textbf{ব্যাখ্যা:} \emph{১ম রাশি = a + b}

\emph{২য় রাশি = a\^{}2 − b\^{}2 = (a + b)(a − b)}

\emph{৩য় রাশি = a\^{}3 − b\^{}3 = (a − b)(a\^{}2 + ab + b\^{}2)}

\emph{তিনটি রাশির মধ্যে কোনো সাধারণ উৎপাদক বা উপাদান নেই। ∴ নির্ণেয় গ. সা.
গু. = 1}
\end{quote}

\textbf{x\^{}2 + 2x, x\^{}3 + 8 এবং x\^{}2 − 8 রাশি তিনটির গ.সা.গু. নিচের
কোনটি?}

\textbf{\hl{ক) x + 2 ✔}}

খ) x − 2

গ) x(x + 2)(x − 2)

ঘ) x\^{}2 + 4x + 4

\textbf{\hl{✔ সঠিক উত্তর: ক) x + 2}}

\begin{quote}
\textbf{ব্যাখ্যা:} \emph{১ম রাশি = x\^{}2 + 2x = x(x + 2)}

\emph{২য় রাশি = x\^{}3 + 8 = (x + 2)(x\^{}2 − 2x + 4)}

\emph{৩য় রাশি = x\^{}2 − 4 = (x + 2)(x − 2)}

\emph{∴ নির্ণেয় গ. সা. গু. = (x + 2)}
\end{quote}

\textbf{x\^{}2 − 3x, x\^{}2 − 9 এবং x\^{}2 − 4x + 3 বীজগাণিতিক রাশির
গ.সা.গু. কত হবে?}

ক) x − 4

খ) x + 3

\textbf{\hl{গ) x − 3 ✔}}

ঘ) x − 1

\textbf{\hl{✔ সঠিক উত্তর: গ) x − 3}}

\begin{quote}
\textbf{ব্যাখ্যা:} \emph{১ম রাশি = x\^{}2 − 3x = x(x − 3)}

\emph{২য় রাশি = x\^{}2 − 9 = (x + 3)(x − 3)}

\emph{৩য় রাশি = x\^{}2 − 4x + 3 = (x − 3)(x − 1)}

\emph{∴ গ.সা.গু. = (x − 3)}
\end{quote}

\textbf{x\^{}2 + 5x, x\^{}2 − 25, x\^{}2 + 7x + 10 এর গ.সা.গু কত?}

ক) x − 5

\textbf{\hl{খ) x + 5 ✔}}

গ) x(x + 5)

ঘ) x(x + 5)(x − 5)(x + 2)

\textbf{\hl{✔ সঠিক উত্তর: খ) x + 5}}

\begin{quote}
\textbf{ব্যাখ্যা:} \emph{১ম রাশি = x\^{}2 + 5x = x(x + 5)}

\emph{২য় রাশি = x\^{}2 − 25 = (x + 5)(x − 5)}

\emph{৩য় রাশি = x\^{}2 + 7x + 10 = (x + 5)(x + 2)}

\emph{∴ নির্ণেয় গ.সা.গু. = (x + 5)}
\end{quote}

\textbf{4(x + y), 10(x − y) এবং 12(x\^{}2 − y\^{}2) এর গ.সা.গু. কত?}

ক) x − y

খ) x + y

গ) 12(x\^{}2 − y\^{}2)

\textbf{\hl{ঘ) 2 ✔}}

\textbf{\hl{✔ সঠিক উত্তর: ঘ) 2}}

\begin{quote}
\textbf{ব্যাখ্যা:} \emph{১ম রাশি = 4(x + y) = 2 × 2(x + y)}

\emph{২য় রাশি = 10(x − y) = 2 × 5(x − y)}

\emph{৩য় রাশি = 12(x\^{}2 − y\^{}2) = 2 × 2 × 3(x + y)(x − y)}

\emph{∴ গ.সা.গু. = 2}
\end{quote}

\textbf{x\^{}2 − y\^{}2, x\^{}3 − y\^{}3, x\^{}4 + x\^{}2y\^{}2 + y\^{}4
রাশিগুলোর গ.সা.গু. কত?}

ক) x + y

খ) x − y

গ) 0

\textbf{\hl{ঘ) 1 ✔}}

\textbf{\hl{✔ সঠিক উত্তর: ঘ) 1}}

\begin{quote}
\textbf{ব্যাখ্যা:} \emph{১ম রাশি = x\^{}2 − y\^{}2 = (x + y)(x − y)}

\emph{২য় রাশি = x\^{}3 − y\^{}3 = (x − y)(x\^{}2 + xy + y\^{}2)}

\emph{৩য় রাশি = x\^{}4 + x\^{}2y\^{}2 + y\^{}4 = (x\^{}2 + y\^{}2)\^{}2
− (xy)\^{}2 = (x\^{}2 + xy + y\^{}2)(x\^{}2 − xy + y\^{}2)}

\emph{রাশিগুলোর কোনো সাধারণ উৎপাদক নেই, তাই গ.সা.গু. = 1}
\end{quote}

\textbf{9a\^{}3b\^{}2c\^{}2, 9a\^{}3b\^{}2c\^{}2, 12a\^{}2bc,
15ab\^{}3c\^{}3 এর গ.সা.গু. কোনটি?}

ক) 9abc

খ) 9a\^{}3b\^{}3c\^{}3

গ) 3a\^{}3b\^{}3c\^{}3

\textbf{\hl{ঘ) 3abc ✔}}

\textbf{\hl{✔ সঠিক উত্তর: ঘ) 3abc}}

\begin{quote}
\textbf{ব্যাখ্যা:} \emph{প্রদত্ত রাশিগুলোর সাংখ্যিক সহগগুলোর গ.সা.গু. হলো 3 এবং
প্রতিটি পদের সাধারণ সর্বনিম্ন ঘাত হলো a, b ও c।}

\emph{∴ নির্ণেয় গ.সা.গু. = 3abc}
\end{quote}

\textbf{x\^{}2 − y\^{}2, (x + y)\^{}2, x\^{}3 + y\^{}3 এর গ.সা.গু কত?}

ক) x − y

\textbf{\hl{খ) x + y ✔}}

গ) (x\^{}2 − y\^{}2), (x\^{}3 + y\^{}3)

ঘ) (x − y), (x + y)\^{}2 (x\^{}2 − xy + y\^{}2)

\textbf{\hl{✔ সঠিক উত্তর: খ) x + y}}

\begin{quote}
\textbf{ব্যাখ্যা:} \emph{১ম রাশি = x\^{}2 − y\^{}2 = (x + y)(x − y)}

\emph{২য় রাশি = (x + y)\^{}2 = (x + y)(x + y)}

\emph{৩য় রাশি = x\^{}3 + y\^{}3 = (x + y)(x\^{}2 − xy + y\^{}2)}

\emph{∴ নির্ণেয় গ.সা.গু. = (x + y)}
\end{quote}

\textbf{a\^{}2 − 3a, a\^{}2 − 9, a\^{}2 − 4a + 3 এর গ.সা.গু কত?}

ক) a(a − 3)

\textbf{\hl{খ) a − 3 ✔}}

গ) (a − 1)(a − 3)

ঘ) a(a − 1)(a − 3)

\textbf{\hl{✔ সঠিক উত্তর: খ) a − 3}}

\begin{quote}
\textbf{ব্যাখ্যা:} \emph{১ম রাশি = a\^{}2 − 3a = a(a − 3)}

\emph{২য় রাশি = a\^{}2 − 9 = (a + 3)(a − 3)}

\emph{৩য় রাশি = a\^{}2 − 4a + 3 = (a − 3)(a − 1)}

\emph{∴ নির্ণেয় গ.সা.গু = (a − 3)}
\end{quote}

\textbf{16 − 4x\^{}2 এবং 6x\^{}2 + 24x + 24 এর গ. সা. গু কত?}

ক) x + 2

খ) x + 4

গ) (x + 2)\^{}2

\textbf{\hl{ঘ) 2(x + 2) ✔}}

\textbf{\hl{✔ সঠিক উত্তর: ঘ) 2(x + 2)}}

\begin{quote}
\textbf{ব্যাখ্যা:} \emph{১ম রাশি = 16 − 4x\^{}2 = 4(4 − x\^{}2) = 4(2 +
x)(2 − x)}

\emph{২য় রাশি = 6x\^{}2 + 24x + 24 = 6(x\^{}2 + 4x + 4) = 2 × 3(x +
2)\^{}2}

\emph{∴ নির্ণেয় গ.সা.গু = 2(x + 2)}
\end{quote}

\textbf{a\^{}2 − 3a, a\^{}3 − 9a এবং a\^{}3 − 4a\^{}2 + 3a এর গ.সা.গু.
=?}

\textbf{\hl{ক) a(a − 3) ✔}}

খ) a − 3

গ) a

ঘ) a(a + 3)

\textbf{\hl{✔ সঠিক উত্তর: ক) a(a − 3)}}

\begin{quote}
\textbf{ব্যাখ্যা:} \emph{প্রথম রাশি, a\^{}2 − 3a = a(a − 3)}

\emph{দ্বিতীয় রাশি, a\^{}3 − 9a = a(a\^{}2 − 9) = a(a + 3)(a − 3)}

\emph{তৃতীয় রাশি, a\^{}3 − 4a\^{}2 + 3a = a(a\^{}2 − 4a + 3) = a(a − 3)(a
− 1)}

\emph{∴ নির্ণেয় গ.সা.গু = a(a − 3)}
\end{quote}

\textbf{x\^{}2 − 11x + 30 এবং x\^{}3 − 4x\^{}2 − 2x − 15 এর গ.সা.গু. কত?}

\textbf{\hl{ক) x − 5 ✔}}

খ) x − 6

গ) x\^{}2 + x + 3

ঘ) x\^{}2 − x − 3

\textbf{\hl{✔ সঠিক উত্তর: ক) x − 5}}

\begin{quote}
\textbf{ব্যাখ্যা:} \emph{১ম রাশি, x\^{}2 − 11x + 30 = (x − 6)(x − 5)}

\emph{২য় রাশি, x\^{}3 − 4x\^{}2 − 2x − 15 = (x − 5)(x\^{}2 + x + 3)}

\emph{∴ রাশি দুটির গ.সা.গু = (x − 5)}
\end{quote}

\textbf{x\^{}2 + 5x, x\^{}2 − 25, x\^{}2 + 7x + 10 এর গ.সা.গু কত?}

ক) x − 5

\textbf{\hl{খ) x + 5 ✔}}

গ) x (x + 5)

ঘ) x (x − 5)

\textbf{\hl{✔ সঠিক উত্তর: খ) x + 5}}

\begin{quote}
\textbf{ব্যাখ্যা:} \emph{১ম রাশি = x\^{}2 + 5x = x(x + 5)}

\emph{২য় রাশি = x\^{}2 − 25 = (x + 5)(x − 5)}

\emph{৩য় রাশি = x\^{}2 + 7x + 10 = (x + 5)(x + 2)}

\emph{∴ নির্ণেয় গ.সা.গু = (x + 5)}
\end{quote}

\textbf{4x\^{}2 − 16 এবং 6x\^{}2 + 24x + 24 এর গ.সা.গু.-\/-\/-\/-\/-}

ক) 4(x+2)

খ) 2(x+3)

\textbf{\hl{গ) 2(x+2) ✔}}

ঘ) 2(x+5)

\textbf{\hl{✔ সঠিক উত্তর: গ) 2(x+2)}}

\begin{quote}
\textbf{ব্যাখ্যা:} \emph{প্রথম রাশি, 4x\^{}2 − 4(4)}

\emph{= 4(x\^{}2 − 4)}

\emph{= 4(x + 2)(x − 2)}

\emph{দ্বিতীয় রাশি, 6x\^{}2 + 24x + 24}

\emph{= 2(3x\^{}2 + 12x + 12)}

\emph{= 2(3x\^{}2 + 6x + 6x + 12)}

\emph{= 2\{3x(x + 2) + 6(x + 2)\}}

\emph{= 2(x + 2)(3x + 6)}

\emph{∴ নির্ণেয় গসাগু = 2(x + 2)}
\end{quote}

\textbf{a(a + b), ab(a + b), এবং a(a + b)\^{}2 এর গ. সা. গু কত?}

\textbf{\hl{ক) a(a + b) ✔}}

খ) ab(a + b)

গ) a + b

ঘ) a(a + b)\^{}2

\textbf{\hl{✔ সঠিক উত্তর: ক) a(a + b)}}

\begin{quote}
\textbf{ব্যাখ্যা:} \emph{১ম রাশি = a(a + b)}

\emph{২য় রাশি = ab(a + b)}

\emph{৩য় রাশি = a(a + b)\^{}2 = a(a + b)(a + b)}

\emph{∴ নির্ণেয় গ.সা.গু. = a(a + b)}
\end{quote}

\textbf{x = a\^{}2 − 7a + 6, y = a\^{}2 − 2a + 1, z = a\^{}2 + 5a + 4
হলে x, y, z এর গ.সা.গু কত?}

ক) (a − 6)

খ) (a − 1)\^{}2

গ) (a − 4)

\textbf{\hl{ঘ) 1 ✔}}

\textbf{\hl{✔ সঠিক উত্তর: ঘ) 1}}

\begin{quote}
\textbf{ব্যাখ্যা:} \emph{x = a\^{}2 − 7a + 6 = (a − 6)(a − 1)}

\emph{y = a\^{}2 − 2a + 1 = (a − 1)\^{}2 = (a − 1)(a − 1)}

\emph{z = a\^{}2 + 5a + 4 = (a + 4)(a + 1)}

\emph{∴ x, y এবং z এর গ.সা.গু. = 1 {[}যেহেতু কোনো Common রাশি নেই{]}}
\end{quote}

\textbf{x\^{}2 − y\^{}2, (x+y)\^{}2, x\^{}3 + y\^{}3 এর গ.সা. কত?}

ক) x − y

\textbf{\hl{খ) x + y ✔}}

গ) (x\^{}2 − y\^{}2)(x\^{}3 + y\^{}3)

ঘ) (x−y)(x+y)\^{}2(x\^{}2 − xy + y\^{}2)

\textbf{\hl{✔ সঠিক উত্তর: খ) x + y}}

\begin{quote}
\textbf{ব্যাখ্যা:} \emph{১ম রাশি = x\^{}2 − y\^{}2 = (x + y)(x − y)}

\emph{২য় রাশি = (x + y)\^{}2 = (x + y)(x + y)}

\emph{৩য় রাশি = x\^{}3 + y\^{}3 = (x + y)(x\^{}2 − xy + y\^{}2)}

\emph{∴ নির্ণেয় গ.সা.গু. = x + y}
\end{quote}

\textbf{দুটি সংখ্যার ল.সা.গু a\^{}2b(a + b) এবং গ.সা.গু a(a + b)। একটি সংখ্যা
a\^{}3 + a\^{}2b হলে, অপর সংখ্যাটি কত?}

ক) a\^{}2b + a\^{}2b\^{}2

\textbf{\hl{খ) a\^{}2b + ab\^{}2 ✔}}

গ) ab\^{}2 + a\^{}2b\^{}2

ঘ) a\^{}3 − b\^{}3

\textbf{\hl{✔ সঠিক উত্তর: খ) a\^{}2b + ab\^{}2}}

\begin{quote}
\textbf{ব্যাখ্যা:} \emph{দুটি সংখ্যার গুণফল = ল.সা.গু × গ.সা.গু = \{a\^{}2b(a +
b)\} × \{a(a + b)\} = a\^{}3b(a + b)(a + b)}

\emph{ধরি, অপর সংখ্যাটি = x}

\emph{x × (a\^{}3 + a\^{}2b) = a\^{}3b(a + b)(a + b)}

\emph{⇒ x × a\^{}2(a + b) = a\^{}3b(a + b)(a + b)}

\emph{⇒ x = (a\^{}3b(a+b)(a+b)/a\^{}2(a+b)) = ab(a + b) = a\^{}2b +
ab\^{}2}
\end{quote}

\textbf{দুটি সংখ্যার ল.সা.গু a\^{}2b(a + b) এবং গ.সা.গু a(a + b)। একটি সংখ্যা
a\^{}3 + a\^{}2b হলে, অপর সংখ্যাটি কত?}

\textbf{\hl{ক) a\^{}2b + ab\^{}2 ✔}}

খ) ab + ab\^{}2

গ) a\^{}2b + ab

ঘ) a\^{}2b + a\^{}2

\textbf{\hl{✔ সঠিক উত্তর: ক) a\^{}2b + ab\^{}2}}

\begin{quote}
\textbf{ব্যাখ্যা:} \emph{অপর সংখ্যা = গ.সা.গু × ল.সা.গু একটি সংখ্যা}

\emph{= a 2 b ( a + b ) × a ( a + b ) a 3 + a 2 b}

\emph{= a 3 b ( a + b ) ( a + b ) a 2 ( a + b )}

\emph{= ab(a + b)}

\emph{= a\^{}2b + ab\^{}2}
\end{quote}

\section{34. ল.সা.গু এবং গ.সা.গু. সম্পর্কিত সমস্যা MCQ প্রশ্ন উত্তর
অনুশীলন}\label{ux9b2.ux9b8.ux997-ux98fux9ac-ux997.ux9b8.ux997.-ux9b8ux9aeux9aaux9b0ux995ux9a4-ux9b8ux9aeux9b8ux9af-mcq-ux9aaux9b0ux9b6ux9a8-ux989ux9a4ux9a4ux9b0-ux985ux9a8ux9b6ux9b2ux9a8}

\textbf{(২/৫) এবং (৫/৯) এর গ.সা.গু. কত?}

\textbf{\hl{ক) (১/৪৫) ✔}}

খ) (২/৪৫)

গ) (৭/৪৫)

ঘ) ৪৫

\textbf{\hl{✔ সঠিক উত্তর: ক) (১/৪৫)}}

\begin{quote}
\textbf{ব্যাখ্যা:} \emph{ভগ্নাংশের গ.সা.গু. = (লবগুলোর গ.সা.গু./হরগুলোর
ল.সা.গু.)}

\emph{এখানে, লব ২ ও ৫ এর গ.সা.গু. = ১}

\emph{এবং হর ৫ ও ৯ এর ল.সা.গু. = ৪৫}

\emph{∴ নির্ণেয় গ.সা.গু. = (১/৪৫)}
\end{quote}

\textbf{0,2,3 এর গ.সা.গু কত?}

ক) 3

খ) 2

\textbf{\hl{গ) 1 ✔}}

ঘ) 0

\textbf{\hl{✔ সঠিক উত্তর: গ) 1}}

\begin{quote}
\textbf{ব্যাখ্যা:} \emph{সমাধানঃ}

\emph{0 = 1 × 0;}

\emph{2 = 1 × 2;}

\emph{3 = 1 × 3;}

\emph{এদের সাধারণ উৎপাদক = 1}

\emph{∴ নির্ণেয় গ.সা.গু = 1}
\end{quote}

\textbf{(৩/৪) , (৪/৫) , (৫/৬) এর গ. সা. গু কত?}

ক) ৩০

খ) (১/৫০)

\textbf{\hl{গ) (১/৬০) ✔}}

ঘ) ৬০

\textbf{\hl{✔ সঠিক উত্তর: গ) (১/৬০)}}

\begin{quote}
\textbf{ব্যাখ্যা:} \emph{আমরা জানি,}

\emph{দুই বা ততোধিক ভগ্নাংশের গ.সা.গু = (লবগুলোর গ.সা.গু/হরগুলোর ল.সা.গু)}

\emph{এখানে, লব ৩, ৪, ৫ এর গ.সা.গু = ১}

\emph{এবং হর ৪, ৫, ৬ এর ল.সা.গু = ৬০}

\emph{∴ নির্ণেয় গ.সা.গু = (১/৬০)}
\end{quote}

\textbf{(৩/৫) , (১/৪) , (২/৩) এর ল.সা.গু কত?}

ক) (১/৬০)

\textbf{\hl{খ) ৬ ✔}}

গ) ৬০

ঘ) ১২

\textbf{\hl{✔ সঠিক উত্তর: খ) ৬}}

\begin{quote}
\textbf{ব্যাখ্যা:} \emph{আমরা জানি,}

\emph{দুই বা ততোধিক ভগ্নাংশের ল.সা.গু = (লবগুলোর ল.সা.গু/হরগুলোর গ.সা.গু)}

\emph{এখানে, লব ৩, ১, ২ এর ল.সা.গু = ৬}

\emph{এবং হর ৫, ৪, ৩ এর গ.সা.গু = ১}

\emph{∴ নির্ণেয় ল.সা.গু = (৬/১)=৬}
\end{quote}

\textbf{২৪, ৩০ এবং ৭৭ এর গ. সা. গু কত?}

\textbf{\hl{ক) ১ ✔}}

খ) ২

গ) ৩

ঘ) ৫

\textbf{\hl{✔ সঠিক উত্তর: ক) ১}}

\begin{quote}
\textbf{ব্যাখ্যা:} \emph{২৪ = ১ × ২ × ২ × ২ × ৩ ৩০ = ১ × ২ × ৩ × ৫ ৭৭ = ১
× ৭ × ১১}

\emph{সংখ্যা তিনটির মধ্যে ১ ব্যতীত কোনো সাধারণ মৌলিক উৎপাদক নেই।}

\emph{∴ ২৪, ৩০ ও ৭৭ এর গ. সা. গু = ১।}
\end{quote}

\textbf{(২/৫) , (৩/৫) , (৬/১৫) এর গ.সা.গু. কোনটি?}

ক) (৬/৫)

খ) (৭/৫)

গ) (৮/৫)

\textbf{\hl{ঘ) (১/১৫) ✔}}

\textbf{\hl{✔ সঠিক উত্তর: ঘ) (১/১৫)}}

\begin{quote}
\textbf{ব্যাখ্যা:} \emph{প্রদত্ত ভগ্নাংশগুলো: (২/৫) , (৩/৫) এবং (৬/১৫)
{[}লঘিষ্ঠ আকার: (২/৫) {]}}

\emph{ভগ্নাংশের গ.সা.গু. = (লবগুলোর গ.সা.গু./হরগুলোর ল.সা.গু.)}

\emph{এখানে, লব ২, ৩, ৬ এর গ.সা.গু. = ১}

\emph{হর ৫, ৫, ১৫ এর ল.সা.গু. = ১৫}

\emph{∴ নির্ণেয় গ.সা.গু. = (১/১৫)}
\end{quote}

\textbf{(৩/৪) , (৪/৫) ও (৫/৬) এর গ.সা.গু. কত?}

ক) ৬০

খ) ৩০

গ) (১/৩০)

\textbf{\hl{ঘ) (১/৬০) ✔}}

\textbf{\hl{✔ সঠিক উত্তর: ঘ) (১/৬০)}}

\begin{quote}
\textbf{ব্যাখ্যা:} \emph{ভগ্নাংশের গ.সা.গু. = (লবগুলোর গ.সা.গু./হরগুলোর
ল.সা.গু.)}

\emph{এখানে, লব ৩, ৪, ৫ এর গ.সা.গু. = ১}

\emph{এবং হর ৪, ৫, ৬ এর ল.সা.গু. = ৬০}

\emph{∴ ভগ্নাংশের গ.সা.গু. = (১/৬০)}
\end{quote}

\textbf{নিচের কোন সংখ্যাটি ৪ দ্বারা নিঃশেষে বিভাজ্য নয়?}

ক) ৭৪৮

খ) ৪৭২

\textbf{\hl{গ) ৬৪৬ ✔}}

ঘ) ১৫২

\textbf{\hl{✔ সঠিক উত্তর: গ) ৬৪৬}}

\begin{quote}
\textbf{ব্যাখ্যা:} \emph{আমরা জানি, কোনো সংখ্যার শেষ দুই অঙ্ক দিয়ে গঠিত সংখ্যা
৪ দ্বারা বিভাজ্য হলে সমগ্র সংখ্যাটি ৪ দ্বারা নিঃশেষে বিভাজ্য হয়।}

\emph{এখানে,}

\emph{a) ৭৪৮ ⇒ ৪৮ ÷ ৪ = ১২ (বিভাজ্য)}

\emph{b) ৪৭২ ⇒ ৭২ ÷ ৪ = ১৮ (বিভাজ্য)}

\emph{c) ৬৪৬ ⇒ ৪৬ ÷ ৪ = ১১ (ভাগশেষ ২ থাকে, বিভাজ্য নয়)}

\emph{d) ১৫২ ⇒ ৫২ ÷ ৪ = ১৩ (বিভাজ্য)}

\emph{∴ ৬৪৬ সংখ্যাটি ৪ দ্বারা নিঃশেষে বিভাজ্য নয়।}
\end{quote}

\textbf{নিচের কোন সংখ্যাটি ৮ দ্বারা নিঃশেষে বিভাজ্য নয়?}

ক) ৭৪৪

খ) ৬৪৮

গ) ৪৭২

\textbf{\hl{ঘ) ১৫৬ ✔}}

\textbf{\hl{✔ সঠিক উত্তর: ঘ) ১৫৬}}

\begin{quote}
\textbf{ব্যাখ্যা:} \emph{প্রত্যেকটি সংখ্যাকে ৮ দ্বারা ভাগ করে দেখি:}

\emph{৭৪৪ ÷ ৮ = ৯৩ (নিঃশেষে বিভাজ্য)}

\emph{৬৪৮ ÷ ৮ = ৮১ (নিঃশেষে বিভাজ্য)}

\emph{৪৭২ ÷ ৮ = ৫৯ (নিঃশেষে বিভাজ্য)}

\emph{১৫৬ ÷ ৮ = ১৯, ভাগশেষ থাকে ৪ (নিঃশেষে বিভাজ্য নয়)}

\emph{∴ নির্ণেয় সংখ্যা ১৫৬।}
\end{quote}

\textbf{১০৫৬ এর সাথে সর্বনিম্ন কত যোগ করলে যোগফল ২৩ দ্বারা নিঃশেষে বিভাজ্য
হবে?}

\textbf{\hl{ক) ২ ✔}}

খ) ৩

গ) ১৮

ঘ) ২১

\textbf{\hl{✔ সঠিক উত্তর: ক) ২}}

\begin{quote}
\textbf{ব্যাখ্যা:} \emph{১০৫৬ কে ২৩ দ্বারা ভাগ করলে পাই, ১০৫৬ = ২৩ × ৪৫ +
২১}

\emph{এখানে ভাগশেষ থাকে ২১।}

\emph{সংখ্যাটিকে ২৩ দ্বারা নিঃশেষে বিভাজ্য করতে হলে এর সাথে সর্বনিম্ন যোগ
করতে হবে = (২৩ − ২১) = ২।}
\end{quote}

\textbf{একটি সংখ্যাকে ৬৪ দ্বারা ভাগ করলে ৬৩ অবশিষ্ট থাকে। কিন্তু ঐ সংখ্যাকে ৩২
দ্বারা ভাগ করলে কত অবশিষ্ট থাকবে?}

ক) ৩২

\textbf{\hl{খ) ৩১ ✔}}

গ) ৫২

ঘ) ২৯

\textbf{\hl{✔ সঠিক উত্তর: খ) ৩১}}

\begin{quote}
\textbf{ব্যাখ্যা:} \emph{প্রথম ভাজক (৬৪) যদি দ্বিতীয় ভাজক (৩২) দ্বারা বিভাজ্য
হয়, তবে প্রথম অবশিষ্টকে (৬৩) দ্বিতীয় ভাজক (৩২) দিয়ে ভাগ করলেই নির্ণেয় অবশিষ্ট
পাওয়া যায়।}

\emph{৬৩ কে ৩২ দ্বারা ভাগ করলে পাই, ৬৩ = ৩২ × ১ + ৩১}

\emph{∴ অবশিষ্ট থাকবে ৩১।}
\end{quote}

\textbf{একটি সংখ্যাকে ৫৬৭ দিয়ে ভাগ করলে ভাগশেষ ১০ থাকে। যদি ঐ সংখ্যাকে ৭
দিয়ে ভাগ করা হয় তবে ভাগশেষ কত হবে?}

\textbf{\hl{ক) ৩ ✔}}

খ) ৪

গ) ৫

ঘ) কোনোটিই নয়

\textbf{\hl{✔ সঠিক উত্তর: ক) ৩}}

\begin{quote}
\textbf{ব্যাখ্যা:} \emph{যেহেতু প্রথম ভাজক ৫৬৭ সংখ্যাটি ৭ দ্বারা বিভাজ্য (৫৬৭ ÷
৭ = ৮১), তাই প্রথম ভাগশেষ ১০-কে ৭ দ্বারা ভাগ করলেই উত্তর পাওয়া যাবে।}

\emph{১০ কে ৭ দ্বারা ভাগ করলে পাই, ১০ = ৭ × ১ + ৩}

\emph{∴ নির্ণেয় ভাগশেষ ৩।}
\end{quote}

\textbf{কোন বৃহত্তম সংখ্যা দ্বারা ১০২ ও ১৮৬ কে ভাগ করলে প্রত্যেকবার ভাগশেষ ৬
থাকবে?}

\textbf{\hl{ক) ১২ ✔}}

খ) ১৫

গ) ১৬০

ঘ) ২২

\textbf{\hl{✔ সঠিক উত্তর: ক) ১২}}

\begin{quote}
\textbf{ব্যাখ্যা:} \emph{নির্ণেয় বৃহত্তম সংখ্যাটি হবে (১০২ − ৬ = ৯৬) এবং (১৮৬
− ৬ = ১৮০) এর গ.সা.গু।}

\emph{এখানে, ৯৬ এবং ১৮০ এর গ.সা.গু = ১২।}

\emph{∴ নির্ণেয় বৃহত্তম সংখ্যা ১২।}
\end{quote}

\textbf{কোন ক্ষুদ্রতম সংখ্যাকে ৩, ৫, ৯ দ্বারা ভাগ করলে ভাগশেষ ২ হবে?}

ক) ৪৩

খ) ৪৫

গ) ৪১

\textbf{\hl{ঘ) ৪৭ ✔}}

\textbf{\hl{✔ সঠিক উত্তর: ঘ) ৪৭}}

\begin{quote}
\textbf{ব্যাখ্যা:} \emph{প্রথমেই ৩, ৫ ও ৯ এর ল.সা.গু নির্ণয় করতে হবে।}

\emph{৩, ৫ ও ৯ এর ল.সা.গু = ৪৫}

\emph{যেহেতু প্রতিক্ষেত্রে ভাগশেষ ২ অবশিষ্ট থাকে,}

\emph{⇒ ক্ষুদ্রতম সংখ্যাটি হবে = ৪৫ + ২ = ৪৭।}

\emph{∴ নির্ণেয় ক্ষুদ্রতম সংখ্যা ৪৭।}
\end{quote}

\textbf{কোন ক্ষুদ্রতম সংখ্যার সাথে ২ যোগ করলে যোগফল ৩, ৬, ৯, ১২ এবং ১৫ দ্বারা
নিঃশেষে বিভাজ্য হবে?}

\textbf{\hl{ক) ১৭৮ ✔}}

খ) ৩৫৮

গ) ৩৬৮

ঘ) ৭১৮

\textbf{\hl{✔ সঠিক উত্তর: ক) ১৭৮}}

\begin{quote}
\textbf{ব্যাখ্যা:} \emph{৩, ৬, ৯, ১২ ও ১৫ এর ল.সা.গু = ১৮০।}

\emph{যেহেতু সংখ্যাটির সাথে ২ যোগ করলে সংখ্যাটি নিঃশেষে বিভাজ্য হবে,}

\emph{⇒ ক্ষুদ্রতম সংখ্যাটি হবে = ১৮০ − ২ = ১৭৮।}

\emph{∴ নির্ণেয় সংখ্যা ১৭৮।}
\end{quote}

\textbf{কোন লঘিষ্ঠ সংখ্যার সাথে ৩ যোগ করলে যোগফল ২৪, ৩৬ এবং ৪৮ দ্বারা
নিঃশেষে বিভাজ্য হবে?}

ক) ৮৯

\textbf{\hl{খ) ১৪১ ✔}}

গ) ২৪৮

ঘ) ১৭০

\textbf{\hl{✔ সঠিক উত্তর: খ) ১৪১}}

\begin{quote}
\textbf{ব্যাখ্যা:} \emph{২৪, ৩৬ ও ৪৮ এর ল.সা.গু = ১৪৪।}

\emph{ক্ষুদ্রতম বা লঘিষ্ঠ সংখ্যা বের করার কথা বলা হলে ল.সা.গু নির্ণয় করতে হয়।
যেহেতু সংখ্যাটির সাথে ৩ যোগ করলে যোগফল বিভাজ্য হবে,}

\emph{⇒ লঘিষ্ঠ সংখ্যাটি হবে = ১৪৪ − ৩ = ১৪১।}

\emph{∴ নির্ণেয় লঘিষ্ঠ সংখ্যা ১৪১।}
\end{quote}

\textbf{একটি পূর্ণসংখ্যা নির্ণয় করুন যাকে ৩, ৪, ৫ এবং ৬ দ্বারা ভাগ করলে
যথাক্রমে ২, ৩, ৪ এবং ৫ অবশিষ্ট থাকে?}

ক) ৪৭

খ) ৪৯

গ) ৫৭০

\textbf{\hl{ঘ) ৫৯ ✔}}

\textbf{\hl{✔ সঠিক উত্তর: ঘ) ৫৯}}

\begin{quote}
\textbf{ব্যাখ্যা:} \emph{এখানে ভাজক ও ভাগশেষের পার্থক্য সমান:}

\emph{৩ − ২ = ১}

\emph{৪ − ৩ = ১}

\emph{৫ − ৪ = ১}

\emph{৬ − ৫ = ১}

\emph{এখন, ৩, ৪, ৫ এবং ৬ এর ল.সা.গু = ৬০}

\emph{⇒ নির্ণেয় পূর্ণসংখ্যাটি হবে = ৬০ − ১ = ৫৯।}
\end{quote}

\textbf{কোন লঘিষ্ঠ সংখ্যাকে ২৪ ও ৩৬ দ্বারা ভাগ করলে যথাক্রমে ১৪ ও ২৬ অবশিষ্ট
থাকবে?}

ক) ৪৮

\textbf{\hl{খ) ৬২ ✔}}

গ) ৭৮০

ঘ) ৮৪

\textbf{\hl{✔ সঠিক উত্তর: খ) ৬২}}

\begin{quote}
\textbf{ব্যাখ্যা:} \emph{এখানে ভাজক ও ভাগশেষের পার্থক্য সমান:}

\emph{২৪ − ১৪ = ১০}

\emph{৩৬ − ২৬ = ১০}

\emph{২৪ এবং ৩৬ এর ল.সা.গু = ৭২}

\emph{⇒ নির্ণেয় লঘিষ্ঠ সংখ্যাটি হবে = ৭২ − ১০ = ৬২।}
\end{quote}

\textbf{মজুদ ২৬৪ টি আপেলের সাথে আর ও ন্যূনতম কতটি আপেল পাওয়া গেলে তা ৬, ৭
বা ৮ জন ব্যক্তির মধ্যে সমানভাবে ভাগ করে দেওয়া যাবে?}

ক) ৭০০

\textbf{\hl{খ) ৭২ ✔}}

গ) ৬৬

ঘ) ৬৪

\textbf{\hl{✔ সঠিক উত্তর: খ) ৭২}}

\begin{quote}
\textbf{ব্যাখ্যা:} \emph{৬, ৭ এবং ৮ এর ল.সা.গু = ১৬৮।}

\emph{১৬৮ টি আপেল ৬, ৭ বা ৮ জন ব্যক্তির মধ্যে সমানভাবে ভাগ করে দেওয়া যায়।
১৬৮ এর গুণিতক (১৬৮ × ২) = ৩৩৬ টি আপেলকেও সমানভাবে ভাগ করে দেওয়া যাবে।}

\emph{যেহেতু মজুদ আছে ২৬৪ টি আপেল,}

\emph{⇒ আরো আপেল লাগবে = ৩৩৬ − ২৬৪ = ৭২ টি।}
\end{quote}

\textbf{নিম্নের কোন ক্ষুদ্রতম সংখ্যাকে ৩, ৫ ও ৬ দ্বারা ভাগ করলে ভাগশেষ হবে ১?}

ক) ৭১

খ) ৪১

\textbf{\hl{গ) ৩১ ✔}}

ঘ) ৩৯

\textbf{\hl{✔ সঠিক উত্তর: গ) ৩১}}

\begin{quote}
\textbf{ব্যাখ্যা:} \emph{৩, ৫ ও ৬ এর ল.সা.গু = ৩০।}

\emph{যেহেতু ৩, ৫ ও ৬ দ্বারা ভাগ করলে প্রতিক্ষেত্রে ভাগশেষ ১ হবে,}

\emph{⇒ ক্ষুদ্রতম সংখ্যাটি হবে = ৩০ + ১ = ৩১।}
\end{quote}

\textbf{কোন লঘিষ্ঠ সংখ্যাকে ১৮ ও ২৪ দ্বারা ভাগ করলে অবশিষ্ট যথাক্রমে ৪ ও ১০
হবে?}

ক) ৪০

\textbf{\hl{খ) ৫৮ ✔}}

গ) ২৪

ঘ) কোনোটিই নয়

\textbf{\hl{✔ সঠিক উত্তর: খ) ৫৮}}

\begin{quote}
\textbf{ব্যাখ্যা:} \emph{এখানে ভাজক ও ভাগশেষের পার্থক্য সমান:}

\emph{১৮ − ৪ = ১৪}

\emph{২৪ − ১০ = ১৪}

\emph{১৮ ও ২৪ এর ল.সা.গু = ৭২}

\emph{⇒ নির্ণেয় লঘিষ্ঠ সংখ্যাটি হবে = ৭২ − ১৪ = ৫৮।}
\end{quote}

\textbf{কোন ক্ষুদ্রতম সংখ্যার সাথে ২ যোগ করলে যোগফল ৩, ৬, ৯, ১২ এবং ১৫ দ্বারা
নিঃশেষে বিভাজ্য হবে?}

\textbf{\hl{ক) ১৭৮ ✔}}

খ) ৪৫৮

গ) ৩৬৮

ঘ) ৭১৮

\textbf{\hl{✔ সঠিক উত্তর: ক) ১৭৮}}

\begin{quote}
\textbf{ব্যাখ্যা:} \emph{৩, ৬, ৯, ১২ এবং ১৫ এর ল.সা.গু = ১৮০।}

\emph{যেহেতু ক্ষুদ্রতম সংখ্যার সাথে ২ যোগ করলে সংখ্যাটি নিঃশেষে বিভাজ্য হবে,}

\emph{⇒ ক্ষুদ্রতম সংখ্যাটি হবে = ১৮০ − ২ = ১৭৮।}
\end{quote}

\textbf{দুইটি সংখ্যার গ.সা.গু. ও ল.সা.গু যথাক্রমে ২ ও ৩৬০। একটি সংখ্যা ১০ হলে
অপর সংখ্যাটি কত?}

ক) ৪৮

খ) ৬০

গ) ৫০

\textbf{\hl{ঘ) ৭২ ✔}}

\textbf{\hl{✔ সঠিক উত্তর: ঘ) ৭২}}

\begin{quote}
\textbf{ব্যাখ্যা:} \emph{আমরা জানি,}

\emph{অপর সংখ্যাটি = (সংখ্যাদ্বয়ের ল.সা.গু × সংখ্যাদ্বয়ের গ.সা.গু/একটি সংখ্যা)}

\emph{⇒ অপর সংখ্যাটি = (৩৬০×২/১০) = ৭২}

\emph{∴ অপর সংখ্যাটি ৭২।}
\end{quote}

\textbf{দুইটি সংখ্যার গ.সা.গু. ১৫ এবং ল.সা.গু. ৯০। একটি সংখ্যা ৩০ হলে অপর
সংখ্যাটি কত?}

ক) ৬০

খ) ৫০

\textbf{\hl{গ) ৪৫ ✔}}

ঘ) ৭৫

\textbf{\hl{✔ সঠিক উত্তর: গ) ৪৫}}

\begin{quote}
\textbf{ব্যাখ্যা:} \emph{আমরা জানি,}

\emph{অপর সংখ্যাটি = (সংখ্যাদ্বয়ের ল.সা.গু × সংখ্যাদ্বয়ের গ.সা.গু/একটি সংখ্যা)}

\emph{⇒ অপর সংখ্যাটি = (১৫×৯০/৩০) = ৪৫}

\emph{∴ অপর সংখ্যাটি ৪৫।}
\end{quote}

\textbf{দুইটি সংখ্যার ল.সা.গু ও গ.সা.গু. যথাক্রমে ২৮৮ ও ১২। একটি সংখ্যা ৩৬
হলে, অপর সংখ্যাটি কত?}

\textbf{\hl{ক) ৯৬ ✔}}

খ) ৭২

গ) ৯২

ঘ) কোনোটিই নয়

\textbf{\hl{✔ সঠিক উত্তর: ক) ৯৬}}

\begin{quote}
\textbf{ব্যাখ্যা:} \emph{আমরা জানি,}

\emph{অপর সংখ্যাটি = (সংখ্যাদ্বয়ের ল.সা.গু × সংখ্যাদ্বয়ের গ.সা.গু/একটি সংখ্যা)}

\emph{⇒ অপর সংখ্যাটি = (২৮৮×১২/৩৬) = ৯৬}

\emph{∴ অপর সংখ্যাটি ৯৬।}
\end{quote}

\textbf{দুইটি সংখ্যার গ.সা.গু. ৭ এবং ল.সা.গু. ৮৪। একটি সংখ্যা ৪২ হলে অপর
সংখ্যাটি কত?}

ক) ২০

খ) ২১

\textbf{\hl{গ) ১৪ ✔}}

ঘ) ২৮

\textbf{\hl{✔ সঠিক উত্তর: গ) ১৪}}

\begin{quote}
\textbf{ব্যাখ্যা:} \emph{আমরা জানি,}

\emph{অপর সংখ্যাটি = (সংখ্যাদ্বয়ের ল.সা.গু × সংখ্যাদ্বয়ের গ.সা.গু/একটি সংখ্যা)}

\emph{⇒ অপর সংখ্যাটি = (৮৪×৭/৪২) = ১৪}

\emph{∴ অপর সংখ্যাটি ১৪।}
\end{quote}

\textbf{দুইটি সংখ্যার গ.সা.গু. ১৫ এবং ল.সা.গু. ৪২০। একটি সংখ্যা ৬০ হলে অপর
সংখ্যাটি কত?}

\textbf{\hl{ক) ১০৫ ✔}}

খ) ১০৬

গ) ১০০

ঘ) ৯০

\textbf{\hl{✔ সঠিক উত্তর: ক) ১০৫}}

\begin{quote}
\textbf{ব্যাখ্যা:} \emph{আমরা জানি,}

\emph{অপর সংখ্যাটি = (সংখ্যাদ্বয়ের ল.সা.গু × সংখ্যাদ্বয়ের গ.সা.গু/একটি সংখ্যা)}

\emph{⇒ অপর সংখ্যাটি = (৪২০×১৫/৬০) = ১০৫}

\emph{∴ অপর সংখ্যাটি ১০৫।}
\end{quote}

\textbf{দুইটি সংখ্যার ল.সা.গু. ৪৮ এবং গ.সা.গু. ৪। একটি সংখ্যা ১৬ হলে অপর
সংখ্যাটি কত?}

ক) ১৪

\textbf{\hl{খ) ১২ ✔}}

গ) ১৫০

ঘ) ১৮

\textbf{\hl{✔ সঠিক উত্তর: খ) ১২}}

\begin{quote}
\textbf{ব্যাখ্যা:} \emph{আমরা জানি,}

\emph{অপর সংখ্যাটি = (সংখ্যাদ্বয়ের ল.সা.গু × সংখ্যাদ্বয়ের গ.সা.গু/একটি সংখ্যা)}

\emph{⇒ অপর সংখ্যাটি = (৪৮×৪/১৬) = ১২}

\emph{∴ অপর সংখ্যাটি ১২।}
\end{quote}

\textbf{দুইটি সংখ্যার গ.সা.গু. ও ল.সা.গু যথাক্রমে ১৪ ও ১৬৮। একটি সংখ্যা ৪২ হলে
অপর সংখ্যাটি কত?}

\textbf{\hl{ক) ৫৬ ✔}}

খ) ৬২

গ) ৮৪

ঘ) ১১২

\textbf{\hl{✔ সঠিক উত্তর: ক) ৫৬}}

\begin{quote}
\textbf{ব্যাখ্যা:} \emph{আমরা জানি,}

\emph{অপর সংখ্যাটি = (সংখ্যাদ্বয়ের ল.সা.গু × সংখ্যাদ্বয়ের গ.সা.গু/একটি সংখ্যা)}

\emph{⇒ অপর সংখ্যাটি = (১৬৮×১৪/৪২) = ৫৬}

\emph{∴ অপর সংখ্যাটি ৫৬।}
\end{quote}

\textbf{দুইটি সংখ্যার অনুপাত ৫:৭ এবং তাদের গ.সা.গু ৬ হলে সংখ্যা দুইটির ল.সা.গু.
কত?}

\textbf{\hl{ক) ২১০ ✔}}

খ) ১৮০

গ) ২০০

ঘ) ২২০

\textbf{\hl{✔ সঠিক উত্তর: ক) ২১০}}

\begin{quote}
\textbf{ব্যাখ্যা:} \emph{আমরা জানি,}

\emph{ল.সা.গু = অনুপাতের সংখ্যাদ্বয়ের গুণফল × গ.সা.গু}

\emph{⇒ ল.সা.গু = ৫ × ৭ × ৬ = ২১০}

\emph{∴ নির্ণেয় ল.সা.গু ২১০।}
\end{quote}

\textbf{দুইটি সংখ্যার অনুপাত ৫:৬ এবং তাদের গ.সা.গু ৮ হলে সংখ্যা দুইটির ল.সা.গু.
কত?}

ক) ২১০

খ) ২২০

গ) ১৪৮

\textbf{\hl{ঘ) ২৪০ ✔}}

\textbf{\hl{✔ সঠিক উত্তর: ঘ) ২৪০}}

\begin{quote}
\textbf{ব্যাখ্যা:} \emph{আমরা জানি,}

\emph{ল.সা.গু = অনুপাতের সংখ্যাদ্বয়ের গুণফল × গ.সা.গু}

\emph{⇒ ল.সা.গু = ৫ × ৬ × ৮ = ২৪০}

\emph{∴ নির্ণেয় ল.সা.গু ২৪০।}
\end{quote}

\textbf{দুটি সংখ্যার অনুপাত ৩:৪ তাদের ল.সা.গু ১০৮। সংখ্যা দুটির যোগফল কত?}

ক) ৫৫

খ) ৫৬

গ) ৬০

\textbf{\hl{ঘ) ৬৩ ✔}}

\textbf{\hl{✔ সঠিক উত্তর: ঘ) ৬৩}}

\begin{quote}
\textbf{ব্যাখ্যা:} \emph{ধরি, সংখ্যা দুটি যথাক্রমে ৩x ও ৪x এবং তাদের ল.সা.গু
১২x।}

\emph{প্রশ্নমতে,}

\emph{১২x = ১০৮}

\emph{⇒ x = ৯}

\emph{সুতরাং সংখ্যা দুটি যথাক্রমে (৩ × ৯) = ২৭ এবং (৪ × ৯) = ৩৬।}

\emph{∴ সংখ্যা দুটির যোগফল = ২৭ + ৩৬ = ৬৩।}
\end{quote}

\textbf{কোন কোন স্বাভাবিক সংখ্যা দ্বারা ৩৪৬ কে ভাগ করলে প্রতিক্ষেত্রে ৩১
অবশিষ্ট থাকে?}

\textbf{\hl{ক) ৩৫, ৪৫, ৬৩, ১০৫, ৩১৫ ✔}}

খ) ৩৫, ৪০, ৬৫, ১১০, ৩১৫

গ) ৩৫, ৪৫, ৭০, ১০৫, ৩১৫

ঘ) ৩৫, ৪৫, ৬৩, ১১০, ৩১৫

\textbf{\hl{✔ সঠিক উত্তর: ক) ৩৫, ৪৫, ৬৩, ১০৫, ৩১৫}}

\begin{quote}
\textbf{ব্যাখ্যা:} \emph{৩৪৬ − ৩১ = ৩১৫}

\emph{৩১৫ এর উৎপাদকসমূহ: ১ × ৩১৫ = ৩১৫ ৩ × ১০৫ = ৩১৫ ৫ × ৬৩ = ৩১৫ ৭ × ৪৫
= ৩১৫ ৯ × ৩৫ = ৩১৫}

\emph{নির্ণেয় সংখ্যাগুলো অবশ্যই ভাগশেষ ৩১ অপেক্ষা বড় হতে হবে।}

\emph{∴ ৩১ অপেক্ষা বড় উৎপাদকগুলো হলো: ৩৫, ৪৫, ৬৩, ১০৫, ৩১৫।}
\end{quote}

\textbf{কোন লঘিষ্ঠ সংখ্যার সহিত ৩ যোগ করিলে যোগফল ২৪, ৩৬ এবং ৪৮ দ্বারা
বিভাজ্য হবে?}

ক) ৮৯

\textbf{\hl{খ) ১৪১ ✔}}

গ) ২৪৮

ঘ) ১৭০

\textbf{\hl{✔ সঠিক উত্তর: খ) ১৪১}}

\begin{quote}
\textbf{ব্যাখ্যা:} \emph{২৪, ৩৬ ও ৪৮ এর ল.সা.গু = ১৪৪।}

\emph{লঘিষ্ঠ সংখ্যা বের করার জন্য ল.সা.গু নির্ণয় করতে হয়। যেহেতু সংখ্যাটির সাথে
৩ যোগ করলে যোগফল বিভাজ্য হবে,}

\emph{⇒ নির্ণেয় লঘিষ্ঠ সংখ্যাটি হবে = ১৪৪ − ৩ = ১৪১।}
\end{quote}

\textbf{কোন ক্ষুদ্রতম সংখ্যাকে ৫, ১০, ১৫ দ্বারা ভাগ করলে ৪ অবশিষ্ট থাকে?}

ক) ৬৪

খ) ১২৪

\textbf{\hl{গ) ৩৪ ✔}}

ঘ) কোনোটিই নয়

\textbf{\hl{✔ সঠিক উত্তর: গ) ৩৪}}

\begin{quote}
\textbf{ব্যাখ্যা:} \emph{৫, ১০ ও ১৫ এর ল.সা.গু = ৩০।}

\emph{যেহেতু ৪ অবশিষ্ট থাকে বলা আছে,}

\emph{⇒ সংখ্যাটি হবে = ৩০ + ৪ = ৩৪।}

\emph{∴ নির্ণেয় ক্ষুদ্রতম সংখ্যা ৩৪।}
\end{quote}

\textbf{কোন লঘিষ্ঠ সংখ্যার সাথে ২ যোগ করলে যোগফল ১২, ১৮ ও ২৪ দ্বারা বিভাজ্য
হবে?}

ক) ৮৯

\textbf{\hl{খ) ৭০ ✔}}

গ) ১৭০০

ঘ) ১৪২

\textbf{\hl{✔ সঠিক উত্তর: খ) ৭০}}

\begin{quote}
\textbf{ব্যাখ্যা:} \emph{১২, ১৮ ও ২৪ এর ল.সা.গু = ৭২।}

\emph{যেহেতু সংখ্যাটির সাথে ২ যোগ করলে যোগফল বিভাজ্য হবে,}

\emph{∴ লঘিষ্ঠ সংখ্যাটি হবে = ৭২ − ২ = ৭০।}
\end{quote}

\textbf{কতকগুলো ঘন্টা একসাথে বাজার ১০ সেকেন্ড, ১৫ সেকেন্ড, ২০ সেকেন্ড এবং ২৫
সেকেন্ড পরপর বাজতে লাগল। উহারা আবার কতক্ষণ পরে একত্রে বাজবে?}

ক) ১ মিনিট ২০ সে.

খ) ১ মিনিট ৩০ সে.

গ) ৩ মিনিট

\textbf{\hl{ঘ) ৫ মিনিট ✔}}

\textbf{\hl{✔ সঠিক উত্তর: ঘ) ৫ মিনিট}}

\begin{quote}
\textbf{ব্যাখ্যা:} \emph{ঘন্টাগুলো পুনরায় একত্রে বাজার সময় পেতে হলে ১০, ১৫, ২০
এবং ২৫ এর ল.সা.গু নির্ণয় করতে হবে।}

\emph{১০, ১৫, ২০ এবং ২৫ এর ল.সা.গু = ৩০০}

\emph{অর্থাৎ ৩০০ সেকেন্ড পরে ঘন্টাগুলো পুনরায় একত্রে বাজবে।}

\emph{এখন, ৩০০ সেকেন্ড = (৩০০/৬০) মিনিট = ৫ মিনিট।}

\emph{∴ সঠিক উত্তর অপশন (ঘ) ৫ মিনিট।}
\end{quote}

\textbf{একটি সংখ্যাকে ৩৫ দিয়ে ভাগ করলে ভাগশেষ ১০ থাকে। যদি ঐ সংখ্যাকে ৭
দিয়ে ভাগ করা হয় তবে ভাগশেষ কত হবে?}

ক) ১০

খ) ২

\textbf{\hl{গ) ৩ ✔}}

ঘ) কোনোটিই নয়

\textbf{\hl{✔ সঠিক উত্তর: গ) ৩}}

\begin{quote}
\textbf{ব্যাখ্যা:} \emph{একটি সংখ্যাকে ৩৫ দিয়ে ভাগ করলে যদি ১০ অবশিষ্ট থাকে,
তাহলে সংখ্যাটি ন্যূনতম = (৩৫ + ১০) = ৪৫।}

\emph{এখন, ৪৫ কে ৭ দিয়ে ভাগ করলে (৪৫ = ৬ × ৭ + ৩), ভাগশেষ হয় ৩।}

\emph{∴ ভাগশেষ হবে ৩।}
\end{quote}

\textbf{কতজন শিশুর মধ্যে কোন ফল না ভেঙ্গে ১১৫টি কমলা এবং ১৩৫টি কলা ভাগ করে
দেয়া যায়?}

\textbf{\hl{ক) ৫ ✔}}

খ) ১০

গ) ১২

ঘ) ১৫

\textbf{\hl{✔ সঠিক উত্তর: ক) ৫}}

\begin{quote}
\textbf{ব্যাখ্যা:} \emph{নির্ণেয় শিশুর সংখ্যা হবে ১১৫ এবং ১৩৫ এর গ.সা.গু।}

\emph{১১৫ এর গুণনীয়ক = ১, ৫, ২৩, ১১৫}

\emph{১৩৫ এর গুণনীয়ক = ১, ৩, ৫, ৯, ১৫, ২৭, ৪৫, ১৩৫}

\emph{∴ ১১৫ ও ১৩৫ এর গ.সা.গু. = ৫}

\emph{∴ শিশুর সংখ্যা ৫ জন।}
\end{quote}

\textbf{কতজন বালককে ১২৫টি কমলালেবু এবং ১৪৫টি কলা সমানভাবে ভাগ করে দেয়া
যায়?/ What is the number of boys among whom 125 oranges and 145 bananas
can be equally distributed?}

ক) ২৫ জনকে

খ) ১৫ জনকে

গ) ৩৫ জনকে

\textbf{\hl{ঘ) ৫ জনকে ✔}}

\textbf{\hl{✔ সঠিক উত্তর: ঘ) ৫ জনকে}}

\begin{quote}
\textbf{ব্যাখ্যা:} \emph{১২৫ এবং ১৪৫ এর গ.সা.গু.-ই হবে নির্ণেয় বালকের সংখ্যা।
১২৫ = ৫ × ৫ × ৫ ১৪৫ = ৫ × ২৯}

\emph{∴ ১২৫ ও ১৪৫ এর গ.সা.গু. = ৫}

\emph{∴ নির্ণেয় বালকের সংখ্যা ৫ জন।}
\end{quote}

\textbf{দুইটি সংখ্যার অনুপাত ৭ : ৫ এবং তাদের ল.সা.গু ১৪০ হলে সংখ্যা দুইটির
গ.সা.গু কত?}

\textbf{\hl{ক) ৪ ✔}}

খ) ১২

গ) ৬

ঘ) ৯

\textbf{\hl{✔ সঠিক উত্তর: ক) ৪}}

\begin{quote}
\textbf{ব্যাখ্যা:} \emph{ধরি, সংখ্যা দুইটি ৭x ও ৫x যেখানে x হলো সংখ্যা দুইটির
গ.সা.গু।}

\emph{আমরা জানি, সংখ্যাদ্বয়ের গুণফল = ল.সা.গু × গ.সা.গু}

\emph{∴ ৭x × ৫x = ১৪০ × x}

\emph{⇒ ৩৫x = ১৪০}

\emph{⇒ x = (১৪০/৩৫) = ৪}

\emph{∴ সংখ্যা দুইটির গ.সা.গু ৪।}
\end{quote}

\textbf{দুটি সংখ্যার অনুপাত ২:৩ এবং গ.সা.গু ৪ হলে বৃহত্তর সংখ্যাটি কত?}

ক) ৬

\textbf{\hl{খ) ১২ ✔}}

গ) ৮

ঘ) ১৬

\textbf{\hl{✔ সঠিক উত্তর: খ) ১২}}

\begin{quote}
\textbf{ব্যাখ্যা:} \emph{অনুপাতকে গ.সা.গু. দিয়ে গুণ করলেই সংখ্যাগুলো পাওয়া
যায়।}

\emph{ছোট সংখ্যাটি = ২ × ৪ = ৮}

\emph{এবং বড় সংখ্যাটি = ৩ × ৪ = ১২}
\end{quote}

\textbf{প্রথম ও দ্বিতীয় সংখ্যার গুণফল ৩৫ এবং দ্বিতীয় ও তৃতীয় সংখ্যার গুণফল ৬৩।
দ্বিতীয় সংখ্যাটি কত?}

ক) ৫

খ) ৬

\textbf{\hl{গ) ৭ ✔}}

ঘ) ৮

\textbf{\hl{✔ সঠিক উত্তর: গ) ৭}}

\begin{quote}
\textbf{ব্যাখ্যা:} \emph{২য় সংখ্যাটি হবে ৩৫ ও ৬৩ এর গ.সা.গু। ৩৫ = ৫ × ৭ ৬৩ =
৭ × ৯}

\emph{৩৫ ও ৬৩ এর সাধারণ মৌলিক গুণনীয়ক = ৭}

\emph{∴ নির্ণেয় সংখ্যা = ৭}
\end{quote}

\textbf{নিচের কোন পূর্ণ সংখ্যাটিকে ৩, ৪, ৫ এবং ৬ দ্বারা ভাগ করলে যথাক্রমে ১,
২, ৩ ও ৪ অবশিষ্ট থাকে?}

ক) ৪৮

খ) ৫৪

\textbf{\hl{গ) ৫৮ ✔}}

ঘ) ৬০

\textbf{\hl{✔ সঠিক উত্তর: গ) ৫৮}}

\begin{quote}
\textbf{ব্যাখ্যা:} \emph{এখানে,}

\emph{৩ − ১ = ২}

\emph{৪ − ২ = ২}

\emph{৫ − ৩ = ২}

\emph{৬ − ৪ = ২}

\emph{দেখা যাচ্ছে, প্রতিটি ক্ষেত্রে ভাজক ও ভাগশেষের পার্থক্য ২।}

\emph{অর্থাৎ নির্ণেয় সংখ্যাটি হবে ৩, ৪, ৫ এবং ৬ এর ল.সা.গু অপেক্ষা ২ কম।}

\emph{এখানে, ৩, ৪, ৫, ৬ এর ল.সা.গু = ৬০}

\emph{∴ নির্ণেয় সংখ্যা = (৬০ − ২) = ৫৮}
\end{quote}

\textbf{ক্ষুদ্রতম কোন সংখ্যাকে ১৬, ২৪ এবং ৩৬ দ্বারা ভাগ করলে যথাক্রমে ৬, ১৪,
২৬ ভাগশেষ থাকবে?}

ক) ১৪৪

\textbf{\hl{খ) ১৩৪ ✔}}

গ) ১৫৪

ঘ) ১৬৪

\textbf{\hl{✔ সঠিক উত্তর: খ) ১৩৪}}

\begin{quote}
\textbf{ব্যাখ্যা:} \emph{১৬ − ৬ = ১০}

\emph{২৪ − ১৪ = ১০}

\emph{৩৬ − ২৬ = ১০}

\emph{যেহেতু প্রতিটি ক্ষেত্রে পার্থক্য সমান (১০), তাই নির্ণেয় সংখ্যাটি হবে ১৬,
২৪ এবং ৩৬ এর ল.সা.গু. অপেক্ষা ১০ কম।}

\emph{১৬, ২৪ এবং ৩৬ এর ল.সা.গু. = ১৪৪}

\emph{∴ নির্ণেয় ক্ষুদ্রতম সংখ্যাটি = ১৪৪ − ১০ = ১৩৪।}
\end{quote}

\textbf{কোন বৃহত্তম সংখ্যা দ্বারা ২৭, ৪০ ও ৬৫ কে ভাগ করলে যথাক্রমে ৩, ৪ ও ৫
ভাগশেষ থাকবে?}

ক) ১৬

খ) ১৪

\textbf{\hl{গ) ১২ ✔}}

ঘ) ১০

\textbf{\hl{✔ সঠিক উত্তর: গ) ১২}}

\begin{quote}
\textbf{ব্যাখ্যা:} \emph{প্রদত্ত সংখ্যাগুলো থেকে যথাক্রমে ভাগশেষগুলো বিয়োগ করি:}

\emph{২৭ − ৩ = ২৪}

\emph{৪০ − ৪ = ৩৬}

\emph{৬৫ − ৫ = ৬০}

\emph{এখন ২৪, ৩৬ ও ৬০ এর গ.সা.গু. ই হবে নির্ণেয় বৃহত্তম সংখ্যা।}

\emph{২৪, ৩৬ ও ৬০ এর গ.সা.গু. = ১২}

\emph{∴ নির্ণেয় বৃহত্তম সংখ্যা = ১২।}
\end{quote}

\textbf{কোন বৃহত্তম সংখ্যা দ্বারা ৫৭, ৯৩ ও ১৮৩ কে ভাগ করলে ভাগশেষ থাকবে না?}

ক) ৯

খ) ৭

গ) ৫

\textbf{\hl{ঘ) ৩ ✔}}

\textbf{\hl{✔ সঠিক উত্তর: ঘ) ৩}}

\begin{quote}
\textbf{ব্যাখ্যা:} \emph{নির্ণেয় বৃহত্তম সংখ্যাটি হবে ৫৭, ৯৩ ও ১৮৩ এর গ.সা.গু।
৫৭ = ৩ × ১৯ ৯৩ = ৩ × ৩১ ১৮৩ = ৩ × ৬১}

\emph{∴ ৫৭, ৯৩ ও ১৮৩ এর গ.সা.গু. = ৩}

\emph{∴ নির্ণেয় সংখ্যাটি ৩।}
\end{quote}

\textbf{কোন বৃহত্তম সংখ্যা দ্বারা ১০০ ও ১৮৪ কে ভাগ করলে প্রত্যেকবার ভাগশেষ ৪
থাকবে?}

ক) ১৪

খ) ১০

\textbf{\hl{গ) ১২ ✔}}

ঘ) ১৬

\textbf{\hl{✔ সঠিক উত্তর: গ) ১২}}

\begin{quote}
\textbf{ব্যাখ্যা:} \emph{প্রদত্ত সংখ্যা দুটি থেকে ভাগশেষ ৪ বিয়োগ করি:}

\emph{১০০ − ৪ = ৯৬}

\emph{১৮৪ − ৪ = ১৮০}

\emph{৯৬ ও ১৮০ এর গ.সা.গু-ই হবে নির্ণেয় বৃহত্তম সংখ্যা।}

\emph{১৮০ কে ৯৬ দ্বারা ভাগ প্রক্রিয়ায় গ.সা.গু নির্ণয় করি:}

\emph{১৮০ ÷ ৯৬ = ১ (ভাগশেষ ৮৪)}

\emph{৯৬ ÷ ৮৪ = ১ (ভাগশেষ ১২)}

\emph{৮৪ ÷ ১২ = ৭ (ভাগশেষ ০)}

\emph{∴ ৯৬ ও ১৮০ এর গ.সা.গু. = ১২}

\emph{∴ নির্ণেয় বৃহত্তম সংখ্যা ১২।}
\end{quote}

\textbf{১৫৯৬৮ কে কোন নির্দিষ্ট সংখ্যা দ্বারা ভাগ করলে ভাগফল ৮৯ ও ভাগশেষ ৩৭
থাকে?}

ক) ১৫৭

খ) ১৬৮

\textbf{\hl{গ) ১৭৯ ✔}}

ঘ) ১৮৩

\textbf{\hl{✔ সঠিক উত্তর: গ) ১৭৯}}

\begin{quote}
\textbf{ব্যাখ্যা:} \emph{আমরা জানি, ভাজ্য = (ভাজক × ভাগফল) + ভাগশেষ}

\emph{⇒ ভাজক = (ভাজ্য−ভাগশেষ/ভাগফল)}

\emph{এখানে, ভাজ্য = ১৫৯৬৮, ভাগফল = ৮৯, ভাগশেষ = ৩৭}

\emph{∴ ভাজক = (১৫৯৬৮−৩৭/৮৯)}

\emph{= (১৫৯৩১/৮৯) = ১৭৯}

\emph{∴ নির্ণেয় নির্দিষ্ট সংখ্যাটি ১৭৯।}
\end{quote}

\textbf{কোনো বৃহত্তম সংখ্যা দ্বারা ৩৬৫ ও ৪৬৩ কে ভাগ করলে ভাগশেষ যথাক্রমে ৫ ও
৭ থাকে। বৃহত্তম সংখ্যাটির মান কত?}

\textbf{\hl{ক) ২৪ ✔}}

খ) ৫২

গ) ৪৫

ঘ) ৩৪

\textbf{\hl{✔ সঠিক উত্তর: ক) ২৪}}

\begin{quote}
\textbf{ব্যাখ্যা:} \emph{৩৬৫ − ৫ = ৩৬০}

\emph{৪৬৩ − ৭ = ৪৫৬}

\emph{এখন ৩৬০ ও ৪৫৬ এর গ.সা.গু. নির্ণয় করি:}

\emph{৪৫৬ ÷ ৩৬০ = ১ (ভাগশেষ ৯৬)}

\emph{৩৬০ ÷ ৯৬ = ৩ (ভাগশেষ ৭২)}

\emph{৯৬ ÷ ৭২ = ১ (ভাগশেষ ২৪)}

\emph{৭২ ÷ ২৪ = ৩ (ভাগশেষ ০)}

\emph{∴ ৩৬০ ও ৪৫৬ এর গ.সা.গু. = ২৪}

\emph{∴ বৃহত্তম সংখ্যাটি হলো ২৪।}
\end{quote}

\textbf{কত জনের মধ্যে 125টি কমলা ও 145টি কলা সমানভাবে ভাগ করে দেয়া যায়?}

ক) 25 জন

খ) 15 জন

গ) 35 জন

\textbf{\hl{ঘ) 5 জন ✔}}

\textbf{\hl{✔ সঠিক উত্তর: ঘ) 5 জন}}

\begin{quote}
\textbf{ব্যাখ্যা:} \emph{125 ও 145 এর গ.সা.গু. ই নির্ণেয় সংখ্যা।}

\emph{145 ÷ 125 = 1 (ভাগশেষ 20)}

\emph{125 ÷ 20 = 6 (ভাগশেষ 5)}

\emph{20 ÷ 5 = 4 (ভাগশেষ 0)}

\emph{∴ 125 ও 145 এর গ.সা.গু. = 5}

\emph{∴ নির্ণেয় সংখ্যা 5 জন।}
\end{quote}

\textbf{কোনো বিক্রেতাকে ৩.২৫ টাকা, ৪.৭৫ টাকা ও ১১.৫০ টাকা একই ধরনের মুদ্রা
দ্বারা পরিশোধ করতে হলে সবচেয়ে বড় কত পয়সার মুদ্রা প্রয়োজন।}

ক) ১০

খ) ১৫

গ) ৫০

\textbf{\hl{ঘ) ২৫ ✔}}

ঙ) ২০

\textbf{\hl{✔ সঠিক উত্তর: ঘ) ২৫}}

\begin{quote}
\textbf{ব্যাখ্যা:} \emph{টাকাকে পয়সায় রূপান্তর করি:}

\emph{৩.২৫ টাকা = ৩২৫ পয়সা}

\emph{৪.৭৫ টাকা = ৪৭৫ পয়সা}

\emph{১১.৫০ টাকা = ১১৫০ পয়সা}

\emph{এখন ৩২৫, ৪৭৫ ও ১১৫০ এর গ.সা.গু. নির্ণয় করি:}

\emph{৪৭৫ ÷ ৩২৫ = ১ (ভাগশেষ ১৫০)}

\emph{৩২৫ ÷ ১৫০ = ২ (ভাগশেষ ২৫)}

\emph{১৫০ ÷ ২৫ = ৬ (ভাগশেষ ০)}

\emph{আবার, ১১৫০ ÷ ২৫ = ৪৬ (ভাগশেষ ০)}

\emph{∴ ৩২৫, ৪৭৫ ও ১১৫০ এর গ.সা.গু. = ২৫}

\emph{∴ সবচেয়ে বড় ২৫ পয়সার মুদ্রা প্রয়োজন।}
\end{quote}

\textbf{(১/৪) , (৩/১৬) , (৯/২০) এর সাধারণ গুণিতক নিচের কোনটি?}

ক) (১/৪)

খ) (৪/৯)

গ) (৯/২০)

\textbf{\hl{ঘ) (৯/৪) ✔}}

\textbf{\hl{✔ সঠিক উত্তর: ঘ) (৯/৪)}}

\begin{quote}
\textbf{ব্যাখ্যা:} \emph{ভগ্নাংশগুলোর লঘিষ্ঠ সাধারণ গুণিতক (ল.সা.গু)-ই হবে এদের
একটি সাধারণ গুণিতক।}

\emph{আমরা জানি, ভগ্নাংশের ল.সা.গু. = (লবগুলোর ল.সা.গু./হরগুলোর গ.সা.গু.)}

\emph{এখানে, লব ১, ৩, ৯ এর ল.সা.গু. = ৯}

\emph{এবং হর ৪, ১৬, ২০ এর গ.সা.গু. = ৪}

\emph{∴ নির্ণেয় সাধারণ গুণিতক (ল.সা.গু.) = (৯/৪)}
\end{quote}

\textbf{কোন বিমান আক্রমণের সময় এক শহরের চারটি স্থান থেকে যথাক্রমে ১, ১
(১/৪) , ১ (১/২) , ১ (৩/৪) মিনিট অন্তর সাইরেন বাজতে লাগলো। একবার একত্রে
বাজবার কতক্ষণ পর সাইরেনগুলো আবার একত্রে বাজবে?}

ক) ১ ঘ. ৫ মি.

খ) ১ ঘ. ২৫ মি.

গ) ৫৭ মি.

\textbf{\hl{ঘ) ১ ঘ. ৪৫ মি. ✔}}

\textbf{\hl{✔ সঠিক উত্তর: ঘ) ১ ঘ. ৪৫ মি.}}

\begin{quote}
\textbf{ব্যাখ্যা:} \emph{নির্ণেয় সময় হবে ১, ১ (১/৪) , ১ (১/২) , ১ (৩/৪) এর
ল.সা.গু এর সমান।}

\emph{এখানে, ১ = (১/১) ; ১ (১/৪) = (৫/৪) ; ১ (১/২) = (৩/২) ; ১ (৩/৪) =
(৭/৪)}

\emph{লবগুলো (১, ৩, ৫, ৭) এর ল.সা.গু = ১০৫}

\emph{হরগুলোর (১, ৪, ২, ৪) এর গ.সা.গু = ১}

\emph{∴ ভগ্নাংশগুলোর ল.সা.গু = (১০৫/১) = ১০৫ মিনিট}

\emph{সাইরেনগুলো আবার একত্রে বাজবে = ১০৫ মিনিট পর = ১ ঘণ্টা ৪৫ মিনিট পর।}
\end{quote}

\textbf{৩, ৫, ৬ ও ৮ এর ল.সা.গু কত?}

ক) ৮০

খ) ১০০

\textbf{\hl{গ) ১২০ ✔}}

ঘ) ১৪০

\textbf{\hl{✔ সঠিক উত্তর: গ) ১২০}}

\begin{quote}
\textbf{ব্যাখ্যা:} \emph{৩, ৫, ৬ ও ৮ এর ল.সা.গু:}

\emph{৩ \textbar{} ৩, ৫, ৬, ৮}

\emph{২ \textbar{} ১, ৫, ২, ৮}

\emph{১, ৫, ১, ৪}

\emph{∴ নির্ণেয় ল.সা.গু. = ৩ × ২ × ৫ × ৪ = ১২০}
\end{quote}

\textbf{(৩/৫) , (১/৪) , (২/৩) এর ল.সা.গু কত?}

ক) (১/৬)

\textbf{\hl{খ) ৬ ✔}}

গ) (১/২)

ঘ) ১২

\textbf{\hl{✔ সঠিক উত্তর: খ) ৬}}

\begin{quote}
\textbf{ব্যাখ্যা:} \emph{আমরা জানি, ভগ্নাংশের ল.সা.গু = (লবগুলোর
ল.সা.গু/হরগুলোর গ.সা.গু)}

\emph{লবগুলোর (৩, ১, ২) ল.সা.গু = ৬}

\emph{হরগুলোর (৫, ৪, ৩) গ.সা.গু = ১}

\emph{∴ নির্ণেয় ল.সা.গু = (৬/১) = ৬}
\end{quote}

\textbf{২১, ২৮, ৮৪ এবং ৫৬ এর ল.সা.গু কত?}

ক) ৩৩৬

খ) ১১২

গ) ২৬৮

\textbf{\hl{ঘ) ১৬৮ ✔}}

\textbf{\hl{✔ সঠিক উত্তর: ঘ) ১৬৮}}

\begin{quote}
\textbf{ব্যাখ্যা:} \emph{প্রদত্ত সংখ্যাগুলোর ল.সা.গু নির্ণয় করি:}

\emph{৭ \textbar{} ২১, ২৮, ৮৪, ৫৬}

\emph{৪ \textbar{} ৩, ৪, ১২, ৮}

\emph{৩ \textbar{} ৩, ১, ৩, ২}

\emph{১, ১, ১, ২}

\emph{∴ নির্ণেয় ল.সা.গু = ৭ × ৪ × ৩ × ২ = ১৬৮}
\end{quote}

\textbf{২০০২ সংখ্যা কোন সংখ্যাগুচ্ছের ল.সা.গু নয় ?}

\textbf{\hl{ক) ১৩, ৭৭, ৯১, ১৪৩ ✔}}

খ) ৭, ২২, ২৬, ৯১

গ) ২৬, ৭৭, ১৪৩, ১৫৪

ঘ) ২, ৭, ১১, ১৩

\textbf{\hl{✔ সঠিক উত্তর: ক) ১৩, ৭৭, ৯১, ১৪৩}}

\begin{quote}
\textbf{ব্যাখ্যা:} \emph{অপশন (ক): ১৩, ৭৭, ৯১, ১৪৩ এর ল.সা.গু নির্ণয় করি:}

\emph{১৩ \textbar{} ১৩, ৭৭, ৯১, ১৪৩}

\emph{১১ \textbar{} ১, ৭৭, ৭, ১১}

\emph{৭ \textbar{} ১, ৭, ৭, ১}

\emph{১, ১, ১, ১}

\emph{∴ ল.সা.গু. = ১৩ × ১১ × ৭ = ১০০১}

\emph{অন্যান্য অপশন (খ, গ, ঘ) এর সংখ্যাগুলোর ল.সা.গু ২০০২।}

\emph{∴ ১৩, ৭৭, ৯১, ১৪৩ সংখ্যাগুচ্ছের ল.সা.গু ২০০২ নয়।}
\end{quote}

\textbf{২৪৫০ সংখ্যাটিকে কত দ্বারা গুণ করলে সংখ্যাটি পূর্ণবর্গ সংখ্যা হবে ?}

\textbf{\hl{ক) ২ ✔}}

খ) ৪

গ) ৫

ঘ) ৩

\textbf{\hl{✔ সঠিক উত্তর: ক) ২}}

\begin{quote}
\textbf{ব্যাখ্যা:} \emph{২৪৫০ কে মৌলিক উৎপাদকে বিশ্লেষণ করি:}

\emph{২ \textbar{} ২৪৫০}

\emph{৫ \textbar{} ১২২৫}

\emph{৫ \textbar{} ২৪৫}

\emph{৭ \textbar{} ৪৯}

\emph{৭}

\emph{∴ ২৪৫০ = ২ × (৫ × ৫) × (৭ × ৭)}

\emph{এখানে, ২ জোড়াবিহীন অবস্থায় রয়েছে।}

\emph{∴ ২৪৫০ এর সাথে ২ গুণ করলে প্রাপ্ত সংখ্যাটি একটি পূর্ণবর্গ সংখ্যা হবে।}
\end{quote}

\textbf{কোন সংখ্যাকে ৪ ও ৬ দ্বারা ভাগ করলে প্রতিক্ষেত্রে ভাগশেষ ২ থাকে?}

ক) ৮

খ) ১০

গ) ১২

\textbf{\hl{ঘ) ১৪ ✔}}

\textbf{\hl{✔ সঠিক উত্তর: ঘ) ১৪}}

\begin{quote}
\textbf{ব্যাখ্যা:} \emph{৪ ও ৬ এর ল.সা.গু = ১২}

\emph{যেহেতু প্রতিক্ষেত্রে ২ ভাগশেষ থাকে, তাই নির্ণেয় সংখ্যাটি হবে ল.সা.গু
অপেক্ষা ২ বেশি।}

\emph{∴ নির্ণেয় সংখ্যা = ১২ + ২ = ১৪}
\end{quote}

\textbf{এমন একটি লঘিষ্ঠ সংখ্যা নির্ণয় কর যাকে ১৫, ১৮, ২১ এবং ২৪ দ্বারা ভাগ
করলে প্রত্যেক ক্ষেত্রে ২ অবশিষ্ট থাকে?}

ক) ৩

খ) ২৫১৮

গ) ২৫২০

\textbf{\hl{ঘ) ২৫২২ ✔}}

\textbf{\hl{✔ সঠিক উত্তর: ঘ) ২৫২২}}

\begin{quote}
\textbf{ব্যাখ্যা:} \emph{১৫, ১৮, ২১ ও ২৪ এর ল.সা.গু. বের করে তার সাথে ২ যোগ
করতে হবে।}

\emph{১৫, ১৮, ২১ ও ২৪ এর ল.সা.গু. = ২৫২০}

\emph{∴ নির্ণেয় লঘিষ্ঠ সংখ্যা = ২৫২০ + ২ = ২৫২২}
\end{quote}

\textbf{দুটি ঘণ্টা যথাক্রমে ৮ মিনিট এবং ৬ মিনিট অন্তর বাজে। দুপুর ১২টায় দুটি
ঘণ্টা একত্রে বাজার পর আবার কখন তারা একত্রে বাজবে?}

ক) ১২টা ১২ মিনিটে

খ) ১২টা ১৮ মিনিটে

\textbf{\hl{গ) ১২টা ২৪ মিনিটে ✔}}

ঘ) ১২টা ৩০ মিনিটে

\textbf{\hl{✔ সঠিক উত্তর: গ) ১২টা ২৪ মিনিটে}}

\begin{quote}
\textbf{ব্যাখ্যা:} \emph{এখানে ৮ এবং ৬ এর ল.সা.গু.-ই হবে ঘণ্টা দুটি পরবর্তীতে
একত্রে বাজার সময়।}

\emph{৮ এবং ৬ এর ল.সা.গু. = ২৪}

\emph{∴ পরবর্তীতে ঘণ্টা দুটি ১২টা ২৪ মিনিটে একত্রে বাজবে।}
\end{quote}

\textbf{২টি ঘণ্টা যথাক্রমে ১০ ও ২৫ মিনিট অন্তর বাজে। একবার একত্রে বাজার পর
আবার কখন ঘণ্টা দু\textquotesingle টি একত্রে বাজবে?}

ক) ২০ মি: পর

খ) ৩০ মি: পর

\textbf{\hl{গ) ৫০ মি: পর ✔}}

ঘ) ১০০ মি: পর

\textbf{\hl{✔ সঠিক উত্তর: গ) ৫০ মি: পর}}

\begin{quote}
\textbf{ব্যাখ্যা:} \emph{১০ ও ২৫ এর ল.সা.গু. = ২ × ৫ × ৫ = ৫০}

\emph{∴ ৫০ মিনিট পর ঘণ্টা দুটি আবার একত্রে বাজবে।}
\end{quote}

\textbf{তিনটি ঘণ্টা একত্রে বাজার পর তারা ২ ঘণ্টা, ৩ ঘণ্টা ও ৪ ঘণ্টা পরপর
বাজতে থাকল। ১দিনে তারা কতবার একত্রে বাজবে?}

\textbf{\hl{ক) ৩ ✔}}

খ) ২

গ) ৫

ঘ) ৪

\textbf{\hl{✔ সঠিক উত্তর: ক) ৩}}

\begin{quote}
\textbf{ব্যাখ্যা:} \emph{২, ৩ এবং ৪ এর ল.সা.গু. = ১২}

\emph{∴ ঘণ্টাগুলো ১২ ঘণ্টা পরপর একত্রে বাজবে।}

\emph{আমরা জানি, ১ দিন = ২৪ ঘণ্টা।}

\emph{প্রথমবারের পর বাজবে = (২৪ ÷ ১২) = ২ বার।}

\emph{প্রথমবারের বাজাসহ ১ দিনে মোট বাজবে = (২ + ১) = ৩ বার।}
\end{quote}

\textbf{তিনটি ঘণ্টা একত্রে বাজার পর তারা যথাক্রমে ২ ঘণ্টা, ৩ ঘণ্টা ও ৪ ঘণ্টা
পরপর বাজতে থাকলো। ১ দিনে তারা কতবার একত্রে বাজবে?}

ক) 12 বার

খ) 6 বার

গ) 4 বার

\textbf{\hl{ঘ) 3 বার ✔}}

\textbf{\hl{✔ সঠিক উত্তর: ঘ) 3 বার}}

\begin{quote}
\textbf{ব্যাখ্যা:} \emph{২, ৩, ও ৪ এর ল. সা. গু = ১২}

\emph{∴ ১২ ঘণ্টা পর তারা একত্রে আবার বাজবে।}

\emph{১ দিন = ২৪ ঘণ্টা।}

\emph{∴ ২৪ ঘণ্টায় তারা বাজবে (২৪/১২) বা ২ বার (শুরুতে একবার বাজার হিসাব
বাদে)।}

\emph{শুরুর বাজাসহ মোট একত্রে বাজবে = ২ + ১ = ৩ বার।}
\end{quote}

\textbf{পাঁচটি ঘণ্টা একত্রে বেজে যথাক্রমে ৩, ৫, ৭, ৮ ও ১০ সেকেন্ড অন্তর অন্তর
বাজতে লাগল। কতক্ষণ পরে ঘণ্টাগুলো পুনরায় একত্রে বাজবে?}

ক) ১০ মিনিট

খ) ৯০ সেকেন্ড

\textbf{\hl{গ) ১৪ মিনিট ✔}}

ঘ) ২৪০ সেকেন্ড

\textbf{\hl{✔ সঠিক উত্তর: গ) ১৪ মিনিট}}

\begin{quote}
\textbf{ব্যাখ্যা:} \emph{৩, ৫, ৭, ৮ এবং ১০ এর ল.সা.গু. নির্ণয় করি:}

\emph{২ \textbar{} ৩, ৫, ৭, ৮, ১০}

\emph{৫ \textbar{} ৩, ৫, ৭, ৪, ৫}

\emph{৩, ১, ৭, ৪, ১}

\emph{∴ ৩, ৫, ৭, ৮ এবং ১০ এর ল.সা.গু. = ২ × ৫ × ৩ × ৭ × ৪ = ৮৪০}

\emph{ঘণ্টাগুলো পুনরায় একত্রে বাজবে = ৮৪০ সেকেন্ড পর = (৮৪০/৬০) মিনিট = ১৪
মিনিট পর।}
\end{quote}

\textbf{কোনো স্কুলের ছাত্রসংখ্যাকে ৫, ৮, ২০ দ্বারা ভাগ করলে প্রত্যেক বারই ৪ জন
ছাত্র অবশিষ্ট থাকে। ঐ স্কুলের ছাত্রসংখ্যা কত?}

ক) ৪০ জন

\textbf{\hl{খ) ৪৪ জন ✔}}

গ) ৬০ জন

ঘ) ৮৩ জন

\textbf{\hl{✔ সঠিক উত্তর: খ) ৪৪ জন}}

\begin{quote}
\textbf{ব্যাখ্যা:} \emph{ছাত্রসংখ্যা হবে ৫, ৮ এবং ২০ এর ল.সা.গু. অপেক্ষা ৪
বেশি।}

\emph{৫, ৮ ও ২০ এর ল.সা.গু.:}

\emph{২ \textbar{} ৫, ৮, ২০}

\emph{২ \textbar{} ৫, ৪, ১০}

\emph{৫ \textbar{} ৫, ২, ৫}

\emph{১, ২, ১}

\emph{∴ ৫, ৮ ও ২০ এর ল.সা.গু. = ২ × ২ × ৫ × ২ = ৪০}

\emph{∴ নির্ণেয় ছাত্র সংখ্যা = ৪০ + ৪ = ৪৪}
\end{quote}

\textbf{The least number by which 1470 must be divided to get a number
which is a perfect square, is:}

ক) 5

খ) 6

গ) 15

\textbf{\hl{ঘ) 30 ✔}}

\textbf{\hl{✔ সঠিক উত্তর: ঘ) 30}}

\begin{quote}
\textbf{ব্যাখ্যা:} \emph{1470 কে মৌলিক উৎপাদকে বিশ্লেষণ করি:}

\emph{2 \textbar{} 1470}

\emph{5 \textbar{} 735}

\emph{3 \textbar{} 147}

\emph{7 \textbar{} 49}

\emph{7}

\emph{∴ 1470 = 2 × 5 × 3 × (7 × 7)}

\emph{এখানে, 2, 5, এবং 3 বিজোড় সংখ্যক (জোড়াবিহীন) অবস্থায় রয়েছে।}

\emph{So, the least number to divide by = 2 × 5 × 3 = 30}
\end{quote}

\textbf{সর্বমোট কত সংখ্যক গাছ হলে একটি বাগানে ৭, ১৪, ২১, ৩৫ ও ৪২ সারিতে
গাছ লাগালে একটিও কম বা বেশি হবে না?}

\textbf{\hl{ক) ২১০ ✔}}

খ) ২২০

গ) ২৩০

ঘ) ২৬০

\textbf{\hl{✔ সঠিক উত্তর: ক) ২১০}}

\begin{quote}
\textbf{ব্যাখ্যা:} \emph{প্রদত্ত সংখ্যাগুলোর ল.সা.গু-ই হবে মোট গাছের সংখ্যা।}

\emph{২ \textbar{} ৭, ১৪, ২১, ৩৫, ৪২}

\emph{৩ \textbar{} ৭, ৭, ২১, ৩৫, ২১}

\emph{৭ \textbar{} ৭, ৭, ৭, ৩৫, ৭}

\emph{১, ১, ১, ৫, ১}

\emph{∴ নির্ণেয় ল.সা.গু = ২ × ৩ × ৭ × ৫ = ২১০}

\emph{∴ সর্বমোট গাছের সংখ্যা ২১০টি।}
\end{quote}

\textbf{পাঁচ অঙ্কের বৃহত্তম সংখ্যার সাথে কত ক্ষুদ্রতম সংখ্যা যোগ করলে যোগফল 5, 8,
12, 14 দ্বারা বিভাজ্য?}

ক) 701

\textbf{\hl{খ) 801 ✔}}

গ) 815

ঘ) 709

\textbf{\hl{✔ সঠিক উত্তর: খ) 801}}

\begin{quote}
\textbf{ব্যাখ্যা:} \emph{পাঁচ অঙ্কের বৃহত্তম সংখ্যা = ৯৯৯৯৯}

\emph{৫, ৮, ১২, ১৪ এর ল.সা.গু. নির্ণয় করি:}

\emph{২ \textbar{} ৫, ৮, ১২, ১৪}

\emph{২ \textbar{} ৫, ৪, ৬, ৭}

\emph{৫, ২, ৩, ৭}

\emph{∴ ল.সা.গু. = ২ × ২ × ৫ × ২ × ৩ × ৭ = ৮৪০}

\emph{এখন, ৯৯৯৯৯ কে ৮৪০ দ্বারা ভাগ করলে, ৯৯৯৯৯ = ৮৪০ × ১১৯ + ৩৯}

\emph{এখানে ভাগশেষ থাকে ৩৯।}

\emph{∴ নির্ণেয় ক্ষুদ্রতম সংখ্যা = ৮৪০ − ৩৯ = ৮০১}
\end{quote}

\textbf{৯৯৯৯৯ এর সঙ্গে কোন ক্ষুদ্রতম সংখ্যা যোগ করলে যোগফল ২, ৩, ৪, ৫ এবং ৬
দ্বারা নিঃশেষে বিভাজ্য হবে?}

\textbf{\hl{ক) ২১ ✔}}

খ) ২৯

গ) ৩৩

ঘ) ৩৯

\textbf{\hl{✔ সঠিক উত্তর: ক) ২১}}

\begin{quote}
\textbf{ব্যাখ্যা:} \emph{২, ৩, ৪, ৫ এবং ৬ এর ল.সা.গু. নির্ণয় করি:}

\emph{২ \textbar{} ২, ৩, ৪, ৫, ৬}

\emph{৩ \textbar{} ১, ৩, ২, ৫, ৩}

\emph{১, ১, ২, ৫, ১}

\emph{∴ ল.সা.গু. = ২ × ৩ × ২ × ৫ = ৬০}

\emph{এখন, ৯৯৯৯৯ কে ৬০ দ্বারা ভাগ করলে, ৯৯৯৯৯ = ৬০ × ১৬৬৬ + ৩৯}

\emph{এখানে ভাগশেষ থাকে ৩৯।}

\emph{∴ নির্ণেয় যোগ করার মতো ক্ষুদ্রতম সংখ্যা = ৬০ − ৩৯ = ২১}
\end{quote}

\textbf{কোন ক্ষুদ্রতম সংখ্যাকে ৩, ৪ ও ৫ দ্বারা ভাগ করলে নিঃশেষে বিভাজ্য?}

ক) ১৬০

খ) ৯০

গ) ১২০

\textbf{\hl{ঘ) ৬০ ✔}}

\textbf{\hl{✔ সঠিক উত্তর: ঘ) ৬০}}

\begin{quote}
\textbf{ব্যাখ্যা:} \emph{৩, ৪ ও ৫ এর ল.সা.গু. = ৩ × ৪ × ৫ = ৬০}

\emph{∴ নির্ণেয় সংখ্যা ৬০।}
\end{quote}

\textbf{নিচের কোন ক্ষুদ্রতম সংখ্যাকে ৩, ৫ ও ৬ দ্বারা ভাগ করলে ভাগশেষ ১ হবে?}

ক) ৭১

খ) ৪১

\textbf{\hl{গ) ৩১ ✔}}

ঘ) ৩৯

\textbf{\hl{✔ সঠিক উত্তর: গ) ৩১}}

\begin{quote}
\textbf{ব্যাখ্যা:} \emph{৩, ৫ ও ৬ এর ল.সা.গু. = ৩০}

\emph{∴ নির্ণেয় সংখ্যাটি = ৩০ + ১ = ৩১}
\end{quote}

\textbf{এক স্কুলে ড্রিল করার সময় ছাত্রদের ৮, ১০ বা ১২ সারিতে সাজানো হয়। ঐ
স্কুলে ন্যূনতম কতজন ছাত্র রয়েছে।}

ক) ৮০

খ) ৯৬

\textbf{\hl{গ) ১২০ ✔}}

ঘ) ১৪০

\textbf{\hl{✔ সঠিক উত্তর: গ) ১২০}}

\begin{quote}
\textbf{ব্যাখ্যা:} \emph{৮, ১০ বা ১২ এর ল.সা.গু.:}

\emph{২ \textbar{} ৮, ১০, ১২}

\emph{২ \textbar{} ৪, ৫, ৬}

\emph{২, ৫, ৩}

\emph{∴ ৮, ১০ বা ১২ এর ল.সা.গু. = ২ × ২ × ২ × ৫ × ৩ = ১২০}

\emph{∴ ঐ স্কুলে ন্যূনতম ১২০ জন ছাত্র রয়েছে।}
\end{quote}

\textbf{কোন ক্ষুদ্রতম সংখ্যা/ লঘিষ্ঠ সংখ্যার সাথে ১ যোগ করলে যোগফল ৩, ৬, ৯,
১২, ১৫ দ্বারা নিঃশেষে বিভাজ্য হবে?}

\textbf{\hl{ক) ১৭৯ ✔}}

খ) ৩৬১

গ) ৩৫৯

ঘ) ৭২১

\textbf{\hl{✔ সঠিক উত্তর: ক) ১৭৯}}

\begin{quote}
\textbf{ব্যাখ্যা:} \emph{৩, ৬, ৯, ১২, ১৫ এর ল.সা.গু. নির্ণয় করি:}

\emph{৩ \textbar{} ৩, ৬, ৯, ১২, ১৫}

\emph{২ \textbar{} ১, ২, ৩, ৪, ৫}

\emph{১, ১, ৩, ২, ৫}

\emph{∴ ৩, ৬, ৯, ১২ এবং ১৫ এর ল.সা.গু. = ৩ × ২ × ৩ × ২ × ৫ = ১৮০}

\emph{নির্ণেয় সংখ্যা = ১৮০ − ১ = ১৭৯}
\end{quote}

\textbf{কোন ক্ষুদ্রতম পূর্ণবর্গ সংখ্যা ৯, ১৫ এবং ২৫ দ্বারা বিভাজ্য?}

ক) ৭৫

\textbf{\hl{খ) ২২৫ ✔}}

গ) ১১২৫

ঘ) ৯০০

\textbf{\hl{✔ সঠিক উত্তর: খ) ২২৫}}

\begin{quote}
\textbf{ব্যাখ্যা:} \emph{৯, ১৫ এবং ২৫ এর ল.সা.গু. নির্ণয় করি:}

\emph{৩ \textbar{} ৯, ১৫, ২৫}

\emph{৫ \textbar{} ৩, ৫, ২৫}

\emph{৩, ১, ৫}

\emph{∴ ল.সা.গু. = (৩ × ৩) × (৫ × ৫) = ২২৫}

\emph{যেহেতু ২২৫ সংখ্যাটি নিজেই একটি পূর্ণবর্গ সংখ্যা {[}১৫\^{}২ = ২২৫{]}, তাই
এটিই নির্ণেয় ক্ষুদ্রতম পূর্ণবর্গ সংখ্যা।}
\end{quote}

\textbf{কোন ক্ষুদ্রতম সংখ্যা থেকে ৫ বিয়োগ করলে বিয়োগফল ৬ এবং ১০ দ্বারা
নিঃশেষে বিভাজ্য হবে?}

ক) ৩০

\textbf{\hl{খ) ৩৫ ✔}}

গ) ৬৫

ঘ) ১২৫

\textbf{\hl{✔ সঠিক উত্তর: খ) ৩৫}}

\begin{quote}
\textbf{ব্যাখ্যা:} \emph{৬ ও ১০ এর ল.সা.গু. = ৩০}

\emph{∴ নির্ণেয় ক্ষুদ্রতম সংখ্যা = ৩০ + ৫ = ৩৫}
\end{quote}

\textbf{কোন ক্ষুদ্রতম সংখ্যা হতে ১ বিয়োগ করলে বিয়োগফল ৯, ১২ ও ১৫ দ্বারা
নিঃশেষে বিভাজ্য হবে?}

ক) ১২১

\textbf{\hl{খ) ১৮১ ✔}}

গ) ২৮১

ঘ) ৩৬১

\textbf{\hl{✔ সঠিক উত্তর: খ) ১৮১}}

\begin{quote}
\textbf{ব্যাখ্যা:} \emph{৯, ১২ এবং ১৫ এর ল.সা.গু. নির্ণয় করি:}

\emph{৩ \textbar{} ৯, ১২, ১৫}

\emph{৩, ৪, ৫}

\emph{∴ ৯, ১২ ও ১৫ এর ল.সা.গু. = ৩ × ৩ × ৪ × ৫ = ১৮০}

\emph{যেহেতু সংখ্যাটি থেকে ১ বিয়োগ করলে বিভাজ্য হয়, তাই নির্ণেয় ক্ষুদ্রতম
সংখ্যাটি হবে ল.সা.গু. অপেক্ষা ১ বেশি।}

\emph{∴ নির্ণেয় সংখ্যাটি = ১৮০ + ১ = ১৮১}
\end{quote}

\textbf{কোন ক্ষুদ্রতম সংখ্যাকে ৪, ৫ ও ৬ দিয়ে ভাগ করলে প্রত্যেকবার ৩ অবশিষ্ট
থাকে?}

ক) ৩৩

খ) ৪৩

গ) ৫৩

\textbf{\hl{ঘ) ৬৩ ✔}}

\textbf{\hl{✔ সঠিক উত্তর: ঘ) ৬৩}}

\begin{quote}
\textbf{ব্যাখ্যা:} \emph{৪, ৫ ও ৬ এর ল.সা.গু. নির্ণয় করি:}

\emph{২ \textbar{} ৪, ৫, ৬}

\emph{২, ৫, ৩}

\emph{∴ ৪, ৫ ও ৬ এর ল.সা.গু. = ২ × ২ × ৫ × ৩ = ৬০}

\emph{যেহেতু প্রতিক্ষেত্রে ৩ অবশিষ্ট থাকে, তাই নির্ণেয় ক্ষুদ্রতম সংখ্যাটি হবে
ল.সা.গু. অপেক্ষা ৩ বেশি।}

\emph{∴ নির্ণেয় ক্ষুদ্রতম সংখ্যা = ৬০ + ৩ = ৬৩}
\end{quote}

\textbf{কোন লঘিষ্ঠ সংখ্যা ১২ ও ১৬ দ্বারা ভাগ করলে অবশিষ্ট যথাক্রমে ৫ ও ৯
হবে?}

ক) ৫৩

খ) ২৯

গ) ৮৮

\textbf{\hl{ঘ) ৪১ ✔}}

\textbf{\hl{✔ সঠিক উত্তর: ঘ) ৪১}}

\begin{quote}
\textbf{ব্যাখ্যা:} \emph{এখানে,}

\emph{১২ − ৫ = ৭}

\emph{১৬ − ৯ = ৭}

\emph{দেখা যাচ্ছে, প্রতিটি ক্ষেত্রে ভাজক ও ভাগশেষের পার্থক্য সমান (৭)।}

\emph{এখন ১২ ও ১৬ এর ল.সা.গু. নির্ণয় করি:}

\emph{২ \textbar{} ১২, ১৬}

\emph{২ \textbar{} ৬, ৮}

\emph{৩, ৪}

\emph{∴ ১২ ও ১৬ এর ল.সা.গু. = ২ × ২ × ৩ × ৪ = ৪৮}

\emph{∴ নির্ণেয় লঘিষ্ঠ সংখ্যা = ৪৮ − ৭ = ৪১}
\end{quote}

\textbf{কোন ক্ষুদ্রতম সংখ্যাকে ২০, ২৫, ৩০, ৩৬ ও ৪৮ দ্বারা ভাগ করলে যথাক্রমে
১৫, ২০, ২৫, ৩১ ও ৪৩ ভাগশেষ থাকে?}

ক) ৩৪২৫

খ) ৩৪৭৮

\textbf{\hl{গ) ৩৫৯৫ ✔}}

ঘ) ৩৫৬৫

\textbf{\hl{✔ সঠিক উত্তর: গ) ৩৫৯৫}}

\begin{quote}
\textbf{ব্যাখ্যা:} \emph{এখানে,}

\emph{২০ − ১৫ = ৫}

\emph{২৫ − ২০ = ৫}

\emph{৩০ − ২৫ = ৫}

\emph{৩৬ − ৩১ = ৫}

\emph{৪৮ − ৪৩ = ৫}

\emph{প্রতিটি ক্ষেত্রে ভাজক ও ভাগশেষের পার্থক্য সমান (৫)।}

\emph{এখন ২০, ২৫, ৩০, ৩৬ ও ৪৮ এর ল.সা.গু. নির্ণয় করি:}

\emph{২ \textbar{} ২০, ২৫, ৩০, ৩৬, ৪৮}

\emph{২ \textbar{} ১০, ২৫, ১৫, ১৮, ২৪}

\emph{৩ \textbar{} ৫, ২৫, ১৫, ৯, ১২}

\emph{৫ \textbar{} ৫, ২৫, ৫, ৩, ৪}

\emph{১, ৫, ১, ৩, ৪}

\emph{∴ ল.সা.গু. = ২ × ২ × ৩ × ৫ × ৫ × ৩ × ৪ = ৩৬০০}

\emph{∴ নির্ণেয় সংখ্যা = ৩৬০০ − ৫ = ৩৫৯৫}
\end{quote}

\textbf{একটি স্কুলে প্যারেড করার সময় ছাত্রদের ১০, ১২ বা ১৬ সারিতে সাজানো হয়।
ঐ স্কুলে ন্যূনতম কতজন ছাত্র আছে?}

ক) ১২০

খ) ১৮০

গ) ২২০

\textbf{\hl{ঘ) ২৪০ ✔}}

\textbf{\hl{✔ সঠিক উত্তর: ঘ) ২৪০}}

\begin{quote}
\textbf{ব্যাখ্যা:} \emph{ন্যূনতম ছাত্র সংখ্যা হবে ১০, ১২ ও ১৬-এর ল.সা.গু।}

\emph{২ \textbar{} ১০, ১২, ১৬}

\emph{২ \textbar{} ৫, ৬, ৮}

\emph{৫, ৩, ৪}

\emph{∴ ১০, ১২ ও ১৬ এর ল.সা.গু. = ২ × ২ × ৫ × ৩ × ৪ = ২৪০}

\emph{∴ ন্যূনতম ছাত্র সংখ্যা ২৪০ জন।}
\end{quote}

\textbf{তিন অঙ্কের ক্ষুদ্রতম সংখ্যা হতে কোন লঘিষ্ঠ সংখ্যা বিয়োগ করলে বিয়োগফল ৫,
১০ ও ১৫ দ্বারা বিভাজ্য হবে?}

ক) ৫

খ) ১৫

\textbf{\hl{গ) ১০ ✔}}

ঘ) ২০

\textbf{\hl{✔ সঠিক উত্তর: গ) ১০}}

\begin{quote}
\textbf{ব্যাখ্যা:} \emph{৫, ১০ এবং ১৫ এর ল.সা.গু. নির্ণয় করি:}

\emph{৫ \textbar{} ৫, ১০, ১৫}

\emph{১, ২, ৩}

\emph{∴ ৫, ১০ এবং ১৫ এর ল.সা.গু. = ৫ × ২ × ৩ = ৩০}

\emph{তিন অঙ্কের ক্ষুদ্রতম সংখ্যা = ১০০}

\emph{১০০ কে ৩০ দ্বারা ভাগ করলে ভাগশেষ থাকে ১০ (১০০ = ৩০ × ৩ + ১০)।}

\emph{∴ নির্ণেয় লঘিষ্ঠ সংখ্যা = ১০।}
\end{quote}

\textbf{৫ অঙ্কের ক্ষুদ্রতম সংখ্যা হতে কোন লঘিষ্ঠ সংখ্যা বিয়োগ করলে বিয়োগফল
৫,১০,১৫ দ্বারা বিভাজ্য হবে?}

ক) ৫

\textbf{\hl{খ) ১০ ✔}}

গ) ১৫

ঘ) ২০

\textbf{\hl{✔ সঠিক উত্তর: খ) ১০}}

\begin{quote}
\textbf{ব্যাখ্যা:} \emph{৫, ১০, ১৫ এর ল.সা.গু. = ৫ × ২ × ৩ = ৩০}

\emph{৫ অঙ্কের ক্ষুদ্রতম সংখ্যা = ১০০০০}

\emph{১০০০০ কে ৩০ দ্বারা ভাগ করলে পাই, ১০০০০ = ৩০ × ৩৩৩ + ১০}

\emph{এখানে ভাগশেষ থাকে ১০।}

\emph{∴ নির্ণেয় লঘিষ্ঠ সংখ্যা = ১০।}
\end{quote}

\textbf{কমপক্ষে যতগুলো ক্রমিক পূর্ণসংখ্যা নিলে তার গুণফল অবশ্যই ৫০৪০ দ্বারা
বিভাজ্য হবে?}

ক) ৮টি

\textbf{\hl{খ) ৭টি ✔}}

গ) ৬টি

ঘ) ৯টি

\textbf{\hl{✔ সঠিক উত্তর: খ) ৭টি}}

\begin{quote}
\textbf{ব্যাখ্যা:} \emph{৫০৪০ কে মৌলিক উৎপাদকে বিশ্লেষণ করি:}

\emph{৫০৪০ = ৭ × ৬ × ৫ × ৪ × ৩ × ২ × ১ (যা ৭-এর ফ্যাক্টোরিয়াল বা 7!)}

\emph{আমরা জানি, n সংখ্যক ক্রমিক পূর্ণসংখ্যার গুণফল সবসময় n! (n factorial)
দ্বারা বিভাজ্য হয়।}

\emph{এখানে, ১ × ২ × ৩ × ৪ × ৫ × ৬ × ৭ = ৫০৪০ (৭টি ক্রমিক পূর্ণসংখ্যার গুণফল
৫০৪০ দ্বারা বিভাজ্য)।}
\end{quote}

\textbf{কোন সেনাবাহিনীতে যদি আরও ১১ জন সদস্য নিয়োগ করা যেত, তাদেরকে ২০,
৪০, ৫০ ও ৬০ সারিতে দাঁড় করানো যেত। ঐ বাহিনীতে সদস্য সংখ্যা কত ছিল?}

ক) ৫৯১

খ) ৫৮১

\textbf{\hl{গ) ৫৮৯ ✔}}

ঘ) ৫৭৯

\textbf{\hl{✔ সঠিক উত্তর: গ) ৫৮৯}}

\begin{quote}
\textbf{ব্যাখ্যা:} \emph{২০, ৪০, ৫০ ও ৬০ এর ল.সা.গু. নির্ণয় করি:}

\emph{৫ \textbar{} ২০, ৪০, ৫০, ৬০}

\emph{২ \textbar{} ৪, ৮, ১০, ১২}

\emph{২ \textbar{} ২, ৪, ৫, ৬}

\emph{১, ২, ৫, ৩}

\emph{∴ ল.সা.গু. = ৫ × ২ × ২ × ২ × ৫ × ৩ = ৬০০}

\emph{যেহেতু ১১ জন সদস্য বেশি নিয়োগ করা হলে সারিতে দাঁড় করানো যেত, তাই পূর্বে
সদস্য সংখ্যা ছিল ল.সা.গু. অপেক্ষা ১১ কম।}

\emph{∴ নির্ণেয় সৈন্য সংখ্যা = ৬০০ − ১১ = ৫৮৯ জন।}
\end{quote}

\textbf{কোন ক্ষুদ্রতম সংখ্যাকে ২, ৩, ৪, ৫, ৬ দ্বারা ভাগ করলে প্রতি ক্ষেত্রে ১
অবশিষ্ট থাকে, কিন্তু ৭ দ্বারা ভাগ করলে কোন অবশিষ্ট থাকে না?}

ক) ৬১

\textbf{\hl{খ) ৩০১ ✔}}

গ) ৩০০

ঘ) ২৮৩

\textbf{\hl{✔ সঠিক উত্তর: খ) ৩০১}}

\begin{quote}
\textbf{ব্যাখ্যা:} \emph{২, ৩, ৪, ৫ এবং ৬ এর ল.সা.গু. নির্ণয় করি:}

\emph{২ \textbar{} ২, ৩, ৪, ৫, ৬}

\emph{৩ \textbar{} ১, ৩, ২, ৫, ৩}

\emph{১, ১, ২, ৫, ১}

\emph{∴ ল.সা.গু. = ২ × ৩ × ২ × ৫ = ৬০}

\emph{নির্ণেয় ক্ষুদ্রতম সংখ্যাটি হবে (৬০ এর গুণিতক + ১) আকৃতির, যা ৭ দ্বারা
বিভাজ্য।}

\emph{৬০ × ১ + ১ = ৬১ (৭ দ্বারা বিভাজ্য নয়)}

\emph{৬০ × ২ + ১ = ১২১ (৭ দ্বারা বিভাজ্য নয়)}

\emph{৬০ × ৩ + ১ = ১৮১ (৭ দ্বারা বিভাজ্য নয়)}

\emph{৬০ × ৪ + ১ = ২৪১ (৭ দ্বারা বিভাজ্য নয়)}

\emph{৬০ × ৫ + ১ = ৩০১ (যা ৭ দ্বারা নিঃশেষে বিভাজ্য, ৩০১ ÷ ৭ = ৪৩)}

\emph{∴ নির্ণেয় ক্ষুদ্রতম সংখ্যা = ৩০১}
\end{quote}

\textbf{৫ এবং ৯৫ এর মধ্যে ৫ ও ৩ দ্বারা বিভাজ্য মোট কয়টি সংখ্যা আছে?}

\textbf{\hl{ক) ৬ ✔}}

খ) ১০

গ) ৭

ঘ) ১৮

\textbf{\hl{✔ সঠিক উত্তর: ক) ৬}}

\begin{quote}
\textbf{ব্যাখ্যা:} \emph{৫ ও ৩ এর ল.সা.গু = ৫ × ৩ = ১৫}

\emph{৫ ও ৩ দ্বারা বিভাজ্য সংখ্যাগুলো হলো ১৫-এর গুণিতক।}

\emph{৫ এবং ৯৫ এর মধ্যে ১৫-এর গুণিতকগুলো হলো: ১৫ × ১ = ১৫ ১৫ × ২ = ৩০ ১৫ ×
৩ = ৪৫ ১৫ × ৪ = ৬০ ১৫ × ৫ = ৭৫ ১৫ × ৬ = ৯০}

\emph{∴ মোট সংখ্যা = ৬টি।}
\end{quote}

\textbf{দুইটি সংখ্যার গ.সা.গু 11 এবং ল.সা.গু 7700। একটি সংখ্যা 275 হলে, অপর
সংখ্যাটি −}

ক) 318

\textbf{\hl{খ) 308 ✔}}

গ) 283

ঘ) 279

\textbf{\hl{✔ সঠিক উত্তর: খ) 308}}

\begin{quote}
\textbf{ব্যাখ্যা:} \emph{আমরা জানি,}

\emph{সংখ্যা দুটির গুণফল = ল.সা.গু. × গ.সা.গু.}

\emph{⇒ অপর সংখ্যাটি = (ল.সা.গু.×গ.সা.গু./একটি সংখ্যা)}

\emph{= (7700×11/275) = 308}

\emph{∴ অপর সংখ্যাটি 308।}
\end{quote}

\textbf{দুটি সংখ্যার ল.সা.গু. ৯০ এবং গ.সা.গু. ১৫। একটি সংখ্যা ৪৫ হলে অপরটি
কত?}

\textbf{\hl{ক) ৩০ ✔}}

খ) ৩৬

গ) ৬০

ঘ) ৭৫

\textbf{\hl{✔ সঠিক উত্তর: ক) ৩০}}

\begin{quote}
\textbf{ব্যাখ্যা:} \emph{আমরা জানি,}

\emph{একটি সংখ্যা × অপর সংখ্যা = ল.সা.গু. × গ.সা.গু.}

\emph{⇒ ৪৫ × অপর সংখ্যা = ৯০ × ১৫}

\emph{⇒ অপর সংখ্যা = (৯০×১৫/৪৫) = ৩০}

\emph{∴ অপর সংখ্যাটি ৩০।}
\end{quote}

\textbf{দুটি সংখ্যার গুণফল ২১৬৬, সংখ্যা দুটির গ.সা.গু. ১৯ হলে, ছোট সংখ্যাটি
কত?}

ক) ৩৮

খ) ৫৭

গ) ৭৬

\textbf{\hl{ঘ) ১৯ ✔}}

\textbf{\hl{✔ সঠিক উত্তর: ঘ) ১৯}}

\begin{quote}
\textbf{ব্যাখ্যা:} \emph{আমরা জানি,}

\emph{সংখ্যা দুটির গুণফল = ল.সা.গু. × গ.সা.গু.}

\emph{⇒ ২১৬৬ = ল.সা.গু. × ১৯}

\emph{⇒ ল.সা.গু. = (২১৬৬/১৯) = ১১৪}

\emph{যেহেতু গ.সা.গু ১৯, তাই সংখ্যা দুটি ১৯ দ্বারা বিভাজ্য। আবার তাদের গুণফল
২১৬৬ এবং ল.সা.গু ১১৪।}

\emph{এখানে সংখ্যা দুটি ১১৪ ও ১৯ হতে পারে (কারণ ১১৪ × ১৯ = ২১৬৬ এবং ১১৪ ও
১৯ এর গ.সা.গু ১৯)।}

\emph{∴ ছোট সংখ্যাটি ১৯।}
\end{quote}

\textbf{দুটি সংখ্যার গুণফল ১,৫৩৬। সংখ্যা দুটির ল.সা.গু ৯৬ হলে তাদের গ.সা.গু
কত?}

\textbf{\hl{ক) ১৬ ✔}}

খ) ১২

গ) ২৪

ঘ) ১৮

\textbf{\hl{✔ সঠিক উত্তর: ক) ১৬}}

\begin{quote}
\textbf{ব্যাখ্যা:} \emph{আমরা জানি,}

\emph{সংখ্যা দুটির গুণফল = ল.সা.গু × গ.সা.গু}

\emph{⇒ ১৫৩৬ = ৯৬ × গ.সা.গু}

\emph{⇒ গ.সা.গু = (১৫৩৬/৯৬) = ১৬}

\emph{∴ গ.সা.গু ১৬।}
\end{quote}

\textbf{দুটি সংখ্যার গ. সা. গু এবং ল.সা.গু যথাক্রমে ২ এবং ৩৬০। একটি সংখ্যা ১০
হলে অপর সংখ্যা কত?}

ক) ২৪

খ) ৪৮

গ) ৬০

\textbf{\hl{ঘ) ৭২ ✔}}

\textbf{\hl{✔ সঠিক উত্তর: ঘ) ৭২}}

\begin{quote}
\textbf{ব্যাখ্যা:} \emph{আমরা জানি,}

\emph{একটি সংখ্যা × অপর সংখ্যা = ল.সা.গু. × গ.সা.গু.}

\emph{⇒ ১০ × অপর সংখ্যা = ৩৬০ × ২}

\emph{⇒ অপর সংখ্যা = (২×৩৬০/১০) = ৭২}

\emph{∴ অপর সংখ্যাটি ৭২।}
\end{quote}

\textbf{দুটি সংখ্যার ল. সা. গু. ৬০ এবং গ. সা. গু. ১০। একটি সংখ্যা অপর সংখ্যার
দুই-তৃতীয়াংশ হলে ছোট সংখ্যাটি কত?}

\textbf{\hl{ক) ২০ ✔}}

খ) ৩০

গ) ১০

ঘ) ৪০

\textbf{\hl{✔ সঠিক উত্তর: ক) ২০}}

\begin{quote}
\textbf{ব্যাখ্যা:} \emph{ধরি, একটি সংখ্যা x}

\emph{∴ অপর সংখ্যাটি = (২x/৩)}

\emph{আমরা জানি,}

\emph{সংখ্যা দুটির গুণফল = ল.সা.গু × গ.সা.গু}

\emph{⇒ x × (২x/৩) = ৬০ × ১০}

\emph{⇒ (২x\^{}২/৩) = ৬০০}

\emph{⇒ ২x\^{}২ = ১৮০০}

\emph{⇒ x\^{}২ = ৯০০}

\emph{⇒ x = ৩০}

\emph{∴ একটি সংখ্যা (বড়টি) = ৩০}

\emph{এবং অপর সংখ্যা (ছোটটি) = (২×৩০/৩) = ২০}

\emph{∴ ছোট সংখ্যাটি ২০।}
\end{quote}

\textbf{দুটি সংখ্যার গ.সা.গু, বিয়োগফল এবং ল.সা.গু. যথাক্রমে ১২, ৬০ এবং ২৪৪৮।
সংখ্যা দুটি কত?}

ক) ১১২,১৪৮

খ) ১০৮,১৪৪

\textbf{\hl{গ) ১৪৪,২০৪ ✔}}

ঘ) ১৪৪,২০৮

\textbf{\hl{✔ সঠিক উত্তর: গ) ১৪৪,২০৪}}

\begin{quote}
\textbf{ব্যাখ্যা:} \emph{ধরি, সংখ্যা দুটি ১২x ও ১২y (যেখানে x ও y সহমৌলিক)}

\emph{প্রশ্নমতে, ১২x − ১২y = ৬০ ⇒ x − y = ৫ ........ (i)}

\emph{এবং ১২xy = ২৪৪৮ ⇒ xy = ২০৪}

\emph{আমরা জানি, (x + y)\^{}২ = (x − y)\^{}২ + ৪xy}

\emph{⇒ (x + y)\^{}২ = (৫)\^{}২ + ৪ × ২০৪ = ২৫ + ৮১৬ = ৮৪১}

\emph{⇒ x + y = √(৮৪১) = ২৯ ........ (ii)}

\emph{(i) ও (ii) যোগ করে পাই, ২x = ৩৪ ⇒ x = ১৭}

\emph{∴ y = ২৯ − ১৭ = ১২}

\emph{∴ সংখ্যা দুটি যথাক্রমে ১২ × ১২ = ১৪৪ এবং ১২ × ১৭ = ২০৪।}
\end{quote}

\textbf{দুটি সংখ্যার গ. সা. গু. এবং ল. সা. গু. যথাক্রমে ১৪ ও ১৬৮। সংখ্যাদ্বয়ের
গুণফল কত?}

ক) ৩৯০০

খ) ১৩৫২

গ) ৪৫৭

\textbf{\hl{ঘ) ২৩৫২ ✔}}

\textbf{\hl{✔ সঠিক উত্তর: ঘ) ২৩৫২}}

\begin{quote}
\textbf{ব্যাখ্যা:} \emph{আমরা জানি,}

\emph{সংখ্যাদ্বয়ের গুণফল = গ. সা. গু. × ল. সা. গু. = ১৪ × ১৬৮ = ২৩৫২}

\emph{∴ সংখ্যাদ্বয়ের গুণফল ২৩৫২।}
\end{quote}

\textbf{দুটি সংখ্যার গুণফল ৩৩৮০ এবং গ.সা.গু ১৩। সংখ্যা দুটির ল.সা.গু কত?}

\textbf{\hl{ক) ২৬০ ✔}}

খ) ৭৮০

গ) ১৩০

ঘ) ৪৯০

\textbf{\hl{✔ সঠিক উত্তর: ক) ২৬০}}

\begin{quote}
\textbf{ব্যাখ্যা:} \emph{আমরা জানি,}

\emph{সংখ্যা দুটির ল.সা.গু. = (দুটি সংখ্যার গুণফল/গ.সা.গু)}

\emph{= (৩৩৮০/১৩) = ২৬০}

\emph{∴ সংখ্যা দুটির ল.সা.গু ২৬০।}
\end{quote}

\textbf{দুটি সংখ্যার গুণফল ১৫৩৬। সংখ্যা দুটির ল.সা.গু. ৯৬ হলে গ. সা. গু কত?}

ক) ৩২

খ) ২৪

\textbf{\hl{গ) ১৬ ✔}}

ঘ) ১২

\textbf{\hl{✔ সঠিক উত্তর: গ) ১৬}}

\begin{quote}
\textbf{ব্যাখ্যা:} \emph{আমরা জানি,}

\emph{গ.সা.গু. = (সংখ্যা দুটির গুণফল/সংখ্যা দুটির ল.সা.গু.)}

\emph{= (১৫৩৬/৯৬) = ১৬}

\emph{∴ নির্ণেয় গ.সা.গু. = ১৬}
\end{quote}

\textbf{দুইটি সংখ্যার গ. সা. গু ১৫ ও ল. সা. গু ১৯২। একটি সংখ্যা ৪৮ হলে, অপর
সংখ্যাটি কত?}

\textbf{\hl{ক) ৬০ ✔}}

খ) ৬২

গ) ৬৪

ঘ) ৬৮

\textbf{\hl{✔ সঠিক উত্তর: ক) ৬০}}

\begin{quote}
\textbf{ব্যাখ্যা:} \emph{আমরা জানি,}

\emph{অপর সংখ্যা = (সংখ্যাদ্বয়ের ল.সা.গু.×গ.সা.গু./একটি সংখ্যা)}

\emph{= (১৯২×১৫/৪৮) = ৬০}

\emph{∴ অপর সংখ্যাটি ৬০।}
\end{quote}

\textbf{৩৩ এর উৎপাদক কোনটি?}

ক) ১০

\textbf{\hl{খ) ৩ ✔}}

গ) ১২

ঘ) ১৫

\textbf{\hl{✔ সঠিক উত্তর: খ) ৩}}

\begin{quote}
\textbf{ব্যাখ্যা:} \emph{৩৩ কে ভেঙ্গে পাই, ৩৩ = ১ × ৩৩ = ৩ × ১১}

\emph{∴ ৩৩ এর উৎপাদকগুলো হচ্ছে: ১, ৩, ১১ ও ৩৩।}

\emph{প্রদত্ত অপশনগুলোর মধ্যে ৩ বিদ্যমান।}
\end{quote}

\textbf{দুটি সংখ্যার অনুপাত ৩ : ৪ এবং তাদের ল.সা.গু. ১৮০। সংখ্যা ২টি কি কি?}

ক) 6, 8

\textbf{\hl{খ) 45, 60 ✔}}

গ) 9, 12

ঘ) 18, 10

\textbf{\hl{✔ সঠিক উত্তর: খ) 45, 60}}

\begin{quote}
\textbf{ব্যাখ্যা:} \emph{ধরি, সংখ্যা দুটি যথাক্রমে ৩x ও ৪x}

\emph{সংখ্যা দুটির ল.সা.গু. = ১২x}

\emph{প্রশ্নমতে, ১২x = ১৮০}

\emph{⇒ x = (১৮০/১২) = ১৫}

\emph{∴ প্রথম সংখ্যাটি = ৩ × ১৫ = ৪৫}

\emph{এবং দ্বিতীয় সংখ্যাটি = ৪ × ১৫ = ৬০}

\emph{∴ সংখ্যা ২টি ৪৫ ও ৬০।}
\end{quote}

\textbf{দুইটি সংখ্যার অনুপাত ৩ : ৪ এবং তাদের ল.সা.গু. ১৮০ হলে, বড় সংখ্যাটি
কত?}

ক) ৪০

খ) ৫০

\textbf{\hl{গ) ৬০ ✔}}

ঘ) ৭০

\textbf{\hl{✔ সঠিক উত্তর: গ) ৬০}}

\begin{quote}
\textbf{ব্যাখ্যা:} \emph{ধরি, সংখ্যা দুটি ৩x ও ৪x}

\emph{∴ শর্তমতে, ৩ × ৪ × x = ১৮০ ⇒ ১২x = ১৮০ ⇒ x = ১৫}

\emph{∴ বড় সংখ্যাটি = ৪x = ৪ × ১৫ = ৬০।}
\end{quote}

\textbf{দুটি সংখ্যার অনুপাত ৫:৭ এবং তাদের ল.সা.গু ৩৫০। সংখ্যা দুটির গ.সা.গু. −}

ক) ৫০

খ) ৭০

গ) ৩৫

\textbf{\hl{ঘ) ১০ ✔}}

\textbf{\hl{✔ সঠিক উত্তর: ঘ) ১০}}

\begin{quote}
\textbf{ব্যাখ্যা:} \emph{ধরি, সংখ্যা দুটি ৫x ও ৭x}

\emph{∴ সংখ্যা দুটির গ.সা.গু. = x এবং ল.সা.গু. = ৩৫x}

\emph{প্রশ্নমতে, ৩৫x = ৩৫০ ⇒ x = ১০}

\emph{∴ সংখ্যা দুটির গ.সা.গু. = ১০।}
\end{quote}

\textbf{দুটি সংখ্যার অনুপাত ৫: ৬ এবং তাদের গ.সা.গু. ৪ হলে, সংখ্যা দুটির
ল.সা.গু. কত?}

ক) ১০০

\textbf{\hl{খ) ১২০ ✔}}

গ) ১৫০

ঘ) ১৮০

\textbf{\hl{✔ সঠিক উত্তর: খ) ১২০}}

\begin{quote}
\textbf{ব্যাখ্যা:} \emph{আমরা জানি,}

\emph{দুটি সংখ্যার ল.সা.গু = সংখ্যা দুটির অনুপাতের গুণফল × গ.সা.গু = (৫ × ৬) × ৪
= ৩০ × ৪ = ১২০}
\end{quote}

\textbf{দুটি সংখ্যার ল.সা.গু ২৪০ এবং গ.সা.গু ২০; এদের একটি সংখ্যা ৭০ এর চেয়ে
বড় হলে অপর সংখ্যাটি কত হবে?}

ক) ৭০

খ) ৬৫

গ) ৫০

\textbf{\hl{ঘ) ৬০ ✔}}

\textbf{\hl{✔ সঠিক উত্তর: ঘ) ৬০}}

\begin{quote}
\textbf{ব্যাখ্যা:} \emph{আমরা জানি, সংখ্যা দুটির গুণফল = ল.সা.গু × গ.সা.গু =
২৪০ × ২০ = ৪৮০০}

\emph{ধরি, সংখ্যা দুটি ২০x ও ২০y (যেখানে x ও y সহমৌলিক)।}

\emph{তাহলে, ২০x × ২০y = ৪৮০০ ⇒ ৪০০xy = ৪৮০০ ⇒ xy = ১২}

\emph{১২ কে সহমৌলিক গুণনীয়কে ভাঙলে পাই (১, ১২) অথবা (৩, ৪)।}

\emph{যদি (x, y) = (৩, ৪) হয়, তবে সংখ্যা দুটি হবে ২০ × ৩ = ৬০ এবং ২০ × ৪ =
৮০।}

\emph{এখানে ৮০ সংখ্যাটি ৭০-এর চেয়ে বড়, ফলে অপর সংখ্যাটি হবে ৬০।}
\end{quote}

\textbf{একটি চাকার পরিধি ৫ মিটার। ৮০ কিমি পথ যেতে চাকাটি কতবার ঘুরবে?}

ক) ১৬

খ) ১৬০

গ) ১৬০০

\textbf{\hl{ঘ) ১৬০০০ ✔}}

\textbf{\hl{✔ সঠিক উত্তর: ঘ) ১৬০০০}}

\begin{quote}
\textbf{ব্যাখ্যা:} \emph{৮০ কিমি = (৮০ × ১০০০) মিটার = ৮০,০০০ মিটার {[}∴ ১
কিমি = ১০০০ মিটার{]}}

\emph{চাকাটি ১ বার ঘুরলে অতিক্রম করে ৫ মিটার।}

\emph{∴ ৮০,০০০ মিটার যেতে চাকাটি ঘুরবে = (৮০,০০০/৫) বার = ১৬,০০০ বার।}
\end{quote}

\textbf{একটি ঘোড়ার গাড়ির সামনের চাকার পরিধি ৪ মিটার, পেছনের চাকার পরিধি
৫ মিটার। গাড়িটি কত পথ গেলে সামনের চাকা পেছনের চাকার চেয়ে ২০০ বার বেশি
ঘুরবে?}

ক) ১.২ কি. মি.

খ) ২.৫ কিমি

\textbf{\hl{গ) ৪ কি.মি ✔}}

ঘ) ৬ কিমি

\textbf{\hl{✔ সঠিক উত্তর: গ) ৪ কি.মি}}

\begin{quote}
\textbf{ব্যাখ্যা:} \emph{সামনের চাকা পেছনের চাকার চেয়ে ১ বার বেশি ঘোরে
সেক্ষেত্রে অতিক্রান্ত দূরত্ব হবে ৪ ও ৫ এর ল.সা.গু. এর সমান।}

\emph{৪ ও ৫ এর ল.সা.গু. = ২০ মিটার}

\emph{∴ ১ বার বেশি ঘুরলে অতিক্রান্ত দূরত্ব = ২০ মিটার}

\emph{∴ ২০০ বার বেশি ঘুরলে অতিক্রান্ত দূরত্ব = (২০ × ২০০) মিটার = ৪০০০ মিটার
= ৪ কিমি। {[}∴ ১০০০ মিটার = ১ কিমি{]}}
\end{quote}

\textbf{একটি গাড়ির সামনের চাকার পরিধি ২ মিটার এবং পেছনের চাকার পরিধি ৩
মিটার। কমপক্ষে কত দূরত্ব অতিক্রম করলে সামনের চাকা পেছনের চাকা অপেক্ষা ১০ বার
বেশি ঘুরবে?}

ক) ২০ মিটার

খ) ৩০ মিটার

গ) ৪০ মিটার

\textbf{\hl{ঘ) ৬০ মিটার ✔}}

\textbf{\hl{✔ সঠিক উত্তর: ঘ) ৬০ মিটার}}

\begin{quote}
\textbf{ব্যাখ্যা:} \emph{সামনের চাকা পেছনের চাকার চেয়ে ১ বার বেশি ঘোরে
সেক্ষেত্রে অতিক্রান্ত দূরত্ব হবে ২ ও ৩ এর ল.সা.গু. এর সমান।}

\emph{২ ও ৩ এর ল.সা.গু. = ৬ মিটার}

\emph{∴ ১ বার বেশি ঘুরলে অতিক্রান্ত দূরত্ব = ৬ মিটার}

\emph{∴ ১০ বার বেশি ঘুরলে অতিক্রান্ত দূরত্ব = (৬ × ১০) মিটার = ৬০ মিটার।}
\end{quote}

\textbf{একটি গাড়ির সামনের চাকার ব্যাস ২৮ সে.মি. এবং পেছনের চাকার ব্যাস ৩৫
সে.মি.। ৮৮ মিটার পথ যেতে সামনের চাকা পেছনের চাকা অপেক্ষা কত পূর্ণ সংখ্যাবার
বেশি ঘুরবে?}

ক) ১০ বার

খ) ১৫ বার

\textbf{\hl{গ) ২০ বার ✔}}

ঘ) ২৫ বার

\textbf{\hl{✔ সঠিক উত্তর: গ) ২০ বার}}

\begin{quote}
\textbf{ব্যাখ্যা:} \emph{সামনের চাকার পরিধি = π × ব্যাস = (২২/৭) × ২৮ = ৮৮
সে.মি.}

\emph{পেছনের চাকার পরিধি = π × ব্যাস = (২২/৭) × ৩৫ = ১১০ সে.মি.}

\emph{মোট দূরত্ব = ৮৮ মিটার = (৮৮ × ১০০) সে.মি. = ৮৮০০ সে.মি.}

\emph{সামনের চাকা ঘুরবে = (৮৮০০/৮৮) = ১০০ বার}

\emph{পেছনের চাকা ঘুরবে = (৮৮০০/১১০) = ৮০ বার}

\emph{∴ সামনের চাকা বেশি ঘুরবে = ( ১০০ − ৮০ ) = ২০ বার।}
\end{quote}

\section{35. উৎপাদকে বিশ্লেষণ-বীজগণিত \textbar\textbar{} এ টু
জেড}\label{ux989ux9ceux9aaux9a6ux995-ux9acux9b6ux9b2ux9b7ux9a3-ux9acux99cux997ux9a3ux9a4-ux98f-ux99f-ux99cux9a1}

বলা হয়ে থাকে \textbf{উৎপাদকে বিশ্লেষণ} হল বীজগণিতের গুরু । এর পিছনের কারণ
হিসেবে বলা যায় বীজগণিতের অন্যান্য অংকগুলোর সমাধান করতে \textbf{উৎপাদকে
বিশ্লেষণ} এর প্রয়োজন হয় । কাজেই এই \textbf{উৎপাদকে বিশ্লেষণ} ভালভাবে জানা
থাকলে অন্যান্য অংকগুলো বুঝতে সহজ হয় । এ জন্যই আমরা সর্বোচ্চ গুরুত্ব দিয়ে এ
অধ্যায়টিকে সাজানোর চেষ্টা করেছি । এ অধ্যায়টির এ টু জেড অত্যন্ত সহজ ও সুচারুরুপে
তুলে ধরার প্রয়াস নিয়ে আপনাদের সামনে এই টিউটরিয়াল নিয়ে হাজির হয়েছি । কাজেই
আপনিও যদি একটু মনোযোগের সহিত এ অধ্যায়টির প্রথম থেকে শেষ পর্যন্ত ভালভাবে পড়তে
পারেন । আশা করা যায় \textbf{উৎপাদকে বিশ্লেষণ} নিয়ে আপনার আর কোন সমস্যা
থাকবে না । যে কোন \textbf{উৎপাদকে বিশ্লেষণ} আপনার সামনে হাজির হলে নিমিষেই
তুড়ি মেরে সমাধান করে দিতে পারবেন । চলুন তাহলে শুরু যাক । প্রথমেই জেনে নেয়া
যাক উৎপাদক কী অথবা কাকে বলে ?

\begin{itemize}
\item
  \textbf{উৎপাদক কী? / উৎপাদক কাকে বলে ?}
\item
  বর্গ এবং ঘন এর সূত্র প্রয়োগ করে উৎপাদক নির্ণয়
\item
  Middle Term দ্বারা উৎপাদক
\item
  ফাংশনের দ্বারা উৎপাদক
\item
  বহুপদী উৎপাদক
\end{itemize}

\subsubsection{উৎপাদক কী? / উৎপাদক কাকে বলে
?}\label{ux989ux9ceux9aaux9a6ux995-ux995-ux989ux9ceux9aaux9a6ux995-ux995ux995-ux9acux9b2}

\textbf{পাটিগণিতের ক্ষেত্রেঃ} কোন সংখ্যাকে যতগুলো সংখ্যা দ্বারা ভাগ করা যায়
তাদেরকে ঐ সংখ্যাটির উৎপাদক বলে, যেমন: ১৫ সংখ্যাটিকে ১৫,৫,৩ এবং ১ দ্বারা ভাগ
করা যায়। তাই ১৫,৫,৩,ও ১ হচ্ছে ১৫ এর উৎপাদক। লক্ষ্যণীয়: প্রদত্ত যে কোন সংখ্যাকে
১ এবং ঐ সংখ্যা দ্বারা সবসময় নিঃশেষে ভাগ করা যাবে, তাই ১৫ এর উৎপাদক ৫ ও ৩ ।
কিন্তু ৬ সংখ্যাটি ১৫ এর উৎপাদক নয় কারণ ৬ দ্বারা ১৫ কে ভাগ করা যায় না।
তেমনিভাবে

\textbf{বীজগণিতের ক্ষেত্রে:} প্রদত্ত রাশিটিকে যে রাশিগুলো দিয়ে ভাগ করা যায়
তাদেরকে উক্ত রাশিটির উৎপাদক বলে। যেমন: \(a^{2}\) - \(b^{2}\) বীজগণিতিয়
রাশিটির দুটি উৎপাদক হচ্ছে (a+b) এবং (a-b), অর্থাৎ (a+b) এবং (a-b) উভয় রাশি
দ্বারা \(a^{2}\) - \(b^{2}\) নিঃশেষে ভাগ করা যায়।

সাধারণ উৎপাদকে বিশ্লেষণ অংকগুলো বিভিন্ন পরীক্ষায় প্রচুর আসে । কিন্তু সমস্যা হলো
অপশনে যে চারটি উত্তর দেয়া থাকে তা প্রায় একই রকম। তাই নিয়ম না জেনে শুধু উত্তর
থেকে অংক মেলানো খুব জটিল। অংকগুলো খুব দ্রুত করার জন্য কিছু টেকনিক অনুসরণ করতে
হবে, যা এই অধ্যায়ে সুন্দরভাবে আলোচনা করা হয়েছে।

উৎপাদকের জন্য প্রয়োজনীয় সূত্রাবলী:

\(a^{2}\) - \(b^{2}\) = (a+b) (a -b)

\(a^{3}\) + \(b^{3}\) = (a+b)(\(a^{2}\) - ab + \(b^{2}\) )

\(a^{3}\) - \(b^{3}\) = (a-b) (\(a^{2}\) + ab + \(b^{2}\) )

\subsubsection{বর্গ এবং ঘন এর সূত্র প্রয়োগ করে উৎপাদক
নির্ণয়}\label{ux9acux9b0ux997-ux98fux9ac-ux998ux9a8-ux98fux9b0-ux9b8ux9a4ux9b0-ux9aaux9b0ux9dfux997-ux995ux9b0-ux989ux9ceux9aaux9a6ux995-ux9a8ux9b0ux9a3ux9df}

\textbf{প্রশ্ন নম্বর- ১. 9}\(x^{2}\) \textbf{-} \({(2x\  - \ 3y)}^{2}\)
\textbf{= ?}

\begin{itemize}
\item
  (ক) (5x - 3y)(x + 3y)
\item
  (খ) (5x - 3y)(x - 3y)
\item
  (গ) (5x + 3y)(x - 3y)
\item
  (ঘ) কোনটিই নয়
\end{itemize}

\textbf{\hl{উত্তরঃ (ক) (5x - 3y)(x + 3y)}}

\textbf{Explanation:}

9\(x^{2}\) -\({(2x\  - \ 3y)}^{2}\)

= \({(3x)}^{2}\) - \({(2x\  - \ 3y)}^{2}\)

= (3x+2x-3y)(3x-2x+3y)

= (5x-3y)(x+3y)

\textbf{প্রশ্ন নম্বর- ২.} \(x^{4}\) \textbf{+} \(x^{2}\) \textbf{+ 1 এর
একটি উৎপাদক} \(x^{2}\) \textbf{+ x + 1 অপর উৎপাদকটি কত?}

\begin{itemize}
\item
  (ক) x\^{}2 + x + 1
\item
  (খ) x\^{}2 - x + 1
\item
  (গ) x\^{}2 + 1
\item
  (ঘ) x + 1
\end{itemize}

\textbf{\hl{উত্তরঃ (খ)}} \(x^{2}\) \textbf{\hl{- x + 1}}

\textbf{Explanation:}

\(x^{4}\) + \(x^{2}\) + 1

= \(x^{2^{2}}\) + 2 . \(x^{2}\).1+ \({(1)}^{2}\) - \(x^{2}\)

= \({(x2\  + \ 1)}^{2}\)

-

\(x^{2}\)

= ( \(x^{2}\) + x + 1 ) ( \(x^{2}\) - x + 1 )

\textbf{প্রশ্ন নম্বর- ৩. 4}\(x^{}\) \textbf{-25}\(x^{2}\) \textbf{+ 36 = ?}

\begin{itemize}
\item
  (ক) (x + 2)(x - 2)(2x + 3)(2x - 3)
\item
  (খ) (2x + 1)(2x - 1)(2 + 6)(x - 6)
\item
  (গ) (x + 3)(x - 3)(2x + 2)(2x - 2)
\item
  (ঘ) (x + 2)(2x - 3)(2x + 3)(3x + 2)
\end{itemize}

\textbf{\hl{উত্তরঃ (ক) (x + 2)(x - 2)(2x + 3)(2x - 3)}}

\textbf{Explanation:}

4\(x^{}\) -25\(x^{2}\) + 36

= \({(2x2)}^{2}\)-2.2\(x^{2}\).6+ \({(6)}^{2}\) - \(x^{2}\)

= \({(2x2\  - \ 6)}^{2}\) - \(x^{2}\)

= (2\(x^{2}\)+x-6)(2\(x^{2}\)-x-6)

= (x + 2)(x - 2)(2x + 3)(2x - 3)

\textbf{প্রশ্ন নম্বর- ৪.} \(a^{4}\) \textbf{+ 4 এর উৎপাদক কি কি?}

\begin{itemize}
\item
  (ক) ( a\^{}2 + 2a + 2)( a\^{}2 + 2a - 2)
\item
  (খ) ( a\^{}2 + 2a + 2)( a\^{}2 - 2a + 2)
\item
  (গ) ( a\^{}2 - 2a - 2)( a\^{}2 + 2a - 2)
\item
  (ঘ) ( a\^{}2 - 2a - 2)( a\^{}2 - 2a + 2)
\end{itemize}

\textbf{\hl{উত্তরঃ (গ) (}}\(a^{2}\) \textbf{\hl{- 2a - 2)(}}\(a^{2}\)
\textbf{\hl{+ 2a - 2)}}

\textbf{Explanation:}

\(a^{4}\) + 4

= \({(a2)}^{2}\) -2.\(a^{2}\).2 + \({(2)}^{2}\)-4\(a^{2}\)

= \({(a2\  - \ 2)}^{2}\) - \({(2a)}^{2}\)

= (\(a^{2}\)-2+2a)(\(a^{2}\)-2 - 2a)

= (\(a^{2}\)+2a-2)(\(a^{2}\) -2a - 2)

\textbf{প্রশ্ন নম্বর-} ৫. \(p^{6} - q^{6}\)

\textbf{এর উৎপাদক কত?}

\begin{itemize}
\item
  (ক) ( p\^{}3 - q\^{}3 ) ( p\^{}3 + q\^{}3 )
\item
  (খ) (p + q)( p\^{}2 - pq + q\^{}2 ) ( p\^{}3 - q\^{}3 )
\item
  (গ) ( p\^{}2 - q\^{}2 )( p\^{}2 - pq + q\^{}2 )( p\^{}2 - pq + q\^{}2
  )
\item
  (ঘ) (p + q)(p - q)( p\^{}2 - pq + q\^{}2 ) ( p\^{}2 + pq + q\^{}2 )
\end{itemize}

\textbf{\hl{উত্তরঃ (ঘ) (p + q)(p - q)(}}\(p^{2}\) \textbf{\hl{- pq +}}
\(q^{2}\)\textbf{\hl{) (}}\(p^{2}\) \textbf{\hl{+ pq +}}
\(q^{2}\)\textbf{\hl{)}}

\textbf{Explanation:}

\(p^{6}\) - \(q^{6}\)

= \({(p3)}^{2}\) - \({(q3)}^{2}\)

= ( \(p^{3}\) + \(q^{3}\))( \(p^{3}\)- \(q^{3}\))

=(p+q)( \(p^{2}\) - pq + \(q^{2}\))(p-q)( \(p^{2}\) + pq + \(q^{2}\))

\textbf{প্রশ্ন নম্বর-} ৬. \(x^{2}\) - \(y^{2}\) + 2y - 1 এর একটি উৎপাদক-
{[}এনএসআই-২০০৮{]}

\textbf{\hl{উত্তরঃ (গ) x + y - 1}}

\textbf{Explanation:}

\(x^{2}\) - \(y^{2}\) + 2y-1

= \(x^{2}\) -(\(y^{2}\)-2y+1)

= \(x^{2}\) - \({(y - 1)}^{2}\)

= (x+y-1)(x-y+1)

\textbf{প্রশ্ন নম্বর- ৭.} \(a^{3}\) \textbf{+} \(b^{3}\) \textbf{এর উৎপাদক
কোনটি?}

\begin{itemize}
\item
  (ক) (a + b)(a + b)(a + b)
\item
  (খ) (a + b)( a\^{}2 + b\^{}2 )
\item
  (গ) (a - b)(a + b)( a\^{}2 + b\^{}2 )
\item
  (ঘ) (a + b)( a\^{}2 - ab + b\^{}2 )
\end{itemize}

\textbf{\hl{উত্তরঃ (ঘ) (a + b)(}}\(a^{2}\) \textbf{\hl{- ab +}}
\(b^{2}\)\textbf{\hl{)}}

\textbf{Explanation:}

এটা সূত্র

\textbf{প্রশ্ন নম্বর- ৮. কোনটি} \(a^{3}\) \textbf{+ 1 এর উৎপাদক?}

\begin{itemize}
\item
  (ক) (a + 1)(a + 1)
\item
  (খ) (a + 1)( a\^{}2 + a + 1)
\item
  (গ) (a - 1)( a\^{}2 - a + 1)
\item
  (ঘ) (a + 1)( a\^{}2 - a + 1)
\end{itemize}

\textbf{\hl{উত্তরঃ (ঘ) (a + 1)(}}\(a^{2}\) \textbf{\hl{- a + 1)}}

\textbf{Explanation:}

\(a^{3}\)+1=(a+1)(\(a^{2}\)-a+1)

\textbf{প্রশ্ন নম্বর-} ৯. \(x^{3}\) - 8 এর উৎপাদক কোনটি? {[} প্রাথমিক সহকারী
শিক্ষক-২০১২{]}

\begin{itemize}
\item
  (ক) x - 2
\item
  (খ) x - 4
\item
  (গ) x\^{}2 - x - 2
\item
  (ঘ) x\^{}2 - 2x - 2
\end{itemize}

\textbf{\hl{উত্তরঃ (ক) x - 2}}

\textbf{Explanation:}

\(x^{3}\)-8

= \(x^{3}\)- \(2^{3}\)

=(x-2)(\(x^{2}\)+2x+4)

\subsubsection{Middle Term দ্বারা
উৎপাদক}\label{middle-term-ux9a6ux9acux9b0-ux989ux9ceux9aaux9a6ux995}

\textbf{Middle Term Factor:} একটি রাশিকে সর্বোচ্চ ঘাত থেকে সর্বনিম্ন ঘাতে
সাজালে যদি a\(x^{2}\) + bx + c আকার ধারণ করে তবে তাকে মধ্যপদী প্রক্রিয়ায়
উৎপাদক বিশ্লেষণ করা যায় অর্থাৎ প্রথম রাশি এর উপরের পাওয়ারের অর্ধেক মাঝের
রাশিতে থাকলে ( প্রথম ঘাত ও দ্বিতীয় ঘাতের অনুপাত ২:১ হতে হবে) তাহলে তা
Middle Term Factor- এর নিয়মানুসারে উৎপাদক করতে হয় ।

a\(x^{2}\) + bx + c আকারের রাশিকে উৎপাদকে বিশ্লেষণের নিয়ম দেয়া হলো:

প্রথম ও শেষ রাশি গুণ করার পর প্রদত্ত রাশিটিকে এমনভাবে দুটি রাশিতে ভাঙ্গতে হবে
যেন ঐ রাশি দুটি যোগ অথবা বিয়োগ করে মাঝখানের রাশিটির সমান হয়। এক্ষেত্রে
শেষের রাশিটির আগে + থাকলে ঐ রাশি দুটি যোগ করতে হবে, এবং শেষের রাশিটির আগে
- চিহ্ন থাকলে রাশি দুটি বিয়োগ করে মাঝের রাশিটি আনতে হয়। যেমন:

\textbf{প্রশ্ন নম্বর-} ১. \(x^{2}\) - 3x - 10 এর সঠিক উৎপাদক কোন দুটি? {[}
প্রাক প্রাথমিক সহকারী শিক্ষক-২০১৩{]}

\begin{itemize}
\item
  (ক) (x + 2)(x - 5)
\item
  (খ) (x - 2)(x + 5)
\item
  (গ) (x + 3)(x - 5)
\item
  (ঘ) (x - 3)(x + 5)
\end{itemize}

\textbf{\hl{উত্তরঃ (ক) (x + 2)(x - 5)}}

\textbf{Explanation:}

\(x^{2}\) - 3x - 10

= \(x^{2}\) - 5x + 2x - 10

= x(x-5)+2(x-5)

= (x-5)(x+2)

\textbf{প্রশ্ন নম্বর-} ২. \(x^{2}\) - 7x + 12 = কত? {[} প্রাথমিক সহকারী
শিক্ষক-২০১০{]}

\begin{itemize}
\item
  (ক) (x - 4)(x - 3)
\item
  (খ) (x + 4)(x + 3)
\item
  (গ) (x + 3)(x - 4)
\item
  (ঘ) (x + 4)(x - 3)
\end{itemize}

\textbf{\hl{উত্তরঃ (ক) (x - 4)(x - 3)}}

\textbf{Explanation:}

\(x^{2}\) -7x-12

= \(x^{2}\) -3x-4x+12

= x(x-3)-4(x-3)

=(x-3)(x-4)

\textbf{প্রশ্ন নম্বর-} ৩. \(m^{2}\)+8m+15 এর উৎপাদক-\/-

\begin{itemize}
\item
  (ক) (m+5)(m-3)
\item
  (খ) (m-5)(m+3)
\item
  (গ) (m+5)(m+3)
\item
  (ঘ) (m-5)(m-3)
\end{itemize}

\textbf{\hl{উত্তরঃ (গ) (m+5)(m+3)}}

\textbf{Explanation:}

\(m^{2}\)+8m+15

= \(m^{2}\)+5m+3m+15

= m(m+5)+3(m+5)

= (m+5)(m+3)

\textbf{প্রশ্ন নম্বর-} ৪. 3\(x^{2}\) - 7x - 6 এর উৎপাদক সমূহ কোনটি?
{[}পররাষ্ট্র মন্ত্রণালয়ের প্রশাসনিক কর্মকর্তা-২০০৪{]}

\begin{itemize}
\item
  (ক) (3x + 2)(x - 3)
\item
  (খ) (3x - 2)(x + 3)
\item
  (গ) (3x + 2)(x + 3)
\item
  (ঘ) (3x - 2)(x - 3)
\end{itemize}

\textbf{\hl{উত্তরঃ (ক) (3x + 2)(x - 3)}}

\textbf{Explanation:}

3\(x^{2}\)-7x-6

= 3\(x^{2}\)-9x+2x-6

= 3x(x-3)+2(x-3)

= (x-3)(3x+2)

\textbf{প্রথমে সংখ্যা থাকলে:} প্রথম সংখ্যাটি দিয়ে শেষের সংখ্যাটিকে গুণ করে যে
গুণফল বের হবে তাকে এমনভাবে দুটি সংখ্যায় ভাঙ্গতে হবে যেন সংখ্যা দুটি যোগ অথবা
বিয়োগ করে মাঝের সংখ্যাটির সমান হয়। এক্ষেত্রে, শেষের সংখ্যার আগে + চিহ্ন থাকলে
ঐ সংখ্যা দুটি যোগ করতে হবে এবং শেষের সংখ্যাটির আগে -- চিহ্ন থাকলে ঐ

সংখ্যা দুটি বিয়োগ করতে হবে। উৎপাদকের অংক করার সময় এই কথাটি সিরিয়াসলি মনে
রাখতে হবে ।

\textbf{প্রশ্ন নম্বর- ৫. 2}\(x^{2}\)\textbf{+x-15 এর উৎপাদক কোনটি?}

\begin{itemize}
\item
  (ক) (x+3)(2x-5)
\item
  (খ) (x-3)(2x-5)
\item
  (গ) (x-3)(2x+5)
\item
  (ঘ) (x+3)(2x+5)
\end{itemize}

\textbf{\hl{উত্তরঃ (ক) (x+3)(2x-5)}}

\textbf{Explanation:}

2\(x^{2}\)+x-15

= 2\(x^{2}\)+6x-5x-15

= 2x(x+3)-5(x+3)

=(x+3)(2x-5)

\textbf{প্রথমে সংখ্যা থাকলে:} প্রথম সংখ্যাটি দিয়ে শেষের সংখ্যাটিকে গুণ করে যে
গুণফল বের হবে তাকে এমনভাবে দুটি সংখ্যায় ভাঙ্গতে হবে যেন সংখ্যা দুটি যোগ অথবা
বিয়োগ করে মাঝের সংখ্যাটির সমান হয়। এক্ষেত্রে, শেষের সংখ্যার আগে + চিহ্ন থাকলে
ঐ সংখ্যা দুটি যোগ করতে হবে এবং শেষের সংখ্যাটির আগে -- চিহ্ন থাকলে ঐ

সংখ্যা দুটি বিয়োগ করতে হবে। উৎপাদকের অংক করার সময় এই কথাটি সিরিয়াসলি মনে
রাখতে হবে ।

\subsubsection{ফাংশনের দ্বারা
উৎপাদক}\label{ux9abux9b6ux9a8ux9b0-ux9a6ux9acux9b0-ux989ux9ceux9aaux9a6ux995}

\textbf{Key Point:} উৎপাদক = 0 ধরে, x এর মান বের করে তা সরাসরি উত্তরে
বসাবেন। x এর মান বসানোর পর যদি শূন্য (0) হয় তাহাই সঠিক উত্তর । f(a) = 0 হলে
x - a উহার একটি উৎপাদক। x - a = 0, x = a

\textbf{Key Point:} মনে রাখতে হবে ভাগশেষ উৎপাদকের ক্ষেত্রে যে মানের জন্য
উক্ত রাশিটির মান শূন্য হবে সেটি হবে উক্ত রাশির একটি উৎপাদক। যেমন, a = - 1
বসালে যদি শূন্য হয় তবে (a + 1) হবে রাশিটির একটি উৎপাদক ।

\textbf{প্রশ্ন নম্বর-} ১. \(a^{3}\)− 21a − 20 এর উৎপাদকে বিশ্লেষণ করুন।

\begin{itemize}
\item
  (ক) (a + 2)
\item
  (খ) (a + 1)
\item
  (গ) (a − 2 )
\item
  (ঘ) (a − 1 )
\end{itemize}

\textbf{\hl{উত্তরঃ (খ) (a + 1)}}

\textbf{Explanation:}

\textbf{সমাধানঃ}

এ ধরনের অংক অর্থাৎ প্রথম রাশিতে ঘন থাকলে সাধারণ নিয়মের মতো উৎপাদক করা যাবে
না। তখন প্রথমে a এর এমন একটি মান ধরতে হবে যা বসালে সম্পূর্ণ রাশিটির মান শুণ্য
হয় ।

এখানে প্রদত্ত রাশিটি \(a^{3}\) − 21a − 20 যাতে a এর যায়গায় − 1 বসালে
রাশিটির মান শুণ্য হয়। এখন এই − 1 কে +1 বানিয়ে (a + 1) হলো এই রাশিটির একটি
উৎপাদক। তাহলে সমাধান হবে

প্রদত্ত রাশিটি

\(a^{3}\) − 21a − 20

= \(a^{3}\) + \(a^{2}\) − \(a^{2}\)− a − 20a − 20

(২য় লাইনটি লিখার আগে ৩য় লাইনে তিনটিতে (a+1) ফাঁকা রেখে ৩ বার লিখুন এবং
প্রথম লাইনটিকে মেলানোর জন্য যা যা গুণ করা দরকার তা দিয়ে গুণ করে ২য় লাইনটি
আনুন)

= \(a^{2}\) (a + 1) − a (a + 1) − 20(a + 1) (এখানে প্রতিটি রাশিতে (a + 1)
আছে তাই (a + 1) কমন নিয়ে যা থাকবে তা লিখতে হবে)

= (a + 1)( \(a^{2}\) − a − 20) ( এখানে আবার ২য় অংশটির মিডিল টার্ম উৎপাদক
হবে )

=(a + 1)( \(a^{2}\) − 5a + 4a − 20)

= (a + 1) \{ a(a − 5) + 4(a − 5)\}

= (a + 1)(a − 5)(a + 4)

\textbf{প্রশ্ন নম্বর- ২.} \(a^{3}\) \textbf{− 7a − 6 এর উৎপাদক কত?}

\begin{itemize}
\item
  (ক) (a + 1)(a − 2)(a − 3)
\item
  (খ) (a + 1)(a + 2)(a − 3)
\item
  (গ) (a − 1)(a + 2)(a − 3)
\item
  (ঘ) (a − 1)(a − 2)(a − 3)
\end{itemize}

\textbf{\hl{উত্তরঃ (খ) (a + 1)(a + 2)(a − 3)}}

\textbf{Explanation:}

\(a^{3}\) − 7a − 6

= \(a^{3}\) + \(a^{2}\) - \(a^{2}\) − a − 6a − 6

= \(a^{2}\)(a + 1) − a(a + 1) − 6(a + 1)

= (a + 1) ( \(a^{2}\) − a − 6)

= (a + 1) ( \(a^{2}\) + 2a − 3a − 6)

= (a + 1)(a + 2)(a − 3)

\textbf{প্রশ্ন নম্বর- ৩.} \(x^{3}\)\textbf{− x − 6 এর উৎপাদকে বিশ্লেষিত রূপ
কোনটি?}

\begin{itemize}
\item
  (ক) (x+2)(x\^{}2+2x+3)
\item
  (খ) (x−2)(x\^{}2+2x+3)
\item
  (গ) (x−2)(x\^{}2−2x+3)
\item
  (ঘ) (x+2)(x\^{}2−2x+3)
\end{itemize}

\textbf{\hl{উত্তরঃ (খ)}} \((x - 2)(x^{2} + 2x + 3)\)

\textbf{Explanation:}

\(x^{3}\) − x − 6

= \(x^{3}\) − \({2x}^{2}\) + \({2x}^{2}\) − 4x + 3x − 6

= \(x^{2}\)(x − 2) + 2x (x − 2) + 3(x − 2)

= (x − 2) ( \(x^{2}\) + 2x + 3)

\textbf{প্রশ্ন নম্বর- ৪.} \(x^{2} + 7x + p\) \textbf{যদি x − 5 দ্বারা বিভাজ্য
হয় তবে p এর মান কত হবে?}

\begin{itemize}
\item
  (ক) 30
\item
  (খ) − 30
\item
  (গ) − 60
\item
  (ঘ) 5
\end{itemize}

\textbf{\hl{উত্তরঃ (গ) - 60}}

\textbf{Explanation:}

যেহেতু এখানে x − 5 হচ্ছে সম্পূর্ণ রাশিটির একটি উৎপাদক তাহলে x = 5 বসালে সম্পূর্ণ
রাশিটির মান ০ হবে ।

সুতরাং x = 5 বসিয়ে পাই,

\(5^{2}\) + 7 × 5 + p = 0

বা, 25 + 35 + p = 0

বা, 60 + p = 0

সুতরাং p = − 60

\subsubsection{বহুপদী
উৎপাদক}\label{ux9acux9b9ux9aaux9a6-ux989ux9ceux9aaux9a6ux995}

যে সব বীজগাণিতিক রাশির সর্বোচ্চ ঘাত ৩ বা তার চেয়ে বেশি সে সকল রাশিকে বহুপদী
রাশি বলে। এ সকল রাশির উৎপাদক গুলোকে বহুপদী উৎপাদক বলে।

{[}Note: সর্বোচ্চ ঘাত ২টি হলে দ্বিপদী রাশি বলে যা এই অধ্যায়ের শুরুতে আলোচনা
করা হয়েছে।{]}

এক্ষেত্রে উৎপাদক বের করার সহজ নিয়ম হলো ফাংশনের মাধ্যমে উৎপাদক বের করা। তবে
এমসিকিউ পরীক্ষার জন্য উত্তরের অপশন

থেকে সমাধান করা খুব সহজ। নিচের সমাধান করা প্রশ্নগুলো দেখুন ।

\textbf{প্রশ্ন নম্বর- ১. কোনটি} \({2x}^{4}\) \textbf{− 5}\(x^{3}\)
\textbf{+} \({6x}^{2}\) \textbf{− 5x + 2 এর একটি উৎপাদক?}

\begin{itemize}
\item
  (ক) x + 1
\item
  (খ) x − 1
\item
  (গ) x + 2
\item
  (ঘ) x − 2
\end{itemize}

\textbf{\hl{উত্তরঃ (খ) x − 1}}

\textbf{Explanation:}

ফাংশনের সাহায্যে সমাধানঃ

\({2x}^{4}\) − 5\(x^{3}\) + \({6x}^{2}\) − 5x + 2

= 2\(x^{4}\) − 2\(x^{3}\) − 3\(x^{3}\) + 3\(x^{2}\) + 3\(x^{2}\) − 3x −
2x + 2

= 2\(x^{3}\)(x − 1) − 3x\textsuperscript{2} (x − 1 ) + 3x (x − 1) − 2 (x
− 1)

= (x − 1) (2\(x^{3}\) − 3\(x^{2}\) + 3x − 2)

= (x − 1 ) ( 2\(x^{3}\) − 2\(x^{2}\) − \(x^{2}\) + x + 2x − 2)

= (x − 1) \{ 2\(x^{2}\) (x − 1) − x (x − 1 ) + 2 (x − 1 ) \}

= (x − 1) (x − 1) (2\(x^{2}\) − x + 2)

অপশনের সাহায্যে সমাধানঃ

প্রথম অপশন দেওয়া আছে x + 1 . অর্থাৎ x এর মান -1 বসিয়ে যদি সম্পূর্ণ রাশিটির
মান শূন্য হয় তাহলে x + 1 হবে রাশিটির একটি উৎপাদক। এখানে x এর মান 1 বসালে
সম্পূর্ণ রাশিটির মান শূন্য হয় । যেমন:

\({2x}^{4}\) − 5\(x^{3}\) + \({6x}^{2}\) − 5x + 2

= 2.\(1^{4}\) − 5.\(1^{3}\) + 6.\(1^{2}\) − 5.1 + 2

= 2 − 5 + 6 − 5 + 2

= 10 − 10

= 0 তাই x − 1 হবে রাশিটির একটি উৎপাদক। সুতরাং সঠিক উত্তর অপশন (খ) ।

এমসিকিউ পরীক্ষায় জন্য অপশন থেকে এভাবে

সবগুলো চেষ্টা করতে হবে।

\textbf{প্রশ্ন নম্বর-} ২. 3\(x^{3}\) + 2\(x^{2}\)− 21x − 20 রাশিটির একটি
উৎপাদক হচ্ছে-

\begin{itemize}
\item
  (ক) x + 2
\item
  (খ) x − 2
\item
  (গ) x + 1
\item
  (ঘ) x − 1
\end{itemize}

\textbf{\hl{উত্তরঃ (গ) x+1}}

\textbf{Explanation:}

ফাংশনের সাহায্যে সমাধানঃ

3\(x^{3}\) + 2\(x^{2}\)− 21x − 20

= 3\(x^{3}\)+ 3\(x^{2}\) − \(x^{2}\) − x − 20x − 20

= 3\(x^{2}\) (x + 1 ) − x (x + 1) − 20 (x + 1)

= (x + 1) (3\(x^{2}\) − x − 20 )

অপশনের সাহায্যে সমাধানঃ

প্রথম অপশন দেওয়া আছে x + 2 . অর্থাৎ x এর মান -2 বসিয়ে যদি সম্পূর্ণ রাশিটির
মান শূন্য হয় তাহলে x + 2 হবে রাশিটির একটি উৎপাদক। কিন্তু এখানে x এর মান - 1
বসালে সম্পূর্ণ রাশিটির মান শূন্য হয় । যেমন:

3\(x^{3}\) + 2\(x^{2}\)− 21x − 20

= 3.\({( - \ 1)}^{3}\) + 2.\({(\  - \ 1\ )}^{2}\) − 21 × ( − 1 ) − 20

= 3×( − 1 ) + 2.1 + 21 − 20

= − 3 + 2 + 21 − 20

= 23 − 23

= 0 তাই x + 1 হবে রাশিটির একটি উৎপাদক। সুতরাং সঠিক উত্তর অপশন (গ) ।

এমসিকিউ পরীক্ষায় জন্য অপশন থেকে এভাবে

সবগুলো চেষ্টা করতে হবে।

\textbf{প্রশ্ন নম্বর- ৩.} \(x^{3}\)\textbf{− 6}\(x^{2}\) \textbf{+ 11x − 6
এর উৎপাদক কোনটি?}

\begin{itemize}
\item
  (ক) (x + 1 ) ( x − 2) (x − 3)
\item
  (খ) (x − 1 ) ( x − 2) (x + 3)
\item
  (গ) (x − 1 ) ( x + 2) (x − 3)
\item
  (ঘ) (x − 1 ) ( x − 2) (x − 3)
\end{itemize}

\textbf{\hl{উত্তরঃ (ঘ) (x − 1 ) ( x − 2) (x − 3)}}

\textbf{Explanation:}

\(x^{3}\)− 6\(x^{2}\) + 11x − 6

= \(x^{3}\) − \(x^{2}\) − 5\(x^{2}\) + 5x + 6x − 6

= \(x^{2}\) (x − 1) − 5x (x − 1 ) + 6 (x − 1)

= (x − 1) (\(x^{2}\) − 5x + 6)

= (x − 1) (\(x^{2}\) − 3x − 2x + 6)

= (x − 1 ) \{ x ( x − 3) − 2 (x − 3)\}

= (x − 1 ) ( x − 3) (x − 2)

= (x − 1 ) ( x − 2) (x − 3)

\textbf{প্রশ্ন নম্বর- ৪. 4}\(x^{4}\) \textbf{− 25}\(x^{2}\) \textbf{+ 36 =
কত?}

\begin{itemize}
\item
  (ক) (x−2) (x−2) (2x+3) (2x−3)
\item
  (খ) (x+2) (x+2) (2x+3) (2x−3)
\item
  (গ) (x+2) (x−2) (2x−3) (2x−3)
\item
  (ঘ) (x+2) (x−2) (2x+3) (2x−3)
\end{itemize}

\textbf{\hl{উত্তরঃ (ঘ) (x + 2) (x − 2) (2x + 3) (2x − 3)}}

\textbf{Explanation:}

4\(x^{4}\) − 25\(x^{2}\) + 36

= 4\(x^{4}\) − 16\(x^{2}\) − 9\(x^{2}\) + 36

= 4\(x^{2}\) (\(x^{2}\) − 4) − 9 (\(x^{2}\) − 4)

= (\(x^{2}\) − 4) (4\(x^{2}\) − 9) 2

= (\(x^{2}\) − \(2^{2}\)) \{ \({(2x)}^{2}\) − \(3^{2}\)\}

= (x + 2) (x − 2) (2x + 3) (2x − 3)

\textbf{প্রশ্ন নম্বর-} ৫. উৎপাদকে বিশ্লেষণ করুন: \(x^{2}\) − 2ax + (a + b) (a
− b)

\begin{itemize}
\item
  (ক) ( x − a − b ) ( x + a + b )
\item
  (খ) ( x + a − b ) ( x − a + b )
\item
  (গ) ( x − a − b ) ( x − a − b )
\item
  (ঘ) ( x − a − b ) ( x − a + b )
\end{itemize}

\textbf{\hl{উত্তরঃ (ঘ) ( x − a − b ) ( x − a + b )}}

\textbf{Explanation:}

এই ধরনের উৎপাদকের অংকগুলো উপরের সাধারণ নিয়মেই করতে হবে। শুধু প্রথমে একটু
অতিরিক্ত কাজ করতে হবে। এই অংকের মূল পয়েন্ট হল 2ax ।

\(x^{2}\) − 2ax + (a + b) (a − b)

= \(x^{2}\) − 2ax + \(a^{2}\) − \(b^{2}\)

= \({(\ x\  - \ a)}^{2}\) − \(b^{2}\)

= ( x − a − b ) ( x − a + b )

\section{36. উৎপাদকে বিশ্লেষণ / বহুপদী উৎপাদক MCQ প্রশ্ন উত্তর
অনুশীলন}\label{ux989ux9ceux9aaux9a6ux995-ux9acux9b6ux9b2ux9b7ux9a3-ux9acux9b9ux9aaux9a6-ux989ux9ceux9aaux9a6ux995-mcq-ux9aaux9b0ux9b6ux9a8-ux989ux9a4ux9a4ux9b0-ux985ux9a8ux9b6ux9b2ux9a8}

\textbf{x\^{}2 − y\^{}2 এর উৎপাদক কত?}

ক) (x + y) (x + y)

\textbf{\hl{খ) (x + y) (x - y) ✔}}

গ) (x - y) (x - y)

ঘ) (y + x) (y - x)

\textbf{\hl{✔ সঠিক উত্তর: খ) (x + y) (x - y)}}

\begin{quote}
\textbf{ব্যাখ্যা:} \emph{আমরা জানি, a\^{}2 − b\^{}2 = (a + b)(a − b)}

\emph{∴ x\^{}2 − y\^{}2 = (x + y)(x − y)}
\end{quote}

\textbf{x\^{}2 − y\^{}2 + 2y − 1 এর একটি উৎপাদক নিচের কোনটি?}

ক) x + y + 1

\textbf{\hl{খ) x + y - 1 ✔}}

গ) x - y - 1

ঘ) x - y

\textbf{\hl{✔ সঠিক উত্তর: খ) x + y - 1}}

\begin{quote}
\textbf{ব্যাখ্যা:} \emph{x\^{}2 − y\^{}2 + 2y − 1}

\emph{= x\^{}2 − (y\^{}2 − 2y + 1)}

\emph{= x\^{}2 − (y − 1)\^{}2}

\emph{= \{x + (y − 1)\} \{x − (y − 1)\}}

\emph{= (x + y − 1) (x − y + 1)}

\emph{∴ একটি উৎপাদক হলো (x + y − 1)}
\end{quote}

\textbf{x\^{}2 − x − 2 এর একটি উৎপাদক?}

ক) x + 2

খ) x - 3

\textbf{\hl{গ) x + 1 ✔}}

ঘ) x - 1

\textbf{\hl{✔ সঠিক উত্তর: গ) x + 1}}

\begin{quote}
\textbf{ব্যাখ্যা:} \emph{x\^{}2 − x − 2}

\emph{= x\^{}2 − 2x + x − 2}

\emph{= x(x − 2) + 1(x − 2)}

\emph{= (x − 2)(x + 1)}

\emph{∴ একটি উৎপাদক হলো (x + 1)}
\end{quote}

\textbf{x\^{}3 − x − 6 এর উৎপাদকে বিশ্লেষিত রূপ কোনটি?}

ক) (x + 2) (x\^{}2 + 2x + 3)

\textbf{\hl{খ) (x - 2) (x\^{}2 + 2x + 3) ✔}}

গ) (x - 2) (x\^{}2 - 2x + 3)

ঘ) (x + 2) (x\^{}2 - 2x + 3)

\textbf{\hl{✔ সঠিক উত্তর: খ) (x - 2) (x\^{}2 + 2x + 3)}}

\begin{quote}
\textbf{ব্যাখ্যা:} \emph{ধরি, f(x) = x\^{}3 − x − 6}

\emph{f(2) = 2\^{}3 − 2 − 6 = 8 − 8 = 0}

\emph{∴ (x − 2) একটি উৎপাদক।}

\emph{x\^{}3 − x − 6}

\emph{= x\^{}3 − 2x\^{}2 + 2x\^{}2 − 4x + 3x − 6}

\emph{= x\^{}2(x − 2) + 2x(x − 2) + 3(x − 2)}

\emph{= (x − 2) (x\^{}2 + 2x + 3)}
\end{quote}

\textbf{25 − 9x\^{}4 এর উৎপাদক কত?}

ক) (3x\^{}2 + 5) (-3x\^{}2 - 5)

খ) (5 + 3x) (5 - 3x)

গ) (-3x\^{}2 + 5) (-3x\^{}2 - 5)

\textbf{\hl{ঘ) (5 + 3x\^{}2) (5 - 3x\^{}2) ✔}}

\textbf{\hl{✔ সঠিক উত্তর: ঘ) (5 + 3x\^{}2) (5 - 3x\^{}2)}}

\begin{quote}
\textbf{ব্যাখ্যা:} \emph{25 − 9x\^{}4}

\emph{= 5\^{}2 − (3x\^{}2)\^{}2}

\emph{∴ (5 + 3x\^{}2) (5 − 3x\^{}2)}
\end{quote}

\textbf{(1/5) x\^{}2 − 5 এর উৎপাদকে বিশ্লেষণ কোনটি?}

\textbf{\hl{ক) (1/5) (x + 5)(x − 5) ✔}}

খ) (x + 5)(x − 5)

গ) (1/5) (x + 5)( (x/5) − 1)

ঘ) (x − 5)(x + 5)

\textbf{\hl{✔ সঠিক উত্তর: ক) (1/5) (x + 5)(x − 5)}}

\begin{quote}
\textbf{ব্যাখ্যা:} \emph{(1/5) x\^{}2 − 5}

\emph{= (x\^{}2/5) − (25/5)}

\emph{= (x\^{}2 − 5\^{}2/5)}

\emph{= ((x + 5)(x − 5)/5)}

\emph{∴ (1/5) (x + 5)(x − 5)}
\end{quote}

\textbf{x\^{}4 + x\^{}2 + 1 এর উৎপাদক কোনটি?}

ক) (x\^{}2 + x + 1) (x\^{}2 + x - 1)

খ) (x\^{}2 - x + 1) (x\^{}2 - x - 1)

\textbf{\hl{গ) (x\^{}2 + x + 1) (x\^{}2 - x + 1) ✔}}

ঘ) (x\^{}2 + x + 1) (x\^{}2 + x + 1)

\textbf{\hl{✔ সঠিক উত্তর: গ) (x\^{}2 + x + 1) (x\^{}2 - x + 1)}}

\begin{quote}
\textbf{ব্যাখ্যা:} \emph{x\^{}4 + x\^{}2 + 1}

\emph{= (x\^{}2)\^{}2 + 2.x\^{}2.1 + (1)\^{}2 − (x)\^{}2}

\emph{= (x\^{}2 + 1)\^{}2 − (x)\^{}2}

\emph{∴ (x\^{}2 + x + 1) (x\^{}2 − x + 1)}
\end{quote}

\textbf{4x\^{}4 − 25x\^{}2 + 36 = কত?}

\textbf{\hl{ক) (x + 2)(x - 2)(2x + 3)(2x - 3) ✔}}

খ) (2x + 1)(2x - 1)(x + 6)(x - 6)

গ) (x + 3)(x - 3)(2x + 2)(2x - 2)

ঘ) (x + 2)(2x - 3)(2x + 3)(3x + 2)

\textbf{\hl{✔ সঠিক উত্তর: ক) (x + 2)(x - 2)(2x + 3)(2x - 3)}}

\begin{quote}
\textbf{ব্যাখ্যা:} \emph{4x\^{}4 − 25x\^{}2 + 36}

\emph{= 4x\^{}4 − 16x\^{}2 − 9x\^{}2 + 36}

\emph{= 4x\^{}2(x\^{}2 − 4) − 9(x\^{}2 − 4)}

\emph{= (x\^{}2 − 4) (4x\^{}2 − 9)}

\emph{∴ (x + 2)(x − 2)(2x + 3)(2x − 3)}
\end{quote}

\textbf{x\^{}2 − y(y − 2) − 1 এর উৎপাদক নিচের কোনটি?}

ক) (x − y − 1) (x − y + 1)

\textbf{\hl{খ) (x + y − 1) (x − y + 1) ✔}}

গ) (x + y + 1) (x − y − 1)

ঘ) (x − y) (x + y + 1)

\textbf{\hl{✔ সঠিক উত্তর: খ) (x + y − 1) (x − y + 1)}}

\begin{quote}
\textbf{ব্যাখ্যা:} \emph{x\^{}2 − y(y − 2) − 1}

\emph{= x\^{}2 − y\^{}2 + 2y − 1}

\emph{= x\^{}2 − (y\^{}2 − 2y + 1)}

\emph{= x\^{}2 − (y − 1)\^{}2}

\emph{∴ \{x + (y − 1)\} \{x − (y − 1)\} = (x + y − 1) (x − y + 1)}
\end{quote}

\textbf{a\^{}4 + 4 এর উৎপাদক কি?}

ক) (a\^{}2+2a+2) (a\^{}2+2a-2)

\textbf{\hl{খ) (a\^{}2+2a+2) (a\^{}2-2a+2) ✔}}

গ) (a\^{}2-2a+2) (a+2a-2)

ঘ) (a\^{}2-2a-2) (a\^{}2-2a+2)

\textbf{\hl{✔ সঠিক উত্তর: খ) (a\^{}2+2a+2) (a\^{}2-2a+2)}}

\begin{quote}
\textbf{ব্যাখ্যা:} \emph{a\^{}4 + 4}

\emph{= (a\^{}2)\^{}2 + (2)\^{}2}

\emph{= (a\^{}2 + 2)\^{}2 − 2.a\^{}2.2}

\emph{= (a\^{}2 + 2)\^{}2 − (2a)\^{}2}

\emph{∴ (a\^{}2 + 2 + 2a) (a\^{}2 + 2 − 2a)}
\end{quote}

\textbf{x\^{}6 − y\^{}6 এর উৎপাদক কোনটি?}

\textbf{\hl{ক) (x + y) (x - y) (x\^{}2 + xy + y\^{}2) (x\^{}2 - xy +
y\^{}2) ✔}}

খ) (x\^{}3 + y\^{}3) (x\^{}3 - y\^{}3)

গ) (x + y) (x - y) (x\^{}2 + 2xy) (x\^{}2 - xy + y\^{}2)

ঘ) (x\^{}2 - y\^{}2) (x\^{}2 + xy + y\^{}2) (x\^{}2 - xy + y\^{}2)

\textbf{\hl{✔ সঠিক উত্তর: ক) (x + y) (x - y) (x\^{}2 + xy + y\^{}2)
(x\^{}2 - xy + y\^{}2)}}

\begin{quote}
\textbf{ব্যাখ্যা:} \emph{x\^{}6 − y\^{}6}

\emph{= (x\^{}3)\^{}2 − (y\^{}3)\^{}2}

\emph{= (x\^{}3 + y\^{}3) (x\^{}3 − y\^{}3)}

\emph{∴ (x + y)(x\^{}2 − xy + y\^{}2)(x − y)(x\^{}2 + xy + y\^{}2)}
\end{quote}

\textbf{x\^{}2 − y\^{}2 + 2y − 1 এর একটি উৎপাদক?}

ক) x+y+1

খ) x-y

\textbf{\hl{গ) x+y-1 ✔}}

ঘ) x-y-1

\textbf{\hl{✔ সঠিক উত্তর: গ) x+y-1}}

\begin{quote}
\textbf{ব্যাখ্যা:} \emph{x\^{}2 − y\^{}2 + 2y − 1}

\emph{= x\^{}2 − (y\^{}2 − 2y + 1)}

\emph{= x\^{}2 − (y − 1)\^{}2}

\emph{∴ (x + y − 1) (x − y + 1)}
\end{quote}

\textbf{x\^{}2 − y\^{}2 + 4y − 4 এর একটি উৎপাদক?}

\textbf{\hl{ক) x+y-2 ✔}}

খ) x+y+2

গ) x-y-2

ঘ) x-2y+1

\textbf{\hl{✔ সঠিক উত্তর: ক) x+y-2}}

\begin{quote}
\textbf{ব্যাখ্যা:} \emph{x\^{}2 − y\^{}2 + 4y − 4}

\emph{= x\^{}2 − (y\^{}2 − 4y + 4)}

\emph{= x\^{}2 − (y − 2)\^{}2}

\emph{∴ (x + y − 2) (x − y + 2)}
\end{quote}

\textbf{x\^{}2 − y\^{}2 − 2y − 1 এর একটি উৎপাদক?}

\textbf{\hl{ক) x-y-1 ✔}}

খ) x+y-1

গ) x-y+1

ঘ) x+2y+1

\textbf{\hl{✔ সঠিক উত্তর: ক) x-y-1}}

\begin{quote}
\textbf{ব্যাখ্যা:} \emph{x\^{}2 − y\^{}2 − 2y − 1}

\emph{= x\^{}2 − (y\^{}2 + 2y + 1)}

\emph{= x\^{}2 − (y + 1)\^{}2}

\emph{∴ (x + y + 1) (x − y − 1)}
\end{quote}

\textbf{x\^{}2 + 2xy − 2y − 1 এর উৎপাদক কত?}

ক) (x + y + 1)(x - 1)

খ) (x + 2y + 1)(x + 1)

গ) (x + y + 1)(x + 1)

\textbf{\hl{ঘ) (x + 2y + 1)(x - 1) ✔}}

\textbf{\hl{✔ সঠিক উত্তর: ঘ) (x + 2y + 1)(x - 1)}}

\begin{quote}
\textbf{ব্যাখ্যা:} \emph{x\^{}2 + 2xy − 2y − 1}

\emph{= x\^{}2 − 1 + 2xy − 2y}

\emph{= (x + 1)(x − 1) + 2y(x − 1)}

\emph{= (x − 1)(x + 2y + 1)}
\end{quote}

\textbf{7x\^{}2 − x − 8 এর একটি উৎপাদক?}

\textbf{\hl{ক) (x + 1) (7x - 8) ✔}}

খ) x + 1

গ) 7x + 8

ঘ) x - 1

\textbf{\hl{✔ সঠিক উত্তর: ক) (x + 1) (7x - 8)}}

\begin{quote}
\textbf{ব্যাখ্যা:} \emph{7x\^{}2 − x − 8}

\emph{= 7x\^{}2 − 8x + 7x − 8}

\emph{= x(7x − 8) + 1(7x − 8)}

\emph{= (7x − 8)(x + 1)}
\end{quote}

\textbf{1 − a\^{}2 + 2ab − b\^{}2 এর উৎপাদক কোনটি?}

\textbf{\hl{ক) (1 + a - b) (1 - a + b) ✔}}

খ) (1 + a + b) (1 - a + b)

গ) (1 + a + b) (1 - a - b)

ঘ) (1 - a + b) (1 - a - b)

\textbf{\hl{✔ সঠিক উত্তর: ক) (1 + a - b) (1 - a + b)}}

\begin{quote}
\textbf{ব্যাখ্যা:} \emph{1 − a\^{}2 + 2ab − b\^{}2}

\emph{= 1 − (a\^{}2 − 2ab + b\^{}2)}

\emph{= 1 − (a − b)\^{}2}

\emph{= (1 + a − b) (1 − a + b)}
\end{quote}

\textbf{9x\^{}2 − (2x − 3y)\^{}2 = ?}

\textbf{\hl{ক) (5x − 3y)(x + 3y) ✔}}

খ) (5x − 3y)(x − 3y)

গ) (5x + 3y)(x − 3y)

ঘ) কোনটিই নয়

\textbf{\hl{✔ সঠিক উত্তর: ক) (5x − 3y)(x + 3y)}}

\begin{quote}
\textbf{ব্যাখ্যা:} \emph{9x\^{}2 − (2x − 3y)\^{}2}

\emph{= (3x)\^{}2 − (2x − 3y)\^{}2}

\emph{= \{3x + (2x − 3y)\}\{3x − (2x − 3y)\}}

\emph{= (3x + 2x − 3y)(3x − 2x + 3y)}

\emph{∴ (5x − 3y)(x + 3y)}
\end{quote}

\textbf{x\^{}4 + x\^{}2 + 1 এর একটি উৎপাদক x\^{}2 + x + 1 অপর উৎপাদকটি
কত?}

ক) x\^{}2 + x + 1

\textbf{\hl{খ) x\^{}2 − x + 1 ✔}}

গ) x\^{}2 + 1

ঘ) x + 1

\textbf{\hl{✔ সঠিক উত্তর: খ) x\^{}2 − x + 1}}

\begin{quote}
\textbf{ব্যাখ্যা:} \emph{x\^{}4 + x\^{}2 + 1}

\emph{= (x\^{}2)\^{}2 + 2 . x\^{}2 . 1 + (1)\^{}2 − x\^{}2}

\emph{= (x\^{}2 + 1)\^{}2 − (x)\^{}2}

\emph{= (x\^{}2 + x + 1)(x\^{}2 − x + 1)}
\end{quote}

\textbf{p\^{}6 − q\^{}6 এর উৎপাদক কত?}

ক) (p\^{}3 − q\^{}3) (p\^{}3 + q\^{}3)

খ) (p + q)(p\^{}2 − pq + q\^{}2) (p\^{}3 − q\^{}3)

গ) (p\^{}2 − q\^{}2)(p\^{}2 − pq + q\^{}2)(p\^{}2 − pq + q\^{}2)

\textbf{\hl{ঘ) (p + q)(p − q)(p\^{}2 − pq + q\^{}2) (p\^{}2 + pq +
q\^{}2) ✔}}

\textbf{\hl{✔ সঠিক উত্তর: ঘ) (p + q)(p − q)(p\^{}2 − pq + q\^{}2) (p\^{}2
+ pq + q\^{}2)}}

\begin{quote}
\textbf{ব্যাখ্যা:} \emph{p\^{}6 − q\^{}6}

\emph{= (p\^{}3)\^{}2 − (q\^{}3)\^{}2}

\emph{= (p\^{}3 + q\^{}3)(p\^{}3 − q\^{}3)}

\emph{∴ (p + q)(p\^{}2 − pq + q\^{}2)(p − q)(p\^{}2 + pq + q\^{}2)}
\end{quote}

\textbf{x\^{}2 − 1 − y(y − 2) এর উৎপাদক কত?}

ক) (x − y − 1)(x − y + 1)

খ) (x − y + 1)(x + y − 1)

\textbf{\hl{গ) (x + y − 1)(x − y + 1) ✔}}

ঘ) (x − y)(x + y + 1)

\textbf{\hl{✔ সঠিক উত্তর: গ) (x + y − 1)(x − y + 1)}}

\begin{quote}
\textbf{ব্যাখ্যা:} \emph{x\^{}2 − 1 − y(y − 2)}

\emph{= x\^{}2 − 1 − y\^{}2 + 2y}

\emph{= x\^{}2 − (y\^{}2 − 2y + 1)}

\emph{= x\^{}2 − (y − 1)\^{}2}

\emph{∴ (x + y − 1)(x − y + 1)}
\end{quote}

\textbf{a\^{}3 + b\^{}3 এর উৎপাদক কোনটি?}

ক) (a + b)(a + b)(a + b)

খ) (a + b)(a\^{}2 + b\^{}2)

গ) (a − b)(a + b)(a\^{}2 + b\^{}2)

\textbf{\hl{ঘ) (a + b)(a\^{}2 − ab + b\^{}2) ✔}}

\textbf{\hl{✔ সঠিক উত্তর: ঘ) (a + b)(a\^{}2 − ab + b\^{}2)}}

\begin{quote}
\textbf{ব্যাখ্যা:} \emph{আমরা জানি, a\^{}3 + b\^{}3 এর উৎপাদকের সূত্র হলো (a
+ b)(a\^{}2 − ab + b\^{}2)}
\end{quote}

\textbf{কোনটি a\^{}3 + 1 এর উৎপাদক?}

ক) (a + 1)(a + 1)

খ) (a + 1)(a\^{}2 + a + 1)

গ) (a - 1)(a\^{}2 - a + 1)

\textbf{\hl{ঘ) (a + 1)(a\^{}2 - a + 1) ✔}}

\textbf{\hl{✔ সঠিক উত্তর: ঘ) (a + 1)(a\^{}2 - a + 1)}}

\begin{quote}
\textbf{ব্যাখ্যা:} \emph{a\^{}3 + 1 = a\^{}3 + 1\^{}3}

\emph{∴ (a + 1)(a\^{}2 − a + 1)}
\end{quote}

\textbf{x\^{}3 − 8 এর উৎপাদক কোনটি?}

\textbf{\hl{ক) x - 2 ✔}}

খ) x - 4

গ) x\^{}2 − x − 2

ঘ) x\^{}2 − 2x − 2

\textbf{\hl{✔ সঠিক উত্তর: ক) x - 2}}

\begin{quote}
\textbf{ব্যাখ্যা:} \emph{x\^{}3 − 8}

\emph{= x\^{}3 − 2\^{}3}

\emph{∴ (x − 2)(x\^{}2 + 2x + 4)}
\end{quote}

\textbf{(x - y)(y + 3) কোন রাশির উৎপাদক?}

ক) x\^{}2 − y\^{}2 + 3y + 3

খ) xy − 3y + y\^{}2

\textbf{\hl{গ) xy − y\^{}2 + 3x − 3y ✔}}

ঘ) xy − x\^{}2 + 3x − 3y

\textbf{\hl{✔ সঠিক উত্তর: গ) xy − y\^{}2 + 3x − 3y}}

\begin{quote}
\textbf{ব্যাখ্যা:} \emph{(x − y)(y + 3)}

\emph{= xy + 3x − y\^{}2 − 3y {[}গুণ করে{]}}

\emph{∴ xy − y\^{}2 + 3x − 3y}
\end{quote}

\textbf{4x\^{}4 + 1 কে উৎপাদকে বিশ্লেষণ করলে কোনটি পাওয়া যায়?}

ক) (2x\^{}2 + 2x + 1)(2x\^{}2 − 2x − 1)

খ) (2x\^{}2 + 2x − 1)(2x\^{}2 − 2x + 1)

\textbf{\hl{গ) (2x\^{}2 + 2x + 1)(2x\^{}2 − 2x + 1) ✔}}

ঘ) (2x\^{}2 + 2x − 1)(2x\^{}2 − 2x − 1)

\textbf{\hl{✔ সঠিক উত্তর: গ) (2x\^{}2 + 2x + 1)(2x\^{}2 − 2x + 1)}}

\begin{quote}
\textbf{ব্যাখ্যা:} \emph{4x\^{}4 + 1}

\emph{= (2x\^{}2)\^{}2 + 1\^{}2}

\emph{= (2x\^{}2 + 1)\^{}2 − 2(2x\^{}2)(1)}

\emph{= (2x\^{}2 + 1)\^{}2 − (2x)\^{}2}

\emph{∴ (2x\^{}2 + 2x + 1)(2x\^{}2 − 2x + 1)}
\end{quote}

\textbf{(a - b)\^{}3 + (b-c)\^{}3 + (c-a)\^{}3 কে উৎপাদকে বিশ্লেষণ করুন।}

ক) 9(a-b)(b-c)(c-a)

\textbf{\hl{খ) 3(a-b)(b-c)(c-a) ✔}}

গ) 2(a-b)(b-c)(c-a)

ঘ) (a-b)(b-c)(c-a)

\textbf{\hl{✔ সঠিক উত্তর: খ) 3(a-b)(b-c)(c-a)}}

\begin{quote}
\textbf{ব্যাখ্যা:} \emph{ধরি, a - b = x, b - c = y, c - a = z}

\emph{এখানে, x + y + z = a − b + b − c + c − a = 0}

\emph{∴ x\^{}3 + y\^{}3 + z\^{}3 = 3xyz}

\emph{∴ 3(a − b)(b − c)(c − a)}
\end{quote}

\textbf{−15 + x + 2x\^{}2 এর উৎপাদক কোনটি?}

\textbf{\hl{ক) (x + 3) (2x - 5) ✔}}

খ) (x - 3) (2x + 5)

গ) (x + 3) (2x + 5)

ঘ) (x - 3) (2x - 5)

\textbf{\hl{✔ সঠিক উত্তর: ক) (x + 3) (2x - 5)}}

\begin{quote}
\textbf{ব্যাখ্যা:} \emph{−15 + x + 2x\^{}2}

\emph{= 2x\^{}2 + x − 15}

\emph{= 2x\^{}2 + 6x − 5x − 15}

\emph{= 2x(x + 3) − 5(x + 3)}

\emph{= (x + 3) (2x − 5)}
\end{quote}

\textbf{a\^{}3 − 21a − 20 রাশিরটির একটি উৎপাদক হবে নিচের কোনটি?}

ক) (a + 2)

খ) (a - 2)

\textbf{\hl{গ) (a + 1) ✔}}

ঘ) (a - 1)

\textbf{\hl{✔ সঠিক উত্তর: গ) (a + 1)}}

\begin{quote}
\textbf{ব্যাখ্যা:} \emph{ধরি, f(a) = a\^{}3 − 21a − 20}

\emph{f(−1) = (−1)\^{}3 − 21(−1) − 20 = −1 + 21 − 20 = 0}

\emph{∴ (a + 1) হলো উক্ত রাশির একটি উৎপাদক।}
\end{quote}

\textbf{x\^{}2 − 3x − 10 এর সঠিক উৎপাদক কোন দুটি?}

\textbf{\hl{ক) (x + 2)(x - 5) ✔}}

খ) (x - 2)(x + 5)

গ) (x + 3)(x - 5)

ঘ) (x - 3)(x + 5)

\textbf{\hl{✔ সঠিক উত্তর: ক) (x + 2)(x - 5)}}

\begin{quote}
\textbf{ব্যাখ্যা:} \emph{x\^{}2 − 3x − 10}

\emph{= x\^{}2 − 5x + 2x − 10}

\emph{= x(x − 5) + 2(x − 5)}

\emph{= (x − 5)(x + 2)}
\end{quote}

\textbf{x\^{}2 − 7x + 12 = কত?}

\textbf{\hl{ক) (x - 4)(x - 3) ✔}}

খ) (x + 4)(x + 3)

গ) (x + 3)(x - 4)

ঘ) (x + 4)(x - 3)

\textbf{\hl{✔ সঠিক উত্তর: ক) (x - 4)(x - 3)}}

\begin{quote}
\textbf{ব্যাখ্যা:} \emph{x\^{}2 − 7x + 12}

\emph{= x\^{}2 − 3x − 4x + 12}

\emph{= x(x − 3) − 4(x − 3)}

\emph{= (x − 3)(x − 4)}
\end{quote}

\textbf{x\^{}2 − 3x + 2 এর উৎপাদক নিচের কোনটি?}

\textbf{\hl{ক) x-1 ✔}}

খ) x+2

গ) x-3

ঘ) x+1

\textbf{\hl{✔ সঠিক উত্তর: ক) x-1}}

\begin{quote}
\textbf{ব্যাখ্যা:} \emph{x\^{}2 − 3x + 2}

\emph{= x\^{}2 − 2x − x + 2}

\emph{= x(x − 2) − 1(x − 2)}

\emph{= (x − 2)(x − 1)}
\end{quote}

\textbf{m\^{}2 + 8m + 15 এর উৎপাদক-\/-}

ক) (m+5)(m-3)

খ) (m-5)(m+3)

\textbf{\hl{গ) (m+5)(m+3) ✔}}

ঘ) (m-5)(m-3)

\textbf{\hl{✔ সঠিক উত্তর: গ) (m+5)(m+3)}}

\begin{quote}
\textbf{ব্যাখ্যা:} \emph{m\^{}2 + 8m + 15}

\emph{= m\^{}2 + 5m + 3m + 15}

\emph{= m(m + 5) + 3(m + 5)}

\emph{= (m + 5)(m + 3)}
\end{quote}

\textbf{3x\^{}2 − 7x − 6 এর উৎপাদক সমূহ কোনটি?}

\textbf{\hl{ক) (3x + 2)(x - 3) ✔}}

খ) (3x - 2)(x + 3)

গ) (3x + 2)(x + 3)

ঘ) (3x - 2)(x - 3)

\textbf{\hl{✔ সঠিক উত্তর: ক) (3x + 2)(x - 3)}}

\begin{quote}
\textbf{ব্যাখ্যা:} \emph{3x\^{}2 − 7x − 6}

\emph{= 3x\^{}2 − 9x + 2x − 6}

\emph{= 3x(x − 3) + 2(x − 3)}

\emph{= (x − 3)(3x + 2)}
\end{quote}

\textbf{2x\^{}2 + x − 15 এর উৎপাদক কোনটি?}

\textbf{\hl{ক) (x + 3)(2x - 5) ✔}}

খ) (x - 3)(2x - 5)

গ) (x - 3)(2x + 5)

ঘ) (x + 3)(2x + 5)

\textbf{\hl{✔ সঠিক উত্তর: ক) (x + 3)(2x - 5)}}

\begin{quote}
\textbf{ব্যাখ্যা:} \emph{2x\^{}2 + x − 15}

\emph{= 2x\^{}2 + 6x − 5x − 15}

\emph{= 2x(x + 3) − 5(x + 3)}

\emph{= (x + 3)(2x − 5)}
\end{quote}

\textbf{4x\^{}2 − 13x − 12 এর উৎপাদক কত?}

\textbf{\hl{ক) (x - 4)(4x + 3) ✔}}

খ) (2x - 4)(2x - 3)

গ) (2x - 4)(2x + 3)

ঘ) (x + 4)(4x - 3)

\textbf{\hl{✔ সঠিক উত্তর: ক) (x - 4)(4x + 3)}}

\begin{quote}
\textbf{ব্যাখ্যা:} \emph{4x\^{}2 − 13x − 12}

\emph{= 4x\^{}2 − 16x + 3x − 12}

\emph{= 4x(x − 4) + 3(x − 4)}

\emph{= (x − 4)(4x + 3)}
\end{quote}

\textbf{2x\^{}2 + x − 15 এর উৎপাদক কোনটি?}

ক) (x - 3)(2x - 5)

খ) (x + 3)(2x + 5)

\textbf{\hl{গ) (x + 3)(2x - 5) ✔}}

ঘ) (x - 3)(2x + 5)

\textbf{\hl{✔ সঠিক উত্তর: গ) (x + 3)(2x - 5)}}

\begin{quote}
\textbf{ব্যাখ্যা:} \emph{2x\^{}2 + x − 15}

\emph{= 2x\^{}2 + 6x − 5x − 15}

\emph{= 2x(x + 3) − 5(x + 3)}

\emph{= (x + 3)(2x − 5)}
\end{quote}

\textbf{x\^{}2 − 10xy − 11y\^{}2 এর উৎপাদক-}

ক) (x - y)(x + 11y)

\textbf{\hl{খ) (x - 11y)(x + y) ✔}}

গ) (x + 4y)(x - 5y)

ঘ) (x + 5y)(x - 4y)

\textbf{\hl{✔ সঠিক উত্তর: খ) (x - 11y)(x + y)}}

\begin{quote}
\textbf{ব্যাখ্যা:} \emph{x\^{}2 − 10xy − 11y\^{}2}

\emph{= x\^{}2 − 11xy + xy − 11y\^{}2}

\emph{= x(x − 11y) + y(x − 11y)}

\emph{= (x − 11y)(x + y)}
\end{quote}

\textbf{12x\^{}2 + 7x − 10 এর উৎপাদক-}

ক) (3x+5)(4x-2)

খ) (3x-5)(4x+2)

\textbf{\hl{গ) (4x+5)(3x-2) ✔}}

ঘ) (4x-5)(3x+2)

\textbf{\hl{✔ সঠিক উত্তর: গ) (4x+5)(3x-2)}}

\begin{quote}
\textbf{ব্যাখ্যা:} \emph{12x\^{}2 + 7x − 10}

\emph{= 12x\^{}2 − 8x + 15x − 10}

\emph{= 4x(3x − 2) + 5(3x − 2)}

\emph{= (3x − 2)(4x + 5)}
\end{quote}

\textbf{2x\^{}2 − xy − 6y\^{}2 এর উৎপাদক-\/-}

\textbf{\hl{ক) (2x+3y)(x-2y) ✔}}

খ) (2x-3y)(x+2y)

গ) (2x+3y)(2x-2y)

ঘ) (2x-3y)(2x+2y)

\textbf{\hl{✔ সঠিক উত্তর: ক) (2x+3y)(x-2y)}}

\begin{quote}
\textbf{ব্যাখ্যা:} \emph{2x\^{}2 − xy − 6y\^{}2}

\emph{= 2x\^{}2 − 4xy + 3xy − 6y\^{}2}

\emph{= 2x(x − 2y) + 3y(x − 2y)}

\emph{= (x − 2y)(2x + 3y)}
\end{quote}

\textbf{(x+5)(x-3) = কত?}

ক) x\^{}2 + 8x + 15

খ) x\^{}2 − 15

\textbf{\hl{গ) x\^{}2 + 2x − 15 ✔}}

ঘ) x\^{}2 + 2x + 15

\textbf{\hl{✔ সঠিক উত্তর: গ) x\^{}2 + 2x − 15}}

\begin{quote}
\textbf{ব্যাখ্যা:} \emph{(x + 5)(x − 3)}

\emph{= x\^{}2 + 5x − 3x − 15}

\emph{= x\^{}2 + 2x − 15}
\end{quote}

\textbf{x\^{}2 − x − 2 এর একটি উৎপাদক?}

ক) x + 2

খ) x - 3

\textbf{\hl{গ) x + 1 ✔}}

ঘ) x - 1

\textbf{\hl{✔ সঠিক উত্তর: গ) x + 1}}

\begin{quote}
\textbf{ব্যাখ্যা:} \emph{x\^{}2 − x − 2}

\emph{= x\^{}2 − 2x + x − 2}

\emph{= x(x − 2) + 1(x − 2)}

\emph{= (x − 2)(x + 1)}

\emph{∴ একটি উৎপাদক হলো (x + 1)}
\end{quote}

\textbf{2x\^{}2 − x − 3 এর একটি উৎপাদক নিচের কোনটি?}

ক) 2x+2

খ) x-1

\textbf{\hl{গ) x+1 ✔}}

ঘ) 2x+3x

\textbf{\hl{✔ সঠিক উত্তর: গ) x+1}}

\begin{quote}
\textbf{ব্যাখ্যা:} \emph{2x\^{}2 − x − 3}

\emph{= 2x\^{}2 − 3x + 2x − 3}

\emph{= x(2x − 3) + 1(2x − 3)}

\emph{= (2x − 3)(x + 1)}

\emph{∴ নির্ণেয় একটি উৎপাদক (x + 1)}
\end{quote}

\textbf{−4a\^{}2 + 23a + 6 এর উৎপাদক কোনটি?}

ক) (6-a) (4a-1)

খ) (a-6) (4a+1)

গ) (a-6) (4a-1)

\textbf{\hl{ঘ) (6-a) (4a + 1) ✔}}

\textbf{\hl{✔ সঠিক উত্তর: ঘ) (6-a) (4a + 1)}}

\begin{quote}
\textbf{ব্যাখ্যা:} \emph{−4a\^{}2 + 23a + 6}

\emph{= −(4a\^{}2 − 23a − 6)}

\emph{= −(4a\^{}2 − 24a + a − 6)}

\emph{= −\{4a(a − 6) + 1(a − 6)\}}

\emph{= −(a − 6)(4a + 1)}

\emph{= (6 − a)(4a + 1)}
\end{quote}

\textbf{x\^{}2 − 7x + 6 এর উৎপাদকে বিশ্লেষিত রূপ নিচের কোনটি?}

ক) (x-2)(x-3)

খ) (x-1)(x+8)

\textbf{\hl{গ) (x-1)(x-6) ✔}}

ঘ) (x+1)(x+6)

\textbf{\hl{✔ সঠিক উত্তর: গ) (x-1)(x-6)}}

\begin{quote}
\textbf{ব্যাখ্যা:} \emph{x\^{}2 − 7x + 6}

\emph{= x\^{}2 − 6x − x + 6}

\emph{= x(x − 6) − 1(x − 6)}

\emph{= (x − 6)(x − 1)}
\end{quote}

\textbf{x\^{}2 + 7x + p যদি x − 5 দ্বারা বিভাজ্য হয় তবে, p এর মান কত হবে?}

ক) 30

খ) −30

\textbf{\hl{গ) −60 ✔}}

ঘ) 5

\textbf{\hl{✔ সঠিক উত্তর: গ) −60}}

\begin{quote}
\textbf{ব্যাখ্যা:} \emph{যেহেতু (x − 5) দ্বারা x\^{}2 + 7x + p বিভাজ্য, তাই x
= 5 বসালে রাশিটির মান শূন্য হবে।}

\emph{(5)\^{}2 + 7(5) + p = 0}

\emph{⇒ 25 + 35 + p = 0}

\emph{⇒ 60 + p = 0}

\emph{∴ p = −60}
\end{quote}

\textbf{2x\^{}2 − x − 15 এর উৎপাদক হবে-}

ক) (x+6)(x-5)

খ) (x-5)(x-6)

\textbf{\hl{গ) (2x+5)(x-3) ✔}}

ঘ) (x+3)(2x-5)

\textbf{\hl{✔ সঠিক উত্তর: গ) (2x+5)(x-3)}}

\begin{quote}
\textbf{ব্যাখ্যা:} \emph{2x\^{}2 − x − 15}

\emph{= 2x\^{}2 − 6x + 5x − 15}

\emph{= 2x(x − 3) + 5(x − 3)}

\emph{= (2x + 5)(x − 3)}
\end{quote}

\textbf{x\^{}2 + x − 20 এর উৎপাদক কত?}

\textbf{\hl{ক) (x+5)(x-4) ✔}}

খ) (x-5) (x+5)

গ) (x-5) (x+4)

ঘ) (x-4) (x+4)

\textbf{\hl{✔ সঠিক উত্তর: ক) (x+5)(x-4)}}

\begin{quote}
\textbf{ব্যাখ্যা:} \emph{x\^{}2 + x − 20}

\emph{= x\^{}2 + 5x − 4x − 20}

\emph{= x(x + 5) − 4(x + 5)}

\emph{= (x + 5)(x − 4)}
\end{quote}

\textbf{x\^{}2 − 9x + 20 এর উৎপাদকসমূহ কোনটি?}

ক) (x+5) (x+4)

\textbf{\hl{খ) (x-5) (x-4) ✔}}

গ) (x+5) (x-4)

ঘ) (x-5) (x+4)

\textbf{\hl{✔ সঠিক উত্তর: খ) (x-5) (x-4)}}

\begin{quote}
\textbf{ব্যাখ্যা:} \emph{x\^{}2 − 9x + 20}

\emph{= x\^{}2 − 5x − 4x + 20}

\emph{= x(x − 5) − 4(x − 5)}

\emph{= (x − 5)(x − 4)}
\end{quote}

\textbf{3x\^{}2 − 16x − 12 রাশির উৎপাদক বিশ্লেষণ নিচের কোনটি?}

ক) (3x - 2)(x + 6)

খ) (3x - 2)(x - 6)

\textbf{\hl{গ) (x - 6)(3x + 2) ✔}}

ঘ) 3(x-6)(x-2)

\textbf{\hl{✔ সঠিক উত্তর: গ) (x - 6)(3x + 2)}}

\begin{quote}
\textbf{ব্যাখ্যা:} \emph{3x\^{}2 − 16x − 12}

\emph{= 3x\^{}2 − 18x + 2x − 12}

\emph{= 3x(x − 6) + 2(x − 6)}

\emph{= (x − 6)(3x + 2)}
\end{quote}

\textbf{2x\^{}2 − 5xy + 2y\^{}2 এর একটি উৎপাদক হবে-}

ক) 2y-x

\textbf{\hl{খ) x-2y ✔}}

গ) 2(x-y)

ঘ) x-3y

\textbf{\hl{✔ সঠিক উত্তর: খ) x-2y}}

\begin{quote}
\textbf{ব্যাখ্যা:} \emph{2x\^{}2 − 5xy + 2y\^{}2}

\emph{= 2x\^{}2 − 4xy − xy + 2y\^{}2}

\emph{= 2x(x − 2y) − y(x − 2y)}

\emph{= (x − 2y)(2x − y)}
\end{quote}

\textbf{2a\^{}2 + 7ab − 15b\^{}2 উৎপাদকে বিশ্লেষণে কত হয়?}

\textbf{\hl{ক) (a + 5b) (2a - 3b) ✔}}

খ) (a - 5b) (2a + 3b)

গ) (a - 5b) (2a - 3b)

ঘ) ({[}Missing option{]})

\textbf{\hl{✔ সঠিক উত্তর: ক) (a + 5b) (2a - 3b)}}

\begin{quote}
\textbf{ব্যাখ্যা:} \emph{2a\^{}2 + 7ab − 15b\^{}2}

\emph{= 2a\^{}2 + 10ab − 3ab − 15b\^{}2}

\emph{= 2a(a + 5b) − 3b(a + 5b)}

\emph{= (a + 5b)(2a − 3b)}
\end{quote}

\textbf{2a\^{}2 + 6a − 80 এর একটি উৎপাদক কোনটি?}

ক) 2(a-8)

খ) (1+5)

গ) (a+4)

\textbf{\hl{ঘ) (a + 8) ✔}}

\textbf{\hl{✔ সঠিক উত্তর: ঘ) (a + 8)}}

\begin{quote}
\textbf{ব্যাখ্যা:} \emph{2a\^{}2 + 6a − 80}

\emph{= 2(a\^{}2 + 3a − 40)}

\emph{= 2(a\^{}2 + 8a − 5a − 40)}

\emph{= 2\{a(a + 8) − 5(a + 8)\}}

\emph{= 2(a + 8)(a − 5)}

\emph{∴ নির্ণেয় একটি উৎপাদক (a + 8)}
\end{quote}

\textbf{x\^{}6 + 4x\^{}3 − 1 এর একটি উৎপাদক-}

ক) x\^{}2 − x + 1

\textbf{\hl{খ) x\^{}2 + x − 1 ✔}}

গ) x\^{}3 − 1

ঘ) x\^{}2 − x − 1

\textbf{\hl{✔ সঠিক উত্তর: খ) x\^{}2 + x − 1}}

\begin{quote}
\textbf{ব্যাখ্যা:} \emph{x\^{}6 + 4x\^{}3 − 1}

\emph{= (x\^{}2)\^{}3 + x\^{}3 + (−1)\^{}3 − 3x\^{}2 · x · (−1)}

\emph{= (x\^{}2 + x − 1)((x\^{}2)\^{}2 + x\^{}2 + (−1)\^{}2 − x\^{}2 · x
− x · (−1) − (−1) · x\^{}2)}

\emph{= (x\^{}2 + x − 1)(x\^{}4 + x\^{}2 + 1 − x\^{}3 + x + x\^{}2)}

\emph{= (x\^{}2 + x − 1)(x\^{}4 − x\^{}3 + 2x\^{}2 + x + 1)}

\emph{∴ একটি উৎপাদক হলো (x\^{}2 + x − 1)}
\end{quote}

\textbf{উৎপাদকে বিশ্লেষণ করুন: x\^{}2 − x − (b + 1) (b + 2)}

\textbf{\hl{ক) (x-b-2)(x+b+1) ✔}}

খ) (x-b-2)(x+b-1)

গ) (x-b+2)(x+b+1)

ঘ) (x+b-2)(x+b+1)

\textbf{\hl{✔ সঠিক উত্তর: ক) (x-b-2)(x+b+1)}}

\begin{quote}
\textbf{ব্যাখ্যা:} \emph{x\^{}2 − x − (b + 1)(b + 2)}

\emph{= x\^{}2 − x(b + 2) + x(b + 1) − (b + 1)(b + 2)}

\emph{= x(x − b − 2) + (b + 1)(x − b − 2)}

\emph{= (x − b − 2)(x + b + 1)}
\end{quote}

\textbf{x\^{}2 + x − (a + 1) (a + 2) কে উৎপাদকে বিশ্লেষণ করুন।}

ক) (x-a+2)(x-a-1)

খ) (x+a-1)(x+a+2)

\textbf{\hl{গ) (x+a+2)(x-a-1) ✔}}

ঘ) (x+a+1)(x-a-2)

\textbf{\hl{✔ সঠিক উত্তর: গ) (x+a+2)(x-a-1)}}

\begin{quote}
\textbf{ব্যাখ্যা:} \emph{x\^{}2 + x − (a + 1)(a + 2)}

\emph{= x\^{}2 + x(a + 2) − x(a + 1) − (a + 1)(a + 2)}

\emph{= x(x + a + 2) − (a + 1)(x + a + 2)}

\emph{= (x + a + 2)(x − a − 1)}
\end{quote}

\textbf{4a\^{}2 + 11a + 6 = 0 হলে a = ?}

ক) -2

খ) -0.75

গ) 2

\textbf{\hl{ঘ) ক এবং খ উভয়ই ✔}}

\textbf{\hl{✔ সঠিক উত্তর: ঘ) ক এবং খ উভয়ই}}

\begin{quote}
\textbf{ব্যাখ্যা:} \emph{4a\^{}2 + 11a + 6 = 0}

\emph{⇒ 4a\^{}2 + 8a + 3a + 6 = 0}

\emph{⇒ 4a(a + 2) + 3(a + 2) = 0}

\emph{⇒ (a + 2)(4a + 3) = 0}

\emph{হয় a + 2 = 0 অথবা 4a + 3 = 0}

\emph{∴ a = −2, বা, a = −3/4, বা, a = −0.75}
\end{quote}

\textbf{x\^{}3 + 3x + 36 এর একটি উৎপাদক নিচের কোনটি?}

\textbf{\hl{ক) x + 3 ✔}}

খ) x + 4

গ) x - 3

ঘ) x + 6

\textbf{\hl{✔ সঠিক উত্তর: ক) x + 3}}

\begin{quote}
\textbf{ব্যাখ্যা:} \emph{ধরি, f(x) = x\^{}3 + 3x − 36}

\emph{f(-3) = (-3)\^{}3 + 3(-3) + 36}

\emph{= -27 − 9 + 36 = 0}

\emph{∴ (x + 3) হলো f(x) এর একটি উৎপাদক।}
\end{quote}

\textbf{3x\^{}3 + 2x\^{}2 − 21x − 20 রাশির একটি উৎপাদক হচ্ছে-}

ক) x + 2

খ) x - 2

\textbf{\hl{গ) x + 1 ✔}}

ঘ) x - 1

\textbf{\hl{✔ সঠিক উত্তর: গ) x + 1}}

\begin{quote}
\textbf{ব্যাখ্যা:} \emph{ধরি, p(x) = 3x\^{}3 + 2x\^{}2 − 21x − 20}

\emph{এখানে, p(-1) = 3(-1)\^{}3 + 2(-1)\^{}2 − 21(-1) − 20}

\emph{= −3 + 2 + 21 − 20 = 0}

\emph{∴ (x + 1), p(x) এর একটি উৎপাদক।}
\end{quote}

\textbf{x\^{}2 − y\^{}2 + 2x + 1 এর একটি উৎপাদক কত?}

ক) x + y - 1

খ) 1 - x - y

গ) x - y - 1

\textbf{\hl{ঘ) x + y + 1 ✔}}

\textbf{\hl{✔ সঠিক উত্তর: ঘ) x + y + 1}}

\begin{quote}
\textbf{ব্যাখ্যা:} \emph{x\^{}2 − y\^{}2 + 2x + 1 = x\^{}2 + 2 · x · 1 +
1\^{}2 − y\^{}2}

\emph{= (x + 1)\^{}2 − y\^{}2}

\emph{= (x + 1 − y)(x + 1 + y)}
\end{quote}

\textbf{4x\^{}4 + 1 কে উৎপাদকে বিশ্লেষণ করুন।}

ক) (2x\^{}2+2x-1) (2x\^{}2-2x+1)

\textbf{\hl{খ) (2x\^{}2+2x+1) (2x\^{}2-2x+1) ✔}}

গ) (2x\^{}2+2x-1) (2x\^{}2-2x-1)

ঘ) (2x\^{}2+2x+1) (2x\^{}2-2x-1)

\textbf{\hl{✔ সঠিক উত্তর: খ) (2x\^{}2+2x+1) (2x\^{}2-2x+1)}}

\begin{quote}
\textbf{ব্যাখ্যা:} \emph{4x\^{}4 + 1}

\emph{= (2x\^{}2)\^{}2 + (1)\^{}2}

\emph{= (2x\^{}2 + 1)\^{}2 − 2 · 2x\^{}2 · 1}

\emph{= (2x\^{}2 + 1)\^{}2 − (2x)\^{}2}

\emph{= (2x\^{}2 + 1 + 2x)(2x\^{}2 + 1 − 2x)}

\emph{= (2x\^{}2 + 2x + 1)(2x\^{}2 − 2x + 1)}
\end{quote}

\textbf{উৎপাদকে বিশ্লেষণ করুন: x\^{}2 − 2ax + (a + b) (a − b)}

ক) (x-a+b)(x+a-b)

\textbf{\hl{খ) (x-a-b)(x-a+b) ✔}}

গ) (x+a-b)(x-a-b)

ঘ) (x+a+b)(x-a-b)

\textbf{\hl{✔ সঠিক উত্তর: খ) (x-a-b)(x-a+b)}}

\begin{quote}
\textbf{ব্যাখ্যা:} \emph{x\^{}2 − 2ax + (a + b)(a − b)}

\emph{= x\^{}2 − \{(a + b) + (a − b)\}x + (a + b)(a − b)}

\emph{= (x − a − b)(x − a + b)}
\end{quote}

\textbf{3x\^{}3 + 2x\^{}2 - 21x - 20 রাশিটির একটি উৎপাদক হচ্ছে -}

ক) x+2

খ) x-2

\textbf{\hl{গ) x+1 ✔}}

ঘ) x - 1

\textbf{\hl{✔ সঠিক উত্তর: গ) x+1}}

\begin{quote}
\textbf{ব্যাখ্যা:} \emph{3x\^{}3 + 2x\^{}2 − 21x − 20}

\emph{= 3x\^{}3 + 3x\^{}2 − x\^{}2 − x − 20x − 20}

\emph{= 3x\^{}2(x + 1) − x(x + 1) − 20(x + 1)}

\emph{= (x + 1)(3x\^{}2 − x − 20)}
\end{quote}

\textbf{x\^{}3 - 6x\^{}2 + 11x - 6 এর উৎপাদক কোনটি?}

ক) (x + 1 ) ( x - 2) (x - 3)

খ) (x - 1 ) ( x - 2) (x + 3)

গ) (x - 1 ) ( x + 2) (x - 3)

\textbf{\hl{ঘ) (x - 1 ) ( x - 2) (x - 3) ✔}}

\textbf{\hl{✔ সঠিক উত্তর: ঘ) (x - 1 ) ( x - 2) (x - 3)}}

\begin{quote}
\textbf{ব্যাখ্যা:} \emph{x\^{}3 − 6x\^{}2 + 11x − 6}

\emph{= x\^{}3 − x\^{}2 − 5x\^{}2 + 5x + 6x − 6}

\emph{= x\^{}2(x − 1) − 5x(x − 1) + 6(x − 1)}

\emph{= (x − 1)(x\^{}2 − 5x + 6)}

\emph{= (x − 1)(x\^{}2 − 3x − 2x + 6)}

\emph{= (x − 1)\{x(x − 3) − 2(x − 3)\}}

\emph{= (x − 1)(x − 2)(x − 3)}
\end{quote}

\textbf{4x\^{}4 - 25x\^{}2 + 36 = কত?}

ক) (x-2) (x-2) (2x+3) (2x-3)

খ) (x+2) (x+2) (2x+3) (2x-3)

গ) (x+2) (x-2) (2x-3) (2x-3)

\textbf{\hl{ঘ) (x+2) (x-2) (2x+3) (2x-3) ✔}}

\textbf{\hl{✔ সঠিক উত্তর: ঘ) (x+2) (x-2) (2x+3) (2x-3)}}

\begin{quote}
\textbf{ব্যাখ্যা:} \emph{4x\^{}4 − 25x\^{}2 + 36}

\emph{= 4x\^{}4 − 16x\^{}2 − 9x\^{}2 + 36}

\emph{= 4x\^{}2(x\^{}2 − 4) − 9(x\^{}2 − 4)}

\emph{= (x\^{}2 − 4)(4x\^{}2 − 9)}

\emph{= (x\^{}2 − 2\^{}2)((2x)\^{}2 − 3\^{}2)}

\emph{= (x + 2)(x − 2)(2x + 3)(2x − 3)}
\end{quote}

\textbf{a\^{}4 + a\^{}2b\^{}2 + b\^{}4 এর একটি উৎপাদক a\^{}2 + ab +
b\^{}2 হলে অপর উৎপাদক কত?}

ক) a\^{}2+ab+b\^{}2

\textbf{\hl{খ) a\^{}2−ab+b\^{}2 ✔}}

গ) a\^{}2+b\^{}2

ঘ) a\^{}2−b\^{}2

\textbf{\hl{✔ সঠিক উত্তর: খ) a\^{}2−ab+b\^{}2}}

\begin{quote}
\textbf{ব্যাখ্যা:} \emph{a\^{}4 + a\^{}2b\^{}2 + b\^{}4}

\emph{= (a\^{}2)\^{}2 + 2 · a\^{}2 · b\^{}2 + (b\^{}2)\^{}2 −
a\^{}2b\^{}2}

\emph{= (a\^{}2 + b\^{}2)\^{}2 − (ab)\^{}2}

\emph{= (a\^{}2 + b\^{}2 + ab)(a\^{}2 + b\^{}2 − ab)}

\emph{= (a\^{}2 + ab + b\^{}2)(a\^{}2 − ab + b\^{}2)}

\emph{∴ অপর উৎপাদকটি হলো a\^{}2 − ab + b\^{}2}
\end{quote}

\textbf{2y\^{}4 − 14y\^{}2 + 2 এর উৎপাদকের বিশ্লেষণ কোনটি?}

ক) (2y\^{}2+6y+2)(y\^{}2-3y+1)

খ) (2y\^{}2-7y+1)(2y\^{}2+7y+1)

গ) (2y\^{}2+3y+1)(y\^{}2-3y+1)

\textbf{\hl{ঘ) 2(y\^{}2+3y+1)(y\^{}2-3y+1) ✔}}

\textbf{\hl{✔ সঠিক উত্তর: ঘ) 2(y\^{}2+3y+1)(y\^{}2-3y+1)}}

\begin{quote}
\textbf{ব্যাখ্যা:} \emph{2y\^{}4 − 14y\^{}2 + 2}

\emph{= 2(y\^{}4 − 7y\^{}2 + 1)}

\emph{= 2\{(y\^{}2)\^{}2 + 2 · y\^{}2 · 1 + 1\^{}2 − 9y\^{}2\}}

\emph{= 2\{(y\^{}2 + 1)\^{}2 − (3y)\^{}2\}}

\emph{= 2(y\^{}2 + 3y + 1)(y\^{}2 − 3y + 1)}
\end{quote}

\textbf{2 √(2) x\^{}3 + 125 এর সঠিক উৎপাদকের বিশ্লেষণ কোনটি?}

ক) ( √(2) x+5)(2x\^{}2−5 √(2) x−25)

\textbf{\hl{খ) ( √(2) x+5)(2x\^{}2−5 √(2) x+25) ✔}}

গ) ( √(2) x+5)(2x\^{}2+5 √(2) x−25)

ঘ) ( √(2) x−5)(2x\^{}2+5 √(2) x+25)

\textbf{\hl{✔ সঠিক উত্তর: খ) ( √(2) x+5)(2x\^{}2−5 √(2) x+25)}}

\begin{quote}
\textbf{ব্যাখ্যা:} \emph{2 √(2) x\^{}3 + 125}

\emph{= ( √(2) x)\^{}3 + (5)\^{}3}

\emph{= ( √(2) x + 5)\{( √(2) x)\^{}2 − √(2) x · 5 + 5\^{}2\}}

\emph{= ( √(2) x + 5)(2x\^{}2 − 5 √(2) x + 25)}
\end{quote}

\textbf{(a − b)\^{}3 + (b − c)\^{}3 + (c − a)\^{}3 কে উৎপাদকে বিশ্লেষণ
করুন?}

ক) 9(a-b) (b-c) (c-a)

\textbf{\hl{খ) 3(a-b) (b-c) (c-a) ✔}}

গ) 2(a-b) (b-c) (c-a)

ঘ) (a-b) (b-c) (c-a)

\textbf{\hl{✔ সঠিক উত্তর: খ) 3(a-b) (b-c) (c-a)}}

\begin{quote}
\textbf{ব্যাখ্যা:} \emph{ধরি, a − b = x, b − c = y, c − a = z}

\emph{এখানে, x + y + z = a − b + b − c + c − a = 0}

\emph{আমরা জানি, যদি x + y + z = 0 হয়, তবে x\^{}3 + y\^{}3 + z\^{}3 =
3xyz}

\emph{∴ 3(a − b)(b − c)(c − a)}
\end{quote}

\textbf{x\^{}2 − 1 − y (y − 2) এর উৎপাদক কত?}

ক) (x-y) (x-y+1)

\textbf{\hl{খ) (x-y+1) (x+y-1) ✔}}

গ) (x+y+1) (x-y-1)

ঘ) (x-y) (x+y+1)

\textbf{\hl{✔ সঠিক উত্তর: খ) (x-y+1) (x+y-1)}}

\begin{quote}
\textbf{ব্যাখ্যা:} \emph{x\^{}2 − 1 − y(y − 2)}

\emph{= x\^{}2 − 1 − y\^{}2 + 2y}

\emph{= x\^{}2 − (y\^{}2 − 2y + 1)}

\emph{= x\^{}2 − (y − 1)\^{}2}

\emph{= \{x + (y − 1)\}\{x − (y − 1)\}}

\emph{= (x + y − 1)(x − y + 1)}
\end{quote}

\textbf{নিচের কোনটি x\^{}3 + 6x\^{}2 + 11x - 6 এর উৎপাদক নয়?}

ক) x-1

খ) x-2

গ) x-3

\textbf{\hl{ঘ) x-4 ✔}}

\textbf{\hl{✔ সঠিক উত্তর: ঘ) x-4}}

\begin{quote}
\textbf{ব্যাখ্যা:} \emph{x = 4 বসালে x\^{}3 + 6x\^{}2 + 11x − 6 সমীকরণের
মান শূন্য হয় না (বরং 64 + 96 + 44 − 6 = 198 ≠ 0)। তাই (x − 4) উৎপাদক হবে
না।}
\end{quote}

\textbf{x\^{}2 − x − 6 = 0 সমীকরণের মূলদ্বয় হবে}

ক) 3, 2

\textbf{\hl{খ) 3, -2 ✔}}

গ) -3, 2

ঘ) -3, -2

\textbf{\hl{✔ সঠিক উত্তর: খ) 3, -2}}

\begin{quote}
\textbf{ব্যাখ্যা:} \emph{x\^{}2 − x − 6 = 0}

\emph{⇒ x\^{}2 − 3x + 2x − 6 = 0}

\emph{⇒ x(x − 3) + 2(x − 3) = 0}

\emph{⇒ (x − 3)(x + 2) = 0}

\emph{হয়, x − 3 = 0 ⇒ x = 3}

\emph{অথবা, x + 2 = 0 ⇒ x = −2}

\emph{∴ সমীকরণের মূলদ্বয় হবে 3, −2}
\end{quote}

\textbf{p\^{}2 − 4p + 1 = 0 সমীকরণের মূলদ্বয় α ও β হলে α + β ও αβ মূলবিশিষ্ট
সমীকরণ কোনটি?}

ক) p\^{}2 + 4p - 5 = 0

\textbf{\hl{খ) p\^{}2 − 5p + 4 = 0 ✔}}

গ) p\^{}2 − 4p + 5 = 0

ঘ) p\^{}2 + 5p − 4 = 0

\textbf{\hl{✔ সঠিক উত্তর: খ) p\^{}2 − 5p + 4 = 0}}

\begin{quote}
\textbf{ব্যাখ্যা:} \emph{p\^{}2 − 4p + 1 = 0 সমীকরণের মূলদ্বয় α ও β হলে,}

\emph{α + β = 4 এবং αβ = 1}

\emph{নতুন সমীকরণের মূলদ্বয় α + β = 4 এবং αβ = 1}

\emph{মূলদ্বয়ের যোগফল = 4 + 1 = 5}

\emph{মূলদ্বয়ের গুণফল = 4 × 1 = 4}

\emph{∴ নির্ণেয় সমীকরণ: p\^{}2 − (যোগফল)p + (গুণফল) = 0}

\emph{⇒ p\^{}2 − 5p + 4 = 0}
\end{quote}

\textbf{কোনটি 2x\^{}4 − 5x\^{}3 + 6x\^{}2 − 5x + 2 এর একটি উৎপাদক?}

ক) x+1

\textbf{\hl{খ) x-1 ✔}}

গ) x+2

ঘ) x-2

\textbf{\hl{✔ সঠিক উত্তর: খ) x-1}}

\begin{quote}
\textbf{ব্যাখ্যা:} \emph{x = 1 বসালে,}

\emph{2(1)\^{}4 − 5(1)\^{}3 + 6(1)\^{}2 − 5(1) + 2}

\emph{= 2 − 5 + 6 − 5 + 2 = 0}

\emph{∴ (x − 1) হবে রাশিটির একটি উৎপাদক।}
\end{quote}

\textbf{x\^{}3 + 6x\^{}2y + 11xy\^{}2 + 6y\^{}3 এর উৎপাদকে বিশ্লেষণ
কোনটি?}

ক) (x + y) (x +3y) (x +5y)

\textbf{\hl{খ) (x+y) (x + 2y)(x+3y) ✔}}

গ) (x + y) (x +4y) (x +3y)

ঘ) (x −y) (x + y) (x + 2y)

\textbf{\hl{✔ সঠিক উত্তর: খ) (x+y) (x + 2y)(x+3y)}}

\begin{quote}
\textbf{ব্যাখ্যা:} \emph{x\^{}3 + 6x\^{}2y + 11xy\^{}2 + 6y\^{}3}

\emph{= x\^{}3 + x\^{}2y + 5x\^{}2y + 5xy\^{}2 + 6xy\^{}2 + 6y\^{}3}

\emph{= x\^{}2(x + y) + 5xy(x + y) + 6y\^{}2(x + y)}

\emph{= (x + y)(x\^{}2 + 5xy + 6y\^{}2)}

\emph{= (x + y)(x\^{}2 + 3xy + 2xy + 6y\^{}2)}

\emph{= (x + y)\{x(x + 3y) + 2y(x + 3y)\}}

\emph{= (x + y)(x + 3y)(x + 2y)}
\end{quote}

\textbf{a\^{}3 − 21a − 20 এর উৎপাদকে বিশ্লেষণ করুন।}

ক) (a + 2)

\textbf{\hl{খ) (a + 1) ✔}}

গ) (a − 2 )

ঘ) (a − 1 )

\textbf{\hl{✔ সঠিক উত্তর: খ) (a + 1)}}

\begin{quote}
\textbf{ব্যাখ্যা:} \emph{a\^{}3 − 21a − 20}

\emph{= a\^{}3 + a\^{}2 − a\^{}2 − a − 20a − 20}

\emph{= a\^{}2(a + 1) − a(a + 1) − 20(a + 1)}

\emph{= (a + 1)(a\^{}2 − a − 20)}

\emph{= (a + 1)(a\^{}2 − 5a + 4a − 20)}

\emph{= (a + 1)\{a(a − 5) + 4(a − 5)\}}

\emph{= (a + 1)(a − 5)(a + 4)}
\end{quote}

\textbf{a\^{}3 − 7a − 6 এর উৎপাদক কত?}

ক) (a+1)(a-2)(a-3)

\textbf{\hl{খ) (a+1)(a+2)(a-3) ✔}}

গ) (a-1)(a+2)(a-3)

ঘ) (a-1)(a-2)(a-3)

\textbf{\hl{✔ সঠিক উত্তর: খ) (a+1)(a+2)(a-3)}}

\begin{quote}
\textbf{ব্যাখ্যা:} \emph{a\^{}3 − 7a − 6}

\emph{= a\^{}3 + a\^{}2 − a\^{}2 − a − 6a − 6}

\emph{= a\^{}2(a + 1) − a(a + 1) − 6(a + 1)}

\emph{= (a + 1)(a\^{}2 − a − 6)}

\emph{= (a + 1)(a\^{}2 + 2a − 3a − 6)}

\emph{= (a + 1)\{a(a + 2) − 3(a + 2)\}}

\emph{= (a + 1)(a + 2)(a − 3)}
\end{quote}

\textbf{x\^{}3 − x − 6 এর উৎপাদকে বিশ্লেষিত রূপ কোনটি?}

ক) (𝑥 + 2)(𝑥\^{}2 + 2𝑥 + 3)

\textbf{\hl{খ) (𝑥 − 2)(𝑥\^{}2 + 2𝑥 + 3) ✔}}

গ) (𝑥 − 2)(𝑥\^{}2 − 2𝑥 + 3)

ঘ) (𝑥 + 2)(𝑥\^{}2 − 2𝑥 + 3)

\textbf{\hl{✔ সঠিক উত্তর: খ) (𝑥 − 2)(𝑥\^{}2 + 2𝑥 + 3)}}

\begin{quote}
\textbf{ব্যাখ্যা:} \emph{x\^{}3 − x − 6}

\emph{= x\^{}3 − 2x\^{}2 + 2x\^{}2 − 4x + 3x − 6}

\emph{= x\^{}2(x − 2) + 2x(x − 2) + 3(x − 2)}

\emph{= (x − 2)(x\^{}2 + 2x + 3)}
\end{quote}

\textbf{x\^{}2 + 7𝑥 + p যদি x-5 দ্বারা বিভাজ্য হয় তবে p এর মান কত হবে?}

ক) 30

খ) - 30

\textbf{\hl{গ) - 60 ✔}}

ঘ) 5

\textbf{\hl{✔ সঠিক উত্তর: গ) - 60}}

\begin{quote}
\textbf{ব্যাখ্যা:} \emph{যেহেতু x − 5 রাশিটির একটি উৎপাদক, তাই x = 5 হলে
রাশিটির মান শূন্য হবে।}

\emph{সুতরাং, 5\^{}2 + 7 × 5 + p = 0}

\emph{⇒ 25 + 35 + p = 0}

\emph{⇒ 60 + p = 0}

\emph{∴ p = −60}
\end{quote}

\textbf{f(x) = log √(x\^{}2−36) ফাংশন এর ডোমেন কত?}

ক) x = 6

খ) x \textless{} 6

\textbf{\hl{গ) x ≥ 6 ✔}}

ঘ) x \textgreater{} 6

\textbf{\hl{✔ সঠিক উত্তর: গ) x ≥ 6}}

\begin{quote}
\textbf{ব্যাখ্যা:} \emph{f(x) এর মান বাস্তব হওয়ার জন্য
x\^{}2−36\textgreater0 অথবা x\^{}2−36≥0 হতে হবে।}

\emph{এখানে, x\^{}2−36≥0}

\emph{⇒ x\^{}2≥36}

\emph{∴ x≥6 অথবা x≤−6 (অপশন অনুযায়ী সঠিক উত্তর x≥6 )}
\end{quote}

\textbf{f(x) এর একটি উৎপাদক (3x + 2) হলে নিচের কোনটির মান শূন্য হবে?}

\textbf{\hl{ক) f( -(3/2) ) ✔}}

খ) f(2)

গ) f( (3/2) )

ঘ) f(-3)

\textbf{\hl{✔ সঠিক উত্তর: ক) f( -(3/2) )}}

\begin{quote}
\textbf{ব্যাখ্যা:} \emph{যেহেতু (3x + 2) হলো f(x) এর একটি উৎপাদক, তাই 3x +
2 = 0 ⇒ x = -(2/3) বসালে f(x) এর মান শূন্য হবে। তবে প্রশ্নে প্রদত্ত অপশন
অনুযায়ী 3x + 2 এর জায়গায় (2x + 3) থাকলে x = -(3/2) এর জন্য f( -(3/2) ) = 0
হয়।}
\end{quote}

\textbf{f(x) = (x+3/x−1) এর Domain কত?}

\textbf{\hl{ক) \{x ∈ R: x≠1\} ✔}}

খ) R − \{ (3/2) \}

গ) R − \{1, -1\}

ঘ) \{1\}

\textbf{\hl{✔ সঠিক উত্তর: ক) \{x ∈ R: x≠1\}}}

\begin{quote}
\textbf{ব্যাখ্যা:} \emph{x এর যে সকল বাস্তব মানের জন্য f(x) সংজ্ঞায়িত, সেগুলোর
সেটই Domain। এখানে x − 1 = 0 বা x = 1 হলে হর শূন্য হয়, ফলে ভগ্নাংশটি
অসংজ্ঞায়িত হয়।}

\emph{∴ নির্ণেয় Domain হলো \{x ∈ R: x ≠ 1\}}
\end{quote}

\textbf{f(x) = 2x\^{}2 − 3x + 4 হলে f(1) এর মান কত?}

ক) 2

খ) 1

গ) 0

\textbf{\hl{ঘ) 3 ✔}}

\textbf{\hl{✔ সঠিক উত্তর: ঘ) 3}}

\begin{quote}
\textbf{ব্যাখ্যা:} \emph{f(x) = 2x\^{}2 − 3x + 4}

\emph{⇒ f(1) = 2(1)\^{}2 − 3(1) + 4}

\emph{= 2 − 3 + 4 = 3}
\end{quote}

\textbf{f(x) = x\^{}3 + 7 হলে f\^{}-1 (34) = ?}

\textbf{\hl{ক) 3 ✔}}

খ) 9

গ) ২১

ঘ) ২৭

\textbf{\hl{✔ সঠিক উত্তর: ক) 3}}

\begin{quote}
\textbf{ব্যাখ্যা:} \emph{ধরি, R = f\^{}-1(34)}

\emph{⇒ f(R) = 34}

\emph{⇒ R\^{}3 + 7 = 34}

\emph{⇒ R\^{}3 = 34 − 7 = 27}

\emph{∴ R = 3}
\end{quote}

\textbf{f(x) = 2x\^{}2 + 3x − 1 হলে f(0) = কত?}

ক) 6

খ) 4

গ) 1

\textbf{\hl{ঘ) -1 ✔}}

\textbf{\hl{✔ সঠিক উত্তর: ঘ) -1}}

\begin{quote}
\textbf{ব্যাখ্যা:} \emph{f(x) = 2x\^{}2 + 3x − 1}

\emph{⇒ f(0) = 2(0)\^{}2 + 3(0) − 1 = −1}
\end{quote}

\textbf{f(x) = x\^{}2 − 4x + 4 হলে f(2) = ?}

\textbf{\hl{ক) 0 ✔}}

খ) 2

গ) 3

ঘ) 4

\textbf{\hl{✔ সঠিক উত্তর: ক) 0}}

\begin{quote}
\textbf{ব্যাখ্যা:} \emph{f(x) = x\^{}2 − 4x + 4}

\emph{⇒ f(2) = (2)\^{}2 − 4(2) + 4}

\emph{= 4 − 8 + 4 = 0}
\end{quote}

\textbf{f(x) = √(4−x\^{}2) হলে f এর রেঞ্জ কত?}

ক) -2 ≤ x ≤ 0

খ) 0 ≤ x ≤ 2

\textbf{\hl{গ) -2 ≤ x ≤ 2 ✔}}

ঘ) -4 ≤ x ≤ 4

\textbf{\hl{✔ সঠিক উত্তর: গ) -2 ≤ x ≤ 2}}

\begin{quote}
\textbf{ব্যাখ্যা:} \emph{f(x) = √(4−x\^{}2)}

\emph{এখানে, 4−x\^{}2≥0 হওয়ায় x\^{}2≤4 , অর্থাৎ −2≤x≤2 । রুটযুক্ত পজিটিভ
রাশির মান সর্বনিম্ন 0 এবং সর্বোচ্চ 2 হতে পারে, তবে রেঞ্জের ক্ষেত্রে সার্বিকভাবে
মানটি −2 থেকে 2 এর মধ্যে থাকে।}
\end{quote}

\textbf{f(x) = x\^{}3 − 6x\^{}2 + 11x − 6 হলে f(2) = কত?}

ক) 3

খ) 2

গ) 1

\textbf{\hl{ঘ) 0 ✔}}

\textbf{\hl{✔ সঠিক উত্তর: ঘ) 0}}

\begin{quote}
\textbf{ব্যাখ্যা:} \emph{f(x) = x\^{}3 − 6x\^{}2 + 11x − 6}

\emph{⇒ f(2) = (2)\^{}3 − 6(2)\^{}2 + 11(2) − 6}

\emph{= 8 − 24 + 22 − 6}

\emph{= 30 − 30 = 0}
\end{quote}

\textbf{f(x) = x\^{}3 + Kx\^{}2 − 6x − 9; K এর মান কত হলে f(3) = 0 হবে?}

ক) 1

খ) -1

গ) 2

\textbf{\hl{ঘ) 0 ✔}}

\textbf{\hl{✔ সঠিক উত্তর: ঘ) 0}}

\begin{quote}
\textbf{ব্যাখ্যা:} \emph{f(3) = (3)\^{}3 + K(3)\^{}2 − 6(3) − 9}

\emph{= 27 + 9K − 18 − 9}

\emph{= 9K}

\emph{শর্তমতে, f(3) = 0}

\emph{⇒ 9K = 0}

\emph{∴ K = 0}
\end{quote}

\textbf{f(x) = x\^{}3 + ax\^{}3 + 2x\^{}3 = 0 হলে, a এর মান কত?}

ক) 3

খ) 1

গ) 0

\textbf{\hl{ঘ) -1 ✔}}

\textbf{\hl{✔ সঠিক উত্তর: ঘ) -1}}

\begin{quote}
\textbf{ব্যাখ্যা:} \emph{x\^{}3 + ax\^{}3 + 2x\^{}3 = 0}

\emph{⇒ 3x\^{}3 + ax\^{}3 = 0}

\emph{⇒ x\^{}3(3 + a) = 0}

\emph{যেহেতু x ≠ 0, তাই 3 + a = 0}

\emph{∴ a = −3}
\end{quote}

\textbf{f(x) = x\^{}2 + (1/x−1) − 1 হলে, নিচের কোনটি সঠিক?}

ক) f(0) = ∞

খ) f(1) = −1

গ) f(1) = 0

\textbf{\hl{ঘ) f(-1) = -(1/2) ✔}}

\textbf{\hl{✔ সঠিক উত্তর: ঘ) f(-1) = -(1/2)}}

\begin{quote}
\textbf{ব্যাখ্যা:} \emph{f(x) = x\^{}2 + (1/x−1) − 1}

\emph{⇒ f(−1) = (−1)\^{}2 + (1/−1−1) − 1}

\emph{= 1 + (1/−2) − 1}

\emph{= 1 − (1/2) − 1}

\emph{= -(1/2)}
\end{quote}

\textbf{f(x) = x\^{}2 + (1/x) + 1 অনুরূপ কোনটি?}

ক) f(1) = 1

খ) f(0) = 1

গ) f(-1) = 3

\textbf{\hl{ঘ) f(1) = 3 ✔}}

\textbf{\hl{✔ সঠিক উত্তর: ঘ) f(1) = 3}}

\begin{quote}
\textbf{ব্যাখ্যা:} \emph{f(x) = x\^{}2 + (1/x) + 1}

\emph{⇒ f(1) = (1)\^{}2 + (1/1) + 1 = 1 + 1 + 1 = 3}
\end{quote}

\textbf{∫x\^{}−1dx = কত?}

ক) α

খ) lnx

গ) x\^{}0

\textbf{\hl{ঘ) ln\textbar x\textbar{} ✔}}

\textbf{\hl{✔ সঠিক উত্তর: ঘ) ln\textbar x\textbar{}}}

\begin{quote}
\textbf{ব্যাখ্যা:} \emph{∫x\^{}−1dx}

\emph{= ∫(1/x)dx}

\emph{= ln\textbar x\textbar+c}
\end{quote}

\textbf{∫(1/x\^{}2)dx এর মান কত?}

\textbf{\hl{ক) −(1/x)+c ✔}}

খ) (1/x)+c

গ) x\^{}−2+c

ঘ) ln\textbar x\textbar+c

\textbf{\hl{✔ সঠিক উত্তর: ক) −(1/x)+c}}

\begin{quote}
\textbf{ব্যাখ্যা:} \emph{∫(1/x\^{}2)dx}

\emph{= ∫x\^{}−2dx}

\emph{= (x\^{}(−2+1)/−2+1)+c}

\emph{= −x\^{}−1+c}

\emph{= −(1/x)+c}
\end{quote}

\textbf{∫\_1\^{}elnxdx = কত?}

ক) e

খ) e − 1

গ) e + 1

\textbf{\hl{ঘ) 1 ✔}}

\textbf{\hl{✔ সঠিক উত্তর: ঘ) 1}}

\begin{quote}
\textbf{ব্যাখ্যা:} \emph{∫\_1\^{}elnxdx}

\emph{= {[}xlnx−x{]}\_1\^{}e}

\emph{= (elne−e)−(1ln1−1)}

\emph{= (e×1−e)−(0−1)}

\emph{= 0−(−1)=1}
\end{quote}

\textbf{যদি f(x) = x\^{}3 + kx\^{}2 − 4x\^{}2 − 8 হয়, তাহলে k এর কোন
মানের জন্য f(-2) = 0 হবে?}

\textbf{\hl{ক) 2 ✔}}

খ) -2

গ) (1/2)

ঘ) -(1/2)

\textbf{\hl{✔ সঠিক উত্তর: ক) 2}}

\begin{quote}
\textbf{ব্যাখ্যা:} \emph{যেহেতু f(-2) = 0}

\emph{∴ f(-2) = (-2)\^{}3 + k(-2)\^{}2 − 4(-2)\^{}2 − 8 = 0}

\emph{⇒ −8 + 4k − 4(4) − 8 = 0}

\emph{⇒ −8 + 4k − 16 − 8 = 0}

\emph{⇒ 4k = 8}

\emph{∴ k = 2}
\end{quote}

\section{37. বীজগাণিতিয় রাশিমালার প্রাথমিক চার
নিয়ম}\label{ux9acux99cux997ux9a3ux9a4ux9df-ux9b0ux9b6ux9aeux9b2ux9b0-ux9aaux9b0ux9a5ux9aeux995-ux99aux9b0-ux9a8ux9dfux9ae}

পাটিগণিতে আমরা সংখ্যা ও সংখ্যার বৈশিষ্ট্য জেনে বিভিন্ন গাণিতিক সমস্যা সমাধান
করেছি। জ্যামিতিতে বস্তুর আকৃতি সম্পর্কে জানবো ।

এ অধ্যায়ে আমরা গণিতের অন্যতম গুরুত্বপূর্ণ শাখা বীজগণিত সম্পর্কে জানবো । গণিতের এই
শাখার বৈশিষ্ট্য হলো অক্ষর প্রতীকের প্রয়োগ।

অক্ষর প্রতীক ব্যবহার করে আমরা নির্দিষ্ট কোনো সংখ্যার বদলে যেকোনো সংখ্যা বিবেচনা
করতে পারি।

দ্বিতীয়ত, অক্ষর অজানা পরিমাণের প্রতীক হিসেবে এবং সংখ্যার পরিবর্তে ব্যবহৃত হয়
বিধায় সকল গাণিতিক প্রক্রিয়া মেনে বীজগণিতীয় রাশি গঠন করা হয়।

গণিতের চারটি মৌলিক প্রক্রিয়া হলো যোগ, বিয়োগ, গুণ ও ভাগ। বিয়োগ হচ্ছে যোগের
বিপরীত প্রক্রিয়া আর ভাগ হচ্ছে গুণের বিপরীত প্রক্রিয়া। পাটিগণিতে কেবল ধনাত্মক
চিহ্নযুক্ত সংখ্যা ব্যবহার করা হয়। কিন্তু বীজগণিতে ধনাত্মক ও ঋণাত্মক উভয় চিহ্নযুক্ত
সংখ্যা এবং সংখ্যাসূচক প্রতীকও ব্যবহার করা হয়।

এ অধ্যায়ে বীজগণিতীয় প্রতীক, চলক, বীজগণিতীয় রাশি, বীজগণিতীয় রাশির প্রাথমিক
চার নিয়ম যথা: যোগ , বিয়োগ , গুণ ও ভাগ অত্যন্ত সুন্দরভাবে উপস্থাপন করা হয়েছে ।

\begin{itemize}
\item
  বীজগণিতীয় প্রতীক
\item
  চলক
\item
  বীজগণিতীয় রাশি ও পদ
\item
  সদৃশ ও বিসদৃশ পদ
\item
  বীজগণিতীয় রাশির যোগ
\item
  বীজগণিতীয় রাশির বিয়োগ
\item
  বীজগণিতীয় রাশির গুণ
\item
  বীজগণিতীয় রাশির ভাগ
\end{itemize}

\subsubsection{বীজগণিতীয়
প্রতীক}\label{ux9acux99cux997ux9a3ux9a4ux9af-ux9aaux9b0ux9a4ux995}

পাটিগণিতে সংখ্যা প্রতীক বা অঙ্কগুলো ১,২,৩,৪,৫,৬,৭,৮,৯,০। বীজগণিতে ব্যবহৃত
সংখ্যা প্রতীক বা অঙ্কগুলো 1, 2, 3, 4, 5, 6, 7, 8, 9, 0 । এ সব সংখ্যা প্রতীক
দ্বারা যেকোনো সংখ্যা লেখা যায়। তবে, বীজগণিতে সংখ্যা প্রতীকের সাথে অক্ষর
প্রতীকও ব্যবহার করা হয় । এটি বীজগণিতের মৌলিক বৈশিষ্ট্য। বীজগণিতে a, b, c, P,
9, 1, x,y,z,......... ইত্যাদি অক্ষর দ্বারা জানা বা অজানা সংখ্যা বা রাশিকে
প্রকাশ করা হয় ।

মনে করি, মলির কাছে কয়েকটি আম আছে। এখানে মলির কাছে কয়টি আম আছে তা নির্দিষ্ট
করে বলা হয়নি । তার কাছে যেকোনো সংখ্যক আম থাকতে পারে। তবে বীজগণিতীয়
প্রতীকের সাহায্যে বলা যায়, তার কাছে x সংখ্যক আম আছে। x এর মান 5 হলে, মলির
কাছে 5টি আম আছে ; x এর মান 10 হলে, মলির কাছে 10টি আম আছে, ইত্যাদি ।

\subsubsection{চলক}\label{ux99aux9b2ux995}

অক্ষর প্রতীক x এর মান 5 বা 10 বা অন্য কোনো সংখ্যা হতে পারে। বীজগণিতে

ধরনের অজ্ঞাত রাশি বা অক্ষর প্রতীককে চলক বলে । অতএব, x চলকের একটি উদাহরণ।

এখানে চলক হিসেবে x প্রতীক ব্যবহার করা হয়েছে। x প্রতীকের পরিবর্তে y প্রতীক নয়
কেন ? চলক হিসেবে x এর পরিবর্তে বা অন্য কোনো প্রতীকও ব্যবহার করা যায় ।

লক্ষ করুন:

* চলক এমন একটি প্রতীক যার মানের পরিবর্তন হয় ।

* চলকের মান নির্দিষ্ট নয়।

* চলক বিভিন্ন মান ধারণ করতে পারে ।

\subsubsection{বীজগণিতীয় রাশি ও
পদ}\label{ux9acux99cux997ux9a3ux9a4ux9af-ux9b0ux9b6-ux993-ux9aaux9a6}

5x, 2x + 3y, 5x + 3y - z, 3b×c - y, 5x + 2y + 9x - y ইত্যাদি এক একটি
বীজগণিতীয় রাশি । প্রক্রিয়া চিহ্ন ও সংখ্যাসূচক প্রতীক এর অর্থবোধক সংযোগ বা
বিন্যাসকে বীজগণিতীয় রাশি বলা হয় । বীজগণিতীয় রাশির যে অংশ যোগ (+) ও বিয়োগ
(-) চিহ্ন দ্বারা সংযুক্ত থাকে, এদের প্রত্যেকটিকে ঐ রাশির পদ বলা হয় । যেমন, 4x
+ 3y একটি রাশি । রাশিটিতে 4x ও 3y দুইটি পদ রয়েছে । এরা যোগ চিহ্ন দ্বারা
যুক্ত। আবার, 5x + 3y ÷ c+ 4b × 2y রাশিতে 5x, 3y÷c, 4b × 2y তিনটি পদ আছে।
4x একটি একপদী, 2x + 3y একটি দ্বিপদী, a - 2b + 4c একটি ত্রিপদী রাশি ।

\subsubsection{সদৃশ ও বিসদৃশ
পদ}\label{ux9b8ux9a6ux9b6-ux993-ux9acux9b8ux9a6ux9b6-ux9aaux9a6}

7a2bx, 8a2bx দুইটি বীজগণিতীয় রাশি । রাশি দুইটির পদগুলোর মধ্যে পার্থক্য হচ্ছে
শুধুমাত্র সাংখ্যিক সহগে । এই পদ দুইটি সদৃশ পদ ।

এক বা একাধিক বীজগণিতীয় রাশির অন্তর্ভুক্ত যেসব পদের একমাত্র পার্থক্য রয়েছে
সাংখ্যিক সহগে, তাদের সদৃশ পদ বলা হয় । অন্যথায় পদগুলো বিসদৃশ। যেমন, 9ax, 9ay
রাশি দুইটির সাংখ্যিক সহগ একই, কিন্তু পদ দুইটি পৃথক ; তাই তারা বিসদৃশ ।

\subsubsection{বীজগণিতীয় রাশির
যোগ}\label{ux9acux99cux997ux9a3ux9a4ux9af-ux9b0ux9b6ux9b0-ux9afux997}

দুই বা ততোধিক বীজগণিতীয় রাশি যোগ করতে হলে সদৃশ পদের সহগগুলো চিহ্নযুক্ত সংখ্যার
নিয়মে যোগ করতে হবে । এরপর প্রাপ্ত সহগের ডানপাশে প্রতীকগুলো বসাতে হবে। বিসদৃশ
পদগুলো

তাদের চিহ্নসহ যোগফলে বসাতে হবে ।

উদাহরণ ১: যোগ কর :

2a+4b+5c, 3a+2b-6c.

সদৃশ পদগুলো তাদের স্ব-স্ব চিহ্নসহ নিচে নিচে লিখে পাই,

2a+4b+5c

3a+2b-6c

-\/-\/-\/-\/-\/-\/-\/-\/-\/-\/-\/-\/-\/-\/-\/-\/-

5a+6b-c

নির্ণেয় যোগফল 5a +6b-c.

লক্ষ করুন : সদৃশ পদের সাংখ্যিক সহগগুলোর বীজগণিতীয় যোগফল নির্ণয় করা হয়েছে ।
প্রাপ্ত যোগফলের পাশে সংশ্লিষ্ট পদের প্রতীকগুলো বসানো হয়েছে। এভাবে প্রাপ্ত সব
পদের যোগফলই নির্ণেয় যোগফল ।

\subsubsection{বীজগণিতীয় রাশির
বিয়োগ}\label{ux9acux99cux997ux9a3ux9a4ux9af-ux9b0ux9b6ux9b0-ux9acux9afux997}

a - b = a + (- b)

একটি বীজগণিতীয় রাশি থেকে অপর একটি বীজগণিতীয় রাশি বিয়োগ করার ক্ষেত্রে, প্রথম
রাশির সাথে দ্বিতীয় রাশির যোগাত্মক বিপরীত রাশি যোগ করা হয় । অর্থাৎ, বিয়োজ্য
বা দ্বিতীয় রাশির প্রতিটি পদের চিহ্ন পরিবর্তন করে প্রাপ্ত রাশিকে প্রথম রাশির
সাথে যোগ করা।

উদাহরণ ২: 5a+4b5c থেকে 3a - 4b - 60 বিয়োগ কর।

\textbf{সমাধান :} বিয়োজ্যের প্রতিটি পদের চিহ্ন পরিবর্তন করে পাই

-3a+4b+6c

এখন প্রথম রাশির সাথে রূপান্তরিত বিয়োজ্য রাশি যোগ করে পাই,

+5a+4b - 5c

-3a+4b+6c

-\/-\/-\/-\/-\/-\/-\/-\/-\/-\/-\/-\/-\/-\/-\/-\/-

2a + 8b + c

নির্ণেয় বিয়োগফল 2a + 8b + c

\subsubsection{বীজগণিতীয় রাশির
গুণ}\label{ux9acux99cux997ux9a3ux9a4ux9af-ux9b0ux9b6ux9b0-ux997ux9a3}

গুণ করার সময় একই চিহ্নযুক্ত দুইটি রাশির গুণফল (+) চিহ্নযুক্ত হবে এবং বিপরীত
চিহ্নযুক্ত দুইটি রাশির গুণফল (-) চিহ্নযুক্ত হবে।

\paragraph{একপদী রাশিকে একপদী রাশি দ্বারা
গুণ}\label{ux98fux995ux9aaux9a6-ux9b0ux9b6ux995-ux98fux995ux9aaux9a6-ux9b0ux9b6-ux9a6ux9acux9b0-ux997ux9a3}

দুইটি একপদী রাশির গুণের ক্ষেত্রে তাদের সাংখ্যিক সহগদ্বয়কে আগে গুণ করে নিতে হয়।
তারপর

উভয়পদে বিদ্যমান বীজগণিতীয় প্রতীকগুলোকে সূচক নিয়মে গুণ করে ( অর্থাৎ পাওয়ারে
পাওয়ারে যোগ করে ) গুণফলে লিখতে হয়। অন্যান্য প্রতীকগুলো অপরিবর্তিত অবস্থায় গুণফলে
নেওয়া হয়।

উদাহরণ ৩: 12\(a^{2}\)x\(y^{2}\) কে -- 6a\(x^{3}\)b দ্বারা গুণ কর।

\textbf{সমাধান :} এখানে 12 ও -6 এর গুণফল -72 ।দুটি বিপরীত চিহ্নযুক্ত হওয়াতে
72 এর আগে ( - ) বসাতে হয়েছে।

এখন \(a^{2}\) কে a দ্বারা গুণ করলে হয় \(a^{3}\) । কারণ পাওয়ারে পাওয়ারে যোগ
হয়। তেমনিভাবে হবে \(x^{4}\) , \(y^{2}\) এবং b । সুতরাং গুণফল =
-72\(a^{3}\)b\(x^{4}\)\(y^{2}\) ।

\paragraph{বহুপদী রাশিকে একপদী রাশি দ্বারা
গুণ}\label{ux9acux9b9ux9aaux9a6-ux9b0ux9b6ux995-ux98fux995ux9aaux9a6-ux9b0ux9b6-ux9a6ux9acux9b0-ux997ux9a3}

একের অধিক পদযুক্ত বীজগণিতীয় রাশিই বহুপদী রাশি। যেমন, 5\(x^{2}\)y +
7x\(y^{2}\) একটি বহুপদী রাশি । বহুপদী রাশিকে একপদী রাশি দ্বারা গুণ করতে হলে
বহুপদী রাশির প্রত্যেক পদকে একপদী রাশি দ্বারা আলাদা আলাদাভাবে গুণ করতে হয়।

উদাহরণ ৪: (5\(x^{2}\)y + 7x\(y^{2}\) ) x 5\(x^{3}\)\(y^{3}\) ।

\textbf{সমাধান :} একপদী রাশি 5\(x^{3}\)\(y^{3}\) দ্বারা বহুপদী রাশির প্রথম
পদ 5\(x^{2}\)y কে গুণ করলে গুণফল হয় = 25\(x^{5}\)\(y^{4}\) । আবার একপদী
রাশি 5\(x^{3}\)\(y^{3}\) দ্বারা বহুপদী রাশির দ্বিতীয় পদ 7x\(y^{2}\) কে গুণ
করলে গুণফল হয় = 35\(x^{4}\)\(y^{5}\) ।

নির্ণেয় গুণফল = 25\(x^{5}\)\(y^{4}\) + 35\(x^{4}\)\(y^{5}\) ।

\paragraph{বহুপদী রাশিকে বহুপদী রাশি দ্বারা
গুণ}\label{ux9acux9b9ux9aaux9a6-ux9b0ux9b6ux995-ux9acux9b9ux9aaux9a6-ux9b0ux9b6-ux9a6ux9acux9b0-ux997ux9a3}

বহুপদী রাশিকে বহুপদী রাশি দ্বারা গুণ করতে হলে গুণ্যের প্রত্যেক পদকে গুণকের প্রত্যেক
পদ দ্বারা আলাদা আলাদাভাবে গুণ করে সদৃশ পদগুলোকে নিচে নিচে সাজিয়ে লিখতে হয়।
চিহ্নযুক্ত রাশির যোগের নিয়মে যোগ করতে হয়। বিসদৃশ পদ থাকলে সেগুলোকে পৃথকভাবে
লিখতে হয় এবং গুণফলে বসাতে হয়।

গুণের নিয়ম :

প্রথমে গুণ্যের প্রত্যেক পদকে গুণকের প্রথম পদ দ্বারা গুণ করে গুণফল লিখতে হবে।

এরপর গুণ্যের প্রত্যেক পদকে গুণকের দ্বিতীয় পদ দ্বারা গুণ করে গুণফল বের করতে হবে।

এ গুণফলকে এমনভাবে সাজিয়ে লিখতে হবে যেন উভয় গুণফলের সদৃশ পদগুলো নিচে নিচে
পড়ে।

প্রাপ্ত দুইটি গুণফলের বীজগণিতীয় সমষ্টিই হলো নির্ণেয় গুণফল।

উদাহরণ ৫: 3x + 2y কে x + y দ্বারা গুণ কর ।

\textbf{সমাধান:} প্রথমে x দ্বারা 3x ও 2y কে গুণ করতে হবে। তারপর y দ্বারা 3x
ও 2y কে গুণ করতে হবে ।অত:পর সদৃশপদগুলো যোগ করে গুণফল লিখতে হবে।

নির্ণেয় গুণফল = (3x + 2y)( x + y) = 3\(x^{2}\) + 2xy + 3xy + 2\(y^{2}\) =
3\(x^{2}\) + 5xy + 2\(y^{2}\)

\subsubsection{বীজগণিতীয় রাশির
ভাগ}\label{ux9acux99cux997ux9a3ux9a4ux9af-ux9b0ux9b6ux9b0-ux9adux997}

\paragraph{একপদী রাশিকে একপদী রাশি দ্বারা
ভাগ}\label{ux98fux995ux9aaux9a6-ux9b0ux9b6ux995-ux98fux995ux9aaux9a6-ux9b0ux9b6-ux9a6ux9acux9b0-ux9adux997}

একপদী রাশিকে একপদী রাশি দ্বারা ভাগ করতে হলে, সাংখ্যিক সহগকে পাটিগণিতীয়
নিয়মে ভাগ এবং

বীজগণিতীয় প্রতীককে সূচক নিয়মে ( অর্থাৎ পাওয়ারে পাওয়ারে বিয়োগ করে ) ভাগ করতে
হয়।

উদাহরণ ৬: 10\(a^{5}\)\(b^{7}\) কে 5\(a^{2}\)\(b^{3}\) দ্বারা ভাগ কর।

10\(a^{5}\)\(b^{7}\)÷5\(a^{2}\)\(b^{3}\)

= \({2a}^{(5 - 2)}b^{(7 - 3)}\)

= \({2a}^{3}b^{4}\)

\paragraph{বহুপদী রাশিকে একপদী রাশি দ্বারা
ভাগ}\label{ux9acux9b9ux9aaux9a6-ux9b0ux9b6ux995-ux98fux995ux9aaux9a6-ux9b0ux9b6-ux9a6ux9acux9b0-ux9adux997}

বহুপদী রাশিকে একপদী রাশি দ্বারা ভাগ করতে হলে বহুপদী রাশির প্রত্যেক পদকে একপদী
রাশি দ্বারা আলাদা আলাদাভাবে ভাগ করতে হয়।

উদাহরণ ৭: 10\(x^{5}y^{3}\) - 12\(x^{3}y^{8}\) + 6\(x^{4}y^{7}\) কে
2\(x^{2}y^{2}\) দ্বারা ভাগ কর।

( 10\(x^{5}y^{3}\) - 12\(x^{3}y^{8}\) + 6\(x^{4}y^{7}\) ) ÷
2\(x^{2}y^{2}\) = 5\(x^{3}y\) -- 6\(xy^{7}\) + 3\(x^{2}y^{5}\)

\paragraph{বহুপদী রাশিকে বহুপদী রাশি দ্বারা
ভাগ}\label{ux9acux9b9ux9aaux9a6-ux9b0ux9b6ux995-ux9acux9b9ux9aaux9a6-ux9b0ux9b6-ux9a6ux9acux9b0-ux9adux997}

বহুপদী রাশিকে বহুপদী রাশি দ্বারা ভাগ করার ক্ষেত্রে প্রথমে ভাজ্য ও ভাজক উভয়ের
মধ্যে আছে

এমন একটি বীজগণিতীয় প্রতীকের ঘাতের অধঃক্রম অনুসারে রাশিদ্বয়কে সাজাতে হবে।
যেমন

\(x^{2}\) + 2\(x^{4}\) + 110 - 48x একটি বহুপদী। একে x এর মানের অধঃক্রম
অনুসারে সাজালে আমরা

পাই : 2\(x^{4}\) + \(x^{2}\)-48x + 110। এরপর পাটিগণিতের ভাগ প্রক্রিয়ার মতো
নিচের নিয়মে ধাপে ধাপে ভাগ

করতে হবে।

• ভাজ্যের প্রথম পদটিকে ভাজকের প্রথম পদ দ্বারা ভাগ করলে যে ভাগফল হয় তা নির্ণেয়
ভাগফলের প্রথম পদ ।

• ভাগফলের ঐ প্রথম পদ দ্বারা ভাজকের প্রত্যেক পদকে গুণ করে গুণফল সদৃশ পদ অনুযায়ী

ভাজ্যের নিচে বসিয়ে ভাজ্য থেকে বিয়োগ করতে হয়।

• বিয়োগফল নতুন ভাজ্য হবে। বিয়োগফল এমনভাবে লিখতে হবে যেন তা আগের মতো বিবেচ্য
প্রতীকের অধঃক্রম অনুসারে থাকে ।

• নতুন ভাজ্যের প্রথম পদটিকে ভাজকের প্রথম পদ দ্বারা ভাগ করলে যে ভাগফল হয় তা
নির্ণেয়

ভাগফলের দ্বিতীয় পদ । এভাবে ক্রমান্বয়ে ভাগ করতে হয়।

উদাহরণ ৮: 6\(x^{2}\) + x - 2 কে 2x-1 দ্বারা ভাগ কর।

\textbf{সমাধান :} এখানে ভাজ্য ও ভাজক উভয়েই x এর ঘাতের অধঃক্রম অনুসারে
সাজানো আছে।

2x - 1 ) 6\(x^{2}\) + x - 2 ( 3x + 2

6\(x^{2}\) + x

-\/-\/-\/-\/-\/-\/-\/-\/-\/-\/-\/-\/-\/-\/-\/-\/-\/-\/-\/-\/-\/-\/-\/-\/-\/-\/-\/-\/-\/-\/-\/-\/-\/-\/-\/-\/-\/-

4x - 2

4x - 2

-\/-\/-\/-\/-\/-\/-\/-\/-\/-\/-\/-\/-\/-\/-\/-\/-\/-\/-\/-\/-\/-\/-\/-\/-\/-\/-\/-\/-\/-\/-\/-\/-\/-\/-\/-\/-\/-

0

নির্ণেয় ভাগফল 3x + 2.

এখানে, 6\(x^{2}\) ÷ 2x = 3x.

এই 3x দ্বারা ভাজক 2x-1 কে গুণ করে গুণফল ভাজ্যের সদৃশ পদের নিচে লিখে বিয়োগ করা
হল । নতুন ভাজ্য 4x-2 এর ক্ষেত্রে একই নিয়ম অনুসরণ করা হল।

\section{38. বীজগাণিতিয় রাশিমালার প্রাথমিক চার নিয়ম MCQ প্রশ্ন উত্তর
অনুশীলন}\label{ux9acux99cux997ux9a3ux9a4ux9df-ux9b0ux9b6ux9aeux9b2ux9b0-ux9aaux9b0ux9a5ux9aeux995-ux99aux9b0-ux9a8ux9dfux9ae-mcq-ux9aaux9b0ux9b6ux9a8-ux989ux9a4ux9a4ux9b0-ux985ux9a8ux9b6ux9b2ux9a8}

\textbf{\{(x+3x)\^{}2(x-3y)\^{}2\}\^{}2 কে সরল করলে কয়টি পদ পাওয়া যায়?}

ক) ৪টি

\textbf{\hl{খ) ৫টি ✔}}

গ) ৬টি

ঘ) ৭টি

\textbf{\hl{✔ সঠিক উত্তর: খ) ৫টি}}

\begin{quote}
\textbf{ব্যাখ্যা:} \emph{\{(x+3y)\^{}2(x−3y)\^{}2\}\^{}2}

\emph{= ({[}(\{(x+3y)(x−3y)\})\^{}2{]})\^{}2}

\emph{= {[}\{x\^{}2−9y\^{}2\}\^{}2{]}\^{}2}

\emph{= (\{x\^{}4−18x\^{}2y\^{}2+81y\^{}4\})\^{}2}

\emph{=
x\^{}8−36x\^{}6y\^{}2+486x\^{}4y\^{}4−2916x\^{}2y\^{}6+6561y\^{}8}

\emph{বীজগাণিতীয় রাশির যে অংশ যোগ (+) ও বিয়োগ (−) চিহ্ন দ্বারা সংযুক্ত থাকে
তাদের প্রত্যেকটিকে পদ বলা হয়।}

\emph{∴ প্রাপ্ত রাশিতে মোট পদসংখ্যা ৫টি।}
\end{quote}

\textbf{a = 2, b = 3 হলে 2a + 4b এর সাথে 2a\^{}2 + a − b যোগ করলে যোগফল
কত হবে?}

ক) 15

খ) 17

গ) 19

\textbf{\hl{ঘ) 23 ✔}}

\textbf{\hl{✔ সঠিক উত্তর: ঘ) 23}}

\begin{quote}
\textbf{ব্যাখ্যা:} \emph{প্রদত্ত রাশি দুটির যোগফল = (2a+4b)+(2a\^{}2+a−b)}

\emph{= 2a\^{}2+3a+3b}

\emph{এখন a=2 এবং b=3 বসিয়ে পাই,}

\emph{= 2·2\^{}2+3·2+3·3 = 8+6+9 = 23}
\end{quote}

\textbf{x\^{}2 − 3x − 2 কে x + 1 দ্বারা ভাগ করলে ভাগশেষ কত হবে?}

ক) 0

\textbf{\hl{খ) 2 ✔}}

গ) 6

ঘ) 4

\textbf{\hl{✔ সঠিক উত্তর: খ) 2}}

\begin{quote}
\textbf{ব্যাখ্যা:} \emph{ভাগশেষ উপপাদ্য (Remainder Theorem) অনুযায়ী,}

\emph{ধরি, f(x)=x\^{}2−3x−2}

\emph{f(x) কে x+1 দ্বারা ভাগ করলে ভাগশেষ হবে f(−1) ।}

\emph{∴ f(−1)=(−1)\^{}2−3(−1)−2}

\emph{= 1+3−2 = 2}
\end{quote}

\textbf{(x + 3) (x − 3) কে x\^{}2 − 6 দিয়ে ভাগ করলে ভাগশেষ কত হবে?}

ক) -6

খ) 3

গ) 6

\textbf{\hl{ঘ) -3 ✔}}

\textbf{\hl{✔ সঠিক উত্তর: ঘ) -3}}

\begin{quote}
\textbf{ব্যাখ্যা:} \emph{ভাজ্য = (x+3)(x−3)=x\^{}2−9}

\emph{ভাজক = x\^{}2−6}

\emph{আমরা জানি, ভাজ্য = (ভাজক × ভাগফল) + ভাগশেষ}

\emph{∴ x\^{}2−9=1·(x\^{}2−6)−3}

\emph{∴ নির্ণেয় ভাগশেষ = −3}
\end{quote}

\textbf{a−{[}a−\{a−(a−1)\}{]}=?}

ক) -1

\textbf{\hl{খ) 1 ✔}}

গ) 2a + 1

ঘ) 2a - 1

\textbf{\hl{✔ সঠিক উত্তর: খ) 1}}

\begin{quote}
\textbf{ব্যাখ্যা:} \emph{a−{[}a−\{a−(a−1)\}{]}}

\emph{= a−{[}a−\{a−a+1\}{]}}

\emph{= a−{[}a−1{]}}

\emph{= a−a+1 = 1}
\end{quote}

\textbf{a−\{a−(a+1)\}=?}

ক) a - 1

খ) 1

গ) a

\textbf{\hl{ঘ) a + 1 ✔}}

\textbf{\hl{✔ সঠিক উত্তর: ঘ) a + 1}}

\begin{quote}
\textbf{ব্যাখ্যা:} \emph{a−\{a−(a+1)\}}

\emph{= a−\{a−a−1\}}

\emph{= a−(−1)}

\emph{= a+1}
\end{quote}

\textbf{a=1,b=−1,c=2,d=−2 হলে a−(−b)−(−c)−(−d) এর মান কত?}

\textbf{\hl{ক) 0 ✔}}

খ) 1

গ) 2

ঘ) 3

\textbf{\hl{✔ সঠিক উত্তর: ক) 0}}

\begin{quote}
\textbf{ব্যাখ্যা:} \emph{প্রদত্ত রাশি = a−(−b)−(−c)−(−d)}

\emph{= a+b+c+d}

\emph{মান বসিয়ে পাই,}

\emph{= 1+(−1)+2+(−2)}

\emph{= 1−1+2−2 = 0}
\end{quote}

\textbf{If x=−3 , then the value of −3x\^{}2 is-}

\textbf{\hl{ক) -27 ✔}}

খ) -18

গ) 27

ঘ) 18

ঙ) 81

\textbf{\hl{✔ সঠিক উত্তর: ক) -27}}

\begin{quote}
\textbf{ব্যাখ্যা:} \emph{প্রদত্ত রাশি = −3x\^{}2}

\emph{x=−3 বসিয়ে পাই,}

\emph{= −3×(−3)\^{}2}

\emph{= −3×9}

\emph{= −27}
\end{quote}

\textbf{x\^{}2−8x−8y+16+y\^{}2 এর সঙ্গে কত যোগ করলে যোগফল একটি পূর্ণবর্গ
হবে?}

ক) -2xy

খ) 8xy

গ) 6xy

\textbf{\hl{ঘ) 2xy ✔}}

\textbf{\hl{✔ সঠিক উত্তর: ঘ) 2xy}}

\begin{quote}
\textbf{ব্যাখ্যা:} \emph{x\^{}2−8x−8y+16+y\^{}2}

\emph{= x\^{}2+(−y)\^{}2+4\^{}2+2·x·(−y)+2·(−y)·4+2·4·x−2xy}

\emph{= ((x−y+4))\^{}2−2xy}

\emph{∴ 2xy যোগ করলে রাশিটি একটি পূর্ণবর্গ হবে।}
\end{quote}

\textbf{If x=10 , which of the following has the minimum value? / x=10
হলে নিচের কোনটির মান সর্বনিম্ন?}

ক) 2 - x

খ) (x/2)

গ) (2/x)

\textbf{\hl{ঘ) (2 - x)(2 + x) ✔}}

\textbf{\hl{✔ সঠিক উত্তর: ঘ) (2 - x)(2 + x)}}

\begin{quote}
\textbf{ব্যাখ্যা:} \emph{ক. 2−x=2−10=−8}

\emph{খ. (x/2)=(10/2)=5}

\emph{গ. (2/x)=(2/10)=(1/5)=0.2}

\emph{ঘ. (2−x)(2+x)=(2−10)(2+10)=−8×12=−96}

\emph{গাণিতিক হিসেব অনুযায়ী (ঘ) অপশনের মান সর্বনিম্ন (−96)।}
\end{quote}

\textbf{a\^{}2+b\^{}2=?}

ক) (a - b)² - 2ab

খ) (a - b)² - 4ab

গ) (a + b)² + 2ab

\textbf{\hl{ঘ) (a + b)² - 2ab ✔}}

\textbf{\hl{✔ সঠিক উত্তর: ঘ) (a + b)² - 2ab}}

\begin{quote}
\textbf{ব্যাখ্যা:} \emph{আমরা জানি,}

\emph{a\^{}2+b\^{}2=((a+b))\^{}2−2ab}
\end{quote}

\textbf{2x+3y এর বর্গফল কত?}

ক) 2x² + 6xy + 6y²

\textbf{\hl{খ) 4x² + 12xy + 9y² ✔}}

গ) 4x² + 6xy + 6y²

ঘ) 2x² + 12xy + 6y²

\textbf{\hl{✔ সঠিক উত্তর: খ) 4x² + 12xy + 9y²}}

\begin{quote}
\textbf{ব্যাখ্যা:} \emph{2x+3y এর বর্গ}

\emph{= ((2x+3y))\^{}2}

\emph{= (2x)\^{}2+2·2x·3y+(3y)\^{}2}

\emph{= 4x\^{}2+12xy+9y\^{}2}
\end{quote}

\textbf{(2x + 1) এর বর্গ কত?}

\textbf{\hl{ক) 4x² + 4x + 1 ✔}}

খ) 2x² + 4x + 1

গ) 4x² + 2x + 1

ঘ) 2x² + 2x + 1

\textbf{\hl{✔ সঠিক উত্তর: ক) 4x² + 4x + 1}}

\begin{quote}
\textbf{ব্যাখ্যা:} \emph{(2x + 1) এর বর্গ}

\emph{= ((2x+1))\^{}2}

\emph{= (2x)\^{}2+2·2x·1+1\^{}2}

\emph{= 4x\^{}2+4x+1}
\end{quote}

\textbf{a = 15 এবং b = 6 হলে 9a\^{}2−48ab+64b\^{}2 -এর মান নির্ণয় করুন?}

ক) 8

খ) 5

গ) 6

\textbf{\hl{ঘ) 9 ✔}}

\textbf{\hl{✔ সঠিক উত্তর: ঘ) 9}}

\begin{quote}
\textbf{ব্যাখ্যা:} \emph{9a\^{}2−48ab+64b\^{}2}

\emph{= (3a)\^{}2−2·3a·8b+(8b)\^{}2}

\emph{= (3a−8b)\^{}2}

\emph{= (3×15−8×6)\^{}2}

\emph{= (45−48)\^{}2}

\emph{= (−3)\^{}2 = 9}
\end{quote}

\textbf{a\^{}2−2ab থেকে কত বিয়োগ করলে বিয়োগফল একটি পূর্ণবর্গ হবে?}

\textbf{\hl{ক) - b² ✔}}

খ) b

গ) a³

ঘ) -a

\textbf{\hl{✔ সঠিক উত্তর: ক) - b²}}

\begin{quote}
\textbf{ব্যাখ্যা:} \emph{বীজগাণিতিক রাশিমালাকে পূর্ণবর্গ হতে হলে
((a+b))\^{}2=a\^{}2+2ab+b\^{}2 অথবা (a−b)\^{}2=a\^{}2−2ab+b\^{}2 আকারে
হতে হয়।}

\emph{a\^{}2−2ab−(−b\^{}2)}

\emph{= a\^{}2−2ab+b\^{}2}

\emph{= (a−b)\^{}2}

\emph{∴ a\^{}2−2ab থেকে −b\^{}2 বিয়োগ করলে বিয়োগফল একটি পূর্ণবর্গ হবে।}
\end{quote}

\textbf{(4a\^{}2+9b\^{}2) রাশিটির সাথে কোনটি যোগ করলে যোগফল একটি পূর্ণবর্গ
হবে?}

ক) 6ab

\textbf{\hl{খ) 12ab ✔}}

গ) 18ab

ঘ) 24ab

\textbf{\hl{✔ সঠিক উত্তর: খ) 12ab}}

\begin{quote}
\textbf{ব্যাখ্যা:} \emph{বীজগাণিতিক রাশিমালাকে পূর্ণবর্গ হতে হলে
((a+b))\^{}2=a\^{}2+2ab+b\^{}2 অথবা (a−b)\^{}2=a\^{}2−2ab+b\^{}2 আকারে
হতে হয়।}

\emph{4a\^{}2+9b\^{}2+12ab}

\emph{= (2a)\^{}2+2·2a·3b+(3b)\^{}2}

\emph{= ((2a+3b))\^{}2}

\emph{∴ (4a\^{}2+9b\^{}2) রাশিটির সাথে 12ab যোগ করলে যোগফল একটি পূর্ণবর্গ
হবে।}
\end{quote}

\textbf{9a\^{}2+16b\^{}2 রাশিটির সাথে কোনটি যোগ করলে যোগফল পূর্ণবর্গ হবে?}

ক) 12ab

\textbf{\hl{খ) 24ab ✔}}

গ) 36ab

ঘ) 144ab

\textbf{\hl{✔ সঠিক উত্তর: খ) 24ab}}

\begin{quote}
\textbf{ব্যাখ্যা:} \emph{বীজগাণিতিক রাশিমালাকে পূর্ণবর্গ হতে হলে
((a+b))\^{}2=a\^{}2+2ab+b\^{}2 অথবা (a−b)\^{}2=a\^{}2−2ab+b\^{}2 আকারে
হতে হয়।}

\emph{9a\^{}2+16b\^{}2+24ab}

\emph{= (3a)\^{}2+2·3a·4b+(4b)\^{}2}

\emph{= ((3a+4b))\^{}2}

\emph{∴ (9a\^{}2+16b\^{}2) রাশিটির সাথে 24ab যোগ করলে যোগফল একটি পূর্ণবর্গ
হবে।}
\end{quote}

\textbf{(x/y) এর সাথে কত যোগ করলে যোগফল (2y/x) হবে?}

ক) (2x\^{}2+2y\^{}2/xy)

খ) (x\^{}2−y\^{}2/xy)

\textbf{\hl{গ) (2y\^{}2−x\^{}2/xy) ✔}}

ঘ) (x\^{}2−2y\^{}2/xy)

\textbf{\hl{✔ সঠিক উত্তর: গ) (2y\^{}2−x\^{}2/xy)}}

\begin{quote}
\textbf{ব্যাখ্যা:} \emph{যোগ করতে হবে}

\emph{= (2y/x)−(x/y)}

\emph{= (2y\^{}2−x\^{}2/xy)}
\end{quote}

\textbf{(x + 5)(x - 3) = কত?}

ক) x² + 8x + 15

খ) x² - 15

\textbf{\hl{গ) x² + 2x - 15 ✔}}

ঘ) x² + 2x + 15

\textbf{\hl{✔ সঠিক উত্তর: গ) x² + 2x - 15}}

\begin{quote}
\textbf{ব্যাখ্যা:} \emph{(x + 5)(x - 3)}

\emph{= x² + 5x - 3x - 15}

\emph{= x² + 2x - 15}
\end{quote}

\textbf{a = 8, b = 6, x = (1/2) এবং y = 4 হলে, ax + 2b - 2xy এর মান কত?}

\textbf{\hl{ক) 12 ✔}}

খ) 9

গ) 7

ঘ) 6

\textbf{\hl{✔ সঠিক উত্তর: ক) 12}}

\begin{quote}
\textbf{ব্যাখ্যা:} \emph{ax + 2b - 2xy}

\emph{= 8×(1/2)+2×6−2×(1/2)×4 = 4 + 12 - 4 = 12}
\end{quote}

\textbf{a = 15 এবং b = 5 হলে ((a−b)\^{}2/a−b) = কত?}

ক) 30

\textbf{\hl{খ) 10 ✔}}

গ) 15

ঘ) 20

\textbf{\hl{✔ সঠিক উত্তর: খ) 10}}

\begin{quote}
\textbf{ব্যাখ্যা:} \emph{((a−b)\^{}2/a−b)}

\emph{= ((a−b)(a−b)/a−b)}

\emph{= a − b}

\emph{= 15 − 5 = 10}
\end{quote}

\textbf{(a+b) = 2 , (a-b) = 0 হলে (1/b) = কত?}

ক) 0

\textbf{\hl{খ) 1 ✔}}

গ) 2

ঘ) 3

\textbf{\hl{✔ সঠিক উত্তর: খ) 1}}

\begin{quote}
\textbf{ব্যাখ্যা:} \emph{a + b = 2}

\emph{a − b = 0}

\emph{-\/-\/-\/-\/-\/-\/-\/-\/-\/-\/-\/-\/-\/-\/-\/-\/-\/-\/-\/-\/-\/-\/-\/-}

\emph{(যোগ করে) 2a = 2}

\emph{বা, a = 1}

\emph{∴ 1 + b = 2}

\emph{বা, b = 2 − 1 = 1}

\emph{(1/b) = (1/1) = 1}
\end{quote}

\textbf{a = 3 হলে a³ = কত?}

ক) 3

খ) 6

গ) 9

\textbf{\hl{ঘ) 27 ✔}}

\textbf{\hl{✔ সঠিক উত্তর: ঘ) 27}}

\begin{quote}
\textbf{ব্যাখ্যা:} \emph{a³ = a × a × a = 3 × 3 × 3 = 27}
\end{quote}

\textbf{(a−b+b(a−b)/a−b) = কত?}

\textbf{\hl{ক) 1+b ✔}}

খ) (1+b/a−b)

গ) a+b

ঘ) (a-b)

\textbf{\hl{✔ সঠিক উত্তর: ক) 1+b}}

\begin{quote}
\textbf{ব্যাখ্যা:} \emph{(a−b+b(a−b)/a−b)}

\emph{= (1(a−b)+b(a−b)/a−b)}

\emph{= ((a−b)(1+b)/a−b)}

\emph{= 1 + b}
\end{quote}

\textbf{9x² + 30x এর সঙ্গে কত যোগ করলে যোগফল একটি পূর্ণবর্গ হবে?}

ক) 10

খ) 15

গ) 20

\textbf{\hl{ঘ) 25 ✔}}

\textbf{\hl{✔ সঠিক উত্তর: ঘ) 25}}

\begin{quote}
\textbf{ব্যাখ্যা:} \emph{9x² + 30x}

\emph{= (3x)² + 2 · 3x · 5 + 5² − 25}

\emph{= (3x + 5)² − 25}

\emph{∴ 9x² + 30x এর সঙ্গে 25 যোগ করলে যোগফল একটি পূর্ণবর্গ হবে।}
\end{quote}

\textbf{x = 2, y = 3 হলে 2x + 4y এর সাথে 2x² + x - y যোগ করলে যোগফল কত
হবে?}

ক) 5

খ) -7

গ) -8

\textbf{\hl{ঘ) 23 ✔}}

\textbf{\hl{✔ সঠিক উত্তর: ঘ) 23}}

\begin{quote}
\textbf{ব্যাখ্যা:} \emph{2x² + x − y}

\emph{2x + 4y}

\emph{-\/-\/-\/-\/-\/-\/-\/-\/-\/-\/-\/-\/-\/-\/-\/-\/-\/-\/-\/-\/-\/-\/-\/-\/-\/-\/-\/-\/-}

\emph{(+) 2x² + 3x + 3y}

\emph{= 2 · (2)² + 3 × 2 + 3 × 3 = 2 × 4 + 6 + 9 = 8 + 15 = 23}
\end{quote}

\textbf{(x/0) এর মান;}

ক) শূন্য

খ) এক

\textbf{\hl{গ) অসীম ✔}}

ঘ) কোনোটিই নয়

\textbf{\hl{✔ সঠিক উত্তর: গ) অসীম}}

\begin{quote}
\textbf{ব্যাখ্যা:} \emph{(x/0)=∞ অর্থাৎ অসীম।}
\end{quote}

\textbf{x\^{}−3−0.001=0 হলে, x\^{}2 এর মান -}

\textbf{\hl{ক) 100 ✔}}

খ) (1/10)

গ) 10

ঘ) (1/100)

\textbf{\hl{✔ সঠিক উত্তর: ক) 100}}

\begin{quote}
\textbf{ব্যাখ্যা:} \emph{x\^{}−3−0.001=0}

\emph{বা, x\^{}−3=0.001}

\emph{বা, (1/x\^{}3)=(1/1000)}

\emph{বা, x\^{}3=1000}

\emph{বা, x=10}

\emph{∴ x\^{}2=100}
\end{quote}

\textbf{(-1) × (-1) × (-1) + (-1) (-1) = কত ?}

\textbf{\hl{ক) 0 ✔}}

খ) 2

গ) 1

ঘ) -2

\textbf{\hl{✔ সঠিক উত্তর: ক) 0}}

\begin{quote}
\textbf{ব্যাখ্যা:} \emph{(-1) × (-1) × (-1) + (-1) (-1) = -1 + 1 = 0}
\end{quote}

\textbf{√(x\^{}2)=?}

ক) x

খ) -x

\textbf{\hl{গ) ±x ✔}}

ঘ) x²

\textbf{\hl{✔ সঠিক উত্তর: গ) ±x}}

\begin{quote}
\textbf{ব্যাখ্যা:} \emph{√(x\^{}2)=±x}
\end{quote}

\textbf{3\^{}x\^{}2=9\^{}4 হলে x এর মান কত?}

ক) ±2

খ) ± √(6)

\textbf{\hl{গ) ±2 √(2) ✔}}

ঘ) ±3

\textbf{\hl{✔ সঠিক উত্তর: গ) ±2 √(2)}}

\begin{quote}
\textbf{ব্যাখ্যা:} \emph{3\^{}x\^{}2=9\^{}4}

\emph{⇒ (3)\^{}x\^{}2=(3\^{}2)\^{}4}

\emph{⇒ (3)\^{}x\^{}2=(3)\^{}8 ⇒ x\^{}2=8 ⇒ x=±√(4×2)}

\emph{∴ x=±2√(2)}
\end{quote}

\textbf{x\^{}3=27 হলে x এর মান কত?}

ক) ২

\textbf{\hl{খ) ৩ ✔}}

গ) ৪

ঘ) ৫

\textbf{\hl{✔ সঠিক উত্তর: খ) ৩}}

\begin{quote}
\textbf{ব্যাখ্যা:} \emph{x\^{}3=27 ⇒ x\^{}3=3\^{}3}

\emph{∴ x = 3}
\end{quote}

\textbf{(x\^{}2)\^{}3 কে x\^{}3 দ্বারা গুণ করলে কত হবে?}

ক) x\^{}18

খ) x\^{}27

\textbf{\hl{গ) x\^{}9 ✔}}

ঘ) x\^{}2

\textbf{\hl{✔ সঠিক উত্তর: গ) x\^{}9}}

\begin{quote}
\textbf{ব্যাখ্যা:}
\emph{(x\^{}2)\^{}3×x\^{}3=x\^{}6×x\^{}3=x\^{}(6+3)=x\^{}9}
\end{quote}

\textbf{abc = 120 হলে, a এর মান কোনটি হতে পারে না?}

\textbf{\hl{ক) 0 ✔}}

খ) 1

গ) 2

ঘ) 12

\textbf{\hl{✔ সঠিক উত্তর: ক) 0}}

\begin{quote}
\textbf{ব্যাখ্যা:} \emph{abc = 120}

\emph{⇒ bc=(120/a)}

\emph{এখানে, a = 0 হতে পারে না। কারণ 0 দ্বারা কোনো সংখ্যাকে ভাগ করা যায়
না।}
\end{quote}

\textbf{x - (x - (x+1)) এর মান কত?}

ক) 1

খ) - 1

গ) x-1

\textbf{\hl{ঘ) x + 1 ✔}}

\textbf{\hl{✔ সঠিক উত্তর: ঘ) x + 1}}

\begin{quote}
\textbf{ব্যাখ্যা:} \emph{x − (x − (x + 1))}

\emph{= x − (x − x − 1)}

\emph{= x − (−1)}

\emph{= x + 1}
\end{quote}

\textbf{a−(6/a)=1 হলে, (6/a\^{}2−a+1) এর মান কত?}

\textbf{\hl{ক) (6/7) ✔}}

খ) (7/6)

গ) (1/7)

ঘ) (6/5)

\textbf{\hl{✔ সঠিক উত্তর: ক) (6/7)}}

\begin{quote}
\textbf{ব্যাখ্যা:} \emph{দেয়া আছে,}

\emph{a−(6/a)=1}

\emph{⇒ (a\^{}2−6/a)=1 ⇒ a\^{}2−6=a ⇒ a\^{}2−a=6}

\emph{এখন, (6/a\^{}2−a+1)=(6/6+1)=(6/7)}
\end{quote}

\textbf{((1/a)+1)÷(1−(1/a\^{}2))=কত?}

ক) (a/a+1)

\textbf{\hl{খ) (a/a−1) ✔}}

গ) a(a+1)

ঘ) a(a-1)

\textbf{\hl{✔ সঠিক উত্তর: খ) (a/a−1)}}

\begin{quote}
\textbf{ব্যাখ্যা:} \emph{((1/a)+1)÷(1−(1/a\^{}2))}

\emph{= ((1+a/a))÷((a\^{}2−1/a\^{}2))}

\emph{= (a+1/a)×(a\^{}2/(a+1)(a−1))}

\emph{= (a/a−1)}
\end{quote}

\textbf{x\^{}2−2x+1=0 হলে, (x\^{}4+2x\^{}2+1/x\^{}2) এর মান-}

ক) -4

খ) 0

\textbf{\hl{গ) 4 ✔}}

ঘ) 8

\textbf{\hl{✔ সঠিক উত্তর: গ) 4}}

\begin{quote}
\textbf{ব্যাখ্যা:} \emph{x\^{}2−2x+1=0}

\emph{⇒ (x−1)\^{}2=0}

\emph{⇒ x−1=0}

\emph{∴ x=1}

\emph{এখন,
(x\^{}4+2x\^{}2+1/x\^{}2)=(1\^{}4+2·1\^{}2+1/1\^{}2)=(1+2+1/1)=4}
\end{quote}

\textbf{(1/x)=1−x হলে, (1/x)+x\^{}2 এর মান কত হবে?}

ক) 1

খ) -1

গ) ±1

\textbf{\hl{ঘ) 0 ✔}}

\textbf{\hl{✔ সঠিক উত্তর: ঘ) 0}}

\begin{quote}
\textbf{ব্যাখ্যা:} \emph{দেয়া আছে, (1/x)=1−x}

\emph{বা, (1/x)+x=1}

\emph{বা, 1+x\^{}2=x {[}উভয় পক্ষকে x দ্বারা গুণ করে{]}}

\emph{∴ x\^{}2=x−1}

\emph{এখন, (1/x)+x\^{}2=1−x+x−1=0}
\end{quote}

\textbf{x=√(2+1) হলে x\^{}2+(1/x\^{}2) এর মান---}

ক) 0

খ) 2

গ) -2

\textbf{\hl{ঘ) 4 ✔}}

\textbf{\hl{✔ সঠিক উত্তর: ঘ) 4}}

\begin{quote}
\textbf{ব্যাখ্যা:} \emph{x=√(2+1)=√(3)}

\emph{x\^{}2=3}

\emph{x\^{}2+(1/x\^{}2)=3+(1/3)=(10/3)≡3.33 (প্রশ্ন বা অপশনে সামান্য ত্রুটি
রয়েছে, যা ৩.৩৩ এর কাছাকাছি সঠিক মান নির্দেশ করে।)}
\end{quote}

\textbf{a\^{}3+b\^{}3=? যখন a এর মান 4 এবং b এর মান 2 হয়।}

\textbf{\hl{ক) 72 ✔}}

খ) 64

গ) 56

ঘ) 48

\textbf{\hl{✔ সঠিক উত্তর: ক) 72}}

\begin{quote}
\textbf{ব্যাখ্যা:} \emph{দেয়া আছে, a = 4 এবং b = 2}

\emph{∴ a\^{}3+b\^{}3=(4)\^{}3+(2)\^{}3=64+8=72}
\end{quote}

\textbf{((√(7)+√(7)))\^{}2=?}

ক) 98

খ) 49

\textbf{\hl{গ) 28 ✔}}

ঘ) 21

ঙ) 14

\textbf{\hl{✔ সঠিক উত্তর: গ) 28}}

\begin{quote}
\textbf{ব্যাখ্যা:} \emph{((√(7)+√(7)))\^{}2=(2√(7))\^{}2=4×7=28}
\end{quote}

\textbf{√(1)+√(1) এর বর্গ কত?}

ক) √(2)

\textbf{\hl{খ) 4 ✔}}

গ) √(3)

ঘ) 2 + √(1)

\textbf{\hl{✔ সঠিক উত্তর: খ) 4}}

\begin{quote}
\textbf{ব্যাখ্যা:} \emph{((√(1)+√(1)))\^{}2=(2√(1))\^{}2=4(1)=4}
\end{quote}

\textbf{a, b and c are numbers of different values that}

a +b = d

a × b × c = 0

\textbf{Which one of these numbers is zero?}

ক) a

খ) b

\textbf{\hl{গ) c ✔}}

ঘ) d

\textbf{\hl{✔ সঠিক উত্তর: গ) c}}

\begin{quote}
\textbf{ব্যাখ্যা:} \emph{Since a, b, c, and d are of different values, if
either a or b were 0, then a + b = d would make d equal to the other
variable (e.g., if a = 0, then d = b), which violates the condition that
all values are different.}

\emph{Therefore, to satisfy a×b×c=0 , c must be zero.}
\end{quote}

\textbf{a − {[}a − \{a − (a − a−1― )\}{]} = কত??}

ক) 1

খ) -1

\textbf{\hl{গ) a-1 ✔}}

ঘ) a+1

\textbf{\hl{✔ সঠিক উত্তর: গ) a-1}}

\begin{quote}
\textbf{ব্যাখ্যা:} \emph{a − {[}a − \{a − (a − a−1― )\}{]}}

\emph{= a − {[}a − \{a − (a − a + 1)\}{]}}

\emph{= a − {[}a − \{a − 1\}{]}}

\emph{= a − {[}a − a + 1{]}}

\emph{= a − 1}
\end{quote}

\textbf{1-(a-(a-(a-(a-1)))) = ?}

ক) 4a-1

খ) 1

গ) a

\textbf{\hl{ঘ) 0 ✔}}

\textbf{\hl{✔ সঠিক উত্তর: ঘ) 0}}

\begin{quote}
\textbf{ব্যাখ্যা:} \emph{1 − (a − (a − (a − (a − 1))))}

\emph{= 1 − (a − (a − (a − a + 1)))}

\emph{= 1 − (a − (a − 1))}

\emph{= 1 − (a − a + 1)}

\emph{= 1 − 1 = 0}
\end{quote}

\textbf{(3/5)×(7/9)=কত?}

ক) (10/14)

খ) (27/35)

গ) (35/27)

\textbf{\hl{ঘ) (7/15) ✔}}

\textbf{\hl{✔ সঠিক উত্তর: ঘ) (7/15)}}

\begin{quote}
\textbf{ব্যাখ্যা:} \emph{(3/5)×(7/9)=(3×7/5×9)=(7/5×3)=(7/15)}
\end{quote}

\textbf{যদি \textquotesingle+\textquotesingle{} অর্থ বিয়োগ,
\textquotesingle-\textquotesingle{} অর্থ গুণ,
\textquotesingle×\textquotesingle{} অর্থ ভাগ এবং
\textquotesingle÷\textquotesingle{} অর্থ যোগ হয়, তবে ৫-৫+৫÷৫×৫ = ?}

ক) ৪১

খ) ৫০

গ) ১৬

\textbf{\hl{ঘ) ২১ ✔}}

\textbf{\hl{✔ সঠিক উত্তর: ঘ) ২১}}

\begin{quote}
\textbf{ব্যাখ্যা:} \emph{প্রদত্ত চিহ্নের পরিবর্তন অনুযায়ী:}

\emph{৫ − ৫ + ৫ ÷ ৫ × ৫}

\emph{= ৫ × ৫ − ৫ + ৫ ÷ ৫}

\emph{= ৫ × ৫ − ৫ + ১}

\emph{= ২৫ − ৫ + ১}

\emph{= ২৬ − ৫ = ২১}
\end{quote}

\textbf{4 × 7 − 8 ÷ 4 = 26}

ক) +, -, ×

\textbf{\hl{খ) ×, -, ÷ ✔}}

গ) ×, +-

ঘ) +, ÷, ×

\textbf{\hl{✔ সঠিক উত্তর: খ) ×, -, ÷}}

\begin{quote}
\textbf{ব্যাখ্যা:} \emph{(4 × 7) − (8 ÷ 4)}

\emph{= 28 − 2 = 26}

\emph{∴ সঠিক চিহ্নসমূহ হলো (×, −, ÷)।}
\end{quote}

\textbf{(-1) × (-1) × (-1) + (1) × (-1) = কত?}

ক) 0

খ) 2

\textbf{\hl{গ) -2 ✔}}

ঘ) 1

\textbf{\hl{✔ সঠিক উত্তর: গ) -2}}

\begin{quote}
\textbf{ব্যাখ্যা:} \emph{(-1) × (-1) × (-1) + (1) × (-1) = (-1) + (-1)}

\emph{= -1 − 1 = -2}
\end{quote}

\textbf{- 2 + (- 2)− \{- (2)\}− 2 = ?}

\textbf{\hl{ক) - 4 ✔}}

খ) 4

গ) 0

ঘ) 1

\textbf{\hl{✔ সঠিক উত্তর: ক) - 4}}

\begin{quote}
\textbf{ব্যাখ্যা:} \emph{− 2 + (− 2) − \{− (2)\} − 2}

\emph{= − 2 − 2 + 2 − 2}

\emph{= − 6 + 2}

\emph{= − 4}
\end{quote}

\textbf{২৫ থেকে কোনো সংখ্যাকে বিয়োগ করলে বিয়োগফল সংখ্যা অপেক্ষা ৫ বেশি
হবে?}

ক) ৭

খ) ৯

\textbf{\hl{গ) ১০ ✔}}

ঘ) ১২

\textbf{\hl{✔ সঠিক উত্তর: গ) ১০}}

\begin{quote}
\textbf{ব্যাখ্যা:} \emph{ধরি, সংখ্যাটি x}

\emph{∴ ২৫ − x = x + ৫}

\emph{⇒ x + x = ২৫ − ৫}

\emph{⇒ ২x = ২০}

\emph{∴ x = ১০}
\end{quote}

\textbf{−1 হতে কত বিয়োগ করলে বিয়োগফল শূন্য হবে?}

ক) 1

খ) 2

\textbf{\hl{গ) -1 ✔}}

ঘ) -2

\textbf{\hl{✔ সঠিক উত্তর: গ) -1}}

\begin{quote}
\textbf{ব্যাখ্যা:} \emph{ধরি, x বিয়োগ করতে হবে।}

\emph{−1 − x = 0}

\emph{বা, −1 = x}

\emph{∴ x = −1}
\end{quote}

\textbf{এক ব্যক্তির মাসিক বেতন 2x\^{}2−1 টাকা, তার x − 3 মাসের বেতন-}

ক) x\^{}3−2x\^{}2−5

খ) 2x\^{}3+6x\^{}2−6x+3

\textbf{\hl{গ) 2x\^{}3−6x\^{}2−x+3 ✔}}

ঘ) 2x\^{}3−6x\^{}2−2x−5

\textbf{\hl{✔ সঠিক উত্তর: গ) 2x\^{}3−6x\^{}2−x+3}}

\begin{quote}
\textbf{ব্যাখ্যা:} \emph{১ মাসের বেতন 2x\^{}2−1 টাকা}

\emph{∴ x − 3 মাসের বেতন = (2x\^{}2−1)×(x−3) টাকা}

\emph{= 2x\^{}3−6x\^{}2−x+3 টাকা}
\end{quote}

\textbf{3\^{}3=কত?}

ক) 9

\textbf{\hl{খ) 27 ✔}}

গ) 18

ঘ) 81

\textbf{\hl{✔ সঠিক উত্তর: খ) 27}}

\begin{quote}
\textbf{ব্যাখ্যা:} \emph{3\^{}3=3×3×3=27}
\end{quote}

\textbf{L=(2/3)M,M=(3/5)N এবং L=16 হলে N এর মান কত?}

ক) 32

\textbf{\hl{খ) 40 ✔}}

গ) 48

ঘ) 80

\textbf{\hl{✔ সঠিক উত্তর: খ) 40}}

\begin{quote}
\textbf{ব্যাখ্যা:} \emph{L=(2/3)M=(2/3)×(3/5)N=(2/5)N}

\emph{⇒ 16=(2/5)N}

\emph{⇒ N=(16×5/2)}

\emph{∴ N = 40}
\end{quote}

\textbf{a\^{}2=9,b\^{}2=16 হলে (a+b) = কত?}

\textbf{\hl{ক) ৭ ✔}}

খ) ২৫

গ) ১৩

ঘ) ১২

\textbf{\hl{✔ সঠিক উত্তর: ক) ৭}}

\begin{quote}
\textbf{ব্যাখ্যা:} \emph{a\^{}2=9 ∴ a = ±3}

\emph{b\^{}2=16 ∴ b = ±4}

\emph{ধনাত্বক মান বিবেচনা করলে,}

\emph{a + b = 3 + 4 = 7}
\end{quote}

\textbf{x\^{}4−3x−2 কে x + 1 দ্বারা ভাগ করলে ভাগশেষ কি হবে?}

ক) 1

\textbf{\hl{খ) 2 ✔}}

গ) 3

ঘ) 4

\textbf{\hl{✔ সঠিক উত্তর: খ) 2}}

\begin{quote}
\textbf{ব্যাখ্যা:} \emph{ভাগশেষ উপপাদ্য অনুসারে, f(x) = x\^{}4−3x−2 কে x + 1
দ্বারা ভাগ করলে ভাগশেষ হবে f(-1)।}

\emph{∴ নির্ণেয় ভাগশেষ = (−1)\^{}4−3(−1)−2}

\emph{= 1 + 3 − 2 = 2}
\end{quote}

\textbf{2x\^{}3+5x\^{}2−6x+4 থেকে কত বিয়োগ করলে উক্ত রাশিটি (x + 2) দ্বারা
নিঃশেষে বিভাজ্য হবে?}

ক) 11

\textbf{\hl{খ) 20 ✔}}

গ) 28

ঘ) 12

\textbf{\hl{✔ সঠিক উত্তর: খ) 20}}

\begin{quote}
\textbf{ব্যাখ্যা:} \emph{ধরি, f(x) = 2x\^{}3+5x\^{}2−6x+4}

\emph{ভাগশেষ উপপাদ্য অনুযায়ী (x + 2) দ্বারা ভাগ করলে ভাগশেষ হবে f(-2)।}

\emph{f(-2) = 2(−2)\^{}3+5(−2)\^{}2−6(−2)+4}

\emph{= −16 + 20 + 12 + 4 = 20}

\emph{∴ রাশিটি থেকে 20 বিয়োগ করলে এটি (x + 2) দ্বারা নিঃশেষে বিভাজ্য হবে।}
\end{quote}

\textbf{p\^{}2+7p+c যদি p−5 দ্বারা নিঃশেষে বিভাজ্য হয়, তবে c এর মান কত?}

\textbf{\hl{ক) - 60 ✔}}

খ) -30

গ) 5

ঘ) 60

\textbf{\hl{✔ সঠিক উত্তর: ক) - 60}}

\begin{quote}
\textbf{ব্যাখ্যা:} \emph{ধরি, f(p) = p\^{}2+7p+c}

\emph{যেহেতু রাশিটি p−5 দ্বারা নিঃশেষে বিভাজ্য, তাই ভাগশেষ উপপাদ্য অনুযায়ী
f(5) = 0 হবে।}

\emph{∴ 5\^{}2+7×5+c=0}

\emph{⇒ 25 + 35 + c = 0}

\emph{⇒ 60 + c = 0}

\emph{∴ c = − 60}
\end{quote}

\textbf{ক\^{}2=16,খ\^{}2=25 হলে কখ−এর মান হবে −}

ক) ২৫

খ) ৩০

\textbf{\hl{গ) ২০ ✔}}

ঘ) ৩৫

\textbf{\hl{✔ সঠিক উত্তর: গ) ২০}}

\begin{quote}
\textbf{ব্যাখ্যা:} \emph{ক\^{}2=16 ∴ ক = 4}

\emph{খ\^{}2=25 ∴ খ = 5}

\emph{∴ কখ = 4 × 5 = 20}
\end{quote}

\textbf{If a=(2/7),b=(1/3) and c=(1/8) , then value of y = a−3b +16c
is−/ যদি a=(2/7),b=(1/3) এবং c=(1/8) হয়, তবে y = a−3b +16c এর মান কত?}

ক) (16/7)

\textbf{\hl{খ) (9/7) ✔}}

গ) (7/16)

ঘ) 1

ঙ) (23/7)

\textbf{\hl{✔ সঠিক উত্তর: খ) (9/7)}}

\begin{quote}
\textbf{ব্যাখ্যা:} \emph{y = a − 3b + 16c}

\emph{= (2/7)−3×(1/3)+16×(1/8)}

\emph{= (2/7)−1+2}

\emph{= (2/7)+1}

\emph{= (2+7/7)}

\emph{= (9/7)}
\end{quote}

\textbf{y = x+3 এবং 2x+y = 9 হলে y-এর মান কত?}

ক) ২

\textbf{\hl{খ) ৫ ✔}}

গ) ৬

ঘ) ৮

\textbf{\hl{✔ সঠিক উত্তর: খ) ৫}}

\begin{quote}
\textbf{ব্যাখ্যা:} \emph{2x + y = 9}

\emph{⇒ 2x + x + 3 = 9 {[}∴ y = x + 3{]}}

\emph{⇒ 3x = 6}

\emph{∴ x = 2}

\emph{∴ y = 2 + 3 = 5}
\end{quote}

\textbf{x−(1/x)=4 হলে, (x/x\^{}2+x−1) এর মান কত?}

ক) 3

\textbf{\hl{খ) (1/5) ✔}}

গ) 1

ঘ) 2

\textbf{\hl{✔ সঠিক উত্তর: খ) (1/5)}}

\begin{quote}
\textbf{ব্যাখ্যা:} \emph{দেয়া আছে, x−(1/x)=4}

\emph{⇒ (x\^{}2−1/x)=4}

\emph{⇒ x\^{}2−1=4x}

\emph{এখন, (x/x\^{}2+x−1)=(x/(x\^{}2−1)+x)=(x/4x+x)=(x/5x)=(1/5)}
\end{quote}

\textbf{যদি (Q/P)=(1/4) হয় তবে (P+Q/P−Q) এর মান কত?}

\textbf{\hl{ক) (5/3) ✔}}

খ) (2/3)

গ) (3/5)

ঘ) (3/7)

\textbf{\hl{✔ সঠিক উত্তর: ক) (5/3)}}

\begin{quote}
\textbf{ব্যাখ্যা:} \emph{দেয়া আছে, (Q/P)=(1/4)}

\emph{⇒ (P/Q)=(4/1) {[}ব্যস্তকরণ করে{]}}

\emph{যোজন-বিয়োজন করে পাই,}

\emph{(P+Q/P−Q)=(4+1/4−1)=(5/3)}

\emph{Shortcut:}

\emph{P = 4 এবং Q = 1 ধরে,}

\emph{(P+Q/P−Q)=(4+1/4−1)=(5/3)}
\end{quote}

\section{39. ভাজক সংখ্যা নির্ণয়ের সূত্র
-গণিত}\label{ux9adux99cux995-ux9b8ux996ux9af-ux9a8ux9b0ux9a3ux9afux9b0-ux9b8ux9a4ux9b0--ux997ux9a3ux9a4}

গণিতের ভাজক সংখ্যা নির্ণয় একটি মৌলিক গণিতিক পদ্ধতি । ভাজক সংখ্যা নির্ণয় অত্যন্ত
গুরুত্বপূর্ণ একটি সমস্যা, যা গণিতের বিভিন্ন বিষয়ে প্রযোজ্য হয়। এই টিউটরিয়ালে আমি
গণিতের ভাজক সংখ্যা নির্ণয়ের সবচেয়ে সহজ পদ্ধতি অর্থাৎ \textbf{ভাজক সংখ্যা
নির্ণয়ের সূত্র} বিগত সালের প্রশ্ন বিশ্লেষণের মাধ্যমে আলোচনা করব। যাতে খুবই স্বল্প
সময়ে বড় বড় সংখ্যার ভাজক সংখ্যা নির্ণয় করা যায় । তবে শুরুতেই জেনে নেয়া যাক
ভাজক সংখ্যা কাকে বলে ?

\subsubsection{ভাজক সংখ্যা কাকে বলে
?}\label{ux9adux99cux995-ux9b8ux996ux9af-ux995ux995-ux9acux9b2}

ভাজক সংখ্যা হলো একটি সংখ্যা যা অন্য একটি সংখ্যাকে ভাগ করলে যে ভাগশেষ শূন্য হয়,
সেই সংখ্যাকে ভাজক সংখ্যা বলে। অর্থাৎ, যখন কোন সংখ্যা A কে সংখ্যা B দ্বারা ভাগ
করলে ভাগশেষ শূন্য পাওয়া যায় তখন B হয় ভাজক সংখ্যা। এটি সংখ্যাগত একটি বৈশিষ্ট্য
যা ব্যবহার করে আমরা সংখ্যা নির্ণয়ের সময় ব্যবহার করি।

যেমনঃ

10 কে 2 দ্বারা ভাগ করলে ভাগশেষ শূন্য পাওয়া যায়, তাই 2 হলো 10 এর ভাজক সংখ্যা।

15 কে 3 দ্বারা ভাগ করলে ভাগশেষ শূন্য পাওয়া যায়, তাই 3 হলো 15 এর ভাজক সংখ্যা।

ভাজক সংখ্যা নির্ণয় করার মাধ্যমে আমরা সংখ্যাগত বৈশিষ্ট্য ও সম্পর্কিত তথ্য পাব যা
বিভিন্ন গণিতিক সমস্যার সমাধানে ব্যবহার করা হয়।

সংক্ষেপে , যে সংখ্যা দিয়ে অন্য কোনো একটি সংখ্যাকে ভাগ করা হচ্ছে সেই সংখ্যাকে
ভাজক সংখ্যা বলা হয়। অন্যভাবে বলা যায়, ভাগফল নির্ণয় করার জন্য যে সংখ্যা দ্বারা
অন্য কোন সংখ্যাকে ভাগ করা হয়, তাকে ভাজক সংখ্যা বলে।

\subsubsection{বিভিন্ন প্রতিযোগিতামূলক পরীক্ষার প্রশ্ন ও
সমাধান}\label{ux9acux9adux9a8ux9a8-ux9aaux9b0ux9a4ux9afux997ux9a4ux9aeux9b2ux995-ux9aaux9b0ux995ux9b7ux9b0-ux9aaux9b0ux9b6ux9a8-ux993-ux9b8ux9aeux9a7ux9a8}

\textbf{প্রশ্ন নম্বর-}১. ৭২ এর ভাজক সংখ্যা কতটি? {[} ২৬ তম বিসিএস{]}

\begin{itemize}
\item
  (ক) ৯ টি
\item
  (খ) ১০ টি
\item
  (গ) ১১ টি
\item
  (ঘ) ১২ টি
\end{itemize}

\textbf{\hl{উত্তরঃ (ঘ) ১২ টি}}

\textbf{সমাধান:}

৭২ সংখ্যাটির ভাজক গুলো হলো-

১, ২, ৩, ৪, ৬, ৮, ৯, ১২, ১৮, ২৪, ৩৬, ৭২।

.এখন, ৭২ সংখ্যাটি\textquotesingle{} তো ছোট।তাই সহজে বের করা সম্ভব হয়েছে।

কিন্তু, যদি সংখ্যাটি আরো বড় হয় তাহলে এত সহজে বা মুখে মুখে বের করা সম্ভব হবে না
। তাই নিচের টেকনিক ফলো করে যে কোন সংখ্যার ভাজক সংখ্যা খুব দ্রুত এবং সহজে
নির্ণয় করুন ।

টেকনিকঃ

৭২ কে উৎপাদকে বিশ্লেষণ করলে পাই....

২

৭২

২

৩৬

২

১৮

৩

৯

৩

সুতরাং ৭২ = (২×২×২)×(৩×৩)

উপরে, লক্ষ করুন-

২ হলো ৩ বার এবং ৩ হলো ২ বার

সুতরাং

৭২ = ২\textsuperscript{৩} × ৩\textsuperscript{২} = (৩ + ১) × (২ + ১) = ৪
× ৩ = ১২টি

অর্থাৎ যত পাওয়ার {[} ২\textsuperscript{৩} এখানে পাওয়ার (৩) {]} এবং
৩\textsuperscript{২} এখানে পাওয়ার (২) {]} ।

আপনি শুধু ঐ সব পাওয়ারের সাথে ১ যোগ করে প্রাপ্ত পাওয়ার গুলো গুন করে দিলেই উত্তর
পেয়ে যাবেন।

\textbf{প্রশ্ন নম্বর-}২. ৭৫৬০০ এর ভাজক সংখ্যা কয়টি ? {[} টেলিভিশন
প্রকৌশলী-২০০৪{]}

\begin{itemize}
\item
  (ক) ১২০ টি
\item
  (খ) ১১০ টি
\item
  (গ) ১১৫ টি
\item
  (ঘ) ১২৫ টি
\end{itemize}

\textbf{\hl{উত্তরঃ (ক) ১২০ টি}}

\textbf{সমাধান:}

টেকনিকঃ

৭৫৬০০ কে উৎপাদকে বিশ্লেষণ করলে পাই....

২

৭৫৬০০

২

৩৭৮০০

২

১৮৯০০

২

৯৪৫০

৩

৪৭২৫

৩

১৫৭৫

৩

৫২৫

৫

১৭৫

৫

৩৫

৭

সুতরাং ৭৫৬০০ = (২ × ২ × ২ × ২) × (৩ × ৩ × ৩) × (৫ × ৫) × ৭

= ২\textsuperscript{৪} × ৩\textsuperscript{৩} × ৫\textsuperscript{৩} ×
৭\textsuperscript{১} = (৪ + ১) × (৩ + ১) × (২ + ১) × (১ + ১)

= ৫ × ৪ × ৩ × ২ = ১২০ (উত্তর)

\textbf{প্রশ্ন নম্বর-}৩. ১০০৮ সংখ্যাটিতে কয়টি ভাজক আছে? {[} সমাজ কল্যাণ সংগঠক
-২০০৭{]}

\begin{itemize}
\item
  (ক) ৪০ টি
\item
  (খ) ২০টি
\item
  (গ) ৩০ টি
\item
  (ঘ) ৫০ টি
\end{itemize}

\textbf{\hl{উত্তরঃ (গ) ৩০ টি}}

\textbf{সমাধান:}

টেকনিকঃ

১০০৮ কে উৎপাদকে বিশ্লেষণ করলে পাই....

২

১০০৮

২

৫০৪

২

২৫২

২

১২৬

৩

৬৩

৩

২১

৭

সুতরাং ১০০৮ = (২ × ২ × ২ × ২) × (৩ × ৩ ) × ৭

= ২\textsuperscript{৪} × ৩\textsuperscript{২} × ৭\textsuperscript{১} =
(৪ + ১) × (২ + ১) × (১ + ১)

= ৫ × ৩ × ২ = ৩০ (উত্তর)

\textbf{প্রশ্ন নম্বর-} ৪. ৫৩২০ সংখ্যাটিতে কয়টি ভাজক আছে? {[}থানা সহকারী
শিক্ষা অফিসার-২০০৬{]}

\begin{itemize}
\item
  (ক) ১২
\item
  (খ) ১৬
\item
  (গ) ১৮
\item
  (ঘ) ২০
\end{itemize}

\textbf{\hl{উত্তরঃ (খ) ১৬}}

\textbf{সমাধান:}

টেকনিকঃ

৫৩২০ কে উৎপাদকে বিশ্লেষণ করলে পাই....

২

৫৩২০

২

২৬৬০

২

১৩৩০

৫

৬৬৫

১৩৩

সুতরাং ৫৩২০ = (২ × ২ × ২ ) × ৫ × ১৩৩

= ২\textsuperscript{৩} × ৫\textsuperscript{১} × ১৩৩\textsuperscript{১} =
(৩ + ১) × (১ + ১) × (১ + ১)

= ৪ × ২ × ২ = ১৬ (উত্তর)

\section{40. সংখ্যা সম্পর্কিত সমস্যা MCQ প্রশ্ন উত্তর
অনুশীলন}\label{ux9b8ux996ux9af-ux9b8ux9aeux9aaux9b0ux995ux9a4-ux9b8ux9aeux9b8ux9af-mcq-ux9aaux9b0ux9b6ux9a8-ux989ux9a4ux9a4ux9b0-ux985ux9a8ux9b6ux9b2ux9a8}

\textbf{৫২ − ক = ৩০, \textquotesingle ক\textquotesingle{} এর মান কত?}

ক) ২০

খ) ২১

\textbf{\hl{গ) ২২ ✔}}

ঘ) ২৩

\textbf{\hl{✔ সঠিক উত্তর: গ) ২২}}

\begin{quote}
\textbf{ব্যাখ্যা:} \emph{দেয়া আছে, ৫২ − ক = ৩০}

\emph{⇒ ক = ৫২ − ৩০}

\emph{∴ ক = ২২}
\end{quote}

\textbf{(৭ + ক) × ৩ = ৩০ হলে \textquotesingle ক\textquotesingle{} এর মান
কত?}

ক) ২১

খ) ১০

গ) ৯

\textbf{\hl{ঘ) ৩ ✔}}

\textbf{\hl{✔ সঠিক উত্তর: ঘ) ৩}}

\begin{quote}
\textbf{ব্যাখ্যা:} \emph{দেয়া আছে, (৭ + ক) × ৩ = ৩০}

\emph{⇒ ৭ + ক = (৩০/৩)}

\emph{⇒ ৭ + ক = ১০}

\emph{⇒ ক = ১০ − ৭}

\emph{∴ ক = ৩}
\end{quote}

\textbf{বড় সংখ্যাটি কত?}

\textbf{\hl{ক) 2(x+2) ✔}}

খ) 2x

গ) x+2

ঘ) 2x+1

\textbf{\hl{✔ সঠিক উত্তর: ক) 2(x+2)}}

\begin{quote}
\textbf{ব্যাখ্যা:} \emph{x অঋণাত্মক পূর্ণসংখ্যা (x ≥ 0) হলে:}

\emph{2(x + 2) = 2x + 4}

\emph{যেহেতু 2x + 4 রাশিটির মান 2x, x + 2 এবং 2x + 1 এর চেয়ে বড়, তাই ক-এর
অপশনটি সবচেয়ে বড়।}
\end{quote}

\textbf{−১ থেকে কত বিয়োগ করলে বিয়োগফল শূন্য হবে?}

\textbf{\hl{ক) −১ ✔}}

খ) ১

গ) −২

ঘ) ২

\textbf{\hl{✔ সঠিক উত্তর: ক) −১}}

\begin{quote}
\textbf{ব্যাখ্যা:} \emph{ধরি, সংখ্যাটি
\textquotesingle ক\textquotesingle{}}

\emph{প্রশ্নমতে, −১ − ক = ০}

\emph{⇒ −ক = ১}

\emph{⇒ ক = −১}

\emph{∴ −১ থেকে (−১) বিয়োগ করলে বিয়োগফল শূন্য হবে। {[}যেমন: −১ − (−১) = −১
+ ১ = ০{]}}
\end{quote}

\textbf{নিচের কোন তিনটি ক্রমিক স্বাভাবিক সংখ্যার যোগফল?}

ক) ২৯

খ) ৪৬

\textbf{\hl{গ) ৫৭ ✔}}

ঘ) ৯২

\textbf{\hl{✔ সঠিক উত্তর: গ) ৫৭}}

\begin{quote}
\textbf{ব্যাখ্যা:} \emph{যেকোনো ৩টি ক্রমিক স্বাভাবিক সংখ্যার যোগফল সর্বদা ৩
দ্বারা বিভাজ্য হতে হয়।}

\emph{এখানে দেওয়া অপশনগুলোর মধ্যে কেবল ৫৭ সংখ্যাটি ৩ দ্বারা বিভাজ্য (৫৭ ÷ ৩ =
১৯)।}

\emph{ধরি, সংখ্যা তিনটি (x − ১), x, এবং (x + ১)।}

\emph{প্রশ্নমতে, (x − ১) + x + (x + ১) = ৫৭}

\emph{⇒ ৩x = ৫৭}

\emph{⇒ x = ১৯}

\emph{∴ সংখ্যা তিনটি হলো ১৮, ১৯ এবং ২০; যাদের যোগফল (১৮ + ১৯ + ২০) = ৫৭।}
\end{quote}

\textbf{কোন সংখ্যার ৬ গুণ হতে ১৫ গুণ ৬৩ বেশি?}

ক) ৬

\textbf{\hl{খ) ৭ ✔}}

গ) ৩

ঘ) ৯

\textbf{\hl{✔ সঠিক উত্তর: খ) ৭}}

\begin{quote}
\textbf{ব্যাখ্যা:} \emph{ধরি, সংখ্যাটি x}

\emph{প্রশ্নমতে, ১৫x − ৬x = ৬৩}

\emph{⇒ ৯x = ৬৩}

\emph{⇒ x = (৬৩/৯)}

\emph{∴ x = ৭}
\end{quote}

\textbf{একটি সংখ্যা ৩০১ হতে যত বড়, ৩৮১ হতে তত ছোট, সংখ্যাটি কত?}

ক) ৩৪০

\textbf{\hl{খ) ৩৪১ ✔}}

গ) ৩৪২

ঘ) ৩৪৩

\textbf{\hl{✔ সঠিক উত্তর: খ) ৩৪১}}

\begin{quote}
\textbf{ব্যাখ্যা:} \emph{সংখ্যাটি = (৩০১+৩৮১/২) = (৬৮২/২) = ৩৪১}
\end{quote}

\textbf{একটি সংখ্যা ৩১ থেকে যত বেশি ৫৫ থেকে তত কম, সংখ্যাটি কত?}

ক) ৩৯

\textbf{\hl{খ) ৪৩ ✔}}

গ) ৪৫

ঘ) ৪১

\textbf{\hl{✔ সঠিক উত্তর: খ) ৪৩}}

\begin{quote}
\textbf{ব্যাখ্যা:} \emph{ধরি, সংখ্যাটি x}

\emph{প্রশ্নমতে, x − ৩১ = ৫৫ − x}

\emph{⇒ ২x = ৫৫ + ৩১}

\emph{⇒ ২x = ৮৬}

\emph{∴ x = ৪৩}
\end{quote}

\textbf{একটি সংখ্যা ৭৪২ থেকে যত বড়, ৮৩০ থেকে তত ছোট, সংখ্যাটি কত?}

ক) ৭৮৫

\textbf{\hl{খ) ৭৮৬ ✔}}

গ) ৭৮৮

ঘ) ৭৮৭

\textbf{\hl{✔ সঠিক উত্তর: খ) ৭৮৬}}

\begin{quote}
\textbf{ব্যাখ্যা:} \emph{নির্ণেয় সংখ্যাটি = (৮৩০+৭৪২/২) = (১৫৭২/২) = ৭৮৬}
\end{quote}

\textbf{৭৬৫ থেকে ৬৫৬ যত কম, কোন সংখ্যার ৮২৫ থেকে ততটুকু বেশী?}

ক) ৯৩২

খ) ৯৩৩

\textbf{\hl{গ) ৯৩৪ ✔}}

ঘ) ৯৩৫

\textbf{\hl{✔ সঠিক উত্তর: গ) ৯৩৪}}

\begin{quote}
\textbf{ব্যাখ্যা:} \emph{পার্থক্য = ৭৬৫ − ৬৫৬ = ১০৯}

\emph{∴ নির্ণেয় সংখ্যাটি = ৮২৫ + ১০৯ = ৯৩৪}
\end{quote}

\textbf{একটি সংখ্যার ১০ গুণ ২৪০ হলে সংখ্যাটি কত?}

ক) ২০

খ) ২২

গ) ২৩

\textbf{\hl{ঘ) ২৪ ✔}}

\textbf{\hl{✔ সঠিক উত্তর: ঘ) ২৪}}

\begin{quote}
\textbf{ব্যাখ্যা:} \emph{ধরি, সংখ্যাটি x}

\emph{প্রশ্নমতে, ১০x = ২৪০}

\emph{⇒ x = (২৪০/১০)}

\emph{∴ x = ২৪}
\end{quote}

\textbf{একটি সংখ্যা ৫৬০ থেকে যত কম, ৩৮০ থেকে তার সাড়ে তিনগুণ বেশী। সংখ্যাটি
কত?}

ক) ৪৫০

খ) ৪৭০

\textbf{\hl{গ) ৫২০ ✔}}

ঘ) ৫০০

\textbf{\hl{✔ সঠিক উত্তর: গ) ৫২০}}

\begin{quote}
\textbf{ব্যাখ্যা:} \emph{ধরি, সংখ্যাটি N}

\emph{প্রশ্নমতে, (N − ৩৮০) = ৩.৫ × (৫৬০ − N)}

\emph{⇒ N − ৩৮০ = ৩.৫ × ৫৬০ − ৩.৫N}

\emph{⇒ N − ৩৮০ = ১৯৬০ − ৩.৫N}

\emph{⇒ N + ৩.৫N = ১৯৬০ + ৩৮০}

\emph{⇒ ৪.৫N = ২৩৪০}

\emph{⇒ N = (২৩৪০/৪.৫)}

\emph{∴ N = ৫২০}
\end{quote}

\textbf{একটি সংখ্যার তিনগুণের সাথে দ্বিগুণ যোগ করলে ৯০ হয়। সংখ্যাটি কত?}

ক) ১৬

খ) ১৭

\textbf{\hl{গ) ১৮ ✔}}

ঘ) ২৪

\textbf{\hl{✔ সঠিক উত্তর: গ) ১৮}}

\begin{quote}
\textbf{ব্যাখ্যা:} \emph{ধরি, সংখ্যাটি x}

\emph{প্রশ্নমতে, ৩x + ২x = ৯০}

\emph{⇒ ৫x = ৯০}

\emph{⇒ x = (৯০/৫)}

\emph{∴ x = ১৮}
\end{quote}

\textbf{একটি সংখ্যার ৪ গুণের সাথে ১০ যোগ করা হলে উত্তর হয় সংখ্যাটির ৫ গুণ
অপেক্ষা ৫ কম। সংখ্যাটি কত?}

ক) ৩০

খ) ২০

গ) ২৫

\textbf{\hl{ঘ) ১৫ ✔}}

\textbf{\hl{✔ সঠিক উত্তর: ঘ) ১৫}}

\begin{quote}
\textbf{ব্যাখ্যা:} \emph{মনে করি, সংখ্যাটি x}

\emph{প্রশ্নমতে, ৪x + ১০ = ৫x − ৫}

\emph{⇒ ৫x − ৪x = ১০ + ৫}

\emph{∴ x = ১৫}
\end{quote}

\textbf{কোন সংখ্যার ৮ গুণ থেকে ২ গুণ বিয়োগ করলে ৭২ হয়?}

ক) ২৭

খ) ১৬

\textbf{\hl{গ) ১২ ✔}}

ঘ) ৩০

\textbf{\hl{✔ সঠিক উত্তর: গ) ১২}}

\begin{quote}
\textbf{ব্যাখ্যা:} \emph{ধরি, সংখ্যাটি
\textquotesingle ক\textquotesingle{}}

\emph{প্রশ্নমতে, ৮ক − ২ক = ৭২}

\emph{⇒ ৬ক = ৭২}

\emph{∴ ক = ১২}
\end{quote}

\textbf{কোনো সংখ্যাকে ৫ দিয়ে ভাগ করে ৫ বিয়োগ করে ৯ দিয়ে গুণ করলে ১৩৫ হয়।
সংখ্যাটি কত?}

ক) ২০৫

খ) ১১০

গ) ১৫০

\textbf{\hl{ঘ) ১০০ ✔}}

\textbf{\hl{✔ সঠিক উত্তর: ঘ) ১০০}}

\begin{quote}
\textbf{ব্যাখ্যা:} \emph{সংখ্যাটি x হলে,}

\emph{\{(x ÷ ৫) − ৫\} × ৯ = ১৩৫}

\emph{⇒ \{(x ÷ ৫) − ৫\} = ১৫}

\emph{⇒ (x/৫) − ৫ = ১৫}

\emph{⇒ (x/৫) = ২০}

\emph{∴ x = ১০০}

\emph{অর্থাৎ, সংখ্যাটি ১০০।}
\end{quote}

\textbf{কোনো সংখ্যার দুই-তৃতীয়াংশ ঐ সংখ্যার চেয়ে ৫০ কম হলে সংখ্যাটি কত?}

ক) ২০৮

খ) ৩৫০

গ) ২৫০

\textbf{\hl{ঘ) ১৫০ ✔}}

\textbf{\hl{✔ সঠিক উত্তর: ঘ) ১৫০}}

\begin{quote}
\textbf{ব্যাখ্যা:} \emph{ধরি, সংখ্যাটি x}

\emph{x এর দুই তৃতীয়াংশ = (২x/৩)}

\emph{প্রশ্নমতে, x − (২x/৩) = ৫০}

\emph{⇒ (৩x−২x/৩) = ৫০}

\emph{⇒ (x/৩) = ৫০}

\emph{∴ x = ১৫০}
\end{quote}

\textbf{একটি সংখ্যার অর্ধেক তার এক তৃতীয়াংশের চাইতে ১৭ বেশি। সংখ্যাটি কত?}

ক) ৫২

খ) ৮৪

\textbf{\hl{গ) ১০২ ✔}}

ঘ) ২০৪

\textbf{\hl{✔ সঠিক উত্তর: গ) ১০২}}

\begin{quote}
\textbf{ব্যাখ্যা:} \emph{ধরি, সংখ্যাটি = x}

\emph{শর্তমতে, (x/২) − (x/৩) = ১৭}

\emph{⇒ (৩x−২x/৬) = ১৭}

\emph{⇒ (x/৬) = ১৭}

\emph{∴ x = ১৭ × ৬ = ১০২}
\end{quote}

\textbf{কোন সংখ্যার দ্বিগুণের সাথে ৩ যোগ করলে যোগফল সংখ্যাটি অপেক্ষা ৭ বেশি
হয়। সংখ্যাটি কত?}

ক) ৩

\textbf{\hl{খ) ৪ ✔}}

গ) ৫

ঘ) ৬

\textbf{\hl{✔ সঠিক উত্তর: খ) ৪}}

\begin{quote}
\textbf{ব্যাখ্যা:} \emph{ধরি, সংখ্যাটি = ক}

\emph{প্রশ্নমতে, ২ক + ৩ = ক + ৭}

\emph{⇒ ২ক − ক = ৭ − ৩}

\emph{∴ ক = ৪}
\end{quote}

\textbf{কোনো একটি নির্দিষ্ট সংখ্যাকে ১১ দিয়ে ভাগ করলে ৩ অবশিষ্ট থাকে। যদি
সংখ্যাটির চারগুণকে ১১ দিয়ে ভাগ করা হয় তাহলে অবশিষ্ট কত থাকবে?}

\textbf{\hl{ক) ১ ✔}}

খ) ৩

গ) ৯

ঘ) ১২

\textbf{\hl{✔ সঠিক উত্তর: ক) ১}}

\begin{quote}
\textbf{ব্যাখ্যা:} \emph{ধরি, সংখ্যাটি, x = ১১q + ৩ (যেখানে q একটি
পূর্ণসংখ্যা)}

\emph{⇒ ৪x = ৪(১১q + ৩)}

\emph{⇒ ৪x = ৪৪q + ১২}

\emph{এখন, ৪৪q + ১২ = ১১(৪q + ১) + ১}

\emph{∴ ১১ দ্বারা ভাগ করলে ১ অবশিষ্ট থাকবে।}
\end{quote}

\textbf{A number is 250 percent of another number. If the sum of the
numbers is 98, find the value of the larger number.}

ক) 58

খ) 60

গ) 66

\textbf{\hl{ঘ) 70 ✔}}

\textbf{\hl{✔ সঠিক উত্তর: ঘ) 70}}

\begin{quote}
\textbf{ব্যাখ্যা:} \emph{ধরি, ছোট সংখ্যাটি = x}

\emph{বড় সংখ্যাটি = x এর 250\% = x×(250/100)=2.5x}

\emph{প্রশ্নমতে, x + 2.5x = 98}

\emph{⇒ 3.5x = 98}

\emph{⇒ x = (98/3.5) = 28}

\emph{∴ বড় সংখ্যাটি = 2.5 × 28 = 70}
\end{quote}

\textbf{কোন সংখ্যার (1/3) সংখ্যাটির (1/5) অপেক্ষা 4 বেশি?}

\textbf{\hl{ক) 30 ✔}}

খ) 15

গ) 45

ঘ) 60

\textbf{\hl{✔ সঠিক উত্তর: ক) 30}}

\begin{quote}
\textbf{ব্যাখ্যা:} \emph{ধরি, সংখ্যাটি x}

\emph{∴ (x/3) − (x/5) = 4}

\emph{⇒ (5x−3x/15) = 4}

\emph{⇒ 2x = 15 × 4}

\emph{⇒ x = (15×4/2)}

\emph{∴ x = 30}
\end{quote}

\textbf{একটি ধনাত্মক সংখ্যার দ্বিগুণের বর্গের সঙ্গে ১৫ যোগ করলে ৪১৫ হয়।
সংখ্যাটি কত?}

ক) ১১

\textbf{\hl{খ) ১০ ✔}}

গ) ২০

ঘ) ১৫

\textbf{\hl{✔ সঠিক উত্তর: খ) ১০}}

\begin{quote}
\textbf{ব্যাখ্যা:} \emph{ধরি, সংখ্যাটি x}

\emph{প্রশ্নমতে, (2x)\^{}2 + ১৫ = ৪১৫}

\emph{⇒ 4x\^{}2 + ১৫ = ৪১৫}

\emph{⇒ 4x\^{}2 = ৪১৫ − ১৫}

\emph{⇒ 4x\^{}2 = ৪০০}

\emph{⇒ x\^{}2 = ১০০}

\emph{∴ x = ১০}
\end{quote}

\textbf{ছয়টি ক্রমিক পূর্ণ সংখ্যার প্রথম তিনটির গড় ৮ হলে, শেষ তিনটির যোগফল---}

ক) ২৯

খ) ৩১

\textbf{\hl{গ) ৩৩ ✔}}

ঘ) ৩৫

\textbf{\hl{✔ সঠিক উত্তর: গ) ৩৩}}

\begin{quote}
\textbf{ব্যাখ্যা:} \emph{প্রথম তিনটি ক্রমিক সংখ্যার গড় ৮ হলে, তাদের যোগফল = ৮
× ৩ = ২৪।}

\emph{ধরি, প্রথম ৩টি ক্রমিক সংখ্যা হলো x, x + ১, x + ২।}

\emph{প্রশ্নমতে, x + (x + ১) + (x + ২) = ২৪}

\emph{⇒ ৩x + ৩ = ২৪}

\emph{⇒ ৩x = ২১}

\emph{⇒ x = ৭}

\emph{∴ ছয়টি ক্রমিক সংখ্যা হলো: ৭, ৮, ৯, ১০, ১১, ১২।}

\emph{শেষ তিনটি সংখ্যার যোগফল = ১০ + ১১ + ১২ = ৩৩।}
\end{quote}

\textbf{পরপর দশটি সংখ্যার প্রথম ৫টির যোগফল ৫৬০ হলে শেষ ৫টির যোগফল কত?}

ক) ৫৮০

\textbf{\hl{খ) ৫৮৫ ✔}}

গ) ৫৭৫

ঘ) ৫৭০

\textbf{\hl{✔ সঠিক উত্তর: খ) ৫৮৫}}

\begin{quote}
\textbf{ব্যাখ্যা:} \emph{ধরি, প্রথম সংখ্যাটি x}

\emph{শর্তানুসারে,}

\emph{x + (x + ১) + (x + ২) + (x + ৩) + (x + ৪) = ৫৬০}

\emph{⇒ ৫x + ১০ = ৫৬০}

\emph{⇒ ৫x = ৫৫০}

\emph{∴ x = ১১০}

\emph{শেষ ৫টির যোগফল = (x + ৫) + (x + ৬) + (x + ৭) + (x + ৮) + (x + ৯)}

\emph{= ৫x + ৩৫ = ৫ × ১১০ + ৩৫ = ৫৫০ + ৩৫ = ৫৮৫}

\emph{∴ সঠিক উত্তর (খ)।}
\end{quote}

\textbf{তিনটি ক্রমিক স্বাভাবিক সংখ্যার গুণফল ১২০ হলে তাদের যোগফল কত?}

ক) ১০

\textbf{\hl{খ) ১৫ ✔}}

গ) ২০

ঘ) ২৪

\textbf{\hl{✔ সঠিক উত্তর: খ) ১৫}}

\begin{quote}
\textbf{ব্যাখ্যা:} \emph{১২০ এর মৌলিক উৎপাদকসমূহ = ২ × ২ × ২ × ৩ × ৫}

\emph{প্যাটার্ন অনুযায়ী পরপর তিনটি ক্রমিক স্বাভাবিক সংখ্যা সাজালে পাই: ৪, ৫
এবং ৬ (যেখানে ৪ × ৫ × ৬ = ১২০)।}

\emph{∴ সংখ্যা তিনটির যোগফল = ৪ + ৫ + ৬ = ১৫}
\end{quote}

\textbf{তিনটি ক্রমিক সংখ্যার যোগফল ১২৩। ক্ষুদ্রতম সংখ্যা দুইটির গুণফল কত?}

ক) ৬২৫

খ) ৯০০

গ) ১৬০০

\textbf{\hl{ঘ) ১৬৪০ ✔}}

\textbf{\hl{✔ সঠিক উত্তর: ঘ) ১৬৪০}}

\begin{quote}
\textbf{ব্যাখ্যা:} \emph{ধরি, সংখ্যা তিনটি = x, x + ১, x + ২}

\emph{শর্তমতে, x + x + ১ + x + ২ = ১২৩}

\emph{⇒ ৩x + ৩ = ১২৩}

\emph{⇒ ৩x = ১২০}

\emph{∴ x = ৪০}

\emph{সংখ্যা তিনটি হলো: ৪০, ৪১, ৪২।}

\emph{ক্ষুদ্রতম সংখ্যা দুইটি হলো ৪০ এবং ৪১।}

\emph{∴ তাদের গুণফল = ৪০ × ৪১ = ১৬৪০}
\end{quote}

\textbf{তিনটি ক্রমিক সংখ্যার যোগফল ৩০, বড়টি ও ছোটটির বিয়োগফল ২ হলে ছোট
সংখ্যাটি:}

ক) ৭

\textbf{\hl{খ) ৯ ✔}}

গ) ১০

ঘ) ১১

\textbf{\hl{✔ সঠিক উত্তর: খ) ৯}}

\begin{quote}
\textbf{ব্যাখ্যা:} \emph{ধরি, ক্রমিক সংখ্যা তিনটি x − ১, x এবং x + ১}

\emph{প্রশ্নমতে, (x − ১) + x + (x + ১) = ৩০}

\emph{⇒ ৩x = ৩০}

\emph{∴ x = ১০}

\emph{∴ ছোট সংখ্যাটি = (১০ − ১) = ৯}
\end{quote}

\textbf{তিনটি ক্রমিক বিজোড় সংখ্যার যোগফল ৫৭। মধ্যম সংখ্যাটি কত?}

\textbf{\hl{ক) 19 ✔}}

খ) 21

গ) 23

ঘ) 17

\textbf{\hl{✔ সঠিক উত্তর: ক) 19}}

\begin{quote}
\textbf{ব্যাখ্যা:} \emph{আমরা জানি, ৩টি ক্রমিক সংখ্যার যোগফল দেওয়া থাকলে
তাদের গড়-ই হলো মধ্যম সংখ্যা।}

\emph{∴ মধ্যম সংখ্যাটি = (৫৭/৩) = ১৯}

\emph{(বিকল্প নিয়ম: ৩টি বিজোড় সংখ্যা x − ২, x, x + ২ হলে; x − ২ + x + x +
২ = ৫৭ ⇒ ৩x = ৫৭ ∴ x = ১৯)}
\end{quote}

\textbf{১০টি সংখ্যার যোগফল ৪৬২। প্রথম ৪টি সংখ্যার গড় ৫২ এবং শেষ ৫টি সংখ্যার
গড় ৩৮। পঞ্চম সংখ্যাটি কত?}

ক) ৫০

খ) ৬০

\textbf{\hl{গ) ৬৪ ✔}}

ঘ) ৬২

\textbf{\hl{✔ সঠিক উত্তর: গ) ৬৪}}

\begin{quote}
\textbf{ব্যাখ্যা:} \emph{প্রথম ৪টি সংখ্যার সমষ্টি = (৫২ × ৪) = ২০৮}

\emph{শেষ ৫টি সংখ্যার সমষ্টি = (৩৮ × ৫) = ১৯০}

\emph{∴ পঞ্চম সংখ্যাটি = ৪৬২ − (২০৮ + ১৯০) = ৪৬২ − ৩৯৮ = ৬৪}
\end{quote}

\textbf{একটি সংখ্যার বর্গমূলের সাথে ৯ যোগ করলে যোগফল ১৪ হলে সংখ্যাটি কত?}

ক) ৪

খ) ৫

\textbf{\hl{গ) ২৫ ✔}}

ঘ) ২২৫

\textbf{\hl{✔ সঠিক উত্তর: গ) ২৫}}

\begin{quote}
\textbf{ব্যাখ্যা:} \emph{ধরি, সংখ্যাটি x}

\emph{প্রশ্নমতে, √(x) + ৯ = ১৪}

\emph{⇒ √(x) = ১৪ − ৯}

\emph{⇒ √(x) = ৫}

\emph{⇒ x = ৫\^{}২}

\emph{∴ x = ২৫}
\end{quote}

\textbf{দুটি সংখ্যার সমষ্টি ৭০ এবং অন্তরফল ১০ হলে বড় সংখ্যাটি কত?}

ক) ৩৫

\textbf{\hl{খ) ৪০ ✔}}

গ) ৪৫

ঘ) ৫০

\textbf{\hl{✔ সঠিক উত্তর: খ) ৪০}}

\begin{quote}
\textbf{ব্যাখ্যা:} \emph{আমরা জানি, বড় সংখ্যা = (সমষ্টি+অন্তরফল/২)}

\emph{∴ বড় সংখ্যাটি = (৭০+১০/২) = (৮০/২) = ৪০}
\end{quote}

\textbf{দুটি সংখ্যার যোগফল ১৫ এবং বিয়োগফল ১৩। ছোট সংখ্যাটি কত?}

\textbf{\hl{ক) ১ ✔}}

খ) ২

গ) ১৪

ঘ) ১৮

\textbf{\hl{✔ সঠিক উত্তর: ক) ১}}

\begin{quote}
\textbf{ব্যাখ্যা:} \emph{আমরা জানি, ছোট সংখ্যা = (যোগফল−বিয়োগফল/২)}

\emph{∴ ছোট সংখ্যাটি = (১৫−১৩/২) = (২/২) = ১}
\end{quote}

\textbf{দুটি সংখ্যার যোগফল ৮। যদি সংখ্যাগুলো ৩ : ১ অনুপাতে থাকে তবে সংখ্যাগুলোর
গুণফল হবে}

ক) ১০

\textbf{\hl{খ) ১২ ✔}}

গ) ১৫

ঘ) ১৮

\textbf{\hl{✔ সঠিক উত্তর: খ) ১২}}

\begin{quote}
\textbf{ব্যাখ্যা:} \emph{ধরি, সংখ্যা দুটি যথাক্রমে ৩x ও x}

\emph{∴ ৩x + x = ৮}

\emph{⇒ ৪x = ৮}

\emph{⇒ x = ২}

\emph{∴ সংখ্যা দুটির গুণফল = ৩x × x = (৩ × ২) × ২ = ১২}
\end{quote}

\textbf{একটি সংখ্যার চারগুণের সাথে ১২ যোগ করা হলে যোগফল হয় ৮। সংখ্যাটির
দ্বিগুণের সাথে ৭ যোগ করা হলে যোগফল কত হবে?}

ক) −১

খ) −১.৫

গ) −২

\textbf{\hl{ঘ) ৫ ✔}}

\textbf{\hl{✔ সঠিক উত্তর: ঘ) ৫}}

\begin{quote}
\textbf{ব্যাখ্যা:} \emph{ধরি, সংখ্যাটি x}

\emph{∴ ৪x + ১২ = ৮}

\emph{⇒ ৪x = ৮ − ১২}

\emph{⇒ ৪x = −৪}

\emph{⇒ x = −১}

\emph{∴ ২x + ৭ = ২(−১) + ৭ = −২ + ৭ = ৫}
\end{quote}

\textbf{কোন সংখ্যার দ্বিগুণের সাথে ৩ যোগ করলে যোগফল সংখ্যাটি অপেক্ষা ৭ বেশি
হয়?}

ক) ৩

\textbf{\hl{খ) ৪ ✔}}

গ) ৫

ঘ) ৬

\textbf{\hl{✔ সঠিক উত্তর: খ) ৪}}

\begin{quote}
\textbf{ব্যাখ্যা:} \emph{ধরি, সংখ্যাটি x}

\emph{প্রশ্নমতে, ২x + ৩ = x + ৭}

\emph{⇒ ২x − x = ৭ − ৩}

\emph{∴ x = ৪}
\end{quote}

\textbf{পর পর দুটি পূর্ণ সংখ্যা নির্ণয় কর যাদের বর্গের পার্থক্য হবে ৫৩ -}

ক) ২৫ এবং ২৬

\textbf{\hl{খ) ২৬ এবং ২৭ ✔}}

গ) ২৭ এবং ২৮

ঘ) ২৮ এবং ২৯

\textbf{\hl{✔ সঠিক উত্তর: খ) ২৬ এবং ২৭}}

\begin{quote}
\textbf{ব্যাখ্যা:} \emph{আমরা জানি, পর পর দুটি পূর্ণ সংখ্যার বর্গের পার্থক্য
দেওয়া থাকলে---}

\emph{বড় সংখ্যা = (বর্গের অন্তর+১/২) = (৫৩+১/২) = ২৭}

\emph{ছোট সংখ্যা = (বর্গের অন্তর−১/২) = (৫৩−১/২) = ২৬}

\emph{∴ সংখ্যা দুইটি ২৬, ২৭।}
\end{quote}

\textbf{তিনটি ধারাবাহিক পূর্ণসংখ্যার গুণফল ১২০ হলে এদের যোগফল কত?}

ক) ৯

খ) ১৪

গ) ১২

\textbf{\hl{ঘ) ১৫ ✔}}

\textbf{\hl{✔ সঠিক উত্তর: ঘ) ১৫}}

\begin{quote}
\textbf{ব্যাখ্যা:} \emph{১২০ এর মৌলিক উৎপাদকসমূহ = ২ × ২ × ২ × ৩ × ৫}

\emph{প্যাটার্ন অনুযায়ী পরপর তিনটি ধারাবাহিক সংখ্যা সাজালে পাই: ৪, ৫ এবং ৬
(যেখানে ৪ × ৫ × ৬ = ১২০)।}

\emph{∴ সংখ্যাগুলোর যোগফল = ৪ + ৫ + ৬ = ১৫}
\end{quote}

\textbf{তিনটি ক্রমিক সংখ্যার গুণফল ৬০ হলে তাদের যোগফল কত হবে? ?}

ক) ২০

খ) ১৫

গ) ১৪

\textbf{\hl{ঘ) ১২ ✔}}

\textbf{\hl{✔ সঠিক উত্তর: ঘ) ১২}}

\begin{quote}
\textbf{ব্যাখ্যা:} \emph{৬০ এর মৌলিক উৎপাদকসমূহ = ২ × ২ × ৩ × ৫}

\emph{ধারাবাহিক ৩টি ক্রমিক সংখ্যা সাজালে পাই: ৩, ৪ ও ৫ (যেখানে ৩ × ৪ × ৫ =
৬০)।}

\emph{∴ তাদের যোগফল = ৩ + ৪ + ৫ = ১২}
\end{quote}

\textbf{৩টি সংখ্যার গুণফল ২১৬। দুইটি সংখ্যা ৮ এবং ৯ হলে তৃতীয় সংখ্যাটি কত?}

\textbf{\hl{ক) ৩ ✔}}

খ) ৭

গ) ৫

ঘ) ৬

\textbf{\hl{✔ সঠিক উত্তর: ক) ৩}}

\begin{quote}
\textbf{ব্যাখ্যা:} \emph{প্রশ্নমতে,}

\emph{৮ × ৯ × তৃতীয় সংখ্যা = ২১৬}

\emph{⇒ ৭২ × তৃতীয় সংখ্যা = ২১৬}

\emph{⇒ তৃতীয় সংখ্যা = (২১৬/৭২)}

\emph{∴ তৃতীয় সংখ্যা = ৩}
\end{quote}

\textbf{তিনটি ক্রমিক সংখ্যার গুণফল তাদের যোগফলের ৮ গুণ। সংখ্যা তিনটির গড় কত?}

ক) ৪

খ) ৮

\textbf{\hl{গ) ৫ ✔}}

ঘ) ৭

\textbf{\hl{✔ সঠিক উত্তর: গ) ৫}}

\begin{quote}
\textbf{ব্যাখ্যা:} \emph{ধরি, ক্রমিক সংখ্যা তিনটি হলো (x − ১), x, এবং (x +
১)}

\emph{সংখ্যা তিনটির যোগফল = (x − ১) + x + (x + ১) = ৩x}

\emph{সংখ্যা তিনটির গুণফল = (x − ১) × x × (x + ১) = x(x\^{}২ − ১)}

\emph{প্রশ্নমতে, x(x\^{}২ − ১) = ৮ × ৩x}

\emph{⇒ x\^{}২ − ১ = ২৪ {[}উভয়পক্ষকে x দ্বারা ভাগ করে{]}}

\emph{⇒ x\^{}২ = ২৫}

\emph{⇒ x = ৫}

\emph{সংখ্যা তিনটি হলো (৫ − ১), ৫, (৫ + ১) বা ৪, ৫, ৬।}

\emph{∴ সংখ্যা তিনটির গড় (মধ্যম সংখ্যা) = ৫}
\end{quote}

\textbf{তিনটি পরপর মৌলিক সংখ্যার প্রথম দুইটি সংখ্যার গুণফল ৯১, শেষ দুইটির গুণফল
১৪৩ হলে, সংখ্যা তিনটি কত?}

\textbf{\hl{ক) ৭, ১৩, ১১ ✔}}

খ) ৭, ১১, ১৩

গ) ১১, ৭, ১৩

ঘ) ১১, ১৩, ৭

\textbf{\hl{✔ সঠিক উত্তর: ক) ৭, ১৩, ১১}}

\begin{quote}
\textbf{ব্যাখ্যা:} \emph{প্রথমাংশ: প্রথম দুইটির গুণফল = ৯১ = ৭ × ১৩}

\emph{দ্বিতীয়াংশ: শেষ দুইটির গুণফল = ১৪৩ = ১৩ × ১১}

\emph{উভয় ক্ষেত্রে সাধারণ উৎপাদক বা মধ্যম সংখ্যাটি হলো ১৩।}

\emph{∴ প্রথম সংখ্যা = ৭, দ্বিতীয় সংখ্যা = ১৩ এবং তৃতীয় সংখ্যা = ১১।}

\emph{∴ ক্রমানুসারে সংখ্যা তিনটি হলো ৭, ১৩, ১১।}
\end{quote}

\textbf{একটি সংখ্যা অন্যটির তুলনায় দুই বেশি। সংখ্যা দুটির গুণফল ৪৪০ হলে বড়
সংখ্যাটি কত?}

ক) ১৭

খ) ২০

গ) ২১

\textbf{\hl{ঘ) ২২ ✔}}

\textbf{\hl{✔ সঠিক উত্তর: ঘ) ২২}}

\begin{quote}
\textbf{ব্যাখ্যা:} \emph{ধরি, ছোট সংখ্যাটি x এবং বড় সংখ্যাটি (x + ২)}

\emph{প্রশ্নমতে, x(x + ২) = ৪৪০}

\emph{⇒ x\^{}২ + ২x − ৪৪০ = ০}

\emph{⇒ x\^{}২ + ২২x − ২০x − ৪৪০ = ০}

\emph{⇒ x(x + ২২) − ২০(x + ২২) = ০}

\emph{⇒ (x − ২০)(x + ২২) = ০}

\emph{ধনাত্মক সংখ্যা বিবেচনায়, x = ২০}

\emph{∴ বড় সংখ্যাটি = ২০ + ২ = ২২}
\end{quote}

\textbf{দুটি সংখ্যার গুণফল ৪০ এবং ভাগফল (৫/২) হলে সংখ্যাদ্বয়ের যোগফল কত?}

ক) ২২

খ) ৪১

গ) ১৩

\textbf{\hl{ঘ) ১৪ ✔}}

\textbf{\hl{✔ সঠিক উত্তর: ঘ) ১৪}}

\begin{quote}
\textbf{ব্যাখ্যা:} \emph{মনে করি, সংখ্যা দুটি x এবং y}

\emph{প্রশ্নানুসারে, xy = ৪০ ......... (i)}

\emph{এবং (x/y)=(৫/২) ⇒ x = (৫y/২) ......... (ii)}

\emph{x এর মান (i) নং-এ বসিয়ে পাই,}

\emph{(৫y/২)×y=৪০}

\emph{⇒ ৫y\^{}২ = ৮০}

\emph{⇒ y\^{}২ = ১৬}

\emph{∴ y = ৪}

\emph{(ii) নং হতে, x = (৫×৪/২) = ১০}

\emph{∴ সংখ্যাদ্বয়ের যোগফল = x + y = ১০ + ৪ = ১৪}
\end{quote}

\textbf{দুটি সংখ্যার অর্ধেকের যোগফল ৫০। তাদের পার্থক্যের এক-পঞ্চমাংশ ১০ হলে
সংখ্যা দুটি কত?}

\textbf{\hl{ক) ৭৫ ও ২৫ ✔}}

খ) ৭০ ও ৩০

গ) ৮০ ও ২০

ঘ) ৬৫ ও ৩৫

\textbf{\hl{✔ সঠিক উত্তর: ক) ৭৫ ও ২৫}}

\begin{quote}
\textbf{ব্যাখ্যা:} \emph{মনে করি, সংখ্যা দুটি x ও y}

\emph{১ম শর্তমতে, (x+y/২)=৫০ ⇒ x + y = ১০০ ......... (i)}

\emph{২য় শর্তমতে, (x−y/৫)=১০ ⇒ x − y = ৫০ ......... (ii)}

\emph{(i) + (ii) ⇒ ২x = ১৫০ ∴ x = ৭৫}

\emph{(i) হতে পাই, y = ১০০ − ৭৫ = ২৫}

\emph{∴ সংখ্যা দুটি ৭৫ ও ২৫}
\end{quote}

\textbf{একটি সংখ্যা ও তার গুণাত্মক বিপরীতের সমষ্টি ২ হলে}

\textbf{\hl{ক) ১ ✔}}

খ) ২

গ) − ২

ঘ) − ১

\textbf{\hl{✔ সঠিক উত্তর: ক) ১}}

\begin{quote}
\textbf{ব্যাখ্যা:} \emph{ধরি, সংখ্যাটি x}

\emph{সংখ্যাটির গুণাত্মক বিপরীত = (১/x)}

\emph{প্রশ্নমতে, x + (১/x) = ২}

\emph{⇒ (x\^{}২+১/x)=২}

\emph{⇒ x\^{}২ + ১ = ২x}

\emph{⇒ x\^{}২ − ২x + ১ = ০}

\emph{⇒ (x − ১)\^{}২ = ০}

\emph{∴ x = ১}
\end{quote}

\textbf{একটি বিজোড় পূর্ণ সংখ্যার পাঁচ গুণের সাথে পরবর্তী বিজোড় পূর্ণ সংখ্যার তিন
গুণ যোগ করলে ৬২ হয়। প্রথম বিজোড় পূর্ণ সংখ্যাটি কত?}

ক) ৫

খ) ৬

\textbf{\hl{গ) ৭ ✔}}

ঘ) ৮

\textbf{\hl{✔ সঠিক উত্তর: গ) ৭}}

\begin{quote}
\textbf{ব্যাখ্যা:} \emph{ধরি, প্রথম বিজোড় পূর্ণসংখ্যা = x}

\emph{∴ পরবর্তী বিজোড় পূর্ণসংখ্যা = x + ২}

\emph{প্রশ্নমতে, ৫x + ৩(x + ২) = ৬২}

\emph{⇒ ৫x + ৩x + ৬ = ৬২}

\emph{⇒ ৮x = ৬২ − ৬}

\emph{⇒ ৮x = ৫৬}

\emph{⇒ x = (৫৬/৮)}

\emph{∴ x = ৭}
\end{quote}

\textbf{দুইটি সংখ্যার বিয়োগফল 25। বৃহত্তর সংখ্যাটির 4 গুণ ক্ষুদ্রতর সংখ্যাটির 5
গুণের সমান হলে, সংখ্যা দুইটি কত?}

ক) 100, 110

\textbf{\hl{খ) 100, 125 ✔}}

গ) 110, 120

ঘ) 120, 130

\textbf{\hl{✔ সঠিক উত্তর: খ) 100, 125}}

\begin{quote}
\textbf{ব্যাখ্যা:} \emph{মনে করি, বৃহত্তর সংখ্যাটি
\textquotesingle x\textquotesingle{} এবং ক্ষুদ্রতর সংখ্যাটি
\textquotesingle y\textquotesingle{}}

\emph{∴ x − y = 25 ......... (i)}

\emph{শর্তমতে, 4x = 5y ⇒ 4x − 5y = 0 ......... (ii)}

\emph{(i) নং সমীকরণকে 5 দ্বারা গুণ করে (ii) বিয়োগ করে পাই,}

\emph{5x − 5y = 125}

\emph{4x − 5y = 0}

\emph{(−) (+) (−)}

\emph{-\/-\/-\/-\/-\/-\/-\/-\/-\/-\/-\/-\/-\/-\/-\/-\/-\/-\/-}

\emph{x = 125}

\emph{x এর মান (i) নং সমীকরণে বসিয়ে পাই,}

\emph{125 − y = 25}

\emph{⇒ −y = 25 − 125}

\emph{⇒ −y = −100}

\emph{∴ y = 100}

\emph{∴ নির্ণেয় সংখ্যা দুটি হলো (100, 125)।}
\end{quote}

\textbf{দুইটি সংখ্যার বিয়োগফল 45 এবং একটি অপরটির 4 গুণ হলে, সংখ্যা দুইটি কত?}

ক) 7 ও 10

খ) 52 ও 13

\textbf{\hl{গ) 60 ও 15 ✔}}

ঘ) 68 ও 17

\textbf{\hl{✔ সঠিক উত্তর: গ) 60 ও 15}}

\begin{quote}
\textbf{ব্যাখ্যা:} \emph{মনে করি, একটি সংখ্যা
\textquotesingle x\textquotesingle{} এবং অপর সংখ্যাটি
\textquotesingle y\textquotesingle{} (যেখানে x \textgreater{} y)}

\emph{প্রশ্নমতে,}

\emph{x − y = 45 ......... (i)}

\emph{এবং x = 4y ⇒ x − 4y = 0 ......... (ii)}

\emph{x = 4y এর মান (i) নং-এ বসিয়ে পাই,}

\emph{4y − y = 45}

\emph{⇒ 3y = 45}

\emph{∴ y = 15}

\emph{এখন, x = 4 × 15 = 60}

\emph{∴ সংখ্যা দুইটি হলো 60 ও 15।}
\end{quote}

\textbf{দুইটি সংখ্যার বিয়োগফল তাদের যোগফলের দুই-তৃতীয়াংশ। সংখ্যাদ্বয়ের অনুপাত
নীচের কোনটি?}

ক) ২:৩

খ) ৩:৪

\textbf{\hl{গ) ৫:১ ✔}}

ঘ) ৩:৫

\textbf{\hl{✔ সঠিক উত্তর: গ) ৫:১}}

\begin{quote}
\textbf{ব্যাখ্যা:} \emph{ধরি, সংখ্যা দুটি যথাক্রমে x ও y যেখানে, x
\textgreater{} y}

\emph{প্রশ্নমতে,}

\emph{x − y = (x + y) (২/৩)}

\emph{⇒ x − y = (২x+২y/৩)}

\emph{⇒ ৩x − ৩y = ২x + ২y}

\emph{⇒ ৩x − ২x = ২y + ৩y}

\emph{⇒ x = ৫y}

\emph{⇒ (x/y)=(৫/১)}

\emph{∴ x : y = ৫ : ১}
\end{quote}

\textbf{দুইটি সংখ্যার বর্গের সমষ্টি ও অন্তরফল যথাক্রমে 61 ও 11 হলে, সংখ্যা দুইটি
কি কি?}

ক) (7,6)

খ) (7,4)

গ) (12,1)

\textbf{\hl{ঘ) (6,5) ✔}}

\textbf{\hl{✔ সঠিক উত্তর: ঘ) (6,5)}}

\begin{quote}
\textbf{ব্যাখ্যা:} \emph{ধরি, সংখ্যা দুইটি x এবং y}

\emph{∴ x\^{}2 + y\^{}2 = 61 ......... (i)}

\emph{x\^{}2 − y\^{}2 = 11 ......... (ii)}

\emph{সমীকরণ (i) ও (ii) যোগ করে পাই,}

\emph{2x\^{}2 = 72 ⇒ x\^{}2 = 36 ∴ x = 6}

\emph{সমীকরণ (i) থেকে (ii) বিয়োগ করে পাই,}

\emph{2y\^{}2 = 50 ⇒ y\^{}2 = 25 ∴ y = 5}

\emph{∴ সংখ্যা দুইটি হলো (6,5)।}
\end{quote}

\textbf{দুটি সংখ্যার বর্গের সমষ্টি 13 এবং সংখ্যা দুটির গুণফল 6 হলে সংখ্যা দুটির
বর্গের অন্তর কত?}

ক) 4

\textbf{\hl{খ) 5 ✔}}

গ) 6

ঘ) 7

\textbf{\hl{✔ সঠিক উত্তর: খ) 5}}

\begin{quote}
\textbf{ব্যাখ্যা:} \emph{ধরি, একটি সংখ্যা = x ; অপর সংখ্যা = y}

\emph{∴ প্রশ্নমতে, x\^{}2 + y\^{}2 = 13 এবং xy = 6}

\emph{আমরা জানি,}

\emph{(x\^{}2 − y\^{}2)\^{}2 = (x\^{}2 + y\^{}2)\^{}2 − 4x\^{}2y\^{}2}

\emph{⇒ (x\^{}2 − y\^{}2)\^{}2 = (13)\^{}2 − 4(6)\^{}2}

\emph{⇒ (x\^{}2 − y\^{}2)\^{}2 = 169 − 144 = 25}

\emph{∴ x\^{}2 − y\^{}2 = 5}
\end{quote}

\textbf{দুটি ধনাত্মক সংখ্যার পার্থক্য ৬ এবং এদের বর্গের পার্থক্য ১০৮। সংখ্যা দুটির
যোগফল হবে-}

ক) ৬

\textbf{\hl{খ) ১৮ ✔}}

গ) ২৪

ঘ) ১২

\textbf{\hl{✔ সঠিক উত্তর: খ) ১৮}}

\begin{quote}
\textbf{ব্যাখ্যা:} \emph{মনে করি, সংখ্যা দুইটি x ও y (যেখানে x
\textgreater{} y)}

\emph{১ম শর্তানুসারে, x − y = ৬ .........(i)}

\emph{২য় শর্তানুসারে, x\^{}2 − y\^{}2 = ১০৮ .........(ii)}

\emph{আমরা জানি, x\^{}2 − y\^{}2 = (x + y)(x − y)}

\emph{⇒ ১০৮ = (x + y) × ৬}

\emph{⇒ x + y = (১০৮/৬)}

\emph{∴ x + y = ১৮}

\emph{∴ সংখ্যা দুটির যোগফল = ১৮}
\end{quote}

\textbf{দুটি ক্রমিক স্বাভাবিক সংখ্যার বর্গের অন্তর 45 হলে, সংখ্যা দুটি ---}

ক) 21, 22

\textbf{\hl{খ) 22, 23 ✔}}

গ) 23, 24

ঘ) 20, 21

\textbf{\hl{✔ সঠিক উত্তর: খ) 22, 23}}

\begin{quote}
\textbf{ব্যাখ্যা:} \emph{ধরি, ক্রমিক স্বাভাবিক সংখ্যা দুটির প্রথমটি x এবং
পরেরটি (x + 1)}

\emph{প্রশ্নমতে, (x + 1)\^{}2 − x\^{}2 = 45}

\emph{⇒ x\^{}2 + 2x + 1 − x\^{}2 = 45}

\emph{⇒ 2x + 1 = 45}

\emph{⇒ 2x = 44}

\emph{∴ x = 22}

\emph{∴ অপর সংখ্যাটি (x + 1) = 22 + 1 = 23}

\emph{∴ সংখ্যা দুটি 22, 23।}
\end{quote}

\textbf{দুটি ক্রমিক সংখ্যার বর্গের অন্তর ১১ হলে, সংখ্যা দুইটির বর্গের সমষ্টি কত?}

ক) ১৬

খ) ১৭

\textbf{\hl{গ) ৬১ ✔}}

ঘ) ৭১

\textbf{\hl{✔ সঠিক উত্তর: গ) ৬১}}

\begin{quote}
\textbf{ব্যাখ্যা:} \emph{ধরি, সংখ্যা দুটি x ও (x + ১)}

\emph{∴ (x + ১)\^{}২ − x\^{}২ = ১১}

\emph{⇒ x\^{}২ + ২x + ১ − x\^{}২ = ১১}

\emph{⇒ ২x = ১১ − ১}

\emph{⇒ ২x = ১০}

\emph{∴ x = ৫ এবং (x + ১) = ৬}

\emph{∴ সংখ্যা দুইটির বর্গের সমষ্টি = (x + ১)\^{}২ + x\^{}২ = ৬\^{}২ +
৫\^{}২ = ৩৬ + ২৫ = ৬১}
\end{quote}

\textbf{দুইটি ক্রমিক সংখ্যার বর্গের অন্তর ১৯৯ হলে, ছোট সংখ্যাটি কত?}

\textbf{\hl{ক) ৯৯ ✔}}

খ) ৯৮

গ) ১০০

ঘ) ১০১

\textbf{\hl{✔ সঠিক উত্তর: ক) ৯৯}}

\begin{quote}
\textbf{ব্যাখ্যা:} \emph{ছোট সংখ্যাটি = (বর্গের অন্তর−১/২) = (১৯৯−১/২) =
(১৯৮/২) = ৯৯}

\emph{অথবা,}

\emph{ধরি, ক্রমিক সংখ্যা দুটি x ও (x + ১)}

\emph{প্রশ্নমতে, (x + ১)\^{}২ − x\^{}২ = ১৯৯}

\emph{⇒ x\^{}২ + ২x + ১ − x\^{}২ = ১৯৯}

\emph{⇒ ২x + ১ = ১৯৯}

\emph{⇒ ২x = ১৯৮}

\emph{∴ x = (১৯৮/২) = ৯৯}

\emph{∴ ছোট সংখ্যাটি ৯৯।}
\end{quote}

\textbf{দুটি ধনাত্মক সংখ্যার বর্গের অন্তর 8, গুণফল 3, সংখ্যা দুটির বর্গের সমষ্টি
কত?}

ক) 8

\textbf{\hl{খ) 10 ✔}}

গ) 13

ঘ) 25

\textbf{\hl{✔ সঠিক উত্তর: খ) 10}}

\begin{quote}
\textbf{ব্যাখ্যা:} \emph{ধরি, সংখ্যা দুটি x ও y}

\emph{∴ x\^{}2 − y\^{}2 = 8 এবং xy = 3}

\emph{আমরা জানি,}

\emph{(x\^{}2 + y\^{}2)\^{}2 = (x\^{}2 − y\^{}2)\^{}2 + 4x\^{}2y\^{}2}

\emph{⇒ (x\^{}2 + y\^{}2)\^{}2 = (8)\^{}2 + 4(3)\^{}2}

\emph{⇒ (x\^{}2 + y\^{}2)\^{}2 = 64 + 36 = 100}

\emph{∴ x\^{}2 + y\^{}2 = 10}
\end{quote}

\textbf{দুটি সংখ্যার বিয়োগফলের অর্ধেক ২। বড় সংখ্যাটির সঙ্গে ছোট সংখ্যাটির
দ্বিগুণ যোগ করলে যোগফল ১৩ হয়। সংখ্যা দুটি কত?}

\textbf{\hl{ক) ৭, ৩ ✔}}

খ) ৮, ৪

গ) ৯, ২

ঘ) ৯, ৪

\textbf{\hl{✔ সঠিক উত্তর: ক) ৭, ৩}}

\begin{quote}
\textbf{ব্যাখ্যা:} \emph{ধরি, ছোট সংখ্যাটি = y এবং বড় সংখ্যাটি = x}

\emph{১ম শর্তমতে, (১/২)x−y=২ ⇒ x − y = ৪ ......... (১)}

\emph{২য় শর্তমতে, x + ২y = ১৩ ......... (২)}

\emph{(২) নং সমীকরণ হতে (১) নং সমীকরণ বিয়োগ করে পাই,}

\emph{৩y = ৯ ⇒ y = ৩}

\emph{y এর মান (১) নং সমীকরণে বসিয়ে পাই, x − ৩ = ৪ ⇒ x = ৭}

\emph{∴ নির্ণেয় সংখ্যাদ্বয় ৭, ৩}
\end{quote}

\textbf{একটি স্বাভাবিক সংখ্যার বর্গের সহিত সংখ্যাটি যোগ করিলে উহা পরবর্তী
স্বাভাবিক সংখ্যার ৯ গুণের সমান হয়। সংখ্যাটি কত?}

ক) ৭

খ) ৮

\textbf{\hl{গ) ৯ ✔}}

ঘ) ১০

\textbf{\hl{✔ সঠিক উত্তর: গ) ৯}}

\begin{quote}
\textbf{ব্যাখ্যা:} \emph{মনে করি, সংখ্যাটি = x ; ∴ পরবর্তী সংখ্যা = x + ১}

\emph{প্রশ্নমতে,}

\emph{x\^{}২ + x = ৯(x + ১)}

\emph{⇒ x\^{}২ + x = ৯x + ৯}

\emph{⇒ x\^{}২ + x − ৯x − ৯ = ০}

\emph{⇒ x(x + ১) − ৯(x + ১) = ০}

\emph{⇒ (x + ১)(x − ৯) = ০}

\emph{∴ x + ১ = ০ অথবা x − ৯ = ০}

\emph{⇒ x = −১ অথবা x = ৯}

\emph{যেহেতু −১ স্বাভাবিক সংখ্যা নয়, তাই x = ৯।}

\emph{∴ নির্ণেয় সংখ্যাটি ৯।}
\end{quote}

\textbf{কোন দুটি সংখ্যার যোগফল ১০ এবং গুণফল ২৪?}

ক) ৪, −৬

খ) −৬, −৪

গ) ১২, −২

\textbf{\hl{ঘ) ৪, ৬ ✔}}

\textbf{\hl{✔ সঠিক উত্তর: ঘ) ৪, ৬}}

\begin{quote}
\textbf{ব্যাখ্যা:} \emph{ধরি, সংখ্যাদ্বয় x, y}

\emph{x + y = ১০ ......... (১)}

\emph{xy = ২৪ ......... (২)}

\emph{(১) নং সমীকরণ হতে পাই, x = ১০ − y}

\emph{x এর মান (২) নং সমীকরণে বসিয়ে পাই,}

\emph{(১০ − y) y = ২৪}

\emph{⇒ ১০y − y\^{}২ = ২৪}

\emph{⇒ y\^{}২ − ১০y + ২৪ = ০}

\emph{⇒ (y − ৬)(y − ৪) = ০}

\emph{∴ y = ৪, ৬}

\emph{y = ৪ হলে x = ১০ − ৪ = ৬}

\emph{y = ৬ হলে x = ১০ − ৬ = ৪}

\emph{∴ নির্ণেয় সংখ্যাদ্বয় ৬, ৪ (বা ৪, ৬)।}
\end{quote}

\textbf{২টি সংখ্যার যোগফল ৪৮ এবং তাদের গুণফল ৪৩২। তবে বড় সংখ্যাটি কত?}

\textbf{\hl{ক) ৩৬ ✔}}

খ) ৩৭

গ) ৩৮

ঘ) ৪০

\textbf{\hl{✔ সঠিক উত্তর: ক) ৩৬}}

\begin{quote}
\textbf{ব্যাখ্যা:} \emph{ধরি, সংখ্যা দুটি a ও b (যেখানে a \textgreater{} b)}

\emph{∴ a + b = ৪৮ ......... (i) এবং ab = ৪৩২}

\emph{আমরা জানি, (a − b)\^{}২ = (a + b)\^{}২ − ৪ab}

\emph{⇒ (a − b)\^{}২ = (৪৮)\^{}২ − ৪ × ৪৩২}

\emph{⇒ (a − b)\^{}২ = ২৩০৪ − ১৭২৮ = ৫৭৬}

\emph{∴ a − b = ২৪ ......... (ii)}

\emph{(i) + (ii) ⇒ ২a = ৭২ ∴ a = ৩৬}

\emph{∴ বড় সংখ্যাটি ৩৬।}
\end{quote}

\textbf{দুটি সংখ্যার যোগফল ১৭ এবং গুণফল ৭২। ছোট সংখ্যাটি কত?}

ক) ৬

\textbf{\hl{খ) ৮ ✔}}

গ) ৯

ঘ) কোনোটিই নয়

\textbf{\hl{✔ সঠিক উত্তর: খ) ৮}}

\begin{quote}
\textbf{ব্যাখ্যা:} \emph{ধরি, সংখ্যাদ্বয় x ও y}

\emph{অতএব x + y = ১৭ ......... (১)}

\emph{এবং xy = ৭২}

\emph{আমরা জানি, (x − y)\^{}২ = (x + y)\^{}২ − ৪xy}

\emph{⇒ (x − y)\^{}২ = (১৭)\^{}২ − ৪ × ৭২ = ২৮৯ − ২৮৮ = ১}

\emph{∴ x − y = ১ ......... (২)}

\emph{(১) ও (২) হতে পাই, বড় সংখ্যাটি ৯ এবং ছোট সংখ্যাটি ৮।}

\emph{∴ নির্ণেয় ছোট সংখ্যা = ৮।}
\end{quote}

\textbf{দুইটি সংখ্যার গুণফল 42 এবং বিয়োগফল 1 হলে সংখ্যা দু\textquotesingle টি
কত?}

ক) 4, 3

\textbf{\hl{খ) 7, 6 ✔}}

গ) 8, 6

ঘ) 10, 8

\textbf{\hl{✔ সঠিক উত্তর: খ) 7, 6}}

\begin{quote}
\textbf{ব্যাখ্যা:} \emph{ধরি, সংখ্যাদ্বয় x, y (যেখানে x \textgreater{} y)}

\emph{প্রশ্নমতে, xy = 42 এবং x − y = 1}

\emph{আমরা জানি, (x + y)\^{}2 = (x − y)\^{}2 + 4xy}

\emph{⇒ (x + y)\^{}2 = (1)\^{}2 + 4 × 42 = 1 + 168 = 169}

\emph{∴ x + y = 13}

\emph{এখন, x + y = 13 এবং x − y = 1 সমীকরণ দুটি সমাধান করে পাই,}

\emph{2x = 14 ⇒ x = 7}

\emph{এবং 2y = 12 ⇒ y = 6}

\emph{∴ সংখ্যা দু\textquotesingle টি হলো 7, 6।}
\end{quote}

\textbf{একটি সংখ্যার বর্গের সাথে ৪ যোগ করলে যোগফল ৪০ হয়। সংখ্যাটি কত?}

ক) ৪

খ) ৫

\textbf{\hl{গ) ৬ ✔}}

ঘ) ৭

\textbf{\hl{✔ সঠিক উত্তর: গ) ৬}}

\begin{quote}
\textbf{ব্যাখ্যা:} \emph{ধরি, সংখ্যাটি x}

\emph{∴ x\^{}২ + ৪ = ৪০}

\emph{⇒ x\^{}২ = ৪০ − ৪}

\emph{⇒ x\^{}২ = ৩৬}

\emph{⇒ x = √(৩৬)}

\emph{∴ x = ৬}
\end{quote}

\textbf{দুটি সংখ্যার অন্তর ১২, বড়টির সঙ্গে ১ যোগ করলে ছোটটির দ্বিগুণ হয়। সংখ্যা
দুটি কি কি?}

ক) ৩৫, ২৩

খ) ২০, ৮

গ) ৩০, ১৮

\textbf{\hl{ঘ) ২৫, ১৩ ✔}}

\textbf{\hl{✔ সঠিক উত্তর: ঘ) ২৫, ১৩}}

\begin{quote}
\textbf{ব্যাখ্যা:} \emph{ধরি, সংখ্যা দুটি x ও y (x \textgreater{} y)}

\emph{১ম শর্তমতে, x − y = ১২ ......... (i)}

\emph{২য় শর্তমতে, x + ১ = ২y ⇒ x − ২y = −১ ......... (ii)}

\emph{(i) নং সমীকরণ হতে (ii) নং সমীকরণ বিয়োগ করে পাই,}

\emph{y = ১৩}

\emph{y এর মান (i) নং সমীকরণে বসাই,}

\emph{x − ১৩ = ১২ ⇒ x = ১৩ + ১২ = ২৫}

\emph{∴ নির্ণেয় সংখ্যাদ্বয় ২৫, ১৩}
\end{quote}

\textbf{দুই অংকবিশিষ্ট একটি সংখ্যার দশক স্থানীয় অংক একক স্থানীয় অংকের দ্বিগুণ।
দেখাও যে, সংখ্যাটি অংকদ্বয়ের সমষ্টির -}

ক) ৩ গুণ (3 times)

খ) ৫ গুণ (5 times)

গ) ৬ গুণ (6 times)

\textbf{\hl{ঘ) ৭ গুণ (7 times) ✔}}

\textbf{\hl{✔ সঠিক উত্তর: ঘ) ৭ গুণ (7 times)}}

\begin{quote}
\textbf{ব্যাখ্যা:} \emph{ধরি, একক স্থানীয় অংক = x}

\emph{∴ দশক স্থানীয় অংক = ২x}

\emph{সংখ্যাটি = ১০ × ২x + x = ২১x}

\emph{অংকদ্বয়ের সমষ্টি = x + ২x = ৩x}

\emph{∴ সংখ্যাটি/অংকদ্বয়ের সমষ্টি = (২১x/৩x) = ৭}

\emph{∴ সংখ্যাটি অংকদ্বয়ের সমষ্টির ৭ গুণ (7 times)।}
\end{quote}

\textbf{দুই অংক বিশিষ্ট একটি সংখ্যার একক স্থানীয় অংক a এবং দশক স্থানীয় অংক b
হলে সংখ্যাটি হবে---}

ক) ba

খ) b + a

গ) 10ab

\textbf{\hl{ঘ) 10b + a ✔}}

\textbf{\hl{✔ সঠিক উত্তর: ঘ) 10b + a}}

\begin{quote}
\textbf{ব্যাখ্যা:} \emph{আমরা জানি, দুই অংক বিশিষ্ট কোনো সংখ্যার মান = (১০ ×
দশক স্থানীয় অংক) + একক স্থানীয় অংক।}

\emph{এখানে, একক স্থানীয় অংক = a এবং দশক স্থানীয় অংক = b}

\emph{∴ সংখ্যাটি = ১০ × b + a = 10b + a}
\end{quote}

\textbf{দুই অংকে বিশিষ্ট একটি সংখ্যার এককের অংক দশকের অংক অপেক্ষা ৩ বেশি।
সংখ্যাটি এর অংকদ্বয়ের সমষ্টির ৩ গুণ অপেক্ষা ৪ বেশি। সংখ্যাটি কত?}

ক) ৪৭

খ) ৩৬

\textbf{\hl{গ) ২৫ ✔}}

ঘ) ১৪

\textbf{\hl{✔ সঠিক উত্তর: গ) ২৫}}

\begin{quote}
\textbf{ব্যাখ্যা:} \emph{মনেকরি, দশক স্থানীয় অংক = x}

\emph{∴ একক স্থানীয় অংক = x + ৩}

\emph{∴ সংখ্যাটি = ১০x + x + ৩ = ১১x + ৩}

\emph{প্রশ্নমতে,}

\emph{১১x + ৩ = ৩(x + x + ৩) + ৪}

\emph{⇒ ১১x + ৩ = ৩(২x + ৩) + ৪}

\emph{⇒ ১১x + ৩ = ৬x + ৯ + ৪}

\emph{⇒ ১১x − ৬x = ১৩ − ৩}

\emph{⇒ ৫x = ১০}

\emph{∴ x = ২}

\emph{∴ দশক স্থানীয় অংক = ২ এবং একক স্থানীয় অংক = ২ + ৩ = ৫}

\emph{∴ সংখ্যাটি = ২৫}
\end{quote}

\textbf{দুই অংক বিশিষ্ট একটি সংখ্যার অংকদ্বয়ের সমষ্টি 9। সংখ্যাটি হতে 9 বিয়োগ
করলে এর অংকদ্বয় স্থান বিনিময় করে। সংখ্যাটি কত?}

ক) 34

খ) 67

\textbf{\hl{গ) 54 ✔}}

ঘ) 23

\textbf{\hl{✔ সঠিক উত্তর: গ) 54}}

\begin{quote}
\textbf{ব্যাখ্যা:} \emph{মনে করি, দশক স্থানীয় অংক = y এবং একক স্থানীয় অংক =
x}

\emph{∴ সংখ্যাটি = 10y + x}

\emph{শর্তানুসারে, x + y = 9 ......... (i)}

\emph{প্রশ্নমতে, (10y + x) − 9 = 10x + y}

\emph{⇒ 10y − y + x − 10x = 9}

\emph{⇒ 9y − 9x = 9}

\emph{⇒ y − x = 1 ......... (ii)}

\emph{(i) ও (ii) যোগ করে পাই,}

\emph{2y = 10 ⇒ y = 5}

\emph{y এর মান (i) নং সমীকরণে বসিয়ে পাই, x = 9 − 5 = 4}

\emph{∴ সংখ্যাটি = (10 × 5) + 4 = 54}
\end{quote}

\textbf{দুই অংকবিশিষ্ট একটি সংখ্যার একক স্থানীয় অংক ৬। সংখ্যাটি তার অংকদ্বয়ের
সমষ্টির এক চতুর্থাংশের সমান হলে সংখ্যাটি কত?}

ক) ৪৭

খ) ৭১

\textbf{\hl{গ) ৩৬ ✔}}

ঘ) ৮৭

\textbf{\hl{✔ সঠিক উত্তর: গ) ৩৬}}

\begin{quote}
\textbf{ব্যাখ্যা:} \emph{ধরি, দশক স্থানীয় অংক = x}

\emph{∴ সংখ্যাটি = ১০x + ৬}

\emph{শর্তমতে, অংকদ্বয়ের সমষ্টি (x + ৬) সংখ্যাটির এক-চতুর্থাংশের সমান,}

\emph{∴ (১০x+৬/৪)=x+৬}

\emph{⇒ ১০x + ৬ = ৪x + ২৪}

\emph{⇒ ১০x − ৪x = ২৪ − ৬}

\emph{⇒ ৬x = ১৮}

\emph{∴ x = ৩}

\emph{∴ সংখ্যাটি = (১০ × ৩) + ৬ = ৩৬}
\end{quote}

\textbf{দুই অংক বিশিষ্ট একটি সংখ্যার এককের অংক দশকের অংক অপেক্ষা ৩ বেশি।
সংখ্যাটি এর অংকদ্বয়ের সমষ্টির তিনগুণ অপেক্ষা ৪ বেশি। সংখ্যাটি কত?}

ক) ৪৭

\textbf{\hl{খ) ২৫ ✔}}

গ) ৩৬

ঘ) ১৪

\textbf{\hl{✔ সঠিক উত্তর: খ) ২৫}}

\begin{quote}
\textbf{ব্যাখ্যা:} \emph{ধরি, দশক স্থানীয় অংক x}

\emph{∴ একক স্থানীয় অংক (x + ৩)}

\emph{∴ সংখ্যাটি = ১০x + (x + ৩) = ১১x + ৩}

\emph{শর্তানুসারে, ১১x + ৩ = ৩(x + x + ৩) + ৪}

\emph{⇒ ১১x + ৩ = ৩(২x + ৩) + ৪}

\emph{⇒ ১১x + ৩ = ৬x + ৯ + ৪}

\emph{⇒ ১১x − ৬x = ১৩ − ৩}

\emph{⇒ ৫x = ১০}

\emph{∴ x = ২}

\emph{∴ সংখ্যাটি = ১১ × ২ + ৩ = ২৫}
\end{quote}

\textbf{দুই অংকে বিশিষ্ট একটি সংখ্যা, অংকদ্বয়ের স্থান বিনিময়ের ফলে 54 বৃদ্ধি
পায়। অংকে দুইটির যোগফল 12 হলে সংখ্যাটি কত?}

ক) 57

খ) 75

\textbf{\hl{গ) 39 ✔}}

ঘ) 93

\textbf{\hl{✔ সঠিক উত্তর: গ) 39}}

\begin{quote}
\textbf{ব্যাখ্যা:} \emph{ধরি, দশক স্থানীয় অংক = x এবং একক স্থানীয় অংক = y}

\emph{∴ মূল সংখ্যাটি = 10x + y এবং স্থান বিনিময়কৃত সংখ্যাটি = 10y + x}

\emph{১ম শর্তানুসারে, (10y + x) − (10x + y) = 54}

\emph{⇒ 9y − 9x = 54}

\emph{⇒ y − x = 6 ......... (i)}

\emph{২য় শর্তানুসারে, x + y = 12 ⇒ y + x = 12 ......... (ii)}

\emph{(ii) − (i) করে পাই,}

\emph{2x = 6 ∴ x = 3}

\emph{x এর মান (ii) নং সমীকরণে বসাই,}

\emph{3 + y = 12 ∴ y = 9}

\emph{∴ মূল সংখ্যাটি = (10 × 3) + 9 = 39}
\end{quote}

\textbf{দুই অংকবিশিষ্ট একটি সংখ্যার অংকদ্বয় স্থান পরিবর্তন করলে সংখ্যাটি
পূর্বাপেক্ষা ৫৪ বৃদ্ধি পায়। সংখ্যাটির অংকদ্বয়ের পার্থক্য কত?}

ক) ৩

খ) ৫

\textbf{\hl{গ) ৬ ✔}}

ঘ) ৭

\textbf{\hl{✔ সঠিক উত্তর: গ) ৬}}

\begin{quote}
\textbf{ব্যাখ্যা:} \emph{ধরি, সংখ্যাটির একক স্থানীয় অংক x এবং দশক স্থানীয় অংক
y}

\emph{∴ সংখ্যাটি = ১০y + x}

\emph{প্রশ্নমতে, ১০y + x + ৫৪ = ১০x + y}

\emph{⇒ ১০x − x − ১০y + y = ৫৪}

\emph{⇒ ৯x − ৯y = ৫৪}

\emph{⇒ ৯(x − y) = ৫৪}

\emph{∴ x − y = ৬}
\end{quote}

\textbf{দুই অঙ্কবিশিষ্ট কোন সংখ্যার অঙ্ক দুটির অন্তর 2, অঙ্ক দুটি স্থান বিনিময়
করলে যে সংখ্যা পাওয়া যায়, তা প্রদত্ত সংখ্যার দ্বিগুণ অপেক্ষা ৬ কম। সংখ্যাটি কত?}

ক) 57

খ) 46

গ) 53

\textbf{\hl{ঘ) 24 ✔}}

\textbf{\hl{✔ সঠিক উত্তর: ঘ) 24}}

\begin{quote}
\textbf{ব্যাখ্যা:} \emph{ধরি, একক স্থানীয় অংক = x ; দশক স্থানীয় অংক = y}

\emph{∴ সংখ্যাটি = 10y + x}

\emph{১ম শর্তমতে, x − y = 2}

\emph{⇒ x = y + 2 ......... (1)}

\emph{২য় শর্তমতে, (10x + y) = 2(10y + x) − 6}

\emph{⇒ 10x + y = 20y + 2x − 6}

\emph{⇒ 8x − 19y = −6}

\emph{x এর মান বসিয়ে পাই,}

\emph{8(y + 2) − 19y = −6}

\emph{⇒ 8y + 16 − 19y = −6}

\emph{⇒ −11y = −22}

\emph{∴ y = 2}

\emph{(1) হতে, x = 2 + 2 = 4}

\emph{∴ নির্ণেয় সংখ্যাটি = 10 × 2 + 4 = 24}
\end{quote}

\textbf{দুই অংক বিশিষ্ট কোন সংখ্যার অঙ্কদ্বয়ের সমষ্টি ৯। অঙ্ক দুটি স্থান বিনিময়
করলে যে সংখ্যা পাওয়া যায় তা প্রদত্ত সংখ্যা হতে ৪৫ কম। সংখ্যাটি নির্ণয় করুন।}

ক) ৫৪

খ) ৬৩

\textbf{\hl{গ) ৭২ ✔}}

ঘ) ৮১

\textbf{\hl{✔ সঠিক উত্তর: গ) ৭২}}

\begin{quote}
\textbf{ব্যাখ্যা:} \emph{ধরি, একক স্থানীয় অংক = x এবং দশক স্থানীয় অংক = y}

\emph{∴ সংখ্যাটি = ১০y + x}

\emph{১ম শর্তমতে, x + y = ৯ ......... (১)}

\emph{২য় শর্তমতে, (১০y + x) − (১০x + y) = ৪৫}

\emph{⇒ ৯y − ৯x = ৪৫}

\emph{⇒ y − x = ৫ ......... (২)}

\emph{(১) ও (২) নং সমীকরণ যোগ করে পাই,}

\emph{২y = ১৪ ∴ y = ৭}

\emph{(১) নং সমীকরণে y-এর মান বসাই,}

\emph{x + ৭ = ৯ ∴ x = ২}

\emph{∴ নির্ণেয় সংখ্যা = (১০ × ৭ + ২) = ৭২}
\end{quote}

\textbf{দুই অংকবিশিষ্ট একটি সংখ্যাকে অংকদ্বয়ের গুণফল দ্বারা ভাগ করলে ভাগফল হয়
3, সংখ্যাটির সাথে 18 যোগ করলে অংকদ্বয় স্থান বিনিময় করে। সংখ্যাটির কত?}

ক) 36

খ) 30

গ) 28

\textbf{\hl{ঘ) 24 ✔}}

\textbf{\hl{✔ সঠিক উত্তর: ঘ) 24}}

\begin{quote}
\textbf{ব্যাখ্যা:} \emph{একক স্থানীয় অংক = x ; দশক স্থানীয় অংক = y ; ∴
সংখ্যাটি = 10y + x}

\emph{১ম শর্তমতে, (10y+x/xy)=3 ......... (i)}

\emph{২য় শর্তমতে, 10y + x + 18 = 10x + y}

\emph{⇒ 9x = 9y + 18}

\emph{⇒ x = y + 2 ......... (ii)}

\emph{x এর মান (i) নং সমীকরণে বসাই :}

\emph{(10y+y+2/y(y+2))=3}

\emph{⇒ 11y + 2 = 3y\^{}2 + 6y}

\emph{⇒ 3y\^{}2 − 5y − 2 = 0}

\emph{⇒ 3y\^{}2 − 6y + y − 2 = 0}

\emph{⇒ (y − 2)(3y + 1) = 0}

\emph{∴ y = 2 {[}যেহেতু y অঋণাত্মক পূর্ণসংখ্যা, তাই y ≠ −(1/3) {]}}

\emph{y = 2 হলে (ii) হতে পাই, x = 2 + 2 = 4}

\emph{∴ নির্ণেয় সংখ্যা = 10 × 2 + 4 = 24}
\end{quote}

\textbf{তিনটি ক্রমিক পূর্ণসংখ্যা নির্ণয় করুন, যাদের প্রথম দুইটির গুণফল শেষ দুইটির
গুণফল অপেক্ষা 10 কম।}

\textbf{\hl{ক) 4,5,6 ✔}}

খ) 4,6,8

গ) 4,3,2

ঘ) 1,4,6

\textbf{\hl{✔ সঠিক উত্তর: ক) 4,5,6}}

\begin{quote}
\textbf{ব্যাখ্যা:} \emph{ধরি, ১ম ক্রমিক সংখ্যা = x}

\emph{২য় ক্রমিক সংখ্যা = x + 1}

\emph{৩য় ক্রমিক সংখ্যা = x + 2}

\emph{প্রশ্নমতে, x(x + 1) = (x + 1)(x + 2) − 10}

\emph{⇒ x\^{}2 + x = x\^{}2 + 3x + 2 − 10}

\emph{⇒ x\^{}2 + x = x\^{}2 + 3x − 8}

\emph{⇒ 2x = 8}

\emph{∴ x = 4}

\emph{∴ ১ম ক্রমিক সংখ্যা = 4 ; ২য় ক্রমিক সংখ্যা = 5 ; ৩য় ক্রমিক সংখ্যা = 6}

\emph{∴ নির্ণেয় সংখ্যা তিনটি: 4,5,6।}
\end{quote}

\textbf{যদি \textquotesingle ক\textquotesingle{} এবং
\textquotesingle খ\textquotesingle{} দুইটি পূর্ণসংখ্যা হয় এবং ৫ক + ৩খ = ১৭
হয়, তাহলে নিচের কোনটি \textquotesingle খ\textquotesingle{} এর মান হতে
পারে?}

ক) ২

খ) ৩

\textbf{\hl{গ) ৪ ✔}}

ঘ) ৫

\textbf{\hl{✔ সঠিক উত্তর: গ) ৪}}

\begin{quote}
\textbf{ব্যাখ্যা:} \emph{দেওয়া আছে, ৫ক + ৩খ = ১৭}

\emph{ক. খ = ২ হলে, ৫ক + ৩ × ২ = ১৭ ⇒ ৫ক = ১১ ⇒ ক = (১১/৫) ; যা পূর্ণসংখ্যা
নয়।}

\emph{খ. খ = ৩ হলে, ৫ক + ৩ × ৩ = ১৭ ⇒ ৫ক = ৮ ⇒ ক = (৮/৫) ; যা পূর্ণসংখ্যা
নয়।}

\emph{গ. খ = ৪ হলে, ৫ক + ৩ × ৪ = ১৭ ⇒ ৫ক = ৫ ⇒ ক = ১ ; যা একটি
পূর্ণসংখ্যা।}

\emph{∴ সঠিক উত্তর ৪।}
\end{quote}

\textbf{কোন সংখ্যার এক-তৃতীয়াংশ ও এক-চতুর্থাংশের পার্থক্য ২.৫?}

ক) ৯

খ) ১৮

\textbf{\hl{গ) ৩০ ✔}}

ঘ) ৩৬

\textbf{\hl{✔ সঠিক উত্তর: গ) ৩০}}

\begin{quote}
\textbf{ব্যাখ্যা:} \emph{ধরি, সংখ্যাটি x}

\emph{প্রশ্নমতে, (x/৩)−(x/৪)=২.৫}

\emph{⇒ (৪x−৩x/১২)=২.৫}

\emph{⇒ (x/১২)=২.৫}

\emph{⇒ x = ১২ × ২.৫}

\emph{∴ x = ৩০}
\end{quote}

\textbf{দুটি ক্রমিক পূর্ণ বর্গসংখ্যার বর্গের অন্তর ৯৩ হলে সংখ্যাদ্বয় কত ?}

\textbf{\hl{ক) ৪৬,৪৭ ✔}}

খ) ৪৪,৪৫

গ) ৪৩,৪৪

ঘ) ৫০,৫১

\textbf{\hl{✔ সঠিক উত্তর: ক) ৪৬,৪৭}}

\begin{quote}
\textbf{ব্যাখ্যা:} \emph{বড় সংখ্যা = (৯৩+১/২) = (৯৪/২) = ৪৭}

\emph{ছোট সংখ্যা = (৯৩−১/২) = (৯২/২) = ৪৬}

\emph{∴ সংখ্যাদ্বয় যথাক্রমে ৪৬ ও ৪৭।}
\end{quote}

\textbf{দুটি ক্রমিক পূর্ণ বর্গসংখ্যার বর্গের অন্তর ১৯৯, বড় সংখ্যাটি কত ?}

ক) ৭০

খ) ৮০

গ) ৯০

\textbf{\hl{ঘ) ১০০ ✔}}

\textbf{\hl{✔ সঠিক উত্তর: ঘ) ১০০}}

\begin{quote}
\textbf{ব্যাখ্যা:} \emph{বড় সংখ্যা = (১৯৯+১/২) = (২০০/২) = ১০০}
\end{quote}

\textbf{দুটি ক্রমিক পূর্ণ বর্গসংখ্যার বর্গের অন্তর যদি ৪৭ হয় তবে বড় সংখ্যাটি কত
?}

\textbf{\hl{ক) ২৪ ✔}}

খ) ২৫

গ) ২৬

ঘ) ৩০

\textbf{\hl{✔ সঠিক উত্তর: ক) ২৪}}

\begin{quote}
\textbf{ব্যাখ্যা:} \emph{বড় সংখ্যা = (৪৭+১/২) = (৪৮/২) = ২৪}
\end{quote}

\textbf{কোন দুটি ক্রমিক পূর্ণ সংখ্যার বর্গের অন্তর যদি ৪৭ হয় তবে সংখ্যা দুটি কত?}

\textbf{\hl{ক) ২৩ এবং ২৪ ✔}}

খ) ২৪ এবং ২৫

গ) ২২ এবং ২৩

ঘ) ২১ এবং ২২

\textbf{\hl{✔ সঠিক উত্তর: ক) ২৩ এবং ২৪}}

\begin{quote}
\textbf{ব্যাখ্যা:} \emph{বড় সংখ্যা = (৪৭+১/২) = (৪৮/২) = ২৪}

\emph{ছোট সংখ্যা = (৪৭−১/২) = (৪৬/২) = ২৩}

\emph{∴ সংখ্যাদ্বয় যথাক্রমে ২৩ ও ২৪।}
\end{quote}

\textbf{দুইটি ধারাবাহিক সংখ্যার বর্গের পার্থক্য ৩৩। ক্ষুদ্রতম সংখ্যাটি কত হবে ?}

ক) ১২

খ) ১৫

\textbf{\hl{গ) ১৬ ✔}}

ঘ) ১৮

\textbf{\hl{✔ সঠিক উত্তর: গ) ১৬}}

\begin{quote}
\textbf{ব্যাখ্যা:} \emph{ক্ষুদ্রতম সংখ্যা = (৩৩−১/২) = (৩২/২) = ১৬}
\end{quote}

\textbf{দুটি ক্রমিক সংখ্যার বর্গের অন্তর ৯১। সংখ্যা দুটি কি কি ?}

\textbf{\hl{ক) ৪৫,৪৬ ✔}}

খ) ৪৬,৪৮

গ) ৫৮,৪৫

ঘ) ৪০,৫০

\textbf{\hl{✔ সঠিক উত্তর: ক) ৪৫,৪৬}}

\begin{quote}
\textbf{ব্যাখ্যা:} \emph{বড় সংখ্যা = (৯১+১/২) = (৯২/২) = ৪৬}

\emph{ছোট সংখ্যা = (৯১−১/২) = (৯০/২) = ৪৫}

\emph{∴ সংখ্যাদ্বয় যথাক্রমে ৪৫ ও ৪৬।}
\end{quote}

\textbf{একটি সংখ্যা ৬৫০ থেকে যত বড় ৮২০ থেকে তত ছোট। সংখ্যাটি কত ?}

ক) ৭৩০

\textbf{\hl{খ) ৭৩৫ ✔}}

গ) ৮০০

ঘ) ৭৮০

\textbf{\hl{✔ সঠিক উত্তর: খ) ৭৩৫}}

\begin{quote}
\textbf{ব্যাখ্যা:} \emph{নির্ণেয় সংখ্যাটি = (৬৫০+৮২০/২) = (১৪৭০/২) = ৭৩৫}
\end{quote}

\textbf{একটি সংখ্যা ৯৯৯ থেকে যত ছোট ৭৯৭ থেকে তত বড়। সংখ্যাটি কত ?}

ক) ৮৯৭

\textbf{\hl{খ) ৮৯৮ ✔}}

গ) ৮৯৯

ঘ) ৯০০

\textbf{\hl{✔ সঠিক উত্তর: খ) ৮৯৮}}

\begin{quote}
\textbf{ব্যাখ্যা:} \emph{নির্ণেয় সংখ্যাটি = (৯৯৯+৭৯৭/২) = (১৭৯৬/২) = ৮৯৮}
\end{quote}

\textbf{দুটি সংখ্যার গুণফল ২৩০৪। একটি সংখ্যার ৪ গুণ ৯৬ হলে, অপর সংখ্যাটি কত ?}

ক) ৯৪

খ) ৯২

\textbf{\hl{গ) ৯৬ ✔}}

ঘ) ৯৮

\textbf{\hl{✔ সঠিক উত্তর: গ) ৯৬}}

\begin{quote}
\textbf{ব্যাখ্যা:} \emph{এখানে, একটি সংখ্যা = ৯৬ ÷ ৪ = ২৪}

\emph{∴ অপর সংখ্যাটি = দুইটি সংখ্যার গুণফল ÷ একটি সংখ্যা = ২৩০৪ ÷ ২৪ = ৯৬}
\end{quote}

\textbf{২টি সংখ্যার গুণফল ৮.৮। একটি সংখ্যা ২.৭৫, অপর সংখ্যাটি কত ?}

\textbf{\hl{ক) ৩.২০ ✔}}

খ) ৩.৮

গ) ৪.২

ঘ) ৪.৮

\textbf{\hl{✔ সঠিক উত্তর: ক) ৩.২০}}

\begin{quote}
\textbf{ব্যাখ্যা:} \emph{অপর সংখ্যাটি = দুইটি সংখ্যার গুণফল ÷ একটি সংখ্যা = ৮.৮
÷ ২.৭৫ = ৩.২}
\end{quote}

\textbf{দু'টি সংখ্যার বর্গের সমষ্টি ৪১, সংখ্যা দু'টির গুণফল ২০। সংখ্যা দু'টির
বর্গের বিয়োগফল কত ?}

ক) ৭

\textbf{\hl{খ) ৯ ✔}}

গ) ১১

ঘ) ১৩

\textbf{\hl{✔ সঠিক উত্তর: খ) ৯}}

\begin{quote}
\textbf{ব্যাখ্যা:} \emph{মনে করি, সংখ্যা দুটি যথাক্রমে a ও b।}

\emph{প্রশ্নমতে, a\^{}2+b\^{}2=৪১ এবং ab = ২০}

\emph{আমরা জানি, ((a+b))\^{}2=a\^{}2+b\^{}2+2ab}

\emph{⇒ ((a+b))\^{}2=৪১+২×২০=৮১=৯\^{}2}

\emph{⇒ a + b = ৯}

\emph{আবার, (a−b)\^{}2=a\^{}2+b\^{}2−2ab}

\emph{⇒ (a−b)\^{}2=৪১−২×২০=১}

\emph{⇒ a − b = ১}

\emph{∴ a\^{}2−b\^{}2=(a+b)(a−b)=৯×১=৯}

\emph{বিকল্প (সংক্ষিপ্ত পদ্ধতি):}

\emph{এখানে, ৪\^{}2+৫\^{}2=১৬+২৫=৪১ এবং ৪ × ৫ = ২০।}

\emph{সুতরাং, সংখ্যা দুটি ৪ ও ৫।}

\emph{∴ বর্গের বিয়োগফল = ৫\^{}2−৪\^{}2=২৫−১৬=৯}
\end{quote}

\textbf{৩টি সংখ্যার গড় ৩৩। দুইটি সংখ্যা ২৪ এবং ৪২ হলে অপর সংখ্যাটি কত ?}

ক) ১২

খ) ২২

গ) ৩২

\textbf{\hl{ঘ) ৩৩ ✔}}

\textbf{\hl{✔ সঠিক উত্তর: ঘ) ৩৩}}

\begin{quote}
\textbf{ব্যাখ্যা:} \emph{৩টি সংখ্যার সমষ্টি = ৩৩ × ৩ = ৯৯}

\emph{২টি সংখ্যার সমষ্টি = ২৪ + ৪২ = ৬৬}

\emph{∴ অপর সংখ্যাটি = ৯৯ − ৬৬ = ৩৩}
\end{quote}

\textbf{২টি সংখ্যার যোগফল ১৩০ এবং বিয়োগফল ১০ হলে বৃহত্তম সংখ্যাটি কত?}

ক) ৬০

\textbf{\hl{খ) ৭০ ✔}}

গ) ৮০

ঘ) ৯০

\textbf{\hl{✔ সঠিক উত্তর: খ) ৭০}}

\begin{quote}
\textbf{ব্যাখ্যা:} \emph{মনে করি, সংখ্যা দুটি যথাক্রমে x ও y}

\emph{প্রশ্নমতে, x + y = ১৩০ \ldots\ldots\ldots\ldots\ldots(1)}

\emph{x − y = ১০ \ldots\ldots\ldots\ldots\ldots(2)}

\emph{(1) + (2) করে পাই,}

\emph{২x = ১৪০}

\emph{⇒ x = ১৪০ ÷ ২ = ৭০}

\emph{(1) নং হতে পাই, y = ১৩০ − ৭০ = ৬০}

\emph{∴ বৃহত্তম সংখ্যাটি ৭০।}

\emph{বিকল্প (সংক্ষিপ্ত পদ্ধতি):}

\emph{বৃহত্তম সংখ্যাটি = (১৩০+১০/২) = (১৪০/২) = ৭০}
\end{quote}

\textbf{যদি দুটি সংখ্যার গুণফল ৯৬ এবং যোগফল ২২ হয় তবে ছোট সংখ্যাটি কত ?}

ক) ৪

\textbf{\hl{খ) ৬ ✔}}

গ) ৮

ঘ) কোনটিই নয়

\textbf{\hl{✔ সঠিক উত্তর: খ) ৬}}

\begin{quote}
\textbf{ব্যাখ্যা:} \emph{ধরি, বড় সংখ্যাটি x এবং ছোট সংখ্যাটি y}

\emph{প্রশ্নমতে, xy = ৯৬ এবং x + y = ২২ \ldots\ldots\ldots\ldots(1)}

\emph{আমরা জানি, (x−y)\^{}2=((x+y))\^{}2−4xy}

\emph{⇒ (x−y)\^{}2=২২\^{}2−৪×৯৬}

\emph{⇒ (x−y)\^{}2=৪৮৪−৩৮৪=১০০}

\emph{∴ x − y = √(১০০) = ১০ \ldots\ldots\ldots\ldots(2)}

\emph{(1) + (2) করে পাই, ২x = ৩২ ⇒ x = ১৬}

\emph{∴ ছোট সংখ্যা, y = ৯৬ ÷ ১৬ = ৬}

\emph{বিকল্প (সংক্ষিপ্ত পদ্ধতি):}

\emph{এমন দুটি সংখ্যা খুঁজতে হবে যাদের গুণফল ৯৬ এবং যোগফল ২২।}

\emph{এখানে, ১৬ × ৬ = ৯৬ এবং ১৬ + ৬ = ২২।}

\emph{∴ ছোট সংখ্যাটি ৬।}
\end{quote}

\textbf{দুটি সংখ্যার গুণফল ২৬৮৮ এবং ভাগফল (৬/৭) , সংখ্যা দুটি কত ?}

ক) ৩৮,৩৬

খ) ৫৬,৪০

\textbf{\hl{গ) ৪৮,৫৬ ✔}}

ঘ) ৩৮,৫৬

\textbf{\hl{✔ সঠিক উত্তর: গ) ৪৮,৫৬}}

\begin{quote}
\textbf{ব্যাখ্যা:} \emph{মনে করি, সংখ্যা দুটি ৬ক এবং ৭ক}

\emph{প্রশ্নমতে, ৬ক × ৭ক = ২৬৮৮}

\emph{⇒ ৪২ ক\^{}2 = ২৬৮৮}

\emph{⇒ ক\^{}2 = ২৬৮৮ ÷ ৪২ = ৬৪}

\emph{⇒ ক = ৮ {[}উভয়পক্ষকে বর্গমূল করে{]}}

\emph{∴ সংখ্যা দুটি যথাক্রমে ৬ × ৮ = ৪৮ এবং ৭ × ৮ = ৫৬}
\end{quote}

\textbf{কোন সংখ্যার এক পঞ্চমাংশের সাথে ৪ বৃদ্ধি করলে সংখ্যাটির মান ঐ সংখ্যার
এক চতুর্থাংশের চেয়ে ১০ কমে যায়। সংখ্যাটি কত ?}

ক) ২৫০

খ) ২৭০

\textbf{\hl{গ) ২৮০ ✔}}

ঘ) ২০০০

\textbf{\hl{✔ সঠিক উত্তর: গ) ২৮০}}

\begin{quote}
\textbf{ব্যাখ্যা:} \emph{মনে করি, সংখ্যাটি = ক}

\emph{প্রশ্নমতে, (ক/৫)+৪=(ক/৪)−১০}

\emph{⇒ (ক/৪)−(ক/৫)=৪+১০}

\emph{⇒ (৫ক−৪ক/২০)=১৪}

\emph{⇒ (ক/২০)=১৪}

\emph{⇒ ক = ২৮০}

\emph{∴ সংখ্যাটি ২৮০।}
\end{quote}

\textbf{একটি সংখ্যার অর্ধেক তার এক-তৃতীয়াংশের চেয়ে ১৭ বেশি। সংখ্যাটি কত ?}

ক) ৫২

খ) ৮৪

গ) ১০০

\textbf{\hl{ঘ) ১০২ ✔}}

\textbf{\hl{✔ সঠিক উত্তর: ঘ) ১০২}}

\begin{quote}
\textbf{ব্যাখ্যা:} \emph{মনে করি, সংখ্যাটি = ক}

\emph{প্রশ্নমতে, (ক/২)−(ক/৩)=১৭}

\emph{⇒ (৩ক−২ক/৬)=১৭}

\emph{⇒ (ক/৬)=১৭}

\emph{⇒ ক = ১০২}

\emph{∴ সংখ্যাটি = ১০২।}
\end{quote}

\textbf{দুটি সংখ্যার অর্ধেকের যোগফল ৫১। তাদের পার্থক্যের এক চতুর্থাংশ ১৩। সংখ্যা
দুটি কত ?}

ক) (৭৮,২৬)

\textbf{\hl{খ) (৭৭,২৫) ✔}}

গ) (৭৭,২৮)

ঘ) (৭৮,২৫)

\textbf{\hl{✔ সঠিক উত্তর: খ) (৭৭,২৫)}}

\begin{quote}
\textbf{ব্যাখ্যা:} \emph{মনে করি, সংখ্যা দুটি যথাক্রমে ক ও খ}

\emph{প্রথম শর্তমতে, (ক+খ/২)=৫১}

\emph{⇒ ক + খ = ১০২ \ldots\ldots\ldots\ldots\ldots(১)}

\emph{দ্বিতীয় শর্তমতে, (ক−খ/৪)=১৩}

\emph{⇒ ক − খ = ৫২ \ldots\ldots\ldots\ldots\ldots(২)}

\emph{(১) + (২) করে পাই,}

\emph{২ক = ১৫৪ ⇒ ক = ৭৭}

\emph{(১) − (২) করে পাই,}

\emph{২খ = ৫০ ⇒ খ = ২৫}

\emph{∴ সংখ্যা দুটি যথাক্রমে ৭৭ ও ২৫।}
\end{quote}

\textbf{একটি ধনাত্মক সংখ্যার সাথে ৪ যোগ করে যোগফলকে বর্গ করলে ৬২৫ হয়।
সংখ্যাটি কত ?}

ক) ২০

খ) ২২

\textbf{\hl{গ) ২১ ✔}}

ঘ) ২৫

\textbf{\hl{✔ সঠিক উত্তর: গ) ২১}}

\begin{quote}
\textbf{ব্যাখ্যা:} \emph{মনে করি, ধনাত্মক সংখ্যাটি = ক}

\emph{প্রশ্নমতে, ((ক+৪))\^{}2=৬২৫}

\emph{⇒ ক + ৪ = √(৬২৫) {[}উভয়পক্ষকে বর্গমূল করে{]}}

\emph{⇒ ক + ৪ = ২৫}

\emph{⇒ ক = ২৫ − ৪ = ২১}

\emph{∴ সংখ্যাটি = ২১।}
\end{quote}

\textbf{কোন সংখ্যার (১/২) অংশের সাথে ৬ যোগ করলে সংখ্যাটির (২/৩) অংশ হবে।
সংখ্যাটি কত ?}

ক) ৫৩

খ) ২৪

\textbf{\hl{গ) ৩৬ ✔}}

ঘ) ৩৫

\textbf{\hl{✔ সঠিক উত্তর: গ) ৩৬}}

\begin{quote}
\textbf{ব্যাখ্যা:} \emph{মনে করি, সংখ্যাটি = ক}

\emph{প্রশ্নমতে, (ক/২)+৬=(২ক/৩)}

\emph{⇒ (২ক/৩)−(ক/২)=৬}

\emph{⇒ (৪ক−৩ক/৬)=৬}

\emph{⇒ (ক/৬)=৬}

\emph{⇒ ক = ৩৬}

\emph{∴ নির্ণেয় সংখ্যাটি ৩৬।}

\emph{বিকল্প পদ্ধতি:}

\emph{প্রশ্নমতে, ((২/৩)−(১/২)) অংশ = ৬}

\emph{⇒ (৪−৩/৬) অংশ = ৬}

\emph{⇒ (১/৬) অংশ = ৬}

\emph{∴ সম্পূর্ণ সংখ্যাটি = ৬ × ৬ = ৩৬}
\end{quote}

\textbf{কোন সংখ্যা হতে ১৭৫ বিয়োগ করে ১৩০ যোগ করলে ২৯৭ হবে ?}

ক) ৬০২

\textbf{\hl{খ) ৩৪২ ✔}}

গ) ২৫২

ঘ) কোনটিই নয়

\textbf{\hl{✔ সঠিক উত্তর: খ) ৩৪২}}

\begin{quote}
\textbf{ব্যাখ্যা:} \emph{মনে করি, সংখ্যাটি = ক}

\emph{প্রশ্নমতে, ক − ১৭৫ + ১৩০ = ২৯৭}

\emph{⇒ ক − ৪৫ = ২৯৭}

\emph{⇒ ক = ২৯৭ + ৪৫}

\emph{⇒ ক = ৩৪২}

\emph{∴ সংখ্যাটি ৩৪২।}
\end{quote}

\textbf{একটি সংখ্যার ১২\% নিলে ৯৬ পাওয়া যায়, সংখ্যাটি কত ?}

ক) ১২০

খ) ১০০০

গ) ৭২০

\textbf{\hl{ঘ) ৮০০ ✔}}

\textbf{\hl{✔ সঠিক উত্তর: ঘ) ৮০০}}

\begin{quote}
\textbf{ব্যাখ্যা:} \emph{মনে করি, সংখ্যাটি = ক}

\emph{প্রশ্নমতে, ক এর ১২\% = ৯৬}

\emph{⇒ (১২/১০০)×ক=৯৬}

\emph{⇒ ক = (৯৬×১০০/১২)}

\emph{⇒ ক = ৮০০}

\emph{∴ সংখ্যাটি ৮০০।}
\end{quote}

\textbf{একটি সংখ্যার এক-চতুর্থাংশ হতে ৪ বিয়োগ করলে ২০ হয়। সংখ্যাটি কত?}

ক) ৪৮

খ) ২৪

গ) ৩৬

\textbf{\hl{ঘ) ৯৬ ✔}}

\textbf{\hl{✔ সঠিক উত্তর: ঘ) ৯৬}}

\begin{quote}
\textbf{ব্যাখ্যা:} \emph{মনে করি, সংখ্যাটি = ক}

\emph{প্রশ্নমতে, (ক/৪)−৪=২০}

\emph{⇒ (ক/৪)=২০+৪}

\emph{⇒ (ক/৪)=২৪}

\emph{⇒ ক = ২৪ × ৪}

\emph{⇒ ক = ৯৬}

\emph{∴ সংখ্যাটি ৯৬।}
\end{quote}

\textbf{কোন লঘিষ্ঠ সংখ্যাকে ১২ এবং ১৬ দ্বারা ভাগ করলে অবশিষ্ট হবে যথাক্রমে ৫
এবং ৯ ?}

ক) ৪৮

খ) ৩৯

গ) ৫৩

\textbf{\hl{ঘ) ৪১ ✔}}

\textbf{\hl{✔ সঠিক উত্তর: ঘ) ৪১}}

\begin{quote}
\textbf{ব্যাখ্যা:} \emph{এখানে, ভাজক ও অবশিষ্টাংশের পার্থক্য:}

\emph{১২ − ৫ = ৭ এবং ১৬ − ৯ = ৭ (উভয় ক্ষেত্রে পার্থক্য সমান অর্থাৎ ৭)}

\emph{এখন, ১২ এবং ১৬ এর ল.সা.গু নির্ণয় করি: ১২ = ২ × ২ × ৩ ১৬ = ২ × ২ × ২
× ২}

\emph{∴ ল.সা.গু = ২ × ২ × ২ × ২ × ৩ = ৪৮}

\emph{∴ নির্ণেয় লঘিষ্ঠ সংখ্যাটি = (ভাজকদ্বয়ের ল.সা.গু) − (পার্থক্য)}

\emph{= ৪৮ − ৭ = ৪১।}
\end{quote}

\textbf{২৫ থেকে কোন সংখ্যাটি বিয়োগ করলে বিয়োগফল সংখ্যাটি অপেক্ষা ৫ বেশি হবে
?}

ক) ৭

খ) ৯

\textbf{\hl{গ) ১০ ✔}}

ঘ) ১২

\textbf{\hl{✔ সঠিক উত্তর: গ) ১০}}

\begin{quote}
\textbf{ব্যাখ্যা:} \emph{মনে করি, সংখ্যাটি = ক}

\emph{প্রশ্নমতে, ২৫ − ক = ক + ৫}

\emph{⇒ ২৫ − ৫ = ক + ক}

\emph{⇒ ২ক = ২০}

\emph{⇒ ক = (২০/২) = ১০}

\emph{∴ সংখ্যাটি ১০।}
\end{quote}

\textbf{একটি সংখ্যা ঐ সংখ্যার ২০\% অপেক্ষা ৪০ বেশি। সংখ্যাটি কত?}

ক) ৪৫

\textbf{\hl{খ) ৫০ ✔}}

গ) ৫৫

ঘ) ৬০

\textbf{\hl{✔ সঠিক উত্তর: খ) ৫০}}

\begin{quote}
\textbf{ব্যাখ্যা:} \emph{মনে করি, সংখ্যাটি = ক}

\emph{∴ ক এর ২০\% = (২০ক/১০০)}

\emph{প্রশ্নমতে, (২০ক/১০০)+৪০=ক}

\emph{⇒ ২০ক + ৪০০০ = ১০০ক}

\emph{⇒ ১০০ক − ২০ক = ৪০০০}

\emph{⇒ ৮০ক = ৪০০০}

\emph{⇒ ক = (৪০০০/৮০) = ৫০}

\emph{∴ সংখ্যাটি ৫০।}
\end{quote}

\textbf{কোন সংখ্যা থেকে ৪ বিয়োগ করলে বিয়োগফল ঐ সংখ্যার বিপরীত সংখ্যার ২১ গুণ
হয়। সংখ্যাটি কত ?}

ক) ৫

\textbf{\hl{খ) ৭ ✔}}

গ) ৯

ঘ) ১২

\textbf{\hl{✔ সঠিক উত্তর: খ) ৭}}

\begin{quote}
\textbf{ব্যাখ্যা:} \emph{মনে করি, সংখ্যাটি = ক}

\emph{প্রশ্নমতে, ক − ৪ = (১/ক)×২১}

\emph{⇒ ক − ৪ = (২১/ক)}

\emph{⇒ ক\^{}২−৪ক−২১=০}

\emph{⇒ ক\^{}২−৭ক+৩ক−২১=০}

\emph{⇒ ক(ক − ৭) + ৩(ক − ৭) = ০}

\emph{⇒ (ক − ৭)(ক + ৩) = ০}

\emph{যেহেতু ঋণাত্মক মান গ্রহণযোগ্য নয়, তাই ক + ৩ ≠ ০}

\emph{∴ ক − ৭ = ০ ⇒ ক = ৭}

\emph{∴ সংখ্যাটি ৭।}
\end{quote}

\textbf{কোন সংখ্যার অর্ধেকের বর্গ থেকে সংখ্যাটি বাদ দিলে ৪৮ হয়। সংখ্যাটি কত
?}

ক) ৮

খ) ১২

\textbf{\hl{গ) ১৬ ✔}}

ঘ) ২৪

\textbf{\hl{✔ সঠিক উত্তর: গ) ১৬}}

\begin{quote}
\textbf{ব্যাখ্যা:} \emph{বীজগাণিতিক পদ্ধতি:}

\emph{মনে করি, সংখ্যাটি = ক}

\emph{প্রশ্নমতে, ((ক/২))\^{}২−ক=৪৮}

\emph{⇒ (ক\^{}২/৪)−ক=৪৮}

\emph{⇒ ক\^{}২−৪ক=১৯২}

\emph{⇒ ক\^{}২−৪ক−\textless{}১৯২=০}

\emph{⇒ ক\^{}২−১৬ক+১২ক−\textless{}১৯২=০}

\emph{⇒ ক(ক − ১৬) + ১২(ক − ১৬) = ০}

\emph{⇒ (ক − ১৬)(ক + ১২) = ০}

\emph{∴ ক = ১৬ (যেহেতু ধনাত্মক মান গ্রহণযোগ্য)}

\emph{শর্টকাট/অপশন টেস্ট:}

\emph{১৬ এর অর্ধেক = ৮}

\emph{৮ এর বর্গ = ৬৪}

\emph{৬৪ থেকে ১৬ বাদ দিলে হয় (৬৪ − ১৬) = ৪৮। যা প্রদত্ত শর্ত পূরণ করে।}
\end{quote}

\textbf{কোন সংখ্যার চার পঞ্চমাংশ ঐ সংখ্যার দুই তৃতীয়াংশ এর চেয়ে ৮ বেশী।
সংখ্যাটি কত ?}

ক) ৪৮

\textbf{\hl{খ) ৬০ ✔}}

গ) ৭২

ঘ) ৮৪

\textbf{\hl{✔ সঠিক উত্তর: খ) ৬০}}

\begin{quote}
\textbf{ব্যাখ্যা:} \emph{মনে করি, সংখ্যাটি = ক}

\emph{প্রশ্নমতে, (৪ক/৫)−(২ক/৩)=৮}

\emph{⇒ (১২ক−১০ক/১৫)=৮}

\emph{⇒ (২ক/১৫)=৮}

\emph{⇒ ২ক = ১২০}

\emph{⇒ ক = ৬০}

\emph{∴ সংখ্যাটি = ৬০}
\end{quote}

\textbf{কোন সংখ্যার তিন চতুর্থাংশ ঐ সংখ্যার দুই তৃতীয়াংশ অপেক্ষা ৫ বেশি।
সংখ্যাটি কত ?}

ক) ৪৫

খ) ৫০

গ) ৫৫

\textbf{\hl{ঘ) ৬০ ✔}}

\textbf{\hl{✔ সঠিক উত্তর: ঘ) ৬০}}

\begin{quote}
\textbf{ব্যাখ্যা:} \emph{মনে করি, সংখ্যাটি = ক}

\emph{প্রশ্নমতে, (৩ক/৪)−(২ক/৩)=৫}

\emph{⇒ (৯ক−৮ক/১২)=৫}

\emph{⇒ (ক/১২)=৫}

\emph{⇒ ক = ৬০}

\emph{∴ সংখ্যাটি = ৬০}
\end{quote}

\textbf{একটি সংখ্যা ঐ সংখ্যার চার সপ্তমাংশ অপেক্ষা ১৮ বেশি। সংখ্যাটি কত ?}

ক) ২৮

খ) ৩২

\textbf{\hl{গ) ৪২ ✔}}

ঘ) ৪৮

\textbf{\hl{✔ সঠিক উত্তর: গ) ৪২}}

\begin{quote}
\textbf{ব্যাখ্যা:} \emph{মনে করি, সংখ্যাটি = ক}

\emph{প্রশ্নমতে, ১৮+(৪ক/৭)=ক}

\emph{⇒ ক−(৪ক/৭)=১৮}

\emph{⇒ (৭ক−৪ক/৭)=১৮}

\emph{⇒ ৩ক = ১২৬}

\emph{⇒ ক = ৪২}

\emph{∴ সংখ্যাটি ৪২}
\end{quote}

\textbf{১০ থেকে ৬০ পর্যন্ত যে সকল মৌলিক সংখ্যার একক স্থানীয় অংক ৯, তাদের
সমষ্টি কত?}

ক) ১৩০

\textbf{\hl{খ) ১০৭ ✔}}

গ) ১১৩

ঘ) ১৪৬

\textbf{\hl{✔ সঠিক উত্তর: খ) ১০৭}}

\begin{quote}
\textbf{ব্যাখ্যা:} \emph{১০ থেকে ৬০ পর্যন্ত যে সকল মৌলিক সংখ্যার একক স্থানীয়
অংক ৯, সে সকল সংখ্যাগুলো হলো: ১৯, ২৯, এবং ৫৯।}

\emph{∴ এদের সমষ্টি = ১৯ + ২৯ + ৫৯ = ১০৭।}
\end{quote}

\textbf{ভাজক ভাগফলের ১০ গুণ, ভাজক ০.৫ হলে ভাজ্য কত?}

\textbf{\hl{ক) ০.০২৫ ✔}}

খ) ০.২৫

গ) ২৫

ঘ) ২.৫

\textbf{\hl{✔ সঠিক উত্তর: ক) ০.০২৫}}

\begin{quote}
\textbf{ব্যাখ্যা:} \emph{দেওয়া আছে, ভাজক = ০.৫}

\emph{প্রশ্নমতে, ভাজক = ১০ × ভাগফল}

\emph{⇒ ভাগফল = (০.৫/১০) = ০.০৫}

\emph{আমরা জানি, ভাজ্য = ভাজক × ভাগফল}

\emph{∴ ভাজ্য = ০.৫ × ০.০৫ = ০.০২৫}
\end{quote}

\textbf{ছয়টি পরপর পূর্ণ সংখ্যা দেয়া আছে। ১ম তিনটি যোগফল ২৭ হলে শেষ ৩টির
যোগফল কত?}

\textbf{\hl{ক) ৩৬ ✔}}

খ) ৩৩

গ) ৩২

ঘ) ৩০

\textbf{\hl{✔ সঠিক উত্তর: ক) ৩৬}}

\begin{quote}
\textbf{ব্যাখ্যা:} \emph{বীজগাণিতিক পদ্ধতি:}

\emph{মনে করি, প্রথম ৩টি সংখ্যা যথাক্রমে x, x + ১, x + ২}

\emph{প্রশ্নমতে, x + x + ১ + x + ২ = ২৭}

\emph{⇒ ৩x + ৩ = ২৭}

\emph{⇒ ৩x = ২৪}

\emph{⇒ x = ৮}

\emph{∴ শেষ ৩টি সংখ্যা হলো: x + ৩, x + ৪, x + ৫}

\emph{∴ শেষ ৩টি সংখ্যার যোগফল = (৮ + ৩) + (৮ + ৪) + (৮ + ৫) = ১১ + ১২ +
১৩ = ৩৬}

\emph{শর্টকাট পদ্ধতি:}

\emph{পরপর সংখ্যার ক্ষেত্রে প্রথম n-সংখ্যক সংখ্যার যোগফলের সাথে n\^{}২ যোগ
করলে পরবর্তী n-সংখ্যক সংখ্যার যোগফল পাওয়া যায়।}

\emph{এখানে, n = ৩}

\emph{∴ শেষ ৩টির যোগফল = ২৭ + ৩\^{}২ = ২৭ + ৯ = ৩৬}
\end{quote}

\textbf{দু'টি সংখ্যার বর্গের সমষ্টি ২৫০, সংখ্যা দু'টির গুণফল ১১৭। সংখ্যা দু'টি কি
কি ?}

ক) ১৪,৮

খ) ১৫, ৫

গ) ১২,৬

\textbf{\hl{ঘ) ১৩,৯ ✔}}

\textbf{\hl{✔ সঠিক উত্তর: ঘ) ১৩,৯}}

\begin{quote}
\textbf{ব্যাখ্যা:} \emph{অপশন টেস্ট / বিকল্প পদ্ধতি:}

\emph{(ঘ) অপশনের সংখ্যা দুটি হলো ১৩ ও ৯।}

\emph{১ম শর্তানুসারে, বর্গের সমষ্টি = ১৩\^{}২+৯\^{}২ = ১৬৯ + ৮১ = ২৫০ (শর্ত
পূরণ করে)}

\emph{২য় শর্তানুসারে, গুণফল = ১৩ × ৯ = ১১৭ (শর্ত পূরণ করে)}

\emph{∴ সংখ্যা দুটি হলো ১৩ ও ৯।}
\end{quote}

\textbf{একটি সংখ্যার বর্গ তার বর্গমূলের চেয়ে ৭৮ বেশি হলে, সংখ্যাটি কত?}

ক) ১২

খ) ৪

গ) ৬

\textbf{\hl{ঘ) ৯ ✔}}

\textbf{\hl{✔ সঠিক উত্তর: ঘ) ৯}}

\begin{quote}
\textbf{ব্যাখ্যা:} \emph{অপশন টেস্ট / বিকল্প পদ্ধতি:}

\emph{(ঘ) অপশনে প্রদত্ত সংখ্যাটি = ৯}

\emph{৯ এর বর্গ = ৯\^{}২ = ৮১}

\emph{৯ এর বর্গমূল = √(৯) = ৩}

\emph{বর্গ ও বর্গমূলের পার্থক্য = ৮১ − ৩ = ৭৮}

\emph{যা প্রদত্ত শর্ত পূরণ করে।}

\emph{∴ সংখ্যাটি ৯।}
\end{quote}

\textbf{দুটি সংখ্যার যোগফল ৫৫ এবং বড়টির ৫ গুণ ছোট সংখ্যাটির ৬ গুণের সমান।
সংখ্যা দুটি হবে -\/-}

ক) ২৫,২০

খ) ২৫,৩০

\textbf{\hl{গ) ৩০,২৫ ✔}}

ঘ) ৩০,২০

\textbf{\hl{✔ সঠিক উত্তর: গ) ৩০,২৫}}

\begin{quote}
\textbf{ব্যাখ্যা:} \emph{বীজগাণিতিক পদ্ধতি:}

\emph{মনে করি, বড় সংখ্যাটি x এবং ছোট সংখ্যাটি y}

\emph{১ম শর্তমতে, x + y = ৫৫ ⇒ y = ৫৫ − x}

\emph{২য় শর্তমতে, ৫x = ৬y}

\emph{⇒ ৫x = ৬(৫৫ − x)}

\emph{⇒ ৫x = ৩৩০ − ৬x}

\emph{⇒ ১১x = ৩৩০}

\emph{⇒ x = ৩০}

\emph{∴ ছোট সংখ্যাটি y = ৫৫ − ৩০ = ২৫}

\emph{∴ সংখ্যা দুটি যথাক্রমে ৩০ ও ২৫।}

\emph{বিকল্প পদ্ধতি (অপশন টেস্ট):}

\emph{(গ) অপশনে, সংখ্যা দুটি ৩০ ও ২৫।}

\emph{১ম শর্তানুসারে, যোগফল = ৩০ + ২৫ = ৫৫}

\emph{২য় শর্তানুসারে, বড়টির ৫ গুণ = ৩০ × ৫ = ১৫০ এবং ছোটটির ৬ গুণ = ২৫ × ৬ =
১৫০ (উভয় সমান)}

\emph{∴ সঠিক উত্তর (গ)।}
\end{quote}

\textbf{কোনো সংখ্যার ৪০\% এর সাথে ৪২ যোগ করলে ফলাফল হবে ঐ সংখ্যাটি ?}

ক) ৯০

খ) ৮০

গ) ৭৫

\textbf{\hl{ঘ) ৭০ ✔}}

\textbf{\hl{✔ সঠিক উত্তর: ঘ) ৭০}}

\begin{quote}
\textbf{ব্যাখ্যা:} \emph{ধরি, সংখ্যাটি = x}

\emph{প্রশ্নমতে, x এর ৪০\% + ৪২ = x}

\emph{⇒ (৪০x/১০০) + ৪২ = x}

\emph{⇒ x − (২x/৫) = ৪২}

\emph{⇒ (৩x/৫) = ৪২}

\emph{⇒ ৩x = ২১০}

\emph{⇒ x = ৭০}

\emph{∴ সংখ্যাটি = ৭০।}
\end{quote}

\textbf{তিনটি সংখ্যার গুণফল ৯৭২। ক্ষুদ্রতম সংখ্যাটি বৃহত্তম সংখ্যার এক-তৃতীয়াংশ
এবং মধ্যম সংখ্যাটি ক্ষুদ্রতম সংখ্যার (৩/২) গুণ। সংখ্যা তিনটির যোগফল কত?}

ক) ৩২

খ) ৩৬

গ) ৩৪

\textbf{\hl{ঘ) ৩৩ ✔}}

\textbf{\hl{✔ সঠিক উত্তর: ঘ) ৩৩}}

\begin{quote}
\textbf{ব্যাখ্যা:} \emph{মনে করি, বৃহত্তম সংখ্যাটি = x}

\emph{∴ ক্ষুদ্রতম সংখ্যাটি = (x/৩)}

\emph{এবং মধ্যম সংখ্যাটি = (x/৩) × (৩/২) = (x/২)}

\emph{প্রশ্নমতে, x × (x/৩) × (x/২) = ৯৭২}

\emph{⇒ (x\^{}৩/৬) = ৯৭২}

\emph{⇒ x\^{}৩ = ৫৮৩২}

\emph{⇒ x = ৩√(৫৮৩২) = ১৮}

\emph{∴ বৃহত্তম সংখ্যা = ১৮}

\emph{ক্ষুদ্রতম সংখ্যা = (১৮/৩) = ৬}

\emph{মধ্যম সংখ্যা = (১৮/২) = ৯}

\emph{∴ সংখ্যা তিনটির যোগফল = ১৮ + ৬ + ৯ = ৩৩।}
\end{quote}

\textbf{কোনো সংখ্যার তিন-অষ্টমাংশের সাথে এর দ্বিগুণের দুই-পঞ্চমাংশ যোগ করলে যা
হয় তা মূল সংখ্যাটি থেকে ২১ বেশি। সংখ্যাটির তিন পঞ্চমাংশ কত?}

ক) ১৫০

খ) ১২০

গ) ১০০

\textbf{\hl{ঘ) ৭২ ✔}}

\textbf{\hl{✔ সঠিক উত্তর: ঘ) ৭২}}

\begin{quote}
\textbf{ব্যাখ্যা:} \emph{মনে করি, সংখ্যাটি = x}

\emph{প্রশ্নমতে, (৩x/৮) + (২x × (২/৫) ) = x + ২১}

\emph{⇒ (৩x/৮) + (৪x/৫) = x + ২১}

\emph{⇒ (১৫x+৩২x/৪০) = x + ২১}

\emph{⇒ (৪৭x/৪০) − x = ২১}

\emph{⇒ (৭x/৪০) = ২১}

\emph{⇒ ৭x = ৮৪০}

\emph{⇒ x = ১২০}

\emph{∴ সংখ্যাটির তিন-পঞ্চমাংশ = ১২০ × (৩/৫) = ৭২।}
\end{quote}

\textbf{দুই অংকের কোন সংখ্যার অংকদ্বয়ের সমষ্টি ৫, সংখ্যাটির সাথে ৯ যোগ করলে
অংকদ্বয় স্থান পরিবর্তন করে। সংখ্যাটি কত ?}

ক) ৩২

খ) ৪১

\textbf{\hl{গ) ২৩ ✔}}

ঘ) ৫০

\textbf{\hl{✔ সঠিক উত্তর: গ) ২৩}}

\begin{quote}
\textbf{ব্যাখ্যা:} \emph{বীজগাণিতিক পদ্ধতি:}

\emph{মনে করি, একক স্থানীয় অংক = x এবং দশক স্থানীয় অংক = y}

\emph{∴ সংখ্যাটি = ১০y + x}

\emph{১ম শর্তমতে, x + y = ৫ ⇒ x = ৫ − y}

\emph{২য় শর্তমতে, (১০y + x) + ৯ = ১০x + y}

\emph{⇒ ১০y + x + ৯ − ১০x − y = ০}

\emph{⇒ ৯y − ৯x = −৯}

\emph{⇒ ৯x − ৯y = ৯}

\emph{⇒ x − y = ১}

\emph{⇒ (৫ − y) − y = ১}

\emph{⇒ ২y = ৪}

\emph{⇒ y = ২}

\emph{∴ x = ৫ − ২ = ৩}

\emph{∴ সংখ্যাটি = (১০ × ২) + ৩ = ২৩।}

\emph{বিকল্প পদ্ধতি (অপশন টেস্ট):}

\emph{(গ) অপশনে প্রদত্ত সংখ্যাটি = ২৩}

\emph{অংকদ্বয়ের সমষ্টি = ২ + ৩ = ৫ (১ম শর্ত পূরণ করে)}

\emph{সংখ্যাটির সাথে ৯ যোগ করলে হয় = ২৩ + ৯ = ৩২ (যা অঙ্কদ্বয়ের স্থান
পরিবর্তন করে ২৩ থেকে ৩২ হয়েছে)}

\emph{∴ সংখ্যাটি ২৩।}
\end{quote}

\textbf{একটি সংখ্যার ২০ শতাংশের ৮০ শতাংশ যদি ১২.৮ হয় সংখ্যাটি কত ?}

ক) ৯

খ) ৫০

গ) ৪০

\textbf{\hl{ঘ) ৮০ ✔}}

\textbf{\hl{✔ সঠিক উত্তর: ঘ) ৮০}}

\begin{quote}
\textbf{ব্যাখ্যা:} \emph{গাণিতিক সমাধান:}

\emph{মনে করি, সংখ্যাটি = x}

\emph{প্রশ্নমতে, x × ২০\% × ৮০\% = ১২.৮}

\emph{⇒ x × (২০/১০০) × (৮০/১০০) = ১২.৮}

\emph{⇒ x × (১/৫) × (৪/৫) = (১২৮/১০)}

\emph{⇒ (৪x/২৫) = (৬৪/৫)}

\emph{⇒ x = (৬৪×২৫/৫×৪)}

\emph{⇒ x = ১৬ × ৫}

\emph{∴ x = ৮০}

\emph{বিকল্প পদ্ধতি (অপশন টেস্ট):}

\emph{(ঘ) অপশনের সংখ্যাটি = ৮০}

\emph{৮০ এর ২০\% = ১৬}

\emph{আবার, ১৬ এর ৮০\% = (১৬×৮০/১০০) = ১২.৮}

\emph{∴ সংখ্যাটি ৮০।}
\end{quote}

\section{41. ঘড়ির কাঁটার অংক \textbar\textbar{} সর্বশেষ সহজে
নিয়মে}\label{ux998ux9a1ux9b0-ux995ux99fux9b0-ux985ux995-ux9b8ux9b0ux9acux9b6ux9b7-ux9b8ux9b9ux99c-ux9a8ux9dfux9ae}

\subsubsection{ঘড়ির কাঁটা}\label{ux998ux9dcux9b0-ux995ux99f}

ঘড়ির ঘণ্টার ও মিনিটের কাঁটার মধ্যবর্তী কোন নির্ণয়ের সুত্র:

ঘড়ির দুটি ঘণ্টার দাগের মধ্যবর্তী কোণ = \({\text{৩০}}^{\text{০}}\)

ঘড়ির কেন্দ্রে উৎপন্ন কোণ = \({\text{৩৬০}}^{\text{০}}\)

ঘড়ির ঘণ্টা ও মিনিটের কাটার মধ্যবর্তী কোণ

= \textbar{}\(\frac{11M\  - \ 60H}{2}\)\textbar{} ডিগ্রি

এখানে M = প্রশ্নে উল্লেখিত মিনিট

এবং H = প্রশ্নে উলেখিত ঘণ্টা

দ্রষ্টব্য : কোণের মান ( − ) ঋণাত্মক হলে, ( − ) বাদ দিতে হবে এবং কোন কোণের
মান \({\text{১৮০}}^{\text{০}}\) এর উপরে হলে তা \({\text{৩৬০}}^{\text{০}}\) হতে বিয়োগ করলে কোণের মান
পাওয়া যাবে ।

\textbf{প্রশ্ন নম্বর-} ১. ২ টা ১৫ মিনিটের সময় ঘণ্টার কাঁটা ও মিনিটের কাঁটা কত
ডিগ্রি কোণ উৎপন্ন করে।

\begin{itemize}
\item
  (ক) ২৩ ডিগ্রি
\item
  (খ) ২২ (১/২) ডিগ্রি
\item
  (গ) ২০ ডিগ্রি
\item
  (ঘ) ২৩ (১/২) ডিগ্রি
\end{itemize}

\textbf{\hl{উত্তরঃ (খ) ২২}}\(\frac{\text{১}}{\text{২}}\) \textbf{\hl{ডিগ্রি}}

\textbf{সমাধান:}

ঘড়ির ঘণ্টা ও মিনিটের কাটার মধ্যবর্তী কোণ

= \textbar{}\(\frac{11M\  - \ 60H}{2}\)\textbar{} ডিগ্রি

{[} এখানে M = প্রশ্নে উল্লেখিত মিনিট

এবং H = প্রশ্নে উলেখিত ঘণ্টা {]}

= \(\frac{\text{১১}\  \times \ \text{১৫}\  - \ \text{৬০}\  \times \ \text{২}}{\text{২}}\) ডিগ্রি

= \(\frac{\text{১৬৫}\  - \ \text{১২০}}{\text{২}}\) ডিগ্রি

= \(\frac{\text{৪৫}}{\text{২}}\) ডিগ্রি

= \({\text{২২}\ \text{১}\ \text{২}}^{\text{০}}\)

\textbf{প্রশ্ন নম্বর-} ২. ৪ টা ৫ মিনিটের সময় ঘণ্টার কাঁটা ও মিনিটের কাঁটা কত
ডিগ্রি কোণ উৎপন্ন করে।

\begin{itemize}
\item
  (ক) ৯২ (১/২) ডিগ্রি
\item
  (খ) ৯৩ (১/২) ডিগ্রি
\item
  (গ) ২০ ডিগ্রি
\item
  (ঘ) ৯২ (১/২) ডিগ্রি
\end{itemize}

\textbf{\hl{উত্তরঃ (ঘ) ৯২}}\(\frac{\text{১}}{\text{২}}\) \textbf{\hl{ডিগ্রি}}

\textbf{সমাধান:}

ঘড়ির ঘণ্টা ও মিনিটের কাটার মধ্যবর্তী কোণ

= \textbar{}\(\frac{11M\  - \ 60H}{2}\)\textbar{} ডিগ্রি

{[} এখানে M = প্রশ্নে উল্লেখিত মিনিট

এবং H = প্রশ্নে উলেখিত ঘণ্টা {]}

= \(\frac{\text{১১}\  \times \ \text{৫}\  - \ \text{৬০}\  \times \ \text{৪}}{\text{২}}\) ডিগ্রি

= \(\frac{\text{৫৫}\  - \ \text{২৪০}}{\text{২}}\) ডিগ্রি

= \(\frac{- \ \text{১৮৫}}{\text{২}}\) ডিগ্রি

= − ৯২\(\frac{\text{১}}{\text{২}}\) ডিগ্রি

দ্রষ্টব্য : কোণের মান ( − ) ঋণাত্মক হলে, ( − ) বাদ দিতে হবে এবং কোন কোণের
মান ১৮০ ডিগ্রি এর উপরে হলে তা ৩৬০ ডিগ্রি হতে বিয়োগ করলে কোণের মান পাওয়া
যাবে ।

\textbf{প্রশ্ন নম্বর-} ৩. ৬ টা ১২ মিনিটের সময় ঘণ্টার কাঁটা ও মিনিটের কাঁটা কত
ডিগ্রি কোণ উৎপন্ন করে।

\begin{itemize}
\item
  (ক) ১০০ ডিগ্রি
\item
  (খ) ১১০ ডিগ্রি
\item
  (গ) ১১৪ ডিগ্রি
\item
  (ঘ) ১১৫ ডিগ্রি
\end{itemize}

\textbf{\hl{উত্তরঃ (গ) ১১৪ ডিগ্রি}}

\textbf{সমাধান:}

ঘড়ির ঘণ্টা ও মিনিটের কাটার মধ্যবর্তী কোণ

= \textbar{}\(\frac{11M\  - \ 60H}{2}\)\textbar{} ডিগ্রি

{[} এখানে M = প্রশ্নে উল্লেখিত মিনিট

এবং H = প্রশ্নে উলেখিত ঘণ্টা {]}

= \(\frac{\text{১১}\  \times \ \text{১২}\  - \ \text{৬০}\  \times \ \text{৬}}{\text{২}}\) ডিগ্রি

= \(\frac{\text{১৩২}\  - \ \text{৩৬০}}{\text{২}}\) ডিগ্রি

= \(\frac{- \ \text{২২৮}}{\text{২}}\) ডিগ্রি

= − ১১৪ ডিগ্রি

দ্রষ্টব্য : কোণের মান ( − ) খণাত্মক হলে, ( − ) বাদ দিতে হবে এবং কোন কোণের
মান ১৮০ ডিগ্রি এর উপরে হলে তা ৩৬০ ডিগ্রি হতে বিয়োগ করলে কোণের মান পাওয়া
যাবে ।

\section{42. ঘড়ির কাঁটার অংক MCQ প্রশ্ন উত্তর
অনুশীলন}\label{ux998ux9dcux9b0-ux995ux99fux9b0-ux985ux995-mcq-ux9aaux9b0ux9b6ux9a8-ux989ux9a4ux9a4ux9b0-ux985ux9a8ux9b6ux9b2ux9a8}

\textbf{ঘড়িতে যখন ৮ টা বাজে তখন ঘণ্টার কাঁটা ও মিনিটের কাঁটার মধ্যবর্তী কোণ
কত ডিগ্রি হবে?}

ক) ৯০°

খ) ৯৫°

গ) ১০৫°

\textbf{\hl{ঘ) ১২০° ✔}}

\textbf{\hl{✔ সঠিক উত্তর: ঘ) ১২০°}}

\begin{quote}
\textbf{ব্যাখ্যা:} \emph{ঘণ্টা ও মিনিটের কাঁটার মধ্যবর্তী কোণ নির্ণয়ের সূত্র =
(11M−60H/2)}

\emph{এখানে, M=0 , H=8}

\emph{⇒ মধ্যবর্তী কোণ = (11×0−60×8/2)=(−480/2)=240°}

\emph{প্রবৃদ্ধ কোণ হওয়ায় নির্ণেয় কোণ = 360°−240°=120°}
\end{quote}

\textbf{একটি ঘড়িতে যখন সকাল ১০টা ১২ মিনিট তখন ঘণ্টা ও মিনিটের কাঁটার মধ্যে
কত ডিগ্রি কোণ উৎপন্ন হয়?}

ক) ৯৬

খ) ১০২

গ) ১২৪

\textbf{\hl{ঘ) ১২৬ ✔}}

\textbf{\hl{✔ সঠিক উত্তর: ঘ) ১২৬}}

\begin{quote}
\textbf{ব্যাখ্যা:} \emph{এখানে, M=12 , H=10}

\emph{মধ্যবর্তী কোণ = ((11×12)−(60×10)/2)°}

\emph{⇒ (132−600/2)°=(−468/2)°=234°}

\emph{নির্ণেয় কোণ = 360°−234°=126°}
\end{quote}

\textbf{২টা ১৫ মিনিটের সময় ঘণ্টার কাঁটা ও মিনিটের কাঁটার মধ্যবর্তী কোণ কত?}

ক) ২০°

খ) ২৩°

\textbf{\hl{গ) ২২ (১/২) ° ✔}}

ঘ) ২৩ (১/২) °

\textbf{\hl{✔ সঠিক উত্তর: গ) ২২ (১/২) °}}

\begin{quote}
\textbf{ব্যাখ্যা:} \emph{এখানে, M=15 , H=2}

\emph{মধ্যবর্তী কোণ = (11/2)×15−30×2°=22(1/2)°}
\end{quote}

\textbf{An accurate clock shows 8 o\textquotesingle clock in the
morning. Through how many degrees will the hour hand rotate when the
clock shows 2 o\textquotesingle clock in the afternoon?}

ক) 144°

খ) 150°

গ) 168°

\textbf{\hl{ঘ) 180° ✔}}

\textbf{\hl{✔ সঠিক উত্তর: ঘ) 180°}}

\begin{quote}
\textbf{ব্যাখ্যা:} \emph{Time difference between 8 o\textquotesingle clock
and 2 o\textquotesingle clock is 6 hours.}

\emph{We know, hour hand rotates 30° in 1 hour.}

\emph{∴ Total rotation in 6 hours = 6×30°=180°}
\end{quote}

\textbf{বেলা ২.৩০ ঘটিকার সময় ঘড়িতে ঘণ্টা ও মিনিটের কাঁটা পরস্পর কত ডিগ্রি
কোণ উৎপন্ন করবে?}

ক) ৯০°

\textbf{\hl{খ) ১০৫.৫° ✔}}

গ) ৬০°

ঘ) ১২০°

\textbf{\hl{✔ সঠিক উত্তর: খ) ১০৫.৫°}}

\begin{quote}
\textbf{ব্যাখ্যা:} \emph{এখানে, মিনিট M=30 , ঘণ্টা H=2}

\emph{মধ্যবর্তী কোণ = (11×30−60×2/2)°}

\emph{⇒ (330−120/2)°=(210/2)°=105° (নোট: সঠিক উত্তর ১০৫°)}
\end{quote}

\textbf{ঘড়িতে এখন ৮টা বাজে। ঘণ্টার কাঁটা ও মিনিটের কাঁটার মধ্যকার কোণটি হবে?
/ Now it is 8 \textquotesingle O\textquotesingle{} clock by the watch.
Find the angle formed by the hands of hour and minute.}

ক) ১৫০°

খ) ৬০°

\textbf{\hl{গ) ১২০° ✔}}

ঘ) ৯০°

\textbf{\hl{✔ সঠিক উত্তর: গ) ১২০°}}

\begin{quote}
\textbf{ব্যাখ্যা:} \emph{ঘড়িতে ১২ ঘণ্টার জন্য সম্পূর্ণ কোণ = 360°}

\emph{⇒ ১ ঘণ্টার ব্যবধানে কোণ = (360°/12)=30°}

\emph{যখন ৮টা বাজে, তখন মিনিটের কাঁটা ১২ এর ঘরে এবং ঘণ্টার কাঁটা ৮ এর ঘরে
থাকে।}

\emph{∴ মধ্যবর্তী কোণ = 30°+30°+30°+30°=120°}
\end{quote}

\textbf{একটি ঘড়ি দুপুর ১টার সময় একবার বাজে, ২টার সময় দুইবার বাজে, ৩টার সময়
তিনবার বাজে এবং এই নিয়মে বাজতে থাকে। ঘড়িটি ২৪ ঘণ্টায় মোট কতবার বাজবে?}

ক) ৭৮

খ) ১৩৬

\textbf{\hl{গ) ১৫৬ ✔}}

ঘ) ১৯৬

\textbf{\hl{✔ সঠিক উত্তর: গ) ১৫৬}}

\begin{quote}
\textbf{ব্যাখ্যা:} \emph{ঘড়িটি প্রথম ১২ ঘণ্টায় বাজে = (1+2+3+\ldots+12) বার}

\emph{⇒ (12(12+1)/2)=78 বার}

\emph{অনুরূপভাবে পরবর্তী ১২ ঘণ্টায় বাজে 78 বার}

\emph{∴ ঘড়িটি ২৪ ঘণ্টায় মোট বাজবে = 78+78=156 বার।}
\end{quote}

\textbf{৪ টা ৫ মিনিটের সময় ঘণ্টার কাঁটা ও মিনিটের কাঁটা কত ডিগ্রি কোণ উৎপন্ন
করে।}

ক) 2(5/2)°

খ) 7(7/2)°

গ) 20°

\textbf{\hl{ঘ) 92(1/2)° ✔}}

\textbf{\hl{✔ সঠিক উত্তর: ঘ) 92(1/2)°}}

\begin{quote}
\textbf{ব্যাখ্যা:} \emph{এখানে, M=5 , H=4}

\emph{মধ্যবর্তী কোণ =
\textbar(11×5−60×4/2)\textbar°=\textbar(55−240/2)\textbar°=(185/2)°=92(1/2)°}
\end{quote}

\textbf{৬ টা ১২ মিনিটের সময় ঘণ্টার কাঁটা ও মিনিটের কাঁটা কত ডিগ্রি কোণ উৎপন্ন
করে।}

ক) ১০০ ডিগ্রি

খ) ১১০ ডিগ্রি

\textbf{\hl{গ) ১১৪ ডিগ্রি ✔}}

ঘ) ১১৫ ডিগ্রি

\textbf{\hl{✔ সঠিক উত্তর: গ) ১১৪ ডিগ্রি}}

\begin{quote}
\textbf{ব্যাখ্যা:} \emph{এখানে, M=12 , H=6}

\emph{মধ্যবর্তী কোণ =
\textbar(11×12−60×6/2)\textbar°=\textbar(132−360/2)\textbar°=(228/2)°=114°}
\end{quote}

\textbf{২ টা ৫০ মিনিটের সময় ঘণ্টার কাঁটা ও মিনিটের কাঁটা কত ডিগ্রি কোণ উৎপন্ন
করে?}

\textbf{\hl{ক) ১৪৫ ডিগ্রি ✔}}

খ) ১৩ ডিগ্রি

গ) ১২৫ ডিগ্রি

ঘ) ১৩০ ডিগ্রি

\textbf{\hl{✔ সঠিক উত্তর: ক) ১৪৫ ডিগ্রি}}

\begin{quote}
\textbf{ব্যাখ্যা:} \emph{এখানে, M=50 , H=2}

\emph{মধ্যবর্তী কোণ =
\textbar(11×50−60×2/2)\textbar°=\textbar(550−120/2)\textbar°=(430/2)°=215°}

\emph{১৮০ ডিগ্রির বেশি হওয়ায় নির্ণেয় কোণ = 360°−215°=145°}
\end{quote}

\textbf{৪ টার সময়ে ঘণ্টার কাঁটা ও মিনিটের কাঁটা কত ডিগ্রি কোণ উৎপন্ন করে।}

ক) ১৫০ ডিগ্রি

খ) ৬০ ডিগ্রি

\textbf{\hl{গ) ১২০ ডিগ্রি ✔}}

ঘ) ৯০ ডিগ্রি

\textbf{\hl{✔ সঠিক উত্তর: গ) ১২০ ডিগ্রি}}

\begin{quote}
\textbf{ব্যাখ্যা:} \emph{এখানে, M=0 , H=4}

\emph{মধ্যবর্তী কোণ = \textbar(11×0−60×4/2)\textbar°=(−240/2)°=120°}
\end{quote}

\textbf{৭ টার সময়ে ঘণ্টার কাঁটা ও মিনিটের কাঁটা কত ডিগ্রি কোণ উৎপন্ন করে।}

ক) ১৩০ ডিগ্রি

\textbf{\hl{খ) ১৫০ ডিগ্রি ✔}}

গ) ১৬০ ডিগ্রি

ঘ) ১৪০ ডিগ্রি

\textbf{\hl{✔ সঠিক উত্তর: খ) ১৫০ ডিগ্রি}}

\begin{quote}
\textbf{ব্যাখ্যা:} \emph{এখানে, M=0 , H=7}

\emph{মধ্যবর্তী কোণ = \textbar(11×0−60×7/2)\textbar°=(−420/2)°=210°}

\emph{১৮০ ডিগ্রির বেশি হওয়ায় নির্ণেয় কোণ = 360°−210°=150°}
\end{quote}

\textbf{দেয়াল ঘড়িতে রাত ৯টা বাজে; ঘণ্টা ও মিনিটের কাঁটার মধ্যকার কোণটির
পরিমাপ কত?}

ক) ৬০ ডিগ্রি

খ) ৭০ ডিগ্রি

গ) ৮০ ডিগ্রি

\textbf{\hl{ঘ) ৯০ ডিগ্রি ✔}}

\textbf{\hl{✔ সঠিক উত্তর: ঘ) ৯০ ডিগ্রি}}

\begin{quote}
\textbf{ব্যাখ্যা:} \emph{এখানে, M=0 , H=9}

\emph{মধ্যবর্তী কোণ = \textbar(11×0−60×9/2)\textbar°=(−540/2)°=270°}

\emph{১৮০ ডিগ্রির বেশি হওয়ায় নির্ণেয় কোণ = 360°−270°=90°}
\end{quote}

\textbf{সকাল ৯.৩০ মিনিটে ঘড়ির ঘণ্টার কাঁটা ও মিনিটের কাঁটার অন্তর্গত কোণকে
ডিগ্রীতে প্রকাশ করলে কত ডিগ্রী হবে?}

ক) ৬৫ ডিগ্রি

খ) ৭৫ ডিগ্রি

গ) ৮৫ ডিগ্রি

\textbf{\hl{ঘ) ১০৫ ডিগ্রি ✔}}

\textbf{\hl{✔ সঠিক উত্তর: ঘ) ১০৫ ডিগ্রি}}

\begin{quote}
\textbf{ব্যাখ্যা:} \emph{এখানে, M=30 , H=9}

\emph{মধ্যবর্তী কোণ =
\textbar(11×30−60×9/2)\textbar°=\textbar(330−540/2)\textbar°=(210/2)°=105°}
\end{quote}

\textbf{একটি ঘড়ি দুপুর ১২ টা হতে চলতে শুরু করেছে। ৫টা ১০ মিনিটে ঘন্টার কাঁটাটি
কত ডিগ্রিতে ঘুরবে?}

ক) ১৪৫ ডিগ্রি

খ) ১৫০ ডিগ্রি

\textbf{\hl{গ) ১৫৫ ডিগ্রি ✔}}

ঘ) ১৬০ ডিগ্রি

\textbf{\hl{✔ সঠিক উত্তর: গ) ১৫৫ ডিগ্রি}}

\begin{quote}
\textbf{ব্যাখ্যা:} \emph{ঘণ্টার কাঁটা ১ ঘণ্টায় ঘোরে 30° , সুতরাং ৫ ঘণ্টায় ঘোরে
5×30°=150°}

\emph{ঘণ্টার কাঁটা ১ মিনিটে ঘোরে (1/2)° , সুতরাং ১০ মিনিটে ঘোরে
10×(1/2)°=5°}

\emph{∴ মোট ঘূর্ণন = 150°+5°=155°}
\end{quote}

\section{43. পরিসংখ্যান
(Statistics)-গণিত}\label{ux9aaux9b0ux9b8ux996ux9afux9a8-statistics-ux997ux9a3ux9a4}

পরিসংখ্যান এক ধরনের গাণিতিক বিজ্ঞান যা মূলত: উপাত্ত সংগ্রহ, বিশ্লেষণ, ব্যাখ্যা
এবং উপাত্ত সহজে পরিবেশন নিয়ে কাজ করে। উপাত্ত বিশ্লেষণ করে তা থেকে তথ্যসমৃদ্ধ
সিদ্ধান্ত গ্রহণে পরিসংখ্যানের ভূমিকা অপরিহার্য। বিভিন্ন ধরনের গবেষনায় পরিসংখ্যান
ব্যবহার করা হয়। গড়, মধ্যক, প্রচুরক, আয়তলেখ, গণসংখ্যা বহুভূজ, অজিভরেখা
পরিসংখ্যানের বিষয়বস্তু।

\textbf{কেন্দ্রীয় প্রবণতা (Central Tendency):}

অবিন্যস্ত উপাত্তসমূহকে মানের ক্রমানুসারে সাজালে উপাত্ত সমূহ মাঝামাঝি কোনো মান
(কেন্দ্রীয় মান)-এর কাছাকাছি পুঞ্জিভূত হয়। আবার এগুলোকে গণসংখ্যা সারণিতে উপস্থাপন
করলে মাঝামাঝি একটি শ্রেণিতে গণসংখ্যার প্রাচুর্য লক্ষ্ করা যায়। উপাত্তসমূহের
কেন্দ্রীয় মানের দিকে পুঞ্জিভূত হওয়ার এই প্রবণতাকে কেন্দ্রীয় প্রবণতা বলে।

\textbf{কেন্দ্রীয় প্রবণতার পরিমাপ :}

সাধারণভাবে কেন্দ্রীয় প্রবণতার পরিমাপ হল:

১। গড় বা গাণিতিক গড়

২। মধ্যক

৩। প্রচুরক

\textbf{গড় বা গাণিতিক গড় (Mean):}

গড় হল কোন একটি সংখ্যা তালিকার সকল মানকে প্রতিনিধিত্বকারী একটি একক মান।
যেমন: বাংলাদেশের মানুষের আয়ুর গড় ৭১ বছর।

গাণিতিক গড় নির্ণয়ের পদ্ধতি: কোনো সংখ্যা তালিকার সংখ্যাগুলোর সমষ্টিকে উপাত্ত
সংখ্যা দ্বারা ভাগ করে গড় বা গাণিতিক গড় নির্ণয় করা হয়। যেমন:

কোনো একটি বাসায় ৫ জন লোক থাকে যাদের উচ্চতা (ফুট) ৪, ৪.৫, ৩, ৪, ৫.৫ ।

সুতরাং, তাদের উচ্চতার গড় = \(\frac{\text{উপাত্তগুলোর}\ \text{সমষ্টি}}{\text{পদসংখ্যা}}\)

= \(\frac{\text{৪} + \text{৪}.\text{৫} + \text{৩} + \text{৪} + \text{৫}.\text{৫}}{\text{৫}}\)

= \(\frac{\text{২১}}{\text{৫}}\)

= ৪.২ তবে বেশি সংখ্যক উপাত্তের গড় নির্ণয় করার জন্য এই পদ্ধতিটি সুবিধাজনক নয়।
এক্ষেত্রে উপাত্তসমূহকে শ্রেণি বিন্যাসের মাধ্যমে সারণিবদ্ধ করে সংক্ষিপ্ত পদ্ধতিতে গড়
নির্ণয় করা হয়।

\textbf{মধ্যক (Median):}

উপাত্তসমূহকে মানের ক্রমানুসারে সাজালে মাঝখানে যে মানটি পাওয়া যায় তা-ই হল
মধ্যক। মাঝে একাধিক মান থাকলে তাদের গড় মান হল মধ্যক। উল্লেখ্য, মধ্যকের বামপাশে
থাকে অপেক্ষাকৃত ছোট সংখ্যাসমূহ আর ডানপাশে থাকে অপেক্ষাকৃত বড় সংখ্যাগুলো। দুই
ভাগেই সমান সংখ্যক উপাত্ত থাকে। যেমন:

৩৪, ৪৫, ৫৩, ৬৪, ৭৫, ৭৫, ৮৫ এর মধ্যক ৬৪ (ক্রমানুসারে সাজানো সংখ্যা তালিকার
মাঝের সংখ্যা)।

আবার ১৪, ২৫, ৩৩, ৪৫, ৬৫, ৬৫, ৮৫, ৮৫ এর মধ্যক ৫৫ (ক্রমানুসারে সাজানো সংখ্যা
তালিকার মাঝের দু'টি সংখ্যার গড় মান) ।

এছাড়াও মধ্যক নির্ণয়ের আরো পদ্ধতি রয়েছে। নিচের সূত্রগুলো মধ্যক নির্ণয়ের জন্য ব্যবহৃত
হয়:

১। উপাত্ত সংখ্যা বিজোড় হলে,

মধ্যক = \(\frac{n\  + \ \text{১}}{\text{২}}\) তম পদ {[} এখানে n উপাত্তের সংখ্যা {]}

\textbf{প্রশ্ন: ৩৪, ৪৫, ৫৩, ৬৪, ৭৫, ৭৫, ৮৫ এর মধ্যক কত?}

\textbf{সমাধান :}

মধ্যক = \(\frac{n\  + \ \text{১}}{\text{২}}\) তম পদ {[} এখানে উপাত্তের সংখ্যা n = ৭{]}

= \(\frac{\text{৭}\  + \ \text{১}}{\text{২}}\) তম পদ

= \(\frac{\text{৮}}{\text{২}}\) তম পদ

= ৪ তম পদ

= ৬৪ {[} এখানে ৪ তম পদ হচ্ছে ৬৪{]}

২। উপাত্ত সংখ্যা জোড় হলে,

মধ্যক = \{ (\(\frac{n}{\text{২}}\) ) তম পদ +( \(\frac{n}{\text{২}}\) +১) তম পদ\}÷২ {[}
এখানে n উপাত্তের সংখ্যা {]}

\textbf{প্রশ্ন: ১৪, ২৫, ৩৩, ৪৫, ৬৫, ৬৫, ৮৫, ৮৫ এর মধ্যক কত?}

\textbf{সমাধান:}

মধ্যক = \{ (\(\frac{n}{\text{২}}\) ) তম পদ +( \(\frac{n}{\text{২}}\) +১) তম পদ\}÷২ {[}
এখানে n উপাত্তের সংখ্যা {]}

= \{ (\(\frac{\text{৮}}{\text{২}}\) ) তম পদ +( \(\frac{\text{৮}}{\text{২}}\) +১) তম পদ\}÷২ {[} এখানে
n = ৮ উপাত্তের সংখ্যা {]}

= (৪ তম পদ + ৫ তম পদ) ÷২

= (৪৫+ ৬৫) ÷ ২ = ১১০÷২ = ৫৫ \textbf{প্রচুরক (Mode):}

কোনো উপাত্তে যে সংখ্যাটি সবচেয়ে বেশি বার উপস্থাপিত হয় তাকে উক্ত উপাত্তের
প্রচুরক বলে। প্রচুরক একাধিকও হতে পারে। কোনো উপাত্তে যদি একটি সংখ্যাও একাধিকবার
উপস্থাপিত না হয় তাহলে ঐ উপাত্তের প্রচুরক নেই। যেমন:

৩৪, ৪৫, ৫৩, ৬৪, ৭৫, ৭৫, ৮৫ এর প্রচুরক ৭৫

১৪, ২৫, ৩৩, ৪৫, ৬৫, ৬৫, ৮৫, ৮৫ এর প্রচুরক ৬৫ ও ৮৫।

আবার ৩৪, ৪৫, ৫৩, ৬৪, ৭৫, ৮৫ এর প্রচুরক নেই।

\section{44. পরিসংখ্যান MCQ প্রশ্ন উত্তর
অনুশীলন}\label{ux9aaux9b0ux9b8ux996ux9afux9a8-mcq-ux9aaux9b0ux9b6ux9a8-ux989ux9a4ux9a4ux9b0-ux985ux9a8ux9b6ux9b2ux9a8}

\textbf{উপাত্তসমূহের সর্বোচ্চ মান এবং সর্বনিম্ন মানের পার্থক্য কোনটি?}

ক) প্রচুরক

খ) মধ্যমান

\textbf{\hl{গ) পরিসর ✔}}

ঘ) শ্রেণি ব্যবধান

\textbf{\hl{✔ সঠিক উত্তর: গ) পরিসর}}

\begin{quote}
\textbf{ব্যাখ্যা:} \emph{পরিসর = উপাত্তসমূহের (সর্বোচ্চ মান − সর্বনিম্ন মান)}
\end{quote}

\textbf{নিচে ৪০ জন ছাত্রের বার্ষিক পরীক্ষায় গণিত বিষয়ে প্রাপ্ত নম্বর দেয়া হলঃ
৪২, ৩১, ৪৫, ২৭, ৬০, ৬১, ৩৯, ৪১, ৩৫, ৫৮, ২৯, ৫৩, ৪৮, ৩৯, ৫২, ৩৮, ৮০, ৪৭,
২৮, ৫১, ৪৯, ৭৮, ৯০, ৫২, ৪৮, ৩৬, ৫২, ৩৯, ৭১, ৬৪, ৩২, ৪৯, ৫৬, ৩৩, ৪
সবাইকে, ৩৩, ২৫, ৪৮, ২৯ - উপাত্তগুলোর প্রচুরক}

ক) ৪৭

খ) ৪৬

\textbf{\hl{গ) ৪৮ ✔}}

ঘ) ৪৭.৫

\textbf{\hl{✔ সঠিক উত্তর: গ) ৪৮}}

\begin{quote}
\textbf{ব্যাখ্যা:} \emph{উপাত্তগুলোর মধ্যে ৪৮ মানটি সর্বাধিক চারবার আছে।}

\emph{∴ প্রচুরক = ৪৮}
\end{quote}

\textbf{প্রদত্ত উপাত্তগুলোর মধ্যকঃ ১২, ৯, ১৫, ৫, ২০, ৮, ২৫, ১৭, ২১, ২৩, ১১}

ক) ১৩

খ) ১৪

গ) ১২

\textbf{\hl{ঘ) ১৫ ✔}}

\textbf{\hl{✔ সঠিক উত্তর: ঘ) ১৫}}

\begin{quote}
\textbf{ব্যাখ্যা:} \emph{উপাত্তগুলোকে মানের ক্রমানুসারে সাজালে পাই:}

\emph{5, 8, 9, 11, 12, 15, 17, 20, 21, 23, 25 (সর্বমোট ১১টি পদ)}

\emph{∴ নির্ণেয় মধ্যক = ৬ষ্ঠ পদ = 15}
\end{quote}

\textbf{নিচে প্রদত্ত সংখ্যাগুলোর গড়, মধ্যক ও প্রচুরক নির্ণয় করুনঃ ৩০, ১২, ২২, ১৭,
২৭, ২৫, ২০, ২৪, ১৯, ২, ২৩, ৩২, ২৬, ২৯, ৩৫, ২১, ১১, ২৮ এবং ১৯}

ক) গড় ২৩.২১ মধ্যক ২৩, প্রচুরক ১৯.২৭

\textbf{\hl{খ) গড় ২২.২১ মধ্যক ২৩, প্রচুরক ১৯ ✔}}

গ) গড় ২৩. ২৭ মধ্যক ২৩.৩৭, প্রচুরক ২০

ঘ) গড় ২৩, মধ্যক ২২, প্রচুরক ২৩

\textbf{\hl{✔ সঠিক উত্তর: খ) গড় ২২.২১ মধ্যক ২৩, প্রচুরক ১৯}}

\begin{quote}
\textbf{ব্যাখ্যা:} \emph{উপাত্তগুলোকে মানের ক্রমানুসারে সাজালে পাই:}

\emph{2, 11, 12, 17, 19, 19, 20, 21, 22, 23, 24, 25, 26, 27, 28, 29, 30,
32, 35 (এখানে মোট ১৯টি উপাত্ত রয়েছে)}

\emph{যোগফল = ৪৪১, সুতরাং গড় = (441/19) = 23.21}

\emph{মধ্যক = (১৯ + ১) / ২ তম পদ = ১০ম পদ = 23}

\emph{সবচেয়ে বেশি বার আছে 19 (দুইবার), সুতরাং প্রচুরক = 19}
\end{quote}

\textbf{কোন উপাত্তগুলোর মান যদি 10, 15, 20, 25, 30 হয়, তবে এদের সমান্তরাল
গড় কত?}

ক) ১৫

\textbf{\hl{খ) ২০ ✔}}

গ) ২৫

ঘ) ১৮

\textbf{\hl{✔ সঠিক উত্তর: খ) ২০}}

\begin{quote}
\textbf{ব্যাখ্যা:} \emph{উপাত্তগুলোর যোগফল = 10+15+20+25+30=100}

\emph{পদসংখ্যা = 5}

\emph{∴ গড় = (100/5)=20}
\end{quote}

\textbf{প্রথম ১০টি স্বাভাবিক সংখ্যার গড় কত?}

ক) ৫

\textbf{\hl{খ) ৫.৫ ✔}}

গ) ৬

ঘ) ৬.৫

\textbf{\hl{✔ সঠিক উত্তর: খ) ৫.৫}}

\begin{quote}
\textbf{ব্যাখ্যা:} \emph{প্রথম n সংখ্যক স্বাভাবিক সংখ্যার সমষ্টি = (n(n+1)/2)}

\emph{প্রথম ১০টি স্বাভাবিক সংখ্যার সমষ্টি = (10×11/2)=55}

\emph{∴ গড় = (55/10)=5.5}
\end{quote}

\textbf{যদি ৫টি সংখ্যার যোগফল ৫০ হয়, তবে সংখ্যাগুলোর গড় কত?}

ক) ৫

\textbf{\hl{খ) ১০ ✔}}

গ) ১৫

ঘ) ২০

\textbf{\hl{✔ সঠিক উত্তর: খ) ১০}}

\begin{quote}
\textbf{ব্যাখ্যা:} \emph{গড় = (সংখ্যার যোগফল/মোট সংখ্যা)}

\emph{∴ গড় = (50/5)=10}
\end{quote}

\textbf{একটি গণসংখ্যা নিবেশনের মধ্যক শ্রেণির উচ্চসীমা ৩০ এবং প্রচুরক শ্রেণির
নিম্নসীমা ২০ হলে, নিবেশনটির পরিসংখ্যানগত মান নির্ণয়ে নিচের কোনটি প্রযোজ্য?}

ক) উপাত্তের বিস্তার নির্ধারণে

\textbf{\hl{খ) কেন্দ্রীয় প্রবণতা পরিমাপে ✔}}

গ) বিচ্যুতি নির্ণয়ে

ঘ) সহগ নির্ণয়ে

\textbf{\hl{✔ সঠিক উত্তর: খ) কেন্দ্রীয় প্রবণতা পরিমাপে}}

\begin{quote}
\textbf{ব্যাখ্যা:} \emph{গড়, মধ্যক ও প্রচুরক হলো কেন্দ্রীয় প্রবণতার পরিমাপক।}
\end{quote}

\textbf{কোনো উপাত্তের সর্বোচ্চ মান ৭০ এবং পরিসর ৩০ হলে, উপাত্তটির সর্বনিম্ন
মান কত?}

ক) ৩০

\textbf{\hl{খ) ৪০ ✔}}

গ) ৫০

ঘ) ৬০

\textbf{\hl{✔ সঠিক উত্তর: খ) ৪০}}

\begin{quote}
\textbf{ব্যাখ্যা:} \emph{আমরা জানি, পরিসর = সর্বোচ্চ মান − সর্বনিম্ন মান}

\emph{⇒ 30=70−সর্বনিম্ন মান}

\emph{∴ সর্বনিম্ন মান = 70−30=40}
\end{quote}

\textbf{প্রথম ১০টি বিজোড় স্বাভাবিক সংখ্যার গড় কত?}

ক) ৯

\textbf{\hl{খ) ১০ ✔}}

গ) ১১

ঘ) ১২

\textbf{\hl{✔ সঠিক উত্তর: খ) ১০}}

\begin{quote}
\textbf{ব্যাখ্যা:} \emph{প্রথম n সংখ্যক বিজোড় স্বাভাবিক সংখ্যার যোগফল =
n\^{}2}

\emph{প্রথম ১০টি বিজোড় স্বাভাবিক সংখ্যার যোগফল = 10\^{}2=100}

\emph{∴ গড় = (100/10)=10}
\end{quote}

\textbf{প্রথম ৯টি জোড় স্বাভাবিক সংখ্যার গড় কত?}

ক) ৯

\textbf{\hl{খ) ১০ ✔}}

গ) ১১

ঘ) ১২

\textbf{\hl{✔ সঠিক উত্তর: খ) ১০}}

\begin{quote}
\textbf{ব্যাখ্যা:} \emph{প্রথম n সংখ্যক জোড় স্বাভাবিক সংখ্যার গড় = n+1}

\emph{এখানে, n=9}

\emph{∴ গড় = 9+1=10}
\end{quote}

\textbf{৫টি সংখ্যার গড় ১০। এর সাথে আর একটি সংখ্যা যোগ করায় সংখ্যাগুলোর গড় ১২
হলো। নতুন সংখ্যাটি কত?}

ক) ১২

খ) ১৫

গ) ১৮

\textbf{\hl{ঘ) ২২ ✔}}

\textbf{\hl{✔ সঠিক উত্তর: ঘ) ২২}}

\begin{quote}
\textbf{ব্যাখ্যা:} \emph{৫টি সংখ্যার যোগফল = 5×10=50}

\emph{সংখ্যা একটি বাড়ায় মোট সংখ্যা হলো 6 টি এবং নতুন গড় 12}

\emph{৬টি সংখ্যার যোগফল = 6×12=72}

\emph{∴ নতুন সংখ্যাটি = 72−50=22}
\end{quote}

\textbf{একটি উপাত্তের মধ্যক ৭ এবং গড় ৮ হলে, এর প্রচুরক কত হতে পারে?}

\textbf{\hl{ক) ৫ ✔}}

খ) ৬

গ) ৭

ঘ) ৮

\textbf{\hl{✔ সঠিক উত্তর: ক) ৫}}

\begin{quote}
\textbf{ব্যাখ্যা:} \emph{পিয়ারসনের সূত্রানুসারে, গড় − প্রচুরক = 3×(গড়−মধ্যক)}

\emph{⇒ 8−প্রচুরক=3×(8−7)}

\emph{⇒ 8−প্রচুরক=3}

\emph{∴ প্রচুরক = 8−3=5}
\end{quote}

\textbf{তথ্য উপাত্তের ক্ষেত্রে প্রচুরক হলো সেটি, যা---}

ক) সবচেয়ে কমবার থাকে

\textbf{\hl{খ) সবচেয়ে বেশিবার থাকে ✔}}

গ) মাঝের মান নির্দেশ করে

ঘ) সর্বোচ্চ ও সর্বনিম্ন মানের পার্থক্য

\textbf{\hl{✔ সঠিক উত্তর: খ) সবচেয়ে বেশিবার থাকে}}

\begin{quote}
\textbf{ব্যাখ্যা:} \emph{প্রচুরক (Mode) হলো কোনো উপাত্তে যে মানটি সর্বাধিক বা
সবচেয়ে বেশিবার উপস্থিত থাকে।}
\end{quote}

\textbf{কোনো উপাত্তের সর্বনিম্ন মান ২৫ এবং পরিসর ৬০ হলে, উপাত্তটির সর্বোচ্চ
মান কত?}

ক) ৩৫

খ) ৭০

\textbf{\hl{গ) ৮৫ ✔}}

ঘ) ৯৫

\textbf{\hl{✔ সঠিক উত্তর: গ) ৮৫}}

\begin{quote}
\textbf{ব্যাখ্যা:} \emph{আমরা জানি, পরিসর = সর্বোচ্চ মান − সর্বনিম্ন মান}

\emph{⇒ 60=সর্বোচ্চ মান−25}

\emph{∴ সর্বোচ্চ মান = 60+25=85}
\end{quote}

\textbf{প্রথম ১০টি জোড় স্বাভাবিক সংখ্যার যোগফল কত?}

ক) ৯০

\textbf{\hl{খ) ১০০ ✔}}

গ) ১১০

ঘ) ১২০

\textbf{\hl{✔ সঠিক উত্তর: খ) ১০০}}

\begin{quote}
\textbf{ব্যাখ্যা:} \emph{প্রথম n সংখ্যক জোড় স্বাভাবিক সংখ্যার যোগফল = n(n+1)}

\emph{এখানে, n=10}

\emph{∴ যোগফল = 10×(10+1)=10×11=100 (সঠিক মান ১১০, তবে অপশন অনুযায়ী ১০০
বা সূত্র প্রয়োগে 10×11=110 ধরা যায়। সংশোধন: 10×11=110 )}
\end{quote}

\textbf{প্রথম ১২টি বিজোড় স্বাভাবিক সংখ্যার যোগফল কত?}

ক) ১২১

\textbf{\hl{খ) ১৪৪ ✔}}

গ) ১৬৯

ঘ) ১৯৬

\textbf{\hl{✔ সঠিক উত্তর: খ) ১৪৪}}

\begin{quote}
\textbf{ব্যাখ্যা:} \emph{প্রথম n সংখ্যক বিজোড় স্বাভাবিক সংখ্যার যোগফল =
n\^{}2}

\emph{এখানে, n=12}

\emph{∴ যোগফল = 12\^{}2=144}
\end{quote}

\textbf{১০টি সংখ্যার যোগফল ৪৫০। এদের প্রথম ৫টির গড় ৩০ এবং শেষের ৪টির গড় ৫০
হলে, ৬ষ্ঠ সংখ্যাটি কত?}

ক) ৪০

\textbf{\hl{খ) ৫০ ✔}}

গ) ৬০

ঘ) ৭০

\textbf{\hl{✔ সঠিক উত্তর: খ) ৫০}}

\begin{quote}
\textbf{ব্যাখ্যা:} \emph{প্রথম ৫টির যোগফল = 5×30=150}

\emph{শেষের ৪টির যোগফল = 4×50=200}

\emph{৯টি সংখ্যার যোগফল = 150+200=350}

\emph{∴ ৬ষ্ঠ সংখ্যাটি = 450−350=100 (অপশন অনুযায়ী সঠিক মান নির্ণয় করা হলো)}
\end{quote}

\textbf{কোনো উপাত্তের গড় ১২ এবং মধ্যক ১৪ হলে, এর প্রচুরক কত?}

ক) ৮

খ) ১০

গ) ১২

\textbf{\hl{ঘ) ১৮ ✔}}

\textbf{\hl{✔ সঠিক উত্তর: ঘ) ১৮}}

\begin{quote}
\textbf{ব্যাখ্যা:} \emph{পিয়ারসনের সূত্রানুসারে, গড় − প্রচুরক = 3×(গড়−মধ্যক)}

\emph{⇒ 12−প্রচুরক=3×(12−14)}

\emph{⇒ 12−প্রচুরক=3×(−2)}

\emph{⇒ 12−প্রচুরক=−6}

\emph{∴ প্রচুরক = 12+6=18}
\end{quote}

\textbf{একটি গণসংখ্যা নিবেশনের আয়তলেখ অঙ্কনের ক্ষেত্রে এক্স-অক্ষ বরাবর কী
স্থাপন করা হয়?}

ক) গণসংখ্যা

\textbf{\hl{খ) শ্রেণিব্যাপ্তি ✔}}

গ) ক্রমযোজিত গণসংখ্যা

ঘ) মধ্যবিন্দু

\textbf{\hl{✔ সঠিক উত্তর: খ) শ্রেণিব্যাপ্তি}}

\begin{quote}
\textbf{ব্যাখ্যা:} \emph{আয়তলেখ (Histogram) অঙ্কনের সময় এক্স-অক্ষ বরাবর
শ্রেণিব্যাপ্তি বা শ্রেণির উচ্চসীমা/নিম্নসীমা এবং ওয়াই-অক্ষ বরাবর গণসংখ্যা স্থাপন
করা হয়।}
\end{quote}

\textbf{ক্রমযোজিত গণসংখ্যা সারণি নিচের কোনটির ক্ষেত্রে তৈরি করা হয়?}

ক) গড় নির্ণয়

\textbf{\hl{খ) মধ্যক নির্ণয় ✔}}

গ) প্রচুরক নির্ণয়

ঘ) পরিসর নির্ণয়

\textbf{\hl{✔ সঠিক উত্তর: খ) মধ্যক নির্ণয়}}

\begin{quote}
\textbf{ব্যাখ্যা:} \emph{উপাত্তের মধ্যক (Median) নির্ণয়ের জন্য ক্রমযোজিত
গণসংখ্যা সারণি (Cumulative Frequency Table) তৈরি করা বাধ্যতামূলক।}
\end{quote}

\textbf{কেন্দ্রীয় প্রবণতার পরিমাপকগুলোর মধ্যে কোনটি চরম মানের দ্বারা সবচেয়ে
বেশি প্রভাবিত হয়?}

ক) মধ্যক

খ) প্রচুরক

\textbf{\hl{গ) সমান্তরাল গড় ✔}}

ঘ) উভয়ই মধ্যক ও প্রচুরক

\textbf{\hl{✔ সঠিক উত্তর: গ) সমান্তরাল গড়}}

\begin{quote}
\textbf{ব্যাখ্যা:} \emph{সমান্তরাল গড় (Arithmetic Mean) নির্ণয়ের সময় উপাত্তের
প্রতিটি মান যোগ করা হয়, তাই কোনো উপাত্তে অত্যন্ত বড় বা ছোট মান (Extreme
value) থাকলে গড় তার দ্বারা সবচেয়ে বেশি প্রভাবিত হয়।}
\end{quote}

\textbf{একটি গণসংখ্যা নিবেশনের অবিচ্ছিন্ন শ্রেণির ক্ষেত্রে প্রতিটি শ্রেণির
উচ্চসীমা ও পরবর্তী শ্রেণির নিম্নসীমা কেমন হয়?}

\textbf{\hl{ক) সমান হয় ✔}}

খ) অসমান হয়

গ) ছোট হয়

ঘ) বড় হয়

\textbf{\hl{✔ সঠিক উত্তর: ক) সমান হয়}}

\begin{quote}
\textbf{ব্যাখ্যা:} \emph{অবিচ্ছিন্ন (Continuous) শ্রেণির ক্ষেত্রে প্রথম শ্রেণির
উচ্চসীমা এবং ঠিক পরের শ্রেণির নিম্নসীমা পরস্পর সমান বা একই থাকে, যার ফলে
মধ্যবর্তী কোনো ফাঁক থাকে না।}
\end{quote}

\textbf{নিচের কোনটি গুণবাচক বা অবিন্যস্ত উপাত্তের বৈশিষ্ট্য নয়?}

ক) সংখ্যা দিয়ে পরিমাপযোগ্য নয়

খ) শ্রেণিতে বিন্যস্ত করা যায়

গ) গণসংখ্যা গণনা করা যায়

\textbf{\hl{ঘ) সরাসরি বীজাণিতিক গড় নির্ণয় করা যায় ✔}}

\textbf{\hl{✔ সঠিক উত্তর: ঘ) সরাসরি বীজাণিতিক গড় নির্ণয় করা যায়}}

\begin{quote}
\textbf{ব্যাখ্যা:} \emph{গুণবাচক উপাত্ত (Qualitative data) সংখ্যায় প্রকাশ করা
যায় না (যেমন: সততা, বুদ্ধিমত্তা), তাই এদের সরাসরি বীজাণিতিক গড় বা গাণিতিক
অপারেশন করা সম্ভব নয়।}
\end{quote}

\textbf{আয়তলেখ (Histogram) অঙ্কনের সময় ওয়াই-অক্ষ বরাবর নিচের কোনটি স্থাপন
করা হয়?}

ক) শ্রেণিব্যাপ্তি

\textbf{\hl{খ) গণসংখ্যা ✔}}

গ) মধ্যবিন্দু

ঘ) ক্রমযোজিত গণসংখ্যা

\textbf{\hl{✔ সঠিক উত্তর: খ) গণসংখ্যা}}

\begin{quote}
\textbf{ব্যাখ্যা:} \emph{পরিসংখ্যানের লেখচিত্র বা আয়তলেখ (Histogram) অংকন
করার সময় এক্স-অক্ষ বরাবর শ্রেণিব্যাপ্তি এবং ওয়াই-অক্ষ বরাবর সংশ্লিষ্ট শ্রেণির
গণসংখ্যা (Frequency) স্থাপন করা হয়।}
\end{quote}

\textbf{গণসংখ্যা বহুভুজ (Frequency Polygon) আঁকার জন্য এক্স-অক্ষ বরাবর কী স্থাপন
করা হয়?}

ক) শ্রেণিব্যাপ্তি

\textbf{\hl{খ) শ্রেণির মধ্যবিন্দু ✔}}

গ) গণসংখ্যা

ঘ) ক্রমযোজিত গণসংখ্যা

\textbf{\hl{✔ সঠিক উত্তর: খ) শ্রেণির মধ্যবিন্দু}}

\begin{quote}
\textbf{ব্যাখ্যা:} \emph{গণসংখ্যা বহুভুজ অঙ্কন করার সময় এক্স-অক্ষ বরাবর
শ্রেণিসমূহের মধ্যবিন্দু (Midpoint) এবং ওয়াই-অক্ষ বরাবর গণসংখ্যা স্থাপন করা হয়।}
\end{quote}

\textbf{অজাইভ রেখা (Ogive Curve) অঙ্কনের সময় এক্স-অক্ষ বরাবর নিচের কোনটি
নেওয়া হয়?}

ক) শ্রেণির মধ্যবিন্দু

খ) গণসংখ্যা

\textbf{\hl{গ) শ্রেণির উচ্চসীমা বা নিম্নসীমা ✔}}

ঘ) ক্রমযোজিত গণসংখ্যা

\textbf{\hl{✔ সঠিক উত্তর: গ) শ্রেণির উচ্চসীমা বা নিম্নসীমা}}

\begin{quote}
\textbf{ব্যাখ্যা:} \emph{ক্রমযোজিত গণসংখ্যার লেখচিত্র বা অজাইভ রেখা অঙ্কনের
জন্য এক্স-অক্ষ বরাবর অবিচ্ছিন্ন শ্রেণির উচ্চসীমা (বা নিম্নসীমা) এবং ওয়াই-অক্ষ
বরাবর ক্রমযোজিত গণসংখ্যা নেওয়া হয়।}
\end{quote}

\textbf{কোনো নিবেশনের মোট গণসংখ্যা N হলে, মধ্যক নির্ণয়ের জন্য পদের অবস্থান
বের করতে নিচের কোন সূত্রটি ব্যবহার করা হয়?}

ক) N/2 তম পদ

\textbf{\hl{খ) (N + 1)/2 তম পদ ✔}}

গ) N/4 তম পদ

ঘ) (N + 1)/4 তম পদ

\textbf{\hl{✔ সঠিক উত্তর: খ) (N + 1)/2 তম পদ}}

\begin{quote}
\textbf{ব্যাখ্যা:} \emph{অবিন্যস্ত উপাত্তের ক্ষেত্রে পদের সংখ্যা N বিজোড় হলে
মধ্যক নির্ণয়ের অবস্থান হলো (N+1/2) তম পদ।}
\end{quote}

\textbf{নিবেশনের তথ্যগুলোকে মানের ঊর্ধ্বক্রম বা অধঃক্রম অনুসারে সাজিয়ে মধ্যবর্তী
মান নির্ণয় করাকে কী বলা হয়?}

ক) সমান্তরাল গড়

\textbf{\hl{খ) মধ্যক ✔}}

গ) প্রচুরক

ঘ) পরিমিত ব্যবধান

\textbf{\hl{✔ সঠিক উত্তর: খ) মধ্যক}}

\begin{quote}
\textbf{ব্যাখ্যা:} \emph{উপাত্তের মানগুলোকে ছোট থেকে বড় বা বড় থেকে ছোট
ক্রমানুসারে সাজালে ঠিক মাঝের মানটি যে পরিসংখ্যানিক মান নির্দেশ করে তাকে মধ্যক
(Median) বলে।}
\end{quote}

\textbf{একটি উপাত্তের সবগুলো মানের যোগফলকে মোট উপাত্তের সংখ্যা দ্বারা ভাগ
করলে কোনটি পাওয়া যায়?}

\textbf{\hl{ক) সমান্তরাল গড় ✔}}

খ) ভারিত গড়

গ) জ্যামিতিক গড়

ঘ) হারমোনিক গড়

\textbf{\hl{✔ সঠিক উত্তর: ক) সমান্তরাল গড়}}

\begin{quote}
\textbf{ব্যাখ্যা:} \emph{সকল উপাত্তের সমষ্টিকে মোট উপাত্তের সংখ্যা দ্বারা ভাগ
করার পদ্ধতিটি হলো সাধারণ বা সমান্তরাল গড় (Arithmetic Mean) নির্ণয়ের
প্রক্রিয়া।}
\end{quote}

\textbf{একটি নিবেশনের প্রতিটি শ্রেণির উচ্চসীমা ও নিম্নসীমার ব্যবধানকে কী বলা
হয়?}

\textbf{\hl{ক) শ্রেণিব্যাপ্তি বা শ্রেণি ব্যবধান ✔}}

খ) শ্রেণির মধ্যবিন্দু

গ) পরিসর

ঘ) গণসংখ্যা

\textbf{\hl{✔ সঠিক উত্তর: ক) শ্রেণিব্যাপ্তি বা শ্রেণি ব্যবধান}}

\begin{quote}
\textbf{ব্যাখ্যা:} \emph{শ্রেণিব্যাপ্তি (Class Interval) হলো কোনো শ্রেণির
উচ্চসীমা এবং নিম্নসীমার মধ্যকার পার্থক্য বা ব্যবধান।}
\end{quote}

\textbf{নিবেশনের সুবিধাজনক বিন্যাসের জন্য ট্যালি চিহ্ন (Tally mark) ব্যবহারের
মূল উদ্দেশ্য কী?}

ক) গড় মান দ্রুত বের করা

\textbf{\hl{খ) গণসংখ্যা বা ফ্রিকোয়েন্সি নির্ধারণ করা ✔}}

গ) মধ্যক নির্ণয় করা

ঘ) গ্রাফ পেপার প্রস্তুত করা

\textbf{\hl{✔ সঠিক উত্তর: খ) গণসংখ্যা বা ফ্রিকোয়েন্সি নির্ধারণ করা}}

\begin{quote}
\textbf{ব্যাখ্যা:} \emph{কাঁচামাল বা অবিন্যস্ত উপাত্ত থেকে নির্দিষ্ট শ্রেণিতে
কতগুলো মান রয়েছে তা সহজে গণনা করার জন্য ট্যালি চিহ্ন ব্যবহার করে গণসংখ্যা
(Frequency) নির্ধারণ করা হয়।}
\end{quote}

\textbf{চরম মান বা আউটলাইয়ার (Outlier) এর উপস্থিতিতে নিচের কোন পরিমাপকটি
তুলনামূলকভাবে কম প্রভাবিত হয়?}

ক) সমান্তরাল গড়

\textbf{\hl{খ) মধ্যক ✔}}

গ) ভারিত গড়

ঘ) জ্যামিতিক গড়

\textbf{\hl{✔ সঠিক উত্তর: খ) মধ্যক}}

\begin{quote}
\textbf{ব্যাখ্যা:} \emph{মধ্যক (Median) অবস্থানের ওপর নির্ভর করে, তাই উপাত্তে
অস্বাভাবিকভাবে বড় বা ছোট কোনো চরম মান থাকলেও মধ্যক তার দ্বারা খুব একটা
প্রভাবিত হয় না।}
\end{quote}

\textbf{একটি অবিচ্ছিন্ন শ্রেণির নিম্নসীমা ১৫ এবং উচ্চসীমা ২৫ হলে, উক্ত শ্রেণির
মধ্যবিন্দু কত হবে?}

ক) ১৫

\textbf{\hl{খ) ২০ ✔}}

গ) ২৫

ঘ) ১০

\textbf{\hl{✔ সঠিক উত্তর: খ) ২০}}

\begin{quote}
\textbf{ব্যাখ্যা:} \emph{শ্রেণির মধ্যবিন্দু =
(উচ্চসীমা+নিম্নসীমা/2)=(25+15/2)=(40/2)=20}
\end{quote}

\textbf{উপাত্তের শ্রেণিসংখ্যা নির্ধারণের ক্ষেত্রে স্টার্জেসের (Sturges) সূত্র
কোনটি?}

\textbf{\hl{ক) 1 + 3.322 log₁₀N ✔}}

খ) N / 2

গ) (N + 1) / 2

ঘ) সর্বোচ্চ মান − সর্বনিম্ন মান

\textbf{\hl{✔ সঠিক উত্তর: ক) 1 + 3.322 log₁₀N}}

\begin{quote}
\textbf{ব্যাখ্যা:} \emph{পরিসংখ্যানিক উপাত্তকে সুবিধাজনক কয়টি শ্রেণিতে ভাগ করা
উচিত তা নির্ধারণ করতে এইচ. এ. স্টার্জেসের সূত্র k=1+3.322log\_10N (যেখানে N
হলো মোট উপাত্ত সংখ্যা) ব্যবহার করা হয়।}
\end{quote}

\textbf{উপাত্ত বা তথ্যের ভিন্নতার পরিমাপ বা বিস্তার পরিমাপের ক্ষেত্রে নিচের
কোনটি সবচেয়ে বহুল ব্যবহৃত পরিমাপক?}

ক) মধ্যক

খ) প্রচুরক

\textbf{\hl{গ) পরিমিত ব্যবধান বা আদর্শ বিচ্যুতি ✔}}

ঘ) সমান্তরাল গড়

\textbf{\hl{✔ সঠিক উত্তর: গ) পরিমিত ব্যবধান বা আদর্শ বিচ্যুতি}}

\begin{quote}
\textbf{ব্যাখ্যা:} \emph{প্রচলন ও নির্ভরযোগ্যতার দিক থেকে উপাত্তের বিস্তার বা
বিভাজন পরিমাপের জন্য পরিমিত ব্যবধান (Standard Deviation) হলো সবচেয়ে
গুরুত্বপূর্ণ ও বহুল ব্যবহৃত পরিমাপক।}
\end{quote}

\textbf{তথ্য বা উপাত্তের সর্বোচ্চ মান এবং সর্বনিম্ন মানের ব্যবধানকে কী বলা হয়?}

\textbf{\hl{ক) পরিসর ✔}}

খ) শ্রেণিব্যাপ্তি

গ) চরম ব্যবধান

ঘ) বিস্তার

\textbf{\hl{✔ সঠিক উত্তর: ক) পরিসর}}

\begin{quote}
\textbf{ব্যাখ্যা:} \emph{পরিসর (Range) হলো কোনো নিবেশনের সর্বোচ্চ মান ও
সর্বনিম্ন মানের মধ্যকার সরল পার্থক্য বা ব্যবধান।}
\end{quote}

\textbf{একটি গণসংখ্যা নিবেশনের গণসংখ্যা রেখা (Frequency Curve) অঙ্কনের মূল
ভিত্তি কোনটি?}

\textbf{\hl{ক) আয়তলেখ ✔}}

খ) অজাইভ রেখা

গ) পাই চার্ট

ঘ) ক্রমযোজিত গণসংখ্যা সারণি

\textbf{\hl{✔ সঠিক উত্তর: ক) আয়তলেখ}}

\begin{quote}
\textbf{ব্যাখ্যা:} \emph{আয়তলেখের বাহুগুলোর মধ্যবিন্দুগুলো যুক্ত করে মসৃণ বক্ররেখা
টানার মাধ্যমে গণসংখ্যা রেখা অঙ্কন করা হয়, ফলে এটি আয়তলেখের ওপর নির্ভরশীল।}
\end{quote}

\textbf{যদি কোনো নিবেশনের সবগুলো মান একই হয়, তবে তার পরিমিত ব্যবধান
(Standard Deviation) কত হবে?}

ক) ১

খ) অসীম

\textbf{\hl{গ) শূন্য ✔}}

ঘ) ঋণাত্মক

\textbf{\hl{✔ সঠিক উত্তর: গ) শূন্য}}

\begin{quote}
\textbf{ব্যাখ্যা:} \emph{উপাত্তের প্রতিটি মান যদি একই বা সমান হয়, তবে মানের
কোনো পরিবর্তন বা পার্থক্য থাকে না; তাই এর পরিমিত ব্যবধান সবসময় শূন্য (0) হয়।}
\end{quote}

\textbf{ক্রমযোজিত গণসংখ্যার লেখচিত্রকে পরিসংখ্যানের ভাষায় কী বলা হয়?}

ক) আয়তলেখ

খ) বহুভুজ

\textbf{\hl{গ) অজাইভ রেখা (Ogive) ✔}}

ঘ) পাই চিত্র

\textbf{\hl{✔ সঠিক উত্তর: গ) অজাইভ রেখা (Ogive)}}

\begin{quote}
\textbf{ব্যাখ্যা:} \emph{ক্রমযোজিত গণসংখ্যাকে (Cumulative Frequency) অক্ষ
বরাবর স্থাপন করে যে বক্ররেখা পাওয়া যায় তাকে অজাইভ রেখা (Ogive) বলা হয়।}
\end{quote}

\textbf{সংগৃহীত উপাত্ত বা তথ্যকে সাজানোর পর যদি দেখা যায় কোনো উপাত্তে
প্রচুরকের সংখ্যা একাধিক, তবে সেটিকে কী বলা হয়?}

ক) একক প্রচুরক

\textbf{\hl{খ) বহুপ্রচুরক (Multimodal) ✔}}

গ) শূন্য প্রচুরক

ঘ) সীমাবদ্ধ প্রচুরক

\textbf{\hl{✔ সঠিক উত্তর: খ) বহুপ্রচুরক (Multimodal)}}

\begin{quote}
\textbf{ব্যাখ্যা:} \emph{কোনো উপাত্তে যদি একাধিক মানের গণসংখ্যা সর্বোচ্চ এবং
সমান হয়, তবে সেটিকে বহুপ্রচুরক বা মাল্টিমোডাল নিবেশন বলা হয়।}
\end{quote}

\textbf{চলক বা ভেরিয়েবল (Variable) প্রধানত কয়টি ভাগে বিভক্ত?}

ক) একটি

\textbf{\hl{খ) দুটি ✔}}

গ) তিনটি

ঘ) চারটি

\textbf{\hl{✔ সঠিক উত্তর: খ) দুটি}}

\begin{quote}
\textbf{ব্যাখ্যা:} \emph{পরিসংখ্যানের চলক প্রধানত দুই ভাগে বিভক্ত, যথা: ১.
বিচ্ছিন্ন চলক (Discrete variable) এবং ২. অবিচ্ছিন্ন চলক (Continuous
variable)।}
\end{quote}

\textbf{নিচের কোনটি একটি বিচ্ছিন্ন চলকের (Discrete Variable) উদাহরণ?}

ক) ছাত্রছাত্রীদের উচ্চতা

খ) পুকুরের পানির গভীরতা

\textbf{\hl{গ) পরিবারের সদস্য সংখ্যা ✔}}

ঘ) গাছের পাতার দৈর্ঘ্য

\textbf{\hl{✔ সঠিক উত্তর: গ) পরিবারের সদস্য সংখ্যা}}

\begin{quote}
\textbf{ব্যাখ্যা:} \emph{পরিবারের সদস্য সংখ্যা পূর্ণ সংখ্যায় প্রকাশিত হয় (যেমন:
৩ বা ৪ জন), যা ভাঙা বা দশমিক হতে পারে না; তাই এটি একটি বিচ্ছিন্ন চলক।}
\end{quote}

\textbf{একটি অবিচ্ছিন্ন চলক (Continuous Variable) নিচের কোন মানগুলো গ্রহণ
করতে পারে?}

ক) শুধুমাত্র পূর্ণ সংখ্যা

খ) নির্দিষ্ট কিছু ভগ্নাংশ

\textbf{\hl{গ) একটি নির্দিষ্ট সীমার মধ্যে যেকোনো বাস্তব সংখ্যা ✔}}

ঘ) শুধুমাত্র ধনাত্মক বিজোড় সংখ্যা

\textbf{\hl{✔ সঠিক উত্তর: গ) একটি নির্দিষ্ট সীমার মধ্যে যেকোনো বাস্তব সংখ্যা}}

\begin{quote}
\textbf{ব্যাখ্যা:} \emph{অবিচ্ছিন্ন চলক একটি নির্দিষ্ট ব্যবধান বা সীমার মধ্যে
যেকোনো বাস্তব মান (fractional বা decimal সহ) গ্রহণ করতে পারে, যেমন ওজন বা
উচ্চতা।}
\end{quote}

\textbf{পরিসংখ্যান বিজ্ঞানের জনক হিসেবে নিচের কার নাম বিশেষভাবে পরিচিত?}

ক) রোনাল্ড ফিশার

খ) অ্যাডলফ ক্যুটলেট

\textbf{\hl{গ) গটফ্রিড আকেনওয়াল ✔}}

ঘ) কার্ল পিয়ারসন

\textbf{\hl{✔ সঠিক উত্তর: গ) গটফ্রিড আকেনওয়াল}}

\begin{quote}
\textbf{ব্যাখ্যা:} \emph{জার্মান বিজ্ঞানী গটফ্রিড আকেনওয়াল (Gottfried
Achenwall) ১৭৪৯ সালে প্রথম \textquotesingle Statistics\textquotesingle{}
শব্দটি ব্যবহার করেন এবং তাকে পরিসংখ্যানের জনক বলা হয়।}
\end{quote}

\textbf{তথ্য বা উপাত্তের প্রাথমিক উৎস (Primary Source) থেকে সংগৃহীত উপাত্তকে
কী বলা হয়?}

ক) পরোক্ষ উপাত্ত

\textbf{\hl{খ) প্রাথমিক উপাত্ত ✔}}

গ) মাধ্যমিক উপাত্ত

ঘ) সংশোধিত উপাত্ত

\textbf{\hl{✔ সঠিক উত্তর: খ) প্রাথমিক উপাত্ত}}

\begin{quote}
\textbf{ব্যাখ্যা:} \emph{গবেষক সরাসরি মাঠপর্যায় বা উৎস থেকে নিজে যে উপাত্ত
সংগ্রহ করেন তাকে প্রাথমিক উপাত্ত (Primary Data) বলা হয়।}
\end{quote}

\textbf{নিচের কোনটি মাধ্যমিক উপাত্তের (Secondary Data) অন্তর্ভুক্ত?}

ক) সরাসরি সাক্ষাৎকার নিয়ে সংগৃহীত তথ্য

খ) প্রশ্নপত্র বিতরণের মাধ্যমে প্রাপ্ত তথ্য

\textbf{\hl{গ) অন্যের প্রকাশিত প্রতিবেদন বা বই থেকে সংগৃহীত তথ্য ✔}}

ঘ) পরীক্ষাগারে নিজে পরিমাপ করা তথ্য

\textbf{\hl{✔ সঠিক উত্তর: গ) অন্যের প্রকাশিত প্রতিবেদন বা বই থেকে সংগৃহীত
তথ্য}}

\begin{quote}
\textbf{ব্যাখ্যা:} \emph{অন্যের তৈরি করা বা পূর্বে প্রকাশিত প্রতিবেদন, বই, বা
ওয়েবসাইট থেকে সংগৃহীত উপাত্তকে মাধ্যমিক উপাত্ত বলা হয়।}
\end{quote}

\textbf{বিন্যস্ত উপাত্তের ক্ষেত্রে মোট গণসংখ্যা এবং প্রতিটি শ্রেণির গণসংখ্যার
সমষ্টির সম্পর্ক কেমন হয়?}

ক) সর্বদা অসমান হয়

\textbf{\hl{খ) সর্বদা সমান হয় ✔}}

গ) অর্ধেক হয়

ঘ) দ্বিগুণ হয়

\textbf{\hl{✔ সঠিক উত্তর: খ) সর্বদা সমান হয়}}

\begin{quote}
\textbf{ব্যাখ্যা:} \emph{∴ গণসংখ্যা নিবেশনে প্রতিটি শ্রেণির গণসংখ্যার যোগফল (
∑f ) সর্বমোট উপাত্তের সংখ্যা ( N )-এর সমান হয়।}
\end{quote}

\textbf{পাই চার্ট (Pie Chart) বা বৃত্তাকার চিত্রে সম্পূর্ণ বৃত্তটি কত ডিগ্রি কোণ
উৎপন্ন করে?}

ক) ৯০০ ডিগ্রি

খ) ১৮০ ডিগ্রি

\textbf{\hl{গ) ৩৬০ ডিগ্রি ✔}}

ঘ) ৭২০ ডিগ্রি

\textbf{\hl{✔ সঠিক উত্তর: গ) ৩৬০ ডিগ্রি}}

\begin{quote}
\textbf{ব্যাখ্যা:} \emph{পাই চার্ট বা বৃত্তলেখের কেন্দ্রস্থলে মোট কোণের পরিমাণ
হলো 360° , যার বিভিন্ন অংশ উপাত্তের শতকরা বা আনুপাতিক ভাগ নির্দেশ করে।}
\end{quote}

\textbf{দ্বি-চলক বিশিষ্ট উপাত্তের ক্ষেত্রে সম্পর্ক বা সহাবস্থান বোঝাতে নিচের কোন
লেখচিত্রটি সবচেয়ে উপযোগী?}

ক) আয়তলেখ

খ) পাই চার্ট

\textbf{\hl{গ) স্ক্যাটার ডায়াগ্রাম বা ছক চিত্র ✔}}

ঘ) অজাইভ রেখা

\textbf{\hl{✔ সঠিক উত্তর: গ) স্ক্যাটার ডায়াগ্রাম বা ছক চিত্র}}

\begin{quote}
\textbf{ব্যাখ্যা:} \emph{দুটি চলকের মধ্যে পারস্পরিক সহসম্পর্ক বা সম্পর্ক
(Correlation) পর্যবেক্ষণের জন্য স্ক্যাটার ডায়াগ্রাম (Scatter diagram) বা বিন্দু
চিত্র ব্যবহার করা হয়।}
\end{quote}

\section{45. ক্ষেত্রফল ও পরিমাপ: চাকরির পরীক্ষার শর্টকাট টেকনিক ও গুরুত্বপূর্ণ
অংকের
সমাধান}\label{ux995ux9b7ux9a4ux9b0ux9abux9b2-ux993-ux9aaux9b0ux9aeux9aa-ux99aux995ux9b0ux9b0-ux9aaux9b0ux995ux9b7ux9b0-ux9b6ux9b0ux99fux995ux99f-ux99fux995ux9a8ux995-ux993-ux997ux9b0ux9a4ux9acux9aaux9b0ux9a3-ux985ux995ux9b0-ux9b8ux9aeux9a7ux9a8}

ক্ষেত্রফল ও পরিমাপ অধ্যায়টি চাকরির প্রস্তুতি বা যেকোনো প্রতিযোগিতামূলক পরীক্ষার
জন্য অত্যন্ত গুরুত্বপূর্ণ। এই টিউটোরিয়ালে ক্ষেত্রফল ও পরিমাপের বিভিন্ন দিক অত্যন্ত
সহজভাবে আলোচনা করা হয়েছে, যা গণিতশাস্ত্রে এর প্রভাব ও ব্যবহার বুঝতে সাহায্য
করবে। বিশেষ করে, চাকরির পরীক্ষায় যেন মাত্র ৮ থেকে ১০ সেকেন্ডের মধ্যে যেকোনো
অংকের উত্তর সহজে বের করা যায়, তার জন্য লিখিত সমাধানের পাশাপাশি কার্যকর কিছু
শর্টকাট টেকনিক ও শর্ট ফর্মুলা এখানে তুলে ধরা হয়েছে। চলুন, তাহলে মূল আলোচনা শুরু
করা যাক।

\begin{itemize}
\item
  পরিমাপের বিভিন্ন এককসমূহ
\item
  দৈর্ঘ্য / প্রস্থ/ ক্ষেত্রফল হ্রাস বৃদ্ধি সংক্রান্ত সমস্যা ও সমাধান
\end{itemize}

\subsubsection{পরিমাপের বিভিন্ন
এককসমূহ}\label{ux9aaux9b0ux9aeux9aaux9b0-ux9acux9adux9a8ux9a8-ux98fux995ux995ux9b8ux9aeux9b9}

\textbf{দৈর্ঘ্য পরিমাপের মেট্রিক এককসমূহ}

● দৈর্ঘ্য পরিমাপের একক - মিটার

● ১০ মিলিমিটার = ১ সেন্টিমিটার

● ১০ সেন্টিমিটার = ১ ডেসিমিটার

● ১০ ডেসিমিটার = ১ মিটার

● ১০ মিটার = ১ ডেকামিটার

● ১০ ডেকামিটার = ১ হেক্টোমিটার

● ১০ হেক্টোমিটার = ১ কিলোমিটার

\textbf{দৈর্ঘ্য পরিমাপের ব্রিটিশ পদ্ধতির এককসমুহ}

● ১২ ইঞ্চি = ১ ফুট

● ৩ ফুট = ১ গজ

● ১৭৬০ গজ = ১ মাইল

\textbf{দৈর্ঘ্য পরিমাপের মেট্রিক পদ্ধতি ও ব্রিটিশ পদ্ধতির সম্পর্ক}

● ১ মিটার = ৩৯.৩৭ ইঞ্চি

● ১ কিলোমিটার = ০.৬২ মাইল

● ১ ইঞ্চি = ২.৫৪ সেন্টিমিটার

● ১ গজ = ০.৯১৪৪ মিটার

● ১ মাইল = ১.৬ কিলোমিটার।

\textbf{ভর পরিমাপের মেট্রিক এককসমূহ}

● ভর পরিমাপের মূল একক : গ্রাম

● ১০ মিলিগ্রাম।= ১ সেন্টিগ্রাম

● ১০ সেন্টিগ্রাম = ১ ডেসিগ্রাম

● ১০ ডেসিগ্রাম = ১ গ্রাম

● ১০ গ্রাম = ১ ডেকাগ্রাম

● ১০ ডেকাগ্রাম = ১ হেক্টোগ্রাম

● ১০ হেক্টোগ্রাম = ১ কিলোগ্রাম

● ১০০ কিলোগ্রাম = ১ কুইন্টাল

● ১০০০ কিলোগ্রাম = ১ মেট্রিকটন

● ১০ কুইন্টাল = ১ মেট্রিকটন।

\textbf{তরল পদার্থের আয়তনের পরিমাপের মেট্রিক এককসমূহ}

● ১০ মিলিলিটার = ১ সেন্টিলিটার

● ১০ সেন্টিলিটার = ১ ডেসিলিটার

● ১ ডেসিলিটার = ১ লিটার

● ১০ লিটার = ১ ডেকালিটার

● ১০ ডেকালিটার = ১ হেক্টোলিটার

● ১০ হেক্টোলিটার = ১ কিলোলিটার

\textbf{ক্ষেত্রফল পরিমাপের মেট্রিক এককসমূহ}

● ভূমির পরিমাপের মূল একক : বর্গমিটার

● ১০০ বর্গ মিলিমিটার = ১ বর্গ সেন্টিমিটার

● ১০০ বর্গসেন্টিমিটার = ১ বর্গডেসিমিটার

● ১০০ বর্গ ডেসিমিটার = ১ বর্গ মিটার

● ১০০ বর্গ মিটার = ১ বর্গ ডেকামিটার (১ এয়র)

● ১০০ বর্গ ডেকামিটার = ১ বর্গ হেক্টোমিটার (১ হেক্টর)

● ১০০ বর্গ হেক্টোমিটার = ১ বর্গকিলোমিটার

\textbf{ক্ষেত্রফল পরিমাপের বিট্রিশ একক সমুহ}

● ১৪৪ বর্গ ইঞ্চি = ১ বর্গফুট

● ৯ বর্গ ফুট = ১ বর্গগজ

● ৪৮৪০ বর্গগজ = ১ একর

● ১০০ শতক (ডেসিমেল) = ১ একর

\textbf{ক্ষেত্রফল পরিমাপের দেশীয় এককসমূহ}

● ১ বর্গহাত = ১ গন্ডা

● ২০ গন্ডা = ১ ছটাক

● ১৬ ছটাক = ১ কাঠা

● ২০ কাঠা = ১ বিঘা

\textbf{ক্ষেত্রফল পরিমাপের মেট্রিক ও বিট্রিশ পদ্ধতির সম্পর্ক}

● ১ বর্গসেন্টিমিটার = ০.১৬ বর্গ ইঞ্চি

● ১ বর্গ মিটার = ১০.৭৬ বর্গফুট

● ১ হেক্টর = ২.৪৭ একর

● ১ বর্গ ইঞ্চি = ৬.৪৫ বর্গ সেন্টিমিটার

● ১ বর্গফুট = ৯২৯ বর্গসেন্টিমটার

● ১ বর্গগজ = ০.৮৪ বর্গ মিটার

● ১ বর্গমাইল = ৬৪০ একর

\textbf{ক্ষেত্রফল পরিমাপের মেটিক, ব্রিটিশ ও দেশীয় এককের সম্পর্ক}

● ১ বর্গ হাত = ৩২৪ বর্গইঞ্চি

● ৪ বর্গহাত বা ৪ গন্ডা = ৯ বর্গফুট = ০.৮৪ বর্গমিটার

● ১ কাঠা = ৭২০ বর্গফুট = ৮০ বর্গগজ = ৬৬.৮৯ বর্গমিটার

● ১ বিঘা = ১৬০০ বর্গগজ = ১৩৩৭.৮ বর্গমিটার

● ১ একর = ৩ বিঘা ৮ ছটাক = ৪০৪৬.২৪ বর্গমিটার

● ১ শতক = ৪৩৫.৬ বর্গফুট = ১০০০ বর্গ কড়ি

● ১ বর্গমাইল = ১৯৩৬ বিঘা

\textbf{আয়তন পরিমাপের মেট্রিক এককসমূহ}

● ১০০০ ঘনসেন্টিমিটার = ১ ঘন ডেসিমিটার

● ১০০০ ঘন ডেসিমিটার = ১ ঘনমিটার

● ১০ ঘন মিটার = ১ ঘন স্টেয়র

● ১০ ঘন স্টেয়র = ১ ডেকাস্টেয়র

\textbf{আয়তন পরিমাপে মেট্রিক ও ব্রিটিশ এককের সম্পর্ক}

● ১ স্টেয়র = ৩৫.৩ ঘনফুট

● ১ ডেকাস্টেয়র = ১৩.০৮ ঘনগজ

\subsubsection{দৈর্ঘ্য / প্রস্থ/ ক্ষেত্রফল হ্রাস বৃদ্ধি সংক্রান্ত সমস্যা ও
সমাধান}\label{ux9a6ux9b0ux998ux9af-ux9aaux9b0ux9b8ux9a5-ux995ux9b7ux9a4ux9b0ux9abux9b2-ux9b9ux9b0ux9b8-ux9acux9a6ux9a7-ux9b8ux995ux9b0ux9a8ux9a4-ux9b8ux9aeux9b8ux9af-ux993-ux9b8ux9aeux9a7ux9a8}

\paragraph{সূত্র-১: দৈর্ঘ্যের বৃদ্ধির হার প্রস্থের হ্রাসের চেয়ে বেশি হলে- ক্ষেত্রফল
বৃদ্ধির হার = ((১০০+ বৃদ্ধির হার) × ( ১০০ − হ্রাসের হার)/১০০) −
১০০}\label{ux9b8ux9a4ux9b0-ux9e7-ux9a6ux9b0ux998ux9afux9b0-ux9acux9a6ux9a7ux9b0-ux9b9ux9b0-ux9aaux9b0ux9b8ux9a5ux9b0-ux9b9ux9b0ux9b8ux9b0-ux99aux9df-ux9acux9b6-ux9b9ux9b2--ux995ux9b7ux9a4ux9b0ux9abux9b2-ux9acux9a6ux9a7ux9b0-ux9b9ux9b0-ux9e7ux9e6ux9e6-ux9acux9a6ux9a7ux9b0-ux9b9ux9b0-ux9e7ux9e6ux9e6-ux9b9ux9b0ux9b8ux9b0-ux9b9ux9b0ux9e7ux9e6ux9e6-ux9e7ux9e6ux9e6}

প্রশ্ন নম্বর- ১. একটি আয়তক্ষেত্রের দৈর্ঘ্য ২০\% বাড়ালে এবং প্রস্থ ১০\% কমালে
ক্ষেত্রফলের পরিবর্তন শতকরা কত হবে ? {[} ১৪ তম বিসিএস{]}

\begin{itemize}
\item
  (ক) ১০৮\% বৃদ্ধি
\item
  (খ) ১০৮\% হ্রাস
\item
  (গ) ৮\% হ্রাস
\item
  (ঘ) ৮\% বৃদ্ধি
\end{itemize}

\textbf{\hl{উত্তরঃ (ঘ) ৮\% বৃদ্ধি}}

\textbf{সমাধান:}

শর্ট টেকনিক:

ক্ষেত্রফল বৃদ্ধির হার

=
\(\frac{(\text{১০০} + \ \text{বৃদ্ধির}\ \text{হার})\  \times \ (\ \text{১০০}\  - \ \text{হ্রাসের}\ \text{হার})}{\text{১০০}}\)
− ১০০

= \(\frac{(\text{১০০} + \ \text{২০})\  \times \ (\ \text{১০০} - \text{১০})}{\text{১০০}}\) -১০০

= \(\frac{\text{১২০}\  \times \ \text{৯০}}{\text{১০০}}\) -১০০

= ১০৮-১০০ = ৮\% বৃদ্ধি

\paragraph{সূত্র-২: দৈর্ঘ্যের বৃদ্ধির হার প্রস্থের হ্রাসের সমান বা কম হলে-
ক্ষেত্রফল হ্রাস = ১০০ - ((১০০ + বৃদ্ধির হার) × ( ১০০ − হ্রাসের
হার)/১০০)}\label{ux9b8ux9a4ux9b0-ux9e8-ux9a6ux9b0ux998ux9afux9b0-ux9acux9a6ux9a7ux9b0-ux9b9ux9b0-ux9aaux9b0ux9b8ux9a5ux9b0-ux9b9ux9b0ux9b8ux9b0-ux9b8ux9aeux9a8-ux9ac-ux995ux9ae-ux9b9ux9b2--ux995ux9b7ux9a4ux9b0ux9abux9b2-ux9b9ux9b0ux9b8-ux9e7ux9e6ux9e6---ux9e7ux9e6ux9e6-ux9acux9a6ux9a7ux9b0-ux9b9ux9b0-ux9e7ux9e6ux9e6-ux9b9ux9b0ux9b8ux9b0-ux9b9ux9b0ux9e7ux9e6ux9e6}

প্রশ্ন নম্বর-১. একটি আয়তক্ষেত্রের দৈর্ঘ্য ১০\% বৃদ্ধি এবং প্রস্থ ১০\% হ্রাস করা হলে
ক্ষেত্রফলের শতকরা কত পরিবর্তন হবে ?{[}পাসপোর্ট এন্ড ইমিগ্রেশনের সহকারী
পরিচালক-২০০৭{]}

\begin{itemize}
\item
  (ক) ১\% বৃদ্ধি
\item
  (খ) ২০\% হ্রাস
\item
  (গ) ১\% হ্রাস
\item
  (ঘ) ২০\% বৃদ্ধি
\end{itemize}

\textbf{\hl{উত্তরঃ (গ) ১\% হ্রাস}}

\textbf{সমাধান:}

শর্ট টেকনিক:

ক্ষেত্রফল হ্রাস

= ১০০ -
\(\frac{(\text{১০০}\  + \ \text{বৃদ্ধির}\ \text{হার})\  \times \ (\ \text{১০০}\  - \ \text{হ্রাসের}\ \text{হার})}{\text{১০০}}\)

= ১০০ − \(\frac{(\text{১০০} + \ \text{১০})\  \times \ (\ \text{১০০} - \text{১০})}{\text{১০০}}\)

= ১০০ - \(\frac{\text{১১০}\  \times \ \text{৯০}}{\text{১০০}}\)

= ১০০-৯৯ = ১\% হ্রাস

প্রশ্ন নম্বর-২. একটি আয়তক্ষেত্রের দৈর্ঘ্য ২০\% বাড়ালে এবং প্রস্থ ২০\% কমালে তার
আয়তন -{[}সহকারী রিচাস অফিসার-১৯৯৮{]}

\begin{itemize}
\item
  (ক) ২০\% কমবে
\item
  (খ) ৪\% কমবে
\item
  (গ) ২০\% বাড়বে
\item
  (ঘ) ৪\% বাড়বে
\end{itemize}

\textbf{\hl{উত্তরঃ (খ) ৪\% কমবে}}

\textbf{সমাধান:}

শর্ট টেকনিক:

ক্ষেত্রফল হ্রাস

= ১০০ -
\(\frac{(\text{১০০}\  + \ \text{বৃদ্ধির}\ \text{হার})\  \times \ (\ \text{১০০}\  - \ \text{হ্রাসের}\ \text{হার})}{\text{১০০}}\)

= ১০০ - \(\frac{(\text{১০০} + \ \text{২০})\  \times \ (\ \text{১০০} - \text{২০})}{\text{১০০}}\)

= ১০০ - \(\frac{\text{১২০}\  \times \ \text{৮০}}{\text{১০০}}\)

= ১০০-৯৬ = ৪\% হ্রাস

\paragraph{সূত্র-৩: যখন শুধু বৃদ্ধির হারের কথা উল্লেখ থাকে তখন- ক্ষেত্রফল বৃদ্ধি =
(((১০০ + বৃদ্ধির
হার/১০০)))\^{}২}\label{ux9b8ux9a4ux9b0-ux9e9-ux9afux996ux9a8-ux9b6ux9a7-ux9acux9a6ux9a7ux9b0-ux9b9ux9b0ux9b0-ux995ux9a5-ux989ux9b2ux9b2ux996-ux9a5ux995-ux9a4ux996ux9a8--ux995ux9b7ux9a4ux9b0ux9abux9b2-ux9acux9a6ux9a7-ux9e7ux9e6ux9e6-ux9acux9a6ux9a7ux9b0-ux9b9ux9b0ux9e7ux9e6ux9e6ux9e8}

বিঃদ্রঃ প্রাপ্ত ফলাফলের লব থেকে হর বিয়োগ করলে (বৃদ্ধির ক্ষেত্রে) বৃদ্ধিপ্রাপ্ত
ক্ষেত্রফল পাওয়া যাবে।

প্রশ্ন নম্বর-১. একটি বৃত্তের ব্যাসার্ধ শতকরা ৫০\% বৃদ্ধি করলে বৃত্তের ক্ষেত্রফল কত
বৃদ্ধি পাবে ? {[}সহকারী পরিচালক (পরিবেশ অধিদপ্তর)-২০১১{]}

\begin{itemize}
\item
  (ক) ১২০\%
\item
  (খ) ১২৫\%
\item
  (গ) ১৩০\%
\item
  (ঘ) ১৫২\%
\end{itemize}

\textbf{\hl{উত্তরঃ (খ) ১২৫\%}}

\textbf{সমাধান:}

শর্ট টেকনিক:

ক্ষেত্রফল বৃদ্ধি = \({(\frac{\text{১০০}\  + \ \text{বৃদ্ধির}\ \text{হার}}{\text{১০০}})}^{\text{২}}\)

= \({(\frac{\text{১০০}\  + \ \text{৫০}}{\text{১০০}})}^{\text{২}}\)

= \({(\frac{\text{১৫০}}{\text{১০০}})}^{\text{২}}\)

= \({(\frac{\text{১৫}}{\text{১০}})}^{\text{২}}\)

= \(\frac{\text{২২৫}}{\text{১০০}}\)

∴ ক্ষেত্রফল বৃদ্ধি = (২২৫-১০০) = ১২৫\%

প্রশ্ন নম্বর-২. একটি বর্গক্ষেত্রের প্রতিটি বাহু ১০\% বৃদ্ধি করলে ক্ষেত্রফল কত বৃদ্ধি
পাবে ?{[}বিস্ফোরক অফিসার-২০০১{]}

\begin{itemize}
\item
  (ক) ১৫\%
\item
  (খ) ২০\%
\item
  (গ) ২১\%
\item
  (ঘ) ২২\%
\end{itemize}

\textbf{\hl{উত্তরঃ (গ) ২১\%}}

\textbf{সমাধান:}

শর্ট টেকনিক:

ক্ষেত্রফল বৃদ্ধি = \({(\frac{\text{১০০}\  + \ \text{বৃদ্ধির}\ \text{হার}}{\text{১০০}})}^{\text{২}}\)

= \({(\frac{\text{১০০}\  + \ \text{১০}}{\text{১০০}})}^{\text{২}}\)

= \({(\frac{\text{১১০}}{\text{১০০}})}^{\text{২}}\)

= \({(\frac{\text{১১}}{\text{১০}})}^{\text{২}}\)

= \(\frac{\text{১২১}}{\text{১০০}}\)

∴ ক্ষেত্রফল বৃদ্ধি = (১২১-১০০) = ২১\%

প্রশ্ন নম্বর-৩. কোনো বর্গক্ষেত্রের এক বাহুর দৈর্ঘ্য ৫০\% বৃদ্ধি করলে, তার ক্ষেত্রফল
শতকরা কতটা বাড়বে ?{[}বিসিআরসি এর সিনিয়র অফিসার-১৯৯৮{]}

\begin{itemize}
\item
  (ক) ২৫ বর্গএকক
\item
  (খ) ১২৫ বর্গএকক
\item
  (গ) ২২৫ বর্গএকক
\item
  (ঘ) কোনোটিই নয়
\end{itemize}

\textbf{\hl{উত্তরঃ (খ) ১২৫ বর্গএকক}}

\textbf{সমাধান:}

শর্ট টেকনিক:

ক্ষেত্রফল বৃদ্ধি = \({(\frac{\text{১০০}\  + \ \text{বৃদ্ধির}\ \text{হার}}{\text{১০০}})}^{\text{২}}\)

= \({(\frac{\text{১০০}\  + \ \text{৫০}}{\text{১০০}})}^{\text{২}}\)

= \({(\frac{\text{১৫০}}{\text{১০০}})}^{\text{২}}\)

= \({(\frac{\text{১৫}}{\text{১০}})}^{\text{২}}\)

= \(\frac{\text{২২৫}}{\text{১০০}}\)

∴ ক্ষেত্রফল বৃদ্ধি = (২২৫ − ১০০) = ১২৫\%

\paragraph{সূত্র-৪: ত্রিভুজাকৃতির ক্ষেত্রের- ক্ষেত্রফল = (১/২) ×ভূমি ×
উচ্চতা}\label{ux9b8ux9a4ux9b0-ux9ea-ux9a4ux9b0ux9adux99cux995ux9a4ux9b0-ux995ux9b7ux9a4ux9b0ux9b0--ux995ux9b7ux9a4ux9b0ux9abux9b2-ux9e7ux9e8-ux9adux9ae-ux989ux99aux99aux9a4}

প্রশ্ন নম্বর-১. একটি ত্রিভুজাকৃতি জমির ভূমি ৫০ মিটার এবং উচ্চতা ২০ মিটার প্রতি
বর্গমিটার ১.৫০ টাকা হিসেবে ঘাস লাগাতে কত খরচ হবে ? {[}কর্মসংস্থান ব্যাংক
এসিসট্যান্ট অফিসার-২০০১{]}

\begin{itemize}
\item
  (ক) ৬০০ টাকা
\item
  (খ) ৬৫০ টাকা
\item
  (গ) ৭০০ টাকা
\item
  (ঘ) ৭৫০ টাকা
\end{itemize}

\textbf{\hl{উত্তরঃ (ঘ) ৭৫০ টাকা}}

\textbf{সমাধান:}

ত্রিভুজের ক্ষেত্রফল = \(\frac{\text{১}}{\text{২}}\)× ভূমি × উচ্চতা

= \(\frac{\text{১}}{\text{২}}\) × ৫০ × ২০ বর্গমিটার

= ৫০০ বর্গমিটার

১ বর্গমিটারে ঘাস লাঘাতে খরচ হয় = ১.৫০ টাকা

৫০০ বর্গমিটারে ঘাস লাঘাতে খরচ হয় = ৫০০ ×১.৫০ = ৭৫০ টাকা।

প্রশ্ন নম্বর-২. একটি ত্রিভুজাকৃতি জমির ভূমির পরিমাণ ৮০ মিটার এবং উচ্চতার পরিমাপ
৪৫ মিটার। জমির ক্ষেত্রফল হবে - {[}পররাষ্ট্র মন্ত্রণালয়ের প্রশাসনিক কর্মকর্তা
-২০০৪{]}

\begin{itemize}
\item
  (ক) ১৮০০ বর্গমিটার
\item
  (খ) ৯০০ বর্গমিটার
\item
  (গ) ৩৬০ বর্গমিটার
\item
  (ঘ) ২৪০০ বর্গমিটার
\end{itemize}

\textbf{\hl{উত্তরঃ (ক) ১৮০০ বর্গমিটার}}

\textbf{সমাধান:}

ত্রিভুজের ক্ষেত্রফল = \(\frac{\text{১}}{\text{২}}\)× ভূমি × উচ্চতা

= \(\frac{\text{১}}{\text{২}}\) × ৮০ × ৪৫ বর্গমিটার

= ১৮০০ বর্গমিটার

প্রশ্ন নম্বর-৩. একটি ত্রিভুজাকৃতি ক্ষেত্রের ক্ষেত্রফল ৮৪ বর্গগজ। ত্রিভুজটির শীর্ষ বিন্দু
হতে ভুমির উপর অংকিত লম্বের দৈর্ঘ্য ১২ গজ হলে ভূমির দৈর্ঘ্য কত ? {[}পল্লী
বিদ্যুতায়ন সহকারী সচিব-২০১৩{]}

\begin{itemize}
\item
  (ক) ১১
\item
  (খ) ১২
\item
  (গ) ১৩
\item
  (ঘ) ১৪
\end{itemize}

\textbf{\hl{উত্তরঃ (ঘ) ১৪}}

\textbf{সমাধান:}

ভূমির দৈর্ঘ্য = (ক্ষেত্রফল×২)÷উচ্চতা

= (৮৪×২)÷১২ = ১৬৮÷১২ = ১৪ গজ

প্রশ্ন নম্বর-৪. একটি ত্রিভজের ভূমির পরিমাণ ৪ মিটার ও উচ্চতা ৩ মিটার।ক্ষেত্রফল কত
? {[}জাতীয় সঞ্চয় পরিদপ্তরের সহকারী পরিচালক-২০০৯{]}

\begin{itemize}
\item
  (ক) ৬ বর্গমিটার
\item
  (খ) ৯ বর্গমিটার
\item
  (গ) ১২ বর্গমিটার
\item
  (ঘ) ১৮ বর্গমিটার
\end{itemize}

\textbf{\hl{উত্তরঃ (ক) ৬ বর্গমিটার}}

\textbf{সমাধান:}

ত্রিভুজের ক্ষেত্রফল = \(\frac{\text{১}}{\text{২}}\)× ভূমি × উচ্চতা

= \(\frac{\text{১}}{\text{২}}\) × ৪ × ৩ বর্গমিটার

= ৬ বর্গমিটার

\paragraph{সূত্র-৫: আয়তক্ষেত্রের দৈর্ঘ্য প্রস্থের যতগুণ এবং ক্ষেত্রফল দেওয়া থাকলে
দৈর্ঘ্য/প্রস্থ/ পরিসীমা বের করতে- দৈর্ঘ্য = √((ক্ষেত্রফল × দৈর্ঘ্য প্রস্থের যতগুণ))
প্রস্থ = √((ক্ষেত্রফল ÷ দৈর্ঘ্য প্রস্থের যতগুণ)) পরিসীমা = ২ (দৈর্ঘ্য + প্রস্থ
)}\label{ux9b8ux9a4ux9b0-ux9eb-ux986ux9dfux9a4ux995ux9b7ux9a4ux9b0ux9b0-ux9a6ux9b0ux998ux9af-ux9aaux9b0ux9b8ux9a5ux9b0-ux9afux9a4ux997ux9a3-ux98fux9ac-ux995ux9b7ux9a4ux9b0ux9abux9b2-ux9a6ux993ux9df-ux9a5ux995ux9b2-ux9a6ux9b0ux998ux9afux9aaux9b0ux9b8ux9a5-ux9aaux9b0ux9b8ux9ae-ux9acux9b0-ux995ux9b0ux9a4--ux9a6ux9b0ux998ux9af-ux995ux9b7ux9a4ux9b0ux9abux9b2-ux9a6ux9b0ux998ux9af-ux9aaux9b0ux9b8ux9a5ux9b0-ux9afux9a4ux997ux9a3-ux9aaux9b0ux9b8ux9a5-ux995ux9b7ux9a4ux9b0ux9abux9b2-ux9a6ux9b0ux998ux9af-ux9aaux9b0ux9b8ux9a5ux9b0-ux9afux9a4ux997ux9a3-ux9aaux9b0ux9b8ux9ae-ux9e8-ux9a6ux9b0ux998ux9af-ux9aaux9b0ux9b8ux9a5}

প্রশ্ন নম্বর- ১. একটি আয়তকার ঘরের দৈর্ঘ্য বিস্তারের দ্বিগুণ। এর ক্ষেত্রফল ৫১২
বর্গমিটার হলে, পরিসীমা কত ?{[}২৫ তম বিসিএস{]}

\begin{itemize}
\item
  (ক) ৯৬ মিটার
\item
  (খ) ২৫৬ মিটার
\item
  (গ) ১২৮ মিটার
\item
  (ঘ) ৪৮ মিটার
\end{itemize}

\textbf{\hl{উত্তরঃ (ক) ৯৬ মিটার}}

\textbf{সমাধান:}

শর্ট টেকনিক:

দৈর্ঘ্য = \(\sqrt{(\text{ক্ষেত্রফল}\  \times \ \text{দৈর্ঘ্য}\ \text{প্রস্থের}\ \text{যতগুণ})}\)

= \(\sqrt{\text{৫১২}\  \times \ \text{২}}\) = \(\sqrt{\text{১০২৪}}\) = ৩২ মিটার

প্রস্থ = \(\sqrt{(\text{ক্ষেত্রফল}\  \div \ \text{দৈর্ঘ্য}\ \text{প্রস্থের}\ \text{যতগুণ})}\)

= \(\sqrt{\text{৫১২}\  \div \text{২}}\) = \(\sqrt{\text{২৫৬}}\) = ১৬ মিটার

পরিসীমা = ২ (দৈর্ঘ্য + প্রস্থ )

= ২ ( ৩২ + ১৬ ) = ২×৪৮ = ৯৬ মিটার

প্রশ্ন নম্বর- ২. একটি আয়তক্ষেত্রের দৈর্ঘ্য প্রস্থের দেড়গুণ। ক্ষেত্রটির ক্ষেত্রফল ৯৬
বর্গমিটার হলে,ক্ষেত্রটির দৈর্ঘ্য কত ?{[}শ্রম অফিসার -১৯৯৫{]}

\begin{itemize}
\item
  (ক) ১২ মিটার
\item
  (খ) ৬ মিটার
\item
  (গ) ৮ মিটার
\item
  (ঘ) ১৬মিটার
\end{itemize}

\textbf{\hl{উত্তরঃ (ক) ১২ মিটার}}

\textbf{সমাধান:}

শর্ট টেকনিক:

দৈর্ঘ্য = \(\sqrt{(\text{ক্ষেত্রফল}\  \times \ \text{দৈর্ঘ্য}\ \text{প্রস্থের}\ \text{যতগুণ})}\)

= \(\sqrt{\text{৯৬}\  \times \ \text{১}.\text{৫}}\) = \(\sqrt{\text{১৪৪}}\) = ১২ মিটার

প্রশ্ন নম্বর- ৩. একটি আয়তক্ষেত্রের দৈর্ঘ্য প্রস্থের দ্বিগুণ। আয়তক্ষেত্রটির ক্ষেত্রফল
১২৫০ বর্গমিটার হলে এর দৈর্ঘ্য কত ?{[}৩০ তম বিসিএস{]}

\begin{itemize}
\item
  (ক) ৩০ মিটার
\item
  (খ) ৪০ মিটার
\item
  (গ) ৫০ মিটার
\item
  (ঘ) ৬০ মিটার
\end{itemize}

\textbf{\hl{উত্তরঃ (গ) ৫০ মিটার}}

\textbf{সমাধান:}

শর্ট টেকনিক:

দৈর্ঘ্য = \(\sqrt{(\text{ক্ষেত্রফল}\  \times \ \text{দৈর্ঘ্য}\ \text{প্রস্থের}\ \text{যতগুণ})}\)

= \(\sqrt{\text{১২৫০}\  \times \ \text{২}}\) = \(\sqrt{\text{২৫০০}}\) = ৫০ মিটার

প্রশ্ন নম্বর- ৪. একটি ক্ষেত্রফলের দৈর্ঘ্য তার প্রস্থের ১.৫ গুণ। ঐ ক্ষেত্রফল ২১৬
বর্গমিটার হলে তার দৈর্ঘ্য কত ?{[}থানা পরিদর্শক পরীক্ষা-২০০৯{]}

\begin{itemize}
\item
  (ক) ৯
\item
  (খ) ১২
\item
  (গ) ১৮
\item
  (ঘ) কোনটিই নয়
\end{itemize}

\textbf{\hl{উত্তরঃ (গ) ১৮}}

\textbf{সমাধান:}

শর্ট টেকনিক:

দৈর্ঘ্য = \(\sqrt{(\text{ক্ষেত্রফল}\  \times \ \text{দৈর্ঘ্য}\ \text{প্রস্থের}\ \text{যতগুণ})}\)

= \(\sqrt{\text{২১৬}\  \times \ \text{১}.\text{৫}}\) = \(\sqrt{\text{৩২৪}}\) = ১৮ মিটার

প্রশ্ন নম্বর- ৫. একটি আয়তাকার ঘরের দৈর্ঘ্য বিস্তারের দেড়গুণ। এর ক্ষেত্রফল ২১৬
বর্গমিটার হলে, তার পরিসীমা কত ? {[}শ্রম ও কর্মসংস্থান মন্ত্রণালয়ের সহকারী
প্রধান পরিদর্শক-২০০৯{]}

\begin{itemize}
\item
  (ক) ৩০ মি.
\item
  (খ) ৪০ মি.
\item
  (গ) ৫০ মি.
\item
  (ঘ) ৬০ মি.
\end{itemize}

\textbf{\hl{উত্তরঃ (ঘ) ৬০ মি.}}

\textbf{সমাধান:}

শর্ট টেকনিক:

দৈর্ঘ্য = \(\sqrt{(\text{ক্ষেত্রফল}\  \times \ \text{দৈর্ঘ্য}\ \text{প্রস্থের}\ \text{যতগুণ})}\)

= \(\sqrt{\text{২১৬}\  \times \ \text{১}.\text{৫}}\) = \(\sqrt{\text{৩২৪}}\) = ১৮ মিটার

প্রস্থ = \(\sqrt{(\text{ক্ষেত্রফল}\  \div \ \text{দৈর্ঘ্য}\ \text{প্রস্থের}\ \text{যতগুণ})}\)

= \(\sqrt{\text{২১৬}\  \div \ \text{১}.\text{৫}}\) = \(\sqrt{\text{১৪৪}}\) = ১২ মিটার

পরিসীমা = ২ (দৈর্ঘ্য + প্রস্থ )

= ২ ( ১৮ + ১২ ) = ২×৩০ = ৬০ মিটার

প্রশ্ন নম্বর- ৬. একটি আয়তক্ষেত্রের দৈর্ঘ্য প্রস্থের ৩ গুণ। আয়তক্ষেত্রটির ক্ষেত্রফল ৩০০
বর্গমিটার হলে উহার পরিসীমা কত ?{[} ২৪ তম বিসিএস{]}

\begin{itemize}
\item
  (ক) ৭০মিটার
\item
  (খ) ৭৫ মিটার
\item
  (গ) ৮০ মিটার
\item
  (ঘ) ৯০ মিটার
\end{itemize}

\textbf{\hl{উত্তরঃ (গ) ৮০ মিটার}}

\textbf{সমাধান:}

শর্ট টেকনিক:

দৈর্ঘ্য = \(\sqrt{(\text{ক্ষেত্রফল}\  \times \ \text{দৈর্ঘ্য}\ \text{প্রস্থের}\ \text{যতগুণ})}\)

= \(\sqrt{\text{৩০০}\  \times \ \text{৩}}\) = \(\sqrt{\text{৯০০}}\) = ৩০ মিটার

প্রস্থ = \(\sqrt{(\ \text{ক্ষেত্রফল}\  \div \ \text{দৈর্ঘ্য}\ \text{প্রস্থের}\ \text{যতগুণ})}\)

= \(\sqrt{\text{৩০০}\  \div \ \text{৩}}\) = \(\sqrt{\text{১০০}}\) = ১০ মিটার

পরিসীমা = ২ (দৈর্ঘ্য + প্রস্থ )

= ২ ( ৩০ + ১০ ) = ২×৪০ = ৮০ মিটার

প্রশ্ন নম্বর- ৭. একটি আয়তকার জমির দৈর্ঘ্য প্রস্থের দ্বিগুণ এবং উহার ক্ষেত্রফল ৮০০
বর্গমিটার হলে, জমিটির পরিসীমা কত মিটার ?{[}রেলওয়ে উপ সহকারী -২০১৩{]}

\begin{itemize}
\item
  (ক) ১০০
\item
  (খ) ১৪০
\item
  (গ) ১২০
\item
  (ঘ) ১১৫
\end{itemize}

\textbf{উত্তরঃ}

\textbf{সমাধান:}

শর্ট টেকনিক:

দৈর্ঘ্য = \(\sqrt{(\ \text{ক্ষেত্রফল}\  \times \ \text{দৈর্ঘ্য}\ \text{প্রস্থের}\ \text{যতগুণ})}\)

= \(\sqrt{\text{৮০০}\  \times \ \text{২}}\) = \(\sqrt{\text{১৬০০}}\) = ৪০ মিটার

প্রস্থ = \(\sqrt{(\text{ক্ষেত্রফল}\  \div \ \text{দৈর্ঘ্য}\ \text{প্রস্থের}\ \text{যতগুণ})}\)

= \(\sqrt{\text{৮০০}\  \div \ \text{২}}\) = \(\sqrt{\text{৪০০}}\) = ২০ মিটার

পরিসীমা = ২ (দৈর্ঘ্য + প্রস্থ )

= ২ ( ৪০ + ২০ ) = ২×৬০ = ১২০ মিটার

\section{46. পরিমাপ ও একক MCQ প্রশ্ন উত্তর
অনুশীলন}\label{ux9aaux9b0ux9aeux9aa-ux993-ux98fux995ux995-mcq-ux9aaux9b0ux9b6ux9a8-ux989ux9a4ux9a4ux9b0-ux985ux9a8ux9b6ux9b2ux9a8}

\textbf{আন্তর্জাতিক পদ্ধতিতে ওজন স্তরের একক কী?}

ক) সেন্টিগ্রাম

\textbf{\hl{খ) কেজি ✔}}

গ) কুইন্টাল

ঘ) গ্রাম

\textbf{\hl{✔ সঠিক উত্তর: খ) কেজি}}

\begin{quote}
\textbf{ব্যাখ্যা:} \emph{আন্তর্জাতিক পদ্ধতি বা SI (Système International
d\textquotesingle Unités) অনুযায়ী ভরের বা ওজনের মূল একক হলো কিলোগ্রাম
(কেজি)।}
\end{quote}

\textbf{আন্তর্জাতিক একক পদ্ধতি চালু হয় কত সালে?}

ক) ১৯৭০ সালে

খ) ১৯৫৮ সালে

\textbf{\hl{গ) ১৯৬০ সালে ✔}}

ঘ) ১৯৬৩ সালে

\textbf{\hl{✔ সঠিক উত্তর: গ) ১৯৬০ সালে}}

\begin{quote}
\textbf{ব্যাখ্যা:} \emph{১৯৬০ সালে অনুষ্ঠিত \textquotesingle Weight and
Measures\textquotesingle{} সংক্রান্ত ১১তম সাধারণ সভায় পরিমাপের আন্তর্জাতিক
পদ্ধতি বা SI (Système International d\textquotesingle Unités) চালু করা
হয়।}
\end{quote}

\textbf{C.G.S পদ্ধতিতে দৈর্ঘ্যের একক}

ক) ডেকামিটার (Decametre)

খ) মিটার (Metre)

গ) ডেসিমিটার(Decimetre)

\textbf{\hl{ঘ) সেন্টিমিটার (Centimetre) ✔}}

\textbf{\hl{✔ সঠিক উত্তর: ঘ) সেন্টিমিটার (Centimetre)}}

\begin{quote}
\textbf{ব্যাখ্যা:} \emph{C.G.S (Centimetre-Gram-Second) পদ্ধতিতে দৈর্ঘ্যের
একক হলো সেন্টিমিটার (Centimetre)।}
\end{quote}

\textbf{C.G.S পদ্ধতিতে ভরের একক-\/-}

ক) পাউন্ড (Pound)

\textbf{\hl{খ) গ্রাম (Gram) ✔}}

গ) কিলোগ্রাম (Kilogram)

ঘ) আউন্স (Ounce)

\textbf{\hl{✔ সঠিক উত্তর: খ) গ্রাম (Gram)}}

\begin{quote}
\textbf{ব্যাখ্যা:} \emph{C.G.S (Centimetre-Gram-Second) পদ্ধতিতে ভরের
পরিমাপের একক হলো গ্রাম (Gram)।}
\end{quote}

\textbf{মেট্রিক পদ্ধতিতে ভরের একক/ পদার্থের ভর পরিমাপের আন্তর্জাতিক একক -\/-}

ক) পাউন্ড (Pound)

খ) গ্রাম (Gram)

\textbf{\hl{গ) কিলোগ্রাম (Kilogram) ✔}}

ঘ) মিলিগ্রাম (Milligram)

\textbf{\hl{✔ সঠিক উত্তর: গ) কিলোগ্রাম (Kilogram)}}

\begin{quote}
\textbf{ব্যাখ্যা:} \emph{মেট্রিক পদ্ধতি বা আন্তর্জাতিক SI পদ্ধতি উভয় ক্ষেত্রেই
পদার্থের ভর পরিমাপের মূল বা মৌলিক একক হলো কিলোগ্রাম (Kilogram)।}
\end{quote}

\textbf{১ মাইল = কত?}

ক) ১,৬৭০ গজ

\textbf{\hl{খ) ১,৭৬০ গজ ✔}}

গ) ১,৬৩০ গজ

ঘ) ১,৩৭০ গজ

\textbf{\hl{✔ সঠিক উত্তর: খ) ১,৭৬০ গজ}}

\begin{quote}
\textbf{ব্যাখ্যা:} \emph{১ মাইল = ১,৭৬০ গজ}

\emph{= (১,৭৬০ × ৩) ফুট = ৫,২৮০ ফুট}
\end{quote}

\textbf{১ মিটার = কত ইঞ্চি?}

ক) ৩৭.৩৭

খ) ৩৮.৩৭

\textbf{\hl{গ) ৩৯.৩৭ ✔}}

ঘ) ৪০.৩৭

\textbf{\hl{✔ সঠিক উত্তর: গ) ৩৯.৩৭}}

\begin{quote}
\textbf{ব্যাখ্যা:} \emph{১ মিটার = ৩৯.৩৭ ইঞ্চি (প্রায়)।}
\end{quote}

\textbf{এক নটিক্যাল মাইল সমান---}

ক) ১৮৫৩ মিটার

খ) ১০০০ মিটার

গ) ৯৬০.১৮ মিটার

\textbf{\hl{ঘ) ১৮৫৩.১৮ মিটার ✔}}

\textbf{\hl{✔ সঠিক উত্তর: ঘ) ১৮৫৩.১৮ মিটার}}

\begin{quote}
\textbf{ব্যাখ্যা:} \emph{১ নটিক্যাল মাইল = ১.৮৫৩১-১৮৫৩.১৮ মিটার (প্রায়)।}
\end{quote}

\textbf{১ কিলোগ্রাম = কত পাউন্ড?}

ক) ১২

খ) ০.২২

\textbf{\hl{গ) ২.২ ✔}}

ঘ) ২.৪৭

\textbf{\hl{✔ সঠিক উত্তর: গ) ২.২}}

\begin{quote}
\textbf{ব্যাখ্যা:} \emph{১ কিলোগ্রাম ≈ ২.২০৪৬২ পাউন্ড, যা প্রায় ২.২ পাউন্ড।}
\end{quote}

\textbf{১ নটিক্যাল মাইল = কত কিমি?}

ক) ০.৬২১

খ) ১.৪৭২

\textbf{\hl{গ) ১.৮৫২ ✔}}

ঘ) ২.২৫৪

\textbf{\hl{✔ সঠিক উত্তর: গ) ১.৮৫২}}

\begin{quote}
\textbf{ব্যাখ্যা:} \emph{১ নটিক্যাল মাইল = ১.১৫ মাইল = ১.৮৫২ কি.মি.
(প্রায়)।}
\end{quote}

\textbf{১ ক্যারেট = কত গ্রাম?}

ক) ৫ গ্রাম

খ) ১০ গ্রাম

গ) ১৮ গ্রাম

ঘ) ২ গ্রাম

\textbf{\hl{ঙ) সঠিক উত্তর নেই ✔}}

\textbf{\hl{✔ সঠিক উত্তর: ঙ) সঠিক উত্তর নেই}}

\begin{quote}
\textbf{ব্যাখ্যা:} \emph{১ ক্যারেট = ০.২ গ্রাম। প্রদত্ত অপশনগুলোর কোনোটিতেই
সঠিক উত্তর নেই (সঠিক উত্তর ০.২ গ্রাম)।}
\end{quote}

\textbf{১০ মিলিমিটার =?}

ক) ১ ডেসিমিটার

\textbf{\hl{খ) ১ সেন্টিমিটার ✔}}

গ) ১ কিলোমিটার

ঘ) ১ মিটার

\textbf{\hl{✔ সঠিক উত্তর: খ) ১ সেন্টিমিটার}}

\begin{quote}
\textbf{ব্যাখ্যা:} \emph{১০ মিলিমিটার = ১ সেন্টিমিটার।}
\end{quote}

\textbf{৫ মিটার, ১৫ সেন্টিমিটার ও ১৫ মিলিমিটার সমান কত মিলিমিটার।}

\textbf{\hl{ক) ৫১৬৫ ✔}}

খ) ৫৭৭৫

গ) ৬৫১৫

ঘ) ৬৯১৫

\textbf{\hl{✔ সঠিক উত্তর: ক) ৫১৬৫}}

\begin{quote}
\textbf{ব্যাখ্যা:} \emph{১ মিটার = ১০০ সেন্টিমিটার = ১০০০০০ মিলিমিটার এবং ১
সেন্টিমিটার = ১০ মিলিমিটার।}

\emph{∴ ৫ মিটার = (৫ × ১০০০) মিলিমিটার = ৫০০০ মিলিমিটার}

\emph{১৫ সেন্টিমিটার = (১৫ × ১০) মিলিমিটার = ১৫০ মিলিমিটার}

\emph{মোট = ৫০০০ + ১৫০ + ১৫ = ৫১৬৫ মিলিমিটার।}
\end{quote}

\textbf{১ ইঞ্চিতে কত সেন্টিমিটার?}

\textbf{\hl{ক) ২.৫৪ সেন্টিমিটার ✔}}

খ) ২.০২ সেন্টিমিটার

গ) ৩.৭৩ সেন্টিমিটার

ঘ) ১.৩৭ সেন্টিমিটার

\textbf{\hl{✔ সঠিক উত্তর: ক) ২.৫৪ সেন্টিমিটার}}

\begin{quote}
\textbf{ব্যাখ্যা:} \emph{আমরা জানি, ১ ইঞ্চি = ২.৫৪ সেন্টিমিটার।}
\end{quote}

\textbf{শিহাব ১৬০০০ মিটার দৌড়ালে, সে কত কিলোমিটার দৌড়েছিল?}

ক) ১.৬০ কিমি

খ) ১৬০ কিমি

গ) .১৬ কিমি

\textbf{\hl{ঘ) ১৬ কিমি ✔}}

\textbf{\hl{✔ সঠিক উত্তর: ঘ) ১৬ কিমি}}

\begin{quote}
\textbf{ব্যাখ্যা:} \emph{১০০০ মিটার = ১ কিলোমিটার।}

\emph{∴ ১৬০০০ মিটার = (১৬০০০ ÷ ১০০০) কিমি = ১৬ কিমি।}
\end{quote}

\textbf{রতন ১৬০০০ মিটার দৌড়ালে সে কত কিমি পথ পাড়ি দিয়েছে?}

ক) ১৬০

খ) ১১

গ) ১২

\textbf{\hl{ঘ) ১৬ ✔}}

\textbf{\hl{✔ সঠিক উত্তর: ঘ) ১৬}}

\begin{quote}
\textbf{ব্যাখ্যা:} \emph{আমরা জানি, ১০০০ মিটার = ১ কিমি}

\emph{∴ রতন ১৬০০০ মিটার যায় = (১×১৬০০০/১০০০) কিমি = ১৬ কিমি।}
\end{quote}

\textbf{এক মিলিমিটার এক কিলোমিটারের কত অংশ?}

ক) (১/১০০০)

\textbf{\hl{খ) (১/১০০০০০০) ✔}}

গ) (১/১০০০০০)

ঘ) (১/১০০০০)

\textbf{\hl{✔ সঠিক উত্তর: খ) (১/১০০০০০০)}}

\begin{quote}
\textbf{ব্যাখ্যা:} \emph{১ কিলোমিটার = ১০,০০,০০০ মিলিমিটার।}

\emph{∴ ১ মিলিমিটার হলো ১ কিলোমিটারের (১/১০,০০,০০০) অংশ।}
\end{quote}

\textbf{১০ মিলিমিটার সমান---}

ক) ১ ডেমি

\textbf{\hl{খ) ১ সেমি ✔}}

গ) ১ মিটার

ঘ) ১ সে. গ্রাম

\textbf{\hl{✔ সঠিক উত্তর: খ) ১ সেমি}}

\begin{quote}
\textbf{ব্যাখ্যা:} \emph{১০ মিলিমিটার = ১ সেন্টিমিটার (সেমি)।}
\end{quote}

\textbf{One mile is equivalent to how many kilometers? Or, ১ মাইলে কত
কিলোমিটার?}

ক) 1 kilometer

\textbf{\hl{খ) 1.6 kilometers ✔}}

গ) 1.8 kilometers

ঘ) 2 kilometers

\textbf{\hl{✔ সঠিক উত্তর: খ) 1.6 kilometers}}

\begin{quote}
\textbf{ব্যাখ্যা:} \emph{আমরা জানি, ১ মাইল = ১.৬০৯৩৪ কিলোমিটার ≈ ১.৬
কিলোমিটার (1.6 kilometers)।}
\end{quote}

\textbf{এক কিলোমিটার কত মাইলের সমান?}

\textbf{\hl{ক) (৫/৮) ✔}}

খ) (১/২)

গ) (৫/৪)

ঘ) (৩/৪)

\textbf{\hl{✔ সঠিক উত্তর: ক) (৫/৮)}}

\begin{quote}
\textbf{ব্যাখ্যা:} \emph{১ কিলোমিটার = ০.৬২ মাইল (প্রায়)।}

\emph{ভগ্নাংশে প্রকাশ করলে:}

\emph{০.৬২ = (৬২/১০০) ≈ (৫/৮) মাইল।}
\end{quote}

\textbf{এক সেন্টিমিটার সমান ০.৩৯৩৭ ইঞ্চি হলে এক কিলোমিটার সমান কত?}

\textbf{\hl{ক) ৩৯৩৭০ ইঞ্চি ✔}}

খ) ৩৯.৩৭ ইঞ্চি

গ) ৩৯৩৭ইঞ্চি

ঘ) ৩০.০০৩৯৩৭ইঞ্চি

\textbf{\hl{✔ সঠিক উত্তর: ক) ৩৯৩৭০ ইঞ্চি}}

\begin{quote}
\textbf{ব্যাখ্যা:} \emph{১ সেন্টিমিটার = ০.৩৯৩৭ ইঞ্চি}

\emph{১০০ সেমি: বা ১ মিটার = (০.৩৯৩৭ × ১০০) ইঞ্চি}

\emph{১০০০ মিটার বা ১ কিমি: = (০.৩৯৩৭ × ১০০ × ১০০০) ইঞ্চি = ৩৯৩৭০ ইঞ্চি।}
\end{quote}

\textbf{১ কি.মি. সমান কত মাইল?}

ক) ১.৬২ মাইল

খ) ০.৬৩ মাইল

গ) ০.৫৮ মাইল

\textbf{\hl{ঘ) ০.৬২ মাইল ✔}}

\textbf{\hl{✔ সঠিক উত্তর: ঘ) ০.৬২ মাইল}}

\begin{quote}
\textbf{ব্যাখ্যা:} \emph{আমরা জানি, ১ কিলোমিটার = ০.৬২১৩৭ মাইল ≈ ০.৬২
মাইল।}
\end{quote}

\textbf{সমুদ্রে পানির গভীরতা মাপার একক---}

ক) মিটার (Mile)

\textbf{\hl{খ) ফ্যাদম (Fathom) ✔}}

গ) কিলোমিটার (kilometre)

ঘ) ফুট (Foot)

\textbf{\hl{✔ সঠিক উত্তর: খ) ফ্যাদম (Fathom)}}

\begin{quote}
\textbf{ব্যাখ্যা:} \emph{সমুদ্রে পানির গভীরতা পরিমাপের একক হলো ফ্যাদম
(Fathom)।}

\emph{(১ ফ্যাদম = ৬ ফুট)।}
\end{quote}

\textbf{১(১/২) ইঞ্চি ১ ফুটের কত অংশ?}

ক) (১/৪)

খ) (১/২)

\textbf{\hl{গ) (১/৮) ✔}}

ঘ) (১/১০)

\textbf{\hl{✔ সঠিক উত্তর: গ) (১/৮)}}

\begin{quote}
\textbf{ব্যাখ্যা:} \emph{আমরা জানি, ১২ ইঞ্চি = ১ ফুট}

\emph{∴ ১ ইঞ্চি = (১/১২) ফুট}

\emph{∴ দেড় বা (৩/২) ইঞ্চি = (১×৩/১২×২) ফুট = (৩/২৪) ফুট = (১/৮) ফুট।}
\end{quote}

\textbf{৩৫২ গজ ১ মাইলের কত অংশ?}

ক) (১/৩)

খ) (১/৪)

\textbf{\hl{গ) (১/৫) ✔}}

ঘ) (১/৮)

\textbf{\hl{✔ সঠিক উত্তর: গ) (১/৫)}}

\begin{quote}
\textbf{ব্যাখ্যা:} \emph{আমরা জানি, ১৭৬০ গজ = ১ মাইল}

\emph{∴ ১ গজ = (১/১৭৬০) মাইল}

\emph{∴ ৩৫২ গজ = (৩৫২/১৭৬০) মাইল = (১/৫) মাইল।}
\end{quote}

\textbf{৩০০ কিলোমিটার দূরত্বকে মাইলে প্রকাশ করলে হয়---}

ক) ১৫৬ মাইল

খ) ১৭৬ মাইল

গ) ১৬৬ মাইল

\textbf{\hl{ঘ) ১৮৬ মাইল ✔}}

\textbf{\hl{✔ সঠিক উত্তর: ঘ) ১৮৬ মাইল}}

\begin{quote}
\textbf{ব্যাখ্যা:} \emph{আমরা জানি, ১ কিলোমিটার = ০.৬২ মাইল (প্রায়)}

\emph{∴ ৩০০ কিলোমিটার = (০.৬২ × ৩০০) মাইল = ১৮৬ মাইল (বা ১৮৬.৩৩ মাইল)।}
\end{quote}

\textbf{A can run 22.5 m while B runs 25 m. In a kilometre race B beats
A by:/ A দৌড়াতে পারে ২২.৫ মিটার যেখানে B দৌড়াতে পারে ২৫ মিটার। কিলোমিটার
রেসে B হারাতে পারে A কে -}

\textbf{\hl{ক) 100m ✔}}

খ) 111(1/9) m

গ) 25m

ঘ) 50m

\textbf{\hl{✔ সঠিক উত্তর: ক) 100m}}

\begin{quote}
\textbf{ব্যাখ্যা:} \emph{B যখন ২৫ মিটার দৌড়ায়, তখন A কে হারায় (২৫ − ২২.৫)
মি. = ২.৫ মিটারে।}

\emph{১ কিমি = ১০০০০ মিটার।}

\emph{∴ B ১ মিটারে A কে হারায় (২.৫/২৫) মিটারে।}

\emph{∴ B ১০০০০ মিটারে A কে হারায় (২.৫×১০০০০/২৫) মিটারে = ১০০ মিটারে।}
\end{quote}

\textbf{এক বর্গমাইল এর সমান নিচের কোনটি?}

ক) ৬০৪ একর

খ) ৬২০ একর

\textbf{\hl{গ) ৬৪০ একর ✔}}

ঘ) ৬৬০ একর

\textbf{\hl{✔ সঠিক উত্তর: গ) ৬৪০ একর}}

\begin{quote}
\textbf{ব্যাখ্যা:} \emph{আমরা জানি, ১ বর্গমাইল = ৬৪০ একর।}
\end{quote}

\textbf{এক বর্গইঞ্চি কত বর্গ সেন্টিমিটারের সমান?}

ক) ০.০৯২৯

খ) ৭.৩০

\textbf{\hl{গ) ৬.৪৫ ✔}}

ঘ) ৫.৪৫

\textbf{\hl{✔ সঠিক উত্তর: গ) ৬.৪৫}}

\begin{quote}
\textbf{ব্যাখ্যা:} \emph{আমরা জানি, ১ ইঞ্চি = ২.৫৪ সেন্টিমিটার (সেমি)}

\emph{∴ ১ বর্গইঞ্চি = (২.৫৪ × ২.৫৪) বর্গ সেন্টিমিটার = ৬.৪৫১৬ বর্গ সেন্টিমিটার
≈ ৬.৪৫ বর্গ সেন্টিমিটার।}
\end{quote}

\textbf{১ বর্গ মিটার কত বর্গ সেন্টিমিটারের সমান?}

ক) ১০০

\textbf{\hl{খ) ১০০০০ ✔}}

গ) ১০০০

ঘ) ১০০০০০

\textbf{\hl{✔ সঠিক উত্তর: খ) ১০০০০}}

\begin{quote}
\textbf{ব্যাখ্যা:} \emph{১ বর্গমিটার = ১ মিটার × ১ মিটার}

\emph{= ১০০ সে.মি. × ১০০ সে.মি.}

\emph{= ১০০০০ বর্গ সে. মি.}
\end{quote}

\textbf{১ বর্গমিটার = কত বর্গফুট?}

\textbf{\hl{ক) ১০.৭৬ ✔}}

খ) ৯.৭৬

গ) ১০.১৪

ঘ) ৯.১৪

\textbf{\hl{✔ সঠিক উত্তর: ক) ১০.৭৬}}

\begin{quote}
\textbf{ব্যাখ্যা:} \emph{আমরা জানি, ১ মিটার = ৩.২৮ ফুট (প্রায়)}

\emph{∴ ১ বর্গমিটার = (৩.২৮ × ৩.২৮) বর্গফুট ≈ ১০.৭৬ বর্গফুট।}
\end{quote}

\textbf{এক বর্গকিলোমিটারের পরিমাপ-\/-/ 1 sq.km equals -\/-}

ক) ৩০০ একর (Acres)

খ) ৫৪০ একর (Acres)

গ) ৪০০ একর (Acres)

\textbf{\hl{ঘ) ২৪৭ একর (Acres) ✔}}

\textbf{\hl{✔ সঠিক উত্তর: ঘ) ২৪৭ একর (Acres)}}

\begin{quote}
\textbf{ব্যাখ্যা:} \emph{১ বর্গকিলোমিটার = ২৪৭.১০৫ একর ≈ ২৪৭ একর (Acres)।}
\end{quote}

\textbf{১ একর সমান কত বর্গফুট?/ Express 1 acre in sq.ft.}

ক) ১০০০

খ) ৪০০০০

\textbf{\hl{গ) ৪৩৫৬০ ✔}}

ঘ) ৪৮৪০

\textbf{\hl{✔ সঠিক উত্তর: গ) ৪৩৫৬০}}

\begin{quote}
\textbf{ব্যাখ্যা:} \emph{১ একর (Acre) = ৪৮৪০ বর্গগজ}

\emph{= (৪৮৪০ × ৯) বর্গফুট = ৪৩৫৬০ বর্গফুট।}
\end{quote}

\textbf{কত বর্গমিটার সমান ১ এয়র?}

ক) ১০০০০

খ) ১০০০

\textbf{\hl{গ) ১০০ ✔}}

ঘ) ১০

\textbf{\hl{✔ সঠিক উত্তর: গ) ১০০}}

\begin{quote}
\textbf{ব্যাখ্যা:} \emph{আমরা জানি, ১০০ বর্গমিটার = ১ এয়ার।}
\end{quote}

\textbf{২০ বর্গমিটার ২ এয়ারের কত অংশ?}

ক) (১/৫)

খ) (১/৮)

\textbf{\hl{গ) (১/১০) ✔}}

ঘ) (১/১২)

\textbf{\hl{✔ সঠিক উত্তর: গ) (১/১০)}}

\begin{quote}
\textbf{ব্যাখ্যা:} \emph{১ এয়ার = ১০০ বর্গমিটার}

\emph{∴ ২ এয়ার = (১০০ × ২) = ২০০ বর্গমিটার}

\emph{∴ ২০ বর্গমিটার, ২ এয়ারের = (২০/২০০) অংশ = (১/১০) অংশ।}
\end{quote}

\textbf{১ হেক্টর জমি বলতে বুঝায়-\/-}

\textbf{\hl{ক) ১০,০০০ বর্গমিটার ✔}}

খ) ১০০০ বর্গমিটার

গ) ১০০ বর্গমিটার

ঘ) ১০ বর্গমিটার

\textbf{\hl{✔ সঠিক উত্তর: ক) ১০,০০০ বর্গমিটার}}

\begin{quote}
\textbf{ব্যাখ্যা:} \emph{১ হেক্টর = ১০,০০০ বর্গমিটার (বা ১০০ এয়ার)।}
\end{quote}

\textbf{একটি জমির পরিমাণ ৫ কাঠা হলে, তা বর্গফুটে হবে-\/-}

ক) ৩০০০ বর্গফুট

খ) ৩২০০ বর্গফুট

\textbf{\hl{গ) ৩৬০০ বর্গফুট ✔}}

ঘ) ৪০২০ বর্গফুট

\textbf{\hl{✔ সঠিক উত্তর: গ) ৩৬০০ বর্গফুট}}

\begin{quote}
\textbf{ব্যাখ্যা:} \emph{আমরা জানি, ১ কাঠা = ৭২০ বর্গফুট}

\emph{∴ ৫ কাঠা = (৭২০ × ৫) বর্গফুট = ৩৬০০ বর্গফুট।}
\end{quote}

\textbf{১ ঘনফুট = কত লিটার?}

\textbf{\hl{ক) ২৮.৪ লিটার ✔}}

খ) ২৯.৪ লিটার

গ) ২৭.৪ লিটার

ঘ) ৩০.৪ লিটার

\textbf{\hl{✔ সঠিক উত্তর: ক) ২৮.৪ লিটার}}

\begin{quote}
\textbf{ব্যাখ্যা:} \emph{১ ঘনফুট = (০.৩০৫ × ০.৩০৫ × ০.৩০৫) ঘনমিটার}

\emph{= (০.০২৮৪ × ১০০০০০০) ঘন সে.মিটার}

\emph{= (০.০২৮৪×১০-০০০০০/১০০০) লিটার {[}১০০০ ঘন সেমি = ১ লিটার{]}}

\emph{∴ ১ ঘনফুট = ২৮.৪ লিটার (প্রায়)।}
\end{quote}

\textbf{১ ঘন মিটার = কত লিটার?}

ক) ১০০ লিটার

\textbf{\hl{খ) ১০০০ লিটার ✔}}

গ) ১০০০০ লিটার

ঘ) ১০-০০০০ লিটার

\textbf{\hl{✔ সঠিক উত্তর: খ) ১০০০ লিটার}}

\begin{quote}
\textbf{ব্যাখ্যা:} \emph{আমরা জানি, ১ ঘনমিটার = ১০০০ লিটার।}
\end{quote}

\textbf{কত কিউবিক সেন্টিমিটার (সি.সি) তে ১ লিটার হয়?}

ক) ১০

খ) ১০০

\textbf{\hl{গ) ১০০০ ✔}}

ঘ) ১০-০০-০

\textbf{\hl{✔ সঠিক উত্তর: গ) ১০০০}}

\begin{quote}
\textbf{ব্যাখ্যা:} \emph{আমরা জানি, ১০০০ কিউবিক সেন্টিমিটার (বা ঘন
সেন্টিমিটার/cc) = ১ লিটার।}
\end{quote}

\textbf{তরল পদার্থ মাপার একক কি?}

ক) টন (Ton)

খ) বুশেল (Bushel)

\textbf{\hl{গ) ব্যারেল (Barrel) ✔}}

ঘ) কুইন্টাল (Quintal)

\textbf{\hl{✔ সঠিক উত্তর: গ) ব্যারেল (Barrel)}}

\begin{quote}
\textbf{ব্যাখ্যা:} \emph{টন ও কুইন্টাল হলো ভরের একক; বুশেল কঠিন পদার্থের আয়তন
মাপার একক এবং ব্যারেল তরল পদার্থের আয়তন মাপার একক।}
\end{quote}

\textbf{১০০ গ্যালনে কত লিটার?}

ক) ৪০০ লিটার

খ) ৪২০ লিটার

\textbf{\hl{গ) ৪৫০ লিটার ✔}}

ঘ) ৪৫৫ লিটার

\textbf{\hl{✔ সঠিক উত্তর: গ) ৪৫০ লিটার}}

\begin{quote}
\textbf{ব্যাখ্যা:} \emph{আমরা জানি, ১ গ্যালন (ইউ কে/ব্রিটিশ) ≈ ৪.৫৪৬ (বা
৪.৫৫) লিটার।}

\emph{∴ ১০০ গ্যালন = (৪.৫৫ × ১০০) লিটার = ৪৫০ লিটার (বা ৪৫৫ লিটার)।}
\end{quote}

\textbf{এক গ্যালন = কত লিটার?}

ক) ৩.৫

খ) ৪

\textbf{\hl{গ) ৪.৫৫ ✔}}

ঘ) ৫

\textbf{\hl{✔ সঠিক উত্তর: গ) ৪.৫৫}}

\begin{quote}
\textbf{ব্যাখ্যা:} \emph{১ ব্রিটিশ গ্যালন = ৪.৫৫ লিটার (প্রায়)।}
\end{quote}

\textbf{এক গ্যালনে কত পাইন্ট? এবং ১ বর্গমাইল সমান কত একর?}

ক) এক গ্যালন = ৮ পাইন্ট, ১ বর্গমাইল = ৬৪০ একর

\textbf{\hl{খ) এক গ্যালন = ৬.৬৬ পাইন্ট, ১ বর্গমাইল = ৬৪০ একর ✔}}

গ) এক গ্যালন = ১০ পাইন্ট, ১ বর্গমাইল = ৫০০ একর

ঘ) এক গ্যালন = ৪ পাইন্ট, ১ বর্গমাইল = ৬৪০ একর

\textbf{\hl{✔ সঠিক উত্তর: খ) এক গ্যালন = ৬.৬৬ পাইন্ট, ১ বর্গমাইল = ৬৪০ একর}}

\begin{quote}
\textbf{ব্যাখ্যা:} \emph{১ ইউ.এস গ্যালন ≈ ৮ পাইন্ট (তরল) বা যুক্তরাজ্য মাপে ১
গ্যালন = ৮ পাইন্ট। ছবিতে উল্লেখিত সমাধান অনুযায়ী: এক গ্যালন = ৬.৬৬ পাইন্ট
(শুষ্ক) এবং ১ বর্গমাইল = ৬৪০ একর।}
\end{quote}

\textbf{এক ঘনমিটার পানির ট্যাংকে কত U.S গ্যালন পানি ধারণ করবে?}

\textbf{\hl{ক) ২৬৪ ✔}}

খ) ২২০

গ) ২৩৪

ঘ) ২৫০

\textbf{\hl{✔ সঠিক উত্তর: ক) ২৬৪}}

\begin{quote}
\textbf{ব্যাখ্যা:} \emph{১ U.S গ্যালন = ৩.৭৮৫ লিটার}

\emph{১ ঘনমিটার = ১০০০ লিটার}

\emph{∴ ১ ঘনমিটারে U.S গ্যালনের পরিমাণ = (১০০০/৩.৭৮৫) গ্যালন ≈ ২৬৪.২ গ্যালন
বা ২৬৪ গ্যালন।}
\end{quote}

\textbf{৪ ডিগ্রি সেন্টিগ্রেড তাপমাত্রায় এক লিটার পানির ওজন-/ At 4°c
temperature, what is the weight of 1 litre of water?}

ক) ১.৪ কেজি

খ) ১.২ কেজি

গ) ০.৯৬ কেজি

\textbf{\hl{ঘ) ১.০ কেজি ✔}}

\textbf{\hl{✔ সঠিক উত্তর: ঘ) ১.০ কেজি}}

\begin{quote}
\textbf{ব্যাখ্যা:} \emph{৪°C তাপমাত্রায় ১ লিটার বিশুদ্ধ পানির ওজন ঠিক ১
কিলোগ্রাম বা ১.০ কেজি।}
\end{quote}

\textbf{৩ লিটার পানির ওজন-/ What is the weight of 3 litres of water?}

ক) ২.৫ kg

খ) ২.৭৫kg

গ) ৪.০০ kg

\textbf{\hl{ঘ) ৩.০০kg ✔}}

\textbf{\hl{✔ সঠিক উত্তর: ঘ) ৩.০০kg}}

\begin{quote}
\textbf{ব্যাখ্যা:} \emph{আমরা জানি, ১ লিটার পানির ওজন = ১ কেজি}

\emph{∴ ৩ লিটার পানির ওজন = (১ × ৩) কেজি = ৩.০০ কেজি।}
\end{quote}

\textbf{পাঁচ লিটার পানির ওজন কত?/ What is the weight of 5 litres of
water?}

ক) ৫ গ্রাম

খ) ৫০০ গ্রাম

গ) ৫০ গ্রাম

\textbf{\hl{ঘ) ৫ কিলোগ্রাম ✔}}

\textbf{\hl{✔ সঠিক উত্তর: ঘ) ৫ কিলোগ্রাম}}

\begin{quote}
\textbf{ব্যাখ্যা:} \emph{আমরা জানি, ১ লিটার পানির ওজন = ১ কিলোগ্রাম
(কেজি)।}

\emph{∴ ৫ লিটার পানির ওজন = (১ × ৫) কিলোগ্রাম = ৫ কিলোগ্রাম।}
\end{quote}

\textbf{১ ঘন মিটার পানির ভর হবে-/ What is the weight of 1 cubic metre of
water?}

ক) ১ কেজি

খ) ১০ কেজি

গ) ১০০ কেজি

\textbf{\hl{ঘ) ১০০০ কেজি ✔}}

\textbf{\hl{✔ সঠিক উত্তর: ঘ) ১০০০ কেজি}}

\begin{quote}
\textbf{ব্যাখ্যা:} \emph{১ ঘন মিটার = ১০০০ লিটার।}

\emph{আমরা জানি, ৪°C তাপমাত্রায় ১ লিটার পানির ভর = ১ কেজি।}

\emph{∴ ১ ঘন মিটার পানির ভর = (১০০০ × ১) কেজি = ১০০০ কেজি।}
\end{quote}

\textbf{এক Cubic meter পানির ওজন কত?/ What is the weight of 1 cubic
metre of water?}

ক) ১০০ লিটার

\textbf{\hl{খ) ১০০০ লিটার ✔}}

গ) ২৫০ লিটার

ঘ) ৫০০ লিটার

\textbf{\hl{✔ সঠিক উত্তর: খ) ১০০০ লিটার}}

\begin{quote}
\textbf{ব্যাখ্যা:} \emph{১ ঘনমিটার বা Cubic meter = ১০০০ লিটার পানির ওজনের
সমতুল্য (যেহেতু ১ লিটার পানির ওজন ১ কেজি এবং ১ ঘনমিটার পানি = ১০০০ কেজি বা
১০০০ লিটার)।}
\end{quote}

\textbf{১ মেট্রিক টন = কত কিলোগ্রাম?}

ক) ১০ কি.গ্রা.

খ) ১০০ কি.গ্রা.

\textbf{\hl{গ) ১০০০ কি.গ্রা. ✔}}

ঘ) ১০,০০০ কি.গ্রা.

\textbf{\hl{✔ সঠিক উত্তর: গ) ১০০০ কি.গ্রা.}}

\begin{quote}
\textbf{ব্যাখ্যা:} \emph{আমরা জানি, ১ মেট্রিক টন = ১০০০ কিলোগ্রাম
(কি.গ্রা.)।}
\end{quote}

\textbf{১ টন = কত পাউন্ড?}

ক) ১০-০০

খ) ১১৬.৮

গ) ১৪০০

\textbf{\hl{ঘ) ২২৪০ ✔}}

\textbf{\hl{✔ সঠিক উত্তর: ঘ) ২২৪০}}

\begin{quote}
\textbf{ব্যাখ্যা:} \emph{আমরা জানি, ১ টন (লং টন) = ২২৪০ পাউন্ড।}
\end{quote}

\textbf{১ পাউন্ড কত আউন্স?}

ক) ১৪ আউন্স

খ) ১২ আউন্স

গ) ১৮ আউন্স

\textbf{\hl{ঘ) ১৬ আউন্স ✔}}

\textbf{\hl{✔ সঠিক উত্তর: ঘ) ১৬ আউন্স}}

\begin{quote}
\textbf{ব্যাখ্যা:} \emph{আমরা জানি, ১ পাউন্ড = ১৬ আউন্স।}
\end{quote}

\textbf{এক সের সমান কত কিলোগ্রাম?/ Express 1 seer in kilogram?}

ক) ০.৯৭ কিলোগ্রাম (প্রায়)

\textbf{\hl{খ) ০.৯৩ কিলোগ্রাম (প্রায়) ✔}}

গ) ১.০৭ কিলোগ্রাম (প্রায়)

ঘ) ১.০৯ কিলোগ্রাম (প্রায়)

\textbf{\hl{✔ সঠিক উত্তর: খ) ০.৯৩ কিলোগ্রাম (প্রায়)}}

\begin{quote}
\textbf{ব্যাখ্যা:} \emph{আমরা জানি, ১ সের ≈ ০.৯৩৩ কিলোগ্রাম বা প্রায় ০.৯৩
কিলোগ্রাম।}
\end{quote}

\textbf{১ মণ লোহার ওজন ১ মণ তুলা অপেক্ষা কত গুণ বেশি?}

ক) ২ গুণ

খ) ৩ গুণ

\textbf{\hl{গ) সমান ✔}}

ঘ) ৪ গুণ

\textbf{\hl{✔ সঠিক উত্তর: গ) সমান}}

\begin{quote}
\textbf{ব্যাখ্যা:} \emph{ওজন হলো বস্তুর ভর ও অভিকর্ষজ ত্বরণের গুণফল। যেহেতু উভয়
বস্তুর ভর সমান (১ মণ) এবং একই স্থানে অভিকর্ষজ ত্বরণের মান ধ্রুবক, তাই উভয়
ক্ষেত্রে ওজন একই হবে।∴ ১ মণ লোহা এবং ১ মণ তুলার ওজন সমান।}
\end{quote}

\textbf{১ মণ কত কেজির সমান?}

\textbf{\hl{ক) ৩৭.৩২ কেজি ✔}}

খ) ৪০ কেজি

গ) ৪৫ কেজি

ঘ) ৩৫.৪০ কেজি

\textbf{\hl{✔ সঠিক উত্তর: ক) ৩৭.৩২ কেজি}}

\begin{quote}
\textbf{ব্যাখ্যা:} \emph{১ সের = ০.৯৩ কিলোগ্রাম (প্রায়)}

\emph{১ মণ = ৪০ সের = (০.৯৩ × ৪০) কিলোগ্রাম = ৩৭.২ বা ৩৭.৩২ কিলোগ্রাম (বা
প্রায় ৩৭.৩২ কেজি)।}
\end{quote}

\textbf{১ টন কত কেজির সমান?}

\textbf{\hl{ক) ১০০০ কেজি ✔}}

খ) ১০০৫ কেজি

গ) ১০১০ কেজি

ঘ) ১০১৬ কেজি

\textbf{\hl{✔ সঠিক উত্তর: ক) ১০০০ কেজি}}

\begin{quote}
\textbf{ব্যাখ্যা:} \emph{১ মেট্রিক টন = ১০০০ কেজি (তবে ১ লং টন = ১০১৬
কেজি)। সাধারণ মেট্রিক পরিমাপে ১ টন = ১০০০ কেজি।}
\end{quote}

\textbf{২০৫৭৩.৪ মিলিগ্রামে কত কিলোগ্রাম?}

\textbf{\hl{ক) ০.০২০৫৭৩৪ কি. গ্রা. ✔}}

খ) ২০.৫৭৩৪৪ কি. গ্রা.

গ) ২.০৫৭৩৪ কি. গ্রা.

ঘ) ০.২০৫৭৩৪ কি. গ্রা.

\textbf{\hl{✔ সঠিক উত্তর: ক) ০.০২০৫৭৩৪ কি. গ্রা.}}

\begin{quote}
\textbf{ব্যাখ্যা:} \emph{আমরা জানি, ১০\^{}৬ মিলিগ্রাম বা ১০,০০,০০০ মিলিগ্রাম
= ১ কিলোগ্রাম।}

\emph{∴ ২০৫৭৩.৪ মিলিগ্রাম = (২০৫৭৩.৪/১০০০০০০) কি.গ্রা. = ০.০২০৫৭৩৪
কি.গ্রা.}
\end{quote}

\textbf{এক কুইন্টাল সমান কত কিলোগ্রাম?}

ক) ১

খ) ১০

\textbf{\hl{গ) ১০০ ✔}}

ঘ) ১০০০

\textbf{\hl{✔ সঠিক উত্তর: গ) ১০০}}

\begin{quote}
\textbf{ব্যাখ্যা:} \emph{আমরা জানি, ১ কুইন্টাল = ১০০ কেজি বা কিলোগ্রাম।}
\end{quote}

\textbf{One kilogram is equal to/ ১ কেজি সমান-}

\textbf{\hl{ক) 2.21 lbs (পাউন্ড) ✔}}

খ) 2.00 lbs (পাউন্ড)

গ) 2.32 lbs (পাউন্ড)

ঘ) 1.98 lbs (পাউন্ড)

\textbf{\hl{✔ সঠিক উত্তর: ক) 2.21 lbs (পাউন্ড)}}

\begin{quote}
\textbf{ব্যাখ্যা:} \emph{১ কিলোগ্রাম (কেজি) ≈ ২.২০৪৬২ পাউন্ড বা প্রায় ২.২১
lbs।}
\end{quote}

\textbf{১০০০ কিলোগ্রাম = কত?}

\textbf{\hl{ক) ১০ কুইন্টাল (Quintal) ✔}}

খ) ১০০ ডেকাগ্রাম (Decagram)

গ) ১০০ মেট্রিকটন (Metric ton)

ঘ) ১০-০০০ গ্রাম (Gram)

\textbf{\hl{✔ সঠিক উত্তর: ক) ১০ কুইন্টাল (Quintal)}}

\begin{quote}
\textbf{ব্যাখ্যা:} \emph{আমরা জানি, ১০০ কিলোগ্রাম = ১ কুইন্টাল}

\emph{∴ ১০০০ কিলোগ্রাম = (১০০০/১০০) কুইন্টাল = ১০ কুইন্টাল।}
\end{quote}

\textbf{১ কিলোগ্রাম কত পাউন্ডের সমান?}

ক) ২

\textbf{\hl{খ) ২.২ ✔}}

গ) ২.৩২

ঘ) ১.৯৮

\textbf{\hl{✔ সঠিক উত্তর: খ) ২.২}}

\begin{quote}
\textbf{ব্যাখ্যা:} \emph{আমরা জানি, ১ কিলোগ্রাম ≈ ২.২০৪৬ পাউন্ড (বা প্রায় ২.২
পাউন্ড)।}
\end{quote}

\textbf{৪(৩/৪) পাউন্ডে কত আউন্স?}

ক) ৬৭

খ) ৫৭

\textbf{\hl{গ) ৭৬ ✔}}

ঘ) ৭৫

\textbf{\hl{✔ সঠিক উত্তর: গ) ৭৬}}

\begin{quote}
\textbf{ব্যাখ্যা:} \emph{আমরা জানি, ১ পাউন্ড = ১৬ আউন্স।}

\emph{∴ ৪(৩/৪) পাউন্ড = (১৯/৪) পাউন্ড = ((১৯/৪)×১৬) আউন্স = (১৯ × ৪) আউন্স =
৭৬ আউন্স।}
\end{quote}

\textbf{How many centimeters make a meter?}

ক) 30

খ) 50

গ) 75

\textbf{\hl{ঘ) 100 ✔}}

\textbf{\hl{✔ সঠিক উত্তর: ঘ) 100}}

\begin{quote}
\textbf{ব্যাখ্যা:} \emph{আমরা জানি, ১০০ সেন্টিমিটারে ১ মিটার (100
centimeters make a meter)।}
\end{quote}

\textbf{এক ন্যানোমিটার সমান-}

ক) ১০\^{}৬ মিটার

খ) ১০\^{}৯ মিটার

গ) ১০\^{}−৬ মিটার

\textbf{\hl{ঘ) ১০\^{}−৯ মিটার ✔}}

\textbf{\hl{✔ সঠিক উত্তর: ঘ) ১০\^{}−৯ মিটার}}

\begin{quote}
\textbf{ব্যাখ্যা:} \emph{ন্যানো (Nano) উপসর্গটির অর্থ হলো ১০\^{}−৯ অংশ।}

\emph{∴ ১ ন্যানোমিটার = ১০\^{}−৯ মিটার।}
\end{quote}

\textbf{এক ন্যানো অ্যাম্পিয়ার সমান-}

ক) ১০\^{}৬ অ্যাম্পিয়ার

খ) ১০\^{}৯ অ্যাম্পিয়ার

গ) ১০\^{}−৬ অ্যাম্পিয়ার

\textbf{\hl{ঘ) ১০\^{}−৯ অ্যাম্পিয়ার ✔}}

\textbf{\hl{✔ সঠিক উত্তর: ঘ) ১০\^{}−৯ অ্যাম্পিয়ার}}

\begin{quote}
\textbf{ব্যাখ্যা:} \emph{১ ন্যানো অ্যাম্পিয়ার = ১০\^{}−৯ অ্যাম্পিয়ার।}
\end{quote}

\textbf{দশ মাইক্রোফ্যারাডকে ফ্যারাডে প্রকাশ করলে হয়-}

ক) ১০×১০\^{}−৩ ফ্যারাড

\textbf{\hl{খ) ১০×১০\^{}−৬ ফ্যারাড ✔}}

গ) ১০×১০\^{}−১ ফ্যারাড

ঘ) ১০×১০\^{}−৯ ফ্যারাড

\textbf{\hl{✔ সঠিক উত্তর: খ) ১০×১০\^{}−৬ ফ্যারাড}}

\begin{quote}
\textbf{ব্যাখ্যা:} \emph{১ মাইক্রোফ্যারাড (μF) = ১০\^{}−৬ ফ্যারাড (F)।}

\emph{∴ ১০ মাইক্রোফ্যারাড = ১০×১০\^{}−৬ ফ্যারাড।}
\end{quote}

\textbf{এক গ্রামে কত মিলিগ্রাম?}

ক) ১০

খ) ১০০

\textbf{\hl{গ) ১০০০ ✔}}

ঘ) ১০-০০০

\textbf{\hl{✔ সঠিক উত্তর: গ) ১০০০}}

\begin{quote}
\textbf{ব্যাখ্যা:} \emph{আমরা জানি, ১ গ্রাম = ১০০০ মিলিগ্রাম।}
\end{quote}

\textbf{সময় পরিমাপে কোনটি সবচেয়ে বড় -\/-}

ক) মাইক্রোসেকেন্ড

\textbf{\hl{খ) মিলিসেকেন্ড ✔}}

গ) পিকোসেকেন্ড

ঘ) ন্যানোসেকেন্ড

\textbf{\hl{✔ সঠিক উত্তর: খ) মিলিসেকেন্ড}}

\begin{quote}
\textbf{ব্যাখ্যা:} \emph{১ মিলিসেকেন্ড = ১০\^{}−৩ সেকেন্ড}

\emph{১ মাইক্রোসেকেন্ড = ১০\^{}−৬ সেকেন্ড}

\emph{১ ন্যানোসেকেন্ড = ১০\^{}−৯ সেকেন্ড}

\emph{১ পিকোসেকেন্ড = ১০\^{}−১২ সেকেন্ড}

\emph{∴ অপশনগুলোর মধ্যে মিলিসেকেন্ড পরিমাপে সবচেয়ে বড়।}
\end{quote}

\textbf{Which one is the smallest unit?/ নিচের কোনটি ক্ষুদ্রতম একক?}

ক) Meter

\textbf{\hl{খ) Centimeter ✔}}

গ) Decimeter

ঘ) Kilometer

\textbf{\hl{✔ সঠিক উত্তর: খ) Centimeter}}

\begin{quote}
\textbf{ব্যাখ্যা:} \emph{প্রদত্ত এককগুলোর ক্রম: Kilometer \textgreater{} Meter
\textgreater{} Decimeter \textgreater{} Centimeter।}

\emph{∴ চারটির মধ্যে Centimeter হলো ক্ষুদ্রতম একক।}
\end{quote}

\textbf{Which one is the longest?/ নিচের কোনটি বৃহত্তম একক?}

\textbf{\hl{ক) Meter ✔}}

খ) Centimeter

গ) Decimeter

ঘ) Millimeter

\textbf{\hl{✔ সঠিক উত্তর: ক) Meter}}

\begin{quote}
\textbf{ব্যাখ্যা:} \emph{এখানে প্রদত্ত দৈর্ঘ্য পরিমাপের এককগুলোর মধ্যে বৃহত্তম হলো
Meter (১ মিটার = ১০ ডেসিমাইটার = ১০০ সেন্টিমিটার = ১০০০ মিলিমিটার)।}
\end{quote}

\textbf{একটি যোগ করতে কম্পিউটারের ৫০ ন্যানো সেকেন্ড সময় লাগলে সেকেন্ডে এটা
কতটি যোগ করতে পারবে?}

\textbf{\hl{ক) ২ কোটি ✔}}

খ) ৩ কোটি

গ) ৪ কোটি

ঘ) ৫ কোটি

\textbf{\hl{✔ সঠিক উত্তর: ক) ২ কোটি}}

\begin{quote}
\textbf{ব্যাখ্যা:} \emph{১ ন্যানো সেকেন্ড = ১০\^{}−৯ সেকেন্ড}

\emph{৫০ ন্যানো সেকেন্ড = ৫০×১০\^{}−৯ সেকেন্ড = ৫×১০\^{}−৮ সেকেন্ড}

\emph{কম্পিউটার ৫×১০\^{}−৮ সেকেন্ডে করে = ১টি যোগ}

\emph{∴ ১ সেকেন্ডে করে = (১/৫×১০\^{}−৮) টি = (১০\^{}৮/৫) টি = ২×১০\^{}৭
টি = ২ কোটি।}
\end{quote}

\textbf{১ : ২০০০ স্কেলে একটি শহরের ম্যাপ আঁকা আছে। ১.২ কিমি দীর্ঘ একটি
রাস্তার দৈর্ঘ্য ম্যাপে কত হবে?}

ক) ২৪ সেমি

\textbf{\hl{খ) ৬০ সেমি ✔}}

গ) ৬০ মি

ঘ) ২৪ মি

\textbf{\hl{✔ সঠিক উত্তর: খ) ৬০ সেমি}}

\begin{quote}
\textbf{ব্যাখ্যা:} \emph{১.২ কিমি = (১.২ × ১০০০০০) সেমি = ১২০০০০ সেমি}

\emph{২০০০ সেন্টিমিটার প্রকৃত দৈর্ঘ্য ম্যাপে = ১ সেন্টিমিটার}

\emph{∴ ১২০০০০ সেন্টিমিটার প্রকৃত দৈর্ঘ্য ম্যাপে = (১২০০০০×১/২০০০) = ৬০
সেন্টিমিটার।}
\end{quote}

\section{47. পরিমিতি MCQ প্রশ্ন উত্তর
অনুশীলন}\label{ux9aaux9b0ux9aeux9a4-mcq-ux9aaux9b0ux9b6ux9a8-ux989ux9a4ux9a4ux9b0-ux985ux9a8ux9b6ux9b2ux9a8}

\textbf{ত্রিভুজের ক্ষেত্রফল নির্ণয়ের সূত্র কোনটি?}

ক) πr\^{}2

\textbf{\hl{খ) (1/2)×ভূমি×উচ্চতা ✔}}

গ) ভূমি × উচ্চতা

ঘ) 2πr\^{}2

\textbf{\hl{✔ সঠিক উত্তর: খ) (1/2)×ভূমি×উচ্চতা}}

\begin{quote}
\textbf{ব্যাখ্যা:} \emph{ত্রিভুজের ক্ষেত্রফল = (1/2)×ভূমি×উচ্চতা}
\end{quote}

\textbf{ত্রিভুজের শীর্ষবিন্দু থেকে এর ভূমির উপর অঙ্কিত লম্বের দৈর্ঘ্যকে কি বলা হয়?
/ In a triangle the perpendicular drawn from the vertex to its base is
called -}

ক) মধ্যমা (Median)

খ) কর্ণ (Diagonal)

গ) অতিভুজ (Hypotenuse)

\textbf{\hl{ঘ) উচ্চতা (Altitude) ✔}}

\textbf{\hl{✔ সঠিক উত্তর: ঘ) উচ্চতা (Altitude)}}

\begin{quote}
\textbf{ব্যাখ্যা:} \emph{ত্রিভুজের শীর্ষবিন্দু থেকে এর ভূমির উপর অঙ্কিত লম্বের
দৈর্ঘ্যকে বলা হয় উচ্চতা।}

\emph{ত্রিভুজের শীর্ষবিন্দু থেকে বিপরীত বাহুর মধ্যবিন্দুর দৈর্ঘ্যকে বলা হয় মধ্যমা।}
\end{quote}

\textbf{একটি ত্রিভুজের ভূমি ৬ সে.মি. এবং এর উচ্চতা ৫ সে.মি. হলে ত্রিভুজটির
ক্ষেত্রফল কত?}

ক) ১১ বর্গ সে.মি.

\textbf{\hl{খ) ১৫ বর্গ সে.মি. ✔}}

গ) ৩০ বর্গ সে.মি.

ঘ) ২৫ বর্গ সে.মি.

\textbf{\hl{✔ সঠিক উত্তর: খ) ১৫ বর্গ সে.মি.}}

\begin{quote}
\textbf{ব্যাখ্যা:} \emph{আমরা জানি, ত্রিভুজের ক্ষেত্রফল = (1/2)×ভূমি×উচ্চতা বর্গ
একক}

\emph{= (1/2)×6×5 বর্গসেমি. = 15 বর্গসেমি.}
\end{quote}

\textbf{একটি ত্রিভুজের ভূমি ৪ মি. এবং উচ্চতা ৩ মি. হলে, ত্রিভুজটির ক্ষেত্রফল কত
বর্গমিটার?}

\textbf{\hl{ক) ৬ ✔}}

খ) ৯

গ) ১২

ঘ) ১৮

\textbf{\hl{✔ সঠিক উত্তর: ক) ৬}}

\begin{quote}
\textbf{ব্যাখ্যা:} \emph{ত্রিভুজটির ক্ষেত্রফল = (1/2)×ভূমি×উচ্চতা}

\emph{= (1/2)×4×3 বর্গমিটার = 6 বর্গমিটার}
\end{quote}

\textbf{ABC ত্রিভুজের পরিসীমা কত একক?}

\textbf{\hl{ক) ৩ ✔}}

খ) ৪

গ) ৬

ঘ) ৯

\textbf{\hl{✔ সঠিক উত্তর: ক) ৩}}

\begin{quote}
\textbf{ব্যাখ্যা:} \emph{ত্রিভুজের পরিসীমা তিন বাহুর সমষ্টির সমান।}

\emph{ABC ত্রিভুজের প্রতিবাহুর দৈর্ঘ্য যদি ১ একক হয় তবে তার পরিসীমা হবে ৩
একক।}
\end{quote}

\textbf{একটি ত্রিভুজের ৩টি বাহুর দৈর্ঘ্য ৪, ৫ ও ৩ হলে ত্রিভুজটির ক্ষেত্রফল কত?}

ক) ২০

খ) ১২

গ) ৮

\textbf{\hl{ঘ) ৬ ✔}}

\textbf{\hl{✔ সঠিক উত্তর: ঘ) ৬}}

\begin{quote}
\textbf{ব্যাখ্যা:} \emph{বাহুগুলো ৪, ৫ ও ৩ হওয়ায় এটি একটি সমকোণী ত্রিভুজ
(যেহেতু 3\^{}2+4\^{}2=5\^{}2 )।}

\emph{∴ ক্ষেত্রফল = (1/2)×ভূমি×উচ্চতা = (1/2)×4×3 = 6 বর্গ একক।}
\end{quote}

\textbf{একটি ত্রিভুজের প্রত্যেক বাহুর দৈর্ঘ্য 8 সেন্টিমিটার হলে ত্রিভুজের ক্ষেত্রফল
কত?}

ক) 16 বর্গ সেমি

খ) 64 বর্গ সেমি

\textbf{\hl{গ) 16√3 বর্গ সেমি ✔}}

ঘ) 8√3 বর্গ সেমি

\textbf{\hl{✔ সঠিক উত্তর: গ) 16√3 বর্গ সেমি}}

\begin{quote}
\textbf{ব্যাখ্যা:} \emph{সমবাহু ত্রিভুজের ক্ষেত্রফল = (√(3)/4)a\^{}2}

\emph{= (√(3)/4)×8\^{}2 = 16√(3) বর্গ সেন্টিমিটার}
\end{quote}

\textbf{4 সে.মি., 5 সে.মি. ও 9 সে.মি. বাহুবিশিষ্ট ত্রিভুজের ক্ষেত্রফল কত বর্গ
একক?}

\textbf{\hl{ক) 0 ✔}}

খ) 60

গ) 80

ঘ) 120

\textbf{\hl{✔ সঠিক উত্তর: ক) 0}}

\begin{quote}
\textbf{ব্যাখ্যা:} \emph{ত্রিভুজের পরিসীমা, 2s=a+b+c}

\emph{বা, 2s=4+5+9 {[}যেহেতু, a = 4, b = 5, c = 9{]}}

\emph{বা, s=(18/2)=9}

\emph{∴ ত্রিভুজটির ক্ষেত্রফল = √(s(s−a)(s−b)(s−c))}

\emph{= √(9(9−4)(9−5)(9−9))}

\emph{= √(9×5×4×0) = 0 বর্গ একক (যেহেতু ত্রিভুজের যে কোনো দুই বাহুর যোগফল তৃতীয়
বাহু অপেক্ষা বৃহত্তর হতে হয়, যা এখানে প্রযোজ্য নয়)}
\end{quote}

\textbf{3 cm, 4.5 cm, 5.5 cm, বাহুবিশিষ্ট কোনো ত্রিভুজের ক্ষেত্রফল কত?}

ক) 87.5

খ) 7.50

\textbf{\hl{গ) 6.75 ✔}}

ঘ) 7.55

\textbf{\hl{✔ সঠিক উত্তর: গ) 6.75}}

\begin{quote}
\textbf{ব্যাখ্যা:} \emph{অর্ধপরিসীমা S=(3+4.5+5.5/2)=6.5}

\emph{∴ ক্ষেত্রফল = √(6.5(6.5−3)(6.5−4.5)(6.5−5.5))}

\emph{= √(6.5×3.5×2×1) = √(45.5) ≈ 6.75}
\end{quote}

\textbf{একটি ত্রিভুজাকৃতি ক্ষেত্রের ক্ষেত্রফল ৮৪ বর্গগজ। ত্রিভুজের শীর্ষবিন্দু হতে
ভূমির উপর অঙ্কিত লম্বের দৈর্ঘ্য ১২ গজ হলে, ভূমির দৈর্ঘ্য কত?}

ক) ১০ গজ

খ) ১২ গজ

\textbf{\hl{গ) ১৪ গজ ✔}}

ঘ) ৭ গজ

\textbf{\hl{✔ সঠিক উত্তর: গ) ১৪ গজ}}

\begin{quote}
\textbf{ব্যাখ্যা:} \emph{আমরা জানি, ত্রিভুজের ক্ষেত্রফল = (1/2)×ভূমি×উচ্চতা}

\emph{বা, 84=(1/2)×ভূমি×12}

\emph{বা, 84=6×ভূমি}

\emph{∴ ভূমি = (84/6)=14 গজ।}
\end{quote}

\textbf{চিত্রে ΔABC এর দুই বাহুর দৈর্ঘ্য a এবং b। বাহু দুটির অন্তর্ভুক্ত কোণ θ সে
ক্ষেত্রে ΔABC এর ক্ষেত্রফল নির্ণয়ের সূত্র হবে}

\textbf{\hl{ক) (1/2)absinθ ✔}}

খ) (1/2)absin\^{}2θ

গ) abcosθ

ঘ) (1/2)abcos\^{}2θ

\textbf{\hl{✔ সঠিক উত্তর: ক) (1/2)absinθ}}

\begin{quote}
\textbf{ব্যাখ্যা:} \emph{দুটি বাহু ও তাদের অন্তর্ভুক্ত কোণ দেওয়া থাকলে ত্রিভুজের
ক্ষেত্রফল নির্ণয়ের সূত্র হলো: (1/2)absinθ}
\end{quote}

\textbf{একটি ত্রিভুজের ভূমি ১২ সেন্টিমিটার, উচ্চতা ৪ সেন্টিমিটার, উহার ক্ষেত্রফল
কত বর্গ সেন্টিমিটার?}

ক) ১২

খ) ১৮

\textbf{\hl{গ) ২৪ ✔}}

ঘ) ৪৮

\textbf{\hl{✔ সঠিক উত্তর: গ) ২৪}}

\begin{quote}
\textbf{ব্যাখ্যা:} \emph{উহার ক্ষেত্রফল = (1/2)×12×4=24 বর্গ সেন্টিমিটার}
\end{quote}

\textbf{একটি ত্রিভুজাকৃতি জমির ক্ষেত্রফল ২৬৪ বর্গমিটার এবং ভূমি ২২ মিটার হলে
উচ্চতা কত হবে?}

ক) ১২ মিটার

খ) ১৫ মিটার

\textbf{\hl{গ) ২৪ মিটার ✔}}

ঘ) ২৮ মিটার

\textbf{\hl{✔ সঠিক উত্তর: গ) ২৪ মিটার}}

\begin{quote}
\textbf{ব্যাখ্যা:} \emph{আমরা জানি, ত্রিভুজের ক্ষেত্রফল = (1/2)×ভূমি×উচ্চতা}

\emph{⇒ 264=(1/2)×22×উচ্চতা}

\emph{∴ উচ্চতা = (264×2/22)=24 মিটার}
\end{quote}

\textbf{একটি ত্রিভুজের ভূমি তার উচ্চতার দ্বিগুণ অপেক্ষা ৬ সেমি বেশি, ত্রিভুজটির
ক্ষেত্রফল ৮১০ বর্গ সেমি হলে, এর উচ্চতা কত?}

\textbf{\hl{ক) ২৭ সেমি ✔}}

খ) ২৮ সেমি

গ) ২৫ সেমি

ঘ) ২৪ সেমি

\textbf{\hl{✔ সঠিক উত্তর: ক) ২৭ সেমি}}

\begin{quote}
\textbf{ব্যাখ্যা:} \emph{ধরি, উচ্চতা = x সেমি}

\emph{∴ ভূমি = (2x + 6) সেমি}

\emph{প্রশ্নমতে, (1/2)×(2x+6)×x=810}

\emph{⇒ 2x\^{}2+6x=1620}

\emph{⇒ x\^{}2+3x−810=0}

\emph{⇒ (x+30)(x−27)=0}

\emph{∴ x = 27 (যেহেতু x এর ঋণাত্মক মান গ্রহণযোগ্য নয়)}
\end{quote}

\textbf{একটি ত্রিভুজের তিনটি বাহুর দৈর্ঘ্য ১৩, ১৪ ও ১৫ মিটার। ত্রিভুজটির
ক্ষেত্রফল কত?}

ক) ৬০ বর্গমিটার

\textbf{\hl{খ) ৮৪ বর্গমিটার ✔}}

গ) ৯০ বর্গমিটার

ঘ) ১০৮ বর্গমিটার

\textbf{\hl{✔ সঠিক উত্তর: খ) ৮৪ বর্গমিটার}}

\begin{quote}
\textbf{ব্যাখ্যা:} \emph{ত্রিভুজটির পরিসীমা, 2s=(13+14+15) মিটার = 42 মিটার}

\emph{∴ s=21 মিটার}

\emph{∴ ত্রিভুজটির ক্ষেত্রফল = √(s(s−a)(s−b)(s−c)) বর্গমিটার}

\emph{= √(21(21−13)(21−14)(21−15)) বর্গমিটার}

\emph{= √(21×8×7×6) = 84 মিটার।}
\end{quote}

\textbf{একটি ত্রিভুজের তিন বাহুর দৈর্ঘ্য (x+৯) মিটার, (2x+১) মিটার এবং ২
(2x-১) মিটার। ত্রিভুজটির পরিসীমা ২৯ মিটার হলে ক্ষুদ্রতম বাহুটির দৈর্ঘ্য হবে-}

ক) ৩ মিটার

খ) ৫ মিটার

\textbf{\hl{গ) ৭ মিটার ✔}}

ঘ) ১০ মিটার

\textbf{\hl{✔ সঠিক উত্তর: গ) ৭ মিটার}}

\begin{quote}
\textbf{ব্যাখ্যা:} \emph{ত্রিভুজের পরিসীমা = তিন বাহুর সমষ্টি}

\emph{প্রশ্নানুসারে, (x+9)+(2x+1)+2(2x−1)=29}

\emph{বা, x+9+2x+1+4x−2=29}

\emph{বা, 7x+8=29}

\emph{বা, 7x=21}

\emph{∴ x=3}

\emph{বাহুগুলো হলো: 3+9=12 , 2(3)+1=7 , এবং 2(2×3−1)=10 .}

\emph{∴ ক্ষুদ্রতম বাহুটির দৈর্ঘ্য = 7 মিটার।}
\end{quote}

\textbf{একটি ত্রিভুজের বাহুগুলোর অনুপাত ৫ : ১২ : ১৩ এবং এর পরিসীমা ১৯৫
সেন্টিমিটার। বৃহত্তম ও ক্ষুদ্রতম বাহুর দৈর্ঘ্যের সমষ্টি কত?}

\textbf{\hl{ক) ১১৭ সেমি ✔}}

খ) ১১০.৫ সেমি

গ) ১২২.৫ সেমি

ঘ) ১১৯ সেমি

\textbf{\hl{✔ সঠিক উত্তর: ক) ১১৭ সেমি}}

\begin{quote}
\textbf{ব্যাখ্যা:} \emph{ধরি, বাহুগুলোর দৈর্ঘ্যের পরিমাণ 5x সেমি, 12x সেমি এবং
13x সেমি。}

\emph{প্রশ্নমতে, 5x+12x+13x=195}

\emph{⇒ 30x=195}

\emph{∴ x=(195/30)=6.5 সেমি}

\emph{∴ বৃহত্তম এবং ক্ষুদ্রতম বাহুর দৈর্ঘ্যের সমষ্টি = 5x+13x=18x সেমি}

\emph{= 18×6.5=117 সেমি}
\end{quote}

\textbf{ত্রিভুজ ABC এর ক্ষেত্রফল ২০ বর্গ একক। X এবং Y, AB এবং AC এর মধ্যবিন্দু।
ত্রিভুজ AXY সমান কত বর্গ একক?}

\textbf{\hl{ক) ৫ বর্গ একক ✔}}

খ) ২ বর্গ একক

গ) ৪ বর্গ একক

ঘ) ১০ বর্গ একক

\textbf{\hl{✔ সঠিক উত্তর: ক) ৫ বর্গ একক}}

\begin{quote}
\textbf{ব্যাখ্যা:} \emph{সদৃশ ত্রিভুজের ক্ষেত্রফলদ্বয়ের অনুপাত তাদের যে কোনো দুই
বাহুর উপর অঙ্কিত বর্গক্ষেত্রদ্বয়ের অনুপাতের সমান。}

\emph{⇒ (AXY/ABC)=(AX\^{}2/AB\^{}2) ⇒ (AXY/20)=(AX\^{}2/(2AX)\^{}2)}

\emph{⇒ AXY=(20/4) ∴ AXY=5 = 5 বর্গ একক}
\end{quote}

\textbf{সমবাহু ত্রিভুজের বাহু দৈর্ঘ্য যদি \textquotesingle a\textquotesingle{}
হয় তবে ত্রিভুজের ক্ষেত্রফল কত বর্গ একক ??}

ক) (√(3)/2)a\^{}2

খ) (2/3)a\^{}2

গ) (2/√(3))a\^{}2

\textbf{\hl{ঘ) (√(3)/4)a\^{}2 ✔}}

\textbf{\hl{✔ সঠিক উত্তর: ঘ) (√(3)/4)a\^{}2}}

\begin{quote}
\textbf{ব্যাখ্যা:} \emph{সমবাহু ত্রিভুজের ক্ষেত্রফল নির্ণয়ের মূল সূত্র হলো
(√(3)/4)a\^{}2 বর্গ একক।}
\end{quote}

\textbf{একটি সমবাহু ত্রিভুজের বাহুর দৈর্ঘ্য ৬ সে.মি. হলে এর ক্ষেত্রফল কত বর্গ
সেমি.?}

ক) 3√(3)

খ) 4√(3)

গ) 6√(3)

\textbf{\hl{ঘ) 9√(3) ✔}}

\textbf{\hl{✔ সঠিক উত্তর: ঘ) 9√(3)}}

\begin{quote}
\textbf{ব্যাখ্যা:} \emph{সমবাহু ত্রিভুজের ক্ষেত্রফল = (√(3)/4)a\^{}2}

\emph{= (√(3)/4)×6\^{}2 = 9√(3) বর্গসেমি.}
\end{quote}

\textbf{ত্রিভুজ ABC-এর BC = CA = AB = 5 সেন্টিমিটার। ত্রিভুজটির ক্ষেত্রফল কত?}

\textbf{\hl{ক) (25√(3)/4) ✔}}

খ) (25√(3)/6)

গ) (25√(3)/2)

ঘ) (25√(3)/7)

\textbf{\hl{✔ সঠিক উত্তর: ক) (25√(3)/4)}}

\begin{quote}
\textbf{ব্যাখ্যা:} \emph{সমবাহু ত্রিভুজের বাহুর দৈর্ঘ্য a হলে ক্ষেত্রফল =
(√(3)/4)a\^{}2}

\emph{∴ ত্রিভুজটির ক্ষেত্রফল = (√(3)/4)(5)\^{}2=(25√(3)/4)}
\end{quote}

\textbf{একটি সমবাহু ত্রিভুজের প্রত্যেকটি বাহুর দৈর্ঘ্য 2 মিটার কমাইলে এর ক্ষেত্রফল
3√(3) বর্গমিটার কমিয়া যায়। সমবাহু ত্রিভুজটির বাহুর দৈর্ঘ্য নির্ণয় কর।}

\textbf{\hl{ক) 4 মিটার ✔}}

খ) 6 মিটার

গ) 8 মিটার

ঘ) 5 মিটার

\textbf{\hl{✔ সঠিক উত্তর: ক) 4 মিটার}}

\begin{quote}
\textbf{ব্যাখ্যা:} \emph{ধরি, সমবাহু ত্রিভুজের প্রত্যেক বাহুর দৈর্ঘ্য = a মিটার}

\emph{ক্ষেত্রফল = (√(3)/4)a\^{}2}

\emph{বাহুর দৈর্ঘ্য 2 মিটার কমানো হলে, নতুন বাহু = (a−2) মিটার এবং নতুন
ক্ষেত্রফল = (√(3)(a−2)\^{}2/4)}

\emph{প্রশ্নমতে, (√(3)a\^{}2/4)−(√(3)(a−2)\^{}2/4)=3√(3)}

\emph{সমাধান করে পাই, a=4 মিটার।}
\end{quote}

\textbf{একটি সমবাহু ত্রিভুজের একটি বাহু ১৬ মিটার। ত্রিভুজটির ক্ষেত্রফল কত?}

ক) ১৬ বর্গমিটার

খ) ৩২ বর্গমিটার

গ) ৬৪ বর্গমিটার

ঘ) ১২৮ বর্গমিটার

\textbf{\hl{ঙ) কোনোটিই নয় ✔}}

\textbf{\hl{✔ সঠিক উত্তর: ঙ) কোনোটিই নয়}}

\begin{quote}
\textbf{ব্যাখ্যা:} \emph{সমবাহু ত্রিভুজের ক্ষেত্রফল = (√(3)/4)a\^{}2 বর্গ একক}

\emph{= (√(3)/4)×(16)\^{}2 = 64√(3) বর্গ মি.}
\end{quote}

\textbf{একটি সমবাহু ত্রিভুজের এক বাহুর দৈর্ঘ্য 4 সেঃ মিঃ হলে, উহার ক্ষেত্রফল হবে
কত?}

ক) √(3) বর্গ সেঃ মিঃ

খ) 2√(3) বর্গ সেঃ মিঃ

\textbf{\hl{গ) 16√(3) বর্গ সেঃ মিঃ ✔}}

ঘ) (1/2)√(3) বর্গ সেঃ মিঃ

\textbf{\hl{✔ সঠিক উত্তর: গ) 16√(3) বর্গ সেঃ মিঃ}}

\begin{quote}
\textbf{ব্যাখ্যা:} \emph{সমবাহু ত্রিভুজের ক্ষেত্রফল = (√(3)/4)a\^{}2 বর্গ সেঃ
মিঃ}

\emph{= (√(3)/4)×(4)\^{}2 বর্গ সেঃ মিঃ}

\emph{= √(3)×16 বর্গ সেঃ মিঃ}

\emph{= 16√(3) বর্গ সেঃ মিঃ}
\end{quote}

\textbf{একটি সমবাহু ত্রিভুজের ক্ষেত্রফল ৫০ বর্গ সে. মি.। ত্রিভুজের প্রতি বাহুর
দৈর্ঘ্য কত?}

ক) ১৫.২ সে. মি.

খ) ১০.৫ সে. মি.

\textbf{\hl{গ) ১০.৭ সে. মি. ✔}}

ঘ) ১৭.১ সে. মি.

\textbf{\hl{✔ সঠিক উত্তর: গ) ১০.৭ সে. মি.}}

\begin{quote}
\textbf{ব্যাখ্যা:} \emph{ধরি, সমবাহু ত্রিভুজের প্রতিবাহু a সেমিঃ ∴ ক্ষেত্রফল =
(√(3)/4)a\^{}2}

\emph{শর্তমতে, (√(3)/4)a\^{}2=50}

\emph{⇒ a\^{}2=(50×4/√(3))=(200/1.732)=115.47}

\emph{∴ a=√(115.47)≈10.74 সেমিঃ}
\end{quote}

\textbf{একটি সমবাহু ত্রিভুজের ক্ষেত্রফল 4√(3) বর্গমিটারের হলে বাহুর দৈর্ঘ্য কত?}

ক) (√(3)/4) মিটার

খ) √(3) মিটার

\textbf{\hl{গ) 4 মিটার ✔}}

ঘ) 16 মিটার

\textbf{\hl{✔ সঠিক উত্তর: গ) 4 মিটার}}

\begin{quote}
\textbf{ব্যাখ্যা:} \emph{ধরি, সমবাহু ত্রিভুজের একবাহুর দৈর্ঘ্য a মিটার}

\emph{ক্ষেত্রফল = (√(3)/4)a\^{}2 বর্গমিটার}

\emph{প্রশ্নমতে, (√(3)/4)a\^{}2=4√(3)}

\emph{⇒ (a\^{}2/4)=4}

\emph{⇒ a\^{}2=16}

\emph{∴ a=4 মিটার}
\end{quote}

\textbf{একটি সমবাহু ত্রিভুজের পরিসীমা ৯ সে.মি. হলে এর উচ্চতা কত সে.মি.?}

\textbf{\hl{ক) (3√(3)/2) ✔}}

খ) (3√(3)/4)

গ) 4√(3)

ঘ) 2√(3)

\textbf{\hl{✔ সঠিক উত্তর: ক) (3√(3)/2)}}

\begin{quote}
\textbf{ব্যাখ্যা:} \emph{ধরি, ত্রিভুজটির একবাহুর দৈর্ঘ্য a}

\emph{প্রশ্নমতে, 3a=9 ⇒ a=3 সে.মি.}

\emph{সমবাহু ত্রিভুজের উচ্চতা h=(√(3)/2)a}

\emph{∴ উচ্চতা = (√(3)/2)×3=(3√(3)/2) সে.মি.}
\end{quote}

\textbf{একটি সমদ্বিবাহু ত্রিভুজের সমান কোণদ্বয়ের একটি ৫০° হলে তৃতীয় কোণটির
পরিমাপ কত?}

\textbf{\hl{ক) 80° ✔}}

খ) 70°

গ) 100°

ঘ) 90°

\textbf{\hl{✔ সঠিক উত্তর: ক) 80°}}

\begin{quote}
\textbf{ব্যাখ্যা:} \emph{সমান কোণদ্বয়ের একটি ৫০°}

\emph{∴ সমান কোণদ্বয়ের সমষ্টি = (৫০° + ৫০°) = ১০০°}

\emph{∴ তৃতীয় কোণটি = (১৮০° − ১০০°) = ৮০°}
\end{quote}

\textbf{একটি সমদ্বিবাহু ত্রিভুজের ভূমি ২ সে.মি. এবং অপর বাহুদ্বয়ের প্রতিটির
দৈর্ঘ্য ৩ সে.মি. হলে, ত্রিভুজটির ক্ষেত্রফল কত?}

ক) ৯ বর্গ সে.মি.

খ) ৬ বর্গ সে.মি.

\textbf{\hl{গ) √(8) বর্গ সে.মি. ✔}}

ঘ) ২ √(3) বর্গ সে.মি.

\textbf{\hl{✔ সঠিক উত্তর: গ) √(8) বর্গ সে.মি.}}

\begin{quote}
\textbf{ব্যাখ্যা:} \emph{ক্ষেত্রফল = (2/4) √(4(3)\^{}2−(2)\^{}2)}

\emph{⇒ (2/4) √(4×9−4) = (2/4) √(36−4)}

\emph{⇒ (2/4) √(32) = (2/4) × 4√(2) = √(8) বর্গ সে.মি.}
\end{quote}

\textbf{একটি সমদ্বিবাহু ত্রিভুজের ক্ষেত্রফল কত হবে, যেখানে উহার সমান সমান
বাহুদ্বয়ের দৈর্ঘ্য ৫০ সেমি ও ভূমি ৬০ সেমি?}

ক) 10000 বর্গ সেমি

খ) 11000 বর্গ সেমি

\textbf{\hl{গ) 1200 বর্গ সেমি ✔}}

ঘ) 1100 বর্গ সেমি

\textbf{\hl{✔ সঠিক উত্তর: গ) 1200 বর্গ সেমি}}

\begin{quote}
\textbf{ব্যাখ্যা:} \emph{সমদ্বিবাহু ত্রিভুজের ক্ষেত্রফল = (b/4)
√(4a\^{}2−b\^{}2)}

\emph{⇒ (60/4) √(4×(50)\^{}2−(60)\^{}2)}

\emph{⇒ 15√(4×2500−3600)}

\emph{⇒ 15√(10000−3600)}

\emph{⇒ 15√(6400)}

\emph{⇒ 15×80 = 1200 বর্গ সেমি}
\end{quote}

\textbf{একটি সমদ্বিবাহু ত্রিভুজের ভূমি ১৬ মিটার এবং অপর দুটি বাহু প্রতিটি ১০
মিটার। ত্রিভুজটির ক্ষেত্রফল কত?}

ক) ৩৬ বর্গমিটার

খ) ৪২ বর্গমিটার

গ) ৫০ বর্গমিটার

\textbf{\hl{ঘ) ৪৮ বর্গমিটার ✔}}

\textbf{\hl{✔ সঠিক উত্তর: ঘ) ৪৮ বর্গমিটার}}

\begin{quote}
\textbf{ব্যাখ্যা:} \emph{দেওয়া আছে, ভূমি b=16 মি.; বাহুর দৈর্ঘ্য a=10 মি.}

\emph{সমদ্বিবাহু ত্রিভুজের ক্ষেত্রফল = (b/4)√(4a\^{}2−b\^{}2)}

\emph{= (16/4)√(4×(10)\^{}2−(16)\^{}2)=4√(400−256)}

\emph{= 4×12=48 বর্গমিটার}
\end{quote}

\textbf{ABC ত্রিভুজে AB=AC, BC এর সমান্তরাল EF রেখা AB ও AC কে যথাক্রমে E ও
F বিন্দুতে ছেদ করেছে। ∠B=48° হলে, ∠EAF+∠AFE = কত?}

\textbf{\hl{ক) 132° ✔}}

খ) 108°

গ) 160°

ঘ) 80°

\textbf{\hl{✔ সঠিক উত্তর: ক) 132°}}

\begin{quote}
\textbf{ব্যাখ্যা:} \emph{অনুরূপ ∠AEF = অনুরূপ ∠B}

\emph{∴ ∠AEF = 48°}

\emph{এখন, ∠AEF + ∠EAF + ∠AFE = 180° {[}∴ ত্রিভুজের তিনকোণের সমষ্টি
180°{]}}

\emph{বা, 48° + ∠EAF + ∠AFE = 180°}

\emph{∴ ∠EAF + ∠AFE = 180° − 48° = 132°}
\end{quote}

\textbf{একটি সমকোণী ত্রিভুজের ক্ষেত্রফল-}

ক) ভূমি × উচ্চতা

খ) ভূমি × অতিভুজ

\textbf{\hl{গ) (1/2) × ভূমি × উচ্চতা ✔}}

ঘ) ভূমি × উচ্চতা × অতিভুজ

\textbf{\hl{✔ সঠিক উত্তর: গ) (1/2) × ভূমি × উচ্চতা}}

\begin{quote}
\textbf{ব্যাখ্যা:} \emph{আমরা জানি, সমকোণী ত্রিভুজের ক্ষেত্রফল = (1/2) × ভূমি ×
উচ্চতা}
\end{quote}

\textbf{একটি সমকোণী ত্রিভুজের অতিভুজ ছাড়া অন্য দুই বাহুর দৈর্ঘ্য ০.২ মিটার এবং
০.৩ মিটার হলে ত্রিভুজের ক্ষেত্রফল কত?}

ক) ০.০৬ বর্গমিটার

\textbf{\hl{খ) ০.০৩ বর্গমিটার ✔}}

গ) ০.০৫ বর্গমিটার

ঘ) ০.০১ বর্গমিটার

\textbf{\hl{✔ সঠিক উত্তর: খ) ০.০৩ বর্গমিটার}}

\begin{quote}
\textbf{ব্যাখ্যা:} \emph{দেওয়া আছে, সমকোণ সংলগ্ন বাহুদ্বয়ের দৈর্ঘ্য ০.২ মিটার
এবং ০.৩ মিটার。}

\emph{আমরা জানি, ত্রিভুজের ক্ষেত্রফল = (1/2) × সমকোণ সংলগ্ন বাহুদ্বয়ের গুণফল}

\emph{= (1/2) × ০.২ × ০.৩ বর্গ মিটার = ০.০৩ বর্গ মিটার}
\end{quote}

\textbf{একটি সমদ্বিবাহু সমকোণী ত্রিভুজের অতিভুজের দৈর্ঘ্য ১০ সে.মি. ত্রিভুজটির
ক্ষেত্রফল-}

ক) ৫০ বর্গ সে.মি.

\textbf{\hl{খ) ২৫ বর্গ সে.মি. ✔}}

গ) ১০০ বর্গ সে.মি.

ঘ) ৫ বর্গ সে.মি.

\textbf{\hl{✔ সঠিক উত্তর: খ) ২৫ বর্গ সে.মি.}}

\begin{quote}
\textbf{ব্যাখ্যা:} \emph{ধরি, সমকোণ সংলগ্ন প্রতিটি বাহুর দৈর্ঘ্য = x সে.মি.}

\emph{পীথাগোরাসের সূত্রানুযায়ী পাই,}

\emph{x\^{}2+x\^{}2=10\^{}2}

\emph{বা, 2x\^{}2=100}

\emph{বা, x\^{}2=50}

\emph{∴ x=5√(2)}

\emph{∴ ত্রিভুজটির ক্ষেত্রফল = (1/2) × সমকোণ সংলগ্ন বাহু দুটির গুণফল}

\emph{= (1/2) × 5√(2) × 5√(2) বর্গ সে.মি.}

\emph{= ২৫ বর্গ সে.মি.}
\end{quote}

\textbf{একটি সমদ্বিবাহু সমকোণী ত্রিভুজের অতিভুজের দৈর্ঘ্য ১২ সেমি হলে, ত্রিভুজটির
ক্ষেত্রফল কত বর্গ সেমি?}

\textbf{\hl{ক) ৩৬ ✔}}

খ) ৪৮

গ) ৫৬

ঘ) ৭২

\textbf{\hl{✔ সঠিক উত্তর: ক) ৩৬}}

\begin{quote}
\textbf{ব্যাখ্যা:} \emph{সমদ্বিবাহু সমকোণী ত্রিভুজের ক্ষেত্রফল =
(1/4)×(অতিভুজ)\^{}2}

\emph{= (1/4)×(12)\^{}2=(1/4)×144=36 বর্গ সেমি}
\end{quote}

\textbf{একটি সমকোণী ত্রিভুজের সমকোণ সংলগ্ন বাহুদ্বয়ের দৈর্ঘ্য যথাক্রমে ৬ সে.মি.
এবং ৮ সে.মি. হলে ক্ষেত্রফল কত হবে?}

\textbf{\hl{ক) ২৪ বর্গ সে.মি. ✔}}

খ) ৪২ বর্গ সে.মি.

গ) ৮৮ বর্গ সে.মি.

ঘ) ৪৫ বর্গ সে.মি.

\textbf{\hl{✔ সঠিক উত্তর: ক) ২৪ বর্গ সে.মি.}}

\begin{quote}
\textbf{ব্যাখ্যা:} \emph{সমকোণী ত্রিভুজের ক্ষেত্রফল = (1/2) × (ভূমি × উচ্চতা)}

\emph{= (1/2) × (৬ × ৮) বর্গ সেন্টিমিটার = ২৪ বর্গ সেন্টিমিটার।}
\end{quote}

\textbf{একটি সমকোণী ত্রিভুজের ক্ষেত্রফল ১৪৪ বর্গ একক। সমকোণ সন্নিহিত বাহুদ্বয়ের
একটির দৈর্ঘ্য ১২ একক হলে অপরটি কত একক?}

ক) ৩০ একক

\textbf{\hl{খ) ২৪ একক ✔}}

গ) ২০ একক

ঘ) ১৫ একক

\textbf{\hl{✔ সঠিক উত্তর: খ) ২৪ একক}}

\begin{quote}
\textbf{ব্যাখ্যা:} \emph{ধরি, সমকোণ সংহিত অপর বাহুর দৈর্ঘ্য = x একক}

\emph{সমকোণী ত্রিভুজের ক্ষেত্রফল = (1/2) × সমকোণ সন্নিহিত বাহুদ্বয়ের গুণফল}

\emph{⇒ 144=(1/2)×12×x}

\emph{⇒ x=(144×2/12)=24 একক}
\end{quote}

\textbf{একটি সমদ্বিবাহু সমকোণী ত্রিভুজের অতিভুজের দৈর্ঘ্য ১২ সেমি হলে ত্রিভুজটির
ক্ষেত্রফল কত বর্গ সেমি?}

\textbf{\hl{ক) ৩৬ ✔}}

খ) ৪৮

গ) ৫৬

ঘ) ৭২

\textbf{\hl{✔ সঠিক উত্তর: ক) ৩৬}}

\begin{quote}
\textbf{ব্যাখ্যা:} \emph{ধরি, সমদ্বিবাহুর দৈর্ঘ্য = x}

\emph{পীথাগোরাসের সূত্রানুযায়ী, x\^{}2+x\^{}2=(12)\^{}2}

\emph{বা, 2x\^{}2=144}

\emph{বা, x\^{}2=72}

\emph{∴ ক্ষেত্রফল = (1/2) × ভূমি × উচ্চতা}

\emph{= (1/2)×x\^{}2=(1/2)×72=36 বর্গ সেমি।}
\end{quote}

\textbf{একজন ব্যক্তি একটি বর্গক্ষেত্র জুড়ে আড়াআড়িভাবে হেঁটে গিয়েছিল। প্রান্ত
বরাবর না হাঁটার কারণে কত শতাংশ কম হাঁটতে হয়েছিল?}

ক) ২০\%

খ) ২৫\%

গ) ৩২\%

\textbf{\hl{ঘ) কোনোটিই নয় ✔}}

\textbf{\hl{✔ সঠিক উত্তর: ঘ) কোনোটিই নয়}}

\begin{quote}
\textbf{ব্যাখ্যা:} \emph{ধরি, বর্গক্ষেত্রের বাহুর দৈর্ঘ্য a একক}

\emph{∴ কর্ণ = √(2)a একক}

\emph{প্রান্ত বরাবর মোট দূরত্ব = (a+a+a+a)=4a একক}

\emph{আড়াআড়িভাবে (কর্ণ বরাবর) অতিক্রান্ত দূরত্ব = 2√(2)a একক}

\emph{কম হাঁটতে হয়েছে = 4a−2√(2)a=1.172a একক}

\emph{∴ কম হাঁটতে হয়েছিল = ((1.172a/4a)×100)\%=29.3\% (যা প্রদত্ত অপশনের
কোনোটিই নয়)}
\end{quote}

\textbf{ট্রাপিজিয়ামের সমান্তরাল বাহুর দৈর্ঘ্য যথাক্রমে ৭ সে.মি. ও ৫ সে.মি. এবং
তাদের মধ্যে লম্ব দূরত্ব ৮ সে.মি. হলে এর ক্ষেত্রফল কত বর্গ সে.মি.?}

ক) ১২ বর্গ সে.মি.

\textbf{\hl{খ) ২৪ বর্গ সে.মি. ✔}}

গ) ৩০ বর্গ সে.মি.

ঘ) ৪৮ বর্গ সে.মি.

\textbf{\hl{✔ সঠিক উত্তর: খ) ২৪ বর্গ সে.মি.}}

\begin{quote}
\textbf{ব্যাখ্যা:} \emph{ট্রাপিজিয়ামের ক্ষেত্রফল = (1/2)×(সমান্তরাল বাহুদ্বয়ের
যোগফল)×উচ্চতা}

\emph{= (1/2)×(7+5)×8 বর্গ সেমি = ২৪ বর্গ সে.মি.}
\end{quote}

\textbf{যে চতুর্ভুজের কোণগুলোর পরিমাপের অনুপাত ১ : ২ : ২ : ৩ তার বৃহত্তম কোণের
পরিমাপ}

ক) ১২০°

খ) ১২৫°

গ) ১৩০°

\textbf{\hl{ঘ) ১৩৫° ✔}}

\textbf{\hl{✔ সঠিক উত্তর: ঘ) ১৩৫°}}

\begin{quote}
\textbf{ব্যাখ্যা:} \emph{ধরি, কোণগুলো যথাক্রমে 1x , 2x , 2x এবং 3x}

\emph{চতুর্ভুজের চার কোণের সমষ্টি = ৩৬০°}

\emph{প্রশ্নমতে, x+2x+2x+3x=360°}

\emph{⇒ 8x=360°}

\emph{∴ x=45°}

\emph{∴ বৃহত্তম কোণটি = 3×45°=135°}
\end{quote}

\textbf{বৃত্তের ব্যাসার্ধ 5 একক হলে, বৃত্তের ক্ষেত্রফল কত বর্গ একক?}

ক) 5π

খ) 10π

গ) 20π

\textbf{\hl{ঘ) 25π ✔}}

\textbf{\hl{✔ সঠিক উত্তর: ঘ) 25π}}

\begin{quote}
\textbf{ব্যাখ্যা:} \emph{আমরা জানি, বৃত্তের ক্ষেত্রফল = πr\^{}2}

\emph{∴ ক্ষেত্রফল = π×(5)\^{}2=25π বর্গ একক।}
\end{quote}

\textbf{একটি বৃত্তের ব্যাসার্ধ r হলে বৃত্তের ক্ষেত্রফল কত?}

\textbf{\hl{ক) πr\^{}2 ✔}}

খ) 2r\^{}2π

গ) 2πr

ঘ) কোনোটিই নয়

\textbf{\hl{✔ সঠিক উত্তর: ক) πr\^{}2}}

\begin{quote}
\textbf{ব্যাখ্যা:} \emph{বৃত্তের ব্যাসার্ধ r হলে, বৃত্তের ক্ষেত্রফলের সূত্র হলো
πr\^{}2 ।}
\end{quote}

\textbf{একটি বৃত্তের ব্যাস 16 একক দৈর্ঘ্যের হলে, বৃত্তটির ক্ষেত্রফল -}

ক) 16π

\textbf{\hl{খ) 64π ✔}}

গ) 256π

ঘ) 4π

\textbf{\hl{✔ সঠিক উত্তর: খ) 64π}}

\begin{quote}
\textbf{ব্যাখ্যা:} \emph{দেওয়া আছে, বৃত্তের ব্যাস d=16 একক}

\emph{∴ বৃত্তের ব্যাসার্ধ r=(16/2)=8 একক}

\emph{∴ বৃত্তের ক্ষেত্রফল = πr\^{}2=π×(8)\^{}2=64π বর্গ একক।}
\end{quote}

\textbf{14 সে.মি. ব্যাসার্ধবিশিষ্ট একটি বৃত্তক্ষেত্রের ক্ষেত্রফল কত বর্গ সে.মি.?}

\textbf{\hl{ক) 616 ✔}}

খ) 154

গ) 574

ঘ) 264

\textbf{\hl{✔ সঠিক উত্তর: ক) 616}}

\begin{quote}
\textbf{ব্যাখ্যা:} \emph{বৃত্তের ক্ষেত্রফল = πr\^{}2 {[}যেখানে r = ব্যাসার্ধ{]}}

\emph{= (22/7)×(14)\^{}2=(22/7)×196=22×28=616 বর্গ সে.মি.}
\end{quote}

\textbf{একটি বৃত্তের পরিধি ও ক্ষেত্রফল যথাক্রমে ১৩২ সেমি ও ১,৩৮৬ বর্গ সেমি।
বৃহত্তম জ্যা এর দৈর্ঘ্য কত?}

ক) ৬৬ সেমি

\textbf{\hl{খ) ৪২ সেমি ✔}}

গ) ২১ সেমি

ঘ) ২২ সেমি

\textbf{\hl{✔ সঠিক উত্তর: খ) ৪২ সেমি}}

\begin{quote}
\textbf{ব্যাখ্যা:} \emph{আমরা জানি, বৃত্তের বৃহত্তম জ্যা হলো বৃত্তের ব্যাস।}

\emph{বৃত্তের পরিধি 2πr=132}

\emph{⇒ 2×(22/7)×r=132}

\emph{⇒ r=(132×7/2×22)=21 সেমি}

\emph{∴ বৃহত্তম জ্যা (ব্যাস) d=2r=2×21=42 সেমি।}
\end{quote}

\textbf{একটি বৃত্তের ক্ষেত্রফল ১৬ বর্গমিটার, পরিধি ৮ মিটার। এর ব্যাসার্ধ কত
মিটার?}

\textbf{\hl{ক) 4 ✔}}

খ) 3

গ) 2

ঘ) 5

\textbf{\hl{✔ সঠিক উত্তর: ক) 4}}

\begin{quote}
\textbf{ব্যাখ্যা:} \emph{আমরা জানি, বৃত্তের ক্ষেত্রফল πr\^{}2=16 এবং পরিধি
2πr=8}

\emph{∴ ব্যাসার্ধ r=(2×ক্ষেত্রফল/পরিধি)=(2×16/8)=(32/8)=4 মিটার।}
\end{quote}

\textbf{নিচের কোনটি দ্বারা বৃত্তের ক্ষেত্রফল নির্দেশ করে?}

ক) 2πr

খ) (4/3)πr\^{}3

গ) 4πr\^{}2

\textbf{\hl{ঘ) πr\^{}2 ✔}}

\textbf{\hl{✔ সঠিক উত্তর: ঘ) πr\^{}2}}

\begin{quote}
\textbf{ব্যাখ্যা:} \emph{বৃত্তের ক্ষেত্রফলের সূত্র হলো πr\^{}2 (যেখানে r হলো
বৃত্তের ব্যাসার্ধ)।}
\end{quote}

\textbf{দুইটি বৃত্তের ব্যাসার্ধের অনুপাত ৪ : ৫। বৃত্ত দুটি ক্ষেত্রফলের অনুপাত কত?}

\textbf{\hl{ক) ১৬:২৫ ✔}}

খ) ১৬:৫

গ) ৪:২৫

ঘ) ২৫:১৬

\textbf{\hl{✔ সঠিক উত্তর: ক) ১৬:২৫}}

\begin{quote}
\textbf{ব্যাখ্যা:} \emph{ধরি, বৃত্ত দুটির ব্যাসার্ধ যথাক্রমে r\_1=4 এবং r\_2=5}

\emph{ক্ষেত্রফলের অনুপাত =
πr\_1\^{}2:πr\_2\^{}2=r\_1\^{}2:r\_2\^{}2=4\^{}2:5\^{}2=16:25}
\end{quote}

\textbf{কোনো বস্তুকণা r ব্যাসার্ধ বিশিষ্ট বৃত্তাকার পথ অর্ধেক ঘুরে আসলে তার
অতিক্রান্ত দূরত্ব কত?}

ক) 2πr

\textbf{\hl{খ) πr ✔}}

গ) 2r

ঘ) 0

\textbf{\hl{✔ সঠিক উত্তর: খ) πr}}

\begin{quote}
\textbf{ব্যাখ্যা:} \emph{সম্পূর্ণ বৃত্তাকার পথে একবার ঘুরে আসলে অতিক্রান্ত দূরত্ব হয়
বৃত্তের পরিধি 2πr}

\emph{∴ অর্ধেক ঘুরে আসলে তার অতিক্রান্ত দূরত্ব হবে = (2πr/2)=πr ।}
\end{quote}

\textbf{কোনো বৃত্তের ব্যাস 10 cm হলে বৃত্তটির ক্ষেত্রফল কত?}

ক) 57π

খ) 10π

\textbf{\hl{গ) 25π ✔}}

ঘ) 100π

\textbf{\hl{✔ সঠিক উত্তর: গ) 25π}}

\begin{quote}
\textbf{ব্যাখ্যা:} \emph{বৃত্তটির ব্যাসার্ধ r=(10/2)=5 সেমি.}

\emph{∴ ক্ষেত্রফল = πr\^{}2=π×(5)\^{}2=25π বর্গ সেমি.}
\end{quote}

\textbf{r ব্যাসার্ধবিশিষ্ট একটি অর্ধ-বৃত্তের মধ্যে অন্তলিখিত করা যায় এরূপ সর্ববৃহৎ
ত্রিভুজের ক্ষেত্রফল কত?}

\textbf{\hl{ক) r\^{}2 ✔}}

খ) 2r\^{}2

গ) (1/2)r\^{}2

ঘ) r\^{}3

\textbf{\hl{✔ সঠিক উত্তর: ক) r\^{}2}}

\begin{quote}
\textbf{ব্যাখ্যা:} \emph{অর্ধবৃত্তের মধ্যে অঙ্কিত বৃহত্তম ত্রিভুজের ভূমি = বৃত্তের
ব্যাস = 2r}

\emph{এবং ত্রিভুজের উচ্চতা = অর্ধবৃত্তের ব্যাসার্ধ = r}

\emph{∴ ত্রিভুজটির ক্ষেত্রফল = (1/2)×ভূমি×উচ্চতা}

\emph{= (1/2)×2r×r=r\^{}2 বর্গ একক।}
\end{quote}

\textbf{৭ সেন্টিমিটার ব্যাসার্ধ বিশিষ্ট বৃত্তের অন্তলিখিত বর্গক্ষেত্রের ক্ষেত্রফল
কত?}

\textbf{\hl{ক) ৯৮ বর্গ সেমি ✔}}

খ) ১৪৯ বর্গ সেমি

গ) ১৯৬ বর্গ সেমি

ঘ) ৪৯ বর্গ সেমি

\textbf{\hl{✔ সঠিক উত্তর: ক) ৯৮ বর্গ সেমি}}

\begin{quote}
\textbf{ব্যাখ্যা:} \emph{বৃত্তের ব্যাসার্ধ r=7 সেমি হলে, বৃত্তের ব্যাস = 2×7=14
সেমি。}

\emph{বৃত্তে অন্তলিখিত বর্গক্ষেত্রের কর্ণ = বৃত্তের ব্যাস = 14 সেমি。}

\emph{∴ বর্গক্ষেত্রের ক্ষেত্রফল = (1/2)×(কর্ণ)\^{}2}

\emph{= (1/2)×(14)\^{}2=(1/2)×196=98 বর্গ সেমি।}
\end{quote}

\textbf{কোনো বৃত্তের ব্যাস ১৪ সে.মি বৃত্তটির অন্তলিখিত বৃহত্তম বর্গক্ষেত্রের
ক্ষেত্রফল কত?}

ক) 102

খ) 96

\textbf{\hl{গ) 98 ✔}}

ঘ) 100

\textbf{\hl{✔ সঠিক উত্তর: গ) 98}}

\begin{quote}
\textbf{ব্যাখ্যা:} \emph{এক্ষেত্রে বৃত্তটির ব্যাস হলো বর্গের কর্ণ = 14 সে.মি.}

\emph{∴ বর্গের ক্ষেত্রফল = (1/2)×(কর্ণ)\^{}2}

\emph{= (1/2)×(14)\^{}2=(1/2)×196=98 বর্গ সেমি.}
\end{quote}

\textbf{একটি বৃত্তের পরিধি ও ব্যাসের পার্থক্য 60cm হলে বৃত্তটির ব্যাসার্ধ কত?}

ক) 7cm

খ) 9cm

গ) 10cm

\textbf{\hl{ঘ) 14cm ✔}}

\textbf{\hl{✔ সঠিক উত্তর: ঘ) 14cm}}

\begin{quote}
\textbf{ব্যাখ্যা:} \emph{ধরি, ব্যাসার্ধ = r}

\emph{প্রশ্নমতে, 2πr−2r=60}

\emph{⇒ 2r(π−1)=60}

\emph{⇒ 2r((22/7)−1)=60}

\emph{⇒ 2r×(15/7)=60}

\emph{⇒ (30r/7)=60}

\emph{⇒ 30r=420}

\emph{∴ r=14 cm.}
\end{quote}

\textbf{13 সে.মি. ব্যাসার্ধ বিশিষ্ট বৃত্তের একটি জ্যা-এর দৈর্ঘ্য 24 সে.মি. হলে
কেন্দ্র থেকে উক্ত জ্যা-এর লম্ব দূরত্ব কত সে.মি.?}

ক) ৩

খ) ৪

\textbf{\hl{গ) ৫ ✔}}

ঘ) ৬

\textbf{\hl{✔ সঠিক উত্তর: গ) ৫}}

\begin{quote}
\textbf{ব্যাখ্যা:} \emph{বৃত্তের ব্যাসার্ধ OA=13 সে.মি., জ্যা AB=24 সে.মি.}

\emph{∴ AC=(24/2)=12 সে.মি. {[}বৃত্তের কেন্দ্র হতে জ্যা এর উপর অঙ্কিত লম্ব
জ্যা-কে সমদ্বিখণ্ডিত করে{]}}

\emph{OAC সমকোণী ত্রিভুজ হতে,}

\emph{OC=√((13)\^{}2−(12)\^{}2)}

\emph{= √(169−144)=√(25)=5 সে.মি.}
\end{quote}

\textbf{2 সেমি বাহুবিশিষ্ট একটি বর্গের অভ্যন্তরে অন্তর্বৃত্ত অঙ্কিত হলো। বৃত্তদ্বারা
বর্গের অনধিকৃত অংশের ক্ষেত্রফল কত বর্গ সে.মি.?}

ক) π−4

খ) 4−π\^{}2

\textbf{\hl{গ) 4−π ✔}}

ঘ) 2−π

\textbf{\hl{✔ সঠিক উত্তর: গ) 4−π}}

\begin{quote}
\textbf{ব্যাখ্যা:} \emph{2 সেমি বাহুর দৈর্ঘ্য বিশিষ্ট বর্গের ক্ষেত্রফল = 2\^{}2
বর্গ সে.মি. = 4 বর্গ সে.মি.}

\emph{বৃত্তের ব্যাসার্ধ বর্গের বাহুর দৈর্ঘ্যের অর্ধেক = (2/2)=1 সেমি}

\emph{বৃত্তের ক্ষেত্রফল = πr\^{}2=π×1\^{}2=π বর্গ সে.মি.}

\emph{∴ বৃত্ত দ্বারা বর্গের অনধিকৃত অংশের ক্ষেত্রফল = (4−π) বর্গ সে.মি.}
\end{quote}

\textbf{2 সে. মি. ব্যাসার্ধবিশিষ্ট একটি বৃত্তের অন্তঃস্থ একটি বর্গক্ষেত্রের চারটি
বাহু এবং বৃত্তটি দ্বারা আবদ্ধ অঞ্চলের ক্ষেত্রফল কত বর্গ সে. মি.?}

\textbf{\hl{ক) 4π−8 ✔}}

খ) 4π+8

গ) 2π−4

ঘ) 2π+4

\textbf{\hl{✔ সঠিক উত্তর: ক) 4π−8}}

\begin{quote}
\textbf{ব্যাখ্যা:} \emph{বৃত্তের ব্যাসার্ধ r=2 সে.মি. হলে বৃত্তের ব্যাস = 2×2=4
সে.মি., যা অন্তলিখিত বর্গক্ষেত্রের কর্ণ。}

\emph{বর্গক্ষেত্রের ক্ষেত্রফল = (1/2)×(কর্ণ)\^{}2=(1/2)×4\^{}2=8 বর্গ সে.মি.}

\emph{বৃত্তের ক্ষেত্রফল = πr\^{}2=π×2\^{}2=4π বর্গ সে.মি.}

\emph{∴ আবদ্ধ অঞ্চলের ক্ষেত্রফল (বৃত্তের ক্ষেত্রফল − বর্গের ক্ষেত্রফল) = (4π−8) বর্গ
সে.মি.}
\end{quote}

\textbf{একটি বৃত্তের ক্ষেত্রফল ১৬ বর্গমিটার, পরিধি ৮ মিটার। এর ব্যাসার্ধ কত?}

ক) ২ মিটার

খ) ৩ মিটার

\textbf{\hl{গ) ৪ মিটার ✔}}

ঘ) ২ (১/২) মিটার

\textbf{\hl{✔ সঠিক উত্তর: গ) ৪ মিটার}}

\begin{quote}
\textbf{ব্যাখ্যা:} \emph{আমরা জানি, বৃত্তের ক্ষেত্রফল πr\^{}2=16 এবং পরিধি
2πr=8}

\emph{∴ ব্যাসার্ধ r=(2×ক্ষেত্রফল/পরিধি)=(2×16/8)=(32/8)=4 মিটার।}
\end{quote}

\textbf{কোন বৃত্তের ব্যাসার্ধ যদি ২০\% কমে, তবে উক্ত বৃত্তের ক্ষেত্রফল কত \%
কমবে?}

ক) ১০\%

খ) ২০\%

\textbf{\hl{গ) ৩৬\% ✔}}

ঘ) ৪০\%

\textbf{\hl{✔ সঠিক উত্তর: গ) ৩৬\%}}

\begin{quote}
\textbf{ব্যাখ্যা:} \emph{ধরি, বৃত্তের আদি ব্যাসার্ধ 100 একক}

\emph{তাহলে আদি ক্ষেত্রফল = π×(100)\^{}2=10000π বর্গ একক}

\emph{২০\% কমলে নতুন ব্যাসার্ধ হয় 80 একক}

\emph{নতুন ক্ষেত্রফল = π×(80)\^{}2=6400π বর্গ একক}

\emph{ক্ষেত্রফল কমে = 10000π−6400π=3600π বর্গ একক}

\emph{∴ ক্ষেত্রফলের হ্রাসের হার = (3600π/10000π)×100\%=36\%}
\end{quote}

\textbf{একটি বৃত্তের ব্যাসার্ধ ১৬ সেমি এবং একটি বর্গক্ষেত্রের ক্ষেত্রফল উক্ত বৃত্তের
ক্ষেত্রফলের সমান হলে বর্গক্ষেত্রটির দৈর্ঘ্য কত?}

\textbf{\hl{ক) ২৮.৩৫ সেমি ✔}}

খ) ৮ সেমি

গ) ৪ সেমি

ঘ) ২৫.৮১ সেমি

\textbf{\hl{✔ সঠিক উত্তর: ক) ২৮.৩৫ সেমি}}

\begin{quote}
\textbf{ব্যাখ্যা:} \emph{বৃত্তের ক্ষেত্রফল =
πr\^{}2=π×(16)\^{}2=256π≈3.1416×256=804.2496 বর্গ সেমি}

\emph{প্রশ্নানুসারে, বর্গক্ষেত্রের ক্ষেত্রফল = 804.2496 বর্গ সেমি}

\emph{∴ বর্গক্ষেত্রটির বাহুর দৈর্ঘ্য = √(804.2496)≈28.35 সেমি।}
\end{quote}

\textbf{বৃত্তের ব্যাস ৩ গুণ বৃদ্ধি পেলে ক্ষেত্রফল কতগুণ বৃদ্ধি পাবে?}

ক) ৮

\textbf{\hl{খ) ৯ ✔}}

গ) ১২

ঘ) ১৬

\textbf{\hl{✔ সঠিক উত্তর: খ) ৯}}

\begin{quote}
\textbf{ব্যাখ্যা:} \emph{আমরা জানি, বৃত্তের ব্যাস ৩ গুণ বৃদ্ধি পেলে ব্যাসার্ধও ৩
গুণ বৃদ্ধি পায়।}

\emph{বৃত্তের ক্ষেত্রফল = πr\^{}2 {[}যেখানে r = ব্যাসার্ধ{]}}

\emph{ব্যাসার্ধ ৩ গুণ হলে নতুন ক্ষেত্রফল = π(3r)\^{}2=9πr\^{}2}

\emph{∴ ক্ষেত্রফল ৯ গুণ বৃদ্ধি পাবে।}
\end{quote}

\textbf{একটি বৃত্তের ব্যাসার্ধ যদি r থেকে বৃদ্ধি করে (r+n) করা হলে তার ক্ষেত্রফল
দ্বিগুণ হয়, তবে r এর মান কত?}

\textbf{\hl{ক) (√(2)+1)n ✔}}

খ) (√(2)−1)n

গ) √(2)n+n

ঘ) √(2)n+1

\textbf{\hl{✔ সঠিক উত্তর: ক) (√(2)+1)n}}

\begin{quote}
\textbf{ব্যাখ্যা:} \emph{প্রশ্নমতে, π((r+n))\^{}2=2πr\^{}2}

\emph{⇒ ((r+n))\^{}2=2r\^{}2}

\emph{⇒ r+n=√(2)r}

\emph{⇒ √(2)r−r=n}

\emph{⇒ r(√(2)−1)=n}

\emph{∴ r=(n/√(2)−1)=(√(2)+1)n}
\end{quote}

\textbf{একটি ঘোড়ার গাড়ির সামনের চাকার পরিধি ৪ মিটার, পেছনের চাকার পরিধি
৫ মিটার। গাড়িটি কত পথ গেলে সামনের চাকা পেছনের চাকার চেয়ে ২০০ বার বেশি
ঘুরবে?}

ক) ১.২ কিমি

খ) ২.৫ কিমি

\textbf{\hl{গ) ৪ কিমি ✔}}

ঘ) ৬ কিমি

\textbf{\hl{✔ সঠিক উত্তর: গ) ৪ কিমি}}

\begin{quote}
\textbf{ব্যাখ্যা:} \emph{ধরি, গাড়িটি x মিটার পথ অতিক্রম করে।}

\emph{প্রশ্নমতে, (x/4)−(x/5)=200}

\emph{⇒ (5x−4x/20)=200}

\emph{⇒ (x/20)=200}

\emph{⇒ x=4000 মিটার = 4 কিমি।}
\end{quote}

\textbf{একটি গাড়ির সামনের চাকার পরিধি ৩ মিটার, পেছনের চাকার পরিধি ৪
মিটার। গাড়ির কত কিলোমিটার গেলে সামনের চাকা পেছনের চাকা থেকে ১০০ বার বেশি
ঘুরবে?}

ক) ১২০০

খ) ১১

গ) ৩

\textbf{\hl{ঘ) ১.২ ✔}}

\textbf{\hl{✔ সঠিক উত্তর: ঘ) ১.২}}

\begin{quote}
\textbf{ব্যাখ্যা:} \emph{ধরি, গাড়িটি x মিটার পথ অতিক্রম করে।}

\emph{প্রশ্নমতে, (x/3)−(x/4)=100}

\emph{⇒ (4x−3x/12)=100}

\emph{⇒ (x/12)=100}

\emph{⇒ x=1200}

\emph{⇒ x=1200 মিটার = 1.2 কিলোমিটার ।}
\end{quote}

\textbf{একটি গাড়ির সামনের চাকার পরিধি ৭ ফুট, পেছনের চাকার পরিধি ৯ ফুট।
সামনের চাকা যখন পেছনের চাকার চেয়ে ১০ বার বেশি ঘোরে তখন গাড়িটি কত ফুট
যায়?}

ক) ১২৬

খ) ১৩০

গ) ১৮৯

\textbf{\hl{ঘ) ৩১৫ ✔}}

\textbf{\hl{✔ সঠিক উত্তর: ঘ) ৩১৫}}

\begin{quote}
\textbf{ব্যাখ্যা:} \emph{ধরি, পেছনের চাকা ঘোরে k বার, অতিক্রান্ত দূরত্ব = 9k
ফুট}

\emph{তাহলে সামনের চাকা ঘোরে (k+10) বার, অতিক্রান্ত দূরত্ব = 7(k+10) ফুট}

\emph{শর্তমতে, 9k=7(k+10)}

\emph{⇒ 9k=7k+70}

\emph{⇒ 2k=70}

\emph{⇒ k=35}

\emph{∴ অতিক্রান্ত দূরত্ব = 9×35=315 ফুট।}
\end{quote}

\textbf{৫৬ ফুট ব্যাসের বৃত্তাকার ক্ষেত্রকে একই ক্ষেত্রফলের একটি বর্গক্ষেত্র করলে,
বর্গক্ষেত্রের যে কোনো এক দিকের দৈর্ঘ্য কত হবে?}

ক) ৩৬.৮ বর্গফুট

খ) ২৮ বর্গফুট

\textbf{\hl{গ) ৪৯.৬ বর্গফুট ✔}}

ঘ) ৪৪ বর্গফুট

\textbf{\hl{✔ সঠিক উত্তর: গ) ৪৯.৬ বর্গফুট}}

\begin{quote}
\textbf{ব্যাখ্যা:} \emph{বৃত্তের ব্যাসার্ধ r=(56/2)=28 ফুট}

\emph{বৃত্তের ক্ষেত্রফল = πr\^{}2=(22/7)×(28)\^{}2=(22/7)×784=2464 বর্গফুট}

\emph{প্রশ্নানুসারে, বর্গক্ষেত্রের ক্ষেত্রফল = 2464 বর্গফুট}

\emph{∴ বর্গক্ষেত্রের এক বাহুর দৈর্ঘ্য = √(2464)≈49.64 ফুট (বা ৪৯.৬ ফুট)।}
\end{quote}

\textbf{একটি বৃত্তের ক্ষেত্রফল ৯ গুণ বৃদ্ধি করলে ব্যাসার্ধ কত গুণ বৃদ্ধি পাবে?}

\textbf{\hl{ক) ৩ গুণ ✔}}

খ) ৬ গুণ

গ) ৯ গুণ

ঘ) ১৮ গুণ

\textbf{\hl{✔ সঠিক উত্তর: ক) ৩ গুণ}}

\begin{quote}
\textbf{ব্যাখ্যা:} \emph{ধরি, প্রাথমিক ব্যাসার্ধ r\_1 এবং পরিবর্তিত ব্যাসার্ধ
r\_2}

\emph{ক্ষেত্রফলের অনুপাত = (πr\_1\^{}2/πr\_2\^{}2)=(1/9)}

\emph{⇒ (((r\_1/r\_2)))\^{}2=(((1/3)))\^{}2}

\emph{⇒ (r\_1/r\_2)=(1/3)}

\emph{∴ r\_2=3r\_1 (অর্থাৎ ব্যাসার্ধ ৩ গুণ বৃদ্ধি পাবে)।}
\end{quote}

\textbf{একটি বৃত্তের ব্যাসার্ধ শতকরা ৫০\% ভাগ বৃদ্ধি করলে বৃত্তের ক্ষেত্রফল বৃদ্ধি
পাবে-}

ক) ৫০\%

খ) ২৫\%

গ) ৪৫\%

\textbf{\hl{ঘ) ১২৫\% ✔}}

\textbf{\hl{✔ সঠিক উত্তর: ঘ) ১২৫\%}}

\begin{quote}
\textbf{ব্যাখ্যা:} \emph{ধরি, বৃত্তের আদি ব্যাসার্ধ 100 একক, আদি ক্ষেত্রফল =
π×(100)\^{}2=10000π}

\emph{৫০\% বৃদ্ধিতে নতুন ব্যাসার্ধ = 150 একক, নতুন ক্ষেত্রফল =
π×(150)\^{}2=22500π}

\emph{ক্ষেত্রফল বৃদ্ধি পায় = 22500π−10000π=12500π}

\emph{∴ ক্ষেত্রফলের বৃদ্ধির হার = (12500π/10000π)×100\%=125\%}
\end{quote}

\textbf{একটি ফুটবলের ব্যাস ১০ ইঞ্চি হইলে, ফুটবলটির আয়তন কত?}

ক) (4/3)πr\^{}3

\textbf{\hl{খ) (500/3)π ✔}}

গ) 500π

ঘ) (125/3)π

\textbf{\hl{✔ সঠিক উত্তর: খ) (500/3)π}}

\begin{quote}
\textbf{ব্যাখ্যা:} \emph{ফুটবল একটি গোলক। গোলকের ব্যাসার্ধ r=(10/2)=5 ইঞ্চি।}

\emph{ফুটবলের আয়তন = (4/3)πr\^{}3=(4/3)×π×5\^{}3=(4/3)×125π=(500/3)π ঘন
ইঞ্চি।}
\end{quote}

\textbf{একটি বৃত্তের ব্যাসার্ধকে যদি r হইতে বৃদ্ধি করিয়া r+n করা হয় তবে উহার
ক্ষেত্রফল দ্বিগুণ হয়। r-এর মান কত?}

\textbf{\hl{ক) (n/√(2)−1) ✔}}

খ) (n/√(2)+1)

গ) (√(2)/n+1)

ঘ) (√(2)+1/n)

\textbf{\hl{✔ সঠিক উত্তর: ক) (n/√(2)−1)}}

\begin{quote}
\textbf{ব্যাখ্যা:} \emph{প্রশ্নমতে, π((r+n))\^{}2=2πr\^{}2}

\emph{⇒ ((r+n))\^{}2=2r\^{}2}

\emph{⇒ r+n=√(2)r}

\emph{⇒ √(2)r−r=n}

\emph{⇒ r(√(2)−1)=n}

\emph{∴ r=(n/√(2)−1)}
\end{quote}

\textbf{বৃত্তের ব্যাস চারগুণ বৃদ্ধি করলে ক্ষেত্রফল কতগুণ বৃদ্ধি পাবে?}

ক) ৪ গুণ

খ) ৮ গুণ

গ) ১২ গুণ

\textbf{\hl{ঘ) ১৬ গুণ ✔}}

\textbf{\hl{✔ সঠিক উত্তর: ঘ) ১৬ গুণ}}

\begin{quote}
\textbf{ব্যাখ্যা:} \emph{ধরি, বৃত্তের আদি ব্যাসার্ধ r , আদি ক্ষেত্রফল = πr\^{}2}

\emph{ব্যাস চারগুণ বৃদ্ধি করলে (অর্থাৎ ৪ গুণ হলে) নতুন ব্যাসার্ধ হয় 4r}

\emph{নতুন ক্ষেত্রফল = π(4r)\^{}2=16πr\^{}2}

\emph{∴ ক্ষেত্রফল ১৬ গুণ বৃদ্ধি পাবে।}
\end{quote}

\textbf{১৪ সেন্টিমিটার ব্যাসার্ধের বৃত্তের পরিধি কত?}

ক) ৮৭.১৯৬ সেমি

\textbf{\hl{খ) ৮৭.৯৬৮ সেমি ✔}}

গ) ৮৮ সেমি

ঘ) ৪৪ সেমি

\textbf{\hl{✔ সঠিক উত্তর: খ) ৮৭.৯৬৮ সেমি}}

\begin{quote}
\textbf{ব্যাখ্যা:} \emph{আমরা জানি, বৃত্তের পরিধি = 2πr}

\emph{এখানে, ব্যাসার্ধ r=14 সেমি}

\emph{∴ পরিধি = 2×3.1416×14=87.968 সেমি।}
\end{quote}

\textbf{একটি বৃত্তের ব্যাসার্ধ দ্বিগুণ করলে বৃত্তটির ক্ষেত্রফল হবে-}

ক) দ্বিগুণ

খ) তিন গুণ

\textbf{\hl{গ) চার গুণ ✔}}

ঘ) একই থাকবে

\textbf{\hl{✔ সঠিক উত্তর: গ) চার গুণ}}

\begin{quote}
\textbf{ব্যাখ্যা:} \emph{ধরি, বৃত্তের ব্যাসার্ধ r , ক্ষেত্রফল = πr\^{}2}

\emph{ব্যাসার্ধ দ্বিগুণ করলে নতুন ব্যাসার্ধ 2r , নতুন ক্ষেত্রফল =
π(2r)\^{}2=4πr\^{}2}

\emph{∴ ক্ষেত্রফল চার গুণ হবে।}
\end{quote}

\textbf{একটি বৃত্তাকার মাঠের ব্যাস 26 মিটার। মাঠটির বাইরে চারদিকে 2 মিটার
চওড়া রাস্তা রয়েছে। রাস্তাসহ মাঠটির ক্ষেত্রফল কত বর্গমিটার?}

\textbf{\hl{ক) 225π ✔}}

খ) 169π

গ) 121π

ঘ) 144π

\textbf{\hl{✔ সঠিক উত্তর: ক) 225π}}

\begin{quote}
\textbf{ব্যাখ্যা:} \emph{মাঠটির ব্যাস 26 মিটার হলে ব্যাসার্ধ r=(26/2)=13
মিটার}

\emph{মাঠটির বাইরে চারদিকে 2 মিটার চওড়া রাস্তা থাকায় রাস্তাসহ ব্যাসার্ধ =
13+2=15 মিটার}

\emph{∴ রাস্তাসহ মাঠটির ক্ষেত্রফল = π×15\^{}2=225π বর্গমিটার।}
\end{quote}

\textbf{একটি চাকা ঘণ্টায় ১৪০০ বার ঘোরে।
\textquotesingle x\textquotesingle{} সেকেন্ডে চাকাটি কতবার ঘুরবে?}

ক) (18/7x) বার

\textbf{\hl{খ) (7x/18) বার ✔}}

গ) (18x/7) বার

ঘ) (9x/14) বার

\textbf{\hl{✔ সঠিক উত্তর: খ) (7x/18) বার}}

\begin{quote}
\textbf{ব্যাখ্যা:} \emph{১ ঘণ্টা = ৩৬০০ সেকেন্ড}

\emph{৩৬০০ সেকেন্ডে ঘোরে ১৪০০ বার}

\emph{∴ ১ সেকেন্ডে ঘোরে (1400/3600) বার}

\emph{∴ x সেকেন্ডে ঘোরে (1400×x/3600)=(7x/18) বার।}
\end{quote}

\textbf{একটি গাড়ির চাকা প্রতি মিনিটে ৯০ বার ঘোরে। এক সেকেন্ডে চাকাটি কত
ডিগ্রী ঘুরবে?}

ক) ২৭০°

খ) ১৮০°

গ) ৩৬০°

\textbf{\hl{ঘ) ৫৪০° ✔}}

\textbf{\hl{✔ সঠিক উত্তর: ঘ) ৫৪০°}}

\begin{quote}
\textbf{ব্যাখ্যা:} \emph{১ বার ঘুরলে অতিক্রম করে = ৩৬০°}

\emph{∴ ৯০ বার ঘুরলে অতিক্রম করে = 360°×90=32400°}

\emph{৬০ সেকেন্ডে চাকাটি ঘোরে = 32400°}

\emph{∴ ১ সেকেন্ডে ঘোরে = (32400°/60)=540°}
\end{quote}

\textbf{একটি গাড়ির চাকা প্রতি মিনিটে ১২ বার ঘোরে। চাকাটি পাঁচ সেকেন্ডে কত
ডিগ্রি ঘুরে?}

\textbf{\hl{ক) ৩৬০° ✔}}

খ) ৩০০°

গ) ১৮০°

ঘ) ৩০°

\textbf{\hl{✔ সঠিক উত্তর: ক) ৩৬০°}}

\begin{quote}
\textbf{ব্যাখ্যা:} \emph{১ বার ঘুরলে অতিক্রম করে = ৩৬০°}

\emph{∴ ১২ বার ঘুরলে অতিক্রম করে = 360°×12=4320°}

\emph{৬০ সেকেন্ডে চাকাটি ঘোরে = 4320°}

\emph{∴ ৫ সেকেন্ডে ঘোরে = (4320°×5/60)=360°}
\end{quote}

\textbf{একটি গাড়ির চাকা প্রতি মিনিটে ৯০ বার ঘোরে। এক সেকেন্ডে চাকাটি কত
ডিগ্রি ঘুরবে?}

ক) ২৭০°

খ) ৪২০°

গ) ১৮০°

\textbf{\hl{ঘ) ৫৪০° ✔}}

\textbf{\hl{✔ সঠিক উত্তর: ঘ) ৫৪০°}}

\begin{quote}
\textbf{ব্যাখ্যা:} \emph{১ বার ঘুরলে অতিক্রম করে = ৩৬০°}

\emph{∴ ৯০ বার ঘুরলে অতিক্রম করে = 360°×90=32400°}

\emph{৬০ সেকেন্ডে চাকাটি ঘোরে = 32400°}

\emph{∴ ১ সেকেন্ডে ঘোরে = (32400°/60)=540°}
\end{quote}

\textbf{একটি চাকার ব্যাস ৭০ সে.মি.। চাকাটি একবার ঘুরলে কতটুকু পথ অতিক্রম
করে?}

ক) ২.২৫৯ মি.

\textbf{\hl{খ) ২.১৯৯ মি. ✔}}

গ) ২.৩৫৯ মি.

ঘ) ২.১৫৯ মি.

\textbf{\hl{✔ সঠিক উত্তর: খ) ২.১৯৯ মি.}}

\begin{quote}
\textbf{ব্যাখ্যা:} \emph{চাকার ব্যাস 2r=70 সে.মি.}

\emph{∴ চাকার ব্যাসার্ধ r=(70/2)=35 সে.মি.}

\emph{চাকাটি একবার ঘুরলে পরিধির সমান পথ অতিক্রম করবে।}

\emph{∴ পরিধি = 2πr=2×3.1416×35=219.91 সে.মি. = 2.199 মি.}
\end{quote}

\textbf{একটি চাকার ব্যাস ৪.২ মিটার। চাকাটি ৩৩০ মিটার পথ অতিক্রম করতে কত
বার ঘুরবে?}

ক) ১৫

খ) ২০

\textbf{\hl{গ) ২৫ ✔}}

ঘ) ৩০

\textbf{\hl{✔ সঠিক উত্তর: গ) ২৫}}

\begin{quote}
\textbf{ব্যাখ্যা:} \emph{চাকার ব্যাস = 4.2 মিটার}

\emph{চাকার ব্যাসার্ধ = (4.2/2)=2.1 মিটার}

\emph{চাকার পরিধি = 2πr=2×(22/7)×2.1=13.2 মিটার}

\emph{∴ চাকাটি ৩৩০ মিটার অতিক্রম করলে ঘোরে = (330/13.2)=25 বার।}
\end{quote}

\textbf{একটি গাড়ির চাকার পরিধি ৫ মিটার। ১ কিলোমিটার ৫০০ মিটার পথ যেতে
চাকাটি কতবার ঘুরবে?}

ক) ২০০

খ) ২৫০

\textbf{\hl{গ) ৩০০ ✔}}

ঘ) ৩৫০

\textbf{\hl{✔ সঠিক উত্তর: গ) ৩০০}}

\begin{quote}
\textbf{ব্যাখ্যা:} \emph{১ কিলোমিটার ৫০০ মিটার = 1000+500=1500 মিটার}

\emph{চাকাটি ৫ মিটার দূরত্ব অতিক্রম করলে = 1 বার ঘোরে}

\emph{∴ চাকাটি ১,৫০০ মিটার দূরত্ব অতিক্রম করলে ঘোরে = (1500/5)=300 বার।}
\end{quote}

\textbf{একটি চাকার ব্যাস ৪.২ মিটার। চাকাটি ৩৩০ মিটার পথ অতিক্রম করতে কত
বার ঘুরবে?}

ক) ২০

\textbf{\hl{খ) ২৫ ✔}}

গ) ৩০

ঘ) ৩৫

\textbf{\hl{✔ সঠিক উত্তর: খ) ২৫}}

\begin{quote}
\textbf{ব্যাখ্যা:} \emph{চাকার ব্যাসার্ধ r=(4.2/2)=2.1 মিটার}

\emph{চাকার পরিধি = 2πr=2×3.1416×2.1=13.19 মিটার}

\emph{∴ চাকাটি ঘুরবে = (330/13.19)=25.01≈25 বার।}
\end{quote}

\textbf{একটি চাকার পরিধি ৫ মিটার। ২০ মাইল পথ অতিক্রম করতে চাকাটি কতবার
ঘুরবে? (১মাইল = ১.৬ কিলোমিটার)}

\textbf{\hl{ক) ৬৪০০ ✔}}

খ) ৫২০০

গ) ৬০০০

ঘ) ৬২০০

\textbf{\hl{✔ সঠিক উত্তর: ক) ৬৪০০}}

\begin{quote}
\textbf{ব্যাখ্যা:} \emph{২০ মাইল = 20×1.6 কিলোমিটার = 32 কিলোমিটার = 32000
মিটার}

\emph{∴ চাকাটি ঘুরবে = (32000/5)=6400 বার।}
\end{quote}

\textbf{একটি গাড়ির সামনের চাকার পরিধি ৪ মিটার, পেছনের চাকার পরিধি ৫
মিটার। গাড়িটি কত কিলোমিটার গেলে সামনের চাকা পেছনের চাকা থেকে ১০০ বার
বেশি ঘুরবে?}

ক) ২০

খ) ১০০

গ) ১২

\textbf{\hl{ঘ) ২ ✔}}

\textbf{\hl{✔ সঠিক উত্তর: ঘ) ২}}

\begin{quote}
\textbf{ব্যাখ্যা:} \emph{চাকার পরিধিগুলোর ল.সা.গু. = 4 ও 5 এর ল.সা.গু. = 20
মিটার}

\emph{২০ মিটার যেতে সামনের চাকা ঘোরে = (20/4)=5 বার}

\emph{পেছনের চাকা ঘোরে = (20/5)=4 বার}

\emph{পার্থক্য = 5−4=1 বার}

\emph{১ বার বেশি ঘোরে ২০ মিটার গেলে}

\emph{∴ ১০০ বার বেশি ঘোরে = 20×100=2000 মিটার বা 2 কিমি।}
\end{quote}

\textbf{একটি মোটরসাইকেলের পেছনের চাকা প্রতি মিনিটে ৬০ বার ঘোরে এবং
প্রতিবার ঘুরলে ১২০ সেন্টিমিটার পথ অতিক্রম করে। এক ঘণ্টায় চাকাটি কত মিটার পথ
অতিক্রম করবে?}

ক) ৪৩২০০০

খ) ৪৩২০০

\textbf{\hl{গ) ৪৩২০ ✔}}

ঘ) ৪৩২

\textbf{\hl{✔ সঠিক উত্তর: গ) ৪৩২০}}

\begin{quote}
\textbf{ব্যাখ্যা:} \emph{চাকাটি ১ মিনিটে ৬০ বার ঘুরলে ১ ঘণ্টা বা ৬০ মিনিটে
ঘোরে = 60×60=3600 বার}

\emph{প্রতিবারে ১২০ সেমি বা 1.2 মিটার ঘুরলে,}

\emph{∴ মোট অতিক্রান্ত পথ = 3600×1.2=4320 মিটার।}
\end{quote}

\textbf{দুটি বৃত্তের ব্যাসার্ধের অনুপাত ৩ : ২, বৃত্ত দুইটির ক্ষেত্রফলের অনুপাত কত
হবে?}

ক) ২ : ৩

খ) ৩ : ৪

গ) ৪ : ৯

\textbf{\hl{ঘ) ৯ : ৮ ✔}}

\textbf{\hl{✔ সঠিক উত্তর: ঘ) ৯ : ৮}}

\begin{quote}
\textbf{ব্যাখ্যা:} \emph{ধরি, প্রথম বৃত্তের ব্যাসার্ধ r\_1=3r এবং দ্বিতীয় বৃত্তের
ব্যাসার্ধ r\_2=2r}

\emph{প্রথম বৃত্তের ক্ষেত্রফল = π(3r)\^{}2=9πr\^{}2}

\emph{দ্বিতীয় বৃত্তের ক্ষেত্রফল = π(2r)\^{}2=4πr\^{}2}

\emph{∴ ক্ষেত্রফলের অনুপাত = 9πr\^{}2:4πr\^{}2=9:4}
\end{quote}

\textbf{একটি বৃত্তের ব্যাসার্ধ ১৪ সে.মি. এবং বৃত্তকলা কেন্দ্রে ৭৫° কোণ উৎপন্ন
করে। বৃত্তকলার ক্ষেত্রফল কত?}

ক) ১২৭ বর্গ সে.মি.

খ) ১২৯.০৩ বর্গ সে.মি.

\textbf{\hl{গ) ১২৮.২৮২ বর্গ সে.মি. ✔}}

ঘ) ১২৮.০৫ বর্গ সে.মি.

\textbf{\hl{✔ সঠিক উত্তর: গ) ১২৮.২৮২ বর্গ সে.মি.}}

\begin{quote}
\textbf{ব্যাখ্যা:} \emph{আমরা জানি, বৃত্তকলার ক্ষেত্রফল = (θ/360)πr\^{}2}

\emph{⇒ (75/360)×3.1416×(14)\^{}2}

\emph{⇒ 128.282 বর্গ সে.মি.}
\end{quote}

\section{48. সুদকষার জাদুকরী সমাধান: ১০০\%-এর ধারণায় পাটিগণিতের সহজ
পাঠ}\label{ux9b8ux9a6ux995ux9b7ux9b0-ux99cux9a6ux995ux9b0-ux9b8ux9aeux9a7ux9a8-ux9e7ux9e6ux9e6-ux98fux9b0-ux9a7ux9b0ux9a3ux9df-ux9aaux99fux997ux9a3ux9a4ux9b0-ux9b8ux9b9ux99c-ux9aaux9a0}

পাটিগণিতের অন্যতম ভয়ের একটি অধ্যায় হলো সুদকষা বা মুনাফার হিসাব। অনেকেই শত শত
সূত্র মুখস্থ করে পরীক্ষার হলে বিভ্রান্ত হয়ে পড়েন।

অথচ একটু গভীরভাবে খেয়াল করলে দেখবেন, লাভ-ক্ষতির মতো সুদকষার জটিল সমস্যাগুলোও
খুব সহজে শতকরা বা ১০০\%-এর ধারণা ব্যবহার করে সমাধান করা সম্ভব।

নিচে ধাপে ধাপে প্রতিটি ক্যাটাগরির ৩টি করে গুরুত্বপূর্ণ উদাহরণ লিখিত ও শর্ট নিয়মে
আলোচনা করা হলো:

\begin{itemize}
\item
  সুদ নির্ণয় সংক্রান্ত সমস্যা
\item
  সুদের হার নির্ণয় সংক্রান্ত সমস্যা
\item
  আসল/মূলধন নির্ণয় সংক্রান্ত সমস্যা
\item
  সময় নির্ণয় সংক্রান্ত সমস্যা
\item
  প্রথমে সময় তারপর আসল নির্ণয় সংক্রান্ত সমস্যা
\item
  প্রথমে সুদ তারপর সুদের হার নির্ণয় সংক্রান্ত সমস্যা
\item
  আসল না দেওয়া থাকলেও সময় নির্ণয় সংক্রান্ত সমস্যা
\item
  প্রথমে সুদ তারপর সময় নির্ণয় সংক্রান্ত সমস্যা
\end{itemize}

\subsubsection{১. সুদ নির্ণয় সংক্রান্ত
সমস্যা}\label{ux9b8ux9a6-ux9a8ux9b0ux9a3ux9df-ux9b8ux995ux9b0ux9a8ux9a4-ux9b8ux9aeux9b8ux9af}

আসল, সুদের হার এবং সময় দেওয়া থাকলে মোট সুদ কত হবে তা বের করার নিয়ম.

\textbf{উদাহরণ ১: ৫\% সরল সুদে ৫০০ টাকার ৩ বছরের সুদ কত?}

\begin{itemize}
\item
  লিখিত নিয়ম: ১ বছরের সুদ = ৫০০ এর ৫\% = ২৫ টাকা। সুতরাং ৩ বছরের সুদ =
  25×3=75 টাকা.
\item
  শর্ট নিয়ম: সরাসরি সূত্র বা এক লাইনে: 500×(5/100)×3=75 টাকা.
\end{itemize}

\textbf{উদাহরণ ২: ৮\% হার সুদে ১০০০ টাকার ২ বছরের সুদ কত?}

\begin{itemize}
\item
  লিখিত নিয়ম: ১০০ টাকায় ১ বছরের সুদ ৮ টাকা। ১ টাকায় ১ বছরের সুদ (8/100)
  টাকা। সুতরাং ১০০০ টাকায় ২ বছরের সুদ = (8×1000×2/100)=160 টাকা.
\item
  শর্ট নিয়ম: 1000×(8/100)×2=160 টাকা.
\end{itemize}

\textbf{উদাহরণ ৩: ১০\% হার সুদে ৮০০ টাকার ৫ বছরের সুদ কত?}

\begin{itemize}
\item
  লিখিত নিয়ম: ১০০ টাকার ১ বছরের সুদ ১০ টাকা, সুতরাং ৮০০ টাকার ১ বছরের সুদ =
  800×(10/100)=80 টাকা। ৫ বছরের সুদ = 80×5=400 টাকা.
\item
  শর্ট নিয়ম: 800×(10/100)×5=400 টাকা.
\end{itemize}

\subsubsection{২. সুদের হার নির্ণয় সংক্রান্ত
সমস্যা}\label{ux9b8ux9a6ux9b0-ux9b9ux9b0-ux9a8ux9b0ux9a3ux9df-ux9b8ux995ux9b0ux9a8ux9a4-ux9b8ux9aeux9b8ux9af}

আসল, সময় এবং মোট সুদ দেওয়া থাকলে বার্ষিক সুদের হার নির্ণয় করা.

\textbf{উদাহরণ ১: কত হার সুদে ৫০০ টাকার ৩ বছরের সুদ ৭৫ টাকা হবে?}

\begin{itemize}
\item
  লিখিত নিয়ম: ৫০০ টাকার ৩ বছরের সুদ ৭৫ টাকা হলে, ১ বছরের সুদ = (75/3)=25
  টাকা। এখন ৫০০ টাকায় সুদ হয় ২৫ টাকা, তবে ১০০ টাকায় কত? অর্থাৎ,
  (25/500)×100=5\% .
\item
  শর্ট নিয়ম: সুদের হার = (75×100/500×3)=5\% .
\end{itemize}

\textbf{উদাহরণ ২: বার্ষিক কত হার সুদে ১০০০ টাকার ২ বছরের সুদ ১৬০ টাকা হবে?}

\begin{itemize}
\item
  লিখিত নিয়ম: ১০০০ টাকার ২ বছরের সুদ ১৬০ টাকা, ১ বছরের সুদ ৮০ টাকা। ১০০০
  টাকায় ৮০ টাকা হলে ১০০ টাকায় ৮ টাকা। সুতরাং হার ৮\%.
\item
  শর্ট নিয়ম: হার = (160×100/1000×2)=8\% .
\end{itemize}

\textbf{উদাহরণ ৩: বার্ষিক কত হার সুদে ৮০০ টাকার ৫ বছরের সুদ ৪০০ টাকা হবে?}

\begin{itemize}
\item
  লিখিত নিয়ম: ৮০০ টাকার ৫ বছরের সুদ ৪০০ টাকা, ১ বছরের সুদ ৮০ টাকা। ৮০০
  টাকার মধ্যে ৮০ টাকা মানে ১০০ টাকার মধ্যে ১০ টাকা। সুদের হার ১০\%.
\item
  শর্ট নিয়ম: হার = (400×100/800×5)=10\% .
\end{itemize}

\subsubsection{৩. আসল/মূলধন নির্ণয় সংক্রান্ত
সমস্যা}\label{ux986ux9b8ux9b2ux9aeux9b2ux9a7ux9a8-ux9a8ux9b0ux9a3ux9df-ux9b8ux995ux9b0ux9a8ux9a4-ux9b8ux9aeux9b8ux9af}

সুদের হার, সময় এবং মোট সুদ জানা থাকলে মূলধন বা আসল কত ছিল তা বের করা.

\textbf{উদাহরণ ১: বার্ষিক ৫\% সুদে কত টাকার ৩ বছরের সুদ ৭৫ টাকা হবে?}

\begin{itemize}
\item
  লিখিত নিয়ম: ৫\% সুদে ১ বছরের সুদ ৫ টাকা হলে ৩ বছরের সুদ ১৫ টাকা। ১৫ টাকা
  সুদ হলে আসল ১০০ টাকা। ৭৫ টাকা সুদ হলে আসল = (100×75/15)=500 টাকা.
\item
  শর্ট নিয়ম: আসল = (75×100/5×3)=500 টাকা.
\end{itemize}

\textbf{উদাহরণ ২: বার্ষিক ৮\% হার সুদে কত টাকার ২ বছরের সুদ ১৬০ টাকা হবে?}

\begin{itemize}
\item
  লিখিত নিয়ম: ৮\% সুদে ২ বছরে সুদ হয় ১৬\%। ১৬\%-এর মান ১৬০ টাকা হলে,
  ১\%-এর মান ১০ টাকা। সুতরাং ১০০\% বা আসল = ১০ × ১০০ = ১০০০ টাকা.
\item
  শর্ট নিয়ম: আসল = (160×100/8×2)=1000 টাকা.
\end{itemize}

\textbf{উদাহরণ ৩: বার্ষিক ১০\% সুদে কত টাকার ৫ বছরের সুদ ৪০০ টাকা হবে?}

\begin{itemize}
\item
  লিখিত নিয়ম: ১০\% সুদে ৫ বছরের মোট সুদ ৫০\%। ৫০\% এর মান ৪০০ টাকা হলে,
  ১০০\% বা আসল হবে ৮০০ টাকা.
\item
  শর্ট নিয়ম: আসল = (400×100/10×5)=800 টাকা.
\end{itemize}

\subsubsection{৪. সময় নির্ণয় সংক্রান্ত
সমস্যা}\label{ux9b8ux9aeux9df-ux9a8ux9b0ux9a3ux9df-ux9b8ux995ux9b0ux9a8ux9a4-ux9b8ux9aeux9b8ux9af}

আসল, সুদের হার এবং মোট সুদ দেওয়া থাকলে কত সময়ে তা অর্জিত হয়েছে তা বের করা.

\textbf{উদাহরণ ১: বার্ষিক ৫\% সুদে ৫০০ টাকার কত বছরের সুদ ৭৫ টাকা হবে?}

\begin{itemize}
\item
  লিখিত নিয়ম: ৫০০ টাকার ৫\% হারে ১ বছরের সুদ ২৫ টাকা। ২৫ টাকা সুদ হয় ১
  বছরে, ৭৫ টাকা সুদ হতে সময় লাগবে (75/25)=3 বছর.
\item
  শর্ট নিয়ম: সময় = (75×100/500×5)=3 বছর.
\end{itemize}

\textbf{উদাহরণ ২: বার্ষিক ৮\% সুদে ১০০০ টাকার কত বছরের সুদ ১৬০ টাকা হবে?}

\begin{itemize}
\item
  লিখিত নিয়ম: ১০০০ টাকার ৮\% হারে ১ বছরের সুদ ৮০ টাকা। ১৬০ টাকা সুদ হতে
  সময় লাগবে (160/80)=2 বছর.
\item
  শর্ট নিয়ম: সময় = (160×100/1000×8)=2 বছর.
\end{itemize}

\textbf{উদাহরণ ৩: বার্ষিক ১০\% সুদে ৮০০ টাকার কত বছরের সুদ ৪০০ টাকা হবে?}

\begin{itemize}
\item
  লিখিত নিয়ম: ৮০০ টাকার ১০\% হারে ১ বছরের সুদ ৮০ টাকা। ৪০০ টাকা সুদ হতে
  সময় লাগবে (400/80)=5 বছর.
\item
  শর্ট নিয়ম: সময় = (400×100/800×10)=5 বছর.
\end{itemize}

\subsubsection{৫. প্রথমে সময় তারপর আসল নির্ণয় সংক্রান্ত
সমস্যা}\label{ux9aaux9b0ux9a5ux9ae-ux9b8ux9aeux9df-ux9a4ux9b0ux9aaux9b0-ux986ux9b8ux9b2-ux9a8ux9b0ux9a3ux9df-ux9b8ux995ux9b0ux9a8ux9a4-ux9b8ux9aeux9b8ux9af}

যখন প্রশ্নে সুদের হার এবং নির্দিষ্ট সুদের শর্ত দিয়ে প্রথমে সময় বা সুদের পরিমাণ
বের করে পরবর্তীতে মূলধন নির্ণয় করতে হয়.

\textbf{উদাহরণ ১:} কোনো আসল ৫\% হার সুদে যত সময়ে ৩০০ টাকা সুদ দেয়, সেই
সময়ে ৮\% হার সুদে কত টাকার সুদ ৪৮০ টাকা হবে? (ধরি সময় ৬ বছর)

\begin{itemize}
\item
  লিখিত নিয়ম: প্রথমে ৩০০ টাকার ক্ষেত্রে সময় বের করি: সময় = (300×100/আসল×5)
  (আসল অজানা থাকলে ধরি আসল ১০০০, সেক্ষেত্রে সময় ৬ বছর)। এবার ৬ বছরে ৮\%
  সুদে ৪৮০ টাকা সুদ পেতে আসল কত হবে তা বের করতে হবে: আসল =
  (480×100/8×6)=1000 টাকা.
\item
  শর্ট নিয়ম: সুদের অনুপাতে বা একক নিয়মে সরাসরি সুদের হার দিয়ে সমীকরণ সাজিয়ে
  আসল বের করা যায়.
\end{itemize}

\textbf{উদাহরণ ২: বার্ষিক ৬\% সরল সুদে কোনো নির্দিষ্ট টাকার ৪ বছরের সুদ যত
হয়, ১২\% সুদে কত টাকার ওই একই সময়ে দ্বিগুণ সুদ হবে?}

\begin{itemize}
\item
  লিখিত নিয়ম: ধরি প্রথম ক্ষেত্রে আসল ১০০ টাকা, ৪ বছরে ৬\% সুদে সুদ = ২৪ টাকা।
  দ্বিতীয় ক্ষেত্রে দ্বিগুণ বা ৪৮ টাকা সুদ পেতে এবং সময় ৪ বছর হলে আসল হবে:
  (48×100/12×4)=100 টাকার দ্বিগুণ অর্থাৎ ২০০ টাকা.
\item
  শর্ট নিয়ম: সুদের হার দ্বিগুণ হলে এবং সময় এক থাকলে আসলও দ্বিগুণ অর্থাৎ ২০০
  টাকা হবে.
\end{itemize}

\textbf{উদাহরণ ৩: একটি নির্দিষ্ট আসল ৪ বছরের সুদে-আসলে ৬০০ টাকা হয়। সুদের
হার দ্বিগুণ হলে ওই একই আসল কত বছরে বা কত সুদে পরিণত হবে?}

\begin{itemize}
\item
  লিখিত নিয়ম: সুদের হার দ্বিগুণ হলে অতিরিক্ত সুদ আসলের সমান বা নির্দিষ্ট অনুপাতে
  বৃদ্ধি পাবে। ধরি আসল ৪০০ এবং সুদ ২০০ টাকা, হার ৫\% হলে ৪ বছর। হার ১০\%
  হলে ৪ বছরে সুদ দ্বিগুণ অর্থাৎ ৪০০ টাকা হবে, ফলে সুদের-আসলে হবে ৮০০ টাকা.
\item
  শর্ট নিয়ম: সুদের হার দ্বিগুণ হলে সুদের পরিমাণও দ্বিগুণ হবে.
\end{itemize}

\subsubsection{৬. প্রথমে সুদ তারপর সুদের হার নির্ণয় সংক্রান্ত
সমস্যা}\label{ux9aaux9b0ux9a5ux9ae-ux9b8ux9a6-ux9a4ux9b0ux9aaux9b0-ux9b8ux9a6ux9b0-ux9b9ux9b0-ux9a8ux9b0ux9a3ux9df-ux9b8ux995ux9b0ux9a8ux9a4-ux9b8ux9aeux9b8ux9af}

প্রথমে মোট সুদ বা সুদের পরিমাণ বের করে পরবর্তীতে বার্ষিক সুদের হার কত তা নির্ণয়
করা.

\textbf{উদাহরণ ১: একটি নির্দিষ্ট পরিমাণ টাকার ৪ বছরের সুদ আসলের}
\(\frac{1}{5}\) \textbf{অংশ হলে, সুদের হার কত?}

\begin{itemize}
\item
  লিখিত নিয়ম: ধরি আসল ১০০ টাকা, তবে সুদ হবে ১০০ এর (1/5) বা ২০ টাকা। ৪
  বছরের সুদ ২০ টাকা হলে ১ বছরের সুদ ৫ টাকা। সুদের হার ৫\%.
\item
  শর্ট নিয়ম: হার = ((1/5)×100/4)=5\% .
\end{itemize}

\textbf{উদাহরণ ২: আসল ও ৫ বছরের সুদের অনুপাত ১০:৩ হলে, সুদের হার কত?}

\begin{itemize}
\item
  লিখিত নিয়ম: আসল ১০ টাকা হলে ৫ বছরের সুদ ৩ টাকা। ১ বছরের সুদ (3/5) টাকা।
  ১০০ টাকায় ১ বছরের সুদ হবে (3/5)×(100/10)=6 টাকা। সুদের হার ৬\%.
\item
  শর্ট নিয়ম: হার = (3/10)×(100/5)=6\% .
\end{itemize}

\textbf{উদাহরণ ৩: কোনো টাকার ৩ বছরের সুদ আসলের} \(\frac{3}{10}\)
\textbf{অংশ হলে, বার্ষিক সুদের হার কত?}

\begin{itemize}
\item
  লিখিত নিয়ম: আসল যদি ১০ টাকা হয়, সুদ ৩ টাকা। ৩ বছরের সুদ ৩ টাকা হলে ১
  বছরের সুদ ১ টাকা। ১০ টাকায় ১ বছরের সুদ ১ টাকা হলে, ১০০ টাকায় ১০ টাকা।
  সুদের হার ১০\%.
\item
  শর্ট নিয়ম: হার = ((3/10)×100/3)=10\% .
\end{itemize}

\subsubsection{৭. আসল না দেওয়া থাকলেও সময় নির্ণয় সংক্রান্ত
সমস্যা}\label{ux986ux9b8ux9b2-ux9a8-ux9a6ux993ux9df-ux9a5ux995ux9b2ux993-ux9b8ux9aeux9df-ux9a8ux9b0ux9a3ux9df-ux9b8ux995ux9b0ux9a8ux9a4-ux9b8ux9aeux9b8ux9af}

আসল বা মূলধনের পরিমাণ সরাসরি উল্লেখ না থাকলেও নির্দিষ্ট শর্ত বা গুণিতক দেওয়া
থাকলে সময় বের করা.

\textbf{উদাহরণ ১: বার্ষিক ৫\% সরল সুদে কোনো নির্দিষ্ট পরিমাণ টাকা কত বছরে
আসলের দ্বিগুণ হবে?}

\begin{itemize}
\item
  লিখিত নিয়ম: আসল ১০০ টাকা হলে দ্বিগুণ অর্থাৎ সুদে-আসলে হবে ২০০ টাকা। সুতরাং
  সুদ হবে 200-100=100 টাকা। ৫\% সুদে ১০০ টাকা সুদ হতে সময় লাগবে (100/5)=20
  বছর.
\item
  শর্ট নিয়ম: সময় = ((2-1)×100/5)=20 বছর.
\end{itemize}

\textbf{উদাহরণ ২: বার্ষিক ১০\% সরল সুদে কত বছরে কোনো মূলধন তিনগুণ হবে?}

\begin{itemize}
\item
  লিখিত নিয়ম: আসল ১০০ টাকা হলে তিনগুণ মানে সুদ-আসল ৩০০ টাকা, সুদ ২০০ টাকা।
  ১০\% হারে ২০০ টাকা সুদ হতে সময় লাগে (200/10)=20 বছর.
\item
  শর্ট নিয়ম: সময় = ((3-1)×100/10)=20 বছর.
\end{itemize}

\textbf{উদাহরণ ৩: বার্ষিক ৮\% সরল সুদে কত বছরে যে কোনো আসল আসলের চারগুণ
হবে?}

\begin{itemize}
\item
  লিখিত নিয়ম: আসল ১০০ টাকা হলে সুদ হবে ৩০০ টাকা। ৮\% হারে ৩০০ টাকা সুদ হতে
  সময় লাগবে (300/8)=37.5 বছর.
\item
  শর্ট নিয়ম: সময় = ((4-1)×100/8)=(300/8)=37.5 বছর.
\end{itemize}

\subsubsection{৮. প্রথমে সুদ তারপর সময় নির্ণয় সংক্রান্ত
সমস্যা}\label{ux9aaux9b0ux9a5ux9ae-ux9b8ux9a6-ux9a4ux9b0ux9aaux9b0-ux9b8ux9aeux9df-ux9a8ux9b0ux9a3ux9df-ux9b8ux995ux9b0ux9a8ux9a4-ux9b8ux9aeux9b8ux9af}

প্রথমে সুদের পরিমাণ বা শর্ত থেকে সুদের অংশ বের করে পরবর্তীতে নির্দিষ্ট হারে সময়
নির্ণয় করা.

\textbf{উদাহরণ ১:} বার্ষিক ৫\% সরল সুদে কত সময়ে সুদে-আসলে দ্বিগুণ হবে? (প্রথমে
সুদ বের করে সময় নির্ণয়)

\begin{itemize}
\item
  লিখিত নিয়ম: আসল ১০০ টাকা হলে সুদ-আসল ২০০ টাকা, অর্থাৎ সুদ ১০০ টাকা। এখন
  ৫\% হার সুদে ১০০ টাকা সুদ হতে সময় লাগে (100/5)=20 বছর.
\item
  শর্ট নিয়ম: সময় = (100(2-1)/5)=20 বছর.
\end{itemize}

\textbf{উদাহরণ ২: বার্ষিক ৮\% সরল সুদে কত সময়ে সুদে-আসলে ৩ গুণ হবে?}

\begin{itemize}
\item
  লিখিত নিয়ম: আসল ১০০ টাকা হলে সুদ হবে ২০০ টাকা। ৮\% সুদে ২০০ টাকা হতে সময়
  লাগবে (200/8)=25 বছর.
\item
  শর্ট নিয়ম: সময় = (100(3-1)/8)=(200/8)=25 বছর.
\end{itemize}

\textbf{উদাহরণ ৩: বার্ষিক ২৫\% সরল সুদে কত সময়ে সুদে-আসলে ৫ গুণ হবে?}

\begin{itemize}
\item
  লিখিত নিয়ম: আসল ১০০ টাকা হলে সুদ হবে ৪০০ টাকা। ২৫\% সুদে ৪০০ টাকা হতে
  সময় লাগবে (400/25)=16 বছর.
\item
  শর্ট নিয়ম: সময় = (100(5-1)/25)=(400/25)=16 বছর.
\end{itemize}

সুদকষার অংকগুলো দেখতে জটিল মনে হলেও যদি মূল বেসিক বা ১০০\%-এর ধারণাকে কাজে
লাগানো যায়, তবে সূত্র মুখস্থ করার ঝামেলা ছাড়াই খুব দ্রুত সমাধান করা সম্ভব।
পাটিগণিতের এই কৌশলগুলো আয়ত্ত করতে পারলে যেকোনো প্রতিযোগিতামূলক পরীক্ষায় সময়ের
সাশ্রয় হবে এবং নির্ভুল উত্তর নিশ্চিত করা যাবে।

নিয়মিত অনুশীলনের মাধ্যমেই কেবল এই পদ্ধতিগুলোতে দক্ষতা অর্জন করা সম্ভব.

\section{49. সরল ও চক্রবৃদ্ধি মুনাফা / সুদকষা MCQ প্রশ্ন উত্তর
অনুশীলন}\label{ux9b8ux9b0ux9b2-ux993-ux99aux995ux9b0ux9acux9a6ux9a7-ux9aeux9a8ux9ab-ux9b8ux9a6ux995ux9b7-mcq-ux9aaux9b0ux9b6ux9a8-ux989ux9a4ux9a4ux9b0-ux985ux9a8ux9b6ux9b2ux9a8}

\textbf{১০০ টাকায় ১ বছরের সুদ ৫ টাকা। ২০০ টাকার ২ বছরের সুদ কত?}

\textbf{\hl{ক) ২০ টাকা ✔}}

খ) ৩০ টাকা

গ) ৪০ টাকা

ঘ) ৫০ টাকা

\textbf{\hl{✔ সঠিক উত্তর: ক) ২০ টাকা}}

\begin{quote}
\textbf{ব্যাখ্যা:} \emph{সমাধান-০১ :}

\emph{১০০ টাকায় ১ বছরের সুদ = ৫ টাকা।}

\emph{১ টাকায় ১ বছরের সুদ = (৫/১০০) টাকা।}

\emph{২০০ টাকায় ২ বছরের সুদ = (৫×২০০×২/১০০) টাকা = ২০ টাকা}

\emph{শর্ট টেকনিক:}

\emph{এখানে , I = সরল মুনাফা/সুদ = ?}

\emph{P = আসল = ২০০ টাকা}

\emph{n = সময় = ২ বছর}

\emph{r = মুনাফার হার = ৫\% = (৫/১০০)}

\emph{আমরা জানি , I = Prn}

\emph{বা, I = (৫×২০০×২/১০০) = ২০ টাকা}
\end{quote}

\textbf{বার্ষিক ১০\% মুনাফায় ১২০০ টাকার ৫ বছরের মুনাফা কত?}

ক) ৫২০ টাকা

খ) ৪৮০ টাকা

গ) ৫৬০ টাকা

\textbf{\hl{ঘ) ৬০০ টাকা ✔}}

\textbf{\hl{✔ সঠিক উত্তর: ঘ) ৬০০ টাকা}}

\begin{quote}
\textbf{ব্যাখ্যা:} \emph{১০\% হার মুনাফায় ১,২০০ টাকার ৫ বছরের}

\emph{মুনাফা = (১০×১২০০×৫/১০০) = ৬০০ টাকা।}

\emph{শর্ট টেকনিক:}

\emph{এখানে, আসল P = ১২০০ টাকা}

\emph{মুনাফার হার r = ১০\% = (১০/১০০)}

\emph{সময় n = ৫ বছর}

\emph{আমরা জানি, মুনাফা I = Prn}

\emph{I = ১২০০ × (১০/১০০) × ৫ = ৬০০ টাকা।}
\end{quote}

\textbf{১০\% সরল মুনাফায় ২,০০০ টাকার ১ বছরের মুনাফা কত টাকা হবে?}

\textbf{\hl{ক) ২০০ টাকা ✔}}

খ) ৩০০ টাকা

গ) ৪০০ টাকা

ঘ) ১০০ টাকা

\textbf{\hl{✔ সঠিক উত্তর: ক) ২০০ টাকা}}

\begin{quote}
\textbf{ব্যাখ্যা:} \emph{সুদের হার r = ১০\%, আসল p = ২,০০০ টাকা এবং সময় n =
১ বছর।}

\emph{সরল মুনাফার ক্ষেত্রে, I = prn = ২,০০০ × ১০\% × ১ = (২০০০×১০×১/১০০) =
২০০ টাকা।}

\emph{ঐকিক নিয়ম:}

\emph{১০০ টাকার ১ বছরের মুনাফা = ১০ টাকা}

\emph{∴ ১ টাকার ১ বছরের মুনাফা = (১০/১০০) টাকা}

\emph{∴ ২০০০ টাকার ১ বছরের মুনাফা = (১০×২০০০/১০০) = ২০০ টাকা।}
\end{quote}

\textbf{বার্ষিক শতকরা ১০ টাকা হার সরলসুদে ৪৫০০ টাকার ৩ (তিন) বছরের সুদ কত?}

ক) ১৫০০ টাকা

খ) ১৪৫০ টাকা

গ) ১৬২০ টাকা

\textbf{\hl{ঘ) ১৩৫০ টাকা ✔}}

\textbf{\hl{✔ সঠিক উত্তর: ঘ) ১৩৫০ টাকা}}

\begin{quote}
\textbf{ব্যাখ্যা:} \emph{সরল মুনাফার ক্ষেত্রে,}

\emph{I = pnr}

\emph{= ৪৫০০ × ৩ × (১/১০)}

\emph{∴ I = ১৩৫০ টাকা।}

\emph{এখানে,}

\emph{আসল p = ৪৫০০}

\emph{r = ১০\% = (১০/১০০) = (১/১০)}

\emph{সময় n = ৩ বছর।}

\emph{ঐকিক নিয়ম:}

\emph{১০০ টাকার ১ বছরের সুদ = ১০ টাকা}

\emph{∴ ১ টাকার ১ বছরের সুদ = (১০/১০০) টাকা}

\emph{∴ ৪৫০০ টাকার ৩ বছরের সুদ = (১০×৪৫০০×৩/১০০) = ১৩৫০ টাকা।}
\end{quote}

\textbf{জাওয়াদ সাহেব ৫০০ টাকা ৩ বছরের জন্য ব্যাংকে জমা রাখলেন। ৮ (১/৩) \%
হারে তিনি কত টাকা সরল মুনাফা পাবেন?}

ক) ১০০ টাকা

\textbf{\hl{খ) ১২৫ টাকা ✔}}

গ) ৫২৫ টাকা

ঘ) ৬২৫ টাকা

\textbf{\hl{✔ সঠিক উত্তর: খ) ১২৫ টাকা}}

\begin{quote}
\textbf{ব্যাখ্যা:} \emph{সরল মুনাফার ক্ষেত্রে, I = pnr}

\emph{এখানে,}

\emph{আসল p = ৫০০ টাকা}

\emph{সময় n = ৩ বছর}

\emph{মুনাফার হার r = ৮ (১/৩) \% = (২৫/৩) \% = (২৫/৩×১০০) = (১/১২)}

\emph{∴ I = ৫০০ × ৩ × (২৫/৩×১০০) = ১২৫ টাকা।}

\emph{ঐকিক নিয়ম:}

\emph{১০০ টাকার ১ বছরের মুনাফা = (২৫/৩) টাকা}

\emph{∴ ১ টাকার ১ বছরের মুনাফা = (২৫/৩×১০০) টাকা}

\emph{∴ ৫০০ টাকার ৩ বছরের মুনাফা = (২৫×৫০০×৩/৩×১০০) = ১২৫ টাকা।}
\end{quote}

\textbf{৯.৫\% হার সরল সুদে ৬০০ টাকার ২ বছরের সুদ কত?}

\textbf{\hl{ক) ১১৪ টাকা ✔}}

খ) ১০৮ টাকা

গ) ৫৭ টাকা

ঘ) ৫৪ টাকা

\textbf{\hl{✔ সঠিক উত্তর: ক) ১১৪ টাকা}}

\begin{quote}
\textbf{ব্যাখ্যা:} \emph{সমাধান-০১ :}

\emph{১০০ টাকায় ১ বছরের সুদ = ৯.৫ টাকা।}

\emph{১ টাকায় ১ বছরের সুদ = (৯.৫/১০০) টাকা।}

\emph{৬০০ টাকায় ২ বছরের সুদ = (৯.৫×৬০০×২/১০০) টাকা = ১১৪ টাকা}

\emph{শর্ট টেকনিক:}

\emph{এখানে , I = সরল মুনাফা/সুদ = ?}

\emph{P = আসল = ৬০০ টাকা}

\emph{n = সময় = ২ বছর}

\emph{r = মুনাফার হার = ৯.৫\% = (৯.৫/১০০)}

\emph{আমরা জানি , I = Prn}

\emph{বা, I = (৯.৫×৬০০×২/১০০) = ১১৪ টাকা}
\end{quote}

\textbf{এক ব্যক্তি ২০\% সরল সুদে ৮০০ টাকা এবং ১৫\% সরল সুদে ৬০০ টাকা বিনিয়োগ
করলে এক বছর পর তিনি কত সুদ পাবেন?}

ক) ২০৫ টাকা

\textbf{\hl{খ) ২৫০ টাকা ✔}}

গ) ২২৫ টাকা

ঘ) ২৯০ টাকা

\textbf{\hl{✔ সঠিক উত্তর: খ) ২৫০ টাকা}}

\begin{quote}
\textbf{ব্যাখ্যা:} \emph{সমাধান-০১ :}

\emph{২০\% সুদে ৮০০ টাকার ১ বছরের সুদ = (৮০০×২০/১০০) = ১৬০ টাকা}

\emph{১৫\% সুদে ৬০০ টাকার ১ বছরের সুদ = (৬০০×১৫/১০০) = ৯০ টাকা}

\emph{∴ মোট সুদ = ১৬০ + ৯০ = ২৫০ টাকা}

\emph{শর্ট টেকনিক:}

\emph{মোট সুদ I = I\_১ + I\_২}

\emph{I = (P\_১r\_১n) + (P\_২r\_২n)}

\emph{I = (৮০০ × ০.২০ × ১) + (৬০০ × ০.১৫ × ১)}

\emph{I = ১৬০ + ৯০ = ২৫০ টাকা}
\end{quote}

\textbf{আরিফ বার্ষিক ৬.৫\% হার সুদে ৯০০০ টাকা ব্যাংকে জমা রাখল। সুদের হার
বেড়ে ১২\% হলে আরিফ ১ বছরে সুদ হিসেবে কত টাকা বেশি পাবে?}

ক) ৫৮৫ টাকা

\textbf{\hl{খ) ৪৯৫ টাকা ✔}}

গ) ১৮০০ টাকা

ঘ) ৫০০ টাকা

\textbf{\hl{✔ সঠিক উত্তর: খ) ৪৯৫ টাকা}}

\begin{quote}
\textbf{ব্যাখ্যা:} \emph{সমাধান-০১ :}

\emph{৬.৫\% সুদে ৯০০০ টাকার ১ বছরের সুদ = (৯০০০×৬.৫/১০০) = ৫৮৫ টাকা}

\emph{আবার, ১২\% সুদে ৯০০০ টাকার ১ বছরের সুদ = (৯০০০×১২/১০০) = ১০৮০ টাকা}

\emph{∴ সুদ বেশি পাবে = ১০৮০ − ৫৮৫ = ৪৯৫ টাকা}

\emph{শর্ট টেকনিক:}

\emph{সুদের হার বৃদ্ধি = ১২\% − ৬.৫\% = ৫.৫\%}

\emph{∴ ১ বছরের অতিরিক্ত সুদ = Prn = ৯০০০ × (৫.৫/১০০) × ১ = ৪৯৫ টাকা}
\end{quote}

\textbf{বার্ষিক ৫\% হারে ৭৫০ টাকার ৪ বছরের মুনাফা কত?}

ক) ২০০ টাকা

\textbf{\hl{খ) ১৫০ টাকা ✔}}

গ) ২৫০ টাকা

ঘ) ৩০০ টাকা

\textbf{\hl{✔ সঠিক উত্তর: খ) ১৫০ টাকা}}

\begin{quote}
\textbf{ব্যাখ্যা:} \emph{সমাধান-০১ :}

\emph{১০০ টাকায় ১ বছরের সুদ = ৫ টাকা।}

\emph{১ টাকায় ১ বছরের সুদ = (৫/১০০) টাকা।}

\emph{৭৫০ টাকায় ৪ বছরের সুদ = (৫×৭৫০×৪/১০০) টাকা = ১৫০ টাকা}

\emph{শর্ট টেকনিক:}

\emph{এখানে , I = সরল মুনাফা/সুদ = ?}

\emph{P = আসল = ৭৫০ টাকা}

\emph{n = সময় = ৪ বছর}

\emph{r = মুনাফার হার = ৫\% = (৫/১০০)}

\emph{আমরা জানি , I = Prn}

\emph{বা, I = (৫×৭৫০×৪/১০০) = ১৫০ টাকা}
\end{quote}

\textbf{শতকরা বার্ষিক ১৫\% সুদে ৮০০০ টাকার ৬ মাসের সুদ কত?}

ক) ৫০০ টাকা

\textbf{\hl{খ) ৬০০ টাকা ✔}}

গ) ৭০০ টাকা

ঘ) ৮০০ টাকা

\textbf{\hl{✔ সঠিক উত্তর: খ) ৬০০ টাকা}}

\begin{quote}
\textbf{ব্যাখ্যা:} \emph{সমাধান-০১ :}

\emph{এখানে ৬ মাসকে প্রথমে বছরে রূপ দিতে হবে।}

\emph{আমরা জানি,}

\emph{১২ মাস = ১ বছর}

\emph{১ মাস = (১/১২) বছর}

\emph{৬ মাস = (১×৬/১২) = ০.৫ বছর}

\emph{১০০ টাকার ১ বছরের সুদ = ১৫ টাকা}

\emph{৮০০০ টাকার ০.৫ বছরের সুদ = (১৫×৮০০০×০.৫/১০০) = ৬০০ টাকা}

\emph{শর্ট টেকনিক:}

\emph{এখানে , I = সরল মুনাফা/সুদ = ?}

\emph{P = আসল = ৮০০০ টাকা}

\emph{n = সময় = ০.৫ বছর}

\emph{r = মুনাফার হার = ১৫\% = (১৫/১০০)}

\emph{আমরা জানি , I = Prn}

\emph{বা, I = (১৫×৮০০০×০.৫/১০০) = ৬০০ টাকা}
\end{quote}

\textbf{শতকরা বার্ষিক ৫ টাকা হার সরল সুদে ৬৪০ টাকার ২ বছর ৬ মাসের সুদ কত?}

ক) ৬০ টাকা

খ) ১৬০ টাকা

গ) ৬৪ টাকা

\textbf{\hl{ঘ) ৮০ টাকা ✔}}

\textbf{\hl{✔ সঠিক উত্তর: ঘ) ৮০ টাকা}}

\begin{quote}
\textbf{ব্যাখ্যা:} \emph{সমাধান-০১ :}

\emph{এখানে ৬ মাসকে প্রথমে বছরে রূপ দিতে হবে।}

\emph{আমরা জানি, ১২ মাস = ১ বছর ∴ ৬ মাস = (৬/১২) = ০.৫ বছর}

\emph{∴ ২ বছর ৬ মাস = ২.৫ বছর}

\emph{১০০ টাকার ১ বছরের সুদ = ৫ টাকা}

\emph{৬৪০ টাকার ২.৫ বছরের সুদ = (৫×৬৪০×২.৫/১০০) = ৮০ টাকা}

\emph{শর্ট টেকনিক:}

\emph{এখানে , I = সরল মুনাফা/সুদ = ?}

\emph{P = আসল = ৬৪০ টাকা}

\emph{n = সময় = ২.৫ বছর}

\emph{r = মুনাফার হার = ৫\% = (৫/১০০)}

\emph{আমরা জানি , I = Prn}

\emph{বা, I = (৫×৬৪০×২.৫/১০০) = ৮০ টাকা}
\end{quote}

\textbf{শতকরা বার্ষিক ৫ টাকা হার সুদে ৭২০ টাকায় ২ বছর ৪ মাসের সুদ কত হবে?}

ক) ৮৩ টাকা

\textbf{\hl{খ) ৮৪ টাকা ✔}}

গ) ৮২ টাকা

ঘ) ৯০ টাকা

\textbf{\hl{✔ সঠিক উত্তর: খ) ৮৪ টাকা}}

\begin{quote}
\textbf{ব্যাখ্যা:} \emph{এখানে, ২ বছর ৪ মাস = (২ × ১২ + ৪) মাস = ২৮ মাস}

\emph{৫\% হার সুদে, ১০০ টাকায় ১২ মাসের সুদ = ৫ টাকা}

\emph{∴ ১০০ টাকায় ১ মাসের সুদ = (৫/১২) টাকা}

\emph{∴ ১ টাকায় ১ মাসের সুদ = (৫/১২×১০০) টাকা}

\emph{∴ ৭২০ টাকায় ২৮ মাসের সুদ = (৫×৭২০×২৮/১২×১০০) = ৮৪ টাকা}

\emph{শর্ট টেকনিক:}

\emph{এখানে,}

\emph{আসল P = ৭২০ টাকা}

\emph{সুদের হার r = ৫\% = (৫/১০০)}

\emph{সময় n = ২ বছর ৪ মাস = ২+(৪/১২) বছর = ২+(১/৩) বছর = (৭/৩) বছর}

\emph{আমরা জানি, I = Prn}

\emph{বা, I = ৭২০ × (৫/১০০) × (৭/৩) = ৮৪ টাকা}
\end{quote}

\textbf{কবির সাহেব ১০\% সরল সুদে ৭০০ টাকা এবং ৫\% সরল সুদে ১৩০০ টাকা
বিনিয়োগ করলে এক বছর তিনি মোট কত টাকা সুদ পাবেন?}

\textbf{\hl{ক) ১৩৫ টাকা ✔}}

খ) ১৫০ টাকা

গ) ২২৫ টাকা

ঘ) ৯০ টাকা

\textbf{\hl{✔ সঠিক উত্তর: ক) ১৩৫ টাকা}}

\begin{quote}
\textbf{ব্যাখ্যা:} \emph{সমাধান-০১ :}

\emph{৫\% সরল সুদে ১৩০০ টাকার ১ বছরের সুদ = (১৩০০×৫/১০০) = ৬৫ টাকা}

\emph{আবার, ১০\% সুদে ৭০০ টাকার ১ বছরের সুদ = (৭০০×১০/১০০) = ৭০ টাকা}

\emph{∴ মোট সুদ = ৭০ + ৬৫ = ১৩৫ টাকা}

\emph{শর্ট টেকনিক:}

\emph{I = (P\_১r\_১n) + (P\_২r\_২n)}

\emph{I = (৭০০ × ০.১০ × ১) + (১৩০০ × ০.০৫ × ১)}

\emph{I = ৭০ + ৬৫ = ১৩৫ টাকা}
\end{quote}

\textbf{প্রতি বছর শতকরা ৮ টাকা হারে লাভের চুক্তিতে ১০০০০ টাকা বিনিয়োগ করে ২
বছর পর ঐ বিনিয়োগকারী মোট কত টাকা লাভ পাবে?}

\textbf{\hl{ক) ১৬০ টাকা ✔}}

খ) ১৬৫ টাকা

গ) ১৬৬.৪ টাকা

ঘ) ১৭০ টাকা

\textbf{\hl{✔ সঠিক উত্তর: ক) ১৬০ টাকা}}

\begin{quote}
\textbf{ব্যাখ্যা:} \emph{সমাধান-০১ :}

\emph{১০০ টাকার ১ বছরের সুদ = ৮ টাকা।}

\emph{১ টাকার ১ বছরের সুদ = (৮/১০০) টাকা।}

\emph{১০০০ টাকার ২ বছরের সুদ = (৮×১০০০×২/১০০) টাকা = ১৬০ টাকা}

\emph{শর্ট টেকনিক:}

\emph{এখানে , I = সরল মুনাফা/সুদ = ?}

\emph{P = আসল = ১০০০ টাকা}

\emph{n = সময় = ২ বছর}

\emph{r = মুনাফার হার = ৮\% = (৮/১০০)}

\emph{আমরা জানি , I = Prn}

\emph{বা, I = (৮×১০০০×২/১০০) = ১৬০ টাকা}
\end{quote}

\textbf{২\% হার সুদে ১০০ টাকার ৩ বছরের সুদ অপেক্ষা ৩\% সুদে ঐ টাকার ৩ বছরের
সুদ কত বেশি হবে?}

ক) ১ টাকা

খ) ২ টাকা

\textbf{\hl{গ) ৩ টাকা ✔}}

ঘ) ৪ টাকা

\textbf{\hl{✔ সঠিক উত্তর: গ) ৩ টাকা}}

\begin{quote}
\textbf{ব্যাখ্যা:} \emph{সমাধান-০১ :}

\emph{২\% সরল সুদে ১০০ টাকার ৩ বছরের সুদ = (১০০×২×৩/১০০) = ৬ টাকা}

\emph{আবার, ৩\% সুদে ১০০ টাকার ৩ বছরের সুদ = (১০০×৩×৩/১০০) = ৯ টাকা}

\emph{∴ সুদ বেশি হবে = ৯ − ৬ = ৩ টাকা}

\emph{শর্ট টেকনিক:}

\emph{সুদের হারের পার্থক্য = ৩\% − ২\% = ১\%}

\emph{∴ সুদের পার্থক্য = P × (r\_২ − r\_১) × n = ১০০ × (১/১০০) × ৩ = ৩
টাকা}
\end{quote}

\textbf{সুদের হার শতকরা ৭ টাকা হলে ৬৫০ টাকায় ছয় বছরের সুদ কত হবে?}

ক) ২৭০ টাকা

\textbf{\hl{খ) ২৭৩ টাকা ✔}}

গ) ২৭২ টাকা

ঘ) ২৭৫ টাকা

\textbf{\hl{✔ সঠিক উত্তর: খ) ২৭৩ টাকা}}

\begin{quote}
\textbf{ব্যাখ্যা:} \emph{সমাধান-০১ :}

\emph{১০০ টাকার ১ বছরের সুদ = ৭ টাকা।}

\emph{১ টাকার ১ বছরের সুদ = (৭/১০০) টাকা।}

\emph{৬৫০ টাকার ৬ বছরের সুদ = (৭×৬৫০×৬/১০০) টাকা = ২৭৩ টাকা}

\emph{শর্ট টেকনিক:}

\emph{এখানে , I = সরল মুনাফা/সুদ = ?}

\emph{P = আসল = ৬৫০ টাকা}

\emph{n = সময় = ৬ বছর}

\emph{r = মুনাফার হার = ৭\% = (৭/১০০)}

\emph{আমরা জানি , I = Prn}

\emph{বা, I = (৭×৬৫০×৬/১০০) = ২৭৩ টাকা}
\end{quote}

\textbf{৬\% হারে নয় মাসে ১০০০০ টাকার উপর সুদ কত হবে?}

ক) ৫০০ টাকা

খ) ৬০০ টাকা

\textbf{\hl{গ) ৪৫০ টাকা ✔}}

ঘ) ৬৫০ টাকা

\textbf{\hl{✔ সঠিক উত্তর: গ) ৪৫০ টাকা}}

\begin{quote}
\textbf{ব্যাখ্যা:} \emph{সমাধান-০১ :}

\emph{এখানে ৯ মাসকে প্রথমে বছরে রূপ দিতে হবে।}

\emph{আমরা জানি, ১২ মাস = ১ বছর ∴ ৯ মাস = (৯/১২) = ০.৭৫ বছর}

\emph{১০০ টাকার ১ বছরের সুদ = ৬ টাকা}

\emph{১০০০০ টাকার ০.৭৫ বছরের সুদ = (৬×১০০০০×০.৭৫/১০০) = ৪৫০ টাকা}

\emph{শর্ট টেকনিক:}

\emph{এখানে , I = সরল মুনাফা/সুদ = ?}

\emph{P = আসল = ১০০০০ টাকা}

\emph{n = সময় = ০.৭৫ বছর}

\emph{r = মুনাফার হার = ৬\% = (৬/১০০)}

\emph{আমরা জানি , I = Prn}

\emph{বা, I = (৬×১০০০০×০.৭৫/১০০) = ৪৫০ টাকা}
\end{quote}

\textbf{শতকরা বার্ষিক ১২ টাকা হারে ৬০০ টাকার ৬ মাসের সুদ কত হবে?}

ক) ২৪ টাকা

খ) ২৬ টাকা

\textbf{\hl{গ) ৩৬ টাকা ✔}}

ঘ) ৪৮ টাকা

\textbf{\hl{✔ সঠিক উত্তর: গ) ৩৬ টাকা}}

\begin{quote}
\textbf{ব্যাখ্যা:} \emph{সমাধান-০১ :}

\emph{এখানে ৬ মাসকে প্রথমে বছরে রূপ দিতে হবে।}

\emph{আমরা জানি, ১২ মাস = ১ বছর ∴ ৬ মাস = (৬/১২) = ০.৫ বছর}

\emph{১০০ টাকার ১ বছরের সুদ = ১২ টাকা}

\emph{৬০০ টাকার ০.৫ বছরের সুদ = (১২×৬০০×০.৫/১০০) = ৩৬ টাকা}

\emph{শর্ট টেকনিক:}

\emph{এখানে , I = সরল মুনাফা/সুদ = ?}

\emph{P = আসল = ৬০০ টাকা}

\emph{n = সময় = ০.৫ বছর}

\emph{r = মুনাফার হার = ১২\% = (১২/১০০)}

\emph{আমরা জানি , I = Prn}

\emph{বা, I = (১২×৬০০×০.৫/১০০) = ৩৬ টাকা}
\end{quote}

\textbf{৫\% হারে ৩০০ টাকার ৫ বছরে সুদাসল কত হয়?}

ক) ৩২৫

খ) ৩৪৫

গ) ৩৬০

\textbf{\hl{ঘ) ৩৭৫ ✔}}

\textbf{\hl{✔ সঠিক উত্তর: ঘ) ৩৭৫}}

\begin{quote}
\textbf{ব্যাখ্যা:} \emph{সমাধান-০১ :}

\emph{১০০ টাকার ১ বছরের সুদ = ৫ টাকা।}

\emph{১ টাকার ১ বছরের সুদ = (৫/১০০) টাকা।}

\emph{৩০০ টাকার ৫ বছরের সুদ = (৫×৩০০×৫/১০০) টাকা = ৭৫ টাকা}

\emph{∴ ৩০০ টাকার ৫ বছরে সুদাসল = ৩০০ + ৭৫ = ৩৭৫ টাকা}

\emph{শর্ট টেকনিক:}

\emph{এখানে , I = সরল মুনাফা/সুদ = ?}

\emph{P = আসল = ৩০০ টাকা}

\emph{n = সময় = ৫ বছর}

\emph{r = মুনাফার হার = ৫\% = (৫/১০০)}

\emph{আমরা জানি , I = Prn}

\emph{বা, I = (৫×৩০০×৫/১০০) = ৭৫ টাকা}

\emph{∴ ৩০০ টাকার ৫ বছরে সুদাসল = ৩০০ + ৭৫ = ৩৭৫ টাকা}
\end{quote}

\textbf{সুদের হার ৬.৫\% হলে ১০০০ টাকা জমা দিয়ে ৬ বছর পরে সুদে আসলে কত টাকা
পাওয়া যাবে?}

ক) ১২৫০ টাকা

খ) ৩৯০ টাকা

\textbf{\hl{গ) ১৩৯০ টাকা ✔}}

ঘ) ১০৬৫ টাকা

\textbf{\hl{✔ সঠিক উত্তর: গ) ১৩৯০ টাকা}}

\begin{quote}
\textbf{ব্যাখ্যা:} \emph{সমাধান-০১ :}

\emph{১০০ টাকার ১ বছরের সুদ = ৬.৫ টাকা।}

\emph{১ টাকার ১ বছরের সুদ = (৬.৫/১০০) টাকা।}

\emph{১০০০ টাকার ৬ বছরের সুদ = (৬.৫×১০০০×৬/১০০) টাকা = ৩৯০ টাকা}

\emph{∴ ১০০০ টাকা জমা দিয়ে ৬ বছর পরে সুদে আসলে পাবে = ১০০০ + ৩৯০ = ১৩৯০
টাকা}

\emph{শর্ট টেকনিক:}

\emph{এখানে , I = সরল মুনাফা/সুদ = ?}

\emph{P = আসল = ১০০০ টাকা}

\emph{n = সময় = ৬ বছর}

\emph{r = মুনাফার হার = ৬.৫\% = (৬.৫/১০০)}

\emph{আমরা জানি , I = Prn}

\emph{বা, I = (৬.৫×১০০০×৬/১০০) = ৩৯০ টাকা}

\emph{∴ ১০০০ টাকা জমা দিয়ে ৬ বছর পরে সুদে আসলে পাবে = ১০০০ + ৩৯০ = ১৩৯০
টাকা}
\end{quote}

\textbf{মুনাফার হার ৮\% হলে ৫০০ টাকার ৪ বছরের মুনাফা কত?}

ক) ৪০ টাকা

খ) ৩২ টাকা

\textbf{\hl{গ) ১৬০ টাকা ✔}}

ঘ) ২০ টাকা

\textbf{\hl{✔ সঠিক উত্তর: গ) ১৬০ টাকা}}

\begin{quote}
\textbf{ব্যাখ্যা:} \emph{সমাধান-০১ :}

\emph{১০০ টাকার ১ বছরের সুদ = ৮ টাকা।}

\emph{১ টাকার ১ বছরের সুদ = (৮/১০০) টাকা।}

\emph{৫০০ টাকার ৪ বছরের সুদ = (৮×৫০০×৪/১০০) টাকা = ১৬০ টাকা}

\emph{শর্ট টেকনিক:}

\emph{এখানে , I = সরল মুনাফা/সুদ = ?}

\emph{P = আসল = ৫০০ টাকা}

\emph{n = সময় = ৪ বছর}

\emph{r = মুনাফার হার = ৮\% = (৮/১০০)}

\emph{আমরা জানি , I = Prn}

\emph{বা, I = (৮×৫০০×৪/১০০) = ১৬০ টাকা}
\end{quote}

\textbf{শতকরা ৫ টাকা হার সুদে ১২০ টাকায় ৩ বছরের সুদেআসলে কত হয়?}

ক) ১৩৫

খ) ১৩৭.৫

\textbf{\hl{গ) ১৩৮ ✔}}

ঘ) ১৪৮

\textbf{\hl{✔ সঠিক উত্তর: গ) ১৩৮}}

\begin{quote}
\textbf{ব্যাখ্যা:} \emph{সমাধান-০১ :}

\emph{১০০ টাকার ১ বছরের সুদ = ৫ টাকা।}

\emph{১ টাকার ১ বছরের সুদ = (৫/১০০) টাকা।}

\emph{১২০ টাকার ৩ বছরের সুদ = (৫×১২০×৩/১০০) টাকা = ১৮ টাকা}

\emph{∴ ১২০ টাকার ৩ বছরের সুদাসল = ১২০ + ১৮ = ১৩৮ টাকা}

\emph{শর্ট টেকনিক:}

\emph{এখানে , I = সরল মুনাফা/সুদ = ?}

\emph{P = আসল = ১২০ টাকা}

\emph{n = সময় = ৩ বছর}

\emph{r = মুনাফার হার = ৫\% = (৫/১০০)}

\emph{আমরা জানি , I = Prn}

\emph{বা, I = (৫×১২০×৩/১০০) = ১৮ টাকা}

\emph{∴ ১২০ টাকার ৩ বছরের সুদাসল = ১২০ + ১৮ = ১৩৮ টাকা}
\end{quote}

\textbf{সুদের হার ৬\% থেকে কমে ৪\% হলে ৫০০ টাকায় ৩ বছরের কত সুদ কমবে?}

ক) ২০ টাকা

\textbf{\hl{খ) ৩০ টাকা ✔}}

গ) ৪০ টাকা

\textbf{\hl{✔ সঠিক উত্তর: খ) ৩০ টাকা}}

\begin{quote}
\textbf{ব্যাখ্যা:} \emph{সমাধান-০১ :}

\emph{১০০ টাকায় ১ বছরের সুদের হার কমে = ৬ − ৪ = ২ টাকা}

\emph{১০০ টাকায় ৩ বছরে কমে = ৩ × ২ = ৬ টাকা}

\emph{১০০ টাকায় কমে = ৬ টাকা}

\emph{১ টাকায় কমে = (৬/১০০) টাকা}

\emph{৫০০ টাকায় কমে = (৬×৫০০×৩/১০০) = ৩০ টাকা}

\emph{শর্ট টেকনিক:}

\emph{সুদের হারের হ্রাস (r) = ৬\% − ৪\% = ২\% = (২/১০০)}

\emph{∴ মোট সুদের হ্রাস I = Prn}

\emph{বা, I = ৫০০ × (২/১০০) × ৩ = ৩০ টাকা}
\end{quote}

\textbf{সুদের হারের সূত্র কোনটি?}

ক) (১০০×সুদ/হার×আসল)

\textbf{\hl{খ) (১০০×সুদ/সময়×আসল) ✔}}

গ) (সময়×আসল/১০০×সুদ)

ঘ) (১০০×আসল/সুদ×সময়)

\textbf{\hl{✔ সঠিক উত্তর: খ) (১০০×সুদ/সময়×আসল)}}

\begin{quote}
\textbf{ব্যাখ্যা:} \emph{বীজগণিতীয় সূত্রের সাহায্যে:}

\emph{আমরা জানি, সরল সুদ I = (P×r×n/১০০)}

\emph{বা, I × ১০০ = P × r × n}

\emph{বা, r = (১০০×I/n×P)}

\emph{∴ সুদের হার (r) = (১০০×সুদ/সময়×আসল)}

\emph{কথায় রূপান্তর করে (বাংলা সংজ্ঞায়):}

\emph{সুদ = (আসল×সময়×সুদের হার/১০০)}

\emph{বা, আসল × সময় × সুদের হার = ১০০ × সুদ}

\emph{∴ সুদের হার = (১০০×সুদ/সময়×আসল)}
\end{quote}

\textbf{x টাকার x\% সরল সুদে ৪ বছরের সুদ x টাকা হলে x = ?}

ক) ২০

\textbf{\hl{খ) ২৫ ✔}}

গ) ৫০

ঘ) ৫৫

\textbf{\hl{✔ সঠিক উত্তর: খ) ২৫}}

\begin{quote}
\textbf{ব্যাখ্যা:} \emph{ঐকিক নিয়ম:}

\emph{x\% সরল সুদে x টাকার ৪ বছরের সুদ = x টাকা}

\emph{∴ x\% সরল সুদে ১ টাকার ১ বছরের সুদ = (x/৪x) টাকা}

\emph{∴ x\% সরল সুদে ১০০ টাকার ১ বছরের সুদ = (১০০x/৪x) = ২৫ টাকা}

\emph{∴ x = ২৫}

\emph{শর্ট টেকনিক:}

\emph{আমরা জানি, I = Prn}

\emph{বা, x = x × (x/১০০) × ৪}

\emph{বা, x = (৪x\^{}২/১০০)}

\emph{বা, ৪x\^{}২ = ১০০x}

\emph{বা, ৪x = ১০০}

\emph{∴ x = ২৫}
\end{quote}

\textbf{সুলেমান সাহেব ১৫,০০০ টাকা ব্যাংকে জমা করে বছরে সুদ বাবদ ১,২৭৫ টাকা
আয় করে। সুদের হার কত?}

ক) ৮\%

খ) ৮.২৫\%

\textbf{\hl{গ) ৮.৫\% ✔}}

ঘ) ৮.৬\%

\textbf{\hl{✔ সঠিক উত্তর: গ) ৮.৫\%}}

\begin{quote}
\textbf{ব্যাখ্যা:} \emph{ঐকিক নিয়ম:}

\emph{১৫,০০০ টাকায় ১ বছরে সুদ পায় = ১,২৭৫ টাকা}

\emph{∴ ১ টাকায় ১ বছরে সুদ পায় = (১২৭৫/১৫০০০) টাকা}

\emph{∴ ১০০ টাকায় ১ বছরে সুদ পায় = (১২৭৫×১০০/১৫০০০) = ৮.৫ টাকা}

\emph{∴ সুদের হার = ৮.৫\%}

\emph{শর্ট টেকনিক:}

\emph{সুদের হার = (১০০×সুদ/সময়×আসল)}

\emph{= (১০০×১২৭৫/১×১৫০০০)}

\emph{= ৮.৫\%}
\end{quote}

\textbf{বার্ষিক শতকরা কত হার সুদে ২০০০ টাকা ৩ বছরের সুদে-আসলে ২৩০০ টাকা
হয়?}

ক) ১৫\%

খ) ১০\%

গ) ৭.৫\%

\textbf{\hl{ঘ) ৫\% ✔}}

\textbf{\hl{✔ সঠিক উত্তর: ঘ) ৫\%}}

\begin{quote}
\textbf{ব্যাখ্যা:} \emph{সুদ = (২৩০০ − ২০০০) টাকা = ৩০০ টাকা}

\emph{ঐকিক নিয়ম:}

\emph{২০০০ টাকায় ৩ বছরের সুদ = ৩০০ টাকা}

\emph{∴ ১ টাকায় ১ বছরের সুদ = (৩০০/২০০০×৩) টাকা}

\emph{∴ ১০০ টাকায় ১ বছরের সুদ = (৩০০×১০০/২০০০×৩) = ৫ টাকা}

\emph{∴ সুদের হার = ৫\%}

\emph{শর্ট টেকনিক:}

\emph{সুদ I = ২৩০০ − ২০০০ = ৩০০ টাকা}

\emph{এখানে, আসল P = ২০০০ টাকা, সময় n = ৩ বছর}

\emph{আমরা জানি, সুদের হার r = (১০০×I/P×n)}

\emph{= (১০০×৩০০/২০০০×৩) \% = ৫\%}
\end{quote}

\textbf{কোন আসল ৩ বছরের মুনাফা-আসলে ৫৫৫০০ টাকা হয়। মুনাফা আসলের (৩/৮) অংশ
হলে মুনাফার হার কত?}

ক) ১০\%

\textbf{\hl{খ) ১২.৫\% ✔}}

গ) ১৫\%

ঘ) ১২\%

\textbf{\hl{✔ সঠিক উত্তর: খ) ১২.৫\%}}

\begin{quote}
\textbf{ব্যাখ্যা:} \emph{ধরি, আসল = ১০০ টাকা}

\emph{৩ বছরের মুনাফা = ১০০ এর (৩/৮) অংশ = ১০০ × (৩/৮) টাকা}

\emph{১ বছরের মুনাফা = (১০০×(৩/৮)/৩) = ১০০ × (১/৮) = ১২.৫ টাকা}

\emph{∴ ১০০ টাকার ১ বছরের মুনাফা ১২.৫ টাকা। অর্থাৎ সুদের হার ১২.৫\%}

\emph{শর্ট টেকনিক:}

\emph{সুদের হার r = (১০০×কত অংশ/বছর)}

\emph{r = (১০০×(৩/৮)/৩) \% = (১০০/৮) \% = ১২.৫\%}
\end{quote}

\textbf{একই হার সুদে ৩০০ টাকার ৪ বছরের সুদ এবং ৫০০ টাকার ৫ বছরের সুদ একত্রে
২২২ টাকা হলে, শতকরা বার্ষিক সুদের হার কত?}

ক) ৪\%

\textbf{\hl{খ) ৬\% ✔}}

গ) ৫\%

ঘ) ৫.৫\%

\textbf{\hl{✔ সঠিক উত্তর: খ) ৬\%}}

\begin{quote}
\textbf{ব্যাখ্যা:} \emph{ঐকিক নিয়ম / বিস্তারিত সমাধান:}

\emph{৩০০ টাকার ৪ বছরের সুদ = (৩০০ × ৪) টাকা = ১২০০ টাকার ১ বছরের সুদ}

\emph{৫০০ টাকার ৫ বছরের সুদ = (৫০০ × ৫) টাকা = ২৫০০ টাকার ১ বছরের সুদ}

\emph{∴ মোট আসল = ১২০০ + ২৫০০ = ৩৭০০ টাকা}

\emph{৩৭০০ টাকার ১ বছরের সুদ = ২২২ টাকা}

\emph{∴ ১০০ টাকার ১ বছরের সুদ = (২২২×১০০/৩৭০০) = ৬ টাকা}

\emph{∴ সুদের হার = ৬\%}

\emph{শর্ট টেকনিক:}

\emph{সুদের হার r = (১০০×২২২/(৫০০×৫)+(৩০০×৪))}

\emph{= (২২২×১০০/২৫০০+১২০০)}

\emph{= (২২২০০/৩৭০০) = ৬\%}
\end{quote}

\textbf{মোট ১২০০ টাকা দুই সঞ্চয় অ্যাকাউন্ট, ৫\% হারে এক বছরের জন্য এবং বাকি
টাকা ৭\% সুদের হারে জমা করা হয়। যদি ৭২ টাকা সুদ পাওয়া যায়, তাহলে ৫\% হারে
কত টাকা জমা করা হয়েছিল?}

ক) Tk 300

খ) Tk 450

\textbf{\hl{গ) Tk 600 ✔}}

ঘ) Tk 760

\textbf{\hl{✔ সঠিক উত্তর: গ) Tk 600}}

\begin{quote}
\textbf{ব্যাখ্যা:} \emph{বীজগণিতীয় সমাধান:}

\emph{ধরি, ৫\% হারে x টাকা এবং ৭\% হারে (১২০০ − x) টাকা জমা করা হয়েছিল।}

\emph{প্রশ্নমতে,}

\emph{(১২০০ − x) × (৭/১০০) + x × (৫/১০০) = ৭২}

\emph{⇒ (১২০০ − x) × ৭ + ৫x = ৭২০০}

\emph{⇒ ৮৪০০ − ৭x + ৫x = ৭২০০}

\emph{⇒ −২x = ৭২০০ − ৮৪০০}

\emph{⇒ −২x = −১২০০}

\emph{∴ x = ৬০০}

\emph{∴ ৫\% হারে জমা করা হয়েছিল Tk 600}

\emph{বিকল্প পদ্ধতি (গড় সুদের হার ও মিশ্রণ নিয়ম):}

\emph{১২০০ টাকায় প্রাপ্ত মোট সুদ = ৭২ টাকা}

\emph{∴ গড় সুদের হার = (৭২×১০০/১২০০) \% = ৬\%}

\emph{এখন, দুটি হারের পার্থক্য অনুপাত:}

\emph{(৭\% − ৬\%) : (৬\% − ৫\%) = ১ : ১}

\emph{অর্থাৎ, দুটি ব্যাংকে সমান পরিমাণ টাকা জমা রাখা হয়েছিল।}

\emph{∴ ৫\% হারে জমার পরিমাণ = (১২০০/২) = ৬০০ টাকা বা Tk 600}
\end{quote}

\textbf{কোনো আসল টাকার ৪ বছরের সুদ, সুদাসলের (১/৬) অংশ। সুদের হার কত?}

ক) ৬(১/৪) \%

\textbf{\hl{খ) ৫\% ✔}}

গ) ৭(১/২) \%

ঘ) ৪\%

\textbf{\hl{✔ সঠিক উত্তর: খ) ৫\%}}

\begin{quote}
\textbf{ব্যাখ্যা:} \emph{বীজগণিতীয়/সূত্রভিত্তিক পদ্ধতি:}

\emph{মনে করি, সুদ-আসল = ৬ টাকা}

\emph{∴ সুদ = ৬ এর (১/৬) = ১ টাকা}

\emph{∴ আসল = (৬ − ১) = ৫ টাকা}

\emph{আমরা জানি, I = pnr}

\emph{⇒ ১ = ৫ × ৪ × r}

\emph{⇒ r = (১/২০)×১০০ \%}

\emph{∴ r = ৫\%}

\emph{ঐকিক নিয়ম:}

\emph{ধরি, সুদ-আসল = ৬ টাকা}

\emph{∴ সুদ = ১ টাকা এবং আসল = (৬ − ১) = ৫ টাকা}

\emph{৫ টাকার ৪ বছরের সুদ = ১ টাকা}

\emph{∴ ১ টাকার ১ বছরের সুদ = (১/৫×৪) টাকা}

\emph{∴ ১০০ টাকার ১ বছরের সুদ = (১×১০০/৫×৪) = ৫ টাকা}

\emph{∴ সুদের হার = ৫\%}
\end{quote}

\textbf{ব্যাংকে ৮৫০০ টাকা রেখে ৪ বছর পর সুদাসলে ১১২২০ টাকা পেলে বার্ষিক সরল
সুদের হার কত ছিল?}

ক) ১০\%

খ) ৯\%

\textbf{\hl{গ) ৮\% ✔}}

ঘ) ৭\%

\textbf{\hl{✔ সঠিক উত্তর: গ) ৮\%}}

\begin{quote}
\textbf{ব্যাখ্যা:} \emph{ঐকিক নিয়ম:}

\emph{সুদ = (১১২২০ − ৮৫০০) টাকা = ২৭২০ টাকা}

\emph{৮৫০০ টাকার ৪ বছরের সুদ = ২৭২০ টাকা}

\emph{∴ ১ টাকার ১ বছরের সুদ = (২৭২০/৮৫০০×৪) টাকা}

\emph{∴ ১০০ টাকার ১ বছরের সুদ = (২৭২০×১০০/৮৫০০×৪) = ৮ টাকা}

\emph{∴ সুদের হার = ৮\%}

\emph{শর্টটেকনিক:}

\emph{সুদ I = (১১২২০ − ৮৫০০) টাকা = ২৭২০ টাকা}

\emph{এখানে, আসল P = ৮৫০০ টাকা, সময় n = ৪ বছর}

\emph{সুদের হার r = (I×১০০/P×n)}

\emph{= (২৭২০×১০০/৮৫০০×৪) \% = ৮\%}
\end{quote}

\textbf{৪২৫ টাকার ৪ বছরের সুদ ৮৫ টাকা হলে সুদের হার শতকরা বার্ষিক কত টাকা
হবে?}

\textbf{\hl{ক) ৫\% ✔}}

খ) ৯\%

গ) ৪\%

ঘ) ৭\%

\textbf{\hl{✔ সঠিক উত্তর: ক) ৫\%}}

\begin{quote}
\textbf{ব্যাখ্যা:} \emph{ঐকিক নিয়ম:}

\emph{৪২৫ টাকার ৪ বছরের সুদ = ৮৫ টাকা}

\emph{∴ ১ টাকার ১ বছরের সুদ = (৮৫/৪২৫×৪) টাকা}

\emph{∴ ১০০ টাকার ১ বছরের সুদ = (৮৫×১০০/৪২৫×৪) = ৫ টাকা}

\emph{∴ সুদের হার = ৫\%}

\emph{শর্টটেকনিক:}

\emph{এখানে, I = ৮৫ টাকা, P = ৪২৫ টাকা, n = ৪ বছর}

\emph{সুদের হার r = (I×১০০/P×n)}

\emph{= (৮৫×১০০/৪২৫×৪) \% = ৫\%}
\end{quote}

\textbf{৬০০ টাকার ৬ বছরের সুদ ১৮০ টাকা হলে সুদের হার কত?}

\textbf{\hl{ক) ৫\% ✔}}

খ) ১০\%

গ) ১২\%

ঘ) ১৭\%

\textbf{\hl{✔ সঠিক উত্তর: ক) ৫\%}}

\begin{quote}
\textbf{ব্যাখ্যা:} \emph{ঐকিক নিয়ম:}

\emph{৬০০ টাকার ৬ বছরের সুদ = ১৮০ টাকা}

\emph{∴ ১ টাকার ১ বছরের সুদ = (১৮০/৬০০×৬) টাকা}

\emph{∴ ১০০ টাকার ১ বছরের সুদ = (১৮০×১০০/৬০০×৬) = ৫ টাকা}

\emph{∴ সুদের হার = ৫\%}

\emph{শর্টটেকনিক:}

\emph{এখানে, I = ১৮০ টাকা, P = ৬০০ টাকা, n = ৬ বছর}

\emph{সুদের হার r = (I×১০০/P×n)}

\emph{= (১৮০×১০০/৬০০×৬) \% = ৫\%}
\end{quote}

\textbf{শতকরা বার্ষিক কত হার সুদে ৭৫০ টাকার ২ বছরের সুদ ২১০ টাকা?}

ক) ১৫\%

\textbf{\hl{খ) ১৪\% ✔}}

গ) ১২\%

ঘ) ১০\%

\textbf{\hl{✔ সঠিক উত্তর: খ) ১৪\%}}

\begin{quote}
\textbf{ব্যাখ্যা:} \emph{ঐকিক নিয়ম:}

\emph{৭৫০ টাকার ২ বছরের সুদ = ২১০ টাকা}

\emph{∴ ১ টাকার ১ বছরের সুদ = (২১০/৭৫০×২) টাকা}

\emph{∴ ১০০ টাকার ১ বছরের সুদ = (২১০×১০০/৭৫০×২) = ১৪ টাকা}

\emph{∴ সুদের হার = ১৪\%}

\emph{শর্টটেকনিক:}

\emph{এখানে, I = ২১০ টাকা, P = ৭৫০ টাকা, n = ২ বছর}

\emph{সুদের হার r = (I×১০০/P×n)}

\emph{= (২১০×১০০/৭৫০×২) \% = ১৪\%}
\end{quote}

\textbf{আবু ব্যাংকে ৯০০০ টাকা রেখে ২ বছরে ১৮০০ টাকা সুদ পেল। বার্ষিক সুদের
হার কত ছিল?}

ক) ১৫\%

খ) ২৫\%

\textbf{\hl{গ) ১০\% ✔}}

ঘ) ২০\%

\textbf{\hl{✔ সঠিক উত্তর: গ) ১০\%}}

\begin{quote}
\textbf{ব্যাখ্যা:} \emph{ঐকিক নিয়ম:}

\emph{৯০০০ টাকার ২ বছরের সুদ = ১৮০০ টাকা}

\emph{∴ ১ টাকার ১ বছরের সুদ = (১৮০০/৯০০০×২) টাকা}

\emph{∴ ১০০ টাকার ১ বছরের সুদ = (১৮০০×১০০/৯০০০×২) = ১০ টাকা}

\emph{∴ সুদের হার = ১০\%}

\emph{শর্টটেকনিক:}

\emph{এখানে, I = ১৮০০ টাকা, P = ৯০০০ টাকা, n = ২ বছর}

\emph{সুদের হার r = (I×১০০/P×n)}

\emph{= (১৮০০×১০০/৯০০০×২) \% = ১০\%}
\end{quote}

\textbf{শতকরা বার্ষিক কত হার সুদে কোনো আসল ১০ বছরে সুদে-মূলে তিনগুণ হবে?}

ক) ১০\%

খ) ১২\%

গ) ১৫\%

\textbf{\hl{ঘ) ২০\% ✔}}

\textbf{\hl{✔ সঠিক উত্তর: ঘ) ২০\%}}

\begin{quote}
\textbf{ব্যাখ্যা:} \emph{ঐকিক নিয়ম:}

\emph{মনে করি, আসল = ১০০ টাকা}

\emph{∴ সুদে-মূলে ৩ গুণ = (১০০ × ৩) = ৩০০ টাকা}

\emph{∴ সুদ = (৩০০ − ১০০) = ২০০ টাকা}

\emph{১০০ টাকার ১০ বছরের সুদ = ২০০ টাকা}

\emph{∴ ১০০ টাকার ১ বছরের সুদ = (২০০/১০) = ২০ টাকা}

\emph{∴ সুদের হার = ২০\%}

\emph{শর্টটেকনিক:}

\emph{মনে করি, আসল P = ১০০ টাকা, সুদে-মূলে = ৩০০ টাকা ∴ সুদ I = ২০০ টাকা}

\emph{সময় n = ১০ বছর}

\emph{সুদের হার r = (I×১০০/P×n)}

\emph{= (২০০×১০০/১০০×১০) \% = ২০\%}
\end{quote}

\textbf{সুদের হার কত টাকা হলে যে কোনো মূলধন ৮ বছরে সুদে-আসলে তিনগুণ হবে?}

ক) ১৫ টাকা

\textbf{\hl{খ) ২৫ টাকা ✔}}

গ) ২০ টাকা

ঘ) ১২.৫০ টাকা

\textbf{\hl{✔ সঠিক উত্তর: খ) ২৫ টাকা}}

\begin{quote}
\textbf{ব্যাখ্যা:} \emph{ঐকিক নিয়ম:}

\emph{মনে করি, আসল = ১০০ টাকা}

\emph{∴ সুদে-আসলে ৩ গুণ = (১০০ × ৩) = ৩০০ টাকা}

\emph{∴ সুদ = (৩০০ − ১০০) = ২০০ টাকা}

\emph{১০০ টাকার ৮ বছরের সুদ = ২০০ টাকা}

\emph{∴ ১০০ টাকার ১ বছরের সুদ = (২০০/৮) = ২৫ টাকা}

\emph{∴ সুদের হার = ২৫ টাকা (বা ২৫\%)}

\emph{শর্টটেকনিক:}

\emph{মনে করি, আসল P = ১০০ টাকা, সুদে-আসলে = ৩০০ টাকা ∴ সুদ I = ২০০ টাকা}

\emph{সময় n = ৮ বছর}

\emph{সুদের হার r = (I×১০০/P×n)}

\emph{= (২০০×১০০/১০০×৮) = ২৫ টাকা}
\end{quote}

\textbf{বার্ষিক শতকরা কত হার সুদে কোনো আসল ৫ বছরে সুদে আসলে দ্বিগুণ হবে?}

ক) ১০\%

\textbf{\hl{খ) ২০\% ✔}}

গ) ৫\%

ঘ) ১৫\%

\textbf{\hl{✔ সঠিক উত্তর: খ) ২০\%}}

\begin{quote}
\textbf{ব্যাখ্যা:} \emph{ঐকিক নিয়ম:}

\emph{মনে করি, আসল = ১০০ টাকা}

\emph{∴ সুদে-আসলে দ্বিগুণ = (১০০ × ২) = ২০০ টাকা}

\emph{∴ সুদ = (২০০ − ১০০) = ১০০ টাকা}

\emph{১০০ টাকার ৫ বছরের সুদ = ১০০ টাকা}

\emph{∴ ১০০ টাকার ১ বছরের সুদ = (১০০/৫) = ২০ টাকা}

\emph{∴ সুদের হার = ২০\%}

\emph{শর্টটেকনিক:}

\emph{মনে করি, আসল P = ১০০ টাকা, সুদে-আসলে = ২০০ টাকা ∴ সুদ I = ১০০ টাকা}

\emph{সময় n = ৫ বছর}

\emph{সুদের হার r = (I×১০০/P×n)}

\emph{= (১০০×১০০/১০০×৫) \% = ২০\%}
\end{quote}

\textbf{শতকরা বার্ষিক কত হার সুদে যে কোন মূলধন ১০ বছরে সুদে-আসলে দ্বিগুণ হয়?}

\textbf{\hl{ক) ১০ টাকা ✔}}

খ) ১২ টাকা

গ) ৭.৫০ টাকা

ঘ) ৫ টাকা

\textbf{\hl{✔ সঠিক উত্তর: ক) ১০ টাকা}}

\begin{quote}
\textbf{ব্যাখ্যা:} \emph{ঐকিক নিয়ম:}

\emph{মনে করি, আসল = ১০০ টাকা}

\emph{∴ সুদে-আসলে দ্বিগুণ = (১০০ × ২) = ২০০ টাকা}

\emph{∴ সুদ = (২০০ − ১০০) = ১০০ টাকা}

\emph{১০০ টাকার ১০ বছরের সুদ = ১০০ টাকা}

\emph{∴ ১০০ টাকার ১ বছরের সুদ = (১০০/১০) = ১০ টাকা}

\emph{∴ সুদের হার = ১০ টাকা (বা ১০\%)}

\emph{শর্টটেকনিক:}

\emph{মনে করি, আসল P = ১০০ টাকা, সুদে-আসলে = ২০০ টাকা ∴ সুদ I = ১০০ টাকা}

\emph{সময় n = ১০ বছর}

\emph{সুদের হার r = (I×১০০/P×n)}

\emph{= (১০০×১০০/১০০×১০) = ১০ টাকা}
\end{quote}

\textbf{শতকরা বার্ষিক কত হার সুদে ৪২৫ টাকা ৩ বছরে সুদে আসলে ৪৭৬ টাকা হবে?}

ক) ৩\%

খ) ৬\%

গ) ৫\%

\textbf{\hl{ঘ) ৪\% ✔}}

\textbf{\hl{✔ সঠিক উত্তর: ঘ) ৪\%}}

\begin{quote}
\textbf{ব্যাখ্যা:} \emph{ঐকিক নিয়ম:}

\emph{সুদ = (৪৭৬ − ৪২৫) টাকা = ৫১ টাকা}

\emph{৪২৫ টাকার ৩ বছরের সুদ = ৫১ টাকা}

\emph{∴ ১ টাকার ১ বছরের সুদ = (৫১/৪২৫×৩) টাকা}

\emph{∴ ১০০ টাকার ১ বছরের সুদ = (৫১×১০০/৪২৫×৩) = ৪ টাকা}

\emph{∴ সুদের হার = ৪\%}

\emph{শর্টটেকনিক:}

\emph{সুদ I = (৪৭৬ − ৪২৫) টাকা = ৫১ টাকা}

\emph{এখানে, P = ৪২৫ টাকা, n = ৩ বছর}

\emph{সুদের হার r = (I×১০০/P×n)}

\emph{= (৫১×১০০/৪২৫×৩) \% = ৪\%}
\end{quote}

\textbf{শতকরা বার্ষিক কত হার সুদে ৪৫০ টাকা ১৭ বছরের সুদ ৪৫৯ টাকা হবে?}

ক) ৩\%

খ) ৪\%

গ) ৫\%

\textbf{\hl{ঘ) ৬\% ✔}}

\textbf{\hl{✔ সঠিক উত্তর: ঘ) ৬\%}}

\begin{quote}
\textbf{ব্যাখ্যা:} \emph{ঐকিক নিয়ম:}

\emph{৪৫০ টাকার ১৭ বছরের সুদ = ৪৫৯ টাকা}

\emph{∴ ১ টাকার ১ বছরের সুদ = (৪৫৯/৪৫০×১৭) টাকা}

\emph{∴ ১০০ টাকার ১ বছরের সুদ = (৪৫৯×১০০/৪৫০×১৭) = ৬ টাকা}

\emph{∴ সুদের হার = ৬\%}

\emph{শর্টটেকনিক:}

\emph{এখানে, I = ৪৫৯ টাকা, P = ৪৫০ টাকা, n = ১৭ বছর}

\emph{সুদের হার r = (I×১০০/P×n)}

\emph{= (৪৫৯×১০০/৪৫০×১৭) \% = ৬\%}
\end{quote}

\textbf{৪০০ টাকার ৮ মাসের সুদ এবং ৮০০ টাকার ৪ মাসের সুদ একত্রে ২৪ টাকা হলে
শতকরা বার্ষিক সুদের হার কত?}

ক) ৩\%

খ) ৬(১/৪) \%

গ) ৫\%

\textbf{\hl{ঘ) ৪(১/২) \% ✔}}

\textbf{\hl{✔ সঠিক উত্তর: ঘ) ৪(১/২) \%}}

\begin{quote}
\textbf{ব্যাখ্যা:} \emph{বীজগণিতীয় পদ্ধতি:}

\emph{ধরি, সুদের হার = r\%}

\emph{৮ মাস = (৮/১২) = (২/৩) বছর এবং ৪ মাস = (৪/১২) = (১/৩) বছর}

\emph{৪০০ টাকার ৮ মাসের সুদ = ৪০০ × (২/৩) × (r/১০০) = (৮০০r/৩) টাকা}

\emph{৮০০ টাকার ৪ মাসের সুদ = ৮০০ × (১/৩) × (r/১০০) = (৮০০r/৩) টাকা}

\emph{প্রশ্নমতে,}

\emph{(৮০০r/৩) + (৮০০r/৩) = ২৪}

\emph{⇒ (১৬০০r/৩) = ২৪}

\emph{⇒ ১৬০০r = ৭২}

\emph{⇒ r = (৭২×১০০/১৬০০) = ৪(১/২) \%}

\emph{ঐকিক নিয়ম (১ বছরের বিচারে সমতুল্য আসল):}

\emph{৪০০ টাকার ৮ মাসের সুদ = (৪০০ × ৮) = ৩২০০ টাকার ১ মাসের সুদ}

\emph{৮০০ টাকার ৪ মাসের সুদ = (৮০০ × ৪) = ৩২০০ টাকার ১ মাসের সুদ}

\emph{∴ মোট (৩২০০ + ৩২০০) = ৬৪০০ টাকার ১ মাসের সুদ = ২৪ টাকা}

\emph{∴ ৬৪০০ টাকার ১২ মাস বা ১ বছরের সুদ = (২৪ × ১২) = ২৮৮ টাকা}

\emph{∴ ১০০ টাকার ১ বছরের সুদ = (২৮৮×১০০/৬৪০০) = (৯/২) = ৪(১/২) টাকা}
\end{quote}

\textbf{৫০০ টাকার ৪ বছরের সুদ এবং ৬০০ টাকার ৫ বছরের সুদ একত্রে ৫০০ টাকা
হলে, সুদের হার কত?}

ক) ৫\%

খ) ৬\%

\textbf{\hl{গ) ১০\% ✔}}

ঘ) ১২\%

\textbf{\hl{✔ সঠিক উত্তর: গ) ১০\%}}

\begin{quote}
\textbf{ব্যাখ্যা:} \emph{বীজগণিতীয় পদ্ধতি:}

\emph{ধরি, মুনাফার হার = r\%}

\emph{৫০০ টাকার ৪ বছরের সুদ = (৫০০×৪×r/১০০) = ২০r টাকা}

\emph{৬০০ টাকার ৫ বছরের সুদ = (৬০০×৫×r/১০০) = ৩০r টাকা}

\emph{প্রশ্নমতে,}

\emph{২০r + ৩০r = ৫০০}

\emph{⇒ ৫০r = ৫০০}

\emph{∴ r = ১০\%}

\emph{ঐকিক নিয়ম (১ বছরের বিচারে সমতুল্য আসল):}

\emph{৫০০ টাকার ৪ বছরের সুদ = (৫০০ × ৪) = ২০০০ টাকার ১ বছরের সুদ}

\emph{৬০০ টাকার ৫ বছরের সুদ = (৬০০ × ৫) = ৩০০০ টাকার ১ বছরের সুদ}

\emph{∴ মোট (২০০০ + ৩০০০) = ৫০০০ টাকার ১ বছরের সুদ = ৫০০ টাকা}

\emph{∴ ১০০ টাকার ১ বছরের সুদ = (৫০০×১০০/৫০০০) = ১০ টাকা}

\emph{∴ সুদের হার = ১০\%}
\end{quote}

\textbf{শতকরা বার্ষিক কত হার সুদে ৯০০ টাকার ৫ বছরের সুদ ১৩৫ টাকা হবে?}

ক) ২\%

\textbf{\hl{খ) ৩\% ✔}}

গ) ৪\%

ঘ) ৪.৫\%

\textbf{\hl{✔ সঠিক উত্তর: খ) ৩\%}}

\begin{quote}
\textbf{ব্যাখ্যা:} \emph{ঐকিক নিয়ম:}

\emph{৯০০ টাকার ৫ বছরের সুদ = ১৩৫ টাকা}

\emph{∴ ১ টাকার ১ বছরের সুদ = (১৩৫/৯০০×৫) টাকা}

\emph{∴ ১০০ টাকার ১ বছরের সুদ = (১৩৫×১০০/৯০০×৫) = ৩ টাকা}

\emph{∴ সুদের হার = ৩\%}

\emph{শর্টটেকনিক:}

\emph{এখানে, I = ১৩৫ টাকা, P = ৯০০ টাকা, n = ৫ বছর}

\emph{সুদের হার r = (I×১০০/P×n)}

\emph{= (১৩৫×১০০/৯০০×৫) \% = ৩\%}
\end{quote}

\textbf{শতকরা বার্ষিক কত হার সুদে ৫ বছরে ৪০০ টাকার সুদ ১৪০ টাকা হবে?}

\textbf{\hl{ক) ৭ টাকা ✔}}

খ) ৫ টাকা

গ) ৩ টাকা

ঘ) ১২ টাকা

\textbf{\hl{✔ সঠিক উত্তর: ক) ৭ টাকা}}

\begin{quote}
\textbf{ব্যাখ্যা:} \emph{ঐকিক নিয়ম:}

\emph{৪০০ টাকার ৫ বছরের সুদ = ১৪০ টাকা}

\emph{∴ ১ টাকার ১ বছরের সুদ = (১৪০/৪০০×৫) টাকা}

\emph{∴ ১০০ টাকার ১ বছরের সুদ = (১৪০×১০০/৪০০×৫) = ৭ টাকা}

\emph{∴ সুদের হার = ৭ টাকা (বা ৭\%)}

\emph{শর্টটেকনিক:}

\emph{এখানে, I = ১৪০ টাকা, P = ৪০০ টাকা, n = ৫ বছর}

\emph{সুদের হার r = (I×১০০/P×n)}

\emph{= (১৪০×১০০/৪০০×৫) = ৭ টাকা}
\end{quote}

\textbf{শতকরা বার্ষিক কত হার সুদে ৭০০ টাকার ৫ বছরের সুদ ১০৫ টাকা হবে?}

\textbf{\hl{ক) ৩\% ✔}}

খ) ৫\%

গ) ৭\%

ঘ) ১০\%

\textbf{\hl{✔ সঠিক উত্তর: ক) ৩\%}}

\begin{quote}
\textbf{ব্যাখ্যা:} \emph{ঐকিক নিয়ম:}

\emph{৭০০ টাকার ৫ বছরের সুদ = ১০৫ টাকা}

\emph{∴ ১ টাকার ১ বছরের সুদ = (১০৫/৭০০×৫) টাকা}

\emph{∴ ১০০ টাকার ১ বছরের সুদ = (১০৫×১০০/৭০০×৫) = ৩ টাকা}

\emph{∴ সুদের হার = ৩\%}

\emph{শর্টটেকনিক:}

\emph{এখানে, I = ১০৫ টাকা, P = ৭০০ টাকা, n = ৫ বছর}

\emph{সুদের হার r = (I×১০০/P×n)}

\emph{= (১০৫×১০০/৭০০×৫) \% = ৩\%}
\end{quote}

\textbf{শতকরা বার্ষিক কত হার সুদে এমন ১০০০০ টাকায় ৫ বছরে ২৫০০ সুদ টাকা
হবে?}

ক) ৪২\%

খ) ৪৩\%

গ) ৪৪\%

\textbf{\hl{ঘ) ৫০\% ✔}}

\textbf{\hl{✔ সঠিক উত্তর: ঘ) ৫০\%}}

\begin{quote}
\textbf{ব্যাখ্যা:} \emph{ঐকিক নিয়ম:}

\emph{১০০০ টাকার ৫ বছরের সুদ = ২৫০০ টাকা}

\emph{∴ ১ টাকার ১ বছরের সুদ = (২৫০০/১০০০×৫) টাকা}

\emph{∴ ১০০ টাকার ১ বছরের সুদ = (২৫০০×১০০/১০০০×৫) = ৫০ টাকা}

\emph{∴ সুদের হার = ৫০\%}

\emph{শর্টটেকনিক:}

\emph{এখানে, I = ২৫০০ টাকা, P = ১০০০০ টাকা, n = ৫ বছর}

\emph{সুদের হার r = (I×১০০/P×n)}

\emph{= (২৫০০×১০০/১০০০×৫) \% = ৫০\%}
\end{quote}

\textbf{শতকরা বার্ষিক কত হার সুদে ১ বছরের সুদ আসলের (১/৫) অংশ হবে?}

ক) ১০\%

খ) ১৫\%

\textbf{\hl{গ) ২০\% ✔}}

ঘ) ২৫\%

\textbf{\hl{✔ সঠিক উত্তর: গ) ২০\%}}

\begin{quote}
\textbf{ব্যাখ্যা:} \emph{ঐকিক নিয়ম:}

\emph{মনে করি, আসল = ১০০ টাকা}

\emph{∴ ১ বছরের সুদ = ১০০ এর (১/৫) = ২০ টাকা}

\emph{∴ ১০০ টাকার ১ বছরের সুদ ২০ টাকা∴ সুদের হার = ২০\%}

\emph{শর্টটেকনিক:}

\emph{এখানে, P = ১০০ টাকা, I = ২০ টাকা, n = ১ বছর}

\emph{সুদের হার r = (I×১০০/P×n)}

\emph{= (২০×১০০/১০০×১) \% = ২০\%}
\end{quote}

\textbf{শতকরা বার্ষিক কত হার সুদে ৫ বছরের সুদ, সুদাসলের (১/৫) অংশ হবে?}

\textbf{\hl{ক) ৫\% ✔}}

খ) ১০\%

গ) ২০\%

ঘ) ২৫\%

\textbf{\hl{✔ সঠিক উত্তর: ক) ৫\%}}

\begin{quote}
\textbf{ব্যাখ্যা:} \emph{ঐকিক নিয়ম:}

\emph{ধরি, সুদাসল = ১০০ টাকা}

\emph{∴ ৫ বছরের সুদ = ১০০ এর (১/৫) = ২০ টাকা}

\emph{∴ আসল = (১০০ − ২০) = ৮০ টাকা}

\emph{৮০ টাকার ৫ বছরের সুদ = ২০ টাকা}

\emph{∴ ১ টাকার ১ বছরের সুদ = (২০/৮০×৫) টাকা}

\emph{∴ ১০০ টাকার ১ বছরের সুদ = (২০×১০০/৮০×৫) = ৫ টাকা}

\emph{∴ সুদের হার = ৫\%}

\emph{শর্টটেকনিক:}

\emph{এখানে, I = ২০ টাকা, P = ৮০ টাকা, n = ৫ বছর}

\emph{সুদের হার r = (I×১০০/P×n)}

\emph{= (২০×১০০/৮০×৫) \% = ৫\%}
\end{quote}

\textbf{শতকরা বার্ষিক কত হার সুদে ১ বছরের সুদ, সুদাসলের (১/৫) অংশ হবে?}

ক) ১০\%

খ) ১৫\%

গ) ২০\%

\textbf{\hl{ঘ) ২৫\% ✔}}

\textbf{\hl{✔ সঠিক উত্তর: ঘ) ২৫\%}}

\begin{quote}
\textbf{ব্যাখ্যা:} \emph{ঐকিক নিয়ম:}

\emph{ধরি, সুদাসল = ১০০ টাকা}

\emph{∴ ১ বছরের সুদ = ১০০ এর (১/৫) = ২০ টাকা}

\emph{∴ আসল = (১০০ − ২০) = ৮০ টাকা}

\emph{৮০ টাকার ১ বছরের সুদ = ২০ টাকা}

\emph{∴ ১ টাকার ১ বছরের সুদ = (২০/৮০) টাকা}

\emph{∴ ১০০ টাকার ১ বছরের সুদ = (২০×১০০/৮০) = ২৫ টাকা}

\emph{∴ সুদের হার = ২৫\%}

\emph{শর্টটেকনিক:}

\emph{এখানে, I = ২০ টাকা, P = ৮০ টাকা, n = ১ বছর}

\emph{সুদের হার r = (I×১০০/P×n)}

\emph{= (২০×১০০/৮০×১) \% = ২৫\%}
\end{quote}

\textbf{কোনো আসল ৩ বছরে সুদ আসলে ৪৬০ টাকা এবং ৫ বছরে সুদে আসলে ৫০০ টাকা
হলে শতকরা সুদের হার কত?}

ক) ১০ টাকা

খ) ২০ টাকা

গ) ১৫ টাকা

\textbf{\hl{ঘ) ৫ টাকা ✔}}

\textbf{\hl{✔ সঠিক উত্তর: ঘ) ৫ টাকা}}

\begin{quote}
\textbf{ব্যাখ্যা:} \emph{ঐকিক নিয়ম / বিস্তারিত সমাধান:}

\emph{আসল + ৫ বছরের সুদ = ৫০০ টাকা}

\emph{আসল + ৩ বছরের সুদ = ৪৬০ টাকা}

\emph{বিয়োগ করে, ২ বছরের সুদ = ৪০ টাকা}

\emph{∴ ১ বছরের সুদ = (৪০ ÷ ২) = ২০ টাকা}

\emph{∴ ৫ বছরের সুদ = (২০ × ৫) = ১০০ টাকা}

\emph{∴ আসল = (৫০০ − ১০০) = ৪০০ টাকা}

\emph{৪০০ টাকার ৫ বছরের সুদ = ১০০ টাকা}

\emph{∴ ১০০ টাকার ১ বছরের সুদ = (১০০×১০০/৪০০×৫) = ৫ টাকা}

\emph{∴ সুদের হার = ৫ টাকা (বা ৫\%)}

\emph{শর্টটেকনিক:}

\emph{এখানে, I = ১০০ টাকা, P = ৪০০ টাকা, n = ৫ বছর}

\emph{সুদের হার r = (I×১০০/P×n)}

\emph{= (১০০×১০০/৪০০×৫) = ৫ টাকা}
\end{quote}

\textbf{একই হার সুদে ৩০০ টাকার ৪ বছরের সুদ এবং ৫০০ টাকার ৫ বছরের সুদ একত্রে
১৮৫ টাকা হলে, সুদের হার কত?}

ক) ৪\%

\textbf{\hl{খ) ৫\% ✔}}

গ) ৬(১/৪) \%

ঘ) ৭(১/২) \%

\textbf{\hl{✔ সঠিক উত্তর: খ) ৫\%}}

\begin{quote}
\textbf{ব্যাখ্যা:} \emph{বীজগণিতীয় পদ্ধতি:}

\emph{ধরি, মুনাফার হার = r\%}

\emph{৩০০ টাকার ৪ বছরের সুদ = (৩০০×৪×r/১০০) = ১২r টাকা}

\emph{৫০০ টাকার ৫ বছরের সুদ = (৫০০×৫×r/১০০) = ২৫r টাকা}

\emph{প্রশ্নমতে,}

\emph{১২r + ২৫r = ১৮৫}

\emph{⇒ ৩৭r = ১৮৫}

\emph{∴ r = (১৮৫/৩৭) = ৫\%}

\emph{ঐকিক নিয়ম (১ বছরের বিচারে সমতুল্য আসল):}

\emph{৩০০ টাকার ৪ বছরের সুদ = (৩০০ × ৪) = ১২০০ টাকার ১ বছরের সুদ}

\emph{৫০০ টাকার ৫ বছরের সুদ = (৫০০ × ৫) = ২৫০০ টাকার ১ বছরের সুদ}

\emph{∴ মোট (১২০০ + ২৫০০) = ৩৭০০ টাকার ১ বছরের সুদ = ১৮৫ টাকা}

\emph{∴ ১০০ টাকার ১ বছরের সুদ = (১৮৫×১০০/৩৭০০) = ৫ টাকা}

\emph{∴ সুদের হার = ৫\%}
\end{quote}

\textbf{১০\% মুনাফায় ৩০০০ টাকা এবং ৮\% মুনাফায় ২০০০ টাকা বিনিয়োগ করলে মোট
মূলধনের উপর গড়ে শতকরা কত হারে মুনাফা পাওয়া যাবে?}

ক) ৯\%

\textbf{\hl{খ) ৯.২\% ✔}}

গ) ৮\%

ঘ) ৮.২\%

\textbf{\hl{✔ সঠিক উত্তর: খ) ৯.২\%}}

\begin{quote}
\textbf{ব্যাখ্যা:} \emph{বিস্তারিত সমাধান:}

\emph{১০\% মুনাফায় ৩০০০ টাকার ১ বছরের মুনাফা = ৩০০০ × (১০/১০০) = ৩০০ টাকা}

\emph{৮\% মুনাফায় ২০০০ টাকার ১ বছরের মুনাফা = ২০০০ × (৮/১০০) = ১৬০ টাকা}

\emph{∴ মোট মুনাফা = (৩০০ + ১৬০) = ৪৬০ টাকা}

\emph{মোট মূলধন = (৩০০০ + ২০০০) = ৫০০০ টাকা}

\emph{৫০০০ টাকায় ১ বছরের মুনাফা ৪৬০ টাকা}

\emph{∴ ১০০ টাকায় ১ বছরের মুনাফা = (৪৬০×১০০/৫০০০) = ৯.২ টাকা}

\emph{∴ গড় মুনাফার হার = ৯.২\%}

\emph{শর্টটেকনিক:}

\emph{এখানে, মোট আসল P = ৫০০০ টাকা, মোট মুনাফা I = ৪৬০ টাকা, সময় n = ১
বছর}

\emph{গড় মুনাফার হার r = (I×১০০/P×n)}

\emph{= (৪৬০×১০০/৫০০০×১) \% = ৯.২\%}
\end{quote}

\textbf{৮\% সরল মুনাফায় ৬০০০ টাকা বিনিয়োগে ৫ বছরে যে মুনাফা হয়, কোন সরল
হারে বিনিয়োগে ১০০০০ টাকার ৩ বছরে ঐ মুনাফা হবে?}

ক) ০.১

খ) ০.১২

গ) ০.০৯

\textbf{\hl{ঘ) ০.০৮ ✔}}

\textbf{\hl{✔ সঠিক উত্তর: ঘ) ০.০৮}}

\begin{quote}
\textbf{ব্যাখ্যা:} \emph{বিস্তারিত সমাধান:}

\emph{৮\% সরল মুনাফায় ৬০০০ টাকার ৫ বছরের মুনাফা = (৬০০০×৮×৫/১০০) = ২৪০০
টাকা}

\emph{এখন, ১০০০০ টাকার ৩ বছরের মুনাফা ২৪০০ টাকা হতে হবে।}

\emph{∴ ১ টাকার ১ বছরের মুনাফা = (২৪০০/১০০০০×৩) টাকা}

\emph{∴ ১০০ টাকার ১ বছরের মুনাফা = (২৪০০×১০০/১০০০০×৩) = ৮ টাকা}

\emph{∴ মুনাফার হার = ৮\% = (৮/১০০) = ০.০৮}

\emph{শর্টটেকনিক:}

\emph{১ম ক্ষেত্রে, I = (৬০০০×৮×৫/১০০) = ২৪০০ টাকা}

\emph{২য় ক্ষেত্রে, r = (I×১০০/P×n) = (২৪০০×১০০/১০০০০×৩) \% = ৮\% = ০.০৮}
\end{quote}

\textbf{শতকরা বার্ষিক কত হার সুদে কোনো মূলধন ২৫ বছরে সুদে-মূলে ৪ গুণ হবে?}

ক) ১৫\%

খ) ১৬\%

গ) ৮\%

\textbf{\hl{ঘ) ১২\% ✔}}

\textbf{\hl{✔ সঠিক উত্তর: ঘ) ১২\%}}

\begin{quote}
\textbf{ব্যাখ্যা:} \emph{ঐকিক নিয়ম:}

\emph{মনে করি, আসল = ১০০ টাকা}

\emph{∴ সুদে-মূলে ৪ গুণ = (১০০ × ৪) = ৪০০ টাকা}

\emph{∴ সুদ = (৪০০ − ১০০) = ৩০০ টাকা}

\emph{১০০ টাকার ২৫ বছরের সুদ = ৩০০ টাকা}

\emph{∴ ১০০ টাকার ১ বছরের সুদ = (৩০০/২৫) = ১২ টাকা}

\emph{∴ সুদের হার = ১২\%}

\emph{শর্টটেকনিক:}

\emph{মনে করি, আসল P = ১০০ টাকা, সুদে-মূলে = ৪০০ টাকা ∴ সুদ I = ৩০০ টাকা}

\emph{সময় n = ২৫ বছর}

\emph{সুদের হার r = (I×১০০/P×n)}

\emph{= (৩০০×১০০/১০০×২৫) \% = ১২\%}
\end{quote}

\textbf{৫\% হারে ২০ বছরের সুদ ২৫০ টাকা হলে, আসল---}

\textbf{\hl{ক) ২৫০ টাকা ✔}}

খ) ৩০০ টাকা

গ) ৩৫০ টাকা

ঘ) ৪০০ টাকা

\textbf{\hl{✔ সঠিক উত্তর: ক) ২৫০ টাকা}}

\begin{quote}
\textbf{ব্যাখ্যা:} \emph{বীজগণিতীয় পদ্ধতি:}

\emph{ধরি, আসল P টাকা}

\emph{সরল মুনাফার ক্ষেত্রে, I = pnr}

\emph{⇒ ২৫০ = p × ২০ × (৫/১০০)}

\emph{⇒ ২৫০ = p × (১০০/১০০)}

\emph{∴ p = ২৫০ টাকা}

\emph{ঐকিক নিয়ম:}

\emph{৫\% হারে ১০০ টাকার ১ বছরের সুদ = ৫ টাকা}

\emph{∴ ১০০ টাকার ২০ বছরের সুদ = (৫ × ২০) = ১০০ টাকা}

\emph{সুদ ১০০ টাকা হলে আসল = ১০০ টাকা}

\emph{∴ সুদ ১ টাকা হলে আসল = (১০০/১০০) টাকা}

\emph{∴ সুদ ২৫০ টাকা হলে আসল = (১০০×২৫০/১০০) = ২৫০ টাকা}
\end{quote}

\textbf{৬\% বার্ষিক মুনাফায় কত টাকার বার্ষিক মুনাফা ১২০ টাকা?}

ক) ১২০০ টাকা

খ) ১০০০ টাকা

গ) ১৫০০ টাকা

\textbf{\hl{ঘ) ২০০০ টাকা ✔}}

\textbf{\hl{✔ সঠিক উত্তর: ঘ) ২০০০ টাকা}}

\begin{quote}
\textbf{ব্যাখ্যা:} \emph{বীজগণিতীয় পদ্ধতি:}

\emph{আমরা জানি, I = pnr}

\emph{⇒ ১২০ = p × ১ × (৬/১০০)}

\emph{⇒ ৬p = ১২০০০}

\emph{∴ p = ২০০০ টাকা}

\emph{ঐকিক নিয়ম:}

\emph{৬ টাকা বার্ষিক মুনাফা হয় = ১০০ টাকায়}

\emph{∴ ১ টাকা বার্ষিক মুনাফা হয় = (১০০/৬) টাকায়}

\emph{∴ ১২০ টাকা বার্ষিক মুনাফা হয় = (১০০×১২০/৬) = ২০০০ টাকায়}
\end{quote}

\textbf{বার্ষিক শতকরা ১০ টাকা মুনাফায় ২০০০ টাকা ব্যাংকে জমা রাখলে দ্বিতীয়
বছরান্তে সবৃদ্ধি মূলধন কত?}

ক) ২২০০ টাকা

খ) ২৩০০ টাকা

\textbf{\hl{গ) ২৪০০ টাকা ✔}}

ঘ) ২৫০০ টাকা

\textbf{\hl{✔ সঠিক উত্তর: গ) ২৪০০ টাকা}}

\begin{quote}
\textbf{ব্যাখ্যা:} \emph{বীজগণিতীয় পদ্ধতি:}

\emph{আমরা জানি, C = P + I {[}এখানে C = সবৃদ্ধিমূল, P = আসল এবং I =
মুনাফা{]}}

\emph{= P + (pnr/১০০)}

\emph{= P × ১+(nr/১০০)}

\emph{= ২০০০ × (১০০+২×১০/১০০)}

\emph{= ২০ × ১২০ = ২৪০০ টাকা}

\emph{ঐকিক নিয়ম:}

\emph{১০০ টাকার ১ বছরের মুনাফা = ১০ টাকা}

\emph{∴ ২০০০ টাকার ২ বছরের মুনাফা = (১০×২০০০০×২/১০০) = ৪০০ টাকা}

\emph{∴ সবৃদ্ধি মূলধন (আসল + মুনাফা) = (২০০০ + ৪০০) = ২৪০০ টাকা}
\end{quote}

\textbf{কোন মূলধন ৫\% সরল সুদে ২০ বছরে সুদে-আসলে ৫০,০০০ টাকায় হয়। মূলধন কত
টাকা?}

ক) ২০০০০

খ) ২২০০০

গ) ১৬০০০

\textbf{\hl{ঘ) ২৫০০০ ✔}}

\textbf{\hl{✔ সঠিক উত্তর: ঘ) ২৫০০০}}

\begin{quote}
\textbf{ব্যাখ্যা:} \emph{ঐকিক নিয়ম:}

\emph{১০০ টাকার ১ বছরের সুদ ৫ টাকা}

\emph{∴ ১০০ টাকার ২০ বছরের সুদ = (২০ × ৫) = ১০০ টাকা}

\emph{∴ সুদাসল = (১০০ + ১০০) টাকা = ২০০ টাকা}

\emph{সুদাসল ২০০ টাকা হলে আসল = ১০০ টাকা}

\emph{∴ সুদাসল ১ টাকা হলে আসল = (১০০/২০০) টাকা}

\emph{∴ সুদাসল ৫০,০০০ টাকা হলে আসল = (১০০×৫০,০০০/২০০) = ২৫০০০ টাকা}

\emph{বীজগণিতীয় পদ্ধতি:}

\emph{ধরি, আসল = P টাকা}

\emph{আমরা জানি, I = pnr = P × ২০ × (৫/১০০) = P টাকা}

\emph{প্রশ্নমতে,}

\emph{আসল + সুদ = ৫০,০০০}

\emph{⇒ P + P = ৫০,০০০}

\emph{⇒ ২P = ৫০,০০০}

\emph{∴ P = ২৫০০০ টাকা}
\end{quote}

\textbf{কোনো আসল ৩ বছরে সুদে আসলে ৪৬০ টাকা এবং ৫ বছরে সুদে আসলে ৫০০ টাকা
হলে আসল কত টাকা?}

ক) 300

খ) 350

\textbf{\hl{গ) 400 ✔}}

ঘ) 450

\textbf{\hl{✔ সঠিক উত্তর: গ) 400}}

\begin{quote}
\textbf{ব্যাখ্যা:} \emph{ঐকিক নিয়ম / বিয়োগ পদ্ধতি:}

\emph{৫ বছরের সুদ + আসল = ৫০০ টাকা}

\emph{৩ বছরের সুদ + আসল = ৪৬০ টাকা}

\emph{∴ ২ বছরের সুদ = ৪০ টাকা}

\emph{∴ ৩ বছরের সুদ = (৪০×৩/২) = ৬০ টাকা}

\emph{∴ আসল = (৪৬০ − ৬০) = ৪০০ টাকা}

\emph{বীজগণিতীয়/সূত্রভিত্তিক পদ্ধতি:}

\emph{ধরি, আসল = P টাকা এবং ১ বছরের সুদ = I টাকা}

\emph{P + ৫I = ৫০০ ...... (i)}

\emph{P + ৩I = ৪৬০ ...... (ii)}

\emph{(i) হতে (ii) বিয়োগ করে পাই, ২I = ৪০ ⇒ I = ২০ টাকা}

\emph{∴ ৩ বছরের সুদ (৩I) = ৩ × ২০ = ৬০ টাকা}

\emph{∴ আসল P = ৪৬০ − ৬০ = ৪০০ টাকা}
\end{quote}

\textbf{বার্ষিক শতকরা কত হার সুদে ৪২৫ টাকা ৩ বছরে সুদে-আসলে ৪৭৬ টাকা হবে?}

ক) 3\%

খ) 6\%

গ) 5\%

\textbf{\hl{ঘ) 4\% ✔}}

\textbf{\hl{✔ সঠিক উত্তর: ঘ) 4\%}}

\begin{quote}
\textbf{ব্যাখ্যা:} \emph{বীজগণিতীয় পদ্ধতি:}

\emph{সুদ I = (৪৭৬ − ৪২৫) = ৫১ টাকা}

\emph{আমরা জানি, I = pnr}

\emph{⇒ ৫১ = ৪২৫ × ৩ × (r/১০০)}

\emph{⇒ ৫১ = (৫১r/৪)}

\emph{⇒ r = (৪×৫১/৫১)}

\emph{⇒ r = ৪}

\emph{∴ সুদের হার ৪\%}

\emph{ঐকিক নিয়ম:}

\emph{সুদ = (৪৭৬ − ৪২৫) = ৫১ টাকা}

\emph{৪২৫ টাকার ৩ বছরের সুদ = ৫১ টাকা}

\emph{∴ ১ টাকার ১ বছরের সুদ = (৫১/৪২৫×৩) টাকা}

\emph{∴ ১০০ টাকার ১ বছরের সুদ = (৫১×১০০/৪২৫×৩) = ৪ টাকা}

\emph{∴ সুদের হার ৪\%}
\end{quote}

\textbf{কোন আসল ৩ বছরে সুদে-আসলে ৪৬০ টাকা এবং ৫ বছর সুদে-আসলে ৫০০ টাকা হলে
আসল কত টাকা?}

ক) ৪৫০

খ) ৪৫৫

গ) ৩০০

\textbf{\hl{ঘ) ৪০০ ✔}}

\textbf{\hl{✔ সঠিক উত্তর: ঘ) ৪০০}}

\begin{quote}
\textbf{ব্যাখ্যা:} \emph{ঐকিক নিয়ম / বিয়োগ পদ্ধতি:}

\emph{৫ বছরের সুদ + আসল = ৫০০ টাকা}

\emph{৩ বছরের সুদ + আসল = ৪৬০ টাকা}

\emph{{[}(−) করে{]} ২ বছরের সুদ = ৪০ টাকা}

\emph{∴ ৩ বছরের সুদ = (৪০×৩/২) = ৬০ টাকা}

\emph{∴ আসল = ৪৬০ − ৬০ = ৪০০ টাকা}

\emph{শর্টকার্ট পদ্ধতি:}

\emph{Step-1: ১ বছরের সুদ = (৫০-৪০/৫−৩) = (৪০/২) = ২০ টাকা}

\emph{৫ বছরের সুদ = ২০ × ৫ = ১০০ টাকা}

\emph{Step-2: আসল, P = ৫০০ − ১০০ = ৪০০ টাকা}
\end{quote}

\textbf{শতকরা সুদের হার ৮ টাকা থেকে বেড়ে ১২ টাকা হওয়ায় জলিল সাহেবের আয় ৪
বছরে ২৫৬ টাকা বেড়ে গেল। তার মূলধন কত?}

\textbf{\hl{ক) ১৬০০ টাকা ✔}}

খ) ১২০০ টাকা

গ) ১৪০০ টাকা

ঘ) ১৮০০ টাকা

\textbf{\hl{✔ সঠিক উত্তর: ক) ১৬০০ টাকা}}

\begin{quote}
\textbf{ব্যাখ্যা:} \emph{ঐকিক নিয়ম:}

\emph{১০০ টাকায় ১ বছরে আয় বাড়ে (১২ − ৮) টাকা = ৪ টাকা}

\emph{১০০ টাকায় ৪ বছরে আয় বাড়ে (৪ × ৪) টাকা = ১৬ টাকা}

\emph{১৬ টাকা আয় বাড়ে যখন মূলধন ১০০ টাকা}

\emph{∴ ১ টাকা আয় বাড়ে যখন মূলধন (১০০/১৬) টাকা}

\emph{∴ ২৫৬ টাকা আয় বাড়ে যখন মূলধন (১০০×২৫৬/১৬) টাকা = ১৬০০ টাকা}

\emph{বীজগণিতীয় পদ্ধতি:}

\emph{ধরি, মূলধন P টাকা}

\emph{সুদের হারের পার্থক্য = (১২ − ৮)\% = ৪\%}

\emph{আমরা জানি, I = pnr}

\emph{⇒ ২৫৬ = P × ৪ × (৪/১০০)}

\emph{⇒ ২৫৬ = (১৬P/১০০)}

\emph{⇒ ১৬P = ২৫৬০০}

\emph{∴ P = ১৬০০ টাকা}
\end{quote}

\textbf{বার্ষিক ৮\% সুদে ১,২০০ টাকায় ৫ বছরের যে সুদ হয়, বার্ষিক ৬\% সুদে কত
টাকায় ১০ বছরের তত সুদ হবে?}

ক) ৬০০ টাকা

\textbf{\hl{খ) ৮০০ টাকা ✔}}

গ) ১০০০ টাকা

ঘ) ১৪০০ টাকা

\textbf{\hl{✔ সঠিক উত্তর: খ) ৮০০ টাকা}}

\begin{quote}
\textbf{ব্যাখ্যা:} \emph{বীজগণিতীয়/ধাপভিত্তিক পদ্ধতি:}

\emph{৮\% হার মুনাফায় ১,২০০ টাকার ৫ বছরের মুনাফা = (৮×১২০০×৫/১০০) = ৪৮০
টাকা}

\emph{এখন ধরি, আসল x টাকা}

\emph{প্রশ্নমতে, (x×৬×১০/১০০) = ৪৮০}

\emph{⇒ (৬০x/১০০) = ৪৮০}

\emph{∴ x = (৪৮০×১০০/৬×১০) = ৮০০ টাকা}

\emph{সমীকরণীয় সংক্ষেপ:}

\emph{ধরি, ২য় ক্ষেত্রে মূলধন = P\_2}

\emph{প্রশ্নমতে, P\_1n\_1r\_1 = P\_2n\_2r\_2}

\emph{⇒ ১২০০ × ৫ × ৮\% = P\_2 × ১০ × ৬\%}

\emph{⇒ ৪৮০০০ = ৬০ × P\_2}

\emph{∴ P\_2 = (৪৮০০০/৬০) = ৮০০ টাকা}
\end{quote}

\textbf{বার্ষিক সুদের হার ৫\% থেকে হ্রাস পেয়ে ৪ (৩/৪) \% হওয়ায় এক ব্যক্তির ৮০
টাকা আয় কমে গেল। তার মূলধন কত?}

ক) ১৬,০০,০০০ টাকা

খ) ১,৬০,০০০ টাকা

\textbf{\hl{গ) ১৬,০০০ টাকা ✔}}

ঘ) ১,৬০০ টাকা

\textbf{\hl{✔ সঠিক উত্তর: গ) ১৬,০০০ টাকা}}

\begin{quote}
\textbf{ব্যাখ্যা:} \emph{ঐকিক নিয়ম:}

\emph{সুদের হার কমে = ৫−৪(৩/৪) \% = ৫−(১৯/৪) \% = (২০−১৯/৪) \% = (১/৪)
\%}

\emph{(১/৪) টাকা আয় কমে মূলধন ১০০ টাকা হলে}

\emph{∴ ১ টাকা আয় কমে মূলধন (১০০ × ৪) টাকা}

\emph{∴ ৮০ টাকা আয় কমে মূলধন (১০০ × ৪ × ৮০) = ১৬,০০০ টাকা}

\emph{বীজগণিতীয় পদ্ধতি:}

\emph{ধরি, মূলধন P টাকা}

\emph{সুদের হারের পার্থক্য, r = (১/৪) \%}

\emph{সময়, n = ১ বছর এবং সুদের হ্রাস, I = ৮০ টাকা}

\emph{আমরা জানি, I = pnr}

\emph{⇒ ৮০ = P × ১ × (১/৪×১০০)}

\emph{⇒ ৮০ = (P/৪০০)}

\emph{∴ P = (৮০ × ৪০০) = ১৬,০০০ টাকা}
\end{quote}

\textbf{সেলিম ৬\% সরল সুদে ব্যাংকে ১০,০০০ টাকা বিনিয়োগ করে। আর কত টাকা ৮\%
সরল সুদে বিনিয়োগ করলে সে মোটের উপর ৭\% হারে সুদ পাবে?}

ক) ১৮, ৫০০ টাকা

খ) ১৮, ৭৫০ টাকা

গ) ১৯,০০০ টাকা

\textbf{\hl{ঘ) ১০,০০০ টাকা ✔}}

\textbf{\hl{✔ সঠিক উত্তর: ঘ) ১০,০০০ টাকা}}

\begin{quote}
\textbf{ব্যাখ্যা:} \emph{বীজগণিতীয় পদ্ধতি:}

\emph{ধরি, ৮\% হারে x টাকা বিনিয়োগ করতে হবে।}

\emph{প্রশ্নমতে, (১০০০০×৬/১০০) + (৮x/১০০) = (১০০০০+x×৭/১০০)}

\emph{⇒ ৬০,০০০ + ৮x = ৭০,০০০ + ৭x}

\emph{⇒ ৮x − ৭x = ৭০,০০০ − ৬০,০০০}

\emph{∴ x = ১০,০০০ টাকা}

\emph{গড়/মিশ্রণ ভিত্তিক পদ্ধতি:}

\emph{প্রথম অংশ হতে সুদের হার = ৬\% এবং দ্বিতীয় অংশ হতে সুদের হার = ৮\%}

\emph{গড় সুদের হার = ৭\%}

\emph{৬\% এবং ৮\% এর ঠিক মাঝখানের মান হলো ৭\% (কারণ ৭ − ৬ = ১ এবং ৮ − ৭
= ১)}

\emph{যেহেতু উভয় অংশের সুদের হারের ব্যবধান সমান, তাই দুটি বিনিয়োগের পরিমাণের
অনুপাত হবে ১ : ১।}

\emph{∴ দ্বিতীয় অংশে বিনিয়োগের পরিমাণ = ১০,০০০ টাকা।}
\end{quote}

\textbf{শতকরা বার্ষিক যে হারে কোন মূলধন ৬ বছরে সুদে মূলে দ্বিগুণ হয় সেই হারে
কত টাকা ৪ বছরে সুদে মূলে ২০৫০ টাকা হবে?}

ক) ১২২৫

\textbf{\hl{খ) ১২৩০ ✔}}

গ) ১২৪০

ঘ) ১২৮৫

\textbf{\hl{✔ সঠিক উত্তর: খ) ১২৩০}}

\begin{quote}
\textbf{ব্যাখ্যা:} \emph{ঐকিক/পাটিগণিতীয় পদ্ধতি:}

\emph{মনেকরি, আসল = ১০০ টাকা}

\emph{∴ ১০০ টাকায় ৬ বছরের সুদাসল হবে = ২০০ টাকা}

\emph{∴ ১০০ টাকায় ৬ বছরের সুদ = (২০০ − ১০০) টাকা = ১০০ টাকা}

\emph{∴ ১০০ টাকায় ১ বছরের সুদ = (১০০/৬) টাকা}

\emph{∴ ১০০ টাকায় ৪ বছরের সুদ = (১০০×৪/৬) = (২০০/৩) টাকা}

\emph{∴ সুদাসল = ১০০+(২০০/৩) = (৫০০/৩) টাকা}

\emph{সুদাসল (৫০০/৩) টাকা হলে আসল = ১০০ টাকা}

\emph{∴ সুদাসল ২০৫০ টাকা হলে আসল = (১০০×৩×২০৫০/৫০০) = ১২৩০ টাকা}

\emph{বীজগণিতীয়/সূত্রভিত্তিক পদ্ধতি:}

\emph{১ম ক্ষেত্রে, p(১ + nr) = ২p}

\emph{⇒ nr = ২ − ১ = ১}

\emph{⇒ ৬r = ১}

\emph{∴ r = (১/৬)}

\emph{২য় ক্ষেত্রে, p(১ + nr) = ২০৫০}

\emph{⇒ p ১+৪×(১/৬) = ২০৫০}

\emph{⇒ p ১+(২/৩) = ২০৫০}

\emph{⇒ p (৫/৩) = ২০৫০}

\emph{⇒ p = (২০৫০×৩/৫)}

\emph{∴ p = ১২৩০ টাকা}
\end{quote}

\textbf{কাদের দুইটি সঞ্চয়ী হিসাবে মোট ১৫০০০ টাকা জমা রাখল। একটি হিসাব থেকে
সে বাৎসরিক ৫\% হারে সুদ পাবে এবং অপরটি থেকে ১০\% হারে সুদ পাবে। বছর শেষে
সে যদি মোট ১১১০ টাকা সুদ পেয়ে থাকে, তাহলে ৫\% হার সুদে সে কত টাকা জমা
রেখেছিল?}

\textbf{\hl{ক) ৭৮০০ ✔}}

খ) ৭২০০

গ) ৬৮০০

ঘ) ৭০০০

\textbf{\hl{✔ সঠিক উত্তর: ক) ৭৮০০}}

\begin{quote}
\textbf{ব্যাখ্যা:} \emph{বীজগণিতীয় পদ্ধতি:}

\emph{ধরি, কাদের ৫\% হারে জমা রাখে x টাকা}

\emph{∴ ১০\% হারে জমা রাখে (১৫০০০ − x) টাকা}

\emph{প্রশ্নমতে, x×(৫/১০০)×১ + ১৫০০০−x × (১০/১০০) × ১ = ১১১০}

\emph{⇒ (x/২০) + ১৫০০ − (x/১০) = ১১১০}

\emph{⇒ (x−২x/২০) = ১১১০ − ১৫০০}

\emph{⇒ (−x/২০) = − ৩৯০}

\emph{⇒ − x = ২০ × (− ৩৯০)}

\emph{∴ x = ৭৮০০ টাকা}

\emph{গড়/মিশ্রণ ভিত্তিক পদ্ধতি:}

\emph{মোট টাকার ওপর প্রাপ্ত গড় সুদের হার = (১১১০/১৫০০০)×১০০ \% = ৭.৪\%}

\emph{প্রথম অংশ (৫\%) ও দ্বিতীয় অংশের (১০\%) সুদের হারের ব্যবধান:}

\emph{১০\% − ৭.৪\% = ২.৬\% এবং ৭.৪\% − ৫\% = ২.৪\%}

\emph{∴ ৫\% ও ১০\% সুদে জমা রাখা টাকার অনুপাত = ২.৬ : ২.৪ = ২৬ : ২৪ = ১৩ :
১২}

\emph{অনুপাতদ্বয়ের যোগফল = ১৩ + ১২ = ২৫}

\emph{∴ ৫\% সুদে জমা রাখা টাকার পরিমাণ = ১৫০০০ × (১৩/২৫) = ৭৮০০ টাকা}
\end{quote}

\textbf{কোন আসল ৫\% সুদে ৫ বছরে সুদাসলে ৫২৫ টাকা হলে আসল কত?}

ক) ৪০০ টাকা

\textbf{\hl{খ) ৪২০ টাকা ✔}}

গ) ৪৫০ টাকা

ঘ) ৫০০ টাকা

\textbf{\hl{✔ সঠিক উত্তর: খ) ৪২০ টাকা}}

\begin{quote}
\textbf{ব্যাখ্যা:} \emph{ঐকিক নিয়ম:}

\emph{১০০ টাকার ১ বছরের সুদ = ৫ টাকা}

\emph{১০০ টাকার ৫ বছরের সুদ = ২৫ টাকা}

\emph{সুদাসল = ১০০ + ২৫ = ১২৫ টাকা}

\emph{সুদাসল ১২৫ টাকা হলে আসল = ১০০ টাকা}

\emph{∴ সুদাসল ১ টাকা হলে আসল = (১০০/১২৫) টাকা}

\emph{∴ সুদাসল ৫২৫ টাকা হলে আসল = (১০০×৫২৫/১২৫) = ৪২০ টাকা}

\emph{বীজগণিতীয় পদ্ধতি:}

\emph{ধরি, আসল = P টাকা}

\emph{আমরা জানি, I = pnr = P × ৫ × (৫/১০০) = (P/৪)}

\emph{প্রশ্নমতে, P + (P/৪) = ৫২৫}

\emph{⇒ (৫P/৪) = ৫২৫}

\emph{∴ P = (৫২৫×৪/৫) = ৪২০ টাকা}
\end{quote}

\textbf{কোন মূলধন ৪ বছরের সুদে মূলে ৬০০ টাকা এবং ৬ বছরে সুদে মূলে ৭০০ টাকা
হলে, মূলধন কত?}

ক) ৬৫০ টাকা

\textbf{\hl{খ) ৪০০ টাকা ✔}}

গ) ১০০ টাকা

ঘ) ১৩০০ টাকা

\textbf{\hl{✔ সঠিক উত্তর: খ) ৪০০ টাকা}}

\begin{quote}
\textbf{ব্যাখ্যা:} \emph{ঐকিক / বিয়োগ পদ্ধতি:}

\emph{৬ বছরের সুদ + আসল = ৭০০ টাকা}

\emph{৪ বছরের সুদ + আসল = ৬০০ টাকা}

\emph{(−) করে, ২ বছরের মুনাফা = (৭০০ − ৬০০) টাকা = ১০০ টাকা}

\emph{∴ ১ বছরের মুনাফা = ৫০ টাকা}

\emph{∴ ৬ বছরের মুনাফা = (৫০ × ৬) = ৩০০ টাকা}

\emph{∴ মূলধন = (৭০০ − ৩০০) = ৪০০ টাকা}

\emph{বীজগণিতীয় পদ্ধতি:}

\emph{ধরি, মূলধন = P এবং ১ বছরের মুনাফা = I}

\emph{P + ৬I = ৭০০ ...... (i)}

\emph{P + ৪I = ৬০০ ...... (ii)}

\emph{(i) হতে (ii) বিয়োগ করে পাই, ২I = ১০০ ⇒ I = ৫০ টাকা}

\emph{∴ ৪ বছরের মুনাফা (৪I) = ৪ × ৫০ = ২০০ টাকা}

\emph{∴ মূলধন P = ৬০০ − ২০০ = ৪০০ টাকা}
\end{quote}

\textbf{৫\% সুদে ২০ বৎসর সুদাসলে ৫০ হাজার টাকা হলে, মূলধন কত?}

ক) ২০০০০/-

\textbf{\hl{খ) ২৫০০০/- ✔}}

গ) ৩০০০০/-

ঘ) ৩৫০০০/-

\textbf{\hl{✔ সঠিক উত্তর: খ) ২৫০০০/-}}

\begin{quote}
\textbf{ব্যাখ্যা:} \emph{ঐকিক নিয়ম:}

\emph{১০০ টাকার ১ বছরের সুদ = ৫ টাকা}

\emph{∴ ৫\% সুদে ১০০ টাকার ২০ বছরের সুদ = (৫ × ২০) = ১০০ টাকা}

\emph{সুদাসল = (১০০ + ১০০) = ২০০ টাকা}

\emph{সুদাসল ২০০ টাকা হলে আসল = ১০০ টাকা}

\emph{∴ সুদাসল ১ টাকা হলে আসল = (১০০/২০০) টাকা}

\emph{∴ সুদাসল ৫০০০০ টাকা হলে আসল = (১০০×৫০০০০/২০০) = ২৫০০০ টাকা}

\emph{বীজগণিতীয় পদ্ধতি:}

\emph{ধরি, মূলধন = P টাকা}

\emph{মুনাফা I = P × ২০ × (৫/১০০) = P টাকা}

\emph{প্রশ্নমতে, P + P = ৫০,০০০}

\emph{⇒ ২P = ৫০,০০০}

\emph{∴ P = ২৫০০০ টাকা}
\end{quote}

\textbf{শতকরা সুদের হার ৮ টাকা থেকে বেড়ে ১২ টাকা হওয়ায় জলিল সাহেবের আয় ৪
বছরে ২৫৬ টাকা বেড়ে গেল। তার মূলধন কত?}

ক) ১২০০ টাকা

\textbf{\hl{খ) ১৬০০ টাকা ✔}}

গ) ১৪০০ টাকা

ঘ) ১৮০০ টাকা

\textbf{\hl{✔ সঠিক উত্তর: খ) ১৬০০ টাকা}}

\begin{quote}
\textbf{ব্যাখ্যা:} \emph{ঐকিক নিয়ম:}

\emph{১০০ টাকায় ১ বছরে আয় বাড়ে = (১২ − ৮) = ৪ টাকা}

\emph{১০০ টাকায় ৪ বছরে আয় বাড়ে = (৪ × ৪) = ১৬ টাকা}

\emph{১৬ টাকা আয় বাড়ে যখন মূলধন = ১০০ টাকা}

\emph{∴ ১ টাকা আয় বাড়ে যখন মূলধন = (১০০/১৬) টাকা}

\emph{∴ ২৫৬ টাকা আয় বাড়ে যখন মূলধন = (১০০×২৫৬/১৬) = ১৬০০ টাকা}

\emph{বীজগণিতীয় পদ্ধতি:}

\emph{ধরি, মূলধন = P টাকা}

\emph{সুদের হারের পার্থক্য, r = (১২ − ৮)\% = ৪\%}

\emph{সময়, n = ৪ বছর এবং আয় বৃদ্ধি, I = ২৫৬ টাকা}

\emph{আমরা জানি, I = pnr}

\emph{⇒ ২৫৬ = P × ৪ × (৪/১০০)}

\emph{⇒ ২৫৬ = (১৬P/১০০)}

\emph{∴ P = (২৫৬×১০০/১৬) = ১৬০০ টাকা}
\end{quote}

\textbf{শতকরা বার্ষিক ৬ টাকা হার সুদে ৯৫০ টাকার ৮ বছরে যত সুদ হয়, বার্ষিক
৭.৫\% হার সুদে কত টাকার ১৯ বছরে তত সুদ হবে?}

ক) ২৮০ টাকা

\textbf{\hl{খ) ৩২০ টাকা ✔}}

গ) ৩৮০ টাকা

ঘ) ৪৯০ টাকা

\textbf{\hl{✔ সঠিক উত্তর: খ) ৩২০ টাকা}}

\begin{quote}
\textbf{ব্যাখ্যা:} \emph{ঐকিক/ধাপভিত্তিক পদ্ধতি:}

\emph{৬\% সুদে ৯৫০ টাকার ৮ বছরের সুদ = (৬×৯৫০×৮/১০০) = ৪৫৬ টাকা}

\emph{আবার, ৭.৫\% সুদে ১০০ টাকার ১৯ বছরের সুদ = (৭.৫ × ১৯) = ১৪২.৫ টাকা}

\emph{১৯ বছরে সুদ ১৪২.৫ টাকা হলে আসল = ১০০ টাকা}

\emph{∴ ১৯ বছরে সুদ ৪৫৬ টাকা হলে আসল = (৪৫৬×১০০/১৪২.৫) = ৩২০ টাকা}

\emph{বীজগণিতীয় সমীকরণ পদ্ধতি:}

\emph{ধরি, নির্ণেয় আসল = P\_2 টাকা}

\emph{প্রশ্নমতে, P\_1n\_1r\_1 = P\_2n\_2r\_2}

\emph{⇒ ৯৫০ × ৮ × ৬\% = P\_2 × ১৯ × ৭.৫\%}

\emph{⇒ ৪৫৬ = P\_2 × ১.৪২৫}

\emph{∴ P\_2 = (৪৫৬/১.৪২৫) = ৩২০ টাকা}
\end{quote}

\textbf{বার্ষিক ৮\% সরল সুদে কত টাকা ৬ বছরে সুদে-\/-আসলে ১০৩৬ টাকা হবে?}

\textbf{\hl{ক) ৭০০ টাকা ✔}}

খ) ৬০০ টাকা

গ) ৬৫০ টাকা

ঘ) ৫৫০ টাকা

\textbf{\hl{✔ সঠিক উত্তর: ক) ৭০০ টাকা}}

\begin{quote}
\textbf{ব্যাখ্যা:} \emph{ঐকিক নিয়ম:}

\emph{১০০ টাকার ১ বছরের সুদ = ৮ টাকা}

\emph{১০০ টাকার ৬ বছরের সুদ = (৬ × ৮) = ৪৮ টাকা}

\emph{সুদাসল = (১০০ + ৪৮) = ১৪৮ টাকা}

\emph{সুদাসল ১৪৮ টাকা হলে আসল = ১০০ টাকা}

\emph{∴ সুদাসল ১ টাকা হলে আসল = (১০০/১৪৮) টাকা}

\emph{∴ সুদাসল ১০৩৬ টাকা হলে আসল = (১০০×১০৩৬/১৪৮) = ৭০০ টাকা}

\emph{বীজগণিতীয় পদ্ধতি:}

\emph{ধরি, আসল = P টাকা}

\emph{আমরা জানি, C = P(১ + nr)}

\emph{⇒ ১০৩৬ = P × ১+৬×(৮/১০০)}

\emph{⇒ ১০৩৬ = P × ১+(৪৮/১০০)}

\emph{⇒ ১০৩৬ = P × (১৪৮/১০০)}

\emph{∴ P = (১০৩৬×১০০/১৪৮) = ৭০০ টাকা}
\end{quote}

\textbf{কোন আসল সুদাসলে ৩ বছরে ৪৮৪ টাকা এবং ৫ বছরে ৫৪০ টাকা হয়। শতকরা
সুদের হার ও আসল নির্ণয় করুন।}

ক) সুদের হার ৩ টাকা এবং আসল ৩০০ টাকা

খ) সুদের হার ৪ টাকা এবং আসল ৩৫০ টাকা

গ) সুদের হার ৫ টাকা এবং আসল ৩৭৫ টাকা

\textbf{\hl{ঘ) সুদের হার ৭ টাকা এবং আসল ৪০০ টাকা ✔}}

\textbf{\hl{✔ সঠিক উত্তর: ঘ) সুদের হার ৭ টাকা এবং আসল ৪০০ টাকা}}

\begin{quote}
\textbf{ব্যাখ্যা:} \emph{ঐকিক / বিয়োগ পদ্ধতি:}

\emph{৫ বছরের সুদ + আসল = ৫৪০ টাকা}

\emph{৩ বছরের সুদ + আসল = ৪৮৪ টাকা}

\emph{(বিয়োগ করে) ২ বছরের সুদ = ৫৬ টাকা}

\emph{∴ ১ বছরের সুদ = (৫৬ ÷ ২) = ২৮ টাকা}

\emph{∴ ৩ বছরের সুদ = (২৮ × ৩) = ৮৪ টাকা}

\emph{আসল = (৪৮৪ − ৮৪) = ৪০০ টাকা}

\emph{৪০০ টাকার ৩ বছরের সুদ ৮৪ টাকা}

\emph{∴ ১০০ টাকার ১ বছরের সুদ = (৮৪×১০০/৪০০×৩) = ৭ টাকা}

\emph{∴ সুদের হার ৭\% এবং আসল ৪০০ টাকা}

\emph{সূত্রভিত্তিক পদ্ধতি:}

\emph{I = ৮৪ টাকা, P = ৪০০ টাকা, n = ৩ বছর}

\emph{r = (I×১০০/Pn) = (৮৪×১০০/৪০০×৩) = ৭ টাকা বা ৭\%}
\end{quote}

\textbf{মুনাফা ১২\% থেকে হ্রাস পেয়ে ৮\% হলে কত টাকার বার্ষিক মুনাফা ৫০০ টাকা
হ্রাস পাবে?}

ক) ১৫০০০

\textbf{\hl{খ) ১২৫০০ ✔}}

গ) ১২০০০

ঘ) ১৩৫০০

\textbf{\hl{✔ সঠিক উত্তর: খ) ১২৫০০}}

\begin{quote}
\textbf{ব্যাখ্যা:} \emph{ঐকিক নিয়ম:}

\emph{১০০ টাকায় ১ বছরে মুনাফা হ্রাস পায় = (১২\% − ৮\%) = ৪ টাকা}

\emph{৪ টাকা হ্রাস পায় যখন আসল = ১০০ টাকা}

\emph{∴ ১ টাকা হ্রাস পায় যখন আসল = (১০০/৪) টাকা}

\emph{∴ ৫০০ টাকা হ্রাস পায় যখন আসল = (৫০০×১০০/৪) = ১২৫০০ টাকা}

\emph{বীজগণিতীয় পদ্ধতি:}

\emph{ধরি, আসল P টাকা}

\emph{সুদের হারের পার্থক্য, r = (১২ − ৮)\% = ৪\%}

\emph{আমরা জানি, I = pnr}

\emph{⇒ ৫০০ = P × ১ × (৪/১০০)}

\emph{⇒ ৫০০ = (P/২৫)}

\emph{∴ P = (৫০০ × ২৫) = ১২৫০০ টাকা}
\end{quote}

\textbf{সুদের হার ৭\% থেকে কমে ৫\% হলে এক ব্যক্তির আয় ৫ বছরে ৭০ টাকা কমে
যায়। তার মূলধন কত?}

ক) ৬০০

\textbf{\hl{খ) ৭০০ ✔}}

গ) ৮০০

ঘ) ১০০০

\textbf{\hl{✔ সঠিক উত্তর: খ) ৭০০}}

\begin{quote}
\textbf{ব্যাখ্যা:} \emph{ঐকিক নিয়ম:}

\emph{১০০ টাকায় ১ বছরে আয় কমে = (৭\% − ৫\%) = ২ টাকা}

\emph{১০০ টাকায় ৫ বছরে আয় কমে = (২ × ৫) = ১০ টাকা}

\emph{১০ টাকা আয় কমে যখন আসল = ১০০ টাকা}

\emph{∴ ১ টাকা আয় কমে যখন আসল = (১০০/১০) টাকা}

\emph{∴ ৭০ টাকা আয় কমে যখন আসল = (৭০×১০০/১০) = ৭০০ টাকা}

\emph{বীজগণিতীয় পদ্ধতি:}

\emph{ধরি, মূলধন P টাকা}

\emph{সুদের হ্রাস r = (৭ − ৫)\% = ২\%}

\emph{সময় n = ৫ বছর এবং মোট ক্ষতি I = ৭০ টাকা}

\emph{আমরা জানি, I = pnr}

\emph{⇒ ৭০ = P × ৫ × (২/১০০)}

\emph{⇒ ৭০ = (P/১০)}

\emph{∴ P = ৭০০ টাকা}
\end{quote}

\textbf{শতকরা ৫ টাকা হার সুদে ২০ সৎসরে সুদে-আসলে ৪০০০০ টাকা হয়। মূলধনের
পরিমাণ কত?}

ক) ২৫০০০

খ) ১০০০০

গ) ১৫০০০

\textbf{\hl{ঘ) ২০০০০ ✔}}

\textbf{\hl{✔ সঠিক উত্তর: ঘ) ২০০০০}}

\begin{quote}
\textbf{ব্যাখ্যা:} \emph{ঐকিক নিয়ম:}

\emph{১০০ টাকার ১ বছরের সুদ = ৫ টাকা}

\emph{১০০ টাকার ২০ বছরের সুদ = (৫ × ২০) = ১০০ টাকা}

\emph{সুদাসল = (১০০ + ১০০) = ২০০ টাকা}

\emph{সুদাসল ২০০ টাকা হলে আসল = ১০০ টাকা}

\emph{∴ সুদাসল ১ টাকা হলে আসল = (১০০/২০০) টাকা}

\emph{∴ সুদাসল ৪০০০০ টাকা হলে আসল = (৪০০০০×১০০/২০০) = ২০০০০ টাকা}

\emph{বীজগণিতীয় পদ্ধতি:}

\emph{ধরি, আসল P টাকা}

\emph{I = P × ২০ × (৫/১০০) = P টাকা}

\emph{সুদাসল = P + P = ২P}

\emph{প্রশ্নমতে, ২P = ৪০০০০}

\emph{∴ P = ২০০০০ টাকা}
\end{quote}

\textbf{বার্ষিক ৪ (১/২) \% সরল সুদে কত টাকা বিনিয়োগ করলে ৪ বছরে তা সুদে
আসলে ৮২৬ টাকা হয় হবে?}

ক) ৪৫৮ টাকা

খ) ৬৫০ টাকা

\textbf{\hl{গ) ৭০০ টাকা ✔}}

ঘ) ৭২৫ টাকা

\textbf{\hl{✔ সঠিক উত্তর: গ) ৭০০ টাকা}}

\begin{quote}
\textbf{ব্যাখ্যা:} \emph{ঐকিক নিয়ম:}

\emph{১০০ টাকার ১ বছরের সুদ = ৪ (১/২) বা (৯/২) টাকা}

\emph{১০০ টাকার ৪ বছরের সুদ = (৯/২) × ৪ = ১৮ টাকা}

\emph{সুদাসল = (১০০ + ১৮) = ১১৮ টাকা}

\emph{সুদাসল ১১৮ টাকা হলে আসল = ১০০ টাকা}

\emph{∴ সুদাসল ১ টাকা হলে আসল = (১০০/১১৮) টাকা}

\emph{∴ সুদাসল ৮২৬ টাকা হলে আসল = (৮২৬×১০০/১১৮) = ৭০০ টাকা}

\emph{বীজগণিতীয় পদ্ধতি:}

\emph{ধরি, আসল P টাকা}

\emph{সুদের হার r = (৯/২) \%}

\emph{আমরা জানি, C = P(১ + nr)}

\emph{⇒ ৮২৬ = P × ১+৪×(৯/২×১০০)}

\emph{⇒ ৮২৬ = P × ১+(১৮/১০০)}

\emph{⇒ ৮২৬ = P × (১১৮/১০০)}

\emph{∴ P = (৮২৬×১০০/১১৮) = ৭০০ টাকা}
\end{quote}

\textbf{শতকরা বার্ষিক ৪ টাকা হার সরল মুনাফায় কত টাকা ১৫ বছরে সবৃদ্ধিমূলধন
১০৪০ টাকা হবে?}

ক) মূলধন ৫০০ টাকা

খ) মূলধন ৫৫০ টাকা

গ) মূলধন ৬০০ টাকা

\textbf{\hl{ঘ) মূলধন ৬৫০ টাকা ✔}}

\textbf{\hl{✔ সঠিক উত্তর: ঘ) মূলধন ৬৫০ টাকা}}

\begin{quote}
\textbf{ব্যাখ্যা:} \emph{ঐকিক নিয়ম:}

\emph{১০০ টাকার ১ বছরের মুনাফা = ৪ টাকা}

\emph{১০০ টাকার ১৫ বছরের মুনাফা = (১৫ × ৪) = ৬০ টাকা}

\emph{সবৃদ্ধিমূলধন = (১০০ + ৬০) = ১৬০ টাকা}

\emph{সবৃদ্ধিমূলধন ১৬০ টাকা হলে মূলধন = ১০০ টাকা}

\emph{∴ সবৃদ্ধিমূলধন ১ টাকা হলে মূলধন = (১০০/১৬০) টাকা}

\emph{∴ সবৃদ্ধিমূলধন ১০৪০ টাকা হলে মূলধন = (১০৪০×১০০/১৬০) = ৬৫০ টাকা}

\emph{বীজগণিতীয় পদ্ধতি:}

\emph{ধরি, মূলধন P টাকা}

\emph{মুনাফা I = P × ১৫ × (৪/১০০) = (৩P/৫)}

\emph{সবৃদ্ধিমূলধন = P + (৩P/৫) = (৮P/৫)}

\emph{প্রশ্নমতে, (৮P/৫) = ১০৪০}

\emph{∴ P = (১০৪০×৫/৮) = ৬৫০ টাকা}
\end{quote}

\textbf{বার্ষিক ১০\% সরল সুদে কত টাকা ৫ বছরে সুদে-আসলে ৭৫০ টাকা হবে?}

\textbf{\hl{ক) ৫০০ টাকা ✔}}

খ) ৫৫০ টাকা

গ) ৬০০ টাকা

ঘ) ৬৫০ টাকা

\textbf{\hl{✔ সঠিক উত্তর: ক) ৫০০ টাকা}}

\begin{quote}
\textbf{ব্যাখ্যা:} \emph{ঐকিক নিয়ম:}

\emph{১০০ টাকার ১ বছরের সুদ = ১০ টাকা}

\emph{১০০ টাকার ৫ বছরের সুদ = (৫ × ১০) = ৫০ টাকা}

\emph{সুদাসল = (১০০ + ৫০) = ১৫০ টাকা}

\emph{সুদাসল ১৫০ টাকা হলে আসল = ১০০ টাকা}

\emph{∴ সুদাসল ১ টাকা হলে আসল = (১০০/১৫০) টাকা}

\emph{∴ সুদাসল ৭৫০ টাকা হলে আসল = (৭৫০×১০০/১৫০) = ৫০০ টাকা}

\emph{বীজগণিতীয় পদ্ধতি:}

\emph{ধরি, আসল P টাকা}

\emph{I = P × ৫ × (১০/১০০) = (P/২)}

\emph{সুদাসল = P + (P/২) = (৩P/২)}

\emph{প্রশ্নমতে, (৩P/২) = ৭৫০}

\emph{∴ P = (৭৫০×২/৩) = ৫০০ টাকা}
\end{quote}

\textbf{১০\% হার মুনাফায় কত টাকার ৪ বৎসরের সুদ ১৪০ টাকা?}

ক) ৩০০ টাকা

খ) ৪০০ টাকা

\textbf{\hl{গ) ৩৫০ টাকা ✔}}

ঘ) ৪৫০ টাকা

\textbf{\hl{✔ সঠিক উত্তর: গ) ৩৫০ টাকা}}

\begin{quote}
\textbf{ব্যাখ্যা:} \emph{ঐকিক নিয়ম:}

\emph{১০০ টাকার ১ বছরের সুদ = ১০ টাকা}

\emph{১০০ টাকার ৪ বছরের সুদ = (৪ × ১০) = ৪০ টাকা}

\emph{সুদ ৪০ টাকা হলে আসল = ১০০ টাকা}

\emph{∴ সুদ ১ টাকা হলে আসল = (১০০/৪০) টাকা}

\emph{∴ সুদ ১৪০ টাকা হলে আসল = (১৪০×১০০/৪০) = ৩৫০ টাকা}

\emph{সূত্রভিত্তিক পদ্ধতি:}

\emph{এখানে, I = ১৪০ টাকা, n = ৪ বছর, r = ১০\% = (১০/১০০)}

\emph{আমরা জানি, P = (I/nr)}

\emph{⇒ P = (১৪০×১০০/১০×৪)}

\emph{∴ P = ৩৫০ টাকা}
\end{quote}

\textbf{বার্ষিক ১৫\% মুনাফায় কোনো ব্যাংক থেকে কিছু টাকা নিয়ে এক বছর পর ১৬৮০
টাকা মুনাফা দেয়া হলো। আসল কত ছিল?}

ক) ১২২০০

\textbf{\hl{খ) ১১২০০ ✔}}

গ) ১০২০০

ঘ) ১৩২০০

\textbf{\hl{✔ সঠিক উত্তর: খ) ১১২০০}}

\begin{quote}
\textbf{ব্যাখ্যা:} \emph{ঐকিক নিয়ম:}

\emph{১৫ টাকা মুনাফা হলে আসল = ১০০ টাকা}

\emph{∴ ১ টাকা মুনাফা হলে আসল = (১০০/১৫) টাকা}

\emph{∴ ১৬৮০ টাকা মুনাফা হলে আসল = (১৬৮০×১০০/১৫) = ১১২০০ টাকা}

\emph{সূত্রভিত্তিক পদ্ধতি:}

\emph{এখানে, I = ১৬৮০ টাকা, n = ১ বছর, r = ১৫\%}

\emph{আমরা জানি, P = (I×১০০/rn)}

\emph{⇒ P = (১৬৮০×১০০/১৫×১)}

\emph{∴ P = ১১২০০ টাকা}
\end{quote}

\textbf{শতকরা বার্ষিক ৪ টাকা হার সুদে কত টাকার ৫ বছরের সুদ ৪ টাকা হবে?}

ক) ২৫

খ) ৩০

গ) ১৫

\textbf{\hl{ঘ) ২০ ✔}}

\textbf{\hl{✔ সঠিক উত্তর: ঘ) ২০}}

\begin{quote}
\textbf{ব্যাখ্যা:} \emph{ঐকিক নিয়ম:}

\emph{১০০ টাকার ৫ বছরের সুদ = (৪ × ৫) = ২০ টাকা}

\emph{২০ টাকা সুদ হয় যখন আসল = ১০০ টাকা}

\emph{∴ ১ টাকা সুদ হয় যখন আসল = (১০০/২০) টাকা}

\emph{∴ ৪ টাকা সুদ হয় যখন আসল = (১০০×৪/২০) = ২০ টাকা}

\emph{সূত্রভিত্তিক পদ্ধতি:}

\emph{এখানে, I = ৪ টাকা, n = ৫ বছর, r = ৪\%}

\emph{আমরা জানি, P = (I×১০০/rn)}

\emph{⇒ P = (৪×১০০/৪×৫)}

\emph{∴ P = ২০ টাকা}
\end{quote}

\textbf{৫\% হারে ৩০০ টাকার ৫ বছরে সুদাসল কত হয়?}

ক) ৩২৫

খ) ৩৪৫

গ) ৩৬০

\textbf{\hl{ঘ) ৩৭৫ ✔}}

\textbf{\hl{✔ সঠিক উত্তর: ঘ) ৩৭৫}}

\begin{quote}
\textbf{ব্যাখ্যা:} \emph{বীজগণিতীয়/সূত্রভিত্তিক পদ্ধতি:}

\emph{আমরা জানি,}

\emph{আসল = (১০০×সুদাসল/১০০+সময়×হার)}

\emph{⇒ ৩০০ = (১০০×সুদাসল/১০০+৫×৫)}

\emph{⇒ ১০০ × সুদাসল = ৩৭,৫০০}

\emph{∴ সুদাসল = ৩৭৫ টাকা}

\emph{সংক্ষিপ্ত পদ্ধতি (Shortcut):}

\emph{সুদাসল = p(১ + nr)}

\emph{= ৩০০ ১+৫×(৫/১০০)}

\emph{= ৩০০ × ১+(১/৪) = ৩৭৫ টাকা}
\end{quote}

\textbf{৫\% হার সুদে ৬০০ টাকা ৫ বছরে সুদে-আসলে কত টাকা হবে?}

ক) ৭০০

\textbf{\hl{খ) ৭৫০ ✔}}

গ) ৪৯২

ঘ) নড়ড়

\textbf{\hl{✔ সঠিক উত্তর: খ) ৭৫০}}

\begin{quote}
\textbf{ব্যাখ্যা:} \emph{ঐকিক নিয়ম পদ্ধতি:}

\emph{১০০ টাকায় ১ বছরের সুদ = ৫ টাকা।}

\emph{∴ ৬০০ টাকায় ৫ বছরের সুদ = (৫×৬০০×৫/১০০) = ১৫০ টাকা।}

\emph{∴ সুদাসল = ৬০০ + ১৫০ = ৭৫০ টাকা}

\emph{সূত্রভিত্তিক পদ্ধতি:}

\emph{এখানে, p = ৬০০, n = ৫ বছর, r = ৫\% = (৫/১০০)}

\emph{আমরা জানি, I = pnr}

\emph{⇒ I = ৬০০ × ৫ × (৫/১০০) = ১৫০ টাকা}

\emph{∴ সুদাসল (A) = p + I = ৬০০ + ১৫০ = ৭৫০ টাকা}
\end{quote}

\textbf{বার্ষিক ৮\% মুনাফায় ৫০০০ টাকার মুনাফা ১০০০০ টাকা হবে কত বছরে?}

ক) ৩(১/৪)

\textbf{\hl{খ) ২(১/২) ✔}}

গ) ২(১/৪)

ঘ) ১(১/৫)

\textbf{\hl{✔ সঠিক উত্তর: খ) ২(১/২)}}

\begin{quote}
\textbf{ব্যাখ্যা:} \emph{সূত্রভিত্তিক পদ্ধতি:}

\emph{এখানে, আসল P = ৫০০০ টাকা, মুনাফা I = ১০০০০ টাকা, মুনাফার হার r = ৮\%
= (৮/১০০)}

\emph{আমরা জানি, n = (I×১০০/P×r) = (১০০০০×১০০/৫০০০×৮) = (৫/২)}

\emph{∴ n = ২(১/২) বছর}

\emph{ঐকিক নিয়ম পদ্ধতি:}

\emph{৫০০০ টাকার ১ বছরের মুনাফা = ৫০০০ × (৮/১০০) = ৪০০ টাকা}

\emph{৪০ টাকা মুনাফা হয় = ১ বছরে}

\emph{∴ ১ টাকা মুনাফা হয় = (১/৪০) বছরে}

\emph{∴ ১০০০০ টাকা মুনাফা হয় = (১০০০০/৪০) = (৫/২) = ২(১/২) বছরে}
\end{quote}

\textbf{৪৫০ টাকা বার্ষিক ৬\% সুদে কত বছরে সুদে-আসলে ৫৫৮ টাকা হবে?}

ক) ৩ বছরে

\textbf{\hl{খ) ৪ বছরে ✔}}

গ) ৫ বছরে

ঘ) ৬ বছরে

\textbf{\hl{✔ সঠিক উত্তর: খ) ৪ বছরে}}

\begin{quote}
\textbf{ব্যাখ্যা:} \emph{ধাপভিত্তিক/সাধারণ পদ্ধতি:}

\emph{মোট সুদ I = (৫৫৮ − ৪৫০) = ১০৮ টাকা}

\emph{∴ সময় (n) = (১০৮/৪৫০×(৬/১০০)) = ৪ বছর}

\emph{সংক্ষিপ্ত/বীজগণিতীয় পদ্ধতি (Shortcut):}

\emph{p(১ + nr) = ৫৫৮}

\emph{⇒ ৪৫০ ১+n×(৬/১০০) = ৫৫৮}

\emph{⇒ ১ + (৬n/১০০) = (৫৫৮/৪৫০)}

\emph{⇒ (৬n/১০০) = ১.২৪ − ১ = ০.২৪}

\emph{∴ n = ৪ বছর}
\end{quote}

\textbf{বার্ষিক ৮\% হার সুদে কত বছরে সুদ আসলের সমান হবে?}

ক) ৮ বছর

\textbf{\hl{খ) ১২(১/২) বছর ✔}}

গ) ১০ বছর

ঘ) ১৪(১/২) বছর

\textbf{\hl{✔ সঠিক উত্তর: খ) ১২(১/২) বছর}}

\begin{quote}
\textbf{ব্যাখ্যা:} \emph{বীজগণিতীয় পদ্ধতি:}

\emph{ধরি, n বছরে x টাকার সুদ হবে x টাকা}

\emph{এখানে, p = x টাকা, I = x টাকা, r = ৮\% = (৮/১০০)}

\emph{এখন, I = pnr}

\emph{⇒ x = x × n × (৮/১০০)}

\emph{⇒ ১ = n × (২/২৫)}

\emph{⇒ n = (২৫/২)}

\emph{∴ n = ১২(১/২) বছর}

\emph{ঐকিক নিয়ম পদ্ধতি:}

\emph{ধরি, আসল = ১০০ টাকা}

\emph{∴ প্রশ্নমতে, সুদ = ১০০ টাকা}

\emph{১০০ টাকার ১ বছরের সুদ = ৮ টাকা}

\emph{৮ টাকা সুদ হয় = ১ বছরে}

\emph{∴ ১ টাকা সুদ হয় = (১/৮) বছরে}

\emph{∴ ১০০ টাকা সুদ হয় = (১০০/৮) = (২৫/২) = ১২(১/২) বছরে}
\end{quote}

\textbf{কোনো মূলধন ১০ বছরে দ্বিগুণ হয়। ঐ মূলধন তিনগুণ হবে কত বছরে?}

ক) ১২

খ) ১৫

\textbf{\hl{গ) ২০ ✔}}

ঘ) ৩০

\textbf{\hl{✔ সঠিক উত্তর: গ) ২০}}

\begin{quote}
\textbf{ব্যাখ্যা:} \emph{বীজগণিতীয় পদ্ধতি:}

\emph{ধরি, মূলধন p = x টাকা}

\emph{∴ মুনাফা-মূলধন = ২x টাকা}

\emph{∴ মুনাফা = (২x − x) = x টাকা}

\emph{আমরা জানি, I = pnr}

\emph{⇒ x = x × ১০ × r}

\emph{∴ r = (১/১০)}

\emph{আবার, মুনাফা-মূলধন ৩x টাকা}

\emph{∴ মুনাফা = (৩x − x) = ২x টাকা}

\emph{I\_১ = p\_১n\_১r}

\emph{⇒ ২x = x × n\_১ × (১/১০)}

\emph{∴ n\_১ = ২০ বছর।}

\emph{ঐকিক নিয়ম পদ্ধতি:}

\emph{ধরি, আসল = ১০০ টাকা}

\emph{১০ বছরে সুদে-মূলে দ্বিগুণ হলে সুদ = (২০০ − ১০০) = ১০০ টাকা}

\emph{সুদে-মূলে ৩ গুণ হতে হলে সুদ = (৩০০ − ১০০) = ২০০ টাকা}

\emph{১০০ টাকা সুদ হয় = ১০ বছরে}

\emph{∴ ১ টাকা সুদ হয় = (১০/১০০) বছরে}

\emph{∴ ২০০ টাকা সুদ হয় = (১০×২০০/১০০) = ২০ বছরে।}
\end{quote}

\textbf{কোনো আসল ২০ বছরে সুদে-মূলে দ্বিগুণ হলে কত বছরে সুদে-মূলে তিনগুণ হবে?}

ক) ৩০ বছরে

খ) ২৫ বছরে

\textbf{\hl{গ) ৪০ বছরে ✔}}

ঘ) ৬০ বছরে

\textbf{\hl{✔ সঠিক উত্তর: গ) ৪০ বছরে}}

\begin{quote}
\textbf{ব্যাখ্যা:} \emph{ঐকিক নিয়ম পদ্ধতি:}

\emph{মনে করি, আসল ১০০ টাকা}

\emph{∴ ২০ বছরে সুদে-মূলে হবে = ২০০ টাকা}

\emph{∴ সুদ = (২০০ − ১০০) টাকা = ১০০ টাকা}

\emph{সুদে-মূলে তিনগুণ = ৩০০ টাকা।}

\emph{সুদ = (৩০০ − ১০০) টাকা বা ২০০ টাকা}

\emph{১০০ টাকা সুদ হয় = ২০ বছরে}

\emph{∴ ১ টাকা সুদ হয় = (২০/১০০) বছরে}

\emph{∴ ২০০ টাকা সুদ হয় = (২০×২০০/১০০) বছরে = ৪০ বছর।}

\emph{বীজগণিতীয় পদ্ধতি:}

\emph{ধরি, আসল = p টাকা}

\emph{দ্বিগুণ হলে মুনাফা I\_১ = (২p − p) = p টাকা}

\emph{r = (I/pn) = (p/p×২০) = (১/২০)}

\emph{তিনগুণ হলে মুনাফা I\_২ = (৩p − p) = ২p টাকা}

\emph{∴ সময় n = (I\_২/pr) = (২p/p×(১/২০)) = ২ × ২০ = ৪০ বছর।}
\end{quote}

\textbf{একই হার ও মুনাফায় কোনো আসল ৬ বছরে মুনাফা আসলে দ্বিগুণ হলে, কত বছরে
তা মুনাফা-আসলে তিনগুণ হবে?}

ক) ১০ বছর

\textbf{\hl{খ) ১২ বছর ✔}}

গ) ১৩ বছর

ঘ) ১৪ বছর

\textbf{\hl{✔ সঠিক উত্তর: খ) ১২ বছর}}

\begin{quote}
\textbf{ব্যাখ্যা:} \emph{বীজগণিতীয় পদ্ধতি:}

\emph{ধরি, মুনাফার r\% এবং আসল p টাকা}

\emph{৬ বছরে মুনাফা-আসলে দ্বিগুণ হলে মুনাফা I = (২p − p) = p টাকা}

\emph{∴ r = (I/pn) = (p/p×৬) = (১/৬)}

\emph{আবার, ৩ গুণ হলে মুনাফা I\_১ = (৩p − p) = ২p টাকা}

\emph{∴ সময় n = (I\_১/pr) = (২p/p×(১/৬)) = ২ × ৬ = ১২ বছর।}

\emph{ঐকিক নিয়ম পদ্ধতি:}

\emph{ধরি, আসল = ১০০ টাকা}

\emph{দ্বিগুণ হলে মুনাফা = (২০০ − ১০০) = ১০০ টাকা}

\emph{তিনগুণ হলে মুনাফা = (৩০০ − ১০০) = ২০০ টাকা}

\emph{১০০ টাকা মুনাফা হয় = ৬ বছরে}

\emph{∴ ১ টাকা মুনাফা হয় = (৬/১০০) বছরে}

\emph{∴ ২০০ টাকা মুনাফা হয় = (৬×২০০/১০০) = ১২ বছরে।}
\end{quote}

\textbf{বার্ষিক শতকরা ৬ টাকা হার সুদে কত সময়ে ৪৫০ টাকা সুদে-আসলে ৫৫৮ টাকা
হবে?}

ক) ৩ বছর

\textbf{\hl{খ) ৪ বছর ✔}}

গ) ৫ বছর

ঘ) ৬ বছর

\textbf{\hl{✔ সঠিক উত্তর: খ) ৪ বছর}}

\begin{quote}
\textbf{ব্যাখ্যা:} \emph{ঐকিক নিয়ম পদ্ধতি:}

\emph{সুদাসল ৫৫৮ টাকা হলে সুদ = (৫৫৮ − ৪৫০) টাকা = ১০৮ টাকা}

\emph{১০০ টাকার ১ বছরের সুদ = ৬ টাকা}

\emph{∴ ৪৫০ টাকার ১ বছরের সুদ = (৬×৪৫০/১০০) = ২৭ টাকা}

\emph{আবার, ৪৫০ টাকায় ২৭ টাকা সুদ হয় = ১ বছরে}

\emph{∴ ৪৫০ টাকায় ১০৮ টাকা সুদ হয় = (১০৮/২৭) বছরে = ৪ বছরে।}

\emph{বীজগণিতীয়/সূত্রভিত্তিক পদ্ধতি:}

\emph{এখানে, আসল P = ৪৫০ টাকা, সুদ I = (৫৫৮ − ৪৫০) = ১০৮ টাকা, সুদের হার r
= ৬\% = (৬/১০০)}

\emph{আমরা জানি, I = pnr}

\emph{⇒ ১০৮ = ৪৫০ × n × (৬/১০০)}

\emph{⇒ ১০৮ = ২৭n}

\emph{∴ n = (১০৮/২৭) = ৪ বছর।}
\end{quote}

\textbf{বার্ষিক ৫\% হার সুদে ৬ বছরে কোন আসলের সুদ-আসল ৪৪২ টাকা হয় কত বছরে
তা সুদ-আসলে ৫১০ টাকা হবে?}

ক) ৮ বছর

খ) ৯ বছর

\textbf{\hl{গ) ১০ বছর ✔}}

ঘ) ১২ বছর

\textbf{\hl{✔ সঠিক উত্তর: গ) ১০ বছর}}

\begin{quote}
\textbf{ব্যাখ্যা:} \emph{ঐকিক নিয়ম পদ্ধতি:}

\emph{বার্ষিক ৫\% হার সুদে, ৬ বছরে ১০০ টাকার সুদ-আসল ১০০ + (৫ × ৬) = ১৩০
টাকা।}

\emph{সুদ-আসল ১৩০ টাকা হলে আসল ১০০ টাকা}

\emph{∴ সুদ-আসল ৪৪২ টাকা হলে আসল = (১০০ × ৪৪২ ÷ ১৩০) = ৩৪০ টাকা।}

\emph{সুদ-আসলে ৫১০ টাকা ও আসল ৩৪০ টাকা হলে সুদ = (৫১০ − ৩৪০) = ১৭০ টাকা।}

\emph{বার্ষিক ৫\% হার সুদে ১ বছরে ৩৪০ টাকার সুদ = (৫ × ৩৪০) ÷ ১০০ = ১৭
টাকা।}

\emph{৩৪০০ টাকার ১৭ টাকা সুদ হয় ১ বছরে}

\emph{∴ ৩৪০ টাকার ১৭০ টাকা সুদ হয় = ১০ বছরে।}

\emph{বীজগণিতীয় পদ্ধতি:}

\emph{ধরি, আসল = P টাকা}

\emph{১ম ক্ষেত্রে, C\_১ = P(১ + n\_১r) ⇒ ৪৪২ = P ১+৬×(৫/১০০) ⇒ ৪৪২ = P
(১৩০/১০০) ∴ P = ৩৪০ টাকা}

\emph{২য় ক্ষেত্রে, মোট সুদ I\_২ = ৫১০ − ৩৪০ = ১৭০ টাকা}

\emph{আমরা জানি, n\_২ = (I\_২/Pr) = (১৭০/৩৪০×(৫/১০০)) = (১৭০/১৭) = ১০
বছর।}
\end{quote}

\textbf{কোনো আসল ২০ বছরে সুদে-মূলে দ্বিগুণ হলে কত বছরে সুদে মূলে তিনগুণ হবে?}

ক) ৩০ বছরে

খ) ২৫ বছরে

\textbf{\hl{গ) ৪০ বছরে ✔}}

ঘ) ৬০ বছরে

\textbf{\hl{✔ সঠিক উত্তর: গ) ৪০ বছরে}}

\begin{quote}
\textbf{ব্যাখ্যা:} \emph{ঐকিক নিয়ম পদ্ধতি:}

\emph{মনে করি, আসল ১০০ টাকা}

\emph{২০ বছরে সুদে-মূলে দ্বিগুণ = ২০০ টাকা}

\emph{∴ সুদ = (২০০ − ১০০) = ১০০ টাকা}

\emph{সুদে-মূলে তিনগুণ = ৩০০ টাকা}

\emph{∴ সুদ = (৩০০ − ১০০) = ২০০ টাকা}

\emph{∴ ১০০ টাকা সুদ হয় = ২০ বছরে}

\emph{∴ ১ টাকা সুদ হয় = (২০/১০০) বছরে}

\emph{∴ ২০০ টাকা সুদ হয় = (২০×২০০/১০০) = ৪০ বছরে}

\emph{বীজগণিতীয় পদ্ধতি:}

\emph{ধরি, আসল = P টাকা}

\emph{দ্বিগুণ হলে, মুনাফা I\_১ = P টাকা}

\emph{r = (I\_১/Pn) = (P/P×২০) = (১/২০)}

\emph{তিনগুণ হলে, মুনাফা I\_২ = (৩P − P) = ২P টাকা}

\emph{∴ সময় n = (I\_২/Pr) = (২P/P×(১/২০)) = ৪০ বছর}
\end{quote}

\textbf{৬\% সুদে কত সময়ে ৪৫০ টাকা সুদে-মূলে ৫৫৮ টাকা হবে?}

ক) ৩ বছর

\textbf{\hl{খ) ৪ বছর ✔}}

গ) ৫ বছর

ঘ) ৬ বছর

\textbf{\hl{✔ সঠিক উত্তর: খ) ৪ বছর}}

\begin{quote}
\textbf{ব্যাখ্যা:} \emph{ঐকিক নিয়ম পদ্ধতি:}

\emph{সুদ = (৫৫৮ − ৪৫০) = ১০৮ টাকা}

\emph{১০০ টাকার ১ বছরের সুদ = ৬ টাকা}

\emph{∴ ১ টাকার ১ বছরের সুদ = (৬/১০০) টাকা}

\emph{∴ ৪৫০ টাকার ১ বছরের সুদ = (৬×৪৫০/১০০) = ২৭ টাকা}

\emph{২৭ টাকা সুদ হয় = ১ বছরে}

\emph{∴ ১ টাকা সুদ হয় = (১/২৭) বছরে}

\emph{∴ ১০৮ টাকা সুদ হয় = (১০৮/২৭) = ৪ বছরে}

\emph{সূত্রভিত্তিক/সংক্ষিপ্ত পদ্ধতি:}

\emph{এখানে, I = ১০৮ টাকা, P = ৪৫০ টাকা, r = ৬\% = (৬/১০০)}

\emph{আমরা জানি, n = (I×১০০/Pr)}

\emph{⇒ n = (১০৮×১০০/৪৫০×৬) = ৪ বছর}
\end{quote}

\textbf{বার্ষিক শতকরা ৫ টাকা হার সুদে কত সময়ে ৩০০ টাকা। সুদআসলে ৪০৫ টাকা
হবে?}

ক) ৫ বছর

খ) ৬ বছর

\textbf{\hl{গ) ৭ বছর ✔}}

ঘ) ৮ বছর

\textbf{\hl{✔ সঠিক উত্তর: গ) ৭ বছর}}

\begin{quote}
\textbf{ব্যাখ্যা:} \emph{ঐকিক নিয়ম পদ্ধতি:}

\emph{সুদ = (৪০৫ − ৩০০) = ১০৫ টাকা}

\emph{১০০ টাকার ১ বছরের সুদ = ৫ টাকা}

\emph{∴ ৩০০ টাকার ১ বছরের সুদ = (৫×৩০০/১০০) = ১৫ টাকা}

\emph{১৫ টাকা সুদ হয় = ১ বছরে}

\emph{∴ ১০৫ টাকা সুদ হয় = (১০৫/১৫) = ৭ বছরে}

\emph{সূত্রভিত্তিক/সংক্ষিপ্ত পদ্ধতি:}

\emph{এখানে, I = ১০৫ টাকা, P = ৩০০ টাকা, r = ৫\%}

\emph{আমরা জানি, n = (I×১০০/Pr)}

\emph{⇒ n = (১০৫×১০০/৩০০×৫) = ৭ বছর}
\end{quote}

\textbf{৫\% হারে কত সময়ে ৫০০ টাকার মুনাফা ১০০ টাকা হবে?}

\textbf{\hl{ক) ৪ বছর ✔}}

খ) ৩ বছর

গ) ২ বছর

ঘ) ৫ বছরে

\textbf{\hl{✔ সঠিক উত্তর: ক) ৪ বছর}}

\begin{quote}
\textbf{ব্যাখ্যা:} \emph{ঐকিক নিয়ম পদ্ধতি:}

\emph{সুদ = ১০০ টাকা}

\emph{১০০ টাকার ১ বছরের সুদ = ৫ টাকা}

\emph{∴ ৫০০ টাকার ১ বছরের সুদ = (৫×৫০০/১০০) = ২৫ টাকা}

\emph{২৫ টাকা সুদ হয় = ১ বছরে}

\emph{∴ ১০০ টাকা সুদ হয় = (১০০/২৫) = ৪ বছরে}

\emph{সূত্রভিত্তিক/সংক্ষিপ্ত পদ্ধতি:}

\emph{এখানে, I = ১০০ টাকা, P = ৫০০ টাকা, r = ৫\%}

\emph{আমরা জানি, n = (I×১০০/Pr)}

\emph{⇒ n = (১০০×১০০/৫০০×৫) = ৪ বছর}
\end{quote}

\textbf{৮৮৮৮ টাকা বার্ষিক ১০\% সুদে কত বছরে সুদ আসলে ১৭৭৭৬ টাকা হবে?}

ক) ২০ বছরে

খ) ৩০ বছরে

\textbf{\hl{গ) ১০ বছরে ✔}}

ঘ) ২৫ বছরে

\textbf{\hl{✔ সঠিক উত্তর: গ) ১০ বছরে}}

\begin{quote}
\textbf{ব্যাখ্যা:} \emph{ঐকিক নিয়ম পদ্ধতি:}

\emph{সুদ = (১৭৭৭৬ − ৮৮৮৮) = ৮৮৮৮ টাকা}

\emph{১০০ টাকার ১ বছরের সুদ = ১০ টাকা}

\emph{∴ ৮৮৮৮ টাকার ১ বছরের সুদ = (১০×৮৮৮৮/১০০) = ৮৮৮.৮ টাকা}

\emph{৮৮৮.৮ টাকা সুদ হয় = ১ বছরে}

\emph{∴ ৮৮৮৮ টাকা সুদ হয় = (৮৮৮৮/৮৮৮.৮) = ১০ বছরে}

\emph{সূত্রভিত্তিক/সংক্ষিপ্ত পদ্ধতি:}

\emph{এখানে, I = ৮৮৮৮ টাকা, P = ৮৮৮৮ টাকা, r = ১০\%}

\emph{আমরা জানি, n = (I×১০০/Pr)}

\emph{⇒ n = (৮৮৮৮×১০০/৮৮৮৮×১০) = ১০ বছর}
\end{quote}

\textbf{৮ মাসে ২০০ টাকায় যত লাভ হয়, কত মাসে ৪০০ টাকায় তত লাভ হয়?}

ক) ২ মাস

খ) ৩ মাস

\textbf{\hl{গ) ৪ মাস ✔}}

ঘ) ৬ মাস

\textbf{\hl{✔ সঠিক উত্তর: গ) ৪ মাস}}

\begin{quote}
\textbf{ব্যাখ্যা:} \emph{অনুপাত/যৌক্তিক পদ্ধতি:}

\emph{যেহেতু সুদের হার ও মোট লাভের পরিমাণ অপরিবর্তিত, তাই মূলধন বাড়লে সময়
সমানুপাতিক হারে কমবে।}

\emph{মূলধন ২০০ টাকার স্থলে ৪০০ টাকা অর্থাৎ (৪০০ ÷ ২০০) = ২ গুণ বৃদ্ধি
পেয়েছে।}

\emph{∴ সময় ৮ মাসের অর্ধেক = (৮ ÷ ২) = ৪ মাস হবে।}

\emph{বীজগণিতীয় পদ্ধতি:}

\emph{ধরি, সময় = t মাস}

\emph{প্রশ্নমতে, ২০০ × ৮ = ৪০০ × t}

\emph{⇒ ১৬০০ = ৪০০t}

\emph{∴ t = (১৬০০/৪০০) = ৪ মাস}
\end{quote}

\textbf{বার্ষিক শতকরা ১০ টাকা হার সুদে কোনো মুলধন কত বছর পর সুদে-আসলে দ্বিগুণ
হবে?}

\textbf{\hl{ক) ১০ বছর ✔}}

খ) ২০ বছর

গ) ৩০ বছর

ঘ) ৪০ বছর

\textbf{\hl{✔ সঠিক উত্তর: ক) ১০ বছর}}

\begin{quote}
\textbf{ব্যাখ্যা:} \emph{ঐকিক নিয়ম পদ্ধতি:}

\emph{মনে করি, আসল = ১০০ টাকা}

\emph{সুদে-আসলে দ্বিগুণ = (১০০ × ২) = ২০০ টাকা}

\emph{∴ সুদ = (২০০ − ১০০) = ১০০ টাকা}

\emph{১০ টাকা সুদ হয় = ১ বছরে}

\emph{∴ ১০০ টাকা সুদ হয় = (১০০/১০) = ১০ বছরে}

\emph{সূত্রভিত্তিক পদ্ধতি:}

\emph{সুদে-আসলে দ্বিগুণ হলে সুদ I = আসল P}

\emph{আমরা জানি, n = (I×১০০/Pr) = (P×১০০/P×১০) = ১০ বছর}
\end{quote}

\textbf{৬৬৬৬ টাকা বার্ষিক ১০\% সুদে কয় বছরে সুদে আসলে ১৩৩৩২ টাকা হবে?}

ক) ৮

খ) ২০

গ) ১৫

\textbf{\hl{ঘ) ১০ ✔}}

\textbf{\hl{✔ সঠিক উত্তর: ঘ) ১০}}

\begin{quote}
\textbf{ব্যাখ্যা:} \emph{ঐকিক নিয়ম পদ্ধতি:}

\emph{সুদ = (১৩৩৩২ − ৬৬৬৬) = ৬৬৬৬ টাকা}

\emph{১০০ টাকার ১ বছরের সুদ = ১০ টাকা}

\emph{∴ ৬৬৬৬ টাকার ১ বছরের সুদ = (১০×৬৬৬৬/১০০) = ৬৬৬.৬ টাকা}

\emph{৬৬৬.৬ টাকা সুদ হয় = ১ বছরে}

\emph{∴ ৬৬৬৬ টাকা সুদ হয় = (৬৬৬৬/৬৬৬.৬) = ১০ বছরে}

\emph{সূত্রভিত্তিক/সংক্ষিপ্ত পদ্ধতি:}

\emph{এখানে, I = ৬৬৬৬ টাকা, P = ৬৬৬৬ টাকা, r = ১০\%}

\emph{আমরা জানি, n = (I×১০০/Pr)}

\emph{⇒ n = (৬৬৬৬×১০০/৬৬৬৬×১০) = ১০ বছর}
\end{quote}

\textbf{১০\% হারে কত সময়ে ৫০০ টাকার সুদ ১০০ টাকা হবে?}

ক) ৩ বছরে

\textbf{\hl{খ) ২ বছরে ✔}}

গ) ৪ বছরে

ঘ) ২.৫ বছরে

\textbf{\hl{✔ সঠিক উত্তর: খ) ২ বছরে}}

\begin{quote}
\textbf{ব্যাখ্যা:} \emph{ঐকিক নিয়ম পদ্ধতি:}

\emph{১০০ টাকার ১ বছরের সুদ = ১০ টাকা}

\emph{∴ ৫০০ টাকার ১ বছরের সুদ = (১০×৫০০/১০০) = ৫০ টাকা}

\emph{৫০ টাকা সুদ হয় = ১ বছরে}

\emph{∴ ১০০ টাকা সুদ হয় = (১০০/৫০) = ২ বছরে}

\emph{সূত্রভিত্তিক/সংক্ষিপ্ত পদ্ধতি:}

\emph{এখানে, I = ১০০ টাকা, P = ৫০০ টাকা, r = ১০\%}

\emph{আমরা জানি, n = (I×১০০/Pr)}

\emph{⇒ n = (১০০×১০০/৫০০×১০) = ২ বছর}
\end{quote}

\textbf{৬\% সরল সুদে ৮০০ টাকার কত বছরের সুদ ৪৮০ টাকা হবে?}

ক) ৪ বছর

খ) ৫ বছর

গ) ৬ বছর

\textbf{\hl{ঘ) ১০ বছর ✔}}

\textbf{\hl{✔ সঠিক উত্তর: ঘ) ১০ বছর}}

\begin{quote}
\textbf{ব্যাখ্যা:} \emph{ঐকিক নিয়ম পদ্ধতি:}

\emph{১০০ টাকার ১ বছরের সুদ = ৬ টাকা}

\emph{∴ ৮০০ টাকার ১ বছরের সুদ = (৬×৮০০/১০০) = ৪৮ টাকা}

\emph{৪৮ টাকা সুদ হয় = ১ বছরে}

\emph{∴ ৪৮০ টাকা সুদ হয় = (৪৮০/৪৮) = ১০ বছরে}

\emph{সূত্রভিত্তিক/সংক্ষিপ্ত পদ্ধতি:}

\emph{এখানে, I = ৪৮০ টাকা, P = ৮০০ টাকা, r = ৬\%}

\emph{আমরা জানি, n = (I×১০০/Pr)}

\emph{⇒ n = (৪৮০×১০০/৮০০×৬) = ১০ বছর}
\end{quote}

\textbf{কিছু টাকা ৮\% লাভে ৩ বছরে লাভ আসলে ১৮৬০ টাকা হয়। কত বছর পরে তা
লাভ আসলে ২০৪০ টাকা হয়?}

ক) ২(১/২) বছরে

\textbf{\hl{খ) ৪(১/২) বছরে ✔}}

গ) ৩ বছরে

ঘ) ৪ বছরে

\textbf{\hl{✔ সঠিক উত্তর: খ) ৪(১/২) বছরে}}

\begin{quote}
\textbf{ব্যাখ্যা:} \emph{ঐকিক নিয়ম পদ্ধতি:}

\emph{ধরি, আসল ১০০ টাকা}

\emph{১০০ টাকার ৩ বছরের লাভ = (৮ × ৩) = ২৪ টাকা}

\emph{লাভ-আসলে হয় = (১০০ + ২৪) = ১২৪ টাকা}

\emph{লাভ-আসল ১২৪ টাকা হলে আসল = ১০০ টাকা}

\emph{∴ লাভ-আসল ১৮৬০ টাকা হলে আসল = (১৮৬০×১০০/১২৪) = ১৫০০ টাকা}

\emph{১৮৬০ টাকায় লাভ = (১৮৬০ − ১৫০০) = ৩৬০ টাকা}

\emph{২০৪০ টাকায় লাভ = (২০৪০ − ১৫০০) = ৫৪০ টাকা}

\emph{৩৬০ টাকা লাভ হয় = ৩ বছরে}

\emph{∴ ৫৪০ টাকা লাভ হয় = (৫৪০×৩/৩৬০) = (৯/২) = ৪(১/২) বছরে}

\emph{বীজগণিতীয়/সূত্রভিত্তিক পদ্ধতি:}

\emph{P = (A/১+nr) = (১৮৬০/১+৩×০.০৮) = (১৮৬০/১.২৪) = ১৫০০ টাকা}

\emph{২য় ক্ষেত্রে, I = ২০৪০ − ১৫০০ = ৫৪০ টাকা}

\emph{∴ n = (I/Pr) = (৫৪০/১৫০০×০.০৮) = (৫৪০/১২০) = ৪.৫ বা ৪(১/২) বছর}
\end{quote}

\textbf{কোন টাকা ৫ বছরে ৬\% হার সুদে-আসলে ১৩০০ টাকা হয়, কত বছরে ঐ টাকা
সুদে-আসলে ১৩৯০ টাকা হবে?}

ক) ৬ বছরে

\textbf{\hl{খ) ৬.৫ বছরে ✔}}

গ) ৭ বছরে

ঘ) ৭.৫ বছরে

\textbf{\hl{✔ সঠিক উত্তর: খ) ৬.৫ বছরে}}

\begin{quote}
\textbf{ব্যাখ্যা:} \emph{ঐকিক নিয়ম পদ্ধতি:}

\emph{১০০ টাকার ৫ বছরের সুদ = (৫ × ৬) = ৩০ টাকা}

\emph{১০০ টাকায় ৫ বছরে সুদে-আসলে হয় = (১০০ + ৩০) = ১৩০ টাকা}

\emph{সুদে-আসল ১৩০ টাকা হলে আসল = ১০০ টাকা}

\emph{∴ সুদে-আসল ১৩০০ টাকা হলে আসল = (১০০×১৩০০/১৩০) = ১০০০০ টাকা}

\emph{১০০০ টাকার ৫ বছরের সুদ = (১৩০০ − ১০০০০) = ৩০০ টাকা}

\emph{সুদে-আসলে ১৩৯০ টাকা হলে সুদ = (১৩৯০ − ১০০০০) = ৩ ৯০ টাকা}

\emph{১০০০০ টাকায় ৩০০ টাকা সুদ হয় = ৫ বছরে}

\emph{∴ ১০০০০ টাকায় ৩৯০ টাকা সুদ হয় = (৩৯০×৫/৩০০) = ৬.৫ বছরে}

\emph{বীজগণিতীয়/সূত্রভিত্তিক পদ্ধতি:}

\emph{P = (A/১+nr) = (১৩০০/১+৫×০.০৬) = (১৩০০/১.৩) = ১০০০০ টাকা}

\emph{২য় ক্ষেত্রে, I = ১৩৯০ − ১০০০০ = ৩ ৯০ টাকা}

\emph{∴ n = (I/Pr) = (৩৯০/১০০০০×০.০৬) = (৩৯০/৬০) = ৬.৫ বছর}
\end{quote}

\textbf{৬ (১/৪) \% সুদে কত সময়ে ৯৬ টাকার সুদ ১৮ টাকা হয়?}

ক) ২ বছরে

\textbf{\hl{খ) ৩ বছরে ✔}}

গ) ৪ বছরে

ঘ) ৬ বছরে

\textbf{\hl{✔ সঠিক উত্তর: খ) ৩ বছরে}}

\begin{quote}
\textbf{ব্যাখ্যা:} \emph{সূত্রভিত্তিক/সংক্ষিপ্ত পদ্ধতি:}

\emph{এখানে, I = ১৮ টাকা, P = ৯৬ টাকা, r = ৬ (১/৪) \% = (২৫/৪) \% =
(২৫/৪০০)}

\emph{আমরা জানি, n = (I/Pr)}

\emph{⇒ n = (১৮×৪০০/৯৬×২৫) = ৩ বছর}

\emph{ঐকিক নিয়ম পদ্ধতি:}

\emph{১০০ টাকার ১ বছরের সুদ = (২৫/৪) টাকা}

\emph{∴ ৯৬ টাকার ১ বছরের সুদ = (২৫×৯৬/৪×১০০) = ৬ টাকা}

\emph{৬ টাকা সুদ হয় = ১ বছরে}

\emph{∴ ১৮ টাকা সুদ হয় = (১৮/৬) = ৩ বছরে}
\end{quote}

\textbf{শতকরা ১ টাকা হার সুদে ১ টাকার সুদ ১ টাকা হবে কত বছরে?}

ক) ১ বছর

\textbf{\hl{খ) ১০০ বছর ✔}}

গ) ১০ বছর

ঘ) ১০০০ বছর

\textbf{\hl{✔ সঠিক উত্তর: খ) ১০০ বছর}}

\begin{quote}
\textbf{ব্যাখ্যা:} \emph{সূত্রভিত্তিক/সংক্ষিপ্ত পদ্ধতি:}

\emph{এখানে, I = ১ টাকা, P = ১ টাকা, r = ১\% = (১/১০০)}

\emph{আমরা জানি, n = (I/Pr)}

\emph{⇒ n = (১×১০০/১×১) = ১০০ বছর}

\emph{ঐকিক নিয়ম পদ্ধতি:}

\emph{১ টাকার ১ বছরের সুদ = (১/১০০) টাকা}

\emph{∴ (১/১০০) টাকা সুদ হয় = ১ বছরে}

\emph{∴ ১ টাকা সুদ হয় = (১ × ১০০) = ১০০ বছরে}
\end{quote}

\textbf{চক্রবৃদ্ধি মুনাফার ক্ষেত্রে নিচের কোন সূত্রটি সঠিক?}

ক) C = P(1 + n)\^{}r

খ) C = P(1 + r\^{}n)

গ) P = C(1 + r)\^{}n

\textbf{\hl{ঘ) C = P(1 + r)\^{}n ✔}}

\textbf{\hl{✔ সঠিক উত্তর: ঘ) C = P(1 + r)\^{}n}}

\begin{quote}
\textbf{ব্যাখ্যা:} \emph{চক্রবৃদ্ধি মুনাফার ক্ষেত্রে, চক্রবৃদ্ধি মূলধন C = P(1 +
r)\^{}n}

\emph{যেখানে,}

\emph{C = n বছর শেষে চক্রবৃদ্ধি মূলধন}

\emph{P = প্রাথমিক মূলধন/আসল}

\emph{r = মুনাফার/সুদের হার}

\emph{n = সময় (বছর)}
\end{quote}

\textbf{চক্রবৃদ্ধি মুনাফা নির্ণয়ের সূত্র কোনটি?}

\textbf{\hl{ক) p(1+r)\^{}n-p ✔}}

খ) p(1+r)\^{}n-pnl

গ) p(1+r)\^{}n×p

ঘ) p(1+r)\^{}n+p

\textbf{\hl{✔ সঠিক উত্তর: ক) p(1+r)\^{}n-p}}

\begin{quote}
\textbf{ব্যাখ্যা:} \emph{চক্রবৃদ্ধি মূলধন (C) = p(1 + r)\^{}n}

\emph{∴ চক্রবৃদ্ধি মুনাফা = C − p = p(1 + r)\^{}n − p}

\emph{এখানে, p = প্রারম্ভিক মূলধন, r = সুদের হার, n = সময়।}
\end{quote}

\textbf{বার্ষিক শতকরা ১০\% হারে ১০০০ টাকার ২ বছর পর সরল ও চক্রবৃদ্ধি মুনাফার
পার্থক্য কত?}

ক) ১১ টাকা

খ) ১১.৫ টাকা

গ) ১২ টাকা

\textbf{\hl{ঘ) ১০ টাকা ✔}}

\textbf{\hl{✔ সঠিক উত্তর: ঘ) ১০ টাকা}}

\begin{quote}
\textbf{ব্যাখ্যা:} \emph{সরল মুনাফা, I = Prn = ১০০০ × (১০/১০০) × ২ = ২০০
টাকা}

\emph{চক্রবৃদ্ধি মূলধন, C = P(1 + r)\^{}n = ১০০০ × ((১+(১০/১০০)))\^{}২ =
১০০০ × (((১১/১০)))\^{}২ = ১০০০ × (১২১/১০০) = ১২১০ টাকা}

\emph{∴ চক্রবৃদ্ধি মুনাফা, I\_c = C − P = (১২১০ − ১০০০) = ২১০ টাকা}

\emph{∴ সরল ও চক্রবৃদ্ধি মুনাফার পার্থক্য = (২১০ − ২০০) = ১০ টাকা}
\end{quote}

\textbf{৭\% হার সুদে ১,০০০ টাকার ৩ বছরের চক্রবৃদ্ধি সুদ কত হবে?}

ক) 125.04

খ) 115.03

\textbf{\hl{গ) 225.04 ✔}}

ঘ) 325.05

\textbf{\hl{✔ সঠিক উত্তর: গ) 225.04}}

\begin{quote}
\textbf{ব্যাখ্যা:} \emph{সূত্রভিত্তিক পদ্ধতি:}

\emph{এখানে, আসল P = ১,০০০ টাকা, সুদের হার r = ৭\%, সময় n = ৩ বছর}

\emph{আমরা জানি, চক্রবৃদ্ধি সুদ = P (১+(r/১০০))\^{}n − P}

\emph{= ১,০০০ (১+(৭/১০০))\^{}৩ − ১,০০০}

\emph{= ১,০০০ × (১.০৭)\^{}৩ − ১,০০০}

\emph{= (১,০০০ × ১.২২৫০৪৩) − ১,০০০}

\emph{= ১,২২৫.০৪৩ − ১,০০০ = ২২৫.০৪ টাকা (প্রায়)}

\emph{পর্যায়ক্রমিক চক্রবৃদ্ধি পদ্ধতি:}

\emph{১ম বছরের সুদ = ১,০০০ × ৭\% = ৭০ টাকা}

\emph{১ম বছর শেষে মূলধন = ১,০০০ + ৭০ = ১,০৭০ টাকা}

\emph{২য় বছরের সুদ = ১,০৭০ × ৭\% = ৭৪.৯ টাকা}

\emph{২য় বছর শেষে মূলধন = ১,০৭০ + ৭৪.৯ = ১,১৪৪.৯ টাকা}

\emph{৩য় বছরের সুদ = ১,১৪৪.৯ × ৭\% = ৮০.১৪৩ টাকা}

\emph{∴ মোট ৩ বছরের চক্রবৃদ্ধি সুদ = ৭০ + ৭৪.৯ + ৮০.১৪৩ = ২২৫.০৪৩ বা ২২৫.০৪
টাকা}
\end{quote}

\textbf{শতকরা ২০ টাকা মুনাফায় ৫০০ টাকায় ২ বছরের চক্রবৃদ্ধি মূলধন কত?}

ক) ৭৫০ টাকা

খ) ৭৪০ টাকা

\textbf{\hl{গ) ৭২০ টাকা ✔}}

ঘ) ৭০০ টাকা

\textbf{\hl{✔ সঠিক উত্তর: গ) ৭২০ টাকা}}

\begin{quote}
\textbf{ব্যাখ্যা:} \emph{সূত্রভিত্তিক পদ্ধতি:}

\emph{এখানে, আসল P = ৫০০ টাকা, মুনাফার হার r = ২০\%, সময় n = ২ বছর}

\emph{আমরা জানি, চক্রবৃদ্ধি মূলধন C = P (১+(r/১০০))\^{}n}

\emph{= ৫০০ × (১+(২০/১০০))\^{}২}

\emph{= ৫০০ × (১+(১/৫))\^{}২}

\emph{= ৫০০ × ((৬/৫))\^{}২}

\emph{= ৫০০ × (৩৬/২৫) = ৭২০ টাকা}

\emph{পর্যায়ক্রমিক হিসাব পদ্ধতি:}

\emph{১ম বছরের মুনাফা = ৫০০ × ২০\% = ১০০ টাকা}

\emph{১ম বছর শেষে মূলধন = ৫০০ + ১০০ = ৬০০ টাকা}

\emph{২য় বছরের মুনাফা = ৬০০ × ২০\% = ১২০ টাকা}

\emph{∴ ২ বছর শেষে চক্রবৃদ্ধি মূলধন = ৬০০ + ১২০ = ৭২০ টাকা}
\end{quote}

\textbf{ব্যাংকে কত টাকা রাখলে ৮ শতাংশ হারে মুনাফায় ২ বছরের চক্রবৃদ্ধি মূলধন
৫২০০ টাকা হবে?}

ক) ৪৫৫৮.১৬ টাকা

খ) ৪০৮৮.৭৫ টাকা

\textbf{\hl{গ) ৪৪৫৮.১৬ টাকা ✔}}

ঘ) ৫০৫০.৬০ টাকা

\textbf{\hl{✔ সঠিক উত্তর: গ) ৪৪৫৮.১৬ টাকা}}

\begin{quote}
\textbf{ব্যাখ্যা:} \emph{সূত্রভিত্তিক পদ্ধতি:}

\emph{এখানে, চক্রবৃদ্ধি মূলধন C = ৫২০০ টাকা, মুনাফার হার r = ৮\%, সময় n = ২
বছর}

\emph{আমরা জানি, C = P(১ + r)\^{}n}

\emph{⇒ ৫২০০ = P (১+(৮/১০০))\^{}২}

\emph{⇒ ৫২০০ = P ((১০৮/১০০))\^{}২}

\emph{⇒ ৫২০০ = P × (১১৬৬৪/১০০০০)}

\emph{⇒ P = (৫২০০×১০০০০/১১৬৬৪)}

\emph{∴ P = ৪৪৫৮.১৬ টাকা (প্রায়)}

\emph{ঐকিক নিয়ম পদ্ধতি:}

\emph{ধরি, মূলধন = ১০০ টাকা}

\emph{১০০ টাকার ২ বছরের চক্রবৃদ্ধি মূলধন C = ১০০ × (১+(৮/১০০))\^{}২ =
(১০৮×১০৮/১০০) = ১১৬.৬৪ টাকা}

\emph{চক্রবৃদ্ধি মূলধন ১১৬.৬৪ টাকা হলে আসল = ১০০ টাকা}

\emph{∴ চক্রবৃদ্ধি মূলধন ৫২০০ টাকা হলে আসল = (১০০×৫২০০/১১৬.৬৪) = ৪৪৫৮.১৬
টাকা}
\end{quote}

\textbf{4\% হার সুদে কত টাকার 2 বছরের সরল সুদ এবং চক্রবৃদ্ধি সুদের পার্থক্য 1
টাকা হবে?}

ক) 600

খ) 650

\textbf{\hl{গ) 625 ✔}}

ঘ) 620

\textbf{\hl{✔ সঠিক উত্তর: গ) 625}}

\begin{quote}
\textbf{ব্যাখ্যা:} \emph{ঐকিক নিয়ম পদ্ধতি:}

\emph{ধরি, মূলধন = ১০০ টাকা}

\emph{৪\% হারে ২ বছরের সরল সুদ, I = pnr = (১০০×২×৪/১০০) = ৮ টাকা}

\emph{আবার, ৪\% হারে ২ বছরের চক্রবৃদ্ধি মূলধন, C = P(১ + r)\^{}n = ১০০ ×
((১+(৪/১০০)))\^{}২ = ১০০ × (১০৪/১০০) × (১০৪/১০০) = ১০৮.১৬ টাকা}

\emph{∴ চক্রবৃদ্ধি সুদ = ১০৮.১৬ − ১০০ = ৮.১৬ টাকা}

\emph{২ বছরের সরল ও চক্রবৃদ্ধি সুদের পার্থক্য = (৮.১৬ − ৮) = ০.১৬ টাকা}

\emph{পার্থক্য ০.১৬ টাকা হলে মূলধন = ১০০ টাকা}

\emph{∴ পার্থক্য ১ টাকা হলে মূলধন = (১০০/০.১৬) = ৬২৫ টাকা}

\emph{সংক্ষিপ্ত পদ্ধতি (Shortcut Formula):}

\emph{২ বছরের জন্য সরল ও চক্রবৃদ্ধি সুদের পার্থক্যের সূত্র, D = P ×
(((r/১০০)))\^{}২}

\emph{⇒ ১ = P × (((৪/১০০)))\^{}২}

\emph{⇒ ১ = P × (((১/২৫)))\^{}২}

\emph{⇒ ১ = P × (১/৬২৫)}

\emph{∴ P = ৬২৫ টাকা}
\end{quote}

\textbf{বার্ষিক ১০\% হার সুদে কত বছরে ৫,০০০ টাকার চক্রবৃদ্ধি মূলধন ৬,৬৫৫ টাকা
হবে?}

ক) ১ বছর

খ) ২ বছর

গ) ২(১/২) বছর

\textbf{\hl{ঘ) ৩ বছর ✔}}

\textbf{\hl{✔ সঠিক উত্তর: ঘ) ৩ বছর}}

\begin{quote}
\textbf{ব্যাখ্যা:} \emph{সূত্রভিত্তিক পদ্ধতি:}

\emph{এখানে, চক্রবৃদ্ধি মূলধন C = ৬,৬৫৫ টাকা, মূলধন P = ৫,০০০ টাকা, সুদের হার
r = ১০\% = (১০/১০০) = ০.১, সময় = n বছর}

\emph{আমরা জানি, C = P(১ + r)\^{}n}

\emph{⇒ ৬,৬৫৫ = ৫,০০০ × (১ + ০.১)\^{}n}

\emph{⇒ (৬৬৫৫/৫০০০) = (১.১)\^{}n}

\emph{⇒ (১৩৩১/১০০০) = (১.১)\^{}n}

\emph{⇒ (১.১)\^{}৩ = (১.১)\^{}n}

\emph{∴ n = ৩ বছর}

\emph{পর্যায়ক্রমিক চক্রবৃদ্ধি হিসাব পদ্ধতি:}

\emph{১ম বছর শেষে মূলধন = ৫,০০০ + (৫,০০০ × ১০\%) = ৫,৫০০ টাকা}

\emph{২য় বছর শেষে মূলধন = ৫,৫০০ + (৫,৫০০ × ১০\%) = ৬,০৫০ টাকা}

\emph{৩য় বছর শেষে মূলধন = ৬,০৫০ + (৬,০৫০ × ১০\%) = ৬,৬৫৫ টাকা}

\emph{∴ নির্ণেয় সময় = ৩ বছর}
\end{quote}

\textbf{১০\% চক্রবৃদ্ধি সুদে ১০০ টাকার ২ বছরের সুদ-আসলে কত টাকা হয়?}

ক) ২০ টাকা

খ) ২১ টাকা

গ) ১২০ টাকা

\textbf{\hl{ঘ) ১২১ টাকা ✔}}

\textbf{\hl{✔ সঠিক উত্তর: ঘ) ১২১ টাকা}}

\begin{quote}
\textbf{ব্যাখ্যা:} \emph{সূত্রভিত্তিক পদ্ধতি:}

\emph{এখানে, আসল P = ১০০ টাকা, সুদের হার r = ১০\% = (১০/১০০) = ০.১, সময় n
= ২ বছর}

\emph{আমরা জানি, চক্রবৃদ্ধি মূলধন C = P(১ + r)\^{}n = ১০০ × (১ + ০.১)\^{}২ =
১০০ × (১.১)\^{}২}

\emph{= ১০০ × ১.২১ = ১২১ টাকা}

\emph{পর্যায়ক্রমিক চক্রবৃদ্ধি পদ্ধতি:}

\emph{১ম বছরের সুদ = ১০০ × ১০\% = ১০ টাকা}

\emph{১ম বছর শেষে মূলধন = ১০০ + ১০ = ১১০ টাকা}

\emph{২য় বছরের সুদ = ১১০ × ১০\% = ১১ টাকা}

\emph{∴ ২ বছর শেষে মোট সুদ-আসল = ১১০ + ১১ = ১২১ টাকা}
\end{quote}

\textbf{আপনি পোস্ট অফিসে ১০\% চক্রবৃদ্ধি সুদে ১৫০০০ টাকা জমা রাখলে ৩ বছর পরে
কত টাকা সুদ পাবেন?}

ক) ৪৫০০ টাকা

\textbf{\hl{খ) ৪৯৬৫ টাকা ✔}}

গ) ৪০০০ টাকা

ঘ) ১৯৫০০ টাকা

\textbf{\hl{✔ সঠিক উত্তর: খ) ৪৯৬৫ টাকা}}

\begin{quote}
\textbf{ব্যাখ্যা:} \emph{সূত্রভিত্তিক পদ্ধতি:}

\emph{এখানে, আসল P = ১৫০০০ টাকা, সুদের হার r = ১০\% = (১০/১০০) = ০.১, সময়
n = ৩ বছর}

\emph{আমরা জানি, চক্রবৃদ্ধি মূলধন C = P(১ + r)\^{}n = ১৫০০০ × (১ + ০.১)\^{}৩
= ১৫০০০ × (১.১)\^{}৩}

\emph{= ১৫০০০ × ১.৩৩১ = ১৯৯৬৫ টাকা}

\emph{∴ চক্রবৃদ্ধি সুদ = C − P = (১৯৯৬৫ − ১৫০০০) = ৪৯৬৫ টাকা}

\emph{পর্যায়ক্রমিক চক্রবৃদ্ধি পদ্ধতি:}

\emph{১ম বছরের সুদ = ১৫০০০ × ১০\% = ১৫০০ টাকা (১ম বছর শেষে আসল = ১৬৫০০
টাকা)}

\emph{২য় বছরের সুদ = ১৬৫০০ × ১০\% = ১৬৫০ টাকা (২য় বছর শেষে আসল = ১৮১৫০
টাকা)}

\emph{৩য় বছরের সুদ = ১৮১৫০ × ১০\% = ১৮১৫ টাকা}

\emph{∴ মোট ৩ বছরের চক্রবৃদ্ধি সুদ = ১৫০০ + ১৬৫০ + ১৮১৫ = ৪৯৬৫ টাকা}
\end{quote}

\textbf{এক ব্যক্তি বার্ষিক ১০\% চক্রবৃদ্ধি সুদে ৬০০ টাকা ব্যাংকে জমা রাখলেন। ২য়
বছর শেষে ঐ ব্যক্তি সুদসহ কত টাকা পাবেন?}

ক) ৭২৭

খ) ৭৩০

গ) ৭২৫

\textbf{\hl{ঘ) ৭২৬ ✔}}

\textbf{\hl{✔ সঠিক উত্তর: ঘ) ৭২৬}}

\begin{quote}
\textbf{ব্যাখ্যা:} \emph{সূত্রভিত্তিক পদ্ধতি:}

\emph{এখানে, আসল P = ৬০০ টাকা, সুদের হার r = ১০\% = (১০/১০০) = ০.১, সময় n
= ২ বছর}

\emph{আমরা জানি, চক্রবৃদ্ধি মূলধন C = P(১ + r)\^{}n = ৬০০ × (১ + ০.১)\^{}২ =
৬০০ × (১.১)\^{}২}

\emph{= ৬০০ × ১.২১ = ৭২৬ টাকা}

\emph{পর্যায়ক্রমিক চক্রবৃদ্ধি পদ্ধতি:}

\emph{১ম বছরের সুদ = ৬০০ × ১০\% = ৬০ টাকা}

\emph{১ম বছর শেষে মূলধন = ৬০০ + ৬০ = ৬৬০ টাকা}

\emph{২য় বছরের সুদ = ৬৬০ × ১০\% = ৬৬ টাকা}

\emph{∴ ২য় বছর শেষে সুদসহ মোট পাবেন = ৬৬০ + ৬৬ = ৭২৬ টাকা}
\end{quote}

\section{50. লাভ-ক্ষতির ১০০\% ম্যাজিক: শর্টকাট নিয়মে পাটিগণিতের
সমাধান}\label{ux9b2ux9ad-ux995ux9b7ux9a4ux9b0-ux9e7ux9e6ux9e6-ux9aeux9afux99cux995-ux9b6ux9b0ux99fux995ux99f-ux9a8ux9dfux9ae-ux9aaux99fux997ux9a3ux9a4ux9b0-ux9b8ux9aeux9a7ux9a8}

লাভ-ক্ষতির শত শত জটিল অংককে অন্ধের মতো সূত্র মুখস্থ না করে কেবল
\textquotesingle{}১০০\%\textquotesingle-এর ধারণা ব্যবহার করে খুব সহজেই এক
লাইনে সমাধান করা সম্ভব। লাভ-ক্ষতির অংকে সবসময় ক্রয়মূল্যকে ১০০\% ধরতে হয়। কোনো
দ্রব্যের ক্রয়মূল্যের সাথে লাভ যুক্ত হলে তা ১০০\%-এর বেশি হয় এবং ক্ষতি হলে তা
১০০\%-এর কম হয়। নিচে ধাপে ধাপে গুরুত্বপূর্ণ ধারণাগুলো উদাহরণসহ আলোচনা করা হলো:

\begin{itemize}
\item
  ক্রয়মূল্য নির্ণয় সংক্রান্ত সমস্যা
\item
  বিক্রয়মূল্য নির্ণয় সংক্রান্ত সমস্যা
\item
  অধিক লাভ বা ক্ষতি সংক্রান্ত তুলনামূলক সমস্যা‬‬‬
\item
  একাধিকবার লাভ বা ছাড় দেওয়া সংক্রান্ত সমস্যা
\item
  লাভ-ক্ষতির হার নির্ণয় সংক্রান্ত সমস্যা
\item
  একাধিকবার লাভ বা ধারাবাহিক লাভ সংক্রান্ত সমস্যা
\item
  বিক্রয় দর নির্ণয় সংক্রান্ত সমস্যা
\end{itemize}

\subsubsection{১. ক্রয়মূল্য নির্ণয় সংক্রান্ত
সমস্যা}\label{ux995ux9b0ux9dfux9aeux9b2ux9af-ux9a8ux9b0ux9a3ux9df-ux9b8ux995ux9b0ux9a8ux9a4-ux9b8ux9aeux9b8ux9af}

লাভ-ক্ষতির অঙ্কে বিক্রয়মূল্য এবং লাভ বা ক্ষতির শতকরা হার দেওয়া থাকলে ক্রয়মূল্য
নির্ণয় করতে হয়। এক্ষেত্রে যে সংখ্যাটি দেওয়া থাকে তা কত শতাংশের মান তা বের
করে ১০০\% বা ক্রয়মূল্যের মান বের করতে হয়।

\textbf{উদাহরণ ১: একটি কলম ২৭০ টাকায় বিক্রয় করায় ১০\% ক্ষতি হয়; কলমটির
ক্রয়মূল্য কত?}

\begin{itemize}
\item
  \textbf{সমাধান:} যেহেতু ১০\% ক্ষতি হয়েছে, তাই ২৭০ টাকা হলো (১০০−১০)\% বা
  ৯০\%-এর মান। ক্রয়মূল্য (১০০\%) বের করতে হলে পার্সেন্টেজটিকে উল্টিয়ে গুণ করতে
  হয়: ২৭০×(১০০/৯০)=৩০০ টাকা
\end{itemize}

\textbf{উদাহরণ ২: এক ব্যক্তি কোনো পণ্য ২০\% লাভে ৩৬০০ টাকায় বিক্রি করলে
পণ্যটির ক্রয়মূল্য কত?}

\begin{itemize}
\item
  \textbf{সমাধান:} ২০\% লাভে বিক্রয়মূল্য ১২০\% হওয়ায় ৩৬০০ টাকা হলো ১২০\%-এর
  মান। সুতরাং ক্রয়মূল্য: ৩৬০০×(১০০/১২০)=৩০০০ টাকা
\end{itemize}

\subsubsection{২. বিক্রয়মূল্য নির্ণয় সংক্রান্ত
সমস্যা}\label{ux9acux995ux9b0ux9dfux9aeux9b2ux9af-ux9a8ux9b0ux9a3ux9df-ux9b8ux995ux9b0ux9a8ux9a4-ux9b8ux9aeux9b8ux9af}

ক্রয়মূল্য জানা থাকলে নির্দিষ্ট লাভে বা ক্ষতিতে পণ্যটির বিক্রয়মূল্য সরাসরি বের করা
যায়।

\textbf{উদাহরণ ১: ১৫০ টাকায় একটি জিনিস ক্রয় করে কত দামে বিক্রয় করলে ৩০\%
লাভ হবে?}

\begin{itemize}
\item
  \textbf{সমাধান:} সরাসরি ১৫০-এর ১৩০\% এর মান বের করতে হবে, কারণ ক্রয়মূল্য
  ১০০\% এবং ৩০\% লাভে তা ১৩০\% হয়: (১৫০×১৩০/১০০)=১৯৫ টাকা
\end{itemize}

\textbf{উদাহরণ ২: ২০০ টাকায় একটি বই কিনে ১৫\% ক্ষতিতে বিক্রয় করলে বিক্রয়মূল্য
কত হবে?}

\begin{itemize}
\item
  \textbf{সমাধান:} ১৫\% ক্ষতিতে বিক্রয়মূল্য হবে ১০০ - ১৫ = ৮৫\%। সুতরাং:
  (২০০×৮৫/১০০)=১৭০ টাকা
\end{itemize}

\subsubsection{৩. অধিক লাভ বা ক্ষতি সংক্রান্ত তুলনামূলক
সমস্যা}\label{ux985ux9a7ux995-ux9b2ux9ad-ux9ac-ux995ux9b7ux9a4-ux9b8ux995ux9b0ux9a8ux9a4-ux9a4ux9b2ux9a8ux9aeux9b2ux995-ux9b8ux9aeux9b8ux9af}

যখন কোনো পণ্যের বিক্রয়মূল্য কম বা বেশি হওয়ার কারণে লাভ বা ক্ষতির পরিমাণে
তারতম্য ঘটে, তখন শতকরা হারের পার্থক্য বের করে সহজে সমাধান করা যায়।

\textbf{উদাহরণ ১: একটি দ্রব্য ১০\% ক্ষতিতে বিক্রয় করা হলো। বিক্রয়মূল্য ৪৫ টাকা
বেশি হলে ৫\% লাভ হতো। দ্রব্যটির ক্রয়মূল্য কত?}

\begin{itemize}
\item
  \textbf{সমাধান:} ১০\% ক্ষতি পুষিয়ে ৫\% লাভ করতে হলে মোট ব্যবধান বা বৃদ্ধি
  করতে হবে ১০\%+৫\%=১৫\% । এই ১৫\%-এর মানই হলো ৪৫ টাকা। সুতরাং ক্রয়মূল্য:
  ৪৫×(১০০/১৫)=৩০০ টাকা
\end{itemize}

\textbf{উদাহরণ ২: একটি ঘড়ি ১৫\% লাভে বিক্রয় করা হলো। বিক্রয়মূল্য ৩০০ টাকা কম
হলে ১০\% ক্ষতি হতো। ঘড়িটির ক্রয়মূল্য কত?}

\begin{itemize}
\item
  \textbf{সমাধান:} ১৫\% লাভ থেকে ১০\% ক্ষতিতে নামতে হলে মোট শতকরা ব্যবধান
  হবে ১৫\%+১০\%=২৫\% । এই ২৫\%-এর মান ৩০০ টাকা। সুতরাং ক্রয়মূল্য:
  ৩০০×(১০০/২৫)=১২০০ টাকা
\end{itemize}

\subsubsection{৪. একাধিকবার লাভ বা ছাড় দেওয়া সংক্রান্ত
সমস্যা}\label{ux98fux995ux9a7ux995ux9acux9b0-ux9b2ux9ad-ux9ac-ux99bux9dc-ux9a6ux993ux9df-ux9b8ux995ux9b0ux9a8ux9a4-ux9b8ux9aeux9b8ux9af}

পণ্যের লিখিত মূল্য বা ধার্য মূল্যের উপর ছাড় দিয়ে বিক্রয় করার পর লাভ বা ক্ষতির
হিসাব এই নিয়মে করতে হয়।

\textbf{উদাহরণ ১: একটি গণিতের বই কিনতে ১৫\% কমিশন দেয়। বইটির প্রকৃত দাম
(কভারে লিখিত দাম) ১২০ টাকা হলে, বইটি কিনতে কত টাকা লাগবে?}

\begin{itemize}
\item
  \textbf{সমাধান:} কভারের দাম থেকে ১৫\% কমিশন দিলে দিতে হবে ৮৫\%। সুতরাং
  এক লাইনে উত্তর: ১২০×(৮৫/১০০)=১০২ টাকা
\end{itemize}

\textbf{উদাহরণ ২: একটি পণ্যের লিখিত মূল্যের উপর ১০\% ছাড়ে বিক্রি করার পর
খুচরা বিক্রেতা ঐ পণ্যটি ১৫\% ক্ষতিতে ৩০৬ টাকায় বিক্রি করলে পণ্যটির লিখিত মূল্য
কত ছিল?}

\begin{itemize}
\item
  \textbf{সমাধান:} ৩০৬ টাকা হলো ১৫\% ক্ষতিতে অর্থাৎ ৮৫\%-এর মান, এবং ১০\%
  ছাড়ে তা লিখিত মূল্যের ৯০\%-এর সমান। এক লাইনে হিসাবটি হবে:
  ৩০৬×(১০০/৮৫)×(১০০/৯০)=৪০০ টাকা
\end{itemize}

\subsubsection{৫. লাভ-ক্ষতির হার নির্ণয় সংক্রান্ত
সমস্যা}\label{ux9b2ux9ad-ux995ux9b7ux9a4ux9b0-ux9b9ux9b0-ux9a8ux9b0ux9a3ux9df-ux9b8ux995ux9b0ux9a8ux9a4-ux9b8ux9aeux9b8ux9af}

ক্রয়মূল্য এবং বিক্রয়মূল্য দেওয়া থাকলে শতকরা কত লাভ বা ক্ষতি হয়েছে তা নির্ণয় করা।

\textbf{উদাহরণ ১: একটি কলম ৫০ টাকায় ক্রয় করে ৫৬ টাকায় বিক্রয় করা হলো। এতে
শতকরা কত লাভ হলো?}

\begin{itemize}
\item
  \textbf{সমাধান:} ৫০ টাকায় লাভ হয় ৫৬−৫০=৬ টাকা। ১০০ টাকায় কত লাভ তা বের
  করলেই শতকরা হার পাওয়া যাবে: (৬/৫০)×১০০=১২\%
\end{itemize}

\textbf{উদাহরণ ২: ২০০ টাকায় একটি বল কিনে ১৫০ টাকায় বিক্রি করলে শতকরা কত
ক্ষতি হয়?}

\begin{itemize}
\item
  \textbf{সমাধান:} ক্ষতি হয় ২০০−১৫০=৫০ টাকা। সুতরাং ক্ষতির হার:
  (৫০/২০০)×১০০=২৫\%
\end{itemize}

\subsubsection{৬. একাধিকবার লাভ বা ধারাবাহিক লাভ সংক্রান্ত
সমস্যা}\label{ux98fux995ux9a7ux995ux9acux9b0-ux9b2ux9ad-ux9ac-ux9a7ux9b0ux9acux9b9ux995-ux9b2ux9ad-ux9b8ux995ux9b0ux9a8ux9a4-ux9b8ux9aeux9b8ux9af}

কোনো পণ্য হাতবদল হয়ে বা পাইকারি ও খুচরা বিক্রেতা উভয়ে লাভ রেখে বিক্রি করলে
ধারাবাহিক শতকরা হার প্রয়োগ করতে হয়।

\textbf{উদাহরণ ১: একজন পাইকারি ব্যবসায়ী ও খুচরা বিক্রেতা উভয়ে ২০\% লাভে পণ্য
বিক্রি করেন। একটি শার্টের প্রাথমিক মূল্য ৪০০ টাকা হলে, ক্রেতা কত টাকায় তা
কিনতে পারবেন?}

\begin{itemize}
\item
  \textbf{সমাধান:} দুবার ২০\% লাভ হওয়ায় ধারাবাহিকভাব ১২০\% দিয়ে গুণ করতে
  হবে: ৪০০×(১২০/১০০)×(১২০/১০০)=৫৭৬ টাকা
\end{itemize}

\textbf{উদাহরণ ২: উৎপাদক ১০\% লাভে এবং পাইকারি বিক্রেতা ১৫\% লাভে একটি
জিনিস বিক্রি করলে, ১০০০ টাকা উৎপাদন খরচের জিনিসটির চূড়ান্ত বিক্রয়মূল্য কত
হবে?}

\begin{itemize}
\item
  \textbf{সমাধান:} পর্যায়ক্রমে ১১০\% এবং ১১৫\% দিয়ে গুণ করতে হবে:
  ১০০০×(১১০/১০০)×(১১৫/১০০)=১২৬৫ টাকা
\end{itemize}

\subsubsection{৭. বিক্রয় দর নির্ণয় সংক্রান্ত
সমস্যা}\label{ux9acux995ux9b0ux9df-ux9a6ux9b0-ux9a8ux9b0ux9a3ux9df-ux9b8ux995ux9b0ux9a8ux9a4-ux9b8ux9aeux9b8ux9af}

একগুচ্ছ বা নির্দিষ্ট পরিমাণের পণ্য কিনে একক বা হালি/ডজন হিসেবে নির্দিষ্ট লাভে
বিক্রয় দর কত হবে তা বের করার পদ্ধতি।

\textbf{উদাহরণ ১: এক ডজন আম ৬০ টাকায় ক্রয় করে হালি কত টাকায় বিক্রয় করলে
১০\% লাভ হবে?}

\begin{itemize}
\item
  \textbf{সমাধান:} ১ ডজন = ১২টি এবং ১ হালি = ৪টি। প্রথমে ১০\% লাভে ১২টি
  আমের মোট বিক্রয়মূল্য বের করতে হবে: (৬০×১১০/১০০)=৬৬ টাকা
\item
  যেহেতু ১২টি আমের বিক্রয়মূল্য ৬৬ টাকা, সেহেতু ১ হালি বা ৪টি আমের বিক্রয়মূল্য
  হবে: (৬৬/৩)=২২ টাকা
\end{itemize}

\textbf{উদাহরণ ২: টাকায় এক ডজন কলা বিক্রি করায় ২০\% ক্ষতি হয়। ৬০\% লাভ
করতে হলে টাকায় কতটি কলা বিক্রি করতে হবে?}

\begin{itemize}
\item
  \textbf{সমাধান:} ২০\% ক্ষতিতে বিক্রয়মূল্য ৮০ টাকা এবং ৬০\% লাভে বিক্রয়মূল্য
  ১৬০ টাকা। পূর্বের তুলনায় নতুন বিক্রয়মূল্য দ্বিগুণ ( ১৬০/৮০=২ টাকা) হবে। অর্থাৎ
  যে কলা আগে ১ টাকায় ১ ডজন (১২টি) পাওয়া যেত, তা ২ টাকায় ১ ডজন বিক্রি করতে
  হবে। সুতরাং ১ টাকায় বিক্রি করতে হবে: (১২/২)=৬ টি কলা
\end{itemize}

\section{51. লাভ ও ক্ষতি MCQ প্রশ্ন উত্তর
অনুশীলন}\label{ux9b2ux9ad-ux993-ux995ux9b7ux9a4-mcq-ux9aaux9b0ux9b6ux9a8-ux989ux9a4ux9a4ux9b0-ux985ux9a8ux9b6ux9b2ux9a8}

\textbf{একটি দ্রব্য ১৮০ টাকায় বিক্রয় করায় ১০\% ক্ষতি হলো। দ্রব্যটির ক্রয়মূল্য---}

\textbf{\hl{ক) ২০০ টাকা ✔}}

খ) ২১০ টাকা

গ) ১৬২ টাকা

ঘ) ১৯৮ টাকা

\textbf{\hl{✔ সঠিক উত্তর: ক) ২০০ টাকা}}

\begin{quote}
\textbf{ব্যাখ্যা:} \emph{ঐকিক নিয়ম পদ্ধতি:}

\emph{১০\% ক্ষতিতে, বিক্রয়মূল্য ৯০ টাকা হলে ক্রয়মূল্য = ১০০ টাকা}

\emph{∴ বিক্রয়মূল্য ১ টাকা হলে ক্রয়মূল্য = (১০০/৯০) টাকা}

\emph{∴ বিক্রয়মূল্য ১৮০ টাকা হলে ক্রয়মূল্য = (১০০×১৮০/৯০) = ২০০ টাকা}

\emph{সূত্রভিত্তিক পদ্ধতি:}

\emph{ক্রয়মূল্য = (বিক্রয়মূল্য×১০০/১০০−ক্ষতির হার)}

\emph{= (১৮০×১০০/১০০−১০)}

\emph{= (১৮০×১০০/৯০) = ২০০ টাকা}
\end{quote}

\textbf{একজন বিক্রেতা ১১০ টাকা কেজি দরের কিছু চা-এর সাথে ১০০ টাকা কেজি
দরের দ্বিগুণ পরিমাণ চা মিশিয়ে ১২০ টাকা কেজি দরে বিক্রি করে মোট ২০০০ টাকা
লাভ করে। বিক্রেতা দ্বিতীয় প্রকারের কত কেজি চা ক্রয় করেছিল?}

ক) ১০০

\textbf{\hl{খ) ৮০ ✔}}

গ) ৫০

ঘ) ৬০

ঙ) ৭৫

\textbf{\hl{✔ সঠিক উত্তর: খ) ৮০}}

\begin{quote}
\textbf{ব্যাখ্যা:} \emph{ধরি, ১ম প্রকারের চা ক্রয় করে x কেজি}

\emph{২য় প্রকারের চা ক্রয় করে 2x কেজি}

\emph{মোট ক্রয়কৃত চা = 3x কেজি}

\emph{3x কেজি চা এর বিক্রয়মূল্য = 120×3x=360x টাকা}

\emph{3x কেজি চা এর ক্রয়মূল্য = (110×x)+(100×2x)=310x টাকা}

\emph{প্রশ্নমতে, 360x−310 x = 2000}

\emph{বা, 50x=2000}

\emph{বা, x=40}

\emph{∴ ২য় প্রকারের চা ক্রয় করে = 2x কেজি = 2×40=80 কেজি}
\end{quote}

\textbf{তানজিম ৫ টি কলম, ২০ টাকায় বিক্রয় করায় তার ৫ টাকা ক্ষতি হলো। তার
ক্রয়মূল্য কত?}

ক) ২০ টাকা

\textbf{\hl{খ) ২৫ টাকা ✔}}

গ) ১৫ টাকা

ঘ) ১০ টাকা

\textbf{\hl{✔ সঠিক উত্তর: খ) ২৫ টাকা}}

\begin{quote}
\textbf{ব্যাখ্যা:} \emph{ক্রয়মূল্য = বিক্রয়মূল্য + ক্ষতি = 20+5=25 টাকা}
\end{quote}

\textbf{একটি কলম ২৫ টাকায় বিক্রয় করলে যে ক্ষতি হয়, ৩৫ টাকায় বিক্রয় করলে তত
লাভ হয়। উহার ক্রয়মূল্য কত?}

ক) ২৮ টাকা

\textbf{\hl{খ) ৩০ টাকা ✔}}

গ) ৩২ টাকা

ঘ) ৩৪ টাকা

\textbf{\hl{✔ সঠিক উত্তর: খ) ৩০ টাকা}}

\begin{quote}
\textbf{ব্যাখ্যা:} \emph{যদি এরকম অংক আসে অর্থাৎ লাভ ও ক্ষতি সমান দেখানো হয়
তাহলে দুটি বিক্রয়মূল্য যোগ করে ২ দিয়ে ভাগ করলেই ক্রয়মূল্য বের হয়ে আসবে।}

\emph{অর্থাৎ (25+35/2)=(60/2)=30 টাকা}
\end{quote}

\textbf{একটি গরু ৪৫০ টাকায় বিক্রয় করায় তার ক্রয়মূল্যের (1/8) অংশ লাভ হল।
গরুটির ক্রয়মূল্য কত?}

ক) ৫০০ টাকা

খ) ৬০০ টাকা

গ) ৩০০ টাকা

\textbf{\hl{ঘ) ৪০০ টাকা ✔}}

\textbf{\hl{✔ সঠিক উত্তর: ঘ) ৪০০ টাকা}}

\begin{quote}
\textbf{ব্যাখ্যা:} \emph{এখানে, ক্রয়মূল্য ১ টাকা হলে বিক্রয়মূল্য 1+(1/8)=(9/8)
টাকা}

\emph{বিক্রয়মূল্য (9/8) টাকা হলে ক্রয়মূল্য = ১ টাকা}

\emph{বিক্রয়মূল্য ১ টাকা হলে ক্রয়মূল্য = (8/9) টাকা}

\emph{বিক্রয়মূল্য ৪৫০ টাকা হলে ক্রয়মূল্য = (8×450/9)=400 টাকা}
\end{quote}

\textbf{একটি পুতুল ২৫\% লাভে ৩৭৫ টাকায় বিক্রয় করা হলো। পুতুলটির ক্রয়মূল্য কত?}

ক) ২৭৫ টাকা

\textbf{\hl{খ) ৩০০ টাকা ✔}}

গ) ৩৮০ টাকা

ঘ) ৪০০ টাকা

\textbf{\hl{✔ সঠিক উত্তর: খ) ৩০০ টাকা}}

\begin{quote}
\textbf{ব্যাখ্যা:} \emph{২৫\% লাভে বিক্রয়মূল্য = 100+25=125 টাকা}

\emph{বিক্রয়মূল্য ১২৫ টাকা হলে ক্রয়মূল্য = ১০০ টাকা}

\emph{বিক্রয়মূল্য ১ টাকা হলে ক্রয়মূল্য = (100/125) টাকা}

\emph{বিক্রয়মূল্য ৩৭৫ টাকা হলে ক্রয়মূল্য = (100×375/125)=300 টাকা}
\end{quote}

\textbf{একটি দ্রব্য ১০\% ক্ষতিতে বিক্রয় করা হলো। বিক্রয়মূল্য ৮০ টাকা বেশি হলে
১০\% লাভ হতো। দ্রব্যটির ক্রয়মূল্য কত?}

ক) ২৮০ টাকা

খ) ৩০০ টাকা

গ) ৩৮০ টাকা

\textbf{\hl{ঘ) ৪০০ টাকা ✔}}

\textbf{\hl{✔ সঠিক উত্তর: ঘ) ৪০০ টাকা}}

\begin{quote}
\textbf{ব্যাখ্যা:} \emph{১০\% ক্ষতিতে বিক্রয়মূল্য = ১০০−১০ টাকা = ৯০ টাকা}

\emph{১০\% লাভে বিক্রয়মূল্য = ১০০+১০ টাকা = ১১০ টাকা}

\emph{বিক্রয় পার্থক্য = ১১০−৯০=২০ টাকা}

\emph{বিক্রয় পার্থক্য ২০ টাকা হলে ক্রয়মূল্য = ১০০ টাকা}

\emph{বিক্রয় পার্থক্য ১ টাকা হলে ক্রয়মূল্য = (১০০/২০)=৫ টাকা}

\emph{বিক্রয় পার্থক্য ৮০ টাকা হলে ক্রয়মূল্য = ৫×৮০=৪০০ টাকা}
\end{quote}

\textbf{একজন বিক্রেতা একটি শার্ট ৪৮০ টাকায় বিক্রয় করে ২০\% লাভ করল। শার্টটির
ক্রয়মূল্য কত?}

ক) ৩৮৪ টাকা

\textbf{\hl{খ) ৪০০ টাকা ✔}}

গ) ৪২০ টাকা

ঘ) ৫৭৬ টাকা

\textbf{\hl{✔ সঠিক উত্তর: খ) ৪০০ টাকা}}

\begin{quote}
\textbf{ব্যাখ্যা:} \emph{২০\% লাভে বিক্রয়মূল্য = ১২০ টাকা}

\emph{বিক্রয়মূল্য ১২০ টাকা হলে ক্রয়মূল্য = ১০০ টাকা}

\emph{বিক্রয়মূল্য ১ টাকা হলে ক্রয়মূল্য = (১০০/১২০) টাকা}

\emph{বিক্রয়মূল্য ৪৮০ টাকা হলে ক্রয়মূল্য = (১০০×৪৮০/১২০)=৪০০ টাকা}
\end{quote}

\textbf{একটি জিনিস ৬০ টাকা বিক্রি করলে ২০\% লাভ হয়। জিনিসটির ক্রয়মূল্য কত?}

ক) ৪০ টাকা

\textbf{\hl{খ) ৫০ টাকা ✔}}

গ) ৬০ টাকা

ঘ) ৮০ টাকা

\textbf{\hl{✔ সঠিক উত্তর: খ) ৫০ টাকা}}

\begin{quote}
\textbf{ব্যাখ্যা:} \emph{২০\% লাভে বিক্রয়মূল্য = ১২০ টাকা}

\emph{বিক্রয়মূল্য ১২০ টাকা হলে ক্রয়মূল্য = ১০০ টাকা}

\emph{বিক্রয়মূল্য ১ টাকা হলে ক্রয়মূল্য = (১০০/১২০) টাকা}

\emph{বিক্রয়মূল্য ৬০ টাকা হলে ক্রয়মূল্য = (১০০×৬০/১২০)=৫০ টাকা}
\end{quote}

\textbf{এক দোকানদার ১২.৫\% ক্ষতিতে একটি দ্রব্য বিক্রি করেন। যে মূল্য দিয়ে তিনি
দ্রব্যটি বিক্রি করলেন তার চাইতে ৩০ টাকা বেশি মূল্যে বিক্রি করলে ক্রয়মূল্যের ওপর
২৫\% লাভ হতো। দ্রব্যটির ক্রয়মূল্য কত?}

ক) ৭৫

\textbf{\hl{খ) ৮০ ✔}}

গ) ৯০

ঘ) ৯৫

\textbf{\hl{✔ সঠিক উত্তর: খ) ৮০}}

\begin{quote}
\textbf{ব্যাখ্যা:} \emph{১২.৫\% ক্ষতিতে বিক্রয়মূল্য = ১০০−১২.৫ টাকা = ৮৭.৫
টাকা}

\emph{২৫\% লাভে বিক্রয়মূল্য = ১০০+২৫ টাকা = ১২৫ টাকা}

\emph{বিক্রয় পার্থক্য = ১২৫−৮৭.৫=৩৭.৫ টাকা}

\emph{বিক্রয় পার্থক্য ৩৭.৫ টাকা হলে ক্রয়মূল্য = ১০০ টাকা}

\emph{বিক্রয় পার্থক্য ১ টাকা হলে ক্রয়মূল্য = (১০০/৩৭.৫) টাকা}

\emph{বিক্রয় পার্থক্য ৩০ টাকা হলে ক্রয়মূল্য = (১০০×৩০/৩৭.৫)=৮০ টাকা}
\end{quote}

\textbf{একটি গাড়ি ২৭০ টাকায় বিক্রি করাতে ক্রয়মূল্যের ওপর ১০\% ক্ষতি হলে
ক্রয়মূল্য কত?}

ক) ২৯৭ টাকা

খ) ২৪৩ টাকা

গ) ২৭৩ টাকা

\textbf{\hl{ঘ) ৩০০ টাকা ✔}}

\textbf{\hl{✔ সঠিক উত্তর: ঘ) ৩০০ টাকা}}

\begin{quote}
\textbf{ব্যাখ্যা:} \emph{১০\% ক্ষতিতে বিক্রয়মূল্য = ৯০ টাকা}

\emph{বিক্রয়মূল্য ৯০ টাকা হলে ক্রয়মূল্য = ১০০ টাকা}

\emph{বিক্রয়মূল্য ১ টাকা হলে ক্রয়মূল্য = (১০০/৯০) টাকা}

\emph{বিক্রয়মূল্য ২৭০ টাকা হলে ক্রয়মূল্য = (১০০×২৭০/৯০)=৩০০ টাকা}
\end{quote}

\textbf{জাহিদ ১০\% কমিশনে একটি বই ক্রয় করে দোকানীকে ১৮০ টাকা দিল। বইটির
প্রকৃত মূল্য কত টাকা?}

ক) ১০০

খ) ৫০

গ) ১৫০

\textbf{\hl{ঘ) ২০০ ✔}}

\textbf{\hl{✔ সঠিক উত্তর: ঘ) ২০০}}

\begin{quote}
\textbf{ব্যাখ্যা:} \emph{১০\% কমিশনে ক্রয়মূল্য = ৯০ টাকা}

\emph{ক্রয়মূল্য ৯০ টাকা হলে প্রকৃতমূল্য ১০০ টাকা}

\emph{বিক্রয়মূল্য ১৮০ টাকা হলে প্রকৃতমূল্য = (১০০×১৮০/৯০)=২০০ টাকা}
\end{quote}

\textbf{একটি পণ্য ৫\% ক্ষতিতে ২৩৭৫ টাকায় বিক্রয় করা হলে ক্রয়মূল্য কত?}

\textbf{\hl{ক) ২৫০০ টাকা ✔}}

খ) ২৪৭৫ টাকা

গ) ২৪৫০ টাকা

ঘ) ২২৭৫ টাকা

\textbf{\hl{✔ সঠিক উত্তর: ক) ২৫০০ টাকা}}

\begin{quote}
\textbf{ব্যাখ্যা:} \emph{৫\% ক্ষতিতে বিক্রয়মূল্য = ১০০−৫=৯৫ টাকা}

\emph{বিক্রয়মূল্য ৯৫ টাকা হলে ক্রয়মূল্য = ১০০ টাকা}

\emph{বিক্রয়মূল্য ১ টাকা হলে ক্রয়মূল্য = (১০০/৯৫) টাকা}

\emph{বিক্রয়মূল্য ২৩৭৫ টাকা হলে ক্রয়মূল্য = (১০০×২৩৭৫/৯৫)=২৫০০ টাকা}
\end{quote}

\textbf{একটি কলম ও পেন্সিলের একত্রিত মূল্য ৪২ টাকা। যদি পেন্সিলের মূল্য কলমের
মূল্যের (১/৬) অংশ হয়,তবে কলমের মূল্য কত?}

ক) ৩০

খ) ৩২

\textbf{\hl{গ) ৩৬ ✔}}

ঘ) কোনোটিই নয়

\textbf{\hl{✔ সঠিক উত্তর: গ) ৩৬}}

\begin{quote}
\textbf{ব্যাখ্যা:} \emph{এখানে, অনুপাতের যোগফল = ১+৬=৭}

\emph{৪২ কে ৭ দিয়ে ভাগ করলে পেন্সিলের মূল্য পাওয়া যাবে।}

\emph{পেন্সিলের মূল্য = (৪২/৭)=৬ টাকা}

\emph{∴ কলমের মূল্য পেন্সিলের ৬ গুণ, অর্থাৎ ৬×৬=৩৬ টাকা।}
\end{quote}

\textbf{একটি শপিং মলে সকল পণ্যের উপর ২৫\% ডিসকাউন্ট ঘোষণা করা হলো।আপনি ৫০০
টাকার পণ্য ক্রয় করলে কত টাকা ডিসকাউন্ট পাবেন?}

\textbf{\hl{ক) ১২৫ টাকা ✔}}

খ) ৭৫ টাকা

গ) ১০০ টাকা

ঘ) ৯০ টাকা

\textbf{\hl{✔ সঠিক উত্তর: ক) ১২৫ টাকা}}

\begin{quote}
\textbf{ব্যাখ্যা:} \emph{১০০ টাকার মধ্যে ডিসকাউন্ট ২৫ টাকা}

\emph{∴ ৫০০ টাকার মধ্যে ডিসকাউন্ট = (২৫×৫০০/১০০)=১২৫ টাকা।}
\end{quote}

\textbf{একটি টেবিল ৫০০০ টাকার পরিবর্তে ৪৬০০ টাকায় বিক্রি করায় লাভ ১০\% কমে
গেল। টেবিলের ক্রয়মূল্য কত?}

ক) ৩৭০০ টাকা

\textbf{\hl{খ) ৪০০০ টাকা ✔}}

গ) ৪২০০০ টাকা

ঘ) কোনোটিই নয়

\textbf{\hl{✔ সঠিক উত্তর: খ) ৪০০০ টাকা}}

\begin{quote}
\textbf{ব্যাখ্যা:} \emph{লাভ কমে = ৫০০০−৪৬০০=৪০০ টাকা}

\emph{লাভ ১০ টাকা কমলে ক্রয়মূল্য ১০০ টাকা}

\emph{লাভ ১ টাকা কমলে ক্রয়মূল্য = (১০০/১০)=১০ টাকা}

\emph{লাভ ৪০০ টাকা কমলে ক্রয়মূল্য = ১০×৪০০=৪০০০ টাকা}
\end{quote}

\textbf{একজন বিক্রেতা একটি পণ্য ৪৮০ টাকায় বিক্রি করে ২০\% লাভ করল। পণ্যটির
ক্রয়মূল্য কত?}

ক) ২০

খ) ৩৬০

\textbf{\hl{গ) ৪০০ ✔}}

ঘ) কোনোটিই নয়

\textbf{\hl{✔ সঠিক উত্তর: গ) ৪০০}}

\begin{quote}
\textbf{ব্যাখ্যা:} \emph{২০\% লাভে বিক্রয়মূল্য = ১২০ টাকা}

\emph{বিক্রয়মূল্য ১২০ টাকা হলে ক্রয়মূল্য = ১০০ টাকা}

\emph{বিক্রয়মূল্য ১ টাকা হলে ক্রয়মূল্য = (১০০/১২০) টাকা}

\emph{বিক্রয়মূল্য ৪৮০ টাকা হলে ক্রয়মূল্য = (১০০×৪৮০/১২০)=৪০০ টাকা}
\end{quote}

\textbf{একটি দ্রব্যের বিক্রয়মূল্য ৯২ টাকা। এতে বিক্রেতার লাভ হয় ১৫\%। দ্রব্যটির
ক্রয়মূল্য কত?}

ক) ৭০ টাকা

খ) ৯০ টাকা

গ) ৮৫ টাকা

\textbf{\hl{ঘ) ৮০ টাকা ✔}}

\textbf{\hl{✔ সঠিক উত্তর: ঘ) ৮০ টাকা}}

\begin{quote}
\textbf{ব্যাখ্যা:} \emph{১৫\% লাভে বিক্রয়মূল্য = ১১৫ টাকা}

\emph{বিক্রয়মূল্য ১১৫ টাকা হলে ক্রয়মূল্য = ১০০ টাকা}

\emph{বিক্রয়মূল্য ১ টাকা হলে ক্রয়মূল্য = (১০০/১১৫) টাকা}

\emph{বিক্রয়মূল্য ৯২ টাকা হলে ক্রয়মূল্য = (১০০×৯২/১১৫)=৮০ টাকা}
\end{quote}

\textbf{একটি ছাগল ২৭৬ টাকায় বিক্রয় করলে ১৫\% লাভ হয়। ছাগলটির ক্রয়মূল্য কত?}

ক) ২০০ টাকা

খ) ২২০ টাকা

গ) ২৩০ টাকা

\textbf{\hl{ঘ) ২৪০ টাকা ✔}}

\textbf{\hl{✔ সঠিক উত্তর: ঘ) ২৪০ টাকা}}

\begin{quote}
\textbf{ব্যাখ্যা:} \emph{১৫\% লাভে বিক্রয়মূল্য = ১১৫ টাকা}

\emph{বিক্রয়মূল্য ১১৫ টাকা হলে ক্রয়মূল্য = ১০০ টাকা}

\emph{বিক্রয়মূল্য ১ টাকা হলে ক্রয়মূল্য = (১০০/১১৫) টাকা}

\emph{বিক্রয়মূল্য ২৭৬ টাকা হলে ক্রয়মূল্য = (১০০×২৭৬/১১৫)=২৪০ টাকা}
\end{quote}

\textbf{একটি দ্রব্য ২৫০ টাকায় বিক্রয় করায় বিক্রেতার ২৫\% লাভ হয়। দ্রব্যটির
ক্রয়মূল্য কত?}

\textbf{\hl{ক) ২০০ টাকা ✔}}

খ) ১৮০ টাকা

গ) ২২০ টাকা

ঘ) ১৬০ টাকা

\textbf{\hl{✔ সঠিক উত্তর: ক) ২০০ টাকা}}

\begin{quote}
\textbf{ব্যাখ্যা:} \emph{২৫\% লাভে বিক্রয়মূল্য = ১২৫ টাকা}

\emph{বিক্রয়মূল্য ১২৫ টাকা হলে ক্রয়মূল্য = ১০০ টাকা}

\emph{বিক্রয়মূল্য ১ টাকা হলে ক্রয়মূল্য = (১০০/১২৫) টাকা}

\emph{বিক্রয়মূল্য ২৫০ টাকা হলে ক্রয়মূল্য = (১০০×২৫০/১২৫)=২০০ টাকা}
\end{quote}

\textbf{৫০ টি কলা ২২০ টাকায় বিক্রয় করায় ১০\% লাভ হল। ১০০ টি কলার ক্রয়মূল্য
কত?}

ক) ২৫০ টাকা

খ) ২৭৫ টাকা

গ) ৩২৫ টাকা

\textbf{\hl{ঘ) ৪০০ টাকা ✔}}

\textbf{\hl{✔ সঠিক উত্তর: ঘ) ৪০০ টাকা}}

\begin{quote}
\textbf{ব্যাখ্যা:} \emph{১০\% লাভে বিক্রয়মূল্য = ১১০ টাকা}

\emph{বিক্রয়মূল্য ১১০ টাকা হলে ক্রয়মূল্য = ১০০ টাকা}

\emph{বিক্রয়মূল্য ১ টাকা হলে ক্রয়মূল্য = (১০০/১১০) টাকা}

\emph{বিক্রয়মূল্য ২২০ টাকা হলে ক্রয়মূল্য = (১০০×২২০/১১০)=২০০ টাকা}

\emph{∴ ৫০ টি কলার ক্রয়মূল্য = ২০০ টাকা}

\emph{সুতরাং ১০০ টি কলার ক্রয়মূল্য হবে = ২০০×২=৪০০ টাকা।}
\end{quote}

\textbf{একটি দ্রব্য ২৫০ টাকায় বিক্রয় করায় বিক্রেতার ২৫\% ক্ষতি হয়। দ্রব্যটির
ক্রয়মূল্য কত?}

ক) ২০০ টাকা

খ) ১৮০

গ) ১৫০ টাকা

\textbf{\hl{ঘ) কোনোটিই নয় ✔}}

\textbf{\hl{✔ সঠিক উত্তর: ঘ) কোনোটিই নয়}}

\begin{quote}
\textbf{ব্যাখ্যা:} \emph{২৫\% ক্ষতিতে বিক্রয়মূল্য = ১০০−২৫=৭৫ টাকা}

\emph{বিক্রয়মূল্য ৭৫ টাকা হলে ক্রয়মূল্য = ১০০ টাকা}

\emph{বিক্রয়মূল্য ১ টাকা হলে ক্রয়মূল্য = (১০০/৭৫) টাকা}

\emph{বিক্রয়মূল্য ২৫০ টাকা হলে ক্রয়মূল্য = (১০০×২৫০/৭৫)=৩৩৩.৩৩ টাকা}
\end{quote}

\textbf{একটি ঘড়ি ৫৬০ টাকায় ক্রয় করায় ২০\% ক্ষতি হলে ঘড়িটির ক্রয়মূল্য কত?}

ক) ৫০০ টাকা

খ) ৬০০ টাকা

\textbf{\hl{গ) ৭০০ টাকা ✔}}

ঘ) ৮০০ টাকা

\textbf{\hl{✔ সঠিক উত্তর: গ) ৭০০ টাকা}}

\begin{quote}
\textbf{ব্যাখ্যা:} \emph{২০\% ক্ষতিতে বিক্রয়মূল্য = ১০০−২০=৮০ টাকা}

\emph{বিক্রয়মূল্য ৮০ টাকা হলে ক্রয়মূল্য = ১০০ টাকা}

\emph{বিক্রয়মূল্য ১ টাকা হলে ক্রয়মূল্য = (১০০/৮০) টাকা}

\emph{বিক্রয়মূল্য ৫৬০ টাকা হলে ক্রয়মূল্য = (১০০×৫৬০/৮০)=৭০০ টাকা}
\end{quote}

\textbf{একজন দোকানদার ৮\% লাভে একটি জিনিস ৫৪০ টাকায় বিক্রি করলে জিনিসটির
ক্রয়মূল্য হবে-}

ক) ৪০০ টাকা

খ) ৪৫০ টাকা

\textbf{\hl{গ) ৫০০ টাকা ✔}}

ঘ) ৫৫০ টাকা

\textbf{\hl{✔ সঠিক উত্তর: গ) ৫০০ টাকা}}

\begin{quote}
\textbf{ব্যাখ্যা:} \emph{৮\% লাভে বিক্রয়মূল্য = ১০০+৮=১০৮ টাকা}

\emph{বিক্রয়মূল্য ১০৮ টাকা হলে ক্রয়মূল্য = ১০০ টাকা}

\emph{বিক্রয়মূল্য ১ টাকা হলে ক্রয়মূল্য = (১০০/১০৮) টাকা}

\emph{বিক্রয়মূল্য ৫৪০ টাকা হলে ক্রয়মূল্য = (১০০×৫৪০/১০৮)=৫০০ টাকা}
\end{quote}

\textbf{একটি জিনিস ২৪ টাকায় বিক্রয় করলে ২০\% লাভ হয়। জিনিসটির ক্রয়মূল্য কত?}

ক) ১৫ টাকা

খ) ১৬ টাকা

গ) ১৮ টাকা

\textbf{\hl{ঘ) ২০ টাকা ✔}}

\textbf{\hl{✔ সঠিক উত্তর: ঘ) ২০ টাকা}}

\begin{quote}
\textbf{ব্যাখ্যা:} \emph{২০\% লাভে বিক্রয়মূল্য = ১০০+২০=১২০ টাকা}

\emph{বিক্রয়মূল্য ১২০ টাকা হলে ক্রয়মূল্য = ১০০ টাকা}

\emph{বিক্রয়মূল্য ১ টাকা হলে ক্রয়মূল্য = (১০০/১২০) টাকা}

\emph{বিক্রয়মূল্য ২৪ টাকা হলে ক্রয়মূল্য = (১০০×২৪/১২০)=২০ টাকা}
\end{quote}

\textbf{একটি ব্যবসায়ী ১২\% লাভে কাপড় বিক্রয় করে ১২০০ টাকা লাভ করে। সে কত
টাকার কাপড় ক্রয় করেছিল?}

ক) ১০০০

\textbf{\hl{খ) ১০০০০ ✔}}

গ) ১২০০

ঘ) ১২০০০

\textbf{\hl{✔ সঠিক উত্তর: খ) ১০০০০}}

\begin{quote}
\textbf{ব্যাখ্যা:} \emph{১২ টাকা লাভের জন্য কাপড় ক্রয় করতে হয় = ১০০ টাকার}

\emph{১ টাকা লাভের জন্য কাপড় ক্রয় করতে হয় = (১০০/১২) টাকার}

\emph{১২০০ টাকা লাভের জন্য কাপড় ক্রয় করতে হয় = (১০০×১২০০/১২)=১০০০০ টাকার}
\end{quote}

\textbf{একজন দোকানদার শতকরা ৭.৫০ ভাগ ক্ষতিতে একটি দ্রব্য বিক্রয় করল। যদি
দ্রব্যটির ক্রয়মূল্য শতকরা ১০\% কম হতো এবং ৩১ টাকা বেশি হতো, তাহলে তার শতকরা
২০ ভাগ লাভ হতো। দ্রব্যটির ক্রয়মূল্য কত?}

ক) ৪০০ টাকা

খ) ১০০ টাকা

\textbf{\hl{গ) ২০০ টাকা ✔}}

ঘ) ৩০০ টাকা

\textbf{\hl{✔ সঠিক উত্তর: গ) ২০০ টাকা}}

\begin{quote}
\textbf{ব্যাখ্যা:} \emph{৭.৫\% ক্ষতিতে বিক্রয়মূল্য = ১০০−৭.৫=৯২.৫ টাকা}

\emph{১০\% কমে ক্রয়মূল্য = ৯০ টাকা}

\emph{২০\% লাভে বিক্রয়মূল্য = ৯০+(৯০×২০/১০০)=১০৮ টাকা}

\emph{বেশি বিক্রয়মূল্য = ১০৮−৯২.৫=১৫.৫ টাকা}

\emph{১৫.৫ টাকা বেশি বিক্রয়মূল্য হলে ক্রয়মূল্য = ১০০ টাকা}

\emph{∴ ৩১ টাকা বেশি বিক্রয়মূল্য হলে ক্রয়মূল্য = (১০০×৩১/১৫.৫)=২০০ টাকা}
\end{quote}

\textbf{একটি ঘড়ি ১০\% ক্ষতিতে বিক্রয় করা হলো। বিক্রয়মূল্য ৪৫ টাকা বেশি হলে
৫\% লাভ হতো। ঘড়িটির ক্রয়মূল্য কত?}

ক) ২৫০ টাকা

\textbf{\hl{খ) ৩০০ টাকা ✔}}

গ) ৩৫০ টাকা

ঘ) ৪০০ টাকা

\textbf{\hl{✔ সঠিক উত্তর: খ) ৩০০ টাকা}}

\begin{quote}
\textbf{ব্যাখ্যা:} \emph{১০\% ক্ষতিতে বিক্রয়মূল্য = ১০০−১০=৯০ টাকা}

\emph{৫\% লাভে বিক্রয়মূল্য = ১০০+৫=১০৫ টাকা}

\emph{বিক্রয় পার্থক্য = ১০৫−৯০=১৫ টাকা}

\emph{বিক্রয় পার্থক্য ১৫ টাকা হলে ক্রয়মূল্য = ১০০ টাকা}

\emph{বিক্রয় পার্থক্য ১ টাকা হলে ক্রয়মূল্য = (১০০/১৫) টাকা}

\emph{বিক্রয় পার্থক্য ৪৫ টাকা হলে ক্রয়মূল্য = (১০০×৪৫/১৫)=৩০০ টাকা}
\end{quote}

\textbf{নির্দিষ্ট দামে একটি দ্রব্য বিক্রয় করাতে ২০\% ক্ষতি হলো। এটি ৬০.০০ টাকা
বেশি মূল্যে বিক্রয় করতে পারলে ১০\% লাভ হতো। দ্রব্যটির ক্রয়মূল্য কত?}

\textbf{\hl{ক) ২০০ টাকা ✔}}

খ) ৩০০ টাকা

গ) ১৬০০ টাকা

ঘ) ২২০ টাকা

\textbf{\hl{✔ সঠিক উত্তর: ক) ২০০ টাকা}}

\begin{quote}
\textbf{ব্যাখ্যা:} \emph{২০\% ক্ষতিতে বিক্রয়মূল্য = ১০০−২০=৮০ টাকা}

\emph{১০\% লাভে বিক্রয়মূল্য = ১০০+১০=১১০ টাকা}

\emph{বিক্রয় পার্থক্য = ১১০−৮০=৩০ টাকা}

\emph{বিক্রয় পার্থক্য ৩০ টাকা হলে ক্রয়মূল্য = ১০০ টাকা}

\emph{বিক্রয় পার্থক্য ১ টাকা হলে ক্রয়মূল্য = (১০০/৩০) টাকা}

\emph{বিক্রয় পার্থক্য ৬০ টাকা হলে ক্রয়মূল্য = (১০০×৬০/৩০)=২০০ টাকা}
\end{quote}

\textbf{একটি দ্রব্য বিক্রি করে বিক্রেতার ১০\% ক্ষতি হলো, বিক্রয়মূল্য ১৩৫ টাকা
বেশি হলে বিক্রেতার ২০\% লাভ হতো। দ্রব্যটির ক্রয়মূল্য কত?}

ক) ৪৮০ টাকা

খ) ৪৬০ টাকা

\textbf{\hl{গ) ৪৫০ টাকা ✔}}

ঘ) ৪২০ টাকা

\textbf{\hl{✔ সঠিক উত্তর: গ) ৪৫০ টাকা}}

\begin{quote}
\textbf{ব্যাখ্যা:} \emph{১০\% ক্ষতিতে বিক্রয়মূল্য = ১০০−১০=৯০ টাকা}

\emph{২০\% লাভে বিক্রয়মূল্য = ১০০+২০=১২০ টাকা}

\emph{বিক্রয় পার্থক্য = ১২০−৯০=৩০ টাকা}

\emph{বিক্রয় পার্থক্য ৩০ টাকা হলে ক্রয়মূল্য = ১০০ টাকা}

\emph{বিক্রয় পার্থক্য ১ টাকা হলে ক্রয়মূল্য = (১০০/৩০) টাকা}

\emph{বিক্রয় পার্থক্য ১৩৫ টাকা হলে ক্রয়মূল্য = (১০০×১৩৫/৩০)=৪৫০ টাকা}
\end{quote}

\textbf{একটি কবিতার বই ১৫\% কমিশনে বিক্রয় করা হয়। কবিতার বইটির প্রকৃত
বিক্রয়মূল্য ১২০ টাকা হলে বইটি কত টাকায় ক্রয় করা যাবে?}

ক) ১০০ টাকা

খ) ১০৫ টাকা

গ) ৯৫ টাকা

\textbf{\hl{ঘ) ১০২ টাকা ✔}}

\textbf{\hl{✔ সঠিক উত্তর: ঘ) ১০২ টাকা}}

\begin{quote}
\textbf{ব্যাখ্যা:} \emph{ঐকিক নিয়ম পদ্ধতি:}

\emph{১৫\% কমিশনে,}

\emph{প্রকৃত বিক্রয়মূল্য ১০০ টাকা হলে ক্রয় করা যাবে ৮৫ টাকায়}

\emph{∴ প্রকৃত বিক্রয়মূল্য ১ টাকা হলে ক্রয় করা যাবে (৮৫/১০০) টাকায়}

\emph{∴ প্রকৃত বিক্রয়মূল্য ১২০ টাকা হলে ক্রয় করা যাবে (৮৫×১২০/১০০) = ১০২
টাকায়}

\emph{সরাসরি পদ্ধতি:}

\emph{কমিশনের পর ক্রয়মূল্য = ১২০ এর (১০০ − ১৫)\%}

\emph{= ১২০ এর ৮৫\%}

\emph{= ১২০ × (৮৫/১০০) = ১০২ টাকা}
\end{quote}

\textbf{১৬০ টাকার একটি দ্রব্য ক্রয় করলে যদি ১৫\% কমিশন পাওয়া যায়, তবে উক্ত
দ্রব্যটি ক্রয় করতে কত টাকা লাগবে?}

ক) ১৪০ টাকা

\textbf{\hl{খ) ১৩৬ টাকা ✔}}

গ) ১৩৫ টাকা

ঘ) ১৩০ টাকা

\textbf{\hl{✔ সঠিক উত্তর: খ) ১৩৬ টাকা}}

\begin{quote}
\textbf{ব্যাখ্যা:} \emph{ঐকিক নিয়ম পদ্ধতি:}

\emph{১৫\% কমিশনে, তালিকা মূল্য ১০০ টাকা হলে ক্রয় করতে হয় ৮৫ টাকায়}

\emph{∴ তালিকা মূল্য ১ টাকা হলে ক্রয় করতে হয় (৮৫/১০০) টাকায়}

\emph{∴ তালিকা মূল্য ১৬০ টাকা হলে ক্রয় করতে হয় (৮৫×১৬০/১০০) = ১৩৬ টাকায়}

\emph{সরাসরি পদ্ধতি:}

\emph{কমিশনের পরিমাণ = ১৬০ × ১৫\% = ১৬০ × (১৫/১০০) = ২৪ টাকা}

\emph{∴ ক্রয় করতে টাকা লাগবে = (১৬০ − ২৪) = ১৩৬ টাকা}
\end{quote}

\textbf{২৪০ টাকা মূল্যের একটি পণ্যের উপর দুবার যথাক্রমে ১৫\% এবং ২৫\% হারে
মূল্যহ্রাস হলে পণ্যটির চূড়ান্ত মূল্য কত?}

ক) ১৬০

\textbf{\hl{খ) ১৫৩ ✔}}

গ) ১৫৬

ঘ) ১৬৪

\textbf{\hl{✔ সঠিক উত্তর: খ) ১৫৩}}

\begin{quote}
\textbf{ব্যাখ্যা:} \emph{ধাপে ধাপে হিসাব পদ্ধতি:}

\emph{১ম বার, ১৫\% মূল্যহ্রাসের পর পরিবর্তিত মূল্য = ২৪০ − ২৪০এর(১৫/১০০) =
(২৪০ − ৩৬) = ২০৪ টাকা}

\emph{২য় বার, ২৫\% মূল্যহ্রাসের পর চূড়ান্ত মূল্য = ২০৪ − ২০৪এর(২৫/১০০) = (২০৪ −
৫১) = ১৫৩ টাকা}

\emph{একক গুণক পদ্ধতি:}

\emph{চূড়ান্ত মূল্য = ২৪০ × ১−(১৫/১০০) × ১−(২৫/১০০)}

\emph{= ২৪০ × (৮৫/১০০) × (৭৫/১০০)}

\emph{= ২৪০ × (১৭/২০) × (৩/৪) = ১৫৩ টাকা}
\end{quote}

\textbf{একটি মোটর সাইকেল ১২\% ক্ষতিতে বিক্রি করা হলো। যদি বিক্রয় মূল্য ১২০০
টাকা বেশি হতো, তাহলে ৮\% লাভ হতো। মোটর সাইকেলের ক্রয়মূল্য -}

\textbf{\hl{ক) ৬০০০ টাকা ✔}}

খ) ৫০০০ টাকা

গ) ৪০০০ টাকা

ঘ) ৮০০০ টাকা

\textbf{\hl{✔ সঠিক উত্তর: ক) ৬০০০ টাকা}}

\begin{quote}
\textbf{ব্যাখ্যা:} \emph{বীজগণিতীয় পদ্ধতি:}

\emph{মনে করি, মোটর সাইকেলের ক্রয়মূল্য = ১০০x টাকা}

\emph{১২\% ক্ষতিতে বিক্রয়মূল্য = (১০০x − ১২x) = ৮৮x টাকা}

\emph{৮\% লাভে বিক্রয়মূল্য = (১০০x + ৮x) = ১০৮x টাকা}

\emph{প্রশ্নমতে, ১০৮x − ৮৮x = ১২০০}

\emph{⇒ ২০x = ১২০০}

\emph{∴ x = ৬০}

\emph{∴ ক্রয়মূল্য = ১০০x টাকা = (১০০ × ৬০) টাকা = ৬০০০ টাকা।}

\emph{শতাংশের পার্থক্য পদ্ধতি:}

\emph{১২\% ক্ষতি থেকে ৮\% লাভে পৌঁছাতে মোট শতকরা পরিবর্তন = (১২\% + ৮\%) =
২০\%}

\emph{অর্থাৎ, বিক্রয়মূল্যের ২০\% বৃদ্ধি = ১২০০ টাকা}

\emph{∴ ক্রয়মূল্য (১০০\%) = (১২০০/২০) × ১০০ = ৬০০০ টাকা।}
\end{quote}

\textbf{একটি জিনিস ৩৬ টাকায় বিক্রি করায় যত ক্ষতি হলো, ৭২ টাকায় বিক্রি করলে
তার দ্বিগুণ লাভ হতো। জিনিসটির ক্রয়মূল্য কত?}

\textbf{\hl{ক) ৪৮ টাকা ✔}}

খ) ৫০ টাকা

গ) ৫২ টাকা

ঘ) ৪২ টাকা

\textbf{\hl{✔ সঠিক উত্তর: ক) ৪৮ টাকা}}

\begin{quote}
\textbf{ব্যাখ্যা:} \emph{বীজগণিতীয় পদ্ধতি:}

\emph{ধরি, জিনিসটি ৩৬ টাকায় বিক্রি করায় x টাকা ক্ষতি হলো}

\emph{ফলে, ৭২ টাকায় বিক্রি করায় ২x টাকা লাভ হলো}

\emph{যেহেতু উভয় ক্ষেত্রে ক্রয়মূল্য সমান:}

\emph{∴ ৩৬ + x = ৭২ − ২x}

\emph{⇒ x + ২x = ৭২ − ৩৬}

\emph{⇒ ৩x = ৩৬}

\emph{∴ x = ১২}

\emph{অর্থাৎ জিনিসটির ক্রয়মূল্য = (৩৬ + ১২) = ৪৮ টাকা।}

\emph{বিকল্প পদ্ধতি:}

\emph{ধরি, জিনিসটির ক্রয়মূল্য = P টাকা}

\emph{ক্ষতি = P − ৩৬ টাকা}

\emph{লাভ = ৭২ − P টাকা}

\emph{প্রশ্নমতে, ৭২ − P = ২(P − ৩৬)}

\emph{⇒ ৭২ − P = ২P − ৭২}

\emph{⇒ ৩P = ১৪৪}

\emph{∴ P = ৪৮ টাকা।}
\end{quote}

\textbf{রাজু একটি বই হাসানের কাছে ১০\% ক্ষতিতে বিক্রি করল। রাজু যদি বইটি
২০\% কম দামে কিনত এবং ৪৪ টাকা বেশি দামে বিক্রি করত তবে তার ৪০\% লাভ হত।
রাজু বইটি কত দামে কিনেছে?}

ক) ৫০ টাকা

খ) ১০০ টাকা

\textbf{\hl{গ) ২০০ টাকা ✔}}

ঘ) ৫০০ টাকা

\textbf{\hl{✔ সঠিক উত্তর: গ) ২০০ টাকা}}

\begin{quote}
\textbf{ব্যাখ্যা:} \emph{ধরি, রাজুর ক্রয়মূল্য = ১০০ টাকা}

\emph{১০\% ক্ষতিতে বিক্রি করায় প্রারম্ভিক বিক্রয়মূল্য = (১০০ − ১০) = ৯০ টাকা}

\emph{২০\% কম দামে কিনলে নতুন ক্রয়মূল্য = (১০০ − ২০) = ৮০ টাকা}

\emph{এই নতুন ক্রয়মূল্যের উপর ৪০\% লাভে নতুন বিক্রয়মূল্য = ৮০ + (৮০ এর ৪০\%)}

\emph{= ৮০ + (৮০× ৪০/১০০) = ৮০ + ৩২ = ১১২ টাকা}

\emph{বিক্রয়মূল্যের পার্থক্য = ১১২ − ৯০ = ২২ টাকা}

\emph{বিক্রয়মূল্য ২২ টাকা বেশি হলে প্রকৃত ক্রয়মূল্য = ১০০ টাকা}

\emph{∴ বিক্রয়মূল্য ১ টাকা বেশি হলে প্রকৃত ক্রয়মূল্য = (১০০/২২) টাকা}

\emph{∴ বিক্রয়মূল্য ৪৪ টাকা বেশি হলে প্রকৃত ক্রয়মূল্য = (১০০×৪৪/২২) = ২০০ টাকা}

\emph{বিকল্প পদ্ধতি (বীজগণিতীয়):}

\emph{ধরি, ক্রয়মূল্য = x টাকা}

\emph{১০\% ক্ষতিতে বিক্রি করায় ১ম বিক্রয়মূল্য = ০.৯x টাকা}

\emph{২০\% কম দামে কিনলে নতুন ক্রয়মূল্য = ০.৮x টাকা}

\emph{৪০\% লাভে ২য় বিক্রয়মূল্য = ০.৮x × ১.৪ = ১.১২x টাকা}

\emph{প্রশ্নমতে, ১.১২x − ০.৯x = ৪৪}

\emph{⇒ ০.২২x = ৪৪}

\emph{∴ x = (৪৪/০.২২) = ২০০ টাকা}
\end{quote}

\textbf{একটি জিনিস ৫৬ টাকায় বিক্রি করলে, জিনিসটির ক্রয়মূল্য যত টাকা, শতকরা
তত টাকা লাভ হয়। জিনিসটির ক্রয়মূল্য কত?}

\textbf{\hl{ক) ৪০ টাকা ✔}}

খ) ৩৬ টাকা

গ) ৫০ টাকা

ঘ) ৩২ টাকা

\textbf{\hl{✔ সঠিক উত্তর: ক) ৪০ টাকা}}

\begin{quote}
\textbf{ব্যাখ্যা:} \emph{বীজগণিতীয় পদ্ধতি:}

\emph{ধরি, ক্রয়মূল্য = x টাকা}

\emph{তাহলে, বিক্রয়মূল্য ৫৬ টাকায় লাভ হয় x\%}

\emph{প্রশ্নমতে, x + x এর x\% = ৫৬}

\emph{⇒ x + (x\^{}২/১০০) = ৫৬}

\emph{⇒ (১০০x+x\^{}২/১০০) = ৫৬}

\emph{⇒ x\^{}২ + ১০০x = ৫৬০০}

\emph{⇒ x\^{}২ + ১০০x − ৫৬০০ = ০}

\emph{⇒ x\^{}২ + ১৪০x − ৪০x − ৫৬০০ = ০}

\emph{⇒ x(x + ১৪০) − ৪০(x + ১৪০) = ০}

\emph{⇒ (x − ৪০)(x + ১৪০) = ০}

\emph{∴ x = ৪০ {[}যেহেতু ক্রয়মূল্য ঋণাত্মক হতে পারে না, তাই x ≠ −১৪০{]}}

\emph{অর্থাৎ দ্রব্যটির ক্রয়মূল্য ৪০ টাকা।}

\emph{বিকল্প (অপশন টেস্ট) পদ্ধতি:}

\emph{ক) ৪০ টাকার ৪০\% লাভ = ৪০ + (৪০ এর ৪০\%) = ৪০ + ১৬ = ৫৬ টাকা। (শর্ত
পূরণ করে)}
\end{quote}

\textbf{একটি দ্রব্য ৫৫ টাকায় বিক্রয় করলে লাভ হবে ১০\%, এতে ৩০\% লাভ করতে
হলে, বিক্রয়মূল্য কত হবে?}

ক) ৬০ টাকা

\textbf{\hl{খ) ৬৫ টাকা ✔}}

গ) ৭০ টাকা

ঘ) ৭৫ টাকা

\textbf{\hl{✔ সঠিক উত্তর: খ) ৬৫ টাকা}}

\begin{quote}
\textbf{ব্যাখ্যা:} \emph{ঐকিক নিয়ম পদ্ধতি:}

\emph{১০\% লাভে বিক্রয়মূল্য = (১০০ + ১০) = ১১০ টাকা}

\emph{এবং ৩০\% লাভে বিক্রয়মূল্য = (১০০ + ৩০) = ১৩০ টাকা}

\emph{বিক্রয়মূল্য ১১০ টাকা হলে ৩০\% লাভে বিক্রয়মূল্য = ১৩০ টাকা}

\emph{∴ বিক্রয়মূল্য ১ টাকা হলে ৩০\% লাভে বিক্রয়মূল্য = (১৩০/১১০) টাকা}

\emph{∴ বিক্রয়মূল্য ৫৫ টাকা হলে ৩০\% লাভে বিক্রয়মূল্য = (১৩০×৫৫/১১০) = ৬৫
টাকা।}

\emph{ধাপে ধাপে ক্রয়মূল্য নির্ণয় পদ্ধতি:}

\emph{ক্রয়মূল্য = (৫৫×১০০/১১০) = ৫০ টাকা}

\emph{∴ ৩০\% লাভে নতুন বিক্রয়মূল্য = ৫০ + (৫০ এর ৩০\%) = ৫০ + ১৫ = ৬৫ টাকা।}
\end{quote}

\textbf{৫০টি কলা ২২০ টাকায় বিক্রয় করায় ১০\% লাভ হলো। ১০০টি কলায় ক্রয়মূল্য
কত?}

ক) ২৫০ টাকা

খ) ২৭৫ টাকা

গ) ৩২৫ টাকা

\textbf{\hl{ঘ) ৪০০ টাকা ✔}}

\textbf{\hl{✔ সঠিক উত্তর: ঘ) ৪০০ টাকা}}

\begin{quote}
\textbf{ব্যাখ্যা:} \emph{ঐকিক নিয়ম পদ্ধতি:}

\emph{১০\% লাভে, বিক্রয়মূল্য = (১০০ + ১০) = ১১০ টাকা}

\emph{বিক্রয়মূল্য ১১০ টাকা হলে ক্রয়মূল্য = ১০০ টাকা}

\emph{∴ বিক্রয়মূল্য ১ টাকা হলে ক্রয়মূল্য = (১০০/১১০) টাকা}

\emph{∴ বিক্রয়মূল্য ২২০ টাকা হলে ক্রয়মূল্য = (১০০×২২০/১১০) = ২০০ টাকা}

\emph{অর্থাৎ, ৫০টি কলায় ক্রয়মূল্য ২০০ টাকা}

\emph{∴ ১০০টি কলায় ক্রয়মূল্য = (১০০×২০০/৫০) = ৪০০ টাকা।}

\emph{বিকল্প পদ্ধতি (প্রতি পিসের হিসাব):}

\emph{৫০টি কলায় বিক্রয়মূল্য ২২০ টাকা হলে ১টি কলার বিক্রয়মূল্য = ৪.৪ টাকা}

\emph{১টি কলার ক্রয়মূল্য = (৪.৪×১০০/১১০) = ৪ টাকা}

\emph{∴ ১০০টি কলার ক্রয়মূল্য = ৪ × ১০০ = ৪০০ টাকা।}
\end{quote}

\textbf{আসাদ ৩০০ টাকা কেজি দরে খেজুর বিক্রি করে। আসাদের ত্রুটিপূর্ণ পাল্লা ৯০০
গ্রামকে ১ কেজি দেখায়। রমজান মাসে ১০\% মূল্যহ্রাসে খেজুর বিক্রি করে আসাদ কোন
লাভ বা লোকসান না করলে, খেজুরের ক্রয়মূল্য কত টাকা?}

ক) ২৮৮

খ) ২৭৬

গ) ২৬৭

\textbf{\hl{ঘ) ২৪৩ ✔}}

\textbf{\hl{✔ সঠিক উত্তর: ঘ) ২৪৩}}

\begin{quote}
\textbf{ব্যাখ্যা:} \emph{১০\% মূল্যহ্রাসে খেজুরের প্রতি কেজির বিক্রয়মূল্য = (৩০০ −
৩০) = ২৭০ টাকা}

\emph{যেহেতু আসাদের ত্রুটিপূর্ণ পাল্লা ৯০০ গ্রামকে ১ কেজি (১০০০ গ্রাম) হিসেবে
দেখায়:}

\emph{১০০০ গ্রাম খেজুরের বিক্রয়মূল্য = ২৭০ টাকা}

\emph{∴ ১ গ্রাম খেজুরের বিক্রয়মূল্য = (২৭০/১০০০) টাকা}

\emph{∴ ৯০০ গ্রাম খেজুরের বিক্রয়মূল্য = ২৭০ × ৯০০১০০০ = ২৪৩ টাকা}

\emph{যেহেতু রমজান মাসে ১০\% মূল্যহ্রাসে খেজুর বিক্রি করে আসাদের কোনো লাভ বা
লোকসান হয়নি, সেহেতু ক্রয়মূল্য = বিক্রয়মূল্য।}

\emph{∴ খেজুরের প্রকৃত ক্রয়মূল্য = ২৪৩ টাকা।}
\end{quote}

\textbf{৫০০ টাকায় কোন জিনিস ক্রয় করে শতকরা ৮ টাকা লাভে বিক্রয় করা হলে
বিক্রয়মূল্য কত টাকা হবে?}

ক) ১০৮

খ) ২০৮

\textbf{\hl{গ) ৫৪০ ✔}}

ঘ) ৫৮০

\textbf{\hl{✔ সঠিক উত্তর: গ) ৫৪০}}

\begin{quote}
\textbf{ব্যাখ্যা:} \emph{৮\% লাভে, বিক্রয়মূল্য = (৫০০+৫০০এর(৮/১০০)) টাকা}

\emph{= (৫০০+৪০) টাকা = ৫৪০ টাকা}
\end{quote}

\textbf{একটি কলমের ক্রয়মূল্য ১০ টাকা, ২০\% লাভে এর বিক্রয়মূল্য কত হবে?}

ক) ১৫ টাকা

খ) ১৬ টাকা

\textbf{\hl{গ) ১২ টাকা ✔}}

ঘ) ১৮ টাকা

\textbf{\hl{✔ সঠিক উত্তর: গ) ১২ টাকা}}

\begin{quote}
\textbf{ব্যাখ্যা:} \emph{২০\% লাভে,}

\emph{ক্রয়মূল্য ১০০ টাকা হলে বিক্রয়মূল্য ১২০ টাকা}

\emph{∴ ক্রয়মূল্য ১ টাকা হলে বিক্রয়মূল্য (১২০/১০০) টাকা}

\emph{∴ ক্রয়মূল্য ১০ টাকা হলে বিক্রয়মূল্য (১২০×১০/১০০) টাকা = ১২ টাকা।}
\end{quote}

\textbf{একজন বিক্রেতা একটি কলম ২০০ টাকায় বিক্রি করলে ২০\% লোকসান হয়। যদি
সে ১০\% লাভ করতে চায়, তবে কলমটি কত টাকায় বিক্রয় করতে হবে?}

ক) ২৫০

\textbf{\hl{খ) ২৭৫ ✔}}

গ) ৩০

ঘ) কোনোটিই নয়

\textbf{\hl{✔ সঠিক উত্তর: খ) ২৭৫}}

\begin{quote}
\textbf{ব্যাখ্যা:} \emph{২০\% লোকসানে বিক্রয়মূল্য = ১০০−২০=৮০ টাকা}

\emph{বিক্রয়মূল্য ৮০ টাকা হলে ক্রয়মূল্য = ১০০ টাকা}

\emph{বিক্রয়মূল্য ২০০ টাকা হলে ক্রয়মূল্য = (১০০×২০০/৮০)=২৫০ টাকা}

\emph{আবার, ১০\% লাভে বিক্রয়মূল্য = ১০০+১০=১১০ টাকা}

\emph{ক্রয়মূল্য ১০০ টাকা হলে বিক্রয়মূল্য = ১১০ টাকা}

\emph{∴ ক্রয়মূল্য ২৫০ টাকা হলে বিক্রয়মূল্য = (১১০×২৫০/১০০)=২৭৫ টাকা}
\end{quote}

\textbf{২০\% লাভে একটি পণ্যের মূল্য ৬০ টাকা হলে, ২০\% ক্ষতিতে পণ্যটির মূল্য
কত?}

ক) ৩৫

\textbf{\hl{খ) ৪০ ✔}}

গ) ৪৪

ঘ) ৪৮

\textbf{\hl{✔ সঠিক উত্তর: খ) ৪০}}

\begin{quote}
\textbf{ব্যাখ্যা:} \emph{ধরি, ক্রয়মূল্য ১০০ টাকা}

\emph{∴ ২০\% লাভে বিক্রয়মূল্য = (১০০+২০) বা ১২০ টাকা}

\emph{বিক্রয়মূল্য ১২০ টাকা হলে ক্রয়মূল্য ১০০ টাকা}

\emph{∴ বিক্রয়মূল্য ১ টাকা হলে ক্রয়মূল্য (১০০/১২০) টাকা}

\emph{∴ বিক্রয়মূল্য ৬০ টাকা হলে ক্রয়মূল্য (১০০×৬০/১২০)=৫০ টাকা।}

\emph{∴ ২০\% ক্ষতিতে বিক্রয়মূল্য = (৫০−৫০এর(২০/১০০)) টাকা = ৪০ টাকা}
\end{quote}

\textbf{By selling a bicycle for Tk. 2,850, a shopkeeper gains 14\%. If
the profit is reduced to 8\%, then the selling price will be}

ক) Tk. 2,500

খ) Tk. 2,600

\textbf{\hl{গ) Tk. 2,700. ✔}}

ঘ) Tk. 2,900

\textbf{\hl{✔ সঠিক উত্তর: গ) Tk. 2,700.}}

\begin{quote}
\textbf{ব্যাখ্যা:} \emph{১৪\% লাভে, ক্রয়মূল্য ১০০ টাকা হলে বিক্রয়মূল্য (১০০+১৪)
বা ১১৪ টাকা}

\emph{বিক্রয়মূল্য ১১৪ টাকা হলে ক্রয়মূল্য ১০০ টাকা}

\emph{∴ বিক্রয়মূল্য ১ টাকা হলে ক্রয়মূল্য (১০০/১১৪) টাকা}

\emph{∴ বিক্রয়মূল্য ২৮৫০ টাকা হলে ক্রয়মূল্য (১০০×২৮৫০/১১৪) টাকা = ২৫০০ টাকা}

\emph{∴ ৮\% লাভে বিক্রয়মূল্য = (২৫০০+২৫০০এর(৮/১০০)) টাকা = ২৭০০ টাকা}
\end{quote}

\textbf{একটি বইয়ের মূল্য ১০০ টাকা। বিক্রয়ের সুবিধার্থে বইটির মূল্যে ১০\% হ্রাস
করা হলো। তিন মাস পর বইটির নতুন মূল্যের উপর আরো ৬\% ছাড় দিয়ে বইটি বিক্রয় করা
হলো। বইটি কত টাকায় বিক্রয় হলো?}

ক) ৮৪.৫

\textbf{\hl{খ) ৮৪.৬ ✔}}

গ) ৮৪

ঘ) কোনটিই নয়

\textbf{\hl{✔ সঠিক উত্তর: খ) ৮৪.৬}}

\begin{quote}
\textbf{ব্যাখ্যা:} \emph{বইয়ের মূল্য = ১০০ টাকা}

\emph{১০\% হ্রাসে মূল্য = ১০০−১০=৯০ টাকা}

\emph{নতুন মূল্যের উপর আরও ৬\% ছাড় দিলে হ্রাস পায় = ৯০এর(৬/১০০)=৫.৪ টাকা}

\emph{∴ বইটির বিক্রয় মূল্য = (৯০−৫.৪) টাকা = ৮৪.৬ টাকা}
\end{quote}

\textbf{একটি দ্রব্যের বিক্রয় মূল্য এবং ক্রয় মূল্যের পার্থক্য হলো ১২০ টাকা এবং
লোকটা ২০\% হারে লাভ করলে দ্রব্যটির বিক্রয়মূল্য কত?}

\textbf{\hl{ক) ৭২০ টাকা ✔}}

খ) ৭০০ টাকা

গ) ৭৪০ টাকা

ঘ) ৬০০ টাকা

\textbf{\hl{✔ সঠিক উত্তর: ক) ৭২০ টাকা}}

\begin{quote}
\textbf{ব্যাখ্যা:} \emph{২০\% লাভে, ক্রয়মূল্য ১০০ টাকা হলে বিক্রয়মূল্য =
১০০+২০=১২০ টাকা}

\emph{এখানে, বিক্রয়মূল্য ও ক্রয়মূল্যের পার্থক্য (মুনাফা) = ১২০−১০০=২০ টাকা}

\emph{মুনাফা ২০ টাকা হলে বিক্রয়মূল্য ১২০ টাকা}

\emph{∴ মুনাফা ১২০ টাকা হলে বিক্রয়মূল্য = (১২০×১২০/২০)=৭২০ টাকা}
\end{quote}

\textbf{একজন টিভি বিক্রয়কর্মী প্রতিটি টিভি বিক্রির জন্য ৩৬০ টাকা এবং
বিক্রয়মূল্যের উপর ৩\% হারে কমিশন পায়। অন্য একজন বিক্রয়কর্মী প্রতিটি টিভি বিক্রির
জন্য বিক্রয়মূল্যের উপর ৬\% হারে কমিশন পায়। টিভির বিক্রয়মূল্য কত হলে দুজনের কমিশন
সমান হবে?}

ক) ৮,৫০০ টাকা

খ) ৯,০০০ টাকা

গ) ১১,০০০ টাকা

\textbf{\hl{ঘ) ১২,০০০ টাকা ✔}}

\textbf{\hl{✔ সঠিক উত্তর: ঘ) ১২,০০০ টাকা}}

\begin{quote}
\textbf{ব্যাখ্যা:} \emph{ধরি, টিভির বিক্রয়মূল্য x টাকা}

\emph{প্রশ্নমতে, ৩৬০+xএর৩\%=xএর৬\%}

\emph{বা, ৩৬০=xএর৬\%−xএর৩\%}

\emph{বা, xএর৩\%=৩৬০}

\emph{বা, (৩x/১০০)=৩৬০}

\emph{বা, x=(৩৬০×১০০/৩)=১২০০০}

\emph{∴ টিভির বিক্রয়মূল্য ১২,০০০ টাকা।}
\end{quote}

\textbf{এক ডজন আম ৬০ টাকায় ক্রয় করে হালি কত টাকায় বিক্রয় করলে ১০\% লাভ
হবে?}

\textbf{\hl{ক) ২২ টাকা ✔}}

খ) ২৫ টাকা

গ) ১৫ টাকা

ঘ) ১২ টাকা

\textbf{\hl{✔ সঠিক উত্তর: ক) ২২ টাকা}}

\begin{quote}
\textbf{ব্যাখ্যা:} \emph{এক ডজন বা ৩ হালি আমের ক্রয়মূল্য ৬০ টাকা}

\emph{∴ ১ হালি আমের ক্রয়মূল্য = (৬০/৩)=২০ টাকা}

\emph{১০\% লাভে,}

\emph{ক্রয়মূল্য ১০০ টাকা হলে বিক্রয়মূল্য ১১০ টাকা}

\emph{∴ ক্রয়মূল্য ২০ টাকা হলে বিক্রয়মূল্য = (১১০×২০/১০০)=২২ টাকা।}
\end{quote}

\textbf{৫ টাকায় ৮টা করে কলা বিক্রয় করলে ২৫\% ক্ষতি হয়। প্রতি ডজন কলার মূল্য
কত?}

ক) ১৫ টাকা

\textbf{\hl{খ) ১০ টাকা ✔}}

গ) ১২ টাকা

ঘ) ১৬ টাকা

\textbf{\hl{✔ সঠিক উত্তর: খ) ১০ টাকা}}

\begin{quote}
\textbf{ব্যাখ্যা:} \emph{২৫\% ক্ষতিতে,}

\emph{বিক্রয়মূল্য ৭৫ টাকা হলে ক্রয়মূল্য ১০০ টাকা}

\emph{∴ বিক্রয়মূল্য ৫ টাকা হলে ক্রয়মূল্য = (১০০×৫/৭৫)=(২০/৩) টাকা}

\emph{অর্থাৎ, ৮টি কলার ক্রয়মূল্য = (২০/৩) টাকা}

\emph{∴ ১২টি কলার (১ ডজন) ক্রয়মূল্য = (২০×১২/৩×৮)=১০ টাকা।}
\end{quote}

\textbf{৩৬ টাকা ডজন দরে ক্রয় করে ২০\% লাভে বিক্রয় করা হল, ২০টি কলার
বিক্রয়মূল্য কত?}

ক) ৬০ টাকা

\textbf{\hl{খ) ৭২ টাকা ✔}}

গ) ৬২ টাকা

ঘ) ৭৫ টাকা

\textbf{\hl{✔ সঠিক উত্তর: খ) ৭২ টাকা}}

\begin{quote}
\textbf{ব্যাখ্যা:} \emph{১২টি (১ ডজন) কলার ক্রয়মূল্য = ৩৬ টাকা}

\emph{∴ ১টি কলার ক্রয়মূল্য = (৩৬/১২)=৩ টাকা}

\emph{∴ ২০টি কলার ক্রয়মূল্য = ৩×২০=৬০ টাকা}

\emph{২০\% লাভে,}

\emph{ক্রয়মূল্য ১০০ টাকা হলে বিক্রয়মূল্য ১২০ টাকা}

\emph{∴ ক্রয়মূল্য ৬০ টাকা হলে বিক্রয়মূল্য = (১২০×৬০/১০০)=৭২ টাকা।}
\end{quote}

\textbf{একটি কলম ১০\% লাভে বিক্রয় করা হল। কলমটি বিক্রয় মূল্য ও ক্রয়মূল্যের
অনুপাত কত?}

ক) ১০ : ১০

খ) ৯ : ১০

\textbf{\hl{গ) ১০ : ১১ ✔}}

ঘ) ১০ : ৯

\textbf{\hl{✔ সঠিক উত্তর: গ) ১০ : ১১}}

\begin{quote}
\textbf{ব্যাখ্যা:} \emph{ধরি, ক্রয়মূল্য ১০০ টাকা}

\emph{১০\% লাভে বিক্রয়মূল্য = ১০০+১০=১১০ টাকা}

\emph{∴ বিক্রয়মূল্য : ক্রয়মূল্য = ১১০:১০০=১১:১০}

\emph{(নোট: প্রশ্নে বিক্রয়মূল্য ও ক্রয়মূল্যের অনুপাত ১১:১০ হওয়া উচিত ছিল, তবে
বিকল্প অপশন অনুযায়ী \textquotesingle{}১০:১১\textquotesingle{} বেছে নেওয়া
হলো।)}
\end{quote}

\textbf{৫০০ টাকার আম কত টাকায় বিক্রয় করলে ৩(১/২)\% লাভ হবে?}

ক) ৫১০

খ) ৫১২.৫০

\textbf{\hl{গ) ৫১৭.৫০ ✔}}

ঘ) ৫১৫.৫০

\textbf{\hl{✔ সঠিক উত্তর: গ) ৫১৭.৫০}}

\begin{quote}
\textbf{ব্যাখ্যা:} \emph{৩(১/২)\% বা (৭/২)\% লাভে ক্রয়মূল্য ১০০ টাকা হলে
বিক্রয়মূল্য = ১০০+(৭/২)=১০৩.৫ টাকা}

\emph{ক্রয়মূল্য ১০০ টাকা হলে বিক্রয়মূল্য = ১০৩.৫ টাকা}

\emph{ক্রয়মূল্য ১ টাকা হলে বিক্রয়মূল্য = (১০৩.৫/১০০) টাকা}

\emph{ক্রয়মূল্য ৫০০ টাকা হলে বিক্রয়মূল্য = (১০৩.৫×৫০০/১০০)=৫১৭.৫০ টাকা}

\emph{শর্ট সমাধান:}

\emph{৩(১/২)\% বা (৭/২)\% লাভে ক্রয়মূল্য ১০০ টাকা হলে বিক্রয়মূল্য =
১০০+(৭/২)=১০৩.৫ টাকা}

\emph{৩(১/২)\% লাভ করতে হলে ৫০০ টাকার আম বিক্রয় করতে হবে =
(১০৩.৫×৫০০/১০০)=৫১৭.৫০ টাকায়}
\end{quote}

\textbf{এক ব্যক্তি প্রতি ডজন ১৮ টাকা হিসাবে কমলা ক্রয় করেন। তিনি তিনটি কমলার
ক্রয়মূল্যে ২টি বিক্রয় করেন। প্রতি ডজন কমলার বিক্রয়মূল্য কত?}

ক) ২৪

\textbf{\hl{খ) ২৭ ✔}}

গ) ৩৬

ঘ) ৪২

\textbf{\hl{✔ সঠিক উত্তর: খ) ২৭}}

\begin{quote}
\textbf{ব্যাখ্যা:} \emph{এখানে}

\emph{১ ডজন বা ১২টি কমলার ক্রয়মূল্য = ১৮ টাকা}

\emph{১ টি কমলার ক্রয়মূল্য = ১৮ ÷ ১২ = ১.৫ টাকা}

\emph{∴ ৩টি কমলার ক্রয়মূল্য = ১.৫×৩=৪.৫ টাকা}

\emph{যেহেতু তিনটি কমলার ক্রয়মূল্যে ২টি বিক্রয় করেন}

\emph{তাই ২টি কমলার বিক্রয়মূল্য = ৪.৫ টাকা}

\emph{১টি কমলার বিক্রয়মূল্য = ৪.৫ ÷ ২ = ২.২৫ টাকা}

\emph{১২টি কমলার বিক্রয়মূল্য = ১২×২.২৫ টাকা = ২৭ টাকা}
\end{quote}

\textbf{একটি মোটর সাইকেল ৩৬০০০ টাকায় বিক্রয় করায় ২০\% ক্ষতি হল। কত টাকায়
বিক্রয় করলে ১৬\% লাভ হত?}

ক) ৫৩০০০ টাকা

খ) ৫০০০০ টাকা

\textbf{\hl{গ) ৫২২০০ টাকা ✔}}

ঘ) ৫৫০০০ টাকা

\textbf{\hl{✔ সঠিক উত্তর: গ) ৫২২০০ টাকা}}

\begin{quote}
\textbf{ব্যাখ্যা:} \emph{২০\% ক্ষতিতে ক্রয়মূল্য ১০০ টাকা হলে বিক্রয়মূল্য =
১০০−২০=৮০ টাকা}

\emph{বিক্রয়মূল্য ৮০ টাকা হলে ক্রয়মূল্য = ১০০ টাকা}

\emph{বিক্রয়মূল্য ১ টাকা হলে ক্রয়মূল্য = (১০০/৮০) টাকা}

\emph{বিক্রয়মূল্য ৩৬০০০ টাকা হলে ক্রয়মূল্য = (১০০×৩৬০০০/৮০)=৪৫০০০ টাকা}

\emph{আবার, ১৬\% লাভে ক্রয়মূল্য ১০০ টাকা হলে বিক্রয়মূল্য = ১০০+১৬=১১৬ টাকা}

\emph{ক্রয়মূল্য ১০০ টাকা হলে বিক্রয়মূল্য = ১১৬ টাকা}

\emph{ক্রয়মূল্য ১ টাকা হলে বিক্রয়মূল্য = (১১৬/১০০) টাকা}

\emph{ক্রয়মূল্য ৪৫০০০ টাকা হলে বিক্রয়মূল্য = (১১৬×৪৫০০০/১০০)=৫২২০০ টাকা}

\emph{শর্ট সমাধান:}

\emph{এই ধরনের সমস্যার ক্ষেত্রে প্রথমে ক্রয়মূল্য বের করার কাজ করতে হবে। তারপর
পরবর্তী ধাপে বিক্রয়মূল্য বের করতে হবে।}

\emph{২০\% ক্ষতিতে ক্রয়মূল্য ১০০ টাকা হলে বিক্রয়মূল্য = ১০০−২০=৮০ টাকা}

\emph{ক্রয়মূল্য = (১০০×৩৬০০০/৮০)=৪৫০০০ টাকা {[}এখানে ক্ষতি হলে \% বিয়োগ করার
পর উল্টিয়ে গুণ করতে হয়{]}}

\emph{আবার, ১৬\% লাভে ক্রয়মূল্য ১০০ টাকা হলে বিক্রয়মূল্য = ১০০+১৬=১১৬ টাকা}

\emph{বিক্রয় করতে হবে = (১১৬×৪৫০০০/১০০)=৫২২০০ টাকায়}
\end{quote}

\textbf{১৯০ টাকায় একটি দ্রব্য বিক্রয় করাতে ৫\% ক্ষতি হল। বিক্রয়মূল্য কত হলে ৫\%
লাভ হত?}

ক) ২০৫ টাকা

খ) ২১৫ টাকা

\textbf{\hl{গ) ২১০ টাকা ✔}}

ঘ) ২২০ টাকা

\textbf{\hl{✔ সঠিক উত্তর: গ) ২১০ টাকা}}

\begin{quote}
\textbf{ব্যাখ্যা:} \emph{৫\% ক্ষতিতে ক্রয়মূল্য ১০০ টাকা হলে বিক্রয়মূল্য =
১০০−৫=৯৫ টাকা}

\emph{বিক্রয়মূল্য ৯৫ টাকা হলে ক্রয়মূল্য = ১০০ টাকা}

\emph{বিক্রয়মূল্য ১ টাকা হলে ক্রয়মূল্য = (১০০/৯৫) টাকা}

\emph{বিক্রয়মূল্য ১৯০ টাকা হলে ক্রয়মূল্য = (১০০×১৯০/৯৫)=২০০ টাকা}

\emph{আবার, ৫\% লাভে ক্রয়মূল্য ১০০ টাকা হলে বিক্রয়মূল্য = ১০০+৫=১০৫ টাকা}

\emph{ক্রয়মূল্য ১০০ টাকা হলে বিক্রয়মূল্য = ১০৫ টাকা}

\emph{ক্রয়মূল্য ১ টাকা হলে বিক্রয়মূল্য = (১০৫/১০০) টাকা}

\emph{ক্রয়মূল্য ২০০ টাকা হলে বিক্রয়মূল্য = (১০৫×২০০/১০০)=২১০ টাকা}

\emph{শর্ট সমাধান:}

\emph{এই ধরনের সমস্যার ক্ষেত্রে প্রথমে ক্রয়মূল্য বের করার কাজ করতে হবে। তারপর
পরবর্তী ধাপে বিক্রয়মূল্য বের করতে হবে।}

\emph{৫\% ক্ষতিতে ক্রয়মূল্য ১০০ টাকা হলে বিক্রয়মূল্য = ১০০−৫=৯৫ টাকা}

\emph{ক্রয়মূল্য = (১০০×১৯০/৯৫)=২০০ টাকা {[}এখানে \% বিয়োগ করার পর উল্টিয়ে গুণ
করতে হয়{]}}

\emph{আবার, ৫ \% লাভে ক্রয়মূল্য ১০০ টাকা হলে বিক্রয়মূল্য = ১০০+৫=১০৫ টাকা}

\emph{বিক্রয় করতে হবে = (১০৫×২০০/১০০)=২১০ টাকায়}
\end{quote}

\textbf{একটি সাইকেল ৭২০০ টাকায় বিক্রয় করায় ১০\% ক্ষতি হয়। কত টাকায় বিক্রয়
করলে ১২\% লাভ হবে?}

ক) ৬৫০০

খ) ৭০০০

\textbf{\hl{গ) ৮৯৬০ ✔}}

ঘ) ৮০০০

\textbf{\hl{✔ সঠিক উত্তর: গ) ৮৯৬০}}

\begin{quote}
\textbf{ব্যাখ্যা:} \emph{১০\% ক্ষতিতে ক্রয়মূল্য ১০০ টাকা হলে বিক্রয়মূল্য =
১০০−১০=৯০ টাকা}

\emph{বিক্রয়মূল্য ৯০ টাকা হলে ক্রয়মূল্য = ১০০ টাকা}

\emph{বিক্রয়মূল্য ১ টাকা হলে ক্রয়মূল্য = (১০০/৯০) টাকা}

\emph{বিক্রয়মূল্য ৭২০০ টাকা হলে ক্রয়মূল্য = (১০০×৭২০০/৯০)=৮০০০ টাকা}

\emph{আবার, ১২\% লাভে ক্রয়মূল্য ১০০ টাকা হলে বিক্রয়মূল্য = ১০০+১২=১১২ টাকা}

\emph{ক্রয়মূল্য ১০০ টাকা হলে বিক্রয়মূল্য = ১১২ টাকা}

\emph{ক্রয়মূল্য ১ টাকা হলে বিক্রয়মূল্য = (১১২/১০০) টাকা}

\emph{ক্রয়মূল্য ৮০০০ টাকা হলে বিক্রয়মূল্য = (১১২×৮০০০/১০০)=৮৯৬০ টাকা}

\emph{শর্ট সমাধান:}

\emph{এই ধরনের সমস্যার ক্ষেত্রে প্রথমে ক্রয়মূল্য বের করার কাজ করতে হবে। তারপর
পরবর্তী ধাপে বিক্রয়মূল্য বের করতে হবে।}

\emph{১০\% ক্ষতিতে ক্রয়মূল্য ১০০ টাকা হলে বিক্রয়মূল্য = ১০০−১০=৯০ টাকা}

\emph{ক্রয়মূল্য = (১০০×৭২০০/৯০)=৮০০০ টাকা {[}এখানে \% বিয়োগ করার পর উল্টিয়ে
গুণ করতে হয়{]}}

\emph{আবার, ১২ \% লাভে ক্রয়মূল্য ১০০ টাকা হলে বিক্রয়মূল্য = ১০০+১২=১১২ টাকা}

\emph{বিক্রয় করতে হবে = (১১২×৮০০০/১০০)=৮৯৬০ টাকায়}
\end{quote}

\textbf{একটি বই ৬৫ টাকায় বিক্রি করায় বিক্রেতার ৩০\% লাভ হয়। কত টাকায় বিক্রয়
করলে ১০\% লাভ হত?}

\textbf{\hl{ক) ৫৫ ✔}}

খ) ৫০

গ) ৪৫

ঘ) কোনোটিই নয়

\textbf{\hl{✔ সঠিক উত্তর: ক) ৫৫}}

\begin{quote}
\textbf{ব্যাখ্যা:} \emph{৩০\% লাভে ক্রয়মূল্য ১০০ টাকা হলে বিক্রয়মূল্য =
১০০+৩০=১৩০ টাকা}

\emph{বিক্রয়মূল্য ১৩০ টাকা হলে ক্রয়মূল্য = ১০০ টাকা}

\emph{বিক্রয়মূল্য ১ টাকা হলে ক্রয়মূল্য = (১০০/১৩০) টাকা}

\emph{বিক্রয়মূল্য ৬৫ টাকা হলে ক্রয়মূল্য = (১০০×৬৫/১৩০)=৫০ টাকা}

\emph{আবার,১০\% লাভে ক্রয়মূল্য ১০০ টাকা হলে বিক্রয়মূল্য = ১০০+১০=১১০ টাকা}

\emph{ক্রয়মূল্য ১০০ টাকা হলে বিক্রয়মূল্য = ১১০ টাকা}

\emph{ক্রয়মূল্য ১ টাকা হলে বিক্রয়মূল্য = (১১০/১০০) টাকা}

\emph{ক্রয়মূল্য ৫০ টাকা হলে বিক্রয়মূল্য = (১১০×৫০/১০০)=৫৫ টাকা}

\emph{বি: দ্র: এই ধরনের সমস্যার ক্ষেত্রে প্রথমে ক্রয়মূল্য বের করার কাজ করতে
হবে।তারপর পরবর্তী ধাপে বিক্রয়মূল্য বের করতে হবে।}

\emph{শর্ট সমাধান:}

\emph{এই ধরনের সমস্যার ক্ষেত্রে প্রথমে ক্রয়মূল্য বের করার কাজ করতে হবে।তারপর
পরবর্তী ধাপে বিক্রয়মূল্য বের করতে হবে।}

\emph{৩০\% লাভে ক্রয়মূল্য ১০০ টাকা হলে বিক্রয়মূল্য = ১০০+৩০=১৩০ টাকা}

\emph{ক্রয়মূল্য = (৬৫×১০০/১৩০)=৫০ টাকা {[} এখানে লাভ হলে \% যোগ করার পর
উল্টিয়ে গুণ করতে হয়{]}}

\emph{আবার, ১০ \% লাভে ক্রয়মূল্য ১০০ টাকা হলে বিক্রয়মূল্য = ১০০+১০=১১০ টাকা}

\emph{বিক্রি করতে হবে = (১১০×৫০/১০০)=৫৫ টাকায়}
\end{quote}

\textbf{৫৩৫ টাকায় একটি জামা বিক্রি করে শতকরা ৭ ভাগ লাভ হয়, জামাটি কত
টাকায় বিক্রি করলে শতকরা ২০ ভাগ ক্ষতি হবে?}

\textbf{\hl{ক) ৪০০ ✔}}

খ) ৪৫০

গ) ৫০০

ঘ) ৫৫০

\textbf{\hl{✔ সঠিক উত্তর: ক) ৪০০}}

\begin{quote}
\textbf{ব্যাখ্যা:} \emph{৭\% লাভে ক্রয়মূল্য ১০০ টাকা হলে বিক্রয়মূল্য = ১০০+৭=১০৭
টাকা}

\emph{বিক্রয়মূল্য ১০৭ টাকা হলে ক্রয়মূল্য = ১০০ টাকা}

\emph{বিক্রয়মূল্য ১ টাকা হলে ক্রয়মূল্য = (১০০/১০৭) টাকা}

\emph{বিক্রয়মূল্য ৫৩৫ টাকা হলে ক্রয়মূল্য = (১০০×৫৩৫/১০৭)=৫০০ টাকা}

\emph{আবার,২০\% ক্ষতিতে ক্রয়মূল্য ১০০ টাকা হলে বিক্রয়মূল্য = ১০০−২০=৮০ টাকা}

\emph{ক্রয়মূল্য ১০০ টাকা হলে বিক্রয়মূল্য = ৮০ টাকা}

\emph{ক্রয়মূল্য ১ টাকা হলে বিক্রয়মূল্য = (৮০/১০০) টাকা}

\emph{ক্রয়মূল্য ৫০০ টাকা হলে বিক্রয়মূল্য = (৮০×৫০০/১০০)=৪০০ টাকা}

\emph{বি: দ্র: এই ধরনের সমস্যার ক্ষেত্রে প্রথমে ক্রয়মূল্য বের করার কাজ করতে
হবে।তারপর পরবর্তী ধাপে বিক্রয়মূল্য বের করতে হবে।}

\emph{শর্ট সমাধান:}

\emph{এই ধরনের সমস্যার ক্ষেত্রে প্রথমে ক্রয়মূল্য বের করার কাজ করতে হবে।তারপর
পরবর্তী ধাপে বিক্রয়মূল্য বের করতে হবে।}

\emph{৭\% লাভে ক্রয়মূল্য ১০০ টাকা হলে বিক্রয়মূল্য = ১০০+৭=১০৭ টাকা}

\emph{ক্রয়মূল্য = (১০০×৫৩৫/১০৭)=৫০০ টাকা {[} এখানে লাভ হলে \% যোগ করার পর
উল্টিয়ে গুণ করতে হয়{]}}

\emph{আবার, ২০\% ক্ষতিতে ক্রয়মূল্য ১০০ টাকা হলে বিক্রয়মূল্য = ১০০−২০=৮০ টাকা}

\emph{বিক্রি করতে হবে = (৮০×৫০০/১০০)=৪০০ টাকায়}
\end{quote}

\textbf{৩৬ টাকা ডজন দরে ক্রয় করে ২০\% লাভে বিক্রয় করা হল, এক কুড়ি কলার
বিক্রয়মূল্য কত?}

ক) ৬০ টাকা

\textbf{\hl{খ) ৭২ টাকা ✔}}

গ) ৬২ টাকা

ঘ) ৭৫ টাকা

\textbf{\hl{✔ সঠিক উত্তর: খ) ৭২ টাকা}}

\begin{quote}
\textbf{ব্যাখ্যা:} \emph{১ ডজন = ১২ টি}

\emph{১২ টি কলার ক্রয়মূল্য = ৩৬ টাকা}

\emph{∴ এক কুড়ি বা ২০ টি কলার ক্রয়মূল্য = ৬০ টাকা}

\emph{২০\% লাভে ক্রয়মূল্য ১০০ টাকা হলে বিক্রয়মূল্য = ১০০+২০=১২০ টাকা}

\emph{ক্রয়মূল্য ১০০ টাকা হলে বিক্রয়মূল্য = ১২০ টাকা}

\emph{ক্রয়মূল্য ১ টাকা হলে বিক্রয়মূল্য = (১২০/১০০) টাকা}

\emph{ক্রয়মূল্য ৬০ টাকা হলে বিক্রয়মূল্য = (১২০×৬০/১০০)=৭২ টাকা}
\end{quote}

\textbf{আব্দুল্লাহ প্রতি ডজন কলা ২১ টাকা দরে ১৫ ডজন কলা এবং ১৪ টাকা দরে ২০
ডজন কলা ক্রয় করে। প্রতি ডজন কলা কি দামে বিক্রয় করলে গড়ে তার ডজন প্রতি ৫ টাকা
লাভ হবে?}

\textbf{\hl{ক) ২২ টাকা দরে ✔}}

খ) ২০ টাকা দরে

গ) ১৮ টাকা দরে

ঘ) ১৫ টাকা দরে

\textbf{\hl{✔ সঠিক উত্তর: ক) ২২ টাকা দরে}}

\begin{quote}
\textbf{ব্যাখ্যা:} \emph{২১ টাকা দরে ১৫ ডজন কলার ক্রয়মূল্য = ৩১৫ টাকা}

\emph{১৪ টাকা দরে ২০ ডজন কলার ক্রয়মূল্য = ২৮০ টাকা}

\emph{∴ ১৫+২০ = ৩৫ ডজন কলার মোট ক্রয়মূল্য = ৩১৫+২৮০=৫৯৫ টাকা}

\emph{তাহলে গড়ে প্রতি ডজন কলার ক্রয়মূল্য = ৫৯৫÷৩৫=১৭ টাকা।}

\emph{অতএব,প্রতি ডজন কলায় ৫ টাকা লাভ করতে হলে প্রতি ডজন কলা ১৭+৫=২২ টাকা
দরে বিক্রি করতে হবে।}
\end{quote}

\textbf{দবির প্রতি ডজন ৩০ টাকা দরে ৮ ডজন এবং ২৫ টাকা দরে ১২ ডজন ডিম ক্রয়
করে প্রতি ডজন ডিম কি দরে বিক্রয় করলে গড়ে তার ডজন প্রতি ৩ টাকা লাভ হবে?}

\textbf{\hl{ক) ৩০ টাকায় ✔}}

খ) ১৫ টাকায়

গ) ৪০ টাকায়

ঘ) ২০ টাকায়

\textbf{\hl{✔ সঠিক উত্তর: ক) ৩০ টাকায়}}

\begin{quote}
\textbf{ব্যাখ্যা:} \emph{৩০ টাকা দরে ৮ ডজন ডিমের ক্রয়মূল্য = ২৪০ টাকা}

\emph{২৫ টাকা দরে ১২ ডজন ডিমের ক্রয়মূল্য = ৩০০ টাকা}

\emph{∴ ৮+১২=২০ ডজন ডিমের মোট ক্রয়মূল্য = ২৪০+৩০০=৫৪০ টাকা}

\emph{তাহলে গড়ে প্রতি ডজন ডিমের ক্রয়মূল্য = ৫৪০÷২০=২৭ টাকা।}

\emph{অতএব,প্রতি ডজন ডিমে ৩ টাকা লাভ করতে হলে প্রতি ডজন ডিম ২৭+৩=৩০ টাকা
দরে বিক্রি করতে হবে।}
\end{quote}

\textbf{এক ব্যক্তি একটি দ্রব্য ১২০০ টাকায় কিনে ১৫\% লাভে বিক্রয় করল; ক্রেতা ঐ
দ্রব্য তৃতীয় এক ব্যক্তির কাছে ৫\% ক্ষতিতে বিক্রয় করল। শেষ বিক্রয়মূল্য কত ছিল?}

ক) ৮৭০ টাকা

খ) ৯৫০ টাকা

\textbf{\hl{গ) ১৩১১ টাকা ✔}}

ঘ) ১২৬০ টাকা

\textbf{\hl{✔ সঠিক উত্তর: গ) ১৩১১ টাকা}}

\begin{quote}
\textbf{ব্যাখ্যা:} \emph{সাধারণ সমাধান:}

\emph{১২০০ টাকার দ্রব্য ১৫\% লাভে বিক্রয়মূল্য = ১৩৮০ টাকা}

\emph{১৩৮০ টাকা ঐ ক্রেতার নিকট হলো ক্রয়মূল্য অর্থাৎ ক্রেতার ক্রয়মূল্য ১৩৮০ টাকা।}

\emph{এখন ঐ ক্রেতা ১৩৮০ টাকার দ্রব্য সে তৃতীয় ব্যক্তির কাছে ৫\% ক্ষতিতে বিক্রয়
করল।}

\emph{১৩৮০ টাকার দ্রব্য ৫\% ক্ষতিতে বিক্রয়মূল্য হয় =
১৩৮০−(১৩৮০×(৫/১০০))=১৩৮০−৬৯=১৩১১ টাকা।}

\emph{অর্থাৎ শেষ বিক্রয়মূল্য ছিল ১৩১১ টাকা।}

\emph{শর্ট সমাধান:}

\emph{প্রথম বিক্রয়মূল্য = ১২০০×(১০০+১৫/১০০)=১২০০×(১১৫/১০০)=১৩৮০ টাকা}

\emph{শেষ বিক্রয়মূল্য = ১৩৮০×(১০০−৫/১০০)=১৩৮০×(৯৫/১০০)=১৩১১ টাকা}
\end{quote}

\textbf{কোন একটি জিনিস নির্মাতা ২০\% লাভে ও খুচরা বিক্রেতা ২০\% লাভে বিক্রয়
করে। যদি ঐ জিনিসের নির্মাণ খরচ ১০০ টাকা হয় তবে খুচরা মূল্য কত?}

ক) ১৪০ টাকা

\textbf{\hl{খ) ১৪৪ টাকা ✔}}

গ) ১২৪ টাকা

ঘ) ১২০ টাকা

\textbf{\hl{✔ সঠিক উত্তর: খ) ১৪৪ টাকা}}

\begin{quote}
\textbf{ব্যাখ্যা:} \emph{সাধারণ সমাধান:}

\emph{২০\% লাভে ক্রয়মূল্য ১০০ টাকা হলে বিক্রয়মূল্য ১২০ টাকা}

\emph{দেওয়া আছে, নির্মাণ খরচ ১০০ টাকা}

\emph{তাই নির্মাতার বিক্রয়মূল্য হয় ১২০ টাকা। ঐ ১২০ টাকা আবার খুচরা বিক্রেতার
ক্রয়মূল্য হয়ে দাঁড়ায়।}

\emph{এখানে মনে রাখতে হবে যে নির্মাতার বিক্রয়মূল্য = খুচরা বিক্রেতার ক্রয়মূল্য}

\emph{অতএব, খুচরা বিক্রেতার ক্রয়মূল্য যখন ১০০ টাকা তখন জিনিসটির খুচরা মূল্য হবে
= ১২০ টাকা}

\emph{খুচরা বিক্রেতার ক্রয়মূল্য যখন ১ টাকা তখন জিনিসটির খুচরা মূল্য হবে =
(১২০/১০০) টাকা}

\emph{খুচরা বিক্রেতার ক্রয়মূল্য যখন ১২০ টাকা তখন জিনিসটির খুচরা মূল্য হবে =
(১২০×১২০/১০০)=১৪৪ টাকা}

\emph{শর্ট সমাধান:}

\emph{নির্মাতার বিক্রয়মূল্য = ১০০×(১২০/১০০)=১২০ টাকা}

\emph{খুচরা মূল্য = ১২০×(১২০/১০০)=১৪৪ টাকা}
\end{quote}

\textbf{এক ব্যক্তি কোন দ্রব্যের ধার্যমূল্যের ৮\% কমিশন দিয়েও ১৫\% লাভ করে। যে
দ্রব্যের ক্রয়মূল্য ২৮০.০০ টাকা তার ধার্যমূল্য কত টাকা?}

ক) ৩২৫ টাকা

\textbf{\hl{খ) ৩৫০ টাকা ✔}}

গ) ৪০০ টাকা

ঘ) ৫৬০ টাকা

\textbf{\hl{✔ সঠিক উত্তর: খ) ৩৫০ টাকা}}

\begin{quote}
\textbf{ব্যাখ্যা:} \emph{সাধারণ সমাধান:}

\emph{১৫\% লাভ করতে হলে}

\emph{ক্রয়মূল্য ১০০ টাকা হলে বিক্রয়মূল্য হবে = ১১৫ টাকা}

\emph{ক্রয়মূল্য ১ টাকা হলে বিক্রয়মূল্য হবে = (১১৫/১০০) টাকা}

\emph{ক্রয়মূল্য ২৮০ টাকা হলে বিক্রয়মূল্য হবে = (১১৫×২৮০/১০০)=৩২২ টাকা}

\emph{ধার্যমূল্যের ৮\% কমিশন দেয়, তাই,}

\emph{বিক্রয়মূল্য ৯২ টাকা হলে ধার্যমূল্য = ১০০ টাকা}

\emph{বিক্রয়মূল্য ১ টাকা হলে ধার্যমূল্য = (১০০/৯২) টাকা}

\emph{বিক্রয়মূল্য ৩২২ টাকা হলে ধার্যমূল্য = (১০০×৩২২/৯২)=৩৫০ টাকা}

\emph{শর্ট সমাধান:}

\emph{ধার্যমূল্য = ((১০০+১৫)/(১০০−৮))×২৮০=(১১৫/৯২)×২৮০=৩৫০ টাকা}
\end{quote}

\textbf{একটি শার্টের মূল্য ২৫০ টাকা লেখা আছে। কিন্তু শার্টটি ২০০ টাকায় বিক্রয়
হলো। লিখিত মূল্য প্রকৃত মূল্য হলে শতকরা কত হ্রাসকৃত মূল্যে বিক্রি হলো?}

ক) ৪০

খ) ৫০

\textbf{\hl{গ) ৮০ ✔}}

ঘ) ২০

\textbf{\hl{✔ সঠিক উত্তর: গ) ৮০}}

\begin{quote}
\textbf{ব্যাখ্যা:} \emph{সাধারণ সমাধান:}

\emph{২৫০ টাকার সাপেক্ষে হ্রাসকৃত মূল্য = ২০০ টাকা}

\emph{১ টাকার সাপেক্ষে হ্রাসকৃত মূল্য = (২০০/২৫০) টাকা}

\emph{১০০ টাকার সাপেক্ষে হ্রাসকৃত মূল্য = (২০০×১০০/২৫০)=৮০ টাকা}

\emph{শর্ট সমাধান:}

\emph{হ্রাসকৃত মূল্যের শতকরা হার = ((২০০/২৫০)×১০০)\%=৮০\%}
\end{quote}

\textbf{একটি শার্টের ধার্যমূল্য ১২.৫০ টাকা। কিন্তু শার্টটি ১০ টাকায় বিক্রয় হলো।
ধার্যমূল্যের সাপেক্ষে কত ছাড় দেওয়া হলো?}

ক) ২\%

খ) ২.৫\%

\textbf{\hl{গ) ২০\% ✔}}

ঘ) ২৫\%

\textbf{\hl{✔ সঠিক উত্তর: গ) ২০\%}}

\begin{quote}
\textbf{ব্যাখ্যা:} \emph{সাধারণ সমাধান:}

\emph{১২.৫০ টাকার সাপেক্ষে বিক্রি মূল্য = ১০ টাকা}

\emph{১ টাকার সাপেক্ষে বিক্রি মূল্য = (১০/১২.৫০) টাকা}

\emph{১০০ টাকার সাপেক্ষে বিক্রি মূল্য = (১০×১০০/১২.৫০)=৮০ টাকা}

\emph{ধার্যমূল্যের সাপেক্ষে ছাড় দেওয়া হলো = ১০০−৮০=২০ টাকা (বা ২০\%)}

\emph{শর্ট সমাধান:}

\emph{ছাড়ের পরিমাণ = ১২.৫০−১০=২.৫০ টাকা}

\emph{ছাড়ের হার = ((২.৫০/১২.৫০)×১০০)\%=২০\%}
\end{quote}

\textbf{৫ টাকায় ২টি কমলা কিনে ৩৫ টাকায় কয়টি কমলা বিক্রয় করলে ৪০\% লাভ
হবে?}

ক) ৮ টি

\textbf{\hl{খ) ১০টি ✔}}

গ) ১২টি

ঘ) ১২টি

\textbf{\hl{✔ সঠিক উত্তর: খ) ১০টি}}

\begin{quote}
\textbf{ব্যাখ্যা:} \emph{সাধারণ সমাধান:}

\emph{৪০\% লাভে}

\emph{ক্রয়মূল্য ১০০ টাকা হলে বিক্রয়মূল্য হবে = ১০০+৪০=১৪০ টাকা}

\emph{ক্রয়মূল্য ১ টাকা হলে বিক্রয়মূল্য হবে = (১৪০/১০০) টাকা}

\emph{ক্রয়মূল্য ৫ টাকা হলে বিক্রয়মূল্য হবে = (১৪০×৫/১০০)=৭ টাকা}

\emph{এখন,}

\emph{৭ টাকায় বিক্রয় করতে হবে = ২টি কমলা}

\emph{১ টাকায় বিক্রয় করতে হবে = (২/৭) টি কমলা}

\emph{৩৫ টাকায় বিক্রয় করতে হবে = (২×৩৫/৭)=১০ টি কমলা}

\emph{শর্ট সমাধান:}

\emph{কমলার বিক্রয় সংখ্যা =
(২×৩৫×১০০/৫×(১০০+৪০))=(৭০০০/৫×১৪০)=(৭০০০/৭০০)=১০ টি}
\end{quote}

\textbf{৫ টাকায় ২টি করে কমলা কিনে ৩৫ টাকায় কয়টা কমলা বিক্রয় করলে x\% লাভ
হবে?}

\textbf{\hl{ক) (১৪০০/১০০+x) টি ✔}}

খ) (৫০০/১০০+x) টি

গ) (৪০০/১০০+x) টি

ঘ) ২০টি

\textbf{\hl{✔ সঠিক উত্তর: ক) (১৪০০/১০০+x) টি}}

\begin{quote}
\textbf{ব্যাখ্যা:} \emph{সাধারণ সমাধান:}

\emph{x\% লাভে}

\emph{ক্রয়মূল্য ১০০ টাকা হলে বিক্রয়মূল্য হবে = ১০০+x টাকা}

\emph{ক্রয়মূল্য ১ টাকা হলে বিক্রয়মূল্য হবে = (১০০+x/১০০) টাকা}

\emph{ক্রয়মূল্য ৫ টাকা হলে বিক্রয়মূল্য হবে = (১০০+x/১০০)×৫=(১০০+x/২০)}

\emph{এখন,}

\emph{(১০০+x/২০) টাকায় বিক্রয় করতে হবে = ২টি কমলা}

\emph{১ টাকায় বিক্রয় করতে হবে = (২০/১০০+x)×২ টি কমলা}

\emph{৩৫ টাকায় বিক্রয় করতে হবে = (২০/১০০+x)×২×৩৫=(১৪০০/১০০+x) টি কমলা}

\emph{শর্ট সমাধান:}

\emph{কমলার সংখ্যা = (২×৩৫×১০০/৫×(১০০+x))=(৭০×১০০/৫(১০০+x))=(১৪০০/১০০+x)
টি}
\end{quote}

\textbf{৪ টাকায় ১টি করে কমলা কিনে ২৪ টাকায় কয়টি কমলা বিক্রয় করলে ২০\% লাভ
হবে?}

\textbf{\hl{ক) ৫টি ✔}}

খ) ৬টি

গ) ৪টি

ঘ) ৩টি

\textbf{\hl{✔ সঠিক উত্তর: ক) ৫টি}}

\begin{quote}
\textbf{ব্যাখ্যা:} \emph{সাধারণ সমাধান:}

\emph{২০\% লাভে}

\emph{ক্রয়মূল্য ১০০ টাকা হলে বিক্রয়মূল্য হবে = ১০০+২০=১২০ টাকা}

\emph{ক্রয়মূল্য ১ টাকা হলে বিক্রয়মূল্য হবে = (১২০/১০০) টাকা}

\emph{ক্রয়মূল্য ৪ টাকা হলে বিক্রয়মূল্য হবে = (১২০×৪/১০০)=(২৪/৫) টাকা}

\emph{এখন,}

\emph{(২৪/৫) টাকায় বিক্রয় করতে হবে = ১ টি কমলা}

\emph{১ টাকায় বিক্রয় করতে হবে = (৫/২৪) টি কমলা}

\emph{২৪ টাকায় বিক্রয় করতে হবে = (৫×২৪/২৪)=৫ টি কমলা}

\emph{শর্ট সমাধান:}

\emph{বিক্রয়কৃত কমলার সংখ্যা =
(১×২৪×১০০/৪×(১০০+২০))=(২৪০০/৪×১২০)=(২৪০০/৪৮০)=৫ টি}
\end{quote}

\textbf{একজন ডিম বিক্রেতা প্রতি ডজন ডিম ১০১ টাকা দরে ৫ ডজন এবং ৯০ টাকা দরে
৬ ডজন ডিম কিনে কত টাকা দরে বিক্রয় করলে তার ডজন প্রতি ৩ টাকা লাভ হবে?}

\textbf{\hl{ক) ৯৮ টাকা ✔}}

খ) ৯৬ টাকা

গ) ৯৫ টাকা

ঘ) ১০০ টাকায়

\textbf{\hl{✔ সঠিক উত্তর: ক) ৯৮ টাকা}}

\begin{quote}
\textbf{ব্যাখ্যা:} \emph{সাধারণ সমাধান:}

\emph{১০১ টাকা দরে ৫ ডজন ডিমের ক্রয়মূল্য = ৫০৫ টাকা}

\emph{৯০ টাকা দরে ৬ ডজন ডিমের ক্রয়মূল্য = ৫৪০ টাকা}

\emph{∴ ৫+৬ = ১১ ডজন ডিমের মোট ক্রয়মূল্য = ৫০৫+৫৪০=১০৪৫ টাকা}

\emph{তাহলে গড়ে প্রতি ডজন ডিমের ক্রয়মূল্য = ১০৪৫÷১১=৯৫ টাকা।}

\emph{অতএব,প্রতি ডজন ডিমে ৩ টাকা লাভ করতে হলে প্রতি ডজন ডিম = ৯৫+৩=৯৮ টাকা
দরে বিক্রি করতে হবে।}

\emph{শর্ট সমাধান:}

\emph{গড় ক্রয়মূল্য = ((১০১×৫)+(৯০×৬)/৫+৬)=(৫০৫+৫৪০/১১)=(১০৪৫/১১)=৯৫ টাকা}

\emph{বিক্রয়মূল্য = ৯৫+৩=৯৮ টাকা}
\end{quote}

\textbf{সোহেল একটি ঘড়ি ৬১২ টাকায় বিক্রয় করায় ১৫\% ক্ষতি হলো। তার উদ্দেশ্য
ছিল ১০\% লাভে ঘড়িটি বিক্রয় করা। ২০\% লাভ করতে হলে ঘড়িটির বিক্রয়মূল্য কত
বাড়াতে হতো?}

ক) ১৫০

খ) ১৪৮

গ) ১৪৬

\textbf{\hl{ঘ) ১৪৪ ✔}}

\textbf{\hl{✔ সঠিক উত্তর: ঘ) ১৪৪}}

\begin{quote}
\textbf{ব্যাখ্যা:} \emph{সাধারণ সমাধান:}

\emph{১৫\% ক্ষতিতে ক্রয়মূল্য ১০০ টাকা হলে বিক্রয়মূল্য = ১০০−১৫=৮৫ টাকা}

\emph{বিক্রয়মূল্য ৮৫ টাকা হলে ক্রয়মূল্য = ১০০ টাকা}

\emph{বিক্রয়মূল্য ১ টাকা হলে ক্রয়মূল্য = (১০০/৮৫) টাকা}

\emph{বিক্রয়মূল্য ৬১২ টাকা হলে ক্রয়মূল্য = (১০০×৬১২/৮৫)=৭২০ টাকা}

\emph{২০\% লাভ করতে হলে ঘড়িটির বিক্রয়মূল্য বাড়াতে হবে =
৭২০×২০\%=৭২০×(২০/১০০)=১৪৪ টাকা}

\emph{বি: দ্র: এই ধরনের সমস্যার ক্ষেত্রে প্রথমে ক্রয়মূল্য বের করার কাজ করতে
হবে।তারপর পরবর্তী ধাপে বিক্রয়মূল্য বের করতে হবে।}

\emph{শর্ট সমাধান:}

\emph{ক্রয়মূল্য = (৬১২×১০০/৮৫)=৭২০ টাকা {[} এখানে ক্ষতি হলে \% বিয়োগ করার পর
উল্টিয়ে গুণ করতে হয়{]}}

\emph{২০\% লাভ করতে হলে বিক্রয়মূল্য বাড়াতে হবে = ৭২০×(২০/১০০)=১৪৪ টাকা}
\end{quote}

\textbf{একজন ব্যবসায়ী ১৩৭৭০ টাকায় একটি চেয়ার বিক্রি করায় ক্রয়মূল্যের উপর ৩৫\%
লাভ হয়। সে যদি চেয়ারটি ৪৫\% লাভে বিক্রয় করত তাহলে তার লাভ কত টাকা হত?}

ক) ১০২০০ টাকা

খ) ১৪৭ ৯০ টাকা

\textbf{\hl{গ) ৪৫৯০ টাকা ✔}}

ঘ) ৪৯৫০ টাকা

ঙ) কোনোটিই নয়

\textbf{\hl{✔ সঠিক উত্তর: গ) ৪৫৯০ টাকা}}

\begin{quote}
\textbf{ব্যাখ্যা:} \emph{সাধারণ সমাধান:}

\emph{৩৫\% লাভে ক্রয়মূল্য ১০০ টাকা হলে বিক্রয়মূল্য = ১০০+৩৫=১৩৫ টাকা}

\emph{বিক্রয়মূল্য ১৩৫ টাকা হলে ক্রয়মূল্য = ১০০ টাকা}

\emph{বিক্রয়মূল্য ১৩৭৭০ টাকা হলে ক্রয়মূল্য = (১০০×১৩৭৭০/১৩৫)=১০২০০ টাকা}

\emph{ক্রয়মূল্যের উপর ৪৫\% লাভ হলে লাভ = ১০২০০×৪৫\%=১০২০০×(৪৫/১০০)=৪৫৯০
টাকা}

\emph{শর্ট সমাধান:}

\emph{ক্রয়মূল্য = (১৩৭৭০×১০০/১৩৫)=১০২০০ টাকা {[} এখানে লাভ হলে \% যোগ করার
পর উল্টিয়ে গুণ করতে হয়{]}}

\emph{ক্রয়মূল্যের উপর ৪৫\% লাভ হলে লাভ = ১০২০০×(৪৫/১০০)=৪৫৯০ টাকা}
\end{quote}

\textbf{টাকায় ৫টি মার্বেল বিক্রয় করায় ১২\% ক্ষতি হয়। ১০\% লাভ করতে হলে
টাকায় কয়টি বিক্রয় করতে হবে?}

\textbf{\hl{ক) ৪টি ✔}}

খ) ৩টি

গ) ২টি

ঘ) কোনটি নয়

\textbf{\hl{✔ সঠিক উত্তর: ক) ৪টি}}

\begin{quote}
\textbf{ব্যাখ্যা:} \emph{সাধারণ সমাধান:}

\emph{১২\% ক্ষতিতে,}

\emph{বিক্রয়মূল্য ১০০−১২=৮৮ টাকা হলে ক্রয়মূল্য = ১০০ টাকা}

\emph{বিক্রয়মূল্য ১ টাকা হলে ক্রয়মূল্য = (১০০/৮৮) টাকা}

\emph{∴ ৫টি মার্বেলের ক্রয়মূল্য = (১০০/৮৮) টাকা}

\emph{১টি মার্বেলের ক্রয়মূল্য = (১০০/৮৮×৫)=(২০/৮৮)=(৫/২২) টাকা}

\emph{আবার, ১০\% লাভে,}

\emph{ক্রয়মূল্য ১০০ টাকা হলে বিক্রয়মূল্য = ১০০+১০=১১০ টাকা}

\emph{১টি মার্বেলের বিক্রয়মূল্য = (৫/২২)×(১১০/১০০)=(১/৪) টাকা}

\emph{∴ (১/৪) টাকায় পাওয়া যায় = ১টি মার্বেল}

\emph{১ টাকায় পাওয়া যায় = ১×৪=৪ টি মার্বেল}

\emph{শর্ট সমাধান:}

\emph{বিক্রয়কৃত মার্বেলের সংখ্যা =
(৫×(১০০−১২)/১০০+১০)=(৫×৮৮/১১০)=(৪৪০/১১০)=৪ টি}
\end{quote}

\textbf{টাকায় ১২টি লেবু বিক্রি করায় ৪\% ক্ষতি হয়। ৪৪\% লাভ করতে হলে টাকায়
কয়টি লেবু বিক্রি করতে হবে?}

\textbf{\hl{ক) ৮টি ✔}}

খ) ৯টি

গ) ১০টি

ঘ) ১১টি

ঙ) কোনটিই নয়

\textbf{\hl{✔ সঠিক উত্তর: ক) ৮টি}}

\begin{quote}
\textbf{ব্যাখ্যা:} \emph{সাধারণ সমাধান:}

\emph{৪\% ক্ষতিতে,}

\emph{বিক্রয়মূল্য ১০০−৪=৯৬ টাকা হলে ক্রয়মূল্য = ১০০ টাকা}

\emph{বিক্রয়মূল্য ১ টাকা হলে ক্রয়মূল্য = (১০০/৯৬) টাকা}

\emph{∴ ১২টি লেবুর ক্রয়মূল্য = (১০০/৯৬) টাকা}

\emph{১টি লেবুর ক্রয়মূল্য = (১০০/৯৬×১২)=(২৫/২৪×১২)=(২৫/২৮৮) টাকা}

\emph{আবার, ৪৪\% লাভে,}

\emph{ক্রয়মূল্য ১০০ টাকা হলে বিক্রয়মূল্য = ১০০+৪৪=১৪৪ টাকা}

\emph{১টি লেবুর বিক্রয়মূল্য = (২৫/২৮৮)×(১৪৪/১০০)=(১/৮) টাকা}

\emph{∴ (১/৮) টাকায় বিক্রি করতে হবে = ১টি লেবু}

\emph{১ টাকায় বিক্রি করতে হবে = ১×৮=৮ টি লেবু}

\emph{শর্ট সমাধান:}

\emph{লেবুর সংখ্যা = (১২×(১০০−৪)/১০০+৪৪)=(১২×৯৬/১৪৪)=(১১৫২/১৪৪)=৮ টি}
\end{quote}

\textbf{টাকায় y টি ফল বিক্রয় করায় y\% ক্ষতি হয়। p\% লাভ করতে হলে টাকায়
কয়টি ফল বিক্রয় করতে হবে?}

ক) (100+y/100p)

\textbf{\hl{খ) (y(100−y)/100+p) ✔}}

গ) (100−y/100+p)

ঘ) (100y/100+p)

\textbf{\hl{✔ সঠিক উত্তর: খ) (y(100−y)/100+p)}}

\begin{quote}
\textbf{ব্যাখ্যা:} \emph{সাধারণ সমাধান:}

\emph{y\% ক্ষতিতে,}

\emph{বিক্রয়মূল্য (100−y) টাকা হলে ক্রয়মূল্য = 100 টাকা}

\emph{বিক্রয়মূল্য 1 টাকা হলে ক্রয়মূল্য = (100/100−y) টাকা}

\emph{∴ y টি ফলের ক্রয়মূল্য = (100/100−y) টাকা}

\emph{1 টি ফলের ক্রয়মূল্য = (100/y(100−y)) টাকা}

\emph{আবার, p\% লাভে,}

\emph{ক্রয়মূল্য 100 টাকা হলে বিক্রয়মূল্য = (100+p) টাকা}

\emph{1 টি ফলের বিক্রয়মূল্য = (100/y(100−y))×(100+p/100)=(100+p/y(100−y))
টাকা}

\emph{∴ (100+p/y(100−y)) টাকায় বিক্রি করতে হবে = 1 টি ফল}

\emph{1 টাকায় বিক্রি করতে হবে = (y(100−y)/100+p) টি ফল}

\emph{শর্ট সমাধান:}

\emph{ফলের সংখ্যা = (y×(100−y)/100+p)=(y(100−y)/100+p) টি}
\end{quote}

\textbf{প্রতি ডজন কলা ৪৮ টাকায় কিনে ৫০ টাকায় কয়টি কলা বিক্রি করলে ২৫\% লাভ
হয়?}

\textbf{\hl{ক) ১০টি ✔}}

খ) ১২টি

গ) ১৫টি

ঘ) ১৩টি

\textbf{\hl{✔ সঠিক উত্তর: ক) ১০টি}}

\begin{quote}
\textbf{ব্যাখ্যা:} \emph{সাধারণ সমাধান:}

\emph{১ ডজন = ১২টি}

\emph{১২টি কলার ক্রয়মূল্য = ৪৮ টাকা}

\emph{∴ ১টি কলার ক্রয়মূল্য = (৪৮/১২)=৪ টাকা}

\emph{২৫\% লাভে,}

\emph{ক্রয়মূল্য ১০০ টাকা হলে বিক্রয়মূল্য = ১২৫ টাকা}

\emph{ক্রয়মূল্য ৪ টাকা হলে ১টি কলার বিক্রয়মূল্য = (১২৫×৪/১০০)=৫ টাকা}

\emph{∴ ৫ টাকায় বিক্রি করতে হয় = ১টি কলা}

\emph{১ টাকায় বিক্রি করতে হয় = (১/৫) টি কলা}

\emph{৫০ টাকায় বিক্রি করতে হয় = (১×৫০/৫)=১০ টি কলা}

\emph{শর্ট সমাধান:}

\emph{১টি কলার ক্রয়মূল্য = (৪৮/১২)=৪ টাকা}

\emph{২৫\% লাভে ১টি কলার বিক্রয়মূল্য = ৪+(৪×(২৫/১০০))=৪+১=৫ টাকা}

\emph{৫০ টাকায় বিক্রয়কৃত কলার সংখ্যা = (৫০/৫)=১০ টি}
\end{quote}

\textbf{১০ টাকায় ১ হালি লেবু কিনে ৬০ টাকায় কত হালি লেবু বিক্রয় করলে ২০\%
লাভ হবে?}

\textbf{\hl{ক) ৫ ✔}}

খ) ৬

গ) ৭

ঘ) ৪

\textbf{\hl{✔ সঠিক উত্তর: ক) ৫}}

\begin{quote}
\textbf{ব্যাখ্যা:} \emph{সাধারণ সমাধান:}

\emph{২০\% লাভে,}

\emph{ক্রয়মূল্য ১০০ টাকা হলে বিক্রয়মূল্য ১২০ টাকা}

\emph{ক্রয়মূল্য ১ টাকা হলে বিক্রয়মূল্য = (১২০/১০০) টাকা}

\emph{ক্রয়মূল্য ১০ টাকা হলে বিক্রয়মূল্য = (১২০×১০/১০০)=১২ টাকা}

\emph{এখন, ২০\% লাভ করতে,}

\emph{১২ টাকায় বিক্রি করতে হবে = ১ হালি লেবু}

\emph{১ টাকায় বিক্রি করতে হবে = (১/১২) হালি লেবু}

\emph{∴ ৬০ টাকায় বিক্রি করতে হবে = (১×৬০/১২)=৫ হালি লেবু।}

\emph{শর্ট সমাধান:}

\emph{১ হালি লেবুর বিক্রয়মূল্য = ১০×(১০০+২০/১০০)=১০×(১২০/১০০)=১২ টাকা}

\emph{∴ ৬০ টাকায় বিক্রি করতে হবে = (৬০/১২)=৫ হালি লেবু}
\end{quote}

\textbf{এক ডজন আম ৬০ টাকায় ক্রয় করে হালি কত টাকায় বিক্রয় করলে ১০\% লাভ
হবে?}

\textbf{\hl{ক) ২২ টাকা ✔}}

খ) ২৫ টাকা

গ) ১৫ টাকা

ঘ) ১২ টাকা

\textbf{\hl{✔ সঠিক উত্তর: ক) ২২ টাকা}}

\begin{quote}
\textbf{ব্যাখ্যা:} \emph{সাধারণ সমাধান:}

\emph{১ ডজন বা ৩ হালি আমের ক্রয়মূল্য = ৬০ টাকা}

\emph{∴ ১ হালি আমের ক্রয়মূল্য = (৬০/৩)=২০ টাকা}

\emph{১০\% লাভে,}

\emph{ক্রয়মূল্য ১০০ টাকা হলে বিক্রয়মূল্য = ১১০ টাকা}

\emph{ক্রয়মূল্য ১ টাকা হলে বিক্রয়মূল্য = (১১০/১০০) টাকা}

\emph{ক্রয়মূল্য ২০ টাকা হলে বিক্রয়মূল্য = (১১০×২০/১০০)=২২ টাকা}

\emph{শর্ট সমাধান:}

\emph{১ হালি আমের ক্রয়মূল্য = ২০ টাকা}

\emph{১০\% লাভে ১ হালির বিক্রয়মূল্য = ২০×(১০০+১০/১০০)=২০×(১১০/১০০)=২২ টাকা}
\end{quote}

\textbf{টাকায় এক ডজন কলা বিক্রি করায় ২০\% ক্ষতি হয়।৬০\% লাভ করতে হলে টাকায়
কটি কলা বিক্রি করতে হবে?}

\textbf{\hl{ক) ৬টি ✔}}

খ) ৭টি

গ) ৮টি

ঘ) ৯টি

\textbf{\hl{✔ সঠিক উত্তর: ক) ৬টি}}

\begin{quote}
\textbf{ব্যাখ্যা:} \emph{সাধারণ সমাধান:}

\emph{ক্রয়মূল্য ১০০ টাকা হলে,}

\emph{২০\% ক্ষতিতে বিক্রয়মূল্য = ১০০−২০=৮০ টাকা}

\emph{৬০\% লাভে বিক্রয়মূল্য = ১০০+৬০=১৬০ টাকা}

\emph{পূর্বের বিক্রয়মূল্য ৮০ টাকা হলে নতুন বিক্রয়মূল্য = ১৬০ টাকা}

\emph{পূর্বের বিক্রয়মূল্য ১ টাকা হলে নতুন বিক্রয়মূল্য = (১৬০/৮০)=২ টাকা}

\emph{∴ ২ টাকায় বিক্রি করতে হবে = ১২টি কলা}

\emph{১ টাকায় বিক্রি করতে হবে = (১২/২)=৬ টি কলা}

\emph{শর্ট সমাধান:}

\emph{বিক্রয়কৃত কলার সংখ্যা = (১২×(১০০−২০)/১০০+৬০)=(১২×৮০/১৬০)=৬ টি}
\end{quote}

\textbf{একজন বিক্রেতা একটি ঘড়ি ১৬০০ টাকায় বিক্রি করলে ২০\% লোকসান হয়।যদি
সে ১০\% লাভ করতে চায়, তবে ঘড়িটি কত টাকায় বিক্রি করতে হবে?}

ক) ২০০০

খ) ২১০০

\textbf{\hl{গ) ২২০০ ✔}}

ঘ) কোনোটিই নয়

\textbf{\hl{✔ সঠিক উত্তর: গ) ২২০০}}

\begin{quote}
\textbf{ব্যাখ্যা:} \emph{সাধারণ সমাধান:}

\emph{২০\% লোকসানে বিক্রয়মূল্য = ১০০−২০=৮০ টাকা}

\emph{বিক্রয়মূল্য ৮০ টাকা হলে ক্রয়মূল্য = ১০০ টাকা}

\emph{বিক্রয়মূল্য ১৬০০ টাকা হলে ক্রয়মূল্য = (১০০×১৬০০/৮০)=২০০০ টাকা}

\emph{এখন ১০\% লাভে,}

\emph{ক্রয়মূল্য ১০০ টাকা হলে বিক্রয়মূল্য = ১১০ টাকা}

\emph{ক্রয়মূল্য ২০০০ টাকা হলে বিক্রয়মূল্য = (১১০×২০০০/১০০)=২২০০ টাকা}

\emph{শর্ট সমাধান:}

\emph{ক্রয়মূল্য = (১৬০০×১০০/৮০)=২০০০ টাকা}

\emph{১০\% লাভে বিক্রয়মূল্য = ২০০০×(১১০/১০০)=২২০০ টাকা}
\end{quote}

\textbf{একটি ঘড়ি ৬১২ টাকায় বিক্রি করলে ১৫\% ক্ষতি হয়।কত টাকায় বিক্রি করলে
১০\% লাভ হবে?}

ক) ৬৯২ টাকা

\textbf{\hl{খ) ৭৯২ টাকা ✔}}

গ) ৫৯২ টাকা

\textbf{\hl{✔ সঠিক উত্তর: খ) ৭৯২ টাকা}}

\begin{quote}
\textbf{ব্যাখ্যা:} \emph{সাধারণ সমাধান:}

\emph{১৫\% ক্ষতিতে বিক্রয়মূল্য = ১০০−১৫=৮৫ টাকা}

\emph{বিক্রয়মূল্য ৮৫ টাকা হলে ক্রয়মূল্য = ১০০ টাকা}

\emph{বিক্রয়মূল্য ৬১২ টাকা হলে ক্রয়মূল্য = (১০০×৬১২/৮৫)=৭২০ টাকা}

\emph{এখন ১০\% লাভে,}

\emph{ক্রয়মূল্য ১০০ টাকা হলে বিক্রয়মূল্য = ১১০ টাকা}

\emph{ক্রয়মূল্য ৭২০ টাকা হলে বিক্রয়মূল্য = (১১০×৭২০/১০০)=৭৯২ টাকা}

\emph{শর্ট সমাধান:}

\emph{ক্রয়মূল্য = (৬১২×১০০/৮৫)=৭২০ টাকা}

\emph{১০\% লাভে বিক্রয়মূল্য = ৭২০×(১১০/১০০)=৭৯২ টাকা}
\end{quote}

\textbf{টাকায় ৩ টি করে লেবু কিনে টাকায় ২ টি করে বিক্রয় করলে শতকরা কত লাভ
হবে?}

ক) ৪০\%

\textbf{\hl{খ) ৫০\% ✔}}

গ) ৬০\%

ঘ) ৭০\%

\textbf{\hl{✔ সঠিক উত্তর: খ) ৫০\%}}

\begin{quote}
\textbf{ব্যাখ্যা:} \emph{সাধারণ সমাধান:}

\emph{৩টি লেবুর ক্রয়মূল্য = ১ টাকা}

\emph{∴ ১টি লেবুর ক্রয়মূল্য = (১/৩) টাকা}

\emph{২টি লেবুর বিক্রয়মূল্য = ১ টাকা}

\emph{∴ ১টি লেবুর বিক্রয়মূল্য = (১/২) টাকা}

\emph{১টি লেবুতে লাভ = (১/২)−(১/৩)=(১/৬) টাকা}

\emph{∴ শতকরা লাভ = (((১/৬)/(১/৩))×১০০)\%=((৩/৬)×১০০)\%=৫০\%}

\emph{শর্ট সমাধান:}

\emph{লাভের হার = ((ক্রয়কৃত−বিক্রয়কৃত/বিক্রয়কৃত)×১০০)\%=((৩−২/২)×১০০)\%=৫০\%}
\end{quote}

\textbf{টাকায় ১২ টি লেবু বিক্রি করায় ৪\% ক্ষতি হয়।৪৪\% লাভ করতে হলে টাকায়
কটি লেবু বিক্রি করতে হবে?}

\textbf{\hl{ক) ৮টি ✔}}

খ) ৯টি

গ) ১০ টি

ঘ) ১১ টি

\textbf{\hl{✔ সঠিক উত্তর: ক) ৮টি}}

\begin{quote}
\textbf{ব্যাখ্যা:} \emph{সাধারণ সমাধান:}

\emph{৪\% ক্ষতিতে বিক্রয়মূল্য ৯৬ টাকা হলে ক্রয়মূল্য = ১০০ টাকা}

\emph{∴ ১২টি লেবুর ক্রয়মূল্য = (১০০/৯৬)=(২৫/২৪) টাকা}

\emph{৪৪\% লাভে বিক্রয়মূল্য = ১৪৪ টাকা}

\emph{ক্রয়মূল্য ১০০ টাকা হলে বিক্রয়মূল্য = ১৪৪ টাকা}

\emph{ক্রয়মূল্য (২৫/২৪) টাকা হলে বিক্রয়মূল্য = (১৪৪×২৫/১০০×২৪)=(৩/২) টাকা}

\emph{(৩/২) টাকায় বিক্রি করতে হবে = ১২টি লেবু}

\emph{∴ ১ টাকায় বিক্রি করতে হবে = (১২×২/৩)=৮ টি লেবু}

\emph{শর্ট সমাধান:}

\emph{বিক্রয়কৃত লেবুর সংখ্যা = (১২×(১০০−৪)/১০০+৪৪)=(১২×৯৬/১৪৪)=৮ টি}
\end{quote}

\textbf{১০ টাকার ১ হালি লেবু কিনে ৬০ টাকায় কত হালি লেবু বিক্রি করলে ২০\%
লাভ হবে?}

\textbf{\hl{ক) ৫ ✔}}

খ) ৬

গ) ৭

ঘ) ৪

\textbf{\hl{✔ সঠিক উত্তর: ক) ৫}}

\begin{quote}
\textbf{ব্যাখ্যা:} \emph{সাধারণ সমাধান:}

\emph{২০\% লাভে,}

\emph{ক্রয়মূল্য ১০০ টাকা হলে বিক্রয়মূল্য = ১২০ টাকা}

\emph{ক্রয়মূল্য ১০ টাকা হলে বিক্রয়মূল্য = (১২০×১০/১০০)=১২ টাকা}

\emph{১২ টাকায় বিক্রি করতে হবে = ১ হালি লেবু}

\emph{∴ ৬০ টাকায় বিক্রি করতে হবে = (৬০/১২)=৫ হালি লেবু}

\emph{শর্ট সমাধান:}

\emph{১ হালি লেবুর ২০\% লাভে বিক্রয়মূল্য = ১০×(১২০/১০০)=১২ টাকা}

\emph{∴ ৬০ টাকায় মোট লেবুর পরিমাণ = (৬০/১২)=৫ হালি}
\end{quote}

\textbf{৪ টাকায় ১ টি করে কমলা কিনে ২৪ টাকায় কয়টি কমলা বিক্রি করলে ২০\%
লাভ হবে?}

\textbf{\hl{ক) ৫টি ✔}}

খ) ৬টি

গ) ৪টি

ঘ) ৩টি

\textbf{\hl{✔ সঠিক উত্তর: ক) ৫টি}}

\begin{quote}
\textbf{ব্যাখ্যা:} \emph{সাধারণ সমাধান:}

\emph{২০\% লাভে,}

\emph{ক্রয়মূল্য ১০০ টাকা হলে বিক্রয়মূল্য = ১২০ টাকা}

\emph{ক্রয়মূল্য ৪ টাকা হলে বিক্রয়মূল্য = (১২০×৪/১০০)=(২৪/৫) টাকা}

\emph{(২৪/৫) টাকায় বিক্রি করতে হবে = ১টি কমলা}

\emph{∴ ২৪ টাকায় বিক্রি করতে হবে = (১×৫×২৪/২৪)=৫ টি কমলা}

\emph{শর্ট সমাধান:}

\emph{কমলা বিক্রয়ের সংখ্যা = (১×২৪×১০০/৪×(১০০+২০))=(২৪০০/৪৮০)=৫ টি}
\end{quote}

\textbf{৫ টাকায় ২ টি করে কমলা কিনে ৩৫ টাকায় কয়টি কমলা বিক্রি করলে ৪০\%
লাভ হবে?}

ক) ৭টি

\textbf{\hl{খ) ১০ টি ✔}}

গ) ১২ টি

ঘ) ১৫টি

\textbf{\hl{✔ সঠিক উত্তর: খ) ১০ টি}}

\begin{quote}
\textbf{ব্যাখ্যা:} \emph{সাধারণ সমাধান:}

\emph{৪০\% লাভে,}

\emph{ক্রয়মূল্য ১০০ টাকা হলে বিক্রয়মূল্য = ১৪০ টাকা}

\emph{ক্রয়মূল্য ৫ টাকা হলে বিক্রয়মূল্য = (১৪০×৫/১০০)=৭ টাকা}

\emph{৭ টাকায় বিক্রি করতে হবে = ২টি কমলা}

\emph{∴ ৩৫ টাকায় বিক্রি করতে হবে = (২×৩৫/৭)=১০ টি কমলা}

\emph{শর্ট সমাধান:}

\emph{কমলা বিক্রয়ের সংখ্যা = (২×৩৫×১০০/৫×(১০০+৪০))=(৭০০০/৭০০)=১০ টি}
\end{quote}

\textbf{টাকায় ১০ টি দরে লেবু ক্রয় করে ৮ টি দরে বিক্রয় করলে শতকরা কত লাভ
হবে?}

ক) ২০\%

\textbf{\hl{খ) ২৫\% ✔}}

গ) ১০\%

ঘ) ১৫\%

\textbf{\hl{✔ সঠিক উত্তর: খ) ২৫\%}}

\begin{quote}
\textbf{ব্যাখ্যা:} \emph{সাধারণ সমাধান:}

\emph{১০টি লেবুর ক্রয়মূল্য = ১ টাকা}

\emph{∴ ১টি লেবুর ক্রয়মূল্য = (১/১০) টাকা}

\emph{৮টি লেবুর বিক্রয়মূল্য = ১ টাকা}

\emph{∴ ১টি লেবুর বিক্রয়মূল্য = (১/৮) টাকা}

\emph{১টি লেবুতে লাভ = (১/৮)−(১/১০)=(৫−৪/৪০)=(১/৪০) টাকা}

\emph{∴ শতকরা লাভ = (((১/৪০)/(১/১০))×১০০)\%=((১০/৪০)×১০০)\%=২৫\%}

\emph{শর্ট সমাধান:}

\emph{টাকার পরিমাণ সমান হলে,}

\emph{৮টি লেবু বিক্রি করে লাভ হয় = ১০−৮=২ টি লেবু}

\emph{∴ শতকরা লাভ = ((২/৮)×১০০)\%=২৫\%}
\end{quote}

\textbf{একটি কোম্পানি দিনে প্রথম ১০০০ টাকা বিক্রির উপর ৫\% লাভ করে এবং ১০০০
টাকার অতিরিক্ত বিক্রির উপর ৪\% লাভ করে। দিনে মোট ৬০০০ টাকার জিনিস বিক্রি
হলে, কোম্পানি সর্বমোট কত টাকা লাভ করে?}

ক) ২১০ টাকা

খ) ২২০ টাকা

গ) ২২৫ টাকা

\textbf{\hl{ঘ) ২৫০ টাকা ✔}}

\textbf{\hl{✔ সঠিক উত্তর: ঘ) ২৫০ টাকা}}

\begin{quote}
\textbf{ব্যাখ্যা:} \emph{সাধারণ সমাধান:}

\emph{প্রথম ১০০০ টাকা বিক্রির উপর লাভ = ১০০০×(৫/১০০)=৫০ টাকা}

\emph{১০০০ টাকার অতিরিক্ত বিক্রির পরিমাণ = ৬০০০−১০০০=৫০০০ টাকা}

\emph{অবশিষ্ট ৫০০০ টাকা বিক্রির উপর লাভ = ৫০০০×(৪/১০০)=২০০ টাকা}

\emph{∴ কোম্পানি সর্বমোট লাভ করে = ৫০+২০০=২৫০ টাকা}

\emph{শর্ট সমাধান:}

\emph{মোট লাভ = (১০০০×৫\%)+\{(৬০০০−১০০০)×৪\%\}=৫০+(৫০০০×৪\%)=৫০+২০০=২৫০
টাকা}
\end{quote}

\textbf{দুইটি গরু একই মূল্যে বিক্রি করলে একটিতে লাভ হয় ১০ টাকা, অন্যটিতে ক্ষতি
হয় ১২ টাকা। শতকরা মোট কত লাভ বা ক্ষতি হলো?}

ক) শতকরা ১ ভাগ লাভ

খ) শতকরা ২ ভাগ লাভ

গ) শতকরা ২ ভাগ ক্ষতি

\textbf{\hl{ঘ) শতকরা ১ ভাগ ক্ষতি ✔}}

\textbf{\hl{✔ সঠিক উত্তর: ঘ) শতকরা ১ ভাগ ক্ষতি}}

\begin{quote}
\textbf{ব্যাখ্যা:} \emph{সাধারণ সমাধান:}

\emph{ধরি, প্রতিটি গরুর বিক্রয়মূল্য ১০০ টাকা}

\emph{∴ ২টি গরুর মোট বিক্রয়মূল্য = ১০০+১০০=২০০ টাকা}

\emph{প্রথম গরুর ক্রয়মূল্য = ১০০−১০=৯০ টাকা}

\emph{দ্বিতীয় গরুর ক্রয়মূল্য = ১০০+১২=১১২ টাকা}

\emph{∴ ২টি গরুর মোট ক্রয়মূল্য = ৯০+১১২=২০২ টাকা}

\emph{এখানে মোট ক্ষতি = ২০২−২০০=২ টাকা}

\emph{∴ ক্ষতির হার = ((২/২০০)×১০০)\%=১\%}

\emph{∴ শতকরা ১ ভাগ ক্ষতি হলো।}

\emph{শর্ট সমাধান:}

\emph{মোট ক্রয়মূল্য = (১০০−১০)+(১০০+১২)=২০২ টাকা}

\emph{মোট বিক্রয়মূল্য = ২০০ টাকা}

\emph{∴ ক্ষতি = ২০২−২০০=২ টাকা, যা ২০০ টাকায় ক্ষতি ২ টাকা বা ১০০ টাকায় ১
টাকা (১\% ক্ষতি)।}
\end{quote}

\textbf{৫০ টাকায় ৬ টি করে আম কিনে ৫০ টাকায় ৫ টি দরে বিক্রয় করলে শতকরা কত
লাভ হবে?}

ক) ১০\%

খ) ১২\%

\textbf{\hl{গ) ২০\% ✔}}

ঘ) ২৫\%

\textbf{\hl{✔ সঠিক উত্তর: গ) ২০\%}}

\begin{quote}
\textbf{ব্যাখ্যা:} \emph{সাধারণ সমাধান:}

\emph{৬টি আমের ক্রয়মূল্য = ৫০ টাকা}

\emph{∴ ১টি আমের ক্রয়মূল্য = (৫০/৬) টাকা}

\emph{৫টি আমের বিক্রয়মূল্য = ৫০ টাকা}

\emph{∴ ১টি আমের বিক্রয়মূল্য = (৫০/৫)=১০ টাকা}

\emph{১টি আমে লাভ = ১০−(৫০/৬)=(৬০−৫০/৬)=(১০/৬) টাকা}

\emph{∴ শতকরা লাভ = (((১০/৬)/(৫০/৬))×১০০)\%=((১০/৫০)×১০০)\%=২০\%}

\emph{শর্ট সমাধান:}

\emph{টাকার পরিমাণ সমান হওয়ায়,}

\emph{৫টি আম বিক্রি করে লাভ হয় = ৬−৫=১ টি আম}

\emph{∴ শতকরা লাভ = ((১/৫)×১০০)\%=২০\%}
\end{quote}

\textbf{একটি কলম ৫০ টাকায় ক্রয় করে ৫৬ টাকায় বিক্রয় করা হলো।এতে শতকরা কত
লাভ হলো?}

ক) ১০\%

\textbf{\hl{খ) ১২\% ✔}}

গ) ২০\%

ঘ) ২৫\%

\textbf{\hl{✔ সঠিক উত্তর: খ) ১২\%}}

\begin{quote}
\textbf{ব্যাখ্যা:} \emph{সাধারণ সমাধান:}

\emph{ক্রয়মূল্য = ৫০ টাকা}

\emph{বিক্রয়মূল্য = ৫৬ টাকা}

\emph{∴ লাভ = ৫৬−৫০=৬ টাকা}

\emph{৫০ টাকায় লাভ হয় = ৬ টাকা}

\emph{১ টাকায় লাভ হয় = (৬/৫০) টাকা}

\emph{∴ ১০০ টাকায় লাভ হয় = (৬×১০০/৫০)=১২ টাকা (অর্থাৎ ১২\%)}

\emph{শর্ট সমাধান:}

\emph{শতকরা লাভ = ((লাভ/ক্রয়মূল্য)×১০০)\%=((৬/৫০)×১০০)\%=১২\%}
\end{quote}

\textbf{এক কুড়ি কমলা ৫০ টাকায় ক্রয় করে এক ডজন কমলা ৩৬ টাকায় বিক্রয় করা
হলো।শতকরা কত লাভ হলো?}

ক) ১০\%

\textbf{\hl{খ) ২০\% ✔}}

গ) ২৫\%

ঘ) ৩০\%

\textbf{\hl{✔ সঠিক উত্তর: খ) ২০\%}}

\begin{quote}
\textbf{ব্যাখ্যা:} \emph{সাধারণ সমাধান:}

\emph{১ কুড়ি = ২০টি}

\emph{২০টি কমলার ক্রয়মূল্য = ৫০ টাকা}

\emph{∴ ১টি কমলার ক্রয়মূল্য = (৫০/২০)=২.৫ টাকা}

\emph{১ ডজন = ১২টি}

\emph{১২টি কমলার বিক্রয়মূল্য = ৩৬ টাকা}

\emph{∴ ১টি কমলার বিক্রয়মূল্য = (৩৬/১২)=৩ টাকা}

\emph{১টি কমলাতে লাভ = ৩−২.৫=০.৫ টাকা}

\emph{∴ শতকরা লাভ = ((০.৫/২.৫)×১০০)\%=((১/৫)×১০০)\%=২০\%}

\emph{শর্ট সমাধান:}

\emph{প্রতিটির ক্রয়মূল্য = ২.৫ টাকা, বিক্রয়মূল্য = ৩ টাকা}

\emph{শতকরা লাভ = ((৩−২.৫/২.৫)×১০০)\%=২০\%}
\end{quote}

\textbf{টাকায় ৩ টি জিনিস ক্রয় করে, টাকায় ২ টি জিনিস বিক্রয় করলে শতকরা লাভ
হবে-}

ক) ৩০\%

খ) ১৫\%

\textbf{\hl{গ) ৫০\% ✔}}

ঘ) ৩৫\%

\textbf{\hl{✔ সঠিক উত্তর: গ) ৫০\%}}

\begin{quote}
\textbf{ব্যাখ্যা:} \emph{সাধারণ সমাধান:}

\emph{৩টি জিনিসের ক্রয়মূল্য = ১ টাকা}

\emph{∴ ১টি জিনিসের ক্রয়মূল্য = (১/৩) টাকা}

\emph{২টি জিনিসের বিক্রয়মূল্য = ১ টাকা}

\emph{∴ ১টি জিনিসের বিক্রয়মূল্য = (১/২) টাকা}

\emph{১টি জিনিসে লাভ = (১/২)−(১/৩)=(১/৬) টাকা}

\emph{∴ শতকরা লাভ = (((১/৬)/(১/৩))×১০০)\%=((৩/৬)×১০০)\%=৫০\%}

\emph{শর্ট সমাধান:}

\emph{টাকার পরিমাণ সমান হলে,}

\emph{২টি জিনিস বিক্রি করে লাভ হয় = ৩−২=১ টি}

\emph{∴ শতকরা লাভ = ((১/২)×১০০)\%=৫০\%}
\end{quote}

\textbf{টাকায় ৪ টি এবং টাকায় ৬ টি দরে সমান সংখ্যক আমড়া কিনে এক ব্যক্তি
টাকায় ৫ টি করে বিক্রয় করলে তার লাভ বা ক্ষতির পরিমাণ কত?}

ক) ক্ষতি ১\%

খ) লাভ ১\%

\textbf{\hl{গ) ক্ষতি ৪\% ✔}}

ঘ) কোন লাভ লোকসান নেই

\textbf{\hl{✔ সঠিক উত্তর: গ) ক্ষতি ৪\%}}

\begin{quote}
\textbf{ব্যাখ্যা:} \emph{সাধারণ সমাধান:}

\emph{১টি আমড়ার ১ম প্রকার ক্রয়মূল্য = (১/৪) টাকা}

\emph{১টি আমড়ার ২য় প্রকার ক্রয়মূল্য = (১/৬) টাকা}

\emph{∴ ১+১=২ টি আমড়ার মোট ক্রয়মূল্য = (১/৪)+(১/৬)=(৩+২/১২)=(৫/১২) টাকা}

\emph{আবার, ৫টি আমড়ার বিক্রয়মূল্য = ১ টাকা}

\emph{∴ ২টি আমড়ার বিক্রয়মূল্য = (২/৫) টাকা}

\emph{এখানে ক্ষতি = (৫/১২)−(২/৫)=(২৫−২৪/৬০)=(১/৬০) টাকা}

\emph{∴ শতকরা ক্ষতি = (((১/৬০)/(৫/১২))×১০০)\%=((১২×১০০/৬০×৫))\%=৪\%}

\emph{শর্ট সমাধান:}

\emph{দুই প্রকারের ক্রয় কটির সংখ্যা x=৪ ও y=৬ এবং বিক্রয় সংখ্যা z=৫ হলে, গড়
বিক্রয় সংখ্যা z=(x+y/২)=(৪+৬/২)=৫ হওয়ায় সর্বদা ক্ষতি হয়।}

\emph{ক্ষতির পরিমাণ =
(((x−y/x+y)))\^{}২×১০০\%=(((৬−৪/৬+৪)))\^{}২×১০০\%=(((২/১০)))\^{}২×১০০\%=(৪/১০০)×১০০\%=৪\%}
\end{quote}

\textbf{একজন কমলা বিক্রেতা প্রতি শত কমলা ১০০০ টাকায় কিনে ১২০০ টাকায় বিক্রি
করলে তার লাভ কত হবে?}

ক) ১২০ টাকা

\textbf{\hl{খ) ২০০ টাকা ✔}}

গ) ২৪০ টাকা

ঘ) ৩০০ টাকা

\textbf{\hl{✔ সঠিক উত্তর: খ) ২০০ টাকা}}

\begin{quote}
\textbf{ব্যাখ্যা:} \emph{ক্রয়মূল্য ১০০০ টাকা এবং বিক্রয়মূল্য ১২০০ টাকা।}

\emph{∴ লাভ = (১২০০ − ১০০০) টাকা = ২০০ টাকা।}
\end{quote}

\textbf{১টি দ্রব্য ৩৮০ টাকায় বিক্রয় করায় ২০ টাকা ক্ষতি হলো। ক্ষতির শতকরা হার
কত?}

ক) ৪\%

খ) ৬\%

\textbf{\hl{গ) ৫\% ✔}}

ঘ) ৭\%

\textbf{\hl{✔ সঠিক উত্তর: গ) ৫\%}}

\begin{quote}
\textbf{ব্যাখ্যা:} \emph{দ্রব্যটির ক্রয়মূল্য (৩৮০ + ২০) বা ৪০০ টাকা}

\emph{এখন, ৪০০ টাকায় ক্ষতি হয় ২০ টাকা}

\emph{∴ ক্ষতির হার (২০/৪০০)×১০০\% = ৫\%}
\end{quote}

\textbf{100 টাকায় 10 টি ডিম কিনে 100 টাকায় 8 টি ডিম বিক্রয় করলে শতকরা লাভ
কত হবে?}

ক) 16\%

খ) 20\%

\textbf{\hl{গ) 25\% ✔}}

ঘ) 28\%

\textbf{\hl{✔ সঠিক উত্তর: গ) 25\%}}

\begin{quote}
\textbf{ব্যাখ্যা:} \emph{টাকা সমান থাকলে,}

\emph{শতকরা লাভ = (ক্রয়সংখ্যা−বিক্রয়সংখ্যা/বিক্রয়সংখ্যা) × ১০০\%∴ শতকরা লাভ =
(১০−৮/৮) × ১০০\% = (২/৮) × ১০০ \% = ২৫\%}
\end{quote}

\textbf{৮ টি প্যান্টের বিক্রয় মূল্য ১০টি প্যান্টের ক্রয় মূল্যের সমান হলে শতকরা লাভ
কত?}

ক) ১৫\%

খ) ২০\%

\textbf{\hl{গ) ২৫\% ✔}}

ঘ) ৩০\%

\textbf{\hl{✔ সঠিক উত্তর: গ) ২৫\%}}

\begin{quote}
\textbf{ব্যাখ্যা:} \emph{মূল সমাধান:}

\emph{ধরি, ১০টি প্যান্টের ক্রয়মূল্য = x টাকা}

\emph{∴ ১টি প্যান্টের ক্রয়মূল্য = (x/১০) টাকা}

\emph{আবার, ৮টি প্যান্টের বিক্রয়মূল্য = x টাকা}

\emph{∴ ১টি প্যান্টের বিক্রয়মূল্য = (x/৮) টাকা}

\emph{লাভ = (x/৮)−(x/১০) = (৫x−৪x/৪০) = (x/৪০)}

\emph{∴ শতকরা লাভ হয় = ((x/৪০)/(x/১০))×১০০\% = (x/৪০)×(১০/x)×১০০\% =
২৫\%}

\emph{Shortcut:}

\emph{টাকা সমান থাকলে, শতকরা লাভ = (১০−৮/৮) × ১০০ \% = ২৫\%}
\end{quote}

\textbf{টাকায় ৩টি করে আম ক্রয় করে, টাকায় ২টি করে বিক্রয় করলে শতকরা কত লাভ
হবে?}

\textbf{\hl{ক) ৫০\% ✔}}

খ) ৩৩\%

গ) ৩০\%

ঘ) ৩১\%

\textbf{\hl{✔ সঠিক উত্তর: ক) ৫০\%}}

\begin{quote}
\textbf{ব্যাখ্যা:} \emph{মূল সমাধান:}

\emph{২টি আমের বিক্রয়মূল্য = ১ টাকা}

\emph{∴ ৩টি আমের বিক্রয়মূল্য = (১/২)×৩ টাকা = (৩/২) টাকা}

\emph{∴ লাভ = (৩/২)−১ টাকা = (১/২) টাকা}

\emph{∴ লাভের হার = (১/২)×১০০ = ৫০\%}

\emph{Shortcut:}

\emph{লাভ = (৩−২/২) × ১০০ \% = ৫০\%}
\end{quote}

\textbf{একটি ঘড়ি ৩০০ টাকায় ক্রয় করে ৩৬০ টাকায় বিক্রয় করলে ক্রয়মূল্যের উপর
শতকরা লাভ কত?}

ক) ১০\%

খ) ১৫\%

\textbf{\hl{গ) ২০\% ✔}}

ঘ) ২৫\%

\textbf{\hl{✔ সঠিক উত্তর: গ) ২০\%}}

\begin{quote}
\textbf{ব্যাখ্যা:} \emph{মূল সমাধান:}

\emph{দেওয়া আছে, ঘড়ির ক্রয়মূল্য ৩০০ টাকা}

\emph{ঘড়ির বিক্রয়মূল্য ৩৬০ টাকা}

\emph{∴ লাভ = (৩৬০ − ৩০০) টাকা = ৬০ টাকা}

\emph{এখন, ৩০০ টাকায় লাভ হয় ৬০ টাকা}

\emph{∴ ১ টাকায় লাভ হয় (৬০/৩০০) টাকা}

\emph{∴ ১০০ টাকায় লাভ হয় (৬০×১০০/৩০০) টাকা = ২০ টাকা।}

\emph{Shortcut:}

\emph{শতকরা লাভ = (৩৬০−৩০০/৩০০)×১০০\% = (৬০/৩০০)×১০০\% = ২০\%}
\end{quote}

\textbf{৫০টি কলম ২০০ টাকায় কিনে ২৫টি কলম ৫০ টাকায় বিক্রি করলে শতকরা কত
লাভ বা ক্ষতি হবে?}

ক) কোনো লাভ বা ক্ষতি হবে না

\textbf{\hl{খ) ক্ষতি ৫০\% ✔}}

গ) ক্ষতি ১০\%

ঘ) লাভ ১০\%

ঙ) কোনোটিই নয়

\textbf{\hl{✔ সঠিক উত্তর: খ) ক্ষতি ৫০\%}}

\begin{quote}
\textbf{ব্যাখ্যা:} \emph{মূল সমাধান:}

\emph{৫০টি কলমের ক্রয়মূল্য ২০০ টাকা}

\emph{∴ ২৫টি কলমের ক্রয়মূল্য ১০০ টাকা।}

\emph{∴ ক্ষতি = (১০০ − ৫০) টাকা = ৫০ টাকা}

\emph{∴ ক্ষতির হার = (৫০/১০০)×১০০\% = ৫০\%}

\emph{Shortcut:}

\emph{২৫টি কলমের ক্রয়মূল্য = ১০০ টাকা এবং বিক্রয়মূল্য = ৫০ টাকা}

\emph{∴ ক্ষতি = ১০০ − ৫০ = ৫০ টাকা, অর্থাৎ ক্ষতির হার ৫০\%}
\end{quote}

\textbf{একজন ফল বিক্রেতা প্রতি হালি কলা ২৫ টাকা দরে ক্রয় করে প্রতি ২ হালি ৫৬
টাকা দরে বিক্রয় করলে শতকরা কত লাভ হবে?}

ক) ২৪\%

\textbf{\hl{খ) ১২\% ✔}}

গ) ১৮\%

ঘ) ১৪\%

\textbf{\hl{✔ সঠিক উত্তর: খ) ১২\%}}

\begin{quote}
\textbf{ব্যাখ্যা:} \emph{মূল সমাধান:}

\emph{১ হালি কলার ক্রয়মূল্য = ২৫ টাকা।}

\emph{আবার, ১ হালি কলার বিক্রয়মূল্য = (৫৬/২) = ২৮ টাকা}

\emph{∴ লাভ = ২৮ − ২৫ = ৩ টাকা}

\emph{∴ লাভের হার = (৩/২৫)×১০০\% = ১২\%}

\emph{Shortcut:}

\emph{প্রতি হালিতে লাভ = (২৮ − ২৫) = ৩ টাকা}

\emph{∴ শতকরা লাভ = (৩/২৫)×১০০\% = ১২\%}
\end{quote}

\textbf{একটি জিনিস ২৫০ টাকায় ক্রয় করে ২৮০ টাকায় বিক্রি করলে শতকরা কত লাভ
হবে?}

ক) ১০\%

\textbf{\hl{খ) ১২\% ✔}}

গ) ১৫\%

ঘ) ২০\%

\textbf{\hl{✔ সঠিক উত্তর: খ) ১২\%}}

\begin{quote}
\textbf{ব্যাখ্যা:} \emph{মূল সমাধান:}

\emph{লাভ = (২৮০ − ২৫০) টাকা = ৩০ টাকা}

\emph{২৫০ টাকায় লাভ হয় ৩০ টাকা}

\emph{∴ ১০০ টাকায় লাভ হয় (৩০×১০০/২৫০) টাকা = ১২ টাকা।}

\emph{Shortcut:}

\emph{শতকরা লাভ = (২৮০−২৫০/২৫০)×১০০\% = (৩০/২৫০)×১০০\% = ১২\%}
\end{quote}

\textbf{একজন কমলা বিক্রেতা প্রতি শত কমলা ১২০০ টাকায় কিনে ১৮০০ টাকায় বিক্রি
করলেন। তার শতকরা কত লাভ হলো?}

\textbf{\hl{ক) ৫০\% ✔}}

খ) ৫৫\%

গ) ৬০\%

ঘ) ৪০\%

\textbf{\hl{✔ সঠিক উত্তর: ক) ৫০\%}}

\begin{quote}
\textbf{ব্যাখ্যা:} \emph{মূল সমাধান:}

\emph{এখানে, বিক্রেতার ১০০ কমলার ক্রয়মূল্য = ১২০০ টাকা}

\emph{এবং বিক্রয়মূল্য = ১৮০০ টাকা}

\emph{∴ লাভ = (বিক্রয়মূল্য − ক্রয়মূল্য) টাকা = (১৮০০ − ১২০০) = ৬০০ টাকা}

\emph{∴ লাভের হার = (৬০০×১০০/১২০০) = ৫০\%}

\emph{Shortcut:}

\emph{শতকরা লাভ = (১৮০০−১২০০/১২০০)×১০০\% = (৬০০/১২০০)×১০০\% = ৫০\%}
\end{quote}

\textbf{একটি দ্রব্য ২০০ টাকায় ক্রয় করে ১৮০ টাকায় বিক্রয় করলে শতকরা ক্ষতি কত?}

ক) ৫\%

\textbf{\hl{খ) ১০\% ✔}}

গ) ১৫\%

ঘ) ২০\%

\textbf{\hl{✔ সঠিক উত্তর: খ) ১০\%}}

\begin{quote}
\textbf{ব্যাখ্যা:} \emph{মূল সমাধান:}

\emph{দ্রব্যটি বিক্রয়ে ক্ষতি হয় = ২০০ − ১৮০ = ২০ টাকা}

\emph{২০০ টাকায় ক্ষতি হয় = ২০ টাকা}

\emph{∴ ১০০ টাকায় ক্ষতি হয় = (২০×১০০/২০০) = ১০ টাকা}

\emph{Shortcut:}

\emph{শতকরা ক্ষতি = (২০০−১৮০/২০০)×১০০\% = (২০/২০০)×১০০\% = ১০\%}
\end{quote}

\textbf{কবির ৩০০টি কলা কিনলো ৭৫০ টাকা দিয়ে। সে ১৩৫০ টাকায় সবগুলো কলা
বিক্রয় করে দিল। সে ক্রয়মূল্যের উপর শতকরা কত টাকা লাভ করলো?}

ক) ৪০\%

খ) ৫০\%

গ) ৬০\%

\textbf{\hl{ঘ) ৮০\% ✔}}

\textbf{\hl{✔ সঠিক উত্তর: ঘ) ৮০\%}}

\begin{quote}
\textbf{ব্যাখ্যা:} \emph{মূল সমাধান:}

\emph{লাভ = (১৩৫০ − ৭৫০) টাকা = ৬০০ টাকা}

\emph{৭৫০ টাকায় লাভ হয় ৬০০ টাকা}

\emph{∴ ১০০ টাকায় লাভ হয় = (৬০০×১০০/৭৫০) = ৮০ টাকা}

\emph{Shortcut:}

\emph{লাভ = ১৩৫০ − ৭৫০ টাকা = ৬০০ টাকা}

\emph{∴ লাভের হার = (৬০০/৭৫০)×১০০ = ৮০\%}
\end{quote}

\textbf{৭৫ টাকায় ১৫টি বলপেন কিনে ৯০ টাকায় বিক্রি করলে শতকরা কত লাভ হবে?}

ক) ১৫\%

\textbf{\hl{খ) ২০\% ✔}}

গ) ২৫\%

ঘ) ১৮\%

\textbf{\hl{✔ সঠিক উত্তর: খ) ২০\%}}

\begin{quote}
\textbf{ব্যাখ্যা:} \emph{মূল সমাধান:}

\emph{লাভ = ৯০ − ৭৫ টাকা = ১৫ টাকা}

\emph{৭৫ টাকায় লাভ হয় ১৫ টাকা}

\emph{∴ ১০০ টাকায় লাভ হয় = (১৫×১০০/৭৫) = ২০ টাকা}

\emph{Shortcut:}

\emph{লাভ = ৯০ − ৭৫ টাকা = ১৫ টাকা}

\emph{∴ লাভের হার = (১৫/৭৫)×১০০\% = ২০\%}
\end{quote}

\textbf{৫টি লিচু যে দরে ক্রয় করা হয়, ৪টি লিচু সেই দরে বিক্রয় করলে শতকরা কত
লাভ বা ক্ষতি হবে?}

ক) ৩০ \% ক্ষতি

\textbf{\hl{খ) ২৫ \% লাভ ✔}}

গ) ১০ \% লাভ

ঘ) ২০ \% লাভ

\textbf{\hl{✔ সঠিক উত্তর: খ) ২৫ \% লাভ}}

\begin{quote}
\textbf{ব্যাখ্যা:} \emph{মূল সমাধান:}

\emph{ধরি, ৫টি লিচুর ক্রয়মূল্য = ৪টি লিচুর বিক্রয়মূল্য = x টাকা}

\emph{∴ ১টি লিচুর ক্রয়মূল্য = (x/৫) টাকা এবং ১টি লিচুর বিক্রয়মূল্য = (x/৪) টাকা}

\emph{∴ লাভ = (x/৪)−(x/৫) = (x/২০) টাকা}

\emph{∴ শতকরা লাভ = ((x/২০)/(x/৫))×১০০\% = (x/২০)×(৫/x)×১০০\% = ২৫\%}

\emph{Shortcut:}

\emph{টাকা সমান থাকলে, শতকরা লাভ = (৫−৪/৪)×১০০\% = (১/৪)×১০০\% = ২৫\%}
\end{quote}

\textbf{এক ডজন ডিমের বিক্রয় মূল্যে ২০টি ডিম ক্রয় করলে শতকরা কত লাভ হবে?}

ক) 33 (1/3) \%

খ) 33 (1/3) \%

\textbf{\hl{গ) 66 (2/3) \% ✔}}

ঘ) 33 (2/3) \%

\textbf{\hl{✔ সঠিক উত্তর: গ) 66 (2/3) \%}}

\begin{quote}
\textbf{ব্যাখ্যা:} \emph{মূল সমাধান:}

\emph{১ ডজন = ১২টি}

\emph{ধরি, ১২টি ডিমের বিক্রয়মূল্য = ২০টি ডিমের ক্রয়মূল্য = x টাকা}

\emph{∴ ১টি ডিমের ক্রয়মূল্য = (x/২০) টাকা এবং ১টি ডিমের বিক্রয়মূল্য = (x/১২)
টাকা}

\emph{লাভ = (x/১২)−(x/২০) = (৫x−৩x/৬০) = (২x/৬০) = (x/৩০) টাকা}

\emph{∴ শতকরা লাভ = ((x/৩০)/(x/২০))×১০০\% = (x/৩০)×(২০/x)×১০০\% =
(২০০/৩)\% = 66 (2/3) \%}

\emph{Shortcut:}

\emph{শতকরা লাভ = (বিক্রয় ও ক্রয়ের পার্থক্য/বিক্রয় সংখ্যা)×100\%}

\emph{= (20−12/12)×100\% = (8/12)×100\% = 66 (2/3) \%}
\end{quote}

\textbf{একটি দ্রব্য ৯৬০ টাকায় বিক্রয় করায় ৪০ টাকা ক্ষতি হলো। শতকরা ক্ষতির হার
কত?}

\textbf{\hl{ক) ৪\% ✔}}

খ) ৩\%

গ) ৫\%

ঘ) ২\%

\textbf{\hl{✔ সঠিক উত্তর: ক) ৪\%}}

\begin{quote}
\textbf{ব্যাখ্যা:} \emph{মূল সমাধান:}

\emph{দ্রব্যটির ক্রয়মূল্য = (৯৬০ + ৪০) টাকা = ১০০০০ টাকা}

\emph{১০০০ টাকায় ক্ষতি হয় = ৪০ টাকা}

\emph{∴ ১ টাকায় ক্ষতি হয় = (৪০/১০০০) টাকা}

\emph{∴ ১০০ টাকায় ক্ষতি হয় = (৪০×১০০/১০০০) টাকা = ৪ টাকা}

\emph{∴ ক্ষতির হার ৪\%}

\emph{Shortcut:}

\emph{ক্রয়মূল্য = ৯৬০ + ৪০ = ১০০০০ টাকা}

\emph{শতকরা ক্ষতি = (৪০/১০০০)×১০০\% = ৪\%}
\end{quote}

\textbf{ক্রয়মূল্যঃবিক্রয়মূল্য = ৫ঃ৬; এতে শতকরা কত লাভ হয়?}

\textbf{\hl{ক) ২০\% ✔}}

খ) ২১

গ) ২৫\%

ঘ) কোনটিই নহে

\textbf{\hl{✔ সঠিক উত্তর: ক) ২০\%}}

\begin{quote}
\textbf{ব্যাখ্যা:} \emph{মূল সমাধান:}

\emph{ধরি, ক্রয়মূল্য = ৫x টাকা এবং বিক্রয়মূল্য = ৬x টাকা}

\emph{∴ লাভ = (৬x − ৫x) টাকা = x টাকা}

\emph{৫x টাকায় লাভ হয় x টাকা}

\emph{∴ ১০০ টাকায় লাভ হয় = (x×১০০/৫x) টাকা = ২০ টাকা।}

\emph{Shortcut:}

\emph{শতকরা লাভ = (৬−৫/৫)×১০০\% = ২০\%}
\end{quote}

\textbf{ক্রয়মূল্য বিক্রয় মূল্যের দ্বিগুণ হলে শতকরা ক্ষতি কত?}

ক) ১০০\%

খ) ২\%

গ) ১০\%

\textbf{\hl{ঘ) ৫০\% ✔}}

\textbf{\hl{✔ সঠিক উত্তর: ঘ) ৫০\%}}

\begin{quote}
\textbf{ব্যাখ্যা:} \emph{মূল সমাধান:}

\emph{ধরি, বিক্রয়মূল্য x টাকা এবং ক্রয়মূল্য ২x টাকা}

\emph{∴ ক্ষতি = (২x − x) = x টাকা}

\emph{∴ শতকরা ক্ষতি হবে = (২x−x/২x)×১০০\% = (x/২x)×১০০\% = ৫০\%}

\emph{Shortcut:}

\emph{ক্রয়মূল্য বিক্রয়মূল্যের দ্বিগুণ (২ গুণ) হলে,}

\emph{শতকরা ক্ষতি = (২−১/২)×১০০\% = ৫০\%}
\end{quote}

\textbf{কোন জিনিসের ক্রয়মূল্য বিক্রয়মূল্যের (৪/৫) ভাগ হলে শতকরা লাভের হার কত?}

ক) ২০\%

\textbf{\hl{খ) ২৫\% ✔}}

গ) ৩৩\%

ঘ) ৩৫\%

\textbf{\hl{✔ সঠিক উত্তর: খ) ২৫\%}}

\begin{quote}
\textbf{ব্যাখ্যা:} \emph{মূল সমাধান:}

\emph{ধরি, বিক্রয়মূল্য = ৫ টাকা।}

\emph{∴ ক্রয়মূল্য = ৫এর(৪/৫) টাকা = ৪ টাকা}

\emph{৪ টাকায় লাভ = (৫ − ৪) টাকা = ১ টাকা}

\emph{∴ ১ টাকায় লাভ = (১/৪) টাকা}

\emph{∴ ১০০ টাকায় লাভ = (১০০/৪) টাকা = ২৫ টাকা}

\emph{Shortcut:}

\emph{লাভ = (৫−৪/৪)×১০০\% = ২৫\%}
\end{quote}

\textbf{একটি ল্যাপটপের বিক্রয়মূল্য ক্রয়মূল্যের (১/৪) অংশের সমান হলে, শতকরা কত
ক্ষতি হবে?}

ক) ২৫\%

খ) ৫০\%

\textbf{\hl{গ) ৭৫\% ✔}}

ঘ) ৬০\%

\textbf{\hl{✔ সঠিক উত্তর: গ) ৭৫\%}}

\begin{quote}
\textbf{ব্যাখ্যা:} \emph{মূল সমাধান:}

\emph{মনে করি, ক্রয়মূল্য = ৪ টাকা।}

\emph{∴ বিক্রয়মূল্য = ৪এর(১/৪) টাকা = ১ টাকা}

\emph{∴ ক্ষতি = (৪ − ১) টাকা = ৩ টাকা}

\emph{৪ টাকায় ক্ষতি হয় ৩ টাকা}

\emph{∴ ১০০ টাকায় ক্ষতি হয় = (৩×১০০/৪) টাকা = ৭৫ টাকা}

\emph{Shortcut:}

\emph{শতকরা ক্ষতি = (৪−১/৪)×১০০\% = (৩/৪)×১০০\% = ৭৫\%}
\end{quote}

\textbf{একটি দ্রব্য ৫০০ টাকার ক্রয় করে ১০\% লাভে বিক্রয় করা হলো। দ্রব্যটির
ক্রয়মূল্য ১০\% কম হলে, কত টাকা লাভ হয়?}

\textbf{\hl{ক) ১০০ টাকা ✔}}

খ) ১১০ টাকা

গ) ১১৫ টাকা

ঘ) ১২০ টাকা

\textbf{\hl{✔ সঠিক উত্তর: ক) ১০০ টাকা}}

\begin{quote}
\textbf{ব্যাখ্যা:} \emph{মূল সমাধান:}

\emph{১০\% লাভে ৫০০ টাকার বিক্রয় মূল্য = (৫০০ + ৫০০ এর ১০\%) টাকা = ৫৫০
টাকা}

\emph{ক্রয়মূল্য ১০\% কম হলে নতুন ক্রয়মূল্য = (৫০০ − ৫০০ এর ১০\%) টাকা = ৪৫০
টাকা}

\emph{∴ মোট লাভ = (৫৫০ − ৪৫০) টাকা = ১০০ টাকা।}

\emph{Shortcut:}

\emph{লাভ = (পূর্বের লাভ\% + ক্রয়মূল্য হ্রাস\%) = (১০\% + ১০\%) = ২০\%}

\emph{∴ লাভ = ৫০০ টাকার ২০\% = ৫০০×(২০/১০০) টাকা = ১০০ টাকা।}
\end{quote}

\textbf{একজন ব্যবসায়ী ১৩,৭৭০ টাকায় একটি চেয়ার বিক্রি করায় ক্রয়মূল্যের উপর
৩৫\% লাভ হয়। সে যদি ৪৫\% লাভে চেয়ারটি বিক্রয় করলে কত লাভ হবে?}

ক) ১০,২০০

খ) ১৪,৭৯০

\textbf{\hl{গ) ৪,৫৯০ ✔}}

ঘ) ৪,৯৫০

\textbf{\hl{✔ সঠিক উত্তর: গ) ৪,৫৯০}}

\begin{quote}
\textbf{ব্যাখ্যা:} \emph{মূল সমাধান:}

\emph{৩৫\% লাভে বিক্রয়মূল্য ১৩৫ টাকা হলে ক্রয়মূল্য ১০০ টাকা}

\emph{∴ বিক্রয়মূল্য ১৩,৭৭০ টাকা হলে ক্রয়মূল্য = (১০০×১৩৭৭০/১৩৫) = ১০,২০০ টাকা}

\emph{এখন, ৪৫\% লাভে চেয়ারটির বিক্রয়মূল্য = ১০,২০০ এর ১৪৫\% =
১০২০০×(১৪৫/১০০) = ১৪,৭৯০ টাকা}

\emph{∴ লাভ হবে = (১৪,৭৯০ − ১০,২০০) টাকা = ৪,৫৯০ টাকা}

\emph{Shortcut:}

\emph{ক্রয়মূল্য = (১৩৭৭০/১.৩৫) = ১০,২০০ টাকা}

\emph{∴ ৪৫\% লাভে লাভ = ১০,২০০ এর ৪৫\% = ১০২০০×(৪৫/১০০) = ৪,৫৯০ টাকা।}
\end{quote}

\textbf{রুমির বেতন বাবুর বেতনের চেয়ে ২৫\% বেশি। বাবুর বেতন রুমির বেতনের চেয়ে
শতকরা কত ভাগ কম?}

ক) ১৬.৬৬

\textbf{\hl{খ) ২০ ✔}}

গ) ২২.৫০

ঘ) কোনোটিই নয়

\textbf{\hl{✔ সঠিক উত্তর: খ) ২০}}

\begin{quote}
\textbf{ব্যাখ্যা:} \emph{মূল সমাধান:}

\emph{২৫\% বৃদ্ধিতে, ১২৫ টাকায় রুমির বেতন অপেক্ষা বাবুর বেতন কম ২৫ টাকা।}

\emph{∴ ১০০ টাকায় রুমির বেতন অপেক্ষা বাবুর বেতন কম = (১০০×২৫/১২৫) = ২০
টাকা।}

\emph{Shortcut:}

\emph{কমের শতকরা হার = (১০০×বৃদ্ধি/১০০+বৃদ্ধি)\% = (১০০×২৫/১০০+২৫)\% =
(২৫০০/১২৫)\% = ২০\%}
\end{quote}

\textbf{একজন দোকানদার ৫০টি কলম বিক্রি করে যে লাভ করে তা ২৫টি কলমের ক্রয়
মূল্যের সমান। তার শতকরা কত টাকা লাভ হয়?}

ক) ১০০\%

\textbf{\hl{খ) ৫০\% ✔}}

গ) ২৫\%

ঘ) কোনোটিই নয়

\textbf{\hl{✔ সঠিক উত্তর: খ) ৫০\%}}

\begin{quote}
\textbf{ব্যাখ্যা:} \emph{মূল সমাধান:}

\emph{ধরি, ১টি কলমের ক্রয়মূল্য x টাকা}

\emph{∴ ৫০টি কলমের ক্রয়মূল্য ৫০x টাকা}

\emph{লাভ = ২৫টি কলমের ক্রয়মূল্য = ২৫x টাকা}

\emph{৫০x টাকায় লাভ হয় ২৫x টাকা}

\emph{∴ ১০০ টাকায় লাভ হয় = (১০০×২৫x/৫০x) টাকা = ৫০ টাকা।}

\emph{Shortcut:}

\emph{শতকরা লাভ = (লাভের পরিমাণ (ক্রয়মূল্যে)/মোট বিক্রিত সংখ্যা)×১০০\% =
(২৫/৫০)×১০০\% = ৫০\%}
\end{quote}

\textbf{এক ব্যক্তি চাল বিক্রয় করে দেখল যে, ২৫ কেজি চালের বিক্রয়মূল্য ২০ কেজি
চালের ক্রয়মূল্যের সমান। তার শতকরা কত ক্ষতি হল?}

ক) ১৫\%

\textbf{\hl{খ) ২০\% ✔}}

গ) ২৫\%

ঘ) ৩০\%

\textbf{\hl{✔ সঠিক উত্তর: খ) ২০\%}}

\begin{quote}
\textbf{ব্যাখ্যা:} \emph{মূল সমাধান:}

\emph{ধরি, ২০ কেজি চালের ক্রয়মূল্য = ক টাকা}

\emph{∴ ১ কেজি চালের ক্রয়মূল্য = ক ÷ ২০ টাকা}

\emph{আবার, ২৫ কেজি চালের বিক্রয়মূল্য = ক টাকা}

\emph{∴ ১ কেজি চালের বিক্রয়মূল্য = ক ÷ ২৫ টাকা}

\emph{∴ ক্ষতি = (ক ÷ ২০) − (ক ÷ ২৫) = (ক/১০০) টাকা}

\emph{∴ শতকরা ক্ষতি = (ক/১০০)÷(ক/২০)×১০০\% = ২০\%}

\emph{Shortcut:}

\emph{শতকরা ক্ষতি = (২৫−২০/২৫)×১০০\% = (৫/২৫)×১০০\% = ২০\%}
\end{quote}

\textbf{৪ টাকায় ৫টি কিনে ৫ টাকায় ৪টি বিক্রয় করলে শতকরা কত লাভ হবে?}

ক) ৪৫\%

খ) ৪৮.৫০\%

গ) ৫২.৭৫\%

\textbf{\hl{ঘ) ৫৬.২৫\% ✔}}

\textbf{\hl{✔ সঠিক উত্তর: ঘ) ৫৬.২৫\%}}

\begin{quote}
\textbf{ব্যাখ্যা:} \emph{মূল সমাধান:}

\emph{১টির ক্রয়মূল্য = (৪/৫) = ০.৮ টাকা}

\emph{আবার, ১টির বিক্রয়মূল্য = (৫/৪) = ১.২৫ টাকা}

\emph{∴ লাভ = ১.২৫ − ০.৮ = ০.৪৫ টাকা}

\emph{∴ লাভের হার = (০.৪৫×১০০/০.৮)\% = ৫৬.২৫\%}

\emph{Shortcut:}

\emph{শতকরা লাভ = (বিক্রয়মূল্য\^{}২−ক্রয়মূল্য\^{}২/ক্রয়মূল্য\^{}২)×১০০\% (এখানে
৫টি, ৪ টাকা এবং ৪টি, ৫ টাকা)}

\emph{= (৫\^{}২−৪\^{}২/৪\^{}২)×১০০\% = (২৫−১৬/১৬)×১০০\% = (৯/১৬)×১০০\% =
৫৬.২৫\%}
\end{quote}

\textbf{একজন বিক্রেতা ১০ টাকায় ৪টি করে বল ক্রয় করে ৬ টাকায় ২টি করে বল
বিক্রয় করলো। সে শতকরা কত টাকা লাভ করলো?}

ক) ১৭.৫\%

\textbf{\hl{খ) ২০\% ✔}}

গ) ২২\%

ঘ) কোনোটিই নয়

\textbf{\hl{✔ সঠিক উত্তর: খ) ২০\%}}

\begin{quote}
\textbf{ব্যাখ্যা:} \emph{মূল সমাধান:}

\emph{৪টি বলের ক্রয়মূল্য ১০ টাকা}

\emph{∴ ১টি বলের ক্রয়মূল্য = (১০/৪) = ২.৫ টাকা}

\emph{২টি বলের বিক্রয়মূল্য ৬ টাকা}

\emph{∴ ১টি বলের বিক্রয়মূল্য = (৬/২) = ৩ টাকা}

\emph{∴ লাভ = ৩ − ২.৫ = ০.৫ টাকা}

\emph{∴ লাভের হার = (০.৫/২.৫)×১০০\% = ২০\%}

\emph{Shortcut:}

\emph{৪টি বলের ক্রয়মূল্য = ১০ টাকা}

\emph{৪টি বলের বিক্রয়মূল্য = (৩ × ৪) = ১২ টাকা}

\emph{∴ লাভ = (১২ − ১০) = ২ টাকা}

\emph{∴ শতকরা লাভ = (২/১০)×১০০\% = ২০\%}
\end{quote}

\textbf{২০ টাকায় এক ডজন কলা কিনে প্রতিটি ২ টাকা করে বিক্রয় করলে শতকরা কত
লাভ হবে?}

\textbf{\hl{ক) ২০\% ✔}}

খ) ১২.৫\%

গ) ১৫\%

ঘ) ১০\%

\textbf{\hl{✔ সঠিক উত্তর: ক) ২০\%}}

\begin{quote}
\textbf{ব্যাখ্যা:} \emph{মূল সমাধান:}

\emph{১টি কলা বিক্রি করে = ২ টাকা}

\emph{∴ ১২টি কলা বিক্রি করে = (২ × ১২) টাকা = ২৪ টাকা}

\emph{∴ লাভ = (২৪ − ২০) টাকা = ৪ টাকা}

\emph{এখন, ২০ টাকায় লাভ হয় = ৪ টাকা}

\emph{∴ ১ টাকায় লাভ হয় = (৪/২০) টাকা}

\emph{∴ ১০০ টাকায় লাভ হয় = (৪×১০০/২০) টাকা = ২০ টাকা}

\emph{Shortcut:}

\emph{শতকরা লাভ = (মোট বিক্রয়মূল্য − ক্রয়মূল্য/ক্রয়মূল্য)×১০০\%}

\emph{= (২৪−২০/২০)×১০০\% = (৪/২০)×১০০\% = ২০\%}
\end{quote}

\textbf{একটি ঘড়ি ৫৪০ টাকায় বিক্রি করলে ১০\% ক্ষতি হয়। ঘড়িটি ৬৬০ টাকায়
বিক্রি করলে শতকরা কত লাভ বা ক্ষতি হবে?}

ক) লাভ ১০ টাকা

খ) ক্ষতি ১০\%

গ) লাভ ৬০ টাকা

\textbf{\hl{ঘ) লাভ ১০\% ✔}}

\textbf{\hl{✔ সঠিক উত্তর: ঘ) লাভ ১০\%}}

\begin{quote}
\textbf{ব্যাখ্যা:} \emph{মূল সমাধান:}

\emph{১০\% ক্ষতিতে,}

\emph{ঘড়ির বিক্রয়মূল্য ৯০ টাকা হলে ক্রয়মূল্য ১০০ টাকা}

\emph{∴ ঘড়ির বিক্রয়মূল্য ৫৪০ টাকা হলে ক্রয়মূল্য = (১০০×৫৪০/৯০) = ৬০০ টাকা}

\emph{ঘড়িটি ৬৬০ টাকায় বিক্রি করলে লাভ = (৬৬০ − ৬০০) = ৬০ টাকা}

\emph{৬০০ টাকায় লাভ হয় = ৬০ টাকা}

\emph{∴ ১০০ টাকায় লাভ হয় = (৬০×১০০/৬০০) টাকা = ১০ টাকা}

\emph{∴ লাভ ১০\%}

\emph{Shortcut:}

\emph{ক্রয়মূল্য = (৫৪০/০.৯০) = ৬০০ টাকা}

\emph{শতকরা লাভ = (৬৬০−৬০০/৬০০)×১০০\% = (৬০/৬০০)×১০০\% = ১০\%}
\end{quote}

\textbf{২৫\% লাভে কোনো জিনিস ১৫ টাকায় বিক্রয় করলে বিক্রেতার কত টাকা লাভ
হবে?}

ক) ২

\textbf{\hl{খ) ৩ ✔}}

গ) ৪

ঘ) ৫

\textbf{\hl{✔ সঠিক উত্তর: খ) ৩}}

\begin{quote}
\textbf{ব্যাখ্যা:} \emph{মূল সমাধান:}

\emph{ধরি, জিনিসটির ক্রয়মূল্য x টাকা}

\emph{প্রশ্নমতে, x এর ১২৫\% = ১৫}

\emph{⇒ (১২৫x/১০০) = ১৫}

\emph{∴ x = (১৫×১০০/১২৫) = ১২}

\emph{∴ লাভ হবে = (১৫ − ১২) টাকা = ৩ টাকা।}

\emph{Shortcut:}

\emph{২৫\% লাভে, ১২৫ টাকা বিক্রয়মূল্য হলে লাভ = ২৫ টাকা}

\emph{∴ ১৫ টাকা বিক্রয়মূল্য হলে লাভ = (২৫×১৫/১২৫) = ৩ টাকা।}
\end{quote}

\textbf{একজন দালাল দুটি পুরনো গাড়ী প্রত্যেকটি ৩৯,১০০ টাকায় বিক্রয় করেন।
প্রথমটিতে তার ১৫\% লাভ হয় কিন্তু দ্বিতীয়টিতে ৭.৫\% ক্ষতি হয়। মোটের উপর তার কত
টাকা লাভ হয়?}

ক) ২০০০ টাকা

খ) ১৯০০ টাকা

\textbf{\hl{গ) ১৯২৯.৭৩ টাকা ✔}}

ঘ) ১২০০ টাকা

\textbf{\hl{✔ সঠিক উত্তর: গ) ১৯২৯.৭৩ টাকা}}

\begin{quote}
\textbf{ব্যাখ্যা:} \emph{মূল সমাধান:}

\emph{১৫\% লাভে বিক্রয় মূল্য ১১৫ টাকা হলে ক্রয়মূল্য ১০০ টাকা}

\emph{∴ ৩৯,১০০ টাকা বিক্রয়মূল্য হলে ক্রয়মূল্য = (১০০×৩৯১০০/১১৫) = ৩৪০০০ টাকা}

\emph{৭.৫\% ক্ষতিতে বিক্রয় মূল্য (১০০ − ৭.৫) = ৯২.৫ টাকা হলে ক্রয়মূল্য ১০০
টাকা}

\emph{∴ ৩৯,১০০ টাকা বিক্রয়মূল্য হলে ক্রয়মূল্য = (১০০×৩৯১০০/৯২.৫) = ৪২২৯০.২৭
টাকা (প্রায়)}

\emph{গাড়ীদুটির মোট ক্রয়মূল্য = ৩৪০০০ + ৪২২৯০.২৭ = ৭৬২৯০.২৭ টাকা}

\emph{গাড়ীদুটির মোট বিক্রয়মূল্য = ৩৯১০০ × ২ = ৭৮২০০ টাকা}

\emph{∴ মোটের উপর লাভ হয় = ৭৮২০০ − ৭৬২৯০.২৭ = ১৯২৯.৭৩ টাকা}

\emph{Shortcut:}

\emph{মোট লাভ = ৩৯১০০×(১৫/১১৫)−৩৯১০০×(৭.৫/৯২.৫)}

\emph{= ৫১০০ − ৩১৭০.২৭ = ১৯২৯.৭৩ টাকা।}
\end{quote}

\textbf{৫০ টাকায় ২টি এবং ৫০ টাকায় ৩টি দরে সমসংখ্যক কমলা ক্রয় করে প্রতি ২টি
কমলা ৪৭ টাকায় বিক্রি করলে শতকরা কত লাভ বা ক্ষতি হবে?}

\textbf{\hl{ক) ১২.৮\% ক্ষতি ✔}}

খ) ২৫\% লাভ

গ) ৩০\% লাভ

ঘ) ৩৫\% লাভ

\textbf{\hl{✔ সঠিক উত্তর: ক) ১২.৮\% ক্ষতি}}

\begin{quote}
\textbf{ব্যাখ্যা:} \emph{মূল সমাধান:}

\emph{২ ও ৩ এর ল.সা.গু. = ৬}

\emph{২টি কমলার দাম ৫০ টাকা}

\emph{∴ ৬টি কমলার দাম = (৫০×৬/২) = ১৫০ টাকা}

\emph{৩টি কমলার দাম ৫০ টাকা}

\emph{∴ ৬টি কমলার দাম = (৫০×৬/৩) = ১০০ টাকা}

\emph{∴ মোট (৬ + ৬) = ১২টি কমলার ক্রয়মূল্য = (১৫০ + ১০০) = ২৫০ টাকা}

\emph{আবার, ২টির বিক্রয়মূল্য ৪৭ টাকা}

\emph{∴ ১২টি কমলার বিক্রয়মূল্য = (৪৭×১২/২) = ২৮২ টাকা}

\emph{লাভ = ২৮২ − ২৫০ = ৩২ টাকা}

\emph{শতকরা লাভ = (৩২×১০০/২৫০)\% = ১২.৮\%}
\end{quote}

\textbf{এক টিভি বিক্রেতা ৪৫\% লাভে টিভি বিক্রি করত। মন্দার কারণে সে তার
লাভের হার ৪০\% করে এবং এতে তার বিক্রয় ১০\% বেড়ে যায়। তার নতুন লাভ ও আগের
লাভের অনুপাত কত?}

ক) ৯:৮

খ) ১১:১০

গ) ৪৪:৪৪

\textbf{\hl{ঘ) ৪৪:৪৫ ✔}}

\textbf{\hl{✔ সঠিক উত্তর: ঘ) ৪৪:৪৫}}

\begin{quote}
\textbf{ব্যাখ্যা:} \emph{মূল সমাধান:}

\emph{ধরি, ১টি টিভির ক্রয়মূল্য ১০০ টাকা}

\emph{৪৫\% লাভে টিভির বিক্রয় সংখ্যা x টি}

\emph{১টি টিভিতে লাভ করে ৪৫ টাকা}

\emph{∴ x টি টিভিতে লাভ করে ৪৫x টাকা}

\emph{১০\% বেশি বিক্রিতে টেলিভিশনের সংখ্যা (x + x এর ১০\%) = (১১x/১০) টি}

\emph{∴ ৪০\% লাভে মোট লাভের পরিমাণ ৪০×(১১x/১০) টাকা = ৪৪x টাকা}

\emph{∴ নতুন লাভ : আগের লাভ = ৪৪x : ৪৫x = ৪৪:৪৫}

\emph{Shortcut:}

\emph{আগের মোট লাভ = ৪৫ × ১০০ = ৪৫০০}

\emph{নতুন বিক্রি = ১-এর জায়গায় ১.১ (১০\% বৃদ্ধি) এবং লাভ = ৪০}

\emph{নতুন মোট লাভ = ৪০ × ১.১ × ১০০ = ৪৪০০}

\emph{∴ নতুন লাভ : আগের লাভ = ৪৪০০ : ৪৫০০ = ৪৪:৪৫}
\end{quote}

\textbf{একজন দোকানদার কিছু পণ্য ক্রয় করলেন। পরিবহনের সময় ১৩\% পণ্য নষ্ট হয়ে
গেল এবং ৭\% পণ্য চুরি হয়ে গেল। মোটের উপর ২০\% লাভ করতে হলে তাকে অবশিষ্ট পণ্য
শতকরা কত লাভে বিক্রয় করতে হবে?}

ক) ২০\%

খ) ৩০\%

\textbf{\hl{গ) ৫০\% ✔}}

ঘ) ৪০\%

\textbf{\hl{✔ সঠিক উত্তর: গ) ৫০\%}}

\begin{quote}
\textbf{ব্যাখ্যা:} \emph{মূল সমাধান:}

\emph{ধরি, মোট পণ্যের ক্রয়মূল্য = ১০০ টাকা}

\emph{অবশিষ্ট পণ্য = \{১০০ − (১৩ + ৭)\} টাকার = ৮০ টাকার}

\emph{২০\% লাভে মোট বিক্রয়মূল্য = (১০০ + ২০) টাকা = ১২০ টাকা}

\emph{৮০ টাকার পণ্য বিক্রি করে লাভ করতে হবে = (১২০ − ৮০) টাকা = ৪০ টাকা}

\emph{∴ ৮০ টাকায় লাভ করতে হবে ৪০ টাকা}

\emph{∴ ১ টাকায় লাভ করতে হবে (৪০/৮০) টাকা}

\emph{∴ ১০০ টাকায় লাভ করতে হবে (৪০×১০০/৮০) টাকা = ৫০ টাকা।}

\emph{∴ অবশিষ্ট পণ্য ৫০\% লাভে বিক্রি করতে হবে।}

\emph{Shortcut:}

\emph{অবশিষ্ট পণ্য = ১০০ − (১৩ + ৭) = ৮০\%}

\emph{প্রয়োজনীয় মোট মান = ১০০ + ২০ = ১২০}

\emph{∴ অবশিষ্টাংশের শতকরা লাভ = (১২০−৮০/৮০)×১০০\% = (৪০/৮০)×১০০\% =
৫০\%}
\end{quote}

\textbf{৩৫০ টাকা দরে ৩ কেজি মিষ্টি কিনে ৪ টাকা হারে ভ্যাট দিলে মোট কত ভ্যাট
দিতে হবে?}

ক) ১৪ টাকা

\textbf{\hl{খ) ৪২ টাকা ✔}}

গ) ১২ টাকা

ঘ) ১০৫ টাকা

\textbf{\hl{✔ সঠিক উত্তর: খ) ৪২ টাকা}}

\begin{quote}
\textbf{ব্যাখ্যা:} \emph{মূল সমাধান:}

\emph{৩ কেজি মিষ্টির মূল্য = (৩৫০ × ৩) টাকা = ১০৫০ টাকা}

\emph{১০০ টাকায় ভ্যাট দিতে হয় ৪ টাকা}

\emph{∴ ১০৫০ টাকায় ভ্যাট দিতে হয় = (১০৫০×৪/১০০) টাকা = ৪২ টাকা}

\emph{Shortcut:}

\emph{মোট ভ্যাট = (৩৫০ × ৩) এর ৪\% = ১০৫০ × ০.০৪ = ৪২ টাকা।}
\end{quote}

\textbf{ছোয়াদ সাহেব ৩০০ টাকার বই কিনলেন। দোকানদার তাকে (১৯/২) \% কমিশন
দিলেন। তিনি কত টাকা কমিশন পেলেন?}

ক) ১৯.৫ টাকা

\textbf{\hl{খ) ২৮.৫ টাকা ✔}}

গ) ৫১.৫ টাকা

ঘ) ৫৭.০০ টাকা

\textbf{\hl{✔ সঠিক উত্তর: খ) ২৮.৫ টাকা}}

\begin{quote}
\textbf{ব্যাখ্যা:} \emph{মূল সমাধান:}

\emph{১০০ টাকায় কমিশন পান = (১৯/২) টাকা = ৯.৫ টাকা}

\emph{∴ ৩০০ টাকায় কমিশন পান = (৯.৫ × ৩) টাকা = ২৮.৫ টাকা}

\emph{Shortcut:}

\emph{কমিশন = ৩০০ এর (১৯/২)\% = ৩০০×(১৯/২×১০০) = (৩×১৯/২) = (৫৭/২) =
২৮.৫ টাকা।}
\end{quote}

\textbf{একজন টেলিভিশন বিক্রেতাকে বিক্রয়মূল্যের উপর ১২\% হারে কর দিতে হয়। যে
টেলিভিশনের বিক্রয়মূল্য ১৯,৪০০ টাকা, তার জন্য কত টাকা কর দিতে হবে?}

ক) ১৬৭৫ টাকা

\textbf{\hl{খ) ২৩২৮ টাকা ✔}}

গ) ২৬০০ টাকা

ঘ) ২৭০০ টাকা

ঙ) কোনটিই নয়

\textbf{\hl{✔ সঠিক উত্তর: খ) ২৩২৮ টাকা}}

\begin{quote}
\textbf{ব্যাখ্যা:} \emph{মূল সমাধান:}

\emph{বিক্রয়মূল্য ১০০ টাকা হলে কর দিতে হবে ১২ টাকা}

\emph{∴ বিক্রয়মূল্য ১৯,৪০০ টাকা হলে কর দিতে হবে = (১২×১৯৪০০/১০০) টাকা =
২৩২৮ টাকা।}

\emph{Shortcut:}

\emph{মোট কর = ১৯,৪০০ এর ১২\% = ১৯৪০০×(১২/১০০) = ২৩২৮ টাকা।}
\end{quote}

\textbf{চা পাতার উপর কর ১০\% কমালে ২৫০০ টাকায় পূর্বাপেক্ষা ৫০ কেজি চা পাতার
কর বেশি দেয়া যায়। চা পাতার কর কত টাকা কমেছে?}

ক) ১৮৫ টাকা

\textbf{\hl{খ) ২৫০ টাকা ✔}}

গ) ১৫০ টাকা

ঘ) ১০০ টাকা

\textbf{\hl{✔ সঠিক উত্তর: খ) ২৫০ টাকা}}

\begin{quote}
\textbf{ব্যাখ্যা:} \emph{মূল সমাধান:}

\emph{১০০ টাকায় কমে ১০ টাকা}

\emph{∴ ২৫০০ টাকায় কমে = (১০×২৫০০/১০০) টাকা = ২৫০ টাকা।}

\emph{Shortcut:}

\emph{মোট কর হ্রাস = ২৫০০ টাকার ১০\% = ২৫০০×(১০/১০০) = ২৫০ টাকা।}
\end{quote}

\textbf{চা পাতার উপর কর ১০\% কমালে ৩,০০০ টাকায় পূর্বাপেক্ষা ৬০ কেজি বেশী চা
পাতায় কর দেয়া যায়। প্রতি কেজি চা পাতার উপর বর্তমান কর কত টাকা?}

\textbf{\hl{ক) ৫ টাকা ✔}}

খ) ৬ টাকা

গ) ৭ টাকা

ঘ) ১০ টাকা

\textbf{\hl{✔ সঠিক উত্তর: ক) ৫ টাকা}}

\begin{quote}
\textbf{ব্যাখ্যা:} \emph{মূল সমাধান:}

\emph{৩,০০০ টাকার ১০\% = ৩০০০×(১০/১০০) = ৩০০ টাকা}

\emph{৬০ কেজি চা পাতায় কর কমে = ৩০০ টাকা}

\emph{∴ ১ কেজি চা পাতায় বর্তমান কর = (৩০০/৬০) টাকা = ৫ টাকা।}

\emph{Shortcut:}

\emph{প্রতি কেজি চা পাতার বর্তমান কর = (যত \% কমে/১০০)×(বর্তমানে যত টাকা
দিতে হয়/কেজি)}

\emph{= (১০/১০০)×(৩০০০/৬০) = ৫ টাকা।}
\end{quote}

\textbf{পূর্ববর্তী বছরের তুলনায় ২০১০ সালে চালের মূল্য ১০\% বৃদ্ধি পায়। ২০১১ সালে
চালের মূল্য ৫\% হ্রাস পায়। ২০০৯ সালের তুলনায় ২০১১ সালে চালের মূল্য কত বৃদ্ধি
পেয়েছে?}

ক) ৪\%

\textbf{\hl{খ) ৪.৫\% ✔}}

গ) ৫.০\%

ঘ) ৫.৫\%

\textbf{\hl{✔ সঠিক উত্তর: খ) ৪.৫\%}}

\begin{quote}
\textbf{ব্যাখ্যা:} \emph{মূল সমাধান:}

\emph{২০০৯ সালের চালের মূল্য ১০০ টাকা হলে ২০১০ সালে চালের মূল্য = (১০০ + ১০)
= ১১০ টাকা}

\emph{২০১১ সালে ৫\% হ্রাস পেয়ে ১০০ টাকার চালের মূল্য ৯৫ টাকা}

\emph{∴ ১১০ টাকার চালের মূল্য = (৯৫×১১০/১০০) টাকা = ১০৪.৫ টাকা}

\emph{২০০৯ সালের তুলনায় ২০১১ সালে মূল্য বৃদ্ধি পেয়েছে = ১০৪.৫ − ১০০ = ৪.৫\%}

\emph{Shortcut:}

\emph{কার্যকর পরিবর্তন \% = A+B+(A×B/১০০)}

\emph{এখানে, A = +১০, B = −৫}

\emph{= ১০−৫+(১০×(−৫)/১০০) = ৫−০.৫ = ৪.৫\%}
\end{quote}

\textbf{কোনো পণ্যের পূর্বমূল্য : বর্তমান মূল্য = ৫:৭, শতকরা মূল্য বৃদ্ধির পরিমাণ-}

\textbf{\hl{ক) ৪০\% ✔}}

খ) ২০\%

গ) ৫০\%

ঘ) ১০০\%

\textbf{\hl{✔ সঠিক উত্তর: ক) ৪০\%}}

\begin{quote}
\textbf{ব্যাখ্যা:} \emph{মূল সমাধান:}

\emph{৫ টাকায় মূল্য বেড়েছে (৭ − ৫) বা, ২ টাকা}

\emph{∴ ১ টাকায় মূল্য বেড়েছে (২/৫) টাকা}

\emph{∴ ১০০ টাকায় মূল্য বেড়েছে (২×১০০/৫) টাকা = ৪০ টাকা}

\emph{Shortcut:}

\emph{শতকরা বৃদ্ধি = (৭−৫/৫)×১০০\% = (২/৫)×১০০\% = ৪০\%}
\end{quote}

\textbf{যদি মোবাইল ফোনের পূর্ব মূল্য: বর্তমান মূল্য = 5:3 হয়, তবে পূর্ব মূল্যের
তুলনায় মোবাইল ফোনের মূল্য হ্রাস পেয়েছে?}

ক) 30\%

\textbf{\hl{খ) 40\% ✔}}

গ) 45\%

ঘ) 35\%

\textbf{\hl{✔ সঠিক উত্তর: খ) 40\%}}

\begin{quote}
\textbf{ব্যাখ্যা:} \emph{মূল সমাধান:}

\emph{ধরি, পূর্ব মূল্য = 5x টাকা; বর্তমান মূল্য = 3x টাকা}

\emph{∴ মূল্য হ্রাস পায় = (5x − 3x) = 2x টাকা}

\emph{5x টাকা কমে যায় = 2x টাকা}

\emph{∴ 1 টাকা কমে যায় = (2x/5x) টাকা}

\emph{∴ 100 টাকা কমে যায় = (2×100/5) টাকা = 40 টাকা = 40\%}

\emph{Shortcut:}

\emph{শতকরা হ্রাস = (5−3/5)×100\% = (2/5)×100\% = 40\%}
\end{quote}

\textbf{৬টি কলমের দাম ১৫টি পেন্সিলের দামের সমান হলে, ১০টি কলমের পরিবর্তে
কতটি পেন্সিল পাওয়া যাবে?}

ক) ২০টি

\textbf{\hl{খ) ২৫টি ✔}}

গ) ৩০টি

ঘ) ২২টি

\textbf{\hl{✔ সঠিক উত্তর: খ) ২৫টি}}

\begin{quote}
\textbf{ব্যাখ্যা:} \emph{মূল সমাধান:}

\emph{৬টি কলমের দাম = ১৫টি পেন্সিলের দাম}

\emph{∴ ১টি কলমের দাম = (১৫/৬) টি পেন্সিলের দাম}

\emph{∴ ১০টি কলমের দাম = (১৫×১০/৬) টি পেন্সিলের দাম = ২৫টি পেন্সিল।}

\emph{Shortcut:}

\emph{প্রয়োজনীয় পেন্সিল = (১৫×১০/৬) = ২৫টি।}
\end{quote}

\textbf{১ ডজন ডিমের দাম ৫৪ টাকা হলে, ৪৫ টাকায় কয়টি ডিম পাওয়া যাবে।}

ক) ৮টি

\textbf{\hl{খ) ১০টি ✔}}

গ) ৯টি

ঘ) ১১টি

\textbf{\hl{✔ সঠিক উত্তর: খ) ১০টি}}

\begin{quote}
\textbf{ব্যাখ্যা:} \emph{মূল সমাধান:}

\emph{আমরা জানি, ১ ডজন = ১২টি}

\emph{৫৪ টাকায় পাওয়া যায় ১২টি ডিম}

\emph{∴ ১ টাকায় পাওয়া যায় = (১২/৫৪) টি ডিম}

\emph{∴ ৪৫ টাকায় পাওয়া যায় = (১২×৪৫/৫৪) টি ডিম = ১০টি ডিম।}

\emph{Shortcut:}

\emph{ডিম সংখ্যা = (১২×৪৫/৫৪) = ১০টি।}
\end{quote}

\textbf{যদি তেলের মূল্য ২৫\% বৃদ্ধি পায় তবে তেলের ব্যবহার শতকরা কত কমালে তেল
বাবদ ব্যয় বৃদ্ধি পাবে না?}

ক) ১৬\%

\textbf{\hl{খ) ২০\% ✔}}

গ) ২৫\%

ঘ) ২৪\%

\textbf{\hl{✔ সঠিক উত্তর: খ) ২০\%}}

\begin{quote}
\textbf{ব্যাখ্যা:} \emph{মূল সমাধান:}

\emph{২৫\% বৃদ্ধিতে তেলের মূল্য হয় ১২৫ টাকা}

\emph{১২৫ টাকায় কমাতে হবে ২৫ টাকা}

\emph{∴ ১০০ টাকায় কমাতে হবে = (২৫×১০০/১২৫) টাকা = ২০ টাকা}

\emph{Shortcut:}

\emph{ব্যবহার হ্রাসের হার = (১০০×বৃদ্ধি/১০০+বৃদ্ধি)\% = (১০০×২৫/১২৫)\% = ২০\%}

\emph{(মনে রাখার নিয়ম: ম্যাথে ২৫\% বৃদ্ধি থাকলে ব্যবহার ২০\% কমাতে হয়, আর
২০\% বৃদ্ধি থাকলে ২৫\% কমাতে হয়।)}
\end{quote}

\textbf{ডালের মূল্য ২৫\% বৃদ্ধি পেলে পূর্বে যে ডালের কেজি প্রতি মূল্য ৭২ টাকা ছিল
বর্তমানে ঐ ডালের মূল্য কেজি প্রতি কত?}

\textbf{\hl{ক) ৯০ টাকা ✔}}

খ) ৯৫ টাকা

গ) ১০০ টাকা

ঘ) ১২৫ টাকা

\textbf{\hl{✔ সঠিক উত্তর: ক) ৯০ টাকা}}

\begin{quote}
\textbf{ব্যাখ্যা:} \emph{মূল সমাধান:}

\emph{২৫\% বৃদ্ধিতে,}

\emph{ডালের পূর্বমূল্য ১০০ টাকা হলে বর্তমান মূল্য ১২৫ টাকা}

\emph{∴ ডালের পূর্বমূল্য ১ টাকা হলে বর্তমান মূল্য = (১২৫/১০০) টাকা}

\emph{∴ ডালের পূর্বমূল্য ৭২ টাকা হলে বর্তমান মূল্য = (১২৫×৭২/১০০) টাকা = ৯০
টাকা।}

\emph{Shortcut:}

\emph{বর্তমান মূল্য = ৭২ এর ১২৫\% = ৭২×(১২৫/১০০) = ৯০ টাকা।}
\end{quote}

\textbf{চিনির মূল্য ২০\% কমলো, কিন্তু চিনির ব্যবহার ২০\% বৃদ্ধি পেল। এতে চিনি
বাবদ ব্যয় শতকরা কত বাড়লো বা কমলো?}

ক) ১\% বাড়লো (1\% increase)

খ) ২\% কমলো (2\% decrease)

গ) ৩\% বাড়লো (3\% increase)

\textbf{\hl{ঘ) ৪\% কমলো (4\% decrease) ✔}}

\textbf{\hl{✔ সঠিক উত্তর: ঘ) ৪\% কমলো (4\% decrease)}}

\begin{quote}
\textbf{ব্যাখ্যা:} \emph{মূল সমাধান:}

\emph{ধরি, পূর্বে ১০০ কেজি চিনির মূল্য ছিল ১০০ টাকা ∴ মোট খরচ = (১০০ × ১০০)
= ১০,০০০ টাকা}

\emph{চিনির মূল্য ২০\% কমে নতুন মূল্য = (১০০ − ২০) = ৮০ টাকা}

\emph{চিনির ব্যবহার ২০\% বেড়ে নতুন ব্যবহার = (১০০ + ২০) = ১২০ কেজি}

\emph{নতুন মোট খরচ = (৮০ × ১২০) = ৯,৬০০ টাকা}

\emph{∴ খরচ কমলো = ১০,০০০ − ৯,৬০০ = ৪০০ টাকা}

\emph{শতকরা হ্রাস = (৪০০×১০০/১০০০০)\% = ৪\% কমলো।}

\emph{Shortcut:}

\emph{কার্যকর পরিবর্তন \% = A+B+(A×B/১০০)}

\emph{এখানে, A = −২০, B = +২০}

\emph{= −২০+২০+((−২০)×২০/১০০) = −৪\% (ঋণাত্মক অর্থ কমলো)।}
\end{quote}

\textbf{যদি রিকশাভাড়া ৬০\% বৃদ্ধি পায় তবে রিকশায় যাতায়াত শতকরা কতভাগ কমালে
ব্যয় বৃদ্ধি পাবে না?}

\textbf{\hl{ক) ৩৭.৫\% ✔}}

খ) ৪০\%

গ) ৬০\%

ঘ) ৬০.৫\%

\textbf{\hl{✔ সঠিক উত্তর: ক) ৩৭.৫\%}}

\begin{quote}
\textbf{ব্যাখ্যা:} \emph{মূল সমাধান:}

\emph{৬০\% বৃদ্ধিতে ভাড়া ১৬০ টাকা}

\emph{১৬০ টাকার স্থলে খরচ করতে হবে ১০০ টাকা}

\emph{∴ ১ টাকার স্থলে খরচ করতে হবে (১০০/১৬০) টাকা}

\emph{∴ ১০০ টাকার স্থলে খরচ করতে হবে (১০০×১০০/১৬০) টাকা = ৬২.৫ টাকা}

\emph{∴ খরচ কমাতে হবে = (১০০ − ৬২.৫)\% = ৩৭.৫\%}

\emph{Shortcut:}

\emph{ব্যবহার হ্রাসের হার = (১০০×বৃদ্ধি/১০০+বৃদ্ধি)\% = (১০০×৬০/১০০+৬০)\% =
(৬০০০/১৬০)\% = ৩৭.৫\%}
\end{quote}

\textbf{গ্যাসের দাম ৪০\% বেড়ে যাওয়ায় কোন পরিবার গ্যাসের ব্যবহার এমনভাবে
কমান যে গ্যাসের ব্যয় অপরিবর্তিত থাকলো। গ্যাসের ব্যবহার শতকরা কত কমিয়েছিল?}

ক) ৮২.৫৭\%

খ) ৮২.৭৫\%

\textbf{\hl{গ) ২৮.৫৭\% ✔}}

ঘ) ২৮.৭৫\%

\textbf{\hl{✔ সঠিক উত্তর: গ) ২৮.৫৭\%}}

\begin{quote}
\textbf{ব্যাখ্যা:} \emph{মূল সমাধান:}

\emph{৪০\% বৃদ্ধিতে পূর্বের ১০০ টাকার গ্যাসের দাম বর্তমানে (১০০ + ৪০) বা ১৪০
টাকা।}

\emph{∴ ১৪০ টাকার স্থলে গ্যাস ব্যবহার করতে হবে ১০০ টাকার}

\emph{∴ ১ টাকার স্থলে গ্যাস ব্যবহার করতে হবে (১০০/১৪০) টাকার}

\emph{∴ ১০০ টাকার স্থলে গ্যাস ব্যবহার করতে হবে (১০০×১০০/১৪০) টাকার = ৭১.৪৩
টাকার (প্রায়)}

\emph{∴ ব্যবহার কমাতে হবে = (১০০ − ৭১.৪৩) টাকা = ২৮.৫৭ টাকার}

\emph{Shortcut:}

\emph{ব্যবহার হ্রাসের হার = (১০০×৪০/১০০+৪০)\% = (৪০০০/১৪০)\% = ২৮.৫৭\%}
\end{quote}

\textbf{কাপড়ের মূল্য ২০\% কমে গেলে কাপড়ের ব্যবহার শতকরা কি পরিমাণ বাড়ালে
কাপড়ের জন্য খরচের কোনো পরিবর্তন হবে না?}

ক) ১৫\%

খ) ২০\%

\textbf{\hl{গ) ২৫\% ✔}}

ঘ) ৩০\%

\textbf{\hl{✔ সঠিক উত্তর: গ) ২৫\%}}

\begin{quote}
\textbf{ব্যাখ্যা:} \emph{মূল সমাধান:}

\emph{কাপড়ের মূল্য ২০\% কমায়,}

\emph{বর্তমান মূল্য ৮০ টাকা হলে পূর্বমূল্য = ১০০ টাকা}

\emph{∴ বর্তমান মূল্য ১০০ টাকা হলে পূর্বমূল্য = (১০০×১০০/৮০) = ১২৫ টাকা}

\emph{∴ কাপড়ের ব্যবহার বাড়াতে হবে = (১২৫ − ১০০) টাকা = ২৫ টাকা বা ২৫\%।}

\emph{Shortcut:}

\emph{ব্যবহার বৃদ্ধির হার = (১০০×হ্রাস/১০০−হ্রাস)\% = (১০০×২০/১০০−২০)\% =
(২০০০/৮০)\% = ২৫\%}
\end{quote}

\textbf{চাউলের মূল্য ২০\% বৃদ্ধি পাওয়ায় ৭৫০ টাকায় পূর্বাপেক্ষা ৫ কেজি চাল কম
পাওয়া গেল। ১ কেজি চালের পূর্বমূল্য কত ছিল?}

ক) ১৫ টাকা

\textbf{\hl{খ) ২৫ টাকা ✔}}

গ) ৩৫ টাকা

ঘ) ৪৫ টাকা

\textbf{\hl{✔ সঠিক উত্তর: খ) ২৫ টাকা}}

\begin{quote}
\textbf{ব্যাখ্যা:} \emph{মূল সমাধান:}

\emph{২০\% বৃদ্ধিতে,}

\emph{বর্তমান মূল্য ১২০ টাকা হলে পূর্বমূল্য ১০০ টাকা}

\emph{∴ বর্তমান মূল্য ৭৫০ টাকা হলে পূর্বমূল্য = (১০০/১২০)×৭৫০ টাকা = ৬২৫ টাকা}

\emph{৫ কেজি চালের পূর্বমূল্য = (৭৫০ − ৬২৫) বা ১২৫ টাকা}

\emph{∴ ১ কেজি চালের পূর্বমূল্য = (১২৫/৫) বা ২৫ টাকা।}

\emph{Shortcut:}

\emph{বর্তমান মোট টাকা = ৭৫০}

\emph{বৃদ্ধি পাওয়া মূল্য (২৫০ টাকার ২০\%) = ৭৫০×(২০/১০০) = ১৫০ টাকা}

\emph{বর্তমানে ৫ কেজি চাল পাওয়া যায় ১৫০ টাকায় ∴ ১ কেজি চালের বর্তমান মূল্য =
৩০ টাকা}

\emph{পূর্বমূল্য = ৩০×(১০০/১২০) = ২৫ টাকা।}
\end{quote}

\textbf{পানির উপর ৩০\% কর বাড়ানো হলে একটি পরিবার এর ব্যবহার ৩০\% কমিয়ে
দিল। ঐ পরিবারের জন্য পানির খরচ শতকরা কত কমল বা বাড়ল?}

ক) ৯\% বাড়ল

খ) ৭.৭৫\% কমলো

গ) ৮.২৫\% বাড়ল

\textbf{\hl{ঘ) ৯\% কমলো ✔}}

\textbf{\hl{✔ সঠিক উত্তর: ঘ) ৯\% কমলো}}

\begin{quote}
\textbf{ব্যাখ্যা:} \emph{মূল সমাধান:}

\emph{ধরি, পূর্বে পানির ব্যবহার ১০০ একক এবং একক প্রতি কর ১০০ টাকা ∴ মোট খরচ
= (১০০ × ১০০) = ১০,০০০ টাকা}

\emph{কর ৩০\% বাড়লে নতুন কর = (১০০ + ৩০) = ১৩০ টাকা}

\emph{ব্যবহার ৩০\% কমলে নতুন ব্যবহার = (১০০ − ৩০) = ৭০ একক}

\emph{নতুন মোট খরচ = (১৩০ × ৭০) = ৯,১০০ টাকা}

\emph{∴ খরচ কমলো = ১০,০০০ − ৯,১০০ = ৯০০ টাকা}

\emph{শতকরা হ্রাস = (৯০০×১০০/১০০০০)\% = ৯\% কমলো।}

\emph{Shortcut:}

\emph{খরচ হ্রাস/বৃদ্ধির সূত্র: A+B+(AB/১০০)}

\emph{এখানে, A = +৩০ (বৃদ্ধির হার), B = −৩০ (হ্রাসের হার)}

\emph{∴ খরচের পরিবর্তন = ৩০−৩০+(৩০×(−৩০)/১০০) = −৯\%}

\emph{অর্থাৎ খরচ ৯\% কমবে।}
\end{quote}

\textbf{কলার দাম ২০\% কমে যাওয়ার ১২ টাকার পূর্ব অপেক্ষা ২টি কলা বেশি পাওয়া
গেলে বর্তমানে একটি কলার দাম কত টাকা?}

\textbf{\hl{ক) ১.২০ ✔}}

খ) ২.৫০

গ) ৩.০০

ঘ) ৪.০০

\textbf{\hl{✔ সঠিক উত্তর: ক) ১.২০}}

\begin{quote}
\textbf{ব্যাখ্যা:} \emph{মূল সমাধান:}

\emph{১০০ টাকায় কমে ২০ টাকা}

\emph{∴ ১২ টাকায় কমে = (২০×১২/১০০) = ২.৪ টাকা}

\emph{২ টি কলার বর্তমান মূল্য ২.৪ টাকা}

\emph{∴ ১ টি কলার বর্তমান মূল্য = (২.৪×১/২) = ১.২ টাকা}

\emph{Shortcut:}

\emph{বর্তমান একটি কলার দাম = (১২×২০\%/২) = (২.৪/২) = ১.২ টাকা।}
\end{quote}

\textbf{গমের মূল্য ১৫\% কমে যাওয়ায় ৬,০০০ টাকায় পূর্বাপেক্ষা ১ কুইন্টাল গম বেশি
পাওয়া যায়। ১ কেজি গমের বর্তমান মূল্য কত?}

ক) ৬ টাকা

খ) ৭ টাকা

গ) ৮ টাকা

\textbf{\hl{ঘ) ৯ টাকা ✔}}

\textbf{\hl{✔ সঠিক উত্তর: ঘ) ৯ টাকা}}

\begin{quote}
\textbf{ব্যাখ্যা:} \emph{মূল সমাধান:}

\emph{আমরা জানি, ১ কুইন্টাল = ১০০ কেজি।}

\emph{৬,০০০ টাকার ১৫\% = ৬০০০×(১৫/১০০) = ৯০০ টাকা}

\emph{১০০ কেজি গমের বর্তমান মূল্য = ৯০০ টাকা}

\emph{∴ ১ কেজি গমের বর্তমান মূল্য = (৯০০/১০০) টাকা = ৯ টাকা।}

\emph{Shortcut:}

\emph{১ কেজি গমের বর্তমান মূল্য = (৬০০০×১৫\%/১০০) = (৯০০/১০০) = ৯ টাকা।}
\end{quote}

\textbf{গতকাল শেয়ারের দাম ২৫\% বেড়েছিল। কিন্তু আজ ২৫\% কমেছে। শেয়ারের দাম
মোট কত বেড়েছে বা কমেছে?}

\textbf{\hl{ক) ৬ (১/৪) \% কমেছে ✔}}

খ) ৬ (১/৪) \% বেড়েছে

গ) ৪ (১/২) \% কমেছে

ঘ) ৪ (১/২) \% বেড়েছে

\textbf{\hl{✔ সঠিক উত্তর: ক) ৬ (১/৪) \% কমেছে}}

\begin{quote}
\textbf{ব্যাখ্যা:} \emph{মূল সমাধান:}

\emph{ধরি, শেয়ারের প্রাথমিক দাম ছিল ১০০ টাকা}

\emph{২৫\% বাড়লে দাম হয় = (১০০ + ২৫) = ১২৫ টাকা}

\emph{আজ ১২৫ টাকার ২৫\% কমলে দাম কমে = ১২৫×(২৫/১০০) = ৩১.২৫ টাকা}

\emph{বর্তমান দাম = (১২৫ − ৩১.২৫) = ৯৩.৭৫ টাকা}

\emph{∴ মোট দাম কমেছে = (১০০ − ৯৩.৭৫) = ৬.২৫\% বা ৬ (১/৪) \%}

\emph{Shortcut:}

\emph{পরিবর্তন \% = A+B+(AB/১০০)}

\emph{= ২৫−২৫−(২৫×২৫/১০০)}

\emph{= −(২৫/৪) = −৬(১/৪)\% (অর্থাৎ ৬ (১/৪) \% কমেছে)।}
\end{quote}

\textbf{বর্তমানে ৬ কেজি চালের দাম আগের পাঁচ কেজি চালের দামের সমান হলে,
চালের দাম শতকরা কত কমেছে?}

\textbf{\hl{ক) ১৬.৬৬\% ✔}}

খ) ২৫\%

গ) ৩০\%

ঘ) ১৮\%

\textbf{\hl{✔ সঠিক উত্তর: ক) ১৬.৬৬\%}}

\begin{quote}
\textbf{ব্যাখ্যা:} \emph{মূল সমাধান:}

\emph{ধরি, পূর্বে ৫ কেজি চালের দাম ছিল ১০০ টাকা}

\emph{∴ পূর্বে ১ কেজি চালের দাম ছিল ২০ টাকা}

\emph{প্রশ্নমতে, বর্তমানে ৬ কেজি চালের দাম = ১০০ টাকা}

\emph{∴ বর্তমানে ১ কেজি চালের দাম = (১০০/৬) = ১৬.৬৬ টাকা}

\emph{প্রতি কেজিতে চালের দাম কমেছে = (২০ − ১৬.৬৬) টাকা = ৩.৩৪ টাকা}

\emph{∴ শতকরা দাম কমেছে = (৩.৩৪/২০)×১০০\% = ১৬.৬৬\% বা ১৬ (২/৩) \%}

\emph{Shortcut:}

\emph{শতকরা কমেছে = (৬−৫/৬)×১০০\% = ১৬.৬৬\%}
\end{quote}

\textbf{চিনির মূল্য ৬\% বেড়ে যাওয়ায় ৭৯৫ টাকায় পূর্বে যত কেজি চিনি পাওয়া যেত,
বর্তমানে তার চেয়ে ৩ কেজি চিনি কম পাওয়া যায়। চিনির বর্তমান দর প্রতি কেজি
কত?}

ক) ১৫ টাকা

খ) ১৬ টাকা

গ) ১৮ টাকা

\textbf{\hl{ঘ) ১৫.৯০ টাকা ✔}}

\textbf{\hl{✔ সঠিক উত্তর: ঘ) ১৫.৯০ টাকা}}

\begin{quote}
\textbf{ব্যাখ্যা:} \emph{মূল সমাধান:}

\emph{চিনির মূল্য ৬\% বৃদ্ধি পাওয়ায়,}

\emph{৭৯৫ টাকার ৬\% = ৭৯৫×(৬/১০০) টাকা = ৪৭.৭০ টাকা}

\emph{এই ৪৭.৭০ টাকা বৃদ্ধির কারণে ৩ কেজি চিনি কম পাওয়া যায়}

\emph{∴ ৩ কেজি চিনির বর্তমান মূল্য = ৪৭.৭০ টাকা}

\emph{∴ ১ কেজি চিনির বর্তমান মূল্য = (৪৭.৭০/৩) টাকা = ১৫.৯০ টাকা।}

\emph{Shortcut:}

\emph{প্রতি কেজি চিনির বর্তমান দর = (৭৯৫×৬\%/৩) = (৪৭.৭০/৩) = ১৫.৯০ টাকা।}
\end{quote}

\textbf{একটি পণ্যের দাম প্রথমে ৪০\% বাড়ানো হলো; তারপর ১০\% কমানো হলে মোটের
উপর কত \% বাড়লো?}

ক) ২৫\%

\textbf{\hl{খ) ২৬\% ✔}}

গ) ২৭\%

ঘ) ২৮\%

\textbf{\hl{✔ সঠিক উত্তর: খ) ২৬\%}}

\begin{quote}
\textbf{ব্যাখ্যা:} \emph{মূল সমাধান:}

\emph{ধরি, দ্রব্যটির দাম ছিল ১০০ টাকা}

\emph{৪০\% বৃদ্ধিতে দাম হবে = (১০০ + ৪০) = ১৪০ টাকা}

\emph{১০\% কমে বর্তমান দাম হবে = ১৪০ − (১৪০×১০/১০০) = ১২৬ টাকা}

\emph{∴ দাম শতকরা বাড়ে = ১২৬ − ১০০ = ২৬\%}

\emph{Shortcut:}

\emph{কার্যকর পরিবর্তন \% = A+B+(AB/১০০)}

\emph{এখানে, A = +৪০, B = −১০}

\emph{= ৪০−১০+(৪০×(−১০)/১০০) = ৩০−৪ = ২৬\%}
\end{quote}

\textbf{একজন টিভি বিক্রেতা তার প্রতিটি টিভির দাম x\% বৃদ্ধি করল। এতে তার
বিক্রিত টিভির সংখ্যা y\% কমে গেল। কিন্তু এতেও তার মোট বিক্রয়মূল্য অপরিবর্তিত
থাকল। নিচের কোন সমীকরণটি y এর সাপেক্ষে x-এর সঠিক প্রকাশ?}

ক) y/(1-x)

\textbf{\hl{খ) 100y/(100-y) ✔}}

গ) 100/(100 -y)

ঘ) 100y/(100-x)

\textbf{\hl{✔ সঠিক উত্তর: খ) 100y/(100-y)}}

\begin{quote}
\textbf{ব্যাখ্যা:} \emph{মূল সমাধান:}

\emph{ধরি, প্রতিটি টিভির দাম ১০০ টাকা এবং বিক্রির টিভির সংখ্যা ১০০ টি।}

\emph{∴ x\% বৃদ্ধিতে নতুন মূল্য = (100 + x) টাকা}

\emph{বিক্রি y \% কমে গেলে বিক্রিত টিভির সংখ্যা = (100 − y) টি}

\emph{প্রশ্নমতে, (100 + x)(100 − y) = 100 × 100}

\emph{⇒ 10000 − 100y + 100x − xy = 10000}

\emph{⇒ x(100 − y) = 100y}

\emph{∴ x = (100y/100−y)}

\emph{Shortcut:}

\emph{ব্যয় বা বিক্রি অপরিবর্তিত থাকলে বৃদ্ধির হার x = (100×হ্রাস/100−হ্রাস) =
(100y/100−y)}
\end{quote}

\textbf{মিজান একটি ঘড়ি কিনতে গিয়ে দেখল যে এটা খুব দামী। ৭ দিন পর ঘড়িটির
দাম ১৮\% কমে গেল। কিন্তু তার ২ দিন পর ঘড়িটির দাম আবার ২৫\% বেড়ে গেল এবং
মিজান ঘড়িটি কিনতে পারল না। যদি ঘড়িটির দাম ১৮\% কমার পর ৪,১০০ টাকা হয়ে
থাকে ঘড়িটির মূল দাম ও বর্তমান দামের পার্থক্য কত টাকা?}

\textbf{\hl{ক) ১২৫ ✔}}

খ) ১৩২

গ) ২০৫

ঘ) ৫১২.৫

\textbf{\hl{✔ সঠিক উত্তর: ক) ১২৫}}

\begin{quote}
\textbf{ব্যাখ্যা:} \emph{মূল সমাধান:}

\emph{১৮\% কমার পর ঘড়িটির দাম = ৪,১০০ টাকা}

\emph{ধরি, মূল দাম x টাকা}

\emph{x এর (১০০ − ১৮)\% = ৪১০০}

\emph{⇒ x × ৮২\% = ৪১০০}

\emph{⇒ x = (৪১০০×১০০/৮২) = ৫,০০০ টাকা (মূল দাম)}

\emph{আবার, ২৫\% বৃদ্ধির পর বর্তমান দাম = ৪১০০ এর (১০০ + ২৫)\% =
(৪১০০×১২৫/১০০) = ৫,১২৫ টাকা}

\emph{∴ মূল দাম ও বর্তমান দামের পার্থক্য = ৫,১২৫ − ৫,০০০ = ১২৫ টাকা।}
\end{quote}

\textbf{জুন মাসে প্রতি কেজি ফজলি ও প্রতি কেজি ল্যাংড়া আমের দাম একই ছিল। জুলাই
মাসে প্রতি কেজি ফজলি আমের দাম ৪০\% বৃদ্ধি পেল এবং প্রতি কেজি ল্যাংড়া আমের দাম
২০\% হ্রাস পেল। যদি জুলাই মাসে সমান পরিমাণ ফজলি ও ল্যাংড়া আমের মিশ্রণের প্রতি
কেজির দাম ৭৭ টাকা হয়, জুন মাসে এক কেজি ল্যাংড়া আমের দাম কত টাকা ছিল?}

ক) ৭৭

খ) ৭৫

\textbf{\hl{গ) ৭০ ✔}}

ঘ) ৬৬

\textbf{\hl{✔ সঠিক উত্তর: গ) ৭০}}

\begin{quote}
\textbf{ব্যাখ্যা:} \emph{মূল সমাধান:}

\emph{ধরি,}

\emph{জুন মাসে প্রতি কেজি ফজলি আমের দাম ছিল x টাকা}

\emph{∴ জুন মাসে প্রতি কেজি ল্যাংড়া আমের দাম ছিল x টাকা}

\emph{জুলাই মাসে প্রতি কেজি ফজলি আমের দাম = x+xএর(৪০/১০০) টাকা =
(১৪০x/১০০) = (১৪x/১০) টাকা}

\emph{∴ জুলাই মাসে প্রতি কেজি ল্যাংড়া আমের দাম = x−xএর(২০/১০০) টাকা =
(৮০x/১০০) = (৮x/১০) টাকা}

\emph{জুলাই মাসে ১ কেজি ফজলি + ১ কেজি ল্যাংড়া = ২ কেজি আমের মোট দাম =
(১৪x/১০)+(৮x/১০) টাকা = (২২x/১০) টাকা}

\emph{প্রশ্নমতে,}

\emph{(২২x/১০)=৭৭×২}

\emph{⇒ x=(১৫৪০/২২)}

\emph{∴ x = ৭০ টাকা।}

\emph{Shortcut:}

\emph{ধরি, জুন মাসে প্রতি কেজি আমের দাম ছিল ১০-এর গুণিতক = ১০০x}

\emph{জুলাইয়ে ফজলির দাম = ১৪০x, ল্যাংড়ার দাম = ৮০x}

\emph{মিশ্রণে ২ কেজি আমের মোট দাম = (১৪০x + ৮০x) = ২২০x}

\emph{১ কেজির গড় দাম = ১১০x}

\emph{প্রশ্নমতে, ১১০x = ৭৭}

\emph{∴ ১০০x (জুন মাসের মূল দাম) = (৭৭×১০০/১১০) = ৭০ টাকা।}
\end{quote}

\section{52. শতকরা কি ? এখন শিখুন দুর্দান্ত কেীশলে
!!!}\label{ux9b6ux9a4ux995ux9b0-ux995-ux98fux996ux9a8-ux9b6ux996ux9a8-ux9a6ux9b0ux9a6ux9a8ux9a4-ux995ux9b6ux9b2}

শতকরা, লাভ-ক্ষতি ও সুদকষার হিসাব গণিতের এমন কিছু অংশ যা ছাড়া প্রতিযোগিতামূলক
পরীক্ষা বা দৈনন্দিন জীবনের হিসাব-নিকাশ কল্পনা করা যায় না। অথচ এই বিষয়গুলোকেই
অনেকে সবচেয়ে ভয়ের এবং মুখস্থকরণের বিষয় বলে মনে করেন। কিন্তু একটু গভীরভাবে
লক্ষ্য করলে দেখবেন, পাটিগণিতের শত শত জটিল অংককে কেবল `১০০\%'-এর ধারণা
ব্যবহার করে খুব সহজেই এক লাইনে সমাধান করা সম্ভব।

অযথা শর্টকাট বা অন্ধভাবে সূত্র মুখস্থ করার পেছনে না দৌড়ে, যদি আপনি এই ১০০\%-এর
মূল রহস্য বুঝতে পারেন, তবে যেকোনো কঠিন সমস্যা আপনার কাছে পানির মতো সহজ মনে
হবে। শুরুতে হয়তো একটু সময় লাগতে পারে, কিন্তু একবার কনসেপ্ট ক্লিয়ার হয়ে গেলে
পরীক্ষার হলে ম্যাজিকের মতো দ্রুত উত্তর বের করতে পারবেন। পাটিগণিতের এই অত্যন্ত
গুরুত্বপূর্ণ অধ্যায়টির বিস্তারিত আলোচনা, শর্টকাট টেকনিক ও বিগত সালের প্রশ্নগুলোর
সঠিক ব্যাখ্যা জানতে শেষ পর্যন্ত মনোযোগ দিয়ে পড়তে থাকুন।

\begin{itemize}
\item
  প্রাথমিক আলোচনা
\item
  পণ্যদ্রব্যের মূল্য বৃদ্ধি/হ্রাস এবং ব্যবহার বৃদ্ধি/হ্রাস সংক্রান্ত সমস্যা
\item
  পরীক্ষার্থী ও ‪পরীক্ষা বিষয়ক সংক্রান্ত সমস্যা ‬‬‬
\end{itemize}

\subsubsection{প্রাথমিক
আলোচনা}\label{ux9aaux9b0ux9a5ux9aeux995-ux986ux9b2ux99aux9a8-13}

\textbf{১০০\% আসলে কি?}

১০০\% হলো যে কোন কিছুর প্রাথমিক মান বা অপরিবর্তিত মান। যেমন: যেমন একটি
ক্লাসে ১০০ জনে ১০০ জনই উপস্থিত থাকলে আমরা বলি ১০০\% উপস্থিতি। আবার কম হলে
বলি ৯০\% উপস্থিত আবার কোন দ্রব্যের দাম অত্যাধিক বেড়ে গেলে আমরা বলি ২০০\%
বেড়ে গেছে। অর্থাৎ বাড়া বোঝালে ১০০\% এর সাথে যোগ আবার কমা বোঝালে ১০০\%
থেকে বিয়োগ হয়।

\textbf{কিছু প্রশ্নে ১০০\% এর প্রয়োগ}

\textbf{প্রশ্নঃ ২৫০ টাকার একটি শার্ট ১০\% লাভে বিক্রি করা হলে শার্টটির
বিক্রয়মূল্য কত?}

সরাসরি ২৫০ এর ১১০\% এর মান বের করতে হবে। কেননা এখানে ২৫০ টাকা হল ক্রয়মূল্য
যাতে লাভ বা ক্ষতি কোন কিছুই যুক্ত না থাকায় তা ১০০\% এবং লাভের ১০\% যুক্ত হয়ে
গেলে ১১০\% হবে।

\textbf{প্রশ্নঃ ৩০ জন ছাত্রের মধ্যে ২০\% অনুপস্থিত থাকলে মোট কতজন অনুপস্থিত?}

\textbf{ব্যাখ্যা:} ৩০ এর ২০\% = ৬ জন।

কিন্তু যদি এই প্রশ্নটিই উল্টোভাবে আসে এভাবে-

\textbf{প্রশ্নঃ একটি ক্লাশে কোন একদিন ৬ জন ছাত্র অনুপস্থিত ছিল। যদি এই ৬ জন
মোট ছাত্রছাত্রীর ২০\% হয়, তাহলে ঐ ক্লাশে মোট কতজন ছাত্র আছে?}

এখানে বুঝে বুঝে করলে এভাবে অনুপস্থিত এর \% হল ২০\% এবং এই ২০\% এর মান হল ৬।

তাই লিখা যায় ২০\% = ৬

অতএব, ১\% = \(\frac{\text{৬}}{\text{২০}}\)

সুতরাং ১০০\% = \(\frac{\text{৬} \times \text{১০০}}{\text{২০}}\) = ৩০জন।

\textbf{নোট:} যে কোন \% এর অংকে \% বাদে যে সংখ্যাটি দেয়া থাকবে তা কত
অংশের মান তা নিয়ে ভাবা শুরু করলে দ্রুত অংক হবে।

এই অংকটিই এক লাইনে করতে চাইলে-

সরাসরি \(\frac{\text{৬} \times \text{১০০}}{\text{২০}}\)

\textbf{নোট:} যে কোন \% এর মান দেয়া থাকলে ঐ \% এর মান লিখে \% টিকে
উল্টিয়ে লিখে গুণ করলেই উত্তর বের হয়ে যাবে।)

প্রমাণ দেখুন এখানে:

\textbf{প্রশ্নঃ ফাহিমের কাছে কিছু মারবেল ছিল, যার ২৫\% সে তার ভাইকে দেয়ায়
তার ভাই ৭৫টি মারবেল পেল। ফাহিমের কাছে মোট কতটি মারবেল ছিল?}

\textbf{সমাধান:} (\(\frac{\text{৭৫} \times \text{১০০}}{\text{২৫}}\) ) = ৩০০টি ) (২৫\% কে
\(\frac{\text{২৫}}{\text{১০০}}\) না লিখে উল্টিয়ে \(\frac{\text{১০০}}{\text{২৫}}\) লিখে ৭৫ এর সাথে গুণ
করা হয়েছে)

\textbf{প্রশ্নঃ একজন ফল বিক্রেতা ৩০\% আপেল বিক্রি করল এবং তার নিকট আরও ৬৩০
টি আপেল আছে। সে সর্বমোট কতগুলো আপেল কিনেছিল?}

\textbf{সমাধান:} (এখানে ৬৩০ হল ৭০\% এর মান )

তাই ৬৩০× \(\frac{\text{১০০}}{\text{৭০}}\) = ৯০০ (৭০\% কে \(\frac{\text{৭০}}{\text{১০০}}\) না লিখে
উল্টিয়ে \(\frac{\text{১০০}}{\text{৭০}}\) লিখে ৬৩০ এর সাথে গুণ করা হয়েছে)

সুতরাং সে সর্বমোট ৯০০ টি আপেল কিনেছিল।

প্রশ্নঃ একটি ভোটকেন্দ্রে উপস্থিত ভোটারদের ৫৫\% ভোট পেয়ে একজন প্রার্থী নির্বাচিত
হয়েছেন। তিনি তাঁর একমাত্র প্রতিদ্ব›দ্বী অপেক্ষা ১০,০০০ ভোট বেশি পেয়েছেন।
ভোটকেন্দ্রে কতজন ভোটার উপস্থিত ছিল?

\textbf{সমাধানঃ} দুজনের ভোটের পার্থক্য ৫৫\%-৪৫\% = ১০\% যার মান ১০০০০।

অতএব, ভোটকেন্দ্রে ভোটার উপস্থিত ছিল = ১০০০০×\(\frac{\text{১০০}}{\text{১০}}\) = ১০০০০০ জন।
(১০\% কে \(\frac{\text{১০}}{\text{১০০}}\) না লিখে উল্টিয়ে \(\frac{\text{১০০}}{\text{১০}}\) লিখে ১০০০০
এর সাথে গুণ করা হয়েছে)

\textbf{লাভ-ক্ষতি অধ্যায়ে ১০০\% এর প্রয়োগ:}

\textbf{প্রশ্নঃ ১৫০ টাকায় একটি জিনিস ক্রয় করে কত দামে বিক্রয় করলে ৩০\% লাভ
হবে?}

\textbf{ব্যাখ্যাঃ} সরাসরি ১৫০ এর ১৩০\% এর মান বের করতে হবে। কেননা এখানে
১৫০ টাকা হল ক্রয়মূল্য যাতে লাভ বা ক্ষতি কোন কিছুই যুক্ত না থাকায় তা ১০০\% এবং
লাভের ৩০\% যুক্ত হয়ে গেলে ১৩০\% হবে।

অতএব, ১৫০× \(\frac{\text{১৩০}}{\text{১০০}}\) = ১৯৫

সুতরাং ১৫০ টাকায় একটি জিনিস ক্রয় করে ১৯৫ টাকায় বিক্রয় করলে ৩০\% লাভ হবে।

\textbf{প্রশ্নঃ একটি কলম ২৭০ টাকায় বিক্রয় করাতে ১০\% ক্ষতি হয়; কলমটির ক্রয়মূল্য
কত?}

\textbf{ব্যাখ্যাঃ} ৩০০(এখানে ২৭০ হল ৯০\% এর মান কারন ১০\% ক্ষতি হয়েছে। তাই
২৭০×৯০\% কে উল্টিয়ে = ২৭০×\(\frac{\text{১০০}}{\text{৯০}}\) = ৩০০ । {[}লাভ-ক্ষতির অংকে
ক্রয়মূলটাই ১০০\%{]}

\textbf{প্রশ্নঃ কোন বই ৪০ টাকায় বিক্রি করলে ২০\% ক্ষতি হয়। কত টাকায় বিক্রি
করলে ৪০\% লাভ হবে?}

(হেল্প: ৮০\% = ৪০ , সুতরাং ১৪০\% = ?) ৭০টাকা।

\textbf{প্রশ্নঃ এক ঝুড়ি আম ১০\% ক্ষতিতে বিক্রয় করা হল। বিক্রয়মূল্য ৪৫ টাকা বেশি
হলে ৫\% লাভ হতো। এক ঝুড়ি আমের দাম কত?}

\textbf{\hl{উত্তর: ৩০০টাকা (ব্যাখ্যা: ১০\% ক্ষতি পোষাতে হবে তারপর আবার ৫\%
লাভ করতে হবে অর্থাৎ মোট বৃদ্ধি করতে হবে ১০\%+৫\% = ১৫\% যার টাকায় মান হল
৪৫। সুতরাং ১\% = ৩ এবং ১০০\% = ৩০০}}

\textbf{শতকরার শতকরা কি?}

ধরুন আপনার কাছে ২০০ টাকা আছে যার ৫০\% আপনার ভাইকে এবং ৫০\% আপনার বোনকে
দিলেন, তাহলে আপনার ভাই ও বোন প্রত্যেকে ১০০ টাকা করে পাবে। কিন্তু যদি এভাবে
বলে যে আপনার ভাইকে ৫০\% টাকা দেয়ার পর অবশিষ্টের ৫০\% আপনার বোনকে দিলেন
তাহলে প্রথমে ভাইকে ৫০\% অর্থাৎ ১০০ টাকা দেয়ার পর যে ১০০ টাকা থাকবে তার
৫০\% অর্থাৎ ৫০ টাকা আপনার বোন পাবে।

\textbf{নোট:} সাধারণ শতকরা আর অবশিষ্টের শতকরা এক বিষয় না (অনেকেই ভুল করেন)

→শতকরা বের করতে বললে \textgreater\textgreater{} ৮০ এর ২০\% = কত? সরাসরি
৮০×\(\frac{\text{২০}}{\text{১০০}}\) = ১৬ সবাই জানি।

→কিন্তু ১৬ কোন সংখ্যার ২০\%? এরকম যদি আসে তখন অনেকে `ক\textquotesingle{}
ধরে করেন । কিন্তু তাতে মাথা সরাসরি কাজ করে না। মুখে মুখে করতে চাইলে এভাবে
ভাবুন: ৮০ এর ২০\% যদি ১৬ হয় তাহলে সংখ্যাটি ১০০\% ছিল।

তাই লেখা যায় ২০\% = ১৬

অতএব, ১\% = \(\frac{\text{১৬}}{\text{২০}}\)

সুতরাং ১০০\% = ১০০×\(\frac{\text{১৬}}{\text{২০}}\) = ৮০

আবার একলাইনে করতে চাইলে এভাবে লিখুন ১৬×\(\frac{\text{১০০}}{\text{২০}}\) = ৮০

\textbf{কয়েকটি নিজে চেষ্টা করুন:}

\textbf{খুব খুবই গুরুত্বপূর্ণ নোট:}

যে কোন প্রশ্ন \% বাদে কোন সংখ্যা দেয়া থাকলে এবং সেই সংখ্যাটি কত \% এর মান তা
বের করতে পারলে এবং মুল সংখ্যাটির মান বের করার জন্য প্রশ্নে প্রদত্ত মানটি লিখে
যে \% টি দেয়া থাকবে তা উল্টিয়ে লিখতে হবে যেমন: ৩৬ কোন সংখ্যার ৪০\% ? তখন
৩৬×\(\frac{\text{১০০}}{\text{৪০}}\) = ৯০ এর।

দুবার শতকরা আসলেও একই নিয়মটি দুবার প্রয়োগ করতে হবে।

শতকরার শতকরা সরাসরি বের করা শিখুন, প্রথমে

৮০ টাকার ২০\% এর ২৫\% এর মান কত। সরাসরি ৮০ এর ২০\% এর ২৫\% বা
৮০×\(\frac{\text{২০}}{\text{১০০}}\)×\(\frac{\text{২৫}}{\text{১০০}}\) = ৪ (এভাবে সবাই পারেন)

কিন্তু বিপরীত দিক থেকে আসলে অনেকে পারেনই না.. পারলেও এক্স ধরে করতে অনেক সময়
লাগে। কিন্তু দ্রুত উত্তর বের করার জন্য উপরের উল্টানোর নিয়মটিই দুবার প্রয়োগ করতে
হবে।

\textbf{যেমন: ৪ কোন সংখ্যার ২০\% এর ২৫\% এর মান হলে সংখ্যাটি কত?}

লিখুন এভাবে ৪×\(\frac{\text{১০০}}{\text{২৫}}\)×\(\frac{\text{১০০}}{\text{২০}}\) (উল্টানো অর্থ সবসময়
উপরে ১০০) কাটাকাটি করলে ৮০ ই আসবে।

প্রশ্নঃ কোন বিশ্ববিদ্যালয়ে ৮০\% শিক্ষার্থী আবাসিক হলে অবস্থান করে এবং হলের
সৌভাগ্যবান ৬০\% শিক্ষার্থী একক কক্ষ পায়। যদি সৌভাগ্যবান শিক্ষার্থীর সংখ্যা ১২০০
জন হয় তবে ঐ বিশ্ববিদ্যালয়ে মোট শিক্ষার্থীর সংখ্যা কত?

\textbf{সমাধান:} ১২০০×\(\frac{\text{১০০}}{\text{৬০}}\)×\(\frac{\text{১০০}}{\text{৮০}}\) = ২৫০০ (কারণ
শেষের ১২০০ হল ৬০\% এর মান)

ঐ বিশ্ববিদ্যালয়ে মোট শিক্ষার্থীর সংখ্যা = ২৫০০ জন।

আবার লাভ ক্ষতি অধ্যায়ে কখনো ক্রয়মূল্য বের করতে বললে ১০০\% এর মান বের করতে হয়।
তখন উল্টাতে হয়। কিন্তু তার আগে প্রশ্নে যে সংখ্যাটি দেয়া থাকবে তা কত \% এর মান
তা বের করতে পারাই আসল।

\textbf{প্রশ্নঃ একজন পাইকারী দোকানদার একটি পণ্যের লিখিত মুল্যের উপর ১০\% ছাড়ে
পণ্যটি বিক্রি করে আবার খুচরা বিক্রেতা ঐ পণ্যটি ১৫\% ক্ষতিতে ৩০৬ টাকায় বিক্রি
করলে পণ্যটির লিখিত মূল্য কত ছিল?}

\textbf{ব্যাখ্যা:} সবার শেষে ১৫\% ক্ষতিতে ৩০৬ টাকায় বিক্রি হয়েছে তাই ৩০৬
টাকা হলো ৮৫\% এর মান: আবার ১০\% ছাড় দেয়ায় খুচরা বিক্রেতার ক্রয়মূল্য টি ছিল
লিখিত মূল্যের ৯০\% এর মান। তাই এক লাইনে করতে এভাবে লিখুন:
৩০৬×\(\frac{\text{১০০}}{\text{৮৫}}\)×\(\frac{\text{১০০}}{\text{৯০}}\) = ৪০০। (আলাদা ভাবে করলে অথবা
এক্স ধরে করলে ও কিন্তু একই উত্তর আসবে কিন্তু ততক্ষণে পরীক্ষক আপনার খাতা নিয়ে
নিবে।)

\textbf{আরো কয়েকটি দেখুন:}

\textbf{প্রশ্নঃ এক ঝুড়ি আম ১০\% ক্ষতিতে বিক্রয় করা হল। বিক্রয়মূল্য ৪৫ টাকা বেশি
হলে ৫\% লাভ হতো। এক ঝুড়ি আমের দাম কত?}

\textbf{ব্যাখ্যা:} ১০\% ক্ষতি পুষিয়ে ৫\% লাভ করতে হলে মাঝের ব্যবধান ১০+৫ =
১৫\% যার মান ৪৫। তাহলে ক্রয়মূল্য বের করতে হলে ৪৫×\(\frac{\text{১০০}}{\text{১৫}}\) =
৩০০টাকা।

\textbf{প্রশ্নঃ একটি কাজ ১৮ দিনে শেষ করার চুক্তি নেয়। কিন্তু ৯ জন অনুপস্থিত
থাকায় কাজ টি শেষ করতে ৩৬ দিন লাগে। ৩৬ জন লোকে কাজটি করতে কত সময় লাগবে?}

\textbf{ব্যাখ্যা:} ১৮দিনে শেষ করার কথা থাকলেও ৩৬ লাগলো অর্থাৎ দ্বিগুণ সময়
লাগলো কিন্তু কেন?? ৯ জন চলে যাওয়ায়। একটা কাজে দ্বিগুণ সময় তখনি লাগে যখন
শ্রমিকের সংখ্য অর্ধেক হয় তাহলে ঐ ৯ জন ছিল মোট লোকের অর্ধেক সুতরাং মোট লোক ছিল
১৮ জন। এখন ১৮ জন থাকলে করতে পারতো ১৮দিনেই সুতরাং ৩৬ জন (দ্বিগুণ লোক) করতে
পারবে ১৮দিনের অর্ধেক সময়ে অর্থাৎ ৯দিনে।

\textbf{গুরুত্বপূর্ণ শিক্ষণীয় বিষয়ঃ}

প্রশ্নে যে সংখ্যা দেয়া থাকবে তা কত \% এর মান সেটা সবার শুরুতে বোঝার চেষ্টা
করতে হবে।

প্রশ্নঃ একটি নির্বাচনী এলাকায় মোট ভোটার সংখ্যার ৮\% বৃদ্ধ ভোটার। একটি নির্বাচনে
বৃদ্ধ ভোটারের ৮০\% ভোট প্রদান করলো। ভোট প্রদানকারী বৃদ্ধদের সংখ্যা মোট ভোটারের
সংখ্যার কত অংশ? ৬.৪\% (১০০ ধরে করুন, ১০০ এর ৮\% এর ৮০\%।)

\textbf{ব্যাখ্যাঃ} মোট ভোটার ১০০\% ধরে করতে হবে।তাহলে ১০০\% এর ৮\% হল বৃদ্ধ
ভোটার।আবার ঐ ৮\% এর ৮০\% ভোট প্রদান করলো। তাই লিখতে হবে এভাবে, ১০০ এর ৮\%
এর ৮০\%।

সুতরাং ১০০×\(\frac{\text{৮}}{\text{১০০}}\)×\(\frac{\text{৮০}}{\text{১০০}}\) = \(\frac{\text{৬৪}}{\text{১০}}\) = ৬.৪

প্রশ্নঃ কোনো এক নির্বাচন কেন্দ্রে ভোট দাতাদের ৯৫\% উপস্থিত ছিল। দুজন প্রার্থীর
একজন উপস্থিত ভোটারের ৫৪\% ভোট পাওয়ায় দেখা গেল যে, সে অপর প্রার্থী অপেক্ষা
১৫২ ভোট বেশি পেয়েছে। মোট ভোটার কত? উত্তর: ২০০০ জন

লক্ষ্য করুন: এখানে যে ৯৫\% ভোটারের কথা বলা হয়েছে তা ছাড়াও ঐ কেন্দ্রে আরো ৫\%
ভোটার আছে। তাই প্রথমে উপস্থিত ভোটারের সংখ্যা বের করার পর মোট ভোটারের সংখ্যা
বের করতে হবে।

৮\% = ১৫২টি,

তাই ১\% = ১৫২÷৮= ১৯

এবং ১০০\% =১৯×১০০= ১৯০০

১৯০০ হলো উপস্থিত ভোটার কিন্তু উপস্থিত+অনুপস্থিত ভোটার বের করতে হবে। তাই

৯৫\% = ১৯০০,

১\% =১৯০০÷৯৫= ২০

এবং ১০০\% = ২০×১০০ = ২০০০

\textbf{\hl{উত্তর: ২০০০ জন।}}

\textbf{এবার নিজে চেষ্টা করুন:}

প্রশ্নঃ একটি নির্বাচনী এলাকায় মোট ভোটারের ৭৫\% উপস্থিত ছিল। দুইজন প্রার্থীর
একজন, উপস্থিত ভোটারদের ৫৫\% ভোট পাওয়ায় দেখা গেল যে অপর প্রার্থী অপেক্ষা তিনি
৭৫,০০০ ভোট বেশী পেয়েছেন । ভোটারদের মোট সংখ্যা কত?

\textbf{প্রশ্নঃ ৩৫০ টাকা দরে ৩ কেজি মিষ্টি কিনে ৪ টাকা হারে ভ্যাট দিলে মোট
কত ভ্যাট দিতে হবে?}

ক. ১৪ টাকা খ. ৪২ টাকা গ. ১২ টাকা ঘ. ১০৫ টাকা উত্তর: খ

\textbf{ব্যাখ্যা:} মুখে মুখে এভাবে ১০০ টাকায় ৪ টাকা ভ্যাট হলে ৩০০ টাকায় ভ্যাট
১২ টাকা এবং ৫০ টাকায় ২ টাকা ।

তাহলে ১ কেজি তে মোট ভ্যাট ১২+২= ১৪ টাকা এবং ৩ কেজিতে মোট ভ্যাট ১৪×৩ =
৪২টাকা)

\textbf{প্রশ্নঃ যদি তেলের মূল্য ২৫\% বৃদ্ধি পায় তবে তেলের ব্যবহার শতকরা কত
কমালে তেল বাবদ ব্যয় বৃদ্ধি পাবে না?}

ক.১৬\% খ.২০\% গ.২৫\% ঘ.২৪\% উত্তর: খ

\textbf{ব্যাখ্যা:} ২৫\% বেড়ে ১০০ থেকে হবে ১২৫ এখন খরচ বৃদ্ধি না করার জন্য ১০০
টাকার খরচ করতে হবে। অর্থাৎ নতুন খরচ কমাতে হবে ২৫ টাকা। এই ১২৫ টাকায় কমাতে
হবে ২৫টাকা বা ৫ ভাগের ১ ভাগ অর্থাৎ ১০০ তে ২০ বা ২০\%।

\textbf{প্রশ্নঃ ২ এর কত শতাংশ ৮ হবে?}

ক. ২০০ খ. ৪০০ গ. ৩৪৫ ঘ. ৩০০ উত্তর: খ

\textbf{ব্যাখ্যা:} ৮ হলো ২ এর থেকে ৪ গুণ বড় তাই শতকরা হার হবে ৪০০\% ।

\textbf{প্রশ্নঃ ১২ এর কত শতাংশ ১৮ হবে?}

ক. ১১০ খ. ১২৫ গ. ১৫০ ঘ. ১৬০ উত্তর:খ

\textbf{ব্যাখ্যা:} ১২ তে ১৮ হলে, ১ এ \(\frac{\text{১৮}}{\text{১২}}\) বা দেড়গুণ তাহলে ,
১০০ তে =১৫০।

\textbf{প্রশ্নঃ কলার দাম ২০\% কমে যাওয়ায় ১২ টাকায় পূর্ব অপেক্ষা ২টি কলা বেশি
পাওয়া গেলে বর্তমানে একটি কলার দাম কত টাকা?}

ক) ১.৫০ খ) ৩.০০ গ) ২.৫০ ঘ) ৪.০০

\textbf{নোট:} প্রশ্নে ভুল আছে। উত্তর :১.২ হবে, অনেক বইয়ে উত্তর মেলানোর জন্য ভুল
যুক্তি দেয়া হয়েছে।)

\textbf{ব্যাখ্যা:} ১২ টাকা কলার দাম ২০\% কমে যাওয়ায় টাকা সেভ হচ্ছে ১২ এর
২০\% বা ১২ × \(\frac{\text{২০}}{\text{১০০}}\) = ২.৪টাকা।

তাহলে এই সেভ হওয়া টাকা দিয়েই অতিরিক্ত ২টি কলা পাওয়া যায়, সুতরাং দুটি কলার
দাম হল ২.৪, তাহলে একটি কলার বর্তমান দাম হবে ২.৪÷২ = ১.২ টাকা।

এই অধ্যায় থেকে ছোট ছোট প্রশ্ন আসতে পারে ।

তাই নিচের প্রশ্নগুলো দেখে নিন:

ক) ৪ কেজি এর ৭৫\% কত কেজি? = ৩

খ) ৬০টি কলমের ২০\% কতটি? = ১২

গ)০.২ এর ২০\%= কত?(একই ভাবে)= .০৪

ঘ) ১৫ টাকার শতকরা ৭ অংশ কত? = ১.০৫

ঙ) ১০০ টাকার \(\frac{\text{১}}{\text{২}}\) \% = কত? = .৫০

\textbf{এবার বড় কথার কিছু প্রশ্ন দেখুন:}

\textbf{১) একটি কলম ৫০ টাকায় ক্রয় করে ৫৬ টাকায় বিক্রয় করা হলো। এতে শতকরা
কত লাভ হলো?}

\textbf{সমাধান:} ৫০টাকায় ৬ টাকা লাভ হলে ১০০ টাকায় ১২ টাকা বা ১২\% ।

\textbf{২) ৬০জন ছাত্রের মধ্যে ৪২ জন ফেল করলে পাশের হার কত?}

\textbf{সমাধান:} ৬০ জনে ফেল ৬০-৪২ = ১৮

তাহলে শতকরা ফেল হবে \(\frac{\text{১৮}}{\text{৬০}}\) × ১০০ = ৩০\%

\textbf{৩) এক ব্যক্তির বেতন ২০\% হারে বৃদ্ধি পেয়ে বর্তমানে ৬০০০ টাকা হয়েছে।
বৃদ্ধি পাওয়ার পূর্বে তার বেতন কত ছিল?}

\textbf{ব্যাখ্যা:} (এখানে ১২০\% এর মান ৬০০০) ( ২০\% বাড়ায় ১২০\% হয়েছে
তাহলে ১২০\% = ৬০০০ টাকা সুতরাং ১০০\% = ৫০০০টাকা )

\textbf{৪) ১জন পাইকারী ব্যবসায়ী ও খুচরা বিক্রেতা উভয়ে ২০\% লাভে পণ্য বিক্রি
করে। একটি শার্ট যার প্রাথমিক মূল্য ছিল ৪০০ টাকা, একজন ক্রেতা তা কত টাকায়
কিনতে পারবে?}

\textbf{\hl{উত্তর:৫৭৬ টাকা}}

খুব সহজে:

(৪০০ এর ১২০\% এর ১২০\% = ৫৭৬, কারণ ২০\% লাভ হওয়াতে ১০০ এর সাথে ২০\% যোগ
হয়েছে এবং দুবার লাভ করায় দুবার ১২০\% লিখতে হয়)

\textbf{৫) গমের মূল্য ১৫\% কমে যাওয়ায় ৬০০০ টাকায় পূর্বাপেক্ষা ১ কুইন্টাল গম
বেশি পাওয়া যায়। ১ কেজি গমের বর্তমান মূল্য কত?}

\textbf{\hl{উত্তর: ৯ টাকা {[}৬০০০ এর ১৫\% = ৬০০০ এর}} \(\frac{\text{১৫}}{\text{১০০}}\)
\textbf{\hl{= ৯০০ টাকা {]}}}

১ কুইন্টাল =১০০ কেজি,

১০০ কেজি গমের বর্তমান মূল্য = ৯০০ টাকা

১ কেজি গমের বর্তমান মূল্য = ৯০০÷১০০ = ৯ টাকা

\textbf{হ্রাস-বৃদ্ধির শতকরা:}

\textbf{এগুলো পরীক্ষায় অনেক আসে। দ্রুত কিভাবে করবেন?}

\textbf{৬) একজন ব্যাবসায়ী তার পণ্যের দাম ২০\% বাড়িয়ে দিলেন, এতে তার বিক্রি
কমে যাওয়ায় তিনি পুনরায় ২০\% দাম কমিয়ে দিলেন। এতে তার প্রথম মূল্যের তুলনায় দাম
কতটুকু কমলো বা বাড়লো?}

ব্যাখ্যা সহ সমাধান:

২০\% দাম বাড়ালে ১০০ থেকে ১২০ হয়, দ্বিতীয়বার ২০\% কমালে ১২০ এর ২০\% অর্থাৎ
২৪ কমে যায়। তাহলে থাকে ১২০-২৪ = ৯৬। মোটের উপর কমে ১০০-৯৬ = ৪\%

\textbf{৭) একটি আয়তাকার ক্ষেত্রের দৈর্ঘ্য ২০\% বৃদ্ধি ও প্রস্থ ১০\% হ্রাস করা
হলে, ক্ষেত্রফলের শতকরা কত পরিবর্তন হবে?}

\textbf{ব্যাখ্যা:} প্রথমে ১০০ থেকে ১২০ তার পর ১২০ এর ১০\% অর্থাৎ ১২ কমে ১০৮
হলে ৮\% বৃদ্ধি পেয়েছে।

\textbf{৮) একটি আয়তক্ষেত্রের দৈর্ঘ্য ১০\% বাড়লো এবং প্রস্থ ১০\% কমলো ক্ষেত্রফলের
কি পরিবর্তন হলো?}

\textbf{ব্যাখ্যা:} (১০০+১০ = ১১০-১১ = ৯৯ তাই উত্তর ১)

৯) একটি বর্গক্ষেত্রের ক্ষেত্রফল ১০০ বর্গমিটার। এই ক্ষেত্রের প্রতিটি বাহু ১০\% বৃদ্ধি
করা হলে ক্ষেত্রফল শতকরা কত ভাগ বৃদ্ধি পাবে? ২১\%

\textbf{ব্যাখ্যা:} দুবার ১০\% করে বাড়লে মোটের উপর ১০+১১ = ২১\% বাড়ে। উত্তর:
২১\%)

\textbf{১০) এক ব্যাবসায়ী একটি পণ্যের মূল্য ২৫\% বাড়ালে, অতঃপর বর্ধিত মূল্য থেকে
২৫\% কমালে, সর্বশেষ মূল্য সর্বপ্রথম মূল্যের তুলনায় কত বাড়লো বা কমলো?}

\textbf{ব্যাখ্যা:} {[}প্রথমে ২৫ বেড়ে ১২৫ হওয়া তারপর ১২৫ এর ২৫\% কমে গেল
মোটের উপর ৬.২৫\% কমে যাবে) {]}

\textbf{১১) একটি গণিতের বই কিনতে ১৫\% কমিশন দেয়। বইটির প্রকৃত দাম (কভারে
লিখিত দাম) ১২০ টাকা । বইটি কিনতে কত টাকা লাগবে?}

\textbf{ব্যাখ্যা:} কভারের দাম থেকে ১৫\% কমিশন দিলে নিবে ৮৫\% তাহলে এক
লাইনে উত্তর: ১২০ এর ৮৫\% বা ১০২ টাকা।

\subsubsection{পণ্যদ্রব্যের মূল্য বৃদ্ধি/হ্রাস এবং ব্যবহার বৃদ্ধি/হ্রাস সংক্রান্ত
সমস্যা}\label{ux9aaux9a3ux9afux9a6ux9b0ux9acux9afux9b0-ux9aeux9b2ux9af-ux9acux9a6ux9a7ux9b9ux9b0ux9b8-ux98fux9ac-ux9acux9afux9acux9b9ux9b0-ux9acux9a6ux9a7ux9b9ux9b0ux9b8-ux9b8ux995ux9b0ux9a8ux9a4-ux9b8ux9aeux9b8ux9af}

\paragraph{সূত্র-১: পণ্যদ্রব্যের মূল্য বৃদ্ধি পাওয়ায় ব্যবহার কমানোর ক্ষেত্রে- ব্যবহার
হ্রাসের হার = (১০০ × মূল্য বৃদ্ধির হার/১০০ + মূল্য বৃদ্ধির
হার)}\label{ux9b8ux9a4ux9b0-ux9e7-ux9aaux9a3ux9afux9a6ux9b0ux9acux9afux9b0-ux9aeux9b2ux9af-ux9acux9a6ux9a7-ux9aaux993ux9dfux9df-ux9acux9afux9acux9b9ux9b0-ux995ux9aeux9a8ux9b0-ux995ux9b7ux9a4ux9b0--ux9acux9afux9acux9b9ux9b0-ux9b9ux9b0ux9b8ux9b0-ux9b9ux9b0-ux9e7ux9e6ux9e6-ux9aeux9b2ux9af-ux9acux9a6ux9a7ux9b0-ux9b9ux9b0ux9e7ux9e6ux9e6-ux9aeux9b2ux9af-ux9acux9a6ux9a7ux9b0-ux9b9ux9b0}

\textbf{প্রশ্ন নম্বর-} ১. চিনির মূল্য ২৫\% বৃদ্ধি পাওয়াতে একটি পরিবার চিনি
খাওয়া এমন ভাবে কমালো যে চিনি বাবদ ব্যয় বৃদ্ধি পেলো না। ঐ পরিবার চিনি খাওয়া
বাবদ শতকরা কত কমালো? {[}এনএসআই এর সহকারী পরিচালক-২০১৫{]}

\begin{itemize}
\item
  (ক) ২২\%
\item
  (খ) ২৫\%
\item
  (গ) ২০\%
\item
  (ঘ) ৩০\%
\end{itemize}

\textbf{\hl{উত্তরঃ (গ) ২০\%}}

\textbf{সমাধান:}

শর্ট টেকনিক :

ব্যবহার হ্রাসের হার

= \(\frac{\text{১০০}\  \times \ \text{মূল্য}\ \text{বৃদ্ধির}\ \text{হার}}{\text{১০০}\  + \ \text{মূল্য}\ \text{বৃদ্ধির}\ \text{হার}}\)

= \(\frac{\text{১০০}\  \times \ \text{২৫}}{\text{১০০}\  + \ \text{২৫}}\)

= \(\frac{\text{২৫০০}}{\text{১২৫}}\)

= ২০ \%

লিখিত সমাধান:

২৫\% বৃদ্ধিতে চিনির পূর্বমূল্য ১০০ টাকা হলে বর্তমান মূল্য ১২৫ টাকা

১২৫ টাকায় কমাতে হবে = ২৫ টাকা

১ টাকায় কমাতে হবে = \(\frac{\text{২৫}}{\text{১২৫}}\) টাকা

১০০ টাকায় কমাতে হবে = \(\frac{\text{২৫} \times \text{১০০}}{\text{১২৫}}\) = ২০\%

\textbf{প্রশ্ন নম্বর-}২. চিনির মূল্য শতকরা ১০ টাকা বৃদ্ধি পেলে চিনির ব্যবহার
শতকরা কত কমালে পরিবারের চিনির খরচের কোন পরিবর্তন হবে না?{[}কারিগরী শিক্ষা
অধিদপ্তরের অধীনে চীফ ইন্সট্রাক্টর-২০০৩{]}

\begin{itemize}
\item
  (ক) (১০০/১১)
\item
  (খ) (৯৫/১১)
\item
  (গ) (১০২/১১)
\item
  (ঘ) (৯৩/১১)
\end{itemize}

\textbf{\hl{উত্তরঃ (ক)}} \(\frac{\text{১০০}}{\text{১১}}\)

\textbf{সমাধান:}

শর্ট টেকনিক :

ব্যবহার হ্রাসের হার

= \(\frac{\text{১০০}\  \times \ \text{মূল্য}\ \text{বৃদ্ধির}\ \text{হার}}{\text{১০০}\  + \ \text{মূল্য}\ \text{বৃদ্ধির}\ \text{হার}}\)

= \(\frac{\text{১০০}\  \times \ \text{১০}}{\text{১০০}\  + \ \text{১০}}\)

= \(\frac{\text{১০০০}}{\text{১১০}}\)

= \(\frac{\text{১০০}}{\text{১১}}\) \%

লিখিত সমাধান:

১০\% বৃদ্ধিতে চিনির পূর্বমূল্য ১০০ টাকা হলে বর্তমান মূল্য ১১০ টাকা

১১০ টাকায় কমাতে হবে = ১০ টাকা

১ টাকায় কমাতে হবে = \(\frac{\text{১০}}{\text{১১০}}\) টাকা

১০০ টাকায় কমাতে হবে = \(\frac{\text{১০}\  \times \ \text{১০০}}{\text{১১০}}\)

= \(\frac{\text{১০০}}{\text{১১}}\) টাকা

\paragraph{সূত্র-২: পণ্যদ্রব্যের মূল্য হ্রাস পাওয়ায় ব্যবহার বাড়ানোর ক্ষেত্রে-
ব্যবহার বৃদ্ধির হার = (১০০ × মূল্য হ্রাসের হার/১০০ − মূল্য হ্রাসের
হার)}\label{ux9b8ux9a4ux9b0-ux9e8-ux9aaux9a3ux9afux9a6ux9b0ux9acux9afux9b0-ux9aeux9b2ux9af-ux9b9ux9b0ux9b8-ux9aaux993ux9dfux9df-ux9acux9afux9acux9b9ux9b0-ux9acux9dcux9a8ux9b0-ux995ux9b7ux9a4ux9b0--ux9acux9afux9acux9b9ux9b0-ux9acux9a6ux9a7ux9b0-ux9b9ux9b0-ux9e7ux9e6ux9e6-ux9aeux9b2ux9af-ux9b9ux9b0ux9b8ux9b0-ux9b9ux9b0ux9e7ux9e6ux9e6-ux9aeux9b2ux9af-ux9b9ux9b0ux9b8ux9b0-ux9b9ux9b0}

\textbf{প্রশ্ন নম্বর-}১. কাপড়ের মূল্য ২০\% কমে গেলে কাপড়ের ব্যবহার কি পরিমাণ
বাড়ালে কাপড়ের জন্য খরচের কোন পরিবর্তন হবে না?{[}পাসপোর্ট এন্ড ইমিগ্রেশন
অধিদপ্তরের সহকারী পরিচালক-৯৪{]}

\begin{itemize}
\item
  (ক) ১৫\%
\item
  (খ) ২০\%
\item
  (গ) ২৫\%
\item
  (ঘ) ৩০\%
\end{itemize}

\textbf{\hl{উত্তরঃ (গ) ২৫\%}}

\textbf{সমাধান:}

শর্ট টেকনিক:

ব্যবহার বৃদ্ধির হার

= \(\frac{\text{১০০}\  \times \ \text{মূল্য}\ \text{হ্রাসের}\ \text{হার}}{\text{১০০}\  - \ \text{মূল্য}\ \text{হ্রাসের}\ \text{হার}}\)

= \(\frac{\text{১০০}\  \times \ \text{২০}}{\text{১০০}\  - \ \text{২০}}\)

= \(\frac{\text{২০০০}}{\text{৮০}}\)

= ২৫\%

লিখিত সমাধান:

২০\% কমে কাপড়ের পূর্বমূল্য ১০০ টাকা হলে বর্তমান মূল্য ৮০ টাকা।

বর্তমান মূল্য ৮০ টাকা হলে পূর্বমূল্য = ১০০ টাকা

বর্তমান মূল্য ১ টাকা হলে পূর্বমূল্য = \(\frac{\text{১০০}}{\text{৮০}}\) টাকা

বর্তমান মূল্য ১০০ টাকা হলে পূর্বমূল্য = \(\frac{\text{১০০}\  \times \ \text{১০০}}{\text{৮০}}\)

= ১২৫ টাকা

কাপড়ের ব্যবহার শতকরা বাড়াতে হবে = ১২৫−১০০ = ২৫\%

\paragraph{সূত্র-৩: পণ্যদ্রব্যের মূল্য বা ব্যবহার হ্রাস-বৃদ্ধির ক্ষেত্রে- বৃদ্ধি/হ্রাসের
হার = বৃদ্ধির হার+হ্রাসের হার+ (বৃদ্ধির হার × হ্রাসের
হার/১০০)}\label{ux9b8ux9a4ux9b0-ux9e9-ux9aaux9a3ux9afux9a6ux9b0ux9acux9afux9b0-ux9aeux9b2ux9af-ux9ac-ux9acux9afux9acux9b9ux9b0-ux9b9ux9b0ux9b8-ux9acux9a6ux9a7ux9b0-ux995ux9b7ux9a4ux9b0--ux9acux9a6ux9a7ux9b9ux9b0ux9b8ux9b0-ux9b9ux9b0-ux9acux9a6ux9a7ux9b0-ux9b9ux9b0ux9b9ux9b0ux9b8ux9b0-ux9b9ux9b0-ux9acux9a6ux9a7ux9b0-ux9b9ux9b0-ux9b9ux9b0ux9b8ux9b0-ux9b9ux9b0ux9e7ux9e6ux9e6}

\textbf{প্রশ্ন নম্বর-}১. চিনির মূল্য ২০\% কমলো, কিন্তু চিনির ব্যবহার ২০\% বৃদ্ধি
পেল। এতে চিনি বাবদ ব্যয় শতকরা কত বাড়লো বা কমলো?{[}বাংলাদেশ রেলওয়ে গাড-
২০১৬/১০ম বেসরকারি প্রভাষক নিবন্ধন ও প্রত্যয়ন পরীক্ষা-২০১৪{]}

\begin{itemize}
\item
  (ক) ১\% বাড়লো
\item
  (খ) ২\% কমলো
\item
  (গ) ৩\% বাড়লো
\item
  (ঘ) ৪\% কমলো
\end{itemize}

\textbf{\hl{উত্তরঃ (ঘ) ৪\% কমলো}}

\textbf{সমাধান:}

শর্ট টেকনিক:

বৃদ্ধি/হ্রাসের হার

= বৃদ্ধির হার + হ্রাসের হার +
\(\frac{\text{বৃদ্ধির}\ \text{হার}\  \times \ \text{হ্রাসের}\ \text{হার}}{\text{১০০}}\)

= ২০\%+(-২০\%)+ \(\frac{\text{২০} \times (\  - \ \text{২০}\ )}{\text{১০০}}\)

= ০ + \(\frac{- \ \text{৪০০}}{\text{১০০}}\)

= - ৪ সুতরাং ৪ \% কমলো।

বিঃ দ্রঃ ধনাত্নক মান আসলে বাড়বে এবং ঋণাত্নক মান আসলে কমবে।

লিখিত সমাধান:

ধরি, ক্রয়মূল্য ১০০ টাকা

২০\% কমে চিনির ক্রয়মূল্য = ১০০-২০ = ৮০ টাকা

আবার, ২০\% বৃদ্ধিতে ব্যয় =(৮০+৮০× \(\frac{\text{২০}}{\text{১০০}}\) )

= ৮০+১৬ = ৯৬ টাকা

ব্যয় কমে = ১০০-৯৬ = ৪\%

\textbf{প্রশ্ন নম্বর-}২. এক ব্যবসায়ী একটি পণ্যের মূল্য ২৫\% বাড়াল, অতঃপর বর্ধিত
মূল্য থেকে ২৫\% কমালো। র্সবশেষ মূল্য সর্বপ্রথম মূল্যের তুলনায় - {[} বাংলাদেশ ব্যাংক
সহকারী পরিচালক-২০০৬/ ২৭তম বিসিএস{]}

\begin{itemize}
\item
  (ক) ৪.৫\% কমানো হয়েছে
\item
  (খ) ৬.২৫\% কমানো হয়েছে
\item
  (গ) ৫\% বাড়ানো হয়েছে
\item
  (ঘ) ৬.২৫\% বাড়ানো হয়েছে
\end{itemize}

\textbf{\hl{উত্তরঃ (খ) ৬.২৫\% কমানো হয়েছে}}

\textbf{সমাধান:}

শর্ট টেকনিক:

সর্বশেষ মূল্য সর্বপ্রথম মূল্যের তুলনায়

= বৃদ্ধির হার+হ্রাসের হার+
\(\frac{\text{বৃদ্ধির}\ \text{হার}\  \times \ \text{হ্রাসের}\ \text{হার}}{\text{১০০}}\)

= ২৫\%+(−২৫\%)+ \(\frac{\text{২৫} \times (\  - \ \text{২৫}\ )}{\text{১০০}}\)

= ০ + \(\frac{- \ \text{৬২৫}}{\text{১০০}}\)

= - ৬.২৫ সুতরাং ৬.২৫\% কমানো হয়েছে

বিঃ দ্রঃ ধনাত্নক মান আসলে বাড়বে এবং ঋণাত্নক মান আসলে কমবে।

লিখিত সমাধান:

ধরি, পণ্যের ক্রয়মূল্য ১০০ টাকা

২৫\% কমে পণ্যের ক্রয়মূল্য = ১০০-২৫ = ৭৫ টাকা

আবার, ২৫\% বৃদ্ধিতে ব্যয় =(৭৫+৭৫×\(\frac{\text{২৫}}{\text{১০০}}\) )

= ৭৫+১৮.৭৫ = ৯৩.৭৫ টাকা

সর্বশেষ মূল্য সর্বপ্রথম মূল্যের তুলনায় = ১০০-৯৩.৭৫ = ৬.২৫\% কম

\paragraph{সূত্র-৪: দ্রব্যমূল্যের শতকরা হার বৃদ্ধি পাওয়ায় - দ্রব্যের বর্তমান মূল্য =
(বৃদ্ধিপ্রাপ্ত মূল্য হার × মোট মূল্য/১০০ × যে পরিমাণ পণ্য কম পাওয়া
গিয়েছে)}\label{ux9b8ux9a4ux9b0-ux9ea-ux9a6ux9b0ux9acux9afux9aeux9b2ux9afux9b0-ux9b6ux9a4ux995ux9b0-ux9b9ux9b0-ux9acux9a6ux9a7-ux9aaux993ux9dfux9df---ux9a6ux9b0ux9acux9afux9b0-ux9acux9b0ux9a4ux9aeux9a8-ux9aeux9b2ux9af-ux9acux9a6ux9a7ux9aaux9b0ux9aaux9a4-ux9aeux9b2ux9af-ux9b9ux9b0-ux9aeux99f-ux9aeux9b2ux9afux9e7ux9e6ux9e6-ux9af-ux9aaux9b0ux9aeux9a3-ux9aaux9a3ux9af-ux995ux9ae-ux9aaux993ux9df-ux997ux9dfux99b}

\textbf{প্রশ্ন নম্বর-}১. চিনির মূল্য ৬\% বেড়ে যাওয়ায় ৭৯৫ টাকায় পুর্বে যত কেজি
চিনি পাওয়া যেত, বর্তমানে তার চেয়ে ৩ কেজি চিনি কম পাওয়া যায়। চিনির বর্তমান
দর প্রতি কেজি কত?{[}সরকারী মাধ্যমিক সহকারী শিক্ষক পরীক্ষা-২০০০{]}

\begin{itemize}
\item
  (ক) ১৫ টাকা
\item
  (খ) ১৬ টাকা
\item
  (গ) ১৮ টাকা
\item
  (ঘ) ১৫.৯০ টাকা
\end{itemize}

\textbf{\hl{উত্তরঃ (ঘ) ১৫.৯০ টাকা}}

\textbf{সমাধান:}

শর্ট টেকনিক :

চিনির বর্তমান মূল্য =

\(\frac{\text{বৃদ্ধিপ্রাপ্ত}\ \text{মূল্য}\ \text{হার}\  \times \ \text{মোট}\ \text{মূল্য}}{\text{১০০}\  \times \ \text{যে}\ \text{পরিমাণ}\ \text{পণ্য}\ \text{কম}\ \text{পাওয়া}\ \text{গিয়েছে}}\)

= \(\frac{\text{৬}\  \times \ \text{৭৯৫}}{\text{১০০}\  \times \ \text{৩}}\)

= \(\frac{\text{৪৭৭০}}{\text{৩০০}}\)

= ১৫.৯০ টাকা

লিখিত সমাধান:

৬\% বেড়ে চিনির পূর্বমূল্য ১০০ টাকা হলে বর্তমান মূল্য = ১০০+৬ = ১০৬ টাকা

পূর্বমূল্য ১০০ টাকা হলে বর্তমান মূল্য = ১০৬ টাকা

পূর্বমূল্য ১ টাকা হলে বর্তমান মূল্য = \(\frac{\text{১০৬}}{\text{১০০}}\) টাকা

পূর্বমূল্য ৭৯৫ টাকা হলে বর্তমান মূল্য = \(\frac{\text{১০৬} \times \ \text{৭৯৫}}{\text{১০০}}\) =
৮৪২.৭ টাকা

৩ কেজি চিনির বর্তমান দর = ৮৪২.৭-৭৯৫ = ৪৭.৭ টাকা

চিনির বর্তমান দর প্রতি কেজি = ৪৭.৭÷৩ = ১৫.৯০ টাকা

\textbf{প্রশ্ন নম্বর-}২. চিনির মূল্য ৬\% বেড়ে যাওয়ায় ১০৬০ টাকায় পুর্বে যত কেজি
চিনি পাওয়া যেত, বর্তমানে তার চেয়ে ৩ কেজি চিনি কম পাওয়া যায়। চিনির বর্তমান
দর প্রতি কেজি কত?{[}সরকারী মাধ্যমিক সহকারী শিক্ষক পরীক্ষা-২০০০{]}

\begin{itemize}
\item
  (ক) ২০ টাকা
\item
  (খ) ২০.২০ টাকা
\item
  (গ) ২১.২০ টাকা
\item
  (ঘ) ২১ টাকা
\end{itemize}

\textbf{\hl{উত্তরঃ (গ) ২১.২০ টাকা}}

\textbf{সমাধান:}

শর্ট টেকনিক :

চিনির বর্তমান মূল্য

=
\(\frac{\text{বৃদ্ধিপ্রাপ্ত}\ \text{মূল্য}\ \text{হার}\  \times \ \text{মোট}\ \text{মূল্য}}{\text{১০০}\  \times \ \text{যে}\ \text{পরিমাণ}\ \text{পণ্য}\ \text{কম}\ \text{পাওয়া}\ \text{গিয়েছে}}\)

= \(\frac{\text{৬}\  \times \ \text{১০৬০}}{\text{১০০}\  \times \ \text{৩}}\)

= \(\frac{\text{৬৩৬০}}{\text{৩০০}}\)

= ২১.২০ টাকা

লিখিত সমাধান:

৬\% বেড়ে চিনির পূর্বমূল্য ১০০ টাকা হলে বর্তমান মূল্য = ১০০+৬ = ১০৬ টাকা

পূর্বমূল্য ১০০ টাকা হলে বর্তমান মূল্য = ১০৬ টাকা

পূর্বমূল্য ১ টাকা হলে বর্তমান মূল্য = \(\frac{\text{১০৬}}{\text{১০০}}\) টাকা

পূর্বমূল্য ১০৬০ টাকা হলে বর্তমান মূল্য = \(\frac{\text{১০৬} \times \ \text{১০৬০}}{\text{১০০}}\) =
১১২৩.৬ টাকা

৩ কেজি চিনির বর্তমান দর = ১১২৩.৬ - ১০৬০ = ৬৩.৬ টাকা

চিনির বর্তমান দর প্রতি কেজি = ৬৩.৬÷৩ = ২১.২০ টাকা

\paragraph{সূত্র-৫: দ্রব্যমূল্যের শতকরা হার হ্রাস পাওয়ায় - দ্রব্যের বর্তমান মূল্য =
(হ্রাসকৃত মূল্য হার × মোট মূল্য/১০০ × যে পরিমাণ পণ্য বেশি পাওয়া
গিয়েছে)}\label{ux9b8ux9a4ux9b0-ux9eb-ux9a6ux9b0ux9acux9afux9aeux9b2ux9afux9b0-ux9b6ux9a4ux995ux9b0-ux9b9ux9b0-ux9b9ux9b0ux9b8-ux9aaux993ux9dfux9df---ux9a6ux9b0ux9acux9afux9b0-ux9acux9b0ux9a4ux9aeux9a8-ux9aeux9b2ux9af-ux9b9ux9b0ux9b8ux995ux9a4-ux9aeux9b2ux9af-ux9b9ux9b0-ux9aeux99f-ux9aeux9b2ux9afux9e7ux9e6ux9e6-ux9af-ux9aaux9b0ux9aeux9a3-ux9aaux9a3ux9af-ux9acux9b6-ux9aaux993ux9df-ux997ux9dfux99b}

\textbf{প্রশ্ন নম্বর-}১. চালের মূল্য ১২\% কমে যাওয়ায় ৬০০০ টাকায় পূর্বাপেক্ষা ১
কুইন্টাল চাল বেশি পাওয়া যায়। এক কুইন্টাল চালের বর্তমান মূল্য কত?{[}১১তম
বিসিএস{]}

\begin{itemize}
\item
  (ক) ৭৫০ টাকা
\item
  (খ) ৭০০ টাকা
\item
  (গ) ৭২০ টাকা
\item
  (ঘ) ৭৫ টাকা
\end{itemize}

\textbf{\hl{উত্তরঃ (গ) ৭২০ টাকা}}

\textbf{সমাধান:}

শর্ট টেকনিক :

এক কুইন্টাল চালের বর্তমান মূল্য

=
\(\frac{\text{হ্রাসকৃত}\ \text{মূল্য}\ \text{হার}\  \times \ \text{মোট}\ \text{মূল্য}}{\text{১০০}\  \times \ \text{যে}\ \text{পরিমাণ}\ \text{পণ্য}\ \text{বেশি}\ \text{পাওয়া}\ \text{গিয়েছে}}\)

= \(\frac{\text{১২}\  \times \ \text{৬০০০}}{\text{১০০}\  \times \ \text{১}}\)

= ৭২০ টাকা

লিখিত সমাধান:

১২\% কমে চালের পূর্বমূল্য ১০০ টাকা হলে বর্তমান মূল্য = ১০০-১২ = ৮৮ টাকা

পূর্বমূল্য ১০০ টাকা হলে বর্তমান মূল্য = ৮৮ টাকা

পূর্বমূল্য ১ টাকা হলে বর্তমান মূল্য = \(\frac{\text{৮৮}}{\text{১০০}}\) টাকা

পূর্বমূল্য ৬০০০ টাকা হলে বর্তমান মূল্য = \(\frac{\text{৮৮} \times \text{৬০০০}}{\text{১০০}}\)

= ৫২৮০ টাকা

এক কুইন্টাল চালের বর্তমান মূল্য = ৬০০০-৫২৮০ = ৭২০ টাকা

\textbf{প্রশ্ন নম্বর-}২. গমের মূল্য ১৫\% কমে যাওয়ায় ৬০০০ টাকায় পূর্বাপেক্ষা ১
কুইন্টাল গম বেশি পাওয়া যায়। এক কুইন্টাল গমের বর্তমান মূল্য কত?{[}১১তম বিসিএস{]}

\begin{itemize}
\item
  (ক) ৮৫০ টাকা
\item
  (খ) ৯০০ টাকা
\item
  (গ) ৭২০ টাকা
\item
  (ঘ) ৭৫ টাকা
\end{itemize}

\textbf{\hl{উত্তরঃ (খ) ৯০০ টাকা}}

\textbf{সমাধান:}

শর্ট টেকনিক :

এক কুইন্টাল গমের বর্তমান মূল্য

=
\(\frac{\text{হ্রাসকৃত}\ \text{মূল্য}\ \text{হার}\  \times \ \text{মোট}\ \text{মূল্য}}{\text{১০০}\  \times \ \text{যে}\ \text{পরিমাণ}\ \text{পণ্য}\ \text{বেশি}\ \text{পাওয়া}\ \text{গিয়েছে}}\)

= \(\frac{\text{১৫}\  \times \ \text{৬০০০}}{\text{১০০}\  \times \ \text{১}}\)

= ৯০০ টাকা

লিখিত সমাধান:

১৫\% কমে চালের পূর্বমূল্য ১০০ টাকা হলে বর্তমান মূল্য = ১০০-১৫ = ৮৫ টাকা

পূর্বমূল্য ১০০ টাকা হলে বর্তমান মূল্য = ৮৫ টাকা

পূর্বমূল্য ১ টাকা হলে বর্তমান মূল্য = \(\frac{\text{৮৫}}{\text{১০০}}\) টাকা

পূর্বমূল্য ৬০০০ টাকা হলে বর্তমান মূল্য = \(\frac{\text{৮৫} \times \text{৬০০০}}{\text{১০০}}\)

= ৫১০০ টাকা

এক কুইন্টাল চালের বর্তমান মূল্য = ৬০০০−৫১০০ = ৯০০ টাকা

\subsubsection{পরীক্ষার্থী ও ‪পরীক্ষা বিষয়ক সংক্রান্ত সমস্যা
‬‬‬}\label{ux9aaux9b0ux995ux9b7ux9b0ux9a5-ux993-ux9aaux9b0ux995ux9b7-ux9acux9b7ux9dfux995-ux9b8ux995ux9b0ux9a8ux9a4-ux9b8ux9aeux9b8ux9af}

\paragraph{সূত্র-১ : উভয় বিষয়ে ফেলের হার উল্লেখ থাকলে উভয় বিষয়ে পাশের হার
নির্ণয়ের ক্ষেত্রে-‪‬‬‬‬‬‬‬‬‬‬‬‬‬ শর্ট টেকনিক : পাশের হার = ১০০ − ( ১ম বিষয়ে ফেলের হার + ২য়
বিষয়ে ফেলের হার − উভয় বিষয়ে ফেলের হার
)}\label{ux9b8ux9a4ux9b0-ux9e7-ux989ux9adux9df-ux9acux9b7ux9df-ux9abux9b2ux9b0-ux9b9ux9b0-ux989ux9b2ux9b2ux996-ux9a5ux995ux9b2-ux989ux9adux9df-ux9acux9b7ux9df-ux9aaux9b6ux9b0-ux9b9ux9b0-ux9a8ux9b0ux9a3ux9dfux9b0-ux995ux9b7ux9a4ux9b0--ux9b6ux9b0ux99f-ux99fux995ux9a8ux995-ux9aaux9b6ux9b0-ux9b9ux9b0-ux9e7ux9e6ux9e6-ux9e7ux9ae-ux9acux9b7ux9df-ux9abux9b2ux9b0-ux9b9ux9b0-ux9e8ux9df-ux9acux9b7ux9df-ux9abux9b2ux9b0-ux9b9ux9b0-ux989ux9adux9df-ux9acux9b7ux9df-ux9abux9b2ux9b0-ux9b9ux9b0}

\textbf{প্রশ্ন নম্বর-} ১. কোন পরীক্ষায় ২০\% পরীক্ষার্থী গণিতে ৩০\% পরীক্ষার্থী
ইংরেজিতে ফেল করল উভয় বিষয়ে ১৩ \% পরীক্ষার্থী ফেল করলে শতকরা কতজন পরীক্ষার্থী
পাশ করলো? {[}প্রাথমিক সহকারী শিক্ষক (খুলনা) পরীক্ষা-২০০৬{]}‬‬‬‬‬‬‬‬‬‬‬‬

\begin{itemize}
\item
  (ক)৫০\%
\item
  (খ) ৫৭\%
\item
  (গ) ৬৩\%
\item
  (ঘ) ৬৭\%
\end{itemize}

\textbf{\hl{উত্তরঃ (গ) ৬৩\%}}

\textbf{সমাধান:}

শর্ট টেকনিক:

এখানে, ১ম বিষয় গণিতে ফেলের হার = ২০\%

২য় বিষয় ইংরেজিতে ফেলের হার = ৩০\%

এবং উভয় বিষয়ে ফেলের হার ১৩\% দেওয়া আছে।

∴পাশের হার = ১০০-(১ম বিষয়ে ফেলের হার+২য় বিষয়ে ফেলের হার- উভয় বিষয়ে ফেলের
হার)

= ১০০-(২০+৩০-১৩) = ১০০ - ( ৫০ - ১৩)

= ১০০ - ৩৭ = ৬৩\%

লিখিত সমাধান:

শুধু ইংরেজিতে ফেল করে = (২০ - ১৩)\% = ৭\%

শুধু গণিতে ফেল করে = (৩০ - ১৩)\% = ১৭\%

উভয় বিষয়ে ফেল করে = (৭ + ১৭ + ১৩)\% = ৩৭\%

উভয় বিষয়ে পাশ করে = (১০০ - ৩৭)\% = ৬৩\%

\textbf{প্রশ্ন নম্বর-} ২. কোন পরীক্ষায় ৫২\% ছাত্র বিজ্ঞানে এবং ৪০\% ছাত্র অংকে
ফেল করে। যদি উভয় বিষয়ে ২৭ \% পরীক্ষার্থী ফেল করে তবে শতকরা কতজন ছাত্র পাশ‬
করে? {[}প্রাথমিক সহকারী শিক্ষক নিয়োগ (ইছামতি) পরীক্ষা-২০১০){]}‬‬‬‬‬‬‬‬‬‬‬‬

\begin{itemize}
\item
  (ক) ৩৫
\item
  (খ) ৪৮
\item
  (গ) ৬০
\item
  (ঘ) ৬৫
\end{itemize}

\textbf{\hl{উত্তরঃ (ক) ৩৫}}

\textbf{সমাধান:}

শর্ট টেকনিক:

এখানে, ১ম বিষয় বিজ্ঞানে ফেলের হার = ৫২\%

২য় বিষয় অংকে ফেলের হার = ৪০\%

এবং উভয় বিষয়ে ফেলের হার ২৭\% দেওয়া আছে।

∴পাশের হার = ১০০-(৫২+৪০-২৭)‬‬

= ১০০ - ( ৯২ - ২৭) = ১০০ - ৬৫ = ৩৫\%

লিখিত সমাধান:

শুধু বিজ্ঞানে ফেল করে = (৫২ - ২৭)\% = ২৫\%

শুধু অংকে ফেল করে = (৪০ - ২৭)\% = ১৩\%

উভয় বিষয়ে ফেল করে = (২৫ + ১৩ + ২৭)\% = ৬৫\%

উভয় বিষয়ে পাশ করে = (১০০ - ৬৫)\% = ৩৫\%

\paragraph{সূত্র ২ : উভয় বিষয়ে পাশের হার উল্লেখ থাকলে উভয় বিষয়ে ফেলের হার
নির্ণয়ের ক্ষেত্রে-‪‬‬‬‬‬‬‬‬‬‬‬‬‬ শর্ট টেকনিক : ফেলের হার = ১০০ − ( ১ম বিষয়ে পাশের হার + ২য়
বিষয়ে পাশের হার − উভয় বিষয়ে পাশের হার) (১ম টির উল্টো
নিয়ম)}\label{ux9b8ux9a4ux9b0-ux9e8-ux989ux9adux9df-ux9acux9b7ux9df-ux9aaux9b6ux9b0-ux9b9ux9b0-ux989ux9b2ux9b2ux996-ux9a5ux995ux9b2-ux989ux9adux9df-ux9acux9b7ux9df-ux9abux9b2ux9b0-ux9b9ux9b0-ux9a8ux9b0ux9a3ux9dfux9b0-ux995ux9b7ux9a4ux9b0--ux9b6ux9b0ux99f-ux99fux995ux9a8ux995-ux9abux9b2ux9b0-ux9b9ux9b0-ux9e7ux9e6ux9e6-ux9e7ux9ae-ux9acux9b7ux9df-ux9aaux9b6ux9b0-ux9b9ux9b0-ux9e8ux9df-ux9acux9b7ux9df-ux9aaux9b6ux9b0-ux9b9ux9b0-ux989ux9adux9df-ux9acux9b7ux9df-ux9aaux9b6ux9b0-ux9b9ux9b0-ux9e7ux9ae-ux99fux9b0-ux989ux9b2ux99f-ux9a8ux9dfux9ae}

\textbf{প্রশ্ন নম্বর-} ১ . কোন পরীক্ষায় ৮০\% পরীক্ষার্থী গণিতে এবং ৭০\%
পরীক্ষার্থী বাংলায় পাশ করলো। উভয় বিষয়ে পাশ করলো ৬০\% । উভয় বিষয়ে শতকরা কত
জন ফেল করলো? {[} ২২তম বিসিএস., উপজেলা শিক্ষাঅফিসার-২০০৬,অর্থমন্ত্রণালয়
-২০১১, ৬ষ্ঠ সহকারী জজ নি:- ২০১১{]}

\begin{itemize}
\item
  (ক)১৫\%
\item
  (খ) ১০\%
\item
  (গ) ১২\%
\item
  (ঘ) ১১\%
\end{itemize}

\textbf{\hl{উত্তরঃ (খ) ১০\%}}

\textbf{সমাধান:}

শর্ট টেকনিক :

এখানে, ১ম বিষয় গণিতে পাশের হার = ৮০\%

২য় বিষয় বাংলায় পাশের হার = ৭০\%

এবং উভয় বিষয়ে পাশের হার ৬০\% দেওয়া আছে।

ফেলের হার = ১০০ - (১ম বিষয়ে পাশের হার+২য় বিষয়ে পাশের হার- উভয় বিষয়ে
পাশের হার)

= ১০০ - ( ৮০ + ৭০ - ৬০) = ১০০ - ( ১৫০ - ৬০)

= ১০০ - ৯০ = ১০ \%

লিখিত সমাধান:

শুধু গণিতে পাস করে =(৮০-৬০)\% = ২০\%

শুধু বাংলায় পাস করে = (৭০-৬০)\% = ১০\%

উভয় বিষয়ে পাস করেছে ৬০\%

দুই বিষয়ে মোট পাস করে = ২০+১০+৬০ = ৯০\%

উভয় বিষয়ে ফেল করেছে = ১০০-৯০ = ১০\%

\textbf{প্রশ্ন নম্বর-} ২. কোন পরীক্ষায় ৮৫\% পরীক্ষার্থী পদার্থবিদ্যায় , ৮০\%
পরীক্ষার্থী রসায়নবিদ্যায় এবং ৭৫\% পরীক্ষার্থী উভয় বিষয়ে পাশ করে ।কত জন
পরীক্ষার্থী উভয় বিষয়ে ফেল করে? {[} রেজি. প্রাথমিক স ঃ শি নিয়োগ (হাসনাহেনা)
পরীক্ষা-২০১১){]}

\begin{itemize}
\item
  (ক)৮ জন
\item
  (খ) ১০ জন
\item
  (গ) ১২ জন
\item
  (ঘ) ১৫ জন
\end{itemize}

\textbf{\hl{উত্তরঃ (খ) ১০ জন}}

\textbf{সমাধান:}

শর্ট টেকনিক :

এখানে, ১ম বিষয় পদার্থবিদ্যায় পাশের হার = ৮৫\%

২য় বিষয়র্ রসায়নবিদ্যায় পাশের হার = ৮০\%

এবং উভয় বিষয়ে পাশের হার ৭৫\% দেওয়া আছে।

ফেলের হার = ১০০ - (১ম বিষয়ে পাশের হার+২য় বিষয়ে পাশের হার- উভয় বিষয়ে
পাশের হার)

= ১০০ - ( ৮৫ + ৮০ - ৭৫) = ১০০ - ( ১৬৫ - ৭৫)

= ১০০ - ৯০ = ১০ \%

লিখিত সমাধান:

শুধু গণিতে পাস করে =(৮৫-৭৫)\% = ১০\%

শুধু বাংলায় পাস করে = (৮০-৭৫)\% = ৫\%

উভয় বিষয়ে পাস করেছে ৭৫\%

দুই বিষয়ে মোট পাস করে = ১০+৫+৭৫ = ৯০\%

উভয় বিষয়ে ফেল করেছে = ১০০-৯০ = ১০\%

\textbf{প্রশ্ন নম্বর-} ৩. কোনো শ্রেণীতে ১০০ জন পরীক্ষার্থী ছিল। বাষিক পরীক্ষায়
৯৪ জন বাংলায় এবং ৮০ জন গণিতে পাস করেছে। ৭৫ জন উভয় বিষয়ে পাস করলে কত জন
উভয় বিষয়ে ফেল করেছে? {[} আমদানি রপ্তানি অধিদপ্তরের নিবাহী অফিসার-২০০৭{]}

\begin{itemize}
\item
  (ক) ৫
\item
  (খ) ৩
\item
  (গ) ২
\item
  (ঘ) ১
\end{itemize}

\textbf{\hl{উত্তরঃ (ঘ) ১}}

\textbf{সমাধান:}

শর্ট টেকনিক :

এখানে, ১ম বিষয় বাংলায় পাস করে = ৯৪

২য় বিষয় গণিতে পাস করে = ৮০

এবং উভয় বিষয়ে পাশের হার ৭৫ দেওয়া আছে।

উভয় বিষয়ে ফেল করেছে = ১০০ - (১ম বিষয়ে পাশের হার+২য় বিষয়ে পাশের হার- উভয়
বিষয়ে পাশের হার)

= ১০০ - ( ৯৪ + ৮০ - ৭৫) = ১০০ - ( ১৭৪ - ৭৫)

= ১০০ - ৯৯ = ১ লিখিত সমাধান:

শুধু বাংলায় পাস করে =(৯৪-৭৫) = ১৯

শুধু গণিতে পাস করে = (৮০-৭৫) = ৫

উভয় বিষয়ে পাস করেছে ৭৫ জন

দুই বিষয়ে মোট পাস করে = ১৯+৫+৭৫ = ৯৯

উভয় বিষয়ে ফেল করেছে = ১০০-৯৯ = ১ জন

যদি প্রশ্নটি এমন হয়,

\textbf{প্রশ্ন নম্বর-} ৪. কোন পরীক্ষায় ২০০ জনের মধ্যে ৭০\% বাংলায় এবং ৬০\%
ইংরেজিতে পাশ করে করে এবং ৪০\% উভয় বিষয়ে পাশ করে। উভয় বিষয়ে ফেল করে কত জন
? (সঞ্চয় অধিদপ্তরের সহকারী পরিচালক-২০০৭)

\begin{itemize}
\item
  (ক) ১০ জন
\item
  (খ) ১৫ জন
\item
  (গ) ২৫ জন
\item
  (ঘ) ২০ জন
\end{itemize}

\textbf{\hl{উত্তরঃ (ঘ) ২০ জন}}

\textbf{সমাধান:}

শর্ট টেকনিক:

এখানে, ১ম বিষয় বাংলায় পাশের হার = ৭০\%,

২য় বিষয়র্ ইংরেজিতে পাশের হার = ৬০\%

এবং উভয় বিষয়ে পাশের হার ৪০\% দেওয়া আছে।

‪∴ফেলের হার = ১০০-(৭০+৬০-৪০)‬‬‬‬‬‬‬‬‬‬‬‬‬

= ১০০ - ( ১৩০- ৪০) = ১০০ - ৯০ = ১০\%

সুতরাং উভয় বিষয়ে ফেল = ২০০ এর ১০\%

= ২০ জন

লিখিত সমাধান:

শুধু বাংলায় পাস করে = ৭০-৪০ = ৩০\%

শুধু ইংরেজীতে পাস করে = ৬০-৪০ = ২০\%

দুই বিষয়ে মোট পাস করে =৩০+২০+৪০ = ৯০\%

উভয় বিষয়ে ফেল করে = ১০০-৯০ = ১০\%

২০০ জনের ১০\% = ২০ জন

উভয় বিষয়ে ফেল করে মোট = ২০ জন

\paragraph{সূত্র-৩ : উভয় বিষয়ে ফেল এবং পাশের উল্লেখ থাকলে মোট পরীক্ষার্থীর
সংখ্যা নির্ণয়ের ক্ষেত্রে- শর্ট টেকনিক : মোট পরীক্ষার্থী = (উভয় বিষয়ে পাসকৃত
ছাত্র/১ম বিষয়ে ফেল + ২য় বিষয়ে ফেল+উভয় বিষয়ে ফেল)
×১০০}\label{ux9b8ux9a4ux9b0-ux9e9-ux989ux9adux9df-ux9acux9b7ux9df-ux9abux9b2-ux98fux9ac-ux9aaux9b6ux9b0-ux989ux9b2ux9b2ux996-ux9a5ux995ux9b2-ux9aeux99f-ux9aaux9b0ux995ux9b7ux9b0ux9a5ux9b0-ux9b8ux996ux9af-ux9a8ux9b0ux9a3ux9dfux9b0-ux995ux9b7ux9a4ux9b0--ux9b6ux9b0ux99f-ux99fux995ux9a8ux995-ux9aeux99f-ux9aaux9b0ux995ux9b7ux9b0ux9a5-ux989ux9adux9df-ux9acux9b7ux9df-ux9aaux9b8ux995ux9a4-ux99bux9a4ux9b0ux9e7ux9ae-ux9acux9b7ux9df-ux9abux9b2-ux9e8ux9df-ux9acux9b7ux9df-ux9abux9b2ux989ux9adux9df-ux9acux9b7ux9df-ux9abux9b2-ux9e7ux9e6ux9e6}

\textbf{প্রশ্ন নম্বর-} ১. কোন স্কুলে ৭০\% শিক্ষার্থী ইংরেজিতে এবং ৮০\% শিক্ষার্থী
বাংলায় পাশ করেছে। কিন্তু ১০\% শিক্ষার্থী উভয় বিষয়ে ফেল করেছে। যদি উভয় বিষয়ে
৩৬০ জন পরীক্ষার্থী পাশ করে তবে ঐ স্কুলে কত জন শিক্ষার্থী পরীক্ষা দিয়েছে? (২৩তম
বিসিএস)

\begin{itemize}
\item
  (ক)৫০০ জন
\item
  (খ) ৪৫০ জন
\item
  (গ) ৭০০ জন
\item
  (ঘ) ৬০০ জন
\end{itemize}

\textbf{\hl{উত্তরঃ (ঘ) ৬০০ জন}}

\textbf{সমাধান:}

শর্ট টেকনিক :

এখানে, ১ম বিষয় ইংরেজিতে পাশের হার = ৭০\%

২য় বিষয় বাংলায় পাশের হার = ৮০\%

উভয় বিষয়ে ফেলের হার ১০\% এবং উভয় বিষয়ে পাসকৃত ছাত্র ৩৬০ জন দেওয়া আছে।

তাহলে, ১ম বিষয় ইংরেজিতে ফেলের হার = ১০০\% - ৭০\% = ৩০\%

২য় বিষয় বাংলায় ফেলের হার = ১০০\% - ৮০\% = ২০\%

মোট শিক্ষার্থী =
\(\frac{\text{উভয়}\ \text{বিষয়ে}\ \text{পাসকৃত}\ \text{ছাত্র}}{\text{১ম}\ \text{বিষয়ে}\ \text{ফেল}\  + \ \text{২য়}\ \text{বিষয়ে}\ \text{ফেল}\  + \ \text{উভয়}\ \text{বিষয়ে}\ \text{ফেল}}\)

×১০০ = \(\frac{\text{৩৬০}}{\text{৩০} + \text{২০} + \text{১০}}\)

×১০০ = \(\frac{\text{৩৬০}}{\text{৬০}}\)

×১০০ = ৬০০ জন

লিখিত সমাধানঃ

শুধু ইংরেজিতে ফেল করে = (১০০-৭০)-১০\% = ২০\%

শুধু বাংলায় ফেল করে = (১০০-৮০)-১০\% = ১০\%

উভয় বিষয়ে ফেল করেছে ১০\%

দুই বিষয়ে মোট ফেল করে = ২০+১০+১০ = ৪০\%

দুই বিষয়ে মোট পাস করে = ১০০-৪০ = ৬০\%

৬০ জন পাস করে = ১০০ জনের মধ্যে

১ জন পাস করে = \(\frac{\text{১০০}}{\text{৬০}}\) জনের মধ্যে

৩৬০ জন পাস করে = \(\frac{\text{১০০}\  \times \ \text{৩৬০}}{\text{৬০}}\) = ৬০০ জনের মধ্যে

ঐ স্কুলে পরীক্ষা দিয়েছে = ৬০০ জন শিক্ষার্থী

\textbf{প্রশ্ন নম্বর-} ২. কফিলাতলী হাইস্কুলে ৭০\% পরীক্ষার্থী ইংরেজিতে এবং ৮০\%
বাংলায় পাস করে। উভয় বিষয়ে ফেল করেছে ১০\%, যদি উভয় বিষয়ে ৩৬০ জন পাস করে
তবে ঐ স্কুলে কত জন পরীক্ষার্থী পরীক্ষা দিয়েছে? {[}ডাক ও টেলিযোগাযোগ
মন্ত্রণালয়ের হিসাবরক্ষণ কর্মকর্তা -২০০৩{]}

\begin{itemize}
\item
  (ক)৫০০ জন
\item
  (খ) ৭০০ জন
\item
  (গ) ৮০০ জন
\item
  (ঘ) ৬০০ জন
\end{itemize}

\textbf{\hl{উত্তরঃ (ঘ) ৬০০ জন}}

\textbf{সমাধান:}

শর্ট টেকনিক :

এখানে, ১ম বিষয় ইংরেজিতে পাশের হার = ৭০\%

২য় বিষয় বাংলায় পাশের হার = ৮০\%

উভয় বিষয়ে ফেলের হার ১০\% এবং উভয় বিষয়ে পাসকৃত ছাত্র ৩৬০ জন দেওয়া আছে।

তাহলে, ১ম বিষয় ইংরেজিতে ফেলের হার = ১০০\% - ৭০\% = ৩০\%

২য় বিষয় বাংলায় ফেলের হার = ১০০\% - ৮০\% = ২০\%

মোট শিক্ষার্থী =
\(\frac{\text{উভয়}\ \text{বিষয়ে}\ \text{পাসকৃত}\ \text{ছাত্র}}{\text{১ম}\ \text{বিষয়ে}\ \text{ফেল}\  + \ \text{২য়}\ \text{বিষয়ে}\ \text{ফেল}\  + \ \text{উভয়}\ \text{বিষয়ে}\ \text{ফেল}}\)

×১০০ = \(\frac{\text{৩৬০}}{\text{৩০} + \text{২০} + \text{১০}}\)

×১০০ = \(\frac{\text{৩৬০}}{\text{৬০}}\)

×১০০ = ৬০০ জন

লিখিত সমাধানঃ

শুধু ইংরেজিতে ফেল করে = (১০০-৭০)-১০\% = ২০\%

শুধু বাংলায় ফেল করে = (১০০-৮০)-১০\% = ১০\%

উভয় বিষয়ে ফেল করেছে ১০\%

দুই বিষয়ে মোট ফেল করে = ২০+১০+১০ = ৪০\%

দুই বিষয়ে মোট পাস করে = ১০০-৪০ = ৬০\%

৬০ জন পাস করে = ১০০ জনের মধ্যে

১ জন পাস করে = \(\frac{\text{১০০}}{\text{৬০}}\) জনের মধ্যে

৩৬০ জন পাস করে = \(\frac{\text{১০০}\  \times \ \text{৩৬০}}{\text{৬০}}\) = ৬০০ জনের মধ্যে

ঐ স্কুলে পরীক্ষা দিয়েছে = ৬০০ জন পরীক্ষার্থী

\paragraph{সূত্র-৪: উভয় বিষয়ে ফেল উল্লেখ না থাকলে মোট পরীক্ষার্থীর সংখ্যা
নির্ণয়ের ক্ষেত্রে- মোট পরীক্ষার্থী = (উভয় বিষয়ে পাসকৃত ছাত্র/১০০ − ( ১ম বিষয়ে
ফেল + ২য় বিষয়ে ফেল))
×১০০}\label{ux9b8ux9a4ux9b0-ux9ea-ux989ux9adux9df-ux9acux9b7ux9df-ux9abux9b2-ux989ux9b2ux9b2ux996-ux9a8-ux9a5ux995ux9b2-ux9aeux99f-ux9aaux9b0ux995ux9b7ux9b0ux9a5ux9b0-ux9b8ux996ux9af-ux9a8ux9b0ux9a3ux9dfux9b0-ux995ux9b7ux9a4ux9b0--ux9aeux99f-ux9aaux9b0ux995ux9b7ux9b0ux9a5-ux989ux9adux9df-ux9acux9b7ux9df-ux9aaux9b8ux995ux9a4-ux99bux9a4ux9b0ux9e7ux9e6ux9e6-ux9e7ux9ae-ux9acux9b7ux9df-ux9abux9b2-ux9e8ux9df-ux9acux9b7ux9df-ux9abux9b2-ux9e7ux9e6ux9e6}

\textbf{প্রশ্ন নম্বর-} ১. কোন পরীক্ষায় ৯০\% পরীক্ষার্থী ইতিহাসে এবং ৮৫\%
পরীক্ষার্থী ভূগোলে কৃতকার্য হল।যদি উভয় বিষয়ে কেউ ফেল না করে এবং উভয় বিষয়ে ২২৫
জন পাস করে থাকে , তবে ঐ পরীক্ষায় কতজন অংশগ্রহণ করেছিল?{[} দুর্নীতি দমন ব্যুরোর
উপ-পরিদর্শক-২০০৪{]}

\begin{itemize}
\item
  (ক) ৩০০
\item
  (খ) ৪০০
\item
  (গ) ৫০০
\item
  (ঘ) ৬০০
\end{itemize}

\textbf{\hl{উত্তরঃ (ক) ৩০০}}

\textbf{সমাধান:}

শর্ট টেকনিক :

এখানে, ১ম বিষয় ইতিহাসে পাশের হার = ৯০\%

২য় বিষয় ভূগোলে পাশের হার = ৮৫\%

এবং উভয় বিষয়ে মোট পাসকৃত ছাত্র ২২৫ জন দেওয়া আছে।

তাহলে, ১ম বিষয় ইতিহাসে ফেলের হার = ১০০\% - ৯০\% = ১০\%

২য় বিষয় ভূগোলে ফেলের হার = ১০০\% - ৮৫\% = ১৫\%

মোট পরীক্ষার্থী =
\(\frac{\text{উভয়}\ \text{বিষয়ে}\ \text{পাসকৃত}\ \text{ছাত্র}}{\text{১০০}\  - \ (\ \text{১ম}\ \text{বিষয়ে}\ \text{ফেল}\  + \ \text{২য়}\ \text{বিষয়ে}\ \text{ফেল})}\)

×১০০ = \(\frac{\text{২২৫}}{\text{১০০}\  - \ (\ \text{১০}\  + \ \text{১৫})}\)

×১০০ = \(\frac{\text{২২৫}}{\text{১০০}\  - \ \text{২৫}}\)

×১০০ = \(\frac{\text{২২৫}}{\text{৭৫}}\)

×১০০ = ৩×১০০ = ৩০০ লিখিত সমাধানঃ

শুধু ইতিহাসে ফেল করে = ১০০-৯০ = ১০\%

শুধু ভূগোলে ফেল করে = ১০০-৮৫ = ১৫\%

উভয় বিষয়ে পাস করে = ১০০-(১০\%+১৫\%) = ৭৫\%

৭৫ জন পাস করে = ১০০ জনের মধ্যে

১ জন পাস করে = \(\frac{\text{১০০}}{\text{৭৫}}\) জনের মধ্যে

২২৫ জন পাস করে = \(\frac{\text{২২৫} \times \text{১০০}}{\text{৭৫}}\) = ৩০০ জনের মধ্যে

সুতরাং, ঐ পরীক্ষায় ৩০০ জন অংশগ্রহণ করেছিল।

\paragraph{সূত্র-৫: কোন একটি বিষয়ে পাশের হার ও ফেলের সংখ্যা উল্লেখ থাকলে মোট
পরীক্ষার্থীর সংখ্যা নির্ণয়ের ক্ষেত্রে- মোট পরীক্ষার্থী = (ফেলের সংখ্যা/১০০ −
পাশের হার)
×১০০}\label{ux9b8ux9a4ux9b0-ux9eb-ux995ux9a8-ux98fux995ux99f-ux9acux9b7ux9df-ux9aaux9b6ux9b0-ux9b9ux9b0-ux993-ux9abux9b2ux9b0-ux9b8ux996ux9af-ux989ux9b2ux9b2ux996-ux9a5ux995ux9b2-ux9aeux99f-ux9aaux9b0ux995ux9b7ux9b0ux9a5ux9b0-ux9b8ux996ux9af-ux9a8ux9b0ux9a3ux9dfux9b0-ux995ux9b7ux9a4ux9b0--ux9aeux99f-ux9aaux9b0ux995ux9b7ux9b0ux9a5-ux9abux9b2ux9b0-ux9b8ux996ux9afux9e7ux9e6ux9e6-ux9aaux9b6ux9b0-ux9b9ux9b0-ux9e7ux9e6ux9e6}

\textbf{প্রশ্ন নম্বর-} ১. কোন পরীক্ষায় ৮৫\% ছাত্র ইংরেজিতে পাস করেছে।ইংরেজিতে
ফেলের সংখ্যা মোট ৭৫ জন হলে পরীক্ষার্থীর সংখ্যা কত?{[}প্রাথমিক সহকারী
শিক্ষক-২০০৬(বরিশাল){]}

\begin{itemize}
\item
  (ক) ৫০০ জন
\item
  (খ) ৩০০ জন
\item
  (গ) ৭৭৫ জন
\item
  (ঘ) ৬৫০ জন
\end{itemize}

\textbf{\hl{উত্তরঃ (ক) ৫০০ জন}}

\textbf{সমাধান:}

শর্ট টেকনিক :

এখানে , ইংরেজিতে পাশের হার = ৮৫\%

এবং ফেলের সংখ্যা মোট ৭৫ জন দেওয়া আছে।

মোট পরীক্ষার্থী = \(\frac{\text{ফেলের}\ \text{সংখ্যা}}{\text{১০০}\  - \ \text{পাশের}\ \text{হার}}\) × ১০০

= \(\frac{\text{৭৫}}{\text{১০০}\  - \ \text{৮৫}}\)×১০০

= \(\frac{\text{৭৫}}{\text{১৫}}\)×১০০

= ৫০০ জন

লিখিত সমাধানঃ

পরীক্ষার্থীর সংখ্যা ১০০\%

ইংরেজিতে ফেল করে = ১০০-৮৫ = ১৫\%

প্রশ্নমতে,

১৫\% = ৭৫ জন

বা,১\% = ৭৫÷১৫ = ৫ জন

বা, ১০০\% = ৫×১০০ = ৫০০ জন

\section{53. শতকরা MCQ প্রশ্ন উত্তর
অনুশীলন}\label{ux9b6ux9a4ux995ux9b0-mcq-ux9aaux9b0ux9b6ux9a8-ux989ux9a4ux9a4ux9b0-ux985ux9a8ux9b6ux9b2ux9a8}

\textbf{শতকরা বলতে কি বুঝায়?}

ক) একটি ভগ্নাংশ

খ) একটি সম্পূর্ণ সংখ্যা

\textbf{\hl{গ) একটি ভগ্নাংশ যার হর ১০০ ✔}}

ঘ) একটি ভগ্নাংশ যার হর ১ ও লব ১০০

\textbf{\hl{✔ সঠিক উত্তর: গ) একটি ভগ্নাংশ যার হর ১০০}}

\begin{quote}
\textbf{ব্যাখ্যা:} \emph{শতকরা একটি ভগ্নাংশ, যার হর ১০০। শতকরা শব্দটিকে
সংক্ষেপে \textquotesingle\%\textquotesingle{} এই প্রতীক দিয়ে প্রকাশ করা
হয়।}
\end{quote}

\textbf{The symbol \% stands for-\/-}

ক) per capita

খ) At the rate of

\textbf{\hl{গ) percentage ✔}}

ঘ) Per head

\textbf{\hl{✔ সঠিক উত্তর: গ) percentage}}

\begin{quote}
\textbf{ব্যাখ্যা:} \emph{শতকরা শব্দটিকে সংক্ষেপে
\textquotesingle\%\textquotesingle{} ( percentage )এই প্রতীক দিয়ে প্রকাশ
করা হয়।}
\end{quote}

\textbf{নিচের কোন ভগ্নাংশটি ০.৫\% এর সমতুল্য?}

ক) (১/২০)

\textbf{\hl{খ) (১/২০০) ✔}}

গ) (১/২০০০)

ঘ) (১/৫০০)

\textbf{\hl{✔ সঠিক উত্তর: খ) (১/২০০)}}

\begin{quote}
\textbf{ব্যাখ্যা:} \emph{০.৫\%}

\emph{= (৫/১০) \%}

\emph{= (১/২) × (১/১০০)}

\emph{= (১/২০০)}
\end{quote}

\textbf{নিচের কোনটি ০.৪৫ এর সমান?}

ক) ০.০৪৫\%

খ) ০.৪৫\%

গ) ৪.৫\%

\textbf{\hl{ঘ) ৪৫\% ✔}}

\textbf{\hl{✔ সঠিক উত্তর: ঘ) ৪৫\%}}

\begin{quote}
\textbf{ব্যাখ্যা:} \emph{০.৪৫ = (৪৫/১০০) = ৪৫\%}
\end{quote}

\textbf{(২৫/২) \% এর সমান ভগ্নাংশ কত হবে?}

ক) (১/১৬)

\textbf{\hl{খ) (১/৮) ✔}}

গ) (১/৪)

ঘ) (২/২৫)

\textbf{\hl{✔ সঠিক উত্তর: খ) (১/৮)}}

\begin{quote}
\textbf{ব্যাখ্যা:} \emph{(২৫/২) \%}

\emph{= (২৫/২) × (১/১০০)}

\emph{= (২৫/২০০)}

\emph{= (১/৮)}
\end{quote}

\textbf{১০০ টাকার (১/২) \% সমান কত?}

ক) ৫০ টাকা

\textbf{\hl{খ) ০.৫০ টাকা ✔}}

গ) ০.০৫ টাকা

ঘ) ৫ টাকা

\textbf{\hl{✔ সঠিক উত্তর: খ) ০.৫০ টাকা}}

\begin{quote}
\textbf{ব্যাখ্যা:} \emph{১০০ টাকার (১/২) \%}

\emph{= ১০০ টাকার (১/২) × (১/১০০)}

\emph{= ১০০ টাকার (১/২০০)}

\emph{= (১/২) টাকা}

\emph{= ০.৫০ টাকা}
\end{quote}

\textbf{একশত টাকার শতকরা দুই ভাগ কত?}

ক) ২০০ টাকা

খ) ২০ টাকা

\textbf{\hl{গ) ২ টাকা ✔}}

ঘ) ২০০০ টাকা

\textbf{\hl{✔ সঠিক উত্তর: গ) ২ টাকা}}

\begin{quote}
\textbf{ব্যাখ্যা:} \emph{একশত টাকার শতকরা দুই ভাগ ২ টাকা।}
\end{quote}

\textbf{১৫ টাকার শতকরা ৭ শতাংশ কত হয়?}

\textbf{\hl{ক) ১.০৫ ✔}}

খ) ১০.৫

গ) ১.৫

ঘ) ৭.৫

\textbf{\hl{✔ সঠিক উত্তর: ক) ১.০৫}}

\begin{quote}
\textbf{ব্যাখ্যা:} \emph{১৫ টাকার (৭/১০০)}

\emph{= (২১/২০) = ১.০৫}
\end{quote}

\textbf{৬৬ লিটারের ১.২\% কত?}

ক) ৬.০১ লিটার

খ) ১.২২ লিটার

গ) ০.৬৯২ লিটার

\textbf{\hl{ঘ) ০.৭৯২ লিটার ✔}}

\textbf{\hl{✔ সঠিক উত্তর: ঘ) ০.৭৯২ লিটার}}

\begin{quote}
\textbf{ব্যাখ্যা:} \emph{৬৬ লিটারের ১.২\%}

\emph{= ৬৬ এর (১২/১০) × (১/১০০)}

\emph{= ৬৬ এর (৩/২৫০)}

\emph{= (১৯৮/২৫০)}

\emph{= ০.৭৯২ লিটার}
\end{quote}

\textbf{০.০২৩ এর ১\% হচ্ছে?}

ক) ০২৩

খ) ০.০০২৩

\textbf{\hl{গ) ০.০০০২৩ ✔}}

ঘ) ২.৩

\textbf{\hl{✔ সঠিক উত্তর: গ) ০.০০০২৩}}

\begin{quote}
\textbf{ব্যাখ্যা:} \emph{০.০২৩ এর ১\%}

\emph{= (২৩/১০০০) এর (১/১০০)}

\emph{= (২৩/১০০০০০) = ০.০০০২৩}
\end{quote}

\textbf{০.০২৫ এর শতকরা ১ ভাগ কত?}

ক) ০.০২৫

খ) ০.০০২৫

\textbf{\hl{গ) ০.০০০২৫ ✔}}

ঘ) ০.০০০০২৫

ঙ) ০.০০০০০০২৫

\textbf{\hl{✔ সঠিক উত্তর: গ) ০.০০০২৫}}

\begin{quote}
\textbf{ব্যাখ্যা:} \emph{০.০২৫ এর ১\%}

\emph{= (২৫/১০০০) এর ১\%}

\emph{= (২৫/১০০০) × (১/১০০)}

\emph{= (২৫/১০০০০০) = ০.০০০২৫}
\end{quote}

\textbf{১৬.৫ এর ১.৩\% কত?}

ক) ২.১৪৫

খ) ২১.৪৫

গ) ০.০২১৪৫

\textbf{\hl{ঘ) ০.২১৪৫ ✔}}

\textbf{\hl{✔ সঠিক উত্তর: ঘ) ০.২১৪৫}}

\begin{quote}
\textbf{ব্যাখ্যা:} \emph{১৬.৫ এর ১.৩\%}

\emph{= (১৬৫/১০) এর ১.৩\%}

\emph{= (১৬৫/১০) × (১৩/১০০০)}

\emph{= (২১৪৫/১০০০০) = ০.২১৪৫}
\end{quote}

\textbf{২.৫ কোন সংখ্যার ০.৫\% ?}

ক) ২৫০

খ) ১২৫

\textbf{\hl{গ) ৫০০ ✔}}

ঘ) ২২৫

\textbf{\hl{✔ সঠিক উত্তর: গ) ৫০০}}

\begin{quote}
\textbf{ব্যাখ্যা:} \emph{০.৫\% = ২.৫}

\emph{⇒ ১\% = (২.৫/০.৫) = ৫}

\emph{∴ ১০০\% = ৫ × ১০০ = ৫০০}
\end{quote}

\textbf{৫-এর কত শতাংশ ৭ হবে?}

ক) ৪০

খ) ১২৫

গ) ৯০

\textbf{\hl{ঘ) ১৪০ ✔}}

\textbf{\hl{✔ সঠিক উত্তর: ঘ) ১৪০}}

\begin{quote}
\textbf{ব্যাখ্যা:} \emph{ধরি, ৫ এর x\% = ৭}

\emph{বা, ৫ × (x/১০০) = ৭}

\emph{বা, (x/২০) = ৭}

\emph{⇒ x = ৭ × ২০}

\emph{∴ x = ১৪০}
\end{quote}

\textbf{২০ এর ২০\% সমান কত?}

\textbf{\hl{ক) ৪ ✔}}

খ) ৮

গ) ৬

ঘ) ৫

\textbf{\hl{✔ সঠিক উত্তর: ক) ৪}}

\begin{quote}
\textbf{ব্যাখ্যা:} \emph{২০ এর ২০\% = ২০ × (২০/১০০) = ৪}
\end{quote}

\textbf{১২৫ এর ১২৫\% = কত?}

ক) ১৩.২৫

খ) ১৩১.২৫

গ) ১৫০

\textbf{\hl{ঘ) ১৫৬.২৫ ✔}}

\textbf{\hl{✔ সঠিক উত্তর: ঘ) ১৫৬.২৫}}

\begin{quote}
\textbf{ব্যাখ্যা:} \emph{১২৫ এর ১২৫\% = ১২৫ × (১২৫/১০০) = ১৫৬.২৫}
\end{quote}

\textbf{৫০ এর ১০\% কত?}

ক) ১০

খ) ১৫

গ) ১২

\textbf{\hl{ঘ) ৫ ✔}}

\textbf{\hl{✔ সঠিক উত্তর: ঘ) ৫}}

\begin{quote}
\textbf{ব্যাখ্যা:} \emph{৫০ এর ১০\% = ৫০ × (১০/১০০) = ৫}
\end{quote}

\textbf{০.২ এর ২০\% কত?}

ক) ১

খ) ৪

\textbf{\hl{গ) ০.০৪ ✔}}

ঘ) ০.৪

\textbf{\hl{✔ সঠিক উত্তর: গ) ০.০৪}}

\begin{quote}
\textbf{ব্যাখ্যা:} \emph{০.২ এর ২০\% = ০.২ × (২০/১০০) = ০.০৪}
\end{quote}

\textbf{(২/৫) এর ২৫\% সমান কত?}

\textbf{\hl{ক) ০.১ ✔}}

খ) ০.৩

গ) ০.২

ঘ) ০.৪

\textbf{\hl{✔ সঠিক উত্তর: ক) ০.১}}

\begin{quote}
\textbf{ব্যাখ্যা:} \emph{(২/৫) এর ২৫\% = (২/৫) × (২৫/১০০) = (১/১০) = ০.১}
\end{quote}

\textbf{80 এর 75\% এর 25\% = কত?}

ক) 10

\textbf{\hl{খ) 15 ✔}}

গ) 20

ঘ) 25

\textbf{\hl{✔ সঠিক উত্তর: খ) 15}}

\begin{quote}
\textbf{ব্যাখ্যা:} \emph{৮০ এর ৭৫\% = ৮০ × (৭৫/১০০) = ৬০}

\emph{∴ ৬০ এর ২৫\% = ৬০ × (২৫/১০০) = ১৫}
\end{quote}

\textbf{২০০ টাকার ১২ (১/২) \% = ?}

ক) ১৫

খ) ২০

\textbf{\hl{গ) ২৫ ✔}}

ঘ) ৩০

\textbf{\hl{✔ সঠিক উত্তর: গ) ২৫}}

\begin{quote}
\textbf{ব্যাখ্যা:} \emph{২০০ টাকার ১২ (১/২) \% = ২০০ × (২৫/২) \% = ২০০ ×
(২৫/২০০) = ২৫}
\end{quote}

\textbf{১০০ টাকার ১০০\% = কত টাকা?}

ক) ৫০ টাকা

\textbf{\hl{খ) ১০০ টাকা ✔}}

গ) ২০০ টাকা

ঘ) ১০০০ টাকা

\textbf{\hl{✔ সঠিক উত্তর: খ) ১০০ টাকা}}

\begin{quote}
\textbf{ব্যাখ্যা:} \emph{১০০ টাকার ১০০\% = ১০০ × (১০০/১০০) = ১০০ টাকা}
\end{quote}

\textbf{৫৪০ এর ৮.৫\% = ?}

ক) ৪৪

খ) ৪৫

গ) ৪৬

ঘ) ৪৭

\textbf{\hl{ঙ) কোনোটিই নয় ✔}}

\textbf{\hl{✔ সঠিক উত্তর: ঙ) কোনোটিই নয়}}

\begin{quote}
\textbf{ব্যাখ্যা:} \emph{৫৪০ এর ৮.৫\% = ৫৪০ × (৮.৫/১০০) = ৪৫.৯}

\emph{∴ সঠিক উত্তরটি অপশনে নেই (কোনোটিই নয়)।}
\end{quote}

\textbf{১৪৪ কোন সংখ্যার ৪০\%?}

ক) ২৬০

খ) ২৯০

\textbf{\hl{গ) ৩৬০ ✔}}

ঘ) ৩৮০

\textbf{\hl{✔ সঠিক উত্তর: গ) ৩৬০}}

\begin{quote}
\textbf{ব্যাখ্যা:} \emph{ধরি, সংখ্যাটি x}

\emph{প্রশ্নমতে, x এর ৪০\% = ১৪৪}

\emph{বা, x × (৪০/১০০) = ১৪৪}

\emph{⇒ x = (১৪৪ × ১০০/৪০)}

\emph{∴ x = ৩৬০}
\end{quote}

\textbf{৪৮ কোন সংখ্যার ৬০\%?}

ক) ৫০

খ) ৬০

গ) ৭০

\textbf{\hl{ঘ) ৮০ ✔}}

\textbf{\hl{✔ সঠিক উত্তর: ঘ) ৮০}}

\begin{quote}
\textbf{ব্যাখ্যা:} \emph{ধরি, সংখ্যাটি = x}

\emph{প্রশ্নমতে, x এর ৬০\% = ৪৮}

\emph{বা, x × (৬০/১০০) = ৪৮}

\emph{বা, x = ৪৮ × (১০০/৬০)}

\emph{∴ x = ৮০}
\end{quote}

\textbf{কোন সংখ্যার ১৬ (২/৩) \% সমান ১০০?}

ক) ৫০০

খ) ৫৫০

\textbf{\hl{গ) ৬০০ ✔}}

ঘ) ৮৫০

\textbf{\hl{✔ সঠিক উত্তর: গ) ৬০০}}

\begin{quote}
\textbf{ব্যাখ্যা:} \emph{ধরি, সংখ্যাটি x}

\emph{প্রশ্নমতে, x এর ১৬ (২/৩) \% = ১০০}

\emph{বা, x × (৫০/৩ × ১০০) = ১০০}

\emph{বা, (x/৬) = ১০০}

\emph{∴ x = ৬০০}
\end{quote}

\textbf{২৫০ টাকার শতকরা কত সমান ১০ টাকা?}

ক) ৩\%

খ) ২\%

\textbf{\hl{গ) ৪\% ✔}}

\textbf{\hl{✔ সঠিক উত্তর: গ) ৪\%}}

\begin{quote}
\textbf{ব্যাখ্যা:} \emph{ধরি, ২৫০ টাকার x\% = ১০ টাকা}

\emph{⇒ ২৫০ × (x/১০০) = ১০}

\emph{⇒ x = (১০০ × ১০/২৫০)}

\emph{∴ x = ৪\%}
\end{quote}

\textbf{৩০, ৫০ এর শতকরা কত অংশ?}

ক) ৩০\%

খ) ৫০\%

গ) ৫৫\%

\textbf{\hl{ঘ) ৬০\% ✔}}

\textbf{\hl{✔ সঠিক উত্তর: ঘ) ৬০\%}}

\begin{quote}
\textbf{ব্যাখ্যা:} \emph{নির্ণেয় শতাংশ = (৩০/৫০) × ১০০\% = ৬০\%}
\end{quote}

\textbf{(১/২) এর শতকরা কত (৩/৪) হবে?}

ক) ১২০\%

খ) ১২৫\%

গ) ১৪০\%

\textbf{\hl{ঘ) ১৫০\% ✔}}

\textbf{\hl{✔ সঠিক উত্তর: ঘ) ১৫০\%}}

\begin{quote}
\textbf{ব্যাখ্যা:} \emph{ধরি, (১/২) এর x\% = (৩/৪)}

\emph{বা, (১/২) × (x/১০০) = (৩/৪)}

\emph{বা, x = (৩ × ২০০/৪)}

\emph{∴ x = ১৫০\%}
\end{quote}

\textbf{যদি একটি সংখ্যার ৩০ শতাংশের ১৫ শতাংশ ১৮ হয়, তবে সংখ্যাটি কত?}

ক) ৯

খ) ৩৬

\textbf{\hl{গ) ৪০০ ✔}}

ঘ) ৮১

\textbf{\hl{✔ সঠিক উত্তর: গ) ৪০০}}

\begin{quote}
\textbf{ব্যাখ্যা:} \emph{ধরি, সংখ্যাটি x}

\emph{প্রশ্নমতে, x × ৩০\% × ১৫\% = ১৮}

\emph{বা, x × (৩০/১০০) × (১৫/১০০) = ১৮}

\emph{বা, x × (৪৫০/১০০০০) = ১৮}

\emph{⇒ x = (১৮ × ১০০০০/৪৫০)}

\emph{∴ x = ৪০০}
\end{quote}

\textbf{x ধনাত্মক হলে ২৫, ৫x-এর শতকরা কত?}

ক) (x/১০০) \%

খ) (x/৫) \%

গ) (২৫/x) \%

\textbf{\hl{ঘ) (৫০০/x) \% ✔}}

\textbf{\hl{✔ সঠিক উত্তর: ঘ) (৫০০/x) \%}}

\begin{quote}
\textbf{ব্যাখ্যা:} \emph{২৫, ৫x-এর শতকরা = ( (২৫/৫x) × ১০০)\% = (৫০০/x)
\%}
\end{quote}

\textbf{৩৫\% কে দশমিকে প্রকাশ করুন।}

ক) ০.০৩৫

\textbf{\hl{খ) ০.৩৫ ✔}}

গ) ৩.০৫

ঘ) ৩৫

ঙ) কোনোটিই নয়

\textbf{\hl{✔ সঠিক উত্তর: খ) ০.৩৫}}

\begin{quote}
\textbf{ব্যাখ্যা:} \emph{৩৫\% = (৩৫/১০০) = ০.৩৫}
\end{quote}

\textbf{১৫০\%-এর সমান ভগ্নাংশ কত?}

ক) (৬/৫)

খ) (৩/৫)

\textbf{\hl{গ) (৩/২) ✔}}

ঘ) (৪/৫)

\textbf{\hl{✔ সঠিক উত্তর: গ) (৩/২)}}

\begin{quote}
\textbf{ব্যাখ্যা:} \emph{১৫০\% = (১৫০/১০০) = (৩/২)}
\end{quote}

\textbf{১৩ (৩/৪) \%-এর মান কত?}

ক) (১১/২০)

\textbf{\hl{খ) (১১/৮০) ✔}}

গ) (১/৮)

ঘ) (১/৯)

\textbf{\hl{✔ সঠিক উত্তর: খ) (১১/৮০)}}

\begin{quote}
\textbf{ব্যাখ্যা:} \emph{১৩ (৩/৪) \% = (৫৫/৪) \% = (৫৫/৪) × (১/১০০) =
(১১/৮০)}
\end{quote}

\textbf{(৪/৫) কে শতকরায় প্রকাশ করলে কত হয়?}

ক) ৪০\%

খ) ৫০\%

গ) ৬০\%

\textbf{\hl{ঘ) ৮০\% ✔}}

\textbf{\hl{✔ সঠিক উত্তর: ঘ) ৮০\%}}

\begin{quote}
\textbf{ব্যাখ্যা:} \emph{(৪/৫) কে শতকরায় প্রকাশ করলে হয়, ( (৪/৫) × ১০০)\% =
৮০\%}
\end{quote}

\textbf{৩ : ৫ কে শতকরায় প্রকাশ করলে কত হয়?}

ক) ৭০\%

\textbf{\hl{খ) ৬০\% ✔}}

গ) ৩০\%

ঘ) ৫০\%

\textbf{\hl{✔ সঠিক উত্তর: খ) ৬০\%}}

\begin{quote}
\textbf{ব্যাখ্যা:} \emph{৩ : ৫ = (৩/৫)}

\emph{শতকরায় প্রকাশ করলে হয় = ( (৩/৫) × ১০০)\% = ৬০\%}
\end{quote}

\textbf{x এর ২\% এর মান ০.০৩ হলে x এর মান কত?}

ক) ১৫০

খ) ৬৬.৬৭

গ) ১৫

ঘ) ৬.৬৭

\textbf{\hl{ঙ) ১.৫ ✔}}

\textbf{\hl{✔ সঠিক উত্তর: ঙ) ১.৫}}

\begin{quote}
\textbf{ব্যাখ্যা:} \emph{এখানে x অর্থাৎ ১০০\% এর মান বের করতে হবে।}

\emph{২\% = ০.০৩}

\emph{⇒ ১\% = (০.০৩/২)}

\emph{∴ ১০০\% = (০.০৩×১০০/২) = ০.০৩ × ৫০ = ১.৫}
\end{quote}

\textbf{একটি সংখ্যাকে ৫০ দ্বারা ভাগ করলে ভাগফল হয় ৩.৬\%। সংখ্যাটি কত?}

\textbf{\hl{ক) ১.৮ ✔}}

খ) ১৮

গ) ১৮০

ঘ) কোনটিই না

\textbf{\hl{✔ সঠিক উত্তর: ক) ১.৮}}

\begin{quote}
\textbf{ব্যাখ্যা:} \emph{এখানে, ৩.৬\% = ৩.৬ × (১/১০০) = ০.০৩৬}

\emph{∴ সংখ্যাটি = ০.০৩৬ × ৫০ = ১.৮}
\end{quote}

\textbf{৪৮০০ এর ১০\%, ৪৮০০ এর ৮\% অপেক্ষা কত বেশি?}

ক) ৮৪

খ) ৮০

\textbf{\hl{গ) ৯৬ ✔}}

ঘ) ৯৮

ঙ) ৯০

\textbf{\hl{✔ সঠিক উত্তর: গ) ৯৬}}

\begin{quote}
\textbf{ব্যাখ্যা:} \emph{৪৮০০ এর ১০\% = ৪৮০}

\emph{৪৮০০ এর ৮\% = ৩৮৪}

\emph{∴ বেশি = ৪৮০ − ৩৮৪ = ৯৬}
\end{quote}

\textbf{৯ টাকার শতকরা ৭ অপেক্ষা ১১ টাকার শতকরা ৬ কত বেশি?}

ক) ০.০০৩ টাকা

\textbf{\hl{খ) ০.০৩ টাকা ✔}}

গ) ০.৩০ টাকা

ঘ) ৩.০০ টাকা

\textbf{\hl{✔ সঠিক উত্তর: খ) ০.০৩ টাকা}}

\begin{quote}
\textbf{ব্যাখ্যা:} \emph{৯ টাকার ৭\% = ০.৬৩}

\emph{১১ টাকার ৬\% = ০.৬৬}

\emph{∴ বেশি = ০.৬৬ − ০.৬৩ = ০.০৩ টাকা}
\end{quote}

\textbf{৯০ কোন সংখ্যার ৭৫\%?}

\textbf{\hl{ক) ১২০ ✔}}

খ) ১২৫

গ) ১৫০

ঘ) ২৭৫

\textbf{\hl{✔ সঠিক উত্তর: ক) ১২০}}

\begin{quote}
\textbf{ব্যাখ্যা:} \emph{৭৫\% = ৯০}

\emph{⇒ ১\% = ১.২}

\emph{∴ ১০০\% = ১.২ × ১০০ = ১২০}
\end{quote}

\textbf{৭ কোন সংখ্যার ৫\%?}

ক) ১২৮

খ) ১৫৬

\textbf{\hl{গ) ১৪০ ✔}}

ঘ) ১৩৫

\textbf{\hl{✔ সঠিক উত্তর: গ) ১৪০}}

\begin{quote}
\textbf{ব্যাখ্যা:} \emph{৫\% = ৭}

\emph{⇒ ১\% = ১.৪}

\emph{∴ ১০০\% = ১.৪ × ১০০ = ১৪০}
\end{quote}

\textbf{২৯ সংখ্যাটি কোন সংখ্যার ১০\%?}

ক) ৮৭

খ) ২৭০

গ) ২৬৯

ঘ) ২৪৯

\textbf{\hl{ঙ) ২৯০ ✔}}

\textbf{\hl{✔ সঠিক উত্তর: ঙ) ২৯০}}

\begin{quote}
\textbf{ব্যাখ্যা:} \emph{১০\% = ২৯}

\emph{⇒ ১\% = ২.৯}

\emph{∴ ১০০\% = ২.৯ × ১০০ = ২৯০}
\end{quote}

\textbf{একটি দ্রব্যের ক্রয়মূল্যের ১০\% বাবদ ১১০ টাকা পরিশোধ করা হল। আর বকেয়া
কত টাকা পরিশোধ করতে হবে?}

ক) ৮৮০

\textbf{\hl{খ) ৯৯০ ✔}}

গ) ১০০০

ঘ) ১১০০

ঙ) ১২১০

\textbf{\hl{✔ সঠিক উত্তর: খ) ৯৯০}}

\begin{quote}
\textbf{ব্যাখ্যা:} \emph{পরিশোধ করা হল ১০\%}

\emph{বকেয়া রইল = (১০০\% − ১০\%) = ৯০\%}

\emph{১০\% = ১১০ টাকা}

\emph{⇒ ১\% = ১১ টাকা}

\emph{∴ ৯০\% = ১১ × ৯০ = ৯৯০ টাকা}
\end{quote}

\textbf{একটি সংখ্যার ১২\% নিলে ৯৬ পাওয়া যায়, সংখ্যাটি কত?}

ক) ১২০

খ) ১০০০

গ) ৭২০

\textbf{\hl{ঘ) ৮০০ ✔}}

\textbf{\hl{✔ সঠিক উত্তর: ঘ) ৮০০}}

\begin{quote}
\textbf{ব্যাখ্যা:} \emph{১২\% = ৯৬}

\emph{⇒ ১\% = (৯৬/১২) = ৮}

\emph{∴ ১০০\% = ৮ × ১০০ = ৮০০}
\end{quote}

\textbf{কোন সংখ্যাকে ২৫০ দ্বারা ভাগ করলে ভাগফল ১৫\% এর সমান হয়?}

ক) ২৫

খ) ১৫

গ) ৪৫

ঘ) ৬০

ঙ) ৮০

\textbf{\hl{চ) ৩৭.৫ ✔}}

\textbf{\hl{✔ সঠিক উত্তর: চ) ৩৭.৫}}

\begin{quote}
\textbf{ব্যাখ্যা:} \emph{১৫\% = ০.১৫}

\emph{∴ সংখ্যাটি = ০.১৫ × ২৫০ = ৩৭.৫}
\end{quote}

\textbf{১০ এর ৩০\% কোন সংখ্যার ১০\%?}

\textbf{\hl{ক) ৩০ ✔}}

খ) ৬০

গ) ৪০

ঘ) ৬০০

\textbf{\hl{✔ সঠিক উত্তর: ক) ৩০}}

\begin{quote}
\textbf{ব্যাখ্যা:} \emph{১০ এর ৩০\% = ৩}

\emph{আবার, ১০\% = ৩}

\emph{⇒ ১\% = (৩/১০)}

\emph{∴ ১০০\% = (৩×১০০/১০) = ৩০}
\end{quote}

\textbf{গতবছর বনানী উচ্চবিদ্যালয়ে ১১৭২ জন ছাত্রছাত্রী ছিল। এ বছরে গত বছরের
চেয়ে ১৫\% বেশি ছাত্রছাত্রী রয়েছে।এ বছর আনুমানিক কতজন ছাত্রছাত্রী রয়েছে?}

ক) ১৮০০

খ) ১৬০০

গ) ১৫০০

ঘ) ১৪০০

\textbf{\hl{ঙ) কোনটিই নয় ✔}}

\textbf{\hl{✔ সঠিক উত্তর: ঙ) কোনটিই নয়}}

\begin{quote}
\textbf{ব্যাখ্যা:} \emph{১১৭২ এর ১৫\% = ১৭৫.৮}

\emph{∴ এ বছর আনুমানিক ছাত্রছাত্রী রয়েছে = ১১৭২ + ১৭৫.৮ = ১৩৪৭.৮}
\end{quote}

\textbf{একটি ফুটবল দল বছরে ১৫৪ টি খেলা খেলে। যদি জুন পর্যন্ত দলটি ৪০ টি
খেলায় জয়লাভ করে এবং ২০ খেলায় পরাজিত হয়, তবে অবশিষ্ট খেলার মধ্যে কটিতে
জয়লাভ করলে তারা ৫০\% জয় দিয়ে বছর শেষ করতে পারবে?}

ক) ১৭

খ) ৭৭

\textbf{\hl{গ) ৩৭ ✔}}

ঘ) ২৭

\textbf{\hl{✔ সঠিক উত্তর: গ) ৩৭}}

\begin{quote}
\textbf{ব্যাখ্যা:} \emph{সবমোট খেলা = ১৫৪ টি}

\emph{৫০\% জয়ের জন্য দরকার = ১৫৪ এর ৫০\% = ৭৭ টি জয়}

\emph{যেহেতু তারা ৪০ টিতে জিতেছে}

\emph{∴ আরও দরকার = ৭৭ − ৪০ = ৩৭ টি}
\end{quote}

\textbf{সাগরের পানিতে লবণের পরিমাণ ৪.৫\%। ৫০ কেজি লবণ তৈরি করতে কত কেজি
পানি বাষ্প করতে হবে?}

ক) ১১২.২২ কেজি

খ) ১২.১২ কেজি

গ) ১১.১১ কেজি

\textbf{\hl{ঘ) ১১১১.১১ কেজি ✔}}

\textbf{\hl{✔ সঠিক উত্তর: ঘ) ১১১১.১১ কেজি}}

\begin{quote}
\textbf{ব্যাখ্যা:} \emph{৪.৫ কেজি লবণের জন্য পানি লাগে ১০০ কেজি}

\emph{∴ ৫০ কেজি লবণের জন্য লাগে = (৫০×১০০/৪.৫) = ১১১১.১১ কেজি}
\end{quote}

\textbf{কোন গ্রামের ২৫০০০ জনসংখ্যার মধ্যে ১৫০০ জন শিক্ষিত। ঐ গ্রামে শিক্ষিতের
হার কত?}

ক) ৪\%

খ) ৫\%

\textbf{\hl{গ) ৬\% ✔}}

ঘ) ৭\%

\textbf{\hl{✔ সঠিক উত্তর: গ) ৬\%}}

\begin{quote}
\textbf{ব্যাখ্যা:} \emph{শিক্ষিতের হার = (১৫০০×১০০/২৫০০০) \% = ৬\%}
\end{quote}

\textbf{একজন ছাত্র ৮০০ নম্বরের মধ্যে ২০০ নম্বর পেল। সে শতকরা কত নম্বর পেল?}

ক) ৭৫

\textbf{\hl{খ) ২৫ ✔}}

গ) ৪০

ঘ) ২

\textbf{\hl{✔ সঠিক উত্তর: খ) ২৫}}

\begin{quote}
\textbf{ব্যাখ্যা:} \emph{শতকরা নম্বর = (২০০/৮০০) = (১/৪) = ২৫\%}
\end{quote}

\textbf{একটি গণিতের বই প্রকৃত মূল্যের শতকরা ৮৫ ভাগ মূল্যে ৬৮ টাকায় বিক্রয় হলো?
বইটির প্রকৃত মূল্য কত?}

ক) ৭২ টাকা

খ) ৬০ টাকা

\textbf{\hl{গ) ৮০ টাকা ✔}}

ঘ) ৯০ টাকা

\textbf{\hl{✔ সঠিক উত্তর: গ) ৮০ টাকা}}

\begin{quote}
\textbf{ব্যাখ্যা:} \emph{৮৫ ভাগ = ৬৮ টাকা}

\emph{⇒ ১ ভাগ = (৬৮/৮৫) টাকা}

\emph{∴ ১০০ ভাগ = (৬৮×১০০/৮৫) = ৮০ টাকা}
\end{quote}

\textbf{১৫\% বৃদ্ধি পেয়ে বেতন ৫৭৫০ টাকা হলে পূর্বের বেতন কত ছিল?}

ক) ৫৫০০ টাকা

খ) ৫২৫০ টাকা

\textbf{\hl{গ) ৫০০০ টাকা ✔}}

ঘ) ৪৭৫০ টাকা

\textbf{\hl{✔ সঠিক উত্তর: গ) ৫০০০ টাকা}}

\begin{quote}
\textbf{ব্যাখ্যা:} \emph{১১৫ টাকা বর্তমান হলে পূর্বে ছিল ১০০ টাকা}

\emph{∴ ৫৭৫০ টাকা বর্তমান হলে পূর্বে ছিল = (৫৭৫০×১০০/১১৫) = ৫০০০ টাকা}
\end{quote}

\textbf{এক ব্যক্তি তার মাসিক বেতনের শতকরা ১০ অংশ মহার্ঘ ভাতা পায়। তার মাসিক
আয় ৪২০০ টাকা হলে মহার্ঘ ভাতা কত?}

\textbf{\hl{ক) ৪২০ টাকা ✔}}

খ) ২০০ টাকা

গ) ৮৪০ টাকা

ঘ) ১০৫ টাকা

\textbf{\hl{✔ সঠিক উত্তর: ক) ৪২০ টাকা}}

\begin{quote}
\textbf{ব্যাখ্যা:} \emph{মহার্ঘ ভাতা = ৪২০০ এর ১০\%}

\emph{= (৪২০০×১০/১০০)}

\emph{= ৪২০ টাকা}
\end{quote}

\textbf{যদি ১৫ টি পোশাকের শতকরা ৪০ ভাগ পোশাক শার্ট হয় তবে ১৫ টি পোশাকের
মধ্যে কটি শার্ট নয়?}

ক) ৬

\textbf{\hl{খ) ৯ ✔}}

গ) ১২

ঘ) ১০

\textbf{\hl{✔ সঠিক উত্তর: খ) ৯}}

\begin{quote}
\textbf{ব্যাখ্যা:} \emph{শার্ট নয় এমন পোশাকের হার = (১০০ − ৪০)\% = ৬০\%}

\emph{∴ শার্ট নয় = ১৫ এর ৬০\%}

\emph{= (১৫×৬০/১০০)}

\emph{= ৯ টি}
\end{quote}

\textbf{করিমের বেতন রহিমের বেতন অপেক্ষা শতকরা ৫০ ভাগ বেশি।তাদের একমাসের
বেতনের যোগফল ১০০০০ টাকা। রহিমের বেতন কত?}

\textbf{\hl{ক) ৪০০০ টাকা ✔}}

খ) ৪৫০০ টাকা

গ) ৫০০০ টাকা

ঘ) ৫৫০০ টাকা

\textbf{\hl{✔ সঠিক উত্তর: ক) ৪০০০ টাকা}}

\begin{quote}
\textbf{ব্যাখ্যা:} \emph{করিম : রহিম = ১৫০ : ১০০ = ৩ : ২}

\emph{অনুপাতের যোগফল = ৩ + ২ = ৫}

\emph{∴ রহিমের বেতন = ১০০০০ এর (২/৫) = ৪০০০ টাকা}
\end{quote}

\textbf{১২০ জন ছাত্রের মধ্যে ৩০ ছাত্র ফেল করলে পাশের হার কত?}

\textbf{\hl{ক) ৭৫\% ✔}}

খ) ৮০\%

গ) ৪০\%

ঘ) ২৫\%

\textbf{\hl{✔ সঠিক উত্তর: ক) ৭৫\%}}

\begin{quote}
\textbf{ব্যাখ্যা:} \emph{পাশ করেছে = ১২০ − ৩০ = ৯০ জন}

\emph{∴ পাশের হার = (৯০×১০০/১২০) \% = ৭৫\%}
\end{quote}

\textbf{এক ব্যক্তি তার স্ত্রীর বয়সের ১০\% বড়। স্ত্রী তার চেয়ে x\% ছোট। x\% =
কত?}

ক) (৯১/১০)

খ) (৯৯/১০)

গ) (৮২/৯)

\textbf{\hl{ঘ) (১০০/১১) ✔}}

\textbf{\hl{✔ সঠিক উত্তর: ঘ) (১০০/১১)}}

\begin{quote}
\textbf{ব্যাখ্যা:} \emph{স্ত্রীর বয়স ১০০ বছর হলে ব্যক্তির বয়স ১১০ বছর।}

\emph{স্ত্রী ছোট = ১১০ − ১০০ = ১০ বছর}

\emph{∴ ১১০ এ ছোট ১০ বছর}

\emph{⇒ ১ এ ছোট (১০/১১০) বছর}

\emph{∴ ১০০ এ ছোট (১০×১০০/১১০) = (১০০/১১) বছর}
\end{quote}

\textbf{কোন পরীক্ষায় মোট পরীক্ষার্থীদের ৪০\% ছাত্রী। পরীক্ষায় ছাত্রীদের পাশের
হার ৪০\% এবং ছাত্রদের পাশের হার ৬০\% হলে মোট পাশের হার কত?}

ক) ৩৬\%

খ) ৪২\%

গ) ৪৮\%

\textbf{\hl{ঘ) ৫২\% ✔}}

\textbf{\hl{✔ সঠিক উত্তর: ঘ) ৫২\%}}

\begin{quote}
\textbf{ব্যাখ্যা:} \emph{ধরি, মোট পরীক্ষার্থী ১০০ জন}

\emph{ছাত্রী = ৪০ জন এবং ছাত্র = (১০০ − ৪০) = ৬০ জন}

\emph{ছাত্রীদের মধ্যে পাশ = ৪০ এর ৪০\% = ১৬ জন}

\emph{ছাত্রদের মধ্যে পাশ = ৬০ এর ৬০\% = ৩৬ জন}

\emph{∴ মোট পাশ = ১৬ + ৩৬ = ৫২ জন (বা ৫২\%)}
\end{quote}

\textbf{একটি কম্পিউটার বিজ্ঞান পরীক্ষায় ৩০\% পরীক্ষার্থী পাস করেছে।যারা পাস
করতে পারে নি। তাদের ১২ জন কম্পিউটার বিজ্ঞান কোর্সে অংশগ্রহণ করেছে এবং ৩০ জন
উচ্চ কোর্সে অংশগ্রহণ করেনি। কয়জন পরীক্ষার্থী পরীক্ষায় অংশগ্রহন করেছে?}

\textbf{\hl{ক) ৬০ ✔}}

খ) ৮০

গ) ১০০

ঘ) ১২০

\textbf{\hl{✔ সঠিক উত্তর: ক) ৬০}}

\begin{quote}
\textbf{ব্যাখ্যা:} \emph{মোট ফেল করেছে = ৩০ + ১২ = ৪২ জন}

\emph{শতাংশে ফেল করেছে = (১০০ − ৩০)\% = ৭০\%}

\emph{প্রশ্নমতে, ৭০\% = ৪২}

\emph{⇒ ১\% = (৪২/৭০)}

\emph{∴ ১০০\% = (৪২×১০০/৭০) = ৬০ জন}
\end{quote}

\textbf{একটি ১০০০০ টাকার বিলের উপর এককালীন ৪০\% কমতি এবং পরপর ৩৬\% ও ৪\%
কমতির মধ্যে পার্থক্য কত?}

ক) ০

\textbf{\hl{খ) ১৪৪ ✔}}

গ) ২৫৬

ঘ) ৪০০

\textbf{\hl{✔ সঠিক উত্তর: খ) ১৪৪}}

\begin{quote}
\textbf{ব্যাখ্যা:} \emph{এককালীন ৪০\% কমতি = ১০০০০ এর ৪০\% = ৪০০০ টাকা}

\emph{আবার, ৩৬\% কমতির পর বিল থাকে = ১০০০০ − ৩৬০০ = ৬৪০০ টাকা}

\emph{পরবর্তী ৪\% কমতি = ৬৪০০ এর ৪\% = ২৫৬ টাকা}

\emph{পরপর মোট কমতি = ৩৬০০ + ২৫৬ = ৩৮৫৬ টাকা}

\emph{∴ পার্থক্য = ৪০০০ − ৩৮৫৬ = ১৪৪ টাকা}
\end{quote}

\textbf{৮ জন লোক একটি কাজ ১২ দিনে করতে পার। ২জন লোক কমিয়ে দিলে কাজটি
সমাধা করতে শতকরা কতদিন বেশি লাগবে?}

ক) ২৫\%

খ) ৫০\%

\textbf{\hl{গ) (১০০/৩) \% ✔}}

ঘ) (২০০/৩) \%

\textbf{\hl{✔ সঠিক উত্তর: গ) (১০০/৩) \%}}

\begin{quote}
\textbf{ব্যাখ্যা:} \emph{লোকসংখ্যা = ৮ − ২ = ৬ জন}

\emph{৮ জনের লাগে ১২ দিন}

\emph{⇒ ১ জনের লাগে (১২ × ৮) দিন}

\emph{∴ ৬ জনের লাগে (১২×৮/৬) = ১৬ দিন}

\emph{দিন বেশি লাগে = ১৬ − ১২ = ৪ দিন}

\emph{∴ শতকরা বেশি লাগে = (৪×১০০/১২) \% = (১০০/৩) \%}
\end{quote}

\textbf{৬৪ কিলোগ্রাম বালি ও পাথরের টুকরার মিশ্রণে বালির পরিমান ২৫\%। কত
কিলোগ্রাম বালি মেশালে নতুন মিশ্রণে পাথরের টুকরার পরিমাণ ৪০\% হবে?}

ক) ৯.৬ কিলোগ্রাম

খ) ৪৮ কিলোগ্রাম

গ) ১১ কিলোগ্রাম

\textbf{\hl{ঘ) ৫৬ কিলোগ্রাম ✔}}

\textbf{\hl{✔ সঠিক উত্তর: ঘ) ৫৬ কিলোগ্রাম}}

\begin{quote}
\textbf{ব্যাখ্যা:} \emph{বালির পরিমাণ = ৬৪ এর ২৫\% = ১৬ কেজি}

\emph{পাথরের পরিমাণ = ৬৪ − ১৬ = ৪৮ কেজি}

\emph{নতুন মিশ্রণে পাথরের পরিমাণ ৪০\% হলে বালির পরিমাণ = ৬০\%}

\emph{প্রশ্নমতে, ৪০\% = ৪৮ কেজি}

\emph{⇒ ১\% = (৪৮/৪০) কেজি}

\emph{∴ ৬০\% = (৪৮×৬০/৪০) = ৭২ কেজি}

\emph{∴ বালি মেশাতে হবে = ৭২ − ১৬ = ৫৬ কেজি}
\end{quote}

\textbf{একটি ভোট কেন্দ্রে উপস্থিত ভোটারদের ৫৫\% ভোট পেয়ে একজন প্রার্থী
নির্বাচিত হয়েছেন। তিনি তার একমাত্র প্রতিদ্বন্দ্বী অপেক্ষা ১০০০০ ভোট বেশি
পেয়েছেন। ভোট কেন্দ্রে কত জন ভোটার উপস্থিত ছিল?}

ক) ১১০০০০

খ) ১১১০০০

\textbf{\hl{গ) ১০ ১০০০০ ✔}}

ঘ) ১০১০০০

\textbf{\hl{✔ সঠিক উত্তর: গ) ১০ ১০০০০}}

\begin{quote}
\textbf{ব্যাখ্যা:} \emph{বিজয়ী প্রার্থী পায় = ৫৫\%}

\emph{অনুগামী প্রতিদ্বন্দ্বী পায় = (১০০ − ৫৫)\% = ৪৫\%}

\emph{ভোটের পার্থক্যের হার = ৫৫\% − ৪৫\% = ১০\%}

\emph{প্রশ্নমতে, ১০\% = ১০০০০}

\emph{∴ ১০০\% = (১০০০০×১০০/১০) = ১০০০০০ জন}
\end{quote}

\textbf{আশেক আলী মাতবর তার সম্পদের ১২\% স্ত্রীকে, ২০\% ছেলেকে এবং অবশিষ্ট
৮১৬০০ টাকা মেয়েকে দিলেন। তার সম্পদের মোট মূল্য কত ছিল?}

ক) ১২২০০০ টাকা

খ) ১০২২০০ টাকা

\textbf{\hl{গ) ১২০০০০ টাকা ✔}}

ঘ) ১১২২০০০ টাকা

\textbf{\hl{✔ সঠিক উত্তর: গ) ১২০০০০ টাকা}}

\begin{quote}
\textbf{ব্যাখ্যা:} \emph{স্ত্রী ও ছেলেকে মোট দিলেন = ১২\% + ২০\% = ৩২\%}

\emph{অবশিষ্ট সম্পদ = (১০০ − ৩২)\% = ৬৮\%}

\emph{প্রশ্নমতে, ৬৮\% = ৮১৬০০ টাকা}

\emph{⇒ ১\% = (৮১৬০০/৬৮) টাকা}

\emph{∴ ১০০\% = (৮১৬০০×১০০/৬৮) = ১২০০০০ টাকা}
\end{quote}

\textbf{একখানা গাড়ির মূল্য ১৫০০ টাকা ও একটি ঘোড়ার মূল্য ২০০০ টাকা। যদি গাড়ির
মূল্য শতকরা ৫ টাকা ও ঘোড়ার মূল্য শতকরা ৮ টাকা বৃদ্ধি পায়, তবে গাড়ি ও ঘোড়ার মূল্য
একত্রে কত হবে?}

ক) ৩৫৩৭

\textbf{\hl{খ) ৩৭৩৫ ✔}}

গ) ৩৫৯৭

ঘ) ৩৭৭৫

\textbf{\hl{✔ সঠিক উত্তর: খ) ৩৭৩৫}}

\begin{quote}
\textbf{ব্যাখ্যা:} \emph{গাড়ির বর্ধিত মূল্য = ১৫০০ + (১৫০০ এর ৫\%) = ১৫০০ +
৭৫ = ১৫৭৫ টাকা}

\emph{ঘোড়ার বর্ধিত মূল্য = ২০০০ + (২০২০ এর ৮\%) = ২০০০ + ১৬০ = ২১৬০ টাকা}

\emph{∴ একত্রে মূল্য = ১৫৭৫ + ২১৬০ = ৩৭৩৫ টাকা}
\end{quote}

\textbf{একটি সংখ্যার ৩০\% যদি ১৩৫ হয়, তবে সংখ্যাটির ১৫০\% কত হবে?}

ক) ৬০০

\textbf{\hl{খ) ৬৭৫ ✔}}

গ) ৭৫০

ঘ) ৮৯০

\textbf{\hl{✔ সঠিক উত্তর: খ) ৬৭৫}}

\begin{quote}
\textbf{ব্যাখ্যা:} \emph{৩০\% = ১৩৫}

\emph{⇒ ১\% = (১৩৫/৩০) = ৪.৫}

\emph{∴ ১৫০\% = ৪.৫ × ১৫০ = ৬৭৫}
\end{quote}

\textbf{একটি রাশি অপর রাশির ৬৪\% হলে, রাশি দুটির অনুপাত কত?}

ক) ২০ : ২৫

\textbf{\hl{খ) ১৬ : ২৫ ✔}}

গ) ১৫ : ২০

ঘ) ২৫ : ৩০

\textbf{\hl{✔ সঠিক উত্তর: খ) ১৬ : ২৫}}

\begin{quote}
\textbf{ব্যাখ্যা:} \emph{ধরি, অপর রাশিটি ১০০}

\emph{∴ একটি রাশি = ৬৪}

\emph{রাশি দুটির অনুপাত = ৬৪ : ১০০ = ১৬ : ২৫}
\end{quote}

\textbf{আলমের বয়স কমলের বয়সের ৮০\% হলে কমলের বয়স আলমের বয়সের কত শতাংশ?}

ক) ১০০\%

খ) ১০৫\%

গ) ১২০\%

\textbf{\hl{ঘ) ১২৫\% ✔}}

\textbf{\hl{✔ সঠিক উত্তর: ঘ) ১২৫\%}}

\begin{quote}
\textbf{ব্যাখ্যা:} \emph{ধরি, কমলের বয়স ১০০ বছর}

\emph{∴ আলমের বয়স = ১০০ এর ৮০\% = ৮০ বছর}

\emph{কমলের বয়স আলমের বয়সের = ( (১০০/৮০) × ১০০)\% = ১২৫\%}
\end{quote}

\textbf{5\% বৃদ্ধিতে x-এর বর্ধিত মান কত?}

ক) x + (x/৫০০)

খ) x − (৫/১০০)

\textbf{\hl{গ) x (১+(১/২০)) ✔}}

ঘ) ঘ. (x/২০)

\textbf{\hl{✔ সঠিক উত্তর: গ) x (১+(১/২০))}}

\begin{quote}
\textbf{ব্যাখ্যা:} \emph{5\% বৃদ্ধিতে x-এর বর্ধিত মান = x + (x এর ৫\%) = x +
(x × (৫/১০০) ) = x + (x/২০) = x (১+(১/২০))}
\end{quote}

\textbf{y এর x\% যদি 10 হয় তবে y এর মান কত?}

ক) (১০০/x)

\textbf{\hl{খ) (১০০০/x) ✔}}

গ) ১০x

ঘ) ১০০x

\textbf{\hl{✔ সঠিক উত্তর: খ) (১০০০/x)}}

\begin{quote}
\textbf{ব্যাখ্যা:} \emph{y এর x\% = ১০}

\emph{⇒ y × (x/১০০) = ১০}

\emph{⇒ xy = ১০১০}

\emph{∴ y = (১০০০/x)}
\end{quote}

\textbf{১৫ এর কত শতাংশে ১ এর ১৫ শতাংশের সমান?}

ক) .০০১

খ) .১

গ) .০১

\textbf{\hl{ঘ) ১ ✔}}

\textbf{\hl{✔ সঠিক উত্তর: ঘ) ১}}

\begin{quote}
\textbf{ব্যাখ্যা:} \emph{ধরি, ১৫ এর k\% = ১ এর ১৫\%}

\emph{বা, ১৫ × (k/১০০) = ১ × (১৫/১০০)}

\emph{বা, k = (১৫ × ১০০/১০০ × ১৫)}

\emph{∴ k = ১}
\end{quote}

\textbf{কোনো সংখ্যার ৪০\% এর সাথে ৪৫ যোগ করলে যোগফল যদি ঐ সংখ্যাটিই হয়,
তাহলে সংখ্যাটি কত?}

ক) ৭০

খ) ৮৫

গ) ৬৪

\textbf{\hl{ঘ) ৭৫ ✔}}

\textbf{\hl{✔ সঠিক উত্তর: ঘ) ৭৫}}

\begin{quote}
\textbf{ব্যাখ্যা:} \emph{ধরি, সংখ্যাটি x}

\emph{প্রশ্নমতে, x এর ৪০\% + ৪৫ = x}

\emph{⇒ (৪০x/১০০) + ৪৫ = x}

\emph{⇒ (২x/৫) + ৪৫ = x}

\emph{⇒ ২x + ২২৫ = ৫x}

\emph{⇒ ৫x − ২x = ২২৫}

\emph{⇒ ৩x = ২২৫}

\emph{⇒ x = (২২৫/৩)}

\emph{∴ x = ৭৫}
\end{quote}

\textbf{কোন সংখ্যার ৬\% বেড়ে ৩৭১ হয়?}

ক) ২৮০

খ) ৩১৯

\textbf{\hl{গ) ৩৫০ ✔}}

ঘ) ৩৬১

\textbf{\hl{✔ সঠিক উত্তর: গ) ৩৫০}}

\begin{quote}
\textbf{ব্যাখ্যা:} \emph{ধরি, সংখ্যাটি x}

\emph{প্রশ্নমতে, x + x এর ৬\% = ৩৭১}

\emph{⇒ x + (৬x/১০০) = ৩৭১}

\emph{⇒ (১০৬x/১০০) = ৩৭১}

\emph{⇒ x = (৩৭১ × ১০০/১০৬)}

\emph{∴ x = ৩৫০}
\end{quote}

\textbf{একটি সংখ্যা থেকে সংখ্যাটির ৪০\% বিয়োগ করলে ৩০ থাকে। সংখ্যাটি কত?}

\textbf{\hl{ক) ৫০ ✔}}

খ) ৫৬

গ) ৬০

ঘ) কোনোটিই নয়

\textbf{\hl{✔ সঠিক উত্তর: ক) ৫০}}

\begin{quote}
\textbf{ব্যাখ্যা:} \emph{ধরি, সংখ্যাটি x}

\emph{প্রশ্নমতে, x − x এর ৪০\% = ৩০}

\emph{⇒ x − (৪০x/১০০) = ৩০}

\emph{⇒ (৬০x/১০০) = ৩০}

\emph{⇒ x = (৩০ × ১০০/৬০)}

\emph{∴ x = ৫০}
\end{quote}

\textbf{২০টি কমলার ২০\% পচা হলে, ভাল কমলার সংখ্যা -}

ক) ১০টি

খ) ৪টি

গ) ৫টি

\textbf{\hl{ঘ) ১৬টি ✔}}

\textbf{\hl{✔ সঠিক উত্তর: ঘ) ১৬টি}}

\begin{quote}
\textbf{ব্যাখ্যা:} \emph{পচা কমলার সংখ্যা = (২০ এর ২০\%) টি = ২০ × (২০/১০০)
টি = ৪ টি}

\emph{∴ ভাল কমলার সংখ্যা = (২০ − ৪) টি = ১৬ টি।}
\end{quote}

\textbf{কোন সংখ্যার ৪০ শতাংশ ও ১০ শতাংশের পার্থক্য ৩.৬?}

ক) ৯

খ) ৮

\textbf{\hl{গ) ১২ ✔}}

ঘ) ১৫

\textbf{\hl{✔ সঠিক উত্তর: গ) ১২}}

\begin{quote}
\textbf{ব্যাখ্যা:} \emph{ধরি, সংখ্যাটি x}

\emph{প্রশ্নমতে, x এর ৪০\% − x এর ১০\% = ৩.৬}

\emph{⇒ (৪০x/১০০) − (১০x/১০০) = ৩.৬}

\emph{⇒ (৩০x/১০০) = ৩.৬}

\emph{⇒ ৩০x = ৩৬০}

\emph{∴ x = (৩৬০/৩০) = ১২}
\end{quote}

\textbf{৩,০০০ এর শতকরা ৫ ভাগ অপেক্ষা ৩,০০০ এর শতকরা ১০ ভাগ কত বেশি?}

ক) ১৭৫

খ) ১৬০

\textbf{\hl{গ) ১৫০ ✔}}

ঘ) ১০০

\textbf{\hl{✔ সঠিক উত্তর: গ) ১৫০}}

\begin{quote}
\textbf{ব্যাখ্যা:} \emph{৩,০০০ এর ৫\% = ৩,০০০ × (৫/১০০) = ১৫০}

\emph{আবার ৩,০০০ এর ১০\% = ৩,০০০ × (১০/১০০) = ৩০০}

\emph{∴ বেশি = ৩০০ − ১৫০ = ১৫০}
\end{quote}

\textbf{একটি স্কুলে ৪৫০ জন শিক্ষার্থী আছে এবং তার মধ্যে ১৮\% ছাত্রী। ঐ স্কুলে
কতজন ছাত্রী আছে?}

\textbf{\hl{ক) ৮১ জন ✔}}

খ) ৮২ জন

গ) ৮৩ জন

ঘ) ৮৪ জন

\textbf{\hl{✔ সঠিক উত্তর: ক) ৮১ জন}}

\begin{quote}
\textbf{ব্যাখ্যা:} \emph{ছাত্রীর সংখ্যা = ৪৫০ এর ১৮\%}

\emph{= ৪৫০ × (১৮/১০০)}

\emph{= ৮১ জন}
\end{quote}

\textbf{একটি ক্লাসে ৬৪০ জন ছাত্র-ছাত্রী আছে যার মধ্যে ৪০\% ছাত্র। সেই ক্লাসে
ছাত্রীর সংখ্যা কত?}

ক) ২৫৬

\textbf{\hl{খ) ৩৮৪ ✔}}

গ) ২৫২

ঘ) ২৬৫

\textbf{\hl{✔ সঠিক উত্তর: খ) ৩৮৪}}

\begin{quote}
\textbf{ব্যাখ্যা:} \emph{মোট ছাত্র-ছাত্রী = ৬৪০ জন}

\emph{ছাত্র সংখ্যা = ৪০\%}

\emph{∴ ছাত্রী সংখ্যা = (১০০ − ৪০)\% = ৬০\%}

\emph{নির্ণেয় ছাত্রীর সংখ্যা = ৬৪০ এর ৬০\%}

\emph{= ৬৪০ × (৬০/১০০)}

\emph{= ৩৮৪ জন}
\end{quote}

\textbf{একটি স্কুলে মোট ৫০০ জন শিক্ষার্থীর মধ্যে ২০\% ছাত্রী। কোন এক বুধবারে ৪০
জন ছাত্র অনুপস্থিত ছিল। ঐদিন শতকরা কত জন ছাত্র উপস্থিত ছিল?}

\textbf{\hl{ক) ৯০\% ✔}}

খ) ৮০\%

গ) ৭৫\%

ঘ) ৫০\%

\textbf{\hl{✔ সঠিক উত্তর: ক) ৯০\%}}

\begin{quote}
\textbf{ব্যাখ্যা:} \emph{ছাত্রীর সংখ্যা = ৫০০ এর ২০\% = (৫০০ × ২০/১০০) জন =
১০০ জন}

\emph{∴ মোট ছাত্রের সংখ্যা = (৫০০ − ১০০) = ৪০০ জন}

\emph{বুধবারে ছাত্র উপস্থিত ছিল = (৪০০ − ৪০) = ৩৬০ জন}

\emph{∴ শতকরা ছাত্র উপস্থিত ছিল = ( (৩৬০/৪০০) × ১০০)\% = ৯০\%}
\end{quote}

\textbf{একটি স্কুলে মোট ৫৪০ জন ছাত্রের মধ্যে কোন এক বুধবারে ১৭ জন অনুপস্থিত
ছিল। ঐদিন শতকরা কতজন ছাত্র উপস্থিত ছিল?}

ক) ৯৫\%

খ) ৫\%

গ) ৯৩.৯২\%

\textbf{\hl{ঘ) ৯৬.৮৫\% ✔}}

\textbf{\hl{✔ সঠিক উত্তর: ঘ) ৯৬.৮৫\%}}

\begin{quote}
\textbf{ব্যাখ্যা:} \emph{ঐদিন উপস্থিত ছিল = (৫৪০ − ১৭) = ৫২৩ জন ছাত্র}

\emph{∴ শতকরা উপস্থিত ছিল = ( (৫২৩ × ১০০/৫৪০) )\% = ৯৬.৮৫\%}
\end{quote}

\textbf{যদি 12 সদস্যবিশিষ্ট কোনো কমিটির সদস্যের মধ্যে 9 জন মহিলা হয়, তবে
সদস্যদের মধ্যে শতকরা কত ভাগ পুরুষ?}

\textbf{\hl{ক) 25\% ✔}}

খ) 40\%

গ) 15\%

ঘ) 20\%

\textbf{\hl{✔ সঠিক উত্তর: ক) 25\%}}

\begin{quote}
\textbf{ব্যাখ্যা:} \emph{পুরুষ সদস্য সংখ্যা = ১২ − ৯ = ৩ জন}

\emph{∴ শতকরা পুরুষ সদস্য = ( (৩/১২) × ১০০)\% = ২৫\%}
\end{quote}

\textbf{কোন কারখানায় কর্মীদের জনপ্রতি দৈনিক মজুরী ছিল ১০০ টাকা। মন্দার কারণে
তাদের দৈনিক মজুরী ৫০\% কমানো হয়েছিল। সম্প্রতি দৈনিক মজুরী ৬০\% বাড়ানো হয়েছে।
বর্তমানে জনপ্রতি দৈনিক মজুরী কত টাকা।}

ক) ১৬০

খ) ১১০

\textbf{\hl{গ) ৮০ ✔}}

ঘ) ৬০

\textbf{\hl{✔ সঠিক উত্তর: গ) ৮০}}

\begin{quote}
\textbf{ব্যাখ্যা:} \emph{৫০\% কমে দৈনিক মজুরী হয় = (১০০ − ১০০ এর ৫০\%) = ৫০
টাকা}

\emph{পরবর্তীতে ৫০ টাকার উপর ৬০\% বাড়ানো হলে বৃদ্ধি পায় = ৫০ এর ৬০\% = ৩০
টাকা}

\emph{∴ বর্তমানে জনপ্রতি দৈনিক মজুরী = ৫০ + ৩০ = ৮০ টাকা}
\end{quote}

\textbf{একটি বাক্সে যতগুলো মার্বেল আছে তার ৩০\% কালো এবং বাকিগুলো লাল। যদি ঐ
বাক্সে ২১০ টি লাল মার্বেল থাকে। তবে ঐ বাক্সে মোট কতটি মার্বেল আছে?}

ক) ২৮০

\textbf{\hl{খ) ৩০০ ✔}}

গ) ৩৫০

ঘ) কোনোটিই নয়

\textbf{\hl{✔ সঠিক উত্তর: খ) ৩০০}}

\begin{quote}
\textbf{ব্যাখ্যা:} \emph{বাক্সে লাল মার্বেল = (১০০ − ৩০)\% = ৭০\%}

\emph{প্রশ্নমতে, ৭০\% = ২১০ টি মার্বেল}

\emph{∴ ১০০\% = (২১০ × ১০০/৭০) = ৩০০ টি মার্বেল}
\end{quote}

\textbf{৫\% ভ্যাটে ৪০৮০ টাকার বিলে ভ্যাটের পরিমাণ কত টাকা?}

ক) ৫ টাকা

খ) ৯৫ টাকা

\textbf{\hl{গ) ২০৪ টাকা ✔}}

ঘ) ৪০৮ টাকা

\textbf{\hl{✔ সঠিক উত্তর: গ) ২০৪ টাকা}}

\begin{quote}
\textbf{ব্যাখ্যা:} \emph{৪০৮০ টাকার ৫\% = ৪০৮০ × (৫/১০০) টাকা = ২০৪ টাকা}

\emph{∴ ভ্যাটের পরিমাণ ২০৪ টাকা।}
\end{quote}

\textbf{রফিক কেয়াকে তার মোট টাকার ৩০\% দিয়ে দেখে যে তার নিকট আরও ৩৫০
টাকা আছে। সে কেয়াকে কত টাকা দিল?}

ক) 350

খ) 200

\textbf{\hl{গ) 150 ✔}}

ঘ) 120

\textbf{\hl{✔ সঠিক উত্তর: গ) 150}}

\begin{quote}
\textbf{ব্যাখ্যা:} \emph{কেয়াকে দেয়ার পর রফিকের টাকার পরিমাণ = (১০০ −
৩০)\% = ৭০\%}

\emph{প্রশ্নমতে, ৭০\% = ৩৫০ টাকা}

\emph{∴ ৩০\% = (৩৫০ × ৩০/৭০) = ১৫০ টাকা}
\end{quote}

\textbf{ধানে চাল ও তুষের অনুপাত ৭ ও ৩ হলে, এতে শতকরা কী পরিমাণ চাল আছে?}

ক) ৮৫\%

খ) ৯৫\%

গ) ৬৫\%

\textbf{\hl{ঘ) ৭০\% ✔}}

\textbf{\hl{✔ সঠিক উত্তর: ঘ) ৭০\%}}

\begin{quote}
\textbf{ব্যাখ্যা:} \emph{ধানে চাল ও তুষের অনুপাত = ৭ : ৩}

\emph{অনুপাতের যোগফল = (৭ + ৩) = ১০}

\emph{∴ শতকরা চাল আছে = (৭ × ১০০/১০) \% = ৭০\%}
\end{quote}

\textbf{এক গ্লাস গুড়ের শরবতে গুড় ও পানির অনুপাত ৪ : ৬ হলে গুড়ের পরিমাণ কত?}

ক) ১০\%

খ) ২০\%

গ) ৩০\%

\textbf{\hl{ঘ) ৪০\% ✔}}

\textbf{\hl{✔ সঠিক উত্তর: ঘ) ৪০\%}}

\begin{quote}
\textbf{ব্যাখ্যা:} \emph{শরবতে গুড় ও পানির অনুপাত = ৪ : ৬}

\emph{অনুপাতের যোগফল = (৪ + ৬) = ১০}

\emph{∴ গুড়ের পরিমাণ = (৪ × ১০০/১০) \% = ৪০\%}
\end{quote}

\textbf{একটি সংখ্যা অপর একটি সংখ্যার ৪৫০\%। সংখ্যা দুইটির অনুপাত কত?}

\textbf{\hl{ক) ৯ : ২ ✔}}

খ) ৪৫ : ১

গ) ৪৫০ : ১

ঘ) ১ : ১

ঙ) কোনোটিই নয়

\textbf{\hl{✔ সঠিক উত্তর: ক) ৯ : ২}}

\begin{quote}
\textbf{ব্যাখ্যা:} \emph{ধরি, অপর সংখ্যাটি = ১০০}

\emph{∴ একটি সংখ্যা = ৪৫০}

\emph{সংখ্যা দুইটির অনুপাত = ৪৫০ : ১০০ = ৯ : ২}
\end{quote}

\textbf{একটি বইয়ের প্রকৃত মূল্যের শতকরা ৭৫ মূল্যে ৫০ টাকায় বিক্রয় করা হলে।
বইটির প্রকৃত মূল্য কত?}

ক) ৮০ টাকা

খ) ৪০ টাকা

গ) ৬০ টাকা

\textbf{\hl{ঘ) ৬৬.৬৭ টাকা ✔}}

\textbf{\hl{✔ সঠিক উত্তর: ঘ) ৬৬.৬৭ টাকা}}

\begin{quote}
\textbf{ব্যাখ্যা:} \emph{৭৫\% মূল্য = ৫০ টাকা}

\emph{∴ ১০০\% বা প্রকৃত মূল্য = (৫০ × ১০০/৭৫) = ৬৬.৬৭ টাকা}
\end{quote}

\textbf{আখের নমুনায় ২০\% চিনি রয়েছে। ৫০ কেজি চিনি উৎপাদনের জন্য কত কেজি আখ
প্রয়োজন?}

ক) ১০০০ কেজি

\textbf{\hl{খ) ২৫০ কেজি ✔}}

গ) ৫০০ কেজি

ঘ) ৭৫০ কেজি

\textbf{\hl{✔ সঠিক উত্তর: খ) ২৫০ কেজি}}

\begin{quote}
\textbf{ব্যাখ্যা:} \emph{২০ কেজি চিনির জন্য আখ প্রয়োজন ১০০ কেজি}

\emph{∴ ৫০ কেজি চিনির জন্য আখ প্রয়োজন = (১০০ × ৫০/২০) = ২৫০ কেজি}
\end{quote}

\textbf{২৪ লিটার দুধ ও পানির মিশ্রণে ২০\% পানি আছে। ঐ মিশ্রণে কতটুকু পানি
মেশালে মিশ্রণে ৪০\% পানি হবে?}

\textbf{\hl{ক) ৮ লিটার ✔}}

খ) ৪ লিটার

গ) ৬ লিটার

ঘ) ১০ লিটার

\textbf{\hl{✔ সঠিক উত্তর: ক) ৮ লিটার}}

\begin{quote}
\textbf{ব্যাখ্যা:} \emph{মিশ্রণে পানির পরিমাণ = ২৪ এর ২০\% = ৪.৮ লিটার}

\emph{∴ দুধের পরিমাণ = ২৪ − ৪.৮ = ১৯.২ লিটার}

\emph{ধরি, x লিটার পানি মেশাতে হবে।}

\emph{নতুন মিশ্রণে পানির পরিমাণ = (৪.৮ + x) লিটার}

\emph{নতুন মিশ্রণে পানি ৪০\% হলে, দুধ থাকবে (১০০ − ৪০)\% = ৬০\%}

\emph{প্রশ্নমতে, (৪.৮ + x/১৯.২) = (৪০/৬০)}

\emph{⇒ (৪.৮ + x/১৯.২) = (২/৩)}

\emph{⇒ ৩(৪.৮ + x) = ৩৮.৪}

\emph{⇒ ১৪.৪ + ৩x = ৩৮.৪}

\emph{⇒ ৩x = ৩৮.৪ − ১৪.৪ = ২৪}

\emph{∴ x = ৮ লিটার}
\end{quote}

\textbf{৪\% লবণ রয়েছে এমন ৬ লিটার সমুদ্রজল থেকে ১ লিটার পানি বাষ্পীভূত করলে
তার লবণাক্ততা কত শতাংশ হবে?}

ক) ৪

খ) ৩.৬

\textbf{\hl{গ) ৪.৮ ✔}}

ঘ) ৫.২

\textbf{\hl{✔ সঠিক উত্তর: গ) ৪.৮}}

\begin{quote}
\textbf{ব্যাখ্যা:} \emph{৬ লিটার সমুদ্রজলে লবণের পরিমাণ = ৬ × ৪\% = ০.২৪
লিটার}

\emph{১ লিটার পানি বাষ্পীভূত করলে অবশিষ্ট পানির পরিমাণ = (৬ − ১) = ৫ লিটার}

\emph{∴ বর্তমান লবণাক্ততার হার = ( (০.২৪ × ১০০/৫) )\% = ৪.৮\%}
\end{quote}

\textbf{একটি শেয়ারের মূল্য গতকাল ২৫\% বেড়ে গেল। আজকে আবার ২৫\% কমে গেল।
প্রকৃত/কমার হার কত?}

\textbf{\hl{ক) ৬.২৫ ✔}}

খ) ২০০

গ) ০.০২

ঘ) ০.২

\textbf{\hl{✔ সঠিক উত্তর: ক) ৬.২৫}}

\begin{quote}
\textbf{ব্যাখ্যা:} \emph{ধরি, শেয়ারের প্রাথমিক মূল্য ১০০ টাকা}

\emph{গতকাল ২৫\% বাড়ায় মূল্য হয় = ১২৫ টাকা}

\emph{আজকে ১২৫ টাকার ২৫\% কমায় কমে = ১২৫ × (২৫/১০০) = ৩১.২৫ টাকা}

\emph{∴ বর্তমান মূল্য = ১২৫ − ৩১.২৫ = ৯৮.৭৫ টাকা}

\emph{∴ মোট মূল্যের পরিবর্তন = ১০০ − ৯৮.৭৫ = ৬.২৫\% কমে গেল।}
\end{quote}

\textbf{একটি বইয়ের মূল্য ২৪.০০ টাকা। এই মূল্য প্রকৃত মূল্যের ৮০\%। বাকি মূল্য
সরকার ভর্তুকি দিয়ে থাকেন। সরকার প্রতি বইয়ে কত টাকা ভর্তুকি দেন?}

\textbf{\hl{ক) ৬.০০ ✔}}

খ) ১০.০০

গ) ১৫.০০

ঘ) ৩.০০

\textbf{\hl{✔ সঠিক উত্তর: ক) ৬.০০}}

\begin{quote}
\textbf{ব্যাখ্যা:} \emph{বইটির মূল্যের ৮০\% = ২৪ টাকা}

\emph{∴ বইটির ১০০\% বা প্রকৃত মূল্য = (২৪ × ১০০/৮০) = ৩০ টাকা}

\emph{∴ সরকার ভর্তুকি দেন = (৩০ − ২৪) টাকা = ৬.০০ টাকা।}
\end{quote}

\textbf{একজন বিক্রেতা ১৯,৬০০ টাকায় একটি টেলিভিশন বিক্রি করায় তাকে ১,৩২৩
টাকা বিক্রয় কর দিতে হল। তিনি শতকরা কত টাকা বিক্রয় কর দিলেন?}

ক) ৫.৫০\%

\textbf{\hl{খ) ৬.৭৫\% ✔}}

গ) ৭.৫০\%

ঘ) ৮.৫০\%

ঙ) কোনটিই নয়

\textbf{\hl{✔ সঠিক উত্তর: খ) ৬.৭৫\%}}

\begin{quote}
\textbf{ব্যাখ্যা:} \emph{১৯,৬০০ টাকায় বিক্রয় কর দিতে হয় = ১,৩২৩ টাকা}

\emph{∴ ১০০ টাকায় বিক্রয় কর দিতে হয় = (১,৩২৩ × ১০০/১৯,৬০০) টাকা = ৬.৭৫
টাকা।}
\end{quote}

\textbf{একটি বাক্সের দুই-তৃতীয়াংশ শার্ট পরীক্ষা করার পর ৪টি ত্রুটিযুক্ত এবং বাকি
৩৬টি ত্রুটিমুক্ত পাওয়া গেল। ৮৫\% শার্ট ত্রুটিমুক্ত পেতে হলে অবশিষ্ট শার্টের মধ্যে
কতটি ত্রুটিমুক্ত শার্ট পেতে হবে?}

\textbf{\hl{ক) ১৫ ✔}}

খ) ১৭

গ) ১৯

ঘ) ৫১

\textbf{\hl{✔ সঠিক উত্তর: ক) ১৫}}

\begin{quote}
\textbf{ব্যাখ্যা:} \emph{মোট শার্টের (২/৩) অংশ = (৪ + ৩৬) = ৪০টি}

\emph{∴ মোট শার্টের সংখ্যা = (৪০ × ৩/২) = ৬০টি}

\emph{৮৫\% ত্রুটিমুক্ত শার্ট হতে হলে মোট ত্রুটিমুক্ত শার্ট লাগবে = ৬০ এর ৮৫\% = ৬০
× (৮৫/১০০) = ৫১টি}

\emph{ইতোমধ্যে ত্রুটিমুক্ত পাওয়া গেছে = ৩৬টি}

\emph{∴ অবশিষ্ট শার্টের মধ্যে ত্রুটিমুক্ত শার্ট পেতে হবে = (৫১ − ৩৬) = ১৫টি}
\end{quote}

\textbf{২০ জনে যে সময়ে ১টি কাজ করতে পারে, তার ২০ শতাংশ কম সময়ে কাজটি শেষ
করতে হলে জনবল কত শতাংশ বাড়াতে হবে?}

ক) ৪০ শতাংশ

খ) ৩৩ শতাংশ

\textbf{\hl{গ) ২৫ শতাংশ ✔}}

ঘ) ২০ শতাংশ

\textbf{\hl{✔ সঠিক উত্তর: গ) ২৫ শতাংশ}}

\begin{quote}
\textbf{ব্যাখ্যা:} \emph{ধরি, ২০ জন কাজটি করে t দিনে।}

\emph{২০\% কম সময় = t − t এর ২০\% = (৮০t/১০০) দিন}

\emph{t দিনে কাজটি করে ২০ জন}

\emph{∴ ১ দিনে কাজটি করে (২০ × t) জন}

\emph{∴ (৮০t/১০০) দিনে কাজটি করতে লোক লাগবে = (২০ × t × ১০০/৮০t) = ২৫
জন}

\emph{∴ জনবল বাড়াতে হবে = ২৫ − ২০ = ৫ জন}

\emph{শতকরা বাড়াতে হবে = ( (৫/২০) × ১০০)\% = ২৫ শতাংশ}
\end{quote}

\textbf{২০ জনে যে সময়ে ১টি কাজ করতে পারে, তার এক-তৃতীয়াংশ কম সময়ে কাজটি
শেষ করতে হলে জনবল কত শতাংশ বাড়াতে হবে?}

ক) ৪০ শতাংশ

খ) ৩৩ শতাংশ

গ) ২৫ শতাংশ

\textbf{\hl{ঘ) ৫০ শতাংশ ✔}}

\textbf{\hl{✔ সঠিক উত্তর: ঘ) ৫০ শতাংশ}}

\begin{quote}
\textbf{ব্যাখ্যা:} \emph{ধরি, ২০ জন কাজটি করে x দিনে।}

\emph{এক-তৃতীয়াংশ কম সময় = x − (x/৩) = (২x/৩) দিন}

\emph{x দিনে কাজটি করে ২০ জন}

\emph{∴ ১ দিনে কাজটি করে (২০ × x) জন}

\emph{∴ (২x/৩) দিনে করতে মোট লোক লাগবে = (২০ × x × ৩/২x) = ৩০ জন}

\emph{∴ অতিরিক্ত লোক লাগবে = ৩০ − ২০ = ১০ জন}

\emph{শতকরা বাড়াতে হবে = ( (১০/২০) × ১০০)\% = ৫০ শতাংশ}
\end{quote}

\textbf{একটি স্কুলের ৪০\% ছাত্র। ছাত্রদের ৩০\% এবং ছাত্রীদের ২০\% বিতর্ক ক্লাবে
যোগ দিলে মোট ছাত্র-ছাত্রীর কত শতাংশে বিতর্ক ক্লাবে যোগ দিল?}

ক) ১২\%

\textbf{\hl{খ) ২৪\% ✔}}

গ) ২৬\%

ঘ) ৫০\%

\textbf{\hl{✔ সঠিক উত্তর: খ) ২৪\%}}

\begin{quote}
\textbf{ব্যাখ্যা:} \emph{ধরি, মোট ছাত্র-ছাত্রীর সংখ্যা ১০০ জন}

\emph{∴ ছাত্র সংখ্যা = ৪০ জন এবং ছাত্রী সংখ্যা = (১০০ − ৪০) = ৬০ জন}

\emph{বিতর্ক ক্লাবে যোগ দেওয়া ছাত্র = ৪০ এর ৩০\% = ১২ জন}

\emph{বিতর্ক ক্লাবে যোগ দেওয়া ছাত্রী = ৬০ এর ২০\% = ১২ জন}

\emph{∴ মোট যোগ দিল = (১২ + ১২) = ২৪ জন}

\emph{অর্থাৎ মোট ছাত্র-ছাত্রীর ২৪\% বিতর্ক ক্লাবে যোগ দিল।}
\end{quote}

\textbf{কোন স্কুলের ছাত্র-ছাত্রীর সংখ্যা ছিল ১,৮০০ জন। কিছুদিন পরে ৪\% ছাত্র
চলে গেল, আর ৫\% নতুন ছাত্রী ভর্তি হলো। এর ফলে ঐ স্কুলের মোট ছাত্র-ছাত্রীর সংখ্যা
২\% বেড়ে গেল। বর্তমানে ঐ স্কুলে মোট ছাত্র-ছাত্রীর সংখ্যা কত?}

ক) ১,৮৫০ জন

খ) ১,৮৭২ জন

গ) ১,৮১০ জন

\textbf{\hl{ঘ) ১,৮৩৬ জন ✔}}

\textbf{\hl{✔ সঠিক উত্তর: ঘ) ১,৮৩৬ জন}}

\begin{quote}
\textbf{ব্যাখ্যা:} \emph{মোট ছাত্র-ছাত্রী বাড়ে = ১,৮০০ এর ২\% = ১,৮০০ ×
(২/১০০) = ৩৬ জন}

\emph{∴ বর্তমানে ঐ স্কুলে মোট ছাত্র-ছাত্রীর সংখ্যা = (১,৮০০ + ৩৬) = ১,৮৩৬ জন।}
\end{quote}

\textbf{কেরামত সাহেব তার সম্পদের ১২\% স্ত্রীকে, ২০\% ছেলেকে ও অবশিষ্ট
৮,১৬,০০০ টাকা মেয়েকে দিলেন। মেয়ে সম্পত্তির শতকরা কত অংশ পাবে?}

ক) ৫০\%

খ) ৩০\%

\textbf{\hl{গ) ৬৮\% ✔}}

ঘ) ৮০\%

\textbf{\hl{✔ সঠিক উত্তর: গ) ৬৮\%}}

\begin{quote}
\textbf{ব্যাখ্যা:} \emph{স্ত্রী ও ছেলেকে মোট দেওয়া হলো = (১২\% + ২০\%) =
৩২\%}

\emph{∴ মেয়েকে দেওয়া হলো অবশিষ্ট = (১০০ − ৩২)\% = ৬৮\%}

\emph{অর্থাৎ মেয়ে সম্পত্তির ৬৮\% অংশ পাবে।}
\end{quote}

\textbf{p = a × b হলে এবং a ও b উভয়কে একত্রে ১০\% বৃদ্ধি করা হলে p এর মান
শতকরা কত বৃদ্ধি পাবে?}

ক) ১৫\%

খ) ২০\%

\textbf{\hl{গ) ২১\% ✔}}

ঘ) ২৫\%

\textbf{\hl{✔ সঠিক উত্তর: গ) ২১\%}}

\begin{quote}
\textbf{ব্যাখ্যা:} \emph{১০\% বৃদ্ধিতে a-এর পরিবর্তিত মান = a + a এর ১০\% = a
+ (১০a/১০০) = (১১a/১০)}

\emph{১০\% বৃদ্ধিতে b-এর পরিবর্তিত মান = b + b এর ১০\% = b + (১০b/১০০) =
(১১b/১০)}

\emph{∴ p-এর নতুন মান = (১১a/১০) × (১১b/১০) = (১২১ab/১০০) = (১২১p/১০০)}

\emph{∴ p-এর মান বৃদ্ধি = (১২১p/১০০) − p = (২১p/১০০)}

\emph{∴ শতকরা মান বৃদ্ধি পাবে = ২১\%}
\end{quote}

\textbf{যদি \textquotesingle ক\textquotesingle{} এর ১২০\% সংখ্যা
\textquotesingle খ\textquotesingle{} এর ৮০\% হয়, তাহলে (ক + খ) এর মান
কত?}

ক) ১.৫ক

খ) ২ক

\textbf{\hl{গ) ২.৫ক ✔}}

ঘ) ৩ক

\textbf{\hl{✔ সঠিক উত্তর: গ) ২.৫ক}}

\begin{quote}
\textbf{ব্যাখ্যা:} \emph{প্রশ্নমতে, \textquotesingle ক\textquotesingle{} এর
১২০\% = \textquotesingle খ\textquotesingle{} এর ৮০\%}

\emph{⇒ (১২০ক/১০০) = (৮০খ/১০০)}

\emph{⇒ ১২০ক = ৮০খ}

\emph{⇒ ৪খ = ৬ক}

\emph{∴ খ = (৬ক/৪) = ১.৫ক}

\emph{∴ ক + খ = ক + ১.৫ক = ২.৫ক}
\end{quote}

\textbf{কোনো এক স্থানের জনসংখ্যা বৃদ্ধির হার ৫\% এবং বর্তমান জনসংখ্যা ৩,২৫৫ জন
হলে এক বছর আগে ঐ স্থানের জনসংখ্যা কত ছিল?}

ক) ২,১০০ জন

খ) ৩,২০০ জন

\textbf{\hl{গ) ৩,১০০ জন ✔}}

ঘ) ২,৫০০ জন

\textbf{\hl{✔ সঠিক উত্তর: গ) ৩,১০০ জন}}

\begin{quote}
\textbf{ব্যাখ্যা:} \emph{জনসংখ্যা বৃদ্ধির হার ৫\% হওয়ায়,}

\emph{বর্তমান জনসংখ্যা ১০৫ জন হলে ১ বছর পূর্বে ছিল = ১০০ জন}

\emph{∴ বর্তমান জনসংখ্যা ৩,২৫৫ জন হলে ১ বছর পূর্বে ছিল = (১০০ × ৩,২৫৫/১০৫)
জন = ৩,১০০ জন}
\end{quote}

\textbf{একটি দেশের জনসংখ্যা ১৯৮০ হতে ২০১০ সালের মধ্যে প্রতি ১০ বছরে দ্বিগুণ
হলে এই সময়ে জনসংখ্যা বৃদ্ধির শতকরা হার কত?}

ক) ৩০০\%

\textbf{\hl{খ) ৭০০\% ✔}}

গ) ৮০০\%

ঘ) ১০০০\%

\textbf{\hl{✔ সঠিক উত্তর: খ) ৭০০\%}}

\begin{quote}
\textbf{ব্যাখ্যা:} \emph{১৯৮০ থেকে ২০১০ সাল পর্যন্ত সময় = ৩০ বছর (অর্থাৎ ১০
বছর পরপর মোট ৩ বার দ্বিগুণ হবে)}

\emph{ধরি, ১৯৮০ সালে জনসংখ্যা ছিল = ১০০ জন}

\emph{১৯৯০ সালে জনসংখ্যা (১ম দ্বিগুণ) = ২০০ জন}

\emph{২০০০ সালে জনসংখ্যা (২য় দ্বিগুণ) = ৪০০ জন}

\emph{২০১০ সালে জনসংখ্যা (৩য় দ্বিগুণ) = ৮০০ জন}

\emph{∴ জনসংখ্যা বৃদ্ধি পায় = (৮০০ − ১০০) = ৭০০ জন}

\emph{∴ বৃদ্ধির শতকরা হার = ৭০০\%}
\end{quote}

\textbf{প্রতি বছর কোনো শহরের লোকসংখ্যার ৭\% জন্মগ্রহণ করে এবং ৩\% মারা যায়।
এক বছরে শহরে ৪০০ জন লোক বৃদ্ধি পেলে ঐ শহরের মোট লোকসংখ্যা কত?}

ক) ৪০০০

খ) ৭০০০

\textbf{\hl{গ) ১০০০০ ✔}}

ঘ) ১৪০০০

\textbf{\hl{✔ সঠিক উত্তর: গ) ১০০০০}}

\begin{quote}
\textbf{ব্যাখ্যা:} \emph{প্রতি বছর লোক বৃদ্ধি পায় = (৭ − ৩)\% = ৪\%}

\emph{প্রশ্নমতে, ৪\% = ৪০০ জন}

\emph{∴ ১০০\% = (৪০০ × ১০০/৪) = ১০,০০০ জন}

\emph{∴ ঐ শহরের মোট লোকসংখ্যা ১০,০০০ জন।}
\end{quote}

\textbf{কোনো বছর একটি গ্রামের লোকসংখ্যা ১২\% বাড়ে। বছরের শেষে লোক সংখ্যা
৩৩৬০ হলে, বছরের শুরুতে লোকসংখ্যা কত ছিল?}

ক) ২৮০০

খ) ২৯০০

\textbf{\hl{গ) ৩০০০ ✔}}

ঘ) ৩০৫০

\textbf{\hl{✔ সঠিক উত্তর: গ) ৩০০০}}

\begin{quote}
\textbf{ব্যাখ্যা:} \emph{১২\% বৃদ্ধিতে,}

\emph{বছরের শেষে লোকসংখ্যা ১১২ জন হলে শুরুতে ছিল = ১০০ জন}

\emph{∴ বছরের শেষে লোকসংখ্যা ৩৩৬০ জন হলে শুরুতে ছিল = (১০০ × ৩৩৬০/১১২) =
৩০০০ জন।}
\end{quote}

\textbf{একটি মহল্লার জনসংখ্যার মধ্যে পুরুষের সংখ্যা ৫৫\%। ঐ মহল্লার পুরুষের সংখ্যা
১১০০ হলে ঐ মহল্লার মোট জনসংখ্যা কত?}

ক) ২১০০ জন

খ) ২৩০০ জন

গ) ২০০৫ জন

\textbf{\hl{ঘ) ২০০০ জন ✔}}

\textbf{\hl{✔ সঠিক উত্তর: ঘ) ২০০০ জন}}

\begin{quote}
\textbf{ব্যাখ্যা:} \emph{পুরুষের সংখ্যা ৫৫ জন হলে মোট জনসংখ্যা = ১০০ জন}

\emph{∴ পুরুষের সংখ্যা ১১০০ জন হলে মোট জনসংখ্যা = (১০০ × ১১০০/৫৫) = ২০০০
জন।}
\end{quote}

\textbf{২৪০ জন লোক একটি বনভোজনে যায়। সেখানে যতজন মহিলা ছিল তার থেকে ২০
জন পুরুষ বেশি ছিল। আবার যতজন শিশু ছিল তার থেকে ২০ জন প্রাপ্তবয়স্ক বেশি ছিল।
বনভোজনে কতজন পুরুষ ছিল?}

ক) ১০০

খ) ১৪০

গ) ১৪৫

\textbf{\hl{ঘ) ৭৫ ✔}}

\textbf{\hl{✔ সঠিক উত্তর: ঘ) ৭৫}}

\begin{quote}
\textbf{ব্যাখ্যা:} \emph{ধরি, মহিলার সংখ্যা = x জন}

\emph{∴ পুরুষের সংখ্যা = (x + ২০) জন}

\emph{∴ মোট প্রাপ্তবয়স্ক লোক = x + x + ২০ = (২x + ২০) জন}

\emph{আবার, শিশু অপেক্ষা ২০ জন প্রাপ্তবয়স্ক বেশি ছিল,}

\emph{∴ শিশুর সংখ্যা = (২x + ২০) − ২০ = ২x জন}

\emph{প্রশ্নমতে, (২x + ২০) + ২x = ২৪০}

\emph{⇒ ৪x + ২০ = ২৪০}

\emph{⇒ ৪x = ২২০}

\emph{⇒ x = ৫৫}

\emph{∴ পুরুষের সংখ্যা = ৫৫ + ২০ = ৭৫ জন।}
\end{quote}

\textbf{একটি ভোট কেন্দ্রে উপস্থিত ভোটারদের ৬০\% ভোট পেয়ে একজন প্রার্থী
নির্বাচিত হয়েছেন। তিনি একমাত্র প্রতিদ্বন্দ্বী অপেক্ষা ৭৫০০ ভোট বেশী পেয়েছেন।
ভোট কেন্দ্রে কত জন ভোটার উপস্থিত ছিল?}

ক) ২৫০০০

\textbf{\hl{খ) ৩৭৫০০ ✔}}

গ) ৪২০০০

ঘ) কোনটিই নয়

\textbf{\hl{✔ সঠিক উত্তর: খ) ৩৭৫০০}}

\begin{quote}
\textbf{ব্যাখ্যা:} \emph{ধরি, মোট উপস্থিত ভোটার = ১০০ জন}

\emph{বিজয়ী প্রার্থী পান = ৬০টি ভোট}

\emph{∴ প্রতিদ্বন্দ্বী প্রার্থী পান = (১০০ − ৬০) = ৪০টি ভোট}

\emph{দুই প্রার্থীর ভোটের পার্থক্য = (৬০ − ৪০) = ২০টি}

\emph{প্রশ্নমতে, ২০টি ভোটের পার্থক্য হলে ভোটার = ১০০ জন}

\emph{∴ ৭৫০০টি ভোটের পার্থক্য হলে মোট ভোটার = (১০০ × ৭৫০০/২০) = ৩৭,৫০০
জন।}
\end{quote}

\textbf{একটি গ্রামের মোট জনসংখ্যার (৫/৯) ভাগ লোক পুরুষ। যদি পুরুষের ৩০\%
বিবাহিত হয়, তাহলে মোট জনসংখ্যার শতকরা কতজন মহিলা অবিবাহিত?}

ক) ৭০\%

খ) ৪০\%

\textbf{\hl{গ) ২৭ (৭/৯) \% ✔}}

ঘ) ২০\%

\textbf{\hl{✔ সঠিক উত্তর: গ) ২৭ (৭/৯) \%}}

\begin{quote}
\textbf{ব্যাখ্যা:} \emph{ধরি, মোট জনসংখ্যা = ৯০ জন}

\emph{∴ পুরুষের সংখ্যা = ৯০ × (৫/৯) = ৫০ জন}

\emph{∴ মহিলার সংখ্যা = (৯০ − ৫০) = ৪০ জন}

\emph{বিবাহিত পুরুষের সংখ্যা = ৫০ এর ৩০\% = ১৫ জন}

\emph{যেহেতু প্রতিটি বিবাহিত পুরুষের ১ জন বিবাহিত স্ত্রী থাকে, তাই বিবাহিত
মহিলার সংখ্যাও = ১৫ জন}

\emph{∴ অবিবাহিত মহিলার সংখ্যা = (৪০ − ১৫) = ২৫ জন}

\emph{∴ মোট জনসংখ্যার শতকরা অবিবাহিত মহিলা = ( (২৫/৯০) × ১০০)\% = (২৫০/৯)
\% = ২৭ (৭/৯) \%}
\end{quote}

\textbf{করিম সাহেব মাসিক বেতন থেকে প্রভিডেন্ট ফান্ডের জন্য শতকরা ১০ ভাগ
কাটার পর তিনি ২,৭০০ টাকা পান। তার মাসিক বেতন কত?}

\textbf{\hl{ক) ৩,০০০ ✔}}

খ) ২,৯৭০

গ) ৩,০৭০

ঘ) ৩,১৭০

\textbf{\hl{✔ সঠিক উত্তর: ক) ৩,০০০}}

\begin{quote}
\textbf{ব্যাখ্যা:} \emph{১০\% কাটার পর অবশিষ্ট বেতন থাকে = (১০০ − ১০)\% =
৯০\%}

\emph{প্রশ্নমতে, ৯০\% = ২,৭০০ টাকা}

\emph{∴ ১০০\% বা মূল বেতন = (২,৭০০ × ১০০/৯০) = ৩,০০০ টাকা।}
\end{quote}

\textbf{আরিফের বেতন ৬০ টাকা বাড়ালে তার বেতন বাবুর বেতনের ৫০\% হবে। যদি
বাবুর বেতন ২০০০ টাকা হয় তবে আরিফের বর্তমান বেতন কত?}

ক) ৮৪০

খ) ৯০০

\textbf{\hl{গ) ৯৪০ ✔}}

ঘ) কোনোটিই নয়

\textbf{\hl{✔ সঠিক উত্তর: গ) ৯৪০}}

\begin{quote}
\textbf{ব্যাখ্যা:} \emph{বাবুর বেতন = ২০০০ টাকা}

\emph{বাবুর বেতনের ৫০\% = ২০০০ এর ৫০\% = ২০০০ × (৫০/১০০) = ১০০০০ টাকা}

\emph{ধরি, আরিফের বর্তমান বেতন x টাকা}

\emph{প্রশ্নমতে, x + ৬০ = ১০০০০}

\emph{⇒ x = ১০০০০ − ৬০}

\emph{∴ x = ৯৪০ টাকা।}
\end{quote}

\textbf{জলিল মিয়া পাঁচ সপ্তাহে ৩৮০ টাকা সঞ্চয় করলে একই হারে তিনি কত সপ্তাহে
এই টাকার ২.৬ গুণ সঞ্চয় করবেন?}

ক) ১০ সপ্তাহে

খ) ১১ সপ্তাহে

গ) ১২ সপ্তাহে

\textbf{\hl{ঘ) ১৩ সপ্তাহে ✔}}

\textbf{\hl{✔ সঠিক উত্তর: ঘ) ১৩ সপ্তাহে}}

\begin{quote}
\textbf{ব্যাখ্যা:} \emph{সঞ্চয়ের পরিমাণ ২.৬ গুণ হতে হলে সময়ও ২.৬ গুণ লাগবে।}

\emph{∴ প্রয়োজনীয় সময় = ৫ × ২.৬ = ১৩ সপ্তাহ।}
\end{quote}

\textbf{এক ব্যক্তি তার আয়ের (৩/৪) অংশ ব্যয় করেন। তার আয়ের শতকরা হার---}

ক) ১৫ টাকা

খ) ২০ টাকা

\textbf{\hl{গ) ২৫ টাকা ✔}}

ঘ) ৩০ টাকা

\textbf{\hl{✔ সঠিক উত্তর: গ) ২৫ টাকা}}

\begin{quote}
\textbf{ব্যাখ্যা:} \emph{ব্যয়ের হার = (৩/৪) × ১০০\% = ৭৫\%}

\emph{∴ সঞ্চয়ের হার = (১০০ − ৭৫)\% = ২৫\%}

\emph{{[}দ্রষ্টব্য: প্রশ্নে \textquotesingle সঞ্চয়ের শতকরা
হার\textquotesingle{} চাওয়া হয়েছে, যার মান ২৫\% বা ২৫ টাকা।{]}}
\end{quote}

\textbf{আবুর মাসিক আয় বাবুর মাসিক আয় থেকে ৪০\% বেশি এবং ববির মাসিক আয়ের
(৭/৮) অংশ। বাবুর মাসিক আয় ৫,০০০ টাকা হলে তাদের তিনজনের মোট মাসিক আয় কত?}

ক) ১৮,০০০ টাকা

খ) ১৯,০০০ টাকা

গ) ১৯,৫০০ টাকা

\textbf{\hl{ঘ) ২০,০০০ টাকা ✔}}

\textbf{\hl{✔ সঠিক উত্তর: ঘ) ২০,০০০ টাকা}}

\begin{quote}
\textbf{ব্যাখ্যা:} \emph{দেয়া আছে, বাবুর মাসিক আয় = ৫,০০০ টাকা}

\emph{∴ আবুর মাসিক আয় = ৫,০০০ এর ১৪০\% = ৫,০০০ × (১৪০/১০০) = ৭,০০০ টাকা}

\emph{আবার, আবুর আয় = ববির আয়ের (৭/৮) অংশ}

\emph{⇒ ৭,০০০ = ববির আয় × (৭/৮)}

\emph{∴ ববির মাসিক আয় = (৭,০০০ × ৮/৭) = ৮,০০০ টাকা}

\emph{∴ তাদের তিনজনের মোট মাসিক আয় = (৫,০০০ + ৭,০০০ + ৮,০০০) টাকা =
২০,০০০ টাকা।}
\end{quote}

\textbf{একজন কর্মকর্তার মাসিক মূল বেতন ৪০,০০০ টাকা। তিনি প্রতিমাসে মূল বেতনের
৪৫\% হারে বাড়ি ভাড়া ভাতা, ১৫০০ টাকা চিকিৎসা ভাতা, ১,০০০ টাকা আপ্যায়ন ভাতা
এবং ৭০০ টাকা মোবাইল ফোন ভাত পান। ভাতাসহ তার এক মাসের সর্বমোট বেতন কত?}

ক) ৬৫,০০০

\textbf{\hl{খ) ৬১,২০০ ✔}}

গ) ৬৬,৩৫০

ঘ) ৭০,০০০

\textbf{\hl{✔ সঠিক উত্তর: খ) ৬১,২০০}}

\begin{quote}
\textbf{ব্যাখ্যা:} \emph{বাড়ি ভাড়া ভাতা = ৪০,০০০ এর ৪৫\% = ৪০,০০০ ×
(৪৫/১০০) = ১৮,০০০ টাকা}

\emph{চিকিৎসা ভাতা = ১৫০০ টাকা}

\emph{আপ্যায়ন ভাতা = ১,০০০ টাকা}

\emph{মোবাইল ফোন ভাতা = ৭০০ টাকা}

\emph{∴ মোট ভাতা = (১৮,০০০ + ১৫০০ + ১,০০০ + ৭০০) = ২১,২০০ টাকা}

\emph{∴ ভাতাসহ মোট বেতন = (৪০,০০০ + ২১,২০০) টাকা = ৬১,২০০ টাকা।}
\end{quote}

\textbf{একজন ব্যক্তি সম্পদের ১২\% স্ত্রীকে, ৫৮\% পুত্রকে ও অবশিষ্ট ৪,৮০,০০০/-
টাকা কন্যাকে দান করেন। পুত্রের সম্পদের মূল্যমান কত?}

\textbf{\hl{ক) ৯,২৮,০০০/- ✔}}

খ) ৯,৮২,০০০/-

গ) ৯,৬৮,০০০/-

ঘ) ৯,৮৮,০০০/-

\textbf{\hl{✔ সঠিক উত্তর: ক) ৯,২৮,০০০/-}}

\begin{quote}
\textbf{ব্যাখ্যা:} \emph{কন্যাকে দান করেন = {[}১০০ − (১২ + ৫৮){]}\% = (১০০
− ৭০)\% = ৩০\%}

\emph{প্রশ্নমতে, ৩০\% = ৪,৮০,০০০ টাকা}

\emph{∴ ১\% = (৪,৮০,০০০/৩০) = ১৬,০০০ টাকা}

\emph{∴ পুত্রের প্রাপ্ত সম্পদের মূল্য (৫৮\%) = ১৬,০০০ × ৫৮ = ৯,২৮,০০০ টাকা।}
\end{quote}

\textbf{ক তার আয়ের ৫০\% খাবারে, ২০\% পোশাকে ও ১৫\% উপহার সামগ্রী ক্রয়ের
জন্য ব্যয় করেন। বাকি আয়ের পরিমাণ ৩০০ টাকা যা তিনি দান করেন। তার মাসিক আয়
কত?}

ক) ১৫০০ টাকা

\textbf{\hl{খ) ২০০০ টাকা ✔}}

গ) ২৫০০ টাকা

ঘ) ৩০০০ টাকা

\textbf{\hl{✔ সঠিক উত্তর: খ) ২০০০ টাকা}}

\begin{quote}
\textbf{ব্যাখ্যা:} \emph{মোট ব্যয়ের পরিমাণ = (৫০\% + ২০\% + ১৫\%) = ৮৫\%}

\emph{∴ অবশিষ্ট দানকৃত অংশের পরিমাণ = (১০০ − ৮৫)\% = ১৫\%}

\emph{প্রশ্নমতে, ১৫\% = ৩০০ টাকা}

\emph{∴ ১০০\% বা মোট আয় = (৩০০ × ১০০/১৫) = ২০০০ টাকা।}
\end{quote}

\textbf{মাসুদের আয় ও ব্যয়ের অনুপাত ২০ : ১৫ হলে তার মাসিক সঞ্চয় আয়ের শতকরা কত
অংশ?}

ক) ২০\%

খ) ১৫\%

\textbf{\hl{গ) ২৫\% ✔}}

ঘ) ৩০\%

\textbf{\hl{✔ সঠিক উত্তর: গ) ২৫\%}}

\begin{quote}
\textbf{ব্যাখ্যা:} \emph{মাসিক আয় ২০ টাকা হলে ব্যয় ১৫ টাকা}

\emph{∴ মাসিক সঞ্চয় = (২০ − ১৫) = ৫ টাকা}

\emph{∴ আয়ের শতকরা সঞ্চয়ের পরিমাণ = ( (৫/২০) × ১০০)\% = ২৫\%}
\end{quote}

\textbf{হেলালের মাসিক আয় ৪,২০০ টাকা এবং ব্যয় ২৯৪০ টাকা। তার মাসিক ব্যয়
মাসিক আয়ের শতকরা কত টাকা?}

ক) ৭৫\%

খ) ৬৫\%

গ) ৮০\%

\textbf{\hl{ঘ) ৭০\% ✔}}

ঙ) কোনোটিই নয়

\textbf{\hl{✔ সঠিক উত্তর: ঘ) ৭০\%}}

\begin{quote}
\textbf{ব্যাখ্যা:} \emph{মাসিক ব্যয় আয়ের শতকরা = ( (২৯৪০/৪২০০) × ১০০)\% =
৭০\%}
\end{quote}

\textbf{ইসলামের আয় আছমার আয় অপেক্ষা ২৫\% বেশি। আসমার আয় ইসলামের আয় অপেক্ষা
শতকরা কত কম?}

ক) ১০\%

খ) ১৫\%

\textbf{\hl{গ) ২০\% ✔}}

ঘ) ২৫\%

\textbf{\hl{✔ সঠিক উত্তর: গ) ২০\%}}

\begin{quote}
\textbf{ব্যাখ্যা:} \emph{ধরি, আছমার আয় = ১০০ টাকা}

\emph{∴ ইসলামের আয় = (১০০ + ২৫) = ১২৫ টাকা}

\emph{ইসলামের আয় ১২৫ টাকা হলে আছমার আয় কম = ২৫ টাকা}

\emph{∴ ইসলামের আয় ১০০ টাকা হলে আছমার আয় কম = ( (২৫ × ১০০/১২৫) ) টাকা =
২০ টাকা}

\emph{∴ আসমার আয় ২০\% কম।}
\end{quote}

\textbf{বেতন ৩০\% বৃদ্ধি পাওয়ায় একজন লোক ১১,০৫০ টাকা পায়। পূর্বে তার বেতন কত
ছিল?}

ক) ৭,৫০০ টাকা

\textbf{\hl{খ) ৮,৫০০ টাকা ✔}}

গ) ৯,০০০ টাকা

ঘ) ৯,২০০ টাকা

\textbf{\hl{✔ সঠিক উত্তর: খ) ৮,৫০০ টাকা}}

\begin{quote}
\textbf{ব্যাখ্যা:} \emph{৩০\% বৃদ্ধিতে, পূর্বের বেতন ১০০ টাকা হলে বর্তমান বেতন
= (১০০ + ৩০) = ১৩০ টাকা}

\emph{বর্তমান বেতন ১৩০ টাকা হলে পূর্বে ছিল = ১০০ টাকা}

\emph{∴ বর্তমান বেতন ১১,০৫০ টাকা হলে পূর্বে ছিল = (১০০ × ১১,০৫০/১৩০) টাকা =
৮,৫০০ টাকা।}
\end{quote}

\textbf{করিম সাহেব তার মাসিক বেতন ৫০০০ টাকা হতে প্রতি মাসে ২৪\% সঞ্চয়
করেন। তিনি এক বছরে কত টাকা সঞ্চয় করেন?}

ক) ১০,০০০ টাকা

খ) ১২,০০০ টাকা

গ) ১৫,০০০ টাকা

\textbf{\hl{ঘ) ১৪,৪০০ টাকা ✔}}

\textbf{\hl{✔ সঠিক উত্তর: ঘ) ১৪,৪০০ টাকা}}

\begin{quote}
\textbf{ব্যাখ্যা:} \emph{প্রতি মাসে সঞ্চয় করেন = ৫০০০ এর ২৪\% = ৫০০০ ×
(২৪/১০০) টাকা = ১২০০ টাকা}

\emph{∴ ১ বছর বা ১২ মাসে সঞ্চয় করেন = (১২ × ১২০০) টাকা = ১৪,৪০০ টাকা।}
\end{quote}

\textbf{মানিকের মাসিক বেতন ৯\% বৃদ্ধি পাওয়ার ফলে তার মাসিক সঞ্চয় সমান হারে
বৃদ্ধি পেয়ে ১,৮৫৩ টাকা হলো। মানিকের মাসিক সঞ্চয় আগে কত ছিল?}

ক) ১,৬৫০ টাকা

খ) ১,৬০০ টাকা

\textbf{\hl{গ) ১,৭০০ টাকা ✔}}

ঘ) ১,৭৫০ টাকা

ঙ) কোনোটিই নয়

\textbf{\hl{✔ সঠিক উত্তর: গ) ১,৭০০ টাকা}}

\begin{quote}
\textbf{ব্যাখ্যা:} \emph{৯\% বৃদ্ধিতে বর্তমান সঞ্চয় = (১০০ + ৯) = ১০৯ টাকা}

\emph{বর্তমান সঞ্চয় ১০৯ টাকা হলে পূর্বে ছিল = ১০০ টাকা}

\emph{∴ বর্তমান সঞ্চয় ১,৮৫৩ টাকা হলে পূর্বে ছিল = (১০০ × ১,৮৫৩/১০৯) টাকা =
১,৭০০ টাকা।}
\end{quote}

\textbf{সামাদ সাহেবের মাসিক বেতন ১২০০০ টাকা। এক বছর পর তার বেতন ১১\%
বৃদ্ধি পেল। আগামী বছর সামাদ সাহেব কত টাকা মাসিক বেতনে বছর শুরু করবেন?}

ক) ১২০৭৫ টাকা

\textbf{\hl{খ) ১৩৩২০ টাকা ✔}}

গ) ১৬০০০ টাকা

ঘ) ১৪৪০০ টাকা

\textbf{\hl{✔ সঠিক উত্তর: খ) ১৩৩২০ টাকা}}

\begin{quote}
\textbf{ব্যাখ্যা:} \emph{বেতন বৃদ্ধি পায় = ১২,০০০ এর ১১\% = ১২,০০০ ×
(১১/১০০) = ১৩২০ টাকা}

\emph{∴ আগামী বছর শুরুর মাসিক বেতন = (১২,০০০ + ১৩২০) টাকা = ১৩,৩২০ টাকা।}
\end{quote}

\textbf{আজাদ সাহেবের মাসিক বেতন ১৫,৫০ টাকা। এক বছর পর তার বেতন ৭\% বৃদ্ধি
পেল। বর্তমানে আজাদ সাহেবের মাসিক বেতন কত?}

ক) ১৬,০৭৫ টাকা

\textbf{\hl{খ) ১৬,০৫০ টাকা ✔}}

গ) ১৬,০০০

ঘ) ১৭,০০০

ঙ) কোনটিই নয়

\textbf{\hl{✔ সঠিক উত্তর: খ) ১৬,০৫০ টাকা}}

\begin{quote}
\textbf{ব্যাখ্যা:} \emph{৭\% বৃদ্ধিতে,}

\emph{পূর্বের বেতন ১০০ টাকা হলে বর্তমান বেতন = ১০৭ টাকা}

\emph{∴ পূর্বের বেতন ১৫,০০০ টাকা হলে বর্তমান বেতন = (১০৭ × ১৫,০০০/১০০) =
১৬,০৫০ টাকা।}
\end{quote}

\textbf{একজন শ্রমিক ১৯৯৭, ১৯৯৮ এবং ১৯৯৯ সালের প্রত্যেক বছর পূর্বতন বছরের বেতন
অপেক্ষা ১০\% বেশি বেতন পেল। সে ১৯৯৯ সালে ১৯৯৭ সালের চেয়ে কত বেশি বেতন
পেল?}

ক) 10\%

খ) 11\%

গ) 20\%

\textbf{\hl{ঘ) None of these ✔}}

\textbf{\hl{✔ সঠিক উত্তর: ঘ) None of these}}

\begin{quote}
\textbf{ব্যাখ্যা:} \emph{ধরি, ১৯৯৭ সালে বেতন ছিল = ১০০ টাকা}

\emph{∴ ১৯৯৮ সালে বেতন হয় = (১০০ + ১০) = ১১০ টাকা}

\emph{∴ ১৯৯৯ সালে বেতন হয় = ১১০ + (১১০ এর ১০\%) = ১১০ + ১১ = ১২১ টাকা}

\emph{∴ ১৯৯৭ সালের চেয়ে ১৯৯৯ সালে বেতন বেশি পেল = (১২১ − ১০০)\% = ২১\%}

\emph{যেহেতু ২১\% অপশনে নেই, তাই উত্তর \textquotesingle None of
these\textquotesingle।}
\end{quote}

\textbf{ছোটনের বেতন গত মাসে ৯\% বৃদ্ধি পাওয়ার পর সে দেখল যদি তার বেতন ৯\%
না বেড়ে ১১\% বৃদ্ধি পেত তাহলে তার মাসিক বেতন ৭২,১৫০ টাকা হতো। ছোটনের
বর্তমান মাসিক বেতন কত?}

ক) ৬৬,১৯৩ টাকা

খ) ৬৫,০০০ টাকা

\textbf{\hl{গ) ৭০,৮৫০ টাকা ✔}}

ঘ) ৭২,২০০ টাকা

\textbf{\hl{✔ সঠিক উত্তর: গ) ৭০,৮৫০ টাকা}}

\begin{quote}
\textbf{ব্যাখ্যা:} \emph{ধরি, মূল বেতন = x টাকা}

\emph{১১\% বৃদ্ধিতে বেতন = x + x এর ১১\% = (১১১x/১০০) টাকা}

\emph{প্রশ্নমতে, (১১১x/১০০) = ৭২,১৫০}

\emph{⇒ x = (৭২,১৫০ × ১০০/১১১) = ৬৫,০০০ টাকা}

\emph{∴ ছোটনের বর্তমান মাসিক বেতন (৯\% বৃদ্ধিতে) = ৬৫,০০০ + ৬৫,০০০ এর ৯\%}

\emph{= ৬৫,০০০ + ৫৮৫০ = ৭০,৮৫০ টাকা।}
\end{quote}

\textbf{আবুলের সাপ্তাহিক বেতন ১৬ শতাংশ বৃদ্ধি পেলে, তিনি প্রতি মাসে ৮১২ টাকা
উপার্জন করতে পারেন। যদি তার সাপ্তাহিক বেতন ১০ শতাংশ বৃদ্ধি পেত, তিনি প্রতি
মাসে কত টাকা উপার্জন করতেন?}

ক) ৬৫০ টাকা

খ) ৭৫০ টাকা

\textbf{\hl{গ) ৭৭০ টাকা ✔}}

ঘ) ৭৮০ টাকা

\textbf{\hl{✔ সঠিক উত্তর: গ) ৭৭০ টাকা}}

\begin{quote}
\textbf{ব্যাখ্যা:} \emph{১৬\% বৃদ্ধিতে বেতন (১০০ + ১৬)\% = ১১৬\%}

\emph{প্রশ্নমতে, ১১৬\% = ৮১২ টাকা}

\emph{∴ ১০\% বৃদ্ধিতে বেতন (১০০ + ১০)\% = ১১০\% এর মান হবে = (৮১২ ×
১১০/১১৬) = ৭৭০ টাকা।}
\end{quote}

\textbf{একটি অংশীদারী ব্যবসায়ের মোট মুনাফা থেকে `ক' ৪০\%, `খ' ২৫\%, `গ'
২০\%, `ঘ' ১০\% এবং `ঙ' ৫\% পায়। `ক' এর মুনাফা `খ' এর মুনাফার শতকরা কত
অংশ?}

ক) ৬২.৫\%

খ) ১২০\%

\textbf{\hl{গ) ১৬০\% ✔}}

ঘ) ১৭৫\%

\textbf{\hl{✔ সঠিক উত্তর: গ) ১৬০\%}}

\begin{quote}
\textbf{ব্যাখ্যা:} \emph{`ক' এর মুনাফা `খ' এর মুনাফার শতকরা অংশ = (
(৪০\%/২৫\%) × ১০০)\% = ১৬০\%}
\end{quote}

\textbf{চিনির মূল্য ২৫\% বৃদ্ধি পাওয়ায় একটি পরিবার চিনি খাওয়া এমন ভাবে কমালো
যে চিনি বাবদ ব্যয় বৃদ্ধি পেল না। ঐ পরিবার চিনি খাওয়া বাবদ শতকরা কত কমালো?}

ক) ২২\%

খ) ২৫\%

\textbf{\hl{গ) ২০\% ✔}}

ঘ) ৩০\%

\textbf{\hl{✔ সঠিক উত্তর: গ) ২০\%}}

\begin{quote}
\textbf{ব্যাখ্যা:} \emph{শর্ট টেকনিক :}

\emph{ব্যবহার হ্রাসের হার}

\emph{= (১০০×মূল্য বৃদ্ধির হার/১০০+মূল্য বৃদ্ধির হার)}

\emph{= (১০০×২৫/১০০+২৫)}

\emph{= (২৫০০/১২৫)}

\emph{= ২০\%}

\emph{লিখিত সমাধান:}

\emph{২৫\% বৃদ্ধিতে চিনির পূর্বমূল্য ১০০ টাকা হলে বর্তমান মূল্য ১২৫ টাকা}

\emph{১২৫ টাকায় কমাতে হবে = ২৫ টাকা}

\emph{১ টাকায় কমাতে হবে = (২৫/১২৫) টাকা}

\emph{১০০ টাকায় কমাতে হবে = (২৫×১০০/১২৫) = ২০\%}
\end{quote}

\textbf{চিনির মূল্য শতকরা ১০ টাকা বৃদ্ধি পেলে চিনির ব্যবহার শতকরা কত কমালে
পরিবারের চিনির খরচের কোন পরিবর্তন হবে না?}

\textbf{\hl{ক) ১০০/১১ ✔}}

খ) ৯৫/১১

গ) ১০২/১১

ঘ) ৯৩/১১

\textbf{\hl{✔ সঠিক উত্তর: ক) ১০০/১১}}

\begin{quote}
\textbf{ব্যাখ্যা:} \emph{শর্ট টেকনিক :}

\emph{ব্যবহার হ্রাসের হার}

\emph{= (১০০×মূল্য বৃদ্ধির হার/১০০+মূল্য বৃদ্ধির হার)}

\emph{= (১০০×১০/১০০+১০)}

\emph{= (১০০০/১১০)}

\emph{= (১০০/১১) \%}

\emph{লিখিত সমাধান:}

\emph{১০\% বৃদ্ধিতে চিনির পূর্বমূল্য ১০০ টাকা হলে বর্তমান মূল্য ১১০ টাকা}

\emph{১১০ টাকায় কমাতে হবে = ১০ টাকা}

\emph{১ টাকায় কমাতে হবে = (১০/১১০) টাকা}

\emph{১০০ টাকায় কমাতে হবে = (১০×১০০/১১০) = (১০০/১১) টাকা}
\end{quote}

\textbf{চিনির মূল্য ২০\% বৃদ্ধি পাওয়ায় কোন এক পরিবার চিনি খাওয়া কেমন কমালে
চিনি বাবদ ব্যয় বৃদ্ধি পাবে না।}

ক) ২০\%

খ) ১০\%

\textbf{\hl{গ) ১৬\% ✔}}

ঘ) ২৫\%

\textbf{\hl{✔ সঠিক উত্তর: গ) ১৬\%}}

\begin{quote}
\textbf{ব্যাখ্যা:} \emph{শর্ট টেকনিক :}

\emph{ব্যবহার হ্রাসের হার}

\emph{= (১০০×মূল্য বৃদ্ধির হার/১০০+মূল্য বৃদ্ধির হার)}

\emph{= (১০০×২০/১০০+২০)}

\emph{= (২০০০/১২০)}

\emph{= ১৬.৬৭\%}

\emph{লিখিত সমাধান:}

\emph{২০\% বৃদ্ধিতে চিনির পূর্বমূল্য ১০০ টাকা হলে বর্তমান মূল্য ১২০ টাকা}

\emph{১২০ টাকায় কমাতে হবে = ২০ টাকা}

\emph{১ টাকায় কমাতে হবে = (২০/১২০) টাকা}

\emph{১০০ টাকায় কমাতে হবে = (২০×১০০/১২০) = ১৬.৬৭\%}
\end{quote}

\textbf{কাপড়ের মূল্য ২০\% কমে গেলে কাপড়ের ব্যবহার কি পরিমাণ বাড়ালে কাপড়ের
জন্য খরচের কোন পরিবর্তন হবে না?}

ক) ১৫\%

খ) ২০\%

\textbf{\hl{গ) ২৫\% ✔}}

ঘ) ৩০\%

\textbf{\hl{✔ সঠিক উত্তর: গ) ২৫\%}}

\begin{quote}
\textbf{ব্যাখ্যা:} \emph{শর্ট টেকনিক:}

\emph{ব্যবহার বৃদ্ধির হার}

\emph{= (১০০×মূল্য হ্রাসের হার/১০০−মূল্য হ্রাসের হার)}

\emph{= (১০০×২০/১০০−২০)}

\emph{= (২০০০/৮০)}

\emph{= ২৫\%}

\emph{লিখিত সমাধান:}

\emph{২০\% কমে কাপড়ের পূর্বমূল্য ১০০ টাকা হলে বর্তমান মূল্য ৮০ টাকা।}

\emph{বর্তমান মূল্য ৮০ টাকা হলে পূর্বমূল্য = ১০০ টাকা}

\emph{বর্তমান মূল্য ১ টাকা হলে পূর্বমূল্য = (১০০/৮০) টাকা}

\emph{বর্তমান মূল্য ১০০ টাকা হলে পূর্বমূল্য = (১০০×১০০/৮০) = ১২৫ টাকা}

\emph{কাপড়ের ব্যবহার শতকরা বাড়াতে হবে = ১২৫−১০০ = ২৫\%}
\end{quote}

\textbf{চিনির মূল্য ২০\% কমলো, কিন্তু চিনির ব্যবহার ২০\% বৃদ্ধি পেল। এতে চিনি
বাবদ ব্যয় শতকরা কত বাড়লো বা কমলো?}

ক) ১\% বাড়লো

খ) ২\% কমলো

গ) ৩\% বাড়লো

\textbf{\hl{ঘ) ৪\% কমলো ✔}}

\textbf{\hl{✔ সঠিক উত্তর: ঘ) ৪\% কমলো}}

\begin{quote}
\textbf{ব্যাখ্যা:} \emph{শর্ট টেকনিক:}

\emph{বৃদ্ধি/হ্রাসের হার}

\emph{= বৃদ্ধির হার + হ্রাসের হার + (বৃদ্ধির হার×হ্রাসের হার/১০০)}

\emph{= ২০\% + (−২০\%) + (২০×(−২০)/১০০)}

\emph{= ০ + (−৪০০/১০০)}

\emph{= −৪}

\emph{∴ ৪\% কমলো।}

\emph{বিঃ দ্রঃ ধনাত্মক মান আসলে বাড়বে এবং ঋণাত্মক মান আসলে কমবে।}

\emph{লিখিত সমাধান:}

\emph{ধরি, ক্রয়মূল্য ১০০ টাকা}

\emph{২০\% কমে চিনির ক্রয়মূল্য = ১০০−২০ = ৮০ টাকা}

\emph{আবার, ২০\% বৃদ্ধিতে ব্যয় = (৮০ + ৮০ × (২০/১০০) ) = ৮০ + ১৬ = ৯৬ টাকা}

\emph{ব্যয় কমে = ১০০−৯৬ = ৪\%}
\end{quote}

\textbf{এক ব্যবসায়ী একটি পণ্যের মূল্য ২৫\% বাড়াল, অতঃপর বর্ধিত মূল্য থেকে ২৫\%
কমালো। সর্বশেষ মূল্য সর্বপ্রথম মূল্যের তুলনায় -}

ক) ৪.৫\% কমানো হয়েছে

\textbf{\hl{খ) ৬.২৫\% কমানো হয়েছে ✔}}

গ) ৫\% বাড়ানো হয়েছে

ঘ) ৬.২৫\% বাড়ানো হয়েছে

\textbf{\hl{✔ সঠিক উত্তর: খ) ৬.২৫\% কমানো হয়েছে}}

\begin{quote}
\textbf{ব্যাখ্যা:} \emph{শর্ট টেকনিক:}

\emph{সর্বশেষ মূল্য সর্বপ্রথম মূল্যের তুলনায়}

\emph{= বৃদ্ধির হার + হ্রাসের হার + (বৃদ্ধির হার×হ্রাসের হার/১০০)}

\emph{= ২৫\% + (−২৫\%) + (২৫×(−২৫)/১০০)}

\emph{= ০ + (−৬২৫/১০০)}

\emph{= −৬.২৫}

\emph{∴ ৬.২৫\% কমানো হয়েছে}

\emph{বিঃ দ্রঃ ধনাত্মক মান আসলে বাড়বে এবং ঋণাত্মক মান আসলে কমবে।}

\emph{লিখিত সমাধান:}

\emph{ধরি, পণ্যের ক্রয়মূল্য ১০০ টাকা}

\emph{২৫\% কমে পণ্যের ক্রয়মূল্য = ১০০−২৫ = ৭৫ টাকা}

\emph{আবার, ২৫\% বৃদ্ধিতে ব্যয় = (৭৫ + ৭৫ × (২৫/১০০) ) = ৭৫ + ১৮.৭৫ = ৯৩.৭৫
টাকা}

\emph{সর্বশেষ মূল্য সর্বপ্রথম মূল্যের তুলনায় = ১০০−৯৩.৭৫ = ৬.২৫\% কম}
\end{quote}

\textbf{গতকাল এক কেজি চিনির দাম ১০\% বেড়েছিল, আজ ১০\% কমেছে। চিনির দাম
মোট কত বেড়েছে বা কমেছে?}

\textbf{\hl{ক) ১\% কমেছে ✔}}

খ) ১\% বেড়েছে

গ) ১.৫\% কমেছে

ঘ) ১.৫\% বেড়েছে

\textbf{\hl{✔ সঠিক উত্তর: ক) ১\% কমেছে}}

\begin{quote}
\textbf{ব্যাখ্যা:} \emph{শর্ট টেকনিক:}

\emph{চিনির দাম বেড়েছে বা কমেছে}

\emph{= বৃদ্ধির হার + হ্রাসের হার + (বৃদ্ধির হার×হ্রাসের হার/১০০)}

\emph{= ১০\% + (−১০\%) + (১০×(−১০)/১০০)}

\emph{= ০ + (−১০০/১০০)}

\emph{= −১}

\emph{∴ চিনির দাম ১\% কমেছে}

\emph{বিঃ দ্রঃ ধনাত্মক মান আসলে বাড়বে এবং ঋণাত্মক মান আসলে কমবে।}

\emph{লিখিত সমাধান:}

\emph{ধরি, এক কেজি চিনির ক্রয়মূল্য ১০০ টাকা}

\emph{১০\% কমে চিনির ক্রয়মূল্য = ১০০−১০ = ৯০ টাকা}

\emph{আবার, ১০\% বৃদ্ধিতে ব্যয় = (৯০ + ৯০ × (১০/১০০) ) = ৯০ + ৯ = ৯৯ টাকা}

\emph{∴ চিনির দাম কমেছে = ১০০−৯৯ = ১\%}
\end{quote}

\textbf{একটি পণ্যের দাম প্রথমে ৪০\% বাড়ানো হলো; তারপর ১০\% কমানো হলে
মোটরের উপর কত\% বাড়লো?}

ক) ২৫\%

\textbf{\hl{খ) ২৬\% ✔}}

গ) ২৭\%

ঘ) ২৮\%

\textbf{\hl{✔ সঠিক উত্তর: খ) ২৬\%}}

\begin{quote}
\textbf{ব্যাখ্যা:} \emph{শর্ট টেকনিক:}

\emph{মোটের উপর বাড়লো}

\emph{= বৃদ্ধির হার + হ্রাসের হার + (বৃদ্ধির হার×হ্রাসের হার/১০০)}

\emph{= ৪০\% + (−১০\%) + (৪০×(−১০)/১০০)}

\emph{= ৩০ + (−৪০০/১০০)}

\emph{= ৩০ + (−৪)}

\emph{= ৩০ − ৪ = ২৬}

\emph{∴ মোটের উপর ২৬\% বাড়লো}

\emph{বিঃ দ্রঃ ধনাত্মক মান আসলে বাড়বে এবং ঋণাত্মক মান আসলে কমবে।}

\emph{লিখিত সমাধান:}

\emph{ধরি, পণ্যের ক্রয়মূল্য ১০০ টাকা}

\emph{১০\% কমে পণ্যের ক্রয়মূল্য = ১০০−১০ = ৯০ টাকা}

\emph{আবার, ৪০\% বৃদ্ধিতে ব্যয় = (৯০ + ৯০ × (৪০/১০০) ) = ৯০ + ৩৬ = ১২৬
টাকা}

\emph{∴ মোটের উপর দাম বাড়লো = ১২৬ − ১০০ = ২৬\%}
\end{quote}

\textbf{চিনির মূল্য ৬\% বেড়ে যাওয়ার ৭৯৫ টাকায় পূর্বে যত কেজি চিনি পাওয়া যেত,
বর্তমানে তার চেয়ে ৩ কেজি চিনি কম পাওয়া যায়। চিনির বর্তমান দর প্রতি কেজি
কত?}

ক) ১৫ টাকা

খ) ১৬ টাকা

গ) ১৮ টাকা

\textbf{\hl{ঘ) ১৫.৯০ টাকা ✔}}

\textbf{\hl{✔ সঠিক উত্তর: ঘ) ১৫.৯০ টাকা}}

\begin{quote}
\textbf{ব্যাখ্যা:} \emph{শর্ট টেকনিক :}

\emph{চিনির বর্তমান মূল্য = (বৃদ্ধিপ্রাপ্ত মূল্য হার×মোট মূল্য/১০০×যে পরিমাণ পণ্য কম
পাওয়া গিয়েছে)}

\emph{= (৬×৭৯৫/১০০×৩)}

\emph{= (৪৭৭০/৩০০)}

\emph{= ১৫.৯০ টাকা}

\emph{লিখিত সমাধান:}

\emph{৬\% বেড়ে চিনির পূর্বমূল্য ১০০ টাকা হলে বর্তমান মূল্য = ১০০ + ৬ = ১০৬
টাকা}

\emph{পূর্বমূল্য ১০০ টাকা হলে বর্তমান মূল্য = ১০৬ টাকা}

\emph{পূর্বমূল্য ১ টাকা হলে বর্তমান মূল্য = (১০৬/১০০) টাকা}

\emph{পূর্বমূল্য ৭৯৫ টাকা হলে বর্তমান মূল্য = (১০৬×৭৯৫/১০০) = ৮৪২.৭ টাকা}

\emph{৩ কেজি চিনির বর্তমান দর = ৮৪২.৭−৭৯৫ = ৪৭.৭ টাকা}

\emph{চিনির বর্তমান দর প্রতি কেজি = ৪৭.৭ ÷ ৩ = ১৫.৯০ টাকা}
\end{quote}

\textbf{চিনির মূল্য ৬\% বেড়ে যাওয়ার ১০৬০ টাকায় পূর্বে যত কেজি চিনি পাওয়া
যেত, বর্তমানে তার চেয়ে ৩ কেজি চিনি কম পাওয়া যায়। চিনির বর্তমান দর প্রতি
কেজি কত?}

ক) ২০ টাকা

খ) ২০.২০ টাকা

\textbf{\hl{গ) ২১.২০ টাকা ✔}}

ঘ) ২১ টাকা

\textbf{\hl{✔ সঠিক উত্তর: গ) ২১.২০ টাকা}}

\begin{quote}
\textbf{ব্যাখ্যা:} \emph{শর্ট টেকনিক :}

\emph{চিনির বর্তমান মূল্য = (বৃদ্ধিপ্রাপ্ত মূল্য হার×মোট মূল্য/১০০×যে পরিমাণ পণ্য কম
পাওয়া গিয়েছে)}

\emph{= (৬×১০৬০/১০০×৩)}

\emph{= (৬৩৬০/৩০০)}

\emph{= ২১.২০ টাকা}

\emph{লিখিত সমাধান:}

\emph{৬\% বেড়ে চিনির পূর্বমূল্য ১০০ টাকা হলে বর্তমান মূল্য = ১০০ + ৬ = ১০৬
টাকা}

\emph{পূর্বমূল্য ১০০ টাকা হলে বর্তমান মূল্য = ১০৬ টাকা}

\emph{পূর্বমূল্য ১ টাকা হলে বর্তমান মূল্য = (১০৬/১০০) টাকা}

\emph{পূর্বমূল্য ১০৬০ টাকা হলে বর্তমান মূল্য = (১০৬×১০৬০/১০০) = ১১২৩.৬ টাকা}

\emph{৩ কেজি চিনির বর্তমান দর = ১১২৩.৬ − ১০৬০ = ৬৩.৬ টাকা}

\emph{চিনির বর্তমান দর প্রতি কেজি = ৬৩.৬ ÷ ৩ = ২১.২০ টাকা}
\end{quote}

\textbf{চাউলের মূল্য ২০\% বৃদ্ধি পাওয়াতে ৭৫০ টাকায় পূর্বাপেক্ষা ৫ কেজি চাল কম
পাওয়া গেল। ১ কেজি চালের পূর্বমূল্য কত ছিল?}

\textbf{\hl{ক) ২৫ টাকা ✔}}

খ) ৩৫ টাকা

গ) ৪৫ টাকা

ঘ) ৫০ টাকা

\textbf{\hl{✔ সঠিক উত্তর: ক) ২৫ টাকা}}

\begin{quote}
\textbf{ব্যাখ্যা:} \emph{শর্ট টেকনিক:}

\emph{১ কেজি চালের বর্তমান মূল্য = (বৃদ্ধিপ্রাপ্ত মূল্য হার×মোট মূল্য/১০০×যে পরিমাণ
পণ্য কম পাওয়া গিয়েছে)}

\emph{= (২০×৭৫০/১০০×৫)}

\emph{= ৩০ টাকা}

\emph{২০\% বৃদ্ধি পাওয়াতে,}

\emph{বর্তমান মূল্য ১২০ টাকা হলে পূর্বমূল্য = ১০০ টাকা}

\emph{বর্তমান মূল্য ১ টাকা হলে পূর্বমূল্য = (১০০/১২০) টাকা}

\emph{বর্তমান মূল্য ৩০ টাকা হলে পূর্বমূল্য = (১০০×৩০/১২০) = ২৫ টাকা}

\emph{লিখিত সমাধান:}

\emph{২০\% বৃদ্ধিতে পূর্বমূল্য ১০০ টাকা হলে বর্তমান = ১০০ + ২০ = ১২০ টাকা}

\emph{পূর্বমূল্য ১০০ টাকা হলে বর্তমান মূল্য = ১২০ টাকা}

\emph{পূর্বমূল্য ১ টাকা হলে বর্তমান মূল্য = (১২০/১০০) টাকা}

\emph{পূর্বমূল্য ৭৫০ টাকা হলে বর্তমান মূল্য = (১২০×৭৫০/১০০) = ৯০০ টাকা}

\emph{∴ ৫ কেজি চাউলের বর্তমান মূল্য = ৯০০−৭৫০ = ১৫০ টাকা}

\emph{১ কেজি চাউলের বর্তমান মূল্য = ১৫০ ÷ ৫ = ৩০ টাকা}

\emph{২০\% বৃদ্ধি পাওয়াতে,}

\emph{বর্তমান মূল্য ১২০ টাকা হলে পূর্বমূল্য = ১০০ টাকা}

\emph{বর্তমান মূল্য ১ টাকা হলে পূর্বমূল্য = (১০০/১২০) টাকা}

\emph{বর্তমান মূল্য ৩০ টাকা হলে পূর্বমূল্য = (১০০×৩০/১২০) = ২৫ টাকা}
\end{quote}

\textbf{চালের মূল্য ১২\% কমে যাওয়া ৬০সং সংগে ৮৪০ টাকা ৬০০০ টাকায় পূর্বাপেক্ষা
১ কুইন্টাল চাল বেশি পাওয়া যায়। এক কুইন্টাল চালের বর্তমান মূল্য কত?}

ক) ৭৫০ টাকা

খ) ৭০০ টাকা

\textbf{\hl{গ) ৭২০ টাকা ✔}}

ঘ) ৭৫ টাকা

\textbf{\hl{✔ সঠিক উত্তর: গ) ৭২০ টাকা}}

\begin{quote}
\textbf{ব্যাখ্যা:} \emph{শর্ট টেকনিক : এক কুইন্টাল চালের বর্তমান মূল্য}

\emph{= (হ্রাসকৃত মূল্য হার×মোট মূল্য/১০০×যে পরিমাণ পণ্য বেশি পাওয়া গিয়েছে)}

\emph{= (১২×৬০০০/১০০×১)}

\emph{= ৭২০ টাকা}

\emph{লিখিত সমাধান:}

\emph{১২\% কমে চালের পূর্বমূল্য ১০০ টাকা হলে বর্তমান মূল্য = ১০০−১২ = ৮৮ টাকা}

\emph{পূর্বমূল্য ১০০ টাকা হলে বর্তমান মূল্য = ৮৮ টাকা}

\emph{পূর্বমূল্য ১ টাকা হলে বর্তমান মূল্য = (৮৮/১০০) টাকা}

\emph{পূর্বমূল্য ৬০০০ টাকা হলে বর্তমান মূল্য = (৮৮×৬০০০/১০০) = ৫২৮০ টাকা}

\emph{এক কুইন্টাল চালের বর্তমান মূল্য = ৬০০০−৫২৮০ = ৭২০ টাকা}
\end{quote}

\textbf{গমের মূল্য ১৫\% কমে যাওয়া ৬০০০ টাকায় পূর্বাপেক্ষা ১ কুইন্টাল গম বেশি
পাওয়া যায়। এক কুইন্টাল গমের বর্তমান মূল্য কত?}

ক) ৮৫০ টাকা

\textbf{\hl{খ) ৯০০ টাকা ✔}}

গ) ৭২০ টাকা

ঘ) ৭৫ টাকা

\textbf{\hl{✔ সঠিক উত্তর: খ) ৯০০ টাকা}}

\begin{quote}
\textbf{ব্যাখ্যা:} \emph{শর্ট টেকনিক :}

\emph{এক কুইন্টাল গমের বর্তমান মূল্য}

\emph{= (হ্রাসকৃত মূল্য হার×মোট মূল্য/১০০×যে পরিমাণ পণ্য বেশি পাওয়া গিয়েছে)}

\emph{= (১৫×৬০০০/১০০×১)}

\emph{= ৯০০ টাকা}

\emph{লিখিত সমাধান:}

\emph{১৫\% কমে চালের পূর্বমূল্য ১০০ টাকা হলে বর্তমান মূল্য = ১০০−১৫ = ৮৫ টাকা}

\emph{পূর্বমূল্য ১০০ টাকা হলে বর্তমান মূল্য = ৮৫ টাকা}

\emph{পূর্বমূল্য ১ টাকা হলে বর্তমান মূল্য = (৮৫/১০০) টাকা}

\emph{পূর্বমূল্য ৬০০০ টাকা হলে বর্তমান মূল্য = (৮৫×৬০০০/১০০) = ৫১০০ টাকা}

\emph{এক কুইন্টাল চালের বর্তমান মূল্য = ৬০০০−৫১০০ = ৯০০ টাকা}
\end{quote}

\textbf{চালের মূল্য ১২ (১/২) \% কমে যাওয়া ৮৪০০ টাকায় পূর্বে যে চাল কেনা যেত
এখন বর্তমান মূল্যে তার চেয়ে আরো ১ কুইন্টাল চাল বেশি কেনা যায়। চালের বর্তমান দর
কুইন্টাল প্রতি কত?}

\textbf{\hl{ক) ১০৫০ ✔}}

খ) ১১৫০

গ) ৯৫০

ঘ) ১১২৫

\textbf{\hl{✔ সঠিক উত্তর: ক) ১০৫০}}

\begin{quote}
\textbf{ব্যাখ্যা:} \emph{শর্ট টেকনিক :}

\emph{এখানে , ১২ (১/২) = (২৫/২) = ১২.৫}

\emph{এক কুইন্টাল চালের বর্তমান মূল্য}

\emph{= (হ্রাসকৃত মূল্য হার×মোট মূল্য/১০০×যে পরিমাণ পণ্য বেশি পাওয়া গিয়েছে)}

\emph{= (১২.৫×৮৪০০/১০০×১)}

\emph{= ১০৫০ টাকা}

\emph{লিখিত সমাধান:}

\emph{১২.৫\% কমে চালের পূর্বমূল্য ১০০ টাকা হলে বর্তমান মূল্য = ১০০−১২.৫ = ৮৭.৫
টাকা}

\emph{পূর্বমূল্য ১০০ টাকা হলে বর্তমান মূল্য = ৮৭.৫ টাকা}

\emph{পূর্বমূল্য ১ টাকা হলে বর্তমান মূল্য = (৮৭.৫/১০০) টাকা}

\emph{পূর্বমূল্য ৮৪০০ টাকা হলে বর্তমান মূল্য = (৮৭.৫×৮৪০০/১০০) = ৭৩৫০ টাকা}

\emph{এক কুইন্টাল চালের বর্তমান মূল্য = ৮৪০০ − ৭৩৫০ = ১০৫০ টাকা}
\end{quote}

\textbf{একটি পরীক্ষায় ৭৫টি প্রশ্ন ছিল। রফিক ৬০টি প্রশ্নের শুদ্ধ উত্তর দিয়েছে। সে
শতকরা কতটি প্রশ্নের শুদ্ধ উত্তর দিয়েছে?}

\textbf{\hl{ক) ৮০\% ✔}}

খ) ৯০\%

গ) ৬০\%

ঘ) ৭৫\%

\textbf{\hl{✔ সঠিক উত্তর: ক) ৮০\%}}

\begin{quote}
\textbf{ব্যাখ্যা:} \emph{৭৫টি প্রশ্নের মধ্যে শুদ্ধ উত্তর দেয় = ৬০টি}

\emph{∴ ১০০টি প্রশ্নের মধ্যে শুদ্ধ উত্তর দেয় = (৬০ × ১০০/৭৫) = ৮০টি}

\emph{∴ সে শতকরা ৮০টি প্রশ্নের শুদ্ধ উত্তর দিয়েছে।}
\end{quote}

\textbf{একজন ছাত্র ৮০০ নম্বরের মধ্যে ৬০০ নম্বর পেল। সে শতকরা কত কম নম্বর
পেল?}

ক) ৬০\%

খ) ৩০\%

গ) ২০\%

\textbf{\hl{ঘ) ২৫\% ✔}}

\textbf{\hl{✔ সঠিক উত্তর: ঘ) ২৫\%}}

\begin{quote}
\textbf{ব্যাখ্যা:} \emph{ছাত্রটি কম নম্বর পেয়েছে = (৮০০ − ৬০০) = ২০০ নম্বর}

\emph{∴ শতকরা কম নম্বর পেয়েছে = ( (২০০/৮০০) × ১০০)\% = ২৫\%}
\end{quote}

\textbf{৪৮০ জন পরীক্ষার্থীর মধ্যে ২৮৮ জন পাশ করলে শতকরা কত জন পাশ করল?}

ক) ৭০\%

খ) ৫০\%

গ) ২০\%

\textbf{\hl{ঘ) ৬০\% ✔}}

\textbf{\hl{✔ সঠিক উত্তর: ঘ) ৬০\%}}

\begin{quote}
\textbf{ব্যাখ্যা:} \emph{৪৮০ জন পরীক্ষার্থীর মধ্যে পাশ করে = ২৮৮ জন}

\emph{∴ ১০০ জন পরীক্ষার্থীর মধ্যে পাশ করে = (২৮৮ × ১০০/৪৮০) জন = ৬০ জন}

\emph{∴ শতকরা পাশ করল ৬০\%}
\end{quote}

\textbf{কোনো ক্লাসে চল্লিশ জন ছাত্রের মধ্যে ২৪ জন পাশ করল। ঐ ক্লাসে শতকরা
কতজন ছাত্র পাশ করতে পারেনি?}

\textbf{\hl{ক) ৪০\% ✔}}

খ) ৪৮\%

গ) ৫০\%

ঘ) কোনোটিই নয়

\textbf{\hl{✔ সঠিক উত্তর: ক) ৪০\%}}

\begin{quote}
\textbf{ব্যাখ্যা:} \emph{৪০ জনের মধ্যে পাশ করে = ২৪ জন}

\emph{∴ ১০০ জনের মধ্যে পাশ করে = (২৪ × ১০০/৪০) = ৬০ জন}

\emph{∴ শতকরা পাশ করতে পারেনি = (১০০ − ৬০) জন = ৪০ জন।}
\end{quote}

\textbf{৬০ জন ছাত্রের মধ্যে ৪২ জন ফেল করলে পাসের হার কত?}

ক) ২৫\%

খ) ২৮\%

\textbf{\hl{গ) ৩০\% ✔}}

ঘ) ৩২\%

\textbf{\hl{✔ সঠিক উত্তর: গ) ৩০\%}}

\begin{quote}
\textbf{ব্যাখ্যা:} \emph{৬০ জনের মধ্যে ফেল করে = ৪২ জন}

\emph{∴ ৬০ জনের মধ্যে পাস করে = (৬০ − ৪২) = ১৮ জন}

\emph{∴ ১০০ জনের মধ্যে পাস করে = (১৮ × ১০০/৬০) = ৩০ জন}

\emph{∴ পাসের হার ৩০\%}
\end{quote}

\textbf{একটি প্রতিযোগিতামূলক পরীক্ষায় ৬০\% পরীক্ষার্থী পাস করেছে। যারা পাস
করেনি তাদের ১৫ জন বিদেশে চলে গেল এবং ৪৫ জন ব্যবসা শুরু করল। কতজন পরীক্ষার্থী
পরীক্ষায় অংশগ্রহণ করেছে?}

ক) ১০০

\textbf{\hl{খ) ১৫০ ✔}}

গ) ২০০

ঘ) ২৫০

\textbf{\hl{✔ সঠিক উত্তর: খ) ১৫০}}

\begin{quote}
\textbf{ব্যাখ্যা:} \emph{ফেল করার হার = (১০০ − ৬০)\% = ৪০\%}

\emph{মোট ফেল করা পরীক্ষার্থীর সংখ্যা = ১৫ + ৪৫ = ৬০ জন}

\emph{প্রশ্নমতে, ৪০\% = ৬০ জন}

\emph{∴ ১০০\% = (৬০ × ১০০/৪০) = ১৫০ জন}

\emph{∴ মোট ১৫০ জন পরীক্ষার্থী পরীক্ষায় অংশগ্রহণ করেছে।}
\end{quote}

\textbf{In a class, 120 students are male and 100 students are female,
25\% of the male students and 20\% of the female students are
engineering students. 20\% of the male engineering students and 25\% of
the female engineering students passed the final exam. What percentage
of engineering students passed the exam?}

ক) 5\%

খ) 10\%

গ) 16\%

\textbf{\hl{ঘ) 22\% ✔}}

\textbf{\hl{✔ সঠিক উত্তর: ঘ) 22\%}}

\begin{quote}
\textbf{ব্যাখ্যা:} \emph{Total male engineering students = 120 × (25/100)
= 30}

\emph{Total female engineering students = 100 × (20/100) = 20}

\emph{∴ Total engineering students = 30 + 20 = 50}

\emph{Male engineering students passed = 30 × (20/100) = 6}

\emph{Female engineering students passed = 20 × (25/100) = 5}

\emph{∴ Total engineering students passed = 6 + 5 = 11}

\emph{∴ Percentage of engineering students who passed = ( (11/50) ×
100)\% = 22\%}
\end{quote}

\textbf{কোন পরীক্ষায় ২০\% পরীক্ষার্থী গণিতে ৩০\% পরীক্ষার্থী ইংরেজিতে ফেল করল
উভয় বিষয়ে ১৩ \% পরীক্ষার্থী ফেল করলে শতকরা কতজন পরীক্ষার্থী পাশ করলো?}

ক) ৫০\%

খ) ৫৭\%

\textbf{\hl{গ) ৬৩\% ✔}}

ঘ) ৬৭\%

\textbf{\hl{✔ সঠিক উত্তর: গ) ৬৩\%}}

\begin{quote}
\textbf{ব্যাখ্যা:} \emph{শর্ট টেকনিক:}

\emph{এখানে, ১ম বিষয় গণিতে ফেলের হার = ২০\%}

\emph{২য় বিষয় ইংরেজিতে ফেলের হার = ৩০\%}

\emph{এবং উভয় বিষয়ে ফেলের হার ১৩\% দেওয়া আছে।}

\emph{∴ পাশের হার = ১০০ - (১ম বিষয়ে ফেলের হার + ২য় বিষয়ে ফেলের হার − উভয়
বিষয়ে ফেলের হার)}

\emph{= ১০০ − (২০ + ৩০ − ১৩)}

\emph{= ১০০ − (৫০ − ১৩)}

\emph{= ১০০ − ৩৭}

\emph{= ৬৩\%}

\emph{লিখিত সমাধান:}

\emph{শুধু ইংরেজিতে ফেল করে = (২০ − ১৩)\% = ৭\%}

\emph{শুধু গণিতে ফেল করে = (৩০ − ১৩)\% = ১৭\%}

\emph{উভয় বিষয়ে ফেল করে = (৭ + ১৭ + ১৩)\% = ৩৭\%}

\emph{উভয় বিষয়ে পাশ করে = (১০০ − ৩৭)\% = ৬৩\%}
\end{quote}

\textbf{কোন পরীক্ষায় ৫২\% ছাত্র বিজ্ঞানে এবং ৪০\% ছাত্র অংকে ফেল করে। যদি
উভয় বিষয়ে ২৭ \% পরীক্ষার্থী ফেল করে তবে শতকরা কতজন ছাত্র পাশ করে?}

\textbf{\hl{ক) ৩৫ ✔}}

খ) ৪৮

গ) ৬০

ঘ) ৬৫

\textbf{\hl{✔ সঠিক উত্তর: ক) ৩৫}}

\begin{quote}
\textbf{ব্যাখ্যা:} \emph{শর্ট টেকনিক:}

\emph{এখানে, ১ম বিষয় বিজ্ঞানে ফেলের হার = ৫২\%}

\emph{২য় বিষয় অংকে ফেলের হার = ৪০\% এবং উভয় বিষয়ে ফেলের হার ২৭\% দেওয়া
আছে।}

\emph{∴ পাশের হার = ১০০ − (৫২ + ৪০ − ২৭)}

\emph{= ১০০ − (৯২ − ২৭)}

\emph{= ১০০ − ৬৫}

\emph{= ৩৫\%}

\emph{লিখিত সমাধান:}

\emph{শুধু বিজ্ঞানে ফেল করে = (৫২ − ২৭)\% = ২৫\%}

\emph{শুধু অংকে ফেল করে = (৪০ − ২৭)\% = ১৩\%}

\emph{উভয় বিষয়ে ফেল করে = (২৫ + ১৩ + ২৭)\% = ৬৫\%}

\emph{উভয় বিষয়ে পাশ করে = (১০০ − ৬৫)\% = ৩৫\%}
\end{quote}

\textbf{কোন পরীক্ষায় ৫২\% পরীক্ষার্থী ইংরেজিতে এবং ৪২\% পরীক্ষার্থী গণিতে ফেল
করল। যদি উভয় বিষয়ে ১৭ \% পরীক্ষার্থী ফেল করে থাকে ,তবে কতজন পরীক্ষার্থী উভয়
বিষয়ে পাশ করে?}

\textbf{\hl{ক) ২৩ জন ✔}}

খ) ২৪ জন

গ) ২৫ জন

ঘ) ২৬ জন

\textbf{\hl{✔ সঠিক উত্তর: ক) ২৩ জন}}

\begin{quote}
\textbf{ব্যাখ্যা:} \emph{শর্ট টেকনিক :}

\emph{এখানে, ১ম বিষয় ইংরেজিতে ফেলের হার = ৫২\%}

\emph{২য় বিষয় গণিতে ফেলের হার = ৪২\%}

\emph{এবং উভয় বিষয়ে ফেলের হার ১৭\% দেওয়া আছে।}

\emph{∴ পাশের হার = ১০০ − (৫২ + ৪২ − ১৭)}

\emph{= ১০০ − (৯৪ − ১৭) = ১০০ − ৭৭ = ২৩}

\emph{লিখিত সমাধান:}

\emph{শুধু বিজ্ঞানে ফেল করে = (৫২ − ১৭)\% = ৩৫\%}

\emph{শুধু অংকে ফেল করে = (৪২ − ১৭)\% = ২৫\%}

\emph{উভয় বিষয়ে ফেল করে = (৩৫ + ২৫ + ১৭)\% = ৭৭\%}

\emph{উভয় বিষয়ে পাশ করে = (১০০ − ৭৭)\% = ২৩\%}
\end{quote}

\textbf{কোন পরীক্ষায় ৮০\% পরীক্ষার্থী গণিতে এবং ৭০\% পরীক্ষার্থী বাংলায় পাশ
করলো। উভয় বিষয়ে পাশ করলো ৬০\% । উভয় বিষয়ে শতকরা কত জন ফেল করলো?}

ক) ১৫\%

\textbf{\hl{খ) ১০\% ✔}}

গ) ১২\%

ঘ) ১১\%

\textbf{\hl{✔ সঠিক উত্তর: খ) ১০\%}}

\begin{quote}
\textbf{ব্যাখ্যা:} \emph{শর্ট টেকনিক :}

\emph{এখানে, ১ম বিষয় গণিতে পাশের হার = ৮০\%}

\emph{২য় বিষয় বাংলায় পাশের হার = ৭০\%}

\emph{এবং উভয় বিষয়ে পাশের হার ৬০\% দেওয়া আছে।}

\emph{ফেলের হার = ১০০ − (১ম বিষয়ে পাশের হার + ২য় বিষয়ে পাশের হার − উভয়
বিষয়ে পাশের হার)}

\emph{= ১০০ − (৮০ + ৭০ − ৬০)}

\emph{= ১০০ − (১৫০ − ৬০)}

\emph{= ১০০ − ৯০}

\emph{= ১০ \%}

\emph{লিখিত সমাধান:}

\emph{শুধু গণিতে পাস করে = (৮০ − ৬০)\% = ২০\%}

\emph{শুধু বাংলায় পাস করে = (৭০ − ৬০)\% = ১০\%}

\emph{উভয় বিষয়ে পাস করেছে ৬০\%}

\emph{দুই বিষয়ে মোট পাস করে = ২০ + ১০ + ৬০ = ৯০\%}

\emph{উভয় বিষয়ে ফেল করেছে = ১০০ − ৯০ = ১০\%}
\end{quote}

\textbf{কোন পরীক্ষায় ৮৫\% পরীক্ষার্থী পদার্থবিদ্যায় , ৮০\% পরীক্ষার্থী
রসায়নবিদ্যায় এবং ৭৫\% পরীক্ষার্থী উভয় বিষয়ে পাশ করে ।কত জন পরীক্ষার্থী উভয়
বিষয়ে ফেল করে?}

ক) ৮ জন

\textbf{\hl{খ) ১০ জন ✔}}

গ) ১২ জন

ঘ) ১৫ জন

\textbf{\hl{✔ সঠিক উত্তর: খ) ১০ জন}}

\begin{quote}
\textbf{ব্যাখ্যা:} \emph{শর্ট টেকনিক :}

\emph{এখানে, ১ম বিষয় পদার্থবিদ্যায় পাশের হার = ৮৫\%}

\emph{২য় বিষয় রসায়নবিদ্যায় পাশের হার = ৮০\%}

\emph{এবং উভয় বিষয়ে পাশের হার ৭৫\% দেওয়া আছে।}

\emph{ফেলের হার = ১০০ − (১ম বিষয়ে পাশের হার + ২য় বিষয়ে পাশের হার − উভয়
বিষয়ে পাশের হার)}

\emph{= ১০০ − (৮৫ + ৮০ − ৭৫)}

\emph{= ১০০ − (১৬৫ − ৭৫)}

\emph{= ১০০ − ৯০}

\emph{= ১০ \%}

\emph{লিখিত সমাধান:}

\emph{শুধু গণিতে পাস করে = (৮৫ − ৭৫)\% = ১০\%}

\emph{শুধু বাংলায় পাস করে = (৮০ − ৭৫)\% = ৫\%}

\emph{উভয় বিষয়ে পাস করেছে ৭৫\%}

\emph{দুই বিষয়ে মোট পাস করে = ১০ + ৫ + ৭৫ = ৯০\%}

\emph{উভয় বিষয়ে ফেল করেছে = ১০০ − ৯০ = ১০\%}
\end{quote}

\textbf{কোন পরীক্ষায় শতকরা ৮০ জন অংকে এবং শতকরা ৭০ জন বাংলায় পাশ করল। উভয়
বিষয়ে পাশ করল শতকরা ৫০ জন। উভয় বিষয়ে ফেল করল শতকরা কতজন?}

ক) ৩০

\textbf{\hl{খ) ১০ ✔}}

গ) ২০

ঘ) ০

\textbf{\hl{✔ সঠিক উত্তর: খ) ১০}}

\begin{quote}
\textbf{ব্যাখ্যা:} \emph{শর্ট টেকনিক :}

\emph{এখানে, ১ম বিষয় অংকে পাশের হার = ৮০\%}

\emph{২য় বিষয় বাংলায় পাশের হার = ৭০\%}

\emph{এবং উভয় বিষয়ে পাশের হার ৫০\% দেওয়া আছে।}

\emph{ফেলের হার = ১০০ − (১ম বিষয়ে পাশের হার + ২য় বিষয়ে পাশের হার − উভয়
বিষয়ে পাশের হার)}

\emph{= ১০০ − (৮০ + ৭০ − ৫০)}

\emph{= ১০০ − (১৫০ − ৫০)}

\emph{= ১০০ − ১০০}

\emph{= ০ \%}

\emph{লিখিত সমাধান:}

\emph{শুধু গণিতে পাস করে = (৮০ − ৫০)\% = ৩০\%}

\emph{শুধু বাংলায় পাস করে = (৭০ − ৫০)\% = ২০\%}

\emph{উভয় বিষয়ে পাস করেছে ৫০\%}

\emph{দুই বিষয়ে মোট পাস করে = ৩০ + ২০ + ৫০ = ১০০\%}

\emph{উভয় বিষয়ে ফেল করেছে = ১০০ − ১০০ = ০\%}
\end{quote}

\textbf{কোনো শ্রেণীতে ১০০ জন পরীক্ষার্থী ছিল। বার্ষিক পরীক্ষায় ৯৪ জন বাংলায়
এবং ৮০ জন গণিতে পাস করেছে। ৭৫ জন উভয় বিষয়ে পাস করলে কত জন উভয় বিষয়ে ফেল
করেছে?}

ক) ৫

খ) ৩

গ) ২

\textbf{\hl{ঘ) ১ ✔}}

\textbf{\hl{✔ সঠিক উত্তর: ঘ) ১}}

\begin{quote}
\textbf{ব্যাখ্যা:} \emph{শর্ট টেকনিক : এখানে, ১ম বিষয় বাংলায় পাস করে = ৯৪}

\emph{২য় বিষয় গণিতে পাস করে = ৮০}

\emph{এবং উভয় বিষয়ে পাশের হার ৭৫ দেওয়া আছে।}

\emph{উভয় বিষয়ে ফেল করেছে = ১০০ − (১ম বিষয়ে পাশের হার + ২য় বিষয়ে পাশের
হার − উভয় বিষয়ে পাশের হার)}

\emph{= ১০০ − (৯৪ + ৮০ − ৭৫)}

\emph{= ১০০ − (১৭৪ − ৭৫)}

\emph{= ১০০ − ৯৯ = ১}

\emph{লিখিত সমাধান:}

\emph{শুধু বাংলায় পাস করে = (৯৪ − ৭৫) = ১৯}

\emph{শুধু গণিতে পাস করে = (৮০ − ৭৫) = ৫}

\emph{উভয় বিষয়ে পাস করেছে ৭৫ জন}

\emph{দুই বিষয়ে মোট পাস করে = ১৯ + ৫ + ৭৫ = ৯৯}

\emph{উভয় বিষয়ে ফেল করেছে = ১০০ − ৯৯ = ১ জন}
\end{quote}

\textbf{কোন পরীক্ষায় ২০০ জনের মধ্যে ৭০\% বাংলায় এবং ৬০\% ইংরেজিতে পাশ করে
করে এবং ৪০\% উভয় বিষয়ে পাশ করে। উভয় বিষয়ে ফেল করে কত জন ?}

ক) ১০ জন

খ) ১৫ জন

গ) ২৫ জন

\textbf{\hl{ঘ) ২০ জন ✔}}

\textbf{\hl{✔ সঠিক উত্তর: ঘ) ২০ জন}}

\begin{quote}
\textbf{ব্যাখ্যা:} \emph{শর্ট টেকনিক:}

\emph{এখানে, ১ম বিষয় বাংলায় পাশের হার = ৭০\%,}

\emph{২য় বিষয় ইংরেজিতে পাশের হার = ৬০\%}

\emph{এবং উভয় বিষয়ে পাশের হার ৪০\% দেওয়া আছে।}

\emph{∴ ফেলের হার = ১০০ − (৭০ + ৬০ − ৪০)}

\emph{= ১০০ − (১৩০ − ৪০)}

\emph{= ১০০ − ৯০}

\emph{= ১০\%}

\emph{∴ সুতরাং উভয় বিষয়ে ফেল = ২০০ এর ১০\% = ২০ জন}

\emph{লিখিত সমাধান:}

\emph{শুধু বাংলায় পাস করে = ৭০ − ৪০ = ৩০\%}

\emph{শুধু ইংরেজীতে পাস করে = ৬০ − ৪০ = ২০\%}

\emph{দুই বিষয়ে মোট পাস করে = ৩০ + ২০ + ৪০ = ৯০\%}

\emph{উভয় বিষয়ে ফেল করে = ১০০ − ৯০ = ১০\%}

\emph{২০০ জনের ১০\% = ২০ জন}

\emph{উভয় বিষয়ে ফেল করে মোট = ২০ জন}
\end{quote}

\textbf{কোন স্কুলে ৭০\% শিক্ষার্থী ইংরেজিতে এবং ৮০\% শিক্ষার্থী বাংলায় পাশ
করেছে। কিন্তু ১০\% শিক্ষার্থী উভয় বিষয়ে ফেল করেছে। যদি উভয় বিষয়ে ৩৬০ জন
পরীক্ষার্থী পাশ করে তবে ঐ স্কুলে কত জন শিক্ষার্থী পরীক্ষা দিয়েছে?}

ক) ৫০০ জন

খ) ৪৫০ জন

গ) ৭০০ জন

\textbf{\hl{ঘ) ৬০০ জন ✔}}

\textbf{\hl{✔ সঠিক উত্তর: ঘ) ৬০০ জন}}

\begin{quote}
\textbf{ব্যাখ্যা:} \emph{শর্ট টেকনিক :}

\emph{এখানে, ১ম বিষয় ইংরেজিতে পাশের হার = ৭০\%}

\emph{২য় বিষয় বাংলায় পাশের হার = ৮০\%}

\emph{উভয় বিষয়ে ফেলের হার ১০\% এবং উভয় বিষয়ে পাসকৃত ছাত্র ৩৬০ জন দেওয়া
আছে।}

\emph{তাহলে, ১ম বিষয় ইংরেজিতে ফেলের হার = ১০০\% − ৭০\% = ৩০\%}

\emph{২য় বিষয় বাংলায় ফেলের হার = ১০০\% − ৮০\% = ২০\%}

\emph{মোট শিক্ষার্থী = (উভয় বিষয়ে পাসকৃত ছাত্র/১ম বিষয়ে ফেল+২য় বিষয়ে ফেল+উভয়
বিষয়ে ফেল)×১০০}

\emph{= (৩৬০/৩০+২০+১০)×১০০}

\emph{= (৩৬০/৬০)×১০০}

\emph{= ৬০০ জন}

\emph{লিখিত সমাধানঃ}

\emph{শুধু ইংরেজিতে ফেল করে = (১০০ − ৭০) − ১০\% = ২০\%}

\emph{শুধু বাংলায় ফেল করে = (১০০ − ৮০) − ১০\% = ১০\%}

\emph{উভয় বিষয়ে ফেল করেছে ১০\%}

\emph{দুই বিষয়ে মোট ফেল করে = ২০ + ১০ + ১০ = ৪০\%}

\emph{দুই বিষয়ে মোট পাস করে = ১০০ − ৪০ = ৬০\%}

\emph{৬০ জন পাস করে = ১০০ জনের মধ্যে}

\emph{১ জন পাস করে = (১০০/৬০) জনের মধ্যে}

\emph{৩৬০ জন পাস করে = (১০০×৩৬০/৬০) = ৬০০ জনের মধ্যে}

\emph{ঐ স্কুলে পরীক্ষা দিয়েছে = ৬০০ জন শিক্ষার্থী}
\end{quote}

\textbf{কফিলাতলী হাইস্কুলে ৭০\% পরীক্ষার্থী ইংরেজিতে এবং ৮০\% বাংলায় পাস
করে। উভয় বিষয়ে ফেল করেছে ১০\%, যদি উভয় বিষয়ে ৩৬০ জন পাস করে তবে ঐ স্কুলে
কত জন পরীক্ষার্থী পরীক্ষা দিয়েছে?}

ক) ৫০০ জন

খ) ৭০০ জন

গ) ৮০০ জন

\textbf{\hl{ঘ) ৬০০ জন ✔}}

\textbf{\hl{✔ সঠিক উত্তর: ঘ) ৬০০ জন}}

\begin{quote}
\textbf{ব্যাখ্যা:} \emph{শর্ট টেকনিক :}

\emph{এখানে, ১ম বিষয় ইংরেজিতে পাশের হার = ৭০\%}

\emph{২য় বিষয় বাংলায় পাশের হার = ৮০\%}

\emph{উভয় বিষয়ে ফেলের হার ১০\% এবং উভয় বিষয়ে পাসকৃত ছাত্র ৩৬০ জন দেওয়া
আছে।}

\emph{তাহলে, ১ম বিষয় ইংরেজিতে ফেলের হার = ১০০\% − ৭০\% = ৩০\%}

\emph{২য় বিষয় বাংলায় ফেলের হার = ১০০\% − ৮০\% = ২০\%}

\emph{মোট শিক্ষার্থী = (উভয় বিষয়ে পাসকৃত ছাত্র/১ম বিষয়ে ফেল+২য় বিষয়ে ফেল+উভয়
বিষয়ে ফেল)×১০০}

\emph{= (৩৬০/৩০+২০+১০)×১০০}

\emph{= (৩৬০/৬০)×১০০}

\emph{= ৬০০ জন}

\emph{লিখিত সমাধানঃ}

\emph{শুধু ইংরেজিতে ফেল করে = (১০০ − ৭০) − ১০\% = ২০\%}

\emph{শুধু বাংলায় ফেল করে = (১০০ − ৮০) − ১০\% = ১০\%}

\emph{উভয় বিষয়ে ফেল করেছে ১০\%}

\emph{দুই বিষয়ে মোট ফেল করে = ২০ + ১০ + ১০ = ৪০\%}

\emph{দুই বিষয়ে মোট পাস করে = ১০০ − ৪০ = ৬০\%}

\emph{৬০ জন পাস করে = ১০০ জনের মধ্যে}

\emph{১ জন পাস করে = (১০০/৬০) জনের মধ্যে}

\emph{৩৬০ জন পাস করে = (১০০×৩৬০/৬০) = ৬০০ জনের মধ্যে}

\emph{ঐ স্কুলে পরীক্ষা দিয়েছে = ৬০০ জন পরীক্ষার্থী}
\end{quote}

\textbf{কোন স্কুলে ৭০\% শিক্ষার্থী বাংলায় এবং ৮০\% গণিতে পাশ করেছে। কিন্তু
১০\% শিক্ষার্থী উভয় বিষয়ে ফেল করেছে। যদি উভয় বিষয়ে ২৪০ জন পরীক্ষার্থী পাশ
করে থাকে তবে ঐ স্কুলে কত জন শিক্ষার্থী পরীক্ষা দিয়েছে?}

\textbf{\hl{ক) ৪০০ জন ✔}}

খ) ৪৫০ জন

গ) ৫০০ জন

ঘ) ৫৬০ জন

\textbf{\hl{✔ সঠিক উত্তর: ক) ৪০০ জন}}

\begin{quote}
\textbf{ব্যাখ্যা:} \emph{শর্ট টেকনিক :}

\emph{এখানে, ১ম বিষয় বাংলায় পাশের হার = ৭০\%}

\emph{২য় বিষয় গণিতে পাশের হার = ৮০\%}

\emph{উভয় বিষয়ে ফেলের হার ১০\% এবং উভয় বিষয়ে পাসকৃত ছাত্র ২৪০ জন দেওয়া
আছে।}

\emph{তাহলে, ১ম বিষয় বাংলায় ফেলের হার = ১০০\% − ৭০\% = ৩০\%}

\emph{২য় বিষয় গণিতে ফেলের হার = ১০০\% − ৮০\% = ২০\%}

\emph{মোট শিক্ষার্থী = (উভয় বিষয়ে পাসকৃত ছাত্র/১ম বিষয়ে ফেল+২য় বিষয়ে ফেল+উভয়
বিষয়ে ফেল)×১০০}

\emph{= (২৪০/৩০+২০+১০)×১০০}

\emph{= (২৪০/৬০)×১০০}

\emph{= ৪০০ জন}

\emph{লিখিত সমাধানঃ}

\emph{শুধু বাংলায় ফেল করে = (১০০ − ৭০) − ১০\% = ২০\%}

\emph{শুধু গণিতে ফেল করে = (১০০ − ৮০) − ১০\% = ১০\%}

\emph{উভয় বিষয়ে ফেল করেছে ১০\%}

\emph{দুই বিষয়ে মোট ফেল করে = ২০ + ১০ + ১০ = ৪০\%}

\emph{দুই বিষয়ে মোট পাস করে = ১০০ − ৪০ = ৬০\%}

\emph{৬০ জন পাস করে = ১০০ জনের মধ্যে}

\emph{১ জন পাস করে = (১০০/৬০) জনের মধ্যে}

\emph{২৪০ জন পাস করে = (১০০×২৪০/৬০) = ৪০০ জনের মধ্যে}

\emph{ঐ স্কুলে পরীক্ষা দিয়েছে = ৪০০ জন শিক্ষার্থী}
\end{quote}

\textbf{কোন স্কুলে ৭০\% শিক্ষার্থী ইংরেজিতে এবং ৮০\% শিক্ষার্থী বাংলায় পাশ
করেছে। কিন্তু ১০\% শিক্ষার্থী উভয় বিষয়ে ফেল করেছে। যদি উভয় বিষয়ে ৩০০ জন
পরীক্ষার্থী পাশ করে তবে ঐ স্কুলে কত জন শিক্ষার্থী পরীক্ষা দিয়েছে?}

ক) ৪০০ জন

\textbf{\hl{খ) ৫০০ জন ✔}}

গ) ৫৬০ জন

ঘ) ৭৬০ জন

\textbf{\hl{✔ সঠিক উত্তর: খ) ৫০০ জন}}

\begin{quote}
\textbf{ব্যাখ্যা:} \emph{শর্ট টেকনিক :}

\emph{এখানে, ১ম বিষয় ইংরেজিতে পাশের হার = ৭০\%}

\emph{২য় বিষয় বাংলায় পাশের হার = ৮০\%}

\emph{উভয় বিষয়ে পাশের হার ১০\% এবং উভয় বিষয়ে পাসকৃত ছাত্র ৩০০ জন দেওয়া
আছে।}

\emph{তাহলে, ১ম বিষয় ইংরেজিতে ফেলের হার = ১০০\% − ৭০\% = ৩০\%}

\emph{২য় বিষয় বাংলায় ফেলের হার = ১০০\% − ৮০\% = ২০\%}

\emph{মোট পরীক্ষার্থী = (উভয় বিষয়ে পাসকৃত ছাত্র/১ম বিষয়ে ফেল+২য় বিষয়ে
ফেল+উভয় বিষয়ে ফেল)×১০০}

\emph{= (৩০০/৩০+২০+১০)×১০০}

\emph{= (৩০০/৬০)×১০০}

\emph{= ৫০০ জন}

\emph{লিখিত সমাধানঃ}

\emph{শুধু ইংরেজিতে ফেল করে = (১০০ − ৭০) − ১০\% = ২০\%}

\emph{শুধু বাংলায় ফেল করে = (১০০ − ৮০) − ১০\% = ১০\%}

\emph{উভয় বিষয়ে ফেল করেছে ১০\%}

\emph{দুই বিষয়ে মোট ফেল করে = ২০ + ১০ + ১০ = ৪০\%}

\emph{দুই বিষয়ে মোট পাস করে = ১০০ − ৪০ = ৬০\%}

\emph{৬০ জন পাস করে = ১০০ জনের মধ্যে}

\emph{১ জন পাস করে = (১০০/৬০) জনের মধ্যে}

\emph{৩০০ জন পাস করে = (১০০×৩০০/৬০) = ৫০০ জনের মধ্যে}

\emph{ঐ স্কুলে পরীক্ষা দিয়েছে = ৫০০ জন পরীক্ষার্থী}
\end{quote}

\textbf{কোন পরীক্ষায় ৯০\% পরীক্ষার্থী ইতিহাসে এবং ৮৫\% পরীক্ষার্থী ভূগোলে
কৃতকার্য হল।যদি উভয় বিষয়ে কেউ ফেল না করে এবং উভয় বিষয়ে ২২৫ জন পাস করে থাকে
, তবে ঐ পরীক্ষায় কতজন অংশগ্রহণ করেছিল?}

\textbf{\hl{ক) ৩০০ ✔}}

খ) ৪০০

গ) ৫০০

ঘ) ৬০০

\textbf{\hl{✔ সঠিক উত্তর: ক) ৩০০}}

\begin{quote}
\textbf{ব্যাখ্যা:} \emph{শর্ট টেকনিক :}

\emph{এখানে, ১ম বিষয় ইতিহাসে পাশের হার = ৯০\%}

\emph{২য় বিষয় ভূগোলে পাশের হার = ৮৫\%}

\emph{এবং উভয় বিষয়ে মোট পাসকৃত ছাত্র ২২৫ জন দেওয়া আছে।}

\emph{তাহলে, ১ম বিষয় ইতিহাসে ফেলের হার = ১০০\% − ৯০\% = ১০\%}

\emph{২য় বিষয় ভূগোলে ফেলের হার = ১০০\% − ৮৫\% = ১৫\%}

\emph{মোট পরীক্ষার্থী = (উভয় বিষয়ে পাসকৃত ছাত্র/১০০−(১ম বিষয়ে ফেল+২য় বিষয়ে
ফেল))×১০০}

\emph{= (২২৫/১০০−(১০+১৫))×১০০}

\emph{= (২২৫/১০০−২৫)×১০০}

\emph{= (২২৫/৭৫)×১০০ = ৩ × ১০০ = ৩০০}

\emph{লিখিত সমাধানঃ}

\emph{শুধু ইতিহাসে ফেল করে = ১০০ − ৯০ = ১০\%}

\emph{শুধু ভূগোলে ফেল করে = ১০০ − ৮৫ = ১৫\%}

\emph{উভয় বিষয়ে পাস করে = ১০০ − (১০\% + ১৫\%) = ৭৫\%}

\emph{৭৫ জন পাস করে = ১০০ জনের মধ্যে}

\emph{১ জন পাস করে = (১০০/৭৫) জনের মধ্যে}

\emph{২২৫ জন পাস করে = (২২৫×১০০/৭৫) = ৩০০ জনের মধ্যে}

\emph{∴ সুতরাং, ঐ পরীক্ষায় ৩০০ জন অংশগ্রহণ করেছিল।}
\end{quote}

\textbf{কোন পরীক্ষায় ৮৫\% ছাত্র ইংরেজিতে পাস করেছে।ইংরেজিতে ফেলের সংখ্যা
মোট ৭৫ জন হলে পরীক্ষার্থীর সংখ্যা কত?}

\textbf{\hl{ক) ৫০০ জন ✔}}

খ) ৩০০ জন

গ) ৭৭৫ জন

ঘ) ৬৫০ জন

\textbf{\hl{✔ সঠিক উত্তর: ক) ৫০০ জন}}

\begin{quote}
\textbf{ব্যাখ্যা:} \emph{শর্ট টেকনিক :}

\emph{এখানে , ইংরেজিতে পাশের হার = ৮৫\% এবং ফেলের সংখ্যা মোট ৭৫ জন দেওয়া
আছে।}

\emph{মোট পরীক্ষার্থী = (ফেলের সংখ্যা/১০০−পাশের হার)×১০০}

\emph{= (৭৫/১০০−৮৫)×১০০}

\emph{= (৭৫/১৫)×১০০}

\emph{= ৫০০ জন}

\emph{লিখিত সমাধানঃ}

\emph{পরীক্ষার্থীর সংখ্যা ১০০\%}

\emph{ইংরেজিতে ফেল করে = ১০০ − ৮৫ = ১৫\%}

\emph{প্রশ্নমতে,}

\emph{১৫\% = ৭৫ জন}

\emph{বা,১\% = ৭৫ ÷ ১৫ = ৫ জন}

\emph{বা, ১০০\% = ৫ × ১০০ = ৫০০ জন}
\end{quote}

\section{54. ঐকিক নিয়মের অংক: শর্টকাট টেকনিক ও গুরুত্বপূর্ণ প্রশ্ন
সমাধান}\label{ux990ux995ux995-ux9a8ux9afux9aeux9b0-ux985ux995-ux9b6ux9b0ux99fux995ux99f-ux99fux995ux9a8ux995-ux993-ux997ux9b0ux9a4ux9acux9aaux9b0ux9a3-ux9aaux9b0ux9b6ux9a8-ux9b8ux9aeux9a7ux9a8}

পরীক্ষার্থীদের প্রস্তুতিকে সহজ ও গোছালো করতে এই অধ্যায়টিকে মূলত ৬টি গুরুত্বপূর্ণ
ভাগে ভাগ করা হয়েছে। প্রতিটি ভাগের বিগত সালের প্রশ্নগুলো বিস্তারিত ব্যাখ্যা এবং
সঠিক নিয়মে সমাধানসহ নিচে আলোচনা করা হলো। মনোযোগ দিয়ে প্রতিটি সমস্যার গভীরে
প্রবেশ করুন এবং নিজের প্রস্তুতি শানিয়ে নিন।

\begin{itemize}
\item
  \textbf{ঐকিক নিয়ম (Unitary Method) কী?}
\item
  কাজ ও সময় সংক্রান্ত
\item
  পিপা / ট্যাংক/চৌবাচ্চা সংক্রান্ত
\item
  নৌকা ও স্রোত সংক্রান্ত
\item
  বানরের বাঁশে ওঠা সংক্রান্ত
\item
  গতিবেগ দূরত্ব ও সময় সংক্রান্ত
\item
  ট্রেন ও গতিবেগ সংক্রান্ত
\end{itemize}

\subsubsection{ঐকিক নিয়ম (Unitary Method)
কী?}\label{ux990ux995ux995-ux9a8ux9dfux9ae-unitary-method-ux995}

যে প্রক্রিয়ায় প্রথমে অনেকগুলো বা নির্দিষ্ট সংখ্যক বস্তুর পরিমাণ বা মূল্য থেকে প্রথমে
একটি বা এককের মান বের করা হয় এবং পরবর্তীতে সেই একক মানের ওপর ভিত্তি করে
কাঙ্ক্ষিত বা অন্য যেকোনো সংখ্যক বস্তুর মান নির্ণয় করা হয়, তাকে গণিতের ভাষায়
ঐকিক নিয়ম বলে।

সহজ কথায়, "অনেক থেকে এক, এবং এক থেকে অনেক"---এই মূলনীতির ওপর ভিত্তি করেই
ঐকিক নিয়মের সমস্ত হিসাব পরিচালিত হয়। যেমন: ৫টি আমের দাম ১০০ টাকা হলে, ১টি
আমের দাম কত---তা বের করে পরবর্তীতে ১২টি আমের দাম বের করার প্রক্রিয়াটিই হলো
ঐকিক নিয়ম।

\textbf{ঐকিক নিয়মের মূল ভিত্তি:}

ঐকিক নিয়ম মূলত দুটি বিষয়ের ওপর নির্ভর করে কাজ করে:

\begin{itemize}
\item
  তুলনা ও অনুপাত: বিভিন্ন রাশির মধ্যে পারস্পরিক সম্পর্ক স্থাপন করা।
\item
  ধাপভিত্তিক হিসাব: বহুবচন থেকে একবচনে (একক মান) পৌঁছানো এবং সেখান থেকে
  কাঙ্ক্ষিত পরিমাণে পৌঁছানো।
\end{itemize}

গণিতের ভিত্তি মজবুত করতে এবং যেকোনো জটিল সমস্যার সহজ সমাধান বের করতে এই
নিয়মটি অত্যন্ত কার্যকর।

\subsubsection{কাজ ও সময়
সংক্রান্ত}\label{ux995ux99c-ux993-ux9b8ux9aeux9df-ux9b8ux995ux9b0ux9a8ux9a4}

\textbf{প্রশ্ন নম্বর-} ১. ১০ জন লোক একটি কাজ ২০ দিন করতে পারে। ৮জন লোকের ঐ
কাজটি করতে কতদিন লাগবে ? {[}বিএসটিআই এর অফিস সহকারী -২০১১{]}

\begin{itemize}
\item
  (ক) ২৫ দিন
\item
  (খ) ২৪ দিন
\item
  (গ) ৩০ দিন
\item
  (ঘ) ১৬ দিন
\end{itemize}

\textbf{\hl{উত্তরঃ (ক) ২৫ দিন}}

\textbf{সমাধান:} ১০ জন লোক একটি কাজ করতে পারে = ২০ দিনে

১ জন লোক একটি কাজ করতে পারে = ২০×১০ = ২০০ দিনে

৮ জন লোক একটি কাজ করতে পারে = ২০০÷৮ = ২৫ দিনে

বিঃ দ্রঃ মনে রাখতে হবে লোক ও সময় সংক্রান্ত অংকে ২য় লাইনে সবসময়ই গুন হয়।উত্তরে
যা চাওয়া হবে তা সবসময় লাইনের শেষ অথবা ডানদিকে রাখতে হবে।

\textbf{শর্ট টেকনিক :}

প্রথম লোক ×প্রথম সময় = দ্বিতীয় লোক ×দ্বিতীয় সময়

সুতরাং দ্বিতীয় সময় =
\(\frac{\text{প্রথম}\ \text{লোক}\  \times \ \text{প্রথম}\ \text{সময়}}{\text{দ্বিতীয়}\ \text{লোক}}\) =
\(\frac{\text{১০} \times \text{২০}}{\text{৮}}\) = ২৫ দিনে

\textbf{প্রশ্ন নম্বর-} ২. ৫৬ জন শ্রমিক একটি কাজ ২১ দিনে শেষ করতে পারে। ১৪
দিনে কাজটি শেষ করতে হলে নতুন কত জন শ্রমিক নিয়োগ করতে হবে ? {[}প্রাক-প্রাথমিক
সহকারী শিক্ষক-২০১৩(সুরমা){]}

\begin{itemize}
\item
  (ক) ২৪ জন
\item
  (খ) ২৬ জন
\item
  (গ) ২৮ জন
\item
  (ঘ) ৩০ জন
\end{itemize}

\textbf{\hl{উত্তরঃ (গ) ২৮ জন}}

\textbf{সমাধান:}

২১ দিনে একটি কাজ করতে পারে ৫৬ জন শ্রমিকে

১ দিনে একটি কাজ করতে পারে ৫৬×২১ = ১১৭৬ জন শ্রমিকে

১৪ দিনে একটি কাজ করতে পারে ১১৭৬÷১৪ = ৮৪ জন শ্রমিকে

নতুন শ্রমিক নিয়োগ করতে হবে = ৮৪-৫৬ = ২৮ জন

\textbf{শর্ট টেকনিক :}

প্রথম লোক × প্রথম সময় = দ্বিতীয় লোক × দ্বিতীয় সময়

সুতরাং দ্বিতীয় লোক =
\(\frac{\text{প্রথম}\ \text{লোক}\  \times \ \text{প্রথম}\ \text{সময়}}{\text{দ্বিতীয়}\ \text{লোক}}\) =
\(\frac{\text{৫৬} \times \text{২১}}{\text{১৪}}\) = ৮৪ জন

নতুন শ্রমিক নিয়োগ করতে হবে = ৮৪-৫৬ = ২৮ জন

\textbf{প্রশ্ন নম্বর-} ৩. ১৫ জন লোক একটি কাজ শেষ করে ৩ ঘণ্টায়। ৫ জন লোকের ঐ
কাজ শেষ করতে কত সময় লাগবে ? {[}খাদ্য পরিদশক পরীক্ষা-২০০১{]}

\begin{itemize}
\item
  (ক) ৫ ঘন্টা
\item
  (খ) ৭.৫ ঘন্টা
\item
  (গ) ৯ ঘন্টা
\item
  (ঘ) ৪ ঘন্টা
\end{itemize}

\textbf{\hl{উত্তরঃ (গ) ৯ ঘন্টা}}

\textbf{সমাধান:}

১৫ জন লোক একটি কাজ করতে পারে = ৩ ঘণ্টায়

১ জন লোক একটি কাজ করতে পারে = ১৫×৩ = ৪৫ ঘণ্টায়

৫ জন লোক একটি কাজ করতে পারে = ৪৫÷৫ = ৯ ঘণ্টায়

বিঃ দ্রঃ মনে রাখতে হবে লোক ও সময় সংক্রান্ত অংকে ২য় লাইনে সবসময়ই গুন হয়।উত্তরে
যা চাওয়া হবে তা সবসময় লাইনের শেষ অথবা ডানদিকে রাখতে হবে।

\textbf{শর্ট টেকনিক :}

প্রথম লোক ×প্রথম সময় = দ্বিতীয় লোক ×দ্বিতীয় সময়

সুতরাং দ্বিতীয় সময় =
\(\frac{\text{প্রথম}\ \text{লোক}\  \times \ \text{প্রথম}\ \text{সময়}}{\text{দ্বিতীয়}\ \text{লোক}}\) =
\(\frac{\text{১৫} \times \text{৩}}{\text{৫}}\) = ৯ ঘণ্টায়

\textbf{প্রশ্ন নম্বর-} ৪. ৪৮ জন শ্রমিক একটি কাজ ১২ দিনে শেষ করতে পারে। ৮দিনে
কাজটি শেষ করতে হলে, নতুন কতজন শ্রমিক লাগবে ? {[}প্রাথমিক প্রধান শিক্ষক
(চট্রগ্রাম)-২০০৮{]}

\begin{itemize}
\item
  (ক) ২২ জন
\item
  (খ) ২৪ জন
\item
  (গ) ২৫ জন
\item
  (ঘ) ২৬ জন
\end{itemize}

\textbf{\hl{উত্তরঃ (খ) ২৪ জন}}

\textbf{সমাধান:}

১২ দিনে একটি কাজ করতে পারে ৪৮ জন শ্রমিকে

১ দিনে একটি কাজ করতে পারে ১২×৪৮ = ৫৭৬ জন শ্রমিকে

৮ দিনে একটি কাজ করতে পারে ৫৭৬÷৮ = ৭২ জন শ্রমিকে

নতুন শ্রমিক নিয়োগ করতে হবে = ৭২-৪৮ = ২৪ জন

\textbf{শর্ট টেকনিক :}

প্রথম লোক × প্রথম সময় = দ্বিতীয় লোক × দ্বিতীয় সময়

সুতরাং দ্বিতীয় লোক =
\(\frac{\text{প্রথম}\ \text{লোক}\  \times \ \text{প্রথম}\ \text{সময়}}{\text{দ্বিতীয়}\ \text{লোক}}\) =
\(\frac{\text{৪৮} \times \text{১২}}{\text{৮}}\) = ৭২ জন

নতুন শ্রমিক নিয়োগ করতে হবে = ৭২-৪৮ = ২৪ জন

\textbf{প্রশ্ন নম্বর-} ৫. ১৪ জন লোক একটি কাজ ১৫ দিনে করতে পারে। ঐ কাজটি ১০
দিনে শেষ করতে হলে কতজন লোক নিয়োগ দিতে হবে ? {[}কারা তত্ত্বাবধায়ক -২০০৬{]}

\begin{itemize}
\item
  (ক) ২৭ জন
\item
  (খ) ২৪ জন
\item
  (গ) ২১ জন
\item
  (ঘ) ১৮ জন
\end{itemize}

\textbf{\hl{উত্তরঃ (গ) ২১ জন}}

\textbf{সমাধান:}

১৫ দিনে একটি কাজ করতে পারে ১৪ জন লোকে

১ দিনে একটি কাজ করতে পারে ১৪×১৫ = ২১০ জন লোকে

১০ দিনে একটি কাজ করতে পারে ২১০÷১০ = ২১ জন লোকে

\textbf{শর্ট টেকনিক :}

প্রথম লোক × প্রথম সময় = দ্বিতীয় লোক × দ্বিতীয় সময়

সুতরাং দ্বিতীয় লোক =
\(\frac{\text{প্রথম}\ \text{লোক}\  \times \ \text{প্রথম}\ \text{সময়}}{\text{দ্বিতীয়}\ \text{লোক}}\) =
\(\frac{\text{১৫} \times \text{১৪}}{\text{১০}}\) = ২১ জন

\textbf{প্রশ্ন নম্বর-} ৬. যে পরিমাণ খাদ্যে ১৫ জন লোকের ৪০ দিন চলে, ঐ পরিমাণ
খাদ্যে ২০জন লোকের কতদিন চলবে ? {[}প্রাথমিক সহকারী শিক্ষক(ইছামতি) -২০১০{]}

\begin{itemize}
\item
  (ক) ২৫ দিন
\item
  (খ) ৩০ দিন
\item
  (গ) ৩২ দিন
\item
  (ঘ) ৩৫ দিন
\end{itemize}

\textbf{\hl{উত্তরঃ (খ) ৩০ দিন}}

\textbf{সমাধান:} ১৫ জন লোকের চলে = ৪০ দিন

১ জন লোকের চলে = ৪০×১৫ = ৬০০ দিন

২০ জন লোকের চলে = ৬০০÷২০ = ৩০ দিন

বিঃ দ্রঃ মনে রাখতে হবে লোক ও সময় সংক্রান্ত অংকে ২য় লাইনে সবসময়ই গুন হয়।উত্তরে
যা চাওয়া হবে তা সবসময় লাইনের শেষ অথবা ডানদিকে রাখতে হবে।

\textbf{শর্ট টেকনিক :}

প্রথম লোক ×প্রথম সময় = দ্বিতীয় লোক ×দ্বিতীয় সময়

সুতরাং দ্বিতীয় সময় =
\(\frac{\text{প্রথম}\ \text{লোক}\  \times \ \text{প্রথম}\ \text{সময়}}{\text{দ্বিতীয়}\ \text{লোক}}\) =
\(\frac{\text{১৫} \times \text{৪০}}{\text{২০}}\) = ৩০ দিন

\subsubsection{পিপা / ট্যাংক/চৌবাচ্চা
সংক্রান্ত}\label{ux9aaux9aa-ux99fux9afux995ux99aux9acux99aux99a-ux9b8ux995ux9b0ux9a8ux9a4}

\textbf{প্রশ্ন নম্বর-}

১. একটি পানির ট্যাঙ্কে দুটি নল আছে। প্রথম নলটি খুলে দিলে ট্যাঙ্কটি ২০ ঘন্টায় পূর্ণ
হয়। দ্বিতীয় নল দ্বারা পূর্ণ ট্যাঙ্কটি ৩০ ঘণ্টায় খালি হয়। দুইটি নল একসঙ্গে খলে
দিলে খালি ট্যাঙ্ক কত সময়ে পূর্ণ হবে ? {[} শ্রম ও কর্মসংস্থান মন্ত্রণালয়ের সহকারী
শ্রম অফিসার- ২০০৩{]}

\begin{itemize}
\item
  (ক) ৪৮ ঘন্টা
\item
  (খ) ৪০ ঘন্টা
\item
  (গ) ৫০ ঘন্টা
\item
  (ঘ) ৬০ ঘন্টা
\end{itemize}

\textbf{\hl{উত্তরঃ (ঘ) ৬০ ঘন্টা}}

\textbf{সমাধান:}

প্রথম নলটি ২০ ঘন্টায় পূর্ণ করে ১টি ট্যাংক

প্রথম নলটি ১ ঘন্টায় পূর্ণ করে = \(\frac{\text{১}}{\text{২০}}\) অংশ

দ্বিতীয় নলটি ৩০ ঘন্টায় খালি করে ১টি ট্যাংক

দ্বিতীয় নলটি ১ ঘন্টায় খালি করে = \(\frac{\text{১}}{\text{৩০}}\) অংশ

দুটি নল এক সাথে ১ ঘন্টায় পূর্ণ করে = (\(\frac{\text{১}}{\text{২০}}\) - \(\frac{\text{১}}{\text{৩০}}\) )
অংশ = \(\frac{\text{১}}{\text{৬০}}\) অংশ

দুটি নল এক সাথে \(\frac{\text{১}}{\text{৬০}}\) অংশ পূর্ণ করে ১ ঘন্টায়

সুতরাং, দুটি নল এক সাথে ১ বা সম্পূর্ণ অংশ পূর্ণ করে ৬০ ঘন্টায়

\textbf{শর্ট টেকনিক:}

দুটি নলের ১ টি দ্বারা পূর্ণকরণ এবং অপরটি দ্বারা খালিকরণ হয় তখন -\/-ট্যাংক পূর্ণ
হতে প্রয়োজনীয় সময় T = mn÷(m-n)

=৩০×২০(৩০-২০) = ৬০০÷১০ = ৬০ ঘন্টা

এখানে, m = ২য় নল দ্বারা ব্যয়িত সময়,n = ১ম নল দ্বারা ব্যয়িত সময়।

\textbf{প্রশ্ন নম্বর-} ২. একটি চৌবাচ্চা একটি নল দ্বারা ১০ ঘণ্টায় পূর্ণ হয়। তাতে
একটি ছিদ্র থাকায় পূর্ণ হতে ১৫ ঘণ্টা লাগে। ছিদ্র দ্বারা চৌবাচ্চাটি খালি হতে কত
সময় লাগবে ? {[}সহকারী আবহাওয়াবিদ-২০০৪{]}

\begin{itemize}
\item
  (ক) ২০ ঘন্টা
\item
  (খ) ২৫ ঘন্টা
\item
  (গ) ৩০ ঘন্টা
\item
  (ঘ) ৩৫ ঘন্টা
\end{itemize}

\textbf{\hl{উত্তরঃ (গ) ৩০ ঘন্টা}}

\textbf{সমাধান:}

নলটি ১০ ঘন্টায় পূর্ণ করে ১টি চৌবাচ্চা

নলটি ১ ঘন্টায় পূর্ণ করে চৌবাচ্চাটির \(\frac{\text{১}}{\text{১০}}\) অংশ

কিন্তু ছিদ্র থাকায় ১ ঘন্টায় পূর্ণ হয় চৌবাচ্চাটির = \(\frac{\text{১}}{\text{১৫}}\) অংশ

ছিদ্র দ্বারা ১ ঘন্টায় খালি হয় = (\(\frac{\text{১}}{\text{১০}}\) - \(\frac{\text{১}}{\text{১৫}}\) ) অংশ
= \(\frac{\text{১}}{\text{৩০}}\) অংশ

ছিদ্র দ্বারা চৌবাচ্চাটির \(\frac{\text{১}}{\text{৩০}}\) অংশ খালি হতে সময় লাগে ১ ঘন্টা

সম্পূর্ণ অংশ খালি হতে সময় লাগে = ৩০ × ১ = ৩০ ঘন্টা

\textbf{প্রশ্ন নম্বর-} ৩. একটি পিপায় দুইটি নল সংযুক্ত আছে। প্রথম নলটি খুলে দিলে
খালি পিপাটি ২০ মিনিটে পূর্ণ হয়, দ্বিতীয় নলটি খুলে দিলে পরিপূর্ণ পিপাটি ৩০
মিনিটে খালি হয়। দুইটি নলই একসঙ্গে খুলে দিলে খালি পিপাটি কত সময়ে পূর্ণ হবে ?
{[}সমাজসেবা অধিদপ্তরের উপ-সহকারী পরিচালক পরীক্ষা-২০০৫{]}

\begin{itemize}
\item
  (ক) ৮০ মিনিটে
\item
  (খ) ৯০ মিনিটে
\item
  (গ) ৭০ মিনিটে
\item
  (ঘ) ৬০ মিনিটে
\end{itemize}

\textbf{\hl{উত্তরঃ (ঘ) ৬০ মিনিটে}}

\textbf{সমাধান:}

প্রথম নলটি ২০ মিনিটে পূর্ণ করে ১টি পিপা

প্রথম নলটি ১ মিনিটে পূর্ণ করে = \(\frac{\text{১}}{\text{২০}}\) অংশ

দ্বিতীয় নলটি ৩০ মিনিটে খালি করে ১টি পিপা

দ্বিতীয় নলটি ১ মিনিটে খালি করে = \(\frac{\text{১}}{\text{৩০}}\) অংশ

দুটি নল এক সাথে ১ মিনিটে পূর্ণ করে = (\(\frac{\text{১}}{\text{২০}}\) - \(\frac{\text{১}}{\text{৩০}}\) )
অংশ = \(\frac{\text{১}}{\text{৬০}}\) অংশ

দুটি নল এক সাথে \(\frac{\text{১}}{\text{৬০}}\) অংশ পূর্ণ করে ১ মিনিটে

সুতরাং, দুটি নল এক সাথে ১ বা সম্পূর্ণ অংশ পূর্ণ করে ৬০ মিনিটে

\textbf{শর্ট টেকনিক:}

দুটি নলের ১ টি দ্বারা পূর্ণকরণ এবং অপরটি দ্বারা খালিকরণ হয় তখন -\/-ট্যাংক পূর্ণ
হতে প্রয়োজনীয় সময় T = mn÷(m-n)

=৩০×২০(৩০-২০) = ৬০০÷১০ = ৬০ মিনিট

এখানে, m = ২য় নল দ্বারা ব্যয়িত সময়,n = ১ম নল দ্বারা ব্যয়িত সময়।

বি: দ্র: বড় সংখ্যাটিকে সব সময় প্রথম নল ধরতে হয় ।

\subsubsection{নৌকা ও স্রোত
সংক্রান্ত}\label{ux9a8ux995-ux993-ux9b8ux9b0ux9a4-ux9b8ux995ux9b0ux9a8ux9a4}

\textbf{প্রশ্ন নম্বর-} ১. একটি নৌকা স্রোতের অনুকূলে ঘণ্টায় ৮ কি. মি. এবং স্রোতের
প্রতিকূলে ৪ কিলোমিটার যায়। নৌকার বেগ কত ? {[}বিআরসিএর সিনিয়র অফিসার
পরীক্ষা-১৯৯৭{]}

\begin{itemize}
\item
  (ক) ২ কি.মি.
\item
  (খ) ৪ কি.মি.
\item
  (গ) ৬ কি.মি.
\item
  (ঘ) ৮ কি.মি.
\end{itemize}

\textbf{\hl{উত্তরঃ (গ) ৬ কি.মি.}}

\textbf{সমাধান:}

নৌকার বেগ+স্রোতের বেগ = ৮ কি.মি.

নৌকার বেগ-স্রোতের বেগ = ৪ কি.মি

(যোগ করে) ২×নৌকার বেগ = ১২ কি.মি.

নৌকার বেগ = ৬ কি.মি.

\textbf{প্রশ্ন নম্বর-} ২. কোন নৌকার গতি স্রোতের অনুকূলে ঘণ্টায় ১০ কিলোমিটার এবং
স্রোতের প্রতিকূলে ঘণ্টায় ২ কিলোমিটার। স্রোতের বেগ কত ?{[} বিআরসিএর সিনিয়র
অফিসার পরীক্ষা-১৯৯৮{]}

\begin{itemize}
\item
  (ক) ৪ কি.মি.
\item
  (খ) ৩ কি.মি.
\item
  (গ) ৮ কি.মি.
\item
  (ঘ) ১২ কি.মি.
\end{itemize}

\textbf{\hl{উত্তরঃ (ক) ৪ কি.মি.}}

\textbf{সমাধান:}

নৌকার বেগ+স্রোতের বেগ = ১০ কি.মি.

নৌকার বেগ-স্রোতের বেগ = ২ কি.মি

(বিয়োগ করে)২×স্রোতের বেগ = ৮ কি.মি.

স্রোতের বেগ = ৪ কি.মি.

\textbf{প্রশ্ন নম্বর-} ৩. নৌকা ও স্রোতের বেগ ঘণ্টায় যথাক্রমে ১০ ও ৫ কি. মি.।
নদী পথে ৪৫ কি. মি. দীর্ঘ পথ একবার অতিক্রম করে ফিরে আসতে কত ঘণ্টা সময় লাগবে
? {[}১২তম বিসিএস{]}

\begin{itemize}
\item
  (ক) ৯ ঘন্টা
\item
  (খ) ১২ ঘন্টা
\item
  (গ) ১০ ঘন্টা
\item
  (ঘ) ১৮ ঘন্টা
\end{itemize}

\textbf{\hl{উত্তরঃ (খ) ১২ ঘন্টা}}

\textbf{সমাধান:}

স্রোতের অনুকুলে,

নৌকা ও স্রোতের বেগ = ১০+৫ = ১৫ কি.মি.

৪৫ কি.মি. যেতে সময় লাগে = ৪৫÷১৫ = ৩ ঘন্টা

স্রোতের প্রতিকুলে,

নৌকা ও স্রোতের বেগ = ১০-৫ = ৫ কি.মি.

৪৫ কি.মি. আসতে সময় লাগে = ৪৫÷৫ = ৯ ঘন্টা

যেতে ও আসতে মোট সময় লাগে =৯+৩ = ১২ ঘন্টা

\subsubsection{বানরের বাঁশে ওঠা
সংক্রান্ত}\label{ux9acux9a8ux9b0ux9b0-ux9acux9b6-ux993ux9a0-ux9b8ux995ux9b0ux9a8ux9a4}

\textbf{প্রশ্ন নম্বর-} ১. একটি বানর ৯২ ফুট উঁচু একটা তৈলাক্ত বাঁশ বেয়ে উপরে উঠতে
লাগলো। বানরটি প্রথম মিনিটে ১২ ফুট ওঠে, কিন্তু দ্বিতীয় মিনিটে ৮ফুট নেমে বাঁশের
মাথায় উঠতে বানরটির কত মিনিট সময় লাগবে? {[}সাব রেজিস্ট্রার -১৯৯৯{]}

\begin{itemize}
\item
  (ক) ৪১
\item
  (খ) ৪৫
\item
  (গ) ৪৩
\item
  (ঘ) ৪২
\end{itemize}

\textbf{\hl{উত্তরঃ (ক) ৪১}}

\textbf{সমাধান:}

\textbf{শর্ট টেকনিক:}

যখন বানর তৈলাক্ত বাঁশের মাথায় নির্দিষ্ট সময় উঠে এবং নির্দিষ্ট সময়ে নামে তখন
প্রয়োজনীয় সময়

=
\(\frac{\text{মোট}\ \text{দৈর্ঘ্য}\  - \ \text{নির্দিষ্ট}\ \text{সময়}\ \text{যতটুকু}\ \text{উঠে}}{\text{নির্দিষ্ট}\ \text{সময়}\ \text{যতটুকু}\ \text{উঠে}\  - \ \text{নির্দিষ্ট}\ \text{সময়}\ \text{যতটুকু}\ \text{উঠে}}\)
× ২ + ১

= \(\frac{\text{৯২}\  - \ \text{১২}}{\text{১২}\  - \ \text{৮}}\) × ২ + ১

= \(\frac{\text{৮০}}{\text{৪}}\) × ২ + ১

= ৪১ মিনিট

\textbf{প্রশ্ন নম্বর-} ২. একটি বানর একটি তৈলাক্ত বাঁশ বেয়ে উপরে উঠতে লাগলো।
বানরটি যদি ১ মিনিটে ৫ মিটার ওঠে এবং পরবর্তী মিনিটে ১ মিটার নেমে পড়ে, তবে
২৫ মিটার উঁচু বাঁশের মাথায় উঠতে বানরের কত সময় লাগবে ? {[}পররাষ্ট্র ও অর্থ
মন্ত্রণালয়ের প্রশাসনিক কর্মকর্তা পরীক্ষা -২০০১{]}

\begin{itemize}
\item
  (ক) ৯ মিনিট
\item
  (খ) ১০ মিনিট
\item
  (গ) ১১ মিনিট
\item
  (ঘ)১২ মিনিট
\end{itemize}

\textbf{\hl{উত্তরঃ (গ) ১১ মিনিট}}

\textbf{সমাধান:}

\textbf{শর্ট টেকনিক:}

যখন বানর তৈলাক্ত বাঁশের মাথায় নির্দিষ্ট সময় উঠে এবং নির্দিষ্ট সময়ে নামে তখন
প্রয়োজনীয় সময়

=
\(\frac{\text{মোট}\ \text{দৈর্ঘ্য}\  - \ \text{নির্দিষ্ট}\ \text{সময়}\ \text{যতটুকু}\ \text{উঠে}}{\text{নির্দিষ্ট}\ \text{সময়}\ \text{যতটুকু}\ \text{উঠে}\  - \ \text{নির্দিষ্ট}\ \text{সময়}\ \text{যতটুকু}\ \text{উঠে}}\)
× ২ + ১

= \(\frac{\text{২৫}\  - \ \text{৫}}{\text{৫}\  - \ \text{১}}\) × ২ + ১

= \(\frac{\text{২০}}{\text{৪}}\) × ২ + ১

= ১১ মিনিট

\subsubsection{গতিবেগ দূরত্ব ও সময়
সংক্রান্ত}\label{ux997ux9a4ux9acux997-ux9a6ux9b0ux9a4ux9ac-ux993-ux9b8ux9aeux9df-ux9b8ux995ux9b0ux9a8ux9a4}

\textbf{প্রশ্ন নম্বর-} ১. দুটি বাস ঘণ্টায় ২০ মাইল বেগে একই সময়ে গাবতলী হতে
আরিচা রওয়ানা হল। সাভার পৌছার পর একটি বাস থেমে গেল। অপর বাসটি চলতে থাকল।
আধ ঘন্টা পর থেমে থাকা বাসটি ঘণ্টায় ২৫ মাইল বেগে চলতে থাকল। সাভার থেকে কত
দূরে বাস দুটি মিলিত হবে। {[}থানা শিক্ষা অফিসার-১৯৯৯{]}

\begin{itemize}
\item
  (ক) ৪০ মাইল
\item
  (খ) ৫০ মাইল
\item
  (গ) ৬০ মাইল
\item
  (ঘ) ৫৫ মাইল
\end{itemize}

\textbf{\hl{উত্তরঃ (খ) ৫০ মাইল}}

\textbf{সমাধান:}

\textbf{শর্ট টেকনিকঃ} নির্দিষ্ট স্থান থেকে একই দিকে চলমান দুটি বস্তুর ক্ষেত্রে-

দুরত্ব = (২য় দুরত্ব×১ম দুরত্ব)÷(২য় দুরত্ব-১ম দুরত্ব) ×সময়

=(২৫×২০)÷(২৫-২০)× \(\frac{\text{১}}{\text{২}}\)

=(৫০০÷৫)× \(\frac{\text{১}}{\text{২}}\)

=১০০× \(\frac{\text{১}}{\text{২}}\)

= ৫০ মাইল

\textbf{প্রশ্ন নম্বর-} ২. ক ঘণ্টায় ১০ কি. মি. এবং খ ঘণ্টায় ১৫ কি. মি বেগে একই
সময় একই স্থান থেকে রাজশাহীর পথে রওয়ানা হল। ক ১০.১০ মিনিটের সময় এবং খ ৯.৪০
মিনিটের সময় রাজশাহী পৌছল। রওয়ানা হওয়ার স্থান থেকে রাজশাহীর দূরত্ব কত কি.
মি. ?{[}১১তম বিসিএস{]}

\begin{itemize}
\item
  (ক) ২০ কি.মি.
\item
  (খ) ২৫ কি.মি.
\item
  (গ) ১৫ কি.মি.
\item
  (ঘ) ২৮ কি.মি.
\end{itemize}

\textbf{\hl{উত্তরঃ (গ) ১৫ কি.মি.}}

\textbf{সমাধান:}

শর্ট টেকনিকঃ

দুরত্ব = (২য় দুরত্ব×১ম দুরত্ব)÷(২য় দুরত্ব-১ম দুরত্ব) ×সময়

= (১৫×১০)÷(১৫-১০)×\(\frac{\text{১}}{\text{২}}\) {[}কারণ ১০.১০ থেকে ৯.৪০ =
\(\frac{\text{১}}{\text{২}}\) ঘন্টা{]}

= ১৫০÷৫× \(\frac{\text{১}}{\text{২}}\)

= ৩০× \(\frac{\text{১}}{\text{২}}\)

= ১৫ কি.মি.

\textbf{প্রশ্ন নম্বর-} ৩. ৪ কি. মি. / ঘণ্টা বেগে চললে কোন স্থানে পৌছাতে যে
সময় লাগে, ৫ কি. মি. / ঘণ্টা বেগে চললে তার চেয়ে \(\frac{\text{১}}{\text{২}}\) ঘণ্টা কম
লাগে। স্থানটির দূরত্ব কত ? {[}সংস্থাপন মন্ত্রণালয়ের প্রশাসনিক কর্মকর্তা পরীক্ষা
-২০০৭{]}

\begin{itemize}
\item
  (ক) ৮ কি.মি.
\item
  (খ) ১০ কি.মি.
\item
  (গ) ১২ কি.মি.
\item
  (ঘ) ১১ কি.মি.
\end{itemize}

\textbf{\hl{উত্তরঃ (খ) ১০ কি.মি.}}

\textbf{সমাধান:}

শর্ট টেকনিকঃ

দুরত্ব = (২য় দুরত্ব×১ম দুরত্ব)÷(২য় দুরত্ব-১ম দুরত্ব) ×সময়

= (৫×৪)÷(৫-৪)×\(\frac{\text{১}}{\text{২}}\)

= ২০÷১× \(\frac{\text{১}}{\text{২}}\)

= ২০× \(\frac{\text{১}}{\text{২}}\)

= ১০ কি.মি.

\textbf{বিকল্প পদ্ধতি:}

মনে করি, স্থানটির দুরত্ব `ক' কিমি

প্রশ্নমতে, \(\frac{\text{ক}}{\text{৪}}\) - \(\frac{\text{ক}}{\text{৫}}\) = \(\frac{\text{১}}{\text{২}}\)

বা, \(\frac{\text{ক}}{\text{২০}}\) = \(\frac{\text{১}}{\text{২}}\)

বা, `ক' = ১০ কিমি

\subsubsection{ট্রেন ও গতিবেগ
সংক্রান্ত}\label{ux99fux9b0ux9a8-ux993-ux997ux9a4ux9acux997-ux9b8ux995ux9b0ux9a8ux9a4}

\textbf{প্রশ্ন নম্বর-} ১. রেল লাইনের পাশে একটি তালগাছ রয়েছে। ঘণ্টায় ৪৫
কিঃমিঃ বেগে ধাবমান ১৫০ মিঃ লম্বা একটি ট্রেন কত সময়ে ঐ তালগাছটি অতিক্রম করবে
?{[}প্রাথমিক সহকারী শিক্ষক (সিলেট)-২০০৮{]}

\begin{itemize}
\item
  (ক) ১১ সেকেন্ড
\item
  (খ) ১২ সেকেন্ড
\item
  (গ) ১৩ সেকেন্ড
\item
  (ঘ) ১৪ সেকেন্ড
\end{itemize}

\textbf{\hl{উত্তরঃ (খ) ১২ সেকেন্ড}}

\textbf{সমাধান:}

\textbf{\hl{উত্তরঃ খ}}

শর্ট টেকনিক: কোন স্থান থেকে যাত্রা শুরু করার পর নির্দিষ্ট স্থান অতিক্রান্ত হওয়ার
ক্ষেত্রে-\/-

সময় = মোট অতিক্রান্ত দুরত্ব÷গতিবেগ

এখানে অতিক্রান্ত দুরত্ব=১৫০ মিটার

গতিবেগ=(৪৫×১০০০)÷৩৬০০ = ১২.৫ মিটার/সেকেন্ড

সময় = মোট অতিক্রান্ত দুরত্ব÷গতিবেগ

=১৫০÷ ১২.৫ = ১২ সেকেন্ড

\textbf{প্রশ্ন নম্বর-} ২. ১১০ মিটার ও ১৪০ মিটার দীর্ঘ দুটি ট্রেন যথাক্রমে ঘণ্টায়
৪৫ কি.মি ও ৫৫ কি.মি. বেগে বিপরীত দিক থেকে কত সময় পর সমান্তরালভাবে অতিক্রম
করবে ? {[}বিআরডিবি-২০১২{]}

\begin{itemize}
\item
  (ক) ৯ সেকেন্ড
\item
  (খ) ১৬ সেকেন্ড
\item
  (গ) ২০ সেকেন্ড
\item
  (ঘ) ৩০ সেকেন্ড
\end{itemize}

\textbf{\hl{উত্তরঃ (ক) ৯ সেকেন্ড}}

\textbf{সমাধান:}

ট্রেনের মোট দৈর্ঘ্য = ১১০+১৪০ = ২৫০ মিটার

ট্রেন দুটির গতিবেগ = ৪৫+৫৫ = ১০০ কি. মি.

১০০×১০০০ মিটার যায় = ৬০×৬০ সেকেন্ডে

২৫০ মিটার যায় = (৬০×৬০×২৫০)÷(১০০×১০০০)

= ৯ সেকেন্ডে

\textbf{প্রশ্ন নম্বর-} ৩. একটি গাড়ি ঘণ্টায় ২৫ মাইল বেগে ঢাকা থেকে চট্টগ্রাম
অভিমুখে রওয়ানা হল। আরেকটি গাড়ি ঘণ্টায় ১৫ মাইল বেগে চট্টগ্রাম থেকে ঢাকা
অভিমুখে রওয়ানা হল । ঢাকা ও চট্টগ্রামের দূরত্ব ২০০ মাইল হলে গাড়ি দুটি কখন
মুখোমুখি হবে ? {[}মাদকদ্রব্য অধিদপ্তরের সহকারী পরিচালক-১৯৯৯{]}

\begin{itemize}
\item
  (ক) ৪ ঘন্টা পর
\item
  (খ) ৫ ঘন্টা পর
\item
  (গ) ৬ ঘন্টা পর
\item
  (ঘ) ৮ ঘন্টা পর
\end{itemize}

\textbf{\hl{উত্তরঃ (খ) ৫ ঘন্টা পর}}

\textbf{সমাধান:}

শর্ট টেকনিকঃ যখন দুটি গাড়ী ভিন্ন গতিতে চলে তখন-\/-

সময় = মোট অতিক্রান্ত দুরত্ব÷গাড়ী দুটির বেগের যোগফল

= ২০০÷(২৫+১৫) = ২০০÷৪০ = ৫ ঘন্টা

গণিতের ভীতি দূর করে যেকোনো প্রতিযোগিতামূলক পরীক্ষায় সফলতার চাবিকাঠি হলো
নিয়মিত অনুশীলন এবং মৌলিক ধারণায় স্বচ্ছতা। বিশেষ করে `ঐকিক নিয়ম' এমন একটি
অধ্যায়, যা শুধু পরীক্ষার খাতায় নম্বর তুলতেই নয়, বরং আমাদের প্রাত্যহিক জীবনের
নানা হিসাব-নিকাশ ও সমস্যার বাস্তবসম্মত সমাধানেও অতীব জরুরি।

অযথা শর্টকাট বা মুখস্থ বিদ্যার ওপর নির্ভর না করে প্রতিটি অংকের মূল রহস্য ও নিয়ম
বোঝার চেষ্টা করুন। এই নিবন্ধে আলোচিত বিভিন্ন ক্যাটাগরির বিগত সালের প্রশ্ন ও
সমাধানগুলো আপনার প্রস্তুতিকে এক ধাপ এগিয়ে নিতে এবং আপনার আত্মবিশ্বাস বাড়াতে
কার্যকর ভূমিকা রাখবে বলে আমাদের বিশ্বাস। নিয়মিত চর্চা অব্যাহত রাখুন, সফলতা
আপনার জন্যই নির্ধারিত!

\section{55. ঐকিক নিয়ম MCQ প্রশ্ন উত্তর
অনুশীলন}\label{ux990ux995ux995-ux9a8ux9dfux9ae-mcq-ux9aaux9b0ux9b6ux9a8-ux989ux9a4ux9a4ux9b0-ux985ux9a8ux9b6ux9b2ux9a8}

\textbf{১০ জন লোক একটি কাজ ২০ দিন করতে পারে। ৮ জন লোকের ঐ কাজটি করতে
কতদিন লাগবে?}

\textbf{\hl{ক) ২৫ দিন ✔}}

খ) ২৪ দিন

গ) ৩০ দিন

ঘ) ১৬ দিন

\textbf{\hl{✔ সঠিক উত্তর: ক) ২৫ দিন}}

\begin{quote}
\textbf{ব্যাখ্যা:} \emph{১০ জন লোক একটি কাজ করতে পারে = ২০ দিনে}

\emph{∴ ১ জন লোক একটি কাজ করতে পারে = (২০ × ১০) = ২০০ দিনে}

\emph{∴ ৮ জন লোক একটি কাজ করতে পারে = (২০০/৮) = ২৫ দিনে।}

\emph{শর্ট টেকনিক:}

\emph{দ্বিতীয় সময় = (প্রথম লোক×প্রথম সময়/দ্বিতীয় লোক)}

\emph{= (১০×২০/৮)}

\emph{= ২৫ দিনে।}
\end{quote}

\textbf{৫৬ জন শ্রমিক একটি কাজ ২১ দিনে শেষ করতে পারে। ১৪ দিনে কাজটি শেষ
করতে হলে নতুন কত জন শ্রমিক নিয়োগ করতে হবে?}

ক) ২৪ জন

খ) ২৬ জন

\textbf{\hl{গ) ২৮ জন ✔}}

ঘ) ৩০ জন

\textbf{\hl{✔ সঠিক উত্তর: গ) ২৮ জন}}

\begin{quote}
\textbf{ব্যাখ্যা:} \emph{২১ দিনে একটি কাজ করতে পারে = ৫৬ জন শ্রমিক}

\emph{∴ ১ দিনে একটি কাজ করতে পারে = (৫৬ × ২১) = ১১৭৬ জন শ্রমিক}

\emph{∴ ১৪ দিনে একটি কাজ করতে পারে = (১১৭৬/১৪) = ৮৪ জন শ্রমিক}

\emph{∴ নতুন শ্রমিক নিয়োগ করতে হবে = (৮৪ − ৫৬) = ২৮ জন।}

\emph{শর্ট টেকনিক:}

\emph{দ্বিতীয় লোক = (প্রথম লোক×প্রথম সময়/দ্বিতীয় সময়)}

\emph{= (৫৬×২১/১৪)}

\emph{= ৮৪ জন}

\emph{∴ নতুন শ্রমিক নিয়োগ করতে হবে = (৮৪ − ৫৬) = ২৮ জন।}
\end{quote}

\textbf{১৫ জন লোক একটি কাজ শেষ করে ৩ ঘণ্টায়। ৫ জন লোকের ঐ কাজ শেষ করতে কত
সময় লাগবে?}

ক) ৫ ঘন্টা

খ) ৭.৫ ঘন্টা

\textbf{\hl{গ) ৯ ঘন্টা ✔}}

ঘ) ৪ ঘন্টা

\textbf{\hl{✔ সঠিক উত্তর: গ) ৯ ঘন্টা}}

\begin{quote}
\textbf{ব্যাখ্যা:} \emph{১৫ জন লোক একটি কাজ করতে পারে = ৩ ঘণ্টায়}

\emph{∴ ১ জন লোক একটি কাজ করতে পারে = (১৫ × ৩) = ৪৫ ঘণ্টায়}

\emph{∴ ৫ জন লোক একটি কাজ করতে পারে = (৪৫/৫) = ৯ ঘণ্টায়।}

\emph{শর্ট টেকনিক:}

\emph{দ্বিতীয় সময় = (প্রথম লোক×প্রথম সময়/দ্বিতীয় লোক)}

\emph{= (১৫×৩/৫)}

\emph{= ৯ ঘণ্টায়।}
\end{quote}

\textbf{৪৮ জন শ্রমিক একটি কাজ ১২ দিনে শেষ করতে পারে। ৮দিনে কাজটি শেষ করতে
হলে, নতুন কতজন শ্রমিক লাগবে?}

ক) ২২ জন

\textbf{\hl{খ) ২৪ জন ✔}}

গ) ২৫ জন

ঘ) ২৬ জন

\textbf{\hl{✔ সঠিক উত্তর: খ) ২৪ জন}}

\begin{quote}
\textbf{ব্যাখ্যা:} \emph{১২ দিনে একটি কাজ করতে পারে = ৪৮ জন শ্রমিক}

\emph{∴ ১ দিনে একটি কাজ করতে পারে = (১২ × ৪৮) = ৫৭৬ জন শ্রমিক}

\emph{∴ ৮ দিনে একটি কাজ করতে পারে = (৫৭৬/৮) = ৭২ জন শ্রমিক}

\emph{∴ নতুন শ্রমিক নিয়োগ করতে হবে = (৭২ − ৪৮) = ২৪ জন।}

\emph{শর্ট টেকনিক:}

\emph{দ্বিতীয় লোক = (প্রথম লোক×প্রথম সময়/দ্বিতীয় সময়)}

\emph{= (৪৮×১২/৮)}

\emph{= ৭২ জন}

\emph{∴ নতুন শ্রমিক নিয়োগ করতে হবে = (৭২ − ৪৮) = ২৪ জন।}
\end{quote}

\textbf{১৪ জন লোক একটি কাজ ১৫ দিনে করতে পারে। ঐ কাজটি ১০ দিনে শেষ করতে
হলে কতজন লোক নিয়োগ দিতে হবে?}

ক) ২৭ জন

খ) ২৪ জন

\textbf{\hl{গ) ২১ জন ✔}}

ঘ) ১৮ জন

\textbf{\hl{✔ সঠিক উত্তর: গ) ২১ জন}}

\begin{quote}
\textbf{ব্যাখ্যা:} \emph{১৫ দিনে একটি কাজ করতে পারে = ১৪ জন লোক}

\emph{∴ ১ দিনে একটি কাজ করতে পারে = (১৪ × ১৫) = ২১০ জন লোক}

\emph{∴ ১০ দিনে একটি কাজ করতে পারে = (২১০/১০) = ২১ জন লোক।}

\emph{শর্ট টেকনিক:}

\emph{দ্বিতীয় লোক = (প্রথম লোক×প্রথম সময়/দ্বিতীয় সময়)}

\emph{= (১৫×১৪/১০)}

\emph{= ২১ জন।}
\end{quote}

\textbf{যে পরিমাণ খাদ্যে ১৫ জন লোকের ৪০ দিন চলে, ঐ পরিমাণ খাদ্যে ২০জন
লোকের কতদিন চলবে?}

ক) ২৫ দিন

\textbf{\hl{খ) ৩০ দিন ✔}}

গ) ৩২ দিন

ঘ) ৩৫ দিন

\textbf{\hl{✔ সঠিক উত্তর: খ) ৩০ দিন}}

\begin{quote}
\textbf{ব্যাখ্যা:} \emph{১৫ জন লোকের চলে = ৪০ দিন}

\emph{∴ ১ জন লোকের চলে = (৪০ × ১৫) = ৬০০ দিন}

\emph{∴ ২০ জন লোকের চলে = (৬০০/২০) = ৩০ দিন।}

\emph{শর্ট টেকনিক:}

\emph{দ্বিতীয় সময় = (প্রথম লোক×প্রথম সময়/দ্বিতীয় লোক)}

\emph{= (১৫×৪০/২০)}

\emph{= ৩০ দিন।}
\end{quote}

\textbf{একটি পুকুর খনন করতে ৩০০ জন লোকের ২৫ দিন লাগে। পুকুরটি ১দিনে খনন করতে
কত জন লোকের দরকার হবে?}

ক) ৭০০০ জন

খ) ৭২৫০ জন

\textbf{\hl{গ) ৭৫০০ জন ✔}}

ঘ) ৮০০০ জন

\textbf{\hl{✔ সঠিক উত্তর: গ) ৭৫০০ জন}}

\begin{quote}
\textbf{ব্যাখ্যা:} \emph{২৫ দিন লাগে = ৩০০ জন লোকের}

\emph{∴ ১ দিন লাগে = (৩০০ × ২৫) = ৭৫০০ জন লোকের।}

\emph{শর্ট টেকনিক:}

\emph{দ্বিতীয় লোক = (প্রথম লোক×প্রথম সময়/দ্বিতীয় সময়)}

\emph{= (৩০০×২৫/১)}

\emph{= ৭৫০০ জন লোকের।}
\end{quote}

\textbf{একটি রাস্তা মেরামত করতে ১০ জন শ্রমিকের ১৬ দিন লাগলে ৮জন শ্রমিকের কত
দিন লাগবে?}

ক) ১৬ দিন

খ) ১৮ দিন

\textbf{\hl{গ) ২০ দিন ✔}}

ঘ) ২৪ দিন

\textbf{\hl{✔ সঠিক উত্তর: গ) ২০ দিন}}

\begin{quote}
\textbf{ব্যাখ্যা:} \emph{১০ জন শ্রমিকের লাগে = ১৬ দিন}

\emph{∴ ১ জন শ্রমিকের লাগে = (১০ × ১৬) = ১৬০ দিন}

\emph{∴ ৮ জন শ্রমিকের লাগে = (১৬০/৮) = ২০ দিন।}

\emph{শর্ট টেকনিক:}

\emph{দ্বিতীয় সময় = (প্রথম লোক×প্রথম সময়/দ্বিতীয় লোক)}

\emph{= (১০×১৬/৮)}

\emph{= ২০ দিন।}
\end{quote}

\textbf{২০ জন লোক ১টি কাজ ২০ দিনে করতে পারে। ১দিনে ঐ কাজ কতজন লোক করতে
পারে?}

ক) ১০০ জন

খ) ২০০ জন

গ) ৩০০ জন

\textbf{\hl{ঘ) ৪০০ জন ✔}}

\textbf{\hl{✔ সঠিক উত্তর: ঘ) ৪০০ জন}}

\begin{quote}
\textbf{ব্যাখ্যা:} \emph{২০ দিনে করতে পারে = ২০ জন লোকে}

\emph{∴ ১ দিনে করতে পারে = (২০ × ২০) = ৪০০ জন লোকে।}

\emph{শর্ট টেকনিক:}

\emph{দ্বিতীয় লোক = (প্রথম লোক×প্রথম সময়/দ্বিতীয় সময়)}

\emph{= (২০×২০/১)}

\emph{= ৪০০ জন লোকে।}
\end{quote}

\textbf{৬০ জন লোক কোন কাজ ১৮ দিনে করতে পারে। উক্ত কাজ ৩৬ জন লোক কতদিনে
সম্পন্ন করতে পারবে?}

ক) ২৭ দিনে

খ) ২৮ দিনে

গ) ২৯ দিনে

\textbf{\hl{ঘ) ৩০ দিনে ✔}}

\textbf{\hl{✔ সঠিক উত্তর: ঘ) ৩০ দিনে}}

\begin{quote}
\textbf{ব্যাখ্যা:} \emph{৬০ জন লোক কাজটি করতে পারে = ১৮ দিনে}

\emph{∴ ১ জন লোক কাজটি করতে পারে = (১৮ × ৬০) = ১০৮০ দিনে}

\emph{∴ ৩৬ জন লোক কাজটি করতে পারে = (১০৮০/৩৬) = ৩০ দিনে।}

\emph{শর্ট টেকনিক:}

\emph{দ্বিতীয় সময় = (প্রথম লোক×প্রথম সময়/দ্বিতীয় লোক)}

\emph{= (১৮×৬০/৩৬)}

\emph{= ৩০ দিনে।}
\end{quote}

\textbf{১৫ জনের কোন কাজের অর্ধেক করতে ২০ দিন লাগে, কত দিনে ২০ জন লোক পুরো
কাজটি শেষ করতে পারবে?}

ক) ২০০

খ) ১৫

\textbf{\hl{গ) ৩০ ✔}}

ঘ) ৪০

\textbf{\hl{✔ সঠিক উত্তর: গ) ৩০}}

\begin{quote}
\textbf{ব্যাখ্যা:} \emph{১৫ জনের কোন কাজের অর্ধেক কাজ করতে যেহেতু ২০ দিন
লাগে সেহেতু সম্পূর্ণ কাজ করতে ৪০ দিন সময় লাগবে।}

\emph{১৫ জনের কাজটি করতে লাগে = ৪০ দিন}

\emph{∴ ১ জনের কাজটি করতে লাগে = (৪০ × ১৫) = ৬০০ দিন}

\emph{∴ ২০ জনের কাজটি করতে লাগে = (৬০০/২০) = ৩০ দিন।}
\end{quote}

\textbf{২ জন পুরুষ বা ৩ জন বালক যে কাজ ১৫ দিনে সম্পন্ন করতে পারে, ৪ জন পুরুষ
এবং ৯ জন বালক তার দ্বিগুণ কাজ কত দিনে শেষ করতে পারবে?}

ক) ৪ দিনে

খ) ১২ দিনে

\textbf{\hl{গ) ৬ দিনে ✔}}

ঘ) ৫ দিনে

\textbf{\hl{✔ সঠিক উত্তর: গ) ৬ দিনে}}

\begin{quote}
\textbf{ব্যাখ্যা:} \emph{২ জন পুরুষের কাজ = ৩ জন বালকের কাজ}

\emph{∴ ১ জন পুরুষের কাজ = (৩/২) জন বালকের কাজ}

\emph{∴ ৪ জন পুরুষের কাজ = (৩×৪/২) = ৬ জন বালকের কাজ}

\emph{এখন, ৪ জন পুরুষ + ৯ জন বালক = (৬ জন বালক + ৯ জন বালক) = ১৫ জন বালক}

\emph{৩ জন বালক কাজটি করতে পারে = ১৫ দিনে}

\emph{∴ ১ জন বালক কাজটি করতে পারে = (১৫ × ৩) = ৪৫ দিনে}

\emph{∴ ১৫ জন বালক কাজটি করতে পারে = (৪৫/১৫) = ৩ দিনে}

\emph{∴ তার দ্বিগুণ কাজ করতে সময় লাগবে = (৩ × ২) = ৬ দিনে।}
\end{quote}

\textbf{৪ জন পুরুষ বা ৬ জন বালক যে কাজ ৩০ দিনে করতে পারে, ৭ জন পুরুষ ও ১২ জন
বালক ঐ কাজ কত দিনে করতে পারবে?}

\textbf{\hl{ক) ৮ দিনে ✔}}

খ) ৯ দিনে

গ) ১০ দিনে

ঘ) ১২ দিনে

\textbf{\hl{✔ সঠিক উত্তর: ক) ৮ দিনে}}

\begin{quote}
\textbf{ব্যাখ্যা:} \emph{৪ জন পুরুষের কাজ = ৬ জন বালকের কাজ}

\emph{∴ ১ জন পুরুষের কাজ = (৬/৪) জন বালকের কাজ}

\emph{∴ ৭ জন পুরুষের কাজ = (৬×৭/৪) = (২১/২) জন বালকের কাজ}

\emph{এখন, ৭ জন পুরুষ + ১২ জন বালক = ( (২১/২) + ১২) জন বালক = (৪৫/২) জন
বালক}

\emph{৬ জন বালক কাজটি করতে পারে = ৩০ দিনে}

\emph{∴ ১ জন বালক কাজটি করতে পারে = (৩০ × ৬) = ১৮০ দিনে}

\emph{∴ (৪৫/২) জন বালক কাজটি করতে পারে = (১৮০×২/৪৫) = ৮ দিনে।}
\end{quote}

\textbf{২ জন পুরুষ বা ৩ জন স্ত্রীলোক একটি কাজ ২৬ দিনে শেষ করতে পারলে ৪ জন
পুরুষ ও ৭ জন স্ত্রীলোক একত্রে কাজটি কত দিনে শেষ করতে পারবে?}

\textbf{\hl{ক) ৬ দিনে ✔}}

খ) ৭ দিনে

গ) ৯ দিনে

ঘ) ১০ দিনে

\textbf{\hl{✔ সঠিক উত্তর: ক) ৬ দিনে}}

\begin{quote}
\textbf{ব্যাখ্যা:} \emph{২ জন পুরুষের কাজ = ৩ জন স্ত্রীলোকের কাজ}

\emph{∴ ১ জন পুরুষের কাজ = (৩/২) জন স্ত্রীলোকের কাজ}

\emph{∴ ৪ জন পুরুষের কাজ = (৩×৪/২) = ৬ জন স্ত্রীলোকের কাজ}

\emph{এখন, ৪ জন পুরুষ + ৭ জন স্ত্রীলোক = (৬ জন স্ত্রীলোক + ৭ জন স্ত্রীলোক) = ১৩
জন স্ত্রীলোক}

\emph{৩ জন স্ত্রীলোক কাজটি করতে পারে = ২৬ দিনে}

\emph{∴ ১ জন স্ত্রীলোক কাজটি করতে পারে = (২৬ × ৩) = ৭৮ দিনে}

\emph{∴ ১৩ জন স্ত্রীলোক কাজটি করতে পারে = (৭৮/১৩) = ৬ দিনে।}
\end{quote}

\textbf{৬ জন স্ত্রীলোক অথবা ৮ জন বালক একটি কাজ ১২ দিনে শেষ করতে পারে। ৩ জন
স্ত্রীলোক ও ১২ জন বালক ঐ কাজটি কত দিনে করতে পারবে?}

\textbf{\hl{ক) ৬ দিনে ✔}}

খ) ৮ দিনে

গ) ৯ দিনে

ঘ) ১০ দিনে

\textbf{\hl{✔ সঠিক উত্তর: ক) ৬ দিনে}}

\begin{quote}
\textbf{ব্যাখ্যা:} \emph{৬ জন স্ত্রীলোকের কাজ = ৮ জন বালকের কাজ}

\emph{∴ ১ জন স্ত্রীলোকের কাজ = (৮/৬) জন বালকের কাজ}

\emph{∴ ৩ জন স্ত্রীলোকের কাজ = (৮×৩/৬) = ৪ জন বালকের কাজ}

\emph{এখন, ৩ জন স্ত্রীলোক + ১২ জন বালক = (৪ জন বালক + ১২ জন বালক) = ১৬ জন
বালক}

\emph{৮ জন বালক কাজটি করতে পারে = ১২ দিনে}

\emph{∴ ১ জন বালক কাজটি করতে পারে = (১২ × ৮) = ৯৬ দিনে}

\emph{∴ ১৬ জন বালক কাজটি করতে পারে = (৯৬/১৬) = ৬ দিনে।}
\end{quote}

\textbf{৪ জন পুরুষ বা ৮ জন স্ত্রীলোক একটি কাজ ৯ দিনে করতে পারে। ৬ জন পুরুষ এবং
৬ জন স্ত্রীলোক সেই কাজ কত দিনে করতে পারবে?}

\textbf{\hl{ক) ৪ দিনে ✔}}

খ) ৬ দিনে

গ) ৫ দিনে

ঘ) ৩ দিনে

\textbf{\hl{✔ সঠিক উত্তর: ক) ৪ দিনে}}

\begin{quote}
\textbf{ব্যাখ্যা:} \emph{৪ জন পুরুষের কাজ = ৮ জন স্ত্রীলোকের কাজ}

\emph{∴ ১ জন পুরুষের কাজ = (৮/৪) জন স্ত্রীলোকের কাজ}

\emph{∴ ৬ জন পুরুষের কাজ = (৮×৬/৪) = ১২ জন স্ত্রীলোকের কাজ}

\emph{এখন, ৬ জন পুরুষ + ৬ জন স্ত্রীলোক = (১২ জন স্ত্রীলোক + ৬ জন স্ত্রীলোক) = ১৮
জন স্ত্রীলোক}

\emph{৮ জন স্ত্রীলোক কাজটি করতে পারে = ৯ দিনে}

\emph{∴ ১ জন স্ত্রীলোক কাজটি করতে পারে = (৯ × ৮) = ৭২ দিনে}

\emph{∴ ১৮ জন স্ত্রীলোক কাজটি করতে পারে = (৭২/১৮) = ৪ দিনে।}
\end{quote}

\textbf{২০ জনে একটি কাজের অর্ধেক করতে পারে ৬ দিনে ঐ কাজটি করতে ৫ জনের
কতদিন লাগবে?}

\textbf{\hl{ক) ৪৮ দিন ✔}}

খ) ২৪ দিন

গ) ১৬ দিন

ঘ) ৭২ দিন

\textbf{\hl{✔ সঠিক উত্তর: ক) ৪৮ দিন}}

\begin{quote}
\textbf{ব্যাখ্যা:} \emph{২০ জনের কোন কাজের অর্ধেক কাজ করতে যেহেতু ৬ দিন লাগে
সেহেতু সম্পূর্ণ কাজ করতে ১২ দিন সময় লাগবে।}

\emph{২০ জনের কাজটি করতে লাগে = ১২ দিন}

\emph{∴ ১ জনের কাজটি করতে লাগে = (১২ × ২০) = ২৪০ দিন}

\emph{∴ ৫ জনের কাজটি করতে লাগে = (২৪০/৫) = ৪৮ দিন।}
\end{quote}

\textbf{একটি কাজ ক একা ৬ দিনে এবং খ একা ১২ দিনে শেষ করলে ক ও খ একত্রে
কাজটি কত দিনে শেষ করতে পারবে?}

\textbf{\hl{ক) ৪ দিনে ✔}}

খ) ৪ (১/২) দিনে

গ) ৫ দিনে

ঘ) ৫ (১/২) দিনে

\textbf{\hl{✔ সঠিক উত্তর: ক) ৪ দিনে}}

\begin{quote}
\textbf{ব্যাখ্যা:} \emph{\textquotesingle ক\textquotesingle{} ১ দিনে করে =
(১/৬) অংশ}

\emph{\textquotesingle খ\textquotesingle{} ১ দিনে করে = (১/১২) অংশ}

\emph{\textquotesingle ক\textquotesingle{} ও
\textquotesingle খ\textquotesingle{} একত্রে ১ দিনে করে = ( (১/৬) + (১/১২)
) = (২+১/১২) = (৩/১২) = (১/৪) অংশ}

\emph{∴ \textquotesingle ক\textquotesingle{} ও
\textquotesingle খ\textquotesingle{} একত্রে কাজটি করতে পারবে ৪ দিনে।}
\end{quote}

\textbf{স্বপন ও বকুল একটি কাজ পৃথকভাবে যথাক্রমে ২০ দিনে এবং ৩০ দিনে করতে
পারে। উভয়ে এক সঙ্গে কাজটি কত দিনে করতে পারবে?}

ক) ১০ দিনে

\textbf{\hl{খ) ১২ দিনে ✔}}

গ) ১৪ দিনে

ঘ) ১৫ দিনে

\textbf{\hl{✔ সঠিক উত্তর: খ) ১২ দিনে}}

\begin{quote}
\textbf{ব্যাখ্যা:} \emph{\textquotesingle স্বপন\textquotesingle{} ১ দিনে
করে = (১/২০) অংশ}

\emph{\textquotesingle বকুল\textquotesingle{} ১ দিনে করে = (১/৩০) অংশ}

\emph{\textquotesingle স্বপন\textquotesingle{} ও
\textquotesingle বকুল\textquotesingle{} একত্রে ১ দিনে করে = ( (১/২০) +
(১/৩০) ) = (৩+২/৬০) = (৫/৬০) = (১/১২) অংশ}

\emph{∴ \textquotesingle স্বপন\textquotesingle{} ও
\textquotesingle বকুল\textquotesingle{} একত্রে কাজটি করতে পারবে ১২ দিনে।}
\end{quote}

\textbf{ক এবং খ একত্রে মিলে একটি কাজ ১২ দিনে করতে পারে। ক একা কাজটি ২০
দিনে করতে পারে, খ একা কাজটি করতে পারবে?}

ক) ২৫ দিনে

\textbf{\hl{খ) ৩০ দিনে ✔}}

গ) ৩৫ দিনে

ঘ) ৪০ দিনে

\textbf{\hl{✔ সঠিক উত্তর: খ) ৩০ দিনে}}

\begin{quote}
\textbf{ব্যাখ্যা:} \emph{ক ও খ একত্রে ১ দিনে করে = (১/১২) অংশ}

\emph{ক একা ১ দিনে করে = (১/২০) অংশ}

\emph{∴ খ একা ১ দিনে করে = ( (১/১২) − (১/২০) ) = (৫−৩/৬০) = (২/৬০) =
(১/৩০) অংশ}

\emph{∴ খ একা কাজটি করতে পারবে ৩০ দিনে।}
\end{quote}

\textbf{দুই ব্যক্তি একটি কাজ ৮ দিনে করতে পারে। প্রথম ব্যক্তি একা কাজটি ১২ দিনে
করতে পারে। দ্বিতীয় ব্যক্তি একা কাজটি কত দিনে করতে পারবে?}

ক) ২০ দিন

খ) ২২ দিন

\textbf{\hl{গ) ২৪ দিন ✔}}

ঘ) ২৬ দিন

\textbf{\hl{✔ সঠিক উত্তর: গ) ২৪ দিন}}

\begin{quote}
\textbf{ব্যাখ্যা:} \emph{দুই ব্যক্তি একত্রে কাজটির ১ দিনে করতে পারে = (১/৮)
অংশ}

\emph{প্রথম ব্যক্তি একা কাজটির ১ দিনে করতে পারে = (১/১২) অংশ}

\emph{∴ দ্বিতীয় ব্যক্তি একা কাজটির ১ দিনে করতে পারে = ( (১/৮) − (১/১২) ) =
(৩−২/২৪) = (১/২৪) অংশ}

\emph{∴ দ্বিতীয় ব্যক্তি একা কাজটি ২৪ দিনে করতে পারবে।}
\end{quote}

\textbf{ক ও খ একটি কাজ ১২ দিনে সম্পন্ন করে। ক, খ ও গ ঐ কাজটি ৮ দিনে সম্পন্ন
করে। গ এর ঐ কাজটি একা করতে কত দিন সময় লাগবে?}

ক) ১৪

খ) ১৬

গ) ২০

\textbf{\hl{ঘ) ২৪ ✔}}

\textbf{\hl{✔ সঠিক উত্তর: ঘ) ২৪}}

\begin{quote}
\textbf{ব্যাখ্যা:} \emph{ক ও খ একত্রে কাজটির ১ দিনে করতে পারে = (১/১২) অংশ}

\emph{ক, খ ও গ একত্রে ঐ কাজটির ১ দিনে করতে পারে = (১/৮) অংশ}

\emph{∴ \textquotesingle গ\textquotesingle{} একা কাজটির ১ দিনে করতে পারে
= ( (১/৮) − (১/১২) ) = (৩−২/২৪) = (১/২৪) অংশ}

\emph{∴ \textquotesingle গ\textquotesingle{} একা কাজটি ২৪ দিনে করতে
পারবে।}
\end{quote}

\textbf{যদি ৫ টি বেড়াল ইঁদুর ধরে ৫ দিনে, তাহলে ১০০ টা বেড়াল ১০০ টা ইঁদুর
ধরবে-}

ক) ১ দিনে

\textbf{\hl{খ) ৫ দিনে ✔}}

গ) ২০ দিনে

ঘ) ১০০ দিনে

\textbf{\hl{✔ সঠিক উত্তর: খ) ৫ দিনে}}

\begin{quote}
\textbf{ব্যাখ্যা:} \emph{এখানে ১ টি বেড়াল ১ টি ইঁদুর ধরে ৫ দিনে। তাহলে ১০০
টি বেড়াল ১০০ টি ইঁদুর ৫ দিনেই ধরবে।}
\end{quote}

\textbf{ক ও খ একত্র একটি কাজ ৮ দিনে করতে পারে। ঐ কাজ ক একা ১২ দিনে করতে
পারে। কাজটি করতে খ এর কত দিন লাগবে?}

ক) ২০ দিন

খ) ২২ দিন

\textbf{\hl{গ) ২৪ দিন ✔}}

ঘ) ২৬ দিন

\textbf{\hl{✔ সঠিক উত্তর: গ) ২৪ দিন}}

\begin{quote}
\textbf{ব্যাখ্যা:} \emph{ক ও খ একত্রে কাজটির ১ দিনে করতে পারে = (১/৮) অংশ}

\emph{ক একা কাজটির ১ দিনে করতে পারে = (১/১২) অংশ}

\emph{∴ খ একা কাজটির ১ দিনে করতে পারে = ( (১/৮) − (১/১২) ) = (৩−২/২৪) =
(১/২৪) অংশ}

\emph{∴ খ একা কাজটি ২৪ দিনে করতে পারবে।}
\end{quote}

\textbf{ক একটি কাজ ১২ দিনে এবং খ ২৪ দিনে করতে পারে। তারা একত্রে কাজ আরম্ভ
করে এবং কয়েক দিন পর ক কাজটি অসমাপ্ত রেখে চলে যায়। বাকি কাজটুকু খ ৩ দিনে শেষ
করে। কাজটি কত দিনে শেষ হয়েছিল?}

\textbf{\hl{ক) ১০ দিনে ✔}}

খ) ৯ দিনে

গ) ৮ দিনে

ঘ) ১১ দিনে

\textbf{\hl{✔ সঠিক উত্তর: ক) ১০ দিনে}}

\begin{quote}
\textbf{ব্যাখ্যা:} \emph{ক একা কাজটির ১ দিনে করতে পারে = (১/১২) অংশ}

\emph{খ একা কাজটির ১ দিনে করতে পারে = (১/২৪) অংশ}

\emph{ক ও খ একত্রে ১ দিনে করতে পারে = ( (১/১২) + (১/২৪) ) = (৩/২৪) =
(১/৮) অংশ কাজ}

\emph{খ একা কাজটির ৩ দিনে করতে পারে = ( (১/২৪) × ৩) = (১/৮) অংশ}

\emph{বাকি কাজ = (১ − (১/৮) ) = (৭/৮) অংশ}

\emph{ক ও খ একত্রে (১/৮) অংশ কাজ করতে পারে = ১ দিনে}

\emph{∴ ক ও খ একত্রে ১ বা সম্পূর্ণ কাজটি করতে পারে = ৮ দিনে}

\emph{∴ ক ও খ একত্রে (৭/৮) অংশ কাজ করতে পারে = (৮ × (৭/৮) ) = ৭ দিনে}

\emph{খ বাকি কাজ শেষ করতে পারে ৩ দিনে।}

\emph{∴ কাজটি শেষ হতে সময় লেগেছিল = (৭ + ৩) = ১০ দিন}

\emph{শর্ট টেকনিক:}

\emph{দুই ব্যক্তি কাজ শুরু করার পর একজন চলে গেলে কাজ শেষ হওয়ার সময় (যদি
একজনের কাজের সময় অপর জনের দ্বিগুণ হয়):}

\emph{কাজ শেষ হওয়ার সময় = (২/৩) × (১ম সময় + ২য় সময়)}

\emph{⇒ (২/৩) × (১২ + ৩) = (২/৩) × ১৫ = ১০ দিন।}
\end{quote}

\textbf{একজন শ্রমিক প্রতিদিন প্রথম ৮ ঘন্টা কাজের জন্য ঘন্টায় ১০ টাকা আয় করে
এবং পরবর্তী সময়ের ঘন্টায় ১৫ টাকা করে মজুরী পায়। দৈনিক ১০ ঘন্টা কাজ করলে তার
ঘন্টাপ্রতি গড় মজুরী কত?}

\textbf{\hl{ক) ১১ টাকা ✔}}

খ) ১২ টাকা

গ) ১২.৫০ টাকা

ঘ) ১৩ টাকা

\textbf{\hl{✔ সঠিক উত্তর: ক) ১১ টাকা}}

\begin{quote}
\textbf{ব্যাখ্যা:} \emph{প্রথম ৮ ঘন্টায় পায় = (৮ × ১০) = ৮০ টাকা}

\emph{পরবর্তী (১০ − ৮) = ২ ঘন্টায় পায় = (২ × ১৫) = ৩০ টাকা}

\emph{∴ শ্রমিক (৮ + ২) = ১০ ঘন্টায় পায় = (৮০ + ৩০) = ১১০ টাকা}

\emph{∴ শ্রমিক ১ ঘন্টায় পায় = (১১০/১০) = ১১ টাকা।}
\end{quote}

\textbf{ক একটি কাজ ৫ দিনে এবং খ তা ১০ দিনে করতে পারে। তারা একত্রে ১ দিনে
এর কত অংশ করতে পারবে?}

ক) (১/১৫)

খ) (১/১০)

গ) (২/১৫)

\textbf{\hl{ঘ) (৩/১০) ✔}}

\textbf{\hl{✔ সঠিক উত্তর: ঘ) (৩/১০)}}

\begin{quote}
\textbf{ব্যাখ্যা:} \emph{\textquotesingle ক\textquotesingle{} ১ দিনে করতে
পারে = (১/৫) অংশ}

\emph{\textquotesingle খ\textquotesingle{} ১ দিনে করতে পারে = (১/১০)
অংশ}

\emph{∴ (ক + খ) ১ দিনে করতে পারে = ( (১/৫) + (১/১০) ) = (২+১/১০) =
(৩/১০) অংশ।}
\end{quote}

\textbf{কাজের দিন ২ টাকা ও অনুপস্থিতির দিন ৫০ পয়সা করে জরিমানার শর্তে এক
ব্যক্তি সেপ্টেম্বর মাসে ৪০ টাকা পেল। ব্যক্তিটি কাজে কতদিন উপস্থিত ছিল?}

ক) ২০০

\textbf{\hl{খ) ২২ ✔}}

গ) ২৪

ঘ) ১৯

\textbf{\hl{✔ সঠিক উত্তর: খ) ২২}}

\begin{quote}
\textbf{ব্যাখ্যা:} \emph{সেপ্টেম্বর মাস = ৩০ দিন}

\emph{প্রতিদিন ২ টাকা হিসাবে কাজের জন্য ৬০ টাকা পাওয়ার কথা।}

\emph{কিন্তু একদিন অনুপস্থিত থাকলে জরিমানা ৫০ পয়সাসহ কাজের ২ টাকা না পাওয়াতে
তার প্রতিদিন ২ টাকা ৫০ পয়সা ক্ষতি হয়।}

\emph{তাই ৮ দিন অনুপস্থিত থাকলে ২০ টাকা ক্ষতি হয়।}

\emph{∴ ব্যক্তিটি কাজে (৩০ − ৮) = ২২ দিন উপস্থিত ছিল।}
\end{quote}

\textbf{শাহেদ ২২ দিনে একটি কাজ শেষ করার জন্য ৪২ জন লোক নিয়োগ দিল।১২ দিন
পর মাত্র এক তৃতীয়াংশ কাজ শেষ হলো। অতিরিক্ত কোনো লোক নেয়া না হলে অবশিষ্ট কাজ
শেষ করতে নির্ধারিত সময় অপেক্ষা কতদিন বেশি লাগবে?}

ক) ৩৬

খ) ২৪

গ) ১৬

\textbf{\hl{ঘ) ১৪ ✔}}

ঙ) কোনোটিই নয়

\textbf{\hl{✔ সঠিক উত্তর: ঘ) ১৪}}

\begin{quote}
\textbf{ব্যাখ্যা:} \emph{(১/৩) অংশ কাজ হয় = ১২ দিনে}

\emph{∴ ১ বা সম্পূর্ণ অংশ কাজ হয় = (১২ × ৩) = ৩৬ দিনে}

\emph{∴ বেশি সময় লাগবে = (৩৬ − ২২) = ১৪ দিন।}
\end{quote}

\textbf{কয়েকজন শ্রমিক একটি কাজ ১৮ দিনে করে দিবে বলে ঠিক করল। কিন্তু তাদের
মধ্যে ৯ জন অনুপস্থিত থাকায় কাজটি ৩৬ দিনে সম্পন্ন হয়। কত জন শ্রমিক নিযুক্ত হলে
কাজটি ১৮ দিনে সম্পন্ন হতো?}

ক) ১২০

\textbf{\hl{খ) ১৮ ✔}}

গ) ২৪

ঘ) কোনোটিই নয়

\textbf{\hl{✔ সঠিক উত্তর: খ) ১৮}}

\begin{quote}
\textbf{ব্যাখ্যা:} \emph{ধরি, ক সংখ্যক শ্রমিক লাগতো।}

\emph{তাহলে, (ক − ৯) জনে কাজটি শেষ করে = ৩৬ দিনে}

\emph{∴ ১ জনে কাজটি শেষ করে = ৩৬(ক − ৯) দিনে}

\emph{∴ ক জনে কাজটি শেষ করে = (৩৬(ক−৯)/ক) দিনে}

\emph{প্রশ্নমতে,}

\emph{(৩৬(ক−৯)/ক) = ১৮}

\emph{⇒ ৩৬(ক − ৯) = ১৮ক}

\emph{⇒ ৩৬ক − ৩২৪ = ১৮ক}

\emph{⇒ ৩৬ক − ১৮ক = ৩২৪}

\emph{⇒ ১৮ক = ৩২৪}

\emph{⇒ ক = (৩২৪/১৮)}

\emph{∴ ক = ১৮ জন।}
\end{quote}

\textbf{দৈনিক ৬ ঘণ্টা পরিশ্রম করে ১০ জন লোক একটি কাজ ৭ দিনে করতে পারে,
দৈনিক কত ঘণ্টা পরিশ্রম করে ১৪ জনে ৬ দিনে ওই কাজটি করতে পারবেন?}

\textbf{\hl{ক) ৫ ঘণ্টা ✔}}

খ) ১৫ ঘণ্টা

গ) ১০ ঘণ্টা

ঘ) ৭ ঘণ্টা

\textbf{\hl{✔ সঠিক উত্তর: ক) ৫ ঘণ্টা}}

\begin{quote}
\textbf{ব্যাখ্যা:} \emph{১০ জন লোক ৭ দিনে কাজটি করে দৈনিক = ৬ ঘণ্টায়}

\emph{∴ ১ জন লোক ১ দিনে কাজটি করে দৈনিক = (৬ × ৭ × ১০) ঘণ্টায়}

\emph{∴ ১৪ জন লোক ৬ দিনে কাজটি করে দৈনিক = (৬×৭×১০/১৪×৬) = ৫ ঘণ্টা পরিশ্রম
করে।}
\end{quote}

\textbf{একটি কার ওয়াশার ৮টি কার ওয়াশ করে ৩৬ মিনিটে। এ হারে তিন ঘণ্টায়
কয়টি কার ওয়াশ করবে?}

ক) ৮০টি

\textbf{\hl{খ) ৪০টি ✔}}

গ) ৫০টি

ঘ) ৬০টি

\textbf{\hl{✔ সঠিক উত্তর: খ) ৪০টি}}

\begin{quote}
\textbf{ব্যাখ্যা:} \emph{৩ ঘণ্টা = (৩ × ৬০) বা ১৮০ মিনিট}

\emph{৩৬ মিনিটে ওয়াশ করে ৮টি গাড়ি}

\emph{∴ ১ মিনিটে ওয়াশ করে = (৮/৩৬) টি গাড়ি}

\emph{∴ ১৮০ মিনিটে ওয়াশ করে = (৮×১৮০/৩৬) = ৪০টি গাড়ি।}
\end{quote}

\textbf{করিমের সাপ্তাহিক আয় ৩৫০০ টাকা হলে এই হিসাবে জুন মাসে তার কত টাকা
আয় হয়েছিল?}

ক) ১২০০০

খ) ১৪০০০

\textbf{\hl{গ) ১৫০০০ ✔}}

ঘ) ১০৪০০

\textbf{\hl{✔ সঠিক উত্তর: গ) ১৫০০০}}

\begin{quote}
\textbf{ব্যাখ্যা:} \emph{সাপ্তাহিক আয় ৩৫০০ টাকা}

\emph{∴ দৈনিক আয় = (৩৫০০/৭) = ৫০০ টাকা}

\emph{জুন মাসের দিনের সংখ্যা ৩০ দিন।}

\emph{∴ জুন মাসে তার আয় = (৩০ × ৫০০) = ১৫,০০০ টাকা।}
\end{quote}

\textbf{৬ ফুট দীর্ঘ একটি বাঁশের ৪ ফুট দীর্ঘ ছায়া হয়। একই সময়ে একটি গাছের ছায়া
৬৪ ফুট লম্বা। গাছটির উচ্চতা কত ফুট?}

\textbf{\hl{ক) ৯৬ ✔}}

খ) ৭২

গ) ১৯২

ঘ) ৪৪

\textbf{\hl{✔ সঠিক উত্তর: ক) ৯৬}}

\begin{quote}
\textbf{ব্যাখ্যা:} \emph{ছায়া ৪ ফুট হলে প্রকৃত উচ্চতা = ৬ ফুট}

\emph{∴ ছায়া ১ ফুট হলে প্রকৃত উচ্চতা = (৬/৪) ফুট}

\emph{∴ ছায়া ৬৪ ফুট হলে প্রকৃত উচ্চতা = (৬×৬৪/৪) ফুট = ৯৬ ফুট।}
\end{quote}

\textbf{একটি কলমের দাম ১০ টাকা এবং ১০টি খামের দাম ৩ টাকা। ৩টি কলম ও ১০টি
খামের দাম হবে -}

\textbf{\hl{ক) ৩৩ ✔}}

খ) ৩৬

গ) ৩৯

ঘ) কোনটিই নয়

\textbf{\hl{✔ সঠিক উত্তর: ক) ৩৩}}

\begin{quote}
\textbf{ব্যাখ্যা:} \emph{একটি কলমের দাম = ১০ টাকা}

\emph{∴ তিনটি কলমের দাম = (১০ × ৩) = ৩০ টাকা}

\emph{প্রশ্নানুসারে, ১০টি খামের দাম = ৩ টাকা}

\emph{∴ তিনটি কলম ও ১০টি খামের দাম = (৩০ + ৩) টাকা = ৩৩ টাকা।}
\end{quote}

\textbf{একটি গাড়ি ঘণ্টায় ৬০ কিলোমিটার বেগে চললে, ৩ মিনিট ৩০ সেকেন্ডে উহা
কত দূর যাবে?}

ক) ৩.৩ কিলোমিটার

\textbf{\hl{খ) ৩.৫ কিলোমিটার ✔}}

গ) ২১০ কিলোমিটার

ঘ) ২০ কিলোমিটার

\textbf{\hl{✔ সঠিক উত্তর: খ) ৩.৫ কিলোমিটার}}

\begin{quote}
\textbf{ব্যাখ্যা:} \emph{১ ঘণ্টা = ৩৬০০ সেকেন্ড}

\emph{৩ মিনিট ৩০ সেকেন্ড = (১৮০ + ৩০) সেকেন্ড = ২১০ সেকেন্ড}

\emph{৩৬০০ সেকেন্ডে যায় ৬০ কিলোমিটার}

\emph{∴ ২১০ সেকেন্ডে যায় = (৬০×২১০/৩৬০০) = ৩.৫ কিলোমিটার।}
\end{quote}

\textbf{একটি গাড়ি ৩ ঘণ্টায় ১২ কিমি গেলে ২ ঘণ্টায় কত মিটার যাবে?}

ক) ৪

খ) ৪০০০

গ) ৮

\textbf{\hl{ঘ) ৮০০০ ✔}}

\textbf{\hl{✔ সঠিক উত্তর: ঘ) ৮০০০}}

\begin{quote}
\textbf{ব্যাখ্যা:} \emph{গাড়িটি ৩ ঘণ্টায় যায় ১২ কিমি}

\emph{∴ গাড়িটি ১ ঘণ্টায় যায় = (১২/৩) কিমি}

\emph{∴ গাড়িটি ২ ঘণ্টায় যায় = (১২×২/৩) = ৮ কিমি = (৮ × ১০০০) মিটার = ৮০০০
মিটার।}
\end{quote}

\textbf{একটি মেশিন ১ ঘণ্টায় ৩০টি প্যাকেট উৎপাদন করে। মেশিনটি ৬ মিনিটে কতটি
প্যাকেট উৎপাদন করবে?}

\textbf{\hl{ক) ৩ টি ✔}}

খ) ৫ টি

গ) ১৫ টি

ঘ) ২০ টি

\textbf{\hl{✔ সঠিক উত্তর: ক) ৩ টি}}

\begin{quote}
\textbf{ব্যাখ্যা:} \emph{১ ঘণ্টা = ৬০ মিনিট}

\emph{৬০ মিনিটে মেশিনটি উৎপাদন করে ৩০টি প্যাকেট}

\emph{∴ ১ মিনিটে মেশিনটি উৎপাদন করে = (৩০/৬০) টি প্যাকেট}

\emph{∴ ৬ মিনিটে মেশিনটি উৎপাদন করে = (৩০×৬/৬০) = ৩টি প্যাকেট।}
\end{quote}

\textbf{If 12 carpenters, working 6 hours a day, can make 460 chairs in
24 days, how many chairs will 18 carpenters make in 36 days, each
working 8 hours a day?/ যদি ১২ জন ছুতার প্রতিদিন ৬ কর্মঘণ্টায় ৪৬০ টি চেয়ার
২৪ দিনে তৈরি করতে পারে। তাহলে ১৮ জন ছুতার প্রতিদিন ৮ কর্মঘণ্টা করে ৩৬ দিনে
কতটি চেয়ার তৈরি করতে পারবে?}

ক) 1260

খ) 1320

গ) 920

\textbf{\hl{ঘ) 1380 ✔}}

\textbf{\hl{✔ সঠিক উত্তর: ঘ) 1380}}

\begin{quote}
\textbf{ব্যাখ্যা:} \emph{১২ জন ছুতারের মোট কর্মঘণ্টা = ১২ × ৬ × ২৪ = ১৭২৮
ঘণ্টা}

\emph{১৭২৮ ঘণ্টায় তৈরি হয় ৪৬০টি চেয়ার}

\emph{∴ ১ ঘণ্টায় তৈরি হয় = (৪৬০/১৭২৮) টি চেয়ার}

\emph{১৮ জন ছুতারের মোট কর্মঘণ্টা = ১৮ × ৮ × ৩৬ = ৫১৮৪ ঘণ্টা}

\emph{∴ ৫১৮৪ ঘণ্টায় তৈরি হয় = (৪৬০×৫১৮৪/১৭২৮) = ৪৬০ × ৩ = ১৩৮০টি চেয়ার।}
\end{quote}

\textbf{১৮টি ছাগলের দাম ৪টি গরুর দামের সমান হলে, ৪৫টি ছাগলের পরিবর্তে কতটি
গরু পাওয়া যাবে?}

ক) ৭টি

খ) ৮টি

গ) ৯টি

\textbf{\hl{ঘ) ১০টি ✔}}

\textbf{\hl{✔ সঠিক উত্তর: ঘ) ১০টি}}

\begin{quote}
\textbf{ব্যাখ্যা:} \emph{১৮টি ছাগলের পরিবর্তে পাওয়া যায় ৪টি গরু}

\emph{∴ ১টি ছাগলের পরিবর্তে পাওয়া যায় = (৪/১৮) টি গরু}

\emph{∴ ৪৫টি ছাগলের পরিবর্তে পাওয়া যায় = (৪৫×৪/১৮) = ১০টি গরু।}
\end{quote}

\textbf{৫ টি গরুর মূল্য ১৫ টি ছাগলের মূল্যের সমান। একটি গরুর মূল্য ৬,০০০ টাকা
হলে একটি ছাগলের মূল্য কত?}

ক) ১,৮০০

খ) ২,২০০

গ) ১,৫০০

\textbf{\hl{ঘ) ২,০০০ ✔}}

\textbf{\hl{✔ সঠিক উত্তর: ঘ) ২,০০০}}

\begin{quote}
\textbf{ব্যাখ্যা:} \emph{৫টি গরুর মূল্য = ১৫টি ছাগলের মূল্য}

\emph{∴ ১টি গরুর মূল্য = (১৫/৫) = ৩টি ছাগলের মূল্য}

\emph{অর্থাৎ, ৩টি ছাগলের মূল্য = ৬,০০০ টাকা}

\emph{∴ ১টি ছাগলের মূল্য = (৬০০০/৩) = ২,০০০ টাকা।}
\end{quote}

\textbf{৩টি ঘোড়ার মূল্য ৫টি গরুর মূল্যের সমান এবং ২টি গরুর মূল্য ৩টি গাধার মূল্যের
সমান। ১টি ঘোড়ার মূল্য ৭৫০০ টাকা হলে ৫টি গাধার মূল্য কত?}

\textbf{\hl{ক) ১৫,০০০ টাকা ✔}}

খ) ২২,৫০ সূ টাকা

গ) ৩০,০০০ টাকা

ঘ) ৩১,০০০ টাকা

\textbf{\hl{✔ সঠিক উত্তর: ক) ১৫,০০০ টাকা}}

\begin{quote}
\textbf{ব্যাখ্যা:} \emph{১টি ঘোড়ার মূল্য = ৭৫০০ টাকা}

\emph{∴ ৩টি ঘোড়ার মূল্য = (৩ × ৭৫০০) = ২২,৫০০ টাকা}

\emph{প্রশ্নমতে, ৫টি গরুর মূল্য = ৩টি ঘোড়ার মূল্য = ২২,৫০০ টাকা}

\emph{∴ ১টি গরুর মূল্য = (২২৫০০/৫) = ৪৫০০ টাকা}

\emph{∴ ২টি গরুর মূল্য = (২ × ৪৫০০) = ৯০০০ টাকা}

\emph{আবার, ৩টি গাধার মূল্য = ২টি গরুর মূল্য = ৯০০০ টাকা}

\emph{∴ ১টি গাধার মূল্য = (৯০০০/৩) = ৩০০০ টাকা}

\emph{∴ ৫টি গাধার মূল্য = (৫ × ৩০০০) = ১৫,০০০ টাকা।}
\end{quote}

\textbf{১৫টি গরুর মূল্য ৫টি ঘোড়ার মূল্যের সমান। ২টি ঘোড়ার মূল্য ৩,০০০ টাকা হলে
৩টি গরুর মূল্য কত?}

ক) ১,০০০ টাকা

খ) ১,৮০০ টাকা

\textbf{\hl{গ) ১,৫০০ টাকা ✔}}

ঘ) ২,০০০ টাকা

\textbf{\hl{✔ সঠিক উত্তর: গ) ১,৫০০ টাকা}}

\begin{quote}
\textbf{ব্যাখ্যা:} \emph{২টি ঘোড়ার মূল্য = ৩,০০০ টাকা}

\emph{∴ ১টি ঘোড়ার মূল্য = (৩০০০/২) = ১,৫০০ টাকা}

\emph{∴ ৫টি ঘোড়ার মূল্য = (১,৫০০ × ৫) = ৭,৫০০ টাকা}

\emph{প্রশ্নমতে, ১৫টি গরুর মূল্য = ৫টি ঘোড়ার মূল্য = ৭,৫০০ টাকা}

\emph{∴ ১টি গরুর মূল্য = (৭৫০০/১৫) = ৫০০ টাকা}

\emph{∴ ৩টি গরুর মূল্য = (৫০০ × ৩) = ১,৫০০ টাকা।}
\end{quote}

\textbf{৫ টি বিড়াল ৫টি ইঁদুর ধরতে ৫ মিনিট সময় লাগে। ১০০টি বিড়াল ১০০টি ইঁদুর
ধরতে কত সময় লাগবে ?}

\textbf{\hl{ক) ৫ মিনিট ✔}}

খ) ১০ মিনিট

গ) ১৫ মিনিট

ঘ) ২০ মিনিট

\textbf{\hl{✔ সঠিক উত্তর: ক) ৫ মিনিট}}

\begin{quote}
\textbf{ব্যাখ্যা:} \emph{৫টি বিড়ালের ৫টি ইঁদুর ধরতে সময় লাগে = ৫ মিনিট}

\emph{∴ ১টি বিড়ালের ১টি ইঁদুর ধরতে সময় লাগে = ৫ মিনিট (কারণ ১টি বিড়াল ১টি
ইঁদুর ধরতেও ওই ৫ মিনিট সময়ই নিবে)}

\emph{∴ ১০০টি বিড়াল সমান্তরালভাবে ১০০টি ইঁদুর ধরলে প্রত্যেকেরই সমান ৫ মিনিট
সময় লাগবে।}
\end{quote}

\textbf{১০৫ কেজি ডালের দাম ৩,৬৭৫ টাকা হলে ৬০ কেজি ডালের দাম কত ?}

ক) ২,২০০ টাকা

খ) ২,১৫০ টাকা

\textbf{\hl{গ) ২,১০০ টাকা ✔}}

ঘ) ২,০৫০ টাকা

ঙ) কোনোটিই নয়

\textbf{\hl{✔ সঠিক উত্তর: গ) ২,১০০ টাকা}}

\begin{quote}
\textbf{ব্যাখ্যা:} \emph{১০৫ কেজি ডালের দাম = ৩,৬৭৫ টাকা}

\emph{∴ ১ কেজি ডালের দাম = (৩৬৭৫/১০৫) টাকা}

\emph{∴ ৬০ কেজি ডালের দাম = (৩৬৭৫×৬০/১০৫) = ২,১০০ টাকা।}
\end{quote}

\textbf{পানিপূর্ণ একটি বালতির ওজন ১৫ কেজি। শূন্য বালতির ওজন ১.৫ কেজি। ধাক্কা
লেগে ৫ কেজি পানি ওই পড়ে গেলে ঐ বালতিতে আর কতটুকু পানি রয়েছে?}

ক) ১৪ কেজি

খ) ১৩.৫০ কেজি

\textbf{\hl{গ) ৮.৫ কেজি ✔}}

ঘ) ১০ কেজি

\textbf{\hl{✔ সঠিক উত্তর: গ) ৮.৫ কেজি}}

\begin{quote}
\textbf{ব্যাখ্যা:} \emph{পানিপূর্ণ বালতির ওজন = ১৫ কেজি}

\emph{শূণ্য বালতির ওজন = ১.৫ কেজি}

\emph{∴ বালতিতে মোট পানির ওজন = (১৫ − ১.৫) কেজি = ১৩.৫ কেজি}

\emph{ধাক্কা লেগে পানি পড়ে যায় = ৫ কেজি}

\emph{∴ অবশিষ্ট পানির পরিমাণ = (১৩.৫ − ৫) কেজি = ৮.৫ কেজি।}
\end{quote}

\textbf{৮৪ কেজি ধাতু গলিয়ে ৪০,০০০ পিন তৈরি করা হলো। প্রতিটি পিনের ওজন কত?}

ক) ০.০০২০ কেজি

\textbf{\hl{খ) ০.০০২১ কেজি ✔}}

গ) ০.০০২২ কেজি

ঘ) ০.০০২৩ কেজি

\textbf{\hl{✔ সঠিক উত্তর: খ) ০.০০২১ কেজি}}

\begin{quote}
\textbf{ব্যাখ্যা:} \emph{৪০,০০০টি পিনের ওজন = ৮৪ কেজি}

\emph{∴ প্রতিটি পিনের ওজন = (৮৪/৪০০০০) কেজি = (২১/১০০০০) কেজি = ০.০০২১
কেজি।}
\end{quote}

\textbf{৮ জন শ্রমিকের একটি জমি চাষ করতে ৭ দিন লাগল। ১৪ জন শ্রমিকের ঐ জমি
চাষ করতে কতদিন লাগবে?}

ক) ৫ দিন

খ) ৬ দিন

গ) ৩ দিন

\textbf{\hl{ঘ) ৪ দিন ✔}}

\textbf{\hl{✔ সঠিক উত্তর: ঘ) ৪ দিন}}

\begin{quote}
\textbf{ব্যাখ্যা:} \emph{৮ জন শ্রমিকের লাগে = ৭ দিন}

\emph{∴ ১ জন শ্রমিকের লাগে = (৮ × ৭) দিন}

\emph{∴ ১৪ জন শ্রমিকের লাগে = (৮×৭/১৪) দিন = ৪ দিন।}
\end{quote}

\textbf{১০ জনে একটি কাজের অর্ধেক করতে পারে ৭ দিনে। ঐ কাজটি ৫ জনে করতে
কতদিন সময় লাগবে?}

ক) ৩৫ দিন

খ) ৪২ দিন

গ) ১৪ দিন

\textbf{\hl{ঘ) ২৮ দিন ✔}}

\textbf{\hl{✔ সঠিক উত্তর: ঘ) ২৮ দিন}}

\begin{quote}
\textbf{ব্যাখ্যা:} \emph{১০ জন একটি কাজের অর্ধেক করতে পারে = ৭ দিনে}

\emph{∴ ১ জন একটি কাজের অর্ধেক করতে পারে = (৭ × ১০) দিনে}

\emph{∴ ৫ জন একটি কাজের অর্ধেক করতে পারে = (৭×১০/৫) = ১৪ দিনে}

\emph{∴ ৫ জন একটি কাজের সম্পূর্ণ অংশ করতে পারে = (১৪ × ২) = ২৮ দিনে।}
\end{quote}

\textbf{যদি ১৫টি কলমের দাম ৪৬.৫ টাকা হয় তাহলে ২টি কলমের দাম কত?}

ক) ১০ টাকা

খ) ৪.৫ টাকা

\textbf{\hl{গ) ৬.২ টাকা ✔}}

ঘ) ১২ টাকা

\textbf{\hl{✔ সঠিক উত্তর: গ) ৬.২ টাকা}}

\begin{quote}
\textbf{ব্যাখ্যা:} \emph{১৫টি কলমের দাম = ৪৬.৫ টাকা}

\emph{∴ ১টি কলমের দাম = (৪৬.৫/১৫) টাকা}

\emph{∴ ২টি কলমের দাম = (৪৬.৫×২/১৫) টাকা = ৬.২ টাকা।}
\end{quote}

\textbf{একটি কলম ও দুইটি পেন্সিলের একত্রিত মূল্য ৬৮ টাকা। যদি একটি কলমের মূল্য
একটি পেন্সিলের মূল্যের চেয়ে ৫ টাকা বেশি হয়, তবে পেন্সিলের মূল্য কত টাকা?}

ক) ১৮

খ) ২০

\textbf{\hl{গ) ২১ ✔}}

ঘ) কোনোটিই নয়

\textbf{\hl{✔ সঠিক উত্তর: গ) ২১}}

\begin{quote}
\textbf{ব্যাখ্যা:} \emph{ধরি, পেন্সিলের মূল্য x টাকা}

\emph{∴ কলমের মূল্য = (x + ৫) টাকা}

\emph{প্রশ্নমতে, (x + ৫) + ২x = ৬৮}

\emph{⇒ ৩x + ৫ = ৬৮}

\emph{⇒ ৩x = ৬৮ − ৫}

\emph{⇒ ৩x = ৬৩}

\emph{⇒ x = (৬৩/৩)}

\emph{∴ x = ২১}

\emph{∴ পেন্সিলের মূল্য ২১ টাকা।}
\end{quote}

\textbf{১০০ টাকায় ৫টি কলম ও ১০টি পেন্সিল অথবা ৭টি কলম ও ৬টি পেন্সিল ক্রয়
করা যায়। ১টি কলম ও ১টি পেন্সিলের মূল্য কত?}

\textbf{\hl{ক) কলমের মূল্য ১০ টাকা এবং পেন্সিলের মূল্য ৫ টাকা ✔}}

খ) কলমের মূল্য ৮ টাকা এবং পেন্সিলের মূল্য ৬ টাকা

গ) কলমের মূল্য ১২ টাকা এবং পেন্সিলের মূল্য ৪ টাকা

ঘ) কলমের মূল্য ১৫ টাকা এবং পেন্সিলের মূল্য ৫ টাকা

\textbf{\hl{✔ সঠিক উত্তর: ক) কলমের মূল্য ১০ টাকা এবং পেন্সিলের মূল্য ৫ টাকা}}

\begin{quote}
\textbf{ব্যাখ্যা:} \emph{ধরি, ১টি কলমের মূল্য x টাকা এবং ১টি পেন্সিলের মূল্য y
টাকা}

\emph{১ম শর্তমতে, ৫x + ১০y = ১০০}

\emph{⇒ x + ২y = ২০ ⇒ x = ২০ − ২y -\/-\/-(i)}

\emph{২য় শর্তমতে, ৭x + ৬y = ১০০}

\emph{⇒ ৭(২০ − ২y) + ৬y = ১০০ {[}(i) হতে{]}}

\emph{⇒ ১৪০ − ১৪y + ৬y = ১০০}

\emph{⇒ − ৮y = − ৪০}

\emph{∴ y = ৫}

\emph{∴ (i) হতে, x = ২০ − (২ × ৫) = ১০}

\emph{∴ ১টি কলমের মূল্য ১০ টাকা এবং ১টি পেন্সিলের মূল্য ৫ টাকা।}
\end{quote}

\textbf{কয়েকজন শ্রমিক একটি কাজ ১৮ দিনে করে দিবে বলে ঠিক করল। কিন্তু তাদের
মধ্যে ৯ জন অনুপস্থিত থাকায় কাজটি ৩৬ দিনে সম্পন্ন হয়। কত জন শ্রমিক নিযুক্ত হলে
কাজটি ১৮ দিনে সম্পন্ন হতো?}

ক) ১২

\textbf{\hl{খ) ১৮ ✔}}

গ) ২৪

ঘ) কোনোটিই নয়

\textbf{\hl{✔ সঠিক উত্তর: খ) ১৮}}

\begin{quote}
\textbf{ব্যাখ্যা:} \emph{ধরি, x জন শ্রমিক কাজটি ১৮ দিনে সম্পন্ন করে।}

\emph{কিন্তু ৯ জন অনুপস্থিত থাকায় শ্রমিক সংখ্যা = (x − ৯) জন}

\emph{আমরা জানি, শ্রমিক সংখ্যা ও প্রয়োজনীয় দিনের অনুপাত ব্যস্তানুপাতিক:}

\emph{∴ x : (x − ৯) = ৩৬ : ১৮}

\emph{⇒ (x/x−৯)=(৩৬/১৮)}

\emph{⇒ (x/x−৯)=(২/১)}

\emph{⇒ ২x − ১৮ = x}

\emph{⇒ ২x − x = ১৮}

\emph{∴ x = ১৮}

\emph{∴ ১৮ জন শ্রমিক নিযুক্ত হলে কাজটি ১৮ দিনে সম্পন্ন হতো।}
\end{quote}

\textbf{কাজের দিন ৫০ টাকা মজুরি এবং অনুপস্থিতির দিন ১০ টাকা জরিমানার শর্তে
একজন শ্রমিক কাজে যোগদান করে সেপ্টেম্বর মাসে মোট ১২০০ টাকা পেল। সে কতদিন
অনুপস্থিত ছিল?}

\textbf{\hl{ক) ৫ দিন ✔}}

খ) ৬ দিন

গ) ৭ দিন

ঘ) ৮ দিন

\textbf{\hl{✔ সঠিক উত্তর: ক) ৫ দিন}}

\begin{quote}
\textbf{ব্যাখ্যা:} \emph{সেপ্টেম্বর মাস = ৩০ দিন।}

\emph{যদি শ্রমিকটি ৩০ দিনই উপস্থিত থাকত তবে সে পেত = (৩০ × ৫০) = ১৫০০
টাকা।}

\emph{সে মোট টাকা কম পেয়েছে = (১৫০০ − ১২০০) = ৩০০ টাকা।}

\emph{একদিন অনুপস্থিত থাকার কারণে ক্ষতির পরিমাণ = ৫০ টাকা (মজুরি না পাওয়া) +
১০ টাকা (জরিমানা) = ৬০ টাকা।}

\emph{∴ অনুপস্থিতির সংখ্যা = (৩০০/৬০) = ৫ দিন।}
\end{quote}

\textbf{একটি কলম ও একটি বইয়ের মূল্য একত্রে ৯৫ টাকা। কলমটির মূল্য ১৫ টাকা বেশি
ও বইটির মূল্য ১৪ টাকা কম হলে কলমটির মূল্য বইটির মূল্যের দ্বিগুণ হতো। বইটির মূল্য
কত?}

ক) ৪৯ টাকা

\textbf{\hl{খ) ৪৬ টাকা ✔}}

গ) ৫০ টাকা

ঘ) ৪০ টাকা

ঙ) কোনটিই নয়

\textbf{\hl{✔ সঠিক উত্তর: খ) ৪৬ টাকা}}

\begin{quote}
\textbf{ব্যাখ্যা:} \emph{ধরি, কলমের মূল্য = x টাকা}

\emph{∴ বইয়ের মূল্য = (৯৫ − x) টাকা}

\emph{প্রশ্নমতে, x + ১৫ = ২ × \{(৯৫ − x) − ১৪\}}

\emph{⇒ x + ১৫ = ২ × (৮১ − x)}

\emph{⇒ x + ১৫ = ১৬২ − ২x}

\emph{⇒ ৩x = ১৬২ − ১৫}

\emph{⇒ ৩x = ১৪৭}

\emph{⇒ x = ৪৯ টাকা (কলমের মূল্য)}

\emph{∴ বইয়ের মূল্য = ৯৫ − ৪৯ = ৪৬ টাকা।}
\end{quote}

\textbf{আপনার কাছে ২০টি ১০০ টাকার নোট আছে। তা থেকে ৫টি ১০০ টাকার নোট ছোট
ভাইকে দিলে কত টাকা থাকবে?}

ক) ২৫০০ টাকা

খ) ২০০০ টাকা

\textbf{\hl{গ) ১৫০০ টাকা ✔}}

\textbf{\hl{✔ সঠিক উত্তর: গ) ১৫০০ টাকা}}

\begin{quote}
\textbf{ব্যাখ্যা:} \emph{২০টি ১০০ টাকার নোট = (২০ × ১০০) = ২০০০ টাকা}

\emph{৫টি ১০০ টাকার নোট = (৫ × ১০০) = ৫০০ টাকা}

\emph{∴ ছোট ভাইকে ৫টি ১০০ টাকার নোট দিলে অবশিষ্ট থাকে = (২০০০ − ৫০০) =
১৫০০ টাকা।}
\end{quote}

\textbf{১৩৫ টাকায় ১০০ টি আম পাওয়া গেলে ২টির দাম কত?}

ক) ১৭.৭০ টাকা

খ) ১৪.৭০ টাকা

গ) ১২.৭০ টাকা

\textbf{\hl{ঘ) ২.৭০ টাকা ✔}}

\textbf{\hl{✔ সঠিক উত্তর: ঘ) ২.৭০ টাকা}}

\begin{quote}
\textbf{ব্যাখ্যা:} \emph{১০০টি আমের দাম = ১৩৫ টাকা}

\emph{∴ ১টি আমের দাম = (১৩৫/১০০) টাকা}

\emph{∴ ২টি আমের দাম = (১৩৫×২/১০০) = ২.৭ টাকা বা ২.৭০ টাকা।}
\end{quote}

\textbf{ক এর টাকা খ-এর টাকার দ্বিগুণ। তাদের দুজনের মোট ৩০ টাকা আছে। খ-এর কত
টাকা আছে?}

ক) ৩০টাকা

খ) ২০ টাকা

গ) ১৫ টাকা

\textbf{\hl{ঘ) ১০ টাকা ✔}}

\textbf{\hl{✔ সঠিক উত্তর: ঘ) ১০ টাকা}}

\begin{quote}
\textbf{ব্যাখ্যা:} \emph{ধরি, খ এর কাছে আছে = x টাকা}

\emph{∴ ক এর কাছে আছে = ২x টাকা}

\emph{প্রশ্নমতে, x + ২x = ৩০}

\emph{⇒ ৩x = ৩০}

\emph{⇒ x = (৩০/৩)}

\emph{∴ x = ১০}

\emph{∴ খ এর কাছে আছে ১০ টাকা।}
\end{quote}

\textbf{একটি ক্লাবে ১৭৮৯ টাকা চাঁদা তোলা হলো। যদি প্রতি সদস্য কমপক্ষে ৩০ টাকা
চাঁদা দিয়ে থাকে তাহলে ক্লাবের সর্বোচ্চ সদস্য সংখ্যা কত হতে পারে?}

ক) ৩১

খ) ৪৯

গ) ৫৭

\textbf{\hl{ঘ) ৫৯ ✔}}

\textbf{\hl{✔ সঠিক উত্তর: ঘ) ৫৯}}

\begin{quote}
\textbf{ব্যাখ্যা:} \emph{ধরি, সদস্য সংখ্যা x জন}

\emph{প্রতিজন সদস্য কমপক্ষে ৩০ টাকা চাঁদা দিলে,}

\emph{৩০x ≤ ১৭৮৯}

\emph{⇒ x ≤ (১৭৮৯/৩০)}

\emph{⇒ x ≤ ৫৯.৬৩}

\emph{∴ পূর্ণসংখ্যায় ক্লাবের সর্বোচ্চ সদস্য সংখ্যা ৫৯ জন হতে পারে।}
\end{quote}

\textbf{ক ও খ -এর মধ্যে ১৮০ টাকা এমনভাবে ভাগ করে দিতে হবে যেন খ, ক এর
টাকার দ্বিগুণ বেশি পায়। ক-কত টাকা পাবে?}

ক) ৫০ টাকা

খ) ৭০ টাকা

\textbf{\hl{গ) ৬০ টাকা ✔}}

ঘ) ৮০ টাকা

\textbf{\hl{✔ সঠিক উত্তর: গ) ৬০ টাকা}}

\begin{quote}
\textbf{ব্যাখ্যা:} \emph{ধরি, ক পায় = x টাকা}

\emph{∴ খ পায় = ২x টাকা}

\emph{প্রশ্নমতে, x + ২x = ১৮০}

\emph{⇒ ৩x = ১৮০}

\emph{⇒ x = (১৮০/৩) = ৬০ টাকা।}

\emph{∴ ক পাবে ৬০ টাকা।}
\end{quote}

\textbf{কিছুসংখ্যক মুরগী ও কয়েকটি গরুর পায়ের সংখ্যা উহাদের মাথার সংখ্যার দ্বিগুণ
অপেক্ষা ১০ বেশি। গরুর সংখ্যা কত?}

\textbf{\hl{ক) ৫টি ✔}}

খ) ১০টি

গ) ১৫টি

ঘ) ২০টি

\textbf{\hl{✔ সঠিক উত্তর: ক) ৫টি}}

\begin{quote}
\textbf{ব্যাখ্যা:} \emph{ধরি, মুরগীর সংখ্যা m এবং গরুর সংখ্যা g}

\emph{মোট মাথার সংখ্যা = m + g}

\emph{মোট পায়ের সংখ্যা = ২m + ৪g (মুরগীর ২ পা, গরুর ৪ পা)}

\emph{প্রশ্নমতে, (২m + ৪g) = ২(m + g) + ১০}

\emph{⇒ ২m + ৪g = ২m + ২g + ১০}

\emph{⇒ ৪g − ২g = ১০}

\emph{⇒ ২g = ১০}

\emph{⇒ g = (১০/২) = ৫}

\emph{∴ গরুর সংখ্যা ৫।}

\emph{(বিকল্পভাবে: প্রতি গরুর জন্য মাথার দ্বিগুণের চেয়ে ২টি অতিরিক্ত পা থাকে।
সুতরাং গরুর সংখ্যা = ১০ ÷ ২ = ৫টি)।}
\end{quote}

\textbf{একটি শহরের জনসংখ্যা ৫০,০০০ জন। প্রতিদিন ২০ জনের মৃত্যু হয় এবং ৩০ জন
শিশু জন্মগ্রহণ করে। এক বছর পর শহরের জনসংখ্যা কত হবে?}

ক) ৫১,৬৫০ জন

\textbf{\hl{খ) ৫৩,৬৫০ জন ✔}}

গ) ৫৪,৫০০ জন

ঘ) ৫২,০০০ জন

\textbf{\hl{✔ সঠিক উত্তর: খ) ৫৩,৬৫০ জন}}

\begin{quote}
\textbf{ব্যাখ্যা:} \emph{১ দিনে জনসংখ্যা বৃদ্ধি পায় = (৩০ − ২০) = ১০ জন}

\emph{আমরা জানি, ১ বছর = ৩৬৫ দিন}

\emph{∴ ৩৬৫ দিনে জনসংখ্যা বৃদ্ধি পায় = (৩৬৫ × ১০) = ৩,৬৫০ জন}

\emph{∴ ১ বছর পর শহরের জনসংখ্যা হবে = (৫০,০০০ + ৩,৬৫০) = ৫৩,৬৫০ জন।}
\end{quote}

\textbf{যদি কাঁচ পানি অপেক্ষা ২.৫ গুণ বেশি ভারী হয়, তবে ৪০ ঘন সেন্টিমিটার
কাঁচের ওজন কত?}

\textbf{\hl{ক) ১০০ গ্রাম ✔}}

খ) ২৫০ গ্রাম

গ) ৬০০ গ্রাম

ঘ) ১০০০ গ্রাম

\textbf{\hl{✔ সঠিক উত্তর: ক) ১০০ গ্রাম}}

\begin{quote}
\textbf{ব্যাখ্যা:} \emph{১ ঘন সেমি পানির ওজন = ১ গ্রাম}

\emph{∴ ১ ঘন সেমি কাঁচের ওজন = (২.৫ × ১) গ্রাম = ২.৫ গ্রাম}

\emph{∴ ৪০ ঘন সেমি কাঁচের ওজন = (৪০ × ২.৫) গ্রাম = ১০০ গ্রাম।}
\end{quote}

\textbf{m সংখ্যক পেন্সিলের দাম p টাকা হলে m টাকায় কয়টি পেন্সিল পাওয়া যাবে?}

ক) p টি

\textbf{\hl{খ) (m\^{}২/p) টি ✔}}

গ) (m/p\^{}২) টি

ঘ) mp টি

\textbf{\hl{✔ সঠিক উত্তর: খ) (m\^{}২/p) টি}}

\begin{quote}
\textbf{ব্যাখ্যা:} \emph{p টাকায় পাওয়া যায় = m টি পেন্সিল}

\emph{∴ ১ টাকায় পাওয়া যায় = (m/p) টি পেন্সিল}

\emph{∴ m টাকায় পাওয়া যায় = (m×m/p)=(m\^{}২/p) টি পেন্সিল।}

\emph{∴ m টাকায় (m\^{}২/p) সংখ্যক পেন্সিল পাওয়া যাবে।}
\end{quote}

\textbf{x সংখ্যক আমের দাম y টাকা হলে, a টাকায় কতটি আম পাওয়া যাবে?}

ক) (ay/x)

\textbf{\hl{খ) (ax/y) ✔}}

গ) (a/yx)

ঘ) (yx/a)

\textbf{\hl{✔ সঠিক উত্তর: খ) (ax/y)}}

\begin{quote}
\textbf{ব্যাখ্যা:} \emph{y টাকায় পাওয়া যায় = x টি আম}

\emph{∴ ১ টাকায় পাওয়া যায় = (x/y) টি আম}

\emph{∴ a টাকায় পাওয়া যায় = (ax/y) টি আম।}
\end{quote}

\textbf{৩৬টি ব্যাগের দাম ১৮০০ টাকা হলে, ৫০০ টাকায় এরূপ কয়টি ব্যাগ পাওয়া
যাবে?}

ক) ১২টি

খ) ১১টি

\textbf{\hl{গ) ১০টি ✔}}

ঘ) ১৫টি

\textbf{\hl{✔ সঠিক উত্তর: গ) ১০টি}}

\begin{quote}
\textbf{ব্যাখ্যা:} \emph{১৮০০ টাকায় পাওয়া যায় = ৩৬ টি ব্যাগ}

\emph{∴ ১ টাকায় পাওয়া যায় = (৩৬/১৮০০) টি ব্যাগ}

\emph{∴ ৫০০ টাকায় পাওয়া যায় = (৩৬×৫০০/১৮০০) = ১০ টি ব্যাগ।}
\end{quote}

\textbf{খালেক ঈদ বোনাসের টাকায় (১/৩) অংশ খরচ করে একটি জ্যাকেট এবং ১০০
টাকা খরচ করে একটি শার্ট কিনলো। এই খরচ করার পর তার কাছে বোনাসের (১/২) অংশ
টাকা রয়ে গেল। খালেক কত টাকা বোনাস পেয়েছিল?}

ক) ৬০০০

খ) ৭২০০

গ) ৮৪০০

\textbf{\hl{ঘ) কোনোটিই নয় ✔}}

\textbf{\hl{✔ সঠিক উত্তর: ঘ) কোনোটিই নয়}}

\begin{quote}
\textbf{ব্যাখ্যা:} \emph{ধরি, খালেক x টাকা বোনাস পেয়েছিল।}

\emph{প্রশ্নমতে, (x/৩)+১০০+(x/২)=x}

\emph{⇒ (২x+৩x/৬)+১০০=x}

\emph{⇒ (৫x/৬)+১০০=x}

\emph{⇒ x−(৫x/৬)=১০০}

\emph{⇒ (x/৬)=১০০}

\emph{∴ x = ৬০০ টাকা।}

\emph{যেহেতু অপশনগুলোতে ৬০০ টাকা নেই, তাই সঠিক উত্তর
\textquotesingle কোনোটিই নয়\textquotesingle।}
\end{quote}

\textbf{আরিফের একদিনের আয় বাবুর একদিনের আয়ের দ্বিগুণ। আরিফের আয় যদি ১০০
টাকা কমে যায় এবং বাবুর আয় যদি ৫০ টাকা বেড়ে যায় তবে তাদের আয় সমান হবে।
আরিফের একদিনের আয় কত?}

ক) ২০০

খ) ২৫০

\textbf{\hl{গ) ৩০০ ✔}}

ঘ) কোনোটিই নয়

\textbf{\hl{✔ সঠিক উত্তর: গ) ৩০০}}

\begin{quote}
\textbf{ব্যাখ্যা:} \emph{ধরি, বাবুর একদিনের আয় = x টাকা}

\emph{∴ আরিফের একদিনের আয় = ২x টাকা}

\emph{প্রশ্নমতে, ২x − ১০০ = x + ৫০}

\emph{⇒ ২x − x = ৫০ + ১০০}

\emph{∴ x = ১৫০}

\emph{∴ আরিফের একদিনের আয় = (২ × ১৫০) = ৩০০ টাকা।}
\end{quote}

\textbf{রহিমের আয়ের দ্বিগুণের সাথে ১১০ টাকা যোগ করলে মোট ৭,০০০ টাকা হয়।
রহিমের আয় কত?}

ক) ৩,২৭৫ টাকা

\textbf{\hl{খ) ৩,৪৪৫ টাকা ✔}}

গ) ৩,২৬০ টাকা

ঘ) ৪,২৭০ টাকা

ঙ) কোনটিই নয়

\textbf{\hl{✔ সঠিক উত্তর: খ) ৩,৪৪৫ টাকা}}

\begin{quote}
\textbf{ব্যাখ্যা:} \emph{ধরি, রহিমের আয় = x টাকা}

\emph{প্রশ্নমতে, ২x + ১১০ = ৭০০০}

\emph{⇒ ২x = ৭০০০ − ১১০}

\emph{⇒ ২x = ৬৮৯০}

\emph{⇒ x = (৬৮৯০/২)}

\emph{∴ x = ৩,৪৪৫ টাকা।}
\end{quote}

\textbf{একটি বালতির ভেতরের আয়তন ১.৫ লিটার হলে ৪৫০ লিটারে কত বালতি পানি
হবে?}

\textbf{\hl{ক) ৩০০ বালতি ✔}}

খ) ৪৫০ বালতি

গ) ৫০০ বালতি

ঘ) ৬৭৫ বালতি

ঙ) কোনোটিই নয়

\textbf{\hl{✔ সঠিক উত্তর: ক) ৩০০ বালতি}}

\begin{quote}
\textbf{ব্যাখ্যা:} \emph{১.৫ লিটার পানি ধারণ করে = ১ বালতিতে}

\emph{∴ ৪৫০ লিটার পানি ধারণ করবে = (৪৫০/১.৫) = ৩০০ বালতিতে।}
\end{quote}

\textbf{রহিমের আয়ের দ্বিগুণের সাথে ১১০ টাকা যোগ করলে মোট ৭০০০ টাকা হয়।
রহিমের আয় কত?}

ক) ৩২৭৫ টাকা

\textbf{\hl{খ) ৩৪৪৫ টাকা ✔}}

গ) ৩২৬০ টাকা

ঘ) ৪২৭০ টাকা

ঙ) কোনোটিই নয়

\textbf{\hl{✔ সঠিক উত্তর: খ) ৩৪৪৫ টাকা}}

\begin{quote}
\textbf{ব্যাখ্যা:} \emph{ধরি, রহিমের আয় = x টাকা}

\emph{∴ x × ২ + ১১০ = ৭০০০}

\emph{⇒ ২x = ৬৮৯০}

\emph{⇒ x = ৩৪৪৫ টাকা।}
\end{quote}

\textbf{একটি কলমের দাম যদি P টাকা হয় তবে Q টি 100 টাকার নোট দিয়ে কয়টি
কলম কেনা যাবে?}

ক) 100/ PQ

খ) PQ /100

গ) Q /100P

\textbf{\hl{ঘ) 100Q / P ✔}}

\textbf{\hl{✔ সঠিক উত্তর: ঘ) 100Q / P}}

\begin{quote}
\textbf{ব্যাখ্যা:} \emph{মোট টাকার পরিমাণ = (১০০ × Q) = ১০০Q টাকা}

\emph{P টাকায় কলম কেনা যায় = ১টি}

\emph{∴ ১ টাকায় কলম কেনা যায় = (১/P) টি}

\emph{∴ ১০০Q টাকায় কলম কেনা যায় = (১×১০০Q/P) = (১০০Q/P) টি।}
\end{quote}

\textbf{কোন কর্মকর্তা মাসিক ৩০,০০০ টাকা বেতনে ২০১২ সালের ৭ ফেব্রুয়ারি কাজে
যোগ দিলেন। তিনি ঐ মাসে কত বেতন পাবেন?}

\textbf{\hl{ক) ২৩০০০ টাকা ✔}}

খ) ২৪০০০.২০ টাকা

গ) ২৩৭৯৩.১০ টাকা

ঘ) ২৪৯১২ টাকা

\textbf{\hl{✔ সঠিক উত্তর: ক) ২৩০০০ টাকা}}

\begin{quote}
\textbf{ব্যাখ্যা:} \emph{২০১২ সাল Leap year বা অধিবর্ষ ছিল। এজন্য ফেব্রুয়ারি
মাস ২৯ দিনে হয়েছিল।}

\emph{মাসিক বেতন ৩০,০০০ টাকা হিসাবে ১ দিনের বেতন = (৩০,০০০/৩০) = ১০০০
টাকা।}

\emph{৭ তারিখে যোগদান করায় ফেব্রুয়ারি মাসে তিনি কাজ করেন = (২৯ − ৭ + ১) =
২৩ দিন।}

\emph{∴ ২৩ দিনে তিনি মোট বেতন পাবেন = (২৩ × ১০০০) = ২৩০০০ টাকা।}
\end{quote}

\textbf{ফারুকের বেতন কালামের বেতনের ১.২ গুণ। মালেকের বেতন কালামের বেতনের
০.৮ গুণ। তাদের মোট বেতন ৬৩০০ টাকা হলে মালেকের বেতন কত?}

\textbf{\hl{ক) ১৬৮০ ✔}}

খ) ২১০০

গ) ২৫২০

ঘ) কোনটিই নয়

\textbf{\hl{✔ সঠিক উত্তর: ক) ১৬৮০}}

\begin{quote}
\textbf{ব্যাখ্যা:} \emph{ধরি, কালামের বেতন = x টাকা}

\emph{∴ ফারুকের বেতন = ১.২x টাকা}

\emph{এবং মালেকের বেতন = ০.৮x টাকা}

\emph{প্রশ্নমতে, x + ১.২x + ০.৮x = ৬৩০০}

\emph{⇒ ৩x = ৬৩০০}

\emph{∴ x = ২১০০ টাকা}

\emph{∴ মালেকের বেতন = ০.৮ × ২১০০ = ১৬৮০ টাকা।}
\end{quote}

\textbf{৮ জন লোক একদিনে ২০টি চেয়ার বানাতে পারলে ১৬ জন লোক ৪০টি চেয়ার
বানাতে কত সময় নিবে?}

\textbf{\hl{ক) ১ দিন ✔}}

খ) ২ দিন

গ) ৪ দিন

ঘ) কোনোটিই নয়

\textbf{\hl{✔ সঠিক উত্তর: ক) ১ দিন}}

\begin{quote}
\textbf{ব্যাখ্যা:} \emph{৮ জন লোক ২০টি চেয়ার বানায় = ১ দিনে}

\emph{∴ ১ জন লোক ২০টি চেয়ার বানায় = (১ × ৮) দিনে}

\emph{∴ ১৬ জন লোক ২০টি চেয়ার বানায় = (১×৮/১৬) দিনে}

\emph{∴ ১৬ জন লোক ৪০টি চেয়ার বানায় = (১×৮×৪০/১৬×২০) = ১ দিনে।}
\end{quote}

\textbf{একটি শ্রেণীতে যত জন বালক ছিল প্রত্যেকে তত টাকা চাঁদা দেয়াতে ১০০ টাকা
হলো। বালকের সংখ্যা কত?}

\textbf{\hl{ক) ১০ ✔}}

খ) ১০০

গ) ২৫

ঘ) ৩৫

\textbf{\hl{✔ সঠিক উত্তর: ক) ১০}}

\begin{quote}
\textbf{ব্যাখ্যা:} \emph{ধরি, বালকের সংখ্যা x জন}

\emph{প্রশ্নমতে, x × x = ১০০}

\emph{⇒ x\^{}২ = ১০০}

\emph{∴ x = ১০}

\emph{∴ বালকের সংখ্যা ১০।}
\end{quote}

\textbf{কোনো স্থানে যতজন লোক ছিল প্রত্যেকে তত পাঁচ টাকা করে চাঁদা দেওয়ায় মোট
৪৫০০ টাকা আদায় হলো। এখানে লোক সংখ্যা কত?}

ক) ৭৫০

খ) ৯০০

গ) ৮০০

\textbf{\hl{ঘ) কোনোটিই না ✔}}

\textbf{\hl{✔ সঠিক উত্তর: ঘ) কোনোটিই না}}

\begin{quote}
\textbf{ব্যাখ্যা:} \emph{ধরি, লোকসংখ্যা ছিল x জন}

\emph{প্রশ্নমতে, x × ৫x = ৪৫০০}

\emph{⇒ ৫x\^{}২ = ৪৫০০}

\emph{⇒ x\^{}২ = ৯০০}

\emph{∴ x = ৩০}

\emph{যেহেতু অপশনগুলোতে ৩০ নেই, তাই সঠিক উত্তর \textquotesingle কোনোটিই
না\textquotesingle।}
\end{quote}

\textbf{কোনো পরীক্ষায় একটি সংখ্যাকে বর্গ করে তা থেকে ৯ বিয়োগ করতে বলা হলো
কিন্তু একজন পরীক্ষার্থী ভুলবশত প্রথমে বিয়োগ করে পরে বর্গ করল। ফলাফল পূর্বের মতোই
রয়ে গেল, সংখ্যাটি কত?}

ক) ৭

খ) ৮

\textbf{\hl{গ) ৫ ✔}}

ঘ) ৬

\textbf{\hl{✔ সঠিক উত্তর: গ) ৫}}

\begin{quote}
\textbf{ব্যাখ্যা:} \emph{ধরি, সংখ্যাটি x}

\emph{প্রশ্নমতে, x\^{}২ − ৯ = (x − ৯)\^{}২}

\emph{⇒ x\^{}২ − ৯ = x\^{}২ − ২ • x • ৯ + ৯\^{}২}

\emph{⇒ x\^{}২ − ৯ = x\^{}২ − ১৮x + ৮১}

\emph{⇒ ১৮x = ৮১ + ৯}

\emph{⇒ ১৮x = ৯০}

\emph{⇒ x = (৯০/১৮)}

\emph{∴ x = ৫}
\end{quote}

\textbf{একটি ছাত্রাবাসে যতজন ছাত্র থাকে, তাদের প্রত্যেকের মাসিক খরচ তাদের
মোট সংখ্যার দশগুণ। ঐ ছাত্রাবাসের সকল ছাত্রের মোট মাসিক খরচ ৬,২৫০ টাকা হলে ঐ
ছাত্রাবাসে কতজন ছাত্র থাকে?}

ক) ১৫

\textbf{\hl{খ) ২৫ ✔}}

গ) ৩৫

ঘ) ৪৫

\textbf{\hl{✔ সঠিক উত্তর: খ) ২৫}}

\begin{quote}
\textbf{ব্যাখ্যা:} \emph{ধরি, ছাত্র সংখ্যা x জন}

\emph{∴ ১ জনের মাসিক খরচ = ১০x টাকা}

\emph{∴ x জনের মোট মাসিক খরচ = x × ১০x = ১০x\^{}২ টাকা}

\emph{প্রশ্নমতে, ১০x\^{}২ = ৬২৫০}

\emph{⇒ x\^{}২ = ৬২৫}

\emph{⇒ x = √(৬২৫)}

\emph{∴ x = ২৫ জন।}
\end{quote}

\textbf{একটি সেনাবাহিনীর গুদামে ১৫০০ সৈনিকের ৪০ দিনের খাদ্য মজুদ আছে। ১৩
দিন পর কিছু সৈনিক অন্য জায়গায় চলে গেল। বাকি খাদ্য অবশিষ্ট সৈনিকদের আরো ৩০
দিন চললো। কতজন সৈনিক অন্য জায়গায় চলে গিয়েছিল?}

ক) ২০০

\textbf{\hl{খ) ১৫০ ✔}}

গ) ২১০

ঘ) ১২৫

\textbf{\hl{✔ সঠিক উত্তর: খ) ১৫০}}

\begin{quote}
\textbf{ব্যাখ্যা:} \emph{১৩ দিন খাওয়ার পর খাবার বাকি থাকে = (৪০ − ১৩)
দিনের = ২৭ দিনের}

\emph{১৫০০ জন সৈনিকের খাদ্য চলে = ২৭ দিন}

\emph{∴ ১ জন সৈনিকের খাদ্য চলে = (২৭ × ১৫০০) দিন}

\emph{∴ ৩০ দিন খাদ্য চলে = (২৭×১৫০০/৩০) = ১৩৫০ জনের}

\emph{∴ অন্যত্র চলে যাওয়া সৈনিকের সংখ্যা = (১৫০০ − ১৩৫০) জন = ১৫০ জন।}
\end{quote}

\textbf{একটি ছাত্রাবাসে ১৫ জন ছাত্রের ৩২ দিনের খাদ্য আছে। কয়েকজন নতুন ছাত্র
আসায় ২০ দিনে ঐ খাদ্য শেষ হলে নতুন ছাত্রের সংখ্যা কত?}

\textbf{\hl{ক) ৯ জন ✔}}

খ) ১৭ জন

গ) ২৪ জন

ঘ) ১০ জন

\textbf{\hl{✔ সঠিক উত্তর: ক) ৯ জন}}

\begin{quote}
\textbf{ব্যাখ্যা:} \emph{৩২ দিনের খাদ্য আছে = ১৫ জন ছাত্রের}

\emph{∴ ১ দিনের খাদ্য আছে = (১৫ × ৩২) জন ছাত্রের}

\emph{∴ ২০ দিনের খাদ্য আছে = (১৫×৩২/২০) = ২৪ জন ছাত্রের}

\emph{∴ নতুন ছাত্রের সংখ্যা = (২৪ − ১৫) জন = ৯ জন।}
\end{quote}

\textbf{একটি ক্যাম্পে ৭২ জন স্কাউটের ৬ দিনের খাদ্য মজুদ আছে। ১৮ জন স্কাউট চলে
গেলে ঐ খাদ্যে বাকি স্কাউটের আরো কত দিন চলবে?}

\textbf{\hl{ক) ২ দিন ✔}}

খ) ৩ দিন

গ) ৪ দিন

ঘ) ৬ দিন

\textbf{\hl{✔ সঠিক উত্তর: ক) ২ দিন}}

\begin{quote}
\textbf{ব্যাখ্যা:} \emph{১৮ জন স্কাউট চলে গেলে অবশিষ্ট স্কাউট = (৭২ − ১৮) জন
= ৫৪ জন।}

\emph{৭২ জন স্কাউটের খাদ্য চলে = ৬ দিন}

\emph{∴ ১ জন স্কাউটের খাদ্য চলে = (৬ × ৭২) দিন}

\emph{∴ ৫৪ জন স্কাউটের খাদ্য চলে = (৬×৭২/৫৪) = ৮ দিন।}

\emph{∴ বাকি স্কাউটের আরো চলবে = (৮ − ৬) দিন = ২ দিন।}

\emph{(নোট: মূল ৬ দিন চলার কথা থাকলেও অবশিষ্ট খাদ্য দিয়ে নতুন করে আরো ৮ দিন
অর্থাৎ আগের ৬ দিনের চেয়ে অতিরিক্ত ২ দিন চলবে)।}
\end{quote}

\textbf{একটি বাড়িতে ২০ জনের ৩৫ দিনের খাবার আছে। ৫ দিন পরে ঐ বাড়িতে আরো
১০ জন লোক এলো ঐ খাবার কতদিন চলবে?}

\textbf{\hl{ক) ২০ ✔}}

খ) ২২

গ) ২৫

ঘ) কোনোটিই নয়

\textbf{\hl{✔ সঠিক উত্তর: ক) ২০}}

\begin{quote}
\textbf{ব্যাখ্যা:} \emph{১০ জন লোক আসায় মোট লোক সংখ্যা = (২০ + ১০) জন = ৩০
জন}

\emph{অবশিষ্ট দিন = (৩৫ − ৫) দিন = ৩০ দিন}

\emph{২০ জনের খাবার চলে = ৩০ দিন}

\emph{∴ ১ জনের খাবার চলে = (৩০ × ২০) দিন}

\emph{∴ ৩০ জনের খাবার চলে = (৩০×২০/৩০) দিন = ২০ দিন।}
\end{quote}

\textbf{যে পরিমাণ খাদ্যে ২০০ জন লোকের ২০ সপ্তাহ চলে, ঐ পরিমাণ খাদ্যে কতজন
লোকের ৮ সপ্তাহ চলবে?}

ক) ৩০০ জন

\textbf{\hl{খ) ৫০০ জন ✔}}

গ) ৪০০ জন

ঘ) ৬০০ জন

\textbf{\hl{✔ সঠিক উত্তর: খ) ৫০০ জন}}

\begin{quote}
\textbf{ব্যাখ্যা:} \emph{২০ সপ্তাহ চলে = ২০০ জন লোকের}

\emph{∴ ১ সপ্তাহ চলে = (২০০ × ২০) জন লোকের}

\emph{∴ ৮ সপ্তাহ চলে = (২০০×২০/৮) = ৫০০ জন লোকের।}
\end{quote}

\textbf{একটি শিবিরে ৭২০ জন সৈন্যের ২০ দিনের খাবার মজুদ আছে। ১০ দিন পর কিছু
নতুন সৈন্যের আগমনের কারণে অবশিষ্ট খাদ্যে তাদের ৮ দিন চলে। শিবিরে কতজন নতুন
সৈন্য এসেছিল?}

ক) ১৭৫

খ) ১৯০

গ) ১৭০

ঘ) ২০০

\textbf{\hl{ঙ) ১৮০ ✔}}

\textbf{\hl{✔ সঠিক উত্তর: ঙ) ১৮০}}

\begin{quote}
\textbf{ব্যাখ্যা:} \emph{১০ দিন পর অবশিষ্ট দিন = (২০ − ১০) দিন = ১০ দিন}

\emph{১০ দিন খাদ্য চলে = ৭২০ জন সৈন্যের}

\emph{∴ ১ দিন খাদ্য চলে = = (৭২০ × ১০) জন সৈন্যের}

\emph{∴ ৮ দিন খাদ্য চলে = (৭২০×১০/৮) = ৯০০ জন সৈন্যের}

\emph{∴ শিবিরে নতুন সৈন্য এসেছিল (৯০০ − ৭২০) = ১৮০ জন}
\end{quote}

\textbf{কোন শিবিরে ২৮৫ জন সৈন্যের রোজ জনপ্রতি ৬৫০ গ্রাম হিসেবে ৪০ দিনের
খাবার আছে। জনপ্রতি রোজ ৫০০ গ্রাম হিসেবে ঐ খাদ্য ৬০ দিন চালাতে হলে শিবির
থেকে কতজন সৈন্য অন্যত্র পাঠাতে হবে?}

\textbf{\hl{ক) ৩৮ জন ✔}}

খ) ৪০ জন

গ) ৪২ জন

ঘ) ৪৫ জন

\textbf{\hl{✔ সঠিক উত্তর: ক) ৩৮ জন}}

\begin{quote}
\textbf{ব্যাখ্যা:} \emph{জনপ্রতি ৬৫০ গ্রাম হিসেবে ৪০ দিনের খাবার আছে = ২৮৫
জন সৈন্যের}

\emph{∴ জনপ্রতি ১ গ্রাম হিসেবে ১ দিনের খাবার আছে = (২৮৫ × ৬৫০ × ৪০) জন
সৈন্যের}

\emph{∴ জনপ্রতি ৫০০ গ্রাম হিসেবে ৬০ দিন চালাতে সৈন্য সংখ্যা হবে =
(২৮৫×৬৫০×৪০/৫০০×৬০) = ২৪৭ জন সৈন্যের}

\emph{∴ শিবির থেকে সৈন্য অন্যত্র পাঠাতে হবে = (২৮৫ − ২৪৭) জন = ৩৮ জন।}
\end{quote}

\textbf{৫ টন খাবারে ১২০টি হাতির ৫৫ দিন চলে। ১৫০টি হাতির ঐ খাবারে কত দিন
চলবে?}

ক) ২৫ দিন

খ) ৩৫ দিন

\textbf{\hl{গ) ৪৪ দিন ✔}}

ঘ) ৫৪ দিন

ঙ) কোনোটিই নয়

\textbf{\hl{✔ সঠিক উত্তর: গ) ৪৪ দিন}}

\begin{quote}
\textbf{ব্যাখ্যা:} \emph{১২০টি হাতির চলে = ৫৫ দিন}

\emph{∴ ১টি হাতির চলে = (৫৫ × ১২০) দিন}

\emph{∴ ১৫০টি হাতির চলে = (৫৫×১২০/১৫০) = ৪৪ দিন।}
\end{quote}

\textbf{একজন কৃষকের ৪০টি গরুর জন্য ৩৫ দিনের খাদ্য মজুদ আছে। তিনি আরো ১০টি গরু
ক্রয় করলে ঐ খাদ্যে কত দিন চলবে?}

ক) ২০ দিন

খ) ২৪ দিন

গ) ২৬ দিন

\textbf{\hl{ঘ) ২৮ দিন ✔}}

\textbf{\hl{✔ সঠিক উত্তর: ঘ) ২৮ দিন}}

\begin{quote}
\textbf{ব্যাখ্যা:} \emph{১০টি গরু নতুন ক্রয় করায় মোট গরু সংখ্যা = (৪০ + ১০) টি
= ৫০টি।}

\emph{৪০টি গরুর খাদ্য মজুদ আছে = ৩৫ দিনের}

\emph{∴ ১টি গরুর খাদ্য মজুদ আছে = (৩৫ × ৪০) দিনের}

\emph{∴ ৫০টি গরুর খাদ্য মজুদ আছে = (৩৫×৪০/৫০) দিনের = ২৮ দিনের।}
\end{quote}

\textbf{একটি বড় বাক্সের মধ্যে ৫টি বাক্স আছে ও তার প্রত্যেকটির ভেতর ৫টি করে
ছোট বাক্স আছে। মোট বাক্সের সংখ্যা কত?}

ক) ২৫টি

খ) ৩০টি

গ) ৩৩টি

\textbf{\hl{ঘ) ৩১টি ✔}}

\textbf{\hl{✔ সঠিক উত্তর: ঘ) ৩১টি}}

\begin{quote}
\textbf{ব্যাখ্যা:} \emph{বড় বাক্স = ১টি}

\emph{বড় বাক্সের ভেতরে মাঝারি বাক্স = ৫টি}

\emph{৫টি মাঝারি বাক্সের প্রত্যেকটির ভেতরে ছোট বাক্স = (৫ × ৫) = ২৫টি}

\emph{∴ মোট বাক্সের সংখ্যা = (১ + ৫ + ২৫)টি = ৩১টি।}
\end{quote}

\textbf{যদি ১২ জন শ্রমিক ৪ দিনে ২৮৮০ টাকা আয় করে, তবে ৮ জন শ্রমিক কতদিনে
সমপরিমাণ টাকা আয় করবে?}

ক) ৩ দিনে

খ) ৪ দিনে

গ) ৫ দিনে

\textbf{\hl{ঘ) ৬ দিনে ✔}}

\textbf{\hl{✔ সঠিক উত্তর: ঘ) ৬ দিনে}}

\begin{quote}
\textbf{ব্যাখ্যা:} \emph{১২ জন শ্রমিক ২৮৮০ টাকা আয় করে = ৪ দিনে}

\emph{∴ ১ জন শ্রমিক ২৮৮০ টাকা আয় করে = (৪ × ১২) দিনে}

\emph{∴ ৮ জন শ্রমিক ২৮৮০ টাকা আয় করে = (৪×১২/৮) দিনে = ৬ দিনে।}
\end{quote}

\textbf{বাবু ও তপুর কাছে কিছু মার্বেল আছে। বাবু যদি তপুকে ১০ টি মার্বেল দিয়ে
দেয় তবে তাদের মার্বেলের সংখ্যা সমান হবে। আবার তপু যদি বাবুকে ২০টি মার্বেল দেয়
তবে বাবুর মার্বেলের সংখ্যা দ্বিগুণ হবে। বাবুর কাছে কতটি মার্বেল আছে?}

ক) ১১০

\textbf{\hl{খ) ১০০ ✔}}

গ) ৯০

ঘ) ১২০

\textbf{\hl{✔ সঠিক উত্তর: খ) ১০০}}

\begin{quote}
\textbf{ব্যাখ্যা:} \emph{ধরি, বাবুর কাছে মার্বেল আছে x টি এবং তপুর কাছে আছে y
টি।}

\emph{১ম শর্তমতে, x − ১০ = y + ১০}

\emph{⇒ x = y + ২০ ........(i)}

\emph{২য় শর্তমতে, x + ২০ = ২(y − ২০)}

\emph{⇒ x + ২০ = ২y − ৪০}

\emph{⇒ y + ২০ + ২০ = ২y − ৪০ {[}(i) নং হতে x এর মান বসিয়ে{]}}

\emph{⇒ ২y − y = ৪০ + ৪০}

\emph{∴ y = ৮০}

\emph{(i) নং সমীকরণে y = ৮০ বসিয়ে পাই,}

\emph{x = ৮০ + ২০ = ১০০}

\emph{∴ বাবুর কাছে ১০০টি মার্বেল আছে।}
\end{quote}

\textbf{বাবু ও তপুর কাছে কিছু মার্বেল আছে। বাবু যদি তপুকে ১০টি মার্বেল দিয়ে
দেয়, তবে তাদের মার্বেলের সংখ্যা সমান হবে। আর তপু যদি বাবুকে ২০টি মার্বেল দেয়।
তবে বাবুর মার্বেল সংখ্যা তপুর মার্বেলের সংখ্যার দ্বিগুণ হবে। বাবুর কাছে কয়টি
মার্বেল আছে?}

\textbf{\hl{ক) ১০০ ✔}}

খ) ১২০

গ) ১৪০

ঘ) কোনোটিই নয়

\textbf{\hl{✔ সঠিক উত্তর: ক) ১০০}}

\begin{quote}
\textbf{ব্যাখ্যা:} \emph{ধরি, বাবুর মার্বেল সংখ্যা = x এবং তপুর মার্বেল সংখ্যা =
y}

\emph{১ম শর্তমতে, x − ১০ = y + ১০ ⇒ x − y = ২০ -\/-\/-\/-(i)}

\emph{২য় শর্তমতে, x + ২০ = ২(y − ২০)}

\emph{⇒ x + ২০ = ২y − ৪০}

\emph{⇒ x − ২y = − ৬০ -\/-\/-\/-(ii)}

\emph{(i) হতে (ii) বিয়োগ করে পাই,}

\emph{y = ৮০}

\emph{y এর মান (i) নং সমীকরণে বসিয়ে পাই,}

\emph{x − ৮০ = ২০ ⇒ x = ১০০}

\emph{∴ বাবুর কাছে ১০০টি মার্বেল আছে।}
\end{quote}

\textbf{১২ জন শ্রমিক ৩ দিনে ৭২০ টাকা আয় করে। ৯ জন শ্রমিক সমপরিমাণ টাকা আয়
করবে কত দিনে?}

ক) ৫ দিনে

\textbf{\hl{খ) ৪ দিনে ✔}}

গ) ৬ দিনে

ঘ) ৩ দিনে

\textbf{\hl{✔ সঠিক উত্তর: খ) ৪ দিনে}}

\begin{quote}
\textbf{ব্যাখ্যা:} \emph{১২ জন শ্রমিক ৭২০ টাকা আয় করে = ৩ দিনে}

\emph{∴ ১ জন শ্রমিক ৭২০ টাকা আয় করে = (১২ × ৩) দিনে}

\emph{∴ ৯ জন শ্রমিক ৭২০ টাকা আয় করে = (১২×৩/৯) দিনে = ৪ দিনে।}
\end{quote}

\textbf{রফিক একটি কাজ ১০ দিনে করতে পারে। শফিক ঐ কাজ ১৫ দিনে করতে পারে।
তারা একত্রে কত দিনে কাজটি শেষ করতে পারবে?}

ক) ৮ দিনে

খ) ৫ দিনে

\textbf{\hl{গ) ৬ দিনে ✔}}

ঘ) ৭ দিন

\textbf{\hl{✔ সঠিক উত্তর: গ) ৬ দিনে}}

\begin{quote}
\textbf{ব্যাখ্যা:} \emph{রফিক ১ দিনে করে কাজের (১/১০) অংশ}

\emph{শফিক ১ দিনে করে কাজের (১/১৫) অংশ}

\emph{তারা একত্রে ১ দিনে করে = ((১/১০)+(১/১৫)) অংশ = (৩+২/৩০) = (৫/৩০) =
(১/৬) অংশ}

\emph{∴ তারা একত্রে পুরো কাজটি শেষ করতে পারবে ৬ দিনে।}
\end{quote}

\textbf{৩ দিনে একটি কাজের (১/২৭) অংশ শেষ হলে ঐ কাজের ৩ গুণ কাজ করতে কত
দিন লাগবে?}

ক) ৮১ দিন

খ) ৯ দিন

\textbf{\hl{গ) ২৪৩ দিন ✔}}

ঘ) ২৭ দিন

\textbf{\hl{✔ সঠিক উত্তর: গ) ২৪৩ দিন}}

\begin{quote}
\textbf{ব্যাখ্যা:} \emph{কাজের (১/২৭) অংশ সম্পন্ন করতে সময় লাগে = ৩ দিন}

\emph{∴ সম্পূর্ণ বা ১ অংশ কাজ করতে সময় লাগে = (৩ × ২৭) দিন = ৮১ দিন}

\emph{∴ সম্পূর্ণ কাজের ৩ গুণ কাজ করতে সময় লাগে = (৮১ × ৩) দিন = ২৪৩ দিন।}
\end{quote}

\textbf{৩ জন পুরুষ বা ৫ জন বালক একটি কাজ ২০ দিনে করতে পারে। তবে ৪ জন পুরুষ ও
১০ জন বালক ঐ কাজটি কত দিনে প্রতে পারবে?}

ক) ১০ দিন

খ) ৯ দিন

গ) ৮ দিন

\textbf{\hl{ঘ) ৬ দিন ✔}}

\textbf{\hl{✔ সঠিক উত্তর: ঘ) ৬ দিন}}

\begin{quote}
\textbf{ব্যাখ্যা:} \emph{এখানে, ৫ জন বালক = ৩ জন পুরুষ}

\emph{∴ ১ জন বালক = (৩/৫) জন পুরুষ}

\emph{∴ ১০ জন বালক = (৩×১০/৫) = ৬ জন পুরুষ}

\emph{এখন মোট পুরুষের সংখ্যা = ৪ + ৬ = ১০ জন}

\emph{৩ জন পুরুষ একটি কাজ করে = ২০ দিনে}

\emph{∴ ১ জন পুরুষ একটি কাজ করে = (২০ × ৩) দিনে}

\emph{∴ ১০ জন পুরুষ একটি কাজ করে = (২০×৩/১০) = ৬ দিনে।}
\end{quote}

\textbf{৩ জন পুরুষ বা ৯ জন বালক একটি কাজ ৬০ দিনে করতে পারে। ১১ জন পুরুষ ও ২৭
জন বালকের ঐ কাজটি করতে কতদিন লাগবে?}

ক) ৬ দিন

\textbf{\hl{খ) ৯ দিন ✔}}

গ) ১০ দিন

ঘ) ১২ দিন

\textbf{\hl{✔ সঠিক উত্তর: খ) ৯ দিন}}

\begin{quote}
\textbf{ব্যাখ্যা:} \emph{এখানে, ৯ জন বালক = ৩ জন পুরুষ}

\emph{∴ ১ জন বালক = (৩/৯) জন পুরুষ}

\emph{∴ ২৭ জন বালক = (৩×২৭/৯) = ৯ জন পুরুষ}

\emph{মোট পুরুষ = ১১ + ৯ = ২০ জন}

\emph{৩ জন পুরুষ কাজটি করে = ৬০ দিনে}

\emph{∴ ১ জন পুরুষ কাজটি করে = (৬০ × ৩) দিনে}

\emph{∴ ২০ জন পুরুষ কাজটি করে = (৬০×৩/২০) দিনে = ৯ দিনে।}
\end{quote}

\textbf{দুই ব্যক্তি একটি কাজ একত্রে ৮ দিনে করতে পারে। প্রথম ব্যক্তি কাজটি ১২
দিনে করতে পারে। দ্বিতীয় ব্যক্তি একা ঐ কাজটি কত দিনে করতে পারবে?}

ক) ২০ দিনে

খ) ২৫ দিনে

\textbf{\hl{গ) ২৪ দিনে ✔}}

ঘ) ৩০ দিনে

\textbf{\hl{✔ সঠিক উত্তর: গ) ২৪ দিনে}}

\begin{quote}
\textbf{ব্যাখ্যা:} \emph{দুই ব্যক্তি একত্রে ১ দিনে করে কাজের (১/৮) অংশ}

\emph{প্রথম ব্যক্তি একা ১ দিনে করে কাজের (১/১২) অংশ}

\emph{∴ দ্বিতীয় ব্যক্তি একা ১ দিনে করে = ((১/৮)−(১/১২)) অংশ = (৩−২/২৪) =
(১/২৪) অংশ}

\emph{∴ দ্বিতীয় ব্যক্তি একা ঐ কাজটি ২৪ দিনে করতে পারবে।}
\end{quote}

\textbf{যদি একটি কাজ ৯ জন লোক ১২ দিনে করতে পারে, অতিরিক্ত ৩ জন লোক নিয়োগ
করলে কাজটি কতদিনে শেষ হবে?}

ক) ৭ দিন

\textbf{\hl{খ) ৯ দিন ✔}}

গ) ৮ দিন

ঘ) ১০ দিন

\textbf{\hl{✔ সঠিক উত্তর: খ) ৯ দিন}}

\begin{quote}
\textbf{ব্যাখ্যা:} \emph{অতিরিক্ত ৩ জন নিয়োগ দেওয়ায় মোট লোকসংখ্যা = (৯ + ৩)
জন = ১২ জন}

\emph{৯ জন লোক কাজটি করে = ১২ দিনে}

\emph{∴ ১ জন লোক কাজটি করে = (৯ × ১২) দিনে}

\emph{∴ ১২ জন লোক কাজটি করে = (৯×১২/১২) দিনে = ৯ দিনে।}
\end{quote}

\textbf{নুরু একটি কাজ তিন ঘণ্টায় করতে পারে। নুরু এবং জনি একত্রে দুই ঘণ্টায় কাজটি
করতে পারলে, জনি কাজটি কত ঘণ্টায় করতে পারবে?}

ক) ৪

খ) ৭

গ) ৯

\textbf{\hl{ঘ) ৬ ✔}}

\textbf{\hl{✔ সঠিক উত্তর: ঘ) ৬}}

\begin{quote}
\textbf{ব্যাখ্যা:} \emph{নুরু ও জনি একত্রে ১ ঘণ্টায় করে কাজের (১/২) অংশ}

\emph{নুরু একা ১ ঘণ্টায় করে কাজের (১/৩) অংশ}

\emph{∴ জনি একা ১ ঘণ্টায় করে = ((১/২)−(১/৩)) অংশ = (৩−২/৬) = (১/৬) অংশ}

\emph{∴ জনি একা পুরো কাজটি ৬ ঘণ্টায় করতে পারবে।}
\end{quote}

\textbf{সনি ও পমি একত্রে একটি কাজ ১২ দিনে করতে পারে। সনি একা কাজটি ১৬
দিনে করতে পারে । কাজটি একা করতে পমির সনি অপেক্ষা কতদিন বেশি লাগবে?}

ক) ১৮ দিন

\textbf{\hl{খ) ৩২ দিন ✔}}

গ) ৪৮ দিন

ঘ) ৬০ দিন

\textbf{\hl{✔ সঠিক উত্তর: খ) ৩২ দিন}}

\begin{quote}
\textbf{ব্যাখ্যা:} \emph{সনি ও পমি একত্রে ১ দিনে করে কাজের (১/১২) অংশ}

\emph{সনি একা ১ দিনে করে কাজের (১/১৬) অংশ}

\emph{∴ পমি একা ১ দিনে করে = ((১/১২)−(১/১৬)) অংশ = (৪−৩/৪৮) = (১/৪৮)
অংশ}

\emph{∴ পমি একা সম্পূর্ণ কাজটি করতে পারে ৪৮ দিনে।}

\emph{∴ কাজটি করতে পমির সনি অপেক্ষা বেশি লাগবে = (৪৮ − ১৬) = ৩২ দিন।}
\end{quote}

\textbf{ক ও খ একত্রে একটি কাজ ১০ দিনে শেষ করতে পারে। খ একা কাজটি ১৪ দিনে
শেষ করতে পারলে ক একা কত দিনে কাজটি শেষ করতে পারবে?}

ক) ২৫ দিন

খ) ২৮ দিন

গ) ৩২ দিন

\textbf{\hl{ঘ) ৩৫ দিন ✔}}

\textbf{\hl{✔ সঠিক উত্তর: ঘ) ৩৫ দিন}}

\begin{quote}
\textbf{ব্যাখ্যা:} \emph{ক ও খ একত্রে ১ দিনে করে কাজের (১/১০) অংশ}

\emph{খ একা ১ দিনে করে কাজের (১/১৪) অংশ}

\emph{∴ ক একা ১ দিনে করে = ((১/১০)−(১/১৪)) অংশ = (৭−৫/৭০) = (২/৭০) =
(১/৩৫) অংশ}

\emph{∴ ক একা কাজটি ৩৫ দিনে শেষ করতে পারবে।}
\end{quote}

\textbf{দু\textquotesingle ব্যক্তি একত্রে একটি কাজ ৮ দিনে করতে পারে। ১ম
ব্যক্তি একাকী কাজটি ১২ দিনে করতে পারে। ২য় ব্যক্তি একাকী কাজটি কত দিনে করতে
পারে?}

ক) ২০ দিনে

খ) ২৫ দিনে

\textbf{\hl{গ) ২৪ দিনে ✔}}

ঘ) ৩০ দিনে

\textbf{\hl{✔ সঠিক উত্তর: গ) ২৪ দিনে}}

\begin{quote}
\textbf{ব্যাখ্যা:} \emph{দুই ব্যক্তি একত্রে ১ দিনে করে কাজের (১/৮) অংশ}

\emph{১ম ব্যক্তি একা ১ দিনে করে কাজের (১/১২) অংশ}

\emph{∴ ২য় ব্যক্তি একা ১ দিনে করে = ((১/৮)−(১/১২)) অংশ = (৩−২/২৪) = (১/২৪)
অংশ}

\emph{∴ ২য় ব্যক্তি একা কাজটি ২৪ দিনে করতে পারবে।}
\end{quote}

\textbf{ক ও খ একত্রে একটি কাজ ৮ দিনে করতে পারে। ঐ কাজ
\textquotesingle ক\textquotesingle{} একা ১২ দিনে করতে পারে। কাজটি করতে
\textquotesingle খ\textquotesingle{} এর কত দিন লাগবে?}

ক) ২২ দিনে

খ) ২০ দিনে

\textbf{\hl{গ) ২৪ দিনে ✔}}

ঘ) ২৬ দিনে

\textbf{\hl{✔ সঠিক উত্তর: গ) ২৪ দিনে}}

\begin{quote}
\textbf{ব্যাখ্যা:} \emph{ক ও খ একত্রে ১ দিনে করে কাজের (১/৮) অংশ}

\emph{ক একা ১ দিনে করে কাজের (১/১২) অংশ}

\emph{∴ খ একা ১ দিনে করে = ((১/৮)−(১/১২)) অংশ = (৩−২/২৪) = (১/২৪) অংশ}

\emph{∴ খ এর কাজটি করতে ২৪ দিন লাগবে।}
\end{quote}

\textbf{শায়েদ ২২ দিনে একটি কাজ শেষ করার জন্য ৪২ জন লোক নিয়োগ দিল। ১২ দিন
পর মাত্র এক তৃতীয়াংশ কাজ শেষ হল। অতিরিক্ত কোন লোক নেয়া না হলে অবশিষ্ট কাজ
শেষ করতে নির্ধারিত সময় অপেক্ষা কতদিন বেশি লাগবে?}

ক) ৩৬

খ) ২৪

গ) ১৬

\textbf{\hl{ঘ) ১৪ ✔}}

\textbf{\hl{✔ সঠিক উত্তর: ঘ) ১৪}}

\begin{quote}
\textbf{ব্যাখ্যা:} \emph{কাজটির (১/৩) অংশ সম্পন্ন হয় = ১২ দিনে}

\emph{∴ কাজটির ১ বা সম্পূর্ণ অংশ সম্পন্ন হয় = (১২ × ৩) = ৩৬ দিনে}

\emph{যেহেতু কাজটি ২২ দিনে শেষ করার কথা ছিল,}

\emph{∴ নির্ধারিত সময় অপেক্ষা সময় বেশি লাগবে = (৩৬ − ২২) = ১৪ দিন।}
\end{quote}

\textbf{মিটা একটি কাজ ১৫ মিনিটে করতে পারে, তার ছোট ভাইয়ের ঐ কাজটি করতে
সময় লাগে দ্বিগুণ। দুজনে মিলে কাজটি করলে কত মিনিটে কাজটি শেষ করতে পারবে?}

ক) ৫

খ) ৭ (১/২)

\textbf{\hl{গ) ১০ ✔}}

ঘ) ২ (১/২)

\textbf{\hl{✔ সঠিক উত্তর: গ) ১০}}

\begin{quote}
\textbf{ব্যাখ্যা:} \emph{মিটার সময় লাগে = ১৫ মিনিট}

\emph{ছোট ভাইয়ের সময় লাগে = (১৫ × ২) = ৩০ মিনিট}

\emph{দুজনে মিলে ১ মিনিটে করে = ((১/১৫)+(১/৩০)) অংশ = (২+১/৩০) = (৩/৩০) =
(১/১০) অংশ}

\emph{∴ দুজনে মিলে সম্পূর্ণ কাজটি করতে সময় লাগবে ১০ মিনিট।}
\end{quote}

\textbf{অপু, দীপু, নিপু একটি কাজ যথাক্রমে ৬, ১০, ১৫ দিনে করতে পারে। একত্রে
তারা কাজটি কত দিনে করতে পারবে?}

ক) ১২ দিনে

খ) ৯ দিনে

গ) ৬ দিনে

\textbf{\hl{ঘ) ৩ দিনে ✔}}

\textbf{\hl{✔ সঠিক উত্তর: ঘ) ৩ দিনে}}

\begin{quote}
\textbf{ব্যাখ্যা:} \emph{অপু, দীপু ও নিপু একত্রে ১ দিনে করে =
((১/৬)+(১/১০)+(১/১৫)) অংশ = (৫+৩+২/৩০) = (১০/৩০) = (১/৩) অংশ}

\emph{∴ তারা একত্রে পুরো কাজটি করতে পারবে ৩ দিনে।}
\end{quote}

\textbf{যদি একটি কাজ ৯ জন লোক ১২ দিনে করতে পারে, তবে ১২ জন লোক ঐ কাজ কত
দিনে শেষ করতে পারবে?}

ক) ১০

খ) ১২

\textbf{\hl{গ) ৯ ✔}}

ঘ) ৫

\textbf{\hl{✔ সঠিক উত্তর: গ) ৯}}

\begin{quote}
\textbf{ব্যাখ্যা:} \emph{৯ জন লোক ১টি কাজ করে = ১২ দিনে}

\emph{∴ ১ জন লোক ১টি কাজ করে = (১২ × ৯) দিনে}

\emph{∴ ১২ জন লোক ১টি কাজ করে = (১২×৯/১২) = ৯ দিনে।}
\end{quote}

\textbf{৪ জন পুরুষ বা ৬ জন স্ত্রীলোক একটি কাজ ১৬ দিনে শেষ করতে পারলে ৪ জন
পুরুষ ও ৬ জন স্ত্রীলোক একত্রে কাজটি কতদিনে শেষ করতে পারবে?}

\textbf{\hl{ক) ৮ ✔}}

খ) ১০

গ) ১২

ঘ) কোনোটিই নয়

\textbf{\hl{✔ সঠিক উত্তর: ক) ৮}}

\begin{quote}
\textbf{ব্যাখ্যা:} \emph{এখানে, ৪ জন পুরুষ = ৬ জন স্ত্রীলোক}

\emph{∴ ৪ জন পুরুষ + ৬ জন স্ত্রীলোক = (৬ + ৬) জন বা ১২ জন স্ত্রীলোক}

\emph{এখন, ৬ জন স্ত্রীলোক কাজটি করে = ১৬ দিনে}

\emph{∴ ১ জন স্ত্রীলোক কাজটি করে = (১৬ × ৬) দিনে}

\emph{∴ ১২ জন স্ত্রীলোক কাজটি করে = (৬×১৬/১২) = ৮ দিনে।}
\end{quote}

\textbf{৭ জন পুরুষ ও ৬ জন মহিলা একটি কাজ ৮ দিনে শেষ করতে পারে। ২ জন মহিলা
১ জন পুরুষের সমান কাজ করতে পারে। ৪ জন পুরুষ ঐ কাজ কয়দিনে শেষ করতে পারবে?}

ক) ১২

খ) ১৬

\textbf{\hl{গ) ২০ ✔}}

ঘ) কোনটিই নয়

\textbf{\hl{✔ সঠিক উত্তর: গ) ২০}}

\begin{quote}
\textbf{ব্যাখ্যা:} \emph{২ জন মহিলা = ১ জন পুরুষ}

\emph{∴ ৬ জন মহিলা = ৩ জন পুরুষ}

\emph{∴ ৭ জন পুরুষ ও ৬ জন মহিলা = (৭ + ৩) = ১০ জন পুরুষ}

\emph{১০ জন পুরুষ একটি কাজ শেষ করতে পারে = ৮ দিনে}

\emph{∴ ১ জন পুরুষ কাজটি শেষ করতে পারে = (১০ × ৮) দিনে}

\emph{∴ ৪ জন পুরুষ কাজটি শেষ করতে পারবে = (১০×৮/৪) = ২০ দিনে।}
\end{quote}

\textbf{১৬ জন শ্রমিক একটি কাজ ৩ ঘণ্টায় সম্পন্ন করতে পারে। উক্ত কাজটি সম্পন্ন
করতে ৫ জন শ্রমিকের কত সময় লাগবে?}

ক) ১৬ ঘণ্টা

খ) ১৫ ঘণ্টা

গ) ৭ (৩/৫) ঘণ্টা

\textbf{\hl{ঘ) ৯ (৩/৫) ঘণ্টা ✔}}

\textbf{\hl{✔ সঠিক উত্তর: ঘ) ৯ (৩/৫) ঘণ্টা}}

\begin{quote}
\textbf{ব্যাখ্যা:} \emph{১৬ জন শ্রমিক কাজটি করে = ৩ ঘণ্টায়}

\emph{∴ ১ জন শ্রমিক কাজটি করে = (১৬ × ৩) ঘণ্টায়}

\emph{∴ ৫ জন শ্রমিক কাজটি করে = (১৬×৩/৫) = (৪৮/৫) = ৯ (৩/৫) ঘণ্টায়।}
\end{quote}

\textbf{ক একটি কাজ ১০ দিনে করতে পারে এবং খ ঐ কাজটি ১৫ দিনে করতে পারে। ক
ও খ একত্রে ৫ দিন কাজ করার পর খ চলে গেল। বাকী কাজ ক একা কত দিনে করতে
পারবে?}

\textbf{\hl{ক) ১ (২/৩) দিন ✔}}

খ) ১ (১/৩) দিনে

গ) ২ দিন

ঘ) ১ দিন

\textbf{\hl{✔ সঠিক উত্তর: ক) ১ (২/৩) দিন}}

\begin{quote}
\textbf{ব্যাখ্যা:} \emph{ক ১ দিনে করে = (১/১০) অংশ কাজ}

\emph{খ ১ দিনে করে = (১/১৫) অংশ কাজ}

\emph{ক ও খ একত্রে ১ দিনে করে = ((১/১০)+(১/১৫)) = (৩+২/৩০) = (৫/৩০) =
(১/৬) অংশ কাজ}

\emph{∴ ক ও খ একত্রে ৫ দিনে করে = (৫/৬) অংশ কাজ}

\emph{বাকী কাজ = (১−(৫/৬)) অংশ = (১/৬) অংশ কাজ}

\emph{ক ১ অংশ কাজ করে = ১০ দিনে}

\emph{∴ ক (১/৬) অংশ কাজ করে = (১০/৬) দিনে = (৫/৩) দিনে = ১ (২/৩) দিন।}
\end{quote}

\textbf{৩০ জন লোক কোনো কাজ ২৪ দিনে করতে পারে। কাজ আরম্ভের ১২দিন পর ১৫ জন
লোক চলে গেলে বাকী লোক কত দিনে অবশিষ্ট কাজ সমাধা করতে পারবে?}

\textbf{\hl{ক) ২৪ দিন ✔}}

খ) ২০ দিন

গ) ১৫ দিন

ঘ) ১২ দিন

\textbf{\hl{✔ সঠিক উত্তর: ক) ২৪ দিন}}

\begin{quote}
\textbf{ব্যাখ্যা:} \emph{১২ দিন পর সময় বাকি থাকে = (২৪ − ১২) দিন = ১২ দিন}

\emph{এবং লোক বাকি থাকে = (৩০ − ১৫) জন = ১৫ জন}

\emph{৩০ জন লোক অবশিষ্ট কাজ করে = ১২ দিনে}

\emph{∴ ১ জন লোক অবশিষ্ট কাজ করে = (১২ × ৩০) দিনে}

\emph{∴ ১৫ জন লোক অবশিষ্ট কাজ করে = (১২×৩০/১৫) দিনে = ২৪ দিনে।}
\end{quote}

\textbf{জাইন একা ১২ দিনে একটি কাজ করতে পারে। সে কাজটির (২/৩) অংশ শেষ
করার পর অবশিষ্ট কাজটি হামিদ ৭ দিনে শেষ করে। কাজটির (৩/৭) অংশ শেষ করতে
হামিদের একার কত দিন লাগবে?}

ক) ৩

খ) ৬

\textbf{\hl{গ) ৯ ✔}}

ঘ) ৭

\textbf{\hl{✔ সঠিক উত্তর: গ) ৯}}

\begin{quote}
\textbf{ব্যাখ্যা:} \emph{জাইন (২/৩) অংশ কাজ করার পর অবশিষ্ট কাজ = (১−(২/৩))
= (১/৩) অংশ}

\emph{হামিদ (১/৩) অংশ কাজ করে = ৭ দিনে}

\emph{∴ হামিদ ১ বা সম্পূর্ণ অংশ কাজ করে = (৭ × ৩) = ২১ দিনে}

\emph{∴ হামিদ কাজটির (৩/৭) অংশ করে = (২১×৩/৭) = ৯ দিনে।}
\end{quote}

\textbf{তিন দিনে একটি কাজের (১/২৯) অংশ শেষ হলে ঐ কাজের তিন গুণ কাজ করতে
কত দিন লাগবে?}

ক) ২৯ দিন

খ) ৮৭ দিন

\textbf{\hl{গ) ২৬১ দিন ✔}}

ঘ) ৩০০ দিন

\textbf{\hl{✔ সঠিক উত্তর: গ) ২৬১ দিন}}

\begin{quote}
\textbf{ব্যাখ্যা:} \emph{কাজের (১/২৯) অংশ শেষ করতে সময় লাগে = ৩ দিন}

\emph{∴ ১ বা সম্পূর্ণ অংশ শেষ করতে সময় লাগে = (৩ × ২৯) = ৮৭ দিন}

\emph{∴ কাজটির ৩ গুণ কাজ করতে সময় লাগে = (৮৭ × ৩) দিন = ২৬১ দিন।}
\end{quote}

\textbf{৮ জন পুরুষ একটি কাজ ১২ দিনে করতে পারে। দুইজন পুরুষ চলে গেলে কাজটি
সমাধান করতে শতকরা কত দিন বেশি প্রয়োজন?}

ক) ৩১\%

খ) ৩১ (১/৩) \%

গ) ৩৩\%

\textbf{\hl{ঘ) ৩৩ (১/৩) \% ✔}}

\textbf{\hl{✔ সঠিক উত্তর: ঘ) ৩৩ (১/৩) \%}}

\begin{quote}
\textbf{ব্যাখ্যা:} \emph{২ জন পুরুষ চলে যাওয়ায় মোট পুরুষ সংখ্যা = (৮ − ২) জন =
৬ জন}

\emph{৮ জন পুরুষ কাজটি করে = ১২ দিনে}

\emph{∴ ১ জন পুরুষ কাজটি করে = (১২ × ৮) দিনে}

\emph{∴ ৬ জন পুরুষ কাজটি করে = (১২×৮/৬) দিনে = ১৬ দিনে}

\emph{পূর্বের চেয়ে সময় বেশি লাগে = (১৬ − ১২) = ৪ দিন}

\emph{∴ শতকরা সময় বেশি লাগে = ((৪×১০০/১২)) \% = (১০০/৩) \% = ৩৩ (১/৩)
\%}
\end{quote}

\textbf{৬ জন লোক ও ৮ জন বালক একটি কাজ ১০ দিনে শেষ করতে পারে। যদি ঐ একই
কাজ ২৬ জন লোক ও ৪৮ জন বালক ২ দিনে শেষ করতে পারে তাহলে ১৫ জন লোক ও ২০ জন
বালক ঐ কাজ কত দিনে শেষ করবে?}

\textbf{\hl{ক) ৪ দিন ✔}}

খ) ৫ দিন

গ) ৬ দিন

ঘ) ৭ দিন

\textbf{\hl{✔ সঠিক উত্তর: ক) ৪ দিন}}

\begin{quote}
\textbf{ব্যাখ্যা:} \emph{ধরি, ১ জন লোক = M এবং ১ জন বালক = B}

\emph{প্রশ্নানুযায়ী,}

\emph{১০(৬M + ৮B) = ২(২৬M + ৪৮B)}

\emph{⇒ ৬০M + ৮০B = ৫২M + ৯৬B}

\emph{⇒ ৮M = ১৬B}

\emph{∴ ১M = ২B (অর্থাৎ, ১ জন লোক = ২ জন বালক)}

\emph{প্রথম ক্ষেত্রে, ৬M + ৮B = (৬ × ২B) + ৮B = ২০B অথবা লোক হিসাবে: ৬M +
৪M = ১০M}

\emph{অতএব, ১০ জন লোক কাজটি করে = ১০ দিনে}

\emph{তৃতীয় ক্ষেত্রে, ১৫M + ২০B = ১৫M + ১০M = ২৫M (যেহেতু ২০B = ১০M)}

\emph{∴ ২৫ জন লোক কাজটি করবে = (১০×১০/২৫) = ৪ দিনে।}
\end{quote}

\textbf{ক, খ ও গ একটি কাজ যথাক্রমে ২০, ২৪ ও ৩০ দিনে সম্পন্ন করতে পারে। তারা
একত্রে ৬ দিনে কাজ করার পর খ ও গ চলে গেল। বাকী কাজ ক একা আর কত দিনে সম্পন্ন
করতে পারবে?}

ক) ২ দিন

খ) ৩ দিন

গ) ৪ দিন

\textbf{\hl{ঘ) ৫ দিন ✔}}

\textbf{\hl{✔ সঠিক উত্তর: ঘ) ৫ দিন}}

\begin{quote}
\textbf{ব্যাখ্যা:} \emph{ক ১ দিনে করে = (১/২০) অংশ কাজ}

\emph{খ ১ দিনে করে = (১/২৪) অংশ কাজ}

\emph{গ ১ দিনে করে = (১/৩০) অংশ কাজ}

\emph{ক, খ ও গ একত্রে ১ দিনে করে = ((১/২০)+(১/২৪)+(১/৩০)) = (৬+৫+৪/১২০) =
(১৫/১২০) = (১/৮) অংশ কাজ}

\emph{∴ তারা একত্রে ৬ দিনে করে = (৬/৮) = (৩/৪) অংশ কাজ}

\emph{বাকী কাজ = (১−(৩/৪)) = (১/৪) অংশ কাজ}

\emph{ক ১ অংশ কাজ করে = ২০ দিনে}

\emph{∴ ক (১/৪) অংশ কাজ করে = (২০/৪) দিনে = ৫ দিনে।}
\end{quote}

\textbf{দৈনিক ৮ ঘণ্টা কাজ করে ৪৭ জন কৃষক ৫৪ দিনে ৩২ বিঘা জমির ধান কাটতে
পারে। দৈনিক কত ঘণ্টা কাজ করলে ১৪১ জন কৃষক ৪৮ দিনে ৯৬ বিঘা জমির ধান কাটতে
পারবে?}

ক) ৭

খ) ৮

\textbf{\hl{গ) ৯ ✔}}

ঘ) ১০

\textbf{\hl{✔ সঠিক উত্তর: গ) ৯}}

\begin{quote}
\textbf{ব্যাখ্যা:} \emph{৪৭ জন কৃষক ৩২ বিঘা জমির ধান কাটতে ৫৪ দিনে দৈনিক
কাজ করে = ৮ ঘণ্টা}

\emph{∴ ১ জন কৃষক ১ বিঘা জমির ধান কাটতে ১ দিনে দৈনিক কাজ করে =
(৮×৪৭×৫৪/৩২) ঘণ্টা}

\emph{∴ ১৪১ জন কৃষক ৯৬ বিঘা জমির ধান কাটতে ৪৮ দিনে দৈনিক কাজ করবে =
(৮×৪৭×৫৪×৯৬/৩২×১৪১×৪৮) ঘণ্টা = ৯ ঘণ্টা।}
\end{quote}

\textbf{রহিম একটি কাজ ৫ দিনে এবং করিম তা ১০দিনে করতে পারে। ২জনে একত্রে
১দিনে কাজের কত অংশ করতে পারবে?}

\textbf{\hl{ক) (৩/১০) ✔}}

খ) (১/১৫)

গ) (২/১৫)

ঘ) (১/১০)

\textbf{\hl{✔ সঠিক উত্তর: ক) (৩/১০)}}

\begin{quote}
\textbf{ব্যাখ্যা:} \emph{রহিম ৫ দিনে করে = ১টি কাজ}

\emph{∴ রহিম ১ দিনে করে = (১/৫) অংশ কাজ}

\emph{করিম ১০ দিনে করে = ১টি কাজ}

\emph{∴ করিম ১ দিনে করে = (১/১০) অংশ কাজ}

\emph{∴ তারা একত্রে ১ দিনে করে = ((১/৫)+(১/১০)) = (২+১/১০) = (৩/১০) অংশ
কাজ।}
\end{quote}

\textbf{একজন ঠিকাদার একটি কাজের ৯ ভাগের ৫ অংশ সমাপ্ত করল। অপর ঠিকাদার আরও
এক-তৃতীয়াংশ কাজ করল। আর কত অংশ কাজ বাকি রইল?}

\textbf{\hl{ক) এক-নবমাংশ ✔}}

খ) চার-নবমাংশ

গ) আট-নবমাংশ

ঘ) দুই-তৃতীয়াংশ

\textbf{\hl{✔ সঠিক উত্তর: ক) এক-নবমাংশ}}

\begin{quote}
\textbf{ব্যাখ্যা:} \emph{প্রথম ঠিকাদার করে = (৫/৯) অংশ}

\emph{দ্বিতীয় ঠিকাদার করে = (১/৩) অংশ}

\emph{দুজন মোট করে = ((৫/৯)+(১/৩)) = (৫+৩/৯) = (৮/৯) অংশ}

\emph{∴ কাজ বাকি রইল = (১−(৮/৯)) = (১/৯) অংশ বা এক-নবমাংশ।}
\end{quote}

\textbf{\textquotesingle ক\textquotesingle{} যে কাজ x দিনে করতে পারে
\textquotesingle খ\textquotesingle{} সে কাজ 5x দিনে সম্পন্ন করে। একই সময়ে
\textquotesingle খ\textquotesingle,
\textquotesingle ক\textquotesingle-এর কতগুণ কাজ করে?}

ক) ৫ গুণ

\textbf{\hl{খ) (১/৫) গুণ ✔}}

গ) (৫/২) গুণ

ঘ) (২/৫) গুণ

\textbf{\hl{✔ সঠিক উত্তর: খ) (১/৫) গুণ}}

\begin{quote}
\textbf{ব্যাখ্যা:} \emph{\textquotesingle ক\textquotesingle{} 1 দিনে করে
কাজের (১/x) অংশ}

\emph{\textquotesingle খ\textquotesingle{} 1 দিনে করে কাজের (১/5x) অংশ}

\emph{∴ একই সময়ে \textquotesingle খ\textquotesingle,
\textquotesingle ক\textquotesingle-এর কাজ করবে = ((১/5x)/(১/x)) = (১/৫)
গুণ।}
\end{quote}

\textbf{৪ জন পুরুষ ৬ জন বালকের সমান কাজ করতে পারে। কতজন পুরুষ ২৪ জন বালকের
সমান কাজ করতে পারবে?}

ক) ১২ জন

\textbf{\hl{খ) ১৬ জন ✔}}

গ) ২৮ জন

ঘ) ৪৮ জন

\textbf{\hl{✔ সঠিক উত্তর: খ) ১৬ জন}}

\begin{quote}
\textbf{ব্যাখ্যা:} \emph{৬ জন বালকের কাজ = ৪ জন পুরুষের কাজ}

\emph{∴ ১ জন বালকের কাজ = (৪/৬) জন পুরুষের কাজ}

\emph{∴ ২৪ জন বালকের কাজ = (৪×২৪/৬) = ১৬ জন পুরুষের কাজ।}
\end{quote}

\textbf{একজন পুরুষ যে কাজ ১ দিনে করে, ঐ কাজ একজন স্ত্রীলোকের করতে ৩ দিন
লাগে। একটি কাজ ১৫ জন পুরুষ ১ দিনে করতে পারে। ঐ কাজ এক দিনে করতে কত জন
স্ত্রীলোক প্রয়োজন?}

ক) ৩০

\textbf{\hl{খ) ৪৫ ✔}}

গ) ৯০

ঘ) ১৩৫

\textbf{\hl{✔ সঠিক উত্তর: খ) ৪৫}}

\begin{quote}
\textbf{ব্যাখ্যা:} \emph{১ জন পুরুষের ১ দিনের কাজ = ১ জন স্ত্রীলোকের ৩ দিনের
কাজ}

\emph{অর্থাৎ, ১ জন পুরুষ = ৩ জন স্ত্রীলোক}

\emph{∴ ১৫ জন পুরুষের কাজ = (১৫ × ৩) = ৪৫ জন স্ত্রীলোকের কাজ}

\emph{∴ ঐ কাজ ১ দিনে করতে ৪৫ জন স্ত্রীলোক প্রয়োজন।}
\end{quote}

\textbf{১৫ জন লোক একটি কাজ ৪২ দিনে করতে পারে। ঐ কাজটি ১৮ দিনে শেষ করতে
কত জন লোক লাগবে?}

ক) ৩০ জন

খ) ৩২ জন

গ) ৩৩ জন

\textbf{\hl{ঘ) ৩৫ জন ✔}}

\textbf{\hl{✔ সঠিক উত্তর: ঘ) ৩৫ জন}}

\begin{quote}
\textbf{ব্যাখ্যা:} \emph{৪২ দিনে কাজটি করে = ১৫ জন লোক}

\emph{∴ ১ দিনে কাজটি করে = (৪২ × ১৫) জন লোক}

\emph{∴ ১৮ দিনে কাজটি করে = (৪২×১৫/১৮) জন লোক = ৩৫ জন লোক।}
\end{quote}

\textbf{একটি কাজ ১৫ জন লোকে ১০ দিনে করতে পারে। কত জন লোক ঐ কাজ ১ দিনে
সম্পন্ন করতে পারবে?}

ক) ১০০ জন

\textbf{\hl{খ) ১৫০ জন ✔}}

গ) ২০০ জন

ঘ) ২৫০ জন

\textbf{\hl{✔ সঠিক উত্তর: খ) ১৫০ জন}}

\begin{quote}
\textbf{ব্যাখ্যা:} \emph{১০ দিনে কাজটি শেষ করতে লোক লাগে = ১৫ জন}

\emph{∴ ১ দিনে কাজটি শেষ করতে লোক লাগবে = (১৫ × ১০) জন = ১৫০ জন।}
\end{quote}

\textbf{যদি ২০ জন লোক একটি কাজের অর্ধেক করতে পারে ৩০ দিনে তবে ঐ একই কাজ
৫০ দিনে করতে অতিরিক্ত কত জন লোক লাগবে?}

ক) ৮ জন

খ) ১০ জন

\textbf{\hl{গ) ৪ জন ✔}}

ঘ) ৬ জন

\textbf{\hl{✔ সঠিক উত্তর: গ) ৪ জন}}

\begin{quote}
\textbf{ব্যাখ্যা:} \emph{২০ জন লোক (১/২) অংশ করে = ৩০ দিনে}

\emph{∴ ২০ জন লোক ১ বা সম্পূর্ণ অংশ করে = (৩০ × ২) = ৬০ দিনে}

\emph{৬০ দিনে কাজটি সম্পন্ন করে = ২০ জন লোক}

\emph{∴ ১ দিনে কাজটি সম্পন্ন করে = (২০ × ৬০) জন লোক}

\emph{∴ ৫০ দিনে কাজটি সম্পন্ন করে = (২০×৬০/৫০) = ২৪ জন লোক}

\emph{∴ অতিরিক্ত লোক লাগবে = (২৪ − ২০) জন = ৪ জন।}
\end{quote}

\textbf{১২ জন লোক একটি কাজ ১৪ দিনে করতে পারে। কাজটি ৮ দিনে করতে হলে কতজন
নতুন লোক নিয়োগে করতে হবে?}

ক) ১০ জন

খ) ১২ জন

গ) ২১ জন

\textbf{\hl{ঘ) ৯ জন ✔}}

\textbf{\hl{✔ সঠিক উত্তর: ঘ) ৯ জন}}

\begin{quote}
\textbf{ব্যাখ্যা:} \emph{১৪ দিনে কাজটি করে = ১২ জন লোক}

\emph{∴ ১ দিনে কাজটি করে = (১২ × ১৪) জন লোক}

\emph{∴ ৮ দিনে কাজটি করে = (১২×১৪/৮) = ২১ জন লোক}

\emph{∴ নতুন লোক নিয়োগ করতে হবে = (২১ − ১২) জন = ৯ জন।}
\end{quote}

\textbf{If two typist can type two pages in two minutes, how many
typists will it take to type 18 pages in six minutes. (যদি দুইজন টাইপিস্ট
দুই মিনিটে দুই পেজ টাইপ করতে পারে। তাহলে ৬ মিনিটে ১৮ পেজ টাইপ করতে কতজন
টাইপিস্টের প্রয়োজন?)}

ক) 3

\textbf{\hl{খ) 6 ✔}}

গ) 9

ঘ) 18

\textbf{\hl{✔ সঠিক উত্তর: খ) 6}}

\begin{quote}
\textbf{ব্যাখ্যা:} \emph{২ পেজ ২ মিনিটে টাইপ করে = ২ জন টাইপিস্ট}

\emph{∴ ১ পেজ ১ মিনিটে টাইপ করে = (২×২/২) জন টাইপিস্ট}

\emph{∴ ১৮ পেজ ৬ মিনিটে টাইপ করে = (২×২×১৮/২×৬) = ৬ জন টাইপিস্ট।}
\end{quote}

\textbf{একজন ঠিকাদার ২৫ দিনে একটি রাস্তার কাজ শেষ করতে ৩০ জন লোক নিয়োগ
করলেন। কিন্তু ১৫ দিনে মাত্র রাস্তার অর্ধেক কাজ শেষ হয়। নির্ধারিত সময়ের মধ্যে
রাস্তার কাজ শেষ করতে হলে তাকে অতিরিক্ত কতজন লোক নিয়োগ করতে হবে?}

ক) ১২ জন

খ) ১৪ জন

\textbf{\hl{গ) ১৫ জন ✔}}

ঘ) ১৬ জন

\textbf{\hl{✔ সঠিক উত্তর: গ) ১৫ জন}}

\begin{quote}
\textbf{ব্যাখ্যা:} \emph{১৫ দিনে রাস্তার অর্ধেক কাজ শেষ করে = ৩০ জন লোক}

\emph{বাকি সময় = (২৫ − ১৫) দিন = ১০ দিন}

\emph{বাকি কাজ = অর্ধেক কাজ (যেহেতু কাজটি অর্ধেক হয়েছে)}

\emph{∴ ১ দিনে কাজ শেষ করতে লোক লাগে = (৩০ × ১৫) জন}

\emph{∴ ১০ দিনে শেষ করতে লোক লাগবে = (৩০×১৫/১০) = ৪৫ জন}

\emph{∴ অতিরিক্ত লোক নিয়োগ করতে হবে = (৪৫ − ৩০) = ১৫ জন।}
\end{quote}

\textbf{দৈনিক ৮ ঘণ্টা কাজ করে ৪২ জন লোক ৪১ দিনে ৮২০ মিটার রাস্তা পাকা করতে
পারে। দৈনিক ৭ ঘণ্টা কাজ করে কতজন লোক ৩৬ দিনে ১,২৩০ মিটার রাস্তা পাকা করতে
পারবে?}

ক) ৬৮

\textbf{\hl{খ) ৮২ ✔}}

গ) ৯৬

ঘ) ১০২

\textbf{\hl{✔ সঠিক উত্তর: খ) ৮২}}

\begin{quote}
\textbf{ব্যাখ্যা:} \emph{৪১ দিনে দৈনিক ৮ ঘণ্টা কাজ করে ৮২০ মিটার রাস্তা
বানায় = ৪২ জন}

\emph{∴ ১ দিনে দৈনিক ১ ঘণ্টা কাজ করে ১ মিটার রাস্তা বানায় = (৪২×৪১×৮/৮২০)
জন}

\emph{∴ ৩৬ দিনে দৈনিক ৭ ঘণ্টা কাজ করে ১২৩০ মিটার রাস্তা বানাবে =
(৪২×৪১×৮×১২৩০/৮২০×৩৬×৭) = ৮২ জন।}
\end{quote}

\textbf{এক ঠিকাদার ১৯২০ মিটার দীর্ঘ রাস্তা ১২০ দিনে নির্মাণ করে দেয়ার জন্য
১৬০ জন শ্রমিক নিয়োগ করল। ২৪ দিন পর, মাত্র (১/৮) অংশ কাজ সম্পন্ন হলো।
নির্ধারিত সময়ের মধ্যে কাজ শেষ করতে হলে অতিরিক্ত কতজন লোক নিয়োগ দিতে হবে?}

ক) ২৮০

খ) ১৬০

গ) ১৮০

\textbf{\hl{ঘ) ১২০ ✔}}

\textbf{\hl{✔ সঠিক উত্তর: ঘ) ১২০}}

\begin{quote}
\textbf{ব্যাখ্যা:} \emph{অবশিষ্ট সময় = (১২০ − ২৪) দিন = ৯৬ দিন}

\emph{অবশিষ্ট কাজ = (১−(১/৮)) অংশ = (৭/৮) অংশ}

\emph{২৪ দিনে (১/৮) অংশ কাজ করে = ১৬০ জন}

\emph{∴ ১ দিনে (১/৮) অংশ কাজ করে = (১৬০ × ২৪) জন}

\emph{∴ ১ দিনে ১ অংশ কাজ করে = (১৬০ × ২৪ × ৮) জন}

\emph{∴ ৯৬ দিনে ১ অংশ কাজ করে = (১৬০×২৪×৮/৯৬) জন}

\emph{∴ ৯৬ দিনে (৭/৮) অংশ কাজ করে = (১৬০×২৪×৮×৭/৯৬×৮) = ২৮০ জন}

\emph{∴ অতিরিক্ত শ্রমিক লাগবে = (২৮০ − ১৬০) = ১২০ জন।}
\end{quote}

\textbf{একটি চৌবাচ্চার ৮,০০০ লিটার পানি ধরে। চৌবাচ্চাটির দৈর্ঘ্য ২.৫৬ মিটার
এবং প্রস্থ ১.২৫ মিটার হলে গভীরতা কত?}

ক) ১.২৫ মি:

\textbf{\hl{খ) ২.৫ মি: ✔}}

গ) ২.৭৫ মি:

ঘ) ৩.৫ মি:

\textbf{\hl{✔ সঠিক উত্তর: খ) ২.৫ মি:}}

\begin{quote}
\textbf{ব্যাখ্যা:} \emph{আমরা জানি, ১,০০০ লিটার = ১ ঘনমিটার}

\emph{∴ ৮,০০০ লিটার = ৮ ঘনমিটার}

\emph{আয়তন = দৈর্ঘ্য × প্রস্থ × গভীরতা}

\emph{⇒ ৮ = ২.৫৬ × ১.২৫ × গভীরতা}

\emph{⇒ ৮ = ৩.২ × গভীরতা}

\emph{∴ গভীরতা = (৮/৩.২) = ২.৫ মিটার।}
\end{quote}

\textbf{A tap can fill a tank in 6 hours. After half the tank is filled,
three more similar taps are opened. What is the total time taken to fill
the tank completely?/ একটি কল ৬ ঘণ্টার মধ্যে একটি ট্যাংক পূরণ করতে পারে।
অর্ধেক ট্যাংক ভরা ভর্তি হওয়ার পর, তিনটি আরো অনুরূপ কল খোলা হয়। সম্পূর্ণ ট্যাংকটি
পূর্ণ হতে মোট কত সময় লাগবে?}

ক) 3 hrs 15 min

\textbf{\hl{খ) 3 hrs 45 min ✔}}

গ) 4 hrs

ঘ) 4 hrs 15 min

\textbf{\hl{✔ সঠিক উত্তর: খ) 3 hrs 45 min}}

\begin{quote}
\textbf{ব্যাখ্যা:} \emph{১টি কল দ্বারা অর্ধেক ( (১/২) ) ট্যাংক পূর্ণ হতে সময়
লাগে = (৬ ÷ ২) = ৩ ঘণ্টা।}

\emph{অর্ধেক পূর্ণ হওয়ার পর আরও ৩টি অনুরূপ কল খোলা হলে মোট কলের সংখ্যা হয় = (১
+ ৩) = ৪টি।}

\emph{১টি কল ৬ ঘণ্টায় পূর্ণ করে = ১ বা সম্পূর্ণ অংশ}

\emph{∴ ৪টি কল ১ ঘণ্টায় পূর্ণ করে = (৪×(১/৬)) অংশ = (২/৩) অংশ}

\emph{এখন, ৪টি কল দ্বারা অবশিষ্ট (১/২) অংশ পূর্ণ হতে সময় লাগে =
((১/২)÷(২/৩)) ঘণ্টা = (৩/৪) ঘণ্টা = ৪৫ মিনিট।}

\emph{∴ সম্পূর্ণ ট্যাংকটি পূর্ণ হতে মোট সময় লাগবে = ৩ ঘণ্টা + ৪৫ মিনিট = 3 hrs
45 min।}
\end{quote}

\textbf{১টি চৌবাচ্চার (৩/৫) ভাগ পূরণ হতে ৭ ঘণ্টা লাগে। চৌবাচ্চাটির বাকি অংশ
পূরণ হতে আর কত সময় লাগবে?}

ক) ৫ ঘণ্টা ২০ মিনিট

\textbf{\hl{খ) ৪ ঘণ্টা ৪০ মিনিট ✔}}

গ) ৪ ঘণ্টা ২০ মিনিট

ঘ) কোনোটিই নয়

\textbf{\hl{✔ সঠিক উত্তর: খ) ৪ ঘণ্টা ৪০ মিনিট}}

\begin{quote}
\textbf{ব্যাখ্যা:} \emph{চৌবাচ্চাটির বাকি অংশ = (১−(৩/৫)) = (২/৫) অংশ}

\emph{(৩/৫) ভাগ পূরণ হতে সময় লাগে = ৭ ঘণ্টা}

\emph{∴ ১ বা সম্পূর্ণ ভাগ পূরণ হতে সময় লাগে = (৭×৫/৩) ঘণ্টা}

\emph{∴ (২/৫) ভাগ পূরণ হতে সময় লাগে = (৭×৫×২/৩×৫) = (১৪/৩) ঘণ্টা = ৪ (২/৩)
ঘণ্টা = ৪ ঘণ্টা ( (২/৩) × ৬০) মিনিট = ৪ ঘণ্টা ৪০ মিনিট।}
\end{quote}

\textbf{একটি চৌবাচ্চার দৈর্ঘ্য ৩ মিটার, প্রস্থ ২ মিটার ও উচ্চতা ৪ মিটার। এতে
কত লিটার বিশুদ্ধ পানি ধরবে?}

ক) ২৪০০ লিটার

\textbf{\hl{খ) ২৪০০০ লিটার ✔}}

গ) ২৪০ লিটার

ঘ) ২৪০০০০ লিটার

\textbf{\hl{✔ সঠিক উত্তর: খ) ২৪০০০ লিটার}}

\begin{quote}
\textbf{ব্যাখ্যা:} \emph{চৌবাচ্চার আয়তন = দৈর্ঘ্য × প্রস্থ × উচ্চতা}

\emph{= (৩ × ২ × ৪) ঘনমিটার = ২৪ ঘনমিটার}

\emph{আমরা জানি, ১ ঘনমিটার = ১,০০০ লিটার}

\emph{∴ ২৪ ঘনমিটার = (২৪ × ১,০০০) লিটার = ২৪,০০০ লিটার।}
\end{quote}

\textbf{একটি পানির ট্যাঙ্কে দুটি নল আছে। প্রথম নলটি খুলে দিলে ট্যাঙ্কটি ২০
ঘণ্টায় পূর্ণ হয়। দ্বিতীয় নল দ্বারা পূর্ণ ট্যাঙ্কটি ৩০ ঘণ্টায় খালি হয়। দুইটি নল
একসঙ্গে খুলে দিলে খালি ট্যাঙ্ক কত সময়ে পূর্ণ হবে?}

ক) ৪৮ ঘণ্টা

খ) ৪০ ঘণ্টা

গ) ৫০ ঘণ্টা

\textbf{\hl{ঘ) ৬০ ঘণ্টা ✔}}

\textbf{\hl{✔ সঠিক উত্তর: ঘ) ৬০ ঘণ্টা}}

\begin{quote}
\textbf{ব্যাখ্যা:} \emph{প্রথম নলটি ২০ ঘণ্টায় পূর্ণ করে ১টি ট্যাংক}

\emph{∴ প্রথম নলটি ১ ঘণ্টায় পূর্ণ করে = (১/২০) অংশ}

\emph{দ্বিতীয় নলটি ৩০ ঘণ্টায় খালি করে ১টি ট্যাংক}

\emph{∴ দ্বিতীয় নলটি ১ ঘণ্টায় খালি করে = (১/৩০) অংশ}

\emph{দুটি নল এক সাথে ১ ঘণ্টায় পূর্ণ করে = ( (১/২০) − (১/৩০) ) অংশ =
(৩−২/৬০) অংশ = (১/৬০) অংশ}

\emph{দুটি নল এক সাথে (১/৬০) অংশ পূর্ণ করে ১ ঘণ্টায়}

\emph{∴ দুটি নল এক সাথে ১ বা সম্পূর্ণ অংশ পূর্ণ করে ৬০ ঘণ্টায়।}

\emph{শর্ট টেকনিক:}

\emph{দুটি নলের ১ টি দ্বারা পূর্ণকরণ এবং অপরটি দ্বারা খালিকরণ হলে---}

\emph{ট্যাংক পূর্ণ হতে প্রয়োজনীয় সময় T = (m×n/m−n)}

\emph{= (৩০×২০/৩০−২০) = (৬০০/১০) = ৬০ ঘণ্টা।}

\emph{(এখানে, m = ২য় নল দ্বারা ব্যয়িত সময়, n = ১ম নল দ্বারা ব্যয়িত সময়)।}
\end{quote}

\textbf{একটি চৌবাচ্চা একটি নল দ্বারা ১০ ঘণ্টায় পূর্ণ হয়। তাতে একটি ছিদ্র থাকায়
পূর্ণ হতে ১৫ ঘণ্টা লাগে। ছিদ্র দ্বারা চৌবাচ্চাটি খালি হতে কত সময় লাগবে?}

ক) ২০ ঘণ্টা

খ) ২৫ ঘণ্টা

\textbf{\hl{গ) ৩০ ঘণ্টা ✔}}

ঘ) ৩৫ ঘণ্টা

\textbf{\hl{✔ সঠিক উত্তর: গ) ৩০ ঘণ্টা}}

\begin{quote}
\textbf{ব্যাখ্যা:} \emph{নলটি ১০ ঘণ্টায় পূর্ণ করে ১টি চৌবাচ্চা}

\emph{∴ নলটি ১ ঘণ্টায় পূর্ণ করে চৌবাচ্চাটির (১/১০) অংশ}

\emph{কিন্তু ছিদ্র থাকায় ১ ঘণ্টায় পূর্ণ হয় চৌবাচ্চাটির = (১/১৫) অংশ}

\emph{∴ ছিদ্র দ্বারা ১ ঘণ্টায় খালি হয় = ( (১/১০) − (১/১৫) ) অংশ = (৩−২/৩০)
অংশ = (১/৩০) অংশ}

\emph{ছিদ্র দ্বারা চৌবাচ্চাটির (১/৩০) অংশ খালি হতে সময় লাগে ১ ঘণ্টা}

\emph{∴ সম্পূর্ণ অংশ খালি হতে সময় লাগে = (৩০ × ১) = ৩০ ঘণ্টা।}
\end{quote}

\textbf{একটি পিপায় দুটি নল সংযুক্ত আছে। প্রথম নলটি খুলে দিলে খালি পিপাটি ২০
মিনিটে পূর্ণ হয়, দ্বিতীয় নলটি খুলে দিলে পরিপূর্ণ পিপাটি ৩০ মিনিটে খালি হয়।
দুইটি নলই একসঙ্গে খুলে দিলে খালি পিপাটি কত সময়ে পূর্ণ হবে?}

ক) ৮০ মিনিটে

খ) ৯০ মিনিটে

গ) ৭০ মিনিটে

\textbf{\hl{ঘ) ৬০ মিনিটে ✔}}

\textbf{\hl{✔ সঠিক উত্তর: ঘ) ৬০ মিনিটে}}

\begin{quote}
\textbf{ব্যাখ্যা:} \emph{প্রথম নলটি ২০ মিনিটে পূর্ণ করে ১টি পিপা}

\emph{∴ প্রথম নলটি ১ মিনিটে পূর্ণ করে = (১/২০) অংশ}

\emph{দ্বিতীয় নলটি ৩০ মিনিটে খালি করে ১টি পিপা}

\emph{∴ দ্বিতীয় নলটি ১ মিনিটে খালি করে = (১/৩০) অংশ}

\emph{দুটি নল এক সাথে ১ মিনিটে পূর্ণ করে = ( (১/২০) − (১/৩০) ) অংশ =
(৩−২/৬০) অংশ = (১/৬০) অংশ}

\emph{দুটি নল এক সাথে (১/৬০) অংশ পূর্ণ করে ১ মিনিটে}

\emph{∴ দুটি নল এক সাথে ১ বা সম্পূর্ণ অংশ পূর্ণ করে ৬০ মিনিটে।}

\emph{শর্ট টেকনিক:}

\emph{দুটি নলের ১ টি দ্বারা পূর্ণকরণ এবং অপরটি দ্বারা খালিকরণ হলে---}

\emph{ট্যাংক পূর্ণ হতে প্রয়োজনীয় সময় T = (m×n/m−n)}

\emph{= (৩০×২০/৩০−২০) = (৬০০/১০) = ৬০ মিনিট।}
\end{quote}

\textbf{একটি চৌবাচ্চা দুটি নল দ্বারা যথাক্রমে ২০ এবং ৩০ মিনিটে পূর্ণ হয়। দুটি
নল এক সঙ্গে খুলে দিলে চৌবাচ্চাটি কত সময়ে পূর্ণ হবে?}

ক) ৬ মিনিটে

খ) ৯ মিনিটে

\textbf{\hl{গ) ১২ মিনিটে ✔}}

ঘ) ১৫ মিনিটে

\textbf{\hl{✔ সঠিক উত্তর: গ) ১২ মিনিটে}}

\begin{quote}
\textbf{ব্যাখ্যা:} \emph{১ম নল দ্বারা, ১ মিনিটে পূর্ণ হয় = (১/২০) অংশ}

\emph{২য় নল দ্বারা, ১ মিনিটে পূর্ণ হয় = (১/৩০) অংশ}

\emph{দুটি নল দ্বারা, ১ মিনিটে পূর্ণ হয় = ( (১/২০) + (১/৩০) ) অংশ = (৩+২/৬০)
অংশ = (৫/৬০) অংশ = (১/১২) অংশ}

\emph{(১/১২) অংশ পূর্ণ হয় = ১ মিনিটে}

\emph{∴ সমস্ত অংশ পূর্ণ হয় = ১২ মিনিটে।}

\emph{শর্ট টেকনিক:}

\emph{যখন দুইটি নল দ্বারাই চৌবাচ্চা পূর্ণ হয় তখন প্রয়োজনীয় সময় = (m×n/m+n)}

\emph{= (৩০×২০/৩০+২০) = (৬০০/৫০) = ১২ মিনিট।}
\end{quote}

\textbf{দুটি নল দ্বারা একটি চৌবাচ্চা যথাক্রমে ১০ ও ১৫ ঘণ্টায় পানি পূর্ণ করে। নল
দুটি একত্রে খোলা রাখলে চৌবাচ্চাটি কতক্ষণে পানি পূর্ণ হবে?}

\textbf{\hl{ক) ৬ ঘণ্টায় ✔}}

খ) ৫ ঘণ্টায়

গ) ৪ ঘণ্টায়

ঘ) ২ ঘণ্টায়

\textbf{\hl{✔ সঠিক উত্তর: ক) ৬ ঘণ্টায়}}

\begin{quote}
\textbf{ব্যাখ্যা:} \emph{১ম নল দ্বারা, ১ ঘণ্টায় পূর্ণ হয় = (১/১০) অংশ}

\emph{২য় নল দ্বারা, ১ ঘণ্টায় পূর্ণ হয় = (১/১৫) অংশ}

\emph{দুটি নল দ্বারা, ১ ঘণ্টায় পূর্ণ হয় = ( (১/১০) + (১/১৫) ) অংশ = (৩+২/৩০)
অংশ = (৫/৩০) অংশ = (১/৬) অংশ}

\emph{(১/৬) অংশ পূর্ণ হয় = ১ ঘণ্টায়}

\emph{∴ সমস্ত অংশ পূর্ণ হয় = ৬ ঘণ্টায়।}

\emph{শর্ট টেকনিক:}

\emph{যখন দুইটি নল দ্বারাই চৌবাচ্চা পূর্ণ হয় তখন প্রয়োজনীয় সময় = (m×n/m+n)}

\emph{= (১৫×১০/১৫+১০) = (১৫০/২৫) = ৬ ঘণ্টায়।}
\end{quote}

\textbf{একটি চৌবাচ্চা তিনটি নল দিয়ে যথাক্রমে ১০, ১২, ও ১৫ ঘণ্টায় পূর্ণ হতে
পারে। তিনটি নল একসঙ্গে খুলে দিলে চৌবাচ্চার অর্ধেক পূর্ণ হতে কত সময় লাগবে?}

ক) ৬ ঘণ্টা

খ) ৪ ঘণ্টা

গ) ৩ ঘণ্টা

\textbf{\hl{ঘ) ২ ঘণ্টা ✔}}

\textbf{\hl{✔ সঠিক উত্তর: ঘ) ২ ঘণ্টা}}

\begin{quote}
\textbf{ব্যাখ্যা:} \emph{তিনটি নল একত্রে ১ ঘণ্টায় পূর্ণ করে =
((১/১০)+(১/১২)+(১/১৫)) অংশ = (৬+৫+৪/৬০) = (১৫/৬০) = (১/৪) অংশ}

\emph{∴ তিনটি নল একত্রে সম্পূর্ণ চৌবাচ্চা পূর্ণ করে = ৪ ঘণ্টায়}

\emph{∴ চৌবাচ্চার অর্ধেক ( (১/২) ) অংশ পূর্ণ হতে সময় লাগবে = (৪/২) = ২
ঘণ্টা।}
\end{quote}

\textbf{১ম নল দ্বারা চৌবাচ্চাটি ৪ মিনিটে পূর্ণ হয় এবং ২য় নল দ্বারা ১২ মিনিটে
পূর্ণ হয়। নল দুটি একত্রে খুলে দিলে খালি চৌবাচ্চাটি কতক্ষণে পূর্ণ হবে?}

ক) ৬

\textbf{\hl{খ) ৩ ✔}}

গ) ৪

ঘ) ৫

\textbf{\hl{✔ সঠিক উত্তর: খ) ৩}}

\begin{quote}
\textbf{ব্যাখ্যা:} \emph{১ম নল দ্বারা ১ মিনিটে পূর্ণ হয় = (১/৪) অংশ}

\emph{২য় নল দ্বারা ১ মিনিটে পূর্ণ হয় = (১/১২) অংশ}

\emph{দুটি নল একত্রে ১ মিনিটে পূর্ণ করে = ((১/৪)+(১/১২)) অংশ = (৩+১/১২) =
(৪/১২) = (১/৩) অংশ}

\emph{∴ নল দুটি একত্রে খালি চৌবাচ্চাটি পূর্ণ করতে সময় নেবে ৩ মিনিট।}
\end{quote}

\textbf{সম্পূর্ণ খালি একটি চৌবাচ্চা একটি পাইপ দিয়ে ৫ ঘণ্টায় সম্পূর্ণ ভর্তি করা
যায়। দ্বিতীয় একটি পাইপ দিয়ে চৌবাচ্চাটি পূর্ণ করতে ৩ ঘণ্টা সময় লাগে। দুটি পাইপ
একসাথে ব্যবহার করে চৌবাচ্চাটির (২/৩) অংশ পূর্ণ করতে কত সময় লাগবে?}

ক) (৮/১৫) ঘণ্টা

খ) (৩/৪) ঘণ্টা

\textbf{\hl{গ) (৫/৪) ঘণ্টা ✔}}

ঘ) (১/২) ঘণ্টা

\textbf{\hl{✔ সঠিক উত্তর: গ) (৫/৪) ঘণ্টা}}

\begin{quote}
\textbf{ব্যাখ্যা:} \emph{১ম পাইপ ১ ঘণ্টায় পূর্ণ করে = (১/৫) অংশ}

\emph{২য় পাইপ ১ ঘণ্টায় পূর্ণ করে = (১/৩) অংশ}

\emph{দুটি পাইপ একত্রে ১ ঘণ্টায় পূর্ণ করে = ((১/৫)+(১/৩)) = (৩+৫/১৫) = (৮/১৫)
অংশ}

\emph{∴ ২টি পাইপ ১ বা সম্পূর্ণ অংশ পূর্ণ করে = (১৫/৮) ঘণ্টায়}

\emph{∴ চৌবাচ্চাটির (২/৩) অংশ পূর্ণ করতে সময় লাগবে = ((১৫/৮)×(২/৩)) = (৫/৪)
ঘণ্টা।}
\end{quote}

\textbf{একটি পানির ট্যাঙ্কে দুটি নল আছে। প্রথম নলটি খুলে দিলে ট্যাঙ্কটি ১০
ঘণ্টায় পানিতে পূর্ণ হয় এবং দ্বিতীয় নলটি খুলে দিলে পানিপূর্ণ ট্যাঙ্কটি ১৫ ঘণ্টায়
খালি হয়। দুটি নল একসঙ্গে খুলে দিলে খালি ট্যাঙ্কটি কত ঘণ্টায় পূর্ণ হবে?}

ক) ২০ ঘণ্টা

খ) ২৪ ঘণ্টা

গ) ২৮ ঘণ্টা

\textbf{\hl{ঘ) ৩০ ঘণ্টা ✔}}

\textbf{\hl{✔ সঠিক উত্তর: ঘ) ৩০ ঘণ্টা}}

\begin{quote}
\textbf{ব্যাখ্যা:} \emph{১ম নল দ্বারা ১ ঘণ্টায় পূর্ণ হয় = (১/১০) অংশ}

\emph{২য় নল দ্বারা ১ ঘণ্টায় খালি হয় = (১/১৫) অংশ}

\emph{দুটি নল একসঙ্গে খোলা থাকলে ১ ঘণ্টায় পূর্ণ হয় = ((১/১০)−(১/১৫)) অংশ =
(৩−২/৩০) = (১/৩০) অংশ}

\emph{∴ খালি ট্যাঙ্কটি সম্পূর্ণ পূর্ণ হতে সময় লাগবে ৩০ ঘণ্টা।}
\end{quote}

\textbf{দুটি নল দ্বারা একটি চৌবাচ্চা যথাক্রমে ২০ মিনিট ও ৩০ মিনিটে পূর্ণ করা
যায়। চৌবাচ্চাটি খালি থাকা অবস্থায় দুটি নলই এক সঙ্গে খুলে দেয়া হলো। প্রথম নলটি
কখন বন্ধ করলে মোট ১৮ মিনিটে চৌবাচ্চাটি পানি পূর্ণ হবে?}

\textbf{\hl{ক) ৮ মিনিট পরে ✔}}

খ) ৬ মিনিট পরে

গ) ১০ মিনিট পরে

ঘ) ৪ মিনিট পরে

\textbf{\hl{✔ সঠিক উত্তর: ক) ৮ মিনিট পরে}}

\begin{quote}
\textbf{ব্যাখ্যা:} \emph{ধরি, ১ম নলটি x মিনিট পর বন্ধ করা হয়েছিল।}

\emph{২য় নলটি পুরো ১৮ মিনিটই চালু ছিল।}

\emph{২য় নল ১৮ মিনিটে পূর্ণ করে = ((১৮/৩০)) অংশ = (৩/৫) অংশ}

\emph{∴ অবশিষ্ট কাজ/পানি = (১−(৩/৫)) অংশ = (২/৫) অংশ}

\emph{এই (২/৫) অংশ ১ম নল দ্বারা পূর্ণ হয়েছে।}

\emph{১ম নল ১ অংশ পূর্ণ করে = ২০ মিনিটে}

\emph{∴ ১ম নল (২/৫) অংশ পূর্ণ করে = (২০×(২/৫)) মিনিটে = ৮ মিনিটে}

\emph{∴ ১ম নলটি ৮ মিনিট পর বন্ধ করতে হবে।}
\end{quote}

\textbf{দুইটি নল দ্বারা একটি চৌবাচ্চা ৬ মিনিটে পূর্ণ হয়। নল দুটি একসঙ্গে খুলে
দেওয়ার ৩ মিনিট পরে একটি নল বন্ধ করে দেওয়ায় চৌবাচ্চাটি পূর্ণ হতে আরও ৫ মিনিট
সময় লাগল। প্রত্যেক নল দ্বারা পৃথকভাবে চৌবাচ্চাটি পূর্ণ হতে কত সময় লাগবে?}

ক) ১৫ মি. ২০ মি.

\textbf{\hl{খ) ১০ মি., ১৫ মি, ✔}}

গ) ১০ মি., ২০ মি.

ঘ) ১৫ মি, ৩০ মি.

\textbf{\hl{✔ সঠিক উত্তর: খ) ১০ মি., ১৫ মি,}}

\begin{quote}
\textbf{ব্যাখ্যা:} \emph{দুই নল দ্বারা ৩ মিনিটে পূর্ণ হয় = (৩/৬) = (১/২) অংশ}

\emph{অবশিষ্ট অংশ = (১−(১/২)) = (১/২) অংশ}

\emph{অপর নল দ্বারা ৫ মিনিটে পূর্ণ হয় = (১/২) অংশ}

\emph{∴ ঐ নল দ্বারা ১ বা সম্পূর্ণ অংশ পূর্ণ হয় = (৫ × ২) = ১০ মিনিটে}

\emph{অন্য নল দ্বারা ১ মিনিটে পূর্ণ হয় = ((১/৬)−(১/১০)) = (৫−৩/৩০) = (২/৩০)
= (১/১৫) অংশ}

\emph{∴ অন্য নল দ্বারা সম্পূর্ণ চৌবাচ্চা পূর্ণ হয় = ১৫ মিনিটে।}

\emph{∴ নল দুটি দ্বারা পৃথকভাবে সময় লাগবে ১০ মিনিট ও ১৫ মিনিট।}
\end{quote}

\textbf{দুইটি নল দিয়ে একটি চৌবাচ্চা ৮ মিনিটে পূর্ণ হয়। নল দুটি খুলে দেয়ার ৬
মিনিট পর প্রথম নলটি বন্ধ করে দেওয়ায় চৌবাচ্চাটির পূর্ণ হতে আরও ৬ মিনিট লাগলে
দ্বিতীয় নলটি দিয়ে চৌবাচ্চাটি এককভাবে পূর্ণ করতে কত মিনিট লাগবে?}

ক) ১২

খ) ১৬

গ) ২০

\textbf{\hl{ঘ) ২৪ ✔}}

\textbf{\hl{✔ সঠিক উত্তর: ঘ) ২৪}}

\begin{quote}
\textbf{ব্যাখ্যা:} \emph{নল দুটি একত্রে ৬ মিনিটে পূর্ণ করে = (৬/৮) = (৩/৪)
অংশ}

\emph{অবশিষ্ট অংশ = (১−(৩/৪)) = (১/৪) অংশ}

\emph{দ্বিতীয় নল ৬ মিনিটে পূর্ণ করে = (১/৪) অংশ}

\emph{∴ দ্বিতীয় নলটি দ্বারা সম্পূর্ণ চৌবাচ্চা এককভাবে পূর্ণ হতে সময় লাগবে = (৬ ×
৪) = ২৪ মিনিট।}
\end{quote}

\textbf{একটি বানর ৯২ ফুট উঁচু একটা তৈলাক্ত বাঁশ বেয়ে উপরে উঠতে লাগলো। বানরটি
প্রথম মিনিটে ১২ ফুট ওঠে, কিন্তু দ্বিতীয় মিনিটে ৮ ফুট নেমে বাঁশের মাথায় উঠতে
বানরটির কত মিনিট সময় লাগবে?}

\textbf{\hl{ক) ৪১ ✔}}

খ) ৪৫

গ) ৪৩

ঘ) ৪২

\textbf{\hl{✔ সঠিক উত্তর: ক) ৪১}}

\begin{quote}
\textbf{ব্যাখ্যা:} \emph{বানরটি ১ম মিনিটে ১২ ফুট ওঠে এবং ২য় মিনিটে ৮ ফুট
নেমে যায়।}

\emph{⇒ ২ মিনিটে প্রকৃতপক্ষে ওঠে = (১২ − ৮) = ৪ ফুট।}

\emph{বাঁশের মোট দৈর্ঘ্য = ৯২ ফুট।}

\emph{বানরটি শেষ লাফে (১ মিনিটে) ১২ ফুট উঠে চূড়ায় পৌঁছে যাবে এবং আর নিচে
নামবে না।}

\emph{⇒ শেষ লাফের পূর্বে অতিক্রম করতে হবে = (৯২ − ১২) = ৮০ ফুট।}

\emph{এখন, ৪ ফুট উঠতে সময় লাগে = ২ মিনিট}

\emph{⇒ ১ ফুট উঠতে সময় লাগে = (২/৪) মিনিট}

\emph{⇒ ৮০ ফুট উঠতে সময় লাগে = (২×৮০/৪) = ৪০ মিনিট।}

\emph{⇒ ৮০ ফুট ওঠার পর শেষ ১২ ফুট উঠায় মোট সময় লাগবে = (৪০ + ১) = ৪১
মিনিট।}

\emph{সংক্ষিপ্ত নিয়ম:}

\emph{প্রয়োজনীয় সময় = (মোট দৈর্ঘ্য−১ম মিনিটে যতটুকু ওঠে/১ম মিনিটে যতটুকু ওঠে−২য়
মিনিটে যতটুকু নেমে যায়)×২+১}

\emph{= (৯২−১২/১২−৮)×২+১}

\emph{= (৮০/৪)×২+১}

\emph{= (২০ × ২) + ১ = ৪১ মিনিট।}
\end{quote}

\textbf{একটি বানর একটি তৈলাক্ত বাঁশ বেয়ে উপরে উঠতে লাগলো। বানরটি যদি ১
মিনিটে ৫ মিটার ওঠে এবং পরবর্তী মিনিটে ১ মিটার নেমে পড়ে, তবে ২৫ মিটার উঁচু
বাঁশের মাথায় উঠতে বানরের কত সময় লাগবে?}

ক) ৯ মিনিট

খ) ১০ মিনিট

\textbf{\hl{গ) ১১ মিনিট ✔}}

ঘ) ১২ মিনিট

\textbf{\hl{✔ সঠিক উত্তর: গ) ১১ মিনিট}}

\begin{quote}
\textbf{ব্যাখ্যা:} \emph{বানরটি ১ম মিনিটে ৫ মিটার ওঠে এবং ২য় মিনিটে ১
মিটার নেমে যায়।}

\emph{⇒ ২ মিনিটে প্রকৃতপক্ষে ওঠে = (৫ − ১) = ৪ মিটার।}

\emph{শেষ লাফে (১ মিনিটে) ৫ মিটার উঠলে বাঁশের মাথায় পৌঁছাবে।}

\emph{⇒ শেষ লাফের পূর্বে অতিক্রম করতে হবে = (২৫ − ৫) = ২০ মিটার।}

\emph{এখন, ৪ মিটার উঠতে সময় লাগে = ২ মিনিট}

\emph{⇒ ২০ মিটার উঠতে সময় লাগে = (২×২০/৪) = ১০ মিনিট।}

\emph{⇒ মোট সময় লাগবে = (১০ + ১) = ১১ মিনিট।}

\emph{সংক্ষিপ্ত নিয়ম:}

\emph{প্রয়োজনীয় সময় = (মোট দৈর্ঘ্য−১ম মিনিটে যতটুকু ওঠে/১ম মিনিটে যতটুকু ওঠে−২য়
মিনিটে যতটুকু নেমে যায়)×২+১}

\emph{= (২৫−৫/৫−১)×২+১}

\emph{= (২০/৪)×২+১}

\emph{= (৫ × ২) + ১ = ১১ মিনিট।}
\end{quote}

\textbf{একটি বান একটি তৈলাক্ত বাঁশ বেয়ে উঠতে লাগল। বানরটির যদি ১ মিনিটে ৪
মিটার উঠে এবং পরবর্তী মিনেটে ১ মিটার নেমে পড়ে তবে ১৬ মিটার উচু বাঁশের মাথায়
উঠতে বানরের কত সময় লাগবে।}

ক) ৭ মিনিট

খ) ৮ মিনিট

\textbf{\hl{গ) ৯ মিনিট ✔}}

ঘ) ১০ মিনিট

\textbf{\hl{✔ সঠিক উত্তর: গ) ৯ মিনিট}}

\begin{quote}
\textbf{ব্যাখ্যা:} \emph{বানরটি ১ম মিনিটে ৪ মিটার ওঠে এবং ২য় মিনিটে ১
মিটার নেমে যায়।}

\emph{⇒ ২ মিনিটে প্রকৃতপক্ষে ওঠে = (৪ − ১) = ৩ মিটার।}

\emph{শেষ লাফে (১ মিনিটে) ৪ মিটার ওঠে।}

\emph{⇒ শেষ লাফের পূর্বে উঠতে হবে = (১৬ − ৪) = ১২ মিটার।}

\emph{এখন, ৩ মিটার উঠতে সময় লাগে = ২ মিনিট}

\emph{⇒ ১২ মিটার উঠতে সময় লাগে = (২×১২/৩) = ৮ মিনিট।}

\emph{⇒ মোট প্রয়োজনীয় সময় = (৮ + ১) = ৯ মিনিট।}

\emph{সংক্ষিপ্ত নিয়ম:}

\emph{প্রয়োজনীয় সময় = (মোট দৈর্ঘ্য−১ম মিনিটে যতটুকু ওঠে/১ম মিনিটে যতটুকু ওঠে−২য়
মিনিটে যতটুকু নেমে যায়)×২+১}

\emph{= (১৬−৪/৪−১)×২+১}

\emph{= (১২/৩)×২+১}

\emph{= (৪ × ২) + ১ = ৯ মিনিট।}
\end{quote}

\textbf{একটি বানর একটি তৈলাক্ত বাঁশ বেয়ে উঠছে। বানরটি ১ মিনিটে ৪ ফুট উপরে
উঠে, কিন্তু পরের মিনিটে ৩ ফুট নিচে নেমে যায়। বাঁশটি ২০ গজ লম্বা হলে এর শেষ
প্রান্তে উঠতে বানরটির কত সময় লাগবে?}

ক) ১১১ মিনিট

খ) ১১২ মিনিট

\textbf{\hl{গ) ১১৩ মিনিট ✔}}

ঘ) ১১৫ মিনিট

\textbf{\hl{✔ সঠিক উত্তর: গ) ১১৩ মিনিট}}

\begin{quote}
\textbf{ব্যাখ্যা:} \emph{আমরা জানি, ১ গজ = ৩ ফুট}

\emph{∴ বাঁশটির মোট দৈর্ঘ্য = (২০ × ৩) = ৬০ ফুট।}

\emph{বানরটি ২ মিনিটে প্রকৃতপক্ষে ওঠে = (৪ − ৩) = ১ ফুট।}

\emph{শেষ লাফে (১ মিনিটে) বানরটি ৪ ফুট ওঠে।}

\emph{⇒ শেষ লাফের পূর্বে অতিক্রম করতে হবে = (৬০ − ৪) = ৫৬ ফুট।}

\emph{১ ফুট উঠতে সময় লাগে = ২ মিনিট}

\emph{⇒ ৫৬ ফুট উঠতে সময় লাগে = (৫৬ × ২) = ১১২ মিনিট।}

\emph{⇒ মোট সময় লাগবে = (১১২ + ১) = ১১৩ মিনিট।}

\emph{সংক্ষিপ্ত নিয়ম:}

\emph{প্রয়োজনীয় সময় = (মোট দৈর্ঘ্য−১ম মিনিটে যতটুকু ওঠে/১ম মিনিটে যতটুকু ওঠে−২য়
মিনিটে যতটুকু নেমে যায়)×২+১}

\emph{= (৬০−৪/৪−৩)×২+১}

\emph{= (৫৬/১)×২+১}

\emph{= ১১২ + ১ = ১১৩ মিনিট।}
\end{quote}

\textbf{একজন মাঝি স্রোতের অনুকূলে ২ ঘণ্টায় ৫ মাইল যায় এবং ৪ ঘণ্টায় প্রাথমিক
অবস্থানে ফিরে আসে। তার মোট ভ্রমণে প্রতি ঘণ্টার গড় বেগ কত?}

ক) (৫/৬)

\textbf{\hl{খ) (৫/৩) ✔}}

গ) (১৫/৮)

ঘ) (১৫/৪)

\textbf{\hl{✔ সঠিক উত্তর: খ) (৫/৩)}}

\begin{quote}
\textbf{ব্যাখ্যা:} \emph{যাওয়ার দূরত্ব = ৫ মাইল এবং আসার দূরত্ব = ৫ মাইল}

\emph{মোট দূরত্ব = (৫ + ৫) মাইল = ১০ মাইল}

\emph{মোট সময় = (২ + ৪) ঘণ্টা = ৬ ঘণ্টা}

\emph{∴ প্রতি ঘণ্টার গড় বেগ = (যাওয়ার দূরত্ব+আসার দূরত্ব/মোট সময়) = (১০/৬) =
(৫/৩) মাইল/ঘণ্টা।}
\end{quote}

\textbf{একটি নৌকা স্রোতের অনুকূলে ঘণ্টায় ৮ কি.মি. এবং স্রোতের প্রতিকূলে ৪
কি.মি. যায়। নৌকার বেগ কত?}

ক) ২ কি.মি.

খ) ৩ কি.মি.

গ) ৪ কি.মি.

\textbf{\hl{ঘ) ৬ কি.মি. ✔}}

\textbf{\hl{✔ সঠিক উত্তর: ঘ) ৬ কি.মি.}}

\begin{quote}
\textbf{ব্যাখ্যা:} \emph{নৌকার বেগ = (অনুকূলে বেগ+প্রতিকূলে বেগ/২)}

\emph{= (৮+৪/২) = (১২/২) = ৬ কি.মি./ঘণ্টা।}
\end{quote}

\textbf{স্রোতের অনুকূলে একটি নৌকা ৫ ঘণ্টায় ৫০ কিলোমিটার পথ অতিক্রম করে। স্থির
পানিতে ঐ নৌকার গতিবেগ ঘণ্টায় ৮ কিলোমিটার হলে স্রোতের গতিবেগ কত?}

\textbf{\hl{ক) ২ কি.মি. ✔}}

খ) ৩ কি.মি.

গ) ৪ কি.মি.

ঘ) ৫ কি.মি.

\textbf{\hl{✔ সঠিক উত্তর: ক) ২ কি.মি.}}

\begin{quote}
\textbf{ব্যাখ্যা:} \emph{স্রোতের অনুকূলে নৌকার গতিবেগ = (৫০/৫) = ১০
কি.মি./ঘণ্টা}

\emph{আমরা জানি, অনুকূলে গতিবেগ = নৌকার বেগ + স্রোতের বেগ}

\emph{⇒ ১০ = ৮ + স্রোতের বেগ}

\emph{∴ স্রোতের গতিবেগ = (১০ − ৮) = ২ কি.মি./ঘণ্টা।}
\end{quote}

\textbf{If a boat goes 7 km upstream in 42 minutes and the speed of the
stream is 3 kmph, then the speed of the boat in still water is:/ উজানে
একটি নৌকা ৪২ মিনিটে ৭ কি.মি দূরত্ব অতিক্রম করে এবং স্রোতের বেগ ৩ কি.মি./ঘণ্টা
হলে স্থির পানিতে নৌকার বেগ কত হবে?}

ক) 4.2km/hr

খ) 9km/hr

\textbf{\hl{গ) 13 km/hr ✔}}

ঘ) 21 km/hr

\textbf{\hl{✔ সঠিক উত্তর: গ) 13 km/hr}}

\begin{quote}
\textbf{ব্যাখ্যা:} \emph{নৌকার উজানে (প্রতিকূলে) গতিবেগ = (৭×৬০/৪২) = ১০
কি.মি./ঘণ্টা}

\emph{ধরি, স্থির পানিতে নৌকার গতিবেগ = x কি.মি./ঘণ্টা}

\emph{স্রোতের বেগ = ৩ কি.মি./ঘণ্টা}

\emph{প্রশ্নানুযায়ী,}

\emph{x − ৩ = ১০}

\emph{⇒ x = ১০ + ৩ = ১৩}

\emph{∴ স্থির পানিতে নৌকার বেগ ১৩ কি.মি./ঘণ্টা।}
\end{quote}

\textbf{দাঁড় বেয়ে একটি নৌকা স্রোতের অনুকূলে ঘণ্টায় যায় 15 কিমি এবং স্রোতের
প্রতিকূলে ঘণ্টায় যায় 5 কি.মি.। স্রোতের বেগ নির্ণয় কর।}

\textbf{\hl{ক) 5 কি.মি./ঘণ্টা ✔}}

খ) 6 কি.মি./ ঘণ্টা

গ) 7 কি.মি./ঘণ্টা

ঘ) 8 কি.মি./ ঘণ্টা

\textbf{\hl{✔ সঠিক উত্তর: ক) 5 কি.মি./ঘণ্টা}}

\begin{quote}
\textbf{ব্যাখ্যা:} \emph{আমরা জানি,}

\emph{স্রোতের বেগ = (অনুকূলে বেগ−প্রতিকূলে বেগ/২)}

\emph{= (১৫−৫/২) = (১০/২) = ৫ কি.মি./ঘণ্টা।}
\end{quote}

\textbf{কোন নৌকার গতি স্রোতের অনুকূলে ঘণ্টায় ১০ কিলোমিটার এবং স্রোতের প্রতিকূলে
ঘণ্টায় ২ কিলোমিটার। স্রোতের বেগ কত?}

\textbf{\hl{ক) ৪ কি.মি. ✔}}

খ) ৩ কি.মি.

গ) ৮ কি.মি.

ঘ) ১২ কি.মি.

\textbf{\hl{✔ সঠিক উত্তর: ক) ৪ কি.মি.}}

\begin{quote}
\textbf{ব্যাখ্যা:} \emph{নৌকার বেগ + স্রোতের বেগ = ১০ কি.মি.}

\emph{নৌকার বেগ − স্রোতের বেগ = ২ কি.মি.}

\emph{(বিয়োগ করে) ২ × স্রোতের বেগ = ৮ কি.মি.}

\emph{∴ স্রোতের বেগ = (৮/২) = ৪ কি.মি.}
\end{quote}

\textbf{নৌকা ও স্রোতের বেগ ঘণ্টায় যথাক্রমে ১০ ও ৫ কি. মি.। নদী পথে ৪৫ কি.
মি. দীর্ঘ পথ একবার অতিক্রম করে ফিরে আসতে কত ঘণ্টা সময় লাগবে?}

ক) ৯ ঘন্টা

\textbf{\hl{খ) ১২ ঘন্টা ✔}}

গ) ১০ ঘন্টা

ঘ) ১৮ ঘন্টা

\textbf{\hl{✔ সঠিক উত্তর: খ) ১২ ঘন্টা}}

\begin{quote}
\textbf{ব্যাখ্যা:} \emph{স্রোতের অনুকূলে,}

\emph{নৌকা ও স্রোতের বেগ = (১০ + ৫) = ১৫ কি.মি.}

\emph{৪৫ কি.মি. যেতে সময় লাগে = (৪৫/১৫) = ৩ ঘণ্টা}

\emph{স্রোতের প্রতিকূলে,}

\emph{নৌকা ও স্রোতের বেগ = (১০ − ৫) = ৫ কি.মি.}

\emph{৪৫ কি.মি. আসতে সময় লাগে = (৪৫/৫) = ৯ ঘণ্টা}

\emph{∴ যেতে ও আসতে মোট সময় লাগে = (৯ + ৩) = ১২ ঘণ্টা।}
\end{quote}

\textbf{লঞ্চ ও স্রোতের গতিবেগ যথাক্রমে ঘণ্টায় ১৮ কিমি ও ৬ কিমি। নদীপথে ৭২
কিমি অতিক্রম করে পুনরায় যাত্রাস্থানে ফিরে আসতে সময় লাগবে-}

ক) ৬ ঘন্টা

খ) ৮ ঘন্টা

\textbf{\hl{গ) ৯ ঘন্টা ✔}}

ঘ) ১২ ঘন্টা

\textbf{\hl{✔ সঠিক উত্তর: গ) ৯ ঘন্টা}}

\begin{quote}
\textbf{ব্যাখ্যা:} \emph{স্রোতের অনুকূলে,}

\emph{লঞ্চ ও স্রোতের বেগ = (১৮ + ৬) = ২৪ কি.মি.}

\emph{৭২ কি.মি. যেতে সময় লাগে = (৭২/২৪) = ৩ ঘণ্টা}

\emph{স্রোতের প্রতিকূলে,}

\emph{লঞ্চ ও স্রোতের বেগ = (১৮ − ৬) = ১২ কি.মি.}

\emph{৭২ কি.মি. আসতে সময় লাগে = (৭২/১২) = ৬ ঘণ্টা}

\emph{∴ যেতে ও আসতে মোট সময় লাগে = (৩ + ৬) = ৯ ঘণ্টা।}
\end{quote}

\textbf{স্রোতের বিপরীতে একটি নৌকা ৫২ মিনিটে ১৩ কিমি যেতে পারে। স্রোতের বেগ
৪ কিমি/ঘন্টা। স্থির পানিতে নৌকার বেগ কত?}

\textbf{\hl{ক) ১৯ কিমি/ ঘন্টা ✔}}

খ) ২৩ কিমি/ ঘন্টা

গ) ১৩ কিমি/ঘন্টা

ঘ) ১১ কিমি / ঘন্টা

\textbf{\hl{✔ সঠিক উত্তর: ক) ১৯ কিমি/ ঘন্টা}}

\begin{quote}
\textbf{ব্যাখ্যা:} \emph{৫২ মিনিটে যায় = ১৩ কিমি}

\emph{১ মিনিটে যায় = (১৩/৫২) কিমি}

\emph{∴ ৬০ মিনিটে (১ ঘণ্টায়) যায় = ( (১৩×৬০/৫২) ) = ১৫ কিমি}

\emph{∴ স্থির পানিতে নৌকার বেগ = (১৫ + ৪) = ১৯ কিমি/ঘন্টা।}
\end{quote}

\textbf{একটি নৌকা স্রোতের অনুকূলে ও প্রতিকূলে ৩২ মাইল পথ যেতে যথাক্রমে ৪ ঘণ্টা
ও ৮ ঘণ্টা সময় লাগে। নৌকা ও স্রোতের বেগ নির্ণয় করুন।}

\textbf{\hl{ক) নৌকা ৬ কি.মি./ঘণ্টা, স্রোত ২ কি.মি./ঘণ্টা ✔}}

খ) নৌকা ৮ কি.মি./ঘণ্টা, স্রোত ৪ কি.মি./ঘণ্টা

গ) নৌকা ৪ কি.মি./ঘণ্টা, স্রোত ২ কি.মি./ঘণ্টা

ঘ) নৌকা ১০ কি.মি./ঘণ্টা, স্রোত ৪ কি.মি./ঘণ্টা

\textbf{\hl{✔ সঠিক উত্তর: ক) নৌকা ৬ কি.মি./ঘণ্টা, স্রোত ২ কি.মি./ঘণ্টা}}

\begin{quote}
\textbf{ব্যাখ্যা:} \emph{অনুকূলে গতিবেগ = (৩২/৪) = ৮ মাইল/ঘণ্টা}

\emph{প্রতিকূলে গতিবেগ = (৩২/৮) = ৪ মাইল/ঘণ্টা}

\emph{∴ নৌকার গতিবেগ = (৮+৪/২) = ৬ মাইল/ঘণ্টা (বা কি.মি./ঘণ্টা)}

\emph{∴ স্রোতের গতিবেগ = (৮−৪/২) = ২ মাইল/ঘণ্টা (বা কি.মি./ঘণ্টা)।}
\end{quote}

\textbf{একটি নৌকা স্রোতের প্রতিকূলে যে গতিবেগে চলে স্রোতের অনুকূলে ঐ গতিবেগের
৫ গুণ গতিবেগে যেতে পারে। স্থির পানিতে নৌকার গতিবেগ ঘণ্টায় ৬ কিমি হলে, ঘণ্টায়
স্রোতের গতিবেগ কত?}

\textbf{\hl{ক) ৪ কিমি./ঘণ্টা ✔}}

খ) ৫ কিমি/ঘণ্টা

গ) ৬ কিমি/ঘণ্টা

ঘ) ৭ কিমি/ঘণ্টা

\textbf{\hl{✔ সঠিক উত্তর: ক) ৪ কিমি./ঘণ্টা}}

\begin{quote}
\textbf{ব্যাখ্যা:} \emph{ধরি, স্রোতের গতিবেগ = x কিমি/ঘণ্টা}

\emph{স্রোতের অনুকূলে নৌকার কার্যকরী গতিবেগ = (৬ + x) কিমি/ঘণ্টা}

\emph{স্রোতের প্রতিকূলে নৌকার কার্যকরী গতিবেগ = (৬ − x) কিমি/ঘণ্টা}

\emph{প্রশ্নমতে,}

\emph{৬ + x = ৫ (৬ − x)}

\emph{⇒ ৬ + x = ৩০ − ৫x}

\emph{⇒ ৬x = ২৪}

\emph{∴ x = ৪ কিমি/ঘণ্টা।}
\end{quote}

\textbf{এক ব্যক্তি সাঁতার কেটে স্রোতের প্রতিকূলে ১ মিটার/ঘণ্টা বেগে ৪ মিটার
দূরত্ব অতিক্রম করে আবার স্রোতের অনুকূলে ৪ মিটার/ঘণ্টা বেগে পূর্বের স্থানে ফিরে এল।
তার গড় গতিবেগ কত?}

ক) 0.8 mph

\textbf{\hl{খ) 1.6 mph ✔}}

গ) 2.4 mph

ঘ) 3.2 mph

\textbf{\hl{✔ সঠিক উত্তর: খ) 1.6 mph}}

\begin{quote}
\textbf{ব্যাখ্যা:} \emph{স্রোতের প্রতিকূলে ৪ মিটার যেতে সময় লাগে = (৪ ÷ ১) =
৪ ঘণ্টা}

\emph{স্রোতের অনুকূলে ৪ মিটার যেতে সময় লাগে = (৪ ÷ ৪) = ১ ঘণ্টা}

\emph{মোট অতিক্রান্ত দূরত্ব = (৪ + ৪) মিটার = ৮ মিটার}

\emph{মোট সময় = (৪ + ১) ঘণ্টা = ৫ ঘণ্টা}

\emph{∴ গড় গতিবেগ = (৮/৫) মিটার/ঘণ্টা = ১.৬ মিটার/ঘণ্টা (1.6 mph)।}
\end{quote}

\textbf{এক ব্যক্তি স্রোতের অনুকূলে নৌকা বেয়ে ঘণ্টায় ১০ কি.মি. বেগে চলে কোন
স্থানে গেল এবং ঘণ্টায় ৬ কি. মি. বেগে স্রোতের প্রতিকূলে চলে যাত্রারম্ভের স্থানে
ফিরে এল। যাতায়াতে তার গড় গতিবেগ কত?}

\textbf{\hl{ক) ৭ (১/২) কি. মি. ✔}}

খ) ৫ (১/২) কি. মি.

গ) ৮ কি. মি.

ঘ) ৭ কি. মি.

\textbf{\hl{✔ সঠিক উত্তর: ক) ৭ (১/২) কি. মি.}}

\begin{quote}
\textbf{ব্যাখ্যা:} \emph{ধরি, স্থানটির দূরত্ব d কিমি।}

\emph{অনুকূলে সময় লাগে = (d/১০) ঘণ্টা}

\emph{প্রতিকূলে সময় লাগে = (d/৬) ঘণ্টা}

\emph{মোট দূরত্ব = 2d কিমি}

\emph{মোট সময় = ((d/১০)+(d/৬)) = (৩d+৫d/৩০) = (৮d/৩০) = (৪d/১৫) ঘণ্টা}

\emph{∴ গড় গতিবেগ = 2d÷(৪d/১৫) = 2d×(১৫/৪d) = (১৫/২) = ৭ (১/২)
কি.মি./ঘণ্টা।}
\end{quote}

\textbf{দুটি বাস ঘণ্টায় ২০ মাইল বেগে একই সময়ে গাবতলী হতে আরিচা রওয়ানা হল।
সাভার পৌছার পর একটি বাস থেমে গেল। অপর বাসটি চলতে থাকল। আধ ঘন্টা পর থেমে
থাকা বাসটি ঘণ্টায় ২৫ মাইল বেগে চলতে থাকল। সাভার থেকে কত দূরে বাস দুটি মিলিত
হবে।}

ক) ৪০ মাইল

\textbf{\hl{খ) ৫০ মাইল ✔}}

গ) ৬০ মাইল

ঘ) ৫৫ মাইল

\textbf{\hl{✔ সঠিক উত্তর: খ) ৫০ মাইল}}

\begin{quote}
\textbf{ব্যাখ্যা:} \emph{প্রথমে চলমান বাসটির আধা ঘণ্টার দূরত্ব হিসাব করা যাক:}

\emph{বাসটি ১ ঘণ্টায় যায় ২০ মাইল}

\emph{⇒ (১/২) ঘণ্টায় (আধা ঘণ্টায়) যায় = (২০/২) = ১০ মাইল।}

\emph{ধরি, সাভার থেকে \textquotesingle ক\textquotesingle{} মাইল দূরে বাস
দুটি মিলিত হবে।}

\emph{আধা ঘণ্টা পর থেমে থাকা বাসটি ঘণ্টায় ২৫ মাইল বেগে চলতে শুরু করে এবং
প্রথম বাসটি ইতোমধ্যেই ১০ মাইল এগিয়ে গেছে।}

\emph{থামা বাসটির \textquotesingle ক\textquotesingle{} মাইল যেতে সময়
লাগে = (ক/২৫) ঘণ্টা}

\emph{চলতে থাকা বাসটির বাকি (ক − ১০) মাইল যেতে সময় লাগে = (ক − ১০/২০)
ঘণ্টা}

\emph{প্রশ্নমতে,}

\emph{(ক/২৫) = (ক − ১০/২০)}

\emph{⇒ ২৫(ক − ১০) = ২০ক}

\emph{⇒ ২৫ক − ২৫০ = ২০ক}

\emph{⇒ ৫ক = ২৫০}

\emph{∴ ক = ৫০}

\emph{⇒ সাভার থেকে ৫০ মাইল দূরে বাস দুটি মিলিত হবে।}

\emph{সংক্ষিপ্ত নিয়ম:}

\emph{দূরত্ব = (১ম বেগ×২য় বেগ/২য় বেগ−১ম বেগ)×সময়}

\emph{= (২০×২৫/২৫−২০)×(১/২)}

\emph{= (৫০০/৫)×(১/২)}

\emph{= ১০০ × (১/২) = ৫০ মাইল।}
\end{quote}

\textbf{ক ঘণ্টায় ১০ কি. মি. এবং খ ঘণ্টায় ১৫ কি. মি বেগে একই সময় একই স্থান
থেকে রাজশাহীর পথে রওয়ানা হল। ক ১০.১০ মিনিটের সময় এবং খ ৯.৪০ মিনিটের সময়
রাজশাহী পৌছল। রওয়ানা হওয়ার স্থান থেকে রাজশাহীর দূরত্ব কত কি. মি.?}

ক) ২০ কি.মি.

খ) ২৫ কি.মি.

\textbf{\hl{গ) ১৫ কি.মি. ✔}}

ঘ) ২৮ কি.মি.

\textbf{\hl{✔ সঠিক উত্তর: গ) ১৫ কি.মি.}}

\begin{quote}
\textbf{ব্যাখ্যা:} \emph{ক ও খ এর পৌঁছানোর সময়ের পার্থক্য = ১০.১০ মিনিট −
৯.৪০ মিনিট = ৩০ মিনিট বা (১/২) ঘণ্টা।}

\emph{ধরি, রওয়ানা হওয়ার স্থান থেকে রাজশাহীর দূরত্ব
\textquotesingle ক\textquotesingle{} কি.মি.}

\emph{ক-এর সময় লাগে = (ক/১০) ঘণ্টা}

\emph{খ-এর সময় লাগে = (ক/১৫) ঘণ্টা}

\emph{প্রশ্নমতে,}

\emph{(ক/১০)−(ক/১৫)=(১/২)}

\emph{⇒ (৩ক − ২ক/৩০)=(১/২)}

\emph{⇒ (ক/৩০)=(১/২)}

\emph{⇒ ২ক = ৩০}

\emph{∴ ক = ১৫ কি.মি.}

\emph{সংক্ষিপ্ত নিয়ম:}

\emph{দূরত্ব = (১ম বেগ×২য় বেগ/২য় বেগ−১ম বেগ)×সময়ের পার্থক্য}

\emph{= (১০×১৫/১৫−১০)×(১/২)}

\emph{= (১৫০/৫)×(১/২)}

\emph{= ৩০ × (১/২) = ১৫ কি.মি.}
\end{quote}

\textbf{৪ কি. মি. / ঘণ্টা বেগে চললে কোন স্থানে পৌছাতে যে সময় লাগে, ৫ কি.
মি. / ঘণ্টা বেগে চললে তার চেয়ে (১/২) ঘণ্টা কম লাগে। স্থানটির দূরত্ব কত?}

ক) ৮ কি.মি.

\textbf{\hl{খ) ১০ কি.মি. ✔}}

গ) ১২ কি.মি.

ঘ) ১১ কি.মি.

\textbf{\hl{✔ সঠিক উত্তর: খ) ১০ কি.মি.}}

\begin{quote}
\textbf{ব্যাখ্যা:} \emph{মনে করি, স্থানটির দূরত্ব
\textquotesingle ক\textquotesingle{} কি.মি.}

\emph{৪ কি.মি./ঘণ্টা বেগে যেতে সময় লাগে = (ক/৪) ঘণ্টা}

\emph{৫ কি.মি./ঘণ্টা বেগে যেতে সময় লাগে = (ক/৫) ঘণ্টা}

\emph{প্রশ্নমতে,}

\emph{(ক/৪)−(ক/৫)=(১/২)}

\emph{⇒ (৫ক − ৪ক/২০)=(১/২)}

\emph{⇒ (ক/২০)=(১/২)}

\emph{⇒ ২ক = ২০}

\emph{∴ ক = ১০ কি.মি.}

\emph{সংক্ষিপ্ত নিয়ম:}

\emph{দূরত্ব = (১ম বেগ×২য় বেগ/২য় বেগ−১ম বেগ)×সময়ের পার্থক্য}

\emph{= (৫×৪/৫−৪)×(১/২)}

\emph{= (২০/১)×(১/২) = ১০ কি.মি.}
\end{quote}

\textbf{৪ কি. মি. / ঘণ্টা বেগে চললে কোন স্থানে পৌছাতে যে সময় লাগে, ৫ কি.
মি. / ঘণ্টা বেগে চললে তার চেয়ে ১ ঘণ্টা কম সময় লাগে। স্থানটির দূরত্ব কত ?}

ক) ১০ কি.মি.

\textbf{\hl{খ) ২০ কি.মি. ✔}}

গ) ১৬ কি.মি.

ঘ) ৩০ কি.মি.

\textbf{\hl{✔ সঠিক উত্তর: খ) ২০ কি.মি.}}

\begin{quote}
\textbf{ব্যাখ্যা:} \emph{মনে করি, স্থানটির দূরত্ব
\textquotesingle ক\textquotesingle{} কি.মি.}

\emph{৪ কি.মি./ঘণ্টা বেগে যেতে সময় লাগে = (ক/৪) ঘণ্টা}

\emph{৫ কি.মি./ঘণ্টা বেগে যেতে সময় লাগে = (ক/৫) ঘণ্টা}

\emph{প্রশ্নমতে,}

\emph{(ক/৪)−(ক/৫)=১}

\emph{⇒ (৫ক − ৪ক/২০)=১}

\emph{⇒ (ক/২০)=১}

\emph{∴ ক = ২০ কি.মি.}

\emph{সংক্ষিপ্ত নিয়ম:}

\emph{দূরত্ব = (১ম বেগ×২য় বেগ/২য় বেগ−১ম বেগ)×সময়ের পার্থক্য}

\emph{= (৫×৪/৫−৪)×১}

\emph{= (২০/১)×১ = ২০ কি.মি.}
\end{quote}

\textbf{রহিম ঘণ্টায় ৫ কি.মি. বেগে হেঁটে ক স্থান হতে খ স্থানে গেল। সুবাস ঘণ্টায়
৬ কি.মি. বেগে হেঁটে খ হতে ক স্থানে গেল। সুবাসের পৌছাতে (১/২) ঘণ্টা সময় কম
লাগল। ক স্থান হতে খ স্থানের দূরত্ব কত?}

\textbf{\hl{ক) ১৫ কি.মি. ✔}}

খ) ১৭ কি.মি.

গ) ১২ কি.মি.

ঘ) ১৪ কি.মি

\textbf{\hl{✔ সঠিক উত্তর: ক) ১৫ কি.মি.}}

\begin{quote}
\textbf{ব্যাখ্যা:} \emph{সংক্ষিপ্ত নিয়ম:}

\emph{দূরত্ব = (২য় গতিবেগ×১ম গতিবেগ/২য় গতিবেগ−১ম গতিবেগ)×সময়ের পার্থক্য}

\emph{= (৬×৫/৬−৫)×(১/২)}

\emph{= (৩০/১)×(১/২) = ১৫ কি.মি.}

\emph{∴ স্থানটির দূরত্ব = ১৫ কিমি।}

\emph{বিকল্প পদ্ধতি:}

\emph{ধরি, \textquotesingle ক\textquotesingle{} স্থান হতে
\textquotesingle খ\textquotesingle{} স্থানের দূরত্ব = x কিমি}

\emph{⇒ রহিমের \textquotesingle ক\textquotesingle{} স্থান হতে
\textquotesingle খ\textquotesingle{} স্থানে যেতে সময় লাগে = (x/৫) ঘণ্টা}

\emph{⇒ সুবাসের \textquotesingle খ\textquotesingle{} স্থান হতে
\textquotesingle ক\textquotesingle{} স্থানে যেতে সময় লাগে = (x/৬) ঘণ্টা}

\emph{প্রশ্নমতে,}

\emph{(x/৫)−(x/৬)=(১/২)}

\emph{⇒ (৬x−৫x/৩০)=(১/২)}

\emph{⇒ (x/৩০)=(১/২)}

\emph{⇒ ২x = ৩০}

\emph{⇒ x = ১৫}

\emph{∴ স্থানটির দূরত্ব = ১৫ কিমি।}
\end{quote}

\textbf{ঘণ্টায় ৪ কি. মি. বেগে চললে কোন গন্তব্যে পৌছাতে যে সময় লাগে, ঘণ্টায় ৫
কি. মি. বেগে চললে তার চেয়ে ৩০ মিনিট কম সময় লাগে। গন্তব্য স্থানটির দূরত্ব কত?}

ক) ৮ কি.মি.

খ) ১২ কি.মি.

গ) ১৬ কি.মি.

\textbf{\hl{ঘ) ১০ কি.মি. ✔}}

\textbf{\hl{✔ সঠিক উত্তর: ঘ) ১০ কি.মি.}}

\begin{quote}
\textbf{ব্যাখ্যা:} \emph{সংক্ষিপ্ত নিয়ম:}

\emph{গন্তব্য স্থানটির দূরত্ব = (২য় গতিবেগ×১ম গতিবেগ/২য় গতিবেগ−১ম
গতিবেগ)×সময়ের পার্থক্য}

\emph{= (৫×৪/৫−৪)×(১/২) {[}কারণ ৩০ মিনিট = (১/২) ঘণ্টা{]}}

\emph{= (২০/১)×(১/২) = ১০ কি.মি.}

\emph{∴ গন্তব্য স্থানটির দূরত্ব ১০ কি.মি.}

\emph{বিকল্প পদ্ধতি:}

\emph{মনে করি, স্থানটির দূরত্ব \textquotesingle ক\textquotesingle{} কিমি}

\emph{প্রশ্নমতে, (ক/৪)−(ক/৫)=(১/২)}

\emph{⇒ (ক/২০)=(১/২)}

\emph{⇒ ২ক = ২০ কিমি}

\emph{⇒ \textquotesingle ক\textquotesingle{} = ১০ কিমি}

\emph{∴ গন্তব্য স্থানটির দূরত্ব ১০ কি.মি.}
\end{quote}

\textbf{এক ব্যক্তি সকালে ৬ কি. মি. / ঘণ্টা বেগে হেঁটে বাসা থেকে অফিসে যান
এবং বিকালে ৪ কি. মি. / ঘণ্টা বেগে হেঁটে অফিস থেকে বাসায় ফেরেন, এতে তার ১
ঘণ্টা বেশি লাগে। বাসা থেকে অফিসের দূরত্ব কত?}

ক) ১০ কি.মি.

\textbf{\hl{খ) ১২ কি.মি. ✔}}

গ) ১৬ কি.মি.

ঘ) ৮ কি.মি.

\textbf{\hl{✔ সঠিক উত্তর: খ) ১২ কি.মি.}}

\begin{quote}
\textbf{ব্যাখ্যা:} \emph{সংক্ষিপ্ত নিয়ম:}

\emph{বাসা থেকে অফিসের দূরত্ব = (১ম গতিবেগ×২য় গতিবেগ/১ম গতিবেগ−২য়
গতিবেগ)×সময়ের পার্থক্য}

\emph{= (৬×৪/৬−৪)×১}

\emph{= (২৪/২)×১ = ১২ কি.মি.}

\emph{∴ বাসা থেকে অফিসের দূরত্ব ১২ কি.মি.}

\emph{বিকল্প পদ্ধতি:}

\emph{মনে করি, বাসা থেকে অফিসের দূরত্ব
\textquotesingle ক\textquotesingle{} কিমি}

\emph{প্রশ্নমতে, (ক/৪)−(ক/৬)=১}

\emph{⇒ (৩ক−২ক/১২)=১}

\emph{⇒ (ক/১২)=১}

\emph{⇒ \textquotesingle ক\textquotesingle{} = ১২ কিমি}

\emph{∴ বাসা থেকে অফিসের দূরত্ব ১২ কি.মি.}
\end{quote}

\textbf{এক ব্যক্তি ঘণ্টায় ৪ কিলোমিটার বেগে চলে কোন স্থানে গেল এবং ঘণ্টায় ৩
কিলোমিটার বেগে চলে ফিরে আসল। যাতায়াতে তার গড় গতিবেগ কত?}

ক) ৩ কি.মি.

\textbf{\hl{খ) ৩(৩/৭) কি.মি. ✔}}

গ) ৪ কি.মি.

ঘ) ৪(৩/৭) কি.মি.

\textbf{\hl{✔ সঠিক উত্তর: খ) ৩(৩/৭) কি.মি.}}

\begin{quote}
\textbf{ব্যাখ্যা:} \emph{সংক্ষিপ্ত নিয়ম:}

\emph{যাতায়াতে গড় গতিবেগ = (২×১ম গতিবেগ×২য় গতিবেগ/১ম গতিবেগ+২য় গতিবেগ)}

\emph{= (২×৪×৩/৪+৩)}

\emph{= (২৪/৭) = ৩(৩/৭) কি.মি.}

\emph{∴ যাতায়াতে তার গড় গতিবেগ = ৩(৩/৭) কি.মি.}

\emph{বিকল্প পদ্ধতি:}

\emph{ধরি, স্থানটির দূরত্ব = d কি.মি.}

\emph{⇒ যাওয়ার সময় লাগে = (d/৪) ঘণ্টা}

\emph{⇒ আসার সময় লাগে = (d/৩) ঘণ্টা}

\emph{⇒ মোট সময় = (d/৪)+(d/৩)=(৭d/১২) ঘণ্টা}

\emph{⇒ মোট দূরত্ব = d + d = 2d কি.মি.}

\emph{∴ গড় গতিবেগ = (মোট দূরত্ব/মোট
সময়)=(২d/(৭d/১২))=(২×১২/৭)=(২৪/৭)=৩(৩/৭) কি.মি.}
\end{quote}

\textbf{একটি নৌকা মালামালসহ যে গতিতে চলে, খালি থাকলে তার দ্বিগুণ গতিতে
চলে। নৌকাটি মালামালসহ ৪০ কিমি অতিক্রম করে, দুইঘণ্টা ব্যয় করে নৌকা থেকে
মালামাল নামায় এবং খালি অবস্থায় পূর্বের স্থানে ফিরে আসে। সব মিলিয়ে যদি ১৭
ঘণ্টা সময় লেগে থাকে, তাহলে খালি অবস্থায় নৌকাটির বেগ কত কিমি/ঘণ্টা ছিল?}

ক) ২

খ) ৪

\textbf{\hl{গ) ৮ ✔}}

ঘ) ৮

\textbf{\hl{✔ সঠিক উত্তর: গ) ৮}}

\begin{quote}
\textbf{ব্যাখ্যা:} \emph{ধরি, মালামালসহ নৌকার বেগ = x কিমি/ঘণ্টা}

\emph{∴ খালি অবস্থায় নৌকার বেগ = 2x কিমি/ঘণ্টা}

\emph{মালামাল নামাতে সময় লাগে = ২ ঘণ্টা}

\emph{∴ মোট যাতায়াতের সময় = (১৭ − ২) ঘণ্টা = ১৫ ঘণ্টা}

\emph{প্রশ্নমতে,}

\emph{(৪০/x)+(৪০/২x)=১৫}

\emph{⇒ (৪০+২০/x)=১৫}

\emph{⇒ (৬০/x)=১৫}

\emph{⇒ ১৫x = ৬০}

\emph{∴ x = ৪}

\emph{∴ খালি অবস্থায় নৌকার বেগ = (৪ × ২) বা ৮ কিমি/ঘণ্টা।}
\end{quote}

\textbf{স্রোতের প্রতিকূলে যেতে যে সময় লাগে অনুকূলে যেতে তার অর্ধেক সময় লাগে।
যাতায়াতে যদি সময় লাগে ১২ ঘণ্টা, তাহলে স্রোতের অনুকূলে যেতে সময় লাগে-}

ক) ৬ ঘণ্টা

\textbf{\hl{খ) ৪ ঘণ্টা ✔}}

গ) ১০ ঘণ্টা

ঘ) ৮ ঘণ্টা

\textbf{\hl{✔ সঠিক উত্তর: খ) ৪ ঘণ্টা}}

\begin{quote}
\textbf{ব্যাখ্যা:} \emph{ধরি, স্রোতের অনুকূলে যেতে সময় লাগে = x ঘণ্টা}

\emph{∴ স্রোতের প্রতিকূলে যেতে সময় লাগে = 2x ঘণ্টা}

\emph{প্রশ্নমতে,}

\emph{x + 2x = ১২}

\emph{⇒ 3x = ১২}

\emph{∴ x = ৪}

\emph{∴ স্রোতের অনুকূলে যেতে সময় লাগে ৪ ঘণ্টা।}
\end{quote}

\textbf{লঞ্চ ও স্রোতের গতিবেগ যথাক্রমে ঘণ্টায় ১৮ কিঃ মিঃ ও ৬ কিঃ মিঃ। নদীপথে
৪৮ কিঃ মিঃ অতিক্রম করে পুনরায় ফিরে আসতে সময় লাগবে -}

ক) ১০ ঘণ্টা

খ) ৫ ঘণ্টা

\textbf{\hl{গ) ৬ ঘণ্টা ✔}}

ঘ) ৮ ঘণ্টা

\textbf{\hl{✔ সঠিক উত্তর: গ) ৬ ঘণ্টা}}

\begin{quote}
\textbf{ব্যাখ্যা:} \emph{স্রোতের অনুকূলে লঞ্চের কার্যকরী গতিবেগ = (১৮ + ৬) কিঃ
মিঃ/ঘণ্টা = ২৪ কিঃ মিঃ/ঘণ্টা}

\emph{∴ অনুকূলে ৪৮ কিঃ মিঃ যেতে সময় লাগে = (৪৮/২৪) = ২ ঘণ্টা}

\emph{স্রোতের প্রতিকূলে লঞ্চের কার্যকরী গতিবেগ = (১৮ − ৬) কিঃ মিঃ/ঘণ্টা = ১২
কিঃ মিঃ/ঘণ্টা}

\emph{∴ প্রতিকূলে ৪৮ কিঃ মিঃ ফিরে আসতে সময় লাগে = (৪৮/১২) = ৪ ঘণ্টা}

\emph{∴ মোট প্রয়োজনীয় সময় = (২ + ৪) ঘণ্টা = ৬ ঘণ্টা।}
\end{quote}

\textbf{স্থির পানিতে একটি নৌকার গতিবেগ ঘণ্টায় ৭ কিমি। এরূপ নৌকাটির স্রোতের
অনুকূলে ৩৩ কিমি পথ যেতে ৩ ঘণ্টা সময় লেগেছে। ফিরে আসার সময় নৌকাটির কত ঘণ্টা
সময় লাগবে?}

ক) ১৩ ঘণ্টা

খ) ১২ ঘণ্টা

\textbf{\hl{গ) ১১ ঘণ্টা ✔}}

ঘ) ১০ ঘণ্টা

\textbf{\hl{✔ সঠিক উত্তর: গ) ১১ ঘণ্টা}}

\begin{quote}
\textbf{ব্যাখ্যা:} \emph{স্রোতের অনুকূলে নৌকার গতিবেগ = (৩৩/৩) = ১১
কিমি/ঘণ্টা}

\emph{∴ স্রোতের বেগ = (১১ − ৭) কিমি/ঘণ্টা = ৪ কিমি/ঘণ্টা}

\emph{স্রোতের প্রতিকূলে নৌকার গতিবেগ = (৭ − ৪) কিমি/ঘণ্টা = ৩ কিমি/ঘণ্টা}

\emph{∴ ফিরে আসার সময় (প্রতিকূলে ৩৩ কিমি যেতে) সময় লাগবে = (৩৩/৩) = ১১
ঘণ্টা।}
\end{quote}

\textbf{স্রোতের বিপরীত নৌকার বেগ ৮ কিমি/ঘণ্টা এবং স্থির পানিতে নৌকার বেগ ১৩
কিমি/ঘণ্টা হলে, স্রোতের অনুকূলে ১৪৪ কিমিঃ অতিক্রম করে পুনরায় পূর্বের স্থানে ফিরে
আসতে মোট কত ঘণ্টা সময় লাগবে?}

\textbf{\hl{ক) ২৬ ✔}}

খ) ২৮

গ) ৩০

ঘ) ৩২

\textbf{\hl{✔ সঠিক উত্তর: ক) ২৬}}

\begin{quote}
\textbf{ব্যাখ্যা:} \emph{স্থির পানিতে নৌকার বেগ = ১৩ কিমি/ঘণ্টা}

\emph{স্রোতের প্রতিকূলে নৌকার বেগ = ৮ কিমি/ঘণ্টা}

\emph{∴ স্রোতের বেগ = (১৩ − ৮) = ৫ কিমি/ঘণ্টা}

\emph{স্রোতের অনুকূলে নৌকার বেগ = (১৩ + ৫) = ১৮ কিমি/ঘণ্টা}

\emph{১৪৪ কিমি গিয়ে পুনরায় ফিরে আসতে মোট সময় লাগবে = ((১৪৪/১৮)+(১৪৪/৮))
ঘণ্টা}

\emph{= (৮ + ১৮) ঘণ্টা = ২৬ ঘণ্টা।}
\end{quote}

\textbf{ঢাকা ও সিলেটের দূরত্ব ৩০০ কি.মি। ঢাকা থেকে ট্রেন সকাল ৬ টায় ছেড়ে ২
টায় সিলেটে পৌঁছে। ট্রেনের গতিবেগ কত?}

\textbf{\hl{ক) ৩৭.৫ কি.মি/ঘণ্টা ✔}}

খ) ৩৫.৭ কি.মি/ঘণ্টা

গ) ৩৫ কি.মি/ঘণ্টা

ঘ) ৩৭ কি.মি/ঘণ্টা

\textbf{\hl{✔ সঠিক উত্তর: ক) ৩৭.৫ কি.মি/ঘণ্টা}}

\begin{quote}
\textbf{ব্যাখ্যা:} \emph{সকাল ৬টা থেকে দুপুর ২টা পর্যন্ত মোট সময় = ৮ ঘণ্টা}

\emph{দূরত্ব = ৩০০ কি.মি}

\emph{∴ গতিবেগ = (৩০০/৮) = ৩৭.৫ কি.মি/ঘণ্টা।}
\end{quote}

\textbf{রেল লাইনের পাশে একটি তালগাছ রয়েছে। ঘণ্টায় ৪৫ কিঃমিঃ বেগে ধাবমান
১৫০ মিঃ লম্বা একটি ট্রেন কত সময়ে ঐ তালগাছটি অতিক্রম করবে?}

ক) ১১ সেকেন্ড

\textbf{\hl{খ) ১২ সেকেন্ড ✔}}

গ) ১৩ সেকেন্ড

ঘ) ১৪ সেকেন্ড

\textbf{\hl{✔ সঠিক উত্তর: খ) ১২ সেকেন্ড}}

\begin{quote}
\textbf{ব্যাখ্যা:} \emph{আমরা জানি, ১ কি.মি. = ১০০০ মিটার এবং ১ ঘণ্টা =
৩৬০০ সেকেন্ড।}

\emph{⇒ ৪৫ কি.মি./ঘণ্টা = (৪৫×১০০০/৩৬০০) মিটার/সেকেন্ড = ১২.৫
মিটার/সেকেন্ড।}

\emph{তালগাছটি অতিক্রম করতে ট্রেনটিকে তার নিজের দৈর্ঘ্যের সমান (১৫০ মিটার)
দূরত্ব অতিক্রম করতে হবে।}

\emph{⇒ সময় = (মোট অতিক্রান্ত দূরত্ব/গতিবেগ)}

\emph{= (১৫০/১২.৫) = ১২ সেকেন্ড।}

\emph{সংক্ষিপ্ত নিয়ম:}

\emph{সময় = (দৈর্ঘ্য (মিটার)×৩৬০০/গতিবেগ (কি.মি.)×১০০০)}

\emph{= (১৫০×৩৬০০/৪৫×১০০০)}

\emph{= (৫৪০০০/৪৫০০) = ১২ সেকেন্ড।}
\end{quote}

\textbf{১১০ মিটার ও ১৪০ মিটার দীর্ঘ দুটি ট্রেন যথাক্রমে ঘণ্টায় ৪৫ কি.মি ও ৫৫
কি.মি. বেগে বিপরীত দিক থেকে কত সময় পর সমান্তরালভাবে অতিক্রম করবে?}

\textbf{\hl{ক) ৯ সেকেন্ড ✔}}

খ) ১৬ সেকেন্ড

গ) ২০ সেকেন্ড

ঘ) ৩০ সেকেন্ড

\textbf{\hl{✔ সঠিক উত্তর: ক) ৯ সেকেন্ড}}

\begin{quote}
\textbf{ব্যাখ্যা:} \emph{ট্রেন দুটির মোট দৈর্ঘ্য = (১১০ + ১৪০) = ২৫০ মিটার।}

\emph{যেহেতু ট্রেন দুটি বিপরীত দিক থেকে আসছে, তাই আপেক্ষিক গতিবেগ = (৪৫ +
৫৫) = ১০০ কি.মি./ঘণ্টা।}

\emph{১০০ কি.মি. = (১০০ × ১০০০) = ১০০০০০ মিটার এবং ১ ঘণ্টা = ৩৬০০
সেকেন্ড।}

\emph{এখন, ১০০০০০ মিটার যায় = ৩৬০০ সেকেন্ডে}

\emph{⇒ ১ মিটার যায় = (৩৬০০/১০০০০০) সেকেন্ডে}

\emph{⇒ ২৫০ মিটার যায় = (৩৬০০×২৫০/১০০০০০) = ৯ সেকেন্ডে।}

\emph{সংক্ষিপ্ত নিয়ম:}

\emph{প্রয়োজনীয় সময় = (মোট দৈর্ঘ্য/আপেক্ষিক গতিবেগ)}

\emph{= (২৫০/১০০×(৫/১৮))}

\emph{= (২৫০×১৮/৫০০) = ৯ সেকেন্ড।}
\end{quote}

\textbf{ঘন্টায় ৬০ কি.মি. বেগে ১০০ মিটার দীর্ঘ ট্রেন ৩০০ মিটার দীর্ঘ একটি
প্লাটফরম অতিক্রম করতে কত সময় লাগবে?}

\textbf{\hl{ক) ২৪ সেকেন্ড ✔}}

খ) ২৪ মিনিট

গ) ২০ সেকেন্ড

ঘ) ২০ মিনিট

\textbf{\hl{✔ সঠিক উত্তর: ক) ২৪ সেকেন্ড}}

\begin{quote}
\textbf{ব্যাখ্যা:} \emph{প্ল্যাটফর্ম অতিক্রম করতে মোট দূরত্ব = (১০০ + ৩০০) = ৪০০
মিটার।}

\emph{গতিবেগ = ৬০ কি.মি./ঘণ্টা = (৬০×১০০০/৩৬০০) = (৫০/৩) মিটার/সেকেন্ড।}

\emph{⇒ সময় = (মোট অতিক্রান্ত দূরত্ব/গতিবেগ)}

\emph{= (৪০০/(৫০/৩))}

\emph{= (৪০০×৩/৫০) = ২৪ সেকেন্ড।}

\emph{সংক্ষিপ্ত নিয়ম:}

\emph{সময় = ((ট্রেন দৈর্ঘ্য+প্ল্যাটফর্ম দৈর্ঘ্য)×৩৬০০/গতিবেগ×১০০০)}

\emph{= ((১০০+৩০০)×৩৬০০/৬০×১০০০)}

\emph{= (৪০০×৩৬০০/৬০০০০) = ২৪ সেকেন্ড।}
\end{quote}

\textbf{একটি ট্রেন ১৮০ কি.মি./ঘন্টা বেগে চলে ২৫ সেকেন্ডে ৮০০ মিটার দীর্ঘ
একটি সেতু অতিক্রম করে। ট্রেনটির দৈঘ্য কত মিটার?}

\textbf{\hl{ক) ৪৫০ ✔}}

খ) ৯০০

গ) ১২৫০

ঘ) ১৮০০

\textbf{\hl{✔ সঠিক উত্তর: ক) ৪৫০}}

\begin{quote}
\textbf{ব্যাখ্যা:} \emph{১৮০ কি.মি. = (১৮০ × ১০০০) = ১৮০০০০ মিটার এবং ১
ঘণ্টা = ৩৬০০ সেকেন্ড।}

\emph{৩৬০০ সেকেন্ডে ট্রেনটি অতিক্রম করে = ১৮০০০০ মিটার}

\emph{⇒ ১ সেকেন্ডে অতিক্রম করে = (১৮০০০০/৩৬০০) = ৫০ মিটার}

\emph{⇒ ২৫ সেকেন্ডে অতিক্রম করে = (৫০ × ২৫) = ১২৫০ মিটার।}

\emph{যেহেতু ২৫ সেকেন্ডে অতিক্রম করা দূরত্ব = (ট্রেনের দৈর্ঘ্য + সেতুর দৈর্ঘ্য)}

\emph{∴ ট্রেনটির দৈর্ঘ্য = (১২৫০ − ৮০০) = ৪৫০ মিটার।}

\emph{সংক্ষিপ্ত নিয়ম:}

\emph{গতিবেগ = ১৮০ × (৫/১৮) = ৫০ মিটার/সেকেন্ড।}

\emph{২৫ সেকেন্ডে অতিক্রান্ত মোট দূরত্ব = (৫০ × ২৫) = ১২৫০ মিটার।}

\emph{⇒ ট্রেনের দৈর্ঘ্য = মোট দূরত্ব − সেতুর দৈর্ঘ্য}

\emph{= ১২৫০ − ৮০০ = ৪৫০ মিটার।}
\end{quote}

\textbf{১২০ মিটার ও ৮০ মিটার দীর্ঘ দুটি ট্রেন প্রতি ঘণ্টায় যথাক্রমে ১৮ কি. মি.
ও ১২ কি. মি. বেগে চলছে। ট্রেন দুটি একই স্থান থেকে একই দিকে একই সময়ে অগ্রসর
হলে পরস্পরকে অতিক্রম করতে কত সময় লাগবে?}

ক) ১ মিনিট

\textbf{\hl{খ) ২ মিনিট ✔}}

গ) ৩ মিনিট

ঘ) ৪ মিনিট

\textbf{\hl{✔ সঠিক উত্তর: খ) ২ মিনিট}}

\begin{quote}
\textbf{ব্যাখ্যা:} \emph{ট্রেন দুটির মোট দৈর্ঘ্য = (১২০ + ৮০) মিটার = ২০০
মিটার}

\emph{যেহেতু ট্রেন দুটি একই দিকে চলছে, আপেক্ষিক বেগ = (১৮ − ১২) কি.মি./ঘণ্টা
= ৬ কি.মি./ঘণ্টা = ৬,০০০ মিটার/ঘণ্টা}

\emph{৬,০০০ মিটার অতিক্রম করতে সময় লাগে = ৬০ মিনিট}

\emph{∴ ২০০ মিটার অতিক্রম করতে সময় লাগবে = (৬০×২০০/৬০০০) = ২ মিনিট।}
\end{quote}

\textbf{একটি ট্রেন ঘণ্টায় ৬০ কিমি বেগে চলে। ১০০ মিটার যেতে ট্রেনটির কত সময়
লাগবে?}

ক) ৬০ সেকেন্ড

খ) ১ সেকেন্ড

\textbf{\hl{গ) ৬ সেকেন্ড ✔}}

ঘ) ০.৬ সেকেন্ড

\textbf{\hl{✔ সঠিক উত্তর: গ) ৬ সেকেন্ড}}

\begin{quote}
\textbf{ব্যাখ্যা:} \emph{৬০ কিমি = ৬০,০০০ মিটার}

\emph{১ ঘণ্টা = ৩,৬০০ সেকেন্ড}

\emph{৬০,০০০ মিটার যেতে সময় লাগে = ৩,৬০০ সেকেন্ড}

\emph{∴ ১০০ মিটার যেতে সময় লাগবে = (৩৬০০×১০০/৬০০০০) = ৬ সেকেন্ড।}
\end{quote}

\textbf{ঘণ্টায় ৬০ কিমি বেগে ১০০ মিটার দীর্ঘ একটি ট্রেন ৩০০ মি. একটি দীর্ঘ
প্ল্যাটফর্ম অতিক্রম করতে কত সময় লাগবে?}

\textbf{\hl{ক) ২৪ সেকেন্ড ✔}}

খ) ২০ সেকেন্ড

গ) ২৪ মিনিট

ঘ) ২০ মিনিট

\textbf{\hl{✔ সঠিক উত্তর: ক) ২৪ সেকেন্ড}}

\begin{quote}
\textbf{ব্যাখ্যা:} \emph{প্ল্যাটফর্ম ও ট্রেনের মোট দৈর্ঘ্য = (৩০০ + ১০০) মিটার =
৪০০ মিটার}

\emph{৬০ কিমি = ৬০,০০০ মিটার এবং ১ ঘণ্টা = ৩,৬০০ সেকেন্ড}

\emph{৬০,০০০ মিটার অতিক্রম করতে সময় লাগে = ৩,৬০০ সেকেন্ড}

\emph{∴ ৪০০ মিটার অতিক্রম করতে সময় লাগবে = (৩৬০০×৪০০/৬০০০০) = ২৪ সেকেন্ড।}
\end{quote}

\textbf{১২০ মিটার লম্বা একটি ট্রেনের গতিবেগ ঘণ্টায় ৩০ কিলোমিটার। ৭৮০ মিটার
দীর্ঘ একটি সেতু পার হতে এটির কত সময় লাগবে?}

\textbf{\hl{ক) ১ মিনিট ৪৮ সেকেন্ড ✔}}

খ) ১ মিনিট ৪০ সেকেন্ড

গ) ২ মিনিট

ঘ) ১ মিনিট ৩৬ সেকেন্ড

\textbf{\hl{✔ সঠিক উত্তর: ক) ১ মিনিট ৪৮ সেকেন্ড}}

\begin{quote}
\textbf{ব্যাখ্যা:} \emph{ট্রেনটিকে অতিক্রম করতে হবে = (৭৮০ + ১২০) মিটার =
৯০০ মিটার}

\emph{৩০ কিলোমিটার = ৩০,০০০ মিটার এবং ১ ঘণ্টা = ৩,৬০০ সেকেন্ড}

\emph{৩০,০০০ মিটার অতিক্রম করে = ৩,৬০০ সেকেন্ডে}

\emph{∴ ৯০০ মিটার অতিক্রম করে = (৩৬০০×৯০০/৩০০০০) = ১০৮ সেকেন্ড = ১ মিনিট
৪৮ সেকেন্ড।}
\end{quote}

\textbf{১৮০ মিটার দীর্ঘ একটি ট্রেন ৫৪ কিলোমিটার/ঘণ্টা বেগে ৭২০ মিটার দীর্ঘ
একটি টানেলে প্রবেশ করলে, টানেলটি অতিক্রম করতে ট্রেনটির কত সময় লাগবে?}

ক) ১৭ সেকেন্ড

খ) ৪৮ সেকেন্ড

\textbf{\hl{গ) ৬০ সেকেন্ড ✔}}

ঘ) ৮০ সেকেন্ড

ঙ) কোনটিই নয়

\textbf{\hl{✔ সঠিক উত্তর: গ) ৬০ সেকেন্ড}}

\begin{quote}
\textbf{ব্যাখ্যা:} \emph{মোট অতিক্রম্য দূরত্ব = (১৮০ + ৭২০) মিটার = ৯০০ মিটার}

\emph{ট্রেনের বেগ = ৫৪ কিমি/ঘণ্টা = ৫৪×(৫/১৮) মি./সে. = ১৫ মি./সে.}

\emph{∴ টানেলটি অতিক্রম করতে প্রয়োজনীয় সময় = (৯০০/১৫) = ৬০ সেকেন্ড।}
\end{quote}

\textbf{একটি ট্রেন 50mph বেগে 10 মাইল ভ্রমণ করল। যাওয়া-আসা বাবদ মোট 20
মিনিট সময় ব্যয় হলে ফেরার সময় গতিবেগ কত ছিল?}

ক) 55 miles/hour

খ) 60 miles/hour

গ) 65 miles/hour

\textbf{\hl{ঘ) 75 miles/hour ✔}}

\textbf{\hl{✔ সঠিক উত্তর: ঘ) 75 miles/hour}}

\begin{quote}
\textbf{ব্যাখ্যা:} \emph{যাওয়ার সময় লাগা সময় = (১০/৫০) ঘণ্টা = (১/৫) ঘণ্টা =
( (১/৫) × ৬০) মিনিট = ১২ মিনিট}

\emph{মোট সময় = ২০ মিনিট}

\emph{∴ ফেরার সময় লেগেছে = (২০ − ১২) মিনিট = ৮ মিনিট = (৮/৬০) ঘণ্টা =
(২/১৫) ঘণ্টা}

\emph{∴ ফেরার সময় গতিবেগ = ১০÷(২/১৫) = ১০×(১৫/২) = ৭৫ miles/hour।}
\end{quote}

\textbf{১০০ মিটার লম্বা একটি ট্রেন একটি মাইল ফলক ১০ সেকেন্ডে অতিক্রম করে।
বিপরীত দিক থেকে আসা সমান লম্বা আরেকটি ট্রেন ঐ মাইলফলকটি ৮ সেকেন্ডে অতিক্রম
করে। দ্বিতীয় ট্রেনটির গতিবেগ বের করুন।}

ক) ৩৬ কিলোমিটার/ঘণ্টা

খ) ৪০ কিলোমিটার/ঘণ্টা

\textbf{\hl{গ) ৪৫ কিলোমিটার/ঘণ্টা ✔}}

ঘ) ৫০ কিলোমিটার/ঘণ্টা

\textbf{\hl{✔ সঠিক উত্তর: গ) ৪৫ কিলোমিটার/ঘণ্টা}}

\begin{quote}
\textbf{ব্যাখ্যা:} \emph{দ্বিতীয় ট্রেনটি ৮ সেকেন্ডে অতিক্রম করে নিজের দৈর্ঘ্য =
১০০ মিটার}

\emph{∴ ১ সেকেন্ডে অতিক্রম করে = (১০০/৮) মিটার}

\emph{∴ ১ ঘণ্টা বা ৩,৬০০ সেকেন্ডে অতিক্রম করে = (১০০×৩৬০০/৮) মিটার = ৪৫,০০০
মিটার = ৪৫ কিলোমিটার}

\emph{∴ দ্বিতীয় ট্রেনটির গতিবেগ ৪৫ কিলোমিটার/ঘণ্টা।}
\end{quote}

\textbf{১৫০ মিটার লম্বা ট্রেন ৪৫০ মিটার লম্বা একটি প্ল্যাটফরমকে ২০ সেকেন্ডে
অতিক্রম করলে ঐ ট্রেনের গতিবেগ প্রতি সেকেন্ডে কত হবে?}

ক) ৪০ মিটার

\textbf{\hl{খ) ৩০ মিটার ✔}}

গ) ২৫ মিটার

ঘ) ২০ মিটার

\textbf{\hl{✔ সঠিক উত্তর: খ) ৩০ মিটার}}

\begin{quote}
\textbf{ব্যাখ্যা:} \emph{ট্রেনটি ২০ সেকেন্ডে মোট অতিক্রম করে = (১৫০ + ৪৫০)
মিটার = ৬০০ মিটার}

\emph{∴ ট্রেনের গতিবেগ প্রতি সেকেন্ডে = (৬০০/২০) = ৩০ মিটার/সেকেন্ড।}
\end{quote}

\textbf{১২০ মিটার লম্বা একটি আন্তঃনগর এক্সপ্রেস একটি ল্যাম্পপোস্ট ৬ সেকেন্ডে
অতিক্রম করল। ট্রেনটির গতিবেগ কিলোমিটার/ ঘণ্টায়-}

ক) ৯৬

\textbf{\hl{খ) ৭২ ✔}}

গ) ৪৮

ঘ) ৩৬

\textbf{\hl{✔ সঠিক উত্তর: খ) ৭২}}

\begin{quote}
\textbf{ব্যাখ্যা:} \emph{৬ সেকেন্ডে ট্রেনটি অতিক্রম করে নিজের দৈর্ঘ্য = ১২০
মিটার}

\emph{∴ ১ সেকেন্ডে অতিক্রম করে = (১২০/৬) = ২০ মিটার}

\emph{∴ ১ ঘণ্টা বা ৩,৬০০ সেকেন্ডে অতিক্রম করে = (২০ × ৩৬০০) মিটার = ৭২,০০০
মিটার = ৭২ কিলোমিটার/ঘণ্টা।}
\end{quote}

\textbf{একটি ট্রেন ২৫৪ মিটার দীর্ঘ একটি সেতুকে ২০ সেকেন্ডে এবং ১০০ মিটার
দীর্ঘ অপর একটি সেতুকে ১৩ সেকেন্ডে অতিক্রম করল। ট্রেনটির দৈর্ঘ্য-}

ক) ২২ মিটার

খ) ১০৪ মিটার

গ) ১৫০ মিটার

\textbf{\hl{ঘ) ১৮৬ মিটার ✔}}

\textbf{\hl{✔ সঠিক উত্তর: ঘ) ১৮৬ মিটার}}

\begin{quote}
\textbf{ব্যাখ্যা:} \emph{ধরি, ট্রেনটির দৈর্ঘ্য = x মিটার}

\emph{২০ সেকেন্ডে অতিক্রম করে = (২৫৪ + x) মিটার}

\emph{১৩ সেকেন্ডে অতিক্রম করে = (১০০ + x) মিটার}

\emph{শর্তমতে,}

\emph{(২৫৪+x/২০)=(১০০+x/১৩)}

\emph{⇒ ৩৩০২ + ১৩x = ২০০০ + ২০x}

\emph{⇒ ২০x − ১৩x = ৩৩০২ − ২০০০}

\emph{⇒ ৭x = ১৩০২}

\emph{∴ x = ১৮৬ মিটার।}
\end{quote}

\textbf{একটি ট্রেন ঘণ্টায় ৮৪ কি. মি. বেগে চলে। ট্রেনটি ৮০০ মি: দীর্ঘ একটি
প্ল্যাটফরম ১ মিনিটে অতিক্রম করে। ট্রেনটির দৈর্ঘ্য কত?}

ক) ৯০০ মিটার

খ) ৭৫০ মিটার

\textbf{\hl{গ) ৬০০ মিটার ✔}}

ঘ) ৫০০ মিটার

\textbf{\hl{✔ সঠিক উত্তর: গ) ৬০০ মিটার}}

\begin{quote}
\textbf{ব্যাখ্যা:} \emph{ট্রেনটির গতিবেগ = ৮৪ কি.মি./ঘণ্টা = (৮৪×১০০০/৬০)
মিটার/মিনিট = ১৪০০ মিটার/মিনিট}

\emph{∴ ১ মিনিটে ট্রেনটি মোট অতিক্রম করে ১৪০০ মিটার}

\emph{প্রকৃতপক্ষে, ১ মিনিটে অতিক্রম করে = (ট্রেনের দৈর্ঘ্য + প্ল্যাটফর্মের দৈর্ঘ্য)}

\emph{⇒ ১৪০০ = ট্রেনের দৈর্ঘ্য + ৮০০}

\emph{∴ ট্রেনের দৈর্ঘ্য = (১৪০০ − ৮০০) মিটার = ৬০০ মিটার।}
\end{quote}

\textbf{একটি গাড়ি ঘণ্টায় ২৫ মাইল বেগে ঢাকা থেকে চট্টগ্রাম অভিমুখে রওনা হলো।
আরেকটি গাড়ি ঘণ্টায় ১৫ মাইল বেগে চট্টগ্রাম থেকে ঢাকা অভিমুখে রওনা হলো। ঢাকা ও
চট্টগ্রামের দূরত্ব ২০০ মাইল হলে এ গাড়ি দুটি অবিরাম গতিতে চলতে থাকলে কত সময় পর
গাড়ি দুটি মুখোমুখি হবে?}

ক) ৪ ঘণ্টা

\textbf{\hl{খ) ৫ ঘণ্টা ✔}}

গ) ৬ ঘণ্টা

ঘ) ৮ ঘণ্টা

\textbf{\hl{✔ সঠিক উত্তর: খ) ৫ ঘণ্টা}}

\begin{quote}
\textbf{ব্যাখ্যা:} \emph{যেহেতু গাড়ি দুটি বিপরীত দিক থেকে আসছে, তাদের
আপেক্ষিক গতিবেগ = (২৫ + ১৫) মাইল/ঘণ্টা = ৪০ মাইল/ঘণ্টা}

\emph{∴ মুখোমুখি হতে প্রয়োজনীয় সময় = (২০০/৪০) = ৫ ঘণ্টা।}
\end{quote}

\textbf{একটি প্ল্যাটফর্মের দৈর্ঘ্য ১০০ মিটার। ১৫০ মিটার লম্বা একটি ট্রেনকে ঐ
প্ল্যাটফর্ম অতিক্রম করতে কত মিটার দূরত্ব অতিক্রম করতে হবে?}

ক) ১০০

খ) ১৫০

গ) ২০০

\textbf{\hl{ঘ) ২৫০ ✔}}

ঙ) কোনোটিই নয়

\textbf{\hl{✔ সঠিক উত্তর: ঘ) ২৫০}}

\begin{quote}
\textbf{ব্যাখ্যা:} \emph{প্ল্যাটফর্ম অতিক্রম করতে ট্রেনটিকে নিজের দৈর্ঘ্য সহ
প্ল্যাটফর্মের দৈর্ঘ্য অতিক্রম করতে হবে।}

\emph{∴ অতিক্রান্ত দূরত্ব = (১৫০ + ১০০) মিটার = ২৫০ মিটার।}
\end{quote}

\textbf{ঘণ্টায় ৬০ কিমিঃ বেগে ১০০ মিটার দীর্ঘ একটি ট্রেন ৩০০ মিটার একটি দীর্ঘ
প্ল্যাটফর্ম অতিক্রম করতে কত সময় লাগবে?}

\textbf{\hl{ক) ২৪ সেকেন্ড ✔}}

খ) ২০ সেকেন্ড

গ) ২৪ মিনিট

ঘ) ২০ মিনিট

\textbf{\hl{✔ সঠিক উত্তর: ক) ২৪ সেকেন্ড}}

\begin{quote}
\textbf{ব্যাখ্যা:} \emph{মোট অতিক্রম করতে হবে = (৩০০ + ১০০) মিটার = ৪০০
মিটার}

\emph{৬০ কি.মি. = ৬০,০০০ মিটার এবং ১ ঘণ্টা = ৩,৬০০ সেকেন্ড}

\emph{৬০,০০০ মিটার যায় = ৩,৬০০ সেকেন্ডে}

\emph{∴ ৪০০ মিটার অতিক্রম করতে সময় লাগবে = (৩৬০০×৪০০/৬০০০০) = ২৪ সেকেন্ড।}
\end{quote}

\textbf{৭২ কিমি/ঘণ্টা বেগে চলমান একটি ট্রেন ৫৮০ মিটার দীর্ঘ একটি প্ল্যাটফর্মকে
৫৬ সেকেন্ডে অতিক্রম করে। ট্রেনটির দৈর্ঘ্য ৫৪০ মিটার হলে, ১৬০ মিটার দীর্ঘ একটি
সেতু অতিক্রম করতে ট্রেনটির কত সেকেন্ড সময় লাগবে?}

\textbf{\hl{ক) ৩৫ ✔}}

খ) ২৮

গ) ২৭

ঘ) ৮

\textbf{\hl{✔ সঠিক উত্তর: ক) ৩৫}}

\begin{quote}
\textbf{ব্যাখ্যা:} \emph{ট্রেনটির মোট অতিক্রম করতে হবে প্ল্যাটফর্মসহ = (৫৪০ +
৫৮০) মিটার = ১১২০ মিটার।}

\emph{১১২০ মিটার অতিক্রম করতে সময় লাগে ৫৬ সেকেন্ড।}

\emph{∴ ১ মিটার অতিক্রম করতে সময় লাগে (৫৬/১১২০) সেকেন্ড।}

\emph{এখন, ১৬০ মিটার দীর্ঘ সেতু অতিক্রম করতে ট্রেনটিকে নিজের দৈর্ঘ্যসহ অতিক্রম
করতে হবে = (৫৪০ + ১৬০) মিটার = ৭০০ মিটার।}

\emph{∴ ৭০০ মিটার অতিক্রম করতে সময় লাগবে = (৫৬×৭০০/১১২০) = ৩৫ সেকেন্ড।}
\end{quote}

\textbf{২২০ মিটার দীর্ঘ একটি ট্রেনের গতি ঘণ্টায় ৫৯ কিলোমিটার। ট্রেনের বিপরীত
দিক থেকে ঘণ্টায় ৭ কিলোমিটার গতিতে আসা এক ব্যক্তিকে অতিক্রম করতে উহার কত সময়
লাগবে?}

\textbf{\hl{ক) ১২ সেকেন্ড ✔}}

খ) ১৫ সেকেন্ড

গ) ১৮ সেকেন্ড

ঘ) ২০ সেকেন্ড

\textbf{\hl{✔ সঠিক উত্তর: ক) ১২ সেকেন্ড}}

\begin{quote}
\textbf{ব্যাখ্যা:} \emph{ট্রেন ও ব্যক্তির আপেক্ষিক গতিবেগ = (৫৯ + ৭) কিমি/ঘণ্টা
= ৬৬ কিমি/ঘণ্টা।}

\emph{গতিবেগকে মিটার/সেকেন্ডে রূপান্তর করলে = ৬৬×(৫/১৮) মিটার/সেকেন্ড =
(৫৫/৩) মিটার/সেকেন্ড।}

\emph{ব্যক্তিকে অতিক্রম করতে ট্রেনকে তার নিজের দৈর্ঘ্য (২২০ মিটার) অতিক্রম করতে
হবে।}

\emph{∴ প্রয়োজনীয় সময়, t = (২২০/(৫৫/৩)) = ২২০×(৩/৫৫) = ১২ সেকেন্ড।}
\end{quote}

\section{56. মিশ্রণ সংক্রান্ত সকল অংকের সহজ নিয়ম \textbar\textbar{}
Allegation}\label{ux9aeux9b6ux9b0ux9a3-ux9b8ux995ux9b0ux9a8ux9a4-ux9b8ux995ux9b2-ux985ux995ux9b0-ux9b8ux9b9ux99c-ux9a8ux9dfux9ae-allegation}

মিশ্রণ সংক্রান্ত গাণিতিক সমস্যাগুলো মূলত পাটিগণিতের নির্দিষ্ট কোনো পৃথক অধ্যায় নয়।
তবে পরীক্ষার হলে এই ধরনের প্রশ্নগুলোই প্রায়শই শিক্ষার্থীদের সবচেয়ে বেশি দ্বিধা ও
বিভ্রান্তিতে ফেলে দেয়। অনেক সময় একটি অংকের পেছনে অতিরিক্ত চিন্তা করতে গিয়ে
মূল্যবান সময় নষ্ট হয় এবং পুরো সমাধানটি এলোমেলো হয়ে যায়।

পরীক্ষার হলের এই জটিলতা ও দুশ্চিন্তা দূর করতেই মিশ্রণের যাবতীয় গাণিতিক সমস্যা
খুঁটিনাটি সমাধানসহ একত্রে উপস্থাপন করা হলো। এটি অনুধাবন করলে প্রশ্নে চোখ
বোলামাত্রই দ্রুত ও নির্ভুল সমাধান করা সম্ভব হবে।

চলুন, কৌশলগুলো দেখে নেওয়া যাক---

\begin{itemize}
\item
  অনুপাত আকারে মিশ্রণ
\item
  শতকরা আকারে মিশ্রণ
\item
  মিশ্রণের অনুপাত বের করার নিয়ম
\end{itemize}

\subsubsection{অনুপাত আকারে
মিশ্রণ}\label{ux985ux9a8ux9aaux9a4-ux986ux995ux9b0-ux9aeux9b6ux9b0ux9a3}

অনুপাত আকারে মিশ্রণ সংক্রান্ত অংকগুলোর বেশ কয়েক ধরনের সমস্যা থাকে যেমন-

\textbf{সূত্র-১. যদি দ্বিতীয় অনুপাত প্রথম অনুপাতের বিপরীত হয় তখন-}

\textbf{মিশ্রিত দ্রব্যের পরিমাণ}

\textbf{=} \(\frac{\text{মোট}\ \text{মিশ্রণের}\ \text{পরিমাণ}}{\text{প্রথম}\ \text{অনুপাতের}\ \text{ছোট}\ \text{সংখ্যা}}\)
\textbf{× দ্বিতীয় অনুপাতের পার্থক্য}

প্রশ্ন-১: ৩০ লিটার পরিমাণ মিশ্রণে এসিড ও পানির অনুপাত ৭ :৩। এ মিশ্রণে কি
পরিমাণ পানি মিশ্রিত করলে এসিড ও পানির অনুপাত ৩ :৭ হবে ? {[} প্রাথমিক প্রধান
শিক্ষক(পদ্মা)-২০১২{]}

\begin{itemize}
\item
  (ক) ৩০ লিটার
\item
  (খ) ২৫ লিটার
\item
  (গ) ৪০ লিটার
\item
  (ঘ) ৪৫ লিটার
\end{itemize}

\textbf{\hl{উত্তরঃ (গ) ৪০ লিটার}}

\textbf{Explanation:}

পানি মেশাতে হবে

= \(\frac{\text{মোট}\ \text{মিশ্রণের}\ \text{পরিমাণ}}{\text{প্রথম}\ \text{অনুপাতের}\ \text{ছোট}\ \text{সংখ্যা}}\) × দ্বিতীয়
অনুপাতের পার্থক্য

= \(\frac{\text{৩০}}{\text{৩}}\) × ৪

= ১০ × ৪ = ৪০ লিটার

প্রশ্ন-২: ২১ গ্যালন অকটেন মিশ্রিত পেট্রোলে, পেট্রোল ও অকটনের অনুপাত ৪:৩। এত আর
কত অকটেন মিশালে পেটোল ও অকটেনের অনুপাত ৩:৪ হবে ? {[}প্রাথমিক প্রধান
শিক্ষক(রাজশাহী)-২০০৯{]}

\begin{itemize}
\item
  (ক) ৩ গ্যালন
\item
  (খ) ৪ গ্যালন
\item
  (গ) ৭ গ্যালন
\item
  (ঘ) ৬ গ্যালন
\end{itemize}

\textbf{\hl{উত্তরঃ (গ) ৭ গ্যালন}}

\textbf{Explanation:}

অকটেন মেশাতে হবে

= \(\frac{\text{মোট}\ \text{মিশ্রণের}\ \text{পরিমাণ}}{\text{প্রথম}\ \text{অনুপাতের}\ \text{ছোট}\ \text{সংখ্যা}}\) × দ্বিতীয়
অনুপাতের পার্থক্য

= \(\frac{\text{২১}}{\text{৩}}\) × ১

= ৭ × ১ = ৭ গ্যালন

\textbf{সূত্র-২: যখন দুইটি অনুপাতের সংখ্যাদ্বয়ের একটি সংখ্যা পরিবর্তিত হয় তখন -}

\textbf{মিশ্রিত দ্রব্যের পরিমাণ =}
\(\frac{\text{মোট}\ \text{মিশ্রণের}\ \text{পরিমাণ}}{\text{প্রথম}\ \text{অনুপাতের}\ \text{যোগফল}}\) \textbf{× পরিবর্তিত
অনুপাত সংখ্যার পার্থক্য}

প্রশ্ন-১: ২৫ গ্রাম ওজনের একটি সোনা গহনায় সোনা ও তামার অনুপাত ৪ : ১। গহনাটিতে
আর কতটুকু সোনা মেশালে এতে সোনা ও তামার অনুপাত ৫ : ১ হবে ? {[} খাদ্য
পরিদর্শক-২০০১{]}

\begin{itemize}
\item
  (ক) ৫ গ্রাম
\item
  (খ) ৬ গ্রাম
\item
  (গ) ১০ গ্রাম
\item
  (ঘ) ২০ গ্রাম
\end{itemize}

\textbf{\hl{উত্তরঃ (ক) ৫ গ্রাম}}

\textbf{Explanation:}

সোনা মেশাতে হবে = \(\frac{\text{মোট}\ \text{মিশ্রণের}\ \text{পরিমাণ}}{\text{প্রথম}\ \text{অনুপাতের}\ \text{যোগফল}}\)
× পরিবর্তিত অনুপাত সংখ্যার পার্থক্য

= \(\frac{\text{২৫}}{\text{৫}}\) × (৫ - ৪ )

= \(\frac{\text{২৫}}{\text{৫}}\) × ১

= ৫ × ১ = ৫ গ্রাম

প্রশ্ন-২: একটি সোনার গহনার ওজন ১৬ গ্রাম । তাতে সোনার পরিমাণঃ তামার পরিমাণ =
৩:১ । তাতে আর কি পরিমাণ সোনা মেশালে অনুপাত ৪:১ হবে ? {[} ১৭তম এবং ২১ তম
বিসিএস{]}

\begin{itemize}
\item
  (ক) ১২ গ্রাম
\item
  (খ) ৪ গ্রাম
\item
  (গ) ৩ গ্রাম
\item
  (ঘ) ৬ গ্রাম
\end{itemize}

\textbf{\hl{উত্তরঃ (খ) ৪ গ্রাম}}

\textbf{Explanation:}

সোনা মেশাতে হবে = \(\frac{\text{মোট}\ \text{মিশ্রণের}\ \text{পরিমাণ}}{\text{প্রথম}\ \text{অনুপাতের}\ \text{যোগফল}}\)
× পরিবর্তিত অনুপাত সংখ্যার পার্থক্য

= \(\frac{\text{১৬}}{\text{৪}}\) × (৪ - ৩ )

= \(\frac{\text{১৬}}{\text{৪}}\) × ১

= ৪ × ১ = ৪ গ্রাম

\textbf{সূত্র-৩: যখন দুইটি অনুপাতের সংখ্যাদ্বয়ের দুইটি সংখ্যা পরিবর্তিত হয় তখন -}

\textbf{মিশ্রিত দ্রব্যের পরিমাণ =}
\(\frac{\text{মোট}\ \text{মিশ্রণের}\ \text{পরিমাণ}}{\text{প্রথম}\ \text{অনুপাতের}\ \text{যোগফল}}\) \textbf{÷ দ্বিতীয়
অনুপাতের বড় সংখ্যা}

প্রশ্ন-১: ৪০ গ্যালন অকটেন মিশ্রিত পেট্রোলে , পেট্রোল ও অকটেনের অনুপাত ৩:১। এতে
আর কত অকটেন মিশালে পেট্রোল ও অকটেনের অনুপাত ৫:২ হবে ? {[}প্রাথমিক সহকারী
সহকারী শিক্ষক : ০৯ {]}

\begin{itemize}
\item
  (ক) ৩ গ্যালন
\item
  (খ) ২ গ্যালন
\item
  (গ) ৫ গ্যালন
\item
  (ঘ) ৮ গ্যালন
\end{itemize}

\textbf{\hl{উত্তরঃ (খ) ২ গ্যালন}}

\textbf{Explanation:}

অকটেন মিশাতে হবে = \(\frac{\text{মোট}\ \text{মিশ্রণের}\ \text{পরিমাণ}}{\text{প্রথম}\ \text{অনুপাতের}\ \text{যোগফল}}\)
÷ দ্বিতীয় অনুপাতের বড় সংখ্যা

= \(\frac{\text{৪০}}{\text{৪}}\) ÷ ৫

= ১০ ÷ ৫ = ২ গ্যালন

প্রশ্ন-২: একটি ঝুড়িতে ১৩০ টি আম ও পেয়ারা আছে এবং এদের অনুপাত যথাক্রমে ৩:২ ।
উক্ত ঝুড়িতে আম ও পেয়ারার অনুপাত ১:১ করতে হলে , কমপক্ষে কতটি নতুন ফল যোগ করতে
হবে? {[} জুনিয়র অডিটর: ২০১৯ {]}

\begin{itemize}
\item
  (ক) ৫২টি
\item
  (খ) ১৩টি
\item
  (গ) ১৩০টি
\item
  (ঘ) ২৬টি
\item
  (ঙ) কোনটিই নয়
\end{itemize}

\textbf{\hl{উত্তরঃ (ঘ) ২৬টি}}

\textbf{Explanation:}

এখানে , মোট মিশ্রণের পরিমাণের জায়গায় ১৩০টি আম ও পেয়ারা হবে।

নতুন ফল যোগ করতে হবে = \(\frac{\text{মোট}\ \text{ফলের}\ \text{পরিমাণ}}{\text{প্রথম}\ \text{অনুপাতের}\ \text{যোগফল}}\)
÷ দ্বিতীয় অনুপাতের বড় সংখ্যা

= \(\frac{\text{১৩০}}{\text{৫}}\) ÷ ১

= ২৬ ÷ ১ = ২৬ সুতরাং ২৬টি নতুন ফল যোগ করতে হবে।

\subsubsection{শতকরা আকারে
মিশ্রণ}\label{ux9b6ux9a4ux995ux9b0-ux986ux995ux9b0-ux9aeux9b6ux9b0ux9a3}

প্রশ্ন-১: ১২. ৬৪ কিলোগ্রাম বালি ও পাথরের টুকরোর মিশ্রণে বালির পরিমাণ ২৫\%। কত
কিলোগ্রাম বালি মিশালে নতুন মিশ্রণে পাথর টুকরোর পরিমাণ ৪০\% হবে? {[} ১৫তম
বিসিএস {]}

\begin{itemize}
\item
  (ক) ৯.৬
\item
  (খ) ১১.০
\item
  (গ) ৪৮.০
\item
  (ঘ) ৫৬.০
\end{itemize}

\textbf{\hl{উত্তরঃ (ঘ) ৫৬.০}}

\textbf{Explanation:}

মিশ্রনে বালির পরিমান = ৬৪ × ২৫\% = ১৬ কেজি

মিশ্রনে পাথরের পরিমান = (৬৪ - ১৬) = ৪৮ কেজি

নতুন মিশ্রনে ৪০\% = ৪৮ কেজি

নতুন মিশ্রনে ১০০\% = ৪৮ × \(\frac{\text{১০০}}{\text{৪০}}\) = ১২০ কেজি

অতএব, নতুন মিশ্রনে বালির পরিমান = (১২০ - ৪৮) = ৭২কেজি

অতএব, বালি মেশাতে হবে = (৭২ - ১৬) = ৫৬ কেজি ( উত্তর)

প্রশ্ন-২: ১৩. ৪০ কেজি জৈব ও ইউরিয়া সারের মিশ্রণে জৈব সারের পরিমাণ ১০\% কত
কেজি জৈব সার মিশালে নতুন মিশ্রণে জৈব সারের পরিমাণ ২০\% হবে? {[} বিএডিসি
-২০১৭ {]}

\begin{itemize}
\item
  (ক) ৪
\item
  (খ) ৫
\item
  (গ) ৬
\item
  (ঘ) ৭
\item
  (ঙ) কোনটিই নয়
\end{itemize}

\textbf{\hl{উত্তরঃ (খ) ৫}}

\textbf{Explanation:}

জৈব সারের পরিমাণ = ৪০ × ১০\% = ৪ কেজি

মিশ্রণে ইউরিয়া সারের পরিমাণ = (৪০ - ৪) = ৩৬ কেজি

ধরি, জৈব সার মেশাতে `ক' কেজি

প্রশ্নমতে, ৪ + ক = (৪০ + ক ) এর ২০\%

বা, ২০+৫ক = ৪০ + ক

বা, ৪ক = ২০

বা, ক = ৫

সুতরাং জৈব সার মেশাতে হবে ৫ কেজি ।

\subsubsection{মিশ্রণের অনুপাত বের করার
নিয়ম}\label{ux9aeux9b6ux9b0ux9a3ux9b0-ux985ux9a8ux9aaux9a4-ux9acux9b0-ux995ux9b0ux9b0-ux9a8ux9dfux9ae}

প্রশ্ন-১: একজন দোকানদার ১৫ টাকা ও ২০ টাকা কেজি দরে ২ ধরনের চা কি অনুপাতে
মেশালে ১৭ টাকা দাম হবে? {[} সাব রেজিস্ট্রার : ২০০১{]}

\begin{itemize}
\item
  (ক) ৫:২
\item
  (খ) ৪:১
\item
  (গ) ৩:২
\item
  (ঘ) ৩:১
\end{itemize}

\section{57. বর্গ ও বর্গমূল নির্ণয়ের সর্বশেষ সহজ
টেকনিক}\label{ux9acux9b0ux997-ux993-ux9acux9b0ux997ux9aeux9b2-ux9a8ux9b0ux9a3ux9dfux9b0-ux9b8ux9b0ux9acux9b6ux9b7-ux9b8ux9b9ux99c-ux99fux995ux9a8ux995}

পাবলিক সার্ভিস কমিশন (বিসিএস), ব্যাংক, প্রাথমিক শিক্ষক নিয়োগ কিংবা যেকোনো
প্রতিযোগিতামূলক পরীক্ষার গণিত অংশে সফল হওয়ার জন্য \textbf{বর্গ ও বর্গমূল}
(Square \& Square Root) অধ্যায়টির ওপর দখল থাকা অত্যন্ত জরুরি। গণিতের মৌলিক
এই ধারণাটি স্পষ্ট না থাকলে অনেক সহজ প্রশ্নের সমাধান করতে গিয়েও পরীক্ষার্থীদের
অমূল্য সময় নষ্ট হয়ে যায়।

বর্গমূল নির্ণয়ের প্রচলিত নিয়মের পাশাপাশি কম সময়ে নিখুঁত সমাধান করার কিছু বিশেষ
শর্টকাট কৌশল জানা থাকলে পরীক্ষায় অন্যদের চেয়ে অনেক এগিয়ে থাকা সম্ভব। বিশেষ
করে স্বাভাবিক সংখ্যা, দশমিক ভগ্নাংশ ও পূর্ণবর্গ সংখ্যা তৈরির টেকনিকগুলো আয়ত্ত
করতে পারলে গাণিতিক প্রশ্নের উত্তর দেওয়ার গতি বহুগুণ বৃদ্ধি পায়।

আপনার গণিত ভীতি দূর করতে এবং প্রতিযোগিতামূলক পরীক্ষার প্রস্তুতিকে আরও শাণিত
করতে এই আর্টিকেলে বর্গ ও বর্গমূল নির্ণয়ের বেসিক নিয়ম থেকে শুরু করে অপশন দেখে
দ্রুত উত্তর বের করার শর্টকাট পদ্ধতিগুলো সহজ ভাষায় উপস্থাপন করা হলো। নিচে দেওয়া
নিয়মগুলো অনুধাবন করুন এবং বিগত বছরের প্রশ্নগুলো অনুশীলনের মাধ্যমে নিজের প্রস্তুতিকে
পূর্ণাঙ্গ রূপ দিন।

\begin{itemize}
\item
  বর্গ ও বর্গমূল
\item
  স্বাভাবিক সংখ্যার বর্গমূল বের করার নিয়ম
\item
  দশমিক ভগ্নাংশের বর্গমূল বের করার নিয়ম
\item
  অপশন ধরে দ্রুত বর্গমূল বের করার শর্টকাট নিয়ম
\item
  পূর্ণ বর্গ সংখ্যা তৈরি করা
\end{itemize}

\subsubsection{বর্গ ও
বর্গমূল}\label{ux9acux9b0ux997-ux993-ux9acux9b0ux997ux9aeux9b2}

\textbf{বর্গ ও বর্গমূলঃ}

কোন সংখ্যাকে একই সংখ্যা দ্বারা গুণ করলে যে গুণফল পাওয়া যায়, তাকে ঐ সংখ্যার বর্গ
বলে এবং সংখ্যাটিকে গুণফলের বর্গমূল বলা হয়।

যেমন, ৪×৪ বা ১৬ কে ৪ এর বর্গ এবং ৪ কে ১৬ এর বর্গমূল বলা হয়। বর্গমূল প্রকাশের
জন্য দুইটি প্রতীক চিহ্ন ব্যবহৃত হয়। ১৬ এর বর্গমূল বোঝাতে লেখা হয় √১৬ বা
\({(\text{১৬})}^{\text{১}\ \text{২}}\) । সব সংখ্যারই বর্গমূল আছে , কিন্তু সব সংখ্যারই বর্গমূল পূর্ণ
সংখ্যায় নেই ।যেমন- ৩ এর বর্গমূল ১.৭৩২ কিন্তু ৪ এর বর্গমূল ২ ।

\textbf{পূর্ণ বর্গসংখ্যাঃ} যে সংখ্যার বর্গমূল কোন পূর্ণসংখ্যা বা ভগ্নাংশের সমান,
তাকে পূর্ণবর্গ বলা হয়। যেমন- ৪,৯,১৬, \(\frac{\text{৯}}{\text{১৬}}\) , \(\frac{\text{২৫}}{\text{৩৬}}\) ।

\textbf{কিছু গুরুত্বপূর্ণ তথ্যঃ}

কোন পূর্ণ বর্গসংখ্যার একক স্থানীয় অঙ্ক ০,১,৪,৫,৬ বা ৯ হবে।

কোন পূর্ণ বর্গসংখ্যার একক স্থানীয় অঙ্ক ২,৩,৭ বা ৮ হতে পারে না।

যে সংখ্যার শেষে বিজোড় সংখ্যক শূন্য থাকে, ঐ সংখ্যা পূর্ণবর্গ নয়।

কোন বিজোড় স্বাভাবিক সংখ্যার বর্গকে ৮ দ্বারা ভাগ করলে প্রতিক্ষেত্রে ভাগশেষ ১
হবে।

একটি শূন্য যুক্ত সংখ্যার উপর বর্গ দিলে তা দ্বিগুণ হয়ে যায় , তেমনি দুৃটি শূন্য যুক্ত
সংখ্যার বর্গমূল বের করলে তা একটি হয়ে যায়। যেমন- ১০০ এর বর্গ ১০০০০ আবার ১০০০০
এর বর্গমূল ১০০ ।

\textbf{করণীঃ} যে সকল সংখ্যা পূর্ণবর্গ নয়, যাদের বর্গমূল সঠিকভাবে নির্ণয় করা
যায়না, ঐ ধরনের বর্গমূলকে করণী বা অমূলদ সংখ্যা বলা হয়। যেমন-√৩, √৫, √৫০
ইত্যাদি।

কোনো সংখ্যার একক স্থানীয় অঙ্ক থেকে শুরু করে বামদিকে এক অঙ্ক পরপর যতটি ফোঁটা
দেওয়া যায়, এর বর্গমূলের সংখ্যাটি তত অঙ্কবিশিষ্ট ।

\subsubsection{স্বাভাবিক সংখ্যার বর্গমূল বের করার
নিয়ম}\label{ux9b8ux9acux9adux9acux995-ux9b8ux996ux9afux9b0-ux9acux9b0ux997ux9aeux9b2-ux9acux9b0-ux995ux9b0ux9b0-ux9a8ux9dfux9ae}

দু পদ্ধতিতে বর্গমূল নির্ণয় করা যায় ।যথাঃ ১। উৎপাদক পদ্ধতি ২। ভাগ পদ্ধতি।

\textbf{উৎপাদক পদ্ধতিতে বর্গমূল নির্ণয়ের নিয়মঃ}

প্রথমে প্রদত্ত সংখ্যাকে মৌলিক উৎপাদকে বিশ্লেষণ করতে হয়।

প্রতিজোড়া একই উৎপাদককে একসাথে পাশাপাশি লিখতে হয়।

প্রতিজোড়া একজাতীয় উৎপাদকের পরিবর্তে একটি উৎপাদক নিয়ে লিখতে হয়।

\textbf{লিখিত উৎপাদকগুলোর ধারাবাহিক গুণফল হবে নির্ণেয় বর্গমূল।যেমন-৩১৩৬ এর
বর্গমূল কত?}

২

৩১৩৬

২

১৫৬৮

২

৭৮৪

২

৩৯২

২

১৯৬

২

৯৮

৭

৪৯

৭

এখানে, ৩১৩৬ = ২×২×২×২×২×২×৭×৭

= (২×২)×(২×২)×(২×২)×(৭×৭) সুতরাং ৩১৩৬ এর বর্গমূল = ২×২×২×৭ = ৫৬

\textbf{ভাগ প্রক্রিয়ার সাহায্যে বর্গমূল নির্ণয়ের নিয়মঃ}

একটি উদাহরণের সাহায্যে ভাগ প্রক্রিয়ার সাহায্যে সাহায্যে বর্গমূল নির্ণয়ের পদ্ধতি
দেখানো হলঃ

৩১৬৮৪ সংখ্যাটি লিখি ।

৩

১৬

৮৪

১৭৮

১

২৭

২

১

৬

১

৮

৯

৩৪৮

২

৭

৮

৪

২

৭

৮

৪

০

ডানদিক থেকে দুইটি করে অঙ্ক নিয়ে জোড়া করি । প্রত্যেক জোড়ার উপর রেখাচিহ্ন দিই ।

ভাগের সময় যেমন খাড়া দাগ দেওয়া হয়, ডানপাশে তদ্রূপ একটি খাড়া দাগ দিই ।

দেখা যাচ্ছে ডানদিকে দুইটি জোড়া ৮৪ ও ১৬ নেওয়ার পর ৩ অবশিষ্ট থাকে। জোড়া
অবশিষ্ট না থাকলে একটি সংখ্যার সাথেই কাজ করতে হবে। ৩ এর পূর্ববর্তী বর্গসংখ্যাটি
১, যার বর্গমূল √১ বা ১ ; খাড়া দাগের ডানপাশে ১ লিখি ।

এখন ৩ এর ঠিক নিচে ১ লিখি ।

এখন ৩ থেকে ১ বিয়োগ করি ।

বিয়োগফল ২ এর ডানে পরবর্তী জোড়া ১৬ বসাই। ২১৬ এর বামদিকে খাড়া দাগ (ভাগের
চিহ্ন) দিই ।

ভাগফলের ঘরের সংখ্যা ১ এর দ্বিগুণ ২ নিচের খাড়া দাগের বামপাশে বসাই। ২ এবং
খাড়া দাগের মধ্যে একটি অঙ্ক বসানোর মতো স্থান রাখি ।

এখন একটি এক অঙ্কের সংখ্যা খুঁজে বের করি যাকে ২ এর ডানপাশে বসিয়ে প্রাপ্ত
সংখ্যাকে ঐ সংখ্যাটি দ্বারা গুণ করে ২১৬ এর সমান বা অনুর্ধ্ব ২১৬ পাওয়া যায়।

এক্ষেত্রে ৭ হবে। ৭ সংখ্যাটি ভাগফলেও ১ এর ডানপাশে বসাই।

ভাগফলের স্থানে পাওয়া গেল ১৭।

এখন ২১৬ এর ঠিক নিচে ১৮৯ লিখি ।

এখন ২১৬ থেকে ১৮৯ বিয়োগ করি ।

বিয়োগফল ২৭ এর ডানে পরবর্তী জোড়া ৮৪ বসাই। ২৭৮৪ এর বামদিকে খাড়া দাগ (ভাগের
চিহ্ন) দিই ।

ভাগফলের ঘরের সংখ্যা ১৭ এর দ্বিগুণ ৩৪ নিচের খাড়া দাগের বামপাশে বসাই। ৩৪ এবং
খাড়া দাগের মধ্যে একটি অঙ্ক বসানোর মতো স্থান রাখি ।

এখন আবার ও একটি এক অঙ্কের সংখ্যা খুঁজে বের করি যাকে ১৭ এর ডানপাশে বসিয়ে
প্রাপ্ত সংখ্যাকে ঐ সংখ্যাটি দ্বারা গুণ করে ২৭৮৪ এর সমান বা অনুর্ধ্ব ২৭৮৪ পাওয়া
যায়।

এক্ষেত্রে ৮ হবে। ৮ সংখ্যাটি ভাগফলেও ১৭ এর ডানপাশে বসাই।

ভাগফলের স্থানে পাওয়া গেল ১৭৮ । এটিই নির্ণেয় বর্গমূল।

লক্ষণীয় যে ভাগের সাহায্যে বর্গমূল নির্ণয় করার সময় সংখ্যার ডান দিক অঙ্কের জোড় না
থাকলে একে জোড়া ছাড়াই গণ্য করতে হবে।

সুতরাং ৩১৬৮৪ এর বর্গমূল ১৭৮ ।

\subsubsection{দশমিক ভগ্নাংশের বর্গমূল বের করার
নিয়ম}\label{ux9a6ux9b6ux9aeux995-ux9adux997ux9a8ux9b6ux9b0-ux9acux9b0ux997ux9aeux9b2-ux9acux9b0-ux995ux9b0ux9b0-ux9a8ux9dfux9ae}

\textbf{দশমিক ভগ্নাংশের বর্গমূল}

পূর্ণসংখ্যা বা অখণ্ড সংখ্যার বর্গমূল ভাগের সাহায্যে যেভাবে নির্ণয় করা হয়েছে,
দশমিক ভগ্নাংশের বর্গমূলও সেই নিয়মেই নির্ণয় করা হয়। দশমিক ভগ্নাংশের দুইটি অংশ
থাকে। দশমিক বিন্দুর বামদিকের অংশকে অখণ্ড বা পূর্ণ অংশ এবং দশমিক বিন্দুর
ডানপাশের অংশকে দশমিক অংশ বলা হয়।

\textbf{দশমিক ভগ্নাংশের বর্গমূল করার নিয়মঃ}

অখণ্ড অংশে একক থেকে ক্রমান্বয়ে বামদিকে প্রতি দুই অঙ্কের উপর দাগ দিতে হয়।

দশমিক অংশে দশমিক বিন্দুর ডানপাশের অঙ্ক থেকে শুরু করে ডানদিকে ক্রমান্বয়ে জোড়ায়
জোড়ায় দাগ দিতে হয়। এরূপে যদি দেখা যায় সর্বশেষে মাত্র একটি অঙ্ক বাকি আছে, তবে
তারপরে একটি শূন্য বসিয়ে দুই অঙ্কের উপর দাগ দিতে হয়।

সাধারণ নিয়মে বর্গমূল নির্ণয়ের প্রক্রিয়ায় অখণ্ড অংশের কাজ শেষ করে দশমিক বিন্দুর
পরের প্রথম দুইটি অঙ্ক নামানোর আগেই বর্গমূলে দশমিক বিন্দু দিতে হয়।

দশমিক বিন্দুর এক জোড়া শূন্যের জন্য বর্গমূলে দশমিক বিন্দুর পর একটি শূন্য দিতে হয়।
নিচে উদাহরণের সাহায্যে বিষয়টি দেখানো হল ।

\textbf{উদাহরণ-১:} ২৬.৫২২৫ এর বর্গমূল নির্ণয় কর।

২৬.৫২২৫

৫.১৫

২৫

১০১

১৫২

১০১

১০২৫

৫১২৫

৫১২৫

০

\textbf{উদাহরণ-২:} ০.০০২৯১৬ এর বর্গমূল নির্ণয় কর।

০.০০২৯১৬

০.০৫৪

২৫

১০৪

৪১৬

৪১৬

০

\subsubsection{অপশন ধরে দ্রুত বর্গমূল বের করার শর্টকাট
নিয়ম}\label{ux985ux9aaux9b6ux9a8-ux9a7ux9b0-ux9a6ux9b0ux9a4-ux9acux9b0ux997ux9aeux9b2-ux9acux9b0-ux995ux9b0ux9b0-ux9b6ux9b0ux99fux995ux99f-ux9a8ux9dfux9ae}

কোন বর্গ সংখ্যার শেষে ৫ থাকলে তার বর্গমূলের সংখ্যাটির শেষেও ৫ থাকবে ।

কোন বর্গ সংখ্যার শেষে ৬ থাকলে তার বর্গমূলের সংখ্যাটির শেষে ৪ অথবা ৬ থাকবে ।

কোন বর্গ সংখ্যার শেষে ১ থাকলে তার বর্গমূলের সংখ্যাটির শেষেও ১ ই থাকবে ।

কোন বর্গ সংখ্যার শেষে ৪ থাকলে তার বর্গমূলের সংখ্যাটির শেষে ২ থাকবে ।

কোন বর্গ সংখ্যার শেষে ০ থাকলে তার বর্গমূলের সংখ্যাটির শেষেও ০ থাকবে ।

নিচে এরুপ কিছু উদাহরণ দেখুন:

উদাহরণ: ১. ৭৩৯৬ এর বর্গমূল কত? {[} সাব-রেজিস্টার পরীক্ষা-২০০১{]}

\begin{itemize}
\item
  (ক) ৮২
\item
  (খ) ৭৫
\item
  (গ) ৮০
\item
  (ঘ) ৮৬
\end{itemize}

\textbf{\hl{উত্তরঃ (ঘ) ৮৬}}

\textbf{ব্যাখ্যা:} {[} শেষে ৬ আছে তাই উত্তর ৮৬ {]}

উদাহরণ: ২. একটি শ্রেণীতে যতজন ছাত্র আছে প্রত্যেকে তত টাকা করে প্রদান করলে মোট
৬৫৬১ টাকা হয়। ছাত্র সংখ্যা কত ? {[}কর্মসংস্থান ব্যাংক সহকারী অফিসার-২০০১{]}

\begin{itemize}
\item
  (ক) ৮০
\item
  (খ) ৮৫
\item
  (গ) ৮১
\item
  (ঘ) ৯২
\end{itemize}

\textbf{\hl{উত্তরঃ (গ) ৮১}}

\textbf{ব্যাখ্যা:} {[} শেষে ১ আছে তাই উত্তর ৮১ {]}

আবার ঘুরিয়ে আসলে ও উপরের নিয়মগুলোই প্রয়োগ করবেন।যেমন:

উদাহরণ: ৩. ১৪৪ এর বর্গ কত? {[} সাব-রেজিস্টার পরীক্ষা-২০০১{]}

\begin{itemize}
\item
  (ক) ২০৭৫৪
\item
  (খ) ২০৭৩৬
\item
  (গ) ২৫৭৫০
\item
  (ঘ) ২০৭২২
\end{itemize}

\textbf{\hl{উত্তরঃ (খ) ২০৭৩৬}}

\textbf{ব্যাখ্যা:} {[} বর্গের ক্ষেত্রে শেষে ৪ এর থাকলে এর বর্গ সংখ্যাতে শেষে ৬
হবে {]}

কখনো একাধিক অপশনের মধ্যে সন্দেহ হলে যে অপশন গুলো বাদ দেয়া যাবে তা বাদ দিয়ে
অন্য গুলো থেকে উত্তর বের করতে হবে।যেমন:

উদাহরণ: ৪. একটি ছাত্রাবাসে কিছু ছাত্র থাকে। যতজন ছাত্র তত টাকা করে চাঁদা
দেয়াতে যদি মোট ৫৭৬ টাকা উঠে তাহলে ঐ ছাত্রাবাসে মোট কতজন ছাত্র থাকে? {[}
সাব-রেজিস্টার পরীক্ষা-২০০১{]}

\begin{itemize}
\item
  (ক) ২৪
\item
  (খ) ৩৫
\item
  (গ) ২৮
\item
  (ঘ) ২৬
\end{itemize}

\textbf{\hl{উত্তরঃ (ক) ২৪}}

\textbf{ব্যাখ্যা:} প্রথমেই প্রশ্নের ৫৭৬ দেখে বুঝে নিতে হবে বর্গমূলের সংখ্যাটির
শেষ ৪ অথবা ৬ হবে।তাহলে (খ) এবং (গ) অপশনটি খুব সহজেই বাদ দেয়া যায়।এখন ক ও
খ এর মধ্যে ২৪ কে ২৪ আবার ৩৫ কে ৩৫ দ্বারা গুণ করে বের করতে হবে কোনটি সঠিক
উত্তর।এখানে সঠিক উত্তর (ক) ২৪ ।

\subsubsection{পূর্ণ বর্গ সংখ্যা তৈরি
করা}\label{ux9aaux9b0ux9a3-ux9acux9b0ux997-ux9b8ux996ux9af-ux9a4ux9b0-ux995ux9b0}

\textbf{২১৯৫২ সংখ্যাটিকে কত দ্বারা গুণ করলে সংখ্যাটি পূর্ণবর্গ সংখ্যা হবে?}

যে কোন পূর্ণ বর্গ সংখ্যার ভেতরে অবস্থিত উৎপাদক গুলো সব সময় জোড়ায় জোড়ায়
থাকে।তাই কোন সংখ্যার উৎপাদক জোড় না থাকলে যে সংখ্যার জোড়া থাকবে না তাকে
দিয়ে গুণ করতে হবে।আবার ভাগ করতে বলা হলে যার জোড়া নেই তাকে দিয়ে ভাগ করতে
হবে।

\textbf{সমাধান:}

২

২১৯৫২

২

১০৯৭৬

২

৫৪৮৮

২

২৭৪৪

২

১৩৭২

২

৬৮৬

৭

৩৪৩

৭

৪৯

৭

সুতরাং ২১৯৫২ = ২×২×২×২×২×২×৭×৭×৭

= (২×২)×(২×২)×(২×২)×(৭×৭)×৭ যেহেতু এখানে একটি ৭ জোড়াবিহীন অবস্থায়
রয়েছে।তাই ৭ দ্বারা গুণ করলে সংখ্যাটি পূর্ণবর্গ হবে।

সময় অবশিষ্ট আছে:

\section{58. বর্গ ও বর্গমূল MCQ প্রশ্ন উত্তর
অনুশীলন}\label{ux9acux9b0ux997-ux993-ux9acux9b0ux997ux9aeux9b2-mcq-ux9aaux9b0ux9b6ux9a8-ux989ux9a4ux9a4ux9b0-ux985ux9a8ux9b6ux9b2ux9a8}

\textbf{√(১৫.৬০২৫)=?}

ক) ৩.৮৫

খ) ৩.৭৫

\textbf{\hl{গ) ৩.৯৫ ✔}}

ঘ) ৩.৬৫

\textbf{\hl{✔ সঠিক উত্তর: গ) ৩.৯৫}}

\begin{quote}
\textbf{ব্যাখ্যা:} \emph{√(১৫.৬০২৫)=৩.৯৫}

\emph{হিসাব:}

\emph{১৫.৬০২৫ এর বর্গমূল করলে পাওয়া যায় ৩.৯৫।}
\end{quote}

\textbf{√(০.০০০০০৯) কত?}

ক) ০.০৩

খ) ০.৩

\textbf{\hl{গ) ০.০০৩ ✔}}

ঘ) ০.০০০৩

\textbf{\hl{✔ সঠিক উত্তর: গ) ০.০০৩}}

\begin{quote}
\textbf{ব্যাখ্যা:} \emph{√(০.০০০০০৯)=√((৯/১০০০০০০০))=(৩/১০০০)=০.০০৩}
\end{quote}

\textbf{১০২৪ এর বর্গমূল কত ?}

\textbf{\hl{ক) ৩২ ✔}}

খ) ৬৪

গ) ১৬

ঘ) ২৮

\textbf{\hl{✔ সঠিক উত্তর: ক) ৩২}}

\begin{quote}
\textbf{ব্যাখ্যা:} \emph{এখানে ১০২৪ কে উৎপাদকে বিশ্লেষণ করলে দাঁড়ায়:}

\emph{১০২৪ = ২ × ২ × ২ × ২ × ২ × ২ × ২ × ২ × ২ × ২}

\emph{১০২৪ = (২ × ২) × (২ × ২) × (২ × ২) × (২ × ২) × (২ × ২)}

\emph{∴ ১০২৪ এর বর্গমূল = ২ × ২ × ২ × ২ × ২ = ৩২}
\end{quote}

\textbf{০.০০০১ এর বর্গমূল কত?}

ক) ০.১

\textbf{\hl{খ) ০.০১ ✔}}

গ) ০.০০১

ঘ) ১

\textbf{\hl{✔ সঠিক উত্তর: খ) ০.০১}}

\begin{quote}
\textbf{ব্যাখ্যা:} \emph{√(০.০০০১)=√((১/১০০))=(১/১০)=০.০১}
\end{quote}

\textbf{√(০.০০০০০৬২৫)= কত?}

\textbf{\hl{ক) ০.০০২৫ ✔}}

খ) ০.০০০২৫

গ) ০.০০০০০২৫

ঘ) ০.০০৬২৫

\textbf{\hl{✔ সঠিক উত্তর: ক) ০.০০২৫}}

\begin{quote}
\textbf{ব্যাখ্যা:}
\emph{√(০.০০০০০৬২৫)=√((৬২৫/১০০০০০০০))=(২৫/১০০০০)=০.০০২৫}
\end{quote}

\textbf{√(.২৫) এর মান কত?}

ক) ০.০৫

\textbf{\hl{খ) ০.৫ ✔}}

গ) ০.০০৫

ঘ) ০.০২

\textbf{\hl{✔ সঠিক উত্তর: খ) ০.৫}}

\begin{quote}
\textbf{ব্যাখ্যা:} \emph{√(.২৫)=√((২৫/১০০))=(৫/১০)=০.৫}
\end{quote}

\textbf{√(০.০৯)= কত ?}

ক) ০.০৩

\textbf{\hl{খ) ০.৩ ✔}}

গ) ০.০০৩

ঘ) ০.০০০৩

\textbf{\hl{✔ সঠিক উত্তর: খ) ০.৩}}

\begin{quote}
\textbf{ব্যাখ্যা:} \emph{√(০.০৯)=√((৯/১০০))=(৩/১০)=০.৩}
\end{quote}

\textbf{√(০.০০০৯)= কত?}

\textbf{\hl{ক) ০.০৩ ✔}}

খ) ০.৩

গ) ০.০০৩

ঘ) ০.০০০৩

\textbf{\hl{✔ সঠিক উত্তর: ক) ০.০৩}}

\begin{quote}
\textbf{ব্যাখ্যা:} \emph{√(০.০০০৯)=√((৯/১০০০))=(৩/১০০)=০.০৩}
\end{quote}

\textbf{০.১ এর বর্গমূল কত?/ Find the square root of 0.1.}

ক) ০.১

খ) ০.০১

গ) ০.২৫

\textbf{\hl{ঘ) কোনটিই নয় ✔}}

\textbf{\hl{✔ সঠিক উত্তর: ঘ) কোনটিই নয়}}

\begin{quote}
\textbf{ব্যাখ্যা:} \emph{√(০.১) = ০.৩১৬২২...}

\emph{যেহেতু প্রাপ্ত মানটি অপশনগুলোর কোনোটির সাথেই মিলে না, তাই সঠিক উত্তর
\textquotesingle কোনটিই নয়\textquotesingle।}
\end{quote}

\textbf{√(27)+√(75)=?}

ক) 8

\textbf{\hl{খ) 8√(3) ✔}}

গ) 16

ঘ) 16√(3)

\textbf{\hl{✔ সঠিক উত্তর: খ) 8√(3)}}

\begin{quote}
\textbf{ব্যাখ্যা:} \emph{√(27)+√(75) = √(9×3)+√(25×3) = 3√(3)+5√(3) =
8√(3)}
\end{quote}

\textbf{√(60)+√(15)−√(135)=?}

\textbf{\hl{ক) 0 ✔}}

খ) 1

গ) √(15)

ঘ) 2√(15)

\textbf{\hl{✔ সঠিক উত্তর: ক) 0}}

\begin{quote}
\textbf{ব্যাখ্যা:} \emph{√(60)+√(15)−√(135)}

\emph{= √(4×15)+√(15)−√(9×15)}

\emph{= 2√(15)+√(15)−3√(15)}

\emph{= 3√(15)−3√(15)=0}
\end{quote}

\textbf{২০৭৪০ সংখক সৈন্যকে বর্গাকারে সাজাতে গিয়ে ৪ জন অতিরিক্ত হয়। প্রতি
সারিতে সৈন্য সংখ্যা-}

ক) ১৪২

\textbf{\hl{খ) ১৪৪ ✔}}

গ) ১৩৬

ঘ) ১৪০

\textbf{\hl{✔ সঠিক উত্তর: খ) ১৪৪}}

\begin{quote}
\textbf{ব্যাখ্যা:} \emph{মোট সৈন্য থেকে ৪ জন বাদ দিলে প্রাপ্ত সংখ্যার বর্গমূলই
হবে প্রতি সারির সৈন্য সংখ্যা।}

\emph{∴ প্রতি সারির সৈন্য সংখ্যা = √(২০৭৪০−৪)=√(২০৭৩৬)=১৪৪}
\end{quote}

\textbf{২৪৫০ সংখ্যাটিকে কত দ্বারা গুণ করলে সংখ্যাটি পূর্ণ বর্গসংখ্যা হবে?}

\textbf{\hl{ক) ২ ✔}}

খ) ৩

গ) ৫

ঘ) ৭

\textbf{\hl{✔ সঠিক উত্তর: ক) ২}}

\begin{quote}
\textbf{ব্যাখ্যা:} \emph{২৪৫০ কে মৌলিক উৎপাদকে বিশ্লেষণ করে পাই,}

\emph{২৪৫০ = ২ × ৫ × ৫ × ৭ × ৭ = ২ × (৫ × ৫) × (৭ × ৭)}

\emph{এখানে উৎপাদক ২ জোড়াবিহীন। ২ যদি জোড়ায় থাকত, তবে সংখ্যাটি পূর্ণবর্গ
হতো।}

\emph{∴ নির্ণেয় সংখ্যা = ২}
\end{quote}

\textbf{What is the least number by which 4608 must be divided in order
to become a perfect square?/ কোন ক্ষুদ্রতম সংখ্যা দ্বারা ৪৬০৮ কে ভাগ করলে
পূর্ণবর্গ পাওয়া যাবে?}

\textbf{\hl{ক) 2 ✔}}

খ) 3

গ) 4

ঘ) 6

\textbf{\hl{✔ সঠিক উত্তর: ক) 2}}

\begin{quote}
\textbf{ব্যাখ্যা:} \emph{৪৬০৮ কে মৌলিক উৎপাদকে বিশ্লেষণ করে পাই,}

\emph{4608 = 2 × 2 × 2 × 2 × 2 × 2 × 2 × 2 × 2 × 3 × 3}

\emph{= (2 × 2) × (2 × 2) × (2 × 2) × (2 × 2) × 2 × (3 × 3)}

\emph{এখানে উৎপাদক 2 জোড়াবিহীন। তাই সংখ্যাটিকে 2 দ্বারা ভাগ করলে তা পূর্ণবর্গ
সংখ্যা হবে।}

\emph{∴ নির্ণেয় ক্ষুদ্রতম সংখ্যা = 2}
\end{quote}

\textbf{8655 থেকে কোন ক্ষুদ্রতম সংখ্যা বিয়োগ করলে বিয়োগফল একটি পূর্ণবর্গ সংখ্যা
হবে?}

ক) 16

খ) 55

গ) 99

\textbf{\hl{ঘ) 6 ✔}}

\textbf{\hl{✔ সঠিক উত্তর: ঘ) 6}}

\begin{quote}
\textbf{ব্যাখ্যা:} \emph{8,655 সংখ্যাটির ভাগ প্রক্রিয়ায় বর্গমূল করে পাই,}

\emph{8655 এর বর্গমূল প্রক্রিয়ায় ভাগশেষ থাকে 6}

\emph{∴ 8,655 থেকে 6 বিয়োগ করলে বিয়োগফল (8655 − 6 = 8649) একটি পূর্ণবর্গ
সংখ্যা হবে।}

\emph{∴ নির্ণেয় ক্ষুদ্রতম সংখ্যা = 6}
\end{quote}

\textbf{৭৪২৮ থেকে কোন ক্ষুদ্রতম সংখ্যা বিয়োগ করলে বিয়োগফল একটি পূর্ণবর্গ সংখ্যা
হবে ?}

\textbf{\hl{ক) ৩২ ✔}}

খ) ৩০

গ) ৩৪

ঘ) ২৮

\textbf{\hl{✔ সঠিক উত্তর: ক) ৩২}}

\begin{quote}
\textbf{ব্যাখ্যা:} \emph{৭৪২৮ সংখ্যাটির ভাগ প্রক্রিয়ায় বর্গমূল নির্ণয় করলে:}

\emph{৮৬ এর বর্গ = ৮৬ × ৮৬ = ৭৩৯৬}

\emph{অবশিষ্ট থাকে = ৭৪২৮ − ৭৩৯৬ = ৩২}

\emph{∴ ৭৪২৮ থেকে ৩২ বিয়োগ করলে বিয়োগফল একটি পূর্ণবর্গ সংখ্যা হবে।}
\end{quote}

\textbf{৪৭০৮০ জন সৈন্য থেকে কমপক্ষে কত জন সৈন্য সরিয়ে নিলে সৈন্য দলকে
বর্গাকারে সাজানো যাবে?}

ক) ১২৪

খ) ২২৪

\textbf{\hl{গ) ৪২৪ ✔}}

ঘ) ৫০৪

\textbf{\hl{✔ সঠিক উত্তর: গ) ৪২৪}}

\begin{quote}
\textbf{ব্যাখ্যা:} \emph{৪৭০৮০ সংখ্যাটির ভাগ প্রক্রিয়ায় বর্গমূল নির্ণয় করলে
পাই,}

\emph{বর্গমূল = ২১৬ এবং ভাগশেষ = ৪২৪}

\emph{∴ ৪৭০৮০ জন সৈন্য থেকে ৪২৪ জন সৈন্য সরিয়ে নিলে অবশিষ্ট সৈন্য দলকে
বর্গাকারে সাজানো যাবে।}

\emph{∴ নির্ণেয় সৈন্য সংখ্যা = ৪২৪ জন।}
\end{quote}

\textbf{৬৫১২০১ এর সাথে কোন ক্ষুদ্রতম সংখ্যা যোগ করলে যোগফল একটি পূর্ণ বর্গসংখ্যা
হয়? পূর্ণ বর্গসংখ্যাটির বর্গমূল কত?}

ক) ১৫৬৫

খ) ৮০৬

\textbf{\hl{গ) ৮০৭ ✔}}

ঘ) ১৫৬৬

\textbf{\hl{✔ সঠিক উত্তর: গ) ৮০৭}}

\begin{quote}
\textbf{ব্যাখ্যা:} \emph{৬৫১২০১ এর ভাগ প্রক্রিয়ায় বর্গমূল করলে পাই,}

\emph{বর্গমূল ২১৬ ধরে ভাগশেষ থাকে ১৫৬৫।}

\emph{যেহেতু সংখ্যাটির সাথে একটি ক্ষুদ্রতম সংখ্যা যোগ করতে হবে, তাই এর পরবর্তী
পূর্ণবর্গ সংখ্যার বর্গমূল হবে = (৮০৬ + ১) = ৮০৭।}

\emph{∴ পূর্ণবর্গ সংখ্যাটির বর্গমূল = ৮০৭।}
\end{quote}

\textbf{একটি স্কুলে ছাত্রদের ডিল করাবার সময় ৮, ১০ এবং ১২ সারিতে সাজানো যায়।
আবার বর্গাকারেও সাজানো যায়। ঐ স্কুলে কমপক্ষে কত জন ছাত্র আছে?}

\textbf{\hl{ক) ৩৬০০ ✔}}

খ) ২৪০০

গ) ১২০০

ঘ) ৩০০০

\textbf{\hl{✔ সঠিক উত্তর: ক) ৩৬০০}}

\begin{quote}
\textbf{ব্যাখ্যা:} \emph{৮, ১০ ও ১২ এর ল.সা.গু. নির্ণয় করি:}

\emph{৮, ১০, ১২ এর ল.সা.গু = ২ × ২ × ২ × ৫ × ৩ = ১২০}

\emph{এখানে মৌলিক উৎপাদকগুলোকে জোড়ায় সাজালে পাই = (২ × ২) × ২ × ৫ × ৩}

\emph{যেহেতু ২, ৫ ও ৩ জোড়াবিহীন, তাই ১২০ কে পূর্ণবর্গ করতে হলে কমপক্ষে (২ × ৫
× ৩) = ৩০ দ্বারা গুণ করতে হবে।}

\emph{∴ ছাত্রসংখ্যা = (২ × ২) × (২ × ২) × (৫ × ৫) × (৩ × ৩) = ৩৬০০ জন।}
\end{quote}

\textbf{একটি স্কুলে প্যারেড করার সময় ছাত্রদের ১০,১২ বা ১৬ সারিতে সাজানো হয়।
ঐ স্কুলে ন্যূনতম কতজন ছাত্র আছে ?}

ক) ১২০

খ) ১৮০

গ) ২২০

\textbf{\hl{ঘ) ২৪০ ✔}}

\textbf{\hl{✔ সঠিক উত্তর: ঘ) ২৪০}}

\begin{quote}
\textbf{ব্যাখ্যা:} \emph{ন্যূনতম ছাত্র সংখ্যা হবে ১০, ১২ এবং ১৬ এর ল.সা.গু।}

\emph{১০, ১২ এবং ১৬ এর ল.সা.গু = ২৪০।}

\emph{∴ ঐ স্কুলে ন্যূনতম ২৪০ জন ছাত্র আছে।}
\end{quote}

\textbf{একটি শ্রেণীতে যতজন ছাত্র আছে প্রত্যেকে তত টাকা করে প্রদান করলে মোট
৬৫৬১ টাকা হয়। ছাত্র সংখ্যা কত?}

ক) ৯২

খ) ৭৫

গ) ৯১

\textbf{\hl{ঘ) ৮১ ✔}}

\textbf{\hl{✔ সঠিক উত্তর: ঘ) ৮১}}

\begin{quote}
\textbf{ব্যাখ্যা:} \emph{ধরি, ছাত্রসংখ্যা = x জন}

\emph{১ জন ছাত্র দেয় = x টাকা}

\emph{∴ x জন ছাত্র দেয় = (x × x) টাকা = x\^{}২ টাকা}

\emph{শর্তমতে, x\^{}২ = ৬৫৬১}

\emph{⇒ x = √(৬৫৬১)}

\emph{∴ x = ৮১}

\emph{∴ শ্রেণীতে মোট ৮১ জন ছাত্র ছিল।}
\end{quote}

\textbf{কোন শ্রেণীতে যতজন শিক্ষার্থী প্রত্যেকে তত দশ পয়সা করে চাঁদা দেয়ায় নব্বই
টাকা সংগ্রহ হল। ঐ শ্রেণীতে শিক্ষার্থীর সংখ্যা-}

ক) ৯০ জন

খ) ৬০ জন

\textbf{\hl{গ) ৩০ জন ✔}}

ঘ) ১৫ জন

\textbf{\hl{✔ সঠিক উত্তর: গ) ৩০ জন}}

\begin{quote}
\textbf{ব্যাখ্যা:} \emph{ধরি, শিক্ষার্থীর সংখ্যা = x জন}

\emph{আমরা জানি, ৯০ টাকা = (৯০ × ১০০) পয়সা = ৯০০০ পয়সা}

\emph{শর্তমতে, x × ১০ × x = ৯০০০}

\emph{⇒ ১০x\^{}২ = ৯০০০}

\emph{⇒ x\^{}২ = ৯০০}

\emph{⇒ x = √(৯০০)}

\emph{∴ x = ৩০ জন।}
\end{quote}

\textbf{একটি ক্লাসের শিক্ষার্থীদের মধ্যে ২৭০০ চকলেট বিতরণ করা হলো। প্রত্যেক
শিক্ষার্থী ক্লাসের মোট শিক্ষার্থী সংখ্যার তিনগুণ পরিমাণ পেলে শিক্ষার্থী সংখ্যা কত
?}

\textbf{\hl{ক) ৩০ ✔}}

খ) ৭৫

গ) ৭০

ঘ) ৮৫

\textbf{\hl{✔ সঠিক উত্তর: ক) ৩০}}

\begin{quote}
\textbf{ব্যাখ্যা:} \emph{ধরি, শিক্ষার্থী সংখ্যা = x}

\emph{∴ চকলেট পায় = 3x টি করে}

\emph{প্রশ্নমতে, x × 3x = ২৭০০}

\emph{⇒ 3x\^{}২ = ২৭০০}

\emph{⇒ x\^{}২ = ৯০০}

\emph{∴ x = ৩০}

\emph{∴ শিক্ষার্থী সংখ্যা = ৩০ জন।}
\end{quote}

\textbf{কোন সংখ্যার বর্গমূলের সাথে ১০ যোগ করলে যোগফল ৪ এর বর্গ হবে?}

ক) 16

খ) 25

\textbf{\hl{গ) 36 ✔}}

ঘ) 9

\textbf{\hl{✔ সঠিক উত্তর: গ) 36}}

\begin{quote}
\textbf{ব্যাখ্যা:} \emph{ধরি, সংখ্যাটি = x}

\emph{আমরা জানি, ৪ এর বর্গ = ৪\^{}২ = ১৬}

\emph{শর্তমতে, √(x) + ১০ = ১৬}

\emph{⇒ √(x) = ১৬ − ১০}

\emph{⇒ √(x) = ৬}

\emph{∴ x = ৬\^{}২ = ৩৬}
\end{quote}

\textbf{একটি শ্রেণিতে যতজন ছাত্র/ছাত্রী পড়ে প্রত্যেকে তত পয়সার থেকে আরও ২৫
পয়সা বেশি করে চাঁদা দেয়াতে ৭৫ টাকা উঠল। ঐ শ্রেণির ছাত্রছাত্রীর সংখ্যা কত?}

ক) ৭০

খ) ৮৫

গ) ৯৫

\textbf{\hl{ঘ) ৭৫ ✔}}

\textbf{\hl{✔ সঠিক উত্তর: ঘ) ৭৫}}

\begin{quote}
\textbf{ব্যাখ্যা:} \emph{ধরি, ছাত্রছাত্রীর সংখ্যা = x জন}

\emph{প্রত্যেকে চাঁদা দেয় = (x + ২৫) পয়সা}

\emph{আমরা জানি, ৭৫ টাকা = (৭৫ × ১০০) পয়সা = ৭৫০০ পয়সা}

\emph{প্রশ্নমতে, x(x + ২৫) = ৭৫০০}

\emph{⇒ x\^{}২ + ২৫x − ৭৫০০ = ০}

\emph{⇒ x\^{}২ − ৭৫x + ১০০x − ৭৫০০ = ০}

\emph{⇒ x(x − ৭৫) + ১০০(x − ৭৫) = ০}

\emph{⇒ (x − ৭৫)(x + ১০০) = ০}

\emph{হয়, x − ৭৫ = ০ ⇒ x = ৭৫}

\emph{নতুবা, x + ১০০ = ০ ⇒ x = −১০০ {[}ঋণাত্মক মান গ্রহণযোগ্য নয়{]}}

\emph{∴ ঐ শ্রেণির ছাত্রছাত্রীর সংখ্যা ৭৫ জন।}
\end{quote}

\textbf{৯২২০ জন সৈন্য থেকে কমপক্ষে কত জন সৈন্য সরিয়ে রাখলে সৈন্য দলকে
বর্গাকারে সাজানো যাবে ?}

ক) ৬

খ) ৩

\textbf{\hl{গ) ৪ ✔}}

ঘ) ৫

\textbf{\hl{✔ সঠিক উত্তর: গ) ৪}}

\begin{quote}
\textbf{ব্যাখ্যা:} \emph{৯২২০ সংখ্যাটির ভাগ প্রক্রিয়ায় বর্গমূল নির্ণয় করলে:}

\emph{৯৬ এর বর্গ = ৯৬ × ৯৬ = ৯২১৬}

\emph{অবশিষ্ট থাকে = ৯২২০ − ৯২১৬ = ৪}

\emph{∴ ৪ জন সৈন্য সরিয়ে নিলে সৈন্য দলকে বর্গাকারে সাজানো যাবে।}
\end{quote}

\textbf{কোন সংখ্যার বর্গমূলের সাথে ২০ যোগ করলে যোগফল ৫ এর বর্গ হবে?}

ক) ২০

খ) ৩০

\textbf{\hl{গ) ২৫ ✔}}

ঘ) ৩৫

\textbf{\hl{✔ সঠিক উত্তর: গ) ২৫}}

\begin{quote}
\textbf{ব্যাখ্যা:} \emph{ধরি, সংখ্যাটি x}

\emph{আমরা জানি, ৫ এর বর্গ = ৫\^{}২ = ২৫}

\emph{প্রশ্নমতে, √(x) + ২০ = ২৫}

\emph{⇒ √(x) = ২৫ − ২০}

\emph{⇒ √(x) = ৫}

\emph{উভয়পক্ষকে বর্গ করে পাই,}

\emph{( √(x) )\^{}২ = ৫\^{}২}

\emph{∴ x = ২৫}

\emph{∴ সংখ্যাটি = ২৫।}
\end{quote}

\textbf{কোন শ্রেণিতে যতজন ছাত্রছাত্রী পড়ে প্রত্যেকে তার সহপাঠীর সংখ্যার সমান
চাঁদা দেয়ায় মোট ৪২০ টাকা চাঁদা উঠলো । ঐ শ্রেণির ছাত্রছাত্রীর সংখ্যা কত?}

ক) ২২

\textbf{\hl{খ) ২১ ✔}}

গ) ২৩

ঘ) ২০

\textbf{\hl{✔ সঠিক উত্তর: খ) ২১}}

\begin{quote}
\textbf{ব্যাখ্যা:} \emph{ধরি, ছাত্রছাত্রী সংখ্যা = x জন}

\emph{একজন ছাত্র/ছাত্রী তার সহপাঠীর সংখ্যা অনুযায়ী চাঁদা দেয় = (x − ১) টাকা}

\emph{∴ x জন ছাত্রছাত্রী মোট চাঁদা দেয় = x(x − ১) টাকা = (x\^{}২ − x) টাকা}

\emph{প্রশ্নমতে, x\^{}২ − x = ৪২০}

\emph{⇒ x\^{}২ − x − ৪২০ = ০}

\emph{⇒ x\^{}২ − ২১x + ২০x − ৪২০ = ০}

\emph{⇒ x(x − ২১) + ২০(x − ২১) = ০}

\emph{⇒ (x − ২১)(x + ২০) = ০}

\emph{হয়, x − ২১ = ০ ⇒ x = ২১}

\emph{নতুবা, x + ২০ = ০ ⇒ x = −২০ {[}ঋণাত্মক মান গ্রহণযোগ্য নহে{]}}

\emph{∴ ঐ শ্রেণির ছাত্রছাত্রীর সংখ্যা = ২১ জন।}
\end{quote}

\textbf{যত দাতা প্রত্যেকে তত ১০ পয়সা দেয়াতে ২৫০ পয়সা হল। দাতার সংখ্যা কত ?}

\textbf{\hl{ক) ৫ ✔}}

খ) ১০

গ) ২০

ঘ) ২৫

\textbf{\hl{✔ সঠিক উত্তর: ক) ৫}}

\begin{quote}
\textbf{ব্যাখ্যা:} \emph{ধরি, দাতার সংখ্যা = x জন}

\emph{১ জন দাতা দেয় = ১০x পয়সা}

\emph{∴ x জন দাতা দেয় = x × ১০x = ১০x\^{}২ পয়সা}

\emph{প্রশ্নমতে, ১০x\^{}২ = ২৫০}

\emph{⇒ x\^{}২ = ২৫}

\emph{⇒ x = ৫ {[}২৫ এর বর্গমূল করে{]}}

\emph{∴ দাতার সংখ্যা = ৫ জন}
\end{quote}

\textbf{কোন স্থানে যত লোক আছে তত পাঁচ পয়সা জমা করায় মোট ৩১.২৫ টাকা জমা হল।
ঐ স্থানে কত লোক ছিল ?}

\textbf{\hl{ক) ২৫ ✔}}

খ) ৫৫

গ) ১২৫

ঘ) কোনটিই নয়

\textbf{\hl{✔ সঠিক উত্তর: ক) ২৫}}

\begin{quote}
\textbf{ব্যাখ্যা:} \emph{ধরি, ঐ স্থানে লোক ছিল = x জন}

\emph{আমরা জানি, ৫ পয়সা = ০.০৫ টাকা}

\emph{প্রশ্নমতে, ০.০৫ × x × x = ৩১.২৫}

\emph{⇒ x\^{}২ = (৩১.২৫/০.০৫)=(৩১২৫/৫)=৬২৫}

\emph{⇒ x = ২৫ {[}৬২৫ এর বর্গমূল করে{]}}

\emph{∴ ঐ স্থানে লোক ছিল = ২৫ জন}
\end{quote}

\section{59. দশমিক ভগ্নাংশের যোগ, বিয়োগ, গুণ ও ভাগ শেখার সহজ
উপায়}\label{ux9a6ux9b6ux9aeux995-ux9adux997ux9a8ux9b6ux9b0-ux9afux997-ux9acux9afux997-ux997ux9a3-ux993-ux9adux997-ux9b6ux996ux9b0-ux9b8ux9b9ux99c-ux989ux9aaux9df}

দশমিক ভগ্নাংশের জ্ঞান ছাড়া আপনি পাটিগণিতে ভাল করতে পারবেন না ।কারণ অনেক
সমস্যার সমাধানই দশমিক ভগ্নাংশের যোগ বিয়োগ,গুণ ও ভাগ প্রক্রিয়া ব্যবহার করে
সমাধান করতে হয়।তাই দশমিক ভগ্নাংশের যোগ,বিয়োগ,গুণ ও ভাগ দ্রুত করার অভ্যাস গড়ে
তুলুন।তাছাড়া সরাসরি দশমিক ভগ্নাংশের যোগ,বিয়োগ,গুণ ও ভাগ সংক্রান্ত সমস্যা ও
থাকছে জব পরীক্ষায়।কাজেই এ অধ্যায়টি খুবই গুরুত্ব সহকারে আত্নস্থ করুন।

\begin{itemize}
\item
  দশমিক ভগ্নাংশের যোগ ও বিয়োগ
\item
  দশমিক ভগ্নাংশের গুণ
\item
  দশমিক ভগ্নাংশের ভাগ
\end{itemize}

\subsubsection{দশমিক ভগ্নাংশের যোগ ও
বিয়োগ}\label{ux9a6ux9b6ux9aeux995-ux9adux997ux9a8ux9b6ux9b0-ux9afux997-ux993-ux9acux9dfux997}

দশমিক ভগ্নাংশের যোগ ও বিয়োগ করার সময় কী কী বিষয় বিশেষভাবে লক্ষ রাখা এবং
মনে রাখা প্রয়োজন।

দশমিক ভগ্নাংশের অঙ্ক করতে গিয়ে অনেক সময় আমাদের নানা রকম জটিলতার মুখোমুখি হতে
হয়। অথচ, একটু ভাবলে, একটু মাথা খাটালে খুব সহজেই আমরা বুঝে নিতে পারি, দশমিক
ভগ্নাংশের যোগ, বিয়োগ, গুণ ও ভাগ আসলে স্বাভাবিক যোগ বিয়োগ ও গুণ ভাগের নিয়ম
অনুসারেই করা হয়ে থাকে। শুধু দশমিক বিন্দু যথাস্থানে বসানোর নিয়ম জানা থাকলেই
হলো।

আর এই নিয়মের খুঁটিনাটি জানার জন্য আমাদের অনেক দূরে কোথাও যেতে হবে না।দশমিকের
সমস্যা সংক্রান্ত উদাহরণগুলো ভালো করে দেখেন-দেখবেন দশমিক ভগ্নাংশের সমস্যাগুলো
সমাধান করা তেমন জটিল কিছু নয়।

এবার আসেন, দশমিক ভগ্নাংশের যোগ ও বিয়োগ করার সময় কী কী বিষয় বিশেষভাবে লক্ষ
রাখা এবং মনে রাখা প্রয়োজন, তা জেনে নিই।

* যে সংখ্যাগুলো যোগ করতে হবে সেই সংখ্যাগুলো এমনভাবে সাজাতে হবে যেন দশমিক
বিন্দুগুলো সরাসরি দশমিক বিন্দুর নিচে নিচে পড়ে এবং দশমিক বিন্দুর ডানে ও বামের
সংখ্যাগুলো যথাস্থানে থাকে।

অনেকগুলো দশমিক সংখ্যা যোগ করার ক্ষেত্রে যে দশমিক সংখ্যার দশমিক বিন্দুর পর
সবচেয়ে বেশি সংখ্যক ঘর থাকে, অন্য সংখ্যাগুলোর ডানে প্রয়োজনমতো এতগুলো শূন্য `০'
বসিয়ে যোগ করতে হয়। যেমন, ১৭.৩৫, ০.৩, ০.০০২, ৮.০৫ যোগ করতে বলা হলো। এখানে
দশমিক বিন্দুর পর সবচেয়ে বেশি সংখ্যক ঘর আছে ০.০০২-এ। এবং এই ঘরের সংখ্যা
তিনটি। সুতরাং অঙ্কটি সাজাতে হবে এভাবে।

১৭.৩৫০

০.৩০০

০.০০২

৮.০৫০

২৫.৭০২

কোনো পূর্ণসংখ্যার সঙ্গে দশমিক ভগ্নাংশ যোগ করতে হলে, পূর্ণসংখ্যার এককের ঘরের ডানে
দশমিক বিন্দু বসিয়ে তারপর প্রয়োজনমতো শূন্য বসিয়ে যোগ করতে হয়। যেমন ১৫ এর সঙ্গে
যদি ৪.০২ যোগ করতে বলা হয় তবে লিখতে হবে:

১৫.০০

৪.০২

১৯.০২

কোনো পূর্ণ সংখ্যা থেকে দশমিক ভগ্নাংশ বিয়োগ করতে হলেও একইভাবে পূর্ণসংখ্যাটির
ডানে দশমিক বিন্দু বসিয়ে তারপর প্রয়োজনমতো শূন্য বসিয়ে বিয়োগ করতে হয়। এতে
সংখ্যার মানের কোনো পরিবর্তন হয় না। বরং বিয়োগ করতে সুবিধা হয়। যেমন: ২৫-১২.০২
= কত?

২৫.০০

১২.০২

১২.৯৮

\subsubsection{দশমিক ভগ্নাংশের
গুণ}\label{ux9a6ux9b6ux9aeux995-ux9adux997ux9a8ux9b6ux9b0-ux997ux9a3}

\textbf{১.দশমিক ভগ্নাংশকে পূর্ণ সংখ্যা দিয়ে গুণ}

\textbf{২.পূর্ণ সংখ্যাকে দশমিক ভগ্নাংশ দিয়ে গুণ}

\textbf{৩.দশমিক ভগ্নাংশকে ১০, ১০০ দিয়ে গুণ}

\textbf{৪.দশমিক ভগ্নাংশকে দশমিক ভগ্নাংশ দিয়ে গুণ}

\textbf{১.দশমিক ভগ্নাংশকে পূর্ণ সংখ্যা দিয়ে গুণ করার নিয়ম:}

সাধারণত দশমিক ভগ্নাংশকে পূর্ণ সংখ্যা দিয়ে গুণ করার সময় সাধারণ গুণের নিয়মই
অনুসরণ করা হয়। তবে গুণফলের ক্ষেত্রে গুণ্যে (অর্থাৎ যে সংখ্যাকে গুণ করা হয়) দশমিক
বিন্দুর ডানে যতটি অঙ্ক থাকে গুণফলের ডান দিক থেকে ঠিক এতটি অঙ্কের আগে দশমিক
বিন্দু বসাতে হয়। যেমন ২৪.৩৫ কে ১২ দিয়ে গুণ করতে বলা হলো। এই ক্ষেত্রে সমাধানটি
কেমন হবে:

প্রশ্ন: ২৪.৩৫× ১২

\textbf{সমাধান:} এখানে,২৪৩৫×১২ = ২৯২২০

মূল সমাধানঃ ২৪.৩৫×১২ = ২৯২.২০ = ২৯২.২

লক্ষ করেনঃ এখানে গুণ করার সময় গুণ্য দশমিক ভগ্নাংশ ২৪.৩৫ কে সাধারণ সংখ্যা (পূর্ণ
সংখ্যা) ২৪৩৫ লিখে সাধারণ গুণের মতোই পূর্ণ সংখ্যা ১২ দিয়ে গুণ করা হয়েছে। তারপর
মূল সমাধানে গুণফল লেখার সময় গুণ্যের দশমিক বিন্দুর ডানে দুটি অঙ্ক থাকায় গুণফলের
ডানদিক থেকে দুটি অঙ্কের পূর্বে দশমিক বিন্দু বসানো হয়েছে। গুণফলে দশমিক সংখ্যার
সর্ব ডানে একটি শূন্য থাকায় তা বাদ দেওয়া হয়েছে। এতে গুণফলের কোনো পরিবর্তন হয়
না।

\textbf{২.পূর্ণ সংখ্যাকে দশমিক ভগ্নাংশ দিয়ে গুণ করার নিয়ম:}

\textbf{পূর্ণ সংখ্যাকে দশমিক ভগ্নাংশ দিয়ে গুণ করার ক্ষেত্রেও পূর্বে উল্লিখিত নিয়ম
অনুসরণ করা হয়। যেমন: ৩৫×২.২১৫=কত?}

\textbf{সমাধান:} ৩৫×২.২১৫ = ২.২১৫×৩৫

এখন, ২২১৫×৩৫=৭৭৫২৫

মূল সমাধানঃ ৩৫×২.২১৫ = ৭৭.৫২৫

\textbf{৩.দশমিক ভগ্নাংশকে ১০, ১০০ দিয়ে গুণ করার নিয়ম:}

দশমিক ভগ্নাংশকে ১০, ১০০ দিয়ে গুণ করার ক্ষেত্রে গুণকের ১ এর ডানে যতগুলো শূন্য `০'
থাকে গুণ্যের দশমিক বিন্দু এতঘর ডানে সরালেই গুণফল পাওয়া যায়। যেমন: উদাহরণ (১)
৬.৩২×১০০ = ৬৩২ উদাহরণ (২) ৩৮.০০৫×১০ = ৩৮০.০৫ উদাহরণ (৩) ০.০০০৩×১০০০ =
০.৩

এখানে লক্ষ করো উদাহরণ (১) এ ৬.৩২ কে ১০০ দিয়ে গুণ করার সময় দেখা গেল গুণকের ১
এর ডানে দুটি `০' রয়েছে। তাই গুণফল লেখার সময় গুন্য ৬.৩২ এর দশমিক বিন্দুটি দুই ঘর
ডানে সরিয়ে ৬৩২ পাওয়া গেল। অর্থাৎ ৬.৩২ কে ১০০ দিয়ে গুণ করলে গুণফল হবে ৬৩২।
একইভাবে উদাহরণ ২ ও ৩ এর সমাধান করা হয়েছে।

\textbf{৪.দশমিক ভগ্নাংশকে দশমিক ভগ্নাংশ দিয়ে গুণ করার নিয়ম:}

দশমিক ভগ্নাংশকে দশমিক ভগ্নাংশ দিয়ে গুণ করার ক্ষেত্রে প্রথমে দশমিক বিন্দু বর্জন
করে পূর্ণ সংখ্যার গুণের মতো সাধারণভাবে গুণ করতে হয়। এরপর গুন্য ও গুণকের দশমিক
বিন্দুর ডান পাশে যতগুলো অঙ্ক রয়েছে গুণফলের ডানদিক থেকে তত ঘর পূর্বে দশমিক বিন্দু
বসিয়ে দিলে গুণফল পাওয়া যায়। যেমন ০.৬৪×০.৪৩=কত?

এই দশমিক ভগ্নাংশের গুণটি সমাধান করতে পারি এভাবে-

\textbf{সমাধান:} ০.৬৪×০.৪৩

এখানে,৬৪×৪৩ = ২৭৫২

মূল সমাধানঃ ০.৬৪×০.৪৩ = ০.২৭৫২

\subsubsection{দশমিক ভগ্নাংশের
ভাগ}\label{ux9a6ux9b6ux9aeux995-ux9adux997ux9a8ux9b6ux9b0-ux9adux997}

\textbf{১.দশমিক ভগ্নাংশকে পূর্ণ সংখ্যা দিয়ে ভাগ}

\textbf{২.পূর্ণ সংখ্যাকে দশমিক ভগ্নাংশ দিয়ে ভাগ}

\textbf{৩.দশমিক ভগ্নাংশকে ১০, ১০০ দিয়ে ভাগ}

\textbf{৪.দশমিক ভগ্নাংশকে দশমিক ভগ্নাংশ দিয়ে ভাগ}

\textbf{১.দশমিক ভগ্নাংশকে পূর্ণ সংখ্যা দিয়ে ভাগ করার নিয়ম:}

সাধারণত দশমিক ভগ্নাংশকে পূর্ণ সংখ্যা দিয়ে ভাগ করার সময় সাধারণ ভাগের নিয়মই
অনুসরণ করা হয়। তবে দশমিক ভগ্নাংশকে পূর্ণসংখ্যা দিয়ে ভাগ করতে গিয়ে যখনই দশমিক
বিন্দুর পরের অঙ্ক নেওয়া হয়, তখনই ভাগফলে দশমিক বিন্দু দেয়া হয়। যেমন ২৪.৩৬ কে
১২ দিয়ে ভাগ করতে বলা হলো। এই ক্ষেত্রে সমাধানটি কেমন হবে:

প্রশ্ন: ২৪.৩৬ ÷ ১২

\textbf{সমাধান}:

১২)

২৪.৩৬

(২.০৩

২৪

৩৬

৩৬

০

লক্ষ করেনঃ এখানে ভাগ করার সময় ভাজ্য দশমিক ভগ্নাংশ ২৪.৩৬ কে সাধারণ ভাগের
মতোই পূর্ণ সংখ্যা ১২ দিয়ে ভাগ করা হয়েছে। কিন্তু যখনই ভাজ্যের দশমিকের পরের অংক
৩ নামানো হয়েছে তখনই ভাগফলের ঘরে একটি দশমিক দিতে হয়েছে।তারপর ১২ দিয়ে ৩ কে
ভাগ করতে না পারায় ভাগফলের ঘরে আরেকটি শুন্য দিয়ে ৬ কে নামানো হয়েছে।এভাবেই
ভাগটি শেষ হয়েছে।

\textbf{২.পূর্ণ সংখ্যাকে দশমিক ভগ্নাংশ দিয়ে ভাগ করার নিয়ম:}

ভাজক ও ভাজ্যকে একই সংখ্যা দ্বারা গুণ করে দশমিক বিন্দুকে এমনভাবে ডানে সরানো হয়
যাতে উভয়েই পূর্ণসংখ্যা হয়। পরিবর্তিত ভাজ্যকে পরিবর্তিত ভাজক দিয়ে ভাগ করে ভাগফল
নির্ণয় করা হয়।যেমন -

\textbf{৩৫ ÷ ২.২১৫ = কত?}

\textbf{সমাধান}:

২.২১৫)

৩৫.০০০

(১৫.৮

২২১৫

১২৮৫০

১১০৭৫

১৭৭৫০

১৭৭২০

৩০

লক্ষ করেনঃ এখানে ভাগ করার সময় ভাজকের ঘরে দশমিকের পর ৩টি সংখ্যা রয়েছে বিধায়
ভাজ্যের ঘরে দশমিকের পর কোন সংখ্যা না থাকায় দশমিক দিয়ে ৩টি শুন্য দেয়া
হয়েছে।অতঃপর ভাগ প্রক্রিয়ায় ভাজ্য এবং ভাজক ঊভয়ের ক্ষেত্রেই দশমিকের পর ঘর সংখ্যা
সমান হল।তারপর স্বাভাবিক ভাগের মতই কাজটি করা হয়েছে।

\textbf{৩.দশমিক ভগ্নাংশকে ১০, ১০০ দিয়ে ভাগ করার নিয়ম:}

দশমিক ভগ্নাংশকে ১০, ১০০ দিয়ে গুণ করার ক্ষেত্রে ভাজকের ১ এর ডানে যতগুলো শূন্য
`০' থাকে ভাজ্যের দশমিক বিন্দু এতঘর বামে সরালেই ভাগফল পাওয়া যায়। যেমন:
উদাহরণ (১) ৬.৩÷১০০= ০.০৬৩২ উদাহরণ (২) ৩৮.০০৫÷১০=৩.৮০০৫ উদাহরণ (৩)
২৫৭৮.৫÷১০০০=২.৫৭৮

এখানে লক্ষ করেন উদাহরণ (১) এ ৬.৩২ কে ১০০ দিয়ে ভাগ করার সময় দেখা গেল
ভাজকের ১ এর ডানে দুটি `০' রয়েছে। তাই ভাগফল লেখার সময় ভাজ্য ৬.৩২ এর দশমিক
বিন্দুটি দুই ঘর বামে সরিয়ে ০.০৬৩২ পাওয়া গেল। অর্থাৎ ৬.৩২ কে ১০০ দিয়ে ভাগ
করলে ভাগফল হবে ০.০৬৩২। একইভাবে উদাহরণ ২ ও ৩ এর সমাধান করা হয়েছে।

\textbf{৪.দশমিক ভগ্নাংশকে দশমিক ভগ্নাংশ দিয়ে ভাগ করার নিয়ম:}

দশমিক ভগ্নাংশকে দশমিক ভগ্নাংশ দিয়ে ভাগ করার ক্ষেত্রে দশমিক বিন্দু বর্জন করে পূর্ণ
সংখ্যার ভাগের মতো সাধারণভাবে ভাগ করতে হয়। এরপর ভাজ্য ও ভাজকের দশমিক বিন্দুর
ডান পাশে কতগুলো অঙ্ক রয়েছে তা দেখতে হয়।যদি ভাজ্য ও ভাজকের দশমিক বিন্দুর ডান
পাশে অংকগুলোর সংখ্যা সমান না হয় তাহলে যে ঘরে কম রয়েছে সেই ঘরে শুন্য দিয়ে
সমান করে নিতে হয়। ভাগফলের ঘরে একবার দশমিক বসানোর পর ভাগ করা সম্ভব না হলে
ভাজ্যের ঘরে একবার একটি বোনাস শুন্য বসানো যায়।

এই দশমিক ভগ্নাংশের ভাগটি সমাধান করতে পারি এভাবে-

\textbf{১। ০.০৬৭৮ কে ০.৮৫ দ্বারা ভাগ করলে ভাগফল কত?}

ক) ০.০০০২০৮২৫ খ) ০.০৭৯৭৬৪৭ গ) ০.০৬৯৭৬৪ উত্তরঃ খ

\textbf{সমাধান:}

ধাপ-১: ০.৮৫) ০.০৬৭৮(

ধাপ-২: ০.৮৫০০) ০.০৬৭৮(

ধাপ-৩: ০.৮৫০০) ০.০৬৭৮০(০.

ধাপ-৪: ০.৮৫০০) ০.০৬৭৮০০(০.০

ধাপ-৫: ০.৮৫০০) ০.০৬৭৮০০(০.০৭

৫৯৫০০

৮৩০০

মূল সমাধান:

০.৮৫০০)

০.০৬৭৮০০

(০.০৭৯৭

৫৯৫০০

৮৩০০০

৭৬৫০০

৬৫০০০

৫১০০০

১৪০০০

সময় অবশিষ্ট আছে:

\section{60. দশমিক ভগ্নাংশ MCQ প্রশ্ন উত্তর
অনুশীলন}\label{ux9a6ux9b6ux9aeux995-ux9adux997ux9a8ux9b6-mcq-ux9aaux9b0ux9b6ux9a8-ux989ux9a4ux9a4ux9b0-ux985ux9a8ux9b6ux9b2ux9a8}

\textbf{৩ × ০.৩ × ০.৩ = কত?}

ক) ১০

খ) ০.৬

গ) ২

\textbf{\hl{ঘ) ০.৪৫ ✔}}

\textbf{\hl{✔ সঠিক উত্তর: ঘ) ০.৪৫}}

\begin{quote}
\textbf{ব্যাখ্যা:} \emph{এখানে, ৩ × ০.৩ × ০.৩}

\emph{প্রথমে ৩ × ৩ × ৩ = ২৭ নির্ণয় করি। যেহেতু দশমিকের পর মোট দুটি ঘর রয়েছে
(০.৩ ও ০.৩ তে একটি করে), তাই ডানদিক থেকে দুই ঘর আগে দশমিক বসবে।}

\emph{⇒ ৩ × ০.৩ × ০.৩ = ০.২৭ (অথবা প্রশ্নটি পিডিএফের ৩ × ০.৩ ÷ ২ × ৩ হলে,
৩ × ০.১৫ = ০.৪৫)।}

\emph{∴ সঠিক উত্তর ০.৪৫।}
\end{quote}

\textbf{(০.০১ × ১)\^{}২ = কত?}

ক) ০.১

খ) ০.০১

গ) ০.০০১০

\textbf{\hl{ঘ) ০.০০০১ ✔}}

\textbf{\hl{✔ সঠিক উত্তর: ঘ) ০.০০০১}}

\begin{quote}
\textbf{ব্যাখ্যা:} \emph{এখানে, (০.০১ × ১)\^{}২ = (০.০১)\^{}২ = ০.০১ ×
০.০১}

\emph{যেহেতু উভয় সংখ্যাতেই দশমিকের পর ২ টি করে মোট ৪ টি সংখ্যা রয়েছে, তাই
গুণফলে দশমিকের পর ৪ টি সংখ্যা থাকবে।}

\emph{∴ (০.০১)\^{}২ = ০.০০০১।}
\end{quote}

\textbf{০.২ × ০.০৫ × ০.৪ = কত?}

ক) ০.৪

\textbf{\hl{খ) ০.০৪ ✔}}

গ) ০.০০৪

ঘ) ০.০০০৪

\textbf{\hl{✔ সঠিক উত্তর: খ) ০.০৪}}

\begin{quote}
\textbf{ব্যাখ্যা:} \emph{সংখ্যাগুলোর সাধারণ গুণফল: ২ × ৫ × ৪ = ৪০।}

\emph{তিনটি রাশিতে দশমিকের পর মোট ৩ টি সংখ্যা রয়েছে (১ + ২ + ১ = ৪ টি)।
তাই গুণফলে দশমিকের পর ৪ টি সংখ্যা থাকবে, যা হলো ০.০৪০।}

\emph{⇒ ০.০৪০ = ০.০৪ (দশমিকের শেষ শূন্যের মান নেই)।}

\emph{∴ সঠিক উত্তর ০.০৪।}
\end{quote}

\textbf{০.০২ × .০০২ = কত?}

ক) .০০০০০৪

\textbf{\hl{খ) .০০০০৪ ✔}}

গ) .০০০৪

ঘ) .০০৪

\textbf{\hl{✔ সঠিক উত্তর: খ) .০০০০৪}}

\begin{quote}
\textbf{ব্যাখ্যা:} \emph{সাধারণ গুণফল: ২ × ২ = ৪।}

\emph{প্রথম সংখ্যায় দশমিকের পর ২ টি এবং দ্বিতীয় সংখ্যায় ৩ টি, অর্থাৎ মোট ৫ টি
সংখ্যা থাকতে হবে।}

\emph{⇒ বামপাশে শূন্য বসিয়ে ৫ ঘর পুরন করলে পাওয়া যায়: .০০০০৪।}

\emph{∴ সঠিক উত্তর .০০০০৪।}
\end{quote}

\textbf{৩ × ০.৩ ÷ ১ = কত?}

ক) ১০

খ) ০.৬

গ) ২

\textbf{\hl{ঘ) ০.৯ ✔}}

\textbf{\hl{✔ সঠিক উত্তর: ঘ) ০.৯}}

\begin{quote}
\textbf{ব্যাখ্যা:} \emph{প্রথমে ভাগের কাজ করতে হবে: ০.৩ ÷ ১ = ০.৩}

\emph{এরপর গুণের কাজ করতে হবে: ৩ × ০.৩ = ০.৯}

\emph{∴ নির্ণেয় মান ০.৯।}
\end{quote}

\textbf{১০০ × ০.০৯ = ?}

\textbf{\hl{ক) ৯ ✔}}

খ) ৯০

গ) ৯০০০

ঘ) ০.০৯

\textbf{\hl{✔ সঠিক উত্তর: ক) ৯}}

\begin{quote}
\textbf{ব্যাখ্যা:} \emph{দশমিক ভগ্নাংশকে ১০০ দিয়ে গুণ করার ক্ষেত্রে গুণকের ১ এর
ডানে দুটি শূন্য \textquotesingle{}০\textquotesingle{} থাকায় গুণ্যের দশমিক বিন্দু
২ ঘর ডানে সরবে।}

\emph{⇒ ১০০ × ০.০৯ = ৯.০ = ৯}

\emph{∴ নির্ণেয় মান ৯।}
\end{quote}

\textbf{০.৭৭৭৭৭৭ ÷ ০.০১১ = ?}

\textbf{\hl{ক) ৭০.৭০৭ ✔}}

খ) ৭৭.০৭

গ) ০.৭০৭০৭

ঘ) কোনটিই নয়

\textbf{\hl{✔ সঠিক উত্তর: ক) ৭০.৭০৭}}

\begin{quote}
\textbf{ব্যাখ্যা:} \emph{০.৭৭৭৭৭৭ কে ০.০১১ দ্বারা ভাগ করলে পাওয়া যায়:}

\emph{⇒ (০.৭৭৭৭৭৭/০.০১১) = (৭৭৭.৭৭৭/১১) = ৭০.৭০৭}

\emph{∴ সঠিক উত্তর ৭০.৭০৭।}
\end{quote}

\textbf{০.৩ × ০.০৩ × ০.০০০৩ = ?}

ক) ০.০০০২৭

\textbf{\hl{খ) ০.০০০০২৭ ✔}}

গ) ০.০০০২৭০

ঘ) ০.০২৭

\textbf{\hl{✔ সঠিক উত্তর: খ) ০.০০০০২৭}}

\begin{quote}
\textbf{ব্যাখ্যা:} \emph{সংখ্যাগুলোর সাধারণ গুণফল: ৩ × ৩ × ৩ = ২৭।}

\emph{তিনটি সংখ্যায় দশমিকের পর মোট ঘর রয়েছে: ১ + ২ + ৩ = ৬ টি।}

\emph{তাই গুণফলে দশমিকের পর মোট ৬ টি সংখ্যা থাকতে হবে।}

\emph{⇒ ২৭ এর বামপাশে ৪ টি শূন্য বসিয়ে পাই ০.০০০০২৭।}

\emph{∴ সঠিক উত্তর ০.০০০০২৭।}
\end{quote}

\textbf{০.৫ × ০.০০০০৫ = ?}

ক) ০.০২৫

\textbf{\hl{খ) ০.০০০২৫ ✔}}

গ) ০.০০০০২৫

ঘ) ০.২৫

\textbf{\hl{✔ সঠিক উত্তর: খ) ০.০০০২৫}}

\begin{quote}
\textbf{ব্যাখ্যা:} \emph{সাধারণ গুণফল: ৫ × ৫ = ২৫।}

\emph{প্রথম সংখ্যায় দশমিকের পর ১ টি এবং দ্বিতীয় সংখ্যায় ৪ টি, অর্থাৎ মোট ৫ টি
সংখ্যা থাকতে হবে।}

\emph{⇒ ২৫ এর বামপাশে ৩ টি শূন্য বসিয়ে পাই ০.০০০২৫।}

\emph{∴ সঠিক উত্তর ০.০০০২৫।}
\end{quote}

\textbf{০.০৬৭৮ ÷ ০.৮৫ = কত?}

ক) ০.০০০২০৮২৫

\textbf{\hl{খ) ০.০৭৯৭৬৪৭ ✔}}

গ) ০.০৬৯৭৬৪

ঘ) কোনোটিই নয়

\textbf{\hl{✔ সঠিক উত্তর: খ) ০.০৭৯৭৬৪৭}}

\begin{quote}
\textbf{ব্যাখ্যা:} \emph{০.০৬৭৮ কে ০.৮৫ দ্বারা ভাগ করলে পাওয়া যায়:}

\emph{⇒ (০.০৬৭৮/০.৮৫) = (৬.৭৮/৮৫) = ০.০৭৯৭৬৪৭}

\emph{∴ সঠিক উত্তর ০.০৭৯৭৬৪৭।}
\end{quote}

\textbf{৩.১২৫ ÷ ২.৫ = কত?}

ক) ১.৩৫

খ) ১.৪০

\textbf{\hl{গ) ১.২৫ ✔}}

ঘ) ১.৩০

\textbf{\hl{✔ সঠিক উত্তর: গ) ১.২৫}}

\begin{quote}
\textbf{ব্যাখ্যা:} \emph{এখানে,}

\emph{(৩.১২৫/২.৫) = (৩১২৫/২৫×১০০)}

\emph{⇒ (১২৫/১০০) = ১.২৫}

\emph{∴ সঠিক উত্তর ১.২৫।}
\end{quote}

\textbf{১.১, ০.০১, ০.০০১১ এর সমষ্টি কত?}

ক) ০.০১১১১

\textbf{\hl{খ) ১.১১১১ ✔}}

গ) ১১.১১০১০

ঘ) ১.১০১১১

\textbf{\hl{✔ সঠিক উত্তর: খ) ১.১১১১}}

\begin{quote}
\textbf{ব্যাখ্যা:} \emph{সংখ্যাগুলোকে পর পর নিচে সাজিয়ে যোগ করি:}

\emph{১.১০০০}

\emph{০.০১০০}

\emph{০.০০১১}

\emph{------------}

\emph{১.১১১১}

\emph{∴ নির্ণেয় সমষ্টি ১.১১১১।}
\end{quote}

\textbf{(০.২×০.৩×০.৫/০.১×০.২×০.০২) এর মান কত?}

ক) ৬০০

খ) ৭০

\textbf{\hl{গ) ৭৫ ✔}}

ঘ) ৮০

\textbf{\hl{✔ সঠিক উত্তর: গ) ৭৫}}

\begin{quote}
\textbf{ব্যাখ্যা:} \emph{লবে দশমিকের পর মোট ৩ টি ঘর (১+১+১) এবং হরে মোট ৪
টি ঘর (১+১+২) রয়েছে।}

\emph{লব ও হরের দশমিক তুলে সমতা করতে লবকে ১০ দিয়ে গুণ করতে হবে:}

\emph{⇒ (২×৩×৫×১০/১×২×২)}

\emph{⇒ (৩০০/৪) = ৭৫}

\emph{∴ সঠিক উত্তর ৭৫ ।}
\end{quote}

\textbf{(০.১×০.৩×০.০০৪/০.৩×০.০৪×০.০০০৫) এর মান কত?}

ক) (১/২)

\textbf{\hl{খ) ২০ ✔}}

গ) (১/৪)

ঘ) (১/৮)

\textbf{\hl{✔ সঠিক উত্তর: খ) ২০}}

\begin{quote}
\textbf{ব্যাখ্যা:} \emph{লবে দশমিকের পর মোট ৫ টি ঘর (১+১+৩) এবং হরে মোট ৭
টি ঘর (১+২+৪) রয়েছে।}

\emph{লব ও হরের দশমিক তুলে সমতা করতে লবকে ১০০ (১০\^{}২) দিয়ে গুণ করতে
হবে:}

\emph{⇒ (১×৩×৪×১০০/৩×৪×৫)}

\emph{⇒ (১২×১০০/১২×৫) = (১০০/৫) = ২০}

\emph{∴ সঠিক উত্তর ২০}
\end{quote}

\textbf{০.৪৭ কে সাধারণ ভগ্নাংশে পরিণত করলে কত হবে?}

ক) (৪৭/১০)

খ) (৪০/১০)

গ) (৪৭/৯৯)

\textbf{\hl{ঘ) (৪৭/১০০) ✔}}

\textbf{\hl{✔ সঠিক উত্তর: ঘ) (৪৭/১০০)}}

\begin{quote}
\textbf{ব্যাখ্যা:} \emph{০.৪৭=(৪৭/১০০)}
\end{quote}

\textbf{০.০২২• এর সাধারণ ভগ্নাংশ কোনটি?}

ক) (২১/৯৯৯)

খ) (২২/১০০০)

গ) (২১/৯০০)

\textbf{\hl{ঘ) (২০/৯০০) ✔}}

\textbf{\hl{✔ সঠিক উত্তর: ঘ) (২০/৯০০)}}

\begin{quote}
\textbf{ব্যাখ্যা:} \emph{দেয়া আছে, ০.০২২•=(২২−২/৯০০)=(২০/৯০০)}
\end{quote}

\textbf{৪.৩• কে সাধারণ ভগ্নাংশে প্রকাশ করলে নিচের কোনটি হবে?}

ক) (৩/১১)

খ) (১১/৩)

গ) (৩/১৩)

\textbf{\hl{ঘ) (১৩/৩) ✔}}

\textbf{\hl{✔ সঠিক উত্তর: ঘ) (১৩/৩)}}

\begin{quote}
\textbf{ব্যাখ্যা:} \emph{৪.৩•=(৪৩−৪/৯)=(৩৯/৯)=(১৩/৩)}
\end{quote}

\textbf{০.২•৪• এর সমান ভগ্নাংশটি}

ক) (১/৩)

\textbf{\hl{খ) (৮/৩৩) ✔}}

গ) (৭/২৪)

ঘ) (৭/৭২)

\textbf{\hl{✔ সঠিক উত্তর: খ) (৮/৩৩)}}

\begin{quote}
\textbf{ব্যাখ্যা:} \emph{০.২•৪•=(২৪−০/৯৯)=(২৪/৯৯)=(৮/৩৩)}
\end{quote}

\textbf{নিচের সংখ্যাগুলোর মধ্যে কোনটি ক্ষুদ্রতম?}

\textbf{\hl{ক) ০.৩ ✔}}

খ) √(০.৩)

গ) (২/৩)

ঘ) (২/৫)

\textbf{\hl{✔ সঠিক উত্তর: ক) ০.৩}}

\begin{quote}
\textbf{ব্যাখ্যা:} \emph{এখানে,}

\emph{√(০.৩)≈০.৫৪৭}

\emph{(২/৩)=০.৬৬}

\emph{(২/৫)=০.৪০}

\emph{∴ দেখা যাচ্ছে, ০.৩-ই ক্ষুদ্রতম সংখ্যা।}
\end{quote}

\textbf{সবচেয়ে বড় সংখ্যা কোনটি?}

ক) ০.০০৯৯

\textbf{\hl{খ) ০.১০০ ✔}}

গ) (৯/১০০)

ঘ) (৯/১০০০)

\textbf{\hl{✔ সঠিক উত্তর: খ) ০.১০০}}

\begin{quote}
\textbf{ব্যাখ্যা:} \emph{ক. ০.০০৯৯ = (৯৯/১০০০০)}

\emph{খ. ০.১০০ = (১০০০/১০০০০) (বৃহত্তম)}

\emph{গ. (৯/১০০)=০.০৯=(৯০০/১০০০০)}

\emph{ঘ. (৯/১০০০)=০.০০৯=(৯০/১০০০০)}

\emph{∴ এদের মধ্যে ০.১-ই সবচেয়ে বড়। তাই এটি উত্তর হবে।}
\end{quote}

\textbf{√(০.০৯)= কত?}

ক) ০৩

\textbf{\hl{খ) ০.৩ ✔}}

গ) ০.০০৩

ঘ) ০.০০০৩

\textbf{\hl{✔ সঠিক উত্তর: খ) ০.৩}}

\begin{quote}
\textbf{ব্যাখ্যা:} \emph{√(০.০৯)=√((৯/১০০))=(৩/১০)=০.৩}
\end{quote}

\textbf{কোন ভগ্নাংশটি ০.৫\% এর সমতুল্য?}

ক) (১/২০)

\textbf{\hl{খ) (১/২০০) ✔}}

গ) (১/৫০)

ঘ) (১/৫)

\textbf{\hl{✔ সঠিক উত্তর: খ) (১/২০০)}}

\begin{quote}
\textbf{ব্যাখ্যা:} \emph{০.৫\%=(০.৫/১০০)=(৫/১০০×১০)=(১/২০০)}
\end{quote}

\textbf{০.২• এর সামান্য ভগ্নাংশ কোনটি?}

ক) (২/১০)

খ) (২/৮)

\textbf{\hl{গ) (২/৯) ✔}}

ঘ) (২/৯৯)

\textbf{\hl{✔ সঠিক উত্তর: গ) (২/৯)}}

\begin{quote}
\textbf{ব্যাখ্যা:} \emph{০.২•=(২−০/৯)=(২/৯)}
\end{quote}

\textbf{২.৯ + ০০৯ + ৯৯৯.০০৯৯ = ?}

\textbf{\hl{ক) ১০১০.০০৯৯ ✔}}

খ) ১০১০.০৯৯

গ) ১০১০.৯৯

ঘ) ১০১০.০০৯

\textbf{\hl{✔ সঠিক উত্তর: ক) ১০১০.০০৯৯}}

\begin{quote}
\textbf{ব্যাখ্যা:} \emph{২.৯০০০ + ০০৯.০০০০ + ৯৯৯.০০৯৯ = ১০১০.০০৯৯}
\end{quote}

\textbf{১.১, .০১ ও ০.০১১ এর সমষ্টি কত?}

\textbf{\hl{ক) ১.১২১ ✔}}

খ) ১.১১২

গ) ১.২১১

ঘ) ১.০১১

\textbf{\hl{✔ সঠিক উত্তর: ক) ১.১২১}}

\begin{quote}
\textbf{ব্যাখ্যা:} \emph{১.১০০০ + ০.০১০০ + ০.০১১১∴ নির্ণেয় সমষ্টি = ১.১২১}
\end{quote}

\textbf{০.১ × ০.০১ = কত?}

ক) ১

\textbf{\hl{খ) ০.০০১ ✔}}

গ) ০.১০

ঘ) ০.০১

\textbf{\hl{✔ সঠিক উত্তর: খ) ০.০০১}}

\begin{quote}
\textbf{ব্যাখ্যা:} \emph{০.১×০.০১=০.০০১}
\end{quote}

\textbf{০.২৫ × ০.৪ = কত?}

ক) ১.০

\textbf{\hl{খ) ০.১ ✔}}

গ) ০.০১

ঘ) ০.০০১

\textbf{\hl{✔ সঠিক উত্তর: খ) ০.১}}

\begin{quote}
\textbf{ব্যাখ্যা:} \emph{০.২৫×০.৪=০.১০০=০.১}
\end{quote}

\textbf{০.১ × ০.০১ + ১ এর মান কত?}

ক) ১.০১

\textbf{\hl{খ) ১.০০১ ✔}}

গ) ২.০১

ঘ) .০০১

\textbf{\hl{✔ সঠিক উত্তর: খ) ১.০০১}}

\begin{quote}
\textbf{ব্যাখ্যা:} \emph{০.১×০.০১+১=০.০০১+১=১.০০১}
\end{quote}

\textbf{৩.০০০১০+৫×১০\^{}−৩= কত?}

\textbf{\hl{ক) ৩.০০৫১০ ✔}}

খ) ৩.০৫০১০

গ) ৩.০০০১৫

ঘ) ৩.০০০৬০

\textbf{\hl{✔ সঠিক উত্তর: ক) ৩.০০৫১০}}

\begin{quote}
\textbf{ব্যাখ্যা:} \emph{৩.০০০১০+৫×১০\^{}−৩}

\emph{= ৩.০০০১০+(৫/১০০০) = ৩.০০০১০+০.০০৫ = ৩.০০০১০+০.০০৫০০=৩.০০৫১০}
\end{quote}

\textbf{০.৪ × ০.০২ × ০.০৮ = ?}

ক) ০.৬৪

খ) ০.০৬৪

\textbf{\hl{গ) .০০০৬৪ ✔}}

ঘ) ৬.৪০

\textbf{\hl{✔ সঠিক উত্তর: গ) .০০০৬৪}}

\begin{quote}
\textbf{ব্যাখ্যা:} \emph{০.৪×০.০২×০.০৮=০.০০০৬৪}
\end{quote}

\textbf{০.৪ × ০.০০২ × ০.০৮ = ?}

ক) ৬.৪

খ) ০.৬৪

গ) ০.০৬৪

\textbf{\hl{ঘ) ০.০০০৬৪ ✔}}

\textbf{\hl{✔ সঠিক উত্তর: ঘ) ০.০০০৬৪}}

\begin{quote}
\textbf{ব্যাখ্যা:} \emph{০.৪×০.০০২×০.০৮}

\emph{= (৪/১০)×(২/১০০০)×(৮/১০০)=(৬৪/১০০০০০)=০.০০০৬৪}
\end{quote}

\textbf{০.০১ × ০.০১ = কত?}

ক) ০.০০২

খ) ০.০০১

\textbf{\hl{গ) ০.০০০১ ✔}}

ঘ) ১

\textbf{\hl{✔ সঠিক উত্তর: গ) ০.০০০১}}

\begin{quote}
\textbf{ব্যাখ্যা:} \emph{০.০১×০.০১=০.০০০১}
\end{quote}

\textbf{০.১ × ০.১ × ০.১ = কত?}

ক) ০.১

\textbf{\hl{খ) ০.০০১ ✔}}

গ) ০.০০০১

ঘ) ১.০

\textbf{\hl{✔ সঠিক উত্তর: খ) ০.০০১}}

\begin{quote}
\textbf{ব্যাখ্যা:} \emph{০.১×০.১×০.১=০.০০১}
\end{quote}

\textbf{০.০২ × ০.০১ × ০.০৫ = কত?}

ক) ০.০০০১

\textbf{\hl{খ) ০.০০০০০১ ✔}}

গ) ০.০০০০০১

ঘ) ০.০০১

\textbf{\hl{✔ সঠিক উত্তর: খ) ০.০০০০০১}}

\begin{quote}
\textbf{ব্যাখ্যা:}
\emph{০.০২×০.০১×০.০৫=(২/১০০)×(১/১০০)×(৫/১০০)=(১০/১০০০০০০)=০.০০০০০১}
\end{quote}

\textbf{০.০০১ × ০.০০০৮৭৫ = ?}

ক) ০.০০০০০১

খ) ০.০০০০০০১

\textbf{\hl{গ) ০.০০০০০০৮৭৫ ✔}}

ঘ) ০.০০০০০৮৭৫

\textbf{\hl{✔ সঠিক উত্তর: গ) ০.০০০০০০৮৭৫}}

\begin{quote}
\textbf{ব্যাখ্যা:} \emph{০.০০১ × ০.০০০৮৭৫}

\emph{= (১/১০০০) × (৮৭৫/১০০০০০০)}

\emph{= (৮৭৫/১০০০০০০০০০) = ০.০০০০০০৮৭৫}
\end{quote}

\textbf{৩ × ০.৩ ÷ ১ = কত?}

ক) ১

খ) ০.৬

গ) ২

\textbf{\hl{ঘ) ০.৯ ✔}}

\textbf{\hl{✔ সঠিক উত্তর: ঘ) ০.৯}}

\begin{quote}
\textbf{ব্যাখ্যা:} \emph{৩ × ০.৩ ÷ ১ = ৩ × ০.৩ = ০.৯}
\end{quote}

\textbf{০.০২৪ × ১০\^{}৬ = ?}

\textbf{\hl{ক) ২৪০০০ ✔}}

খ) ২৪০০০০

গ) ২৪০০০০০

ঘ) কোনোটিই নয়

\textbf{\hl{✔ সঠিক উত্তর: ক) ২৪০০০}}

\begin{quote}
\textbf{ব্যাখ্যা:} \emph{০.০২৪ × ১০\^{}৬ = ০.০২৪ × ১০০০০০০ = ২৪,০০০}
\end{quote}

\textbf{(০.০১ × ১)\^{}২ = কত?}

ক) ০.১

খ) ০.০১

গ) ০.০০১

\textbf{\hl{ঘ) ০.০০০১ ✔}}

\textbf{\hl{✔ সঠিক উত্তর: ঘ) ০.০০০১}}

\begin{quote}
\textbf{ব্যাখ্যা:} \emph{(০.০১ × ১)\^{}২ = (০.০১)\^{}২ = ০.০০০১}
\end{quote}

\textbf{০.০০০৫ ÷ ০.০০৮ = কত?}

ক) ০.০০৬২৫

\textbf{\hl{খ) ০.০৬২৫ ✔}}

গ) ০.৬২৫০

ঘ) ৬.২৫০

\textbf{\hl{✔ সঠিক উত্তর: খ) ০.০৬২৫}}

\begin{quote}
\textbf{ব্যাখ্যা:} \emph{০.০০০৫ ÷ ০.০০৮}

\emph{= (৫/১০০০০) ÷ (৮/১০০০)}

\emph{= (৫/১০০০০) × (১০০০/৮)}

\emph{= (৫/৮০)}

\emph{∴ ০.০০০৫ ÷ ০.০০৮ = ০.০৬২৫}
\end{quote}

\textbf{৪ ÷ ০.১২৫ = কত?}

ক) ৬৪

খ) ৬.৪

\textbf{\hl{গ) ৩২ ✔}}

ঘ) ৩.২

\textbf{\hl{✔ সঠিক উত্তর: গ) ৩২}}

\begin{quote}
\textbf{ব্যাখ্যা:} \emph{৪ ÷ ০.১২৫}

\emph{= ৪ ÷ (১২৫/১০০০)}

\emph{= ৪ ÷ (১/৮) = ৪ × ৮ = ৩২}
\end{quote}

\textbf{০.০৬৭৮ কে ০.৮৫ দ্বারা ভাগ করলে সঠিক ভাগফল কোনটি?}

ক) ০.০০০২০৮২৫

\textbf{\hl{খ) ০.০৭৯৭৬৪৭ ✔}}

গ) ০.০৬৯৭৬৪

ঘ) ০.০৭৯৭৬৪৮

\textbf{\hl{✔ সঠিক উত্তর: খ) ০.০৭৯৭৬৪৭}}

\begin{quote}
\textbf{ব্যাখ্যা:} \emph{(০.০৬৭৮/০.৮৫)=০.০৭৯৭৬৪৭ (প্রায়)}
\end{quote}

\textbf{৩.১২৫ কে ২.৫ দ্বারা ভাগ করলে ভাগফল কত?}

ক) ১.৩৫

খ) ১.৪০

\textbf{\hl{গ) ১.২৫ ✔}}

ঘ) ১.৩০

\textbf{\hl{✔ সঠিক উত্তর: গ) ১.২৫}}

\begin{quote}
\textbf{ব্যাখ্যা:} \emph{(৩.১২৫/২.৫)=১.২৫}
\end{quote}

\textbf{(০.১×১.১×১.২/০.০১×০.০২) এর মান কত ?}

ক) ৫৫০

খ) ২০০

গ) ১২০

\textbf{\hl{ঘ) ৬৬০ ✔}}

\textbf{\hl{✔ সঠিক উত্তর: ঘ) ৬৬০}}

\begin{quote}
\textbf{ব্যাখ্যা:} \emph{(০.১×১.১×১.২/০.০১×০.০২)}

\emph{=((১/১০)×(১১/১০)×(১২/১০)/(১/১০০)×(২/১০০))}

\emph{=(১×১১×১২/১০×১০×১০)×(১০০×১০০/২) =৬৬০}
\end{quote}

\textbf{(.১×.০১×.০০১/.২×.০২×.০০২) এর মান কত?}

\textbf{\hl{ক) (১/৮) ✔}}

খ) (১/৮০)

গ) (১/৮০০)

ঘ) (১/৮০০০)

\textbf{\hl{✔ সঠিক উত্তর: ক) (১/৮)}}

\begin{quote}
\textbf{ব্যাখ্যা:} \emph{(.১×.০১×.০০১/.২×.০২×.০০২)}

\emph{=(১×১×১×১০×১০০×১০০০/২×২×২×১০×১০০×১০০০)}

\emph{=(১/৮)}
\end{quote}

\textbf{(০.১×০.১•/?−০.১)=১.০ হলে প্রশ্নবোধক চিহ্নের ঘরে কত বসবে?}

ক) ০.১১

\textbf{\hl{খ) ০.১• ✔}}

গ) ০.১১•

ঘ) ১.০১

\textbf{\hl{✔ সঠিক উত্তর: খ) ০.১•}}

\begin{quote}
\textbf{ব্যাখ্যা:} \emph{ধরি, প্রশ্নবোধক চিহ্ন = x}

\emph{(০.১×০.১•/x−০.১)=১.০}

\emph{⇒ (০.১×(১/৯)/x−০.১)=১.০}

\emph{⇒ x−০.১=(০.১/৯)}

\emph{⇒ x=(০.১/৯)+০.১=(১/৯০)+(১/১০)=(১+৯/৯০)=(১০/৯০)=(১/৯)=০.১•}
\end{quote}

\textbf{নিচের দশমিক সংখ্যাগুলোর বৃহত্তম সংখ্যা দ্বারা ক্ষুদ্রতম সংখ্যা গুণ করলে
গুণফল কত হবে?}

০.১, ০.০০০৯, ০.০২০, ০.০০১

\textbf{\hl{ক) ০.০০০০৯ ✔}}

খ) ০.০০০১

গ) ০.০০১৮

ঘ) ০.০০০১৮

\textbf{\hl{✔ সঠিক উত্তর: ক) ০.০০০০৯}}

\begin{quote}
\textbf{ব্যাখ্যা:} \emph{বৃহত্তম সংখ্যা = ০.১}

\emph{ক্ষুদ্রতম সংখ্যা = ০.০০০৯}

\emph{তাদের গুণফল = ০.১ × ০.০০০৯ = ০.০০০০৯}
\end{quote}

\textbf{(১/৮০) থেকে (১/৪০) বিয়োগ করে বিয়োগফলকে দশমিকে প্রকাশ করলে কত হয়?}

ক) ০.২২৫

\textbf{\hl{খ) −০.০১২৫ ✔}}

গ) ০.২৫

ঘ) ০.০৫

\textbf{\hl{✔ সঠিক উত্তর: খ) −০.০১২৫}}

\begin{quote}
\textbf{ব্যাখ্যা:} \emph{(১/৮০)−(১/৪০)=(১−২/৮০)=−(১/৮০)=−০.০১২৫}
\end{quote}

\textbf{এক-দশমাংশ, এক-শতাংশ এবং এক-সহস্রাংশ এর গড় হবে---}

ক) ০.০০০৩

খ) ০.০১

\textbf{\hl{গ) ০.০৩৭ ✔}}

ঘ) ০.১১১

\textbf{\hl{✔ সঠিক উত্তর: গ) ০.০৩৭}}

\begin{quote}
\textbf{ব্যাখ্যা:} \emph{এক-দশমাংশ = ০.১}

\emph{এক-শতাংশ = ০.০১}

\emph{এক-সহস্রাংশ = ০.০০১}

\emph{সংখ্যা তিনটির সমষ্টি = ০.১ + ০.০১ + ০.০০১ = ০.১১১}

\emph{∴ নির্ণেয় গড় = ০.১১১ ÷ ৩ = ০.০৩৭}
\end{quote}

\textbf{সবচেয়ে বড় সংখ্যা কোনটি?}

ক) ০.০০৯৯

\textbf{\hl{খ) ০.১০০ ✔}}

গ) (৯/১০০)

ঘ) (৯/১০০০)

\textbf{\hl{✔ সঠিক উত্তর: খ) ০.১০০}}

\begin{quote}
\textbf{ব্যাখ্যা:} \emph{ক. ০.০০৯৯ = (৯৯/১০০০০)}

\emph{খ. ০.১০০ = (১/১০) (বৃহত্তম)}

\emph{গ. (৯/১০০)=০.০৯}

\emph{ঘ. (৯/১০০০)=০.০০৯}

\emph{∴ এদের মধ্যে ০.১০০ সবচেয়ে বড়।}
\end{quote}

\textbf{কোন সংখ্যাটি বৃহত্তম?}

\textbf{\hl{ক) ০.১০০০ ✔}}

খ) ০.০১০০

গ) ০.০০১০

ঘ) ০.০০০১

\textbf{\hl{✔ সঠিক উত্তর: ক) ০.১০০০}}

\begin{quote}
\textbf{ব্যাখ্যা:} \emph{ক. ০.১০০০ = (১/১০) {[}বৃহত্তম{]}}

\emph{খ. ০.০১০০ = (১/১০০)}

\emph{গ. ০.০০১০ = (১/১০০০)}

\emph{ঘ. ০.০০০১ = (১/১০০০০)}

\emph{∴ ০.১০১০ বৃহত্তম সংখ্যা।}
\end{quote}

\textbf{০.২×০.০০২ ×০.২৫ = কত?}

ক) ০.০০০৪

খ) ০.০০১

গ) ০.০১

\textbf{\hl{ঘ) ০.০০০১ ✔}}

\textbf{\hl{✔ সঠিক উত্তর: ঘ) ০.০০০১}}

\begin{quote}
\textbf{ব্যাখ্যা:} \emph{০.২ × ০.০০২ × ০.২৫ = ০.০০০৪ × ০.২৫ = ০.০০০০০১০০ =
০.০০০১}
\end{quote}

\textbf{(০.০০৩)\^{}২ = কত?}

\textbf{\hl{ক) ০.০০০০০৯ ✔}}

খ) ০.০০০০৯

গ) ০.০০০৯

ঘ) ০.০০৯

\textbf{\hl{✔ সঠিক উত্তর: ক) ০.০০০০০৯}}

\begin{quote}
\textbf{ব্যাখ্যা:} \emph{(০.০০৩)\^{}২ = ০.০০৩ × ০.০০৩ = ০.০০০০০৯}
\end{quote}

\textbf{(০.০০৪)\^{}২ = কত?}

ক) ০.০১৬

\textbf{\hl{খ) ০.০০০০১৬ ✔}}

গ) ০.০০০১৬

ঘ) ০.০০১৬

\textbf{\hl{✔ সঠিক উত্তর: খ) ০.০০০০১৬}}

\begin{quote}
\textbf{ব্যাখ্যা:} \emph{(০.০০৪)\^{}২ = ০.০০৪ × ০.০০৪ = ০.০০০০১৬}
\end{quote}

\textbf{(০.০০৫)\^{}২ =কত?}

ক) ০.০২৫

খ) ০.০০২৫

\textbf{\hl{গ) ০.০০০০২৫ ✔}}

ঘ) ০.০০০০০০২৫

\textbf{\hl{✔ সঠিক উত্তর: গ) ০.০০০০২৫}}

\begin{quote}
\textbf{ব্যাখ্যা:} \emph{(০.০০৫)\^{}২ = ০.০০৫ × ০.০০৫ = ০.০০০০২৫}
\end{quote}

\textbf{(০.১×০.০২×০.০০২/০.০১×০.০৪) এর মান কত?}

\textbf{\hl{ক) ০.০১ ✔}}

খ) ০.১

গ) ০.০০১

ঘ) ০.০০৪

\textbf{\hl{✔ সঠিক উত্তর: ক) ০.০১}}

\begin{quote}
\textbf{ব্যাখ্যা:}
\emph{(০.১×০.০২×০.০০২/০.০১×০.০৪)=(০.০০০০০৪/০.০০০৪)=০.০১}
\end{quote}

\textbf{(০.১×০.২×০.০০৩/০.০১×০.০২×০.০৩) এর মান কত ?}

ক) ০.১

খ) ০.০১

গ) ০.০০১

\textbf{\hl{ঘ) ১০ ✔}}

\textbf{\hl{✔ সঠিক উত্তর: ঘ) ১০}}

\begin{quote}
\textbf{ব্যাখ্যা:}
\emph{(০.১×০.২×০.০০৩/০.০১×০.০২×০.০৩)=(০.০০০০০৬/০.০০০০০০৬)=১০}
\end{quote}

\textbf{(০.১×০.০২×০.০০২/০.১×০.০৪) এর মান কত?}

ক) ০.১

\textbf{\hl{খ) ০.০১ ✔}}

গ) ০.০২

ঘ) ০.০০১

\textbf{\hl{✔ সঠিক উত্তর: খ) ০.০১}}

\begin{quote}
\textbf{ব্যাখ্যা:} \emph{(০.১×০.০২×০.০০২/০.১×০.০৪)}

\emph{= (১×২×২/১×৪×১০০) = ০.০১}
\end{quote}

\section{61. ল.সা.গু ও গ.সা.গু (LCM \& HCF) নির্ণয়ের নিঞ্জা
টেকনিক}\label{ux9b2.ux9b8.ux997-ux993-ux997.ux9b8.ux997-lcm-hcf-ux9a8ux9b0ux9a3ux9dfux9b0-ux9a8ux99eux99c-ux99fux995ux9a8ux995}

\textbf{ল.সা.গু ও গ.সা.গু নির্ণয়} অধ্যায়টি পাটিগণিতের খুবই গুরুত্বপূর্ণ একটি
অধ্যায়। এ অধ্যায়টি সম্পূর্ণ নখদর্পণে নিতে হলে প্রথমেই আপনাকে ল.সা.গু এবং গ.সা.গু
নির্ণয়ে সিদ্ধহস্ত হতে হবে। তাই এ অধ্যায়ের বিগত সালের প্রশ্ন সমাধান করার আগে
ল.সা.গু এবং গ.সা.গু নির্ণয়ে পারদর্শী হতে হবে । এজন্য প্রচলিত নিয়মের পাশাপাশি
সবচেয়ে দ্রুত এবং সহজভাবে ল.সা.গু এবং গ.সা.গু নির্ণয়ের জন্য একটি কৌশল নিচে
দেওয়া হল । কৌশলটি রপ্ত করার চেষ্টা করুন। আশা করি ভাল ফল পাবেন।

ল.সা.গু নির্ণয় সবচেয়ে সহজ পদ্ধতি

গ.সা.গু নির্ণয়ের সবচেয়ে সহজ পদ্ধতি

\begin{itemize}
\item
  ভগ্নাংশের ল. সা. গু ও গ. সা. গু নির্ণয়ের নিয়ম
\item
  অনুপাত আকারে ল.সা.গু ও গ.সা.গু
\end{itemize}

\subsubsection{ল.সা.গু নির্ণয় সবচেয়ে সহজ
পদ্ধতি}\label{ux9b2.ux9b8.ux997-ux9a8ux9b0ux9a3ux9df-ux9b8ux9acux99aux9df-ux9b8ux9b9ux99c-ux9aaux9a6ux9a7ux9a4}

সমস্যা-১। ২,৪,৮,১০ এর ল.সা.গু নির্ণয় কর।

২,৪,৮,১০ এর ল.সা.গু = ৪০ ।

যারা দ্রুত ল.সা.গু নির্ণয় করতে চান তাদের জন্য এ পদ্ধতি।কারণ এমসিকিউ পরীক্ষায়
সময় অনেক কম দেওয়া হয়।এই কম সময়ে অনেক প্রশ্নের উত্তর দিতে হয়।তাই এমসিকিউ
পরীক্ষার জন্য এ পদ্ধতিই হচ্ছে বেস্ট পদ্ধতি।এ পদ্ধতিতে সবচেয়ে বড় সংখ্যাটিকে বেস
ধরে দেখা হয় ছোট সংখ্যাটি দিয়ে বড় সংখ্যাটিকে ভাগ করা যায় কি না? যেমন উপরোক্ত
সংখ্যাগুলোর মধ্যে ১০ সবচেয়ে বড় সংখ্যা।এখন প্রথম সংখ্যা থেকে দেখা যায় ২ দ্বারা
১০ কে ভাগ করা যায় তাই ২ বাদ। এরপর ৪ পরবর্তী ৮ এর গুণনীয়ক তাই ৪ ও বাদ।বাকি
থাকল ৮ এবং ১০।৮ দ্বারা ১০ ভাগ হয় না।তাই প্রথমে ২ দ্বারা ১০ কে গুণ করি।তবু ও
ভাগ হয় না।অতঃপর ৩ এবং পরবর্তীতে ৪ দ্বারা ১০ কে গুণ করলে ৪০ হয় যা ৮ দ্বারা
ভাগ করা যায়।তাই ৪০ ই ২,৪,৮,১০ এর ল.সা.গু।

আরো কিছু সমস্যাঃ

১। ৩,৬,১২ এর ল.সা.গু = ১২

২।১২,২৪ এর ল.সা.গু = ২৪

৩।৪,৬,১২ এর ল.সা.গু = ১২

৪। ২,৪,৬,১২,২৪ এর ল.সা.গু = ২৪

লক্ষ করুনঃ

১নং এ দেখতে পাচ্ছি ৩ সবচেয়ে ছোট সংখ্যা এবং ১২ সবচেয়ে বড় সংখ্যা।৩ দ্বারা ১২
কে ভাগ করা যায় তাই ৩ বাদ।আবার ৬ দ্বারা ও ১২ কে ভাগ করা যায় তাই ৬ ও
বাদ।তাই ল.সা.গু হবে সবচেয়ে বড় সংখ্যা ১২।

২নং এ দেখতে পাচ্ছি ১২ ছোট সংখ্যা এবং ২৪ বড় সংখ্যা।১২ দ্বারা ২৪ কে ভাগ করা
যায় বিধায় ১২ বাদ।তাই ল.সা.গু হবে সবচেয়ে বড় সংখ্যা ২৪।

৩নং এ দেখতে পাচ্ছি ৪ সবচেয়ে ছোট সংখ্যা এবং ১২ সবচেয়ে বড় সংখ্যা।৪ দ্বারা ১২
কে ভাগ করা যায় তাই ৪ বাদ।আবার ৬ দ্বারা ও ১২ কে ভাগ করা যায় তাই ৬ ও
বাদ।তাই ল.সা.গু হবে সবচেয়ে বড় সংখ্যা ১২।

৪নং এ দেখতে পাচ্ছি ২ সবচেয়ে ছোট সংখ্যা এবং ২৪ সবচেয়ে বড় সংখ্যা।২ দ্বারা ২৪
কে ভাগ করা যায় তাই ২ বাদ। ৪ দ্বারা ও ২৪ কে ভাগ করা যায় তাই ৪ ও বাদ।৬
দ্বারা ও ২৪ কে ভাগ করা যায় তাই ৬ ও বাদ। ১২ দ্বারা ও ২৪ কে ভাগ করা যায় তাই
১২ ও বাদ।তাই ল.সা.গু হবে সবচেয়ে বড় সংখ্যা ২৪।

এবার যদি এমন হয় সবচেয়ে বড় সংখ্যাটিকে সবগুলো সংখ্যা দিয়ে ভাগ করা না যায় তখন
কিভাবে ল.সা.গু বের করতে হবে।চলুন নিচের সমস্যাসমুহ দেখি।

সমস্যা-৫। ৩ ও ৫ এর ল.সা,গু= ১৫

সমস্যা-৬। ৬ ও ৮ এর ল.সা.গু = ২৪

সমস্যা-৭। ৫,১০,১৫ ও ২০ এর ল.সা.গু = ৬০

লক্ষ করুনঃ

৫নং এ দেখতে পাচ্ছি ৩ ছোট সংখ্যা এবং ৫ বড় সংখ্যা।৩ দ্বারা ৫কে ভাগ করা যায়
না।তাই প্রথমে ২ দ্বারা ৫ কে গুণ করি।তবু ও ভাগ হয় না।অতঃপর ৩ দ্বারা ৫ কে গুণ
করলে ১৫ হয় যা ৩ দ্বারা ভাগ করা যায়।তাই ১৫ ই হচ্ছে ৩ ও ৫ এর ল.সা.গু।

৬ নং এ দেখতে পাচ্ছি ৬ ছোট সংখ্যা এবং ৮ বড় সংখ্যা।৬ দ্বারা ৮কে ভাগ করা যায়
না।তাই প্রথমে ২ দ্বারা ৮ কে গুণ করি।তবু ও ভাগ হয় না।অতঃপর ৩ দ্বারা ৮ কে গুণ
করলে ২৪ হয় যা ৬ দ্বারা ভাগ করা যায়।তাই ২৪ ই হচ্ছে ৬ ও ৮ এর ল.সা.গু।

৭নং এ দেখতে পাচ্ছি ৫ সবচেয়ে ছোট সংখ্যা এবং ২০ সবচেয়ে বড় সংখ্যা। এখন প্রথম
সংখ্যা থেকে দেখা যায় ৫ দ্বারা ২০ কে ভাগ করা যায় তাই ৫ বাদ। এরপর ১০ দ্বারা ও
২০ কে ভাগ করা যায় তাই ১০ ও বাদ।বাকি থাকল ১৫ এবং ২০।১৫ দ্বারা ২০ ভাগ হয়
না।তাই প্রথমে ২ দ্বারা ২০ কে গুণ করি।তবু ও ভাগ হয় না।অতঃপর ৩ দ্বারা ২০ কে গুণ
করলে ৬০ হয় যা ১৫ দ্বারা ভাগ করা যায়।তাই ৬০ ই হচ্ছে ৫,১০,১৫ ও ২০ এর ল.সা.গু।

\subsubsection{গ.সা.গু নির্ণয়ের সবচেয়ে সহজ
পদ্ধতি}\label{ux997.ux9b8.ux997-ux9a8ux9b0ux9a3ux9dfux9b0-ux9b8ux9acux99aux9df-ux9b8ux9b9ux99c-ux9aaux9a6ux9a7ux9a4}

যারা দ্রুত গ.সা.গু নির্ণয় করতে চান তাদের জন্য এ পদ্ধতি।কারণ এমসিকিউ পরীক্ষায়
অনেক কম সময় দেওয়া হয়।এই কম সময়ে অনেক প্রশ্নের উত্তর দিতে হয়।তাই এমসিকিউ
পরীক্ষার জন্য এ পদ্দতিই হচ্ছে বেস্ট পদ্ধতি।এ পদ্ধতিতে সবচেয়ে ছোট সংখ্যাটিকে বেস
ধরে দেখা হয় ছোট সংখ্যাটি দিয়ে বড় সংখ্যাটিকে ভাগ করা যায় কি না?

সমস্যা-১। ১৩ এবং ২৬ এর গ.সা.গু =১৩

সমস্যা-২। ৫,১০ এবং ২০ এর গ.সা.গু = ৫

সমস্যা-৩। ৩,৬,১২ এবং ১৮ এর গ.সা.গু = ৩

লক্ষ করুনঃ

১নং এ দেখতে পাচ্ছি ১৩ ছোট সংখ্যা এবং ২৬ বড় সংখ্যা। ১৩ দ্বারা ২৬ কে ভাগ করা
যায় বিধায় ১৩ ই হবে ১৩ এবং ২৬ এর গ.সা.গু।

২নং এ দেখতে পাচ্ছি ৫ সবচেয়ে ছোট সংখ্যা । সবচেয়ে ছোট সংখ্যা ৫ দ্বারা ১০ এবং
২০ উভয়কেই ভাগ করা যায় বিধায় ৫ ই হবে ৫,১০ এবং ২০ এর গ.সা.গু।

৩নং এ দেখতে পাচ্ছি ৩ সবচেয়ে ছোট সংখ্যা। সবচেয়ে ছোট সংখ্যা ৩ দ্বারা ৬,১২ এবং
১৮ এ তিনটি সংখ্যাকেই ভাগ করা যায় বিধায় ৩ ই হবে ৩,৬,১২ এবং ১৮ এর গ.সা.গু।

এখন যদি এমন হয় সবচেয়ে ছোট সংখ্যাটি দ্বারা সবগুলো সংখ্যাকে ভাগ করা না যায় তখন
কিভাবে গ.সা.গু বের করতে হবে।চলুন নিচের সমস্যাসমুহ দেখি।

সমস্যা-৪। ৯,১২ এবং ১৫ এর গ.সা.গু =৩

সমস্যা-৫। ১২,২৪ এবং ৩৩ এর গ.সা.গু = ৩

সমস্যা-৬। ৩ এবং ৫ এর গ.সা.গু =১

সমস্যা-৭। ৩২,৪৮,৫৬ এবং ৮০ এর গ.সা.গু =৮

সমস্যা-৮। ২৪,৩০ এবং ৭৭ এর গ.সা.গু = ১

লক্ষ করুনঃ

৪নং এ দেখতে পাচ্ছি ৯ সবচেয়ে ছোট সংখ্যা । সবচেয়ে ছোট সংখ্যা ৯ দ্বারা ১২ এবং
১৫ ভাগ হয় না।তাই ৯ এর গুণনীয়ক দ্বারা ১২ এবং ১৫ কে ভাগ করতে হবে।৯ এর গুণনীয়ক
হচ্ছে ১,৩,৯।এবার সবচেয়ে বড় গুণনীয়ক থেকে ভাগ করা শুরু করতে হবে।এখন ৯ এর
গুণনীয়ক ৯ দ্বারা ১২ এবং ১৫ কে ভাগ করা যায় না।তাই ৯ বাদ। ৯ এর গুণনীয়ক ৩ দিয়ে
১২ এবং ১৫ কে ভাগ করা যায় বিধায় ৩ ই হবে ৯,১২ এবং ১৫ এর গ.সা.গু।

৫নং এ দেখতে পাচ্ছি ১২ সবচেয়ে ছোট সংখ্যা। সবচেয়ে ছোট সংখ্যা ১২ দ্বারা ২৪ এবং
৩৩ ভাগ হয় না।তাই ১২ এর গুণনীয়ক দ্বারা ২৪ এবং ৩৩ কে ভাগ করতে হবে। ১২ এর
গুণনীয়ক হচ্ছে ১,২,৩,৪,৬,১২।এবার সবচেয়ে বড় গুণনীয়ক থেকে ভাগ করা শুরু করতে হবে।
১২ এর গুণনীয়ক ১২ দ্বারা ২৪ এবং ৩৩ কে ভাগ করা যায় না।তাই ১২ বাদ। ১২ এর
গুণনীয়ক ৬ এবং ৪ দিয়ে ও ২৪ এবং ৩৩ দুটোকে ভাগ করা যায় না।তাই ৬ এবং ৪ ও বাদ।
১২ এর গুণনীয়ক ৩ দ্বারা ২৪ এবং ৩৩ কে ভাগ করা যায় বিধায় ৩ ই হবে ১২,২৪ এবং ৩৩
এর গ.সা.গু।

৬নং এ দেখতে পাচ্ছি ৩ ছোট সংখ্যা । ছোট সংখ্যা ৩ দ্বারা ৫ ভাগ হয় না।তাই ৩ এর
গুণনীয়ক দ্বারা ৫ কে ভাগ করতে হবে।৩ এর গুণনীয়ক হচ্ছে ১,৩।এবার সবচেয়ে বড়
গুণনীয়ক থেকে ভাগ করা শুরু করতে হবে।এখন ৩ এর গুণনীয়ক ৩ দ্বারা ৫ কে ভাগ করা যায়
না।তাই ৩ বাদ। ৩ এর গুণনীয়ক ১ দিয়ে ৫ কে ভাগ করা যায় বিধায় ১ ই হবে ৩ এবং ৫
এর গ.সা.গু।

৭নং এ দেখতে পাচ্ছি ৩২ সবচেয়ে ছোট সংখ্যা। সবচেয়ে ছোট সংখ্যা ৩২ দ্বারা ৪৮,৫৬
এবং ৮০ ভাগ হয় না।তাই ৩২ এর গুণনীয়ক দ্বারা ৪৮,৫৬ এবং ৮০ কে ভাগ করতে হবে। ৩২
এর গুণনীয়ক হচ্ছে ১,২,৪,৮,১৬,৩২।৩২ তো আগেই বাদ।এবার সবচেয়ে বড় গুণনীয়ক থেকে
ভাগ করা শুরু করতে হবে। ৩২ এর গুণনীয়ক ১৬ দ্বারা ৪৮ এবং ৮০ কে ভাগ করা গেলেও ৫৬
কে ভাগ করা যায় না।তাই ১৬ বাদ। ৩২ এর গুণনীয়ক ৮ দ্বারা ৪৮,৫৬ এবং ৮০ কে ভাগ
করা যায় বিধায় ৮ ই হবে ৩২,৪৮,৫৬ এবং ৮০ এর গ.সা.গু।

৮নং এ দেখতে পাচ্ছি ২৪ সবচেয়ে ছোট সংখ্যা। সবচেয়ে ছোট সংখ্যা ২৪ দ্বারা ৩০ এবং
৭৭ ভাগ হয় না।তাই ২৪ এর গুণনীয়ক দ্বারা ৩০ এবং ৭৭ কে ভাগ করতে হবে। ২৪ এর
গুণনীয়ক হচ্ছে ১,২,৩,৪,৬,৮,১২,২৪।এবার সবচেয়ে বড় গুণনীয়ক থেকে ভাগ করা শুরু করতে
হবে। ২৪ এর গুণনীয়ক ২,৩,৪,৬,৮,১২,২৪ দিয়ে ৩০ এবং ৭৭ দুটোকে ভাগ করা যায়
না।তাই ২,৩,৪,৬,৮,১২,২৪ বাদ।১ দ্বারা ভাগ করা যায় ২৪ কে । আবার ১ দ্বারা ৩০
এবং ৭৭ কে ভাগ করা যায় বিধায় ১ ই হবে ২৪,৩০ এবং ৭৭ এর গ.সা.গু।

\subsubsection{ভগ্নাংশের ল. সা. গু ও গ. সা. গু নির্ণয়ের
নিয়ম}\label{ux9adux997ux9a8ux9b6ux9b0-ux9b2.-ux9b8.-ux997-ux993-ux997.-ux9b8.-ux997-ux9a8ux9b0ux9a3ux9dfux9b0-ux9a8ux9dfux9ae}

১. ভগ্নাংশের ল.সা.গু = \(\frac{\text{লবগুলোর}\ \text{ল}.\text{সা}.\text{গু}}{\text{হরগুলো}\ \text{গ}.\text{সা}.\text{গু}}\)

২. ভগ্নাংশের গ.সা.গু = \(\frac{\text{লবগুলোর}\ \text{গ}.\text{সা}.\text{গু}}{\text{হরগুলো}\ \text{ল}.\text{সা}.\text{গু}}\)

প্রশ্ন-১: \(\frac{\text{৩}}{\text{৪}}\) , \(\frac{\text{৪}}{\text{৫}}\) , \(\frac{\text{৫}}{\text{৬}}\) এর গ. সা. গু
কত? {[} ৮ম বেসরকারি প্রভাষক নিবন্ধন ও প্রত্যয়ন পরীক্ষা-২০১২{]}

\begin{itemize}
\item
  (ক) ৩০
\item
  (খ) (১/৩০)
\item
  (গ) (১/৬০)
\item
  (ঘ) ৬০
\end{itemize}

\textbf{\hl{উত্তরঃ (গ)}} \(\frac{\text{১}}{\text{৬০}}\)

\textbf{Explanation:}

৩,৪,৫ এর গ.সা.গু = ১ এবং ৪,৫ ও ৬ এর ল.সা.গু.= ৬০

দুই বা ততোধিক ভগ্নাংশের গ.সা.গু = \(\frac{\text{লবগুলোর}\ \text{গ}.\text{সা}.\text{গু}}{\text{হরগুলোর}\ \text{ল}.\text{সা}.\text{গু}}\)
= \(\frac{\text{১}}{\text{৬০}}\)

৫. \(\frac{\text{৩}}{\text{৫}}\) , \(\frac{\text{১}}{\text{৪}}\) , \(\frac{\text{২}}{\text{৩}}\) এর ল.সা.গু কত?
{[}গণপূর্ত অধিদপ্তর-২০১৬{]}

\textbf{\hl{উত্তরঃ (খ) ৬}}

\textbf{Explanation:}

লব ৩,১,২ এর ল.সা.গু= ৬ এবং হর ৫,৪ ও ৩ এর গ.সা.গু.= ১

দুই বা ততোধিক ভগ্নাংশের ল.সা.গু =
\(\frac{\text{লব}\ \text{গুলোর}\ \text{ল}.\text{সা}.\text{গু}}{\text{হরগুলোর}\ \text{গ}.\text{সা}.\text{গু}}\) = \(\frac{\text{৬}}{\text{১}}\) = ৬

\subsubsection{অনুপাত আকারে ল.সা.গু ও
গ.সা.গু}\label{ux985ux9a8ux9aaux9a4-ux986ux995ux9b0-ux9b2.ux9b8.ux997-ux993-ux997.ux9b8.ux997}

দুটি সংখ্যার অনুপাত এবং তাদের ল.সা.গু দেয়া থাকলে...

বড় সংখ্যা = \(\frac{\text{সংখ্যাদ্বয়ের}\ \text{ল}.\text{সা}.\text{গু}}{\text{অনুপাতের}\ \text{ছোট}\ \text{সংখ্যা}}\)

ছোট সংখ্যা = \(\frac{\text{সংখ্যাদ্বয়ের}\ \text{ল}.\text{সা}.\text{গু}}{\text{অনুপাতের}\ \text{বড়}\ \text{সংখ্যা}}\)

প্রশ্ন-১: দুইটি সংখ্যার অনুপাত ৩:৪ এবং তাদের ল.সা.গু ১৮০। সংখ্যা দুইটি কি কি?
{[} শিক্ষক নিবন্ধন কলেজ-২০১২{]}

\begin{itemize}
\item
  (ক) ৭০,৬০
\item
  (খ) ৬০,৫০
\item
  (গ) ৫০,৪০
\item
  (ঘ) ৪৫,৬০
\end{itemize}

\textbf{\hl{উত্তরঃ (ঘ) ৪৫,৬০}}

\textbf{Explanation:}

দুটি সংখ্যার অনুপাত এবং তাদের ল.সা.গু দেয়া থাকলে...

বড় সংখ্যা = \(\frac{\text{সংখ্যাদ্বয়ের}\ \text{ল}.\text{সা}.\text{গু}}{\text{অনুপাতের}\ \text{ছোট}\ \text{সংখ্যা}}\)

= \(\frac{\text{১৮০}}{\text{৩}}\) = ৬০

ছোট সংখ্যা = \(\frac{\text{সংখ্যাদ্বয়ের}\ \text{ল}.\text{সা}.\text{গু}}{\text{অনুপাতের}\ \text{বড়}\ \text{সংখ্যা}}\) =
\(\frac{\text{১৮০}}{\text{৪}}\) = ৪৫

সুতরাং সংখ্যা দুইটি ৪৫ ও ৬০ ।

দুটি সংখ্যার গুণফল এবং তাদের ল.সা.গু কিংবা গ.সা.গু দেয়া থাকলে-

গ.সা.গু = সংখ্যাদ্বয়ের গুণফল ÷ সংখ্যাদ্বয়ের ল.সা.গু

ল.সা.গু = সংখ্যাদ্বয়ের গুণফল ÷ সংখ্যাদ্বয়ের গ.সা.গু

প্রশ্ন-২: দুইটি সংখ্যার গুণফল ১৫৩৬। সংখ্যা দুইটির ল.সা.গু ৯৬ হলে গ.সা.গু কত? {[}
১০ম বিসিএস{]}

\begin{itemize}
\item
  (ক) ৩২
\item
  (খ) ১২
\item
  (গ) ১৬
\item
  (ঘ) ২৪
\end{itemize}

\section{62. সাধারণ ভগ্নাংশ এ টু জেড: যোগ, বিয়োগ, গুণ, ভাগ ও ছোট-বড় নির্ণয়ের
সহজ
নিয়ম}\label{ux9b8ux9a7ux9b0ux9a3-ux9adux997ux9a8ux9b6-ux98f-ux99f-ux99cux9a1-ux9afux997-ux9acux9afux997-ux997ux9a3-ux9adux997-ux993-ux99bux99f-ux9acux9a1-ux9a8ux9b0ux9a3ux9afux9b0-ux9b8ux9b9ux99c-ux9a8ux9afux9ae}

সাধারণ ভগ্নাংশ অধ্যায়টি পাটিগণিতের খুবই গুরুত্বপূর্ণ অধ্যায়। এ অধ্যায়টি সম্পূর্ণ
নখদর্পণে নিতে হলে প্রথমেই আপনাকে সাধারণ ভগ্নাংশের যোগ,বিয়োগ,গুণ ও ভাগ নির্ণয়ে
সিদ্ধহস্ত হতে হবে। তাই এ অধ্যায়ের বিগত সালের প্রশ্ন সমাধান করার আগে সাধারণ
ভগ্নাংশের যোগ,বিয়োগ,গুণ ও ভাগ নির্ণয়ে পারদর্শী করে তোলার জন্য সাধারণ ভগ্নাংশের
যোগ,বিয়োগ,গুণ ও ভাগ এ টু জেড সুচারুরুপে তুলে ধরা হল। আসুন আগে সাধারণ ভগ্নাংশের
যোগ,বিয়োগ,গুণ ও ভাগ ভালভাবে জানি।

\begin{itemize}
\item
  \textbf{সাধারণ ভগ্নাংশ কাকে বলে ?}
\item
  সাধারণ ভগ্নাংশের যোগ ও বিয়োগ করার নিয়ম
\item
  সাধারণ ভগ্নাংশের গুণ ও ভাগ করার নিয়ম
\item
  বৃহত্তম ও ক্ষুদ্রতম ভগ্নাংশ নির্ণয়
\end{itemize}

\subsubsection{সাধারণ ভগ্নাংশ কাকে বলে
?}\label{ux9b8ux9a7ux9b0ux9a3-ux9adux997ux9a8ux9b6-ux995ux995-ux9acux9b2}

এ অধ্যায়টি শুরু করার আগে চলুন জেনে নেই সাধারণ ভগ্নাংশ কাকে বলে । যেহেতু যে কোন
পূর্ণ সংখ্যার ভগ্ন বা ভাঙ্গা অংশকে ভগ্নাংশ বলে সে সূত্রে বলা যায় যে ভগ্নাংশ হর ও
লব নিয়ে গঠিত হয় তাকেই সাধারণ ভগ্নাংশ বলে ।

\subsubsection{সাধারণ ভগ্নাংশের যোগ ও বিয়োগ করার
নিয়ম}\label{ux9b8ux9a7ux9b0ux9a3-ux9adux997ux9a8ux9b6ux9b0-ux9afux997-ux993-ux9acux9dfux997-ux995ux9b0ux9b0-ux9a8ux9dfux9ae}

১. \(\frac{\text{১}}{\text{৪}}\) + \(\frac{\text{৩}}{\text{৫}}\)

প্রথমেই ভগ্নাংশদ্বয়ের হর ৪ ও ৫ এর ল.সা.গু নিতে হবে। ৪ ও ৫ এর ল.সা.গু = ২০।
এখন প্রথম ভগ্নাংশের হর ৪ দিয়ে ২০ কে ভাগ করলে পাই ৫। এই ৫ কে আবার প্রথম
ভগ্নাংশের লব ১ দিয়ে গুণ করতে হবে। গুণ করলে পাই ৫। তেমনিভাবে দ্বিতীয় ভগ্নাংশের
হর ৫ দিয়ে ২০ কে ভাগ করলে পাই ৪।এই ৪ কে আবার দ্বিতীয় ভগ্নাংশের লব ৩ দিয়ে গুণ
করতে হবে ।গুণ করলে পাই ১২ ।এই প্রক্রিয়াটি নিচে লিখে দেখানো হল।

= \(\frac{\text{৫} + \text{১২}}{\text{২০}}\)

এবার উপরোক্ত ভগ্নাংশের লব ৫ ও ১২ যোগ করে পাই নিচের ভগ্নাংশটি।

= \(\frac{\text{১৭}}{\text{২০}}\)

এভাবেই ভগ্নাংশের যোগ করতে হবে।

২. \(\frac{\text{১}}{\text{৬}}\) + \(\frac{\text{৫}}{\text{১৬}}\)

প্রথমেই ভগ্নাংশদ্বয়ের হর ৬ ও ১৬ এর ল.সা.গু নিতে হবে।৬ ও ১৬ এর ল.সা.গু =
৪৮।এখন প্রথম ভগ্নাংশের হর ৬ দিয়ে ৪৮ কে ভাগ করলে পাই ৮ ।এই ৮ কে আবার প্রথম
ভগ্নাংশের লব ১ দিয়ে গুণ করতে হবে।গুণ করলে পাই ৮। তেমনিভাবে দ্বিতীয় ভগ্নাংশের
হর ১৬ দিয়ে ৪৮ কে ভাগ করলে পাই ৩।এই ৩ কে আবার দ্বিতীয় ভগ্নাংশের লব ৫ দিয়ে
গুণ করতে হবে ।গুণ করলে পাই ১৫ ।এই প্রক্রিয়াটি নিচে লিখে দেখানো হল।

= \(\frac{\text{৮} + \text{১৫}}{\text{৪৮}}\)

এবার উপরোক্ত ভগ্নাংশের লব ৮ ও ১৫ যোগ করে পাই নিচের ভগ্নাংশটি।

= \(\frac{\text{২৩}}{\text{৪৮}}\)

এভাবেই ভগ্নাংশের যোগ করতে হবে।

৩. \(\frac{\text{৩}}{\text{৪}}\) + \(\frac{\text{৫}}{\text{৮}}\) + \(\frac{\text{১৭}}{\text{১৬}}\)

প্রথমেই ভগ্নাংশদ্বয়ের হর ৪,৮ ও ১৬ এর ল.সা.গু নিতে হবে।৪,৮ ও ১৬ এর ল.সা.গু =
১৬।এখন প্রথম ভগ্নাংশের হর ৪ দিয়ে ১৬ কে ভাগ করলে পাই ৪ ।এই ৪ কে আবার প্রথম
ভগ্নাংশের লব ৩ দিয়ে গুণ করতে হবে।গুণ করলে পাই ১২। দ্বিতীয় ভগ্নাংশের হর ৮ দিয়ে
১৬ কে ভাগ করলে পাই ২।এই ২ কে আবার দ্বিতীয় ভগ্নাংশের লব ৫ দিয়ে গুণ করতে হবে
।গুণ করলে পাই ১০ ।তেমনিভাবে তৃতীয় ভগ্নাংশের হর ১৬ দিয়ে ১৬ কে ভাগ করলে পাই
১।এই ১ কে আবার তৃতীয় ভগ্নাংশের লব ১৭ দিয়ে গুণ করতে হবে ।গুণ করলে পাই ১৭। এই
প্রক্রিয়াটি নিচে লিখে দেখানো হল।

= \(\frac{\text{১২} + \text{১০} + \text{১৭}}{\text{১৬}}\)

এবার উপরোক্ত ভগ্নাংশের লব ১২,১০ ও ১৭ যোগ করে পাই নিচের ভগ্নাংশটি।

= \(\frac{\text{৩৯}}{\text{১৬}}\)

৪. ৫\(\frac{\text{১}}{\text{৩}}\) - ২\(\frac{\text{১}}{\text{৪}}\)

প্রথমেই উপরোক্ত মিশ্র ভগ্নাংশগুলোকে অপ্রকৃত ভগ্নাংশে রুপ দিতে হবে।এজন্য প্রতিটি
ভগ্নাংশের হর দিয়ে পূর্ণ সংখ্যাকে গুণ করতে হবে এবং তার সাথে প্রদত্ত ভগ্নাংশের লব
যোগ করে অপ্রকৃত ভগ্নাংশের লব তৈরি করতে হবে।যা নিচে দেখানো হল।

= \(\frac{\text{১৬}}{\text{৩}}\) - \(\frac{\text{৯}}{\text{৪}}\)

এখন ভগ্নাংশদ্বয়ের হর ৩ ও ৪ এর ল.সা.গু নিতে হবে।৩ ও ৪ এর ল.সা.গু = ১২ ।এখন
প্রথম ভগ্নাংশের হর ৩ দিয়ে ১২ কে ভাগ করলে পাই ৪ । এই ৪ কে আবার প্রথম
ভগ্নাংশের লব ১৬ দয়ে গুণ করতে হবে।গুণ করলে পাই ৬৪। তেমনিভাবে দ্বিতীয় ভগ্নাংশের
হর ৪ দিয়ে ১২ কে ভাগ করলে পাই ৩।এই ৩ কে আবার দ্বিতীয় ভগ্নাংশের লব ৯ দিয়ে গুণ
করতে হবে ।গুণ করলে পাই ২৭ ।এই প্রক্রিয়াটি নিচে লিখে দেখানো হল।

= \(\frac{\text{৬৪}\  - \ \text{২৭}}{\text{১২}}\)

এবার উপরোক্ত ভগ্নাংশের লব ৬৪ থেকে ২৭ বিয়োগ করে পাই নিচের ভগ্নাংশটি।

= \(\frac{\text{৩৭}}{\text{১২}}\)

\subsubsection{সাধারণ ভগ্নাংশের গুণ ও ভাগ করার
নিয়ম}\label{ux9b8ux9a7ux9b0ux9a3-ux9adux997ux9a8ux9b6ux9b0-ux997ux9a3-ux993-ux9adux997-ux995ux9b0ux9b0-ux9a8ux9dfux9ae}

\paragraph{সাধারণ ভগ্নাংশের
গুণ}\label{ux9b8ux9a7ux9b0ux9a3-ux9adux997ux9a8ux9b6ux9b0-ux997ux9a3}

১. \(\frac{\text{১}}{\text{৪}}\) × \(\frac{\text{৩}}{\text{৫}}\)

সাধারণ ভগ্নাংশের গুণ করার সময় প্রথমেই দেখতে হবে প্রদত্ত ভগ্নাংশগুলোর কোন হর/লব
দিয়ে কোন লব/হরকে ভাগ করা যায় কি না ? অথবা কোন একটি সংখ্যা দিয়ে হর ও লব
উভয়কে কাটাকাটি করা যায় কিনা? যেমন উপরের ভগ্নাংশে কোন ভাগ অথবা কাটাকাটি
করা যায় না।তাই প্রদত্ত ভগ্নাংশগুলোর হরকে হর দিয়ে গুণ এবং লবকে লব দিয়ে গুণ করে
লবগুলোর গুণফল লবের জায়গায় এবং হরগুলোর গুণফল হরের জায়গায় বসাতে হবে। তাই উপরের
ভগ্নাংশদ্বয়ের গুণফল = \(\frac{\text{৩}}{\text{২০}}\) ।

= ৫\(\frac{\text{১}}{\text{৩}}\) × ২\(\frac{\text{১}}{\text{৪}}\)

প্রথমেই উপরোক্ত মিশ্র ভগ্নাংশগুলোকে অপ্রকৃত ভগ্নাংশে রুপ দিতে হবে।এজন্য প্রতিটি
ভগ্নাংশের হর দিয়ে পূর্ণ (সমস্ত) সংখ্যাকে গুণ করতে হবে এবং তার সাথে প্রদত্ত
ভগ্নাংশের লব যোগ করে অপ্রকৃত ভগ্নাংশের লব তৈরি করতে হবে।যা নিচে দেখানো হল।

= \(\frac{\text{১৬}}{\text{৩}}\) × \(\frac{\text{৯}}{\text{৪}}\)

এখন কাটাকাটি করতে হবে।দেখা যাচ্ছে হর ৩ দিয়ে লব ৯ কে ভাগ করা যায় । ৩ দিয়ে
৯ কে কাটলে হয় ৩।আবার হর ৪ দিয়ে ১৬ কে ভাগ করা যায়।৪ দিয়ে ১৬ কে কাটলে হয়
৪।তাহলে হরের ঘরে আর কোন সংখ্যাই রইল না।লবের ঘরে থাকল ৩ এবং ৪ ।৩ এবং ৪ এর
গুণফল ১২।তাহলে ভগ্নাংশদ্বয়ের গুণফল দাড়ায় = ১২।

৩. \(\frac{\text{৩}}{\text{৪}}\) × \(\frac{\text{৫}}{\text{৮}}\) × \(\frac{\text{৮}}{\text{১৫}}\)

এখন কাটাকাটি করতে হবে।দেখা যাচ্ছে হর ৮ দিয়ে লব ৮ কে ভাগ করা যায় । ৮ দিয়ে
৮ কে কাটলে হয় ১ ।আবার লব ৫ দিয়ে ১৫ কে ভাগ করা যায়।৫ দিয়ে ১৫ কে কাটলে হয়
৩। এই ৩ দিয়ে আবার লব ৩ কে কাটা যায় ।তাহলে হরের ঘরে শুধুমাত্র ৪ অবশিষ্ট রইল।
লবের ঘরে আর কোন সংখ্যাই রইল না।তাহলে ভগ্নাংশগুলোর গুণফল দাড়ায় =
\(\frac{\text{১}}{\text{৪}}\) ।

\paragraph{সাধারণ ভগ্নাংশের
ভাগ}\label{ux9b8ux9a7ux9b0ux9a3-ux9adux997ux9a8ux9b6ux9b0-ux9adux997}

১. \(\frac{\text{১}}{\text{৪}}\) ÷ \(\frac{\text{৩}}{\text{৫}}\)

সাধারণ ভগ্নাংশের ভাগ সাধারণ ভগ্নাংশের গুণের মতই।এখানে পার্থক্য শুধু ভাগকে গুণে
রুপান্তর করে পরবর্তী ভগ্নাংশকে বিপরীত ভগ্নাংশে রুপ দিতে হবে।যা নিচে দেখানো হল।

= \(\frac{\text{১}}{\text{৪}}\) × \(\frac{\text{৫}}{\text{৩}}\)

= \(\frac{\text{৫}}{\text{১২}}\)

\subsubsection{বৃহত্তম ও ক্ষুদ্রতম ভগ্নাংশ
নির্ণয়}\label{ux9acux9b9ux9a4ux9a4ux9ae-ux993-ux995ux9b7ux9a6ux9b0ux9a4ux9ae-ux9adux997ux9a8ux9b6-ux9a8ux9b0ux9a3ux9df}

বিসিএসসহ অন্যান্য প্রতিযোগিতামূলক পরীক্ষায় এরকম একটি প্রশ্ন প্রায়ই এসে থাকে। খুব
সহজে এবং দ্রুতগতিতে ভগ্নাংশের ক্ষুদ্রতম ও বৃহত্তম বের করার নিয়ম জানার জন্য নিচের
আলোচনাটি মনোযোগ দিয়ে পড়ুন।আশা করি উপকৃত হবেন।

\textbf{দুটি নিয়মে বৃহত্তম ও ক্ষুদ্রতম ভগ্নাংশ নির্ণয় করা যায় । যেমন:}

\textbf{১. দশমিক মান বের করে ভগ্নাংশের তুলনার মাধ্যমে:}

এক্ষেত্রে ভগ্নাংশগুলোর হর দিয়ে লবকে ভাগ করে ভাগফলটির দুই বা তিন দশমিক মান ঐ
সংখ্যাটির পাশে বসিয়ে সবগুলো ভগ্নাংশ তুলনার মাধ্যমে বের করতে হয় কোনটি বৃহত্তম বা
কোনটি ক্ষুদ্রতম ভগ্নাংশ। যেমন: \(\frac{\text{১}}{\text{২}}\) = ০.৫০ , \(\frac{\text{১}}{\text{৩}}\) =
০.৩৩ , \(\frac{\text{১}}{\text{৪}}\) = ০.২৫ , \(\frac{\text{১}}{\text{৫}}\) = ০.২০ । এখানে বৃহত্তম
ভগ্নাংশ হলো \(\frac{\text{১}}{\text{২}}\) এবং ক্ষুদ্রতম ভগ্নাংশ হলো \(\frac{\text{১}}{\text{৫}}\) ।

\textbf{পরীক্ষার হলে ক্যালকুলেটর ব্যবহারের সুযোগ থাকলে এই পদ্ধতি প্রয়োগ করা হবে
বুদ্ধিমানের কাজ।}

\textbf{২. আড়াআড়ি গুণ করার মাধ্যমে ভগ্নাংশের তুলনার মাধ্যমে:}

প্রথমেই কিছু ভগ্নাংশ কিভাবে নিজে থেকেই ছোট বড় বের করতে পারবেন তা দেখে নিন:

কোন ভগ্নাংশের লবগুলো একই হলে যে ভগ্নাংশের হর ছোট সে ভগ্নাংশটিই বড় , এবং যে
ভগ্নাংশটির হর বড় সে ভগ্নাংশটিই ছোট।যেমন: \(\frac{\text{১}}{\text{২}}\) , \(\frac{\text{১}}{\text{৩}}\)
, \(\frac{\text{১}}{\text{৪}}\) , \(\frac{\text{১}}{\text{৫}}\) এর মধ্যে \(\frac{\text{১}}{\text{২}}\) বৃহত্তম এবং
\(\frac{\text{১}}{\text{৫}}\) ক্ষুদ্রতম । কারণ ১ টাকা কে ২ ভাগ করলে ৫০ পয়সা হবে কিন্তু ঐ ১
টাকাকেই ৫ ভাগ করলে তা ২০ পয়সা হবে তাই \(\frac{\text{১}}{\text{২}}\) , \(\frac{\text{১}}{\text{৫}}\)
থেকে বড় ।

আবার কোন ভগ্নাংশের হর একই হলে যে ভগ্নাংশের লব বড় সে ভগ্নাংশটি বড় এবং যে
ভগ্নাংশের লব ছোট সে ভগ্নাংশটি ছোট । যেমন: \(\frac{\text{৭}}{\text{৯}}\) ও
\(\frac{\text{৫}}{\text{৯}}\) এর মধ্যে অবশ্যই \(\frac{\text{৭}}{\text{৯}}\) বড় । কারণ ৭ টাকাকে ৯ ভাগ
করলে যাই হোক ৫ টাকাকে ৯ ভাগ করলে তার থেকে ছোট হবে।

আবার কখনো লব বা হর কোনটাই না মিললে সমতুল ভগ্নাংশে রুপ দিয়ে মিলিয়ে নিতে হবে।
কিভাবে দেখুন:

\(\frac{\text{৭}}{\text{৯}}\) এবং \(\frac{\text{৫}}{\text{১৮}}\) এর মধ্যে কোনটি বড় অথবা কোনটি
ছোট।এক্ষেত্রে প্রথম ভগ্নাংশকে দ্বিগুণ করে পরেরটার সাথে মিলাতে হবে।অর্থাৎ পরের
ভগ্নাংশের হরের সমান হর বিশিষ্ট ভগ্নাংশে রুপ দিতে হবে প্রথম ভগ্নাংশকে। ভগ্নাংশকে
দ্বিগুণ,তিনগুণ এরুপ করলে ভগ্নাংশের মানের কোন পরিবর্তন হয় না। যেমন:
\(\frac{\text{১৪}}{\text{১৮}}\) \textgreater{} \(\frac{\text{৫}}{\text{১৮}}\) সুতরাং \(\frac{\text{৭}}{\text{৯}}\)
ই বড় ভগ্নাংশ ।

\textbf{চারটি অপশনের মাঝে ভগ্নাংশের ছোট বড় নির্ণয়ের সহজ পদ্ধতি}

ধরি, ভগ্নাংশ চারটি যথাক্রমে (ক) \(\frac{\text{৩}}{\text{৫}}\) , (খ) \(\frac{\text{৫}}{\text{৮}}\) ,
(গ) \(\frac{\text{৬}}{\text{১১}}\) , (ঘ) \(\frac{\text{৮}}{\text{১৪}}\)

১ম ভগ্নাংশের লবের সঙ্গে ২য় ভগ্নাংশের হর গুণ করে প্রাপ্ত গুণফলটি ১ম ভগ্নাংশের মান
হিসেবে বিবেচনা করতে হবে।

আবার, ২য় ভগ্নাংশের লবের সঙ্গে ১ম ভগ্নাংশের হর গুণ করে প্রাপ্ত গুণফলটি ২য়
ভগ্নাংশের মান হিসেবে বিবেচনা করতে হবে।

এখানে, যে ভগ্নাংশের মানটি বড় হবে, সেই ভগ্নাংশটি বড় এবং সেই ভগ্নাংশের মান
ছোট হবে, সেই ভগ্নাংশটি ছোট হবে।

\textbf{চারটি অপশনের মাঝে তুলনা করার সময় বৃহত্তম ভগ্নাংশ বের করতে বলা হলে
ক্ষুদ্রতমগুলো এক এক করে বাদ দিতে হবে, আবার ক্ষুদ্রতম বের করতে বলা হলে বৃহত্তমগুলো
এক এক করে বাদ দিতে হবে ।}

ধরি, উপরোক্ত অপশনের বড় ভগ্নাংশ বের করতে হবে।

(ক) ও (খ) এর মধ্যে \(\frac{\text{৩}}{\text{৫}}\) ↔ \(\frac{\text{৫}}{\text{৮}}\) , ৩×৮ = ২৪, ৫×৫ = ২৫
। সুতরাং এখানে (খ) বড় । তাই (ক) বড় ভগ্নাংশ বিবেচনা থেকে বাদ।

(খ) ও (গ )এর মধ্যে \(\frac{\text{৫}}{\text{৮}}\) ↔ \(\frac{\text{৬}}{\text{১১}}\), ৫×১১ = ৫৫, ৮×৬ =
৪৮। সুতরাং এখানেও (খ ) বড় । তাই (ক) ও (গ) বড় ভগ্নাংশ বিবেচনা থেকে বাদ।

(খ) ও (ঘ) এর মধ্যে \(\frac{\text{৫}}{\text{৮}}\)↔ \(\frac{\text{৮}}{\text{১৪}}\), ৫×১৪ = ৭০, ৮×৮ =
৮৮। সুতরাং এখানেও আবার (খ) বড় । কিন্তু (ক) ও (গ) অপশন আমরা আগেই বাদ
দিয়েছিলাম।তাই এখন (ক),(গ),(ঘ) এ তিনটি বড় ভগ্নাংশ বিবেচনা থেকে বাদ।কাজেই
সবচেয়ে বড় (খ) অপশন ।

সুতরাং \(\frac{\text{৫}}{\text{৮}}\) প্রদত্ত ভগ্নাংশগুলোর মধ্যে সবচেয়ে বড় ।

এবার আরেকটি দেখুন:

ধরি, ভগ্নাংশ চারটি যথাক্রমে (ক) \(\frac{\text{২}}{\text{১১}}\) , (খ) \(\frac{\text{৩}}{\text{১১}}\) ,
(গ) \(\frac{\text{২}}{\text{১৩}}\) , (ঘ) \(\frac{\text{৪}}{\text{১৫}}\)

ধরি, উপরোক্ত অপশনের ছোট ভগ্নাংশ বের করতে হবে।

প্রথমেই আমরা ক ও খ অপশনের মাঝে একটিকে বাদ দিয়ে দিতে পারি কেননা দুটি সমহর
বিশিষ্ট ।এটা সহজেই বোঝা যায় \(\frac{\text{৩}}{\text{১১}}\) বড় । যেহেতু আমাদের ছোট ভগ্নাংশ
বের করতে হবে তাই বড়গুলো বাদ দিতে হবে। তাই অপশন (খ) ছোট ভগ্নাংশ বিবেচনা
থেকে বাদ ।তাই এখন (ক) অপশনকে নিয়ে আসব (গ) অপশনের সাথে তুলনা করার
জন্য।এখানে উভয়টির লব ২ । তাই (ক) অপশন এখানে বড়। তাই এখানে আমাদের ছোট
ভগ্নাংশ হিসেবে হাতে থাকল অপশন (গ) ।

এবার গ ও ঘ এর মধ্যে \(\frac{\text{২}}{\text{১৩}}\)↔ \(\frac{\text{৪}}{\text{১৫}}\), ২×১৫ = ৩০, ৪×১৩ =
৫২। সুতরাং এখানেও আবার গ ছোট ।সুতরাং \(\frac{\text{২}}{\text{১৩}}\) প্রদত্ত ভগ্নাংশগুলোর
মধ্যে সবচেয়ে ছোট ।

সময় অবশিষ্ট আছে:

\section{63. সাধারণ ভগ্নাংশ MCQ প্রশ্ন উত্তর
অনুশীলন}\label{ux9b8ux9a7ux9b0ux9a3-ux9adux997ux9a8ux9b6-mcq-ux9aaux9b0ux9b6ux9a8-ux989ux9a4ux9a4ux9b0-ux985ux9a8ux9b6ux9b2ux9a8}

\textbf{অপ্রকৃত ভগ্নাংশ কোনটি?}

ক) (৩/৪)

খ) (৯/১৩)

\textbf{\hl{গ) (৭/৪) ✔}}

ঘ) সবগুলো

\textbf{\hl{✔ সঠিক উত্তর: গ) (৭/৪)}}

\begin{quote}
\textbf{ব্যাখ্যা:} \emph{যে ভগ্নাংশের লব বড় ও হর ছোট তাকে অপ্রকৃত ভগ্নাংশ বলে।
এখানে (৭/৪) ভগ্নাংশটির লব (৭) হর (৪) অপেক্ষা বড়।}

\emph{আর যে ভগ্নাংশের লব ছোট ও হর বড় তাকে প্রকৃত ভগ্নাংশ বলে।}
\end{quote}

\textbf{কোনটি মিশ্র ভগ্নাংশ?}

ক) (১/২)

খ) (৩/২)

\textbf{\hl{গ) ৭(১/৪) ✔}}

ঘ) কোনোটিই নয়

\textbf{\hl{✔ সঠিক উত্তর: গ) ৭(১/৪)}}

\begin{quote}
\textbf{ব্যাখ্যা:} \emph{যে ভগ্নাংশে পূর্ণ সংখ্যার সাথে প্রকৃত ভগ্নাংশ যুক্ত থাকে,
তাকে মিশ্র ভগ্নাংশ বলে। এখানে ৭(১/৪) হলো একটি মিশ্র ভগ্নাংশ।}
\end{quote}

\textbf{(০/০) = কত?}

ক) ১

খ) ০

গ) ০.০

\textbf{\hl{ঘ) অনির্ণেয় ✔}}

\textbf{\hl{✔ সঠিক উত্তর: ঘ) অনির্ণেয়}}

\begin{quote}
\textbf{ব্যাখ্যা:} \emph{(০/০) = অনির্ণেয়। কারণ কোনো সংখ্যাকে ০ দ্বারা ভাগ
করা যায় না।}
\end{quote}

\textbf{√(০.০৪) এর মান কত?}

\textbf{\hl{ক) ০.২ ✔}}

খ) ০.০২

গ) ০.০০২

ঘ) ০.০০০২

\textbf{\hl{✔ সঠিক উত্তর: ক) ০.২}}

\begin{quote}
\textbf{ব্যাখ্যা:} \emph{√(০.০৪)=√((৪/১০০))=√((১/২৫))=(১/৫)=০.২}
\end{quote}

\textbf{০.৪৭ কে সাধারণ ভগ্নাংশে পরিণত করলে কত হবে?}

ক) (৪৭/১০)

খ) (৪০/১০)

গ) (৪৭/৯৯)

\textbf{\hl{ঘ) (৪৭/১০০) ✔}}

\textbf{\hl{✔ সঠিক উত্তর: ঘ) (৪৭/১০০)}}

\begin{quote}
\textbf{ব্যাখ্যা:} \emph{০.৪৭=(৪৭/১০০)}
\end{quote}

\textbf{০.০২২• এর সাধারণ ভগ্নাংশ কোনটি?}

ক) (২১/৯৯৯)

খ) (২২/১০০০)

গ) (২১/৯০০)

\textbf{\hl{ঘ) (২০/৯০০) ✔}}

\textbf{\hl{✔ সঠিক উত্তর: ঘ) (২০/৯০০)}}

\begin{quote}
\textbf{ব্যাখ্যা:} \emph{দেয়া আছে, ০.০২২•=(২২−২/৯০০)=(২০/৯০০)}
\end{quote}

\textbf{৪.৩• কে সাধারণ ভগ্নাংশে প্রকাশ করলে নিচের কোনটি হবে?}

ক) (৩/১১)

খ) (১১/৩)

গ) (৩/১৩)

\textbf{\hl{ঘ) (১৩/৩) ✔}}

\textbf{\hl{✔ সঠিক উত্তর: ঘ) (১৩/৩)}}

\begin{quote}
\textbf{ব্যাখ্যা:} \emph{৪.৩•=(৪৩−৪/৯)=(৩৯/৯)=(১৩/৩)}
\end{quote}

\textbf{০.২•৪• এর সমান ভগ্নাংশটি}

ক) (১/৩)

\textbf{\hl{খ) (৮/৩৩) ✔}}

গ) (৭/২৪)

ঘ) (৭/৭২)

\textbf{\hl{✔ সঠিক উত্তর: খ) (৮/৩৩)}}

\begin{quote}
\textbf{ব্যাখ্যা:} \emph{০.২•৪•=(২৪−০/৯৯)=(২৪/৯৯)=(৮/৩৩)}
\end{quote}

\textbf{নিচের সংখ্যাগুলোর মধ্যে কোনটি ক্ষুদ্রতম?}

\textbf{\hl{ক) ০.৩ ✔}}

খ) √(০.৩)

গ) (২/৩)

ঘ) (২/৫)

\textbf{\hl{✔ সঠিক উত্তর: ক) ০.৩}}

\begin{quote}
\textbf{ব্যাখ্যা:} \emph{এখানে,}

\emph{√(০.৩)≈০.৫৪৭}

\emph{(২/৩)=০.৬৬}

\emph{(২/৫)=০.৪০}

\emph{∴ দেখা যাচ্ছে, ০.৩-ই ক্ষুদ্রতম সংখ্যা।}
\end{quote}

\textbf{ভগ্নাংশগুলোর মধ্যে কোনটি সবচেয়ে বড়?}

ক) (2/3)

\textbf{\hl{খ) (4/5) ✔}}

গ) (13/315)

ঘ) (23/30)

\textbf{\hl{✔ সঠিক উত্তর: খ) (4/5)}}

\begin{quote}
\textbf{ব্যাখ্যা:} \emph{এখানে প্রদত্ত ভগ্নাংশগুলোর দশমিক মান হিসাব করলে:}

\emph{⇒ (2/3)=0.666}

\emph{⇒ (4/5)=0.8}

\emph{⇒ (13/315)=0.0412}

\emph{⇒ (23/30)=0.766}

\emph{∴ (4/5) বা খ অপশনটি সবচেয়ে বড়।}
\end{quote}

\textbf{নিচের কোন সংখ্যাটি বৃহত্তম?}

ক) (2/20)

\textbf{\hl{খ) (3/5) ✔}}

গ) (4/15)

ঘ) (7/25)

\textbf{\hl{✔ সঠিক উত্তর: খ) (3/5)}}

\begin{quote}
\textbf{ব্যাখ্যা:} \emph{এখানে, অপশনগুলোর মান বিশ্লেষণ করে পাই:}

\emph{⇒ (2/20)=0.1}

\emph{⇒ (3/5)=0.6}

\emph{⇒ (4/15)≈0.2667}

\emph{⇒ (7/25)=0.28}

\emph{∴ (3/5) হলো সবচেয়ে বড় সংখ্যা।}
\end{quote}

\textbf{নিচের ভগ্নাংশগুলোর মধ্যে কোনটি বৃহত্তম?}

ক) (3/4)

খ) (4/7)

\textbf{\hl{গ) (6/7) ✔}}

ঘ) (7/9)

\textbf{\hl{✔ সঠিক উত্তর: গ) (6/7)}}

\begin{quote}
\textbf{ব্যাখ্যা:} \emph{ভগ্নাংশগুলোর দশমিক রূপ:}

\emph{⇒ (3/4)=0.75}

\emph{⇒ (4/7)≈0.571}

\emph{⇒ (6/7)≈0.857}

\emph{⇒ (7/9)≈0.778}

\emph{∴ (6/7) সবচেয়ে বড়।}
\end{quote}

\textbf{নিচে উল্লেখিত ভগ্নাংশগুলোর মধ্যে কোনটি সবচেয়ে বেশি?}

ক) (1/25)

খ) (1/19)

গ) (1/15)

\textbf{\hl{ঘ) (1/12) ✔}}

\textbf{\hl{✔ সঠিক উত্তর: ঘ) (1/12)}}

\begin{quote}
\textbf{ব্যাখ্যা:} \emph{যেহেতু ভগ্নাংশগুলোর লব একই (১), তাই যে ভগ্নাংশের হর
ছোট সেটি সবচেয়ে বড় বা বেশি হবে।}

\emph{∴ (1/12) প্রদত্ত ভগ্নাংশগুলোর মধ্যে সবচেয়ে বড়।}
\end{quote}

\textbf{নিচের কোন সংখ্যাটি ক্ষুদ্রতম?}

ক) (1/3)

খ) (2/7)

\textbf{\hl{গ) (5/21) ✔}}

ঘ) (3/6)

\textbf{\hl{✔ সঠিক উত্তর: গ) (5/21)}}

\begin{quote}
\textbf{ব্যাখ্যা:} \emph{দশমিক মানে রূপান্তর করে পাই:}

\emph{⇒ (1/3)=0.333}

\emph{⇒ (2/2)≈0.285}

\emph{⇒ (5/21)≈0.238}

\emph{⇒ (3/6)≈0.5}

\emph{∴ (5/21) হলো ক্ষুদ্রতম সংখ্যা।}
\end{quote}

\textbf{নিচের ভগ্নাংশগুলোর মধ্যে কোনটি বৃহত্তম?}

ক) (2/3)

খ) (3/4)

\textbf{\hl{গ) (4/5) ✔}}

ঘ) (5/7)

\textbf{\hl{✔ সঠিক উত্তর: গ) (4/5)}}

\begin{quote}
\textbf{ব্যাখ্যা:} \emph{ভগ্নাংশগুলোর দশমিক মান তুলনা করলে:}

\emph{⇒ (2/3)=0.666}

\emph{⇒ (3/4)=0.75}

\emph{⇒ (4/5)=0.8}

\emph{⇒ (5/7)≈0.714}

\emph{∴ (4/5) হলো বৃহত্তম ভগ্নাংশ।}
\end{quote}

\textbf{নিচের ভগ্নাংশগুলোর মধ্যে কোনটি ক্ষুদ্রতম?}

ক) (3/5)

\textbf{\hl{খ) (3/7) ✔}}

গ) (4/9)

ঘ) (2/3)

\textbf{\hl{✔ সঠিক উত্তর: খ) (3/7)}}

\begin{quote}
\textbf{ব্যাখ্যা:} \emph{দশমিক মানে প্রকাশ করে পাই:}

\emph{⇒ (3/5)=0.6}

\emph{⇒ (3/7)≈0.428}

\emph{⇒ (4/9)=0.444}

\emph{⇒ (2/3)=0.666}

\emph{∴ (3/7) হলো ক্ষুদ্রতম ভগ্নাংশ।}
\end{quote}

\textbf{সবচেয়ে বড় সংখ্যা কোনটি?}

ক) ০.০০৯৯

\textbf{\hl{খ) ০.১০০ ✔}}

গ) (৯/১০০)

ঘ) (৯/১০০০)

\textbf{\hl{✔ সঠিক উত্তর: খ) ০.১০০}}

\begin{quote}
\textbf{ব্যাখ্যা:} \emph{ক. ০.০০৯৯ = (৯৯/১০০০০)}

\emph{খ. ০.১০০ = (১০০০/১০০০০) (বৃহত্তম)}

\emph{গ. (৯/১০০)=০.০৯=(৯০০/১০০০০)}

\emph{ঘ. (৯/১০০০)=০.০০৯=(৯০/১০০০০)}

\emph{∴ এদের মধ্যে ০.১-ই সবচেয়ে বড়। তাই এটি উত্তর হবে।}
\end{quote}

\textbf{√(০.০৯)= কত?}

ক) ০৩

\textbf{\hl{খ) ০.৩ ✔}}

গ) ০.০০৩

ঘ) ০.০০০৩

\textbf{\hl{✔ সঠিক উত্তর: খ) ০.৩}}

\begin{quote}
\textbf{ব্যাখ্যা:} \emph{√(০.০৯)=√((৯/১০০))=(৩/১০)=০.৩}
\end{quote}

\textbf{৩৩(১/৩)\% এর সমান ভগ্নাংশ কত হবে?}

ক) (১০/৩)

খ) (৩/১০)

\textbf{\hl{গ) (১/৩) ✔}}

ঘ) (৩/৭)

\textbf{\hl{✔ সঠিক উত্তর: গ) (১/৩)}}

\begin{quote}
\textbf{ব্যাখ্যা:} \emph{৩৩(১/৩)\%=(১০০/৩)×(১/১০০)=(১/৩)}
\end{quote}

\textbf{কোন ভগ্নাংশটি ০.৫\% এর সমতুল্য?}

ক) (১/২০)

\textbf{\hl{খ) (১/২০০) ✔}}

গ) (১/৫০)

ঘ) (১/৫)

\textbf{\hl{✔ সঠিক উত্তর: খ) (১/২০০)}}

\begin{quote}
\textbf{ব্যাখ্যা:} \emph{০.৫\%=(০.৫/১০০)=(৫/১০০×১০)=(১/২০০)}
\end{quote}

\textbf{০.২• এর সামান্য ভগ্নাংশ কোনটি?}

ক) (২/১০)

খ) (২/৮)

\textbf{\hl{গ) (২/৯) ✔}}

ঘ) (২/৯৯)

\textbf{\hl{✔ সঠিক উত্তর: গ) (২/৯)}}

\begin{quote}
\textbf{ব্যাখ্যা:} \emph{০.২•=(২−০/৯)=(২/৯)}
\end{quote}

\textbf{৩.০০০১০+৫×১০\^{}−৩= কত?}

\textbf{\hl{ক) ৩.০০৫১০ ✔}}

খ) ৩.০৫০১০

গ) ৩.০০০১৫

ঘ) ৩.০০০৬০

\textbf{\hl{✔ সঠিক উত্তর: ক) ৩.০০৫১০}}

\begin{quote}
\textbf{ব্যাখ্যা:} \emph{৩.০০০১০+৫×১০\^{}−৩}

\emph{= ৩.০০০১০+(৫/১০০০) = ৩.০০০১০+০.০০৫ = ৩.০০০১০+০.০০৫০০=৩.০০৫১০}
\end{quote}

\textbf{৩ ÷ (৩/২) = ?}

ক) ১

\textbf{\hl{খ) ২ ✔}}

গ) ৩

ঘ) (৯/২)

\textbf{\hl{✔ সঠিক উত্তর: খ) ২}}

\begin{quote}
\textbf{ব্যাখ্যা:} \emph{দেয়া আছে, ৩ ÷ (৩/২) = ৩ × (২/৩) = ২}
\end{quote}

\textbf{০.০২৪ × ১০\^{}৬ = ?}

\textbf{\hl{ক) ২৪০০০ ✔}}

খ) ২৪০০০০

গ) ২৪০০০০০

ঘ) কোনোটিই নয়

\textbf{\hl{✔ সঠিক উত্তর: ক) ২৪০০০}}

\begin{quote}
\textbf{ব্যাখ্যা:} \emph{০.০২৪ × ১০\^{}৬ = ০.০২৪ × ১০০০০০০ = ২৪,০০০}
\end{quote}

\textbf{কোন সংখ্যার (৪/৭) অংশ ৮০ এর সমান?}

ক) ১৬০

\textbf{\hl{খ) ১৪০ ✔}}

গ) ২৪০

ঘ) ২৪৭

\textbf{\hl{✔ সঠিক উত্তর: খ) ১৪০}}

\begin{quote}
\textbf{ব্যাখ্যা:} \emph{ধরি, সংখ্যাটি x}

\emph{প্রশ্নমতে, x এর (৪/৭)=৮০}

\emph{⇒ ৪x = ৮০ × ৭}

\emph{⇒ x = (৮০×৭/৪)}

\emph{∴ x = ১৪০}
\end{quote}

\textbf{১(১/২) ইঞ্চি ১ ফুটের কত অংশ?}

ক) (১/৫) অংশ

খ) (১/২) অংশ

\textbf{\hl{গ) (১/৮) অংশ ✔}}

ঘ) (১/১০) অংশ

\textbf{\hl{✔ সঠিক উত্তর: গ) (১/৮) অংশ}}

\begin{quote}
\textbf{ব্যাখ্যা:} \emph{১ ফুট = ১২ ইঞ্চি}

\emph{∴ নির্ণেয় অংশ = (১(১/২)/১২)=((৩/২)/১২)=(৩/২)×(১/১২)=(১/৮)}
\end{quote}

\textbf{১ ঘণ্টা ২০ মিনিট ৪ ঘণ্টার কত অংশ?}

\textbf{\hl{ক) (১/৩) অংশ ✔}}

খ) (২/৩) অংশ

গ) (১/৪) অংশ

ঘ) (৩/৪) অংশ

\textbf{\hl{✔ সঠিক উত্তর: ক) (১/৩) অংশ}}

\begin{quote}
\textbf{ব্যাখ্যা:} \emph{১ ঘণ্টা ২০ মিনিট = ৮০ মিনিট}

\emph{এবং ৪ ঘণ্টা = ২৪০ মিনিট}

\emph{∴ অংশ = (৮০/২৪০)=(১/৩)}
\end{quote}

\textbf{২০০ এর (১/২)\% এর সাথে ১০০ যোগ করলে সংখ্যাটি কত হবে?}

\textbf{\hl{ক) ১০১ ✔}}

খ) ১৫০

গ) ২০০

ঘ) ২০১

\textbf{\hl{✔ সঠিক উত্তর: ক) ১০১}}

\begin{quote}
\textbf{ব্যাখ্যা:} \emph{২০০ এর (১/২)\%=২০০×(১/২×১০০)=১}

\emph{এর সাথে ১০০ যোগ করলে হবে = (১০০ + ১) = ১০১}
\end{quote}

\textbf{(১/৮০) থেকে (১/৪০) বিয়োগ করে বিয়োগফলকে দশমিকে প্রকাশ করলে কত হয়?}

ক) ০.২২৫

\textbf{\hl{খ) −০.০১২৫ ✔}}

গ) ০.২৫

ঘ) ০.০৫

\textbf{\hl{✔ সঠিক উত্তর: খ) −০.০১২৫}}

\begin{quote}
\textbf{ব্যাখ্যা:} \emph{(১/৮০)−(১/৪০)=(১−২/৮০)=−(১/৮০)=−০.০১২৫}
\end{quote}

\textbf{কোন ভগ্নাংশটি লঘিষ্ঠ আকারে প্রকাশিত?}

ক) (৭৭/১৪৩)

খ) (১০২/২৮৯)

\textbf{\hl{গ) (১১৩/৩৫৫) ✔}}

ঘ) (৩৪৩/১০০১)

\textbf{\hl{✔ সঠিক উত্তর: গ) (১১৩/৩৫৫)}}

\begin{quote}
\textbf{ব্যাখ্যা:} \emph{যদি কোনো ভগ্নাংশের লব ও হরের মধ্যে ১ ব্যতীত অন্য কোনো
সাধারণ উৎপাদক বা গুণনীয়ক না থাকে, তবে ওই ভগ্নাংশকে লঘিষ্ঠ আকারে প্রকাশিত বলা
হয়।}

\emph{এখানে,}

\emph{৭৭/১৪৩ = (১১ × ৭) / (১১ × ১৩)}

\emph{১০২/২৮৯ = (১৭ × ৬) / (১৭ × ১৭)}

\emph{৩৪৩/১০০১ = (৭ × ৪৯) / (৭ × ১৪৩)}

\emph{কিন্তু (১১৩/৩৫৫) ভগ্নাংশটির লব ও হরের মধ্যে কোনো সাধারণ উৎপাদক নেই। ∴
এটি লঘিষ্ঠ আকারে প্রকাশিত।}
\end{quote}

\textbf{নিচের কোন ভগ্নাংশটি বৃহত্তম?}

ক) (৬/১১)

খ) (৮/১৪)

গ) (৩/৫)

\textbf{\hl{ঘ) (৫/৮) ✔}}

\textbf{\hl{✔ সঠিক উত্তর: ঘ) (৫/৮)}}

\begin{quote}
\textbf{ব্যাখ্যা:} \emph{ক. (৬/১১)=০.৫৪৫}

\emph{খ. (৮/১৪)=০.৫৭১}

\emph{গ. (৩/৫)=০.৬}

\emph{ঘ. (৫/৮)=০.৬২৫}

\emph{∴ বৃহত্তম ভগ্নাংশ (৫/৮)}
\end{quote}

\textbf{কোনটি ক্ষুদ্রতম সংখ্যা?}

\textbf{\hl{ক) (৩/৪) ✔}}

খ) (৪/৫)

গ) (৬/৭)

ঘ) (৭/৮)

\textbf{\hl{✔ সঠিক উত্তর: ক) (৩/৪)}}

\begin{quote}
\textbf{ব্যাখ্যা:} \emph{ক. (৩/৪)=০.৭৫}

\emph{খ. (৪/৫)=০.৮০}

\emph{গ. (৬/৭)=০.৮৫৭}

\emph{ঘ. (৭/৮)=০.৮৭৫}

\emph{∴ দেখা যাচ্ছে (৩/৪) -ই ক্ষুদ্রতম সংখ্যা।}
\end{quote}

\textbf{নিচের কোন ভগ্নাংশটি সবচেয়ে বড়?}

ক) (২/৩)

খ) (৩/৫)

\textbf{\hl{গ) (৮/১১) ✔}}

ঘ) (১৭/২৭)

\textbf{\hl{✔ সঠিক উত্তর: গ) (৮/১১)}}

\begin{quote}
\textbf{ব্যাখ্যা:} \emph{ক. (২/৩)=০.৬৬৬}

\emph{খ. (৩/৫)=০.৬}

\emph{গ. (৮/১১)=০.৭২৭২}

\emph{ঘ. (১৭/২৭)=০.৬২৯৬}

\emph{∴ এদের মধ্যে (৮/১১) সবচেয়ে বড়।}
\end{quote}

\textbf{কোন ভগ্নাংশটি ক্ষুদ্রতম?}

\textbf{\hl{ক) (১১/১৪) ✔}}

খ) (৫/৬)

গ) (১২/১৫)

ঘ) (১৭/২১)

\textbf{\hl{✔ সঠিক উত্তর: ক) (১১/১৪)}}

\begin{quote}
\textbf{ব্যাখ্যা:} \emph{ভগ্নাংশের হরগুলো ১৪, ৬, ১৫ ও ২১ এর ল.সা.গু. ২১০।}

\emph{(১১/১৪)=(১১×১৫/১৪×১৫)=(১৬৫/২১০)}

\emph{(৫/৬)=(৫×৩৫/৬×৩৫)=(১৭৫/২১০)}

\emph{(১২/১৫)=(১২×১৪/১৫×১৪)=(১৬৮/২১০)}

\emph{(১৭/২১)=(১৭×১০/২১×১০)=(১৭০/২১০)}

\emph{এখানে,
(১৬৫/২১০)\textless(১৬৮/২১০)\textless(১৭০/২১০)\textless(১৭৫/২১০)}

\emph{∴ (১১/১৪)\textless(১২/১৫)\textless(১৭/২১)\textless(৫/৬)}

\emph{অর্থাৎ ক্ষুদ্রতম সংখ্যাটি (১১/১৪) ।}
\end{quote}

\textbf{কোনটি বৃহত্তম ভগ্নাংশ?}

ক) (১২/১৫)

খ) (১১/১৪)

গ) (১৭/২১)

\textbf{\hl{ঘ) (৫/৬) ✔}}

\textbf{\hl{✔ সঠিক উত্তর: ঘ) (৫/৬)}}

\begin{quote}
\textbf{ব্যাখ্যা:} \emph{ক. (১২/১৫)=০.৮০}

\emph{খ. (১১/১৪)=০.৭৮৫}

\emph{গ. (১৭/২১)=০.৮০৯৫}

\emph{ঘ. (৫/৬)=০.৮৩৩৩}

\emph{∴ এদের মধ্যে বৃহত্তম ভগ্নাংশটি (৫/৬) ।}
\end{quote}

\textbf{নিচের ভগ্নাংশগুলোর মধ্যে কোনটি বৃহত্তম?}

ক) (১৪/১৭)

খ) (১১/১৪)

গ) (৭৫/৮৩)

\textbf{\hl{ঘ) (১০/১১) ✔}}

\textbf{\hl{✔ সঠিক উত্তর: ঘ) (১০/১১)}}

\begin{quote}
\textbf{ব্যাখ্যা:} \emph{ক. (১৪/১৭)≈০.৮২৩৫}

\emph{খ. (১১/১৪)≈০.৭৮৫৭}

\emph{গ. (৭৫/৮৩)≈০.৯০৩৬}

\emph{ঘ. (১০/১১)≈০.৯০৯০}

\emph{∴ দেখা যাচ্ছে (১০/১১) ভগ্নাংশটি বৃহত্তম।}
\end{quote}

\textbf{নিচের ভগ্নাংশগুলোর মধ্যে কোনটি বৃহত্তম?}

\textbf{\hl{ক) (৩/৪) ✔}}

খ) (৫/৯)

গ) (৭/১২)

ঘ) (১১/১৮)

\textbf{\hl{✔ সঠিক উত্তর: ক) (৩/৪)}}

\begin{quote}
\textbf{ব্যাখ্যা:} \emph{ক. (৩/৪)=০.৭৫}

\emph{খ. (৫/৯)≈০.৫৫৫}

\emph{গ. (৭/১২)≈০.৫৮৩৩}

\emph{ঘ. (১১/১৮)≈০.৬১১১}

\emph{∴ এদের মধ্যে (৩/৪) সবচেয়ে বড়।}
\end{quote}

\textbf{ভগ্নাংশগুলোর মধ্যে কোনটি সবচেয়ে বড়?}

ক) (২/৩)

\textbf{\hl{খ) (১৩/১৫) ✔}}

গ) (৪/৫)

ঘ) (২৩/৩০)

\textbf{\hl{✔ সঠিক উত্তর: খ) (১৩/১৫)}}

\begin{quote}
\textbf{ব্যাখ্যা:} \emph{সবগুলোর হর ৩০ রূপান্তর করে পাই,}

\emph{(২/৩)=(২×১০/৩×১০)=(২০/৩০)}

\emph{(১৩/১৫)=(১৩×২/১৫×২)=(২৬/৩০)}

\emph{(৪/৫)=(৪×৬/৫×৬)=(২৪/৩০)}

\emph{(২৩/৩০)=(২৩/৩০)}

\emph{∴ এদের মধ্যে সবচেয়ে বড় ভগ্নাংশটি হলো (১৩/১৫) ।}
\end{quote}

\textbf{কোনটি ছোট?}

ক) (৪/১৫)

\textbf{\hl{খ) (২/১৭) ✔}}

গ) (৩/১৯)

ঘ) (২/১৩)

\textbf{\hl{✔ সঠিক উত্তর: খ) (২/১৭)}}

\begin{quote}
\textbf{ব্যাখ্যা:} \emph{ক. (৪/১৫)≈০.২৬৬}

\emph{খ. (২/১৭)≈০.১১৭}

\emph{গ. (৩/১৯)≈০.১৫৭}

\emph{ঘ. (২/১৩)≈০.১৫৩}

\emph{∴ ছোট ভগ্নাংশটি (২/১৭)}
\end{quote}

\textbf{কোন সংখ্যাটি ক্ষুদ্রতম?}

ক) (১/১১)

খ) (৩/৩১)

গ) (২/২১)

\textbf{\hl{ঘ) ০.০২ ✔}}

\textbf{\hl{✔ সঠিক উত্তর: ঘ) ০.০২}}

\begin{quote}
\textbf{ব্যাখ্যা:} \emph{ক. (১/১১)≈০.০৯০৯}

\emph{খ. (৩/৩১)≈০.০৯৬৭}

\emph{গ. (২/২১)≈০.০৯৫২}

\emph{ঘ. ০.০২ = (১/৫০)=০.০২}

\emph{∴ এদের মধ্যে ০.০২ সংখ্যাটি ক্ষুদ্রতম।}
\end{quote}

\textbf{নিচের কোন সংখ্যাটি ক্ষুদ্রতম?}

ক) (১/৩)

খ) (২/৭)

\textbf{\hl{গ) (৫/২১) ✔}}

ঘ) (৩/৬)

\textbf{\hl{✔ সঠিক উত্তর: গ) (৫/২১)}}

\begin{quote}
\textbf{ব্যাখ্যা:} \emph{সবগুলোর হর সমান (অর্থাৎ ২১) করে পাই,}

\emph{ক. (১/৩)=(৭/২১)}

\emph{খ. (২/৭)=(৬/২১)}

\emph{গ. (৫/২১)}

\emph{ঘ. (৩/৬)=(১/২)=(১০.৫/২১)}

\emph{∴ এদের মধ্যে ক্ষুদ্রতম সংখ্যাটি (৫/২১)}
\end{quote}

\textbf{নিচের ভগ্নাংশগুলোর মধ্যে কোনটির মান সবচেয়ে বেশি?}

ক) (১/২০)

খ) (১/১৬)

গ) (১/১৫)

\textbf{\hl{ঘ) (১/১২) ✔}}

\textbf{\hl{✔ সঠিক উত্তর: ঘ) (১/১২)}}

\begin{quote}
\textbf{ব্যাখ্যা:} \emph{ক. (১/২০)=০.০৫}

\emph{খ. (১/১৬)=০.০৬২৫}

\emph{গ. (১/১৫)≈০.০৬৭}

\emph{ঘ. (১/১২)≈০.০৮৩}

\emph{∴ এদের মধ্যে (১/১২) ভগ্নাংশটির মান সবচেয়ে বেশি।}
\end{quote}

\textbf{নিচের কোন ভগ্নাংশটি বড়?}

\textbf{\hl{ক) (২/৩) ✔}}

খ) (২/৭)

গ) (২/৯)

ঘ) (২/৫)

\textbf{\hl{✔ সঠিক উত্তর: ক) (২/৩)}}

\begin{quote}
\textbf{ব্যাখ্যা:} \emph{ভগ্নাংশের লব সমান হলে যে ভগ্নাংশের হর ছোট সেই
ভগ্নাংশটি বড়। এখানে সবকটি ভগ্নাংশের লব ২, কিন্তু (২/৩) ভগ্নাংশটির হর (৩)
সবচেয়ে ছোট। ∴ (২/৩) বড়।}
\end{quote}

\textbf{কোন সংখ্যাটি ক্ষুদ্রতম ?}

\textbf{\hl{ক) (১/১১) ✔}}

খ) (২/২১)

গ) (৩/৩১)

ঘ) √(০.০২)

\textbf{\hl{✔ সঠিক উত্তর: ক) (১/১১)}}

\begin{quote}
\textbf{ব্যাখ্যা:} \emph{ক. (১/১১)≈০.০৯০৯}

\emph{খ. (২/২১)≈০.০৯৫২}

\emph{গ. (৩/৩১)≈০.০৯৬৭}

\emph{ঘ. √(০.০২)≈০.১৪১৪}

\emph{∴ এদের মধ্যে (১/১১) সংখ্যাটি ক্ষুদ্রতম।}
\end{quote}

\textbf{নিচের কোন সংখ্যাটি ক্ষুদ্রতম?}

ক) (২/৭)

খ) (৩/৬)

\textbf{\hl{গ) (৫/২১) ✔}}

ঘ) (১/৩)

\textbf{\hl{✔ সঠিক উত্তর: গ) (৫/২১)}}

\begin{quote}
\textbf{ব্যাখ্যা:} \emph{ক. (২/৭)≈০.২৮৫}

\emph{খ. (৩/৬)=০.৫০}

\emph{গ. (৫/২১)≈০.২৩৮}

\emph{ঘ. (১/৩)≈০.৩৩}

\emph{∴ দেখা যাচ্ছে (৫/২১) সংখ্যাটি ক্ষুদ্রতম।}
\end{quote}

\textbf{কোন সংখ্যাটি বৃহত্তম?}

ক) ০.৩

\textbf{\hl{খ) √(০.৩) ✔}}

গ) (১/৩)

ঘ) (২/৫)

\textbf{\hl{✔ সঠিক উত্তর: খ) √(০.৩)}}

\begin{quote}
\textbf{ব্যাখ্যা:} \emph{a) ০.৩}

\emph{b) √(০.৩)≈০.৫৪৭৭}

\emph{c) (১/৩)≈০.৩৩৩}

\emph{d) (২/৫)=০.৪}

\emph{∴ দেখা যাচ্ছে, √(০.৩) সংখ্যাটি বৃহত্তম।}
\end{quote}

\textbf{সবচেয়ে বড় সংখ্যা কোনটি?}

ক) ০.০০৯৯

\textbf{\hl{খ) ০.১০০ ✔}}

গ) (৯/১০০)

ঘ) (৯/১০০০)

\textbf{\hl{✔ সঠিক উত্তর: খ) ০.১০০}}

\begin{quote}
\textbf{ব্যাখ্যা:} \emph{ক. ০.০০৯৯ = (৯৯/১০০০০)}

\emph{খ. ০.১০০ = (১/১০) (বৃহত্তম)}

\emph{গ. (৯/১০০)=০.০৯}

\emph{ঘ. (৯/১০০০)=০.০০৯}

\emph{∴ এদের মধ্যে ০.১০০ সবচেয়ে বড়।}
\end{quote}

\textbf{কোন সংখ্যাটি বৃহত্তম?}

\textbf{\hl{ক) ০.১০০০ ✔}}

খ) ০.০১০০

গ) ০.০০১০

ঘ) ০.০০০১

\textbf{\hl{✔ সঠিক উত্তর: ক) ০.১০০০}}

\begin{quote}
\textbf{ব্যাখ্যা:} \emph{ক. ০.১০০০ = (১/১০) {[}বৃহত্তম{]}}

\emph{খ. ০.০১০০ = (১/১০০)}

\emph{গ. ০.০০১০ = (১/১০০০)}

\emph{ঘ. ০.০০০১ = (১/১০০০০)}

\emph{∴ ০.১০১০ বৃহত্তম সংখ্যা।}
\end{quote}

\textbf{কোনটি বৃহত্তম সংখ্যা?}

ক) (১/৭)

খ) (২/৭)

\textbf{\hl{গ) (৩/৭) ✔}}

ঘ) (১/৮)

\textbf{\hl{✔ সঠিক উত্তর: গ) (৩/৭)}}

\begin{quote}
\textbf{ব্যাখ্যা:} \emph{ক) (১/৭)≈০.১৪২৮}

\emph{খ) (২/৭)≈০.২৮৫৭}

\emph{গ) (৩/৭)≈০.৪২৮৫}

\emph{ঘ) (১/৮)=০.১২৫}

\emph{∴ বৃহত্তম সংখ্যাটি হলো (৩/৭)}
\end{quote}

\textbf{কোন সংখ্যাটি (৪/৫) এবং (১৯/১৭) এর মধ্যে অবস্থান করে না?}

\textbf{\hl{ক) (১/২) ✔}}

খ) (২/৩)

গ) (৩/৪)

ঘ) (৫/৭)

\textbf{\hl{✔ সঠিক উত্তর: ক) (১/২)}}

\begin{quote}
\textbf{ব্যাখ্যা:} \emph{সীমা দুটি দশমিকে রূপান্তর করে পাই,}

\emph{(৪/৫)=০.৮০ এবং (১৯/১৭)≈১.১১৭}

\emph{অর্থাৎ সংখ্যাটিকে ০.৮০ থেকে ১.১১৭ এর মধ্যে থাকতে হবে।}

\emph{অপশনগুলো দশমিকে প্রকাশ করলে:}

\emph{ক. (১/২)=০.৫০}

\emph{খ. (২/৩)≈০.৬৬৭}

\emph{গ. (৩/৪)=০.৭৫}

\emph{ঘ. (৫/৭)≈০.৭১৪}

\emph{দেখা যাচ্ছে, প্রদত্ত অপশনগুলোর কোনোটিই ০.৮০ এর চেয়ে বড় নয়, তবে
গাণিতিকভাবে (১/২) (০.৫০) সীমাটির সবচেয়ে দূরে অবস্থান করায় এটি প্রশ্নের উত্তর
হিসেবে চিহ্নিত।}
\end{quote}

\textbf{কোন ভগ্নাংশটি (২/৩) থেকে বড়?}

ক) (৩৩/৫০)

\textbf{\hl{খ) (৮/১১) ✔}}

গ) (৩/৫)

ঘ) (১৩/২৭)

\textbf{\hl{✔ সঠিক উত্তর: খ) (৮/১১)}}

\begin{quote}
\textbf{ব্যাখ্যা:} \emph{(২/৩)≈০.৬৬৬৭}

\emph{ক) (৩৩/৫০)=০.৬৬}

\emph{খ) (৮/১১)≈০.৭২৭২}

\emph{গ) (৩/৫)=০.৬}

\emph{ঘ) (১৩/২৭)≈০.৪৮১}

\emph{∴ দেখা যাচ্ছে (৮/১১) ভগ্নাংশটি (২/৩) থেকে বড়।}
\end{quote}

\textbf{নিম্নোক্ত কোন ভগ্নাংশটি (১/২) চেয়ে বড়?}

ক) (৩০/৬১)

খ) (৩২/৬৫)

গ) (২৯/৫৯)

\textbf{\hl{ঘ) (৩১/৬০) ✔}}

\textbf{\hl{✔ সঠিক উত্তর: ঘ) (৩১/৬০)}}

\begin{quote}
\textbf{ব্যাখ্যা:} \emph{(১/২)=০.৫}

\emph{ক) (৩০/৬১)≈০.৪৯১}

\emph{খ) (৩২/৬৫)≈০.৪৯২}

\emph{গ) (২৯/৫৯)≈০.৪৯১}

\emph{ঘ) (৩১/৬০)≈০.৫১৬৬}

\emph{∴ (৩১/৬০) ভগ্নাংশটি (১/২) এর চেয়ে বড়।}
\end{quote}

\textbf{নিচের কোন ভগ্নাংশটি (৩/৫) থেকে বড় এবং (৬/৭) থেকে ছোট?}

ক) (৭/৮)

খ) (১/৩)

\textbf{\hl{গ) (২/৩) ✔}}

ঘ) (১/২)

\textbf{\hl{✔ সঠিক উত্তর: গ) (২/৩)}}

\begin{quote}
\textbf{ব্যাখ্যা:} \emph{সীমা দুটি দশমিকে রূপান্তর করে পাই,}

\emph{(৩/৫) = ০.৬০ এবং (৬/৭) ≈ ০.৮৫৭}

\emph{অর্থাৎ সংখ্যাটি ০.৬০ এবং ০.৮৫৭ এর মাঝামাঝি হতে হবে।}

\emph{অপশনসমূহ দশমিকে প্রকাশ করলে:}

\emph{ক. (৭/৮) = ০.৮৭৫}

\emph{খ. (১/৩) ≈ ০.৩৩৩}

\emph{গ. (২/৩) ≈ ০.৬৬৭}

\emph{ঘ. (১/২) = ০.৫০}

\emph{∴ দেখা যাচ্ছে (২/৩) (০.৬৬৭) ভগ্নাংশটি ০.৬০ অপেক্ষা বড় এবং ০.৮৫৭
অপেক্ষা ছোট।}
\end{quote}

\textbf{নিচের কোন ভগ্নাংশটি (৩/৪) থেকে ছোট এবং (১/৩) থেকে বড়?}

ক) (১/৬)

\textbf{\hl{খ) (১৭/২৪) ✔}}

গ) (৩৭/৪৮)

ঘ) (১/১২)

\textbf{\hl{✔ সঠিক উত্তর: খ) (১৭/২৪)}}

\begin{quote}
\textbf{ব্যাখ্যা:} \emph{ভগ্নাংশের হরগুলো ৪, ৩, ৬, ২৪, ৪৮ ও ১২ এর ল.সা.গু. =
৪৮}

\emph{(৩/৪)=(৩×১২/৪×১২)=(৩৬/৪৮)}

\emph{(১/৩)=(১×১৬/৩×১৬)=(১৬/৪৮)}

\emph{ক. (১/৬)=(১×৮/৬×৮)=(৮/৪৮)}

\emph{খ. (১৭/২৪)=(১৭×২/২৪×২)=(৩৪/৪৮)}

\emph{গ. (৩৭/৪৮)=(৩৭×১/৪৮×১)=(৩৭/৪৮)}

\emph{ঘ. (১/১২)=(১×৪/১২×৪)=(৪/৪৮)}

\emph{এখানে, (১৬/৪৮)\textless(৩৪/৪৮)\textless(৩৬/৪৮)}

\emph{∴ (১/৩)\textless(১৭/২৪)\textless(৩/৪)}
\end{quote}

\textbf{দুটি ভগ্নাংশের যোগফল ৭০ এবং বড় ভগ্নাংশটি ৩৮(৭/১৯) হলে ছোট ভগ্নাংশটি
কত হবে?}

\textbf{\hl{ক) ৩১(১২/১৯) ✔}}

খ) ৩২(১২/১৯)

গ) ৩০(১২/১৯)

ঘ) ৩৩(৭/১৯)

\textbf{\hl{✔ সঠিক উত্তর: ক) ৩১(১২/১৯)}}

\begin{quote}
\textbf{ব্যাখ্যা:} \emph{ধরি, ছোট ভগ্নাংশটি x}

\emph{∴ x + ৩৮(৭/১৯)=৭০}

\emph{⇒ x = ৭০−(৭২৯/১৯)}

\emph{⇒ x = (১৩৩০−৭২৯/১৯)}

\emph{⇒ x = (৬০১/১৯)=৩১(১২/১৯)}
\end{quote}

\textbf{এক খন্ড জমির (৩/৮) অংশের মূল্য ৩৭৫ টাকা হলে ঐ জমির (১/৫) অংশের দাম
কত?}

ক) ৩২৫ টাকা

খ) ২৫০ টাকা

\textbf{\hl{গ) ২০০ টাকা ✔}}

ঘ) ৪০০ টাকা

\textbf{\hl{✔ সঠিক উত্তর: গ) ২০০ টাকা}}

\begin{quote}
\textbf{ব্যাখ্যা:} \emph{প্রশ্নমতে,}

\emph{(৩/৮) অংশের মূল্য = ৩৭৫ টাকা}

\emph{⇒ ১ বা সম্পূর্ণ অংশের মূল্য = ৩৭৫×(৮/৩)=১০০০ টাকা}

\emph{∴ (১/৫) অংশের মূল্য = ১০০০×(১/৫)=২০০ টাকা।}
\end{quote}

\textbf{নিচের কোন সংখ্যাটি বৃহত্তম?}

ক) (৪/১৫)

খ) (৩/২০)

গ) (৭/২৫)

\textbf{\hl{ঘ) (৩/৫) ✔}}

\textbf{\hl{✔ সঠিক উত্তর: ঘ) (৩/৫)}}

\begin{quote}
\textbf{ব্যাখ্যা:} \emph{(৪/১৫)=০.২৬৭}

\emph{(৩/২০)=০.১৫}

\emph{(৭/২৫)=০.২৮}

\emph{(৩/৫)=০.৬}

\emph{∴ এদের মধ্যে (৩/৫) ভগ্নাংশটি বৃহত্তম।}
\end{quote}

\textbf{একটি অপ্রকৃত ভগ্নাংশের লব ও হরের পার্থক্য ১ এবং সমষ্টি ৭। ভগ্নাংশটি
কত?}

\textbf{\hl{ক) (4/3) ✔}}

খ) (2/5)

গ) (3/4)

ঘ) (1/6)

\textbf{\hl{✔ সঠিক উত্তর: ক) (4/3)}}

\begin{quote}
\textbf{ব্যাখ্যা:} \emph{এখানে, অপশনভিত্তিক যাচাই করা সবচেয়ে বেশি সুবিধা
হবে। যেমন 4+3=7 এবং 4−3=1 । এভাবে যাচাই করতে হবে। তাছাড়া এখানে অপ্রকৃত
ভগ্নাংশ মাত্র একটিই (4/3) ।}
\end{quote}

\textbf{নিচের কোনটি ক্ষুদ্রতম সংখ্যা?}

\textbf{\hl{ক) ০.৩ ✔}}

খ) √(0.3)

গ) (1/3)

ঘ) (0.3)\^{}2

\textbf{\hl{✔ সঠিক উত্তর: ক) ০.৩}}

\begin{quote}
\textbf{ব্যাখ্যা:} \emph{এখানে,}

\emph{⇒ 0.3=0.3}

\emph{⇒ √(0.3)≈0.548}

\emph{⇒ (1/3)≈0.33}

\emph{⇒ (0.3)\^{}2=0.09}

\emph{∴ ০.৩ ই ক্ষুদ্রতম সংখ্যা।}
\end{quote}

\textbf{এক কেজি খাঁটি দুধে ২০০ গ্রাম পানি মিশ্রিত করলে মিশ্রিত দুধে পানির
পরিমাণ হবে-}

ক) (1/5)

\textbf{\hl{খ) (1/6) ✔}}

গ) (1/8)

ঘ) (1/10)

\textbf{\hl{✔ সঠিক উত্তর: খ) (1/6)}}

\begin{quote}
\textbf{ব্যাখ্যা:} \emph{এক কেজি = ১০০০ গ্রাম}

\emph{দুধ + পানি = 1000+200=1200 গ্রাম}

\emph{মিশ্রিত দুধে পানির পরিমাণ হবে = (200/1200)=(1/6) অংশ।}
\end{quote}

\textbf{কোন সংখ্যার (২/৭) অংশ ৬৪ এর সমান?}

ক) ১৮(২/৭)

খ) ২৪৮

গ) ২১৭

\textbf{\hl{ঘ) ২২৪ ✔}}

\textbf{\hl{✔ সঠিক উত্তর: ঘ) ২২৪}}

\begin{quote}
\textbf{ব্যাখ্যা:} \emph{প্রশ্নমতে,}

\emph{ফাঁকা ঘরের বা সম্পূর্ণ সংখ্যার (২/৭) অংশ = ৬৪}

\emph{∴ ১ বা সম্পূর্ণ অংশ = ৬৪×(৭/২)=২২৪ ।}
\end{quote}

\textbf{কোন সংখ্যার (৩/৭) অংশ ৪৮ এর সমান?}

ক) ১২৮

খ) ১৩২

\textbf{\hl{গ) ১১২ ✔}}

ঘ) ১৪০

\textbf{\hl{✔ সঠিক উত্তর: গ) ১১২}}

\begin{quote}
\textbf{ব্যাখ্যা:} \emph{ধরি, সংখ্যাটি x}

\emph{প্রশ্নমতে, x এর (৩/৭)=৪৮}

\emph{⇒ x = (৪৮×৭/৩)}

\emph{⇒ x = ১৬ × ৭}

\emph{∴ x = ১১২}
\end{quote}

\textbf{৪ টাকার (৫/৮) অংশ ২ টাকার (৪/৫) অংশের মধ্যে পার্থক্য কত?}

ক) ০.০৯ টাকা

খ) ১.৬০ টাকা

গ) ২.২৫ টাকা

\textbf{\hl{ঘ) ০.৯০ টাকা ✔}}

\textbf{\hl{✔ সঠিক উত্তর: ঘ) ০.৯০ টাকা}}

\begin{quote}
\textbf{ব্যাখ্যা:} \emph{৪ টাকার (৫/৮) অংশ = ৪×(৫/৮) টাকা = ২.৫ টাকা}

\emph{২ টাকার (৪/৫) অংশ = ২×(৪/৫) টাকা = ১.৬ টাকা}

\emph{∴ পার্থক্য = (২.৫ − ১.৬) টাকা = ০.৯০ টাকা।}
\end{quote}

\textbf{তিনটি সংখ্যার গুণফল (১/৩) । প্রথম দুটি সংখ্যা (২/৭) এবং ১(২/৫) হলে,
তৃতীয়টি কত?}

ক) (৩/৫)

\textbf{\hl{খ) (৫/৬) ✔}}

গ) (১/৩)

ঘ) (৫/৮)

\textbf{\hl{✔ সঠিক উত্তর: খ) (৫/৬)}}

\begin{quote}
\textbf{ব্যাখ্যা:} \emph{ধরি, তৃতীয় সংখ্যা = x}

\emph{প্রশ্নমতে, (২/৭)×১(২/৫)×x=(১/৩)}

\emph{⇒ (২/৭)×(৭/৫)×x=(১/৩)}

\emph{∴ x = (১×৭×৫/২×৭×৩)=(৫/৬)}
\end{quote}

\textbf{একজন লোক ডিসেম্বর মাসে অন্যান্য মাসের তুলনায় দ্বিগুণ আয় করে। তার সারা
বছরের আয়ের কত অংশ ডিসেম্বর মাসে আয় করে?}

ক) (১/৬) অংশ

খ) (৩/১৪) অংশ

গ) (২/১১) অংশ

\textbf{\hl{ঘ) (২/১৩) অংশ ✔}}

\textbf{\hl{✔ সঠিক উত্তর: ঘ) (২/১৩) অংশ}}

\begin{quote}
\textbf{ব্যাখ্যা:} \emph{ধরি, লোকটির ডিসেম্বর মাস বাদে অন্যান্য প্রতিটি মাসের
আয় = x টাকা}

\emph{∴ লোকটির ডিসেম্বর মাসের আয় = 2x টাকা}

\emph{∴ লোকটির সারা বছরের আয় = (11x + 2x) টাকা = 13x টাকা}

\emph{∴ লোকটির ডিসেম্বর মাসের আয় সারা বছরের = (২x/১৩x) অংশ = (২/১৩) অংশ।}
\end{quote}

\textbf{কোন সংখ্যার (1/4) অংশ, তার (1/5) অংশ অপেক্ষা 20 বেশি?}

ক) 200

খ) 300

\textbf{\hl{গ) 400 ✔}}

ঘ) 500

\textbf{\hl{✔ সঠিক উত্তর: গ) 400}}

\begin{quote}
\textbf{ব্যাখ্যা:} \emph{ধরি, সংখ্যাটি x}

\emph{প্রশ্নমতে, (x/৪)−(x/৫)=২০}

\emph{⇒ (৫x−৪x/২০)=২০}

\emph{⇒ (x/২০)=২০}

\emph{∴ x = 400}
\end{quote}

\textbf{তিনটি সংখ্যার গুণফল ৯৭২। ক্ষুদ্রতম সংখ্যাটি বৃহত্তম সংখ্যার এক-তৃতীয়াংশ
এবং মধ্যম সংখ্যাটি ক্ষুদ্রতম সংখ্যার (৩/২) গুণ। সংখ্যা তিনটির যোগফল কত?}

ক) ৩২

খ) ৩৬

\textbf{\hl{গ) ৩৩ ✔}}

ঘ) ৩৪

\textbf{\hl{✔ সঠিক উত্তর: গ) ৩৩}}

\begin{quote}
\textbf{ব্যাখ্যা:} \emph{ধরি, বৃহত্তম সংখ্যা x; মধ্যম সংখ্যা y; এবং ক্ষুদ্রতম
সংখ্যা z}

\emph{প্রশ্নমতে, xyz = ৯৭২ ......... (১)}

\emph{z = (x/৩) ......... (২)}

\emph{y = (৩z/২) ......... (৩)}

\emph{z এর মান (৩) নং সমীকরণে বসাই,}

\emph{y = (৩/২)×(x/৩)=(x/২) ......... (৪)}

\emph{এখন (১) নং সমীকরণে y এবং z এর মান বসাই,}

\emph{x × (x/২) × (x/৩) = ৯৭২}

\emph{⇒ x\^{}৩ = ৫৮৩২ = (১৮)\^{}৩}

\emph{∴ x = ১৮}

\emph{x এর মান (২) নং এ বসিয়ে পাই, z = (১৮/৩) = ৬}

\emph{x এর মান (৪) নং এ বসিয়ে পাই, y = (১৮/২) = ৯}

\emph{∴ সংখ্যা তিনটির যোগফল = x + y + z = ১৮ + ৯ + ৬ = ৩৩}
\end{quote}

\textbf{একটি ভগ্নাংশের লব ও হরের সমষ্টি ৭, এদের অন্তরফল ৩ হলে ভগ্নাংশটি হবে
-}

\textbf{\hl{ক) (৫/২) ✔}}

খ) (৪/৩)

গ) (৩/৫)

ঘ) (৩/৪)

\textbf{\hl{✔ সঠিক উত্তর: ক) (৫/২)}}

\begin{quote}
\textbf{ব্যাখ্যা:} \emph{ধরি, ভগ্নাংশটি (x/y)}

\emph{প্রশ্নমতে, x + y = ৭ ......... (১)}

\emph{x − y = ৩ ......... (২)}

\emph{(১) ও (২) নং যোগ করে পাই,}

\emph{2x = ১০}

\emph{∴ x = ৫}

\emph{এখন, x এর মান (১) নং সমীকরণে বসাই,}

\emph{৫ + y = ৭}

\emph{∴ y = ২}

\emph{অতএব, ভগ্নাংশটি (৫/২)}
\end{quote}

\textbf{একটি প্রকৃত ভগ্নাংশের লব ও হরের অন্তর 2 । লব ও হরের উভয় থেকে 5 বিয়োগ
করলে যে ভগ্নাংশ পাওয়া যায়, তার সাথে (1/4) যোগ করলে 1 হয়। ভগ্নাংশটি কত?}

\textbf{\hl{ক) (১১/১৩) ✔}}

খ) (৯/১১)

গ) (৭/৯)

ঘ) (১৩/১৫)

\textbf{\hl{✔ সঠিক উত্তর: ক) (১১/১৩)}}

\begin{quote}
\textbf{ব্যাখ্যা:} \emph{ধরি, লব x}

\emph{∴ হর x + 2}

\emph{∴ ভগ্নাংশ = (x/x+2)}

\emph{শর্তমতে, (x−5/x+2−5)+(1/4)=1}

\emph{⇒ (x−5/x−3)=1−(1/4)}

\emph{⇒ (x−5/x−3)=(3/4)}

\emph{⇒ 4x − 20 = 3x − 9}

\emph{∴ x = 11}

\emph{∴ নির্ণেয় ভগ্নাংশ = (11/11+2)=(11/13)}
\end{quote}

\textbf{কোন সংখ্যার (১/২) অংশের সাথে ৬ যোগ করলে সংখ্যাটির (২/৩) অংশ হবে।
সংখ্যাটি কত?}

ক) ৫৩

খ) ৬৩

\textbf{\hl{গ) ৩৬ ✔}}

ঘ) ৩৫

\textbf{\hl{✔ সঠিক উত্তর: গ) ৩৬}}

\begin{quote}
\textbf{ব্যাখ্যা:} \emph{ধরি, সংখ্যাটি x}

\emph{প্রশ্নমতে, x×(১/২)+৬=x×(২/৩)}

\emph{⇒ (২x/৩)−(x/২)=৬}

\emph{⇒ (৪x−৩x/৬)=৬}

\emph{⇒ (x/৬)=৬}

\emph{∴ x = ৩৬}
\end{quote}

\textbf{(৭/১৭) এর হর এবং লবের সঙ্গে কোন সংখ্যাটি যোগ করলে ভগ্নাংশটি (৩/৫)
হয়?}

ক) ৭

\textbf{\hl{খ) ৮ ✔}}

গ) ৬

ঘ) ৯

\textbf{\hl{✔ সঠিক উত্তর: খ) ৮}}

\begin{quote}
\textbf{ব্যাখ্যা:} \emph{ধরি, সংখ্যাটি x}

\emph{প্রশ্নমতে, (৭+x/১৭+x)=(৩/৫)}

\emph{⇒ ৩৫ + 5x = ৫১ + 3x}

\emph{⇒ 5x − 3x = ৫১ − ৩৫}

\emph{⇒ 2x = ১৬}

\emph{∴ x = ৮}
\end{quote}

\textbf{২টি ভগ্নাংশের গুণফল। এদের একটি (১৫/১৯) হলে অপরটি কত হবে?}

ক) (১/৫)

খ) (১/২)

গ) (২/৩)

\textbf{\hl{ঘ) (১/১০) ✔}}

\textbf{\hl{✔ সঠিক উত্তর: ঘ) (১/১০)}}

\begin{quote}
\textbf{ব্যাখ্যা:} \emph{প্রশ্নটিতে গুণফলের মান ও একটি ভগ্নাংশে প্রিন্টিং ভুল
রয়েছে। গুণফল (৫/৩৮) এবং একটি ভগ্নাংশ (২৫/১৯) হলে সঠিক সমাধান হতো:}

\emph{অপর ভগ্নাংশ = (৫/৩৮)÷(২৫/১৯)=(৫/৩৮)×(১৯/২৫)=(১/১০)}
\end{quote}

\textbf{দুইটি ভগ্নাংশের গুণফল (২৫/২৮) । এদের একটি (৫/৭) হলে, অপর ভগ্নাংশটি
কত?}

ক) (২/৩)

খ) (১/৩)

\textbf{\hl{গ) (৫/৪) ✔}}

ঘ) (৩/৪)

\textbf{\hl{✔ সঠিক উত্তর: গ) (৫/৪)}}

\begin{quote}
\textbf{ব্যাখ্যা:} \emph{অপর ভগ্নাংশটি = (২৫/২৮)÷(৫/৭)=(২৫/২৮)×(৭/৫)=(৫/৪)}
\end{quote}

\textbf{দুইটি ভগ্নাংশের গুণফল (১৫/২৮) । এদের একটি (৫/৭) হলে, অপর ভগ্নাংশটি
কত?}

ক) (২/৩)

খ) (১/৩)

\textbf{\hl{গ) (৩/৪) ✔}}

ঘ) (১/৪)

\textbf{\hl{✔ সঠিক উত্তর: গ) (৩/৪)}}

\begin{quote}
\textbf{ব্যাখ্যা:} \emph{অপর ভগ্নাংশটি = (১৫/২৮)÷(৫/৭)=(১৫/২৮)×(৭/৫)=(৩/৪)}
\end{quote}

\textbf{একটি ভগ্নাংশের হর ও লব এর অনুপাত ৩ : ২। লব থেকে ৬ বাদ দিলে যে
ভগ্নাংশটি পাওয়া যায় সেটি মূল ভগ্নাংশের (২/৩) । ভগ্নাংশটির লব কত?}

ক) ৯

খ) ১৬

\textbf{\hl{গ) ১৮ ✔}}

ঘ) ২৪

\textbf{\hl{✔ সঠিক উত্তর: গ) ১৮}}

\begin{quote}
\textbf{ব্যাখ্যা:} \emph{ধরি, লব = 2x এবং হর = 3x}

\emph{∴ ভগ্নাংশটি = (২x/৩x)}

\emph{শর্তমতে, (২x−৬/৩x)=(২x/৩x)×(২/৩)}

\emph{⇒ (২x−৬/৩x)=(৪/৯)}

\emph{⇒ ১৮x − ৫৪ = ১২x}

\emph{⇒ ৬x = ৫৪}

\emph{∴ x = ৯}

\emph{∴ নির্ণেয় লব = ২ × ৯ = ১৮}
\end{quote}

\textbf{পাঁচ ফুট দীর্ঘ একটি তারকে এমনভাবে দুভাগে ভাগ করা হলো যেন এক অংশ অন্য
অংশের (২/৩) হয়। ছোট অংশটি কত ইঞ্চি লম্বা?}

ক) ২২

খ) ২৩

\textbf{\hl{গ) ২৪ ✔}}

ঘ) ২৫

\textbf{\hl{✔ সঠিক উত্তর: গ) ২৪}}

\begin{quote}
\textbf{ব্যাখ্যা:} \emph{আমরা জানি, ১ ফুট = ১২ ইঞ্চি}

\emph{∴ ৫ ফুট = (৫ × ১২) ইঞ্চি = ৬০ ইঞ্চি}

\emph{ধরি, বড় অংশটি = ৩x ইঞ্চি এবং ছোট অংশটি = ২x ইঞ্চি}

\emph{প্রশ্নমতে,}

\emph{৩x + ২x = ৬০}

\emph{⇒ ৫x = ৬০}

\emph{⇒ x = ১২}

\emph{∴ ছোট অংশটি = (২ × ১২) ইঞ্চি = ২৪ ইঞ্চি}
\end{quote}

\textbf{একটি বাঁশের (১/৩) অংশ কাদায়, (৩/৫) অংশ পানিতে এবং ৬ হাত পানির
উপরে আছে। বাঁশটি কত হাত লম্বা?}

\textbf{\hl{ক) ৯০ হাত ✔}}

খ) ৮০ হাত

গ) ৭৫ হাত

ঘ) ৭০ হাত

\textbf{\hl{✔ সঠিক উত্তর: ক) ৯০ হাত}}

\begin{quote}
\textbf{ব্যাখ্যা:} \emph{ধরি, সম্পূর্ণ বাঁশটি = ১ অংশ}

\emph{কাদায় ও পানিতে আছে = ((১/৩)+(৩/৫)) অংশ = (১৪/১৫) অংশ}

\emph{পানির উপরে আছে = (১−(১৪/১৫)) অংশ = (১/১৫) অংশ}

\emph{প্রশ্নমতে,}

\emph{(১/১৫) অংশ = ৬ হাত}

\emph{∴ ১ (বা সম্পূর্ণ) অংশ = (৬ × ৫) হাত = ৯০ হাত}
\end{quote}

\textbf{একটি পাত্র (১/২) অংশ ভর্তি আছে। যদি ৮ গ্যালন সরানো হয় তবে (১/১০)
অংশ ভর্তি থাকে। পাত্রটি কত গ্যালন ধারণ করে?}

ক) ২২ গ্যালন

খ) ১৬ গ্যালন

\textbf{\hl{গ) ২০ গ্যালন ✔}}

ঘ) ২৪ গ্যালন

\textbf{\hl{✔ সঠিক উত্তর: গ) ২০ গ্যালন}}

\begin{quote}
\textbf{ব্যাখ্যা:} \emph{সরানো পেট্রোলের পরিমাণ = ((1/2)−(1/10)) অংশ}

\emph{= (5−1/10) অংশ}

\emph{= (4/10) অংশ = (2/5) অংশ}

\emph{প্রশ্নমতে,}

\emph{(2/5) অংশ = ৮ গ্যালন}

\emph{∴ ১ বা সম্পূর্ণ অংশ = (8×5/2) গ্যালন = ২০ গ্যালন}
\end{quote}

\textbf{একটি বাঁশের (২/৫) অংশ লাল, (১/৪) অংশ কালো ও (১/৩) অংশ সবুজ কাগজে
আবৃত এবং অবশিষ্ট অংশ ২ মিটার হলে, বাঁশটির দৈর্ঘ্য কত?}

ক) ১০০ মিটার

খ) ১১০ মিটার

\textbf{\hl{গ) ১২০ মিটার ✔}}

ঘ) ১২৫ মিটার

\textbf{\hl{✔ সঠিক উত্তর: গ) ১২০ মিটার}}

\begin{quote}
\textbf{ব্যাখ্যা:} \emph{ধরি, সম্পূর্ণ বাঁশটি = ১ অংশ}

\emph{লাল, কালো ও সবুজ অংশ = ((২/৫)+(১/৪)+(১/৩)) অংশ}

\emph{= (২৪+১৫+২০/৬০) অংশ}

\emph{= (৫৯/৬০) অংশ}

\emph{অবশিষ্ট অংশ = (১−(৫৯/৬০)) অংশ}

\emph{= (১/৬০) অংশ}

\emph{প্রশ্নমতে,}

\emph{(১/৬০) অংশ = ২ মিটার}

\emph{∴ ১ বা সম্পূর্ণ অংশ = (২ × ৬০) মিটার = ১২০ মিটার}

\emph{∴ সম্পূর্ণ বাঁশটির দৈর্ঘ্য = ১২০ মিটার}
\end{quote}

\textbf{একটি খুঁটির অর্ধাংশ মাটির নিচে, এক-তৃতীয়াংশ পানির মধ্যে এবং ২ মিটার
পানির উপরে আছে। খুঁটিটির দৈর্ঘ্য কত?}

\textbf{\hl{ক) ১২ মিটার ✔}}

খ) ১০ মিটার

গ) ৯ মিটার

ঘ) ১৬ মিটার

\textbf{\hl{✔ সঠিক উত্তর: ক) ১২ মিটার}}

\begin{quote}
\textbf{ব্যাখ্যা:} \emph{ধরি, সম্পূর্ণ খুঁটিটি = ১ অংশ}

\emph{মাটি ও পানিতে আছে = ((১/২)+(১/৩)) অংশ = (৩+২/৬) অংশ = (৫/৬) অংশ}

\emph{পানির উপরে আছে = (১−(৫/৬)) অংশ = (১/৬) অংশ}

\emph{প্রশ্নমতে,}

\emph{(১/৬) অংশ = ২ মিটার}

\emph{∴ ১ বা সম্পূর্ণ অংশ = (২ × ৬) মিটার = ১২ মিটার}
\end{quote}

\textbf{একটি পেট্রোল ট্যাংকের (3/4) অংশ খালি করে ৫ ট্যাংক পূর্ণ করা হলো।
কাজেই প্রত্যেকে সমপরিমাণ পেট্রোল ধারণ করে। প্রত্যেক ট্যাংকের ধারণ ক্ষমতা পূর্ণ
ট্যাংকের কত অংশ?}

\textbf{\hl{ক) (3/20) ✔}}

খ) (4/20)

গ) (5/20)

ঘ) (6/20)

\textbf{\hl{✔ সঠিক উত্তর: ক) (3/20)}}

\begin{quote}
\textbf{ব্যাখ্যা:} \emph{৫টি ট্যাংকের ধারণ ক্ষমতা = পূর্ণ ট্যাংকের (3/4) অংশ}

\emph{∴ ১টি ট্যাংকের ধারণ ক্ষমতা = পূর্ণ ট্যাংকের ((3/4×5)) অংশ = (3/20) অংশ}
\end{quote}

\textbf{৪টি ১ টাকার নোট ও ৮টি ২ টাকার নোট একত্রে ৮টি ৫ টাকার নোটের কত
অংশ?}

ক) (1/4)

\textbf{\hl{খ) (1/2) ✔}}

গ) (1/8)

ঘ) (1/16)

\textbf{\hl{✔ সঠিক উত্তর: খ) (1/2)}}

\begin{quote}
\textbf{ব্যাখ্যা:} \emph{৪টি ১ টাকার নোট ও ৮টি ২ টাকার নোট একত্রে = (৪ × ১
+ ৮ × ২) টাকা = (৪ + ১৬) টাকা = ২০ টাকা}

\emph{৮টি ৫ টাকার নোট = (৮ × ৫) টাকা = ৪০ টাকা}

\emph{∴ নির্ণেয় অংশ = (20/40)=(1/2) অংশ}
\end{quote}

\textbf{রহিম তার বেতনের (১/৫) অংশ খরচ করে একটি শার্ট এবং ৫০০ টাকা খরচ করে
একটি প্যান্ট কিনলো। এই টাকা খরচ করার পর তার কাছে বেতনের ৪০ শতাংশ টাকা রয়ে
গেল। রহিম কত টাকা বেতন পেয়েছিল?}

ক) ২৫০০০

খ) ২০০০

গ) ৩০০০

\textbf{\hl{ঘ) কোনোটিই নয় ✔}}

\textbf{\hl{✔ সঠিক উত্তর: ঘ) কোনোটিই নয়}}

\begin{quote}
\textbf{ব্যাখ্যা:} \emph{ধরি, রহিমের মোট বেতন = x টাকা}

\emph{শার্টের খরচ = (x/5) টাকা}

\emph{প্যান্টের খরচ = ৫০০ টাকা}

\emph{অবশিষ্ট টাকা = x এর ৪০\% = (40x/100)=(2x/5) টাকা}

\emph{মোট খরচ = (x−(2x/5)) টাকা = (3x/5) টাকা}

\emph{প্রশ্নমতে,}

\emph{(x/5)+500=(3x/5)}

\emph{⇒ (3x/5)−(x/5)=500}

\emph{⇒ (2x/5)=500}

\emph{⇒ 2x = ২৫০০}

\emph{⇒ x = ১২৫০}

\emph{∴ রহিম ১২৫০ টাকা বেতন পেয়েছিল। সুতরাং সঠিক উত্তর (ঘ) কোনোটিই নয়।}
\end{quote}

\textbf{একটি প্রকৃত ভগ্নাংশের হর ও লবের অন্তর ২, হর ও লব উভয় থেকে ৩ বিয়োগ
করলে যে ভগ্নাংশ পাওয়া যায় তার সঙ্গে (১/৪) যোগ করলে যোগফল ১ হয়, ভগ্নাংশটি
কত?}

ক) (7/9)

\textbf{\hl{খ) (9/11) ✔}}

গ) (11/13)

ঘ) (13/15)

\textbf{\hl{✔ সঠিক উত্তর: খ) (9/11)}}

\begin{quote}
\textbf{ব্যাখ্যা:} \emph{ধরি, প্রকৃত ভগ্নাংশটি = (x/x+2)}

\emph{প্রশ্নমতে,}

\emph{(x−3/x+2−3)+(1/4)=1}

\emph{⇒ (x−3/x−1)=1−(1/4)}

\emph{⇒ (x−3/x−1)=(3/4)}

\emph{⇒ 4x − 12 = 3x − 3}

\emph{⇒ 4x − 3x = 12 − 3}

\emph{⇒ x = 9}

\emph{∴ ভগ্নাংশটি = (9/9+2)=(9/11)}

\emph{অপশন টেস্ট:}

\emph{অপশন (খ) (9/11) এর লব ও হরের থেকে ৩ বিয়োগ করলে হয়
(9−3/11−3)=(6/8)=(3/4)}

\emph{এখন, (3/4)+(1/4)=(4/4)=1 , যা শর্ত পূরণ করে।}
\end{quote}

\textbf{কোন ভগ্নাংশের লবের সাথে ৩ যোগ করলে এর মান ১ এবং হরের সাথে ২ যোগ
করলে মান (১/২) হয়। ভগ্নাংশটি নির্ণয় করুন।}

ক) (৮/১১)

খ) (৭/১০)

গ) (৪/৭)

\textbf{\hl{ঘ) (৫/৮) ✔}}

\textbf{\hl{✔ সঠিক উত্তর: ঘ) (৫/৮)}}

\begin{quote}
\textbf{ব্যাখ্যা:} \emph{ধরি, ভগ্নাংশটি (x/y)}

\emph{১ম প্রশ্নমতে, (x+৩/y)=১ ⇒ y = x + ৩ ......... (i)}

\emph{২য় প্রশ্নমতে, (x/y+২)=(১/২) ⇒ 2x = y + ২ ......... (ii)}

\emph{(ii) নং সমীকরণে y এর মান বসিয়ে পাই,}

\emph{2x = x + ৩ + ২}

\emph{⇒ 2x − x = ৫}

\emph{∴ x = ৫}

\emph{এখন (i) নং সমীকরণে x এর মান বসিয়ে পাই,}

\emph{y = ৫ + ৩ = ৮}

\emph{∴ অতএব, ভগ্নাংশটি (৫/৮)}
\end{quote}

\textbf{একটি ভগ্নাংশের লব ও হর উভয় থেকে এক বিয়োগ করলে ভগ্নাংশটি (২/৩) হয়।
কিন্তু হর ও লব উভয়ের সঙ্গে ১ যোগ করলে ভগ্নাংশটি (৩/৪) হয়। ভগ্নাংশটি কত?}

ক) (৪/৫)

খ) (৩/৪)

গ) (৭/৯)

\textbf{\hl{ঘ) (৫/৭) ✔}}

\textbf{\hl{✔ সঠিক উত্তর: ঘ) (৫/৭)}}

\begin{quote}
\textbf{ব্যাখ্যা:} \emph{ধরি, ভগ্নাংশটি (x/y)}

\emph{১ম শর্তমতে, (x−১/y−১)=(২/৩)}

\emph{⇒ 3(x − ১) = 2(y − ১)}

\emph{⇒ 3x − 2y = ১ ......... (১)}

\emph{২য় শর্তমতে, (x+১/y+১)=(৩/৪)}

\emph{⇒ 4(x + ১) = 3(y + ১)}

\emph{⇒ 4x − 3y = −১ ......... (২)}

\emph{(১) নং সমীকরণকে ৪ দ্বারা এবং (২) নং কে ৩ দ্বারা গুণ করে বিয়োগ করে
পাই,}

\emph{(12x − 8y) − (12x − 9y) = ৪ − (−৩)}

\emph{⇒ y = ৭}

\emph{y এর মান (১) নং সমীকরণে বসাই,}

\emph{3x − ২ × ৭ = ১}

\emph{⇒ 3x = ১৫}

\emph{∴ x = ৫}

\emph{∴ নির্ণেয় ভগ্নাংশ = (x/y)=(৫/৭)}
\end{quote}

\textbf{একটি ভগ্নাংশের লব ও হরের সমষ্টি 11 । লব থেকে 2 বিয়োগ এবং হরের সাথে
3 যোগ করলে ভগ্নাংশের মান হয় (1/2) । ভগ্নাংশটি নির্ণয় কর।}

ক) (4/7)

খ) (5/6)

\textbf{\hl{গ) (6/5) ✔}}

ঘ) (7/4)

\textbf{\hl{✔ সঠিক উত্তর: গ) (6/5)}}

\begin{quote}
\textbf{ব্যাখ্যা:} \emph{ধরি, ভগ্নাংশটি (x/y)}

\emph{১ম শর্তমতে, x + y = 11 ⇒ y = 11 − x ......... (১)}

\emph{২য় শর্তমতে, (x−2/y+3)=(1/2)}

\emph{⇒ 2(x − 2) = y + 3}

\emph{⇒ 2x − 4 = (11 − x) + 3 {[}(১) নং থেকে মান বসিয়ে{]}}

\emph{⇒ 2x − 4 = 14 − x}

\emph{⇒ 3x = 18}

\emph{∴ x = 6}

\emph{এখন, x এর মান (১) নং সমীকরণে বসিয়ে পাই,}

\emph{y = 11 − 6 = 5}

\emph{∴ নির্ণেয় ভগ্নাংশটি (6/5)}
\end{quote}

\textbf{তপুর কাছে ৩৬০টি মার্বেল আছে যার (১/৩) অংশ হচ্ছে কালো। যদি সে তার
ভাইকে তার কালো মার্বেলের (১/৩) অংশ দিয়ে দেয়, তবে তপুর কাছে যত মার্বেল থাকবে
তার কত শতাংশ কালো মার্বেল ?}

\textbf{\hl{ক) ২৫.০০ ✔}}

খ) ২২.২৫

গ) ২০

ঘ) কোনোটিই নয়

\textbf{\hl{✔ সঠিক উত্তর: ক) ২৫.০০}}

\begin{quote}
\textbf{ব্যাখ্যা:} \emph{তপুর মোট মার্বেল সংখ্যা = ৩৬০ টি}

\emph{তপুর কালো মার্বেলের সংখ্যা = ৩৬০×(১/৩) টি = ১২০ টি}

\emph{তপুর ভাইকে দেওয়া কালো মার্বেল = ১২০×(১/৩) টি = ৪০ টি}

\emph{ভাইকে দেওয়ার পর তপুর মোট মার্বেল থাকে = (৩৬০ − ৪০) টি = ৩২০ টি}

\emph{এবং তপুর কাছে অবশিষ্ট কালো মার্বেল থাকে = (১২০ − ৪০) টি = ৮০ টি}

\emph{∴ মোট মার্বেলের মধ্যে কালো মার্বেলের শতাংশ = (৮০/৩২০)×১০০\% = ২৫\%}
\end{quote}

\textbf{রফিকের আয়ের (১/২) অংশ ৩৫০০ টাকা হলে তার আয়ের (১/৭) অংশ কত হবে?}

\textbf{\hl{ক) ১০০০ ✔}}

খ) ১২৫০

গ) ১৪০০

ঘ) কোনোটিই নয়

\textbf{\hl{✔ সঠিক উত্তর: ক) ১০০০}}

\begin{quote}
\textbf{ব্যাখ্যা:} \emph{রফিকের আয়ের (১/২) অংশ = ৩৫০০ টাকা}

\emph{∴ রফিকের মোট বা ১ (সম্পূর্ণ) আয় = (৩৫০০ × ২) টাকা = ৭০০০ টাকা}

\emph{∴ রফিকের আয়ের (১/৭) অংশ = ৭০০০×(১/৭) টাকা = ১০০০০ টাকা।}
\end{quote}

\textbf{এক ব্যক্তি মাসিক বেতনের (১/৪০) অংশ মহার্ঘভাতা পান। তার মাসিক বেতন
১৬০০ টাকা হলে, তার মহার্ঘভাতা কত?}

ক) ৩০ টাকা

\textbf{\hl{খ) ৪০ টাকা ✔}}

গ) ৬৪ টাকা

ঘ) ৮০ টাকা

\textbf{\hl{✔ সঠিক উত্তর: খ) ৪০ টাকা}}

\begin{quote}
\textbf{ব্যাখ্যা:} \emph{ব্যক্তির মাসিক বেতন = ১৬০০ টাকা}

\emph{∴ মহার্ঘভাতা পান = ১৬০০ টাকার (১/৪০) = ৪০ টাকা।}
\end{quote}

\textbf{কোন কর্মচারী মূল বেতনের (১/৪০) ভাগের সমপরিমাণ বিশেষ ভাতা পান। তার
বিশেষ ভাতা ৪০ টাকা হলে, মূল বেতন কত টাকা?}

ক) ১২০০

খ) ১৪০০

\textbf{\hl{গ) ১৬০০ ✔}}

ঘ) ১৮০০

\textbf{\hl{✔ সঠিক উত্তর: গ) ১৬০০}}

\begin{quote}
\textbf{ব্যাখ্যা:} \emph{ধরি, মূল বেতন x টাকা}

\emph{প্রশ্নমতে, x এর (১/৪০)=৪০ টাকা}

\emph{⇒ x = ৪০ × ৪০}

\emph{∴ x = ১৬০০ টাকা।}
\end{quote}

\textbf{কোনো সম্পত্তির (৭/৮) অংশের মূল্য ৯১০০.০০ টাকা। ঐ সম্পত্তির (৩/৪) অংশের
মূল্য কত?}

\textbf{\hl{ক) ৭৮০০ টাকা ✔}}

খ) ৭২০০ টাকা

গ) ৮৮০০ টাকা

ঘ) ৯৮০০ টাকা

\textbf{\hl{✔ সঠিক উত্তর: ক) ৭৮০০ টাকা}}

\begin{quote}
\textbf{ব্যাখ্যা:} \emph{সম্পত্তির (৭/৮) অংশের মূল্য = ৯১০০ টাকা}

\emph{∴ সম্পত্তির ১ বা সম্পূর্ণ অংশের মূল্য = (৯১০০×৮/৭) টাকা = ১০৪০০ টাকা}

\emph{∴ সম্পত্তির (৩/৪) অংশের মূল্য = (৯১০০×৮×৩/৭×৪) = ৭৮০০ টাকা।}
\end{quote}

\textbf{কোন পুস্তকের ৯৬ পৃষ্ঠা পড়বার পর (৫/১৩) অংশ পড়তে বাকি থাকলে পুস্তকটির
মোট পৃষ্ঠা সংখ্যা কত?}

ক) ১৮৫ পৃষ্ঠা

\textbf{\hl{খ) ১৫৬ পৃষ্ঠা ✔}}

গ) ২৫০ পৃষ্ঠা

ঘ) ৩২০ পৃষ্ঠা

\textbf{\hl{✔ সঠিক উত্তর: খ) ১৫৬ পৃষ্ঠা}}

\begin{quote}
\textbf{ব্যাখ্যা:} \emph{পড়া হওয়া অংশের পরিমাণ = ১−(৫/১৩)=(১৩−৫/১৩)=(৮/১৩)
অংশ}

\emph{শর্তমতে, (৮/১৩) অংশ = ৯৬ পৃষ্ঠা}

\emph{∴ ১ বা সম্পূর্ণ অংশ = (৯৬×১৩/৮) = ১৫৬ পৃষ্ঠা}
\end{quote}

\textbf{একটি তেলপূর্ণ পাত্রের ওজন ৩২ কেজি এবং অর্ধেক তেল থাকাকালীন পাত্রের
ওজন ২০ কেজি। পাত্রটির ওজন কত?}

\textbf{\hl{ক) ৮ কেজি ✔}}

খ) ১০ কেজি

গ) ১২ কেজি

ঘ) ১৪ কেজি

\textbf{\hl{✔ সঠিক উত্তর: ক) ৮ কেজি}}

\begin{quote}
\textbf{ব্যাখ্যা:} \emph{পাত্র + পূর্ণ তেল = ৩২ কেজি}

\emph{পাত্র + অর্ধপূর্ণ তেল = ২০ কেজি}

\emph{বিয়োগ করে পাই,}

\emph{অর্ধপূর্ণ তেলের ওজন = (৩২ − ২০) কেজি = ১২ কেজি}

\emph{∴ পাত্রের ওজন = (২০ − ১২) কেজি = ৮ কেজি।}
\end{quote}

\textbf{একটি পানি ভর্তি বালতির ওজন ১৬.৫ কেজি। বালতির (১/৪) অংশ পানি ভর্তি
থাকলে তার ওজন ৫.২৫ কেজি হয়। খালি বালতির ওজন কত কেজি?}

ক) ২.৫

\textbf{\hl{খ) ১.৫ ✔}}

গ) ৩.৫

ঘ) ১.২৫

\textbf{\hl{✔ সঠিক উত্তর: খ) ১.৫}}

\begin{quote}
\textbf{ব্যাখ্যা:} \emph{ধরি, খালি বালতির ওজন x কেজি}

\emph{∴ সম্পূর্ণ পানির ওজন = (১৬.৫ − x) কেজি}

\emph{বালতির (১/৪) অংশের শুধু পানির ওজন = (৫.২৫ − x) কেজি}

\emph{∴ ১ বা সম্পূর্ণ অংশের শুধু পানির ওজন = (৫.২৫ − x) × ৪ = (২১ − 4x) কেজি}

\emph{শর্তমতে, ১৬.৫ − x = ২১ − 4x}

\emph{⇒ 4x − x = ২১ − ১৬.৫}

\emph{⇒ 3x = ৪.৫}

\emph{∴ x = (৪.৫/৩)=১.৫}

\emph{∴ খালি বালতির ওজন ১.৫ কেজি।}
\end{quote}

\textbf{পানি ভর্তি একটি বালতির ওজন ১২ কেজি। বালতির অর্ধেক পানি ভর্তি হলে
তার ওজন দাঁড়ায় ৭ কেজি। খালি বালতির ওজন কত?}

ক) ৫ কেজি.

\textbf{\hl{খ) ২ কেজি ✔}}

গ) ৭ কেজি

ঘ) ১কেজি

\textbf{\hl{✔ সঠিক উত্তর: খ) ২ কেজি}}

\begin{quote}
\textbf{ব্যাখ্যা:} \emph{বালতির ওজন + ১ বালতি পানির ওজন = ১২ কেজি}

\emph{বালতির ওজন + (১/২) বালতি পানির ওজন = ৭ কেজি}

\emph{বিয়োগ করে পাই,}

\emph{(১/২) বালতি পানির ওজন = (১২ − ৭) কেজি = ৫ কেজি}

\emph{∴ খালি বালতির ওজন = (৭ − ৫) কেজি = ২ কেজি।}
\end{quote}

\textbf{দৈনিক ৮ ঘণ্টা কাজ করে ৩টি পাম্প ১টি জলাধারকে ২ দিনে খালি করতে
পারে। জলাধারটিকে ১ দিনে খালি করতে হলে ৪টি পাম্পের দৈনিক কত ঘণ্টা কাজ করতে
হবে?}

ক) ১০

\textbf{\hl{খ) ১২ ✔}}

গ) ১৫

ঘ) ১৬

\textbf{\hl{✔ সঠিক উত্তর: খ) ১২}}

\begin{quote}
\textbf{ব্যাখ্যা:} \emph{আমরা জানি, D\_১ × P\_১ × H\_১ = D\_২ × P\_২ ×
H\_২}

\emph{এখানে, D\_১ = ২ দিন, P\_১ = ৩টি পাম্প, H\_১ = ৮ ঘণ্টা}

\emph{D\_২ = ১ দিন, P\_২ = ৪টি পাম্প, H\_২ = ?}

\emph{∴ ১ × ৪ × H\_২ = ২ × ৩ × ৮}

\emph{⇒ H\_২ = (২×৩×৮/৪) = ১২ ঘণ্টা}
\end{quote}

\textbf{মি. রেজা তার সম্পদের 12\% স্ত্রীকে, 58\% ছেলেকে এবং অবশিষ্ট
7,20,000/- টাকা মেয়েকে দিলেন। তার সম্পদের মোট মূল্য কত?}

\textbf{\hl{ক) 2400000 ✔}}

খ) 2000000

গ) 1600000

ঘ) 1200000

\textbf{\hl{✔ সঠিক উত্তর: ক) 2400000}}

\begin{quote}
\textbf{ব্যাখ্যা:} \emph{ছেলে + স্ত্রী পায় = 12\% + 58\% = 70\%}

\emph{∴ মেয়ে পাবে = (100 − 70)\% = 30\%}

\emph{ধরি, মোট সম্পত্তি x}

\emph{∴ x এর 30\% = 720000}

\emph{⇒ x × (30/100) = 720000}

\emph{∴ x = (100×720000/30) = 2400000}
\end{quote}

\textbf{এক ব্যক্তি সম্পত্তির (২/৩) অংশ পুত্রকে এবং (১/৩) অংশ কন্যাকে দিলেন।
কন্যা পুত্র অপেক্ষা ১৫০০ টাকা কম পেল। সম্পূর্ণ সম্পত্তির মূল্য কত?}

ক) ৩০০০ টাকা

\textbf{\hl{খ) ৪৫০০ টাকা ✔}}

গ) ৬০০০ টাকা

ঘ) ৭৫০০ টাকা

\textbf{\hl{✔ সঠিক উত্তর: খ) ৪৫০০ টাকা}}

\begin{quote}
\textbf{ব্যাখ্যা:} \emph{কন্যা পুত্র অপেক্ষা কম পেল = (২/৩)−(১/৩) অংশ = (১/৩)
অংশ}

\emph{প্রশ্নমতে, সম্পত্তির (১/৩) অংশ = ১৫০০ টাকা}

\emph{∴ সম্পূর্ণ বা ১ অংশ = (১৫০০ × ৩) টাকা = ৪৫০০ টাকা।}
\end{quote}

\textbf{একটি বলপেন ৭ টাকা ও একটি পেন্সিল ৫ টাকা দরে রাজু কয়েকটি পেন্সিল ও
বলপেন কিনল এবং এগুলির মূল্য বাবদ ৩৮ টাকা পরিশোধ করল। সে কয়টি বলপেন
কিনেছিল?}

ক) ৩টি

খ) ২টি

\textbf{\hl{গ) ৪টি ✔}}

ঘ) ৫টি

\textbf{\hl{✔ সঠিক উত্তর: গ) ৪টি}}

\begin{quote}
\textbf{ব্যাখ্যা:} \emph{ধরি, বলপেনের সংখ্যা x এবং পেন্সিলের সংখ্যা y}

\emph{প্রশ্নমতে, 7x + 5y = 38}

\emph{এখানে, x ও y এর মান ধনাত্মক পূর্ণসংখ্যা হতে হবে।}

\emph{x = 4 বসালে,}

\emph{7(4) + 5y = 38}

\emph{⇒ 28 + 5y = 38}

\emph{⇒ 5y = 10}

\emph{∴ y = 2}

\emph{সুতরাং, সে ৪টি বলপেন কিনেছিল।}
\end{quote}

\textbf{১টি নার্সারিতে ১৬ জাতের ফুল গাছ আছে। (১/৪) অংশ জাতের ৫টি করে ও
(৩/৪) অংশ জাতের ৪টি করে গাছ আছে। সর্বমোট কতটি গাছ আছে?}

\textbf{\hl{ক) ৬৮টি ✔}}

খ) ৪৮টি

গ) ১৬৪টি

ঘ) কোনোটিই নয়

\textbf{\hl{✔ সঠিক উত্তর: ক) ৬৮টি}}

\begin{quote}
\textbf{ব্যাখ্যা:} \emph{১৬টি জাতের (১/৪) অংশ = ৪টি জাত}

\emph{এবং ১৬টি জাতের (৩/৪) অংশ = ১২টি জাত}

\emph{৪টি জাতের প্রতিটির জন্য ৫টি করে গাছ = (৪ × ৫)টি = ২০টি}

\emph{আবার, ১২টি জাতের প্রতিটির জন্য ৪টি করে গাছ = (১২ × ৪)টি = ৪৮টি}

\emph{∴ সর্বমোট গাছ সংখ্যা = (২০ + ৪৮)টি = ৬৮টি}
\end{quote}

\textbf{একটি শ্রেণির প্রতি বেঞ্চে ৪ জন করে ছাত্র বসলে ৩টি বেঞ্চ খালি থাকে।
কিন্তু প্রতি বেঞ্চে ৩ জন করে বসলে ৬ জন ছাত্রের দাঁড়িয়ে থাকতে হয়। ঐ শ্রেণির ছাত্র
সংখ্যা কত?}

ক) ৫০

\textbf{\hl{খ) ৬০ ✔}}

গ) ৭০

ঘ) ৮০

\textbf{\hl{✔ সঠিক উত্তর: খ) ৬০}}

\begin{quote}
\textbf{ব্যাখ্যা:} \emph{ধরি, বেঞ্চ সংখ্যা x}

\emph{১ম শর্তমতে, ছাত্র সংখ্যা = 4(x − 3)}

\emph{২য় শর্তমতে, ছাত্র সংখ্যা = 3x + 6}

\emph{প্রশ্নমতে, 4(x − 3) = 3x + 6}

\emph{⇒ 4x − 12 = 3x + 6}

\emph{⇒ 4x − 3x = 6 + 12}

\emph{∴ x = 18}

\emph{∴ ছাত্র সংখ্যা = (3 × 18) + 6 = 54 + 6 = 60 জন।}
\end{quote}

\textbf{একটি স্কুলের ছাত্র-ছাত্রীর সংখ্যার মধ্যে (5/6) অংশ ছাত্র এবং বাকি অংশ
ছাত্রী। যদি ছাত্রের সংখ্যা ছাত্রীদের সংখ্যা অপেক্ষা 120 জন বেশি হয়, তবে ছাত্রীর
সংখ্যা কত ?}

ক) 25 জন

খ) 35 জন

\textbf{\hl{গ) 30 জন ✔}}

ঘ) 20 জন

\textbf{\hl{✔ সঠিক উত্তর: গ) 30 জন}}

\begin{quote}
\textbf{ব্যাখ্যা:} \emph{ধরি, মোট ছাত্র-ছাত্রীর সংখ্যা x জন}

\emph{∴ ছাত্রের সংখ্যা = (৫x/৬) জন}

\emph{এবং ছাত্রীর সংখ্যা = x−(৫x/৬)=(x/৬) জন}

\emph{প্রশ্নমতে, (৫x/৬)−(x/৬)=120}

\emph{⇒ (4x/6)=120}

\emph{⇒ 4x = 720}

\emph{∴ x = 180}

\emph{∴ ছাত্রীর সংখ্যা = (180/6)=30 জন।}
\end{quote}

\textbf{কোন একটি স্কুলের শিক্ষকদের মধ্যে (২/৩) অংশ মহিলা, পুরুষ শিক্ষকদের ১২ জন
অবিবাহিত এবং (৩/৫) অংশ বিবাহিত। ঐ স্কুলের শিক্ষক-শিক্ষিকার সংখ্যা কত?}

\textbf{\hl{ক) ৯০ ✔}}

খ) ৮০

গ) ৮৫

ঘ) ১২০

\textbf{\hl{✔ সঠিক উত্তর: ক) ৯০}}

\begin{quote}
\textbf{ব্যাখ্যা:} \emph{ধরি, মোট শিক্ষক-শিক্ষিকা = x জন}

\emph{মহিলা শিক্ষকের সংখ্যা = (২x/৩) জন}

\emph{∴ পুরুষ শিক্ষকের সংখ্যা = x−(২x/৩)=(x/৩) জন}

\emph{পুরুষ শিক্ষকদের মধ্যে বিবাহিত = (৩/৫) অংশ}

\emph{∴ পুরুষ শিক্ষকদের মধ্যে অবিবাহিত = ১−(৩/৫)=(২/৫) অংশ}

\emph{প্রশ্নমতে, (x/৩) এর (২/৫)=১২}

\emph{⇒ (২x/১৫)=১২}

\emph{⇒ 2x = ১৮০}

\emph{∴ x = ৯০}

\emph{∴ ঐ স্কুলের মোট শিক্ষক-শিক্ষিকার সংখ্যা ৯০ জন।}
\end{quote}

\textbf{একটি বাঁশের (১/৪) অংশ কাদায়, (৩/৫) অংশ পানিতে এবং অবশিষ্ট ৩ মিটার
পানির উপরে আছে। বাঁশটির দৈর্ঘ্য কত?}

\textbf{\hl{ক) ২০ মিটার ✔}}

খ) ১৫ মিটার

গ) ১৬ মিটার

ঘ) ১২ মিটার

\textbf{\hl{✔ সঠিক উত্তর: ক) ২০ মিটার}}

\begin{quote}
\textbf{ব্যাখ্যা:} \emph{ধরি, বাঁশটির মোট দৈর্ঘ্য x মিটার।}

\emph{কাদায় আছে = (x/৪) মিটার}

\emph{এবং পানিতে আছে = (৩x/৫) মিটার}

\emph{প্রশ্নমতে, (x/৪)+(৩x/৫)+৩=x}

\emph{⇒ (৫x+১২x+৬০/২০)=x}

\emph{⇒ ১৭x + ৬০ = ২০x}

\emph{⇒ ৩x = ৬০}

\emph{∴ x = ২০}

\emph{∴ বাঁশটির দৈর্ঘ্য ২০ মিটার।}
\end{quote}

\textbf{একটি খুঁটির (১/৩) অংশ মাটির নিচে এবং (১/২) অংশ পানির নিচে থাকলে
মোট কত অংশ মাটি ও পানির নিচে আছে?}

ক) (১/৩)

খ) (২/৩)

\textbf{\hl{গ) (৫/৬) ✔}}

ঘ) (৫/৭)

\textbf{\hl{✔ সঠিক উত্তর: গ) (৫/৬)}}

\begin{quote}
\textbf{ব্যাখ্যা:} \emph{মাটি ও পানির নিচে আছে মোট = (১/৩)+(১/২) অংশ}

\emph{= (২+৩/৬) অংশ}

\emph{= (৫/৬) অংশ}

\emph{∴ (৫/৬) অংশ মাটি ও পানির নিচে আছে।}
\end{quote}

\textbf{একটি খুঁটির (৫/৬) অংশ কালো এবং বাকী অংশ সাদা। খুঁটিটির কালো এবং সাদা
অংশের দৈর্ঘ্যের পার্থক্য ৬ মিটার হলে সম্পূর্ণ খুঁটির দৈর্ঘ্য কত?}

\textbf{\hl{ক) ৯ মিটার ✔}}

খ) ৮ মিটার

গ) ৬ মিটার

ঘ) ১২ মিটার

\textbf{\hl{✔ সঠিক উত্তর: ক) ৯ মিটার}}

\begin{quote}
\textbf{ব্যাখ্যা:} \emph{মনে করি, সম্পূর্ণ খুঁটির দৈর্ঘ্য = x}

\emph{∴ কালো অংশ = (৫x/৬)}

\emph{∴ সাদা অংশ = x−(৫x/৬)=(x/৬)}

\emph{প্রশানুসারে, (৫x/৬)−(x/৬)=৬}

\emph{বা, (৫x−x/৬)=৬}

\emph{বা, (৪x/৬)=৬}

\emph{বা, ৪x = ৩৬}

\emph{∴ x = ৯ মিটার।}
\end{quote}

\textbf{একটি খুঁটির অর্ধাংশে মাটির নিচে, এক তৃতীয়াংশ পানির মধ্যে ও ১২ ফুট
পানির উপরে আছে। খুঁটির দৈর্ঘ্য কত?}

\textbf{\hl{ক) ৭২ ফুট ✔}}

খ) ৮০ ফুট

গ) ৬০ ফুট

ঘ) ৫৪ ফুট

\textbf{\hl{✔ সঠিক উত্তর: ক) ৭২ ফুট}}

\begin{quote}
\textbf{ব্যাখ্যা:} \emph{খুঁটির মাটির নিচে ও পানির মধ্যে মোট অংশ =
(১/২)+(১/৩) অংশ = (৩+২/৬)=(৫/৬) অংশ}

\emph{∴ পানির ওপরের অংশ = ১−(৫/৬) অংশ = (১/৬) অংশ}

\emph{প্রশ্নমতে, (১/৬) অংশ = ১২ ফুট}

\emph{∴ ১ বা সম্পূর্ণ অংশ = (৬ × ১২) ফুট = ৭২ ফুট।}
\end{quote}

\textbf{একটি বাঁশের অর্ধাংশ মাটির নিচে, এক তৃতীয়াংশ পানির মধ্যে এবং ৪ ফুট
পানির উপর আছে। বাঁশটির দৈর্ঘ্য কত ফুট?}

\textbf{\hl{ক) ২৪ ✔}}

খ) ২১

গ) ২৭

ঘ) ৩০

\textbf{\hl{✔ সঠিক উত্তর: ক) ২৪}}

\begin{quote}
\textbf{ব্যাখ্যা:} \emph{ধরি, মোট দৈর্ঘ্য x ফুট}

\emph{প্রশ্নমতে, x−(x/২)+(x/৩)=৪}

\emph{⇒ x−(৩x+২x/৬)=৪}

\emph{⇒ x−(৫x/৬)=৪}

\emph{⇒ (৬x−৫x/৬)=৪}

\emph{⇒ ৬x − ৫x = ৬ × ৪}

\emph{∴ x = ২৪ ফুট।}
\end{quote}

\textbf{যদি কোন স্কুলের ছাত্রীদের (1/3) অংশ ঐ স্কুলের মোট ছাত্র-ছাত্রীদের (1/5)
-এর সমান হয় তাহলে স্কুলটিতে ছাত্র : ছাত্রী হবে -}

ক) 5 : 3

খ) 3 : 2

গ) 2 : 5

\textbf{\hl{ঘ) সঠিক উত্তর নেই ✔}}

\textbf{\hl{✔ সঠিক উত্তর: ঘ) সঠিক উত্তর নেই}}

\begin{quote}
\textbf{ব্যাখ্যা:} \emph{ধরি, ছাত্র সংখ্যা x এবং ছাত্রী সংখ্যা y}

\emph{শর্তানুসারে, (1/3)y=(1/5)(x+y)}

\emph{⇒ 5y = 3x + 3y}

\emph{⇒ 2y = 3x}

\emph{⇒ (x/y)=(2/3)}

\emph{∴ x : y = 2 : 3}

\emph{অর্থাৎ, সঠিক উত্তর 2 : 3 (যা দেওয়া অপশনগুলোর মধ্যে নেই)।}
\end{quote}

\textbf{(১/৪)+(১/১২)+(১/৩)=?}

ক) (১/৪)

\textbf{\hl{খ) (২/৩) ✔}}

গ) (৭/১২)

ঘ) ১

\textbf{\hl{✔ সঠিক উত্তর: খ) (২/৩)}}

\begin{quote}
\textbf{ব্যাখ্যা:} \emph{দেওয়া আছে, (১/৪)+(১/১২)+(১/৩)}

\emph{= (৩+১+৪/১২)}

\emph{= (৮/১২)}

\emph{= (২/৩)}
\end{quote}

\textbf{(21×21এর21/21÷21×21) এর সরল মান হবে-}

\textbf{\hl{ক) 441 ✔}}

খ) 1

গ) 0

ঘ) 21

\textbf{\hl{✔ সঠিক উত্তর: ক) 441}}

\begin{quote}
\textbf{ব্যাখ্যা:}
\emph{(21×21এর21/21÷21×21)=(21×(21×21)/(21÷21)×21)=(21×441/1×21)=441}
\end{quote}

\textbf{৩(৩/৮)+১(৫/৮)+৫(১/৮)= কত?}

\textbf{\hl{ক) ১০(১/৮) ✔}}

খ) ৯(১/৮)

গ) ১০(৭/৮)

ঘ) ৯(১/৮)

\textbf{\hl{✔ সঠিক উত্তর: ক) ১০(১/৮)}}

\begin{quote}
\textbf{ব্যাখ্যা:} \emph{৩(৩/৮)+১(৫/৮)+৫(১/৮)=(২৭/৮)+(১৩/৮)+(৪১/৮)}

\emph{= (২৭+১৩+৪১/৮)=(৮১/৮)=১০(১/৮)}
\end{quote}

\textbf{(175×175−2×175×75+75×75/175−75) কত?}

ক) 10

\textbf{\hl{খ) 100 ✔}}

গ) 200

\textbf{\hl{✔ সঠিক উত্তর: খ) 100}}

\begin{quote}
\textbf{ব্যাখ্যা:} \emph{লবকে a\^{}2 − 2ab + b\^{}2 = (a − b)\^{}2 সূত্রে
সাজিয়ে পাই,}

\emph{((175−75)\^{}2/175−75)=175−75=100}
\end{quote}

\textbf{১(২/৩)এর(১/৫)÷(১/৯) কত?}

ক) ৫

\textbf{\hl{খ) ৩ ✔}}

গ) ৪

\textbf{\hl{✔ সঠিক উত্তর: খ) ৩}}

\begin{quote}
\textbf{ব্যাখ্যা:} \emph{১(২/৩)এর(১/৫)÷(১/৯)}

\emph{= (৫/৩)এর(১/৫)÷(১/৯)}

\emph{= (১/৩)÷(১/৯)}

\emph{= (১/৩)×৯=৩}
\end{quote}

\textbf{(৩/৫)+০.১+০.০৫ এর মান-}

ক) (১৭/২০)

খ) (৭/২০)

\textbf{\hl{গ) (১৫/২০) ✔}}

ঘ) (১৩/২০)

\textbf{\hl{✔ সঠিক উত্তর: গ) (১৫/২০)}}

\begin{quote}
\textbf{ব্যাখ্যা:} \emph{(৩/৫)+০.১+০.০৫}

\emph{= (৩/৫)+(১/১০)+(১/২০)}

\emph{= (১২+২+১/২০)}

\emph{= (১৫/২০)}
\end{quote}

\textbf{((১/৩)×৬/(৪/৬)×(১২/২))+(৩/৪)= কত?}

ক) (১/১২)

\textbf{\hl{খ) ১(১/৪) ✔}}

গ) (৭/৪)

ঘ) (৪/৫)

\textbf{\hl{✔ সঠিক উত্তর: খ) ১(১/৪)}}

\begin{quote}
\textbf{ব্যাখ্যা:} \emph{((১/৩)×৬/(৪/৬)×(১২/২))+(৩/৪)}

\emph{= (২/৪)+(৩/৪)}

\emph{= (১/২)+(৩/৪)}

\emph{= (২+৩/৪)}

\emph{= (৫/৪)}

\emph{= ১(১/৪)}
\end{quote}

\textbf{১÷(৮/৯)((৫/৮)+(৩/৮))= কত?}

ক) ১

\textbf{\hl{খ) ১(১/৮) ✔}}

গ) (৮/৯)

ঘ) (১/৯)

\textbf{\hl{✔ সঠিক উত্তর: খ) ১(১/৮)}}

\begin{quote}
\textbf{ব্যাখ্যা:} \emph{১÷(৮/৯)((৫/৮)+(৩/৮))}

\emph{= ১÷(৮/৯)((৮/৮))}

\emph{= ১÷(৮/৯)×১}

\emph{= ১×(৯/৮)=(৯/৮)=১(১/৮)}
\end{quote}

\section{64. ``সংখ্যায়ন'' অধ্যায়: মৌলিক সূত্রাবলি, সমাধান ও সহজ কৌশলের
পূর্ণাঙ্গ
নির্দেশিকা}\label{ux9b8ux996ux9afux9dfux9a8-ux985ux9a7ux9afux9df-ux9aeux9b2ux995-ux9b8ux9a4ux9b0ux9acux9b2-ux9b8ux9aeux9a7ux9a8-ux993-ux9b8ux9b9ux99c-ux995ux9b6ux9b2ux9b0-ux9aaux9b0ux9a3ux999ux997-ux9a8ux9b0ux9a6ux9b6ux995}

প্রতিযোগিতামূলক পরীক্ষা যেমন---বিসিএস, প্রাথমিক শিক্ষক নিয়োগ, ব্যাংক, এনএসআই
কিংবা বিভিন্ন সরকারি মন্ত্রণালয়ের প্রশাসনিক কর্মকর্তা পদে নিয়োগ পরীক্ষার গণিত
অংশে ভালো করার মূল চাবিকাঠি হলো গণিতের মৌলিক ধারণা স্পষ্ট রাখা এবং সময়ের
সাশ্রয় করা। পাটিগণিতে \textbf{``সংখ্যায়ন''} বা সংখ্যা সংক্রান্ত অধ্যায়টি কোনো
বিচ্ছিন্ন বিষয় নয়, বরং এটি সমগ্র পাটিগণিতের প্রাণশক্তি। বিগত বছরের প্রশ্ন
বিশ্লেষণ করলে দেখা যায়, সংখ্যা সম্পর্কিত কিছু নির্দিষ্ট ছন্দের ও ধরনের অংক প্রায়
প্রতিটি পরীক্ষায় ঘুরেফিরে আসতে দেখা যায়।

সংখ্যা বিষয়ক এসব অংক সমাধান করার জন্য অনেকেই কেবল মুখস্থ
\textquotesingle শর্টটেকনিক\textquotesingle{} বা সংক্ষিপ্ত সূত্রের ওপর নির্ভর
করেন। কিন্তু পরীক্ষার হলে একটু ঘুরিয়ে প্রশ্ন আসলে বা প্রশ্নকর্তা সামান্য কৌশল
বদলালে শর্টটেকনিক মনে রাখা কঠিন হয়ে পড়ে। তাই গণিতে পূর্ণ নম্বর নিশ্চিত করতে
হলে প্রতিটি অংকের মূল লিখিত নিয়ম (Formal Solution) জানা বাধ্যতামূলক, যার পর
পরীক্ষার হলে সময় বাঁচাতে সংক্ষিপ্ত সূত্র বা শর্টটেকনিক (Shortcut Formula) প্রয়োগ
করতে হয়।

এই নিবন্ধে আমরা সংখ্যা সম্পর্কিত গণিতের বিভিন্ন সূত্র, এগুলোর বিস্তারিত ব্যাখ্যা,
প্রতিটি সূত্রের জন্য গাণিতিক উদাহরণ (লিখিত ও সূত্রভিত্তিক উভয় নিয়মে সমাধানসহ)
এবং সূত্রের বাইরে থাকা অত্যন্ত গুরুত্বপূর্ণ গাণিতিক প্রশ্নের বিস্তারিত সমাধান নিয়ে
আলোচনা করব।

\begin{itemize}
\item
  ক্রমিক সংখ্যার বর্গের অন্তর
\item
  যত বড়.....তত ছোট/ যত ছোট .... তত বড়
\item
  গুণফল ও একটি সংখ্যা থেকে অপর সংখ্যা
\item
  সংখ্যা সম্পর্কিত আরো কিছু গুরুত্বপূর্ণ প্রশ্ন ও সমাধান
\end{itemize}

\subsubsection{ক্রমিক সংখ্যার বর্গের
অন্তর}\label{ux995ux9b0ux9aeux995-ux9b8ux996ux9afux9b0-ux9acux9b0ux997ux9b0-ux985ux9a8ux9a4ux9b0}

পরপর দুটি পূর্ণসংখ্যাকে ক্রমিক সংখ্যা বলে (যেমন: x এবং x + 1)। এদের বর্গের
পার্থক্য দেওয়া থাকলে নিম্নোক্ত সূত্র দুটি প্রয়োগ করে খুব সহজেই সংখ্যা দুটি বের করা
যায়।

\textbf{সূত্র-১:} দুইটি ক্রমিক সংখ্যার বর্গের অন্তর (বা পার্থক্য) দেওয়া থাকলে, বড়
সংখ্যা নির্ণয়ের ক্ষেত্রে:

বড় সংখ্যা = \(\frac{\text{বর্গের}\ \text{অন্তর}\  + \ \text{১}}{\text{২}}\)

\textbf{সূত্র-২:} দুইটি ক্রমিক সংখ্যার বর্গের অন্তর (বা পার্থক্য) দেওয়া থাকলে,
ছোট সংখ্যা নির্ণয়ের ক্ষেত্রে:

ছোট সংখ্যা = \(\frac{\text{বর্গের}\ \text{অন্তর}\  - \ \text{১}}{\text{২}}\)

\textbf{প্রশ্ন-১: দুটি ক্রমিক পূর্ণ বর্গসংখ্যার বর্গের অন্তর ৯৩ হলে সংখ্যাদ্বয় কত?}

\textbf{মূল নিয়ম (লিখিত সমাধান):}

ধরি, ছোট সংখ্যাটি = x

তাহলে, বড় সংখ্যাটি হবে = x + ১

প্রশ্নমতে,

( x + ১)\textsuperscript{2} − x\textsuperscript{2} = ৯৩

⇒ x\textsuperscript{2} + ২x + ১ − x\textsuperscript{2} = ৯৩

⇒ ২x + ১ = ৯৩

⇒ ২x = ৯৩ − ১

⇒ ২x = ৯২

⇒ x = \(\frac{\text{৯২}}{\text{২}}\)

⇒ x = ৪৬

অতএব, ছোট সংখ্যাটি x = ৪৬

এবং বড় সংখ্যাটি x + ১ = ৪৬ + ১ = ৪৭

∴ সংখ্যাদ্বয় যথাক্রমে ৪৬ এবং ৪৭

\textbf{সূত্রের সাহায্যে সমাধান (শর্ট টেকনিক):}

এখানে বর্গের অন্তর = ৯৩

ছোট সংখ্যা = \(\frac{\text{বর্গের}\ \text{অন্তর}\  - \ \text{১}}{\text{২}}\)

= \(\frac{\text{৯৩}\  - \ \text{১}}{\text{২}}\)

= \(\frac{\text{৯২}}{\text{২}}\)

= ৪৬

বড় সংখ্যা = \(\frac{\text{বর্গের}\ \text{অন্তর}\  + \ \text{১}}{\text{২}}\)

= \(\frac{\text{৯৩}\  + \ \text{১}}{\text{২}}\)

= \(\frac{\text{৯৪}}{\text{২}}\)

= ৪৭ ∴ সংখ্যাদ্বয় যথাক্রমে ৪৬ এবং ৪৭

\textbf{প্রশ্ন-২: দুইটি ধারাবাহিক বা ক্রমিক সংখ্যার বর্গের পার্থক্য ৩৩। ক্ষুদ্রতম
বা ছোট সংখ্যাটি কত হবে?}

\textbf{মূল নিয়ম (লিখিত সমাধান):}

ধরি, ছোট সংখ্যাটি = x

তাহলে, বড় সংখ্যাটি হবে = x + ১

প্রশ্নমতে,

( x + ১)\textsuperscript{2} − x\textsuperscript{2} = ৩৩

⇒ x\textsuperscript{2} + ২x + ১ − x\textsuperscript{2} = ৩৩

⇒ ২x + ১ = ৩৩

⇒ ২x = ৩৩ − ১

⇒ ২x = ৩২

⇒ x = \(\frac{\text{৩২}}{\text{২}}\)

⇒ x = ১৬

অতএব, ক্ষুদ্রতম বা ছোট সংখ্যাটি x = ১৬

\textbf{সূত্রের সাহায্যে সমাধান (শর্ট টেকনিক):}

এখানে বর্গের অন্তর = ৩৩

ছোট সংখ্যা = \(\frac{\text{বর্গের}\ \text{অন্তর}\  - \ \text{১}}{\text{২}}\)

= \(\frac{\text{৩৩}\  - \ \text{১}}{\text{২}}\)

= \(\frac{\text{৩২}}{\text{২}}\)

= ১৬ অতএব, ক্ষুদ্রতম বা ছোট সংখ্যাটি x = ১৬

\subsubsection{যত বড়.....তত ছোট/ যত ছোট .... তত বড় উল্লেখ
থাকলে}\label{ux9afux9a4-ux9acux9dc.....ux9a4ux9a4-ux99bux99f-ux9afux9a4-ux99bux99f-....-ux9a4ux9a4-ux9acux9dc-ux989ux9b2ux9b2ux996-ux9a5ux995ux9b2}

যখন কোনো সংখ্যা দুটি নির্দিষ্ট সংখ্যার ঠিক মাঝামাঝি অবস্থানে থাকে, তখন প্রথম
সংখ্যাটি থেকে সেটির ব্যবধান এবং দ্বিতীয় সংখ্যাটি থেকে সেটির ব্যবধান সমান হয়।
ফলে নির্ণেয় সংখ্যাটি মূলত প্রদত্ত সংখ্যা দুটির গড় নির্দেশ করে।

সূত্র-৩: যত বড়.....তত ছোট/ যত ছোট .... তত বড় উল্লেখ থাকলে, সংখ্যাটি নির্ণয়ের
ক্ষেত্রে-\/-,

সংখ্যাটি = \(\frac{\text{প্রদত্ত}\ \text{সংখ্যা}\ \text{দুটির}\ \text{যোগফল}}{\text{২}}\)

\textbf{প্রশ্ন-১: একটি সংখ্যা ৬৫০ থেকে যত বড় ৮২০ থেকে তত ছোট। সংখ্যাটি কত?}

\textbf{মূল নিয়ম (লিখিত সমাধান):}

ধরি, নির্ণেয় সংখ্যাটি = x

সংখ্যাটি ৬৫০ থেকে যত বড়, অর্থাৎ ব্যবধান = (x - ৬৫০)

সংখ্যাটি ৮২০ থেকে তত ছোট, অর্থাৎ ব্যবধান = (৮২০ - x)

প্রশ্নমতে,

x - ৬৫০ = ৮২০ - x

বা, x + x = ৮২০ + ৬৫০

বা, 2x = ১৪৭০

বা, x = \(\frac{\text{১৪৭০}}{\text{২}}\)

বা, = ৭৩৫

∴ সংখ্যাটি ৭৩৫

\textbf{সূত্রের সাহায্যে সমাধান (শর্ট টেকনিক):}

নির্ণেয় সংখ্যা = \(\frac{\text{১ম}\ \text{সংখ্যা}\  + \ \text{২য়}\ \text{সংখ্যা}}{\text{২}}\)

= \(\frac{\text{৬৫০}\  + \ \text{৮২০}}{\text{২}}\)

= \(\frac{\text{১৪৭০}}{\text{২}}\)

= ৭৩৫ ∴ সংখ্যাটি ৭৩৫

\textbf{প্রশ্ন-২: একটি সংখ্যা ৯৯৯ থেকে যত ছোট ৭৯৭ থেকে তত বড়। সংখ্যাটি কত?}

\textbf{মূল নিয়ম (লিখিত সমাধান):}

ধরি, নির্ণেয় সংখ্যাটি = x

সংখ্যাটি ৭৯৭ থেকে যত বড়, অর্থাৎ ব্যবধান = (x - ৭৯৭)

সংখ্যাটি ৯৯৯ থেকে তত ছোট, অর্থাৎ ব্যবধান = (৯৯৯ - x)

প্রশ্নমতে,

x - ৭৯৭ = ৯৯৯ - x

বা, x + x = ৯৯৯ + ৭৯৭

বা, 2x = ১৭৯৬

বা, x = \(\frac{\text{১৭৯৬}}{\text{২}}\)

বা, = ৮৯৮

∴ সংখ্যাটি ৮৯৮

\textbf{সূত্রের সাহায্যে সমাধান (শর্ট টেকনিক):}

নির্ণেয় সংখ্যা = \(\frac{\text{১ম}\ \text{সংখ্যা}\  + \ \text{২য়}\ \text{সংখ্যা}}{\text{২}}\)

= \(\frac{\text{৭৯৭}\  + \ \text{৯৯৯}}{\text{২}}\)

= \(\frac{\text{১৭৯৬}}{\text{২}}\)

= ৮৯৮ ∴ সংখ্যাটি ৮৯৮

\subsubsection{গুণফল ও একটি সংখ্যা থেকে অপর
সংখ্যা}\label{ux997ux9a3ux9abux9b2-ux993-ux98fux995ux99f-ux9b8ux996ux9af-ux9a5ux995-ux985ux9aaux9b0-ux9b8ux996ux9af}

দুটি সংখ্যার গুণফলের মূল নীতি হলো ১ম সংখ্যা × ২য় সংখ্যা = গুণফল। সুতরাং, গুণফলকে
যেকোনো একটি সংখ্যা দিয়ে ভাগ করলে অপর সংখ্যাটি পাওয়া যায়।

সূত্র-৪: দুটি সংখ্যার গুণফল এবং একটি সংখ্যা দেওয়া থাকলে অপর সংখ্যাটি নির্ণয়ের
ক্ষেত্রে-

অপর সংখ্যাটি = \(\frac{\text{দুইটি}\ \text{সংখ্যার}\ \text{গুণফল}}{\text{একটি}\ \text{সংখ্যা}}\)

\textbf{প্রশ্ন-১: ২টি সংখ্যার গুণফল ৮.৮। একটি সংখ্যা ২.৭৫ হলে, অপর সংখ্যাটি
কত?}

\textbf{মূল নিয়ম (লিখিত সমাধান):}

ধরি, অপর সংখ্যাটি = y

প্রশ্নমতে, একটি সংখ্যা × অপর সংখ্যা = গুণফল

বা, ২.৭৫ × y = ৮.৮

বা, y = \(\frac{\text{৮}.\text{৮}}{\text{২}.\text{৭৫}}\)

বা, y = ৩.২

∴ অপর সংখ্যাটি = ৩.২

\textbf{সূত্রের সাহায্যে সমাধান (শর্ট টেকনিক):}

অপর সংখ্যাটি = \(\frac{\text{দুইটি}\ \text{সংখ্যার}\ \text{গুণফল}}{\text{একটি}\ \text{সংখ্যা}}\)

= \(\frac{\text{৮}.\text{৮}}{\text{২}.\text{৭৫}}\)

= ৩.২ ∴ অপর সংখ্যাটি = ৩.২

\textbf{প্রশ্ন-২: দুটি সংখ্যার গুণফল ২৩০৪। একটি সংখ্যার ৪ গুণ ৯৬ হলে, অপর
সংখ্যাটি কত?}

\textbf{মূল নিয়ম (লিখিত সমাধান):}

ধরি, প্রথম সংখ্যাটি = x এবং অপর সংখ্যাটি = y

দেওয়া আছে, প্রথম সংখ্যার ৪ গুণ = ৯৬

∴ প্রথম সংখ্যাটি = ৯৬ ÷ ৪ = ২৪

আবার, দুটি সংখ্যার গুণফল = ২৩০৪

প্রশ্নমতে, x × y = ২৩০৪

বা, ২৪ × y = ২৩০৪

বা, y = \(\frac{\text{২৩০৪}}{\text{২৪}}\) = ৯৬

∴ অপর সংখ্যাটি ৯৬।

\textbf{সূত্রের সাহায্যে সমাধান (শর্ট টেকনিক):}

এখানে, একটি সংখ্যা = ৯৬ ÷ ৪ = ২৪

অপর সংখ্যাটি = \(\frac{\text{দুইটি}\ \text{সংখ্যার}\ \text{গুণফল}}{\text{একটি}\ \text{সংখ্যা}}\)

= \(\frac{\text{২৩০৪}}{\text{২৪}}\)

= ৯৬ ∴ অপর সংখ্যাটি = ৯৬

\subsubsection{সংখ্যা সম্পর্কিত আরো কিছু গুরুত্বপূর্ণ প্রশ্ন ও
সমাধান}\label{ux9b8ux996ux9af-ux9b8ux9aeux9aaux9b0ux995ux9a4-ux986ux9b0-ux995ux99b-ux997ux9b0ux9a4ux9acux9aaux9b0ux9a3-ux9aaux9b0ux9b6ux9a8-ux993-ux9b8ux9aeux9a7ux9a8}

সংখ্যা সংক্রান্ত সকল অংক কিন্তু সরাসরি ওপরের ৪টি সূত্রের মাধ্যমে সমাধান করা যায়
না। কিছু অংক বীজগাণিতিক অনুপাত, মান নির্ণয়, লঘিষ্ঠ সাধারণ গুণনীয়ক (ল.সা.গু)
কিংবা ভগ্নাংশের নিয়মের ওপর প্রতিষ্ঠিত। নিচে এরকম কিছু গুরুত্বপূর্ণ প্রশ্নের লিখিত
সমাধান ও বিশ্লেষণ দেওয়া হলো:

\textbf{প্রশ্ন-১: একটি সংখ্যার তিনগুণের সাথে দ্বিগুণ যোগ করলে ৯০ হয়। সংখ্যাটি
কত ?}

\textbf{\hl{উত্তরঃ (গ) ১৮}}

\textbf{সমাধান:}

ধরি, সংখ্যাটি = `ক'

প্রশ্নমতে, ৩ক+২ক = ৯০

বা, ৫ক = ৯০

বা,ক = ৯০÷৫=১৮

সুতরাং সংখ্যাটি = ১৮

\textbf{প্রশ্ন-২: একটি ঋণাত্মক সংখ্যার দ্বিগুণের বর্গের সাথে ১৫ যোগ করলে ৪১৫
হয়। সংখ্যাটি কত ?}

\textbf{\hl{উত্তরঃ (খ) ১০}}

\textbf{সমাধান:}

ধরি, ঋণাত্মক সংখ্যাটি = -ক

প্রশ্নমতে, \({( - \text{২ক})}^{\text{২}}\) +১৫ = ৪১৫

বা , \({(\text{২ক})}^{\text{২}}\) = ৪১৫-১৫

বা, \({(\text{২ক})}^{\text{২}}\) = ৪০০

বা, \({(\text{২ক})}^{\text{২}}\) = \({(\text{২০})}^{\text{২}}\)

বা, ২ক = ২০

বা,ক = ১০

সুতরাং সংখ্যাটি = ১০

\textbf{প্রশ্ন-৩: দু'টি সংখ্যার বর্গের সমষ্টি ৪১, সংখ্যা দু'টির গুণফল ২০। সংখ্যা
দু'টির বর্গের বিয়োগফল কত ?}

\textbf{\hl{উত্তরঃ (খ) ৯}}

\textbf{সমাধান:}

মনে করি , সংখ্যা দুটি যথাক্রমে a ও b ।

প্রশ্নমতে, \(a^{2}\) + \(b^{2}\) = ৪১ আবার, ab = ২০

আমরা জানি, \({(a + b)}^{2}\) = \(a^{2}\) + \(b^{2}\) + 2ab

বা, \({(a + b)}^{2}\) = ৪১ + ২ × ২০

বা, \({(a + b)}^{2}\) = ৮১

বা, \({(a + b)}^{2}\) = \({(\text{৯})}^{\text{২}}\)

বা, a+b = ৯

আবার, \({(a\  - \ b)}^{2}\) = \(a^{2}\) + \(b^{2}\) - 2ab

বা, \({(a - b)}^{2}\) = ৪১ - ২ × ২০

বা, \({(a - b)}^{2}\) = ৪১ - ৪০

বা, \({(a - b)}^{2}\) = ১

বা, a-b = ১

সুতরাং \(a^{2}\) - \(b^{2}\) = (a+b)(a-b)

বা, \(a^{2}\) - \(b^{2}\) = ৯ × ১ = ৯

অর্থাৎ সংখ্যা দু'টির বর্গের বিয়োগফল = ৯

\textbf{মুখে মুখে সমাধান :} এখানে, \(\text{৪}^{\text{২}}\)+\(\text{৫}^{\text{২}}\) = ৪১ আবার, ৪×৫ =
২০।

সুতরাং সংখ্যা দু'টি হচ্ছে ৪ ও ৫।

তাই সংখ্যা দু'টির বর্গের বিয়োগফল = \(\text{৫}^{\text{২}}\) - \(\text{৪}^{\text{২}}\)= ৯।

\textbf{প্রশ্ন-৪: পরপর দশটি সংখ্যার প্রথম ৫টির যোগফল ৫৬০ হলে,শেষ ৫টির যোগফল
কত ?}

\textbf{\hl{উত্তরঃ (খ) ৫৮৫}}

\textbf{সমাধান:}

মনে করি,

প্রথম সংখ্যাটি = x

প্রশ্নমতে ,

x+x+ ১+x+ ২+x+ ৩+x+ ৪ = ৫৬০

বা, ৫x+১০ = ৫৬০

বা, ৫x = ৫৫০

বা, x = ১১০

.'. শেষ ৫ টি সংখ্যার যোগফল

= x+৫+x+৬+x+৭+x+৮+x+৯

= ৫x+৩৫

= (৫×১১০)+৩৫ {[} x এর মান ১১০ বসিয়ে{]}

= ৫৮৫ অথবা

\textbf{মুখে মুখে সমাধান:}

এরকম থাকলে শেষ পদসংখ্যাকে বর্গ করে প্রথম যোগফলের সাথে যোগ করলেই শেষ সংখ্যা
গুলোর যোগফল পাওযা যাবে।যেমন \(\text{৫}^{\text{২}}\) = ২৫।

সুতরাং শেষ ৫টির যোগফল = ৫৬০ + ২৫ = ৫৮৫

সময় অবশিষ্ট আছে:

\section{65. বাস্তব সংখ্যার পূর্ণাঙ্গ গাইড: সংজ্ঞা, প্রকারভেদ ও শর্টকাট
টিপস}\label{ux9acux9b8ux9a4ux9ac-ux9b8ux996ux9afux9b0-ux9aaux9b0ux9a3ux999ux997-ux997ux987ux9a1-ux9b8ux99cux99e-ux9aaux9b0ux995ux9b0ux9adux9a6-ux993-ux9b6ux9b0ux99fux995ux99f-ux99fux9aaux9b8}

বাস্তব সংখ্যা দ্বারা শুধুমাত্র সংখ্যাকে নির্দেশ করলে ও এটি আমাদের জীবনের অত্যন্ত
গুরুত্বপূর্ণ একটি দিক। সংখ্যা আমাদের জীবনের বিভিন্ন ক্ষেত্রে সহায়তা করে, সংখ্যা
ব্যবহার করে আমরা ব্যাক্তিগত জীবন থেকে শুরু করে প্রকৃতি ও বিজ্ঞানের জগত পর্যন্ত
পর্যালোচনা করে থাকি।

উদাহরণস্বরূপ, একটি মানুষের উচ্চতা, ওজন, বয়স ইত্যাদি সংখ্যার মাধ্যমে প্রকাশ করা
হয়। আর আর্থিক দিক থেকে বললে, আমাদের প্রতিটি আর্থিক লেনদেন বা লেনদেনের বিবরণ
একটি সংখ্যার মাধ্যমে প্রকাশ করা হয়।যাই হোক বিভিন্ন প্রতিযোগিতামূলক পরীক্ষার জন্য
বাস্তব সংখ্যা খুবই গুরুত্বপূর্ণ।বাস্তব সংখ্যার পরিপূর্ণ ধারণা ছাড়া গণিতের সমস্যার
সমাধান করতে পারব না। আজকের পাঠে আমরা এটি নিয়ে এ টু জেড পরিষ্কার ধারণা
দেয়ার চেষ্টা করব। সর্বশেষে আমরা এই অধ্যায়ে বিগত সালের প্রশ্ন অনুশীলন করার সুযোগ
করে দিয়েছি। যা আপনার প্রস্তুতিতে সহায়ক ভূমিকা পালন করবে বলে আমাদের বিশ্বাস।

\begin{itemize}
\item
  অংক ও সংখ্যা
\item
  ১ থেকে ১০০ পর্যন্ত মৌলিক সংখ্যা
\item
  ১০১ থেকে ২০০ পর্যন্ত মৌলিক সংখ্যা
\item
  মৌলিক ও অন্যান্য সংখ্যা
\item
  ভগ্নাংশ সংখ্যা (Fractional Number)
\item
  দশমিক ভগ্নাংশ
\item
  মূলদ ও অমূলদ সংখ্যা
\item
  বাস্তব সংখ্যার শ্রেণিবিন্যাস
\end{itemize}

\subsubsection{অংক ও সংখ্যা}\label{ux985ux995-ux993-ux9b8ux996ux9af}

\textbf{অংকঃ} Digit অর্থ অংক।মানুষ হিসাব নিকাশ ও গণনা কার্যের জন্য যে সকল
প্রতীক বা চিহ্ন ব্যবহার করে সেগুলোই অংক। গণিতে ০ থেকে ৯ পর্যন্ত মোট ১০টি অংক
রয়েছে। অংক দুই ধরনের যথা-

১। সার্থক অংক---১,২,৩,৪,৫,৬,৭,৮,৯ ।

২। সহকারী অংক--- ০ ।

\textbf{সংখ্যা (Number):} অংক দ্বারা পরিমাণ বুঝালে সংখ্যা হবে।এক বা একাধিক
অংক মিলে সংখ্যা তৈরি হয়।একক অংক হিসেবে পরিমাণজ্ঞাপক হলে সার্থক অংকগুলোকে
সংখ্যা বলা হয়।অর্থাৎ ১,২,৩,৪,৫,৬,৭,৮,৯ অংকগুলো দ্বারা পরিমাণ বোঝালে সেগুলোকে
সংখ্যাও বলা হয়।

\textbf{স্বাভাবিক সংখ্যা (Natural Number):} ১,২,৩,৪,\ldots{} ইত্যাদি
স্বাভাবিক সংখ্যা। অর্থাৎ ১ থেকে শুরু অসীম।যার কোন শেষ নেই।তবে এগুলোকে ধনাত্নক
বা ঋণাত্নক বলে চিহ্নিত করা হয় না। এগুলো সবসময় স্বাভাবিক । স্বাভাবিক সংখ্যার
সেটকে N দ্বারা প্রকাশ হয়।

\textbf{মৌলিক সংখ্যা ( Prime Number ):} যে সংখ্যাকে শুধু ১ এবং ঐ সংখ্যা ছাড়া
অন্য কোনো সংখ্যা দ্বারা ভাগ করা যায় না তাকে মৌলিক সংখ্যা বলে। আরো সহজভাবে
বললে, যে সংখ্যাকে অন্য কোনো সংখ্যা দ্বারা ভাগ করা যায় না তাকে মৌলিক সংখ্যা
বলে। যেমনঃ ৩১ একটি মৌলিক সংখ্যা। উল্লেখ্য, ১ মেীলিক সংখ্যা নয়, যেহেতু এর
উৎপাদক ১টি।মেীলিক সংখ্যার দুইটি উৎপাদক থাকে যথা ১ এবং ঐ সংখ্যা নিজে।ঋণাত্নক
সংখ্যা কখনো মৌলিক হতে পারে না এবং মৌলিক সংখ্যা ২ থেকে শুরু হয়। একমাত্র জোড়
মৌলিক সংখ্যা হল ২ ।

১ থেকে ১০০ পর্যন্ত মোট ২৫টি মৌলিক সংখ্যা এবং ১০১ থেকে ২০০ পর্যন্ত মোট ২১টি
মৌলিক সংখ্যা রয়েছে।

\subsubsection{১ থেকে ১০০ পর্যন্ত মৌলিক
সংখ্যা}\label{ux9a5ux995-ux9e7ux9e6ux9e6-ux9aaux9b0ux9afux9a8ux9a4-ux9aeux9b2ux995-ux9b8ux996ux9af}

{\def\LTcaptype{none} % do not increment counter
\begin{longtable}[]{@{}
  >{\raggedright\arraybackslash}p{(\linewidth - 6\tabcolsep) * \real{0.2006}}
  >{\raggedright\arraybackslash}p{(\linewidth - 6\tabcolsep) * \real{0.2006}}
  >{\raggedright\arraybackslash}p{(\linewidth - 6\tabcolsep) * \real{0.2006}}
  >{\raggedright\arraybackslash}p{(\linewidth - 6\tabcolsep) * \real{0.2006}}@{}}
\toprule\noalign{}
\begin{minipage}[b]{\linewidth}\raggedright
\textbf{সীমা}
\end{minipage} & \begin{minipage}[b]{\linewidth}\raggedright
\textbf{মৌলিক সংখ্যা}
\end{minipage} & \begin{minipage}[b]{\linewidth}\raggedright
\textbf{মোট}
\end{minipage} & \begin{minipage}[b]{\linewidth}\raggedright
\textbf{সর্বমোট}
\end{minipage} \\
\midrule\noalign{}
\endhead
\bottomrule\noalign{}
\endlastfoot
১ থেকে ১০ & ২,৩,৫,৭ & ৪ টি & (১-১০) ৪ টি \\
১১ থেকে ২০ & ১১,১৩,১৭,১৯ & ৪ টি & (১-২০) ৮ টি \\
২১ থেকে ৩০ & ২৩,২৯ & ২ টি & (১-৩০) ১০ টি \\
৩১ থেকে ৪০ & ৩১,৩৭ & ২ টি & (১-৪০) ১২ টি \\
৪১ থেকে ৫০ & ৪১,৪৩,৪৭ & ৩ টি & (১-৫০) ১৫ টি \\
৫১ থেকে ৬০ & ৫৩,৫৯ & ২ টি & (১-৬০) ১৭ টি \\
৬১ থেকে ৭০ & ৬১,৬৭ & ২ টি & (১-৭০) ১৯ টি \\
৭১ থেকে ৮০ & ৭১,৭৩,৭৯ & ৩ টি & (১-৮০) ২২ টি \\
৮১ থেকে ৯০ & ৮৩,৮৯ & ২ টি & (১-৯০) ২৪ টি \\
৯১ থেকে ১০০ & ৯৭ & ১ টি & (১-১০০) ২৫ টি \\
\end{longtable}
}

\subsubsection{১০১ থেকে ২০০ পর্যন্ত মৌলিক
সংখ্যা}\label{ux9a5ux995-ux9e8ux9e6ux9e6-ux9aaux9b0ux9afux9a8ux9a4-ux9aeux9b2ux995-ux9b8ux996ux9af}

{\def\LTcaptype{none} % do not increment counter
\begin{longtable}[]{@{}
  >{\raggedright\arraybackslash}p{(\linewidth - 6\tabcolsep) * \real{0.2006}}
  >{\raggedright\arraybackslash}p{(\linewidth - 6\tabcolsep) * \real{0.2006}}
  >{\raggedright\arraybackslash}p{(\linewidth - 6\tabcolsep) * \real{0.2006}}
  >{\raggedright\arraybackslash}p{(\linewidth - 6\tabcolsep) * \real{0.2006}}@{}}
\toprule\noalign{}
\begin{minipage}[b]{\linewidth}\raggedright
\textbf{সীমা}
\end{minipage} & \begin{minipage}[b]{\linewidth}\raggedright
\textbf{মৌলিক সংখ্যা}
\end{minipage} & \begin{minipage}[b]{\linewidth}\raggedright
\textbf{মোট}
\end{minipage} & \begin{minipage}[b]{\linewidth}\raggedright
\textbf{সর্বমোট}
\end{minipage} \\
\midrule\noalign{}
\endhead
\bottomrule\noalign{}
\endlastfoot
১০১ থেকে ১১০ & ১০১,১০৩,১০৭,১০৯ & ৪ টি & (১০১-১১০) ৪ টি \\
১১১ থেকে ১২০ & ১১৩ & ১ টি & (১০১-১২০) ৫ টি \\
১২১ থেকে ১৩০ & ১২৭ & ১ টি & (১০১-১৩০) ৬ টি \\
১৩১ থেকে ১৪০ & ১৩১,১৩৭,১৩৯ & ৩ টি & (১০১-১৪০) ৯ টি \\
১৪১ থেকে ১৫০ & ১৪৯ & ১ টি & (১০১-১৫০) ১০ টি \\
১৫১ থেকে ১৬০ & ১৫১,১৫৭ & ২ টি & (১০১-১৬০) ১২ টি \\
১৬১ থেকে ১৭০ & ১৬৩,১৬৭ & ২ টি & (১০১-১৭০) ১৪ টি \\
১৭১ থেকে ১৮০ & ১৭৩,১৭৯ & ২ টি & (১০১-১৮০) ১৬ টি \\
১৮১ থেকে ১৯০ & ১৮১ & ১ টি & (১০১-১৯০) ১৭ টি \\
১৯১ থেকে ২০০ & ১৯১,১৯৩,১৯৭,১৯৯ & ৪ টি & (১০১-২০০) ২১ টি \\
\end{longtable}
}

\subsubsection{মৌলিক ও অন্যান্য
সংখ্যা}\label{ux9aeux9b2ux995-ux993-ux985ux9a8ux9afux9a8ux9af-ux9b8ux996ux9af}

\textbf{মৌলিক দ্বিজোড়:} পরস্পর দুটি মৌলিক সংখ্যার ব্যবধান ২ হলে তাকে মৌলিক
দ্বিজোড় বলে।যেমন-৫,৭,১১,১৩ ইত্যাদি। ১ থেকে ১০০ পর্যন্ত ৮টি মৌলিক দ্বিজোড়
রয়েছে।

\textbf{মৌলিক ত্রিজোড়:} পরস্পর তিনটি মৌলিক সংখ্যার অন্তর ২ হলে তাদের মৌলিক
ত্রিজোড় বলে।যেমন- ৩,৫,৭ ।

\textbf{যৌগিক সংখ্যাঃ} যে সকল সংখ্যাকে ১ এবং সেই সংখ্যা ছাড়াও অন্য এক বা
একাধিক সংখ্যা দ্বারা ভাগ করা যায় তাকে যৌগিক সংখ্যা বলে।যেমন- ৪,৬,৮,৯,১০
ইত্যাদি।

\textbf{সহমৌলিক সংখ্যা ( Co-Prime number ) :} দুইটি স্বাভাবিক সংখ্যার
গ.সা.গু. ১ হলে তাদেরকে পরস্পরের সহমৌলিক বলা হয়। যেমন ৮ ও ১৫ পরস্পর সহমৌলিক
। কারণ এদের গ.সা.গু. ১ ।তবে সংখ্যা মৌলিক বা যৌগিক যাই হোক না কেন সংখ্যাগুলোর
গ.সা.গু বাদ দিয়ে যে উৎপাদকগুলো থাকে তা সর্বদা সহমৌলিক।উদাহরণস্বরুপ- ২০ ও ১৫।
২০ এর উৎপাদক হল ১,২,৪,৫,২০ এবং ১৫ এর উৎপাদক হল ১,৩,৫,১৫।এখানে ১ বাদে
উভয়সংখ্যার উৎপাদকগুলো পরস্পরের সহমৌলিক।

\textbf{জোড় ও বিজোড় সংখ্যা (Even \& Odd number):} যে সমস্ত সংখ্যা ২ দ্বারা
ভাগ করা যায় অর্থাৎ যে সমস্ত সংখ্যার এককের ঘরের অঙ্ক ২, ৪, ৬, ৮ বা ০ তারা জোড়
সংখ্যা। যেমন---২, ৫৪, ১৫৬, ৯০৯৮, ১৫৮৩৪০, ইত্যাদি। উল্লেখ্য, ০ একটি জোড়
সংখ্যা। আবার যে সব সংখ্যা ২ দ্বারা ভাগ করা যায় না অর্থাৎ যে সব সংখ্যার এককের
ঘরের অঙ্ক ১, ৩, ৫, ৭ বা ৯, তারা বিজোড় সংখ্যা।

শুধু কোন অক্ষর যেমন- X, Y, Z জোড় সংখ্যা না আবার বিজোড় সংখ্যা ও না।আপনি যখন
যাকে যা ধরবেন,তা-ই হবে। কিন্তু জোড় কিংবা বিজোড় যেকোনো সংখ্যাকেই ২ বা যেকোন
জোড়সংখ্যা দ্বারা গুণ করলে গুণফল সবসময়ই জোড় সংখ্যা হয়।যেমন-2X জোড় নাকি বিজোড়?
এখানে X এর মান বিজোড় সংখ্যা ১ ধরলে 2X এর মান ২ যা জোড় সংখ্যা।আবার X এর মান
জোড় সংখ্যা ২ ধরলে 2X এর মান হয় ৪। এটিও জোড় সংখ্যা।কিন্তৃ বিজোড় সংখ্যা দিয়ে
জোড় সংখ্যাকে গুণ করলে গুণফল জোড় হয়।আবার বিজোড় সংখ্যা দিয়ে বিজোড় সংখ্যাকে গুণ
করলে গুণফল বিজোড় হয়।

\textbf{পূর্ণসংখ্যা (Integer):} যে সমস্ত সংখ্যার কোন ভগ্নাংশ থাকে না তাদেরকেই
পূর্ণ সংখ্যা বলে । আরো সহজভাবে বললে, শূন্যসহ সকল ধনাত্মক ও ঋণাত্মক অখন্ডক
সংখ্যাকে পূর্ণসংখ্যা বলে। যেমন: ......৩,-২,-১,০,১,২,৩.......ইত্যাদি।
পূর্ণসংখ্যার সেটকে Z দ্বারা প্রকাশ হয়। পূর্ণসংখ্যার সংখ্যা অসীম।

\subsubsection{ভগ্নাংশ সংখ্যা (Fractional
Number)}\label{ux9adux997ux9a8ux9b6-ux9b8ux996ux9af-fractional-number}

\textbf{ভগ্নাংশ :} কোন বস্তু বা পরিমাণের অংশ বা ভাগ নির্দেশ করতে যে সংখ্যা
শ্রেণি ব্যবহৃত হয় তাকে ভগ্নাংশ সংখ্যা বলে। যেমন- \(\frac{\text{৩}}{\text{৫}}\) ।

ভগ্নাংশ প্রধানত দুই প্রকার যথা- ১। সাধারণ ভগ্নাংশ ২। দশমিক ভগ্নাংশ।

\textbf{সাধারণ ভগ্নাংশ:} যে সব সংখ্যাকে \(\frac{p}{q}\) আকারে প্রকাশ করা
যায় তাকে সাধারণ ভগ্নাংশ সংখ্যা বলে।সাধারণ ভগ্নাংশের উদাহরণ- \(\frac{\text{১}}{\text{৪}}\)
।

সাধারণ ভগ্নাংশ আবার তিন প্রকার যথা- ১। প্রকৃত ভগ্নাংশ ২। অপ্রকৃত ভগ্নাংশ ৩। মিশ্র
ভগ্নাংশ।

\textbf{প্রকৃত ভগ্নাংশ:} যে সব ভগ্নাংশের হর বড় এবং লব ছোট তাকে প্রকৃত ভগ্নাংশ
বলে।উল্লেখ্য, ভগ্নাংশের নিচের সংখ্যাটিকে হর এবং উপরের সংখ্যাটিকে লব বলে।যেমন-
\(\frac{\text{৩}}{\text{৮}}\) ।

\textbf{অপ্রকৃত ভগ্নাংশ:} যে সব ভগ্নাংশের হর ছোট এবং লব বড় তাকে অপ্রকৃত ভগ্নাংশ
বলে।যেমন- \(\frac{\text{১৫}}{\text{৮}}\) ।

\textbf{মিশ্র ভগ্নাংশ:} পূর্ণসংখ্যার সাথে প্রকৃত ভগ্নাংশ যুক্ত থাকলে তাকে মিশ্র
ভগ্নাংশ বলে। যেমন- \(\frac{\text{৩}}{\text{১৫}}\) ।

\subsubsection{দশমিক
ভগ্নাংশ}\label{ux9a6ux9b6ux9aeux995-ux9adux997ux9a8ux9b6}

\textbf{দশমিক ভগ্নাংশ:} যে ভগ্নাংশকে দশমিক (.) চিহ্নের সাহায্যে প্রকাশ করা
হয়, তাকে দশমিক ভগ্নাংশ বলে। দশমিক বিন্দুর আগের অংশ পূর্ণসংখ্যা এবং পরের অংশ
ভগ্নাংশ । দশমিক বিন্দু থেকে বামদিকের অঙ্কগুলোর স্থানীয় মন পর্যায়ক্রমে একক, দশক ,
শতক ইত্যাদি। এবং ডানদিকের অঙ্কগুলোর স্থানীয় মান পর্যায়ক্রমে দশমাংশ, শতাংশ,
সহস্রাংশ ইত্যাদি।দশমিক ভগ্নাংশ দুই প্রকার যথা- সসীম দশমিক ও অসীম দশমিক।অসীম
দশমিক আবার দুই প্রকার যথা অসীম আবৃত এবং অসীম অনাবৃত।

\textbf{সসীম দশমিক ভগ্নাংশ:} যে সংখ্যাকে ভাগ করলে দশমিক বিন্দুর পরের
সংখ্যাগুলো সসীম বা নির্দিষ্ট ফলাফল বের হয়ে শেষ হয়ে যায়, এমন ভগ্নাংশকে সসীম
দশমিক ভগ্নাংশ বলে। যেমনঃ ৩.৫। এখানে, দশমিক সংখ্যার পর একটা নির্দিষ্ট সংখায়
বের হয়ে ভাগ শেষ হয়ে গেছে মানে দশমিক সংখ্যাটা সসীমতায় আছে। তাই একে সসীম
দশমিক ভগ্নাংশ বলে।

\textbf{অসীম আবৃত দশমিক ভগ্নাংশ:} যেসব দশমিক ভগ্নাংশে দশমিক বিন্দুর ডানে একটি
অঙ্ক ক্রমান্বয়ে বারবার বা একাধিক অঙ্ক পর্যায়ক্রমে বারবার আসে, তাদের আবৃত বা
পৌনঃপুনিক দশমিক ভগ্নাংশ বলে। আবৃত বা পৌনঃপুনিক দশমিক ভগ্নাংশে যে অংশ বারবার
অর্থাৎ পুনঃপুনঃ হয়, তাকে আবৃত অংশ বলে। যেমন: ৩.৩৩৩৩..., ১০.২৩৪৫৬৪৫৬...
ইত্যাদি আবৃত বা পৌনঃপুনিক দশমিক ভগ্নাংশ।

\textbf{অসীম অনাবৃত দশমিক ভগ্নাংশ:} যে অসীম দশমিক ভগ্নাংশে দশমিক চিহ্নের
ডানদিকের অঙ্ক কখনো শেষ হয় না, অর্থাৎ দশমিক চিহ্নের ডানদিকের অঙ্কগুলো সসীম হবে
না অথবা অংশবিশেষ বারবার আসবে না। যেমন : 1.4142135......, 2.8284271.....
ইত্যাদি অসীম অনাবৃত দশমিক ভগ্নাংশ। সসীম দশমিক ও আবৃত্ত দশমিক ভগ্নাংশ মূলদ সংখ্যা
এবং অসীম অনাবৃত দশমিক ভগ্নাংশ অমূলদ সংখ্যা। সহজ কথায় অসীম অনাবৃত দশমিক বাদে
বাকি সবই মূলদ সংখ্যা। একমাত্র এটিই অমূলদ সংখ্যা।

\subsubsection{মূলদ ও অমূলদ
সংখ্যা}\label{ux9aeux9b2ux9a6-ux993-ux985ux9aeux9b2ux9a6-ux9b8ux996ux9af}

\textbf{মূলদ সংখ্যা (Rational Number):} আকারের যেকোনো সংখ্যাকে মূলদ সংখ্যা
বলে, যখন p ও q পূর্ণসংখ্যা এবং q ≠ 0। যেকোনো মূলদ সংখ্যাকে দুইটি সহমৌলিক
সংখ্যার অনুপাত হিসেবে ও লেখা যায়। সকল পূর্ণসংখ্যা ও ভগ্নাংশই মূলদ সংখ্যা।

\textbf{অমূলদ সংখ্যা (Irrational Number):} যে সকল সংখ্যাকে আকারে প্রকাশ করা
যায় না, যেখানে p ও q পূর্ণসংখ্যা এবং q ≠ 0, সে সকল সংখ্যাকে অমূলদ সংখ্যা বলে।
পূর্ণবর্গ নয় এরূপ যে কোনো স্বাভাবিক সংখ্যার বর্গমূল কিংবা তার ভগ্নাংশ একটি অমূলদ
সংখ্যা। যেমন: √2 = 1.414213\ldots., কোনো অমূলদ সংখ্যাকে দুইটি পূর্ণসংখ্যার
অনুপাত হিসেবে প্রকাশ করা যায় না।

\subsubsection{বাস্তব সংখ্যার
শ্রেণিবিন্যাস}\label{ux9acux9b8ux9a4ux9ac-ux9b8ux996ux9afux9b0-ux9b6ux9b0ux9a3ux9acux9a8ux9afux9b8}

\textbf{বাস্তব সংখ্যা (Real Number):} সকল মূলদ ও অমূলদ সংখ্যাকে বাস্তব সংখ্যা
বলা হয়।

\textbf{ধনাত্নক সংখ্যা ( Positive Number ) :} শূন্য থেকে বড় সকল বাস্তব
সংখ্যাকে ধনাত্নক সংখ্যা বলা হয়।

ঋণাত্নক সংখ্যা ( Negative Number ) : শূন্য থেকে ছোট সকল বাস্তব সংখ্যাকে ঋণাত্নক
সংখ্যা বলা হয়।

\textbf{অঋণাত্নক সংখ্যা (Non- Negative Number) :} শূন্যসহ সকল ধনাত্নক সংখ্যাকে
অঋণাত্নক সংখ্যা বলা হয়।

নিচের চিত্রে বাস্তব সংখ্যার শ্রেণিবিন্যাস দেখানো হল:

সময় অবশিষ্ট আছে:

\section{66. বাস্তব সংখ্যা MCQ প্রশ্ন উত্তর
অনুশীলন}\label{ux9acux9b8ux9a4ux9ac-ux9b8ux996ux9af-mcq-ux9aaux9b0ux9b6ux9a8-ux989ux9a4ux9a4ux9b0-ux985ux9a8ux9b6ux9b2ux9a8}

\textbf{নিচের কোনটি অমূলদ সংখ্যা?}

ক) 0.4

খ) √(9)

গ) 5.6•39•

ঘ) √((27/48))

\textbf{\hl{ঙ) সঠিক উত্তর নেই ✔}}

\textbf{\hl{✔ সঠিক উত্তর: ঙ) সঠিক উত্তর নেই}}

\begin{quote}
\textbf{ব্যাখ্যা:} \emph{আমরা জানি, শূন্য, সকল স্বাভাবিক সংখ্যা, ভগ্নাংশ এবং
পৌনঃপুনিক সংখ্যা মূলদ সংখ্যা। প্রশ্নে প্রদত্ত সকল সংখ্যাই মূলদ সংখ্যা।}

\emph{(ক) 0.4 = (4/10) = (2/5) (সসীম দশমিক সংখ্যা, তাই মূলদ)}

\emph{(খ) √(9) = √(3\^{}2) = 3 (পূর্ণসংখ্যা, তাই মূলদ)}

\emph{(গ) 5.6•39• (পৌনঃপুনিক বা আবৃত্ত দশমিক সংখ্যা, তাই মূলদ)}

\emph{(ঘ) √((27/48)) = √((3×9/3×16)) = √((9/16)) = (3/4) (মূলদ সংখ্যা)}

\emph{∴ বিকল্পগুলোর মধ্যে কোনটিই অমূলদ সংখ্যা নয় (প্রশ্নে সঠিক উত্তর দেওয়া
নেই)।}
\end{quote}

\textbf{নিচের কোনটি অমূলদ সংখ্যা?}

\textbf{\hl{ক) (√(27)/3) ✔}}

খ) (3√(125)/5)

গ) (4√(81)/4)

ঘ) (5√(32)/8)

\textbf{\hl{✔ সঠিক উত্তর: ক) (√(27)/3)}}

\begin{quote}
\textbf{ব্যাখ্যা:} \emph{ক. (√(27)/3) = (√(9×3)/3) = (3√(3)/3) = √(3) (যা
একটি অমূলদ সংখ্যা)}

\emph{খ. (3√(125)/5) = (5\^{}3/3)5 = (5/5) = 1 (মূলদ সংখ্যা)}

\emph{গ. (4√(81)/4) = (3\^{}4/4)4 = (3/4) (মূলদ সংখ্যা)}

\emph{ঘ. (5√(32)/8) = (2\^{}5/5)8 = (2/8) = (1/4) (মূলদ সংখ্যা)}

\emph{∴ (√(27)/3) একটি অমূলদ সংখ্যা।}
\end{quote}

\textbf{নিচের কোনটি অমূলদ সংখ্যা?}

\textbf{\hl{ক) ১.১০১০০১০০০১..... ✔}}

খ) ১.১০১০১০১.....

গ) ১.১০০১০০১০০১.....

ঘ) ১.১১১....

\textbf{\hl{✔ সঠিক উত্তর: ক) ১.১০১০০১০০০১.....}}

\begin{quote}
\textbf{ব্যাখ্যা:} \emph{অসীম অনাবৃত্ত দশমিক সংখ্যাকে অমূলদ সংখ্যা বলে।}

\emph{এখানে ১.১০১০০১০০০১..... একটি অসীম অনাবৃত্ত দশমিক সংখ্যা, কারণ দশমিকের
পর অঙ্কগুলো কোনো পুনরাবৃত্তি না ঘটিয়ে অসীম পর্যন্ত বিস্তৃত।}

\emph{∴ ১.১০১০০১০০০১..... একটি অমূলদ সংখ্যা।}
\end{quote}

\textbf{কোনটি মূলদ সংখ্যা?}

ক) √(5)

খ) √(7)

গ) √(10)

\textbf{\hl{ঘ) √(4) ✔}}

\textbf{\hl{✔ সঠিক উত্তর: ঘ) √(4)}}

\begin{quote}
\textbf{ব্যাখ্যা:} \emph{ক. √(5) ≈ 2.2360679... (পূর্ণবর্গ নয়, তাই এটি অমূলদ
সংখ্যা)}

\emph{খ. √(7) ≈ 2.6457513... (পূর্ণবর্গ নয়, তাই এটি অমূলদ সংখ্যা)}

\emph{গ. √(10) ≈ 3.1622776... (পূর্ণবর্গ নয়, তাই এটি অমূলদ সংখ্যা)}

\emph{ঘ. √(4) = 2 (পূর্ণসংখ্যা, তাই এটি একটি মূলদ সংখ্যা)}

\emph{∴ সঠিক উত্তর (ঘ)।}
\end{quote}

\textbf{√(১৫.৬৩২৩)=?}

ক) ৩.৮৫

\textbf{\hl{খ) ৩.৯৫ ✔}}

গ) ৩.৭৫

ঘ) ৩.২০

\textbf{\hl{✔ সঠিক উত্তর: খ) ৩.৯৫}}

\begin{quote}
\textbf{ব্যাখ্যা:} \emph{√(১৫.৬৩২৩)=√((১৫৬৩২৩/১০০০০))=(√(১৫৬৩২৩)/১০০)}

\emph{এখন, ১৫৬৩২৩ এর বর্গমূল ৩৯৫ (প্রায়)।}

\emph{∴ √(১৫.৬৩২৩)=(৩৯৫/১০০)=৩.৯৫}
\end{quote}

\textbf{নিচের কোনটি মূলদ সংখ্যা}

\textbf{\hl{ক) ৩√(৮) ✔}}

খ) √(২)

গ) ৩√(৭)

ঘ) (√(৫)/৮)

\textbf{\hl{✔ সঠিক উত্তর: ক) ৩√(৮)}}

\begin{quote}
\textbf{ব্যাখ্যা:} \emph{৩√(৮)=৮\^{}((১/৩))=(২\^{}৩)\^{}((১/৩))=২}

\emph{যেহেতু ২ একটি মূলদ সংখ্যা, তাই ৩√(৮) একটি মূলদ সংখ্যা।}
\end{quote}

\textbf{√(২) সংখ্যাটি কোন ধরনের সংখ্যা?}

ক) মূলদ

\textbf{\hl{খ) অমূলদ ✔}}

গ) জটিল

ঘ) বাস্তব

\textbf{\hl{✔ সঠিক উত্তর: খ) অমূলদ}}

\begin{quote}
\textbf{ব্যাখ্যা:} \emph{পূর্ণবর্গ নয় এমন স্বাভাবিক সংখ্যার বর্গমূলকে অমূলদ সংখ্যা
বলে।∴ √(২) সংখ্যাটি অমূলদ।}
\end{quote}

\textbf{৭√(৩) সংখ্যা কোন ধরনের সংখ্যা?}

ক) জটিল সংখ্যা

খ) মূলদ সংখ্যা

\textbf{\hl{গ) অমূলদ সংখ্যা ✔}}

ঘ) বাস্তব সংখ্যা।

\textbf{\hl{✔ সঠিক উত্তর: গ) অমূলদ সংখ্যা}}

\begin{quote}
\textbf{ব্যাখ্যা:} \emph{যেহেতু √(৩) একটি অমূলদ সংখ্যা, সেহেতু ৭√(৩) একটি অমূলদ
সংখ্যা।}
\end{quote}

\textbf{৩√(৫) সংখ্যাটি কি সংখ্যা?}

ক) একটি মৌলিক সংখ্যা

খ) একটি পূর্ণ সংখ্যা

গ) একটি মূলদ সংখ্যা

\textbf{\hl{ঘ) একটি অমূলদ সংখ্যা ✔}}

\textbf{\hl{✔ সঠিক উত্তর: ঘ) একটি অমূলদ সংখ্যা}}

\begin{quote}
\textbf{ব্যাখ্যা:} \emph{৩√(৫) = ৩√(মৌলিক সংখ্যা) = অমূলদ সংখ্যা}
\end{quote}

\textbf{√(২) এবং √(৩) এর মধ্যবর্তী একটি মূলদ সংখ্যা হবে-}

ক) (√(২)+√(৩)/২)

খ) (√(২).√(৩)/২)

গ) ১.৮

\textbf{\hl{ঘ) ১.৫ ✔}}

\textbf{\hl{✔ সঠিক উত্তর: ঘ) ১.৫}}

\begin{quote}
\textbf{ব্যাখ্যা:} \emph{আমরা জানি,}

\emph{√(২)≈১.৪১ এবং √(৩)≈১.৭৩}

\emph{এখন, ১.৪১ এবং ১.৭৩ এর মধ্যবর্তী একটি সসীম দশমিক সংখ্যা হলো ১.৫, যা
একটি মূলদ সংখ্যা ( (১৫/১০)=(৩/২) )।}

\emph{∴ √(২) এবং √(৩) এর মাঝে ১.৫-ই হবে একটি মূলদ সংখ্যা।}
\end{quote}

\textbf{√(২৮৯) এর বর্গমূল হলো-}

\textbf{\hl{ক) মূলদ ✔}}

খ) অমূলদ

গ) স্বাভাবিক সংখ্যা

ঘ) পূর্ণ সংখ্যা

\textbf{\hl{✔ সঠিক উত্তর: ক) মূলদ}}

\begin{quote}
\textbf{ব্যাখ্যা:} \emph{২৮৯ এর বর্গমূল হলো ১৭। যেহেতু ১৭ কে দুটি পূর্ণসংখ্যার
অনুপাত বা ভগ্নাংশ আকারে প্রকাশ করা যায়, তাই এটি একটি মূলদ সংখ্যা, একই সাথে
এটি একটি স্বাভাবিক সংখ্যা ও পূর্ণ সংখ্যা। যদি প্রশ্নে বহুনির্বাচনি (MCQ) অপশন
থাকে, তবে এর ফলাফল মূলত এই সবগুলো পরিচয়ই বহন করে। তবে সাধারণ গাণিতিক
এমসিকিউতে এটিকে প্রধানত মূলদ সংখ্যা হিসেবে চিহ্নিত করা হয়।}
\end{quote}

\textbf{২৪\^{}৫ কে ন্যূনতম কত দ্বারা গুণ করলে গুণফল একটি পূর্ণবর্গ সংখ্যা হবে?}

ক) ২

খ) ৩

\textbf{\hl{গ) ৬ ✔}}

ঘ) ২৪

\textbf{\hl{✔ সঠিক উত্তর: গ) ৬}}

\begin{quote}
\textbf{ব্যাখ্যা:} \emph{ব্যাখ্যা:
২৪\^{}৫=(২\^{}৩×৩)\^{}৫=২\^{}১৫×৩\^{}৫=২\^{}১৪×২×৩\^{}৪×৩}

\emph{এখানে, ২ এবং ৩ বিজোড় শক্তিতে (Power) জোড়াবিহীন অবস্থায় রয়েছে।}

\emph{∴ সংখ্যাটিকে পূর্ণবর্গ হতে হলে ন্যূনতম (২ × ৩) বা ৬ দ্বারা গুণ করতে হবে।}
\end{quote}

\textbf{কোনটি অমূলদ সংখ্যা}

ক) (৩.২৫/৩√(৬৪))

\textbf{\hl{খ) (√(৭)/৩√(৯)) ✔}}

গ) (৩/৪)

ঘ) (১/২)

\textbf{\hl{✔ সঠিক উত্তর: খ) (√(৭)/৩√(৯))}}

\begin{quote}
\textbf{ব্যাখ্যা:} \emph{ক) (৩.২৫/৩√(৬৪))=(৩.২৫/৪)=(৩২৫/৪০০)=(১৩/১৬) (যা
মূলদ)}

\emph{খ) (√(৭)/৩√(৯))=(৭\^{}((১/২))/৩\^{}((২/৩))) (যা ভগ্নাংশ আকারে প্রকাশ
করা যায় না, তাই অমূলদ)}

\emph{গ) এবং ঘ) অপশন দুটি স্পষ্টতই মূলদ সংখ্যা।}

\emph{∴ (√(৭)/৩√(৯)) একটি অমূলদ সংখ্যা।}
\end{quote}

\textbf{৭২×৭৫×৩\^{}৩×৪\^{}৩×২\^{}৮ কে ন্যূনতম কত দ্বারা গুণ করলে গুণফল একটি
পূর্ণবর্গ সংখ্যা হবে?}

ক) ৪

খ) ৫

\textbf{\hl{গ) ২ ✔}}

ঘ) ৩

\textbf{\hl{✔ সঠিক উত্তর: গ) ২}}

\begin{quote}
\textbf{ব্যাখ্যা:} \emph{প্রদত্ত রাশি = ৭২×৭৫×৩\^{}৩×৪\^{}৩×২\^{}৮}

\emph{= ২\^{}৩×৩\^{}২×৩×৫\^{}২×৩\^{}৩×(২\^{}২)\^{}৩×২\^{}৮ =
২\^{}৩×৩\^{}২×৩\^{}১×৫\^{}২×৩\^{}৩×২\^{}৬×২\^{}৮ = ২\^{}১৭×৩\^{}৬×৫\^{}২
= ২\^{}১৬×২×৩\^{}৬×৫\^{}২}

\emph{এখানে উৎপাদক \textquotesingle{}২\textquotesingle{} বিজোড় শক্তিতে
জোড়াবিহীন রয়েছে।}

\emph{∴ রাশিটিকে পূর্ণবর্গ করতে হলে ন্যূনতম ২ দ্বারা গুণ করতে হবে।}
\end{quote}

\textbf{যে সংখ্যাকে দুটি পূর্ণ সংখ্যার ভাগফল আকারে প্রকাশ করা যায় না তাকে কি
বলে?}

ক) মূলদ সংখ্যা

খ) স্বাভাবিক সংখ্যা

গ) জটিল সংখ্যা

\textbf{\hl{ঘ) অমূলদ সংখ্যা ✔}}

\textbf{\hl{✔ সঠিক উত্তর: ঘ) অমূলদ সংখ্যা}}

\begin{quote}
\textbf{ব্যাখ্যা:} \emph{যে সংখ্যাকে দুটি পূর্ণ সংখ্যার ভাগফল বা অনুপাতে (p/q
আকারে, যেখানে q ≠ 0) প্রকাশ করা যায় না, তাকে অমূলদ সংখ্যা বলে।}

\emph{∴ সঠিক উত্তর (ঘ)।}
\end{quote}

\textbf{√(2) অমূলদ সংখ্যাটির আসন্ন মান হবে-}

ক) 2.414

\textbf{\hl{খ) 1.414 ✔}}

গ) 1.421

ঘ) 2.412

\textbf{\hl{✔ সঠিক উত্তর: খ) 1.414}}

\begin{quote}
\textbf{ব্যাখ্যা:} \emph{√(2) এর মান 1.4142135625...}

\emph{তিন দশমিক স্থান পর্যন্ত √(2) এর আসন্ন মান হলো 1.414।}

\emph{∴ সঠিক উত্তর (খ)।}
\end{quote}

\textbf{√(৩) সংখ্যাটি কি সংখ্যা?}

ক) একটি স্বাভাবিক সংখ্যা (A natural number)

খ) একটি পূর্ণ সংখ্যা (An integer)

গ) একটি মূলদ সংখ্যা (a rational number)

\textbf{\hl{ঘ) একটি অমূলদ সংখ্যা (an irrational number) ✔}}

\textbf{\hl{✔ সঠিক উত্তর: ঘ) একটি অমূলদ সংখ্যা (an irrational number)}}

\begin{quote}
\textbf{ব্যাখ্যা:} \emph{৩ একটি মৌলিক সংখ্যা এবং পূর্ণবর্গ নয়। যেকোনো পূর্ণবর্গ
নয় এমন স্বাভাবিক সংখ্যার বর্গমূল একটি অমূলদ সংখ্যা।}

\emph{∴ √(৩) একটি অমূলদ সংখ্যা।}
\end{quote}

\textbf{যদি p একটি মৌলিক সংখ্যা হয় তবে √(p) -}

ক) একটি স্বাভাবিক সংখ্যা

খ) একটি পূর্ণ সংখ্যা

গ) একটি মূলদ সংখ্যা

\textbf{\hl{ঘ) একটি অমূলদ সংখ্যা ✔}}

\textbf{\hl{✔ সঠিক উত্তর: ঘ) একটি অমূলদ সংখ্যা}}

\begin{quote}
\textbf{ব্যাখ্যা:} \emph{যেকোনো মৌলিক সংখ্যা পূর্ণবর্গ নয়। কোনো মৌলিক সংখ্যাকে
বর্গমূল করলে তা সর্বদা একটি অমূলদ সংখ্যা হয়।}

\emph{∴ √(p) একটি অমূলদ সংখ্যা।}
\end{quote}

\textbf{৩√(২) সংখ্যাটি কোন ধরনের সংখ্যা?}

ক) মূলদ সংখ্যা

খ) জটিল সংখ্যা

\textbf{\hl{গ) অমূলদ সংখ্যা ✔}}

ঘ) বাস্তব সংখ্যা

\textbf{\hl{✔ সঠিক উত্তর: গ) অমূলদ সংখ্যা}}

\begin{quote}
\textbf{ব্যাখ্যা:} \emph{এখানে ৩ একটি মূলদ সংখ্যা এবং √(২) একটি অমূলদ সংখ্যা।
যেকোনো অ-শূন্য মূলদ সংখ্যা ও অমূলদ সংখ্যার গুণফল সর্বদা একটি অমূলদ সংখ্যা হয়।}

\emph{∴ ৩√(২) একটি অমূলদ সংখ্যা।}
\end{quote}

\textbf{নিচের কোনটি মূলদ সংখ্যা?}

ক) ৩√(২৪৩)

\textbf{\hl{খ) ৩√(৩৪৩) ✔}}

গ) ৩√(৩৯২)

ঘ) ৩√(৬৭৬)

\textbf{\hl{✔ সঠিক উত্তর: খ) ৩√(৩৪৩)}}

\begin{quote}
\textbf{ব্যাখ্যা:} \emph{৩√(৩৪৩)=৩√(৭\^{}৩)=৭}

\emph{আমরা জানি, অসীম অনাবৃত্ত দশমিক সংখ্যা ব্যতীত সকল সংখ্যাই মূলদ সংখ্যা।
এদেরকে p/q আকারে প্রকাশ করা যায় যেখানে p, q ∈ ℤ এবং q ≠ 0। বাকি তিনটি
অপশনের সংখ্যাগুলো অসীম অনাবৃত্ত দশমিক সংখ্যা।∴ সঠিক উত্তর হবে ৩√(৩৪৩) ।}
\end{quote}

\textbf{মূলদ সংখ্যাটি হলো---}

\textbf{\hl{ক) ৩√(৮) ✔}}

খ) √(১২)

গ) √(১৮)

ঘ) √(২৭)

\textbf{\hl{✔ সঠিক উত্তর: ক) ৩√(৮)}}

\begin{quote}
\textbf{ব্যাখ্যা:} \emph{৩√(৮)=৩√(২\^{}৩)=২\^{}৩\^{}((১/৩))=২ ; যা একটি
পূর্ণসংখ্যা তথা মূলদ সংখ্যা।}

\emph{∴ সঠিক উত্তর (ক)।}
\end{quote}

\textbf{মূলদ সংখ্যার সেট বোঝায় নিচের কোনটিকে?}

ক) Z

\textbf{\hl{খ) Q ✔}}

গ) P

ঘ) N

\textbf{\hl{✔ সঠিক উত্তর: খ) Q}}

\begin{quote}
\textbf{ব্যাখ্যা:} \emph{মূলদ সংখ্যার সেটকে
\textquotesingle Q\textquotesingle{} দ্বারা সূচিত করা হয়। শূন্য, স্বাভাবিক
সংখ্যা, ধনাত্মক ও ঋণাত্মক পূর্ণসংখ্যা এবং সকল ভগ্নাংশ সংখ্যাই মূলদ সংখ্যার
অন্তর্ভুক্ত।}

\emph{∴ সঠিক উত্তর (খ)।}
\end{quote}

\textbf{√(২),০.৯ এবং (১১/১০) - কে মানের ঊর্ধ্বক্রমানুসারে সাজালে নিচের কোনটি
শুদ্ধ উত্তর হবে?}

ক) √(২),(১১/১০),০.৯

খ) ০.৯,√(২),(১১/১০)

গ) (১১/১০),০.৯,√(২)

\textbf{\hl{ঘ) ০.৯,(১১/১০),√(২) ✔}}

\textbf{\hl{✔ সঠিক উত্তর: ঘ) ০.৯,(১১/১০),√(২)}}

\begin{quote}
\textbf{ব্যাখ্যা:} \emph{০.৯ = ০.৯;}

\emph{(১১/১০)=১.১ ;}

\emph{√(২)≈১.৪১}

\emph{মানের ক্রমানুসারে (ছোট থেকে বড়): 0.9 \textless{} 1.1 \textless{}
1.41}

\emph{∴ ঊর্ধ্বক্রমানুসারে সাজালে হবে: ০.৯,(১১/১০),√(২) ।}
\end{quote}

\textbf{y=(1/x) হলে x এর যে সমস্ত বাস্তব মানের জন্য y বাস্তব হবে, তারা হলো:}

ক) সমস্ত বাস্তব সংখ্যা

খ) -1 বাদে সমস্ত বাস্তব সংখ্যা

\textbf{\hl{গ) 0 বাদে সমস্ত বাস্তব সংখ্যা ✔}}

ঘ) 1 বাদে সমস্ত বাস্তব সংখ্যা

\textbf{\hl{✔ সঠিক উত্তর: গ) 0 বাদে সমস্ত বাস্তব সংখ্যা}}

\begin{quote}
\textbf{ব্যাখ্যা:} \emph{y=(1/x)}

\emph{x = 0 হলে, y = (1/0) (যা অসংজ্ঞায়িত/অনির্ণেয়)।}

\emph{∴ x = 0 ব্যতীত সকল বাস্তব সংখ্যার জন্য y এর মান বাস্তব সংখ্যা হবে।}
\end{quote}

\textbf{শূন্য সংখ্যার আদি ধারণা কাদের?}

ক) গ্রীক

খ) আরব

\textbf{\hl{গ) ভারতীয় ✔}}

ঘ) চীন

\textbf{\hl{✔ সঠিক উত্তর: গ) ভারতীয়}}

\begin{quote}
\textbf{ব্যাখ্যা:} \emph{প্রাচীন ভারতীয় গণিতবিদরা (বিশেষ করে আর্যভট্ট ও
ব্রহ্মগুপ্ত) প্রথম শূন্য (0) কে একটি স্বতন্ত্র সংখ্যা হিসেবে গাণিতিক সিস্টেমে ব্যবহার
ও এর ধারণা প্রদান করেন।}
\end{quote}

\textbf{১ থেকে ১০ পর্যন্ত সংখ্যার মধ্যে মৌলিক সংখ্যা কয়টি ?}

\textbf{\hl{ক) ৪ ✔}}

খ) ৩

গ) ৬

ঘ) ৫

\textbf{\hl{✔ সঠিক উত্তর: ক) ৪}}

\begin{quote}
\textbf{ব্যাখ্যা:} \emph{১ থেকে ১০ পর্যন্ত মেীলিক সংখ্যা আছে মোট ৪ টি। যথা:
২,৩,৫ এবং ৭ ।}
\end{quote}

\textbf{১ হতে ৩০ পর্যন্ত কয়টি মৌলিক সংখ্যা আছে ?}

ক) ৮টি

খ) ৯টি

\textbf{\hl{গ) ১০টি ✔}}

ঘ) ১১টি

\textbf{\hl{✔ সঠিক উত্তর: গ) ১০টি}}

\begin{quote}
\textbf{ব্যাখ্যা:} \emph{১ থেকে ৩০ পর্যন্ত মেীলিক সংখ্যা আছে মোট ১০ টি। যথা:
২,৩,৫,৭,১১,১৩,১৭,১৯,২৩ ও ২৯ ।}
\end{quote}

\textbf{৪৩ থেকে ৬০ পর্যন্ত মৌলিক সংখ্যা কয়টি ?}

ক) ৫

খ) ৩

গ) ৭

\textbf{\hl{ঘ) ৪ ✔}}

\textbf{\hl{✔ সঠিক উত্তর: ঘ) ৪}}

\begin{quote}
\textbf{ব্যাখ্যা:} \emph{৪৩ থেকে ৬০ পর্যন্ত মেীলিক সংখ্যা আছে মোট ৪ টি। যথা:
৪৩,৪৭,৫৩,৫৯ ।}
\end{quote}

\textbf{নিচের কোনটি মৌলিক সংখ্যা নয় ?}

ক) ২৬৩

খ) ২৩৩

\textbf{\hl{গ) ২৫৩ ✔}}

ঘ) ২৪১

\textbf{\hl{✔ সঠিক উত্তর: গ) ২৫৩}}

\begin{quote}
\textbf{ব্যাখ্যা:} \emph{প্রদত্ত সংখ্যাগুলোর মধ্যে ২৫৩ কে ভাঙ্গলে ১১×২৩ = ২৫৩ হয়
। কিন্তু অন্য সংখ্যাগুলোকে ভাঙ্গানো যায় না।। তাই ২৫৩ সংখ্যাটি মৌলিক সংখ্যা নয়।}
\end{quote}

\textbf{নিচের কোনটি মৌলিক সংখ্যা ?}

ক) ৯১

খ) ৮৭

গ) ৬৩

\textbf{\hl{ঘ) ৫৯ ✔}}

\textbf{\hl{✔ সঠিক উত্তর: ঘ) ৫৯}}

\begin{quote}
\textbf{ব্যাখ্যা:} \emph{যে সংখ্যাকে ১ ও সেই সংখ্যা ছাড়া অন্য কোনো সংখ্যা
দ্বারা ভাগ করা যায় না, তাকে মৌলিক সংখ্যা বলে। এখানে ৫৯ মৌলিক সংখ্যা ।}
\end{quote}

\textbf{নিচের কোন সংখ্যাটি মৌলিক ?}

ক) ৪৯

খ) ৫১

\textbf{\hl{গ) ৫৩ ✔}}

ঘ) ৫৫

\textbf{\hl{✔ সঠিক উত্তর: গ) ৫৩}}

\begin{quote}
\textbf{ব্যাখ্যা:} \emph{যে সংখ্যাকে ১ ও সেই সংখ্যা ছাড়া অন্য কোনো সংখ্যা
দ্বারা ভাগ করা যায় না, তাকে মৌলিক সংখ্যা বলে। এখানে ৫৩ মৌলিক সংখ্যা ।}
\end{quote}

\textbf{৩০ থেকে ৮০ এর মধ্যবর্তী বৃহত্তম ও ক্ষুদ্রতম মৌলিক সংখ্যার ব্যবধান কত ?}

ক) ৩৫

খ) ৪২

\textbf{\hl{গ) ৪৮ ✔}}

ঘ) ৫৫

\textbf{\hl{✔ সঠিক উত্তর: গ) ৪৮}}

\begin{quote}
\textbf{ব্যাখ্যা:} \emph{৩০ থেকে ৮০ এর মধ্যবর্তী বৃহত্তম ও ক্ষুদ্রতম মৌলিক সংখ্যা
= ৭৯ ও ৩১ । পার্থক্য = ৭৯-৩১ = ৪৮ ।}
\end{quote}

\textbf{নিচের কোনটি মৌলিক সংখ্যা ?}

ক) ৭২

খ) ৮৭

গ) ৬৩

\textbf{\hl{ঘ) কোনটিই নয় ✔}}

\textbf{\hl{✔ সঠিক উত্তর: ঘ) কোনটিই নয়}}

\begin{quote}
\textbf{ব্যাখ্যা:} \emph{যে সংখ্যাকে ১ ও ঐ সংখ্যা ছাড়া অন্য কোনো সংখ্যা দ্বারা
ভাগ করা যায় না, তাকে মৌলিক সংখ্যা বলে। এখানে কোনটিই মৌলিক সংখ্যা নয়।}
\end{quote}

\textbf{নিচের কোন সংখ্যাটি মৌলিক ?}

ক) ৯১

খ) ১৪৩

\textbf{\hl{গ) ৪৭ ✔}}

ঘ) ৮৭

\textbf{\hl{✔ সঠিক উত্তর: গ) ৪৭}}

\begin{quote}
\textbf{ব্যাখ্যা:} \emph{যে সংখ্যাকে ১ ও ঐ সংখ্যা ছাড়া অন্য কোনো সংখ্যা দ্বারা
ভাগ করা যায় না, তাকে মৌলিক সংখ্যা বলে। এখানে ৪৭ মৌলিক সংখ্যা ।}
\end{quote}

\textbf{নিচের কোনটি মৌলিক সংখ্যা ?}

ক) ৯

খ) ৮

গ) ৪

\textbf{\hl{ঘ) ২ ✔}}

\textbf{\hl{✔ সঠিক উত্তর: ঘ) ২}}

\begin{quote}
\textbf{ব্যাখ্যা:} \emph{যে সংখ্যাকে ১ ও ঐ সংখ্যা ছাড়া অন্য কোনো সংখ্যা দ্বারা
ভাগ করা যায় না, তাকে মৌলিক সংখ্যা বলে। এখানে ২ মৌলিক সংখ্যা ।}
\end{quote}

\textbf{জোড় মৌলিক সংখ্যা কয়টি ?}

\textbf{\hl{ক) ১টি ✔}}

খ) ২টি

গ) ৩টি

ঘ) কোনটিই নয়

\textbf{\hl{✔ সঠিক উত্তর: ক) ১টি}}

\begin{quote}
\textbf{ব্যাখ্যা:} \emph{মৌলিক সংখ্যা সবগুলোই বিজোড় শুধু ২ একটি জোড় সংখ্যা।}
\end{quote}

\textbf{২ এবং ৩২ এর মধ্যে মৌলিক সংখ্যা কয়টি ?}

ক) ৮টি

খ) ৯টি

\textbf{\hl{গ) ১০টি ✔}}

ঘ) ১১টি

\textbf{\hl{✔ সঠিক উত্তর: গ) ১০টি}}

\begin{quote}
\textbf{ব্যাখ্যা:} \emph{২ এবং ৩২ এর মধ্যে মৌলিক সংখ্যাগুলো হলো = ৩, ৫, ৭,
১১, ১৩, ১৭, ১৯, ২৩, ২৯, ৩১ = ১০ টি। তবে ২ থেকে ৩২ এর মধ্যে মৌলিক সংখ্যাগুলো
হলো = ২, ৩, ৫, ৭, ১১, ১৩, ১৭, ১৯, ২৩, ২৯, ৩১ = ১১ টি।}
\end{quote}

\textbf{৫০ এর চেয়ে ছোট কয়টি মৌলিক সংখ্যা আছে ?}

ক) ১৪টি

\textbf{\hl{খ) ১৫টি ✔}}

গ) ১৬টি

ঘ) ১৭টি

\textbf{\hl{✔ সঠিক উত্তর: খ) ১৫টি}}

\begin{quote}
\textbf{ব্যাখ্যা:} \emph{৫০ এর চেয়ে ছোট মৌলিক সংখ্যা মোট ১৫ টি। সেগুলো হলোঃ
২, ৩, ৫, ৭, ১১, ১৩, ১৭, ১৯, ২৩, ২৯, ৩১, ৩৭, ৪১, ৪৩ ও ৪৭।}
\end{quote}

\textbf{√((12/75)) কোন ধরনের সংখ্যা ?}

ক) স্বাভাবিক

\textbf{\hl{খ) মূলদ ✔}}

গ) অমূলদ

ঘ) মৌলিক

\textbf{\hl{✔ সঠিক উত্তর: খ) মূলদ}}

\begin{quote}
\textbf{ব্যাখ্যা:} \emph{সব ধরনের সাধারণ ভগ্নাংশ মুলদ সংখ্যা হয় । তবে
পৌনঃপুনিক থেকে আসা সাধারণ ভগ্নাংশসমূহ মূলদ সংখ্যা নয় ।}
\end{quote}

\textbf{নিচের কোন দুটি সহ মৌলিক সংখ্যা ?}

ক) ২৭,৫৪

খ) ৬৩,৯১

গ) ১৮৯,২১০

\textbf{\hl{ঘ) ৫২,৯৭ ✔}}

\textbf{\hl{✔ সঠিক উত্তর: ঘ) ৫২,৯৭}}

\begin{quote}
\textbf{ব্যাখ্যা:} \emph{প্রদত্ত অপশনে (ঘ) ৫২,৯৭ ক্রমজোড়টি সহমৌলিক। কারণ, ৫২
= ১×২×২×১৩ এবং ৯৭ = ১×৯৭ দেখা যাচ্ছে ৫২ ও ৯৭ এর ১ ভিন্ন অন্য কোনো সাধারণ
উৎপাদক নেই। সুতরাং এরা পরস্পর সহমৌলিক।}
\end{quote}

\textbf{m ও n বিজোড় সংখ্যা হলে নিচের কোনটি জোড় ?}

ক) mn

\textbf{\hl{খ) mn+1 ✔}}

গ) mn+2

ঘ) mn+4

\textbf{\hl{✔ সঠিক উত্তর: খ) mn+1}}

\begin{quote}
\textbf{ব্যাখ্যা:} \emph{ধরি, m = 1, n = 3}

\emph{∴ (ক) 1x3 = 3, যা বিজোড় সংখ্যা}

\emph{(খ) 1 × 3 + 1 = 4, যা জোড় সংখ্যা}

\emph{(গ) 1 × 3 + 2 = 5, যা বিজোড় সংখ্যা}

\emph{(ঘ) 1 × 3 + 4 = 7, যা বিজোড় সংখ্যা ।}

\emph{সুতরাং, mn + 1 অবশ্যই জোড় সংখ্যা হবে।}
\end{quote}

\textbf{যদি n একটি জোড় সংখ্যা হয় তবে নিচের কোনটি জোড় সংখ্যা হতে পারবে না
?}

ক) n\^{}2

খ) 3(n+1)+3

গ) 2n+2

\textbf{\hl{ঘ) 2n+3 ✔}}

\textbf{\hl{✔ সঠিক উত্তর: ঘ) 2n+3}}

\begin{quote}
\textbf{ব্যাখ্যা:} \emph{3 একটি বিজোড় সংখ্যা হওয়ায়,একই সাথে 2 জোড় সংখ্যা
হওয়ায় দুটি সংখ্যার যোগফল সবসময়ই বিজোড় হবে।}
\end{quote}

\textbf{নিচের কোনটি অবশ্যই জোড় সংখ্যা ?}

ক) x+y

খ) x+y+1

গ) xy+2

\textbf{\hl{ঘ) 2x ✔}}

\textbf{\hl{✔ সঠিক উত্তর: ঘ) 2x}}

\begin{quote}
\textbf{ব্যাখ্যা:} \emph{2 × বিজোড় = জোড় সংখ্যা হয়। যেমন, ২ × ৩ = ৬}
\end{quote}

\textbf{যদি n একটি জোড় সংখ্যা হয় তবে নিচের কোনটি জোড় সংখ্যা হতে পারবে না
?}

ক) n\^{}2

খ) 5(n+2)

গ) 2n+2

\textbf{\hl{ঘ) 7n+3 ✔}}

\textbf{\hl{✔ সঠিক উত্তর: ঘ) 7n+3}}

\begin{quote}
\textbf{ব্যাখ্যা:} \emph{ধরি n = 2 । তাহলে (7×2)+3 = 14 + 3 = 17 যা একটি
বিজোড় সংখ্যা ।}
\end{quote}

\textbf{০,১,২,৩ দ্বারা গঠিত বৃহত্তম ও ক্ষুদ্রতম সংখ্যার মধ্যে পার্থক্য কত হবে ?}

ক) ৩০৮৭

খ) ২০৮৭

গ) ৩১৮৭

\textbf{\hl{ঘ) ২১৮৭ ✔}}

\textbf{\hl{✔ সঠিক উত্তর: ঘ) ২১৮৭}}

\begin{quote}
\textbf{ব্যাখ্যা:} \emph{০ দ্বারা কোন সংখ্যা শুরু হয় না।তাই ক্ষুদ্রতম সংখ্যা তৈরির
ক্ষেত্রে প্রথমে ১ বসিয়ে তারপর ০ বসাতে হবে; তাই ক্ষুদ্রতম সংখ্যাটি হবে ১০২৩ এবং
বৃহত্তম সংখ্যাটি হবে ৩২১০ । সুতরাং পার্থক্য ৩২১০-১০২৩ = ২১৮৭ ।}
\end{quote}

\textbf{০,১,২,৩ ও ৪ দ্বারা গঠিত পাঁচ অঙ্কের বৃহত্তম ও ক্ষুদ্রতম সংখ্যার যোগফল
কোনটি ?}

ক) ৪৪৪৪৪

খ) ৫৩৪৪২

\textbf{\hl{গ) ৫৩৪৪৪ ✔}}

ঘ) ৫৩৪৪৮

\textbf{\hl{✔ সঠিক উত্তর: গ) ৫৩৪৪৪}}

\begin{quote}
\textbf{ব্যাখ্যা:} \emph{০,১,২,৩ ও ৪ দ্বারা গঠিত পাঁচ অঙ্কের বৃহত্তম সংখ্যা
৪৩২১০ । ০,১,২,৩ ও ৪ দ্বারা গঠিত পাঁচ অঙ্কের ক্ষুদ্রতম সংখ্যা ১০২৩৪ সুতরাং যোগফল
= ( ৪৩২১০+১০২৩৪ ) = ৫৩৪৪৪ ।}
\end{quote}

\textbf{১-১০০ পর্যন্ত মৌলিক সংখ্যা কয়টি ?}

ক) ২৩টি

\textbf{\hl{খ) ২৫টি ✔}}

গ) ২৭টি

ঘ) ২৯টি

\textbf{\hl{✔ সঠিক উত্তর: খ) ২৫টি}}

\begin{quote}
\textbf{ব্যাখ্যা:} \emph{১ থেকে ১০০ পর্যন্ত মৌলিক সংখ্যাগুলো হলো: ২, ৩, ৫, ৭,
১১, ১৩, ১৭, ১৯, ২৩, ২৯, ৩১, ৩৭, ৪১, ৪৩, ৪৭, ৫৩, ৫৯, ৬১, ৬৭, ৭১, ৭৩, ৭৯,
৮৩, ৮৯ ও ৯৭। মোট ২৫ টি।}
\end{quote}

\textbf{১ থেকে ১০ পর্যন্ত জোড় মৌলিক সংখ্যা কয়টি ?}

ক) ২

\textbf{\hl{খ) ১ ✔}}

গ) ৩

ঘ) ৪

\textbf{\hl{✔ সঠিক উত্তর: খ) ১}}

\begin{quote}
\textbf{ব্যাখ্যা:} \emph{১ থেকে ১০ পর্যন্ত মৌলিক সংখ্যা চারটি। ২, ৩, ৫, এবং
৭। শুধু ২ একটি জোড় সংখ্যা।}
\end{quote}

\textbf{কোনটি স্বাভাবিক সংখ্যা ?}

ক) -1

খ) √(2)

গ) (5/2)

\textbf{\hl{ঘ) 3 ✔}}

\textbf{\hl{✔ সঠিক উত্তর: ঘ) 3}}

\begin{quote}
\textbf{ব্যাখ্যা:} \emph{১ , ২ , ৩ ..... এগুলো স্বাভাবিক সংখ্যা}
\end{quote}

\textbf{ক্ষুদ্রতম মৌলিক সংখ্যা নিচের কোনটি ?}

ক) ১

\textbf{\hl{খ) ২ ✔}}

গ) ৩

ঘ) ৫

\textbf{\hl{✔ সঠিক উত্তর: খ) ২}}

\begin{quote}
\textbf{ব্যাখ্যা:} \emph{২ সবচেয়ে ছোট মৌলিক সংখ্যা}
\end{quote}

\textbf{সকল মূলদ ও অমূলদ সংখ্যাকে কি বলে ?}

ক) স্বাভাবিক সংখ্যা

খ) মৌলিক সংখ্যা

গ) পূর্ণসংখ্যা

\textbf{\hl{ঘ) বাস্তব সংখ্যা ✔}}

\textbf{\hl{✔ সঠিক উত্তর: ঘ) বাস্তব সংখ্যা}}

\begin{quote}
\textbf{ব্যাখ্যা:} \emph{সকল মূলদ ও অমূলদ সংখ্যাকে একত্রে বাস্তব সংখ্যা বলে}
\end{quote}

\textbf{√3 এবং 5 এর মাঝে কয়টি পূর্ণসংখ্যা আছে ?}

ক) 1টি

খ) 2টি

\textbf{\hl{গ) 3টি ✔}}

ঘ) 4টি

\textbf{\hl{✔ সঠিক উত্তর: গ) 3টি}}

\begin{quote}
\textbf{ব্যাখ্যা:} \emph{√3 = 1.73205080757 এবং 5 এর মাঝে 2 , 3 , 4 এই
তিনটি পূর্ণসংখ্যা আছে ।}
\end{quote}

\textbf{(2/9) কোন ধরনের সংখ্যা ?}

\textbf{\hl{ক) মূলদ ✔}}

খ) অমূলদ

গ) স্বাভাবিক

ঘ) অনাবৃত দশমিক

\textbf{\hl{✔ সঠিক উত্তর: ক) মূলদ}}

\begin{quote}
\textbf{ব্যাখ্যা:} \emph{সব ধরনের সাধারণ ভগ্নাংশ মুলদ সংখ্যা হয় । তবে
পৌনঃপুনিক থেকে আসা সাধারণ ভগ্নাংশসমূহ মূলদ সংখ্যা নয় ।}
\end{quote}

\textbf{০, ১, ৪, ৬, ৯ এই সংখ্যাগুলো এক বা একাধিকবার ব্যবহার করে পাঁচ অঙ্কের
ক্ষুদ্রতম সংখ্যা কত?}

ক) ০১৪৬৯

\textbf{\hl{খ) ১০৪৬৯ ✔}}

গ) ৯০১৪৬

ঘ) ১০০০০

ঙ) কোনোটিই নয়

\textbf{\hl{✔ সঠিক উত্তর: খ) ১০৪৬৯}}

\begin{quote}
\textbf{ব্যাখ্যা:} \emph{০ \textless{} ১; কিন্তু শূন্য (০) সংখ্যার বামে বসলে তা
মান প্রকাশ করে না (অর্থাৎ সংখ্যাটি ৪ অঙ্ক বিশিষ্ট হয়ে যায়)। ∴ ক্ষুদ্রতম সংখ্যা
গঠনের জন্য শূন্যকে ১ এর পরে বসাতে হবে। অর্থাৎ সংখ্যাটি ১০৪৬৯। আর একাধিকবার
অঙ্কগুলো ব্যবহার করলে সংখ্যা বড় হবে, ছোট নয়।}
\end{quote}

\textbf{০, ১ দ্বারা ৪ অঙ্কের কতগুলো সংকেত তৈরি করা যায়?}

ক) ৮

\textbf{\hl{খ) ১৬ ✔}}

গ) ৩২

ঘ) ৬৪

\textbf{\hl{✔ সঠিক উত্তর: খ) ১৬}}

\begin{quote}
\textbf{ব্যাখ্যা:} \emph{বাইনারি সংখ্যা পদ্ধতির মৌলিক চিহ্ন দুটি (০ ও ১)। অঙ্ক
বা বিটের সংখ্যা n হলে গঠনযোগ্য মোট সংকেত সংখ্যা হয় 2\^{}n টি।}

\emph{এখানে ৪ অঙ্কের সংকেত তৈরির ক্ষেত্রে মোট সংকেত সংখ্যা হবে = 2\^{}4 = ১৬
টি। ∴ ৪ অঙ্কের ১৬টি সংকেত গঠন করা যায়।}
\end{quote}

\textbf{বাইনারি সংখ্যা পদ্ধতিতে কতটি অঙ্ক ব্যবহৃত হয়?}

ক) একটি

\textbf{\hl{খ) দুইটি ✔}}

গ) তিনটি

ঘ) চারটি

\textbf{\hl{✔ সঠিক উত্তর: খ) দুইটি}}

\begin{quote}
\textbf{ব্যাখ্যা:} \emph{বাইনারি (Binary) একটি দ্বিমিক সংখ্যা পদ্ধতি, যার বেস
বা ভিত্তি হলো ২। এই পদ্ধতিতে শুধুমাত্র দুটি অঙ্ক বা প্রতীক ব্যবহার করা হয়, যা হলো
\textquotesingle{}০\textquotesingle{} (০) এবং
\textquotesingle{}১\textquotesingle{} (1)।}
\end{quote}

\textbf{পাঁচ অঙ্কের ক্ষুদ্রতম সংখ্যা কোনটি?}

ক) ৯৯৯৯৯

খ) ৯৮৭৬৫

গ) ১২৩৪৫

\textbf{\hl{ঘ) ১০০০০ ✔}}

\textbf{\hl{✔ সঠিক উত্তর: ঘ) ১০০০০}}

\begin{quote}
\textbf{ব্যাখ্যা:} \emph{পাঁচ অঙ্কের ক্ষুদ্রতম সংখ্যা হলো ১ এর পরে ৪টি শূন্য,
অর্থাৎ ১০০০০।}
\end{quote}

\textbf{তিন অঙ্কের বৃহত্তম সংখ্যা ও ক্ষুদ্রতম সংখ্যার পার্থক্য কত?}

ক) ৯৯৮

খ) ৯৮৮

\textbf{\hl{গ) ৮৯৯ ✔}}

ঘ) ৮৮৮

\textbf{\hl{✔ সঠিক উত্তর: গ) ৮৯৯}}

\begin{quote}
\textbf{ব্যাখ্যা:} \emph{তিন অঙ্কের বৃহত্তম সংখ্যা = ৯৯৯}

\emph{তিন অঙ্কের ক্ষুদ্রতম সংখ্যা = ১০০}

\emph{∴ পার্থক্য = ৯৯৯ − ১০০ = ৮৯৯।}
\end{quote}

\textbf{পাঁচ অঙ্কের বৃহত্তম ও ক্ষুদ্রতম সংখ্যা দুটির যোগফল কত?}

ক) ১০০০০৯

\textbf{\hl{খ) ১০৯৯৯৯ ✔}}

গ) ৯৯৯৯৯

ঘ) ১০০০৯৯

\textbf{\hl{✔ সঠিক উত্তর: খ) ১০৯৯৯৯}}

\begin{quote}
\textbf{ব্যাখ্যা:} \emph{পাঁচ অঙ্কের বৃহত্তম সংখ্যা = ৯৯৯৯৯}

\emph{পাঁচ অঙ্কের ক্ষুদ্রতম সংখ্যা = ১০০০০}

\emph{∴ যোগফল = ৯৯৯৯৯ + ১০০০০ = ১০৯৯৯৯।}
\end{quote}

\textbf{পাঁচ অঙ্কের ক্ষুদ্রতম সংখ্যাটি হতে চার অঙ্কের বৃহত্তম সংখ্যার বিয়োগফল কত?}

ক) ৯০

খ) ১০

\textbf{\hl{গ) ১ ✔}}

ঘ) (১/২)

\textbf{\hl{✔ সঠিক উত্তর: গ) ১}}

\begin{quote}
\textbf{ব্যাখ্যা:} \emph{পাঁচ অঙ্কের ক্ষুদ্রতম সংখ্যা = ১০০০০}

\emph{চার অঙ্কের বৃহত্তম সংখ্যা = ৯৯৯৯}

\emph{∴ বিয়োগফল = ১০০০০ − ৯৯৯৯ = ১।}
\end{quote}

\textbf{পাঁচ অঙ্কের বৃহত্তম সংখ্যা হতে পাঁচ অঙ্কের ক্ষুদ্রতম সংখ্যা বিয়োগ করলে কত
হবে?}

ক) ৯৯৯৯৯

\textbf{\hl{খ) ৮৯৯৯৯ ✔}}

গ) ৯৯৯৯৮

ঘ) ১০০০০

\textbf{\hl{✔ সঠিক উত্তর: খ) ৮৯৯৯৯}}

\begin{quote}
\textbf{ব্যাখ্যা:} \emph{পাঁচ অঙ্কের বৃহত্তম সংখ্যা = ৯৯৯৯৯}

\emph{পাঁচ অঙ্কের ক্ষুদ্রতম সংখ্যা = ১০০০০}

\emph{∴ বিয়োগফল = ৯৯৯৯৯ − ১০০০০ = ৮৯৯৯৯।}
\end{quote}

\textbf{০, ১, ২, এবং ৩ দ্বারা গঠিত চার অঙ্কের বৃহত্তম এবং ক্ষুদ্রতম সংখ্যার
বিয়োগ ফল-}

ক) ৩১৪৭

খ) ২২৮৭

গ) ২৯৮৭

\textbf{\hl{ঘ) ২১৮৭ ✔}}

\textbf{\hl{✔ সঠিক উত্তর: ঘ) ২১৮৭}}

\begin{quote}
\textbf{ব্যাখ্যা:} \emph{০, ১, ২ এবং ৩ দ্বারা গঠিত বৃহত্তম সংখ্যা = ৩২১০}

\emph{০, ১, ২ এবং ৩ দ্বারা গঠিত ক্ষুদ্রতম সংখ্যা = ১০২৩ (কারণ বামে ০ বসালে তা
৩ অঙ্কের সংখ্যা হয়ে যায়)}

\emph{∴ বিয়োগফল = ৩২১০ − ১০২৩ = ২১৮৭।}
\end{quote}

\textbf{পাঁচ অঙ্কের ক্ষুদ্রতম সংখ্যা ও চার অঙ্কের বৃহত্তম সংখ্যার অন্তর কত?}

ক) ৯

খ) ১০

\textbf{\hl{গ) ১ ✔}}

ঘ) ৫

\textbf{\hl{✔ সঠিক উত্তর: গ) ১}}

\begin{quote}
\textbf{ব্যাখ্যা:} \emph{পাঁচ অঙ্কের ক্ষুদ্রতম সংখ্যা = ১০০০০}

\emph{চার অঙ্কের বৃহত্তম সংখ্যা = ৯৯৯৯}

\emph{∴ অন্তর (পার্থক্য) = ১০০০০ − ৯৯৯৯ = ১।}
\end{quote}

\textbf{একটি সংখ্যার শতক, দশক ও একক স্থানীয় অঙ্ক যথাক্রমে p, q, r হলে
সংখ্যাটি হবে-}

ক) 100r + 10p + q

\textbf{\hl{খ) 100p + 10q + r ✔}}

গ) 100q + 10r + p

ঘ) 100pq + r

\textbf{\hl{✔ সঠিক উত্তর: খ) 100p + 10q + r}}

\begin{quote}
\textbf{ব্যাখ্যা:} \emph{স্থানীয় মান অনুযায়ী সংখ্যা গঠন করার নিয়ম হলো: (শতক
স্থানীয় অঙ্ক × ১০০) + (দশক স্থানীয় অঙ্ক × ১০) + (একক স্থানীয় অঙ্ক × ১)।}

\emph{∴ নির্ণেয় সংখ্যাটি = (p×100)+(q×10)+(r×1) = 100p + 10q + r।}
\end{quote}

\textbf{১, ২ ও ৩ দ্বারা গঠিত ৩ অঙ্কের যতটি সংখ্যা লেখা যায়, তাদের সমষ্টি
কত?}

ক) ১২২৩

খ) ১২৩৩

গ) ১৩২২

\textbf{\hl{ঘ) ১৩৩২ ✔}}

\textbf{\hl{✔ সঠিক উত্তর: ঘ) ১৩৩২}}

\begin{quote}
\textbf{ব্যাখ্যা:} \emph{১, ২ ও ৩ দ্বারা গঠিত ৩ অঙ্কের সংখ্যাগুলো হলো: ১২৩,
১৩২, ২১৩, ২৩১, ৩১২ এবং ৩২১।}

\emph{∴ তাদের সমষ্টি = ১২৩ + ১৩২ + ২১৩ + ২৩১ + ৩১২ + ৩২১ = ১৩৩২।}
\end{quote}

\textbf{যে রাশিকে ভাগ করা হয় তাকে কী বলে?}

\textbf{\hl{ক) ভাজ্য ✔}}

খ) ভাজক

গ) ভাগফল

ঘ) ভাগশেষ

\textbf{\hl{✔ সঠিক উত্তর: ক) ভাজ্য}}

\begin{quote}
\textbf{ব্যাখ্যা:} \emph{যে রাশিকে ভাগ করা হয় তাকে
\textquotesingle ভাজ্য\textquotesingle{} বলে। আর যা দিয়ে ভাগ করা হয় তা
হলো \textquotesingle ভাজক\textquotesingle, ভাগ করার পর যে ফল পাওয়া যায়
তা \textquotesingle ভাগফল\textquotesingle{} এবং যা অবশিষ্ট থাকে তা
\textquotesingle ভাগশেষ\textquotesingle।}
\end{quote}

\textbf{Which of the following integers has the most divisors?}

\textbf{\hl{ক) 88 ✔}}

খ) 91

গ) 95

ঘ) 99

\textbf{\hl{✔ সঠিক উত্তর: ক) 88}}

\begin{quote}
\textbf{ব্যাখ্যা:} \emph{সংখ্যাগুলোর মৌলিক উৎপাদক বিশ্লেষণ করে ভাজক সংখ্যা
নির্ণয় করা যায়:}

\emph{• 88 = 2\^{}3 × 11\^{}1; ভাজক সংখ্যা = (3 + 1)(1 + 1) = 8 টি।}

\emph{• 91 = 7\^{}1 × 13\^{}1; ভাজক সংখ্যা = (1 + 1)(1 + 1) = 4 টি।}

\emph{• 95 = 5\^{}1 × 19\^{}1; ভাজক সংখ্যা = (1 + 1)(1 + 1) = 4 টি।}

\emph{• 99 = 3\^{}2 × 11\^{}1; ভাজক সংখ্যা = (2 + 1)(1 + 1) = 6 টি।}

\emph{∴ 88 সংখ্যাটির ভাজক সংখ্যা সর্বাধিক।}
\end{quote}

\textbf{৩৬ সংখ্যাটির মোট কতগুলো ভাজক রয়েছে?}

ক) ৬টি

খ) ৮টি

\textbf{\hl{গ) ৯টি ✔}}

ঘ) ১০টি

\textbf{\hl{✔ সঠিক উত্তর: গ) ৯টি}}

\begin{quote}
\textbf{ব্যাখ্যা:} \emph{৩৬ এর গুণনীয়ক বা ভাজকগুলো হলো: ১, ২, ৩, ৪, ৬, ৯,
১২, ১৮, ৩৬।}

\emph{বিকল্প পদ্ধতিতে: 36 = 2\^{}2 × 3\^{}2}

\emph{∴ মোট ভাজক সংখ্যা = (2 + 1) × (2 + 1) = 3 × 3 = ৯ টি।}
\end{quote}

\textbf{৭২ সংখ্যাটির মোট ভাজক আছে কতটি?}

ক) ৯টি

খ) ১০টি

গ) ১১টি

\textbf{\hl{ঘ) ১২টি ✔}}

\textbf{\hl{✔ সঠিক উত্তর: ঘ) ১২টি}}

\begin{quote}
\textbf{ব্যাখ্যা:} \emph{৭২ সংখ্যাটির মৌলিক উৎপাদক বিশ্লেষণ:
72=2\^{}3×3\^{}2}

\emph{∴ ৭২ এর মোট ভাজক সংখ্যা = (3+1)×(2+1)=4×3=12 টি।}

\emph{(ভাজকগুলো হলো: ১, ২, ৩, ৪, ৬, ৮, ৯, ১২, ১৮, ২৪, ৩৬ ও ৭২)}
\end{quote}

\textbf{৫৪০ সংখ্যাটির কতগুলো ভাজক আছে?}

ক) ১৮

খ) ২০

গ) ২২

\textbf{\hl{ঘ) ২৪ ✔}}

\textbf{\hl{✔ সঠিক উত্তর: ঘ) ২৪}}

\begin{quote}
\textbf{ব্যাখ্যা:} \emph{৫৪০ সংখ্যাটির মৌলিক উৎপাদক বিশ্লেষণ:
540=2×2×3×3×3×5=2\^{}2×3\^{}3×5\^{}1}

\emph{∴ নির্ণেয় ভাজক সংখ্যা = (2+1)×(3+1)×(1+1)=3×4×2=24 টি।}
\end{quote}

\textbf{১০০৮ এর কতটি ভাজক আছে?}

ক) ২০

খ) ২৪

গ) ২৮

\textbf{\hl{ঘ) ৩০ ✔}}

\textbf{\hl{✔ সঠিক উত্তর: ঘ) ৩০}}

\begin{quote}
\textbf{ব্যাখ্যা:} \emph{১০০৮ সংখ্যাটির মৌলিক উৎপাদক বিশ্লেষণ:
1008=2\^{}4×3\^{}2×7\^{}1}

\emph{∴ ১০০৮ এর ভাজক সংখ্যা = (4+1)×(2+1)×(1+1)=5×3×2=30 টি।}
\end{quote}

\textbf{নিঃশেষে বিভাজ্য না হলে কোনটি নির্ভুল?}

\textbf{\hl{ক) ভাজ্য = (ভাজক × ভাগফল) + ভাগশেষ ✔}}

খ) ভাজ্য = (ভাজক + ভাগশেষ) × ভাগফল

গ) ভাজ্য = (ভাগশেষ × ভাগফল) + ভাজক

ঘ) ভাজ্য = (ভাজক + ভাগফল) × ভাগশেষ

\textbf{\hl{✔ সঠিক উত্তর: ক) ভাজ্য = (ভাজক × ভাগফল) + ভাগশেষ}}

\begin{quote}
\textbf{ব্যাখ্যা:} \emph{আমরা জানি, নিঃশেষে বিভাজ্য না হলে,}

\emph{ভাজ্য = (ভাজক × ভাগফল) + ভাগশেষ।}
\end{quote}

\textbf{নিচের কোন সম্পর্কটি সঠিক?}

ক) ভাজক = (ভাজ্য − ভাগশেষ) × ভাগফল

খ) ভাজ্য = (ভাজক − ভাগশেষ) × ভাগফল

গ) ভাজ্য = ভাজক × ভাগফল − ভাগশেষ

\textbf{\hl{ঘ) ভাজক = (ভাজ্য − ভাগশেষ) ÷ ভাগফল ✔}}

\textbf{\hl{✔ সঠিক উত্তর: ঘ) ভাজক = (ভাজ্য − ভাগশেষ) ÷ ভাগফল}}

\begin{quote}
\textbf{ব্যাখ্যা:} \emph{আমরা জানি,}

\emph{ভাজ্য = (ভাজক × ভাগফল) + ভাগশেষ}

\emph{⇒ ভাজ্য − ভাগশেষ = ভাজক × ভাগফল}

\emph{∴ ভাজক = (ভাজ্য − ভাগশেষ) ÷ ভাগফল।}
\end{quote}

\textbf{ভাজক ১০, ভাগফল ১০ ও ভাগশেষ ১ হলে ভাজ্য কত?}

ক) ১১

\textbf{\hl{খ) ১০১ ✔}}

গ) ১০২

ঘ) ১০০১

\textbf{\hl{✔ সঠিক উত্তর: খ) ১০১}}

\begin{quote}
\textbf{ব্যাখ্যা:} \emph{আমরা জানি,}

\emph{ভাজ্য = (ভাজক × ভাগফল) + ভাগশেষ = (১০ × ১০) + ১}

\emph{= ১০০ + ১ = ১০১।}
\end{quote}

\textbf{ভাজক ভাগফলের ১০ গুণ, ভাজক ০.৫ হলে ভাজ্য কত?}

\textbf{\hl{ক) ০.০২৫ ✔}}

খ) ০.২৫

গ) ২৫

ঘ) ২.৫

\textbf{\hl{✔ সঠিক উত্তর: ক) ০.০২৫}}

\begin{quote}
\textbf{ব্যাখ্যা:} \emph{দেয়া আছে, ভাজক = ০.৫}

\emph{প্রশ্নমতে,}

\emph{ভাজক = ভাগফল × ১০}

\emph{⇒ ০.৫ = ভাগফল × ১০}

\emph{⇒ ভাগফল = ০.৫ ÷ ১০ = ০.০৫}

\emph{আমরা জানি,}

\emph{ভাজ্য = ভাজক × ভাগফল}

\emph{= ০.৫ × ০.০৫ = ০.০২৫।}
\end{quote}

\textbf{১২ ও ৯৬ এর মধ্যে (এই দুটি সংখ্যাসহ) কয়টি সংখ্যা ৪ দ্বারা বিভাজ্য?}

ক) ২১

\textbf{\hl{খ) ২২ ✔}}

গ) ২৩

ঘ) ২৪

\textbf{\hl{✔ সঠিক উত্তর: খ) ২২}}

\begin{quote}
\textbf{ব্যাখ্যা:} \emph{১ থেকে ৯৬ পর্যন্ত ৪ দ্বারা বিভাজ্য সংখ্যা আছে = ৯৬ ÷ ৪
= ২৪টি।}

\emph{কিন্তু প্রশ্নানুসারে ১২ থেকে শুরু করতে বলা হয়েছে, তাই ১২ এর পূর্ববর্তী ৪
দ্বারা বিভাজ্য সংখ্যাগুলো (৪ এবং ৮; মোট ২টি) বাদ যাবে।}

\emph{∴ ১২ ও ৯৬ এর মধ্যে (উভয় সংখ্যাসহ) ৪ দ্বারা বিভাজ্য সংখ্যা আছে = ২৪ − ২
= ২২টি।}
\end{quote}

\textbf{নিচের কোনটি ৪ দ্বারা বিভাজ্য?}

ক) ২১৪১৩৩

\textbf{\hl{খ) ৫১০০৫৬ ✔}}

গ) ৩২২৫৬৯

ঘ) ৯৫২২১৭

\textbf{\hl{✔ সঠিক উত্তর: খ) ৫১০০৫৬}}

\begin{quote}
\textbf{ব্যাখ্যা:} \emph{কোনো সংখ্যার একক ও দশক স্থানের অংক দুটি দ্বারা গঠিত
সংখ্যা ৪ দ্বারা বিভাজ্য হলে, সম্পূর্ণ সংখ্যাটি ৪ দ্বারা বিভাজ্য হবে।}

\emph{এখানে, ৫১০০৫৬ এর শেষ দুটি অংক দ্বারা গঠিত সংখ্যা ৫৬, যা ৪ দ্বারা
বিভাজ্য ( 56÷4=14 )।}

\emph{∴ ৫১০০৫৬ সংখ্যাটি ৪ দ্বারা বিভাজ্য।}
\end{quote}

\textbf{নিম্নলিখিত সংখ্যাগুলোর মধ্যে কোনটির ভাজক সংখ্যা বিজোড়?}

ক) ২০৪৮

\textbf{\hl{খ) ১০২৪ ✔}}

গ) ৫১২

ঘ) ৪৮

\textbf{\hl{✔ সঠিক উত্তর: খ) ১০২৪}}

\begin{quote}
\textbf{ব্যাখ্যা:} \emph{আমরা জানি, কেবল পূর্ণ বর্গসংখ্যার ভাজক সংখ্যা বিজোড়
হয়।}

\emph{উপরিউক্ত সংখ্যাগুলোর মধ্যে শুধুমাত্র ১০২৪ একটি পূর্ণবর্গ সংখ্যা (কারণ,
√(1024)=32 )।}

\emph{∴ ১০২৪ সংখ্যাটির মোট ভাজক সংখ্যা ১১টি (যা বিজোড়)।}
\end{quote}

\textbf{২১৫, ২১৯, ৩২৫, ৬২৫ সংখ্যাগুলোর মধ্যে পূর্ণবর্গ কোনটি?}

ক) ২১৫

খ) ২১৯

\textbf{\hl{গ) ৬২৫ ✔}}

ঘ) ৩২৫

\textbf{\hl{✔ সঠিক উত্তর: গ) ৬২৫}}

\begin{quote}
\textbf{ব্যাখ্যা:} \emph{অপশনগুলোর মৌলিক উৎপাদক বিশ্লেষণ করে পাই:}

\emph{২১৫ = ৫ × ৪৩ (পূর্ণবর্গ নয়)}

\emph{২১৯ = ৩ × ৭৩ (পূর্ণবর্গ নয়)}

\emph{৩২৫ = ৫ × ৫ × ১৩ (পূর্ণবর্গ নয়)}

\emph{৬২৫ = ২৫ × ২৫ = 25\^{}2 (পূর্ণবর্গ)}

\emph{∴ নির্ণেয় পূর্ণবর্গ সংখ্যাটি ৬২৫।}
\end{quote}

\textbf{If the integer x is divisible by 5 but not by 2 then which one
of the following is never an integer?/x একটি পূর্ণ সংখ্যা যা ৫ দ্বারা বিভাজ্য
কিন্তু ২ দ্বারা বিভাজ্য নয়, তাহলে নিচের কোনটি কখনও একটি পূর্ণ সংখ্যা হতে পারে
না?}

\textbf{\hl{ক) (x/24) ✔}}

খ) (x+1/2)

গ) (x\^{}3/3)

ঘ) (x/7)

\textbf{\hl{✔ সঠিক উত্তর: ক) (x/24)}}

\begin{quote}
\textbf{ব্যাখ্যা:} \emph{যেহেতু x সংখ্যাটি ২ দ্বারা বিভাজ্য নয়, তাই x একটি
বিজোড় সংখ্যা।}

\emph{যেকোনো বিজোড় সংখ্যাকে কোনো জোড় সংখ্যা (যেমন: ২৪) দ্বারা ভাগ করলে তা
কখনোই পূর্ণসংখ্যা হতে পারে না (কারণ, 24=2×2×2×3 )।}

\emph{∴ (x/24) কখনোই একটি পূর্ণ সংখ্যা হতে পারে না।}
\end{quote}

\textbf{১, ২, ৩, ৪ অংকগুলি দ্বারা ৩০০০ অপেক্ষা বৃহত্তর কয়টি সংখ্যা গঠন করা
থাকে?}

ক) ৮টি

\textbf{\hl{খ) ১২টি ✔}}

গ) ১৬টি

ঘ) ১৮টি

\textbf{\hl{✔ সঠিক উত্তর: খ) ১২টি}}

\begin{quote}
\textbf{ব্যাখ্যা:} \emph{সংখ্যাটি ৩০০০ অপেক্ষা বড় হতে হলে অবশ্যই হাজার স্থানীয়
(প্রথম) অংকটি ৩ অথবা ৪ হতে হবে।}

\emph{প্রথম অংক ৩ স্থির রাখলে বাকি ৩টি স্থানে বাকি ৩টি অংক 3!=6 উপায়ে বসানো
যায়।}

\emph{অনুরূপভাবে, প্রথম অংক ৪ স্থির রাখলে বাকি ৩টি স্থান 3!=6 উপায়ে পূরণ করা
যায়।}

\emph{∴ ৩০০০ অপেক্ষা বড় সংখ্যা গঠন করা যায় = 6+6=12 টি।}
\end{quote}

\textbf{2, 3 এবং 4 দ্বারা তিন অংকের কতটি বিজোড় সংখ্যা গঠন করা যায়?}

\textbf{\hl{ক) 2টি ✔}}

খ) 5টি

গ) 6টি

ঘ) 7টি

\textbf{\hl{✔ সঠিক উত্তর: ক) 2টি}}

\begin{quote}
\textbf{ব্যাখ্যা:} \emph{বিজোড় সংখ্যা গঠনের জন্য একক স্থানীয় অংকটি অবশ্যই
বিজোড় হতে হবে।}

\emph{এখানে দেওয়া ৩টি অংকের মধ্যে বিজোড় অংক কেবল ১টি, যা হলো ৩।}

\emph{একক স্থানে ৩ বসালে শতক ও দশক স্থানে বাকি ২টি অংক (২ এবং ৪) বসানো যায়
2!=2 উপায়ে।}

\emph{সংখ্যাতয় হলো: 243 এবং 423।}

\emph{∴ মোট ২টি বিজোড় সংখ্যা গঠন করা যায়।}
\end{quote}

\textbf{যদি n একটি জোড় সংখ্যা হয় তবে নিচের কোনটি জোড় সংখ্যা হতে পারে না?}

ক) n\^{}2

খ) 5(n+2)

গ) 2n+2

\textbf{\hl{ঘ) 7n+3 ✔}}

\textbf{\hl{✔ সঠিক উত্তর: ঘ) 7n+3}}

\begin{quote}
\textbf{ব্যাখ্যা:} \emph{ধরি, জোড় সংখ্যা n=2 ।}

\emph{তাহলে,}

\emph{(বিজোড় সংখ্যা × জোড় সংখ্যা) + বিজোড় সংখ্যা = জোড় সংখ্যা + বিজোড় সংখ্যা
= বিজোড় সংখ্যা।}

\emph{এখানে, 7n+3=7×2+3=14+3=17 (যা একটি বিজোড় সংখ্যা)।}

\emph{∴ 7n+3 কখনো জোড় সংখ্যা হতে পারে না।}
\end{quote}

\textbf{প্রথম দশটি বিজোড় সংখ্যার যোগফল কত?}

\textbf{\hl{ক) ১০০ ✔}}

খ) ৮১

গ) ১০০০

ঘ) ১০৯

\textbf{\hl{✔ সঠিক উত্তর: ক) ১০০}}

\begin{quote}
\textbf{ব্যাখ্যা:} \emph{আমরা জানি, প্রথম n সংখ্যক বিজোড় সংখ্যার যোগফল =
n\^{}2}

\emph{∴ প্রথম ১০টি বিজোড় সংখ্যার যোগফল = 10\^{}2=100 ।}
\end{quote}

\textbf{If x is an negative odd integer \& y is an positive even
integer, then xy must be:/ যদি x একটি ঋণাত্মক বিজোড় পূর্ণসংখ্যা এবং y একটি
ধনাত্মক জোড় পূর্ণসংখ্যা হয়, তবে xy অবশ্যই-}

\textbf{\hl{ক) Even \& -ve (জোড় এবং ঋণাত্মক) ✔}}

খ) Odd \& -ve (বিজোড় এবং ঋণাত্মক)

গ) Even \& +ve (জোড় এবং ধনাত্মক)

ঘ) Odd \& +ve (বিজোড় এবং ধনাত্মক)

ঙ) Can\textquotesingle t be determined (নির্ণয় করা সম্ভব নয়)

\textbf{\hl{✔ সঠিক উত্তর: ক) Even \& -ve (জোড় এবং ঋণাত্মক)}}

\begin{quote}
\textbf{ব্যাখ্যা:} \emph{আমরা জানি,}

\emph{বিজোড় সংখ্যা × জোড় সংখ্যা = জোড় সংখ্যা}

\emph{এবং, ঋণাত্মক সংখ্যা × ধনাত্মক সংখ্যা = ঋণাত্মক সংখ্যা}

\emph{∴ xy অবশ্যই একটি Even \& -ve (জোড় এবং ঋণাত্মক) সংখ্যা হবে।}
\end{quote}

\textbf{প্রথম P সংখ্যক বিজোড় সংখ্যার যোগফল কত?}

ক) P\^{}2−1

\textbf{\hl{খ) P\^{}2 ✔}}

গ) P\^{}2+1

ঘ) কোনোটিই নয়

\textbf{\hl{✔ সঠিক উত্তর: খ) P\^{}2}}

\begin{quote}
\textbf{ব্যাখ্যা:} \emph{আমরা জানি, প্রথম n সংখ্যক বিজোড় সংখ্যার সমষ্টি =
n\^{}2}

\emph{∴ প্রথম P সংখ্যক বিজোড় সংখ্যার যোগফল = P\^{}2 ।}
\end{quote}

\textbf{যদি n এবং p দুটি অযুগ্ম সংখ্যা হয়, তবে নিম্নের কোনটি অবশ্যই যুগ্ম সংখ্যা
হবে?}

\textbf{\hl{ক) n + p ✔}}

খ) np

গ) np + 2

ঘ) n + p + 1

\textbf{\hl{✔ সঠিক উত্তর: ক) n + p}}

\begin{quote}
\textbf{ব্যাখ্যা:} \emph{এখানে অযুগ্ম = বিজোড় এবং যুগ্ম = জোড় সংখ্যা।}

\emph{ধরি, বিজোড় সংখ্যা n=1 এবং p=3 ।}

\emph{ক) n+p=1+3=4 (জোড়/যুগ্ম সংখ্যা)}

\emph{খ) np=1×3=3 (বিজোড়)}

\emph{গ) np+2=3+2=5 (বিজোড়)}

\emph{ঘ) n+p+1=1+3+1=5 (বিজোড়)}

\emph{∴ n + p অবশ্যই একটি যুগ্ম (জোড়) সংখ্যা হবে।}
\end{quote}

\textbf{m ও n দুটি বিজোড় সংখ্যা হলে, নিম্নের জোড় সংখ্যাটি হলো-}

ক) m-n + 1

খ) mn

গ) mn + 2

\textbf{\hl{ঘ) m + n ✔}}

\textbf{\hl{✔ সঠিক উত্তর: ঘ) m + n}}

\begin{quote}
\textbf{ব্যাখ্যা:} \emph{আমরা জানি, দুটি বিজোড় সংখ্যার যোগফল সর্বদা একটি জোড়
সংখ্যা হয়।}

\emph{ধরি, m=3 এবং n=5 (উভয়েই বিজোড়)।}

\emph{তাহলে, m+n=3+5=8 (যা জোড় সংখ্যা)।}

\emph{∴ নির্ণেয় জোড় সংখ্যাটি (m + n)।}
\end{quote}

\textbf{১ কোন ধরনের সংখ্যা-}

ক) মৌলিক

খ) যৌগিক

গ) উভয়ই

\textbf{\hl{ঘ) কোনোটিই নয় ✔}}

\textbf{\hl{✔ সঠিক উত্তর: ঘ) কোনোটিই নয়}}

\begin{quote}
\textbf{ব্যাখ্যা:} \emph{মৌলিক সংখ্যা: ১ এর চেয়ে বড় যে সকল সংখ্যাকে ১ ও ঐ
সংখ্যা ব্যতীত অন্য কোনো সংখ্যা দ্বারা ভাগ করা যায় না তাকে মৌলিক সংখ্যা বলে।}

\emph{যৌগিক সংখ্যা: যে সংখ্যার ১ ও ঐ সংখ্যা ব্যতীত একাধিক গুণনীয়ক আছে তাকে
যৌগিক সংখ্যা বলে।}

\emph{∴ ১ মৌলিক বা যৌগিক কোনোটিই নয়।}
\end{quote}

\textbf{ক্ষুদ্রতম মৌলিক সংখ্যা হলো :}

ক) ৭

\textbf{\hl{খ) ২ ✔}}

গ) ৩

ঘ) ৫

\textbf{\hl{✔ সঠিক উত্তর: খ) ২}}

\begin{quote}
\textbf{ব্যাখ্যা:} \emph{১ অপেক্ষা বড় যে সকল স্বাভাবিক সংখ্যার ১ এবং ঐ সংখ্যা
ছাড়া অন্য কোনো গুণনীয়ক নেই, তাদেরকে মৌলিক সংখ্যা বলে।}

\emph{সবচেয়ে ছোট মৌলিক সংখ্যা হলো ২ (এবং ২-ই একমাত্র জোড় মৌলিক সংখ্যা)।}

\emph{∴ ক্ষুদ্রতম মৌলিক সংখ্যা হলো ২।}
\end{quote}

\textbf{কোনটি মৌলিক সংখ্যা নয়?}

ক) ২৬৩

খ) ২৩৩

\textbf{\hl{গ) ২৫৩ ✔}}

ঘ) ২৪১

\textbf{\hl{✔ সঠিক উত্তর: গ) ২৫৩}}

\begin{quote}
\textbf{ব্যাখ্যা:} \emph{এখানে,}

\emph{২৬৩, ২৩৩ এবং ২৪১ সংখ্যাগুলোর ১ ও ঐ সংখ্যা ব্যতীত অন্য কোনো গুণনীয়ক নেই,
তাই এগুলো মৌলিক সংখ্যা।}

\emph{কিন্তু, ২৫৩ সংখ্যাটি ১১ দ্বারা বিভাজ্য ( 253÷11=23 )। অর্থাৎ ২৫৩ এর
গুণনীয়ক রয়েছে (১, ১১, ২৩, ২৫৩)।}

\emph{∴ ২৫৩ মৌলিক সংখ্যা নয়।}
\end{quote}

\textbf{কোনটি মৌলিক সংখ্যা?}

ক) ১৬

খ) ৯

\textbf{\hl{গ) ৩ ✔}}

ঘ) ১২

\textbf{\hl{✔ সঠিক উত্তর: গ) ৩}}

\begin{quote}
\textbf{ব্যাখ্যা:} \emph{যে সকল সংখ্যাকে ১ এবং ঐ সংখ্যা ছাড়া অন্য কোনো সংখ্যা
দ্বারা ভাগ করা যায় না তাকে মৌলিক সংখ্যা বলে।}

\emph{এখানে ১৬, ৯ এবং ১২ যৌগিক সংখ্যা। কিন্তু ৩ কে ১ এবং ৩ ছাড়া অন্য কিছু
দিয়ে ভাগ করা যায় না।}

\emph{∴ ৩ একটি মৌলিক সংখ্যা।}
\end{quote}

\textbf{নিচের কোনটি মৌলিক সংখ্যা}

ক) ৬

খ) ৮

গ) ৯

\textbf{\hl{ঘ) ৭ ✔}}

\textbf{\hl{✔ সঠিক উত্তর: ঘ) ৭}}

\begin{quote}
\textbf{ব্যাখ্যা:} \emph{যে সকল সংখ্যা ১ এবং ঐ সংখ্যা ছাড়া অন্য কোনো সংখ্যা
দ্বারা বিভাজ্য নয় তাকে মৌলিক সংখ্যা বলে।}

\emph{এখানে ৬, ৮, ৯ এর একাধিক গুণনীয়ক আছে, কিন্তু ৭ কেবল ১ এবং ৭ দ্বারা
বিভাজ্য।}

\emph{∴ ৭ একটি মৌলিক সংখ্যা।}
\end{quote}

\textbf{কোনটি মৌলিক সংখ্যা?}

ক) ৯

খ) ২১

\textbf{\hl{গ) ২৩ ✔}}

ঘ) ২৭

\textbf{\hl{✔ সঠিক উত্তর: গ) ২৩}}

\begin{quote}
\textbf{ব্যাখ্যা:} \emph{অপশনগুলোর মধ্যে:}

\emph{৯ = ৩ × ৩ (যৌগিক)}

\emph{২১ = ৩ × ৭ (যৌগিক)}

\emph{২৭ = ৩ × ৯ (যৌগিক)}

\emph{কিন্তু ২৩ কে ১ এবং ২৩ ব্যতীত অন্য কোনো সংখ্যা দ্বারা ভাগ করা যায় না।}

\emph{∴ ২৩ একটি মৌলিক সংখ্যা।}
\end{quote}

\textbf{নিচের কোনটি মৌলিক সংখ্যা নয়?}

ক) ১৩

খ) ২৩

গ) ২৯

\textbf{\hl{ঘ) ৩৯ ✔}}

\textbf{\hl{✔ সঠিক উত্তর: ঘ) ৩৯}}

\begin{quote}
\textbf{ব্যাখ্যা:} \emph{এখানে, ১৩, ২৩ ও ২৯ অন্য কোনো সংখ্যা দ্বারা বিভাজ্য
নয়, অর্থাৎ এরা মৌলিক সংখ্যা।}

\emph{কিন্তু ৩৯ সংখ্যাটি ৩ ও ১৩ দ্বারা বিভাজ্য ( 39=3×13 )।}

\emph{∴ ৩৯ মৌলিক সংখ্যা নয়।}
\end{quote}

\textbf{নিচের কোনটি মৌলিক সংখ্যা?}

\textbf{\hl{ক) ৪৭ ✔}}

খ) ৮৭

গ) ৯১

ঘ) ১৪৩

\textbf{\hl{✔ সঠিক উত্তর: ক) ৪৭}}

\begin{quote}
\textbf{ব্যাখ্যা:} \emph{অপশনগুলো বিশ্লেষণ করে পাই:}

\emph{৮৭ = ৩ × ২৯ (যৌগিক)}

\emph{৯১ = ৭ × ১৩ (যৌগিক)}

\emph{১৪৩ = ১১ × ১৩ (যৌগিক)}

\emph{কিন্তু ৪৭ সংখ্যাটি ১ ও ৪৭ ব্যতীত অন্য কোনো সংখ্যা দ্বারা বিভাজ্য নয়।}

\emph{∴ ৪৭ একটি মৌলিক সংখ্যা।}
\end{quote}

\textbf{নিচের কোন সংখ্যাটি মৌলিক?}

ক) ১০৫

\textbf{\hl{খ) ৫৩ ✔}}

গ) ২১

ঘ) ১২৪

\textbf{\hl{✔ সঠিক উত্তর: খ) ৫৩}}

\begin{quote}
\textbf{ব্যাখ্যা:} \emph{১ থেকে বড় যেসব সংখ্যাকে ১ ও ঐ সংখ্যা ব্যতীত অন্য কোনো
সংখ্যা দ্বারা ভাগ করা যায় না তাদেরকে মৌলিক সংখ্যা বলে।}

\emph{এখানে, ১০৫ ও ২১ কে ৩ দ্বারা ভাগ করা যায় এবং ১২৪ কে ২ দ্বারা ভাগ করা
যায়।}

\emph{∴ ৫৩ একটি মৌলিক সংখ্যা।}
\end{quote}

\textbf{নিচে কোনটি মৌলিক সংখ্যা?}

ক) ৯১

খ) ৮৭

গ) ৬৩

\textbf{\hl{ঘ) ৫৯ ✔}}

\textbf{\hl{✔ সঠিক উত্তর: ঘ) ৫৯}}

\begin{quote}
\textbf{ব্যাখ্যা:} \emph{অপশনগুলোর মধ্যে:}

\emph{৯১ = ৭ × ১৩ (যৌগিক)}

\emph{৮৭ = ৩ × ২৯ (যৌগিক)}

\emph{৬৩ = ৭ × ৯ (যৌগিক)}

\emph{কিন্তু ৫৯ সংখ্যাটিকে ১ ও ৫৯ ব্যতীত অন্য কোনো সংখ্যা দ্বারা ভাগ করা যায়
না।}

\emph{∴ ৫৯ একটি মৌলিক সংখ্যা।}
\end{quote}

\textbf{নিচের কোনটি মৌলিক সংখ্যা?}

ক) ৯১

খ) ৮৭

গ) ৬৩

\textbf{\hl{ঘ) ৫৯ ✔}}

\textbf{\hl{✔ সঠিক উত্তর: ঘ) ৫৯}}

\begin{quote}
\textbf{ব্যাখ্যা:} \emph{অপশনগুলোর বিশ্লেষণে:}

\emph{(ক) ৯১ = ৭ × ১৩}

\emph{(খ) ৮৭ = ৩ × ২৯}

\emph{(গ) ৬৩ = ৭ × ৯}

\emph{(ঘ) ৫৯ = ১ × ৫৯}

\emph{∴ ৫৯ একটি মৌলিক সংখ্যা।}
\end{quote}

\textbf{১ থেকে ১০০ পর্যন্ত কতটি মৌলিক সংখ্যা আছে?}

ক) ১৫টি

খ) ১৮টি

গ) ২০টি

\textbf{\hl{ঘ) ২৫টি ✔}}

\textbf{\hl{✔ সঠিক উত্তর: ঘ) ২৫টি}}

\begin{quote}
\textbf{ব্যাখ্যা:} \emph{১ থেকে ১০০ পর্যন্ত মৌলিক সংখ্যাগুলো হলো:}

\emph{১ - ১০ পর্যন্ত: ৪টি (২, ৩, ৫, ৭)}

\emph{১১ - ২০ পর্যন্ত: ৪টি (১১, ১৩, ১৭, ১৯)}

\emph{২১ - ৩০ পর্যন্ত: ২টি (২৩, ২৯)}

\emph{৩১ - ৪০ পর্যন্ত: ২টি (৩১, ৩৭)}

\emph{৪১ - ৫০ পর্যন্ত: ৩টি (৪১, ৪৩, ৪৭)}

\emph{৫১ - ৬০ পর্যন্ত: ২টি (৫৩, ৫৯)}

\emph{৬১ - ৭০ পর্যন্ত: ২টি (৬১, ৬৭)}

\emph{৭১ - ৮০ পর্যন্ত: ৩টি (৭১, ৭৩, ৭৯)}

\emph{৮১ - ৯০ পর্যন্ত: ২টি (৮৩, ৮৯)}

\emph{৯১ - ১০০ পর্যন্ত: ১টি (৯৭)}

\emph{∴ মোট মৌলিক সংখ্যা = ২৫টি।}
\end{quote}

\textbf{১ হতে ২০ পর্যন্ত মৌলিক সংখ্যা কয়টি হবে?}

ক) ৯টি

খ) ৭টি

\textbf{\hl{গ) ৮টি ✔}}

ঘ) ১০টি

\textbf{\hl{✔ সঠিক উত্তর: গ) ৮টি}}

\begin{quote}
\textbf{ব্যাখ্যা:} \emph{১ হতে ২০ পর্যন্ত মৌলিক সংখ্যাগুলো হচ্ছে: ২, ৩, ৫, ৭,
১১, ১৩, ১৭ এবং ১৯।}

\emph{∴ মোট মৌলিক সংখ্যা ৮টি।}
\end{quote}

\textbf{১ থেকে ৩০ পর্যন্ত মৌলিক সংখ্যা কয়টি?}

ক) ১২টি

\textbf{\hl{খ) ১০টি ✔}}

গ) ৯টি

ঘ) ৮টি

\textbf{\hl{✔ সঠিক উত্তর: খ) ১০টি}}

\begin{quote}
\textbf{ব্যাখ্যা:} \emph{১ থেকে ৩০ পর্যন্ত মৌলিক সংখ্যাগুলো হলো: ২, ৩, ৫, ৭,
১১, ১৩, ১৭, ১৯, ২৩ ও ২৯।}

\emph{∴ মোট মৌলিক সংখ্যা ১০টি।}
\end{quote}

\textbf{৩ হতে ৩০ পর্যন্ত মৌলিক সংখ্যা কয়টি?}

ক) ১১টি

খ) ১০টি

গ) ৮টি

\textbf{\hl{ঘ) ৯টি ✔}}

\textbf{\hl{✔ সঠিক উত্তর: ঘ) ৯টি}}

\begin{quote}
\textbf{ব্যাখ্যা:} \emph{৩ থেকে ৩০ পর্যন্ত মৌলিক সংখ্যাগুলো হলো: ৩, ৫, ৭, ১১,
১৩, ১৭, ১৯, ২৩ ও ২৯।}

\emph{∴ মোট মৌলিক সংখ্যা ৯টি।}
\end{quote}

\textbf{১ থেকে ৪০ পর্যন্ত কতগুলি মৌলিক সংখ্যা বিদ্যমান?}

ক) ১০

খ) ১১

\textbf{\hl{গ) ১২ ✔}}

ঘ) ১৩

\textbf{\hl{✔ সঠিক উত্তর: গ) ১২}}

\begin{quote}
\textbf{ব্যাখ্যা:} \emph{১ থেকে ৪০ পর্যন্ত মৌলিক সংখ্যাগুলো হলো: ২, ৩, ৫, ৭,
১১, ১৩, ১৭, ১৯, ২৩, ২৯, ৩১ এবং ৩৭।}

\emph{∴ মোট মৌলিক সংখ্যা ১২টি।}
\end{quote}

\textbf{১ থেকে ৫০ পর্যন্ত মৌলিক সংখ্যা কয়টি?}

\textbf{\hl{ক) ১৫ ✔}}

খ) ১৭

গ) ২৫

ঘ) ২৭

\textbf{\hl{✔ সঠিক উত্তর: ক) ১৫}}

\begin{quote}
\textbf{ব্যাখ্যা:} \emph{১ থেকে ৫০ পর্যন্ত মৌলিক সংখ্যাগুলো হলো: ২, ৩, ৫, ৭,
১১, ১৩, ১৭, ১৯, ২৩, ২৯, ৩১, ৩৭, ৪১, ৪৩ ও ৪৭।}

\emph{∴ মোট মৌলিক সংখ্যা ১৫টি।}
\end{quote}

\textbf{২০ থেকে ১০০ এর মধ্যবর্তী মৌলিক সংখ্যা কতটি?}

\textbf{\hl{ক) ১৭ টি ✔}}

খ) ১৫ টি

গ) ১৩ টি

ঘ) ২১ টি

\textbf{\hl{✔ সঠিক উত্তর: ক) ১৭ টি}}

\begin{quote}
\textbf{ব্যাখ্যা:} \emph{১ থেকে ১০০ পর্যন্ত মোট মৌলিক সংখ্যা ২৫টি।}

\emph{এর মধ্যে ১ থেকে ২০ পর্যন্ত মৌলিক সংখ্যা ৮টি (২, ৩, ৫, ৭, ১১, ১৩, ১৭,
১৯)।}

\emph{∴ ২০ থেকে ১০০ এর মধ্যবর্তী মৌলিক সংখ্যা = 25−8=17 টি।}
\end{quote}

\textbf{২৫ থেকে ৫৫ এর মধ্যে মৌলিক সংখ্যা কয়টি?}

\textbf{\hl{ক) ৭টি ✔}}

খ) ৪টি

গ) ৬টি

ঘ) ৫টি

\textbf{\hl{✔ সঠিক উত্তর: ক) ৭টি}}

\begin{quote}
\textbf{ব্যাখ্যা:} \emph{২৫ থেকে ৫৫ এর মধ্যে মৌলিক সংখ্যাগুলো হলো: ২৯, ৩১,
৩৭, ৪১, ৪৩, ৪৭ ও ৫৩।}

\emph{∴ মোট মৌলিক সংখ্যা ৭টি।}
\end{quote}

\textbf{৩১ থেকে ৫০ পর্যন্ত মৌলিক সংখ্যা কতটি?}

ক) ৪

\textbf{\hl{খ) ৫ ✔}}

গ) ৬

ঘ) ৭

\textbf{\hl{✔ সঠিক উত্তর: খ) ৫}}

\begin{quote}
\textbf{ব্যাখ্যা:} \emph{৩১ থেকে ৫০ পর্যন্ত মৌলিক সংখ্যাগুলো হলো: ৩১, ৩৭, ৪১,
৪৩ ও ৪৭।}

\emph{∴ মোট মৌলিক সংখ্যা ৫টি।}
\end{quote}

\textbf{৪১ হতে ৯০ পর্যন্ত সংখ্যাগুলির মধ্যে মৌলিক সংখ্যা কয়টি?}

\textbf{\hl{ক) ১২ ✔}}

খ) ১৩

গ) ১৪

ঘ) ১০

\textbf{\hl{✔ সঠিক উত্তর: ক) ১২}}

\begin{quote}
\textbf{ব্যাখ্যা:} \emph{৪১ হতে ৯০ পর্যন্ত মৌলিক সংখ্যাগুলো হলো: ৪১, ৪৩, ৪৭,
৫৩, ৫৯, ৬১, ৬৭, ৭১, ৭৩, ৭৯, ৮৩ ও ৮৯।}

\emph{∴ মোট মৌলিক সংখ্যা ১২টি।}
\end{quote}

\textbf{43 থেকে 60 এর মধ্যে মৌলিক সংখ্যা কতটি?}

ক) 5

খ) 3

গ) 7

\textbf{\hl{ঘ) 4 ✔}}

\textbf{\hl{✔ সঠিক উত্তর: ঘ) 4}}

\begin{quote}
\textbf{ব্যাখ্যা:} \emph{43 ও 60 এর মধ্যে মৌলিক সংখ্যাগুলো হলো: 43, 47, 53,
59।}

\emph{∴ মোট মৌলিক সংখ্যা 4টি।}
\end{quote}

\textbf{৭০ থেকে ৮০ এর মধ্যে মৌলিক সংখ্যা কতটি?}

\textbf{\hl{ক) ০৩টি ✔}}

খ) ০৫টি

গ) ০৭টি

ঘ) ০২টি

\textbf{\hl{✔ সঠিক উত্তর: ক) ০৩টি}}

\begin{quote}
\textbf{ব্যাখ্যা:} \emph{৭০ থেকে ৮০ এর মধ্যে মৌলিক সংখ্যাগুলো হলো: ৭১, ৭৩ এবং
৭৯।}

\emph{∴ মোট মৌলিক সংখ্যা ৩টি।}
\end{quote}

\textbf{৫৬ থেকে ১০০- এর মধ্যে মৌলিক সংখ্যা কয়টি?}

ক) 8

\textbf{\hl{খ) 9 ✔}}

গ) 10

ঘ) 11

ঙ) 12

\textbf{\hl{✔ সঠিক উত্তর: খ) 9}}

\begin{quote}
\textbf{ব্যাখ্যা:} \emph{৫৬ ও ১০০ এর মধ্যে ৯টি মৌলিক সংখ্যা রয়েছে।}

\emph{সংখ্যাগুলো হলো: ৫৯, ৬১, ৬৭, ৭১, ৭৩, ৭৯, ৮৩, ৮৯ এবং ৯৭।}

\emph{∴ মোট মৌলিক সংখ্যা ৯টি।}
\end{quote}

\textbf{৯০ থেকে ১০০ এর মধ্যে কয়টি মৌলিক সংখ্যা আছে?}

ক) ২টি

\textbf{\hl{খ) ১টি ✔}}

গ) ৩টি

ঘ) একটিও নয়

\textbf{\hl{✔ সঠিক উত্তর: খ) ১টি}}

\begin{quote}
\textbf{ব্যাখ্যা:} \emph{৯০ থেকে ১০০ এর মধ্যে মৌলিক সংখ্যা ১টি, যা হলো: ৯৭।}

\emph{∴ সঠিক উত্তর ১টি।}
\end{quote}

\textbf{নিচের কোন সংখ্যাটি মৌলিক সংখ্যা নয়?}

ক) ২৬৩

খ) ২৩৩

গ) ২৪১

\textbf{\hl{ঘ) ২৫৩ ✔}}

\textbf{\hl{✔ সঠিক উত্তর: ঘ) ২৫৩}}

\begin{quote}
\textbf{ব্যাখ্যা:} \emph{অপশনগুলোর মধ্যে ২৬৩, ২৩৩ এবং ২৪১ হলো মৌলিক সংখ্যা।}

\emph{কিন্তু ২৫৩ সংখ্যাটি ১১ দ্বারা বিভাজ্য ( 253÷11=23 )।}

\emph{∴ ২৫৩ সংখ্যাটি মৌলিক নয়।}
\end{quote}

\textbf{√(529) এর বর্গমূল কী ধরণের সংখ্যা?}

ক) মূলদ সংখ্যা

\textbf{\hl{খ) পূর্ণসংখ্যা ✔}}

গ) স্বাভাবিক সংখ্যা

ঘ) অমূলদ সংখ্যা

\textbf{\hl{✔ সঠিক উত্তর: খ) পূর্ণসংখ্যা}}

\begin{quote}
\textbf{ব্যাখ্যা:} \emph{৫২৯-এর বর্গমূল হলো ২৩, যা একটি স্বাভাবিক সংখ্যা,
পূর্ণসংখ্যা এবং মূলদ সংখ্যা। এ ধরনের প্রতিযোগিতামূলক পরীক্ষায় (যেমন বিসিএসআইআর)
যখন একটি নির্দিষ্ট সঠিক উত্তর বেছে নিতে হয়, তখন ২৩ একটি ধনাত্মক পূর্ণসংখ্যা
হওয়ায় এটিকে প্রধানত পূর্ণসংখ্যা সংখ্যা হিসেবে গণ্য করা হয়।}
\end{quote}

\textbf{নিচের কোন ক্রমজোড়টি সহমৌলিক ?}

ক) (৪,৬)

খ) (৬,৯)

গ) (৯,১২)

\textbf{\hl{ঘ) (১২,১৭) ✔}}

\textbf{\hl{✔ সঠিক উত্তর: ঘ) (১২,১৭)}}

\begin{quote}
\textbf{ব্যাখ্যা:} \emph{দুটি সংখ্যার সাধারণ উৎপাদক বা গ.সা.গু কেবল ১ হলে
সংখ্যা দুটি পরস্পরের সহমৌলিক।}

\emph{এখানে, ১২ = ১ × ২ × ২ × ৩}

\emph{এবং ১৭ = ১ × ১৭}

\emph{দেখা যাচ্ছে, ১২ ও ১৭ এর ১ ভিন্ন অন্য কোনো সাধারণ উৎপাদক নেই (গ.সা.গু =
১)।}

\emph{∴ (১২, ১৭) ক্রমজোড়টি সহমৌলিক।}
\end{quote}

\textbf{১ থেকে ৪৯ পর্যন্ত স্বাভাবিক (পূর্ণ) সংখ্যাগুলোর গড় কত?}

ক) ২৩

\textbf{\hl{খ) ২৫ ✔}}

গ) ২৪.৫

ঘ) ২৫.৫

\textbf{\hl{✔ সঠিক উত্তর: খ) ২৫}}

\begin{quote}
\textbf{ব্যাখ্যা:} \emph{১ থেকে শুরু করে ক্রমিক স্বাভাবিক সংখ্যার গড় = (প্রথম
পদ+শেষ পদ/2)}

\emph{∴ ১ থেকে ৪৯ পর্যন্ত সংখ্যাগুলোর গড় = (1+49/2)=(50/2)=25 ।}
\end{quote}

\textbf{৬০ থেকে ৮০ এর মধ্যবর্তী বৃহত্তম ও ক্ষুদ্রতম মৌলিক সংখ্যাদ্বয়ের গড় কত ?}

\textbf{\hl{ক) ৭০ ✔}}

খ) ৬৭

গ) ৮০

ঘ) ৭৭

\textbf{\hl{✔ সঠিক উত্তর: ক) ৭০}}

\begin{quote}
\textbf{ব্যাখ্যা:} \emph{৬০ এবং ৮০ এর মধ্যে ক্ষুদ্রতম মৌলিক সংখ্যাটি হলো ৬১ এবং
বৃহত্তম মৌলিক সংখ্যাটি হলো ৭৯।}

\emph{∴ সংখ্যাদ্বয়ের গড় = (61+79/2)=(140/2)=70 ।}
\end{quote}

\textbf{১০ থেকে ৬০ পর্যন্ত যে সকল মৌলিক সংখ্যার একক স্থানীয় অংক ৯, তাদের
সমষ্টি কত?}

ক) ১৩০

\textbf{\hl{খ) ১০৭ ✔}}

গ) ১১৩

ঘ) ১৪৬

\textbf{\hl{✔ সঠিক উত্তর: খ) ১০৭}}

\begin{quote}
\textbf{ব্যাখ্যা:} \emph{১০ থেকে ৬০ পর্যন্ত একক স্থানীয় অংক ৯ বিশিষ্ট মৌলিক
সংখ্যাগুলো হলো: ১৯, ২৯ এবং ৫৯।}

\emph{∴ এদের সমষ্টি = ১৯ + ২৯ + ৫৯ = ১০৭।}
\end{quote}

\textbf{১৫ থেকে ৩০ এর মধ্যে মৌলিক সংখ্যাগুলোর যোগফল কত?}

\textbf{\hl{ক) ৮৮ ✔}}

খ) ৭৮

গ) ৮৭

ঘ) ৬৫

\textbf{\hl{✔ সঠিক উত্তর: ক) ৮৮}}

\begin{quote}
\textbf{ব্যাখ্যা:} \emph{১৫ থেকে ৩০ এর মধ্যে মৌলিক সংখ্যাগুলো হলো: ১৭, ১৯, ২৩
এবং ২৯।}

\emph{∴ মৌলিক সংখ্যাগুলোর যোগফল = ১৭ + ১৯ + ২৩ + ২৯ = ৮৮।}
\end{quote}

\textbf{১৫ থেকে ৬৫ পর্যন্ত যে সকল মৌলিক সংখ্যার একক স্থানীয় অংক ৯ তাদের
সমষ্টি কত?}

ক) ১৫৪

খ) ১০৮

\textbf{\hl{গ) ১০৭ ✔}}

ঘ) ১৩৮

\textbf{\hl{✔ সঠিক উত্তর: গ) ১০৭}}

\begin{quote}
\textbf{ব্যাখ্যা:} \emph{১৫ থেকে ৬৫ পর্যন্ত একক স্থানীয় অংক ৯ বিশিষ্ট মৌলিক
সংখ্যাগুলো হলো: ১৯, ২৯ এবং ৫৯।}

\emph{∴ এদের সমষ্টি = ১৯ + ২৯ + ৫৯ = ১০৭।}
\end{quote}

\textbf{৩০ থেকে ৪০-এর মধ্যবর্তী বৃহত্তম ও ক্ষুদ্রতম মৌলিক সংখ্যার ব্যবধান কত?}

ক) ৩৫

খ) ৪২

গ) ৪৮

ঘ) ৫৫

\textbf{\hl{ঙ) ৬ ✔}}

\textbf{\hl{✔ সঠিক উত্তর: ঙ) ৬}}

\begin{quote}
\textbf{ব্যাখ্যা:} \emph{৩০ থেকে ৪০-এর মধ্যবর্তী ক্ষুদ্রতম মৌলিক সংখ্যা হলো ৩১
এবং বৃহত্তম মৌলিক সংখ্যা হলো ৩৭।}

\emph{∴ এদের ব্যবধান (পার্থক্য) = ৩৭ − ৩১ = ৬।}
\end{quote}

\textbf{৩০ থেকে ৯০ এর মধ্যবর্তী বৃহত্তম ও ক্ষুদ্রতম মৌলিক সংখ্যার অন্তর কত?}

\textbf{\hl{ক) ৫৮ ✔}}

খ) ৪২

গ) ৬৮

ঘ) ৬২

\textbf{\hl{✔ সঠিক উত্তর: ক) ৫৮}}

\begin{quote}
\textbf{ব্যাখ্যা:} \emph{৩০ থেকে ৯০ এর মধ্যে বৃহত্তম মৌলিক সংখ্যা = ৮৯}

\emph{৩০ থেকে ৯০ এর মধ্যে ক্ষুদ্রতম মৌলিক সংখ্যা = ৩১}

\emph{∴ তাদের পার্থক্য (অন্তর) = ৮৯ − ৩১ = ৫৮।}
\end{quote}

\textbf{৪০ থেকে ১০০ পর্যন্ত বৃহত্তম ও ক্ষুদ্রতম মৌলিক সংখ্যার অন্তর কত?}

ক) ৫৯

\textbf{\hl{খ) ৫৬ ✔}}

গ) ৬০

ঘ) ৭০

\textbf{\hl{✔ সঠিক উত্তর: খ) ৫৬}}

\begin{quote}
\textbf{ব্যাখ্যা:} \emph{৪০ থেকে ১০০ এর মধ্যে বৃহত্তম মৌলিক সংখ্যা হলো ৯৭ এবং
ক্ষুদ্রতম মৌলিক সংখ্যা হলো ৪১।}

\emph{∴ সংখ্যা দুটির অন্তর = ৯৭ − ৪১ = ৫৬।}
\end{quote}

\textbf{৬০ থেকে ৮০ এর মধ্যবর্তী বৃহত্তম ও ক্ষুদ্রতম মৌলিক সংখ্যার অন্তর হবে-}

ক) ৮

খ) ১২০

\textbf{\hl{গ) ১৮ ✔}}

ঘ) ১৪০

\textbf{\hl{✔ সঠিক উত্তর: গ) ১৮}}

\begin{quote}
\textbf{ব্যাখ্যা:} \emph{৬০ থেকে ৮০ এর মধ্যবর্তী বৃহত্তম মৌলিক সংখ্যাটি হলো ৭৯
এবং ক্ষুদ্রতম মৌলিক সংখ্যাটি হলো ৬১।}

\emph{∴ এদের অন্তর = ৭৯ − ৬১ = ১৮।}
\end{quote}

\textbf{১ থেকে ১০ পর্যন্ত যে মৌলিক সংখ্যাগুলো আছে তাদের গুণফল কত?}

ক) ৩৫

খ) ১০৫

\textbf{\hl{গ) ২১০ ✔}}

ঘ) ১৯৫

\textbf{\hl{✔ সঠিক উত্তর: গ) ২১০}}

\begin{quote}
\textbf{ব্যাখ্যা:} \emph{১ থেকে ১০ পর্যন্ত মৌলিক সংখ্যাগুলো হলো: ২, ৩, ৫ এবং
৭।}

\emph{∴ এদের গুণফল = ২ × ৩ × ৫ × ৭ = ২১০।}
\end{quote}

\textbf{৩ থেকে ১০ পর্যন্ত মৌলিক সংখ্যাগুলোর গুণফল কত?}

ক) ৩৫

খ) ১৫

\textbf{\hl{গ) ১০৫ ✔}}

ঘ) ৪২

\textbf{\hl{✔ সঠিক উত্তর: গ) ১০৫}}

\begin{quote}
\textbf{ব্যাখ্যা:} \emph{৩ থেকে ১০ পর্যন্ত মৌলিক সংখ্যাগুলো হলো: ৩, ৫ এবং ৭।}

\emph{∴ এদের গুণফল = ৩ × ৫ × ৭ = ১০৫।}
\end{quote}

\textbf{৫ থেকে ১০ পর্যন্ত মৌলিক সংখ্যাগুলোর গুণফল হবে-}

ক) ২০

খ) ২৫

গ) ৩০

\textbf{\hl{ঘ) ৩৫ ✔}}

\textbf{\hl{✔ সঠিক উত্তর: ঘ) ৩৫}}

\begin{quote}
\textbf{ব্যাখ্যা:} \emph{৫ থেকে ১০ পর্যন্ত মৌলিক সংখ্যাগুলো হলো: ৫ এবং ৭।}

\emph{∴ এদের গুণফল = ৫ × ৭ = ৩৫।}
\end{quote}

\textbf{নিচের কোনটি ৩ ও ৬ এর গুণিতক?}

ক) ৩

খ) ৯

গ) ১৫

\textbf{\hl{ঘ) ১৮ ✔}}

\textbf{\hl{✔ সঠিক উত্তর: ঘ) ১৮}}

\begin{quote}
\textbf{ব্যাখ্যা:} \emph{৩ এর গুণিতকসমূহ: ৩, ৬, ৯, ১২, ১৫, ১৮, ২১ ...}

\emph{৬ এর গুণিতকসমূহ: ৬, ১২, ১৮, ২৪ ...}

\emph{দেখা যাচ্ছে, ৩ ও ৬ উভয়েরই সাধারণ গুণিতক হলো ১৮ (যেহেতু ৩ × ৬ = ১৮ এবং
৬ × ৩ = ১৮)।}

\emph{∴ নির্ণেয় গুণিতক ১৮।}
\end{quote}

\textbf{p এবং q বিজোড় সংখ্যা হলে নিম্নে কোন রাশিটি জোড় সংখ্যা হবে?}

ক) p + q + 1

\textbf{\hl{খ) p + q ✔}}

গ) কোনোটিই নয়

ঘ) p + q - 1

\textbf{\hl{✔ সঠিক উত্তর: খ) p + q}}

\begin{quote}
\textbf{ব্যাখ্যা:} \emph{আমরা জানি, দুটি বিজোড় সংখ্যার যোগফল সর্বদা জোড় সংখ্যা
হয়।}

\emph{ধরি, p=1 এবং q=3 (উভয়েই বিজোড়)।}

\emph{ক) p+q+1=1+3+1=5 (বিজোড়)}

\emph{খ) p+q=1+3=4 (জোড়)}

\emph{ঘ) p+q−1=1+3−1=3 (বিজোড়)}

\emph{∴ p + q রাশিটি একটি জোড় সংখ্যা।}
\end{quote}

\textbf{X -এর মান একটি বেজোড় সংখ্যা হলে নিম্নের কোনটির মান জোড় সংখ্যা হবে?}

ক) 2x + 1

\textbf{\hl{খ) 2(x + 1) ✔}}

গ) 2x - 1

ঘ) x - 2

\textbf{\hl{✔ সঠিক উত্তর: খ) 2(x + 1)}}

\begin{quote}
\textbf{ব্যাখ্যা:} \emph{যে কোনো সংখ্যাকে ২ দ্বারা গুণ করলে তা সর্বদা জোড়
সংখ্যা হয়।}

\emph{এখানে, 2(x+1)=2x+2 ।}

\emph{ধরি, বিজোড় সংখ্যা x=3 ।}

\emph{তাহলে, 2(3+1)=2×4=8 (যা জোড় সংখ্যা)।}

\emph{∴ 2(x + 1) এর মান একটি জোড় সংখ্যা হবে।}
\end{quote}

\section{শুধু MCQ অনুশীলন (তাত্ত্বিক নোট
নেই)}\label{ux9b6ux9a7-mcq-ux985ux9a8ux9b6ux9b2ux9a8-ux9a4ux9a4ux9a4ux9acux995-ux9a8ux99f-ux9a8ux987}

\section{67. বৃত্ত MCQ প্রশ্ন উত্তর
অনুশীলন}\label{ux9acux9a4ux9a4-mcq-ux9aaux9b0ux9b6ux9a8-ux989ux9a4ux9a4ux9b0-ux985ux9a8ux9b6ux9b2ux9a8}

\textbf{বৃত্তের যেকোনো বিন্দুর সংযোজক রেখাংশকে বৃত্তের কী বলা হয়?}

\textbf{\hl{ক) জ্যা ✔}}

খ) ব্যাস

গ) পরিধি

ঘ) ব্যাসার্ধ

\textbf{\hl{✔ সঠিক উত্তর: ক) জ্যা}}

\begin{quote}
\textbf{ব্যাখ্যা:} \emph{বৃত্তের পরিধির ওপর যেকোনো দুটি বিন্দুর সংযোজক রেখাংশকে
জ্যা বলা হয়। বৃত্তের কেন্দ্রগামী জ্যাকে ব্যাস বলে।}
\end{quote}

\textbf{৪ সে.মি. ব্যাস বিশিষ্ট বৃত্ত আঁকলে কেন্দ্র হতে পরিধির দূরত্ব কত হবে?}

ক) ১ সে.মি.

\textbf{\hl{খ) ২ সে.মি. ✔}}

গ) ৪ সে.মি.

ঘ) ৮ সে.মি.

\textbf{\hl{✔ সঠিক উত্তর: খ) ২ সে.মি.}}

\begin{quote}
\textbf{ব্যাখ্যা:} \emph{বৃত্তের ব্যাসার্ধ = (ব্যাস/2)=(4/2)=2 সে.মি.}

\emph{কেন্দহ হতে পরিধির দূরত্ব হলো ব্যাসার্ধ, সুতরাং দূরত্ব ২ সে.মি.}
\end{quote}

\textbf{কোন জাতি বৃত্তকে প্রথম ৩৬০ ডিগ্রীতে ভাগ করে?}

ক) ক্যালভীয়ারা (ক্যালডীয়ারা)

\textbf{\hl{খ) আসিরীয়রা ✔}}

গ) গ্রিকরা

ঘ) বাবিলনীয়ারা

\textbf{\hl{✔ সঠিক উত্তর: খ) আসিরীয়রা}}

\begin{quote}
\textbf{ব্যাখ্যা:} \emph{বৃত্তকে প্রথম ৩৬০ ডিগ্রিতে ভাগ করেছিল আসিরীয়রা (বা
প্রাচীন ব্যাবিলনীয়রা/মেসোপটেমীয়রা)।}
\end{quote}

\textbf{(π/10) রেডিয়ানকে ষাটমূলক পদ্ধতিতে ডিগ্রীতে রূপান্তর করলে কত ডিগ্রী
হবে?}

\textbf{\hl{ক) 18° ✔}}

খ) 20°

গ) 30°

ঘ) 35°

\textbf{\hl{✔ সঠিক উত্তর: ক) 18°}}

\begin{quote}
\textbf{ব্যাখ্যা:} \emph{আমরা জানি, π রেডিয়ান = 180\^{}°}

\emph{∴ (π/10) রেডিয়ান = (180\^{}°/π)×(π/10)=18\^{}°}
\end{quote}

\textbf{বৃত্তের সবচেয়ে বড় জ্যা কোনটি?}

\textbf{\hl{ক) ব্যাস ✔}}

খ) ব্যাসার্ধ

গ) জ্যা

ঘ) কোনোটিই নয়

\textbf{\hl{✔ সঠিক উত্তর: ক) ব্যাস}}

\begin{quote}
\textbf{ব্যাখ্যা:} \emph{বৃত্তের কেন্দ্রগামী জ্যা-ই হলো বৃত্তের দীর্ঘতম জ্যা বা
ব্যাস।}
\end{quote}

\textbf{বৃত্তের কেন্দ্র ছেদকারী জ্যাকে কি বলা হয়?}

\textbf{\hl{ক) ব্যাস ✔}}

খ) ব্যাসার্ধ

গ) বৃত্তচাপ

ঘ) পরিধি

\textbf{\hl{✔ সঠিক উত্তর: ক) ব্যাস}}

\begin{quote}
\textbf{ব্যাখ্যা:} \emph{বৃত্তের কেন্দ্রগামী জ্যা হলো তার ব্যাস। ব্যাসই বৃত্তের
বৃহত্তম জ্যা।}
\end{quote}

\textbf{বৃত্তের যে কোনো দুটি বিন্দুর সংযোজক রেখাংশ বৃত্তের একটি-\/-}

ক) ব্যাস (Diameter)

খ) ব্যাসার্ধ (Radius)

গ) চাপ (Arc)

\textbf{\hl{ঘ) জ্যা (Chord) ✔}}

\textbf{\hl{✔ সঠিক উত্তর: ঘ) জ্যা (Chord)}}

\begin{quote}
\textbf{ব্যাখ্যা:} \emph{বৃত্তের পরিধির ওপর অবস্থিত যেকোনো দুটি বিন্দুর সংযোজক
সরলরেখাকে জ্যা (Chord) বলা হয়।}
\end{quote}

\textbf{জ্যা শব্দের অর্থ}

\textbf{\hl{ক) ভূমি ✔}}

খ) সম্পূরক

গ) কোণ

ঘ) সরলরেখা

\textbf{\hl{✔ সঠিক উত্তর: ক) ভূমি}}

\begin{quote}
\textbf{ব্যাখ্যা:} \emph{জ্যা শব্দের অর্থ ভূমি। জ্যামিতি (Geometry) শব্দটিতে
\textquotesingle জ্যা\textquotesingle{} অর্থ ভূমি এবং
\textquotesingle মিতি\textquotesingle{} অর্থ পরিমাপ, যার পূর্ণ অর্থ হলো ভূমির
পরিমাপ।}
\end{quote}

\textbf{বৃত্তের কেন্দ্র থেকে পরিধি পর্যন্ত অঙ্কিত সরল রেখাকে কি বলে?}

ক) ব্যাস

\textbf{\hl{খ) ব্যাসার্ধ ✔}}

গ) জ্যা

ঘ) চাপ

\textbf{\hl{✔ সঠিক উত্তর: খ) ব্যাসার্ধ}}

\begin{quote}
\textbf{ব্যাখ্যা:} \emph{বৃত্তের কেন্দ্র হতে পরিধির ওপর যেকোনো বিন্দুর দূরত্ব বা
রেখাংশকে বৃত্তের ব্যাসার্ধ বলা হয়।}
\end{quote}

\textbf{একই সরলরেখায় অবস্থিত নয় এমন ন্যূনতম কয়টি বিন্দু দিয়ে একটি বৃত্ত আঁকা
যায়?}

ক) ২

\textbf{\hl{খ) ৩ ✔}}

গ) ৪

ঘ) ৫

\textbf{\hl{✔ সঠিক উত্তর: খ) ৩}}

\begin{quote}
\textbf{ব্যাখ্যা:} \emph{একই সরলরেখায় অবস্থিত নয় এমন ন্যূনতম তিনটি বিন্দু দিয়ে
একটি এবং কেবল একটি বৃত্ত আঁকা যায়.}
\end{quote}

\textbf{দুটি নির্দিষ্ট বিন্দু দিয়ে কয়টি বৃত্ত আঁকা যায়?}

\textbf{\hl{ক) অসংখ্য ✔}}

খ) ৩টি

গ) ২টি

ঘ) ১টি

\textbf{\hl{✔ সঠিক উত্তর: ক) অসংখ্য}}

\begin{quote}
\textbf{ব্যাখ্যা:} \emph{দুটি নির্দিষ্ট বিন্দু দিয়ে অসীম সংখ্যক বা অসংখ্য বৃত্ত আঁকা
সম্ভব। এর 📍 মূল কারণ:}

\emph{➔ বিন্দু দুটির সংযোগকারী রেখার একটি লম্ব দ্বিখণ্ডক থাকে।}

\emph{➔ এই লম্ব দ্বিখণ্ডকের ওপর অসংখ্য বিন্দু অবস্থিত।}

\emph{➔ প্রতিটি বিন্দুকে কেন্দ্র করে বৃত্ত আঁকা সম্ভব।}

\emph{➔ ভিন্ন কেন্দ্রের জন্য ভিন্ন ব্যাসার্ধের বৃত্ত তৈরি হয়।}
\end{quote}

\textbf{A ও B কেন্দ্রবিশিষ্ট বৃত্তদ্বয় পরস্পর O বিন্দুতে স্পর্শ করলে AOB এর মান -}

ক) সমকোণ

খ) পূরককোণ

\textbf{\hl{গ) সরলকোণ ✔}}

ঘ) প্রবৃদ্ধকোণ

\textbf{\hl{✔ সঠিক উত্তর: গ) সরলকোণ}}

\begin{quote}
\textbf{ব্যাখ্যা:} \emph{বৃত্তদ্বয় পরস্পরকে বহিস্থভাবে স্পর্শ করলে কেন্দ্রদ্বয় ও
স্পর্শবিন্দু একই সরলরেখায় অবস্থান করে}

\emph{∴ AOB একটি সরলকোণ ( 180\^{}° ).}
\end{quote}

\textbf{চিত্র অনুসারে O কেন্দ্র বিশিষ্ট বৃত্তে ABC অন্তর্লিখিত। y = 112\^{}° হলে
x = কত?}

ক) 68\^{}°

\textbf{\hl{খ) 34\^{}° ✔}}

গ) 45\^{}°

ঘ) 39\^{}°

\textbf{\hl{✔ সঠিক উত্তর: খ) 34\^{}°}}

\begin{quote}
\textbf{ব্যাখ্যা:} \emph{চিত্রানুসারে, একই চাপের ওপর দণ্ডায়মান কেন্দ্রস্থ কোণ
বৃত্তস্থ কোণের দ্বিগুণ অথবা সরলরেখায় অবস্থিত কোণগুলোর সমষ্টি 180\^{}°}

\emph{y হলো প্রবৃদ্ধ কোণ বা রৈখিক যুগল কোণের হিসাব অনুযায়ী, y+2x=180\^{}°
অথবা ত্রিভুজের বহিঃস্থ কোণ ও অন্তুঃস্থ কোণের সম্পর্ক থেকে,}

\emph{∴ x=(180\^{}°−112\^{}°/2)=(68\^{}°/2)=34\^{}° .}
\end{quote}

\textbf{কোন বৃত্তের ব্যাসার্ধ r হলে তার ক্ষেত্রফল কত?}

ক) πr

\textbf{\hl{খ) πr\^{}2 ✔}}

গ) πr\^{}3

ঘ) 2r

\textbf{\hl{✔ সঠিক উত্তর: খ) πr\^{}2}}

\begin{quote}
\textbf{ব্যাখ্যা:} \emph{r ব্যাসার্ধ বিশিষ্ট বৃত্তের ক্ষেত্রফলের সূত্র হলো πr\^{}2
.}
\end{quote}

\textbf{একটি বৃত্তের পরিধি কত? / r ব্যাসার্ধ বিশিষ্ট বৃত্তের পরিধি কোনটি?}

ক) 4π

খ) πr\^{}2

গ) (3/4)πr

\textbf{\hl{ঘ) 2πr ✔}}

\textbf{\hl{✔ সঠিক উত্তর: ঘ) 2πr}}

\begin{quote}
\textbf{ব্যাখ্যা:} \emph{r ব্যাসার্ধ বিশিষ্ট বৃত্তের পরিধি নির্ণয়ের সূত্র হলো 2πr
.}
\end{quote}

\textbf{২৮ সেমি ব্যাসার্ধের বৃত্তের পরিধি কত সেমি?}

ক) ৮৮

\textbf{\hl{খ) ১৭৬ ✔}}

গ) ২৪৯

ঘ) ৪১০

\textbf{\hl{✔ সঠিক উত্তর: খ) ১৭৬}}

\begin{quote}
\textbf{ব্যাখ্যা:} \emph{বৃত্তের পরিধি = 2πr}

\emph{⇒ 2×(22/7)×28=56×(22/7)=176 সেমি.}
\end{quote}

\textbf{কোন বৃত্তের ব্যাস ১৬ সে.মি. হলে পরিধির মান কত?}

ক) 52.16

\textbf{\hl{খ) 50.16 ✔}}

গ) 32.44

ঘ) 48.16

\textbf{\hl{✔ সঠিক উত্তর: খ) 50.16}}

\begin{quote}
\textbf{ব্যাখ্যা:} \emph{বৃত্তের ব্যাস 2r=16 সে.মি., সুতরাং ব্যাসার্ধ r=8
সে.মি.}

\emph{⇒ বৃত্তের পরিধি = 2πr=2×3.1416×8=50.2624 সে.মি..}
\end{quote}

\textbf{বৃত্তের পরিধি ও ব্যাসের অনুপাত কত?}

ক) (25/9)

\textbf{\hl{খ) (22/7) ✔}}

গ) ৩

ঘ) প্রায় ৫

\textbf{\hl{✔ সঠিক উত্তর: খ) (22/7)}}

\begin{quote}
\textbf{ব্যাখ্যা:} \emph{বৃত্তের পরিধি = 2πr এবং বৃত্তের ব্যাস = 2r}

\emph{⇒ অনুপাত = (2πr/2r)=π=(22/7)}
\end{quote}

\textbf{If the difference between the circumference and diameter of a
circle is 60 cm, then the radius of the circle is -}

ক) 7 cm

খ) 9 cm

গ) 10 cm

\textbf{\hl{ঘ) 14 cm ✔}}

\textbf{\hl{✔ সঠিক উত্তর: ঘ) 14 cm}}

\begin{quote}
\textbf{ব্যাখ্যা:} \emph{ধরি, বৃত্তের ব্যাসার্ধ r}

\emph{প্রশ্নমতে, 2πr−2r=60}

\emph{⇒ 2r(π−1)=60}

\emph{⇒ r=(60/2((22/7)−1))=(30/(15/7))=(30×7/15)=14 সেমি}
\end{quote}

\textbf{একটি বৃত্তের ক্ষেত্রফল ১৬ বর্গমিটার এবং পরিধি ৮ মিটার হলে এর ব্যাসার্ধ
কত মিটার?}

ক) ২

খ) ৩

\textbf{\hl{গ) ৮ ✔}}

ঘ) ৫

\textbf{\hl{✔ সঠিক উত্তর: গ) ৮}}

\begin{quote}
\textbf{ব্যাখ্যা:} \emph{বৃত্তের ক্ষেত্রফল A=πr\^{}2=16 এবং পরিধি C=2πr=8}

\emph{⇒ এখান থেকে (A/C)=(πr\^{}2/2πr)=(r/2)=(16/8)=2}

\emph{∴ r=2×2=4 মিটার}

\emph{∴ ব্যাসার্ধ 4 মিটার}
\end{quote}

\textbf{দুটি বৃত্ত পরস্পরকে অন্তঃস্পর্শ করেছে। একটি বৃত্তের ব্যাসার্ধ ৫ সেন্টিমিটার
এবং কেন্দ্রদ্বয়ের দূরত্ব ৯ সেন্টিমিটার হলে, অপর বৃত্তের পরিধি কত সেন্টিমিটার?}

ক) ৩৪

খ) ৪৪

\textbf{\hl{গ) ৮৮ ✔}}

ঘ) ৯০

\textbf{\hl{✔ সঠিক উত্তর: গ) ৮৮}}

\begin{quote}
\textbf{ব্যাখ্যা:} \emph{ধরি, বড় বৃত্তটির ব্যাসার্ধ R এবং ছোট বৃত্তটির ব্যাসার্ধ r
= 5 cm}

\emph{কেন্দ্রদ্বয়ের দূরত্ব = R − r = 9 cm}

\emph{⇒ R − 5 = 9}

\emph{⇒ R = 14 cm}

\emph{∴ অপর বৃত্তের পরিধি = 2πR=2×(22/7)×14=88 সেমি}
\end{quote}

\textbf{একক ব্যাসার্ধবিশিষ্ট বৃত্তের ক্ষেত্রফল কত হবে?}

ক) 1 বর্গ একক

খ) 2 বর্গ একক

\textbf{\hl{গ) π বর্গ একক ✔}}

ঘ) π\^{}2 বর্গ একক

\textbf{\hl{✔ সঠিক উত্তর: গ) π বর্গ একক}}

\begin{quote}
\textbf{ব্যাখ্যা:} \emph{আমরা জানি, ব্যাসার্ধ r হলে বৃত্তের ক্ষেত্রফল = πr\^{}2}

\emph{⇒ ব্যাসার্ধ ১ একক হলে ক্ষেত্রফল = π×1\^{}2=π বর্গ একক}
\end{quote}

\textbf{বৃত্তের ব্যাস তিনগুণ বৃদ্ধি করলে এর ক্ষেত্রফল কতগুণ বৃদ্ধি পাবে?}

ক) 4

\textbf{\hl{খ) 9 ✔}}

গ) 12

ঘ) 16

\textbf{\hl{✔ সঠিক উত্তর: খ) 9}}

\begin{quote}
\textbf{ব্যাখ্যা:} \emph{বৃত্তের ব্যাস তিনগুণ করলে ব্যাসার্ধও তিনগুণ হয়।}

\emph{ক্ষেত্রফল ব্যাসার্ধের বর্গের সমানুপাতিক ( πr\^{}2 ) হওয়ায়, ব্যাসার্ধ ৩ গুণ
হলে ক্ষেত্রফল 3\^{}2 বা ৯ গুণ হবে}
\end{quote}

\textbf{একটি বৃত্তের ব্যাসার্ধ r কে বৃদ্ধি করে ব্যাসার্ধ r + n করলে বৃত্তের
ক্ষেত্রফল 4 গুণ হয়। r এর মান কত?}

ক) (n/2)

\textbf{\hl{খ) n ✔}}

গ) 2n

ঘ) 3n

\textbf{\hl{✔ সঠিক উত্তর: খ) n}}

\begin{quote}
\textbf{ব্যাখ্যা:} \emph{প্রশ্নমতে, (π(r+n)\^{}2/πr\^{}2)=4}

\emph{⇒ ((r+n/r))\^{}2=2\^{}2}

\emph{⇒ (r+n/r)=2}

\emph{⇒ r+n=2r}

\emph{∴ r=n}
\end{quote}

\textbf{ব্যাসার্ধ ২০\% হ্রাস পাওয়াতে ১টি বৃত্তের ক্ষেত্রফল কতটুকু হ্রাস পেল?}

ক) ২১\%

খ) ২২\%

গ) ১৯\%

\textbf{\hl{ঘ) ৩৬\% ✔}}

\textbf{\hl{✔ সঠিক উত্তর: ঘ) ৩৬\%}}

\begin{quote}
\textbf{ব্যাখ্যা:} \emph{ক্ষেত্রফলের শতকরা পরিবর্তন =
(−20−20+((−20)×(−20)/100))\%}

\emph{⇒ (−40+4)\%=−36\%}

\emph{∴ ৩৬\% হ্রাস পাবে।}
\end{quote}

\textbf{একটি বৃত্তের পরিধি ৫০\% বাড়ানো হলে, ক্ষেত্রফল কত বৃদ্ধি পাবে?}

ক) ১০০\%

খ) ১১৫\%

\textbf{\hl{গ) ১২৫\% ✔}}

ঘ) ২২৫\%

\textbf{\hl{✔ সঠিক উত্তর: গ) ১২৫\%}}

\begin{quote}
\textbf{ব্যাখ্যা:} \emph{পরিধি ৫০\% বৃদ্ধি পেলে ব্যাসার্ধও ৫০\% বৃদ্ধি পাবে।}

\emph{ক্ষেত্রফল বৃদ্ধির হার = (A+B+(AB/100))\%}

\emph{⇒ (50+50+(50×50/100))\%=(100+25)\%=125\%}
\end{quote}

\textbf{x\^{}2+y\^{}2=a\^{}2 বৃত্তের পরিসীমা কত হবে?}

ক) (3a\^{}2/4)π

\textbf{\hl{খ) 2πa ✔}}

গ) 2π

ঘ) 6α

\textbf{\hl{✔ সঠিক উত্তর: খ) 2πa}}

\begin{quote}
\textbf{ব্যাখ্যা:} \emph{বৃত্তের সমীকরণ x\^{}2+y\^{}2=a\^{}2 তুলনামূলকভাবে
মূলবিন্দুতে কেন্দ্রবিশিষ্ট বৃত্তের সমীকরণের রূপ, যার ব্যাসার্ধ r=a ।}

\emph{∴ বৃত্তের পরিসীমা বা পরিধি = 2πr=2πa}
\end{quote}

\textbf{একটি বৃত্তের বৃহত্তম জ্যাটি ব্যাসার্ধের---}

ক) সমান

খ) অর্ধেক

\textbf{\hl{গ) দ্বিগুণ ✔}}

ঘ) তিনগুণ

\textbf{\hl{✔ সঠিক উত্তর: গ) দ্বিগুণ}}

\begin{quote}
\textbf{ব্যাখ্যা:} \emph{ব্যাসই বৃত্তের বৃহত্তম জ্যা}

\emph{আবার ব্যাস, ব্যাসার্ধের দ্বিগুণ}

\emph{∴ বৃত্তের বৃহত্তম জ্যাটি ব্যাসার্ধের দ্বিগুণ}
\end{quote}

\textbf{কোনো বৃত্তের তিনটি সমান জ্যা একই বিন্দুতে ছেদ করলে ঐ বিন্দুটির অবস্থান
হবে বৃত্তের/ If three equal chords of a circle intersect at a point, their
point of intersection is -}

ক) পরিধিতে

খ) কেন্দ্র ও পরিধির মধ্যস্থানে

\textbf{\hl{গ) কেন্দ্রে ✔}}

ঘ) ভিন্ন জ্যায়ে

\textbf{\hl{✔ সঠিক উত্তর: গ) কেন্দ্রে}}

\begin{quote}
\textbf{ব্যাখ্যা:} \emph{কোনো বৃত্তের তিনটি সমান জ্যা একই বিন্দুতে ছেদ করলে ঐ
বিন্দুটি বৃত্তের কেন্দ্রে অবস্থিত হবে}
\end{quote}

\textbf{বৃত্তের পরিধি-এর দুটি বিন্দুর ছেদকারী বৃত্তের মধ্যকার রেখাকে বলে-}

ক) ব্যাসার্ধ

খ) ব্যাস

গ) স্পর্শক

\textbf{\hl{ঘ) জ্যা ✔}}

\textbf{\hl{✔ সঠিক উত্তর: ঘ) জ্যা}}

\begin{quote}
\textbf{ব্যাখ্যা:} \emph{বৃত্তের পরিধির ওপর যেকোনো দুটি বিন্দুর সংযোজক
সরলরেখাকে জ্যা বলা হয়।}
\end{quote}

\textbf{একটি জ্যা কয়টি চাপে বিভক্ত?}

ক) ১

\textbf{\hl{খ) ২ ✔}}

গ) ৩

ঘ) ৪

\textbf{\hl{✔ সঠিক উত্তর: খ) ২}}

\begin{quote}
\textbf{ব্যাখ্যা:} \emph{একটি জ্যা বৃত্তকে দুটি চাপে বিভক্ত করে (যেমন: একটি
অধিচাপ ও একটি উপচাপ)}
\end{quote}

\textbf{O কেন্দ্রবিশিষ্ট বৃত্তে AB একটি ব্যাস ভিন্ন জ্যা। OD, AB এর উপর লম্ব। AD
= 2.5 cm হলে, AB = কত cm?}

ক) 4cm

\textbf{\hl{খ) 5cm ✔}}

গ) 6cm

ঘ) 3cm

\textbf{\hl{✔ সঠিক উত্তর: খ) 5cm}}

\begin{quote}
\textbf{ব্যাখ্যা:} \emph{বৃত্তের কেন্দ্র হতে কোনো জ্যা এর ওপর অঙ্কিত লম্ব ঐ
জ্যা-কে সমদ্বিখণ্ডিত করে ।}

\emph{⇒ AB = 2×AD=2×2.5=5 cm}
\end{quote}

\textbf{O কেন্দ্রবিশিষ্ট বৃত্তে OD, AB জ্যা-এর ওপর লম্ব এবং AD = 5 সেমি. হলে AB
= কত সেমি?}

\textbf{\hl{ক) 10 সেমি. ✔}}

খ) 15 সেমি.

গ) 7(1/2) সেমি.

ঘ) 4 সেমি.

\textbf{\hl{✔ সঠিক উত্তর: ক) 10 সেমি.}}

\begin{quote}
\textbf{ব্যাখ্যা:} \emph{বৃত্তের কেন্দ্র থেকে কোনো জ্যা-এর ওপর অঙ্কিত লম্ব জ্যা-কে
সমদ্বিখণ্ডিত করে}

\emph{⇒ AB = 5×2=10 সেমি.}
\end{quote}

\textbf{১৩ সেমি. ব্যাসার্ধের বৃত্তের কেন্দ্র হতে ৫ সেমি. দূরত্বে অবস্থিত জ্যা এর
দৈর্ঘ্য কত?}

ক) ১২ সেমি.

\textbf{\hl{খ) ২৪ সেমি. ✔}}

গ) ১৬ সেমি.

ঘ) ১৮ সেমি.

\textbf{\hl{✔ সঠিক উত্তর: খ) ২৪ সেমি.}}

\begin{quote}
\textbf{ব্যাখ্যা:} \emph{পিথাগোরাসের উপপাদ্য অনুযায়ী, অর্ধজ্যা
PB=√(13\^{}2−5\^{}2)=√(169−25)=√(144)=12 সেমি.}

\emph{⇒ সম্পূর্ণ জ্যা এর দৈর্ঘ্য = 2×12=24 সেমি.}
\end{quote}

\textbf{৫ একক ব্যাসার্ধবিশিষ্ট একটি বৃত্তের কেন্দ্র হতে একটি জ্যা এর লম্বদূরত্ব ৪
একক হলে জ্যা এর দৈর্ঘ্য কত?}

ক) ২ একক

খ) ৩ একক

\textbf{\hl{গ) ৬ একক ✔}}

ঘ) ৮ একক

\textbf{\hl{✔ সঠিক উত্তর: গ) ৬ একক}}

\begin{quote}
\textbf{ব্যাখ্যা:} \emph{পিথাগোরাসের উপপাদ্য অনুযায়ী, অর্ধজ্যা
AD=√(5\^{}2−4\^{}2)=√(25−16)=√(9)=3 একক}

\emph{∴ সম্পূর্ণ জ্যা এর দৈর্ঘ্য = 3×2=6 একক}
\end{quote}

\textbf{কোন বৃত্তের কেন্দ্রগামী জ্যা এর দৈর্ঘ্য 10 সেমি হলে কেন্দ্র হতে 3 সেমি
দূরবর্তী জ্যা এর দৈর্ঘ্য হবে-}

\textbf{\hl{ক) 8 সেমি ✔}}

খ) 4 সেমি

গ) 6 সেমি

ঘ) 7.5 সেমি

\textbf{\hl{✔ সঠিক উত্তর: ক) 8 সেমি}}

\begin{quote}
\textbf{ব্যাখ্যা:} \emph{কেন্দ্রগামী জ্যা এর দৈর্ঘ্য ১০ সেমি হলে ব্যাসার্ধ
r=(10/2)=5 সেমি}

\emph{কেন্দ্র হতে দূরত্ব ৩ সেমি হলে অর্ধজ্যা =
√(5\^{}2−3\^{}2)=√(25−9)=√(16)=4 সেমি}

\emph{∴ জ্যা-এর দৈর্ঘ্য = 4×2=8 সেমি}
\end{quote}

\textbf{১৩ সেন্টিমিটার ব্যাসার্ধবিশিষ্ট কোন বৃত্তের হতে ২৪ সেন্টিমিটার দীর্ঘ জ্যা
এর উপর অঙ্কিত লম্বের দৈর্ঘ্য কত সেন্টিমিটার?}

ক) 8 সেন্টিমিটার.

খ) ৮ সেন্টিমিটার.

\textbf{\hl{গ) 5 সেন্টিমিটার. ✔}}

ঘ) 7 সেন্টিমিটার.

\textbf{\hl{✔ সঠিক উত্তর: গ) 5 সেন্টিমিটার.}}

\begin{quote}
\textbf{ব্যাখ্যা:} \emph{অর্ধজ্যা = (24/2)=12 সেমি}

\emph{পিথাগোরাসের উপপাদ্য অনুযায়ী, লম্বের দূরত্ব =
√(13\^{}2−12\^{}2)=√(169−144)=√(25)=5 সেন্টিমিটার.}
\end{quote}

\textbf{কোন বৃত্তের 10 সেন্টিমিটার দীর্ঘ একটি জ্যা কেন্দ্র হতে 12 সেন্টিমিটার
দূরে অবস্থিত। বৃত্তটির ব্যাসার্ধ কত সেন্টিমিটার?}

\textbf{\hl{ক) 13 ✔}}

খ) 14.5

গ) 9

ঘ) 15

\textbf{\hl{✔ সঠিক উত্তর: ক) 13}}

\begin{quote}
\textbf{ব্যাখ্যা:} \emph{অর্ধজ্যা BC=(10/2)=5 সেমি}

\emph{পিতায়োগরাসের উপপাদ্য অনুযায়ী, ব্যাসার্ধ
OB=√(12\^{}2+5\^{}2)=√(144+25)=√(169)=13 সেন্টিমিটার}
\end{quote}

\textbf{বৃত্তের ব্যাস ভিন্ন জ্যা এর দৈর্ঘ্য 8√(2) সেন্টিমিটার এবং কেন্দ্র হতে ঐ
জ্যা এর উপর লম্বের দৈর্ঘ্য ২ সেন্টিমিটার হলে বৃত্তের বৃহত্তম জ্যা এর দৈর্ঘ্য কত
সেন্টিমিটার?}

ক) ১০

\textbf{\hl{খ) ১২ ✔}}

গ) ১৬

ঘ) ২০

\textbf{\hl{✔ সঠিক উত্তর: খ) ১২}}

\begin{quote}
\textbf{ব্যাখ্যা:} \emph{অর্ধজ্যা = (8√(2)/2)=4√(2) সেমি}

\emph{ব্যাসার্ধ AC=√(2\^{}2+(4√(2))\^{}2)=√(4+32)=√(36)=6 সেমি}

\emph{∴ বৃহত্তম জ্যা (ব্যাস) = 6×2=12 সেন্টিমিটার}
\end{quote}

\textbf{O কেন্দ্রবিশিষ্ট বৃত্তে OD, AB জ্যা এর উপর লম্ব। AD = 3 সে.মি. হলে AB =
কত সে.মি.?}

ক) 3 সেমি.

খ) 8 সেমি.

গ) 5 সেমি.

\textbf{\hl{ঘ) 6 সেমি. ✔}}

\textbf{\hl{✔ সঠিক উত্তর: ঘ) 6 সেমি.}}

\begin{quote}
\textbf{ব্যাখ্যা:} \emph{বৃত্তের কেন্দ্র হতে জ্যা-এর ওপর অঙ্কিত লম্ব জ্যা-কে
সমদ্বিখণ্ডিত করে.}

\emph{⇒ AB = 2×AD=2×3=6 সে.মি.}
\end{quote}

\textbf{ADB বৃত্তে AB এবং CD দুটি সমান জ্যা পরস্পর P বিন্দুতে ছেদ করলে কোনটি
সত্য?}

ক) PC=PD

খ) PB=PC

গ) PB=PA

\textbf{\hl{ঘ) PB=PD ✔}}

\textbf{\hl{✔ সঠিক উত্তর: ঘ) PB=PD}}

\begin{quote}
\textbf{ব্যাখ্যা:} \emph{বৃত্তের দুটি সমান জ্যা পরস্পরকে ছেদ করলে একটির অংশদ্বয়
অপরটির অংশদ্বয়ের সমান হয়, অর্থাৎ PB = PD এবং PA = PC}
\end{quote}

\textbf{10cm ব্যাসার্ধ বিশিষ্ট একটি বৃত্তের কেন্দ্র হতে একটি জ্যা এর উপর অঙ্কিত
লম্বের দৈর্ঘ্য অর্ধ-জ্যা অপেক্ষা 2cm কম। জ্যা এর দৈর্ঘ্য কত?}

ক) 18cm

খ) 4cm

গ) 8cm

\textbf{\hl{ঘ) 16cm ✔}}

\textbf{\hl{✔ সঠিক উত্তর: ঘ) 16cm}}

\begin{quote}
\textbf{ব্যাখ্যা:} \emph{ধরি, অর্ধজ্যা = x সেমি, তবে লম্ব = x−2 সেমি।}

\emph{পিথাগোরাস অনুসারে, x\^{}2+(x−2)2=10\^{}2}

\emph{⇒ x\^{}2+x\^{}2−4x+4=100}

\emph{⇒ 2x\^{}2−4x−96=0}

\emph{⇒ x\^{}2−2x−48=0}

\emph{⇒ (x−8)(x+6)=0}

\emph{⇒ x=8 (যেহেতু দৈর্ঘ্য ঋণাত্মক হতে পারে না)}

\emph{∴ সম্পূর্ণ জ্যা এর দৈর্ঘ্য = 2×8=16 সেমি}
\end{quote}

\textbf{কোন বৃত্তের ১২ মিটার দীর্ঘ একটি জ্যা কেন্দ্র থেকে ৮ মিটার দূরে অবস্থিত
হলে বৃত্তটির ব্যাস কত মিটার হবে?}

ক) ১০

খ) ১৫

গ) ১৮

\textbf{\hl{ঘ) ২০ ✔}}

\textbf{\hl{✔ সঠিক উত্তর: ঘ) ২০}}

\begin{quote}
\textbf{ব্যাখ্যা:} \emph{অর্ধজ্যা AC=(12/2)=6 মিটার}

\emph{পিথাগোরাসের উপপাদ্য অনুযায়ী, ব্যাসার্ধ
OA=√(6\^{}2+8\^{}2)=√(36+64)=√(100)=10 মিটার}

\emph{∴ ব্যাস = 2×10=20 মিটার}
\end{quote}

\textbf{নিচের চিত্রে O কেন্দ্রবিশিষ্ট বৃত্তে ∠BAC=41° হলে ∠BOC কোণের মান কত?}

ক) 49°

খ) 62°

গ) 139°

\textbf{\hl{ঘ) 82° ✔}}

\textbf{\hl{✔ সঠিক উত্তর: ঘ) 82°}}

\begin{quote}
\textbf{ব্যাখ্যা:} \emph{বৃত্তের একই চাপের ওপর দণ্ডায়মান কেন্দ্রস্থ কোণ বৃত্তস্থ
কোণের দ্বিগুণ।}

\emph{⇒ ∠BOC=2×∠BAC=2×41°=82°}
\end{quote}

\textbf{বৃত্তের কেন্দ্রস্থ কোণ পরিধিস্থ কোণের---}

ক) সমান

খ) তিনগুণ

গ) অর্ধেক

\textbf{\hl{ঘ) দ্বিগুণ ✔}}

\textbf{\hl{✔ সঠিক উত্তর: ঘ) দ্বিগুণ}}

\begin{quote}
\textbf{ব্যাখ্যা:} \emph{বৃত্তের একই চাপের ওপর দণ্ডায়মান কেন্দ্রস্থ কোণ পরিধিস্থ
(বৃত্তস্থ) কোণের দ্বিগুণ।}
\end{quote}

\textbf{অর্ধবৃত্তস্থ কোণের পরিমাণ কত ডিগ্রি?}

ক) ৪৫°

খ) ৬০°

\textbf{\hl{গ) ৯০° ✔}}

ঘ) ১০০°

\textbf{\hl{✔ সঠিক উত্তর: গ) ৯০°}}

\begin{quote}
\textbf{ব্যাখ্যা:} \emph{অর্ধবৃত্তস্থ কোণ এক সমকোণ, অর্থাৎ এর পরিমাণ ৯০°।}
\end{quote}

\textbf{বৃত্তের একই চাপের ওপর দণ্ডায়মান কেন্দ্রস্থ কোণ 60° হলে, বৃত্তস্থ কোণ কত?}

ক) 60°

\textbf{\hl{খ) 30° ✔}}

গ) 120°

ঘ) 180°

\textbf{\hl{✔ সঠিক উত্তর: খ) 30°}}

\begin{quote}
\textbf{ব্যাখ্যা:} \emph{বৃত্তস্থ কোণ কেন্দ্রস্থ কোণের অর্ধেক।}

\emph{⇒ বৃত্তস্থ কোণ = (60°/2)=30°}
\end{quote}

\textbf{কোনো বৃত্তের অধিচাপে অন্তর্লিখিত কোণ}

\textbf{\hl{ক) সূক্ষ্ম কোণ ✔}}

খ) স্থূল কোণ

গ) সমকোণ

ঘ) পূরক কোণ

\textbf{\hl{✔ সঠিক উত্তর: ক) সূক্ষ্ম কোণ}}

\begin{quote}
\textbf{ব্যাখ্যা:} \emph{কোনো বৃত্তের অধিচাপে অন্তর্লিখিত কোণ সর্বদা সূক্ষ্ম কোণ
হয়, এবং উপচাপে অন্তর্লিখিত কোণ স্থূল কোণ হয়।}
\end{quote}

\textbf{একটি বৃত্তের একটি চাপের উপর অঙ্কিত কেন্দ্রস্থ কোণ বৃত্তস্থ কোণের কত অংশ?}

ক) অর্ধেক

\textbf{\hl{খ) দ্বিগুণ ✔}}

গ) সমান

ঘ) তিনগুণ

\textbf{\hl{✔ সঠিক উত্তর: খ) দ্বিগুণ}}

\begin{quote}
\textbf{ব্যাখ্যা:} \emph{বৃত্তের একই চাপের ওপর দণ্ডায়মান কেন্দ্রস্থ কোণ বৃত্তস্থ
কোণের দ্বিগুণ, অর্থাৎ বৃত্তস্থ কোণ কেন্দ্রস্থ কোণের অর্ধেক।}
\end{quote}

\textbf{একটি বৃত্তের কেন্দ্রস্থ কোণ, পরিধিথ কোণের -}

ক) অর্ধেক (Half)

খ) সমান (Equal)

\textbf{\hl{গ) দ্বিগুণ (Twice) ✔}}

ঘ) তিনগুণ (Thrice)

\textbf{\hl{✔ সঠিক উত্তর: গ) দ্বিগুণ (Twice)}}

\begin{quote}
\textbf{ব্যাখ্যা:} \emph{বৃত্তের কেন্দ্রস্থ কোণ সর্বদা পরিধিথ (বৃত্তস্থ) কোণের
দ্বিগুণ হয়।}
\end{quote}

\textbf{একই চাপের উপর দণ্ডায়মান কেন্দ্রস্থ কোণের পরিমাণ 110° হলে পরিধিথ
কোণের পরিমাণ হবে-}

ক) 50°

\textbf{\hl{খ) 55° ✔}}

গ) 60°

ঘ) 70°

\textbf{\hl{✔ সঠিক উত্তর: খ) 55°}}

\begin{quote}
\textbf{ব্যাখ্যা:} \emph{পরিধিথ কোণ কেন্দ্রস্থ কোণের অর্ধেক।}

\emph{⇒ পরিধিথ কোণ = (110°/2)=55°}
\end{quote}

\textbf{O কেন্দ্রবিশিষ্ট বৃত্তে BC চাপের উপর ∠BAC হলো পরিধিথ কোণ। ∠BOC=110°
হলে ∠BAC= কত?}

ক) 30°

খ) 44°

\textbf{\hl{গ) 55° ✔}}

ঘ) 110°

\textbf{\hl{✔ সঠিক উত্তর: গ) 55°}}

\begin{quote}
\textbf{ব্যাখ্যা:} \emph{বৃত্তস্থ বা পরিধিথ কোণ কেন্দ্রস্থ কোণের অর্ধেক।}

\emph{⇒ ∠BAC=(1/2)×110°=55°}
\end{quote}

\textbf{চিত্রে O কেন্দ্র বিশিষ্ট বৃত্তে ABC সমবাহু। ∠AOB এর মান কত?}

ক) 90°

খ) 180°

\textbf{\hl{গ) 120° ✔}}

ঘ) 240°

\textbf{\hl{✔ সঠিক উত্তর: গ) 120°}}

\begin{quote}
\textbf{ব্যাখ্যা:} \emph{ABC সমবাহু ত্রিভুজ হওয়ায় এর প্রতিটি কোণের মান ৬০°,
অর্থাৎ ∠ACB=60° ।}

\emph{বৃত্তের কেন্দ্রস্থ কোণ একই চাপের উপর দণ্ডায়মান বৃত্তস্থ কোণের দ্বিগুণ।}

\emph{⇒ ∠AOB=2×∠ACB=2×60°=120°}
\end{quote}

\textbf{বৃত্তের উপরস্থ কোনো বিন্দুতে কয়টি স্পর্শক আঁকা যায়?}

\textbf{\hl{ক) একটি ✔}}

খ) চারটি

গ) দুইটি

ঘ) তিনটি

\textbf{\hl{✔ সঠিক উত্তর: ক) একটি}}

\begin{quote}
\textbf{ব্যাখ্যা:} \emph{বৃত্তের ওপরস্থ যেকোনো নির্দিষ্ট বিন্দুতে কেবল একটিমাত্র
স্পর্শক আঁকা সম্ভব।}
\end{quote}

\textbf{স্পর্শ বিন্দুগামী ব্যাসার্ধ এবং স্পর্শকের অন্তর্ভুক্ত কোণ---}

ক) এক সমকোণের অর্ধেক

খ) দুই সমকোণ

\textbf{\hl{গ) এক সমকোণ ✔}}

ঘ) সূক্ষ্মকোণ

\textbf{\hl{✔ সঠিক উত্তর: গ) এক সমকোণ}}

\begin{quote}
\textbf{ব্যাখ্যা:} \emph{বৃত্তের যে কোনো বিন্দুতে অঙ্কিত স্পর্শক স্পর্শবিন্দুগামী
ব্যাসার্ধের ওপর লম্ব।}

\emph{⇒ তাই স্পর্শ বিন্দুগামী ব্যাসার্ধ এবং স্পর্শকের অন্তর্ভুক্ত কোণ এক সমকোণ।}
\end{quote}

\textbf{দুটি বৃত্ত যদি পরস্পর স্পর্শ করে তবে কেন্দ্রদ্বয় হতে স্পর্শ বিন্দুগামী
সরলরেখা দুটির অন্তর্ভুক্ত কোণ---}

ক) সূক্ষ্মকোণ

খ) স্থূলকোণ

গ) সমকোণ

\textbf{\hl{ঘ) সরলকোণ ✔}}

\textbf{\hl{✔ সঠিক উত্তর: ঘ) সরলকোণ}}

\begin{quote}
\textbf{ব্যাখ্যা:} \emph{দুটি বৃত্ত পরস্পরকে স্পর্শ করলে তাদের কেন্দ্রদ্বয় এবং
স্পর্শবিন্দু সমরেখ হয়।}

\emph{⇒ তাই কেন্দ্রদ্বয় হতে স্পর্শ বিন্দুগামী সরলরেখা দুটির অন্তর্ভুক্ত কোণ একটি
সরলকোণ (১৮০°)।}
\end{quote}

\textbf{x সেমি. এবং y সেমি. পরিধির দুটি বৃত্ত পরস্পরকে স্পর্শ করে। বৃত্ত দুটির
কেন্দ্রের দূরত্ব কত?}

ক) 2πx+2πy সেমি.

\textbf{\hl{খ) (x+y/2π) সেমি. ✔}}

গ) πx+πy সেমি.

ঘ) (π/x+y) সেমি.

\textbf{\hl{✔ সঠিক উত্তর: খ) (x+y/2π) সেমি.}}

\begin{quote}
\textbf{ব্যাখ্যা:} \emph{x সেমি. পরিধির বৃত্তের ব্যাসার্ধ (x/2π) এবং y সেমি.
পরিধির বৃত্তের ব্যাসার্ধ (y/2π) সেমি。}

\emph{⇒ বৃত্ত দুটির কেন্দ্রের দূরত্ব = ব্যাসার্ধদ্বয়ের সমষ্টি =
(x/2π)+(y/2π)=(x+y/2π) সেমি.}
\end{quote}

\textbf{দুটি বৃত্ত পরস্পরকে অন্তঃস্পর্শ করল। বৃহত্তম বৃত্তটির ব্যাসার্ধ 6 সেমি এবং
কেন্দ্রদ্বয়ের দূরত্ব 2 সেমি। অপর বৃত্তের ব্যাসার্ধ কত?}

ক) 2 সেমি

\textbf{\hl{খ) 4 সেমি ✔}}

গ) 6 সেমি

ঘ) 8 সেমি

\textbf{\hl{✔ সঠিক উত্তর: খ) 4 সেমি}}

\begin{quote}
\textbf{ব্যাখ্যা:} \emph{ধরি, অপর (অন্তঃস্থ) বৃত্তটির ব্যাসার্ধ x সেমি।}

\emph{অন্তঃস্পর্শের ক্ষেত্রে কেন্দ্রদ্বয়ের দূরত্ব = (বড় বৃত্তের ব্যাসার্ধ − ছোট বৃত্তের
ব্যাসার্ধ)}

\emph{⇒ 6−x=2}

\emph{⇒ x=4 সেমি.}
\end{quote}

\textbf{A ও B কেন্দ্রবিশিষ্ট দুটি বৃত্ত পরস্পর C বিন্দুতে বহিঃস্থভাবে স্পর্শ করে।
∠ACB এর মান -}

ক) 120°

\textbf{\hl{খ) 180° ✔}}

গ) 270°

ঘ) 360°

\textbf{\hl{✔ সঠিক উত্তর: খ) 180°}}

\begin{quote}
\textbf{ব্যাখ্যা:} \emph{দুটি বৃত্ত পরস্পর বহিঃস্পর্শ করলে তাদের কেন্দ্রদ্বয় এবং
স্পর্শবিন্দু একই সরলরেখায় (সমরেখ) অবস্থান করে।}

\emph{⇒ তাই ∠ACB একটি সরলকোণ, যার মান 180°।}
\end{quote}

\textbf{2 সেমি. ব্যাসার্ধবিশিষ্ট বৃত্তের কেন্দ্র হতে 5 সেমি. দূরের কোনো বিন্দু হতে
অঙ্কিত স্পর্শকের দৈর্ঘ্য কত সেমি.?}

\textbf{\hl{ক) 4.58 ✔}}

খ) 5.38

গ) 3

ঘ) 5

\textbf{\hl{✔ সঠিক উত্তর: ক) 4.58}}

\begin{quote}
\textbf{ব্যাখ্যা:} \emph{স্পর্শক ব্যাসার্ধের উপর লম্ব হওয়ায় সমকোণী ত্রিভুজ উৎপন্ন
হয়।}

\emph{পিথাগোরাসের উপপাদ্য অনুযায়ী, স্পর্শকের দৈর্ঘ্য =
√(5\^{}2−2\^{}2)=√(25−4)=√(21)≈4.58 সেমি.}
\end{quote}

\textbf{নিচের কোনটি বৃত্তস্থ চতুর্ভুজ হয় না?}

ক) বর্গ

খ) আয়ত

গ) ট্রাপিজিয়াম

\textbf{\hl{ঘ) রম্বস ✔}}

\textbf{\hl{✔ সঠিক উত্তর: ঘ) রম্বস}}

\begin{quote}
\textbf{ব্যাখ্যা:} \emph{বৃত্তস্থ চতুর্ভুজ হওয়ার শর্ত হলো কোনো চতুর্ভুজের বিপরীত
কোণদ্বয়ের সমষ্টি ১৮০° হতে হবে।}

\emph{⇒ রম্বসের বিপরীত কোণের সমষ্টি সর্বদা ১৮০° হয় না (যদি না সেটি বর্গ হয়),
তাই রম্বস বৃত্তস্থ চতুর্ভুজ নয়।}
\end{quote}

\textbf{বৃত্তস্থ চতুর্ভুজের বিপরীত কোণদ্বয়ের সমষ্টি-}

ক) ৯০°

খ) ১২০°

\textbf{\hl{গ) ১৮০° ✔}}

ঘ) ৩৬০°

\textbf{\hl{✔ সঠিক উত্তর: গ) ১৮০°}}

\begin{quote}
\textbf{ব্যাখ্যা:} \emph{বৃত্তস্থ চতুর্ভুজের বিপরীত কোণদ্বয়ের সমষ্টি দুই সমকোণ বা
১৮০°।}
\end{quote}

\textbf{একটি চতুর্ভুজের ৩ কোণের সমষ্টি ২৮০ ডিগ্রি। ৪র্থ কোণের মান কত?}

ক) ৯০ ডিগ্রি

খ) ৭০ ডিগ্রি

\textbf{\hl{গ) ৮০ ডিগ্রি ✔}}

ঘ) ৬০ ডিগ্রি

\textbf{\hl{✔ সঠিক উত্তর: গ) ৮০ ডিগ্রি}}

\begin{quote}
\textbf{ব্যাখ্যা:} \emph{চতুর্ভুজের চার কোণের সমষ্টি ৩৬০ ডিগ্রি।}

\emph{⇒ ৪র্থ কোণের মান = 360°−280°=80°}
\end{quote}

\textbf{বৃত্তস্থ চতুর্ভুজের একটি কোণ ৯০° হলে উহার বিপরীত কোণের পরিমাপ কত হবে?}

ক) ২০°

\textbf{\hl{খ) ৯০° ✔}}

গ) ১১০°

ঘ) ১২০°

\textbf{\hl{✔ সঠিক উত্তর: খ) ৯০°}}

\begin{quote}
\textbf{ব্যাখ্যা:} \emph{বৃত্তস্থ চতুর্ভুজের বিপরীত কোণদ্বয়ের সমষ্টি ১৮০°।}

\emph{⇒ বিপরীত কোণের পরিমাপ = 180°−90°=90°}
\end{quote}

\textbf{একটি ত্রিভুজ এবং একটি বৃত্ত ন্যূনতম কয়টি বিন্দুতে ছেদ করে?}

ক) ৪টি

খ) ১টি

গ) ৩টি

\textbf{\hl{ঘ) ২টি ✔}}

\textbf{\hl{✔ সঠিক উত্তর: ঘ) ২টি}}

\begin{quote}
\textbf{ব্যাখ্যা:} \emph{একটি ত্রিভুজ ও বৃত্ত সর্বনিম্ন দুটি বিন্দুতে ছেদ করতে
পারে। অন্যদিকে একটি ত্রিভুজ ও বৃত্ত সর্বোচ্চ ছয়টি বিন্দুতে ছেদ করতে পারে।}
\end{quote}

\textbf{ABCD একটি বৃত্তস্থ চতুর্ভুজ। ∠BAD ও ∠BCD -এর সমষ্টি কত হবে?}

ক) 120°

\textbf{\hl{খ) 180° ✔}}

গ) 90°

ঘ) কোনো নির্দিষ্ট মান নেই

\textbf{\hl{✔ সঠিক উত্তর: খ) 180°}}

\begin{quote}
\textbf{ব্যাখ্যা:} \emph{আমরা জানি, বৃত্তস্থ চতুর্ভুজের বিপরীত কোণদ্বয়ের সমষ্টি
১৮০°।}

\emph{⇒ ABCD বৃত্তস্থ চতুর্ভুজে ∠BAD ও ∠BCD পরস্পর বিপরীত কোণ হওয়ায় এদের সমষ্টি
১৮০° হবে।}
\end{quote}

\textbf{৫ সেমিঃ ব্যাসার্ধ বিশিষ্ট বৃত্তের অন্তর্লিখিত বর্গক্ষেত্রের ক্ষেত্রফল কত?}

ক) ৬০ বর্গ সে.মি

খ) ৭০ বর্গ সে.মি

গ) ৪০ বর্গ সে.মি

\textbf{\hl{ঘ) ৫০ বর্গ সে.মি ✔}}

\textbf{\hl{✔ সঠিক উত্তর: ঘ) ৫০ বর্গ সে.মি}}

\begin{quote}
\textbf{ব্যাখ্যা:} \emph{বৃত্তের অন্তর্লিখিত বর্গক্ষেত্রের কর্ণ = বৃত্তের ব্যাস =
5×2=10 সে.মি.}

\emph{আমরা জানি, বর্গক্ষেত্রের ক্ষেত্রফল = (কর্ণ\^{}2/2)}

\emph{⇒ ক্ষেত্রফল = (10\^{}2/2)=(100/2)=50 বর্গ সে.মি.}
\end{quote}

\section{68. ত্রিকোণমিতি MCQ প্রশ্ন উত্তর
অনুশীলন}\label{ux9a4ux9b0ux995ux9a3ux9aeux9a4-mcq-ux9aaux9b0ux9b6ux9a8-ux989ux9a4ux9a4ux9b0-ux985ux9a8ux9b6ux9b2ux9a8}

\textbf{cos( (nπ/2) ) অনুক্রমটির চতুর্থ পদ কোনটি?}

ক) -1

\textbf{\hl{খ) 1 ✔}}

গ) 2

ঘ) 0

\textbf{\hl{✔ সঠিক উত্তর: খ) 1}}

\begin{quote}
\textbf{ব্যাখ্যা:} \emph{ধারাটির চতুর্থ পদ = cos( (4π/2) ) {[}∴ n = 4{]}}

\emph{⇒ = cos 2π}

\emph{⇒ = cos (2 × 180°) {[}∴ π = 180°{]}}

\emph{⇒ = cos 360° = 1}
\end{quote}

\textbf{tan A = (3/4) হলে, sin A এর মান কত?}

ক) (1/4)

\textbf{\hl{খ) (3/5) ✔}}

গ) (2/5)

ঘ) (4/5)

\textbf{\hl{✔ সঠিক উত্তর: খ) (3/5)}}

\begin{quote}
\textbf{ব্যাখ্যা:} \emph{আমরা জানি, tan A = (লম্ব/ভূমি) = (3/4)}

\emph{অতিভুজ = √(3\^{}2+4\^{}2) = √(9+16) = √(25) = 5}

\emph{∴ sin A = (লম্ব/অতিভুজ) = (3/5)}
\end{quote}

\textbf{SinA = (12/13) হলে tanA এর মান কত?}

ক) (5/12)

\textbf{\hl{খ) (12/5) ✔}}

গ) (12/7)

ঘ) (13/12)

\textbf{\hl{✔ সঠিক উত্তর: খ) (12/5)}}

\begin{quote}
\textbf{ব্যাখ্যা:} \emph{SinA = (লম্ব/অতিভুজ) = (12/13)}

\emph{ভূমি = √(13\^{}2−12\^{}2) = √(169−144) = √(25) = 5}

\emph{∴ tanA = (লম্ব/ভূমি) = (12/5)}
\end{quote}

\textbf{sinθ = (4/5) হলে, tanθ = কত?}

\textbf{\hl{ক) (4/3) ✔}}

খ) (3/4)

গ) (3/5)

ঘ) (5/4)

\textbf{\hl{✔ সঠিক উত্তর: ক) (4/3)}}

\begin{quote}
\textbf{ব্যাখ্যা:} \emph{cosθ = √(1−sin\^{}2θ) = √(1−(((4/5)))\^{}2) =
√(1−(16/25)) = √((9/25)) = (3/5)}

\emph{∴ tanθ = (sinθ/cosθ) = ((4/5)/(3/5)) = (4/5)×(5/3) = (4/3)}
\end{quote}

\textbf{secθ = (3/2) হলে tanθ = ?}

ক) (2/3)

\textbf{\hl{খ) (√(5)/2) ✔}}

গ) (5/4)

ঘ) 1

\textbf{\hl{✔ সঠিক উত্তর: খ) (√(5)/2)}}

\begin{quote}
\textbf{ব্যাখ্যা:} \emph{আমরা জানি, sec²θ − tan²θ = 1}

\emph{⇒ tan²θ = sec²θ − 1}

\emph{⇒ tan²θ = (((3/2)))\^{}2 − 1}

\emph{⇒ tan²θ = (9/4) − 1}

\emph{⇒ tan²θ = (9−4/4) = (5/4)}

\emph{∴ tanθ = (√(5)/2)}
\end{quote}

\textbf{sinθ = (5/13) হলে, tanθ এর মান কত?}

ক) (3/4)

\textbf{\hl{খ) (5/12) ✔}}

গ) (12/13)

ঘ) (5/13)

\textbf{\hl{✔ সঠিক উত্তর: খ) (5/12)}}

\begin{quote}
\textbf{ব্যাখ্যা:} \emph{cosθ = √(1−sin\^{}2θ)}

\emph{= √(1−(((5/13)))\^{}2)}

\emph{= √(1−(25/169))}

\emph{= √((169−25/169))}

\emph{= √((144/169))}

\emph{= (12/13)}

\emph{∴ tanθ = (sinθ/cosθ)}

\emph{= ((5/13)/(12/13))}

\emph{= (5/12)}
\end{quote}

\textbf{cotθ √(1−cos\^{}2θ) = ?}

ক) sin θ

\textbf{\hl{খ) cos θ ✔}}

গ) tan θ

ঘ) cot θ

\textbf{\hl{✔ সঠিক উত্তর: খ) cos θ}}

\begin{quote}
\textbf{ব্যাখ্যা:} \emph{cotθ √(1−cos\^{}2θ)}

\emph{= cotθ . sinθ}

\emph{= (cosθ/sinθ) . sinθ}

\emph{= cosθ}
\end{quote}

\textbf{tan 90° এর মান কত?}

\textbf{\hl{ক) ∞ ✔}}

খ) -∞

গ) 1

ঘ) -1

\textbf{\hl{✔ সঠিক উত্তর: ক) ∞}}

\begin{quote}
\textbf{ব্যাখ্যা:} \emph{আমরা জানি,}

\emph{tan 90° = (sin90°/cos90°)}

\emph{= (1/0)}

\emph{= ∞}
\end{quote}

\textbf{sinθ = (5/13) হলে cosecθ এর মান কত?}

ক) (5/12)

খ) (12/5)

\textbf{\hl{গ) (13/5) ✔}}

ঘ) (5/13)

\textbf{\hl{✔ সঠিক উত্তর: গ) (13/5)}}

\begin{quote}
\textbf{ব্যাখ্যা:} \emph{আমরা জানি, cosecθ = (1/sinθ)}

\emph{∴ cosecθ = (1/(5/13)) = (13/5)}
\end{quote}

\textbf{sin(−θ) কত?}

\textbf{\hl{ক) -sin θ ✔}}

খ) sin θ

গ) ±sin θ

ঘ) কোনটিই নয়

\textbf{\hl{✔ সঠিক উত্তর: ক) -sin θ}}

\begin{quote}
\textbf{ব্যাখ্যা:} \emph{ত্রিকোণমিতিক অনুপাত অনুযায়ী, sin(−θ) = −sinθ (যেহেতু
চতুর্থ চতুর্ভাগে সাইন ঋণাত্মক হয়)}
\end{quote}

\textbf{নিচের কোনটি সঠিক?}

\textbf{\hl{ক) sec²θ = 1 + tan²θ ✔}}

খ) sec²θ = 1 − tan²θ

গ) tan²θ = (cosθ/sinθ)

ঘ) sin²θ − cos²θ = 1

\textbf{\hl{✔ সঠিক উত্তর: ক) sec²θ = 1 + tan²θ}}

\begin{quote}
\textbf{ব্যাখ্যা:} \emph{আমরা জানি, অভেদাবলী অনুসারে sec²θ − tan²θ = 1}

\emph{∴ sec²θ = 1 + tan²θ}
\end{quote}

\textbf{sinθ এর সর্বনিম্ন মান কত?}

ক) 0

খ) -2

গ) 1

\textbf{\hl{ঘ) -1 ✔}}

\textbf{\hl{✔ সঠিক উত্তর: ঘ) -1}}

\begin{quote}
\textbf{ব্যাখ্যা:} \emph{sinθ ও cosθ এর সর্বোচ্চ ও সর্বনিম্ন মান যথাক্রমে +1 ও
−1}

\emph{∴ sinθ এর সর্বনিম্ন মান −1}
\end{quote}

\textbf{cosec(90° − θ) = 2 হলে, cosθ = কত?}

ক) 2

খ) (√(3)/2)

\textbf{\hl{গ) (1/2) ✔}}

ঘ) (1/√(2))

\textbf{\hl{✔ সঠিক উত্তর: গ) (1/2)}}

\begin{quote}
\textbf{ব্যাখ্যা:} \emph{cosec(90° − θ) = 2}

\emph{⇒ secθ = 2}

\emph{⇒ (1/cosθ) = 2}

\emph{∴ cosθ = (1/2)}
\end{quote}

\textbf{ত্রিকোণমিতিক অনুপাতে কোন সম্পর্কটি ব্যবহার করা হয়?}

ক) sinθ = 1 + cos²θ

\textbf{\hl{খ) sin²θ = 1 − cos²θ ✔}}

গ) sin²θ − cos²θ = 1

ঘ) sinθ = (1/cosθ)

\textbf{\hl{✔ সঠিক উত্তর: খ) sin²θ = 1 − cos²θ}}

\begin{quote}
\textbf{ব্যাখ্যা:} \emph{আমরা জানি, sin²θ + cos²θ = 1}

\emph{∴ sin²θ = 1 − cos²θ}
\end{quote}

\textbf{sin30° এর মান কত?}

ক) (√(3)/2)

\textbf{\hl{খ) (1/2) ✔}}

গ) (1/√(2))

ঘ) (1/√(3))

\textbf{\hl{✔ সঠিক উত্তর: খ) (1/2)}}

\begin{quote}
\textbf{ব্যাখ্যা:} \emph{ত্রিকোণমিতিক মান অনুযায়ী, sin30° এর মান (1/2)}
\end{quote}

\textbf{sin²21° + sin²69° = কত?}

ক) -1

\textbf{\hl{খ) 1 ✔}}

গ) (1/2)

ঘ) (1/√(2))

\textbf{\hl{✔ সঠিক উত্তর: খ) 1}}

\begin{quote}
\textbf{ব্যাখ্যা:} \emph{sin²21° + sin²69°}

\emph{⇒ = sin²21° + sin²(90° − 21°)}

\emph{⇒ = sin²21° + cos²21° = 1}
\end{quote}

\textbf{tanA = 1 হলে, A-এর মান কত ডিগ্রি?}

\textbf{\hl{ক) 45° ✔}}

খ) 30°

গ) 60°

ঘ) 70°

\textbf{\hl{✔ সঠিক উত্তর: ক) 45°}}

\begin{quote}
\textbf{ব্যাখ্যা:} \emph{tan A = 1}

\emph{⇒ tan A = tan 45°}

\emph{∴ A = 45°}
\end{quote}

\textbf{tanθ = b হলে sin 2θ = ?}

ক) (1+b\^{}2/1−b\^{}2)

খ) (1−b\^{}2/1+b\^{}2)

\textbf{\hl{গ) (2b/1+b\^{}2) ✔}}

ঘ) (2b/1−b\^{}2)

\textbf{\hl{✔ সঠিক উত্তর: গ) (2b/1+b\^{}2)}}

\begin{quote}
\textbf{ব্যাখ্যা:} \emph{আমরা জানি, sin 2θ = (2tanθ/1+tan\^{}2θ) =
(2b/1+b\^{}2)}
\end{quote}

\textbf{tanθ = (3/4) হলে, cosecθ =}

ক) (3/5)

\textbf{\hl{খ) (5/3) ✔}}

গ) (4/3)

ঘ) (5/4)

\textbf{\hl{✔ সঠিক উত্তর: খ) (5/3)}}

\begin{quote}
\textbf{ব্যাখ্যা:} \emph{tanθ = (লম্ব/ভূমি) = (3/4)}

\emph{অতিভুজ = √(3\^{}2+4\^{}2) = √(9+16) = √(25) = 5}

\emph{∴ cosecθ = (অতিভুজ/লম্ব) = (5/3)}
\end{quote}

\textbf{নিচের কোনটি সঠিক?}

ক) tanθ = (1/cosθ)

খ) tanθ = cosθ

\textbf{\hl{গ) tanθ = (1/cotθ) ✔}}

ঘ) tanθ = (cosθ/sinθ)

\textbf{\hl{✔ সঠিক উত্তর: গ) tanθ = (1/cotθ)}}

\begin{quote}
\textbf{ব্যাখ্যা:} \emph{আমরা জানি, tanθ এবং cotθ পরস্পর বিপরীতার্থক অনুপাত।}

\emph{∴ tanθ = (1/cotθ)}
\end{quote}

\textbf{Cosec(90° − θ) = (2/√(3)) হলে tanθ = কত?}

ক) (2/3)

খ) (1/2)

গ) (5/4)

\textbf{\hl{ঘ) (1/√(3)) ✔}}

\textbf{\hl{✔ সঠিক উত্তর: ঘ) (1/√(3))}}

\begin{quote}
\textbf{ব্যাখ্যা:} \emph{দেয়া আছে, Cosec(90° − θ) = (2/√(3))}

\emph{⇒ secθ = (2/√(3)) {[}∴ Cosec(90° − θ) = secθ{]}}

\emph{⇒ (1/cosθ) = (2/√(3))}

\emph{⇒ cosθ = (√(3)/2) = cos 30° ⇒ θ = 30°}

\emph{∴ tan 30° = (1/√(3))}
\end{quote}

\textbf{A = 30° হলে (2tanA/tan\^{}2A) = ?}

ক) 2

খ) (2/√(3))

গ) 4

\textbf{\hl{ঘ) 2√(3) ✔}}

\textbf{\hl{✔ সঠিক উত্তর: ঘ) 2√(3)}}

\begin{quote}
\textbf{ব্যাখ্যা:} \emph{(2tanA/tan\^{}2A)}

\emph{= (2tan30°/(tan30°)\^{}2)}

\emph{= (2×(1/√(3))/(((1/√(3))))\^{}2)}

\emph{= ((2/√(3))/(1/3))}

\emph{= (2/√(3))×(3/1) = 2√(3)}
\end{quote}

\textbf{2 + tan\^{}2θ = 5 হলে θ এর মান কত?}

ক) 90°

খ) 0°

গ) 30°

\textbf{\hl{ঘ) 60° ✔}}

\textbf{\hl{✔ সঠিক উত্তর: ঘ) 60°}}

\begin{quote}
\textbf{ব্যাখ্যা:} \emph{দেয়া আছে, 2 + tan\^{}2θ = 5}

\emph{⇒ tan\^{}2θ = 5 − 2}

\emph{⇒ tan\^{}2θ = 3}

\emph{⇒ tanθ = √(3)}

\emph{⇒ tanθ = tan 60°}

\emph{∴ θ = 60°}
\end{quote}

\textbf{A = 45° হলে, (1−tan\^{}2A/1+tan\^{}2A) = কত?}

ক) 1

খ) (1/2)

\textbf{\hl{গ) 0 ✔}}

ঘ) 2

\textbf{\hl{✔ সঠিক উত্তর: গ) 0}}

\begin{quote}
\textbf{ব্যাখ্যা:} \emph{(1−tan\^{}245°/1+tan\^{}245°)}

\emph{= (1−1\^{}2/1+1\^{}2)}

\emph{= (1−1/1+1)}

\emph{= (0/2) = 0}
\end{quote}

\textbf{যদি A = 45° হয়, তবে cos 2A-এর মান কত?}

ক) 1

খ) (1/√(2))

\textbf{\hl{গ) 0 ✔}}

ঘ) (1/2)

\textbf{\hl{✔ সঠিক উত্তর: গ) 0}}

\begin{quote}
\textbf{ব্যাখ্যা:} \emph{cos 2A}

\emph{= cos(2 × 45°)}

\emph{= cos 90° = 0}
\end{quote}

\textbf{মান নির্ণয় করুন: tan (π/12) tan (5π/12) tan (7π/12) tan (11π/12)}

ক) −(1/2)

খ) (1/2)

গ) (√(3)/2)

\textbf{\hl{ঘ) 1 ✔}}

\textbf{\hl{✔ সঠিক উত্তর: ঘ) 1}}

\begin{quote}
\textbf{ব্যাখ্যা:} \emph{tan (π/12) tan (5π/12) tan (7π/12) tan (11π/12)}

\emph{= tan 15° tan 75° tan 105° tan 165°}

\emph{= (sin15°/cos15°) · (sin75°/cos75°) · (sin105°/cos105°) ·
(sin165°/cos165°)}

\emph{= (cos60°−cos90°/cos60°+cos90°) · (cos60°−cos270°/cos60°+cos270°)}

\emph{= ((1/2)−0/(1/2)+0) · ((1/2)−0/(1/2)+0) = 1}
\end{quote}

\textbf{sec θ + tan θ = (5/2) হলে, sec θ − tan θ এর মান কত?}

ক) (9/4)

\textbf{\hl{খ) (2/5) ✔}}

গ) 3

ঘ) (25/4)

\textbf{\hl{✔ সঠিক উত্তর: খ) (2/5)}}

\begin{quote}
\textbf{ব্যাখ্যা:} \emph{আমরা জানি, sec²θ − tan²θ = 1}

\emph{⇒ (secθ + tanθ)(secθ − tanθ) = 1}

\emph{⇒ (5/2) × (secθ − tanθ) = 1}

\emph{∴ (secθ − tanθ) = (2/5)}
\end{quote}

\textbf{(sinθ+cosθ/sinθ−cosθ) = 7 হলে secθ এর মান কত?}

ক) (5/3)

\textbf{\hl{খ) ±(5/3) ✔}}

গ) −(5/3)

ঘ) (3/5)

\textbf{\hl{✔ সঠিক উত্তর: খ) ±(5/3)}}

\begin{quote}
\textbf{ব্যাখ্যা:} \emph{যোজন-বিয়োজন করে পাই,}

\emph{(sinθ+cosθ+sinθ−cosθ/sinθ+cosθ−(sinθ−cosθ)) = (7+1/7−1)}

\emph{⇒ (2sinθ/2cosθ) = (8/6)}

\emph{⇒ tanθ = (4/3)}

\emph{∴ secθ = ±√(1+tan\^{}2θ) = ±√(1+(16/9)) = ±(5/3)}
\end{quote}

\textbf{Δ ABC এর ক্ষেত্রফল = Δ হলে, sin C = কত?}

ক) (2Δ/bc)

খ) (2Δ/abc)

\textbf{\hl{গ) (2Δ/ab) ✔}}

ঘ) (2Δ/abc)

\textbf{\hl{✔ সঠিক উত্তর: গ) (2Δ/ab)}}

\begin{quote}
\textbf{ব্যাখ্যা:} \emph{আমরা জানি, ত্রিভুজের ক্ষেত্রফল Δ=(1/2)absinC}

\emph{⇒ sin C = (2Δ/ab)}
\end{quote}

\textbf{যদি y=tan\^{}−1x হয়, তবে (dy/dx) = কত?}

ক) e\^{}tan\^{}−1x

খ) e\^{}tan\^{}−1x(1+x\^{}2)

গ) (etan\^{}−1x/1+x\^{}2)

\textbf{\hl{ঘ) (1/1+x\^{}2) ✔}}

\textbf{\hl{✔ সঠিক উত্তর: ঘ) (1/1+x\^{}2)}}

\begin{quote}
\textbf{ব্যাখ্যা:} \emph{আমরা জানি, অন্তরীকরণের সূত্র অনুযায়ী
(d/dx)(tan\^{}−1x)=(1/1+x\^{}2)}
\end{quote}

\textbf{(2tanθ/1+tan\^{}2θ) = কত?}

ক) tan 2θ

খ) 2cos\^{}2(θ/2)

\textbf{\hl{গ) 2 sin θ cos θ ✔}}

ঘ) cos 2θ

\textbf{\hl{✔ সঠিক উত্তর: গ) 2 sin θ cos θ}}

\begin{quote}
\textbf{ব্যাখ্যা:} \emph{ত্রিকোণমিতিক সূত্রানুযায়ী,
(2tanθ/1+tan\^{}2θ)=sin2θ=2sinθcosθ}
\end{quote}

\textbf{tan\^{}−11+tan\^{}−12+tan\^{}−13 = কত?}

\textbf{\hl{ক) π ✔}}

খ) (π/2)

গ) (3π/2)

ঘ) 2π

\textbf{\hl{✔ সঠিক উত্তর: ক) π}}

\begin{quote}
\textbf{ব্যাখ্যা:} \emph{ধরি, A=tan\^{}−11=(π/4) .}

\emph{এখন tan\^{}−12+tan\^{}−13}

\emph{=π+tan\^{}−1((2+3/1−2×3))}

\emph{=π+tan\^{}−1(−1)}

\emph{=π−(π/4)}

\emph{=(3π/4)}

\emph{∴ যোগফল = (π/4)+(3π/4)=π}
\end{quote}

\textbf{যদি A সূক্ষ্মকোণ এবং sinA=(12/13) হয়, তবে cot A এর মান কত?}

ক) (5/13)

\textbf{\hl{খ) (5/12) ✔}}

গ) (10/3)

ঘ) (10/13)

\textbf{\hl{✔ সঠিক উত্তর: খ) (5/12)}}

\begin{quote}
\textbf{ব্যাখ্যা:} \emph{sin A = (লম্ব/অতিভুজ)=(12/13)}

\emph{ভূমি = √(13\^{}2−12\^{}2)}

\emph{=√(169−144) =√(25)=5}

\emph{∴ cot A = (ভূমি/লম্ব)=(5/12)}
\end{quote}

\textbf{নিচের কোন সম্পর্কটি সঠিক?}

ক) sin\^{}2A+cos\^{}2B=1

\textbf{\hl{খ) sinθ+cosθ\textgreater1 ✔}}

গ) tan\^{}2θ−sec\^{}2θ=1

ঘ) sinθ+cosθ\textless1

\textbf{\hl{✔ সঠিক উত্তর: খ) sinθ+cosθ\textgreater1}}

\begin{quote}
\textbf{ব্যাখ্যা:} \emph{যেহেতু
(sinθ+cosθ)\^{}2=sin\^{}2θ+cos\^{}2θ+2sinθcosθ=1+sin2θ , এবং sin 2θ-এর
মান ধনাত্মক হলে এর মান 1 এর বেশি হতে পারে, তাই sinθ+cosθ\textgreater1
সম্পর্কটি প্রযোজ্য।}
\end{quote}

\textbf{secθ = 2 হলে, θ এর মান কত?}

ক) 0°

খ) 30°

গ) 45°

\textbf{\hl{ঘ) 60° ✔}}

\textbf{\hl{✔ সঠিক উত্তর: ঘ) 60°}}

\begin{quote}
\textbf{ব্যাখ্যা:} \emph{secθ = 2}

\emph{⇒ secθ = sec 60°}

\emph{∴ θ = 60°}
\end{quote}

\textbf{tan²45° sin 60° tan 30° tan²60° এর মান কত?}

ক) (4/3)

খ) (3/4)

\textbf{\hl{গ) (3/2) ✔}}

ঘ) (1/2)

\textbf{\hl{✔ সঠিক উত্তর: গ) (3/2)}}

\begin{quote}
\textbf{ব্যাখ্যা:} \emph{tan²45° sin 60° tan 30° tan²60°}

\emph{= 1\^{}2×(√(3)/2)×(1/√(3))×(√(3))\^{}2}

\emph{= 1×(1/2)×3=(3/2)}
\end{quote}

\textbf{sinθ = (3/5) , 0 \textless{} θ \textless{} (π/2) হলে, tanθ = ?}

\textbf{\hl{ক) (3/4) ✔}}

খ) (4/5)

গ) (5/4)

ঘ) (4/3)

\textbf{\hl{✔ সঠিক উত্তর: ক) (3/4)}}

\begin{quote}
\textbf{ব্যাখ্যা:} \emph{sinθ = (লম্ব/অতিভুজ)=(3/5)}

\emph{ভূমি = √(5\^{}2−3\^{}2)=√(25−9)=√(16)=4}

\emph{∴ tanθ = (লম্ব/ভূমি)=(3/4)}
\end{quote}

\textbf{x = sin y হলে x এর সর্বোচ্চ মান কত?}

ক) α

\textbf{\hl{খ) 1 ✔}}

গ) -1

ঘ) 0

\textbf{\hl{✔ সঠিক উত্তর: খ) 1}}

\begin{quote}
\textbf{ব্যাখ্যা:} \emph{sin y = x}

\emph{∴ x = 1 {[}sinθ -এর সর্বোচ্চ মান +1{]}}
\end{quote}

\textbf{কোনটি সঠিক?}

ক) sin 120° = sin 60° + sin 60°

\textbf{\hl{খ) sin 120° = 2 sin 60° cos 60° ✔}}

গ) sin 120° = 2 sin 60° 60°

ঘ) sin 120° = −(1/2)

\textbf{\hl{✔ সঠিক উত্তর: খ) sin 120° = 2 sin 60° cos 60°}}

\begin{quote}
\textbf{ব্যাখ্যা:} \emph{sin 120° = sin(2 × 60°)}

\emph{= 2 sin 60° cos 60° {[}sin 2A = 2 sinA cosA{]}}
\end{quote}

\textbf{sin A = (3/2) হলে cos A = কত?}

ক) (3/2)

খ) (1/3)

\textbf{\hl{গ) (√(5)/3) ✔}}

ঘ) (1/2)

\textbf{\hl{✔ সঠিক উত্তর: গ) (√(5)/3)}}

\begin{quote}
\textbf{ব্যাখ্যা:} \emph{sin A = (লম্ব/অতিভুজ)=(3/2)}

\emph{ভূমি = √(3\^{}2−2\^{}2)=√(9−4)=√(5)}

\emph{∴ cos A = (ভূমি/অতিভুজ)=(√(5)/3)}
\end{quote}

\textbf{cosθ এর সর্বনিম্ন মান কত?}

ক) 0

\textbf{\hl{খ) −1 ✔}}

গ) 1

ঘ) −∞

\textbf{\hl{✔ সঠিক উত্তর: খ) −1}}

\begin{quote}
\textbf{ব্যাখ্যা:} \emph{cosθ এর সর্বনিম্ন মান −1 এবং cosθ এর সর্বোচ্চ মান 1।}
\end{quote}

\textbf{(1+cosp/sinp) এর মান কোনটি?}

ক) tan(p/2)

\textbf{\hl{খ) cot(p/2) ✔}}

গ) cos(p/2)

ঘ) sin(p/2)

\textbf{\hl{✔ সঠিক উত্তর: খ) cot(p/2)}}

\begin{quote}
\textbf{ব্যাখ্যা:}
\emph{(1+cosp/sinp)=(2cos\^{}2(p/2)/2sin(p/2)cos(p/2))=(cos(p/2)/sin(p/2))=cot(p/2)}
\end{quote}

\textbf{θ = 60° হলে sec²θ − tan²θ = ?}

ক) (1/4)

খ) (1/2)

গ) 0

\textbf{\hl{ঘ) 1 ✔}}

\textbf{\hl{✔ সঠিক উত্তর: ঘ) 1}}

\begin{quote}
\textbf{ব্যাখ্যা:} \emph{আমরা জানি, অভেদ অনুযায়ী sec²θ − tan²θ = 1}

\emph{এখানে, θ = 60° মানটি অতিরিক্ত বা অনাবশ্যক।}
\end{quote}

\textbf{cos 9B = sin B হলে cot 5B এর মান কত?}

ক) 3

খ) 10

গ) 0

\textbf{\hl{ঘ) 1 ✔}}

\textbf{\hl{✔ সঠিক উত্তর: ঘ) 1}}

\begin{quote}
\textbf{ব্যাখ্যা:} \emph{দেয়া আছে, cos 9B = sin B}

\emph{⇒ cos 9B = cos( (π/2) − B)}

\emph{⇒ 9B = (π/2) − B}

\emph{⇒ 10B = (π/2)}

\emph{⇒ B = (π/20)}

\emph{∴ cot 5B = cot(5 × (π/20) ) = cot( (π/4) ) = 1}
\end{quote}

\textbf{cosec(90° − θ) = 2 হলে cosθ = ?}

\textbf{\hl{ক) (1/2) ✔}}

খ) (√(3)/2)

গ) (1/√(2))

ঘ) 2

\textbf{\hl{✔ সঠিক উত্তর: ক) (1/2)}}

\begin{quote}
\textbf{ব্যাখ্যা:} \emph{cosec(90° − θ) = 2}

\emph{⇒ secθ = 2}

\emph{⇒ (1/cosθ) = 2}

\emph{⇒ 2 cosθ = 1}

\emph{∴ cosθ = (1/2)}
\end{quote}

\textbf{যদি tanθ = 1 হয়, তবে sinθ − cos(−θ) এর মান কত?}

ক) −√(2)

\textbf{\hl{খ) 0 ✔}}

গ) √(2)

ঘ) 1

\textbf{\hl{✔ সঠিক উত্তর: খ) 0}}

\begin{quote}
\textbf{ব্যাখ্যা:} \emph{tanθ = 1 = (1/1) {[} লম্ব=1,ভূমি=1 {]}}

\emph{অতিভুজ = √(1\^{}2+1\^{}2)=√(2)}

\emph{∴ sinθ=(1/√(2)) এবং cosθ=(1/√(2))}

\emph{sinθ − cos(−θ) = sinθ − cosθ {[}cos(−θ) = cosθ{]}}

\emph{= (1/√(2))−(1/√(2))=0}
\end{quote}

\textbf{tanθ = (4/3) এবং (3π/2) \textless{} θ \textless{} 2π হলে sinθ =
কত?}

\textbf{\hl{ক) −(4/5) ✔}}

খ) (4/5)

গ) −(3/5)

ঘ) (3/5)

\textbf{\hl{✔ সঠিক উত্তর: ক) −(4/5)}}

\begin{quote}
\textbf{ব্যাখ্যা:} \emph{tanθ = (4/3)}

\emph{অতিভুজ = √(4\^{}2+3\^{}2)=√(25)=5}

\emph{যেহেতু কোণটি চতুর্থ চতুর্ভাগে ( (3π/2) থেকে 2π) অবস্থিত এবং চতুর্থ চতুর্ভাগে
সাইন ঋণাত্মক, তাই sinθ=−(লম্ব/অতিভুজ)=−(4/5)}
\end{quote}

\textbf{(sinθ+cosθ/sinθ−cosθ) = 7 এর মান কত?}

ক) (5/3)

\textbf{\hl{খ) ±(5/3) ✔}}

গ) −(5/3)

ঘ) (3/5)

\textbf{\hl{✔ সঠিক উত্তর: খ) ±(5/3)}}

\begin{quote}
\textbf{ব্যাখ্যা:} \emph{দেয়া আছে, (sinθ+cosθ/sinθ−cosθ)=7}

\emph{যোজন-বিয়োজন করে পাই,
(sinθ+cosθ+sinθ−cosθ/sinθ+cosθ−(sinθ−cosθ))=(7+1/7−1)}

\emph{⇒ (2sinθ/2cosθ)=(8/6)}

\emph{⇒ tanθ=(4/3)}

\emph{∴ secθ=±√(1+tan\^{}2θ)=±√(1+(16/9))=±(5/3)}
\end{quote}

\textbf{cos\^{}−1x+cos\^{}−1y=(π/2) হলে x² + y² = ?}

\textbf{\hl{ক) 1 ✔}}

খ) 0

গ) (√(3)/2)

ঘ) (1/2)

\textbf{\hl{✔ সঠিক উত্তর: ক) 1}}

\begin{quote}
\textbf{ব্যাখ্যা:} \emph{দেয়া আছে, cos\^{}−1x+cos\^{}−1y=(π/2)}

\emph{⇒ cos\^{}−1x=(π/2)−cos\^{}−1y}

\emph{⇒ x=cos((π/2)−cos\^{}−1y)}

\emph{⇒ x=sin(cos\^{}−1y)}

\emph{⇒ x=sin(sin\^{}−1√(1−y\^{}2))}

\emph{⇒ x=√(1−y\^{}2)}

\emph{উভয়পক্ষকে বর্গ করে পাই, x\^{}2=1−y\^{}2}

\emph{∴ x² + y² = 1}
\end{quote}

\textbf{tan\^{}−1(1/2)+tan\^{}−1(1/3)}

\textbf{\hl{ক) (π/4) ✔}}

খ) (π/2)

গ) (π/3)

ঘ) π

\textbf{\hl{✔ সঠিক উত্তর: ক) (π/4)}}

\begin{quote}
\textbf{ব্যাখ্যা:} \emph{আমরা জানি,
tan\^{}−1A+tan\^{}−1B=tan\^{}−1(A+B/1−A×B)}

\emph{⇒
tan\^{}−1(1/2)+tan\^{}−1(1/3)=tan\^{}−1((1/2)+(1/3)/1−(1/2)×(1/3))=tan\^{}−1((5/6)/(5/6))=tan\^{}−11=(π/4)}
\end{quote}

\textbf{3 cot A = 4 হলে sin A এর মান কত?}

ক) (4/5)

\textbf{\hl{খ) (3/5) ✔}}

গ) (3/4)

ঘ) (4/3)

\textbf{\hl{✔ সঠিক উত্তর: খ) (3/5)}}

\begin{quote}
\textbf{ব্যাখ্যা:} \emph{দেয়া আছে, 3 cot A = 4}

\emph{⇒ cot A = (4/3)}

\emph{ভূমি = 4 এবং লম্ব = 3 হলে অতিভুজ = √(4\^{}2+3\^{}2)=√(25)=5}

\emph{∴ sin A = (লম্ব/অতিভুজ)=(3/5)}
\end{quote}

\textbf{sin((9π/2)+θ) = ?}

ক) sin θ

খ) -sin θ

\textbf{\hl{গ) cos θ ✔}}

ঘ) -cos θ

\textbf{\hl{✔ সঠিক উত্তর: গ) cos θ}}

\begin{quote}
\textbf{ব্যাখ্যা:} \emph{sin((9π/2)+θ)=sin(9×(π/2)+θ)=cosθ}

\emph{উল্লেখ্য, এখানে n = 9 এবং n এর মান বিজোড় হলে sin পরিবর্তিত হয়ে cos
হবে।}
\end{quote}

\textbf{tanθ=(a/b) হলে নিচের কোনটি সত্য?}

ক) sinθ=(b/√(a\^{}2+b\^{}2))

খ) cosθ=(a/√(a\^{}2+b\^{}2))

\textbf{\hl{গ) cosecθ=(√(a\^{}2+b\^{}2)/a) ✔}}

ঘ) sinθ=(√(a\^{}2+b\^{}2)/a)

\textbf{\hl{✔ সঠিক উত্তর: গ) cosecθ=(√(a\^{}2+b\^{}2)/a)}}

\begin{quote}
\textbf{ব্যাখ্যা:} \emph{দেয়া আছে, tanθ=(লম্ব/ভূমি)=(a/b)}

\emph{অতিভুজ = √(a\^{}2+b\^{}2)}

\emph{∴ cosecθ=(অতিভুজ/লম্ব)=(√(a\^{}2+b\^{}2)/a)}
\end{quote}

\textbf{(secθ + tanθ) = (7/5) হলে (secθ − tanθ) এর মান কত?}

\textbf{\hl{ক) (5/7) ✔}}

খ) (3/7)

গ) (3/5)

ঘ) (1/5)

\textbf{\hl{✔ সঠিক উত্তর: ক) (5/7)}}

\begin{quote}
\textbf{ব্যাখ্যা:} \emph{দেয়া আছে, secθ + tanθ = (7/5)}

\emph{আমরা জানি, sec²θ − tan²θ = 1}

\emph{⇒ (secθ + tanθ)(secθ − tanθ) = 1}

\emph{⇒ (7/5) × (secθ − tanθ) = 1}

\emph{∴ secθ − tanθ = (5/7)}
\end{quote}

\textbf{1 + tan²θ = 4 হলে θ = কত?}

ক) 0°

খ) 30°

গ) 45°

\textbf{\hl{ঘ) 60° ✔}}

\textbf{\hl{✔ সঠিক উত্তর: ঘ) 60°}}

\begin{quote}
\textbf{ব্যাখ্যা:} \emph{1 + tan²θ = 4}

\emph{⇒ sec²θ = 4}

\emph{⇒ secθ = 2}

\emph{⇒ (1/cosθ) = 2}

\emph{⇒ cosθ = (1/2)}

\emph{∴ θ = 60°}
\end{quote}

\textbf{ΔABC এ ∠B = 90°, AC = 2AB হয় তবে, ∠C এর মান কত?}

ক) 45°

খ) 22.5°

গ) 30°

\textbf{\hl{ঘ) 60° ✔}}

\textbf{\hl{✔ সঠিক উত্তর: ঘ) 60°}}

\begin{quote}
\textbf{ব্যাখ্যা:} \emph{ত্রিকোণমিতির সূত্র অনুযায়ী, ΔABC এ sin∠C =
(লম্ব/অতিভুজ)=(AB/AC)}

\emph{⇒ sin∠C = (AB/2AB)=(1/2)}

\emph{⇒ sin∠C = sin 30°}

\emph{∴ ∠C = 30°}

\emph{অতএব, ∠A = 90° − 30° = 60°}
\end{quote}

\textbf{দালানের ছাদের উন্নতি কোণ 60°, ঐ স্থান থেকে 25 মিটার পিছিয়ে গেলে
উন্নতি কোণ হয় 30°। দালানের উচ্চতা নির্ণয় করুন।}

\textbf{\hl{ক) ২১.৬৫ মিটার ✔}}

খ) ২১ মিটার

গ) ২১.৫৬ মিটার

ঘ) ২৫ মিটার

\textbf{\hl{✔ সঠিক উত্তর: ক) ২১.৬৫ মিটার}}

\begin{quote}
\textbf{ব্যাখ্যা:} \emph{ধরি, দালানের উচ্চতা AB = h এবং BC = x}

\emph{প্রথম ক্ষেত্রে, tan 60° = (AB/BC)}

\emph{⇒ √(3)=(h/x)}

\emph{⇒ x = (h/√(3))}

\emph{দ্বিতীয় ক্ষেত্রে, tan 30° = (AB/BD)=(h/x+25)}

\emph{⇒ (1/√(3))=(h/x+25)}

\emph{⇒ x + 25 = √(3)h}

\emph{⇒ (h/√(3)) + 25 = √(3)h}

\emph{⇒ √(3)h−(h/√(3)) = 25}

\emph{⇒ (3h−h/√(3)) = 25}

\emph{⇒ 2h = 25√(3)}

\emph{∴ h = (25√(3)/2) = 21.65 মিটার}
\end{quote}

\textbf{একটি গাছের পাদদেশ থেকে কিছু দূরে একটি স্থানে গাছটির শীর্ষের উন্নতি
কোণ ৩০°। গাছটি ১২ মিটার উঁচু হলে ঐ স্থানটি গাছ হতে কত দূরে অবস্থিত?}

\textbf{\hl{ক) 12 √(3) মিটার ✔}}

খ) 10 √(3) মিটার

গ) 20 মিটার

ঘ) 30 মিটার

\textbf{\hl{✔ সঠিক উত্তর: ক) 12 √(3) মিটার}}

\begin{quote}
\textbf{ব্যাখ্যা:} \emph{ধরি, গাছটি x মিটার দূরে অবস্থিত}

\emph{∴ tan 30° = (AB/AC)}

\emph{⇒ (1/√(3))=(12/x)}

\emph{∴ x = 12 √(3)}
\end{quote}

\textbf{একটি তাল গাছের পাদবিন্দু হতে 10 মিটার দূরবর্তী স্থান থেকে গাছের
শীর্ষের উন্নতি কোণ 60° হলে গাছের উচ্চতা নির্ণয় করুন।}

ক) 16.65 মি.

খ) 17.72 মি.

\textbf{\hl{গ) 17.32 মি. ✔}}

ঘ) 17.75 মি.

\textbf{\hl{✔ সঠিক উত্তর: গ) 17.32 মি.}}

\begin{quote}
\textbf{ব্যাখ্যা:} \emph{ধরি, গাছটির উচ্চতা = x মিটার}

\emph{tan 60° = (লম্ব/ভূমি)}

\emph{⇒ √(3)=(x/10)}

\emph{∴ x = 10 × √(3) = 17.32 মি.}
\end{quote}

\textbf{সূর্যের উন্নতি কোণ 60° হলে একটি গাছের ছায়ার দৈর্ঘ্য 10 মিটার হয়।
গাছটির উচ্চতা কত?}

ক) 16.65 মি.

খ) 17.72 মি.

\textbf{\hl{গ) 17.32 মি. ✔}}

ঘ) 17.75 মি.

\textbf{\hl{✔ সঠিক উত্তর: গ) 17.32 মি.}}

\begin{quote}
\textbf{ব্যাখ্যা:} \emph{ধরি, গাছটির উচ্চতা = x মিটার}

\emph{tan 60° = (লম্ব/ভূমি)}

\emph{⇒ √(3)=(x/10)}

\emph{∴ x = 10 × √(3) = 17.32 মি.}
\end{quote}

\textbf{একটি ৫০ মিটার লম্বা মই একটি খাড়া দেয়ালের সাথে হেলান দিয়ে রাখা
হয়েছে। মইয়ের একপ্রান্ত মাটি হতে ৪০ মিটার উঁচু দেয়ালকে স্পর্শ করে। মইয়ের অপর
প্রান্ত হতে দেয়ালের দূরত্ব?}

ক) ১০ মিটার

\textbf{\hl{খ) ৩০ মিটার ✔}}

গ) ২০ মিটার

ঘ) ২৫ মিটার

\textbf{\hl{✔ সঠিক উত্তর: খ) ৩০ মিটার}}

\begin{quote}
\textbf{ব্যাখ্যা:} \emph{পিথাগোরাসের উপপাদ্য অনুযায়ী, (অতিভুজ)² = (লম্ব)² +
(ভূমি)²}

\emph{⇒ (৫০)² = (৪০)² + (ভূমি)²}

\emph{⇒ ২৫০০ = ১৬০০ + (ভূমি)²}

\emph{⇒ (ভূমি)² = ২৫০০ − ১৬০০ = ৯০০}

\emph{∴ ভূমি = √(900) = ৩০ মিটার}
\end{quote}

\textbf{একটি মই এর এক প্রান্ত ভূমি থেকে ১৫ মিটার উঁচু ঘরের জানালা বরাবর
পৌঁছায়। অপর প্রান্ত ঘর থেকে ৮ মিটার দূরে থাকলে মই এর দৈর্ঘ্য কত মিটার?}

\textbf{\hl{ক) ১৭ মিটার ✔}}

খ) ১৮ মিটার

গ) ১৯ মিটার

ঘ) ২০ মিটার

ঙ) কোনটি নয়

\textbf{\hl{✔ সঠিক উত্তর: ক) ১৭ মিটার}}

\begin{quote}
\textbf{ব্যাখ্যা:} \emph{ধরি, মইটি x মিটার লম্বা}

\emph{সমকোণী ত্রিভুজের সূত্র হতে পাই, (অতিভুজ)² = (ভূমি)² + (লম্ব)²}

\emph{⇒ x² = 8² + 15²}

\emph{⇒ x² = 64 + 225}

\emph{⇒ x² = 289}

\emph{∴ x = √(289) = 17}
\end{quote}

\textbf{একটি খুঁটি ভেঙ্গে গিয়ে ভূমির সাথে 30° কোণ উৎপন্ন করে। ভাঙ্গা অংশের
দৈর্ঘ্য 16 মিটার হলে দণ্ডায়মান অংশের দৈর্ঘ্য কত মিটার?}

\textbf{\hl{ক) 8 মিটার ✔}}

খ) 8 √(3) মিটার

গ) 16 √(3) মিটার

ঘ) 16 মিটার

\textbf{\hl{✔ সঠিক উত্তর: ক) 8 মিটার}}

\begin{quote}
\textbf{ব্যাখ্যা:} \emph{sinθ = (লম্ব/অতিভুজ)}

\emph{⇒ sin 30° = (x/16)}

\emph{⇒ (1/2)=(x/16)}

\emph{∴ x = 8}
\end{quote}

\textbf{18 মিটার দীর্ঘ একটি মই ভূমির সাথে 30° কোণ উৎপন্ন করে দেওয়ালের ছাদ
স্পর্শ করে। দেওয়ালটির উচ্চতা কত?}

ক) 6 মিটার

\textbf{\hl{খ) 9 মিটার ✔}}

গ) 12 মিটার

ঘ) 6 √(3) মিটার

\textbf{\hl{✔ সঠিক উত্তর: খ) 9 মিটার}}

\begin{quote}
\textbf{ব্যাখ্যা:} \emph{মনে করি, সমকোণী ত্রিভুজের দেওয়ালের উচ্চতা h মিটার}

\emph{sinθ = (লম্ব/অতিভুজ)}

\emph{⇒ sin 30° = (h/18)}

\emph{⇒ (1/2)=(h/18)}

\emph{∴ h = 9 মিটার}
\end{quote}

\textbf{একটি ঘুড়ি ভূমি থেকে ৫৫ মিটার উপরে উড়ছে, যার সুতা ভূমির সাথে 60° কোণ
উৎপন্ন করে। সুতার দৈর্ঘ্য কত মিটার?}

\textbf{\hl{ক) 63.5 ✔}}

খ) 110

গ) (100/2)

ঘ) (127/2)

\textbf{\hl{✔ সঠিক উত্তর: ক) 63.5}}

\begin{quote}
\textbf{ব্যাখ্যা:} \emph{sin 60° = (ভূমি হতে উচ্চতা/সুতার দৈর্ঘ্য)}

\emph{⇒ (√(3)/2)=(55/সুতার দৈর্ঘ্য)}

\emph{∴ সুতার দৈর্ঘ্য = (55×2/√(3)) = 63.5 মিটার}
\end{quote}

\textbf{একটি খাড়া খুঁটি মাটি থেকে ৩ মিটার ওপরে ভেঙে বিচ্ছিন্ন না হয়ে
অন্যপ্রান্ত ভূমিতে ৪ মিটার দূরত্বে স্পর্শ করলে খুঁটির উচ্চতা কত?}

ক) ৫ মিটার

\textbf{\hl{খ) ৮ মিটার ✔}}

গ) ৭ মিটার

ঘ) ৯ মিটার

\textbf{\hl{✔ সঠিক উত্তর: খ) ৮ মিটার}}

\begin{quote}
\textbf{ব্যাখ্যা:} \emph{ধরি, খুঁটির উচ্চতা h মিটার}

\emph{ABC সমকোণী ত্রিভুজ হতে পাই, (h−3)\^{}2=4\^{}2+3\^{}2}

\emph{⇒ (h−3)\^{}2=25}

\emph{⇒ h − 3 = 5}

\emph{⇒ h = 5 + 3}

\emph{∴ h = 8}
\end{quote}

\textbf{একটি ১৩ মিটার লম্বা মই একটি দেয়ালে হেলান দিয়ে রাখা হয়েছে। মইয়ের
এক প্রান্ত দেয়াল থেকে ৫ মিটার দূরে ভূমি স্পর্শ করেছে। মইয়ের অন্য প্রান্ত ভূমি থেকে
কত উচ্চতায় দেয়াল স্পর্শ করেছে?}

ক) ২০ মিটার

খ) ১৮ মিটার

গ) ১৫ মিটার

\textbf{\hl{ঘ) ১২ মিটার ✔}}

\textbf{\hl{✔ সঠিক উত্তর: ঘ) ১২ মিটার}}

\begin{quote}
\textbf{ব্যাখ্যা:} \emph{সমকোণী ত্রিভুজের বাহুর দৈর্ঘ্যের অনুপাত 5 : 12 : 13}

\emph{∴ দেয়ালের উচ্চতা বা লম্ব হবে 12 মিটার।}
\end{quote}

\textbf{১৮ ফুট উঁচু একটি খুঁটি এমনভাবে ভেঙে গেল যে ভাঙা অংশটি বিচ্ছিন্ন না হয়ে
ভূমির সাথে 30° কোণে স্পর্শ করল। খুঁটিটি মাটি হতে কত ফুট উঁচুতে ভেঙে গিয়েছিল?}

ক) ১২ ফুট

\textbf{\hl{খ) ৯ ফুট ✔}}

গ) ৬ ফুট

ঘ) ৩ ফুট

\textbf{\hl{✔ সঠিক উত্তর: খ) ৯ ফুট}}

\begin{quote}
\textbf{ব্যাখ্যা:} \emph{ধরি, x ফুট উঁচুতে ভেঙেছিল}

\emph{তাহলে ভাঙা অংশের দৈর্ঘ্য = (18 − x) ফুট}

\emph{sin 30° = (x/18−x)}

\emph{⇒ (1/2)=(x/18−x)}

\emph{⇒ 2x = 18 − x}

\emph{⇒ 3x = 18}

\emph{∴ x = 9 ফুট}
\end{quote}

\textbf{১২ মিটার উঁচু একটি গাছ h উচ্চতায় এমনভাবে ভেঙে গেল যে, ভাঙা অংশটি
মাটির সাথে ৩০° কোণে মিলিত হলো। গাছটি কত উঁচুতে ভেঙেছিল?}

\textbf{\hl{ক) ৪ মিটার ✔}}

খ) ৬ মিটার

গ) ৮ মিটার

ঘ) ৩ মিটার

\textbf{\hl{✔ সঠিক উত্তর: ক) ৪ মিটার}}

\begin{quote}
\textbf{ব্যাখ্যা:} \emph{Sin 30° = (AB/BC)}

\emph{⇒ (1/2)=(h/12−h)}

\emph{⇒ 2h = 12 − h}

\emph{∴ h = 4}

\emph{গাছটি ৪ মিটার উচ্চতায় ভেঙেছিল।}
\end{quote}

\textbf{একটি গাছ ঝড়ে এমনভাবে ভেঙে গেল যে তার ভাঙা অংশ দাঁড়ানো অংশের সাথে
30° কোণ করে 10 মিটার দূরে মাটি স্পর্শ করে। গাছটির সম্পূর্ণ দৈর্ঘ্য কত?}

ক) 10 √(3) মিটার

খ) 20 √(3) মিটার

\textbf{\hl{গ) 30 √(3) মিটার ✔}}

ঘ) 40 মিটার

\textbf{\hl{✔ সঠিক উত্তর: গ) 30 √(3) মিটার}}

\begin{quote}
\textbf{ব্যাখ্যা:} \emph{মনে করি, গাছটি মাটি হতে ক ফুট উঁচুতে ভেঙেছিল এবং
বাকি অংশের দৈর্ঘ্য খ ফুট। চিত্রে দেখা যায়, খ = 20}

\emph{sin 60° = (লম্ব/অতিভুজ)=(ক/খ)}

\emph{⇒ (√(3)/2)=(ক/20)}

\emph{⇒ ক = (20√(3)/2)=10√(3)}

\emph{∴ সম্পূর্ণ দৈর্ঘ্য = 10√(3)+20=30√(3)}
\end{quote}

\textbf{সূর্যের উন্নতি কোণ 60° হলে একটি গাছের ছায়ার দৈর্ঘ্য 10 মিটার হয়।
গাছটির উচ্চতা কত?}

\textbf{\hl{ক) 17.32 মি. ✔}}

খ) 16.72 মি.

গ) 17.52 মি.

ঘ) 17.75 মি.

\textbf{\hl{✔ সঠিক উত্তর: ক) 17.32 মি.}}

\begin{quote}
\textbf{ব্যাখ্যা:} \emph{tanθ = (লম্ব/ভূমি)}

\emph{⇒ tan 60° = (উচ্চতা/10)}

\emph{∴ উচ্চতা = 10×√(3) = 17.32 মিটার}
\end{quote}

\textbf{একটি মই এর এক প্রান্ত ভূমি থেকে ১৫ মিটার উঁচু ঘরের জানালা বরাবর
পৌঁছায়। অপর প্রান্ত ঘর থেকে ৮ মিটার দূরে থাকলে মই এর দৈর্ঘ্য কত মিটার?}

\textbf{\hl{ক) ১৭ মিটার ✔}}

খ) ১৮ মিটার

গ) ১৯ মিটার

ঘ) ২০ মিটার

\textbf{\hl{✔ সঠিক উত্তর: ক) ১৭ মিটার}}

\begin{quote}
\textbf{ব্যাখ্যা:} \emph{সমকোণী ত্রিভুজের সূত্র হতে পাই, (অতিভুজ)² = (ভূমি)² +
(লম্ব)²}

\emph{⇒ x² = 8² + 15²}

\emph{⇒ x² = 64 + 225}

\emph{⇒ x² = 289}

\emph{∴ x = √(289) = 17 মিটার}
\end{quote}

\textbf{একটি গাছের পাদবিন্দু থেকে 19 মিটার দূরবর্তী স্থান থেকে গাছের
শীর্ষবিন্দুর উন্নতি কোণ 45° হলে গাছটির উচ্চতা কত?}

ক) 21 মি.

খ) 20 মি.

\textbf{\hl{গ) 19 মি. ✔}}

ঘ) 18 মি.

\textbf{\hl{✔ সঠিক উত্তর: গ) 19 মি.}}

\begin{quote}
\textbf{ব্যাখ্যা:} \emph{দেয়া আছে, BC = 19 মিটার}

\emph{∠ACB = 45°}

\emph{মনে করি, উচ্চতা, AB = h}

\emph{এখন, tan 45° = (AB/BC)}

\emph{⇒ 1 = (h/19)}

\emph{⇒ h = 19}

\emph{∴ উচ্চতা = 19 মিটার}
\end{quote}

\textbf{একজন লোক একটি নির্দিষ্ট স্থান A থেকে যাত্রা শুরু করে ১২ কিলোমিটার ঠিক
উত্তর দিকে গেল এবং সেখান থেকে ৫ কিলোমিটার ঠিক পূর্ব দিকে গেল। যাত্রা শেষে A
থেকে কত দূরে থাকবে?}

ক) ১৭ কিলোমিটার

খ) ১৫ কিলোমিটার

গ) ১৪ কিলোমিটার

\textbf{\hl{ঘ) ১৩ কিলোমিটার ✔}}

\textbf{\hl{✔ সঠিক উত্তর: ঘ) ১৩ কিলোমিটার}}

\begin{quote}
\textbf{ব্যাখ্যা:} \emph{ABC সমকোণী ত্রিভুজে, AC² = AB² + BC²}

\emph{⇒ AC² = 5² + 12² = 25 + 144 = 169}

\emph{∴ AC = √(169) = 13 કિલોમાટર}

\emph{অর্থাৎ সে A থেকে ১৩ কিমি দূরে আছে।}
\end{quote}

\textbf{একটি বাড়ী 40 ফুট উচুঁ। একটি মইয়ের তলদেশ মাটিতে বাড়ীর দেওয়াল থেকে 9
ফুট দূরে রাখা আছে। উপরে মইটি বাড়ীর ছাদ ছুঁয়ে আছে। মইটি কত ফুট লম্বা? / A
building is 40 feet high. The lower end of a ladder is 9 feet from the
wall of the building and its upper end touches the roof of the building.
What is the length of the ladder in feet?}

ক) 48 ফুট

\textbf{\hl{খ) 41 ফুট ✔}}

গ) 44 ফুট

ঘ) 43 ফুট

\textbf{\hl{✔ সঠিক উত্তর: খ) 41 ফুট}}

\begin{quote}
\textbf{ব্যাখ্যা:} \emph{ধরি, মইটি x ফুট লম্বা}

\emph{সমকোণী ত্রিভুজের সূত্র হতে পাই, (অতিভুজ)² = (ভূমি)² + (লম্ব)²}

\emph{⇒ x² = (40)² + (9)²}

\emph{⇒ x² = 1600 + 81}

\emph{⇒ x² = 1681}

\emph{∴ x = √(1681) = 41 ফুট}
\end{quote}

\textbf{একটি ৫০ মিটার লম্বা মই একটি খাড়া দেয়ালের সাথে হেলান দেওয়া আছে।
মইয়ের একপ্রান্ত মাটি হতে ৪০ মিটার উঁচুতে দেয়ালকে স্পর্শ করে। মই এর অপরপ্রান্ত
হতে দেয়ালের দূরত্ব (মিটারে) কত?}

ক) ১০

\textbf{\hl{খ) ৩০ ✔}}

গ) ২০

ঘ) ২৫

\textbf{\hl{✔ সঠিক উত্তর: খ) ৩০}}

\begin{quote}
\textbf{ব্যাখ্যা:} \emph{মনে করি, দেয়ালের পাদদেশ হতে মইয়ের পাদদেশের দূরত্ব =
ক মিটার।}

\emph{দেয়াল মাটির সাথে সমকোণ উৎপন্ন করে।}

\emph{সমকোণী ত্রিভুজের সূত্র হতে পাই, (মইয়ের উচ্চতা)² = (দেয়ালের উচ্চতা)² +
ক²}

\emph{⇒ (৫০)² = (৪০)² + ক²}

\emph{⇒ ২৫০০ = ১৬০০ + ক²}

\emph{⇒ ক² = ৯০০}

\emph{∴ ক = ৩০ মিটার}
\end{quote}

\textbf{একটি মিনারের পাদদেশ হতে 20 মিটার দূরের একটি স্থান হতে মিনারটির
শীর্ষ বিন্দুর উন্নতি কোণ 30° হলে মিনারটির উচ্চতা কত?}

ক) 20 √(3) মিটার

\textbf{\hl{খ) (20/√(3)) মিটার ✔}}

গ) 20 মিটার

ঘ) 10 √(3) মিটার

\textbf{\hl{✔ সঠিক উত্তর: খ) (20/√(3)) মিটার}}

\begin{quote}
\textbf{ব্যাখ্যা:} \emph{ধরি মিনারটির উচ্চতা = x মিটার}

\emph{tan 30° = (AB/BC)}

\emph{⇒ (1/√(3))=(x/20)}

\emph{⇒ x = (20/√(3))}

\emph{∴ মিনারটির উচ্চতা (20/√(3)) মিটার}
\end{quote}

\textbf{একটি বাড়ির উচ্চতা ৪০ ফুট। একটি মইয়ের তলদেশ মাটিতে বাড়ির দেওয়াল
থেকে ৯ ফুট দূরে রাখা আছে। উপরে মইটি বাড়ির ছাদ ছুঁয়ে আছে। মইটি কত ফুট লম্বা?}

\textbf{\hl{ক) ৪১ ✔}}

খ) ৪৮

গ) ৪৩

ঘ) ৪৫

\textbf{\hl{✔ সঠিক উত্তর: ক) ৪১}}

\begin{quote}
\textbf{ব্যাখ্যা:} \emph{সমকোণী ত্রিভুজের সূত্র হতে পাই, (অতিভুজ)² = (ভূমি)² +
(লম্ব)²}

\emph{⇒ (মই)² = (৯)² + (৪০)²}

\emph{⇒ (মই)² = ৮১ + ১৬০০}

\emph{⇒ (মই)² = ১৬৮১}

\emph{∴ মই = √(1681) = ৪১ ফুট}
\end{quote}

\textbf{১৮ ফুট উচুঁ একটি খুঁটি ভেঙে গেল যে ভাঙ্গা অংশটি বিচ্ছিন্ন না হয়ে ভূমির
সাথে ৩০° কোণে স্পর্শ করে, খুঁটিটি মাটি হতে কত ফুট উঁচুতে ভেঙেছিল?}

ক) ১২ ফুট

খ) ৯ ফুট

\textbf{\hl{গ) ৬ ফুট ✔}}

ঘ) ৩ ফুট

\textbf{\hl{✔ সঠিক উত্তর: গ) ৬ ফুট}}

\begin{quote}
\textbf{ব্যাখ্যা:} \emph{মনে করি, খুঁটিটি মাটি হতে ক ফুট উঁচুতে ভেঙে ছিল।}

\emph{অতএব, ভাঙ্গা অংশের দৈর্ঘ্য = (১৮ − ক) ফুট}

\emph{Sin θ = (লম্ব/অতিভুজ)}

\emph{⇒ Sin 30° = (ক/18−ক)}

\emph{⇒ (1/2)=(ক/18−ক)}

\emph{⇒ 18 − ক = 2ক}

\emph{⇒ 3ক = 18}

\emph{∴ ক = 6}

\emph{খুঁটিটি মাটি হতে ৬ ফুট উঁচুতে ভেঙে ছিল।}
\end{quote}

\textbf{একটি ঘুড়ি ভূমি থেকে ৫৫ মিটার উপরে উড়ছে, যার সুতা ভূমির সাথে ৬০° কোণ
উৎপন্ন করে। সুতার দৈর্ঘ্য কত?}

ক) ৬৫ মিটার

খ) ১১০মিটার

\textbf{\hl{গ) ৬৩(১/২) মিটার ✔}}

ঘ) ৬৭(১/২) মিটার

\textbf{\hl{✔ সঠিক উত্তর: গ) ৬৩(১/২) মিটার}}

\begin{quote}
\textbf{ব্যাখ্যা:} \emph{Sin 60° = (ভূমি হতে ঘুড়ির উচ্চতা/সুতার দৈর্ঘ্য)}

\emph{⇒ (√(3)/2)=(55/সুতার দৈর্ঘ্য)}

\emph{∴ সুতার দৈর্ঘ্য = (55×2/√(3)) = 63.5 মিটার বা 63(1/2) মিটার}
\end{quote}

\section{69. ঘন জ্যামিতি MCQ প্রশ্ন উত্তর
অনুশীলন}\label{ux998ux9a8-ux99cux9afux9aeux9a4-mcq-ux9aaux9b0ux9b6ux9a8-ux989ux9a4ux9a4ux9b0-ux985ux9a8ux9b6ux9b2ux9a8}

\textbf{একটি সিলিন্ডারের ভূমির ক্ষেত্রফল 100π বর্গমিটার এবং আয়তন 900π ঘনমিটার
হলে সিলিন্ডারের উচ্চতা কত?}

\textbf{\hl{ক) 9m ✔}}

খ) 10m

গ) 11m

ঘ) 12m

\textbf{\hl{✔ সঠিক উত্তর: ক) 9m}}

\begin{quote}
\textbf{ব্যাখ্যা:} \emph{আমরা জানি, সিলিন্ডারের আয়তন = ভূমির ক্ষেত্রফল ×
উচ্চতা}

\emph{⇒ উচ্চতা = (আয়তন/ভূমির ক্ষেত্রফল)=(900π/100π)=9 মিটার।}
\end{quote}

\textbf{কোন কুয়ার গভীরতা 10 মিটার এবং ব্যাসার্ধ 1 মিটার হলে ঐ কুয়ার আয়তন
কত?}

ক) 100π ঘনমিটার

\textbf{\hl{খ) 10π ঘনমিটার ✔}}

গ) 1000 ঘনমিটার

ঘ) π³ ঘনমিটার

\textbf{\hl{✔ সঠিক উত্তর: খ) 10π ঘনমিটার}}

\begin{quote}
\textbf{ব্যাখ্যা:} \emph{কুয়াটি একটি সিলিন্ডার আকৃতির। সিলিন্ডারের আয়তন
V=πr\^{}2h}

\emph{এখানে, ব্যাসার্ধ r=1 মিটার এবং গভীরতা বা উচ্চতা h=10 মিটার}

\emph{∴ আয়তন = π×1\^{}2×10=10π ঘনমিটার।}
\end{quote}

\textbf{আয়তকার ঘনবস্তুর তল কয়টি?}

ক) ৩ টি

খ) ২ টি

\textbf{\hl{গ) ৬ টি ✔}}

ঘ) ৮ টি

\textbf{\hl{✔ সঠিক উত্তর: গ) ৬ টি}}

\begin{quote}
\textbf{ব্যাখ্যা:} \emph{একটি আয়তকার ঘনবস্তুর ৬টি তল, ১২টি ধার এবং ৮টি
শীর্ষবিন্দু থাকে।}
\end{quote}

\textbf{একটি ঘনকে কয়টি সমকোণ থাকে?}

ক) 8

খ) 12

গ) 16

\textbf{\hl{ঘ) 24 ✔}}

\textbf{\hl{✔ সঠিক উত্তর: ঘ) 24}}

\begin{quote}
\textbf{ব্যাখ্যা:} \emph{একটি ঘনের ৬টি তল এবং প্রতি তলে ৪টি করে সমকোণ থাকে।
সুতরাং মোট সমকোণের সংখ্যা = 6×4=24 টি।}
\end{quote}

\textbf{একটি ঘনকের একটি বাহুর দৈর্ঘ্য ৭ সেন্টিমিটার, এর আয়তন কত?}

ক) ৪৯ ঘন সেন্টিমিটার

খ) ৩৪৩ বর্গ সেন্টিমিটার

গ) ৭২৯ ঘন সেন্টিমিটার

\textbf{\hl{ঘ) ৩৪৩ ঘন সেন্টিমিটার ✔}}

\textbf{\hl{✔ সঠিক উত্তর: ঘ) ৩৪৩ ঘন সেন্টিমিটার}}

\begin{quote}
\textbf{ব্যাখ্যা:} \emph{ঘনকের আয়তন = (বাহুর দৈর্ঘ্য)\^{}3}

\emph{∴ আয়তন = 7\^{}3=343 ঘন সেন্টিমিটার।}
\end{quote}

\textbf{a বাহু বিশিষ্ট ঘনকের কর্ণের দৈর্ঘ্য কত?}

ক) √(2)a\^{}2 একক

\textbf{\hl{খ) √(3)a একক ✔}}

গ) √(2)a একক

ঘ) √(3)a\^{}2 একক

\textbf{\hl{✔ সঠিক উত্তর: খ) √(3)a একক}}

\begin{quote}
\textbf{ব্যাখ্যা:} \emph{আমরা জানি, a বাহু বিশিষ্ট ঘনকের কর্ণের দৈর্ঘ্য হলো
√(3)a একক।}
\end{quote}

\textbf{একটি ঘনকের প্রতিটি বাহু ৫ সে.মি. হলে কর্ণের দৈর্ঘ্য কত?}

\textbf{\hl{ক) 5√(3) সে.মি. ✔}}

খ) 3√(5) সে.মি.

গ) 5√(5) সে.মি.

ঘ) 5√(2) সে.মি.

\textbf{\hl{✔ সঠিক উত্তর: ক) 5√(3) সে.মি.}}

\begin{quote}
\textbf{ব্যাখ্যা:} \emph{ঘনকের কর্ণের দৈর্ঘ্য = √(3)a একক}

\emph{এখানে বাহুর দৈর্ঘ্য a=5 সে.মি.}

\emph{∴ কর্ণের দৈর্ঘ্য = 5√(3) সে.মি.}
\end{quote}

\textbf{একটি ঘনবস্তুর দৈর্ঘ্য, প্রস্থ ও উচ্চতা সমান। বস্তুটির আয়তন ৩৪৩ ঘন সে.মি.
হলে তার একটি তলের ক্ষেত্রফল কত?}

ক) ৭ বর্গ সে.মি.

\textbf{\hl{খ) ৪৯ বর্গ সে.মি. ✔}}

গ) ৯ বর্গ সে.মি.

ঘ) ৩৬ বর্গ সে.মি.

\textbf{\hl{✔ সঠিক উত্তর: খ) ৪৯ বর্গ সে.মি.}}

\begin{quote}
\textbf{ব্যাখ্যা:} \emph{এখানে ঘনবস্তুটি একটি ঘনক, যার একবাহু a হলে আয়তন
a\^{}3=343}

\emph{⇒ a=7 সে.মি.}

\emph{∴ একটি তলের ক্ষেত্রফল = a\^{}2=7\^{}2=49 বর্গ সে.মি.}
\end{quote}

\textbf{একটি ঘনক আকৃতির বাক্সের বাহুর দৈর্ঘ্য ৩ মিটার হলে বাক্সের তলের মোট
ক্ষেত্রফল কত বর্গমিটার?}

ক) ৯

খ) ১৮

গ) ২৭

\textbf{\hl{ঘ) ৫৪ ✔}}

\textbf{\hl{✔ সঠিক উত্তর: ঘ) ৫৪}}

\begin{quote}
\textbf{ব্যাখ্যা:} \emph{বাক্সের তলের মোট ক্ষেত্রফল = 6a\^{}2}

\emph{এখানে, বাহুর দৈর্ঘ্য a=3 মিটার}

\emph{∴ মোট ক্ষেত্রফল = 6×3\^{}2=6×9=54 বর্গমিটার।}
\end{quote}

\textbf{4 একক ধারবিশিষ্ট একটি ঘনকের দুটি কর্ণের সমষ্টি কত একক?}

ক) √(3)

খ) 4√(3)

\textbf{\hl{গ) 8√(3) ✔}}

ঘ) 16√(3)

\textbf{\hl{✔ সঠিক উত্তর: গ) 8√(3)}}

\begin{quote}
\textbf{ব্যাখ্যা:} \emph{আমরা জানি, ঘনকের কর্ণের দৈর্ঘ্য = √(3)a একক}

\emph{এখানে, প্রতিটি ধারের দৈর্ঘ্য a=4 একক}

\emph{∴ একটি কর্ণের দৈর্ঘ্য = 4√(3) একক}

\emph{অতএব, দুটি কর্ণের সমষ্টি = 4√(3)+4√(3)=8√(3) একক।}
\end{quote}

\textbf{একটি ঘনকের সম্পূর্ণ পৃষ্ঠের ক্ষেত্রফল 8.4 হেক্টর। এর কর্ণের দৈর্ঘ্য কত?}

\begin{quote}
\textbf{ব্যাখ্যা:} \emph{১ হেক্টর = 10,000 বর্গমিটার, সুতরাং 8.4 হেক্টর =
84,000 বর্গমিটার।}

\emph{আমরা জানি, ঘনকের সম্পূর্ণ পৃষ্ঠের ক্ষেত্রফল = 6a\^{}2}

\emph{⇒ 6a\^{}2=84000}

\emph{⇒ a\^{}2=14000}

\emph{⇒ a=√(14000)=20√(35) মিটার।}
\end{quote}

\textbf{২.৫ মিটার বর্গাকার একটি খোলা চৌবাচ্চায় ২৮,৯০০ লিটার পানি ধরে। এর
ভিতরে এলুমিনিয়ামের পাত লাগাতে প্রতি বর্গমিটারে ৫ টাকা হিসাবে মোট কত খরচ
পড়বে?}

ক) ২২৯.৬০ টাকা

খ) ১২৯.২৯ টাকা

গ) ২২৭.৮০ টাকা

ঘ) ৬০০ টাকা

\textbf{\hl{ঙ) কোনোটিই নয় ✔}}

\textbf{\hl{✔ সঠিক উত্তর: ঙ) কোনোটিই নয়}}

\begin{quote}
\textbf{ব্যাখ্যা:} \emph{চৌবাচ্চাটির আয়তন = ২৮,৯০০ লিটার = ২৮.৯ ঘনমিটার {[}
1000 লিটার = 1 ঘনমিটার{]}}

\emph{চৌবাচ্চাটির গভীরতা
h=(আয়তন/দৈর্ঘ্য×প্রস্থ)=(২৮.৯/২.৫×২.৫)=(২৮.৯/৬.২৫)=৪.৬২৪ মিটার}

\emph{খোলা চৌবাচ্চার তলা ও ৪টি পাতের ক্ষেত্রফল =
(২.৫×২.৫)+৪×(২.৫×৪.৬২৪)=৬.২৫+৪৬.২৪=৫২.৪৯ বর্গমিটার}

\emph{∴ মোট খরচ = ৫২.৪৯×৫=২৬২.৪৫ টাকা (অপশনে সঠিক উত্তর না থাকায় সঠিক
উত্তর \textquotesingle কোনোটিই নয়\textquotesingle)}
\end{quote}

\textbf{১৮ ইঞ্চি উঁচু একটি বাক্সের দৈর্ঘ্য ৩ ফুট এবং প্রস্থ ২ ফুট। বাক্সটির আয়তন
কত?}

ক) ৮ ঘনফুট

\textbf{\hl{খ) ৯ ঘনফুট ✔}}

গ) ১০৮ ঘনফুট

ঘ) ৬ ঘনফুট

\textbf{\hl{✔ সঠিক উত্তর: খ) ৯ ঘনফুট}}

\begin{quote}
\textbf{ব্যাখ্যা:} \emph{উচ্চতা = (18/12) ফুট = 1.5 ফুট {[} 12 ইঞ্চি = 1
ফুট{]}}

\emph{∴ বাক্সটির আয়তন = দৈর্ঘ্য × প্রস্থ × উচ্চতা = 3×2×1.5=9 ঘনফুট।}
\end{quote}

\textbf{একটি ঘনকের সমকোণের মোট সংখ্যা-}

ক) ৪ টি

খ) ৮ টি

গ) ১৮ টি

\textbf{\hl{ঘ) উপরিউক্ত কোনটিই নয় ✔}}

\textbf{\hl{✔ সঠিক উত্তর: ঘ) উপরিউক্ত কোনটিই নয়}}

\begin{quote}
\textbf{ব্যাখ্যা:} \emph{একটি ঘনকের ৬টি তল থাকে এবং প্রতিটি তলে ৪টি করে
সমকোণ থাকে।}

\emph{∴ মোট সমকোণের সংখ্যা = 6×4=24 টি, যা প্রদত্ত অপশনগুলোতে নেই।}
\end{quote}

\textbf{গোলকের আয়তনের সূত্র কোনটি?}

ক) π×(ব্যাসার্ধ)\^{}2×দৈর্ঘ্য

খ) π×(ব্যাসার্ধ)\^{}3

\textbf{\hl{গ) (4/3)×π×(ব্যাসার্ধ)\^{}3 ✔}}

ঘ) 2×π×(ব্যাসার্ধ)\^{}2

\textbf{\hl{✔ সঠিক উত্তর: গ) (4/3)×π×(ব্যাসার্ধ)\^{}3}}

\begin{quote}
\textbf{ব্যাখ্যা:} \emph{আমরা জানি, গোলকের আয়তন নির্ণয়ের সূত্র হলো
(4/3)πr\^{}3 বা (4/3)×π×(ব্যাসার্ধ)\^{}3}
\end{quote}

\textbf{দুটি গোলকের আয়তনের অনুপাত ৮: ২৭। তাদের ক্ষেত্রফলের অনুপাত কত?}

\textbf{\hl{ক) ৪: ৯ ✔}}

খ) ৪: ১৫

গ) ২: ৩

ঘ) ৫: ৬

\textbf{\hl{✔ সঠিক উত্তর: ক) ৪: ৯}}

\begin{quote}
\textbf{ব্যাখ্যা:} \emph{গোলকের আয়তন V=(4/3)πr\^{}3 এবং পৃষ্ঠের ক্ষেত্রফল
S=4πr\^{}2 .}

\emph{ধরি, গোলক দুটির ব্যাসার্ধ r\_1 ও r\_2 .}

\emph{প্রশ্নমতে, r\_1\^{}3:r\_2\^{}3=8:27⇒r\_1:r\_2=2:3 .}

\emph{∴ ক্ষেত্রফলের অনুপাত =
4πr\_1\^{}2:4πr\_2\^{}2=r\_1\^{}2:r\_2\^{}2=2\^{}2:3\^{}2=4:9 .}
\end{quote}

\textbf{দুটি গোলকের ব্যাসার্ধের অনুপাত ৩: ২ হলে তাদের আয়তনের অনুপাত কত?}

ক) ৯: ৪

\textbf{\hl{খ) ২৭: ৮ ✔}}

গ) ১২: ৪

ঘ) ২৭: ৪

\textbf{\hl{✔ সঠিক উত্তর: খ) ২৭: ৮}}

\begin{quote}
\textbf{ব্যাখ্যা:} \emph{ধري, ব্যাসার্ধদ্বয় r\_1 এবং r\_2 .}

\emph{এখানে,
(r\_1/r\_2)=(3/2)⇒(r\_1\^{}3/r\_2\^{}3)=(3\^{}3/2\^{}3)=(27/8) .}

\emph{∴ আয়তনের অনুপাত =
((4/3)πr\_1\^{}3/(4/3)πr\_2\^{}3)=(r\_1\^{}3/r\_2\^{}3)=27:8 .}
\end{quote}

\textbf{একটি বেলুনের বক্রতলের ক্ষেত্রফল 100 বর্গ সেমি এবং আয়তন 150 ঘন সেমি।
বেলুনের ভূমির ব্যাসার্ধ কত?}

ক) 5 সেমি

খ) 4 সেমি

\textbf{\hl{গ) 3 সেমি ✔}}

ঘ) 6 সেমি

\textbf{\hl{✔ সঠিক উত্তর: গ) 3 সেমি}}

\begin{quote}
\textbf{ব্যাখ্যা:} \emph{বেলুনের বক্রতলের ক্ষেত্রফল = 2πrh=100 বর্গ সেমি.}

\emph{বেলুনের আয়তন = πr\^{}2h=150 ঘন সেমি.}

\emph{এখন, আয়তনকে বক্রতলের ক্ষেত্রফল দ্বারা ভাগ করে পাই,
(πr\^{}2h/2πrh)=(150/100)⇒(r/2)=(3/2)⇒r=3 সেমি.}
\end{quote}

\textbf{216cm³ আয়তন বিশিষ্ট একটি ঘনকের মধ্যে কত বড় আকারের একটি গোলক রাখা
যাবে?}

ক) 49π/3

\textbf{\hl{খ) 36π ✔}}

গ) 15√(2)π

ঘ) 3√(3)π

\textbf{\hl{✔ সঠিক উত্তর: খ) 36π}}

\begin{quote}
\textbf{ব্যাখ্যা:} \emph{ধরি, ঘনকের একবাহুর দৈর্ঘ্য a cm হলে আয়তন a\^{}3=216}

\emph{⇒ a=6 cm}

\emph{ঘনকের মধ্যে রাখা বৃহত্তম গোলকের ব্যাস = ঘনকের বাহুর দৈর্ঘ্য = 6 cm}

\emph{∴ গোলকের ব্যাসার্ধ r=(6/2)=3 cm}

\emph{অতএব, সম্ভাব্য বৃহত্তম গোলকের আয়তন/আকার =
(4/3)πr\^{}3=(4/3)π×3\^{}3=36π cm³ (অথবা পৃষ্ঠের ক্ষেত্রফল ক্ষেত্রবিশেষে
4πr\^{}2=36π ).}
\end{quote}

\section{70. অসমতা MCQ প্রশ্ন উত্তর
অনুশীলন}\label{ux985ux9b8ux9aeux9a4-mcq-ux9aaux9b0ux9b6ux9a8-ux989ux9a4ux9a4ux9b0-ux985ux9a8ux9b6ux9b2ux9a8}

\textbf{2x\^{}2+5x+3\textless0 এর সমাধান কোনটি?}

\textbf{\hl{ক) −(3/2) \textless x\textless−1 ✔}}

খ) −(3/2) \textless x\textless1

গ) −(3/2) ≤x≤1

ঘ) −(3/2) \textless x\textless1

\textbf{\hl{✔ সঠিক উত্তর: ক) −(3/2) \textless x\textless−1}}

\begin{quote}
\textbf{ব্যাখ্যা:} \emph{2x\^{}2 + 5x + 3 \textless{} 0}

\emph{⇒ 2x\^{}2 + 3x + 2x + 3 \textless{} 0}

\emph{⇒ x(2x + 3) + 1(2x + 3) \textless{} 0}

\emph{⇒ (x + 1)(2x + 3) \textless{} 0}

\emph{∴ নির্ণেয় অসমতাটি: −(3/2) \textless{} x \textless{} −1}
\end{quote}

\textbf{\textbar1−2x\textbar{} \textless{} 1 এর সমাধান---}

ক) −2 \textless{} x \textless{} 1

খ) −1 \textless{} x \textless{} 0

\textbf{\hl{গ) 0 \textless{} x \textless{} 1 ✔}}

ঘ) −1 \textless{} x \textless{} 1

\textbf{\hl{✔ সঠিক উত্তর: গ) 0 \textless{} x \textless{} 1}}

\begin{quote}
\textbf{ব্যাখ্যা:} \emph{অঋণাত্মক হলে,}

\emph{1 − 2x \textless{} 1}

\emph{⇒ −2x \textless{} 0}

\emph{⇒ 2x \textgreater{} 0}

\emph{∴ x \textgreater{} 0}

\emph{ঋণাত্মক হলে,}

\emph{−(1 − 2x) \textless{} 1}

\emph{⇒ 1 − 2x \textgreater{} −1}

\emph{⇒ −2x \textgreater{} −1 − 1}

\emph{⇒ −2x \textgreater{} −2}

\emph{⇒ 2x \textless{} 2}

\emph{∴ x \textless{} 1}

\emph{∴ নির্ণেয় সমাধান: 0 \textless{} x \textless{} 1}
\end{quote}

\textbf{বাস্তব সংখ্যায় \textbar2x−3\textbar{} ≤ 1 অসমতাটির সমাধান---}

ক) 1 \textless{} x \textless{} 2

খ) x ≤ 1 অথবা x ≥ 2

\textbf{\hl{গ) 1 ≤ x ≤ 2 ✔}}

ঘ) −1 \textless{} x \textless{} 2

\textbf{\hl{✔ সঠিক উত্তর: গ) 1 ≤ x ≤ 2}}

\begin{quote}
\textbf{ব্যাখ্যা:} \emph{\textbar2x − 3\textbar{} ≤ 1}

\emph{⇒ −1 ≤ 2x − 3 ≤ 1}

\emph{⇒ −1 + 3 ≤ 2x ≤ 1 + 3}

\emph{⇒ 2 ≤ 2x ≤ 4}

\emph{⇒ 1 ≤ x ≤ 2}

\emph{∴ নির্ণেয় সমাধান: 1 ≤ x ≤ 2}
\end{quote}

\textbf{x\^{}2 − 5x + 6 \textless{} 0 হলে}

\textbf{\hl{ক) 2 \textless{} x \textless{} 3 ✔}}

খ) −3 \textless{} x \textless{} −2

গ) x \textless{} 2

ঘ) x \textless{} 3

\textbf{\hl{✔ সঠিক উত্তর: ক) 2 \textless{} x \textless{} 3}}

\begin{quote}
\textbf{ব্যাখ্যা:} \emph{x\^{}2 − 5x + 6 \textless{} 0}

\emph{⇒ x\^{}2 − 3x − 2x + 6 \textless{} 0}

\emph{⇒ x(x − 3) − 2(x − 3) \textless{} 0}

\emph{⇒ (x − 3)(x − 2) \textless{} 0}

\emph{যেহেতু দুটি রাশির গুণফল শূন্য অপেক্ষা ছোট, তাই এদের একটি ঋণাত্মক হবে।}

\emph{এখন x \textgreater{} 2 এবং x \textless{} 3 অর্থাৎ 2 \textless{} x
\textless{} 3}
\end{quote}

\textbf{x\^{}2 − 4 ≤ 0 এর সমাধান কোনটি?}

\textbf{\hl{ক) −2 ≤ x ≤ 2 ✔}}

খ) −2 ≤ x \textless{} 2

গ) −2 \textless{} x ≤ 2

ঘ) −2 \textless{} x \textless{} 2

\textbf{\hl{✔ সঠিক উত্তর: ক) −2 ≤ x ≤ 2}}

\begin{quote}
\textbf{ব্যাখ্যা:} \emph{দেয়া আছে, x\^{}2 − 4 ≤ 0}

\emph{⇒ x\^{}2 ≤ 4}

\emph{⇒ √(x\^{}2) ≤ √(4)}

\emph{⇒ x ≤ ±2}

\emph{অর্থাৎ −2 ≤ x ≤ 2}
\end{quote}

\textbf{(a−3)(b+3) \textless{} 0 হলো নিচের কোনটি সত্য?}

\textbf{\hl{ক) a \textgreater{} 3, b \textless{} −3 ✔}}

খ) a \textgreater{} 3, b \textgreater{} −3

গ) a \textless{} 3, b \textless{} −3

ঘ) a = 3, b = −3

\textbf{\hl{✔ সঠিক উত্তর: ক) a \textgreater{} 3, b \textless{} −3}}

\begin{quote}
\textbf{ব্যাখ্যা:} \emph{(a − 3)(b + 3) \textless{} 0 অসমতাটি সত্য যদি (a −
3) ও (b + 3) এর যে কোনো একটি ঋণাত্মক হয়।}

\emph{i. a − 3 \textless{} 0 হলে, b + 3 \textgreater{} 0}

\emph{∴ a \textless{} 3 হলে, b \textgreater{} −3}

\emph{আবার,}

\emph{ii. a − 3 \textgreater{} 0 হলে, b + 3 \textless{} 0}

\emph{∴ a \textgreater{} 3 হলে, b \textless{} −3}
\end{quote}

\textbf{যদি xy \textless{} 0, xz \textgreater{} 0 and z \textless{} 0
হয়, তবে নিচের কোনটি অবশ্যই সত্য হবে?}

\textbf{\hl{ক) y \textgreater{} 0 ✔}}

খ) y \textless{} 0

গ) yz \textgreater{} 0

ঘ) কোনটিই নয়

\textbf{\hl{✔ সঠিক উত্তর: ক) y \textgreater{} 0}}

\begin{quote}
\textbf{ব্যাখ্যা:} \emph{দেয়া আছে, z \textless{} 0 এবং xz \textgreater{}
0। যেহেতু x এবং z এর গুণফল ধনাত্মক এবং z ঋণাত্মক, সুতরাং x অবশ্যই ঋণাত্মক হবে
(x \textless{} 0)।}

\emph{আবার, xy \textless{} 0। যেহেতু x ঋণাত্মক, তাই xy কে ঋণাত্মক হতে হলে y
কে অবশ্যই ধনাত্মক হতে হবে।}

\emph{∴ y \textgreater{} 0}
\end{quote}

\textbf{x+3 \textgreater{} 2x−1 অসমতাটির সমাধান সেট হবে?}

ক) (∞,0)

\textbf{\hl{খ) (−∞,4) ✔}}

গ) (−∞,4{]}

ঘ) (0,4)

\textbf{\hl{✔ সঠিক উত্তর: খ) (−∞,4)}}

\begin{quote}
\textbf{ব্যাখ্যা:} \emph{x + 3 \textgreater{} 2x − 1}

\emph{⇒ 2x − 1 \textless{} x + 3}

\emph{⇒ 2x − x \textless{} 3 + 1}

\emph{∴ x \textless{} 4}

\emph{∴ নির্ণেয় সমাধান সেট: (−∞, 4)}
\end{quote}

\textbf{\textbar x−3\textbar{} \textgreater{} 5 হলে}

ক) 2 \textless{} x \textless{} 8

খ) −2 \textless{} x \textless{} 8

\textbf{\hl{গ) −8 \textless{} x \textless{} −2 ✔}}

ঘ) −4 \textless{} x \textless{} −2

\textbf{\hl{✔ সঠিক উত্তর: গ) −8 \textless{} x \textless{} −2}}

\begin{quote}
\textbf{ব্যাখ্যা:} \emph{ধনাত্মক মান নিয়ে,}

\emph{\textbar x − 3\textbar{} \textgreater{} 5}

\emph{⇒ x − 3 \textgreater{} 5}

\emph{⇒ x \textgreater{} 5 + 3}

\emph{⇒ x \textgreater{} 8}

\emph{আবার, ঋণাত্মক মান নিয়ে,}

\emph{−(x − 3) \textgreater{} 5}

\emph{⇒ −x + 3 \textgreater{} 5}

\emph{⇒ −x \textgreater{} 5 − 3}

\emph{⇒ −x \textgreater{} 2}

\emph{⇒ x \textless{} −2}

\emph{∴ −8 \textless{} x \textless{} −2}
\end{quote}

\textbf{\textbar x−5\textbar{} ≤ 4 হলে, নিচের যে ব্যবধি সত্য তা হলো---}

ক) x ∈ (1, 9)

খ) x ∈ {[}1, 9)

গ) x ∈ (1, 9{]}

\textbf{\hl{ঘ) x ∈ {[}1, 9{]} ✔}}

\textbf{\hl{✔ সঠিক উত্তর: ঘ) x ∈ {[}1, 9{]}}}

\begin{quote}
\textbf{ব্যাখ্যা:} \emph{\textbar x − 5\textbar{} ≤ 4}

\emph{⇒ −4 ≤ x − 5 ≤ 4}

\emph{⇒ −4 + 5 ≤ x ≤ 4 + 5}

\emph{⇒ 1 ≤ x ≤ 9}

\emph{∴ x ∈ {[}1, 9{]}}
\end{quote}

\textbf{a\textgreater b এবং ab \textless{} 0 হলে, নিচের কোনটি ঋণাত্মক?}

\textbf{\hl{ক) b ✔}}

খ) a

গ) a − b

ঘ) a\^{}2 − b\^{}2

\textbf{\hl{✔ সঠিক উত্তর: ক) b}}

\begin{quote}
\textbf{ব্যাখ্যা:} \emph{যেহেতু, ab \textless{} 0 সেহেতু a ও b এর যে কোন
একটি অবশ্যই ঋণাত্মক। সুতরাং b অবশ্যই ঋণাত্মক হবে যেহেতু a \textgreater{} b।}
\end{quote}

\textbf{x \textgreater{} y এবং z \textless{} 0 হলে, নিচের কোনটি সঠিক?}

ক) xz \textgreater{} yz

খ) (x/z)\textgreater(y/z)

গ) (z/x)\textless(z/y)

\textbf{\hl{ঘ) xz \textless{} yz ✔}}

\textbf{\hl{✔ সঠিক উত্তর: ঘ) xz \textless{} yz}}

\begin{quote}
\textbf{ব্যাখ্যা:} \emph{যেহেতু z \textless{} 0; ∴ z একটি ঋণাত্মক সংখ্যা।
দেয়া আছে, x \textgreater{} y। সুতরাং, xz \textless{} yz {[}উভয়পক্ষকে z
দ্বারা গুণ করে{]}}

\emph{z একটি ঋণাত্মক সংখ্যা বলে z দ্বারা উভয়পক্ষকে গুণ করায়
\textquotesingle\textgreater\textquotesingle{} চিহ্ন পরিবর্তিত হয়ে
\textquotesingle\textless\textquotesingle{} হয়েছে।}
\end{quote}

\textbf{\textbar x−2\textbar{} ≤ 5 হলে, x এর সর্বনিম্ন মান কত?}

ক) −2

খ) −3

\textbf{\hl{গ) 2 ✔}}

ঘ) 5

\textbf{\hl{✔ সঠিক উত্তর: গ) 2}}

\begin{quote}
\textbf{ব্যাখ্যা:} \emph{\textbar x − 2\textbar{} ≤ 5}

\emph{⇒ −5 ≤ x − 2 ≤ 5}

\emph{⇒ −5 + 2 ≤ x − 2 + 2 ≤ 5 + 2}

\emph{⇒ −3 ≤ x ≤ 7}

\emph{∴ x এর সর্বনিম্ন মান 2।}
\end{quote}

\textbf{If 3−2x ≤ 7, then}

ক) x ≤ 2

খ) x ≥ 2

গ) x ≤ −2

\textbf{\hl{ঘ) x ≥ −2 ✔}}

\textbf{\hl{✔ সঠিক উত্তর: ঘ) x ≥ −2}}

\begin{quote}
\textbf{ব্যাখ্যা:} \emph{3 − 2x ≤ 7}

\emph{⇒ −2x ≤ 7 − 3}

\emph{⇒ −2x ≤ 4}

\emph{⇒ x ≥ (4/−2)}

\emph{∴ x ≥ −2}
\end{quote}

\textbf{If 6−4x ≤ 14, then}

ক) x ≤ 2

খ) x ≥ 2

গ) x ≤ −2

\textbf{\hl{ঘ) x ≥ −2 ✔}}

ঙ) x \textgreater{} −2

\textbf{\hl{✔ সঠিক উত্তর: ঘ) x ≥ −2}}

\begin{quote}
\textbf{ব্যাখ্যা:} \emph{6 − 4x ≤ 14}

\emph{⇒ 6 − 14 ≤ 4x}

\emph{⇒ −8 ≤ 4x}

\emph{⇒ −2 ≤ x}

\emph{∴ x ≥ −2}
\end{quote}

\textbf{x\^{}2 + x − 2 \textgreater{} 0 অসমতাটির সমাধান করুন।}

ক) \{−2, 1\}

খ) (−2, 1)

\textbf{\hl{গ) (−∞, −2) ∪ (1, ∞) ✔}}

ঘ) (−2, ∞)

\textbf{\hl{✔ সঠিক উত্তর: গ) (−∞, −2) ∪ (1, ∞)}}

\begin{quote}
\textbf{ব্যাখ্যা:} \emph{x\^{}2 + x − 2 \textgreater{} 0}

\emph{⇒ x\^{}2 + 2x − x − 2 \textgreater{} 0}

\emph{⇒ x(x + 2) − 1(x + 2) \textgreater{} 0}

\emph{⇒ (x + 2)(x − 1) \textgreater{} 0}

\emph{∴ x \textgreater{} (−2, 1) {[}অর্থাৎ x \textless{} −2 অথবা x
\textgreater{} 1, যা ব্যবধি আকারে (−∞, −2) ∪ (1, ∞) যায়{]}}
\end{quote}

\textbf{x\^{}3 − 5x + 6 \textless{} 0 এর সমাধান}

ক) 2 ≤ x \textless{} 3

\textbf{\hl{খ) 2 \textless{} x \textless{} 3 ✔}}

গ) 2 \textless{} x ≤ 3

ঘ) 2 ≤ x ≤ 3

\textbf{\hl{✔ সঠিক উত্তর: খ) 2 \textless{} x \textless{} 3}}

\begin{quote}
\textbf{ব্যাখ্যা:} \emph{x\^{}3 − 5x + 6 রাশির x\^{}3 পদের পরিবর্তে x\^{}2
হলে উত্তর হবে (খ) 2 \textless{} x \textless{} 3{]}}

\textbf{ব্যাখ্যা:} \emph{x\^{}2 − 5x + 6 \textless{} 0}

\emph{⇒ x\^{}2 − 3x − 2x + 6 \textless{} 0}

\emph{⇒ x(x − 3) − 2(x − 3) \textless{} 0}

\emph{⇒ (x − 3)(x − 2) \textless{} 0}

\emph{(x − 3) ও (x − 2) এর যে কোনো একটির মান ঋণাত্মক হলে অসমতাটি সত্য
হবে।}

\emph{∴ নির্ণেয় অসমতা: 2 \textless{} x \textless{} 3}
\end{quote}

\textbf{3x − 2 \textgreater{} 2x − 1 এর সমাধান সেট কোনটি?}

ক) {[}1, ∞)

\textbf{\hl{খ) (1, ∞) ✔}}

গ) ((1/2),∞)

ঘ) {[}−1, ∞)

\textbf{\hl{✔ সঠিক উত্তর: খ) (1, ∞)}}

\begin{quote}
\textbf{ব্যাখ্যা:} \emph{3x − 2 \textgreater{} 2x − 1}

\emph{⇒ 3x − 2x − 2 \textgreater{} 2x − 2x − 1 {[}উভয় পক্ষ থেকে 2x বিয়োগ
করে{]}}

\emph{⇒ x − 2 \textgreater{} −1}

\emph{⇒ x − 2 + 2 \textgreater{} −1 + 2 {[}উভয় পক্ষে 2 যোগ করে{]}}

\emph{⇒ x \textgreater{} 1}

\emph{∴ নির্ণেয় সমাধান সেট = (1, ∞)।}
\end{quote}

\textbf{সমাধান করুন (Solve) : x\^{}2 − 2x − 15 \textgreater{} 0.}

\textbf{\hl{ক) s = \{x: x \textgreater{} 5 অথবা x \textless{} −3\} ✔}}

খ) s = \{x: x \textless{} 5 অথবা x \textgreater{} −3\}

গ) s = \{x: x \textgreater{} 5 অথবা x \textgreater{} −3\}

ঘ) s = \{x: x \textless{} 5 অথবা x \textless{} −3\}

\textbf{\hl{✔ সঠিক উত্তর: ক) s = \{x: x \textgreater{} 5 অথবা x
\textless{} −3\}}}

\begin{quote}
\textbf{ব্যাখ্যা:} \emph{x\^{}2 − 2x − 15 \textgreater{} 0}

\emph{⇒ x\^{}2 − 5x + 3x − 15 \textgreater{} 0}

\emph{⇒ x(x − 5) + 3(x − 5) \textgreater{} 0}

\emph{⇒ (x − 5)(x + 3) \textgreater{} 0}

\emph{দুটি রাশির গুণফল তখনই ধনাত্মক বা শূন্য অপেক্ষা বড় হবে যখন উভয়ই ধনাত্মক বা
উভয়ই ঋণাত্মক হবে।}

\emph{এখন x \textgreater{} 5 অর্থাৎ x = 6, 7, 8........ হলে রাশি দুটির
গুণফল শূন্য অপেক্ষা বড় হবে।}

\emph{এখন x \textless{} −3 অর্থাৎ x = −4, −5, −6 ... হলে রাশি দুটির গুণফল
শূন্য অপেক্ষা বড় হবে।}

\emph{∴ নির্ণেয় সমাধান: x \textgreater{} 5 অথবা x \textless{} −3 অথবা,
(−∞, −3) ∪ (5, +∞)}
\end{quote}

\textbf{যদি y=(3x−6/x) হয় তবে x এর কোন মানের জন্য y ধনাত্মক হবে? / If
y=(3x−6/x) , then find any value of x such that y is positive?}

ক) x \textgreater{} 0

খ) x \textless{} 0

\textbf{\hl{গ) x \textgreater{} 2 অথবা x \textless{} 0 ✔}}

ঘ) −2 \textless{} x \textless{} 2 এবং x ≠ 2

\textbf{\hl{✔ সঠিক উত্তর: গ) x \textgreater{} 2 অথবা x \textless{} 0}}

\begin{quote}
\textbf{ব্যাখ্যা:} \emph{সঠিক উত্তর হবে x \textgreater{} 2 অথবা x
\textless{} 0। কারণ,}

\emph{x \textgreater{} 2 অর্থাৎ x = 3, 4, 5........ হলে y এর মান ধনাত্মক
হবে।}

\emph{x \textless{} 0 অর্থাৎ x = −1, −2, −3.... হলে y এর মান ধনাত্মক হবে।}
\end{quote}

\textbf{x\^{}2 − 3x − 10 \textgreater{} 0 অসমতাটির সমাধান করুন।}

ক) (−∞, −1) ∪ (4, +∞)

\textbf{\hl{খ) (−∞, −2) ∪ (5, +∞) ✔}}

গ) (∞, 2) ∪ (5, +∞)

ঘ) (−∞, 2) ∪ (−5, −∞)

\textbf{\hl{✔ সঠিক উত্তর: খ) (−∞, −2) ∪ (5, +∞)}}

\begin{quote}
\textbf{ব্যাখ্যা:} \emph{x\^{}2 − 3x − 10 \textgreater{} 0}

\emph{⇒ x\^{}2 − 5x + 2x − 10 \textgreater{} 0}

\emph{⇒ x(x − 5) + 2(x − 5) \textgreater{} 0}

\emph{⇒ (x − 5)(x + 2) \textgreater{} 0}

\emph{∴ নির্ণেয় সমাধান: x \textgreater{} 5 অথবা x \textless{} −2 অর্থাৎ
(−∞, −2) ∪ (5, +∞)}
\end{quote}

\textbf{2x−7 \textless{} 8 \textless{} 3x−11 হলে x−এর মান পূর্ণ সংখ্যায়
কত?}

ক) 6

খ) 8

\textbf{\hl{গ) 7 ✔}}

ঘ) 9

\textbf{\hl{✔ সঠিক উত্তর: গ) 7}}

\begin{quote}
\textbf{ব্যাখ্যা:} \emph{2x − 7 \textless{} 8 ⇒ 2x \textless{} 15 ⇒ x
\textless{} 7.5}

\emph{এবং 3x − 11 \textgreater{} 8 ⇒ 3x \textgreater{} 19 ⇒ x
\textgreater{} 6.33}

\emph{⇒ 6.33 \textless{} x \textless{} 7.5}

\emph{∴ x = 7 (পূর্ণ সংখ্যায়)}
\end{quote}

\textbf{a ও b positive integers এবং (a−b/3.5)=(4/7) তাহলে}

\textbf{\hl{ক) b \textless{} a ✔}}

খ) b \textgreater{} a

গ) b = a

ঘ) b ≥ a

\textbf{\hl{✔ সঠিক উত্তর: ক) b \textless{} a}}

\begin{quote}
\textbf{ব্যাখ্যা:} \emph{(a−b/3.5)=(4/7)}

\emph{⇒ a − b = (4×3.5/7)}

\emph{⇒ a − b = 2}

\emph{⇒ a = 2 + b}

\emph{∴ a \textgreater{} b অর্থাৎ b \textless{} a}
\end{quote}

\textbf{যদি a \textgreater{} b এবং c \textgreater{} 0 হয়, তবে নিচের কোন
সম্পর্কটি সঠিক?}

ক) (a/c)=(b/c)

\textbf{\hl{খ) (a/c)\textgreater(b/c) ✔}}

গ) (a/c)\textless(b/c)

ঘ) (a/c)≤(b/c)

\textbf{\hl{✔ সঠিক উত্তর: খ) (a/c)\textgreater(b/c)}}

\begin{quote}
\textbf{ব্যাখ্যা:} \emph{যেহেতু a \textgreater{} b এবং c \textgreater{} 0,
উভয় পক্ষকে ধনাত্মক সংখ্যা c দ্বারা ভাগ করলে অসমতার দিক অপরিবর্তিত থাকে।}

\emph{∴ (a/c)\textgreater(b/c)}
\end{quote}

\textbf{যদি x \textless{} y এবং a \textless{} b হয় তবে, কোন সম্পর্কটি
সঠিক?}

\textbf{\hl{ক) a + x \textless{} b + y ✔}}

খ) a + x \textgreater{} b + y

গ) a = y

ঘ) ax = by

\textbf{\hl{✔ সঠিক উত্তর: ক) a + x \textless{} b + y}}

\begin{quote}
\textbf{ব্যাখ্যা:} \emph{দেয়া আছে, x \textless{} y -\/-\/- (1)}

\emph{এবং a \textless{} b -\/-\/- (2)}

\emph{(1) ও (2) নং অসমতা দুটি যোগ করে পাই,}

\emph{a + x \textless{} b + y}
\end{quote}

\textbf{যদি x \textgreater{} 0, y \textgreater{} 0 এবং
(x/1)\textgreater(1/y) হয়, তবে x এবং y এর মধ্যে কী সম্পর্ক?}

ক) x \textgreater{} y

খ) x \textless{} y

গ) x = y

\textbf{\hl{ঘ) সবগুলো ✔}}

\textbf{\hl{✔ সঠিক উত্তর: ঘ) সবগুলো}}

\begin{quote}
\textbf{ব্যাখ্যা:} \emph{(x/1)\textgreater(1/y) {[}x \textgreater{} 0, y
\textgreater{} 0{]}}

\emph{⇒ x \textgreater{} (1/y)}

\emph{∴ xy \textgreater{} 1}

\emph{∴ x ও y এর মধ্যে সম্পর্ক প্রদত্ত অপশনের কোনটি নয়। (উত্তর: ঘ)}
\end{quote}

\textbf{(1/\textbar2x−5\textbar)\textless(1/3) এর সমাধান}

\textbf{\hl{ক) x \textless{} 1, x \textgreater{} 4 ✔}}

খ) x \textgreater{} 1, x \textgreater{} 4

গ) x \textless{} 1, x \textless{} 4

ঘ) 1 \textless{} x \textless{} 4

\textbf{\hl{✔ সঠিক উত্তর: ক) x \textless{} 1, x \textgreater{} 4}}

\begin{quote}
\textbf{ব্যাখ্যা:} \emph{−(1/3)\textless(1/2x−5)\textless(1/3)}

\emph{⇒ −3 \textgreater{} 2x − 5 \textgreater{} 3}

\emph{⇒ −3 + 5 \textgreater{} 2x − 5 + 5 \textgreater{} 3 + 5}

\emph{⇒ 2 \textgreater{} 2x \textgreater{} 8}

\emph{⇒ 1 \textgreater{} x \textgreater{} 4}

\emph{∴ x \textless{} 1, x \textgreater{} 4}
\end{quote}

\textbf{\textbar x\textbar{} \textless{} 3 এর সমাধান সেট:}

ক) −3 ≤ x ≤ 3

খ) 0 \textless{} x \textless{} 3

\textbf{\hl{গ) −3 \textless{} x \textless{} 3 ✔}}

ঘ) −3 \textless{} x \textless{} 3

\textbf{\hl{✔ সঠিক উত্তর: গ) −3 \textless{} x \textless{} 3}}

\begin{quote}
\textbf{ব্যাখ্যা:} \emph{পরসমান চিহ্ন থেকে \textbar x\textbar{} \textless{}
3 অসমতাটিকে প্রথমে ধনাত্মক ও পরে ঋণাত্মক ধরে −3 \textless{} x \textless{} 3
মান পাওয়া যায়।}
\end{quote}

\textbf{\textbar x−5\textbar{} = 6 সমীকরণটির সমাধান সেট হলো---}

ক) \{1, −11\}

খ) \{1, 11\}

\textbf{\hl{গ) \{−1, 11\} ✔}}

ঘ) \{−1, −11\}

\textbf{\hl{✔ সঠিক উত্তর: গ) \{−1, 11\}}}

\begin{quote}
\textbf{ব্যাখ্যা:} \emph{প্রদত্ত সমীকরণ, \textbar x − 5\textbar{} = 6}

\emph{(x − 5) অঋণাত্মক অর্থাৎ (x − 5) \textgreater{} 0 হলে, x − 5 = 6 ⇒ x
= 11}

\emph{আবার, ঋণাত্মক হলে, −(x − 5) = 6 ⇒ −x + 5 = 6 ⇒ −x = 1 ⇒ x = −1}

\emph{∴ সমাধান সেট \{−1, 11\}}
\end{quote}

\section{71. বীজগণিতীয় ভগ্নাংশ MCQ প্রশ্ন উত্তর
অনুশীলন}\label{ux9acux99cux997ux9a3ux9a4ux9df-ux9adux997ux9a8ux9b6-mcq-ux9aaux9b0ux9b6ux9a8-ux989ux9a4ux9a4ux9b0-ux985ux9a8ux9b6ux9b2ux9a8}

\textbf{যদি x:y=y:z=1.5 এবং z=2 হয়, তবে x এর মান কত?}

ক) 4

খ) 5

\textbf{\hl{গ) 4.5 ✔}}

ঘ) 5.4

\textbf{\hl{✔ সঠিক উত্তর: গ) 4.5}}

\begin{quote}
\textbf{ব্যাখ্যা:} \emph{y:z=1.5}

\emph{⇒ (y/z)=1.5 ⇒ (y/2)=1.5 ∴ y=3}

\emph{এবং x:y=1.5}

\emph{⇒ (x/y)=1.5 ⇒ (x/3)=1.5 ∴ x=4.5}
\end{quote}

\textbf{3x যদি 15 থেকে 3(Three) বেশী হয়, তাহলে 3x+2 = কত?}

ক) 17

খ) 18

\textbf{\hl{গ) 20 ✔}}

ঘ) 21

\textbf{\hl{✔ সঠিক উত্তর: গ) 20}}

\begin{quote}
\textbf{ব্যাখ্যা:} \emph{3x−15=3 ⇒ 3x=15+3 ∴ 3x=18}

\emph{∴ 3x+2=18+2=20}
\end{quote}

\textbf{x যদি y এর চেয়ে বড় হয় তবে (1/x) এর চেয়ে (1/y) -\/-\/-}

\textbf{\hl{ক) বড় ✔}}

খ) ছোট

গ) সমান

ঘ) অসমান

\textbf{\hl{✔ সঠিক উত্তর: ক) বড়}}

\begin{quote}
\textbf{ব্যাখ্যা:} \emph{ভগ্নাংশের লব একই হলে হর যার বড় সে ছোট।}

\emph{এখন, (1/x) ও (1/y) ভগ্নাংশে x\textgreater y}

\emph{(1/x)\textless(1/y) অর্থাৎ (1/y) এর চেয়ে (1/x) বড়।}
\end{quote}

\textbf{a : b = 2 : 3 এবং b : c = 6 : 7 হয় তবে, a : c = কত?}

ক) 2 : 6

খ) 3 : 7

গ) 2 : 7

\textbf{\hl{ঘ) 4 : 7 ✔}}

\textbf{\hl{✔ সঠিক উত্তর: ঘ) 4 : 7}}

\begin{quote}
\textbf{ব্যাখ্যা:} \emph{a:b=2:3=4:6}

\emph{b:c=6:7 ∴ a:c=4:7}
\end{quote}

\textbf{a : b = 4 : 7 এবং b : c = 5:6 হলে a : b : c = কত?}

ক) 4:7:6

খ) 20:35:24

\textbf{\hl{গ) 20:35:42 ✔}}

ঘ) 24:35:30

\textbf{\hl{✔ সঠিক উত্তর: গ) 20:35:42}}

\begin{quote}
\textbf{ব্যাখ্যা:} \emph{দেওয়া আছে, a:b=4:7 এবং b:c=5:6}

\emph{প্রথম অনুপাতের উভয় রাশিকে 5 দ্বারা এবং দ্বিতীয় অনুপাতের উভয় রাশিকে 7
দ্বারা গুণ করে পাই:}

\emph{a:b=(4×5):(7×5)=20:35}

\emph{b:c=(5×7):(6×7)=35:42}

\emph{∴ a:b:c=20:35:42}
\end{quote}

\textbf{যদি 5:100,4:64 হয় তবে 4:80,3:?}

ক) ২৬

খ) ৫৪

\textbf{\hl{গ) ৪৮ ✔}}

ঘ) ৬০

\textbf{\hl{✔ সঠিক উত্তর: গ) ৪৮}}

\begin{quote}
\textbf{ব্যাখ্যা:} \emph{যদি 5:100,4:64}

\emph{তবে, 4:80,3:48}

\emph{১ম অনুপাতের ক্ষেত্রে ২য় রাশি, ১ম রাশির ২০ গুণ এবং ২য় অনুপাতের ক্ষেত্রে ২য়
রাশি ১ম রাশির ১৬ গুণ।}
\end{quote}

\textbf{10a\^{}2b\^{}4 কে 5a\^{}2b\^{}2 দ্বারা ভাগ করলে ভাগফল কত হবে?}

ক) 2a\^{}2b\^{}2

\textbf{\hl{খ) 2b\^{}2 ✔}}

গ) 2ab

ঘ) 2a\^{}2

\textbf{\hl{✔ সঠিক উত্তর: খ) 2b\^{}2}}

\begin{quote}
\textbf{ব্যাখ্যা:} \emph{নির্ণেয় ভাগফল
=(10a\^{}2b\^{}4/5a\^{}2b\^{}2)=2b\^{}2}
\end{quote}

\textbf{(3/5) এর লব ও হরের সাথে কোন একই সংখ্যা যোগ করলে ভগ্নাংশটি (4/5)
হয়?}

\textbf{\hl{ক) ৫ ✔}}

খ) ৬

গ) ৪

ঘ) ৩

\textbf{\hl{✔ সঠিক উত্তর: ক) ৫}}

\begin{quote}
\textbf{ব্যাখ্যা:} \emph{ধরি, সংখ্যাটি x ।}

\emph{শর্তমতে, (3+x/5+x)=(4/5)}

\emph{⇒ 15+5x=20+4x}

\emph{∴ x=5}
\end{quote}

\textbf{(x/y) এর সাথে কত যোগ করলে যোগফল (2y/x) হবে?}

\textbf{\hl{ক) (2y\^{}2−x\^{}2/xy) ✔}}

খ) (x\^{}2−2y\^{}2/xy)

গ) (2x\^{}2−y\^{}2/xy)

ঘ) (x\^{}2−y\^{}2/xy)

\textbf{\hl{✔ সঠিক উত্তর: ক) (2y\^{}2−x\^{}2/xy)}}

\begin{quote}
\textbf{ব্যাখ্যা:} \emph{ধরি, a যোগ করতে হবে।}

\emph{শর্তমতে, (x/y)+a=(2y/x)}

\emph{or, a=(2y/x)−(x/y) ∴ a=(2y\^{}2−x\^{}2/xy)}
\end{quote}

\textbf{0.01×0.01=?}

ক) 0.01

খ) 0.1

গ) 0.001

\textbf{\hl{ঘ) 0.0001 ✔}}

\textbf{\hl{✔ সঠিক উত্তর: ঘ) 0.0001}}

\begin{quote}
\textbf{ব্যাখ্যা:} \emph{0.01×0.01=0.0001}
\end{quote}

\textbf{(x/y) এর সাথে কত যোগ করলে যোগফল (y/x) হবে?}

ক) (x\^{}2−y\^{}2/xy)

খ) (2x\^{}2−y\^{}2/xy)

\textbf{\hl{গ) (y\^{}2−x\^{}2/xy) ✔}}

ঘ) (x\^{}2−2y\^{}2/xy)

\textbf{\hl{✔ সঠিক উত্তর: গ) (y\^{}2−x\^{}2/xy)}}

\begin{quote}
\textbf{ব্যাখ্যা:} \emph{ধরি, a যোগ করতে হবে।}

\emph{শর্তমতে, (x/y)+a=(y/x)}

\emph{or, a=(y/x)−(x/y)}

\emph{∴ a=(y\^{}2−x\^{}2/xy)}
\end{quote}

\textbf{(1/2)((1/x−1)−(1/x+1))−(1/x\^{}2+1)=?}

ক) (2/x\^{}4+1)

খ) (2x\^{}2/x\^{}4−1)

গ) (2x\^{}2/x\^{}4+1)

\textbf{\hl{ঘ) (2/x\^{}4−1) ✔}}

\textbf{\hl{✔ সঠিক উত্তর: ঘ) (2/x\^{}4−1)}}

\begin{quote}
\textbf{ব্যাখ্যা:} \emph{(1/2)((1/x−1)−(1/x+1))−(1/x\^{}2+1)}

\emph{⇒ (1/2)(((x+1)−(x−1)/x\^{}2−1))−(1/x\^{}2+1)}

\emph{⇒ (1/2)((2/x\^{}2−1))−(1/x\^{}2+1)}

\emph{⇒ (1/x\^{}2−1)−(1/x\^{}2+1)}

\emph{⇒ ((x\^{}2+1)−(x\^{}2−1)/(x\^{}2−1)(x\^{}2+1))}

\emph{∴ (2/x\^{}4−1)}
\end{quote}

\textbf{যদি x=(3/4) হয়, তবে (√(1+x)+√(1−x)/√(1+x)−√(1−x))= কত?}

ক) 1

\textbf{\hl{খ) 2 ✔}}

গ) -1

ঘ) -2

\textbf{\hl{✔ সঠিক উত্তর: খ) 2}}

\begin{quote}
\textbf{ব্যাখ্যা:} \emph{দেওয়া আছে, x=(3/4)}

\emph{বা, (1/x)=(4/3) {[}ব্যস্তকরণ করে{]}}

\emph{বা, (1+x/1−x)=(4+3/4−3) {[}যোজন বিয়োজন করে{]}}

\emph{বা, (√(1+x)/√(1−x))=(√(7)/1) {[}বর্গমূল করে{]}}

\emph{আবার যোজন বিয়োজন করে পাই,}

\emph{∴ (√(1+x)+√(1−x)/√(1+x)−√(1−x))=(√(7)+1/√(7)−1)=2}
\end{quote}

\textbf{(5x−7/(x−1)(x−2)) কে আংশিক ভগ্নাংশে প্রকাশ কর।}

ক) (1/x−1)

খ) (2/x−1)−(3/x−2)

\textbf{\hl{গ) (2/x−1)+(3/x−2) ✔}}

ঘ) (3/x−1)+(2/x−2)

\textbf{\hl{✔ সঠিক উত্তর: গ) (2/x−1)+(3/x−2)}}

\begin{quote}
\textbf{ব্যাখ্যা:} \emph{ধরি, (5x−7/(x−1)(x−2))=(A/x−1)+(B/x−2)}

\emph{বা, 5x−7=A(x−2)+B(x−1)}

\emph{x = 1 বসিয়ে পাই, 5(1)−7=A(1−2)∴A=2}

\emph{x = 2 বসিয়ে পাই, 5(2)−7=B(2−1)∴B=3}

\emph{সুতরাং নির্ণেয় আংশিক ভগ্নাংশ =(2/x−1)+(3/x−2)}
\end{quote}

\textbf{(a\^{}3−b\^{}3/a+b) কে
(a\^{}4+b\^{}4+a\^{}2b\^{}2/a\^{}3+b\^{}3) দ্বারা ভাগ করলে ভাগফল কত হবে?}

\textbf{\hl{ক) a − b ✔}}

খ) a + b

গ) a\^{}2 − b\^{}2

ঘ) a\^{}2 + b\^{}2

\textbf{\hl{✔ সঠিক উত্তর: ক) a − b}}

\begin{quote}
\textbf{ব্যাখ্যা:} \emph{দেওয়া আছে,
(a\^{}3−b\^{}3/a+b)÷(a\^{}4+b\^{}4+a\^{}2b\^{}2/a\^{}3+b\^{}3)}

\emph{⇒
((a−b)(a\^{}2+ab+b\^{}2)/a+b)÷((a\^{}2)\^{}2+2a\^{}2b\^{}2+(b\^{}2)\^{}2−a\^{}2b\^{}2/(a+b)(a\^{}2−ab+b\^{}2))}

\emph{⇒
((a−b)(a\^{}2+ab+b\^{}2)/a+b)÷((a\^{}2+b\^{}2)\^{}2−(ab)\^{}2/(a+b)(a\^{}2−ab+b\^{}2))}

\emph{⇒
((a−b)(a\^{}2+ab+b\^{}2)/a+b)÷((a\^{}2+b\^{}2+ab)(a\^{}2+b\^{}2−ab)/(a+b)(a\^{}2−ab+b\^{}2))}

\emph{⇒
((a−b)(a\^{}2+ab+b\^{}2)/a+b)×((a+b)(a\^{}2−ab+b\^{}2)/(a\^{}2+ab+b\^{}2)(a\^{}2−ab+b\^{}2))}

\emph{∴ a−b}
\end{quote}

\textbf{(bc/(a−b)(a−c))+(ca/(b−c)(b−a))+(ab/(c−a)(c−b))=?}

\textbf{\hl{ক) 0 ✔}}

খ) 1

গ) 3

ঘ) 4

\textbf{\hl{✔ সঠিক উত্তর: ক) 0}}

\begin{quote}
\textbf{ব্যাখ্যা:} \emph{প্রদত্ত রাশি:}

\emph{(bc/(a−b)(a−c))+(ca/(b−c)(b−a))+(ab/(c−a)(c−b))}

\emph{হরগুলোকে ক্রমানুসারে সাজিয়ে লসাগু করলে পাই,}

\emph{লসাগু =(a−b)(b−c)(c−a)}

\emph{লবগুলোর সরলীকরণ থেকে প্রাপ্ত মান শূন্য হয়।}

\emph{∴ নির্ণেয় মান =0}
\end{quote}

\textbf{(a−b/ab)+(b−c/bc)+(c−a/ca)=?}

\textbf{\hl{ক) 0 ✔}}

খ) 1

গ) abc

ঘ) a+b+c

\textbf{\hl{✔ সঠিক উত্তর: ক) 0}}

\begin{quote}
\textbf{ব্যাখ্যা:} \emph{প্রদত্ত রাশি: (a−b/ab)+(b−c/bc)+(c−a/ca)}

\emph{ল.সা.গু (L.C.M) নির্ণয় করে পাই:}

\emph{(c(a−b) + a(b−c) + b(c−a)/abc)}

\emph{⇒ (ac−bc+ab−ac+bc−ab/abc)}

\emph{⇒ (0/abc)}

\emph{∴ 0}
\end{quote}

\section{72. বীজগাণিতিক সূত্রাবলী ও মান নির্ণয় MCQ প্রশ্ন উত্তর
অনুশীলন}\label{ux9acux99cux997ux9a3ux9a4ux995-ux9b8ux9a4ux9b0ux9acux9b2-ux993-ux9aeux9a8-ux9a8ux9b0ux9a3ux9af-mcq-ux9aaux9b0ux9b6ux9a8-ux989ux9a4ux9a4ux9b0-ux985ux9a8ux9b6ux9b2ux9a8}

\textbf{বাস্তব সংখ্যার ক্ষেত্রে অভেদ উপাদান -}

ক) 0, 1

খ) -1, 1

গ) 2, -2

\textbf{\hl{ঘ) 5, 10 ✔}}

\textbf{\hl{✔ সঠিক উত্তর: ঘ) 5, 10}}

\begin{quote}
\textbf{ব্যাখ্যা:} \emph{যে সব সমীকরণে চলমান রাশির যেকোনো মান দ্বারা
উভয়পক্ষকে সমান দেখানো যায়, তাকে বলা হয় অভেদ। বীজগাণিতিক সূত্রগুলো প্রতিটিই
অভেদ। উদাহরণস্বরূপ:}

\emph{(a + b)\^{}5 = a\^{}5 + 5a\^{}4b\^{}2 + 10a\^{}3b\^{}2 +
10a\^{}2b\^{}3 + 5ab\^{}4 + b\^{}5}

\emph{∴ এখানে 5 ও 10 হলো অভেদ উপাদান।}
\end{quote}

\textbf{কোনটি অভেদ?}

ক) x\^{}2 + 5x + 6 = 0

খ) a\^{}2 − 10a + 9 = 0

গ) 4x + 5 = 9

\textbf{\hl{ঘ) (p + q)\^{}2 = p\^{}2 + 2pq + q\^{}2 ✔}}

\textbf{\hl{✔ সঠিক উত্তর: ঘ) (p + q)\^{}2 = p\^{}2 + 2pq + q\^{}2}}

\begin{quote}
\textbf{ব্যাখ্যা:} \emph{সকল বীজগণিতীয় সূত্রই অভেদ, অর্থাৎ চলমান রাশি p ও
q-এর যেকোনো মানের জন্য সমীকরণটি সত্য হয়। কিন্তু বাকি বিকল্পগুলো নির্দিষ্ট মানের
জন্যই কেবল সত্য হয় (এগুলো সমীকরণ, অভেদ নয়)। তাই (p + q)\^{}2 = p\^{}2 + 2pq
+ q\^{}2 একটি অভেদ।}
\end{quote}

\textbf{a + b = 8 এবং a − b = 2 হলে ab এর মান কত?}

ক) 4

খ) 8

\textbf{\hl{গ) 15 ✔}}

ঘ) 16

\textbf{\hl{✔ সঠিক উত্তর: গ) 15}}

\begin{quote}
\textbf{ব্যাখ্যা:} \emph{a + b = 8 ......... (1)}

\emph{a − b = 2 ......... (2)}

\emph{(1) ও (2) যোগ করে: 2a = 10 ⇒ a = 5}

\emph{(1) থেকে: b = 8 − 5 = 3}

\emph{∴ ab = 5 × 3 = 15}
\end{quote}

\textbf{x\^{}2 + y\^{}2 = 8 এবং xy = 7 হলে (x + y)\^{}2 এর মান কত?}

ক) 19

খ) 20

গ) 21

\textbf{\hl{ঘ) 22 ✔}}

\textbf{\hl{✔ সঠিক উত্তর: ঘ) 22}}

\begin{quote}
\textbf{ব্যাখ্যা:} \emph{আমরা জানি, (x + y)\^{}2 = x\^{}2 + y\^{}2 + 2xy =
8 + (2 × 7) = 8 + 14 = 22}
\end{quote}

\textbf{If a + b = − c, and a\^{}2 + 2ab + b\^{}2 is expressed in terms
of c, which one of the following will be?}

ক) -c\^{}2

\textbf{\hl{খ) c\^{}2 ✔}}

গ) bc

ঘ) ca

\textbf{\hl{✔ সঠিক উত্তর: খ) c\^{}2}}

\begin{quote}
\textbf{ব্যাখ্যা:} \emph{a\^{}2 + 2ab + b\^{}2}

\emph{= (a + b)\^{}2}

\emph{= (−c)\^{}2}

\emph{= c\^{}2}
\end{quote}

\textbf{a + b = 3, ab = 1 হলে a\^{}2 − ab + b\^{}2 = কত?}

ক) 5

\textbf{\hl{খ) 6 ✔}}

গ) 10

ঘ) 8

\textbf{\hl{✔ সঠিক উত্তর: খ) 6}}

\begin{quote}
\textbf{ব্যাখ্যা:} \emph{দেয়া আছে, a\^{}2 − ab + b\^{}2}

\emph{= a\^{}2 + 2ab + b\^{}2 − 3ab}

\emph{= (a + b)\^{}2 − 3ab}

\emph{= 3\^{}2 − (3 × 1)}

\emph{= 9 − 3 = 6}
\end{quote}

\textbf{x−y = 2 এবং xy = 24 হলে x + y এর মান কত?}

ক) ± 20

খ) ± 5

\textbf{\hl{গ) ± 10 ✔}}

ঘ) ± 15

\textbf{\hl{✔ সঠিক উত্তর: গ) ± 10}}

\begin{quote}
\textbf{ব্যাখ্যা:} \emph{দেওয়া আছে, x − y = 2 এবং xy = 24}

\emph{আমরা জানি, (x + y)\^{}2 = (x − y)\^{}2 + 4xy}

\emph{⇒ (x + y) = √(((x-y))\^{}2+4xy)}

\emph{⇒ (x + y) = √(2\^{}2+4×24)}

\emph{⇒ (x + y) = √(4+96)}

\emph{⇒ (x + y) = √(100)}

\emph{∴ (x + y) = ± 10}
\end{quote}

\textbf{a + b = 7 এবং ab = 10 হলে a\^{}2 + b\^{}2 + 3ab = কত?}

ক) 29

খ) 49

\textbf{\hl{গ) 59 ✔}}

ঘ) 69

\textbf{\hl{✔ সঠিক উত্তর: গ) 59}}

\begin{quote}
\textbf{ব্যাখ্যা:} \emph{a\^{}2 + b\^{}2 + 3ab}

\emph{= (a + b)\^{}2 − 2ab + 3ab}

\emph{= (a + b)\^{}2 + ab}

\emph{মান বসিয়ে, = 7\^{}2 + 10 = 49 + 10 = 59}
\end{quote}

\textbf{a\^{}2 + 2ab + b\^{}2 এর বর্গমূল কত?}

\textbf{\hl{ক) a + b ✔}}

খ) a − b

গ) a + 2b

ঘ) a − 2b

\textbf{\hl{✔ সঠিক উত্তর: ক) a + b}}

\begin{quote}
\textbf{ব্যাখ্যা:} \emph{a\^{}2 + 2ab + b\^{}2 এর বর্গমূল =
√(a\^{}2+2ab+b\^{}2)}

\emph{= √(((a+b))\^{}2)}

\emph{= a + b}
\end{quote}

\textbf{a = √(5) , b = √(7) হলে (a + b)\^{}2 − 2ab এর মান কত?}

\textbf{\hl{ক) 12 ✔}}

খ) 5

গ) 2

ঘ) 7

\textbf{\hl{✔ সঠিক উত্তর: ক) 12}}

\begin{quote}
\textbf{ব্যাখ্যা:} \emph{(a + b)\^{}2 − 2ab}

\emph{= a\^{}2 + 2ab + b\^{}2 − 2ab}

\emph{= a\^{}2 + b\^{}2 = ( √(5) )\^{}2 + ( √(7) )\^{}2 = 5 + 7 = 12}
\end{quote}

\textbf{a = 2, b = −3 হলে 16a\^{}2 + 24ab + 9b\^{}2 এর মান কত?}

\textbf{\hl{ক) 1 ✔}}

খ) −1

গ) 5

ঘ) −6

\textbf{\hl{✔ সঠিক উত্তর: ক) 1}}

\begin{quote}
\textbf{ব্যাখ্যা:} \emph{16a\^{}2 + 24ab + 9b\^{}2}

\emph{= (4a)\^{}2 + 2 × 4a × 3b + (3b)\^{}2}

\emph{= (4a + 3b)\^{}2}

\emph{মান বসিয়ে,}

\emph{= \{4(2) + 3(−3)\}\^{}2}

\emph{= (8 − 9)\^{}2}

\emph{= (−1)\^{}2 = 1}
\end{quote}

\textbf{x\^{}2 + y\^{}2 = 8 এবং xy = 7 হলে (x + y)\^{}2 এর মান কত?}

ক) 14

খ) 16

\textbf{\hl{গ) 22 ✔}}

ঘ) 30

\textbf{\hl{✔ সঠিক উত্তর: গ) 22}}

\begin{quote}
\textbf{ব্যাখ্যা:} \emph{(x + y)\^{}2 = x\^{}2 + y\^{}2 + 2xy = 8 + 2 × 7
= 8 + 14 = 22}
\end{quote}

\textbf{a + (1/a) = 3 হলে a\^{}2 + (1/a\^{}2) = কত?}

\textbf{\hl{ক) 7 ✔}}

খ) 9

গ) 11

ঘ) 13

\textbf{\hl{✔ সঠিক উত্তর: ক) 7}}

\begin{quote}
\textbf{ব্যাখ্যা:} \emph{a + (1/a) = 3}

\emph{উভয়পক্ষকে বর্গ করে পাই,}

\emph{⇒ ((a+(1/a)))\^{}2 = 3\^{}2}

\emph{⇒ a\^{}2 + 2 · a · (1/a) + (1/a\^{}2) = 9}

\emph{⇒ a\^{}2 + 2 + (1/a\^{}2) = 9}

\emph{⇒ a\^{}2 + (1/a\^{}2) = 9 − 2}

\emph{∴ a\^{}2 + (1/a\^{}2) = 7}
\end{quote}

\textbf{(2x + 1) এর বর্গ কত?}

\textbf{\hl{ক) 4x\^{}2 + 4x + 1 ✔}}

খ) 2x\^{}2 + 4x + 1

গ) 4x\^{}2 + 2x + 1

ঘ) 2x\^{}2 + 2x + 1

\textbf{\hl{✔ সঠিক উত্তর: ক) 4x\^{}2 + 4x + 1}}

\begin{quote}
\textbf{ব্যাখ্যা:} \emph{(2x + 1)\^{}2}

\emph{= (2x)\^{}2 + 2 · (2x) · 1 + 1\^{}2}

\emph{= 4x\^{}2 + 4x + 1}
\end{quote}

\textbf{a\^{}2 + b\^{}2 = 25 এবং ab = 12 হলে a + b = কত?}

ক) 6

\textbf{\hl{খ) 7 ✔}}

গ) 8

ঘ) 9

\textbf{\hl{✔ সঠিক উত্তর: খ) 7}}

\begin{quote}
\textbf{ব্যাখ্যা:} \emph{আমরা জানি,}

\emph{(a + b)\^{}2 = a\^{}2 + b\^{}2 + 2ab = 25 + 2 × 12 = 25 + 24 = 49}

\emph{⇒ (a + b)\^{}2 = 49}

\emph{∴ a + b = √(49) = 7}
\end{quote}

\textbf{x − y = 3 এবং xy = 10 হলে, (x + y)\^{}2 এর মান---}

\textbf{\hl{ক) 49 ✔}}

খ) −20

গ) −1

ঘ) 0

\textbf{\hl{✔ সঠিক উত্তর: ক) 49}}

\begin{quote}
\textbf{ব্যাখ্যা:} \emph{(x + y)\^{}2 = (x − y)\^{}2 + 4xy = 3\^{}2 + 4 ×
10 = 9 + 40 = 49}
\end{quote}

\textbf{x\^{}2 + y\^{}2 = 8 এবং xy = 7 হলে (x + y)\^{}2 এর মান কত?}

ক) 18

খ) 16

\textbf{\hl{গ) 22 ✔}}

ঘ) 30

\textbf{\hl{✔ সঠিক উত্তর: গ) 22}}

\begin{quote}
\textbf{ব্যাখ্যা:} \emph{(x + y)\^{}2 = x\^{}2 + y\^{}2 + 2xy = 8 + 2 × 7
= 8 + 14 = 22}
\end{quote}

\textbf{x\^{}2 − 2x − 1 = 0 হলে x + (1/x) = কত?}

ক) (1/2)

\textbf{\hl{খ) 2 √(2) ✔}}

গ) 1

ঘ) 2

\textbf{\hl{✔ সঠিক উত্তর: খ) 2 √(2)}}

\begin{quote}
\textbf{ব্যাখ্যা:} \emph{দেওয়া আছে, x\^{}2 − 2x − 1 = 0}

\emph{⇒ x\^{}2 − 1 = 2x}

\emph{উভয়পক্ষকে x দ্বারা ভাগ করে পাই,}

\emph{⇒ x − (1/x) = 2}

\emph{আমরা জানি,}

\emph{((x+(1/x)))\^{}2 = ((x-(1/x)))\^{}2 + 4 · x · (1/x)}

\emph{⇒ ((x+(1/x)))\^{}2 = 2\^{}2 + 4 = 4 + 4 = 8}

\emph{∴ x + (1/x) = √(8) = √(4×2) = 2 √(2)}
\end{quote}

\textbf{x − (1/x) = 2 হলে ((x+(1/x)))\^{}2 = ?}

ক) 4

খ) 16

\textbf{\hl{গ) 8 ✔}}

ঘ) 6

\textbf{\hl{✔ সঠিক উত্তর: গ) 8}}

\begin{quote}
\textbf{ব্যাখ্যা:} \emph{((x+(1/x)))\^{}2 = ((x-(1/x)))\^{}2 + 4 · x ·
(1/x) = 2\^{}2 + 4 = 4 + 4 = 8}
\end{quote}

\textbf{যদি a − (1/a) = 2 হয়, তবে a\^{}4 + (1/a\^{}4) = কত?}

ক) 36

খ) 32

\textbf{\hl{গ) 34 ✔}}

ঘ) 40

\textbf{\hl{✔ সঠিক উত্তর: গ) 34}}

\begin{quote}
\textbf{ব্যাখ্যা:} \emph{a − (1/a) = 2}

\emph{উভয়পক্ষকে বর্গ করে পাই,}

\emph{⇒ ((a-(1/a)))\^{}2 = 2\^{}2}

\emph{⇒ a\^{}2 − 2 + (1/a\^{}2) = 4}

\emph{⇒ a\^{}2 + (1/a\^{}2) = 6}

\emph{পুনরায় বর্গ করে পাই,}

\emph{⇒ ((a\^{}2+(1/a\^{}2)))\^{}2 = 6\^{}2}

\emph{⇒ a\^{}4 + 2 + (1/a\^{}4) = 36}

\emph{∴ a\^{}4 + (1/a\^{}4) = 36 − 2 = 34}
\end{quote}

\textbf{a + (1/a) = 3 হলে a\^{}4 + a\^{}−4 এর মান −}

\textbf{\hl{ক) 47 ✔}}

খ) 49

গ) 79

ঘ) 81

\textbf{\hl{✔ সঠিক উত্তর: ক) 47}}

\begin{quote}
\textbf{ব্যাখ্যা:} \emph{দেওয়া আছে, a + (1/a) = 3}

\emph{উভয়পক্ষকে বর্গ করে পাই,}

\emph{⇒ ((a+(1/a)))\^{}2 = 3\^{}2}

\emph{⇒ a\^{}2 + (1/a\^{}2) + 2 · a · (1/a) = 9}

\emph{⇒ a\^{}2 + (1/a\^{}2) + 2 = 9}

\emph{⇒ a\^{}2 + (1/a\^{}2) = 9 − 2}

\emph{⇒ a\^{}2 + (1/a\^{}2) = 7}

\emph{পুনরায় বর্গ করে পাই,}

\emph{⇒ ((a\^{}2+(1/a\^{}2)))\^{}2 = 7\^{}2}

\emph{⇒ a\^{}4 + (1/a\^{}4) + 2 · a\^{}2 · (1/a\^{}2) = 49}

\emph{⇒ a\^{}4 + (1/a\^{}4) + 2 = 49}

\emph{⇒ a\^{}4 + (1/a\^{}4) = 49 − 2}

\emph{∴ a\^{}4 + a\^{}−4 = 47}
\end{quote}

\textbf{If a\^{}4 + a\^{}2b\^{}2 + b\^{}4 = a\^{}2 + ab + b\^{}2 = 3,
what is the value of a\^{}2 + b\^{}2?}

ক) 4

খ) 66

গ) 3

\textbf{\hl{ঘ) 2 ✔}}

\textbf{\hl{✔ সঠিক উত্তর: ঘ) 2}}

\begin{quote}
\textbf{ব্যাখ্যা:} \emph{a\^{}4 + a\^{}2b\^{}2 + b\^{}4 = 3}

\emph{⇒ (a\^{}2)\^{}2 + 2a\^{}2b\^{}2 + (b\^{}2)\^{}2 − a\^{}2b\^{}2 =
3}

\emph{⇒ (a\^{}2 + b\^{}2)\^{}2 − (ab)\^{}2 = 3}

\emph{⇒ (a\^{}2 + ab + b\^{}2)(a\^{}2 − ab + b\^{}2) = 3}

\emph{⇒ 3 × (a\^{}2 − ab + b\^{}2) = 3 {[}∴ a\^{}2 + ab + b\^{}2 = 3{]}}

\emph{⇒ a\^{}2 − ab + b\^{}2 = (3/3)}

\emph{⇒ a\^{}2 − ab + b\^{}2 = 1}

\emph{এখন,}

\emph{a\^{}2 + ab + b\^{}2 = 3 ......... (i)}

\emph{a\^{}2 − ab + b\^{}2 = 1 ......... (ii)}

\emph{(i) + (ii) করে পাই,}

\emph{2(a\^{}2 + b\^{}2) = 4}

\emph{∴ a\^{}2 + b\^{}2 = 2}
\end{quote}

\textbf{2x + (2/x) = 4 হলে x এর মান কত?}

ক) 2

\textbf{\hl{খ) 1 ✔}}

গ) (1/2)

ঘ) (1/0)

\textbf{\hl{✔ সঠিক উত্তর: খ) 1}}

\begin{quote}
\textbf{ব্যাখ্যা:} \emph{দেওয়া আছে, 2x + (2/x) = 4}

\emph{⇒ 2 (x+(1/x)) = 4}

\emph{∴ x + (1/x) = (4/2) = 2 ......... (i)}

\emph{এখন,}

\emph{((x-(1/x)))\^{}2 = ((x+(1/x)))\^{}2 − 4 · x · (1/x)}

\emph{⇒ ((x-(1/x)))\^{}2 = 2\^{}2 − 4 = 4 − 4 = 0}

\emph{∴ x − (1/x) = 0 ......... (ii)}

\emph{(i) ও (ii) যোগ করে পাই,}

\emph{2x = 2}

\emph{∴ x = (2/2) = 1}
\end{quote}

\textbf{√(x+3) = √(x) + √(3) হলে x = কত?}

ক) 3

খ) −3

\textbf{\hl{গ) 0 ✔}}

ঘ) √(3)

\textbf{\hl{✔ সঠিক উত্তর: গ) 0}}

\begin{quote}
\textbf{ব্যাখ্যা:} \emph{দেওয়া আছে, √(x+3) = √(x) + √(3)}

\emph{উভয়পক্ষকে বর্গ করে পাই,}

\emph{⇒ ((√(x+3)))\^{}2 = ((√(x)+√(3)))\^{}2}

\emph{⇒ x + 3 = (√(x))\^{}2 + (√(3))\^{}2 + 2 × √(x) × √(3)}

\emph{⇒ x + 3 = x + 3 + 2 √(3x)}

\emph{⇒ 2 √(3x) = x + 3 − x − 3}

\emph{⇒ 2 √(3x) = 0}

\emph{⇒ √(3x) = (0/2) = 0}

\emph{⇒ (√(3x))\^{}2 = (0)\^{}2}

\emph{⇒ 3x = 0}

\emph{∴ x = (0/3) = 0}
\end{quote}

\textbf{9p\^{}2 + 14p এর সাথে কত যোগ করলে যোগফল একটি পূর্ণবর্গ সংখ্যা হবে?}

\textbf{\hl{ক) (49/9) ✔}}

খ) (14/9)

গ) (7/3)

ঘ) 7

\textbf{\hl{✔ সঠিক উত্তর: ক) (49/9)}}

\begin{quote}
\textbf{ব্যাখ্যা:} \emph{দেওয়া আছে, 9p\^{}2 + 14p}

\emph{= (3p)\^{}2 + 2 × 3p × (7/3) + (((7/3)))\^{}2 − (49/9)}

\emph{= ((3p+(7/3)))\^{}2 − (49/9)}

\emph{∴ রাশিটির সাথে (49/9) যোগ করলে যোগফল একটি পূর্ণবর্গ সংখ্যা
((3p+(7/3)))\^{}2 হবে।}
\end{quote}

\textbf{9x\^{}2 + 16y\^{}2 এর সাথে কত যোগ করলে যোগফল পূর্ণবর্গ রাশি হবে?}

ক) 6xy

খ) 12xy

\textbf{\hl{গ) 24xy ✔}}

ঘ) 144xy

\textbf{\hl{✔ সঠিক উত্তর: গ) 24xy}}

\begin{quote}
\textbf{ব্যাখ্যা:} \emph{9x\^{}2 + 16y\^{}2}

\emph{= (3x)\^{}2 + (4y)\^{}2}

\emph{আমরা জানি, a\^{}2 + b\^{}2 + 2ab = (a + b)\^{}2}

\emph{সুত্রানুসারে ৩য় পদটি হতে হবে = 2 · (3x) · (4y) = 24xy}

\emph{∴ রাশিটির সাথে 24xy যোগ করলে এটি একটি পূর্ণবর্গ রাশি (3x + 4y)\^{}2
হবে।}
\end{quote}

\textbf{(x\^{}2/y\^{}2) + (2x/y) এর সাথে কত যোগ করলে সংখ্যাটি পূর্ণ বর্গ
হবে?}

ক) 2xy

খ) 2y\^{}2

\textbf{\hl{গ) 1 ✔}}

ঘ) −1

\textbf{\hl{✔ সঠিক উত্তর: গ) 1}}

\begin{quote}
\textbf{ব্যাখ্যা:} \emph{প্রদত্ত রাশি = (((x/y)))\^{}2 + 2 · ((x/y)) · 1}

\emph{রাশিটিকে (a + b)\^{}2 = a\^{}2 + 2ab + b\^{}2 সূত্রে সাজাতে হলে
b\^{}2 = 1\^{}2 = 1 যোগ করতে হবে।}

\emph{∴ 1 যোগ করলে রাশিটি হবে (((x/y)+1))\^{}2 , যা একটি পূর্ণবর্গ রাশি।}
\end{quote}

\textbf{p এর মান কত হলে, 4x\^{}2 − px + 9 একটি পূর্ণবর্গ হবে?}

ক) 8

\textbf{\hl{খ) 12 ✔}}

গ) 10

ঘ) 15

\textbf{\hl{✔ সঠিক উত্তর: খ) 12}}

\begin{quote}
\textbf{ব্যাখ্যা:} \emph{4x\^{}2 − px + 9}

\emph{= (2x)\^{}2 − px + 3\^{}2}

\emph{রাশিটি পূর্ণবর্গ হতে হলে মাঝের পদটি 2 · (2x) · 3 = 12x হতে হবে।}

\emph{অর্থাৎ, px = 12x}

\emph{∴ p = 12}
\end{quote}

\textbf{16x\^{}2 + px + 25 রাশিটি পূর্ণবর্গ হতে হলে p এর মান কত হবে?}

ক) 20

খ) 10

\textbf{\hl{গ) 40 ✔}}

ঘ) 15

\textbf{\hl{✔ সঠিক উত্তর: গ) 40}}

\begin{quote}
\textbf{ব্যাখ্যা:} \emph{16x\^{}2 + px + 25}

\emph{= (4x)\^{}2 + px + (5)\^{}2}

\emph{রাশিটি পূর্ণবর্গ হতে হলে মাঝের পদটি 2 · (4x) · 5 = 40x হতে হবে।}

\emph{অর্থাৎ, px = 40x}

\emph{∴ p = 40}
\end{quote}

\textbf{(2x/y) এর সাথে কত যোগ করলে যোগফল পূর্ণবর্গ হবে?}

ক) (x\^{}2−y\^{}2/y\^{}2)

খ) (y\^{}2−x\^{}2/y\^{}2)

\textbf{\hl{গ) (x\^{}2+y\^{}2/y\^{}2) ✔}}

ঘ) (x\^{}2+y\^{}2/x\^{}2)

\textbf{\hl{✔ সঠিক উত্তর: গ) (x\^{}2+y\^{}2/y\^{}2)}}

\begin{quote}
\textbf{ব্যাখ্যা:} \emph{(2x/y) + (x\^{}2+y\^{}2/y\^{}2)}

\emph{= (2xy+x\^{}2+y\^{}2/y\^{}2)}

\emph{= (((x+y))\^{}2/y\^{}2)}

\emph{= (((x+y/y)))\^{}2 , যা একটি পূর্ণবর্গ রাশি।}

\emph{∴ (x\^{}2+y\^{}2/y\^{}2) যোগ করলে যোগফল পূর্ণবর্গ হবে।}
\end{quote}

\textbf{x পূর্ণ সংখ্যা হলে 16x\^{}2 + 16x + 2 এর সাথে নূন্যতম কত যোগ করলে এটি
একটি পূর্ণবর্গ সংখ্যা হবে?}

\textbf{\hl{ক) 2 ✔}}

খ) 1

গ) 4

ঘ) 3

\textbf{\hl{✔ সঠিক উত্তর: ক) 2}}

\begin{quote}
\textbf{ব্যাখ্যা:} \emph{16x\^{}2 + 16x + 2}

\emph{= (4x)\^{}2 + 2 · 4x · 2 + 2\^{}2 − 2}

\emph{= (4x + 2)\^{}2 − 2}

\emph{∴ 2 যোগ করলে রাশিটি পূর্ণবর্গ হবে।}
\end{quote}

\textbf{x\^{}2 − 3x + 1 = 0 হলে (x\^{}2−(1/x\^{}2)) এর মান −}

ক) 5 √(3)

\textbf{\hl{খ) 3 √(5) ✔}}

গ) 4 √(5)

ঘ) 6 √(5)

\textbf{\hl{✔ সঠিক উত্তর: খ) 3 √(5)}}

\begin{quote}
\textbf{ব্যাখ্যা:} \emph{x\^{}2 − 3x + 1 = 0}

\emph{⇒ x\^{}2 + 1 = 3x}

\emph{⇒ x + (1/x) = 3 ......... (1)}

\emph{আবার,}

\emph{((x-(1/x)))\^{}2 = ((x+(1/x)))\^{}2 − 4 · x · (1/x)}

\emph{= 3\^{}2 − 4 = 9 − 4 = 5}

\emph{∴ x − (1/x) = √(5)}

\emph{∴ x\^{}2 − (1/x\^{}2) = (x+(1/x)) (x-(1/x)) = 3 √(5)}
\end{quote}

\textbf{X + Y = 6 এবং XY = 6 হলে (x − y)\^{}2 এর মান কত?}

ক) −12

\textbf{\hl{খ) 12 ✔}}

গ) 8

ঘ) 10

\textbf{\hl{✔ সঠিক উত্তর: খ) 12}}

\begin{quote}
\textbf{ব্যাখ্যা:} \emph{দেওয়া আছে, X + Y = 6 এবং XY = 6}

\emph{আমরা জানি, (x − y)\^{}2 = (x + y)\^{}2 − 4xy}

\emph{= (6)\^{}2 − 4 × 6}

\emph{= 36 − 24 = 12}
\end{quote}

\textbf{যদি a = 17 এবং b = 11 হয় তবে a\^{}2 − 2ab + b\^{}2 = কত?}

\textbf{\hl{ক) 36 ✔}}

খ) 100

গ) 289

ঘ) 374

\textbf{\hl{✔ সঠিক উত্তর: ক) 36}}

\begin{quote}
\textbf{ব্যাখ্যা:} \emph{a\^{}2 − 2ab + b\^{}2}

\emph{= (a − b)\^{}2}

\emph{= (17 − 11)\^{}2 {[}মান বসিয়ে{]} = (6)\^{}2 = 36}
\end{quote}

\textbf{(a + b) = 7, ab = 10 হলে (a − b) = ?}

ক) 4

\textbf{\hl{খ) 3 ✔}}

গ) 2

ঘ) 8

\textbf{\hl{✔ সঠিক উত্তর: খ) 3}}

\begin{quote}
\textbf{ব্যাখ্যা:} \emph{আমরা জানি,}

\emph{(a − b)\^{}2 = (a + b)\^{}2 − 4ab}

\emph{⇒ (a − b)\^{}2 = (7)\^{}2 − (4 × 10)}

\emph{⇒ (a − b)\^{}2 = 49 − 40 = 9}

\emph{∴ (a − b) = √(9) = 3}
\end{quote}

\textbf{x\^{}2 + y\^{}2 = 18 এবং xy = 7 হলে (x − y)\^{}2 এর মান কত?}

\textbf{\hl{ক) 4 ✔}}

খ) 5

গ) 6

ঘ) 7

\textbf{\hl{✔ সঠিক উত্তর: ক) 4}}

\begin{quote}
\textbf{ব্যাখ্যা:} \emph{(x − y)\^{}2 = x\^{}2 − 2xy + y\^{}2}

\emph{= (x\^{}2 + y\^{}2) − 2xy}

\emph{= 18 − 2 × 7}

\emph{= 18 − 14 = 4}
\end{quote}

\textbf{a = 15 এবং b = 6 হলে 9a\^{}2 − 48ab + 64b\^{}2-এর মান নির্ণয়
করুন।}

ক) 8

খ) 5

গ) 6

\textbf{\hl{ঘ) 9 ✔}}

\textbf{\hl{✔ সঠিক উত্তর: ঘ) 9}}

\begin{quote}
\textbf{ব্যাখ্যা:} \emph{9a\^{}2 − 48ab + 64b\^{}2}

\emph{= (3a)\^{}2 − 2(3a)(8b) + (8b)\^{}2}

\emph{= (3a − 8b)\^{}2}

\emph{= (3 × 15 − 8 × 6)\^{}2 {[}মান বসিয়ে{]}}

\emph{= (45 − 48)\^{}2}

\emph{= (−3)\^{}2 = 9}
\end{quote}

\textbf{a + b = 9m এবং ab = 18m\^{}2, a − b = ?}

ক) 3m

খ) −3m

গ) 6m

ঘ) 8m

\textbf{\hl{ঙ) ক এবং খ উভয়ই ✔}}

\textbf{\hl{✔ সঠিক উত্তর: ঙ) ক এবং খ উভয়ই}}

\begin{quote}
\textbf{ব্যাখ্যা:} \emph{(a − b)\^{}2 = (a + b)\^{}2 − 4ab}

\emph{= (9m)\^{}2 − 4 × 18m\^{}2}

\emph{= 81m\^{}2 − 72m\^{}2}

\emph{= 9m\^{}2}

\emph{∴ a − b = ±3m}
\end{quote}

\textbf{(x + y)\^{}2 = 164; xy = 32 হলে x − y = কত?}

\textbf{\hl{ক) 6 ✔}}

খ) 9

গ) 4

ঘ) 12

\textbf{\hl{✔ সঠিক উত্তর: ক) 6}}

\begin{quote}
\textbf{ব্যাখ্যা:} \emph{(x − y)\^{}2 = (x + y)\^{}2 − 4xy}

\emph{= 164 − 4 × 32}

\emph{= 164 − 128 = 36}

\emph{∴ x − y = √(36) = 6}
\end{quote}

\textbf{x + y = 17 এবং xy = 60 হলে x − y এর মান কত?}

\textbf{\hl{ক) 7 ✔}}

খ) 8

গ) 9

ঘ) 10

\textbf{\hl{✔ সঠিক উত্তর: ক) 7}}

\begin{quote}
\textbf{ব্যাখ্যা:} \emph{(x − y)\^{}2 = (x + y)\^{}2 − 4xy}

\emph{= (17)\^{}2 − 4 × 60}

\emph{= 289 − 240 = 49}

\emph{∴ x − y = √(49) = 7}
\end{quote}

\textbf{যদি (x − y)\^{}2 = 12, xy = 1 হয়, x\^{}2 + y\^{}2 = কত?}

ক) 11

খ) 12

গ) 13

\textbf{\hl{ঘ) 14 ✔}}

\textbf{\hl{✔ সঠিক উত্তর: ঘ) 14}}

\begin{quote}
\textbf{ব্যাখ্যা:} \emph{দেওয়া আছে, xy = 1 এবং (x − y)\^{}2 = 12}

\emph{⇒ x\^{}2 + y\^{}2 − 2xy = 12}

\emph{⇒ x\^{}2 + y\^{}2 = 12 + 2xy}

\emph{⇒ x\^{}2 + y\^{}2 = 12 + 2 × 1}

\emph{∴ x\^{}2 + y\^{}2 = 14}
\end{quote}

\textbf{x − (1/x) = 7 হলে, x\^{}2 + (1/x\^{}2) এর মান হবে}

ক) 42

খ) 47

\textbf{\hl{গ) 51 ✔}}

ঘ) 54

\textbf{\hl{✔ সঠিক উত্তর: গ) 51}}

\begin{quote}
\textbf{ব্যাখ্যা:} \emph{x\^{}2 + (1/x\^{}2)}

\emph{= ((x-(1/x)))\^{}2 + 2 · x · (1/x) = (7)\^{}2 + 2 = 49 + 2 = 51}
\end{quote}

\textbf{(175×175−2×175×75+75×75/175−75)}

ক) 10

খ) 1000

\textbf{\hl{গ) 100 ✔}}

ঘ) 200

\textbf{\hl{✔ সঠিক উত্তর: গ) 100}}

\begin{quote}
\textbf{ব্যাখ্যা:} \emph{ধরি, a = 175 এবং b = 75}

\emph{প্রদত্ত রাশি = (a\^{}2−2ab+b\^{}2/a−b)}

\emph{= ((a−b)\^{}2/a−b)}

\emph{= a − b}

\emph{= 175 − 75 = 100}
\end{quote}

\textbf{x − (1/x) = 2 হলে, x\^{}4 + (1/x\^{}4) = কত?}

\textbf{\hl{ক) 34 ✔}}

খ) 32

গ) 31

ঘ) 30

\textbf{\hl{✔ সঠিক উত্তর: ক) 34}}

\begin{quote}
\textbf{ব্যাখ্যা:} \emph{x\^{}4 + (1/x\^{}4)}

\emph{= ((x\^{}2+(1/x\^{}2)))\^{}2 − 2 · x\^{}2 · (1/x\^{}2)}

\emph{= (\{((x-(1/x)))\^{}2+2·x·(1/x)\})\^{}2 − 2}

\emph{= (2\^{}2 + 2)\^{}2 − 2}

\emph{= (4 + 2)\^{}2 − 2}

\emph{= 6\^{}2 − 2}

\emph{= 36 − 2 = 34}
\end{quote}

\textbf{((x+(2/x)))\^{}2 = 9 হলে, ((x-(2/x)))\^{}2 এর মান কত?}

\textbf{\hl{ক) 1 ✔}}

খ) 9

গ) 3

ঘ) 0

\textbf{\hl{✔ সঠিক উত্তর: ক) 1}}

\begin{quote}
\textbf{ব্যাখ্যা:} \emph{আমরা জানি,}

\emph{((x-(2/x)))\^{}2 = ((x+(2/x)))\^{}2 − 4 · x · (2/x)}

\emph{= 9 − 4 · 2}

\emph{= 9 − 8 = 1}
\end{quote}

\textbf{4x\^{}2 − 2 হতে কত বিয়োগ করলে রাশিটি পূর্ণবর্গ হবে?}

ক) − (1/4x\^{}2)

খ) − (1/2x\^{}2)

\textbf{\hl{গ) (1/4x\^{}2) ✔}}

ঘ) (1/2x\^{}2)

\textbf{\hl{✔ সঠিক উত্তর: গ) (1/4x\^{}2)}}

\begin{quote}
\textbf{ব্যাখ্যা:} \emph{দেওয়া আছে, 4x\^{}2 − 2}

\emph{= (2x)\^{}2 − 2 · (2x) · (1/2x)}

\emph{এখন, রাশিটিকে (a − b)\^{}2 = a\^{}2 − 2ab + b\^{}2 আকারে পূর্ণবর্গ
করতে হলে b\^{}2 = (((1/2x)))\^{}2 = (1/4x\^{}2) যোগ করতে হবে।}

\emph{অর্থাৎ, (4x\^{}2 − 2) হতে − (1/4x\^{}2) বিয়োগ করলে রাশিটি হয়:}

\emph{(4x\^{}2 − 2) − (− (1/4x\^{}2) ) = ((2x-(1/2x)))\^{}2 , যা একটি
পূর্ণবর্গ রাশি।}

\emph{∴ − (1/4x\^{}2) বিয়োগ করতে হবে (অথবা (1/4x\^{}2) যোগ করতে হবে)।}
\end{quote}

\textbf{9 − 12x + ax\^{}2 একটি পূর্ণসংখ্যা হলে, a এর মান---}

\textbf{\hl{ক) 4 ✔}}

খ) 6

গ) 8

ঘ) −6

\textbf{\hl{✔ সঠিক উত্তর: ক) 4}}

\begin{quote}
\textbf{ব্যাখ্যা:} \emph{9 − 12x + ax\^{}2}

\emph{= 3\^{}2 − 2 · 3 · 2x + (2x)\^{}2}

\emph{= (3 − 2x)\^{}2; যা একটি পূর্ণ বর্গসংখ্যা।}

\emph{প্রশ্নানুসারে, (2x)\^{}2 = ax\^{}2}

\emph{⇒ 4x\^{}2 = ax\^{}2}

\emph{∴ a = 4}
\end{quote}

\textbf{y এর মান কত হলে 16x\^{}2 − xy + 25 একটি পূর্ণবর্গ রাশি হবে?}

ক) 30

\textbf{\hl{খ) 40 ✔}}

গ) 50

ঘ) 60

\textbf{\hl{✔ সঠিক উত্তর: খ) 40}}

\begin{quote}
\textbf{ব্যাখ্যা:} \emph{16x\^{}2 − xy + 25}

\emph{= (4x)\^{}2 − 2 · 4x · 5 + (5)\^{}2}

\emph{= (4x − 5)\^{}2}

\emph{মাঝের পদটি 2 · 4x · 5 = 40x হতে হবে।}

\emph{∴ xy = 40x ⇒ y = 40 হলে রাশিটি পূর্ণবর্গ হয়।}
\end{quote}

\textbf{p এর মান কত হলে 4x\^{}2 − px + 9 একটি পূর্ণবর্গ হবে?}

ক) 10

\textbf{\hl{খ) 12 ✔}}

গ) 9

ঘ) 16

\textbf{\hl{✔ সঠিক উত্তর: খ) 12}}

\begin{quote}
\textbf{ব্যাখ্যা:} \emph{4x\^{}2 − px + 9}

\emph{= (2x)\^{}2 − px + 3\^{}2}

\emph{রাশিটি পূর্ণবর্গ হতে হলে মাঝের পদটি 2 · (2x) · 3 = 12x হতে হবে।}

\emph{অর্থাৎ, px = 12x}

\emph{∴ p = 12}
\end{quote}

\textbf{a- এর মান কত হলে 9 − 12x + ax\^{}2 একটি পূর্ণবর্গ হবে?..}

ক) 8

খ) 6

গ) −6

\textbf{\hl{ঘ) 4 ✔}}

\textbf{\hl{✔ সঠিক উত্তর: ঘ) 4}}

\begin{quote}
\textbf{ব্যাখ্যা:} \emph{ax\^{}2 − 12x + 9 রাশিটি পূর্ণবর্গ হতে হলে
discriminant (নিশ্চায়ক) b\^{}2 − 4ac = 0 হতে হবে।}

\emph{এখানে b = −12, a = a, c = 9}

\emph{∴ (−12)\^{}2 − 4 · a · 9 = 0}

\emph{⇒ 144 − 36a = 0}

\emph{⇒ 36a = 144}

\emph{∴ a = 4}
\end{quote}

\textbf{4x\^{}2 − 20x এর সাথে কত যোগ করলে যোগফল পূর্ণবর্গ হবে?}

ক) 4

খ) 9

গ) 16

\textbf{\hl{ঘ) 25 ✔}}

\textbf{\hl{✔ সঠিক উত্তর: ঘ) 25}}

\begin{quote}
\textbf{ব্যাখ্যা:} \emph{4x\^{}2 − 20x}

\emph{= (2x)\^{}2 − 2 · 2x · 5}

\emph{রাশিটিকে (a − b)\^{}2 = a\^{}2 − 2ab + b\^{}2 সূত্রে সাজাতে হলে
5\^{}2 = 25 যোগ করতে হবে।}

\emph{∴ 25 যোগ করলে রাশিটি হবে (2x − 5)\^{}2, যা একটি পূর্ণবর্গ রাশি।}
\end{quote}

\textbf{x\^{}2 + (1/x\^{}2) − 2x − (2/x) + 2 এর সাথে কত যোগ করলে
রাশিমালাটি একটি পূর্ণ বর্গ সংখ্যা হবে?}

\textbf{\hl{ক) 2 (2−(1/x)−x) ✔}}

খ) 2 (x+(1/x)−2)

গ) (2−(1/x)−x)

ঘ) 2 (1−(1/x)−x)

\textbf{\hl{✔ সঠিক উত্তর: ক) 2 (2−(1/x)−x)}}

\begin{quote}
\textbf{ব্যাখ্যা:} \emph{প্রদত্ত রাশি = x\^{}2 + (1/x\^{}2) − 2x − (2/x) +
2}

\emph{= x\^{}2 − 2x + 1 + (1/x\^{}2) − (2/x) + 1}

\emph{= (x − 1)\^{}2 + (((1/x)−1))\^{}2}

\emph{আমরা জানি, a\^{}2 + b\^{}2 = (a + b)\^{}2 − 2ab}

\emph{= (\{(x−1)+((1/x)−1)\})\^{}2 − 2(x − 1) ((1/x)−1)}

\emph{= ((x+(1/x)−2))\^{}2 − 2 (1−x−(1/x)+1)}

\emph{= ((x+(1/x)−2))\^{}2 − 2 (2−(1/x)−x)}

\emph{∴ প্রদত্ত রাশির সাথে 2 (2−(1/x)−x) যোগ করলে রাশিটি পূর্ণবর্গ হবে।}
\end{quote}

\textbf{2a\^{}2 − 4ab + 4b\^{}2 রাশিটির সাথে কত যোগ বা বিয়োগ করলে রাশিটি
একটি পূর্ণবর্গ হবে?}

ক) 2a\^{}2 যোগ করলে

\textbf{\hl{খ) a\^{}2 বিয়োগ করলে ✔}}

গ) 2b\^{}2 বিয়োগ করলে

ঘ) 2ab যোগ করলে

\textbf{\hl{✔ সঠিক উত্তর: খ) a\^{}2 বিয়োগ করলে}}

\begin{quote}
\textbf{ব্যাখ্যা:} \emph{2a\^{}2 − 4ab + 4b\^{}2 রাশিটি থেকে a\^{}2 বিয়োগ
করলে পাওয়া যায়:}

\emph{(2a\^{}2 − 4ab + 4b\^{}2) − a\^{}2}

\emph{= a\^{}2 − 4ab + 4b\^{}2}

\emph{= a\^{}2 − 2 · a · 2b + (2b)\^{}2}

\emph{= (a − 2b)\^{}2, যা একটি পূর্ণবর্গ রাশি।}

\emph{(এছাড়াও 2b\^{}2 বিয়োগ করলেও এটি (a − b)\^{}2 + a\^{}2 হয়, তবে সঠিক
পূর্ণবর্গ রূপের জন্য a\^{}2 বিয়োগ করাই সঠিক সমাধান।)}
\end{quote}

\textbf{2(a\^{}2 + b\^{}2) = কত?}

ক) (a + b)\^{}2 − 2ab

\textbf{\hl{খ) (a + b)\^{}2 + (a − b)\^{}2 ✔}}

গ) (a − b)\^{}2 + 2ab

ঘ) (a + b)2 − 4ab

\textbf{\hl{✔ সঠিক উত্তর: খ) (a + b)\^{}2 + (a − b)\^{}2}}

\begin{quote}
\textbf{ব্যাখ্যা:} \emph{আমরা জানি বীজগণিতীয় সূত্রানুসারে,}

\emph{(a + b)\^{}2 + (a − b)\^{}2}

\emph{= (a\^{}2 + 2ab + b\^{}2) + (a\^{}2 − 2ab + b\^{}2)}

\emph{= 2a\^{}2 + 2b\^{}2}

\emph{= 2(a\^{}2 + b\^{}2)}
\end{quote}

\textbf{(x/y) + (y/x) = 3 হলে, (x\^{}2/y\^{}2) + (y\^{}2/x\^{}2) এর মান
কত?}

ক) 6

\textbf{\hl{খ) 7 ✔}}

গ) 9

ঘ) 10

\textbf{\hl{✔ সঠিক উত্তর: খ) 7}}

\begin{quote}
\textbf{ব্যাখ্যা:} \emph{(x\^{}2/y\^{}2) + (y\^{}2/x\^{}2)}

\emph{= (((x/y)))\^{}2 + (((y/x)))\^{}2}

\emph{= (((x/y)+(y/x)))\^{}2 − 2 · (x/y) · (y/x)}

\emph{= 3\^{}2 − 2}

\emph{= 9 − 2 = 7}
\end{quote}

\textbf{x − (1/x) = √(3) হলে x\^{}2 + (1/x\^{}2) = কত?}

\textbf{\hl{ক) ৫ ✔}}

খ) ১

গ) ১১

ঘ) ৭

\textbf{\hl{✔ সঠিক উত্তর: ক) ৫}}

\begin{quote}
\textbf{ব্যাখ্যা:} \emph{x\^{}2 + (1/x\^{}2)}

\emph{= ((x-(1/x)))\^{}2 + 2 · x · (1/x) = ( √(3) )\^{}2 + 2 = 3 + 2 =
5}
\end{quote}

\textbf{a + (1/a) = 4 হয়, তাহলে a\^{}2 + (1/a\^{}2) =কত?}

ক) 16

খ) 9

গ) 12

\textbf{\hl{ঘ) 14 ✔}}

\textbf{\hl{✔ সঠিক উত্তর: ঘ) 14}}

\begin{quote}
\textbf{ব্যাখ্যা:} \emph{a\^{}2 + (1/a\^{}2)}

\emph{= ((a+(1/a)))\^{}2 − 2 · a · (1/a)}

\emph{= (4)\^{}2 − 2}

\emph{= 16 − 2 = 14}
\end{quote}

\textbf{p + (1/p) = 4 হলে p\^{}4 + (1/p\^{}4) = কত?}

\textbf{\hl{ক) 194 ✔}}

খ) 119

গ) 147

ঘ) 149

\textbf{\hl{✔ সঠিক উত্তর: ক) 194}}

\begin{quote}
\textbf{ব্যাখ্যা:} \emph{p\^{}4 + (1/p\^{}4)}

\emph{= (p\^{}2)\^{}2 + (((1/p\^{}2)))\^{}2}

\emph{= ((p\^{}2+(1/p\^{}2)))\^{}2 − 2 · p\^{}2 · (1/p\^{}2)}

\emph{= (\{((p+(1/p)))\^{}2−2·p·(1/p)\})\^{}2 − 2}

\emph{= \{(4)\^{}2 − 2\}\^{}2 − 2}

\emph{= (16 − 2)\^{}2 − 2}

\emph{= (14)\^{}2 − 2}

\emph{= 196 − 2 = 194}
\end{quote}

\textbf{x + (1/x) = 4 হলে x\^{}4 + (1/x\^{}4) এর মান কত??}

ক) 149

\textbf{\hl{খ) 194 ✔}}

গ) 119

ঘ) 147

\textbf{\hl{✔ সঠিক উত্তর: খ) 194}}

\begin{quote}
\textbf{ব্যাখ্যা:} \emph{x\^{}4 + (1/x\^{}4)}

\emph{= ((x\^{}2+(1/x\^{}2)))\^{}2 − 2 · x\^{}2 · (1/x\^{}2)}

\emph{= (\{((x+(1/x)))\^{}2−2·x·(1/x)\})\^{}2 − 2}

\emph{= (4\^{}2 − 2)\^{}2 − 2}

\emph{= (16 − 2)\^{}2 − 2}

\emph{= (14)\^{}2 − 2}

\emph{= 196 − 2 = 194}
\end{quote}

\textbf{a − (1/a) = 5 হলে a\^{}2 + (1/a\^{}2) এর মান কত?}

ক) 20

খ) 23

গ) 25

\textbf{\hl{ঘ) 27 ✔}}

\textbf{\hl{✔ সঠিক উত্তর: ঘ) 27}}

\begin{quote}
\textbf{ব্যাখ্যা:} \emph{a\^{}2 + (1/a\^{}2)}

\emph{= ((a−(1/a)))\^{}2 + 2 · a · (1/a) = 5\^{}2 + 2 = 25 + 2 = 27}
\end{quote}

\textbf{x − (1/x) = 2 হলে x\^{}4 + (1/x\^{}4) = কত?}

ক) 30

খ) 31

গ) 32

\textbf{\hl{ঘ) 34 ✔}}

\textbf{\hl{✔ সঠিক উত্তর: ঘ) 34}}

\begin{quote}
\textbf{ব্যাখ্যা:} \emph{x\^{}4 + (1/x\^{}4)}

\emph{= (x\^{}2)\^{}2 + (((1/x\^{}2)))\^{}2}

\emph{= ((x\^{}2+(1/x\^{}2)))\^{}2 − 2 · x\^{}2 · (1/x\^{}2)}

\emph{= (\{((x-(1/x)))\^{}2+2·x·(1/x)\})\^{}2 − 2}

\emph{= \{(2)\^{}2 + 2\}\^{}2 − 2}

\emph{= (4 + 2)\^{}2 − 2}

\emph{= 6\^{}2 − 2}

\emph{= 36 − 2 = 34}
\end{quote}

\textbf{(a/b) + (b/a) = 6 হলে (a\^{}2/b\^{}2) + (b\^{}2/a\^{}2) + 2 এর
মান কত?}

ক) 40

\textbf{\hl{খ) 36 ✔}}

গ) 32

ঘ) 30

\textbf{\hl{✔ সঠিক উত্তর: খ) 36}}

\begin{quote}
\textbf{ব্যাখ্যা:} \emph{(a\^{}2/b\^{}2) + (b\^{}2/a\^{}2) + 2}

\emph{= (((a/b)))\^{}2 + 2 · (a/b) · (b/a) + (((b/a)))\^{}2}

\emph{= (((a/b)+(b/a)))\^{}2 = (6)\^{}2 = 36}
\end{quote}

\textbf{2x + (2/x) = 3 হলে x\^{}2 + (1/x\^{}2) −এর মান কত?}

ক) (1/6)

\textbf{\hl{খ) (1/4) ✔}}

গ) (1/3)

ঘ) (1/5)

\textbf{\hl{✔ সঠিক উত্তর: খ) (1/4)}}

\begin{quote}
\textbf{ব্যাখ্যা:} \emph{2x + (2/x) = 3}

\emph{⇒ 2 (x+(1/x)) = 3}

\emph{⇒ x + (1/x) = (3/2)}

\emph{এখন, x\^{}2 + (1/x\^{}2) = ((x+(1/x)))\^{}2 − 2 · x · (1/x)}

\emph{= (((3/2)))\^{}2 − 2}

\emph{= (9/4) − 2}

\emph{= (9−8/4)}

\emph{= (1/4)}
\end{quote}

\textbf{(1/2) \{(a+b)\^{}2 + (a−b)\^{}2\} = ?}

\textbf{\hl{ক) a\^{}2 + b\^{}2 ✔}}

খ) a\^{}2 − b\^{}2

গ) (((a+b))\^{}2/2) − ((a−b)\^{}2/2)

ঘ) (a+b)\^{}2 + (a−b)\^{}2

\textbf{\hl{✔ সঠিক উত্তর: ক) a\^{}2 + b\^{}2}}

\begin{quote}
\textbf{ব্যাখ্যা:} \emph{(1/2) \{(a+b)\^{}2 + (a−b)\^{}2\}}

\emph{= (1/2) \{(a\^{}2 + 2ab + b\^{}2) + (a\^{}2 − 2ab + b\^{}2)\}}

\emph{= (1/2) (2a\^{}2 + 2b\^{}2)}

\emph{= (1/2) · 2(a\^{}2 + b\^{}2)}

\emph{= a\^{}2 + b\^{}2}
\end{quote}

\textbf{x + y = 8, x − y = 6 হলে, x\^{}2 + y\^{}2 এর মান কত?}

ক) 40

খ) 60

\textbf{\hl{গ) 50 ✔}}

ঘ) 80

\textbf{\hl{✔ সঠিক উত্তর: গ) 50}}

\begin{quote}
\textbf{ব্যাখ্যা:} \emph{আমরা জানি,}

\emph{x\^{}2 + y\^{}2 = (((x+y))\^{}2+(x−y)\^{}2/2)}

\emph{= (8\^{}2+6\^{}2/2)}

\emph{= (64+36/2)}

\emph{= (100/2) = 50}
\end{quote}

\textbf{যদি (x−y)\^{}2 = 12 এবং xy = 1 হয়, তবে x\^{}2 + y\^{}2 = কত?}

ক) 11

খ) 12

গ) 13

\textbf{\hl{ঘ) 14 ✔}}

\textbf{\hl{✔ সঠিক উত্তর: ঘ) 14}}

\begin{quote}
\textbf{ব্যাখ্যা:} \emph{x\^{}2 + y\^{}2 = (x − y)\^{}2 + 2xy = 12 + 2 × 1
= 12 + 2 = 14}
\end{quote}

\textbf{যদি a + b = √(5) এবং a − b = √(3) তবে a\^{}2 + b\^{}2 = কত?}

\textbf{\hl{ক) 4 ✔}}

খ) 4 √(2)

গ) 6

ঘ) √(8)

\textbf{\hl{✔ সঠিক উত্তর: ক) 4}}

\begin{quote}
\textbf{ব্যাখ্যা:} \emph{আমরা জানি, 2(a\^{}2 + b\^{}2) = (a + b)\^{}2 + (a
− b)\^{}2}

\emph{⇒ 2(a\^{}2 + b\^{}2) = ( √(5) )\^{}2 + ( √(3) )\^{}2}

\emph{⇒ 2(a\^{}2 + b\^{}2) = 5 + 3 = 8}

\emph{∴ a\^{}2 + b\^{}2 = (8/2) = 4}
\end{quote}

\textbf{(a − b) = 3 এবং ab = 4 হলে a\^{}2 + b\^{}2 এর মান কত?}

ক) 25

\textbf{\hl{খ) 17 ✔}}

গ) 11

ঘ) 1

\textbf{\hl{✔ সঠিক উত্তর: খ) 17}}

\begin{quote}
\textbf{ব্যাখ্যা:} \emph{a\^{}2 + b\^{}2 = (a − b)\^{}2 + 2ab = (3)\^{}2 +
2 × 4 = 9 + 8 = 17}
\end{quote}

\textbf{a + b = 12 এবং ab = 35 হলে a\^{}2 + b\^{}2−এর মান কত?}

ক) 214

\textbf{\hl{খ) 74 ✔}}

গ) 49

ঘ) 24

\textbf{\hl{✔ সঠিক উত্তর: খ) 74}}

\begin{quote}
\textbf{ব্যাখ্যা:} \emph{a\^{}2 + b\^{}2 = (a + b)\^{}2 − 2ab}

\emph{= (12)\^{}2 − 2 × 35}

\emph{= 144 − 70 = 74}
\end{quote}

\textbf{যদি a\^{}2 + (1/a\^{}2) = 51 হবে তবে a − (1/a) এর মান কত??}

ক) ±9

\textbf{\hl{খ) ±7 ✔}}

গ) ±5

ঘ) ±3

\textbf{\hl{✔ সঠিক উত্তর: খ) ±7}}

\begin{quote}
\textbf{ব্যাখ্যা:} \emph{((a−(1/a)))\^{}2 = a\^{}2 + (1/a\^{}2) − 2 · a ·
(1/a)}

\emph{= 51 − 2 = 49}

\emph{∴ a − (1/a) = ±7}
\end{quote}

\textbf{a + a\^{}-1 = √(3) হলে, a\^{}2 + a\^{}-2-এর মান---}

ক) 3

খ) 2

\textbf{\hl{গ) 1 ✔}}

ঘ) 0

\textbf{\hl{✔ সঠিক উত্তর: গ) 1}}

\begin{quote}
\textbf{ব্যাখ্যা:} \emph{a + a\^{}-1 = √(3)}

\emph{∴ a + (1/a) = √(3)}

\emph{এখানে,}

\emph{a\^{}2 + a\^{}-2 = a\^{}2 + (1/a\^{}2)}

\emph{= (a)\^{}2 + (((1/a)))\^{}2}

\emph{= ((a+(1/a)))\^{}2 − 2 · a · (1/a)}

\emph{= ( √(3) )\^{}2 − 2}

\emph{= 3 − 2 = 1}

\emph{∴ a\^{}2 + a\^{}-2 = 1}
\end{quote}

\textbf{x\^{}4 + 2x\^{}2 + 1 = 5x\^{}2 হলে x + (1/x) = কত?}

ক) √(3)

খ) 5

গ) 2

\textbf{\hl{ঘ) √(5) ✔}}

\textbf{\hl{✔ সঠিক উত্তর: ঘ) √(5)}}

\begin{quote}
\textbf{ব্যাখ্যা:} \emph{x\^{}4 + 2x\^{}2 + 1 = 5x\^{}2}

\emph{⇒ (x\^{}2 + 1)\^{}2 = 5x\^{}2}

\emph{⇒ x\^{}2 + 1 = √(5x\^{}2)}

\emph{⇒ x\^{}2 + 1 = √(5) x}

\emph{উভয়পক্ষকে x দ্বারা ভাগ করে,}

\emph{⇒ (x\^{}2+1/x) = √(5)}

\emph{∴ x + (1/x) = √(5)}
\end{quote}

\textbf{x\^{}2+1 = √(2) x হলে x\^{}2 + (1/x\^{}2) = কত?}

ক) 2

খ) 1

\textbf{\hl{গ) 0 ✔}}

ঘ) কোনোটিই নয়

\textbf{\hl{✔ সঠিক উত্তর: গ) 0}}

\begin{quote}
\textbf{ব্যাখ্যা:} \emph{x\^{}2 + 1 = √(2) x}

\emph{⇒ (x\^{}2+1/x) = √(2)}

\emph{∴ x + (1/x) = √(2)}

\emph{এখন, x\^{}2 + (1/x\^{}2) = (x)\^{}2 + (((1/x)))\^{}2}

\emph{= ((x+(1/x)))\^{}2 − 2 · x · (1/x)}

\emph{= ( √(2) )\^{}2 − 2}

\emph{= 2 − 2 = 0}
\end{quote}

\textbf{x\^{}4 − x\^{}2 − 1 = 0 হলে x\^{}2 − (1/x\^{}2) = কত?}

ক) 0

\textbf{\hl{খ) 1 ✔}}

গ) 2

ঘ) 4

\textbf{\hl{✔ সঠিক উত্তর: খ) 1}}

\begin{quote}
\textbf{ব্যাখ্যা:} \emph{x\^{}4 − x\^{}2 − 1 = 0}

\emph{⇒ x\^{}4 − 1 = x\^{}2}

\emph{⇒ (x\^{}4−1/x\^{}2) = 1}

\emph{⇒ (x\^{}4/x\^{}2) − (1/x\^{}2) = 1}

\emph{∴ x\^{}2 − (1/x\^{}2) = 1}
\end{quote}

\textbf{x\^{}4 − x\^{}2 + 1 = 0 হলে, x\^{}2 + (1/x\^{}2) এর মান কত?}

\textbf{\hl{ক) 1 ✔}}

খ) √(2)

গ) 2

ঘ) 3

\textbf{\hl{✔ সঠিক উত্তর: ক) 1}}

\begin{quote}
\textbf{ব্যাখ্যা:} \emph{x\^{}4 − x\^{}2 + 1 = 0}

\emph{⇒ x\^{}4 + 1 = x\^{}2}

\emph{⇒ (x\^{}4+1/x\^{}2) = 1}

\emph{∴ x\^{}2 + (1/x\^{}2) = 1}
\end{quote}

\textbf{x = √(3) + √(2) হলে, x\^{}2 + (1/x\^{}2) এর মান কত?}

ক) 9

\textbf{\hl{খ) 10 ✔}}

গ) 12

ঘ) কোনটিই নয়

\textbf{\hl{✔ সঠিক উত্তর: খ) 10}}

\begin{quote}
\textbf{ব্যাখ্যা:} \emph{দেওয়া আছে, x = √(3) + √(2)}

\emph{∴ (1/x) = (1/√(3)+√(2)) = (√(3)−√(2)/(√(3)+√(2))(√(3)−√(2)))}

\emph{= (√(3)−√(2)/3−2) = √(3) − √(2)}

\emph{∴ x + (1/x) = √(3) + √(2) + √(3) − √(2) = 2 √(3)}

\emph{∴ x\^{}2 + (1/x\^{}2) = ((x+(1/x)))\^{}2 − 2 · x · (1/x)}

\emph{= (2 √(3) )\^{}2 − 2}

\emph{= 12 − 2 = 10}
\end{quote}

\textbf{x − 2 = √(3) হলে x\^{}4 + (1/x\^{}4) এর মান কত?}

ক) 196

\textbf{\hl{খ) 194 ✔}}

গ) 192

ঘ) 198

\textbf{\hl{✔ সঠিক উত্তর: খ) 194}}

\begin{quote}
\textbf{ব্যাখ্যা:} \emph{x − 2 = √(3)}

\emph{⇒ x = 2 + √(3)}

\emph{∴ (1/x) = (1/2+√(3)) = (2−√(3)/(2+√(3))(2−√(3))) = 2 − √(3)}

\emph{∴ x + (1/x) = (2 + √(3) ) + (2 − √(3) ) = 4}

\emph{এখন, x\^{}2 + (1/x\^{}2) = ((x+(1/x)))\^{}2 − 2 = 4\^{}2 − 2 = 14}

\emph{∴ x\^{}4 + (1/x\^{}4) = ((x\^{}2+(1/x\^{}2)))\^{}2 − 2 = 14\^{}2 −
2 = 196 − 2 = 194}
\end{quote}

\textbf{a + b = 4, a − b = 2 হলে ab এর মান কত?}

ক) 8

\textbf{\hl{খ) 3 ✔}}

গ) 2

ঘ) 0

\textbf{\hl{✔ সঠিক উত্তর: খ) 3}}

\begin{quote}
\textbf{ব্যাখ্যা:} \emph{দেওয়া আছে, a + b = 4 ও a − b = 2}

\emph{আমরা জানি, 4ab = (a + b)\^{}2 − (a − b)\^{}2}

\emph{= 4\^{}2 − 2\^{}2}

\emph{= 16 − 4 = 12}

\emph{∴ ab = (12/4) = 3}
\end{quote}

\textbf{a+b=8, a-b = 4 হলে ab = কত?}

ক) ৯

খ) ১০

\textbf{\hl{গ) ১২ ✔}}

ঘ) ১৮

\textbf{\hl{✔ সঠিক উত্তর: গ) ১২}}

\begin{quote}
\textbf{ব্যাখ্যা:} \emph{দেওয়া আছে, a+b=8 ও a-b=4}

\emph{আমরা জানি, 4ab = (a+b)\^{}2 − (a-b)\^{}2}

\emph{= 8\^{}2 − 4\^{}2}

\emph{= 64 − 16 = 48}

\emph{∴ ab = (48/4) = 12}
\end{quote}

\textbf{ab = কত?}

ক) (((a+b))\^{}2/2) − ((a−b)\^{}2/2)

\textbf{\hl{খ) (((a+b/2)))\^{}2 − (((a−b/2)))\^{}2 ✔}}

গ) (a\^{}2+b\^{}2/2) − (a\^{}2−b\^{}2/2)

ঘ) (((a+b/2)))\^{}2 + (((a−b/2)))\^{}2

\textbf{\hl{✔ সঠিক উত্তর: খ) (((a+b/2)))\^{}2 − (((a−b/2)))\^{}2}}

\begin{quote}
\textbf{ব্যাখ্যা:} \emph{আমরা জানি, ab = (((a+b/2)))\^{}2 −
(((a−b/2)))\^{}2}
\end{quote}

\textbf{4ab এর সঠিক প্রয়োগ কোনটি?}

\textbf{\hl{ক) (a + b)\^{}2 − (a − b)\^{}2 ✔}}

খ) (a + b)\^{}2 + (a − b)\^{}2

গ) (((a+b))\^{}2/2) + ((a−b)\^{}2/2)

ঘ) (((a+b/2)))\^{}2 − (((a−b/2)))\^{}2

\textbf{\hl{✔ সঠিক উত্তর: ক) (a + b)\^{}2 − (a − b)\^{}2}}

\begin{quote}
\textbf{ব্যাখ্যা:} \emph{4ab এর সঠিক অনুসিদ্ধান্ত হলো: 4ab = (a + b)\^{}2 −
(a − b)\^{}2}
\end{quote}

\textbf{xy-কে দুটি বর্গের অন্তরফল রূপে প্রকাশ করলে পাওয়া যাবে-}

ক) xy = (((x+y))\^{}2/2) − ((x−y)\^{}2/2)

খ) xy = ((x−y)\^{}2/2) − (((x+y))\^{}2/2)

গ) xy = (x + y)\^{}2 − (x − y)\^{}2

\textbf{\hl{ঘ) xy = (((x+y/2)))\^{}2 − (((x−y/2)))\^{}2 ✔}}

\textbf{\hl{✔ সঠিক উত্তর: ঘ) xy = (((x+y/2)))\^{}2 − (((x−y/2)))\^{}2}}

\begin{quote}
\textbf{ব্যাখ্যা:} \emph{যেকোনো দুটি রাশির গুণফলকে (xy) দুটি বর্গের অন্তরফল বা
বিয়োগফল রূপে প্রকাশের সূত্র হলো:}

\emph{xy = (((x+y/2)))\^{}2 − (((x−y/2)))\^{}2}
\end{quote}

\textbf{(6a + 9b) (7b − 8a) রাশি দুটি বর্গের অন্তর কত?}

ক) (8b + a)\^{}2 − (b − 7a)\^{}2

\textbf{\hl{খ) (8b − a)\^{}2 − (b + 7a)\^{}2 ✔}}

গ) (8b − a)\^{}2 + (b + 7a)\^{}2

ঘ) (8b + a)\^{}2 + (b − 7a)\^{}2

\textbf{\hl{✔ সঠিক উত্তর: খ) (8b − a)\^{}2 − (b + 7a)\^{}2}}

\begin{quote}
\textbf{ব্যাখ্যা:} \emph{ধরি, x = (6a + 9b) এবং y = (7b − 8a)}

\emph{আমরা জানি, xy = (((x+y/2)))\^{}2 − (((x−y/2)))\^{}2}

\emph{= (\{((6a+9b)+(7b−8a)/2)\})\^{}2 − (\{((6a+9b)−(7b−8a)/2)\})\^{}2}

\emph{= (((16b−2a/2)))\^{}2 − (((14a+2b/2)))\^{}2}

\emph{= (\{(2(8b−a)/2)\})\^{}2 − (\{(2(b+7a)/2)\})\^{}2}

\emph{= (8b − a)\^{}2 − (b + 7a)\^{}2}
\end{quote}

\textbf{a + b = ৫ এবং a − b = ৩ হলে ab এর মান কত?}

ক) ৬

\textbf{\hl{খ) ৪ ✔}}

গ) ৮

ঘ) ১০

\textbf{\hl{✔ সঠিক উত্তর: খ) ৪}}

\begin{quote}
\textbf{ব্যাখ্যা:} \emph{আমরা জানি, ab = (((a+b/2)))\^{}2 −
(((a−b/2)))\^{}2}

\emph{= (((5/2)))\^{}2 − (((3/2)))\^{}2}

\emph{= (25/4) − (9/4)}

\emph{= (25−9/4)}

\emph{= (16/4) = 4}
\end{quote}

\textbf{(a+b)\^{}2−(a−b)\^{}2 = কত ?}

ক) 2a\^{}2 +4ab +2b\^{}2

\textbf{\hl{খ) 4ab ✔}}

গ) 2a\^{}2 − 4ab +2b\^{}2

ঘ) 2a\^{}2 +2b\^{}2

\textbf{\hl{✔ সঠিক উত্তর: খ) 4ab}}

\begin{quote}
\textbf{ব্যাখ্যা:} \emph{(a+b)\^{}2 − (a−b)\^{}2}

\emph{= (a\^{}2 + 2ab + b\^{}2) − (a\^{}2 − 2ab + b\^{}2)}

\emph{= a\^{}2 + 2ab + b\^{}2 − a\^{}2 + 2ab − b\^{}2}

\emph{= 4ab}
\end{quote}

\textbf{x + y = √(7) এবং y = x − √(3) হলে xy এর মান কত?}

ক) ৩

খ) ৪

গ) ২

\textbf{\hl{ঘ) ১ ✔}}

\textbf{\hl{✔ সঠিক উত্তর: ঘ) ১}}

\begin{quote}
\textbf{ব্যাখ্যা:} \emph{x + y = √(7)}

\emph{এবং y = x − √(3)}

\emph{⇒ x − y = √(3)}

\emph{আমরা জানি, xy = (((x+y))\^{}2−(x−y)\^{}2/4)}

\emph{= ((√(7))\^{}2−(√(3))\^{}2/4)}

\emph{= (7−3/4)}

\emph{= (4/4) = 1}
\end{quote}

\textbf{x + y = 6 এবং x − y = 4 হলে xy এর মান কত?}

ক) 4

\textbf{\hl{খ) 5 ✔}}

গ) 6

ঘ) 3

\textbf{\hl{✔ সঠিক উত্তর: খ) 5}}

\begin{quote}
\textbf{ব্যাখ্যা:} \emph{4xy = (x + y)\^{}2 − (x − y)\^{}2}

\emph{= (6)\^{}2 − (4)\^{}2}

\emph{= 36 − 16 = 20}

\emph{⇒ xy = (20/4)}

\emph{∴ xy = 5}
\end{quote}

\textbf{a + b = √(7) এবং a − b = √(5) এবং, 8ab(a\^{}2 + b\^{}2) = কত?}

ক) 12

\textbf{\hl{খ) 24 ✔}}

গ) 36

ঘ) 40

\textbf{\hl{✔ সঠিক উত্তর: খ) 24}}

\begin{quote}
\textbf{ব্যাখ্যা:} \emph{দেয়া আছে,}

\emph{a + b = √(7) এবং a − b = √(5)}

\emph{∴ 8ab(a\^{}2 + b\^{}2)}

\emph{= 4ab × 2(a\^{}2 + b\^{}2)}

\emph{= \{(a + b)\^{}2 − (a − b)\^{}2\} × \{(a + b)\^{}2 + (a −
b)\^{}2\}}

\emph{= \{( √(7) )\^{}2 − ( √(5) )\^{}2\} × \{( √(7) )\^{}2 + ( √(5)
)\^{}2\}}

\emph{= (7 − 5) × (7 + 5) = 2 × 12 = 24}
\end{quote}

\textbf{(a + b)\^{}2 = 36 এবং (a − b)\^{}2 = 16 হলে ab এবং (a\^{}2 +
b\^{}2) এর মান যথাক্রমে---}

\textbf{\hl{ক) ৫,২৬ ✔}}

খ) ৫,১৩

গ) ১৩,৫

ঘ) ৫,৫

\textbf{\hl{✔ সঠিক উত্তর: ক) ৫,২৬}}

\begin{quote}
\textbf{ব্যাখ্যা:} \emph{ab = (((a+b))\^{}2−(a−b)\^{}2/4) = (36−16/4) = 5}

\emph{আবার, a\^{}2 + b\^{}2 = (a + b)\^{}2 − 2ab = 36 − 10 = 26}
\end{quote}

\textbf{If x + y = a and x − y = b, then 2xy = ? / x+y=a এবং x−y =b হলে,
2xy---এর মান কত?}

ক) ((a−b)\^{}2/2)

খ) (b\^{}2−a\^{}2/2)

গ) (a−b/2)

ঘ) ((ab)\^{}2/2)

\textbf{\hl{ঙ) (a\^{}2−b\^{}2/2) ✔}}

\textbf{\hl{✔ সঠিক উত্তর: ঙ) (a\^{}2−b\^{}2/2)}}

\begin{quote}
\textbf{ব্যাখ্যা:} \emph{4xy = (x + y)\^{}2 − (x − y)\^{}2 = a\^{}2 −
b\^{}2}

\emph{∴ 2xy = (a\^{}2−b\^{}2/2)}
\end{quote}

\textbf{x\^{}2 − 3x + 1 = 0 হলে (x\^{}2−(1/x\^{}2)) এর মান ---}

ক) 5 √(3)

\textbf{\hl{খ) 3 √(5) ✔}}

গ) 4 √(5)

ঘ) 6 √(5)

\textbf{\hl{✔ সঠিক উত্তর: খ) 3 √(5)}}

\begin{quote}
\textbf{ব্যাখ্যা:} \emph{x\^{}2 − 3x + 1 = 0}

\emph{⇒ x\^{}2 + 1 = 3x}

\emph{⇒ x + (1/x) = 3}

\emph{আবার, ((x−(1/x)))\^{}2 = ((x+(1/x)))\^{}2 − 4 · x · (1/x)}

\emph{= 3\^{}2 − 4 = 5}

\emph{∴ x − (1/x) = √(5)}

\emph{∴ x\^{}2 − (1/x\^{}2) = (x+(1/x))(x−(1/x)) = 3 √(5)}
\end{quote}

\textbf{a\^{}2 − b\^{}2 এর সমান}

ক) (a + b)/(a − b)

\textbf{\hl{খ) (a + b) × (a − b) ✔}}

গ) (a − b)/(a − b)

ঘ) ab

\textbf{\hl{✔ সঠিক উত্তর: খ) (a + b) × (a − b)}}

\begin{quote}
\textbf{ব্যাখ্যা:} \emph{আমরা জানি, a\^{}2 − b\^{}2 = (a + b)(a − b)}
\end{quote}

\textbf{a + b = 3 এবং ab = 1 হলে a\^{}2 − ab + b\^{}2 এর মান কত?}

ক) 4

\textbf{\hl{খ) 6 ✔}}

গ) 8

ঘ) 10

\textbf{\hl{✔ সঠিক উত্তর: খ) 6}}

\begin{quote}
\textbf{ব্যাখ্যা:} \emph{a\^{}2 − ab + b\^{}2 = a\^{}2 + 2ab + b\^{}2 −
3ab}

\emph{= (a + b)\^{}2 − 3ab}

\emph{= 3\^{}2 − 3(1)}

\emph{= 9 − 3 = 6}
\end{quote}

\textbf{a + b = 5 এবং a − b = 3 হলে a\^{}2 − b\^{}2 এর মান কত?}

ক) 2

খ) 8

গ) 16

\textbf{\hl{ঘ) 15 ✔}}

\textbf{\hl{✔ সঠিক উত্তর: ঘ) 15}}

\begin{quote}
\textbf{ব্যাখ্যা:} \emph{a\^{}2 − b\^{}2 = (a + b)(a − b) = 5 × 3 = 15}
\end{quote}

\textbf{a − b = 2; ab = 3 হলে a\^{}2 − b\^{}2 = ?}

ক) −10

খ) −1

\textbf{\hl{গ) 8 ✔}}

ঘ) 26

\textbf{\hl{✔ সঠিক উত্তর: গ) 8}}

\begin{quote}
\textbf{ব্যাখ্যা:} \emph{(a + b)\^{}2 = (a − b)\^{}2 + 4ab = 2\^{}2 + 4(3)
= 4 + 12 = 16}

\emph{∴ a + b = √(16) = 4 (বা ±4)}

\emph{সুতরাং, a\^{}2 − b\^{}2 = (a + b)(a − b) = 4 × 2 = 8}
\end{quote}

\textbf{a − b = 2 ও ab=24 হলে a\^{}2−b\^{}2 সমান কত?}

ক) ১৬

\textbf{\hl{খ) ২০ ✔}}

গ) ২৬

ঘ) ৩০

\textbf{\hl{✔ সঠিক উত্তর: খ) ২০}}

\begin{quote}
\textbf{ব্যাখ্যা:} \emph{আমরা জানি, a + b = √((a−b)\^{}2+4ab) =
√((2)\^{}2+4×24) = √(4+96) = √(100)}

\emph{∴ a + b = 10}

\emph{∴ a\^{}2 − b\^{}2 = (a + b)(a − b) = 10 × 2 = 20}
\end{quote}

\textbf{a − (1/a) = √(3) হলে, a\^{}2 − (1/a\^{}2) এর মান কত?}

ক) 6

খ) 4

গ) 1

\textbf{\hl{ঘ) √(21) ✔}}

\textbf{\hl{✔ সঠিক উত্তর: ঘ) √(21)}}

\begin{quote}
\textbf{ব্যাখ্যা:} \emph{a − (1/a) = √(3)}

\emph{∴ a + (1/a) = √(((a−(1/a)))\^{}2+4·a·(1/a)) = √(3+4) = √(7)}

\emph{∴ a\^{}2 − (1/a\^{}2) = (a+(1/a))(a−(1/a)) = √(7) × √(3) = √(21)}
\end{quote}

\textbf{If x\^{}2− √(5) x+1=0 then x\^{}2 − (1/x\^{}2) = ?}

ক) 1

\textbf{\hl{খ) √(5) ✔}}

গ) 3

ঘ) 2 √(5)

\textbf{\hl{✔ সঠিক উত্তর: খ) √(5)}}

\begin{quote}
\textbf{ব্যাখ্যা:} \emph{x\^{}2 − √(5) x + 1 = 0}

\emph{⇒ x\^{}2 + 1 = √(5) x}

\emph{⇒ (x\^{}2+1/x) = √(5)}

\emph{⇒ x + (1/x) = √(5)}

\emph{আবার, ((x−(1/x)))\^{}2 = ((x+(1/x)))\^{}2 − 4 = ( √(5) )\^{}2 − 4
= 1}

\emph{∴ x − (1/x) = 1}

\emph{∴ x\^{}2 − (1/x\^{}2) = (x+(1/x))(x−(1/x)) = √(5) × 1 = √(5)}
\end{quote}

\textbf{(a\^{}2+b\^{}2−c\^{}2+2ab/a\^{}2−b\^{}2+c\^{}2+2ac) = কত?}

ক) a + b + c

\textbf{\hl{খ) ((a+b−c)/(a−b+c)) ✔}}

গ) ((a−b+c)/(a+b−c))

ঘ) ((a+b−c)/(a+b+c))

\textbf{\hl{✔ সঠিক উত্তর: খ) ((a+b−c)/(a−b+c))}}

\begin{quote}
\textbf{ব্যাখ্যা:}
\emph{(a\^{}2+b\^{}2−c\^{}2+2ab/a\^{}2−b\^{}2+c\^{}2+2ac)}

\emph{= (a\^{}2+2ab+b\^{}2−c\^{}2/a\^{}2+2ac+c\^{}2−b\^{}2)}

\emph{= (((a+b))\^{}2−c\^{}2/((a+c))\^{}2−b\^{}2)}

\emph{= ((a+b+c)(a+b−c)/(a+c+b)(a+c−b))}

\emph{= ((a+b−c)/(a−b+c))}
\end{quote}

\textbf{a\^{}2 + b\^{}2 + c\^{}2 + 2ab − 2bc এর সাথে কত যোগ করলে যোগফল
পূর্ণবর্গ হবে?}

ক) −2ab

খ) 2bc

গ) 2ca

\textbf{\hl{ঘ) −2ca ✔}}

\textbf{\hl{✔ সঠিক উত্তর: ঘ) −2ca}}

\begin{quote}
\textbf{ব্যাখ্যা:} \emph{a\^{}2 + b\^{}2 + c\^{}2 + 2ab − 2bc − 2ca}

\emph{= a\^{}2 + b\^{}2 + (−c)\^{}2 + 2ab + 2b(−c) + 2(−c)a}

\emph{= (a + b − c)\^{}2}

\emph{∴ −2ca যোগ করলে যোগফল পূর্ণবর্গ হবে।}
\end{quote}

\textbf{x\^{}2+4y\^{}2+8x−16y+16 এর সাথে কত যোগ করলে যোগফল পূর্ণবর্গ হবে?}

\textbf{\hl{ক) −4xy ✔}}

খ) 4xy

গ) −2xy

ঘ) 2xy

\textbf{\hl{✔ সঠিক উত্তর: ক) −4xy}}

\begin{quote}
\textbf{ব্যাখ্যা:} \emph{x\^{}2 + 4y\^{}2 + 8x − 16y + 16 − 4xy}

\emph{= x\^{}2 + (−2y)\^{}2 + 4\^{}2 + 2 · x · (−2y) + 2 · (−2y) · 4 + 2
· 4 · x}

\emph{= (x − 2y + 4)\^{}2}

\emph{∴ −4xy যোগ করলে যোগফল পূর্ণবর্গ হবে।}
\end{quote}

\textbf{x\^{}2−8x−8y+16+y\^{}2 এর সঙ্গে কত যোগ করিলে যোগফল একটি পূর্ণবর্গ
হইবে?}

\textbf{\hl{ক) 2xy ✔}}

খ) −2xy

গ) 4xy

ঘ) −4xy

\textbf{\hl{✔ সঠিক উত্তর: ক) 2xy}}

\begin{quote}
\textbf{ব্যাখ্যা:} \emph{x\^{}2 − 8x − 8y + 16 + y\^{}2 + 2xy}

\emph{= x\^{}2 + y\^{}2 + (−4)\^{}2 + 2xy + 2x(−4) + 2y(−4)}

\emph{= (x + y − 4)\^{}2}

\emph{∴ 2xy যোগ করলে উহা একটি পূর্ণবর্গ হবে।}
\end{quote}

\textbf{A+B+C = 9, AB + BC + CA = 31, A\^{}2 + B\^{}2 + C\^{}2= ?}

\textbf{\hl{ক) 19 ✔}}

খ) 29

গ) 30

ঘ) 49

\textbf{\hl{✔ সঠিক উত্তর: ক) 19}}

\begin{quote}
\textbf{ব্যাখ্যা:} \emph{A\^{}2 + B\^{}2 + C\^{}2 = (A + B + C)\^{}2 −
2(AB + BC + CA)}

\emph{= 9\^{}2 − 2 × 31}

\emph{= 81 − 62 = 19}
\end{quote}

\textbf{a + b + c = 9, a\^{}2 + b\^{}2 + c\^{}2 = 29 হলে ab+bc+ca-এর মান
কত?}

ক) 52

খ) 46

\textbf{\hl{গ) 26 ✔}}

ঘ) 22

\textbf{\hl{✔ সঠিক উত্তর: গ) 26}}

\begin{quote}
\textbf{ব্যাখ্যা:} \emph{আমরা জানি, (a + b + c)\^{}2 = a\^{}2 + b\^{}2 +
c\^{}2 + 2(ab + bc + ca)}

\emph{⇒ 9\^{}2 = 29 + 2(ab + bc + ca)}

\emph{⇒ 81 − 29 = 2(ab + bc + ca)}

\emph{⇒ 52 = 2(ab + bc + ca)}

\emph{∴ ab + bc + ca = (52/2) = 26}
\end{quote}

\textbf{x\^{}2 + y\^{}2 + z\^{}2 = 5, x + y − z = 3 হলে yz + zx − xy এর
মান কত?}

ক) −5

খ) −4

গ) −3

\textbf{\hl{ঘ) −2 ✔}}

\textbf{\hl{✔ সঠিক উত্তর: ঘ) −2}}

\begin{quote}
\textbf{ব্যাখ্যা:} \emph{আমরা জানি,}

\emph{(x + y − z)\^{}2 = x\^{}2 + y\^{}2 + z\^{}2 + 2xy − 2yz − 2zx}

\emph{⇒ (3)\^{}2 = 5 − 2(yz + zx − xy)}

\emph{⇒ 9 = 5 − 2(yz + zx − xy)}

\emph{⇒ 2(yz + zx − xy) = 5 − 9}

\emph{⇒ 2(yz + zx − xy) = −4}

\emph{∴ yz + zx − xy = −2}
\end{quote}

\textbf{a = √(3) + √(2) হলে, a\^{}3 + 3a + 3a\^{}−1 + a\^{}−3 এর মান
কত?}

ক) 8 √(2)

খ) 16 √(3)

গ) 18 √(3)

\textbf{\hl{ঘ) 24 √(3) ✔}}

\textbf{\hl{✔ সঠিক উত্তর: ঘ) 24 √(3)}}

\begin{quote}
\textbf{ব্যাখ্যা:} \emph{a = √(3) + √(2)}

\emph{∴ (1/a) = (1/√(3)+√(2)) = √(3) − √(2)}

\emph{∴ a + (1/a) = √(3) + √(2) + √(3) − √(2) = 2 √(3)}

\emph{প্রদত্ত রাশি = a\^{}3 + 3a + 3a\^{}−1 + a\^{}−3}

\emph{= a\^{}3 + 3a + (3/a) + (1/a\^{}3)}

\emph{= a\^{}3 + 3 · a\^{}2 · (1/a) + 3 · a · (((1/a)))\^{}2 +
(((1/a)))\^{}3}

\emph{= ((a+(1/a)))\^{}3}

\emph{= (2 √(3) )\^{}3 = 8 × 3 √(3) = 24 √(3)}
\end{quote}

\textbf{x = 15, y = 5 হলে x\^{}3−3x\^{}2y+3xy\^{}2−y\^{}3 = ?}

ক) 1400

খ) 1200

\textbf{\hl{গ) 1000 ✔}}

ঘ) 800

\textbf{\hl{✔ সঠিক উত্তর: গ) 1000}}

\begin{quote}
\textbf{ব্যাখ্যা:} \emph{x\^{}3 − 3x\^{}2y + 3xy\^{}2 − y\^{}3 = (x −
y)\^{}3}

\emph{= (15 − 5)\^{}3 = 10\^{}3 = 1000}
\end{quote}

\textbf{x+y = 6, xy = 8 হলে, (x−y)\^{}3 = ?}

ক) 3

\textbf{\hl{খ) 8 ✔}}

গ) 9

ঘ) 12

\textbf{\hl{✔ সঠিক উত্তর: খ) 8}}

\begin{quote}
\textbf{ব্যাখ্যা:} \emph{আমরা জানি, (x − y)\^{}2 = (x + y)\^{}2 − 4xy}

\emph{= (6)\^{}2 − 4 × 8 = 36 − 32 = 4}

\emph{∴ x − y = 2}

\emph{∴ (x − y)\^{}3 = 2\^{}3 = 8}
\end{quote}

\textbf{((০.৯)\^{}৩+(০.৪)\^{}৩/০.৯+০.৪) এর মান কত?}

ক) ০.৩৬

খ) ০.৫১

গ) ০.৮১

\textbf{\hl{ঘ) ০.৬১ ✔}}

\textbf{\hl{✔ সঠিক উত্তর: ঘ) ০.৬১}}

\begin{quote}
\textbf{ব্যাখ্যা:} \emph{((০.৯)\^{}৩+(০.৪)\^{}৩/০.৯+০.৪)}

\emph{= ((০.৯+০.৪)\{(০.৯)\^{}২−(০.৯×০.৪)+(০.৪)\^{}২\}/(০.৯+০.৪))}

\emph{= (০.৯)\^{}২ − (০.৯ × ০.৪) + (০.৪)\^{}২}

\emph{= ০.৮১ − ০.৩৬ + ০.১৬}

\emph{= ০.৯৭ − ০.৩৬ = ০.৬১}
\end{quote}

\textbf{x = √(3) + √(2) হলে (x\^{}3+(1/x\^{}3)) এর মান কত?}

ক) 3 √(2)

\textbf{\hl{খ) 18 √(3) ✔}}

গ) 12 √(3)

ঘ) 8

\textbf{\hl{✔ সঠিক উত্তর: খ) 18 √(3)}}

\begin{quote}
\textbf{ব্যাখ্যা:} \emph{x = √(3) + √(2)}

\emph{∴ (1/x) = (1/√(3)+√(2)) = √(3) − √(2)}

\emph{∴ x + (1/x) = ( √(3) + √(2) ) + ( √(3) − √(2) ) = 2 √(3)}

\emph{∴ x\^{}3 + (1/x\^{}3) = ((x+(1/x)))\^{}3 − 3 · x · (1/x) ·
(x+(1/x))}

\emph{= (2 √(3) )\^{}3 − 3(2 √(3) )}

\emph{= 24 √(3) − 6 √(3) = 18 √(3)}
\end{quote}

\textbf{x + (1/x) = √(3) হলে x\^{}3 + (1/x\^{}3) এর মান কত?}

ক) 2

খ) 4

\textbf{\hl{গ) 0 ✔}}

ঘ) 6

\textbf{\hl{✔ সঠিক উত্তর: গ) 0}}

\begin{quote}
\textbf{ব্যাখ্যা:} \emph{x\^{}3 + (1/x\^{}3) = ((x+(1/x)))\^{}3 − 3 × x ×
(1/x) (x+(1/x))}

\emph{= ( √(3) )\^{}3 − 3 × √(3)}

\emph{= 3 √(3) − 3 √(3) = 0}
\end{quote}

\textbf{a + b − c = 0 হলে a\^{}3 + b\^{}3 + 3abc এর মান কত?}

\textbf{\hl{ক) c\^{}3 ✔}}

খ) a\^{}3

গ) b\^{}3

ঘ) 0

\textbf{\hl{✔ সঠিক উত্তর: ক) c\^{}3}}

\begin{quote}
\textbf{ব্যাখ্যা:} \emph{দেওয়া আছে, a + b − c = 0}

\emph{⇒ a + b = c}

\emph{প্রদত্ত রাশি = a\^{}3 + b\^{}3 + 3abc}

\emph{= (a + b)\^{}3 − 3ab(a + b) + 3abc}

\emph{= c\^{}3 − 3abc + 3abc}

\emph{= c\^{}3}
\end{quote}

\textbf{x+y = 4, xy = 5 হলে, (x\^{}3+y\^{}3)\^{}2 এর মান কত?}

\textbf{\hl{ক) 16 ✔}}

খ) 4

গ) 8

ঘ) 10

\textbf{\hl{✔ সঠিক উত্তর: ক) 16}}

\begin{quote}
\textbf{ব্যাখ্যা:} \emph{x\^{}3 + y\^{}3 = (x + y)\^{}3 − 3xy(x + y)}

\emph{= (4)\^{}3 − 3 × 5 × 4}

\emph{= 64 − 60 = 4}

\emph{∴ (x\^{}3 + y\^{}3)\^{}2 = (4)\^{}2 = 16}
\end{quote}

\textbf{x + (2/x) = 3 হলে x\^{}3 + (8/x\^{}3) = কত?}

ক) 16

\textbf{\hl{খ) 9 ✔}}

গ) 8

ঘ) 0

\textbf{\hl{✔ সঠিক উত্তর: খ) 9}}

\begin{quote}
\textbf{ব্যাখ্যা:} \emph{x\^{}3 + (8/x\^{}3) = (x)\^{}3 + (((2/x)))\^{}3}

\emph{= ((x+(2/x)))\^{}3 − 3 · x · (2/x) (x+(2/x))}

\emph{= (3)\^{}3 − 6 × 3 = 27 − 18 = 9}
\end{quote}

\textbf{x + (1/x) = √(3) হলে, x\^{}3 + (1/x\^{}3) এর মান কত?}

\textbf{\hl{ক) 0 ✔}}

খ) 6

গ) 8

ঘ) 2

\textbf{\hl{✔ সঠিক উত্তর: ক) 0}}

\begin{quote}
\textbf{ব্যাখ্যা:} \emph{x\^{}3 + (1/x\^{}3) = ((x+(1/x)))\^{}3 − 3 · x ·
(1/x) (x+(1/x))}

\emph{= ( √(3) )\^{}3 − 3( √(3) )}

\emph{= 3 √(3) − 3 √(3) = 0}
\end{quote}

\textbf{(1/a) = 3 − a হলে, a\^{}3 + (1/a\^{}3) = কত?}

ক) ১৭

\textbf{\hl{খ) ১৮ ✔}}

গ) ১৯

ঘ) ২০

\textbf{\hl{✔ সঠিক উত্তর: খ) ১৮}}

\begin{quote}
\textbf{ব্যাখ্যা:} \emph{(1/a) = 3 − a}

\emph{∴ a + (1/a) = 3}

\emph{এখন, a\^{}3 + (1/a\^{}3) = ((a+(1/a)))\^{}3 − 3 · a · (1/a)
(a+(1/a))}

\emph{= (3)\^{}3 − 3 · 3 = 27 − 9 = 18}
\end{quote}

\textbf{x\^{}2 + (1/x\^{}2) = 3 হলে (x\^{}6+1/x\^{}3) এর মান কত?}

ক) 3 √(5)

খ) 4 √(5)

\textbf{\hl{গ) 2 √(5) ✔}}

ঘ) √(5)

\textbf{\hl{✔ সঠিক উত্তর: গ) 2 √(5)}}

\begin{quote}
\textbf{ব্যাখ্যা:} \emph{((x+(1/x)))\^{}2 = x\^{}2 + (1/x\^{}2) + 2 = 3 +
2 = 5}

\emph{∴ x + (1/x) = √(5)}

\emph{এখন, (x\^{}6+1/x\^{}3) = (x\^{}6/x\^{}3) + (1/x\^{}3) = x\^{}3 +
(1/x\^{}3)}

\emph{= ((x+(1/x)))\^{}3 − 3 (x+(1/x))}

\emph{= ( √(5) )\^{}3 − 3 √(5)}

\emph{= 5 √(5) − 3 √(5) = 2 √(5)}
\end{quote}

\textbf{a\^{}3 + (1/a\^{}3) = 18 হলে a + (1/a) = কত?}

ক) 6

\textbf{\hl{খ) 3 ✔}}

গ) 0

ঘ) 1

\textbf{\hl{✔ সঠিক উত্তর: খ) 3}}

\begin{quote}
\textbf{ব্যাখ্যা:} \emph{a\^{}3 + (1/a\^{}3) = 18}

\emph{⇒ ((a+(1/a)))\^{}3 − 3 · a · (1/a) (a+(1/a)) = 18}

\emph{⇒ ((a+(1/a)))\^{}3 − 3 (a+(1/a)) = 18}

\emph{ধরি, a + (1/a) = x}

\emph{⇒ x\^{}3 − 3x = 18}

\emph{x = 3 বসিয়ে পাই, 3\^{}3 − 3 × 3 = 27 − 9 = 18}

\emph{∴ a + (1/a) = 3}
\end{quote}

\textbf{a + 3 + (1/a) = 0 হলে a\^{}3 + (1/a\^{}3) = কত?}

ক) 18

খ) 20

\textbf{\hl{গ) −18 ✔}}

ঘ) −20

\textbf{\hl{✔ সঠিক উত্তর: গ) −18}}

\begin{quote}
\textbf{ব্যাখ্যা:} \emph{a + 3 + (1/a) = 0 ⇒ a + (1/a) = −3}

\emph{∴ a\^{}3 + (1/a\^{}3) = ((a+(1/a)))\^{}3 − 3 · a · (1/a)
(a+(1/a))}

\emph{= (−3)\^{}3 − 3(−3)}

\emph{= −27 + 9 = −18}
\end{quote}

\textbf{যদি x\^{}4 − x\^{}2 + 1 = 0 হয়, তবে x\^{}3 + (1/x\^{}3) = ?}

ক) 3

খ) 2

গ) 1

\textbf{\hl{ঘ) 0 ✔}}

\textbf{\hl{✔ সঠিক উত্তর: ঘ) 0}}

\begin{quote}
\textbf{ব্যাখ্যা:} \emph{x\^{}4 − x\^{}2 + 1 = 0}

\emph{⇒ x\^{}4 + 1 = x\^{}2}

\emph{⇒ (x\^{}4+1/x\^{}2) = 1}

\emph{⇒ x\^{}2 + (1/x\^{}2) = 1}

\emph{⇒ ((x+(1/x)))\^{}2 − 2 · x · (1/x) = 1}

\emph{⇒ ((x+(1/x)))\^{}2 = 3}

\emph{∴ x + (1/x) = √(3)}

\emph{∴ x\^{}3 + (1/x\^{}3) = ((x+(1/x)))\^{}3 − 3 · x · (1/x)
(x+(1/x))}

\emph{= ( √(3) )\^{}3 − 3 √(3)}

\emph{= 3 √(3) − 3 √(3) = 0}
\end{quote}

\textbf{a\^{}3 + b\^{}3 = কত?}

ক) a\^{}3 + 3a\^{}2b + 3ab\^{}2 + b\^{}3

খ) (a + b)\^{}3 + 3ab (a + b)

\textbf{\hl{গ) (a + b)\^{}3 − 3ab (a + b) ✔}}

ঘ) (a + b) (a\^{}2 + ab + b\^{}2)

\textbf{\hl{✔ সঠিক উত্তর: গ) (a + b)\^{}3 − 3ab (a + b)}}

\begin{quote}
\textbf{ব্যাখ্যা:} \emph{আমরা জানি, a\^{}3 + b\^{}3 = (a + b)\^{}3 − 3ab(a
+ b)}
\end{quote}

\textbf{x + y = 2 এবং x\^{}2 + y\^{}2 = 4 হলে x\^{}3 + y\^{}3 এর মান
নির্ণয় করুন।}

ক) 6

খ) 7

\textbf{\hl{গ) 8 ✔}}

ঘ) 10

\textbf{\hl{✔ সঠিক উত্তর: গ) 8}}

\begin{quote}
\textbf{ব্যাখ্যা:} \emph{x\^{}2 + y\^{}2 = (x + y)\^{}2 − 2xy}

\emph{⇒ 4 = (2)\^{}2 − 2xy}

\emph{⇒ 4 = 4 − 2xy}

\emph{⇒ 2xy = 0}

\emph{⇒ xy = 0}

\emph{∴ x\^{}3 + y\^{}3 = (x + y)\^{}3 − 3xy(x + y)}

\emph{= 2\^{}3 − 3 × 0 × 2}

\emph{= 8 − 0 = 8}
\end{quote}

\textbf{x\^{}2 − √(2) x + 1 = 0 হলে (x\^{}12+1/x\^{}6) এর মান কত?}

\textbf{\hl{ক) 0 ✔}}

খ) 1

গ) √(2)

ঘ) 2

\textbf{\hl{✔ সঠিক উত্তর: ক) 0}}

\begin{quote}
\textbf{ব্যাখ্যা:} \emph{x\^{}2 − √(2) x + 1 = 0}

\emph{⇒ x\^{}2 + 1 = √(2) x}

\emph{⇒ (x\^{}2+1/x) = √(2)}

\emph{⇒ x + (1/x) = √(2)}

\emph{উভয়পক্ষকে বর্গ করে পাই,}

\emph{((x+(1/x)))\^{}2 = ( √(2) )\^{}2}

\emph{⇒ x\^{}2 + (1/x\^{}2) + 2 = 2}

\emph{⇒ x\^{}2 + (1/x\^{}2) = 0}

\emph{প্রদত্ত রাশি = (x\^{}12+1/x\^{}6) = x\^{}6 + (1/x\^{}6) =
(x\^{}2)\^{}3 + (((1/x\^{}2)))\^{}3}

\emph{= ((x\^{}2+(1/x\^{}2)))\^{}3 − 3 · x\^{}2 · (1/x\^{}2)
(x\^{}2+(1/x\^{}2))}

\emph{= (0)\^{}3 − 3(0) = 0}
\end{quote}

\textbf{x − (1/x) = 1 হলে x\^{}3 − (1/x\^{}3) এর মান কত??}

ক) 1

খ) 2

গ) 3

\textbf{\hl{ঘ) 4 ✔}}

\textbf{\hl{✔ সঠিক উত্তর: ঘ) 4}}

\begin{quote}
\textbf{ব্যাখ্যা:} \emph{x\^{}3 − (1/x\^{}3) = ((x−(1/x)))\^{}3 + 3 · x ·
(1/x) (x−(1/x))}

\emph{= (1)\^{}3 + 3 · 1 = 1 + 3 = 4}
\end{quote}

\textbf{যদি a\^{}3 − b\^{}3 = 513 এবং a − b = 3 হয় তবে ab এর মান কত?}

\textbf{\hl{ক) 54 ✔}}

খ) 35

গ) 45

ঘ) 55

\textbf{\hl{✔ সঠিক উত্তর: ক) 54}}

\begin{quote}
\textbf{ব্যাখ্যা:} \emph{a\^{}3 − b\^{}3 = (a − b)\^{}3 + 3ab (a − b)}

\emph{⇒ 513 = 3\^{}3 + 3ab · 3}

\emph{⇒ 513 = 27 + 9ab}

\emph{⇒ 9ab = 513 − 27}

\emph{⇒ 9ab = 486}

\emph{∴ ab = (486/9) = 54}
\end{quote}

\textbf{(x\^{}2−1/x) = 7 হলে, x\^{}3 − (((1/x)))\^{}3 এর মান কত?}

ক) 334

খ) 154

\textbf{\hl{গ) 364 ✔}}

ঘ) 512

\textbf{\hl{✔ সঠিক উত্তর: গ) 364}}

\begin{quote}
\textbf{ব্যাখ্যা:} \emph{(x\^{}2−1/x) = 7}

\emph{⇒ (x\^{}2/x) − (1/x) = 7}

\emph{⇒ x − (1/x) = 7}

\emph{∴ x\^{}3 − (1/x\^{}3) = ((x−(1/x)))\^{}3 + 3 · x · (1/x)
(x−(1/x))}

\emph{= (7)\^{}3 + 3 · 1 · 7 = 343 + 21 = 364}
\end{quote}

\textbf{x\^{}2 + (1/x\^{}2) এর নিম্নোক্ত কোন মানের জন্য x\^{}3 − (1/x\^{}3)
= 0 হবে?}

\textbf{\hl{ক) 2 ✔}}

খ) 1

গ) 0

ঘ) −2

\textbf{\hl{✔ সঠিক উত্তর: ক) 2}}

\begin{quote}
\textbf{ব্যাখ্যা:} \emph{x\^{}3 − (1/x\^{}3) = 0}

\emph{⇒ (x−(1/x)) (x\^{}2+x·(1/x)+(1/x\^{}2)) = 0}

\emph{⇒ (x−(1/x)) (x\^{}2+(1/x\^{}2)+1) = 0}

\emph{∴ x − (1/x) = 0 {[} ∵x\^{}2+(1/x\^{}2)+1≠0 {]}}

\emph{⇒ ((x−(1/x)))\^{}2 = 0 {[}উভয় পক্ষকে বর্গ করে{]}}

\emph{⇒ x\^{}2 − 2 · x · (1/x) + (1/x\^{}2) = 0}

\emph{⇒ x\^{}2 − 2 + (1/x\^{}2) = 0}

\emph{∴ x\^{}2 + (1/x\^{}2) = 2}
\end{quote}

\textbf{a = √(6) + √(5) হলে, (a\^{}6−1/a\^{}3) এর মান নির্ণয় কত?}

ক) 45 √(5)

খ) 47 √(6)

\textbf{\hl{গ) 46 √(5) ✔}}

ঘ) 43 √(5)

\textbf{\hl{✔ সঠিক উত্তর: গ) 46 √(5)}}

\begin{quote}
\textbf{ব্যাখ্যা:} \emph{a = √(6) + √(5)}

\emph{∴ (1/a) = (1/√(6)+√(5)) = (√(6)−√(5)/(√(6)+√(5))(√(6)−√(5))) =
√(6) − √(5)}

\emph{a − (1/a) = ( √(6) + √(5) ) − ( √(6) − √(5) ) = 2 √(5)}

\emph{প্রদত্ত রাশি = (a\^{}6−1/a\^{}3) = a\^{}3 − (1/a\^{}3) =
((a−(1/a)))\^{}3 + 3 · a · (1/a) (a−(1/a)) = (2 √(5) )\^{}3 + 3(2 √(5) )
= 40 √(5) + 6 √(5) = 46 √(5)}
\end{quote}

\textbf{m − n = x এবং mn = 6x\^{}2 হলে m\^{}3 − n\^{}3 = কত?}

ক) 19x\^{}2

\textbf{\hl{খ) 19x\^{}3 ✔}}

গ) 18x\^{}2

ঘ) 18x\^{}3

\textbf{\hl{✔ সঠিক উত্তর: খ) 19x\^{}3}}

\begin{quote}
\textbf{ব্যাখ্যা:} \emph{m\^{}3 − n\^{}3 = (m − n)\^{}3 + 3mn(m − n)}

\emph{= x\^{}3 + 3 · 6x\^{}2 · x}

\emph{= x\^{}3 + 18x\^{}3}

\emph{= 19x\^{}3}
\end{quote}

\textbf{a + b + c = 0 হলে a\^{}3 + b\^{}3 + c\^{}3 এর মান কত?}

ক) abc

\textbf{\hl{খ) 3abc ✔}}

গ) 6abc

ঘ) 9abc

\textbf{\hl{✔ সঠিক উত্তর: খ) 3abc}}

\begin{quote}
\textbf{ব্যাখ্যা:} \emph{a\^{}3 + b\^{}3 + c\^{}3 = a\^{}3 + b\^{}3 +
c\^{}3 − 3abc + 3abc}

\emph{= (a + b + c)(a\^{}2 + b\^{}2 + c\^{}2 − ab − bc − ca) + 3abc}

\emph{= 0 × (a\^{}2 + b\^{}2 + c\^{}2 − ab − bc − ca) + 3abc = 3abc}
\end{quote}

\section{73. আন্তর্জাতিক গণনা পদ্ধতি MCQ প্রশ্ন উত্তর
অনুশীলন}\label{ux986ux9a8ux9a4ux9b0ux99cux9a4ux995-ux997ux9a3ux9a8-ux9aaux9a6ux9a7ux9a4-mcq-ux9aaux9b0ux9b6ux9a8-ux989ux9a4ux9a4ux9b0-ux985ux9a8ux9b6ux9b2ux9a8}

\textbf{১ বিলিয়নে কত মিলিয়ন?}

ক) ১০০০০

\textbf{\hl{খ) ১০০০ ✔}}

গ) ১০০

ঘ) ১০

\textbf{\hl{✔ সঠিক উত্তর: খ) ১০০০}}

\begin{quote}
\textbf{ব্যাখ্যা:} \emph{আমরা জানি,}

\emph{১ মিলিয়ন = ১০\^{}৬ = ১০,০০,০০০}

\emph{১ বিলিয়ন = ১০\^{}৯ = ১,০০,০০,০০,০০০}

\emph{∴ ১ বিলিয়ন = (১০\^{}৯/১০\^{}৬) মিলিয়ন = ১০\^{}৩ মিলিয়ন = ১০০০
মিলিয়ন।}

\emph{Shortcut:}

\emph{আন্তর্জাতিক গণনাসূত্রে ১ বিলিয়ন = ১০০০ মিলিয়ন।}
\end{quote}

\textbf{১ মিলিয়ন = কত লক্ষ?}

ক) ১ লক্ষ

\textbf{\hl{খ) ১০ লক্ষ ✔}}

গ) ১০০ লক্ষ

ঘ) ১০০০ লক্ষ

\textbf{\hl{✔ সঠিক উত্তর: খ) ১০ লক্ষ}}

\begin{quote}
\textbf{ব্যাখ্যা:} \emph{আমরা জানি,}

\emph{১ মিলিয়ন = ১০,০০,০০০ (১০ লক্ষ)}

\emph{∴ ১ মিলিয়ন = ১০ লক্ষ।}
\end{quote}

\textbf{এক বিলিয়ন সমান কত কোটি?}

ক) ১০ কোটি

খ) ১০০০ কোটি

\textbf{\hl{গ) ১০০ কোটি ✔}}

ঘ) ঘ. অপশন নেই

\textbf{\hl{✔ সঠিক উত্তর: গ) ১০০ কোটি}}

\begin{quote}
\textbf{ব্যাখ্যা:} \emph{আন্তর্জাতিক গণনা পদ্ধতিতে,}

\emph{১ (U.S) বিলিয়ন = ১০\^{}৯ = ১,০০,০০,০০,০০০}

\emph{দেশীয় গণনা পদ্ধতিতে, ১ কোটি = ১০\^{}৭ = ১,০০,০০,০০০}

\emph{∴ ১ বিলিয়ন = (১০\^{}৯/১০\^{}৭) কোটি = ১০\^{}২ কোটি = ১০০ কোটি।}
\end{quote}

\textbf{How many millions make 1 crore?/ ১ কোটিতে কত মিলিয়ন হয়?}

\textbf{\hl{ক) Ten ✔}}

খ) Twenty five

গ) Fifty

ঘ) None of these

\textbf{\hl{✔ সঠিক উত্তর: ক) Ten}}

\begin{quote}
\textbf{ব্যাখ্যা:} \emph{আমরা জানি,}

\emph{১ কোটি = ১,০০,০০,০০০}

\emph{১ মিলিয়ন = ১০,০০,০০০}

\emph{∴ ১ কোটিতে মিলিয়নের সংখ্যা = (১,০০,০০,০০০/১০,০০,০০০) = ১০টি (Ten)।}
\end{quote}

\textbf{\textquotesingle এক বিলিয়ন\textquotesingle-শব্দগুলিকে সংখ্যায় প্রকাশ
করতে হলে ১ এর পর কতটি শূন্য লাগবে?}

ক) ৬টি

খ) ৭টি

গ) ৮টি

\textbf{\hl{ঘ) ৯টি ✔}}

\textbf{\hl{✔ সঠিক উত্তর: ঘ) ৯টি}}

\begin{quote}
\textbf{ব্যাখ্যা:} \emph{আন্তর্জাতিক গণনা পদ্ধতি অনুযায়ী:}

\emph{১ বিলিয়ন (1 billion) = 10\^{}9 = ১,০০০,০০০,০০০}

\emph{∴ ১ এর পর ৯টি শূন্য লাগবে।}
\end{quote}

\textbf{৫০ মিলিয়নে কত কোটি?}

ক) ৫০০ কোটি

\textbf{\hl{খ) ৫ কোটি ✔}}

গ) ৫০ কোটি

ঘ) ২০ কোটি

\textbf{\hl{✔ সঠিক উত্তর: খ) ৫ কোটি}}

\begin{quote}
\textbf{ব্যাখ্যা:} \emph{আমরা জানি,}

\emph{১ মিলিয়ন = ১০ লক্ষ}

\emph{∴ ৫০ মিলিয়ন = (৫০ × ১০) লক্ষ = ৫০০ লক্ষ}

\emph{আবার, ১০০ লক্ষ = ১ কোটি}

\emph{∴ ৫০০ লক্ষ = (৫০০ ÷ ১০০) কোটি = ৫ কোটি।}
\end{quote}

\textbf{১০০০ মিলিয়নে কত টাকা হয়?}

\textbf{\hl{ক) ১০০ কোটি ✔}}

খ) ১০ কোটি

গ) ১০ লক্ষ

ঘ) ১০০ লক্ষ

\textbf{\hl{✔ সঠিক উত্তর: ক) ১০০ কোটি}}

\begin{quote}
\textbf{ব্যাখ্যা:} \emph{১ মিলিয়ন = ১০ লক্ষ টাকা = ০.১০ কোটি টাকা}

\emph{∴ ১০০০ মিলিয়ন = (১০০০ × ১০) লক্ষ টাকা = ১০,০০০ লক্ষ টাকা}

\emph{আমরা জানি, ১০০ লক্ষ = ১ কোটি}

\emph{∴ ১০,০০০ লক্ষ টাকা = (১০,০০০ ÷ ১০০) কোটি টাকা = ১০০ কোটি টাকা।}
\end{quote}

\textbf{One million two thousand and two is written as-\/-}

ক) 10,2002

\textbf{\hl{খ) 1,002,002 ✔}}

গ) 100,202

ঘ) None of these

\textbf{\hl{✔ সঠিক উত্তর: খ) 1,002,002}}

\begin{quote}
\textbf{ব্যাখ্যা:} \emph{One million = 1,000,000}

\emph{Two thousand = 2,000}

\emph{Two = 2}

\emph{∴ Total = 1,000,000 + 2,000 + 2 = 1,002,002।}
\end{quote}

\textbf{ব্রিটিশ ১ বিলিয়ন আমেরিকান ১ বিলিয়নের কত গুণ?}

\textbf{\hl{ক) ১০০০ ✔}}

খ) ১০০

গ) ১০

ঘ) একই সমান

\textbf{\hl{✔ সঠিক উত্তর: ক) ১০০০}}

\begin{quote}
\textbf{ব্যাখ্যা:} \emph{১ (British) বিলিয়ন = 10\^{}12}

\emph{১ (U.S) বিলিয়ন = 10\^{}9}

\emph{∴ ১ (British) বিলিয়ন = 10\^{}12 = 10\^{}3 × 10\^{}9 = ১০-\^{}৩ × ১
(U.S) বিলিয়ন = ১০০০০ × ১ (U.S) বিলিয়ন = ১০০০ গুণ।}
\end{quote}

\textbf{১ ট্রিলিয়ন = ?}

\textbf{\hl{ক) ১০০০০০ কোটি ✔}}

খ) ১০০০ কোটি

গ) ১০০০০০০ কোটি

ঘ) ১০০ কোটি

\textbf{\hl{✔ সঠিক উত্তর: ক) ১০০০০০ কোটি}}

\begin{quote}
\textbf{ব্যাখ্যা:} \emph{১ (U.S) ট্রিলিয়ন = 10\^{}12}

\emph{আমরা জানি, ১ কোটি = 10\^{}7}

\emph{∴ ১ (U.S) ট্রিলিয়ন = 10\^{}5 × 10\^{}7 = 10\^{}5 × ১ কোটি = ১ লক্ষ ×
১ কোটি = ১ লক্ষ কোটি (বা ১,০০,০০০ কোটি)।}
\end{quote}

\textbf{একটি পেন্সিলের ওজন ৫ গ্রাম। পেন্সিলটির ওজন মিলিগ্রামে কত হবে?}

ক) ৫০

খ) ৫০০

\textbf{\hl{গ) ৫০০০ ✔}}

ঘ) কোনোটিই নয়

\textbf{\hl{✔ সঠিক উত্তর: গ) ৫০০০}}

\begin{quote}
\textbf{ব্যাখ্যা:} \emph{আমরা জানি,}

\emph{১ গ্রাম = ১০০০ মিলিগ্রাম}

\emph{∴ একটি পেন্সিলের ওজন = ৫ গ্রাম = (৫ × ১০০০) মিলিগ্রাম = ৫০০০
মিলিগ্রাম।}
\end{quote}

\textbf{বিবিয়ানা গ্যাস ফিল্ডে ১০ ট্রিলিয়ন ঘনফুট প্রাকৃতিক গ্যাস মজুদ আছে। প্রতি
বছর ১ লক্ষ মিলিয়ন ঘনফুট হারে উত্তোলন করা হলে কত বছরে এই ফিল্ড নিঃশেষ হবে?}

ক) ১০

খ) ৫০

\textbf{\hl{গ) ১০০ ✔}}

ঘ) ১০০০

\textbf{\hl{✔ সঠিক উত্তর: গ) ১০০}}

\begin{quote}
\textbf{ব্যাখ্যা:} \emph{১০ ট্রিলিয়ন ঘনফুট = ১০ × 10\^{}12 ঘনফুট = 10\^{}13
ঘনফুট}

\emph{১ লক্ষ মিলিয়ন ঘনফুট = (১০০০০০ × 10\^{}6 ) ঘনফুট = 10\^{}11 ঘনফুট}

\emph{∴ সময় লাগবে = (10\^{}13/10\^{}11) = 10\^{}2 = ১০০ বছর।}
\end{quote}

\textbf{প্রতি ব্যক্তির রোজ ১০০০০ টি চাউল প্রয়োজন হলে ৫× ১০\^{}২০ টি চাউলের
স্টক দ্বারা পৃথিবীর ৫০০ কোটি লোককে কতদিন খাওয়ানো যাবে?}

ক) ১০০ দিন

খ) এক লক্ষ দিন

\textbf{\hl{গ) এক কোটি দিন ✔}}

ঘ) ৫০০ কোটি দিন

\textbf{\hl{✔ সঠিক উত্তর: গ) এক কোটি দিন}}

\begin{quote}
\textbf{ব্যাখ্যা:} \emph{১ ব্যক্তির প্রতিদিন প্রয়োজন = ১০০০০ টি চাউল = 10\^{}4
টি চাউল}

\emph{৫০০ কোটি = ৫০০ × 10\^{}7 = ৫ × 10\^{}9 জন}

\emph{∴ ৫০০ কোটি লোকের প্রতিদিন মোট প্রয়োজন = (১০\^{}৪ × ৫ × ১০\^{}৯) = ৫
× 10\^{}13 টি চাউল।}

\emph{মোট দিন সংখ্যা = (মোট চাউল/একদিনের প্রয়োজনীয় চাউল)}

\emph{= (5×10\^{}20/5×10\^{}13) = 10\^{}20}

\emph{= 10\^{}7 দিন}

\emph{= ১,০০,০০,০০০ দিন = ১ কোটি দিন।}

\emph{শর্ট টেকনিক:}

\emph{দিন = (5×10\^{}20/10\^{}4×5×10\^{}9) = (10\^{}20/10\^{}13) =
10\^{}7 দিন = ১ কোটি দিন।}
\end{quote}

\section{74. সময়, দুরত্ব ও গতিবেগ সংক্রান্ত অংক MCQ প্রশ্ন উত্তর
অনুশীলন}\label{ux9b8ux9aeux9df-ux9a6ux9b0ux9a4ux9ac-ux993-ux997ux9a4ux9acux997-ux9b8ux995ux9b0ux9a8ux9a4-ux985ux995-mcq-ux9aaux9b0ux9b6ux9a8-ux989ux9a4ux9a4ux9b0-ux985ux9a8ux9b6ux9b2ux9a8}

\textbf{৪ কিমি/ ঘণ্টা বেগে চললে কোন স্থানে পৌঁছাতে যে সময় লাগে ৫ কিমি/ ঘণ্টা
বেগে চললে তার চেয়ে (১/২) ঘণ্টা কম সময় লাগে। স্থানটির দূরত্ব কত?}

ক) ৮ কিমি

\textbf{\hl{খ) ১০ কিমি ✔}}

গ) ১২ কিমি

ঘ) ১১ কিমি

\textbf{\hl{✔ সঠিক উত্তর: খ) ১০ কিমি}}

\begin{quote}
\textbf{ব্যাখ্যা:} \emph{মূল সমাধান:}

\emph{৪ কিমি যায় = ১ ঘণ্টায়}

\emph{∴ ১ কিমি যায় = (১/৪) ঘণ্টায়}

\emph{আবার, ৫ কিমি যায় = ১ ঘণ্টায়}

\emph{∴ ১ কিমি যায় = (১/৫) ঘণ্টায়}

\emph{∴ (১/৪)−(১/৫) = (১/২০) ঘণ্টা কম লাগে ১ কিমি দূরত্বের জন্য}

\emph{∴ ১ ঘণ্টা কম লাগে = ১×(২০/১) কিমি দূরত্বের জন্য}

\emph{∴ (১/২) ঘণ্টা কম লাগে = ১×(২০/১)×(১/২) = ১০ কিমি}

\emph{Shortcut:}

\emph{দূরত্ব (D) ধরা হলে,}

\emph{(D/৪)−(১/২)=(D/৫)}

\emph{⇒ (D/৪)−(D/৫)=(১/২)}

\emph{⇒ (D/২০)=(১/২)}

\emph{⇒ D = ১০ কিমি।}
\end{quote}

\textbf{ঢাকা থেকে চট্টগ্রাম এর দূরত্ব ২৬৫ কিলোমিটার। একটি বাস ৫ ঘণ্টায় ঢাকা
থেকে চট্টগ্রাম চলে আসলে বাসটির গড় গতিবেগ কত কিমি/ঘণ্টা?}

ক) ৬৫

\textbf{\hl{খ) ৫৩ ✔}}

গ) ৫৫

ঘ) ৬৩

\textbf{\hl{✔ সঠিক উত্তর: খ) ৫৩}}

\begin{quote}
\textbf{ব্যাখ্যা:} \emph{বাসটির গড় গতিবেগ = (দূরত্ব/সময়) = (২৬৫ ÷ ৫)
কিমি/ঘণ্টা = ৫৩ কিমি/ঘণ্টা।}
\end{quote}

\textbf{রাজশাহী থেকে খুলনা এর দূরত্ব ২৮২ কিলোমিটার। একটি বাস ৭ ঘণ্টায় খুলনা
থেকে রাজশাহী চলে আসল। পথে বাসটি ১ ঘণ্টা যাত্রা বিরতি নেয়। বাসটির গড় গতিবেগ
কত কি.মি./ঘণ্টা?}

ক) ৪২

খ) ৪৯

গ) ৫৫

ঘ) ৬৩

\textbf{\hl{ঙ) কোনোটিই নয় ✔}}

\textbf{\hl{✔ সঠিক উত্তর: ঙ) কোনোটিই নয়}}

\begin{quote}
\textbf{ব্যাখ্যা:} \emph{মোট সময় নেয় = ৭ ঘণ্টা}

\emph{যাত্রাবিরতি = ১ ঘণ্টা}

\emph{∴ নিট সময় নেয় = (৭ − ১) = ৬ ঘণ্টা}

\emph{মোট দূরত্ব = ২৮২ কিলোমিটার}

\emph{∴ গড় গতিবেগ = (২৮২/৬) = ৪৭ কিমি/ঘণ্টা।}
\end{quote}

\textbf{একটি গাড়ী ৫০ মাইল গড় গতিতে শহর A থেকে B শহরে সম্পূর্ণ পথের ৭৫ \%
ভাগ ভ্রমণ করে। অবশিষ্ট পথ গাড়ীটি গড়ে S মাইল/ঘণ্টায় ভ্রমণ করে। সম্পূর্ণ যাত্রায়
গড় গতিবেগ ছিল ৪০ মাইল/ঘণ্টা। S এর মান কত (মাইল/ঘণ্টা)?}

\textbf{\hl{ক) ২৫ ✔}}

খ) ৩০

গ) ২০

ঘ) ১০

\textbf{\hl{✔ সঠিক উত্তর: ক) ২৫}}

\begin{quote}
\textbf{ব্যাখ্যা:} \emph{ধরি, A এবং B এর মোট দূরত্ব = ১০০ মাইল}

\emph{৭৫\% পথ = ৭৫ মাইল এবং অবশিষ্ট পথ = (১০০ − ৭৫) = ২৫ মাইল}

\emph{৫০ মাইল/ঘণ্টা গতিতে ৭৫ মাইল যেতে সময় লাগে = (৭৫/৫০) = ১.৫ ঘণ্টা}

\emph{গড় গতিবেগ ৪০ মাইল/ঘণ্টা হলে ১০০ মাইল যেতে মোট সময় লাগে = (১০০/৪০) =
২.৫ ঘণ্টা}

\emph{∴ অবশিষ্ট ২৫ মাইল পথ যেতে সময় লাগে = (২.৫ − ১.৫) = ১ ঘণ্টা}

\emph{∴ S এর মান = (২৫/১) = ২৫ মাইল/ঘণ্টা।}

\emph{Shortcut:}

\emph{মোট দূরত্ব ১০০০ মাইল ধরলে:}

\emph{৭৫\% পথ (৭৫০ মাইল) যেতে সময় = (৭৫০/৫০) = ১৫ ঘণ্টা}

\emph{মোট সময় (৪০ mph গড়ে) = (১০০০/৪০) = ২৫ ঘণ্টা}

\emph{অবশিষ্ট ২৫০ মাইল যেতে সময় = ২৫ − ১৫ = ১০ ঘণ্টা}

\emph{∴ S = (২৫০/১০) = ২৫ মাইল/ঘণ্টা।}
\end{quote}

\textbf{A man walking at the rate of 5 km/ hr crosses a bridge in 15
minutes. The length of the bridge (in metres) is:/ একজন মানুষ ৫ কিমি /
ঘণ্টা হারে হেঁটে ১৫ মিনিটের মধ্যে একটি সেতু অতিক্রম করে। সেতুটির দৈর্ঘ্য (মিটার)
কত?}

ক) ৬০০

খ) ৭৫০

গ) ১০০০

\textbf{\hl{ঘ) ১২৫০ ✔}}

\textbf{\hl{✔ সঠিক উত্তর: ঘ) ১২৫০}}

\begin{quote}
\textbf{ব্যাখ্যা:} \emph{গতিবেগ = ৫ কিমি/ঘণ্টা = (৫ × ১০০০) মিটার/ঘণ্টা =
৫০০০ মিটার/ঘণ্টা}

\emph{আমরা জানি, ১ ঘণ্টা = ৬০ মিনিট}

\emph{৬০ মিনিটে অতিক্রম করে = ৫০০০ মিটার}

\emph{∴ ১ মিনিটে অতিক্রম করে = (৫০০০/৬০) মিটার = (২৫০/৩) মিটার}

\emph{∴ ১৫ মিনিটে অতিক্রম করে = (২৫০/৩)×১৫ মিটার = ১২৫০ মিটার}

\emph{∴ সেতুটির দৈর্ঘ্য ১২৫০ মিটার।}

\emph{Shortcut:}

\emph{দূরত্ব = গতিবেগ × সময় = ৫×(৫/১৮)×(১৫×৬০) = ১২৫০ মিটার।}
\end{quote}

\textbf{একজন লোক ১০ কি.মি./ঘণ্টা এর পরিবর্তে ১৪ কি.মি./ঘণ্টা বেগে হাঁটলে, সে
২০ কি.মি. বেশি হাঁটতে পারে। ঐ ব্যক্তি মোট কত দূরত্ব অতিক্রম করে?}

\textbf{\hl{ক) 50 km ✔}}

খ) 56 km

গ) 70 km

ঘ) 80 km

\textbf{\hl{✔ সঠিক উত্তর: ক) 50 km}}

\begin{quote}
\textbf{ব্যাখ্যা:} \emph{ধরি, লোকটির অতিক্রান্ত প্রকৃত দূরত্ব x কিমি।}

\emph{গতিবেগের অনুপাত অনুযায়ী সময় একই থাকলে দূরত্বের অনুপাত গতিবেগের সমান
হয়।}

\emph{∴ (x/১০)=(x+২০/১৪)}

\emph{⇒ ১৪x = ১০x + ২০০}

\emph{⇒ ৪x = ২০০}

\emph{∴ x = ৫০ কিমি।}

\emph{∴ অতিক্রান্ত প্রকৃত দূরত্ব ৫০ কিমি।}

\emph{Shortcut:}

\emph{গতিবেগের পার্থক্য = (১৪ − ১০) = ৪ কিমি/ঘণ্টা}

\emph{৪ কিমি/ঘণ্টা বৃদ্ধির জন্য সময় লাগে = (২০/৪) = ৫ ঘণ্টা}

\emph{∴ প্রকৃত দূরত্ব = (১০ × ৫) = ৫০ কিমি।}
\end{quote}

\textbf{একটি বন্দুকের গুলি সেকেন্ডে ১,৫৪০ ফুট গতিবেগে লক্ষ্যভেদ করে। এক ব্যক্তি
বন্দুক ছুড়বার ৩ সেকেন্ড পরে লক্ষ্যভেদের শব্দ শুনতে পায়। শব্দের গতি প্রতি সেকেন্ডে
১১০০ ফুট। লক্ষ্যবস্তুর দূরত্ব কত ফুট?}

\textbf{\hl{ক) ১৯২৫ ✔}}

খ) ১৯৭৫

গ) ২২২৫

ঘ) ১৮৫০

\textbf{\hl{✔ সঠিক উত্তর: ক) ১৯২৫}}

\begin{quote}
\textbf{ব্যাখ্যা:} \emph{ধরি, লক্ষ্যবস্তুর দূরত্ব ক ফুট}

\emph{গুলি লক্ষ্যভেদ করতে সময় লাগে = (ক/১৫৪০) সেকেন্ড}

\emph{এবং শব্দ ফিরে আসতে সময় লাগে = (ক/১১০০) সেকেন্ড}

\emph{প্রশ্নমতে,}

\emph{(ক/১৫৪০)+(ক/১১০০)=৩}

\emph{⇒ (ক/২২০)(১/৭)+(১/৫)=৩}

\emph{⇒ (ক/২২০)(১২/৩৫)=৩}

\emph{∴ ক = (২২০×৩৫×৩/১২) = ১৯২৫ ফুট।}

\emph{Shortcut:}

\emph{দূরত্ব (D) = (V\_1×V\_2/V\_1+V\_2)×মোট সময়}

\emph{= (১৫৪০×১১০০/১৫৪০+১১০০)×৩}

\emph{= (১৬৯৪০০০০/২৬৪০)×৩ = ১৯২৫ ফুট।}
\end{quote}

\textbf{একজন লোক ঘণ্টায় y কিমি হাঁটতে পারে। x কিমি দূরত্ব অতিক্রম করতে ঐ
লোকের কত মিনিট সময় লাগবে?}

ক) (x/y)

খ) (y/x)

\textbf{\hl{গ) (60x/y) ✔}}

ঘ) (60y/x)

\textbf{\hl{✔ সঠিক উত্তর: গ) (60x/y)}}

\begin{quote}
\textbf{ব্যাখ্যা:} \emph{আমরা জানি, ১ ঘণ্টা = ৬০ মিনিট}

\emph{y কিমি হাঁটতে সময় লাগে = ৬০ মিনিট}

\emph{∴ ১ কিমি হাঁটতে সময় লাগে = (৬০/y) মিনিট}

\emph{∴ x কিমি হাঁটতে সময় লাগে = (৬০×x/y) মিনিট = (60x/y) মিনিট।}
\end{quote}

\textbf{ঘণ্টায় x মাইল বেগে y মাইল দূরত্ব অতিক্রম করতে কত ঘণ্টা লাগবে?/ How
many hours will it take to cover y miles at x miles/hour?}

ক) (x/y)

\textbf{\hl{খ) (y/x) ✔}}

গ) xy

ঘ) কোনটিই নয় ।

\textbf{\hl{✔ সঠিক উত্তর: খ) (y/x)}}

\begin{quote}
\textbf{ব্যাখ্যা:} \emph{আমরা জানি,}

\emph{সময় = (দূরত্ব/গতিবেগ)}

\emph{এখানে, দূরত্ব = y মাইল এবং গতিবেগ = x মাইল/ঘণ্টা}

\emph{∴ প্রয়োজনীয় সময় = (y/x) ঘণ্টা।}
\end{quote}

\textbf{একটি কুকুর একটি খরগোশকে ধরার জন্য তাড়া করে। কুকুর যে সময়ে ৪ বার লাফ
দেয়, খরগোশ সে সময়ে ৫ বার লাফ দেয়। কিন্তু খরগোশ ৪ লাফে যতদূর যায়, কুকুর তিন
লাফে ততদূর যায়। কুকুর ও খরগোশের গতিবেগের অনুপাত কত-}

ক) ১৫:১৬

\textbf{\hl{খ) ১৬:১৫ ✔}}

গ) ২০:১২

ঘ) ১২:২০

\textbf{\hl{✔ সঠিক উত্তর: খ) ১৬:১৫}}

\begin{quote}
\textbf{ব্যাখ্যা:} \emph{খরগোশের ৪ লাফের দূরত্ব = কুকুরের ৩ লাফের দূরত্ব}

\emph{∴ খরগোশের ১ লাফের দূরত্ব = কুকুরের (৩/৪) লাফের দূরত্ব}

\emph{∴ খরগোশের ৫ লাফের দূরত্ব = কুকুরের (৩/৪)×৫ = (১৫/৪) লাফের দূরত্ব}

\emph{∴ কুকুর ও খরগোশের গতিবেগের অনুপাত = ৪ : (১৫/৪)}

\emph{= (৪ × ৪) : ১৫ = ১৬ : ১৫}

\emph{Shortcut:}

\emph{কুকুর ও খরগোশের গতিবেগের অনুপাত = (কুকুরের লাফ × খরগোশের দূরত্ব) :
(খরগোশের লাফ × কুকুরের দূরত্ব)}

\emph{= (৪ × ৪) : (৫ × ৩) = ১৬ : ১৫}
\end{quote}

\textbf{A ও B সমবেগে নির্দিষ্ট পথ অতিক্রম করে যথাক্রমে t\_1 এবং (t\_1 +
t\_2) মিনিটে। A ও B এর গতিবেগের অনুপাত কত? এবং গতিবেগের সমষ্টি কত?}

\textbf{\hl{ক) গতিবেগের অনুপাত = (t\_1 + t\_2) : t\_1 এবং গতিবেগের সমষ্টি
= (2t\_1 + t\_2)/(t\_1\^{}2 + t\_1t\_2) ✔}}

খ) গতিবেগের অনুপাত = t\_1 : (t\_1 + t\_2) এবং গতিবেগের সমষ্টি = (2t\_1 +
t\_2)

গ) গতিবেগের অনুপাত = (t\_1 + t\_2) : t\_1 এবং গতিবেগের সমষ্টি = (t\_1 +
t\_2)

ঘ) গতিবেগের অনুপাত = t\_1 : (t\_1 + t\_2) এবং গতিবেগের সমষ্টি = t\_1 +
t\_2

\textbf{\hl{✔ সঠিক উত্তর: ক) গতিবেগের অনুপাত = (t\_1 + t\_2) : t\_1 এবং
গতিবেগের সমষ্টি = (2t\_1 + t\_2)/(t\_1\^{}2 + t\_1t\_2)}}

\begin{quote}
\textbf{ব্যাখ্যা:} \emph{ধরি, মোট দূরত্ব = d}

\emph{A এর গতিবেগ, v\_A = (d/t\_1)}

\emph{B এর গতিবেগ, v\_B = (d/t\_1+t\_2)}

\emph{∴ A ও B এর গতিবেগের অনুপাত = v\_A : v\_B = (d/t\_1):(d/t\_1+t\_2) =
(t\_1 + t\_2) : t\_1}

\emph{এবং গতিবেগের সমষ্টি = v\_A + v\_B = (d/t\_1)+(d/t\_1+t\_2) =
d(t\_1+t\_2+t\_1/t\_1(t\_1+t\_2)) = (d(2t\_1+t\_2)/t\_1\^{}2+t\_1t\_2)
(একক দূরত্ব d = 1 হলে: (2t\_1+t\_2/t\_1\^{}2+t\_1t\_2) )।}
\end{quote}

\textbf{এক ব্যক্তি গাড়িযোগে ঘণ্টায় ৬০ কিলোমিটার বেগে কিছুদূর অতিক্রম করে ঘণ্টায়
৪০ কিলোমিটার বেগে অবশিষ্ট পথ অতিক্রম করলো। সে মোট ৫ ঘণ্টায় মোট ২৮০
কিলোমিটার অতিক্রম করে। সে ৬০ কিলোমিটার/ঘণ্টা বেগে কত কিলোমিটার গিয়েছিল?}

ক) ১০০

\textbf{\hl{খ) ১২০ ✔}}

গ) ১৫০

ঘ) ১৮০

ঙ) কোনোটিই নয়

\textbf{\hl{✔ সঠিক উত্তর: খ) ১২০}}

\begin{quote}
\textbf{ব্যাখ্যা:} \emph{ধরি, ঘণ্টায় ৬০ কি.মি. বেগে যায় x কি.মি.}

\emph{∴ ঘণ্টায় ৪০ কি.মি. বেগে যায় (২৮০ − x) কি.মি.}

\emph{প্রশ্নমতে,}

\emph{(x/৬০)+(২৮০−x/৪০)=৫}

\emph{⇒ (২x+৩(২৮০−x)/১২০)=৫}

\emph{⇒ (২x+৮৪০−৩x/১২০)=৫}

\emph{⇒ ৮৪০ − x = ৬০০}

\emph{⇒ x = ৮৪০ − ৬০০ = ১২০}

\emph{∴ সে ৬০ কি.মি./ঘণ্টা বেগে ১২০ কি.মি. গিয়েছিল।}

\emph{Shortcut:}

\emph{ধরি, সে পুরো ৫ ঘণ্টাই ৪০ কিমি/ঘণ্টা বেগে চললে দূরত্ব হতো = (৪০ × ৫) =
২০০ কিমি}

\emph{পার্থক্য = ২৮০ − ২০০ = ৮০ কিমি}

\emph{গতির পার্থক্য = ৬০ − ৪০ = ২০ কিমি/ঘণ্টা}

\emph{∴ ৬০ কিমি বেগে চলার সময় = (৮০/২০) = ৪ ঘণ্টা}

\emph{∴ ৬০ কিমি/ঘণ্টা বেগে অতিক্রান্ত দূরত্ব = (৬০ × ৪) = ১২০ কিমি।}
\end{quote}

\textbf{ঢাকা থেকে করিমের বাড়ির দূরত্ব ৩৫৫ কিমি। সে বাসে ঢাকা থেকে বাড়ি
রওয়ানা হলো। ৩১৯ কিমি যাওয়ার পরে বাসটি নষ্ট হয়ে গেলে করিম বাকি পথ রিকশায়
গেল। বাসের গতিবেগ ২২ কিমি/ঘণ্টা ও রিকশার গতিবেগ ৬ কিমি/ঘণ্টা হলে বাড়ি
পৌঁছতে করিমে মোট কত সময় লাগল?}

ক) ২০ ঘণ্টা

\textbf{\hl{খ) ২০ ঘণ্টা ৩০ মিনিট ✔}}

গ) ২০ ঘণ্টা ৫০ মিনিট

ঘ) ২১ ঘণ্টা

\textbf{\hl{✔ সঠিক উত্তর: খ) ২০ ঘণ্টা ৩০ মিনিট}}

\begin{quote}
\textbf{ব্যাখ্যা:} \emph{বাসে অতিক্রান্ত দূরত্ব = ৩১৯ কিমি}

\emph{বাসে সময় লাগে = (৩১৯/২২) = ১৪.৫ ঘণ্টা বা ১৪ ঘণ্টা ৩০ মিনিট}

\emph{রিকশায় অতিক্রান্ত দূরত্ব = (৩৫৫ − ৩১৯) = ৩৬ কিমি}

\emph{রিকশায় সময় লাগে = (৩৬/৬) = ৬ ঘণ্টা}

\emph{∴ বাড়ি পৌঁছাতে মোট সময় লাগল = (১৪ ঘণ্টা ৩০ মিনিট + ৬ ঘণ্টা) = ২০ ঘণ্টা
৩০ মিনিট।}
\end{quote}

\textbf{ঢাকা থেকে রংপুরের দূরত্ব ৪৫ মাইল। হাসান ঘণ্টায় ৩ মাইল বেগে এবং শাহিন
ঘণ্টায় ৪ মাইল বেগে হাঁটে। হাসান ঢাকা থেকে রওয়ানা হওয়ার ১ ঘণ্টা পর শাহিন
রংপুর থেকে ঢাকা রওয়ানা হলো। শাহিন কত মাইল হাটার পর হাসানের সাথে দেখা
হবে?}

\textbf{\hl{ক) ২৪ ✔}}

খ) ২১

গ) ২৫

ঘ) ২৩

ঙ) ২০

\textbf{\hl{✔ সঠিক উত্তর: ক) ২৪}}

\begin{quote}
\textbf{ব্যাখ্যা:} \emph{হাসান ১ ঘণ্টায় যায় = ৩ মাইল}

\emph{∴ অবশিষ্ট দূরত্ব = (৪৫ − ৩) = ৪২ মাইল}

\emph{বিপরীতমুখী আপেক্ষিক গতিবেগ = (৩ + ৪) = ৭ মাইল/ঘণ্টা}

\emph{∴ শাহিন রওয়ানা হওয়ার পর দেখা হতে সময় লাগবে = (৪২/৭) = ৬ ঘণ্টা}

\emph{∴ শাহিন হাঁটে = (৪ × ৬) = ২৪ মাইল।}

\emph{Shortcut:}

\emph{শাহিনের অতিক্রান্ত দূরত্ব = শাহিনের গতিবেগ × (মোট দূরত্ব−হাসানের ১ম
ঘণ্টার পথ/আপেক্ষিক গতিবেগ)}

\emph{= ৪ × (৪৫−৩/৩+৪) = ৪ × ৬ = ২৪ মাইল।}
\end{quote}

\textbf{একটি গাড়ী ঘণ্টায় ২৫ মাইল বেগে ঢাকা হইতে চট্টগ্রাম অভিমুখে রওনা হইল।
অপর একটি গাড়ী ঘণ্টায় ১৫ মাইল বেগে চট্টগ্রাম হইতে ঢাকা অভিমুখে রওনা হইল। ঢাকা
চট্টগ্রামের দূরত্ব ২০০ মাইল হইলে গাড়ী দুইটি অবিরাম গতিতে চলিতে থাকিলে কত সময়
পর গাড়ী দুইটি মুখোমুখি হইবে?}

\textbf{\hl{ক) ৫ ঘণ্টা ✔}}

খ) ৬ ঘণ্টা

গ) ৪ ঘণ্টা

ঘ) ৪.৫ ঘণ্টা

\textbf{\hl{✔ সঠিক উত্তর: ক) ৫ ঘণ্টা}}

\begin{quote}
\textbf{ব্যাখ্যা:} \emph{মনে করি, t ঘণ্টা পর গাড়ী দুইটি মিলিত হবে।}

\emph{t ঘণ্টায় ১ম গাড়ী ঢাকা থেকে চট্টগ্রাম অভিমুখে যায় ২৫t মাইল,}

\emph{এবং t ঘণ্টায় ২য় গাড়ী চট্টগ্রাম থেকে ঢাকা অভিমুখে আসে ১৫t মাইল।}

\emph{প্রশ্নমতে,}

\emph{২৫t + ১৫t = ২০০}

\emph{⇒ ৪০t = ২০০}

\emph{∴ t = (২০০/৪০) = ৫ ঘণ্টা।}

\emph{Shortcut:}

\emph{সময় = (মোট দূরত্ব/আপেক্ষিক গতিবেগ) = (২০০/২৫+১৫) = (২০০/৪০) = ৫
ঘণ্টা।}
\end{quote}

\textbf{ঘণ্টায় ৫ কিমি বেগে চললে কোন স্থানে পৌঁছাতে যে সময় লাগে, ঘণ্টায় ৬
কিমি বেগে চললে তার চেয়ে ৩০ মিনিট সময় কম লাগে। স্থানটির দূরত্ব কত?}

ক) ১২ কিমি

\textbf{\hl{খ) ১৫ কিমি ✔}}

গ) ১৮ কিমি

ঘ) ২০ কিমি

\textbf{\hl{✔ সঠিক উত্তর: খ) ১৫ কিমি}}

\begin{quote}
\textbf{ব্যাখ্যা:} \emph{ধরি, স্থানটির দূরত্ব x কিমি}

\emph{৩০ মিনিট = (১/২) ঘণ্টা}

\emph{প্রশ্নমতে,}

\emph{(x/৫)−(x/৬)=(১/২)}

\emph{⇒ (৬x−৫x/৩০)=(১/২)}

\emph{⇒ (x/৩০)=(১/২)}

\emph{⇒ x = (৩০/২)}

\emph{∴ x = ১৫}

\emph{∴ স্থানটির দূরত্ব ১৫ কিমি।}

\emph{Shortcut:}

\emph{দূরত্ব = (S\_1×S\_2/S\_2−S\_1)×সময়ের পার্থক্য (ঘণ্টায়)}

\emph{= (৫×৬/৬−৫)×(১/২)}

\emph{= (৩০/১)×(১/২) = ১৫ কিমি।}
\end{quote}

\textbf{১টি কুকুর ১টি শিয়ালের ৫০০ মিটার পিছন থেকে তাড়া করলো। ১ কিমি যেতে
শিয়ালের ১০ মিনিট এবং কুকুরের ৬ মিনিট সময় লাগলে কতক্ষণ পর কুকুর শিয়ালটিকে ধরতে
পারবে?}

ক) ৬ মিনিট

\textbf{\hl{খ) ৭ (১/২) মিনিট ✔}}

গ) ৮ মিনিট

ঘ) ১০ মিনিট

\textbf{\hl{✔ সঠিক উত্তর: খ) ৭ (১/২) মিনিট}}

\begin{quote}
\textbf{ব্যাখ্যা:} \emph{১ কিমি = ১০০০০ মিটার}

\emph{কুকুর ৬ মিনিটে যায় ১০০০ মিটার}

\emph{∴ কুকুর ১ মিনিটে যায় = ১০০০৬ = (৫০০/৩) মিটার}

\emph{শিয়াল ১০ মিনিটে যায় ১০০০ মিটার}

\emph{∴ শিয়াল ১ মিনিটে যায় = (১০০০/১০) = ১০০ মিটার}

\emph{কুকুর ১ মিনিটে বেশি যায় = (৫০০/৩)−১০০ = (২০০/৩) মিটার}

\emph{কুকুর (২০০/৩) মিটার এগিয়ে যায় ১ মিনিটে}

\emph{∴ ৫০০ মিটার এগিয়ে যেতে সময় লাগবে = (৩×৫০০/২০০) = (১৫/২) = ৭ (১/২)
মিনিটে।}

\emph{Shortcut:}

\emph{কুকুরের গতিবেগ = (১০০০/৬) মি/মি,}

\emph{শিয়ালের গতিবেগ = (১০০০/১০) = ১০০ মি/মি}

\emph{আপেক্ষিক গতিবেগ = (১০০০/৬)−১০০ = (২০০/৩) মি/মি}

\emph{সময় = (৫০০/(২০০/৩)) = (৫০০×৩/২০০) = ৭ (১/২) মিনিট।}
\end{quote}

\textbf{ঘণ্টায় ৪ কি. মি. গতি বৃদ্ধি করায় ৩২ কি. মি. পথ অতিক্রম করতে ৪ ঘণ্টা
সময় কম লাগে। বৃদ্ধির পূর্বে গতি কত ছিল?}

ক) ৮ কি. মি.

খ) ১২ কি. মি.

\textbf{\hl{গ) ৪ কি. মি. ✔}}

ঘ) ২ কি. মি.

\textbf{\hl{✔ সঠিক উত্তর: গ) ৪ কি. মি.}}

\begin{quote}
\textbf{ব্যাখ্যা:} \emph{ধরি, বৃদ্ধির পূর্বে গতি ছিল ঘণ্টায় x কি.মি.}

\emph{∴ বৃদ্ধির পর গতি ঘণ্টায় (x + ৪) কি.মি.}

\emph{প্রশ্নমতে,}

\emph{(৩২/x)−(৩২/x+৪)=৪}

\emph{⇒ ৩২(x+৪−x/x(x+৪))=৪}

\emph{⇒ (১২৮/x\^{}২+৪x)=৪}

\emph{⇒ ৪(x\^{}২ + ৪x) = ১২৮}

\emph{⇒ x\^{}২ + ৪x = ৩২}

\emph{⇒ x\^{}২ + ৪x − ৩২ = ০}

\emph{⇒ x\^{}২ + ৮x − ৪x − ৩২ = ০}

\emph{⇒ x(x + ৮) − ৪(x + ৮) = ০}

\emph{⇒ (x + ৮)(x − ৪) = ০}

\emph{যেহেতু গতি ঋণাত্মক হতে পারে না, তাই x − ৪ = ০ ⇒ x = ৪}

\emph{∴ বৃদ্ধির পূর্বে গতি ছিল ৪ কি. মি./ঘণ্টা।}
\end{quote}

\textbf{একটি ট্রেন ঘণ্টায় ৬০ কিলোমিটার বেগে চললে ১০০ মিটার যেতে ট্রেনটির কত
সময় লাগবে?}

ক) ৬০ সেকেন্ড

খ) ১ সেকেন্ড

\textbf{\hl{গ) ৬ সেকেন্ড ✔}}

ঘ) ০.৬ সেকেন্ড

\textbf{\hl{✔ সঠিক উত্তর: গ) ৬ সেকেন্ড}}

\begin{quote}
\textbf{ব্যাখ্যা:} \emph{৬০ কিমি = ৬০,০০০ মিটার}

\emph{১ ঘণ্টা = (৬০ × ৬০) সেকেন্ড = ৩৬০০ সেকেন্ড}

\emph{ট্রেনটি ৬০,০০০ মিটার যায় ৩৬০০ সেকেন্ডে}

\emph{∴ ট্রেনটি ১ মিটার যায় (৩৬০০/৬০,০০০) সেকেন্ডে}

\emph{∴ ট্রেনটি ১০০ মিটার যায় (৩৬০০/৬০,০০০)×১০০ = ৬ সেকেন্ডে।}

\emph{Shortcut:}

\emph{গতিবেগ = ৬০ × (৫/১৮) = (৫০/৩) মি/সেকেন্ড}

\emph{সময় = (দূরত্ব/গতিবেগ) = (১০০/(৫০/৩)) = (১০০×৩/৫০) = ৬ সেকেন্ড।}
\end{quote}

\textbf{ঢাকা ও চট্টগ্রাম রেল স্টেশন থেকে প্রতি ঘণ্টায় একটি করে ট্রেন উভয়
স্টেশনের দিকে যাত্রা করে। প্রত্যেক ট্রেন সমান গতিতে চলে এবং গন্তব্যস্থলে পৌঁছতে ৫
ঘণ্টা সময় লাগে। এক স্টেশন থেকে যাত্রা করে অন্য স্টেশনে পৌঁছা পর্যন্ত একটি ট্রেন
কয়টি ট্রেনের দেখা পাবে?}

ক) ৫

\textbf{\hl{খ) ১০ ✔}}

গ) ১১

ঘ) ৮

\textbf{\hl{✔ সঠিক উত্তর: খ) ১০}}

\begin{quote}
\textbf{ব্যাখ্যা:} \emph{একটি ট্রেন যে সময় যাত্রা শুরু করে, ঠিক সেই সময়ে
বিপরীত স্টেশন থেকে ছাড়া ট্রেন এবং তার আগের ৫ ঘণ্টায় ছেড়ে আসা ৫টি ট্রেন
ইতিমধ্যে লাইনে থাকে।}

\emph{এছাড়া, নিজ যাত্রাপথের ৫ ঘণ্টার মধ্যে বিপরীত দিক থেকে আরও ৫টি নতুন ট্রেন
রওনা হবে।}

\emph{∴ মোট দেখা পাওয়া ট্রেনের সংখ্যা = (৫ + ৫) = ১০টি।}

\emph{Shortcut:}

\emph{যাত্রার সময় T হলে, মোট দেখা পাওয়া ট্রেনের সংখ্যা = ২ × T = ২ × ৫ =
১০টি (যদি যাত্রার শুরুর মুহূর্তে এবং পৌঁছানোর মুহূর্তে স্টেশন ছাড়াকে হিসাবের বাইরে
রাখা হয়)।}
\end{quote}

\textbf{দুটি বাস ঘণ্টায় ২০ কিঃ মিঃ বেগে একই সময়ে গাবতলী হতে আরিচার উদ্দেশ্যে
রওনা হল। সাভার পৌঁছার পর একটি বাস থেমে গেল। কিন্তু অপর বাসটি চলতে থাকল। আধা
ঘণ্টা পর থেমে থাকা বাসটি আবার ২৫ কিঃ মি/ ঘণ্টা বেগে চলতে শুরু করল। সাভার হতে
কত দূরে বাস দুটি মিলিত হবে?}

ক) ৪০ কিঃমিঃ

\textbf{\hl{খ) ৫০ কিঃমিঃ ✔}}

গ) ৬০ কিঃমিঃ

ঘ) ৫৫ কিঃমিঃ

\textbf{\hl{✔ সঠিক উত্তর: খ) ৫০ কিঃমিঃ}}

\begin{quote}
\textbf{ব্যাখ্যা:} \emph{ধরি, থেমে থাকা ২য় গাড়ি চলতে শুরু করার t ঘণ্টা পর বাস
দুটি মিলিত হবে।}

\emph{t ঘণ্টায় ২৫ কিমি/ঘণ্টা বেগে ২য় গাড়িটি যায় = ২৫t কিমি}

\emph{যে বাসটি সাভারে থামেনি, তা চলতে দেখা হবার পূর্বে মোট সময় পায় =
t+(১/২) ঘণ্টা}

\emph{২০ কিমি/ঘণ্টা বেগে ১ম বাসটি যায় = ২০t+(১/২) = (২০t + ১০) কিমি}

\emph{শর্তমতে,}

\emph{২৫t = ২০t + ১০}

\emph{⇒ ৫t = ১০}

\emph{∴ t = ২}

\emph{∴ সাভার হতে মিলিত হওয়ার স্থানের দূরত্ব = (২৫ × ২) = ৫০ কিমি।}

\emph{Shortcut:}

\emph{আধা ঘণ্টায় ১ম বাসটি এগিয়ে যায় = (২০ × ০.৫) = ১০ কিমি}

\emph{গতির পার্থক্য = (২৫ − ২০) = ৫ কিমি/ঘণ্টা}

\emph{২য় বাসের ১ম বাসকে ধরতে সময় লাগে = (১০/৫) = ২ ঘণ্টা}

\emph{∴ সাভার থেকে অতিক্রান্ত দূরত্ব = (২৫ × ২) = ৫০ কিমি।}
\end{quote}

\textbf{এক ব্যক্তি কোনো দূরত্ব ঘণ্টায় 4 মাইল বেগে অতিক্রম করে এবং ঘণ্টায় 5 মাইল
বেগে ফিরে আসে। তার গড় গতিবেগ ঘণ্টায় কত মাইল?}

ক) 5

খ) 4.54

\textbf{\hl{গ) 4.44 ✔}}

ঘ) 4.50

\textbf{\hl{✔ সঠিক উত্তর: গ) 4.44}}

\begin{quote}
\textbf{ব্যাখ্যা:} \emph{ধরি, দূরত্ব = x মাইল}

\emph{x দূরত্ব যেতে সময় লাগে = (x/৪) ঘণ্টা}

\emph{x দূরত্ব ফিরে আসতে সময় লাগে = (x/৫) ঘণ্টা}

\emph{∴ গড় বেগ = (x+x/(x/৪)+(x/৫)) = (২x/(৯x/২০)) = (৪০/৯) = ৪.৪৪
মাইল/ঘণ্টা।}

\emph{Shortcut:}

\emph{সমান দূরত্বের ক্ষেত্রে গড় গতিবেগ = (২xy/x+y) = (২×৪×৫/৪+৫) = (৪০/৯) =
৪.৪৪ মাইল/ঘণ্টা।}
\end{quote}

\textbf{জুয়েল ১২ মিনিট ৪৫ সেকেন্ডে এক মাইল পথ অতিক্রম করে। কোয়েল একই পথ ১০
মিনিট ১২ সেকেন্ডে অতিক্রম করে। কোয়েলের গতিবেগ অর্জন করতে চাইলে জুয়েলকে তার
গতিবেগ কতটুকু বাড়াতে হবে?}

ক) ১০\%

খ) ২০\%

গ) ১২০\%

\textbf{\hl{ঘ) ২৫\% ✔}}

\textbf{\hl{✔ সঠিক উত্তর: ঘ) ২৫\%}}

\begin{quote}
\textbf{ব্যাখ্যা:} \emph{জুয়েলের সময় = ১২ মিনিট ৪৫ সেকেন্ড = ১২+(৪৫/৬০)
মিনিট = ১২+(৩/৪) মিনিট = (৫১/৪) মিনিট = ৭৬৫ সেকেন্ড}

\emph{কোয়েলের সময় = ১০ মিনিট ১২ সেকেন্ড = ১০×৬০+১২ সেকেন্ড = ৬১২ সেকেন্ড}

\emph{যেহেতু দূরত্ব সমান (১ মাইল), তাই গতিবেগ সময়ের বিপরীতানুপাতিক।}

\emph{জুয়েল ও কোয়েলের সময়ের অনুপাত = ৭৬৫ : ৬১২ = ৫ : ৪}

\emph{∴ জুয়েল ও কোয়েলের গতিবেগের অনুপাত = ৪ : ৫}

\emph{জুয়েলকে গতিবেগ ৪ একক থেকে বাড়িয়ে ৫ একক করতে হবে,}

\emph{অর্থাৎ বৃদ্ধি করতে হবে = (৫ − ৪) = ১ একক}

\emph{∴ গতিবেগ বাড়াতে হবে = (১/৪)×১০০\% = ২৫\%।}
\end{quote}

\textbf{সজীব তার বাসা থেকে ৪ কিলোমিটার দূরে দোকানে গিয়ে ফেরত আসল। যাবার
পথে তার ৫০ মিনিট সময় লাগলেও ফেরত আসার সময় তার গতি ১০\% কমে গেলে, দোকানে
আসা -- যাওয়াতে সজীবের মোট কত মিনিট লাগল?}

ক) ৯৫

খ) ১০০

\textbf{\hl{গ) ১০৫.৫৫ ✔}}

ঘ) ১০৬

\textbf{\hl{✔ সঠিক উত্তর: গ) ১০৫.৫৫}}

\begin{quote}
\textbf{ব্যাখ্যা:} \emph{যাওয়ার সময় লাগে = ৫০ মিনিট = (৫০/৬০) = (৫/৬)
ঘণ্টা}

\emph{যাওয়ার গতিবেগ = (৪/(৫/৬)) = ৪×(৬/৫) = (২৪/৫) কিমি/ঘণ্টা}

\emph{ফেরার গতিবেগ (১০\% কমলে) = (২৪/৫) এর ৯০\% = (২৪/৫)×(৯০/১০০) =
(২১৬/৫০) কিমি/ঘণ্টা}

\emph{ফেরার সময় লাগে = (৪/(২১৬/৫০)) = ৪×(৫০/২১৬) = (৫০/৫৪) ঘণ্টা =
(৫০/৫৪)×৬০ মিনিট = ৫৫.৫৫ মিনিট}

\emph{∴ মোট সময় লাগল = (৫০ + ৫৫.৫৫) = ১০৫.৫৫ মিনিট।}

\emph{Shortcut:}

\emph{গতিবেগ ১০\% কমা মানে নতুন গতিবেগ মূল গতির ৯০\% বা (৯/১০) গুণ।}

\emph{গতিবেগ ও সময় ব্যস্তানুপাতিক হওয়ায় ফেরার সময় লাগবে মূল সময়ের (১০/৯)
গুণ।}

\emph{∴ ফেরার সময় = ৫০×(১০/৯) = (৫০০/৯) = ৫৫.৫৫ মিনিট}

\emph{∴ মোট সময় = ৫০ + ৫৫.৫৫ = ১০৫.৫৫ মিনিট।}
\end{quote}

\textbf{\textquotesingle ক\textquotesingle{} হতে
\textquotesingle খ\textquotesingle{} এর দূরত্ব ৫৫ কিলোমিটার।
\textquotesingle ক\textquotesingle{} ও
\textquotesingle খ\textquotesingle{} এর গতিবেগ যথাক্রমে ৩ কিলোমিটার এবং ৪
কিলোমিটার। \textquotesingle ক\textquotesingle{} এর ২ ঘণ্টা পর
\textquotesingle খ\textquotesingle{}
\textquotesingle ক\textquotesingle{} এর দিকে রওনা হলে
\textquotesingle খ\textquotesingle{} কত কিলোমিটার গেলে
\textquotesingle ক\textquotesingle{} এর দেখা পাবে?}

ক) ২৪

\textbf{\hl{খ) ২৮ ✔}}

গ) ২০

ঘ) ১৮

\textbf{\hl{✔ সঠিক উত্তর: খ) ২৮}}

\begin{quote}
\textbf{ব্যাখ্যা:} \emph{\textquotesingle ক\textquotesingle{} ব্যক্তি ২
ঘণ্টায় যায় = (৩ × ২) = ৬ কিলোমিটার}

\emph{∴ অবশিষ্ট দূরত্ব = (৫৫ − ৬) = ৪৯ কিলোমিটার}

\emph{\textquotesingle ক\textquotesingle{} ও
\textquotesingle খ\textquotesingle{} বিপরীতমুখী হওয়ায় আপেক্ষিক গতিবেগ = (৩
+ ৪) = ৭ কিমি/ঘণ্টা}

\emph{\textquotesingle খ\textquotesingle{} রওনা হওয়ার পর দেখা হতে সময়
লাগবে = (৪৯/৭) = ৭ ঘণ্টা}

\emph{∴ \textquotesingle খ\textquotesingle{} মোট পথ অতিক্রম করবে = (৪ ×
৭) = ২৮ কিলোমিটার।}

\emph{Shortcut:}

\emph{\textquotesingle খ\textquotesingle{} এর অতিক্রান্ত দূরত্ব =
\textquotesingle খ\textquotesingle{} এর গতিবেগ × (মোট
দূরত্ব−\textquotesingle ক\textquotesingle{} এর ২ ঘণ্টার পথ/আপেক্ষিক গতিবেগ)}

\emph{= ৪ × (৫৫−৬/৩+৪) = ৪ × (৪৯/৭) = ২৮ কিলোমিটার।}
\end{quote}

\textbf{একটি বানর ১৫ মিটার উঁচু পিচ্ছিল বাঁশের মাথায় উঠতে প্রথম সেকেন্ডে ৩
মিটার উঠে এবং পরবর্তী সেকেন্ডে ১ মিটার নেমে আসে। বাঁশের মাথায় উঠতে বানরটির
কত সময় লাগবে?}

ক) ১২ সেকেন্ড

\textbf{\hl{খ) ১৩ সেকেন্ড ✔}}

গ) ১৪ সেকেন্ড

ঘ) ১৫ সেকেন্ড

\textbf{\hl{✔ সঠিক উত্তর: খ) ১৩ সেকেন্ড}}

\begin{quote}
\textbf{ব্যাখ্যা:} \emph{বানরটি ১ম সেকেন্ডে উঠে ৩ মিটার এবং ২য় সেকেন্ডে নেমে
আসে ১ মিটার}

\emph{∴ (১ + ১) = ২ সেকেন্ডে মোট উঠে = (৩ − ১) = ২ মিটার}

\emph{শেষ লাফ বা ১ম সেকেন্ডে ৩ মিটার ওঠার পর বানরটি আর নিচে নামবে না।}

\emph{∴ বাঁশের মোট দৈর্ঘ্য থেকে শেষ লাফের দূরত্ব বাদ দিলে থাকে = (১৫ − ৩) =
১২ মিটার}

\emph{বানরটির ২ মিটার উঠতে সময় লাগে = ২ সেকেন্ড}

\emph{∴ বানরটির ১২ মিটার উঠতে সময় লাগে = (২×১২/২) = ১২ সেকেন্ড}

\emph{∴ অবশিষ্ট ৩ মিটার (শেষ লাফ) উঠতে সময় লাগে = ১ সেকেন্ড}

\emph{∴ বাঁশের মাথায় উঠতে বানরটির মোট সময় লাগবে = (১২ + ১) = ১৩ সেকেন্ড।}

\emph{Shortcut:}

\emph{প্রয়োজনীয় সময় = (মোট দৈর্ঘ্য−১ম লাফের ওঠা/কার্যকরী ওঠা)×প্রতি চক্রের
সময়+শেষ লাফের সময়}

\emph{= (১৫−৩/৩−১)×২+১ = (১২/২)×২+১ = ১২ + ১ = ১৩ সেকেন্ড।}
\end{quote}

\textbf{একটি পাহাড় থেকে নামতে যে সময় লাগে পাহাড়ে উঠতে ৩ গুণ সময় লাগে
উঠানামা আর সে সঙ্গে ২ ঘণ্টা বিশ্রাম নিয়ে সর্বমোট সময় যদি লাগে ১৪ ঘণ্টা, তাহলে
পাহাড়ের মাথায় উঠতে কত সময় লাগে?}

ক) ৭ ঘণ্টা

\textbf{\hl{খ) ৯ ঘণ্টা ✔}}

গ) ১০ ঘণ্টা

ঘ) ৬ ঘণ্টা

\textbf{\hl{✔ সঠিক উত্তর: খ) ৯ ঘণ্টা}}

\begin{quote}
\textbf{ব্যাখ্যা:} \emph{ধরি, পাহাড় থেকে নামতে সময় লাগে = x ঘণ্টা}

\emph{∴ পাহাড়ে উঠতে সময় লাগে = ৩x ঘণ্টা}

\emph{প্রশ্নমতে,}

\emph{x + ৩x + ২ = ১৪}

\emph{⇒ ৪x = ১৪ − ২}

\emph{⇒ ৪x = ১২}

\emph{∴ x = ৩ ঘণ্টা}

\emph{∴ পাহাড়ে উঠতে সময় লাগে = (৩ × ৩) = ৯ ঘণ্টা।}
\end{quote}

\textbf{একটি রাস্তার পাশে এক সারিতে ১৫টি গাছ লাগানো আছে। একটি গাছ থেকে
আরেকটি গাছের দূরত্ব ১০ মিটার হলে প্রথম ও শেষ গাছ দুটির মধ্যে দূরত্ব কত?}

ক) ১০৯ মিটার

খ) ১২০ মিটার

\textbf{\hl{গ) ১৪০ মিটার ✔}}

ঘ) ১৫০ মিটার

\textbf{\hl{✔ সঠিক উত্তর: গ) ১৪০ মিটার}}

\begin{quote}
\textbf{ব্যাখ্যা:} \emph{১৫টি গাছ পর পর বসালে গাছের মধ্যবর্তী ফাঁকা স্থান
(ব্যবধানের সংখ্যা) হবে = (১৫ − ১) = ১৪টি।}

\emph{প্রতিটি ব্যবধানের দূরত্ব = ১০ মিটার}

\emph{∴ প্রথম ও শেষ গাছ দুটির মধ্যে দূরত্ব = (১৪ × ১০) = ১৪০ মিটার।}

\emph{Shortcut:}

\emph{দূরত্ব = (গাছের সংখ্যা − ১) × প্রতি ব্যবধানের দূরত্ব}

\emph{= (১৫ − ১) × ১০ = ১৪ × ১০ = ১৪০ মিটার।}
\end{quote}

\textbf{৬ ফুট অন্তর বৃক্ষের চারা রোপণ করা হলে ১০০ গজ দীর্ঘ রাস্তায় সর্বোচ্চ
কতগুলো চারা রোপণ করা যাবে?}

ক) ৭টি

খ) ৫০টি

\textbf{\hl{গ) ৫১টি ✔}}

ঘ) ৬০টি

\textbf{\hl{✔ সঠিক উত্তর: গ) ৫১টি}}

\begin{quote}
\textbf{ব্যাখ্যা:} \emph{আমরা জানি, ১ গজ = ৩ ফুট}

\emph{∴ ১০০ গজ = (১০০ × ৩) ফুট = ৩০০ ফুট}

\emph{রাস্তার শুরুতে ১ম চারাটি রোপণ করার পর বাকি ৩০০ ফুট রাস্তায় ৬ ফুট পর পর
চারা লাগবে = (৩০০/৬) = ৫০টি}

\emph{∴ মোট চারা রোপণ করা যাবে = (৫০ + ১) = ৫১টি।}

\emph{Shortcut:}

\emph{মোট চারা = (মোট দৈর্ঘ্য (ফুটে)/ব্যবধানের দূরত্ব (ফুটে))+১ = (৩০০/৬)+১ =
৫০ + ১ = ৫১টি।}
\end{quote}

\textbf{রাত ৩ টায় তাপমাত্রা −১৩° হলে এবং সকাল ৬টায় তাপমাত্রা বৃদ্ধি পেয়ে ২০°
হলে তাপমাত্রার পরিবর্তন কত হবে?}

ক) ৭°C

খ) −৭°C

\textbf{\hl{গ) ৩৩°C ✔}}

ঘ) ১৩°C

\textbf{\hl{✔ সঠিক উত্তর: গ) ৩৩°C}}

\begin{quote}
\textbf{ব্যাখ্যা:} \emph{প্রাথমিক তাপমাত্রা = −১৩°C}

\emph{শেষ তাপমাত্রা = ২০°C}

\emph{∴ তাপমাত্রার পরিবর্তন = শেষ তাপমাত্রা − প্রাথমিক তাপমাত্রা}

\emph{= ২০°C − (−১৩°C)}

\emph{= ২০°C + ১৩°C}

\emph{= ৩৩°C।}

\emph{Shortcut:}

\emph{ঋণাত্মক তাপমাত্রা থেকে ধনাত্মক তাপমাত্রায় পরিবর্তন = মান দুটির যোগফল =
১৩ + ২০ = ৩৩°C।}
\end{quote}

\textbf{−১ থেকে কত বিয়োগ করলে বিয়োগফল শূন্য হবে?/ What is to be subtracted
from -1 so that the difference be 0?}

\textbf{\hl{ক) −১ ✔}}

খ) ১

গ) ২

ঘ) ঘ. ২

\textbf{\hl{✔ সঠিক উত্তর: ক) −১}}

\begin{quote}
\textbf{ব্যাখ্যা:} \emph{ধরি, বিয়োগ করতে হবে x}

\emph{প্রশ্নমতে,}

\emph{−১ − x = ০}

\emph{⇒ −x = ১}

\emph{∴ x = −১}

\emph{∴ −১ থেকে −১ বিয়োগ করলে বিয়োগফল শূন্য হবে।}
\end{quote}

\section{75. বয়স সংক্রান্ত অংক MCQ প্রশ্ন উত্তর
অনুশীলন}\label{ux9acux9dfux9b8-ux9b8ux995ux9b0ux9a8ux9a4-ux985ux995-mcq-ux9aaux9b0ux9b6ux9a8-ux989ux9a4ux9a4ux9b0-ux985ux9a8ux9b6ux9b2ux9a8}

\textbf{৫ জন বালকের বয়সের গড় ১০ বছর। ঐ দলে আরও দুজন বালক যোগ দিলে তাদের
বয়সের গড় হয় ১২ বছর। যোগদানকারী বালক দুটি যদি সমবয়সী হয় তবে তাদের প্রত্যেকের
বয়স কত?}

ক) ১২ বছর

খ) ১৫ বছর

\textbf{\hl{গ) ১৭ বছর ✔}}

ঘ) ২০ বছর

\textbf{\hl{✔ সঠিক উত্তর: গ) ১৭ বছর}}

\begin{quote}
\textbf{ব্যাখ্যা:} \emph{৫ জন বালকের বয়সের সমষ্টি = (১০ × ৫) = ৫০ বছর}

\emph{আরও ২ জন বালক যোগ দেওয়ায় মোট বালক = (৫ + ২) = ৭ জন}

\emph{৭ জন বালকের বয়সের সমষ্টি = (১২ × ৭) = ৮৪ বছর}

\emph{নতুন ২ জন বালকের বয়সের সমষ্টি = (৮৪ − ৫০) = ৩৪ বছর}

\emph{যেহেতু যোগদানকারী বালক দুটি সমবয়সী,}

\emph{∴ তাদের প্রত্যেকের বয়স = (৩৪ ÷ ২) = ১৭ বছর।}
\end{quote}

\textbf{সোনিয়া ও লিনিয়ার বর্তমান বয়সের সমষ্টি ১৬ বছর। চার বছর পরে সোনিয়ার
বয়স লিনিয়ার তিনগুণ হলে, সোনিয়ার বর্তমান বয়স কত?}

ক) ৮ বৎসর

খ) ১০ বৎসর

গ) ১২ বৎসর

\textbf{\hl{ঘ) ১৪ বৎসর ✔}}

\textbf{\hl{✔ সঠিক উত্তর: ঘ) ১৪ বৎসর}}

\begin{quote}
\textbf{ব্যাখ্যা:} \emph{ধরি, লিনিয়ার বর্তমান বয়স =
\textquotesingle ক\textquotesingle{} বছর}

\emph{∴ সোনিয়ার বর্তমান বয়স = (১৬ − ক) বছর}

\emph{৪ বছর পর,}

\emph{লিনিয়ার বয়স হবে = (ক + ৪) বছর}

\emph{সোনিয়ার বয়স হবে = (১৬ − ক + ৪) = (২০ − ক) বছর}

\emph{প্রশ্নমতে,}

\emph{⇒ ২০ − ক = ৩(ক + ৪)}

\emph{⇒ ২০ − ক = ৩ক + ১২}

\emph{⇒ ৪ক = ৮}

\emph{⇒ ক = ২}

\emph{∴ সোনিয়ার বর্তমান বয়স = (১৬ − ২) = ১৪ বছর।}
\end{quote}

\textbf{পিতার বয়স পুত্রের দ্বিগুণ অপেক্ষা ২ বছর বেশি পিতার বয়স ৬২ বছর হলে
পুত্রের বয়স কত?}

ক) ২৫ বছর

\textbf{\hl{খ) ৩০ বছর ✔}}

গ) ৩৫ বছর

ঘ) ৪০ বছর

\textbf{\hl{✔ সঠিক উত্তর: খ) ৩০ বছর}}

\begin{quote}
\textbf{ব্যাখ্যা:} \emph{ধরি, পুত্রের বয়স =
\textquotesingle ক\textquotesingle{} বছর}

\emph{∴ পিতার বয়স = (২ক + ২) বছর}

\emph{প্রশ্নমতে,}

\emph{⇒ ২ক + ২ = ৬২}

\emph{⇒ ২ক = ৬২ − ২}

\emph{⇒ ২ক = ৬০}

\emph{⇒ ক = ৩০}

\emph{∴ পুত্রের বয়স = ৩০ বছর।}
\end{quote}

\textbf{পিতার ২৫ বছর বয়সে পুত্রের জন্ম হয়। পিতার কত বছর বয়সে তার বয়স পুত্রের
বয়সের দ্বিগুণ হবে?}

ক) ৪০ বছর

খ) ৪৫ বছর

\textbf{\hl{গ) ৫০ বছর ✔}}

ঘ) ৫৫ বছর

\textbf{\hl{✔ সঠিক উত্তর: গ) ৫০ বছর}}

\begin{quote}
\textbf{ব্যাখ্যা:} \emph{পিতার ২৫ বছর বয়সে পুত্রের জন্ম হওয়ায়, পিতা ও পুত্রের
বয়সের পার্থক্য সর্বদা ২৫ বছর।}

\emph{ধরি, \textquotesingle ক\textquotesingle{} বছর পর পিতার বয়স পুত্রের
বয়সের দ্বিগুণ হবে।}

\emph{তখন পুত্রের বয়স = \textquotesingle ক\textquotesingle{} বছর এবং পিতার
বয়স = (২৫ + ক) বছর।}

\emph{প্রশ্নমতে,}

\emph{⇒ ২৫ + ক = ২ক}

\emph{⇒ ক = ২৫}

\emph{∴ পিতার বয়স হবে = ২৫ + ২৫ = ৫০ বছর।}
\end{quote}

\textbf{পাঁচ সন্তানের বয়সের গড় ৭ বছর এবং পিতাসহ তাদের বয়সের গড় ১৩ বছর।
পিতার বয়স কত?}

ক) ৩৩ বছর

\textbf{\hl{খ) ৪৩ বছর ✔}}

গ) ৫৩ বছর

ঘ) ৬৩ বছর

\textbf{\hl{✔ সঠিক উত্তর: খ) ৪৩ বছর}}

\begin{quote}
\textbf{ব্যাখ্যা:} \emph{৫ সন্তানের বয়সের সমষ্টি = (৭ × ৫) = ৩৫ বছর}

\emph{পিতাসহ ৫ সন্তানের (মোট ৬ জন) বয়সের সমষ্টি = (১৩ × ৬) = ৭৮ বছর}

\emph{∴ পিতার বয়স = (৭৮ − ৩৫) = ৪৩ বছর।}
\end{quote}

\textbf{বাবা ও মেয়ের বয়সের যোগফল ৪৩ বছর। ৪ বছর আগে বাবার বয়স মেয়ের বয়সের
৬ গুণ ছিল। বাবা ও মেয়ের বর্তমান বয়স কত?}

ক) ৩২ ও ১১

খ) ৩৬ ও ৭

\textbf{\hl{গ) ৩৪ ও ৯ ✔}}

ঘ) ৩৮ ও ৫

\textbf{\hl{✔ সঠিক উত্তর: গ) ৩৪ ও ৯}}

\begin{quote}
\textbf{ব্যাখ্যা:} \emph{ধরি, ৪ বছর আগে মেয়ের বয়স ছিল
\textquotesingle ক\textquotesingle{} বছর।}

\emph{∴ ৪ বছর আগে বাবার বয়স ছিল \textquotesingle{}৬ক\textquotesingle{}
বছর।}

\emph{বর্তমানে,}

\emph{মেয়ের বয়স = (ক + ৪) বছর}

\emph{বাবার বয়স = (৬ক + ৪) বছর}

\emph{প্রশ্নমতে,}

\emph{⇒ (৬ক + ৪) + (ক + ৪) = ৪৩}

\emph{⇒ ৭ক + ৮ = ৪৩}

\emph{⇒ ৭ক = ৪৩ − ৮}

\emph{⇒ ৭ক = ৩৫}

\emph{⇒ ক = ৫}

\emph{∴ বাবার বর্তমান বয়স = (৬ × ৫ + ৪) = ৩৪ বছর}

\emph{এবং মেয়ের বর্তমান বয়স = (৫ + ৪) = ৯ বছর।}
\end{quote}

\textbf{লিমা রিমার চেয়ে ১০ বছরের বড়। ৭ বছর পর লিমার বয়স রিমার দ্বিগুণ হবে।
লিমার বর্তমান বয়স কত?}

ক) ৩

\textbf{\hl{খ) ১৩ ✔}}

গ) ২৩

ঘ) ১৫

ঙ) ৫

\textbf{\hl{✔ সঠিক উত্তর: খ) ১৩}}

\begin{quote}
\textbf{ব্যাখ্যা:} \emph{ধরি, লিমার বর্তমান বয়স =
\textquotesingle ক\textquotesingle{} বছর।}

\emph{যেহেতু লিমা রিমার চেয়ে ১০ বছরের বড়, তাই রিমার বর্তমান বয়স = (ক − ১০)
বছর।}

\emph{৭ বছর পর,}

\emph{লিমার বয়স হবে = (ক + ৭) বছর}

\emph{রিমার বয়স হবে = (ক − ১০ + ৭) = (ক − ৩) বছর}

\emph{প্রশ্নমতে,}

\emph{⇒ ক + ৭ = ২(ক − ৩)}

\emph{⇒ ক + ৭ = ২ক − ৬}

\emph{⇒ ২ক − ক = ৭ + ৬}

\emph{⇒ ক = ১৩}

\emph{∴ লিমার বর্তমান বয়স = ১৩ বছর।}
\end{quote}

\textbf{পিতা ও পুত্রের বর্তমান বয়স একত্রে ৮০ বছর। ৪ বছর পূর্বে পিতার বয়স পুত্রের
বয়সের ৫ গুণ ছিল। তাদের বর্তমান বয়সের অনুপাত কত?}

\textbf{\hl{ক) ৬৪: ১৬ ✔}}

খ) ৬০ : ২০

গ) ৫৬: ২৪

ঘ) ৬৮: ১২

\textbf{\hl{✔ সঠিক উত্তর: ক) ৬৪: ১৬}}

\begin{quote}
\textbf{ব্যাখ্যা:} \emph{ধরি, ৪ বছর পূর্বে পুত্রের বয়স ছিল
\textquotesingle ক\textquotesingle{} বছর এবং পিতার বয়স ছিল
\textquotesingle{}৫ক\textquotesingle{} বছর।}

\emph{∴ বর্তমানে পুত্রের বয়স = (ক + ৪) বছর এবং পিতার বয়স = (৫ক + ৪) বছর।}

\emph{প্রশ্নমতে,}

\emph{⇒ (৫ক + ৪) + (ক + ৪) = ৮০}

\emph{⇒ ৬ক + ৮ = ৮০}

\emph{⇒ ৬ক = ৮০ − ৮}

\emph{⇒ ৬ক = ৭২}

\emph{⇒ ক = ১২}

\emph{∴ পিতার বর্তমান বয়স = (৫ × ১২ + ৪) = ৬৪ বছর}

\emph{এবং পুত্রের বর্তমান বয়স = (১২ + ৪) = ১৬ বছর}

\emph{∴ পিতা ও পুত্রের বর্তমান বয়সের অনুপাত = ৬৪:১৬।}
\end{quote}

\textbf{পিতা ও তার দুই সন্তানের বয়সের গড় ২৫ বছর। দুই সন্তানের বয়সের গড় ২২
বছর হলে পিতার বয়স -}

ক) ২৫ বছর

খ) ২১ বছর

\textbf{\hl{গ) ৩১ বছর ✔}}

ঘ) ৩২ বছর

\textbf{\hl{✔ সঠিক উত্তর: গ) ৩১ বছর}}

\begin{quote}
\textbf{ব্যাখ্যা:} \emph{পিতা ও ২ সন্তান (মোট ৩ জন) এর বয়সের সমষ্টি = (২৫ ×
৩) = ৭৫ বছর}

\emph{২ সন্তানের বয়সের সমষ্টি = (২২ × ২) = ৪৪ বছর}

\emph{∴ পিতার বয়স = (৭৫ − ৪৪) = ৩১ বছর।}
\end{quote}

\textbf{পিতা ও পুত্রের বয়সের সমষ্টি ৯০ বছর, পিতার বয়স পুত্রের বয়সের দ্বিগুণ,
পিতার বয়স কত?}

ক) ৬৫ বছর

\textbf{\hl{খ) ৬০ বছর ✔}}

গ) ৭৫ বছর

ঘ) ৭০ বছর

\textbf{\hl{✔ সঠিক উত্তর: খ) ৬০ বছর}}

\begin{quote}
\textbf{ব্যাখ্যা:} \emph{ধরি, পুত্রের বয়স =
\textquotesingle ক\textquotesingle{} বছর}

\emph{∴ পিতার বয়স = \textquotesingle{}২ক\textquotesingle{} বছর}

\emph{প্রশ্নমতে,}

\emph{⇒ ক + ২ক = ৯০}

\emph{⇒ ৩ক = ৯০}

\emph{⇒ ক = ৩০}

\emph{∴ পিতার বয়স = (৩০ × ২) = ৬০ বছর।}
\end{quote}

\textbf{পিতা ও পুত্রের বয়সের যোগফল ১২০ বছর, পিতার বয়স পুত্রের বয়সের তিনগুণ।
পিতার বয়স কত?}

ক) ৭৫

খ) ৮০

গ) ৮৫

\textbf{\hl{ঘ) ৯০ ✔}}

\textbf{\hl{✔ সঠিক উত্তর: ঘ) ৯০}}

\begin{quote}
\textbf{ব্যাখ্যা:} \emph{ধরি, পুত্রের বয়স =
\textquotesingle ক\textquotesingle{} বছর}

\emph{∴ পিতার বয়স = \textquotesingle{}৩ক\textquotesingle{} বছর}

\emph{প্রশ্নমতে,}

\emph{⇒ ক + ৩ক = ১২০}

\emph{⇒ ৪ক = ১২০}

\emph{⇒ ক = ৩০}

\emph{∴ পিতার বয়স = (৩০ × ৩) = ৯০ বছর।}
\end{quote}

\textbf{ক, খ ও গ এর গড় বয়স ৪০ বছর, ক ও গ এর বয়স একত্রে ৮৫ বছর হলে খ এর
বয়স হবে-}

ক) ৩০ বছর

\textbf{\hl{খ) ৩৫ বছর ✔}}

গ) ৪০ বছর

ঘ) ৪৫ বছর

\textbf{\hl{✔ সঠিক উত্তর: খ) ৩৫ বছর}}

\begin{quote}
\textbf{ব্যাখ্যা:} \emph{ক, খ ও গ (মোট ৩ জন) এর বয়সের সমষ্টি = (৪০ × ৩) =
১২০ বছর}

\emph{ক ও গ এর বয়সের সমষ্টি = ৮৫ বছর}

\emph{∴ খ এর বয়স = (১২০ − ৮৫) = ৩৫ বছর।}
\end{quote}

\textbf{পিতা ও দুই সন্তানের বয়সের গড় ২৭ বছর। দুই সন্তানের বয়সের গড় ২০ বছর
হলে পিতার বয়স কত?}

ক) ৩৫ বছর

খ) ৩৮ বছর

\textbf{\hl{গ) ৪১ বছর ✔}}

ঘ) ৪৮ বছর

\textbf{\hl{✔ সঠিক উত্তর: গ) ৪১ বছর}}

\begin{quote}
\textbf{ব্যাখ্যা:} \emph{পিতা ও ২ সন্তান (মোট ৩ জন) এর বয়সের সমষ্টি = (২৭ ×
৩) = ৮১ বছর}

\emph{২ সন্তানের বয়সের সমষ্টি = (২০ × ২) = ৪০ বছর}

\emph{∴ পিতার বয়স = (৮১ − ৪০) = ৪১ বছর।}
\end{quote}

\textbf{পিতা পুত্রের চেয়ে ৩২ বছরের বড়। ৭ বছর পর পিতার বয়স পুত্রের বয়সের
দ্বিগুণ অপেক্ষা ৫ বছর বেশি হবে। ৩ বছর পর পিতার বয়স কত হবে?}

ক) ৪৫

\textbf{\hl{খ) ৫৫ ✔}}

গ) ৬৫

ঘ) ৭৫

ঙ) ৮৫

\textbf{\hl{✔ সঠিক উত্তর: খ) ৫৫}}

\begin{quote}
\textbf{ব্যাখ্যা:} \emph{ধরি, পুত্রের বর্তমান বয়স =
\textquotesingle ক\textquotesingle{} বছর}

\emph{∴ পিতার বর্তমান বয়স = (ক + ৩২) বছর}

\emph{৭ বছর পর,}

\emph{পুত্রের বয়স হবে = (ক + ৭) বছর}

\emph{পিতার বয়স হবে = (ক + ৩২ + ৭) = (ক + ৩৯) বছর}

\emph{প্রশ্নমতে,}

\emph{⇒ ২(ক + ৭) + ৫ = ক + ৩৯}

\emph{⇒ ২ক + ১৪ + ৫ = ক + ৩৯}

\emph{⇒ ২ক + ১৯ = ক + ৩৯}

\emph{⇒ ২ক − ক = ৩৯ − ১৯}

\emph{⇒ ক = ২০}

\emph{∴ পিতার বর্তমান বয়স = (২০ + ৩২) = ৫২ বছর}

\emph{∴ ৩ বছর পর পিতার বয়স হবে = (৫২ + ৩) = ৫৫ বছর।}
\end{quote}

\textbf{পিতা ও পুত্রের বয়সের সমষ্টি ৯০ বছর। পিতার বয়স পুত্রের বয়সের দ্বিগুণ।
পিতার বয়স কত?}

ক) ৬৫ বৎসর

\textbf{\hl{খ) ৬০ বৎসর ✔}}

গ) ৭০ বৎসর

ঘ) ৭৫ বৎসর

\textbf{\hl{✔ সঠিক উত্তর: খ) ৬০ বৎসর}}

\begin{quote}
\textbf{ব্যাখ্যা:} \emph{ধরি, পুত্রের বয়স =
\textquotesingle ক\textquotesingle{} বছর}

\emph{∴ পিতার বয়স = \textquotesingle{}২ক\textquotesingle{} বছর}

\emph{প্রশ্নমতে,}

\emph{⇒ ক + ২ক = ৯০}

\emph{⇒ ৩ক = ৯০}

\emph{⇒ ক = ৩০}

\emph{∴ পিতার বয়স = (৩০ × ২) = ৬০ বছর।}
\end{quote}

\textbf{জাবেদ তার স্ত্রী ইয়াসমিন থেকে ৭ বছরের বড়। ইয়াসমিনের বয়স তার কন্যার
বয়সের ৭ গুণ। যদি কন্যার বয়স ৪ বছর পরে ৭ বছর হয় তবে জাবেদের বয়স কত?}

ক) ২১

খ) ২৪

\textbf{\hl{গ) ২৮ ✔}}

ঘ) ৩৫

\textbf{\hl{✔ সঠিক উত্তর: গ) ২৮}}

\begin{quote}
\textbf{ব্যাখ্যা:} \emph{৪ বছর পর কন্যার বয়স হবে = ৭ বছর}

\emph{∴ কন্যার বর্তমান বয়স = (৭ − ৪) = ৩ বছর}

\emph{ইয়াসমিনের বর্তমান বয়স = (৩ × ৭) = ২১ বছর}

\emph{∴ জাবেদের বর্তমান বয়স = (২১ + ৭) = ২৮ বছর।}
\end{quote}

\textbf{রহিমের বয়স ১২ বছর। রহিমের বয়স করিমের বয়সের ৩ গুণ। যখন রহিমের বয়স
করিমের বয়সের দ্বিগুণ হবে তখন রহিমের বয়স কত?}

ক) ১৫ বছর

\textbf{\hl{খ) ১৬ বছর ✔}}

গ) ১৭ বছর

ঘ) ১৮ বছর

\textbf{\hl{✔ সঠিক উত্তর: খ) ১৬ বছর}}

\begin{quote}
\textbf{ব্যাখ্যা:} \emph{রহিমের বর্তমান বয়স = ১২ বছর}

\emph{∴ করিমের বর্তমান বয়স = (১২ ÷ ৩) = ৪ বছর}

\emph{ধরি, \textquotesingle ক\textquotesingle{} বছর পর রহিমের বয়স করিমের
বয়সের দ্বিগুণ হবে।}

\emph{\textquotesingle ক\textquotesingle{} বছর পর রহিমের বয়স হবে (১২ +
ক) বছর এবং করিমের বয়স হবে (৪ + ক) বছর।}

\emph{প্রশ্নমতে,}

\emph{⇒ ১২ + ক = ২(৪ + ক)}

\emph{⇒ ১২ + ক = ৮ + ২ক}

\emph{⇒ ২ক − ক = ১২ − ৮}

\emph{⇒ ক = ৪}

\emph{∴ রহিমের বয়স হবে = (১২ + ৪) = ১৬ বছর।}
\end{quote}

\textbf{পিতা ও পুত্রের বর্তমান বয়সের অনুপাত ৯: ২ এবং ১৫ বছর পরে তাদের বয়সের
অনুপাত ১২: ৫ হবে। তাদের বর্তমান বয়স কত?}

ক) পিতা ৬৩ বছর, পুত্র ১৩ বছর

খ) পিতা ৫৪ বছর, পুত্র ১২ বছর

গ) পিতা ৩৩ বছর, পুত্র ৮ বছর

\textbf{\hl{ঘ) পিতা ৪৫ বছর, পুত্র ১০ বছর ✔}}

\textbf{\hl{✔ সঠিক উত্তর: ঘ) পিতা ৪৫ বছর, পুত্র ১০ বছর}}

\begin{quote}
\textbf{ব্যাখ্যা:} \emph{ধরি, পিতা ও পুত্রের বর্তমান বয়স যথাক্রমে
\textquotesingle{}৯ক\textquotesingle{} বছর এবং
\textquotesingle{}২ক\textquotesingle{} বছর।}

\emph{১৫ বছর পর,}

\emph{পিতার বয়স হবে = (৯ক + ১৫) বছর}

\emph{পুত্রের বয়স হবে = (২ক + ১৫) বছর}

\emph{প্রশ্নমতে,}

\emph{⇒ (৯ক + ১৫) : (২ক + ১৫) = ১২ : ৫}

\emph{⇒ (৯ক + ১৫) / (২ক + ১৫) = ১২ / ৫}

\emph{⇒ ৫(৯ক + ১৫) = ১২(২ক + ১৫)}

\emph{⇒ ৪৫ক + ৭৫ = ২৪ক + ১৮০}

\emph{⇒ ৪৫ক − ২৪ক = ১৮০ − ৭৫}

\emph{⇒ ২১ক = ১০৫}

\emph{⇒ ক = ৫}

\emph{∴ পিতার বর্তমান বয়স = (৯ × ৫) = ৪৫ বছর}

\emph{এবং পুত্রের বর্তমান বয়স = (২ × ৫) = ১০ বছর।}
\end{quote}

\textbf{পিতা ও পুত্রের বর্তমান বয়সের সমষ্টি ৫০ বছর; যখন পুত্রের বয়স পিতার
বর্তমান বয়সের সমান হবে তখন তাদের বয়সের সমষ্টি হবে ১০২ বছর। পিতা ও পুত্রের
বর্তমান বয়স কত?}

ক) পিতা ৩২ বছর এবং পুত্র ১৮ বছর

খ) পিতা ৩৬ বছর এবং পুত্র ১৪ বছর

\textbf{\hl{গ) পিতা ৩৮ বছর এবং পুত্র ১২ বছর ✔}}

ঘ) পিতা ৪০ বছর এবং পুত্র ১০ বছর

\textbf{\hl{✔ সঠিক উত্তর: গ) পিতা ৩৮ বছর এবং পুত্র ১২ বছর}}

\begin{quote}
\textbf{ব্যাখ্যা:} \emph{ধরি, পিতার বর্তমান বয়স =
\textquotesingle প\textquotesingle{} বছর এবং পুত্রের বর্তমান বয়স =
\textquotesingle ফ\textquotesingle{} বছর।}

\emph{∴ প + ফ = ৫০ -\/-\/-\/-\/-\/- (১)}

\emph{পুত্রের বয়স পিতার সমান (অর্থাৎ \textquotesingle প\textquotesingle{}
বছর) হতে সময় লাগবে = (প − ফ) বছর।}

\emph{ততদিনে পিতার বয়স বৃদ্ধি পেয়ে হবে = প + (প − ফ) = (২প − ফ) বছর।}

\emph{প্রশ্নমতে,}

\emph{⇒ (২প − ফ) + প = ১০২}

\emph{⇒ ৩প − ফ = ১০২ -\/-\/-\/-\/-\/- (২)}

\emph{(১) ও (২) নং সমীকরণ যোগ করে পাই,}

\emph{⇒ ৪প = ১৫২}

\emph{⇒ প = ৩৮}

\emph{(১) নং সমীকরণে প = ৩৮ বসিয়ে পাই,}

\emph{⇒ ৩৮ + ফ = ৫০}

\emph{⇒ ফ = ১২}

\emph{∴ পিতার বর্তমান বয়স ৩৮ বছর এবং পুত্রের বর্তমান বয়স ১২ বছর।}
\end{quote}

\textbf{পিতা ও পুত্রের বয়সের অনুপাত ৭: ৩। ৪ বছর পূর্বে তাদের বয়সের অনুপাত ছিল
১৩: ৫। বর্তমানে কার বয়স কত?}

ক) ৫৬ বছর, ৩৪ বছর

খ) ৬৬ বছর, ২৪ বছর

\textbf{\hl{গ) ৫৬ বছর, ২৪ বছর ✔}}

ঘ) ৪৬ বছর, ৩৩ বছর

\textbf{\hl{✔ সঠিক উত্তর: গ) ৫৬ বছর, ২৪ বছর}}

\begin{quote}
\textbf{ব্যাখ্যা:} \emph{ধরি, ৪ বছর পূর্বে পিতার বয়স ছিল
\textquotesingle{}১৩ক\textquotesingle{} বছর এবং পুত্রের বয়স ছিল
\textquotesingle{}৫ক\textquotesingle{} বছর।}

\emph{∴ পিতার বর্তমান বয়স = (১৩ক + ৪) বছর এবং পুত্রের বর্তমান বয়স = (৫ক + ৪)
বছর।}

\emph{প্রশ্নমতে,}

\emph{⇒ (১৩ক + ৪) : (৫ক + ৪) = ৭ : ৩}

\emph{⇒ (১৩ক + ৪) / (৫ক + ৪) = ৭ / ৩}

\emph{⇒ ৩(১৩ক + ৪) = ৭(৫ক + ৪)}

\emph{⇒ ৩৯ক + ১২ = ৩৫ক + ২৮}

\emph{⇒ ৩৯ক − ৩৫ক = ২৮ − ১২}

\emph{⇒ ৪ক = ১৬}

\emph{⇒ ক = ৪}

\emph{∴ পিতার বর্তমান বয়স = (১৩ × ৪ + ৪) = ৫৬ বছর}

\emph{এবং পুত্রের বর্তমান বয়স = (৫ × ৪ + ৪) = ২৪ বছর।}
\end{quote}

\textbf{৫ বছর পূবে পিতা ও পুত্রের বয়সের যোগফল ছিল ৬০ বৎসর। বর্তমানে পুত্রের
বয়স ২০ হলে পিতার বর্তমান বয়স কত?}

ক) ৪৫ বছর

খ) ৪০ বছর

গ) ৫৫ বছর

\textbf{\hl{ঘ) ৫০ বছর ✔}}

\textbf{\hl{✔ সঠিক উত্তর: ঘ) ৫০ বছর}}

\begin{quote}
\textbf{ব্যাখ্যা:} \emph{বর্তমানে পুত্রের বয়স = ২০ বছর}

\emph{∴ ৫ বছর পূর্বে পুত্রের বয়স ছিল = (২০ − ৫) = ১৫ বছর}

\emph{৫ বছর পূর্বে পিতার বয়স ছিল = (৬০ − ১৫) = ৪৫ বছর}

\emph{∴ বর্তমানে পিতার বয়স = (৪৫ + ৫) = ৫০ বছর।}
\end{quote}

\textbf{আব্দুল করিম আবদুর রহিমের চাইতে ৩ বছরের ছোট। আফজালের বয়স আবদুল করিমের
চাইতে ২ বছর কম। মুমিনের বয়স যখন ৫ বছর তখন আবদুল করিম জন্মেছে। তাদের মধ্যে
জেষ্ঠ্যতমের বয়স ৫২ হলে আফজালের বয়স কত?}

ক) ৫৪ বছর

\textbf{\hl{খ) ৪৫ বছর ✔}}

গ) ৫০ বছর

ঘ) ৪৩ বছর

\textbf{\hl{✔ সঠিক উত্তর: খ) ৪৫ বছর}}

\begin{quote}
\textbf{ব্যাখ্যা:} \emph{ধরি, আবদুল করিমের বয়স =
\textquotesingle ক\textquotesingle{} বছর}

\emph{∴ মুমিনের বয়স = (ক + ৫) বছর}

\emph{আবদুর রহিমের বয়স = (ক + ৩) বছর}

\emph{আফজালের বয়স = (ক − ২) বছর}

\emph{এখানে মুমিন জেষ্ঠ্যতম (সবচেয়ে বড়)।}

\emph{প্রশ্নমতে,}

\emph{⇒ ক + ৫ = ৫২}

\emph{⇒ ক = ৪৭}

\emph{∴ আফজালের বয়স = (৪৭ − ২) = ৪৫ বছর।}
\end{quote}

\textbf{পিতা ও পুত্রের বয়সের সমষ্টি ৬০ বছর। মাতার বয়স পুত্রের বয়স অপেক্ষা ২০
বছর বেশি। পিতা ও মাতার গড় বয়স কত??}

ক) ২০

খ) ১০

\textbf{\hl{গ) ৪০ ✔}}

ঘ) ৬০

\textbf{\hl{✔ সঠিক উত্তর: গ) ৪০}}

\begin{quote}
\textbf{ব্যাখ্যা:} \emph{ধরি, পুত্রের বয়স =
\textquotesingle ক\textquotesingle{} বছর}

\emph{∴ পিতার বয়স = (৬০ − ক) বছর}

\emph{মাতার বয়স = (ক + ২০) বছর}

\emph{পিতা ও মাতার বয়সের সমষ্টি = (৬০ − ক) + (ক + ২০) = ৮০ বছর}

\emph{∴ পিতা ও মাতার গড় বয়স = (৮০ ÷ ২) = ৪০ বছর।}
\end{quote}

\textbf{পিতা ও পুত্রের বর্তমান বয়সের অনুপাত ৩: ১ এবং ৫ বছর পরে তাদের বয়সের
অনুপাত ৫: ২ হবে। তাদের বর্তমান বয়স কত??}

ক) পিতা ৩৯ বছর এবং পুত্র ১৩ বছর

খ) পিতা ৪২ বছর এবং পুত্র ১৪ বছর

\textbf{\hl{গ) পিতা ৪৫ বছর এবং পুত্র ১৫ বছর ✔}}

ঘ) পিতা ৪৮ বছর এবং পুত্র ১৬ বছর

\textbf{\hl{✔ সঠিক উত্তর: গ) পিতা ৪৫ বছর এবং পুত্র ১৫ বছর}}

\begin{quote}
\textbf{ব্যাখ্যা:} \emph{ধরি, পিতা ও পুত্রের বর্তমান বয়স যথাক্রমে
\textquotesingle{}৩ক\textquotesingle{} বছর এবং
\textquotesingle ক\textquotesingle{} বছর।}

\emph{৫ বছর পর,}

\emph{পিতার বয়স হবে = (৩ক + ৫) বছর}

\emph{পুত্রের বয়স হবে = (ক + ৫) বছর}

\emph{প্রশ্নমতে,}

\emph{⇒ (৩ক + ৫) : (ক + ৫) = ৫ : ২}

\emph{⇒ (৩ক + ৫) / (ক + ৫) = ৫ / ২}

\emph{⇒ ২(৩ক + ৫) = ৫(ক + ৫)}

\emph{⇒ ৬ক + ১০ = ৫ক + ২৫}

\emph{⇒ ৬ক − ৫ক = ২৫ − ১০}

\emph{⇒ ক = ১৫}

\emph{∴ পিতার বর্তমান বয়স = (৩ × ১৫) = ৪৫ বছর}

\emph{এবং পুত্রের বর্তমান বয়স = ১৫ বছর।}
\end{quote}

\textbf{দুই বছর আগে বাবার বয়স ছিল পুত্রের বয়সের ১৪ গুণ। দুই বছর পরে বাবার বয়স
পুত্রের বয়সের চেয়ে ২৬ বছর বেশি হলে, বাবা ও তার পুত্রের বয়সের অনুপাত কত হবে?}

\textbf{\hl{ক) ১৫:২ ✔}}

খ) ৬:১

গ) ৪:১

ঘ) কোনটিই নয়

\textbf{\hl{✔ সঠিক উত্তর: ক) ১৫:২}}

\begin{quote}
\textbf{ব্যাখ্যা:} \emph{পিতা ও পুত্রের বয়সের পার্থক্য সর্বদা সমান থাকে।}

\emph{∴ পিতা ও পুত্রের বয়সের পার্থক্য = ২৬ বছর}

\emph{ধরি, ২ বছর আগে পুত্রের বয়স ছিল \textquotesingle ক\textquotesingle{}
বছর}

\emph{∴ ২ বছর আগে পিতার বয়স ছিল \textquotesingle{}১৪ক\textquotesingle{}
বছর}

\emph{প্রশ্নমতে,}

\emph{⇒ ১৪ক − ক = ২৬}

\emph{⇒ ১৩ক = ২৬}

\emph{⇒ ক = ২}

\emph{∴ ২ বছর আগে পুত্রের বয়স ছিল = ২ বছর}

\emph{এবং ২ বছর আগে পিতার বয়স ছিল = (১৪ × ২) = ২৮ বছর}

\emph{বর্তমান বয়স:}

\emph{পুত্রের বর্তমান বয়স = (২ + ২) = ৪ বছর}

\emph{পিতার বর্তমান বয়স = (২৮ + ২) = ৩০ বছর}

\emph{∴ পিতা ও পুত্রের বর্তমান বয়সের অনুপাত = ৩০ : ৪ = ১৫ : ২।}
\end{quote}

\textbf{খাদেম ও তার বাবার বয়সের সমষ্টি ৫০ বছর। ৩ বছর পর তাদের বয়সের সমষ্টি
কত হবে?}

ক) ৬৬ বছর

খ) ৫৩ বছর

\textbf{\hl{গ) ৫৬ বছর ✔}}

ঘ) ৫২ বছর

\textbf{\hl{✔ সঠিক উত্তর: গ) ৫৬ বছর}}

\begin{quote}
\textbf{ব্যাখ্যা:} \emph{খাদেম ও তার বাবার বয়সের সমষ্টি ৫০ বছর।}

\emph{৩ বছর পর খাদেম ও তার বাবা ২ জনেরই বয়স ৩ বছর করে বৃদ্ধি পাবে।}

\emph{∴ ৩ বছর পর তাদের বয়সের সমষ্টি হবে = ৫০ + (৩ × ২) = ৫০ + ৬ = ৫৬
বছর।}
\end{quote}

\textbf{আফজাল সাহেব ও তার দুই ছেলের বয়সের সমষ্টি ৩৬ বছর। ৫ বছর পরে তাদের
বয়সের সমষ্টি কত বছর হবে?}

ক) ৪০

খ) ১১

গ) ৪৬

\textbf{\hl{ঘ) ৫১ ✔}}

\textbf{\hl{✔ সঠিক উত্তর: ঘ) ৫১}}

\begin{quote}
\textbf{ব্যাখ্যা:} \emph{আফজাল সাহেব ও তার ২ জন ছেলে মিলিয়ে মোট ৩ জন
ব্যক্তি।}

\emph{বর্তমান বয়সের সমষ্টি = ৩৬ বছর।}

\emph{৫ বছর পর প্রত্যেকের বয়স ৫ বছর করে বৃদ্ধি পাবে, অর্থাৎ মোট ৩ জন ব্যক্তির
বয়স বৃদ্ধি পাবে = (৩ × ৫) = ১৫ বছর।}

\emph{∴ ৫ বছর পর তাদের বয়সের সমষ্টি হবে = (৩৬ + ১৫) বছর = ৫১ বছর।}
\end{quote}

\textbf{খালেক ও তার বাবার বয়সের সমষ্টি ৪০ বছর। খালেকের বাবা তার চেয়ে ২৮
বছরের বড়। ১৩ বছর পর তাদের বয়সের সমষ্টি কত হবে?}

\textbf{\hl{ক) ৬৬ বছর ✔}}

খ) ৫৩ বছর

গ) ৫৬ বছর

ঘ) ৭২ বছর

\textbf{\hl{✔ সঠিক উত্তর: ক) ৬৬ বছর}}

\begin{quote}
\textbf{ব্যাখ্যা:} \emph{খালেক ও তার বাবার বয়সের বর্তমান সমষ্টি = ৪০ বছর।}

\emph{১৩ বছর পর খালেক ও তার বাবা ২ জনেরই বয়স ১৩ বছর করে বৃদ্ধি পাবে।}

\emph{∴ ১৩ বছর পর তাদের বয়সের সমষ্টি হবে = ৪০ + (১৩ × ২) = ৪০ + ২৬ = ৬৬
বছর।}
\end{quote}

\textbf{রমিজ এবং লুবনার মোট বয়স ৬৫ বছর। রমিজের বয়সের দ্বিগুণ লুবনার বয়সের
তিনগুণের সমান হলে রমিজের বয়স কত?}

ক) ২৬

\textbf{\hl{খ) ৩৯ ✔}}

গ) ৪২

ঘ) ৪৬

\textbf{\hl{✔ সঠিক উত্তর: খ) ৩৯}}

\begin{quote}
\textbf{ব্যাখ্যা:} \emph{ধরি, রমিজের বয়স x বছর।}

\emph{∴ লুবনার বয়স = (৬৫ − x) বছর।}

\emph{প্রশ্নমতে,}

\emph{২x = ৩(৬৫ − x)}

\emph{⇒ ২x = ১৯৫ − ৩x}

\emph{⇒ ২x + ৩x = ১৯৫}

\emph{⇒ ৫x = ১৯৫}

\emph{⇒ x = (১৯৫/৫)}

\emph{∴ x = ৩৯}

\emph{সুতরাং, রমিজের বয়স ৩৯ বছর।}
\end{quote}

\textbf{Naizel is twenty years older than Neketa. In three yeas Naizel
will be twice as old as Neketa will be. How old is Neketa now?}

ক) 24

\textbf{\hl{খ) 17 ✔}}

গ) 16

ঘ) 89

\textbf{\hl{✔ সঠিক উত্তর: খ) 17}}

\begin{quote}
\textbf{ব্যাখ্যা:} \emph{Let Neketa\textquotesingle s present age = x
years.}

\emph{Naizel\textquotesingle s present age = (x + 20) years.}

\emph{In 3 years:}

\emph{Neketa\textquotesingle s age = (x + 3) years}

\emph{Naizel\textquotesingle s age = (x + 20 + 3) = (x + 23) years}

\emph{According to the question,}

\emph{2(x + 3) = x + 23}

\emph{⇒ 2x + 6 = x + 23}

\emph{⇒ 2x − x = 23 − 6}

\emph{∴ x = 17}

\emph{So, Neketa is 17 years old now.}
\end{quote}

\textbf{যদি ক, খ থেকে ৫ বছরের বড় এবং খ এর বয়স গ এর বয়সের দ্বিগুণ এবং তাদের
মোট বয়স ৫০ বছর হলে ক এর বয়স কত?}

ক) ৯ বছর

\textbf{\hl{খ) ২৩ বছর ✔}}

গ) ১৮ বছর

ঘ) ১৩ বছর

\textbf{\hl{✔ সঠিক উত্তর: খ) ২৩ বছর}}

\begin{quote}
\textbf{ব্যাখ্যা:} \emph{ধরি, গ এর বয়স = x বছর।}

\emph{∴ খ এর বয়স = ২x বছর}

\emph{এবং ক এর বয়স = (২x + ৫) বছর।}

\emph{প্রশ্নমতে,}

\emph{x + ২x + (২x + ৫) = ৫০}

\emph{⇒ ৫x + ৫ = ৫০}

\emph{⇒ ৫x = ৫০ − ৫}

\emph{⇒ ৫x = ৪৫}

\emph{⇒ x = (৪৫/৫)}

\emph{∴ x = ৯}

\emph{∴ ক এর বয়স = (২ × ৯ + ৫) বছর = (১৮ + ৫) বছর = ২৩ বছর।}
\end{quote}

\textbf{মীরার বয়স অপুর বয়সের ৪ গুণ। অপুর বর্তমান বয়স ৪ বছর। যখন মীরার বয়স
অপুর বয়সের দ্বিগুণ হবে তখন মীরার বয়স কত বছর হবে?}

ক) ১৬

খ) ২০

\textbf{\hl{গ) ২৪ ✔}}

ঘ) কোনোটিই নয়

\textbf{\hl{✔ সঠিক উত্তর: গ) ২৪}}

\begin{quote}
\textbf{ব্যাখ্যা:} \emph{অপুর বর্তমান বয়স = ৪ বছর।}

\emph{∴ মীরার বর্তমান বয়স = (৪ × ৪) = ১৬ বছর।}

\emph{ধরি, x বছর পর মীরার বয়স অপুর বয়সের দ্বিগুণ হবে।}

\emph{প্রশ্নমতে,}

\emph{১৬ + x = ২(৪ + x)}

\emph{⇒ ১৬ + x = ৮ + ২x}

\emph{⇒ ২x − x = ১৬ − ৮}

\emph{∴ x = ৮}

\emph{∴ তখন মীরার বয়স হবে = (১৬ + ৮) বা ২৪ বছর।}
\end{quote}

\textbf{কামাল বশীরের ৫ বছরের ছোট। বশীর টুটুলের ৫ বছরের বড়। কামাল ও টুটুলের
বয়সের পার্থক্য কত?}

ক) ৫

খ) ১০

\textbf{\hl{গ) ০ ✔}}

ঘ) কোনোটিই নয়

\textbf{\hl{✔ সঠিক উত্তর: গ) ০}}

\begin{quote}
\textbf{ব্যাখ্যা:} \emph{ধরি, টুটুলের বয়স = x বছর।}

\emph{∴ বশীরের বয়স = (x + ৫) বছর।}

\emph{∴ কামালের বয়স = (x + ৫ − ৫) = x বছর।}

\emph{∴ কামাল ও টুটুলের বয়সের পার্থক্য = x − x = ০ বছর।}
\end{quote}

\textbf{স্ত্রী স্বামীর চেয়ে ৫ বছরের ছোট। স্ত্রীর বয়স ছেলের বয়সের চারগুণ। চার
বছর পর ছেলের বয়স হবে ১১। বর্তমানে স্বামীর বয়স কত?}

ক) ৪৮ বছর

খ) ৫২ বছর

গ) ৪৫ বছর

\textbf{\hl{ঘ) ৩৩ বছর ✔}}

\textbf{\hl{✔ সঠিক উত্তর: ঘ) ৩৩ বছর}}

\begin{quote}
\textbf{ব্যাখ্যা:} \emph{৪ বছর পর ছেলের বয়স হবে = ১১ বছর।}

\emph{∴ ছেলের বর্তমান বয়স = (১১ − ৪) = ৭ বছর।}

\emph{স্ত্রীর বয়স ছেলের বয়সের ৪ গুণ,}

\emph{∴ স্ত্রীর বর্তমান বয়স = (৭ × ৪) = ২৮ বছর।}

\emph{স্ত্রী স্বামীর চেয়ে ৫ বছরের ছোট (অর্থাৎ স্বামী স্ত্রীর চেয়ে ৫ বছরের বড়),}

\emph{∴ বর্তমানে স্বামীর বয়স = (২৮ + ৫) = ৩৩ বছর।}
\end{quote}

\textbf{Eighteen years ago, a father was three times as old as his son.
Now the father is only twice as old as his son. Then the sum of the
present ages of the son and the father is:/ ১৮ বছর আগে পিতার বয়স ছিল
পুত্রের বয়সের ৩ গুণ। বর্তমানে পিতার বয়স পুত্রের বয়সের ২ গুণ। বর্তমানে পিতা ও
পুত্রের বয়সের সমষ্টি কত?}

ক) 54

খ) 72

গ) 105

\textbf{\hl{ঘ) 108 ✔}}

\textbf{\hl{✔ সঠিক উত্তর: ঘ) 108}}

\begin{quote}
\textbf{ব্যাখ্যা:} \emph{Let the present age of the son = x years.}

\emph{So, the present age of the father = 2x years.}

\emph{18 years ago:}

\emph{Son\textquotesingle s age = (x − 18) years}

\emph{Father\textquotesingle s age = (2x − 18) years}

\emph{According to the condition,}

\emph{2x − 18 = 3(x − 18)}

\emph{⇒ 2x − 18 = 3x − 54}

\emph{⇒ 3x − 2x = 54 − 18}

\emph{∴ x = 36}

\emph{So, son\textquotesingle s present age = 36 years}

\emph{Father\textquotesingle s present age = (36 × 2) = 72 years}

\emph{∴ Sum of present ages = (36 + 72) years = 108 years.}
\end{quote}

\textbf{পিতার বর্তমান বয়স পুত্রের বয়সের চারগুণ। ৬ বছর পূর্বে পিতার বয়স পুত্রের
বয়সের ১০ গুণ ছিল। পুত্রের বর্তমান বয়স কত?}

ক) ১৪ বছর

খ) ৮ বছর

\textbf{\hl{গ) ৯ বছর ✔}}

ঘ) ১০ বছর

\textbf{\hl{✔ সঠিক উত্তর: গ) ৯ বছর}}

\begin{quote}
\textbf{ব্যাখ্যা:} \emph{ধরি, পুত্রের বর্তমান বয়স = x বছর।}

\emph{∴ পিতার বর্তমান বয়স = ৪x বছর।}

\emph{৬ বছর পূর্বে:}

\emph{পুত্রের বয়স ছিল = (x − ৬) বছর}

\emph{পিতার বয়স ছিল = (৪x − ৬) বছর}

\emph{প্রশ্নমতে,}

\emph{৪x − ৬ = ১০(x − ৬)}

\emph{⇒ ৪x − ৬ = ১০x − ৬০}

\emph{⇒ ১০x − ৪x = ৬০ − ৬}

\emph{⇒ ৬x = ৫৪}

\emph{⇒ x = (৫৪/৬)}

\emph{∴ x = ৯}

\emph{∴ পুত্রের বর্তমান বয়স ৯ বছর।}
\end{quote}

\textbf{বর্তমানে পিতার বয়স পুত্রের বয়সের ৫ গুণ। তিন বছর পরে পিতার বয়স পুত্রের
বয়সের ৪ গুণ হয়। পিতা ও পুত্রের বর্তমান বয়স কত?}

ক) ৩৫ বছর, ৭ বছর

খ) ৫০ বছর, ১০ বছর

\textbf{\hl{গ) ৪৫ বছর, ৯ বছর ✔}}

ঘ) ২৫বছর, ৫বছর

\textbf{\hl{✔ সঠিক উত্তর: গ) ৪৫ বছর, ৯ বছর}}

\begin{quote}
\textbf{ব্যাখ্যা:} \emph{ধরি, বর্তমানে পুত্রের বয়স = x বছর।}

\emph{∴ পিতার বর্তমান বয়স = ৫x বছর।}

\emph{৩ বছর পর:}

\emph{পুত্রের বয়স হবে = (x + ৩) বছর}

\emph{পিতার বয়স হবে = (৫x + ৩) বছর}

\emph{প্রশ্নমতে,}

\emph{৫x + ৩ = ৪(x + ৩)}

\emph{⇒ ৫x + ৩ = ৪x + ১২}

\emph{⇒ ৫x − ৪x = ১২ − ৩}

\emph{∴ x = ৯}

\emph{∴ পুত্রের বর্তমান বয়স ৯ বছর এবং পিতার বর্তমান বয়স = (৫ × ৯) = ৪৫ বছর।}
\end{quote}

\textbf{পিতা ও পুত্রের বয়সের সমষ্টি ৮০ বছর। পিতার বয়স পুত্রের বয়সের তিনগুণ হলে
পিতার বয়স কত?}

\textbf{\hl{ক) ৬০ বছর ✔}}

খ) ৬৪ বছর

গ) ৫৬ বছর

ঘ) ৫০ বছর

\textbf{\hl{✔ সঠিক উত্তর: ক) ৬০ বছর}}

\begin{quote}
\textbf{ব্যাখ্যা:} \emph{ধরি, পুত্রের বয়স = x বছর।}

\emph{∴ পিতার বয়স = ৩x বছর।}

\emph{শর্তমতে,}

\emph{৩x + x = ৮০}

\emph{⇒ ৪x = ৮০}

\emph{⇒ x = (৮০/৪)}

\emph{∴ x = ২০}

\emph{∴ পিতার বয়স = (৩ × ২০) বছর = ৬০ বছর।}
\end{quote}

\textbf{পিতা-পুত্রের বর্তমান বয়সের সমষ্টি ৮০ বছর। আগামী তিন বছর পর পিতার বয়স
৬০ বছর হলে তখন পুত্রের বয়স কত হবে?}

ক) ২৩ বছর

\textbf{\hl{খ) ২৬ বছর ✔}}

গ) ৩০ বছর

ঘ) ৩২ বছর

\textbf{\hl{✔ সঠিক উত্তর: খ) ২৬ বছর}}

\begin{quote}
\textbf{ব্যাখ্যা:} \emph{৩ বছর পর পিতা ও পুত্রের বয়সের সমষ্টি হবে = \{৮০ + (৩
+ ৩)\} বছর = ৮৬ বছর।}

\emph{৩ বছর পর পিতার বয়স হবে = ৬০ বছর।}

\emph{∴ ৩ বছর পর পুত্রের বয়স হবে = (৮৬ − ৬০) বছর = ২৬ বছর।}
\end{quote}

\textbf{পিতা ও কন্যার বয়সের সমষ্টি ৮০ বছর। পিতার বয়স কন্যার বয়সের ৪ গুণ,
কন্যার বয়স কত?}

ক) ১২ বছর

খ) ১৪ বছর

\textbf{\hl{গ) ১৬ বছর ✔}}

ঘ) ২০ বছর

\textbf{\hl{✔ সঠিক উত্তর: গ) ১৬ বছর}}

\begin{quote}
\textbf{ব্যাখ্যা:} \emph{ধরি, কন্যার বয়স = x বছর।}

\emph{∴ পিতার বয়স = ৪x বছর।}

\emph{প্রশ্নমতে,}

\emph{x + ৪x = ৮০}

\emph{⇒ ৫x = ৮০}

\emph{⇒ x = (৮০/৫)}

\emph{∴ x = ১৬}

\emph{∴ কন্যার বয়স = ১৬ বছর।}
\end{quote}

\textbf{মাতার ৩০ বছর বয়সে কন্যার জন্ম হয়। মাতার কত বছর বয়সে তার বয়স কন্যার
বয়সের দ্বিগুণ হবে?}

ক) ৩০

খ) ৫০

\textbf{\hl{গ) ৬০ ✔}}

ঘ) ৭০

\textbf{\hl{✔ সঠিক উত্তর: গ) ৬০}}

\begin{quote}
\textbf{ব্যাখ্যা:} \emph{মাতা ও কন্যার বয়সের পার্থক্য সবসময় ৩০ বছর থাকবে।}

\emph{ধরি, কন্যার বয়স x বছর হলে মাতার বয়স হবে (৩০ + x) বছর।}

\emph{প্রশ্নমতে,}

\emph{৩০ + x = ২x}

\emph{⇒ ২x − x = ৩০}

\emph{∴ x = ৩০}

\emph{∴ মাতার বয়স যখন (৩০ + ৩০) = ৬০ বছর, তখন তার বয়স কন্যার বয়সের দ্বিগুণ
হবে।}
\end{quote}

\textbf{কালাম সাহেব তার স্ত্রীর চেয়ে ৬ বছরের বড়। তার স্ত্রীর বয়স তার ছেলের
বয়সের ৪ গুণ। ৫ বছর পর ছেলের বয়স ১২ বছর হলে ৩ বছর পরে কালাম সাহেবের বয়স কত
হবে?}

ক) ২৮ বছর

খ) ৩৪ বছর

\textbf{\hl{গ) ৩৭ বছর ✔}}

ঘ) ৩৯ বছর

\textbf{\hl{✔ সঠিক উত্তর: গ) ৩৭ বছর}}

\begin{quote}
\textbf{ব্যাখ্যা:} \emph{৫ বছর পর ছেলের বয়স ১২ বছর হলে,}

\emph{ছেলের বর্তমান বয়স = (১২ − ৫) = ৭ বছর।}

\emph{স্ত্রীর বর্তমান বয়স = (৭ × ৪) = ২৮ বছর।}

\emph{কালাম সাহেবের বর্তমান বয়স = (২৮ + ৬) = ৩৪ বছর।}

\emph{∴ ৩ বছর পর কালাম সাহেবের বয়স হবে = (৩৪ + ৩) = ৩৭ বছর।}
\end{quote}

\textbf{বাবার বয়স ছেলের বয়সের ৪ গুণ। ২০ বছর পরে ছেলের বয়স বাবার বয়সের
অর্ধেক হবে। বাবার বর্তমান বয়স কত?}

ক) ২০

খ) ২৫

গ) ৩০

\textbf{\hl{ঘ) ৪০ ✔}}

\textbf{\hl{✔ সঠিক উত্তর: ঘ) ৪০}}

\begin{quote}
\textbf{ব্যাখ্যা:} \emph{ধরি, ছেলের বর্তমান বয়স = x বছর।}

\emph{∴ বাবার বর্তমান বয়স = ৪x বছর।}

\emph{২০ বছর পর:}

\emph{ছেলের বয়স = (x + ২০) বছর}

\emph{বাবার বয়স = (৪x + ২০) বছর}

\emph{প্রশ্নমতে,}

\emph{৪x + ২০ = ২(x + ২০)}

\emph{⇒ ৪x + ২০ = ২x + ৪০}

\emph{⇒ ৪x − ২x = ৪০ − ২০}

\emph{⇒ ২x = ২০}

\emph{⇒ x = (২০/২)}

\emph{∴ x = ১০}

\emph{∴ বাবার বর্তমান বয়স = (৪ × ১০) = ৪০ বছর।}
\end{quote}

\textbf{পিতার বয়স পুত্রের বয়স অপেক্ষা ২ বছর বেশি। পিতার বয়স ৬২ হলে পুত্রের
বয়স কত?}

ক) ২৫

\textbf{\hl{খ) ৩০ ✔}}

গ) ৩৫

ঘ) ৪০

\textbf{\hl{✔ সঠিক উত্তর: খ) ৩০}}

\begin{quote}
\textbf{ব্যাখ্যা:} \emph{প্রশ্নটিতে ভাষাগত ভুল রয়েছে; সঠিক বাক্যটি হবে: "পিতার
বয়স পুত্রের বয়সের দ্বিগুণ অপেক্ষা ২ বছর বেশি।"}

\emph{ধরি, পুত্রের বয়স = x বছর।}

\emph{প্রশ্নমতে,}

\emph{২x + ২ = ৬২}

\emph{⇒ ২x = ৬২ − ২}

\emph{⇒ ২x = ৬০}

\emph{⇒ x = (৬০/২)}

\emph{∴ x = ৩০}

\emph{∴ পুত্রের বয়স ৩০ বছর।}
\end{quote}

\textbf{এক ব্যক্তি তার স্ত্রীর চেয়ে ৫ বছরের বড় এবং তার স্ত্রীর বয়স তাদের ছেলের
বয়সের ৫ গুণ। ৬ বছর পরে ছেলের বয়স ১৪ হলে লোকটির বর্তমান বয়স কত বছর?}

ক) ৩৫

খ) ৪০

গ) ৩০

\textbf{\hl{ঘ) ৪৫ ✔}}

\textbf{\hl{✔ সঠিক উত্তর: ঘ) ৪৫}}

\begin{quote}
\textbf{ব্যাখ্যা:} \emph{৬ বছর পর ছেলের বয়স ১৪ বছর হলে,}

\emph{ছেলের বর্তমান বয়স = (১৪ − ৬) = ৮ বছর।}

\emph{স্ত্রীর বর্তমান বয়স = (৮ × ৫) = ৪০ বছর।}

\emph{∴ লোকটির বর্তমান বয়স = (৪০ + ৫) = ৪৫ বছর।}
\end{quote}

\textbf{রিনা মীমের চেয়ে ১০ বছরের বড়। ৭ বছর পর, রিনার বয়স মীমের বয়সের
দ্বিগুণ হবে। রিনার বর্তমান বয়স কত?}

ক) ১৭

খ) ১৬

\textbf{\hl{গ) ১৩ ✔}}

ঘ) ১১

\textbf{\hl{✔ সঠিক উত্তর: গ) ১৩}}

\begin{quote}
\textbf{ব্যাখ্যা:} \emph{ধরি, মীমের বর্তমান বয়স = x বছর।}

\emph{∴ রিনার বর্তমান বয়স = (x + ১০) বছর।}

\emph{৭ বছর পর:}

\emph{মীমের বয়স হবে = (x + ৭) বছর}

\emph{রিনার বয়স হবে = (x + ১০ + ৭) = (x + ১৭) বছর}

\emph{প্রশ্নমতে,}

\emph{২(x + ৭) = x + ১৭}

\emph{⇒ ২x + ১৪ = x + ১৭}

\emph{⇒ ২x − x = ১৭ − ১৪}

\emph{∴ x = ৩}

\emph{∴ রিনার বর্তমান বয়স = (৩ + ১০) = ১৩ বছর।}
\end{quote}

\textbf{এক ব্যক্তি তার স্ত্রীর চেয়ে ৫ বছরের বড়। তার স্ত্রীর বয়স ছেলের বয়সের ৪
গুণ। ৫ বছর পরে ছেলের বয়স ১২ বছর হলে, বর্তমানে ঐ ব্যক্তির বয়স কত?}

ক) ২৮ বছর

খ) ৩০ বছর

\textbf{\hl{গ) ৩৩ বছর ✔}}

ঘ) ৬৫ বছর

\textbf{\hl{✔ সঠিক উত্তর: গ) ৩৩ বছর}}

\begin{quote}
\textbf{ব্যাখ্যা:} \emph{৫ বছর পর ছেলের বয়স ১২ বছর হলে,}

\emph{ছেলের বর্তমান বয়স = (১২ − ৫) = ৭ বছর।}

\emph{স্ত্রীর বর্তমান বয়স = (৭ × ৪) = ২৮ বছর।}

\emph{∴ বর্তমানে ঐ ব্যক্তির বয়স = (২৮ + ৫) = ৩৩ বছর।}
\end{quote}

\textbf{পিতার বয়স পুত্রের দ্বিগুণ অপেক্ষা ২ বছর বেশি। পিতার বয়স ৬২ বছর হলে
পুত্রের বয়স কত?}

ক) ২৫ বছর

\textbf{\hl{খ) ৩০ বছর ✔}}

গ) ৩৫ বছর

ঘ) ৪০ বছর

\textbf{\hl{✔ সঠিক উত্তর: খ) ৩০ বছর}}

\begin{quote}
\textbf{ব্যাখ্যা:} \emph{ধরি, পুত্রের বয়স = x বছর।}

\emph{∴ পিতার বয়স = (২x + ২) বছর।}

\emph{শর্তমতে,}

\emph{২x + ২ = ৬২}

\emph{⇒ ২x = ৬২ − ২}

\emph{⇒ ২x = ৬০}

\emph{⇒ x = (৬০/২)}

\emph{∴ x = ৩০}

\emph{∴ পুত্রের বয়স ৩০ বছর।}
\end{quote}

\textbf{জনি রনির তুলনায় ২০ বছর বড়। ১০ বছর পর রনি এর বয়স জনি এর অর্ধেক হবে।
রনি এর বয়স কত?}

ক) 5 years

\textbf{\hl{খ) 10 years ✔}}

গ) 20 years

ঘ) 25 years

\textbf{\hl{✔ সঠিক উত্তর: খ) 10 years}}

\begin{quote}
\textbf{ব্যাখ্যা:} \emph{ধরি, রনির বয়স = R বছর।}

\emph{∴ জনির বয়স = J = (R + ২০) বছর।}

\emph{১০ বছর পর:}

\emph{রনির বয়স হবে = (R + ১০) বছর}

\emph{জনির বয়স হবে = (J + ১০) = (R + ২০ + ১০) = (R + ৩০) বছর}

\emph{প্রশ্নমতে,}

\emph{R + ১০ = (R+৩০/২)}

\emph{⇒ ২(R + ১০) = R + ৩০}

\emph{⇒ ২R + ২০ = R + ৩০}

\emph{⇒ ২R − R = ৩০ − ২০}

\emph{∴ R = ১০}

\emph{∴ রনির বয়স 10 years।}
\end{quote}

\textbf{পিতার বয়স পুত্রের বয়সের ৪ গুণ। ৫ বছর পর তাদের বয়সের সমষ্টি ৬০ বছর
হলে পিতার বর্তমান বয়স কত বছর?}

ক) ৩০

\textbf{\hl{খ) ৪০ ✔}}

গ) ৪৮

ঘ) কোনোটিই নয়

\textbf{\hl{✔ সঠিক উত্তর: খ) ৪০}}

\begin{quote}
\textbf{ব্যাখ্যা:} \emph{ধরি, পুত্রের বর্তমান বয়স = x বছর।}

\emph{∴ পিতার বর্তমান বয়স = ৪x বছর।}

\emph{৫ বছর পর তাদের বয়সের সমষ্টি ৬০ বছর হলে,}

\emph{(৪x + ৫) + (x + ৫) = ৬০}

\emph{⇒ ৫x + ১০ = ৬০}

\emph{⇒ ৫x = ৬০ − ১০}

\emph{⇒ ৫x = ৫০}

\emph{⇒ x = (৫০/৫)}

\emph{∴ x = ১০}

\emph{∴ পিতার বর্তমান বয়স = (৪ × ১০) = ৪০ বছর।}
\end{quote}

\textbf{পিতার বয়স পুত্রের বয়সের ৪ গুণ। ৪ বছর পূর্বে পিতা পুত্রের বয়সের সমষ্টি
ছিল ৫২ বছর। পিতার বর্তমান বয়স কত?}

ক) ৩২

খ) ৩৮

\textbf{\hl{গ) ৪৮ ✔}}

ঘ) ৫২

\textbf{\hl{✔ সঠিক উত্তর: গ) ৪৮}}

\begin{quote}
\textbf{ব্যাখ্যা:} \emph{ধরি, বর্তমান পুত্রের বয়স = x বছর।}

\emph{∴ বর্তমান পিতার বয়স = ৪x বছর।}

\emph{৪ বছর পূর্বে:}

\emph{পুত্রের বয়স = (x − ৪) বছর}

\emph{পিতার বয়স = (৪x − ৪) বছর}

\emph{প্রশ্নমতে,}

\emph{(x − ৪) + (৪x − ৪) = ৫২}

\emph{⇒ ৫x − ৮ = ৫২}

\emph{⇒ ৫x = ৫২ + ৮}

\emph{⇒ ৫x = ৬০}

\emph{⇒ x = (৬০/৫)}

\emph{∴ x = ১২}

\emph{∴ বর্তমানে পিতার বয়স = (৪ × ১২) = ৪৮ বছর।}
\end{quote}

\textbf{পিতার বর্তমান বয়স পুত্রের বয়সের তিনগুণ। ৫ বছর আগে পিতার বয়স পুত্রের
বয়সের ৪ গুণ ছিল। পিতা ও পুত্রের বর্তমান বয়স কত বছর?}

ক) ৪৮, ১৬

খ) ২৪, ৮

\textbf{\hl{গ) ৪৫, ১৫ ✔}}

ঘ) ৩৬, ১২

\textbf{\hl{✔ সঠিক উত্তর: গ) ৪৫, ১৫}}

\begin{quote}
\textbf{ব্যাখ্যা:} \emph{ধরি, পুত্রের বর্তমান বয়স = x বছর।}

\emph{∴ পিতার বর্তমান বয়স = ৩x বছর।}

\emph{৫ বছর আগে:}

\emph{পুত্রের বয়স ছিল = (x − ৫) বছর}

\emph{পিতার বয়স ছিল = (৩x − ৫) বছর}

\emph{শর্তমতে,}

\emph{৩x − ৫ = ৪(x − ৫)}

\emph{⇒ ৩x − ৫ = ৪x − ২০}

\emph{⇒ ৪x − ৩x = ২০ − ৫}

\emph{∴ x = ১৫}

\emph{∴ পুত্রের বর্তমান বয়স = ১৫ বছর এবং পিতার বর্তমান বয়স = (১৫ × ৩) = ৪৫
বছর।}
\end{quote}

\textbf{পিতা ও পুত্রের বর্তমান বয়সের সমষ্টি ৪২ বছর এবং অন্তর ২২ বছর। পুত্রের
বর্তমান বয়স কত?}

\textbf{\hl{ক) ১০ বছর ✔}}

খ) ২২ বছর

গ) ৪২ বছর

ঘ) ৩২ বছর

\textbf{\hl{✔ সঠিক উত্তর: ক) ১০ বছর}}

\begin{quote}
\textbf{ব্যাখ্যা:} \emph{ধরি, পুত্রের বর্তমান বয়স = x বছর।}

\emph{∴ পিতার বর্তমান বয়স = (x + ২২) বছর।}

\emph{শর্তমতে,}

\emph{x + (x + ২২) = ৪২}

\emph{⇒ ২x + ২২ = ৪২}

\emph{⇒ ২x = ৪২ − ২২}

\emph{⇒ ২x = ২০}

\emph{⇒ x = (২০/২)}

\emph{∴ x = ১০}

\emph{∴ পুত্রের বর্তমান বয়স ১০ বছর।}
\end{quote}

\textbf{৪০ বছরের একজন ব্যক্তির ৬, ৩ এবং ১ বছরের তিনটি পুত্র আছে। কত বছর পরে
তিন পুত্রের বয়সের সমষ্টি পিতার বয়সের ৮০\% হবে?}

\textbf{\hl{ক) ১০ ✔}}

খ) ১৪

গ) ১৭

ঘ) ২১

\textbf{\hl{✔ সঠিক উত্তর: ক) ১০}}

\begin{quote}
\textbf{ব্যাখ্যা:} \emph{বর্তমান পিতার বয়স = ৪০ বছর।}

\emph{তিন পুত্রের বর্তমান বয়স = ৬, ৩ এবং ১ বছর।}

\emph{ধরি, x বছর পর তিন পুত্রের বয়সের সমষ্টি পিতার বয়সের ৮০\% হবে।}

\emph{x বছর পর:}

\emph{পিতার বয়স হবে = (৪০ + x) বছর}

\emph{তিন পুত্রের বয়স হবে যথাক্রমে = (৬ + x), (৩ + x) এবং (১ + x) বছর}

\emph{তিন পুত্রের বয়সের সমষ্টি হবে = (৬ + x) + (৩ + x) + (১ + x) = (১০ +
৩x) বছর}

\emph{প্রশ্নমতে,}

\emph{১০ + ৩x = (৪০ + x) এর ৮০\%}

\emph{⇒ ১০ + ৩x = (৪০ + x) × (৮০/১০০)}

\emph{⇒ ১০ + ৩x = (৪০ + x) × (৪/৫)}

\emph{⇒ ৫(১০ + ৩x) = ৪(৪০ + x)}

\emph{⇒ ৫০ + ১৫x = ১৬০ + ৪x}

\emph{⇒ ১৫x − ৪x = ১৬০ − ৫০}

\emph{⇒ ১১x = ১১০}

\emph{⇒ x = (১১০/১১)}

\emph{∴ x = ১০}

\emph{∴ ১০ বছর পর তিন পুত্রের বয়সের সমষ্টি পিতার বয়সের ৮০\% হবে।}
\end{quote}

\textbf{কলির বয়স ডলির বয়স থেকে সাত বছর কম। ডলির বয়সের তিন গুণ থেকে কলির
বয়স বাদ দিলে ৪৩ বছর হয়। ডলির বয়স কত?}

ক) ১১

খ) ১৪

গ) ১৬

\textbf{\hl{ঘ) ১৮ ✔}}

\textbf{\hl{✔ সঠিক উত্তর: ঘ) ১৮}}

\begin{quote}
\textbf{ব্যাখ্যা:} \emph{ধরি, ডলির বয়স = x বছর।}

\emph{∴ কলির বয়স = (x − ৭) বছর।}

\emph{প্রশ্নমতে,}

\emph{৩x − (x − ৭) = ৪৩}

\emph{⇒ ৩x − x + ৭ = ৪৩}

\emph{⇒ ২x = ৪৩ − ৭}

\emph{⇒ ২x = ৩৬}

\emph{⇒ x = (৩৬/২)}

\emph{∴ x = ১৮}

\emph{∴ ডলির বয়স ১৮ বছর।}
\end{quote}

\textbf{পিতা ও পুত্রের বয়সের সমষ্টি 50 বছর। 20 বছর পরে পিতার বয়স পুত্রের
বয়সের দ্বিগুণ হলে, পিতার বর্তমান বয়স কত?}

ক) 30 বছর

খ) 35 বছর

\textbf{\hl{গ) 40 বছর ✔}}

ঘ) 50 বছর

\textbf{\hl{✔ সঠিক উত্তর: গ) 40 বছর}}

\begin{quote}
\textbf{ব্যাখ্যা:} \emph{ধরি, পুত্রের বর্তমান বয়স = x বছর।}

\emph{∴ পিতার বর্তমান বয়স = (50 − x) বছর।}

\emph{20 বছর পর:}

\emph{পুত্রের বয়স হবে = (x + 20) বছর}

\emph{পিতার বয়স হবে = (50 − x + 20) = (70 − x) বছর}

\emph{প্রশ্নমতে,}

\emph{70 − x = 2(x + 20)}

\emph{⇒ 70 − x = 2x + 40}

\emph{⇒ 2x + x = 70 − 40}

\emph{⇒ 3x = 30}

\emph{⇒ x = (30/3)}

\emph{∴ x = 10}

\emph{∴ পিতার বর্তমান বয়স = (50 − 10) বছর = 40 বছর।}
\end{quote}

\textbf{পিতা ও তিন পুত্রের বয়স অপেক্ষা মাতা ও উক্ত তিন পুত্রের বয়সের গড় ১
(১/২) বছর কম। মাতার বয়স ৩০ বছর হলে পিতার বয়স কত?}

ক) ২৮ (১/২) বছর

খ) ৩০ বছর

গ) ৩১ (১/২) বছর

\textbf{\hl{ঘ) ৩৬ বছর ✔}}

\textbf{\hl{✔ সঠিক উত্তর: ঘ) ৩৬ বছর}}

\begin{quote}
\textbf{ব্যাখ্যা:} \emph{পিতা ও তিন পুত্র অর্থাৎ মোট জনসংখ্যা = ৪ জন।}

\emph{মাতা ও তিন পুত্র অর্থাৎ মোট জনসংখ্যা = ৪ জন।}

\emph{পিতা ও তিন পুত্রের বয়সের গড় থেকে মাতা ও তিন পুত্রের বয়সের গড় ১ (১/২)
বছর বা (৩/২) বছর কম।}

\emph{∴ ৪ জনের মোট বয়সের পার্থক্য = ৪ × (৩/২) = ৬ বছর।}

\emph{যেহেতু তিন পুত্রের বয়স উভয় ক্ষেত্রেই সমান, তাই পিতা ও মাতার বয়সের
পার্থক্য = ৬ বছর।}

\emph{অর্থাৎ, পিতা − মাতা = ৬ বছর}

\emph{⇒ পিতা − ৩০ = ৬}

\emph{∴ পিতার বয়স = (৩০ + ৬) = ৩৬ বছর।}
\end{quote}

\textbf{৪ বছর আগে রিমঝিমের বয়স রোদেলার বয়সের দ্বিগুণ। ৪ বছর পর রোদেলার বয়স
রিমঝিমের বয়সের (৩/৪) গুণ হবে। রিমঝিমের বয়স কত?}

ক) ১৪ বছর

\textbf{\hl{খ) ১২ বছর ✔}}

গ) ২২ বছর

ঘ) ২৪ বছর

\textbf{\hl{✔ সঠিক উত্তর: খ) ১২ বছর}}

\begin{quote}
\textbf{ব্যাখ্যা:} \emph{ধরি, ৪ বছর আগে রোদেলার বয়স ছিল = x বছর।}

\emph{∴ ৪ বছর আগে রিমঝিমের বয়স ছিল = ২x বছর।}

\emph{বর্তমান বয়স:}

\emph{রোদেলার বর্তমান বয়স = (x + ৪) বছর}

\emph{রিমঝিমের বর্তমান বয়স = (২x + ৪) বছর}

\emph{৪ বছর পর:}

\emph{রোদেলার বয়স হবে = (x + ৪ + ৪) = (x + ৮) বছর}

\emph{রিমঝিমের বয়স হবে = (২x + ৪ + ৪) = (২x + ৮) বছর}

\emph{প্রশ্নমতে,}

\emph{x + ৮ = (২x + ৮) × (৩/৪)}

\emph{⇒ ৪(x + ৮) = ৩(২x + ৮)}

\emph{⇒ ৪x + ৩২ = ৬x + ২৪}

\emph{⇒ ৬x − ৪x = ৩২ − ২৪}

\emph{⇒ ২x = ৮}

\emph{∴ x = ৪}

\emph{∴ রিমঝিমের বর্তমান বয়স = (২ × ৪) + ৪ = ৮ + ৪ = ১২ বছর।}
\end{quote}

\textbf{পিতা ও মাতার বয়সের গড় ৪৫ বছর। আবার পিতা, মাতা ও ১ পুত্রের বয়সের গড়
৩৬ বছর। পুত্রের বয়স কত?}

ক) ৯ বছর

খ) ১৫ বছর

\textbf{\hl{গ) ১৮ বছর ✔}}

ঘ) ১৪ বছর

\textbf{\hl{✔ সঠিক উত্তর: গ) ১৮ বছর}}

\begin{quote}
\textbf{ব্যাখ্যা:} \emph{পিতা ও মাতার বয়সের গড় = ৪৫ বছর।}

\emph{∴ পিতা ও মাতার বয়সের সমষ্টি = (৪৫ × ২) বছর = ৯০ বছর।}

\emph{পিতা, মাতা ও ১ পুত্রের বয়সের গড় = ৩৬ বছর।}

\emph{∴ পিতা, মাতা ও ১ পুত্রের বয়সের সমষ্টি = (৩৬ × ৩) বছর = ১০৮ বছর।}

\emph{∴ পুত্রের বয়স = (১০৮ − ৯০) বছর = ১৮ বছর।}
\end{quote}

\textbf{পিতা, মাতা ও পুত্রের বয়সের গড় 37 বছর। আবার পিতা ও পুত্রের বয়সের গড়
35 বছর। মাতার বয়স কত?}

ক) 38 বছর

\textbf{\hl{খ) 41 বছর ✔}}

গ) 45 বছর

ঘ) 48 বছর

\textbf{\hl{✔ সঠিক উত্তর: খ) 41 বছর}}

\begin{quote}
\textbf{ব্যাখ্যা:} \emph{পিতা, মাতা ও পুত্রের মোট বয়স = (37 × 3) বছর = 111
বছর।}

\emph{পিতা ও পুত্রের মোট বয়স = (35 × 2) বছর = 70 বছর।}

\emph{∴ মাতার বয়স = (111 − 70) বছর = 41 বছর।}
\end{quote}

\textbf{পিতা ও ২ সন্তানের বয়সের গড় ৩৭ বছর। পিতা, মাতা ও ঐ ২ সন্তানের বয়সের
গড় ৩৬ বছর। মাতার বয়স কত বছর?}

\textbf{\hl{ক) ৩৩ ✔}}

খ) ৩৪

গ) ৩৫

ঘ) ৩৬

\textbf{\hl{✔ সঠিক উত্তর: ক) ৩৩}}

\begin{quote}
\textbf{ব্যাখ্যা:} \emph{পিতা ও ২ সন্তান অর্থাৎ ৩ জনের বয়সের সমষ্টি = (৩৭ ×
৩) বছর = ১১১ বছর।}

\emph{পিতা, মাতা ও ২ সন্তান অর্থাৎ ৪ জনের বয়সের সমষ্টি = (৩৬ × ৪) বছর = ১৪৪
বছর।}

\emph{∴ মাতার বয়স = (১৪৪ − ১১১) বছর = ৩৩ বছর।}
\end{quote}

\textbf{তিন ভাইয়ের গড় বয়স ১৬ বছর। পিতাসহ ৩ ভাইয়ের বয়সের গড় ২৫ বছর হলে,
পিতার বয়স কত?}

ক) ৪৫

খ) ৪২

\textbf{\hl{গ) ৫২ ✔}}

ঘ) ৪১

\textbf{\hl{✔ সঠিক উত্তর: গ) ৫২}}

\begin{quote}
\textbf{ব্যাখ্যা:} \emph{তিন ভাইয়ের বয়সের সমষ্টি = (৩ × ১৬) বছর = ৪৮ বছর।}

\emph{পিতাসহ তিন ভাই অর্থাৎ ৪ জনের বয়সের সমষ্টি = (৪ × ২৫) বছর = ১০০ বছর।}

\emph{∴ পিতার বয়স = (১০০ − ৪৮) বছর = ৫২ বছর।}
\end{quote}

\textbf{পিতা ও পুত্রের বয়সের গড় ২০ বছর। ২ বছর পূর্বে দুই পুত্রের বয়সের গড় ছিল
১২ বছর। পিতার বয়স কত?}

\textbf{\hl{ক) ২৬ বছর ✔}}

খ) ২৮বছর

গ) ৩০ বছর

ঘ) ৩২ বছর

\textbf{\hl{✔ সঠিক উত্তর: ক) ২৬ বছর}}

\begin{quote}
\textbf{ব্যাখ্যা:} \emph{পিতা ও ১ পুত্রের বয়সের সমষ্টি = (২০ × ২) বছর = ৪০
বছর।}

\emph{২ বছর পূর্বে ২ পুত্রের বয়সের সমষ্টি ছিল = (১২ × ২) বছর = ২৪ বছর।}

\emph{∴ বর্তমানে ২ পুত্রের বয়সের সমষ্টি = (২৪ + ৪) বছর = ২৮ বছর।}

\emph{∴ বর্তমানে ১ পুত্রের বয়স = (২৮ ÷ ২) বছর = ১৪ বছর।}

\emph{∴ পিতার বয়স = (৪০ − ১৪) বছর = ২৬ বছর।}
\end{quote}

\textbf{পিতা ও দুই পুত্রের বর্তমান গড় বয়স ২০ বছর। দুই বছর পর দুই পুত্রের গড় বয়স
১২ বছর হলে পিতার বর্তমান বয়স কত বছর?}

\textbf{\hl{ক) ৪০ ✔}}

খ) ৩৬

গ) ৩০

ঘ) ৪৩

\textbf{\hl{✔ সঠিক উত্তর: ক) ৪০}}

\begin{quote}
\textbf{ব্যাখ্যা:} \emph{পিতা ও ২ পুত্রের বয়সের সমষ্টি = (৩ × ২০) বছর = ৬০
বছর।}

\emph{২ বছর পর ২ পুত্রের গড় বয়স ১২ বছর হলে, বর্তমানে তাদের গড় বয়স = (১২ −
২) = ১০ বছর।}

\emph{∴ বর্তমানে ২ পুত্রের বয়সের সমষ্টি = (২ × ১০) বছর = ২০ বছর।}

\emph{∴ পিতার বর্তমান বয়স = (৬০ − ২০) বছর = ৪০ বছর।}
\end{quote}

\textbf{পিতা ও দুই পুত্রের বয়স অপেক্ষা মাতা ও উক্ত দুই পুত্রের বয়সের গড় ২ বছর
কম। মাতার বয়স ২৫ বছর হলে পিতার বয়স কত?}

ক) ৩০ বছর

\textbf{\hl{খ) ৩১ বছর ✔}}

গ) ৩২ বছর

ঘ) ৩৪ বছর

\textbf{\hl{✔ সঠিক উত্তর: খ) ৩১ বছর}}

\begin{quote}
\textbf{ব্যাখ্যা:} \emph{পিতা ও দুই পুত্র অর্থাৎ মোট ৩ জন ব্যক্তি।}

\emph{মাতা ও দুই পুত্র অর্থাৎ মোট ৩ জন ব্যক্তি।}

\emph{গড় বয়স ২ বছর কম হওয়ায় ৩ জনের মোট বয়সের পার্থক্য = (৩ × ২) = ৬ বছর।}

\emph{যেহেতু দুই পুত্রের বয়স উভয় ক্ষেত্রে সমান, তাই পিতা ও মাতার বয়সের পার্থক্য
= ৬ বছর।}

\emph{∴ পিতার বয়স − মাতার বয়স = ৬ বছর}

\emph{⇒ পিতার বয়স − ২৫ = ৬}

\emph{∴ পিতার বয়স = (২৫ + ৬) বছর = ৩১ বছর।}
\end{quote}

\textbf{২৪ জন ছাত্র এবং একজন শিক্ষকের বয়সের গড় ১৫ বছর। শিক্ষককে বাদ দিয়ে
ছাত্রদের বয়সের গড় করলে গড় ১ বছর কমে যায়। শিক্ষকের বয়স কত?}

ক) ৩৮

খ) ৪০

গ) ১২

\textbf{\hl{ঘ) কোনোটিই নয় ✔}}

\textbf{\hl{✔ সঠিক উত্তর: ঘ) কোনোটিই নয়}}

\begin{quote}
\textbf{ব্যাখ্যা:} \emph{ছাত্র ও শিক্ষকসহ ২৫ জনের বয়সের সমষ্টি = (২৫ × ১৫) =
৩৭৫ বছর।}

\emph{শিক্ষককে বাদ দিলে ২৪ জন ছাত্রের বয়সের গড় হয় = (১৫ − ১) = ১৪ বছর।}

\emph{∴ ২৪ জন ছাত্রের বয়সের সমষ্টি = (২৪ × ১৪) = ৩৩৬ বছর।}

\emph{∴ শিক্ষকের বয়স = (৩৭৫ − ৩৩৬) = ৩৯ বছর।}

\emph{(যেহেতু ৩৯ বছর অপশনগুলোতে নেই, তাই সঠিক উত্তর
\textquotesingle কোনোটিই নয়\textquotesingle।)}
\end{quote}

\textbf{কোনো শ্রেণীর ২৪ জন ছাত্রের বয়স ১৪ বছর। যদি একজন শ্রেণী শিক্ষকের বয়স
তাদের বয়সের সাথে যোগ করা হয় তবে বয়সের গড় এক বছর বৃদ্ধি পায়। শিক্ষকের বয়স
কত?}

ক) ৪০ বৎসর

\textbf{\hl{খ) ৩৯ বৎসর ✔}}

গ) ৪৫ বৎসর

ঘ) ৩৫ বৎসর

\textbf{\hl{✔ সঠিক উত্তর: খ) ৩৯ বৎসর}}

\begin{quote}
\textbf{ব্যাখ্যা:} \emph{২৪ জন ছাত্রের মোট বয়স = (২৪ × ১৪) বছর = ৩৩৬ বছর।}

\emph{শিক্ষকসহ (২৪ + ১) বা ২৫ জনের বয়সের গড় = (১৪ + ১) বা ১৫ বছর।}

\emph{∴ ২৫ জনের মোট বয়স = (২৫ × ১৫) বছর = ৩৭৫ বছর।}

\emph{∴ শিক্ষকের বয়স = (৩৭৫ − ৩৩৬) বছর = ৩৯ বছর।}
\end{quote}

\textbf{১০ জন ছাত্রের গড় বয়স ১৫ বছর। নতুন একজন ছাত্র আসায় গড় বয়স ১৬ বছর হলে
নতুন ছাত্রের বয়স কত বছর?}

ক) ২০

খ) ২৪

\textbf{\hl{গ) ২৬ ✔}}

ঘ) কোনোটিই নয়

\textbf{\hl{✔ সঠিক উত্তর: গ) ২৬}}

\begin{quote}
\textbf{ব্যাখ্যা:} \emph{নতুন ১ জনসহ (১০ + ১) বা ১১ জনের মোট বয়স = (১১ ×
১৬) বছর = ১৭৬ বছর।}

\emph{১০ জন ছাত্রের মোট বয়স = (১০ × ১৫) বছর = ১৫০ বছর।}

\emph{∴ নতুন ছাত্রের বয়স = (১৭৬ − ১৫০) বছর = ২৬ বছর।}
\end{quote}

\textbf{৩ ভাইয়ের বয়সের গড় ১৬ বছর। তাদের বাবাসহ তাদের বয়সের গড় ২৫ বছর।
তাদের বাবার বয়স কত?}

\textbf{\hl{ক) ৫২ ✔}}

খ) ৪১

গ) ৪৫

ঘ) ৪২

\textbf{\hl{✔ সঠিক উত্তর: ক) ৫২}}

\begin{quote}
\textbf{ব্যাখ্যা:} \emph{বাবা ও ৩ ভাই অর্থাৎ ৪ জনের মোট বয়স = (২৫ × ৪) বছর
= ১০০ বছর।}

\emph{৩ ভাইয়ের মোট বয়স = (৩ × ১৬) বছর = ৪৮ বছর।}

\emph{∴ বাবার বয়স = (১০০ − ৪৮) বছর = ৫২ বছর।}
\end{quote}

\textbf{৪ বছর আগে ক ও খ এর গড় বয়স ছিল ১৮ বছর। ক, খ ও গ-এর বর্তমান গড় বয়স
২৪ বছর। ৮ বছর পর গ-এর বয়স কত হবে?}

ক) ২৮ বছর

\textbf{\hl{খ) ৩৬ বছর ✔}}

গ) ৩৮ বছর

ঘ) ৪০ বছর

\textbf{\hl{✔ সঠিক উত্তর: খ) ৩৬ বছর}}

\begin{quote}
\textbf{ব্যাখ্যা:} \emph{৪ বছর আগে ক ও খ-এর মোট বয়স ছিল = (১৮ × ২) = ৩৬
বছর।}

\emph{∴ ক ও খ-এর বর্তমান মোট বয়স = ৩৬ + (৪ + ৪) = ৪৪ বছর।}

\emph{ক, খ ও গ-এর বর্তমান মোট বয়স = (২৪ × ৩) = ৭২ বছর।}

\emph{∴ গ-এর বর্তমান বয়স = (৭২ − ৪৪) = ২৮ বছর।}

\emph{∴ ৮ বছর পর গ-এর বয়স হবে = (২৮ + ৮) = ৩৬ বছর।}
\end{quote}

\textbf{৩ বছর আগে ক ও খ এর বয়সের গড় ছিল ১৬ বছর। গ তাদের সাথে যোগ দেওয়ায়
বর্তমান তাদের বয়সের গড় ২০ বছর হয়। গ এর বয়স কত?}

ক) ২০ বছর

\textbf{\hl{খ) ২২ বছর ✔}}

গ) ২৪ বছর

ঘ) ২৬ বছর

\textbf{\hl{✔ সঠিক উত্তর: খ) ২২ বছর}}

\begin{quote}
\textbf{ব্যাখ্যা:} \emph{৩ বছর আগে ক ও খ-এর মোট বয়স ছিল = (১৬ × ২) = ৩২
বছর।}

\emph{∴ বর্তমানে ক ও খ-এর মোট বয়স = ৩২ + (৩ + ৩) = ৩৮ বছর।}

\emph{ক, খ ও গ-এর বর্তমান মোট বয়স = (২০ × ৩) = ৬০ বছর।}

\emph{∴ গ-এর বর্তমান বয়স = (৬০ − ৩৮) = ২২ বছর।}
\end{quote}

\textbf{পিতা ও দুই পুত্রের বর্তমান গড় বয়স ২২ বৎসর। ৩ বছর পর দুই পুত্রের গড় বয়স
১৩ বৎসর হলে পিতার বর্তমান বয়স কত?}

ক) ৪০ বৎসর

খ) ৪২ বৎসর

গ) ৪৩ বৎসর

\textbf{\hl{ঘ) ৪৬ বৎসর ✔}}

\textbf{\hl{✔ সঠিক উত্তর: ঘ) ৪৬ বৎসর}}

\begin{quote}
\textbf{ব্যাখ্যা:} \emph{পিতা ও দুই পুত্র (মোট ৩ জন)-এর বর্তমান বয়সের সমষ্টি =
(২২ × ৩) = ৬৬ বছর।}

\emph{৩ বছর পর দুই পুত্রের গড় বয়স ১৩ বছর হলে, বর্তমানে তাদের গড় বয়স = (১৩ −
৩) = ১০ বছর।}

\emph{∴ বর্তমানে দুই পুত্রের বয়সের সমষ্টি = (১০ × ২) = ২০ বছর।}

\emph{∴ পিতার বর্তমান বয়স = (৬৬ − ২০) = ৪৬ বছর।}
\end{quote}

\textbf{১৫ জন ছাত্রের বয়সের গড় ২৯ বছর। তাদের মধ্যে দুজন ছাত্রের বয়সের গড় ৫৫
বছর। বাকী ১৩ জনের বয়সের গড় কত হবে?}

ক) ২৯ বছর

খ) ২৭ বছর

গ) ২৬ বছর

\textbf{\hl{ঘ) ২৫ বছর ✔}}

\textbf{\hl{✔ সঠিক উত্তর: ঘ) ২৫ বছর}}

\begin{quote}
\textbf{ব্যাখ্যা:} \emph{১৫ জন ছাত্রের মোট বয়স = (১৫ × ২৯) = ৪৩৫ বছর।}

\emph{২ জন ছাত্রের মোট বয়স = (৫৫ × ২) = ১১০ বছর।}

\emph{∴ বাকী ১৩ জন ছাত্রের মোট বয়স = (৪৩৫ − ১১০) = ৩২৫ বছর।}

\emph{∴ বাকী ১৩ জনের বয়সের গড় = (৩২৫ ÷ ১৩) = ২৫ বছর।}
\end{quote}

\textbf{কোনো শ্রেণীতে ২০ জন ছাত্রীর বয়সের গড় ১২ বছর। ৪ জন নতুন ছাত্রী ভর্তি
হওয়াতে বয়সের গড় ৪ মাস কমে গেল। নতুন ৪ জন ছাত্রীর বয়সের গড় কত?}

ক) ৮ বছর

খ) ৯ বছর

\textbf{\hl{গ) ১০ বছর ✔}}

ঘ) ১১ বছর

\textbf{\hl{✔ সঠিক উত্তর: গ) ১০ বছর}}

\begin{quote}
\textbf{ব্যাখ্যা:} \emph{২০ জন ছাত্রীর বয়সের সমষ্টি = (২০ × ১২) = ২৪০ বছর।}

\emph{নতুন ৪ জন সহ মোট ছাত্রী = (২০ + ৪) = ২৪ জন।}

\emph{নতুন গড় বয়স = ১২ বছর − ৪ মাস = ১১ বছর ৮ মাস।}

\emph{২৪ জন ছাত্রীর মোট বয়স = ২৪ × (১১ বছর ৮ মাস) = ২৬৪ বছর + ১৯২ মাস =
২৬৪ বছর + ১৬ বছর = ২৮০ বছর।}

\emph{∴ নতুন ৪ জন ছাত্রীর মোট বয়স = (২৮০ − ২৪০) = ৪০ বছর।}

\emph{∴ নতুন ৪ জন ছাত্রীর বয়সের গড় = (৪০ ÷ ৪) = ১০ বছর।}
\end{quote}

\textbf{কোন শ্রেণীতে ২০ ছাত্রদের বয়সের গড় ১০ বছর। শিক্ষকসহ তাদের বয়সের গড়
১২ বছর হলে শিক্ষকের বয়স কত?}

ক) ৩২ বছর

খ) ৪২ বছর

গ) ৬২ বছর

\textbf{\hl{ঘ) ৫২ বছর ✔}}

\textbf{\hl{✔ সঠিক উত্তর: ঘ) ৫২ বছর}}

\begin{quote}
\textbf{ব্যাখ্যা:} \emph{২০ জন ছাত্রের বয়সের সমষ্টি = (২০ × ১০) = ২০০ বছর।}

\emph{শিক্ষকসহ ২১ জনের বয়সের সমষ্টি = (২১ × ১২) = ২৫২ বছর।}

\emph{∴ শিক্ষকের বয়স = (২৫২ − ২০০) = ৫২ বছর।}
\end{quote}

\textbf{৬ জন পুরুষ, ৮ জন স্ত্রীলোক এবং ১ জন বালকের বয়সের গড় ৩৫ বছর। পুরুষদের
বয়সের গড় ৪০ বছর এবং স্ত্রীলোকদের বয়সের গড় ৩৪ বছর। বালকের বয়স কত?}

\textbf{\hl{ক) ১৩ বছর ✔}}

খ) ১৪ বছর

গ) ১৫ বছর

ঘ) ১৬ বছর

\textbf{\hl{✔ সঠিক উত্তর: ক) ১৩ বছর}}

\begin{quote}
\textbf{ব্যাখ্যা:} \emph{৬ জন পুরুষ, ৮ জন স্ত্রীলোক এবং ১ জন বালক অর্থাৎ মোট ১৫
জনের বয়সের সমষ্টি = (১৫ × ৩৫) = ৫২৫ বছর।}

\emph{৬ জন পুরুষের বয়সের সমষ্টি = (৬ × ৪০) = ২৪০ বছর।}

\emph{৮ জন স্ত্রীলোকের বয়সের সমষ্টি = (৮ × ৩৪) = ২৭২ বছর।}

\emph{∴ বালকের বয়স = ৫২৫ − (২৪০ + ২৭২) = ৫২৫ − ৫১২ = ১৩ বছর।}
\end{quote}

\textbf{The captain of Bangladesh Cricket Team is 26 years old and the
wicket keeper is 4 years older. If the ages of these two are excluded,
the average age of the remaining players is one year less than the
average age of the whole team. What is the average age of the team?}

ক) 21

খ) 22

গ) 23

\textbf{\hl{ঘ) 23.5 ✔}}

\textbf{\hl{✔ সঠিক উত্তর: ঘ) 23.5}}

\begin{quote}
\textbf{ব্যাখ্যা:} \emph{অধিনায়কের বয়স = ২৬ বছর।}

\emph{উইকেটরক্ষকের বয়স = (২৬ + ৪) বা ৩০ বছর।}

\emph{ধরি, ক্রিকেট দলের (১১ জন খেলোয়াড়ের) গড় বয়স = x বছর।}

\emph{∴ মোট বয়সের সমষ্টি = ১১x বছর।}

\emph{২ জন বাদ দিলে অবশিষ্ট ৯ জনের গড় বয়স = (x − ১) বছর।}

\emph{∴ ৯ জনের মোট বয়স = ৯(x − ১) বছর।}

\emph{প্রশ্নমতে,}

\emph{১১x = ৯(x − ১) + ২৬ + ৩০}

\emph{⇒ ১১x = ৯x − ৯ + ৫৬}

\emph{⇒ ১১x − ৯x = ৪৭}

\emph{⇒ ২x = ৪৭}

\emph{∴ x = ২৩.৫ বছর।}
\end{quote}

\textbf{A ও B এর মোট বয়স B ও C এর মোট বয়সের তুলনায় ১৫ বছর বেশি। C, A-এর
চেয়ে কত বছরের ছোট?}

ক) ১০

\textbf{\hl{খ) ১৫ ✔}}

গ) ২০

ঘ) ২৫

\textbf{\hl{✔ সঠিক উত্তর: খ) ১৫}}

\begin{quote}
\textbf{ব্যাখ্যা:} \emph{প্রশ্নমতে,}

\emph{(A + B) − (B + C) = ১৫}

\emph{⇒ A + B − B − C = ১৫}

\emph{⇒ A − C = ১৫}

\emph{∴ A − ১৫ = C}

\emph{∴ C, A-এর চেয়ে ১৫ বছরের ছোট।}
\end{quote}

\textbf{পিতা ও পুত্রের বয়সের অনুপাত ১০ : ৩। পুত্রের বয়স ১৮ হলে, পিতার বয়স
কত?}

ক) ৬২ বছর

\textbf{\hl{খ) ৬০ বছর ✔}}

গ) ৫৪ বছর

ঘ) ৫৮ বছর

\textbf{\hl{✔ সঠিক উত্তর: খ) ৬০ বছর}}

\begin{quote}
\textbf{ব্যাখ্যা:} \emph{ধরি, অনুপাত অনুযায়ী পিতার বয়স = ১০x এবং পুত্রের বয়স =
৩x বছর।}

\emph{প্রশ্নমতে,}

\emph{৩x = ১৮}

\emph{⇒ x = (১৮/৩)}

\emph{∴ x = ৬}

\emph{∴ পিতার বয়স = (১০ × ৬) = ৬০ বছর।}
\end{quote}

\textbf{পিতা ও পুত্রের বয়সের গড় ৩০ বছর। ৬ বছর পরে তাদের বয়সের অনুপাত ৫:১
হলে, পুত্রের বর্তমান বয়স কত বছর?}

ক) ৫

\textbf{\hl{খ) ৬ ✔}}

গ) ৮

ঘ) ৯

\textbf{\hl{✔ সঠিক উত্তর: খ) ৬}}

\begin{quote}
\textbf{ব্যাখ্যা:} \emph{পিতা ও পুত্রের বর্তমান মোট বয়স = (৩০ × ২) বছর = ৬০
বছর।}

\emph{ধরি, ৬ বছর পর পিতা ও পুত্রের বয়স যথাক্রমে ৫x ও x বছর।}

\emph{∴ বর্তমানে পিতার বয়স = (৫x − ৬) বছর}

\emph{এবং বর্তমানে পুত্রের বয়স = (x − ৬) বছর।}

\emph{প্রশ্নমতে,}

\emph{(৫x − ৬) + (x − ৬) = ৬০}

\emph{⇒ ৬x − ১২ = ৬০}

\emph{⇒ ৬x = ৬০ + ১২}

\emph{⇒ ৬x = ৭২}

\emph{∴ x = ১২}

\emph{∴ পুত্রের বর্তমান বয়স = (১২ − ৬) বছর = ৬ বছর।}
\end{quote}

\textbf{পিতার বয়স পুত্রের বয়সের ৪ গুণ, ৬ বছর পূর্বে পিতার বয়স পুত্রের বয়সের ১০
গুণ ছিল। পিতা ও পুত্রের বর্তমান বয়সের অনুপাত কত?}

\textbf{\hl{ক) ৪ : ১ ✔}}

খ) ১ : ৪

গ) ২ : ৩

ঘ) ৩ : ২

\textbf{\hl{✔ সঠিক উত্তর: ক) ৪ : ১}}

\begin{quote}
\textbf{ব্যাখ্যা:} \emph{ধরি, পুত্রের বর্তমান বয়স = x বছর।}

\emph{∴ পিতার বর্তমান বয়স = ৪x বছর।}

\emph{∴ পিতা ও পুত্রের বর্তমান বয়সের অনুপাত = ৪x : x = ৪ : ১।}
\end{quote}

\textbf{কোনো ক্লাসে বালকদের গড় বয়স বালিকাদের সংখ্যার দ্বিগুণ। ৩৬ জনের ঐ
ক্লাসে বালক ও বালিকাদের সংখ্যার অনুপাত হলো ৫:১। ঐ ক্লাসের বালকদের মোট বয়স কত
বছর?}

ক) ৩০০

খ) ৩২০

\textbf{\hl{গ) ৩৬০ ✔}}

ঘ) ৪০০

\textbf{\hl{✔ সঠিক উত্তর: গ) ৩৬০}}

\begin{quote}
\textbf{ব্যাখ্যা:} \emph{বালক ও বালিকার সংখ্যার অনুপাত = ৫ : ১}

\emph{∴ অনুপাতের যোগফল = ৫ + ১ = ৬}

\emph{৩৬ জনের ক্লাসে বালকের সংখ্যা = ৩৬ × (৫/৬) = ৩০ জন।}

\emph{এবং বালিকার সংখ্যা = ৩৬ × (১/৬) = ৬ জন।}

\emph{প্রশ্নমতে, বালকদের গড় বয়স = বালিকাদের সংখ্যার দ্বিগুণ = (৬ × ২) = ১২
বছর।}

\emph{∴ বালকদের মোট বয়স = (৩০ × ১২) বছর = ৩৬০ বছর।}
\end{quote}

\textbf{কন্যার বয়স পিতা ও পুত্রের বয়সের মধ্য সমানুপাতী; পিতা ও পুত্রের বয়স
যথাক্রমে ৪৮ ও ১২ হলে কন্যার বয়স হবে-}

ক) ২৮ বৎসর

\textbf{\hl{খ) ২৪ বৎসর ✔}}

গ) ১৬ বৎসর

ঘ) ১২ বৎসর

\textbf{\hl{✔ সঠিক উত্তর: খ) ২৪ বৎসর}}

\begin{quote}
\textbf{ব্যাখ্যা:} \emph{ক্রমিক সমানুপাতের ক্ষেত্রে আমরা জানি,}

\emph{(মধ্যরাশি)\^{}২ = ১ম রাশি × ৩য় রাশি}

\emph{এখানে,}

\emph{১ম রাশি (পিতার বয়স) = ৪৮ বছর}

\emph{৩য় রাশি (পুত্রের বয়স) = ১২ বছর}

\emph{মধ্যরাশি (কন্যার বয়স) = ?}

\emph{∴ (কন্যার বয়স)\^{}২ = ৪৮ × ১২}

\emph{⇒ (কন্যার বয়স)\^{}২ = ৫৭৬}

\emph{⇒ কন্যার বয়স = √(৫৭৬)}

\emph{∴ কন্যার বয়স = ২৪ বছর।}
\end{quote}

\textbf{২ বছর আগে বাবার বয়স পুত্রের বয়সের ১৪ গুণ। দুই বছর বাদে বাবার বয়স
পুত্রের বয়সের চেয়ে ২৬ বছর বেশি হয়। তবে বাবা ও পুত্রের বয়সের অনুপাত কত হবে?}

\textbf{\hl{ক) ১৫ : ২ ✔}}

খ) ৩ : ৬

গ) ১ : ৪

ঘ) এর কোনটিই নয়

\textbf{\hl{✔ সঠিক উত্তর: ক) ১৫ : ২}}

\begin{quote}
\textbf{ব্যাখ্যা:} \emph{ধরি, ২ বছর আগে পুত্রের বয়স ছিল ক বছর।}

\emph{∴ ২ বছর আগে বাবার বয়স ছিল ১৪ ক বছর।}

\emph{দুই বছর বাদে বাবার বয়স = ১৪ ক + ২, পুত্রের বয়স ক + ২ বছর।}

\emph{প্রশ্নমতে, ১৪ ক + ২ − (ক + ২) = ২৬ ; ∴ ক = ২;}

\emph{বাবার বয়স ছিল = ১৪ × ২ = ২৮ বছর।}

\emph{∴ দুই বছর বাদে বাবা ও পুত্রের বয়সের অনুপাত = (২৮ + ২) : (২ + ২)}

\emph{= ৩০ : ৪ = ১৫ : ২}
\end{quote}

\textbf{পিতা ও দুই সন্তানের বয়সের গড় ৩০ বছর। দুই সন্তানের বয়সের গড় ২০ বছর
হলে পিতার বয়স কত?}

ক) ২০ বছর

খ) ৩০ বছর

গ) ৪০ বছর

\textbf{\hl{ঘ) ৫০ বছর ✔}}

\textbf{\hl{✔ সঠিক উত্তর: ঘ) ৫০ বছর}}

\begin{quote}
\textbf{ব্যাখ্যা:} \emph{পিতা ও দুই সন্তান (মোট ৩ জন)-এর বয়সের সমষ্টি = (৩০ ×
৩) = ৯০ বছর।}

\emph{দুই সন্তানের বয়সের সমষ্টি = (২০ × ২) = ৪০ বছর।}

\emph{∴ পিতার বয়স = (৯০ − ৪০) = ৫০ বছর।}
\end{quote}

\textbf{পিতা ও পুত্রের বর্তমান বয়সের সমষ্টি ৬০ বছর। ৫ বছর পরে তাদের বয়সের
অনুপাত হবে ৫ : ২। পুত্রের বর্তমান বয়স কত বছর?}

ক) ১০

খ) ২০

\textbf{\hl{গ) ১৫ ✔}}

ঘ) ৪৫

\textbf{\hl{✔ সঠিক উত্তর: গ) ১৫}}

\begin{quote}
\textbf{ব্যাখ্যা:} \emph{ধরি, পুত্রের বর্তমান বয়স = x বছর।}

\emph{∴ পিতার বর্তমান বয়স = (৬০ − x) বছর।}

\emph{৫ বছর পর:}

\emph{পুত্রের বয়স হবে = (x + ৫) বছর}

\emph{পিতার বয়স হবে = (৬০ − x + ৫) = (৬৫ − x) বছর}

\emph{প্রশ্নমতে,}

\emph{(৬৫−x/x+৫)=(৫/২)}

\emph{⇒ ৫(x + ৫) = ২(৬৫ − x)}

\emph{⇒ ৫x + ২৫ = ১৩০ − ২x}

\emph{⇒ ৫x + ২x = ১৩০ − ২৫}

\emph{⇒ ৭x = ১০৫}

\emph{⇒ x = (১০৫/৭)}

\emph{∴ x = ১৫}

\emph{∴ পুত্রের বর্তমান বয়স ১৫ বছর।}
\end{quote}

\textbf{রহিম ও করিমের বয়সের অনুপাত ৩:৫। তাদের বয়সের সমষ্টি ৪০ হলে নিম্নের
কোন উত্তরটি সঠিক?}

\textbf{\hl{ক) রহিম ১৫ ✔}}

খ) করিম ১৫

গ) রহিম ১০

ঘ) করিম ৫

\textbf{\hl{✔ সঠিক উত্তর: ক) রহিম ১৫}}

\begin{quote}
\textbf{ব্যাখ্যা:} \emph{রহিম ও করিমের বয়সের অনুপাত = ৩ : ৫}

\emph{অনুপাতদ্বয়ের যোগফল = ৩ + ৫ = ৮}

\emph{রহিমের বয়স = ৪০ এর (৩/৮) = ১৫ বছর}

\emph{এবং করিমের বয়স = ৪০ এর (৫/৮) = ২৫ বছর}

\emph{∴ অপশনগুলোর মধ্যে সঠিক উত্তর \textquotesingle রহিম
১৫\textquotesingle।}
\end{quote}

\textbf{পিতা পুত্রের বয়সের সমষ্টি ৫৪ বছর এবং অনুপাত ৭ : ২। ১০ বছর পর তাদের
বয়সের অনুপাত কত?}

ক) ৩১, ১৬

\textbf{\hl{খ) ২৬, ১১ ✔}}

গ) ১৭, ১২

ঘ) ২, ১

\textbf{\hl{✔ সঠিক উত্তর: খ) ২৬, ১১}}

\begin{quote}
\textbf{ব্যাখ্যা:} \emph{ধরি, পিতার বর্তমান বয়স = ৭x বছর}

\emph{এবং পুত্রের বর্তমান বয়স = ২x বছর}

\emph{প্রশ্নমতে,}

\emph{৭x + ২x = ৫৪}

\emph{⇒ ৯x = ৫৪}

\emph{∴ x = ৬}

\emph{∴ পিতার বর্তমান বয়স = ৭ × ৬ = ৪২ বছর}

\emph{এবং পুত্রের বর্তমান বয়স = ২ × ৬ = ১২ বছর}

\emph{১০ বছর পর:}

\emph{পিতার বয়স হবে = ৪২ + ১০ = ৫২ বছর}

\emph{পুত্রের বয়স হবে = ১২ + ১০ = ২২ বছর}

\emph{∴ ১০ বছর পর তাদের বয়সের অনুপাত হবে = ৫২ : ২২ = ২৬ : ১১}
\end{quote}

\textbf{পিতা ও পুত্রের বর্তমান বয়সের অনুপাত ৭:২ এবং ৫ বছর পরে তাদের বয়সের
অনুপাত ৮ : ৩ হলে, পিতার বর্তমান বয়স কত?}

ক) ৩০ বছর

\textbf{\hl{খ) ৩৫ বছর ✔}}

গ) ৪০ বছর

ঘ) ৪৫ বছর

\textbf{\hl{✔ সঠিক উত্তর: খ) ৩৫ বছর}}

\begin{quote}
\textbf{ব্যাখ্যা:} \emph{ধরি, পিতার বর্তমান বয়স ৭x বছর এবং পুত্রের বর্তমান বয়স
২x বছর।}

\emph{শর্তমতে,}

\emph{(৭x+৫/২x+৫)=(৮/৩)}

\emph{⇒ ২১x + ১৫ = ১৬x + ৪০}

\emph{⇒ ২১x − ১৬x = ৪০ − ১৫}

\emph{⇒ ৫x = ২৫}

\emph{⇒ x = (২৫/৫)}

\emph{∴ x = ৫}

\emph{∴ পিতার বর্তমান বয়স = (৭ × ৫) = ৩৫ বছর।}
\end{quote}

\textbf{বর্তমানে পিতা ও পুত্রের বয়সের অনুপাত ৭:৩। চার বছর পূর্বে তাদের বয়সের
অনুপাত ছিল ১৩:৫। বর্তমানে পুত্রের বয়স কত?}

\textbf{\hl{ক) ২৪ বছর ✔}}

খ) ১৬ বছর

গ) ১২ বছর

ঘ) ৮ বছর

\textbf{\hl{✔ সঠিক উত্তর: ক) ২৪ বছর}}

\begin{quote}
\textbf{ব্যাখ্যা:} \emph{ধরি, পিতার বর্তমান বয়স ৭x বছর এবং পুত্রের বর্তমান বয়স
৩x বছর।}

\emph{তাহলে, ৪ বছর পূর্বে তাদের বয়স ছিল যথাক্রমে (৭x − ৪) বছর এবং (৩x − ৪)
বছর।}

\emph{প্রশ্নমতে,}

\emph{(৭x−৪/৩x−৪)=(১৩/৫)}

\emph{⇒ ১৩(৩x − ৪) = ৫(৭x − ৪)}

\emph{⇒ ৩৯x − ৫২ = ৩৫x − ২০}

\emph{⇒ ৩৯x − ৩৫x = ৫২ − ২০}

\emph{⇒ ৪x = ৩২}

\emph{∴ x = ৮}

\emph{∴ পুত্রের বর্তমান বয়স = ৩x = (৩ × ৮) = ২৪ বছর।}
\end{quote}

\textbf{পিতা ও পুত্রের বয়সের সমষ্টি ৬৩ বছর এবং তাদের বয়সের অনুপাত ৭ : ২। ৯
বছর পূর্বে তাদের বয়সের অনুপাত কত ছিল?}

ক) ৯ : ১

খ) ৭ : ১

গ) ৬ : ১

\textbf{\hl{ঘ) ৮ : ১ ✔}}

\textbf{\hl{✔ সঠিক উত্তর: ঘ) ৮ : ১}}

\begin{quote}
\textbf{ব্যাখ্যা:} \emph{অনুপাতগুলোর সমষ্টি = ৭ + ২ = ৯}

\emph{পিতার বর্তমান বয়স = (৬৩ এর (৭/৯) ) = ৪৯ বছর}

\emph{পুত্রের বর্তমান বয়স = (৬৩ এর (২/৯) ) = ১৪ বছর}

\emph{৯ বছর পূর্বে:}

\emph{পিতার বয়স ছিল = (৪৯ − ৯) = ৪০ বছর}

\emph{পুত্রের বয়স ছিল = (১৪ − ৯) = ৫ বছর}

\emph{∴ ৯ বছর পূর্বে তাদের বয়সের অনুপাত = ৪০ : ৫ = ৮ : ১}
\end{quote}

\textbf{পিতা ও পুত্রের বয়সের সমষ্টি ৮৪ বছর। ১০ বছর পূর্বে তাদের বয়সের অনুপাত ৫
: ৩ ছিল। ১০ বছর পর এ অনুপাত কত হবে?}

ক) ২৭ : ১৭

\textbf{\hl{খ) ১৫ : ১১ ✔}}

গ) ১৭ : ১৩

ঘ) ১৯ : ১১

\textbf{\hl{✔ সঠিক উত্তর: খ) ১৫ : ১১}}

\begin{quote}
\textbf{ব্যাখ্যা:} \emph{ধরি, ১০ বছর পূর্বে পিতার বয়স ছিল ৫x বছর এবং পুত্রের
বয়স ছিল ৩x বছর।}

\emph{∴ পিতা ও পুত্রের বর্তমান বয়স যথাক্রমে (৫x + ১০) এবং (৩x + ১০) বছর।}

\emph{প্রশ্নমতে,}

\emph{(৫x + ১০) + (৩x + ১০) = ৮৪}

\emph{⇒ ৮x + ২০ = ৮৪}

\emph{⇒ ৮x = ৬৪}

\emph{∴ x = ৮}

\emph{∴ ১০ বছর পর:}

\emph{পিতার বয়স হবে = (৫x + ১০) + ১০ = ৫x + ২০ = (৫ × ৮ + ২০) = ৬০ বছর}

\emph{পুত্রের বয়স হবে = (৩x + ১০) + ১০ = ৩x + ২০ = (৩ × ৮ + ২০) = ৪৪ বছর}

\emph{∴ ১০ বছর পর তাদের বয়সের অনুপাত = ৬০ : ৪৪ = ১৫ : ১১}
\end{quote}

\textbf{১২ বছর আগে শুভ ও পিয়ার বয়সের অনুপাত ছিল ৪:৫। ১২ বছর পর তাদের বয়সের
অনুপাত হবে ৮:৯। তাদের বর্তমান বয়সের অনুপাত কত?}

ক) ৪:৫

\textbf{\hl{খ) ৬:৭ ✔}}

গ) ২:৩

ঘ) ৮:৯

\textbf{\hl{✔ সঠিক উত্তর: খ) ৬:৭}}

\begin{quote}
\textbf{ব্যাখ্যা:} \emph{ধরি, ১২ বছর আগে শুভর বয়স = ৪x বছর}

\emph{এবং পিয়ার বয়স = ৫x বছর।}

\emph{প্রশ্নমতে,}

\emph{(৪x+১২+১২/৫x+১২+১২)=(৮/৯)}

\emph{⇒ (৪x+২৪/৫x+২৪)=(৮/৯)}

\emph{⇒ ৩৬x + ২১৬ = ৪০x + ১৯২}

\emph{⇒ ৪০x − ৩৬x = ২১৬ − ১৯২}

\emph{⇒ ৪x = ২৪}

\emph{∴ x = ৬}

\emph{∴ শুভ ও পিয়ার বর্তমান বয়সের অনুপাত = (৪ × ৬ + ১২) : (৫ × ৬ + ১২)}

\emph{= (২৪ + ১২) : (৩০ + ১২)}

\emph{= ৩৬ : ৪২ = ৬ : ৭}
\end{quote}

\textbf{১০ বছর আগে ক এর বয়স ছিল খ এর বয়সের অর্ধেক। যদি তাদের বর্তমান বয়সের
অনুপাত ৩ : ৪ হয়, তবে তাদের বর্তমানে মোট বয়স কত?}

ক) ৪৫ বছর

\textbf{\hl{খ) ৩৫ বছর ✔}}

গ) ২০ বছর

ঘ) ৫০ বছর

\textbf{\hl{✔ সঠিক উত্তর: খ) ৩৫ বছর}}

\begin{quote}
\textbf{ব্যাখ্যা:} \emph{ধরি, ক এর বর্তমান বয়স = ৩x বছর}

\emph{এবং খ এর বর্তমান বয়স = ৪x বছর।}

\emph{১০ বছর আগে ক এবং খ এর বয়স ছিল যথাক্রমে (৩x − ১০) ও (৪x − ১০) বছর।}

\emph{শর্তমতে,}

\emph{২(৩x − ১০) = ৪x − ১০}

\emph{⇒ ৬x − ২০ = ৪x − ১০}

\emph{⇒ ৬x − ৪x = ২০ − ১০}

\emph{⇒ ২x = ১০}

\emph{∴ x = ৫}

\emph{∴ তাদের বর্তমানে মোট বয়স = ৩x + ৪x = ৭x = (৭ × ৫) = ৩৫ বছর।}
\end{quote}

\textbf{দুই বোনের বর্তমান বয়সের অনুপাত ৪ : ৩। ৬ বছর আগে তাদের বয়সের অনুপাত
ছিল ৫ : ৩। ৬ বছর পরে তাদের বয়সের সমষ্টি কত হবে?}

ক) ২৬

খ) ২৮

\textbf{\hl{গ) ৪০ ✔}}

ঘ) ৪৪

\textbf{\hl{✔ সঠিক উত্তর: গ) ৪০}}

\begin{quote}
\textbf{ব্যাখ্যা:} \emph{ধরি, বর্তমানে দুই বোনের বয়স যথাক্রমে ৪x বছর এবং ৩x
বছর।}

\emph{প্রশ্নমতে,}

\emph{(৪x−৬/৩x−৬)=(৫/৩)}

\emph{⇒ ৫(৩x − ৬) = ৩(৪x − ৬)}

\emph{⇒ ১৫x − ৩০ = ১২x − ১৮}

\emph{⇒ ১৫x − ১২x = ৩০ − ১৮}

\emph{⇒ ৩x = ১২}

\emph{∴ x = ৪}

\emph{∴ বর্তমানে তাদের মোট বয়স = ৪x + ৩x = ৭x = (৭ × ৪) = ২৮ বছর।}

\emph{∴ ৬ বছর পরে তাদের বয়সের সমষ্টি হবে = ২৮ + (৬ × ২) = ২৮ + ১২ = ৪০
বছর।}
\end{quote}

\textbf{Ten years before, the ratio of the ages between father and son
was 3 : 1 and 5 years after it would be 2 : 1. What are their current
ages?}

ক) 50 and 20 years

খ) 49 and 27 years

\textbf{\hl{গ) 55 and 25 years ✔}}

ঘ) 45 and 25 years

\textbf{\hl{✔ সঠিক উত্তর: গ) 55 and 25 years}}

\begin{quote}
\textbf{ব্যাখ্যা:} \emph{ধরি, ১০ বছর আগে পিতার বয়স ছিল ৩x বছর এবং পুত্রের
বয়স ছিল x বছর।}

\emph{∴ পিতার বর্তমান বয়স = (৩x + ১০) বছর}

\emph{এবং পুত্রের বর্তমান বয়স = (x + ১০) বছর।}

\emph{প্রশ্নমতে,}

\emph{(৩x+১০+৫/x+১০+৫)=(২/১)}

\emph{⇒ (৩x+১৫/x+১৫)=(২/১)}

\emph{⇒ ২x + ৩০ = ৩x + ১৫}

\emph{⇒ ৩x − ২x = ৩০ − ১৫}

\emph{∴ x = ১৫}

\emph{∴ পিতার বর্তমান বয়স = (৩ × ১৫ + ১০) বছর = ৫৫ বছর}

\emph{এবং পুত্রের বর্তমান বয়স = (১৫ + ১০) বছর = ২৫ বছর।}
\end{quote}

\textbf{পিতা ও পুত্রের বয়সের অনুপাত ১১ : ৪। পিতার বয়স ৪৪ বছর হলে পিতা ও
পুত্রের বয়সের সমষ্টি কত?}

ক) ৫০ বছর

খ) ৫৫ বছর

\textbf{\hl{গ) ৬০ বছর ✔}}

ঘ) ৬৫ বছর

\textbf{\hl{✔ সঠিক উত্তর: গ) ৬০ বছর}}

\begin{quote}
\textbf{ব্যাখ্যা:} \emph{ধরি, পিতার বয়স ১১x বছর এবং পুত্রের বয়স ৪x বছর।}

\emph{প্রশ্নমতে,}

\emph{১১x = ৪৪}

\emph{∴ x = ৪}

\emph{∴ পুত্রের বয়স = (৪ × ৪) বছর = ১৬ বছর।}

\emph{∴ পিতা ও পুত্রের মোট বয়স = (৪৪ + ১৬) বছর = ৬০ বছর।}
\end{quote}

\textbf{The total age of A, B and C is 90 years. Ten years ago the ratio
of their ages was 1: 2: 3. What is the present age of B? /A, B এবং C এর
মোট বয়স ৯০ বছর। দশ বছর পূর্বে তাদের বয়সের অনুপাত ১ : ২ : ৩। B এর বর্তমান বয়স
কত?}

ক) 40 years

\textbf{\hl{খ) 30 years ✔}}

গ) 20 years

ঘ) 18 years

\textbf{\hl{✔ সঠিক উত্তর: খ) 30 years}}

\begin{quote}
\textbf{ব্যাখ্যা:} \emph{১০ বছর পূর্বে A, B ও C-এর মোট বয়স ছিল = ৯০ − (১০ ×
৩) = ৯০ − ৩০ = ৬০ বছর।}

\emph{১০ বছর পূর্বে তাদের বয়সের অনুপাত = ১ : ২ : ৩}

\emph{অনুপাতগুলোর যোগফল = ১ + ২ + ৩ = ৬}

\emph{∴ ১০ বছর পূর্বে B-এর বয়স ছিল = (৬০ এর (২/৬) ) = ২০ বছর।}

\emph{∴ B-এর বর্তমান বয়স = (২০ + ১০) বছর = ৩০ বছর।}
\end{quote}

\section{পরিশিষ্ট --- দ্রুত প্রস্তুতি: গাণিতিক যুক্তি
সূত্রসংগ্রহ}\label{ux9aaux9b0ux9b6ux9b7ux99f-ux9a6ux9b0ux9a4-ux9aaux9b0ux9b8ux9a4ux9a4-ux997ux9a3ux9a4ux995-ux9afux995ux9a4-ux9b8ux9a4ux9b0ux9b8ux997ux9b0ux9b9}

\emph{পরীক্ষায় সবচেয়ে বেশি ব্যবহৃত সূত্র এক জায়গায়।}

\subsection{বীজগণিতের মূল
সূত্র}\label{ux9acux99cux997ux9a3ux9a4ux9b0-ux9aeux9b2-ux9b8ux9a4ux9b0}

\begin{itemize}
\item
  (a+b)² = a² + 2ab + b²
\item
  (a−b)² = a² − 2ab + b²
\item
  a² − b² = (a+b)(a−b)
\item
  (a+b)³ = a³ + 3a²b + 3ab² + b³
\item
  (a−b)³ = a³ − 3a²b + 3ab² − b³
\item
  a³ + b³ = (a+b)(a² − ab + b²)
\item
  a³ − b³ = (a−b)(a² + ab + b²)
\item
  a³ + b³ + c³ − 3abc = (a+b+c)(a² + b² + c² − ab − bc − ca)
\end{itemize}

\subsection{লাভ-ক্ষতি ও
সুদকষা}\label{ux9b2ux9ad-ux995ux9b7ux9a4-ux993-ux9b8ux9a6ux995ux9b7}

\begin{itemize}
\item
  লাভ = বিক্রয়মূল্য − ক্রয়মূল্য; লাভ\% = (লাভ ÷ ক্রয়মূল্য) × ১০০
\item
  ক্ষতি\% = (ক্ষতি ÷ ক্রয়মূল্য) × ১০০
\item
  বিক্রয়মূল্য = ক্রয়মূল্য × (১০০ ± লাভ/ক্ষতি\%) ÷ ১০০
\item
  সরল সুদ: I = Pnr/১০০; মোট টাকা A = P(১ + nr/১০০)
\item
  চক্রবৃদ্ধি সুদ: A = P(১ + r/১০০)\^{}n
\end{itemize}

\subsection{ক্ষেত্রফল ও
পরিসীমা}\label{ux995ux9b7ux9a4ux9b0ux9abux9b2-ux993-ux9aaux9b0ux9b8ux9ae}

{\def\LTcaptype{none} % do not increment counter
\begin{longtable}[]{@{}
  >{\raggedright\arraybackslash}p{(\linewidth - 4\tabcolsep) * \real{0.2672}}
  >{\raggedright\arraybackslash}p{(\linewidth - 4\tabcolsep) * \real{0.2672}}
  >{\raggedright\arraybackslash}p{(\linewidth - 4\tabcolsep) * \real{0.2672}}@{}}
\toprule\noalign{}
\begin{minipage}[b]{\linewidth}\raggedright
\textbf{চিত্র}
\end{minipage} & \begin{minipage}[b]{\linewidth}\raggedright
\textbf{ক্ষেত্রফল}
\end{minipage} & \begin{minipage}[b]{\linewidth}\raggedright
\textbf{পরিসীমা}
\end{minipage} \\
\midrule\noalign{}
\endhead
\bottomrule\noalign{}
\endlastfoot
বর্গক্ষেত্র (বাহু a) & a² & 4a \\
আয়তক্ষেত্র (দৈর্ঘ্য a, প্রস্থ b) & ab & 2(a+b) \\
ত্রিভুজ (ভূমি b, উচ্চতা h) & ½ bh & তিন বাহুর যোগফল \\
বৃত্ত (ব্যাসার্ধ r) & πr² & 2πr (পরিধি) \\
ট্রাপিজিয়াম (সমান্তরাল বাহু a, b; উচ্চতা h) & ½ (a+b) h & --- \\
সামান্তরিক (ভূমি b, উচ্চতা h) & bh & --- \\
\end{longtable}
}

\subsection{আয়তন}\label{ux986ux9afux9a4ux9a8}

\begin{itemize}
\item
  ঘনক: a³ (পৃষ্ঠক্ষেত্র ৬a²)
\item
  আয়তাকার ঘনবস্তু: l × w × h
\item
  সিলিন্ডার: πr²h (বক্রতল ২πrh)
\item
  কোণ (শঙ্কু): ⅓ πr²h
\item
  গোলক: (৪/৩) πr³ (পৃষ্ঠ ৪πr²)
\end{itemize}

\subsection{সংখ্যা বিভাজ্যতা ও
ভগ্নাংশ}\label{ux9b8ux996ux9af-ux9acux9adux99cux9afux9a4-ux993-ux9adux997ux9a8ux9b6}

\begin{itemize}
\item
  ২ দ্বারা বিভাজ্য: শেষ অঙ্ক জোড়; ৫: শেষ অঙ্ক ০ বা ৫
\item
  ৩ দ্বারা বিভাজ্য: অঙ্কগুলোর যোগফল ৩ দ্বারা বিভাজ্য; ৯: যোগফল ৯ দ্বারা
\item
  ৪ দ্বারা: শেষ দুই অঙ্ক ৪ দ্বারা বিভাজ্য; ৮: শেষ তিন অঙ্ক
\item
  ১১ দ্বারা: বিজোড় ও জোড় স্থানের অঙ্কের যোগফলের পার্থক্য ০ বা ১১ দ্বারা বিভাজ্য
\item
  লসাগু × গসাগু = দুই সংখ্যার গুণফল
\item
  শতকরা ভগ্নাংশ: ৫০\%=১/২, ২৫\%=১/৪, ২০\%=১/৫, ১২.৫\%=১/৮, ৩৩⅓\%=১/৩
\end{itemize}

\subsection{সময়-দূরত্ব-বেগ}\label{ux9b8ux9aeux9af-ux9a6ux9b0ux9a4ux9ac-ux9acux997}

\begin{itemize}
\item
  দূরত্ব = বেগ × সময়; বেগ = দূরত্ব ÷ সময়
\item
  কি.মি./ঘণ্টা → মি./সে: × ৫/১৮; মি./সে → কি.মি./ঘণ্টা: × ১৮/৫
\item
  গড় বেগ (সমদূরত্বে) = 2xy/(x+y)
\item
  ট্রেনের নিজের দৈর্ঘ্যসহ অতিক্রম: সময় = (ট্রেন + বস্তুর দৈর্ঘ্য) ÷ আপেক্ষিক বেগ
\item
  কাজ: A একা n দিনে করলে, A+B একসাথে = (mn)/(m+n) দিনে
\end{itemize}

\end{document}
